% Options for packages loaded elsewhere
% Options for packages loaded elsewhere
\PassOptionsToPackage{unicode}{hyperref}
\PassOptionsToPackage{hyphens}{url}
\PassOptionsToPackage{dvipsnames,svgnames,x11names}{xcolor}
%
\documentclass[
  letterpaper,
  DIV=11,
  numbers=noendperiod]{scrreprt}
\usepackage{xcolor}
\usepackage{amsmath,amssymb}
\setcounter{secnumdepth}{5}
\usepackage{iftex}
\ifPDFTeX
  \usepackage[T1]{fontenc}
  \usepackage[utf8]{inputenc}
  \usepackage{textcomp} % provide euro and other symbols
\else % if luatex or xetex
  \usepackage{unicode-math} % this also loads fontspec
  \defaultfontfeatures{Scale=MatchLowercase}
  \defaultfontfeatures[\rmfamily]{Ligatures=TeX,Scale=1}
\fi
\usepackage{lmodern}
\ifPDFTeX\else
  % xetex/luatex font selection
\fi
% Use upquote if available, for straight quotes in verbatim environments
\IfFileExists{upquote.sty}{\usepackage{upquote}}{}
\IfFileExists{microtype.sty}{% use microtype if available
  \usepackage[]{microtype}
  \UseMicrotypeSet[protrusion]{basicmath} % disable protrusion for tt fonts
}{}
\makeatletter
\@ifundefined{KOMAClassName}{% if non-KOMA class
  \IfFileExists{parskip.sty}{%
    \usepackage{parskip}
  }{% else
    \setlength{\parindent}{0pt}
    \setlength{\parskip}{6pt plus 2pt minus 1pt}}
}{% if KOMA class
  \KOMAoptions{parskip=half}}
\makeatother
% Make \paragraph and \subparagraph free-standing
\makeatletter
\ifx\paragraph\undefined\else
  \let\oldparagraph\paragraph
  \renewcommand{\paragraph}{
    \@ifstar
      \xxxParagraphStar
      \xxxParagraphNoStar
  }
  \newcommand{\xxxParagraphStar}[1]{\oldparagraph*{#1}\mbox{}}
  \newcommand{\xxxParagraphNoStar}[1]{\oldparagraph{#1}\mbox{}}
\fi
\ifx\subparagraph\undefined\else
  \let\oldsubparagraph\subparagraph
  \renewcommand{\subparagraph}{
    \@ifstar
      \xxxSubParagraphStar
      \xxxSubParagraphNoStar
  }
  \newcommand{\xxxSubParagraphStar}[1]{\oldsubparagraph*{#1}\mbox{}}
  \newcommand{\xxxSubParagraphNoStar}[1]{\oldsubparagraph{#1}\mbox{}}
\fi
\makeatother

\usepackage{color}
\usepackage{fancyvrb}
\newcommand{\VerbBar}{|}
\newcommand{\VERB}{\Verb[commandchars=\\\{\}]}
\DefineVerbatimEnvironment{Highlighting}{Verbatim}{commandchars=\\\{\}}
% Add ',fontsize=\small' for more characters per line
\usepackage{framed}
\definecolor{shadecolor}{RGB}{241,243,245}
\newenvironment{Shaded}{\begin{snugshade}}{\end{snugshade}}
\newcommand{\AlertTok}[1]{\textcolor[rgb]{0.68,0.00,0.00}{#1}}
\newcommand{\AnnotationTok}[1]{\textcolor[rgb]{0.37,0.37,0.37}{#1}}
\newcommand{\AttributeTok}[1]{\textcolor[rgb]{0.40,0.45,0.13}{#1}}
\newcommand{\BaseNTok}[1]{\textcolor[rgb]{0.68,0.00,0.00}{#1}}
\newcommand{\BuiltInTok}[1]{\textcolor[rgb]{0.00,0.23,0.31}{#1}}
\newcommand{\CharTok}[1]{\textcolor[rgb]{0.13,0.47,0.30}{#1}}
\newcommand{\CommentTok}[1]{\textcolor[rgb]{0.37,0.37,0.37}{#1}}
\newcommand{\CommentVarTok}[1]{\textcolor[rgb]{0.37,0.37,0.37}{\textit{#1}}}
\newcommand{\ConstantTok}[1]{\textcolor[rgb]{0.56,0.35,0.01}{#1}}
\newcommand{\ControlFlowTok}[1]{\textcolor[rgb]{0.00,0.23,0.31}{\textbf{#1}}}
\newcommand{\DataTypeTok}[1]{\textcolor[rgb]{0.68,0.00,0.00}{#1}}
\newcommand{\DecValTok}[1]{\textcolor[rgb]{0.68,0.00,0.00}{#1}}
\newcommand{\DocumentationTok}[1]{\textcolor[rgb]{0.37,0.37,0.37}{\textit{#1}}}
\newcommand{\ErrorTok}[1]{\textcolor[rgb]{0.68,0.00,0.00}{#1}}
\newcommand{\ExtensionTok}[1]{\textcolor[rgb]{0.00,0.23,0.31}{#1}}
\newcommand{\FloatTok}[1]{\textcolor[rgb]{0.68,0.00,0.00}{#1}}
\newcommand{\FunctionTok}[1]{\textcolor[rgb]{0.28,0.35,0.67}{#1}}
\newcommand{\ImportTok}[1]{\textcolor[rgb]{0.00,0.46,0.62}{#1}}
\newcommand{\InformationTok}[1]{\textcolor[rgb]{0.37,0.37,0.37}{#1}}
\newcommand{\KeywordTok}[1]{\textcolor[rgb]{0.00,0.23,0.31}{\textbf{#1}}}
\newcommand{\NormalTok}[1]{\textcolor[rgb]{0.00,0.23,0.31}{#1}}
\newcommand{\OperatorTok}[1]{\textcolor[rgb]{0.37,0.37,0.37}{#1}}
\newcommand{\OtherTok}[1]{\textcolor[rgb]{0.00,0.23,0.31}{#1}}
\newcommand{\PreprocessorTok}[1]{\textcolor[rgb]{0.68,0.00,0.00}{#1}}
\newcommand{\RegionMarkerTok}[1]{\textcolor[rgb]{0.00,0.23,0.31}{#1}}
\newcommand{\SpecialCharTok}[1]{\textcolor[rgb]{0.37,0.37,0.37}{#1}}
\newcommand{\SpecialStringTok}[1]{\textcolor[rgb]{0.13,0.47,0.30}{#1}}
\newcommand{\StringTok}[1]{\textcolor[rgb]{0.13,0.47,0.30}{#1}}
\newcommand{\VariableTok}[1]{\textcolor[rgb]{0.07,0.07,0.07}{#1}}
\newcommand{\VerbatimStringTok}[1]{\textcolor[rgb]{0.13,0.47,0.30}{#1}}
\newcommand{\WarningTok}[1]{\textcolor[rgb]{0.37,0.37,0.37}{\textit{#1}}}

\usepackage{longtable,booktabs,array}
\usepackage{calc} % for calculating minipage widths
% Correct order of tables after \paragraph or \subparagraph
\usepackage{etoolbox}
\makeatletter
\patchcmd\longtable{\par}{\if@noskipsec\mbox{}\fi\par}{}{}
\makeatother
% Allow footnotes in longtable head/foot
\IfFileExists{footnotehyper.sty}{\usepackage{footnotehyper}}{\usepackage{footnote}}
\makesavenoteenv{longtable}
\usepackage{graphicx}
\makeatletter
\newsavebox\pandoc@box
\newcommand*\pandocbounded[1]{% scales image to fit in text height/width
  \sbox\pandoc@box{#1}%
  \Gscale@div\@tempa{\textheight}{\dimexpr\ht\pandoc@box+\dp\pandoc@box\relax}%
  \Gscale@div\@tempb{\linewidth}{\wd\pandoc@box}%
  \ifdim\@tempb\p@<\@tempa\p@\let\@tempa\@tempb\fi% select the smaller of both
  \ifdim\@tempa\p@<\p@\scalebox{\@tempa}{\usebox\pandoc@box}%
  \else\usebox{\pandoc@box}%
  \fi%
}
% Set default figure placement to htbp
\def\fps@figure{htbp}
\makeatother


% definitions for citeproc citations
\NewDocumentCommand\citeproctext{}{}
\NewDocumentCommand\citeproc{mm}{%
  \begingroup\def\citeproctext{#2}\cite{#1}\endgroup}
\makeatletter
 % allow citations to break across lines
 \let\@cite@ofmt\@firstofone
 % avoid brackets around text for \cite:
 \def\@biblabel#1{}
 \def\@cite#1#2{{#1\if@tempswa , #2\fi}}
\makeatother
\newlength{\cslhangindent}
\setlength{\cslhangindent}{1.5em}
\newlength{\csllabelwidth}
\setlength{\csllabelwidth}{3em}
\newenvironment{CSLReferences}[2] % #1 hanging-indent, #2 entry-spacing
 {\begin{list}{}{%
  \setlength{\itemindent}{0pt}
  \setlength{\leftmargin}{0pt}
  \setlength{\parsep}{0pt}
  % turn on hanging indent if param 1 is 1
  \ifodd #1
   \setlength{\leftmargin}{\cslhangindent}
   \setlength{\itemindent}{-1\cslhangindent}
  \fi
  % set entry spacing
  \setlength{\itemsep}{#2\baselineskip}}}
 {\end{list}}
\usepackage{calc}
\newcommand{\CSLBlock}[1]{\hfill\break\parbox[t]{\linewidth}{\strut\ignorespaces#1\strut}}
\newcommand{\CSLLeftMargin}[1]{\parbox[t]{\csllabelwidth}{\strut#1\strut}}
\newcommand{\CSLRightInline}[1]{\parbox[t]{\linewidth - \csllabelwidth}{\strut#1\strut}}
\newcommand{\CSLIndent}[1]{\hspace{\cslhangindent}#1}



\setlength{\emergencystretch}{3em} % prevent overfull lines

\providecommand{\tightlist}{%
  \setlength{\itemsep}{0pt}\setlength{\parskip}{0pt}}



 


\KOMAoption{captions}{tableheading}
\makeatletter
\@ifpackageloaded{tcolorbox}{}{\usepackage[skins,breakable]{tcolorbox}}
\@ifpackageloaded{fontawesome5}{}{\usepackage{fontawesome5}}
\definecolor{quarto-callout-color}{HTML}{909090}
\definecolor{quarto-callout-note-color}{HTML}{0758E5}
\definecolor{quarto-callout-important-color}{HTML}{CC1914}
\definecolor{quarto-callout-warning-color}{HTML}{EB9113}
\definecolor{quarto-callout-tip-color}{HTML}{00A047}
\definecolor{quarto-callout-caution-color}{HTML}{FC5300}
\definecolor{quarto-callout-color-frame}{HTML}{acacac}
\definecolor{quarto-callout-note-color-frame}{HTML}{4582ec}
\definecolor{quarto-callout-important-color-frame}{HTML}{d9534f}
\definecolor{quarto-callout-warning-color-frame}{HTML}{f0ad4e}
\definecolor{quarto-callout-tip-color-frame}{HTML}{02b875}
\definecolor{quarto-callout-caution-color-frame}{HTML}{fd7e14}
\makeatother
\makeatletter
\@ifpackageloaded{bookmark}{}{\usepackage{bookmark}}
\makeatother
\makeatletter
\@ifpackageloaded{caption}{}{\usepackage{caption}}
\AtBeginDocument{%
\ifdefined\contentsname
  \renewcommand*\contentsname{Table of contents}
\else
  \newcommand\contentsname{Table of contents}
\fi
\ifdefined\listfigurename
  \renewcommand*\listfigurename{List of Figures}
\else
  \newcommand\listfigurename{List of Figures}
\fi
\ifdefined\listtablename
  \renewcommand*\listtablename{List of Tables}
\else
  \newcommand\listtablename{List of Tables}
\fi
\ifdefined\figurename
  \renewcommand*\figurename{Figure}
\else
  \newcommand\figurename{Figure}
\fi
\ifdefined\tablename
  \renewcommand*\tablename{Table}
\else
  \newcommand\tablename{Table}
\fi
}
\@ifpackageloaded{float}{}{\usepackage{float}}
\floatstyle{ruled}
\@ifundefined{c@chapter}{\newfloat{codelisting}{h}{lop}}{\newfloat{codelisting}{h}{lop}[chapter]}
\floatname{codelisting}{Listing}
\newcommand*\listoflistings{\listof{codelisting}{List of Listings}}
\makeatother
\makeatletter
\makeatother
\makeatletter
\@ifpackageloaded{caption}{}{\usepackage{caption}}
\@ifpackageloaded{subcaption}{}{\usepackage{subcaption}}
\makeatother
\usepackage{bookmark}
\IfFileExists{xurl.sty}{\usepackage{xurl}}{} % add URL line breaks if available
\urlstyle{same}
\hypersetup{
  pdftitle={Advanced Cognitive Modeling Notes},
  pdfauthor={Riccardo Fusaroli},
  colorlinks=true,
  linkcolor={blue},
  filecolor={Maroon},
  citecolor={Blue},
  urlcolor={Blue},
  pdfcreator={LaTeX via pandoc}}


\title{Advanced Cognitive Modeling Notes}
\author{Riccardo Fusaroli}
\date{2026-04-29}
\begin{document}
\maketitle

\renewcommand*\contentsname{Table of contents}
{
\hypersetup{linkcolor=}
\setcounter{tocdepth}{2}
\tableofcontents
}

\bookmarksetup{startatroot}

\chapter{Advanced Cognitive Modeling
Notes}\label{advanced-cognitive-modeling-notes}

\bookmarksetup{startatroot}

\chapter{Advanced Cognitive Modeling
Notes}\label{advanced-cognitive-modeling-notes-1}

These course notes support the Advanced Cognitive Modeling course taught
in the Master's program in Cognitive Science at Aarhus University. The
course represents a journey into how we can understand cognitive
processes through the formalization and implementation of hypothesized
mechanisms, their testing and validation.

\section{Course Philosophy and
Approach}\label{course-philosophy-and-approach}

Advanced cognitive modeling focuses on three interrelated objectives
that shape how we approach the modeling of cognitive processes:

The first objective centers on understanding the thought process behind
model development. Rather than simply providing a toolbox of existing
scripts, we explore how cognitive models are conceptualized and
constructed from the ground up. This approach ensures you develop the
skills to create novel models for unique research questions.

The second objective emphasizes mastering the Bayesian workflow
essential for robust model development. This workflow encompasses
simulation design, prior assessment, parameter recovery testing, and
thorough model fit evaluation. These skills ensure your models are not
just theoretically sound but also practically reliable and generalize
way beyond cognitive modeling.

The third objective focuses on developing advanced probabilistic
modeling capabilities. Through hands-on experience with Stan, you will
learn to implement increasingly sophisticated models while maintaining
scientific rigor.

\section{Course Structure and Learning
Path}\label{course-structure-and-learning-path}

The course follows a carefully structured progression that builds your
modeling capabilities step by step:

After a deepdive into the physics of pizza ovens, we begin with simple
scenarios that introduce fundamental modeling concepts. Each subsequent
chapter introduces new modeling techniques while building upon previous
knowledge. This cumulative approach ensures you develop a deep
understanding of both basic principles and advanced applications.

The chapters include theoretical discussions paired with practical
coding exercises. During practical sessions, we work with real datasets,
design models collaboratively, and implement them using modern
statistical tools. This hands-on approach provides ample opportunity for
questions and exploration.

The course schedule maintains flexibility to adapt to the collective
learning pace of each cohort. While we have clear learning objectives,
we ensure everyone develops a solid foundation before moving to more
advanced topics.

\section{Prerequisites and
Preparation}\label{prerequisites-and-preparation}

To make the most of this course, students should prepare their technical
environment and review fundamental concepts:

Software Requirements: - R (version 4.4 or above) - RStudio (version
2024.12.0 or above) - brms package with proper configuration - cmdstanr
package with complete installation

Technical Prerequisites: - Working knowledge of R programming - Basic
understanding of Bayesian statistics - Familiarity with cognitive
science fundamentals

Additional Resources: - Introduction to R and tidyverse:
https://r4ds.had.co.nz/ - A condensed Bayesian statistics primer (by
Chris Cox and me): https://4ccoxau.github.io/PriorsWorkshop/ - A
condensed primer on all my wisdom on priors:
https://youtu.be/kMSrRd4f2As

\section{Course Resources}\label{course-resources}

The course materials include: - Lecture notes and presentations -
Practical exercise guides - Example code and solutions - Additional
readings and references

For comprehensive information: - Course syllabus: {[}TBA{]} - Lecture
videos:
https://www.youtube.com/playlist?list=PL\_f3yDs5oZx5Ff-rRIZtghpE2tPfyAzMk

\section{About These Notes}\label{about-these-notes}

These notes represent an evolving resource that builds upon previous
iterations of the course while incorporating new developments in the
field. They are designed to serve both as a learning guide during the
course and as a reference for your future research endeavors.

\bookmarksetup{startatroot}

\chapter{The Pizza Experiment}\label{the-pizza-experiment}

\begin{quote}
\textbf{📍 Where we are in the Bayesian modeling workflow:} In these
course notes we will introduce and consistently use a Bayesian Workflow
Before we model cognition, we need to understand what modeling \emph{is}
--- building a mechanism-based account of a system, fitting it to data,
and validating it. Every move we make here (theory → math → code →
fitting → checking) will repeat in every chapter that follows, now
applied to minds instead of pizza stones. The full course roadmap is in
the section below.
\end{quote}

\section{From Pizza to Cognitive Models: An
Introduction}\label{from-pizza-to-cognitive-models-an-introduction}

This chapter introduces core modeling concepts through a delicious lens:
the physics of pizza stone heating. While this might seem far removed
from cognitive science, it provides an insightful introduction to the
challenges and methodologies of modeling complex phenomena.

\section{Why Start with Pizza?}\label{why-start-with-pizza}

Do I even need to answer that question? Because pizza, obviously.

Beyond the universal appeal of pizza, starting with a physical system
like a heating pizza stone offers a tangible way to grasp core
principles fundamental to cognitive modeling. Understanding human
cognition involves tackling immense complexity -- often dealing with
hidden mental states, noisy behavioral data, and competing theories.
Modeling a pizza stone allows us to isolate and understand key
methodological challenges in a more constrained environment.

First, it demonstrates how to bridge theory and math. We rely on
\emph{Prior Knowledge} --- using Newton's Law of Cooling to structure
our equations --- just as cognitive scientists use theories of memory or
attention to guide their models. This forces us to critically determine
our \emph{Levels of Analysis}: just as we simplify complex
thermodynamics into a basic cooling formula, cognitive scientists must
decide whether to model neurons, beliefs, or behavioral choices.

It also highlights the problem of \emph{Latent Variables}. Just as we
cannot ``see'' a heating coefficient---only the noisy temperatures it
produces---we cannot see a ``learning rate'' in the brain, only the
choices it drives.

Finally, this example acts as a sandbox for \emph{Model Complexity} and
\emph{Validation}. We will explicitly compare detailed physics-based
models against simpler statistical approximations to see which performs
better on empirical data. This mirrors the central trade-off in
cognitive science: balancing theoretical completeness with the practical
utility of generating testable predictions from available data.

\section{Learning Objectives}\label{learning-objectives}

This chapter uses the seemingly simple example of a heating pizza stone
to introduce fundamental principles of computational modeling that are
directly applicable to cognitive science. While the specific physics and
code are illustrative, the core goals are to understand the
\emph{process} of modeling.

By completing this chapter, you will be able to:

\begin{itemize}
\item
  \textbf{Appreciate Generative Modeling:} Understand the value of
  building models based on underlying mechanisms (like physics or
  cognitive processes) rather than just describing data patterns.
\item
  \textbf{Connect Theory and Data:} See how theoretical knowledge
  (physics laws) informs model structure and how data is used to
  estimate unknown model parameters (like the heating coefficient).
\item
  \textbf{Implement and Fit Models:} Gain initial experience
  implementing mathematical models in code (R and Stan) and using
  Bayesian inference (\texttt{brms}, \texttt{cmdstanr}) to fit them to
  empirical data.
\item
  \textbf{Evaluate Model Fit:} Learn basic techniques for visualizing
  model predictions against data and understand the limitations of
  purely statistical models compared to theory-driven ones.
\item
  \textbf{Recognize Modeling Trade-offs:} Understand the balance between
  model complexity, theoretical grounding, and practical utility through
  comparing statistical vs.~physics-based approaches.
\item
  \textbf{Apply Modeling Principles:} Recognize how the steps taken here
  -- data exploration, model building, parameter estimation, validation
  -- form a general workflow applicable to cognitive modeling problems
  explored later in the course.
\end{itemize}

Oh, and you'll probably get hungry as well!

Required Packages

\begin{Shaded}
\begin{Highlighting}[]
\CommentTok{\# 1. Ensure \textquotesingle{}pacman\textquotesingle{} is installed}
\ControlFlowTok{if}\NormalTok{ (}\SpecialCharTok{!}\FunctionTok{require}\NormalTok{(}\StringTok{"pacman"}\NormalTok{)) }\FunctionTok{install.packages}\NormalTok{(}\StringTok{"pacman"}\NormalTok{)}

\CommentTok{\# 2. Load packages (pacman will verify, install if missing, and load them)}
\NormalTok{pacman}\SpecialCharTok{::}\FunctionTok{p\_load}\NormalTok{(}
\NormalTok{  tidyverse,  }\CommentTok{\# Core package for data manipulation and visualization}
\NormalTok{  brms,       }\CommentTok{\# Interface for Bayesian regression models using Stan}
\NormalTok{  bayesplot,  }\CommentTok{\# Convenient functions for plotting Bayesian model results}
\NormalTok{  tidybayes,  }\CommentTok{\# Utilities for working with Bayesian model outputs in a tidy format}
\NormalTok{  cmdstanr,    }\CommentTok{\# Backend interface for running Stan models efficiently}
\NormalTok{  here        }\CommentTok{\# For constructing file paths relative to project root}
\NormalTok{)}

\CommentTok{\# Set to TRUE to rerun everything (takes time!), FALSE to load saved results.}
\NormalTok{regenerate\_simulations }\OtherTok{\textless{}{-}} \ConstantTok{FALSE} 
\end{Highlighting}
\end{Shaded}

\section{Part 1: Exploring the Pizza Stone Temperature
Data}\label{part-1-exploring-the-pizza-stone-temperature-data}

In this study, we collected temperature measurements from a pizza stone
in a gas-fired oven using an infrared temperature gun. Three different
raters (N, TR, and R) took measurements over time to track how the stone
heated up. Understanding how pizza stones heat up is crucial for
achieving the perfect pizza crust, as consistent and sufficient stone
temperature is essential for proper baking.

The measurements were taken as follows:

\begin{Shaded}
\begin{Highlighting}[]
\CommentTok{\# Load and examine the data}
\NormalTok{data }\OtherTok{\textless{}{-}} \FunctionTok{tibble}\NormalTok{(}
  \AttributeTok{Order =} \FunctionTok{rep}\NormalTok{(}\DecValTok{0}\SpecialCharTok{:}\DecValTok{18}\NormalTok{, }\DecValTok{3}\NormalTok{),}
  \AttributeTok{Seconds =} \FunctionTok{rep}\NormalTok{(}\FunctionTok{c}\NormalTok{(}\DecValTok{0}\NormalTok{, }\DecValTok{175}\NormalTok{, }\DecValTok{278}\NormalTok{, }\DecValTok{333}\NormalTok{, }\DecValTok{443}\NormalTok{, }\DecValTok{568}\NormalTok{, }\DecValTok{731}\NormalTok{, }\DecValTok{773}\NormalTok{, }\DecValTok{851}\NormalTok{, }\DecValTok{912}\NormalTok{, }\DecValTok{980}\NormalTok{, }
                  \DecValTok{1040}\NormalTok{, }\DecValTok{1074}\NormalTok{, }\DecValTok{1124}\NormalTok{, }\DecValTok{1175}\NormalTok{, }\DecValTok{1237}\NormalTok{, }\DecValTok{1298}\NormalTok{, }\DecValTok{1359}\NormalTok{, }\DecValTok{1394}\NormalTok{), }\DecValTok{3}\NormalTok{),}
  \AttributeTok{Temperature =} \FunctionTok{c}\NormalTok{(}\FloatTok{15.1}\NormalTok{, }\DecValTok{233}\NormalTok{, }\DecValTok{244}\NormalTok{, }\DecValTok{280}\NormalTok{, }\DecValTok{289}\NormalTok{, }\DecValTok{304}\NormalTok{, }\DecValTok{343}\NormalTok{, }\ConstantTok{NA}\NormalTok{, }\DecValTok{333}\NormalTok{, }\DecValTok{341}\NormalTok{, }\DecValTok{320}\NormalTok{, }
                  \DecValTok{370}\NormalTok{, }\DecValTok{325}\NormalTok{, }\DecValTok{362}\NormalTok{, }\DecValTok{363}\NormalTok{, }\DecValTok{357}\NormalTok{, }\DecValTok{380}\NormalTok{, }\DecValTok{376}\NormalTok{, }\DecValTok{380}\NormalTok{,}
                  \FloatTok{14.5}\NormalTok{, }\FloatTok{139.9}\NormalTok{, }\DecValTok{153}\NormalTok{, }\FloatTok{36.1}\NormalTok{, }\DecValTok{254}\NormalTok{, }\DecValTok{459}\NormalTok{, }\DecValTok{263}\NormalTok{, }\DecValTok{369}\NormalTok{, }\FunctionTok{rep}\NormalTok{(}\ConstantTok{NA}\NormalTok{, }\DecValTok{11}\NormalTok{),}
                  \FloatTok{12.9}\NormalTok{, }\FloatTok{149.5}\NormalTok{, }\DecValTok{159}\NormalTok{, }\FloatTok{179.4}\NormalTok{, }\FloatTok{191.7}\NormalTok{, }\DecValTok{201}\NormalTok{, }\DecValTok{210}\NormalTok{, }\ConstantTok{NA}\NormalTok{, }\DecValTok{256}\NormalTok{, }\DecValTok{257}\NormalTok{, }
                  \DecValTok{281}\NormalTok{, }\DecValTok{293}\NormalTok{, }\DecValTok{297}\NormalTok{, }\DecValTok{309}\NormalTok{, }\DecValTok{318}\NormalTok{, }\DecValTok{321}\NormalTok{, }\FunctionTok{rep}\NormalTok{(}\ConstantTok{NA}\NormalTok{, }\DecValTok{3}\NormalTok{)),}
  \AttributeTok{Rater =} \FunctionTok{rep}\NormalTok{(}\FunctionTok{c}\NormalTok{(}\StringTok{"N"}\NormalTok{, }\StringTok{"TR"}\NormalTok{, }\StringTok{"R"}\NormalTok{), }\AttributeTok{each =} \DecValTok{19}\NormalTok{)}
\NormalTok{)}



\CommentTok{\# Create summary statistics}
\NormalTok{summary\_stats }\OtherTok{\textless{}{-}}\NormalTok{ data }\SpecialCharTok{\%\textgreater{}\%}
  \FunctionTok{group\_by}\NormalTok{(Rater) }\SpecialCharTok{\%\textgreater{}\%}
  \FunctionTok{summarize}\NormalTok{(}
    \AttributeTok{n\_measurements =} \FunctionTok{sum}\NormalTok{(}\SpecialCharTok{!}\FunctionTok{is.na}\NormalTok{(Temperature)),}
    \AttributeTok{mean\_temp =} \FunctionTok{mean}\NormalTok{(Temperature, }\AttributeTok{na.rm =} \ConstantTok{TRUE}\NormalTok{),}
    \AttributeTok{sd\_temp =} \FunctionTok{sd}\NormalTok{(Temperature, }\AttributeTok{na.rm =} \ConstantTok{TRUE}\NormalTok{),}
    \AttributeTok{min\_temp =} \FunctionTok{min}\NormalTok{(Temperature, }\AttributeTok{na.rm =} \ConstantTok{TRUE}\NormalTok{),}
    \AttributeTok{max\_temp =} \FunctionTok{max}\NormalTok{(Temperature, }\AttributeTok{na.rm =} \ConstantTok{TRUE}\NormalTok{)}
\NormalTok{  )}

\CommentTok{\# Display summary statistics}
\NormalTok{knitr}\SpecialCharTok{::}\FunctionTok{kable}\NormalTok{(summary\_stats, }\AttributeTok{digits =} \DecValTok{1}\NormalTok{)}
\end{Highlighting}
\end{Shaded}

\begin{longtable}[]{@{}lrrrrr@{}}
\toprule\noalign{}
Rater & n\_measurements & mean\_temp & sd\_temp & min\_temp &
max\_temp \\
\midrule\noalign{}
\endhead
\bottomrule\noalign{}
\endlastfoot
N & 18 & 312.0 & 86.4 & 15.1 & 380 \\
R & 15 & 229.0 & 83.9 & 12.9 & 321 \\
TR & 8 & 211.1 & 155.2 & 14.5 & 459 \\
\end{longtable}

\subsection{Initial Data
Visualization}\label{initial-data-visualization}

Let's visualize how the temperature evolves over time for each rater:

\begin{Shaded}
\begin{Highlighting}[]
\FunctionTok{ggplot}\NormalTok{(data, }\FunctionTok{aes}\NormalTok{(}\AttributeTok{x =}\NormalTok{ Seconds}\SpecialCharTok{/}\DecValTok{60}\NormalTok{, }\AttributeTok{y =}\NormalTok{ Temperature, }\AttributeTok{color =}\NormalTok{ Rater)) }\SpecialCharTok{+}
  \FunctionTok{geom\_point}\NormalTok{(}\AttributeTok{size =} \DecValTok{3}\NormalTok{, }\AttributeTok{alpha =} \FloatTok{0.7}\NormalTok{) }\SpecialCharTok{+}
  \FunctionTok{geom\_line}\NormalTok{(}\AttributeTok{alpha =} \FloatTok{0.5}\NormalTok{) }\SpecialCharTok{+}
  \FunctionTok{labs}\NormalTok{(}
    \AttributeTok{title =} \StringTok{"Pizza Stone Temperature Evolution"}\NormalTok{,}
    \AttributeTok{subtitle =} \StringTok{"Measurements by three different raters"}\NormalTok{,}
    \AttributeTok{x =} \StringTok{"Time (minutes)"}\NormalTok{,}
    \AttributeTok{y =} \StringTok{"Temperature (°C)"}\NormalTok{,}
    \AttributeTok{color =} \StringTok{"Rater"}
\NormalTok{  ) }\SpecialCharTok{+}
  \FunctionTok{theme\_bw}\NormalTok{() }\SpecialCharTok{+}
  \FunctionTok{scale\_color\_brewer}\NormalTok{(}\AttributeTok{palette =} \StringTok{"Set1"}\NormalTok{)}
\end{Highlighting}
\end{Shaded}

\begin{center}
\includegraphics[width=0.8\linewidth,height=\textheight,keepaspectratio]{01-PizzaExperiment_files/figure-pdf/ch02_visualize_data-1.pdf}
\end{center}

\subsection{Key Observations}\label{key-observations}

Several interesting patterns emerge from our data:

\emph{Heating Patterns}: The temperature generally increases over time,
but not uniformly. We observe some fluctuations that might be due to:

\begin{itemize}
\tightlist
\item
  Variation in gas flame intensity
\item
  Different measurement locations on the stone
\item
  Measurement technique differences between raters
\end{itemize}

\emph{Measurement Patterns by Rater}

\begin{itemize}
\tightlist
\item
  Rater N maintained consistent measurements throughout the experiment
\item
  Rater TR shows more variability and fewer total measurements
\item
  Rater R shows a more gradual temperature increase pattern
\end{itemize}

\emph{Missing Data}: Some measurements are missing (NA values),
particularly in the later time points for Rater TR. This is common in
real-world data collection and needs to be considered in our analysis.

Let's examine the rate of temperature change:

\begin{Shaded}
\begin{Highlighting}[]
\CommentTok{\# {-}{-}{-} Calculate Rate of Temperature Change {-}{-}{-}}
\CommentTok{\# To understand how quickly the stone heats up, calculate the change rate per minute.}
\NormalTok{data\_with\_rate }\OtherTok{\textless{}{-}}\NormalTok{ data }\SpecialCharTok{\%\textgreater{}\%}
  \FunctionTok{group\_by}\NormalTok{(Rater) }\SpecialCharTok{\%\textgreater{}\%} \CommentTok{\# Perform calculation separately for each rater}
  \FunctionTok{arrange}\NormalTok{(Seconds) }\SpecialCharTok{\%\textgreater{}\%} \CommentTok{\# Ensure data is sorted chronologically for lag function}
  \FunctionTok{mutate}\NormalTok{(}
    \CommentTok{\# Calculate difference in temperature from previous measurement}
    \AttributeTok{delta\_temp =}\NormalTok{ Temperature }\SpecialCharTok{{-}} \FunctionTok{lag}\NormalTok{(Temperature),}
    \CommentTok{\# Calculate time elapsed since previous measurement}
    \AttributeTok{delta\_secs =}\NormalTok{ Seconds }\SpecialCharTok{{-}} \FunctionTok{lag}\NormalTok{(Seconds),}
    \CommentTok{\# Calculate rate in degrees C per second, then convert to degrees C per minute}
    \AttributeTok{temp\_change\_per\_min =}\NormalTok{ (delta\_temp }\SpecialCharTok{/}\NormalTok{ delta\_secs) }\SpecialCharTok{*} \DecValTok{60}\NormalTok{,}
    \CommentTok{\# Convert time to minutes for plotting}
    \AttributeTok{minutes =}\NormalTok{ Seconds }\SpecialCharTok{/} \DecValTok{60}
\NormalTok{  ) }\SpecialCharTok{\%\textgreater{}\%}
  \CommentTok{\# Remove rows where rate couldn\textquotesingle{}t be calculated (i.e., the first row for each rater)}
  \FunctionTok{filter}\NormalTok{(}\SpecialCharTok{!}\FunctionTok{is.na}\NormalTok{(temp\_change\_per\_min))}

\CommentTok{\# Visualize temperature change rate}
\FunctionTok{ggplot}\NormalTok{(data\_with\_rate, }\FunctionTok{aes}\NormalTok{(}\AttributeTok{x =}\NormalTok{ minutes, }\AttributeTok{y =}\NormalTok{ temp\_change\_per\_min, }\AttributeTok{color =}\NormalTok{ Rater)) }\SpecialCharTok{+}
  \FunctionTok{geom\_point}\NormalTok{() }\SpecialCharTok{+}
  \FunctionTok{geom\_line}\NormalTok{() }\SpecialCharTok{+}
  \FunctionTok{labs}\NormalTok{(}
    \AttributeTok{title =} \StringTok{"Rate of Temperature Change Over Time"}\NormalTok{,}
    \AttributeTok{subtitle =} \StringTok{"Degrees Celsius per minute"}\NormalTok{,}
    \AttributeTok{x =} \StringTok{"Time (minutes)"}\NormalTok{,}
    \AttributeTok{y =} \StringTok{"Temperature Change Rate (°C/min)"}\NormalTok{,}
    \AttributeTok{color =} \StringTok{"Rater"}
\NormalTok{  ) }\SpecialCharTok{+}
  \FunctionTok{theme\_bw}\NormalTok{() }\SpecialCharTok{+}
  \FunctionTok{scale\_color\_brewer}\NormalTok{(}\AttributeTok{palette =} \StringTok{"Set1"}\NormalTok{)}
\end{Highlighting}
\end{Shaded}

\begin{center}
\includegraphics[width=0.8\linewidth,height=\textheight,keepaspectratio]{01-PizzaExperiment_files/figure-pdf/ch02_temp_change_rate-1.pdf}
\end{center}

This visualization reveals that the heating rate is (likely) highest in
the first few minutes and gradually decreases as the stone temperature
approaches the oven temperature. This aligns with Newton's Law of
Cooling/Heating, which we will explore in the next section.

\section{Part 2: Initial Statistical
Modeling}\label{part-2-initial-statistical-modeling}

Before developing our physics-based model, let's explore how standard
statistical approaches perform in modeling our temperature data. We'll
implement two types of models using the brms package: a linear
mixed-effects model and a lognormal mixed-effects model. Both models
will account for variations between raters.

\subsection{Model Setup and Priors}\label{model-setup-and-priors}

First, let's ensure we have a directory for our models and set up our
computational parameters:

\begin{Shaded}
\begin{Highlighting}[]
\CommentTok{\# {-}{-}{-} Model Fitting Setup {-}{-}{-}}
\CommentTok{\# Define relative paths for saving model objects and Stan code.}
\NormalTok{linear\_model\_path }\OtherTok{\textless{}{-}} \FunctionTok{here}\NormalTok{(}\StringTok{"simmodels"}\NormalTok{, }\StringTok{"ch02\_pizza\_linear\_model"}\NormalTok{)}
\NormalTok{lognormal\_model\_path }\OtherTok{\textless{}{-}} \FunctionTok{here}\NormalTok{(}\StringTok{"simmodels"}\NormalTok{, }\StringTok{"ch02\_pizza\_lognormal\_model"}\NormalTok{)}
\NormalTok{stan\_model\_path }\OtherTok{\textless{}{-}} \FunctionTok{here}\NormalTok{(}\StringTok{"stan"}\NormalTok{, }\StringTok{"ch02\_pizza\_physics.stan"}\NormalTok{)}

\CommentTok{\# Create the required directories if they don\textquotesingle{}t already exist}
\ControlFlowTok{if}\NormalTok{ (}\SpecialCharTok{!}\FunctionTok{dir.exists}\NormalTok{(}\StringTok{"simmodels"}\NormalTok{)) }\FunctionTok{dir.create}\NormalTok{(}\StringTok{"simmodels"}\NormalTok{)}
\ControlFlowTok{if}\NormalTok{ (}\SpecialCharTok{!}\FunctionTok{dir.exists}\NormalTok{(}\StringTok{"stan"}\NormalTok{)) }\FunctionTok{dir.create}\NormalTok{(}\StringTok{"stan"}\NormalTok{)}

\CommentTok{\# Define computational parameters}
\NormalTok{mc\_settings }\OtherTok{\textless{}{-}} \FunctionTok{list}\NormalTok{(}
  \AttributeTok{chains =} \DecValTok{2}\NormalTok{,}
  \AttributeTok{iter =} \DecValTok{6000}\NormalTok{,}
  \AttributeTok{seed =} \DecValTok{123}\NormalTok{,}
  \AttributeTok{refresh =} \DecValTok{0}\NormalTok{,}
  \AttributeTok{backend =} \StringTok{"cmdstanr"}
\NormalTok{)}
\end{Highlighting}
\end{Shaded}

\subsection{Linear Mixed-Effects
Model}\label{linear-mixed-effects-model}

We begin with a linear mixed-effects model, which assumes that
temperature increases linearly with time but allows for different
patterns across raters. This model includes both fixed effects (overall
time trend) and random effects (rater-specific variations).

\begin{Shaded}
\begin{Highlighting}[]
\NormalTok{data\_scaled }\OtherTok{\textless{}{-}}\NormalTok{ data }\SpecialCharTok{\%\textgreater{}\%}
  \FunctionTok{mutate}\NormalTok{(}\AttributeTok{Minutes =}\NormalTok{ Seconds }\SpecialCharTok{/} \DecValTok{60}\NormalTok{)}

\CommentTok{\# Define priors for linear model}
\NormalTok{linear\_priors }\OtherTok{\textless{}{-}} \FunctionTok{c}\NormalTok{(}
  \FunctionTok{prior}\NormalTok{(}\FunctionTok{normal}\NormalTok{(}\DecValTok{15}\NormalTok{, }\DecValTok{20}\NormalTok{), }\AttributeTok{class =} \StringTok{"Intercept"}\NormalTok{),  }\CommentTok{\# Centered around room temperature}
  \FunctionTok{prior}\NormalTok{(}\FunctionTok{normal}\NormalTok{(}\DecValTok{0}\NormalTok{, }\DecValTok{20}\NormalTok{), }\AttributeTok{class =} \StringTok{"b"}\NormalTok{),          }\CommentTok{\# Expected temperature change per minute}
  \FunctionTok{prior}\NormalTok{(}\FunctionTok{normal}\NormalTok{(}\DecValTok{0}\NormalTok{, }\DecValTok{50}\NormalTok{), }\AttributeTok{class =} \StringTok{"sigma"}\NormalTok{),         }\CommentTok{\# Residual variation}
  \FunctionTok{prior}\NormalTok{(}\FunctionTok{normal}\NormalTok{(}\DecValTok{0}\NormalTok{, }\DecValTok{50}\NormalTok{), }\AttributeTok{class =} \StringTok{"sd"}\NormalTok{),            }\CommentTok{\# Random effects variation}
  \FunctionTok{prior}\NormalTok{(}\FunctionTok{lkj}\NormalTok{(}\DecValTok{3}\NormalTok{), }\AttributeTok{class =} \StringTok{"cor"}\NormalTok{)                    }\CommentTok{\# Random effects correlation}
\NormalTok{)}

\CommentTok{\# {-}{-}{-} Linear Model {-}{-}{-}}
\NormalTok{linear\_model }\OtherTok{\textless{}{-}} \FunctionTok{brm}\NormalTok{(}
\NormalTok{  Temperature }\SpecialCharTok{\textasciitilde{}}\NormalTok{ Minutes }\SpecialCharTok{+}\NormalTok{ (}\DecValTok{1} \SpecialCharTok{+}\NormalTok{ Minutes }\SpecialCharTok{|}\NormalTok{ Rater),}
  \AttributeTok{data =}\NormalTok{ data\_scaled,}
  \AttributeTok{family =}\NormalTok{ gaussian,}
  \AttributeTok{prior =}\NormalTok{ linear\_priors,}
  \AttributeTok{chains =}\NormalTok{ mc\_settings}\SpecialCharTok{$}\NormalTok{chains,}
  \AttributeTok{iter =}\NormalTok{ mc\_settings}\SpecialCharTok{$}\NormalTok{iter,}
  \AttributeTok{refresh =}\NormalTok{ mc\_settings}\SpecialCharTok{$}\NormalTok{refresh,}
  \AttributeTok{silent =} \DecValTok{2}\NormalTok{,}
  \AttributeTok{seed =}\NormalTok{ mc\_settings}\SpecialCharTok{$}\NormalTok{seed, }\CommentTok{\# Ensures reproducibility of MCMC sampling}
  \AttributeTok{backend =}\NormalTok{ mc\_settings}\SpecialCharTok{$}\NormalTok{backend,}
  \AttributeTok{file =}\NormalTok{ linear\_model\_path, }\CommentTok{\# Use relative path}
  \AttributeTok{file\_refit =} \FunctionTok{ifelse}\NormalTok{(regenerate\_simulations, }\StringTok{"always"}\NormalTok{, }\StringTok{"on\_change"}\NormalTok{), }\CommentTok{\# Refit if model code or data changes}
  \AttributeTok{cores =} \DecValTok{2}\NormalTok{, }\CommentTok{\# Specify cores used for parallel chains}
  \CommentTok{\# Increased adapt\_delta and max\_treedepth help prevent divergent transitions}
  \AttributeTok{adapt\_delta =} \FloatTok{0.99}\NormalTok{,}
  \AttributeTok{max\_treedepth =} \DecValTok{20}
\NormalTok{)}

\CommentTok{\# Display model summary}
\FunctionTok{summary}\NormalTok{(linear\_model)}
\end{Highlighting}
\end{Shaded}

\begin{verbatim}
 Family: gaussian 
  Links: mu = identity 
Formula: Temperature ~ Minutes + (1 + Minutes | Rater) 
   Data: data_scaled (Number of observations: 41) 
  Draws: 2 chains, each with iter = 6000; warmup = 3000; thin = 1;
         total post-warmup draws = 6000

Multilevel Hyperparameters:
~Rater (Number of levels: 3) 
                       Estimate Est.Error l-95% CI u-95% CI Rhat Bulk_ESS
sd(Intercept)             57.01     28.21     7.42   120.86 1.00     2014
sd(Minutes)               31.39     15.76    11.83    71.91 1.00     1655
cor(Intercept,Minutes)     0.00      0.35    -0.65     0.65 1.00     2075
                       Tail_ESS
sd(Intercept)              1104
sd(Minutes)                2789
cor(Intercept,Minutes)     3209

Regression Coefficients:
          Estimate Est.Error l-95% CI u-95% CI Rhat Bulk_ESS Tail_ESS
Intercept    84.92     45.64   -13.93   166.45 1.00     1809     2567
Minutes      -5.54      4.21   -13.23     3.41 1.00     1816     2506

Further Distributional Parameters:
      Estimate Est.Error l-95% CI u-95% CI Rhat Bulk_ESS Tail_ESS
sigma    57.45      6.87    45.63    72.47 1.00     3734     3283

Draws were sampled using sample(hmc). For each parameter, Bulk_ESS
and Tail_ESS are effective sample size measures, and Rhat is the potential
scale reduction factor on split chains (at convergence, Rhat = 1).
\end{verbatim}

\begin{Shaded}
\begin{Highlighting}[]
\CommentTok{\# Generate predictions}
\NormalTok{linear\_preds }\OtherTok{\textless{}{-}} \FunctionTok{fitted}\NormalTok{(}
\NormalTok{  linear\_model,}
  \AttributeTok{newdata =}\NormalTok{ data\_scaled,}
  \AttributeTok{probs =} \FunctionTok{c}\NormalTok{(}\FloatTok{0.025}\NormalTok{, }\FloatTok{0.975}\NormalTok{)}
\NormalTok{) }\SpecialCharTok{\%\textgreater{}\%}
  \FunctionTok{as\_tibble}\NormalTok{() }\SpecialCharTok{\%\textgreater{}\%}
  \FunctionTok{bind\_cols}\NormalTok{(data)}
\end{Highlighting}
\end{Shaded}

\subsection{Lognormal Mixed-Effects
Model}\label{lognormal-mixed-effects-model}

The lognormal model accounts for the fact that temperature changes might
be proportional rather than additive, and ensures predictions cannot go
below zero (I don't bring my oven out in the freezing cold!).

\begin{Shaded}
\begin{Highlighting}[]
\CommentTok{\# Define priors for lognormal model}
\NormalTok{lognormal\_priors }\OtherTok{\textless{}{-}} \FunctionTok{c}\NormalTok{(}
  \FunctionTok{prior}\NormalTok{(}\FunctionTok{normal}\NormalTok{(}\FloatTok{2.7}\NormalTok{, }\FloatTok{0.5}\NormalTok{), }\AttributeTok{class =} \StringTok{"Intercept"}\NormalTok{),     }\CommentTok{\# Log scale for room temperature}
  \FunctionTok{prior}\NormalTok{(}\FunctionTok{normal}\NormalTok{(}\DecValTok{0}\NormalTok{, }\FloatTok{0.3}\NormalTok{), }\AttributeTok{class =} \StringTok{"b"}\NormalTok{),     }\CommentTok{\# Expected log{-}scale change per minute}
  \FunctionTok{prior}\NormalTok{(}\FunctionTok{normal}\NormalTok{(}\DecValTok{0}\NormalTok{, }\DecValTok{1}\NormalTok{), }\AttributeTok{class =} \StringTok{"sigma"}\NormalTok{),        }\CommentTok{\# Log{-}scale residual variation}
  \FunctionTok{prior}\NormalTok{(}\FunctionTok{normal}\NormalTok{(}\DecValTok{0}\NormalTok{, }\FloatTok{0.5}\NormalTok{), }\AttributeTok{class =} \StringTok{"sd"}\NormalTok{),           }\CommentTok{\# Random effects variation}
  \FunctionTok{prior}\NormalTok{(}\FunctionTok{lkj}\NormalTok{(}\DecValTok{3}\NormalTok{), }\AttributeTok{class =} \StringTok{"cor"}\NormalTok{)                   }\CommentTok{\# Random effects correlation}
\NormalTok{)}

\CommentTok{\# {-}{-}{-} Lognormal Model {-}{-}{-}}
\NormalTok{lognormal\_model }\OtherTok{\textless{}{-}} \FunctionTok{brm}\NormalTok{(}
\NormalTok{  Temperature }\SpecialCharTok{\textasciitilde{}}\NormalTok{ Minutes }\SpecialCharTok{+}\NormalTok{ (}\DecValTok{1} \SpecialCharTok{+}\NormalTok{ Minutes }\SpecialCharTok{|}\NormalTok{ Rater),}
  \AttributeTok{data =}\NormalTok{ data\_scaled,}
  \AttributeTok{family =}\NormalTok{ lognormal, }\CommentTok{\# Use lognormal distribution}
  \AttributeTok{prior =}\NormalTok{ lognormal\_priors,}
  \AttributeTok{chains =}\NormalTok{ mc\_settings}\SpecialCharTok{$}\NormalTok{chains,}
  \AttributeTok{cores =} \DecValTok{2}\NormalTok{,}
  \AttributeTok{adapt\_delta =} \FloatTok{0.99}\NormalTok{,}
  \AttributeTok{max\_treedepth =} \DecValTok{20}\NormalTok{,}
  \AttributeTok{iter =}\NormalTok{ mc\_settings}\SpecialCharTok{$}\NormalTok{iter,}
  \AttributeTok{refresh =}\NormalTok{ mc\_settings}\SpecialCharTok{$}\NormalTok{refresh,}
  \AttributeTok{silent =} \DecValTok{2}\NormalTok{,}
  \AttributeTok{seed =}\NormalTok{ mc\_settings}\SpecialCharTok{$}\NormalTok{seed, }\CommentTok{\# Use the same seed for comparison if desired}
  \AttributeTok{backend =}\NormalTok{ mc\_settings}\SpecialCharTok{$}\NormalTok{backend,}
  \AttributeTok{file =}\NormalTok{ lognormal\_model\_path, }\CommentTok{\# Use relative path}
  \AttributeTok{file\_refit =} \FunctionTok{ifelse}\NormalTok{(regenerate\_simulations, }\StringTok{"always"}\NormalTok{, }\StringTok{"on\_change"}\NormalTok{)}
\NormalTok{)}

\CommentTok{\# Generate predictions}
\NormalTok{lognormal\_preds }\OtherTok{\textless{}{-}} \FunctionTok{fitted}\NormalTok{(}
\NormalTok{  lognormal\_model,}
  \AttributeTok{newdata =}\NormalTok{ data\_scaled,}
  \AttributeTok{probs =} \FunctionTok{c}\NormalTok{(}\FloatTok{0.025}\NormalTok{, }\FloatTok{0.975}\NormalTok{)}
\NormalTok{) }\SpecialCharTok{\%\textgreater{}\%}
  \FunctionTok{as\_tibble}\NormalTok{() }\SpecialCharTok{\%\textgreater{}\%}
  \FunctionTok{bind\_cols}\NormalTok{(data)}
\end{Highlighting}
\end{Shaded}

\subsection{Model Comparison and
Visualization}\label{model-comparison-and-visualization}

Let's compare how these models fit our data:

\begin{Shaded}
\begin{Highlighting}[]
\CommentTok{\# Compare models using LOO}
\NormalTok{model\_comparison }\OtherTok{\textless{}{-}} \FunctionTok{loo\_compare}\NormalTok{(}
  \FunctionTok{loo}\NormalTok{(linear\_model),}
  \FunctionTok{loo}\NormalTok{(lognormal\_model)}
\NormalTok{)}

\CommentTok{\# Create comparison plot}
\FunctionTok{ggplot}\NormalTok{() }\SpecialCharTok{+}
  \CommentTok{\# Raw data points}
  \FunctionTok{geom\_point}\NormalTok{(}\AttributeTok{data =}\NormalTok{ data, }
            \FunctionTok{aes}\NormalTok{(}\AttributeTok{x =}\NormalTok{ Seconds}\SpecialCharTok{/}\DecValTok{60}\NormalTok{, }\AttributeTok{y =}\NormalTok{ Temperature, }\AttributeTok{color =}\NormalTok{ Rater),}
            \AttributeTok{alpha =} \FloatTok{0.5}\NormalTok{) }\SpecialCharTok{+}
  \CommentTok{\# Linear model predictions}
  \FunctionTok{geom\_line}\NormalTok{(}\AttributeTok{data =}\NormalTok{ linear\_preds,}
            \FunctionTok{aes}\NormalTok{(}\AttributeTok{x =}\NormalTok{ Seconds}\SpecialCharTok{/}\DecValTok{60}\NormalTok{, }\AttributeTok{y =}\NormalTok{ Estimate, }\AttributeTok{linetype =} \StringTok{"Linear"}\NormalTok{),}
            \AttributeTok{color =} \StringTok{"blue"}\NormalTok{) }\SpecialCharTok{+}
  \FunctionTok{geom\_ribbon}\NormalTok{(}\AttributeTok{data =}\NormalTok{ linear\_preds,}
              \FunctionTok{aes}\NormalTok{(}\AttributeTok{x =}\NormalTok{ Seconds}\SpecialCharTok{/}\DecValTok{60}\NormalTok{, }\AttributeTok{ymin =}\NormalTok{ Q2}\FloatTok{.5}\NormalTok{, }\AttributeTok{ymax =}\NormalTok{ Q97}\FloatTok{.5}\NormalTok{),}
              \AttributeTok{fill =} \StringTok{"blue"}\NormalTok{, }\AttributeTok{alpha =} \FloatTok{0.1}\NormalTok{) }\SpecialCharTok{+}
  \CommentTok{\# Lognormal model predictions}
  \FunctionTok{geom\_line}\NormalTok{(}\AttributeTok{data =}\NormalTok{ lognormal\_preds,}
            \FunctionTok{aes}\NormalTok{(}\AttributeTok{x =}\NormalTok{ Seconds}\SpecialCharTok{/}\DecValTok{60}\NormalTok{, }\AttributeTok{y =}\NormalTok{ Estimate, }\AttributeTok{linetype =} \StringTok{"Lognormal"}\NormalTok{),}
            \AttributeTok{color =} \StringTok{"red"}\NormalTok{) }\SpecialCharTok{+}
  \FunctionTok{geom\_ribbon}\NormalTok{(}\AttributeTok{data =}\NormalTok{ lognormal\_preds,}
              \FunctionTok{aes}\NormalTok{(}\AttributeTok{x =}\NormalTok{ Seconds}\SpecialCharTok{/}\DecValTok{60}\NormalTok{, }\AttributeTok{ymin =}\NormalTok{ Q2}\FloatTok{.5}\NormalTok{, }\AttributeTok{ymax =}\NormalTok{ Q97}\FloatTok{.5}\NormalTok{),}
              \AttributeTok{fill =} \StringTok{"red"}\NormalTok{, }\AttributeTok{alpha =} \FloatTok{0.1}\NormalTok{) }\SpecialCharTok{+}
  \CommentTok{\# Formatting}
  \FunctionTok{facet\_wrap}\NormalTok{(}\SpecialCharTok{\textasciitilde{}}\NormalTok{Rater) }\SpecialCharTok{+}
  \FunctionTok{labs}\NormalTok{(}
    \AttributeTok{title =} \StringTok{"Comparison of Statistical Models"}\NormalTok{,}
    \AttributeTok{subtitle =} \StringTok{"Linear vs Lognormal Mixed{-}Effects Models"}\NormalTok{,}
    \AttributeTok{x =} \StringTok{"Time (minutes)"}\NormalTok{,}
    \AttributeTok{y =} \StringTok{"Temperature (°C)"}\NormalTok{,}
    \AttributeTok{linetype =} \StringTok{"Model Type"}
\NormalTok{  ) }\SpecialCharTok{+}
  \FunctionTok{theme\_bw}\NormalTok{()}
\end{Highlighting}
\end{Shaded}

\begin{center}
\includegraphics[width=0.8\linewidth,height=\textheight,keepaspectratio]{01-PizzaExperiment_files/figure-pdf/ch02_model_comparison-1.pdf}
\end{center}

\begin{Shaded}
\begin{Highlighting}[]
\CommentTok{\# Create comparison plot but capping the y axis}
\FunctionTok{ggplot}\NormalTok{() }\SpecialCharTok{+}
  \CommentTok{\# Raw data points}
  \FunctionTok{geom\_point}\NormalTok{(}\AttributeTok{data =}\NormalTok{ data, }
            \FunctionTok{aes}\NormalTok{(}\AttributeTok{x =}\NormalTok{ Seconds}\SpecialCharTok{/}\DecValTok{60}\NormalTok{, }\AttributeTok{y =}\NormalTok{ Temperature, }\AttributeTok{color =}\NormalTok{ Rater),}
            \AttributeTok{alpha =} \FloatTok{0.5}\NormalTok{) }\SpecialCharTok{+}
  \CommentTok{\# Linear model predictions}
  \FunctionTok{geom\_line}\NormalTok{(}\AttributeTok{data =}\NormalTok{ linear\_preds,}
            \FunctionTok{aes}\NormalTok{(}\AttributeTok{x =}\NormalTok{ Seconds}\SpecialCharTok{/}\DecValTok{60}\NormalTok{, }\AttributeTok{y =}\NormalTok{ Estimate, }\AttributeTok{linetype =} \StringTok{"Linear"}\NormalTok{),}
            \AttributeTok{color =} \StringTok{"blue"}\NormalTok{) }\SpecialCharTok{+}
  \FunctionTok{geom\_ribbon}\NormalTok{(}\AttributeTok{data =}\NormalTok{ linear\_preds,}
              \FunctionTok{aes}\NormalTok{(}\AttributeTok{x =}\NormalTok{ Seconds}\SpecialCharTok{/}\DecValTok{60}\NormalTok{, }\AttributeTok{ymin =}\NormalTok{ Q2}\FloatTok{.5}\NormalTok{, }\AttributeTok{ymax =}\NormalTok{ Q97}\FloatTok{.5}\NormalTok{),}
              \AttributeTok{fill =} \StringTok{"blue"}\NormalTok{, }\AttributeTok{alpha =} \FloatTok{0.1}\NormalTok{) }\SpecialCharTok{+}
  \CommentTok{\# Lognormal model predictions}
  \FunctionTok{geom\_line}\NormalTok{(}\AttributeTok{data =}\NormalTok{ lognormal\_preds,}
            \FunctionTok{aes}\NormalTok{(}\AttributeTok{x =}\NormalTok{ Seconds}\SpecialCharTok{/}\DecValTok{60}\NormalTok{, }\AttributeTok{y =}\NormalTok{ Estimate, }\AttributeTok{linetype =} \StringTok{"Lognormal"}\NormalTok{),}
            \AttributeTok{color =} \StringTok{"red"}\NormalTok{) }\SpecialCharTok{+}
  \FunctionTok{geom\_ribbon}\NormalTok{(}\AttributeTok{data =}\NormalTok{ lognormal\_preds,}
              \FunctionTok{aes}\NormalTok{(}\AttributeTok{x =}\NormalTok{ Seconds}\SpecialCharTok{/}\DecValTok{60}\NormalTok{, }\AttributeTok{ymin =}\NormalTok{ Q2}\FloatTok{.5}\NormalTok{, }\AttributeTok{ymax =}\NormalTok{ Q97}\FloatTok{.5}\NormalTok{),}
              \AttributeTok{fill =} \StringTok{"red"}\NormalTok{, }\AttributeTok{alpha =} \FloatTok{0.1}\NormalTok{) }\SpecialCharTok{+}
  \FunctionTok{ylim}\NormalTok{(}\DecValTok{0}\NormalTok{, }\DecValTok{1000}\NormalTok{) }\SpecialCharTok{+}
  \CommentTok{\# Formatting}
  \FunctionTok{facet\_wrap}\NormalTok{(}\SpecialCharTok{\textasciitilde{}}\NormalTok{Rater) }\SpecialCharTok{+}
  \FunctionTok{labs}\NormalTok{(}
    \AttributeTok{title =} \StringTok{"Comparison of Statistical Models"}\NormalTok{,}
    \AttributeTok{subtitle =} \StringTok{"Linear vs Lognormal Mixed{-}Effects Models"}\NormalTok{,}
    \AttributeTok{x =} \StringTok{"Time (minutes)"}\NormalTok{,}
    \AttributeTok{y =} \StringTok{"Temperature (°C)"}\NormalTok{,}
    \AttributeTok{linetype =} \StringTok{"Model Type"}
\NormalTok{  ) }\SpecialCharTok{+}
  \FunctionTok{theme\_bw}\NormalTok{()}
\end{Highlighting}
\end{Shaded}

\begin{center}
\includegraphics[width=0.8\linewidth,height=\textheight,keepaspectratio]{01-PizzaExperiment_files/figure-pdf/ch02_model_comparison-2.pdf}
\end{center}

\subsection{Model Assessment}\label{model-assessment}

I have seen worse models in my time, but they do seem to have important
issues:

\begin{itemize}
\item
  The linear mixed-effects model assumes a constant rate of temperature
  change, which we can see is not at all accurate. The actual
  temperature increase is fast at the beginning and appears to slow down
  over time, particularly at higher temperatures. While this model has
  the advantage of simplicity, it is not likely to produce accurate
  predictions as it seem to fail to capture the underlying physics of
  heat transfer.
\item
  The lognormal mixed-effects model is completely off.
\end{itemize}

Further, the lognormal model produce some divergences. In a Bayesian
workflow, divergences indicate that the Hamiltonian Monte Carlo sampler
is exploring pathological posterior geometry. Because purely statistical
models lack mechanistic asymptotes, they explore parameter spaces that
lead to mathematically unstable extrapolations. The core issue is that
neither model incorporates our knowledge of heat transfer physics: e.g
the flame temperature limiting the temperature the stone can reach and
the exponential approach to equilibrium temperature. This limitation
motivates our next section, where we'll develop a physics-based model.

\section{Part 3: Understanding the Physics
Model}\label{part-3-understanding-the-physics-model}

Temperature evolution in a pizza stone follows Newton's Law of
Cooling/Heating. We'll start by exploring this physical model before
applying it to real data.

\subsection{The Basic Temperature Evolution
Equation}\label{the-basic-temperature-evolution-equation}

The temperature evolution of a pizza stone in a gas-fired oven is
governed by the heat diffusion equation, which describes how heat flows
through solid materials:

\[\rho c_p \frac{\partial T}{\partial t} = k\nabla^2T + Q\]

where: \(\rho\) represents the stone's density (kg/m³) \(c_p\) denotes
specific heat capacity (J/kg·K) \(T\) is temperature (K) \(t\)
represents time (s) \(k\) is thermal conductivity (W/m·K) \(\nabla^2\)
is the Laplacian operator \(Q\) represents heat input from the oven
(W/m³)

While this equation provides a complete description of heat flow, we can
significantly simplify our analysis by applying the lumped capacitance
model. his simplification assumes that the stone conducts heat
internally much faster than it absorbs it from the surface. This allows
us to ignore the complex internal gradients (\(\nabla^2T\)) and focus
purely on the surface heat transfer, reducing our model to:

\[\frac{dT}{dt} = \frac{hA}{mc_p}(T_{\infty} - T)\]

where: \(h\) is the heat transfer coefficient (W/m²·K) \(A\) is the
surface area exposed to heat (m²) \(m\) is the stone's mass (kg)
\(T_{\infty}\) is the oven temperature (K)

This simplified equation relates the rate of temperature change to the
difference between the current stone temperature T and the ambient
temperature T∞ (as determined by the flame). The coefficient h
represents the heat transfer coefficient between the ambient and stone,
A is the stone's surface area exposed to heat, m is its mass, and cp
remains the specific heat capacity.

To solve this differential equation, we begin by separating variables:

\[\frac{dT}{T_{\infty} - T} = \left(\frac{hA}{mc_p}\right)dt\]

Integration of both sides yields:

\[-\ln|T_{\infty} - T| = \left(\frac{hA}{mc_p}\right)t + C\]

where C is an integration constant.

Using the initial condition \(T = T_i\) at \(t = 0\), we can determine
the integration constant:

\[C = -\ln|T_{\infty} - T_i|\]

Substituting this back and solving for temperature gives us:

\[T = T_{\infty} + (T_i - T_{\infty})\exp\left(-\frac{hA}{mc_p}t\right)\]

For practical reasons, we combine physical parameters into a single
coefficient, which we will call HOT:

\[HOT = \frac{hA}{mc_p}\] Giving our working equation:
\[T = T_{\infty} + (T_i - T_{\infty})\exp(-HOT \cdot t)\]

This equation retains the essential physics while providing a practical
model for analyzing our experimental data. The HOT coefficient
encapsulates the combined effects of heat transfer efficiency, stone
geometry, and material properties into a single parameter that
determines how quickly the stone approaches the flame temperature.

\section{Part 4: Implementing the Physics-Based
Model}\label{part-4-implementing-the-physics-based-model}

Having established the theoretical foundation, we now face a practical
challenge: we cannot measure \(h\), \(A\), or \(c_p\) directly during
the bake. Instead, we must infer the aggregate parameter HOT from the
noisy temperature data collected by our raters. This is where Stan comes
in, enabling us to create a Bayesian implementation of our physics-based
model, and to account for measurement uncertainty and variation between
raters. First, we prepare our data for the Stan model. Our model
requires initial temperatures, time measurements, and observed
temperatures from each rater:

\begin{Shaded}
\begin{Highlighting}[]
\CommentTok{\# Create data structure for Stan}
\NormalTok{stan\_data }\OtherTok{\textless{}{-}} \FunctionTok{list}\NormalTok{(}
  \AttributeTok{N =} \FunctionTok{nrow}\NormalTok{(data }\SpecialCharTok{\%\textgreater{}\%} \FunctionTok{filter}\NormalTok{(}\SpecialCharTok{!}\FunctionTok{is.na}\NormalTok{(Temperature))),}
  \AttributeTok{time =}\NormalTok{ data }\SpecialCharTok{\%\textgreater{}\%} \FunctionTok{filter}\NormalTok{(}\SpecialCharTok{!}\FunctionTok{is.na}\NormalTok{(Temperature)) }\SpecialCharTok{\%\textgreater{}\%} \FunctionTok{pull}\NormalTok{(Seconds),}
  \AttributeTok{temp =}\NormalTok{ data }\SpecialCharTok{\%\textgreater{}\%} \FunctionTok{filter}\NormalTok{(}\SpecialCharTok{!}\FunctionTok{is.na}\NormalTok{(Temperature)) }\SpecialCharTok{\%\textgreater{}\%} \FunctionTok{pull}\NormalTok{(Temperature),}
  \AttributeTok{n\_raters =} \DecValTok{3}\NormalTok{,}
  \AttributeTok{rater =} \FunctionTok{as.numeric}\NormalTok{(}\FunctionTok{factor}\NormalTok{(data }\SpecialCharTok{\%\textgreater{}\%} 
                           \FunctionTok{filter}\NormalTok{(}\SpecialCharTok{!}\FunctionTok{is.na}\NormalTok{(Temperature)) }\SpecialCharTok{\%\textgreater{}\%} 
                           \FunctionTok{pull}\NormalTok{(Rater))),}
  \AttributeTok{Ti =} \FunctionTok{c}\NormalTok{(}\FloatTok{15.1}\NormalTok{, }\FloatTok{14.5}\NormalTok{, }\FloatTok{12.9}\NormalTok{),  }\CommentTok{\# Initial temperature estimates}
  \AttributeTok{Tinf =} \DecValTok{450}              \CommentTok{\# Flame temperature estimate}
\NormalTok{)}
\end{Highlighting}
\end{Shaded}

Next, we implement our physics-based model in Stan. The model
incorporates our derived equation while allowing for rater-specific
heating coefficients:

\begin{Shaded}
\begin{Highlighting}[]
\NormalTok{stan\_code }\OtherTok{\textless{}{-}} \StringTok{"}
\StringTok{data \{}
\StringTok{  int\textless{}lower=0\textgreater{} N;                   // Number of observations}
\StringTok{  vector[N] time;                   // Time points}
\StringTok{  vector[N] temp;                   // Observed temperatures}
\StringTok{  int\textless{}lower=0\textgreater{} n\_raters;           // Number of raters}
\StringTok{  array[N] int\textless{}lower=1,upper=n\_raters\textgreater{} rater;  // Rater indices}
\StringTok{  vector[n\_raters] Ti;             // Initial temperatures}
\StringTok{  real Tinf;                       // Flame temperature}
\StringTok{\}}

\StringTok{parameters \{}
\StringTok{  vector\textless{}lower=0\textgreater{}[n\_raters] HOT;   // Heating coefficients}
\StringTok{  vector\textless{}lower=0\textgreater{}[n\_raters] sigma; // Measurement error}
\StringTok{\}}

\StringTok{model \{}
\StringTok{  // Physics{-}based temperature prediction}
\StringTok{  vector[N] mu = Tinf + (Ti[rater] {-} Tinf) .* exp({-}HOT[rater] .* time);}
\StringTok{  }
\StringTok{  // Prior distributions}
\StringTok{  target += normal\_lpdf(HOT | 0.005, 0.005);    // Prior for heating rate}
\StringTok{  target += exponential\_lpdf(sigma | 1);         // Prior for measurement error}
\StringTok{  }
\StringTok{  // Likelihood}
\StringTok{  target += normal\_lpdf(temp | mu, sigma[rater]);}
\StringTok{\}}
\StringTok{"}

\CommentTok{\# Save the Stan model code using the relative path}
\FunctionTok{writeLines}\NormalTok{(stan\_code, stan\_model\_path)}

\CommentTok{\# Compile and fit the Stan model using the relative path}
\CommentTok{\# Compilation happens once; subsequent runs load the compiled object if unchanged.}
\NormalTok{mod }\OtherTok{\textless{}{-}} \FunctionTok{cmdstan\_model}\NormalTok{(}\StringTok{"stan/ch02\_pizza\_physics.stan"}\NormalTok{, }\AttributeTok{quiet =} \ConstantTok{TRUE}\NormalTok{)}
\end{Highlighting}
\end{Shaded}

\subsection{Prior Predictive Check}\label{prior-predictive-check}

Before fitting our model to empirical data, we must ensure our priors
generate physically plausible heating curves (Simulation Validation). If
our \texttt{normal(0.005,\ 0.005)} prior on \texttt{HOT} is appropriate,
simulated temperatures should gradually approach \(T_{\infty}\) without
exceeding it or dropping below absolute zero.

\begin{Shaded}
\begin{Highlighting}[]
\CommentTok{\# Simulate HOT parameters from the prior}
\FunctionTok{set.seed}\NormalTok{(}\DecValTok{123}\NormalTok{)}
\NormalTok{prior\_hot }\OtherTok{\textless{}{-}} \FunctionTok{rnorm}\NormalTok{(}\DecValTok{100}\NormalTok{, }\AttributeTok{mean =} \FloatTok{0.005}\NormalTok{, }\AttributeTok{sd =} \FloatTok{0.005}\NormalTok{)}
\NormalTok{prior\_hot }\OtherTok{\textless{}{-}}\NormalTok{ prior\_hot[prior\_hot }\SpecialCharTok{\textgreater{}} \DecValTok{0}\NormalTok{] }\CommentTok{\# Truncate at 0 as per Stan lower=0 bound}

\CommentTok{\# Simulate curves}
\NormalTok{time\_seq }\OtherTok{\textless{}{-}} \FunctionTok{seq}\NormalTok{(}\DecValTok{0}\NormalTok{, }\FunctionTok{max}\NormalTok{(stan\_data}\SpecialCharTok{$}\NormalTok{time), }\AttributeTok{length.out =} \DecValTok{100}\NormalTok{)}
\NormalTok{prior\_preds }\OtherTok{\textless{}{-}} \FunctionTok{expand\_grid}\NormalTok{(}\AttributeTok{sim\_id =} \DecValTok{1}\SpecialCharTok{:}\FunctionTok{length}\NormalTok{(prior\_hot), }\AttributeTok{time =}\NormalTok{ time\_seq) }\SpecialCharTok{\%\textgreater{}\%}
  \FunctionTok{mutate}\NormalTok{(}
    \AttributeTok{hot =}\NormalTok{ prior\_hot[sim\_id],}
    \AttributeTok{temp =}\NormalTok{ stan\_data}\SpecialCharTok{$}\NormalTok{Tinf }\SpecialCharTok{+}\NormalTok{ (}\FloatTok{15.0} \SpecialCharTok{{-}}\NormalTok{ stan\_data}\SpecialCharTok{$}\NormalTok{Tinf) }\SpecialCharTok{*} \FunctionTok{exp}\NormalTok{(}\SpecialCharTok{{-}}\NormalTok{hot }\SpecialCharTok{*}\NormalTok{ time)}
\NormalTok{  )}

\FunctionTok{ggplot}\NormalTok{(prior\_preds, }\FunctionTok{aes}\NormalTok{(}\AttributeTok{x =}\NormalTok{ time}\SpecialCharTok{/}\DecValTok{60}\NormalTok{, }\AttributeTok{y =}\NormalTok{ temp, }\AttributeTok{group =}\NormalTok{ sim\_id)) }\SpecialCharTok{+}
  \FunctionTok{geom\_line}\NormalTok{(}\AttributeTok{alpha =} \FloatTok{0.2}\NormalTok{, }\AttributeTok{color =} \StringTok{"blue"}\NormalTok{) }\SpecialCharTok{+}
  \FunctionTok{labs}\NormalTok{(}\AttributeTok{title =} \StringTok{"Prior Predictive Check: Heating Curves"}\NormalTok{,}
       \AttributeTok{x =} \StringTok{"Time (minutes)"}\NormalTok{, }\AttributeTok{y =} \StringTok{"Temperature (°C)"}\NormalTok{) }\SpecialCharTok{+}
  \FunctionTok{theme\_bw}\NormalTok{()}
\end{Highlighting}
\end{Shaded}

\begin{center}
\includegraphics[width=0.8\linewidth,height=\textheight,keepaspectratio]{01-PizzaExperiment_files/figure-pdf/ch02_prior_predictive-1.pdf}
\end{center}

\begin{Shaded}
\begin{Highlighting}[]
\NormalTok{fit\_physics\_path }\OtherTok{\textless{}{-}} \FunctionTok{here}\NormalTok{(}\StringTok{"simmodels"}\NormalTok{, }\StringTok{"ch02\_pizza\_physics\_fit.rds"}\NormalTok{)}
\ControlFlowTok{if}\NormalTok{ (regenerate\_simulations }\SpecialCharTok{||} \SpecialCharTok{!}\FunctionTok{file.exists}\NormalTok{(fit\_physics\_path)) \{}
\NormalTok{  fit }\OtherTok{\textless{}{-}}\NormalTok{ mod}\SpecialCharTok{$}\FunctionTok{sample}\NormalTok{(}
    \AttributeTok{data =}\NormalTok{ stan\_data,}
    \AttributeTok{seed =} \DecValTok{123}\NormalTok{,}
    \AttributeTok{chains =} \DecValTok{2}\NormalTok{,}
    \AttributeTok{parallel\_chains =} \DecValTok{2}\NormalTok{,}
    \AttributeTok{refresh =} \DecValTok{0}\NormalTok{,}
    \AttributeTok{show\_messages =} \ConstantTok{FALSE}
\NormalTok{  )}
\NormalTok{  fit}\SpecialCharTok{$}\FunctionTok{save\_object}\NormalTok{(fit\_physics\_path)}
\NormalTok{\} }\ControlFlowTok{else}\NormalTok{ \{}
\NormalTok{  fit }\OtherTok{\textless{}{-}} \FunctionTok{readRDS}\NormalTok{(fit\_physics\_path)}
\NormalTok{\}}
\end{Highlighting}
\end{Shaded}

The Stan implementation translates our mathematical model into a
computational framework. We assign informative priors to our parameters
based on physical understanding: the heating coefficient (HOT) is
expected to be small but positive, while measurement error (sigma)
follows an exponential distribution to ensure positivity while allowing
for varying levels of uncertainty between raters. To visualize our
model's predictions and assess its performance, we extract posterior
samples and generate predictions across our time range:

\begin{Shaded}
\begin{Highlighting}[]
\CommentTok{\# Extract draws}
\NormalTok{post }\OtherTok{\textless{}{-}} \FunctionTok{as\_draws\_df}\NormalTok{(fit}\SpecialCharTok{$}\FunctionTok{draws}\NormalTok{()) }\SpecialCharTok{\%\textgreater{}\%}
\NormalTok{  dplyr}\SpecialCharTok{::}\FunctionTok{select}\NormalTok{(}\FunctionTok{starts\_with}\NormalTok{(}\StringTok{"HOT"}\NormalTok{), }\FunctionTok{starts\_with}\NormalTok{(}\StringTok{"sigma"}\NormalTok{)) }\SpecialCharTok{\%\textgreater{}\%}
  \FunctionTok{slice\_sample}\NormalTok{(}\AttributeTok{n =} \DecValTok{100}\NormalTok{)}

\CommentTok{\# Create prediction grid}
\NormalTok{pred\_data }\OtherTok{\textless{}{-}}\NormalTok{ tidyr}\SpecialCharTok{::}\FunctionTok{crossing}\NormalTok{(}
  \AttributeTok{time =} \FunctionTok{seq}\NormalTok{(}\DecValTok{0}\NormalTok{, }\FunctionTok{max}\NormalTok{(stan\_data}\SpecialCharTok{$}\NormalTok{time), }\AttributeTok{length.out =} \DecValTok{100}\NormalTok{),}
  \AttributeTok{rater =} \DecValTok{1}\SpecialCharTok{:}\NormalTok{stan\_data}\SpecialCharTok{$}\NormalTok{n\_raters}
\NormalTok{) }\SpecialCharTok{\%\textgreater{}\%}
  \FunctionTok{mutate}\NormalTok{(}
    \AttributeTok{Ti =}\NormalTok{ stan\_data}\SpecialCharTok{$}\NormalTok{Ti[rater],}
    \AttributeTok{Tinf =}\NormalTok{ stan\_data}\SpecialCharTok{$}\NormalTok{Tinf}
\NormalTok{  )}

\CommentTok{\# Generate predictions}
\NormalTok{pred\_matrix }\OtherTok{\textless{}{-}} \FunctionTok{matrix}\NormalTok{(}\ConstantTok{NA}\NormalTok{, }\AttributeTok{nrow =} \FunctionTok{nrow}\NormalTok{(pred\_data), }\AttributeTok{ncol =} \DecValTok{100}\NormalTok{)}
\ControlFlowTok{for}\NormalTok{ (i }\ControlFlowTok{in} \DecValTok{1}\SpecialCharTok{:}\FunctionTok{nrow}\NormalTok{(pred\_data)) \{}
\NormalTok{  pred\_matrix[i,] }\OtherTok{\textless{}{-}} \FunctionTok{with}\NormalTok{(pred\_data[i,], }
\NormalTok{                         Tinf }\SpecialCharTok{+}\NormalTok{ (Ti }\SpecialCharTok{{-}}\NormalTok{ Tinf) }\SpecialCharTok{*} \FunctionTok{exp}\NormalTok{(}\SpecialCharTok{{-}}\FunctionTok{as.matrix}\NormalTok{(post)[,rater] }\SpecialCharTok{*}\NormalTok{ time))}
\NormalTok{\}}

\CommentTok{\# Summarize predictions}
\NormalTok{predictions }\OtherTok{\textless{}{-}}\NormalTok{ pred\_data }\SpecialCharTok{\%\textgreater{}\%}
  \FunctionTok{mutate}\NormalTok{(}
    \AttributeTok{mean =} \FunctionTok{rowMeans}\NormalTok{(pred\_matrix),}
    \AttributeTok{lower =} \FunctionTok{apply}\NormalTok{(pred\_matrix, }\DecValTok{1}\NormalTok{, quantile, }\FloatTok{0.025}\NormalTok{),}
    \AttributeTok{upper =} \FunctionTok{apply}\NormalTok{(pred\_matrix, }\DecValTok{1}\NormalTok{, quantile, }\FloatTok{0.975}\NormalTok{)}
\NormalTok{  )}

\CommentTok{\# Create visualization}
\FunctionTok{ggplot}\NormalTok{(predictions, }\FunctionTok{aes}\NormalTok{(}\AttributeTok{x =}\NormalTok{ time}\SpecialCharTok{/}\DecValTok{60}\NormalTok{)) }\SpecialCharTok{+}
  \FunctionTok{geom\_ribbon}\NormalTok{(}\FunctionTok{aes}\NormalTok{(}\AttributeTok{ymin =}\NormalTok{ lower, }\AttributeTok{ymax =}\NormalTok{ upper), }\AttributeTok{alpha =} \FloatTok{0.2}\NormalTok{) }\SpecialCharTok{+}
  \FunctionTok{geom\_line}\NormalTok{(}\FunctionTok{aes}\NormalTok{(}\AttributeTok{y =}\NormalTok{ mean)) }\SpecialCharTok{+}
  \FunctionTok{geom\_point}\NormalTok{(}
    \AttributeTok{data =}\NormalTok{ data }\SpecialCharTok{\%\textgreater{}\%} 
      \FunctionTok{filter}\NormalTok{(}\SpecialCharTok{!}\FunctionTok{is.na}\NormalTok{(Temperature)) }\SpecialCharTok{\%\textgreater{}\%}
      \FunctionTok{mutate}\NormalTok{(}\AttributeTok{rater =} \FunctionTok{case\_when}\NormalTok{(}
\NormalTok{        Rater }\SpecialCharTok{==} \StringTok{"N"} \SpecialCharTok{\textasciitilde{}} \DecValTok{1}\NormalTok{,}
\NormalTok{        Rater }\SpecialCharTok{==} \StringTok{"TR"} \SpecialCharTok{\textasciitilde{}} \DecValTok{2}\NormalTok{,}
\NormalTok{        Rater }\SpecialCharTok{==} \StringTok{"R"} \SpecialCharTok{\textasciitilde{}} \DecValTok{3}
\NormalTok{      )), }
    \FunctionTok{aes}\NormalTok{(}\AttributeTok{x =}\NormalTok{ Seconds}\SpecialCharTok{/}\DecValTok{60}\NormalTok{, }\AttributeTok{y =}\NormalTok{ Temperature)}
\NormalTok{  ) }\SpecialCharTok{+}
  \FunctionTok{facet\_wrap}\NormalTok{(}\SpecialCharTok{\textasciitilde{}}\NormalTok{rater, }\AttributeTok{labeller =} \FunctionTok{labeller}\NormalTok{(}\AttributeTok{rater =} \FunctionTok{c}\NormalTok{(}
    \StringTok{"1"} \OtherTok{=} \StringTok{"Rater N"}\NormalTok{, }
    \StringTok{"2"} \OtherTok{=} \StringTok{"Rater TR"}\NormalTok{, }
    \StringTok{"3"} \OtherTok{=} \StringTok{"Rater R"}
\NormalTok{  ))) }\SpecialCharTok{+}
  \FunctionTok{labs}\NormalTok{(}
    \AttributeTok{title =} \StringTok{"Physics{-}Based Model Predictions"}\NormalTok{,}
    \AttributeTok{x =} \StringTok{"Time (minutes)"}\NormalTok{,}
    \AttributeTok{y =} \StringTok{"Temperature (°C)"}
\NormalTok{  ) }\SpecialCharTok{+}
  \FunctionTok{theme\_bw}\NormalTok{()}
\end{Highlighting}
\end{Shaded}

\begin{center}
\includegraphics[width=0.8\linewidth,height=\textheight,keepaspectratio]{01-PizzaExperiment_files/figure-pdf/ch02_physics_model_preds-1.pdf}
\end{center}

Our implementation combines the theoretical understanding developed in
Part 3 with practical considerations for real-world data analysis. The
model accounts for measurement uncertainty while maintaining the
fundamental physics of heat transfer, providing a robust framework for
understanding pizza stone temperature evolution.

\section{Part 5: Model Analysis and Practical
Applications}\label{part-5-model-analysis-and-practical-applications}

Having fitted our physics-based model, we must first verify that it
actually describes the data well (Validation) before we can use it to
make predictions. Once we trust our estimates of the HOT parameter, we
can develop practical insights for pizza stone management. A key
question for any pizzaiolo is: How long must I wait before the stone is
ready? To answer this, we don't need a new model; we simply need to look
at our physics equation from a different angle. By rearranging our
original cooling law to solve for time (\(t\)) instead of temperature
(\(T\)), we create a predictive tool:

\[t = -\frac{1}{HOT} \ln\left( \frac{T_{\text{target}} - T_{\infty}}{T_{\text{initial}} - T_{\infty}} \right)\]

We will now implement this function in R to calculate precise waiting
times under different conditions.

\begin{Shaded}
\begin{Highlighting}[]
\NormalTok{time\_to\_temp }\OtherTok{\textless{}{-}} \ControlFlowTok{function}\NormalTok{(target\_temp, HOT, Ti, Tinf) \{}
  \CommentTok{\# Solve: target = Tinf + (Ti {-} Tinf) * exp({-}HOT * t)}
  \CommentTok{\# for t}
\NormalTok{  t }\OtherTok{=} \SpecialCharTok{{-}}\DecValTok{1}\SpecialCharTok{/}\NormalTok{HOT }\SpecialCharTok{*} \FunctionTok{log}\NormalTok{((target\_temp }\SpecialCharTok{{-}}\NormalTok{ Tinf)}\SpecialCharTok{/}\NormalTok{(Ti }\SpecialCharTok{{-}}\NormalTok{ Tinf))}
  \FunctionTok{return}\NormalTok{(t}\SpecialCharTok{/}\DecValTok{60}\NormalTok{)  }\CommentTok{\# Convert seconds to minutes}
\NormalTok{\}}
\end{Highlighting}
\end{Shaded}

To understand heating times across different oven conditions, we examine
how varying flame temperatures affect the time needed to reach
pizza-making temperatures. We extract the heating coefficients from our
fitted model and analyze temperature scenarios:

\begin{Shaded}
\begin{Highlighting}[]
\CommentTok{\# Extract HOT samples from our posterior}
\NormalTok{hot\_samples }\OtherTok{\textless{}{-}} \FunctionTok{as\_draws\_df}\NormalTok{(fit}\SpecialCharTok{$}\FunctionTok{draws}\NormalTok{()) }\SpecialCharTok{\%\textgreater{}\%}
\NormalTok{  dplyr}\SpecialCharTok{::}\FunctionTok{select}\NormalTok{(}\FunctionTok{starts\_with}\NormalTok{(}\StringTok{"HOT"}\NormalTok{))}

\CommentTok{\# Create prediction grid for different flame temperatures}
\NormalTok{pred\_data }\OtherTok{\textless{}{-}}\NormalTok{ tidyr}\SpecialCharTok{::}\FunctionTok{crossing}\NormalTok{(}
  \AttributeTok{Tinf =} \FunctionTok{seq}\NormalTok{(}\DecValTok{450}\NormalTok{, }\DecValTok{1200}\NormalTok{, }\AttributeTok{by =} \DecValTok{50}\NormalTok{),  }\CommentTok{\# Range of flame temperatures}
  \AttributeTok{rater =} \DecValTok{1}\SpecialCharTok{:}\DecValTok{3}
\NormalTok{) }\SpecialCharTok{\%\textgreater{}\%}
  \FunctionTok{mutate}\NormalTok{(}
    \AttributeTok{Ti =}\NormalTok{ stan\_data}\SpecialCharTok{$}\NormalTok{Ti[rater],}
    \AttributeTok{target\_temp =} \DecValTok{400}  \CommentTok{\# Target temperature for pizza cooking}
\NormalTok{  )}

\CommentTok{\# Calculate heating times across conditions}
\NormalTok{n\_samples }\OtherTok{\textless{}{-}} \DecValTok{100}
\NormalTok{time\_preds }\OtherTok{\textless{}{-}} \FunctionTok{map\_dfr}\NormalTok{(}\DecValTok{1}\SpecialCharTok{:}\FunctionTok{nrow}\NormalTok{(pred\_data), }\ControlFlowTok{function}\NormalTok{(i) \{}
\NormalTok{  times }\OtherTok{\textless{}{-}} \FunctionTok{sapply}\NormalTok{(}\DecValTok{1}\SpecialCharTok{:}\NormalTok{n\_samples, }\ControlFlowTok{function}\NormalTok{(j) \{}
\NormalTok{    hot }\OtherTok{\textless{}{-}}\NormalTok{ hot\_samples[j, }\FunctionTok{paste0}\NormalTok{(}\StringTok{"HOT["}\NormalTok{, pred\_data}\SpecialCharTok{$}\NormalTok{rater[i], }\StringTok{"]"}\NormalTok{)][[}\DecValTok{1}\NormalTok{]]}
    \FunctionTok{time\_to\_temp}\NormalTok{(}
\NormalTok{      pred\_data}\SpecialCharTok{$}\NormalTok{target\_temp[i], }
\NormalTok{      hot, }
\NormalTok{      pred\_data}\SpecialCharTok{$}\NormalTok{Ti[i], }
\NormalTok{      pred\_data}\SpecialCharTok{$}\NormalTok{Tinf[i]}
\NormalTok{    )}
\NormalTok{  \})}
  
  \FunctionTok{data.frame}\NormalTok{(}
    \AttributeTok{rater =}\NormalTok{ pred\_data}\SpecialCharTok{$}\NormalTok{rater[i],}
    \AttributeTok{Tinf =}\NormalTok{ pred\_data}\SpecialCharTok{$}\NormalTok{Tinf[i],}
    \AttributeTok{mean\_time =} \FunctionTok{mean}\NormalTok{(times),}
    \AttributeTok{lower =} \FunctionTok{quantile}\NormalTok{(times, }\FloatTok{0.025}\NormalTok{),}
    \AttributeTok{upper =} \FunctionTok{quantile}\NormalTok{(times, }\FloatTok{0.975}\NormalTok{)}
\NormalTok{  )}
\NormalTok{\})}

\CommentTok{\# Visualize heating time predictions}
\FunctionTok{ggplot}\NormalTok{(time\_preds, }\FunctionTok{aes}\NormalTok{(}\AttributeTok{x =}\NormalTok{ Tinf)) }\SpecialCharTok{+}
  \FunctionTok{geom\_ribbon}\NormalTok{(}\FunctionTok{aes}\NormalTok{(}\AttributeTok{ymin =}\NormalTok{ lower, }\AttributeTok{ymax =}\NormalTok{ upper), }\AttributeTok{alpha =} \FloatTok{0.2}\NormalTok{) }\SpecialCharTok{+}
  \FunctionTok{geom\_line}\NormalTok{(}\FunctionTok{aes}\NormalTok{(}\AttributeTok{y =}\NormalTok{ mean\_time)) }\SpecialCharTok{+}
  \FunctionTok{facet\_wrap}\NormalTok{(}\SpecialCharTok{\textasciitilde{}}\NormalTok{rater, }\AttributeTok{labeller =} \FunctionTok{labeller}\NormalTok{(}\AttributeTok{rater =} \FunctionTok{c}\NormalTok{(}
    \StringTok{"1"} \OtherTok{=} \StringTok{"Rater N"}\NormalTok{, }
    \StringTok{"2"} \OtherTok{=} \StringTok{"Rater TR"}\NormalTok{, }
    \StringTok{"3"} \OtherTok{=} \StringTok{"Rater R"}
\NormalTok{  ))) }\SpecialCharTok{+}
  \FunctionTok{labs}\NormalTok{(}
    \AttributeTok{title =} \StringTok{"Time Required to Reach Pizza{-}Making Temperature"}\NormalTok{,}
    \AttributeTok{subtitle =} \StringTok{"Target temperature: 400°C"}\NormalTok{,}
    \AttributeTok{x =} \StringTok{"Flame Temperature (°C)"}\NormalTok{, }
    \AttributeTok{y =} \StringTok{"Minutes to reach target"}
\NormalTok{  ) }\SpecialCharTok{+}
  \FunctionTok{theme\_bw}\NormalTok{()}
\end{Highlighting}
\end{Shaded}

\begin{center}
\includegraphics[width=0.8\linewidth,height=\textheight,keepaspectratio]{01-PizzaExperiment_files/figure-pdf/ch02_heating_time_preds-1.pdf}
\end{center}

Our analysis reveals several important insights for practical pizza
making. First, the heating time decreases non-linearly with flame
temperature, showing diminishing returns at very high temperatures. We
also observed differences between raters in their measured heating
times. These variations likely stem from differences in measurement
technique and location on the stone, highlighting the importance of
consistent temperature monitoring practices. E.g. N measure the
temperature at the boundary between flame and stone, which matches the
model's assumptions better than R, who measures at the center of the
stone.

For practical application, we can provide specific heating guidelines
based on our model. At a typical flame temperature of 800°C, the model
predicts it will take approximately 20-30 minutes to reach optimal
pizza-making temperature, assuming room temperature start. However, this
time can vary significantly based on:

\begin{itemize}
\tightlist
\item
  Initial stone temperature
\item
  Flame temperature and consistency
\item
  Environmental conditions.
\end{itemize}

However, we should also consider that the model best fits the
measurements from Rater N, who measured at the boundary between flame
and stone. This suggests that our predictions are most applicable to the
border of the stone, and less to the center (measured by R). Yet, it is
the center of the stone that is a better index of the temperature all of
our pizza will receive. Therefore we should add a few minutes to the
model estimates to ensure best real-world outcomes.

Can we really wait that long? Well, there are no shortcuts to the
perfect crust! Just like there is no shortcut to understanding human
cognition, and models need always to be critically contextualized in the
real world.

\section{Conclusion: From Pizza to Cognitive
Principles}\label{conclusion-from-pizza-to-cognitive-principles}

The journey from modeling a heating pizza stone to understanding
cognitive processes might seem unusual, but it illustrates fundamental
principles that will guide us throughout this course.

Through this seemingly simple physics problem, we have encountered core
challenges faced daily by cognitive modelers:

\begin{itemize}
\item
  \emph{Levels of Analysis}: We saw how simplifying complex equations
  (heat diffusion) mirrors the need to choose the right level of
  abstraction when modeling cognition (e.g., neural vs.~behavioral).
\item
  \emph{Prior Knowledge}: We used physics laws to build a generative
  model, much like cognitive theories guide model development.
\item
  \emph{Validation}: We confronted the necessity of comparing model
  predictions to real data, a process central to testing theories.
\item
  \emph{Parameter Estimation}: We used the model to quantify an
  unobservable property (HOT) from noisy data, analogous to estimating
  latent cognitive parameters like learning rates or decision
  thresholds.
\end{itemize}

Crucially, the comparison between purely statistical models
(linear/lognormal) and the physics-based Stan model highlighted the
value of mechanistic understanding. While the statistical models
struggled to capture the saturation curve --- often predicting
impossible temperatures via extrapolation --- the physics model offered
deeper insight and robust predictions. This is a primary goal in
cognitive modeling: we do not simply to describe \emph{which} decisions
were made, but we want to capture the underlying mechanism that can
increase our understanding of \emph{how} decisions were made, and
generalize to novel unseen decisions.

The pizza stone experiment also demonstrated the importance of rigorous
methodology. We saw how multiple measurements from different raters
revealed variability in our data, leading us to consider models
assumptions, measurement error and individual differences---themes that
will become crucial when studying human behavior.

\section{The Course Roadmap}\label{the-course-roadmap}

As we progress through the course, the principles from this chapter will
resurface at every step: a generative model driven by theory (Ch. 3),
parameter estimation via Bayesian inference (Ch. 5), validation of that
inference (Ch. 6), and scaling to many participants (Ch. 7). The pizza
stone is not a detour --- it is the template.

Every chapter is one step in a single pipeline {[}NOW COVERING UP TO
MODEL COMPARISON. MORE TO COME{]}. Use this table to orient yourself
whenever you feel lost.

\begin{longtable}[]{@{}
  >{\raggedright\arraybackslash}p{(\linewidth - 4\tabcolsep) * \real{0.2812}}
  >{\raggedright\arraybackslash}p{(\linewidth - 4\tabcolsep) * \real{0.1875}}
  >{\raggedright\arraybackslash}p{(\linewidth - 4\tabcolsep) * \real{0.5312}}@{}}
\toprule\noalign{}
\begin{minipage}[b]{\linewidth}\raggedright
Chapter
\end{minipage} & \begin{minipage}[b]{\linewidth}\raggedright
Step
\end{minipage} & \begin{minipage}[b]{\linewidth}\raggedright
Central question
\end{minipage} \\
\midrule\noalign{}
\endhead
\bottomrule\noalign{}
\endlastfoot
\textbf{1 --- Pizza Experiment} & Why model? & What is a generative
model and what does it mean to fit one to data? \\
\textbf{2 --- Building Models} & Observe \& theorize & What behaviors do
we see, and what verbal strategies could explain them? \\
\textbf{3 --- Formal Models} & Formalize \& simulate & Can we turn a
verbal strategy into code that \emph{generates} realistic data? \\
\textbf{4 --- Inferring Rates} & Fit a single dataset & Given observed
data, can we infer the parameters that produced it? \\
\textbf{5 --- Model Quality} & Validate the engine & Are our parameter
estimates actually trustworthy? (Prior checks, recovery, SBC, posterior
checks) \\
\textbf{6 --- Multilevel Models} & Scale to populations & How do
individuals differ, and how do we pool information across them? \\
\textbf{7 --- Model Comparison} & Choose between theories & Of several
plausible models, which one best predicts unseen data? \\
\end{longtable}

\textbf{The key insight:} steps 1--3 build models; steps 4--7 test them.
A model only earns the right to touch real human data after passing the
validation battery in Chapter 6.

\textbf{Coming up in Chapter 2:} We shift focus from stones to minds.
You have already played the Matching Pennies game in class. Chapter 2
helps you translate what you observed --- your own strategies, your
opponent's patterns --- into candidate verbal models that Chapter 3 will
then formalize in code.

I hope this has whetted your appetite -- both for pizza and for
computational modeling!

\bookmarksetup{startatroot}

\chapter{Building Models of Strategic
Decision-Making}\label{building-models-of-strategic-decision-making}

\begin{quote}
\textbf{📍 Where we are in the Bayesian modeling workflow:} Ch. 1
established \emph{what} modeling means. Now we need something to model.
This chapter gives us our cognitive task (Matching Pennies) and has us
build theories \emph{in words} --- no code yet. These verbal models are
the raw material that Ch. 3 translates into simulatable algorithms and
Ch. 5 fits to real data.
\end{quote}

You've just played Matching Pennies and discussed strategies with your
classmates. You probably noticed patterns in your opponent's play, tried
to be unpredictable, and maybe even changed your strategy mid-game. This
chapter helps you translate those observations and intuitions into the
language of cognitive modeling---preparing you to implement formal
models in Chapter 3.

\section{Learning Objectives}\label{learning-objectives-1}

This chapter bridges the gap between observing behavior and developing
testable theories. By the end of this chapter, using the Matching
Pennies game as a case study, you will be able to:

\begin{itemize}
\tightlist
\item
  \textbf{Identify Key Modeling Steps:} Understand the process of moving
  from behavioral observations and participant reflections to
  formulating initial verbal theories of underlying cognitive
  strategies.
\item
  \textbf{Appreciate Theory Building Challenges:} Recognize common
  issues in theory development, such as the participant vs.~researcher
  perspective, the need for simplification, and incorporating known
  cognitive constraints.
\item
  \textbf{Conceptualize Learning Mechanisms:} Propose distinct models
  (e.g., random choice, heuristics, RL) and organize them into a unified
  framework based on how they process prediction errors.
\item
  \textbf{Connect to Formalization:} Understand the ``Update Rule''
  (\(\text{New} = \text{Old} + \text{Learning Rate} \times \text{Error}\))
  as a universal grammar of learning that connects simple heuristics to
  complex cognitive strategies.
\end{itemize}

\section{Introduction: Observing Behavior to Theorize
Mechanisms}\label{introduction-observing-behavior-to-theorize-mechanisms}

Chapter 1 emphasized the importance of modeling underlying generative
mechanisms. To do this for cognition, we first need a behavior to
explain. This chapter uses the \textbf{Matching Pennies game} as our
initial cognitive phenomenon. It's a simple strategic interaction, yet
rich enough to illustrate the process of developing and refining
cognitive models.

Our goal here is not yet to build the final computational models, but to
practice the crucial preceding steps:

\begin{enumerate}
\def\labelenumi{\arabic{enumi}.}
\item
  Observing behavior in a specific task (through experiments and data
  exploration).
\item
  Reflecting on potential cognitive strategies and constraints (drawing
  on observations, participant reports, and cognitive science
  principles).
\item
  Formulating initial \emph{verbal} theories or candidate models that
  describe the potential underlying mechanisms.
\end{enumerate}

This process lays the groundwork for Chapter 3, where we will translate
these verbal ideas into precise, formal models ready for simulation and
testing.

\begin{quote}
\emph{Why build verbal models before writing any code?} Jumping straight
to code forces you to make dozens of implicit decisions (What counts as
a win? Does memory last 1 trial or 10?) without realizing you're making
them. Working through the verbal model first surfaces those hidden
assumptions so you can make them deliberately. Every ambiguity you
resolve here saves you a debugging session in Ch. 3.
\end{quote}

\section{The Matching Pennies Game}\label{the-matching-pennies-game}

In the matching pennies game, two players engage in a series of choices.
One player attempts to match the other's choice, while the other player
aims to achieve a mismatch, and they repeatedly play with each other.
This is a prototypical example of interacting behaviors that are usually
tackled by game theory, and bring up issues of theory of mind and
recursivity.

For an introduction see the paper: Waade, Peter T., et al.~``Introducing
tomsup: Theory of mind simulations using Python.'' Behavior Research
Methods 55.5 (2023): 2197-2231.

For fun data involving different kinds of primates playing the game,
see: Devaine, M., San-Galli, A., Trapanese, C., Bardino, G., Hano, C.,
Saint Jalme, M., \ldots{} \& Daunizeau, J. (2017). Reading wild minds: a
computational assay of theory of mind sophistication across seven
primate species. PLoS computational biology, 13(11), e1005833. The data
is available in Assignment 2 (for the students at AU/CogSci).

\section{Game Structure}\label{game-structure}

The game proceeds as follows:

\begin{enumerate}
\def\labelenumi{\arabic{enumi}.}
\tightlist
\item
  Two players sit facing each other
\item
  Each round, both players choose either ``left'' or ``right'' to
  indicate where they believe a penny is hidden
\item
  The matcher wins by choosing the same hand as their opponent
\item
  The hider wins by choosing the opposite hand
\item
  Points are awarded: +1 for winning, -1 (or 0, depending on the
  version) for losing
\item
  Repeat
\end{enumerate}

This simple structure creates a rich environment for studying
decision-making strategies, learning, and adaptation.

\section{Empirical Investigation}\label{empirical-investigation}

\subsection{Data Collection Protocol}\label{data-collection-protocol}

{[}NOTES FOR FUTURE YEARS: MAYBE REPLACE WITH THE PRIMATE DATA?{]}

If you are attending my class you have been (or will be) asked to
participate in a matching pennies game. This game provides the
foundation for our modeling efforts. By observing gameplay and
collecting data, we can develop models that capture the cognitive
processes underlying decision-making in strategic situations.

Participants play 30 rounds as the matcher and 30 rounds as the hider,
allowing us to observe behavior in both roles. While playing,
participants track their scores, which can provide quantitative data for
later analysis. Participants are also asked to reflect on their
strategies and the strategies they believe their opponents are using, as
that provides valuable materials to build models on.

\subsection{Initial Observations}\label{initial-observations}

Through the careful observation and discussion of gameplay we do in
class, several patterns typically emerge. For instance, players often
demonstrate strategic adaptation, adjusting their choices based on their
opponent's previous moves. They may attempt to identify patterns in
their opponent's behavior while trying to make their own choices less
predictable. The tension between exploitation of perceived patterns and
maintenance of unpredictability creates fascinating dynamics for
modeling.

\section{Empirical explorations}\label{empirical-explorations}

Below you can observe how a previous year of CogSci did against bots
(computational agents) playing according to different strategies. Look
at the plots below, where the x axes indicate trial, the y axes how many
points the CogSci'ers scored (0 being chance, negative means being
completely owned by the bots, positive owning the bot) and the different
colors indicate different strategies employed by the bots. Strategy
``-2'' was a Win-Stay-Lose-Shift bot: when it got a +1, it repeated its
previous move (e.g.~right if it had just played right), otherwise it
would perform the opposite move (e.g.~left if it had just played right).
Strategy ``-1'' was a biased Nash both, playing ``right'' 80\% of the
time. Strategy ``0'' indicates a reinforcement learning bot; ``1'' a bot
assuming you were playing according to a reinforcement learning strategy
and trying to infer your learning and temperature parameters; ``2'' a
bot assuming you were following strategy ``1'' and trying to accordingly
infer your parameters.

\begin{Shaded}
\begin{Highlighting}[]
\FunctionTok{library}\NormalTok{(tidyverse)}

\CommentTok{\# {-}{-}{-} 1. Data Loading / Generation {-}{-}{-}}
\NormalTok{data\_path }\OtherTok{\textless{}{-}} \FunctionTok{file.path}\NormalTok{(}\StringTok{"data"}\NormalTok{, }\StringTok{"MP\_MSc\_CogSci22.csv"}\NormalTok{)}

\ControlFlowTok{if}\NormalTok{ (}\FunctionTok{file.exists}\NormalTok{(data\_path)) \{}
\NormalTok{  d }\OtherTok{\textless{}{-}} \FunctionTok{read\_csv}\NormalTok{(data\_path)}
\NormalTok{\} }\ControlFlowTok{else}\NormalTok{ \{}
  \CommentTok{\# Generate synthetic data if file is missing (Reproducibility check)}
  \FunctionTok{set.seed}\NormalTok{(}\DecValTok{42}\NormalTok{)}
\NormalTok{  n\_students }\OtherTok{\textless{}{-}} \DecValTok{20}
\NormalTok{  n\_trials }\OtherTok{\textless{}{-}} \DecValTok{30}
\NormalTok{  strategies }\OtherTok{\textless{}{-}} \FunctionTok{c}\NormalTok{(}\SpecialCharTok{{-}}\DecValTok{2}\NormalTok{, }\SpecialCharTok{{-}}\DecValTok{1}\NormalTok{, }\DecValTok{0}\NormalTok{, }\DecValTok{1}\NormalTok{, }\DecValTok{2}\NormalTok{)}
\NormalTok{  d }\OtherTok{\textless{}{-}} \FunctionTok{expand\_grid}\NormalTok{(}
    \AttributeTok{ID =} \FunctionTok{factor}\NormalTok{(}\DecValTok{1}\SpecialCharTok{:}\NormalTok{n\_students),}
    \AttributeTok{BotStrategy =}\NormalTok{ strategies,}
    \AttributeTok{Role =} \FunctionTok{c}\NormalTok{(}\DecValTok{0}\NormalTok{, }\DecValTok{1}\NormalTok{), }\CommentTok{\# 0=Matcher, 1=Hider}
    \AttributeTok{Trial =} \DecValTok{1}\SpecialCharTok{:}\NormalTok{n\_trials}
\NormalTok{  ) }\SpecialCharTok{\%\textgreater{}\%}
    \FunctionTok{mutate}\NormalTok{(}
      \CommentTok{\# Random payoffs for demonstration}
      \AttributeTok{Payoff =} \FunctionTok{sample}\NormalTok{(}\FunctionTok{c}\NormalTok{(}\SpecialCharTok{{-}}\DecValTok{1}\NormalTok{, }\DecValTok{1}\NormalTok{), }\FunctionTok{n}\NormalTok{(), }\AttributeTok{replace =} \ConstantTok{TRUE}\NormalTok{, }\AttributeTok{prob =} \FunctionTok{c}\NormalTok{(}\FloatTok{0.45}\NormalTok{, }\FloatTok{0.55}\NormalTok{))}
\NormalTok{    )}
  \FunctionTok{warning}\NormalTok{(}\StringTok{"Using synthetic data for demonstration."}\NormalTok{)}
\NormalTok{\}}

\CommentTok{\# {-}{-}{-} 2. Data Cleaning (Crucial Step!) {-}{-}{-}}
\CommentTok{\# Map cryptic codes to human{-}readable labels}
\NormalTok{bot\_labels }\OtherTok{\textless{}{-}} \FunctionTok{c}\NormalTok{(}
  \StringTok{"{-}2"} \OtherTok{=} \StringTok{"WSLS Bot"}\NormalTok{,}
  \StringTok{"{-}1"} \OtherTok{=} \StringTok{"Bias Bot (80\%)"}\NormalTok{,}
  \StringTok{"0"}  \OtherTok{=} \StringTok{"RL Bot"}\NormalTok{,}
  \StringTok{"1"}  \OtherTok{=} \StringTok{"ToM{-}1 Bot"}\NormalTok{,}
  \StringTok{"2"}  \OtherTok{=} \StringTok{"ToM{-}2 Bot"}
\NormalTok{)}

\NormalTok{d\_clean }\OtherTok{\textless{}{-}}\NormalTok{ d }\SpecialCharTok{\%\textgreater{}\%}
  \FunctionTok{mutate}\NormalTok{(}
    \CommentTok{\# Make BotStrategy a factor with meaningful names}
    \AttributeTok{BotStrategy =} \FunctionTok{factor}\NormalTok{(BotStrategy, }
                         \AttributeTok{levels =} \FunctionTok{names}\NormalTok{(bot\_labels), }
                         \AttributeTok{labels =}\NormalTok{ bot\_labels),}
    \CommentTok{\# Make Role a factor}
    \AttributeTok{Role =} \FunctionTok{factor}\NormalTok{(Role, }\AttributeTok{levels =} \FunctionTok{c}\NormalTok{(}\DecValTok{0}\NormalTok{, }\DecValTok{1}\NormalTok{), }\AttributeTok{labels =} \FunctionTok{c}\NormalTok{(}\StringTok{"Matcher"}\NormalTok{, }\StringTok{"Hider"}\NormalTok{))}
\NormalTok{  )}

\CommentTok{\# {-}{-}{-} 3. Plot Collective Performance {-}{-}{-}}
\FunctionTok{ggplot}\NormalTok{(d\_clean, }\FunctionTok{aes}\NormalTok{(}\AttributeTok{x =}\NormalTok{ Trial, }\AttributeTok{y =}\NormalTok{ Payoff, }\AttributeTok{color =}\NormalTok{ BotStrategy)) }\SpecialCharTok{+}
  \FunctionTok{geom\_smooth}\NormalTok{(}\AttributeTok{se =} \ConstantTok{FALSE}\NormalTok{, }\AttributeTok{method =} \StringTok{"loess"}\NormalTok{, }\AttributeTok{span =} \FloatTok{0.5}\NormalTok{) }\SpecialCharTok{+}
  \FunctionTok{geom\_hline}\NormalTok{(}\AttributeTok{yintercept =} \DecValTok{0}\NormalTok{, }\AttributeTok{linetype =} \StringTok{"dashed"}\NormalTok{, }\AttributeTok{alpha =} \FloatTok{0.5}\NormalTok{) }\SpecialCharTok{+}
  \FunctionTok{facet\_wrap}\NormalTok{(}\SpecialCharTok{\textasciitilde{}}\NormalTok{Role) }\SpecialCharTok{+}
  \FunctionTok{labs}\NormalTok{(}
    \AttributeTok{title =} \StringTok{"Human vs. Machine: Average Performance"}\NormalTok{,}
    \AttributeTok{subtitle =} \StringTok{"Positive values indicate humans winning against bots"}\NormalTok{,}
    \AttributeTok{y =} \StringTok{"Average Payoff"}\NormalTok{,}
    \AttributeTok{color =} \StringTok{"Opponent Strategy"}
\NormalTok{  ) }\SpecialCharTok{+}
  \FunctionTok{theme\_minimal}\NormalTok{() }\SpecialCharTok{+}
  \FunctionTok{theme}\NormalTok{(}\AttributeTok{legend.position =} \StringTok{"bottom"}\NormalTok{)}
\end{Highlighting}
\end{Shaded}

\begin{center}
\includegraphics[width=0.8\linewidth,height=\textheight,keepaspectratio]{02-BuildingModels_files/figure-pdf/ch03_load_and_clean_data-1.pdf}
\end{center}

That doesn't look too good, ah? What about individual variability? In
the plot below we indicate the score of each of the former students,
against the different bots.

\begin{Shaded}
\begin{Highlighting}[]
\CommentTok{\# {-}{-}{-} Plot 2: Individual Variability in Scores {-}{-}{-}}
\CommentTok{\# Calculate the total score for each student (ID) against each bot strategy.}
\NormalTok{d\_summary }\OtherTok{\textless{}{-}}\NormalTok{ d }\SpecialCharTok{\%\textgreater{}\%}
  \FunctionTok{group\_by}\NormalTok{(ID, BotStrategy, Role) }\SpecialCharTok{\%\textgreater{}\%} \CommentTok{\# Group by student, bot, and role}
  \FunctionTok{summarize}\NormalTok{(}\AttributeTok{TotalScore =} \FunctionTok{sum}\NormalTok{(Payoff), }\AttributeTok{.groups =} \StringTok{"drop"}\NormalTok{) }\CommentTok{\# Calculate total score}

\CommentTok{\# Visualize the distribution of total scores for each bot strategy.}
\CommentTok{\# geom\_boxplot shows the distribution, geom\_point shows individual student scores.}
\FunctionTok{print}\NormalTok{(}
  \FunctionTok{ggplot}\NormalTok{(d\_summary, }\FunctionTok{aes}\NormalTok{(}\AttributeTok{x =}\NormalTok{ BotStrategy, }\AttributeTok{y =}\NormalTok{ TotalScore)) }\SpecialCharTok{+}
    \FunctionTok{geom\_boxplot}\NormalTok{(}\FunctionTok{aes}\NormalTok{(}\AttributeTok{fill =}\NormalTok{ Role), }\AttributeTok{alpha =} \FloatTok{0.3}\NormalTok{, }\AttributeTok{outlier.shape =} \ConstantTok{NA}\NormalTok{) }\SpecialCharTok{+} \CommentTok{\# Boxplot showing distribution}
    \FunctionTok{geom\_jitter}\NormalTok{(}\FunctionTok{aes}\NormalTok{(}\AttributeTok{color =}\NormalTok{ ID), }\AttributeTok{width =} \FloatTok{0.2}\NormalTok{, }\AttributeTok{alpha =} \FloatTok{0.7}\NormalTok{) }\SpecialCharTok{+} \CommentTok{\# Individual student points}
    \FunctionTok{facet\_wrap}\NormalTok{(}\SpecialCharTok{\textasciitilde{}}\NormalTok{Role) }\SpecialCharTok{+} \CommentTok{\# Separate plots for Matcher and Hider}
    \FunctionTok{labs}\NormalTok{(}
      \AttributeTok{title =} \StringTok{"Distribution of Total Scores Against Different Bots"}\NormalTok{,}
      \AttributeTok{subtitle =} \StringTok{"Shows individual student variability"}\NormalTok{,}
      \AttributeTok{x =} \StringTok{"Bot Strategy"}\NormalTok{,}
      \AttributeTok{y =} \StringTok{"Total Score (Sum of Payoffs)"}\NormalTok{,}
      \AttributeTok{fill =} \StringTok{"Player Role"}
\NormalTok{    ) }\SpecialCharTok{+}
    \FunctionTok{theme\_classic}\NormalTok{() }\SpecialCharTok{+}
    \FunctionTok{theme}\NormalTok{(}\AttributeTok{legend.position =} \StringTok{"none"}\NormalTok{) }\CommentTok{\# Hide legend for individual IDs}
\NormalTok{)}
\end{Highlighting}
\end{Shaded}

\begin{center}
\includegraphics[width=0.8\linewidth,height=\textheight,keepaspectratio]{02-BuildingModels_files/figure-pdf/ch03_plot_individual_performance-1.pdf}
\end{center}

\subsection{From Observation to Theory: Identifying Potential
Mechanisms}\label{from-observation-to-theory-identifying-potential-mechanisms}

The plots above reveal patterns: average performance changes over time,
varies by opponent, and differs across individuals. Gameplay
observations and participant reflections (from class discussion or
collected data) add qualitative insights -- perhaps players mention
trying to be unpredictable, guessing opponent biases, or repeating
winning moves.

The crucial next step is to distill these rich, complex observations
into simplified, plausible \textbf{mechanisms} or \textbf{strategies}.
This involves abstraction:

\begin{itemize}
\item
  \textbf{Identifying Core Patterns:} What recurring behaviors seem most
  important? (e.g., reacting to wins/losses, tracking opponent
  frequencies).
\item
  \textbf{Simplifying:} Can we capture the essence of a strategy without
  modeling every detail of a player's thought process or interaction?
  (e.g., modeling just previous behaviors as possible inputs, instead of
  complex patterns in face expressions or movements that real humans
  likely use when trying to guess which hand has the penny).
\item
  \textbf{Drawing on Cognitive Principles:} How do known cognitive
  constraints (like limited memory or processing errors) shape plausible
  strategies?
\end{itemize}

For instance, observing that players often change their choice after a
loss might lead us to propose a ``If Lose Then Shift'' component as part
of a candidate model. Observing that performance differs against biased
vs.~adaptive bots suggests players might be trying to learn or adapt,
leading to memory-based or learning models.

This process generates \emph{verbal models} -- initial hypotheses about
the strategies at play. Key modeling considerations guide this
translation:

\begin{itemize}
\item
  What information do players likely use? (Own past choices? Opponent's
  choices? Payoffs?)
\item
  How far back does memory plausibly extend? (Last trial? Last 5 trials?
  Exponential decay?)
\item
  What is the role of randomness? (True randomness? Exploration?
  Implementation errors?)
\item
  How might strategies adapt over time or differ between individuals?
\end{itemize}

Answering these helps refine our verbal models, paving the way for
formalization. The goal isn't to capture everything, but to propose
distinct, testable mechanisms.

\subsection{The distinction between participant and researcher
perspectives}\label{the-distinction-between-participant-and-researcher-perspectives}

As participants we might not be aware of the strategy we use, or we
might believe something erroneous. The exercise here is to act as
researchers: what are the principles underlying the participants'
behaviors, no matter what the participants know or believe? Note that
talking to participants and being participants helps developing ideas,
but it's not the end point of the process. Also note that as cognitive
scientists we can rely on what we have learned about cognitive processes
(e.g.~memory).

Another important component of the distinction is that participants
leave in a rich world: they rely on facial expressions and bodily
posture, the switch strategies, etc. On the other hand, the researcher
is trying to identify one or few at most ``simple'' strategies. Rich
bodily interactions and mixtures or sequences of multiple strategies are
not a good place to start modeling. These aspects are a poor starting
point for building your first model, and are often pretty difficult to
fit to empirical data. Nevertheless, they are important intuitions that
the researcher should (eventually?) accommodate.

\section{Candidate Models: A Hierarchy of Cognitive
Complexity}\label{candidate-models-a-hierarchy-of-cognitive-complexity}

Based on behavioral observations and cognitive principles, the many
discussions in class during the years have advanced a variety of models,
which I'm here organizing into a hierarchy based on their conception of
learning. We start with simple agents that ignore history, move to
learning agents that track first average outcomes, then also uncertainty
and volatility, and conclude with strategic agents that model the
opponent's mind. To make their conceptual continuity more explicit, I
couldn't resist adding some discussion of the math behind the verbal
intuitions and pointing to how each model can be recast in terms of an
update rule:
\[\text{New Belief} = \text{Old Belief} + \text{Learning Rate} \cdot \text{Predictive Error}\].
This format unifies the different models under a common mathematical
framework, making it easier to understand their relationships and
differences.

Also note that in the formula below, `Belief' is a placeholder. If we
are modeling a coin flipper, the belief is a probability (\(\theta\)).
If we are modeling betting, it might be an expected reward value
(\(V\)). If we are modeling playing strategies, it might be a prediction
of the opponent's move. As it should be obvious, those are partially
overlapping conceptualizations.

\subsection{Level 0: The Static Agent (No
Learning)}\label{level-0-the-static-agent-no-learning}

The simplest assumption is that the agent does not learn or adapt to the
opponent at all. They simply act according to a fixed internal
preference.

\subsection{The Biased Agent (Random
Choice):}\label{the-biased-agent-random-choice}

\begin{itemize}
\item
  \emph{Concept}: The agent is a biased coin flipper. They have a static
  preference (e.g., ``I generally prefer Right'') that does not change
  over time.
\item
  \emph{Traditional Formulation (Bernoulli Process)}: We model the
  choice \(y_t\) at trial \(t\) as a draw from a Bernoulli distribution
  with a fixed rate parameter \(\theta\):
  \[y_t \sim \text{Bernoulli}(\theta)\]
\item
  \emph{The ``Update'' Rule}: Because this agent is static, there is no
  update rule, in other words, a learning of 0. The parameter \(\theta\)
  remains constant regardless of wins or losses.
  \[\theta_{t+1} = \theta_{t} + \text{Learning Rate} \cdot \text{Predictive Error}\]
  where \(\text{Learning Rate} = 0\) and therefore the formula reduces
  to \[\theta_{t+1} = \theta_{t}\].
\item
  \emph{Variable Definitions}:

  \begin{itemize}
  \tightlist
  \item
    \(y_t \in \{0, 1\}\): The choice (e.g., Left/Right).
  \item
    \(\theta \in [0, 1]\): The fixed probability of choosing 1 (bias).
  \end{itemize}
\item
  \emph{Looking Ahead}: This model might seem too simple to be useful,
  but never underestimate the power of null models! In a study using
  k-ToM models (see below), we realized that for k = 6 or above, the
  model behaved in a way that was practically indistinguishabe from a
  biased agent. Also, a simple model allows us to build skills in
  reasonable steps. In the next chapter, we will write code to simulate
  data according to these models and in the following chapter, we will
  then learn how to take a sequence of choices and mathematically work
  backwards to find the most likely value of the model parameters. In
  all these cases, being able to start with a simple model is crucial to
  build intuition and skills before moving to more complex models.
\end{itemize}

{[}FIND REF OF PAPER USING THIS{]}

\subsection{Level 1: Heuristics \& The Immediate
Past}\label{level-1-heuristics-the-immediate-past}

The next step up is an agent that learns minimally: it reacts to
feedback but has a ``memory'' of only one trial.

\subsubsection{Deterministic Win-Stay-Lose-Shift
(WSLS):}\label{deterministic-win-stay-lose-shift-wsls}

\begin{itemize}
\item
  \emph{Concept}: A rigid heuristic: ``If I won, I do the same thing. If
  I lost, I switch.'' This can be viewed as an agent with zero memory
  who effectively ``re-learns'' the world from scratch after every
  single feedback.
\item
  \emph{Traditional Formulation (Conditional Probability)}:
  \[P(y_t = y_{t-1} | outcome_{t-1}) = \begin{cases} 1 & \text{if } outcome_{t-1} = \text{win} \\ 0 & \text{if } outcome_{t-1} = \text{loss} \end{cases}\]
\item
  \emph{The ``Update'' Rule}: We can treat the ``Probability of
  Staying'' (\(P_{stay}\)) as a value that is updated
  trial-by-trial.\[P_{stay, t+1} = P_{stay, t} + \text{Learning Rate} \cdot (Outcome_t - P_{stay, t})\]
  Where \(\text{Learning Rate} = 1\). Because the agent learns instantly
  from the immediate past, the formula reduces to:
  \[P_{stay, t+1} = P_{stay, t} + 1 \cdot (Outcome_t - P_{stay, t})\]\[P_{stay, t+1} = Outcome_t\]

  \begin{itemize}
  \tightlist
  \item
    Implication: If the outcome was a Win (\(1\)), the probability of
    staying becomes \(1\). If the outcome was a Loss (\(0\)), the
    probability of staying becomes \(0\). This mathematically formalizes
    the idea of ``instant forgetting.''
  \end{itemize}
\item
  \emph{Variable Definitions}:\(outcome_{t-1}\):
  \(outcome_{t-1} \in \{0, 1\}\): The feedback from the previous trial,
  where 1 = Win and 0 = Loss
\item
  \emph{Note: The Physics View (Ising Formulation)} {[}only for
  nerds{]}: While we typically code choices as 0/1 for logistic
  regression, physicists and network theorists often code binary states
  as ``spins'': -1 (Left) and +1 (Right). If we also code the outcome as
  -1 (Loss) and +1 (Win), the WSLS update rule becomes mathematically
  elegant: \[y_t = outcome_{t-1} \cdot y_{t-1}\] Why does this work? If
  Win (+1): The equation becomes \(y_t = 1 \cdot y_{t-1}\). The sign
  stays the same (Stay). If Loss (-1): The equation becomes
  \(y_t = -1 \cdot y_{t-1}\). The sign flips (Shift). This formulation
  highlights that ``shifting'' is mathematically equivalent to sign
  inversion. This connects cognitive shifting to ``spin glass'' or ising
  models in physics (e.g., Stephens \& Bialek, 2010), and is useful to
  know of it as it sometimes simplifies the mathematical equations of
  your models.
\end{itemize}

{[}FIND REF OF PAPERS USING THIS MODES{]}

\subsubsection{Probabilistic Win-Stay-Lose-Shift (Logistic
Regression):}\label{probabilistic-win-stay-lose-shift-logistic-regression}

\begin{itemize}
\item
  \emph{Concept}: We relax the strict rule. The agent is more likely to
  stay after a win and shift after a loss, but real behavior is noisy
  (stochastic, probabilistic).
\item
  \emph{Traditional Formulation (GLM)}: We define the probability of
  ``staying'' (\(P_{stay}\)) as a function of the previous outcome using
  a logistic curve.
  \[P(stay_t) = \text{logit}^{-1}(\beta_0 + \beta_1 \cdot outcome_{t-1})\].
  Note that we can think of the reactions to feedback as symmetric
  (probability of staying when winning equal to the possibility of
  shifting when losing), or asymmetric (e.g.~the agent might be more
  sensitive to losses than wins, or vice versa). In the latter case,
  we'd have an equation like:
  \[P(stay_t) = \text{logit}^{-1}(\beta_0 + \beta_1 \cdot win_{t-1} + \beta_2 \cdot loss_{t-1})\]
\item
  \emph{The ``Update'' Rule}: To fit our hierarchy, we treat the agent's
  internal Decision Value (\(Q\)) as the state being
  updated.\[Q_{t+1} = Q_{t} + \text{Learning Rate} \cdot (\text{Target}_t - Q_{t})\]Where
  \(\text{Learning Rate} = 1\) (Instant Update) and the Target is the
  weighted outcome
  (\(\beta_0 + \beta_1 \cdot outcome_t\)).Reduction:\[Q_{t+1} = \beta_0 + \beta_1 \cdot outcome_t\]Implication:
  Just like the Deterministic agent, this agent ``forgets'' the distant
  past instantly. However, a ``Win'' does not guarantee a Stay
  (\(P=1\)); it simply pushes the decision value to a high---but not
  infinite---number.
\item
  \emph{Variable Definitions}:

  \begin{itemize}
  \item
    \(Q\): The internal log-odds of staying.
  \item
    \(\text{logit}^{-1}(x) = \frac{1}{1+e^{-x}}\): The sigmoid function
    mapping value to probability.
  \item
    \(\beta_1\): Sensitivity to feedback. High \(\beta_1\) approximates
    the deterministic heuristic; \(\beta_1 \approx 0\) implies the agent
    ignores feedback.
  \end{itemize}
\end{itemize}

{[}FIND REF OF PAPER USING THIS e.g., Zhang \& Lee, 2010; Worthy et al.,
2013{]}

\subsection{Level 2: Integrating History (The Moving
Average)}\label{level-2-integrating-history-the-moving-average}

Real agents usually care about more than just the last trial. They
integrate history over time to form an expectation.

\subsubsection{The Moving Average (Windowed
Memory):}\label{the-moving-average-windowed-memory}

\begin{itemize}
\item
  \emph{Concept}: ``I think the future will look like the average of the
  recent past.''
\item
  \emph{Traditional Formulation (Batch Average)}:
  \[V_{t} = \frac{1}{N} \sum_{i=0}^{N-1} Outcome_{t-i}\]
\item
  \emph{Transformation to Update Rule}: To avoid recalculating the sum
  from scratch, we can update the average by adding the new outcome and
  removing the oldest
  outcome.\[V_{t} = V_{t-1} + \frac{1}{N} \cdot (Outcome_{t} - Outcome_{t-N})\]
\item
  \emph{Critique (The Memory Problem)}: This formulation reveals a
  massive cognitive cost. To perform this update, the agent cannot just
  store the current average (\(V_{t-1}\)); they must remember exactly
  what happened \(N\) trials ago (\(Outcome_{t-N}\)) to ``remove'' it
  from the sum. Further, a memory from 10 trials ago is recalled
  perfectly, but a memory from 11 trials ago vanishes instantly. This is
  biologically implausible.

  \begin{itemize}
  \tightlist
  \item
    The Mismatch: This does not fit the standard ``Prediction Error''
    update rule. It compares the new reality (\(Outcome_t\)) to the
    oldest reality (\(Outcome_{t-N}\)), rather than comparing reality to
    expectation.
  \end{itemize}
\item
  \emph{Variable Definitions}:

  \begin{itemize}
  \tightlist
  \item
    \(V_t\): The estimated value (probability of winning/Right) at trial
    \(t\).
  \item
    \(N\): The window size (e.g., 10 trials).
  \item
    \(Outcome_t \in \{0, 1\}\): The outcome at trial \(t\).
  \end{itemize}
\end{itemize}

\subsection{Level 3: A more elegant memory decay
mechanism}\label{level-3-a-more-elegant-memory-decay-mechanism}

Reinforcement Learning (RL) is mathematically the elegant solution to
the Moving Average's memory problem. It replaces the ``perfect memory of
\(N\) trials'' with a ``fading memory of everything''.

\subsubsection{From Moving Average to Reinforcement Learning
(Rescorla-Wagner):}\label{from-moving-average-to-reinforcement-learning-rescorla-wagner}

How do we fix the memory problem in Level 2? The agent cannot remember
the specific outcome \(Outcome_{t-N}\). So, they make a guess: ``The
outcome \(N\) trials ago was probably similar to my current average
value.'' Mathematically, we substitute
\(Outcome_{t-N} \approx V_{t-1}\).
\[V_{t} = V_{t-1} + \frac{1}{N} (Outcome_{t} - \mathbf{V_{t-1}})\] If we
rename \(\frac{1}{N}\) to \(\alpha\) (alpha), we get the standard RL
rule.

Note that the update rule only applies to the value of the action that
was actually taken. For instance, if we are modeling the value of
picking ``Right'', the outcome should be 1 if we picked ``Right'' and
won, 0 if we picked ``Right'' and lost, and it should not be updated at
all if we picked ``Left''. In the context of the matching pennies we can
assume that increased expected value for left will go hand in hand with
decreased expected value for right (and so when the participant chooses
``Left'', use \(1-Outcome\) to update the value of ``Right'').

\begin{itemize}
\item
  \emph{``Traditional Formulation (Exponential Weighted Average)}:
  Instead of a sum over \(N\), value is a weighted sum of all past
  history, where weights decay geometrically.
  \[V_{t} = (1-\alpha) \cdot V_{t-1} + \alpha \cdot Reward_{t}\]
\item
  \emph{Transformation to Update Rule (The Delta Rule)}: By rearranging
  terms, we get the update format I used as template for all models:
  \[V_{t} = V_{t-1} + \alpha \cdot (Reward_{t} - V_{t-1})\]
\item
  \emph{Why is this ``Learning''?} The term \((Reward_{t} - V_{t-1})\)
  is the Prediction Error (PE). The agent compares reality (\(Reward\))
  to their expectation (\(V\)) and nudges their belief by a step size
  \(\alpha\).
\item
  \emph{Variable Definitions}:

  \begin{itemize}
  \tightlist
  \item
    \(V_t\): The estimated value (e.g., probability of win) at trial
    \(t\).
  \item
    \(\alpha \in [0, 1]\): The Learning Rate. High \(\alpha\) = fast
    forgetting; Low \(\alpha\) = long memory.
  \item
    \(PE\): The difference between what happened and what was expected.
  \end{itemize}
\end{itemize}

\begin{tcolorbox}[enhanced jigsaw, toprule=.15mm, toptitle=1mm, bottomtitle=1mm, leftrule=.75mm, rightrule=.15mm, colframe=quarto-callout-note-color-frame, left=2mm, colback=white, titlerule=0mm, coltitle=black, colbacktitle=quarto-callout-note-color!10!white, opacityback=0, title=\textcolor{quarto-callout-note-color}{\faInfo}\hspace{0.5em}{Historical Context: From Behaviorism to Bayesian Belief --- The
Prediction-Error Lineage}, arc=.35mm, opacitybacktitle=0.6, breakable, bottomrule=.15mm]

The simple update rule
\(V_t = V_{t-1} + \alpha \cdot (Outcome_t - V_{t-1})\) looks almost
trivially obvious once written down. Getting there took cognitive
science roughly four decades and required overturning an entire
paradigm.

\subsubsection{The Behaviorist
Prohibition}\label{the-behaviorist-prohibition}

For the first half of the twentieth century, scientific psychology
largely banned the internal state variable \(V_t\) from its vocabulary.
The dominant framework --- Hullian behaviorism (Hull 1943) and its
variants --- insisted that a rigorous account of learning had to be
stated purely in terms of observable stimuli and responses, with no
reference to unobservable mental states. Tolman's (Tolman 1948) proposal
that rats form ``cognitive maps'' --- internal representations of their
environment --- was widely regarded as an embarrassing lapse into
mentalism. The prediction error \((Outcome - V)\) cannot exist without
\(V\), and \(V\) is a latent belief. The entire framework of the models
in this chapter was conceptually prohibited.

\subsubsection{The Rescorla-Wagner
Pivot}\label{the-rescorla-wagner-pivot}

The rehabilitation came from classical conditioning research, and
specifically from the realisation that conditioning depends not on the
raw presence of the unconditioned stimulus but on \emph{how surprising}
it is. Rescorla and Wagner (1972) introduced the prediction-error rule
into learning theory to explain blocking (Kamin 1969): if one cue
already perfectly predicts an outcome, adding a second cue produces no
conditioning to it, because the outcome is not surprising. This
``blocking effect'' is exactly what the prediction-error term captures
--- when \(V_{t-1}\) already equals \(Outcome_t\), the error is zero and
no updating occurs. Rescorla and Wagner's model was the first to make
the internal expectation \(V_t\) an explicit, load-bearing quantity in a
learning equation.

The impact was immediate and lasting. The prediction-error rule turned
out to be the correct description of dopamine neuron firing (Schultz,
Dayan, and Montague 1997): dopamine neurons respond to unpredicted
rewards (positive prediction error), are suppressed by
predicted-but-absent rewards (negative prediction error), and are silent
when rewards are fully predicted (zero error). The same equation that
Rescorla and Wagner derived from behavioral blocking data in 1972 became
the neural currency of reinforcement learning two decades later.

\subsubsection{From Fixed to Adaptive Learning Rates: the Kalman
Filter}\label{from-fixed-to-adaptive-learning-rates-the-kalman-filter}

Rescorla-Wagner's \(\alpha\) is fixed --- the same step size applies on
trial 1 as on trial 1000. This is cognitively unrealistic: a rational
agent should trust new observations more when they are uncertain about
the current state, and less when they are confident. Dayan, Kakade, and
Montague (2000) showed that Rescorla-Wagner is a special case of the
\textbf{Kalman filter} (Kalman 1960) obtained by setting the process
noise to zero (a static environment). In the full Kalman filter, the
learning rate becomes the Kalman Gain
\(K_t = \sigma^2_{t-1} / (\sigma^2_{t-1} + \sigma^2_{noise})\) --- a
quantity that starts high (large prior uncertainty) and shrinks as
evidence accumulates. The same update rule, but now the step size is
endogenous and principled rather than fixed and arbitrary.

\subsubsection{Hierarchical Uncertainty: the
HGF}\label{hierarchical-uncertainty-the-hgf}

The Kalman filter still assumes that \(\sigma^2_{noise}\) --- the
volatility of the environment --- is known and fixed. Real environments
are not: sometimes the rules change slowly, sometimes rapidly. Mathys et
al. (2011) introduced the \textbf{Hierarchical Gaussian Filter (HGF)} to
handle this: a hierarchy of Kalman filters where the volatility at one
level is itself tracked and updated at the level above. The learning
rate for beliefs about the world now adapts continuously to inferred
volatility, and there is a separate meta-learning rate \(\kappa\)
governing how quickly volatility estimates are revised. The HGF is,
formally, just another application of the prediction-error principle ---
but applied recursively across levels of abstraction.

The three-decade arc from R-W (1972) to the Kalman filter connection
(2000) to the HGF (2011) is a story about the same core idea ---
\emph{update in proportion to surprise} --- being successively liberated
from artificial constraints: first the prohibition on latent states,
then the constraint of a fixed learning rate, then the assumption of
known and constant volatility.

\end{tcolorbox}

\subsection{Level 4: The Bayesian Update (Static
Uncertainty)}\label{level-4-the-bayesian-update-static-uncertainty}

In RL, the learning rate \(\alpha\) is a fixed number. In Bayesian
inference, the ``learning rate'' becomes adaptive based on uncertainty.

\begin{itemize}
\item
  \emph{Concept}: The agent tracks not just the value, but the
  uncertainty in estimating that value. They weigh new evidence against
  their prior belief based on how noisy the evidence is versus how solid
  their prior is.
\item
  \emph{Traditional Formulation}: We combine a Prior distribution
  (\(\mathcal{N}(\mu_{prior}, \sigma^2_{prior})\)) with Likelihood
  (\(\mathcal{N}(Outcome, \sigma^2_{noise})\)) to get a
  Posterior.\[\text{Posterior Mean} = \frac{\sigma^2_{noise}\mu_{t-1} + \sigma^2_{t-1}Outcome_t}{\sigma^2_{noise} + \sigma^2_{t-1}}\]
\item
  \emph{Transformation to Update Rule (Kalman Gain)}: With some algebra,
  we can rearrange the posterior mean into our standard ``Update +
  Error''
  format.\[\mu_{t} = \mu_{t-1} + K_t \cdot (Outcome_t - \mu_{t-1})\]
  Here, \(K\), what we called learning rate is often called
  \(Kalman Gain\):
  \[K_t = \frac{\sigma^{2}_{t-1}}{\sigma^{2}_{t-1} + \sigma^{2}_{noise}}\]
\item
  \emph{Variable Definitions}:

  \begin{itemize}
  \tightlist
  \item
    \(\mu_t\): The estimated value (mean belief) at trial \(t\).
  \item
    \(\sigma^2_{t-1}\): The uncertainty of the agent's belief before
    seeing the outcome (Prior Variance).
  \item
    \(\sigma^2_{noise}\): The uncertainty/noise of the environment
    (Likelihood Variance).
  \item
    \(Outcome_t\): The observation at trial \(t\).
  \end{itemize}
\item
  \emph{Critique (The Stopping Problem)}:

  \begin{itemize}
  \tightlist
  \item
    The Match: This looks exactly like RL, where \(K_t\) replaces
    \(\alpha\).
  \item
    The Difference: In this model, the uncertainty \(\sigma^2_{t-1}\)
    decreases after every observation (the agent gets more confident).
    As \(\sigma^2_{t-1} \to 0\), the gain \(K_t \to 0\).
  \item
    Implication: The agent eventually stops learning because they become
    ``sure'' they know the truth. This is optimal for static worlds, but
    disastrous if the opponent changes strategy.
  \end{itemize}
\end{itemize}

\subsection{Level 5: The Kalman Filter (Dynamic
Uncertainty)}\label{level-5-the-kalman-filter-dynamic-uncertainty}

Real environments change (volatility). The Kalman Filter extends the
Bayesian update by adding a Prediction Step that prevents the agent from
becoming ``too sure.''

\begin{itemize}
\item
  \emph{Concept}: ``The world is not static; it drifts. Even if I was
  sure yesterday, my knowledge has degraded by today.''
\item
  \emph{Traditional Formulation (State Space Model)}: We assume the true
  hidden value (\(x\)) drifts over time, and that our observations
  (\(y\)) are noisy versions of that drifting value.
  \[x_t = x_{t-1} + w_t, \quad w_t \sim \mathcal{N}(0, Q)\]\[y_t = x_t + v_t, \quad v_t \sim \mathcal{N}(0, R)\]
\item
  \emph{Transformation to Update Rule (Two-Step Process)}: Unlike the
  static Bayesian agent, the Kalman agent ``inflates'' their uncertainty
  before looking at the data.

  \begin{itemize}
  \tightlist
  \item
    Predict (Inflate Uncertainty): Before seeing data, the agent assumes
    the world might have changed. They add Process Noise (\(Q\)) to
    their existing uncertainty.
    \[\sigma^2_{prediction} = \sigma^2_{previous} + Q\]
  \item
    Calculate Gain: The ``Learning Rate'' (\(K_t\)) is calculated using
    this inflated
    uncertainty.\[K_{t} = \frac{\sigma^2_{prediction}}{\sigma^2_{prediction} + R}\]
  \item
    Update (The Standard Rule): Finally, we apply the standard
    error-correction rule.
    \[\mu_{t} = \mu_{t-1} + K_t \cdot (Outcome_t - \mu_{t-1})\]
  \end{itemize}
\item
  \emph{Variable Definitions}:

  \begin{itemize}
  \item
    \(Q\): Process Noise (Volatility). The variance of the drift
    \(w_t\). High \(Q\) means the world changes largely/frequently.
  \item
    \(R\): Measurement Noise (Observation Uncertainty). Equivalent to
    \(\sigma^2_{noise}\) in Level 4.
  \item
    \(\mu_t\): The agent's estimate of the true state \(x_t\).
  \item
    \(\sigma^2_{prediction}\): The agent's uncertainty after accounting
    for volatility but before seeing the new outcome.
  \end{itemize}
\item
  \emph{Why this fixes the ``Stopping Problem''}: In Level 4,
  \(\sigma^2\) shrank to zero, causing learning to stop. Here, because
  we add \(Q\) at every step, \(\sigma^2_{prediction}\) never hits zero.
  Consequently, the learning rate \(K_t\) never hits zero. It settles at
  a dynamic equilibrium---a sweet spot where the agent stays permanently
  ``alert'' to changes without over-reacting to noise.
\end{itemize}

\subsection{Level 6: The Hierarchical Gaussian Filter
(Meta-Learning)}\label{level-6-the-hierarchical-gaussian-filter-meta-learning}

The Kalman Filter assumes the volatility (\(Q\)) is constant. But what
if the opponent plays steadily for 20 rounds, then suddenly starts
switching wildly? The HGF introduces a hierarchy to track this.

\begin{itemize}
\item
  \emph{Concept}:

  \begin{itemize}
  \item
    Level 1 (\(x_1\)): The value of the stimulus (e.g., probability of
    ``Right'').
  \item
    Level 2 (\(x_2\)): The volatility of Level 1.
  \item
    N.B. we could build more levels (e.g.~level 3 would be the
    meta-volatility, the model belief of how fast to update estimates of
    volatility)
  \end{itemize}
\item
  \emph{Transformation to Update Rule (Volatile Learning Rates)}: The
  update rule for Level 1 looks standard, but the Learning Rate is now a
  function of Level
  2.\[\mu_{1,t} = \mu_{1,t-1} + K_1 \cdot (Outcome_t - \mu_{1,t-1})\]
  Where the Gain \(K_1\) depends on the volatility estimate from Level
  2:\[K_1 \approx \frac{\exp(\mu_{2})}{\exp(\mu_{2}) + \sigma^2_{noise}}\]
\item
  \emph{Variable Definitions}:

  \begin{itemize}
  \tightlist
  \item
    \(\mu_{1}\): Estimated value (Probability).
  \item
    \(\mu_{2}\): Estimated volatility (Log-volatility).
  \item
    \(\exp(\mu_2)\): The ``phasic volatility''---how much the agent
    expects the world to change right now.
  \end{itemize}
\item
  \emph{Implication}:

  \begin{itemize}
  \tightlist
  \item
    If Level 2 detects high volatility (\(x_2 \uparrow\)), the learning
    rate \(K_1\) spikes \(\to 1\).
  \item
    If Level 2 detects stability (\(x_2 \downarrow\)), the learning rate
    \(K_1\) drops \(\to 0\).
  \item
    Meta-Learning: The agent is not just learning the value; they are
    learning how fast to learn.
  \end{itemize}
\item
  \emph{Looking Ahead}: This model represents the current
  state-of-the-art in ``Ideal Observer'' models and has had many
  interesting applications in cognitive neurosciences (REFS). While the
  math for inverting this model (Variational Bayes) is complex, the
  intuition is simple: it is a dynamic learning rate manager. Future
  chapters (not yet written) might explore the model and its
  implications and limitations.
\end{itemize}

\subsection{Level 7: Recursive Strategies (Theory of
Mind)}\label{level-7-recursive-strategies-theory-of-mind}

Here is the revised and expanded text for the final sections. I have
aligned Level 7 (ToM) with the ``Update Rule'' framework by clarifying
that while the decision is recursive, the learning is a Bayesian update
of the opponent's parameters. I also kept the Mixture Models section as
a necessary pedagogical conclusion.

Level 7: Recursive Strategies (Theory of Mind)

Finally, we move from learning about the environment to learning about
the agent. This is Theory of Mind (ToM). In our framework, this is
formally framed as Recursive Bayesian Learning.

\begin{itemize}
\item
  \emph{Concept}: The agent (\(k\)-ToM) does not simply track the
  pattern of outcomes; they attempt to track the opponent's hidden
  strategy.

  \begin{itemize}
  \tightlist
  \item
    A 0-ToM agent is a simple learner (e.g., the Biased Agent or RL).
  \item
    A 1-ToM agent simulates a 0-ToM opponent. They ask: ``If I were a
    simple learner observing this history, what would I play next?''
  \item
    A 2-ToM agent simulates a 1-ToM opponent. They ask: ``What does the
    opponent think I will do?''
  \end{itemize}
\item
  \emph{Traditional Formulation (Recursive Simulation)}: The agent holds
  a ``shadow model'' of the opponent. If I am Level \(k\), I assume my
  opponent is Level \(k-1\). I feed the game history into my internal
  simulation of them to predict their next move probability
  (\(\hat{p}_{op}\)).
\item
  \emph{Update Rule (Bayesian Meta-Learning)}: Unlike RL, where we
  update a simple value \(V\), here we update the parameters of the
  simulated opponent.
  \[\text{OpponentModel}_{new} = \text{OpponentModel}_{old} + \text{Learning Rate} \cdot (\text{Prediction Error})\]

  \begin{itemize}
  \item
    Specifics: If I am 1-ToM modeling a 0-ToM opponent (Bias), I use a
    Bayesian update to refine my estimate of their bias \(\theta\).
  \item
    Decision Value: Once the opponent model is updated, the value of my
    action \(a\) is derived from their predicted move:
    \[V_{k}(a) \propto P(\text{Opponent plays } \hat{a}_{op} | \text{Opponent is } k-1)\]
  \end{itemize}
\item
  \emph{Variable Definitions}:

  \begin{itemize}
  \tightlist
  \item
    \(k\): The sophistication level. \(k=0\) is the baseline; \(k=1\) is
    strategic; \(k=2\) is meta-strategic.
  \item
    \(\hat{a}_{op}\): The predicted action of the opponent.(Ref:
    Devaine, et al., 2014; Waade et al., 2023)
  \end{itemize}
\end{itemize}

\subsection{Handling Heterogeneity: Mixture
Models}\label{handling-heterogeneity-mixture-models}

Real behavior is rarely pure. A player might be ``Random'' when tired
but ``Strategic'' when focused. Or they might switch strategies to
confuse the opponent.

\begin{itemize}
\item
  \emph{Concept}: We assume the data is not generated by one single
  model, but by a probabilistic blend of several.
\item
  \emph{Traditional Formulation (The Weighted Likelihood)}: The
  likelihood of the observed data \(D\) is the weighted sum of the
  likelihoods from \(M\) different candidate strategies.
  \[P(D | \Theta) = \sum_{m=1}^{M} w_m \cdot P(D | \text{Model}_m, \theta_m)\]
\item
  \emph{Variable Definitions}:\(w_m\): The mixing weight (e.g., 70\% RL,
  30\% Random). Note that \(\sum w_m = 1\).\(P(D|\text{Model}_m)\): How
  well the specific strategy (e.g., RL) explains the data.
\item
  \emph{Critique} (The Complexity Trap): While realistic, these models
  are hard to fit. If we allow an agent to switch between 5 different
  strategies at any moment, the possible combinations of parameter
  values explaining the same patterns of behavior could explode, making
  the model underdefined for the data. Sometimes one can be lucky and
  integrate the ``mixed'' models into one equation (see an example here:
  https://betanalpha.github.io/assets/chapters\_html/reading\_times.html)
\end{itemize}

\subsection{Plausibility Check: Cognitive
Constraints}\label{plausibility-check-cognitive-constraints}

You probably noticed you couldn't remember every single move your
opponent made. Maybe you lost track of older trials, and recent trials
felt more important. These aren't failures: they're features of human
cognition. Our models should reflect these constraints to be cognitively
realistic. Let's see what difference these constraints make for
predictions.

\subsubsection{Memory Limitations}\label{memory-limitations}

\begin{itemize}
\tightlist
\item
  \emph{Constraint:} Humans have limited working memory and exhibit
  forgetting, often approximated by exponential decay. Perfect recall of
  long trial sequences is unrealistic.
\item
  \emph{Modeling Implication:} This favors models incorporating memory
  decay or finite history windows (like imperfect memory models or RL
  with a learning rate \textless{} 1) over perfect memory models. It
  suggests that even bias-tracking models should discount older
  information.
\end{itemize}

\subsubsection{Perseveration Tendencies}\label{perseveration-tendencies}

\begin{itemize}
\tightlist
\item
  \emph{Constraint:} People sometimes exhibit perseveration -- repeating
  a previous action, especially if it was recently successful or chosen,
  even if a different strategy might suggest otherwise. This can be
  distinct from rational ``win-stay''.
\item
  \emph{Modeling Implication:} This might be incorporated as an
  additional bias parameter influencing the choice probability (e.g., a
  small added probability of repeating the last action \(a_{t-1}\)
  regardless of outcome) or interact with feedback processing (e.g.,
  strengthening the `stay' tendency after wins).
\end{itemize}

\subsubsection{Noise and Errors}\label{noise-and-errors}

\begin{itemize}
\tightlist
\item
  \emph{Constraint:} Human behavior is inherently noisy. People make
  mistakes, have attentional lapses, press the wrong button, or
  misunderstand feedback. Behavior rarely perfectly matches a
  deterministic strategy.
\item
  \emph{Modeling Implication:} Models should almost always include a
  ``noise'' component. This can be implemented in several ways:

  \begin{itemize}
  \tightlist
  \item
    \textbf{Lapse Rate:} A probability (e.g., \(\epsilon\)) that on any
    given trial, the agent makes a random choice instead of following
    their primary strategy (as used in the Mixture Model chapter).
  \item
    \textbf{Decision Noise (Softmax):} In models where choices are based
    on comparing values (like RL), a `temperature' parameter can control
    the stochasticity. High temperature leads to more random choices,
    low temperature leads to more deterministic choices based on values.
  \item
    \textbf{Imperfect Heuristics:} Parameters within a strategy might
    reflect imperfect application (e.g., in WSLS, \(p_{stay\_win} < 1\)
    or \(p_{shift\_loss} < 1\)). This can also capture asymmetric
    responses to feedback (e.g., being more likely to shift after a loss
    than stay after a win).

    \begin{itemize}
    \tightlist
    \item
      \emph{Exploration:} Note that we talked about noise and errors,
      but random deviations can also be framed as adaptive exploration,
      allowing the agent to test actions that their current strategy
      deems suboptimal.
    \end{itemize}
  \end{itemize}
\end{itemize}

\subsubsection{A Note on Ideal
Observers}\label{a-note-on-ideal-observers}

It is important to distinguish between \textbf{Heuristic Models} (like
WSLS) and \textbf{Ideal Observer Models} (like the Kalman Filter or
HGF). Heuristics attempt to describe the \emph{process} a human uses.
Ideal observers describe the \emph{optimal computation} given the
uncertainty. When we fit models like the HGF to human data, we are
effectively asking: ``In what specific ways does the human deviate from
optimality?'' (e.g., do they overestimate how volatile the opponent
is?).

\subsection{Relationships Between
Models}\label{relationships-between-models}

It's useful to note that these candidate models aren't always entirely
distinct. Often, simpler models emerge as special cases of more complex
ones:

\begin{itemize}
\item
  A Random Choice model is like a Memory-Based model where the influence
  of memory is zero.
\item
  WSLS can be seen as a specific type of RL model with a very high
  learning rate and sensitivity only to the immediately preceding
  trial's outcome.
\item
  A 0-ToM model might resemble a Bias Tracking (moving average or RL
  model).
\end{itemize}

Recognizing these connections can guide a principled modeling approach,
starting simple and adding complexity only as needed and justified by
data or theory.

This is what we called ``model nesting'' in a future chapter. It's very
elegant in that it allows to directly compare several models in a
non-exclusive fashion, simply based on parameter values. A good example
is also provided in Betancourt's case study on reading times:
https://betanalpha.github.io/assets/chapters\_html/reading\_times.html

\subsection{Handling Heterogeneity: Mixture
Models}\label{handling-heterogeneity-mixture-models-1}

Sometimes the models cannot be nested, or we might want to capture the
possibility that different participants (or even the same participant at
different times) use different strategies. For instance, some players
might predominantly use WSLS, while others rely more on bias tracking
(moving average or RL). This is where \emph{mixture models} become
relevant (explored in detail in a future chapter).

\begin{itemize}
\item
  \emph{Concept:} Instead of assuming \emph{one} model generated all the
  data, a mixture model assumes the data is a probabilistic blend from
  \emph{multiple} candidate models (e.g., 70\% of choices from WSLS,
  30\% from Random Bias).
\item
  \emph{Purpose:} Allows capturing heterogeneity within or across
  individuals without needing to know \emph{a priori} which strategy was
  used on which trial or by which person. The model estimates the
  \emph{probability} that each data point came from each component
  strategy.
\item
  \emph{Challenge:} Mixture models often require substantial data to
  reliably distinguish between components and estimate their mixing
  proportions.
\end{itemize}

\subsection{Cognitive Modeling vs.~Traditional Statistical Approaches
(e.g.,
GLM)}\label{cognitive-modeling-vs.-traditional-statistical-approaches-e.g.-glm}

How does this modeling approach differ from standard statistical
analyses you might have learned, like ANOVAs or the General Linear Model
- GLM?

\begin{itemize}
\tightlist
\item
  \textbf{Focus:} GLM approaches typically focus on identifying
  \emph{statistical effects}: Does factor X significantly influence
  outcome Y? (e.g., Does the opponent's strategy affect the player's win
  rate?). Cognitive modeling focuses on identifying the underlying
  \emph{process or mechanism}: \emph{How} does the opponent's strategy
  lead to changes in the player's choices via specific computations
  (like learning, memory updating, or strategic reasoning)?
\item
  \textbf{Theory:} Cognitive models are usually derived from theories
  about mental processes. GLMs are more general statistical tools, often
  used agnostically regarding the specific cognitive mechanism.
\item
  \textbf{Parameters:} Cognitive models estimate parameters that often
  have direct psychological interpretations (e.g., learning rate, memory
  decay, decision threshold, bias weight). GLM parameters represent
  statistical associations (e.g., regression coefficients).
\item
  \textbf{Data Level:} Cognitive models often predict behavior at the
  trial level (e.g., predicting the choice on trial \emph{t} based on
  history up to \emph{t-1}). GLM analyses often aggregate data (e.g.,
  comparing average win rates across conditions).
\item
  \textbf{Prediction vs.~Explanation:} While both aim to explain data,
  cognitive modeling often places a stronger emphasis on generating the
  observed behavior pattern from the hypothesized mechanism, allowing
  for simulation and prediction of fine-grained details.
\end{itemize}

\emph{Example Revisited:} In the Matching Pennies game: * A GLM approach
might test if
\texttt{Payoff\ \textasciitilde{}\ BotStrategy\ *\ Role\ +\ (1\textbar{}ID)}
shows a significant effect of \texttt{BotStrategy}. * A cognitive
modeling approach would fit different strategy models (WSLS, RL, etc.)
to the choice data and compare them (using methods from Ch 7) to see
which \emph{mechanism} best explains the choices made against different
bots, potentially revealing \emph{why} performance differs (e.g., due to
changes in estimated learning rates or strategy weights).

Both approaches are valuable, but cognitive modeling aims for a deeper,
mechanistic level of explanation about the underlying cognitive
processes.

\section{Conclusion: From Observations to Verbal
Theories}\label{conclusion-from-observations-to-verbal-theories}

This chapter took us from observing behavior in a specific task -- the
Matching Pennies game -- to the crucial stage of formulating initial
theories about the cognitive processes involved. We explored how
analyzing gameplay data, considering participant reports, applying
cognitive principles (like memory limits and error proneness), and
contrasting different potential strategies (Random, WSLS, Memory-based,
RL, k-ToM) helps us generate plausible \emph{verbal models}.

We saw that the path from raw behavior to a testable model involves
significant abstraction and simplification. We also highlighted the
importance of distinguishing between the participant's experience and
the researcher's theoretical stance, and how cognitive modeling differs
from traditional statistical approaches by focusing on underlying
mechanisms.

You now have a conceptual map of candidate models and understand why
cognitive constraints matter. But verbal descriptions like
`win-stay-lose-shift' hide crucial ambiguities: Does `stay' mean always
stay or usually stay? How do we handle the first trial? In Chapter 3,
you'll implement these models in code, forcing you to make every
assumption explicit. This is where modeling becomes rigorous and often
reveals that our verbal intuitions were vaguer than we thought.

The next chapter, ``From verbal descriptions to formal models,'' tackles
exactly this challenge. We will take some of the candidate models
discussed here (like Random Choice and WSLS) and translate them into
precise mathematical algorithms and R functions. This formalization will
force us to be explicit about our assumptions and enable us to simulate
agent behavior. That simulation is itself a prerequisite: in Ch. 5 we
flip the direction of the pipeline --- instead of generating data from
known parameters, we will \emph{infer} parameters from observed data.
And in Ch. 6, before we trust any of those estimates, we will put the
entire inference procedure through a rigorous validation battery.

\bookmarksetup{startatroot}

\chapter{From verbal descriptions to formal
models}\label{from-verbal-descriptions-to-formal-models}

\begin{quote}
\textbf{📍 Where we are in the Bayesian modeling workflow:} We have
verbal theories (Ch. 2). This chapter turns them into runnable code ---
\textbf{forward models} that generate data \emph{from} parameters. This
is a prerequisite for Ch. 5, where we reverse the direction: given
observed data, we will \emph{infer} the parameters. You cannot run that
inversion without first knowing how to run the model forward.
\end{quote}

In Chapter 1, we explored how the ``physics'' of a pizza
stone---Newton's Law of Cooling---could generate temperature data. We
called this a \textbf{generative model}. In Chapter 2, you observed your
own behavior in the Matching Pennies game and proposed verbal theories
like ``Win-Stay-Lose-Shift'' or ``Random Choice.''

Now, we bridge the gap. In this chapter, we will build ``models of the
mind.'' We will translate the verbal strategies from Chapter 2 into
precise algorithms---creating agents, or \textbf{Synthetic Subjects}.

Just as we used a formula to simulate the pizza stone heating up, we
will write code to simulate a cognitive agent playing a game. If our
code can generate choice data that looks like human data, we have a
candidate model for how the mind works. If it cannot, our theory is
wrong.

\section{Learning Objectives}\label{learning-objectives-2}

This chapter focuses on translating the verbal theories from Chapter 2
into concrete, testable computational models. After completing this
chapter, you will be able to:

\begin{itemize}
\tightlist
\item
  \textbf{Formalize Verbal Models:} Translate abstract strategy
  descriptions (like ``Random Choice'' or ``Win-Stay-Lose-Shift'') into
  precise algorithms.
\item
  \textbf{Implement Agent-Based Simulations:} Create simple software
  ``agents'' that execute these algorithms within a simulated
  environment (the Matching Pennies game).
\item
  \textbf{Run Virtual Experiments:} Use simulation to see how different
  strategies perform against different opponents \emph{before} we ever
  touch real data.
\item
  \textbf{Analyze Simulated Behavior:} Generate and visualize simulation
  data to understand how different strategies perform and interact under
  controlled conditions.
\item
  \textbf{Prepare for Model Fitting:} Appreciate why simulating models
  is a crucial step prior to estimating parameters from empirical data
  (covered in Chapter 5).
\end{itemize}

\section{The Value of Formalization and
Simulation}\label{the-value-of-formalization-and-simulation}

Verbal models are comfortable but dangerous. In Chapter 2, we might have
said: \emph{``The agent stays after a win.''}

But formalization forces us to answer uncomfortable questions, for
instance:

\begin{itemize}
\item
  \textbf{The Persistence Question:} If they win 10 times, do they stay
  forever? Or is there a limit?
\item
  \textbf{The Error Question:} Do they \emph{always} stay, or do they
  sometimes press the wrong button by mistake?
\item
  \textbf{The First Trial Question:} What do they do before they have
  any feedback?
\end{itemize}

Writing code forces us to strip away ambiguity. It reveals hidden
assumptions and forces us to be explicit about the mechanisms driving
behavior. However, the value of this process extends beyond just
clarity. Once formalized, we can create computational agents that act
out the strategy. Simulation allows us to:

\begin{itemize}
\item
  \textbf{Understand Emergent Behavior}: See the consequences of our
  model's rules over many trials or interactions. Does the strategy lead
  to expected patterns? Are there surprising outcomes that weren't
  obvious from the verbal description?
\item
  \textbf{Explore Parameter Space}: How does behavior change if we alter
  a model parameter (e.g., the bias in a random agent, or the stay
  probability in WSLS)?
\item
  \textbf{Generate Testable Predictions}: What specific behavioral
  patterns (e.g., choice sequences) does the model predict under
  different conditions? These predictions can then be compared to real
  data.
\item
  \textbf{Debug Model Logic}: Does the simulated behavior make sense? If
  not, it might indicate a flaw in our formalization of the verbal
  theory.
\end{itemize}

In this chapter, we will:

\begin{enumerate}
\def\labelenumi{\arabic{enumi}.}
\item
  \textbf{Formalize:} Define precise algorithms for two key strategies
  (Random Choice and WSLS).
\item
  \textbf{Implement:} Code these algorithms as reusable R functions.
\item
  \textbf{Simulate:} Run these agents in a ``Virtual Arena'' against
  different opponents.
\item
  \textbf{Scale:} Simulate many agents to understand average behavior
  and visualize the results to compare strategy effectiveness.
\end{enumerate}

\section{Defining general conditions}\label{defining-general-conditions}

\begin{Shaded}
\begin{Highlighting}[]
\CommentTok{\# Load necessary packages}
\CommentTok{\# patchwork helps combine multiple plots}
\NormalTok{pacman}\SpecialCharTok{::}\FunctionTok{p\_load}\NormalTok{(tidyverse, patchwork, here)}

\CommentTok{\# Number of trials per simulation}
\NormalTok{trials }\OtherTok{\textless{}{-}} \DecValTok{120}  

\CommentTok{\# Number of agents to simulate}
\NormalTok{agents }\OtherTok{\textless{}{-}} \DecValTok{100}  

\CommentTok{\# Set to TRUE to rerun everything (takes time!), FALSE to load saved results.}
\NormalTok{regenerate\_simulations }\OtherTok{\textless{}{-}} \ConstantTok{FALSE} 

\CommentTok{\# Optional: Set random seed for reproducibility}
\CommentTok{\# set.seed(123)}
\end{Highlighting}
\end{Shaded}

\section{Implementing a Random Agent}\label{implementing-a-random-agent}

Our simplest model from Chapter 2 was the Random Choice agent. In
cognitive science, this serves as our \textbf{Null Hypothesis}.

A Random Agent is not just ``acting crazily''; it represents a specific,
strong theoretical claim: \textbf{The agent has no memory and no
learning.} It does not care about the past (history) or the future
(planning). It only has a bias.

\textbf{Algorithm:} On every trial \(t\), independently of what happened
at \(t-1\), choose `right' (1) with probability \(\theta\) and `left'
(0) with probability \(1-\theta\).

Let's implement this synthetic subject.

\begin{Shaded}
\begin{Highlighting}[]
\CommentTok{\# {-}{-}{-} Simulation Parameters {-}{-}{-}}
\NormalTok{trials }\OtherTok{\textless{}{-}} \DecValTok{120} \CommentTok{\# Number of trials for a single simulation run}

\CommentTok{\# {-}{-}{-} Basic Implementation {-}{-}{-}}
\CommentTok{\# Simulate an unbiased random agent (rate = 0.5)}
\NormalTok{rate\_unbiased }\OtherTok{\textless{}{-}} \FloatTok{0.5}
\NormalTok{choices\_unbiased }\OtherTok{\textless{}{-}} \FunctionTok{rbinom}\NormalTok{(trials, }\AttributeTok{size =} \DecValTok{1}\NormalTok{, }\AttributeTok{prob =}\NormalTok{ rate\_unbiased)}
\CommentTok{\# rbinom generates random draws from a binomial distribution.}
\CommentTok{\# size = 1 makes it a Bernoulli trial (0 or 1 outcome).}
\CommentTok{\# prob = rate specifies the probability of getting a 1.}

\CommentTok{\# Simulate a biased random agent (rate = 0.8, chooses \textquotesingle{}right\textquotesingle{} 80\% of the time)}
\NormalTok{rate\_biased }\OtherTok{\textless{}{-}} \FloatTok{0.8}
\NormalTok{choices\_biased }\OtherTok{\textless{}{-}} \FunctionTok{rbinom}\NormalTok{(trials, }\AttributeTok{size =} \DecValTok{1}\NormalTok{, }\AttributeTok{prob =}\NormalTok{ rate\_biased)}

\CommentTok{\# {-}{-}{-} Visualization 1: Raw Choices {-}{-}{-}}
\CommentTok{\# Create data frames for plotting}
\NormalTok{d\_unbiased }\OtherTok{\textless{}{-}} \FunctionTok{tibble}\NormalTok{(}\AttributeTok{trial =} \DecValTok{1}\SpecialCharTok{:}\NormalTok{trials, }\AttributeTok{choice =}\NormalTok{ choices\_unbiased, }\AttributeTok{rate =}\NormalTok{ rate\_unbiased)}
\NormalTok{d\_biased }\OtherTok{\textless{}{-}} \FunctionTok{tibble}\NormalTok{(}\AttributeTok{trial =} \DecValTok{1}\SpecialCharTok{:}\NormalTok{trials, }\AttributeTok{choice =}\NormalTok{ choices\_biased, }\AttributeTok{rate =}\NormalTok{ rate\_biased)}

\CommentTok{\# Plot raw choice sequence (0s and 1s)}
\NormalTok{p1 }\OtherTok{\textless{}{-}} \FunctionTok{ggplot}\NormalTok{(d\_unbiased, }\FunctionTok{aes}\NormalTok{(}\AttributeTok{x =}\NormalTok{ trial, }\AttributeTok{y =}\NormalTok{ choice)) }\SpecialCharTok{+}
  \FunctionTok{geom\_line}\NormalTok{(}\AttributeTok{color =} \StringTok{"blue"}\NormalTok{) }\SpecialCharTok{+} \FunctionTok{geom\_point}\NormalTok{(}\AttributeTok{color =} \StringTok{"blue"}\NormalTok{, }\AttributeTok{size =} \DecValTok{1}\NormalTok{) }\SpecialCharTok{+}
  \FunctionTok{labs}\NormalTok{(}\AttributeTok{title =} \StringTok{"Unbiased Random Agent (Rate = 0.5)"}\NormalTok{, }\AttributeTok{y =} \StringTok{"Choice (0/1)"}\NormalTok{) }\SpecialCharTok{+} \FunctionTok{theme\_classic}\NormalTok{() }\SpecialCharTok{+} \FunctionTok{ylim}\NormalTok{(}\SpecialCharTok{{-}}\FloatTok{0.1}\NormalTok{, }\FloatTok{1.1}\NormalTok{)}

\NormalTok{p2 }\OtherTok{\textless{}{-}} \FunctionTok{ggplot}\NormalTok{(d\_biased, }\FunctionTok{aes}\NormalTok{(}\AttributeTok{x =}\NormalTok{ trial, }\AttributeTok{y =}\NormalTok{ choice)) }\SpecialCharTok{+}
  \FunctionTok{geom\_line}\NormalTok{(}\AttributeTok{color =} \StringTok{"red"}\NormalTok{) }\SpecialCharTok{+} \FunctionTok{geom\_point}\NormalTok{(}\AttributeTok{color =} \StringTok{"red"}\NormalTok{, }\AttributeTok{size =} \DecValTok{1}\NormalTok{) }\SpecialCharTok{+}
  \FunctionTok{labs}\NormalTok{(}\AttributeTok{title =} \StringTok{"Biased Random Agent (Rate = 0.8)"}\NormalTok{, }\AttributeTok{y =} \StringTok{"Choice (0/1)"}\NormalTok{) }\SpecialCharTok{+} \FunctionTok{theme\_classic}\NormalTok{() }\SpecialCharTok{+} \FunctionTok{ylim}\NormalTok{(}\SpecialCharTok{{-}}\FloatTok{0.1}\NormalTok{, }\FloatTok{1.1}\NormalTok{)}

\CommentTok{\# Show plots side{-}by{-}side}
\FunctionTok{print}\NormalTok{(p1 }\SpecialCharTok{+}\NormalTok{ p2 }\SpecialCharTok{+} \FunctionTok{plot\_layout}\NormalTok{(}\AttributeTok{guides =} \StringTok{"collect"}\NormalTok{))}
\end{Highlighting}
\end{Shaded}

\begin{center}
\includegraphics[width=0.8\linewidth,height=\textheight,keepaspectratio]{03-BuildingFormalModels_files/figure-pdf/ch04_implement_random_agent-1.pdf}
\end{center}

This first visualization shows the behavior of a purely random agent -
one that chooses between options with equal probability (rate = 0.5).
Looking at the jagged line jumping between 0 and 1, we can see that the
agent's choices appear truly random, with no discernible pattern. This
represents what we might expect from a player who is deliberately trying
to be unpredictable in the matching pennies game. However, this raw
choice plot can be hard to interpret. A more informative way to look at
the agent's behavior is to examine how its average rate of choosing
option 1 evolves over time:

\begin{Shaded}
\begin{Highlighting}[]
\CommentTok{\# {-}{-}{-} Visualization 2: Cumulative Rate {-}{-}{-}}
\CommentTok{\# Calculate the cumulative proportion of \textquotesingle{}right\textquotesingle{} choices (1s)}
\NormalTok{d\_unbiased }\OtherTok{\textless{}{-}}\NormalTok{ d\_unbiased }\SpecialCharTok{\%\textgreater{}\%} \FunctionTok{mutate}\NormalTok{(}\AttributeTok{cumulative\_rate =} \FunctionTok{cumsum}\NormalTok{(choice) }\SpecialCharTok{/} \FunctionTok{row\_number}\NormalTok{())}
\NormalTok{d\_biased }\OtherTok{\textless{}{-}}\NormalTok{ d\_biased }\SpecialCharTok{\%\textgreater{}\%} \FunctionTok{mutate}\NormalTok{(}\AttributeTok{cumulative\_rate =} \FunctionTok{cumsum}\NormalTok{(choice) }\SpecialCharTok{/} \FunctionTok{row\_number}\NormalTok{())}

\CommentTok{\# Combine data for comparison plot}
\NormalTok{d\_combined }\OtherTok{\textless{}{-}} \FunctionTok{bind\_rows}\NormalTok{(d\_unbiased, d\_biased) }\SpecialCharTok{\%\textgreater{}\%}
                \FunctionTok{mutate}\NormalTok{(}\AttributeTok{Agent =} \FunctionTok{paste}\NormalTok{(}\StringTok{"Rate ="}\NormalTok{, rate))}

\CommentTok{\# Plot cumulative rates}
\NormalTok{p3 }\OtherTok{\textless{}{-}} \FunctionTok{ggplot}\NormalTok{(d\_combined, }\FunctionTok{aes}\NormalTok{(}\AttributeTok{x =}\NormalTok{ trial, }\AttributeTok{y =}\NormalTok{ cumulative\_rate, }\AttributeTok{color =}\NormalTok{ Agent)) }\SpecialCharTok{+}
  \FunctionTok{geom\_line}\NormalTok{(}\AttributeTok{size =} \DecValTok{1}\NormalTok{) }\SpecialCharTok{+}
  \FunctionTok{geom\_hline}\NormalTok{(}\AttributeTok{yintercept =} \FloatTok{0.5}\NormalTok{, }\AttributeTok{linetype =} \StringTok{"dashed"}\NormalTok{, }\AttributeTok{color =} \StringTok{"black"}\NormalTok{) }\SpecialCharTok{+} \CommentTok{\# Reference line for 50\%}
  \FunctionTok{ylim}\NormalTok{(}\DecValTok{0}\NormalTok{, }\DecValTok{1}\NormalTok{) }\SpecialCharTok{+}
  \FunctionTok{labs}\NormalTok{(}
    \AttributeTok{title =} \StringTok{"Cumulative Choice Rate Over Time"}\NormalTok{,}
    \AttributeTok{subtitle =} \StringTok{"Shows convergence towards the agent\textquotesingle{}s underlying bias"}\NormalTok{,}
    \AttributeTok{x =} \StringTok{"Trial Number"}\NormalTok{,}
    \AttributeTok{y =} \StringTok{"Cumulative Proportion of Choosing \textquotesingle{}Right\textquotesingle{}"}\NormalTok{,}
    \AttributeTok{color =} \StringTok{"Agent Type"}
\NormalTok{  ) }\SpecialCharTok{+}
  \FunctionTok{theme\_classic}\NormalTok{()}

\FunctionTok{print}\NormalTok{(p3)}
\end{Highlighting}
\end{Shaded}

\begin{center}
\includegraphics[width=0.8\linewidth,height=\textheight,keepaspectratio]{03-BuildingFormalModels_files/figure-pdf/ch04_cumulative_viz-1.pdf}
\end{center}

We can clearly see how bias affects choice behavior. The unbiased agent
(rate = 0.5) stabilizes around choosing each option equally often, while
the biased agent (rate = 0.8) shows a strong preference for option 1,
choosing it approximately 80\% of the time. This comparison helps us
understand how we might detect biases in real players' behavior -
consistent deviation from 50-50 choice proportions could indicate an
underlying preference or strategy.

\subsection{Encapsulating the Agent in a
Function}\label{encapsulating-the-agent-in-a-function}

To make our agent reusable, we'll wrap the logic in a function. This is
good practice for building more complex simulations.

Note that the function is preceded by a number of commented lines. They
are called Roxygen comments, a special documentation format used in R
packages. By using these comments, we keep the code and its manual in
the same place (and if you build a package, the format allows for
automatically generating the help pages).

\subsubsection{How to read it:}\label{how-to-read-it}

\begin{itemize}
\item
  Description: The first lines describe what the function does.
\item
  \texttt{@param}: Describes the inputs (arguments) the function expects
  (e.g., n\_trials, rate).
\item
  \texttt{@return}: Describes the output the function produces (e.g., a
  vector of 0s and 1s).
\end{itemize}

\begin{Shaded}
\begin{Highlighting}[]
\CommentTok{\#\textquotesingle{} Random Agent Function}
\CommentTok{\#\textquotesingle{}}
\CommentTok{\#\textquotesingle{} Generates a sequence of choices based on a fixed probability (rate).}
\CommentTok{\#\textquotesingle{} Includes optional noise causing random 50/50 choices.}
\CommentTok{\#\textquotesingle{}}
\CommentTok{\#\textquotesingle{} @param n\_trials Integer, the number of choices to generate.}
\CommentTok{\#\textquotesingle{} @param rate Numeric, the probability (0{-}1) of choosing option 1 (\textquotesingle{}right\textquotesingle{}).}
\CommentTok{\#\textquotesingle{} @param noise Numeric, the probability (0{-}1) of making a random 50/50 choice}
\CommentTok{\#\textquotesingle{}   instead of following the rate. Default is 0 (no noise).}
\CommentTok{\#\textquotesingle{}}
\CommentTok{\#\textquotesingle{} @return A numeric vector of choices (0s and 1s).}
\CommentTok{\#\textquotesingle{}}
\NormalTok{RandomAgent\_f }\OtherTok{\textless{}{-}} \ControlFlowTok{function}\NormalTok{(n\_trials, }\AttributeTok{rate =} \FloatTok{0.5}\NormalTok{, }\AttributeTok{noise =} \DecValTok{0}\NormalTok{) \{}
  \CommentTok{\# Input validation}
  \ControlFlowTok{if}\NormalTok{ (}\SpecialCharTok{!}\FunctionTok{is.numeric}\NormalTok{(rate) }\SpecialCharTok{||}\NormalTok{ rate }\SpecialCharTok{\textless{}} \DecValTok{0} \SpecialCharTok{||}\NormalTok{ rate }\SpecialCharTok{\textgreater{}} \DecValTok{1}\NormalTok{) \{}
    \FunctionTok{stop}\NormalTok{(}\StringTok{"Rate must be a probability between 0 and 1."}\NormalTok{)}
\NormalTok{  \}}
  \ControlFlowTok{if}\NormalTok{ (}\SpecialCharTok{!}\FunctionTok{is.numeric}\NormalTok{(noise) }\SpecialCharTok{||}\NormalTok{ noise }\SpecialCharTok{\textless{}} \DecValTok{0} \SpecialCharTok{||}\NormalTok{ noise }\SpecialCharTok{\textgreater{}} \DecValTok{1}\NormalTok{) \{}
    \FunctionTok{stop}\NormalTok{(}\StringTok{"Noise must be a probability between 0 and 1."}\NormalTok{)}
\NormalTok{  \}}

  \CommentTok{\# Generate base choices according to the rate}
\NormalTok{  choices }\OtherTok{\textless{}{-}} \FunctionTok{rbinom}\NormalTok{(n\_trials, }\AttributeTok{size =} \DecValTok{1}\NormalTok{, }\AttributeTok{prob =}\NormalTok{ rate)}

  \CommentTok{\# Apply noise: identify trials where noise occurs}
\NormalTok{  noise\_trials }\OtherTok{\textless{}{-}} \FunctionTok{rbinom}\NormalTok{(n\_trials, }\AttributeTok{size =} \DecValTok{1}\NormalTok{, }\AttributeTok{prob =}\NormalTok{ noise) }\SpecialCharTok{==} \DecValTok{1}

  \CommentTok{\# Replace choices with random 50/50 on noise trials}
  \ControlFlowTok{if}\NormalTok{ (}\FunctionTok{any}\NormalTok{(noise\_trials)) \{}
\NormalTok{    choices[noise\_trials] }\OtherTok{\textless{}{-}} \FunctionTok{rbinom}\NormalTok{(}\FunctionTok{sum}\NormalTok{(noise\_trials), }\AttributeTok{size =} \DecValTok{1}\NormalTok{, }\AttributeTok{prob =} \FloatTok{0.5}\NormalTok{)}
\NormalTok{  \}}

  \FunctionTok{return}\NormalTok{(choices)}
\NormalTok{\}}

\CommentTok{\# Example usage:}
\NormalTok{example\_choices }\OtherTok{\textless{}{-}} \FunctionTok{RandomAgent\_f}\NormalTok{(}\AttributeTok{n\_trials =} \DecValTok{10}\NormalTok{, }\AttributeTok{rate =} \FloatTok{0.7}\NormalTok{, }\AttributeTok{noise =} \FloatTok{0.1}\NormalTok{)}

\FunctionTok{print}\NormalTok{(}\FunctionTok{paste}\NormalTok{(}\StringTok{"Example choices (rate=0.7, noise=0.1):"}\NormalTok{, }\FunctionTok{paste}\NormalTok{(example\_choices, }\AttributeTok{collapse =} \StringTok{" "}\NormalTok{)))}
\end{Highlighting}
\end{Shaded}

\begin{verbatim}
[1] "Example choices (rate=0.7, noise=0.1): 1 1 1 1 0 1 1 1 0 1"
\end{verbatim}

\section{Implementing a Win-Stay-Lose-Shift (WSLS)
Agent}\label{implementing-a-win-stay-lose-shift-wsls-agent}

The Random Agent was memoryless. The Win-Stay-Lose-Shift (WSLS) agent is
our first step into \textbf{strategic behavior}: it reacts to the
immediate past.

\textbf{Algorithm:}

\begin{enumerate}
\def\labelenumi{\arabic{enumi}.}
\item
  \textbf{Initialize:} Make a random guess on the first trial.
\item
  \textbf{Loop (Trial \(t > 1\)):}

  \begin{itemize}
  \item
    Check feedback from trial \(t-1\).
  \item
    If \textbf{Win}: Repeat the previous choice.
  \item
    If \textbf{Loss}: Switch to the opposite choice.
  \end{itemize}
\item
  \textbf{Noise (The ``Trembling Hand''):} With a small probability
  \texttt{noise}, ignore the strategy and pick randomly. This accounts
  for attention lapses or motor errors (often called ``trembling hand''
  in game theory).
\end{enumerate}

\subsection{Implementing the WSLS
Function}\label{implementing-the-wsls-function}

\emph{Note on Programming:} In R, we usually prefer vectorization (doing
everything at once) over loops. However, many cognitive models are
inherently \textbf{sequential}: what I do now depends on what I just
saw. Therefore, we use an explicit \texttt{if/else} logic here to mirror
the cognitive process.

\begin{Shaded}
\begin{Highlighting}[]
\CommentTok{\#\textquotesingle{} Win{-}Stay{-}Lose{-}Shift Agent Function}
\CommentTok{\#\textquotesingle{}}
\CommentTok{\#\textquotesingle{} Determines the next choice based on the previous choice and its outcome (feedback).}
\CommentTok{\#\textquotesingle{} Includes optional noise.}
\CommentTok{\#\textquotesingle{}}
\CommentTok{\#\textquotesingle{} @param prevChoice Numeric, the agent\textquotesingle{}s choice on the previous trial (0 or 1).}
\CommentTok{\#\textquotesingle{} @param feedback Numeric, the outcome of the previous trial (1 for win, 0 for loss).}
\CommentTok{\#\textquotesingle{} @param noise Numeric, the probability (0{-}1) of making a random 50/50 choice. Default is 0.}
\CommentTok{\#\textquotesingle{}}
\CommentTok{\#\textquotesingle{} @return Numeric, the agent\textquotesingle{}s next choice (0 or 1).}
\CommentTok{\#\textquotesingle{}}
\NormalTok{WSLSAgent\_f }\OtherTok{\textless{}{-}} \ControlFlowTok{function}\NormalTok{(prevChoice, feedback, }\AttributeTok{noise =} \DecValTok{0}\NormalTok{) \{}
  \CommentTok{\# Input validation}
  \ControlFlowTok{if}\NormalTok{ (}\SpecialCharTok{!}\NormalTok{prevChoice }\SpecialCharTok{\%in\%} \FunctionTok{c}\NormalTok{(}\DecValTok{0}\NormalTok{, }\DecValTok{1}\NormalTok{)) }\FunctionTok{stop}\NormalTok{(}\StringTok{"Previous choice must be 0 or 1."}\NormalTok{)}
  \ControlFlowTok{if}\NormalTok{ (}\SpecialCharTok{!}\NormalTok{feedback }\SpecialCharTok{\%in\%} \FunctionTok{c}\NormalTok{(}\DecValTok{0}\NormalTok{, }\DecValTok{1}\NormalTok{)) }\FunctionTok{stop}\NormalTok{(}\StringTok{"Feedback must be 0 or 1."}\NormalTok{)}
  \ControlFlowTok{if}\NormalTok{ (}\SpecialCharTok{!}\FunctionTok{is.numeric}\NormalTok{(noise) }\SpecialCharTok{||}\NormalTok{ noise }\SpecialCharTok{\textless{}} \DecValTok{0} \SpecialCharTok{||}\NormalTok{ noise }\SpecialCharTok{\textgreater{}} \DecValTok{1}\NormalTok{) }\FunctionTok{stop}\NormalTok{(}\StringTok{"Noise must be a probability between 0 and 1."}\NormalTok{)}

  \CommentTok{\# Core WSLS logic:}
  \CommentTok{\# If feedback is 1 (win), stay: choice = prevChoice}
  \CommentTok{\# If feedback is 0 (loss), shift: choice = 1 {-} prevChoice}
\NormalTok{  choice }\OtherTok{\textless{}{-}} \FunctionTok{ifelse}\NormalTok{(feedback }\SpecialCharTok{==} \DecValTok{1}\NormalTok{, prevChoice, }\DecValTok{1} \SpecialCharTok{{-}}\NormalTok{ prevChoice)}

  \CommentTok{\# Apply noise if specified}
  \ControlFlowTok{if}\NormalTok{ (noise }\SpecialCharTok{\textgreater{}} \DecValTok{0} \SpecialCharTok{\&\&} \FunctionTok{runif}\NormalTok{(}\DecValTok{1}\NormalTok{) }\SpecialCharTok{\textless{}}\NormalTok{ noise) \{}
    \CommentTok{\# Override with a random 50/50 choice}
\NormalTok{    choice }\OtherTok{\textless{}{-}} \FunctionTok{sample}\NormalTok{(}\FunctionTok{c}\NormalTok{(}\DecValTok{0}\NormalTok{, }\DecValTok{1}\NormalTok{), }\DecValTok{1}\NormalTok{)}
\NormalTok{  \}}

  \FunctionTok{return}\NormalTok{(choice)}
\NormalTok{\}}

\CommentTok{\# Example usage:}
\CommentTok{\# Won previous trial (feedback=1) after choosing 1:}
\FunctionTok{print}\NormalTok{(}\FunctionTok{paste}\NormalTok{(}\StringTok{"Next choice after win (chose 1):"}\NormalTok{, }\FunctionTok{WSLSAgent\_f}\NormalTok{(}\AttributeTok{prevChoice =} \DecValTok{1}\NormalTok{, }\AttributeTok{feedback =} \DecValTok{1}\NormalTok{)))}
\end{Highlighting}
\end{Shaded}

\begin{verbatim}
[1] "Next choice after win (chose 1): 1"
\end{verbatim}

\begin{Shaded}
\begin{Highlighting}[]
\CommentTok{\# Lost previous trial (feedback=0) after choosing 1:}
\FunctionTok{print}\NormalTok{(}\FunctionTok{paste}\NormalTok{(}\StringTok{"Next choice after loss (chose 1):"}\NormalTok{, }\FunctionTok{WSLSAgent\_f}\NormalTok{(}\AttributeTok{prevChoice =} \DecValTok{1}\NormalTok{, }\AttributeTok{feedback =} \DecValTok{0}\NormalTok{)))}
\end{Highlighting}
\end{Shaded}

\begin{verbatim}
[1] "Next choice after loss (chose 1): 0"
\end{verbatim}

\subsection{Simulating WSLS
vs.~Opponents}\label{simulating-wsls-vs.-opponents}

Now, let's simulate our WSLS agent playing against different opponents
to see how it performs.

\begin{Shaded}
\begin{Highlighting}[]
\CommentTok{\# {-}{-}{-} Simulation Setup {-}{-}{-}}
\NormalTok{trials }\OtherTok{\textless{}{-}} \DecValTok{120}

\CommentTok{\# {-}{-}{-} Simulation 1: WSLS vs. Biased Random Agent {-}{-}{-}}
\FunctionTok{cat}\NormalTok{(}\StringTok{"Simulating WSLS vs. Biased Random (Rate = 0.8)...}\SpecialCharTok{\textbackslash{}n}\StringTok{"}\NormalTok{)}
\end{Highlighting}
\end{Shaded}

\begin{verbatim}
Simulating WSLS vs. Biased Random (Rate = 0.8)...
\end{verbatim}

\begin{Shaded}
\begin{Highlighting}[]
\CommentTok{\# Opponent always chooses \textquotesingle{}right\textquotesingle{} with 80\% probability}
\NormalTok{opponent\_choices\_biased }\OtherTok{\textless{}{-}} \FunctionTok{RandomAgent\_f}\NormalTok{(}\AttributeTok{n\_trials =}\NormalTok{ trials, }\AttributeTok{rate =} \FloatTok{0.8}\NormalTok{, }\AttributeTok{noise =} \DecValTok{0}\NormalTok{)}

\CommentTok{\# Initialize vectors for WSLS agent}
\NormalTok{wsls\_choices\_vs\_biased }\OtherTok{\textless{}{-}} \FunctionTok{rep}\NormalTok{(}\ConstantTok{NA}\NormalTok{, trials)}
\NormalTok{feedback\_vs\_biased }\OtherTok{\textless{}{-}} \FunctionTok{rep}\NormalTok{(}\ConstantTok{NA}\NormalTok{, trials)}

\CommentTok{\# First choice is random}
\NormalTok{wsls\_choices\_vs\_biased[}\DecValTok{1}\NormalTok{] }\OtherTok{\textless{}{-}} \FunctionTok{sample}\NormalTok{(}\FunctionTok{c}\NormalTok{(}\DecValTok{0}\NormalTok{, }\DecValTok{1}\NormalTok{), }\DecValTok{1}\NormalTok{)}

\CommentTok{\# Simulation loop}
\ControlFlowTok{for}\NormalTok{ (t }\ControlFlowTok{in} \DecValTok{2}\SpecialCharTok{:}\NormalTok{trials) \{}
  \CommentTok{\# Determine feedback from previous trial (WSLS wins if choices match opponent)}
  \CommentTok{\# Note: In Matching Pennies, win condition depends on role (Matcher/Hider).}
  \CommentTok{\# Here we assume the WSLS agent is the \textquotesingle{}Matcher\textquotesingle{} for simplicity.}
\NormalTok{  prev\_feedback }\OtherTok{\textless{}{-}} \FunctionTok{ifelse}\NormalTok{(wsls\_choices\_vs\_biased[t }\SpecialCharTok{{-}} \DecValTok{1}\NormalTok{] }\SpecialCharTok{==}\NormalTok{ opponent\_choices\_biased[t }\SpecialCharTok{{-}} \DecValTok{1}\NormalTok{], }\DecValTok{1}\NormalTok{, }\DecValTok{0}\NormalTok{)}
\NormalTok{  feedback\_vs\_biased[t }\SpecialCharTok{{-}} \DecValTok{1}\NormalTok{] }\OtherTok{\textless{}{-}}\NormalTok{ prev\_feedback }\CommentTok{\# Store feedback for analysis}

  \CommentTok{\# WSLS agent makes choice based on previous feedback}
\NormalTok{  wsls\_choices\_vs\_biased[t] }\OtherTok{\textless{}{-}} \FunctionTok{WSLSAgent\_f}\NormalTok{(}
    \AttributeTok{prevChoice =}\NormalTok{ wsls\_choices\_vs\_biased[t }\SpecialCharTok{{-}} \DecValTok{1}\NormalTok{],}
    \AttributeTok{feedback =}\NormalTok{ prev\_feedback,}
    \AttributeTok{noise =} \DecValTok{0} \CommentTok{\# Assuming no noise for this simulation}
\NormalTok{  )}
\NormalTok{\}}
\CommentTok{\# Record feedback for the last trial}
\NormalTok{feedback\_vs\_biased[trials] }\OtherTok{\textless{}{-}} \FunctionTok{ifelse}\NormalTok{(wsls\_choices\_vs\_biased[trials] }\SpecialCharTok{==}\NormalTok{ opponent\_choices\_biased[trials], }\DecValTok{1}\NormalTok{, }\DecValTok{0}\NormalTok{)}

\CommentTok{\# Create dataframe for analysis}
\NormalTok{df\_vs\_biased }\OtherTok{\textless{}{-}} \FunctionTok{tibble}\NormalTok{(}
  \AttributeTok{trial =} \DecValTok{1}\SpecialCharTok{:}\NormalTok{trials,}
  \AttributeTok{Self\_WSLS =}\NormalTok{ wsls\_choices\_vs\_biased,}
  \AttributeTok{Opponent\_Random =}\NormalTok{ opponent\_choices\_biased,}
  \AttributeTok{Feedback =}\NormalTok{ feedback\_vs\_biased }\CommentTok{\# 1 if WSLS won}
\NormalTok{) }\SpecialCharTok{\%\textgreater{}\%} \FunctionTok{mutate}\NormalTok{(}
  \AttributeTok{Cumulative\_Performance =} \FunctionTok{cumsum}\NormalTok{(Feedback) }\SpecialCharTok{/} \FunctionTok{row\_number}\NormalTok{()}
\NormalTok{)}

\CommentTok{\# {-}{-}{-} Simulation 2: WSLS vs. WSLS {-}{-}{-}}
\CommentTok{\# (Opponent plays WSLS but tries to mismatch {-} feedback is inverted)}
\FunctionTok{cat}\NormalTok{(}\StringTok{"Simulating WSLS vs. WSLS...}\SpecialCharTok{\textbackslash{}n}\StringTok{"}\NormalTok{)}
\end{Highlighting}
\end{Shaded}

\begin{verbatim}
Simulating WSLS vs. WSLS...
\end{verbatim}

\begin{Shaded}
\begin{Highlighting}[]
\NormalTok{wsls1\_choices }\OtherTok{\textless{}{-}} \FunctionTok{rep}\NormalTok{(}\ConstantTok{NA}\NormalTok{, trials)}
\NormalTok{wsls2\_choices }\OtherTok{\textless{}{-}} \FunctionTok{rep}\NormalTok{(}\ConstantTok{NA}\NormalTok{, trials)}
\NormalTok{feedback\_wsls1 }\OtherTok{\textless{}{-}} \FunctionTok{rep}\NormalTok{(}\ConstantTok{NA}\NormalTok{, trials) }\CommentTok{\# Feedback for agent 1}

\NormalTok{wsls1\_choices[}\DecValTok{1}\NormalTok{] }\OtherTok{\textless{}{-}} \FunctionTok{sample}\NormalTok{(}\FunctionTok{c}\NormalTok{(}\DecValTok{0}\NormalTok{, }\DecValTok{1}\NormalTok{), }\DecValTok{1}\NormalTok{)}
\NormalTok{wsls2\_choices[}\DecValTok{1}\NormalTok{] }\OtherTok{\textless{}{-}} \FunctionTok{sample}\NormalTok{(}\FunctionTok{c}\NormalTok{(}\DecValTok{0}\NormalTok{, }\DecValTok{1}\NormalTok{), }\DecValTok{1}\NormalTok{)}

\ControlFlowTok{for}\NormalTok{ (t }\ControlFlowTok{in} \DecValTok{2}\SpecialCharTok{:}\NormalTok{trials) \{}
  \CommentTok{\# Feedback for agent 1 from previous trial (wins if match)}
\NormalTok{  prev\_feedback1 }\OtherTok{\textless{}{-}} \FunctionTok{ifelse}\NormalTok{(wsls1\_choices[t }\SpecialCharTok{{-}} \DecValTok{1}\NormalTok{] }\SpecialCharTok{==}\NormalTok{ wsls2\_choices[t }\SpecialCharTok{{-}} \DecValTok{1}\NormalTok{], }\DecValTok{1}\NormalTok{, }\DecValTok{0}\NormalTok{)}
\NormalTok{  feedback\_wsls1[t }\SpecialCharTok{{-}} \DecValTok{1}\NormalTok{] }\OtherTok{\textless{}{-}}\NormalTok{ prev\_feedback1}

  \CommentTok{\# Feedback for agent 2 (opponent) is the opposite (wins if mismatch)}
\NormalTok{  prev\_feedback2 }\OtherTok{\textless{}{-}} \DecValTok{1} \SpecialCharTok{{-}}\NormalTok{ prev\_feedback1}

  \CommentTok{\# Both agents choose based on their own previous feedback}
\NormalTok{  wsls1\_choices[t] }\OtherTok{\textless{}{-}} \FunctionTok{WSLSAgent\_f}\NormalTok{(wsls1\_choices[t }\SpecialCharTok{{-}} \DecValTok{1}\NormalTok{], prev\_feedback1, }\AttributeTok{noise =} \DecValTok{0}\NormalTok{)}
\NormalTok{  wsls2\_choices[t] }\OtherTok{\textless{}{-}} \FunctionTok{WSLSAgent\_f}\NormalTok{(wsls2\_choices[t }\SpecialCharTok{{-}} \DecValTok{1}\NormalTok{], prev\_feedback2, }\AttributeTok{noise =} \DecValTok{0}\NormalTok{)}
\NormalTok{\}}
\NormalTok{feedback\_wsls1[trials] }\OtherTok{\textless{}{-}} \FunctionTok{ifelse}\NormalTok{(wsls1\_choices[trials] }\SpecialCharTok{==}\NormalTok{ wsls2\_choices[trials], }\DecValTok{1}\NormalTok{, }\DecValTok{0}\NormalTok{)}

\CommentTok{\# Create dataframe}
\NormalTok{df\_vs\_wsls }\OtherTok{\textless{}{-}} \FunctionTok{tibble}\NormalTok{(}
  \AttributeTok{trial =} \DecValTok{1}\SpecialCharTok{:}\NormalTok{trials,}
  \AttributeTok{Self\_WSLS =}\NormalTok{ wsls1\_choices,}
  \AttributeTok{Opponent\_WSLS =}\NormalTok{ wsls2\_choices,}
  \AttributeTok{Feedback =}\NormalTok{ feedback\_wsls1 }\CommentTok{\# 1 if Self\_WSLS won}
\NormalTok{) }\SpecialCharTok{\%\textgreater{}\%} \FunctionTok{mutate}\NormalTok{(}
  \AttributeTok{Cumulative\_Performance =} \FunctionTok{cumsum}\NormalTok{(Feedback) }\SpecialCharTok{/} \FunctionTok{row\_number}\NormalTok{()}
\NormalTok{)}

\CommentTok{\# {-}{-}{-} Visualize Simulation Results {-}{-}{-}}
\CommentTok{\# Plot choices over time}
\NormalTok{p\_choices\_biased }\OtherTok{\textless{}{-}} \FunctionTok{ggplot}\NormalTok{(df\_vs\_biased, }\FunctionTok{aes}\NormalTok{(}\AttributeTok{x =}\NormalTok{ trial)) }\SpecialCharTok{+}
  \FunctionTok{geom\_line}\NormalTok{(}\FunctionTok{aes}\NormalTok{(}\AttributeTok{y =}\NormalTok{ Self\_WSLS, }\AttributeTok{color =} \StringTok{"WSLS Agent"}\NormalTok{)) }\SpecialCharTok{+}
  \FunctionTok{geom\_line}\NormalTok{(}\FunctionTok{aes}\NormalTok{(}\AttributeTok{y =}\NormalTok{ Opponent\_Random }\SpecialCharTok{+} \FloatTok{0.05}\NormalTok{, }\AttributeTok{color =} \StringTok{"Biased Opponent"}\NormalTok{), }\AttributeTok{linetype =} \StringTok{"dashed"}\NormalTok{) }\SpecialCharTok{+} \CommentTok{\# Offset slightly}
  \FunctionTok{labs}\NormalTok{(}\AttributeTok{title =} \StringTok{"WSLS vs. Biased Random"}\NormalTok{, }\AttributeTok{y =} \StringTok{"Choice (0/1)"}\NormalTok{) }\SpecialCharTok{+} \FunctionTok{theme\_classic}\NormalTok{() }\SpecialCharTok{+} \FunctionTok{ylim}\NormalTok{(}\SpecialCharTok{{-}}\FloatTok{0.1}\NormalTok{, }\FloatTok{1.1}\NormalTok{)}

\NormalTok{p\_choices\_wsls }\OtherTok{\textless{}{-}} \FunctionTok{ggplot}\NormalTok{(df\_vs\_wsls, }\FunctionTok{aes}\NormalTok{(}\AttributeTok{x =}\NormalTok{ trial)) }\SpecialCharTok{+}
  \FunctionTok{geom\_line}\NormalTok{(}\FunctionTok{aes}\NormalTok{(}\AttributeTok{y =}\NormalTok{ Self\_WSLS, }\AttributeTok{color =} \StringTok{"WSLS Agent 1"}\NormalTok{)) }\SpecialCharTok{+}
  \FunctionTok{geom\_line}\NormalTok{(}\FunctionTok{aes}\NormalTok{(}\AttributeTok{y =}\NormalTok{ Opponent\_WSLS }\SpecialCharTok{+} \FloatTok{0.05}\NormalTok{, }\AttributeTok{color =} \StringTok{"WSLS Agent 2"}\NormalTok{), }\AttributeTok{linetype =} \StringTok{"dashed"}\NormalTok{) }\SpecialCharTok{+} \CommentTok{\# Offset slightly}
  \FunctionTok{labs}\NormalTok{(}\AttributeTok{title =} \StringTok{"WSLS vs. WSLS"}\NormalTok{, }\AttributeTok{y =} \StringTok{"Choice (0/1)"}\NormalTok{) }\SpecialCharTok{+} \FunctionTok{theme\_classic}\NormalTok{() }\SpecialCharTok{+} \FunctionTok{ylim}\NormalTok{(}\SpecialCharTok{{-}}\FloatTok{0.1}\NormalTok{, }\FloatTok{1.1}\NormalTok{)}

\CommentTok{\# Plot cumulative performance}
\NormalTok{p\_perf\_biased }\OtherTok{\textless{}{-}} \FunctionTok{ggplot}\NormalTok{(df\_vs\_biased, }\FunctionTok{aes}\NormalTok{(}\AttributeTok{x =}\NormalTok{ trial, }\AttributeTok{y =}\NormalTok{ Cumulative\_Performance)) }\SpecialCharTok{+}
  \FunctionTok{geom\_line}\NormalTok{(}\AttributeTok{color =} \StringTok{"blue"}\NormalTok{, }\AttributeTok{size =} \DecValTok{1}\NormalTok{) }\SpecialCharTok{+} \FunctionTok{geom\_hline}\NormalTok{(}\AttributeTok{yintercept =} \FloatTok{0.5}\NormalTok{, }\AttributeTok{linetype =} \StringTok{"dashed"}\NormalTok{) }\SpecialCharTok{+}
  \FunctionTok{labs}\NormalTok{(}\AttributeTok{title =} \StringTok{"WSLS Performance vs. Biased Random"}\NormalTok{, }\AttributeTok{y =} \StringTok{"Proportion Wins"}\NormalTok{) }\SpecialCharTok{+} \FunctionTok{theme\_classic}\NormalTok{() }\SpecialCharTok{+} \FunctionTok{ylim}\NormalTok{(}\DecValTok{0}\NormalTok{, }\DecValTok{1}\NormalTok{)}

\NormalTok{p\_perf\_wsls }\OtherTok{\textless{}{-}} \FunctionTok{ggplot}\NormalTok{(df\_vs\_wsls, }\FunctionTok{aes}\NormalTok{(}\AttributeTok{x =}\NormalTok{ trial, }\AttributeTok{y =}\NormalTok{ Cumulative\_Performance)) }\SpecialCharTok{+}
  \FunctionTok{geom\_line}\NormalTok{(}\AttributeTok{color =} \StringTok{"purple"}\NormalTok{, }\AttributeTok{size =} \DecValTok{1}\NormalTok{) }\SpecialCharTok{+} \FunctionTok{geom\_hline}\NormalTok{(}\AttributeTok{yintercept =} \FloatTok{0.5}\NormalTok{, }\AttributeTok{linetype =} \StringTok{"dashed"}\NormalTok{) }\SpecialCharTok{+}
  \FunctionTok{labs}\NormalTok{(}\AttributeTok{title =} \StringTok{"WSLS Performance vs. WSLS"}\NormalTok{, }\AttributeTok{y =} \StringTok{"Proportion Wins"}\NormalTok{) }\SpecialCharTok{+} \FunctionTok{theme\_classic}\NormalTok{() }\SpecialCharTok{+} \FunctionTok{ylim}\NormalTok{(}\DecValTok{0}\NormalTok{, }\DecValTok{1}\NormalTok{)}

\CommentTok{\# Arrange plots}
\NormalTok{(p\_choices\_biased }\SpecialCharTok{|}\NormalTok{ p\_choices\_wsls) }\SpecialCharTok{/}\NormalTok{ (p\_perf\_biased }\SpecialCharTok{|}\NormalTok{ p\_perf\_wsls) }\SpecialCharTok{+}
  \FunctionTok{plot\_layout}\NormalTok{(}\AttributeTok{guides =} \StringTok{"collect"}\NormalTok{) }\SpecialCharTok{\&} \FunctionTok{theme}\NormalTok{(}\AttributeTok{legend.position =} \StringTok{"bottom"}\NormalTok{)}
\end{Highlighting}
\end{Shaded}

\begin{center}
\includegraphics[width=0.8\linewidth,height=\textheight,keepaspectratio]{03-BuildingFormalModels_files/figure-pdf/ch04_wsls_vs_biased-1.pdf}
\end{center}

This cumulative performance plot reveals the overall effectiveness of
the WSLS strategy. By tracking the running average of successes, we can
see whether the strategy leads to above-chance performance in the long
run. When playing against a biased random opponent, the WSLS agent can
potentially exploit the opponent's predictable tendencies, though
success depends on how strong and consistent the opponent's bias is.
When we pit the WSLS agent against another WSLS agent, the dynamics
become more complex. Both agents are now trying to adapt to each other's
adaptations, creating a more sophisticated strategic interaction. The
resulting behavior often shows interesting patterns of mutual
adaptation, where each agent's attempts to exploit the other's strategy
leads to evolving patterns of play.

\section{Scaling Up: The Virtual
Experiment}\label{scaling-up-the-virtual-experiment}

Now that we have our agents defined as functions, we can run
\textbf{virtual Experiments}, to better understand our models.

We will build a ``Simulation Engine'' function that:

\begin{enumerate}
\def\labelenumi{\arabic{enumi}.}
\item
  Takes two agents (e.g., WSLS vs.~Random).
\item
  Makes them play for \texttt{N} trials.
\item
  Records every move.
\end{enumerate}

To understand the average behavior and robustness of the WSLS strategy,
we need to simulate many agents playing against opponents with varying
biases.

\begin{quote}
\emph{Why run a virtual experiment instead of just testing one agent?} A
single simulation is like flipping a coin once --- you cannot tell if
the coin is fair from one flip. Running 100 simulated agents lets us
separate the signal (the model's actual tendency) from the noise (random
variation in any one run). It also lets us sweep across parameter values
(e.g., varying opponent bias from 0.5 to 1.0) to understand \emph{when}
WSLS works and when it fails --- building intuition before we ever see
real human data.
\end{quote}

First, let's create a function that runs one simulation (one WSLS agent
vs.~one opponent) and returns the results. This avoids repeating code.

\begin{Shaded}
\begin{Highlighting}[]
\CommentTok{\#\textquotesingle{} Run Matching Pennies Simulation}
\CommentTok{\#\textquotesingle{}}
\CommentTok{\#\textquotesingle{} Simulates a full game between an Agent (Subject) and an Opponent.}
\CommentTok{\#\textquotesingle{}}
\CommentTok{\#\textquotesingle{} @param n\_trials Integer, number of trials.}
\CommentTok{\#\textquotesingle{} @param agent\_strategy String, "WSLS" or "Random".}
\CommentTok{\#\textquotesingle{} @param agent\_params List, parameters for the agent (e.g., list(noise=0.1)).}
\CommentTok{\#\textquotesingle{} @param opponent\_strategy String, "WSLS" or "Random".}
\CommentTok{\#\textquotesingle{} @param opponent\_params List, parameters for the opponent (e.g., list(rate=0.8)).}
\CommentTok{\#\textquotesingle{}}
\CommentTok{\#\textquotesingle{} @return A tibble with trial{-}by{-}trial choices and outcomes.}
\CommentTok{\#\textquotesingle{}}
\NormalTok{run\_simulation }\OtherTok{\textless{}{-}} \ControlFlowTok{function}\NormalTok{(n\_trials, agent\_strategy, agent\_params, opponent\_strategy, opponent\_params) \{}

  \CommentTok{\# Pre{-}allocate vectors for speed}
\NormalTok{  agent\_choices }\OtherTok{\textless{}{-}} \FunctionTok{rep}\NormalTok{(}\ConstantTok{NA}\NormalTok{, n\_trials)}
\NormalTok{  opponent\_choices }\OtherTok{\textless{}{-}} \FunctionTok{rep}\NormalTok{(}\ConstantTok{NA}\NormalTok{, n\_trials)}
\NormalTok{  agent\_feedback }\OtherTok{\textless{}{-}} \FunctionTok{rep}\NormalTok{(}\ConstantTok{NA}\NormalTok{, n\_trials) }

  \CommentTok{\# {-}{-}{-} Trial 1 (Initialization) {-}{-}{-}}
  \CommentTok{\# Agents behave randomly or use a specific start bias on trial 1}
\NormalTok{  agent\_choices[}\DecValTok{1}\NormalTok{] }\OtherTok{\textless{}{-}} \FunctionTok{rbinom}\NormalTok{(}\DecValTok{1}\NormalTok{, }\DecValTok{1}\NormalTok{, }\FloatTok{0.5}\NormalTok{)}
\NormalTok{  opponent\_choices[}\DecValTok{1}\NormalTok{] }\OtherTok{\textless{}{-}} \FunctionTok{rbinom}\NormalTok{(}\DecValTok{1}\NormalTok{, }\DecValTok{1}\NormalTok{, }\FloatTok{0.5}\NormalTok{)}
  
  \CommentTok{\# Check winner (Assuming Agent is Matcher: Win if choices match)}
\NormalTok{  agent\_feedback[}\DecValTok{1}\NormalTok{] }\OtherTok{\textless{}{-}} \FunctionTok{ifelse}\NormalTok{(agent\_choices[}\DecValTok{1}\NormalTok{] }\SpecialCharTok{==}\NormalTok{ opponent\_choices[}\DecValTok{1}\NormalTok{], }\DecValTok{1}\NormalTok{, }\DecValTok{0}\NormalTok{)}

  \CommentTok{\# {-}{-}{-} Loop through Game {-}{-}{-}}
  \ControlFlowTok{for}\NormalTok{ (t }\ControlFlowTok{in} \DecValTok{2}\SpecialCharTok{:}\NormalTok{n\_trials) \{}
    
    \CommentTok{\# 1. Agent makes a choice based on history (t{-}1)}
    \ControlFlowTok{if}\NormalTok{ (agent\_strategy }\SpecialCharTok{==} \StringTok{"WSLS"}\NormalTok{) \{}
\NormalTok{      agent\_choices[t] }\OtherTok{\textless{}{-}} \FunctionTok{WSLSAgent\_f}\NormalTok{(agent\_choices[t}\DecValTok{{-}1}\NormalTok{], agent\_feedback[t}\DecValTok{{-}1}\NormalTok{], agent\_params}\SpecialCharTok{$}\NormalTok{noise)}
\NormalTok{    \} }\ControlFlowTok{else} \ControlFlowTok{if}\NormalTok{ (agent\_strategy }\SpecialCharTok{==} \StringTok{"Random"}\NormalTok{) \{}
\NormalTok{      agent\_choices[t] }\OtherTok{\textless{}{-}} \FunctionTok{RandomAgent\_f}\NormalTok{(}\DecValTok{1}\NormalTok{, agent\_params}\SpecialCharTok{$}\NormalTok{rate, agent\_params}\SpecialCharTok{$}\NormalTok{noise)}
\NormalTok{    \}}
    
    \CommentTok{\# 2. Opponent makes a choice}
    \ControlFlowTok{if}\NormalTok{ (opponent\_strategy }\SpecialCharTok{==} \StringTok{"Random"}\NormalTok{) \{}
\NormalTok{      opponent\_choices[t] }\OtherTok{\textless{}{-}} \FunctionTok{RandomAgent\_f}\NormalTok{(}\DecValTok{1}\NormalTok{, opponent\_params}\SpecialCharTok{$}\NormalTok{rate, opponent\_params}\SpecialCharTok{$}\NormalTok{noise)}
\NormalTok{    \} }\ControlFlowTok{else} \ControlFlowTok{if}\NormalTok{ (opponent\_strategy }\SpecialCharTok{==} \StringTok{"WSLS"}\NormalTok{) \{}
      \CommentTok{\# Opponent is Hider: They win if choices MISMATCH. }
      \CommentTok{\# So their feedback is the opposite of the Agent\textquotesingle{}s feedback.}
\NormalTok{      op\_feedback }\OtherTok{\textless{}{-}} \DecValTok{1} \SpecialCharTok{{-}}\NormalTok{ agent\_feedback[t}\DecValTok{{-}1}\NormalTok{]}
\NormalTok{      opponent\_choices[t] }\OtherTok{\textless{}{-}} \FunctionTok{WSLSAgent\_f}\NormalTok{(opponent\_choices[t}\DecValTok{{-}1}\NormalTok{], op\_feedback, opponent\_params}\SpecialCharTok{$}\NormalTok{noise)}
\NormalTok{    \}}

    \CommentTok{\# 3. Determine Outcome}
\NormalTok{    agent\_feedback[t] }\OtherTok{\textless{}{-}} \FunctionTok{ifelse}\NormalTok{(agent\_choices[t] }\SpecialCharTok{==}\NormalTok{ opponent\_choices[t], }\DecValTok{1}\NormalTok{, }\DecValTok{0}\NormalTok{)}
\NormalTok{  \}}

  \CommentTok{\# Return data}
  \FunctionTok{return}\NormalTok{(}\FunctionTok{tibble}\NormalTok{(}
    \AttributeTok{trial =} \DecValTok{1}\SpecialCharTok{:}\NormalTok{n\_trials,}
    \AttributeTok{choice =}\NormalTok{ agent\_choices,}
    \AttributeTok{opponent =}\NormalTok{ opponent\_choices,}
    \AttributeTok{win =}\NormalTok{ agent\_feedback}
\NormalTok{  ))}
\NormalTok{\}}

\CommentTok{\# {-}{-}{-} Running Scaled Simulations {-}{-}{-}}
\NormalTok{sim\_data\_file }\OtherTok{\textless{}{-}} \FunctionTok{here}\NormalTok{(}\StringTok{"simdata"}\NormalTok{, }\StringTok{"ch03\_scaled\_simulations.rds"}\NormalTok{)}

\ControlFlowTok{if}\NormalTok{ (regenerate\_simulations }\SpecialCharTok{||} \SpecialCharTok{!}\FunctionTok{file.exists}\NormalTok{(sim\_data\_file)) \{}
  \FunctionTok{cat}\NormalTok{(}\StringTok{"Running scaled simulations...}\SpecialCharTok{\textbackslash{}n}\StringTok{"}\NormalTok{)}

  \CommentTok{\# Default WSLS parameters (Only \textquotesingle{}noise\textquotesingle{} is used by our simple WSLS agent)}
\NormalTok{  wsls\_params\_default }\OtherTok{\textless{}{-}} \FunctionTok{list}\NormalTok{(}\AttributeTok{noise =} \DecValTok{0}\NormalTok{)}

  \CommentTok{\# Simulation 1: WSLS vs. Random agents with varying biases}
\NormalTok{  opponent\_rates }\OtherTok{\textless{}{-}} \FunctionTok{seq}\NormalTok{(}\AttributeTok{from =} \FloatTok{0.5}\NormalTok{, }\AttributeTok{to =} \DecValTok{1}\NormalTok{, }\AttributeTok{by =} \FloatTok{0.05}\NormalTok{)}
\NormalTok{  results\_vs\_random }\OtherTok{\textless{}{-}} \FunctionTok{list}\NormalTok{() }

  \ControlFlowTok{for}\NormalTok{ (rate }\ControlFlowTok{in}\NormalTok{ opponent\_rates) \{}
    \CommentTok{\# Create a list to store runs for this specific rate}
\NormalTok{    rate\_runs }\OtherTok{\textless{}{-}} \FunctionTok{list}\NormalTok{()}
    
    \ControlFlowTok{for}\NormalTok{ (agent\_i }\ControlFlowTok{in} \DecValTok{1}\SpecialCharTok{:}\NormalTok{agents) \{}
\NormalTok{       opponent\_params\_rand }\OtherTok{\textless{}{-}} \FunctionTok{list}\NormalTok{(}\AttributeTok{rate =}\NormalTok{ rate, }\AttributeTok{noise =} \DecValTok{0}\NormalTok{)}
       
       \CommentTok{\# CORRECTED CALL: Added "WSLS" as the second argument}
\NormalTok{       sim\_run }\OtherTok{\textless{}{-}} \FunctionTok{run\_simulation}\NormalTok{(trials, }\StringTok{"WSLS"}\NormalTok{, wsls\_params\_default, }\StringTok{"Random"}\NormalTok{, opponent\_params\_rand)}
       
\NormalTok{       sim\_run}\SpecialCharTok{$}\NormalTok{agent }\OtherTok{\textless{}{-}}\NormalTok{ agent\_i}
\NormalTok{       sim\_run}\SpecialCharTok{$}\NormalTok{opponent\_rate }\OtherTok{\textless{}{-}}\NormalTok{ rate}
\NormalTok{       rate\_runs[[agent\_i]] }\OtherTok{\textless{}{-}}\NormalTok{ sim\_run}
\NormalTok{    \}}
    \CommentTok{\# Combine all agents for this rate and store in main list}
\NormalTok{    results\_vs\_random[[}\FunctionTok{as.character}\NormalTok{(rate)]] }\OtherTok{\textless{}{-}} \FunctionTok{bind\_rows}\NormalTok{(rate\_runs)}
\NormalTok{  \}}
  \CommentTok{\# Combine everything into one big dataframe}
\NormalTok{  df\_scaled\_vs\_random }\OtherTok{\textless{}{-}} \FunctionTok{bind\_rows}\NormalTok{(results\_vs\_random)}
  \FunctionTok{saveRDS}\NormalTok{(df\_scaled\_vs\_random, sim\_data\_file)}
\NormalTok{\} }\ControlFlowTok{else}\NormalTok{ \{}
\NormalTok{  df\_scaled\_vs\_random }\OtherTok{\textless{}{-}} \FunctionTok{readRDS}\NormalTok{(sim\_data\_file)}
\NormalTok{\}}
\end{Highlighting}
\end{Shaded}

\subsection{Visualizing Results}\label{visualizing-results}

Now let's visualize the average performance of the WSLS agent against
random opponents with different biases.

\begin{Shaded}
\begin{Highlighting}[]
\CommentTok{\# Calculate cumulative performance for each simulation run}
\NormalTok{df\_scaled\_vs\_random }\OtherTok{\textless{}{-}}\NormalTok{ df\_scaled\_vs\_random }\SpecialCharTok{\%\textgreater{}\%}
  \FunctionTok{group\_by}\NormalTok{(agent, opponent\_rate) }\SpecialCharTok{\%\textgreater{}\%}
  \CommentTok{\# CHANGED \textquotesingle{}Feedback\textquotesingle{} to \textquotesingle{}win\textquotesingle{} here:}
  \FunctionTok{mutate}\NormalTok{(}\AttributeTok{Cumulative\_Performance =} \FunctionTok{cumsum}\NormalTok{(win) }\SpecialCharTok{/} \FunctionTok{row\_number}\NormalTok{()) }\SpecialCharTok{\%\textgreater{}\%}
  \FunctionTok{ungroup}\NormalTok{()}

\CommentTok{\# Plot average performance over time, grouped by opponent bias rate}
\NormalTok{p\_scaled\_perf }\OtherTok{\textless{}{-}} \FunctionTok{ggplot}\NormalTok{(df\_scaled\_vs\_random,}
       \FunctionTok{aes}\NormalTok{(}\AttributeTok{x =}\NormalTok{ trial, }\AttributeTok{y =}\NormalTok{ Cumulative\_Performance,}
           \AttributeTok{group =}\NormalTok{ opponent\_rate, }\AttributeTok{color =} \FunctionTok{as.factor}\NormalTok{(opponent\_rate))) }\SpecialCharTok{+} \CommentTok{\# made color a factor for cleaner legend}
  \FunctionTok{stat\_summary}\NormalTok{(}\AttributeTok{fun =}\NormalTok{ mean, }\AttributeTok{geom =} \StringTok{"line"}\NormalTok{, }\AttributeTok{size =} \DecValTok{1}\NormalTok{) }\SpecialCharTok{+}
  \FunctionTok{geom\_hline}\NormalTok{(}\AttributeTok{yintercept =} \FloatTok{0.5}\NormalTok{, }\AttributeTok{linetype =} \StringTok{"dashed"}\NormalTok{) }\SpecialCharTok{+}
  \FunctionTok{scale\_color\_viridis\_d}\NormalTok{() }\SpecialCharTok{+} \CommentTok{\# Discrete scale since we made rate a factor}
  \FunctionTok{labs}\NormalTok{(}
    \AttributeTok{title =} \StringTok{"Average WSLS Performance vs. Random Opponents"}\NormalTok{,}
    \AttributeTok{subtitle =} \StringTok{"Performance improves as opponent\textquotesingle{}s bias increases"}\NormalTok{,}
    \AttributeTok{x =} \StringTok{"Trial Number"}\NormalTok{,}
    \AttributeTok{y =} \StringTok{"Average Proportion Wins"}\NormalTok{,}
    \AttributeTok{color =} \StringTok{"Opponent Bias"}
\NormalTok{  ) }\SpecialCharTok{+}
  \FunctionTok{ylim}\NormalTok{(}\DecValTok{0}\NormalTok{, }\DecValTok{1}\NormalTok{) }\SpecialCharTok{+}
  \FunctionTok{theme\_classic}\NormalTok{()}

\FunctionTok{print}\NormalTok{(p\_scaled\_perf)}
\end{Highlighting}
\end{Shaded}

\begin{center}
\includegraphics[width=0.8\linewidth,height=\textheight,keepaspectratio]{03-BuildingFormalModels_files/figure-pdf/ch04_visualizing_results-1.pdf}
\end{center}

This plot clearly shows that the WSLS strategy becomes more effective
(higher average win rate) as the opponent becomes more predictable
(i.e., their bias rate moves further away from 0.5 towards 1.0). Against
a purely random opponent (rate=0.5), WSLS performs at chance level on
average. This highlights WSLS's ability to exploit predictability.

\section{Conclusion: The Forward
Model}\label{conclusion-the-forward-model}

In this chapter, we successfully built a \textbf{Forward Model}. * We
started with a verbal theory (WSLS). * We translated it into code
(parameters + algorithm). * We ran the code to \textbf{generate}
artificial data.

This is the ``models of the Mind'' in action. We demonstrated that a
simple mechanism (Win-Stay-Lose-Shift) can generate complex success
rates depending on the environment (the opponent).

\subsection{The Problem of Inference}\label{the-problem-of-inference}

However, in real science, we don't start with the parameters. We start
with the data (the choices participants made) and we want to find the
parameters (did they use WSLS? How much noise?).

We need to run our machine in reverse.

\begin{itemize}
\item
  \textbf{Forward (Simulation)}: Parameters \(\rightarrow\) Data
\item
  \textbf{Inverse (Fitting)}: Data \(\rightarrow\) Parameters
\end{itemize}

\textbf{Coming up in Chapter 5:} We reverse the arrow. The simulated
data we just created --- sequences of choices generated by agents with
\emph{known} parameters --- become our test cases. We will learn to
infer those parameters back from the choices using Bayesian inference in
Stan. This single-dataset fit is the foundation; in Chapter 6, we will
then stress-test the entire inference procedure (prior predictive
checks, parameter recovery, SBC, posterior predictive checks) before the
model is allowed anywhere near real human data.

\begin{tcolorbox}[enhanced jigsaw, toprule=.15mm, toptitle=1mm, bottomtitle=1mm, leftrule=.75mm, rightrule=.15mm, colframe=quarto-callout-note-color-frame, left=2mm, colback=white, titlerule=0mm, coltitle=black, colbacktitle=quarto-callout-note-color!10!white, opacityback=0, title=\textcolor{quarto-callout-note-color}{\faInfo}\hspace{0.5em}{Open Science and the Case for Transparent Implementations}, arc=.35mm, opacitybacktitle=0.6, breakable, bottomrule=.15mm]

You have just written your first formal cognitive model: a function that
takes parameters, runs an algorithm, and produces data. It is tempting
to treat this as a private computational tool --- something you run,
extract a result from, and move on. The history of cognitive modeling
suggests this would be a mistake.

A persistent challenge in the field is what we might call the
\textbf{hidden implementation problem}. Two research groups both claim
to implement the same model --- say, a Win-Stay-Lose-Shift agent, or
later in the course, an exemplar model or a reinforcement learner. Their
verbal descriptions match. Their equations look identical on paper. But
their code differs in subtle ways: how they handle the first trial,
whether they normalise probabilities before or after applying a noise
parameter, how they deal with edge cases. The result is that ``the same
model'' produces different predictions, and conflicting empirical
conclusions are blamed on the data rather than the implementation. The
field accumulates apparent contradictions that are really implementation
artifacts.

The solution is not to standardise everyone onto a single package ---
though shared, collaboratively verified implementations are genuinely
valuable. A pioneering example is the \texttt{catlearn} R package (Wills
et al. 2017), which provides formally verified implementations of
classic categorization models alongside canonical benchmark datasets, so
that results from different labs are directly comparable. The broader
principle extends to every domain of cognitive modeling: when
implementations are open and documented, errors can be found and fixed;
when they are private, the same errors propagate silently.

These course notes are designed as a contribution to this open-science
ethos. Rather than handing you pre-built scripts to run, we ask you to
implement each model yourself, document every assumption explicitly,
validate the implementation against known inputs and outputs, and
produce code you could share. You will not struggle with an equation in
silence and then quietly shelve the script because you are not sure it
is correct. You will build the habits --- transparent code, documented
choices, validated outputs --- that make science cumulative rather than
siloed. Every chapter that follows applies this principle: forward
simulator, Stan model, diagnostic battery, and validation checks, all in
open, reproducible code.

\end{tcolorbox}

\bookmarksetup{startatroot}

\chapter{From simulation to model fitting}\label{sec-inferring-rates}

\begin{quote}
\textbf{📍 Where we are in the Bayesian modeling workflow:} Ch. 3 built
forward models --- code that generates data \emph{from} known
parameters. This chapter reverses the arrow: given observed choices, we
infer the parameters that most plausibly produced them. This is
\textbf{Bayesian parameter estimation}, and it is the engine at the
heart of all remaining chapters. We fit to \emph{simulated} data here
precisely because we know the ground truth --- a sanity check before the
engine is handed real human data in Ch. 6.
\end{quote}

In Chapter 3, we formalized verbal models into computational agents and
explored their behavior through simulation by \emph{setting} their
parameters (like bias \texttt{rate} or WSLS rules). However, when
analyzing real experimental data, we don't know the true parameter
values underlying a participant's behavior. The core task shifts from
simulation to \textbf{parameter estimation}: inferring the plausible
values of model parameters \emph{given} the observed data.

\section{Learning Objectives}\label{learning-objectives-3}

This chapter transitions from simulating predefined models (Chapter 3)
to inferring model parameters from data using Bayesian methods. By the
end of this chapter, you will be able to:

\begin{itemize}
\tightlist
\item
  \textbf{Model Cognitive Processes:} Apply these techniques to estimate
  parameters for simple cognitive strategies, including biased choice
  and different types of memory models.
\item
  \textbf{Master Reparameterization:} Learn to rethink the mathematical
  form of models. You will appreciate the utility of the log-odds
  transformation for mapping bounded probability parameters into
  unbounded spaces.
\item
  \textbf{Read MCMC Diagnostics:} Learn the mandatory diagnostic battery
  for Bayesian inference, specifically focusing on chain convergence
  (\(\hat{R}\)) and Effective Sample Size (ESS).
\end{itemize}

In a rigorous Bayesian workflow, we \emph{never} fit a model to
empirical data without validating it first. Therefore, throughout this
chapter, we will fit our models to a single \textbf{simulated dataset}
where we know the ground truth. This allows us to learn the mechanics of
Stan and verify our single-fit posterior. In the next chapter, we will
scale this up into a full rigorous validation pipeline (Parameter
Recovery and predictive checks).

\section{The Challenge: Inferring Latent
Parameters}\label{the-challenge-inferring-latent-parameters}

Why is this a challenge?

\begin{enumerate}
\def\labelenumi{\arabic{enumi}.}
\item
  \textbf{Parameters are Unseen:} Cognitive parameters (like bias
  strength, memory decay, learning rate) are not directly observable; we
  must infer them from behavior (choices, RTs, etc.).
\item
  \textbf{Behavior is Noisy:} Human behavior is variable. Our models
  need to account for this noise while extracting the underlying
  parameter signal.
\item
  \textbf{Model Plausibility:} We need to determine not just \emph{if} a
  model fits, but \emph{how well} it fits, what are its strengths and
  limitations, and whether the estimated parameter values are
  theoretically meaningful.
\end{enumerate}

Consider the biased agent from Chapter 3. We simulated agents with a
known \texttt{rate}. Now, imagine you only have the sequence of choices
(the \texttt{h} data) from an unknown agent. How can you estimate the
underlying \texttt{rate} that most likely generated those choices?

\textbf{Our Approach: Bayesian Inference with Stan}

This chapter introduces \textbf{Bayesian inference} as our primary tool
for parameter estimation. The Bayesian approach allows us to combine
prior knowledge about parameters with information from the observed data
(via the likelihood function) to arrive at an updated understanding
(posterior distribution) of the parameters. We will use \textbf{Stan}, a
powerful probabilistic programming language, implemented via the
\texttt{cmdstanr} R package, to:

\begin{enumerate}
\def\labelenumi{\arabic{enumi}.}
\tightlist
\item
  \textbf{Specify Models:} Define our cognitive models formally,
  including parameters and their prior distributions.
\item
  \textbf{Fit Models:} Use Stan's algorithms (primarily Hamiltonian
  Monte Carlo - HMC, and a faster but more risky approximation method
  called pathfinder) to sample from the posterior distribution of
  parameters given the data.
\item
  \textbf{Evaluate Fits:} Examine the results to understand parameter
  estimates and their uncertainty.
\end{enumerate}

As a crucial validation step, we will first apply these methods to
\textbf{simulated data}, where we know the true parameters. This allows
us to check if our models and fitting procedures work correctly -- a
process called \textbf{parameter recovery}.

\begin{quote}
\textbf{Terminology note --- ``parameter recovery'' used carefully:}
Here and in the conclusion of this chapter, we use ``parameter
recovery'' loosely to mean: \emph{fitting a model to one simulated
dataset and checking whether the posterior lands near the true value.}
This is a useful first sanity check. In Ch. 6 we will use the term more
precisely to mean a \emph{systematic} study that repeats this across
dozens of datasets at many different true parameter values and noise
levels. That is a qualitatively different (and stronger) validation. The
even stricter test --- Simulation-Based Calibration Checks (SBC) ---
checks whether posterior \emph{intervals} are correctly calibrated
across the entire prior space. These distinctions will matter a lot in
Ch. 6, so don't worry about them too much now, you'll hear more about
them.
\end{quote}

\section{Simulating data}\label{simulating-data}

As usual we start with simulated data, where we know the underlying
mechanisms and parameter values.

Simulated data are rarely enough (empirical data often offer unexpected
challenges), but they are a great starting point to stress test your
model: does the model reconstruct the right parameter values? Does it
reproduce the overall patterns in the data?

Here we build a new simulation of random agents with bias and noise. The
code and visualization is really nothing different from last chapter.

\begin{Shaded}
\begin{Highlighting}[]
\CommentTok{\# Flag to control whether to regenerate simulation/fitting results for the knitted files}
\CommentTok{\# Set to TRUE to rerun everything (takes time!), FALSE to load saved results.}
\NormalTok{regenerate\_simulations }\OtherTok{\textless{}{-}} \ConstantTok{FALSE} 

\CommentTok{\# {-}{-}{-} Simulate Data for Fitting {-}{-}{-}}
\CommentTok{\# Goal: Generate data from known models/parameters to test our fitting procedures.}

\NormalTok{trials }\OtherTok{\textless{}{-}} \DecValTok{120} \CommentTok{\# Number of trials per simulated agent}

\CommentTok{\# Define the agent function (copied from previosu chapter)}
\CommentTok{\# This agent chooses \textquotesingle{}1\textquotesingle{} with probability \textquotesingle{}rate\textquotesingle{}, potentially adding noise.}
\NormalTok{RandomAgentNoise\_f }\OtherTok{\textless{}{-}} \ControlFlowTok{function}\NormalTok{(rate, noise) \{}
  \CommentTok{\# Ensure rate is probability (in case log{-}odds are passed, though not here)}
\NormalTok{  rate\_prob }\OtherTok{\textless{}{-}} \FunctionTok{plogis}\NormalTok{(rate) }\CommentTok{\# Assumes rate might be log{-}odds; use identity if rate is already probability}
\NormalTok{  choice }\OtherTok{\textless{}{-}} \FunctionTok{rbinom}\NormalTok{(}\DecValTok{1}\NormalTok{, }\DecValTok{1}\NormalTok{, rate\_prob) }\CommentTok{\# Base choice}
  \ControlFlowTok{if}\NormalTok{ (noise }\SpecialCharTok{\textgreater{}} \DecValTok{0} \SpecialCharTok{\&} \FunctionTok{rbinom}\NormalTok{(}\DecValTok{1}\NormalTok{, }\DecValTok{1}\NormalTok{, noise) }\SpecialCharTok{==} \DecValTok{1}\NormalTok{) \{ }\CommentTok{\# Check if noise occurs}
\NormalTok{    choice }\OtherTok{\textless{}{-}} \FunctionTok{rbinom}\NormalTok{(}\DecValTok{1}\NormalTok{, }\DecValTok{1}\NormalTok{, }\FloatTok{0.5}\NormalTok{) }\CommentTok{\# Override with random 50/50 choice}
\NormalTok{  \}}
  \FunctionTok{return}\NormalTok{(choice)}
\NormalTok{\}}

\CommentTok{\# {-}{-}{-} Generate Data Across Different Conditions {-}{-}{-}}
\CommentTok{\# Define relative paths for saving/loading simulation data}
\NormalTok{sim\_data\_dir }\OtherTok{\textless{}{-}} \StringTok{"simdata"} \CommentTok{\# Assumes a \textquotesingle{}simdata\textquotesingle{} subdirectory}
\NormalTok{sim\_data\_file }\OtherTok{\textless{}{-}} \FunctionTok{here}\NormalTok{(sim\_data\_dir, }\StringTok{"ch06\_randomnoise.csv"}\NormalTok{) }

\CommentTok{\# Create directory if it doesn\textquotesingle{}t exist}
\ControlFlowTok{if}\NormalTok{ (}\SpecialCharTok{!}\FunctionTok{dir.exists}\NormalTok{(sim\_data\_dir)) }\FunctionTok{dir.create}\NormalTok{(sim\_data\_dir)}

\ControlFlowTok{if}\NormalTok{ (regenerate\_simulations }\SpecialCharTok{||} \SpecialCharTok{!}\FunctionTok{file.exists}\NormalTok{(sim\_data\_file)) \{}
  \FunctionTok{cat}\NormalTok{(}\StringTok{"Generating new simulation data...}\SpecialCharTok{\textbackslash{}n}\StringTok{"}\NormalTok{)}
  \CommentTok{\# Use expand\_grid for cleaner parameter combinations}
\NormalTok{  param\_grid }\OtherTok{\textless{}{-}} \FunctionTok{expand\_grid}\NormalTok{(}
    \AttributeTok{noise =} \FunctionTok{seq}\NormalTok{(}\DecValTok{0}\NormalTok{, }\FloatTok{0.5}\NormalTok{, }\FloatTok{0.1}\NormalTok{), }\CommentTok{\# Noise levels}
    \AttributeTok{rate =} \FunctionTok{seq}\NormalTok{(}\DecValTok{0}\NormalTok{, }\DecValTok{1}\NormalTok{, }\FloatTok{0.1}\NormalTok{)     }\CommentTok{\# Bias rate levels (probability scale)}
\NormalTok{  )}

  \CommentTok{\# Use map\_dfr for efficient simulation across parameters}
\NormalTok{  d }\OtherTok{\textless{}{-}} \FunctionTok{pmap\_dfr}\NormalTok{(param\_grid, }\ControlFlowTok{function}\NormalTok{(noise, rate) \{}
    \CommentTok{\# Simulate choices for one agent/condition}
\NormalTok{    choices }\OtherTok{\textless{}{-}} \FunctionTok{map\_int}\NormalTok{(}\DecValTok{1}\SpecialCharTok{:}\NormalTok{trials, }\SpecialCharTok{\textasciitilde{}}\FunctionTok{RandomAgentNoise\_f}\NormalTok{(}\FunctionTok{qlogis}\NormalTok{(rate), noise)) }\CommentTok{\# Convert rate to log{-}odds for function if needed}
    \CommentTok{\# Create tibble for this condition}
    \FunctionTok{tibble}\NormalTok{(}
      \AttributeTok{trial =} \DecValTok{1}\SpecialCharTok{:}\NormalTok{trials,}
      \AttributeTok{choice =}\NormalTok{ choices,}
      \AttributeTok{true\_rate =}\NormalTok{ rate, }\CommentTok{\# Store the true rate used for generation}
      \AttributeTok{true\_noise =}\NormalTok{ noise }\CommentTok{\# Store the true noise used}
\NormalTok{    )}
\NormalTok{  \}) }\SpecialCharTok{\%\textgreater{}\%}
  \CommentTok{\# Calculate cumulative rate for visualization}
  \FunctionTok{group\_by}\NormalTok{(true\_rate, true\_noise) }\SpecialCharTok{\%\textgreater{}\%}
  \FunctionTok{mutate}\NormalTok{(}\AttributeTok{cumulative\_rate =} \FunctionTok{cumsum}\NormalTok{(choice) }\SpecialCharTok{/} \FunctionTok{row\_number}\NormalTok{()) }\SpecialCharTok{\%\textgreater{}\%}
  \FunctionTok{ungroup}\NormalTok{()}

  \CommentTok{\# Save the generated data}
  \FunctionTok{write\_csv}\NormalTok{(d, sim\_data\_file)}
  \FunctionTok{cat}\NormalTok{(}\StringTok{"Generated new simulation data and saved to"}\NormalTok{, sim\_data\_file, }\StringTok{"}\SpecialCharTok{\textbackslash{}n}\StringTok{"}\NormalTok{)}
\NormalTok{\} }\ControlFlowTok{else}\NormalTok{ \{}
  \CommentTok{\# Load existing simulation data}
\NormalTok{  d }\OtherTok{\textless{}{-}} \FunctionTok{read\_csv}\NormalTok{(sim\_data\_file)}
  \FunctionTok{cat}\NormalTok{(}\StringTok{"Loaded existing simulation data from"}\NormalTok{, sim\_data\_file, }\StringTok{"}\SpecialCharTok{\textbackslash{}n}\StringTok{"}\NormalTok{)}
\NormalTok{\}}
\end{Highlighting}
\end{Shaded}

\begin{verbatim}
Loaded existing simulation data from /Users/au209589/Dropbox/Teaching/AdvancedCognitiveModeling23_book/simdata/ch06_randomnoise.csv 
\end{verbatim}

\begin{Shaded}
\begin{Highlighting}[]
\CommentTok{\# {-}{-}{-} Visualize Simulated Data {-}{-}{-}}
\CommentTok{\# Plot cumulative rate, faceted by noise level}
\NormalTok{p1 }\OtherTok{\textless{}{-}} \FunctionTok{ggplot}\NormalTok{(d, }\FunctionTok{aes}\NormalTok{(}\AttributeTok{x =}\NormalTok{ trial, }\AttributeTok{y =}\NormalTok{ cumulative\_rate, }\AttributeTok{group =}\NormalTok{ true\_rate, }\AttributeTok{color =}\NormalTok{ true\_rate)) }\SpecialCharTok{+}
  \FunctionTok{geom\_line}\NormalTok{(}\AttributeTok{alpha =} \FloatTok{0.8}\NormalTok{) }\SpecialCharTok{+}
  \FunctionTok{geom\_hline}\NormalTok{(}\AttributeTok{yintercept =} \FloatTok{0.5}\NormalTok{, }\AttributeTok{linetype =} \StringTok{"dashed"}\NormalTok{) }\SpecialCharTok{+}
  \FunctionTok{scale\_color\_viridis\_c}\NormalTok{() }\SpecialCharTok{+} \CommentTok{\# Use perceptually uniform color scale}
  \FunctionTok{facet\_wrap}\NormalTok{(}\SpecialCharTok{\textasciitilde{}}\FunctionTok{paste}\NormalTok{(}\StringTok{"Noise ="}\NormalTok{, true\_noise)) }\SpecialCharTok{+} \CommentTok{\# Facet by noise level}
  \FunctionTok{labs}\NormalTok{(}
    \AttributeTok{title =} \StringTok{"Simulated Choice Behavior (Cumulative Rate)"}\NormalTok{,}
    \AttributeTok{subtitle =} \StringTok{"Agents with different bias rates and noise levels"}\NormalTok{,}
    \AttributeTok{x =} \StringTok{"Trial Number"}\NormalTok{,}
    \AttributeTok{y =} \StringTok{"Cumulative Proportion Choosing \textquotesingle{}Right\textquotesingle{}"}\NormalTok{,}
    \AttributeTok{color =} \StringTok{"True Bias Rate"}
\NormalTok{  ) }\SpecialCharTok{+}
  \FunctionTok{theme\_minimal}\NormalTok{() }\SpecialCharTok{+}
  \FunctionTok{ylim}\NormalTok{(}\DecValTok{0}\NormalTok{, }\DecValTok{1}\NormalTok{)}

\FunctionTok{print}\NormalTok{(p1)}
\end{Highlighting}
\end{Shaded}

\begin{center}
\includegraphics[width=0.8\linewidth,height=\textheight,keepaspectratio]{04-InferringRates_files/figure-pdf/ch05_simulating data-1.pdf}
\end{center}

\section{How the Posterior Is Actually
Computed}\label{how-the-posterior-is-actually-computed}

Before we hand this problem to Stan, it is worth understanding what Stan
is actually \emph{doing} --- because ``posterior ∝ prior × likelihood''
is easy to say and hard to see. This section builds a posterior by hand
using a \textbf{grid approximation}: a brute-force method that works
when there is only one parameter and we are willing to evaluate every
possible value of it.

\textbf{The setup.} Suppose an agent made 3 choices, picking ``right''
on trials 1 and 3 and ``left'' on trial 2 (\texttt{h\ =\ c(1,\ 0,\ 1)}).
We want to know: what is the most plausible value of θ --- the agent's
underlying bias toward ``right''?

\textbf{Bayes' rule} tells us:

\[P(\theta \mid \text{data}) \propto P(\theta) \times P(\text{data} \mid \theta)\]

In words: the posterior is proportional to the prior (what we believed
before seeing the data) multiplied by the likelihood (how probable the
data is under each value of θ). We then re-normalize so the result sums
to 1.

\begin{Shaded}
\begin{Highlighting}[]
\CommentTok{\# 1. Define a coarse grid of 10 candidate values for theta}
\NormalTok{theta\_grid }\OtherTok{\textless{}{-}} \FunctionTok{seq}\NormalTok{(}\FloatTok{0.1}\NormalTok{, }\FloatTok{1.0}\NormalTok{, }\AttributeTok{by =} \FloatTok{0.1}\NormalTok{)}

\CommentTok{\# 2. Prior: uniform — Beta(1,1) — every value equally plausible before seeing data}
\NormalTok{prior }\OtherTok{\textless{}{-}} \FunctionTok{rep}\NormalTok{(}\DecValTok{1}\NormalTok{, }\FunctionTok{length}\NormalTok{(theta\_grid))}
\NormalTok{prior }\OtherTok{\textless{}{-}}\NormalTok{ prior }\SpecialCharTok{/} \FunctionTok{sum}\NormalTok{(prior)   }\CommentTok{\# Normalize to sum to 1}

\CommentTok{\# 3. Observed data}
\NormalTok{h }\OtherTok{\textless{}{-}} \FunctionTok{c}\NormalTok{(}\DecValTok{1}\NormalTok{, }\DecValTok{0}\NormalTok{, }\DecValTok{1}\NormalTok{)   }\CommentTok{\# 2 "right" out of 3 trials}

\CommentTok{\# 4. Likelihood: P(data | theta)}
\CommentTok{\# Each choice is Bernoulli(theta), so the joint likelihood is the product}
\CommentTok{\# across trials}
\NormalTok{likelihood }\OtherTok{\textless{}{-}} \FunctionTok{sapply}\NormalTok{(theta\_grid, }\ControlFlowTok{function}\NormalTok{(th) \{}
  \FunctionTok{prod}\NormalTok{(}\FunctionTok{dbinom}\NormalTok{(h, }\AttributeTok{size =} \DecValTok{1}\NormalTok{, }\AttributeTok{prob =}\NormalTok{ th))}
\NormalTok{\})}

\CommentTok{\# 5. Unnormalized posterior: prior * likelihood (pointwise)}
\NormalTok{posterior\_unnorm }\OtherTok{\textless{}{-}}\NormalTok{ prior }\SpecialCharTok{*}\NormalTok{ likelihood}

\CommentTok{\# 6. Normalize so it sums to 1}
\NormalTok{posterior }\OtherTok{\textless{}{-}}\NormalTok{ posterior\_unnorm }\SpecialCharTok{/} \FunctionTok{sum}\NormalTok{(posterior\_unnorm)}

\CommentTok{\# {-}{-}{-} Display as a table {-}{-}{-}}
\NormalTok{grid\_table }\OtherTok{\textless{}{-}} \FunctionTok{tibble}\NormalTok{(}
  \StringTok{\textasciigrave{}}\AttributeTok{θ}\StringTok{\textasciigrave{}}           \OtherTok{=}\NormalTok{ theta\_grid,}
  \StringTok{\textasciigrave{}}\AttributeTok{Prior}\StringTok{\textasciigrave{}}       \OtherTok{=} \FunctionTok{round}\NormalTok{(prior, }\DecValTok{3}\NormalTok{),}
  \StringTok{\textasciigrave{}}\AttributeTok{Likelihood}\StringTok{\textasciigrave{}}  \OtherTok{=} \FunctionTok{round}\NormalTok{(likelihood, }\DecValTok{4}\NormalTok{),}
  \StringTok{\textasciigrave{}}\AttributeTok{Posterior}\StringTok{\textasciigrave{}}   \OtherTok{=} \FunctionTok{round}\NormalTok{(posterior, }\DecValTok{3}\NormalTok{)}
\NormalTok{)}

\NormalTok{knitr}\SpecialCharTok{::}\FunctionTok{kable}\NormalTok{(grid\_table,}
             \AttributeTok{caption =} \StringTok{"Grid approximation: prior × likelihood → posterior (3 observations)"}\NormalTok{)}
\end{Highlighting}
\end{Shaded}

\begin{longtable}[]{@{}rrrr@{}}
\caption{Grid approximation: prior × likelihood → posterior (3
observations)}\tabularnewline
\toprule\noalign{}
θ & Prior & Likelihood & Posterior \\
\midrule\noalign{}
\endfirsthead
\toprule\noalign{}
θ & Prior & Likelihood & Posterior \\
\midrule\noalign{}
\endhead
\bottomrule\noalign{}
\endlastfoot
0.1 & 0.1 & 0.009 & 0.011 \\
0.2 & 0.1 & 0.032 & 0.039 \\
0.3 & 0.1 & 0.063 & 0.076 \\
0.4 & 0.1 & 0.096 & 0.116 \\
0.5 & 0.1 & 0.125 & 0.152 \\
0.6 & 0.1 & 0.144 & 0.175 \\
0.7 & 0.1 & 0.147 & 0.178 \\
0.8 & 0.1 & 0.128 & 0.155 \\
0.9 & 0.1 & 0.081 & 0.098 \\
1.0 & 0.1 & 0.000 & 0.000 \\
\end{longtable}

Grid approximation with 3 observed choices. The posterior (blue) is the
pointwise product of the prior (red) and the likelihood (green),
re-normalized. Notice that the prior is flat --- every value of θ was
equally plausible beforehand --- and the posterior peaks near θ = 0.7,
which best explains seeing 2 `right' choices out of 3.

\begin{Shaded}
\begin{Highlighting}[]
\CommentTok{\# {-}{-}{-} Plot all three distributions {-}{-}{-}}
\NormalTok{grid\_long }\OtherTok{\textless{}{-}}\NormalTok{ grid\_table }\SpecialCharTok{\%\textgreater{}\%}
  \FunctionTok{rename}\NormalTok{(}\AttributeTok{theta =} \StringTok{\textasciigrave{}}\AttributeTok{θ}\StringTok{\textasciigrave{}}\NormalTok{, }\AttributeTok{prior =}\NormalTok{ Prior, }\AttributeTok{likelihood =}\NormalTok{ Likelihood, }\AttributeTok{posterior =}\NormalTok{ Posterior) }\SpecialCharTok{\%\textgreater{}\%}
  \FunctionTok{pivot\_longer}\NormalTok{(}\FunctionTok{c}\NormalTok{(prior, likelihood, posterior),}
               \AttributeTok{names\_to =} \StringTok{"distribution"}\NormalTok{, }\AttributeTok{values\_to =} \StringTok{"mass"}\NormalTok{) }\SpecialCharTok{\%\textgreater{}\%}
  \FunctionTok{mutate}\NormalTok{(}\AttributeTok{distribution =} \FunctionTok{factor}\NormalTok{(distribution,}
                                \AttributeTok{levels =} \FunctionTok{c}\NormalTok{(}\StringTok{"prior"}\NormalTok{, }\StringTok{"likelihood"}\NormalTok{, }\StringTok{"posterior"}\NormalTok{),}
                                \AttributeTok{labels =} \FunctionTok{c}\NormalTok{(}\StringTok{"Prior"}\NormalTok{, }\StringTok{"Likelihood"}\NormalTok{, }\StringTok{"Posterior"}\NormalTok{)))}

\FunctionTok{ggplot}\NormalTok{(grid\_long, }\FunctionTok{aes}\NormalTok{(}\AttributeTok{x =}\NormalTok{ theta, }\AttributeTok{y =}\NormalTok{ mass, }\AttributeTok{fill =}\NormalTok{ distribution)) }\SpecialCharTok{+}
  \FunctionTok{geom\_col}\NormalTok{(}\AttributeTok{position =} \StringTok{"dodge"}\NormalTok{, }\AttributeTok{alpha =} \FloatTok{0.8}\NormalTok{, }\AttributeTok{width =} \FloatTok{0.08}\NormalTok{) }\SpecialCharTok{+}
  \FunctionTok{scale\_fill\_manual}\NormalTok{(}\AttributeTok{values =} \FunctionTok{c}\NormalTok{(}\StringTok{"Prior"} \OtherTok{=} \StringTok{"red"}\NormalTok{, }\StringTok{"Likelihood"} \OtherTok{=} \StringTok{"darkgreen"}\NormalTok{,}
                                \StringTok{"Posterior"} \OtherTok{=} \StringTok{"blue"}\NormalTok{)) }\SpecialCharTok{+}
  \FunctionTok{scale\_x\_continuous}\NormalTok{(}\AttributeTok{breaks =}\NormalTok{ theta\_grid) }\SpecialCharTok{+}
  \FunctionTok{labs}\NormalTok{(}
    \AttributeTok{title    =} \StringTok{"Bayes by Hand: Prior × Likelihood → Posterior"}\NormalTok{,}
    \AttributeTok{subtitle =} \StringTok{"Data: h = c(1, 0, 1) — two \textquotesingle{}right\textquotesingle{} choices out of three trials"}\NormalTok{,}
    \AttributeTok{x        =} \StringTok{"θ  (probability of choosing \textquotesingle{}right\textquotesingle{})"}\NormalTok{,}
    \AttributeTok{y        =} \StringTok{"Probability mass"}\NormalTok{,}
    \AttributeTok{fill     =} \ConstantTok{NULL}
\NormalTok{  ) }\SpecialCharTok{+}
  \FunctionTok{theme\_classic}\NormalTok{() }\SpecialCharTok{+}
  \FunctionTok{theme}\NormalTok{(}\AttributeTok{legend.position =} \StringTok{"top"}\NormalTok{)}
\end{Highlighting}
\end{Shaded}

\begin{figure}[H]

{\centering \includegraphics[width=0.8\linewidth,height=\textheight,keepaspectratio]{04-InferringRates_files/figure-pdf/04-bayes-by-hand-1.pdf}

}

\caption{Grid approximation with 3 observed choices. The posterior
(blue) is the pointwise product of the prior (red) and the likelihood
(green), re-normalized. Notice that the prior is flat --- every value of
θ was equally plausible beforehand --- and the posterior peaks near θ =
0.7, which best explains seeing 2 `right' choices out of 3.}

\end{figure}%

Three things to notice from the table and the plot:

\begin{enumerate}
\def\labelenumi{\arabic{enumi}.}
\item
  \textbf{The prior is flat.} Before seeing any data, every value of θ
  is equally plausible. This is what \texttt{Beta(1,\ 1)} means ---
  maximal prior ignorance.
\item
  \textbf{The likelihood peaks near θ = 0.7.} The agent chose ``right''
  2 out of 3 times, so θ = 0.7 makes the observed sequence most
  probable. Values near 0 or 1 make it very improbable.
\item
  \textbf{The posterior is the pointwise product, re-normalized.} It is
  narrower than the prior (the data ruled out extreme values) and
  centred on the likelihood peak. This is exactly the ``update'': the
  data shifted our beliefs from flat to concentrated.
\end{enumerate}

\textbf{Why Stan instead of a grid?} This worked with 1 parameter and 10
grid points. With 2 parameters and 100 grid points each, we need 100 ×
100 = 10,000 evaluations. With 5 parameters: 100⁵ = \textbf{10 billion}.
The grid grows exponentially --- this is the \emph{curse of
dimensionality}. Stan's Hamiltonian Monte Carlo (HMC) sampler sidesteps
this by \emph{sampling} from regions of high posterior probability
rather than evaluating every possible combination. The samples it
returns are the blue histogram you will see in the next section --- same
idea, just no grid required.

\section{Building our First Stan Model: Inferring Bias
Rate}\label{building-our-first-stan-model-inferring-bias-rate}

N.B. Refer to the video and slides for the step by step build-up of the
Stan code.

Let's use Stan to estimate the underlying bias \texttt{rate} (which
we'll call \texttt{theta} in the model) from just the sequence of
choices (\texttt{h}) generated by one of our simulated agents. This
model is represented in the

\begin{figure}[H]

{\centering \includegraphics[width=0.8\linewidth,height=\textheight,keepaspectratio]{04-InferringRates_files/figure-pdf/unnamed-chunk-1-1.png}

}

\caption{Binomial Model with Formulas}

\end{figure}%

\begin{figure}[H]

{\centering \includegraphics[width=0.8\linewidth,height=\textheight,keepaspectratio]{04-InferringRates_files/figure-pdf/unnamed-chunk-2-1.png}

}

\caption{Kruschke-Style Binomial Model}

\end{figure}%

We'll start with a simple case: an agent with \texttt{rate\ =\ 0.8} and
\texttt{noise\ =\ 0}.

\textbf{Preparing Data for Stan} Why a list? Well, dataframes (now
tibbles) are amazing. But they have a big drawback: they require each
variable to have the same length. Lists do not have that limitation,
they are more flexible. So, lists. We'll have to learn how to live with
them.

\begin{Shaded}
\begin{Highlighting}[]
\CommentTok{\# Subset data for one specific agent simulation (rate=0.8, noise=0)}
\NormalTok{d1 }\OtherTok{\textless{}{-}}\NormalTok{ d }\SpecialCharTok{\%\textgreater{}\%} \FunctionTok{filter}\NormalTok{(true\_rate }\SpecialCharTok{==} \FloatTok{0.8}\NormalTok{, true\_noise }\SpecialCharTok{==} \FloatTok{0.0}\NormalTok{) }\CommentTok{\# Make sure noise level matches}

\CommentTok{\# Stan requires data in a list format}
\NormalTok{data\_for\_stan }\OtherTok{\textless{}{-}} \FunctionTok{list}\NormalTok{(}
  \AttributeTok{n =}\NormalTok{ trials,     }\CommentTok{\# Number of trials (integer)}
  \AttributeTok{h =}\NormalTok{ d1}\SpecialCharTok{$}\NormalTok{choice   }\CommentTok{\# Vector of choices (0s and 1s) for this agent}
\NormalTok{)}

\FunctionTok{str}\NormalTok{(data\_for\_stan) }\CommentTok{\# Show the structure of the list}
\end{Highlighting}
\end{Shaded}

\begin{verbatim}
List of 2
 $ n: num 120
 $ h: num [1:120] 1 1 0 1 1 0 1 1 1 1 ...
\end{verbatim}

\subsection{Model}\label{model}

We write the stan code within the R code (so I can show it to you more
easily), then we save it as a stan file, which can be loaded at a later
stage in order to compile it. {[}Missing: more info on compiling etc.{]}

Remember that the minimal Stan model requires 3 chunks, one specifying
the data it will need as input; one specifying the parameters to be
estimated; one specifying the model within which the parameters appear,
and the priors for those parameters.

\begin{Shaded}
\begin{Highlighting}[]
\CommentTok{\# This R chunk defines the Stan code as a string and writes it to a file.}
\CommentTok{\# It\textquotesingle{}s marked eval=FALSE because we don\textquotesingle{}t run the R code defining the string here,}
\CommentTok{\# but the Stan code itself IS the important content.}

\NormalTok{stan\_model\_code }\OtherTok{\textless{}{-}} \StringTok{"}
\StringTok{// Stan Model: Simple Bernoulli for Bias Estimation}
\StringTok{// Goal: Estimate the underlying probability (theta) of choosing 1 (\textquotesingle{}right\textquotesingle{})}
\StringTok{//       given a sequence of binary choices.}

\StringTok{// 1. Data Block: Declares the data Stan expects from R}
\StringTok{data \{}
\StringTok{  int\textless{}lower=1\textgreater{} n;        // Number of trials (must be at least 1)}
\StringTok{  array[n] int\textless{}lower=0, upper=1\textgreater{} h; // Array \textquotesingle{}h\textquotesingle{} of length \textquotesingle{}n\textquotesingle{} containing choices (0 or 1)}
\StringTok{\}}

\StringTok{// 2. Parameters Block: Declares the parameters the model will estimate}
\StringTok{parameters \{}
\StringTok{  real\textless{}lower=0, upper=1\textgreater{} theta; // The bias parameter (probability), constrained between 0 and 1}
\StringTok{\}}

\StringTok{// 3. Model Block: Defines the priors and the data model (used as likelihood during inference)}
\StringTok{model \{}
\StringTok{  // Prior: Our belief about theta *before* seeing the data.}
\StringTok{  // We use a Beta(1, 1) prior, which is equivalent to a Uniform(0, 1) distribution.}
\StringTok{  // This represents maximal prior ignorance about the bias.}
\StringTok{  // \textquotesingle{}target +=\textquotesingle{} adds the log{-}probability density to the overall model log{-}probability.}
\StringTok{  target += beta\_lpdf(theta | 1, 1); // lpdf = log probability density function}

\StringTok{  // Likelihood: How the data \textquotesingle{}h\textquotesingle{} depend on the parameter \textquotesingle{}theta\textquotesingle{}.}
\StringTok{  // We model each choice \textquotesingle{}h\textquotesingle{} as a Bernoulli trial with success probability \textquotesingle{}theta\textquotesingle{}.}
\StringTok{  // The model assesses how likely the observed sequence \textquotesingle{}h\textquotesingle{} is given a value of \textquotesingle{}theta\textquotesingle{}.}
\StringTok{  target += bernoulli\_lpmf(h | theta); // lpmf = log probability mass function (for discrete data)}
\StringTok{\}}

\StringTok{// 4. Generated Quantities Block (Optional but useful)}
\StringTok{// Code here is executed *after* sampling, using the estimated parameter values.}
\StringTok{// Useful for calculating derived quantities or predictions.}
\StringTok{generated quantities \{}
\StringTok{  // let\textquotesingle{}s generate the prior for visualization}
\StringTok{  // first we define it}
\StringTok{  real\textless{}lower = 0, upper = 1\textgreater{} prior\_theta; }
\StringTok{  // then we generate it}
\StringTok{  prior\_theta = beta\_rng(1, 1); // Draw a random sample from the prior distribution for theta}
\StringTok{\}}
\StringTok{"}
\CommentTok{\# Define relative path for Stan model file}
\NormalTok{stan\_model\_dir }\OtherTok{\textless{}{-}} \StringTok{"stan"}
\NormalTok{stan\_file\_bernoulli }\OtherTok{\textless{}{-}} \FunctionTok{here}\NormalTok{(stan\_model\_dir, }\StringTok{"ch06\_SimpleBernoulli.stan"}\NormalTok{)}

\CommentTok{\# Create directory if needed}
\ControlFlowTok{if}\NormalTok{ (}\SpecialCharTok{!}\FunctionTok{dir.exists}\NormalTok{(stan\_model\_dir)) }\FunctionTok{dir.create}\NormalTok{(stan\_model\_dir)}

\CommentTok{\# Write the Stan code to the file}
\FunctionTok{writeLines}\NormalTok{(stan\_model\_code, stan\_file\_bernoulli)}

\FunctionTok{cat}\NormalTok{(}\StringTok{"Stan model written to:"}\NormalTok{, stan\_file\_bernoulli, }\StringTok{"}\SpecialCharTok{\textbackslash{}n}\StringTok{"}\NormalTok{)}
\end{Highlighting}
\end{Shaded}

\begin{verbatim}
Stan model written to: /Users/au209589/Dropbox/Teaching/AdvancedCognitiveModeling23_book/stan/ch06_SimpleBernoulli.stan 
\end{verbatim}

Compiling and Fitting the Stan Model

Now we use cmdstanr to compile this model (translating it into efficient
C++ code) and then run the MCMC sampler to get posterior estimates for
theta.

\begin{Shaded}
\begin{Highlighting}[]
\DocumentationTok{\#\# {-}{-}{-} Compile and Fit {-}{-}{-}}
\CommentTok{\# Specify model file path}
\NormalTok{stan\_file\_bernoulli }\OtherTok{\textless{}{-}} \FunctionTok{here}\NormalTok{(}\StringTok{"stan"}\NormalTok{, }\StringTok{"ch06\_SimpleBernoulli.stan"}\NormalTok{)}
\CommentTok{\# Define path for saving fitted model object}
\NormalTok{stan\_results\_dir }\OtherTok{\textless{}{-}} \StringTok{"simmodels"}
\NormalTok{model\_file\_bernoulli }\OtherTok{\textless{}{-}} \FunctionTok{here}\NormalTok{(stan\_results\_dir, }\StringTok{"ch06\_SimpleBernoulli.rds"}\NormalTok{)}

\ControlFlowTok{if}\NormalTok{ (}\SpecialCharTok{!}\FunctionTok{dir.exists}\NormalTok{(stan\_results\_dir)) }\FunctionTok{dir.create}\NormalTok{(stan\_results\_dir)}

\ControlFlowTok{if}\NormalTok{ (regenerate\_simulations }\SpecialCharTok{||} \SpecialCharTok{!}\FunctionTok{file.exists}\NormalTok{(model\_file\_bernoulli)) \{}
  \FunctionTok{cat}\NormalTok{(}\StringTok{"Compiling and fitting Stan model (Bernoulli)...}\SpecialCharTok{\textbackslash{}n}\StringTok{"}\NormalTok{)}
  \CommentTok{\# Compile the Stan model (only needs to be done once unless code changes)}
\NormalTok{  mod\_bernoulli }\OtherTok{\textless{}{-}} \FunctionTok{cmdstan\_model}\NormalTok{(stan\_file\_bernoulli,}
                                 \AttributeTok{cpp\_options =} \FunctionTok{list}\NormalTok{(}\AttributeTok{stan\_threads =} \ConstantTok{FALSE}\NormalTok{), }\CommentTok{\# Enable threading}
                                 \AttributeTok{stanc\_options =} \FunctionTok{list}\NormalTok{(}\StringTok{"O1"}\NormalTok{)) }\CommentTok{\# Basic optimization}

  \CommentTok{\# Sample from the posterior distribution using MCMC}
\NormalTok{  fit\_bernoulli }\OtherTok{\textless{}{-}}\NormalTok{ mod\_bernoulli}\SpecialCharTok{$}\FunctionTok{sample}\NormalTok{(}
    \AttributeTok{data =}\NormalTok{ data\_for\_stan,          }\CommentTok{\# The data list we prepared}
    \AttributeTok{seed =} \DecValTok{123}\NormalTok{,                    }\CommentTok{\# For reproducible MCMC sampling}
    \AttributeTok{chains =} \DecValTok{4}\NormalTok{,                    }\CommentTok{\# Number of parallel Markov chains (recommend 4)}
    \AttributeTok{parallel\_chains =} \DecValTok{4}\NormalTok{, }\CommentTok{\# Run chains in parallel}
    \CommentTok{\#threads\_per\_chain = 1,         \# For within{-}chain parallelization (usually 1 is fine)}
    \AttributeTok{iter\_warmup =} \DecValTok{1000}\NormalTok{,            }\CommentTok{\# Number of warmup iterations per chain (discarded)}
    \AttributeTok{iter\_sampling =} \DecValTok{2000}\NormalTok{,          }\CommentTok{\# Number of sampling iterations per chain (kept)}
    \AttributeTok{refresh =} \DecValTok{0}\NormalTok{,                   }\CommentTok{\# How often to print progress updates. 0 means never since the model is so fast}
    \AttributeTok{max\_treedepth =} \DecValTok{10}\NormalTok{,            }\CommentTok{\# Controls complexity of MCMC steps (adjust if needed)}
    \AttributeTok{adapt\_delta =} \FloatTok{0.8}              \CommentTok{\# Target acceptance rate (adjust if divergences occur)}
\NormalTok{  )}

  \CommentTok{\# Save the fitted model object}
\NormalTok{  fit\_bernoulli}\SpecialCharTok{$}\FunctionTok{save\_object}\NormalTok{(}\AttributeTok{file =}\NormalTok{ model\_file\_bernoulli)}
  \FunctionTok{cat}\NormalTok{(}\StringTok{"Model fit completed and saved to:"}\NormalTok{, model\_file\_bernoulli, }\StringTok{"}\SpecialCharTok{\textbackslash{}n}\StringTok{"}\NormalTok{)}
\NormalTok{\} }\ControlFlowTok{else}\NormalTok{ \{}
  \CommentTok{\# Load existing fitted model object}
\NormalTok{  fit\_bernoulli }\OtherTok{\textless{}{-}} \FunctionTok{readRDS}\NormalTok{(model\_file\_bernoulli)}
  \FunctionTok{cat}\NormalTok{(}\StringTok{"Loaded existing model fit from:"}\NormalTok{, model\_file\_bernoulli, }\StringTok{"}\SpecialCharTok{\textbackslash{}n}\StringTok{"}\NormalTok{)}
\NormalTok{\}}
\end{Highlighting}
\end{Shaded}

\begin{verbatim}
Loaded existing model fit from: /Users/au209589/Dropbox/Teaching/AdvancedCognitiveModeling23_book/simmodels/ch06_SimpleBernoulli.rds 
\end{verbatim}

\subsection{Assessing MCMC Health (The Mandatory
Diagnostics)}\label{assessing-mcmc-health-the-mandatory-diagnostics}

\begin{quote}
\emph{Why check diagnostics before looking at the estimates?} The
posterior distribution Stan returns is only valid if the sampler
explored the parameter space correctly. If it did not --- due to
divergent transitions, poor mixing, or insufficient samples --- the
numbers in the posterior are artifacts of the algorithm, not properties
of the model. Looking at estimates from a broken sampler is like reading
a map drawn by a robot that kept falling over: the map looks plausible
but cannot be trusted. Diagnostics are the check that the robot walked
properly before you navigate with its map.
\end{quote}

Before we look at the parameter estimates, we MUST check whether the
model was able to sample the posterior space: the health of the MCMC
chains. Fitting a model is like sending a robotic explorer (the sampler)
to map a landscape (the posterior). If the robot's tracks are broken,
the map is mathematically invalid. We must check:

\begin{enumerate}
\def\labelenumi{\arabic{enumi}.}
\tightlist
\item
  \textbf{Divergences:} Did the simulation crash? (Zero tolerance).
\item
  \textbf{\(\hat{R}\) (R-hat):} Did our 4 independent chains map the
  exact same space? (Must be strictly \(< 1.01\)).
\item
  \textbf{Effective Sample Size (ESS):} Do we have enough \emph{truly
  independent} draws? (Bulk and Tail ESS should be \(> 400\)).
\item
  \textbf{Visual Diagnostics:} Trace plots (should look like a ``hairy
  caterpillar'') and Rank Histograms (should be uniform across chains).
\end{enumerate}

Let's run this mandatory diagnostic battery on our first fit:

\begin{Shaded}
\begin{Highlighting}[]
\CommentTok{\# 1. Check for Divergent Transitions}
\NormalTok{diagnostics }\OtherTok{\textless{}{-}}\NormalTok{ fit\_bernoulli}\SpecialCharTok{$}\FunctionTok{diagnostic\_summary}\NormalTok{()}
\FunctionTok{cat}\NormalTok{(}\StringTok{"Divergent transitions per chain:"}\NormalTok{, }\FunctionTok{paste}\NormalTok{(diagnostics}\SpecialCharTok{$}\NormalTok{num\_divergent, }\AttributeTok{collapse =} \StringTok{", "}\NormalTok{), }\StringTok{"}\SpecialCharTok{\textbackslash{}n}\StringTok{"}\NormalTok{)}
\end{Highlighting}
\end{Shaded}

\begin{verbatim}
Divergent transitions per chain: 0, 0, 0, 0 
\end{verbatim}

\begin{Shaded}
\begin{Highlighting}[]
\CommentTok{\# 2. Check R{-}hat and ESS metrics}
\NormalTok{fit\_summary }\OtherTok{\textless{}{-}}\NormalTok{ fit\_bernoulli}\SpecialCharTok{$}\FunctionTok{summary}\NormalTok{(}\AttributeTok{variables =} \StringTok{"theta"}\NormalTok{)}
\FunctionTok{print}\NormalTok{(fit\_summary }\SpecialCharTok{\%\textgreater{}\%}\NormalTok{ dplyr}\SpecialCharTok{::}\FunctionTok{select}\NormalTok{(variable, mean, rhat, ess\_bulk, ess\_tail))}
\end{Highlighting}
\end{Shaded}

\begin{verbatim}
# A tibble: 1 x 5
  variable  mean  rhat ess_bulk ess_tail
  <chr>    <dbl> <dbl>    <dbl>    <dbl>
1 theta    0.836  1.00    2661.    3744.
\end{verbatim}

\begin{Shaded}
\begin{Highlighting}[]
\CommentTok{\# 3. Visual Checks (Trace and Rank Plots)}
\NormalTok{draws\_df }\OtherTok{\textless{}{-}} \FunctionTok{as\_draws\_df}\NormalTok{(fit\_bernoulli}\SpecialCharTok{$}\FunctionTok{draws}\NormalTok{())}

\CommentTok{\# A healthy trace plot is well{-}mixed with no wandering}
\NormalTok{p\_trace }\OtherTok{\textless{}{-}} \FunctionTok{mcmc\_trace}\NormalTok{(draws\_df, }\AttributeTok{pars =} \StringTok{"theta"}\NormalTok{) }\SpecialCharTok{+} 
  \FunctionTok{theme\_classic}\NormalTok{() }\SpecialCharTok{+} 
  \FunctionTok{ggtitle}\NormalTok{(}\StringTok{"Trace Plot"}\NormalTok{, }
          \AttributeTok{subtitle =} \StringTok{"Healthy = all chains overlap and fluctuate around the same stable band"}\NormalTok{)}

\CommentTok{\# A healthy rank histogram has uniform bars across all chains}
\NormalTok{p\_rank }\OtherTok{\textless{}{-}} \FunctionTok{mcmc\_rank\_overlay}\NormalTok{(draws\_df, }\AttributeTok{pars =} \StringTok{"theta"}\NormalTok{, }\AttributeTok{ref\_line =} \ConstantTok{TRUE}\NormalTok{) }\SpecialCharTok{+} 
  \FunctionTok{theme\_classic}\NormalTok{() }\SpecialCharTok{+} 
  \FunctionTok{ggtitle}\NormalTok{(}\StringTok{"Trank plot"}\NormalTok{, }\AttributeTok{subtitle =} \StringTok{"Healthy = chains interleave uniformly around the reference line — no chain dominates any rank band"}\NormalTok{)}

\NormalTok{p\_trace }\SpecialCharTok{/}\NormalTok{ p\_rank}
\end{Highlighting}
\end{Shaded}

\begin{center}
\includegraphics[width=0.8\linewidth,height=\textheight,keepaspectratio]{04-InferringRates_files/figure-pdf/04 biased model quality-1.pdf}
\end{center}

\begin{Shaded}
\begin{Highlighting}[]
\CommentTok{\# 4. Look at the Posterior Estimate}
\FunctionTok{ggplot}\NormalTok{(draws\_df) }\SpecialCharTok{+}
  \FunctionTok{geom\_density}\NormalTok{(}\FunctionTok{aes}\NormalTok{(theta, }\AttributeTok{fill =} \StringTok{"Posterior"}\NormalTok{), }\AttributeTok{alpha =} \FloatTok{0.6}\NormalTok{) }\SpecialCharTok{+}
  \FunctionTok{geom\_density}\NormalTok{(}\FunctionTok{aes}\NormalTok{(prior\_theta, }\AttributeTok{fill =} \StringTok{"Prior"}\NormalTok{), }\AttributeTok{alpha =} \FloatTok{0.6}\NormalTok{) }\SpecialCharTok{+}
  \FunctionTok{geom\_vline}\NormalTok{(}\AttributeTok{xintercept =} \FloatTok{0.8}\NormalTok{, }\AttributeTok{linetype =} \StringTok{"dashed"}\NormalTok{, }\AttributeTok{color =} \StringTok{"black"}\NormalTok{, }\AttributeTok{linewidth =} \FloatTok{1.2}\NormalTok{) }\SpecialCharTok{+}
  \FunctionTok{scale\_fill\_manual}\NormalTok{(}\AttributeTok{values =} \FunctionTok{c}\NormalTok{(}\StringTok{"Posterior"} \OtherTok{=} \StringTok{"blue"}\NormalTok{, }\StringTok{"Prior"} \OtherTok{=} \StringTok{"red"}\NormalTok{)) }\SpecialCharTok{+}
  \FunctionTok{labs}\NormalTok{(}
    \AttributeTok{title    =} \StringTok{"Prior{-}Posterior Update: Bias Parameter (θ)"}\NormalTok{,}
    \AttributeTok{subtitle =} \StringTok{"Healthy = posterior (blue) is narrower than prior (red) and shifted toward the true value (dashed)"}\NormalTok{,}
    \AttributeTok{x        =} \StringTok{"θ  (Probability of Choosing Right)"}\NormalTok{,}
    \AttributeTok{y        =} \StringTok{"Density"}\NormalTok{,}
    \AttributeTok{fill     =} \StringTok{"Distribution"}
\NormalTok{  ) }\SpecialCharTok{+}
  \FunctionTok{theme\_classic}\NormalTok{()}
\end{Highlighting}
\end{Shaded}

\begin{center}
\includegraphics[width=0.8\linewidth,height=\textheight,keepaspectratio]{04-InferringRates_files/figure-pdf/04 biased model quality-2.pdf}
\end{center}

The posterior (blue) is substantially narrower than the prior (red) and
has shifted toward the true value of 0.8 (dashed line). The posterior
mean is approximately 0.755, not exactly 0.8 --- this is expected with
only 120 trials and is not a sign that something is wrong. What matters
is that (a) the posterior is much more informative than the prior,
confirming the data were indeed informative, and (b) the true value 0.8
sits comfortably within the credible interval. If the posterior were
identical to the prior, the data would be telling us nothing; if the
posterior were centered far from the true value and excluded it, the
model would be misspecified.

Now we build the same model, but using the log odds scale for the theta
parameter, which will become useful later when we condition theta on
variables and build multilevel models (as we can do what we want in a
log odds space and it will always be bound between 0 and 1).

\textbf{Alternative Parameterization: Log-Odds Scale}

While estimating theta directly as a probability (0-1) is intuitive, it
can sometimes be computationally advantageous or necessary for more
complex models (like multilevel models in a later chapter) to estimate
parameters on an unbounded scale. The logit (log-odds) transformation
converts a probability \(p\) (from 0 to 1) to log-odds (from \(-\infty\)
to \(+\infty\)): \[\text{logit}(p) = \log\left(\frac{p}{1-p}\right)\]
The inverse transformation (logistic function or inv\_logit) converts
log-odds back to probability:
\[\text{inv\_logit}(x) = \frac{1}{1 + \exp(-x)}\] We can rewrite our
Stan model to estimate theta on the log-odds scale. However, we must be
careful with our priors, as the non-linear transformation distorts the
probability mass.

\begin{itemize}
\item
  Too Narrow: A standard normal prior
  \(\lambda \sim \text{Normal}(0, 1)\) concentrates mass around
  \(p=0.5\), biasing the model toward chance.
\item
  Too Wide: A wide prior (e.g., \(\lambda \sim \text{Normal}(0, 5)\))
  pushes mass to the extremes (0 and 1), implying we expect
  deterministic behavior.
\item
  Just Right: To approximate a uniform distribution on the probability
  scale (like our Beta(1,1)), we typically use
  \(\lambda \sim \text{Normal}(0, 1.5)\). This is the specific ``sweet
  spot'' that is essentially flat, avoiding both the center-bias of
  narrow priors and the extreme-bias of wide priors.
\end{itemize}

\begin{Shaded}
\begin{Highlighting}[]
\CommentTok{\# Stan code using logit parameterization}
\NormalTok{stan\_model\_logit\_code }\OtherTok{\textless{}{-}} \StringTok{"}
\StringTok{// Stan Model: Simple Bernoulli (Logit Parameterization)}
\StringTok{// Estimate theta on the log{-}odds scale}

\StringTok{data \{}
\StringTok{  int\textless{}lower=1\textgreater{} n;}
\StringTok{  array[n] int\textless{}lower=0, upper=1\textgreater{} h;}
\StringTok{\}}

\StringTok{parameters \{}
\StringTok{  real theta\_logit; // Parameter is now on the unbounded log{-}odds scale}
\StringTok{\}}

\StringTok{model \{}
\StringTok{  // Prior on log{-}odds scale (e.g., Normal(0, 1.5))}
\StringTok{  // Normal(0, 1.5) on log{-}odds corresponds roughly to a diffuse prior on probability scale}
\StringTok{  target += normal\_lpdf(theta\_logit | 0, 1.5);}

\StringTok{  // Likelihood using the logit version of the Bernoulli PMF}
\StringTok{  // This tells Stan that theta\_logit is on the log{-}odds scale}
\StringTok{  target += bernoulli\_logit\_lpmf(h | theta\_logit);}
\StringTok{\}}

\StringTok{generated quantities \{}
\StringTok{  // Generate the prior for viz purposes}
\StringTok{  real prior\_theta = inv\_logit(normal\_rng(0, 1.5)); // Sample from the log{-}odds prior}
\StringTok{  // Convert estimate back to probability scale for easier interpretation}
\StringTok{  real\textless{}lower=0, upper=1\textgreater{} theta = inv\_logit(theta\_logit);}
\StringTok{\}}
\StringTok{"}
\CommentTok{\# Define file path}
\NormalTok{stan\_file\_logit }\OtherTok{\textless{}{-}} \FunctionTok{here}\NormalTok{(stan\_model\_dir, }\StringTok{"ch06\_SimpleBernoulli\_logodds.stan"}\NormalTok{)}
\CommentTok{\# Write the Stan code to the file}
\FunctionTok{writeLines}\NormalTok{(stan\_model\_logit\_code, stan\_file\_logit)}
\FunctionTok{cat}\NormalTok{(}\StringTok{"Stan model (logit) written to:"}\NormalTok{, stan\_file\_logit, }\StringTok{"}\SpecialCharTok{\textbackslash{}n}\StringTok{"}\NormalTok{)}
\end{Highlighting}
\end{Shaded}

\begin{verbatim}
Stan model (logit) written to: /Users/au209589/Dropbox/Teaching/AdvancedCognitiveModeling23_book/stan/ch06_SimpleBernoulli_logodds.stan 
\end{verbatim}

\begin{tcolorbox}[enhanced jigsaw, toprule=.15mm, toptitle=1mm, bottomtitle=1mm, leftrule=.75mm, rightrule=.15mm, colframe=quarto-callout-tip-color-frame, left=2mm, colback=white, titlerule=0mm, coltitle=black, colbacktitle=quarto-callout-tip-color!10!white, opacityback=0, title=\textcolor{quarto-callout-tip-color}{\faLightbulb}\hspace{0.5em}{Stan Efficiency --- Sufficient Statistics: Aggregate IID Trials}, arc=.35mm, opacitybacktitle=0.6, breakable, bottomrule=.15mm]

When all trials are \textbf{independently and identically distributed}
(no memory, no sequential dependency, one \(\theta\) for the whole
session), you can replace a loop of \(N\) Bernoulli calls with a single
Binomial call on the \textbf{sufficient statistic} --- the total number
of successes.

\begin{Shaded}
\begin{Highlighting}[]
\CommentTok{// Slower: N separate Bernoulli evaluations (one autodiff node each)}
\ControlFlowTok{for}\NormalTok{ (i }\ControlFlowTok{in} \DecValTok{1}\NormalTok{:N)}
  \KeywordTok{target +=}\NormalTok{ bernoulli\_logit\_lpmf(h[i] | theta\_logit);}

\CommentTok{// Faster: one Binomial evaluation — mathematically identical for IID data}
\KeywordTok{target +=}\NormalTok{ binomial\_logit\_lpmf(sum(h) | N, theta\_logit);}
\end{Highlighting}
\end{Shaded}

Pass \texttt{sum(h)} and \texttt{N} as integers rather than the full
array --- Stan never sees the individual trial values.

\textbf{When this applies:} single-subject models with one rate
parameter per session (Chapters 4--5).

\textbf{When it does not apply:} any model where the likelihood varies
by trial --- memory models, multilevel models with per-subject
parameters indexed by \texttt{subj\_id}, sequential belief-updating
models. For those, the per-trial loop is necessary.

\end{tcolorbox}

\begin{Shaded}
\begin{Highlighting}[]
\DocumentationTok{\#\# With the logit format}
\DocumentationTok{\#\# Specify where the model is}
\NormalTok{file }\OtherTok{\textless{}{-}} \FunctionTok{here}\NormalTok{(}\StringTok{"stan"}\NormalTok{, }\StringTok{"ch06\_SimpleBernoulli\_logodds.stan"}\NormalTok{)}

\CommentTok{\# File path for saved model}
\NormalTok{model\_file }\OtherTok{\textless{}{-}} \FunctionTok{here}\NormalTok{(}\StringTok{"simmodels"}\NormalTok{,}\StringTok{"ch06\_SimpleBernoulli\_logodds.rds"}\NormalTok{)}

\CommentTok{\# Check if we need to rerun the simulation}
\ControlFlowTok{if}\NormalTok{ (regenerate\_simulations }\SpecialCharTok{||} \SpecialCharTok{!}\FunctionTok{file.exists}\NormalTok{(model\_file)) \{}
  \CommentTok{\# Compile the model}
\NormalTok{  mod }\OtherTok{\textless{}{-}} \FunctionTok{cmdstan\_model}\NormalTok{(file, }
                       \AttributeTok{cpp\_options =} \FunctionTok{list}\NormalTok{(}\AttributeTok{stan\_threads =} \ConstantTok{FALSE}\NormalTok{),}
                       \AttributeTok{stanc\_options =} \FunctionTok{list}\NormalTok{(}\StringTok{"O1"}\NormalTok{))}
  
  \CommentTok{\# The following command calls Stan with specific options.}
\NormalTok{  samples\_biased\_logodds }\OtherTok{\textless{}{-}}\NormalTok{ mod}\SpecialCharTok{$}\FunctionTok{sample}\NormalTok{(}
    \AttributeTok{data =}\NormalTok{ data\_for\_stan,}
    \AttributeTok{seed =} \DecValTok{123}\NormalTok{,}
    \AttributeTok{chains =} \DecValTok{2}\NormalTok{,}
    \AttributeTok{parallel\_chains =} \DecValTok{2}\NormalTok{,}
    \CommentTok{\#threads\_per\_chain = 1,}
    \AttributeTok{iter\_warmup =} \DecValTok{1000}\NormalTok{,}
    \AttributeTok{iter\_sampling =} \DecValTok{2000}\NormalTok{,}
    \AttributeTok{refresh =} \DecValTok{0}\NormalTok{,}
    \AttributeTok{max\_treedepth =} \DecValTok{20}\NormalTok{,}
    \AttributeTok{adapt\_delta =} \FloatTok{0.99}\NormalTok{,}
\NormalTok{  )}
  
  \CommentTok{\# Save the fitted model}
\NormalTok{  samples\_biased\_logodds}\SpecialCharTok{$}\FunctionTok{save\_object}\NormalTok{(}\AttributeTok{file =}\NormalTok{ model\_file)}
  \FunctionTok{cat}\NormalTok{(}\StringTok{"Generated new model fit and saved to"}\NormalTok{, model\_file, }\StringTok{"}\SpecialCharTok{\textbackslash{}n}\StringTok{"}\NormalTok{)}
\NormalTok{\} }\ControlFlowTok{else}\NormalTok{ \{}
  \CommentTok{\# Load existing results}
  \FunctionTok{cat}\NormalTok{(}\StringTok{"Loading biased model (log{-}odds) samples...}\SpecialCharTok{\textbackslash{}n}\StringTok{"}\NormalTok{)}
\NormalTok{  samples\_biased\_logodds }\OtherTok{\textless{}{-}} \FunctionTok{readRDS}\NormalTok{(model\_file)}
  \FunctionTok{cat}\NormalTok{(}\StringTok{"Available parameters:"}\NormalTok{, }\FunctionTok{paste}\NormalTok{(}\FunctionTok{colnames}\NormalTok{(}\FunctionTok{as\_draws\_df}\NormalTok{(samples\_biased\_logodds}\SpecialCharTok{$}\FunctionTok{draws}\NormalTok{())), }\AttributeTok{collapse =} \StringTok{", "}\NormalTok{), }\StringTok{"}\SpecialCharTok{\textbackslash{}n}\StringTok{"}\NormalTok{)}
  \FunctionTok{cat}\NormalTok{(}\StringTok{"Loaded existing model fit from"}\NormalTok{, model\_file, }\StringTok{"}\SpecialCharTok{\textbackslash{}n}\StringTok{"}\NormalTok{)}
\NormalTok{\}}
\end{Highlighting}
\end{Shaded}

\begin{verbatim}
Loading biased model (log-odds) samples...
Available parameters: lp__, theta_logit, prior_theta, theta, .chain, .iteration, .draw 
Loaded existing model fit from /Users/au209589/Dropbox/Teaching/AdvancedCognitiveModeling23_book/simmodels/ch06_SimpleBernoulli_logodds.rds 
\end{verbatim}

\subsection{Summarizing the results}\label{summarizing-the-results}

\begin{Shaded}
\begin{Highlighting}[]
\CommentTok{\# Summary}
\NormalTok{samples\_biased\_logodds}\SpecialCharTok{$}\FunctionTok{summary}\NormalTok{()}
\end{Highlighting}
\end{Shaded}

\begin{verbatim}
# A tibble: 4 x 10
  variable       mean  median     sd    mad       q5     q95  rhat ess_bulk
  <chr>         <dbl>   <dbl>  <dbl>  <dbl>    <dbl>   <dbl> <dbl>    <dbl>
1 lp__        -54.9   -54.6   0.725  0.330  -56.2    -54.4   1.000    1234.
2 theta_logit   1.65    1.65  0.247  0.249    1.25     2.04  1.00      884.
3 prior_theta   0.506   0.507 0.274  0.358    0.0790   0.925 1.00     3947.
4 theta         0.836   0.838 0.0337 0.0334   0.777    0.885 1.00      884.
# i 1 more variable: ess_tail <dbl>
\end{verbatim}

\begin{Shaded}
\begin{Highlighting}[]
\CommentTok{\# Extract posterior samples and include sampling of the prior:}
\NormalTok{draws\_df\_biased\_logodds }\OtherTok{\textless{}{-}} \FunctionTok{as\_draws\_df}\NormalTok{(samples\_biased\_logodds}\SpecialCharTok{$}\FunctionTok{draws}\NormalTok{()) }

\CommentTok{\# Explicitly extract theta parameter}
\NormalTok{theta\_param\_logodds }\OtherTok{\textless{}{-}}\NormalTok{ draws\_df\_biased\_logodds}\SpecialCharTok{$}\NormalTok{theta}
\FunctionTok{cat}\NormalTok{(}\StringTok{"Successfully extracted theta parameter with"}\NormalTok{, }\FunctionTok{length}\NormalTok{(theta\_param\_logodds), }\StringTok{"values}\SpecialCharTok{\textbackslash{}n}\StringTok{"}\NormalTok{)}
\end{Highlighting}
\end{Shaded}

\begin{verbatim}
Successfully extracted theta parameter with 4000 values
\end{verbatim}

\begin{Shaded}
\begin{Highlighting}[]
\FunctionTok{ggplot}\NormalTok{(draws\_df\_biased\_logodds, }\FunctionTok{aes}\NormalTok{(.iteration, theta, }\AttributeTok{group =}\NormalTok{ .chain, }\AttributeTok{color =}\NormalTok{ .chain)) }\SpecialCharTok{+}
  \FunctionTok{geom\_line}\NormalTok{() }\SpecialCharTok{+}
  \FunctionTok{theme\_classic}\NormalTok{()}
\end{Highlighting}
\end{Shaded}

\begin{center}
\includegraphics[width=0.8\linewidth,height=\textheight,keepaspectratio]{04-InferringRates_files/figure-pdf/04 biased model log-odds quality assessment-1.pdf}
\end{center}

\begin{Shaded}
\begin{Highlighting}[]
\CommentTok{\# Now let\textquotesingle{}s plot the density for theta (prior and posterior)}
\FunctionTok{ggplot}\NormalTok{(draws\_df\_biased\_logodds) }\SpecialCharTok{+}
  \FunctionTok{geom\_density}\NormalTok{(}\FunctionTok{aes}\NormalTok{(theta), }\AttributeTok{fill =} \StringTok{"blue"}\NormalTok{, }\AttributeTok{alpha =} \FloatTok{0.8}\NormalTok{) }\SpecialCharTok{+}
  \FunctionTok{geom\_density}\NormalTok{(}\FunctionTok{aes}\NormalTok{(prior\_theta), }\AttributeTok{fill =} \StringTok{"red"}\NormalTok{, }\AttributeTok{alpha =} \FloatTok{0.8}\NormalTok{) }\SpecialCharTok{+}
  \FunctionTok{geom\_vline}\NormalTok{(}\AttributeTok{xintercept =} \FloatTok{0.8}\NormalTok{, }\AttributeTok{linetype =} \StringTok{"dashed"}\NormalTok{, }\AttributeTok{color =} \StringTok{"black"}\NormalTok{, }\AttributeTok{size =} \FloatTok{1.5}\NormalTok{) }\SpecialCharTok{+}
  \FunctionTok{xlab}\NormalTok{(}\StringTok{"Rate"}\NormalTok{) }\SpecialCharTok{+}
  \FunctionTok{ylab}\NormalTok{(}\StringTok{"Posterior Density"}\NormalTok{) }\SpecialCharTok{+}
  \FunctionTok{theme\_classic}\NormalTok{()}
\end{Highlighting}
\end{Shaded}

\begin{center}
\includegraphics[width=0.8\linewidth,height=\textheight,keepaspectratio]{04-InferringRates_files/figure-pdf/04 biased model log-odds quality assessment-2.pdf}
\end{center}

We can see that the results are very similar.

\section{Bonus: Pathfinder Variational Inference vs.~Full
MCMC}\label{bonus-pathfinder-variational-inference-vs.-full-mcmc}

\begin{tcolorbox}[enhanced jigsaw, toprule=.15mm, toptitle=1mm, bottomtitle=1mm, leftrule=.75mm, rightrule=.15mm, colframe=quarto-callout-tip-color-frame, left=2mm, colback=white, titlerule=0mm, coltitle=black, colbacktitle=quarto-callout-tip-color!10!white, opacityback=0, title=\textcolor{quarto-callout-tip-color}{\faLightbulb}\hspace{0.5em}{Stan Optimization: Pathfinder as Warm-Start Initialization}, arc=.35mm, opacitybacktitle=0.6, breakable, bottomrule=.15mm]

Stan's default NUTS-HMC sampler draws posterior samples by simulating a
Hamiltonian dynamical system. During \textbf{warm-up}, the sampler
adapts its step size and metric --- a phase that can take hundreds of
iterations if chains start far from the posterior mass.
\textbf{Pathfinder} (Zhang et al. 2022) is a fast variational
approximation that traces the L-BFGS optimization path from an initial
point to the posterior mode and fits a series of multivariate normal
approximations along the way. It is not a gold-standard posterior, but
it tells HMC \emph{roughly where the posterior lives} before NUTS
begins.

Three strategies, in order of increasing accuracy and cost:

\begin{longtable}[]{@{}
  >{\raggedright\arraybackslash}p{(\linewidth - 4\tabcolsep) * \real{0.3333}}
  >{\raggedright\arraybackslash}p{(\linewidth - 4\tabcolsep) * \real{0.3333}}
  >{\raggedright\arraybackslash}p{(\linewidth - 4\tabcolsep) * \real{0.3333}}@{}}
\toprule\noalign{}
\begin{minipage}[b]{\linewidth}\raggedright
Strategy
\end{minipage} & \begin{minipage}[b]{\linewidth}\raggedright
When to use
\end{minipage} & \begin{minipage}[b]{\linewidth}\raggedright
Typical speed
\end{minipage} \\
\midrule\noalign{}
\endhead
\bottomrule\noalign{}
\endlastfoot
\textbf{Pathfinder alone} & Quick exploration, prior sensitivity checks
& \textasciitilde10--50× faster than NUTS \\
\textbf{NUTS with default init} & Standard analysis & Baseline \\
\textbf{Pathfinder → NUTS warm-start} & Hierarchical or complex models
where warm-up is slow & 20--60\% faster than default NUTS \\
\end{longtable}

The warm-start trick: run \texttt{mod\$pathfinder(...)} to get a
posterior approximation, then pass it as \texttt{init} to
\texttt{mod\$sample(...)}. This replaces the random Uniform(−2, 2)
starting points with draws from the Pathfinder approximation, letting
the sampler skip most of the warm-up trajectory. We will use this
pattern throughout the book for models with challenging geometry or many
parameters.

\textbf{When Pathfinder fails:} Near-degenerate posteriors (e.g., hard
funnels, flat ridges) can cause Pathfinder's variational approximation
to collapse. Always wrap it in
\texttt{tryCatch(...,\ error\ =\ function(e)\ NULL)} and fall back to
explicit per-chain inits centered on prior medians if it returns
\texttt{NULL}.

\end{tcolorbox}

\begin{Shaded}
\begin{Highlighting}[]
\DocumentationTok{\#\# Bonus content: Benchmarking NUTS, Pathfinder, and Pathfinder{-}initialized NUTS}

\CommentTok{\# We will compare three methods:}
\CommentTok{\# 1. Standard NUTS (MCMC)}
\CommentTok{\# 2. Pathfinder (Variational Inference)}
\CommentTok{\# 3. NUTS initialized with Pathfinder}

\CommentTok{\# Define file paths for all three models}
\NormalTok{file\_nuts\_std  }\OtherTok{\textless{}{-}} \FunctionTok{here}\NormalTok{(}\StringTok{"simmodels"}\NormalTok{, }\StringTok{"ch06\_benchmark\_nuts\_std.rds"}\NormalTok{)}
\NormalTok{file\_pathfinder }\OtherTok{\textless{}{-}} \FunctionTok{here}\NormalTok{(}\StringTok{"simmodels"}\NormalTok{, }\StringTok{"ch06\_benchmark\_pathfinder.rds"}\NormalTok{)}
\NormalTok{file\_nuts\_init  }\OtherTok{\textless{}{-}} \FunctionTok{here}\NormalTok{(}\StringTok{"simmodels"}\NormalTok{, }\StringTok{"ch06\_benchmark\_nuts\_init.rds"}\NormalTok{)}

\NormalTok{stan\_file\_logit }\OtherTok{\textless{}{-}}\NormalTok{ here}\SpecialCharTok{::}\FunctionTok{here}\NormalTok{(}\StringTok{"stan"}\NormalTok{, }\StringTok{"ch06\_SimpleBernoulli\_logodds.stan"}\NormalTok{)}
\NormalTok{mod }\OtherTok{\textless{}{-}} \FunctionTok{cmdstan\_model}\NormalTok{(}
\NormalTok{  stan\_file\_logit, }
  \AttributeTok{cpp\_options =} \FunctionTok{list}\NormalTok{(}\AttributeTok{stan\_threads =} \ConstantTok{FALSE}\NormalTok{),}
  \AttributeTok{stanc\_options =} \FunctionTok{list}\NormalTok{(}\StringTok{"O1"}\NormalTok{)}
\NormalTok{)}

\CommentTok{\# {-}{-}{-} 1. Run Standard NUTS (Baseline) {-}{-}{-}}
\ControlFlowTok{if}\NormalTok{ (regenerate\_simulations }\SpecialCharTok{||} \SpecialCharTok{!}\FunctionTok{file.exists}\NormalTok{(file\_nuts\_std)) \{}
  \FunctionTok{cat}\NormalTok{(}\StringTok{"Running Standard NUTS...}\SpecialCharTok{\textbackslash{}n}\StringTok{"}\NormalTok{)}
\NormalTok{  fit\_nuts\_std }\OtherTok{\textless{}{-}}\NormalTok{ mod}\SpecialCharTok{$}\FunctionTok{sample}\NormalTok{(}
    \AttributeTok{data =}\NormalTok{ data\_for\_stan,}
    \AttributeTok{seed =} \DecValTok{123}\NormalTok{,}
    \AttributeTok{chains =} \DecValTok{4}\NormalTok{,}
    \AttributeTok{parallel\_chains =} \DecValTok{4}\NormalTok{,}
    \AttributeTok{iter\_warmup =} \DecValTok{1000}\NormalTok{, }\CommentTok{\# Standard warmup}
    \AttributeTok{iter\_sampling =} \DecValTok{2000}\NormalTok{,}
    \AttributeTok{refresh =} \DecValTok{0}
\NormalTok{  )}
\NormalTok{  fit\_nuts\_std}\SpecialCharTok{$}\FunctionTok{save\_object}\NormalTok{(file\_nuts\_std)}
\NormalTok{\} }\ControlFlowTok{else}\NormalTok{ \{}
\NormalTok{  fit\_nuts\_std }\OtherTok{\textless{}{-}} \FunctionTok{readRDS}\NormalTok{(file\_nuts\_std)}
\NormalTok{\}}

\CommentTok{\# {-}{-}{-} 2. Run Pathfinder {-}{-}{-}}
\ControlFlowTok{if}\NormalTok{ (regenerate\_simulations }\SpecialCharTok{||} \SpecialCharTok{!}\FunctionTok{file.exists}\NormalTok{(file\_pathfinder)) \{}
  \FunctionTok{cat}\NormalTok{(}\StringTok{"Running Pathfinder...}\SpecialCharTok{\textbackslash{}n}\StringTok{"}\NormalTok{)}
\NormalTok{  fit\_pathfinder }\OtherTok{\textless{}{-}}\NormalTok{ mod}\SpecialCharTok{$}\FunctionTok{pathfinder}\NormalTok{(}
    \AttributeTok{data =}\NormalTok{ data\_for\_stan,}
    \AttributeTok{seed =} \DecValTok{123}\NormalTok{,}
    \AttributeTok{draws =} \DecValTok{4000}\NormalTok{, }
    \AttributeTok{num\_paths =} \DecValTok{4}\NormalTok{,}
    \AttributeTok{single\_path\_draws =} \DecValTok{1000}\NormalTok{,}
    \AttributeTok{history\_size =} \DecValTok{50}
\NormalTok{  )}
\NormalTok{  fit\_pathfinder}\SpecialCharTok{$}\FunctionTok{save\_object}\NormalTok{(file\_pathfinder)}
\NormalTok{\} }\ControlFlowTok{else}\NormalTok{ \{}
\NormalTok{  fit\_pathfinder }\OtherTok{\textless{}{-}} \FunctionTok{readRDS}\NormalTok{(file\_pathfinder)}
\NormalTok{\}}

\CommentTok{\# {-}{-}{-} 3. Run NUTS initialized with Pathfinder {-}{-}{-}}
\ControlFlowTok{if}\NormalTok{ (regenerate\_simulations }\SpecialCharTok{||} \SpecialCharTok{!}\FunctionTok{file.exists}\NormalTok{(file\_nuts\_init)) \{}
  \FunctionTok{cat}\NormalTok{(}\StringTok{"Running NUTS initialized with Pathfinder...}\SpecialCharTok{\textbackslash{}n}\StringTok{"}\NormalTok{)}
\NormalTok{  fit\_nuts\_init }\OtherTok{\textless{}{-}}\NormalTok{ mod}\SpecialCharTok{$}\FunctionTok{sample}\NormalTok{(}
    \AttributeTok{data =}\NormalTok{ data\_for\_stan,}
    \AttributeTok{seed =} \DecValTok{123}\NormalTok{,}
    \AttributeTok{chains =} \DecValTok{4}\NormalTok{,}
    \AttributeTok{parallel\_chains =} \DecValTok{4}\NormalTok{,}
    \AttributeTok{iter\_warmup =} \DecValTok{500}\NormalTok{,    }\CommentTok{\# Reduced warmup allowed by better init!}
    \AttributeTok{iter\_sampling =} \DecValTok{2000}\NormalTok{,}
    \AttributeTok{init =}\NormalTok{ fit\_pathfinder, }\CommentTok{\# \textless{}{-}{-}{-} The Magic Line}
    \AttributeTok{refresh =} \DecValTok{0}
\NormalTok{  )}
\NormalTok{  fit\_nuts\_init}\SpecialCharTok{$}\FunctionTok{save\_object}\NormalTok{(file\_nuts\_init)}
\NormalTok{\} }\ControlFlowTok{else}\NormalTok{ \{}
\NormalTok{  fit\_nuts\_init }\OtherTok{\textless{}{-}} \FunctionTok{readRDS}\NormalTok{(file\_nuts\_init)}
\NormalTok{\}}


\CommentTok{\# {-}{-}{-} COMPARISON: TIMING {-}{-}{-}}

\CommentTok{\# Extract execution times (in seconds)}
\CommentTok{\# Note: $time()$total gives the wall clock time for the computation}
\NormalTok{time\_nuts\_std  }\OtherTok{\textless{}{-}}\NormalTok{ fit\_nuts\_std}\SpecialCharTok{$}\FunctionTok{time}\NormalTok{()}\SpecialCharTok{$}\NormalTok{total}
\NormalTok{time\_pathfinder }\OtherTok{\textless{}{-}}\NormalTok{ fit\_pathfinder}\SpecialCharTok{$}\FunctionTok{time}\NormalTok{()}\SpecialCharTok{$}\NormalTok{total}
\CommentTok{\# For the initialized model, the total workflow time is Pathfinder + NUTS}
\NormalTok{time\_nuts\_init }\OtherTok{\textless{}{-}}\NormalTok{ time\_pathfinder }\SpecialCharTok{+}\NormalTok{ fit\_nuts\_init}\SpecialCharTok{$}\FunctionTok{time}\NormalTok{()}\SpecialCharTok{$}\NormalTok{total }

\CommentTok{\# Create a summary table}
\NormalTok{benchmarks }\OtherTok{\textless{}{-}} \FunctionTok{tibble}\NormalTok{(}
  \AttributeTok{Method =} \FunctionTok{c}\NormalTok{(}\StringTok{"Standard NUTS"}\NormalTok{, }\StringTok{"Pathfinder"}\NormalTok{, }\StringTok{"Pathfinder + NUTS"}\NormalTok{),}
  \AttributeTok{Time\_Sec =} \FunctionTok{c}\NormalTok{(time\_nuts\_std, time\_pathfinder, time\_nuts\_init),}
  \AttributeTok{Speedup\_vs\_Std =}\NormalTok{ time\_nuts\_std }\SpecialCharTok{/}\NormalTok{ Time\_Sec}
\NormalTok{)}

\FunctionTok{print}\NormalTok{(benchmarks)}
\end{Highlighting}
\end{Shaded}

\begin{verbatim}
# A tibble: 3 x 3
  Method            Time_Sec Speedup_vs_Std
  <chr>                <dbl>          <dbl>
1 Standard NUTS        0.312          1    
2 Pathfinder           0.270          1.16 
3 Pathfinder + NUTS    0.798          0.391
\end{verbatim}

\begin{Shaded}
\begin{Highlighting}[]
\CommentTok{\# {-}{-}{-} COMPARISON: POSTERIORS {-}{-}{-}}

\CommentTok{\# Extract draws}
\NormalTok{draws\_std  }\OtherTok{\textless{}{-}} \FunctionTok{as\_draws\_df}\NormalTok{(fit\_nuts\_std}\SpecialCharTok{$}\FunctionTok{draws}\NormalTok{(}\StringTok{"theta"}\NormalTok{)) }\SpecialCharTok{\%\textgreater{}\%} \FunctionTok{mutate}\NormalTok{(}\AttributeTok{Method =} \StringTok{"Standard NUTS"}\NormalTok{)}
\NormalTok{draws\_path }\OtherTok{\textless{}{-}} \FunctionTok{as\_draws\_df}\NormalTok{(fit\_pathfinder}\SpecialCharTok{$}\FunctionTok{draws}\NormalTok{(}\StringTok{"theta"}\NormalTok{)) }\SpecialCharTok{\%\textgreater{}\%} \FunctionTok{mutate}\NormalTok{(}\AttributeTok{Method =} \StringTok{"Pathfinder"}\NormalTok{)}
\NormalTok{draws\_init }\OtherTok{\textless{}{-}} \FunctionTok{as\_draws\_df}\NormalTok{(fit\_nuts\_init}\SpecialCharTok{$}\FunctionTok{draws}\NormalTok{(}\StringTok{"theta"}\NormalTok{)) }\SpecialCharTok{\%\textgreater{}\%} \FunctionTok{mutate}\NormalTok{(}\AttributeTok{Method =} \StringTok{"Pathfinder + NUTS"}\NormalTok{)}

\CommentTok{\# Combine}
\NormalTok{all\_draws }\OtherTok{\textless{}{-}} \FunctionTok{bind\_rows}\NormalTok{(draws\_std, draws\_path, draws\_init)}

\CommentTok{\# Plot}
\FunctionTok{ggplot}\NormalTok{(all\_draws, }\FunctionTok{aes}\NormalTok{(}\AttributeTok{x =}\NormalTok{ theta, }\AttributeTok{fill =}\NormalTok{ Method)) }\SpecialCharTok{+}
  \FunctionTok{geom\_density}\NormalTok{(}\AttributeTok{alpha =} \FloatTok{0.4}\NormalTok{) }\SpecialCharTok{+}
  \FunctionTok{geom\_vline}\NormalTok{(}\AttributeTok{xintercept =} \FloatTok{0.8}\NormalTok{, }\AttributeTok{linetype =} \StringTok{"dashed"}\NormalTok{) }\SpecialCharTok{+}
  \FunctionTok{scale\_fill\_viridis\_d}\NormalTok{() }\SpecialCharTok{+}
  \FunctionTok{labs}\NormalTok{(}
    \AttributeTok{title =} \StringTok{"Posterior Comparison"}\NormalTok{,}
    \AttributeTok{subtitle =} \StringTok{"Pathfinder approximates well; Initialized NUTS refines it faster than Standard NUTS"}\NormalTok{,}
    \AttributeTok{x =} \StringTok{"Theta (Rate)"}\NormalTok{,}
    \AttributeTok{y =} \StringTok{"Density"}
\NormalTok{  ) }\SpecialCharTok{+}
  \FunctionTok{theme\_classic}\NormalTok{()}
\end{Highlighting}
\end{Shaded}

\begin{center}
\includegraphics[width=0.8\linewidth,height=\textheight,keepaspectratio]{04-InferringRates_files/figure-pdf/ch05_intro to pathfinded-1.pdf}
\end{center}

\section{Moving Beyond Simple Bias: Memory
Models}\label{sec-memory-models}

The simple biased agent model assumes choices are independent over time,
influenced only by a fixed \texttt{theta}. However, behavior is often
history-dependent. Let's explore models where the choice probability
\texttt{theta} on trial \emph{t} depends on previous events.

\subsection{Memory Model 1: GLM-like Approach (External
Predictor)}\label{memory-model-1-glm-like-approach-external-predictor}

One way to incorporate memory is to treat a summary of past events as an
external predictor influencing the current choice probability, similar
to a predictor in a Generalized Linear Model (GLM).

Let's assume the choice probability \texttt{theta} depends on the
\textbf{cumulative rate of the opponent's `right' choices} observed up
to the \emph{previous} trial.To make the variable more intuitive we code
previous rate - which is bound to a probability 0-1 space - into
log-odds via a logit link/transformation. In this way a previous rate
with more left than right choices will result in a negative value,
thereby decreasing our propensity to choose right; and one with more
right than left choices will result in a positive value, thereby
increasing our propensity to choose right.

\begin{Shaded}
\begin{Highlighting}[]
\CommentTok{\# We subset to only include no noise and a specific rate}
\NormalTok{d1 }\OtherTok{\textless{}{-}}\NormalTok{ d }\SpecialCharTok{\%\textgreater{}\%} 
  \FunctionTok{subset}\NormalTok{(true\_noise }\SpecialCharTok{==} \DecValTok{0} \SpecialCharTok{\&}\NormalTok{ true\_rate }\SpecialCharTok{==} \FloatTok{0.8}\NormalTok{) }\SpecialCharTok{\%\textgreater{}\%} 
  \FunctionTok{rename}\NormalTok{(}\AttributeTok{Other =}\NormalTok{ choice) }\SpecialCharTok{\%\textgreater{}\%} 
  \FunctionTok{mutate}\NormalTok{(}\AttributeTok{cumulativerate =} \FunctionTok{lag}\NormalTok{(cumulative\_rate, }\DecValTok{1}\NormalTok{))}

\NormalTok{d1}\SpecialCharTok{$}\NormalTok{cumulativerate[}\DecValTok{1}\NormalTok{] }\OtherTok{\textless{}{-}} \FloatTok{0.5} \CommentTok{\# no prior info at first trial}
\NormalTok{d1}\SpecialCharTok{$}\NormalTok{cumulativerate[d1}\SpecialCharTok{$}\NormalTok{cumulativerate }\SpecialCharTok{==} \DecValTok{0}\NormalTok{] }\OtherTok{\textless{}{-}} \FloatTok{0.01}
\NormalTok{d1}\SpecialCharTok{$}\NormalTok{cumulativerate[d1}\SpecialCharTok{$}\NormalTok{cumulativerate }\SpecialCharTok{==} \DecValTok{1}\NormalTok{] }\OtherTok{\textless{}{-}} \FloatTok{0.99}

\CommentTok{\# Now we create the memory agent with a coefficient of 1 (in log odds)}
\NormalTok{MemoryAgent\_f }\OtherTok{\textless{}{-}} \ControlFlowTok{function}\NormalTok{(bias, beta, cumulativerate)\{}
\NormalTok{    choice }\OtherTok{=} \FunctionTok{rbinom}\NormalTok{(}\DecValTok{1}\NormalTok{, }\DecValTok{1}\NormalTok{, }\FunctionTok{plogis}\NormalTok{(bias }\SpecialCharTok{+}\NormalTok{ beta }\SpecialCharTok{*} \FunctionTok{qlogis}\NormalTok{(cumulativerate)))}
  \FunctionTok{return}\NormalTok{(choice)}
\NormalTok{\}}

\NormalTok{d1}\SpecialCharTok{$}\NormalTok{Self[}\DecValTok{1}\NormalTok{] }\OtherTok{\textless{}{-}} \FunctionTok{RandomAgentNoise\_f}\NormalTok{(}\DecValTok{0}\NormalTok{, }\DecValTok{0}\NormalTok{) }\CommentTok{\# 0.5 in the transformed space}

\ControlFlowTok{for}\NormalTok{ (i }\ControlFlowTok{in} \DecValTok{2}\SpecialCharTok{:}\NormalTok{trials) \{}
\NormalTok{  d1}\SpecialCharTok{$}\NormalTok{Self[i] }\OtherTok{\textless{}{-}} \FunctionTok{MemoryAgent\_f}\NormalTok{(}\AttributeTok{bias =} \DecValTok{0}\NormalTok{, }\AttributeTok{beta =} \DecValTok{1}\NormalTok{, d1}\SpecialCharTok{$}\NormalTok{cumulativerate[i])}
\NormalTok{\}}

\DocumentationTok{\#\# Create the data}
\NormalTok{data\_memory }\OtherTok{\textless{}{-}} \FunctionTok{list}\NormalTok{(}
  \AttributeTok{n =} \DecValTok{120}\NormalTok{,}
  \AttributeTok{h =}\NormalTok{ d1}\SpecialCharTok{$}\NormalTok{Self,}
  \AttributeTok{memory =}\NormalTok{ d1}\SpecialCharTok{$}\NormalTok{cumulativerate }\CommentTok{\# this creates the new parameter: the rate of right hands so far in log{-}odds}
\NormalTok{)}
\end{Highlighting}
\end{Shaded}

\begin{Shaded}
\begin{Highlighting}[]
\NormalTok{stan\_model }\OtherTok{\textless{}{-}} \StringTok{"}
\StringTok{// The input (data) for the model. n of trials and h for (right and left) hand}
\StringTok{data \{}
\StringTok{ int\textless{}lower=1\textgreater{} n;}
\StringTok{ array[n] int h;}
\StringTok{ vector[n] memory; // here we add the new variable between 0.01 and .99}
\StringTok{\}}

\StringTok{// The parameters accepted by the model. }
\StringTok{parameters \{}
\StringTok{  real bias; // how likely is the agent to pick right when the previous rate has no information (50{-}50)?}
\StringTok{  real beta; // how strongly is previous rate impacting the decision?}
\StringTok{\}}

\StringTok{// The model to be estimated. }
\StringTok{model \{}
\StringTok{  // priors}
\StringTok{  target += normal\_lpdf(bias | 0, .3);}
\StringTok{  target += normal\_lpdf(beta | 0, .5);}
\StringTok{  }
\StringTok{  // model}
\StringTok{  target += bernoulli\_logit\_lpmf(h | bias + beta * logit(memory));}
\StringTok{\}}

\StringTok{generated quantities \{}
\StringTok{  // Generate prior samples for Prior{-}Posterior Update checks}
\StringTok{  real bias\_prior = normal\_rng(0, 0.3);}
\StringTok{  real beta\_prior = normal\_rng(0, 0.5);}
\StringTok{\}}
\StringTok{"}
\FunctionTok{write\_stan\_file}\NormalTok{(}
\NormalTok{  stan\_model,}
  \AttributeTok{dir =} \StringTok{"stan/"}\NormalTok{,}
  \AttributeTok{basename =} \StringTok{"ch06\_MemoryBernoulli.stan"}\NormalTok{)}
\end{Highlighting}
\end{Shaded}

\begin{verbatim}
[1] "/Users/au209589/Dropbox/Teaching/AdvancedCognitiveModeling23_book/stan/ch06_MemoryBernoulli.stan"
\end{verbatim}

\begin{Shaded}
\begin{Highlighting}[]
\DocumentationTok{\#\# Specify where the model is}
\NormalTok{file }\OtherTok{\textless{}{-}} \FunctionTok{here}\NormalTok{(}\StringTok{"stan"}\NormalTok{,}\StringTok{"ch06\_MemoryBernoulli.stan"}\NormalTok{)}

\CommentTok{\# File path for saved model}
\NormalTok{model\_file\_memory }\OtherTok{\textless{}{-}} \FunctionTok{here}\NormalTok{(}\StringTok{"simmodels"}\NormalTok{, }\StringTok{"ch06\_MemoryBernoulli.rds"}\NormalTok{)}

\CommentTok{\# Check if we need to rerun the simulation}
\ControlFlowTok{if}\NormalTok{ (regenerate\_simulations }\SpecialCharTok{||} \SpecialCharTok{!}\FunctionTok{file.exists}\NormalTok{(model\_file\_memory)) \{}
  \CommentTok{\# Compile the model}
\NormalTok{  mod\_memory }\OtherTok{\textless{}{-}} \FunctionTok{cmdstan\_model}\NormalTok{(file)}
  
  \CommentTok{\# The following command calls Stan with specific options.}
\NormalTok{  samples\_memory }\OtherTok{\textless{}{-}}\NormalTok{ mod\_memory}\SpecialCharTok{$}\FunctionTok{sample}\NormalTok{(}
    \AttributeTok{data =}\NormalTok{ data\_memory,}
    \AttributeTok{seed =} \DecValTok{123}\NormalTok{,}
    \AttributeTok{chains =} \DecValTok{2}\NormalTok{,}
    \AttributeTok{parallel\_chains =} \DecValTok{2}\NormalTok{,}
    \AttributeTok{threads\_per\_chain =} \DecValTok{1}\NormalTok{,}
    \AttributeTok{iter\_warmup =} \DecValTok{1000}\NormalTok{,}
    \AttributeTok{iter\_sampling =} \DecValTok{1000}\NormalTok{,}
    \AttributeTok{refresh =} \DecValTok{0}\NormalTok{,}
    \AttributeTok{max\_treedepth =} \DecValTok{20}\NormalTok{,}
    \AttributeTok{adapt\_delta =} \FloatTok{0.99}\NormalTok{,}
\NormalTok{  )}
  
  \CommentTok{\# Save the fitted model}
\NormalTok{  samples\_memory}\SpecialCharTok{$}\FunctionTok{save\_object}\NormalTok{(}\AttributeTok{file =}\NormalTok{ model\_file\_memory)}
  \FunctionTok{cat}\NormalTok{(}\StringTok{"Generated new model fit and saved to"}\NormalTok{, model\_file\_memory, }\StringTok{"}\SpecialCharTok{\textbackslash{}n}\StringTok{"}\NormalTok{)}
\NormalTok{\} }\ControlFlowTok{else}\NormalTok{ \{}
  \CommentTok{\# Load existing results}
  \FunctionTok{cat}\NormalTok{(}\StringTok{"Loading memory model samples...}\SpecialCharTok{\textbackslash{}n}\StringTok{"}\NormalTok{)}
\NormalTok{  samples\_memory }\OtherTok{\textless{}{-}} \FunctionTok{readRDS}\NormalTok{(model\_file\_memory)}
  \FunctionTok{cat}\NormalTok{(}\StringTok{"Available parameters:"}\NormalTok{, }\FunctionTok{paste}\NormalTok{(}\FunctionTok{colnames}\NormalTok{(}\FunctionTok{as\_draws\_df}\NormalTok{(samples\_memory}\SpecialCharTok{$}\FunctionTok{draws}\NormalTok{())), }\AttributeTok{collapse=}\StringTok{", "}\NormalTok{), }\StringTok{"}\SpecialCharTok{\textbackslash{}n}\StringTok{"}\NormalTok{)}
  \FunctionTok{cat}\NormalTok{(}\StringTok{"Loaded existing model fit from"}\NormalTok{, model\_file\_memory, }\StringTok{"}\SpecialCharTok{\textbackslash{}n}\StringTok{"}\NormalTok{)}
\NormalTok{\}}
\end{Highlighting}
\end{Shaded}

\begin{verbatim}
Loading memory model samples...
Available parameters: lp__, bias, beta, bias_prior, beta_prior, .chain, .iteration, .draw 
Loaded existing model fit from /Users/au209589/Dropbox/Teaching/AdvancedCognitiveModeling23_book/simmodels/ch06_MemoryBernoulli.rds 
\end{verbatim}

\subsection{Summarizing the results}\label{summarizing-the-results-1}

\begin{Shaded}
\begin{Highlighting}[]
\CommentTok{\# Check if samples\_memory exists}
\ControlFlowTok{if}\NormalTok{ (}\SpecialCharTok{!}\FunctionTok{exists}\NormalTok{(}\StringTok{"samples\_memory"}\NormalTok{)) \{}
  \FunctionTok{cat}\NormalTok{(}\StringTok{"Loading memory model samples...}\SpecialCharTok{\textbackslash{}n}\StringTok{"}\NormalTok{)}
\NormalTok{  samples\_memory }\OtherTok{\textless{}{-}} \FunctionTok{readRDS}\NormalTok{(}\FunctionTok{here}\NormalTok{(}\StringTok{"simmodels"}\NormalTok{, }\StringTok{"ch06\_MemoryBernoulli.rds"}\NormalTok{))}
  \FunctionTok{cat}\NormalTok{(}\StringTok{"Available parameters:"}\NormalTok{, }\FunctionTok{paste}\NormalTok{(}\FunctionTok{colnames}\NormalTok{(}\FunctionTok{as\_draws\_df}\NormalTok{(samples\_memory}\SpecialCharTok{$}\FunctionTok{draws}\NormalTok{())), }\AttributeTok{collapse =} \StringTok{", "}\NormalTok{), }\StringTok{"}\SpecialCharTok{\textbackslash{}n}\StringTok{"}\NormalTok{)}
\NormalTok{\}}

\CommentTok{\# Extract posterior samples and include sampling of the prior:}
\NormalTok{draws\_df\_memory }\OtherTok{\textless{}{-}} \FunctionTok{as\_draws\_df}\NormalTok{(samples\_memory}\SpecialCharTok{$}\FunctionTok{draws}\NormalTok{()) }

\CommentTok{\# Explicitly extract parameters}
\NormalTok{bias\_param }\OtherTok{\textless{}{-}}\NormalTok{ draws\_df\_memory}\SpecialCharTok{$}\NormalTok{bias}
\NormalTok{beta\_param }\OtherTok{\textless{}{-}}\NormalTok{ draws\_df\_memory}\SpecialCharTok{$}\NormalTok{beta}
\FunctionTok{cat}\NormalTok{(}\StringTok{"Successfully extracted"}\NormalTok{, }\FunctionTok{length}\NormalTok{(bias\_param), }\StringTok{"values for bias parameter}\SpecialCharTok{\textbackslash{}n}\StringTok{"}\NormalTok{)}
\end{Highlighting}
\end{Shaded}

\begin{verbatim}
Successfully extracted 2000 values for bias parameter
\end{verbatim}

\begin{Shaded}
\begin{Highlighting}[]
\FunctionTok{cat}\NormalTok{(}\StringTok{"Successfully extracted"}\NormalTok{, }\FunctionTok{length}\NormalTok{(beta\_param), }\StringTok{"values for beta parameter}\SpecialCharTok{\textbackslash{}n}\StringTok{"}\NormalTok{)}
\end{Highlighting}
\end{Shaded}

\begin{verbatim}
Successfully extracted 2000 values for beta parameter
\end{verbatim}

\begin{Shaded}
\begin{Highlighting}[]
\CommentTok{\# Trace plot for bias}
\FunctionTok{ggplot}\NormalTok{(draws\_df\_memory, }\FunctionTok{aes}\NormalTok{(.iteration, bias, }\AttributeTok{group =}\NormalTok{ .chain, }\AttributeTok{color =}\NormalTok{ .chain)) }\SpecialCharTok{+}
  \FunctionTok{geom\_line}\NormalTok{() }\SpecialCharTok{+}
  \FunctionTok{labs}\NormalTok{(}\AttributeTok{title =} \StringTok{"Trace plot for bias parameter"}\NormalTok{) }\SpecialCharTok{+}
  \FunctionTok{theme\_classic}\NormalTok{()}
\end{Highlighting}
\end{Shaded}

\begin{center}
\includegraphics[width=0.8\linewidth,height=\textheight,keepaspectratio]{04-InferringRates_files/figure-pdf/04 memory model log-odds quality assessment-1.pdf}
\end{center}

\begin{Shaded}
\begin{Highlighting}[]
\CommentTok{\# Trace plot for beta}
\FunctionTok{ggplot}\NormalTok{(draws\_df\_memory, }\FunctionTok{aes}\NormalTok{(.iteration, beta, }\AttributeTok{group =}\NormalTok{ .chain, }\AttributeTok{color =}\NormalTok{ .chain)) }\SpecialCharTok{+}
  \FunctionTok{geom\_line}\NormalTok{() }\SpecialCharTok{+}
  \FunctionTok{labs}\NormalTok{(}\AttributeTok{title =} \StringTok{"Trace plot for beta parameter"}\NormalTok{) }\SpecialCharTok{+}
  \FunctionTok{theme\_classic}\NormalTok{()}
\end{Highlighting}
\end{Shaded}

\begin{center}
\includegraphics[width=0.8\linewidth,height=\textheight,keepaspectratio]{04-InferringRates_files/figure-pdf/04 memory model log-odds quality assessment-2.pdf}
\end{center}

\begin{Shaded}
\begin{Highlighting}[]
\CommentTok{\# Now let\textquotesingle{}s plot the density for bias (prior and posterior)}
\FunctionTok{ggplot}\NormalTok{(draws\_df\_memory) }\SpecialCharTok{+}
  \FunctionTok{geom\_density}\NormalTok{(}\FunctionTok{aes}\NormalTok{(bias), }\AttributeTok{fill =} \StringTok{"blue"}\NormalTok{, }\AttributeTok{alpha =} \FloatTok{0.3}\NormalTok{) }\SpecialCharTok{+}
  \FunctionTok{geom\_density}\NormalTok{(}\FunctionTok{aes}\NormalTok{(bias\_prior), }\AttributeTok{fill =} \StringTok{"red"}\NormalTok{, }\AttributeTok{alpha =} \FloatTok{0.3}\NormalTok{) }\SpecialCharTok{+}
  \FunctionTok{geom\_vline}\NormalTok{(}\AttributeTok{xintercept =} \DecValTok{0}\NormalTok{, }\AttributeTok{linetype =} \StringTok{"dashed"}\NormalTok{, }\AttributeTok{color =} \StringTok{"black"}\NormalTok{, }\AttributeTok{size =} \FloatTok{1.5}\NormalTok{) }\SpecialCharTok{+}
  \FunctionTok{labs}\NormalTok{(}\AttributeTok{title =} \StringTok{"Prior{-}Posterior Update for Bias Parameter"}\NormalTok{,}
       \AttributeTok{subtitle =} \StringTok{"Blue: posterior, Red: prior, Dashed: true value"}\NormalTok{) }\SpecialCharTok{+}
  \FunctionTok{xlab}\NormalTok{(}\StringTok{"Bias"}\NormalTok{) }\SpecialCharTok{+}
  \FunctionTok{ylab}\NormalTok{(}\StringTok{"Posterior Density"}\NormalTok{) }\SpecialCharTok{+}
  \FunctionTok{theme\_classic}\NormalTok{()}
\end{Highlighting}
\end{Shaded}

\begin{center}
\includegraphics[width=0.8\linewidth,height=\textheight,keepaspectratio]{04-InferringRates_files/figure-pdf/04 memory model log-odds quality assessment-3.pdf}
\end{center}

\begin{Shaded}
\begin{Highlighting}[]
\CommentTok{\# Now let\textquotesingle{}s plot the density for beta (prior and posterior)}
\FunctionTok{ggplot}\NormalTok{(draws\_df\_memory) }\SpecialCharTok{+}
  \FunctionTok{geom\_density}\NormalTok{(}\FunctionTok{aes}\NormalTok{(beta), }\AttributeTok{fill =} \StringTok{"blue"}\NormalTok{, }\AttributeTok{alpha =} \FloatTok{0.3}\NormalTok{) }\SpecialCharTok{+}
  \FunctionTok{geom\_density}\NormalTok{(}\FunctionTok{aes}\NormalTok{(beta\_prior), }\AttributeTok{fill =} \StringTok{"red"}\NormalTok{, }\AttributeTok{alpha =} \FloatTok{0.3}\NormalTok{) }\SpecialCharTok{+}
  \FunctionTok{geom\_vline}\NormalTok{(}\AttributeTok{xintercept =} \DecValTok{1}\NormalTok{, }\AttributeTok{linetype =} \StringTok{"dashed"}\NormalTok{, }\AttributeTok{color =} \StringTok{"black"}\NormalTok{, }\AttributeTok{size =} \FloatTok{1.5}\NormalTok{) }\SpecialCharTok{+}
  \FunctionTok{labs}\NormalTok{(}\AttributeTok{title =} \StringTok{"Prior{-}Posterior Update for Beta Parameter"}\NormalTok{,}
       \AttributeTok{subtitle =} \StringTok{"Blue: posterior, Red: prior, Dashed: true value"}\NormalTok{) }\SpecialCharTok{+}
  \FunctionTok{xlab}\NormalTok{(}\StringTok{"Beta"}\NormalTok{) }\SpecialCharTok{+}
  \FunctionTok{ylab}\NormalTok{(}\StringTok{"Posterior Density"}\NormalTok{) }\SpecialCharTok{+}
  \FunctionTok{theme\_classic}\NormalTok{()}
\end{Highlighting}
\end{Shaded}

\begin{center}
\includegraphics[width=0.8\linewidth,height=\textheight,keepaspectratio]{04-InferringRates_files/figure-pdf/04 memory model log-odds quality assessment-4.pdf}
\end{center}

\begin{Shaded}
\begin{Highlighting}[]
\CommentTok{\# Print summary}
\NormalTok{samples\_memory}\SpecialCharTok{$}\FunctionTok{summary}\NormalTok{()}
\end{Highlighting}
\end{Shaded}

\begin{verbatim}
# A tibble: 5 x 10
  variable     mean   median    sd   mad      q5     q95  rhat ess_bulk ess_tail
  <chr>       <dbl>    <dbl> <dbl> <dbl>   <dbl>   <dbl> <dbl>    <dbl>    <dbl>
1 lp__     -70.3    -70.0    0.971 0.736 -72.2   -69.4   1.000     644.     733.
2 bias       0.304    0.311  0.252 0.244  -0.120   0.717 1.01      383.     555.
3 beta       0.518    0.513  0.222 0.225   0.160   0.882 1.01      427.     474.
4 bias_pr~  -0.0138  -0.0106 0.300 0.292  -0.509   0.466 1.000    2201.    1758.
5 beta_pr~   0.0302   0.0395 0.491 0.496  -0.770   0.814 1.000    1861.    1940.
\end{verbatim}

The trace plots should show the same ``hairy caterpillar'' pattern as
before --- well-mixed chains with no systematic drift. The
prior-posterior update plots tell the story: the posterior for
\texttt{bias} should be narrow and centred near 0 (log-odds of 0.5 =
unbiased), and the posterior for \texttt{beta} should be centred near 1
(the true generating value). The data have meaningfully updated both
parameters away from their priors. Formal quality checks --- prior
predictive checks, systematic parameter recovery, and SBC --- come in
Ch. 6.

\subsection{Memory Model 2: Internal State
Variable}\label{memory-model-2-internal-state-variable}

Instead of feeding memory as external data, we can model it as an
internal state that updates within the model.

For this specific model --- where memory is simply a running average of
the opponent's choices --- the update rule depends only on the observed
data (\texttt{other}), not on any parameter being estimated. We should
therefore calculate it in Stan's \texttt{transformed\ data} block.

\begin{tcolorbox}[enhanced jigsaw, toprule=.15mm, toptitle=1mm, bottomtitle=1mm, leftrule=.75mm, rightrule=.15mm, colframe=quarto-callout-tip-color-frame, left=2mm, colback=white, titlerule=0mm, coltitle=black, colbacktitle=quarto-callout-tip-color!10!white, opacityback=0, title=\textcolor{quarto-callout-tip-color}{\faLightbulb}\hspace{0.5em}{Stan Efficiency Rule 1 --- Block Placement: Run it Once, Not Thousands
of Times}, arc=.35mm, opacitybacktitle=0.6, breakable, bottomrule=.15mm]

Every Stan model has blocks that execute at very different frequencies.
Choosing the \textbf{right block} is the most impactful single
optimization:

\begin{longtable}[]{@{}
  >{\raggedright\arraybackslash}p{(\linewidth - 4\tabcolsep) * \real{0.3333}}
  >{\raggedright\arraybackslash}p{(\linewidth - 4\tabcolsep) * \real{0.3333}}
  >{\raggedright\arraybackslash}p{(\linewidth - 4\tabcolsep) * \real{0.3333}}@{}}
\toprule\noalign{}
\begin{minipage}[b]{\linewidth}\raggedright
Block
\end{minipage} & \begin{minipage}[b]{\linewidth}\raggedright
Executes
\end{minipage} & \begin{minipage}[b]{\linewidth}\raggedright
Use for
\end{minipage} \\
\midrule\noalign{}
\endhead
\bottomrule\noalign{}
\endlastfoot
\texttt{transformed\ data} & \textbf{Once}, before sampling & Any
quantity derived \textbf{only from data} --- precomputed distances,
normalized inputs, log-transforms of fixed constants \\
\texttt{transformed\ parameters} & \textbf{Every leapfrog step} &
Deterministic functions of parameters whose gradients must flow backward
--- recursive states, reparameterizations, predicted values used in the
likelihood \\
\texttt{model} (local \texttt{\{\}} block) & Every leapfrog step,
\textbf{not saved} & Scratch variables for accumulating
\texttt{target\ +=}; no overhead from the saved-draw tape \\
\texttt{generated\ quantities} & \textbf{Once per saved draw}
(post-warmup only) & PPCs, log-likelihoods, prior samples --- expensive
quantities you only need in the final posterior \\
\end{longtable}

\textbf{Concrete cost of misplacement.} A typical fit runs 4 chains ×
1000 warmup × \textasciitilde6 leapfrog steps per draw =
\textasciitilde24 000 leapfrog evaluations. Any computation in
\texttt{transformed\ parameters} runs \textasciitilde24 000 times. The
same computation in \texttt{transformed\ data} runs exactly once. For a
model with 500 trials, precomputing 500 memory values in
\texttt{transformed\ data} rather than \texttt{transformed\ parameters}
saves \textasciitilde12 million floating-point operations per fit.

\textbf{The decision tree:}

\begin{enumerate}
\def\labelenumi{\arabic{enumi}.}
\tightlist
\item
  Does this depend \textbf{only on data} (not on any parameter)? →
  \texttt{transformed\ data}
\item
  Does this depend on \textbf{parameters} and does the likelihood or
  another parameter need it? → \texttt{transformed\ parameters}
\item
  Is this needed \textbf{only after sampling} (PPC draws, LOO
  log-likelihoods)? → \texttt{generated\ quantities}, gated by
  \texttt{run\_diagnostics} (introduced properly in Ch. 5)
\item
  Is this just a temporary accumulator inside a loop? → Local variable
  inside \texttt{\{\}} in \texttt{model} or
  \texttt{transformed\ parameters}
\end{enumerate}

In the memory model below, \texttt{memory} depends only on
\texttt{other} (data) → \texttt{transformed\ data}. In the forgetting
model that follows, \texttt{memory} depends on the free parameter
\texttt{forgetting} → local variable inside \texttt{model}.

\textbf{Corollary --- pre-compute in R when possible.} If a quantity
depends only on data and will never change between fits, computing it in
R and passing it as a data field is even faster than
\texttt{transformed\ data}: Stan never sees the raw computation at all.
For example, the multilevel memory model in Chapter 6 pre-computes the
running opponent rate in R and passes \texttt{opp\_rate\_prev} as a
vector, rather than recomputing it inside Stan on every chain.

\end{tcolorbox}

\begin{tcolorbox}[enhanced jigsaw, toprule=.15mm, toptitle=1mm, bottomtitle=1mm, leftrule=.75mm, rightrule=.15mm, colframe=quarto-callout-tip-color-frame, left=2mm, colback=white, titlerule=0mm, coltitle=black, colbacktitle=quarto-callout-tip-color!10!white, opacityback=0, title=\textcolor{quarto-callout-tip-color}{\faLightbulb}\hspace{0.5em}{Stan Efficiency --- Vectorize the Likelihood Call}, arc=.35mm, opacitybacktitle=0.6, breakable, bottomrule=.15mm]

Stan evaluates each \texttt{\textasciitilde{}} or \texttt{target\ +=}
statement as a separate autodiff node. Replacing a trial-level loop with
a \textbf{vectorized distribution call} collapses \(N\) nodes into one,
reducing autodiff graph size and improving cache locality.

\begin{Shaded}
\begin{Highlighting}[]
\CommentTok{// Slower: N separate nodes in the autodiff graph}
\ControlFlowTok{for}\NormalTok{ (i }\ControlFlowTok{in} \DecValTok{1}\NormalTok{:N)}
  \KeywordTok{target +=}\NormalTok{ bernoulli\_logit\_lpmf(h[i] | theta\_logit);}

\CommentTok{// Faster: one vectorized call — h and theta\_logit are both length{-}N vectors}
\KeywordTok{target +=}\NormalTok{ bernoulli\_logit\_lpmf(h | theta\_logit);}
\end{Highlighting}
\end{Shaded}

The vectorized form works whenever \texttt{theta\_logit} is either a
\textbf{scalar} (broadcasts to all trials) or a \textbf{vector of the
same length as \texttt{h}}. In the multilevel model of Chapter 6,
\texttt{target\ +=\ bernoulli\_logit\_lpmf(y\ \textbar{}\ theta\_logit{[}subj\_id{]})}
vectorizes the full N\_total trial loop by scatter-indexing the
per-subject parameters.

\textbf{Caveat:} vectorization is not available for mixture likelihoods
(\texttt{log\_mix}) because each observation requires its own separate
\texttt{log\_mix} call. The loop is unavoidable there (Ch. 8).

\end{tcolorbox}

\begin{Shaded}
\begin{Highlighting}[]
\DocumentationTok{\#\# Create the data — run\_diagnostics = 1 for the single validation fit}
\NormalTok{data }\OtherTok{\textless{}{-}} \FunctionTok{list}\NormalTok{(}
  \AttributeTok{n =} \DecValTok{120}\NormalTok{,}
  \AttributeTok{h =}\NormalTok{ d1}\SpecialCharTok{$}\NormalTok{Self,}
  \AttributeTok{other =}\NormalTok{ d1}\SpecialCharTok{$}\NormalTok{Other,}
  \AttributeTok{run\_diagnostics =} \DecValTok{1}\NormalTok{L}
\NormalTok{)}

\NormalTok{stan\_model }\OtherTok{\textless{}{-}} \StringTok{"}
\StringTok{// Memory{-}based choice model with prior and posterior predictions.}
\StringTok{// memory[] is data{-}derived and lives in transformed data (computed once).}
\StringTok{// run\_diagnostics gates expensive generated quantities for SBC speed.}

\StringTok{data \{}
\StringTok{ int\textless{}lower=1\textgreater{} n;}
\StringTok{ array[n] int h;}
\StringTok{ array[n] int other;}
\StringTok{ int\textless{}lower=0, upper=1\textgreater{} run\_diagnostics; // 1 = compute PPC + log\_lik; 0 = skip (use for SBC)}
\StringTok{\}}

\StringTok{transformed data \{}
\StringTok{  // memory depends only on other[] (data), so it runs ONCE here.}
\StringTok{  // Moving this to transformed parameters would recompute it \textasciitilde{}24 000 times per fit.}
\StringTok{  vector[n] memory;}
\StringTok{  memory[1] = 0.5;}
\StringTok{  for (trial in 1:(n{-}1)) \{}
\StringTok{     memory[trial + 1] = memory[trial] + ((other[trial] {-} memory[trial]) / (trial + 1));}
\StringTok{     if (memory[trial + 1] \textless{} 0.01) memory[trial + 1] = 0.01;}
\StringTok{     if (memory[trial + 1] \textgreater{} 0.99) memory[trial + 1] = 0.99;}
\StringTok{  \}}
\StringTok{\}}

\StringTok{parameters \{}
\StringTok{  real bias;}
\StringTok{  real beta;}
\StringTok{\}}

\StringTok{model \{}
\StringTok{  target += normal\_lpdf(bias | 0, .3);}
\StringTok{  target += normal\_lpdf(beta | 0, .5);}
\StringTok{  // Vectorized likelihood — no loop needed because memory is precomputed.}
\StringTok{  target += bernoulli\_logit\_lpmf(h | bias + beta * logit(memory));}
\StringTok{\}}

\StringTok{generated quantities \{}
\StringTok{  // Prior samples are always generated (cheap).}
\StringTok{  real bias\_prior = normal\_rng(0, 0.3);}
\StringTok{  real beta\_prior = normal\_rng(0, 0.5);}
\StringTok{  // Expensive diagnostics are gated — set run\_diagnostics = 0 during SBC.}
\StringTok{  vector[n] log\_lik;}
\StringTok{  array[n] int y\_rep;}
\StringTok{  if (run\_diagnostics) \{}
\StringTok{    for (trial in 1:n) \{}
\StringTok{      log\_lik[trial] = bernoulli\_logit\_lpmf(h[trial] | bias + beta * logit(memory[trial]));}
\StringTok{      y\_rep[trial]   = bernoulli\_logit\_rng(bias + beta * logit(memory[trial]));}
\StringTok{    \}}
\StringTok{  \}}
\StringTok{\}}
\StringTok{"}

\FunctionTok{write\_stan\_file}\NormalTok{(}
\NormalTok{  stan\_model,}
  \AttributeTok{dir =} \StringTok{"stan/"}\NormalTok{,}
  \AttributeTok{basename =} \StringTok{"ch06\_InternalMemory.stan"}\NormalTok{)}
\end{Highlighting}
\end{Shaded}

\begin{verbatim}
[1] "/Users/au209589/Dropbox/Teaching/AdvancedCognitiveModeling23_book/stan/ch06_InternalMemory.stan"
\end{verbatim}

\begin{Shaded}
\begin{Highlighting}[]
\DocumentationTok{\#\# Specify where the model is}
\NormalTok{file }\OtherTok{\textless{}{-}} \FunctionTok{here}\NormalTok{(}\StringTok{"stan"}\NormalTok{, }\StringTok{"ch06\_InternalMemory.stan"}\NormalTok{)}

\CommentTok{\# File path for saved model}
\NormalTok{model\_file }\OtherTok{\textless{}{-}} \FunctionTok{here}\NormalTok{(}\StringTok{"simmodels"}\NormalTok{, }\StringTok{"ch06\_InternalMemory.rds"}\NormalTok{)}

\CommentTok{\# Check if we need to rerun the simulation}
\ControlFlowTok{if}\NormalTok{ (regenerate\_simulations }\SpecialCharTok{||} \SpecialCharTok{!}\FunctionTok{file.exists}\NormalTok{(model\_file)) \{}
  \CommentTok{\# Compile the model}
\NormalTok{  mod }\OtherTok{\textless{}{-}} \FunctionTok{cmdstan\_model}\NormalTok{(file, }
                       \AttributeTok{cpp\_options =} \FunctionTok{list}\NormalTok{(}\AttributeTok{stan\_threads =} \ConstantTok{FALSE}\NormalTok{),}
                       \AttributeTok{stanc\_options =} \FunctionTok{list}\NormalTok{(}\StringTok{"O1"}\NormalTok{))}
  
  \CommentTok{\# The following command calls Stan with specific options.}
\NormalTok{  samples\_memory\_internal }\OtherTok{\textless{}{-}}\NormalTok{ mod}\SpecialCharTok{$}\FunctionTok{sample}\NormalTok{(}
    \AttributeTok{data =}\NormalTok{ data,}
    \AttributeTok{seed =} \DecValTok{123}\NormalTok{,}
    \AttributeTok{chains =} \DecValTok{1}\NormalTok{,}
    \AttributeTok{parallel\_chains =} \DecValTok{2}\NormalTok{,}
    \AttributeTok{threads\_per\_chain =} \DecValTok{1}\NormalTok{,}
    \AttributeTok{iter\_warmup =} \DecValTok{1000}\NormalTok{,}
    \AttributeTok{iter\_sampling =} \DecValTok{1000}\NormalTok{,}
    \AttributeTok{refresh =} \DecValTok{0}\NormalTok{,}
    \AttributeTok{max\_treedepth =} \DecValTok{20}\NormalTok{,}
    \AttributeTok{adapt\_delta =} \FloatTok{0.99}\NormalTok{,}
\NormalTok{  )}
  
  \CommentTok{\# Save the fitted model}
\NormalTok{  samples\_memory\_internal}\SpecialCharTok{$}\FunctionTok{save\_object}\NormalTok{(}\AttributeTok{file =}\NormalTok{ model\_file)}
  \FunctionTok{cat}\NormalTok{(}\StringTok{"Generated new model fit and saved to"}\NormalTok{, model\_file, }\StringTok{"}\SpecialCharTok{\textbackslash{}n}\StringTok{"}\NormalTok{)}
\NormalTok{\} }\ControlFlowTok{else}\NormalTok{ \{}
  \CommentTok{\# Load existing results}
\NormalTok{  samples\_memory\_internal }\OtherTok{\textless{}{-}} \FunctionTok{readRDS}\NormalTok{(model\_file)}
  \FunctionTok{cat}\NormalTok{(}\StringTok{"Loaded existing model fit from"}\NormalTok{, model\_file, }\StringTok{"}\SpecialCharTok{\textbackslash{}n}\StringTok{"}\NormalTok{)}
\NormalTok{\}}
\end{Highlighting}
\end{Shaded}

\begin{verbatim}
Loaded existing model fit from /Users/au209589/Dropbox/Teaching/AdvancedCognitiveModeling23_book/simmodels/ch06_InternalMemory.rds 
\end{verbatim}

\begin{Shaded}
\begin{Highlighting}[]
\NormalTok{draws\_df }\OtherTok{\textless{}{-}} \FunctionTok{as\_draws\_df}\NormalTok{(samples\_memory\_internal}\SpecialCharTok{$}\FunctionTok{draws}\NormalTok{())}

\CommentTok{\# 1. Check chain convergence}
\CommentTok{\# Plot traces for main parameters}
\FunctionTok{mcmc\_trace}\NormalTok{(draws\_df, }\AttributeTok{pars =} \FunctionTok{c}\NormalTok{(}\StringTok{"bias"}\NormalTok{, }\StringTok{"beta"}\NormalTok{)) }\SpecialCharTok{+}
  \FunctionTok{theme\_minimal}\NormalTok{() }\SpecialCharTok{+}
  \FunctionTok{ggtitle}\NormalTok{(}\StringTok{"Parameter Traces Across Chains"}\NormalTok{)}
\end{Highlighting}
\end{Shaded}

\begin{center}
\includegraphics[width=0.8\linewidth,height=\textheight,keepaspectratio]{04-InferringRates_files/figure-pdf/ch05_internal memory stan model-1.pdf}
\end{center}

\begin{Shaded}
\begin{Highlighting}[]
\CommentTok{\# Plot rank histograms to check mixing}
\FunctionTok{mcmc\_rank\_hist}\NormalTok{(draws\_df, }\AttributeTok{pars =} \FunctionTok{c}\NormalTok{(}\StringTok{"bias"}\NormalTok{, }\StringTok{"beta"}\NormalTok{))}
\end{Highlighting}
\end{Shaded}

\begin{center}
\includegraphics[width=0.8\linewidth,height=\textheight,keepaspectratio]{04-InferringRates_files/figure-pdf/ch05_internal memory stan model-2.pdf}
\end{center}

\begin{Shaded}
\begin{Highlighting}[]
\CommentTok{\# 2. Prior{-}Posterior Update Check}
\NormalTok{p1 }\OtherTok{\textless{}{-}} \FunctionTok{ggplot}\NormalTok{() }\SpecialCharTok{+}
  \FunctionTok{geom\_density}\NormalTok{(}\AttributeTok{data =}\NormalTok{ draws\_df, }\FunctionTok{aes}\NormalTok{(bias, }\AttributeTok{fill =} \StringTok{"Posterior"}\NormalTok{), }\AttributeTok{alpha =} \FloatTok{0.5}\NormalTok{) }\SpecialCharTok{+}
  \FunctionTok{geom\_density}\NormalTok{(}\AttributeTok{data =}\NormalTok{ draws\_df, }\FunctionTok{aes}\NormalTok{(bias\_prior, }\AttributeTok{fill =} \StringTok{"Prior"}\NormalTok{), }\AttributeTok{alpha =} \FloatTok{0.5}\NormalTok{) }\SpecialCharTok{+}
  \FunctionTok{geom\_vline}\NormalTok{(}\AttributeTok{xintercept =} \DecValTok{0}\NormalTok{, }\AttributeTok{linetype =} \StringTok{"dashed"}\NormalTok{) }\SpecialCharTok{+}
  \FunctionTok{scale\_fill\_manual}\NormalTok{(}\AttributeTok{values =} \FunctionTok{c}\NormalTok{(}\StringTok{"Prior"} \OtherTok{=} \StringTok{"red"}\NormalTok{, }\StringTok{"Posterior"} \OtherTok{=} \StringTok{"blue"}\NormalTok{)) }\SpecialCharTok{+}
  \FunctionTok{theme\_minimal}\NormalTok{() }\SpecialCharTok{+}
  \FunctionTok{ggtitle}\NormalTok{(}\StringTok{"Prior{-}Posterior Update: Bias Parameter"}\NormalTok{)}

\NormalTok{p2 }\OtherTok{\textless{}{-}} \FunctionTok{ggplot}\NormalTok{() }\SpecialCharTok{+}
  \FunctionTok{geom\_density}\NormalTok{(}\AttributeTok{data =}\NormalTok{ draws\_df, }\FunctionTok{aes}\NormalTok{(beta, }\AttributeTok{fill =} \StringTok{"Posterior"}\NormalTok{), }\AttributeTok{alpha =} \FloatTok{0.5}\NormalTok{) }\SpecialCharTok{+}
  \FunctionTok{geom\_density}\NormalTok{(}\AttributeTok{data =}\NormalTok{ draws\_df, }\FunctionTok{aes}\NormalTok{(beta\_prior, }\AttributeTok{fill =} \StringTok{"Prior"}\NormalTok{), }\AttributeTok{alpha =} \FloatTok{0.5}\NormalTok{) }\SpecialCharTok{+}
  \FunctionTok{geom\_vline}\NormalTok{(}\AttributeTok{xintercept =} \DecValTok{1}\NormalTok{, }\AttributeTok{linetype =} \StringTok{"dashed"}\NormalTok{) }\SpecialCharTok{+}
  \FunctionTok{scale\_fill\_manual}\NormalTok{(}\AttributeTok{values =} \FunctionTok{c}\NormalTok{(}\StringTok{"Prior"} \OtherTok{=} \StringTok{"red"}\NormalTok{, }\StringTok{"Posterior"} \OtherTok{=} \StringTok{"blue"}\NormalTok{)) }\SpecialCharTok{+}
  \FunctionTok{theme\_minimal}\NormalTok{() }\SpecialCharTok{+}
  \FunctionTok{ggtitle}\NormalTok{(}\StringTok{"Prior{-}Posterior Update: Beta Parameter"}\NormalTok{)}

\NormalTok{p3 }\OtherTok{\textless{}{-}} \FunctionTok{ggplot}\NormalTok{() }\SpecialCharTok{+}
  \FunctionTok{geom\_point}\NormalTok{(}\AttributeTok{data =}\NormalTok{ draws\_df, }\FunctionTok{aes}\NormalTok{(bias, beta), }\AttributeTok{alpha =} \FloatTok{0.5}\NormalTok{) }\SpecialCharTok{+}
  \FunctionTok{theme\_minimal}\NormalTok{() }\SpecialCharTok{+}
  \FunctionTok{ggtitle}\NormalTok{(}\StringTok{"Correlation"}\NormalTok{)}

\NormalTok{p1 }\SpecialCharTok{+}\NormalTok{ p2 }\SpecialCharTok{+}\NormalTok{ p3}
\end{Highlighting}
\end{Shaded}

\begin{center}
\includegraphics[width=0.8\linewidth,height=\textheight,keepaspectratio]{04-InferringRates_files/figure-pdf/ch05_internal memory stan model-3.pdf}
\end{center}

Now that we know how to model memory as an internal state, we can play
with making the update discount the past, setting a parameter that
indicates after how many trials memory is lost, etc.

\subsection{Memory Model 3: Exponential Forgetting (similar to
RL)}\label{memory-model-3-exponential-forgetting-similar-to-rl}

A more cognitively plausible memory model incorporates forgetting, where
recent events have more influence than distant ones. This can be
implemented by adding a \texttt{forgetting} parameter (often called a
learning rate, α, in RL) that controls the weight of the most recent
outcome versus the previous memory state.

\begin{Shaded}
\begin{Highlighting}[]
\NormalTok{stan\_model }\OtherTok{\textless{}{-}} \StringTok{"}
\StringTok{// Memory model with exponential forgetting (learning{-}rate style).}
\StringTok{// memory[] depends on the free parameter \textquotesingle{}forgetting\textquotesingle{}, so it cannot go in}
\StringTok{// transformed data. It lives as a local vector inside the model block —}
\StringTok{// computed at every leapfrog step but never saved to the posterior draws.}

\StringTok{data \{}
\StringTok{  int\textless{}lower=1\textgreater{} n;}
\StringTok{  array[n] int h;}
\StringTok{  array[n] int other;}
\StringTok{  int\textless{}lower=0, upper=1\textgreater{} run\_diagnostics;}
\StringTok{\}}

\StringTok{parameters \{}
\StringTok{  real bias;}
\StringTok{  real beta;}
\StringTok{  real\textless{}lower=0, upper=1\textgreater{} forgetting;}
\StringTok{\}}

\StringTok{model \{}
\StringTok{  target += beta\_lpdf(forgetting | 1, 1);}
\StringTok{  target += normal\_lpdf(bias | 0, .3);}
\StringTok{  target += normal\_lpdf(beta | 0, .5);}

\StringTok{  // memory is a LOCAL variable — declared inside model, so it lives on the}
\StringTok{  // stack and is discarded after each leapfrog step without touching the saved{-}}
\StringTok{  // draw tape. This is the right pattern when memory depends on a parameter.}
\StringTok{  vector[n] memory;}
\StringTok{  for (trial in 1:n) \{}
\StringTok{    if (trial == 1) memory[trial] = 0.5;}
\StringTok{    target += bernoulli\_logit\_lpmf(h[trial] | bias + beta * logit(memory[trial]));}
\StringTok{    if (trial \textless{} n) \{}
\StringTok{      memory[trial + 1] = (1 {-} forgetting) * memory[trial] + forgetting * other[trial];}
\StringTok{      if (memory[trial + 1] \textless{}= 0) memory[trial + 1] = 0.01;}
\StringTok{      if (memory[trial + 1] \textgreater{}= 1) memory[trial + 1] = 0.99;}
\StringTok{    \}}
\StringTok{  \}}
\StringTok{\}}

\StringTok{generated quantities \{}
\StringTok{  real bias\_prior      = normal\_rng(0, 0.3);}
\StringTok{  real beta\_prior      = normal\_rng(0, 0.5);}
\StringTok{  real forgetting\_prior = beta\_rng(1, 1);}
\StringTok{  vector[n] log\_lik;}
\StringTok{  array[n] int y\_rep;}
\StringTok{  if (run\_diagnostics) \{}
\StringTok{    // Recompute memory for the saved draw — runs once per posterior sample.}
\StringTok{    vector[n] memory\_gq;}
\StringTok{    for (trial in 1:n) \{}
\StringTok{      if (trial == 1) memory\_gq[trial] = 0.5;}
\StringTok{      real lin = bias + beta * logit(memory\_gq[trial]);}
\StringTok{      log\_lik[trial] = bernoulli\_logit\_lpmf(h[trial] | lin);}
\StringTok{      y\_rep[trial]   = bernoulli\_logit\_rng(lin);}
\StringTok{      if (trial \textless{} n) \{}
\StringTok{        memory\_gq[trial + 1] = (1 {-} forgetting) * memory\_gq[trial] + forgetting * other[trial];}
\StringTok{        if (memory\_gq[trial + 1] \textless{}= 0) memory\_gq[trial + 1] = 0.01;}
\StringTok{        if (memory\_gq[trial + 1] \textgreater{}= 1) memory\_gq[trial + 1] = 0.99;}
\StringTok{      \}}
\StringTok{    \}}
\StringTok{  \}}
\StringTok{\}}
\StringTok{"}
\FunctionTok{write\_stan\_file}\NormalTok{(}
\NormalTok{  stan\_model,}
  \AttributeTok{dir =} \StringTok{"stan/"}\NormalTok{,}
  \AttributeTok{basename =} \StringTok{"ch06\_InternalMemory2.stan"}\NormalTok{)}
\end{Highlighting}
\end{Shaded}

\begin{verbatim}
[1] "/Users/au209589/Dropbox/Teaching/AdvancedCognitiveModeling23_book/stan/ch06_InternalMemory2.stan"
\end{verbatim}

\begin{Shaded}
\begin{Highlighting}[]
\DocumentationTok{\#\# Specify where the model is}
\NormalTok{file }\OtherTok{\textless{}{-}} \FunctionTok{here}\NormalTok{(}\StringTok{"stan"}\NormalTok{,}\StringTok{"ch06\_InternalMemory2.stan"}\NormalTok{)}

\CommentTok{\# File path for saved model}
\NormalTok{model\_file }\OtherTok{\textless{}{-}} \FunctionTok{here}\NormalTok{(}\StringTok{"simmodels"}\NormalTok{,}\StringTok{"ch06\_InternalMemory2.rds"}\NormalTok{)}

\CommentTok{\# Check if we need to rerun the simulation}
\ControlFlowTok{if}\NormalTok{ (regenerate\_simulations }\SpecialCharTok{||} \SpecialCharTok{!}\FunctionTok{file.exists}\NormalTok{(model\_file)) \{}
  \CommentTok{\# Compile the model}
\NormalTok{  mod }\OtherTok{\textless{}{-}} \FunctionTok{cmdstan\_model}\NormalTok{(file, }
                       \AttributeTok{cpp\_options =} \FunctionTok{list}\NormalTok{(}\AttributeTok{stan\_threads =} \ConstantTok{FALSE}\NormalTok{),}
                       \AttributeTok{stanc\_options =} \FunctionTok{list}\NormalTok{(}\StringTok{"O1"}\NormalTok{))}
  
\NormalTok{  data\_forgetting }\OtherTok{\textless{}{-}} \FunctionTok{modifyList}\NormalTok{(data, }\FunctionTok{list}\NormalTok{(}\AttributeTok{run\_diagnostics =} \DecValTok{1}\NormalTok{L))}
  \CommentTok{\# The following command calls Stan with specific options.}
\NormalTok{  samples\_memory\_forgetting }\OtherTok{\textless{}{-}}\NormalTok{ mod}\SpecialCharTok{$}\FunctionTok{sample}\NormalTok{(}
    \AttributeTok{data =}\NormalTok{ data\_forgetting,}
    \AttributeTok{seed =} \DecValTok{123}\NormalTok{,}
    \AttributeTok{chains =} \DecValTok{4}\NormalTok{,}
    \AttributeTok{parallel\_chains =} \DecValTok{4}\NormalTok{,}
    \AttributeTok{iter\_warmup =} \DecValTok{1000}\NormalTok{,}
    \AttributeTok{iter\_sampling =} \DecValTok{1000}\NormalTok{,}
    \AttributeTok{refresh =} \DecValTok{0}\NormalTok{,}
    \AttributeTok{adapt\_delta =} \FloatTok{0.95}
\NormalTok{  )}
  
  \CommentTok{\# Save the fitted model}
\NormalTok{  samples\_memory\_forgetting}\SpecialCharTok{$}\FunctionTok{save\_object}\NormalTok{(}\AttributeTok{file =}\NormalTok{ model\_file)}
  \FunctionTok{cat}\NormalTok{(}\StringTok{"Generated new model fit and saved to"}\NormalTok{, model\_file, }\StringTok{"}\SpecialCharTok{\textbackslash{}n}\StringTok{"}\NormalTok{)}
\NormalTok{\} }\ControlFlowTok{else}\NormalTok{ \{}
  \CommentTok{\# Load existing results}
\NormalTok{  samples\_memory\_forgetting }\OtherTok{\textless{}{-}} \FunctionTok{readRDS}\NormalTok{(model\_file)}
  \FunctionTok{cat}\NormalTok{(}\StringTok{"Loaded existing model fit from"}\NormalTok{, model\_file, }\StringTok{"}\SpecialCharTok{\textbackslash{}n}\StringTok{"}\NormalTok{)}
\NormalTok{\}}
\end{Highlighting}
\end{Shaded}

\begin{verbatim}
Loaded existing model fit from /Users/au209589/Dropbox/Teaching/AdvancedCognitiveModeling23_book/simmodels/ch06_InternalMemory2.rds 
\end{verbatim}

\begin{Shaded}
\begin{Highlighting}[]
\NormalTok{samples\_memory\_forgetting}\SpecialCharTok{$}\FunctionTok{summary}\NormalTok{()}
\end{Highlighting}
\end{Shaded}

\begin{verbatim}
# A tibble: 4 x 10
  variable    mean   median    sd   mad       q5     q95  rhat ess_bulk ess_tail
  <chr>      <dbl>    <dbl> <dbl> <dbl>    <dbl>   <dbl> <dbl>    <dbl>    <dbl>
1 lp__     -75.3   -74.9    1.42  1.18  -78.2    -73.5    1.00     219.     283.
2 bias       0.668   0.674  0.256 0.278   0.242    1.08   1.01     157.     168.
3 beta       0.108   0.0603 0.199 0.162  -0.136    0.478  1.01     109.     322.
4 forgett~   0.243   0.141  0.241 0.157   0.0108   0.802  1.02     165.     262.
\end{verbatim}

The memory model we've implemented can be seen as part of a broader
family of models that track and update beliefs based on incoming
evidence. Let's explore how it relates to some key frameworks.

\section{Memory Model 4: Bayesian Agent (Optimal
Updating)}\label{memory-model-4-bayesian-agent-optimal-updating}

A fully Bayesian agent wouldn't just track the rate but maintain a full
probability distribution (a Beta distribution) over the likely bias of
the opponent, updating it optimally according to Bayes' rule after each
observation.

\begin{Shaded}
\begin{Highlighting}[]
\NormalTok{stan\_model }\OtherTok{\textless{}{-}} \StringTok{"}
\StringTok{data \{}
\StringTok{  int\textless{}lower=1\textgreater{} n;  // number of trials}
\StringTok{  array[n] int h;  // agent\textquotesingle{}s choices (0 or 1)}
\StringTok{  array[n] int other;  // other player\textquotesingle{}s choices (0 or 1)}
\StringTok{\}}

\StringTok{parameters \{}
\StringTok{  real\textless{}lower=0\textgreater{} alpha\_prior;  // Prior alpha parameter}
\StringTok{  real\textless{}lower=0\textgreater{} beta\_prior;   // Prior beta parameter}
\StringTok{\}}

\StringTok{transformed parameters \{}
\StringTok{  vector[n] alpha;  // Alpha parameter at each trial}
\StringTok{  vector[n] beta;   // Beta parameter at each trial}
\StringTok{  vector[n] rate;   // Expected rate at each trial}
\StringTok{  }
\StringTok{  // Initialize with prior}
\StringTok{  alpha[1] = alpha\_prior;}
\StringTok{  beta[1] = beta\_prior;}
\StringTok{  rate[1] = alpha[1] / (alpha[1] + beta[1]);}
\StringTok{  }
\StringTok{  // Sequential updating of Beta distribution}
\StringTok{  for(t in 2:n) \{}
\StringTok{    // Update Beta parameters based on previous observation}
\StringTok{    alpha[t] = alpha[t{-}1] + other[t{-}1];}
\StringTok{    beta[t] = beta[t{-}1] + (1 {-} other[t{-}1]);}
\StringTok{    }
\StringTok{    // Calculate expected rate}
\StringTok{    rate[t] = alpha[t] / (alpha[t] + beta[t]);}
\StringTok{  \}}
\StringTok{\}}

\StringTok{model \{}
\StringTok{  // Priors on hyperparameters}
\StringTok{  target += gamma\_lpdf(alpha\_prior | 2, 1);}
\StringTok{  target += gamma\_lpdf(beta\_prior | 2, 1);}
\StringTok{  }
\StringTok{  // Agent\textquotesingle{}s choices follow current rate estimates}
\StringTok{  for(t in 1:n) \{}
\StringTok{    target += bernoulli\_lpmf(h[t] | rate[t]);}
\StringTok{  \}}
\StringTok{\}}

\StringTok{generated quantities \{}
\StringTok{  // We strictly reserve posterior predictive checks for Chapter 6.}
\StringTok{  // Here, we just calculate the initial expected rate for reference.}
\StringTok{  real initial\_rate = alpha\_prior / (alpha\_prior + beta\_prior);}
\StringTok{\}}
\StringTok{"}
\FunctionTok{write\_stan\_file}\NormalTok{(}
\NormalTok{  stan\_model,}
  \AttributeTok{dir =} \StringTok{"stan/"}\NormalTok{,}
  \AttributeTok{basename =} \StringTok{"ch06\_BayesianMemory.stan"}\NormalTok{)}
\end{Highlighting}
\end{Shaded}

\begin{verbatim}
[1] "/Users/au209589/Dropbox/Teaching/AdvancedCognitiveModeling23_book/stan/ch06_BayesianMemory.stan"
\end{verbatim}

\begin{Shaded}
\begin{Highlighting}[]
\DocumentationTok{\#\# Specify where the model is}
\NormalTok{file }\OtherTok{\textless{}{-}} \FunctionTok{here}\NormalTok{(}\StringTok{"stan"}\NormalTok{,}\StringTok{"ch06\_BayesianMemory.stan"}\NormalTok{)}

\CommentTok{\# File path for saved model}
\NormalTok{model\_file }\OtherTok{\textless{}{-}} \FunctionTok{here}\NormalTok{(}\StringTok{"simmodels"}\NormalTok{,}\StringTok{"ch06\_BayesianMemory.rds"}\NormalTok{)}

\CommentTok{\# Check if we need to rerun the simulation}
\ControlFlowTok{if}\NormalTok{ (regenerate\_simulations }\SpecialCharTok{||} \SpecialCharTok{!}\FunctionTok{file.exists}\NormalTok{(model\_file)) \{}
  \CommentTok{\# Compile the model}
\NormalTok{  mod }\OtherTok{\textless{}{-}} \FunctionTok{cmdstan\_model}\NormalTok{(file)}
  
  \CommentTok{\# The following command calls Stan with specific options.}
\NormalTok{  samples\_memory\_bayes }\OtherTok{\textless{}{-}}\NormalTok{ mod}\SpecialCharTok{$}\FunctionTok{sample}\NormalTok{(}
    \AttributeTok{data =}\NormalTok{ data,}
    \AttributeTok{seed =} \DecValTok{123}\NormalTok{,}
    \AttributeTok{chains =} \DecValTok{1}\NormalTok{,}
    \AttributeTok{parallel\_chains =} \DecValTok{2}\NormalTok{,}
    \AttributeTok{threads\_per\_chain =} \DecValTok{1}\NormalTok{,}
    \AttributeTok{iter\_warmup =} \DecValTok{1000}\NormalTok{,}
    \AttributeTok{iter\_sampling =} \DecValTok{1000}\NormalTok{,}
    \AttributeTok{refresh =} \DecValTok{0}\NormalTok{,}
    \AttributeTok{max\_treedepth =} \DecValTok{20}\NormalTok{,}
    \AttributeTok{adapt\_delta =} \FloatTok{0.99}\NormalTok{,}
\NormalTok{  )}
  
  \CommentTok{\# Save the fitted model}
\NormalTok{  samples\_memory\_bayes}\SpecialCharTok{$}\FunctionTok{save\_object}\NormalTok{(}\AttributeTok{file =}\NormalTok{ model\_file)}
  \FunctionTok{cat}\NormalTok{(}\StringTok{"Generated new model fit and saved to"}\NormalTok{, model\_file, }\StringTok{"}\SpecialCharTok{\textbackslash{}n}\StringTok{"}\NormalTok{)}
\NormalTok{\} }\ControlFlowTok{else}\NormalTok{ \{}
  \CommentTok{\# Load existing results}
\NormalTok{  samples\_memory\_bayes }\OtherTok{\textless{}{-}} \FunctionTok{readRDS}\NormalTok{(model\_file)}
  \FunctionTok{cat}\NormalTok{(}\StringTok{"Loaded existing model fit from"}\NormalTok{, model\_file, }\StringTok{"}\SpecialCharTok{\textbackslash{}n}\StringTok{"}\NormalTok{)}
\NormalTok{\}}
\end{Highlighting}
\end{Shaded}

\begin{verbatim}
Loaded existing model fit from /Users/au209589/Dropbox/Teaching/AdvancedCognitiveModeling23_book/simmodels/ch06_BayesianMemory.rds 
\end{verbatim}

\begin{Shaded}
\begin{Highlighting}[]
\NormalTok{samples\_memory\_bayes}\SpecialCharTok{$}\FunctionTok{summary}\NormalTok{() }
\end{Highlighting}
\end{Shaded}

\begin{verbatim}
# A tibble: 364 x 10
   variable      mean median    sd   mad      q5    q95  rhat ess_bulk ess_tail
   <chr>        <dbl>  <dbl> <dbl> <dbl>   <dbl>  <dbl> <dbl>    <dbl>    <dbl>
 1 lp__        -73.3  -73.0  1.07  0.855 -75.4   -72.2  1.00      370.     379.
 2 alpha_prior   2.11   1.76 1.52  1.31    0.361   5.18 1.000     383.     446.
 3 beta_prior    1.25   1.07 0.819 0.651   0.299   2.90 1.00      407.     399.
 4 alpha[1]      2.11   1.76 1.52  1.31    0.361   5.18 1.000     383.     446.
 5 alpha[2]      3.11   2.76 1.52  1.31    1.36    6.18 1.000     383.     446.
 6 alpha[3]      4.11   3.76 1.52  1.31    2.36    7.18 1.000     383.     446.
 7 alpha[4]      4.11   3.76 1.52  1.31    2.36    7.18 1.000     383.     446.
 8 alpha[5]      5.11   4.76 1.52  1.31    3.36    8.18 1.000     383.     446.
 9 alpha[6]      6.11   5.76 1.52  1.31    4.36    9.18 1.000     383.     446.
10 alpha[7]      6.11   5.76 1.52  1.31    4.36    9.18 1.000     383.     446.
# i 354 more rows
\end{verbatim}

\begin{Shaded}
\begin{Highlighting}[]
\CommentTok{\# Extract draws}
\NormalTok{draws\_df }\OtherTok{\textless{}{-}} \FunctionTok{as\_draws\_df}\NormalTok{(samples\_memory\_bayes}\SpecialCharTok{$}\FunctionTok{draws}\NormalTok{())}

\CommentTok{\# First let\textquotesingle{}s look at the priors}
\FunctionTok{ggplot}\NormalTok{(draws\_df) }\SpecialCharTok{+}
  \FunctionTok{geom\_density}\NormalTok{(}\FunctionTok{aes}\NormalTok{(alpha\_prior), }\AttributeTok{fill =} \StringTok{"blue"}\NormalTok{, }\AttributeTok{alpha =} \FloatTok{0.3}\NormalTok{) }\SpecialCharTok{+}
  \FunctionTok{geom\_density}\NormalTok{(}\FunctionTok{aes}\NormalTok{(beta\_prior), }\AttributeTok{fill =} \StringTok{"red"}\NormalTok{, }\AttributeTok{alpha =} \FloatTok{0.3}\NormalTok{) }\SpecialCharTok{+}
  \FunctionTok{theme\_classic}\NormalTok{() }\SpecialCharTok{+}
  \FunctionTok{labs}\NormalTok{(}\AttributeTok{title =} \StringTok{"Prior Distributions"}\NormalTok{,}
       \AttributeTok{x =} \StringTok{"Parameter Value"}\NormalTok{,}
       \AttributeTok{y =} \StringTok{"Density"}\NormalTok{)}
\end{Highlighting}
\end{Shaded}

\begin{center}
\includegraphics[width=0.8\linewidth,height=\textheight,keepaspectratio]{04-InferringRates_files/figure-pdf/ch05_bayesian memory stan model-1.pdf}
\end{center}

\begin{Shaded}
\begin{Highlighting}[]
\CommentTok{\# Now let\textquotesingle{}s look at how the rate evolves over trials}
\CommentTok{\# First melt the rate values across trials into long format}
\NormalTok{rate\_df }\OtherTok{\textless{}{-}}\NormalTok{ draws\_df }\SpecialCharTok{\%\textgreater{}\%}
\NormalTok{  dplyr}\SpecialCharTok{::}\FunctionTok{select}\NormalTok{(}\FunctionTok{starts\_with}\NormalTok{(}\StringTok{"rate["}\NormalTok{)) }\SpecialCharTok{\%\textgreater{}\%}
  \FunctionTok{pivot\_longer}\NormalTok{(}\FunctionTok{everything}\NormalTok{(), }
               \AttributeTok{names\_to =} \StringTok{"trial"}\NormalTok{,}
               \AttributeTok{values\_to =} \StringTok{"rate"}\NormalTok{,}
               \AttributeTok{names\_pattern =} \StringTok{"rate}\SpecialCharTok{\textbackslash{}\textbackslash{}}\StringTok{[(}\SpecialCharTok{\textbackslash{}\textbackslash{}}\StringTok{d+)}\SpecialCharTok{\textbackslash{}\textbackslash{}}\StringTok{]"}\NormalTok{) }\SpecialCharTok{\%\textgreater{}\%}
  \FunctionTok{mutate}\NormalTok{(}\AttributeTok{trial =} \FunctionTok{as.numeric}\NormalTok{(trial))}

\CommentTok{\# Calculate summary statistics for each trial}
\NormalTok{rate\_summary }\OtherTok{\textless{}{-}}\NormalTok{ rate\_df }\SpecialCharTok{\%\textgreater{}\%}
  \FunctionTok{group\_by}\NormalTok{(trial) }\SpecialCharTok{\%\textgreater{}\%}
  \FunctionTok{summarise}\NormalTok{(}
    \AttributeTok{mean\_rate =} \FunctionTok{mean}\NormalTok{(rate),}
    \AttributeTok{lower =} \FunctionTok{quantile}\NormalTok{(rate, }\FloatTok{0.025}\NormalTok{),}
    \AttributeTok{upper =} \FunctionTok{quantile}\NormalTok{(rate, }\FloatTok{0.975}\NormalTok{)}
\NormalTok{  )}


\NormalTok{plot\_data }\OtherTok{\textless{}{-}} \FunctionTok{tibble}\NormalTok{(}\AttributeTok{trial =} \FunctionTok{seq}\NormalTok{(}\DecValTok{120}\NormalTok{), }\AttributeTok{choices =}\NormalTok{ data}\SpecialCharTok{$}\NormalTok{other)}

\CommentTok{\# Plot the evolution of rate estimates}
\FunctionTok{ggplot}\NormalTok{(rate\_summary, }\FunctionTok{aes}\NormalTok{(}\AttributeTok{x =}\NormalTok{ trial)) }\SpecialCharTok{+}
  \FunctionTok{geom\_ribbon}\NormalTok{(}\FunctionTok{aes}\NormalTok{(}\AttributeTok{ymin =}\NormalTok{ lower, }\AttributeTok{ymax =}\NormalTok{ upper), }\AttributeTok{alpha =} \FloatTok{0.2}\NormalTok{) }\SpecialCharTok{+}
  \FunctionTok{geom\_line}\NormalTok{(}\FunctionTok{aes}\NormalTok{(}\AttributeTok{y =}\NormalTok{ mean\_rate), }\AttributeTok{color =} \StringTok{"blue"}\NormalTok{) }\SpecialCharTok{+}
  \CommentTok{\# Add true data points}
  \FunctionTok{geom\_line}\NormalTok{(}\AttributeTok{data =}\NormalTok{ plot\_data, }
             \FunctionTok{aes}\NormalTok{(}\AttributeTok{x =}\NormalTok{ trial, }\AttributeTok{y =}\NormalTok{ choices), }\AttributeTok{color =} \StringTok{"orange"}\NormalTok{, }\AttributeTok{alpha =} \FloatTok{0.5}\NormalTok{) }\SpecialCharTok{+}
  \FunctionTok{theme\_classic}\NormalTok{() }\SpecialCharTok{+}
  \FunctionTok{labs}\NormalTok{(}\AttributeTok{title =} \StringTok{"Evolution of Rate Estimates"}\NormalTok{,}
       \AttributeTok{x =} \StringTok{"Trial"}\NormalTok{,}
       \AttributeTok{y =} \StringTok{"Rate"}\NormalTok{,}
       \AttributeTok{subtitle =} \StringTok{"Blue line: posterior mean, Gray band: 95\% CI"}\NormalTok{) }\SpecialCharTok{+}
  \FunctionTok{ylim}\NormalTok{(}\DecValTok{0}\NormalTok{, }\DecValTok{1}\NormalTok{)}
\end{Highlighting}
\end{Shaded}

\begin{center}
\includegraphics[width=0.8\linewidth,height=\textheight,keepaspectratio]{04-InferringRates_files/figure-pdf/ch05_bayesian memory stan model-2.pdf}
\end{center}

\begin{Shaded}
\begin{Highlighting}[]
\CommentTok{\# Let\textquotesingle{}s also look at the correlation between alpha and beta parameters}
\FunctionTok{ggplot}\NormalTok{(draws\_df) }\SpecialCharTok{+}
  \FunctionTok{geom\_point}\NormalTok{(}\FunctionTok{aes}\NormalTok{(alpha\_prior, beta\_prior), }\AttributeTok{alpha =} \FloatTok{0.1}\NormalTok{) }\SpecialCharTok{+}
  \FunctionTok{theme\_classic}\NormalTok{() }\SpecialCharTok{+}
  \FunctionTok{labs}\NormalTok{(}\AttributeTok{title =} \StringTok{"Correlation between Alpha and Beta Parameters"}\NormalTok{,}
       \AttributeTok{x =} \StringTok{"Alpha"}\NormalTok{,}
       \AttributeTok{y =} \StringTok{"Beta"}\NormalTok{)}
\end{Highlighting}
\end{Shaded}

\begin{center}
\includegraphics[width=0.8\linewidth,height=\textheight,keepaspectratio]{04-InferringRates_files/figure-pdf/ch05_bayesian memory stan model-3.pdf}
\end{center}

\section{Relationship to
Rescorla-Wagner}\label{relationship-to-rescorla-wagner}

The Rescorla-Wagner model of learning---which drives the Reinforcement
Learning agents we will build in in a later chapter follows the form:

\[V_{t+1} = V_{t} + \alpha (R_{t} - V_{t})\]

where:

\begin{itemize}
\item
  \(V_{t}\) is the current Value estimate (in our code: memory{[}t{]}).
\item
  \(\alpha\) is the Learning Rate (in our code: forgetting).
\item
  \(R_{t}\) is the Reward or Outcome (in our code: other{[}t{]}).
\item
  \((R_{t} - V_{t})\) is the Prediction Error.
\end{itemize}

Our memory model with a forgetting parameter follows a very similar
structure:\[\text{memory}_{t+1} = (1 - \text{forgetting}) \times \text{memory}_{t} + \text{forgetting} \times \text{outcome}_{t}\]If
we expand this, we can see it is mathematically identical to the
Rescorla-Wagner update rule:\[\begin{aligned}
\text{memory}_{t+1} &= \text{memory}_{t} - \text{forgetting} \times \text{memory}_{t} + \text{forgetting} \times \text{outcome}_{t} \\
\text{memory}_{t+1} &= \text{memory}_{t} + \text{forgetting} \times (\text{outcome}_{t} - \text{memory}_{t})
\end{aligned}\]

This reveals that our ``exponential forgetting'' memory model is
actually a standard Reinforcement Learning agent with a big fake
moustache. The forgetting parameter dictates how quickly the agent
updates their beliefs (value estimate) based on new prediction errors.

\subsection{Connection to Kalman
Filters}\label{connection-to-kalman-filters}

Our memory model updates beliefs about the probability of right-hand
choices using a weighted average of past observations. This is
conceptually similar to how a Kalman filter works, though simpler:

\begin{itemize}
\item
  Kalman filters maintain both an estimate and uncertainty about that
  estimate
\item
  They optimally weight new evidence based on relative uncertainty
\item
  Our forgetting model uses a fixed weighting scheme (1/trial or the
  forgetting parameter)
\item
  The Bayesian Agent (\texttt{ch06\_BayesianMemory.stan}) captures
  uncertainty via the Beta distribution, similar in spirit but simpler
  than a Kalman filter.
\end{itemize}

\subsection{Connection to Hierarchical Gaussian Filter
(HGF)}\label{connection-to-hierarchical-gaussian-filter-hgf}

The HGF extends these ideas by:

\begin{itemize}
\item
  Tracking beliefs at multiple levels
\item
  Allowing learning rates to vary over time
\item
  Explicitly modeling environmental volatility
\end{itemize}

Our model could be seen as the simplest case of an HGF where:

\begin{itemize}
\item
  We only track one level (probability of right-hand choice)
\item
  Have a fixed learning rate (forgetting parameter)
\item
  Don't explicitly model environmental volatility
\end{itemize}

\subsection{Implications for Model
Development}\label{implications-for-model-development}

Understanding these relationships helps us think about how models relate
to each other and to extend our model:

\begin{itemize}
\item
  We could add uncertainty estimates to get Kalman-like behavior
\item
  We could make the forgetting parameter dynamic to capture changing
  learning rates
\item
  We could add multiple levels to track both immediate probabilities and
  longer-term trends
\end{itemize}

Each extension would make the model more flexible but also more complex
to fit to data. The choice depends on our specific research questions
and available data.

\section{Conclusion: Estimating Parameters and Exploring
Memory}\label{conclusion-estimating-parameters-and-exploring-memory}

This chapter marked a crucial transition from simulating models with
known parameters to the core task of \textbf{parameter estimation}:
inferring plausible parameter values from observed data using Bayesian
inference and Stan.

We learned how to:

\begin{itemize}
\item
  \textbf{Specify Bayesian models} in Stan, defining data, parameters,
  priors, and data models.
\item
  \textbf{Fit models} using \texttt{cmdstanr} to obtain posterior
  distributions for parameters like choice bias (\texttt{theta}).
\item
  \textbf{Utilize transformations} like log-odds for computational
  benefits and flexible prior specification.
\item
  \textbf{Validate our fitting procedure} through \textbf{parameter
  recovery}, ensuring our models can retrieve known values from
  simulated data.
\item
  \textbf{Implement different cognitive models} for history-dependent
  choice, exploring various ways to represent \textbf{memory}: as an
  external predictor, an internal updating state, incorporating
  exponential forgetting (linking to Reinforcement Learning principles),
  and even as a fully Bayesian belief update.
\item
  \textbf{Recognize connections} between these models and broader
  frameworks like Kalman filters and HGF.
\end{itemize}

By fitting these models, we moved beyond simply describing behavior
(like the cumulative rates in Chapter 3) to quantifying underlying
latent parameters (like \texttt{bias}, \texttt{beta},
\texttt{forgetting}).

However, a single fit to a single simulated dataset is only a
preliminary sanity check --- the equivalent of flipping a coin once to
check if it's fair. It does not tell us whether the model behaves well
across a wide range of parameter values, whether our credible intervals
are correctly calibrated, or whether the fitted model can actually
reproduce the patterns in the data.

\textbf{Coming up in Chapter 6:} We take the Stan models built in this
chapter and subject them to a four-phase validation battery before they
are allowed near real human data:

\begin{enumerate}
\def\labelenumi{\arabic{enumi}.}
\tightlist
\item
  \textbf{Prior Predictive Check} --- do the priors allow only
  cognitively plausible behavior?
\item
  \textbf{Parameter Recovery} --- does the model find the right answer
  across many simulated datasets and noise levels?
\item
  \textbf{Simulation-Based Calibration Checks (SBC)} --- are our
  posterior intervals actually calibrated (does a 95\% CI contain the
  truth 95\% of the time)?
\item
  \textbf{Posterior Predictive Check} --- does the fitted model
  reproduce key patterns in the observed data?
\end{enumerate}

Only after passing all four tests does a model earn the right to touch
real empirical data.

\bookmarksetup{startatroot}

\chapter{Assessing Model Quality}\label{sec-model-quality-checks}

\begin{quote}
\textbf{📍 Where we are in the Bayesian modeling workflow:} Ch. 5 fit a
Stan model to one simulated dataset. Obtaining a posterior is not the
same as trusting it. This chapter is the validation gate: a six-phase
battery the model must pass before touching real human data. These
techniques reappear in every subsequent chapter --- learn them once
here, apply them everywhere.
\end{quote}

\section{Learning Objectives}\label{learning-objectives-4}

In Chapter 5, we learned how to write Stan code and infer parameters for
a single simulated dataset. However, obtaining a posterior distribution
does not mean the model is correct or the inference is reliable.

This chapter introduces the strict diagnostic battery required to
validate any cognitive model. After completing this chapter, you will be
able to:

\begin{itemize}
\item
  \textbf{Assess Generative Plausibility (Prior Predictive Checks):}
  Simulate data from your model's priors \emph{before} fitting to ensure
  the mathematical bounds and prior assumptions are physically and
  cognitively plausible, that is, generate plausible behaviors.
\item
  \textbf{Validate Identifiability (Parameter Recovery):} Show that your
  model can accurately recover known ground-truth parameters across
  different noise levels.
\item
  \textbf{Calibrate Uncertainty (Simulation-Based Calibration Checks):}
  Use SBC to guarantee your computational engine is unbiased and your
  posterior uncertainties are mathematically sound.
\item
  \textbf{Perform Posterior Predictive Checks:} Simulate data from your
  \emph{fitted} model to assess whether it can reproduce key patterns
  observed in the actual data.
\item
  \textbf{Visualize Prior-Posterior Updates:} Compare prior and
  posterior distributions to understand how the data informed parameter
  estimates.
\item
  \textbf{Test Robustness (Sensitivity Analysis):} Evaluate how
  sensitive your parameter estimates are to the choice of prior
  distributions using both manual methods and the \texttt{priorsense}
  package.
\end{itemize}

Before going into each phase in detail, here is a one-row summary of
what each check does and what a failure looks like. Return to this table
whenever you feel lost.

\begin{longtable}[]{@{}
  >{\raggedright\arraybackslash}p{(\linewidth - 6\tabcolsep) * \real{0.2500}}
  >{\raggedright\arraybackslash}p{(\linewidth - 6\tabcolsep) * \real{0.2500}}
  >{\raggedright\arraybackslash}p{(\linewidth - 6\tabcolsep) * \real{0.2500}}
  >{\raggedright\arraybackslash}p{(\linewidth - 6\tabcolsep) * \real{0.2500}}@{}}
\toprule\noalign{}
\begin{minipage}[b]{\linewidth}\raggedright
Phase
\end{minipage} & \begin{minipage}[b]{\linewidth}\raggedright
Check
\end{minipage} & \begin{minipage}[b]{\linewidth}\raggedright
Core question
\end{minipage} & \begin{minipage}[b]{\linewidth}\raggedright
What failure looks like
\end{minipage} \\
\midrule\noalign{}
\endhead
\bottomrule\noalign{}
\endlastfoot
1 & Mathematical formulation & Is the Stan code doing what we think? &
Model block and generated quantities disagree \\
2 & Prior Predictive & Do the priors allow only plausible behaviour? &
Simulated data saturates impossible boundaries \\
3 & Posterior Predictive & Does the fitted model reproduce observed
patterns? & Observed statistic falls in the tail of replicates \\
4 & Prior-Posterior Update & Did the data teach the model anything? &
Posterior ≈ prior; or posterior far from truth \\
5 & Sensitivity & Do conclusions change if we tweak the prior? & Large
posterior shift under small prior perturbation \\
6 & Parameter Recovery + SBC & Can the engine find the truth? Are
credible intervals calibrated? & Recovery off the diagonal; rank
histogram U-shaped or skewed \\
\end{longtable}

\textbf{Note on terminology:} In Ch. 5 we used ``parameter recovery''
loosely to mean fitting one simulated dataset and checking whether the
posterior landed near the true value. Phases 6a--b here are the
\emph{rigorous} version: systematic recovery across many parameter
values (6a) and formal calibration of the posterior intervals across the
full prior space (SBC, 6b).

\section{Introduction: Why Parameter Estimates Aren't
Enough}\label{introduction-why-parameter-estimates-arent-enough}

In the previous chapters, we journeyed from observing behavior (Chapter
2), formalizing strategies (Chapter 3), to fitting models and estimating
parameters (Chapter 5).

But how confident are we about the model, its assumptions, the way we
coded it? The model itself, the structure we assumed could be
mathematically misspecified (not do what we think it does),
unidentifiable (produce similar behaviors with an infinite set of
combinations of parameter values), and/or simply unable to capture the
data (sometimes even the data that we simulated from it!).

While rigorous Bayesian workflows can take many forms, here we structure
our validations in \textbf{six sequential phases}, each answering a
different question before we are allowed to proceed:

\begin{itemize}
\tightlist
\item
  \textbf{Phase 1} confirms the Stan code is mathematically well-formed.
\item
  \textbf{Phase 2} (Prior Predictive) checks that our priors generate
  plausible data before we see any observations.
\item
  \textbf{Phase 3} (Posterior Predictive) checks that the fitted model
  can reproduce the key patterns in the data it was trained on.
\item
  \textbf{Phase 4} (Prior-Posterior Update) visualises how much the data
  actually moved our beliefs.
\item
  \textbf{Phase 5} (Sensitivity) tests whether conclusions hold under
  small changes to the prior.
\item
  \textbf{Phase 6} (Parameter Recovery + SBC) proves the computational
  engine is unbiased and that credible intervals are correctly
  calibrated.
\end{itemize}

This chapter executes all six phases on the simple biased agent, then
repeats the workflow on the memory model --- where we will discover a
real identifiability problem and fix it.

\section{Loading Packages and Previous
Results}\label{loading-packages-and-previous-results}

Let's start by loading the necessary packages and the fitted model
object from Chapter 5. This avoids re-running the fitting process
unnecessarily and ensures consistency.

\begin{Shaded}
\begin{Highlighting}[]
\CommentTok{\# Load required packages using pacman for efficient package management}
\CommentTok{\# pacman::p\_load ensures packages are installed if missing, then loads them}
\NormalTok{pacman}\SpecialCharTok{::}\FunctionTok{p\_load}\NormalTok{(}
\NormalTok{  tidyverse,    }\CommentTok{\# Core suite for data manipulation (dplyr) and plotting (ggplot2)}
\NormalTok{  here,         }\CommentTok{\# For creating robust file paths relative to the project root (optional but good practice)}
\NormalTok{  posterior,    }\CommentTok{\# Provides tools for working with posterior distributions}
\NormalTok{  cmdstanr,     }\CommentTok{\# The R interface to Stan (via CmdStan)}
\NormalTok{  bayesplot,    }\CommentTok{\# Specialized plotting functions for Bayesian models}
\NormalTok{  patchwork,    }\CommentTok{\# For easily combining multiple ggplot plots}
\NormalTok{  priorsense,   }\CommentTok{\# For automated prior sensitivity analysis}
\NormalTok{  SBC,          }\CommentTok{\# Simulation{-}Based Calibration Checks}
\NormalTok{  future,       }\CommentTok{\# Parallel processing}
\NormalTok{  furrr         }\CommentTok{\# Parallel mapping}
\NormalTok{)}

\CommentTok{\# Set a default theme for ggplot for consistency}
\FunctionTok{theme\_set}\NormalTok{(}\FunctionTok{theme\_classic}\NormalTok{()) }\CommentTok{\# Use a classic theme with white background}

\CommentTok{\# Define relative paths }
\NormalTok{stan\_model\_dir }\OtherTok{\textless{}{-}} \StringTok{"stan"}
\NormalTok{stan\_results\_dir }\OtherTok{\textless{}{-}} \StringTok{"simmodels"}
\NormalTok{sim\_data\_dir }\OtherTok{\textless{}{-}} \StringTok{"simdata"}

\CommentTok{\# Create directories if they don\textquotesingle{}t exist (important for first run)}
\ControlFlowTok{if}\NormalTok{ (}\SpecialCharTok{!}\FunctionTok{dir.exists}\NormalTok{(stan\_model\_dir)) }\FunctionTok{dir.create}\NormalTok{(stan\_model\_dir)}
\ControlFlowTok{if}\NormalTok{ (}\SpecialCharTok{!}\FunctionTok{dir.exists}\NormalTok{(stan\_results\_dir)) }\FunctionTok{dir.create}\NormalTok{(stan\_results\_dir)}
\ControlFlowTok{if}\NormalTok{ (}\SpecialCharTok{!}\FunctionTok{dir.exists}\NormalTok{(sim\_data\_dir)) }\FunctionTok{dir.create}\NormalTok{(sim\_data\_dir)}
\end{Highlighting}
\end{Shaded}

\section{Phase 1: Mathematical Formulation and Stan
implementation}\label{phase-1-mathematical-formulation-and-stan-implementation}

Just as a reminder, for our Simple Biased Agent, the generative process
across \(n\) trials is defined as: \[h_i \sim \text{Bernoulli}(\theta)\]
\[\theta = \text{logit}^{-1}(\theta_{\text{logit}})\]
\[\theta_{\text{logit}} \sim \text{Normal}(0, 1.5)\]

Notice our parameterization. By operating on the unconstrained space
(\(\theta_{\text{logit}} \in [-\infty, \infty]\)) rather than the
constrained probability space (\(\theta \in [0, 1]\)), we avoid boundary
effects that can cause the No-U-Turn Sampler (NUTS) to pathologically
bounce against parameter limits.

Stan's default initialization draws uniformly from \([-2, 2]\) on the
unconstrained (logit) space. When our models will grow more complex,
this arbitrary initialization might be less than optimal and we'll have
to explicitly initialize from the prior or use multi-path variational
inference (\texttt{Pathfinder}) if the chains struggle to start.

We will now reproduce the Stan code for the model, and modify it to
include the generated quantities block. This block executes after the
model has been fitted, and is very useful in our checks as we can
generate and save several variables to be used in our model assessment:

\begin{itemize}
\item
  the prior (e.g.~for visualization in prior-posterior update checks),
\item
  simulations of model expectations from the prior only (before exposure
  to data) for prior predictive checks
\item
  simulations of model expectations from the posteriors (after exposure
  to data), for posterior predictive checks
\item
  the log likelihood for each data point, to assess how well they are
  captured by the model
\item
  the joint log prior, which tells us how likely the posterior is given
  the prior and will be used to assess the sensitivity of the posterior
  to the choice of priors
\end{itemize}

\begin{tcolorbox}[enhanced jigsaw, toprule=.15mm, toptitle=1mm, bottomtitle=1mm, leftrule=.75mm, rightrule=.15mm, colframe=quarto-callout-tip-color-frame, left=2mm, colback=white, titlerule=0mm, coltitle=black, colbacktitle=quarto-callout-tip-color!10!white, opacityback=0, title=\textcolor{quarto-callout-tip-color}{\faLightbulb}\hspace{0.5em}{Stan Design Pattern --- Priors as Data}, arc=.35mm, opacitybacktitle=0.6, breakable, bottomrule=.15mm]

Notice that we declare the prior hyperparameters
(\texttt{prior\_theta\_mu}, \texttt{prior\_theta\_sigma}) in the
\textbf{\texttt{data} block} rather than hardcoding them in the
\texttt{model} block. This separation has two concrete benefits:

\begin{enumerate}
\def\labelenumi{\arabic{enumi}.}
\tightlist
\item
  \textbf{No recompilation when exploring priors.} Changing a hardcoded
  \texttt{normal(0,\ 1.5)} to \texttt{normal(0,\ 2.0)} in Stan requires
  recompiling the C++ model binary (\textasciitilde10--30 seconds).
  Passing them as data means you change a single R list entry and re-run
  instantly.
\item
  \textbf{SBC priors can differ from the fitted model's priors.} When
  SBC simulates from the prior, it draws hyperparameters from the same
  distribution specified in the data list. Keeping this in sync is
  trivial when priors are data; it is easy to introduce accidental
  mismatches when they are hardcoded.
\end{enumerate}

We use this pattern for every model with tunable priors from this
chapter onward.

\end{tcolorbox}

Note that we are adding data for the prior parameters, so we can quickly
change them without recompiling the model, and that we will explain the
different additional variables in detail throughout the chapter.

\begin{tcolorbox}[enhanced jigsaw, toprule=.15mm, toptitle=1mm, bottomtitle=1mm, leftrule=.75mm, rightrule=.15mm, colframe=quarto-callout-tip-color-frame, left=2mm, colback=white, titlerule=0mm, coltitle=black, colbacktitle=quarto-callout-tip-color!10!white, opacityback=0, title=\textcolor{quarto-callout-tip-color}{\faLightbulb}\hspace{0.5em}{Stan Efficiency Rule 2 --- Conditional Diagnostics: The
\texttt{run\_diagnostics} Flag}, arc=.35mm, opacitybacktitle=0.6, breakable, bottomrule=.15mm]

The six-phase validation pipeline requires \textbf{thousands of Stan
fits} during SBC (Simulation-Based Calibration Checks). Each SBC
iteration only needs to verify posterior calibration --- it does
\textbf{not} need:

\begin{itemize}
\tightlist
\item
  PPC draws (\texttt{y\_rep}) --- these require one
  \texttt{bernoulli\_rng()} call per observation per draw
\item
  Pointwise log-likelihoods (\texttt{log\_lik}) --- needed for LOO-CV,
  but LOO is not run on SBC fits
\item
  Prior predictive samples --- only needed for the visual prior check,
  not for calibration
\end{itemize}

Generating these quantities inside \texttt{generated\ quantities} on
every SBC iteration wastes substantial compute. We gate them with a
single integer flag in the data block:

\begin{Shaded}
\begin{Highlighting}[]
\DataTypeTok{int}\NormalTok{\textless{}}\KeywordTok{lower}\NormalTok{=}\DecValTok{0}\NormalTok{, }\KeywordTok{upper}\NormalTok{=}\DecValTok{1}\NormalTok{\textgreater{} run\_diagnostics;  }\CommentTok{// 1 = full diagnostics; 0 = calibration only}
\end{Highlighting}
\end{Shaded}

\begin{Shaded}
\begin{Highlighting}[]
\KeywordTok{generated quantities}\NormalTok{ \{}
  \DataTypeTok{vector}\NormalTok{[N\_trials] log\_lik;}
  \DataTypeTok{array}\NormalTok{[N\_trials] }\DataTypeTok{int}\NormalTok{ y\_rep;}
  \ControlFlowTok{if}\NormalTok{ (run\_diagnostics) \{}
    \ControlFlowTok{for}\NormalTok{ (i }\ControlFlowTok{in} \DecValTok{1}\NormalTok{:N\_trials) \{}
\NormalTok{      log\_lik[i] = bernoulli\_logit\_lpmf(y[i] | theta\_logit);}
\NormalTok{      y\_rep[i]   = bernoulli\_logit\_rng(theta\_logit);}
\NormalTok{    \}}
\NormalTok{  \}}
\NormalTok{\}}
\end{Highlighting}
\end{Shaded}

Then in R, the main validation fit uses \texttt{run\_diagnostics\ =\ 1}
and SBC backends use \texttt{run\_diagnostics\ =\ 0}:

\begin{Shaded}
\begin{Highlighting}[]
\NormalTok{stan\_data\_main }\OtherTok{\textless{}{-}} \FunctionTok{list}\NormalTok{(..., }\AttributeTok{run\_diagnostics =} \DecValTok{1}\NormalTok{L)   }\CommentTok{\# single model check}
\NormalTok{stan\_data\_sbc  }\OtherTok{\textless{}{-}} \FunctionTok{list}\NormalTok{(..., }\AttributeTok{run\_diagnostics =} \DecValTok{0}\NormalTok{L)   }\CommentTok{\# every SBC iteration}
\end{Highlighting}
\end{Shaded}

\textbf{Why it matters at scale.} For a model with 120 trials and 100
SBC datasets × 4 chains × 1000 draws: gating \texttt{log\_lik} and
\texttt{y\_rep} eliminates 100 × 4 × 1000 × 120 ≈ \textbf{48 million}
log-probability evaluations, typically cutting SBC wall-clock time by
20--40\%.

Note that \textbf{prior samples} (e.g.,
\texttt{bias\_prior\ =\ normal\_rng(0,\ 0.3)}) are cheap scalar draws
and can remain ungated --- they don't loop over observations.

We introduced this flag in simplified form in Chapter 4's memory model.
From this chapter onward it appears in every Stan model in the book.

\end{tcolorbox}

\begin{tcolorbox}[enhanced jigsaw, toprule=.15mm, toptitle=1mm, bottomtitle=1mm, leftrule=.75mm, rightrule=.15mm, colframe=quarto-callout-tip-color-frame, left=2mm, colback=white, titlerule=0mm, coltitle=black, colbacktitle=quarto-callout-tip-color!10!white, opacityback=0, title=\textcolor{quarto-callout-tip-color}{\faLightbulb}\hspace{0.5em}{Stan Efficiency --- Post-Hoc Generated Quantities with
\texttt{\$generate\_quantities()}}, arc=.35mm, opacitybacktitle=0.6, breakable, bottomrule=.15mm]

The \texttt{run\_diagnostics} flag trades off between fitting speed and
diagnostic richness. A complementary pattern avoids the trade-off
entirely: \textbf{fit once with \texttt{run\_diagnostics\ =\ 0}, then
replay the saved draws through the \texttt{generated\ quantities} block
on demand} using \texttt{\$generate\_quantities()}.

\begin{Shaded}
\begin{Highlighting}[]
\CommentTok{\# Fast fit — no log\_lik or y\_rep computed during sampling}
\NormalTok{fit\_fast }\OtherTok{\textless{}{-}}\NormalTok{ mod}\SpecialCharTok{$}\FunctionTok{sample}\NormalTok{(}\AttributeTok{data =} \FunctionTok{list}\NormalTok{(..., }\AttributeTok{run\_diagnostics =} \DecValTok{0}\NormalTok{L), ...)}

\CommentTok{\# Later: compute log\_lik and y\_rep from the saved draws without re{-}running MCMC}
\NormalTok{fit\_gq }\OtherTok{\textless{}{-}}\NormalTok{ mod}\SpecialCharTok{$}\FunctionTok{generate\_quantities}\NormalTok{(}
  \AttributeTok{fitted\_params =}\NormalTok{ fit\_fast,}
  \AttributeTok{data          =} \FunctionTok{list}\NormalTok{(..., }\AttributeTok{run\_diagnostics =} \DecValTok{1}\NormalTok{L)}
\NormalTok{)}

\CommentTok{\# Extract as usual}
\NormalTok{log\_lik }\OtherTok{\textless{}{-}}\NormalTok{ fit\_gq}\SpecialCharTok{$}\FunctionTok{draws}\NormalTok{(}\StringTok{"log\_lik"}\NormalTok{)}
\NormalTok{y\_rep   }\OtherTok{\textless{}{-}}\NormalTok{ fit\_gq}\SpecialCharTok{$}\FunctionTok{draws}\NormalTok{(}\StringTok{"y\_rep"}\NormalTok{)}
\end{Highlighting}
\end{Shaded}

\texttt{\$generate\_quantities()} runs only the
\texttt{generated\ quantities} block --- it never touches the sampler.
This means:

\begin{itemize}
\tightlist
\item
  \textbf{LOO and PPC can be deferred} until you actually need them,
  without re-running MCMC.
\item
  \textbf{SBC backends always run with \texttt{run\_diagnostics\ =\ 0}}
  and you call \texttt{\$generate\_quantities()} only on the subset of
  fits you want to inspect.
\item
  The saved \texttt{CmdStanFit} object from a previous session can be
  replayed by loading it with \texttt{readRDS()} and passing it as
  \texttt{fitted\_params}.
\end{itemize}

The only requirement is that the \texttt{data} list passed to
\texttt{\$generate\_quantities()} matches the data block of the compiled
model --- the same model object must be used for both calls.

\end{tcolorbox}

\begin{Shaded}
\begin{Highlighting}[]
\CommentTok{\# Stan code of the model, adding generated quantities for predictive checks}
\NormalTok{stan\_model\_gq\_code }\OtherTok{\textless{}{-}} \StringTok{"}
\StringTok{data \{}
\StringTok{  int\textless{}lower=1\textgreater{} N\_trials;               // Total number of trials}
\StringTok{  array[N\_trials] int\textless{}lower=0, upper=1\textgreater{} y; // Observed choices}
\StringTok{  real prior\_theta\_mu;                 // Mean of the prior for theta\_logit}
\StringTok{  real\textless{}lower=0\textgreater{} prior\_theta\_sigma;     // SD of the prior for theta\_logit}
\StringTok{  int\textless{}lower=0, upper=1\textgreater{} run\_diagnostics; // Flag: 1 to compute PPC/log\_lik, 0 to skip}
\StringTok{\}}
\StringTok{parameters \{}
\StringTok{  real theta\_logit; // Represents the log{-}odds of choosing option 1 (right)}
\StringTok{\}}
\StringTok{model \{}
\StringTok{  target += normal\_lpdf(theta\_logit | prior\_theta\_mu, prior\_theta\_sigma); // prior}
\StringTok{  target += bernoulli\_logit\_lpmf(y | theta\_logit); // likelihood}
\StringTok{\}}

\StringTok{generated quantities \{}
\StringTok{  // Create the prior for theta}
\StringTok{  real theta\_logit\_prior = normal\_rng(prior\_theta\_mu, prior\_theta\_sigma); }
\StringTok{  // {-}{-}{-} Parameter Transformation {-}{-}{-}}
\StringTok{  real theta\_prior = inv\_logit(theta\_logit\_prior); }
\StringTok{  real\textless{}lower=0, upper=1\textgreater{} theta = inv\_logit(theta\_logit);}

\StringTok{  // {-}{-}{-} Conditional Diagnostics {-}{-}{-}}
\StringTok{  // These blocks are expensive to run and store during SBC.}
\StringTok{  // We wrap them in a conditional flag.}
\StringTok{  array[N\_trials] int y\_prior\_rep;}
\StringTok{  int\textless{}lower=0, upper=N\_trials\textgreater{} prior\_rep\_sum;}
\StringTok{  array[N\_trials] int y\_post\_rep;}
\StringTok{  int\textless{}lower=0, upper=N\_trials\textgreater{} post\_rep\_sum;}
\StringTok{  vector[N\_trials] log\_lik;}
\StringTok{  real lprior;}

\StringTok{  if (run\_diagnostics) \{}
\StringTok{    // {-}{-}{-} Prior Predictive Check Simulation {-}{-}{-}}
\StringTok{    y\_prior\_rep = bernoulli\_logit\_rng(rep\_vector(theta\_logit\_prior, N\_trials));}
\StringTok{    prior\_rep\_sum = sum(y\_prior\_rep);}

\StringTok{    // {-}{-}{-} Posterior Predictive Check Simulation {-}{-}{-}}
\StringTok{    y\_post\_rep = bernoulli\_logit\_rng(rep\_vector(theta\_logit, N\_trials));}
\StringTok{    post\_rep\_sum = sum(y\_post\_rep); }
\StringTok{    }
\StringTok{    // {-}{-}{-} Log{-}likelihood {-}{-}{-}}
\StringTok{    for (i in 1:N\_trials) log\_lik[i] = bernoulli\_logit\_lpmf(y[i] | theta\_logit);}
\StringTok{  \}}
\StringTok{  // joint log prior (useful for sensitivity even if run\_diagnostics=0)}
\StringTok{  lprior = normal\_lpdf(theta\_logit | prior\_theta\_mu, prior\_theta\_sigma);}
\StringTok{\}}
\StringTok{"}
\NormalTok{stan\_file\_gq }\OtherTok{\textless{}{-}} \FunctionTok{file.path}\NormalTok{(stan\_model\_dir, }\StringTok{"ch07\_SimpleBernoulli\_GQ.stan"}\NormalTok{)}
\FunctionTok{writeLines}\NormalTok{(stan\_model\_gq\_code, stan\_file\_gq)}

\CommentTok{\# Define path for saving/loading the fitted model object that includes the generated quantities}
\NormalTok{model\_file\_gq }\OtherTok{\textless{}{-}} \FunctionTok{file.path}\NormalTok{(stan\_results\_dir, }\StringTok{"ch07\_SimpleBernoulli\_GQ.rds"}\NormalTok{)}

\CommentTok{\# Generate the required data}

\CommentTok{\# Load the simulation data used in Chapter 5}
\NormalTok{sim\_data\_file }\OtherTok{\textless{}{-}} \FunctionTok{file.path}\NormalTok{(sim\_data\_dir, }\StringTok{"ch06\_randomnoise.csv"}\NormalTok{) }
\NormalTok{d }\OtherTok{\textless{}{-}} \FunctionTok{read\_csv}\NormalTok{(sim\_data\_file, }\AttributeTok{show\_col\_types =} \ConstantTok{FALSE}\NormalTok{) }

\CommentTok{\# Subset data for the specific agent fitted in Chapter 5 (rate=0.8, noise=0)}
\NormalTok{d1 }\OtherTok{\textless{}{-}}\NormalTok{ d }\SpecialCharTok{\%\textgreater{}\%} \FunctionTok{filter}\NormalTok{(true\_rate }\SpecialCharTok{==} \FloatTok{0.8}\NormalTok{, true\_noise }\SpecialCharTok{==} \FloatTok{0.0}\NormalTok{)}

\NormalTok{N\_trials }\OtherTok{\textless{}{-}} \FunctionTok{nrow}\NormalTok{(d1)}

\CommentTok{\# Prepare the observed data list in the format Stan expects}
\NormalTok{observed\_data\_list }\OtherTok{\textless{}{-}} \FunctionTok{list}\NormalTok{(}
  \AttributeTok{N\_trials =}\NormalTok{ N\_trials,}
  \AttributeTok{y =}\NormalTok{ d1}\SpecialCharTok{$}\NormalTok{choice,}
  \AttributeTok{prior\_theta\_mu =} \DecValTok{0}\NormalTok{,}
  \AttributeTok{prior\_theta\_sigma =} \FloatTok{1.5}\NormalTok{,}
  \AttributeTok{run\_diagnostics =} \DecValTok{1}
\NormalTok{)}


\CommentTok{\# {-}{-}{-}{-}{-}{-}{-}{-}{-}{-}{-}{-}{-}{-}{-}{-}{-}{-}{-}{-}{-}{-}{-}{-}{-}{-}{-}{-}{-}{-}{-}{-}{-}{-}{-}{-}{-}{-}{-}{-}{-}{-}{-}{-}{-}{-}{-}{-}{-}{-}{-}{-}{-}{-}{-}{-}{-}{-}{-}{-}{-}{-}{-}{-}{-}{-}{-}{-}{-}{-}{-}{-}{-}{-}{-}}
\CommentTok{\# mod\_gq is compiled OUTSIDE the cache conditional so it is always available —}
\CommentTok{\# for fitting, for the sensitivity analysis, and for the SBC backend.}
\CommentTok{\# Compilation is fast (C++ is cached by CmdStan after the first run).}
\CommentTok{\# {-}{-}{-}{-}{-}{-}{-}{-}{-}{-}{-}{-}{-}{-}{-}{-}{-}{-}{-}{-}{-}{-}{-}{-}{-}{-}{-}{-}{-}{-}{-}{-}{-}{-}{-}{-}{-}{-}{-}{-}{-}{-}{-}{-}{-}{-}{-}{-}{-}{-}{-}{-}{-}{-}{-}{-}{-}{-}{-}{-}{-}{-}{-}{-}{-}{-}{-}{-}{-}{-}{-}{-}{-}{-}{-}}
\NormalTok{mod\_gq }\OtherTok{\textless{}{-}} \FunctionTok{cmdstan\_model}\NormalTok{(stan\_file\_gq)}

\NormalTok{.run\_sampling }\OtherTok{\textless{}{-}}\NormalTok{ regenerate\_simulations }\SpecialCharTok{||} \SpecialCharTok{!}\FunctionTok{file.exists}\NormalTok{(model\_file\_gq)}

\ControlFlowTok{if}\NormalTok{ (}\SpecialCharTok{!}\NormalTok{.run\_sampling) \{}
\NormalTok{  fit\_gq }\OtherTok{\textless{}{-}} \FunctionTok{readRDS}\NormalTok{(model\_file\_gq)}
  \CommentTok{\# Detect stale cache: if y\_post\_rep is missing, the model was saved before PPC was added}
\NormalTok{  .has\_ppc }\OtherTok{\textless{}{-}} \FunctionTok{tryCatch}\NormalTok{(\{ fit\_gq}\SpecialCharTok{$}\FunctionTok{draws}\NormalTok{(}\StringTok{"y\_post\_rep"}\NormalTok{); }\ConstantTok{TRUE}\NormalTok{ \}, }\AttributeTok{error =} \ControlFlowTok{function}\NormalTok{(e) }\ConstantTok{FALSE}\NormalTok{)}
  \ControlFlowTok{if}\NormalTok{ (}\SpecialCharTok{!}\NormalTok{.has\_ppc) .run\_sampling }\OtherTok{\textless{}{-}} \ConstantTok{TRUE}
\NormalTok{\}}

\ControlFlowTok{if}\NormalTok{ (.run\_sampling) \{}
\NormalTok{  fit\_gq }\OtherTok{\textless{}{-}}\NormalTok{ mod\_gq}\SpecialCharTok{$}\FunctionTok{sample}\NormalTok{(}
    \AttributeTok{data             =}\NormalTok{ observed\_data\_list,}
    \AttributeTok{seed             =} \DecValTok{123}\NormalTok{,}
    \AttributeTok{chains           =} \DecValTok{4}\NormalTok{,}
    \AttributeTok{parallel\_chains  =} \FunctionTok{min}\NormalTok{(}\DecValTok{4}\NormalTok{, future}\SpecialCharTok{::}\FunctionTok{availableCores}\NormalTok{()),}
    \AttributeTok{iter\_warmup      =} \DecValTok{1000}\NormalTok{,}
    \AttributeTok{iter\_sampling    =} \DecValTok{2000}\NormalTok{,}
    \AttributeTok{refresh          =} \DecValTok{500}\NormalTok{,}
    \AttributeTok{max\_treedepth    =} \DecValTok{10}\NormalTok{,}
    \AttributeTok{adapt\_delta      =} \FloatTok{0.8}
\NormalTok{  )}
\NormalTok{  fit\_gq}\SpecialCharTok{$}\FunctionTok{save\_object}\NormalTok{(}\AttributeTok{file =}\NormalTok{ model\_file\_gq)}
\NormalTok{\}}
\end{Highlighting}
\end{Shaded}

\subsection{The Mandatory MCMC Diagnostic
Battery}\label{the-mandatory-mcmc-diagnostic-battery}

Before we even look at a posterior distribution, we must prove the
geometry of the posterior was successfully explored. We have zero
tolerance for divergences. If a divergence occurs, the NUTS trajectory
has departed from the true posterior manifold, indicating a biased
approximation.

\begin{Shaded}
\begin{Highlighting}[]
\CommentTok{\# Extract the summary and diagnostics from our fitted model object}
\NormalTok{fit\_summary }\OtherTok{\textless{}{-}}\NormalTok{ fit\_gq}\SpecialCharTok{$}\FunctionTok{summary}\NormalTok{()}
\NormalTok{diagnostic\_summary }\OtherTok{\textless{}{-}}\NormalTok{ fit\_gq}\SpecialCharTok{$}\FunctionTok{diagnostic\_summary}\NormalTok{()}

\CommentTok{\# 1. Check for Divergent Transitions (Zero Tolerance)}
\NormalTok{divergences }\OtherTok{\textless{}{-}} \FunctionTok{sum}\NormalTok{(diagnostic\_summary}\SpecialCharTok{$}\NormalTok{num\_divergent)}
\ControlFlowTok{if}\NormalTok{ (divergences }\SpecialCharTok{\textgreater{}} \DecValTok{0}\NormalTok{) \{}
  \FunctionTok{stop}\NormalTok{(}\FunctionTok{paste}\NormalTok{(}\StringTok{"CRITICAL FAILURE:"}\NormalTok{, divergences, }\StringTok{"divergences detected. The posterior geometry is pathological. Reparameterization required. Do not interpret these results."}\NormalTok{))}
\NormalTok{\}}

\CommentTok{\# 2. Check E{-}BFMI (Energy{-}Bayesian Fraction of Missing Information)}
\CommentTok{\# Threshold is typically 0.3. Values below this indicate poor exploration.}
\NormalTok{min\_ebfmi }\OtherTok{\textless{}{-}} \FunctionTok{min}\NormalTok{(diagnostic\_summary}\SpecialCharTok{$}\NormalTok{ebfmi)}
\ControlFlowTok{if}\NormalTok{ (min\_ebfmi }\SpecialCharTok{\textless{}} \FloatTok{0.3}\NormalTok{) \{}
  \FunctionTok{warning}\NormalTok{(}\FunctionTok{paste}\NormalTok{(}\StringTok{"WARNING: Low E{-}BFMI detected (Min ="}\NormalTok{, }\FunctionTok{round}\NormalTok{(min\_ebfmi, }\DecValTok{3}\NormalTok{), }
                \StringTok{"). The sampler may struggle to explore the tails of the posterior."}\NormalTok{))}
\NormalTok{\}}

\CommentTok{\# 3. Check for Convergence (Rhat \textless{} 1.01)}
\NormalTok{max\_rhat }\OtherTok{\textless{}{-}} \FunctionTok{max}\NormalTok{(fit\_summary}\SpecialCharTok{$}\NormalTok{rhat, }\AttributeTok{na.rm =} \ConstantTok{TRUE}\NormalTok{)}
\ControlFlowTok{if}\NormalTok{ (max\_rhat }\SpecialCharTok{\textgreater{}} \FloatTok{1.01}\NormalTok{) \{}
  \FunctionTok{stop}\NormalTok{(}\FunctionTok{paste}\NormalTok{(}\StringTok{"CRITICAL FAILURE: Chains did not converge. Max Rhat ="}\NormalTok{, }\FunctionTok{round}\NormalTok{(max\_rhat, }\DecValTok{3}\NormalTok{)))}
\NormalTok{\}}

\CommentTok{\# 4. Check Effective Sample Size (ESS \textgreater{} 400 for both bulk and tail)}
\NormalTok{min\_ess\_bulk }\OtherTok{\textless{}{-}} \FunctionTok{min}\NormalTok{(fit\_summary}\SpecialCharTok{$}\NormalTok{ess\_bulk, }\AttributeTok{na.rm =} \ConstantTok{TRUE}\NormalTok{)}
\NormalTok{min\_ess\_tail }\OtherTok{\textless{}{-}} \FunctionTok{min}\NormalTok{(fit\_summary}\SpecialCharTok{$}\NormalTok{ess\_tail, }\AttributeTok{na.rm =} \ConstantTok{TRUE}\NormalTok{)}
\ControlFlowTok{if}\NormalTok{ (min\_ess\_bulk }\SpecialCharTok{\textless{}} \DecValTok{400} \SpecialCharTok{||}\NormalTok{ min\_ess\_tail }\SpecialCharTok{\textless{}} \DecValTok{400}\NormalTok{) \{}
  \FunctionTok{warning}\NormalTok{(}\StringTok{"WARNING: Insufficient Effective Sample Size. Autocorrelation is too high. Run longer chains."}\NormalTok{)}
\NormalTok{\}}
\end{Highlighting}
\end{Shaded}

\section{Phase 2: Generative Plausibility (Prior Predictive
Checks)}\label{phase-2-generative-plausibility-prior-predictive-checks}

\begin{quote}
\emph{Why run a prior predictive check before seeing any data?} Priors
are not just regularizers --- they are part of the generative model. A
prior that assigns high probability to θ = 0.999 (choosing ``right''
almost every single trial) might corrupt the posterior even with 1,000
observations, and/or make the MCMC much slower. Catching an implausible
prior here costs minutes; catching it throughout the fitting process
might waste you much longer.
\end{quote}

Before fitting the model, we specified prior distributions for our
parameters (e.g., \(\theta_{\text{logit}} \sim \text{Normal}(0, 1.5)\)).
Priors do not merely encode our initial uncertainty before seeing the
data; they are fundamental constraints on the data-generating process.
Prior predictive checks help us understand the implications of these
choices. If we simulate data using only parameters drawn from our
priors, what kind of data does the model expect before seeing any actual
observations?

\subsection{The Generative Algorithm:}\label{the-generative-algorithm}

To assess generative plausibility, we explicitly simulate from the joint
distribution of the model:

\begin{itemize}
\item
  Draw from the Prior: Sample \(S\) parameter values from their
  respective prior distributions. For example:
  \(\theta_{\text{logit}}^{(s)} \sim \text{Normal}(0, 1.5)\).
\item
  Transform to Constrained Space: Push these unconstrained values
  through the inverse-link function to map them to the cognitive
  parameter space (probability):
  \(\theta^{(s)} = \text{logit}^{-1}(\theta_{\text{logit}}^{(s)})\).
\item
  Simulate the Likelihood: Using these prior parameter values, simulate
  a replicated dataset from the data-generating function:
  \(y_{\text{prior}}^{(s)} \sim \text{Bernoulli}(\theta^{(s)})\).
\item
  Evaluate Summary Statistics: Calculate theoretically relevant summary
  statistics, \(T(y_{\text{prior}}^{(s)})\), and plot their
  distribution.
\end{itemize}

This protocol reveals the limits of our assumptions. We must interrogate
this space: Are we inadvertently generating boundary-saturating
behaviors? If our prior allows the agent to generate 120 ``left''
choices and 0 ``right'' choices with high frequency, we have likely
specified a prior that is too wide on the logit scale, pushing the
probability mass into the extreme tails of 0 and 1.

\subsection{Performing the Prior Predictive
Check:}\label{performing-the-prior-predictive-check}

We have already embedded the generative simulation (h\_prior\_rep) and
the summary statistic calculation (prior\_rep\_sum) directly into the
Stan generated quantities block. We will now extract these draws and
visualize the geometry of our prior predictions.

\begin{Shaded}
\begin{Highlighting}[]
\CommentTok{\# Extract the summary statistic calculated from prior predictive simulations}
\CommentTok{\# \textquotesingle{}prior\_rep\_sum\textquotesingle{} was calculated within the generated quantities block}
\NormalTok{draws\_gq }\OtherTok{\textless{}{-}} \FunctionTok{as\_draws\_df}\NormalTok{(fit\_gq}\SpecialCharTok{$}\FunctionTok{draws}\NormalTok{()) }\CommentTok{\# Ensure we have the draws as a dataframe}
\NormalTok{observed\_sum }\OtherTok{\textless{}{-}} \FunctionTok{sum}\NormalTok{(observed\_data\_list}\SpecialCharTok{$}\NormalTok{y) }\CommentTok{\# Total number of \textquotesingle{}right\textquotesingle{} choices in the actual observed data}

\NormalTok{ppc\_prior\_ggplot }\OtherTok{\textless{}{-}} \FunctionTok{ggplot}\NormalTok{(draws\_gq, }\FunctionTok{aes}\NormalTok{(}\AttributeTok{x =}\NormalTok{ prior\_rep\_sum)) }\SpecialCharTok{+}
  \CommentTok{\# Draw the histogram of the model\textquotesingle{}s prior predictions (light red)}
  \FunctionTok{geom\_histogram}\NormalTok{(}\AttributeTok{fill =} \StringTok{"\#ffbaba"}\NormalTok{, }\AttributeTok{color =} \StringTok{"\#ff5252"}\NormalTok{, }\AttributeTok{bins =} \DecValTok{30}\NormalTok{, }\AttributeTok{alpha =} \FloatTok{0.8}\NormalTok{) }\SpecialCharTok{+}
  
  \CommentTok{\# Add the vertical line for the actual observed data (dark red)}
  \FunctionTok{geom\_vline}\NormalTok{(}\AttributeTok{xintercept =}\NormalTok{ observed\_sum, }\AttributeTok{color =} \StringTok{"darkred"}\NormalTok{, }\AttributeTok{linewidth =} \FloatTok{1.2}\NormalTok{) }\SpecialCharTok{+}
  
  \CommentTok{\# Labels and theming}
  \FunctionTok{ggtitle}\NormalTok{(}\StringTok{"Prior Predictive Check: Total \textquotesingle{}Right\textquotesingle{} Choices"}\NormalTok{,}
          \AttributeTok{subtitle =} \StringTok{"Healthy = histogram spans the full range; observed value (dashed) is not in the tails"}\NormalTok{) }\SpecialCharTok{+}
  \FunctionTok{xlab}\NormalTok{(}\FunctionTok{paste}\NormalTok{(}\StringTok{"Total Number of \textquotesingle{}Right\textquotesingle{} Choices (out of"}\NormalTok{, N\_trials, }\StringTok{"trials)"}\NormalTok{)) }\SpecialCharTok{+}
  \FunctionTok{ylab}\NormalTok{(}\StringTok{"Frequency"}\NormalTok{) }\SpecialCharTok{+}
  \FunctionTok{theme\_classic}\NormalTok{()}

\NormalTok{ppc\_prior\_ggplot}
\end{Highlighting}
\end{Shaded}

\begin{center}
\includegraphics[width=0.8\linewidth,height=\textheight,keepaspectratio]{05-ModelQualityChecks_files/figure-pdf/ch06_prior_predictive_check-1.pdf}
\end{center}

\textbf{Interpretation:} The histogram covers the full range of possible
trial counts, centred near 60 (chance). This is what we want: the prior
does not pre-commit the model to an extreme outcome before seeing any
data. If the histogram were piled up near 0 or 120, that would mean our
prior is effectively deciding the answer for us --- a broken prior. The
dashed line marks our observed data. It sits well inside the bulk of the
distribution, confirming the prior does not rule out the kind of data we
actually saw.

\section{Phase 3: Posterior Predictive Checks - Does the Fitted Model
Resemble the
Data?}\label{phase-3-posterior-predictive-checks---does-the-fitted-model-resemble-the-data}

\begin{quote}
\emph{Why check posterior predictions after fitting?} A model can
produce a beautiful, well-converged posterior and still be completely
wrong about the data. The posterior predictive check asks: if we used
our fitted model to generate new datasets, would those datasets look
like the one we actually observed? If not, the model's generative
structure is misspecified --- regardless of what the posteriors say.
\end{quote}

Now that we've examined the prior implications, let's perform the
posterior predictive check. The fundamental question of model adequacy
is: Conditional on the data we observed, can the inferred
data-generating process replicate the key structural features of our
data?

\subsection{The Generative Algorithm:}\label{the-generative-algorithm-1}

The algorithm mirrors the prior predictive check, but instead of
sampling from the prior, we condition on the learned posterior
\(p(\theta \mid y)\):

\begin{itemize}
\item
  Draw from the Posterior: For each MCMC sample \(s\), take the inferred
  parameter \(\theta^{(s)}\).
\item
  Simulate Replications: Generate a synthetic dataset
  \(y_{\text{rep}}^{(s)}\) of the exact same dimensions as the empirical
  data, using the data
  model:\[y_{\text{rep}}^{(s)} \sim \text{Bernoulli}(\theta^{(s)})\]
\item
  Evaluate Test Statistics: Calculate theoretically relevant summary
  statistics \(T(y_{\text{rep}}^{(s)})\) and compare them against the
  empirical statistic \(T(y)\).
\end{itemize}

We must interrogate these predictions critically. Our model assumes
independent, identically distributed (i.i.d.) Bernoulli trials. But what
if the participant exhibited sequential dependencies, such as a
``win-stay/lose-shift'' strategy? If so, our model might successfully
capture the overall choice proportion (the sum), but it will fail
catastrophically to reproduce the sequential structure (e.g., the
maximum run-length of identical choices).

\subsection{Performing the Posterior Predictive
Check:}\label{performing-the-posterior-predictive-check}

We extract the replicated datasets (y\_post\_rep) generated
simultaneously within the Stan generated quantities block. We will then
systematically compare the test statistic distributions
\(T(y_{\text{rep}})\) against our empirical ground truth \(T(y)\).

\begin{Shaded}
\begin{Highlighting}[]
\CommentTok{\# Extract posterior predictive simulations from the fitted object}
\CommentTok{\# This matrix has dimensions: [number\_of\_draws, number\_of\_trials]}
\NormalTok{y\_post\_rep }\OtherTok{\textless{}{-}}\NormalTok{ fit\_gq}\SpecialCharTok{$}\FunctionTok{draws}\NormalTok{(}\StringTok{"y\_post\_rep"}\NormalTok{, }\AttributeTok{format =} \StringTok{"matrix"}\NormalTok{)}

\CommentTok{\# Calculate the test statistic T(y) for the *actual* empirical data}
\NormalTok{observed\_sum }\OtherTok{\textless{}{-}} \FunctionTok{sum}\NormalTok{(observed\_data\_list}\SpecialCharTok{$}\NormalTok{y) }

\CommentTok{\# Calculate the test statistic T(y\_rep) for *each* replicated dataset}
\CommentTok{\# We compute the row{-}wise sum across the posterior draws}
\NormalTok{post\_rep\_sums }\OtherTok{\textless{}{-}} \FunctionTok{rowSums}\NormalTok{(y\_post\_rep)}

\CommentTok{\# Bundle into a tibble for strictly structured plotting}
\NormalTok{ppc\_df }\OtherTok{\textless{}{-}} \FunctionTok{tibble}\NormalTok{(}
  \AttributeTok{Replicated\_Sum =}\NormalTok{ post\_rep\_sums}
\NormalTok{)}

\CommentTok{\# {-}{-}{-} Graphical Posterior Predictive Check using ggplot2 {-}{-}{-}}
\NormalTok{ppc\_posterior\_ggplot }\OtherTok{\textless{}{-}} \FunctionTok{ggplot}\NormalTok{(ppc\_df, }\FunctionTok{aes}\NormalTok{(}\AttributeTok{x =}\NormalTok{ Replicated\_Sum)) }\SpecialCharTok{+}
  \FunctionTok{geom\_histogram}\NormalTok{(}\AttributeTok{fill =} \StringTok{"skyblue"}\NormalTok{, }\AttributeTok{color =} \StringTok{"white"}\NormalTok{, }\AttributeTok{alpha =} \FloatTok{0.7}\NormalTok{, }\AttributeTok{bins =} \DecValTok{30}\NormalTok{) }\SpecialCharTok{+}
  \FunctionTok{geom\_vline}\NormalTok{(}\AttributeTok{xintercept =}\NormalTok{ observed\_sum, }\AttributeTok{color =} \StringTok{"black"}\NormalTok{, }\AttributeTok{linetype =} \StringTok{"dashed"}\NormalTok{, }\AttributeTok{linewidth =} \DecValTok{1}\NormalTok{) }\SpecialCharTok{+}
  \FunctionTok{ggtitle}\NormalTok{(}\StringTok{"Posterior Predictive Check: Total \textquotesingle{}Right\textquotesingle{} Choices"}\NormalTok{,}
          \AttributeTok{subtitle =} \StringTok{"Healthy = observed value (dashed) sits inside the bulk of replicated datasets"}\NormalTok{) }\SpecialCharTok{+}
  \FunctionTok{xlab}\NormalTok{(}\FunctionTok{paste}\NormalTok{(}\StringTok{"Total Number of \textquotesingle{}Right\textquotesingle{} Choices (out of"}\NormalTok{, N\_trials, }\StringTok{"trials)"}\NormalTok{)) }\SpecialCharTok{+}
  \FunctionTok{ylab}\NormalTok{(}\StringTok{"Frequency (Replicated Datasets)"}\NormalTok{) }\SpecialCharTok{+}
  \CommentTok{\# Add annotation for the empirical ground truth}
  \FunctionTok{annotate}\NormalTok{(}
    \StringTok{"text"}\NormalTok{, }\AttributeTok{x =}\NormalTok{ observed\_sum, }\AttributeTok{y =} \ConstantTok{Inf}\NormalTok{, }
    \AttributeTok{label =} \StringTok{"Empirical data"}\NormalTok{, }
    \AttributeTok{vjust =} \DecValTok{2}\NormalTok{, }\AttributeTok{hjust =} \ControlFlowTok{if}\NormalTok{(observed\_sum }\SpecialCharTok{\textgreater{}}\NormalTok{ N\_trials}\SpecialCharTok{/}\DecValTok{2}\NormalTok{) }\FloatTok{1.1} \ControlFlowTok{else} \SpecialCharTok{{-}}\FloatTok{0.1}\NormalTok{, }
    \AttributeTok{size =} \DecValTok{4}\NormalTok{, }\AttributeTok{fontface =} \StringTok{"bold"}
\NormalTok{  ) }\SpecialCharTok{+}
  \FunctionTok{theme\_classic}\NormalTok{()}

\NormalTok{ppc\_posterior\_ggplot}
\end{Highlighting}
\end{Shaded}

\begin{center}
\includegraphics[width=0.8\linewidth,height=\textheight,keepaspectratio]{05-ModelQualityChecks_files/figure-pdf/ch06_posterior_predictive_checks-1.pdf}
\end{center}

\begin{Shaded}
\begin{Highlighting}[]
\NormalTok{ppc\_overlay\_df }\OtherTok{\textless{}{-}} \FunctionTok{tibble}\NormalTok{(}
  \AttributeTok{Replicated\_Sum =} \FunctionTok{c}\NormalTok{(draws\_gq}\SpecialCharTok{$}\NormalTok{prior\_rep\_sum, post\_rep\_sums),}
  \AttributeTok{Distribution =} \FunctionTok{factor}\NormalTok{(}
    \FunctionTok{rep}\NormalTok{(}\FunctionTok{c}\NormalTok{(}\StringTok{"Prior Predictive"}\NormalTok{, }\StringTok{"Posterior Predictive"}\NormalTok{), }\AttributeTok{each =} \FunctionTok{length}\NormalTok{(post\_rep\_sums)),}
    \AttributeTok{levels =} \FunctionTok{c}\NormalTok{(}\StringTok{"Prior Predictive"}\NormalTok{, }\StringTok{"Posterior Predictive"}\NormalTok{)}
\NormalTok{  )}
\NormalTok{)}

\CommentTok{\# 2. Generate the overlay plot}
\NormalTok{ppc\_overlay\_ggplot }\OtherTok{\textless{}{-}} \FunctionTok{ggplot}\NormalTok{(ppc\_overlay\_df, }\FunctionTok{aes}\NormalTok{(}\AttributeTok{x =}\NormalTok{ Replicated\_Sum, }\AttributeTok{fill =}\NormalTok{ Distribution)) }\SpecialCharTok{+}
  \CommentTok{\# Use position = "identity" and alpha to ensure both distributions are visible}
  \FunctionTok{geom\_histogram}\NormalTok{(}\AttributeTok{position =} \StringTok{"identity"}\NormalTok{, }\AttributeTok{alpha =} \FloatTok{0.6}\NormalTok{, }\AttributeTok{color =} \StringTok{"white"}\NormalTok{, }\AttributeTok{bins =} \DecValTok{30}\NormalTok{) }\SpecialCharTok{+}
  \FunctionTok{geom\_vline}\NormalTok{(}\AttributeTok{xintercept =}\NormalTok{ observed\_sum, }\AttributeTok{color =} \StringTok{"black"}\NormalTok{, }\AttributeTok{linetype =} \StringTok{"dashed"}\NormalTok{, }\AttributeTok{linewidth =} \DecValTok{1}\NormalTok{) }\SpecialCharTok{+}
  \FunctionTok{scale\_fill\_manual}\NormalTok{(}\AttributeTok{values =} \FunctionTok{c}\NormalTok{(}\StringTok{"Prior Predictive"} \OtherTok{=} \StringTok{"\#ffbaba"}\NormalTok{, }\StringTok{"Posterior Predictive"} \OtherTok{=} \StringTok{"skyblue"}\NormalTok{)) }\SpecialCharTok{+}
  \FunctionTok{ggtitle}\NormalTok{(}\StringTok{"Prior vs. Posterior Predictive: Narrowing of Predictions After Fitting"}\NormalTok{,}
          \AttributeTok{subtitle =} \StringTok{"Posterior (blue) should be much tighter than prior (red) and centred on observed data"}\NormalTok{) }\SpecialCharTok{+}
  \FunctionTok{xlab}\NormalTok{(}\FunctionTok{paste}\NormalTok{(}\StringTok{"Total Number of \textquotesingle{}Right\textquotesingle{} Choices (out of"}\NormalTok{, N\_trials, }\StringTok{"trials)"}\NormalTok{)) }\SpecialCharTok{+}
  \FunctionTok{ylab}\NormalTok{(}\StringTok{"Frequency (Replicated Datasets)"}\NormalTok{) }\SpecialCharTok{+}
  \CommentTok{\# Add annotation for the empirical ground truth}
  \FunctionTok{annotate}\NormalTok{(}
    \StringTok{"text"}\NormalTok{, }\AttributeTok{x =}\NormalTok{ observed\_sum, }\AttributeTok{y =} \ConstantTok{Inf}\NormalTok{, }
    \AttributeTok{label =} \StringTok{"Empirical data"}\NormalTok{, }
    \AttributeTok{vjust =} \DecValTok{2}\NormalTok{, }\AttributeTok{hjust =} \ControlFlowTok{if}\NormalTok{(observed\_sum }\SpecialCharTok{\textgreater{}}\NormalTok{ N\_trials}\SpecialCharTok{/}\DecValTok{2}\NormalTok{) }\FloatTok{1.1} \ControlFlowTok{else} \SpecialCharTok{{-}}\FloatTok{0.1}\NormalTok{, }
    \AttributeTok{size =} \DecValTok{4}\NormalTok{, }\AttributeTok{fontface =} \StringTok{"bold"}
\NormalTok{  ) }\SpecialCharTok{+}
  \FunctionTok{theme\_classic}\NormalTok{() }\SpecialCharTok{+}
  \FunctionTok{theme}\NormalTok{(}
    \AttributeTok{legend.position =} \StringTok{"bottom"}\NormalTok{, }
    \AttributeTok{legend.title =} \FunctionTok{element\_blank}\NormalTok{(),}
    \AttributeTok{legend.text =} \FunctionTok{element\_text}\NormalTok{(}\AttributeTok{size =} \DecValTok{12}\NormalTok{)}
\NormalTok{  )}

\FunctionTok{print}\NormalTok{(ppc\_overlay\_ggplot)}
\end{Highlighting}
\end{Shaded}

\begin{center}
\includegraphics[width=0.8\linewidth,height=\textheight,keepaspectratio]{05-ModelQualityChecks_files/figure-pdf/ch06_posterior_predictive_checks-2.pdf}
\end{center}

\begin{Shaded}
\begin{Highlighting}[]
\CommentTok{\# {-}{-}{-} Example of checking another statistic (e.g., max run length) {-}{-}{-}}
\CommentTok{\# Define a function to calculate the maximum run length in a vector}
\NormalTok{max\_run\_length }\OtherTok{\textless{}{-}} \ControlFlowTok{function}\NormalTok{(x) \{}
  \ControlFlowTok{if}\NormalTok{ (}\FunctionTok{length}\NormalTok{(x) }\SpecialCharTok{==} \DecValTok{0} \SpecialCharTok{||} \FunctionTok{all}\NormalTok{(}\FunctionTok{is.na}\NormalTok{(x))) }\FunctionTok{return}\NormalTok{(}\DecValTok{0}\NormalTok{) }
\NormalTok{  rle\_x }\OtherTok{\textless{}{-}} \FunctionTok{rle}\NormalTok{(x) }\CommentTok{\# Run Length Encoding calculates consecutive identical choices}
  \ControlFlowTok{if}\NormalTok{ (}\FunctionTok{length}\NormalTok{(rle\_x}\SpecialCharTok{$}\NormalTok{lengths) }\SpecialCharTok{==} \DecValTok{0}\NormalTok{) }\FunctionTok{return}\NormalTok{(}\DecValTok{0}\NormalTok{) }
  \FunctionTok{max}\NormalTok{(rle\_x}\SpecialCharTok{$}\NormalTok{lengths) }
\NormalTok{\}}

\CommentTok{\# 2. Calculate the empirical test statistic T(y)}
\NormalTok{observed\_max\_run }\OtherTok{\textless{}{-}} \FunctionTok{max\_run\_length}\NormalTok{(observed\_data\_list}\SpecialCharTok{$}\NormalTok{y)}

\CommentTok{\# 3. Extract prior replicated matrices and compute T(y\_rep) for the prior}
\CommentTok{\# y\_prior\_rep is an S x N matrix generated from the prior distributions}
\NormalTok{y\_prior\_rep }\OtherTok{\textless{}{-}}\NormalTok{ fit\_gq}\SpecialCharTok{$}\FunctionTok{draws}\NormalTok{(}\StringTok{"y\_prior\_rep"}\NormalTok{, }\AttributeTok{format =} \StringTok{"matrix"}\NormalTok{)}
\NormalTok{prior\_max\_runs }\OtherTok{\textless{}{-}} \FunctionTok{apply}\NormalTok{(y\_prior\_rep, }\DecValTok{1}\NormalTok{, max\_run\_length)}

\CommentTok{\# 4. Extract posterior replicated matrices and compute T(y\_rep) for the posterior}
\CommentTok{\# y\_post\_rep is an S x N matrix generated from the posterior distributions}
\NormalTok{y\_post\_rep }\OtherTok{\textless{}{-}}\NormalTok{ fit\_gq}\SpecialCharTok{$}\FunctionTok{draws}\NormalTok{(}\StringTok{"y\_post\_rep"}\NormalTok{, }\AttributeTok{format =} \StringTok{"matrix"}\NormalTok{)}
\NormalTok{post\_max\_runs }\OtherTok{\textless{}{-}} \FunctionTok{apply}\NormalTok{(y\_post\_rep, }\DecValTok{1}\NormalTok{, max\_run\_length)}

\CommentTok{\# 5. Bundle into a strictly structured tibble for overlay plotting}
\NormalTok{ppc\_run\_df }\OtherTok{\textless{}{-}} \FunctionTok{tibble}\NormalTok{(}
  \AttributeTok{Replicated\_Max\_Run =} \FunctionTok{c}\NormalTok{(prior\_max\_runs, post\_max\_runs),}
  \AttributeTok{Distribution =} \FunctionTok{factor}\NormalTok{(}
    \FunctionTok{rep}\NormalTok{(}\FunctionTok{c}\NormalTok{(}\StringTok{"Prior Predictive"}\NormalTok{, }\StringTok{"Posterior Predictive"}\NormalTok{), }\AttributeTok{each =} \FunctionTok{length}\NormalTok{(post\_max\_runs)),}
    \AttributeTok{levels =} \FunctionTok{c}\NormalTok{(}\StringTok{"Prior Predictive"}\NormalTok{, }\StringTok{"Posterior Predictive"}\NormalTok{)}
\NormalTok{  )}
\NormalTok{)}

\CommentTok{\# 6. Generate the overlay plot}
\NormalTok{p\_ppc\_max\_run }\OtherTok{\textless{}{-}} \FunctionTok{ggplot}\NormalTok{(ppc\_run\_df, }\FunctionTok{aes}\NormalTok{(}\AttributeTok{x =}\NormalTok{ Replicated\_Max\_Run, }\AttributeTok{fill =}\NormalTok{ Distribution)) }\SpecialCharTok{+}
  \FunctionTok{geom\_histogram}\NormalTok{(}\AttributeTok{position =} \StringTok{"identity"}\NormalTok{, }\AttributeTok{alpha =} \FloatTok{0.6}\NormalTok{, }\AttributeTok{color =} \StringTok{"white"}\NormalTok{, }\AttributeTok{bins =} \DecValTok{20}\NormalTok{) }\SpecialCharTok{+}
  \FunctionTok{geom\_vline}\NormalTok{(}\AttributeTok{xintercept =}\NormalTok{ observed\_max\_run, }\AttributeTok{color =} \StringTok{"black"}\NormalTok{, }\AttributeTok{linetype =} \StringTok{"dashed"}\NormalTok{, }\AttributeTok{linewidth =} \DecValTok{1}\NormalTok{) }\SpecialCharTok{+}
  \FunctionTok{scale\_fill\_manual}\NormalTok{(}\AttributeTok{values =} \FunctionTok{c}\NormalTok{(}\StringTok{"Prior Predictive"} \OtherTok{=} \StringTok{"\#ffbaba"}\NormalTok{, }\StringTok{"Posterior Predictive"} \OtherTok{=} \StringTok{"skyblue"}\NormalTok{)) }\SpecialCharTok{+}
  \FunctionTok{ggtitle}\NormalTok{(}\StringTok{"Posterior Predictive: Maximum Run Length of Consecutive Choices"}\NormalTok{,}
          \AttributeTok{subtitle =} \StringTok{"Tests sequential structure — a model assuming i.i.d. trials should still pass this"}\NormalTok{) }\SpecialCharTok{+}
  \FunctionTok{xlab}\NormalTok{(}\StringTok{"Test Statistic T(y\_rep): Maximum Run Length of Consecutive Choices"}\NormalTok{) }\SpecialCharTok{+}
  \FunctionTok{ylab}\NormalTok{(}\StringTok{"Frequency (Replicated Datasets)"}\NormalTok{) }\SpecialCharTok{+}
  \CommentTok{\# Add annotation for the empirical ground truth}
  \FunctionTok{annotate}\NormalTok{(}
    \StringTok{"text"}\NormalTok{, }\AttributeTok{x =}\NormalTok{ observed\_max\_run, }\AttributeTok{y =} \ConstantTok{Inf}\NormalTok{, }
    \AttributeTok{label =} \StringTok{"Empirical T(y)"}\NormalTok{, }
    \AttributeTok{vjust =} \DecValTok{2}\NormalTok{, }\AttributeTok{hjust =} \SpecialCharTok{{-}}\FloatTok{0.1}\NormalTok{, }
    \AttributeTok{size =} \DecValTok{4}\NormalTok{, }\AttributeTok{fontface =} \StringTok{"bold"}
\NormalTok{  ) }\SpecialCharTok{+}
  \FunctionTok{theme\_classic}\NormalTok{() }\SpecialCharTok{+}
  \FunctionTok{theme}\NormalTok{(}
    \AttributeTok{legend.position =} \StringTok{"bottom"}\NormalTok{, }
    \AttributeTok{legend.title =} \FunctionTok{element\_blank}\NormalTok{(),}
    \AttributeTok{legend.text =} \FunctionTok{element\_text}\NormalTok{(}\AttributeTok{size =} \DecValTok{12}\NormalTok{)}
\NormalTok{  )}

\FunctionTok{print}\NormalTok{(p\_ppc\_max\_run)}
\end{Highlighting}
\end{Shaded}

\begin{center}
\includegraphics[width=0.8\linewidth,height=\textheight,keepaspectratio]{05-ModelQualityChecks_files/figure-pdf/ch06_posterior_predictive_checks-3.pdf}
\end{center}

\textbf{Interpretation:} Both test statistics --- the total count of
``right'' choices and the maximum run length --- fall comfortably inside
their posterior predictive distributions. This means the model passes on
two fronts: it captures the average tendency (how often the agent picks
``right'') \emph{and} the sequential texture (how long the agent stays
on the same choice).

The maximum run length test is the more important one here. Our model
assumes each trial is independent --- it has no memory of what came
before. If the data showed very long runs (suggesting a
win-stay/lose-shift strategy), the model would fail on that statistic
even while passing on the count. It does not fail. This confirms that an
i.i.d. Bernoulli model is a defensible description of this particular
dataset.

\textbf{So what?} We have passed two out of six phases. The model is
plausibly specified and fits the data it was trained on. Phase 4 now
asks whether the data were actually informative enough to move our
beliefs.

\section{Phase 4: Epistemic Updating (Prior-Posterior Update Checks) -
How Much Did the Data Change our
Beliefs?}\label{phase-4-epistemic-updating-prior-posterior-update-checks---how-much-did-the-data-change-our-beliefs}

\begin{quote}
\emph{Why plot the prior and posterior together?} Two things can go
wrong that only this plot reveals: (1) the posterior as uncertain
(broad) as the prior, or even more, meaning that it wasn't able to learn
from the data; (2) the posterior is far from the true value and narrow
--- meaning the model is confidently wrong. Neither failure is visible
from the posterior alone.
\end{quote}

This check visualizes the core engine of Bayesian inference: the formal
mathematical updating of our structural assumptions by the empirical
evidence. Conceptually, we are observing Bayes' theorem in action as we
transition from the prior \(p(\theta)\) to the posterior
\(p(\theta \mid y) \propto p(y \mid \theta) p(\theta)\).

By explicitly plotting the marginal prior distribution against the
marginal posterior distribution for our parameters, we must interrogate
two critical epistemic properties:

\begin{itemize}
\item
  \textbf{Information Gain}: How aggressively did the likelihood
  \(p(y \mid \theta)\) compress the variance of the prior? If the
  posterior distribution is substantially narrower (higher precision)
  than the prior, the empirical data possessed high informational value
  and strongly constrained our parameter estimates.
\item
  \textbf{Prior Influence}: Did the central tendency of the posterior
  shift significantly away from the prior mean? If the posterior sort of
  mirrors the prior, we must ask a deep structural question: Did the
  data provide too little informational value, or was our prior so
  strong that it overwhelmed the likelihood?
\end{itemize}

We will inspect this transformation on both the unconstrained space
defined by our model (\(\theta_{\text{logit}} \in \mathbb{R}\)) and the
conceptually intuitive probability space (\(\theta \in [0, 1]\)).

\begin{Shaded}
\begin{Highlighting}[]
\CommentTok{\# Extract MCMC draws into a structured dataframe}
\NormalTok{draws\_gq }\OtherTok{\textless{}{-}} \FunctionTok{as\_draws\_df}\NormalTok{(fit\_gq}\SpecialCharTok{$}\FunctionTok{draws}\NormalTok{())}

\CommentTok{\# 2. Define the empirical ground truth for reference}
\NormalTok{true\_theta\_prob }\OtherTok{\textless{}{-}} \FloatTok{0.8}
\NormalTok{true\_theta\_logit }\OtherTok{\textless{}{-}} \FunctionTok{qlogis}\NormalTok{(true\_theta\_prob)}

\CommentTok{\# {-}{-}{-} The Computational (Unconstrained) Space {-}{-}{-}}

\NormalTok{update\_logit\_df }\OtherTok{\textless{}{-}} \FunctionTok{tibble}\NormalTok{(}
  \AttributeTok{Parameter\_Value =} \FunctionTok{c}\NormalTok{(draws\_gq}\SpecialCharTok{$}\NormalTok{theta\_logit\_prior, draws\_gq}\SpecialCharTok{$}\NormalTok{theta\_logit),}
  \AttributeTok{Distribution =} \FunctionTok{factor}\NormalTok{(}
    \FunctionTok{rep}\NormalTok{(}\FunctionTok{c}\NormalTok{(}\StringTok{"Prior"}\NormalTok{, }\StringTok{"Posterior"}\NormalTok{), }\AttributeTok{each =} \FunctionTok{nrow}\NormalTok{(draws\_gq)),}
    \AttributeTok{levels =} \FunctionTok{c}\NormalTok{(}\StringTok{"Prior"}\NormalTok{, }\StringTok{"Posterior"}\NormalTok{)}
\NormalTok{  )}
\NormalTok{)}

\CommentTok{\# Generate the overlay density plot for the logit space}
\NormalTok{p\_update\_logit }\OtherTok{\textless{}{-}} \FunctionTok{ggplot}\NormalTok{(update\_logit\_df, }\FunctionTok{aes}\NormalTok{(}\AttributeTok{x =}\NormalTok{ Parameter\_Value, }\AttributeTok{fill =}\NormalTok{ Distribution)) }\SpecialCharTok{+}
  \FunctionTok{geom\_density}\NormalTok{(}\AttributeTok{alpha =} \FloatTok{0.6}\NormalTok{, }\AttributeTok{color =} \StringTok{"white"}\NormalTok{) }\SpecialCharTok{+}
  \FunctionTok{geom\_vline}\NormalTok{(}\AttributeTok{xintercept =}\NormalTok{ true\_theta\_logit, }\AttributeTok{color =} \StringTok{"black"}\NormalTok{, }\AttributeTok{linetype =} \StringTok{"dashed"}\NormalTok{, }\AttributeTok{linewidth =} \DecValTok{1}\NormalTok{) }\SpecialCharTok{+}
  \FunctionTok{scale\_fill\_manual}\NormalTok{(}\AttributeTok{values =} \FunctionTok{c}\NormalTok{(}\StringTok{"Prior"} \OtherTok{=} \StringTok{"\#ffbaba"}\NormalTok{, }\StringTok{"Posterior"} \OtherTok{=} \StringTok{"skyblue"}\NormalTok{)) }\SpecialCharTok{+}
  \FunctionTok{ggtitle}\NormalTok{(}\StringTok{"Prior{-}Posterior Update: Log{-}Odds Scale"}\NormalTok{,}
          \AttributeTok{subtitle =} \StringTok{"Healthy = posterior (blue) narrower than prior (red), shifted toward true value (dashed)"}\NormalTok{) }\SpecialCharTok{+}
  \FunctionTok{xlab}\NormalTok{(}\FunctionTok{expression}\NormalTok{(theta[logit] }\SpecialCharTok{\textasciitilde{}} \StringTok{" (Log{-}Odds)"}\NormalTok{)) }\SpecialCharTok{+}
  \FunctionTok{ylab}\NormalTok{(}\StringTok{"Marginal Density"}\NormalTok{) }\SpecialCharTok{+}
  \FunctionTok{annotate}\NormalTok{(}
    \StringTok{"text"}\NormalTok{, }\AttributeTok{x =}\NormalTok{ true\_theta\_logit, }\AttributeTok{y =} \ConstantTok{Inf}\NormalTok{, }
    \AttributeTok{label =} \StringTok{"Ground Truth"}\NormalTok{, }
    \AttributeTok{vjust =} \DecValTok{2}\NormalTok{, }\AttributeTok{hjust =} \ControlFlowTok{if}\NormalTok{ (true\_theta\_logit }\SpecialCharTok{\textgreater{}} \DecValTok{0}\NormalTok{) }\FloatTok{1.1} \ControlFlowTok{else} \SpecialCharTok{{-}}\FloatTok{0.1}\NormalTok{, }
    \AttributeTok{size =} \DecValTok{4}\NormalTok{, }\AttributeTok{fontface =} \StringTok{"bold"}
\NormalTok{  ) }\SpecialCharTok{+}
  \FunctionTok{theme\_classic}\NormalTok{() }\SpecialCharTok{+}
  \FunctionTok{theme}\NormalTok{(}
    \AttributeTok{legend.position =} \StringTok{"bottom"}\NormalTok{, }
    \AttributeTok{legend.title =} \FunctionTok{element\_blank}\NormalTok{(),}
    \AttributeTok{legend.text =} \FunctionTok{element\_text}\NormalTok{(}\AttributeTok{size =} \DecValTok{12}\NormalTok{)}
\NormalTok{  )}

\FunctionTok{print}\NormalTok{(p\_update\_logit)}
\end{Highlighting}
\end{Shaded}

\begin{center}
\includegraphics[width=0.8\linewidth,height=\textheight,keepaspectratio]{05-ModelQualityChecks_files/figure-pdf/ch06_prior_posterior_update-1.pdf}
\end{center}

\begin{Shaded}
\begin{Highlighting}[]
\CommentTok{\# {-}{-}{-} The Cognitive (Constrained) Space {-}{-}{-}}

\CommentTok{\# We use plogis() to push the unconstrained prior simulations through the inverse{-}link function.}
\NormalTok{update\_prob\_df }\OtherTok{\textless{}{-}} \FunctionTok{tibble}\NormalTok{(}
  \AttributeTok{Parameter\_Value =} \FunctionTok{c}\NormalTok{(}\FunctionTok{plogis}\NormalTok{(draws\_gq}\SpecialCharTok{$}\NormalTok{theta\_logit\_prior), draws\_gq}\SpecialCharTok{$}\NormalTok{theta),}
  \AttributeTok{Distribution =} \FunctionTok{factor}\NormalTok{(}
    \FunctionTok{rep}\NormalTok{(}\FunctionTok{c}\NormalTok{(}\StringTok{"Prior"}\NormalTok{, }\StringTok{"Posterior"}\NormalTok{), }\AttributeTok{each =} \FunctionTok{nrow}\NormalTok{(draws\_gq)),}
    \AttributeTok{levels =} \FunctionTok{c}\NormalTok{(}\StringTok{"Prior"}\NormalTok{, }\StringTok{"Posterior"}\NormalTok{)}
\NormalTok{  )}
\NormalTok{)}

\CommentTok{\# Generate the overlay density plot for the probability space}
\NormalTok{p\_update\_prob }\OtherTok{\textless{}{-}} \FunctionTok{ggplot}\NormalTok{(update\_prob\_df, }\FunctionTok{aes}\NormalTok{(}\AttributeTok{x =}\NormalTok{ Parameter\_Value, }\AttributeTok{fill =}\NormalTok{ Distribution)) }\SpecialCharTok{+}
  \FunctionTok{geom\_density}\NormalTok{(}\AttributeTok{alpha =} \FloatTok{0.6}\NormalTok{, }\AttributeTok{color =} \StringTok{"white"}\NormalTok{) }\SpecialCharTok{+}
  \FunctionTok{geom\_vline}\NormalTok{(}\AttributeTok{xintercept =}\NormalTok{ true\_theta\_prob, }\AttributeTok{color =} \StringTok{"black"}\NormalTok{, }\AttributeTok{linetype =} \StringTok{"dashed"}\NormalTok{, }\AttributeTok{linewidth =} \DecValTok{1}\NormalTok{) }\SpecialCharTok{+}
  \FunctionTok{scale\_fill\_manual}\NormalTok{(}\AttributeTok{values =} \FunctionTok{c}\NormalTok{(}\StringTok{"Prior"} \OtherTok{=} \StringTok{"\#ffbaba"}\NormalTok{, }\StringTok{"Posterior"} \OtherTok{=} \StringTok{"skyblue"}\NormalTok{)) }\SpecialCharTok{+}
  \FunctionTok{coord\_cartesian}\NormalTok{(}\AttributeTok{xlim =} \FunctionTok{c}\NormalTok{(}\DecValTok{0}\NormalTok{, }\DecValTok{1}\NormalTok{)) }\SpecialCharTok{+} \CommentTok{\# Mathematically enforce the [0, 1] bounds}
  \FunctionTok{ggtitle}\NormalTok{(}\StringTok{"Prior{-}Posterior Update: Probability Scale"}\NormalTok{,}
          \AttributeTok{subtitle =} \StringTok{"Healthy = posterior narrower and closer to true θ = 0.8 than the flat prior"}\NormalTok{) }\SpecialCharTok{+}
  \FunctionTok{xlab}\NormalTok{(}\FunctionTok{expression}\NormalTok{(theta }\SpecialCharTok{\textasciitilde{}} \StringTok{" (Probability of \textquotesingle{}Right\textquotesingle{} Choice)"}\NormalTok{)) }\SpecialCharTok{+}
  \FunctionTok{ylab}\NormalTok{(}\StringTok{"Marginal Density"}\NormalTok{) }\SpecialCharTok{+}
  \FunctionTok{annotate}\NormalTok{(}
    \StringTok{"text"}\NormalTok{, }\AttributeTok{x =}\NormalTok{ true\_theta\_prob, }\AttributeTok{y =} \ConstantTok{Inf}\NormalTok{, }
    \AttributeTok{label =} \StringTok{"Ground Truth"}\NormalTok{, }
    \AttributeTok{vjust =} \DecValTok{2}\NormalTok{, }\AttributeTok{hjust =} \ControlFlowTok{if}\NormalTok{ (true\_theta\_prob }\SpecialCharTok{\textgreater{}} \FloatTok{0.5}\NormalTok{) }\FloatTok{1.1} \ControlFlowTok{else} \SpecialCharTok{{-}}\FloatTok{0.1}\NormalTok{, }
    \AttributeTok{size =} \DecValTok{4}\NormalTok{, }\AttributeTok{fontface =} \StringTok{"bold"}
\NormalTok{  ) }\SpecialCharTok{+}
  \FunctionTok{theme\_classic}\NormalTok{() }\SpecialCharTok{+}
  \FunctionTok{theme}\NormalTok{(}
    \AttributeTok{legend.position =} \StringTok{"bottom"}\NormalTok{, }
    \AttributeTok{legend.title =} \FunctionTok{element\_blank}\NormalTok{(),}
    \AttributeTok{legend.text =} \FunctionTok{element\_text}\NormalTok{(}\AttributeTok{size =} \DecValTok{12}\NormalTok{)}
\NormalTok{  )}

\FunctionTok{print}\NormalTok{(p\_update\_prob)}
\end{Highlighting}
\end{Shaded}

\begin{center}
\includegraphics[width=0.8\linewidth,height=\textheight,keepaspectratio]{05-ModelQualityChecks_files/figure-pdf/ch06_prior_posterior_update-2.pdf}
\end{center}

\textbf{Interpretation:} The posterior (blue) is substantially narrower
than the prior (red) and has moved toward the true value of 0.8 (dashed
line). The data were informative: 120 trials were enough to meaningfully
constrain our estimate. If the posterior had been identical to the
prior, the data would have been uninformative for this parameter --- a
sign that the model structure or the experimental design is broken. If
the posterior were narrow but centred far from the true value, the model
would be confidently wrong. Neither failure occurred here. Phase 5 now
tests whether this conclusion is sensitive to our specific choice of
prior.

Now that we have a better feel for the model, it is high time to go
beyond one dataset and one prior, and be systematic about testing the
model across a wide range of parameter values and prior specifications.
We will first assess the sensitivity of our model to prior values to
then proceed onto parameter recovery.

\section{Phase 5: Prior Sensitivity Analysis - Do Priors Matter Too Too
Much?}\label{phase-5-prior-sensitivity-analysis---do-priors-matter-too-too-much}

\begin{quote}
\emph{Why test sensitivity to prior choice?} We chose
\texttt{Normal(0,\ 1.5)} for the logit-scale prior, but that was a
judgment call. If a researcher using \texttt{Normal(0,\ 0.5)} instead
gets a completely different posterior, we need to more carefully discuss
and motivate the assumptions. Sensitivity analysis quantifies exactly
how much the prior choice matters.
\end{quote}

The next validation check addresses a critical epistemic vulnerability:
Are our structural conclusions pathologically dependent on the specific
mathematical bounds of our chosen prior? If we had specified a tighter
regularizing prior or a completely diffuse one, would our estimate of
the cognitive parameter \(\theta\) change drastically? We must
explicitly investigate the robustness of our inference by perturbing the
prior space and observing the resulting posterior geometry.

\subsection{5a. Empirical Prior Perturbation (Manual Sensitivity
Check)}\label{a.-empirical-prior-perturbation-manual-sensitivity-check}

\textbf{The Mathematical Algorithm}:

\begin{enumerate}
\def\labelenumi{\arabic{enumi}.}
\item
  \emph{Define a Perturbation Grid}: Specify a set of alternative,
  scientifically plausible prior distributions for the parameter of
  interest. For example, if our baseline is
  \(\theta_{\text{logit}} \sim \text{Normal}(0, 1.0)\), we might test a
  highly restrictive prior, \(\text{Normal}(0, 0.5)\), and a diffuse
  prior, \(\text{Normal}(0, 2.0)\).
\item
  \emph{Re-compute the Posterior}: Fit the generative model to the data
  under each alternative prior specification to obtain
  \(p(\theta \mid y, \text{prior}_k)\).
\item
  \emph{Evaluate Invariance}: Overlay and compare the resulting marginal
  posterior distributions.
\end{enumerate}

\textbf{Decision Rule}: If the marginal posteriors are mathematically
invariant (nearly identical) across the perturbation grid, our inference
is computationally robust. The likelihood function contains sufficient
information to overwhelm the prior constraints. Conversely, if the
posteriors shift drastically, our inference is fragile. A fragile model
demands that we either mathematically justify our exact prior
specification, acknowledge the severe uncertainty in our conclusions, or
collect a larger empirical sample. However, it's also important to
remember that in many contexts influential priors are a good choice: to
regularize inference, to incorporate strong domain knowledge and/or
assess the distance between our results and the previous literature,
etc. See this talk for more:
https://youtu.be/kMSrRd4f2As?si=-7bbx3wFWobgxebQ

\subsubsection{Computational
Implementation:}\label{computational-implementation}

To execute this efficiently, we have modified our Stan model to accept
the prior's mean and standard deviation directly via the data block.
This parameterization allows us to loop through our grid of alternative
\(\sigma\) values in R, iteratively updating the posterior without
forcing a costly C++ recompilation of the model at each step.

\begin{Shaded}
\begin{Highlighting}[]
\CommentTok{\# We test 3 location parameters (means) and 10 scale parameters (SDs)}
\NormalTok{prior\_grid }\OtherTok{\textless{}{-}}\NormalTok{ tidyr}\SpecialCharTok{::}\FunctionTok{expand\_grid}\NormalTok{(}
  \AttributeTok{prior\_mu =} \FunctionTok{c}\NormalTok{(}\SpecialCharTok{{-}}\FloatTok{1.5}\NormalTok{, }\DecValTok{0}\NormalTok{, }\FloatTok{1.5}\NormalTok{), }
  \AttributeTok{prior\_sigma =} \FunctionTok{seq}\NormalTok{(}\FloatTok{0.5}\NormalTok{, }\FloatTok{3.0}\NormalTok{, }\AttributeTok{length.out =} \DecValTok{10}\NormalTok{)}
\NormalTok{)}

\CommentTok{\# Standardized filename variable}
\NormalTok{sensitivity\_draws\_file }\OtherTok{\textless{}{-}} \FunctionTok{file.path}\NormalTok{(}\FunctionTok{here}\NormalTok{(stan\_results\_dir), }\StringTok{"ch07\_sensitivity\_draws.rds"}\NormalTok{)}

\CommentTok{\# Check if we need to re{-}run the sensitivity analysis (fixed variable name)}
\ControlFlowTok{if}\NormalTok{ (regenerate\_simulations }\SpecialCharTok{||} \SpecialCharTok{!}\FunctionTok{file.exists}\NormalTok{(sensitivity\_draws\_file)) \{}
  
  \CommentTok{\# Removed \textquotesingle{}seed\_val\textquotesingle{} from the function arguments}
\NormalTok{  fit\_and\_extract }\OtherTok{\textless{}{-}} \ControlFlowTok{function}\NormalTok{(prior\_mu, prior\_sigma) \{}
    \CommentTok{\# Dynamically inject the perturbed parameters into the data list}
\NormalTok{    current\_data\_list }\OtherTok{\textless{}{-}}\NormalTok{ observed\_data\_list}
\NormalTok{    current\_data\_list}\SpecialCharTok{$}\NormalTok{prior\_theta\_mu }\OtherTok{\textless{}{-}}\NormalTok{ prior\_mu}
\NormalTok{    current\_data\_list}\SpecialCharTok{$}\NormalTok{prior\_theta\_sigma }\OtherTok{\textless{}{-}}\NormalTok{ prior\_sigma }

    \CommentTok{\# Fit using the original compiled model}
\NormalTok{    fit }\OtherTok{\textless{}{-}}\NormalTok{ mod\_gq}\SpecialCharTok{$}\FunctionTok{sample}\NormalTok{(}
      \AttributeTok{data =}\NormalTok{ current\_data\_list,}
      \AttributeTok{chains =} \DecValTok{1}\NormalTok{,}
      \AttributeTok{parallel\_chains =} \DecValTok{1}\NormalTok{, }
      \AttributeTok{iter\_warmup =} \DecValTok{1000}\NormalTok{,}
      \AttributeTok{iter\_sampling =} \DecValTok{3000}\NormalTok{,}
      \AttributeTok{refresh =} \DecValTok{0}\NormalTok{, }
      \AttributeTok{adapt\_delta =} \FloatTok{0.8}\NormalTok{,}
      \AttributeTok{show\_messages =} \ConstantTok{FALSE}\NormalTok{,}
      \AttributeTok{show\_exceptions =} \ConstantTok{FALSE} \CommentTok{\# Added to prevent C++ errors from polluting stdout}
\NormalTok{    )}
    
\NormalTok{    draws\_df }\OtherTok{\textless{}{-}}\NormalTok{ fit}\SpecialCharTok{$}\FunctionTok{draws}\NormalTok{(}\StringTok{"theta"}\NormalTok{, }\AttributeTok{format =} \StringTok{"df"}\NormalTok{)}
    
    \CommentTok{\# Return a strictly structured tibble, discarding the fragile fit object}
\NormalTok{    tibble}\SpecialCharTok{::}\FunctionTok{tibble}\NormalTok{(}
      \AttributeTok{theta\_prob =}\NormalTok{ draws\_df}\SpecialCharTok{$}\NormalTok{theta,}
      \AttributeTok{prior\_mu =}\NormalTok{ prior\_mu,}
      \AttributeTok{prior\_sigma =}\NormalTok{ prior\_sigma}
\NormalTok{    )}
\NormalTok{  \}}

  \CommentTok{\# Establish the parallel backend }
\NormalTok{  future}\SpecialCharTok{::}\FunctionTok{plan}\NormalTok{(future}\SpecialCharTok{::}\NormalTok{multisession, }\AttributeTok{workers =} \DecValTok{2}\NormalTok{)}
  
  \CommentTok{\# Execute the parallel mapping}
\NormalTok{  draws\_list }\OtherTok{\textless{}{-}}\NormalTok{ furrr}\SpecialCharTok{::}\FunctionTok{future\_pmap}\NormalTok{(}
\NormalTok{    prior\_grid,}
\NormalTok{    fit\_and\_extract,}
    \AttributeTok{.options =}\NormalTok{ furrr}\SpecialCharTok{::}\FunctionTok{furrr\_options}\NormalTok{(}\AttributeTok{seed =} \ConstantTok{TRUE}\NormalTok{, }\AttributeTok{stdout =} \ConstantTok{FALSE}\NormalTok{) }
\NormalTok{  )}
  
  \CommentTok{\# Collapse the list of dataframes into a single tidy dataframe}
\NormalTok{  sensitivity\_draws }\OtherTok{\textless{}{-}}\NormalTok{ dplyr}\SpecialCharTok{::}\FunctionTok{bind\_rows}\NormalTok{(draws\_list)}
  
  \CommentTok{\# Shut down the parallel backend to free up RAM}
\NormalTok{  future}\SpecialCharTok{::}\FunctionTok{plan}\NormalTok{(future}\SpecialCharTok{::}\NormalTok{sequential)}
  
  \CommentTok{\# Save the indestructible, pure data (fixed variable name)}
  \FunctionTok{saveRDS}\NormalTok{(sensitivity\_draws, sensitivity\_draws\_file)}
  
\NormalTok{\} }\ControlFlowTok{else}\NormalTok{ \{}
  \CommentTok{\# Fixed variable names to match what was saved}
\NormalTok{  sensitivity\_draws }\OtherTok{\textless{}{-}} \FunctionTok{readRDS}\NormalTok{(sensitivity\_draws\_file)}
  \FunctionTok{cat}\NormalTok{(}\StringTok{"Loaded existing perturbation grid from:"}\NormalTok{, sensitivity\_draws\_file, }\StringTok{"}\SpecialCharTok{\textbackslash{}n}\StringTok{"}\NormalTok{)}
\NormalTok{\}}
\end{Highlighting}
\end{Shaded}

\begin{verbatim}
Loaded existing perturbation grid from: /Users/au209589/Dropbox/Teaching/AdvancedCognitiveModeling23_book/simmodels/ch07_sensitivity_draws.rds 
\end{verbatim}

\begin{Shaded}
\begin{Highlighting}[]
\NormalTok{p\_sensitivity\_2d }\OtherTok{\textless{}{-}} \FunctionTok{ggplot}\NormalTok{(sensitivity\_draws, }\FunctionTok{aes}\NormalTok{(}\AttributeTok{x =}\NormalTok{ theta\_prob, }\AttributeTok{group =} \FunctionTok{factor}\NormalTok{(prior\_sigma), }\AttributeTok{color =}\NormalTok{ prior\_sigma)) }\SpecialCharTok{+}
  \FunctionTok{geom\_line}\NormalTok{(}\AttributeTok{stat =} \StringTok{"density"}\NormalTok{, }\AttributeTok{alpha =} \FloatTok{0.8}\NormalTok{, }\AttributeTok{linewidth =} \FloatTok{0.8}\NormalTok{) }\SpecialCharTok{+} 
  \FunctionTok{geom\_vline}\NormalTok{(}\AttributeTok{xintercept =}\NormalTok{ true\_theta\_prob, }\AttributeTok{color =} \StringTok{"black"}\NormalTok{, }\AttributeTok{linetype =} \StringTok{"dashed"}\NormalTok{, }\AttributeTok{linewidth =} \DecValTok{1}\NormalTok{) }\SpecialCharTok{+} 
  \FunctionTok{facet\_wrap}\NormalTok{(}\SpecialCharTok{\textasciitilde{}}\NormalTok{ prior\_mu, }\AttributeTok{labeller =} \FunctionTok{labeller}\NormalTok{(}\AttributeTok{prior\_mu =} \ControlFlowTok{function}\NormalTok{(x) }\FunctionTok{paste0}\NormalTok{(}\StringTok{"Prior Mean (Logit): "}\NormalTok{, x))) }\SpecialCharTok{+}
  \FunctionTok{scale\_color\_viridis\_c}\NormalTok{(}\AttributeTok{option =} \StringTok{"mako"}\NormalTok{, }\AttributeTok{name =} \FunctionTok{expression}\NormalTok{(Prior }\SpecialCharTok{\textasciitilde{}}\NormalTok{ sigma)) }\SpecialCharTok{+} 
  \FunctionTok{coord\_cartesian}\NormalTok{(}\AttributeTok{xlim =} \FunctionTok{c}\NormalTok{(}\FloatTok{0.4}\NormalTok{, }\FloatTok{1.0}\NormalTok{)) }\SpecialCharTok{+} 
  \FunctionTok{ggtitle}\NormalTok{(}\StringTok{"Epistemic Robustness: 2D Prior Perturbation Grid"}\NormalTok{) }\SpecialCharTok{+}
  \FunctionTok{xlab}\NormalTok{(}\FunctionTok{expression}\NormalTok{(theta }\SpecialCharTok{\textasciitilde{}} \StringTok{" (Probability scale)"}\NormalTok{)) }\SpecialCharTok{+}
  \FunctionTok{ylab}\NormalTok{(}\StringTok{"Marginal Density"}\NormalTok{) }\SpecialCharTok{+}
  \FunctionTok{theme\_classic}\NormalTok{() }\SpecialCharTok{+}
  \FunctionTok{theme}\NormalTok{(}
    \AttributeTok{legend.position =} \StringTok{"bottom"}\NormalTok{,}
    \AttributeTok{strip.background =} \FunctionTok{element\_rect}\NormalTok{(}\AttributeTok{fill =} \StringTok{"gray95"}\NormalTok{, }\AttributeTok{color =} \StringTok{"NA"}\NormalTok{),}
    \AttributeTok{strip.text =} \FunctionTok{element\_text}\NormalTok{(}\AttributeTok{face =} \StringTok{"bold"}\NormalTok{, }\AttributeTok{size =} \DecValTok{11}\NormalTok{)}
\NormalTok{  )}

\NormalTok{p\_sensitivity\_2d}
\end{Highlighting}
\end{Shaded}

\begin{center}
\includegraphics[width=0.8\linewidth,height=\textheight,keepaspectratio]{05-ModelQualityChecks_files/figure-pdf/ch06_manual_sensitivity_analysis-1.pdf}
\end{center}

\begin{Shaded}
\begin{Highlighting}[]
\NormalTok{sensitivity\_summary }\OtherTok{\textless{}{-}}\NormalTok{ sensitivity\_draws }\SpecialCharTok{|\textgreater{}} 
\NormalTok{  dplyr}\SpecialCharTok{::}\FunctionTok{group\_by}\NormalTok{(prior\_mu, prior\_sigma) }\SpecialCharTok{|\textgreater{}} 
\NormalTok{  dplyr}\SpecialCharTok{::}\FunctionTok{summarize}\NormalTok{(}
    \AttributeTok{post\_mean =} \FunctionTok{mean}\NormalTok{(theta\_prob),}
    \CommentTok{\# The 95\% CrI (Outer Ribbon)}
    \AttributeTok{lower\_95 =} \FunctionTok{quantile}\NormalTok{(theta\_prob, }\FloatTok{0.025}\NormalTok{),}
    \AttributeTok{upper\_95 =} \FunctionTok{quantile}\NormalTok{(theta\_prob, }\FloatTok{0.975}\NormalTok{),}
    \CommentTok{\# The 68\% CrI (\textasciitilde{}1 Standard Deviation, Inner Ribbon)}
    \AttributeTok{lower\_68 =} \FunctionTok{quantile}\NormalTok{(theta\_prob, }\FloatTok{0.16}\NormalTok{),}
    \AttributeTok{upper\_68 =} \FunctionTok{quantile}\NormalTok{(theta\_prob, }\FloatTok{0.84}\NormalTok{),}
    \AttributeTok{.groups =} \StringTok{"drop"}
\NormalTok{  ) }\SpecialCharTok{|\textgreater{}} 
  \CommentTok{\# Create strict, mathematically explicit facet labels}
\NormalTok{  dplyr}\SpecialCharTok{::}\FunctionTok{mutate}\NormalTok{(}
    \AttributeTok{prior\_mu\_label =} \FunctionTok{factor}\NormalTok{(}
\NormalTok{      prior\_mu, }
      \AttributeTok{levels =} \FunctionTok{c}\NormalTok{(}\SpecialCharTok{{-}}\FloatTok{1.5}\NormalTok{, }\DecValTok{0}\NormalTok{, }\FloatTok{1.5}\NormalTok{), }
      \AttributeTok{labels =} \FunctionTok{c}\NormalTok{(}\StringTok{"Prior \textbackslash{}u03bc = {-}1.5 (Strong Left Bias)"}\NormalTok{, }
                 \StringTok{"Prior \textbackslash{}u03bc = 0 (Neutral)"}\NormalTok{, }
                 \StringTok{"Prior \textbackslash{}u03bc = 1.5 (Strong Right Bias)"}\NormalTok{)}
\NormalTok{    )}
\NormalTok{  )}

\CommentTok{\# 2. Generate the bounded trajectory plot with graded topological ribbons}
\NormalTok{p\_sensitivity\_trajectory }\OtherTok{\textless{}{-}} \FunctionTok{ggplot}\NormalTok{(sensitivity\_summary, }\FunctionTok{aes}\NormalTok{(}\AttributeTok{x =}\NormalTok{ prior\_sigma, }\AttributeTok{y =}\NormalTok{ post\_mean)) }\SpecialCharTok{+}
  \CommentTok{\# The generative ground truth is our invariant reference line}
  \FunctionTok{geom\_hline}\NormalTok{(}\AttributeTok{yintercept =}\NormalTok{ true\_theta\_prob, }\AttributeTok{color =} \StringTok{"black"}\NormalTok{, }\AttributeTok{linetype =} \StringTok{"dashed"}\NormalTok{, }\AttributeTok{linewidth =} \DecValTok{1}\NormalTok{) }\SpecialCharTok{+}
  
  \CommentTok{\# The 95\% Credible Interval (Outer Ribbon, Highly Transparent)}
  \FunctionTok{geom\_ribbon}\NormalTok{(}\FunctionTok{aes}\NormalTok{(}\AttributeTok{ymin =}\NormalTok{ lower\_95, }\AttributeTok{ymax =}\NormalTok{ upper\_95, }\AttributeTok{fill =}\NormalTok{ prior\_mu\_label), }\AttributeTok{alpha =} \FloatTok{0.15}\NormalTok{) }\SpecialCharTok{+}
  
  \CommentTok{\# The 68\% Credible Interval (\textasciitilde{}1 SD, Inner Ribbon, More Opaque)}
  \FunctionTok{geom\_ribbon}\NormalTok{(}\FunctionTok{aes}\NormalTok{(}\AttributeTok{ymin =}\NormalTok{ lower\_68, }\AttributeTok{ymax =}\NormalTok{ upper\_68, }\AttributeTok{fill =}\NormalTok{ prior\_mu\_label), }\AttributeTok{alpha =} \FloatTok{0.35}\NormalTok{) }\SpecialCharTok{+}
  
  \CommentTok{\# The posterior mean trajectory}
  \FunctionTok{geom\_line}\NormalTok{(}\FunctionTok{aes}\NormalTok{(}\AttributeTok{color =}\NormalTok{ prior\_mu\_label), }\AttributeTok{linewidth =} \DecValTok{1}\NormalTok{) }\SpecialCharTok{+}
  \FunctionTok{geom\_point}\NormalTok{(}\FunctionTok{aes}\NormalTok{(}\AttributeTok{color =}\NormalTok{ prior\_mu\_label), }\AttributeTok{size =} \DecValTok{2}\NormalTok{) }\SpecialCharTok{+}
  
  \CommentTok{\# Structural faceting}
  \FunctionTok{facet\_wrap}\NormalTok{(}\SpecialCharTok{\textasciitilde{}}\NormalTok{ prior\_mu\_label) }\SpecialCharTok{+}
  \FunctionTok{scale\_color\_viridis\_d}\NormalTok{(}\AttributeTok{option =} \StringTok{"mako"}\NormalTok{, }\AttributeTok{begin =} \FloatTok{0.3}\NormalTok{, }\AttributeTok{end =} \FloatTok{0.7}\NormalTok{) }\SpecialCharTok{+}
  \FunctionTok{scale\_fill\_viridis\_d}\NormalTok{(}\AttributeTok{option =} \StringTok{"mako"}\NormalTok{, }\AttributeTok{begin =} \FloatTok{0.3}\NormalTok{, }\AttributeTok{end =} \FloatTok{0.7}\NormalTok{) }\SpecialCharTok{+}
  \FunctionTok{coord\_cartesian}\NormalTok{(}\AttributeTok{ylim =} \FunctionTok{c}\NormalTok{(}\FloatTok{0.5}\NormalTok{, }\FloatTok{1.0}\NormalTok{)) }\SpecialCharTok{+} \CommentTok{\# Zoom to the scientifically relevant probability bounds}
  
  \CommentTok{\# Pedagogical labeling}
  \FunctionTok{ggtitle}\NormalTok{(}\StringTok{"Epistemic Trajectory: Graded Posterior Estimates across Prior Perturbations"}\NormalTok{) }\SpecialCharTok{+}
  \FunctionTok{xlab}\NormalTok{(}\FunctionTok{expression}\NormalTok{(}\StringTok{"Prior Standard Deviation ("} \SpecialCharTok{*}\NormalTok{ sigma }\SpecialCharTok{*} \StringTok{") {-}\textgreater{} Increasing Diffuseness"}\NormalTok{)) }\SpecialCharTok{+}
  \FunctionTok{ylab}\NormalTok{(}\FunctionTok{expression}\NormalTok{(}\StringTok{"Posterior Estimate of "} \SpecialCharTok{*}\NormalTok{ theta }\SpecialCharTok{*} \StringTok{" (Mean, 68\%, and 95\% CrI)"}\NormalTok{)) }\SpecialCharTok{+}
  
  \FunctionTok{theme\_classic}\NormalTok{() }\SpecialCharTok{+}
  \FunctionTok{theme}\NormalTok{(}
    \AttributeTok{legend.position =} \StringTok{"none"}\NormalTok{, }\CommentTok{\# Redundant with facet labels}
    \AttributeTok{strip.background =} \FunctionTok{element\_rect}\NormalTok{(}\AttributeTok{fill =} \StringTok{"gray95"}\NormalTok{, }\AttributeTok{color =} \StringTok{"NA"}\NormalTok{),}
    \AttributeTok{strip.text =} \FunctionTok{element\_text}\NormalTok{(}\AttributeTok{face =} \StringTok{"bold"}\NormalTok{, }\AttributeTok{size =} \DecValTok{11}\NormalTok{),}
    \CommentTok{\# Add light y{-}grid lines to help the eye track the shrinkage}
    \AttributeTok{panel.grid.major.y =} \FunctionTok{element\_line}\NormalTok{(}\AttributeTok{color =} \StringTok{"gray90"}\NormalTok{, }\AttributeTok{linetype =} \StringTok{"dotted"}\NormalTok{) }
\NormalTok{  )}

\FunctionTok{print}\NormalTok{(p\_sensitivity\_trajectory)}
\end{Highlighting}
\end{Shaded}

\begin{center}
\includegraphics[width=0.8\linewidth,height=\textheight,keepaspectratio]{05-ModelQualityChecks_files/figure-pdf/ch06_manual_sensitivity_analysis-2.pdf}
\end{center}

\subsection{5b. Sensitivity Analysis using
priorsense}\label{b.-sensitivity-analysis-using-priorsense}

Executing a computationally heavy 30-model grid is acceptable for a
single-parameter Bernoulli process, but what happens when we transition
to high-dimensional hierarchical structures in later chapter? The
computational cost of manual grid searches becomes prohibitive.

To solve this, we rely on the priorsense package, which elegantly
estimates prior sensitivity directly from the original model fit without
requiring a single MCMC re-run. It achieves this via power-scaling, a
mathematical technique that analytically down-weights alternatively the
likelihood and the prior to simulate how the posterior would deform
under different epistemic conditions. Conceptually, it evaluates the
perturbed posterior:

\[p(\theta \mid y) \propto p(y \mid \theta)^{\alpha} p(\theta)^{\beta}\]

where \(\alpha\) and \(\beta\) are scaling factors.

\subsubsection{Conceptual Idea:}\label{conceptual-idea}

\begin{enumerate}
\def\labelenumi{\arabic{enumi}.}
\item
  Baseline Fit: Fit the model once using your principled, generative
  priors.
\item
  Importance Sampling: Use priorsense::powerscale\_sensitivity to apply
  Pareto-smoothed importance sampling (PSIS). This analytically
  estimates the posterior geometry under perturbed scaling factors
  (\(\alpha \neq 1, \beta \neq 1\)).
\item
  Compute Sensitivity (\(\psi\)): The algorithm calculates specific
  diagnostic metrics, such as sensitivity\_psi. This metric quantifies
  the gradient of change: how much does the posterior mean shift
  (relative to its standard deviation) when the prior is perturbed?
\item
  Epistemic Diagnosis: High sensitivity values (e.g., \(\psi > 0.1\))
  indicate structural fragility. This suggests either that the prior is
  pathologically dominating the likelihood, or that the prior and
  likelihood are fundamentally in conflict (prior-data conflict).
\end{enumerate}

\begin{Shaded}
\begin{Highlighting}[]
\CommentTok{\# We run prior sense}
\NormalTok{sensitivity\_summary }\OtherTok{\textless{}{-}}\NormalTok{ priorsense}\SpecialCharTok{::}\FunctionTok{powerscale\_sensitivity}\NormalTok{(fit\_gq, }
    \AttributeTok{variable =} \FunctionTok{c}\NormalTok{(}\StringTok{"theta"}\NormalTok{))}

\CommentTok{\# we extract the variable of interest}
\NormalTok{theta\_metrics }\OtherTok{\textless{}{-}}\NormalTok{ sensitivity\_summary }\SpecialCharTok{|\textgreater{}}
\NormalTok{  dplyr}\SpecialCharTok{::}\FunctionTok{filter}\NormalTok{(variable }\SpecialCharTok{==} \StringTok{"theta"}\NormalTok{)}

\FunctionTok{cat}\NormalTok{(}\StringTok{"}\SpecialCharTok{\textbackslash{}n}\StringTok{{-}{-}{-} Epistemic Diagnosis {-}{-}{-}}\SpecialCharTok{\textbackslash{}n}\StringTok{"}\NormalTok{)}
\end{Highlighting}
\end{Shaded}

\begin{verbatim}

--- Epistemic Diagnosis ---
\end{verbatim}

\begin{Shaded}
\begin{Highlighting}[]
\FunctionTok{print}\NormalTok{(theta\_metrics)}
\end{Highlighting}
\end{Shaded}

\begin{verbatim}
Sensitivity based on cjs_dist
Prior selection: all priors
Likelihood selection: all data

 variable prior likelihood diagnosis
    theta 0.027      0.104         -
\end{verbatim}

\begin{Shaded}
\begin{Highlighting}[]
\CommentTok{\# Visualize the geometric deformation of the posterior}
\NormalTok{p\_priorsense\_dens }\OtherTok{\textless{}{-}}\NormalTok{ priorsense}\SpecialCharTok{::}\FunctionTok{powerscale\_plot\_dens}\NormalTok{(fit\_gq, }\AttributeTok{variables =} \StringTok{"theta"}\NormalTok{) }\SpecialCharTok{+}
  \FunctionTok{ggtitle}\NormalTok{(}\StringTok{"Analytical Epistemic Fragility: Power{-}Scaled Density Deformation"}\NormalTok{) }\SpecialCharTok{+}
  \FunctionTok{theme\_classic}\NormalTok{() }\SpecialCharTok{+}
  \FunctionTok{theme}\NormalTok{(}
    \AttributeTok{strip.background =} \FunctionTok{element\_rect}\NormalTok{(}\AttributeTok{fill =} \StringTok{"gray95"}\NormalTok{, }\AttributeTok{color =} \StringTok{"NA"}\NormalTok{),}
    \AttributeTok{strip.text =} \FunctionTok{element\_text}\NormalTok{(}\AttributeTok{face =} \StringTok{"bold"}\NormalTok{, }\AttributeTok{size =} \DecValTok{11}\NormalTok{),}
    \AttributeTok{legend.position =} \StringTok{"bottom"}
\NormalTok{  )}

\NormalTok{p\_priorsense\_dens}
\end{Highlighting}
\end{Shaded}

\begin{center}
\includegraphics[width=0.8\linewidth,height=\textheight,keepaspectratio]{05-ModelQualityChecks_files/figure-pdf/ch06_priorsense_analysis-1.pdf}
\end{center}

\section{Phase 6: Computational
Calibration}\label{phase-6-computational-calibration}

\begin{quote}
\emph{Why run parameter recovery and SBC after the model already fits
the data?} Phases 2--5 tested the model's \emph{structure}. Phase 6
tests the \emph{engine}. Even a perfectly specified model produces
garbage if the HMC sampler cannot explore the posterior correctly. And
viceversa, testing the model across many parameter values is a waste of
computation if there are obvious issues with the model that could have
been solved looking at one simulation only. Parameter recovery (6a)
checks whether the engine finds the right answer at specific known
values. SBC (6b) checks whether the posterior \emph{intervals} are
calibrated across the entire prior space --- a 95\% CI should contain
the truth 95\% of the time, not 70\% or 99\%. These are different tests
that can fail independently.
\end{quote}

We must investigate two structural properties of our inference before we
can trust the posterior:

\begin{itemize}
\item
  Data Adequacy (Parameter Recovery): Does our specific experimental
  design (e.g., \(N=120\) trials) contain enough informational density
  to shrink the posterior around the true parameter?
\item
  Algorithmic Unbiasedness (SBC): Is the Hamiltonian Monte Carlo engine
  mathematically well-calibrated across the entire prior space, or is it
  structurally overconfident or biased?
\end{itemize}

\subsection{6a. Parameter Recovery and Posterior
Shrinkage}\label{a.-parameter-recovery-and-posterior-shrinkage}

We will simulate empirical data across an escalating grid of trial
lengths (30, 60, 120, 240) for three distinct generative agents: a
neutral agent (\(\theta = 0.5\)), a weakly right-biased agent
(\(\theta = 0.7\)), and a strongly-biased agen (\(\theta = 0.9\)). We
then fit our Stan model to each dataset and track the changes in mean
and variance of the estimates.

\begin{Shaded}
\begin{Highlighting}[]
\NormalTok{recovery\_grid }\OtherTok{\textless{}{-}}\NormalTok{ tidyr}\SpecialCharTok{::}\FunctionTok{expand\_grid}\NormalTok{(}
  \AttributeTok{N\_trials =} \FunctionTok{c}\NormalTok{(}\DecValTok{30}\NormalTok{, }\DecValTok{60}\NormalTok{, }\DecValTok{120}\NormalTok{, }\DecValTok{240}\NormalTok{),}
  \AttributeTok{true\_theta =} \FunctionTok{c}\NormalTok{(}\FloatTok{0.5}\NormalTok{, }\FloatTok{0.7}\NormalTok{, }\FloatTok{0.9}\NormalTok{)}
\NormalTok{)}

\NormalTok{recovery\_file }\OtherTok{\textless{}{-}} \FunctionTok{file.path}\NormalTok{(stan\_results\_dir, }\StringTok{"ch07\_recovery\_draws.rds"}\NormalTok{)}


\ControlFlowTok{if}\NormalTok{ (regenerate\_simulations }\SpecialCharTok{||} \SpecialCharTok{!}\FunctionTok{file.exists}\NormalTok{(recovery\_file)) \{}
  
  \CommentTok{\# 2. Define the isolated recovery pipeline}
\NormalTok{  run\_recovery }\OtherTok{\textless{}{-}} \ControlFlowTok{function}\NormalTok{(N\_trials, true\_theta) \{}
    \CommentTok{\# Generate synthetic empirical data}
\NormalTok{    sim\_y }\OtherTok{\textless{}{-}} \FunctionTok{rbinom}\NormalTok{(}\AttributeTok{n =}\NormalTok{ N\_trials, }\AttributeTok{size =} \DecValTok{1}\NormalTok{, }\AttributeTok{prob =}\NormalTok{ true\_theta)}
    
    \CommentTok{\# Structure the data block for Stan}
\NormalTok{    rec\_data\_list }\OtherTok{\textless{}{-}} \FunctionTok{list}\NormalTok{(}
      \AttributeTok{N\_trials =}\NormalTok{ N\_trials,}
      \AttributeTok{y =}\NormalTok{ sim\_y,}
      \AttributeTok{prior\_theta\_mu =} \DecValTok{0}\NormalTok{,}
      \AttributeTok{prior\_theta\_sigma =} \FloatTok{1.5}\NormalTok{,}
      \AttributeTok{run\_diagnostics =} \DecValTok{0}
\NormalTok{    )}
    
    \CommentTok{\# Fit the generative model}
\NormalTok{    fit }\OtherTok{\textless{}{-}}\NormalTok{ mod\_gq}\SpecialCharTok{$}\FunctionTok{sample}\NormalTok{(}
      \AttributeTok{data =}\NormalTok{ rec\_data\_list,}
      \AttributeTok{chains =} \DecValTok{1}\NormalTok{,}
      \AttributeTok{iter\_warmup =} \DecValTok{1000}\NormalTok{,}
      \AttributeTok{iter\_sampling =} \DecValTok{2000}\NormalTok{,}
      \AttributeTok{refresh =} \DecValTok{0}\NormalTok{,}
      \AttributeTok{show\_messages =} \ConstantTok{FALSE}
\NormalTok{    )}
    
    \CommentTok{\# Extract structural draws immediately to avoid temporary file fragility}
\NormalTok{    draws\_df }\OtherTok{\textless{}{-}}\NormalTok{ fit}\SpecialCharTok{$}\FunctionTok{draws}\NormalTok{(}\StringTok{"theta"}\NormalTok{, }\AttributeTok{format =} \StringTok{"df"}\NormalTok{)}
    
\NormalTok{    tibble}\SpecialCharTok{::}\FunctionTok{tibble}\NormalTok{(}
      \AttributeTok{trials =}\NormalTok{ N\_trials,}
      \AttributeTok{true\_theta =}\NormalTok{ true\_theta,}
      \AttributeTok{theta\_prob =}\NormalTok{ draws\_df}\SpecialCharTok{$}\NormalTok{theta}
\NormalTok{    )}
\NormalTok{  \}}

  \CommentTok{\# 3. Establish the parallel backend and execute mapping}
\NormalTok{  future}\SpecialCharTok{::}\FunctionTok{plan}\NormalTok{(future}\SpecialCharTok{::}\NormalTok{multisession, }\AttributeTok{workers =} \DecValTok{2}\NormalTok{)}
  
\NormalTok{  recovery\_draws }\OtherTok{\textless{}{-}}\NormalTok{ furrr}\SpecialCharTok{::}\FunctionTok{future\_pmap}\NormalTok{(}
\NormalTok{    recovery\_grid,}
\NormalTok{    run\_recovery,}
    \AttributeTok{.options =}\NormalTok{ furrr}\SpecialCharTok{::}\FunctionTok{furrr\_options}\NormalTok{(}\AttributeTok{seed =} \ConstantTok{TRUE}\NormalTok{, }\AttributeTok{stdout =} \ConstantTok{FALSE}\NormalTok{)}
\NormalTok{  ) }\SpecialCharTok{|\textgreater{}} 
\NormalTok{    dplyr}\SpecialCharTok{::}\FunctionTok{bind\_rows}\NormalTok{()}
    
\NormalTok{  future}\SpecialCharTok{::}\FunctionTok{plan}\NormalTok{(future}\SpecialCharTok{::}\NormalTok{sequential)}
  
  \FunctionTok{saveRDS}\NormalTok{(recovery\_draws, recovery\_file)}
  
\NormalTok{\} }\ControlFlowTok{else}\NormalTok{ \{}
\NormalTok{  recovery\_draws }\OtherTok{\textless{}{-}} \FunctionTok{readRDS}\NormalTok{(recovery\_file)}
  \FunctionTok{cat}\NormalTok{(}\StringTok{"Loaded existing 2D recovery draws from:"}\NormalTok{, recovery\_file, }\StringTok{"}\SpecialCharTok{\textbackslash{}n}\StringTok{"}\NormalTok{)}
\NormalTok{\}}
\end{Highlighting}
\end{Shaded}

\begin{verbatim}
Loaded existing 2D recovery draws from: simmodels/ch07_recovery_draws.rds 
\end{verbatim}

\begin{Shaded}
\begin{Highlighting}[]
\CommentTok{\# 4. Compress draws into rigorous posterior summaries}
\NormalTok{recovery\_summary }\OtherTok{\textless{}{-}}\NormalTok{ recovery\_draws }\SpecialCharTok{|\textgreater{}} 
\NormalTok{  dplyr}\SpecialCharTok{::}\FunctionTok{group\_by}\NormalTok{(trials, true\_theta) }\SpecialCharTok{|\textgreater{}} 
\NormalTok{  dplyr}\SpecialCharTok{::}\FunctionTok{summarize}\NormalTok{(}
    \AttributeTok{post\_mean =} \FunctionTok{mean}\NormalTok{(theta\_prob),}
    \AttributeTok{lower\_95 =} \FunctionTok{quantile}\NormalTok{(theta\_prob, }\FloatTok{0.025}\NormalTok{),}
    \AttributeTok{upper\_95 =} \FunctionTok{quantile}\NormalTok{(theta\_prob, }\FloatTok{0.975}\NormalTok{),}
    \AttributeTok{.groups =} \StringTok{"drop"}
\NormalTok{  ) }\SpecialCharTok{|\textgreater{}} 
\NormalTok{  dplyr}\SpecialCharTok{::}\FunctionTok{mutate}\NormalTok{(}
    \AttributeTok{facet\_label =} \FunctionTok{paste0}\NormalTok{(}\StringTok{"Generative Ground Truth: \textbackslash{}u03b8 = "}\NormalTok{, true\_theta)}
\NormalTok{  )}

\CommentTok{\# 5. Visualize the geometric shrinkage across the parameter space}

\NormalTok{p\_recovery\_2d }\OtherTok{\textless{}{-}} \FunctionTok{ggplot}\NormalTok{(recovery\_summary, }\FunctionTok{aes}\NormalTok{(}\AttributeTok{x =} \FunctionTok{factor}\NormalTok{(trials), }\AttributeTok{y =}\NormalTok{ post\_mean)) }\SpecialCharTok{+}
  \CommentTok{\# Map the true value to the hline dynamically per facet}
  \FunctionTok{geom\_hline}\NormalTok{(}\FunctionTok{aes}\NormalTok{(}\AttributeTok{yintercept =}\NormalTok{ true\_theta), }\AttributeTok{color =} \StringTok{"black"}\NormalTok{, }\AttributeTok{linetype =} \StringTok{"dashed"}\NormalTok{, }\AttributeTok{linewidth =} \DecValTok{1}\NormalTok{) }\SpecialCharTok{+}
  \FunctionTok{geom\_pointrange}\NormalTok{(}\FunctionTok{aes}\NormalTok{(}\AttributeTok{ymin =}\NormalTok{ lower\_95, }\AttributeTok{ymax =}\NormalTok{ upper\_95, }\AttributeTok{color =} \FunctionTok{factor}\NormalTok{(true\_theta)), }\AttributeTok{linewidth =} \DecValTok{1}\NormalTok{, }\AttributeTok{size =} \FloatTok{0.8}\NormalTok{) }\SpecialCharTok{+}
  
  \FunctionTok{facet\_wrap}\NormalTok{(}\SpecialCharTok{\textasciitilde{}}\NormalTok{ facet\_label) }\SpecialCharTok{+}
  \CommentTok{\# Adjust palette range so colors remain distinct}
  \FunctionTok{scale\_color\_viridis\_d}\NormalTok{(}\AttributeTok{option =} \StringTok{"mako"}\NormalTok{, }\AttributeTok{begin =} \FloatTok{0.2}\NormalTok{, }\AttributeTok{end =} \FloatTok{0.8}\NormalTok{) }\SpecialCharTok{+}
  
  \FunctionTok{ggtitle}\NormalTok{(}\StringTok{"Parameter Recovery: Posterior Mean vs. True θ Across N"}\NormalTok{,}
          \AttributeTok{subtitle =} \StringTok{"Healthy = points near dashed true{-}value line; intervals shrink as N increases"}\NormalTok{) }\SpecialCharTok{+}
  \FunctionTok{xlab}\NormalTok{(}\StringTok{"Experimental Design (Number of Trials, N)"}\NormalTok{) }\SpecialCharTok{+}
  \FunctionTok{ylab}\NormalTok{(}\FunctionTok{expression}\NormalTok{(}\StringTok{"Posterior Estimate of "} \SpecialCharTok{*}\NormalTok{ theta }\SpecialCharTok{*} \StringTok{" (Mean \& 95\% CrI)"}\NormalTok{)) }\SpecialCharTok{+}
  
  \FunctionTok{coord\_cartesian}\NormalTok{(}\AttributeTok{ylim =} \FunctionTok{c}\NormalTok{(}\FloatTok{0.35}\NormalTok{, }\FloatTok{1.0}\NormalTok{)) }\SpecialCharTok{+} \CommentTok{\# Bound visually to the relevant domain}
  \FunctionTok{theme\_classic}\NormalTok{() }\SpecialCharTok{+}
  \FunctionTok{theme}\NormalTok{(}
    \AttributeTok{legend.position =} \StringTok{"none"}\NormalTok{,}
    \AttributeTok{strip.background =} \FunctionTok{element\_rect}\NormalTok{(}\AttributeTok{fill =} \StringTok{"gray95"}\NormalTok{, }\AttributeTok{color =} \StringTok{"NA"}\NormalTok{),}
    \AttributeTok{strip.text =} \FunctionTok{element\_text}\NormalTok{(}\AttributeTok{face =} \StringTok{"bold"}\NormalTok{, }\AttributeTok{size =} \DecValTok{11}\NormalTok{),}
    \AttributeTok{panel.grid.major.y =} \FunctionTok{element\_line}\NormalTok{(}\AttributeTok{color =} \StringTok{"gray90"}\NormalTok{, }\AttributeTok{linetype =} \StringTok{"dotted"}\NormalTok{)}
\NormalTok{  )}

\FunctionTok{print}\NormalTok{(p\_recovery\_2d)}
\end{Highlighting}
\end{Shaded}

\begin{center}
\includegraphics[width=0.8\linewidth,height=\textheight,keepaspectratio]{05-ModelQualityChecks_files/figure-pdf/ch06_parameter_recovery_grid-1.pdf}
\end{center}

\begin{Shaded}
\begin{Highlighting}[]
\NormalTok{p\_recovery\_2d\_alt }\OtherTok{\textless{}{-}} \FunctionTok{ggplot}\NormalTok{(recovery\_summary, }\FunctionTok{aes}\NormalTok{(}\AttributeTok{x =}\NormalTok{ true\_theta, }\AttributeTok{y =}\NormalTok{ post\_mean)) }\SpecialCharTok{+}
  \CommentTok{\# The Identity Line: y = x (Estimate perfectly matches Ground Truth)}
  \FunctionTok{geom\_abline}\NormalTok{(}\AttributeTok{intercept =} \DecValTok{0}\NormalTok{, }\AttributeTok{slope =} \DecValTok{1}\NormalTok{, }\AttributeTok{color =} \StringTok{"black"}\NormalTok{, }\AttributeTok{linetype =} \StringTok{"dashed"}\NormalTok{, }\AttributeTok{linewidth =} \DecValTok{1}\NormalTok{) }\SpecialCharTok{+}
  
  \CommentTok{\# The Posterior Estimates and 95\% CrI}
  \FunctionTok{geom\_pointrange}\NormalTok{(}\FunctionTok{aes}\NormalTok{(}\AttributeTok{ymin =}\NormalTok{ lower\_95, }\AttributeTok{ymax =}\NormalTok{ upper\_95, }\AttributeTok{color =} \FunctionTok{factor}\NormalTok{(true\_theta)), }\AttributeTok{linewidth =} \DecValTok{1}\NormalTok{, }\AttributeTok{size =} \FloatTok{0.8}\NormalTok{) }\SpecialCharTok{+}
  
  \CommentTok{\# Facet strictly by the Experimental Design (N trials)}
  \FunctionTok{facet\_wrap}\NormalTok{(}
    \SpecialCharTok{\textasciitilde{}}\NormalTok{ trials, }
    \AttributeTok{labeller =} \FunctionTok{as\_labeller}\NormalTok{(}\ControlFlowTok{function}\NormalTok{(x) }\FunctionTok{paste0}\NormalTok{(x, }\StringTok{" Trials"}\NormalTok{))}
\NormalTok{  ) }\SpecialCharTok{+}
  
  \FunctionTok{scale\_color\_viridis\_d}\NormalTok{(}\AttributeTok{option =} \StringTok{"mako"}\NormalTok{, }\AttributeTok{begin =} \FloatTok{0.2}\NormalTok{, }\AttributeTok{end =} \FloatTok{0.8}\NormalTok{) }\SpecialCharTok{+}
  
  \CommentTok{\# Ensure the coordinate space is perfectly square to preserve the visual geometry of y=x}
  \FunctionTok{coord\_fixed}\NormalTok{(}\AttributeTok{xlim =} \FunctionTok{c}\NormalTok{(}\FloatTok{0.4}\NormalTok{, }\FloatTok{1.0}\NormalTok{), }\AttributeTok{ylim =} \FunctionTok{c}\NormalTok{(}\FloatTok{0.4}\NormalTok{, }\FloatTok{1.0}\NormalTok{)) }\SpecialCharTok{+} 
  
  \FunctionTok{xlab}\NormalTok{(}\FunctionTok{expression}\NormalTok{(}\StringTok{"Generative Ground Truth ("} \SpecialCharTok{*}\NormalTok{ theta }\SpecialCharTok{*} \StringTok{")"}\NormalTok{)) }\SpecialCharTok{+}
  \FunctionTok{ylab}\NormalTok{(}\FunctionTok{expression}\NormalTok{(}\StringTok{"Posterior Estimate ("} \SpecialCharTok{*} \FunctionTok{hat}\NormalTok{(theta) }\SpecialCharTok{*} \StringTok{" Mean \& 95\% CrI)"}\NormalTok{)) }\SpecialCharTok{+}
  \FunctionTok{ggtitle}\NormalTok{(}\StringTok{"Parameter Recovery Identity Plot"}\NormalTok{,}
          \AttributeTok{subtitle =} \StringTok{"Healthy = points on the diagonal (y = x); bars are 95\% CrI. N grows left → right."}\NormalTok{) }\SpecialCharTok{+}
  \FunctionTok{theme\_classic}\NormalTok{() }\SpecialCharTok{+}
  \FunctionTok{theme}\NormalTok{(}
    \AttributeTok{legend.position =} \StringTok{"none"}\NormalTok{,}
    \AttributeTok{strip.background =} \FunctionTok{element\_rect}\NormalTok{(}\AttributeTok{fill =} \StringTok{"gray95"}\NormalTok{, }\AttributeTok{color =} \StringTok{"NA"}\NormalTok{),}
    \AttributeTok{strip.text =} \FunctionTok{element\_text}\NormalTok{(}\AttributeTok{face =} \StringTok{"bold"}\NormalTok{, }\AttributeTok{size =} \DecValTok{11}\NormalTok{),}
    \CommentTok{\# Add a faint grid to help the eye assess the distance from the diagonal}
    \AttributeTok{panel.grid.major =} \FunctionTok{element\_line}\NormalTok{(}\AttributeTok{color =} \StringTok{"gray90"}\NormalTok{, }\AttributeTok{linetype =} \StringTok{"dotted"}\NormalTok{)}
\NormalTok{  )}

\FunctionTok{print}\NormalTok{(p\_recovery\_2d\_alt)}
\end{Highlighting}
\end{Shaded}

\begin{center}
\includegraphics[width=0.8\linewidth,height=\textheight,keepaspectratio]{05-ModelQualityChecks_files/figure-pdf/ch06_parameter_recovery_grid-2.pdf}
\end{center}

\emph{Interpretation}

By faceting across our experimental design (\(N\)) and plotting the
generative ground truth on the x-axis, we have constructed a parameter
recovery identity plot. Perfect inference dictates that our posterior
estimates (\(y\)) fall exactly on the true parameters (\(x\)), tracing
the dashed diagonal line (\(y=x\)).

\begin{enumerate}
\def\labelenumi{\arabic{enumi}.}
\item
  Epistemic Evolution: Read the facets from left to right. At \(N=30\),
  the model is fundamentally underpowered. The 95\% Credible Intervals
  (vertical bars) are massive, and the mean estimates may drift away
  from the diagonal due to prior shrinkage. The likelihood is too weak
  to fully constrain the parameter space.
\item
  Asymptotic Consistency: As we scale to \(N=240\), the geometry snaps
  into focus. The estimates lock perfectly onto the diagonal identity
  line, and the uncertainty bounds collapse. We have mathematically
  proven that if we run an experiment with 120 to 240 trials, our
  cognitive architecture is very likely to identify the true underlying
  bias of the agent.
\end{enumerate}

\subsection{Phase 6b. Simulation-Based Calibration Checks
(SBC)}\label{phase-6b.-simulation-based-calibration-checks-sbc}

This is the final crucible for our Simple Biased Agent. Parameter
recovery (Phase 6a) proved that the mean of our posterior is
asymptotically consistent. However, we want a more direct test of
whether the uncertainty of the model is well calibrated.

If our inference systematically generates posterior distributions that
are too narrow (overconfident) or too wide (underconfident), our 95\%
Credible Intervals are fictions. To formally investigate whether our
posterior geometry is calibrated across the entire prior space, we must
execute Simulation-Based Calibration Checks (SBC).

\subsubsection{The Mathematical Theorem of
SBC}\label{the-mathematical-theorem-of-sbc}

SBC relies on a fundamental self-consistency property of Bayesian
inference. If we:

\begin{enumerate}
\def\labelenumi{\arabic{enumi}.}
\item
  Draw a ground-truth parameter from the exact prior:
  \(\tilde{\theta} \sim p(\theta)\)
\item
  Simulate a dataset from the data model:
  \(\tilde{y} \sim p(y \mid \tilde{\theta})\)
\item
  Sample the posterior using our Stan engine:
  \(p(\theta \mid \tilde{y})\)
\end{enumerate}

Then the rank of the true parameter \(\tilde{\theta}\) within the
posterior draws must be strictly Uniformly distributed. If we repeat
this process 1,000 times, any deviation from uniformity mathematically
proves that our computational engine is structurally biased or our model
is misspecified.

\begin{Shaded}
\begin{Highlighting}[]
\NormalTok{sbc\_results\_file }\OtherTok{\textless{}{-}} \FunctionTok{file.path}\NormalTok{(stan\_results\_dir, }\StringTok{"ch07\_sbc\_results.rds"}\NormalTok{)}

\CommentTok{\# {-}{-}{-}{-}{-}{-}{-}{-}{-}{-}{-}{-}{-}{-}{-}{-}{-}{-}{-}{-}{-}{-}{-}{-}{-}{-}{-}{-}{-}{-}{-}{-}{-}{-}{-}{-}{-}{-}{-}{-}{-}{-}{-}{-}{-}{-}{-}{-}{-}{-}{-}{-}{-}{-}{-}{-}{-}{-}{-}{-}{-}{-}{-}{-}{-}{-}{-}{-}{-}{-}{-}{-}{-}{-}{-}}
\CommentTok{\# Backend is defined OUTSIDE the cache conditional so it is always available —}
\CommentTok{\# both for Phase 6b itself and for the power analysis in the next section.}
\CommentTok{\# {-}{-}{-}{-}{-}{-}{-}{-}{-}{-}{-}{-}{-}{-}{-}{-}{-}{-}{-}{-}{-}{-}{-}{-}{-}{-}{-}{-}{-}{-}{-}{-}{-}{-}{-}{-}{-}{-}{-}{-}{-}{-}{-}{-}{-}{-}{-}{-}{-}{-}{-}{-}{-}{-}{-}{-}{-}{-}{-}{-}{-}{-}{-}{-}{-}{-}{-}{-}{-}{-}{-}{-}{-}{-}{-}}
\NormalTok{sbc\_backend }\OtherTok{\textless{}{-}}\NormalTok{ SBC}\SpecialCharTok{::}\FunctionTok{SBC\_backend\_cmdstan\_sample}\NormalTok{(}
\NormalTok{  mod\_gq,}
  \AttributeTok{chains        =} \DecValTok{2}\NormalTok{,}
  \AttributeTok{iter\_warmup   =} \DecValTok{500}\NormalTok{,}
  \AttributeTok{iter\_sampling =} \DecValTok{1000}\NormalTok{,}
  \AttributeTok{refresh       =} \DecValTok{0}\NormalTok{,}
  \AttributeTok{show\_messages =} \ConstantTok{FALSE}
\NormalTok{)}

\ControlFlowTok{if}\NormalTok{ (regenerate\_simulations }\SpecialCharTok{||} \SpecialCharTok{!}\FunctionTok{file.exists}\NormalTok{(sbc\_results\_file)) \{}

  \CommentTok{\# 1. Define the exact mathematical Generative Function}
\NormalTok{  sbc\_generator\_fn }\OtherTok{\textless{}{-}} \ControlFlowTok{function}\NormalTok{() \{}
    \CommentTok{\# Draw from the structural prior (on the unconstrained space)}
\NormalTok{    theta\_logit\_true }\OtherTok{\textless{}{-}} \FunctionTok{rnorm}\NormalTok{(}\DecValTok{1}\NormalTok{, }\AttributeTok{mean =} \DecValTok{0}\NormalTok{, }\AttributeTok{sd =} \FloatTok{1.5}\NormalTok{)}
\NormalTok{    theta\_true       }\OtherTok{\textless{}{-}} \FunctionTok{plogis}\NormalTok{(theta\_logit\_true)}

    \CommentTok{\# Simulate empirical reality (N = 120 trials)}
\NormalTok{    N\_trials }\OtherTok{\textless{}{-}} \DecValTok{120}
\NormalTok{    y\_sim    }\OtherTok{\textless{}{-}} \FunctionTok{rbinom}\NormalTok{(N\_trials, }\AttributeTok{size =} \DecValTok{1}\NormalTok{, }\AttributeTok{prob =}\NormalTok{ theta\_true)}

    \CommentTok{\# Format strictly required by the SBC package}
    \FunctionTok{list}\NormalTok{(}
      \AttributeTok{variables =} \FunctionTok{list}\NormalTok{(}\AttributeTok{theta\_logit =}\NormalTok{ theta\_logit\_true, }\AttributeTok{theta =}\NormalTok{ theta\_true),}
      \AttributeTok{generated =} \FunctionTok{list}\NormalTok{(}
        \AttributeTok{N\_trials          =}\NormalTok{ N\_trials,}
        \AttributeTok{y                 =}\NormalTok{ y\_sim,}
        \AttributeTok{prior\_theta\_mu    =} \DecValTok{0}\NormalTok{,}
        \AttributeTok{prior\_theta\_sigma =} \FloatTok{1.5}\NormalTok{,}
        \AttributeTok{run\_diagnostics   =} \DecValTok{0}
\NormalTok{      )}
\NormalTok{    )}
\NormalTok{  \}}

  \CommentTok{\# Cast the raw function into a formal SBC generator object}
\NormalTok{  sbc\_generator\_obj }\OtherTok{\textless{}{-}}\NormalTok{ SBC}\SpecialCharTok{::}\FunctionTok{SBC\_generator\_function}\NormalTok{(sbc\_generator\_fn)}

  \CommentTok{\# 2. Generate datasets and compute rank statistics}
\NormalTok{  future}\SpecialCharTok{::}\FunctionTok{plan}\NormalTok{(future}\SpecialCharTok{::}\NormalTok{multisession, }\AttributeTok{workers =} \DecValTok{2}\NormalTok{)}

\NormalTok{  sbc\_datasets }\OtherTok{\textless{}{-}}\NormalTok{ SBC}\SpecialCharTok{::}\FunctionTok{generate\_datasets}\NormalTok{(sbc\_generator\_obj, }\AttributeTok{n\_sims =} \DecValTok{200}\NormalTok{)}
\NormalTok{  sbc\_results  }\OtherTok{\textless{}{-}}\NormalTok{ SBC}\SpecialCharTok{::}\FunctionTok{compute\_SBC}\NormalTok{(sbc\_datasets, sbc\_backend, }\AttributeTok{keep\_fits =} \ConstantTok{FALSE}\NormalTok{)}

\NormalTok{  future}\SpecialCharTok{::}\FunctionTok{plan}\NormalTok{(future}\SpecialCharTok{::}\NormalTok{sequential)}

  \FunctionTok{saveRDS}\NormalTok{(sbc\_results, sbc\_results\_file)}
  \FunctionTok{cat}\NormalTok{(}\StringTok{"SBC completed. Results saved to:"}\NormalTok{, sbc\_results\_file, }\StringTok{"}\SpecialCharTok{\textbackslash{}n}\StringTok{"}\NormalTok{)}

\NormalTok{\} }\ControlFlowTok{else}\NormalTok{ \{}
\NormalTok{  sbc\_results }\OtherTok{\textless{}{-}} \FunctionTok{readRDS}\NormalTok{(sbc\_results\_file)}
  \FunctionTok{cat}\NormalTok{(}\StringTok{"Loaded existing SBC results from:"}\NormalTok{, sbc\_results\_file, }\StringTok{"}\SpecialCharTok{\textbackslash{}n}\StringTok{"}\NormalTok{)}
\NormalTok{\}}
\end{Highlighting}
\end{Shaded}

\begin{verbatim}
Loaded existing SBC results from: simmodels/ch07_sbc_results.rds 
\end{verbatim}

\begin{Shaded}
\begin{Highlighting}[]
\CommentTok{\# 4. Visualize the Calibration Geometry}

\CommentTok{\# Plot A: Rank Histogram}

\NormalTok{p\_sbc\_hist }\OtherTok{\textless{}{-}}\NormalTok{ SBC}\SpecialCharTok{::}\FunctionTok{plot\_rank\_hist}\NormalTok{(sbc\_results, }\AttributeTok{variables =} \StringTok{"theta"}\NormalTok{) }\SpecialCharTok{+} 
  \FunctionTok{ggtitle}\NormalTok{(}\StringTok{"SBC Rank Histogram: Algorithmic Calibration"}\NormalTok{,}
          \AttributeTok{subtitle =} \StringTok{"Healthy = bars flat within grey band. U{-}shape = overconfident. Hump = underconfident. Skew = biased."}\NormalTok{) }\SpecialCharTok{+}
  \FunctionTok{theme\_classic}\NormalTok{() }\SpecialCharTok{+}
  \FunctionTok{theme}\NormalTok{(}\AttributeTok{strip.background =} \FunctionTok{element\_rect}\NormalTok{(}\AttributeTok{fill =} \StringTok{"gray95"}\NormalTok{, }\AttributeTok{color =} \StringTok{"NA"}\NormalTok{))}

\CommentTok{\# Plot B: Empirical Cumulative Distribution Function (ECDF)}

\NormalTok{p\_sbc\_ecdf }\OtherTok{\textless{}{-}}\NormalTok{ SBC}\SpecialCharTok{::}\FunctionTok{plot\_ecdf}\NormalTok{(sbc\_results, }\AttributeTok{variables =} \StringTok{"theta"}\NormalTok{) }\SpecialCharTok{+}
  \FunctionTok{ggtitle}\NormalTok{(}\StringTok{"SBC ECDF: Deviation from Uniformity"}\NormalTok{,}
          \AttributeTok{subtitle =} \StringTok{"Healthy = blue line stays inside the grey ellipse at all points"}\NormalTok{) }\SpecialCharTok{+}
  \FunctionTok{theme\_classic}\NormalTok{() }\SpecialCharTok{+}
  \FunctionTok{theme}\NormalTok{(}\AttributeTok{strip.background =} \FunctionTok{element\_rect}\NormalTok{(}\AttributeTok{fill =} \StringTok{"gray95"}\NormalTok{, }\AttributeTok{color =} \StringTok{"NA"}\NormalTok{))}

\CommentTok{\# Combine plots  using patchwork}
\NormalTok{p\_sbc\_combined }\OtherTok{\textless{}{-}}\NormalTok{ p\_sbc\_hist }\SpecialCharTok{/}\NormalTok{ p\_sbc\_ecdf}
\NormalTok{p\_sbc\_combined}
\end{Highlighting}
\end{Shaded}

\begin{center}
\includegraphics[width=0.8\linewidth,height=\textheight,keepaspectratio]{05-ModelQualityChecks_files/figure-pdf/ch06_sbc_calibration-1.pdf}
\end{center}

\textbf{How to read these two plots:}

\emph{Rank histogram:} Under perfect calibration, every rank band is
equally likely, so the bars should be flat within the grey 95\%
confidence band.

\begin{longtable}[]{@{}
  >{\raggedright\arraybackslash}p{(\linewidth - 4\tabcolsep) * \real{0.3333}}
  >{\raggedright\arraybackslash}p{(\linewidth - 4\tabcolsep) * \real{0.3333}}
  >{\raggedright\arraybackslash}p{(\linewidth - 4\tabcolsep) * \real{0.3333}}@{}}
\toprule\noalign{}
\begin{minipage}[b]{\linewidth}\raggedright
Shape
\end{minipage} & \begin{minipage}[b]{\linewidth}\raggedright
Diagnosis
\end{minipage} & \begin{minipage}[b]{\linewidth}\raggedright
Action
\end{minipage} \\
\midrule\noalign{}
\endhead
\bottomrule\noalign{}
\endlastfoot
Flat (within grey band) & Calibrated ✓ & Proceed \\
U-shape (high at edges) & Overconfident --- posteriors too narrow &
Widen priors or reparameterise \\
Hump (high in centre) & Underconfident --- posteriors too wide & Tighten
priors \\
Skewed left or right & Directional bias in the sampler & Check model
code; try non-centred parameterisation \\
\end{longtable}

\emph{ECDF plot:} The blue line is the empirical CDF of the ranks; the
grey ellipse is the region consistent with uniformity. Any excursion
outside the ellipse at a given rank level is a calibration failure at
that part of the parameter space.

\textbf{So what?} A flat histogram and an in-bounds ECDF line mean that
when we report a 95\% credible interval, it actually contains the true
parameter 95\% of the time. This is the formal guarantee that makes
Bayesian inference meaningful.

\subsubsection{Parameter Recovery via
SBC}\label{parameter-recovery-via-sbc}

The rank histogram and ECDF confirm \emph{algorithmic calibration} ---
that the sampler is unbiased. But we have been handed something more for
free. The \texttt{sbc\_results\$stats} dataframe, which the SBC package
computes as a byproduct of the rank calculation, contains for every
simulated universe:

\begin{itemize}
\item
  \texttt{simulated\_value}: the ground-truth \(\theta\) drawn from the
  prior --- the quantity we passed in \texttt{variables}.
\item
  \texttt{median}, \texttt{q5}, \texttt{q95}: the posterior median and
  the bounds of the 90\% credible interval from the fitted posterior.
\end{itemize}

This is precisely the data we need for a \textbf{systematic parameter
recovery study}, and we already have it. No extra computation is
required. The crucial distinction from our grid recovery in Phase 6a is
architectural:

\begin{longtable}[]{@{}
  >{\raggedright\arraybackslash}p{(\linewidth - 2\tabcolsep) * \real{0.5000}}
  >{\raggedright\arraybackslash}p{(\linewidth - 2\tabcolsep) * \real{0.5000}}@{}}
\toprule\noalign{}
\begin{minipage}[b]{\linewidth}\raggedright
Phase 6a (grid recovery)
\end{minipage} & \begin{minipage}[b]{\linewidth}\raggedright
SBC-based recovery
\end{minipage} \\
\midrule\noalign{}
\endhead
\bottomrule\noalign{}
\endlastfoot
Fixed ground truths: \(\theta \in \{0.5, 0.7, 0.9\}\) & 200 values
spanning \(\text{Normal}(0, 1.5)\) on logit scale \\
Asks ``does it work at these specific values?'' & Asks ``does it work
everywhere the prior places mass?'' \\
Controlled, pedagogically interpretable & Exhaustive, scientifically
complete \\
\end{longtable}

The grid recovery shows the shrinkage structure around specific anchor
points with interpretable coordinates. The SBC recovery shows the full
recovery profile across the entire plausible parameter space (if we have
set our priors correctly :-)). I set up the first to give you a concrete
idea of what happens in the code, but the latter is more systematic and
rigorous.

\begin{Shaded}
\begin{Highlighting}[]
\CommentTok{\# Extract the lightweight stats dataframe — already computed as part of SBC above.}
\NormalTok{sbc\_stats }\OtherTok{\textless{}{-}}\NormalTok{ sbc\_results}\SpecialCharTok{$}\NormalTok{stats}

\CommentTok{\# We plot on the probability scale for interpretability, matching Phase 6a.}
\CommentTok{\# The \textquotesingle{}theta\textquotesingle{} variable was stored (on probability scale) alongside \textquotesingle{}theta\_logit\textquotesingle{}}
\CommentTok{\# in our generator\textquotesingle{}s variables list.}
\NormalTok{recovery\_sbc\_df }\OtherTok{\textless{}{-}}\NormalTok{ sbc\_stats }\SpecialCharTok{|\textgreater{}}
\NormalTok{  dplyr}\SpecialCharTok{::}\FunctionTok{filter}\NormalTok{(variable }\SpecialCharTok{==} \StringTok{"theta"}\NormalTok{)}

\CommentTok{\# {-}{-}{-}{-}{-}{-}{-}{-}{-}{-}{-}{-}{-}{-}{-}{-}{-}{-}{-}{-}{-}{-}{-}{-}{-}{-}{-}{-}{-}{-}{-}{-}{-}{-}{-}{-}{-}{-}{-}{-}{-}{-}{-}{-}{-}{-}{-}{-}{-}{-}{-}{-}{-}{-}{-}{-}{-}{-}{-}{-}{-}{-}{-}{-}{-}{-}{-}{-}{-}{-}{-}{-}{-}{-}{-}}
\CommentTok{\# Identity plot: true value (x) vs. posterior median (y)}
\CommentTok{\# {-}{-}{-}{-}{-}{-}{-}{-}{-}{-}{-}{-}{-}{-}{-}{-}{-}{-}{-}{-}{-}{-}{-}{-}{-}{-}{-}{-}{-}{-}{-}{-}{-}{-}{-}{-}{-}{-}{-}{-}{-}{-}{-}{-}{-}{-}{-}{-}{-}{-}{-}{-}{-}{-}{-}{-}{-}{-}{-}{-}{-}{-}{-}{-}{-}{-}{-}{-}{-}{-}{-}{-}{-}{-}{-}}
\CommentTok{\# Under perfect recovery, every point lies on the dashed y = x diagonal.}
\CommentTok{\# Scatter around the diagonal is unavoidable estimation uncertainty.}
\CommentTok{\# Systematic curvature (bowing above or below) would reveal directional bias.}
\CommentTok{\# {-}{-}{-}{-}{-}{-}{-}{-}{-}{-}{-}{-}{-}{-}{-}{-}{-}{-}{-}{-}{-}{-}{-}{-}{-}{-}{-}{-}{-}{-}{-}{-}{-}{-}{-}{-}{-}{-}{-}{-}{-}{-}{-}{-}{-}{-}{-}{-}{-}{-}{-}{-}{-}{-}{-}{-}{-}{-}{-}{-}{-}{-}{-}{-}{-}{-}{-}{-}{-}{-}{-}{-}{-}{-}{-}}
\NormalTok{p\_sbc\_recovery }\OtherTok{\textless{}{-}} \FunctionTok{ggplot}\NormalTok{(recovery\_sbc\_df, }\FunctionTok{aes}\NormalTok{(}\AttributeTok{x =}\NormalTok{ simulated\_value, }\AttributeTok{y =}\NormalTok{ median)) }\SpecialCharTok{+}

  \CommentTok{\# Perfect{-}recovery reference line}
  \FunctionTok{geom\_abline}\NormalTok{(}\AttributeTok{intercept =} \DecValTok{0}\NormalTok{, }\AttributeTok{slope =} \DecValTok{1}\NormalTok{,}
              \AttributeTok{linetype =} \StringTok{"dashed"}\NormalTok{, }\AttributeTok{color =} \StringTok{"gray50"}\NormalTok{, }\AttributeTok{linewidth =} \FloatTok{0.8}\NormalTok{) }\SpecialCharTok{+}

  \CommentTok{\# 90\% posterior credible intervals}
  \CommentTok{\# Under correct calibration, \textasciitilde{}90\% of bars should contain the true value —}
  \CommentTok{\# which we already verified via the rank histogram above.}
  \FunctionTok{geom\_errorbar}\NormalTok{(}\FunctionTok{aes}\NormalTok{(}\AttributeTok{ymin =}\NormalTok{ q5, }\AttributeTok{ymax =}\NormalTok{ q95),}
                \AttributeTok{width =} \DecValTok{0}\NormalTok{, }\AttributeTok{alpha =} \FloatTok{0.25}\NormalTok{, }\AttributeTok{color =} \StringTok{"steelblue"}\NormalTok{) }\SpecialCharTok{+}

  \CommentTok{\# Posterior median point estimates}
  \FunctionTok{geom\_point}\NormalTok{(}\AttributeTok{color =} \StringTok{"midnightblue"}\NormalTok{, }\AttributeTok{alpha =} \FloatTok{0.6}\NormalTok{, }\AttributeTok{size =} \FloatTok{1.8}\NormalTok{) }\SpecialCharTok{+}

  \FunctionTok{coord\_fixed}\NormalTok{(}\AttributeTok{xlim =} \FunctionTok{c}\NormalTok{(}\DecValTok{0}\NormalTok{, }\DecValTok{1}\NormalTok{), }\AttributeTok{ylim =} \FunctionTok{c}\NormalTok{(}\DecValTok{0}\NormalTok{, }\DecValTok{1}\NormalTok{)) }\SpecialCharTok{+}

  \FunctionTok{labs}\NormalTok{(}
    \AttributeTok{title    =} \StringTok{"Parameter Recovery via SBC: Simple Biased Agent"}\NormalTok{,}
    \AttributeTok{subtitle =} \FunctionTok{paste0}\NormalTok{(}
      \StringTok{"N = 200 simulated universes. True \textbackslash{}u03b8 drawn from the prior predictive.}\SpecialCharTok{\textbackslash{}n}\StringTok{"}\NormalTok{,}
      \StringTok{"Each point is one universe. Bars are 90\% posterior credible intervals."}
\NormalTok{    ),}
    \AttributeTok{x =} \FunctionTok{expression}\NormalTok{(}\StringTok{"True "} \SpecialCharTok{*}\NormalTok{ theta }\SpecialCharTok{*} \StringTok{" (probability scale, drawn from prior)"}\NormalTok{),}
    \AttributeTok{y =} \StringTok{"Posterior Median"}
\NormalTok{  ) }\SpecialCharTok{+}
  \FunctionTok{theme\_classic}\NormalTok{()}

\FunctionTok{print}\NormalTok{(p\_sbc\_recovery)}
\end{Highlighting}
\end{Shaded}

\begin{center}
\includegraphics[width=0.8\linewidth,height=\textheight,keepaspectratio]{05-ModelQualityChecks_files/figure-pdf/ch06_sbc_recovery_plot-1.pdf}
\end{center}

\begin{Shaded}
\begin{Highlighting}[]
\CommentTok{\# {-}{-}{-}{-}{-}{-}{-}{-}{-}{-}{-}{-}{-}{-}{-}{-}{-}{-}{-}{-}{-}{-}{-}{-}{-}{-}{-}{-}{-}{-}{-}{-}{-}{-}{-}{-}{-}{-}{-}{-}{-}{-}{-}{-}{-}{-}{-}{-}{-}{-}{-}{-}{-}{-}{-}{-}{-}{-}{-}{-}{-}{-}{-}{-}{-}{-}{-}{-}{-}{-}{-}{-}{-}{-}{-}}
\CommentTok{\# Quantitative recovery metrics}
\CommentTok{\# {-}{-}{-}{-}{-}{-}{-}{-}{-}{-}{-}{-}{-}{-}{-}{-}{-}{-}{-}{-}{-}{-}{-}{-}{-}{-}{-}{-}{-}{-}{-}{-}{-}{-}{-}{-}{-}{-}{-}{-}{-}{-}{-}{-}{-}{-}{-}{-}{-}{-}{-}{-}{-}{-}{-}{-}{-}{-}{-}{-}{-}{-}{-}{-}{-}{-}{-}{-}{-}{-}{-}{-}{-}{-}{-}}
\CommentTok{\# Pearson r   : linear correlation between truth and estimate (1.0 = perfect)}
\CommentTok{\# RMSE        : average estimation error on the probability scale (0 = perfect)}
\CommentTok{\# coverage90  : empirical 90\% CI coverage — should be near 0.90 under calibration}
\CommentTok{\# {-}{-}{-}{-}{-}{-}{-}{-}{-}{-}{-}{-}{-}{-}{-}{-}{-}{-}{-}{-}{-}{-}{-}{-}{-}{-}{-}{-}{-}{-}{-}{-}{-}{-}{-}{-}{-}{-}{-}{-}{-}{-}{-}{-}{-}{-}{-}{-}{-}{-}{-}{-}{-}{-}{-}{-}{-}{-}{-}{-}{-}{-}{-}{-}{-}{-}{-}{-}{-}{-}{-}{-}{-}{-}{-}}

\NormalTok{recovery\_metrics\_sbc }\OtherTok{\textless{}{-}}\NormalTok{ recovery\_sbc\_df }\SpecialCharTok{|\textgreater{}}
\NormalTok{  dplyr}\SpecialCharTok{::}\FunctionTok{summarize}\NormalTok{(}
    \AttributeTok{Pearson\_r  =} \FunctionTok{cor}\NormalTok{(simulated\_value, median),}
    \AttributeTok{RMSE       =} \FunctionTok{sqrt}\NormalTok{(}\FunctionTok{mean}\NormalTok{((simulated\_value }\SpecialCharTok{{-}}\NormalTok{ median)}\SpecialCharTok{\^{}}\DecValTok{2}\NormalTok{)),}
    \AttributeTok{coverage90 =} \FunctionTok{mean}\NormalTok{(simulated\_value }\SpecialCharTok{\textgreater{}=}\NormalTok{ q5 }\SpecialCharTok{\&}\NormalTok{ simulated\_value }\SpecialCharTok{\textless{}=}\NormalTok{ q95)}
\NormalTok{  )}

\FunctionTok{cat}\NormalTok{(}\StringTok{"=== SBC PARAMETER RECOVERY METRICS (theta, probability scale) ===}\SpecialCharTok{\textbackslash{}n}\StringTok{"}\NormalTok{)}
\end{Highlighting}
\end{Shaded}

\begin{verbatim}
=== SBC PARAMETER RECOVERY METRICS (theta, probability scale) ===
\end{verbatim}

\begin{Shaded}
\begin{Highlighting}[]
\FunctionTok{print}\NormalTok{(}
\NormalTok{  knitr}\SpecialCharTok{::}\FunctionTok{kable}\NormalTok{(}
\NormalTok{    recovery\_metrics\_sbc,}
    \AttributeTok{digits    =} \DecValTok{3}\NormalTok{,}
    \AttributeTok{col.names =} \FunctionTok{c}\NormalTok{(}\StringTok{"Pearson r"}\NormalTok{, }\StringTok{"RMSE"}\NormalTok{, }\StringTok{"Empirical 90\% CI Coverage"}\NormalTok{)}
\NormalTok{  )}
\NormalTok{)}
\end{Highlighting}
\end{Shaded}

\begin{verbatim}


| Pearson r|  RMSE| Empirical 90% CI Coverage|
|---------:|-----:|-------------------------:|
|     0.988| 0.041|                     0.915|
\end{verbatim}

\subsubsection{A Precision (Pseudo-Power) Analysis via
SBC}\label{a-precision-pseudo-power-analysis-via-sbc}

Our Phase 6a grid recovery showed how precision evolves at three
specific \(\theta\) values as \(N\) grows. We can now generalise that
question systematically: across the \emph{full} prior space, at what
\(N\) does the model become reliable?

Because the SBC generator is a plain R function, we can pass different
\texttt{N\_trials} values and re-run the calibration+recovery pipeline
at each design point. The faceted identity plots below show the entire
prior-weighted recovery profile for \(N \in \{30, 60, 120, 240\}\),
directly comparable to the grid facets in Phase 6a but without the
restriction to three anchor values.

This is a Precision Analysis, the Bayesian analogue of a frequentist
power analysis, with three important distinctions:

\begin{enumerate}
\def\labelenumi{\arabic{enumi}.}
\item
  \textbf{No asymptotic approximations.} We simulate from the exact
  model.
\item
  \textbf{No point null.} The ``true effect'' is drawn from the prior,
  characterising performance over the full realistic range.
\item
  \textbf{Precision, not rejection rate.} The outcome metric is credible
  interval width, not whether a \(p\)-value crosses a threshold.
\end{enumerate}

\begin{Shaded}
\begin{Highlighting}[]
\CommentTok{\# {-}{-}{-}{-}{-}{-}{-}{-}{-}{-}{-}{-}{-}{-}{-}{-}{-}{-}{-}{-}{-}{-}{-}{-}{-}{-}{-}{-}{-}{-}{-}{-}{-}{-}{-}{-}{-}{-}{-}{-}{-}{-}{-}{-}{-}{-}{-}{-}{-}{-}{-}{-}{-}{-}{-}{-}{-}{-}{-}{-}{-}{-}{-}{-}{-}{-}{-}{-}{-}{-}{-}{-}{-}{-}{-}}
\CommentTok{\# precision{-}analysis generator: same logic as sbc\_generator\_fn above,}
\CommentTok{\# but parameterised on N\_trials so we can loop over the design grid.}
\CommentTok{\# {-}{-}{-}{-}{-}{-}{-}{-}{-}{-}{-}{-}{-}{-}{-}{-}{-}{-}{-}{-}{-}{-}{-}{-}{-}{-}{-}{-}{-}{-}{-}{-}{-}{-}{-}{-}{-}{-}{-}{-}{-}{-}{-}{-}{-}{-}{-}{-}{-}{-}{-}{-}{-}{-}{-}{-}{-}{-}{-}{-}{-}{-}{-}{-}{-}{-}{-}{-}{-}{-}{-}{-}{-}{-}{-}}
\NormalTok{make\_sbc\_precision\_generator }\OtherTok{\textless{}{-}} \ControlFlowTok{function}\NormalTok{(N\_trials) \{}
  \ControlFlowTok{function}\NormalTok{() \{}
\NormalTok{    theta\_logit\_true }\OtherTok{\textless{}{-}} \FunctionTok{rnorm}\NormalTok{(}\DecValTok{1}\NormalTok{, }\AttributeTok{mean =} \DecValTok{0}\NormalTok{, }\AttributeTok{sd =} \FloatTok{1.5}\NormalTok{)}
\NormalTok{    theta\_true       }\OtherTok{\textless{}{-}} \FunctionTok{plogis}\NormalTok{(theta\_logit\_true)}
\NormalTok{    y\_sim            }\OtherTok{\textless{}{-}} \FunctionTok{rbinom}\NormalTok{(N\_trials, }\AttributeTok{size =} \DecValTok{1}\NormalTok{, }\AttributeTok{prob =}\NormalTok{ theta\_true)}

    \FunctionTok{list}\NormalTok{(}
      \AttributeTok{variables =} \FunctionTok{list}\NormalTok{(}
        \AttributeTok{theta\_logit =}\NormalTok{ theta\_logit\_true,}
        \AttributeTok{theta       =}\NormalTok{ theta\_true}
\NormalTok{      ),}
      \AttributeTok{generated =} \FunctionTok{list}\NormalTok{(}
        \AttributeTok{N\_trials          =}\NormalTok{ N\_trials,}
        \AttributeTok{y                 =}\NormalTok{ y\_sim,}
        \AttributeTok{prior\_theta\_mu    =} \DecValTok{0}\NormalTok{,}
        \AttributeTok{prior\_theta\_sigma =} \FloatTok{1.5}\NormalTok{,}
        \AttributeTok{run\_diagnostics   =} \DecValTok{0}
\NormalTok{      )}
\NormalTok{    )}
\NormalTok{  \}}
\NormalTok{\}}

\CommentTok{\# Grid of experimental designs to evaluate}
\NormalTok{trial\_grid\_precision }\OtherTok{\textless{}{-}} \FunctionTok{c}\NormalTok{(}\DecValTok{30}\NormalTok{, }\DecValTok{60}\NormalTok{, }\DecValTok{120}\NormalTok{, }\DecValTok{240}\NormalTok{)}

\CommentTok{\# Accumulate recovery stats across all design conditions}
\NormalTok{precision\_stats }\OtherTok{\textless{}{-}}\NormalTok{ purrr}\SpecialCharTok{::}\FunctionTok{map\_dfr}\NormalTok{(trial\_grid\_precision, }\ControlFlowTok{function}\NormalTok{(n\_t) \{}

  \FunctionTok{cat}\NormalTok{(}\FunctionTok{sprintf}\NormalTok{(}\StringTok{"  precision analysis: N\_trials = \%d ...}\SpecialCharTok{\textbackslash{}n}\StringTok{"}\NormalTok{, n\_t))}

\NormalTok{  precision\_file }\OtherTok{\textless{}{-}} \FunctionTok{file.path}\NormalTok{(stan\_results\_dir,}
                          \FunctionTok{sprintf}\NormalTok{(}\StringTok{"ch07\_sbc\_precision\_n\%d.rds"}\NormalTok{, n\_t))}

  \ControlFlowTok{if}\NormalTok{ (regenerate\_simulations }\SpecialCharTok{||} \SpecialCharTok{!}\FunctionTok{file.exists}\NormalTok{(precision\_file)) \{}

\NormalTok{    gen\_precision }\OtherTok{\textless{}{-}}\NormalTok{ SBC}\SpecialCharTok{::}\FunctionTok{SBC\_generator\_function}\NormalTok{(}\FunctionTok{make\_sbc\_precision\_generator}\NormalTok{(n\_t))}

    \CommentTok{\# 200 datasets per condition — enough to characterise the recovery envelope}
\NormalTok{    dsets\_precision  }\OtherTok{\textless{}{-}}\NormalTok{ SBC}\SpecialCharTok{::}\FunctionTok{generate\_datasets}\NormalTok{(gen\_precision, }\AttributeTok{n\_sims =} \DecValTok{200}\NormalTok{)}
\NormalTok{    result\_precision }\OtherTok{\textless{}{-}}\NormalTok{ SBC}\SpecialCharTok{::}\FunctionTok{compute\_SBC}\NormalTok{(}
\NormalTok{      dsets\_precision,}
\NormalTok{      sbc\_backend,   }\CommentTok{\# reuse the compiled backend from Phase 6b}
      \AttributeTok{keep\_fits =} \ConstantTok{FALSE}
\NormalTok{    )}
    \FunctionTok{saveRDS}\NormalTok{(}\FunctionTok{list}\NormalTok{(}\AttributeTok{datasets =}\NormalTok{ dsets\_precision, }\AttributeTok{results =}\NormalTok{ result\_precision), precision\_file)}

\NormalTok{  \} }\ControlFlowTok{else}\NormalTok{ \{}
\NormalTok{    obj\_precision    }\OtherTok{\textless{}{-}} \FunctionTok{readRDS}\NormalTok{(precision\_file)}
\NormalTok{    result\_precision }\OtherTok{\textless{}{-}}\NormalTok{ obj\_precision}\SpecialCharTok{$}\NormalTok{results}
\NormalTok{  \}}

\NormalTok{  result\_precision}\SpecialCharTok{$}\NormalTok{stats }\SpecialCharTok{|\textgreater{}}
\NormalTok{    dplyr}\SpecialCharTok{::}\FunctionTok{filter}\NormalTok{(variable }\SpecialCharTok{==} \StringTok{"theta"}\NormalTok{) }\SpecialCharTok{|\textgreater{}}
\NormalTok{    dplyr}\SpecialCharTok{::}\FunctionTok{mutate}\NormalTok{(}\AttributeTok{n\_trials =}\NormalTok{ n\_t)}
\NormalTok{\})}
\end{Highlighting}
\end{Shaded}

\begin{verbatim}
  precision analysis: N_trials = 30 ...
  precision analysis: N_trials = 60 ...
  precision analysis: N_trials = 120 ...
  precision analysis: N_trials = 240 ...
\end{verbatim}

\begin{Shaded}
\begin{Highlighting}[]
\CommentTok{\# {-}{-}{-}{-}{-}{-}{-}{-}{-}{-}{-}{-}{-}{-}{-}{-}{-}{-}{-}{-}{-}{-}{-}{-}{-}{-}{-}{-}{-}{-}{-}{-}{-}{-}{-}{-}{-}{-}{-}{-}{-}{-}{-}{-}{-}{-}{-}{-}{-}{-}{-}{-}{-}{-}{-}{-}{-}{-}{-}{-}{-}{-}{-}{-}{-}{-}{-}{-}{-}{-}{-}{-}{-}{-}{-}}
\CommentTok{\# Summary table: key precision metrics as a function of N}
\CommentTok{\# {-}{-}{-}{-}{-}{-}{-}{-}{-}{-}{-}{-}{-}{-}{-}{-}{-}{-}{-}{-}{-}{-}{-}{-}{-}{-}{-}{-}{-}{-}{-}{-}{-}{-}{-}{-}{-}{-}{-}{-}{-}{-}{-}{-}{-}{-}{-}{-}{-}{-}{-}{-}{-}{-}{-}{-}{-}{-}{-}{-}{-}{-}{-}{-}{-}{-}{-}{-}{-}{-}{-}{-}{-}{-}{-}}
\NormalTok{precision\_summary }\OtherTok{\textless{}{-}}\NormalTok{ precision\_stats }\SpecialCharTok{|\textgreater{}}
\NormalTok{  dplyr}\SpecialCharTok{::}\FunctionTok{group\_by}\NormalTok{(n\_trials) }\SpecialCharTok{|\textgreater{}}
\NormalTok{  dplyr}\SpecialCharTok{::}\FunctionTok{summarize}\NormalTok{(}
    \AttributeTok{median\_CI\_width =} \FunctionTok{median}\NormalTok{(q95 }\SpecialCharTok{{-}}\NormalTok{ q5),}
    \AttributeTok{coverage90      =} \FunctionTok{mean}\NormalTok{(simulated\_value }\SpecialCharTok{\textgreater{}=}\NormalTok{ q5 }\SpecialCharTok{\&}\NormalTok{ simulated\_value }\SpecialCharTok{\textless{}=}\NormalTok{ q95),}
    \AttributeTok{RMSE            =} \FunctionTok{sqrt}\NormalTok{(}\FunctionTok{mean}\NormalTok{((simulated\_value }\SpecialCharTok{{-}}\NormalTok{ median)}\SpecialCharTok{\^{}}\DecValTok{2}\NormalTok{)),}
    \AttributeTok{.groups         =} \StringTok{"drop"}
\NormalTok{  )}

\FunctionTok{cat}\NormalTok{(}\StringTok{"=== PRECISION ANALYSIS: RECOVERY QUALITY VS. N\_TRIALS ===}\SpecialCharTok{\textbackslash{}n}\StringTok{"}\NormalTok{)}
\end{Highlighting}
\end{Shaded}

\begin{verbatim}
=== PRECISION ANALYSIS: RECOVERY QUALITY VS. N_TRIALS ===
\end{verbatim}

\begin{Shaded}
\begin{Highlighting}[]
\FunctionTok{print}\NormalTok{(}
\NormalTok{  knitr}\SpecialCharTok{::}\FunctionTok{kable}\NormalTok{(}
\NormalTok{    precision\_summary,}
    \AttributeTok{digits    =} \DecValTok{3}\NormalTok{,}
    \AttributeTok{col.names =} \FunctionTok{c}\NormalTok{(}\StringTok{"N Trials"}\NormalTok{, }\StringTok{"Median 90\% CI Width"}\NormalTok{,}
                  \StringTok{"Empirical 90\% Coverage"}\NormalTok{, }\StringTok{"RMSE"}\NormalTok{)}
\NormalTok{  )}
\NormalTok{)}
\end{Highlighting}
\end{Shaded}

\begin{verbatim}


| N Trials| Median 90% CI Width| Empirical 90% Coverage|  RMSE|
|--------:|-------------------:|----------------------:|-----:|
|       30|               0.251|                  0.890| 0.079|
|       60|               0.183|                  0.835| 0.060|
|      120|               0.134|                  0.945| 0.031|
|      240|               0.095|                  0.930| 0.026|
\end{verbatim}

\begin{Shaded}
\begin{Highlighting}[]
\CommentTok{\# {-}{-}{-}{-}{-}{-}{-}{-}{-}{-}{-}{-}{-}{-}{-}{-}{-}{-}{-}{-}{-}{-}{-}{-}{-}{-}{-}{-}{-}{-}{-}{-}{-}{-}{-}{-}{-}{-}{-}{-}{-}{-}{-}{-}{-}{-}{-}{-}{-}{-}{-}{-}{-}{-}{-}{-}{-}{-}{-}{-}{-}{-}{-}{-}{-}{-}{-}{-}{-}{-}{-}{-}{-}{-}{-}}
\CommentTok{\# Faceted identity plot: one panel per trial count}
\CommentTok{\# {-}{-}{-}{-}{-}{-}{-}{-}{-}{-}{-}{-}{-}{-}{-}{-}{-}{-}{-}{-}{-}{-}{-}{-}{-}{-}{-}{-}{-}{-}{-}{-}{-}{-}{-}{-}{-}{-}{-}{-}{-}{-}{-}{-}{-}{-}{-}{-}{-}{-}{-}{-}{-}{-}{-}{-}{-}{-}{-}{-}{-}{-}{-}{-}{-}{-}{-}{-}{-}{-}{-}{-}{-}{-}{-}}
\FunctionTok{ggplot}\NormalTok{(precision\_stats, }\FunctionTok{aes}\NormalTok{(}\AttributeTok{x =}\NormalTok{ simulated\_value, }\AttributeTok{y =}\NormalTok{ median)) }\SpecialCharTok{+}
  \FunctionTok{geom\_abline}\NormalTok{(}\AttributeTok{intercept =} \DecValTok{0}\NormalTok{, }\AttributeTok{slope =} \DecValTok{1}\NormalTok{,}
              \AttributeTok{linetype =} \StringTok{"dashed"}\NormalTok{, }\AttributeTok{color =} \StringTok{"gray50"}\NormalTok{, }\AttributeTok{linewidth =} \FloatTok{0.7}\NormalTok{) }\SpecialCharTok{+}
  \FunctionTok{geom\_errorbar}\NormalTok{(}\FunctionTok{aes}\NormalTok{(}\AttributeTok{ymin =}\NormalTok{ q5, }\AttributeTok{ymax =}\NormalTok{ q95),}
                \AttributeTok{width =} \DecValTok{0}\NormalTok{, }\AttributeTok{alpha =} \FloatTok{0.20}\NormalTok{, }\AttributeTok{color =} \StringTok{"steelblue"}\NormalTok{) }\SpecialCharTok{+}
  \FunctionTok{geom\_point}\NormalTok{(}\AttributeTok{color =} \StringTok{"midnightblue"}\NormalTok{, }\AttributeTok{alpha =} \FloatTok{0.45}\NormalTok{, }\AttributeTok{size =} \FloatTok{0.9}\NormalTok{) }\SpecialCharTok{+}
  \FunctionTok{coord\_fixed}\NormalTok{(}\AttributeTok{xlim =} \FunctionTok{c}\NormalTok{(}\DecValTok{0}\NormalTok{, }\DecValTok{1}\NormalTok{), }\AttributeTok{ylim =} \FunctionTok{c}\NormalTok{(}\DecValTok{0}\NormalTok{, }\DecValTok{1}\NormalTok{)) }\SpecialCharTok{+}
  \FunctionTok{facet\_wrap}\NormalTok{(}\SpecialCharTok{\textasciitilde{}}\FunctionTok{paste0}\NormalTok{(}\StringTok{"N = "}\NormalTok{, n\_trials), }\AttributeTok{nrow =} \DecValTok{1}\NormalTok{) }\SpecialCharTok{+}
  \FunctionTok{labs}\NormalTok{(}
    \AttributeTok{title    =} \StringTok{"Precision Analysis: Recovery Quality vs. Number of Trials"}\NormalTok{,}
    \AttributeTok{subtitle =} \FunctionTok{paste0}\NormalTok{(}
      \StringTok{"Each panel: 200 independent SBC universes. "}\NormalTok{,}
      \StringTok{"Bars are 90\% credible intervals.}\SpecialCharTok{\textbackslash{}n}\StringTok{"}\NormalTok{,}
      \StringTok{"Tighter scatter around the diagonal = higher estimation precision."}
\NormalTok{    ),}
    \AttributeTok{x =} \FunctionTok{expression}\NormalTok{(}\StringTok{"True "} \SpecialCharTok{*}\NormalTok{ theta }\SpecialCharTok{*} \StringTok{" (drawn from prior)"}\NormalTok{),}
    \AttributeTok{y =} \StringTok{"Posterior Median"}
\NormalTok{  ) }\SpecialCharTok{+}
  \FunctionTok{theme\_classic}\NormalTok{() }\SpecialCharTok{+}
  \FunctionTok{theme}\NormalTok{(}\AttributeTok{strip.background =} \FunctionTok{element\_rect}\NormalTok{(}\AttributeTok{fill =} \StringTok{"gray95"}\NormalTok{, }\AttributeTok{color =} \ConstantTok{NA}\NormalTok{))}
\end{Highlighting}
\end{Shaded}

\begin{center}
\includegraphics[width=0.8\linewidth,height=\textheight,keepaspectratio]{05-ModelQualityChecks_files/figure-pdf/ch06_sbc_precision_analysis-1.pdf}
\end{center}

Read the precision table as follows. The \textbf{median CI width} column
measures how precisely a typical experiment estimates \(\theta\) at a
given \(N\): a width of \(0.3\) means the 90\% credible interval spans
roughly \(0.3\) probability units, which may or may not be adequate
depending on the scientific question. The \textbf{empirical 90\%
coverage} column should remain near \(0.90\) at every \(N\) --- if
coverage collapses at small \(N\), the posterior is becoming
overconfident relative to the calibration guarantee we verified above.
The \textbf{RMSE} column captures average estimation error; halving it
requires roughly quadrupling \(N\), consistent with the classical
\(1/\sqrt{N}\) rate.

Comparing the faceted identity plots here to the pointrange plots in
Phase 6a, the two visualisations are complementary. Phase 6a shows the
shrinkage behaviour at three \emph{specific} anchor values with clearly
labelled coordinates --- pedagogically invaluable for understanding how
prior pull operates at the boundary (\(\theta = 0.9\)) versus at the
centre (\(\theta = 0.5\)). The SBC precision plots show the full
recovery \emph{envelope} across the prior's probability mass, making it
impossible to cherry-pick a favourable test point. Together they
constitute a complete characterisation of the model's identifiability as
a function of experimental design.

\subsection{Phase 6c: Posterior SBC}\label{phase-6c-posterior-sbc}

However, if you paid attention, you've realized that the calibration was
based on the prior. In other words, it's strongly bound by our
assumptions (and we should be careful about that). But what if the
posterior brings us to a slightly different region of the parameter
space? We can also check the calibration of the posterior by running SBC
with a different generator function that samples from the posterior
instead of the prior. This answers the question: ``Given the constraints
the actual data has placed on my parameters, is my model-sampler
combination still unbiased in this specific high-probability manifold?''

\subsubsection{The Mathematical Logic of
PSBC}\label{the-mathematical-logic-of-psbc}

The self-consistency property of Bayes' theorem extends to any valid
distribution. In PSBC, we treat our empirical posterior as a ``local
prior'':

\begin{enumerate}
\def\labelenumi{\arabic{enumi}.}
\item
  Draw a ``true'' parameter \(\tilde{\theta}\) from the fitted
  posterior: \(\tilde{\theta} \sim p(\theta \mid y_{obs})\).
\item
  Simulate a new synthetic dataset \(\tilde{y}\) from the data model:
  \(\tilde{y} \sim p(y \mid \tilde{\theta})\).
\item
  Fit the model to this synthetic data to get a ``new'' posterior
  \(p(\theta \mid \tilde{y}, y_{obs})\).
\item
  Check Ranks: The rank of \(\tilde{\theta}\) within the new posterior
  must still be Uniform.
\end{enumerate}

If PSBC passes but Prior SBC fails, it means your model is only reliable
for data similar to what you observed. If PSBC fails, you cannot trust
the specific estimates for your current experiment, even if the model
works ``in general''

\begin{Shaded}
\begin{Highlighting}[]
\CommentTok{\# 1. Extract posterior draws as the required \textquotesingle{}draws\_matrix\textquotesingle{} type}
\NormalTok{post\_draws\_theta\_raw }\OtherTok{\textless{}{-}}\NormalTok{ posterior}\SpecialCharTok{::}\FunctionTok{as\_draws\_matrix}\NormalTok{(fit\_gq}\SpecialCharTok{$}\FunctionTok{draws}\NormalTok{(}\FunctionTok{c}\NormalTok{(}\StringTok{"theta"}\NormalTok{, }\StringTok{"theta\_logit"}\NormalTok{)))}

\CommentTok{\# 2. Sequential Data Generation}
\CommentTok{\# We manually draw \textquotesingle{}true\textquotesingle{} values from the posterior instead of a formula}
\NormalTok{generate\_psbc\_data\_theta }\OtherTok{\textless{}{-}} \ControlFlowTok{function}\NormalTok{(}\AttributeTok{n\_sims =} \DecValTok{1000}\NormalTok{) \{}
\NormalTok{  indices }\OtherTok{\textless{}{-}} \FunctionTok{sample}\NormalTok{(}\DecValTok{1}\SpecialCharTok{:}\FunctionTok{nrow}\NormalTok{(post\_draws\_theta\_raw), n\_sims, }\AttributeTok{replace =} \ConstantTok{TRUE}\NormalTok{)}
\NormalTok{  sampled\_vars }\OtherTok{\textless{}{-}}\NormalTok{ post\_draws\_theta\_raw[indices, , drop }\OtherTok{=} \ConstantTok{FALSE}\NormalTok{]}
  
\NormalTok{  gen\_list }\OtherTok{\textless{}{-}} \FunctionTok{list}\NormalTok{()}
  \ControlFlowTok{for}\NormalTok{ (i }\ControlFlowTok{in} \DecValTok{1}\SpecialCharTok{:}\NormalTok{n\_sims) \{}
\NormalTok{    true\_logit }\OtherTok{\textless{}{-}} \FunctionTok{as.numeric}\NormalTok{(sampled\_vars[i, }\StringTok{"theta\_logit"}\NormalTok{])}
\NormalTok{    true\_prob }\OtherTok{\textless{}{-}} \FunctionTok{plogis}\NormalTok{(true\_logit) }\CommentTok{\# Logistic transformation}
    
    \CommentTok{\# Simulate data based on these posterior{-}conditioned parameters }
\NormalTok{    y\_sim }\OtherTok{\textless{}{-}} \FunctionTok{rbinom}\NormalTok{(observed\_data\_list}\SpecialCharTok{$}\NormalTok{N\_trials, }\AttributeTok{size =} \DecValTok{1}\NormalTok{, }\AttributeTok{prob =}\NormalTok{ true\_prob)}
    
\NormalTok{    gen\_list[[i]] }\OtherTok{\textless{}{-}} \FunctionTok{list}\NormalTok{(}
      \AttributeTok{N\_trials          =}\NormalTok{ observed\_data\_list}\SpecialCharTok{$}\NormalTok{N\_trials,}
      \AttributeTok{y                 =}\NormalTok{ y\_sim,}
      \AttributeTok{prior\_theta\_mu    =} \DecValTok{0}\NormalTok{,}
      \AttributeTok{prior\_theta\_sigma =} \FloatTok{1.5}\NormalTok{,}
      \AttributeTok{run\_diagnostics   =} \DecValTok{0}
\NormalTok{    )}
\NormalTok{  \}}
  
  \CommentTok{\# PSBC requires parameters in \textquotesingle{}variables\textquotesingle{} and Stan data in \textquotesingle{}generated\textquotesingle{}}
  \FunctionTok{return}\NormalTok{(SBC}\SpecialCharTok{::}\FunctionTok{SBC\_datasets}\NormalTok{(}\AttributeTok{variables =}\NormalTok{ sampled\_vars, }\AttributeTok{generated =}\NormalTok{ gen\_list))}
\NormalTok{\}}

\NormalTok{psbc\_results\_file }\OtherTok{\textless{}{-}} \FunctionTok{file.path}\NormalTok{(stan\_results\_dir, }\StringTok{"ch07\_psbc\_results.rds"}\NormalTok{)}

\ControlFlowTok{if}\NormalTok{ (regenerate\_simulations }\SpecialCharTok{||} \SpecialCharTok{!}\FunctionTok{file.exists}\NormalTok{(psbc\_results\_file)) \{}
\NormalTok{  psbc\_datasets\_theta }\OtherTok{\textless{}{-}} \FunctionTok{generate\_psbc\_data\_theta}\NormalTok{(}\AttributeTok{n\_sims =} \DecValTok{100}\NormalTok{)}
\NormalTok{  future}\SpecialCharTok{::}\FunctionTok{plan}\NormalTok{(future}\SpecialCharTok{::}\NormalTok{sequential)}
\NormalTok{  psbc\_results\_theta }\OtherTok{\textless{}{-}}\NormalTok{ SBC}\SpecialCharTok{::}\FunctionTok{compute\_SBC}\NormalTok{(psbc\_datasets\_theta, sbc\_backend)}
  \FunctionTok{saveRDS}\NormalTok{(psbc\_results\_theta, psbc\_results\_file)}
\NormalTok{\} }\ControlFlowTok{else}\NormalTok{ \{}
\NormalTok{  psbc\_results\_theta }\OtherTok{\textless{}{-}} \FunctionTok{readRDS}\NormalTok{(psbc\_results\_file)}
  \FunctionTok{cat}\NormalTok{(}\StringTok{"Loaded existing PSBC results from:"}\NormalTok{, psbc\_results\_file, }\StringTok{"}\SpecialCharTok{\textbackslash{}n}\StringTok{"}\NormalTok{)}
\NormalTok{\}}
\end{Highlighting}
\end{Shaded}

\begin{verbatim}
Loaded existing PSBC results from: simmodels/ch07_psbc_results.rds 
\end{verbatim}

\begin{Shaded}
\begin{Highlighting}[]
\CommentTok{\# Visualize Rank Histograms for local calibration}
\NormalTok{SBC}\SpecialCharTok{::}\FunctionTok{plot\_rank\_hist}\NormalTok{(psbc\_results\_theta, }\AttributeTok{variables =} \StringTok{"theta"}\NormalTok{) }\SpecialCharTok{+}
  \FunctionTok{ggtitle}\NormalTok{(}\StringTok{"Posterior SBC: Simple Biased Agent"}\NormalTok{) }\SpecialCharTok{+}
  \FunctionTok{theme\_classic}\NormalTok{()}
\end{Highlighting}
\end{Shaded}

\begin{center}
\includegraphics[width=0.8\linewidth,height=\textheight,keepaspectratio]{05-ModelQualityChecks_files/figure-pdf/ch06_posterior_sbc-1.pdf}
\end{center}

Once we have done all this work and we are satisfied it's time to either
move onto empirical data, or to increase the model complexity and make
sure that the new model also works. We'll leave empirical data for a
future chapter, so now let's move to a more complex model.

\section{Extending Checks to the Memory
Model}\label{extending-checks-to-the-memory-model}

The biased agent is a clean single-parameter model. Real cognitive
models are messier. This section applies the same six-phase workflow to
the memory model from Ch. 5 --- a model with two parameters
(\texttt{bias} and \texttt{beta}) where the choice probability on each
trial depends on the opponent's cumulative history.

The workflow will expose a genuine identifiability problem:
\texttt{bias} and \texttt{beta} are correlated in the posterior when the
opponent plays with a constant bias. We will diagnose why this happens,
show that it is a property of the \emph{experimental design} rather than
the model, and fix it by simulating data from a reversal learning
paradigm --- an opponent who switches preference halfway through the
session. The fixed dataset is then saved for use in Chapters 6 and 7.

This is the normal cycle of model development: build → validate →
discover failure → understand its cause → fix the design → validate
again.

When a model incorporates recursive memory --- the agent's internal
state at time \(t\) depends on time \(t-1\) --- the posterior geometry
becomes nonlinear and HMC engines can struggle if the priors are
misspecified or parameters are weakly identified. The checks below catch
all of this.

Before we can subject the Memory Model to our diagnostic battery, we
must remind ourself of its mathematical architecture and build its
isolated generative simulator in R. This simulator will serve as the
invariant engine for our Prior Predictives, Parameter Recovery, and SBC.

\begin{Shaded}
\begin{Highlighting}[]
\CommentTok{\# {-}{-}{-} The Social Memory Generative Engine {-}{-}{-}}
\CommentTok{\# }
\CommentTok{\#\textquotesingle{} Simulate a Social Memory Agent}
\CommentTok{\#\textquotesingle{} }
\CommentTok{\#\textquotesingle{} @param bias Numeric. Baseline log{-}odds of choosing Right.}
\CommentTok{\#\textquotesingle{} @param beta Numeric. Sensitivity to the opponent\textquotesingle{}s cumulative rate.}
\CommentTok{\#\textquotesingle{} @param opponent\_choices Integer Vector. The sequence of choices made by the \textquotesingle{}Other\textquotesingle{}.}
\CommentTok{\#\textquotesingle{} @return A tibble of the agent\textquotesingle{}s internal states and choices.}
\NormalTok{simulate\_memory }\OtherTok{\textless{}{-}} \ControlFlowTok{function}\NormalTok{(bias, beta, opponent\_choices) \{}
\NormalTok{  n\_trials }\OtherTok{\textless{}{-}} \FunctionTok{length}\NormalTok{(opponent\_choices)}
  
  \CommentTok{\# Calculate cumulative rate up to trial t{-}1 (Memory Input)}
  \CommentTok{\# We start with a neutral prior of 0.5}
\NormalTok{  opp\_rate }\OtherTok{\textless{}{-}} \FunctionTok{cumsum}\NormalTok{(opponent\_choices) }\SpecialCharTok{/} \FunctionTok{seq\_along}\NormalTok{(opponent\_choices)}
\NormalTok{  memory\_input }\OtherTok{\textless{}{-}}\NormalTok{ dplyr}\SpecialCharTok{::}\FunctionTok{lag}\NormalTok{(opp\_rate, }\AttributeTok{default =} \FloatTok{0.5}\NormalTok{)}
  
  \CommentTok{\# Enforce numerical stability for the logit transformation}
\NormalTok{  memory\_input }\OtherTok{\textless{}{-}} \FunctionTok{pmin}\NormalTok{(}\FunctionTok{pmax}\NormalTok{(memory\_input, }\FloatTok{0.01}\NormalTok{), }\FloatTok{0.99}\NormalTok{)}
  
  \CommentTok{\# Compute choice probabilities: logit(p) = bias + beta * logit(memory)}
  \CommentTok{\# We use qlogis() for the logit transformation}
\NormalTok{  theta\_probs }\OtherTok{\textless{}{-}} \FunctionTok{plogis}\NormalTok{(bias }\SpecialCharTok{+}\NormalTok{ beta }\SpecialCharTok{*} \FunctionTok{qlogis}\NormalTok{(memory\_input))}
  
  \CommentTok{\# Generate choices}
\NormalTok{  self\_choices }\OtherTok{\textless{}{-}} \FunctionTok{rbinom}\NormalTok{(n\_trials, }\DecValTok{1}\NormalTok{, theta\_probs)}
  
\NormalTok{  tibble}\SpecialCharTok{::}\FunctionTok{tibble}\NormalTok{(}
    \AttributeTok{trial =} \DecValTok{1}\SpecialCharTok{:}\NormalTok{n\_trials,}
    \AttributeTok{other\_choice =}\NormalTok{ opponent\_choices,}
    \AttributeTok{memory\_state =}\NormalTok{ memory\_input,}
    \AttributeTok{theta =}\NormalTok{ theta\_probs,}
    \AttributeTok{choice =}\NormalTok{ self\_choices}
\NormalTok{  )}
\NormalTok{\}}
\end{Highlighting}
\end{Shaded}

\subsection{2. Defining the Stan Model with Generated
Quantities:}\label{defining-the-stan-model-with-generated-quantities}

Now, we define the Stan model for the memory agent. The memory state is
calculated internally within the transformed parameters block based on
the opponent's choices (other). Crucially, we include simulations in the
generated quantities block for both prior and posterior predictive
checks, and calculate log-likelihood and joint log prior for sensitivity
analyses.

\begin{Shaded}
\begin{Highlighting}[]
\CommentTok{\# Stan model for Memory Agent with Generated Quantities}
\NormalTok{stan\_memory\_gq\_code }\OtherTok{\textless{}{-}} \StringTok{"}
\StringTok{// ch07\_MemoryBernoulli\_GQ.stan}
\StringTok{// Social Memory Agent: Estimates bias and memory sensitivity.}

\StringTok{data \{}
\StringTok{  int\textless{}lower=1\textgreater{} N\_trials;               // Number of trials}
\StringTok{  array[N\_trials] int y;               // Agent\textquotesingle{}s choices (0, 1)}
\StringTok{  array[N\_trials] int other;           // Opponent\textquotesingle{}s choices (0, 1)}
\StringTok{  }
\StringTok{  // Dynamic Prior Controls}
\StringTok{  real prior\_bias\_mu;}
\StringTok{  real\textless{}lower=0\textgreater{} prior\_bias\_sigma;}
\StringTok{  real prior\_beta\_mu;}
\StringTok{  real\textless{}lower=0\textgreater{} prior\_beta\_sigma;}
\StringTok{  int\textless{}lower=0, upper=1\textgreater{} run\_diagnostics;}
\StringTok{\}}

\StringTok{parameters \{}
\StringTok{  real bias; }
\StringTok{  real beta; }
\StringTok{\}}

\StringTok{transformed parameters \{}
\StringTok{  vector[N\_trials] memory;}
\StringTok{  memory[1] = 0.5; // Neutral start}
\StringTok{  }
\StringTok{  for (t in 2:N\_trials) \{}
\StringTok{    // Recursive update: running average of \textquotesingle{}other\textquotesingle{} choices}
\StringTok{    memory[t] = memory[t{-}1] + ((other[t{-}1] {-} memory[t{-}1]) / (t{-}1.0));}
\StringTok{    }
\StringTok{    // Boundary clipping for numerical stability of the logit link}
\StringTok{    memory[t] = fmax(0.01, fmin(0.99, memory[t]));}
\StringTok{  \}}
\StringTok{\}}

\StringTok{model \{}
\StringTok{  // Principled regularizing priors}
\StringTok{  target += normal\_lpdf(bias | prior\_bias\_mu, prior\_bias\_sigma);}
\StringTok{  target += normal\_lpdf(beta | prior\_beta\_mu, prior\_beta\_sigma);}

\StringTok{  // Likelihood function}
\StringTok{  // We use the vectorized bernoulli\_logit for computational efficiency}
\StringTok{  target += bernoulli\_logit\_lpmf(y | bias + beta * logit(memory));}
\StringTok{\}}

\StringTok{generated quantities \{}
\StringTok{  // {-}{-}{-} Prior Predictive Checks {-}{-}{-}}
\StringTok{  real bias\_prior = normal\_rng(prior\_bias\_mu, prior\_bias\_sigma);}
\StringTok{  real beta\_prior = normal\_rng(prior\_beta\_mu, prior\_beta\_sigma);}
\StringTok{  }
\StringTok{  // {-}{-}{-} Conditional Diagnostics {-}{-}{-}}
\StringTok{  array[N\_trials] int y\_prior\_rep;}
\StringTok{  int prior\_sum;}
\StringTok{  array[N\_trials] int y\_post\_rep;}
\StringTok{  int post\_sum;}
\StringTok{  vector[N\_trials] log\_lik;}

\StringTok{  if (run\_diagnostics) \{}
\StringTok{    for (t in 1:N\_trials) \{}
\StringTok{      y\_prior\_rep[t] = bernoulli\_logit\_rng(bias\_prior + beta\_prior * logit(memory[t]));}
\StringTok{    \}}
\StringTok{    // Summary Statistic: Cumulative Choice Rate}
\StringTok{    prior\_sum = sum(y\_prior\_rep);}

\StringTok{    // {-}{-}{-} Posterior Predictive Checks {-}{-}{-}}
\StringTok{    for (t in 1:N\_trials) \{}
\StringTok{      y\_post\_rep[t] = bernoulli\_logit\_rng(bias + beta * logit(memory[t]));}
\StringTok{    \}}
\StringTok{    post\_sum = sum(y\_post\_rep);}
\StringTok{    }
\StringTok{    for (t in 1:N\_trials) \{}
\StringTok{      log\_lik[t] = bernoulli\_logit\_lpmf(y[t] | bias + beta * logit(memory[t]));}
\StringTok{    \}}
\StringTok{  \}}

\StringTok{  // 2. Total Log{-}Prior (for Sensitivity/Update analysis)}
\StringTok{  real lprior = normal\_lpdf(bias | prior\_bias\_mu, prior\_bias\_sigma) + }
\StringTok{                normal\_lpdf(beta | prior\_beta\_mu, prior\_beta\_sigma);}
\StringTok{\}}
\StringTok{"}
\NormalTok{stan\_file\_mem\_gq }\OtherTok{\textless{}{-}} \FunctionTok{file.path}\NormalTok{(stan\_model\_dir, }\StringTok{"ch07\_Memory\_GQ.stan"}\NormalTok{)}
\FunctionTok{writeLines}\NormalTok{(stan\_memory\_gq\_code, stan\_file\_mem\_gq)}
\end{Highlighting}
\end{Shaded}

\subsection{Fitting and Generative
Validation}\label{fitting-and-generative-validation}

We execute the fit using the data generated by our invariant R
simulator. Notice we keep the priors relatively narrow to prevent the
``Beta'' parameter (memory weight) from creating pathological
deterministic choices during the warmup phase.

\begin{Shaded}
\begin{Highlighting}[]
\NormalTok{trials }\OtherTok{\textless{}{-}} \DecValTok{120}
\NormalTok{opponent\_choices }\OtherTok{\textless{}{-}} \FunctionTok{rbinom}\NormalTok{(trials, }\DecValTok{1}\NormalTok{, }\FloatTok{0.8}\NormalTok{) }\CommentTok{\# Opponent prefers \textquotesingle{}Right\textquotesingle{} 80\% of the time}

\NormalTok{d\_memory\_obs }\OtherTok{\textless{}{-}} \FunctionTok{simulate\_memory}\NormalTok{(}
  \AttributeTok{bias =} \DecValTok{0}\NormalTok{, }
  \AttributeTok{beta =} \FloatTok{0.9}\NormalTok{, }
  \AttributeTok{opponent\_choices =}\NormalTok{ opponent\_choices}
\NormalTok{)}

\CommentTok{\# 2. Prepare the data list for Stan}
\NormalTok{observed\_data\_list\_mem }\OtherTok{\textless{}{-}} \FunctionTok{list}\NormalTok{(}
  \AttributeTok{N\_trials =}\NormalTok{ N\_trials,}
  \AttributeTok{y =}\NormalTok{ d\_memory\_obs}\SpecialCharTok{$}\NormalTok{choice,}
  \AttributeTok{other =}\NormalTok{ d\_memory\_obs}\SpecialCharTok{$}\NormalTok{other\_choice,}
  \AttributeTok{prior\_bias\_mu =} \DecValTok{0}\NormalTok{, }\AttributeTok{prior\_bias\_sigma =} \FloatTok{0.5}\NormalTok{,}
  \AttributeTok{prior\_beta\_mu =} \DecValTok{0}\NormalTok{, }\AttributeTok{prior\_beta\_sigma =} \FloatTok{1.0}\NormalTok{,}
  \AttributeTok{run\_diagnostics =} \DecValTok{1}
\NormalTok{)}

\NormalTok{model\_file\_mem\_gq }\OtherTok{\textless{}{-}} \FunctionTok{file.path}\NormalTok{(stan\_results\_dir, }\StringTok{"ch07\_Memory\_GQ.rds"}\NormalTok{)}

\CommentTok{\# 3. Fit the model}
\ControlFlowTok{if}\NormalTok{ (regenerate\_simulations }\SpecialCharTok{||} \SpecialCharTok{!}\FunctionTok{file.exists}\NormalTok{(model\_file\_mem\_gq)) \{}
\NormalTok{  mod\_mem\_gq }\OtherTok{\textless{}{-}} \FunctionTok{cmdstan\_model}\NormalTok{(stan\_file\_mem\_gq)}

\NormalTok{  fit\_mem\_gq }\OtherTok{\textless{}{-}}\NormalTok{ mod\_mem\_gq}\SpecialCharTok{$}\FunctionTok{sample}\NormalTok{(}
    \AttributeTok{data =}\NormalTok{ observed\_data\_list\_mem,}
    \AttributeTok{chains =} \DecValTok{4}\NormalTok{,}
    \AttributeTok{parallel\_chains =} \FunctionTok{max}\NormalTok{(}\DecValTok{1}\NormalTok{, parallel}\SpecialCharTok{::}\FunctionTok{detectCores}\NormalTok{() }\SpecialCharTok{{-}} \DecValTok{2}\NormalTok{),}
    \AttributeTok{iter\_warmup =} \DecValTok{1000}\NormalTok{,}
    \AttributeTok{iter\_sampling =} \DecValTok{2000}\NormalTok{,}
    \AttributeTok{refresh =} \DecValTok{0}\NormalTok{,}
    \AttributeTok{adapt\_delta =} \FloatTok{0.95}\NormalTok{,     }\CommentTok{\# Increased for the nonlinear memory geometry}
    \AttributeTok{max\_treedepth =} \DecValTok{12}
\NormalTok{  )}
  \CommentTok{\#fit\_mem\_gq$save\_object(file = model\_file\_mem\_gq)}
\NormalTok{\} }\ControlFlowTok{else}\NormalTok{ \{}
\NormalTok{  fit\_mem\_gq }\OtherTok{\textless{}{-}} \FunctionTok{readRDS}\NormalTok{(model\_file\_mem\_gq)}
\NormalTok{\}}
\end{Highlighting}
\end{Shaded}

\begin{verbatim}
Running MCMC with 4 chains, at most 8 in parallel...

Chain 1 finished in 0.4 seconds.
Chain 2 finished in 0.4 seconds.
Chain 3 finished in 0.4 seconds.
Chain 4 finished in 0.4 seconds.

All 4 chains finished successfully.
Mean chain execution time: 0.4 seconds.
Total execution time: 0.6 seconds.
\end{verbatim}

\begin{Shaded}
\begin{Highlighting}[]
\CommentTok{\# {-}{-}{-} Visualizing the Cognitive State {-}{-}{-}}
\CommentTok{\# Let\textquotesingle{}s observe how the model \textquotesingle{}thinks\textquotesingle{} the memory state evolved}

\NormalTok{plot\_memory\_trace }\OtherTok{\textless{}{-}}\NormalTok{ d\_memory\_obs }\SpecialCharTok{\%\textgreater{}\%}
  \FunctionTok{ggplot}\NormalTok{(}\FunctionTok{aes}\NormalTok{(}\AttributeTok{x =}\NormalTok{ trial, }\AttributeTok{y =}\NormalTok{ memory\_state)) }\SpecialCharTok{+}
  \FunctionTok{geom\_line}\NormalTok{(}\AttributeTok{color =} \StringTok{"darkred"}\NormalTok{, }\AttributeTok{linewidth =} \DecValTok{1}\NormalTok{) }\SpecialCharTok{+}
  \FunctionTok{geom\_point}\NormalTok{(}\FunctionTok{aes}\NormalTok{(}\AttributeTok{y =}\NormalTok{ other\_choice), }\AttributeTok{alpha =} \FloatTok{0.2}\NormalTok{) }\SpecialCharTok{+}
  \FunctionTok{ggtitle}\NormalTok{(}\StringTok{"The Memory Trace: Tracked Opponent Choice Rate"}\NormalTok{) }\SpecialCharTok{+}
  \FunctionTok{theme\_classic}\NormalTok{()}

\FunctionTok{print}\NormalTok{(plot\_memory\_trace)}
\end{Highlighting}
\end{Shaded}

\begin{center}
\includegraphics[width=0.8\linewidth,height=\textheight,keepaspectratio]{05-ModelQualityChecks_files/figure-pdf/05-fit-memory-model-1.pdf}
\end{center}

The transition from a static bias to a dynamic memory sensitivity
(\(\beta\)) introduces nonlinearities in the posterior. Our first step
is to verify that the NUTS sampler navigated this geometry without
pathology.

\begin{Shaded}
\begin{Highlighting}[]
\NormalTok{fit\_mem\_gq}\SpecialCharTok{$}\FunctionTok{summary}\NormalTok{(}\FunctionTok{c}\NormalTok{(}\StringTok{"bias"}\NormalTok{, }\StringTok{"beta"}\NormalTok{))}
\end{Highlighting}
\end{Shaded}

\begin{verbatim}
# A tibble: 2 x 10
  variable  mean median    sd   mad     q5   q95  rhat ess_bulk ess_tail
  <chr>    <dbl>  <dbl> <dbl> <dbl>  <dbl> <dbl> <dbl>    <dbl>    <dbl>
1 bias     0.291  0.294 0.384 0.382 -0.346 0.926  1.00    1850.    2522.
2 beta     0.375  0.370 0.334 0.330 -0.174 0.933  1.00    1845.    2543.
\end{verbatim}

\begin{Shaded}
\begin{Highlighting}[]
\NormalTok{diagnostic\_summary }\OtherTok{\textless{}{-}}\NormalTok{ fit\_mem\_gq}\SpecialCharTok{$}\FunctionTok{diagnostic\_summary}\NormalTok{()}

\CommentTok{\# 1. Check for Divergent Transitions (Zero Tolerance)}
\NormalTok{divergences }\OtherTok{\textless{}{-}} \FunctionTok{sum}\NormalTok{(diagnostic\_summary}\SpecialCharTok{$}\NormalTok{num\_divergent)}
\ControlFlowTok{if}\NormalTok{ (divergences }\SpecialCharTok{\textgreater{}} \DecValTok{0}\NormalTok{) \{}
  \FunctionTok{stop}\NormalTok{(}\FunctionTok{paste}\NormalTok{(}\StringTok{"CRITICAL FAILURE:"}\NormalTok{, divergences, }\StringTok{"divergences detected. The posterior geometry is pathological. Reparameterization required. Do not interpret these results."}\NormalTok{))}
\NormalTok{\}}

\CommentTok{\# 2. Check E{-}BFMI (Energy{-}Bayesian Fraction of Missing Information)}
\CommentTok{\# Threshold is typically 0.3. Values below this indicate poor exploration.}
\NormalTok{min\_ebfmi }\OtherTok{\textless{}{-}} \FunctionTok{min}\NormalTok{(diagnostic\_summary}\SpecialCharTok{$}\NormalTok{ebfmi)}
\ControlFlowTok{if}\NormalTok{ (min\_ebfmi }\SpecialCharTok{\textless{}} \FloatTok{0.3}\NormalTok{) \{}
  \FunctionTok{warning}\NormalTok{(}\FunctionTok{paste}\NormalTok{(}\StringTok{"WARNING: Low E{-}BFMI detected (Min ="}\NormalTok{, }\FunctionTok{round}\NormalTok{(min\_ebfmi, }\DecValTok{3}\NormalTok{), }
                \StringTok{"). The sampler may struggle to explore the tails of the posterior."}\NormalTok{))}
\NormalTok{\}}

\CommentTok{\# 3. Check for Convergence (Rhat \textless{} 1.01)}
\NormalTok{max\_rhat }\OtherTok{\textless{}{-}} \FunctionTok{max}\NormalTok{(fit\_summary}\SpecialCharTok{$}\NormalTok{rhat, }\AttributeTok{na.rm =} \ConstantTok{TRUE}\NormalTok{)}
\ControlFlowTok{if}\NormalTok{ (max\_rhat }\SpecialCharTok{\textgreater{}} \FloatTok{1.01}\NormalTok{) \{}
  \FunctionTok{stop}\NormalTok{(}\FunctionTok{paste}\NormalTok{(}\StringTok{"CRITICAL FAILURE: Chains did not converge. Max Rhat ="}\NormalTok{, }\FunctionTok{round}\NormalTok{(max\_rhat, }\DecValTok{3}\NormalTok{)))}
\NormalTok{\}}

\CommentTok{\# 4. Check Effective Sample Size (ESS \textgreater{} 400 for both bulk and tail)}
\NormalTok{min\_ess\_bulk }\OtherTok{\textless{}{-}} \FunctionTok{min}\NormalTok{(fit\_summary}\SpecialCharTok{$}\NormalTok{ess\_bulk, }\AttributeTok{na.rm =} \ConstantTok{TRUE}\NormalTok{)}
\NormalTok{min\_ess\_tail }\OtherTok{\textless{}{-}} \FunctionTok{min}\NormalTok{(fit\_summary}\SpecialCharTok{$}\NormalTok{ess\_tail, }\AttributeTok{na.rm =} \ConstantTok{TRUE}\NormalTok{)}
\ControlFlowTok{if}\NormalTok{ (min\_ess\_bulk }\SpecialCharTok{\textless{}} \DecValTok{400} \SpecialCharTok{||}\NormalTok{ min\_ess\_tail }\SpecialCharTok{\textless{}} \DecValTok{400}\NormalTok{) \{}
  \FunctionTok{warning}\NormalTok{(}\StringTok{"WARNING: Insufficient Effective Sample Size. Autocorrelation is too high. Run longer chains."}\NormalTok{)}
\NormalTok{\}}
\end{Highlighting}
\end{Shaded}

\subsubsection{Prior and Posterior Predictive
Check:}\label{prior-and-posterior-predictive-check}

We now compare what the model expected before seeing the data versus
what it ``learned'' after seeing the agent's choices. We will use the
cumulative choice rate as our primary summary statistic.

\begin{Shaded}
\begin{Highlighting}[]
\CommentTok{\# 1. Extract Prior and Posterior Replications}
\NormalTok{draws\_mem }\OtherTok{\textless{}{-}} \FunctionTok{as\_draws\_df}\NormalTok{(fit\_mem\_gq}\SpecialCharTok{$}\FunctionTok{draws}\NormalTok{())}

\CommentTok{\# 2. Wrangle for Comparison}
\NormalTok{predictive\_df }\OtherTok{\textless{}{-}} \FunctionTok{tibble}\NormalTok{(}
  \AttributeTok{Value =} \FunctionTok{c}\NormalTok{(draws\_mem}\SpecialCharTok{$}\NormalTok{prior\_sum, draws\_mem}\SpecialCharTok{$}\NormalTok{post\_sum),}
  \AttributeTok{Type =} \FunctionTok{rep}\NormalTok{(}\FunctionTok{c}\NormalTok{(}\StringTok{"Prior Predictive"}\NormalTok{, }\StringTok{"Posterior Predictive"}\NormalTok{), }\AttributeTok{each =} \FunctionTok{nrow}\NormalTok{(draws\_mem))}
\NormalTok{)}

\CommentTok{\# 3. Visualize Predictive Overlap}

\FunctionTok{ggplot}\NormalTok{(predictive\_df, }\FunctionTok{aes}\NormalTok{(}\AttributeTok{x =}\NormalTok{ Value, }\AttributeTok{fill =}\NormalTok{ Type)) }\SpecialCharTok{+}
  \FunctionTok{geom\_histogram}\NormalTok{(}\AttributeTok{alpha =} \FloatTok{0.6}\NormalTok{, }\AttributeTok{color =} \StringTok{"white"}\NormalTok{) }\SpecialCharTok{+}
  \FunctionTok{geom\_vline}\NormalTok{(}\AttributeTok{xintercept =} \FunctionTok{sum}\NormalTok{(observed\_data\_list\_mem}\SpecialCharTok{$}\NormalTok{h), }
             \AttributeTok{linetype =} \StringTok{"dashed"}\NormalTok{, }\AttributeTok{linewidth =} \DecValTok{1}\NormalTok{) }\SpecialCharTok{+}
  \FunctionTok{scale\_fill\_manual}\NormalTok{(}\AttributeTok{values =} \FunctionTok{c}\NormalTok{(}\StringTok{"Prior Predictive"} \OtherTok{=} \StringTok{"gray70"}\NormalTok{, }\StringTok{"Posterior Predictive"} \OtherTok{=} \StringTok{"skyblue"}\NormalTok{)) }\SpecialCharTok{+}
  \FunctionTok{labs}\NormalTok{(}\AttributeTok{title =} \StringTok{"Predictive Check: Total \textquotesingle{}Right\textquotesingle{} Choices"}\NormalTok{,}
       \AttributeTok{x =} \StringTok{"Cumulative Count of Choices"}\NormalTok{, }\AttributeTok{y =} \StringTok{"Density"}\NormalTok{) }\SpecialCharTok{+}
  \FunctionTok{theme\_classic}\NormalTok{()}
\end{Highlighting}
\end{Shaded}

\begin{center}
\includegraphics[width=0.8\linewidth,height=\textheight,keepaspectratio]{05-ModelQualityChecks_files/figure-pdf/05-memory-predictive-checks-1.pdf}
\end{center}

We must investigate how our findings are shaped by our priors. We'll
first run a prior-posterior update check and then use the analytical
priorsense power-scaling to check both bias and beta.

\begin{Shaded}
\begin{Highlighting}[]
\NormalTok{true\_vals }\OtherTok{\textless{}{-}} \FunctionTok{tibble}\NormalTok{(}
  \AttributeTok{Parameter =} \FunctionTok{c}\NormalTok{(}\StringTok{"Bias"}\NormalTok{, }\StringTok{"Memory Strength (Beta)"}\NormalTok{),}
  \AttributeTok{Value =} \FunctionTok{c}\NormalTok{(}\DecValTok{0}\NormalTok{, }\FloatTok{0.9}\NormalTok{) }\CommentTok{\# true\_bias\_mem and true\_beta\_mem}
\NormalTok{)}

\CommentTok{\# Wrangle draws for comparison}
\NormalTok{update\_long }\OtherTok{\textless{}{-}}\NormalTok{ draws\_mem }\SpecialCharTok{\%\textgreater{}\%}
\NormalTok{  dplyr}\SpecialCharTok{::}\FunctionTok{select}\NormalTok{(bias, beta, bias\_prior, beta\_prior) }\SpecialCharTok{\%\textgreater{}\%}
  \FunctionTok{pivot\_longer}\NormalTok{(}\FunctionTok{everything}\NormalTok{(), }\AttributeTok{names\_to =} \StringTok{"param\_raw"}\NormalTok{, }\AttributeTok{values\_to =} \StringTok{"Value"}\NormalTok{) }\SpecialCharTok{\%\textgreater{}\%}
  \FunctionTok{mutate}\NormalTok{(}
    \AttributeTok{Distribution =} \FunctionTok{factor}\NormalTok{(}\FunctionTok{ifelse}\NormalTok{(}\FunctionTok{str\_detect}\NormalTok{(param\_raw, }\StringTok{"prior"}\NormalTok{), }\StringTok{"Prior"}\NormalTok{, }\StringTok{"Posterior"}\NormalTok{),}
                          \AttributeTok{levels =} \FunctionTok{c}\NormalTok{(}\StringTok{"Prior"}\NormalTok{, }\StringTok{"Posterior"}\NormalTok{)),}
    \AttributeTok{Parameter =} \FunctionTok{factor}\NormalTok{(}\FunctionTok{ifelse}\NormalTok{(}\FunctionTok{str\_detect}\NormalTok{(param\_raw, }\StringTok{"bias"}\NormalTok{), }\StringTok{"Bias"}\NormalTok{, }\StringTok{"Memory Strength (Beta)"}\NormalTok{),}
                       \AttributeTok{levels =} \FunctionTok{c}\NormalTok{(}\StringTok{"Bias"}\NormalTok{, }\StringTok{"Memory Strength (Beta)"}\NormalTok{))}
\NormalTok{  )}

\CommentTok{\# Plot with Ground Truth markers}
\NormalTok{p\_mem\_update }\OtherTok{\textless{}{-}} \FunctionTok{ggplot}\NormalTok{(update\_long, }\FunctionTok{aes}\NormalTok{(}\AttributeTok{x =}\NormalTok{ Value, }\AttributeTok{fill =}\NormalTok{ Distribution)) }\SpecialCharTok{+}
  \FunctionTok{geom\_density}\NormalTok{(}\AttributeTok{alpha =} \FloatTok{0.6}\NormalTok{, }\AttributeTok{color =} \StringTok{"white"}\NormalTok{) }\SpecialCharTok{+}
  \FunctionTok{geom\_vline}\NormalTok{(}\AttributeTok{data =}\NormalTok{ true\_vals, }\FunctionTok{aes}\NormalTok{(}\AttributeTok{xintercept =}\NormalTok{ Value), }
             \AttributeTok{linetype =} \StringTok{"dashed"}\NormalTok{, }\AttributeTok{linewidth =} \DecValTok{1}\NormalTok{, }\AttributeTok{color =} \StringTok{"black"}\NormalTok{) }\SpecialCharTok{+}
  \FunctionTok{facet\_wrap}\NormalTok{(}\SpecialCharTok{\textasciitilde{}}\NormalTok{Parameter, }\AttributeTok{scales =} \StringTok{"free"}\NormalTok{) }\SpecialCharTok{+}
  \FunctionTok{scale\_fill\_manual}\NormalTok{(}\AttributeTok{values =} \FunctionTok{c}\NormalTok{(}\StringTok{"Prior"} \OtherTok{=} \StringTok{"\#ffbaba"}\NormalTok{, }\StringTok{"Posterior"} \OtherTok{=} \StringTok{"skyblue"}\NormalTok{)) }\SpecialCharTok{+}
  \FunctionTok{labs}\NormalTok{(}\AttributeTok{title =} \StringTok{"Epistemic Update: Parameter Convergence"}\NormalTok{,}
       \AttributeTok{subtitle =} \StringTok{"Dashed line represents the true generative parameter"}\NormalTok{,}
       \AttributeTok{x =} \StringTok{"Parameter Value"}\NormalTok{, }\AttributeTok{y =} \StringTok{"Density"}\NormalTok{) }\SpecialCharTok{+}
  \FunctionTok{theme\_classic}\NormalTok{() }\SpecialCharTok{+} \FunctionTok{theme}\NormalTok{(}\AttributeTok{legend.position =} \StringTok{"bottom"}\NormalTok{)}

\FunctionTok{print}\NormalTok{(p\_mem\_update)}
\end{Highlighting}
\end{Shaded}

\begin{center}
\includegraphics[width=0.8\linewidth,height=\textheight,keepaspectratio]{05-ModelQualityChecks_files/figure-pdf/05 -prior-sensitivity-memory-1.pdf}
\end{center}

\begin{Shaded}
\begin{Highlighting}[]
\NormalTok{sensitivity\_mem }\OtherTok{\textless{}{-}}\NormalTok{ priorsense}\SpecialCharTok{::}\FunctionTok{powerscale\_sensitivity}\NormalTok{(fit\_mem\_gq, }\AttributeTok{variable =} \FunctionTok{c}\NormalTok{(}\StringTok{"bias"}\NormalTok{, }\StringTok{"beta"}\NormalTok{))}

\NormalTok{priorsense}\SpecialCharTok{::}\FunctionTok{powerscale\_plot\_dens}\NormalTok{(fit\_mem\_gq, }\AttributeTok{variables =} \FunctionTok{c}\NormalTok{(}\StringTok{"bias"}\NormalTok{, }\StringTok{"beta"}\NormalTok{)) }\SpecialCharTok{+}
  \FunctionTok{theme\_classic}\NormalTok{()}
\end{Highlighting}
\end{Shaded}

\begin{center}
\includegraphics[width=0.8\linewidth,height=\textheight,keepaspectratio]{05-ModelQualityChecks_files/figure-pdf/05 -prior-sensitivity-memory-2.pdf}
\end{center}

While the marginal updates show that we recovered the parameters, we
must also inspect the joint posterior distribution. If Bias and Beta are
highly correlated, our experimental design may not be strong enough to
pull these two cognitive strategies apart.

\begin{Shaded}
\begin{Highlighting}[]
\CommentTok{\# {-}{-}{-} Phase 10b: Parameter Correlation \& Joint Posterior {-}{-}{-}}

\CommentTok{\# 1. Extract the posterior draws (excluding prior simulations)}
\NormalTok{post\_draws\_mem }\OtherTok{\textless{}{-}}\NormalTok{ draws\_mem }\SpecialCharTok{\%\textgreater{}\%}
\NormalTok{  dplyr}\SpecialCharTok{::}\FunctionTok{select}\NormalTok{(bias, beta)}

\CommentTok{\# 2. Calculate the Pearson correlation coefficient}
\NormalTok{post\_cor }\OtherTok{\textless{}{-}} \FunctionTok{cor}\NormalTok{(post\_draws\_mem}\SpecialCharTok{$}\NormalTok{bias, post\_draws\_mem}\SpecialCharTok{$}\NormalTok{beta)}

\CommentTok{\# 3. Visualize the Joint Posterior Surface}
\NormalTok{p\_joint\_posterior }\OtherTok{\textless{}{-}} \FunctionTok{ggplot}\NormalTok{(post\_draws\_mem, }\FunctionTok{aes}\NormalTok{(}\AttributeTok{x =}\NormalTok{ bias, }\AttributeTok{y =}\NormalTok{ beta)) }\SpecialCharTok{+}
  \CommentTok{\# Add a 2D density surface to see the "shape" of the uncertainty}
  \FunctionTok{stat\_density\_2d}\NormalTok{(}\FunctionTok{aes}\NormalTok{(}\AttributeTok{fill =} \FunctionTok{after\_stat}\NormalTok{(level)), }\AttributeTok{geom =} \StringTok{"polygon"}\NormalTok{, }\AttributeTok{alpha =} \FloatTok{0.3}\NormalTok{) }\SpecialCharTok{+}
  \CommentTok{\# Add the individual MCMC points (thinned for clarity)}
  \FunctionTok{geom\_point}\NormalTok{(}\AttributeTok{alpha =} \FloatTok{0.1}\NormalTok{, }\AttributeTok{color =} \StringTok{"skyblue"}\NormalTok{, }\AttributeTok{size =} \FloatTok{0.5}\NormalTok{) }\SpecialCharTok{+}
  \CommentTok{\# Add the ground truth marker (Bias=0, Beta=0.9)}
  \FunctionTok{geom\_point}\NormalTok{(}\FunctionTok{aes}\NormalTok{(}\AttributeTok{x =} \DecValTok{0}\NormalTok{, }\AttributeTok{y =} \FloatTok{0.9}\NormalTok{), }\AttributeTok{color =} \StringTok{"black"}\NormalTok{, }\AttributeTok{shape =} \DecValTok{3}\NormalTok{, }\AttributeTok{size =} \DecValTok{5}\NormalTok{, }\AttributeTok{stroke =} \FloatTok{1.5}\NormalTok{) }\SpecialCharTok{+}
  
  \FunctionTok{scale\_fill\_viridis\_c}\NormalTok{(}\AttributeTok{option =} \StringTok{"mako"}\NormalTok{, }\AttributeTok{guide =} \StringTok{"none"}\NormalTok{) }\SpecialCharTok{+}
  
  \FunctionTok{labs}\NormalTok{(}
    \AttributeTok{title =} \StringTok{"Joint Posterior: Bias vs. Memory Strength"}\NormalTok{,}
    \AttributeTok{subtitle =} \FunctionTok{paste0}\NormalTok{(}\StringTok{"Posterior Correlation (r) = "}\NormalTok{, }\FunctionTok{round}\NormalTok{(post\_cor, }\DecValTok{3}\NormalTok{)),}
    \AttributeTok{x =} \FunctionTok{expression}\NormalTok{(}\StringTok{"Baseline Bias ("} \SpecialCharTok{*}\NormalTok{ bias }\SpecialCharTok{*} \StringTok{")"}\NormalTok{),}
    \AttributeTok{y =} \FunctionTok{expression}\NormalTok{(}\StringTok{"Memory Strength ("} \SpecialCharTok{*}\NormalTok{ beta }\SpecialCharTok{*} \StringTok{")"}\NormalTok{)}
\NormalTok{  ) }\SpecialCharTok{+}
  \FunctionTok{annotate}\NormalTok{(}\StringTok{"text"}\NormalTok{, }\AttributeTok{x =} \FloatTok{0.05}\NormalTok{, }\AttributeTok{y =} \FloatTok{0.95}\NormalTok{, }\AttributeTok{label =} \StringTok{"Ground Truth"}\NormalTok{, }\AttributeTok{fontface =} \StringTok{"bold"}\NormalTok{) }\SpecialCharTok{+}
  \FunctionTok{theme\_classic}\NormalTok{()}

\FunctionTok{print}\NormalTok{(p\_joint\_posterior)}
\end{Highlighting}
\end{Shaded}

\begin{center}
\includegraphics[width=0.8\linewidth,height=\textheight,keepaspectratio]{05-ModelQualityChecks_files/figure-pdf/ch06_memory_correlation-1.pdf}
\end{center}

This is a problem. How can we make this correlation disappear? We have
at least two options: 1. Play with the priors, make them stricter 2.
Increase the amount of data

As we kinda like the priors, we will now explicitly manipulate the
experimental design (\(N \in \{30, 60, 120, 240, 480\}\)). We want to
observe the ``cigar'' shaped posterior mathematically decouple into a
tight, independent sphere.

\begin{Shaded}
\begin{Highlighting}[]
\CommentTok{\# {-}{-}{-} Phase 12b: Asymptotic Identifiability \& Covariance Collapse {-}{-}{-}}

\FunctionTok{cat}\NormalTok{(}\StringTok{"Initiating Asymptotic Identifiability Grid...}\SpecialCharTok{\textbackslash{}n}\StringTok{"}\NormalTok{)}
\end{Highlighting}
\end{Shaded}

\begin{verbatim}
Initiating Asymptotic Identifiability Grid...
\end{verbatim}

\begin{Shaded}
\begin{Highlighting}[]
\CommentTok{\# 1. Define the experimental bounds}
\NormalTok{trial\_grid }\OtherTok{\textless{}{-}} \FunctionTok{c}\NormalTok{(}\DecValTok{30}\NormalTok{, }\DecValTok{60}\NormalTok{, }\DecValTok{120}\NormalTok{, }\DecValTok{240}\NormalTok{, }\DecValTok{480}\NormalTok{)}
\NormalTok{true\_bias\_mem }\OtherTok{\textless{}{-}} \FloatTok{0.0}
\NormalTok{true\_beta\_mem }\OtherTok{\textless{}{-}} \FloatTok{0.9}

\NormalTok{cov\_collapse\_file }\OtherTok{\textless{}{-}} \FunctionTok{file.path}\NormalTok{(stan\_results\_dir, }\StringTok{"ch07\_cov\_collapse\_draws.rds"}\NormalTok{)}

\ControlFlowTok{if}\NormalTok{ (regenerate\_simulations }\SpecialCharTok{||} \SpecialCharTok{!}\FunctionTok{file.exists}\NormalTok{(cov\_collapse\_file)) \{}
  
  \CommentTok{\# 2. Define the isolated recovery pipeline}
\NormalTok{  run\_cov\_recovery }\OtherTok{\textless{}{-}} \ControlFlowTok{function}\NormalTok{(n\_trials) \{}
    \CommentTok{\# Simulate the opponent (80\% Right preference)}
\NormalTok{    opp\_choices }\OtherTok{\textless{}{-}} \FunctionTok{rbinom}\NormalTok{(n\_trials, }\DecValTok{1}\NormalTok{, }\FloatTok{0.8}\NormalTok{)}
    
    \CommentTok{\# Simulate our agent using the invariant generator}
\NormalTok{    sim\_data }\OtherTok{\textless{}{-}} \FunctionTok{simulate\_memory}\NormalTok{(}
      \AttributeTok{bias =}\NormalTok{ true\_bias\_mem, }
      \AttributeTok{beta =}\NormalTok{ true\_beta\_mem, }
      \AttributeTok{opponent\_choices =}\NormalTok{ opp\_choices}
\NormalTok{    )}
    
    \CommentTok{\# Structure the data block for Stan}
\NormalTok{    rec\_data\_list }\OtherTok{\textless{}{-}} \FunctionTok{list}\NormalTok{(}
      \AttributeTok{n =}\NormalTok{ n\_trials,}
      \AttributeTok{h =}\NormalTok{ sim\_data}\SpecialCharTok{$}\NormalTok{choice,}
      \AttributeTok{other =}\NormalTok{ sim\_data}\SpecialCharTok{$}\NormalTok{other\_choice,}
      \AttributeTok{prior\_bias\_m =} \DecValTok{0}\NormalTok{, }\AttributeTok{prior\_bias\_s =} \FloatTok{0.5}\NormalTok{,}
      \AttributeTok{prior\_beta\_m =} \DecValTok{0}\NormalTok{, }\AttributeTok{prior\_beta\_s =} \FloatTok{1.0}
\NormalTok{    )}
    
    \CommentTok{\# Fit the generative model}
\NormalTok{    fit }\OtherTok{\textless{}{-}}\NormalTok{ mod\_mem\_gq}\SpecialCharTok{$}\FunctionTok{sample}\NormalTok{(}
      \AttributeTok{data =}\NormalTok{ rec\_data\_list,}
      \AttributeTok{chains =} \DecValTok{1}\NormalTok{,}
      \AttributeTok{parallel\_chains =} \DecValTok{1}\NormalTok{, }\CommentTok{\# Prevent core oversubscription}
      \AttributeTok{iter\_warmup =} \DecValTok{1000}\NormalTok{,}
      \AttributeTok{iter\_sampling =} \DecValTok{2000}\NormalTok{,}
      \AttributeTok{adapt\_delta =} \FloatTok{0.95}\NormalTok{,}
      \AttributeTok{refresh =} \DecValTok{0}\NormalTok{,}
      \AttributeTok{show\_messages =} \ConstantTok{FALSE}
\NormalTok{    )}
    
    \CommentTok{\# Extract strictly the parameters of interest}
\NormalTok{    draws\_df }\OtherTok{\textless{}{-}}\NormalTok{ fit}\SpecialCharTok{$}\FunctionTok{draws}\NormalTok{(}\FunctionTok{c}\NormalTok{(}\StringTok{"bias"}\NormalTok{, }\StringTok{"beta"}\NormalTok{), }\AttributeTok{format =} \StringTok{"df"}\NormalTok{)}
    
\NormalTok{    tibble}\SpecialCharTok{::}\FunctionTok{tibble}\NormalTok{(}
      \AttributeTok{trials =}\NormalTok{ n\_trials,}
      \AttributeTok{bias =}\NormalTok{ draws\_df}\SpecialCharTok{$}\NormalTok{bias,}
      \AttributeTok{beta =}\NormalTok{ draws\_df}\SpecialCharTok{$}\NormalTok{beta}
\NormalTok{    )}
\NormalTok{  \}}

  \CommentTok{\# 3. Execute mapping using parallelization}
\NormalTok{  future}\SpecialCharTok{::}\FunctionTok{plan}\NormalTok{(future}\SpecialCharTok{::}\NormalTok{multisession, }\AttributeTok{workers =} \DecValTok{2}\NormalTok{)}
  
\NormalTok{  cov\_draws }\OtherTok{\textless{}{-}}\NormalTok{ furrr}\SpecialCharTok{::}\FunctionTok{future\_map\_dfr}\NormalTok{(}
\NormalTok{    trial\_grid, }
\NormalTok{    run\_cov\_recovery, }
    \AttributeTok{.options =}\NormalTok{ furrr}\SpecialCharTok{::}\FunctionTok{furrr\_options}\NormalTok{(}\AttributeTok{seed =} \ConstantTok{TRUE}\NormalTok{, }\AttributeTok{stdout =} \ConstantTok{FALSE}\NormalTok{)}
\NormalTok{  )}
  
\NormalTok{  future}\SpecialCharTok{::}\FunctionTok{plan}\NormalTok{(future}\SpecialCharTok{::}\NormalTok{sequential)}
  
  \FunctionTok{saveRDS}\NormalTok{(cov\_draws, cov\_collapse\_file)}
  \FunctionTok{cat}\NormalTok{(}\StringTok{"Covariance collapse grid completed. Draws saved to:"}\NormalTok{, cov\_collapse\_file, }\StringTok{"}\SpecialCharTok{\textbackslash{}n}\StringTok{"}\NormalTok{)}
  
\NormalTok{\} }\ControlFlowTok{else}\NormalTok{ \{}
\NormalTok{  cov\_draws }\OtherTok{\textless{}{-}} \FunctionTok{readRDS}\NormalTok{(cov\_collapse\_file)}
  \FunctionTok{cat}\NormalTok{(}\StringTok{"Loaded existing covariance draws from:"}\NormalTok{, cov\_collapse\_file, }\StringTok{"}\SpecialCharTok{\textbackslash{}n}\StringTok{"}\NormalTok{)}
\NormalTok{\}}
\end{Highlighting}
\end{Shaded}

\begin{verbatim}
Loaded existing covariance draws from: simmodels/ch07_cov_collapse_draws.rds 
\end{verbatim}

\begin{Shaded}
\begin{Highlighting}[]
\CommentTok{\# 4. Compute the correlation coefficient (r) for each trial count}
\NormalTok{cor\_summary }\OtherTok{\textless{}{-}}\NormalTok{ cov\_draws }\SpecialCharTok{|\textgreater{}} 
\NormalTok{  dplyr}\SpecialCharTok{::}\FunctionTok{group\_by}\NormalTok{(trials) }\SpecialCharTok{|\textgreater{}} 
\NormalTok{  dplyr}\SpecialCharTok{::}\FunctionTok{summarize}\NormalTok{(}
    \AttributeTok{r\_val =} \FunctionTok{cor}\NormalTok{(bias, beta),}
    \AttributeTok{label =} \FunctionTok{paste0}\NormalTok{(}\StringTok{"r = "}\NormalTok{, }\FunctionTok{round}\NormalTok{(r\_val, }\DecValTok{2}\NormalTok{)),}
    \AttributeTok{.groups =} \StringTok{"drop"}
\NormalTok{  ) }\SpecialCharTok{|\textgreater{}} 
  \CommentTok{\# Create a highly structured facet label}
\NormalTok{  dplyr}\SpecialCharTok{::}\FunctionTok{mutate}\NormalTok{(}\AttributeTok{facet\_label =} \FunctionTok{factor}\NormalTok{(trials, }\AttributeTok{levels =}\NormalTok{ trial\_grid, }\AttributeTok{labels =} \FunctionTok{paste0}\NormalTok{(}\StringTok{"N = "}\NormalTok{, trial\_grid)))}

\CommentTok{\# Merge facet labels back into the main draws for ggplot}
\NormalTok{cov\_draws }\OtherTok{\textless{}{-}}\NormalTok{ cov\_draws }\SpecialCharTok{|\textgreater{}} 
\NormalTok{  dplyr}\SpecialCharTok{::}\FunctionTok{left\_join}\NormalTok{(cor\_summary }\SpecialCharTok{|\textgreater{}}\NormalTok{ dplyr}\SpecialCharTok{::}\FunctionTok{select}\NormalTok{(trials, facet\_label), }\AttributeTok{by =} \StringTok{"trials"}\NormalTok{)}

\CommentTok{\# 5. Visualize the geometric collapse of the covariance}

\NormalTok{p\_cov\_collapse }\OtherTok{\textless{}{-}} \FunctionTok{ggplot}\NormalTok{(cov\_draws, }\FunctionTok{aes}\NormalTok{(}\AttributeTok{x =}\NormalTok{ bias, }\AttributeTok{y =}\NormalTok{ beta)) }\SpecialCharTok{+}
  \CommentTok{\# Draw the density topology}
  \FunctionTok{stat\_density\_2d}\NormalTok{(}\FunctionTok{aes}\NormalTok{(}\AttributeTok{fill =} \FunctionTok{after\_stat}\NormalTok{(level)), }\AttributeTok{geom =} \StringTok{"polygon"}\NormalTok{, }\AttributeTok{alpha =} \FloatTok{0.6}\NormalTok{) }\SpecialCharTok{+}
  
  \CommentTok{\# The invariant generative ground truth}
  \FunctionTok{geom\_point}\NormalTok{(}\AttributeTok{x =}\NormalTok{ true\_bias\_mem, }\AttributeTok{y =}\NormalTok{ true\_beta\_mem, }
             \AttributeTok{color =} \StringTok{"black"}\NormalTok{, }\AttributeTok{shape =} \DecValTok{3}\NormalTok{, }\AttributeTok{size =} \DecValTok{3}\NormalTok{, }\AttributeTok{stroke =} \DecValTok{1}\NormalTok{) }\SpecialCharTok{+}
  
  \CommentTok{\# Add the correlation text dynamically to each facet}
  \FunctionTok{geom\_text}\NormalTok{(}\AttributeTok{data =}\NormalTok{ cor\_summary, }\FunctionTok{aes}\NormalTok{(}\AttributeTok{x =} \SpecialCharTok{{-}}\FloatTok{0.8}\NormalTok{, }\AttributeTok{y =} \FloatTok{2.5}\NormalTok{, }\AttributeTok{label =}\NormalTok{ label), }
            \AttributeTok{fontface =} \StringTok{"bold"}\NormalTok{, }\AttributeTok{color =} \StringTok{"darkred"}\NormalTok{, }\AttributeTok{inherit.aes =} \ConstantTok{FALSE}\NormalTok{) }\SpecialCharTok{+}
  
  \FunctionTok{facet\_wrap}\NormalTok{(}\SpecialCharTok{\textasciitilde{}}\NormalTok{ facet\_label, }\AttributeTok{nrow =} \DecValTok{1}\NormalTok{) }\SpecialCharTok{+}
  \FunctionTok{scale\_fill\_viridis\_c}\NormalTok{(}\AttributeTok{option =} \StringTok{"mako"}\NormalTok{, }\AttributeTok{guide =} \StringTok{"none"}\NormalTok{) }\SpecialCharTok{+}
  
  \FunctionTok{coord\_cartesian}\NormalTok{(}\AttributeTok{xlim =} \FunctionTok{c}\NormalTok{(}\SpecialCharTok{{-}}\FloatTok{1.0}\NormalTok{, }\FloatTok{1.0}\NormalTok{), }\AttributeTok{ylim =} \FunctionTok{c}\NormalTok{(}\SpecialCharTok{{-}}\FloatTok{1.0}\NormalTok{, }\FloatTok{3.0}\NormalTok{)) }\SpecialCharTok{+}
  
  \FunctionTok{xlab}\NormalTok{(}\FunctionTok{expression}\NormalTok{(}\StringTok{"Baseline Bias ("} \SpecialCharTok{*}\NormalTok{ bias }\SpecialCharTok{*} \StringTok{")"}\NormalTok{)) }\SpecialCharTok{+}
  \FunctionTok{ylab}\NormalTok{(}\FunctionTok{expression}\NormalTok{(}\StringTok{"Memory Strength ("} \SpecialCharTok{*}\NormalTok{ beta }\SpecialCharTok{*} \StringTok{")"}\NormalTok{)) }\SpecialCharTok{+}
  
  \FunctionTok{theme\_classic}\NormalTok{() }\SpecialCharTok{+}
  \FunctionTok{theme}\NormalTok{(}
    \AttributeTok{strip.background =} \FunctionTok{element\_rect}\NormalTok{(}\AttributeTok{fill =} \StringTok{"gray95"}\NormalTok{, }\AttributeTok{color =} \StringTok{"NA"}\NormalTok{),}
    \AttributeTok{strip.text =} \FunctionTok{element\_text}\NormalTok{(}\AttributeTok{face =} \StringTok{"bold"}\NormalTok{, }\AttributeTok{size =} \DecValTok{11}\NormalTok{),}
    \AttributeTok{panel.grid.major =} \FunctionTok{element\_line}\NormalTok{(}\AttributeTok{color =} \StringTok{"gray90"}\NormalTok{, }\AttributeTok{linetype =} \StringTok{"dotted"}\NormalTok{)}
\NormalTok{  )}

\FunctionTok{print}\NormalTok{(p\_cov\_collapse)}
\end{Highlighting}
\end{Shaded}

\begin{center}
\includegraphics[width=0.8\linewidth,height=\textheight,keepaspectratio]{05-ModelQualityChecks_files/figure-pdf/05-memory-covariance-collapse-1.pdf}
\end{center}

\textbf{Diagnosis:} The correlation between \texttt{bias} and
\texttt{beta} is not a numerical accident --- it is a structural
consequence of the experimental design. When the opponent plays with a
\emph{constant} bias (e.g.~always 80\% right), both parameters can
explain the agent's rightward tendency equally well:

\begin{itemize}
\tightlist
\item
  A high \texttt{beta} with neutral \texttt{bias} says ``the agent
  tracks the opponent's pattern and follows it.''
\item
  A high \texttt{bias} with low \texttt{beta} says ``the agent just
  always prefers right.''
\end{itemize}

Both stories are consistent with the data when the opponent never
changes. The model cannot tell them apart because the data provide no
discriminating evidence. Adding more trials of the same constant-bias
opponent does reduce the correlation somewhat (as the covariance
collapse grid showed), but not enough to make the parameters cleanly
identifiable.

\textbf{The fix --- reversal learning:} If the opponent \emph{switches}
preference halfway through the session (e.g.~80\% right for the first
120 trials, then 20\% right for the next 120), a pure-bias agent and a
memory agent make completely different predictions in the second half.
The pure-bias agent continues right; the memory agent flips. This
contrast is what discriminates \texttt{bias} from \texttt{beta}. We call
this a \textbf{reversal learning} paradigm.

We now simulate this design, fit the same memory model, and verify that
the \texttt{bias}/\texttt{beta} correlation collapses. We also save the
reversal dataset for use in Chapters 6 and 7.

\begin{Shaded}
\begin{Highlighting}[]
\FunctionTok{set.seed}\NormalTok{(}\DecValTok{42}\NormalTok{)}
\NormalTok{n\_reversal }\OtherTok{\textless{}{-}} \DecValTok{240}  \CommentTok{\# 120 trials per phase}

\CommentTok{\# {-}{-}{-} 1. Simulate opponent with preference reversal {-}{-}{-}}
\NormalTok{opponent\_reversal }\OtherTok{\textless{}{-}} \FunctionTok{c}\NormalTok{(}
  \FunctionTok{rbinom}\NormalTok{(}\DecValTok{120}\NormalTok{, }\DecValTok{1}\NormalTok{, }\FloatTok{0.8}\NormalTok{),  }\CommentTok{\# Phase 1: strongly right{-}biased}
  \FunctionTok{rbinom}\NormalTok{(}\DecValTok{120}\NormalTok{, }\DecValTok{1}\NormalTok{, }\FloatTok{0.2}\NormalTok{)   }\CommentTok{\# Phase 2: reversal — strongly left{-}biased}
\NormalTok{)}

\CommentTok{\# {-}{-}{-} 2. Simulate agent (true bias = 0, true beta = 0.9) {-}{-}{-}}
\NormalTok{d\_reversal }\OtherTok{\textless{}{-}} \FunctionTok{simulate\_memory}\NormalTok{(}
  \AttributeTok{bias             =} \DecValTok{0}\NormalTok{,}
  \AttributeTok{beta             =} \FloatTok{0.9}\NormalTok{,}
  \AttributeTok{opponent\_choices =}\NormalTok{ opponent\_reversal}
\NormalTok{)}

\CommentTok{\# {-}{-}{-} 3. Visualise the reversal to confirm the design {-}{-}{-}}
\NormalTok{p\_reversal\_design }\OtherTok{\textless{}{-}} \FunctionTok{tibble}\NormalTok{(}
  \AttributeTok{trial         =} \DecValTok{1}\SpecialCharTok{:}\NormalTok{n\_reversal,}
  \AttributeTok{opponent\_rate =} \FunctionTok{cumsum}\NormalTok{(opponent\_reversal) }\SpecialCharTok{/} \FunctionTok{seq\_along}\NormalTok{(opponent\_reversal),}
  \AttributeTok{memory\_state  =}\NormalTok{ d\_reversal}\SpecialCharTok{$}\NormalTok{memory\_state}
\NormalTok{) }\SpecialCharTok{\%\textgreater{}\%}
  \FunctionTok{pivot\_longer}\NormalTok{(}\FunctionTok{c}\NormalTok{(opponent\_rate, memory\_state),}
               \AttributeTok{names\_to =} \StringTok{"series"}\NormalTok{, }\AttributeTok{values\_to =} \StringTok{"value"}\NormalTok{) }\SpecialCharTok{\%\textgreater{}\%}
  \FunctionTok{mutate}\NormalTok{(}\AttributeTok{series =} \FunctionTok{recode}\NormalTok{(series,}
    \AttributeTok{opponent\_rate =} \StringTok{"Opponent cumulative rate"}\NormalTok{,}
    \AttributeTok{memory\_state  =} \StringTok{"Agent memory state"}
\NormalTok{  )) }\SpecialCharTok{\%\textgreater{}\%}
  \FunctionTok{ggplot}\NormalTok{(}\FunctionTok{aes}\NormalTok{(}\AttributeTok{x =}\NormalTok{ trial, }\AttributeTok{y =}\NormalTok{ value, }\AttributeTok{colour =}\NormalTok{ series)) }\SpecialCharTok{+}
  \FunctionTok{geom\_line}\NormalTok{(}\AttributeTok{linewidth =} \FloatTok{0.9}\NormalTok{) }\SpecialCharTok{+}
  \FunctionTok{geom\_vline}\NormalTok{(}\AttributeTok{xintercept =} \DecValTok{120}\NormalTok{, }\AttributeTok{linetype =} \StringTok{"dashed"}\NormalTok{, }\AttributeTok{colour =} \StringTok{"grey40"}\NormalTok{) }\SpecialCharTok{+}
  \FunctionTok{annotate}\NormalTok{(}\StringTok{"text"}\NormalTok{, }\AttributeTok{x =} \DecValTok{121}\NormalTok{, }\AttributeTok{y =} \FloatTok{0.95}\NormalTok{, }\AttributeTok{label =} \StringTok{"Reversal"}\NormalTok{, }\AttributeTok{hjust =} \DecValTok{0}\NormalTok{,}
           \AttributeTok{colour =} \StringTok{"grey40"}\NormalTok{, }\AttributeTok{size =} \FloatTok{3.5}\NormalTok{) }\SpecialCharTok{+}
  \FunctionTok{scale\_colour\_manual}\NormalTok{(}\AttributeTok{values =} \FunctionTok{c}\NormalTok{(}\StringTok{"steelblue"}\NormalTok{, }\StringTok{"darkred"}\NormalTok{)) }\SpecialCharTok{+}
  \FunctionTok{labs}\NormalTok{(}\AttributeTok{title  =} \StringTok{"Reversal Learning Design"}\NormalTok{,}
       \AttributeTok{subtitle =} \StringTok{"Opponent switches from 80\% right to 20\% right at trial 120"}\NormalTok{,}
       \AttributeTok{x =} \StringTok{"Trial"}\NormalTok{, }\AttributeTok{y =} \StringTok{"Rate / Memory"}\NormalTok{, }\AttributeTok{colour =} \ConstantTok{NULL}\NormalTok{) }\SpecialCharTok{+}
  \FunctionTok{theme\_classic}\NormalTok{() }\SpecialCharTok{+} \FunctionTok{theme}\NormalTok{(}\AttributeTok{legend.position =} \StringTok{"top"}\NormalTok{)}

\FunctionTok{print}\NormalTok{(p\_reversal\_design)}
\end{Highlighting}
\end{Shaded}

\begin{center}
\includegraphics[width=0.8\linewidth,height=\textheight,keepaspectratio]{05-ModelQualityChecks_files/figure-pdf/05-reversal-learning-fix-1.pdf}
\end{center}

\begin{Shaded}
\begin{Highlighting}[]
\CommentTok{\# {-}{-}{-} 4. Fit the memory model to the reversal data {-}{-}{-}}
\NormalTok{data\_reversal }\OtherTok{\textless{}{-}} \FunctionTok{list}\NormalTok{(}
  \AttributeTok{n            =}\NormalTok{ n\_reversal,}
  \AttributeTok{h            =}\NormalTok{ d\_reversal}\SpecialCharTok{$}\NormalTok{choice,}
  \AttributeTok{other        =}\NormalTok{ d\_reversal}\SpecialCharTok{$}\NormalTok{other\_choice,}
  \AttributeTok{prior\_bias\_m =} \DecValTok{0}\NormalTok{, }\AttributeTok{prior\_bias\_s =} \FloatTok{0.5}\NormalTok{,}
  \AttributeTok{prior\_beta\_m =} \DecValTok{0}\NormalTok{, }\AttributeTok{prior\_beta\_s =} \FloatTok{1.0}
\NormalTok{)}
\end{Highlighting}
\end{Shaded}

\begin{Shaded}
\begin{Highlighting}[]
\NormalTok{model\_file\_reversal }\OtherTok{\textless{}{-}} \FunctionTok{file.path}\NormalTok{(stan\_results\_dir, }\StringTok{"ch06\_Memory\_Reversal.rds"}\NormalTok{)}

\ControlFlowTok{if}\NormalTok{ (regenerate\_simulations }\SpecialCharTok{||} \SpecialCharTok{!}\FunctionTok{file.exists}\NormalTok{(model\_file\_reversal)) \{}
  \CommentTok{\# mod\_mem\_gq was compiled earlier in this chapter}
\NormalTok{  fit\_reversal }\OtherTok{\textless{}{-}}\NormalTok{ mod\_mem\_gq}\SpecialCharTok{$}\FunctionTok{sample}\NormalTok{(}
    \AttributeTok{data            =}\NormalTok{ data\_reversal,}
    \AttributeTok{chains          =} \DecValTok{4}\NormalTok{,}
    \AttributeTok{parallel\_chains =}\NormalTok{ parallel}\SpecialCharTok{::}\FunctionTok{detectCores}\NormalTok{() }\SpecialCharTok{{-}} \DecValTok{2}\NormalTok{,}
    \AttributeTok{iter\_warmup     =} \DecValTok{1000}\NormalTok{,}
    \AttributeTok{iter\_sampling   =} \DecValTok{2000}\NormalTok{,}
    \AttributeTok{refresh         =} \DecValTok{0}\NormalTok{,}
    \AttributeTok{adapt\_delta     =} \FloatTok{0.95}\NormalTok{,}
    \AttributeTok{max\_treedepth   =} \DecValTok{12}
\NormalTok{  )}
\NormalTok{  fit\_reversal}\SpecialCharTok{$}\FunctionTok{save\_object}\NormalTok{(}\AttributeTok{file =}\NormalTok{ model\_file\_reversal)}
\NormalTok{\} }\ControlFlowTok{else}\NormalTok{ \{}
\NormalTok{  fit\_reversal }\OtherTok{\textless{}{-}} \FunctionTok{readRDS}\NormalTok{(model\_file\_reversal)}
\NormalTok{\}}

\CommentTok{\# {-}{-}{-} 5. Joint posterior: confirm bias/beta correlation collapses {-}{-}{-}}
\NormalTok{draws\_reversal }\OtherTok{\textless{}{-}} \FunctionTok{as\_draws\_df}\NormalTok{(fit\_reversal}\SpecialCharTok{$}\FunctionTok{draws}\NormalTok{(}\FunctionTok{c}\NormalTok{(}\StringTok{"bias"}\NormalTok{, }\StringTok{"beta"}\NormalTok{)))}
\NormalTok{cor\_reversal   }\OtherTok{\textless{}{-}} \FunctionTok{cor}\NormalTok{(draws\_reversal}\SpecialCharTok{$}\NormalTok{bias, draws\_reversal}\SpecialCharTok{$}\NormalTok{beta)}

\NormalTok{p\_joint\_reversal }\OtherTok{\textless{}{-}} \FunctionTok{ggplot}\NormalTok{(draws\_reversal, }\FunctionTok{aes}\NormalTok{(}\AttributeTok{x =}\NormalTok{ bias, }\AttributeTok{y =}\NormalTok{ beta)) }\SpecialCharTok{+}
  \FunctionTok{stat\_density\_2d}\NormalTok{(}\FunctionTok{aes}\NormalTok{(}\AttributeTok{fill =} \FunctionTok{after\_stat}\NormalTok{(level)), }\AttributeTok{geom =} \StringTok{"polygon"}\NormalTok{,}
                  \AttributeTok{alpha =} \FloatTok{0.6}\NormalTok{) }\SpecialCharTok{+}
  \FunctionTok{geom\_point}\NormalTok{(}\AttributeTok{x =} \DecValTok{0}\NormalTok{, }\AttributeTok{y =} \FloatTok{0.9}\NormalTok{, }\AttributeTok{colour =} \StringTok{"black"}\NormalTok{, }\AttributeTok{shape =} \DecValTok{3}\NormalTok{,}
             \AttributeTok{size =} \DecValTok{5}\NormalTok{, }\AttributeTok{stroke =} \FloatTok{1.5}\NormalTok{) }\SpecialCharTok{+}
  \FunctionTok{scale\_fill\_viridis\_c}\NormalTok{(}\AttributeTok{option =} \StringTok{"mako"}\NormalTok{, }\AttributeTok{guide =} \StringTok{"none"}\NormalTok{) }\SpecialCharTok{+}
  \FunctionTok{coord\_cartesian}\NormalTok{(}\AttributeTok{xlim =} \FunctionTok{c}\NormalTok{(}\SpecialCharTok{{-}}\DecValTok{1}\NormalTok{, }\DecValTok{1}\NormalTok{), }\AttributeTok{ylim =} \FunctionTok{c}\NormalTok{(}\SpecialCharTok{{-}}\FloatTok{0.5}\NormalTok{, }\DecValTok{2}\NormalTok{)) }\SpecialCharTok{+}
  \FunctionTok{labs}\NormalTok{(}
    \AttributeTok{title    =} \StringTok{"Joint Posterior After Reversal Learning Design"}\NormalTok{,}
    \AttributeTok{subtitle =} \FunctionTok{paste0}\NormalTok{(}\StringTok{"Correlation r = "}\NormalTok{, }\FunctionTok{round}\NormalTok{(cor\_reversal, }\DecValTok{2}\NormalTok{),}
                      \StringTok{"   (compare to r ≈ 0.9 with constant{-}bias design)"}\NormalTok{),}
    \AttributeTok{x =} \StringTok{"Baseline Bias"}\NormalTok{,}
    \AttributeTok{y =} \StringTok{"Memory Strength (Beta)"}
\NormalTok{  ) }\SpecialCharTok{+}
  \FunctionTok{annotate}\NormalTok{(}\StringTok{"text"}\NormalTok{, }\AttributeTok{x =} \FloatTok{0.05}\NormalTok{, }\AttributeTok{y =} \FloatTok{0.95}\NormalTok{, }\AttributeTok{label =} \StringTok{"Ground Truth"}\NormalTok{,}
           \AttributeTok{fontface =} \StringTok{"bold"}\NormalTok{) }\SpecialCharTok{+}
  \FunctionTok{theme\_classic}\NormalTok{()}

\FunctionTok{print}\NormalTok{(p\_joint\_reversal)}
\end{Highlighting}
\end{Shaded}

\begin{center}
\includegraphics[width=0.8\linewidth,height=\textheight,keepaspectratio]{05-ModelQualityChecks_files/figure-pdf/ch06_fit_memory_reversal-1.pdf}
\end{center}

\begin{Shaded}
\begin{Highlighting}[]
\FunctionTok{cat}\NormalTok{(}\FunctionTok{sprintf}\NormalTok{(}
  \StringTok{"}\SpecialCharTok{\textbackslash{}n}\StringTok{Correlation with constant{-}bias design: \textasciitilde{}0.9}\SpecialCharTok{\textbackslash{}n}\StringTok{"}\NormalTok{,}
  \StringTok{"Correlation with reversal design:       \%.3f}\SpecialCharTok{\textbackslash{}n}\StringTok{"}\NormalTok{,}
\NormalTok{  cor\_reversal}
\NormalTok{))}
\end{Highlighting}
\end{Shaded}

\begin{verbatim}

Correlation with constant-bias design: ~0.9
\end{verbatim}

The joint posterior is now compact and approximately circular around the
true parameter values --- \texttt{bias} and \texttt{beta} are
identified. The reversal in the opponent's preference provides the
critical contrast that lets the model distinguish ``the agent is
generally biased right'' from ``the agent tracks the opponent and
follows them.''

\textbf{Saving the reversal dataset for Chapters 6 and 7.} Chapters 6
(multilevel models) and 7 (model comparison) both need a well-identified
memory dataset as their running example. We save the full reversal
simulation object here so those chapters can load it directly without
regenerating it.

\begin{Shaded}
\begin{Highlighting}[]
\NormalTok{reversal\_data\_file }\OtherTok{\textless{}{-}} \FunctionTok{file.path}\NormalTok{(sim\_data\_dir, }\StringTok{"ch07\_reversal\_data.rds"}\NormalTok{)}

\FunctionTok{saveRDS}\NormalTok{(}
  \FunctionTok{list}\NormalTok{(}
    \AttributeTok{stan\_data        =}\NormalTok{ data\_reversal,    }\CommentTok{\# Ready{-}to{-}use Stan data list}
    \AttributeTok{sim\_tibble       =}\NormalTok{ d\_reversal,       }\CommentTok{\# Full trial{-}level tibble}
    \AttributeTok{opponent\_choices =}\NormalTok{ opponent\_reversal,}\CommentTok{\# Raw opponent sequence}
    \AttributeTok{true\_bias        =} \DecValTok{0}\NormalTok{,}
    \AttributeTok{true\_beta        =} \FloatTok{0.9}\NormalTok{,}
    \AttributeTok{n\_trials         =}\NormalTok{ n\_reversal,}
    \AttributeTok{description      =} \FunctionTok{paste0}\NormalTok{(}
      \StringTok{"Reversal learning paradigm: opponent plays 80\% right (trials 1{-}120) "}\NormalTok{,}
      \StringTok{"then 20\% right (trials 121{-}240). Agent: bias=0, beta=0.9."}
\NormalTok{    )}
\NormalTok{  ),}
\NormalTok{  reversal\_data\_file}
\NormalTok{)}

\FunctionTok{cat}\NormalTok{(}\StringTok{"Reversal dataset saved to:"}\NormalTok{, reversal\_data\_file, }\StringTok{"}\SpecialCharTok{\textbackslash{}n}\StringTok{"}\NormalTok{)}
\end{Highlighting}
\end{Shaded}

\begin{verbatim}
Reversal dataset saved to: simdata/ch07_reversal_data.rds 
\end{verbatim}

\begin{Shaded}
\begin{Highlighting}[]
\FunctionTok{cat}\NormalTok{(}\StringTok{"Load in Ch. 7 and Ch. 8 with: readRDS(here(\textquotesingle{}simdata\textquotesingle{}, \textquotesingle{}ch07\_reversal\_data.rds\textquotesingle{}))}\SpecialCharTok{\textbackslash{}n}\StringTok{"}\NormalTok{)}
\end{Highlighting}
\end{Shaded}

\begin{verbatim}
Load in Ch. 7 and Ch. 8 with: readRDS(here('simdata', 'ch07_reversal_data.rds'))
\end{verbatim}

Finally, we run Simulation-Based Calibration Checks. We use the SBC
package to guarantee that our logit(memory) parameterization doesn't
introduce hidden biases in the sampler's rank statistics.

\begin{Shaded}
\begin{Highlighting}[]
\CommentTok{\# 1. Define the Generator using our simulate\_function}
\NormalTok{sbc\_memory\_generator\_fn }\OtherTok{\textless{}{-}} \ControlFlowTok{function}\NormalTok{() \{}
  \CommentTok{\# Draw true params from priors}
\NormalTok{  t\_bias }\OtherTok{\textless{}{-}} \FunctionTok{rnorm}\NormalTok{(}\DecValTok{1}\NormalTok{, }\DecValTok{0}\NormalTok{, }\FloatTok{0.5}\NormalTok{)}
\NormalTok{  t\_beta }\OtherTok{\textless{}{-}} \FunctionTok{rnorm}\NormalTok{(}\DecValTok{1}\NormalTok{, }\DecValTok{0}\NormalTok{, }\FloatTok{1.0}\NormalTok{)}
  
  \CommentTok{\# Use fixed opponent choices for simplicity in this check}
\NormalTok{  opp\_choices }\OtherTok{\textless{}{-}} \FunctionTok{rbinom}\NormalTok{(}\DecValTok{120}\NormalTok{, }\DecValTok{1}\NormalTok{, }\FloatTok{0.8}\NormalTok{)}
  
\NormalTok{  sim }\OtherTok{\textless{}{-}} \FunctionTok{simulate\_memory}\NormalTok{(}\AttributeTok{bias =}\NormalTok{ t\_bias, }\AttributeTok{beta =}\NormalTok{ t\_beta, }\AttributeTok{opponent\_choices =}\NormalTok{ opp\_choices)}
  
  \FunctionTok{list}\NormalTok{(}
    \AttributeTok{variables =} \FunctionTok{list}\NormalTok{(}\AttributeTok{bias =}\NormalTok{ t\_bias, }\AttributeTok{beta =}\NormalTok{ t\_beta),}
    \AttributeTok{generated =} \FunctionTok{list}\NormalTok{(}\AttributeTok{N\_trials =} \DecValTok{120}\NormalTok{, }\AttributeTok{y =}\NormalTok{ sim}\SpecialCharTok{$}\NormalTok{choice, }\AttributeTok{other =}\NormalTok{ opp\_choices,}
                     \AttributeTok{prior\_bias\_mu =} \DecValTok{0}\NormalTok{, }\AttributeTok{prior\_bias\_sigma =} \FloatTok{0.5}\NormalTok{,}
                     \AttributeTok{prior\_beta\_mu =} \DecValTok{0}\NormalTok{, }\AttributeTok{prior\_beta\_sigma =} \FloatTok{1.0}\NormalTok{,}
                     \AttributeTok{run\_diagnostics =} \DecValTok{0}\NormalTok{L)}
\NormalTok{  )}
\NormalTok{\}}

\CommentTok{\# 2. Execute SBC (1000 simulations)}
\NormalTok{sbc\_mem\_obj }\OtherTok{\textless{}{-}}\NormalTok{ SBC}\SpecialCharTok{::}\FunctionTok{SBC\_generator\_function}\NormalTok{(sbc\_memory\_generator\_fn)}
\NormalTok{sbc\_mem\_backend }\OtherTok{\textless{}{-}}\NormalTok{ SBC}\SpecialCharTok{::}\FunctionTok{SBC\_backend\_cmdstan\_sample}\NormalTok{(mod\_mem\_gq, }\AttributeTok{chains =} \DecValTok{1}\NormalTok{, }\AttributeTok{refresh =} \DecValTok{0}\NormalTok{)}

\CommentTok{\# Run parallel fits}
\NormalTok{future}\SpecialCharTok{::}\FunctionTok{plan}\NormalTok{(future}\SpecialCharTok{::}\NormalTok{sequential)}
\NormalTok{sbc\_mem\_results }\OtherTok{\textless{}{-}}\NormalTok{ SBC}\SpecialCharTok{::}\FunctionTok{compute\_SBC}\NormalTok{(SBC}\SpecialCharTok{::}\FunctionTok{generate\_datasets}\NormalTok{(sbc\_mem\_obj, }\AttributeTok{n\_sims =} \DecValTok{200}\NormalTok{), sbc\_mem\_backend)}
\NormalTok{future}\SpecialCharTok{::}\FunctionTok{plan}\NormalTok{(future}\SpecialCharTok{::}\NormalTok{sequential)}

\CommentTok{\# 3. Visualize calibration}
\NormalTok{p\_mem\_sbc }\OtherTok{\textless{}{-}}\NormalTok{ SBC}\SpecialCharTok{::}\FunctionTok{plot\_rank\_hist}\NormalTok{(sbc\_mem\_results) }\SpecialCharTok{+} \FunctionTok{theme\_classic}\NormalTok{()}
\FunctionTok{print}\NormalTok{(p\_mem\_sbc)}
\end{Highlighting}
\end{Shaded}

\begin{center}
\includegraphics[width=0.8\linewidth,height=\textheight,keepaspectratio]{05-ModelQualityChecks_files/figure-pdf/ch06_memory_sbc-1.pdf}
\end{center}

\section{Conclusion: The Validation Workflow as a Reusable
Checklist}\label{conclusion-the-validation-workflow-as-a-reusable-checklist}

Fitting a model and obtaining parameter estimates is just the beginning.
This chapter introduced a six-phase validation battery that every model
must pass before its estimates can be trusted:

\begin{enumerate}
\def\labelenumi{\arabic{enumi}.}
\tightlist
\item
  \textbf{Mathematical formulation} --- confirm the Stan code does what
  you think.
\item
  \textbf{Prior predictive check} --- priors must generate plausible
  data.
\item
  \textbf{Posterior predictive check} --- fitted model must reproduce
  observed patterns.
\item
  \textbf{Prior-posterior update} --- data must genuinely update
  beliefs.
\item
  \textbf{Sensitivity analysis} --- conclusions must survive small prior
  perturbations.
\item
  \textbf{Parameter recovery + SBC} --- the engine must find the truth
  and intervals must be calibrated.
\end{enumerate}

Failing any check is not a setback --- it is information. The memory
model's \texttt{bias}/\texttt{beta} correlation failure in Phase 6 told
us exactly what was wrong with the experimental design and how to fix
it. That is the validation workflow working correctly: it caught a real
problem before the model touched empirical data.

The reversal learning dataset generated in this chapter
(\texttt{simdata/ch07\_reversal\_data.rds}) is the running example for
the next two chapters. It is well-identified, validated against all six
checks, and ready to be extended.

\textbf{Coming up in Chapter 7:} We take the memory model --- now
validated on the reversal paradigm --- and scale it to many simulated
participants using a hierarchical (multilevel) model. The central
question shifts from ``can we estimate parameters for one agent?'' to
``how do individuals vary, and how do we share information across
them?''

\bookmarksetup{startatroot}

\chapter{Individual Differences in Cognitive Strategies: The
Hierarchical Architecture}\label{sec-multilevel-models}

\begin{quote}
\textbf{📍 Where we are in the Bayesian modeling workflow:} Chs. 1--5
fitted models to a single agent and validated the inference engine. Now
we scale up: real experiments involve many participants who differ from
each other. This chapter builds hierarchical (multilevel) models that
estimate both individual parameters \emph{and} the population
distribution they are drawn from --- and validates the whole pipeline
with the same six-phase battery from Ch. 6. The memory model fitted here
uses the \textbf{reversal learning dataset from Ch. 6}, which was the
first dataset that cleanly separated \texttt{bias} from \texttt{beta}.
The fitted model and its LOO score are saved for Ch. 8.
\end{quote}

\section{Introduction}\label{introduction}

Our exploration of decision-making models has so far focused on single
agents or averaged behavior. However, individuals differ systematically
in how they approach tasks and process information. Some people may be
more risk-averse, have better memory, learn faster, or employ entirely
different strategies than others. Ignoring these individual differences
can lead to misleading conclusions about cognitive processes. From a
statistical perspective, if we assume a single set of parameters
generates all data, our model is mathematically misspecified, which
leads to pathological posterior geometries and misleading inferential
results.

This chapter introduces hierarchical modeling (also called multilevel or
mixed effects modeling) as a powerful framework for capturing these
individual differences while still identifying population-level
patterns.

\subsection{The Epistemic Challenge: Pooling
Information}\label{the-epistemic-challenge-pooling-information}

Traditional statistical approaches to handling data from multiple
individuals force a false dichotomy, imposing rigid, unjustified
structural assumptions on the data:

\begin{enumerate}
\def\labelenumi{\arabic{enumi}.}
\tightlist
\item
  \textbf{Complete Pooling (Underfitting)}: Treats all participants as
  identical: it estimates a single set of parameters for the entire
  group, completely ignoring potentially meaningful individual
  variation. This assumes the population variance (the average error in
  predicting an individual from the population) is strictly zero
  (\(\sigma \to 0\)).

  \begin{itemize}
  \tightlist
  \item
    \emph{The Pathology}: Systematically fails to capture cognitive
    heterogeneity, describing an ``average'' participant that does not
    necessarily exist.
  \end{itemize}
\item
  \textbf{No Pooling (Overfitting)}: Analyzes each participant in
  complete isolation, estimating independent parameters (\(\theta_j\))
  for each individual. No pooling corresponds to fitting \(J\) separate
  models to \(J\) individual participants. This assumes the population
  variance is infinite (\(\sigma \to \infty\)).

  \begin{itemize}
  \tightlist
  \item
    \emph{The Pathology}: Severely inefficient and vulnerable to noise.
    For individuals with sparse or noisy data, the posterior estimate
    will be excessively wide or completely dominated by the random
    fluctuations of their specific likelihood.
  \end{itemize}
\end{enumerate}

\begin{quote}
\emph{Why does this dichotomy matter in practice?} With 120 trials per
agent, complete pooling produces one parameter estimate that represents
nobody in particular. No pooling produces 50 estimates, many of which
are dominated by noise for agents with only 20--30 trials. Partial
pooling is not a philosophical compromise --- it produces estimates that
are demonstrably \emph{closer to the truth} for low-data individuals, as
the shrinkage plot at the end of this chapter will show directly.
\end{quote}

\subsection{The Solution: Partial Pooling (Multilevel or Hierarchical
Modeling)}\label{the-solution-partial-pooling-multilevel-or-hierarchical-modeling}

Hierarchical modeling offers a principled, data-driven middle ground.
Instead of forcing \(\sigma\) to be \(0\) or \(\infty\), we treat the
population variance as a free parameter estimated from the data. The
model is structured hierarchically:

\begin{itemize}
\item
  \textbf{Individual Level}: Each individual agent \(j\) has a specific
  parameter vector (e.g., their intrinsic bias \(\theta_j\) or memory
  decay \(\beta_j\)).
\item
  \textbf{Population Level}: These individual parameters are explicitly
  modeled as draws from an overarching continuous population
  distribution (e.g., a Gaussian characterized by a population mean
  \(\mu\) and standard deviation \(\sigma\)).
\item
  \textbf{The Bayesian Compromise}: The estimation of an individual's
  parameter (\(\theta_j\)) is a precision-weighted combination of their
  own empirical data (the likelihood) and the population distribution
  (the hierarchical prior).
\item
  \textbf{Shrinkage}: Estimates for individuals with sparse or noisy
  data are ``shrunk'' toward the population mean \(\mu\). This is a
  mathematical mechanism that regularizes inference, preventing
  overfitting and yielding more robust out-of-sample estimates.
\end{itemize}

\begin{tcolorbox}[enhanced jigsaw, toprule=.15mm, toptitle=1mm, bottomtitle=1mm, leftrule=.75mm, rightrule=.15mm, colframe=quarto-callout-note-color-frame, left=2mm, colback=white, titlerule=0mm, coltitle=black, colbacktitle=quarto-callout-note-color!10!white, opacityback=0, title=\textcolor{quarto-callout-note-color}{\faInfo}\hspace{0.5em}{Historical Context: Shrinkage, Partial Pooling, and the Recovery of
Individual Differences}, arc=.35mm, opacitybacktitle=0.6, breakable, bottomrule=.15mm]

The hierarchical model is now a standard tool, but the intellectual path
to it was surprisingly contentious. Accepting partial pooling required
overturning the intuition --- deeply held by statisticians trained in
classical estimation theory --- that combining information from
different units is always a form of cheating.

\subsubsection{Mixed Models in the
Fields}\label{mixed-models-in-the-fields}

The mathematical foundations were laid not in cognitive science but in
animal breeding. Henderson (1950) needed to estimate the genetic merit
of individual bulls from data on their offspring, spread across farms
that differed in management and feeding. The complete-pooling solution
(assume all farms are identical) was agronomically wrong; the no-pooling
solution (analyse each farm in isolation) wasted the information that
herds from the same breed share a common distribution of genetic merit.
Henderson's \textbf{Best Linear Unbiased Prediction} (BLUP) formalised
the middle path: individual estimates are shrunk toward the population
mean by an amount that depends on how much data each individual
contributes. The same equations appear, under different names, as the
core of the hierarchical models in this chapter.

\subsubsection{Stein's Paradox and the Optimality of
Shrinkage}\label{steins-paradox-and-the-optimality-of-shrinkage}

The theoretical justification arrived from a different direction, and it
was deeply counterintuitive. James and Stein (1961) proved that the
ordinary (no-pooling) estimator of three or more means is
\emph{inadmissible} under squared-error loss: there always exists a
shrinkage estimator --- one that pulls all estimates toward a common
value --- with uniformly lower expected loss. This is Stein's paradox.
It implies that, regardless of whether the true means are related to
each other, you are statistically better off treating them as
exchangeable draws from a common distribution than estimating each
independently. The result was controversial because it seemed to say
that knowing the speed of light should improve your estimate of the
boiling point of water --- which is absurd --- but the paradox dissolves
once you recognise that the admissibility comparison is over the
\emph{joint} vector of estimates, not individual components.

\subsubsection{Bayesian Formalisation}\label{bayesian-formalisation}

Lindley and Smith (1972) showed that the James-Stein estimator is a
Bayesian posterior mean under the assumption that the true parameters
are drawn from a common normal distribution --- exactly the hierarchical
prior structure used in this chapter. This connected the frequentist
inadmissibility result to a principled Bayesian generative story:
shrinkage is not a trick but the natural consequence of treating
individuals as exchangeable members of a population. Gelman and Rubin
(1992) and Gelman et al. (1995) developed the computational
infrastructure (via MCMC) that made this formalism applicable to
realistic psychological datasets.

\subsubsection{Into Cognitive Science}\label{into-cognitive-science}

The adoption of hierarchical models in cognitive science was slow. For
much of the 1990s and early 2000s, individual-differences research
either fitted separate models to each participant and then performed a
second-stage group analysis --- the ``two-stage'' procedure that
discards uncertainty --- or fitted a single population model that
ignored individual variation entirely. Rouder et al. (2005) and Rouder
and Lu (2008) demonstrated that hierarchical models substantially
improve parameter recovery in reaction-time and signal-detection
contexts: individual estimates derived from partial pooling are closer
to the true generating values than either the complete-pooling or
no-pooling alternatives, especially for participants who contribute few
trials. The shrinkage plot at the end of this chapter makes this
advantage visible directly: you can watch the model borrow strength from
the population to regularise noisy individuals.

\subsubsection{Non-Centred Parameterisation and Neal's
Funnel}\label{non-centred-parameterisation-and-neals-funnel}

The one purely computational contribution to this story came from Neal
(2003) and was adopted as standard practice by the Stan community via
Betancourt and Girolami (2013). When individual parameters are
parameterised as \(\theta_j = \mu + \sigma \cdot z_j\) with
\(z_j \sim \mathcal{N}(0,1)\) --- the non-centred form --- the geometry
of the posterior becomes approximately Gaussian and Hamiltonian Monte
Carlo can traverse it efficiently. In the centred form
\(\theta_j \sim \mathcal{N}(\mu, \sigma^2)\), the posterior develops a
funnel geometry near \(\sigma \approx 0\) that causes HMC to diverge.
The non-centred parameterisation is the single most important practical
technique in hierarchical Bayesian computation, and the reason the
models in this chapter spend time on it.

\end{tcolorbox}

\section{Learning Objectives}\label{learning-objectives-5}

After completing this chapter, you will be able to:

\begin{itemize}
\item
  \textbf{Explain Partial Pooling}: Understand how multilevel modeling
  balances individual and group-level information through shrinkage.
\item
  \textbf{Distinguish Pooling Approaches}: Articulate the differences
  and trade-offs between complete pooling (\(\sigma=0\)), no pooling
  (\(\sigma=\infty\)), and partial pooling.
\item
  \textbf{Implement Hierarchical Models}: Write Stan code for
  hierarchical models with varying intercepts and/or slopes.
\item
  \textbf{Reparameterize a model}: Understand why and how to
  reparameterize a model and in particular how to use non-centered
  parameterization to improve sampling efficiency in hierarchical
  models.
\item
  \textbf{Model Parameter Correlations}: Implement and interpret models
  that estimate correlations between individual-level parameters (via
  Cholesky-factorized correlation matrices \(\mathbf{L}_\Omega\)).
\item
  \textbf{Perform Comprehensive Model Checks}: Apply a rigorous Bayesian
  workflow to hierarchical models (Predictive Checks, Prior-Posterior
  Update Checks, Sensitivity Analysis, Parameter Recovery,
  Simulation-Based Calibration Checks, etc).
\item
  \textbf{Apply to Cognitive Questions}: Use multilevel modeling to
  investigate individual differences in cognitive strategies (e.g.,
  bias, memory, WSLS).
\end{itemize}

This chapter has two parts. The first applies the full Bayesian workflow
to the \textbf{Hierarchical Biased Agent} --- one parameter per person,
one population distribution --- and uses it to expose and fix the
Centered Parameterization's pathological geometry (Neal's Funnel). The
second applies the same workflow to the \textbf{Hierarchical Memory
Agent}, which has two correlated parameters per person, requires a
bivariate population distribution, and uses the reversal learning
dataset from Ch. 6 that we validated there. The LOO score from the
fitted memory model is saved for Ch. 8's model comparison.

\subsection{Biased Agent Model
(Multilevel)}\label{biased-agent-model-multilevel}

In this architecture, each agent \(j \in \{1, \dots, J\}\) possesses an
unobserved, individual-level bias parameter \(\theta_j\). This
individual parameter is generated by a latent population distribution,
parameterized by a global mean \(\mu_\theta\) and global standard
deviation \(\sigma_\theta\). Finally, the empirical choices \(h_{ij}\)
for trial \(i\) are generated conditionally on the individual parameter
\(\theta_j\).

Mathematically, this translates to the following generative structure on
the unconstrained (logit)
space:\[h_{ij} \sim \text{Bernoulli}(\text{logit}^{-1}(\theta_{\text{logit}, j}))\]\[\theta_{\text{logit}, j} \sim \text{Normal}(\mu_{\theta}, \sigma_{\theta})\]\[\mu_{\theta} \sim \text{Normal}(0, 1.5)\]\[\sigma_{\theta} \sim \text{Exponential}(1)\]

By mapping this DAG, we have strictly defined our generative
assumptions. We are now ready to translate this architecture into Stan
code and subject it to Prior Predictive Checks.

\begin{figure}[H]

{\centering \includegraphics[width=0.8\linewidth,height=\textheight,keepaspectratio]{06-MultilevelModels_files/figure-pdf/hierarchical_biased_dags-1.png}

}

\caption{Standard Plate Notation (Left) and Kruschke-Style Generative
Notation (Right) for the Hierarchical Biased Agent.}

\end{figure}%

\begin{quote}
\textbf{Computational Note: The Geometry of Hierarchical Models} The
mathematical definition provided above---specifically
\(\theta_{\text{logit}, j} \sim \text{Normal}(\mu_{\theta}, \sigma_{\theta})\)---is
known as the \textbf{Centered Parameterization}. It is the most direct
conceptual translation of our theory. However, Hamiltonian Monte Carlo
(NUTS) often struggles with this exact form. As \(\sigma_\theta\)
approaches zero, the target geometry constrains into a region of extreme
curvature known as ``Neal's Funnel.'' In the next phase, we will
purposefully implement this naive form to diagnose its pathological
failures before introducing the coordinate-transform solution
(Non-Centered Parameterization).
\end{quote}

\section{Phase 1: Generative Plausibility of the Hierarchical Biased
Agent}\label{phase-1-generative-plausibility-of-the-hierarchical-biased-agent}

Before we write the inference algorithm in Stan, we must construct the
forward generative model in R. This code must strictly mirror the
mathematical dependencies formalized in our DAG.

We model the parameters on the logit (log-odds) scale. As discussed in
previous chapters, probabilities are bounded between \(0\) and \(1\). If
we attempt to place a hierarchical Normal distribution directly on the
probability space, the probability mass will inevitably spill over the
bounds, creating mathematically impossible values (e.g., a probability
of \(1.2\)). The logit transform
\(\text{logit}(p) = \log(\frac{p}{1-p})\) maps the bounded \([0, 1]\)
interval to the unbounded real line \((-\infty, \infty)\). The HMC
sampler traverses this unbounded space smoothly, and we use the
inverse-logit function (plogis() in R, inv\_logit() in Stan) to map the
inferences back to probabilities at the likelihood level.

Let us build a unified generative function. In accordance with the
rigorous Bayesian workflow, this function is designed to either:

\begin{itemize}
\item
  Draw hyperparameters directly from our stated priors (for Prior
  Predictive Checks).
\item
  Accept fixed ``ground-truth'' hyperparameters (for Parameter Recovery
  in Phase 2).
\end{itemize}

\begin{Shaded}
\begin{Highlighting}[]
\CommentTok{\#\textquotesingle{} Simulate Hierarchical Biased Agents}
\CommentTok{\#\textquotesingle{}}
\CommentTok{\#\textquotesingle{} @param n\_agents Integer, number of agents (J).}
\CommentTok{\#\textquotesingle{} @param n\_trials\_min Integer, minimum number of trials per agent.}
\CommentTok{\#\textquotesingle{} @param n\_trials\_max Integer, maximum number of trials per agent.}
\CommentTok{\#\textquotesingle{} @param mu\_theta Numeric, population mean on the logit scale. If NULL, draws from prior.}
\CommentTok{\#\textquotesingle{} @param sigma\_theta Numeric, population SD on the logit scale. If NULL, draws from prior.}
\CommentTok{\#\textquotesingle{} @param seed Integer, for exact computational reproducibility.}
\CommentTok{\#\textquotesingle{}}
\CommentTok{\#\textquotesingle{} @return A tidy tibble containing both latent parameters and observed choices.}
\NormalTok{simulate\_hierarchical\_biased }\OtherTok{\textless{}{-}} \ControlFlowTok{function}\NormalTok{(}\AttributeTok{n\_agents =} \DecValTok{100}\NormalTok{, }
                                         \AttributeTok{n\_trials\_min =} \DecValTok{20}\NormalTok{, }
                                         \AttributeTok{n\_trials\_max =} \DecValTok{200}\NormalTok{, }
                                         \AttributeTok{mu\_theta =} \ConstantTok{NULL}\NormalTok{, }
                                         \AttributeTok{sigma\_theta =} \ConstantTok{NULL}\NormalTok{,}
                                         \AttributeTok{seed =} \ConstantTok{NULL}\NormalTok{) \{}
  
  \ControlFlowTok{if}\NormalTok{ (}\SpecialCharTok{!}\FunctionTok{is.null}\NormalTok{(seed)) }\FunctionTok{set.seed}\NormalTok{(seed)}
  
  \CommentTok{\# 1. Sample Population Hyperparameters (if not provided)}
  \ControlFlowTok{if}\NormalTok{ (}\FunctionTok{is.null}\NormalTok{(mu\_theta)) mu\_theta }\OtherTok{\textless{}{-}} \FunctionTok{rnorm}\NormalTok{(}\DecValTok{1}\NormalTok{, }\AttributeTok{mean =} \DecValTok{0}\NormalTok{, }\AttributeTok{sd =} \FloatTok{1.5}\NormalTok{)}
  \ControlFlowTok{if}\NormalTok{ (}\FunctionTok{is.null}\NormalTok{(sigma\_theta)) sigma\_theta }\OtherTok{\textless{}{-}} \FunctionTok{rexp}\NormalTok{(}\DecValTok{1}\NormalTok{, }\AttributeTok{rate =} \DecValTok{1}\NormalTok{)}
  
  \CommentTok{\# 2. Sample Individual Latent Parameters and Trial Counts}
  \CommentTok{\# We purposefully vary N\_j to demonstrate shrinkage later.}
\NormalTok{  agent\_params }\OtherTok{\textless{}{-}} \FunctionTok{tibble}\NormalTok{(}
    \AttributeTok{agent\_id =} \DecValTok{1}\SpecialCharTok{:}\NormalTok{n\_agents,}
    \AttributeTok{n\_trials =} \FunctionTok{sample}\NormalTok{(n\_trials\_min}\SpecialCharTok{:}\NormalTok{n\_trials\_max, n\_agents, }\AttributeTok{replace =} \ConstantTok{TRUE}\NormalTok{),}
    \AttributeTok{theta\_logit =} \FunctionTok{rnorm}\NormalTok{(n\_agents, }\AttributeTok{mean =}\NormalTok{ mu\_theta, }\AttributeTok{sd =}\NormalTok{ sigma\_theta),}
    \AttributeTok{theta\_prob =} \FunctionTok{plogis}\NormalTok{(theta\_logit) }
\NormalTok{  )}
  
  \CommentTok{\# 3. Generate Observed Data (Likelihood)}
  \CommentTok{\# h\_ij \textasciitilde{} Bernoulli(inv\_logit(theta\_logit\_j))}
\NormalTok{  sim\_data }\OtherTok{\textless{}{-}}\NormalTok{ agent\_params }\SpecialCharTok{|\textgreater{}} 
    \FunctionTok{mutate}\NormalTok{(}
      \AttributeTok{trial\_data =} \FunctionTok{map2}\NormalTok{(theta\_prob, n\_trials, }\ControlFlowTok{function}\NormalTok{(prob, n) \{}
        \FunctionTok{tibble}\NormalTok{(}
          \AttributeTok{trial =} \DecValTok{1}\SpecialCharTok{:}\NormalTok{n,}
          \AttributeTok{choice =} \FunctionTok{rbinom}\NormalTok{(n, }\AttributeTok{size =} \DecValTok{1}\NormalTok{, }\AttributeTok{prob =}\NormalTok{ prob)}
\NormalTok{        )}
\NormalTok{      \})}
\NormalTok{    ) }\SpecialCharTok{|\textgreater{}} 
    \FunctionTok{unnest}\NormalTok{(trial\_data)}
  
  \FunctionTok{attr}\NormalTok{(sim\_data, }\StringTok{"hyperparameters"}\NormalTok{) }\OtherTok{\textless{}{-}} \FunctionTok{list}\NormalTok{(}
    \AttributeTok{mu\_theta =}\NormalTok{ mu\_theta,}
    \AttributeTok{sigma\_theta =}\NormalTok{ sigma\_theta}
\NormalTok{  )}
  
  \FunctionTok{return}\NormalTok{(sim\_data)}
\NormalTok{\}}
\end{Highlighting}
\end{Shaded}

\subsection{Executing the Prior Predictive Check: The Logit-Normal
Artifact}\label{executing-the-prior-predictive-check-the-logit-normal-artifact}

When placing priors on an unbounded scale (logit) that will be
transformed to a bounded scale (probability), our intuition often fails.
A seemingly ``uninformative'' prior like
\(\mu \sim \text{Normal}(0, 10)\) on the logit scale actually forces
almost all probability mass to exactly \(0\) or \(1\) on the outcome
scale, implying that all agents are entirely deterministic.

Let us generate 12 datasets drawn entirely from our chosen priors
(\(\mu_\theta \sim N(0, 1.5), \sigma_\theta \sim \text{Exp}(1)\)) to
verify they generate scientifically plausible, heterogeneous
populations.

\begin{Shaded}
\begin{Highlighting}[]
\CommentTok{\# Generate 12 distinct "universes" from the priors}
\NormalTok{prior\_sims }\OtherTok{\textless{}{-}} \FunctionTok{map\_dfr}\NormalTok{(}\DecValTok{1}\SpecialCharTok{:}\DecValTok{12}\NormalTok{, }\ControlFlowTok{function}\NormalTok{(sim\_id) \{}
  \CommentTok{\# We use a fixed seed per universe for stable workflow mapping}
\NormalTok{  sim\_data }\OtherTok{\textless{}{-}} \FunctionTok{simulate\_hierarchical\_biased}\NormalTok{(}
    \AttributeTok{n\_agents =} \DecValTok{100}\NormalTok{, }
    \AttributeTok{n\_trials\_min =} \DecValTok{1}\NormalTok{, }\AttributeTok{n\_trials\_max =} \DecValTok{1}\NormalTok{, }\CommentTok{\# Only need 1 trial to extract parameters}
    \AttributeTok{seed =}\NormalTok{ sim\_id }\SpecialCharTok{*} \DecValTok{10} 
\NormalTok{  )}
  
\NormalTok{  sim\_data }\SpecialCharTok{|\textgreater{}} 
    \FunctionTok{distinct}\NormalTok{(agent\_id, theta\_prob) }\SpecialCharTok{|\textgreater{}} 
    \FunctionTok{mutate}\NormalTok{(}\AttributeTok{simulation\_id =} \FunctionTok{paste}\NormalTok{(}\StringTok{"Prior Universe"}\NormalTok{, sim\_id))}
\NormalTok{\})}

\CommentTok{\# Visualize the generative plausibility}
\FunctionTok{ggplot}\NormalTok{(prior\_sims, }\FunctionTok{aes}\NormalTok{(}\AttributeTok{x =}\NormalTok{ theta\_prob)) }\SpecialCharTok{+}
  \FunctionTok{geom\_density}\NormalTok{(}\AttributeTok{fill =} \StringTok{"steelblue"}\NormalTok{, }\AttributeTok{alpha =} \FloatTok{0.5}\NormalTok{, }\AttributeTok{color =} \ConstantTok{NA}\NormalTok{) }\SpecialCharTok{+}
  \FunctionTok{facet\_wrap}\NormalTok{(}\SpecialCharTok{\textasciitilde{}}\NormalTok{simulation\_id, }\AttributeTok{ncol =} \DecValTok{4}\NormalTok{) }\SpecialCharTok{+}
  \FunctionTok{scale\_x\_continuous}\NormalTok{(}\AttributeTok{limits =} \FunctionTok{c}\NormalTok{(}\DecValTok{0}\NormalTok{, }\DecValTok{1}\NormalTok{), }\AttributeTok{breaks =} \FunctionTok{c}\NormalTok{(}\DecValTok{0}\NormalTok{, }\FloatTok{0.5}\NormalTok{, }\DecValTok{1}\NormalTok{)) }\SpecialCharTok{+}
  \FunctionTok{labs}\NormalTok{(}
    \AttributeTok{title =} \StringTok{"Prior Predictive Check: Implied Population Distributions"}\NormalTok{,}
    \AttributeTok{subtitle =} \StringTok{"Healthy = populations spread across (0,1), not piled at 0 or 1. Each panel is one prior universe."}\NormalTok{,}
    \AttributeTok{x =} \StringTok{"Agent Bias (Probability Scale)"}\NormalTok{,}
    \AttributeTok{y =} \StringTok{"Density"}
\NormalTok{  )}
\end{Highlighting}
\end{Shaded}

\begin{center}
\includegraphics[width=0.8\linewidth,height=\textheight,keepaspectratio]{06-MultilevelModels_files/figure-pdf/ch07_prior_predictive_check_viz-1.pdf}
\end{center}

\section{Phase 2: Data Preparation and the Centered Stan
Model}\label{phase-2-data-preparation-and-the-centered-stan-model}

Before we unleash the Hamiltonian Monte Carlo (HMC) engine, we must
format our simulated data structure to match the static typing of the
Stan compiler. Because our agents have varying numbers of trials
(\(N_j\)), we cannot pass a simple \(J \times N\) matrix. Instead, we
use a ``long format'' array indexing strategy, where a specific vector
maps every trial \(n \in \{1, \dots, N_{\text{total}}\}\) to its
generating agent \(j\).

\begin{Shaded}
\begin{Highlighting}[]
\CommentTok{\#\textquotesingle{} Prepare Data for Stan}
\CommentTok{\#\textquotesingle{}}
\CommentTok{\#\textquotesingle{} Maps a tidy tibble into the strictly typed list required by CmdStanR.}
\CommentTok{\#\textquotesingle{} }
\CommentTok{\#\textquotesingle{} @param sim\_data Tidy tibble generated by simulate\_hierarchical\_biased}
\CommentTok{\#\textquotesingle{} @return A named list formatted for CmdStanR}
\NormalTok{prep\_stan\_data }\OtherTok{\textless{}{-}} \ControlFlowTok{function}\NormalTok{(sim\_data) \{}
  \FunctionTok{list}\NormalTok{(}
    \AttributeTok{N\_total                  =} \FunctionTok{nrow}\NormalTok{(sim\_data),}
    \AttributeTok{N\_subjects               =} \FunctionTok{length}\NormalTok{(}\FunctionTok{unique}\NormalTok{(sim\_data}\SpecialCharTok{$}\NormalTok{agent\_id)),}
    \AttributeTok{subj\_id                  =}\NormalTok{ sim\_data}\SpecialCharTok{$}\NormalTok{agent\_id,}
    \AttributeTok{y                        =}\NormalTok{ sim\_data}\SpecialCharTok{$}\NormalTok{choice,}
    \AttributeTok{prior\_mu\_theta\_mu        =} \DecValTok{0}\NormalTok{,}
    \AttributeTok{prior\_mu\_theta\_sigma     =} \FloatTok{1.5}\NormalTok{,}
    \AttributeTok{prior\_sigma\_theta\_lambda =} \DecValTok{1}\NormalTok{,}
    \AttributeTok{run\_diagnostics          =} \DecValTok{1}
\NormalTok{  )}
\NormalTok{\}}

\CommentTok{\# Execute the preparation on a specifically generated universe for recovery}
\NormalTok{sim\_data\_centered }\OtherTok{\textless{}{-}} \FunctionTok{simulate\_hierarchical\_biased}\NormalTok{(}
  \AttributeTok{n\_agents =} \DecValTok{50}\NormalTok{, }
  \AttributeTok{n\_trials\_min =} \DecValTok{30}\NormalTok{,  }\CommentTok{\# Intentional sparsity to force shrinkage}
  \AttributeTok{n\_trials\_max =} \DecValTok{120}\NormalTok{,}
  \AttributeTok{mu\_theta =} \FloatTok{1.38}\NormalTok{,}\CommentTok{\# slightly below .8 in prob}
  \AttributeTok{sigma\_theta =} \FloatTok{0.8}\NormalTok{,  }
  \AttributeTok{seed =} \DecValTok{1981}
\NormalTok{)}
\NormalTok{stan\_data\_centered }\OtherTok{\textless{}{-}} \FunctionTok{prep\_stan\_data}\NormalTok{(sim\_data\_centered)}
\end{Highlighting}
\end{Shaded}

\subsection{Stan Model: Centered Parameterization
(CP)}\label{stan-model-centered-parameterization-cp}

We will now implement the Centered Parameterization (CP) of the
hierarchical biased agent. In CP, we explicitly model the individual
latent parameters (\(\theta_{\text{logit}, j}\)) - determining their
probability of choosing ``right'' versus ``left'' - as being drawn
directly from a population distribution with mean \texttt{thetaM} and
standard deviation \texttt{thetaSD}.

Mathematically, this translates to the following hierarchical structure,
which strictly mirrors our Data Generating Process (DGP) from Phase 1:

\begin{itemize}
\item
  \textbf{Population Level}: \[\mu_\theta \sim \text{Normal}(0, 1.5)\]
  \[\sigma_\theta \sim \text{Exponential}(1)\]
\item
  \textbf{Individual Level}:
  \[\theta_{\text{logit}, j} \sim \text{Normal}(\mu_\theta, \sigma_\theta)\]
\item
  \textbf{Data Level (Likelihood)}:
  \[h_{ij} \sim \text{Bernoulli}(\text{logit}^{-1}(\theta_{\text{logit}, j}))\]
\end{itemize}

This architecture mathematically instantiates partial pooling: the
estimation of an individual's bias balances their specific choice
history (\(h_{ij}\)) with the population's overarching distributional
geometry (\(\mu_\theta, \sigma_\theta\)).

\begin{quote}
The Anatomy of a Potential Pathology: Look closely at the
individual-level formulation:
\(\theta_{\text{logit}, j} \sim \text{Normal}(\mu_\theta, \sigma_\theta)\).
The permissible geometric space for \(\theta_{\text{logit}, j}\) is
dynamically dictated by \(\sigma_\theta\). If the data suggests the
population is highly homogenous, the sampler will push
\(\sigma_\theta \to 0\). As it does, the parameter space for
\(\theta_{\text{logit}, j}\) is violently crushed toward \(\mu_\theta\),
creating a funnel of infinite curvature. We are fitting this model
precisely to watch the NUTS sampler fail in this region.
\end{quote}

Let us write the Stan code.

\begin{Shaded}
\begin{Highlighting}[]
\NormalTok{stan\_biased\_ml\_cp\_code }\OtherTok{\textless{}{-}} \StringTok{"}
\StringTok{// Multilevel Biased Agent Model: Centered Parameterization (CP)}
\StringTok{data \{}
\StringTok{  int\textless{}lower=1\textgreater{} N\_total;                 // Total number of observations}
\StringTok{  int\textless{}lower=1\textgreater{} N\_subjects;              // Total number of agents}
\StringTok{  array[N\_total] int\textless{}lower=1, upper=N\_subjects\textgreater{} subj\_id; // Agent ID index}
\StringTok{  array[N\_total] int\textless{}lower=0, upper=1\textgreater{} y;     // Observed choices}
\StringTok{  real prior\_mu\_theta\_mu;}
\StringTok{  real\textless{}lower=0\textgreater{} prior\_mu\_theta\_sigma;}
\StringTok{  real\textless{}lower=0\textgreater{} prior\_sigma\_theta\_lambda;}
\StringTok{  int\textless{}lower=0, upper=1\textgreater{} run\_diagnostics;}
\StringTok{\}}

\StringTok{parameters \{}
\StringTok{  // Population{-}level hyperparameters}
\StringTok{  real mu\_theta;                        // Population mean bias (logit scale)}
\StringTok{  real\textless{}lower=0\textgreater{} sigma\_theta;            // Population SD of bias (logit scale)}

\StringTok{  // Individual{-}level parameters (CENTERED)}
\StringTok{  vector[N\_subjects] theta\_logit;       // Agent{-}specific biases }
\StringTok{\}}

\StringTok{model \{}
\StringTok{  // 1. Population{-}level hyperpriors}
\StringTok{  target += normal\_lpdf(mu\_theta | prior\_mu\_theta\_mu, prior\_mu\_theta\_sigma);}
\StringTok{  target += exponential\_lpdf(sigma\_theta | prior\_sigma\_theta\_lambda);}

\StringTok{  // 2. Hierarchical prior (CENTERED)}
\StringTok{  target += normal\_lpdf(theta\_logit | mu\_theta, sigma\_theta);}

\StringTok{  // 3. Likelihood}
\StringTok{  target += bernoulli\_logit\_lpmf(y | theta\_logit[subj\_id]); }
\StringTok{\}}

\StringTok{generated quantities \{}
\StringTok{  // Transform hyperparameters back to the probability (scale)}
\StringTok{  real\textless{}lower=0, upper=1\textgreater{} mu\_theta\_prob = inv\_logit(mu\_theta);}
\StringTok{  vector\textless{}lower=0, upper=1\textgreater{}[N\_subjects] theta\_prob = inv\_logit(theta\_logit);}

\StringTok{  // Prior samples — cheap scalar draws; declared at top level so they appear in draws output}
\StringTok{  real mu\_theta\_prior    = normal\_rng(prior\_mu\_theta\_mu, prior\_mu\_theta\_sigma);}
\StringTok{  real sigma\_theta\_prior = exponential\_rng(prior\_sigma\_theta\_lambda);}

\StringTok{  // Initialize arrays for conditional diagnostics}
\StringTok{  vector[N\_total] log\_lik;}
\StringTok{  array[N\_total] int y\_post\_rep;}
\StringTok{  array[N\_total] int y\_prior\_rep;}

\StringTok{  // Accumulator for joint prior log{-}density}
\StringTok{  real lprior;}
\StringTok{  lprior = normal\_lpdf(mu\_theta | prior\_mu\_theta\_mu, prior\_mu\_theta\_sigma) +}
\StringTok{           exponential\_lpdf(sigma\_theta | prior\_sigma\_theta\_lambda) +}
\StringTok{           normal\_lpdf(theta\_logit | mu\_theta, sigma\_theta);}

\StringTok{  if (run\_diagnostics) \{}
\StringTok{    for (i in 1:N\_total) \{}
\StringTok{      log\_lik[i] = bernoulli\_logit\_lpmf(y[i] | theta\_logit[subj\_id[i]]);}
\StringTok{      y\_post\_rep[i] = bernoulli\_logit\_rng(theta\_logit[subj\_id[i]]);}
\StringTok{      real theta\_logit\_prior\_i = normal\_rng(mu\_theta\_prior, sigma\_theta\_prior);}
\StringTok{      y\_prior\_rep[i] = bernoulli\_logit\_rng(theta\_logit\_prior\_i);}
\StringTok{    \}}
\StringTok{  \}}
\StringTok{\}}
\StringTok{"}

\CommentTok{\# Save}
\NormalTok{stan\_file\_biased\_ml\_cp }\OtherTok{\textless{}{-}}\NormalTok{ here}\SpecialCharTok{::}\FunctionTok{here}\NormalTok{(}\StringTok{"stan"}\NormalTok{, }\StringTok{"ch07\_multilevel\_biased\_cp.stan"}\NormalTok{)}
\FunctionTok{writeLines}\NormalTok{(stan\_biased\_ml\_cp\_code, stan\_file\_biased\_ml\_cp)}

\CommentTok{\# Compile the CP model utilizing the CmdStan C++ backend}
\NormalTok{mod\_biased\_ml\_cp }\OtherTok{\textless{}{-}} \FunctionTok{cmdstan\_model}\NormalTok{(stan\_file\_biased\_ml\_cp)}
\end{Highlighting}
\end{Shaded}

\subsection{Fitting the model}\label{fitting-the-model}

This not really necessary yet, but a rigorous workflow demands an
explicit initialization strategy, so we have all the pieces in place for
more complex models. We will initialize the NUTS algorithm by drawing
starting values directly from our generative priors. This ensures the
chains begin in scientifically plausible regions of the parameter space,
rather than Stan's arbitrary default \(U[-2, 2]\) on the unconstrained
space.

\begin{Shaded}
\begin{Highlighting}[]
\NormalTok{fit\_cp\_path }\OtherTok{\textless{}{-}}\NormalTok{ here}\SpecialCharTok{::}\FunctionTok{here}\NormalTok{(}\StringTok{"simmodels"}\NormalTok{, }\StringTok{"ch07\_fit\_cp.rds"}\NormalTok{)}

\NormalTok{init\_cp\_list }\OtherTok{\textless{}{-}} \FunctionTok{lapply}\NormalTok{(}\DecValTok{1}\SpecialCharTok{:}\DecValTok{4}\NormalTok{, }\ControlFlowTok{function}\NormalTok{(i) \{}
  \FunctionTok{list}\NormalTok{(}
    \AttributeTok{mu\_theta    =} \FunctionTok{rnorm}\NormalTok{(}\DecValTok{1}\NormalTok{, }\DecValTok{0}\NormalTok{, }\FloatTok{1.5}\NormalTok{),}
    \AttributeTok{sigma\_theta =} \FunctionTok{rexp}\NormalTok{(}\DecValTok{1}\NormalTok{, }\DecValTok{1}\NormalTok{),}
    \AttributeTok{theta\_logit =} \FunctionTok{rnorm}\NormalTok{(stan\_data\_centered}\SpecialCharTok{$}\NormalTok{N\_subjects, }\DecValTok{0}\NormalTok{, }\FloatTok{0.5}\NormalTok{)}
\NormalTok{  )}
\NormalTok{\})}

\ControlFlowTok{if}\NormalTok{ (regenerate\_simulations }\SpecialCharTok{||} \SpecialCharTok{!}\FunctionTok{file.exists}\NormalTok{(fit\_cp\_path)) \{}
  
\NormalTok{  fit\_cp }\OtherTok{\textless{}{-}}\NormalTok{ mod\_biased\_ml\_cp}\SpecialCharTok{$}\FunctionTok{sample}\NormalTok{(}
    \AttributeTok{data            =}\NormalTok{ stan\_data\_centered,}
    \AttributeTok{seed            =} \DecValTok{123}\NormalTok{,}
    \AttributeTok{chains          =} \DecValTok{4}\NormalTok{,}
    \AttributeTok{parallel\_chains =} \DecValTok{4}\NormalTok{,}
    \AttributeTok{iter\_warmup     =} \DecValTok{1000}\NormalTok{,}
    \AttributeTok{iter\_sampling   =} \DecValTok{1000}\NormalTok{,}
    \AttributeTok{init            =}\NormalTok{ init\_cp\_list, }
    \AttributeTok{adapt\_delta     =} \FloatTok{0.85}
\NormalTok{  )}
  
\NormalTok{  fit\_cp}\SpecialCharTok{$}\FunctionTok{save\_object}\NormalTok{(fit\_cp\_path)}
  
\NormalTok{\} }\ControlFlowTok{else}\NormalTok{ \{}
\NormalTok{  fit\_cp }\OtherTok{\textless{}{-}} \FunctionTok{readRDS}\NormalTok{(fit\_cp\_path)}
\NormalTok{\}}
\end{Highlighting}
\end{Shaded}

\subsection{Assessing the fitting
process}\label{assessing-the-fitting-process}

Even if the Stan compiler finishes without throwing explicit warning
messages, our workflow demands that we systematically verify the
geometry of the posterior and the efficiency of the sampler. Hamiltonian
Monte Carlo (HMC) is a physical simulation---it explores the posterior
by simulating a particle gliding over a frictionless manifold dictated
by the joint log-probability density. We must prove the numerical
integration of this physics engine remained stable.

Let us systematically unpack the computational diagnostics of our
Centered Parameterization fit.

\begin{enumerate}
\def\labelenumi{\arabic{enumi}.}
\tightlist
\item
  The Diagnostic Summary: Integration Stability First, we ask cmdstanr
  for a high-level summary of the NUTS sampler's performance.
\end{enumerate}

\begin{Shaded}
\begin{Highlighting}[]
\CommentTok{\#fit\_cp$cmdstan\_diagnose()}
\NormalTok{fit\_cp}\SpecialCharTok{$}\FunctionTok{diagnostic\_summary}\NormalTok{()}
\end{Highlighting}
\end{Shaded}

\begin{verbatim}
$num_divergent
[1] 0 0 0 0

$num_max_treedepth
[1] 0 0 0 0

$ebfmi
[1] 0.9499418 1.0854016 0.9723112 0.9770850
\end{verbatim}

We are looking for strict zeros in num\_divergent. A divergent
transition occurs when the numerical integration process (the leapfrog
steps) fails to accurately map the curvature of the posterior manifold,
causing the simulated trajectory to fly off the energy level. Even one
divergence compromises the mathematical guarantee of the posterior.

\subsubsection{\texorpdfstring{Chain Mixing and Stationarity
(\(\hat{R}\))}{Chain Mixing and Stationarity (\textbackslash hat\{R\})}}\label{chain-mixing-and-stationarity-hatr}

Did our independent Markov chains converge to the same stationary
distribution? We visualize this using trace plots and quantify it using
\(\hat{R}\) (R-hat).

\begin{Shaded}
\begin{Highlighting}[]
\CommentTok{\# Extract the posterior array}
\NormalTok{draws\_cp }\OtherTok{\textless{}{-}}\NormalTok{ fit\_cp}\SpecialCharTok{$}\FunctionTok{draws}\NormalTok{()}

\CommentTok{\# Visual inspection of chain mixing}
\FunctionTok{mcmc\_trace}\NormalTok{(draws\_cp, }\AttributeTok{pars =} \FunctionTok{c}\NormalTok{(}\StringTok{"mu\_theta"}\NormalTok{, }\StringTok{"sigma\_theta"}\NormalTok{), }\AttributeTok{np =} \FunctionTok{nuts\_params}\NormalTok{(fit\_cp)) }\SpecialCharTok{+}
  \FunctionTok{labs}\NormalTok{(}
    \AttributeTok{title =} \StringTok{"Trace Plots: Hyperparameters (CP Model)"}\NormalTok{,}
    \AttributeTok{subtitle =} \StringTok{"Healthy = chains overlap and fluctuate around a stable band. Divergent red triangles mark sampler failures."}
\NormalTok{  )}
\end{Highlighting}
\end{Shaded}

\begin{center}
\includegraphics[width=0.8\linewidth,height=\textheight,keepaspectratio]{06-MultilevelModels_files/figure-pdf/ch07_cp_diagnostics_trace-1.pdf}
\end{center}

\begin{Shaded}
\begin{Highlighting}[]
\CommentTok{\# Quantitative assessment of convergence and efficiency}
\NormalTok{cp\_summary }\OtherTok{\textless{}{-}}\NormalTok{ draws\_cp }\SpecialCharTok{|\textgreater{}} 
  \FunctionTok{subset\_draws}\NormalTok{(}\AttributeTok{variable =} \FunctionTok{c}\NormalTok{(}\StringTok{"mu\_theta"}\NormalTok{, }\StringTok{"sigma\_theta"}\NormalTok{, }\StringTok{"theta\_logit"}\NormalTok{)) }\SpecialCharTok{|\textgreater{}} 
  \FunctionTok{summarise\_draws}\NormalTok{(}\FunctionTok{default\_convergence\_measures}\NormalTok{())}

\CommentTok{\# We require R{-}hat \textless{} 1.01 for all parameters}
\NormalTok{cp\_summary }\SpecialCharTok{|\textgreater{}} 
  \FunctionTok{filter}\NormalTok{(variable }\SpecialCharTok{\%in\%} \FunctionTok{c}\NormalTok{(}\StringTok{"mu\_theta"}\NormalTok{, }\StringTok{"sigma\_theta"}\NormalTok{, }\StringTok{"theta\_logit[1]"}\NormalTok{, }\StringTok{"theta\_logit[2]"}\NormalTok{)) }\SpecialCharTok{|\textgreater{}}
\NormalTok{  knitr}\SpecialCharTok{::}\FunctionTok{kable}\NormalTok{(}\AttributeTok{digits =} \DecValTok{3}\NormalTok{)}
\end{Highlighting}
\end{Shaded}

\begin{longtable}[]{@{}lrrr@{}}
\toprule\noalign{}
variable & rhat & ess\_bulk & ess\_tail \\
\midrule\noalign{}
\endhead
\bottomrule\noalign{}
\endlastfoot
mu\_theta & 1.000 & 4407.162 & 2910.797 \\
sigma\_theta & 1.000 & 3472.594 & 3048.045 \\
theta\_logit{[}1{]} & 1.001 & 4648.021 & 2779.811 \\
theta\_logit{[}2{]} & 1.001 & 4760.858 & 3107.729 \\
\end{longtable}

The potential scale reduction factor (\(\hat{R}\)) compares the variance
within a single chain to the variance between all chains. If the chains
are exploring the same typical set, this ratio converges to 1. Our
strict threshold is \(\hat{R} < 1.01\).

\subsubsection{Effective Sample Size
(ESS)}\label{effective-sample-size-ess}

Markov chains are inherently autocorrelated. Stan might draw 4000
samples, but we need to know the Effective Sample Size---the equivalent
number of completely independent draws.In your summary table, look at
ess\_bulk (efficiency of the posterior mean estimate) and ess\_tail
(efficiency of the posterior interval estimates). We strictly require
ESS \textgreater{} 400 per parameter. If the geometry is difficult, ESS
plummets, indicating the sampler is taking highly correlated,
inefficient steps.

\subsubsection{Energy and E-BFMI}\label{energy-and-e-bfmi}

Finally, we evaluate the Energy-Bayesian Fraction of Missing Information
(E-BFMI). HMC relies on injecting kinetic energy (momentum) into the
system to explore the tails of the target distribution.

\begin{Shaded}
\begin{Highlighting}[]
\CommentTok{\# Assess the momentum resampling efficiency}
\FunctionTok{mcmc\_nuts\_energy}\NormalTok{(}\FunctionTok{nuts\_params}\NormalTok{(fit\_cp)) }\SpecialCharTok{+}
  \FunctionTok{labs}\NormalTok{(}
    \AttributeTok{title =} \StringTok{"Energy Transitions"}\NormalTok{,}
    \AttributeTok{subtitle =} \StringTok{"Marginal and transition energy distributions must align."}
\NormalTok{  )}
\end{Highlighting}
\end{Shaded}

\begin{center}
\includegraphics[width=0.8\linewidth,height=\textheight,keepaspectratio]{06-MultilevelModels_files/figure-pdf/ch07_cp_diagnostics_energy-1.pdf}
\end{center}

If the model geometry is well-behaved, the energy transition
distribution (the energy proposed by the momentum resampling) will
neatly match the marginal energy distribution (the actual energy of the
posterior). An E-BFMI below 0.2 indicates the sampler is taking tiny,
localized steps and failing to explore the global geometry.

\subsection{Prior-Posterior Update Checks: Quantifying
Learning}\label{prior-posterior-update-checks-quantifying-learning}

Bayesian inference is the mathematical process of allocating credibility
across parameter values. The posterior is simply the prior updated by
the likelihood. If the posterior is identical to the prior, the data
contained no information. If the posterior is violently pushed against
the boundaries of the prior, our priors were likely misspecified.

Let us visualize how much the data narrowed our uncertainty about the
population parameters. Because we explicitly generated prior draws in
our Stan generated quantities block, this comparison is trivial.

\begin{Shaded}
\begin{Highlighting}[]
\CommentTok{\# Extract prior and posterior draws for the population mean and variance}
\NormalTok{draws\_df }\OtherTok{\textless{}{-}} \FunctionTok{as\_draws\_df}\NormalTok{(draws\_cp)}

\CommentTok{\# Reshape for ggplot}
\NormalTok{pp\_update\_data }\OtherTok{\textless{}{-}}\NormalTok{ draws\_df }\SpecialCharTok{|\textgreater{}} 
\NormalTok{  dplyr}\SpecialCharTok{::}\FunctionTok{select}\NormalTok{(}
    \AttributeTok{mu\_posterior =}\NormalTok{ mu\_theta, }\AttributeTok{mu\_prior =}\NormalTok{ mu\_theta\_prior,}
    \AttributeTok{sigma\_posterior =}\NormalTok{ sigma\_theta, }\AttributeTok{sigma\_prior =}\NormalTok{ sigma\_theta\_prior}
\NormalTok{  ) }\SpecialCharTok{|\textgreater{}} 
  \FunctionTok{pivot\_longer}\NormalTok{(}\FunctionTok{everything}\NormalTok{(), }\AttributeTok{names\_to =} \FunctionTok{c}\NormalTok{(}\StringTok{"parameter"}\NormalTok{, }\StringTok{"distribution"}\NormalTok{), }\AttributeTok{names\_sep =} \StringTok{"\_"}\NormalTok{)}

\CommentTok{\# Extract the true generative hyperparameters we saved as attributes}
\NormalTok{true\_mu }\OtherTok{\textless{}{-}} \FunctionTok{attr}\NormalTok{(sim\_data\_centered, }\StringTok{"hyperparameters"}\NormalTok{)}\SpecialCharTok{$}\NormalTok{mu\_theta}
\NormalTok{true\_sigma }\OtherTok{\textless{}{-}} \FunctionTok{attr}\NormalTok{(sim\_data\_centered, }\StringTok{"hyperparameters"}\NormalTok{)}\SpecialCharTok{$}\NormalTok{sigma\_theta}

\CommentTok{\# Create a reference dataframe for the ground truth matching the facet names}
\NormalTok{true\_pop\_df }\OtherTok{\textless{}{-}} \FunctionTok{tibble}\NormalTok{(}
  \AttributeTok{parameter =} \FunctionTok{c}\NormalTok{(}\StringTok{"mu"}\NormalTok{, }\StringTok{"sigma"}\NormalTok{),}
  \AttributeTok{true\_value =} \FunctionTok{c}\NormalTok{(true\_mu, true\_sigma)}
\NormalTok{)}

\CommentTok{\# Visualize the update}
\FunctionTok{ggplot}\NormalTok{(pp\_update\_data, }\FunctionTok{aes}\NormalTok{(}\AttributeTok{x =}\NormalTok{ value)) }\SpecialCharTok{+}
  \FunctionTok{geom\_histogram}\NormalTok{(}\FunctionTok{aes}\NormalTok{(}\AttributeTok{fill =}\NormalTok{ distribution), }\AttributeTok{alpha =} \FloatTok{0.6}\NormalTok{, }\AttributeTok{color =} \ConstantTok{NA}\NormalTok{) }\SpecialCharTok{+}
  \CommentTok{\# Add the true values as dashed lines}
  \FunctionTok{geom\_vline}\NormalTok{(}\AttributeTok{data =}\NormalTok{ true\_pop\_df, }\FunctionTok{aes}\NormalTok{(}\AttributeTok{xintercept =}\NormalTok{ true\_value, }\AttributeTok{linetype =} \StringTok{"True Value"}\NormalTok{), }
             \AttributeTok{color =} \StringTok{"darkred"}\NormalTok{, }\AttributeTok{linewidth =} \DecValTok{1}\NormalTok{) }\SpecialCharTok{+}
  \FunctionTok{facet\_wrap}\NormalTok{(}
    \SpecialCharTok{\textasciitilde{}}\NormalTok{parameter, }
    \AttributeTok{scales =} \StringTok{"free"}\NormalTok{,}
    \CommentTok{\# Clean up the facet labels for publication quality}
    \AttributeTok{labeller =} \FunctionTok{as\_labeller}\NormalTok{(}\FunctionTok{c}\NormalTok{(}\StringTok{"mu"} \OtherTok{=} \StringTok{"Population Mean (\textbackslash{}u03BC)"}\NormalTok{, }\StringTok{"sigma"} \OtherTok{=} \StringTok{"Population SD (\textbackslash{}u03C3)"}\NormalTok{))}
\NormalTok{  ) }\SpecialCharTok{+}
  \FunctionTok{scale\_fill\_manual}\NormalTok{(}
    \AttributeTok{name =} \StringTok{"Distribution"}\NormalTok{,}
    \AttributeTok{values =} \FunctionTok{c}\NormalTok{(}\StringTok{"prior"} \OtherTok{=} \StringTok{"gray70"}\NormalTok{, }\StringTok{"posterior"} \OtherTok{=} \StringTok{"steelblue"}\NormalTok{)}
\NormalTok{  ) }\SpecialCharTok{+}
  \FunctionTok{scale\_linetype\_manual}\NormalTok{(}
    \AttributeTok{name =} \StringTok{"Ground Truth"}\NormalTok{, }
    \AttributeTok{values =} \FunctionTok{c}\NormalTok{(}\StringTok{"True Value"} \OtherTok{=} \StringTok{"dashed"}\NormalTok{)}
\NormalTok{  ) }\SpecialCharTok{+}
  \FunctionTok{labs}\NormalTok{(}
    \AttributeTok{title =} \StringTok{"Prior{-}Posterior Update Checks: Population Hyperparameters"}\NormalTok{,}
    \AttributeTok{subtitle =} \StringTok{"Healthy = posterior (blue) narrower than prior (grey) and centred on true value (dashed red)."}\NormalTok{,}
    \AttributeTok{x =} \StringTok{"Parameter Value (Logit Scale)"}\NormalTok{,}
    \AttributeTok{y =} \StringTok{"Density"}\NormalTok{)}
\end{Highlighting}
\end{Shaded}

\begin{center}
\includegraphics[width=0.8\linewidth,height=\textheight,keepaspectratio]{06-MultilevelModels_files/figure-pdf/ch07_cp_pp_update_viz-1.pdf}
\end{center}

A healthy update shows a posterior that is significantly narrower than
the prior (indicating learning) and comfortably contained within the
prior's mass (indicating the prior did not inappropriately censor the
likelihood).

\subsection{Posterior Parameter
Correlations}\label{posterior-parameter-correlations}

Parameters in hierarchical models are structurally correlated. By
plotting the joint posterior geometry of multiple parameters, we can
detect pathological collinearity or explicitly visualize the boundary of
possible Neal's Funnel.

We will use a pairs plot to examine the joint distribution of the
population mean (\(\mu_\theta\)), population standard deviation
(\(\sigma_\theta\)), and a single agent's latent bias
(\(\theta_{\text{logit}, 1}\)).

\begin{Shaded}
\begin{Highlighting}[]
\CommentTok{\# Visualize joint posterior geometries and identify structural correlations}
\FunctionTok{mcmc\_pairs}\NormalTok{(}
\NormalTok{  draws\_cp, }
  \AttributeTok{pars =} \FunctionTok{c}\NormalTok{(}\StringTok{"mu\_theta"}\NormalTok{, }\StringTok{"sigma\_theta"}\NormalTok{, }\StringTok{"theta\_logit[1]"}\NormalTok{),}
  \AttributeTok{np =} \FunctionTok{nuts\_params}\NormalTok{(fit\_cp), }\CommentTok{\# Injects divergence data (red crosses) if any exist}
  \AttributeTok{diag\_fun =} \StringTok{"dens"}\NormalTok{,}
  \AttributeTok{off\_diag\_fun =} \StringTok{"hex"}
\NormalTok{)}
\end{Highlighting}
\end{Shaded}

\begin{center}
\includegraphics[width=0.8\linewidth,height=\textheight,keepaspectratio]{06-MultilevelModels_files/figure-pdf/ch07_cp_pairs_plot-1.pdf}
\end{center}

\subsection{Posterior Predictive Checks (PPC): Generative
Adequacy}\label{posterior-predictive-checks-ppc-generative-adequacy}

Can our fitted model actually generate the reality we observed? If we
simulate data from the posterior distribution, does it look like the
empirical data? If not, the model is structurally misspecified,
regardless of how well the chains mixed.

We will compare the observed choice rate for each agent against the
distribution of choice rates predicted by the posterior.

\begin{Shaded}
\begin{Highlighting}[]
\CommentTok{\# Safely extract vectors}
\NormalTok{y\_obs }\OtherTok{\textless{}{-}} \FunctionTok{as.numeric}\NormalTok{(stan\_data\_centered}\SpecialCharTok{$}\NormalTok{y)}
\NormalTok{agent\_idx }\OtherTok{\textless{}{-}} \FunctionTok{as.numeric}\NormalTok{(stan\_data\_centered}\SpecialCharTok{$}\NormalTok{subj\_id)}

\CommentTok{\# Use cmdstanr\textquotesingle{}s native extraction to guarantee a strict S x N matrix}
\NormalTok{y\_rep }\OtherTok{\textless{}{-}}\NormalTok{ fit\_cp}\SpecialCharTok{$}\FunctionTok{draws}\NormalTok{(}\StringTok{"y\_post\_rep"}\NormalTok{, }\AttributeTok{format =} \StringTok{"matrix"}\NormalTok{)}

\CommentTok{\# Select a random subset of agents to avoid overcrowded, illegible plots}
\NormalTok{n\_sample }\OtherTok{\textless{}{-}} \DecValTok{12}
\NormalTok{sampled\_agents }\OtherTok{\textless{}{-}} \FunctionTok{sample}\NormalTok{(}\FunctionTok{unique}\NormalTok{(agent\_idx), n\_sample)}

\CommentTok{\# Create a logical mask mapping which trials belong to the sampled agents}
\NormalTok{keep\_mask }\OtherTok{\textless{}{-}}\NormalTok{ agent\_idx }\SpecialCharTok{\%in\%}\NormalTok{ sampled\_agents}

\CommentTok{\# Subset the data exactly}
\NormalTok{y\_obs\_sub }\OtherTok{\textless{}{-}}\NormalTok{ y\_obs[keep\_mask]}
\NormalTok{agent\_idx\_sub }\OtherTok{\textless{}{-}}\NormalTok{ agent\_idx[keep\_mask]}
\CommentTok{\# For y\_rep, we subset the columns (trials)}
\NormalTok{y\_rep\_sub }\OtherTok{\textless{}{-}}\NormalTok{ y\_rep[ , keep\_mask] }

\CommentTok{\# Define a custom statistic to explicitly state our cognitive metric }
\CommentTok{\# (and simultaneously silence the bayesplot warning)}
\NormalTok{prop\_right }\OtherTok{\textless{}{-}} \ControlFlowTok{function}\NormalTok{(x) }\FunctionTok{mean}\NormalTok{(x)}

\CommentTok{\# Execute the grouped Posterior Predictive Check}
\FunctionTok{ppc\_stat\_grouped}\NormalTok{(}
  \AttributeTok{y =}\NormalTok{ y\_obs\_sub, }
  \AttributeTok{yrep =}\NormalTok{ y\_rep\_sub, }
  \AttributeTok{group =}\NormalTok{ agent\_idx\_sub, }
  \AttributeTok{stat =} \StringTok{"prop\_right"}
\NormalTok{) }\SpecialCharTok{+}
  \FunctionTok{labs}\NormalTok{(}
    \AttributeTok{title =} \StringTok{"Posterior Predictive Check: Agent{-}Level Choice Rates"}\NormalTok{,}
    \AttributeTok{subtitle =} \FunctionTok{paste0}\NormalTok{(}\StringTok{"Healthy = dark line (observed) inside histogram (predicted). Showing "}\NormalTok{, n\_sample, }\StringTok{" agents."}\NormalTok{),}
    \AttributeTok{x =} \StringTok{"Proportion of \textquotesingle{}Right\textquotesingle{} Choices"}
\NormalTok{  )}
\end{Highlighting}
\end{Shaded}

\begin{center}
\includegraphics[width=0.8\linewidth,height=\textheight,keepaspectratio]{06-MultilevelModels_files/figure-pdf/ch07_cp_ppc_viz-1.pdf}
\end{center}

\subsection{Calibration and Recovery}\label{calibration-and-recovery}

We have proven the model fits the data and generates plausible
predictions. Now, we must prove that the parameters we extract are
actually the true parameters that generated the data. Because this is a
hierarchical model, we must assess recovery at both the population and
individual levels.

\subsubsection{Individual-Level
Recovery}\label{individual-level-recovery}

Having validated the top of the hierarchy (in the prior-posterior update
check), we move to the individual agents. We plot the posterior
estimates against the ground truth. We expect the estimates to align
with the diagonal identity line, but crucially, we expect agents with
fewer trials to deviate from the raw truth, demonstrating shrinkage
toward the estimated population mean.

\begin{Shaded}
\begin{Highlighting}[]
\CommentTok{\# 1. Extract the full summary (which defaults to including median, q5, and q95)}
\NormalTok{recovery\_summary }\OtherTok{\textless{}{-}}\NormalTok{ draws\_cp }\SpecialCharTok{|\textgreater{}} 
\NormalTok{  posterior}\SpecialCharTok{::}\FunctionTok{subset\_draws}\NormalTok{(}\AttributeTok{variable =} \StringTok{"theta\_logit"}\NormalTok{) }\SpecialCharTok{|\textgreater{}} 
\NormalTok{  posterior}\SpecialCharTok{::}\FunctionTok{summarise\_draws}\NormalTok{() }\SpecialCharTok{|\textgreater{}} 
  \FunctionTok{mutate}\NormalTok{(}\AttributeTok{agent\_id =} \FunctionTok{as.integer}\NormalTok{(}\FunctionTok{str\_extract}\NormalTok{(variable, }\StringTok{"(?\textless{}=}\SpecialCharTok{\textbackslash{}\textbackslash{}}\StringTok{[)}\SpecialCharTok{\textbackslash{}\textbackslash{}}\StringTok{d+(?=}\SpecialCharTok{\textbackslash{}\textbackslash{}}\StringTok{])"}\NormalTok{)))}

\CommentTok{\# 2. Extract the true generative parameters directly from the simulated dataset}
\NormalTok{true\_agent\_params }\OtherTok{\textless{}{-}}\NormalTok{ sim\_data\_centered }\SpecialCharTok{|\textgreater{}} 
  \FunctionTok{distinct}\NormalTok{(agent\_id, n\_trials, theta\_logit, theta\_prob)}

\CommentTok{\# 3. Join the posterior estimates with the ground truth}
\NormalTok{recovery\_df }\OtherTok{\textless{}{-}}\NormalTok{ true\_agent\_params }\SpecialCharTok{|\textgreater{}} 
  \FunctionTok{left\_join}\NormalTok{(recovery\_summary, }\AttributeTok{by =} \StringTok{"agent\_id"}\NormalTok{)}

\CommentTok{\# 4. Plot Recovery with Shrinkage Dynamics}
\FunctionTok{ggplot}\NormalTok{(recovery\_df, }\FunctionTok{aes}\NormalTok{(}\AttributeTok{x =}\NormalTok{ theta\_logit, }\AttributeTok{y =}\NormalTok{ median)) }\SpecialCharTok{+}
  \FunctionTok{geom\_abline}\NormalTok{(}\AttributeTok{intercept =} \DecValTok{0}\NormalTok{, }\AttributeTok{slope =} \DecValTok{1}\NormalTok{, }\AttributeTok{linetype =} \StringTok{"dashed"}\NormalTok{, }\AttributeTok{color =} \StringTok{"gray50"}\NormalTok{) }\SpecialCharTok{+}
  \FunctionTok{geom\_errorbar}\NormalTok{(}\FunctionTok{aes}\NormalTok{(}\AttributeTok{ymin =}\NormalTok{ q5, }\AttributeTok{ymax =}\NormalTok{ q95), }\AttributeTok{width =} \DecValTok{0}\NormalTok{, }\AttributeTok{alpha =} \FloatTok{0.5}\NormalTok{, }\AttributeTok{color =} \StringTok{"steelblue"}\NormalTok{) }\SpecialCharTok{+}
  \FunctionTok{geom\_point}\NormalTok{(}\FunctionTok{aes}\NormalTok{(}\AttributeTok{size =}\NormalTok{ n\_trials), }\AttributeTok{color =} \StringTok{"midnightblue"}\NormalTok{, }\AttributeTok{alpha =} \FloatTok{0.8}\NormalTok{) }\SpecialCharTok{+}
  \FunctionTok{labs}\NormalTok{(}
    \AttributeTok{title =} \StringTok{"Parameter Recovery: Hierarchical Biased Agent"}\NormalTok{,}
    \AttributeTok{subtitle =} \StringTok{"Healthy = points near diagonal (y = x). Small dots shrink more — that is partial pooling working correctly."}\NormalTok{,}
    \AttributeTok{x =} \StringTok{"True Parameter Value (Logit Scale)"}\NormalTok{,}
    \AttributeTok{y =} \StringTok{"Posterior Median Estimate (Logit Scale)"}\NormalTok{,}
    \AttributeTok{size =} \StringTok{"Number of Trials (N\_j)"}
\NormalTok{  )}
\end{Highlighting}
\end{Shaded}

\begin{center}
\includegraphics[width=0.8\linewidth,height=\textheight,keepaspectratio]{06-MultilevelModels_files/figure-pdf/ch07_cp_recovery_shrinkage-1.pdf}
\end{center}

\subsection{Prior sensitivity}\label{prior-sensitivity}

\begin{Shaded}
\begin{Highlighting}[]
\CommentTok{\# 1. Calculate the sensitivity diagnostics}
\CommentTok{\# We target the hyperparameters, as they govern the hierarchical shrinkage}
\NormalTok{sens\_results }\OtherTok{\textless{}{-}} \FunctionTok{powerscale\_sensitivity}\NormalTok{(}
\NormalTok{  fit\_cp, }\CommentTok{\# or fit\_ncp, provided lprior is defined}
  \AttributeTok{variable =} \FunctionTok{c}\NormalTok{(}\StringTok{"mu\_theta"}\NormalTok{, }\StringTok{"sigma\_theta"}\NormalTok{)}
\NormalTok{)}

\FunctionTok{cat}\NormalTok{(}\StringTok{"=== PRIOR SENSITIVITY DIAGNOSTICS ===}\SpecialCharTok{\textbackslash{}n}\StringTok{"}\NormalTok{)}
\end{Highlighting}
\end{Shaded}

\begin{verbatim}
=== PRIOR SENSITIVITY DIAGNOSTICS ===
\end{verbatim}

\begin{Shaded}
\begin{Highlighting}[]
\FunctionTok{print}\NormalTok{(sens\_results)}
\end{Highlighting}
\end{Shaded}

\begin{verbatim}
Sensitivity based on cjs_dist
Prior selection: all priors
Likelihood selection: all data

    variable prior likelihood                                diagnosis
    mu_theta 0.165      0.050 potential strong prior / weak likelihood
 sigma_theta 0.354      0.091            potential prior-data conflict
\end{verbatim}

\begin{Shaded}
\begin{Highlighting}[]
\NormalTok{ps\_seq }\OtherTok{\textless{}{-}}\NormalTok{ priorsense}\SpecialCharTok{::}\FunctionTok{powerscale\_sequence}\NormalTok{(}
\NormalTok{  fit\_cp, }
  \AttributeTok{variable =} \FunctionTok{c}\NormalTok{(}\StringTok{"mu\_theta"}\NormalTok{, }\StringTok{"sigma\_theta"}\NormalTok{)}
\NormalTok{)}

\CommentTok{\# Plot the density shifts}
\NormalTok{priorsense}\SpecialCharTok{::}\FunctionTok{powerscale\_plot\_dens}\NormalTok{(ps\_seq) }\SpecialCharTok{+}
  \FunctionTok{theme\_classic}\NormalTok{() }\SpecialCharTok{+}
  \FunctionTok{labs}\NormalTok{(}
    \AttributeTok{title =} \StringTok{"Prior Sensitivity: Density Perturbation"}\NormalTok{,}
    \AttributeTok{subtitle =} \StringTok{"Does scaling the prior/likelihood drastically alter our geometric beliefs?"}
\NormalTok{  )}
\end{Highlighting}
\end{Shaded}

\begin{center}
\includegraphics[width=0.8\linewidth,height=\textheight,keepaspectratio]{06-MultilevelModels_files/figure-pdf/ch07_cp_prior_sensitivity-1.pdf}
\end{center}

\begin{Shaded}
\begin{Highlighting}[]
\CommentTok{\# Plot the quantitative estimate shifts}
\NormalTok{priorsense}\SpecialCharTok{::}\FunctionTok{powerscale\_plot\_quantities}\NormalTok{(ps\_seq) }\SpecialCharTok{+}
  \FunctionTok{theme\_classic}\NormalTok{() }\SpecialCharTok{+}
  \FunctionTok{labs}\NormalTok{(}
    \AttributeTok{title =} \StringTok{"Prior Sensitivity: Parameter Estimates"}\NormalTok{,}
    \AttributeTok{subtitle =} \StringTok{"Horizontal lines indicate the baseline posterior bounds. Deviations indicate sensitivity."}
\NormalTok{  )}
\end{Highlighting}
\end{Shaded}

\begin{center}
\includegraphics[width=0.8\linewidth,height=\textheight,keepaspectratio]{06-MultilevelModels_files/figure-pdf/ch07_cp_prior_sensitivity-2.pdf}
\end{center}

The sensitivity analysis reveals that our population variance
(\(\sigma_\theta\)) exhibits a prior-data conflict. This is a feature,
not a bug, of partial pooling: our Exponential prior applies
regularizing pressure toward zero variance (complete pooling), while the
empirical data resists this pull to account for individual differences.

\subsection{Simulation-Based Calibration Checks
(SBC)}\label{simulation-based-calibration-checks-sbc}

Parameter recovery on a single dataset is anecdotal. A single success
might just mean we got lucky with a benign region of the parameter
space. To mathematically guarantee that our NUTS sampler is unbiased and
our posterior uncertainties are perfectly calibrated across the entire
prior space, we must perform Simulation-Based Calibration Checks (SBC).
The logic of SBC is structurally elegant: if we generate a ground-truth
parameter \(\theta_{\text{true}}\) from the prior, generate data from
the data model conditional on \(\theta_{\text{true}}\), and compute the
posterior, then the posterior must not systematically over- or
under-estimate \(\theta_{\text{true}}\). Under a well-calibrated
sampler, the rank of the true parameter value within the sorted
posterior draws will be uniformly
distributed:\[r \sim \text{Uniform}(0, S)\]where \(S\) is the number of
posterior samples. If the model is overconfident (the posterior is too
narrow and misses the truth), the rank histogram will be U-shaped. If it
is underconfident (the posterior is too wide), it will form an inverted
U-shape. Asymmetries indicate directional bias.

\subsection{Setting up the SBC
Pipeline}\label{setting-up-the-sbc-pipeline}

The SBC package requires a strict data structure. It expects a generator
function that returns a list of two components:

\begin{itemize}
\item
  variables: A named list of the true parameter values.
\item
  generated: The formatted data list required by our Stan model.
\end{itemize}

Let us build a strict wrapper around our simulation and data preparation
functions to satisfy this requirement.

\begin{Shaded}
\begin{Highlighting}[]
\CommentTok{\# Define the strict SBC generator wrapper}
\NormalTok{generate\_sbc\_dataset }\OtherTok{\textless{}{-}} \ControlFlowTok{function}\NormalTok{(}\AttributeTok{n\_agents =} \DecValTok{50}\NormalTok{, }\AttributeTok{n\_trials\_min =} \DecValTok{20}\NormalTok{, }\AttributeTok{n\_trials\_max =} \DecValTok{120}\NormalTok{) \{}
  
  \CommentTok{\# 1. Simulate the universe (draws hyperparams and agent params from priors)}
\NormalTok{  sim\_data }\OtherTok{\textless{}{-}} \FunctionTok{simulate\_hierarchical\_biased}\NormalTok{(}
    \AttributeTok{n\_agents =}\NormalTok{ n\_agents, }
    \AttributeTok{n\_trials\_min =}\NormalTok{ n\_trials\_min, }
    \AttributeTok{n\_trials\_max =}\NormalTok{ n\_trials\_max}
\NormalTok{  )}
  
  \CommentTok{\# 2. Extract the true generative hyperparameters}
\NormalTok{  mu\_true }\OtherTok{\textless{}{-}} \FunctionTok{attr}\NormalTok{(sim\_data, }\StringTok{"hyperparameters"}\NormalTok{)}\SpecialCharTok{$}\NormalTok{mu\_theta}
\NormalTok{  sigma\_true }\OtherTok{\textless{}{-}} \FunctionTok{attr}\NormalTok{(sim\_data, }\StringTok{"hyperparameters"}\NormalTok{)}\SpecialCharTok{$}\NormalTok{sigma\_theta}
  
  \CommentTok{\# 3. Extract the true individual latent parameters (must preserve agent order)}
\NormalTok{  theta\_true }\OtherTok{\textless{}{-}}\NormalTok{ sim\_data }\SpecialCharTok{|\textgreater{}} 
    \FunctionTok{distinct}\NormalTok{(agent\_id, theta\_logit) }\SpecialCharTok{|\textgreater{}} 
    \FunctionTok{arrange}\NormalTok{(agent\_id) }\SpecialCharTok{|\textgreater{}} 
    \FunctionTok{pull}\NormalTok{(theta\_logit)}
  
  \CommentTok{\# 4. Prepare the strictly typed list for CmdStanR}
\NormalTok{  stan\_data }\OtherTok{\textless{}{-}} \FunctionTok{prep\_stan\_data}\NormalTok{(sim\_data)}
  
  \CommentTok{\# 5. Return the exact structure required by the SBC package}
  \FunctionTok{list}\NormalTok{(}
    \AttributeTok{variables =} \FunctionTok{list}\NormalTok{(}
      \AttributeTok{mu\_theta =}\NormalTok{ mu\_true,}
      \AttributeTok{sigma\_theta =}\NormalTok{ sigma\_true,}
      \AttributeTok{theta\_logit =}\NormalTok{ theta\_true}
\NormalTok{    ),}
    \AttributeTok{generated =}\NormalTok{ stan\_data}
\NormalTok{  )}
\NormalTok{\}}
\end{Highlighting}
\end{Shaded}

A production-grade SBC requires \(N_{\text{sim}} \ge 1000\) to smooth
the rank histograms. For computational feasibility in this chapter, we
will run 200 simulations. Because we established our future parallel
processing plan in the setup chunk, the SBC package will automatically
distribute these models across your CPU cores.

\begin{Shaded}
\begin{Highlighting}[]
\NormalTok{sbc\_cp\_path }\OtherTok{\textless{}{-}}\NormalTok{ here}\SpecialCharTok{::}\FunctionTok{here}\NormalTok{(}\StringTok{"simmodels"}\NormalTok{, }\StringTok{"ch07\_sbc\_cp.rds"}\NormalTok{)}

\NormalTok{generator\_cp }\OtherTok{\textless{}{-}} \FunctionTok{SBC\_generator\_function}\NormalTok{(}
\NormalTok{  generate\_sbc\_dataset,}
  \AttributeTok{n\_agents =} \DecValTok{50}\NormalTok{, }\AttributeTok{n\_trials\_min =} \DecValTok{30}\NormalTok{, }\AttributeTok{n\_trials\_max =} \DecValTok{120}
\NormalTok{)}

\NormalTok{backend\_cp }\OtherTok{\textless{}{-}} \FunctionTok{SBC\_backend\_cmdstan\_sample}\NormalTok{(}
\NormalTok{  mod\_biased\_ml\_cp,}
  \AttributeTok{iter\_warmup =} \DecValTok{1000}\NormalTok{, }\AttributeTok{iter\_sampling =} \DecValTok{1000}\NormalTok{, }\AttributeTok{chains =} \DecValTok{1}\NormalTok{, }\AttributeTok{refresh =} \DecValTok{0}
\NormalTok{)}

\ControlFlowTok{if}\NormalTok{ (regenerate\_simulations }\SpecialCharTok{||} \SpecialCharTok{!}\FunctionTok{file.exists}\NormalTok{(sbc\_cp\_path)) \{}
\NormalTok{  future}\SpecialCharTok{::}\FunctionTok{plan}\NormalTok{(future}\SpecialCharTok{::}\NormalTok{multisession, }\AttributeTok{workers =} \DecValTok{8}\NormalTok{)}
\NormalTok{  sbc\_datasets\_cp }\OtherTok{\textless{}{-}} \FunctionTok{generate\_datasets}\NormalTok{(generator\_cp, }\DecValTok{200}\NormalTok{)}
\NormalTok{  sbc\_results\_cp  }\OtherTok{\textless{}{-}} \FunctionTok{compute\_SBC}\NormalTok{(}\AttributeTok{datasets =}\NormalTok{ sbc\_datasets\_cp, }\AttributeTok{backend =}\NormalTok{ backend\_cp,}
  \AttributeTok{keep\_fits =} \ConstantTok{FALSE}\NormalTok{)}
  \FunctionTok{saveRDS}\NormalTok{(}
    \FunctionTok{list}\NormalTok{(}\AttributeTok{datasets =}\NormalTok{ sbc\_datasets\_cp, }\AttributeTok{results =}\NormalTok{ sbc\_results\_cp),}
\NormalTok{    sbc\_cp\_path}
\NormalTok{  )}
  
\NormalTok{\} }\ControlFlowTok{else}\NormalTok{ \{}
\NormalTok{  sbc\_obj        }\OtherTok{\textless{}{-}} \FunctionTok{readRDS}\NormalTok{(sbc\_cp\_path)}
\NormalTok{  sbc\_datasets\_cp }\OtherTok{\textless{}{-}}\NormalTok{ sbc\_obj}\SpecialCharTok{$}\NormalTok{datasets}
\NormalTok{  sbc\_results\_cp  }\OtherTok{\textless{}{-}}\NormalTok{ sbc\_obj}\SpecialCharTok{$}\NormalTok{results}
\NormalTok{\}}

\CommentTok{\# Visualize}
\FunctionTok{plot\_rank\_hist}\NormalTok{(sbc\_results\_cp, }\AttributeTok{variables =} \FunctionTok{c}\NormalTok{(}\StringTok{"mu\_theta"}\NormalTok{, }\StringTok{"sigma\_theta"}\NormalTok{)) }\SpecialCharTok{+}
  \FunctionTok{labs}\NormalTok{(}
    \AttributeTok{title    =} \StringTok{"SBC Rank Histograms: Centered Parameterization"}\NormalTok{,}
    \AttributeTok{subtitle =} \StringTok{"Healthy = bars flat within grey band. U{-}shape = overconfident. Here, U{-}shapes reveal the funnel pathology."}
\NormalTok{  )}
\end{Highlighting}
\end{Shaded}

\begin{center}
\includegraphics[width=0.8\linewidth,height=\textheight,keepaspectratio]{06-MultilevelModels_files/figure-pdf/ch07_sbc_cp_execution-1.pdf}
\end{center}

\begin{Shaded}
\begin{Highlighting}[]
\FunctionTok{plot\_ecdf\_diff}\NormalTok{(sbc\_results\_cp, }\AttributeTok{variables =} \FunctionTok{c}\NormalTok{(}\StringTok{"mu\_theta"}\NormalTok{, }\StringTok{"sigma\_theta"}\NormalTok{)) }\SpecialCharTok{+}
  \FunctionTok{labs}\NormalTok{(}
    \AttributeTok{title    =} \StringTok{"SBC ECDF Difference: Centered Parameterization"}\NormalTok{,}
    \AttributeTok{subtitle =} \StringTok{"Healthy = blue line stays near zero. Excursions here correlate with the divergence failures shown below."}
\NormalTok{  )}
\end{Highlighting}
\end{Shaded}

\begin{center}
\includegraphics[width=0.8\linewidth,height=\textheight,keepaspectratio]{06-MultilevelModels_files/figure-pdf/ch07_sbc_cp_execution-2.pdf}
\end{center}

\begin{Shaded}
\begin{Highlighting}[]
\CommentTok{\# 1. Extract the lightweight stats dataframe from the SBC object}
\NormalTok{sbc\_stats\_df }\OtherTok{\textless{}{-}}\NormalTok{ sbc\_results\_cp}\SpecialCharTok{$}\NormalTok{stats}

\CommentTok{\# 2. Filter for the population{-}level parameters you care about}
\CommentTok{\# (e.g., mu\_alpha, mu\_beta, rho, or even sigma[1])}
\NormalTok{recovery\_df }\OtherTok{\textless{}{-}}\NormalTok{ sbc\_stats\_df }\SpecialCharTok{|\textgreater{}}
  \FunctionTok{filter}\NormalTok{(variable }\SpecialCharTok{\%in\%} \FunctionTok{c}\NormalTok{(}\StringTok{"mu\_theta"}\NormalTok{))}

\CommentTok{\# 3. Create the Parameter Recovery Identity Plot}
\NormalTok{p\_sbc\_recovery }\OtherTok{\textless{}{-}} \FunctionTok{ggplot}\NormalTok{(recovery\_df, }\FunctionTok{aes}\NormalTok{(}\AttributeTok{x =}\NormalTok{ simulated\_value, }\AttributeTok{y =}\NormalTok{ median)) }\SpecialCharTok{+}
  \CommentTok{\# The Identity Line (Perfect Recovery)}
  \FunctionTok{geom\_abline}\NormalTok{(}\AttributeTok{intercept =} \DecValTok{0}\NormalTok{, }\AttributeTok{slope =} \DecValTok{1}\NormalTok{, }\AttributeTok{linetype =} \StringTok{"dashed"}\NormalTok{, }\AttributeTok{color =} \StringTok{"gray50"}\NormalTok{, }\AttributeTok{linewidth =} \DecValTok{1}\NormalTok{) }\SpecialCharTok{+}
  
  \CommentTok{\# The Posterior Estimates with 90\% Credible Intervals}
  \FunctionTok{geom\_errorbar}\NormalTok{(}\FunctionTok{aes}\NormalTok{(}\AttributeTok{ymin =}\NormalTok{ q5, }\AttributeTok{ymax =}\NormalTok{ q95), }\AttributeTok{width =} \DecValTok{0}\NormalTok{, }\AttributeTok{alpha =} \FloatTok{0.3}\NormalTok{, }\AttributeTok{color =} \StringTok{"steelblue"}\NormalTok{) }\SpecialCharTok{+}
  \FunctionTok{geom\_point}\NormalTok{(}\AttributeTok{color =} \StringTok{"midnightblue"}\NormalTok{, }\AttributeTok{alpha =} \FloatTok{0.7}\NormalTok{, }\AttributeTok{size =} \DecValTok{2}\NormalTok{) }\SpecialCharTok{+}
  
  \CommentTok{\# Facet by parameter}
  \FunctionTok{facet\_wrap}\NormalTok{(}\SpecialCharTok{\textasciitilde{}}\NormalTok{variable, }\AttributeTok{scales =} \StringTok{"free"}\NormalTok{) }\SpecialCharTok{+}
  
  \FunctionTok{labs}\NormalTok{(}
    \AttributeTok{title =} \StringTok{"Parameter Recovery (Extracted Directly from SBC)"}\NormalTok{,}
    \AttributeTok{subtitle =} \StringTok{"Points represent the posterior median. Bars are 90\% Credible Intervals."}\NormalTok{,}
    \AttributeTok{x =} \StringTok{"True Generative Value (Drawn from Prior)"}\NormalTok{,}
    \AttributeTok{y =} \StringTok{"Posterior Median Estimate"}
\NormalTok{  ) }\SpecialCharTok{+}
  \FunctionTok{theme\_classic}\NormalTok{() }\SpecialCharTok{+}
  \FunctionTok{theme}\NormalTok{(}\AttributeTok{strip.background =} \FunctionTok{element\_rect}\NormalTok{(}\AttributeTok{fill =} \StringTok{"gray95"}\NormalTok{, }\AttributeTok{color =} \ConstantTok{NA}\NormalTok{))}

\FunctionTok{print}\NormalTok{(p\_sbc\_recovery)}
\end{Highlighting}
\end{Shaded}

\begin{center}
\includegraphics[width=0.8\linewidth,height=\textheight,keepaspectratio]{06-MultilevelModels_files/figure-pdf/ch07_sbc_cp_execution-3.pdf}
\end{center}

\begin{Shaded}
\begin{Highlighting}[]
\CommentTok{\# 4. Quantify the Recovery (Correlation and RMSE)}
\NormalTok{recovery\_metrics }\OtherTok{\textless{}{-}}\NormalTok{ recovery\_df }\SpecialCharTok{|\textgreater{}}
  \FunctionTok{group\_by}\NormalTok{(variable) }\SpecialCharTok{|\textgreater{}}
  \FunctionTok{summarize}\NormalTok{(}
    \AttributeTok{Pearson\_r =} \FunctionTok{cor}\NormalTok{(simulated\_value, median),}
    \AttributeTok{Spearman\_rho =} \FunctionTok{cor}\NormalTok{(simulated\_value, median, }\AttributeTok{method =} \StringTok{"spearman"}\NormalTok{),}
    \AttributeTok{RMSE =} \FunctionTok{sqrt}\NormalTok{(}\FunctionTok{mean}\NormalTok{((simulated\_value }\SpecialCharTok{{-}}\NormalTok{ median)}\SpecialCharTok{\^{}}\DecValTok{2}\NormalTok{)),}
    \AttributeTok{.groups =} \StringTok{"drop"}
\NormalTok{  )}

\FunctionTok{cat}\NormalTok{(}\StringTok{"=== PARAMETER RECOVERY QUANTIFICATION ===}\SpecialCharTok{\textbackslash{}n}\StringTok{"}\NormalTok{)}
\end{Highlighting}
\end{Shaded}

\begin{verbatim}
=== PARAMETER RECOVERY QUANTIFICATION ===
\end{verbatim}

\begin{Shaded}
\begin{Highlighting}[]
\FunctionTok{print}\NormalTok{(knitr}\SpecialCharTok{::}\FunctionTok{kable}\NormalTok{(recovery\_metrics, }\AttributeTok{digits =} \DecValTok{3}\NormalTok{))}
\end{Highlighting}
\end{Shaded}

\begin{verbatim}


|variable | Pearson_r| Spearman_rho|  RMSE|
|:--------|---------:|------------:|-----:|
|mu_theta |      0.99|        0.989| 0.212|
\end{verbatim}

Let's explore the warnings we got.

\begin{Shaded}
\begin{Highlighting}[]
\CommentTok{\# Extract diagnostics now so they are available for the targeted refit below}
\NormalTok{diags\_cp }\OtherTok{\textless{}{-}}\NormalTok{ sbc\_results\_cp}\SpecialCharTok{$}\NormalTok{backend\_diagnostics}
\NormalTok{cp\_fail\_rate }\OtherTok{\textless{}{-}} \FunctionTok{mean}\NormalTok{(diags\_cp}\SpecialCharTok{$}\NormalTok{n\_divergent }\SpecialCharTok{\textgreater{}} \DecValTok{0}\NormalTok{) }\SpecialCharTok{*} \DecValTok{100}

\CommentTok{\# Divergence count distribution across all 200 universes}
\NormalTok{diags\_cp }\SpecialCharTok{|\textgreater{}} 
  \FunctionTok{ggplot}\NormalTok{(}\FunctionTok{aes}\NormalTok{(}\AttributeTok{x =}\NormalTok{ n\_divergent)) }\SpecialCharTok{+}
  \FunctionTok{geom\_histogram}\NormalTok{(}\AttributeTok{binwidth =} \DecValTok{1}\NormalTok{, }\AttributeTok{fill =} \StringTok{"darkred"}\NormalTok{, }\AttributeTok{alpha =} \FloatTok{0.8}\NormalTok{) }\SpecialCharTok{+}
  \FunctionTok{labs}\NormalTok{(}
    \AttributeTok{title =} \StringTok{"CP Divergence Counts Across 200 SBC Universes"}\NormalTok{,}
    \AttributeTok{subtitle =} \FunctionTok{sprintf}\NormalTok{(}\StringTok{"\%.0f\%\% of universes produced at least one divergence. }
\StringTok{Even a single divergence invalidates the geometric guarantee of NUTS."}\NormalTok{, cp\_fail\_rate),}
    \AttributeTok{x =} \StringTok{"Number of Divergent Transitions"}\NormalTok{,}
    \AttributeTok{y =} \StringTok{"Number of SBC Universes"}
\NormalTok{  )}
\end{Highlighting}
\end{Shaded}

\begin{center}
\includegraphics[width=0.8\linewidth,height=\textheight,keepaspectratio]{06-MultilevelModels_files/figure-pdf/ch07_sbc_cp_divergence_viz-1.pdf}
\end{center}

\begin{Shaded}
\begin{Highlighting}[]
\CommentTok{\# Locate the divergence{-}prone region: do they concentrate at small sigma\_theta?}
\NormalTok{divergence\_profile }\OtherTok{\textless{}{-}}\NormalTok{ posterior}\SpecialCharTok{::}\FunctionTok{as\_draws\_df}\NormalTok{(sbc\_datasets\_cp}\SpecialCharTok{$}\NormalTok{variables) }\SpecialCharTok{|\textgreater{}} 
  \FunctionTok{as\_tibble}\NormalTok{() }\SpecialCharTok{|\textgreater{}} 
  \FunctionTok{mutate}\NormalTok{(}\AttributeTok{sim\_id =} \FunctionTok{row\_number}\NormalTok{()) }\SpecialCharTok{|\textgreater{}} 
  \FunctionTok{left\_join}\NormalTok{(}
\NormalTok{    diags\_cp }\SpecialCharTok{|\textgreater{}} \FunctionTok{mutate}\NormalTok{(}\AttributeTok{sim\_id =} \FunctionTok{row\_number}\NormalTok{()) }\SpecialCharTok{|\textgreater{}}\NormalTok{ dplyr}\SpecialCharTok{::}\FunctionTok{select}\NormalTok{(sim\_id, n\_divergent),}
    \AttributeTok{by =} \StringTok{"sim\_id"}
\NormalTok{  ) }\SpecialCharTok{|\textgreater{}} 
  \FunctionTok{mutate}\NormalTok{(}\AttributeTok{had\_divergence =}\NormalTok{ n\_divergent }\SpecialCharTok{\textgreater{}} \DecValTok{0}\NormalTok{)}

\FunctionTok{ggplot}\NormalTok{(divergence\_profile, }\FunctionTok{aes}\NormalTok{(}\AttributeTok{x =}\NormalTok{ sigma\_theta, }\AttributeTok{y =}\NormalTok{ n\_divergent, }\AttributeTok{color =}\NormalTok{ had\_divergence)) }\SpecialCharTok{+}
  \FunctionTok{geom\_point}\NormalTok{(}\AttributeTok{alpha =} \FloatTok{0.5}\NormalTok{, }\AttributeTok{size =} \FloatTok{1.5}\NormalTok{) }\SpecialCharTok{+}
  \FunctionTok{scale\_color\_manual}\NormalTok{(}
    \AttributeTok{values =} \FunctionTok{c}\NormalTok{(}\StringTok{"FALSE"} \OtherTok{=} \StringTok{"steelblue"}\NormalTok{, }\StringTok{"TRUE"} \OtherTok{=} \StringTok{"darkred"}\NormalTok{),}
    \AttributeTok{labels =} \FunctionTok{c}\NormalTok{(}\StringTok{"FALSE"} \OtherTok{=} \StringTok{"Clean"}\NormalTok{, }\StringTok{"TRUE"} \OtherTok{=} \StringTok{"Divergent"}\NormalTok{),}
    \AttributeTok{name =} \ConstantTok{NULL}
\NormalTok{  ) }\SpecialCharTok{+}
  \FunctionTok{labs}\NormalTok{(}
    \AttributeTok{title =} \StringTok{"CP Divergences Concentrate at Small True sigma\_theta"}\NormalTok{,}
    \AttributeTok{subtitle =} \StringTok{"The funnel geometry only manifests when population variance is small."}\NormalTok{,}
    \AttributeTok{x =} \FunctionTok{expression}\NormalTok{(}\StringTok{"True "} \SpecialCharTok{*}\NormalTok{ sigma[theta]),}
    \AttributeTok{y =} \StringTok{"Number of Divergent Transitions"}
\NormalTok{  )}
\end{Highlighting}
\end{Shaded}

\begin{center}
\includegraphics[width=0.8\linewidth,height=\textheight,keepaspectratio]{06-MultilevelModels_files/figure-pdf/ch07_sbc_cp_divergence_viz-2.pdf}
\end{center}

\begin{Shaded}
\begin{Highlighting}[]
\FunctionTok{ggplot}\NormalTok{(divergence\_profile, }\FunctionTok{aes}\NormalTok{(}\AttributeTok{x =}\NormalTok{ mu\_theta, }\AttributeTok{y =}\NormalTok{ sigma\_theta, }\AttributeTok{color =}\NormalTok{ had\_divergence)) }\SpecialCharTok{+}
  \FunctionTok{geom\_point}\NormalTok{(}\AttributeTok{alpha =} \FloatTok{0.7}\NormalTok{, }\AttributeTok{size =} \DecValTok{2}\NormalTok{) }\SpecialCharTok{+}
  \FunctionTok{scale\_color\_manual}\NormalTok{(}
    \AttributeTok{name =} \StringTok{"Sampler Status"}\NormalTok{,}
    \AttributeTok{values =} \FunctionTok{c}\NormalTok{(}\StringTok{"FALSE"} \OtherTok{=} \StringTok{"steelblue"}\NormalTok{, }\StringTok{"TRUE"} \OtherTok{=} \StringTok{"darkred"}\NormalTok{),}
    \AttributeTok{labels =} \FunctionTok{c}\NormalTok{(}\StringTok{"Healthy"}\NormalTok{, }\StringTok{"Diverged"}\NormalTok{)}
\NormalTok{  ) }\SpecialCharTok{+}
  \FunctionTok{labs}\NormalTok{(}
    \AttributeTok{title =} \StringTok{"The Anatomy of Failure: Where does CP break?"}\NormalTok{,}
    \AttributeTok{x =} \StringTok{"True Population Mean (mu\_theta)"}\NormalTok{,}
    \AttributeTok{y =} \StringTok{"True Population Variance (sigma\_theta)"}
\NormalTok{  ) }\SpecialCharTok{+}
  \FunctionTok{theme\_classic}\NormalTok{()}
\end{Highlighting}
\end{Shaded}

\begin{center}
\includegraphics[width=0.8\linewidth,height=\textheight,keepaspectratio]{06-MultilevelModels_files/figure-pdf/ch07_sbc_cp_divergence_viz-3.pdf}
\end{center}

\begin{Shaded}
\begin{Highlighting}[]
\NormalTok{hard\_universe\_path }\OtherTok{\textless{}{-}}\NormalTok{ here}\SpecialCharTok{::}\FunctionTok{here}\NormalTok{(}\StringTok{"simmodels"}\NormalTok{, }\StringTok{"ch07\_fit\_cp\_hard.rds"}\NormalTok{)}

\NormalTok{hard\_universe\_idx }\OtherTok{\textless{}{-}}\NormalTok{ divergence\_profile }\SpecialCharTok{|\textgreater{}} 
  \FunctionTok{filter}\NormalTok{(had\_divergence) }\SpecialCharTok{|\textgreater{}} 
  \FunctionTok{slice\_min}\NormalTok{(sigma\_theta, }\AttributeTok{n =} \DecValTok{1}\NormalTok{) }\SpecialCharTok{|\textgreater{}} 
  \FunctionTok{pull}\NormalTok{(sim\_id)}

\NormalTok{hard\_stan\_data   }\OtherTok{\textless{}{-}}\NormalTok{ sbc\_datasets\_cp}\SpecialCharTok{$}\NormalTok{generated[[hard\_universe\_idx]]}
\CommentTok{\# Ensure compatibility with updated Stan model if cache is stale}
\ControlFlowTok{if}\NormalTok{ (}\SpecialCharTok{!}\StringTok{"prior\_mu\_theta\_mu"} \SpecialCharTok{\%in\%} \FunctionTok{names}\NormalTok{(hard\_stan\_data)) \{}
\NormalTok{  hard\_stan\_data}\SpecialCharTok{$}\NormalTok{prior\_mu\_theta\_mu        }\OtherTok{\textless{}{-}} \DecValTok{0}
\NormalTok{  hard\_stan\_data}\SpecialCharTok{$}\NormalTok{prior\_mu\_theta\_sigma     }\OtherTok{\textless{}{-}} \FloatTok{1.5}
\NormalTok{  hard\_stan\_data}\SpecialCharTok{$}\NormalTok{prior\_sigma\_theta\_lambda }\OtherTok{\textless{}{-}} \DecValTok{1}
\NormalTok{  hard\_stan\_data}\SpecialCharTok{$}\NormalTok{run\_diagnostics          }\OtherTok{\textless{}{-}} \DecValTok{1}
\NormalTok{\}}
\NormalTok{hard\_true\_sigma  }\OtherTok{\textless{}{-}}\NormalTok{ divergence\_profile}\SpecialCharTok{$}\NormalTok{sigma\_theta[hard\_universe\_idx]}

\FunctionTok{cat}\NormalTok{(}\FunctionTok{sprintf}\NormalTok{(}
  \StringTok{"Selected universe \%d: true sigma\_theta = \%.3f}\SpecialCharTok{\textbackslash{}n}\StringTok{"}\NormalTok{,}
\NormalTok{  hard\_universe\_idx, hard\_true\_sigma}
\NormalTok{))}
\end{Highlighting}
\end{Shaded}

\begin{verbatim}
Selected universe 198: true sigma_theta = 0.016
\end{verbatim}

\begin{Shaded}
\begin{Highlighting}[]
\ControlFlowTok{if}\NormalTok{ (regenerate\_simulations }\SpecialCharTok{||} \SpecialCharTok{!}\FunctionTok{file.exists}\NormalTok{(hard\_universe\_path)) \{}
\NormalTok{  fit\_cp\_hard }\OtherTok{\textless{}{-}}\NormalTok{ mod\_biased\_ml\_cp}\SpecialCharTok{$}\FunctionTok{sample}\NormalTok{(}
    \AttributeTok{data            =}\NormalTok{ hard\_stan\_data,}
    \AttributeTok{seed            =} \DecValTok{999}\NormalTok{,}
    \AttributeTok{chains          =} \DecValTok{4}\NormalTok{,}
    \AttributeTok{parallel\_chains =} \DecValTok{4}\NormalTok{,}
    \AttributeTok{iter\_warmup     =} \DecValTok{1000}\NormalTok{,}
    \AttributeTok{iter\_sampling   =} \DecValTok{1000}\NormalTok{,}
    \AttributeTok{init            =}\NormalTok{ init\_cp\_list,}
    \AttributeTok{adapt\_delta     =} \FloatTok{0.85}\NormalTok{, }
    \AttributeTok{refresh         =} \DecValTok{0}
\NormalTok{  )}
\NormalTok{  fit\_cp\_hard}\SpecialCharTok{$}\FunctionTok{save\_object}\NormalTok{(hard\_universe\_path)}
\NormalTok{\} }\ControlFlowTok{else}\NormalTok{ \{}
\NormalTok{  fit\_cp\_hard }\OtherTok{\textless{}{-}} \FunctionTok{readRDS}\NormalTok{(hard\_universe\_path)}
\NormalTok{\}}

\CommentTok{\# 1. Generate the pairs plot }
\NormalTok{p\_funnel }\OtherTok{\textless{}{-}} \FunctionTok{mcmc\_pairs}\NormalTok{(}
\NormalTok{  fit\_cp\_hard}\SpecialCharTok{$}\FunctionTok{draws}\NormalTok{(),}
  \AttributeTok{pars     =} \FunctionTok{c}\NormalTok{(}\StringTok{"mu\_theta"}\NormalTok{, }\StringTok{"sigma\_theta"}\NormalTok{, }\StringTok{"theta\_logit[1]"}\NormalTok{),}
  \AttributeTok{np       =} \FunctionTok{nuts\_params}\NormalTok{(fit\_cp\_hard),}
  \AttributeTok{np\_style =} \FunctionTok{pairs\_style\_np}\NormalTok{(}\AttributeTok{div\_color =} \StringTok{"darkred"}\NormalTok{, }\AttributeTok{div\_alpha =} \FloatTok{0.8}\NormalTok{, }\AttributeTok{div\_size =} \DecValTok{2}\NormalTok{),}
  \AttributeTok{diag\_fun     =} \StringTok{"dens"}\NormalTok{,}
  \AttributeTok{off\_diag\_fun =} \StringTok{"hex"}
\NormalTok{)}

\FunctionTok{wrap\_elements}\NormalTok{(p\_funnel) }\SpecialCharTok{+} 
  \FunctionTok{plot\_annotation}\NormalTok{(}
    \AttributeTok{title =} \FunctionTok{sprintf}\NormalTok{(}\StringTok{"Neal\textquotesingle{}s Funnel: CP Geometry at True sigma\_theta = \%.3f"}\NormalTok{, hard\_true\_sigma),}
    \AttributeTok{subtitle =} \StringTok{"Red points are divergent transitions crashing at the funnel neck."}
\NormalTok{  )}
\end{Highlighting}
\end{Shaded}

\begin{center}
\includegraphics[width=0.8\linewidth,height=\textheight,keepaspectratio]{06-MultilevelModels_files/figure-pdf/ch07_cp_hard_universe_fit-1.pdf}
\end{center}

Here we see a funnel. As sigma goes to zero, theta\_logit{[}1{]} can
only have one value: the value of the population (mu\_theta). The
curvature becomes infinitely steep, and the NUTS sampler fails to
explore this region, resulting in divergent transitions (red points).
This is the geometric pathology of the Centered Parameterization.

\subsection{Non-Centered Parameterization (NCP) as the
Solution}\label{non-centered-parameterization-ncp-as-the-solution}

\begin{quote}
\emph{Why does a coordinate transform fix a geometry problem?} Neal's
Funnel is not a flaw in the model --- it is a flaw in the \emph{space
the sampler must navigate}. In the CP, individual parameters
\(\theta_j\) live in a space whose width is controlled by
\(\sigma_\theta\): as \(\sigma_\theta
\to 0\) the space collapses. In the NCP, we sample standardized offsets
\(z_j \sim N(0,1)\) whose geometry is always the same --- the funnel
disappears because the sampler no longer navigates the narrowing space
directly. The parameters \(\theta_j\) are reconstructed \emph{after}
sampling, not during it.
\end{quote}

We have reached a critical juncture in our Bayesian workflow. We showed
computationally (via SBC and diagnostic pairs plots) that the Centered
Parameterization (CP) creates a pathological geometry---Neal's
Funnel---that makes it difficult for the NUTS sampler to efficiently
traverse the posterior space when data is sparse or population variance
is small. When the model is misbehaving, sometimes we need to reconsider
our theory and/or our data. Yet, in this case, we can do something to
the math underlying the model to topologically transform the posterior
space the model is sampling: a coordinate transformation known as the
Non-Centered Parameterization (NCP).

\subsubsection{The Geometry of the
Transformation}\label{the-geometry-of-the-transformation}

In the Centered Parameterization, the hierarchical dependency is
explicit and
direct:\[\theta_{\text{logit}, j} \sim \text{Normal}(\mu_\theta, \sigma_\theta)\]
This forces the sampler to navigate a space where the permissible bounds
of \(\theta\) are dynamically restricted by \(\sigma\), creating the
funnel. In the Non-Centered Parameterization, we decouple this
dependency, in a way that is not dissimilar from rescaling your
predictors - a common procedure you probably learned in your statistical
course to help model convergence and interpretation.

We introduce an auxiliary, standardized latent variable \(z_j\) for each
agent. We force the sampler to explore this \(z\)-space, which is simply
a perfectly behaved Normal distribution:\[z_j \sim \text{Normal}(0, 1)\]
Then, we deterministically shift and scale \(z_j\) to reconstruct our
target cognitive
parameter:\[\theta_{\text{logit}, j} = \mu_\theta + z_j \times \sigma_\theta\]

Mathematically, this is identical to the CP
(\(z \sim N(0,1) \implies \mu + z\sigma \sim N(\mu, \sigma)\)). However,
the computational geometry is entirely transformed and when the sampler
traverses it, the problems it meets are likely different.

\subsection{Stan Model: Multilevel Biased Agent
(NCP)}\label{stan-model-multilevel-biased-agent-ncp}

Here is the NCP version of the model. Note the changes in the
parameters, transformed parameters, and model blocks compared to the CP
version.

\begin{Shaded}
\begin{Highlighting}[]
\NormalTok{stan\_biased\_ml\_ncp\_code }\OtherTok{\textless{}{-}} \StringTok{"}
\StringTok{// Multilevel Biased Agent Model (Non{-}Centered Parameterization)}
\StringTok{// Estimates individual biases drawn from a population distribution.}

\StringTok{data \{}
\StringTok{  int\textless{}lower=1\textgreater{} N\_total;                 // Total number of observations}
\StringTok{  int\textless{}lower=1\textgreater{} N\_subjects;              // Number of agents}
\StringTok{  array[N\_total] int\textless{}lower=1, upper=N\_subjects\textgreater{} subj\_id; // Agent ID}
\StringTok{  array[N\_total] int\textless{}lower=0, upper=1\textgreater{} y;     // Observed choices}
\StringTok{  real prior\_mu\_theta\_mu;}
\StringTok{  real\textless{}lower=0\textgreater{} prior\_mu\_theta\_sigma;}
\StringTok{  real\textless{}lower=0\textgreater{} prior\_sigma\_theta\_lambda;}
\StringTok{  int\textless{}lower=0, upper=1\textgreater{} run\_diagnostics;}
\StringTok{\}}

\StringTok{parameters \{}
\StringTok{  // Population{-}level parameters}
\StringTok{  real mu\_theta;                     // Population mean bias (logit scale)}
\StringTok{  real\textless{}lower=0\textgreater{} sigma\_theta;         // Population SD of bias (logit scale)}

\StringTok{  // Individual{-}level parameters (standardized deviations, z{-}scores of the agents)}
\StringTok{  vector[N\_subjects] z\_theta;        // Non{-}centered individual effects}
\StringTok{\}}

\StringTok{transformed parameters \{}
\StringTok{  // Deterministic reconstruction of the cognitive parameters}
\StringTok{  vector[N\_subjects] theta\_logit = mu\_theta + z\_theta * sigma\_theta;}
\StringTok{\}}

\StringTok{model \{}
\StringTok{  // Priors for population{-}level parameters}
\StringTok{  target += normal\_lpdf(mu\_theta | prior\_mu\_theta\_mu, prior\_mu\_theta\_sigma);}
\StringTok{  target += exponential\_lpdf(sigma\_theta | prior\_sigma\_theta\_lambda);}
\StringTok{  }
\StringTok{  /// Individual level prior (NON{-}CENTERED)}
\StringTok{  target += std\_normal\_lpdf(z\_theta);}

\StringTok{  // Likelihood}
\StringTok{  target += bernoulli\_logit\_lpmf(y | theta\_logit[subj\_id]);}
\StringTok{\}}

\StringTok{generated quantities \{}
\StringTok{  // Transform hyperparameters back to the probability (scale)}
\StringTok{  real\textless{}lower=0, upper=1\textgreater{} mu\_theta\_prob = inv\_logit(mu\_theta);}
\StringTok{  vector\textless{}lower=0, upper=1\textgreater{}[N\_subjects] theta\_prob = inv\_logit(theta\_logit);}
\StringTok{  }
\StringTok{  // Initialize containers for trial{-}level metrics}
\StringTok{  vector[N\_total] log\_lik;}
\StringTok{  array[N\_total] int y\_post\_rep;}
\StringTok{  array[N\_total] int y\_prior\_rep;}

\StringTok{  // {-}{-}{-} Prior Log{-}Density  {-}{-}{-}}
\StringTok{  real lprior;}
\StringTok{  lprior = normal\_lpdf(mu\_theta | prior\_mu\_theta\_mu, prior\_mu\_theta\_sigma) + }
\StringTok{           exponential\_lpdf(sigma\_theta | prior\_sigma\_theta\_lambda) + }
\StringTok{           std\_normal\_lpdf(z\_theta);}

\StringTok{  // {-}{-}{-} Prior Predictive Checks: Generative Baseline {-}{-}{-}}
\StringTok{  real mu\_theta\_prior = normal\_rng(prior\_mu\_theta\_mu, prior\_mu\_theta\_sigma);}
\StringTok{  real sigma\_theta\_prior = exponential\_rng(prior\_sigma\_theta\_lambda);}

\StringTok{  if (run\_diagnostics) \{}
\StringTok{    // 2. Draw individual agent parameters from the non{-}centered prior}
\StringTok{    vector[N\_subjects] theta\_logit\_prior;}
\StringTok{    for (j in 1:N\_subjects) \{}
\StringTok{      real z\_theta\_prior = normal\_rng(0, 1); }
\StringTok{      theta\_logit\_prior[j] = mu\_theta\_prior + z\_theta\_prior * sigma\_theta\_prior;}
\StringTok{    \}}

\StringTok{    // {-}{-}{-} Trial{-}Level Computations {-}{-}{-}}
\StringTok{    for (i in 1:N\_total) \{}
\StringTok{      log\_lik[i] = bernoulli\_logit\_lpmf(y[i] | theta\_logit[subj\_id[i]]);}
\StringTok{      y\_post\_rep[i] = bernoulli\_logit\_rng(theta\_logit[subj\_id[i]]);}
\StringTok{      y\_prior\_rep[i] = bernoulli\_logit\_rng(theta\_logit\_prior[subj\_id[i]]);}
\StringTok{    \}}
\StringTok{  \}}
\StringTok{\}"}

\CommentTok{\# Write Stan code to file}
\NormalTok{stan\_file\_ncp }\OtherTok{\textless{}{-}}\NormalTok{ here}\SpecialCharTok{::}\FunctionTok{here}\NormalTok{(}\StringTok{"stan"}\NormalTok{, }\StringTok{"ch07\_multilevel\_biased\_ncp.stan"}\NormalTok{)}
\FunctionTok{writeLines}\NormalTok{(stan\_biased\_ml\_ncp\_code, stan\_file\_ncp)}

\NormalTok{mod\_biased\_ml\_ncp }\OtherTok{\textless{}{-}} \FunctionTok{cmdstan\_model}\NormalTok{(stan\_file\_ncp)}

\NormalTok{fit\_ncp\_path }\OtherTok{\textless{}{-}}\NormalTok{ here}\SpecialCharTok{::}\FunctionTok{here}\NormalTok{(}\StringTok{"simmodels"}\NormalTok{, }\StringTok{"ch07\_fit\_ncp.rds"}\NormalTok{)}

\CommentTok{\# Extract N\_subjects safely}
\NormalTok{  N\_subj\_val }\OtherTok{\textless{}{-}}\NormalTok{ stan\_data\_centered}\SpecialCharTok{$}\NormalTok{N\_subjects}
  
  \CommentTok{\# The dumb list (4 chains)}
\NormalTok{  init\_ncp\_list }\OtherTok{\textless{}{-}} \FunctionTok{list}\NormalTok{(}
    \FunctionTok{list}\NormalTok{(}\AttributeTok{mu\_theta =} \FunctionTok{rnorm}\NormalTok{(}\DecValTok{1}\NormalTok{, }\DecValTok{0}\NormalTok{, }\FloatTok{1.5}\NormalTok{), }\AttributeTok{sigma\_theta =} \FunctionTok{rexp}\NormalTok{(}\DecValTok{1}\NormalTok{, }\DecValTok{1}\NormalTok{), }\AttributeTok{z\_theta =} \FunctionTok{rnorm}\NormalTok{(N\_subj\_val, }\DecValTok{0}\NormalTok{, }\DecValTok{1}\NormalTok{)),}
    \FunctionTok{list}\NormalTok{(}\AttributeTok{mu\_theta =} \FunctionTok{rnorm}\NormalTok{(}\DecValTok{1}\NormalTok{, }\DecValTok{0}\NormalTok{, }\FloatTok{1.5}\NormalTok{), }\AttributeTok{sigma\_theta =} \FunctionTok{rexp}\NormalTok{(}\DecValTok{1}\NormalTok{, }\DecValTok{1}\NormalTok{), }\AttributeTok{z\_theta =} \FunctionTok{rnorm}\NormalTok{(N\_subj\_val, }\DecValTok{0}\NormalTok{, }\DecValTok{1}\NormalTok{)),}
    \FunctionTok{list}\NormalTok{(}\AttributeTok{mu\_theta =} \FunctionTok{rnorm}\NormalTok{(}\DecValTok{1}\NormalTok{, }\DecValTok{0}\NormalTok{, }\FloatTok{1.5}\NormalTok{), }\AttributeTok{sigma\_theta =} \FunctionTok{rexp}\NormalTok{(}\DecValTok{1}\NormalTok{, }\DecValTok{1}\NormalTok{), }\AttributeTok{z\_theta =} \FunctionTok{rnorm}\NormalTok{(N\_subj\_val, }\DecValTok{0}\NormalTok{, }\DecValTok{1}\NormalTok{)),}
    \FunctionTok{list}\NormalTok{(}\AttributeTok{mu\_theta =} \FunctionTok{rnorm}\NormalTok{(}\DecValTok{1}\NormalTok{, }\DecValTok{0}\NormalTok{, }\FloatTok{1.5}\NormalTok{), }\AttributeTok{sigma\_theta =} \FunctionTok{rexp}\NormalTok{(}\DecValTok{1}\NormalTok{, }\DecValTok{1}\NormalTok{), }\AttributeTok{z\_theta =} \FunctionTok{rnorm}\NormalTok{(N\_subj\_val, }\DecValTok{0}\NormalTok{, }\DecValTok{1}\NormalTok{))}
\NormalTok{  )}

  \ControlFlowTok{if}\NormalTok{ (regenerate\_simulations }\SpecialCharTok{||} \SpecialCharTok{!}\FunctionTok{file.exists}\NormalTok{(fit\_ncp\_path)) \{}
\NormalTok{    fit\_ncp }\OtherTok{\textless{}{-}}\NormalTok{ mod\_biased\_ml\_ncp}\SpecialCharTok{$}\FunctionTok{sample}\NormalTok{(}
      \AttributeTok{data            =}\NormalTok{ stan\_data\_centered,}
      \AttributeTok{seed            =} \DecValTok{123}\NormalTok{,}
      \AttributeTok{chains          =} \DecValTok{4}\NormalTok{,}
      \AttributeTok{parallel\_chains =} \DecValTok{4}\NormalTok{,}
      \AttributeTok{iter\_warmup     =} \DecValTok{1000}\NormalTok{,}
      \AttributeTok{iter\_sampling   =} \DecValTok{1000}\NormalTok{,}
      \AttributeTok{init            =}\NormalTok{ init\_ncp\_list,}
      \AttributeTok{adapt\_delta     =} \FloatTok{0.85}\NormalTok{,}
      \AttributeTok{refresh         =} \DecValTok{0}
\NormalTok{    )}
\NormalTok{    fit\_ncp}\SpecialCharTok{$}\FunctionTok{save\_object}\NormalTok{(fit\_ncp\_path)}

\NormalTok{\} }\ControlFlowTok{else}\NormalTok{ \{}
\NormalTok{  fit\_ncp }\OtherTok{\textless{}{-}} \FunctionTok{readRDS}\NormalTok{(fit\_ncp\_path)}
\NormalTok{\}}
\end{Highlighting}
\end{Shaded}

\begin{quote}
\textbf{Why we skip the full diagnostic battery here:} The NCP model is
not our production model for this dataset --- it is a \emph{pedagogical
fix}. We introduced the CP model specifically to watch it fail via SBC,
and the NCP to demonstrate the geometric solution. Running the full
six-phase battery (prior predictive, posterior predictive,
prior-posterior update, sensitivity, recovery, SBC) on both
parameterizations would be instructive but redundant for this chapter's
purpose. When you apply the NCP to your \emph{own} data, run all six
phases. Here we focus on the two diagnostics that directly reveal the
CP's failure and the NCP's fix: sampling efficiency and the SBC rank
histograms.
\end{quote}

\begin{tcolorbox}[enhanced jigsaw, toprule=.15mm, toptitle=1mm, bottomtitle=1mm, leftrule=.75mm, rightrule=.15mm, colframe=quarto-callout-tip-color-frame, left=2mm, colback=white, titlerule=0mm, coltitle=black, colbacktitle=quarto-callout-tip-color!10!white, opacityback=0, title=\textcolor{quarto-callout-tip-color}{\faLightbulb}\hspace{0.5em}{Stan Efficiency Rule 3 --- Subject Indexing: \texttt{subj\_id}
vs.~Contiguous Blocks}, arc=.35mm, opacitybacktitle=0.6, breakable, bottomrule=.15mm]

The models in this chapter index subject membership via an integer
array: \texttt{array{[}N\_total{]}\ int\ subj\_id}. On each trial row,
\texttt{subj\_id{[}i{]}} says which subject that row belongs to. This
\textbf{scatter-index} pattern is simple and works perfectly for models
where trial order does not matter --- like the biased-agent model here,
where the likelihood is
\texttt{bernoulli\_logit\_lpmf(y\ \textbar{}\ theta{[}subj\_id{]})}.

However, from Chapter 9 onward, all models are \textbf{path-dependent}:
the predicted probability on trial \(t\) is a function of everything
that happened on trials \(1, \ldots, t-1\) within the same subject. For
those models, we need to process each subject's trials in strict
sequence, which means we must know exactly where each subject's block of
rows starts and ends in the data array.

The \textbf{contiguous-block} pattern achieves this:

\begin{Shaded}
\begin{Highlighting}[]
\CommentTok{\# Sort data by subject, then by trial within subject}
\NormalTok{data\_sorted }\OtherTok{\textless{}{-}}\NormalTok{ data }\SpecialCharTok{|\textgreater{}} \FunctionTok{arrange}\NormalTok{(subject\_id, trial)}
\NormalTok{bounds }\OtherTok{\textless{}{-}}\NormalTok{ data\_sorted }\SpecialCharTok{|\textgreater{}}
  \FunctionTok{mutate}\NormalTok{(}\AttributeTok{row =} \FunctionTok{row\_number}\NormalTok{()) }\SpecialCharTok{|\textgreater{}}
  \FunctionTok{group\_by}\NormalTok{(subject\_id) }\SpecialCharTok{|\textgreater{}}
  \FunctionTok{summarise}\NormalTok{(}\AttributeTok{subj\_start =} \FunctionTok{min}\NormalTok{(row), }\AttributeTok{subj\_end =} \FunctionTok{max}\NormalTok{(row))}
\end{Highlighting}
\end{Shaded}

\begin{Shaded}
\begin{Highlighting}[]
\KeywordTok{data}\NormalTok{ \{}
  \DataTypeTok{array}\NormalTok{[N\_subjects] }\DataTypeTok{int}\NormalTok{ subj\_start;  }\CommentTok{// first row index for each subject}
  \DataTypeTok{array}\NormalTok{[N\_subjects] }\DataTypeTok{int}\NormalTok{ subj\_end;    }\CommentTok{// last row index for each subject}
\NormalTok{\}}
\KeywordTok{model}\NormalTok{ \{}
  \ControlFlowTok{for}\NormalTok{ (j }\ControlFlowTok{in} \DecValTok{1}\NormalTok{:N\_subjects) \{}
    \ControlFlowTok{for}\NormalTok{ (i }\ControlFlowTok{in}\NormalTok{ subj\_start[j]:subj\_end[j]) \{}
      \CommentTok{// process trials in order — history is correct}
\NormalTok{    \}}
\NormalTok{  \}}
\NormalTok{\}}
\end{Highlighting}
\end{Shaded}

This layout also enables the \textbf{\texttt{reduce\_sum} parallelism}
introduced in Chapter 11: because each subject's trials form a
self-contained contiguous block, independent subject loops can be
dispatched to separate CPU threads without any shared mutable state.

We keep \texttt{subj\_id} here for simplicity. You will see
\texttt{subj\_start/subj\_end} throughout Chapters 9--11.

\end{tcolorbox}

\begin{Shaded}
\begin{Highlighting}[]
\CommentTok{\# Timing, ESS, and Pathfinder warm{-}start comparison for CP vs NCP}

\CommentTok{\# ── 1. Wall{-}clock and ESS ─────────────────────────────────────────────────────}
\NormalTok{time\_cp  }\OtherTok{\textless{}{-}}\NormalTok{ fit\_cp}\SpecialCharTok{$}\FunctionTok{time}\NormalTok{()}\SpecialCharTok{$}\NormalTok{total}
\NormalTok{time\_ncp }\OtherTok{\textless{}{-}}\NormalTok{ fit\_ncp}\SpecialCharTok{$}\FunctionTok{time}\NormalTok{()}\SpecialCharTok{$}\NormalTok{total}

\NormalTok{ess\_cp  }\OtherTok{\textless{}{-}} \FunctionTok{summarise\_draws}\NormalTok{(fit\_cp}\SpecialCharTok{$}\FunctionTok{draws}\NormalTok{(}\FunctionTok{c}\NormalTok{(}\StringTok{"mu\_theta"}\NormalTok{, }\StringTok{"sigma\_theta"}\NormalTok{)),  }\StringTok{"ess\_bulk"}\NormalTok{) }\SpecialCharTok{|\textgreater{}}
  \FunctionTok{rename}\NormalTok{(}\AttributeTok{ess\_cp  =}\NormalTok{ ess\_bulk)}
\NormalTok{ess\_ncp }\OtherTok{\textless{}{-}} \FunctionTok{summarise\_draws}\NormalTok{(fit\_ncp}\SpecialCharTok{$}\FunctionTok{draws}\NormalTok{(}\FunctionTok{c}\NormalTok{(}\StringTok{"mu\_theta"}\NormalTok{, }\StringTok{"sigma\_theta"}\NormalTok{)), }\StringTok{"ess\_bulk"}\NormalTok{) }\SpecialCharTok{|\textgreater{}}
  \FunctionTok{rename}\NormalTok{(}\AttributeTok{ess\_ncp =}\NormalTok{ ess\_bulk)}

\NormalTok{efficiency\_comparison }\OtherTok{\textless{}{-}}\NormalTok{ ess\_cp }\SpecialCharTok{|\textgreater{}}
  \FunctionTok{left\_join}\NormalTok{(ess\_ncp, }\AttributeTok{by =} \StringTok{"variable"}\NormalTok{) }\SpecialCharTok{|\textgreater{}}
  \FunctionTok{mutate}\NormalTok{(}
    \AttributeTok{ess\_per\_sec\_cp  =}\NormalTok{ ess\_cp  }\SpecialCharTok{/}\NormalTok{ time\_cp,}
    \AttributeTok{ess\_per\_sec\_ncp =}\NormalTok{ ess\_ncp }\SpecialCharTok{/}\NormalTok{ time\_ncp}
\NormalTok{  )}

\FunctionTok{cat}\NormalTok{(}\StringTok{"=== COMPUTATIONAL SPEED (WALL{-}CLOCK) ===}\SpecialCharTok{\textbackslash{}n}\StringTok{"}\NormalTok{)}
\end{Highlighting}
\end{Shaded}

\begin{verbatim}
=== COMPUTATIONAL SPEED (WALL-CLOCK) ===
\end{verbatim}

\begin{Shaded}
\begin{Highlighting}[]
\FunctionTok{cat}\NormalTok{(}\StringTok{"CP  Total Time:"}\NormalTok{, }\FunctionTok{round}\NormalTok{(time\_cp,  }\DecValTok{2}\NormalTok{), }\StringTok{"seconds}\SpecialCharTok{\textbackslash{}n}\StringTok{"}\NormalTok{)}
\end{Highlighting}
\end{Shaded}

\begin{verbatim}
CP  Total Time: 4.22 seconds
\end{verbatim}

\begin{Shaded}
\begin{Highlighting}[]
\FunctionTok{cat}\NormalTok{(}\StringTok{"NCP Total Time:"}\NormalTok{, }\FunctionTok{round}\NormalTok{(time\_ncp, }\DecValTok{2}\NormalTok{), }\StringTok{"seconds}\SpecialCharTok{\textbackslash{}n}\StringTok{"}\NormalTok{)}
\end{Highlighting}
\end{Shaded}

\begin{verbatim}
NCP Total Time: 6.23 seconds
\end{verbatim}

\begin{Shaded}
\begin{Highlighting}[]
\FunctionTok{cat}\NormalTok{(}\StringTok{"Speedup Factor:"}\NormalTok{, }\FunctionTok{round}\NormalTok{(time\_cp }\SpecialCharTok{/}\NormalTok{ time\_ncp, }\DecValTok{2}\NormalTok{), }\StringTok{"x}\SpecialCharTok{\textbackslash{}n\textbackslash{}n}\StringTok{"}\NormalTok{)}
\end{Highlighting}
\end{Shaded}

\begin{verbatim}
Speedup Factor: 0.68 x
\end{verbatim}

\begin{Shaded}
\begin{Highlighting}[]
\FunctionTok{cat}\NormalTok{(}\StringTok{"=== SAMPLING EFFICIENCY (ESS / SECOND) ===}\SpecialCharTok{\textbackslash{}n}\StringTok{"}\NormalTok{)}
\end{Highlighting}
\end{Shaded}

\begin{verbatim}
=== SAMPLING EFFICIENCY (ESS / SECOND) ===
\end{verbatim}

\begin{Shaded}
\begin{Highlighting}[]
\FunctionTok{print}\NormalTok{(knitr}\SpecialCharTok{::}\FunctionTok{kable}\NormalTok{(efficiency\_comparison, }\AttributeTok{digits =} \DecValTok{1}\NormalTok{))}
\end{Highlighting}
\end{Shaded}

\begin{verbatim}


|variable    | ess_cp| ess_ncp| ess_per_sec_cp| ess_per_sec_ncp|
|:-----------|------:|-------:|--------------:|---------------:|
|mu_theta    | 4407.2|   603.3|         1044.8|            96.8|
|sigma_theta | 3472.6|   600.7|          823.2|            96.4|
\end{verbatim}

\begin{Shaded}
\begin{Highlighting}[]
\CommentTok{\# ── 2. Pathfinder warm{-}start timing for NCP ───────────────────────────────────}
\NormalTok{fit\_ncp\_pf\_path }\OtherTok{\textless{}{-}}\NormalTok{ here}\SpecialCharTok{::}\FunctionTok{here}\NormalTok{(}\StringTok{"simmodels"}\NormalTok{, }\StringTok{"ch07\_fit\_ncp\_pathfinder\_init.rds"}\NormalTok{)}

\ControlFlowTok{if}\NormalTok{ (regenerate\_simulations }\SpecialCharTok{||} \SpecialCharTok{!}\FunctionTok{file.exists}\NormalTok{(fit\_ncp\_pf\_path)) \{}
\NormalTok{  pf\_ncp }\OtherTok{\textless{}{-}} \FunctionTok{tryCatch}\NormalTok{(}
\NormalTok{    mod\_biased\_ml\_ncp}\SpecialCharTok{$}\FunctionTok{pathfinder}\NormalTok{(}\AttributeTok{data =}\NormalTok{ stan\_data\_centered, }\AttributeTok{num\_paths =} \DecValTok{4}\NormalTok{,}
                                  \AttributeTok{seed =} \DecValTok{42}\NormalTok{, }\AttributeTok{refresh =} \DecValTok{0}\NormalTok{),}
    \AttributeTok{error =} \ControlFlowTok{function}\NormalTok{(e) }\ConstantTok{NULL}
\NormalTok{  )}
\NormalTok{  init\_for\_ncp }\OtherTok{\textless{}{-}} \ControlFlowTok{if}\NormalTok{ (}\SpecialCharTok{!}\FunctionTok{is.null}\NormalTok{(pf\_ncp)) pf\_ncp }\ControlFlowTok{else}\NormalTok{ init\_ncp\_list}

\NormalTok{  t\_pf\_start }\OtherTok{\textless{}{-}} \FunctionTok{proc.time}\NormalTok{()[}\StringTok{"elapsed"}\NormalTok{]}
\NormalTok{  fit\_ncp\_pf }\OtherTok{\textless{}{-}}\NormalTok{ mod\_biased\_ml\_ncp}\SpecialCharTok{$}\FunctionTok{sample}\NormalTok{(}
    \AttributeTok{data            =}\NormalTok{ stan\_data\_centered,}
    \AttributeTok{seed            =} \DecValTok{124}\NormalTok{,}
    \AttributeTok{chains          =} \DecValTok{4}\NormalTok{,}
    \AttributeTok{parallel\_chains =} \DecValTok{4}\NormalTok{,}
    \AttributeTok{iter\_warmup     =} \DecValTok{500}\NormalTok{,        }\CommentTok{\# shorter warmup thanks to Pathfinder init}
    \AttributeTok{iter\_sampling   =} \DecValTok{1000}\NormalTok{,}
    \AttributeTok{init            =}\NormalTok{ init\_for\_ncp,}
    \AttributeTok{adapt\_delta     =} \FloatTok{0.85}\NormalTok{,}
    \AttributeTok{refresh         =} \DecValTok{0}
\NormalTok{  )}
\NormalTok{  time\_pf\_total }\OtherTok{\textless{}{-}}\NormalTok{ (}\FunctionTok{proc.time}\NormalTok{()[}\StringTok{"elapsed"}\NormalTok{] }\SpecialCharTok{{-}}\NormalTok{ t\_pf\_start) }\SpecialCharTok{+}
                   \ControlFlowTok{if}\NormalTok{ (}\SpecialCharTok{!}\FunctionTok{is.null}\NormalTok{(pf\_ncp)) pf\_ncp}\SpecialCharTok{$}\FunctionTok{time}\NormalTok{()}\SpecialCharTok{$}\NormalTok{total }\ControlFlowTok{else} \DecValTok{0}
\NormalTok{  fit\_ncp\_pf}\SpecialCharTok{$}\FunctionTok{save\_object}\NormalTok{(fit\_ncp\_pf\_path)}
\NormalTok{\} }\ControlFlowTok{else}\NormalTok{ \{}
\NormalTok{  fit\_ncp\_pf }\OtherTok{\textless{}{-}} \FunctionTok{readRDS}\NormalTok{(fit\_ncp\_pf\_path)}
\NormalTok{  time\_pf\_total }\OtherTok{\textless{}{-}}\NormalTok{ fit\_ncp\_pf}\SpecialCharTok{$}\FunctionTok{time}\NormalTok{()}\SpecialCharTok{$}\NormalTok{total   }\CommentTok{\# approximate (Pathfinder time not stored)}
\NormalTok{\}}

\NormalTok{timing\_summary }\OtherTok{\textless{}{-}}\NormalTok{ tibble}\SpecialCharTok{::}\FunctionTok{tribble}\NormalTok{(}
  \SpecialCharTok{\textasciitilde{}}\NormalTok{Method,                       }\SpecialCharTok{\textasciitilde{}}\StringTok{\textasciigrave{}}\AttributeTok{Time (s)}\StringTok{\textasciigrave{}}\NormalTok{,  }\SpecialCharTok{\textasciitilde{}}\StringTok{\textasciigrave{}}\AttributeTok{Warmup iter}\StringTok{\textasciigrave{}}\NormalTok{,}
  \StringTok{"CP (standard)"}\NormalTok{,                }\FunctionTok{round}\NormalTok{(time\_cp, }\DecValTok{1}\NormalTok{),    }\DecValTok{1000}\NormalTok{L,}
  \StringTok{"NCP (standard init)"}\NormalTok{,          }\FunctionTok{round}\NormalTok{(time\_ncp, }\DecValTok{1}\NormalTok{),   }\DecValTok{1000}\NormalTok{L,}
  \StringTok{"NCP (Pathfinder warm{-}start)"}\NormalTok{,  }\FunctionTok{round}\NormalTok{(time\_pf\_total, }\DecValTok{1}\NormalTok{), }\DecValTok{500}\NormalTok{L}
\NormalTok{)}
\NormalTok{knitr}\SpecialCharTok{::}\FunctionTok{kable}\NormalTok{(timing\_summary, }\AttributeTok{caption =} \StringTok{"Wall{-}clock time: CP vs NCP, with and without Pathfinder warm{-}start"}\NormalTok{)}
\end{Highlighting}
\end{Shaded}

\begin{longtable}[]{@{}lrr@{}}
\caption{Wall-clock time: CP vs NCP, with and without Pathfinder
warm-start}\tabularnewline
\toprule\noalign{}
Method & Time (s) & Warmup iter \\
\midrule\noalign{}
\endfirsthead
\toprule\noalign{}
Method & Time (s) & Warmup iter \\
\midrule\noalign{}
\endhead
\bottomrule\noalign{}
\endlastfoot
CP (standard) & 4.2 & 1000 \\
NCP (standard init) & 6.2 & 1000 \\
NCP (Pathfinder warm-start) & 5.0 & 500 \\
\end{longtable}

Here we see that the NCP model is sometimes slightly faster, sometimes
slightly slower than CP in raw wall-clock time, but always less
\emph{sampling-efficient} (lower ESS/second). This is the trade-off of
NCP: it eliminates the pathological geometry that causes SBC failures in
CP, at the cost of slightly less efficient sampling in data-rich
regimes. The Pathfinder warm-start partially recovers that efficiency by
cutting warmup iterations.

Yet, NCP might still end up exploring the space better. Let's see how it
fares on the SBC assessment that was giving issues to CP.

We will generate 200 random universes from the priors, fit the NCP model
to each, and plot the rank histograms.

\begin{Shaded}
\begin{Highlighting}[]
\NormalTok{sbc\_ncp\_path }\OtherTok{\textless{}{-}}\NormalTok{ here}\SpecialCharTok{::}\FunctionTok{here}\NormalTok{(}\StringTok{"simmodels"}\NormalTok{, }\StringTok{"ch07\_sbc\_ncp.rds"}\NormalTok{)}

\NormalTok{backend\_ncp }\OtherTok{\textless{}{-}} \FunctionTok{SBC\_backend\_cmdstan\_sample}\NormalTok{(}
\NormalTok{  mod\_biased\_ml\_ncp,}
  \AttributeTok{iter\_warmup =} \DecValTok{1000}\NormalTok{, }\AttributeTok{iter\_sampling =} \DecValTok{1000}\NormalTok{, }\AttributeTok{chains =} \DecValTok{1}\NormalTok{, }\AttributeTok{refresh =} \DecValTok{0}
\NormalTok{)}

\ControlFlowTok{if}\NormalTok{ (regenerate\_simulations }\SpecialCharTok{||} \SpecialCharTok{!}\FunctionTok{file.exists}\NormalTok{(sbc\_ncp\_path)) \{}
\NormalTok{  future}\SpecialCharTok{::}\FunctionTok{plan}\NormalTok{(future}\SpecialCharTok{::}\NormalTok{multisession, }\AttributeTok{workers =} \DecValTok{4}\NormalTok{)}
  
  \CommentTok{\# Add keep\_fits = FALSE right here}
\NormalTok{  sbc\_results\_ncp }\OtherTok{\textless{}{-}} \FunctionTok{compute\_SBC}\NormalTok{(}
    \AttributeTok{datasets =}\NormalTok{ sbc\_datasets\_cp, }
    \AttributeTok{backend =}\NormalTok{ backend\_ncp, }
    \AttributeTok{keep\_fits =} \ConstantTok{FALSE}
\NormalTok{  )}
  
  \FunctionTok{saveRDS}\NormalTok{(sbc\_results\_ncp, sbc\_ncp\_path)}
\NormalTok{\} }\ControlFlowTok{else}\NormalTok{ \{}
\NormalTok{  sbc\_results\_ncp }\OtherTok{\textless{}{-}} \FunctionTok{readRDS}\NormalTok{(sbc\_ncp\_path)}
\NormalTok{\}}

\FunctionTok{plot\_rank\_hist}\NormalTok{(sbc\_results\_ncp, }\AttributeTok{variables =} \FunctionTok{c}\NormalTok{(}\StringTok{"mu\_theta"}\NormalTok{, }\StringTok{"sigma\_theta"}\NormalTok{)) }\SpecialCharTok{+}
  \FunctionTok{labs}\NormalTok{(}\AttributeTok{title =} \StringTok{"SBC Rank Histograms: Non{-}Centered Parameterization"}\NormalTok{)}
\end{Highlighting}
\end{Shaded}

\begin{center}
\includegraphics[width=0.8\linewidth,height=\textheight,keepaspectratio]{06-MultilevelModels_files/figure-pdf/ch07_sbc_ncp_execution-1.pdf}
\end{center}

\begin{Shaded}
\begin{Highlighting}[]
\NormalTok{sbc\_stats\_df }\OtherTok{\textless{}{-}}\NormalTok{ sbc\_results\_cp}\SpecialCharTok{$}\NormalTok{stats}

\CommentTok{\# 2. Filter for the population{-}level parameters you care about}
\CommentTok{\# (e.g., mu\_alpha, mu\_beta, rho, or even sigma[1])}
\NormalTok{recovery\_df }\OtherTok{\textless{}{-}}\NormalTok{ sbc\_stats\_df }\SpecialCharTok{|\textgreater{}}
  \FunctionTok{filter}\NormalTok{(variable }\SpecialCharTok{\%in\%} \FunctionTok{c}\NormalTok{(}\StringTok{"mu\_theta"}\NormalTok{, }\StringTok{"sigma\_theta"}\NormalTok{))}

\CommentTok{\# 3. Create the Parameter Recovery Identity Plot}
\NormalTok{p\_sbc\_recovery }\OtherTok{\textless{}{-}} \FunctionTok{ggplot}\NormalTok{(recovery\_df, }\FunctionTok{aes}\NormalTok{(}\AttributeTok{x =}\NormalTok{ simulated\_value, }\AttributeTok{y =}\NormalTok{ median)) }\SpecialCharTok{+}
  \CommentTok{\# The Identity Line (Perfect Recovery)}
  \FunctionTok{geom\_abline}\NormalTok{(}\AttributeTok{intercept =} \DecValTok{0}\NormalTok{, }\AttributeTok{slope =} \DecValTok{1}\NormalTok{, }\AttributeTok{linetype =} \StringTok{"dashed"}\NormalTok{, }\AttributeTok{color =} \StringTok{"gray50"}\NormalTok{, }\AttributeTok{linewidth =} \DecValTok{1}\NormalTok{) }\SpecialCharTok{+}
  
  \CommentTok{\# The Posterior Estimates with 90\% Credible Intervals}
  \FunctionTok{geom\_errorbar}\NormalTok{(}\FunctionTok{aes}\NormalTok{(}\AttributeTok{ymin =}\NormalTok{ q5, }\AttributeTok{ymax =}\NormalTok{ q95), }\AttributeTok{width =} \DecValTok{0}\NormalTok{, }\AttributeTok{alpha =} \FloatTok{0.3}\NormalTok{, }\AttributeTok{color =} \StringTok{"steelblue"}\NormalTok{) }\SpecialCharTok{+}
  \FunctionTok{geom\_point}\NormalTok{(}\AttributeTok{color =} \StringTok{"midnightblue"}\NormalTok{, }\AttributeTok{alpha =} \FloatTok{0.7}\NormalTok{, }\AttributeTok{size =} \DecValTok{2}\NormalTok{) }\SpecialCharTok{+}
  
  \CommentTok{\# Facet by parameter}
  \FunctionTok{facet\_wrap}\NormalTok{(}\SpecialCharTok{\textasciitilde{}}\NormalTok{variable, }\AttributeTok{scales =} \StringTok{"free"}\NormalTok{) }\SpecialCharTok{+}
  
  \FunctionTok{labs}\NormalTok{(}
    \AttributeTok{title =} \StringTok{"Parameter Recovery (Extracted Directly from SBC)"}\NormalTok{,}
    \AttributeTok{subtitle =} \StringTok{"Points represent the posterior median. Bars are 90\% Credible Intervals."}\NormalTok{,}
    \AttributeTok{x =} \StringTok{"True Generative Value (Drawn from Prior)"}\NormalTok{,}
    \AttributeTok{y =} \StringTok{"Posterior Median Estimate"}
\NormalTok{  ) }\SpecialCharTok{+}
  \FunctionTok{theme\_classic}\NormalTok{() }\SpecialCharTok{+}
  \FunctionTok{theme}\NormalTok{(}\AttributeTok{strip.background =} \FunctionTok{element\_rect}\NormalTok{(}\AttributeTok{fill =} \StringTok{"gray95"}\NormalTok{, }\AttributeTok{color =} \ConstantTok{NA}\NormalTok{))}

\FunctionTok{print}\NormalTok{(p\_sbc\_recovery)}
\end{Highlighting}
\end{Shaded}

\begin{center}
\includegraphics[width=0.8\linewidth,height=\textheight,keepaspectratio]{06-MultilevelModels_files/figure-pdf/ch07_sbc_ncp_execution-2.pdf}
\end{center}

\begin{Shaded}
\begin{Highlighting}[]
\CommentTok{\# 4. Quantify the Recovery (Correlation and RMSE)}
\NormalTok{recovery\_metrics }\OtherTok{\textless{}{-}}\NormalTok{ recovery\_df }\SpecialCharTok{|\textgreater{}}
  \FunctionTok{group\_by}\NormalTok{(variable) }\SpecialCharTok{|\textgreater{}}
  \FunctionTok{summarize}\NormalTok{(}
    \AttributeTok{Pearson\_r =} \FunctionTok{cor}\NormalTok{(simulated\_value, median),}
    \AttributeTok{Spearman\_rho =} \FunctionTok{cor}\NormalTok{(simulated\_value, median, }\AttributeTok{method =} \StringTok{"spearman"}\NormalTok{),}
    \AttributeTok{RMSE =} \FunctionTok{sqrt}\NormalTok{(}\FunctionTok{mean}\NormalTok{((simulated\_value }\SpecialCharTok{{-}}\NormalTok{ median)}\SpecialCharTok{\^{}}\DecValTok{2}\NormalTok{)),}
    \AttributeTok{.groups =} \StringTok{"drop"}
\NormalTok{  )}

\FunctionTok{cat}\NormalTok{(}\StringTok{"=== PARAMETER RECOVERY QUANTIFICATION ===}\SpecialCharTok{\textbackslash{}n}\StringTok{"}\NormalTok{)}
\end{Highlighting}
\end{Shaded}

\begin{verbatim}
=== PARAMETER RECOVERY QUANTIFICATION ===
\end{verbatim}

\begin{Shaded}
\begin{Highlighting}[]
\FunctionTok{print}\NormalTok{(knitr}\SpecialCharTok{::}\FunctionTok{kable}\NormalTok{(recovery\_metrics, }\AttributeTok{digits =} \DecValTok{3}\NormalTok{))}
\end{Highlighting}
\end{Shaded}

\begin{verbatim}


|variable    | Pearson_r| Spearman_rho|  RMSE|
|:-----------|---------:|------------:|-----:|
|mu_theta    |     0.990|        0.989| 0.212|
|sigma_theta |     0.986|        0.982| 0.191|
\end{verbatim}

\begin{Shaded}
\begin{Highlighting}[]
\CommentTok{\# Extract diagnostics now so they are available for the targeted refit below}
\NormalTok{diags\_ncp }\OtherTok{\textless{}{-}}\NormalTok{ sbc\_results\_ncp}\SpecialCharTok{$}\NormalTok{backend\_diagnostics}
\NormalTok{ncp\_fail\_rate }\OtherTok{\textless{}{-}} \FunctionTok{mean}\NormalTok{(diags\_ncp}\SpecialCharTok{$}\NormalTok{n\_divergent }\SpecialCharTok{\textgreater{}} \DecValTok{0}\NormalTok{) }\SpecialCharTok{*} \DecValTok{100}

\CommentTok{\# Divergence count distribution across all 200 universes}
\NormalTok{diags\_ncp }\SpecialCharTok{|\textgreater{}} 
  \FunctionTok{ggplot}\NormalTok{(}\FunctionTok{aes}\NormalTok{(}\AttributeTok{x =}\NormalTok{ n\_divergent)) }\SpecialCharTok{+}
  \FunctionTok{geom\_histogram}\NormalTok{(}\AttributeTok{binwidth =} \DecValTok{1}\NormalTok{, }\AttributeTok{fill =} \StringTok{"darkred"}\NormalTok{, }\AttributeTok{alpha =} \FloatTok{0.8}\NormalTok{) }\SpecialCharTok{+}
  \FunctionTok{labs}\NormalTok{(}
    \AttributeTok{title =} \StringTok{"NCP Divergence Counts Across 200 SBC Universes"}\NormalTok{,}
    \AttributeTok{subtitle =} \FunctionTok{sprintf}\NormalTok{(}\StringTok{"\%.0f\%\% of universes produced at least one divergence. }
\StringTok{Even a single divergence invalidates the geometric guarantee of NUTS."}\NormalTok{, cp\_fail\_rate),}
    \AttributeTok{x =} \StringTok{"Number of Divergent Transitions"}\NormalTok{,}
    \AttributeTok{y =} \StringTok{"Number of SBC Universes"}
\NormalTok{  )}
\end{Highlighting}
\end{Shaded}

\begin{center}
\includegraphics[width=0.8\linewidth,height=\textheight,keepaspectratio]{06-MultilevelModels_files/figure-pdf/ch07_sbc_ncp_execution-3.pdf}
\end{center}

\begin{Shaded}
\begin{Highlighting}[]
\CommentTok{\# Locate the divergence{-}prone region: do they concentrate at small sigma\_theta?}
\NormalTok{divergence\_profile }\OtherTok{\textless{}{-}}\NormalTok{ posterior}\SpecialCharTok{::}\FunctionTok{as\_draws\_df}\NormalTok{(sbc\_datasets\_cp}\SpecialCharTok{$}\NormalTok{variables) }\SpecialCharTok{|\textgreater{}} 
  \FunctionTok{as\_tibble}\NormalTok{() }\SpecialCharTok{|\textgreater{}} 
  \FunctionTok{mutate}\NormalTok{(}\AttributeTok{sim\_id =} \FunctionTok{row\_number}\NormalTok{()) }\SpecialCharTok{|\textgreater{}} 
  \FunctionTok{left\_join}\NormalTok{(}
\NormalTok{    diags\_ncp }\SpecialCharTok{|\textgreater{}} \FunctionTok{mutate}\NormalTok{(}\AttributeTok{sim\_id =} \FunctionTok{row\_number}\NormalTok{()) }\SpecialCharTok{|\textgreater{}}\NormalTok{ dplyr}\SpecialCharTok{::}\FunctionTok{select}\NormalTok{(sim\_id, n\_divergent),}
    \AttributeTok{by =} \StringTok{"sim\_id"}
\NormalTok{  ) }\SpecialCharTok{|\textgreater{}} 
  \FunctionTok{mutate}\NormalTok{(}\AttributeTok{had\_divergence =}\NormalTok{ n\_divergent }\SpecialCharTok{\textgreater{}} \DecValTok{0}\NormalTok{)}

\FunctionTok{ggplot}\NormalTok{(divergence\_profile, }\FunctionTok{aes}\NormalTok{(}\AttributeTok{x =}\NormalTok{ sigma\_theta, }\AttributeTok{y =}\NormalTok{ n\_divergent, }\AttributeTok{color =}\NormalTok{ had\_divergence)) }\SpecialCharTok{+}
  \FunctionTok{geom\_point}\NormalTok{(}\AttributeTok{alpha =} \FloatTok{0.5}\NormalTok{, }\AttributeTok{size =} \FloatTok{1.5}\NormalTok{) }\SpecialCharTok{+}
  \FunctionTok{scale\_color\_manual}\NormalTok{(}
    \AttributeTok{values =} \FunctionTok{c}\NormalTok{(}\StringTok{"FALSE"} \OtherTok{=} \StringTok{"steelblue"}\NormalTok{, }\StringTok{"TRUE"} \OtherTok{=} \StringTok{"darkred"}\NormalTok{),}
    \AttributeTok{labels =} \FunctionTok{c}\NormalTok{(}\StringTok{"FALSE"} \OtherTok{=} \StringTok{"Clean"}\NormalTok{, }\StringTok{"TRUE"} \OtherTok{=} \StringTok{"Divergent"}\NormalTok{),}
    \AttributeTok{name =} \ConstantTok{NULL}
\NormalTok{  ) }\SpecialCharTok{+}
  \FunctionTok{labs}\NormalTok{(}
    \AttributeTok{title =} \StringTok{"NCP Divergences Concentrate at Small True sigma\_theta"}\NormalTok{,}
    \AttributeTok{subtitle =} \StringTok{"The funnel geometry only manifests when population variance is small."}\NormalTok{,}
    \AttributeTok{x =} \FunctionTok{expression}\NormalTok{(}\StringTok{"True "} \SpecialCharTok{*}\NormalTok{ sigma[theta]),}
    \AttributeTok{y =} \StringTok{"Number of Divergent Transitions"}
\NormalTok{  )}
\end{Highlighting}
\end{Shaded}

\begin{center}
\includegraphics[width=0.8\linewidth,height=\textheight,keepaspectratio]{06-MultilevelModels_files/figure-pdf/ch07_sbc_ncp_execution-4.pdf}
\end{center}

\begin{Shaded}
\begin{Highlighting}[]
\FunctionTok{ggplot}\NormalTok{(divergence\_profile, }\FunctionTok{aes}\NormalTok{(}\AttributeTok{x =}\NormalTok{ mu\_theta, }\AttributeTok{y =}\NormalTok{ sigma\_theta, }\AttributeTok{color =}\NormalTok{ had\_divergence)) }\SpecialCharTok{+}
  \FunctionTok{geom\_point}\NormalTok{(}\AttributeTok{alpha =} \FloatTok{0.7}\NormalTok{, }\AttributeTok{size =} \DecValTok{2}\NormalTok{) }\SpecialCharTok{+}
  \FunctionTok{scale\_color\_manual}\NormalTok{(}
    \AttributeTok{name =} \StringTok{"Sampler Status"}\NormalTok{,}
    \AttributeTok{values =} \FunctionTok{c}\NormalTok{(}\StringTok{"FALSE"} \OtherTok{=} \StringTok{"steelblue"}\NormalTok{, }\StringTok{"TRUE"} \OtherTok{=} \StringTok{"darkred"}\NormalTok{),}
    \AttributeTok{labels =} \FunctionTok{c}\NormalTok{(}\StringTok{"Healthy"}\NormalTok{, }\StringTok{"Diverged"}\NormalTok{)}
\NormalTok{  ) }\SpecialCharTok{+}
  \FunctionTok{labs}\NormalTok{(}
    \AttributeTok{title =} \StringTok{"The Anatomy of Failure: Where does NCP break?"}\NormalTok{,}
    \AttributeTok{x =} \StringTok{"True Population Mean (mu\_theta)"}\NormalTok{,}
    \AttributeTok{y =} \StringTok{"True Population Variance (sigma\_theta)"}
\NormalTok{  ) }\SpecialCharTok{+}
  \FunctionTok{theme\_classic}\NormalTok{()}
\end{Highlighting}
\end{Shaded}

\begin{center}
\includegraphics[width=0.8\linewidth,height=\textheight,keepaspectratio]{06-MultilevelModels_files/figure-pdf/ch07_sbc_ncp_execution-5.pdf}
\end{center}

\subsubsection{Posterior Simulation-Based Calibration Checks: Probing
Local
Geometry}\label{posterior-simulation-based-calibration-checks-probing-local-geometry}

Prior SBC tested our sampler across the entire prior space --- it asked
whether NUTS is unbiased when the true parameters could be anything the
prior assigns non-negligible mass to. This is a global guarantee, and it
is necessary. But it is not sufficient. A sampler can be globally
unbiased yet struggle precisely in the regions that matter most: the
posterior neighborhoods that the actual observed data force us into.
When the likelihood is highly informative --- as it typically is once we
have a few hundred trials per agent --- the posterior concentrates
sharply in a small corner of the prior space. The sampler must now take
large numbers of tiny, precise steps in a geometrically restricted
region. Prior SBC never directly stresses this regime, because most of
its simulated universes place the true parameter somewhere across the
full prior, and only a fraction happen to fall in the data-rich
posterior neighborhood.

Posterior SBC (PSBC) is a complementary local diagnostic that directly
targets this regime. The logic inverts the generative direction. Instead
of drawing true parameters from the prior, we draw them from the
\emph{fitted posterior} --- a distribution that has already been
conditioned on our actual data and therefore concentrates in precisely
the region the sampler will need to navigate in practice. Concretely: we
sample a single draw \(\tilde{\theta}^{(s)}\) from the posterior
(\(\theta \mid y_{\text{obs}}\)), treat it as a ground truth, generate a
synthetic dataset
\(\tilde{y}^{(s)} \sim p(y \mid \tilde{\theta}^{(s)})\) from the data
model, fit the model to \(\tilde{y}^{(s)}\), and record where
\(\tilde{\theta}^{(s)}\) falls in the resulting posterior. If the
sampler is locally well-calibrated, these ranks are again uniform ---
but now over the concentrated geometry of the posterior, not the diffuse
geometry of the prior. A U-shaped rank histogram here means the sampler
is overconfident \emph{specifically when the data are informative},
which is exactly the failure mode we need to detect. Implementing this
requires us to write a \emph{factory function}: a function that closes
over a fitted model object and a data structure, and returns a new
generator function that the \texttt{SBC} package can call repeatedly to
produce fresh \(({\tilde{\theta}}, \tilde{y})\) pairs. If this pattern
is unfamiliar, read it as a function that \emph{builds} generators
rather than \emph{being} one ---
\texttt{make\_psbc\_generator(fit\_cp,\ stan\_data\_centered)} does not
itself generate a dataset; it returns a self-contained callable that
will do so every time the SBC machinery invokes it.

\begin{Shaded}
\begin{Highlighting}[]
\CommentTok{\# 1. Define a robust factory function for the PSBC generator}
\NormalTok{make\_psbc\_generator }\OtherTok{\textless{}{-}} \ControlFlowTok{function}\NormalTok{(fit, original\_data) \{}
  \CommentTok{\# Extract the full posterior array as a flat data frame}
\NormalTok{  draws\_df }\OtherTok{\textless{}{-}}\NormalTok{ posterior}\SpecialCharTok{::}\FunctionTok{as\_draws\_df}\NormalTok{(fit}\SpecialCharTok{$}\FunctionTok{draws}\NormalTok{())}
  
  \CommentTok{\# Pre{-}compute the exact column names to guarantee flawless extraction}
\NormalTok{  theta\_cols }\OtherTok{\textless{}{-}} \FunctionTok{paste0}\NormalTok{(}\StringTok{"theta\_logit["}\NormalTok{, }\DecValTok{1}\SpecialCharTok{:}\NormalTok{original\_data}\SpecialCharTok{$}\NormalTok{N\_subjects, }\StringTok{"]"}\NormalTok{)}
\NormalTok{  y\_cols }\OtherTok{\textless{}{-}} \FunctionTok{paste0}\NormalTok{(}\StringTok{"y\_post\_rep["}\NormalTok{, }\DecValTok{1}\SpecialCharTok{:}\NormalTok{original\_data}\SpecialCharTok{$}\NormalTok{N\_total, }\StringTok{"]"}\NormalTok{)}
  
  \CommentTok{\# The generator function required by the SBC package}
  \ControlFlowTok{function}\NormalTok{() \{}
    \CommentTok{\# Step A: Randomly sample one specific "state of the world" from the posterior}
\NormalTok{    idx }\OtherTok{\textless{}{-}} \FunctionTok{sample}\NormalTok{(}\FunctionTok{seq\_len}\NormalTok{(}\FunctionTok{nrow}\NormalTok{(draws\_df)), }\DecValTok{1}\NormalTok{)}
    
    \CommentTok{\# Step B: Package the true parameters for this specific simulation}
\NormalTok{    variables }\OtherTok{\textless{}{-}} \FunctionTok{list}\NormalTok{(}
      \AttributeTok{mu\_theta =} \FunctionTok{as.numeric}\NormalTok{(draws\_df}\SpecialCharTok{$}\NormalTok{mu\_theta[idx]),}
      \AttributeTok{sigma\_theta =} \FunctionTok{as.numeric}\NormalTok{(draws\_df}\SpecialCharTok{$}\NormalTok{sigma\_theta[idx]),}
      \AttributeTok{theta\_logit =} \FunctionTok{as.numeric}\NormalTok{(draws\_df[idx, theta\_cols])}
\NormalTok{    )}
    
    \CommentTok{\# Step C: Package the dataset generated by this specific state of the world}
\NormalTok{    generated }\OtherTok{\textless{}{-}} \FunctionTok{list}\NormalTok{(}
      \AttributeTok{N\_total                  =}\NormalTok{ original\_data}\SpecialCharTok{$}\NormalTok{N\_total,}
      \AttributeTok{N\_subjects               =}\NormalTok{ original\_data}\SpecialCharTok{$}\NormalTok{N\_subjects,}
      \AttributeTok{subj\_id                  =}\NormalTok{ original\_data}\SpecialCharTok{$}\NormalTok{subj\_id,}
      \AttributeTok{y                        =} \FunctionTok{as.numeric}\NormalTok{(draws\_df[idx, y\_cols]),}
      \AttributeTok{prior\_mu\_theta\_mu        =}\NormalTok{ original\_data}\SpecialCharTok{$}\NormalTok{prior\_mu\_theta\_mu,}
      \AttributeTok{prior\_mu\_theta\_sigma     =}\NormalTok{ original\_data}\SpecialCharTok{$}\NormalTok{prior\_mu\_theta\_sigma,}
      \AttributeTok{prior\_sigma\_theta\_lambda =}\NormalTok{ original\_data}\SpecialCharTok{$}\NormalTok{prior\_sigma\_theta\_lambda,}
      \AttributeTok{run\_diagnostics          =}\NormalTok{ original\_data}\SpecialCharTok{$}\NormalTok{run\_diagnostics}
\NormalTok{    )}
    
    \FunctionTok{list}\NormalTok{(}\AttributeTok{variables =}\NormalTok{ variables, }\AttributeTok{generated =}\NormalTok{ generated)}
\NormalTok{  \}}
\NormalTok{\}}

\NormalTok{psbc\_path }\OtherTok{\textless{}{-}}\NormalTok{ here}\SpecialCharTok{::}\FunctionTok{here}\NormalTok{(}\StringTok{"simmodels"}\NormalTok{, }\StringTok{"ch07\_psbc.rds"}\NormalTok{)}

\ControlFlowTok{if}\NormalTok{ (regenerate\_simulations }\SpecialCharTok{||} \SpecialCharTok{!}\FunctionTok{file.exists}\NormalTok{(psbc\_path)) \{}
  
\NormalTok{  gen\_psbc\_cp  }\OtherTok{\textless{}{-}} \FunctionTok{SBC\_generator\_function}\NormalTok{(}\FunctionTok{make\_psbc\_generator}\NormalTok{(fit\_cp,  stan\_data\_centered))}
\NormalTok{  gen\_psbc\_ncp }\OtherTok{\textless{}{-}} \FunctionTok{SBC\_generator\_function}\NormalTok{(}\FunctionTok{make\_psbc\_generator}\NormalTok{(fit\_ncp, stan\_data\_centered))}
  
  \FunctionTok{set.seed}\NormalTok{(}\DecValTok{8675309}\NormalTok{)}
\NormalTok{  psbc\_data\_cp  }\OtherTok{\textless{}{-}} \FunctionTok{generate\_datasets}\NormalTok{(gen\_psbc\_cp,  }\DecValTok{100}\NormalTok{)}
\NormalTok{  psbc\_data\_ncp }\OtherTok{\textless{}{-}} \FunctionTok{generate\_datasets}\NormalTok{(gen\_psbc\_ncp, }\DecValTok{100}\NormalTok{)}
  
\NormalTok{  backend\_cp\_psbc  }\OtherTok{\textless{}{-}} \FunctionTok{SBC\_backend\_cmdstan\_sample}\NormalTok{(mod\_biased\_ml\_cp,  }\AttributeTok{iter\_warmup =} \DecValTok{1000}\NormalTok{, }\AttributeTok{iter\_sampling =} \DecValTok{1000}\NormalTok{, }\AttributeTok{chains =} \DecValTok{1}\NormalTok{, }\AttributeTok{refresh =} \DecValTok{0}\NormalTok{)}
\NormalTok{  backend\_ncp\_psbc }\OtherTok{\textless{}{-}} \FunctionTok{SBC\_backend\_cmdstan\_sample}\NormalTok{(mod\_biased\_ml\_ncp, }\AttributeTok{iter\_warmup =} \DecValTok{1000}\NormalTok{, }\AttributeTok{iter\_sampling =} \DecValTok{1000}\NormalTok{, }\AttributeTok{chains =} \DecValTok{1}\NormalTok{, }\AttributeTok{refresh =} \DecValTok{0}\NormalTok{)}
\NormalTok{  future}\SpecialCharTok{::}\FunctionTok{plan}\NormalTok{(future}\SpecialCharTok{::}\NormalTok{multisession, }\AttributeTok{workers =} \DecValTok{4}\NormalTok{)}
\NormalTok{  psbc\_res\_cp  }\OtherTok{\textless{}{-}} \FunctionTok{compute\_SBC}\NormalTok{(psbc\_data\_cp,  backend\_cp\_psbc)}
\NormalTok{  psbc\_res\_ncp }\OtherTok{\textless{}{-}} \FunctionTok{compute\_SBC}\NormalTok{(psbc\_data\_ncp, backend\_ncp\_psbc)}
  
  \FunctionTok{saveRDS}\NormalTok{(}
    \FunctionTok{list}\NormalTok{(}
      \AttributeTok{res\_cp  =}\NormalTok{ psbc\_res\_cp,}
      \AttributeTok{res\_ncp =}\NormalTok{ psbc\_res\_ncp}
\NormalTok{    ),}
\NormalTok{    psbc\_path}
\NormalTok{  )}
  
\NormalTok{\} }\ControlFlowTok{else}\NormalTok{ \{}
\NormalTok{  psbc\_obj    }\OtherTok{\textless{}{-}} \FunctionTok{readRDS}\NormalTok{(psbc\_path)}
\NormalTok{  psbc\_res\_cp  }\OtherTok{\textless{}{-}}\NormalTok{ psbc\_obj}\SpecialCharTok{$}\NormalTok{res\_cp}
\NormalTok{  psbc\_res\_ncp }\OtherTok{\textless{}{-}}\NormalTok{ psbc\_obj}\SpecialCharTok{$}\NormalTok{res\_ncp}
\NormalTok{\}}

\FunctionTok{plot\_rank\_hist}\NormalTok{(psbc\_res\_cp,  }\AttributeTok{variables =} \FunctionTok{c}\NormalTok{(}\StringTok{"mu\_theta"}\NormalTok{, }\StringTok{"sigma\_theta"}\NormalTok{)) }\SpecialCharTok{+}
  \FunctionTok{labs}\NormalTok{(}\AttributeTok{title =} \StringTok{"PSBC Rank Histograms: Centered Parameterization"}\NormalTok{)}
\end{Highlighting}
\end{Shaded}

\begin{center}
\includegraphics[width=0.8\linewidth,height=\textheight,keepaspectratio]{06-MultilevelModels_files/figure-pdf/ch07_psbc_execution-1.pdf}
\end{center}

\begin{Shaded}
\begin{Highlighting}[]
\FunctionTok{plot\_rank\_hist}\NormalTok{(psbc\_res\_ncp, }\AttributeTok{variables =} \FunctionTok{c}\NormalTok{(}\StringTok{"mu\_theta"}\NormalTok{, }\StringTok{"sigma\_theta"}\NormalTok{)) }\SpecialCharTok{+}
  \FunctionTok{labs}\NormalTok{(}\AttributeTok{title =} \StringTok{"PSBC Rank Histograms: Non{-}Centered Parameterization"}\NormalTok{)}
\end{Highlighting}
\end{Shaded}

\begin{center}
\includegraphics[width=0.8\linewidth,height=\textheight,keepaspectratio]{06-MultilevelModels_files/figure-pdf/ch07_psbc_execution-2.pdf}
\end{center}

\begin{Shaded}
\begin{Highlighting}[]
\CommentTok{\# 1. Extract the lightweight stats dataframe from the SBC object}
\NormalTok{sbc\_stats\_df }\OtherTok{\textless{}{-}}\NormalTok{ psbc\_res\_ncp}\SpecialCharTok{$}\NormalTok{stats}

\CommentTok{\# 2. Filter for the population{-}level parameters you care about}
\CommentTok{\# (e.g., mu\_alpha, mu\_beta, rho, or even sigma[1])}
\NormalTok{recovery\_df }\OtherTok{\textless{}{-}}\NormalTok{ sbc\_stats\_df }\SpecialCharTok{|\textgreater{}}
  \FunctionTok{filter}\NormalTok{(variable }\SpecialCharTok{\%in\%} \FunctionTok{c}\NormalTok{(}\StringTok{"mu\_theta"}\NormalTok{, }\StringTok{"sigma\_theta"}\NormalTok{))}

\CommentTok{\# 3. Create the Parameter Recovery Identity Plot}
\NormalTok{p\_sbc\_recovery }\OtherTok{\textless{}{-}} \FunctionTok{ggplot}\NormalTok{(recovery\_df, }\FunctionTok{aes}\NormalTok{(}\AttributeTok{x =}\NormalTok{ simulated\_value, }\AttributeTok{y =}\NormalTok{ median)) }\SpecialCharTok{+}
  \CommentTok{\# The Identity Line (Perfect Recovery)}
  \FunctionTok{geom\_abline}\NormalTok{(}\AttributeTok{intercept =} \DecValTok{0}\NormalTok{, }\AttributeTok{slope =} \DecValTok{1}\NormalTok{, }\AttributeTok{linetype =} \StringTok{"dashed"}\NormalTok{, }\AttributeTok{color =} \StringTok{"gray50"}\NormalTok{, }\AttributeTok{linewidth =} \DecValTok{1}\NormalTok{) }\SpecialCharTok{+}
  
  \CommentTok{\# The Posterior Estimates with 90\% Credible Intervals}
  \FunctionTok{geom\_errorbar}\NormalTok{(}\FunctionTok{aes}\NormalTok{(}\AttributeTok{ymin =}\NormalTok{ q5, }\AttributeTok{ymax =}\NormalTok{ q95), }\AttributeTok{width =} \DecValTok{0}\NormalTok{, }\AttributeTok{alpha =} \FloatTok{0.3}\NormalTok{, }\AttributeTok{color =} \StringTok{"steelblue"}\NormalTok{) }\SpecialCharTok{+}
  \FunctionTok{geom\_point}\NormalTok{(}\AttributeTok{color =} \StringTok{"midnightblue"}\NormalTok{, }\AttributeTok{alpha =} \FloatTok{0.7}\NormalTok{, }\AttributeTok{size =} \DecValTok{2}\NormalTok{) }\SpecialCharTok{+}
  
  \CommentTok{\# Facet by parameter}
  \FunctionTok{facet\_wrap}\NormalTok{(}\SpecialCharTok{\textasciitilde{}}\NormalTok{variable, }\AttributeTok{scales =} \StringTok{"free"}\NormalTok{) }\SpecialCharTok{+}
  
  \FunctionTok{labs}\NormalTok{(}
    \AttributeTok{title =} \StringTok{"Parameter Recovery (Extracted Directly from SBC)"}\NormalTok{,}
    \AttributeTok{subtitle =} \StringTok{"Points represent the posterior median. Bars are 90\% Credible Intervals."}\NormalTok{,}
    \AttributeTok{x =} \StringTok{"True Generative Value (Drawn from Prior)"}\NormalTok{,}
    \AttributeTok{y =} \StringTok{"Posterior Median Estimate"}
\NormalTok{  ) }\SpecialCharTok{+}
  \FunctionTok{theme\_classic}\NormalTok{() }\SpecialCharTok{+}
  \FunctionTok{theme}\NormalTok{(}\AttributeTok{strip.background =} \FunctionTok{element\_rect}\NormalTok{(}\AttributeTok{fill =} \StringTok{"gray95"}\NormalTok{, }\AttributeTok{color =} \ConstantTok{NA}\NormalTok{))}

\FunctionTok{print}\NormalTok{(p\_sbc\_recovery)}
\end{Highlighting}
\end{Shaded}

\begin{center}
\includegraphics[width=0.8\linewidth,height=\textheight,keepaspectratio]{06-MultilevelModels_files/figure-pdf/ch07_psbc_execution-3.pdf}
\end{center}

\begin{Shaded}
\begin{Highlighting}[]
\CommentTok{\# 4. Quantify the Recovery (Correlation and RMSE)}
\NormalTok{recovery\_metrics }\OtherTok{\textless{}{-}}\NormalTok{ recovery\_df }\SpecialCharTok{|\textgreater{}}
  \FunctionTok{group\_by}\NormalTok{(variable) }\SpecialCharTok{|\textgreater{}}
  \FunctionTok{summarize}\NormalTok{(}
    \AttributeTok{Pearson\_r =} \FunctionTok{cor}\NormalTok{(simulated\_value, median),}
    \AttributeTok{RMSE =} \FunctionTok{sqrt}\NormalTok{(}\FunctionTok{mean}\NormalTok{((simulated\_value }\SpecialCharTok{{-}}\NormalTok{ median)}\SpecialCharTok{\^{}}\DecValTok{2}\NormalTok{)),}
    \AttributeTok{.groups =} \StringTok{"drop"}
\NormalTok{  )}

\FunctionTok{cat}\NormalTok{(}\StringTok{"=== PARAMETER RECOVERY QUANTIFICATION ===}\SpecialCharTok{\textbackslash{}n}\StringTok{"}\NormalTok{)}
\end{Highlighting}
\end{Shaded}

\begin{verbatim}
=== PARAMETER RECOVERY QUANTIFICATION ===
\end{verbatim}

\begin{Shaded}
\begin{Highlighting}[]
\FunctionTok{print}\NormalTok{(knitr}\SpecialCharTok{::}\FunctionTok{kable}\NormalTok{(recovery\_metrics, }\AttributeTok{digits =} \DecValTok{3}\NormalTok{))}
\end{Highlighting}
\end{Shaded}

\begin{verbatim}


|variable    | Pearson_r|  RMSE|
|:-----------|---------:|-----:|
|mu_theta    |     0.257| 0.142|
|sigma_theta |     0.358| 0.107|
\end{verbatim}

\begin{Shaded}
\begin{Highlighting}[]
\NormalTok{diags\_psbc\_cp  }\OtherTok{\textless{}{-}}\NormalTok{ psbc\_res\_cp}\SpecialCharTok{$}\NormalTok{backend\_diagnostics}
\NormalTok{diags\_psbc\_ncp }\OtherTok{\textless{}{-}}\NormalTok{ psbc\_res\_ncp}\SpecialCharTok{$}\NormalTok{backend\_diagnostics}

\NormalTok{psbc\_cp\_fail\_rate  }\OtherTok{\textless{}{-}} \FunctionTok{mean}\NormalTok{(diags\_psbc\_cp}\SpecialCharTok{$}\NormalTok{n\_divergent  }\SpecialCharTok{\textgreater{}} \DecValTok{0}\NormalTok{) }\SpecialCharTok{*} \DecValTok{100}
\NormalTok{psbc\_ncp\_fail\_rate }\OtherTok{\textless{}{-}} \FunctionTok{mean}\NormalTok{(diags\_psbc\_ncp}\SpecialCharTok{$}\NormalTok{n\_divergent }\SpecialCharTok{\textgreater{}} \DecValTok{0}\NormalTok{) }\SpecialCharTok{*} \DecValTok{100}

\FunctionTok{cat}\NormalTok{(}\StringTok{"=== POSTERIOR SBC (LOCAL) DIVERGENCE RATES ===}\SpecialCharTok{\textbackslash{}n}\StringTok{"}\NormalTok{)}
\end{Highlighting}
\end{Shaded}

\begin{verbatim}
=== POSTERIOR SBC (LOCAL) DIVERGENCE RATES ===
\end{verbatim}

\begin{Shaded}
\begin{Highlighting}[]
\FunctionTok{cat}\NormalTok{(}\StringTok{"CP  Failure Rate:"}\NormalTok{, psbc\_cp\_fail\_rate,  }\StringTok{"\%}\SpecialCharTok{\textbackslash{}n}\StringTok{"}\NormalTok{)}
\end{Highlighting}
\end{Shaded}

\begin{verbatim}
CP  Failure Rate: 0 %
\end{verbatim}

\begin{Shaded}
\begin{Highlighting}[]
\FunctionTok{cat}\NormalTok{(}\StringTok{"NCP Failure Rate:"}\NormalTok{, psbc\_ncp\_fail\_rate, }\StringTok{"\%}\SpecialCharTok{\textbackslash{}n}\StringTok{"}\NormalTok{)}
\end{Highlighting}
\end{Shaded}

\begin{verbatim}
NCP Failure Rate: 0 %
\end{verbatim}

\subsection{Comparing CP and NCP: Global Fragility vs.~Local
Efficiency}\label{comparing-cp-and-ncp-global-fragility-vs.-local-efficiency}

The CP did not seem to perform too bad, but it broke down when doing
SBC. We therefore reparameterized the model by decoupling the sampling
of individual and population level parameters (from CP to NCP). This
fundamentally altered how the Hamiltonian Monte Carlo (NUTS) algorithm
explores the posterior manifold.

We saw that NCP took slightly longer to run and was less efficient in
sampling (\textless{} ESS). We then ran extensive SBC both sampling the
simulations from the priors (global calibration) and sampling the
simulations from the fitted posteriors (local calibration).

By comparing global calibration (Prior SBC) against local calibration
(Posterior SBC), we have mathematically mapped the boundary conditions
of these parameterizations.

\begin{itemize}
\item
  Integration Stability and Global Fragility (Divergences): In CP,
  sparse data forces the sampler into the infinite curvature of Neal's
  Funnel. As proven by our Prior SBC failure rates, CP is globally
  fragile; it regularly diverged across the broad generative prior space
  whenever the true population variance (\(\sigma_\theta\)) was small.
  NCP, conversely, decouples the parameters. The sampler explores a
  standard, isotropic Gaussian space (\(z \sim \text{Normal}(0, 1)\))
  where the curvature is constant, completely immunizing the model
  against Neal's Funnel and dropping global divergences to zero.
\item
  Sampling Efficiency and Local Vulnerability (ESS and Speed): Because
  the NCP geometry is smooth and well-behaved in data-sparse regimes,
  the Hamiltonian integrator can take massive, efficient step sizes,
  causing the Effective Sample Size (ESS) to skyrocket. However, as our
  Posterior SBC proved, NCP is locally vulnerable to the Data Funnel (or
  ``Reverse Funnel''). In the specific neighborhood of our simulated
  data, the likelihood was highly informative and tightly constrained
  the posterior. CP navigated this sharp likelihood peak effortlessly.
  NCP, however, was forced to blindly explore a broad standard normal
  prior, only to find that the data violently rejected almost all of it,
  causing its ESS to plummet.
\item
  Stationarity (\(\hat{R}\)): While both models can achieve
  \(\hat{R} < 1.01\) under ideal local conditions, NCP achieves it
  reliably across the entire prior space, ensuring all Markov chains mix
  seamlessly regardless of the underlying data-generating universe.
\item
  The Geometric Trade-off: When CP Wins Given the global safety of NCP,
  it is tempting to conclude that we should always use the Non-Centered
  Parameterization. This is mathematically incorrect. The choice between
  CP and NCP is a strict geometric trade-off governed by the density of
  your data.

  ** The Data-Poor Regime (Use NCP): When trial counts (\(N_j\)) are low
  or population variance (\(\sigma_\theta\)) is small, the prior
  dominates the likelihood. This creates Neal's Funnel in the CP space.
  NCP thrives here because it decouples the parameters, allowing the
  sampler to easily and safely explore the prior geometry.

  ** The Data-Rich Regime (Use CP): Consider a scenario where every
  agent completes \(10,000\) trials. The likelihood becomes massively
  informative, crushing the uncertainty around
  \(\theta_{\text{logit}, j}\) into a tiny, tight peak.If we use CP, the
  sampler effortlessly navigates this tight, data-driven likelihood
  space. If we use NCP, the sampler attempts to explore the broad
  \(z \sim \text{Normal}(0,1)\) prior, only to find that almost the
  entire space is immediately rejected by the highly restrictive
  likelihood. This makes NCP highly inefficient, resulting in massive
  autocorrelation, low ESS, and high computational time.
\end{itemize}

In other words, there is no universal panacea in cognitive modeling. We
do not engage in blind ``tinkering'' to get our models to fit. Instead,
we adhere to the Iterative Bayesian Workflow: we fit the model, deploy a
strict diagnostic battery to identify the exact geometric failure
(prior-dominated funnels vs.~data-dominated funnels), and apply the
mathematically appropriate parameterization.

Because typical cognitive tasks (like our Matching Pennies game) involve
relatively sparse data (\(N_j \approx 100\)) and substantial
population-level shrinkage, we will utilize the Non-Centered
Parameterization as our default architecture as we proceed to construct
the more complex Memory and WSLS agents.

\section{The Epistemic Showdown: Visualizing
Shrinkage}\label{the-epistemic-showdown-visualizing-shrinkage}

We have theorized about the dangers of ignoring population structure
(Complete Pooling) and the dangers of ignoring noise in small samples
(No Pooling). In a rigorous Bayesian workflow, we do not rely on
intuition; we compute the alternatives and visualize the geometric
differences.We will quickly define and fit the two extreme models using
our strictly typed stan\_data\_centered. The complete and no-pooling
models are included here solely as visual contrasts --- they serve as
the ``wrong answer'' endpoints that make the partial pooling solution
legible. A full diagnostic battery on these two models would be
redundant: complete pooling is trivially identifiable (one parameter,
well-conditioned), and no-pooling is structurally identical to J
independent Ch. 6 models we already validated. The NCP model, which
actually matters, passed its full battery above.

\subsection{The Complete Pooling Model
(Underfitting)}\label{the-complete-pooling-model-underfitting}

This model assumes all agents are identical clones. The population
variance is forced to exactly \(0\). We estimate a single, universal
bias parameter.

\begin{Shaded}
\begin{Highlighting}[]
\NormalTok{stan\_full\_pooling\_code }\OtherTok{\textless{}{-}} \StringTok{"}
\StringTok{data \{}
\StringTok{  int\textless{}lower=1\textgreater{} N\_total;}
\StringTok{  array[N\_total] int\textless{}lower=0, upper=1\textgreater{} y;}
\StringTok{\}}
\StringTok{parameters \{}
\StringTok{  real theta\_logit; // A single parameter for the entire universe}
\StringTok{\}}
\StringTok{model \{}
\StringTok{  target += normal\_lpdf(theta\_logit | 0, 1.5);}
\StringTok{  target += bernoulli\_logit\_lpmf(y | theta\_logit);}
\StringTok{\}}
\StringTok{"}
\NormalTok{stan\_file\_full }\OtherTok{\textless{}{-}}\NormalTok{ here}\SpecialCharTok{::}\FunctionTok{here}\NormalTok{(}\StringTok{"stan"}\NormalTok{, }\StringTok{"ch07\_biased\_full\_pooling.stan"}\NormalTok{)}
\FunctionTok{writeLines}\NormalTok{(stan\_full\_pooling\_code, stan\_file\_full)}
\NormalTok{mod\_full }\OtherTok{\textless{}{-}} \FunctionTok{cmdstan\_model}\NormalTok{(stan\_file\_full)}

\NormalTok{fit\_full\_path }\OtherTok{\textless{}{-}}\NormalTok{ here}\SpecialCharTok{::}\FunctionTok{here}\NormalTok{(}\StringTok{"simmodels"}\NormalTok{, }\StringTok{"ch07\_fit\_full\_pool.rds"}\NormalTok{)}

\ControlFlowTok{if}\NormalTok{ (regenerate\_simulations }\SpecialCharTok{||} \SpecialCharTok{!}\FunctionTok{file.exists}\NormalTok{(fit\_full\_path)) \{}
\NormalTok{  fit\_full }\OtherTok{\textless{}{-}}\NormalTok{ mod\_full}\SpecialCharTok{$}\FunctionTok{sample}\NormalTok{(}
    \AttributeTok{data            =} \FunctionTok{list}\NormalTok{(}\AttributeTok{N\_total =}\NormalTok{ stan\_data\_centered}\SpecialCharTok{$}\NormalTok{N\_total, }\AttributeTok{y =}\NormalTok{ stan\_data\_centered}\SpecialCharTok{$}\NormalTok{y),}
    \AttributeTok{seed            =} \DecValTok{123}\NormalTok{, }\AttributeTok{chains =} \DecValTok{4}\NormalTok{, }\AttributeTok{parallel\_chains =} \DecValTok{4}\NormalTok{, }\AttributeTok{refresh =} \DecValTok{0}
\NormalTok{  )}
\NormalTok{  fit\_full}\SpecialCharTok{$}\FunctionTok{save\_object}\NormalTok{(fit\_full\_path)}
\NormalTok{\} }\ControlFlowTok{else}\NormalTok{ \{}
\NormalTok{  fit\_full }\OtherTok{\textless{}{-}} \FunctionTok{readRDS}\NormalTok{(fit\_full\_path)}
\NormalTok{\}}

\NormalTok{full\_pool\_est }\OtherTok{\textless{}{-}}\NormalTok{ posterior}\SpecialCharTok{::}\FunctionTok{summarise\_draws}\NormalTok{(fit\_full}\SpecialCharTok{$}\FunctionTok{draws}\NormalTok{(}\StringTok{"theta\_logit"}\NormalTok{), }\StringTok{"median"}\NormalTok{)}\SpecialCharTok{$}\NormalTok{median}
\end{Highlighting}
\end{Shaded}

\subsection{The No Pooling Model
(Overfitting)}\label{the-no-pooling-model-overfitting}

This model assumes every agent is a completely isolated island. The
population variance is implicitly assumed to be \(\infty\). We estimate
an independent parameter for each agent, equipped only with a static,
weakly informative prior.

\begin{Shaded}
\begin{Highlighting}[]
\NormalTok{stan\_no\_pooling\_code }\OtherTok{\textless{}{-}} \StringTok{"}
\StringTok{data \{}
\StringTok{  int\textless{}lower=1\textgreater{} N\_total;}
\StringTok{  int\textless{}lower=1\textgreater{} N\_subjects;}
\StringTok{  array[N\_total] int\textless{}lower=1, upper=N\_subjects\textgreater{} subj\_id;}
\StringTok{  array[N\_total] int\textless{}lower=0, upper=1\textgreater{} y;}
\StringTok{\}}
\StringTok{parameters \{}
\StringTok{  vector[N\_subjects] theta\_logit; // N\_subjects completely independent parameters}
\StringTok{\}}
\StringTok{model \{}
\StringTok{  // Static prior: No hierarchical learning}
\StringTok{  target += normal\_lpdf(theta\_logit | 0, 1.5);}
\StringTok{  target += bernoulli\_logit\_lpmf(y | theta\_logit[subj\_id]);}
\StringTok{\}}
\StringTok{"}
\NormalTok{stan\_file\_no }\OtherTok{\textless{}{-}}\NormalTok{ here}\SpecialCharTok{::}\FunctionTok{here}\NormalTok{(}\StringTok{"stan"}\NormalTok{, }\StringTok{"ch07\_biased\_no\_pooling.stan"}\NormalTok{)}
\FunctionTok{writeLines}\NormalTok{(stan\_no\_pooling\_code, stan\_file\_no)}
\NormalTok{mod\_no }\OtherTok{\textless{}{-}} \FunctionTok{cmdstan\_model}\NormalTok{(stan\_file\_no)}

\CommentTok{\# Fit the No Pooling model}
\NormalTok{fit\_no\_path }\OtherTok{\textless{}{-}}\NormalTok{ here}\SpecialCharTok{::}\FunctionTok{here}\NormalTok{(}\StringTok{"simmodels"}\NormalTok{, }\StringTok{"ch07\_fit\_no\_pool.rds"}\NormalTok{)}

\ControlFlowTok{if}\NormalTok{ (regenerate\_simulations }\SpecialCharTok{||} \SpecialCharTok{!}\FunctionTok{file.exists}\NormalTok{(fit\_no\_path)) \{}
\NormalTok{  fit\_no }\OtherTok{\textless{}{-}}\NormalTok{ mod\_no}\SpecialCharTok{$}\FunctionTok{sample}\NormalTok{(}
    \AttributeTok{data            =}\NormalTok{ stan\_data\_centered,}
    \AttributeTok{seed            =} \DecValTok{123}\NormalTok{, }\AttributeTok{chains =} \DecValTok{4}\NormalTok{, }\AttributeTok{parallel\_chains =} \DecValTok{4}\NormalTok{, }\AttributeTok{refresh =} \DecValTok{0}
\NormalTok{  )}
\NormalTok{  fit\_no}\SpecialCharTok{$}\FunctionTok{save\_object}\NormalTok{(fit\_no\_path)}
\NormalTok{\} }\ControlFlowTok{else}\NormalTok{ \{}
\NormalTok{  fit\_no }\OtherTok{\textless{}{-}} \FunctionTok{readRDS}\NormalTok{(fit\_no\_path)}
\NormalTok{\}}

\NormalTok{no\_pool\_summary }\OtherTok{\textless{}{-}}\NormalTok{ posterior}\SpecialCharTok{::}\FunctionTok{summarise\_draws}\NormalTok{(fit\_no}\SpecialCharTok{$}\FunctionTok{draws}\NormalTok{(}\StringTok{"theta\_logit"}\NormalTok{), }\StringTok{"median"}\NormalTok{) }\SpecialCharTok{|\textgreater{}} 
  \FunctionTok{mutate}\NormalTok{(}\AttributeTok{agent\_id =} \FunctionTok{as.integer}\NormalTok{(}\FunctionTok{str\_extract}\NormalTok{(variable, }\StringTok{"(?\textless{}=}\SpecialCharTok{\textbackslash{}\textbackslash{}}\StringTok{[)}\SpecialCharTok{\textbackslash{}\textbackslash{}}\StringTok{d+(?=}\SpecialCharTok{\textbackslash{}\textbackslash{}}\StringTok{])"}\NormalTok{))) }\SpecialCharTok{|\textgreater{}} 
  \FunctionTok{rename}\NormalTok{(}\AttributeTok{theta\_no\_pool =}\NormalTok{ median)}
\end{Highlighting}
\end{Shaded}

\subsection{The Geometry of Shrinkage}\label{the-geometry-of-shrinkage}

We now have three distinct sets of inferences: * No Pooling: The raw,
isolated estimates (highly vulnerable to the noise of low \(N_j\)). *
Complete Pooling: The rigid, universal average (blind to genuine
individual differences). * Partial Pooling: Our hierarchical NCP model
estimates (the precision-weighted Bayesian compromise).

To truly evaluate the epistemological value of these models, we must
plot them not just against each other, but against the True Generative
Parameters (the hidden reality we simulated in R). We will also extract
the naive arithmetic mean of the No Pooling estimates to contrast it
with the precision-weighted Partial Pooling mean.

\begin{Shaded}
\begin{Highlighting}[]
\CommentTok{\# 1. Extract the Partial Pooling Population and Individual Estimates (from the NCP model)}
\NormalTok{mu\_partial\_pool\_est }\OtherTok{\textless{}{-}}\NormalTok{ posterior}\SpecialCharTok{::}\FunctionTok{summarise\_draws}\NormalTok{(fit\_ncp}\SpecialCharTok{$}\FunctionTok{draws}\NormalTok{(}\StringTok{"mu\_theta"}\NormalTok{), }\StringTok{"median"}\NormalTok{)}\SpecialCharTok{$}\NormalTok{median}

\NormalTok{partial\_pool\_summary }\OtherTok{\textless{}{-}}\NormalTok{ posterior}\SpecialCharTok{::}\FunctionTok{summarise\_draws}\NormalTok{(fit\_ncp}\SpecialCharTok{$}\FunctionTok{draws}\NormalTok{(}\StringTok{"theta\_logit"}\NormalTok{), }\StringTok{"median"}\NormalTok{) }\SpecialCharTok{|\textgreater{}} 
  \FunctionTok{mutate}\NormalTok{(}\AttributeTok{agent\_id =} \FunctionTok{as.integer}\NormalTok{(}\FunctionTok{str\_extract}\NormalTok{(variable, }\StringTok{"(?\textless{}=}\SpecialCharTok{\textbackslash{}\textbackslash{}}\StringTok{[)}\SpecialCharTok{\textbackslash{}\textbackslash{}}\StringTok{d+(?=}\SpecialCharTok{\textbackslash{}\textbackslash{}}\StringTok{])"}\NormalTok{))) }\SpecialCharTok{|\textgreater{}} 
  \FunctionTok{rename}\NormalTok{(}\AttributeTok{theta\_partial\_pool =}\NormalTok{ median)}

\CommentTok{\# 2. Calculate the No Pooling Mean (simple arithmetic average of the isolated estimates)}
\NormalTok{no\_pool\_mean\_est }\OtherTok{\textless{}{-}} \FunctionTok{mean}\NormalTok{(no\_pool\_summary}\SpecialCharTok{$}\NormalTok{theta\_no\_pool)}

\CommentTok{\# 3. Calculate the True Generative Mean (The actual center of our simulated universe)}
\NormalTok{true\_mean\_est }\OtherTok{\textless{}{-}} \FunctionTok{mean}\NormalTok{(true\_agent\_params}\SpecialCharTok{$}\NormalTok{theta\_logit)}

\CommentTok{\# 4. Combine into a single dataframe, including the TRUE generative parameters}
\NormalTok{comparison\_df }\OtherTok{\textless{}{-}}\NormalTok{ no\_pool\_summary }\SpecialCharTok{|\textgreater{}} 
\NormalTok{  dplyr}\SpecialCharTok{::}\FunctionTok{select}\NormalTok{(agent\_id, theta\_no\_pool) }\SpecialCharTok{|\textgreater{}} 
  \FunctionTok{left\_join}\NormalTok{(}
\NormalTok{    partial\_pool\_summary }\SpecialCharTok{|\textgreater{}}\NormalTok{ dplyr}\SpecialCharTok{::}\FunctionTok{select}\NormalTok{(agent\_id, theta\_partial\_pool), }
    \AttributeTok{by =} \StringTok{"agent\_id"}
\NormalTok{  ) }\SpecialCharTok{|\textgreater{}} 
  \CommentTok{\# Bring in the ground truth}
  \FunctionTok{left\_join}\NormalTok{(}
\NormalTok{    true\_agent\_params }\SpecialCharTok{|\textgreater{}}\NormalTok{ dplyr}\SpecialCharTok{::}\FunctionTok{select}\NormalTok{(agent\_id, n\_trials, }\AttributeTok{theta\_true =}\NormalTok{ theta\_logit),}
    \AttributeTok{by =} \StringTok{"agent\_id"}
\NormalTok{  ) }\SpecialCharTok{|\textgreater{}} 
  \CommentTok{\# Sort agents by their No Pooling estimate for a cleaner visual curve}
  \FunctionTok{arrange}\NormalTok{(theta\_no\_pool) }\SpecialCharTok{|\textgreater{}} 
  \FunctionTok{mutate}\NormalTok{(}\AttributeTok{rank =} \FunctionTok{row\_number}\NormalTok{())}

\CommentTok{\# The Ultimate Epistemic Plot}
\FunctionTok{ggplot}\NormalTok{(comparison\_df) }\SpecialCharTok{+}
  \CommentTok{\# {-}{-}{-} POPULATION AVERAGES (Mapped to linetype for a clean legend) {-}{-}{-}}
  \FunctionTok{geom\_hline}\NormalTok{(}\FunctionTok{aes}\NormalTok{(}\AttributeTok{yintercept =}\NormalTok{ no\_pool\_mean\_est, }\AttributeTok{linetype =} \StringTok{"No Pooling Mean"}\NormalTok{), }
             \AttributeTok{color =} \StringTok{"darkorange"}\NormalTok{, }\AttributeTok{linewidth =} \FloatTok{0.8}\NormalTok{) }\SpecialCharTok{+}
  \FunctionTok{geom\_hline}\NormalTok{(}\FunctionTok{aes}\NormalTok{(}\AttributeTok{yintercept =}\NormalTok{ full\_pool\_est, }\AttributeTok{linetype =} \StringTok{"Complete Pooling Mean"}\NormalTok{), }
             \AttributeTok{color =} \StringTok{"darkred"}\NormalTok{, }\AttributeTok{linewidth =} \FloatTok{0.8}\NormalTok{) }\SpecialCharTok{+}
  \FunctionTok{geom\_hline}\NormalTok{(}\FunctionTok{aes}\NormalTok{(}\AttributeTok{yintercept =}\NormalTok{ mu\_partial\_pool\_est, }\AttributeTok{linetype =} \StringTok{"Partial Pooling Mean (\textbackslash{}U03BC)"}\NormalTok{), }
             \AttributeTok{color =} \StringTok{"steelblue"}\NormalTok{, }\AttributeTok{linewidth =} \FloatTok{1.2}\NormalTok{) }\SpecialCharTok{+}
  \FunctionTok{geom\_hline}\NormalTok{(}\FunctionTok{aes}\NormalTok{(}\AttributeTok{yintercept =}\NormalTok{ true\_mean\_est, }\AttributeTok{linetype =} \StringTok{"True Generative Mean"}\NormalTok{), }
             \AttributeTok{color =} \StringTok{"black"}\NormalTok{, }\AttributeTok{linewidth =} \FloatTok{0.8}\NormalTok{) }\SpecialCharTok{+}
  
  \CommentTok{\# {-}{-}{-} INDIVIDUAL ESTIMATES {-}{-}{-}}
  \CommentTok{\# Shrinkage Vectors (from No Pool to Partial Pool)}
  \FunctionTok{geom\_segment}\NormalTok{(}\FunctionTok{aes}\NormalTok{(}\AttributeTok{x =}\NormalTok{ rank, }\AttributeTok{xend =}\NormalTok{ rank, }
                   \AttributeTok{y =}\NormalTok{ theta\_no\_pool, }\AttributeTok{yend =}\NormalTok{ theta\_partial\_pool), }
               \AttributeTok{color =} \StringTok{"gray65"}\NormalTok{, }\AttributeTok{arrow =} \FunctionTok{arrow}\NormalTok{(}\AttributeTok{length =} \FunctionTok{unit}\NormalTok{(}\FloatTok{0.1}\NormalTok{, }\StringTok{"cm"}\NormalTok{))) }\SpecialCharTok{+}
  
  \CommentTok{\# True Generative Parameters (Black Crosses)}
  \FunctionTok{geom\_point}\NormalTok{(}\FunctionTok{aes}\NormalTok{(}\AttributeTok{x =}\NormalTok{ rank, }\AttributeTok{y =}\NormalTok{ theta\_true, }\AttributeTok{color =} \StringTok{"True Value"}\NormalTok{), }
             \AttributeTok{shape =} \DecValTok{4}\NormalTok{, }\AttributeTok{size =} \DecValTok{2}\NormalTok{, }\AttributeTok{stroke =} \DecValTok{1}\NormalTok{) }\SpecialCharTok{+}
             
  \CommentTok{\# No Pooling Estimates (Open Circles)}
  \FunctionTok{geom\_point}\NormalTok{(}\FunctionTok{aes}\NormalTok{(}\AttributeTok{x =}\NormalTok{ rank, }\AttributeTok{y =}\NormalTok{ theta\_no\_pool, }\AttributeTok{color =} \StringTok{"No Pooling"}\NormalTok{, }\AttributeTok{size =}\NormalTok{ n\_trials), }
             \AttributeTok{shape =} \DecValTok{1}\NormalTok{, }\AttributeTok{stroke =} \FloatTok{1.2}\NormalTok{) }\SpecialCharTok{+}
             
  \CommentTok{\# Partial Pooling Estimates (Solid Points)}
  \FunctionTok{geom\_point}\NormalTok{(}\FunctionTok{aes}\NormalTok{(}\AttributeTok{x =}\NormalTok{ rank, }\AttributeTok{y =}\NormalTok{ theta\_partial\_pool, }\AttributeTok{color =} \StringTok{"Partial Pooling"}\NormalTok{, }\AttributeTok{size =}\NormalTok{ n\_trials), }
             \AttributeTok{alpha =} \FloatTok{0.8}\NormalTok{) }\SpecialCharTok{+}
  
  \CommentTok{\# {-}{-}{-} SCALES \& THEMES {-}{-}{-}}
  \FunctionTok{scale\_linetype\_manual}\NormalTok{(}
    \AttributeTok{name =} \StringTok{"Population Averages"}\NormalTok{,}
    \AttributeTok{values =} \FunctionTok{c}\NormalTok{(}
      \StringTok{"True Generative Mean"} \OtherTok{=} \StringTok{"solid"}\NormalTok{,}
      \StringTok{"Partial Pooling Mean (\textbackslash{}U03BC)"} \OtherTok{=} \StringTok{"dotted"}\NormalTok{,}
      \StringTok{"Complete Pooling Mean"} \OtherTok{=} \StringTok{"dashed"}\NormalTok{,}
      \StringTok{"No Pooling Mean"} \OtherTok{=} \StringTok{"dotdash"}
\NormalTok{    )}
\NormalTok{  ) }\SpecialCharTok{+}
  \FunctionTok{scale\_color\_manual}\NormalTok{(}
    \AttributeTok{name =} \StringTok{"Individual Estimates"}\NormalTok{,}
    \AttributeTok{values =} \FunctionTok{c}\NormalTok{(}
      \StringTok{"True Value"} \OtherTok{=} \StringTok{"black"}\NormalTok{,}
      \StringTok{"No Pooling"} \OtherTok{=} \StringTok{"darkorange"}\NormalTok{, }
      \StringTok{"Partial Pooling"} \OtherTok{=} \StringTok{"steelblue"}
\NormalTok{    )}
\NormalTok{  ) }\SpecialCharTok{+}
  \FunctionTok{labs}\NormalTok{(}
    \AttributeTok{title =} \StringTok{"Bayesian Shrinkage: Partial Pooling vs. the Extremes"}\NormalTok{,}
    \AttributeTok{subtitle =} \StringTok{"Arrows show shrinkage from No Pooling → Partial Pooling. Longest arrows = fewest trials. Partial pooling lands closest to truth (black solid line)."}\NormalTok{,}
    \AttributeTok{x =} \StringTok{"Agents (Sorted by No{-}Pooling Estimate)"}\NormalTok{,}
    \AttributeTok{y =} \StringTok{"Estimated Bias Parameter (Logit Scale)"}\NormalTok{,}
    \AttributeTok{size =} \StringTok{"Trial Count (N\_j)"}
\NormalTok{  ) }\SpecialCharTok{+}
  \FunctionTok{theme}\NormalTok{(}
    \AttributeTok{axis.text.x =} \FunctionTok{element\_blank}\NormalTok{(), }
    \AttributeTok{axis.ticks.x =} \FunctionTok{element\_blank}\NormalTok{(),}
    \AttributeTok{legend.position =} \StringTok{"right"}
\NormalTok{  )}
\end{Highlighting}
\end{Shaded}

\begin{center}
\includegraphics[width=0.8\linewidth,height=\textheight,keepaspectratio]{06-MultilevelModels_files/figure-pdf/ch07_shrinkage_viz-1.pdf}
\end{center}

Study this plot carefully. It contains the entire epistemological
justification for hierarchical modeling, visible across two dimensions:
the horizontal population averages and the individual agent estimates.

\subsection{The Horizontal Lines: The Quartet of Population
Means}\label{the-horizontal-lines-the-quartet-of-population-means}

\begin{itemize}
\item
  True Generative Mean (Solid Black): The actual center of our simulated
  universe. This is the structural reality we are trying to recover.
\item
  No Pooling Mean (Dot-dash Orange): The unweighted arithmetic average
  of the isolated estimates. Because it treats all agents equally, it is
  highly susceptible to being dragged away from the truth by extreme
  outliers with very little data.
\item
  Complete Pooling Mean (Dashed Red): The unweighted average of all
  trials. Because it ignores agent boundaries, it allows participants
  who completed massive numbers of trials to disproportionately pull the
  center of gravity toward themselves.
\item
  Partial Pooling Mean (\(\mu_\theta\), Dotted Blue): The
  precision-weighted compromise. By mathematically down-weighting the
  noisy, low-\(N_j\) agents, it reliably lands closest to the true
  generative mean.
\end{itemize}

\subsection{The Points and Arrows: Information-Weighted
Regularization}\label{the-points-and-arrows-information-weighted-regularization}

\begin{itemize}
\item
  The Overfitting Extremes: Look at the tails of the graph. The orange
  open circles (No Pooling) are extreme. If we trusted No Pooling, we
  would conclude these agents have massive, almost deterministic biases.
\item
  The Power of the Prior: The solid blue dots (Partial Pooling) are
  systematically pulled inward. The hierarchical model ``knows'' that
  extremely high or low values are mathematically unlikely given the
  behavior of the rest of the cohort.
\item
  The Calibration of Skepticism: Look at the size of the dots,
  representing the trial count (\(N_j\)). The longest arrows (most
  aggressive shrinkage) almost exclusively belong to the smallest dots.
  When an agent has sparse data, the posterior relies heavily on the
  population prior. The largest dots barely move; when an agent has
  massive amounts of data, the likelihood overwhelms the prior, allowing
  them to remain safely extreme.
\end{itemize}

\section{The Memory Agent: Two Correlated Individual
Differences}\label{the-memory-agent-two-correlated-individual-differences}

\subsection{Introduction}\label{introduction-1}

The Biased Agent established the full Bayesian workflow for a single
individual-level parameter: simulate, fit, diagnose, recover, calibrate.
It also taught us, the hard way, that the Centered Parameterization
fails when population variance is small, and that the Non-Centered
Parameterization resolves this by decoupling individual and population
parameters geometrically. We proceed to the Memory Agent with both of
those lessons in hand.

The Memory Agent is structurally richer. Recall from Chapter 3 its
decision rule. On each trial \(i\), agent \(j\) maintains a running
average of the opponent's past choices \(m_{ij}\) and uses it as a
predictive signal:

\[h_{ij} \sim \text{Bernoulli}(\text{logit}^{-1}(\alpha_j + \beta_j \cdot \text{logit}(m_{ij})))\]

\[m_{ij} = \frac{1}{i-1}\sum_{k=1}^{i-1} o_{kj}\]

There are now two individual-level parameters per agent: \(\alpha_j\),
the baseline directional bias, and \(\beta_j\), the sensitivity to
opponent history. They need not vary independently across the
population. An agent who has a strong directional tendency might also
track the opponent more aggressively --- or might not. Fitting two
separate univariate hierarchical models forecloses this question before
the data can answer it. We model the parameters jointly.

\textbf{Experimental design: the reversal paradigm.} In Ch. 6 we
discovered that fitting the memory model to a \emph{constant-bias}
opponent produces a strong positive correlation between \texttt{bias}
and \texttt{beta} in the posterior --- both parameters explain the same
rightward tendency and cannot be told apart. The reversal learning
design (opponent plays 80\% right for trials 1--120, then 20\% right for
trials 121--240) breaks this confound: in the second phase, a pure-bias
agent continues right while a memory agent flips left, providing the
discriminating contrast the two-parameter model needs. We therefore use
the reversal schedule saved in Ch. 6 as the shared opponent sequence for
all J agents here.

\subsubsection{The Bivariate Population
Distribution}\label{the-bivariate-population-distribution}

\[\begin{pmatrix} \alpha_j \\ \beta_j \end{pmatrix} \sim \mathcal{N}_2\!\left(\boldsymbol{\mu},\, \boldsymbol{\Sigma}\right), \quad \boldsymbol{\Sigma} = \operatorname{diag}(\boldsymbol{\sigma})\, \Omega\, \operatorname{diag}(\boldsymbol{\sigma}), \quad \Omega = \begin{pmatrix} 1 & \rho \\ \rho & 1 \end{pmatrix}\]

This decomposition separates how much individuals differ on each
cognitive dimension (\(\sigma_\alpha\), \(\sigma_\beta\)) from how those
dimensions covary (\(\rho\)). The scalar \(\rho \in (-1, 1)\) is the key
new inferential target.

\subsubsection{The LKJ Prior}\label{the-lkj-prior}

We need a prior on the correlation matrix \(\Omega\). Stan provides the
LKJ distribution (Lewandowski, Kurowicka, and Joe, 2009), with a single
shape parameter \(\eta > 0\):

\[\Omega \sim \text{LKJ}(\eta) \quad \propto \quad (\det \Omega)^{\eta - 1}\]

When \(\eta = 1\) the distribution is uniform over valid correlation
matrices. When \(\eta > 1\) mass concentrates toward the identity,
encoding the prior belief that parameters tend to vary independently
unless the data argue otherwise. For a \(2 \times 2\) matrix, the
marginal on \(\rho\) reduces to a rescaled Beta:
\(\rho \sim 2 \cdot \text{Beta}(\eta, \eta) - 1\). With our default
\(\eta = 2\), the prior is unimodal at zero and weakly regularizes
against extreme correlations.

Stan samples the lower-triangular \textbf{Cholesky factor}
\(\mathbf{L}_\Omega\) (declared as
\texttt{cholesky\_factor\_corr{[}2{]}}) under the
\texttt{lkj\_corr\_cholesky(eta)} prior, recovering
\(\Omega = \mathbf{L}_\Omega \mathbf{L}_\Omega^\top\) in
\texttt{generated\ quantities}. This parameterization is unconstrained
on its natural manifold.

\subsubsection{Non-Centered
Parameterization}\label{non-centered-parameterization}

The scalar NCP from the biased agent
(\(\theta_j = \mu + z_j \cdot \sigma\)) generalizes directly. We
introduce a \(2 \times J\) matrix of independent standard normals:

\[z_{kj} \sim \mathcal{N}(0, 1) \text{ i.i.d.}, \qquad \begin{pmatrix}\alpha_j \\ \beta_j\end{pmatrix} = \begin{pmatrix}\mu_\alpha \\ \mu_\beta\end{pmatrix} + \operatorname{diag}(\boldsymbol{\sigma})\,\mathbf{L}_\Omega\,\mathbf{z}_j\]

The sampler traverses the flat, isotropic space of \(\mathbf{Z}\), whose
geometry is entirely independent of \(\boldsymbol{\sigma}\) and
\(\Omega\). Neal's Funnel is structurally impossible.

\begin{figure}[H]

{\centering \includegraphics[width=0.8\linewidth,height=\textheight,keepaspectratio]{06-MultilevelModels_files/figure-pdf/hierarchical_memory_dag-1.png}

}

\caption{Standard Plate Notation (left) and Kruschke-Style Notation
(right) for the Hierarchical Memory Agent. The correlation node Omega
and the bivariate MVN distinguish this model from the biased agent.}

\end{figure}%

\subsection{Phase 1: Generative Plausibility of the Memory
Agent}\label{phase-1-generative-plausibility-of-the-memory-agent}

Before writing any Stan code, we construct the forward generative model
in R. This code must strictly mirror the mathematical hierarchy above.

One design decision deserves explicit justification. The memory signal
\(m_{ij}\) is a deterministic function of the opponent's past choices.
We could compute it inside Stan's \texttt{transformed\ data} block, but
this executes a nested loop on every leapfrog step of every chain.
Instead, we pre-compute the vector \texttt{opp\_rate\_prev} in R once
and pass it as data. This reduces the memory calculation to a single
vectorized pass before fitting begins, at zero loss of model fidelity.

The simulation function must also mirror the NCP Stan model exactly. We
draw \(\mathbf{z}_j \sim \mathcal{N}(\mathbf{0}, \mathbf{I}_2)\) and
apply the Cholesky transform, exactly as the Stan model will. For a
\(2 \times 2\) correlation matrix with off-diagonal \(\rho\), the lower
Cholesky factor is:

\[\mathbf{L}_\Omega = \begin{pmatrix} 1 & 0 \\ \rho & \sqrt{1-\rho^2} \end{pmatrix}\]

We sample \(\rho\) from the LKJ(2) marginal:
\(\rho \sim 2 \cdot \text{Beta}(2, 2) - 1\).

In accordance with the rigorous Bayesian workflow, the simulation
function is designed to either:

\begin{itemize}
\item
  Draw all hyperparameters directly from the stated priors (for Prior
  Predictive Checks).
\item
  Accept fixed ground-truth hyperparameters (for Parameter Recovery).
\end{itemize}

\begin{Shaded}
\begin{Highlighting}[]
\CommentTok{\#\textquotesingle{} Simulate Hierarchical Memory Agents}
\CommentTok{\#\textquotesingle{}}
\CommentTok{\#\textquotesingle{} @param n\_agents Integer, number of agents (J).}
\CommentTok{\#\textquotesingle{} @param n\_trials\_min Integer, minimum trials per agent.}
\CommentTok{\#\textquotesingle{} @param n\_trials\_max Integer, maximum trials per agent.}
\CommentTok{\#\textquotesingle{} @param mu\_alpha Numeric, population mean baseline bias (logit scale). NULL draws from prior.}
\CommentTok{\#\textquotesingle{} @param mu\_beta Numeric, population mean memory sensitivity. NULL draws from prior.}
\CommentTok{\#\textquotesingle{} @param sigma\_alpha Numeric, population SD of alpha. NULL draws from prior.}
\CommentTok{\#\textquotesingle{} @param sigma\_beta Numeric, population SD of beta. NULL draws from prior.}
\CommentTok{\#\textquotesingle{} @param rho Numeric, population correlation alpha{-}beta. NULL draws from LKJ(2).}
\CommentTok{\#\textquotesingle{} @param opponent\_choices Integer vector of length \textgreater{}= max(n\_trials\_max).}
\CommentTok{\#\textquotesingle{}   If provided, every agent plays a consecutive prefix of this fixed}
\CommentTok{\#\textquotesingle{}   sequence (realistic: all participants face the same opponent).}
\CommentTok{\#\textquotesingle{}   If NULL, generates independent Bernoulli draws using opponent\_rate.}
\CommentTok{\#\textquotesingle{} @param opponent\_rate Numeric in (0,1). Used only when opponent\_choices}
\CommentTok{\#\textquotesingle{}   is NULL. Default 0.7.}
\CommentTok{\#\textquotesingle{} @param seed Integer, for exact computational reproducibility.}
\CommentTok{\#\textquotesingle{}}
\CommentTok{\#\textquotesingle{} @return A tidy tibble containing both latent parameters and observed choices.}
\NormalTok{simulate\_hierarchical\_memory }\OtherTok{\textless{}{-}} \ControlFlowTok{function}\NormalTok{(}
  \AttributeTok{n\_agents      =} \DecValTok{50}\NormalTok{,}
  \AttributeTok{n\_trials\_min  =} \DecValTok{20}\NormalTok{,}
  \AttributeTok{n\_trials\_max  =} \DecValTok{120}\NormalTok{,}
  \AttributeTok{mu\_alpha      =} \ConstantTok{NULL}\NormalTok{,}
  \AttributeTok{mu\_beta       =} \ConstantTok{NULL}\NormalTok{,}
  \AttributeTok{sigma\_alpha   =} \ConstantTok{NULL}\NormalTok{,}
  \AttributeTok{sigma\_beta    =} \ConstantTok{NULL}\NormalTok{,}
  \AttributeTok{rho           =} \ConstantTok{NULL}\NormalTok{,}
  \AttributeTok{opponent\_choices =} \ConstantTok{NULL}\NormalTok{,}
  \AttributeTok{opponent\_rate    =} \FloatTok{0.7}\NormalTok{,}
  \AttributeTok{seed          =} \ConstantTok{NULL}
\NormalTok{) \{}
  \ControlFlowTok{if}\NormalTok{ (}\SpecialCharTok{!}\FunctionTok{is.null}\NormalTok{(seed)) }\FunctionTok{set.seed}\NormalTok{(seed)}

  \CommentTok{\# 1. Sample hyperparameters from priors if not provided}
  \ControlFlowTok{if}\NormalTok{ (}\FunctionTok{is.null}\NormalTok{(mu\_alpha))    mu\_alpha    }\OtherTok{\textless{}{-}} \FunctionTok{rnorm}\NormalTok{(}\DecValTok{1}\NormalTok{, }\DecValTok{0}\NormalTok{, }\DecValTok{1}\NormalTok{)}
  \ControlFlowTok{if}\NormalTok{ (}\FunctionTok{is.null}\NormalTok{(mu\_beta))     mu\_beta     }\OtherTok{\textless{}{-}} \FunctionTok{rnorm}\NormalTok{(}\DecValTok{1}\NormalTok{, }\DecValTok{0}\NormalTok{, }\DecValTok{1}\NormalTok{)}
  \ControlFlowTok{if}\NormalTok{ (}\FunctionTok{is.null}\NormalTok{(sigma\_alpha)) sigma\_alpha }\OtherTok{\textless{}{-}} \FunctionTok{rexp}\NormalTok{(}\DecValTok{1}\NormalTok{, }\DecValTok{1}\NormalTok{)}
  \ControlFlowTok{if}\NormalTok{ (}\FunctionTok{is.null}\NormalTok{(sigma\_beta))  sigma\_beta  }\OtherTok{\textless{}{-}} \FunctionTok{rexp}\NormalTok{(}\DecValTok{1}\NormalTok{, }\DecValTok{1}\NormalTok{)}
  \CommentTok{\# LKJ(2) marginal on rho for 2x2 = 2*Beta(2,2) {-} 1}
  \ControlFlowTok{if}\NormalTok{ (}\FunctionTok{is.null}\NormalTok{(rho))         rho         }\OtherTok{\textless{}{-}} \DecValTok{2} \SpecialCharTok{*} \FunctionTok{rbeta}\NormalTok{(}\DecValTok{1}\NormalTok{, }\DecValTok{2}\NormalTok{, }\DecValTok{2}\NormalTok{) }\SpecialCharTok{{-}} \DecValTok{1}

  \CommentTok{\# 2. Construct the lower{-}triangular Cholesky factor L\_Omega}
\NormalTok{  L\_Omega }\OtherTok{\textless{}{-}} \FunctionTok{matrix}\NormalTok{(}\FunctionTok{c}\NormalTok{(}\DecValTok{1}\NormalTok{, rho, }\DecValTok{0}\NormalTok{, }\FunctionTok{sqrt}\NormalTok{(}\DecValTok{1} \SpecialCharTok{{-}}\NormalTok{ rho}\SpecialCharTok{\^{}}\DecValTok{2}\NormalTok{)), }\AttributeTok{nrow =} \DecValTok{2}\NormalTok{, }\AttributeTok{ncol =} \DecValTok{2}\NormalTok{)}

  \CommentTok{\# 3. Sample individual parameters via the NCP transform (mirrors Stan exactly)}
  \CommentTok{\# Z is 2 x J; Scale = diag(sigma) \%*\% L\_Omega}
\NormalTok{  Z            }\OtherTok{\textless{}{-}} \FunctionTok{matrix}\NormalTok{(}\FunctionTok{rnorm}\NormalTok{(}\DecValTok{2} \SpecialCharTok{*}\NormalTok{ n\_agents), }\AttributeTok{nrow =} \DecValTok{2}\NormalTok{)}
\NormalTok{  Scale        }\OtherTok{\textless{}{-}} \FunctionTok{diag}\NormalTok{(}\FunctionTok{c}\NormalTok{(sigma\_alpha, sigma\_beta)) }\SpecialCharTok{\%*\%}\NormalTok{ L\_Omega}
  \CommentTok{\# indiv\_params is J x 2: column 1 = alpha\_j, column 2 = beta\_j}
\NormalTok{  indiv\_params }\OtherTok{\textless{}{-}} \FunctionTok{t}\NormalTok{(}\FunctionTok{c}\NormalTok{(mu\_alpha, mu\_beta) }\SpecialCharTok{+}\NormalTok{ Scale }\SpecialCharTok{\%*\%}\NormalTok{ Z)}

  \CommentTok{\# 4. Simulate trial{-}level data}
\NormalTok{  n\_trials\_vec }\OtherTok{\textless{}{-}} \FunctionTok{sample}\NormalTok{(n\_trials\_min}\SpecialCharTok{:}\NormalTok{n\_trials\_max, n\_agents, }\AttributeTok{replace =} \ConstantTok{TRUE}\NormalTok{)}

\NormalTok{  sim\_data }\OtherTok{\textless{}{-}} \FunctionTok{map\_dfr}\NormalTok{(}\FunctionTok{seq\_len}\NormalTok{(n\_agents), }\ControlFlowTok{function}\NormalTok{(j) \{}
\NormalTok{    n\_t     }\OtherTok{\textless{}{-}}\NormalTok{ n\_trials\_vec[j]}
\NormalTok{    alpha\_j }\OtherTok{\textless{}{-}}\NormalTok{ indiv\_params[j, }\DecValTok{1}\NormalTok{]}
\NormalTok{    beta\_j  }\OtherTok{\textless{}{-}}\NormalTok{ indiv\_params[j, }\DecValTok{2}\NormalTok{]}

    \ControlFlowTok{if}\NormalTok{ (}\SpecialCharTok{!}\FunctionTok{is.null}\NormalTok{(opponent\_choices)) \{}
      \CommentTok{\# Shared reversal schedule: agent plays the first n\_t steps of the}
      \CommentTok{\# fixed sequence. Agents who leave early see only the first phase.}
      \FunctionTok{stopifnot}\NormalTok{(}\FunctionTok{length}\NormalTok{(opponent\_choices) }\SpecialCharTok{\textgreater{}=}\NormalTok{ n\_t)}
\NormalTok{      opp\_choices }\OtherTok{\textless{}{-}}\NormalTok{ opponent\_choices[}\FunctionTok{seq\_len}\NormalTok{(n\_t)]}
\NormalTok{    \} }\ControlFlowTok{else}\NormalTok{ \{}
\NormalTok{      opp\_choices }\OtherTok{\textless{}{-}} \FunctionTok{rbinom}\NormalTok{(n\_t, }\DecValTok{1}\NormalTok{, opponent\_rate)}
\NormalTok{    \}}

    \CommentTok{\# Pre{-}compute running opponent rate BEFORE each trial (strict t{-}1 history).}
    \CommentTok{\# Trial 1 has no history: initialize to 0.5 (uninformative).}
    \CommentTok{\# Clip to [0.01, 0.99] so logit(opp\_rate\_prev) is always finite.}
\NormalTok{    opp\_rate\_prev    }\OtherTok{\textless{}{-}} \FunctionTok{numeric}\NormalTok{(n\_t)}
\NormalTok{    opp\_rate\_prev[}\DecValTok{1}\NormalTok{] }\OtherTok{\textless{}{-}} \FloatTok{0.5}
    \ControlFlowTok{if}\NormalTok{ (n\_t }\SpecialCharTok{\textgreater{}} \DecValTok{1}\NormalTok{) \{}
\NormalTok{      raw\_rates         }\OtherTok{\textless{}{-}} \FunctionTok{cumsum}\NormalTok{(opp\_choices[}\SpecialCharTok{{-}}\NormalTok{n\_t]) }\SpecialCharTok{/} \FunctionTok{seq\_len}\NormalTok{(n\_t }\SpecialCharTok{{-}} \DecValTok{1}\NormalTok{)}
\NormalTok{      opp\_rate\_prev[}\SpecialCharTok{{-}}\DecValTok{1}\NormalTok{] }\OtherTok{\textless{}{-}} \FunctionTok{pmax}\NormalTok{(}\FloatTok{0.01}\NormalTok{, }\FunctionTok{pmin}\NormalTok{(}\FloatTok{0.99}\NormalTok{, raw\_rates))}
\NormalTok{    \}}

\NormalTok{    logit\_p       }\OtherTok{\textless{}{-}}\NormalTok{ alpha\_j }\SpecialCharTok{+}\NormalTok{ beta\_j }\SpecialCharTok{*} \FunctionTok{qlogis}\NormalTok{(opp\_rate\_prev)}
\NormalTok{    agent\_choices }\OtherTok{\textless{}{-}} \FunctionTok{rbinom}\NormalTok{(n\_t, }\DecValTok{1}\NormalTok{, }\FunctionTok{plogis}\NormalTok{(logit\_p))}

    \FunctionTok{tibble}\NormalTok{(}
      \AttributeTok{agent\_id      =}\NormalTok{ j,}
      \AttributeTok{trial         =} \FunctionTok{seq\_len}\NormalTok{(n\_t),}
      \AttributeTok{n\_trials      =}\NormalTok{ n\_t,}
      \AttributeTok{choice        =}\NormalTok{ agent\_choices,}
      \AttributeTok{opp\_rate\_prev =}\NormalTok{ opp\_rate\_prev,}
      \AttributeTok{true\_alpha    =}\NormalTok{ alpha\_j,}
      \AttributeTok{true\_beta     =}\NormalTok{ beta\_j}
\NormalTok{    )}
\NormalTok{  \})}

  \FunctionTok{attr}\NormalTok{(sim\_data, }\StringTok{"hyperparameters"}\NormalTok{) }\OtherTok{\textless{}{-}} \FunctionTok{list}\NormalTok{(}
    \AttributeTok{mu\_alpha    =}\NormalTok{ mu\_alpha,    }\AttributeTok{mu\_beta     =}\NormalTok{ mu\_beta,}
    \AttributeTok{sigma\_alpha =}\NormalTok{ sigma\_alpha, }\AttributeTok{sigma\_beta  =}\NormalTok{ sigma\_beta,}
    \AttributeTok{rho         =}\NormalTok{ rho}
\NormalTok{  )}

\NormalTok{  sim\_data}
\NormalTok{\}}
\end{Highlighting}
\end{Shaded}

\subsubsection{Executing the Prior Predictive Check: The Joint
Population
Cloud}\label{executing-the-prior-predictive-check-the-joint-population-cloud}

For a bivariate model, the most revealing prior predictive display is a
scatter of individual \((\alpha_j, \beta_j)\) pairs across multiple
prior universes. This simultaneously reveals how dispersed individual
differences are (\(\sigma_\alpha\), \(\sigma_\beta\)), how correlated
the two cognitive dimensions tend to be (\(\rho\)), and whether any
prior universe generates implausible populations (e.g.,
\(|\beta_j| > 8\), meaning the memory signal completely overrides
choice).

\begin{Shaded}
\begin{Highlighting}[]
\CommentTok{\# Generate 12 prior universes}
\NormalTok{prior\_sims\_mem }\OtherTok{\textless{}{-}} \FunctionTok{map\_dfr}\NormalTok{(}\DecValTok{1}\SpecialCharTok{:}\DecValTok{12}\NormalTok{, }\ControlFlowTok{function}\NormalTok{(sim\_id) \{}
\NormalTok{  sim }\OtherTok{\textless{}{-}} \FunctionTok{simulate\_hierarchical\_memory}\NormalTok{(}
    \AttributeTok{n\_agents =} \DecValTok{80}\NormalTok{, }\AttributeTok{n\_trials\_min =} \DecValTok{1}\NormalTok{, }\AttributeTok{n\_trials\_max =} \DecValTok{1}\NormalTok{,}
    \AttributeTok{seed =}\NormalTok{ sim\_id }\SpecialCharTok{*} \DecValTok{13}
\NormalTok{  )}
\NormalTok{  hp }\OtherTok{\textless{}{-}} \FunctionTok{attr}\NormalTok{(sim, }\StringTok{"hyperparameters"}\NormalTok{)}
\NormalTok{  sim }\SpecialCharTok{|\textgreater{}}
    \FunctionTok{distinct}\NormalTok{(agent\_id, true\_alpha, true\_beta) }\SpecialCharTok{|\textgreater{}}
    \FunctionTok{mutate}\NormalTok{(}\AttributeTok{simulation\_id =} \FunctionTok{paste}\NormalTok{(}\StringTok{"Prior Universe"}\NormalTok{, sim\_id))}
\NormalTok{\})}

\FunctionTok{ggplot}\NormalTok{(prior\_sims\_mem, }\FunctionTok{aes}\NormalTok{(}\AttributeTok{x =}\NormalTok{ true\_alpha, }\AttributeTok{y =}\NormalTok{ true\_beta)) }\SpecialCharTok{+}
  \FunctionTok{geom\_point}\NormalTok{(}\AttributeTok{alpha =} \FloatTok{0.4}\NormalTok{, }\AttributeTok{size =} \FloatTok{0.9}\NormalTok{, }\AttributeTok{color =} \StringTok{"steelblue"}\NormalTok{) }\SpecialCharTok{+}
  \FunctionTok{geom\_vline}\NormalTok{(}\AttributeTok{xintercept =} \DecValTok{0}\NormalTok{, }\AttributeTok{linetype =} \StringTok{"dashed"}\NormalTok{, }\AttributeTok{color =} \StringTok{"gray60"}\NormalTok{, }\AttributeTok{linewidth =} \FloatTok{0.4}\NormalTok{) }\SpecialCharTok{+}
  \FunctionTok{geom\_hline}\NormalTok{(}\AttributeTok{yintercept =} \DecValTok{0}\NormalTok{, }\AttributeTok{linetype =} \StringTok{"dashed"}\NormalTok{, }\AttributeTok{color =} \StringTok{"gray60"}\NormalTok{, }\AttributeTok{linewidth =} \FloatTok{0.4}\NormalTok{) }\SpecialCharTok{+}
  \FunctionTok{facet\_wrap}\NormalTok{(}\SpecialCharTok{\textasciitilde{}}\NormalTok{simulation\_id, }\AttributeTok{ncol =} \DecValTok{4}\NormalTok{) }\SpecialCharTok{+}
  \FunctionTok{labs}\NormalTok{(}
    \AttributeTok{title    =} \StringTok{"Prior Predictive Checks: Implied Joint Population Distributions"}\NormalTok{,}
    \AttributeTok{subtitle =} \StringTok{"Populations generated by mu \textasciitilde{} N(0,1), sigma \textasciitilde{} Exp(1), Omega \textasciitilde{} LKJ(2)"}\NormalTok{,}
    \AttributeTok{x        =} \FunctionTok{expression}\NormalTok{(}\StringTok{"Baseline bias  "} \SpecialCharTok{*}\NormalTok{ alpha[j] }\SpecialCharTok{*} \StringTok{"  (logit scale)"}\NormalTok{),}
    \AttributeTok{y        =} \FunctionTok{expression}\NormalTok{(}\StringTok{"Memory sensitivity  "} \SpecialCharTok{*}\NormalTok{ beta[j])}
\NormalTok{  )}
\end{Highlighting}
\end{Shaded}

\begin{center}
\includegraphics[width=0.8\linewidth,height=\textheight,keepaspectratio]{06-MultilevelModels_files/figure-pdf/ch07_memory_prior_predictive_scatter-1.pdf}
\end{center}

A healthy prior generates a range of cloud shapes --- some circular (low
\(|\rho|\)), some elongated and tilted (high \(|\rho|\)), some compact
(low \(\sigma\)), some diffuse (high \(\sigma\)). No prior universe
should collapse all agents to a point or place all agents at implausibly
extreme values.

\subsection{Phase 2: Data Preparation and the NCP Stan
Model}\label{phase-2-data-preparation-and-the-ncp-stan-model}

Having established from the biased agent workflow that NCP is our
default architecture for cognitive tasks with sparse per-agent data, we
implement it directly. We skip the Centered Parameterization entirely.

\begin{Shaded}
\begin{Highlighting}[]
\CommentTok{\#\textquotesingle{} Prepare Memory Agent Data for Stan}
\CommentTok{\#\textquotesingle{}}
\CommentTok{\#\textquotesingle{} @param sim\_data Tibble from simulate\_hierarchical\_memory().}
\CommentTok{\#\textquotesingle{} @return A named list formatted for CmdStanR.}
\NormalTok{prep\_stan\_data\_memory }\OtherTok{\textless{}{-}} \ControlFlowTok{function}\NormalTok{(sim\_data) \{}
  \FunctionTok{list}\NormalTok{(}
    \AttributeTok{N\_total                  =} \FunctionTok{nrow}\NormalTok{(sim\_data),}
    \AttributeTok{N\_subjects               =} \FunctionTok{length}\NormalTok{(}\FunctionTok{unique}\NormalTok{(sim\_data}\SpecialCharTok{$}\NormalTok{agent\_id)),}
    \AttributeTok{subj\_id                  =}\NormalTok{ sim\_data}\SpecialCharTok{$}\NormalTok{agent\_id,}
    \AttributeTok{y                        =}\NormalTok{ sim\_data}\SpecialCharTok{$}\NormalTok{choice,}
    \AttributeTok{opp\_rate\_prev            =}\NormalTok{ sim\_data}\SpecialCharTok{$}\NormalTok{opp\_rate\_prev,}
    \AttributeTok{prior\_mu\_alpha\_mu        =} \DecValTok{0}\NormalTok{,}
    \AttributeTok{prior\_mu\_alpha\_sigma     =} \FloatTok{1.0}\NormalTok{,}
    \AttributeTok{prior\_mu\_beta\_mu         =} \DecValTok{0}\NormalTok{,}
    \AttributeTok{prior\_mu\_beta\_sigma      =} \FloatTok{1.0}\NormalTok{,}
    \AttributeTok{prior\_sigma\_lambda       =} \FloatTok{1.0}\NormalTok{,}
    \AttributeTok{prior\_Omega\_shape        =} \FloatTok{2.0}\NormalTok{,}
    \AttributeTok{run\_diagnostics          =} \DecValTok{1}
\NormalTok{  )}
\NormalTok{\}}


\CommentTok{\# {-}{-}{-} Load the reversal learning schedule from Ch. 6 {-}{-}{-}}
\CommentTok{\# This 240{-}trial sequence (80\% right → reversal → 20\% right) was validated}
\CommentTok{\# in Ch. 6 to cleanly identify bias and beta. All agents play the same}
\CommentTok{\# opponent schedule, with variable trial counts (80–240 trials).}
\NormalTok{reversal\_obj      }\OtherTok{\textless{}{-}} \FunctionTok{readRDS}\NormalTok{(}\FunctionTok{here}\NormalTok{(}\StringTok{"simdata"}\NormalTok{, }\StringTok{"ch06\_reversal\_data.rds"}\NormalTok{))}
\NormalTok{opponent\_reversal }\OtherTok{\textless{}{-}}\NormalTok{ reversal\_obj}\SpecialCharTok{$}\NormalTok{opponent\_choices  }\CommentTok{\# 240{-}element integer vector}

\FunctionTok{cat}\NormalTok{(}\StringTok{"Loaded reversal schedule from Ch. 6:"}\NormalTok{,}
    \FunctionTok{length}\NormalTok{(opponent\_reversal), }\StringTok{"trials |"}\NormalTok{,}
    \FunctionTok{sum}\NormalTok{(opponent\_reversal[}\DecValTok{1}\SpecialCharTok{:}\DecValTok{120}\NormalTok{]), }\StringTok{"right in phase 1 |"}\NormalTok{,}
    \FunctionTok{sum}\NormalTok{(opponent\_reversal[}\DecValTok{121}\SpecialCharTok{:}\DecValTok{240}\NormalTok{]), }\StringTok{"right in phase 2}\SpecialCharTok{\textbackslash{}n}\StringTok{"}\NormalTok{)}
\end{Highlighting}
\end{Shaded}

\begin{verbatim}
Loaded reversal schedule from Ch. 6: 240 trials | 93 right in phase 1 | 23 right in phase 2
\end{verbatim}

\begin{Shaded}
\begin{Highlighting}[]
\CommentTok{\# {-}{-}{-} Generate hierarchical memory data using the reversal schedule {-}{-}{-}}
\NormalTok{sim\_data\_memory }\OtherTok{\textless{}{-}} \FunctionTok{simulate\_hierarchical\_memory}\NormalTok{(}
  \AttributeTok{n\_agents         =} \DecValTok{50}\NormalTok{,}
  \AttributeTok{n\_trials\_min     =} \DecValTok{80}\NormalTok{,    }\CommentTok{\# Minimum: agents see at least the first phase}
  \AttributeTok{n\_trials\_max     =} \DecValTok{240}\NormalTok{,   }\CommentTok{\# Maximum: agents complete both phases}
  \AttributeTok{mu\_alpha         =}  \FloatTok{0.0}\NormalTok{,  }\CommentTok{\# No systematic directional baseline}
  \AttributeTok{mu\_beta          =}  \FloatTok{1.5}\NormalTok{,  }\CommentTok{\# Moderate positive memory tracking}
  \AttributeTok{sigma\_alpha      =}  \FloatTok{0.3}\NormalTok{,}
  \AttributeTok{sigma\_beta       =}  \FloatTok{0.5}\NormalTok{,}
  \AttributeTok{rho              =}  \FloatTok{0.3}\NormalTok{,  }\CommentTok{\# Mild positive correlation}
  \AttributeTok{opponent\_choices =}\NormalTok{ opponent\_reversal,  }\CommentTok{\# \textless{}{-}{-} reversal schedule from Ch. 6}
  \AttributeTok{seed             =} \DecValTok{1981}
\NormalTok{)}

\NormalTok{stan\_data\_memory }\OtherTok{\textless{}{-}} \FunctionTok{prep\_stan\_data\_memory}\NormalTok{(sim\_data\_memory)}

\FunctionTok{cat}\NormalTok{(}\FunctionTok{sprintf}\NormalTok{(}
  \StringTok{"Hierarchical memory data: \%d agents, \%d total trials, \%d–\%d trials per agent}\SpecialCharTok{\textbackslash{}n}\StringTok{"}\NormalTok{,}
\NormalTok{  stan\_data\_memory}\SpecialCharTok{$}\NormalTok{N\_subjects,}
\NormalTok{  stan\_data\_memory}\SpecialCharTok{$}\NormalTok{N\_total,}
  \FunctionTok{min}\NormalTok{(sim\_data\_memory}\SpecialCharTok{$}\NormalTok{n\_trials),}
  \FunctionTok{max}\NormalTok{(sim\_data\_memory}\SpecialCharTok{$}\NormalTok{n\_trials)}
\NormalTok{))}
\end{Highlighting}
\end{Shaded}

\begin{verbatim}
Hierarchical memory data: 50 agents, 8058 total trials, 84–237 trials per agent
\end{verbatim}

\subsubsection{Stan Model: Hierarchical Memory Agent
(NCP)}\label{stan-model-hierarchical-memory-agent-ncp}

The only new structural element relative to the biased agent NCP is the
bivariate extension. Instead of a scalar \(z_j\), we have a
\(2 \times J\) matrix \(\mathbf{Z}\). Instead of scaling by a single
\(\sigma\), we apply \texttt{diag\_pre\_multiply(sigma,\ L\_Omega)} ---
which computes
\(\operatorname{diag}(\boldsymbol{\sigma})\mathbf{L}_\Omega\) and
produces the full Cholesky scale matrix encoding both marginal SDs and
the correlation structure.

\begin{Shaded}
\begin{Highlighting}[]
\NormalTok{stan\_memory\_ncp\_code }\OtherTok{\textless{}{-}} \StringTok{"}
\StringTok{// Hierarchical Memory Agent Model (Non{-}Centered Parameterization)}
\StringTok{// alpha\_j (baseline bias) and beta\_j (memory sensitivity) are jointly}
\StringTok{// drawn from a bivariate Normal. NCP decouples sampling geometry from}
\StringTok{// the population parameters, preventing Neal\textquotesingle{}s Funnel.}

\StringTok{data \{}
\StringTok{  int\textless{}lower=1\textgreater{} N\_total;                    // Total observations}
\StringTok{  int\textless{}lower=1\textgreater{} N\_subjects;                 // Number of agents}
\StringTok{  array[N\_total] int\textless{}lower=1, upper=N\_subjects\textgreater{}  subj\_id;   // Agent index}
\StringTok{  array[N\_total] int\textless{}lower=0, upper=1\textgreater{}  y;       // Observed choices}
\StringTok{  vector\textless{}lower=0.01, upper=0.99\textgreater{}[N\_total] opp\_rate\_prev;}
\StringTok{  }
\StringTok{  real prior\_mu\_alpha\_mu;}
\StringTok{  real\textless{}lower=0\textgreater{} prior\_mu\_alpha\_sigma;}
\StringTok{  real prior\_mu\_beta\_mu;}
\StringTok{  real\textless{}lower=0\textgreater{} prior\_mu\_beta\_sigma;}
\StringTok{  real\textless{}lower=0\textgreater{} prior\_sigma\_lambda;}
\StringTok{  real\textless{}lower=0\textgreater{} prior\_Omega\_shape;}
\StringTok{  int\textless{}lower=0, upper=1\textgreater{} run\_diagnostics;}
\StringTok{\}}

\StringTok{parameters \{}
\StringTok{  // Population{-}level parameters}
\StringTok{  real mu\_alpha;              // Population mean baseline bias (logit scale)}
\StringTok{  real mu\_beta;               // Population mean memory sensitivity}
\StringTok{  vector\textless{}lower=0\textgreater{}[2] sigma;   // Marginal SDs}

\StringTok{  // Cholesky factor of the correlation matrix}
\StringTok{  cholesky\_factor\_corr[2] L\_Omega;}

\StringTok{  // Individual{-}level parameters (NCP)}
\StringTok{  matrix[2, N\_subjects] z;}
\StringTok{\}}

\StringTok{transformed parameters \{}
\StringTok{  matrix[N\_subjects, 2] indiv\_params;}
\StringTok{  \{}
\StringTok{    matrix[2, N\_subjects] deviations = diag\_pre\_multiply(sigma, L\_Omega) * z;}
\StringTok{    for (j in 1:N\_subjects) \{}
\StringTok{      indiv\_params[j, 1] = mu\_alpha + deviations[1, j];}
\StringTok{      indiv\_params[j, 2] = mu\_beta  + deviations[2, j];}
\StringTok{    \}}
\StringTok{  \}}
\StringTok{\}}

\StringTok{model \{}
\StringTok{  // Hyperpriors}
\StringTok{  target += normal\_lpdf(mu\_alpha | prior\_mu\_alpha\_mu, prior\_mu\_alpha\_sigma);}
\StringTok{  target += normal\_lpdf(mu\_beta  | prior\_mu\_beta\_mu,  prior\_mu\_beta\_sigma);}
\StringTok{  target += exponential\_lpdf(sigma | prior\_sigma\_lambda);}
\StringTok{  target += lkj\_corr\_cholesky\_lpdf(L\_Omega | prior\_Omega\_shape);}

\StringTok{  // Individual{-}level prior (NCP)}
\StringTok{  target += std\_normal\_lpdf(to\_vector(z));}

\StringTok{  // Likelihood}
\StringTok{  vector[N\_total] logit\_p;}
\StringTok{  for (i in 1:N\_total) \{}
\StringTok{    logit\_p[i] = indiv\_params[subj\_id[i], 1]}
\StringTok{               + indiv\_params[subj\_id[i], 2] * logit(opp\_rate\_prev[i]);}
\StringTok{  \}}
\StringTok{  target += bernoulli\_logit\_lpmf(y | logit\_p);}
\StringTok{\}}

\StringTok{generated quantities \{}
\StringTok{  matrix[2, 2] Omega = multiply\_lower\_tri\_self\_transpose(L\_Omega);}
\StringTok{  real rho = Omega[1, 2];}

\StringTok{  vector[N\_subjects] alpha    = indiv\_params[, 1];}
\StringTok{  vector[N\_subjects] beta\_mem = indiv\_params[, 2];}

\StringTok{  // Prior log{-}density}
\StringTok{  real lprior = normal\_lpdf(mu\_alpha | prior\_mu\_alpha\_mu, prior\_mu\_alpha\_sigma)}
\StringTok{              + normal\_lpdf(mu\_beta  | prior\_mu\_beta\_mu,  prior\_mu\_beta\_sigma)}
\StringTok{              + exponential\_lpdf(sigma | prior\_sigma\_lambda)}
\StringTok{              + lkj\_corr\_cholesky\_lpdf(L\_Omega | prior\_Omega\_shape)}
\StringTok{              + std\_normal\_lpdf(to\_vector(z));}

\StringTok{  // Predictive checks and pointwise log{-}likelihood}
\StringTok{  vector[N\_total] log\_lik;}
\StringTok{  array[N\_total] int y\_post\_rep;}
\StringTok{  array[N\_total] int y\_prior\_rep;}

\StringTok{  // Prior draws for prior{-}posterior update plots}
\StringTok{  real mu\_alpha\_prior    = normal\_rng(prior\_mu\_alpha\_mu, prior\_mu\_alpha\_sigma);}
\StringTok{  real mu\_beta\_prior     = normal\_rng(prior\_mu\_beta\_mu, prior\_mu\_beta\_sigma);}
\StringTok{  real sigma\_alpha\_prior = exponential\_rng(prior\_sigma\_lambda);}
\StringTok{  real sigma\_beta\_prior  = exponential\_rng(prior\_sigma\_lambda);}
\StringTok{  // LKJ(2) marginal on rho for 2x2: 2*Beta(2,2) {-} 1}
\StringTok{  real rho\_prior         = 2.0 * beta\_rng(2.0, 2.0) {-} 1.0;}

\StringTok{  if (run\_diagnostics) \{}
\StringTok{    matrix[2, 2] L\_Omega\_prior = lkj\_corr\_cholesky\_rng(2, prior\_Omega\_shape);}
\StringTok{    vector[2] sigma\_prior = [sigma\_alpha\_prior, sigma\_beta\_prior]\textquotesingle{};}
\StringTok{    matrix[2, N\_subjects] z\_prior;}
\StringTok{    for (j in 1:N\_subjects) \{ }
\StringTok{      z\_prior[1,j] = normal\_rng(0,1); }
\StringTok{      z\_prior[2,j] = normal\_rng(0,1); }
\StringTok{    \}}
\StringTok{    matrix[N\_subjects, 2] indiv\_prior;}
\StringTok{    \{}
\StringTok{      matrix[2, N\_subjects] dev\_prior = diag\_pre\_multiply(sigma\_prior, L\_Omega\_prior) * z\_prior;}
\StringTok{      for (j in 1:N\_subjects) \{}
\StringTok{        indiv\_prior[j, 1] = mu\_alpha\_prior + dev\_prior[1, j];}
\StringTok{        indiv\_prior[j, 2] = mu\_beta\_prior  + dev\_prior[2, j];}
\StringTok{      \}}
\StringTok{    \}}

\StringTok{    for (i in 1:N\_total) \{}
\StringTok{      int j = subj\_id[i];}
\StringTok{      real lp\_post  = indiv\_params[j,1] + indiv\_params[j,2] * logit(opp\_rate\_prev[i]);}
\StringTok{      real lp\_prior = indiv\_prior[j,1]  + indiv\_prior[j,2]  * logit(opp\_rate\_prev[i]);}

\StringTok{      log\_lik[i]     = bernoulli\_logit\_lpmf(y[i] | lp\_post);}
\StringTok{      y\_post\_rep[i]  = bernoulli\_logit\_rng(lp\_post);}
\StringTok{      y\_prior\_rep[i] = bernoulli\_logit\_rng(lp\_prior);}
\StringTok{    \}}
\StringTok{  \}}
\StringTok{\}}
\StringTok{"}

\NormalTok{stan\_file\_memory\_ncp }\OtherTok{\textless{}{-}}\NormalTok{ here}\SpecialCharTok{::}\FunctionTok{here}\NormalTok{(}\StringTok{"stan"}\NormalTok{, }\StringTok{"ch07\_multilevel\_memory\_ncp.stan"}\NormalTok{)}
\FunctionTok{writeLines}\NormalTok{(stan\_memory\_ncp\_code, stan\_file\_memory\_ncp)}
\NormalTok{mod\_memory\_ncp }\OtherTok{\textless{}{-}} \FunctionTok{cmdstan\_model}\NormalTok{(stan\_file\_memory\_ncp)}
\end{Highlighting}
\end{Shaded}

Two Stan details are worth noting. First, the local block
\texttt{\{\ matrix{[}2,J{]}\ deviations\ =\ ...;\ \}} in
\texttt{transformed\ parameters} confines the intermediate matrix to
local scope, preventing it from appearing in the parameter output.
Second, \texttt{lkj\_corr\_cholesky\_rng(2,\ 2.0)} in
\texttt{generated\ quantities} draws a full Cholesky factor from the
LKJ(2) prior. This is the correct way to generate prior predictive data
--- unlike the Beta marginal trick, it generalizes to the \(3 \times 3\)
case we will need for the WSLS agent.

\begin{tcolorbox}[enhanced jigsaw, toprule=.15mm, toptitle=1mm, bottomtitle=1mm, leftrule=.75mm, rightrule=.15mm, colframe=quarto-callout-tip-color-frame, left=2mm, colback=white, titlerule=0mm, coltitle=black, colbacktitle=quarto-callout-tip-color!10!white, opacityback=0, title=\textcolor{quarto-callout-tip-color}{\faLightbulb}\hspace{0.5em}{Stan Efficiency --- Cholesky Factorization for Correlation Matrices}, arc=.35mm, opacitybacktitle=0.6, breakable, bottomrule=.15mm]

Whenever a model includes a \textbf{correlation matrix} \(\Omega\) over
individual-level parameters, use the Cholesky-factorized
parameterization rather than sampling \(\Omega\) directly.

\begin{longtable}[]{@{}
  >{\raggedright\arraybackslash}p{(\linewidth - 4\tabcolsep) * \real{0.3333}}
  >{\raggedright\arraybackslash}p{(\linewidth - 4\tabcolsep) * \real{0.3333}}
  >{\raggedright\arraybackslash}p{(\linewidth - 4\tabcolsep) * \real{0.3333}}@{}}
\toprule\noalign{}
\begin{minipage}[b]{\linewidth}\raggedright
Parameterization
\end{minipage} & \begin{minipage}[b]{\linewidth}\raggedright
Type
\end{minipage} & \begin{minipage}[b]{\linewidth}\raggedright
Why avoid
\end{minipage} \\
\midrule\noalign{}
\endhead
\bottomrule\noalign{}
\endlastfoot
\texttt{corr\_matrix{[}K{]}\ Omega} + \texttt{lkj\_corr\_lpdf} & Direct
& Requires \(O(K^3)\) matrix square root and positive-definiteness
checks at every leapfrog step \\
\texttt{cholesky\_factor\_corr{[}K{]}\ L\_Omega} +
\texttt{lkj\_corr\_cholesky\_lpdf} & Cholesky & \(L_\Omega\) is already
lower-triangular; lies directly on its natural constrained manifold;
\(O(K^2)\) operations \\
\end{longtable}

\begin{Shaded}
\begin{Highlighting}[]
\CommentTok{// Avoid}
\KeywordTok{parameters}\NormalTok{ \{ }\DataTypeTok{corr\_matrix}\NormalTok{[}\DecValTok{2}\NormalTok{] Omega; \}}
\KeywordTok{model}\NormalTok{     \{ Omega \textasciitilde{} lkj\_corr(}\FloatTok{2.0}\NormalTok{); \}}

\CommentTok{// Prefer}
\KeywordTok{parameters}\NormalTok{ \{ }\DataTypeTok{cholesky\_factor\_corr}\NormalTok{[}\DecValTok{2}\NormalTok{] L\_Omega; \}}
\KeywordTok{model}\NormalTok{      \{ L\_Omega \textasciitilde{} lkj\_corr\_cholesky(}\FloatTok{2.0}\NormalTok{); \}}
\KeywordTok{transformed parameters}\NormalTok{ \{}
  \CommentTok{// Apply the Cholesky scale to NCP offsets in O(K\^{}2) — no matrix inversion}
  \DataTypeTok{matrix}\NormalTok{[}\DecValTok{2}\NormalTok{, J] deviations = diag\_pre\_multiply(sigma, L\_Omega) * z;}
\NormalTok{\}}
\KeywordTok{generated quantities}\NormalTok{ \{}
  \CommentTok{// Recover the full correlation matrix once per saved draw (not every leapfrog step)}
  \DataTypeTok{matrix}\NormalTok{[}\DecValTok{2}\NormalTok{, }\DecValTok{2}\NormalTok{] Omega = multiply\_lower\_tri\_self\_transpose(L\_Omega);}
  \DataTypeTok{real}\NormalTok{ rho = Omega[}\DecValTok{1}\NormalTok{, }\DecValTok{2}\NormalTok{];}
\NormalTok{\}}
\end{Highlighting}
\end{Shaded}

\textbf{Key point:} the matrix multiply
\texttt{diag\_pre\_multiply(sigma,\ L\_Omega)\ *\ z} gives the
individual-level deviations directly --- no matrix square-root needed at
run time. The correlation matrix \(\Omega\) is only reconstructed in
\texttt{generated\ quantities}, where it runs once per saved draw
instead of once per leapfrog step.

For \(K = 2\) parameters the saving is modest; for \(K = 5\) or more
(e.g., five per-subject RL parameters in a full model) it becomes
substantial.

\end{tcolorbox}

\subsubsection{Fitting the model}\label{fitting-the-model-1}

\begin{Shaded}
\begin{Highlighting}[]
\NormalTok{fit\_memory\_path }\OtherTok{\textless{}{-}}\NormalTok{ here}\SpecialCharTok{::}\FunctionTok{here}\NormalTok{(}\StringTok{"simmodels"}\NormalTok{, }\StringTok{"ch07\_fit\_memory\_ncp.rds"}\NormalTok{)}

\CommentTok{\# Initialize at prior means.}
\NormalTok{N\_subj\_val\_mem }\OtherTok{\textless{}{-}}\NormalTok{ stan\_data\_memory}\SpecialCharTok{$}\NormalTok{N\_subjects}

\NormalTok{init\_memory\_list }\OtherTok{\textless{}{-}} \FunctionTok{lapply}\NormalTok{(}\DecValTok{1}\SpecialCharTok{:}\DecValTok{4}\NormalTok{, }\ControlFlowTok{function}\NormalTok{(i) \{}
  \FunctionTok{list}\NormalTok{(}
    \AttributeTok{mu\_alpha =} \FunctionTok{rnorm}\NormalTok{(}\DecValTok{1}\NormalTok{, }\DecValTok{0}\NormalTok{, }\FloatTok{0.3}\NormalTok{), }
    \AttributeTok{mu\_beta =} \FunctionTok{rnorm}\NormalTok{(}\DecValTok{1}\NormalTok{, }\DecValTok{0}\NormalTok{, }\FloatTok{0.3}\NormalTok{), }
    \AttributeTok{sigma =} \FunctionTok{abs}\NormalTok{(}\FunctionTok{rnorm}\NormalTok{(}\DecValTok{2}\NormalTok{, }\DecValTok{0}\NormalTok{, }\FloatTok{0.3}\NormalTok{)) }\SpecialCharTok{+} \FloatTok{0.1}\NormalTok{, }
    \AttributeTok{L\_Omega =} \FunctionTok{matrix}\NormalTok{(}\FunctionTok{c}\NormalTok{(}\DecValTok{1}\NormalTok{, }\DecValTok{0}\NormalTok{, }\DecValTok{0}\NormalTok{, }\DecValTok{1}\NormalTok{), }\AttributeTok{nrow=}\DecValTok{2}\NormalTok{), }
    \AttributeTok{z =} \FunctionTok{matrix}\NormalTok{(}\DecValTok{0}\NormalTok{, }\AttributeTok{nrow=}\DecValTok{2}\NormalTok{, }\AttributeTok{ncol=}\NormalTok{N\_subj\_val\_mem)}
\NormalTok{  )}
\NormalTok{\})}

\ControlFlowTok{if}\NormalTok{ (regenerate\_simulations }\SpecialCharTok{||} \SpecialCharTok{!}\FunctionTok{file.exists}\NormalTok{(fit\_memory\_path)) \{}
\NormalTok{  fit\_memory }\OtherTok{\textless{}{-}}\NormalTok{ mod\_memory\_ncp}\SpecialCharTok{$}\FunctionTok{sample}\NormalTok{(}
    \AttributeTok{data            =}\NormalTok{ stan\_data\_memory,}
    \AttributeTok{seed            =} \DecValTok{123}\NormalTok{,}
    \AttributeTok{chains          =} \DecValTok{4}\NormalTok{,}
    \AttributeTok{parallel\_chains =} \DecValTok{4}\NormalTok{,}
    \AttributeTok{iter\_warmup     =} \DecValTok{1000}\NormalTok{,}
    \AttributeTok{iter\_sampling   =} \DecValTok{1000}\NormalTok{,}
    \AttributeTok{init            =}\NormalTok{ init\_memory\_list, }
    \AttributeTok{adapt\_delta     =} \FloatTok{0.90}\NormalTok{,   }
    \AttributeTok{max\_treedepth   =} \DecValTok{12}\NormalTok{,}
    \AttributeTok{refresh         =} \DecValTok{0}
\NormalTok{  )}
\NormalTok{  fit\_memory}\SpecialCharTok{$}\FunctionTok{save\_object}\NormalTok{(fit\_memory\_path)}
\NormalTok{\} }\ControlFlowTok{else}\NormalTok{ \{}
\NormalTok{  fit\_memory }\OtherTok{\textless{}{-}} \FunctionTok{readRDS}\NormalTok{(fit\_memory\_path)}
\NormalTok{\}}
\end{Highlighting}
\end{Shaded}

We raise \texttt{adapt\_delta} to 0.90. The bivariate NCP is
well-behaved globally, but the correlation parameter \(\rho\) introduces
a curved direction that can require longer leapfrog trajectories. The
higher target acceptance rate forces the dual-averaging algorithm to
find a smaller step size, reducing integration error along this
direction. If \texttt{max\_treedepth} warnings appear (distinct from
divergences), raise \texttt{max\_treedepth} further --- do not raise
\texttt{adapt\_delta} in response to those.

\subsubsection{Assessing the fitting
process}\label{assessing-the-fitting-process-1}

Even if the Stan compiler finishes without explicit warnings, our
workflow demands that we systematically verify the geometry of the
posterior and the efficiency of the sampler.

\begin{Shaded}
\begin{Highlighting}[]
\CommentTok{\#fit\_memory$cmdstan\_diagnose()}
\NormalTok{fit\_memory}\SpecialCharTok{$}\FunctionTok{diagnostic\_summary}\NormalTok{()}
\end{Highlighting}
\end{Shaded}

\begin{verbatim}
$num_divergent
[1] 0 0 0 0

$num_max_treedepth
[1] 0 0 0 0

$ebfmi
[1] 0.7753663 0.6990873 0.7302008 0.7121430
\end{verbatim}

We are looking for strict zeros in \texttt{num\_divergent}. With the
NCP, this should hold. Any divergences here are a signal that
\texttt{adapt\_delta} needs further adjustment, the priors are
misspecified, or there is a structural model issue.

\subsubsection{\texorpdfstring{Chain Mixing and Stationarity
(\(\hat{R}\))}{Chain Mixing and Stationarity (\textbackslash hat\{R\})}}\label{chain-mixing-and-stationarity-hatr-1}

\begin{Shaded}
\begin{Highlighting}[]
\NormalTok{draws\_memory }\OtherTok{\textless{}{-}}\NormalTok{ fit\_memory}\SpecialCharTok{$}\FunctionTok{draws}\NormalTok{()}

\FunctionTok{mcmc\_trace}\NormalTok{(}
\NormalTok{  draws\_memory,}
  \AttributeTok{pars =} \FunctionTok{c}\NormalTok{(}\StringTok{"mu\_alpha"}\NormalTok{, }\StringTok{"mu\_beta"}\NormalTok{, }\StringTok{"sigma[1]"}\NormalTok{, }\StringTok{"sigma[2]"}\NormalTok{, }\StringTok{"rho"}\NormalTok{),}
  \AttributeTok{np   =} \FunctionTok{nuts\_params}\NormalTok{(fit\_memory)}
\NormalTok{) }\SpecialCharTok{+}
  \FunctionTok{labs}\NormalTok{(}
    \AttributeTok{title    =} \StringTok{"Trace Plots: Population Parameters (Memory Agent NCP)"}\NormalTok{,}
    \AttributeTok{subtitle =} \StringTok{"All five population{-}level parameters, including rho. Stationary chains expected throughout."}
\NormalTok{  )}
\end{Highlighting}
\end{Shaded}

\begin{center}
\includegraphics[width=0.8\linewidth,height=\textheight,keepaspectratio]{06-MultilevelModels_files/figure-pdf/ch07_memory_chain_mixing-1.pdf}
\end{center}

\begin{Shaded}
\begin{Highlighting}[]
\NormalTok{draws\_memory }\SpecialCharTok{|\textgreater{}}
  \FunctionTok{subset\_draws}\NormalTok{(}\AttributeTok{variable =} \FunctionTok{c}\NormalTok{(}\StringTok{"mu\_alpha"}\NormalTok{, }\StringTok{"mu\_beta"}\NormalTok{, }\StringTok{"sigma[1]"}\NormalTok{, }\StringTok{"sigma[2]"}\NormalTok{, }\StringTok{"rho"}\NormalTok{,}
                            \StringTok{"alpha[1]"}\NormalTok{, }\StringTok{"alpha[2]"}\NormalTok{, }\StringTok{"beta\_mem[1]"}\NormalTok{, }\StringTok{"beta\_mem[2]"}\NormalTok{)) }\SpecialCharTok{|\textgreater{}}
  \FunctionTok{summarise\_draws}\NormalTok{(}\FunctionTok{default\_convergence\_measures}\NormalTok{()) }\SpecialCharTok{|\textgreater{}}
\NormalTok{  knitr}\SpecialCharTok{::}\FunctionTok{kable}\NormalTok{(}\AttributeTok{digits =} \DecValTok{3}\NormalTok{)}
\end{Highlighting}
\end{Shaded}

\begin{longtable}[]{@{}lrrr@{}}
\toprule\noalign{}
variable & rhat & ess\_bulk & ess\_tail \\
\midrule\noalign{}
\endhead
\bottomrule\noalign{}
\endlastfoot
mu\_alpha & 1.003 & 4536.130 & 3219.616 \\
mu\_beta & 1.000 & 5530.353 & 3365.868 \\
sigma{[}1{]} & 1.001 & 3094.326 & 2956.688 \\
sigma{[}2{]} & 1.000 & 1998.533 & 2940.560 \\
rho & 1.000 & 1646.126 & 2270.844 \\
alpha{[}1{]} & 1.002 & 5728.964 & 3161.287 \\
alpha{[}2{]} & 1.002 & 6427.232 & 3574.001 \\
beta\_mem{[}1{]} & 1.001 & 5303.794 & 2973.740 \\
beta\_mem{[}2{]} & 1.000 & 5467.761 & 3269.354 \\
\end{longtable}

The potential scale reduction factor (\(\hat{R}\)) must be below 1.01
for all parameters. The correlation parameter \texttt{rho} is informed
only by covariation \emph{across} agents, not within any individual
agent's trial sequence, so its ESS will be lower than the marginal
parameters. ESS \textgreater{} 400 remains our threshold. An ESS below
200 for \texttt{rho} warrants investigation of the correlation manifold
geometry.

\paragraph{Effective Sample Size
(ESS)}\label{effective-sample-size-ess-1}

Monitor \texttt{ess\_tail} for \texttt{sigma{[}1{]}} and
\texttt{sigma{[}2{]}}. Population SDs can develop mild funnel-like
behavior even under NCP when the likelihood is highly informative about
individual parameters, because the \(z_j\) draws must simultaneously
explain many agents. If \texttt{ess\_tail} is low for the sigmas but
\texttt{ess\_bulk} is acceptable, raising \texttt{adapt\_delta} further
or running more warmup iterations is the correct response.

\paragraph{Energy (E-BFMI)}\label{energy-e-bfmi}

\begin{Shaded}
\begin{Highlighting}[]
\FunctionTok{mcmc\_nuts\_energy}\NormalTok{(}\FunctionTok{nuts\_params}\NormalTok{(fit\_memory)) }\SpecialCharTok{+}
  \FunctionTok{labs}\NormalTok{(}
    \AttributeTok{title    =} \StringTok{"Energy Transitions: Memory Agent NCP"}\NormalTok{,}
    \AttributeTok{subtitle =} \StringTok{"Marginal and transition energy distributions must align. E{-}BFMI \textgreater{} 0.2 required."}
\NormalTok{  )}
\end{Highlighting}
\end{Shaded}

\begin{center}
\includegraphics[width=0.8\linewidth,height=\textheight,keepaspectratio]{06-MultilevelModels_files/figure-pdf/ch07_memory_energy-1.pdf}
\end{center}

\subsubsection{Posterior Parameter
Correlations}\label{posterior-parameter-correlations-1}

\begin{Shaded}
\begin{Highlighting}[]
\FunctionTok{mcmc\_pairs}\NormalTok{(}
\NormalTok{  draws\_memory,}
  \AttributeTok{pars         =} \FunctionTok{c}\NormalTok{(}\StringTok{"mu\_alpha"}\NormalTok{, }\StringTok{"mu\_beta"}\NormalTok{, }\StringTok{"sigma[1]"}\NormalTok{, }\StringTok{"sigma[2]"}\NormalTok{, }\StringTok{"rho"}\NormalTok{),}
  \AttributeTok{np           =} \FunctionTok{nuts\_params}\NormalTok{(fit\_memory),}
  \AttributeTok{diag\_fun     =} \StringTok{"dens"}\NormalTok{,}
  \AttributeTok{off\_diag\_fun =} \StringTok{"hex"}
\NormalTok{)}
\end{Highlighting}
\end{Shaded}

\begin{center}
\includegraphics[width=0.8\linewidth,height=\textheight,keepaspectratio]{06-MultilevelModels_files/figure-pdf/ch07_memory_posterior_correlations-1.pdf}
\end{center}

The critical panels are \((\sigma_\alpha, \mu_\alpha)\) and
\((\sigma_\beta, \mu_\beta)\). Under the CP biased agent, these showed
triangular wedges --- the Neal's Funnel signature. Under NCP, they
should be roughly elliptical. Any triangular shape here would be
unexpected and would require investigation.

\subsection{Prior-Posterior Update Checks: Quantifying
Learning}\label{prior-posterior-update-checks-quantifying-learning-1}

\begin{Shaded}
\begin{Highlighting}[]
\NormalTok{draws\_df\_memory }\OtherTok{\textless{}{-}} \FunctionTok{as\_draws\_df}\NormalTok{(draws\_memory)}

\CommentTok{\# Extract true generative hyperparameters saved as attributes}
\NormalTok{true\_hp }\OtherTok{\textless{}{-}} \FunctionTok{attr}\NormalTok{(sim\_data\_memory, }\StringTok{"hyperparameters"}\NormalTok{)}

\NormalTok{true\_pop\_mem }\OtherTok{\textless{}{-}} \FunctionTok{tibble}\NormalTok{(}
  \AttributeTok{parameter  =} \FunctionTok{c}\NormalTok{(}\StringTok{"mu\_alpha"}\NormalTok{, }\StringTok{"mu\_beta"}\NormalTok{, }\StringTok{"sigma\_alpha"}\NormalTok{, }\StringTok{"sigma\_beta"}\NormalTok{, }\StringTok{"rho"}\NormalTok{),}
  \AttributeTok{true\_value =} \FunctionTok{c}\NormalTok{(true\_hp}\SpecialCharTok{$}\NormalTok{mu\_alpha, true\_hp}\SpecialCharTok{$}\NormalTok{mu\_beta,}
\NormalTok{                 true\_hp}\SpecialCharTok{$}\NormalTok{sigma\_alpha, true\_hp}\SpecialCharTok{$}\NormalTok{sigma\_beta, true\_hp}\SpecialCharTok{$}\NormalTok{rho)}
\NormalTok{)}

\CommentTok{\# Reshape: split on the underscore immediately before "posterior" or "prior"}
\NormalTok{pp\_update\_mem }\OtherTok{\textless{}{-}}\NormalTok{ draws\_df\_memory }\SpecialCharTok{|\textgreater{}}
\NormalTok{  dplyr}\SpecialCharTok{::}\FunctionTok{select}\NormalTok{(}
    \AttributeTok{mu\_alpha\_posterior  =}\NormalTok{ mu\_alpha,        mu\_alpha\_prior,}
    \AttributeTok{mu\_beta\_posterior   =}\NormalTok{ mu\_beta,         mu\_beta\_prior,}
    \AttributeTok{sigma\_alpha\_posterior =} \StringTok{\textasciigrave{}}\AttributeTok{sigma[1]}\StringTok{\textasciigrave{}}\NormalTok{,    sigma\_alpha\_prior,}
    \AttributeTok{sigma\_beta\_posterior  =} \StringTok{\textasciigrave{}}\AttributeTok{sigma[2]}\StringTok{\textasciigrave{}}\NormalTok{,    sigma\_beta\_prior,}
    \AttributeTok{rho\_posterior       =}\NormalTok{ rho,             rho\_prior}
\NormalTok{  ) }\SpecialCharTok{|\textgreater{}}
  \FunctionTok{pivot\_longer}\NormalTok{(}
    \FunctionTok{everything}\NormalTok{(),}
    \AttributeTok{names\_to  =} \FunctionTok{c}\NormalTok{(}\StringTok{"parameter"}\NormalTok{, }\StringTok{"distribution"}\NormalTok{),}
    \AttributeTok{names\_sep =} \StringTok{"\_(?=posterior|prior)"}
\NormalTok{  )}

\FunctionTok{ggplot}\NormalTok{(pp\_update\_mem, }\FunctionTok{aes}\NormalTok{(}\AttributeTok{x =}\NormalTok{ value)) }\SpecialCharTok{+}
  \FunctionTok{geom\_histogram}\NormalTok{(}\FunctionTok{aes}\NormalTok{(}\AttributeTok{fill =}\NormalTok{ distribution), }\AttributeTok{alpha =} \FloatTok{0.6}\NormalTok{, }\AttributeTok{color =} \ConstantTok{NA}\NormalTok{, }\AttributeTok{bins =} \DecValTok{60}\NormalTok{) }\SpecialCharTok{+}
  \FunctionTok{geom\_vline}\NormalTok{(}
    \AttributeTok{data =}\NormalTok{ true\_pop\_mem,}
    \FunctionTok{aes}\NormalTok{(}\AttributeTok{xintercept =}\NormalTok{ true\_value, }\AttributeTok{linetype =} \StringTok{"True Value"}\NormalTok{),}
    \AttributeTok{color =} \StringTok{"darkred"}\NormalTok{, }\AttributeTok{linewidth =} \DecValTok{1}
\NormalTok{  ) }\SpecialCharTok{+}
  \FunctionTok{facet\_wrap}\NormalTok{(}\SpecialCharTok{\textasciitilde{}}\NormalTok{parameter, }\AttributeTok{scales =} \StringTok{"free"}\NormalTok{) }\SpecialCharTok{+}
  \FunctionTok{scale\_fill\_manual}\NormalTok{(}
    \AttributeTok{name   =} \StringTok{"Distribution"}\NormalTok{,}
    \AttributeTok{values =} \FunctionTok{c}\NormalTok{(}\StringTok{"prior"} \OtherTok{=} \StringTok{"gray70"}\NormalTok{, }\StringTok{"posterior"} \OtherTok{=} \StringTok{"steelblue"}\NormalTok{)}
\NormalTok{  ) }\SpecialCharTok{+}
  \FunctionTok{scale\_linetype\_manual}\NormalTok{(}
    \AttributeTok{name   =} \StringTok{"Ground Truth"}\NormalTok{,}
    \AttributeTok{values =} \FunctionTok{c}\NormalTok{(}\StringTok{"True Value"} \OtherTok{=} \StringTok{"dashed"}\NormalTok{)}
\NormalTok{  ) }\SpecialCharTok{+}
  \FunctionTok{labs}\NormalTok{(}
    \AttributeTok{title =} \StringTok{"Prior{-}Posterior Update Checks: Memory Agent"}\NormalTok{,}
    \AttributeTok{x     =} \StringTok{"Parameter Value"}\NormalTok{,}
    \AttributeTok{y     =} \StringTok{"Count"}
\NormalTok{  )}
\end{Highlighting}
\end{Shaded}

\begin{center}
\includegraphics[width=0.8\linewidth,height=\textheight,keepaspectratio]{06-MultilevelModels_files/figure-pdf/ch07_memory_pp_update-1.pdf}
\end{center}

A healthy update shows a posterior significantly narrower than the prior
and comfortably contained within the prior's mass. The \texttt{rho}
panel carries the central new scientific information. If the posterior
on \(\rho\) is essentially indistinguishable from the LKJ(2) prior, the
data do not contain sufficient information to identify the
population-level correlation --- this can happen with small \(J\) or
short trial sequences. In that case, report the parameter as
unidentified, not the prior as the conclusion. A clearly narrowed
posterior constitutes genuine evidence about the cognitive architecture
of the population.

\subsection{Posterior Predictive Checks (PPC): Generative
Adequacy}\label{posterior-predictive-checks-ppc-generative-adequacy-1}

\begin{Shaded}
\begin{Highlighting}[]
\NormalTok{y\_obs\_mem     }\OtherTok{\textless{}{-}} \FunctionTok{as.numeric}\NormalTok{(stan\_data\_memory}\SpecialCharTok{$}\NormalTok{y)}
\NormalTok{agent\_idx\_mem }\OtherTok{\textless{}{-}} \FunctionTok{as.numeric}\NormalTok{(stan\_data\_memory}\SpecialCharTok{$}\NormalTok{subj\_id)}
\NormalTok{y\_rep\_mem     }\OtherTok{\textless{}{-}}\NormalTok{ fit\_memory}\SpecialCharTok{$}\FunctionTok{draws}\NormalTok{(}\StringTok{"y\_post\_rep"}\NormalTok{, }\AttributeTok{format =} \StringTok{"matrix"}\NormalTok{)}

\NormalTok{n\_sample\_mem       }\OtherTok{\textless{}{-}} \DecValTok{12}
\NormalTok{sampled\_agents\_mem }\OtherTok{\textless{}{-}} \FunctionTok{sample}\NormalTok{(}\FunctionTok{unique}\NormalTok{(agent\_idx\_mem), n\_sample\_mem)}
\NormalTok{keep\_mask\_mem      }\OtherTok{\textless{}{-}}\NormalTok{ agent\_idx\_mem }\SpecialCharTok{\%in\%}\NormalTok{ sampled\_agents\_mem}

\NormalTok{prop\_right }\OtherTok{\textless{}{-}} \ControlFlowTok{function}\NormalTok{(x) }\FunctionTok{mean}\NormalTok{(x)}

\FunctionTok{ppc\_stat\_grouped}\NormalTok{(}
  \AttributeTok{y     =}\NormalTok{ y\_obs\_mem[keep\_mask\_mem],}
  \AttributeTok{yrep  =}\NormalTok{ y\_rep\_mem[, keep\_mask\_mem],}
  \AttributeTok{group =}\NormalTok{ agent\_idx\_mem[keep\_mask\_mem],}
  \AttributeTok{stat  =} \StringTok{"prop\_right"}
\NormalTok{) }\SpecialCharTok{+}
  \FunctionTok{labs}\NormalTok{(}
    \AttributeTok{title    =} \StringTok{"Posterior Predictive Checks: Agent{-}Level Choice Rates (Subset)"}\NormalTok{,}
    \AttributeTok{subtitle =} \FunctionTok{paste}\NormalTok{(}\StringTok{"Showing"}\NormalTok{, n\_sample\_mem, }\StringTok{"randomly sampled agents. Dark line: Observed. Histogram: Predicted."}\NormalTok{),}
    \AttributeTok{x        =} \StringTok{"Proportion of \textquotesingle{}Right\textquotesingle{} Choices"}
\NormalTok{  )}
\end{Highlighting}
\end{Shaded}

\begin{center}
\includegraphics[width=0.8\linewidth,height=\textheight,keepaspectratio]{06-MultilevelModels_files/figure-pdf/ch07_memory_ppc-1.pdf}
\end{center}

\subsubsection{Calibration and
Recovery}\label{calibration-and-recovery-1}

We have proven the model fits the data and generates plausible
predictions. Now we must prove the parameters we extract are actually
the true parameters that generated the data. For a bivariate model,
recovery must be assessed at both levels: all five population
parameters, and the individual \((\alpha_j, \beta_j)\) pairs jointly.

\paragraph{Individual-Level Recovery}\label{individual-level-recovery-1}

\begin{Shaded}
\begin{Highlighting}[]
\NormalTok{recovery\_alpha }\OtherTok{\textless{}{-}}\NormalTok{ draws\_memory }\SpecialCharTok{|\textgreater{}}
  \FunctionTok{subset\_draws}\NormalTok{(}\AttributeTok{variable =} \StringTok{"alpha"}\NormalTok{) }\SpecialCharTok{|\textgreater{}}
  \FunctionTok{summarise\_draws}\NormalTok{() }\SpecialCharTok{|\textgreater{}}
  \FunctionTok{mutate}\NormalTok{(}\AttributeTok{agent\_id =} \FunctionTok{as.integer}\NormalTok{(}\FunctionTok{str\_extract}\NormalTok{(variable, }\StringTok{"(?\textless{}=}\SpecialCharTok{\textbackslash{}\textbackslash{}}\StringTok{[)}\SpecialCharTok{\textbackslash{}\textbackslash{}}\StringTok{d+(?=}\SpecialCharTok{\textbackslash{}\textbackslash{}}\StringTok{])"}\NormalTok{)))}

\NormalTok{recovery\_beta }\OtherTok{\textless{}{-}}\NormalTok{ draws\_memory }\SpecialCharTok{|\textgreater{}}
  \FunctionTok{subset\_draws}\NormalTok{(}\AttributeTok{variable =} \StringTok{"beta\_mem"}\NormalTok{) }\SpecialCharTok{|\textgreater{}}
  \FunctionTok{summarise\_draws}\NormalTok{() }\SpecialCharTok{|\textgreater{}}
  \FunctionTok{mutate}\NormalTok{(}\AttributeTok{agent\_id =} \FunctionTok{as.integer}\NormalTok{(}\FunctionTok{str\_extract}\NormalTok{(variable, }\StringTok{"(?\textless{}=}\SpecialCharTok{\textbackslash{}\textbackslash{}}\StringTok{[)}\SpecialCharTok{\textbackslash{}\textbackslash{}}\StringTok{d+(?=}\SpecialCharTok{\textbackslash{}\textbackslash{}}\StringTok{])"}\NormalTok{)))}

\NormalTok{true\_agent\_params\_mem }\OtherTok{\textless{}{-}}\NormalTok{ sim\_data\_memory }\SpecialCharTok{|\textgreater{}}
  \FunctionTok{distinct}\NormalTok{(agent\_id, n\_trials, true\_alpha, true\_beta)}

\NormalTok{recovery\_mem\_df }\OtherTok{\textless{}{-}}\NormalTok{ true\_agent\_params\_mem }\SpecialCharTok{|\textgreater{}}
  \FunctionTok{left\_join}\NormalTok{(recovery\_alpha }\SpecialCharTok{|\textgreater{}}\NormalTok{ dplyr}\SpecialCharTok{::}\FunctionTok{select}\NormalTok{(agent\_id, }\AttributeTok{alpha\_est =}\NormalTok{ median, }\AttributeTok{alpha\_q5 =}\NormalTok{ q5, }\AttributeTok{alpha\_q95 =}\NormalTok{ q95), }\AttributeTok{by =} \StringTok{"agent\_id"}\NormalTok{) }\SpecialCharTok{|\textgreater{}}
  \FunctionTok{left\_join}\NormalTok{(recovery\_beta  }\SpecialCharTok{|\textgreater{}}\NormalTok{ dplyr}\SpecialCharTok{::}\FunctionTok{select}\NormalTok{(agent\_id, }\AttributeTok{beta\_est  =}\NormalTok{ median, }\AttributeTok{beta\_q5  =}\NormalTok{ q5, }\AttributeTok{beta\_q95  =}\NormalTok{ q95),  }\AttributeTok{by =} \StringTok{"agent\_id"}\NormalTok{)}

\NormalTok{p\_rec\_alpha }\OtherTok{\textless{}{-}} \FunctionTok{ggplot}\NormalTok{(recovery\_mem\_df, }\FunctionTok{aes}\NormalTok{(}\AttributeTok{x =}\NormalTok{ true\_alpha, }\AttributeTok{y =}\NormalTok{ alpha\_est)) }\SpecialCharTok{+}
  \FunctionTok{geom\_abline}\NormalTok{(}\AttributeTok{intercept =} \DecValTok{0}\NormalTok{, }\AttributeTok{slope =} \DecValTok{1}\NormalTok{, }\AttributeTok{linetype =} \StringTok{"dashed"}\NormalTok{, }\AttributeTok{color =} \StringTok{"gray50"}\NormalTok{) }\SpecialCharTok{+}
  \FunctionTok{geom\_errorbar}\NormalTok{(}\FunctionTok{aes}\NormalTok{(}\AttributeTok{ymin =}\NormalTok{ alpha\_q5, }\AttributeTok{ymax =}\NormalTok{ alpha\_q95), }\AttributeTok{width =} \DecValTok{0}\NormalTok{, }\AttributeTok{alpha =} \FloatTok{0.5}\NormalTok{, }\AttributeTok{color =} \StringTok{"steelblue"}\NormalTok{) }\SpecialCharTok{+}
  \FunctionTok{geom\_point}\NormalTok{(}\FunctionTok{aes}\NormalTok{(}\AttributeTok{size =}\NormalTok{ n\_trials), }\AttributeTok{color =} \StringTok{"midnightblue"}\NormalTok{, }\AttributeTok{alpha =} \FloatTok{0.8}\NormalTok{) }\SpecialCharTok{+}
  \FunctionTok{labs}\NormalTok{(}
    \AttributeTok{title =} \FunctionTok{expression}\NormalTok{(}\StringTok{"Recovery: Baseline Bias "} \SpecialCharTok{*}\NormalTok{ alpha[j]),}
    \AttributeTok{x =} \StringTok{"True alpha (logit scale)"}\NormalTok{, }\AttributeTok{y =} \StringTok{"Posterior Median"}\NormalTok{, }\AttributeTok{size =} \StringTok{"N trials"}
\NormalTok{  )}

\NormalTok{p\_rec\_beta }\OtherTok{\textless{}{-}} \FunctionTok{ggplot}\NormalTok{(recovery\_mem\_df, }\FunctionTok{aes}\NormalTok{(}\AttributeTok{x =}\NormalTok{ true\_beta, }\AttributeTok{y =}\NormalTok{ beta\_est)) }\SpecialCharTok{+}
  \FunctionTok{geom\_abline}\NormalTok{(}\AttributeTok{intercept =} \DecValTok{0}\NormalTok{, }\AttributeTok{slope =} \DecValTok{1}\NormalTok{, }\AttributeTok{linetype =} \StringTok{"dashed"}\NormalTok{, }\AttributeTok{color =} \StringTok{"gray50"}\NormalTok{) }\SpecialCharTok{+}
  \FunctionTok{geom\_errorbar}\NormalTok{(}\FunctionTok{aes}\NormalTok{(}\AttributeTok{ymin =}\NormalTok{ beta\_q5, }\AttributeTok{ymax =}\NormalTok{ beta\_q95), }\AttributeTok{width =} \DecValTok{0}\NormalTok{, }\AttributeTok{alpha =} \FloatTok{0.5}\NormalTok{, }\AttributeTok{color =} \StringTok{"steelblue"}\NormalTok{) }\SpecialCharTok{+}
  \FunctionTok{geom\_point}\NormalTok{(}\FunctionTok{aes}\NormalTok{(}\AttributeTok{size =}\NormalTok{ n\_trials), }\AttributeTok{color =} \StringTok{"midnightblue"}\NormalTok{, }\AttributeTok{alpha =} \FloatTok{0.8}\NormalTok{) }\SpecialCharTok{+}
  \FunctionTok{labs}\NormalTok{(}
    \AttributeTok{title =} \FunctionTok{expression}\NormalTok{(}\StringTok{"Recovery: Memory Sensitivity "} \SpecialCharTok{*}\NormalTok{ beta[j]),}
    \AttributeTok{x =} \StringTok{"True beta"}\NormalTok{, }\AttributeTok{y =} \StringTok{"Posterior Median"}\NormalTok{, }\AttributeTok{size =} \StringTok{"N trials"}
\NormalTok{  )}

\NormalTok{p\_rec\_alpha }\SpecialCharTok{+}\NormalTok{ p\_rec\_beta }\SpecialCharTok{+} \FunctionTok{plot\_layout}\NormalTok{(}\AttributeTok{guides =} \StringTok{"collect"}\NormalTok{)}
\end{Highlighting}
\end{Shaded}

\begin{center}
\includegraphics[width=0.8\linewidth,height=\textheight,keepaspectratio]{06-MultilevelModels_files/figure-pdf/ch07_memory_recovery-1.pdf}
\end{center}

Recovery for \(\beta_j\) is typically noisier than for \(\alpha_j\) at
any given trial count. The memory signal \(\text{logit}(m_{ij})\) starts
at zero (the \(\text{logit}(0.5) = 0\) initialization) and becomes
informative only as the cumulative opponent history diverges from 0.5.
Agents with few trials show stronger shrinkage toward
\(\mu_\beta = 1.5\) --- not toward zero, but toward the estimated
population mean, which is the appropriate regularization given what the
rest of the cohort tells us about the population.

\subsubsection{Prior sensitivity}\label{prior-sensitivity-1}

\begin{Shaded}
\begin{Highlighting}[]
\NormalTok{sens\_results\_mem }\OtherTok{\textless{}{-}} \FunctionTok{powerscale\_sensitivity}\NormalTok{(}
\NormalTok{  fit\_memory,}
  \AttributeTok{variable =} \FunctionTok{c}\NormalTok{(}\StringTok{"mu\_alpha"}\NormalTok{, }\StringTok{"mu\_beta"}\NormalTok{, }\StringTok{"sigma[1]"}\NormalTok{, }\StringTok{"sigma[2]"}\NormalTok{, }\StringTok{"rho"}\NormalTok{)}
\NormalTok{)}
\FunctionTok{cat}\NormalTok{(}\StringTok{"=== PRIOR SENSITIVITY DIAGNOSTICS: MEMORY AGENT ===}\SpecialCharTok{\textbackslash{}n}\StringTok{"}\NormalTok{)}
\end{Highlighting}
\end{Shaded}

\begin{verbatim}
=== PRIOR SENSITIVITY DIAGNOSTICS: MEMORY AGENT ===
\end{verbatim}

\begin{Shaded}
\begin{Highlighting}[]
\FunctionTok{print}\NormalTok{(sens\_results\_mem)}
\end{Highlighting}
\end{Shaded}

\begin{verbatim}
Sensitivity based on cjs_dist
Prior selection: all priors
Likelihood selection: all data

 variable prior likelihood                     diagnosis
 mu_alpha 0.033      0.075                             -
  mu_beta 0.058      0.076 potential prior-data conflict
 sigma[1] 0.361      0.208 potential prior-data conflict
 sigma[2] 0.517      0.124 potential prior-data conflict
      rho 0.034      0.197                             -
\end{verbatim}

\begin{Shaded}
\begin{Highlighting}[]
\NormalTok{ps\_seq\_mem }\OtherTok{\textless{}{-}} \FunctionTok{powerscale\_sequence}\NormalTok{(}
\NormalTok{  fit\_memory,}
  \AttributeTok{variable =} \FunctionTok{c}\NormalTok{(}\StringTok{"mu\_alpha"}\NormalTok{, }\StringTok{"mu\_beta"}\NormalTok{, }\StringTok{"sigma[1]"}\NormalTok{, }\StringTok{"sigma[2]"}\NormalTok{, }\StringTok{"rho"}\NormalTok{)}
\NormalTok{)}

\FunctionTok{powerscale\_plot\_dens}\NormalTok{(ps\_seq\_mem) }\SpecialCharTok{+}
  \FunctionTok{labs}\NormalTok{(}
    \AttributeTok{title    =} \StringTok{"Prior Sensitivity: Density Perturbation (Memory Agent)"}\NormalTok{,}
    \AttributeTok{subtitle =} \StringTok{"Does scaling the prior or likelihood drastically alter our beliefs about the population?"}
\NormalTok{  )}
\end{Highlighting}
\end{Shaded}

\begin{center}
\includegraphics[width=0.8\linewidth,height=\textheight,keepaspectratio]{06-MultilevelModels_files/figure-pdf/ch07_memory_prior_sensitivity-1.pdf}
\end{center}

\begin{Shaded}
\begin{Highlighting}[]
\FunctionTok{powerscale\_plot\_quantities}\NormalTok{(ps\_seq\_mem) }\SpecialCharTok{+}
  \FunctionTok{labs}\NormalTok{(}
    \AttributeTok{title    =} \StringTok{"Prior Sensitivity: Parameter Estimates (Memory Agent)"}\NormalTok{,}
    \AttributeTok{subtitle =} \StringTok{"Horizontal lines: baseline posterior bounds. Deviations indicate sensitivity."}
\NormalTok{  )}
\end{Highlighting}
\end{Shaded}

\begin{center}
\includegraphics[width=0.8\linewidth,height=\textheight,keepaspectratio]{06-MultilevelModels_files/figure-pdf/ch07_memory_prior_sensitivity-2.pdf}
\end{center}

Pay particular attention to the \texttt{rho} panel. Because the
correlation is identified only from across-agent covariation, it has
fewer effective observations than the marginal population parameters. If
\texttt{rho} shows strong sensitivity to prior scaling, the data contain
limited information about the correlation structure and the LKJ prior is
doing meaningful regularization.

\subsubsection{Simulation-Based Calibration Checks
(SBC)}\label{simulation-based-calibration-checks-sbc-1}

Parameter recovery on a single dataset is anecdotal. We must prove that
the sampler is globally unbiased across the entire prior space.

\begin{Shaded}
\begin{Highlighting}[]
\NormalTok{sbc\_memory\_path }\OtherTok{\textless{}{-}}\NormalTok{ here}\SpecialCharTok{::}\FunctionTok{here}\NormalTok{(}\StringTok{"simmodels"}\NormalTok{, }\StringTok{"ch07\_sbc\_memory.rds"}\NormalTok{)}

\NormalTok{generate\_sbc\_dataset\_memory }\OtherTok{\textless{}{-}} \ControlFlowTok{function}\NormalTok{(}\AttributeTok{n\_agents =} \DecValTok{50}\NormalTok{, }\AttributeTok{n\_trials\_min =} \DecValTok{80}\NormalTok{, }\AttributeTok{n\_trials\_max =} \DecValTok{240}\NormalTok{) \{}
  \CommentTok{\# Use the same reversal schedule as the main fit.}
  \CommentTok{\# opponent\_reversal is available in the chapter environment from Edit 6.16.}
\NormalTok{  sim  }\OtherTok{\textless{}{-}} \FunctionTok{simulate\_hierarchical\_memory}\NormalTok{(}
    \AttributeTok{n\_agents         =}\NormalTok{ n\_agents,}
    \AttributeTok{n\_trials\_min     =}\NormalTok{ n\_trials\_min,}
    \AttributeTok{n\_trials\_max     =}\NormalTok{ n\_trials\_max,}
    \AttributeTok{opponent\_choices =}\NormalTok{ opponent\_reversal}
\NormalTok{  )}
\NormalTok{  hp   }\OtherTok{\textless{}{-}} \FunctionTok{attr}\NormalTok{(sim, }\StringTok{"hyperparameters"}\NormalTok{)}

  \FunctionTok{list}\NormalTok{(}
    \AttributeTok{variables =} \FunctionTok{list}\NormalTok{(}
      \AttributeTok{mu\_alpha   =}\NormalTok{ hp}\SpecialCharTok{$}\NormalTok{mu\_alpha,}
      \AttributeTok{mu\_beta    =}\NormalTok{ hp}\SpecialCharTok{$}\NormalTok{mu\_beta,}
      \StringTok{\textasciigrave{}}\AttributeTok{sigma[1]}\StringTok{\textasciigrave{}} \OtherTok{=}\NormalTok{ hp}\SpecialCharTok{$}\NormalTok{sigma\_alpha,}
      \StringTok{\textasciigrave{}}\AttributeTok{sigma[2]}\StringTok{\textasciigrave{}} \OtherTok{=}\NormalTok{ hp}\SpecialCharTok{$}\NormalTok{sigma\_beta,}
      \AttributeTok{rho        =}\NormalTok{ hp}\SpecialCharTok{$}\NormalTok{rho,}
      \AttributeTok{alpha      =}\NormalTok{ sim }\SpecialCharTok{|\textgreater{}} \FunctionTok{distinct}\NormalTok{(agent\_id, true\_alpha) }\SpecialCharTok{|\textgreater{}} \FunctionTok{arrange}\NormalTok{(agent\_id) }\SpecialCharTok{|\textgreater{}} \FunctionTok{pull}\NormalTok{(true\_alpha),}
      \AttributeTok{beta\_mem   =}\NormalTok{ sim }\SpecialCharTok{|\textgreater{}} \FunctionTok{distinct}\NormalTok{(agent\_id, true\_beta)  }\SpecialCharTok{|\textgreater{}} \FunctionTok{arrange}\NormalTok{(agent\_id) }\SpecialCharTok{|\textgreater{}} \FunctionTok{pull}\NormalTok{(true\_beta)}
\NormalTok{    ),}
    \AttributeTok{generated =} \FunctionTok{prep\_stan\_data\_memory}\NormalTok{(sim)}
\NormalTok{  )}
\NormalTok{\}}

\NormalTok{generator\_memory }\OtherTok{\textless{}{-}} \FunctionTok{SBC\_generator\_function}\NormalTok{(}
\NormalTok{  generate\_sbc\_dataset\_memory,}
  \AttributeTok{n\_agents =} \DecValTok{50}\NormalTok{, }\AttributeTok{n\_trials\_min =} \DecValTok{30}\NormalTok{, }\AttributeTok{n\_trials\_max =} \DecValTok{120}
\NormalTok{)}

\NormalTok{backend\_memory }\OtherTok{\textless{}{-}} \FunctionTok{SBC\_backend\_cmdstan\_sample}\NormalTok{(}
\NormalTok{  mod\_memory\_ncp,}
  \AttributeTok{iter\_warmup =} \DecValTok{1000}\NormalTok{, }\AttributeTok{iter\_sampling =} \DecValTok{1000}\NormalTok{, }\AttributeTok{chains =} \DecValTok{1}\NormalTok{, }\AttributeTok{refresh =} \DecValTok{0}
\NormalTok{)}

\ControlFlowTok{if}\NormalTok{ (regenerate\_simulations }\SpecialCharTok{||} \SpecialCharTok{!}\FunctionTok{file.exists}\NormalTok{(sbc\_memory\_path)) \{}
\NormalTok{  sbc\_datasets\_memory }\OtherTok{\textless{}{-}} \FunctionTok{generate\_datasets}\NormalTok{(generator\_memory, }\DecValTok{100}\NormalTok{)}
  \FunctionTok{plan}\NormalTok{(sequential)}
  \CommentTok{\#future::plan(future::multisession, workers = 4)}
\NormalTok{  sbc\_results\_memory  }\OtherTok{\textless{}{-}} \FunctionTok{compute\_SBC}\NormalTok{(}\AttributeTok{datasets =}\NormalTok{ sbc\_datasets\_memory, }\AttributeTok{backend =}\NormalTok{ backend\_memory,}
    \AttributeTok{keep\_fits =} \ConstantTok{FALSE}\NormalTok{)}
  \FunctionTok{saveRDS}\NormalTok{(}
    \FunctionTok{list}\NormalTok{(}\AttributeTok{datasets =}\NormalTok{ sbc\_datasets\_memory, }\AttributeTok{results =}\NormalTok{ sbc\_results\_memory),}
\NormalTok{    sbc\_memory\_path}
\NormalTok{  )}
\NormalTok{\} }\ControlFlowTok{else}\NormalTok{ \{}
\NormalTok{  sbc\_mem\_obj         }\OtherTok{\textless{}{-}} \FunctionTok{readRDS}\NormalTok{(sbc\_memory\_path)}
\NormalTok{  sbc\_datasets\_memory }\OtherTok{\textless{}{-}}\NormalTok{ sbc\_mem\_obj}\SpecialCharTok{$}\NormalTok{datasets}
\NormalTok{  sbc\_results\_memory  }\OtherTok{\textless{}{-}}\NormalTok{ sbc\_mem\_obj}\SpecialCharTok{$}\NormalTok{results}
\NormalTok{\}}

\FunctionTok{plot\_rank\_hist}\NormalTok{(sbc\_results\_memory, }\AttributeTok{variables =} \FunctionTok{c}\NormalTok{(}\StringTok{"mu\_alpha"}\NormalTok{, }\StringTok{"mu\_beta"}\NormalTok{, }\StringTok{"sigma[1]"}\NormalTok{, }\StringTok{"sigma[2]"}\NormalTok{, }\StringTok{"rho"}\NormalTok{)) }\SpecialCharTok{+}
  \FunctionTok{labs}\NormalTok{(}\AttributeTok{title =} \StringTok{"SBC Rank Histograms: Memory Agent NCP"}\NormalTok{)}
\end{Highlighting}
\end{Shaded}

\begin{center}
\includegraphics[width=0.8\linewidth,height=\textheight,keepaspectratio]{06-MultilevelModels_files/figure-pdf/ch07_sbc_memory-1.pdf}
\end{center}

\begin{Shaded}
\begin{Highlighting}[]
\FunctionTok{plot\_ecdf\_diff}\NormalTok{(sbc\_results\_memory, }\AttributeTok{variables =} \FunctionTok{c}\NormalTok{(}\StringTok{"mu\_alpha"}\NormalTok{, }\StringTok{"mu\_beta"}\NormalTok{, }\StringTok{"sigma[1]"}\NormalTok{, }\StringTok{"sigma[2]"}\NormalTok{, }\StringTok{"rho"}\NormalTok{)) }\SpecialCharTok{+}
  \FunctionTok{labs}\NormalTok{(}
    \AttributeTok{title    =} \StringTok{"SBC ECDF Difference: Memory Agent NCP"}\NormalTok{,}
    \AttributeTok{subtitle =} \StringTok{"Deviation from the diagonal indicates miscalibration."}
\NormalTok{  )}
\end{Highlighting}
\end{Shaded}

\begin{center}
\includegraphics[width=0.8\linewidth,height=\textheight,keepaspectratio]{06-MultilevelModels_files/figure-pdf/ch07_sbc_memory-2.pdf}
\end{center}

\begin{Shaded}
\begin{Highlighting}[]
\CommentTok{\# 1. Extract the lightweight stats dataframe from the SBC object}
\NormalTok{sbc\_stats\_df }\OtherTok{\textless{}{-}}\NormalTok{ sbc\_results\_memory}\SpecialCharTok{$}\NormalTok{stats}

\CommentTok{\# 2. Filter for the population{-}level parameters you care about}
\CommentTok{\# (e.g., mu\_alpha, mu\_beta, rho, or even sigma[1])}
\NormalTok{recovery\_df }\OtherTok{\textless{}{-}}\NormalTok{ sbc\_stats\_df }\SpecialCharTok{|\textgreater{}}
  \FunctionTok{filter}\NormalTok{(variable }\SpecialCharTok{\%in\%} \FunctionTok{c}\NormalTok{(}\StringTok{"mu\_alpha"}\NormalTok{, }\StringTok{"mu\_beta"}\NormalTok{, }\StringTok{"sigma[1]"}\NormalTok{, }\StringTok{"sigma[2]"}\NormalTok{, }\StringTok{"rho"}\NormalTok{))}

\CommentTok{\# 3. Create the Parameter Recovery Identity Plot}
\NormalTok{p\_sbc\_recovery }\OtherTok{\textless{}{-}} \FunctionTok{ggplot}\NormalTok{(recovery\_df, }\FunctionTok{aes}\NormalTok{(}\AttributeTok{x =}\NormalTok{ simulated\_value, }\AttributeTok{y =}\NormalTok{ median)) }\SpecialCharTok{+}
  \CommentTok{\# The Identity Line (Perfect Recovery)}
  \FunctionTok{geom\_abline}\NormalTok{(}\AttributeTok{intercept =} \DecValTok{0}\NormalTok{, }\AttributeTok{slope =} \DecValTok{1}\NormalTok{, }\AttributeTok{linetype =} \StringTok{"dashed"}\NormalTok{, }\AttributeTok{color =} \StringTok{"gray50"}\NormalTok{, }\AttributeTok{linewidth =} \DecValTok{1}\NormalTok{) }\SpecialCharTok{+}
  
  \CommentTok{\# The Posterior Estimates with 90\% Credible Intervals}
  \FunctionTok{geom\_errorbar}\NormalTok{(}\FunctionTok{aes}\NormalTok{(}\AttributeTok{ymin =}\NormalTok{ q5, }\AttributeTok{ymax =}\NormalTok{ q95), }\AttributeTok{width =} \DecValTok{0}\NormalTok{, }\AttributeTok{alpha =} \FloatTok{0.3}\NormalTok{, }\AttributeTok{color =} \StringTok{"steelblue"}\NormalTok{) }\SpecialCharTok{+}
  \FunctionTok{geom\_point}\NormalTok{(}\AttributeTok{color =} \StringTok{"midnightblue"}\NormalTok{, }\AttributeTok{alpha =} \FloatTok{0.7}\NormalTok{, }\AttributeTok{size =} \DecValTok{2}\NormalTok{) }\SpecialCharTok{+}
  
  \CommentTok{\# Facet by parameter}
  \FunctionTok{facet\_wrap}\NormalTok{(}\SpecialCharTok{\textasciitilde{}}\NormalTok{variable, }\AttributeTok{scales =} \StringTok{"free"}\NormalTok{) }\SpecialCharTok{+}
  
  \FunctionTok{labs}\NormalTok{(}
    \AttributeTok{title =} \StringTok{"Parameter Recovery (Extracted Directly from SBC)"}\NormalTok{,}
    \AttributeTok{subtitle =} \StringTok{"Points represent the posterior median. Bars are 90\% Credible Intervals."}\NormalTok{,}
    \AttributeTok{x =} \StringTok{"True Generative Value (Drawn from Prior)"}\NormalTok{,}
    \AttributeTok{y =} \StringTok{"Posterior Median Estimate"}
\NormalTok{  ) }\SpecialCharTok{+}
  \FunctionTok{theme\_classic}\NormalTok{() }\SpecialCharTok{+}
  \FunctionTok{theme}\NormalTok{(}\AttributeTok{strip.background =} \FunctionTok{element\_rect}\NormalTok{(}\AttributeTok{fill =} \StringTok{"gray95"}\NormalTok{, }\AttributeTok{color =} \ConstantTok{NA}\NormalTok{))}

\FunctionTok{print}\NormalTok{(p\_sbc\_recovery)}
\end{Highlighting}
\end{Shaded}

\begin{center}
\includegraphics[width=0.8\linewidth,height=\textheight,keepaspectratio]{06-MultilevelModels_files/figure-pdf/ch07_sbc_memory-3.pdf}
\end{center}

\begin{Shaded}
\begin{Highlighting}[]
\CommentTok{\# 4. Quantify the Recovery (Correlation and RMSE)}
\NormalTok{recovery\_metrics }\OtherTok{\textless{}{-}}\NormalTok{ recovery\_df }\SpecialCharTok{|\textgreater{}}
  \FunctionTok{group\_by}\NormalTok{(variable) }\SpecialCharTok{|\textgreater{}}
  \FunctionTok{summarize}\NormalTok{(}
    \AttributeTok{Pearson\_r =} \FunctionTok{cor}\NormalTok{(simulated\_value, median),}
    \AttributeTok{RMSE =} \FunctionTok{sqrt}\NormalTok{(}\FunctionTok{mean}\NormalTok{((simulated\_value }\SpecialCharTok{{-}}\NormalTok{ median)}\SpecialCharTok{\^{}}\DecValTok{2}\NormalTok{)),}
    \AttributeTok{.groups =} \StringTok{"drop"}
\NormalTok{  )}

\FunctionTok{cat}\NormalTok{(}\StringTok{"=== PARAMETER RECOVERY QUANTIFICATION ===}\SpecialCharTok{\textbackslash{}n}\StringTok{"}\NormalTok{)}
\end{Highlighting}
\end{Shaded}

\begin{verbatim}
=== PARAMETER RECOVERY QUANTIFICATION ===
\end{verbatim}

\begin{Shaded}
\begin{Highlighting}[]
\FunctionTok{print}\NormalTok{(knitr}\SpecialCharTok{::}\FunctionTok{kable}\NormalTok{(recovery\_metrics, }\AttributeTok{digits =} \DecValTok{3}\NormalTok{))}
\end{Highlighting}
\end{Shaded}

\begin{verbatim}


|variable | Pearson_r|  RMSE|
|:--------|---------:|-----:|
|mu_alpha |     0.988| 0.162|
|mu_beta  |     0.977| 0.220|
|rho      |     0.673| 0.334|
|sigma[1] |     0.988| 0.147|
|sigma[2] |     0.969| 0.237|
\end{verbatim}

\begin{Shaded}
\begin{Highlighting}[]
\NormalTok{diags\_memory  }\OtherTok{\textless{}{-}}\NormalTok{ sbc\_results\_memory}\SpecialCharTok{$}\NormalTok{backend\_diagnostics}
\NormalTok{mem\_fail\_rate }\OtherTok{\textless{}{-}} \FunctionTok{mean}\NormalTok{(diags\_memory}\SpecialCharTok{$}\NormalTok{n\_divergent }\SpecialCharTok{\textgreater{}} \DecValTok{0}\NormalTok{) }\SpecialCharTok{*} \DecValTok{100}

\NormalTok{diags\_memory }\SpecialCharTok{|\textgreater{}}
  \FunctionTok{ggplot}\NormalTok{(}\FunctionTok{aes}\NormalTok{(}\AttributeTok{x =}\NormalTok{ n\_divergent)) }\SpecialCharTok{+}
  \FunctionTok{geom\_histogram}\NormalTok{(}\AttributeTok{binwidth =} \DecValTok{1}\NormalTok{, }\AttributeTok{fill =} \StringTok{"darkred"}\NormalTok{, }\AttributeTok{alpha =} \FloatTok{0.8}\NormalTok{) }\SpecialCharTok{+}
  \FunctionTok{labs}\NormalTok{(}
    \AttributeTok{title    =} \StringTok{"Memory NCP Divergence Counts Across 200 SBC Universes"}\NormalTok{,}
    \AttributeTok{subtitle =} \FunctionTok{sprintf}\NormalTok{(}\StringTok{"\%.0f\%\% of universes produced at least one divergence."}\NormalTok{, mem\_fail\_rate),}
    \AttributeTok{x        =} \StringTok{"Number of Divergent Transitions"}\NormalTok{,}
    \AttributeTok{y        =} \StringTok{"Number of SBC Universes"}
\NormalTok{  )}
\end{Highlighting}
\end{Shaded}

\begin{center}
\includegraphics[width=0.8\linewidth,height=\textheight,keepaspectratio]{06-MultilevelModels_files/figure-pdf/ch07_sbc_memory_diagnostics-1.pdf}
\end{center}

\begin{Shaded}
\begin{Highlighting}[]
\CommentTok{\# Does any failure concentrate at small sigma values?}
\NormalTok{posterior}\SpecialCharTok{::}\FunctionTok{as\_draws\_df}\NormalTok{(sbc\_datasets\_memory}\SpecialCharTok{$}\NormalTok{variables) }\SpecialCharTok{|\textgreater{}}
  \FunctionTok{as\_tibble}\NormalTok{() }\SpecialCharTok{|\textgreater{}}
  \FunctionTok{mutate}\NormalTok{(}\AttributeTok{sim\_id =} \FunctionTok{row\_number}\NormalTok{()) }\SpecialCharTok{|\textgreater{}}
  \FunctionTok{left\_join}\NormalTok{(}
\NormalTok{    diags\_memory }\SpecialCharTok{|\textgreater{}} \FunctionTok{mutate}\NormalTok{(}\AttributeTok{sim\_id =} \FunctionTok{row\_number}\NormalTok{()) }\SpecialCharTok{|\textgreater{}}\NormalTok{ dplyr}\SpecialCharTok{::}\FunctionTok{select}\NormalTok{(sim\_id, n\_divergent),}
    \AttributeTok{by =} \StringTok{"sim\_id"}
\NormalTok{  ) }\SpecialCharTok{|\textgreater{}}
  \FunctionTok{mutate}\NormalTok{(}\AttributeTok{had\_divergence =}\NormalTok{ n\_divergent }\SpecialCharTok{\textgreater{}} \DecValTok{0}\NormalTok{) }\SpecialCharTok{|\textgreater{}}
  \FunctionTok{ggplot}\NormalTok{(}\FunctionTok{aes}\NormalTok{(}\AttributeTok{x =} \StringTok{\textasciigrave{}}\AttributeTok{sigma[1]}\StringTok{\textasciigrave{}}\NormalTok{, }\AttributeTok{y =} \StringTok{\textasciigrave{}}\AttributeTok{sigma[2]}\StringTok{\textasciigrave{}}\NormalTok{, }\AttributeTok{color =}\NormalTok{ had\_divergence)) }\SpecialCharTok{+}
  \FunctionTok{geom\_point}\NormalTok{(}\AttributeTok{alpha =} \FloatTok{0.6}\NormalTok{, }\AttributeTok{size =} \FloatTok{1.5}\NormalTok{) }\SpecialCharTok{+}
  \FunctionTok{scale\_color\_manual}\NormalTok{(}
    \AttributeTok{values =} \FunctionTok{c}\NormalTok{(}\StringTok{"FALSE"} \OtherTok{=} \StringTok{"steelblue"}\NormalTok{, }\StringTok{"TRUE"} \OtherTok{=} \StringTok{"darkred"}\NormalTok{),}
    \AttributeTok{labels =} \FunctionTok{c}\NormalTok{(}\StringTok{"FALSE"} \OtherTok{=} \StringTok{"Clean"}\NormalTok{, }\StringTok{"TRUE"} \OtherTok{=} \StringTok{"Divergent"}\NormalTok{),}
    \AttributeTok{name   =} \ConstantTok{NULL}
\NormalTok{  ) }\SpecialCharTok{+}
  \FunctionTok{labs}\NormalTok{(}
    \AttributeTok{title    =} \StringTok{"Memory NCP: Where Do Divergences Occur?"}\NormalTok{,}
    \AttributeTok{subtitle =} \StringTok{"Any divergences should not concentrate at small sigma values (that would indicate residual funnel geometry)."}\NormalTok{,}
    \AttributeTok{x        =} \FunctionTok{expression}\NormalTok{(sigma[alpha]), }\AttributeTok{y =} \FunctionTok{expression}\NormalTok{(sigma[beta])}
\NormalTok{  )}
\end{Highlighting}
\end{Shaded}

\begin{center}
\includegraphics[width=0.8\linewidth,height=\textheight,keepaspectratio]{06-MultilevelModels_files/figure-pdf/ch07_sbc_memory_diagnostics-2.pdf}
\end{center}

The \texttt{rho} rank histogram will typically have wider confidence
bands than the marginal mean and SD histograms at \(N = 200\). This is
expected: \(\rho\) is identified only from covariation across \(J = 50\)
agents, so each SBC universe contributes one effective draw of \(\rho\)
to the calibration check. Mild irregularity within the confidence
envelope is statistical noise. Systematic U-shape or strong asymmetry
would indicate a genuine calibration failure.

\subsubsection{\texorpdfstring{Cognitive Interpretation of
\(\hat{\rho}\)}{Cognitive Interpretation of \textbackslash hat\{\textbackslash rho\}}}\label{cognitive-interpretation-of-hatrho}

The posterior on \(\rho\) is the inferential result that distinguishes
this model from two independent univariate hierarchical models.

\begin{itemize}
\item
  \textbf{\(\rho \approx 0\)}: Baseline tendency and memory-based
  opponent tracking are orthogonal individual differences. Studying one
  tells us nothing about the other.
\item
  \textbf{\(\rho > 0\)}: Agents with a stronger directional baseline
  also tend to weight memory more heavily --- consistent with a general
  exploitation tendency, where both parameters index willingness to
  depart from 50/50 in response to a signal.
\item
  \textbf{\(\rho < 0\)}: Agents with a strong directional baseline tend
  to suppress memory use --- a compensatory account in which prior
  commitment crowds out sensitivity to the opponent's history.
\end{itemize}

None of these interpretations follow from the model alone. The model
provides a calibrated posterior over \(\rho\). Cognitive theory provides
the framework for deciding what that value means.

\subsubsection{LOO for Model Comparison}\label{loo-for-model-comparison}

We retain the LOO-CV score here for the model comparison section at the
end of the chapter.

\begin{Shaded}
\begin{Highlighting}[]
\NormalTok{loo\_memory }\OtherTok{\textless{}{-}} \FunctionTok{loo}\NormalTok{(fit\_memory}\SpecialCharTok{$}\FunctionTok{draws}\NormalTok{(}\StringTok{"log\_lik"}\NormalTok{, }\AttributeTok{format =} \StringTok{"matrix"}\NormalTok{))}
\FunctionTok{print}\NormalTok{(loo\_memory)}
\end{Highlighting}
\end{Shaded}

\begin{verbatim}

Computed from 4000 by 8058 log-likelihood matrix.

         Estimate   SE
elpd_loo  -4233.0 44.3
p_loo        55.3  1.3
looic      8465.9 88.6
------
MCSE of elpd_loo is 0.1.
MCSE and ESS estimates assume independent draws (r_eff=1).

All Pareto k estimates are good (k < 0.7).
See help('pareto-k-diagnostic') for details.
\end{verbatim}

\begin{Shaded}
\begin{Highlighting}[]
\CommentTok{\# Save everything Ch. 8 needs for model comparison.}
\CommentTok{\# Ch. 8 will load this object to avoid refitting.}
\NormalTok{ch07\_exports\_file }\OtherTok{\textless{}{-}} \FunctionTok{here}\NormalTok{(}\StringTok{"simmodels"}\NormalTok{, }\StringTok{"ch07\_memory\_exports\_for\_ch7.rds"}\NormalTok{)}

\FunctionTok{saveRDS}\NormalTok{(}
  \FunctionTok{list}\NormalTok{(}
    \AttributeTok{fit\_memory       =}\NormalTok{ fit\_memory,        }\CommentTok{\# Fitted CmdStanR object}
    \AttributeTok{loo\_memory       =}\NormalTok{ loo\_memory,        }\CommentTok{\# Pre{-}computed LOO score}
    \AttributeTok{stan\_data\_memory =}\NormalTok{ stan\_data\_memory,  }\CommentTok{\# Stan data list (reversal design)}
    \AttributeTok{sim\_data\_memory  =}\NormalTok{ sim\_data\_memory,   }\CommentTok{\# Full trial{-}level tibble}
    \AttributeTok{reversal\_schedule =}\NormalTok{ opponent\_reversal }\CommentTok{\# 240{-}trial opponent sequence}
\NormalTok{  ),}
\NormalTok{  ch07\_exports\_file}
\NormalTok{)}

\FunctionTok{cat}\NormalTok{(}\StringTok{"Ch. 7 exports saved to:"}\NormalTok{, ch07\_exports\_file, }\StringTok{"}\SpecialCharTok{\textbackslash{}n}\StringTok{"}\NormalTok{)}
\end{Highlighting}
\end{Shaded}

\begin{verbatim}
Ch. 7 exports saved to: /Users/au209589/Dropbox/Teaching/AdvancedCognitiveModeling23_book/simmodels/ch07_memory_exports_for_ch7.rds 
\end{verbatim}

\begin{Shaded}
\begin{Highlighting}[]
\FunctionTok{cat}\NormalTok{(}\StringTok{"Load in Ch. 8 with: readRDS(here(\textquotesingle{}simmodels\textquotesingle{}, \textquotesingle{}ch07\_memory\_exports\_for\_ch7.rds\textquotesingle{}))}\SpecialCharTok{\textbackslash{}n}\StringTok{"}\NormalTok{)}
\end{Highlighting}
\end{Shaded}

\begin{verbatim}
Load in Ch. 8 with: readRDS(here('simmodels', 'ch07_memory_exports_for_ch7.rds'))
\end{verbatim}

\section{Conclusion: From One Agent to a
Population}\label{conclusion-from-one-agent-to-a-population}

This chapter extended the Bayesian modeling workflow to multiple
participants. The core ideas are:

\begin{itemize}
\tightlist
\item
  \textbf{Partial pooling} is not a compromise --- it is the
  statistically correct way to combine information across individuals
  with unequal data, and it provably outperforms both extremes (as the
  shrinkage plot showed).
\item
  \textbf{Centered vs.~non-centered parameterization} is a geometric
  choice, not a philosophical one. CP fails when population variance is
  small (Neal's Funnel); NCP fixes the global geometry but can be
  locally inefficient when data are very informative. SBC revealed which
  regime we are in.
\item
  \textbf{The reversal learning design} is essential for the memory
  model. A constant-bias opponent makes \texttt{bias} and \texttt{beta}
  statistically inseparable; the reversal schedule provides the contrast
  that identifies them.
\end{itemize}

\textbf{Coming up in Chapter 8:} We have now validated two models ---
the hierarchical biased agent and the hierarchical memory agent --- and
computed LOO scores for both. Chapter 8 introduces model comparison:
given that both models have passed validation, which one better predicts
held-out data? The answer will use the LOO scores and the
\texttt{loo\_compare()} function from the \texttt{loo} package.

\bookmarksetup{startatroot}

\chapter{Deciding Between Competing Theories: Model
Comparison}\label{deciding-between-competing-theories-model-comparison}

\begin{quote}
\textbf{📍 Where we are in the Bayesian modeling workflow:} Chs. 1--6
built models, validated them, and scaled them to populations. Every
model passed a six-phase quality battery before fitting real data. This
chapter asks the next question: when two validated models both fit the
data, which one better \emph{predicts} data it has never seen? The
answer uses PSIS-LOO cross-validation and requires the validated fits
from Ch. 7. The reversal learning design introduced in Ch. 6 and used in
Ch. 7 is essential here too --- it is the design that makes the two
models distinguishable in the first place.
\end{quote}

By Chapter 7 we had a working Bayesian workflow: specify a generative
model, validate it with SBC and parameter recovery, fit it to data, and
diagnose the posterior geometry, both in the case of single agent and
hierarchical models. We now face a different question: we rarely have
only one possible model of the cognitive processes underlying the
observed behaviors.

\begin{itemize}
\item
  Sometimes we have different models representing competing theoretical
  accounts of cognitive processes. For instance, a Gaussian filter model
  of learning regularities in one's opponent's behavior in the Matching
  Pennies game vs a Theory of Mind model of learning the model used by
  the opponent vs.~an active inference model trying to actively deceive
  the opponent.
\item
  Sometimes we have simpler and more complex versions of the same model:
  should we assume we equally learn from all feedback (using a
  reinforcement learning model) or should we assume different learning
  from positive and negative feedback (using a dual learning rate
  model)?
\end{itemize}

In either case, once we have ensured that all models involved pass their
individual quality checks, how do we decide which better accounts for
the observed data?

There are many ways to compare models, but here we develop a Bayesian
predictive approach consistent with our approach in previous chapters.
The core argument is that a model should be judged not by how well it
explains the data it was \emph{trained} on, but by how well it predicts
data it has \emph{never seen}. We formalize this as the Expected Log
Predictive Density (ELPD), build an efficient computational
approximation via PSIS-LOO, and apply the full workflow to the biased
and memory agent models from Chapter 7. This approach also gives us a
principled way to identify difficult to predict data points (sometimes
called influential observations), which can tell us much about the
limits of our models and different mechanisms at work in different
portions of the data. However, note that focusing on predictive
abilities does not directly equate to model most likely to be true! It
is a useful pragmatic measure and we should keep it in mind.

As in previous chapters, we heavily rely on simulations to understand
these techniques and validate our model comparison procedure. Because we
control the data-generating mechanism, we can evaluate not just whether
model comparison produces a preference, but whether it produces the
\emph{correct} preference --- the model recovery test.

Remember, model comparison is not a fail-safe procedure to determine the
absolute ``truth''. All models are simplifications. Instead, it helps us
identify which of our candidate models provides the most useful or most
plausible account of the data, given our current understanding and the
available evidence. Model comparison should always be part of a larger
principled scientific workflow using models as stepping stones to deeper
understanding, not as final arbiters of truth.

\begin{quote}
CRITICAL PRECONDITION: Model comparison is only scientifically
meaningful if all candidate models have already passed the individual
quality checks established in Chapter 6. If Model A converges perfectly
but Model B fails its SBC calibration or has persistent divergences, an
ELPD comparison is a ``garbage-in, garbage-out'' exercise. We only
compare models that have earned the right to be taken seriously.
\end{quote}

\section{Learning Objectives}\label{learning-objectives-6}

After completing this chapter, you will be able to:

\begin{itemize}
\item
  Explain the failure mode of in-sample fit: Understand why simply
  choosing the model that fits the training data best is misleading and
  leads to the ``memorization'' of noise.
\item
  Understand Cross-Validation: Grasp the core concept of partitioning
  data into training and test sets to estimate generalization
  performance.
\item
  Define ELPD rigorously: State the Expected Log Predictive Density as
  the formal Bayesian target for predictive accuracy.
\item
  Implement PSIS-LOO: Compute ELPD estimates from Stan fits and use
  Pareto-\(\hat{k}\) to identify when the importance-sampling
  approximation is unreliable.
\item
  Characterize influential observations: Link high \(\hat{k}\) values
  back to the cognitive structure of the task, distinguishing between
  person-level outliers and trial-level surprises.
\item
  Perform exact agent-level K-Fold CV: Implement a parallelized
  cross-validation pipeline to test how models generalize to completely
  new participants.
\item
  Diagnose LOO-PIT calibration: Evaluate whether a model's predictive
  uncertainty matches reality using randomized PIT values and the formal
  POT-C test.
\item
  Interpret loo\_compare() and Bayesian Stacking: Reason about the
  \(|elpd\_diff|/se\_diff\) ratio and use stacking weights to find the
  optimal predictive mixture of models.
\item
  Conduct Power Analysis for Model Recovery: Use 2x2 confusion matrices
  and 2D trade-off plots (Trials vs.~Participants) to determine if an
  experimental design is discriminative enough to identify the true
  cognitive mechanism.
\item
  Evaluate Model Calibration via SBC: Move beyond binary ``wins'' to
  assess the probabilistic calibration of model weights using
  Simulation-Based Calibration Checks for Models.
\item
  Identify sequentially dependent data: Recognize when trial-level
  exchangeability is violated and describe the Leave-Future-Out (LFO-CV)
  remedy.
\end{itemize}

\begin{tcolorbox}[enhanced jigsaw, toprule=.15mm, toptitle=1mm, bottomtitle=1mm, leftrule=.75mm, rightrule=.15mm, colframe=quarto-callout-note-color-frame, left=2mm, colback=white, titlerule=0mm, coltitle=black, colbacktitle=quarto-callout-note-color!10!white, opacityback=0, title=\textcolor{quarto-callout-note-color}{\faInfo}\hspace{0.5em}{Historical Context: From Fit to Prediction --- The History of Model
Selection}, arc=.35mm, opacitybacktitle=0.6, breakable, bottomrule=.15mm]

How to choose among competing models is one of the oldest unsolved
problems in statistics. The modern answer --- estimate out-of-sample
predictive accuracy --- is less obvious than it sounds, and the path to
it passed through several influential wrong turns.

\subsubsection{The In-Sample Era: Likelihood and
R²}\label{the-in-sample-era-likelihood-and-ruxb2}

The classical approach was to ask which model assigns the highest
probability to the observed data: the maximum likelihood. This criterion
was the backbone of hypothesis testing (Fisher, 1922; Neyman \& Pearson,
1933) and dominated empirical science throughout the twentieth century.
Its failure mode was known --- a model with enough free parameters can
fit any finite dataset perfectly --- but was treated as a matter of
scientific judgment rather than a formal problem. The coefficient of
determination \(R^2\) shared the same flaw: it increases monotonically
with model complexity and therefore cannot serve as a principled
selection criterion.

\subsubsection{Information Criteria: AIC and
BIC}\label{information-criteria-aic-and-bic}

The first principled solution came from information theory. Akaike
(1974) derived the \textbf{Akaike Information Criterion} (AIC) by asking
a precise question: what is the expected Kullback-Leibler divergence
between the true data-generating distribution and the distribution
implied by a fitted model, evaluated on \emph{new} data from the same
process? The answer, to leading order, is \(-2\log\hat{L} + 2k\), where
\(\hat{L}\) is the maximized likelihood and \(k\) is the number of free
parameters. The penalty term \(2k\) is a correction for the optimism of
in-sample fit: a model that uses \(k\) extra parameters gains
approximately \(2k\) units of log-likelihood purely by fitting noise,
and AIC subtracts this back out.

Schwarz (1978) derived the \textbf{Bayesian Information Criterion} (BIC)
from a different starting point --- approximating the log marginal
likelihood under a diffuse prior --- and obtained a penalty of
\(k \log n\) that grows with sample size and more aggressively penalises
complexity. AIC and BIC dominated applied statistics through the 1980s
and 1990s, but both inherit a key limitation: they are computed on the
training data and require asymptotic approximations whose validity in
small, heterogeneous cognitive datasets is difficult to verify.

\subsubsection{Cross-Validation and the Equivalence
Theorem}\label{cross-validation-and-the-equivalence-theorem}

Stone (1977) proved a result that initially seemed surprising: AIC and
leave-one-out cross-validation (LOO-CV) are \textbf{asymptotically
equivalent}. Cross-validation --- holding out each observation in turn,
refitting on the rest, and evaluating predictive accuracy on the
held-out point --- gives the same answer as AIC in large samples, but
without requiring asymptotic approximations or a count of free
parameters. This equivalence is important because cross-validation is
model-agnostic: it works on any model that can generate predictions,
whether or not its parameter count is well-defined (as is the case for
hierarchical Bayesian models where effective complexity is
data-dependent).

\subsubsection{PSIS-LOO: Making Cross-Validation
Affordable}\label{psis-loo-making-cross-validation-affordable}

The practical obstacle to LOO-CV was computational: for a model fit to
\(N\) observations, exact LOO requires \(N\) separate refits. Vehtari,
Gelman, and Gabry (2017) solved this with \textbf{Pareto-Smoothed
Importance Sampling LOO} (PSIS-LOO). The key insight is that a model fit
to all \(N\) observations already contains, implicitly, the posterior
for each leave-one-out model --- the left-out observation simply
receives zero weight. Importance sampling reweights the full-data
posterior to approximate the leave-one-out posterior, and
Pareto-smoothing of the importance weights stabilises the estimator and,
crucially, provides a per-observation diagnostic (the \(\hat{k}\)
statistic) that flags when the approximation is unreliable. This
diagnostic turned model comparison from a black-box number into an
interpretable workflow: high \(\hat{k}\) observations are not just
nuisances but windows into the structure of the data that the model is
struggling to predict.

The arc from AIC (1974) to PSIS-LOO (2017) is a story about a single
conceptual shift --- from \emph{how well does the model fit what we have
seen?} to \emph{how well does the model predict what we have not yet
seen?} --- being made progressively more rigorous, more computationally
feasible, and more transparent about its own failure modes.

\end{tcolorbox}

\section{Why In-Sample Fit Fails: The Overfitting
Trap}\label{why-in-sample-fit-fails-the-overfitting-trap}

The most naive approach to comparing two models is to ask which fits the
training data better. The problem with this criterion is that it is
systematically biased toward more complex models. A model with enough
free parameters can memorize any finite dataset. But memorization is not
understanding --- a model that has memorized the noise in its training
data will make poor predictions on new data it has never seen.

To make this concrete, we generate data from a simple quadratic curve
with added noise, then fit three polynomial models of increasing
complexity.

\begin{Shaded}
\begin{Highlighting}[]
\FunctionTok{set.seed}\NormalTok{(}\DecValTok{101}\NormalTok{)}
\NormalTok{n      }\OtherTok{\textless{}{-}} \DecValTok{30}
\NormalTok{x      }\OtherTok{\textless{}{-}} \FunctionTok{seq}\NormalTok{(}\SpecialCharTok{{-}}\DecValTok{3}\NormalTok{, }\DecValTok{3}\NormalTok{, }\AttributeTok{length.out =}\NormalTok{ n)}
\NormalTok{y\_true }\OtherTok{\textless{}{-}} \DecValTok{1} \SpecialCharTok{+} \FloatTok{0.5} \SpecialCharTok{*}\NormalTok{ x }\SpecialCharTok{+} \FloatTok{0.8} \SpecialCharTok{*}\NormalTok{ x}\SpecialCharTok{\^{}}\DecValTok{2}
\NormalTok{y\_obs  }\OtherTok{\textless{}{-}}\NormalTok{ y\_true }\SpecialCharTok{+} \FunctionTok{rnorm}\NormalTok{(n, }\DecValTok{0}\NormalTok{, }\FloatTok{2.5}\NormalTok{)}
\NormalTok{df\_fit }\OtherTok{\textless{}{-}} \FunctionTok{tibble}\NormalTok{(}\AttributeTok{x =}\NormalTok{ x, }\AttributeTok{y =}\NormalTok{ y\_obs)}

\NormalTok{fit\_p1  }\OtherTok{\textless{}{-}} \FunctionTok{lm}\NormalTok{(y }\SpecialCharTok{\textasciitilde{}}\NormalTok{ x,           }\AttributeTok{data =}\NormalTok{ df\_fit)}
\NormalTok{fit\_p2  }\OtherTok{\textless{}{-}} \FunctionTok{lm}\NormalTok{(y }\SpecialCharTok{\textasciitilde{}} \FunctionTok{poly}\NormalTok{(x, }\DecValTok{2}\NormalTok{),  }\AttributeTok{data =}\NormalTok{ df\_fit)}
\NormalTok{fit\_p10 }\OtherTok{\textless{}{-}} \FunctionTok{lm}\NormalTok{(y }\SpecialCharTok{\textasciitilde{}} \FunctionTok{poly}\NormalTok{(x, }\DecValTok{10}\NormalTok{), }\AttributeTok{data =}\NormalTok{ df\_fit)}

\NormalTok{x\_pred  }\OtherTok{\textless{}{-}} \FunctionTok{seq}\NormalTok{(}\SpecialCharTok{{-}}\DecValTok{3}\NormalTok{, }\DecValTok{3}\NormalTok{, }\AttributeTok{length.out =} \DecValTok{200}\NormalTok{)}
\NormalTok{pred\_df }\OtherTok{\textless{}{-}} \FunctionTok{tibble}\NormalTok{(}
  \AttributeTok{x    =}\NormalTok{ x\_pred,}
  \AttributeTok{lin  =} \FunctionTok{predict}\NormalTok{(fit\_p1,  }\AttributeTok{newdata =} \FunctionTok{data.frame}\NormalTok{(}\AttributeTok{x =}\NormalTok{ x\_pred)),}
  \AttributeTok{quad =} \FunctionTok{predict}\NormalTok{(fit\_p2,  }\AttributeTok{newdata =} \FunctionTok{data.frame}\NormalTok{(}\AttributeTok{x =}\NormalTok{ x\_pred)),}
  \AttributeTok{high =} \FunctionTok{predict}\NormalTok{(fit\_p10, }\AttributeTok{newdata =} \FunctionTok{data.frame}\NormalTok{(}\AttributeTok{x =}\NormalTok{ x\_pred))}
\NormalTok{) }\SpecialCharTok{|\textgreater{}}
  \FunctionTok{pivot\_longer}\NormalTok{(}\SpecialCharTok{{-}}\NormalTok{x, }\AttributeTok{names\_to =} \StringTok{"model"}\NormalTok{, }\AttributeTok{values\_to =} \StringTok{"y\_pred"}\NormalTok{) }\SpecialCharTok{|\textgreater{}}
  \FunctionTok{mutate}\NormalTok{(}\AttributeTok{model =} \FunctionTok{factor}\NormalTok{(model,}
    \AttributeTok{levels =} \FunctionTok{c}\NormalTok{(}\StringTok{"lin"}\NormalTok{, }\StringTok{"quad"}\NormalTok{, }\StringTok{"high"}\NormalTok{),}
    \AttributeTok{labels =} \FunctionTok{c}\NormalTok{(}\StringTok{"Degree 1 (underfit)"}\NormalTok{, }\StringTok{"Degree 2 (true DGP)"}\NormalTok{,}
               \StringTok{"Degree 10 (overfit)"}\NormalTok{)))}

\FunctionTok{ggplot}\NormalTok{(df\_fit, }\FunctionTok{aes}\NormalTok{(}\AttributeTok{x =}\NormalTok{ x, }\AttributeTok{y =}\NormalTok{ y)) }\SpecialCharTok{+}
  \FunctionTok{geom\_point}\NormalTok{(}\AttributeTok{size =} \DecValTok{2}\NormalTok{, }\AttributeTok{alpha =} \FloatTok{0.8}\NormalTok{) }\SpecialCharTok{+}
  \FunctionTok{geom\_line}\NormalTok{(}\AttributeTok{data =}\NormalTok{ pred\_df, }\FunctionTok{aes}\NormalTok{(}\AttributeTok{y =}\NormalTok{ y\_pred, }\AttributeTok{color =}\NormalTok{ model), }\AttributeTok{linewidth =} \DecValTok{1}\NormalTok{) }\SpecialCharTok{+}
  \FunctionTok{scale\_color\_brewer}\NormalTok{(}\AttributeTok{palette =} \StringTok{"Set1"}\NormalTok{, }\AttributeTok{name =} \StringTok{"Model"}\NormalTok{) }\SpecialCharTok{+}
  \FunctionTok{labs}\NormalTok{(}\AttributeTok{title    =} \StringTok{"The Overfitting Trap: In{-}Sample Fit Favours Complexity"}\NormalTok{,}
       \AttributeTok{subtitle =} \StringTok{"Degree{-}10 polynomial fits training data perfectly but would fail on new data. Degree{-}2 (true DGP) generalizes best."}\NormalTok{,}
       \AttributeTok{x =} \StringTok{"Predictor"}\NormalTok{, }\AttributeTok{y =} \StringTok{"Outcome"}\NormalTok{) }\SpecialCharTok{+}
  \FunctionTok{theme}\NormalTok{(}\AttributeTok{legend.position =} \StringTok{"bottom"}\NormalTok{)}
\end{Highlighting}
\end{Shaded}

\begin{center}
\includegraphics[width=0.8\linewidth,height=\textheight,keepaspectratio]{07-ModelComparison_files/figure-pdf/07-overfitting-demo-1.pdf}
\end{center}

In this plot, we generated data from a quadratic curve with added noise.
A linear model (red) underfits, missing a crucial feature of the data:
the curvature. A degree-10 polynomial (green) overfits the data: it
captures the curvature, but also the noise, deviating here and there to
capture each stray datapoint. The quadratic model (blue) captures the
underlying trend without fitting every wiggle, representing a better
balance.

The error these models make on the training data (e.g.~root mean square
error) monotonically decreases as we increase the complexity of the
model (and the variance explained, \(R^2\), monotonically increases),
and therefore we cannot use it to choose between models. We need a
criterion that rewards genuine generalization: models that capture the
signal without overfitting the noise.

\subsection{The Concept of
Cross-Validation}\label{the-concept-of-cross-validation}

To find this balance between fitting and generalizing, we need a method
to test how well our model performs on data it hasn't seen yet. Since we
usually only have one dataset, we artificially create ``unseen'' data by
hiding parts of our dataset from the model during the training phase.
This is the core concept of \textbf{Cross-Validation (CV)}.

The general process is:

\begin{enumerate}
\def\labelenumi{\arabic{enumi}.}
\item
  Split your data into a \emph{training set} and a \emph{test set},
  respecting the structure in the data. E.g. if you have multiple
  participants, you might want to hold out entire participants for
  testing, rather than random trials, to test generalization to new
  individuals.
\item
  Fit (train) your model using only the training data.
\item
  Evaluate the model's predictions on the hidden test data.
\item
  Repeat this process with different splits and average the performance.
\end{enumerate}

\subsubsection{K-Fold Cross-Validation}\label{k-fold-cross-validation}

In K-Fold cross-validation, we divide the data into \(K\) equally sized
subsets (folds). We train the model on \(K-1\) folds and test it on the
remaining fold, repeating this until every fold has served as the test
set exactly once.

\begin{Shaded}
\begin{Highlighting}[]
\CommentTok{\# Visualization of k{-}fold cross{-}validation}
\NormalTok{n\_data }\OtherTok{\textless{}{-}} \DecValTok{20}
\NormalTok{cv\_data }\OtherTok{\textless{}{-}} \FunctionTok{tibble}\NormalTok{(}\AttributeTok{index =} \DecValTok{1}\SpecialCharTok{:}\NormalTok{n\_data, }\AttributeTok{value =} \FunctionTok{rnorm}\NormalTok{(n\_data))}
\NormalTok{cv\_data}\SpecialCharTok{$}\NormalTok{fold }\OtherTok{\textless{}{-}} \FunctionTok{sample}\NormalTok{(}\FunctionTok{rep}\NormalTok{(}\DecValTok{1}\SpecialCharTok{:}\DecValTok{5}\NormalTok{, }\AttributeTok{length.out =}\NormalTok{ n\_data))}

\NormalTok{cv\_viz\_data }\OtherTok{\textless{}{-}} \FunctionTok{tibble}\NormalTok{(}
  \AttributeTok{iteration =} \FunctionTok{rep}\NormalTok{(}\DecValTok{1}\SpecialCharTok{:}\DecValTok{5}\NormalTok{, }\AttributeTok{each =}\NormalTok{ n\_data),}
  \AttributeTok{data\_point =} \FunctionTok{rep}\NormalTok{(}\DecValTok{1}\SpecialCharTok{:}\NormalTok{n\_data, }\DecValTok{5}\NormalTok{),}
  \AttributeTok{role =} \FunctionTok{ifelse}\NormalTok{(cv\_data}\SpecialCharTok{$}\NormalTok{fold[data\_point] }\SpecialCharTok{==}\NormalTok{ iteration, }\StringTok{"test"}\NormalTok{, }\StringTok{"train"}\NormalTok{)}
\NormalTok{)}

\FunctionTok{ggplot}\NormalTok{(cv\_viz\_data, }\FunctionTok{aes}\NormalTok{(}\AttributeTok{x =}\NormalTok{ data\_point, }\AttributeTok{y =}\NormalTok{ iteration, }\AttributeTok{fill =}\NormalTok{ role)) }\SpecialCharTok{+}
  \FunctionTok{geom\_tile}\NormalTok{(}\AttributeTok{color =} \StringTok{"white"}\NormalTok{, }\AttributeTok{linewidth =} \FloatTok{0.5}\NormalTok{) }\SpecialCharTok{+}
  \FunctionTok{scale\_fill\_manual}\NormalTok{(}\AttributeTok{values =} \FunctionTok{c}\NormalTok{(}\StringTok{"train"} \OtherTok{=} \StringTok{"steelblue"}\NormalTok{, }\StringTok{"test"} \OtherTok{=} \StringTok{"firebrick"}\NormalTok{), }
                    \AttributeTok{name =} \StringTok{"Data Role"}\NormalTok{) }\SpecialCharTok{+}
  \FunctionTok{labs}\NormalTok{(}\AttributeTok{title =} \StringTok{"K{-}Fold Cross{-}Validation (k=5)"}\NormalTok{,}
       \AttributeTok{subtitle =} \StringTok{"Each row is one iteration; the model must predict the red tiles."}\NormalTok{,}
       \AttributeTok{x =} \StringTok{"Data Point"}\NormalTok{, }\AttributeTok{y =} \StringTok{"Iteration"}\NormalTok{) }\SpecialCharTok{+}
  \FunctionTok{theme}\NormalTok{(}\AttributeTok{panel.grid =} \FunctionTok{element\_blank}\NormalTok{(), }\AttributeTok{legend.position =} \StringTok{"bottom"}\NormalTok{)}
\end{Highlighting}
\end{Shaded}

\begin{center}
\includegraphics[width=0.8\linewidth,height=\textheight,keepaspectratio]{07-ModelComparison_files/figure-pdf/07-kfold-viz-1.pdf}
\end{center}

\subsubsection{Leave-One-Out Cross-Validation
(LOO-CV)}\label{leave-one-out-cross-validation-loo-cv}

Leave-One-Out Cross-Validation (LOO-CV) is the extreme limit of K-Fold,
where \(K\) equals the total number of data points (\(N\)). In each
iteration, we hold out exactly one observation for testing and train the
model on the remaining \(N-1\) observations.

\begin{Shaded}
\begin{Highlighting}[]
\NormalTok{loo\_viz\_data }\OtherTok{\textless{}{-}} \FunctionTok{tibble}\NormalTok{(}
  \AttributeTok{iteration =} \FunctionTok{rep}\NormalTok{(}\DecValTok{1}\SpecialCharTok{:}\NormalTok{n\_data, }\AttributeTok{each =}\NormalTok{ n\_data),}
  \AttributeTok{data\_point =} \FunctionTok{rep}\NormalTok{(}\DecValTok{1}\SpecialCharTok{:}\NormalTok{n\_data, n\_data),}
  \AttributeTok{role =} \FunctionTok{ifelse}\NormalTok{(iteration }\SpecialCharTok{==}\NormalTok{ data\_point, }\StringTok{"test"}\NormalTok{, }\StringTok{"train"}\NormalTok{)}
\NormalTok{)}

\FunctionTok{ggplot}\NormalTok{(loo\_viz\_data, }\FunctionTok{aes}\NormalTok{(}\AttributeTok{x =}\NormalTok{ data\_point, }\AttributeTok{y =}\NormalTok{ iteration, }\AttributeTok{fill =}\NormalTok{ role)) }\SpecialCharTok{+}
  \FunctionTok{geom\_tile}\NormalTok{(}\AttributeTok{color =} \StringTok{"white"}\NormalTok{, }\AttributeTok{linewidth =} \FloatTok{0.5}\NormalTok{) }\SpecialCharTok{+}
  \FunctionTok{scale\_fill\_manual}\NormalTok{(}\AttributeTok{values =} \FunctionTok{c}\NormalTok{(}\StringTok{"train"} \OtherTok{=} \StringTok{"steelblue"}\NormalTok{, }\StringTok{"test"} \OtherTok{=} \StringTok{"firebrick"}\NormalTok{), }
                    \AttributeTok{name =} \StringTok{"Data Role"}\NormalTok{) }\SpecialCharTok{+}
  \FunctionTok{labs}\NormalTok{(}\AttributeTok{title =} \StringTok{"Leave{-}One{-}Out Cross{-}Validation"}\NormalTok{,}
       \AttributeTok{subtitle =} \StringTok{"Maximum training data used, but requires fitting the model N times."}\NormalTok{,}
       \AttributeTok{x =} \StringTok{"Data Point"}\NormalTok{, }\AttributeTok{y =} \StringTok{"Iteration"}\NormalTok{) }\SpecialCharTok{+}
  \FunctionTok{theme}\NormalTok{(}\AttributeTok{panel.grid =} \FunctionTok{element\_blank}\NormalTok{(), }\AttributeTok{legend.position =} \StringTok{"bottom"}\NormalTok{)}
\end{Highlighting}
\end{Shaded}

\begin{center}
\includegraphics[width=0.8\linewidth,height=\textheight,keepaspectratio]{07-ModelComparison_files/figure-pdf/07-loo-cv-viz-1.pdf}
\end{center}

Cross-validation is especially important in cognitive Bayesian modeling
for several reasons:

\begin{enumerate}
\def\labelenumi{\arabic{enumi}.}
\item
  Bayesian models can be highly parameterized (especially hierarchical
  ones), making them prone to overfitting.
\item
  Prior choices can significantly influence a model's predictive
  performance.
\item
  We often need to compare competing theoretical accounts of the same
  behavior.
\end{enumerate}

However, implementing cross-validation for Bayesian models presents two
key challenges that we must solve:

\begin{enumerate}
\def\labelenumi{\arabic{enumi}.}
\item
  \textbf{Proper Scoring}: In standard machine learning, we might just
  calculate the mean squared error or accuracy on the test fold. But
  Bayesian models output \emph{probability distributions}, not just
  point estimates. We need an appropriate metric for evaluating the
  quality of the full predictive distribution.
\item
  \textbf{Computational Cost}: For a dataset with 6,000 observations,
  true LOO-CV would require fitting our Stan model 6,000 separate times.
  Because MCMC sampling can take hours, exact LOO-CV is usually
  practically impossible.
\end{enumerate}

In the next sections, we will solve the first challenge using the
\textbf{Expected Log Predictive Density (ELPD)}, and the second
challenge using an elegant mathematical shortcut called
\textbf{PSIS-LOO}.

\section{Formalizing Predictive Accuracy: The
ELPD}\label{formalizing-predictive-accuracy-the-elpd}

\subsection{The Log-Score Proper Scoring
Rule}\label{the-log-score-proper-scoring-rule}

Any model \(M\) can generate a predictive distribution: how likely
different values of the outcome variable are estimated to be given the
data it has seen and the model's assumptions.

When an actual new outcome \(\tilde{y}_0\) is observed, we evaluate the
model by checking exactly how much probability mass it assigned to the
event that actually happened: \(p(\tilde{y}_0 \mid y, M)\).

We take the natural logarithm of this probability to compute the
log-score:
\[\text{log-score}(M, \tilde{y}_0) = \log p(\tilde{y}_0 \mid y, M)\]

We use the log scale for two primary reasons:

\begin{itemize}
\item
  \textbf{Additivity}: The joint probability of independent future
  observations is the product of their individual probabilities. On the
  log scale, this becomes a sum, which is mathematically tractable and
  avoids computational underflow: probabilities being smaller than one
  being multiplied with each other become small very quickly and run
  into approximation issues, because computers use only a fixed number
  of decimals.
\item
  \textbf{Information Theory}: The negative log-score,
  \(-\log p(\tilde{y}_0 \mid y, M)\), is the \emph{surprisal} or Shannon
  information content of the outcome under the model. Maximizing the
  log-score is equivalent to minimizing the model's surprise.
\end{itemize}

We call the log-score a \emph{scoring rule}: a function that provides a
summary measure for the evaluation of probabilistic forecasts. A scoring
rule is strictly \emph{proper} if the model maximizes its expected score
if and only if the probability distribution it outputs exactly matches
the true data-generating distribution. This is a desirable property: a
model cannot ``game'' the score by being strategically overconfident or
underconfident. Overconfident models - those that assign near-zero
probability to outcomes that then occur - incur unbounded penalties. A
model that says ``I am 99\% certain this participant will choose right''
and the participant chooses left will be penalized far more than a model
that honestly said ``I think it's about 60\%.''

You may have encountered Bayes Factors as a method for model comparison.
Bayes Factors rely on the marginal likelihood---the probability of the
data integrated over the entire prior. While theoretically elegant,
Bayes Factors are notoriously sensitive to the exact width of priors and
are often computationally impossible to calculate for the complex
hierarchical models used in cognitive science. In this course, we rely
on ELPD and Cross-Validation because they focus on predictive
generalization, are more robust to prior choices, and provide
trial-level diagnostics that Bayes Factors lack.

\subsection{The Expected Log Predictive
Density}\label{the-expected-log-predictive-density}

The log-score evaluates a model on a single new observation. But what we
really want is a summary over \emph{all possible} future observations -
a measure of how good the model would be in general, not just on one
lucky or unlucky draw.

The Expected Log Predictive Density (ELPD) does exactly this: it
averages the log-score over all possible future observations, weighted
by how likely those observations are under the true data-generating
distribution \(p_t\):

\[\text{ELPD} = \int p_t(\tilde{y}) \log p(\tilde{y} \mid y, M) \, d\tilde{y}\]

where
\(p(\tilde{y} \mid y, M) = \int p(\tilde{y} \mid \theta) \, p(\theta \mid y) \, d\theta\)
is the posterior predictive distribution - the model's prediction for a
new observation, averaged over the uncertainty in its parameters. The
model with the highest ELPD makes the best average predictions across
the full range of possible future data.

There is one crucial catch: we do not know \(p_t\) - the true
distribution of future observations. If we did, we would not need a
model in the first place! So we have to \emph{estimate} the ELPD from
the data we have. Leave-one-out cross-validation (LOO-CV) provides an
approach that is asymptotically unbiased:

\[\widehat{\text{ELPD}}_{\text{LOO}} = \sum_{i=1}^{N} \log p(y_i \mid y_{-i}, M)\]

The notation \(y_{-i}\) means ``all observations except observation
\(i\)''. So for each data point, we ask: if the model had never seen
this observation, how well would it predict it? A model that memorized
\(y_i\) to reduce its training loss will be penalized here, because the
LOO posterior is computed \emph{without} \(y_i\). The prediction must
generalize - it cannot cheat by peeking at the answer.

But how do we calculate it?

\subsection{The Importance-Sampling
Shortcut}\label{the-importance-sampling-shortcut}

Computing \(\widehat{\text{ELPD}}_{\text{LOO}}\) exactly would require
fitting the model \(N\) times - once with each observation left out. For
our hierarchical models with 50 agents and 120 trials each, that is
6,000 separate MCMC runs. This is not feasible in practice, especially
when we get to more complex models.

Fortunately, there is an elegant mathematical shortcut, which is
analogous to what we have seen when assessing prior sensitivity. The
idea is to use the \emph{single} full-data posterior we already have and
to \emph{reweight} its draws to approximate what the posterior would
have looked like had we omitted observation \(i\). .

The mathematical weights that accomplish this reweighting are the
\emph{inverse} of the likelihood contributions: observations that the
model fits well (high likelihood) get down-weighted when computing what
the model ``would have'' predicted without them. These log-scale
importance weights are simply the negative pointwise log-likelihoods:

\[\log r_i^{(s)} = -\log p(y_i \mid \theta^{(s)})\]

These are already stored in our Stan \texttt{log\_lik} block from
Chapter 7. So no additional computation is required: everything we need
is already sitting in our existing model fits.

However, we should be careful. If the model's estimate would barely
change when we removing observation \(i\), then \(i\) is not very
influential, and we can approximate the LOO estimate cheaply. If the
model's estimate would change a lot, then \(i\) is influential, and our
approximation is less reliable. We therefore need to use a diagnostic to
flag such cases: the Pareto-\(\hat{k}\), which we cover in detail in a
later paragraph.

This combination of LOO cross-validation with Pareto-Smoothed Importance
Sampling is called \textbf{PSIS-LOO} (Vehtari, Gelman \& Gabry, 2017),
and it is our primary tool for model comparison, so we'll go through it
step by step.

\begin{quote}
\emph{Why learn the mechanics before using \texttt{loo()}?} The
\texttt{loo()} function is one line of R. But knowing what it is doing
--- reweighting posterior draws to approximate leave-one-out
predictions, then checking whether those weights are stable via the
Pareto tail --- lets you diagnose failures and interpret the
Pareto-\(\hat{k}\) diagnostic as a \emph{cognitive} signal (which
observations is the model surprised by?) rather than just a warning
light. The under-the-hood demonstration on the next page is worth
reading once carefully; after that, \texttt{loo()} will never be a black
box.
\end{quote}

To really understand what the \texttt{loo()} function will do for us
later, let's look under the hood. We can simulate the PSIS-LOO process
directly using the \texttt{loo} package's internal engine. We will
generate 1,000 simulated log-likelihood values, inject some extreme
outliers to simulate a ``shocked'' model, and feed them directly into
the actual \texttt{loo::psis()} algorithm.

We can then extract the smoothed weights and compare them to the raw
weights to see exactly how the algorithm protects us from extreme MCMC
draws.

\begin{Shaded}
\begin{Highlighting}[]
\NormalTok{n\_draws }\OtherTok{\textless{}{-}} \DecValTok{1000}

\NormalTok{n\_draws }\OtherTok{\textless{}{-}} \DecValTok{1000}

\CommentTok{\# 1. Simulate realistic raw log{-}importance weights}
\CommentTok{\# 995 well{-}behaved draws, and 5 massive breakaway outliers}
\NormalTok{log\_w\_raw }\OtherTok{\textless{}{-}} \FunctionTok{c}\NormalTok{(}\FunctionTok{rnorm}\NormalTok{(}\DecValTok{995}\NormalTok{, }\AttributeTok{mean =} \DecValTok{0}\NormalTok{, }\AttributeTok{sd =} \DecValTok{1}\NormalTok{), }\DecValTok{10}\NormalTok{, }\DecValTok{12}\NormalTok{, }\DecValTok{14}\NormalTok{, }\DecValTok{15}\NormalTok{, }\DecValTok{18}\NormalTok{)}

\CommentTok{\# 2. Apply the ACTUAL PSIS algorithm}
\NormalTok{psis\_result }\OtherTok{\textless{}{-}}\NormalTok{ loo}\SpecialCharTok{::}\FunctionTok{psis}\NormalTok{(log\_w\_raw, }\AttributeTok{r\_eff =} \DecValTok{1}\NormalTok{)}

\CommentTok{\# Print the diagnostic to prove we triggered the smoothing}
\FunctionTok{cat}\NormalTok{(}\StringTok{"Pareto{-}k diagnostic from the real algorithm:"}\NormalTok{, }\FunctionTok{round}\NormalTok{(psis\_result}\SpecialCharTok{$}\NormalTok{diagnostics}\SpecialCharTok{$}\NormalTok{pareto\_k, }\DecValTok{2}\NormalTok{), }\StringTok{"}\SpecialCharTok{\textbackslash{}n}\StringTok{"}\NormalTok{)}
\end{Highlighting}
\end{Shaded}

\begin{verbatim}
Pareto-k diagnostic from the real algorithm: 1.76 
\end{verbatim}

\begin{Shaded}
\begin{Highlighting}[]
\CommentTok{\# Print the diagnostic to prove we triggered the smoothing}
\FunctionTok{cat}\NormalTok{(}\StringTok{"Pareto{-}k diagnostic from the real algorithm:"}\NormalTok{, }\FunctionTok{round}\NormalTok{(psis\_result}\SpecialCharTok{$}\NormalTok{diagnostics}\SpecialCharTok{$}\NormalTok{pareto\_k, }\DecValTok{2}\NormalTok{), }\StringTok{"}\SpecialCharTok{\textbackslash{}n}\StringTok{"}\NormalTok{)}
\end{Highlighting}
\end{Shaded}

\begin{verbatim}
Pareto-k diagnostic from the real algorithm: 1.76 
\end{verbatim}

\begin{Shaded}
\begin{Highlighting}[]
\CommentTok{\# 3. Extract and normalize weights for visualization}
\CommentTok{\# Raw weights}
\NormalTok{w\_raw }\OtherTok{\textless{}{-}} \FunctionTok{exp}\NormalTok{(log\_w\_raw)}
\NormalTok{w\_raw\_norm }\OtherTok{\textless{}{-}}\NormalTok{ w\_raw }\SpecialCharTok{/} \FunctionTok{sum}\NormalTok{(w\_raw)}

\CommentTok{\# Smoothed weights (must also be explicitly normalized)}
\NormalTok{w\_smooth\_raw }\OtherTok{\textless{}{-}} \FunctionTok{exp}\NormalTok{(}\FunctionTok{as.numeric}\NormalTok{(psis\_result}\SpecialCharTok{$}\NormalTok{log\_weights))}
\NormalTok{w\_smooth\_norm }\OtherTok{\textless{}{-}}\NormalTok{ w\_smooth\_raw }\SpecialCharTok{/} \FunctionTok{sum}\NormalTok{(w\_smooth\_raw)}

\CommentTok{\# 4. Sort the weights from smallest to largest to visualize the "tail"}
\CommentTok{\# We sort by the raw weights so we can see how the specific extreme indices change}
\NormalTok{sort\_order }\OtherTok{\textless{}{-}} \FunctionTok{order}\NormalTok{(w\_raw\_norm)}
\NormalTok{weight\_data }\OtherTok{\textless{}{-}} \FunctionTok{data.frame}\NormalTok{(}
  \AttributeTok{raw =}\NormalTok{ w\_raw\_norm[sort\_order],}
  \AttributeTok{smoothed =}\NormalTok{ w\_smooth\_norm[sort\_order]}
\NormalTok{)}

\CommentTok{\# The PSIS algorithm smooths the top 3 * sqrt(S) weights. }
\CommentTok{\# For S=1000, that is the top \textasciitilde{}95 weights.}
\NormalTok{weight\_data\_top }\OtherTok{\textless{}{-}} \FunctionTok{tail}\NormalTok{(weight\_data, }\DecValTok{150}\NormalTok{)}
\NormalTok{weight\_data\_top}\SpecialCharTok{$}\NormalTok{rank }\OtherTok{\textless{}{-}} \DecValTok{1}\SpecialCharTok{:}\DecValTok{150}
\NormalTok{tail\_start\_x }\OtherTok{\textless{}{-}} \DecValTok{150} \SpecialCharTok{{-}} \DecValTok{95} \SpecialCharTok{+} \FloatTok{0.5}

\CommentTok{\# PLOT: Importance weights before and after Pareto smoothing}
\NormalTok{p\_weights }\OtherTok{\textless{}{-}} \FunctionTok{ggplot}\NormalTok{(weight\_data\_top, }\FunctionTok{aes}\NormalTok{(}\AttributeTok{x =}\NormalTok{ rank)) }\SpecialCharTok{+}
  \CommentTok{\# Arrows showing the shift computed by the real algorithm}
  \FunctionTok{geom\_segment}\NormalTok{(}\FunctionTok{aes}\NormalTok{(}\AttributeTok{xend =}\NormalTok{ rank, }\AttributeTok{y =}\NormalTok{ raw, }\AttributeTok{yend =}\NormalTok{ smoothed),}
               \AttributeTok{arrow =} \FunctionTok{arrow}\NormalTok{(}\AttributeTok{length =} \FunctionTok{unit}\NormalTok{(}\FloatTok{0.2}\NormalTok{, }\StringTok{"cm"}\NormalTok{), }\AttributeTok{type =} \StringTok{"closed"}\NormalTok{), }
               \AttributeTok{color =} \StringTok{"purple"}\NormalTok{, }\AttributeTok{linewidth =} \FloatTok{0.7}\NormalTok{, }\AttributeTok{alpha =} \FloatTok{0.7}\NormalTok{) }\SpecialCharTok{+}
               
  \CommentTok{\# Points drawn on top}
  \FunctionTok{geom\_point}\NormalTok{(}\FunctionTok{aes}\NormalTok{(}\AttributeTok{y =}\NormalTok{ raw), }\AttributeTok{color =} \StringTok{"steelblue"}\NormalTok{, }\AttributeTok{size =} \DecValTok{2}\NormalTok{, }\AttributeTok{alpha =} \FloatTok{0.8}\NormalTok{) }\SpecialCharTok{+}
  \FunctionTok{geom\_point}\NormalTok{(}\FunctionTok{aes}\NormalTok{(}\AttributeTok{y =}\NormalTok{ smoothed), }\AttributeTok{color =} \StringTok{"purple"}\NormalTok{, }\AttributeTok{size =} \DecValTok{2}\NormalTok{, }\AttributeTok{alpha =} \FloatTok{1.0}\NormalTok{) }\SpecialCharTok{+}
  
  \CommentTok{\# Add a vertical line to show where the algorithm began smoothing}
  \FunctionTok{geom\_vline}\NormalTok{(}\AttributeTok{xintercept =}\NormalTok{ tail\_start\_x, }\AttributeTok{linetype =} \StringTok{"dashed"}\NormalTok{, }\AttributeTok{color =} \StringTok{"gray50"}\NormalTok{) }\SpecialCharTok{+}
  \FunctionTok{annotate}\NormalTok{(}\StringTok{"text"}\NormalTok{, }\AttributeTok{x =}\NormalTok{ tail\_start\_x }\SpecialCharTok{+} \DecValTok{2}\NormalTok{, }\AttributeTok{y =} \FunctionTok{max}\NormalTok{(weight\_data\_top}\SpecialCharTok{$}\NormalTok{raw) }\SpecialCharTok{*} \FloatTok{0.9}\NormalTok{, }
           \AttributeTok{label =} \StringTok{"Algorithm begins}\SpecialCharTok{\textbackslash{}n}\StringTok{smoothing here"}\NormalTok{, }\AttributeTok{hjust =} \DecValTok{0}\NormalTok{, }\AttributeTok{color =} \StringTok{"gray30"}\NormalTok{) }\SpecialCharTok{+}
  
  \CommentTok{\# Labels and styling}
  \FunctionTok{labs}\NormalTok{(}\AttributeTok{title =} \StringTok{"Actual PSIS Algorithm Smoothing (Top 150 Weights)"}\NormalTok{,}
       \AttributeTok{subtitle =} \StringTok{"The loo::psis() function compresses the extreme tail; remaining weights shift up to absorb the mass."}\NormalTok{,}
       \AttributeTok{x =} \StringTok{"Rank of Weight (Zoomed in on the top 15\%)"}\NormalTok{, }\AttributeTok{y =} \StringTok{"Normalized Weight"}\NormalTok{) }\SpecialCharTok{+}
  \FunctionTok{theme\_minimal}\NormalTok{()}

\CommentTok{\# PLOT: Conceptual diagram of PSIS{-}LOO process}
\NormalTok{p\_flow }\OtherTok{\textless{}{-}} \FunctionTok{ggplot}\NormalTok{() }\SpecialCharTok{+}
  \FunctionTok{annotate}\NormalTok{(}\StringTok{"rect"}\NormalTok{, }\AttributeTok{xmin =} \FloatTok{0.5}\NormalTok{, }\AttributeTok{xmax =} \FloatTok{3.5}\NormalTok{, }\AttributeTok{ymin =} \DecValTok{4}\NormalTok{, }\AttributeTok{ymax =} \DecValTok{5}\NormalTok{, }\AttributeTok{fill =} \StringTok{"lightblue"}\NormalTok{, }\AttributeTok{alpha =} \FloatTok{0.7}\NormalTok{) }\SpecialCharTok{+}
  \FunctionTok{annotate}\NormalTok{(}\StringTok{"text"}\NormalTok{, }\AttributeTok{x =} \DecValTok{2}\NormalTok{, }\AttributeTok{y =} \FloatTok{4.5}\NormalTok{, }\AttributeTok{label =} \StringTok{"1. Fit Model to All Data}\SpecialCharTok{\textbackslash{}n}\StringTok{Obtain posterior samples p(θ|y)"}\NormalTok{) }\SpecialCharTok{+}
  
  \FunctionTok{annotate}\NormalTok{(}\StringTok{"rect"}\NormalTok{, }\AttributeTok{xmin =} \FloatTok{0.5}\NormalTok{, }\AttributeTok{xmax =} \FloatTok{3.5}\NormalTok{, }\AttributeTok{ymin =} \FloatTok{2.5}\NormalTok{, }\AttributeTok{ymax =} \FloatTok{3.5}\NormalTok{, }\AttributeTok{fill =} \StringTok{"lightgreen"}\NormalTok{, }\AttributeTok{alpha =} \FloatTok{0.7}\NormalTok{) }\SpecialCharTok{+}
  \FunctionTok{annotate}\NormalTok{(}\StringTok{"text"}\NormalTok{, }\AttributeTok{x =} \DecValTok{2}\NormalTok{, }\AttributeTok{y =} \DecValTok{3}\NormalTok{, }\AttributeTok{label =} \StringTok{"2. Calculate Log{-}Likelihoods}\SpecialCharTok{\textbackslash{}n}\StringTok{ℓᵢ(θ) = log p(yᵢ|θ)"}\NormalTok{) }\SpecialCharTok{+}
  
  \FunctionTok{annotate}\NormalTok{(}\StringTok{"rect"}\NormalTok{, }\AttributeTok{xmin =} \FloatTok{0.5}\NormalTok{, }\AttributeTok{xmax =} \FloatTok{3.5}\NormalTok{, }\AttributeTok{ymin =} \DecValTok{1}\NormalTok{, }\AttributeTok{ymax =} \DecValTok{2}\NormalTok{, }\AttributeTok{fill =} \StringTok{"lightyellow"}\NormalTok{, }\AttributeTok{alpha =} \FloatTok{0.7}\NormalTok{) }\SpecialCharTok{+}
  \FunctionTok{annotate}\NormalTok{(}\StringTok{"text"}\NormalTok{, }\AttributeTok{x =} \DecValTok{2}\NormalTok{, }\AttributeTok{y =} \FloatTok{1.5}\NormalTok{, }\AttributeTok{label =} \StringTok{"3. Compute Importance Weights}\SpecialCharTok{\textbackslash{}n}\StringTok{wᵢ(θ) ∝ 1/p(yᵢ|θ)"}\NormalTok{) }\SpecialCharTok{+}
  
  \FunctionTok{annotate}\NormalTok{(}\StringTok{"rect"}\NormalTok{, }\AttributeTok{xmin =} \FloatTok{0.5}\NormalTok{, }\AttributeTok{xmax =} \FloatTok{3.5}\NormalTok{, }\AttributeTok{ymin =} \SpecialCharTok{{-}}\FloatTok{0.5}\NormalTok{, }\AttributeTok{ymax =} \FloatTok{0.5}\NormalTok{, }\AttributeTok{fill =} \StringTok{"mistyrose"}\NormalTok{, }\AttributeTok{alpha =} \FloatTok{0.7}\NormalTok{) }\SpecialCharTok{+}
  \FunctionTok{annotate}\NormalTok{(}\StringTok{"text"}\NormalTok{, }\AttributeTok{x =} \DecValTok{2}\NormalTok{, }\AttributeTok{y =} \DecValTok{0}\NormalTok{, }\AttributeTok{label =} \StringTok{"4. Apply Pareto Smoothing}\SpecialCharTok{\textbackslash{}n}\StringTok{Stabilize extreme weights via loo::psis()"}\NormalTok{) }\SpecialCharTok{+}
  
  \FunctionTok{annotate}\NormalTok{(}\StringTok{"rect"}\NormalTok{, }\AttributeTok{xmin =} \FloatTok{0.5}\NormalTok{, }\AttributeTok{xmax =} \FloatTok{3.5}\NormalTok{, }\AttributeTok{ymin =} \SpecialCharTok{{-}}\DecValTok{2}\NormalTok{, }\AttributeTok{ymax =} \SpecialCharTok{{-}}\DecValTok{1}\NormalTok{, }\AttributeTok{fill =} \StringTok{"lavender"}\NormalTok{, }\AttributeTok{alpha =} \FloatTok{0.7}\NormalTok{) }\SpecialCharTok{+}
  \FunctionTok{annotate}\NormalTok{(}\StringTok{"text"}\NormalTok{, }\AttributeTok{x =} \DecValTok{2}\NormalTok{, }\AttributeTok{y =} \SpecialCharTok{{-}}\FloatTok{1.5}\NormalTok{, }\AttributeTok{label =} \StringTok{"5. Compute ELPD Using Weights}\SpecialCharTok{\textbackslash{}n}\StringTok{Approximate leave{-}one{-}out prediction"}\NormalTok{) }\SpecialCharTok{+}
  
  \FunctionTok{annotate}\NormalTok{(}\StringTok{"segment"}\NormalTok{, }\AttributeTok{x =} \DecValTok{2}\NormalTok{, }\AttributeTok{y =} \FunctionTok{c}\NormalTok{(}\DecValTok{4}\NormalTok{, }\FloatTok{2.5}\NormalTok{, }\DecValTok{1}\NormalTok{, }\SpecialCharTok{{-}}\FloatTok{0.5}\NormalTok{), }\AttributeTok{xend =} \DecValTok{2}\NormalTok{, }\AttributeTok{yend =} \FunctionTok{c}\NormalTok{(}\FloatTok{3.5}\NormalTok{, }\DecValTok{2}\NormalTok{, }\FloatTok{0.5}\NormalTok{, }\SpecialCharTok{{-}}\DecValTok{1}\NormalTok{), }
           \AttributeTok{arrow =} \FunctionTok{arrow}\NormalTok{(}\AttributeTok{length =} \FunctionTok{unit}\NormalTok{(}\FloatTok{0.2}\NormalTok{, }\StringTok{"cm"}\NormalTok{))) }\SpecialCharTok{+}
  
  \FunctionTok{annotate}\NormalTok{(}\StringTok{"rect"}\NormalTok{, }\AttributeTok{xmin =} \FloatTok{3.7}\NormalTok{, }\AttributeTok{xmax =} \FloatTok{5.7}\NormalTok{, }\AttributeTok{ymin =} \SpecialCharTok{{-}}\FloatTok{0.2}\NormalTok{, }\AttributeTok{ymax =} \FloatTok{1.5}\NormalTok{, }\AttributeTok{fill =} \StringTok{"white"}\NormalTok{, }\AttributeTok{color =} \StringTok{"black"}\NormalTok{, }\AttributeTok{linetype =} \StringTok{"dashed"}\NormalTok{) }\SpecialCharTok{+}
  \FunctionTok{annotate}\NormalTok{(}\StringTok{"text"}\NormalTok{, }\AttributeTok{x =} \FloatTok{4.7}\NormalTok{, }\AttributeTok{y =} \FloatTok{1.2}\NormalTok{, }\AttributeTok{label =} \StringTok{"Pareto k Diagnostics"}\NormalTok{, }\AttributeTok{fontface =} \StringTok{"bold"}\NormalTok{) }\SpecialCharTok{+}
  \FunctionTok{annotate}\NormalTok{(}\StringTok{"text"}\NormalTok{, }\AttributeTok{x =} \FloatTok{4.7}\NormalTok{, }\AttributeTok{y =} \FloatTok{0.8}\NormalTok{, }\AttributeTok{label =} \StringTok{"k \textless{} 0.5: Reliable"}\NormalTok{, }\AttributeTok{color =} \StringTok{"darkgreen"}\NormalTok{) }\SpecialCharTok{+}
  \FunctionTok{annotate}\NormalTok{(}\StringTok{"text"}\NormalTok{, }\AttributeTok{x =} \FloatTok{4.7}\NormalTok{, }\AttributeTok{y =} \FloatTok{0.4}\NormalTok{, }\AttributeTok{label =} \StringTok{"0.5 \textless{} k \textless{} 0.7: Use caution"}\NormalTok{, }\AttributeTok{color =} \StringTok{"darkorange"}\NormalTok{) }\SpecialCharTok{+}
  \FunctionTok{annotate}\NormalTok{(}\StringTok{"text"}\NormalTok{, }\AttributeTok{x =} \FloatTok{4.7}\NormalTok{, }\AttributeTok{y =} \DecValTok{0}\NormalTok{, }\AttributeTok{label =} \StringTok{"k \textgreater{} 0.7: Unreliable"}\NormalTok{, }\AttributeTok{color =} \StringTok{"darkred"}\NormalTok{) }\SpecialCharTok{+}
  \FunctionTok{theme\_void}\NormalTok{() }\SpecialCharTok{+}
  \FunctionTok{labs}\NormalTok{(}\AttributeTok{title =} \StringTok{"PSIS{-}LOO Process Flow"}\NormalTok{)}

\NormalTok{p\_weights }\SpecialCharTok{/}\NormalTok{ p\_flow}
\end{Highlighting}
\end{Shaded}

\begin{center}
\includegraphics[width=0.8\linewidth,height=\textheight,keepaspectratio]{07-ModelComparison_files/figure-pdf/07-psis-under-the-hood-1.pdf}
\end{center}

Notice how a few random MCMC draws initially monopolize the raw
importance weights (the high blue dots). If we calculated ELPD using
those raw weights, our cross-validation score would be highly erratic
and dictated by statistical noise. The Pareto smoothing (purple dots)
gently compresses the extreme tail of these weights into a more stable
distribution.

Mathematically, an extreme raw importance weight means the approximation
is unstable. Cognitively, an extreme importance weight means the model
is utterly shocked by this specific participant's choice. We will use
this dual mathematical/cognitive interpretation to flag influential
observations later.

But first let's get to the thick of it and start simulating data.

\section{Simulation Setup for Model
Recovery}\label{simulation-setup-for-model-recovery}

We simulate two populations of agents playing against a shared opponent
using the \textbf{reversal design} first validated in Ch. 6 and used
throughout Ch. 7: the opponent plays 80\% right for the first half of
the session, then switches to 20\% right for the second half. We use the
exact sequence loaded from Ch. 7 above (\texttt{opponent\_reversal}),
taking the first 120 trials (60 + 60) for this chapter's shorter
sessions.

Why this design? If the opponent plays a constant bias, the memory
agent's internal running average quickly settles near that value and
stays flat. Behaviorally, an agent responding to a flat memory signal is
mathematically indistinguishable from a biased agent --- both produce a
fixed right-leaning rate. The reversal forces the memory signal to
change sign at trial 61: memory agents reverse their behavior while
biased agents do not. This contrast is what makes the two models
theoretically separable, and why Ch. 6 discovered that a constant-bias
design made \texttt{bias} and \texttt{beta} unidentifiable.

\emph{Note on recovery failure:} We simulate a population with strong
memory sensitivity (\(\mu_\beta = 1.5\)). With weak memory
(\(\mu_\beta \approx 0.3\)), the behavioral difference would be subtle
enough that the simpler biased model would win the ELPD comparison by
default due to its lower complexity penalty --- a correct finding, not a
method failure.

The key insight for model recovery is that we fit \emph{both} models to
\emph{both} datasets, giving a 2×2 matrix of fits:

\begin{itemize}
\item
  \textbf{Biased model on biased data} --- the correctly-specified
  model.
\item
  \textbf{Memory model on biased data} --- the over-specified model; it
  has an extra parameter (\(\beta\)) it does not need.
\item
  \textbf{Biased model on memory data} --- the under-specified model; it
  lacks the flexibility to capture the reversal.
\item
  \textbf{Memory model on memory data} --- the correctly-specified
  model.
\end{itemize}

PSIS-LOO should prefer the correct model in each column. This is
\textbf{model recovery}: the analogue of parameter recovery one level
higher --- instead of asking ``can we recover the right parameter
values?'', we ask ``can we identify the right model?''

\begin{Shaded}
\begin{Highlighting}[]
\NormalTok{simulate\_model\_recovery\_data }\OtherTok{\textless{}{-}} \ControlFlowTok{function}\NormalTok{(}
  \AttributeTok{n\_agents         =} \DecValTok{50}\NormalTok{,  }\AttributeTok{n\_trials      =} \DecValTok{120}\NormalTok{,}
  \AttributeTok{mu\_theta         =} \FloatTok{1.0}\NormalTok{, }\AttributeTok{sigma\_theta   =} \FloatTok{0.5}\NormalTok{,}
  \AttributeTok{mu\_alpha         =} \FloatTok{0.0}\NormalTok{, }\AttributeTok{mu\_beta       =} \FloatTok{1.5}\NormalTok{,}
  \AttributeTok{sigma\_alpha      =} \FloatTok{0.3}\NormalTok{, }\AttributeTok{sigma\_beta    =} \FloatTok{0.5}\NormalTok{,}
  \AttributeTok{rho              =} \FloatTok{0.3}\NormalTok{,}
  \AttributeTok{opponent\_choices =} \ConstantTok{NULL}\NormalTok{,   }
  \AttributeTok{seed             =} \DecValTok{42}
\NormalTok{) \{}
  \FunctionTok{set.seed}\NormalTok{(seed)}

  \CommentTok{\# Reversal block design: 80\% right for t 1:60, 20\% right for t 61:120.}
  \CommentTok{\# If a fixed sequence is supplied (from Ch. 7), use it directly so that}
  \CommentTok{\# all chapters share the same opponent history.}
  \ControlFlowTok{if}\NormalTok{ (}\SpecialCharTok{!}\FunctionTok{is.null}\NormalTok{(opponent\_choices)) \{}
    \FunctionTok{stopifnot}\NormalTok{(}\FunctionTok{length}\NormalTok{(opponent\_choices) }\SpecialCharTok{\textgreater{}=}\NormalTok{ n\_trials)}
\NormalTok{    opp\_choices }\OtherTok{\textless{}{-}}\NormalTok{ opponent\_choices[}\FunctionTok{seq\_len}\NormalTok{(n\_trials)]}
\NormalTok{  \} }\ControlFlowTok{else}\NormalTok{ \{}
\NormalTok{    opp\_choices }\OtherTok{\textless{}{-}} \FunctionTok{c}\NormalTok{(}
      \FunctionTok{rbinom}\NormalTok{(}\DecValTok{60}\NormalTok{, }\DecValTok{1}\NormalTok{, }\FloatTok{0.8}\NormalTok{),}
      \FunctionTok{rbinom}\NormalTok{(n\_trials }\SpecialCharTok{{-}} \DecValTok{60}\NormalTok{, }\DecValTok{1}\NormalTok{, }\FloatTok{0.2}\NormalTok{)}
\NormalTok{    )}
\NormalTok{  \}}
  
  \CommentTok{\# Perfect unweighted memory (cumulative average)}
\NormalTok{  opp\_rate\_prev    }\OtherTok{\textless{}{-}} \FunctionTok{numeric}\NormalTok{(n\_trials)}
\NormalTok{  opp\_rate\_prev[}\DecValTok{1}\NormalTok{] }\OtherTok{\textless{}{-}} \FloatTok{0.5}
  \ControlFlowTok{if}\NormalTok{ (n\_trials }\SpecialCharTok{\textgreater{}} \DecValTok{1}\NormalTok{) \{}
\NormalTok{    raw\_rates         }\OtherTok{\textless{}{-}} \FunctionTok{cumsum}\NormalTok{(opp\_choices[}\SpecialCharTok{{-}}\NormalTok{n\_trials]) }\SpecialCharTok{/} \FunctionTok{seq\_len}\NormalTok{(n\_trials }\SpecialCharTok{{-}} \DecValTok{1}\NormalTok{)}
\NormalTok{    opp\_rate\_prev[}\SpecialCharTok{{-}}\DecValTok{1}\NormalTok{] }\OtherTok{\textless{}{-}} \FunctionTok{pmax}\NormalTok{(}\FloatTok{0.01}\NormalTok{, }\FunctionTok{pmin}\NormalTok{(}\FloatTok{0.99}\NormalTok{, raw\_rates)) }\CommentTok{\# to avoid singularities}
\NormalTok{  \}}

  \CommentTok{\# Generating the hierarchical biased behaviors}
  \CommentTok{\# Generating the hierarchical biased behaviors}
\NormalTok{  theta\_j }\OtherTok{\textless{}{-}} \FunctionTok{rnorm}\NormalTok{(n\_agents, mu\_theta, sigma\_theta)}
  \CommentTok{\# Inject extreme outliers to naturally trigger high Pareto{-}k values}
\NormalTok{  theta\_j[}\DecValTok{1}\NormalTok{] }\OtherTok{\textless{}{-}} \FloatTok{4.5}  \CommentTok{\# Agent 1 is extremely biased right}
\NormalTok{  theta\_j[}\DecValTok{2}\NormalTok{] }\OtherTok{\textless{}{-}} \SpecialCharTok{{-}}\FloatTok{4.5} \CommentTok{\# Agent 2 is extremely biased left}
\NormalTok{  biased\_data }\OtherTok{\textless{}{-}} \FunctionTok{map\_dfr}\NormalTok{(}\FunctionTok{seq\_len}\NormalTok{(n\_agents), }\ControlFlowTok{function}\NormalTok{(j) \{}
    \FunctionTok{tibble}\NormalTok{(}\AttributeTok{agent\_id =}\NormalTok{ j, }\AttributeTok{trial =} \FunctionTok{seq\_len}\NormalTok{(n\_trials),}
           \AttributeTok{choice =} \FunctionTok{rbinom}\NormalTok{(n\_trials, }\DecValTok{1}\NormalTok{, }\FunctionTok{plogis}\NormalTok{(theta\_j[j])),}
           \AttributeTok{opp\_rate\_prev =}\NormalTok{ opp\_rate\_prev, }\AttributeTok{agent\_type =} \StringTok{"biased"}\NormalTok{,}
           \AttributeTok{true\_theta =}\NormalTok{ theta\_j[j])}
\NormalTok{  \})}

  \CommentTok{\# Generating the parameters for the hierarchical memory behaviors with correlation}
  \DocumentationTok{\#\# We follow the NCP parametrization to show you how, but also because it\textquotesingle{}s more}
  \DocumentationTok{\#\# efficient for sampling correlated bias and beta.}
  
  \DocumentationTok{\#\# First we define the Cholesky factor for the correlation matrix:}
\NormalTok{  L\_Omega  }\OtherTok{\textless{}{-}} \FunctionTok{matrix}\NormalTok{(}\FunctionTok{c}\NormalTok{(}\DecValTok{1}\NormalTok{, rho, }\DecValTok{0}\NormalTok{, }\FunctionTok{sqrt}\NormalTok{(}\DecValTok{1} \SpecialCharTok{{-}}\NormalTok{ rho}\SpecialCharTok{\^{}}\DecValTok{2}\NormalTok{)), }\DecValTok{2}\NormalTok{, }\DecValTok{2}\NormalTok{)}
  \DocumentationTok{\#\# Then we need to convert it to the covariance matrix by pre{-}multiplying by }
  \DocumentationTok{\#\# the scale parameters (correlations are z{-}scored, so to get the actual }
  \DocumentationTok{\#\# values we need to put the scales back in):}
\NormalTok{  Scale    }\OtherTok{\textless{}{-}} \FunctionTok{diag}\NormalTok{(}\FunctionTok{c}\NormalTok{(sigma\_alpha, sigma\_beta)) }\SpecialCharTok{\%*\%}\NormalTok{ L\_Omega}
  \DocumentationTok{\#\# Then we create a matrix of independent, standard normal draws}
\NormalTok{  Z        }\OtherTok{\textless{}{-}} \FunctionTok{matrix}\NormalTok{(}\FunctionTok{rnorm}\NormalTok{(}\DecValTok{2} \SpecialCharTok{*}\NormalTok{ n\_agents), }\DecValTok{2}\NormalTok{, n\_agents)}
  \DocumentationTok{\#\# Then we apply our rescaling to the standard normal draws and add the group{-}level means to get the individual{-}level parameters}
\NormalTok{  indiv    }\OtherTok{\textless{}{-}} \FunctionTok{t}\NormalTok{(}\FunctionTok{matrix}\NormalTok{(}\FunctionTok{c}\NormalTok{(mu\_alpha, mu\_beta), }\AttributeTok{ncol =} \DecValTok{1}\NormalTok{) }\SpecialCharTok{\%*\%}
                  \FunctionTok{matrix}\NormalTok{(}\DecValTok{1}\NormalTok{, }\DecValTok{1}\NormalTok{, n\_agents) }\SpecialCharTok{+}\NormalTok{ Scale }\SpecialCharTok{\%*\%}\NormalTok{ Z)}

  \DocumentationTok{\#\# Then we generate the data for the memory model}
\NormalTok{  memory\_data }\OtherTok{\textless{}{-}} \FunctionTok{map\_dfr}\NormalTok{(}\FunctionTok{seq\_len}\NormalTok{(n\_agents), }\ControlFlowTok{function}\NormalTok{(j) \{}
\NormalTok{    logit\_p }\OtherTok{\textless{}{-}}\NormalTok{ indiv[j, }\DecValTok{1}\NormalTok{] }\SpecialCharTok{+}\NormalTok{ indiv[j, }\DecValTok{2}\NormalTok{] }\SpecialCharTok{*} \FunctionTok{qlogis}\NormalTok{(opp\_rate\_prev)}
    \FunctionTok{tibble}\NormalTok{(}\AttributeTok{agent\_id =}\NormalTok{ j, }\AttributeTok{trial =} \FunctionTok{seq\_len}\NormalTok{(n\_trials),}
           \AttributeTok{choice =} \FunctionTok{rbinom}\NormalTok{(n\_trials, }\DecValTok{1}\NormalTok{, }\FunctionTok{plogis}\NormalTok{(logit\_p)),}
           \AttributeTok{opp\_rate\_prev =}\NormalTok{ opp\_rate\_prev, }\AttributeTok{agent\_type =} \StringTok{"memory"}\NormalTok{,}
           \AttributeTok{true\_alpha =}\NormalTok{ indiv[j, }\DecValTok{1}\NormalTok{], }\AttributeTok{true\_beta =}\NormalTok{ indiv[j, }\DecValTok{2}\NormalTok{])}
\NormalTok{  \})}

  \FunctionTok{list}\NormalTok{(}\AttributeTok{biased =}\NormalTok{ biased\_data, }\AttributeTok{memory =}\NormalTok{ memory\_data,}
       \AttributeTok{hyperparams =} \FunctionTok{list}\NormalTok{(}\AttributeTok{mu\_theta =}\NormalTok{ mu\_theta, }\AttributeTok{sigma\_theta =}\NormalTok{ sigma\_theta,}
                          \AttributeTok{mu\_alpha =}\NormalTok{ mu\_alpha, }\AttributeTok{mu\_beta =}\NormalTok{ mu\_beta,}
                          \AttributeTok{sigma\_alpha =}\NormalTok{ sigma\_alpha, }\AttributeTok{sigma\_beta =}\NormalTok{ sigma\_beta,}
                          \AttributeTok{rho =}\NormalTok{ rho))}
\NormalTok{\}}

\NormalTok{sim\_data\_path }\OtherTok{\textless{}{-}}\NormalTok{ here}\SpecialCharTok{::}\FunctionTok{here}\NormalTok{(}\StringTok{"simdata"}\NormalTok{, }\StringTok{"ch08\_model\_recovery\_data.rds"}\NormalTok{)}

\ControlFlowTok{if}\NormalTok{ (regenerate\_simulations }\SpecialCharTok{||} \SpecialCharTok{!}\FunctionTok{file.exists}\NormalTok{(sim\_data\_path)) \{}
\NormalTok{  sim\_data }\OtherTok{\textless{}{-}} \FunctionTok{simulate\_model\_recovery\_data}\NormalTok{(}
    \AttributeTok{opponent\_choices =}\NormalTok{ opponent\_reversal   }\CommentTok{\# Use Ch. 7 reversal schedule}
\NormalTok{  )}
  \FunctionTok{saveRDS}\NormalTok{(sim\_data, sim\_data\_path)}
\NormalTok{\} }\ControlFlowTok{else}\NormalTok{ \{}
\NormalTok{  sim\_data }\OtherTok{\textless{}{-}} \FunctionTok{readRDS}\NormalTok{(sim\_data\_path)}
\NormalTok{\}}

\NormalTok{d\_biased }\OtherTok{\textless{}{-}}\NormalTok{ sim\_data}\SpecialCharTok{$}\NormalTok{biased}
\NormalTok{d\_memory }\OtherTok{\textless{}{-}}\NormalTok{ sim\_data}\SpecialCharTok{$}\NormalTok{memory}
\FunctionTok{cat}\NormalTok{(}\StringTok{"Biased data: N ="}\NormalTok{, }\FunctionTok{nrow}\NormalTok{(d\_biased), }\StringTok{"from"}\NormalTok{, }\FunctionTok{n\_distinct}\NormalTok{(d\_biased}\SpecialCharTok{$}\NormalTok{agent\_id), }\StringTok{"agents}\SpecialCharTok{\textbackslash{}n}\StringTok{"}\NormalTok{)}
\end{Highlighting}
\end{Shaded}

\begin{verbatim}
Biased data: N = 6000 from 50 agents
\end{verbatim}

\begin{Shaded}
\begin{Highlighting}[]
\FunctionTok{cat}\NormalTok{(}\StringTok{"Memory data: N ="}\NormalTok{, }\FunctionTok{nrow}\NormalTok{(d\_memory), }\StringTok{"from"}\NormalTok{, }\FunctionTok{n\_distinct}\NormalTok{(d\_memory}\SpecialCharTok{$}\NormalTok{agent\_id), }\StringTok{"agents}\SpecialCharTok{\textbackslash{}n}\StringTok{"}\NormalTok{)}
\end{Highlighting}
\end{Shaded}

\begin{verbatim}
Memory data: N = 6000 from 50 agents
\end{verbatim}

Let us visualize a sample of agents from both populations before fitting
any models. This is a useful sanity check, and part of any sensible
workflow. Check your data and simulations, folks!

\begin{Shaded}
\begin{Highlighting}[]
\NormalTok{sampled\_agents }\OtherTok{\textless{}{-}} \FunctionTok{bind\_rows}\NormalTok{(d\_biased, d\_memory) }\SpecialCharTok{|\textgreater{}}
  \FunctionTok{distinct}\NormalTok{(agent\_type, agent\_id) }\SpecialCharTok{|\textgreater{}}
  \FunctionTok{group\_by}\NormalTok{(agent\_type) }\SpecialCharTok{|\textgreater{}}
  \FunctionTok{slice\_sample}\NormalTok{(}\AttributeTok{n =} \DecValTok{8}\NormalTok{) }\SpecialCharTok{|\textgreater{}}
  \FunctionTok{ungroup}\NormalTok{()}

\FunctionTok{bind\_rows}\NormalTok{(d\_biased, d\_memory) }\SpecialCharTok{|\textgreater{}}
  \FunctionTok{semi\_join}\NormalTok{(sampled\_agents, }\AttributeTok{by =} \FunctionTok{c}\NormalTok{(}\StringTok{"agent\_type"}\NormalTok{, }\StringTok{"agent\_id"}\NormalTok{)) }\SpecialCharTok{|\textgreater{}}
  \FunctionTok{arrange}\NormalTok{(agent\_type, agent\_id, trial) }\SpecialCharTok{|\textgreater{}}
  \FunctionTok{group\_by}\NormalTok{(agent\_type, agent\_id) }\SpecialCharTok{|\textgreater{}}
  \FunctionTok{mutate}\NormalTok{(}\AttributeTok{cum\_rate =} \FunctionTok{cumsum}\NormalTok{(choice) }\SpecialCharTok{/}\NormalTok{ trial) }\SpecialCharTok{|\textgreater{}}
  \FunctionTok{ungroup}\NormalTok{() }\SpecialCharTok{|\textgreater{}}
  \FunctionTok{ggplot}\NormalTok{(}\FunctionTok{aes}\NormalTok{(}\AttributeTok{x =}\NormalTok{ trial, }\AttributeTok{y =}\NormalTok{ cum\_rate, }\AttributeTok{group =}\NormalTok{ agent\_id, }\AttributeTok{color =}\NormalTok{ agent\_type)) }\SpecialCharTok{+}
  \FunctionTok{geom\_hline}\NormalTok{(}\AttributeTok{yintercept =} \FloatTok{0.5}\NormalTok{, }\AttributeTok{linetype =} \StringTok{"dashed"}\NormalTok{, }\AttributeTok{color =} \StringTok{"gray60"}\NormalTok{,}
             \AttributeTok{linewidth =} \FloatTok{0.4}\NormalTok{) }\SpecialCharTok{+}
  \FunctionTok{geom\_line}\NormalTok{(}\AttributeTok{alpha =} \FloatTok{0.7}\NormalTok{, }\AttributeTok{linewidth =} \FloatTok{0.7}\NormalTok{) }\SpecialCharTok{+}
  \FunctionTok{facet\_wrap}\NormalTok{(}\SpecialCharTok{\textasciitilde{}}\NormalTok{agent\_type, }\AttributeTok{labeller =} \FunctionTok{as\_labeller}\NormalTok{(}
    \FunctionTok{c}\NormalTok{(}\AttributeTok{biased =} \StringTok{"Biased Agents (fixed bias)"}\NormalTok{,}
      \AttributeTok{memory =} \StringTok{"Memory Agents (track opponent)"}\NormalTok{))) }\SpecialCharTok{+}
  \FunctionTok{scale\_color\_brewer}\NormalTok{(}\AttributeTok{palette =} \StringTok{"Set1"}\NormalTok{, }\AttributeTok{guide =} \StringTok{"none"}\NormalTok{) }\SpecialCharTok{+}
  \FunctionTok{scale\_y\_continuous}\NormalTok{(}\AttributeTok{limits =} \FunctionTok{c}\NormalTok{(}\DecValTok{0}\NormalTok{, }\DecValTok{1}\NormalTok{), }\AttributeTok{labels =}\NormalTok{ scales}\SpecialCharTok{::}\NormalTok{percent) }\SpecialCharTok{+}
  \FunctionTok{labs}\NormalTok{(}\AttributeTok{title    =} \StringTok{"Cumulative Choice Rates: Biased vs. Memory Agents"}\NormalTok{,}
       \AttributeTok{subtitle =} \StringTok{"Both populations play the reversal design (80\% right → 20\% right at trial 61). Memory agents shift; biased agents don\textquotesingle{}t."}\NormalTok{,}
       \AttributeTok{x =} \StringTok{"Trial"}\NormalTok{, }\AttributeTok{y =} \StringTok{"Cumulative proportion \textquotesingle{}right\textquotesingle{}"}\NormalTok{)}
\end{Highlighting}
\end{Shaded}

\begin{center}
\includegraphics[width=0.8\linewidth,height=\textheight,keepaspectratio]{07-ModelComparison_files/figure-pdf/07-visual-inspection-1.pdf}
\end{center}

We should see a difference: Both populations initially drift upwards
(rightwards in a sense) one because it is biased that way, the other
because it initially plays against a 80\%-right opponent. Memory agents
then start drifting downwards as they track the reversal in the
opponent's behavior, while biased agents stabilize. But it is not always
easy to tell them apart by eye, especially if we look at individual
agents rather than averages. This is why we need formal model comparison
techniques.

Now we prepare the four Stan data lists for our 2x2 comparison. Note
that when fitting the memory model we need to add the
\texttt{opp\_rate\_prev} even to the biased-generated data, as the model
requires it.

\begin{Shaded}
\begin{Highlighting}[]
\NormalTok{make\_stan\_biased }\OtherTok{\textless{}{-}} \ControlFlowTok{function}\NormalTok{(d)}
  \FunctionTok{list}\NormalTok{(}\AttributeTok{N =} \FunctionTok{nrow}\NormalTok{(d), }\AttributeTok{J =} \FunctionTok{n\_distinct}\NormalTok{(d}\SpecialCharTok{$}\NormalTok{agent\_id),}
       \AttributeTok{agent =}\NormalTok{ d}\SpecialCharTok{$}\NormalTok{agent\_id, }\AttributeTok{h =} \FunctionTok{as.integer}\NormalTok{(d}\SpecialCharTok{$}\NormalTok{choice))}

\NormalTok{make\_stan\_memory }\OtherTok{\textless{}{-}} \ControlFlowTok{function}\NormalTok{(d)}
  \FunctionTok{list}\NormalTok{(}\AttributeTok{N =} \FunctionTok{nrow}\NormalTok{(d), }\AttributeTok{J =} \FunctionTok{n\_distinct}\NormalTok{(d}\SpecialCharTok{$}\NormalTok{agent\_id),}
       \AttributeTok{agent =}\NormalTok{ d}\SpecialCharTok{$}\NormalTok{agent\_id, }\AttributeTok{h =} \FunctionTok{as.integer}\NormalTok{(d}\SpecialCharTok{$}\NormalTok{choice),}
       \AttributeTok{opp\_rate\_prev =}\NormalTok{ d}\SpecialCharTok{$}\NormalTok{opp\_rate\_prev)}

\NormalTok{sd\_biased\_on\_biased }\OtherTok{\textless{}{-}} \FunctionTok{make\_stan\_biased}\NormalTok{(d\_biased)}
\NormalTok{sd\_memory\_on\_biased }\OtherTok{\textless{}{-}} \FunctionTok{make\_stan\_memory}\NormalTok{(d\_biased)}
\NormalTok{sd\_biased\_on\_memory }\OtherTok{\textless{}{-}} \FunctionTok{make\_stan\_biased}\NormalTok{(d\_memory)}
\NormalTok{sd\_memory\_on\_memory }\OtherTok{\textless{}{-}} \FunctionTok{make\_stan\_memory}\NormalTok{(d\_memory)}
\end{Highlighting}
\end{Shaded}

We reuse the NCP Stan models from Chapter 7. The key feature we need for
model comparison is the \texttt{log\_lik} vector in the
\texttt{generated\ quantities} block - this is the pointwise
log-likelihood for each observation that PSIS-LOO requires. This is why
we designed those models with \texttt{log\_lik} from the beginning: the
Chapter 7 workflow was already leading here.

\begin{Shaded}
\begin{Highlighting}[]
\NormalTok{stan\_file\_biased }\OtherTok{\textless{}{-}}\NormalTok{ here}\SpecialCharTok{::}\FunctionTok{here}\NormalTok{(}\StringTok{"stan"}\NormalTok{, }\StringTok{"ch07\_multilevel\_biased\_ncp.stan"}\NormalTok{)}

\ControlFlowTok{if}\NormalTok{ (}\SpecialCharTok{!}\FunctionTok{file.exists}\NormalTok{(stan\_file\_biased)) \{}
  \FunctionTok{writeLines}\NormalTok{(}\StringTok{\textquotesingle{}}
\StringTok{data \{ int\textless{}lower=1\textgreater{} N; int\textless{}lower=1\textgreater{} J;}
\StringTok{  array[N] int\textless{}lower=1,upper=J\textgreater{} agent;}
\StringTok{  array[N] int\textless{}lower=0,upper=1\textgreater{} h; \}}
\StringTok{parameters \{ real mu\_theta; real\textless{}lower=0\textgreater{} sigma\_theta; vector[J] z\_theta; \}}
\StringTok{transformed parameters \{ vector[J] theta\_logit = mu\_theta + z\_theta*sigma\_theta; \}}
\StringTok{model \{}
\StringTok{  target += normal\_lpdf(mu\_theta|0,1.5);}
\StringTok{  target += exponential\_lpdf(sigma\_theta|1);}
\StringTok{  target += std\_normal\_lpdf(z\_theta);}
\StringTok{  target += bernoulli\_logit\_lpmf(h|theta\_logit[agent]); \}}
\StringTok{generated quantities \{}
\StringTok{  real lprior = normal\_lpdf(mu\_theta|0,1.5)+exponential\_lpdf(sigma\_theta|1)}
\StringTok{              + std\_normal\_lpdf(z\_theta);}
\StringTok{  vector[N] log\_lik; array[N] int h\_post\_rep;}
\StringTok{  for (i in 1:N) \{}
\StringTok{    log\_lik[i]    = bernoulli\_logit\_lpmf(h[i]|theta\_logit[agent[i]]);}
\StringTok{    h\_post\_rep[i] = bernoulli\_logit\_rng(theta\_logit[agent[i]]); \} \}\textquotesingle{}}\NormalTok{,}
\NormalTok{    stan\_file\_biased)}
\NormalTok{\}}

\NormalTok{stan\_file\_memory }\OtherTok{\textless{}{-}}\NormalTok{ here}\SpecialCharTok{::}\FunctionTok{here}\NormalTok{(}\StringTok{"stan"}\NormalTok{, }\StringTok{"ch07\_multilevel\_memory\_ncp.stan"}\NormalTok{)}
\ControlFlowTok{if}\NormalTok{ (}\SpecialCharTok{!}\FunctionTok{file.exists}\NormalTok{(stan\_file\_memory)) \{}
  \FunctionTok{writeLines}\NormalTok{(}\StringTok{\textquotesingle{}}
\StringTok{data \{ int\textless{}lower=1\textgreater{} N; int\textless{}lower=1\textgreater{} J;}
\StringTok{  array[N] int\textless{}lower=1,upper=J\textgreater{} agent; array[N] int\textless{}lower=0,upper=1\textgreater{} h;}
\StringTok{  vector\textless{}lower=0.01,upper=0.99\textgreater{}[N] opp\_rate\_prev; \}}
\StringTok{parameters \{}
\StringTok{  real mu\_alpha; real mu\_beta; vector\textless{}lower=0\textgreater{}[2] sigma;}
\StringTok{  cholesky\_factor\_corr[2] L\_Omega; matrix[2,J] z; \}}
\StringTok{transformed parameters \{}
\StringTok{  matrix[J,2] indiv\_params;}
\StringTok{  \{ matrix[2,J] dev = diag\_pre\_multiply(sigma,L\_Omega)*z;}
\StringTok{    for (j in 1:J) \{}
\StringTok{      indiv\_params[j,1] = mu\_alpha+dev[1,j];}
\StringTok{      indiv\_params[j,2] = mu\_beta+dev[2,j]; \} \} \}}
\StringTok{model \{}
\StringTok{  target += normal\_lpdf(mu\_alpha|0,1); target += normal\_lpdf(mu\_beta|0,1);}
\StringTok{  target += exponential\_lpdf(sigma|1);}
\StringTok{  target += lkj\_corr\_cholesky\_lpdf(L\_Omega|2);}
\StringTok{  target += std\_normal\_lpdf(to\_vector(z));}
\StringTok{  vector[N] logit\_p;}
\StringTok{  for (i in 1:N)}
\StringTok{    logit\_p[i] = indiv\_params[agent[i],1]}
\StringTok{               + indiv\_params[agent[i],2]*logit(opp\_rate\_prev[i]);}
\StringTok{  target += bernoulli\_logit\_lpmf(h|logit\_p); \}}
\StringTok{generated quantities \{}
\StringTok{  matrix[2,2] Omega = multiply\_lower\_tri\_self\_transpose(L\_Omega);}
\StringTok{  real rho = Omega[1,2];}
\StringTok{  real lprior = normal\_lpdf(mu\_alpha|0,1)+normal\_lpdf(mu\_beta|0,1)}
\StringTok{              + exponential\_lpdf(sigma|1)+lkj\_corr\_cholesky\_lpdf(L\_Omega|2)}
\StringTok{              + std\_normal\_lpdf(to\_vector(z));}
\StringTok{  vector[N] log\_lik; array[N] int h\_post\_rep;}
\StringTok{  for (i in 1:N) \{}
\StringTok{    real lp = indiv\_params[agent[i],1]}
\StringTok{            + indiv\_params[agent[i],2]*logit(opp\_rate\_prev[i]);}
\StringTok{    log\_lik[i] = bernoulli\_logit\_lpmf(h[i]|lp);}
\StringTok{    h\_post\_rep[i] = bernoulli\_logit\_rng(lp); \} \}\textquotesingle{}}\NormalTok{,}
\NormalTok{    stan\_file\_memory)}
\NormalTok{\}}

\NormalTok{mod\_biased\_ncp }\OtherTok{\textless{}{-}} \FunctionTok{cmdstan\_model}\NormalTok{(stan\_file\_biased)}
\NormalTok{mod\_memory\_ncp }\OtherTok{\textless{}{-}} \FunctionTok{cmdstan\_model}\NormalTok{(stan\_file\_memory)}
\FunctionTok{cat}\NormalTok{(}\StringTok{"Both NCP models compiled.}\SpecialCharTok{\textbackslash{}n}\StringTok{"}\NormalTok{)}
\end{Highlighting}
\end{Shaded}

\begin{verbatim}
Both NCP models compiled.
\end{verbatim}

\section{Fitting the Four Model-Data
Combinations}\label{fitting-the-four-model-data-combinations}

\begin{quote}
\emph{Why fit both models to both datasets?} Fitting only the true model
to its own data tells us the model works --- but it cannot tell us
whether ELPD actually discriminates. For model comparison to be
meaningful, the wrong model must score \emph{worse} than the right one.
Fitting all four combinations gives us a 2×2 confusion matrix: if
PSIS-LOO correctly selects the generating model in each column, the
method has recovered the truth. If it selects the wrong model, that is
scientifically informative too --- it tells us the experimental design
is not discriminative enough at this sample size.
\end{quote}

We now fit all four combinations. This is the computationally expensive
step, but it only needs to be run once --- thereafter, results are
loaded from disk. The \texttt{fit\_or\_load()} helper encapsulates this
logic.

\begin{Shaded}
\begin{Highlighting}[]
\NormalTok{fit\_or\_load }\OtherTok{\textless{}{-}} \ControlFlowTok{function}\NormalTok{(label, model, data\_list, file\_path, init\_fn, }\AttributeTok{seed =} \DecValTok{123}\NormalTok{) \{}
  \ControlFlowTok{if}\NormalTok{ (regenerate\_simulations }\SpecialCharTok{||} \SpecialCharTok{!}\FunctionTok{file.exists}\NormalTok{(file\_path)) \{}
    \FunctionTok{message}\NormalTok{(}\StringTok{"Fitting: "}\NormalTok{, label)}
\NormalTok{    fit }\OtherTok{\textless{}{-}}\NormalTok{ model}\SpecialCharTok{$}\FunctionTok{sample}\NormalTok{(}
      \AttributeTok{data =}\NormalTok{ data\_list, }\AttributeTok{seed =}\NormalTok{ seed, }\AttributeTok{chains =} \DecValTok{4}\NormalTok{, }\AttributeTok{parallel\_chains =} \DecValTok{4}\NormalTok{,}
      \AttributeTok{iter\_warmup =} \DecValTok{1000}\NormalTok{, }\AttributeTok{iter\_sampling =} \DecValTok{1000}\NormalTok{, }\AttributeTok{init =} \FunctionTok{init\_fn}\NormalTok{(data\_list),}
      \AttributeTok{adapt\_delta =} \FloatTok{0.90}\NormalTok{, }\AttributeTok{max\_treedepth =} \DecValTok{12}\NormalTok{, }\AttributeTok{refresh =} \DecValTok{0}\NormalTok{)}
\NormalTok{    fit}\SpecialCharTok{$}\FunctionTok{save\_object}\NormalTok{(file\_path)}
\NormalTok{  \} }\ControlFlowTok{else}\NormalTok{ \{}
\NormalTok{    fit }\OtherTok{\textless{}{-}} \FunctionTok{readRDS}\NormalTok{(file\_path)}
    \FunctionTok{message}\NormalTok{(}\StringTok{"Loaded: "}\NormalTok{, label)}
\NormalTok{  \}}
\NormalTok{  fit}
\NormalTok{\}}

\NormalTok{init\_biased }\OtherTok{\textless{}{-}} \ControlFlowTok{function}\NormalTok{(data\_list) \{}
\NormalTok{  J }\OtherTok{\textless{}{-}}\NormalTok{ data\_list}\SpecialCharTok{$}\NormalTok{J}
  \FunctionTok{lapply}\NormalTok{(}\DecValTok{1}\SpecialCharTok{:}\DecValTok{4}\NormalTok{, }\ControlFlowTok{function}\NormalTok{(i) }\FunctionTok{list}\NormalTok{(}\AttributeTok{mu\_theta =} \FunctionTok{rnorm}\NormalTok{(}\DecValTok{1}\NormalTok{,}\DecValTok{0}\NormalTok{,}\FloatTok{1.5}\NormalTok{),}
    \AttributeTok{sigma\_theta =} \FunctionTok{rexp}\NormalTok{(}\DecValTok{1}\NormalTok{,}\DecValTok{1}\NormalTok{), }\AttributeTok{z\_theta =} \FunctionTok{rnorm}\NormalTok{(J,}\DecValTok{0}\NormalTok{,}\DecValTok{1}\NormalTok{)))}
\NormalTok{\}}

\NormalTok{init\_memory }\OtherTok{\textless{}{-}} \ControlFlowTok{function}\NormalTok{(data\_list) \{}
\NormalTok{  J }\OtherTok{\textless{}{-}}\NormalTok{ data\_list}\SpecialCharTok{$}\NormalTok{J}
  \FunctionTok{lapply}\NormalTok{(}\DecValTok{1}\SpecialCharTok{:}\DecValTok{4}\NormalTok{, }\ControlFlowTok{function}\NormalTok{(i) }\FunctionTok{list}\NormalTok{(}
    \AttributeTok{mu\_alpha =} \FunctionTok{rnorm}\NormalTok{(}\DecValTok{1}\NormalTok{,}\DecValTok{0}\NormalTok{,}\FloatTok{0.3}\NormalTok{), }\AttributeTok{mu\_beta =} \FunctionTok{rnorm}\NormalTok{(}\DecValTok{1}\NormalTok{,}\DecValTok{0}\NormalTok{,}\FloatTok{0.3}\NormalTok{),}
    \AttributeTok{sigma =} \FunctionTok{abs}\NormalTok{(}\FunctionTok{rnorm}\NormalTok{(}\DecValTok{2}\NormalTok{,}\DecValTok{0}\NormalTok{,}\FloatTok{0.3}\NormalTok{)) }\SpecialCharTok{+} \FloatTok{0.1}\NormalTok{,}
    \AttributeTok{L\_Omega =} \FunctionTok{matrix}\NormalTok{(}\FunctionTok{c}\NormalTok{(}\DecValTok{1}\NormalTok{,}\DecValTok{0}\NormalTok{,}\DecValTok{0}\NormalTok{,}\DecValTok{1}\NormalTok{), }\AttributeTok{nrow =} \DecValTok{2}\NormalTok{),}
    \AttributeTok{z =} \FunctionTok{matrix}\NormalTok{(}\DecValTok{0}\NormalTok{, }\AttributeTok{nrow =} \DecValTok{2}\NormalTok{, }\AttributeTok{ncol =}\NormalTok{ J)))}
\NormalTok{\}}

\NormalTok{fit\_b\_on\_b }\OtherTok{\textless{}{-}} \FunctionTok{fit\_or\_load}\NormalTok{(}\StringTok{"Biased{-}\textgreater{}Biased"}\NormalTok{, mod\_biased\_ncp, sd\_biased\_on\_biased,}
\NormalTok{  here}\SpecialCharTok{::}\FunctionTok{here}\NormalTok{(}\StringTok{"simmodels"}\NormalTok{,}\StringTok{"ch08\_fit\_b\_on\_b.rds"}\NormalTok{), init\_biased)}
\NormalTok{fit\_m\_on\_b }\OtherTok{\textless{}{-}} \FunctionTok{fit\_or\_load}\NormalTok{(}\StringTok{"Memory{-}\textgreater{}Biased"}\NormalTok{, mod\_memory\_ncp, sd\_memory\_on\_biased,}
\NormalTok{  here}\SpecialCharTok{::}\FunctionTok{here}\NormalTok{(}\StringTok{"simmodels"}\NormalTok{,}\StringTok{"ch08\_fit\_m\_on\_b.rds"}\NormalTok{), init\_memory)}
\NormalTok{fit\_b\_on\_m }\OtherTok{\textless{}{-}} \FunctionTok{fit\_or\_load}\NormalTok{(}\StringTok{"Biased{-}\textgreater{}Memory"}\NormalTok{, mod\_biased\_ncp, sd\_biased\_on\_memory,}
\NormalTok{  here}\SpecialCharTok{::}\FunctionTok{here}\NormalTok{(}\StringTok{"simmodels"}\NormalTok{,}\StringTok{"ch08\_fit\_b\_on\_m.rds"}\NormalTok{), init\_biased)}
\NormalTok{fit\_m\_on\_m }\OtherTok{\textless{}{-}} \FunctionTok{fit\_or\_load}\NormalTok{(}\StringTok{"Memory{-}\textgreater{}Memory"}\NormalTok{, mod\_memory\_ncp, sd\_memory\_on\_memory,}
\NormalTok{  here}\SpecialCharTok{::}\FunctionTok{here}\NormalTok{(}\StringTok{"simmodels"}\NormalTok{,}\StringTok{"ch08\_fit\_m\_on\_m.rds"}\NormalTok{), init\_memory)}
\end{Highlighting}
\end{Shaded}

Before computing any model comparison statistics, we verify that all
four fits pass the standard geometry diagnostics from Chapter 6. This
step matters especially for the misspecified fits. When we fit the
memory model to biased data, we are asking a model to find a
memory-tracking signal where none exists - it may strain to do so, which
can sometimes produce divergences or poor mixing.

If we find issues, we need to think critically. On the one hand,
comparing a well fitted model to a poorly fitted one is not a strong
test of their relative performance. On the other hand, if the
misspecified model is so difficult to fit that it fails basic
diagnostics, that is itself a sign of poor performance. So as we
investigate the sources of the issues, we need to critically evaluate
whether they are a sign of a fundamental mismatch between the model and
the data-generating mechanism or just a technical issue that can be
resolved with more careful fitting.

Remember that we already run the full bayesian workflow on these models
in the previous chapters, so here we only briefly check whether these
new fitting processes show any sign of issues. If these were new models,
we would need to do a thorough validation first!

\begin{Shaded}
\begin{Highlighting}[]
\NormalTok{fits }\OtherTok{\textless{}{-}} \FunctionTok{list}\NormalTok{(}\StringTok{"Biased{-}\textgreater{}Biased"} \OtherTok{=}\NormalTok{ fit\_b\_on\_b, }\StringTok{"Memory{-}\textgreater{}Biased"} \OtherTok{=}\NormalTok{ fit\_m\_on\_b,}
             \StringTok{"Biased{-}\textgreater{}Memory"} \OtherTok{=}\NormalTok{ fit\_b\_on\_m, }\StringTok{"Memory{-}\textgreater{}Memory"} \OtherTok{=}\NormalTok{ fit\_m\_on\_m)}

\FunctionTok{map\_dfr}\NormalTok{(}\FunctionTok{names}\NormalTok{(fits), }\ControlFlowTok{function}\NormalTok{(label) \{}
\NormalTok{  f }\OtherTok{\textless{}{-}}\NormalTok{ fits[[label]]}
\NormalTok{  sm }\OtherTok{\textless{}{-}}\NormalTok{ f}\SpecialCharTok{$}\FunctionTok{diagnostic\_summary}\NormalTok{(}\AttributeTok{quiet =} \ConstantTok{TRUE}\NormalTok{)}
  
  \CommentTok{\# Compute both metrics in a single pass}
\NormalTok{  post\_summary }\OtherTok{\textless{}{-}} \FunctionTok{summarise\_draws}\NormalTok{(f}\SpecialCharTok{$}\FunctionTok{draws}\NormalTok{(), }\StringTok{"rhat"}\NormalTok{, }\StringTok{"ess\_bulk"}\NormalTok{)}
  
  \FunctionTok{tibble}\NormalTok{(}
    \AttributeTok{Combination  =}\NormalTok{ label,}
    \AttributeTok{Divergences  =} \FunctionTok{sum}\NormalTok{(sm}\SpecialCharTok{$}\NormalTok{num\_divergent),}
    \AttributeTok{Max\_Rhat     =} \FunctionTok{round}\NormalTok{(}\FunctionTok{max}\NormalTok{(post\_summary}\SpecialCharTok{$}\NormalTok{rhat, }\AttributeTok{na.rm =} \ConstantTok{TRUE}\NormalTok{), }\DecValTok{3}\NormalTok{),}
    \AttributeTok{Min\_ESS\_bulk =} \FunctionTok{round}\NormalTok{(}\FunctionTok{min}\NormalTok{(post\_summary}\SpecialCharTok{$}\NormalTok{ess\_bulk, }\AttributeTok{na.rm =} \ConstantTok{TRUE}\NormalTok{), }\DecValTok{0}\NormalTok{)}
\NormalTok{  )}
\NormalTok{\}) }\SpecialCharTok{|\textgreater{}}
\NormalTok{  knitr}\SpecialCharTok{::}\FunctionTok{kable}\NormalTok{(}\AttributeTok{caption =} \FunctionTok{paste}\NormalTok{(}
    \StringTok{"Pre{-}comparison diagnostics. Required: Divergences = 0, R{-}hat \textless{} 1.01, ESS \textgreater{} 400."}\NormalTok{,}
    \StringTok{"Misspecified fits (Memory{-}\textgreater{}Biased, Biased{-}\textgreater{}Memory) may show slightly worse ESS"}\NormalTok{,}
    \StringTok{"because the model is searching for structure that does not exist in the data."}\NormalTok{,}
    \StringTok{"This is informative: severe degradation here would itself be evidence against the model."}
\NormalTok{  ))}
\end{Highlighting}
\end{Shaded}

\begin{longtable}[]{@{}lrrr@{}}
\caption{Pre-comparison diagnostics. Required: Divergences = 0, R-hat
\textless{} 1.01, ESS \textgreater{} 400. Misspecified fits
(Memory-\textgreater Biased, Biased-\textgreater Memory) may show
slightly worse ESS because the model is searching for structure that
does not exist in the data. This is informative: severe degradation here
would itself be evidence against the model.}\tabularnewline
\toprule\noalign{}
Combination & Divergences & Max\_Rhat & Min\_ESS\_bulk \\
\midrule\noalign{}
\endfirsthead
\toprule\noalign{}
Combination & Divergences & Max\_Rhat & Min\_ESS\_bulk \\
\midrule\noalign{}
\endhead
\bottomrule\noalign{}
\endlastfoot
Biased-\textgreater Biased & 0 & 1.006 & 868 \\
Memory-\textgreater Biased & 0 & 1.006 & 963 \\
Biased-\textgreater Memory & 0 & 1.006 & 929 \\
Memory-\textgreater Memory & 0 & 1.006 & 707 \\
\end{longtable}

\section{Computing PSIS-LOO}\label{computing-psis-loo}

\subsection{\texorpdfstring{The Pareto-\(\hat{k}\) Diagnostic: Which
Observations Are
Influential?}{The Pareto-\textbackslash hat\{k\} Diagnostic: Which Observations Are Influential?}}\label{the-pareto-hatk-diagnostic-which-observations-are-influential}

The \texttt{loo()} function performs two critical operations: it
estimates the Expected Log Predictive Density (ELPD) via Pareto-smoothed
importance sampling, and it quantifies the mathematical reliability of
this approximation using the Pareto-\(\hat{k}\) diagnostic.

Recall that PSIS-LOO approximates exact leave-one-out cross-validation
by reweighting draws from the full-data posterior. The raw importance
weights are the inverse of the pointwise likelihoods. If an observation
\(y_i\) is highly unexpected by the model, its importance weight becomes
astronomically large. If a few posterior draws capture most of the
weight, the Monte Carlo approximation collapses (cannot represent the
full dataset).

To stabilize this, the algorithm fits a Generalized Pareto Distribution
to the right tail of these raw weights. The \(\hat{k}\) statistic is the
estimated shape parameter of this distribution, directly measuring the
heaviness of the tail. Conceptually, if removing a single trial would
drastically shift the posterior, extreme importance weights are required
to compensate, causing \(\hat{k}\) to blow up. Thus, \(\hat{k}\) serves
two vital purposes:

\begin{itemize}
\item
  Statistical diagnostic: It flags specific observations where the
  reweighting shortcut is mathematically unreliable.
\item
  Cognitive diagnostic: It identifies highly influential or surprising
  trials that violate the model's structural expectations.
\end{itemize}

We execute the function with save\_psis = TRUE to retain the smoothed
weights matrix, which is strictly required for calibration checks later
in the workflow (LOO-PIT, which will be explained then).

\begin{Shaded}
\begin{Highlighting}[]
\NormalTok{get\_raw\_loo }\OtherTok{\textless{}{-}} \ControlFlowTok{function}\NormalTok{(fit\_obj, file\_path) \{}
  \ControlFlowTok{if}\NormalTok{ (regenerate\_simulations }\SpecialCharTok{||} \SpecialCharTok{!}\FunctionTok{file.exists}\NormalTok{(file\_path)) \{}
    \FunctionTok{message}\NormalTok{(}\StringTok{"Computing and saving raw LOO to "}\NormalTok{, file\_path)}
\NormalTok{    loo\_obj }\OtherTok{\textless{}{-}}\NormalTok{ fit\_obj}\SpecialCharTok{$}\FunctionTok{loo}\NormalTok{(}\AttributeTok{save\_psis =} \ConstantTok{TRUE}\NormalTok{)}
    \FunctionTok{saveRDS}\NormalTok{(loo\_obj, file\_path)}
    \FunctionTok{return}\NormalTok{(loo\_obj)}
\NormalTok{  \} }\ControlFlowTok{else}\NormalTok{ \{}
    \FunctionTok{message}\NormalTok{(}\StringTok{"Loading cached raw LOO from "}\NormalTok{, file\_path)}
    \FunctionTok{return}\NormalTok{(}\FunctionTok{readRDS}\NormalTok{(file\_path))}
\NormalTok{  \}}
\NormalTok{\}}

\NormalTok{loo\_b\_on\_b }\OtherTok{\textless{}{-}} \FunctionTok{get\_raw\_loo}\NormalTok{(fit\_b\_on\_b, here}\SpecialCharTok{::}\FunctionTok{here}\NormalTok{(}\StringTok{"simmodels"}\NormalTok{, }\StringTok{"ch08\_loo\_raw\_b\_on\_b.rds"}\NormalTok{))}
\NormalTok{loo\_m\_on\_b }\OtherTok{\textless{}{-}} \FunctionTok{get\_raw\_loo}\NormalTok{(fit\_m\_on\_b, here}\SpecialCharTok{::}\FunctionTok{here}\NormalTok{(}\StringTok{"simmodels"}\NormalTok{, }\StringTok{"ch08\_loo\_raw\_m\_on\_b.rds"}\NormalTok{))}
\NormalTok{loo\_b\_on\_m }\OtherTok{\textless{}{-}} \FunctionTok{get\_raw\_loo}\NormalTok{(fit\_b\_on\_m, here}\SpecialCharTok{::}\FunctionTok{here}\NormalTok{(}\StringTok{"simmodels"}\NormalTok{, }\StringTok{"ch08\_loo\_raw\_b\_on\_m.rds"}\NormalTok{))}
\NormalTok{loo\_m\_on\_m }\OtherTok{\textless{}{-}} \FunctionTok{get\_raw\_loo}\NormalTok{(fit\_m\_on\_m, here}\SpecialCharTok{::}\FunctionTok{here}\NormalTok{(}\StringTok{"simmodels"}\NormalTok{, }\StringTok{"ch08\_loo\_raw\_m\_on\_m.rds"}\NormalTok{))}

\NormalTok{summarise\_pareto\_k }\OtherTok{\textless{}{-}} \ControlFlowTok{function}\NormalTok{(loo\_obj, label) \{}
\NormalTok{  k     }\OtherTok{\textless{}{-}}\NormalTok{ loo\_obj}\SpecialCharTok{$}\NormalTok{diagnostics}\SpecialCharTok{$}\NormalTok{pareto\_k}
\NormalTok{  p\_loo }\OtherTok{\textless{}{-}}\NormalTok{ loo\_obj}\SpecialCharTok{$}\NormalTok{estimates[}\StringTok{"p\_loo"}\NormalTok{,}\StringTok{"Estimate"}\NormalTok{]}
\NormalTok{  n     }\OtherTok{\textless{}{-}} \FunctionTok{length}\NormalTok{(k)}
  \FunctionTok{tibble}\NormalTok{(}\AttributeTok{Combination=}\NormalTok{label, }\AttributeTok{N=}\NormalTok{n,}
         \AttributeTok{k\_below\_05=}\FunctionTok{sum}\NormalTok{(k}\SpecialCharTok{\textless{}}\FloatTok{0.5}\NormalTok{), }\AttributeTok{k\_ch06\_07=}\FunctionTok{sum}\NormalTok{(k}\SpecialCharTok{\textgreater{}=}\FloatTok{0.5}\SpecialCharTok{\&}\NormalTok{k}\SpecialCharTok{\textless{}}\FloatTok{0.7}\NormalTok{),}
         \AttributeTok{k\_ch08\_10=}\FunctionTok{sum}\NormalTok{(k}\SpecialCharTok{\textgreater{}=}\FloatTok{0.7}\SpecialCharTok{\&}\NormalTok{k}\SpecialCharTok{\textless{}}\DecValTok{1}\NormalTok{), }\AttributeTok{k\_above\_10=}\FunctionTok{sum}\NormalTok{(k}\SpecialCharTok{\textgreater{}=}\DecValTok{1}\NormalTok{),}
         \AttributeTok{Max\_k=}\FunctionTok{round}\NormalTok{(}\FunctionTok{max}\NormalTok{(k),}\DecValTok{3}\NormalTok{),}
         \AttributeTok{p\_hat\_loo=}\FunctionTok{round}\NormalTok{(p\_loo,}\DecValTok{1}\NormalTok{), }\AttributeTok{ratio=}\FunctionTok{round}\NormalTok{(p\_loo}\SpecialCharTok{/}\NormalTok{n,}\DecValTok{3}\NormalTok{),}
         \AttributeTok{ELPD=}\FunctionTok{round}\NormalTok{(loo\_obj}\SpecialCharTok{$}\NormalTok{estimates[}\StringTok{"elpd\_loo"}\NormalTok{,}\StringTok{"Estimate"}\NormalTok{],}\DecValTok{1}\NormalTok{),}
         \AttributeTok{SE=}\FunctionTok{round}\NormalTok{(loo\_obj}\SpecialCharTok{$}\NormalTok{estimates[}\StringTok{"elpd\_loo"}\NormalTok{,}\StringTok{"SE"}\NormalTok{],}\DecValTok{1}\NormalTok{))}
\NormalTok{\}}

\NormalTok{pareto\_k\_table }\OtherTok{\textless{}{-}} \FunctionTok{bind\_rows}\NormalTok{(}
  \FunctionTok{summarise\_pareto\_k}\NormalTok{(loo\_b\_on\_b,}\StringTok{"Biased{-}\textgreater{}Biased"}\NormalTok{),}
  \FunctionTok{summarise\_pareto\_k}\NormalTok{(loo\_m\_on\_b,}\StringTok{"Memory{-}\textgreater{}Biased"}\NormalTok{),}
  \FunctionTok{summarise\_pareto\_k}\NormalTok{(loo\_b\_on\_m,}\StringTok{"Biased{-}\textgreater{}Memory"}\NormalTok{),}
  \FunctionTok{summarise\_pareto\_k}\NormalTok{(loo\_m\_on\_m,}\StringTok{"Memory{-}\textgreater{}Memory"}\NormalTok{))}

\NormalTok{knitr}\SpecialCharTok{::}\FunctionTok{kable}\NormalTok{(pareto\_k\_table,}
  \AttributeTok{caption =} \FunctionTok{paste}\NormalTok{(}\StringTok{"Pareto{-}k diagnostics. ratio = p\_hat\_loo/N measures model"}\NormalTok{,}
    \StringTok{"complexity relative to dataset size {-} we will use this for LOO{-}PIT below."}\NormalTok{))}
\end{Highlighting}
\end{Shaded}

\begin{longtable}[]{@{}
  >{\raggedright\arraybackslash}p{(\linewidth - 20\tabcolsep) * \real{0.1546}}
  >{\raggedleft\arraybackslash}p{(\linewidth - 20\tabcolsep) * \real{0.0515}}
  >{\raggedleft\arraybackslash}p{(\linewidth - 20\tabcolsep) * \real{0.1134}}
  >{\raggedleft\arraybackslash}p{(\linewidth - 20\tabcolsep) * \real{0.1031}}
  >{\raggedleft\arraybackslash}p{(\linewidth - 20\tabcolsep) * \real{0.1031}}
  >{\raggedleft\arraybackslash}p{(\linewidth - 20\tabcolsep) * \real{0.1134}}
  >{\raggedleft\arraybackslash}p{(\linewidth - 20\tabcolsep) * \real{0.0619}}
  >{\raggedleft\arraybackslash}p{(\linewidth - 20\tabcolsep) * \real{0.1031}}
  >{\raggedleft\arraybackslash}p{(\linewidth - 20\tabcolsep) * \real{0.0619}}
  >{\raggedleft\arraybackslash}p{(\linewidth - 20\tabcolsep) * \real{0.0825}}
  >{\raggedleft\arraybackslash}p{(\linewidth - 20\tabcolsep) * \real{0.0515}}@{}}
\caption{Pareto-k diagnostics. ratio = p\_hat\_loo/N measures model
complexity relative to dataset size - we will use this for LOO-PIT
below.}\tabularnewline
\toprule\noalign{}
\begin{minipage}[b]{\linewidth}\raggedright
Combination
\end{minipage} & \begin{minipage}[b]{\linewidth}\raggedleft
N
\end{minipage} & \begin{minipage}[b]{\linewidth}\raggedleft
k\_below\_05
\end{minipage} & \begin{minipage}[b]{\linewidth}\raggedleft
k\_ch06\_07
\end{minipage} & \begin{minipage}[b]{\linewidth}\raggedleft
k\_ch08\_10
\end{minipage} & \begin{minipage}[b]{\linewidth}\raggedleft
k\_above\_10
\end{minipage} & \begin{minipage}[b]{\linewidth}\raggedleft
Max\_k
\end{minipage} & \begin{minipage}[b]{\linewidth}\raggedleft
p\_hat\_loo
\end{minipage} & \begin{minipage}[b]{\linewidth}\raggedleft
ratio
\end{minipage} & \begin{minipage}[b]{\linewidth}\raggedleft
ELPD
\end{minipage} & \begin{minipage}[b]{\linewidth}\raggedleft
SE
\end{minipage} \\
\midrule\noalign{}
\endfirsthead
\toprule\noalign{}
\begin{minipage}[b]{\linewidth}\raggedright
Combination
\end{minipage} & \begin{minipage}[b]{\linewidth}\raggedleft
N
\end{minipage} & \begin{minipage}[b]{\linewidth}\raggedleft
k\_below\_05
\end{minipage} & \begin{minipage}[b]{\linewidth}\raggedleft
k\_ch06\_07
\end{minipage} & \begin{minipage}[b]{\linewidth}\raggedleft
k\_ch08\_10
\end{minipage} & \begin{minipage}[b]{\linewidth}\raggedleft
k\_above\_10
\end{minipage} & \begin{minipage}[b]{\linewidth}\raggedleft
Max\_k
\end{minipage} & \begin{minipage}[b]{\linewidth}\raggedleft
p\_hat\_loo
\end{minipage} & \begin{minipage}[b]{\linewidth}\raggedleft
ratio
\end{minipage} & \begin{minipage}[b]{\linewidth}\raggedleft
ELPD
\end{minipage} & \begin{minipage}[b]{\linewidth}\raggedleft
SE
\end{minipage} \\
\midrule\noalign{}
\endhead
\bottomrule\noalign{}
\endlastfoot
Biased-\textgreater Biased & 6000 & 6000 & 0 & 0 & 0 & 0.162 & 41.1 &
0.007 & -3393.8 & 36.0 \\
Memory-\textgreater Biased & 6000 & 5998 & 2 & 0 & 0 & 0.521 & 47.0 &
0.008 & -3395.1 & 36.0 \\
Biased-\textgreater Memory & 6000 & 6000 & 0 & 0 & 0 & 0.190 & 35.6 &
0.006 & -4038.3 & 15.4 \\
Memory-\textgreater Memory & 6000 & 5969 & 25 & 6 & 0 & 0.759 & 61.7 &
0.010 & -3700.3 & 27.3 \\
\end{longtable}

The action we take depends on how large \(\hat{k}\) is. The thresholds
below correspond roughly to how well the Pareto distribution fits the
tail of the importance weights --- which determines how reliable the
approximation is:

\begin{longtable}[]{@{}lll@{}}
\toprule\noalign{}
\(\hat{k}\) & Interpretation & Action \\
\midrule\noalign{}
\endhead
\bottomrule\noalign{}
\endlastfoot
\(< 0.5\) & Reliable & Accept PSIS-LOO \\
\([0.5, 0.7)\) & Acceptable & But proceed with caution \\
\([0.7, 1.0)\) & Problematic & Apply \texttt{loo\_moment\_match()}
first \\
\(\geq 1.0\) & Very problematic & Use actual K-fold \\
\end{longtable}

The \texttt{loo\_moment\_match()} refits the model excluding the
influential observations one by one, and assesses their out-of-sample
error under the refitted model. This brings two concerns: i)
computational feasibility: can we actually afford to refit the model
multiple times? ii) if there are multiple influential observations, they
might not be independent and only removing one would not capture the
full influence of whatever is causing the underlying issue. So once
again, we need to be critical and do our detective work. However, this
time we are lucky and it seems all is in place.

Notice also the \texttt{p\_hat\_loo} and \texttt{ratio} columns. These
are the effective number of parameters estimated by LOO and its ratio to
total observations - we will need these shortly when checking LOO-PIT
calibration.

The Effective Number of Parameters (\(\hat{p}_{\text{loo}}\)): You may
be familiar with Information Criteria like AIC or BIC, which penalize
models based on their raw parameter count (\(k\)). In Bayesian
hierarchical models, the ``true'' number of parameters is ambiguous
because of shrinkage. Instead, \(\hat{p}_{\text{loo}}\) measures the
model's optimism: the difference between how well the model fits the
training data vs.~how well it predicts held-out
data.\(\hat{p}_{\text{loo}} \approx\) total parameters: The model is
learning very little from the population
prior.\(\hat{p}_{\text{loo}} <\) total parameters: The hierarchical
prior is effectively ``narrowing'' the model's flexibility.PSIS-LOO is
essentially a more robust version of WAIC (Widely Applicable Information
Criterion) that includes the Pareto-\(\hat{k}\) diagnostic to tell you
when the penalty is unreliable.

The ratio \(\hat{p}_{\text{loo}}/N\) tells us the effective flexibility
per observation. This metric is critical for our next step: checking
predictive calibration via LOO-PIT. If this ratio is non-negligible,
removing a single observation drastically shifts the posterior, meaning
our \(N\) leave-one-out predictive distributions are highly correlated.
As Tesso and Vehtari (2026) demonstrated, this dependency invalidates
standard statistical tests for uniformity. We will explicitly rely on
this ratio to justify using specialized dependent-PIT tests (like POT-C)
to accurately diagnose model calibration.

\subsection{\texorpdfstring{Characterizing Influential Observations:
\emph{Why} Are They
Influential?}{Characterizing Influential Observations: Why Are They Influential?}}\label{characterizing-influential-observations-why-are-they-influential}

Finding that some observations have high \(\hat{k}\) is only the
beginning of the story. The scientifically interesting question is
\emph{why} those particular observations are influential. In a
hierarchical cognitive model, there are two structurally distinct
reasons an observation might be hard to predict, with different
theoretical implications:

\begin{enumerate}
\def\labelenumi{\arabic{enumi}.}
\item
  \textbf{The agent is an outlier from the population.} Their individual
  parameters (say, a very high memory sensitivity \(\beta_j\)) are far
  from the population mean. The model is surprised by their behavior not
  because anything unusual happened on that trial, but because the
  hierarchical prior, which uses all other agents to constrain this
  agent's parameters, is pulling the estimate toward the center. Finding
  this pattern suggests the population-level variance may be
  underestimated, or that some participants genuinely follow a different
  strategy than the majority.
\item
  \textbf{The trial itself is unusual for that agent.} Given the agent's
  estimated parameters, this specific observation is surprising, perhaps
  the memory signal was strong (opponent rate 90\%) and the model
  predicted the agent would strongly track it, but instead the agent
  chose randomly. This is a cognitively interesting failure: it
  identifies a trial where our model's assumptions break down,
  potentially revealing strategic flexibility, lapses of attention, or
  simply that agents sometimes deviate from their average behavior
  (e.g.~becoming distracted).
\end{enumerate}

To distinguish these two causes, we join the \(\hat{k}\) values back to
the raw data and look for structure across three dimensions: agent
identity, trial position, and the strength of the memory signal.

\begin{Shaded}
\begin{Highlighting}[]
\NormalTok{k\_mem\_on\_mem }\OtherTok{\textless{}{-}} \FunctionTok{tibble}\NormalTok{(}
  \AttributeTok{k             =}\NormalTok{ loo\_m\_on\_m}\SpecialCharTok{$}\NormalTok{diagnostics}\SpecialCharTok{$}\NormalTok{pareto\_k,}
  \AttributeTok{agent\_id      =}\NormalTok{ d\_memory}\SpecialCharTok{$}\NormalTok{agent\_id,}
  \AttributeTok{trial         =}\NormalTok{ d\_memory}\SpecialCharTok{$}\NormalTok{trial,}
  \AttributeTok{opp\_rate\_prev =}\NormalTok{ d\_memory}\SpecialCharTok{$}\NormalTok{opp\_rate\_prev,}
  \AttributeTok{fit           =} \StringTok{"Memory {-}\textgreater{} Memory"}\NormalTok{)}

\NormalTok{k\_bias\_on\_mem }\OtherTok{\textless{}{-}} \FunctionTok{tibble}\NormalTok{(}
  \AttributeTok{k             =}\NormalTok{ loo\_b\_on\_m}\SpecialCharTok{$}\NormalTok{diagnostics}\SpecialCharTok{$}\NormalTok{pareto\_k,}
  \AttributeTok{agent\_id      =}\NormalTok{ d\_memory}\SpecialCharTok{$}\NormalTok{agent\_id,}
  \AttributeTok{trial         =}\NormalTok{ d\_memory}\SpecialCharTok{$}\NormalTok{trial,}
  \AttributeTok{opp\_rate\_prev =}\NormalTok{ d\_memory}\SpecialCharTok{$}\NormalTok{opp\_rate\_prev,}
  \AttributeTok{fit           =} \StringTok{"Biased {-}\textgreater{} Memory"}\NormalTok{)}

\NormalTok{k\_annotated }\OtherTok{\textless{}{-}} \FunctionTok{bind\_rows}\NormalTok{(k\_mem\_on\_mem, k\_bias\_on\_mem)}

\CommentTok{\# {-}{-} Panel 1: Does influence cluster within specific agents? {-}{-}{-}{-}{-}{-}{-}{-}{-}{-}{-}{-}{-}{-}{-}{-}{-}{-}{-}}
\NormalTok{p\_by\_agent }\OtherTok{\textless{}{-}}\NormalTok{ k\_annotated }\SpecialCharTok{|\textgreater{}}
  \FunctionTok{group\_by}\NormalTok{(fit, agent\_id) }\SpecialCharTok{|\textgreater{}}
  \FunctionTok{summarise}\NormalTok{(}\AttributeTok{n\_high =} \FunctionTok{sum}\NormalTok{(k }\SpecialCharTok{\textgreater{}=} \FloatTok{0.5}\NormalTok{), }\AttributeTok{.groups =} \StringTok{"drop"}\NormalTok{) }\SpecialCharTok{|\textgreater{}}
  \FunctionTok{ggplot}\NormalTok{(}\FunctionTok{aes}\NormalTok{(}\AttributeTok{x =} \FunctionTok{reorder}\NormalTok{(}\FunctionTok{factor}\NormalTok{(agent\_id), }\SpecialCharTok{{-}}\NormalTok{n\_high),}
             \AttributeTok{y =}\NormalTok{ n\_high, }\AttributeTok{fill =}\NormalTok{ fit)) }\SpecialCharTok{+}
  \FunctionTok{geom\_col}\NormalTok{(}\AttributeTok{position =} \StringTok{"dodge"}\NormalTok{) }\SpecialCharTok{+}
  \FunctionTok{scale\_fill\_brewer}\NormalTok{(}\AttributeTok{palette =} \StringTok{"Set1"}\NormalTok{, }\AttributeTok{name =} \ConstantTok{NULL}\NormalTok{) }\SpecialCharTok{+}
  \FunctionTok{labs}\NormalTok{(}\AttributeTok{title    =} \StringTok{"High{-}k Observations per Agent (k \textgreater{}= 0.5)"}\NormalTok{,}
       \AttributeTok{subtitle =} \FunctionTok{paste}\NormalTok{(}\StringTok{"Concentration in specific agents suggests they deviate from"}\NormalTok{,}
                        \StringTok{"the population distribution."}\NormalTok{),}
       \AttributeTok{x =} \StringTok{"Agent ID (sorted by count)"}\NormalTok{, }\AttributeTok{y =} \StringTok{"Number of high{-}k trials"}\NormalTok{) }\SpecialCharTok{+}
  \FunctionTok{theme}\NormalTok{(}\AttributeTok{axis.text.x =} \FunctionTok{element\_text}\NormalTok{(}\AttributeTok{angle =} \DecValTok{90}\NormalTok{, }\AttributeTok{size =} \DecValTok{6}\NormalTok{),}
        \AttributeTok{legend.position =} \StringTok{"bottom"}\NormalTok{)}

\CommentTok{\# {-}{-} Panel 2: Does influence cluster at early trials? {-}{-}{-}{-}{-}{-}{-}{-}{-}{-}{-}{-}{-}{-}{-}{-}{-}{-}{-}{-}{-}{-}{-}{-}{-}}
\CommentTok{\# Early trials have uninformative memory (opp\_rate\_prev \textasciitilde{}= 0.5) so the}
\CommentTok{\# memory model has nothing to predict from, making any outcome surprising.}
\NormalTok{p\_by\_trial }\OtherTok{\textless{}{-}} \FunctionTok{ggplot}\NormalTok{(k\_annotated, }\FunctionTok{aes}\NormalTok{(}\AttributeTok{x =}\NormalTok{ trial, }\AttributeTok{y =}\NormalTok{ k, }\AttributeTok{color =}\NormalTok{ fit)) }\SpecialCharTok{+}
  \FunctionTok{geom\_point}\NormalTok{(}\AttributeTok{size =} \FloatTok{0.6}\NormalTok{, }\AttributeTok{alpha =} \FloatTok{0.3}\NormalTok{) }\SpecialCharTok{+}
  \FunctionTok{geom\_smooth}\NormalTok{(}\AttributeTok{method =} \StringTok{"gam"}\NormalTok{, }\AttributeTok{formula =}\NormalTok{ y }\SpecialCharTok{\textasciitilde{}} \FunctionTok{s}\NormalTok{(x, }\AttributeTok{bs =} \StringTok{"cs"}\NormalTok{),}
              \AttributeTok{linewidth =} \DecValTok{1}\NormalTok{, }\AttributeTok{se =} \ConstantTok{FALSE}\NormalTok{) }\SpecialCharTok{+}
  \FunctionTok{geom\_hline}\NormalTok{(}\AttributeTok{yintercept =} \FloatTok{0.5}\NormalTok{, }\AttributeTok{linetype =} \StringTok{"dashed"}\NormalTok{, }\AttributeTok{color =} \StringTok{"gray40"}\NormalTok{) }\SpecialCharTok{+}
  \FunctionTok{scale\_color\_brewer}\NormalTok{(}\AttributeTok{palette =} \StringTok{"Set1"}\NormalTok{, }\AttributeTok{name =} \ConstantTok{NULL}\NormalTok{) }\SpecialCharTok{+}
  \FunctionTok{labs}\NormalTok{(}\AttributeTok{title    =} \StringTok{"Pareto{-}k by Trial Position"}\NormalTok{,}
       \AttributeTok{subtitle =} \FunctionTok{paste}\NormalTok{(}\StringTok{"Both models should show higher k early: the memory signal"}\NormalTok{,}
                        \StringTok{"is uninformative (opp\_rate \textasciitilde{}= 0.5) so predictions are uncertain."}\NormalTok{),}
       \AttributeTok{x =} \StringTok{"Trial"}\NormalTok{, }\AttributeTok{y =} \FunctionTok{expression}\NormalTok{(}\FunctionTok{hat}\NormalTok{(k))) }\SpecialCharTok{+}
  \FunctionTok{theme}\NormalTok{(}\AttributeTok{legend.position =} \StringTok{"bottom"}\NormalTok{)}

\CommentTok{\# {-}{-} Panel 3: Does influence cluster at extreme opponent rates? {-}{-}{-}{-}{-}{-}{-}{-}{-}{-}{-}{-}{-}{-}{-}{-}}
\CommentTok{\# When opp\_rate\_prev is extreme (near 0 or 1), the memory model makes a}
\CommentTok{\# strong, confident prediction. If the agent does not follow the signal,}
\CommentTok{\# k will be high {-}{-} these are the most cognitively informative failures.}
\NormalTok{p\_by\_opp }\OtherTok{\textless{}{-}} \FunctionTok{ggplot}\NormalTok{(k\_annotated, }\FunctionTok{aes}\NormalTok{(}\AttributeTok{x =}\NormalTok{ opp\_rate\_prev, }\AttributeTok{y =}\NormalTok{ k, }\AttributeTok{color =}\NormalTok{ fit)) }\SpecialCharTok{+}
  \FunctionTok{geom\_point}\NormalTok{(}\AttributeTok{size =} \FloatTok{0.6}\NormalTok{, }\AttributeTok{alpha =} \FloatTok{0.3}\NormalTok{) }\SpecialCharTok{+}
  \FunctionTok{geom\_smooth}\NormalTok{(}\AttributeTok{method =} \StringTok{"gam"}\NormalTok{, }\AttributeTok{formula =}\NormalTok{ y }\SpecialCharTok{\textasciitilde{}} \FunctionTok{s}\NormalTok{(x, }\AttributeTok{bs =} \StringTok{"cs"}\NormalTok{),}
              \AttributeTok{linewidth =} \DecValTok{1}\NormalTok{, }\AttributeTok{se =} \ConstantTok{FALSE}\NormalTok{) }\SpecialCharTok{+}
  \FunctionTok{geom\_hline}\NormalTok{(}\AttributeTok{yintercept =} \FloatTok{0.5}\NormalTok{, }\AttributeTok{linetype =} \StringTok{"dashed"}\NormalTok{, }\AttributeTok{color =} \StringTok{"gray40"}\NormalTok{) }\SpecialCharTok{+}
  \FunctionTok{scale\_color\_brewer}\NormalTok{(}\AttributeTok{palette =} \StringTok{"Set1"}\NormalTok{, }\AttributeTok{name =} \ConstantTok{NULL}\NormalTok{) }\SpecialCharTok{+}
  \FunctionTok{labs}\NormalTok{(}\AttributeTok{title    =} \StringTok{"Pareto{-}k by Opponent\textquotesingle{}s Running Rate"}\NormalTok{,}
       \AttributeTok{subtitle =} \FunctionTok{paste}\NormalTok{(}\StringTok{"At extreme rates the memory model makes confident predictions."}\NormalTok{,}
                        \StringTok{"High k here = the agent did not follow the signal."}\NormalTok{),}
       \AttributeTok{x =} \StringTok{"Opponent running choice rate (the memory signal)"}\NormalTok{,}
       \AttributeTok{y =} \FunctionTok{expression}\NormalTok{(}\FunctionTok{hat}\NormalTok{(k))) }\SpecialCharTok{+}
  \FunctionTok{theme}\NormalTok{(}\AttributeTok{legend.position =} \StringTok{"bottom"}\NormalTok{)}

\NormalTok{p\_by\_agent }\SpecialCharTok{/}\NormalTok{ p\_by\_trial }\SpecialCharTok{/}\NormalTok{ p\_by\_opp }\SpecialCharTok{+}
  \FunctionTok{plot\_annotation}\NormalTok{(}\AttributeTok{title =} \StringTok{"Characterizing Influential Observations"}\NormalTok{,}
    \AttributeTok{subtitle =} \FunctionTok{paste}\NormalTok{(}\StringTok{"Reading all three panels together reveals the cognitive structure"}\NormalTok{,}
      \StringTok{"of model influence:}\SpecialCharTok{\textbackslash{}n}\StringTok{"}\NormalTok{,}
      \StringTok{"Panel 1: person{-}level outliers | Panel 2: temporal learning effects |"}\NormalTok{,}
      \StringTok{"Panel 3: overconfident predictions"}\NormalTok{))}
\end{Highlighting}
\end{Shaded}

\begin{center}
\includegraphics[width=0.8\linewidth,height=\textheight,keepaspectratio]{07-ModelComparison_files/figure-pdf/07-characterize-influential-1.pdf}
\end{center}

Reading these three panels together, we can build a cognitive story
about model influence.

\begin{itemize}
\item
  Panel 1 tells us whether specific agents are outliers from the group
  --- if a handful of agents account for most of the high-\(\hat{k}\)
  observations, those individuals may be genuinely using a different
  strategy from the majority.
\item
  Panel 2 tells us whether the model struggles early in a session,
  before it has accumulated useful opponent history: at trial 1, the
  running opponent rate is exactly 0.5 (the model knows nothing about
  the opponent), so any outcome is equally surprising --- this is
  expected and not a sign of model failure.
\item
  Panel 3 is the most cognitively interesting: when the opponent's
  running rate is extreme (say, 90\% right), the memory model predicts
  the agent will strongly track this. If instead the agent chose
  randomly or in the opposite direction, the model is very surprised ---
  \(\hat{k}\) spikes.
\end{itemize}

These are cognitively informative failures: they identify trials where
our model's assumption that agents track opponent history breaks down.
In our concrete case, we see that early trials are the most problematic.
Given the relatively high bias in the opponent and the few trials,
memory can quickly go to the extreme values and the model has a hard
time dealing with that.

But let's imagine we saw some agents being particularly influential ---
say, they had 10+ high-\(\hat{k}\) observations. This would suggest that
those agents are outliers from the population distribution --- perhaps
they have a very high memory sensitivity \(\beta_j\) that the
hierarchical prior is struggling to accommodate. In this case, we would
want to look at the posterior estimates for those agents and see if they
indeed have extreme parameter values. We could also consider refitting
the model with a more flexible hierarchical structure (e.g.,
heavier-tailed distributions) or even excluding those agents if we
believe they represent a qualitatively different strategy.

\begin{Shaded}
\begin{Highlighting}[]
\CommentTok{\# 1. Extract estimated parameters and transform to probability space (Median and 90\% HDI)}
\NormalTok{agent\_estimates }\OtherTok{\textless{}{-}}\NormalTok{ fit\_m\_on\_m}\SpecialCharTok{$}\FunctionTok{draws}\NormalTok{(}\StringTok{"indiv\_params"}\NormalTok{) }\SpecialCharTok{|\textgreater{}}
  \FunctionTok{spread\_draws}\NormalTok{(indiv\_params[agent\_id, param\_idx]) }\SpecialCharTok{|\textgreater{}}
  \CommentTok{\# Pivot to have alpha and beta in the same row per draw}
  \FunctionTok{pivot\_wider}\NormalTok{(}\AttributeTok{names\_from =}\NormalTok{ param\_idx, }\AttributeTok{values\_from =}\NormalTok{ indiv\_params, }\AttributeTok{names\_prefix =} \StringTok{"param\_"}\NormalTok{) }\SpecialCharTok{|\textgreater{}}
  \FunctionTok{rename}\NormalTok{(}\AttributeTok{alpha =}\NormalTok{ param\_1, }\AttributeTok{beta =}\NormalTok{ param\_2) }\SpecialCharTok{|\textgreater{}}
  \FunctionTok{mutate}\NormalTok{(}
    \CommentTok{\# Probability transformations evaluated at experimental block extremes}
    \AttributeTok{p\_bias   =} \FunctionTok{plogis}\NormalTok{(alpha),                           }\CommentTok{\# Baseline probability (opp = 0.5)}
    \AttributeTok{p\_block1 =} \FunctionTok{plogis}\NormalTok{(alpha }\SpecialCharTok{+}\NormalTok{ beta }\SpecialCharTok{*} \FunctionTok{qlogis}\NormalTok{(}\FloatTok{0.8}\NormalTok{)),      }\CommentTok{\# Prob of Right during 80\% block}
    \AttributeTok{p\_block2 =} \FunctionTok{plogis}\NormalTok{(alpha }\SpecialCharTok{+}\NormalTok{ beta }\SpecialCharTok{*} \FunctionTok{qlogis}\NormalTok{(}\FloatTok{0.2}\NormalTok{))       }\CommentTok{\# Prob of Right during 20\% block}
\NormalTok{  ) }\SpecialCharTok{|\textgreater{}}
  \FunctionTok{pivot\_longer}\NormalTok{(}
    \AttributeTok{cols =} \FunctionTok{c}\NormalTok{(alpha, beta, p\_bias, p\_block1, p\_block2), }
    \AttributeTok{names\_to =} \StringTok{"metric"}\NormalTok{, }
    \AttributeTok{values\_to =} \StringTok{"value"}
\NormalTok{  ) }\SpecialCharTok{|\textgreater{}}
  \FunctionTok{group\_by}\NormalTok{(agent\_id, metric) }\SpecialCharTok{|\textgreater{}}
  \FunctionTok{median\_hdi}\NormalTok{(value, }\AttributeTok{.width =} \FloatTok{0.90}\NormalTok{) }\SpecialCharTok{|\textgreater{}}
  \CommentTok{\# Format as "Median [Lower, Upper]"}
  \FunctionTok{mutate}\NormalTok{(}\AttributeTok{estimate\_label =} \FunctionTok{sprintf}\NormalTok{(}\StringTok{"\%.2f [\%.2f, \%.2f]"}\NormalTok{, value, .lower, .upper)) }\SpecialCharTok{|\textgreater{}}
\NormalTok{  dplyr}\SpecialCharTok{::}\FunctionTok{select}\NormalTok{(agent\_id, metric, estimate\_label) }\SpecialCharTok{|\textgreater{}}
  \FunctionTok{pivot\_wider}\NormalTok{(}\AttributeTok{names\_from =}\NormalTok{ metric, }\AttributeTok{values\_from =}\NormalTok{ estimate\_label, }\AttributeTok{names\_prefix =} \StringTok{"est\_"}\NormalTok{)}

\CommentTok{\# 2. Extract true parameters and compute true probabilities}
\NormalTok{agent\_truths }\OtherTok{\textless{}{-}}\NormalTok{ d\_memory }\SpecialCharTok{|\textgreater{}}
  \FunctionTok{distinct}\NormalTok{(agent\_id, true\_alpha, true\_beta) }\SpecialCharTok{|\textgreater{}}
  \FunctionTok{mutate}\NormalTok{(}
    \AttributeTok{true\_p\_bias   =} \FunctionTok{round}\NormalTok{(}\FunctionTok{plogis}\NormalTok{(true\_alpha), }\DecValTok{2}\NormalTok{),}
    \AttributeTok{true\_p\_block1 =} \FunctionTok{round}\NormalTok{(}\FunctionTok{plogis}\NormalTok{(true\_alpha }\SpecialCharTok{+}\NormalTok{ true\_beta }\SpecialCharTok{*} \FunctionTok{qlogis}\NormalTok{(}\FloatTok{0.8}\NormalTok{)), }\DecValTok{2}\NormalTok{),}
    \AttributeTok{true\_p\_block2 =} \FunctionTok{round}\NormalTok{(}\FunctionTok{plogis}\NormalTok{(true\_alpha }\SpecialCharTok{+}\NormalTok{ true\_beta }\SpecialCharTok{*} \FunctionTok{qlogis}\NormalTok{(}\FloatTok{0.2}\NormalTok{)), }\DecValTok{2}\NormalTok{),}
    \AttributeTok{true\_alpha    =} \FunctionTok{round}\NormalTok{(true\_alpha, }\DecValTok{2}\NormalTok{),}
    \AttributeTok{true\_beta     =} \FunctionTok{round}\NormalTok{(true\_beta, }\DecValTok{2}\NormalTok{)}
\NormalTok{  )}

\CommentTok{\# 3. Join with the Pareto{-}k summary}
\NormalTok{k\_annotated }\SpecialCharTok{|\textgreater{}}
  \FunctionTok{filter}\NormalTok{(fit }\SpecialCharTok{==} \StringTok{"Memory {-}\textgreater{} Memory"}\NormalTok{) }\SpecialCharTok{|\textgreater{}}
  \FunctionTok{group\_by}\NormalTok{(agent\_id) }\SpecialCharTok{|\textgreater{}}
  \FunctionTok{summarise}\NormalTok{(}
    \AttributeTok{n\_high\_k =} \FunctionTok{sum}\NormalTok{(k }\SpecialCharTok{\textgreater{}=} \FloatTok{0.5}\NormalTok{), }
    \AttributeTok{max\_k    =} \FunctionTok{round}\NormalTok{(}\FunctionTok{max}\NormalTok{(k), }\DecValTok{3}\NormalTok{),}
    \AttributeTok{.groups  =} \StringTok{"drop"}
\NormalTok{  ) }\SpecialCharTok{|\textgreater{}}
  \FunctionTok{arrange}\NormalTok{(}\FunctionTok{desc}\NormalTok{(n\_high\_k)) }\SpecialCharTok{|\textgreater{}}
  \FunctionTok{slice\_head}\NormalTok{(}\AttributeTok{n =} \DecValTok{8}\NormalTok{) }\SpecialCharTok{|\textgreater{}}
  \FunctionTok{left\_join}\NormalTok{(agent\_truths, }\AttributeTok{by =} \StringTok{"agent\_id"}\NormalTok{) }\SpecialCharTok{|\textgreater{}}
  \FunctionTok{left\_join}\NormalTok{(agent\_estimates, }\AttributeTok{by =} \StringTok{"agent\_id"}\NormalTok{) }\SpecialCharTok{|\textgreater{}}
  \CommentTok{\# Select columns in a logical display order}
\NormalTok{  dplyr}\SpecialCharTok{::}\FunctionTok{select}\NormalTok{(}
\NormalTok{    agent\_id, n\_high\_k, max\_k, }
    \FunctionTok{starts\_with}\NormalTok{(}\StringTok{"true\_alpha"}\NormalTok{), }\FunctionTok{starts\_with}\NormalTok{(}\StringTok{"est\_alpha"}\NormalTok{),}
    \FunctionTok{starts\_with}\NormalTok{(}\StringTok{"true\_beta"}\NormalTok{), }\FunctionTok{starts\_with}\NormalTok{(}\StringTok{"est\_beta"}\NormalTok{),}
    \FunctionTok{starts\_with}\NormalTok{(}\StringTok{"true\_p\_block1"}\NormalTok{), }\FunctionTok{starts\_with}\NormalTok{(}\StringTok{"est\_p\_block1"}\NormalTok{),}
    \FunctionTok{starts\_with}\NormalTok{(}\StringTok{"true\_p\_block2"}\NormalTok{), }\FunctionTok{starts\_with}\NormalTok{(}\StringTok{"est\_p\_block2"}\NormalTok{)}
\NormalTok{  ) }\SpecialCharTok{|\textgreater{}}
\NormalTok{  knitr}\SpecialCharTok{::}\FunctionTok{kable}\NormalTok{(}\AttributeTok{caption =} \FunctionTok{paste}\NormalTok{(}
    \StringTok{"Top 8 agents by number of high{-}k observations."}\NormalTok{,}
    \StringTok{"Probability space metrics (p\_block1, p\_block2) demonstrate agent determinism at environmental extremes."}\NormalTok{,}
    \StringTok{"Values near 1.00 or 0.00 with narrow HDIs indicate near{-}deterministic behavioral expectations,"}\NormalTok{,}
    \StringTok{"which might trigger high Pareto{-}k values upon any contradictory response."}
\NormalTok{  ))}
\end{Highlighting}
\end{Shaded}

\begin{longtable}[]{@{}
  >{\raggedleft\arraybackslash}p{(\linewidth - 20\tabcolsep) * \real{0.0608}}
  >{\raggedleft\arraybackslash}p{(\linewidth - 20\tabcolsep) * \real{0.0608}}
  >{\raggedleft\arraybackslash}p{(\linewidth - 20\tabcolsep) * \real{0.0405}}
  >{\raggedleft\arraybackslash}p{(\linewidth - 20\tabcolsep) * \real{0.0743}}
  >{\raggedright\arraybackslash}p{(\linewidth - 20\tabcolsep) * \real{0.1419}}
  >{\raggedleft\arraybackslash}p{(\linewidth - 20\tabcolsep) * \real{0.0676}}
  >{\raggedright\arraybackslash}p{(\linewidth - 20\tabcolsep) * \real{0.1216}}
  >{\raggedleft\arraybackslash}p{(\linewidth - 20\tabcolsep) * \real{0.0946}}
  >{\raggedright\arraybackslash}p{(\linewidth - 20\tabcolsep) * \real{0.1216}}
  >{\raggedleft\arraybackslash}p{(\linewidth - 20\tabcolsep) * \real{0.0946}}
  >{\raggedright\arraybackslash}p{(\linewidth - 20\tabcolsep) * \real{0.1216}}@{}}
\caption{Top 8 agents by number of high-k observations. Probability
space metrics (p\_block1, p\_block2) demonstrate agent determinism at
environmental extremes. Values near 1.00 or 0.00 with narrow HDIs
indicate near-deterministic behavioral expectations, which might trigger
high Pareto-k values upon any contradictory response.}\tabularnewline
\toprule\noalign{}
\begin{minipage}[b]{\linewidth}\raggedleft
agent\_id
\end{minipage} & \begin{minipage}[b]{\linewidth}\raggedleft
n\_high\_k
\end{minipage} & \begin{minipage}[b]{\linewidth}\raggedleft
max\_k
\end{minipage} & \begin{minipage}[b]{\linewidth}\raggedleft
true\_alpha
\end{minipage} & \begin{minipage}[b]{\linewidth}\raggedright
est\_alpha
\end{minipage} & \begin{minipage}[b]{\linewidth}\raggedleft
true\_beta
\end{minipage} & \begin{minipage}[b]{\linewidth}\raggedright
est\_beta
\end{minipage} & \begin{minipage}[b]{\linewidth}\raggedleft
true\_p\_block1
\end{minipage} & \begin{minipage}[b]{\linewidth}\raggedright
est\_p\_block1
\end{minipage} & \begin{minipage}[b]{\linewidth}\raggedleft
true\_p\_block2
\end{minipage} & \begin{minipage}[b]{\linewidth}\raggedright
est\_p\_block2
\end{minipage} \\
\midrule\noalign{}
\endfirsthead
\toprule\noalign{}
\begin{minipage}[b]{\linewidth}\raggedleft
agent\_id
\end{minipage} & \begin{minipage}[b]{\linewidth}\raggedleft
n\_high\_k
\end{minipage} & \begin{minipage}[b]{\linewidth}\raggedleft
max\_k
\end{minipage} & \begin{minipage}[b]{\linewidth}\raggedleft
true\_alpha
\end{minipage} & \begin{minipage}[b]{\linewidth}\raggedright
est\_alpha
\end{minipage} & \begin{minipage}[b]{\linewidth}\raggedleft
true\_beta
\end{minipage} & \begin{minipage}[b]{\linewidth}\raggedright
est\_beta
\end{minipage} & \begin{minipage}[b]{\linewidth}\raggedleft
true\_p\_block1
\end{minipage} & \begin{minipage}[b]{\linewidth}\raggedright
est\_p\_block1
\end{minipage} & \begin{minipage}[b]{\linewidth}\raggedleft
true\_p\_block2
\end{minipage} & \begin{minipage}[b]{\linewidth}\raggedright
est\_p\_block2
\end{minipage} \\
\midrule\noalign{}
\endhead
\bottomrule\noalign{}
\endlastfoot
1 & 2 & 0.518 & -0.02 & 0.17 {[}-0.13, 0.48{]} & 1.59 & 1.33 {[}0.77,
1.93{]} & 0.90 & 0.88 {[}0.80, 0.95{]} & 0.10 & 0.16 {[}0.05, 0.30{]} \\
7 & 2 & 0.620 & -0.16 & -0.12 {[}-0.40, 0.16{]} & 1.56 & 1.27 {[}0.74,
1.82{]} & 0.88 & 0.84 {[}0.74, 0.92{]} & 0.09 & 0.13 {[}0.05, 0.25{]} \\
9 & 2 & 0.651 & -0.12 & -0.40 {[}-0.70, -0.08{]} & 1.58 & 1.58 {[}0.97,
2.17{]} & 0.89 & 0.86 {[}0.77, 0.94{]} & 0.09 & 0.07 {[}0.01, 0.14{]} \\
14 & 2 & 0.673 & -0.01 & -0.20 {[}-0.50, 0.10{]} & 2.12 & 1.91 {[}1.29,
2.51{]} & 0.95 & 0.92 {[}0.86, 0.97{]} & 0.05 & 0.06 {[}0.01, 0.11{]} \\
16 & 2 & 0.714 & 0.07 & 0.19 {[}-0.12, 0.48{]} & 1.24 & 1.78 {[}1.17,
2.40{]} & 0.86 & 0.93 {[}0.88, 0.98{]} & 0.16 & 0.09 {[}0.03, 0.19{]} \\
18 & 2 & 0.511 & -0.29 & 0.13 {[}-0.16, 0.42{]} & 1.65 & 1.38 {[}0.81,
1.97{]} & 0.88 & 0.89 {[}0.81, 0.95{]} & 0.07 & 0.14 {[}0.05, 0.27{]} \\
19 & 2 & 0.759 & 0.37 & -0.35 {[}-0.69, -0.03{]} & 0.97 & 2.20 {[}1.58,
2.91{]} & 0.85 & 0.94 {[}0.89, 0.98{]} & 0.28 & 0.03 {[}0.01, 0.07{]} \\
25 & 2 & 0.702 & -0.29 & 0.26 {[}-0.06, 0.57{]} & 1.02 & 2.04 {[}1.36,
2.66{]} & 0.75 & 0.96 {[}0.92, 0.99{]} & 0.15 & 0.07 {[}0.02, 0.16{]} \\
\end{longtable}

\subsection{Moment Matching for Problematic
Observations}\label{moment-matching-for-problematic-observations}

When some observations have \(\hat{k} \geq 0.7\), the standard PSIS
approximation is unreliable for those specific data points. The first
line of defense is \textbf{moment matching} (Paananen et al., 2021): a
method that refines the importance weights by applying a transformation
to the posterior draws so that the moments of the reweighted
distribution better match what the LOO posterior should look like.

Think of it as improving the reweighting without running a full MCMC
from scratch. Instead of asking ``what would the posterior look like if
we had never seen observation \(i\)?'', moment matching asks ``can we
adjust the existing draws so they better approximate that LOO
posterior?'' This is substantially cheaper than refitting and works well
for moderately influential observations. Only if moment matching fails
to bring \(\hat{k}\) below 1.0 should we resort to agent-level K-fold
cross-validation (which we implement next).

\begin{Shaded}
\begin{Highlighting}[]
\NormalTok{apply\_moment\_match\_if\_needed }\OtherTok{\textless{}{-}} \ControlFlowTok{function}\NormalTok{(loo\_obj, fit\_obj, label, file\_path) \{}
  \CommentTok{\# 1. If the cached file exists, load it immediately}
  \ControlFlowTok{if}\NormalTok{ (}\SpecialCharTok{!}\NormalTok{regenerate\_simulations }\SpecialCharTok{\&\&} \FunctionTok{file.exists}\NormalTok{(file\_path)) \{}
    \FunctionTok{message}\NormalTok{(label, }\StringTok{": Loading cached LOO..."}\NormalTok{)}
    \FunctionTok{return}\NormalTok{(}\FunctionTok{readRDS}\NormalTok{(file\_path))}
\NormalTok{  \}}
  
\NormalTok{  n\_bad }\OtherTok{\textless{}{-}} \FunctionTok{sum}\NormalTok{(loo\_obj}\SpecialCharTok{$}\NormalTok{diagnostics}\SpecialCharTok{$}\NormalTok{pareto\_k }\SpecialCharTok{\textgreater{}=} \FloatTok{0.7}\NormalTok{)}
  
  \ControlFlowTok{if}\NormalTok{ (n\_bad }\SpecialCharTok{==} \DecValTok{0}\NormalTok{) \{}
    \FunctionTok{message}\NormalTok{(label, }\StringTok{": All k \textless{} 0.7. PSIS{-}LOO is reliable."}\NormalTok{)}
    \FunctionTok{saveRDS}\NormalTok{(loo\_obj, file\_path)}
    \FunctionTok{return}\NormalTok{(loo\_obj)}
\NormalTok{  \} }
  
  \FunctionTok{message}\NormalTok{(label, }\StringTok{": WARNING! "}\NormalTok{, n\_bad, }\StringTok{" observations with k \textgreater{}= 0.7."}\NormalTok{)}
  \FunctionTok{message}\NormalTok{(}\StringTok{"          Running CmdStanR moment matching to refine importance weights..."}\NormalTok{)}
  
  \CommentTok{\# 2. Run the actual moment matching algorithm}
  \CommentTok{\# (Requires cmdstanr \textgreater{}= 0.7.0 and models that can be unconstrained)}
\NormalTok{  loo\_mm }\OtherTok{\textless{}{-}}\NormalTok{ loo}\SpecialCharTok{::}\FunctionTok{loo\_moment\_match}\NormalTok{(}
\NormalTok{    fit\_obj, }
    \AttributeTok{loo =}\NormalTok{ loo\_obj}
\NormalTok{  )}
  
  \FunctionTok{message}\NormalTok{(}\StringTok{"          Max k reduced from "}\NormalTok{, }\FunctionTok{round}\NormalTok{(}\FunctionTok{max}\NormalTok{(loo\_obj}\SpecialCharTok{$}\NormalTok{diagnostics}\SpecialCharTok{$}\NormalTok{pareto\_k), }\DecValTok{2}\NormalTok{), }
          \StringTok{" to "}\NormalTok{, }\FunctionTok{round}\NormalTok{(}\FunctionTok{max}\NormalTok{(loo\_mm}\SpecialCharTok{$}\NormalTok{diagnostics}\SpecialCharTok{$}\NormalTok{pareto\_k), }\DecValTok{2}\NormalTok{))}
          
  \FunctionTok{saveRDS}\NormalTok{(loo\_mm, file\_path)}
  \FunctionTok{return}\NormalTok{(loo\_mm)}
\NormalTok{\}}

\CommentTok{\# Execute the caching function }
\NormalTok{loo\_b\_on\_b }\OtherTok{\textless{}{-}} \FunctionTok{apply\_moment\_match\_if\_needed}\NormalTok{(loo\_b\_on\_b, fit\_b\_on\_b, }\StringTok{"Biased{-}\textgreater{}Biased"}\NormalTok{, here}\SpecialCharTok{::}\FunctionTok{here}\NormalTok{(}\StringTok{"simmodels"}\NormalTok{, }\StringTok{"ch08\_loo\_mm\_b\_on\_b.rds"}\NormalTok{))}
\NormalTok{loo\_m\_on\_b }\OtherTok{\textless{}{-}} \FunctionTok{apply\_moment\_match\_if\_needed}\NormalTok{(loo\_m\_on\_b, fit\_m\_on\_b, }\StringTok{"Memory{-}\textgreater{}Biased"}\NormalTok{, here}\SpecialCharTok{::}\FunctionTok{here}\NormalTok{(}\StringTok{"simmodels"}\NormalTok{, }\StringTok{"ch08\_loo\_mm\_m\_on\_b.rds"}\NormalTok{))}
\NormalTok{loo\_b\_on\_m }\OtherTok{\textless{}{-}} \FunctionTok{apply\_moment\_match\_if\_needed}\NormalTok{(loo\_b\_on\_m, fit\_b\_on\_m, }\StringTok{"Biased{-}\textgreater{}Memory"}\NormalTok{, here}\SpecialCharTok{::}\FunctionTok{here}\NormalTok{(}\StringTok{"simmodels"}\NormalTok{, }\StringTok{"ch08\_loo\_mm\_b\_on\_m.rds"}\NormalTok{))}
\NormalTok{loo\_m\_on\_m }\OtherTok{\textless{}{-}} \FunctionTok{apply\_moment\_match\_if\_needed}\NormalTok{(loo\_m\_on\_m, fit\_m\_on\_m, }\StringTok{"Memory{-}\textgreater{}Memory"}\NormalTok{, here}\SpecialCharTok{::}\FunctionTok{here}\NormalTok{(}\StringTok{"simmodels"}\NormalTok{, }\StringTok{"ch08\_loo\_mm\_m\_on\_m.rds"}\NormalTok{))}
\end{Highlighting}
\end{Shaded}

\subsection{The Limits of PSIS: Agent-Level
LOO}\label{the-limits-of-psis-agent-level-loo}

If cognitive scientists are primarily interested in generalizing to new
participants, you might wonder: why can't we just calculate PSIS-LOO at
the participant level?

We actually can. To do this, we simply take our trial-level
log-likelihood matrix and sum the log-likelihoods of all 120 trials for
each agent \emph{before} passing the matrix to the \texttt{loo()}
function. This asks the PSIS algorithm to approximate what the posterior
would look like if an entire agent's dataset had been excluded.

Let's try this on our well-fitting
\texttt{Memory\ -\textgreater{}\ Memory} model.

\begin{Shaded}
\begin{Highlighting}[]
\CommentTok{\# 1. Extract the trial{-}level log{-}likelihood matrix (S draws x N trials)}
\NormalTok{ll\_trial\_m }\OtherTok{\textless{}{-}}\NormalTok{ fit\_m\_on\_m}\SpecialCharTok{$}\FunctionTok{draws}\NormalTok{(}\StringTok{"log\_lik"}\NormalTok{, }\AttributeTok{format =} \StringTok{"matrix"}\NormalTok{)}

\CommentTok{\# 2. Create an empty matrix for agent{-}level log{-}likelihoods (S draws x J agents)}
\NormalTok{J }\OtherTok{\textless{}{-}} \FunctionTok{n\_distinct}\NormalTok{(d\_memory}\SpecialCharTok{$}\NormalTok{agent\_id)}
\NormalTok{S }\OtherTok{\textless{}{-}} \FunctionTok{nrow}\NormalTok{(ll\_trial\_m)}
\NormalTok{ll\_agent\_m }\OtherTok{\textless{}{-}} \FunctionTok{matrix}\NormalTok{(}\DecValTok{0}\NormalTok{, }\AttributeTok{nrow =}\NormalTok{ S, }\AttributeTok{ncol =}\NormalTok{ J)}

\CommentTok{\# 3. Sum the trial log{-}likelihoods for each agent}
\ControlFlowTok{for}\NormalTok{ (j }\ControlFlowTok{in} \DecValTok{1}\SpecialCharTok{:}\NormalTok{J) \{}
  \CommentTok{\# Find which columns in the trial matrix belong to agent j}
\NormalTok{  trial\_cols }\OtherTok{\textless{}{-}} \FunctionTok{which}\NormalTok{(d\_memory}\SpecialCharTok{$}\NormalTok{agent\_id }\SpecialCharTok{==}\NormalTok{ j)}
  \CommentTok{\# Sum across those columns for every MCMC draw}
\NormalTok{  ll\_agent\_m[, j] }\OtherTok{\textless{}{-}} \FunctionTok{rowSums}\NormalTok{(ll\_trial\_m[, trial\_cols])}
\NormalTok{\}}

\CommentTok{\# 4. Run PSIS{-}LOO at the agent level}
\NormalTok{loo\_agent\_m }\OtherTok{\textless{}{-}}\NormalTok{ loo}\SpecialCharTok{::}\FunctionTok{loo}\NormalTok{(ll\_agent\_m)}
\FunctionTok{print}\NormalTok{(loo\_agent\_m)}
\end{Highlighting}
\end{Shaded}

\begin{verbatim}

Computed from 4000 by 50 log-likelihood matrix.

         Estimate   SE
elpd_loo  -3711.2 48.6
p_loo        57.1  4.1
looic      7422.4 97.2
------
MCSE of elpd_loo is NA.
MCSE and ESS estimates assume independent draws (r_eff=1).

Pareto k diagnostic values:
                         Count Pct.    Min. ESS
(-Inf, 0.7]   (good)     12    24.0%   136     
   (0.7, 1]   (bad)      30    60.0%   <NA>    
   (1, Inf)   (very bad)  8    16.0%   <NA>    
See help('pareto-k-diagnostic') for details.
\end{verbatim}

I f you run this code, the console should scream at you with warnings.
You will see that a massive percentage of the Pareto-\(\hat{k}\) values
are \(> 0.7\), and many are \(> 1.0\).This is a spectacular mathematical
failure, and it teaches us a crucial lesson about the limits of
Importance Sampling. The PSIS shortcut only works when the
``Leave-One-Out'' posterior is relatively similar to the full-data
posterior. Removing a single trial from a dataset of 6,000 barely shifts
the parameters. But removing an entire agent (and their 120 unique
trials) radically changes the population-level estimates. The shortcut
completely breaks down because the full-data posterior simply does not
contain enough information about what the population would look like
without that agent. This failure provides the exact justification for
our next step. When we want to test generalization to new agents, we are
forced to abandon the PSIS shortcut entirely and perform exact
cross-validation.

\section{Agent-Level K-Fold: Generalizing to New
Participants}\label{agent-level-k-fold-generalizing-to-new-participants}

Cross-validation requires us to make a fundamental choice: what defines
a ``data point'' when we hold data out?

Up until now, we have performed \textbf{trial-level} cross-validation
using the PSIS-LOO shortcut. Trial-level CV asks: \emph{``How well does
the model predict a held-out trial for a participant whose other 119
trials we have already observed and learned from?''}

But cognitive scientists usually want to answer a different question:
\emph{``How well does the model generalize to completely new
participants?''}. To answer this, we must perform \textbf{agent-level}
cross-validation, holding out entire individuals from the training set.

It is crucial to understand that the \emph{level} of cross-validation
(trial vs.~agent) is entirely separate from the \emph{computational
method} we use (PSIS approximation vs.~exact refitting).

In theory, we could perform agent-level LOO-CV (Leave-One-Agent-Out). We
could even try to use our PSIS mathematical shortcut to do it without
refitting the model. However, leaving out an entire agent's worth of
data (120 trials) radically changes the shape of the posterior
distribution. When the LOO posterior is that different from the
full-data posterior, the PSIS importance weights become wildly unstable
(\(\hat{k} \gg 1.0\)), and the shortcut completely fails.

Therefore, when we want to test generalization to new agents, we are
forced to abandon the PSIS shortcut and perform \textbf{exact
cross-validation}, where we actually recompile and refit the Stan model
on the training splits. Because exact Leave-One-Agent-Out would require
fitting the model \(J=50\) times, we compromise to save computation time
and use \textbf{exact agent-level K-fold CV}.

In agent-level K-fold, we hold out \(1/K\) of the agents entirely from
the training set, fit the model to the remaining agents, and then ask:
given \emph{only} the population parameters estimated from the training
group, how well can we predict the held-out strangers' behavior?

\begin{Shaded}
\begin{Highlighting}[]
\CommentTok{\# 1. Biased K{-}Fold Model}
\CommentTok{\# Draws ONE parameter per test agent from the population distribution.}
\FunctionTok{writeLines}\NormalTok{(}\StringTok{"}
\StringTok{data \{}
\StringTok{  int\textless{}lower=1\textgreater{} N\_train; }
\StringTok{  int\textless{}lower=1\textgreater{} J\_train;}
\StringTok{  array[N\_train] int\textless{}lower=1, upper=J\_train\textgreater{} agent\_train;}
\StringTok{  array[N\_train] int\textless{}lower=0, upper=1\textgreater{} h\_train;}
\StringTok{  }
\StringTok{  int\textless{}lower=1\textgreater{} N\_test; }
\StringTok{  int\textless{}lower=1\textgreater{} J\_test;}
\StringTok{  array[N\_test] int\textless{}lower=1, upper=J\_test\textgreater{} agent\_test;}
\StringTok{  array[N\_test] int\textless{}lower=0, upper=1\textgreater{} h\_test;}
\StringTok{\}}

\StringTok{parameters \{}
\StringTok{  real mu\_theta; }
\StringTok{  real\textless{}lower=0\textgreater{} sigma\_theta; }
\StringTok{  vector[J\_train] z\_theta;}
\StringTok{\}}

\StringTok{transformed parameters \{}
\StringTok{  vector[J\_train] theta\_logit = mu\_theta + z\_theta * sigma\_theta;}
\StringTok{\}}

\StringTok{model \{}
\StringTok{  target += normal\_lpdf(mu\_theta | 0, 1.5);}
\StringTok{  target += exponential\_lpdf(sigma\_theta | 1);}
\StringTok{  target += std\_normal\_lpdf(z\_theta);}
\StringTok{  target += bernoulli\_logit\_lpmf(h\_train | theta\_logit[agent\_train]);}
\StringTok{\}}

\StringTok{generated quantities \{}
\StringTok{  vector[J\_test] theta\_new;}
\StringTok{  vector[N\_test] log\_lik\_test;}
\StringTok{  }
\StringTok{  for (j in 1:J\_test) \{}
\StringTok{    theta\_new[j] = normal\_rng(mu\_theta, sigma\_theta);}
\StringTok{  \}}
\StringTok{  }
\StringTok{  for (n in 1:N\_test) \{}
\StringTok{    log\_lik\_test[n] = bernoulli\_logit\_lpmf(h\_test[n] | theta\_new[agent\_test[n]]);}
\StringTok{  \}}
\StringTok{\}"}\NormalTok{, here}\SpecialCharTok{::}\FunctionTok{here}\NormalTok{(}\StringTok{"stan"}\NormalTok{, }\StringTok{"ch08\_biased\_kfold.stan"}\NormalTok{))}

\CommentTok{\# 2. Memory K{-}Fold Model}
\CommentTok{\# Samples correlated parameter pairs for each test agent from the population prior.}
\FunctionTok{writeLines}\NormalTok{(}\StringTok{"}
\StringTok{data \{}
\StringTok{  int\textless{}lower=1\textgreater{} N\_train; }
\StringTok{  int\textless{}lower=1\textgreater{} J\_train;}
\StringTok{  array[N\_train] int\textless{}lower=1, upper=J\_train\textgreater{} agent\_train;}
\StringTok{  array[N\_train] int\textless{}lower=0, upper=1\textgreater{} h\_train;}
\StringTok{  vector\textless{}lower=0.01, upper=0.99\textgreater{}[N\_train] opp\_rate\_prev\_train;}

\StringTok{  int\textless{}lower=1\textgreater{} N\_test; }
\StringTok{  int\textless{}lower=1\textgreater{} J\_test;}
\StringTok{  array[N\_test] int\textless{}lower=1, upper=J\_test\textgreater{} agent\_test;}
\StringTok{  array[N\_test] int\textless{}lower=0, upper=1\textgreater{} h\_test;}
\StringTok{  vector\textless{}lower=0.01, upper=0.99\textgreater{}[N\_test] opp\_rate\_prev\_test;}
\StringTok{\}}

\StringTok{parameters \{}
\StringTok{  real mu\_alpha; }
\StringTok{  real mu\_beta; }
\StringTok{  vector\textless{}lower=0\textgreater{}[2] sigma;}
\StringTok{  cholesky\_factor\_corr[2] L\_Omega; }
\StringTok{  matrix[2, J\_train] z;}
\StringTok{\}}

\StringTok{transformed parameters \{}
\StringTok{  matrix[J\_train, 2] indiv\_params;}
\StringTok{  \{}
\StringTok{    matrix[2, J\_train] dev = diag\_pre\_multiply(sigma, L\_Omega) * z;}
\StringTok{    for (j in 1:J\_train) \{}
\StringTok{      indiv\_params[j, 1] = mu\_alpha + dev[1, j];}
\StringTok{      indiv\_params[j, 2] = mu\_beta + dev[2, j];}
\StringTok{    \}}
\StringTok{  \}}
\StringTok{\}}

\StringTok{model \{}
\StringTok{  target += normal\_lpdf(mu\_alpha | 0, 1); }
\StringTok{  target += normal\_lpdf(mu\_beta | 0, 1);}
\StringTok{  target += exponential\_lpdf(sigma | 1);}
\StringTok{  target += lkj\_corr\_cholesky\_lpdf(L\_Omega | 2);}
\StringTok{  target += std\_normal\_lpdf(to\_vector(z));}
\StringTok{  }
\StringTok{  vector[N\_train] logit\_p;}
\StringTok{  for (i in 1:N\_train) \{}
\StringTok{    logit\_p[i] = indiv\_params[agent\_train[i], 1] + }
\StringTok{                 indiv\_params[agent\_train[i], 2] * logit(opp\_rate\_prev\_train[i]);}
\StringTok{  \}}
\StringTok{  target += bernoulli\_logit\_lpmf(h\_train | logit\_p);}
\StringTok{\}}

\StringTok{generated quantities \{}
\StringTok{  matrix[2, J\_test] new\_indiv;}
\StringTok{  vector[N\_test] log\_lik\_test;}
\StringTok{  }
\StringTok{  for (j in 1:J\_test) \{}
\StringTok{    vector[2] z\_new;}
\StringTok{    z\_new[1] = std\_normal\_rng();}
\StringTok{    z\_new[2] = std\_normal\_rng();}
\StringTok{    vector[2] param\_new = [mu\_alpha, mu\_beta]\textquotesingle{} + diag\_pre\_multiply(sigma, L\_Omega) * z\_new;}
\StringTok{    new\_indiv[1, j] = param\_new[1];}
\StringTok{    new\_indiv[2, j] = param\_new[2];}
\StringTok{  \}}
\StringTok{  }
\StringTok{  for (n in 1:N\_test) \{}
\StringTok{    real lp = new\_indiv[1, agent\_test[n]] + }
\StringTok{              new\_indiv[2, agent\_test[n]] * logit(opp\_rate\_prev\_test[n]);}
\StringTok{    log\_lik\_test[n] = bernoulli\_logit\_lpmf(h\_test[n] | lp);}
\StringTok{  \}}
\StringTok{\}"}\NormalTok{, here}\SpecialCharTok{::}\FunctionTok{here}\NormalTok{(}\StringTok{"stan"}\NormalTok{, }\StringTok{"ch08\_memory\_kfold.stan"}\NormalTok{))}

\NormalTok{mod\_biased\_kfold }\OtherTok{\textless{}{-}} \FunctionTok{cmdstan\_model}\NormalTok{(here}\SpecialCharTok{::}\FunctionTok{here}\NormalTok{(}\StringTok{"stan"}\NormalTok{, }\StringTok{"ch08\_biased\_kfold.stan"}\NormalTok{))}
\NormalTok{mod\_memory\_kfold }\OtherTok{\textless{}{-}} \FunctionTok{cmdstan\_model}\NormalTok{(here}\SpecialCharTok{::}\FunctionTok{here}\NormalTok{(}\StringTok{"stan"}\NormalTok{, }\StringTok{"ch08\_memory\_kfold.stan"}\NormalTok{))}
\FunctionTok{cat}\NormalTok{(}\StringTok{"Both K{-}fold Stan models compiled.}\SpecialCharTok{\textbackslash{}n}\StringTok{"}\NormalTok{)}
\end{Highlighting}
\end{Shaded}

\begin{verbatim}
Both K-fold Stan models compiled.
\end{verbatim}

\begin{Shaded}
\begin{Highlighting}[]
\NormalTok{kfold\_path }\OtherTok{\textless{}{-}}\NormalTok{ here}\SpecialCharTok{::}\FunctionTok{here}\NormalTok{(}\StringTok{"simmodels"}\NormalTok{,}\StringTok{"ch08\_kfold\_comparison.rds"}\NormalTok{)}
\NormalTok{K }\OtherTok{\textless{}{-}} \DecValTok{5}\NormalTok{; }\FunctionTok{set.seed}\NormalTok{(}\DecValTok{21}\NormalTok{)}

\CommentTok{\# We will run this on the Memory data to see the true model win out of sample}
\NormalTok{agent\_ids   }\OtherTok{\textless{}{-}} \FunctionTok{sort}\NormalTok{(}\FunctionTok{unique}\NormalTok{(d\_memory}\SpecialCharTok{$}\NormalTok{agent\_id))}
\NormalTok{fold\_assign }\OtherTok{\textless{}{-}} \FunctionTok{sample}\NormalTok{(}\FunctionTok{rep}\NormalTok{(}\DecValTok{1}\SpecialCharTok{:}\NormalTok{K, }\AttributeTok{length.out=}\FunctionTok{length}\NormalTok{(agent\_ids)))}
\NormalTok{agent\_folds }\OtherTok{\textless{}{-}} \FunctionTok{tibble}\NormalTok{(}\AttributeTok{agent\_id=}\NormalTok{agent\_ids, }\AttributeTok{fold=}\NormalTok{fold\_assign)}

\ControlFlowTok{if}\NormalTok{ (regenerate\_kfold }\SpecialCharTok{||} \SpecialCharTok{!}\FunctionTok{file.exists}\NormalTok{(kfold\_path)) \{}
  \FunctionTok{message}\NormalTok{(}\StringTok{"Running 5{-}fold CV for Biased and Memory models. This takes a moment..."}\NormalTok{)}
  
\NormalTok{  kfold\_list }\OtherTok{\textless{}{-}} \FunctionTok{lapply}\NormalTok{(}\DecValTok{1}\SpecialCharTok{:}\NormalTok{K, }\ControlFlowTok{function}\NormalTok{(k) \{}
\NormalTok{    test\_ids  }\OtherTok{\textless{}{-}}\NormalTok{ agent\_folds}\SpecialCharTok{$}\NormalTok{agent\_id[agent\_folds}\SpecialCharTok{$}\NormalTok{fold }\SpecialCharTok{==}\NormalTok{ k]}
\NormalTok{    train\_ids }\OtherTok{\textless{}{-}}\NormalTok{ agent\_folds}\SpecialCharTok{$}\NormalTok{agent\_id[agent\_folds}\SpecialCharTok{$}\NormalTok{fold }\SpecialCharTok{!=}\NormalTok{ k]}
    
\NormalTok{    d\_train }\OtherTok{\textless{}{-}}\NormalTok{ d\_memory }\SpecialCharTok{|\textgreater{}} \FunctionTok{filter}\NormalTok{(agent\_id }\SpecialCharTok{\%in\%}\NormalTok{ train\_ids) }\SpecialCharTok{|\textgreater{}}
               \FunctionTok{mutate}\NormalTok{(}\AttributeTok{agent\_train =} \FunctionTok{as.integer}\NormalTok{(}\FunctionTok{factor}\NormalTok{(agent\_id)))}
\NormalTok{    d\_test  }\OtherTok{\textless{}{-}}\NormalTok{ d\_memory }\SpecialCharTok{|\textgreater{}} \FunctionTok{filter}\NormalTok{(agent\_id }\SpecialCharTok{\%in\%}\NormalTok{ test\_ids) }\SpecialCharTok{|\textgreater{}}
               \FunctionTok{mutate}\NormalTok{(}\AttributeTok{agent\_test =} \FunctionTok{as.integer}\NormalTok{(}\FunctionTok{factor}\NormalTok{(agent\_id)))}
    
    \CommentTok{\# Stan data lists}
\NormalTok{    sdata\_b }\OtherTok{\textless{}{-}} \FunctionTok{list}\NormalTok{(}
      \AttributeTok{N\_train =} \FunctionTok{nrow}\NormalTok{(d\_train), }\AttributeTok{J\_train =} \FunctionTok{length}\NormalTok{(train\_ids),}
      \AttributeTok{agent\_train =}\NormalTok{ d\_train}\SpecialCharTok{$}\NormalTok{agent\_train, }\AttributeTok{h\_train =} \FunctionTok{as.integer}\NormalTok{(d\_train}\SpecialCharTok{$}\NormalTok{choice),}
      \AttributeTok{N\_test =} \FunctionTok{nrow}\NormalTok{(d\_test), }\AttributeTok{J\_test =} \FunctionTok{length}\NormalTok{(test\_ids),}
      \AttributeTok{agent\_test =}\NormalTok{ d\_test}\SpecialCharTok{$}\NormalTok{agent\_test, }\AttributeTok{h\_test =} \FunctionTok{as.integer}\NormalTok{(d\_test}\SpecialCharTok{$}\NormalTok{choice)}
\NormalTok{    )}
    
\NormalTok{    sdata\_m }\OtherTok{\textless{}{-}} \FunctionTok{list}\NormalTok{(}
      \AttributeTok{N\_train =} \FunctionTok{nrow}\NormalTok{(d\_train), }\AttributeTok{J\_train =} \FunctionTok{length}\NormalTok{(train\_ids),}
      \AttributeTok{agent\_train =}\NormalTok{ d\_train}\SpecialCharTok{$}\NormalTok{agent\_train, }\AttributeTok{h\_train =} \FunctionTok{as.integer}\NormalTok{(d\_train}\SpecialCharTok{$}\NormalTok{choice),}
      \AttributeTok{opp\_rate\_prev\_train =}\NormalTok{ d\_train}\SpecialCharTok{$}\NormalTok{opp\_rate\_prev,}
      \AttributeTok{N\_test =} \FunctionTok{nrow}\NormalTok{(d\_test), }\AttributeTok{J\_test =} \FunctionTok{length}\NormalTok{(test\_ids),}
      \AttributeTok{agent\_test =}\NormalTok{ d\_test}\SpecialCharTok{$}\NormalTok{agent\_test, }\AttributeTok{h\_test =} \FunctionTok{as.integer}\NormalTok{(d\_test}\SpecialCharTok{$}\NormalTok{choice),}
      \AttributeTok{opp\_rate\_prev\_test =}\NormalTok{ d\_test}\SpecialCharTok{$}\NormalTok{opp\_rate\_prev}
\NormalTok{    )}
    
    \CommentTok{\# Fit models (using 2 chains to save rendering time)}
\NormalTok{    fit\_b }\OtherTok{\textless{}{-}}\NormalTok{ mod\_biased\_kfold}\SpecialCharTok{$}\FunctionTok{sample}\NormalTok{(}\AttributeTok{data =}\NormalTok{ sdata\_b, }\AttributeTok{seed =} \DecValTok{200}\SpecialCharTok{+}\NormalTok{k, }\AttributeTok{chains =} \DecValTok{2}\NormalTok{, }
                                     \AttributeTok{parallel\_chains =} \DecValTok{2}\NormalTok{, }\AttributeTok{iter\_warmup =} \DecValTok{1000}\NormalTok{, }
                                     \AttributeTok{iter\_sampling =} \DecValTok{1000}\NormalTok{, }\AttributeTok{adapt\_delta =} \FloatTok{0.90}\NormalTok{, }\AttributeTok{refresh =} \DecValTok{0}\NormalTok{)}
\NormalTok{    fit\_m }\OtherTok{\textless{}{-}}\NormalTok{ mod\_memory\_kfold}\SpecialCharTok{$}\FunctionTok{sample}\NormalTok{(}\AttributeTok{data =}\NormalTok{ sdata\_m, }\AttributeTok{seed =} \DecValTok{300}\SpecialCharTok{+}\NormalTok{k, }\AttributeTok{chains =} \DecValTok{2}\NormalTok{, }
                                     \AttributeTok{parallel\_chains =} \DecValTok{2}\NormalTok{, }\AttributeTok{iter\_warmup =} \DecValTok{1000}\NormalTok{, }
                                     \AttributeTok{iter\_sampling =} \DecValTok{1000}\NormalTok{, }\AttributeTok{adapt\_delta =} \FloatTok{0.90}\NormalTok{, }\AttributeTok{refresh =} \DecValTok{0}\NormalTok{)}
    
\NormalTok{    ll\_b }\OtherTok{\textless{}{-}}\NormalTok{ fit\_b}\SpecialCharTok{$}\FunctionTok{draws}\NormalTok{(}\StringTok{"log\_lik\_test"}\NormalTok{, }\AttributeTok{format =} \StringTok{"matrix"}\NormalTok{)}
\NormalTok{    ll\_m }\OtherTok{\textless{}{-}}\NormalTok{ fit\_m}\SpecialCharTok{$}\FunctionTok{draws}\NormalTok{(}\StringTok{"log\_lik\_test"}\NormalTok{, }\AttributeTok{format =} \StringTok{"matrix"}\NormalTok{)}
\NormalTok{    idx  }\OtherTok{\textless{}{-}} \FunctionTok{which}\NormalTok{(d\_memory}\SpecialCharTok{$}\NormalTok{agent\_id }\SpecialCharTok{\%in\%}\NormalTok{ test\_ids)}
    
    \FunctionTok{list}\NormalTok{(}\AttributeTok{ll\_b =}\NormalTok{ ll\_b, }\AttributeTok{ll\_m =}\NormalTok{ ll\_m, }\AttributeTok{indices =}\NormalTok{ idx, }\AttributeTok{k =}\NormalTok{ k)}
\NormalTok{  \})}
  
\NormalTok{  S\_b }\OtherTok{\textless{}{-}} \FunctionTok{nrow}\NormalTok{(kfold\_list[[}\DecValTok{1}\NormalTok{]]}\SpecialCharTok{$}\NormalTok{ll\_b)}
\NormalTok{  S\_m }\OtherTok{\textless{}{-}} \FunctionTok{nrow}\NormalTok{(kfold\_list[[}\DecValTok{1}\NormalTok{]]}\SpecialCharTok{$}\NormalTok{ll\_m)}
\NormalTok{  ll\_kfold\_b }\OtherTok{\textless{}{-}} \FunctionTok{matrix}\NormalTok{(}\ConstantTok{NA\_real\_}\NormalTok{, S\_b, }\FunctionTok{nrow}\NormalTok{(d\_memory))}
\NormalTok{  ll\_kfold\_m }\OtherTok{\textless{}{-}} \FunctionTok{matrix}\NormalTok{(}\ConstantTok{NA\_real\_}\NormalTok{, S\_m, }\FunctionTok{nrow}\NormalTok{(d\_memory))}
  
  \ControlFlowTok{for}\NormalTok{ (res }\ControlFlowTok{in}\NormalTok{ kfold\_list) \{}
\NormalTok{    ll\_kfold\_b[, res}\SpecialCharTok{$}\NormalTok{indices] }\OtherTok{\textless{}{-}}\NormalTok{ res}\SpecialCharTok{$}\NormalTok{ll\_b}
\NormalTok{    ll\_kfold\_m[, res}\SpecialCharTok{$}\NormalTok{indices] }\OtherTok{\textless{}{-}}\NormalTok{ res}\SpecialCharTok{$}\NormalTok{ll\_m}
\NormalTok{  \}}
  \FunctionTok{saveRDS}\NormalTok{(}\FunctionTok{list}\NormalTok{(}\AttributeTok{biased =}\NormalTok{ ll\_kfold\_b, }\AttributeTok{memory =}\NormalTok{ ll\_kfold\_m), kfold\_path)}
\NormalTok{\} }\ControlFlowTok{else}\NormalTok{ \{}
\NormalTok{  kfold\_res  }\OtherTok{\textless{}{-}} \FunctionTok{readRDS}\NormalTok{(kfold\_path)}
\NormalTok{  ll\_kfold\_b }\OtherTok{\textless{}{-}}\NormalTok{ kfold\_res}\SpecialCharTok{$}\NormalTok{biased}
\NormalTok{  ll\_kfold\_m }\OtherTok{\textless{}{-}}\NormalTok{ kfold\_res}\SpecialCharTok{$}\NormalTok{memory}
\NormalTok{\}}

\ControlFlowTok{if}\NormalTok{ (}\SpecialCharTok{!}\FunctionTok{any}\NormalTok{(}\FunctionTok{is.na}\NormalTok{(ll\_kfold\_b))) \{}
\NormalTok{  elpd\_kf\_b }\OtherTok{\textless{}{-}}\NormalTok{ loo}\SpecialCharTok{::}\FunctionTok{elpd}\NormalTok{(ll\_kfold\_b)}
\NormalTok{  elpd\_kf\_m }\OtherTok{\textless{}{-}}\NormalTok{ loo}\SpecialCharTok{::}\FunctionTok{elpd}\NormalTok{(ll\_kfold\_m)}
  
  \FunctionTok{cat}\NormalTok{(}\StringTok{"=== Trial{-}level LOO vs. Agent{-}level K{-}fold (Memory Data) ===}\SpecialCharTok{\textbackslash{}n\textbackslash{}n}\StringTok{"}\NormalTok{)}
  
  \FunctionTok{cat}\NormalTok{(}\FunctionTok{sprintf}\NormalTok{(}\StringTok{"Biased Model {-} PSIS{-}LOO:  \%.1f (SE=\%.1f)}\SpecialCharTok{\textbackslash{}n}\StringTok{"}\NormalTok{,}
\NormalTok{    loo\_b\_on\_m}\SpecialCharTok{$}\NormalTok{estimates[}\StringTok{"elpd\_loo"}\NormalTok{,}\StringTok{"Estimate"}\NormalTok{], loo\_b\_on\_m}\SpecialCharTok{$}\NormalTok{estimates[}\StringTok{"elpd\_loo"}\NormalTok{,}\StringTok{"SE"}\NormalTok{]))}
  \FunctionTok{cat}\NormalTok{(}\FunctionTok{sprintf}\NormalTok{(}\StringTok{"Biased Model {-} K{-}fold:    \%.1f (SE=\%.1f)}\SpecialCharTok{\textbackslash{}n\textbackslash{}n}\StringTok{"}\NormalTok{,}
\NormalTok{    elpd\_kf\_b}\SpecialCharTok{$}\NormalTok{estimates[}\StringTok{"elpd"}\NormalTok{,}\StringTok{"Estimate"}\NormalTok{], elpd\_kf\_b}\SpecialCharTok{$}\NormalTok{estimates[}\StringTok{"elpd"}\NormalTok{,}\StringTok{"SE"}\NormalTok{]))}
    
  \FunctionTok{cat}\NormalTok{(}\FunctionTok{sprintf}\NormalTok{(}\StringTok{"Memory Model {-} PSIS{-}LOO:  \%.1f (SE=\%.1f)}\SpecialCharTok{\textbackslash{}n}\StringTok{"}\NormalTok{,}
\NormalTok{    loo\_m\_on\_m}\SpecialCharTok{$}\NormalTok{estimates[}\StringTok{"elpd\_loo"}\NormalTok{,}\StringTok{"Estimate"}\NormalTok{], loo\_m\_on\_m}\SpecialCharTok{$}\NormalTok{estimates[}\StringTok{"elpd\_loo"}\NormalTok{,}\StringTok{"SE"}\NormalTok{]))}
  \FunctionTok{cat}\NormalTok{(}\FunctionTok{sprintf}\NormalTok{(}\StringTok{"Memory Model {-} K{-}fold:    \%.1f (SE=\%.1f)}\SpecialCharTok{\textbackslash{}n\textbackslash{}n}\StringTok{"}\NormalTok{,}
\NormalTok{    elpd\_kf\_m}\SpecialCharTok{$}\NormalTok{estimates[}\StringTok{"elpd"}\NormalTok{,}\StringTok{"Estimate"}\NormalTok{], elpd\_kf\_m}\SpecialCharTok{$}\NormalTok{estimates[}\StringTok{"elpd"}\NormalTok{,}\StringTok{"SE"}\NormalTok{]))}
    
  \FunctionTok{cat}\NormalTok{(}\StringTok{"=== K{-}Fold Model Comparison ===}\SpecialCharTok{\textbackslash{}n}\StringTok{"}\NormalTok{)}
  \FunctionTok{print}\NormalTok{(loo}\SpecialCharTok{::}\FunctionTok{loo\_compare}\NormalTok{(elpd\_kf\_m, elpd\_kf\_b))}
\NormalTok{\}}
\end{Highlighting}
\end{Shaded}

\begin{verbatim}
=== Trial-level LOO vs. Agent-level K-fold (Memory Data) ===

Biased Model - PSIS-LOO:  -4038.3 (SE=15.4)
Biased Model - K-fold:    -4075.0 (SE=12.8)

Memory Model - PSIS-LOO:  -3700.3 (SE=27.3)
Memory Model - K-fold:    -3773.2 (SE=25.2)

=== K-Fold Model Comparison ===
       elpd_diff se_diff
model1    0.0       0.0 
model2 -301.7      22.2 
\end{verbatim}

The Agent-Level K-fold ELPD will almost always be lower (worse) than the
trial-level LOO ELPD for both models. This is mathematically expected
and cognitively informative. Trial-level LOO asks: ``How well can I
predict trial 120 for an agent whose personal bias and memory
sensitivity I have already estimated from 119 trials?'' Agent-level
K-fold asks: ``How well can I predict behavior for a complete stranger,
relying purely on population averages?''

The gap between the two ELPDs quantifies exactly how much
individual-level variation matters. If the gap is massive, human
individuality dominates the cognitive process.

\section{LOO-PIT: Calibration Checks for Predictive
Distributions}\label{loo-pit-calibration-checks-for-predictive-distributions}

\subsection{What LOO-PIT Measures}\label{what-loo-pit-measures}

The Pareto-\(\hat{k}\) values tell us whether the \emph{importance
weights} are reliable: whether the PSIS approximation is numerically
stable. LOO-PIT tells us something different and complementary: whether
the model's \emph{predictive distributions} are well-calibrated: whether
the model's uncertainty matches reality.

A model can have perfectly reliable \(\hat{k}\) values and still be
systematically miscalibrated. For example, a biased model fitted to
memory data might produce low \(\hat{k}\) for every observation (the
model is not dramatically wrong on any single trial) while still being
consistently wrong in a systematic, diffuse way across all trials. The
Pareto-\(\hat{k}\) would not flag this; LOO-PIT will.

The Probability Integral Transform (PIT) builds on a simple but powerful
idea. Consider the LOO predictive distribution for observation \(y_i\):
the distribution the model assigns to \(y_i\) before seeing it, having
been fitted on all other observations. We ask: where does the actual
observed \(y_i\) fall within this distribution? Specifically, we compute
\(u_i\), the cumulative probability up to the observed value:

\[u_i = \int_{-\infty}^{y_i} p(\tilde{y} \mid y_{-i}) \, d\tilde{y}\]

This is the \emph{quantile} of the observation within its own LOO
predictive distribution. If the model assigned most of its predictive
probability below \(y_i\), then \(u_i\) is close to 1 --- the model was
surprised, expecting a lower value but getting a higher one. If the
model assigned most of its probability above \(y_i\), then \(u_i\) is
close to 0.

For a perfectly calibrated model, the model's predictions should be
neither systematically too high nor too low. Half the observations
should fall in the lower half of their predicted distributions, half in
the upper half. More precisely, the \(u_i\) values should be
\emph{uniformly distributed} on \([0, 1]\). Departures from uniformity
reveal specific failure modes:

\begin{itemize}
\item
  \textbf{U-shaped distribution} (too many values near 0 and 1): The
  model is \emph{overconfident} --- its predictive distributions are too
  narrow. Observations consistently fall in the tails, surprising the
  model.
\item
  \textbf{Bell-shaped distribution} (too many values near 0.5): The
  model is \emph{underconfident} --- its predictive distributions are
  too wide, spreading probability mass over values that never occur.
\item
  \textbf{Skewed distribution} (more mass on one side): The model is
  \emph{biased}, systematically predicting too high or too low.
\end{itemize}

If this sounds familiar, you have been paying attention! It's the same
principle as in simulation-based-calibration, but applied to
observations instead of parameters. The difference is that in SBC we ask
``where does the true parameter value fall within the posterior
distribution?'' In LOO-PIT we ask ``where does the observed data point
fall within the predictive distribution?'' Both are checks of
calibration --- whether the model's uncertainty matches reality --- but
they apply to different distributions and reveal different failure
modes.

While SBC catches computational bugs and prior-likelihood conflicts,
LOO-PIT catches structural misspecification in the predictions
themselves.

\subsection{The Discrete Data Problem: Randomization for Bernoulli
Outcomes}\label{the-discrete-data-problem-randomization-for-bernoulli-outcomes}

Our outcome is a binary choice: \(h_{it} \in \{0, 1\}\). The PIT as
defined above requires a continuous predictive distribution, but
Bernoulli predictions are discrete. The CDF of a Bernoulli distribution
can only jump between two values: the predicted probability of \(h = 0\)
is \(1 - \hat{p}_i\), and the predicted probability of \(h = 1\) is
\(1.0\). So the raw PIT value is either \(1 - \hat{p}_i\) (when
\(h_i = 0\)) or \(1.0\) (when \(h_i = 1\)). These two families of values
cannot be uniformly distributed on \([0, 1]\), even for a perfectly
correct model.

The standard solution (Czado et al., 2009) is \textbf{randomization}.
Instead of using the CDF value directly, we spread each observation's
PIT value uniformly across the interval the CDF jumps over at that
observation. We draw \(v_i \sim \text{Uniform}(0,1)\) and compute:

\[u_i = F(y_i - 1 \mid y_{-i}) + v_i \cdot \left[F(y_i \mid y_{-i}) - F(y_i - 1 \mid y_{-i})\right]\]

For \(h_i = 1\): \(u_i = (1 - \hat{p}_i) + v_i \cdot \hat{p}_i\), which
spreads uniformly across \([1-\hat{p}_i, 1]\). For \(h_i = 0\):
\(u_i = v_i \cdot (1-\hat{p}_i)\), which spreads uniformly across
\([0, 1-\hat{p}_i]\).

Under a calibrated model, these randomized \(u_i\) values end up
marginally uniform on \([0,1]\). The \texttt{bayesplot} package applies
this randomization automatically when it detects binary data. But
understanding the reasoning matters: without randomization, a PIT plot
for Bernoulli data would look non-uniform even for a perfectly correct
model, which would be deeply misleading.

\subsection{A Critical Warning: The Standard Visual May Miss
Miscalibration}\label{a-critical-warning-the-standard-visual-may-miss-miscalibration}

Here is a subtlety that matters specifically for hierarchical cognitive
models, established recently by Tesso \& Vehtari (2026).

The standard LOO-PIT visualization - the ECDF (empirical cumulative
distribution function) envelope from Sailynoja et al.~(2022) -
constructs its confidence band assuming the \(u_1, u_2, \ldots, u_N\)
values are \emph{independent}. In a simple model with few parameters,
this is approximately true. But in a hierarchical model fitted to
hundreds or thousands of observations, the LOO-PIT values are
\textbf{correlated with one another}.

Why? Each LOO predictive distribution \(p(y_i \mid y_{-i})\) is
determined by the model fitted to the other \(N-1\) observations. Those
\(N-1\) datasets differ from each other by only one data point --- they
are nearly identical. All the LOO predictive distributions therefore
share nearly the same information, and the resulting \(u_i\) values
inherit this shared structure.

Tesso \& Vehtari (2026) prove that this correlation scales with model
complexity: the more effective parameters a model has relative to the
dataset size, the more correlated the LOO-PIT values are, and the
narrower the true null distribution of the ECDF becomes. When we use the
independence-based confidence envelope, that envelope is \emph{too wide}
- it accommodates more deviation than a properly-calibrated null
distribution would allow. The practical consequence is that the standard
visual has \textbf{reduced statistical power}: it fails to flag
miscalibration that is genuinely present. For simple models this is not
a serious problem. For our hierarchical models it can be.

\subsection{\texorpdfstring{Step 1: The \(\hat{p}_{\text{loo}} / N\)
Complexity
Ratio}{Step 1: The \textbackslash hat\{p\}\_\{\textbackslash text\{loo\}\} / N Complexity Ratio}}\label{step-1-the-hatp_textloo-n-complexity-ratio}

Before interpreting any LOO-PIT result, we compute one diagnostic
number: the ratio of the effective number of parameters
(\(\hat{p}_{\text{loo}}\)) to the total number of observations (\(N\)).
Tesso \& Vehtari (2026) show that:

\begin{itemize}
\tightlist
\item
  When \(\hat{p}_{\text{loo}} / N\) is close to zero, the LOO-PIT values
  are nearly independent and the standard visual is reliable.
\item
  When the ratio is non-negligible --- say, above 0.05 --- the
  independence assumption inflates the envelope width and a formal test
  is needed to supplement the visual.
\end{itemize}

Conveniently, \(\hat{p}_{\text{loo}}\) is already reported in our
\texttt{loo()} output. For a hierarchical model with \(J = 50\) agents,
\(\hat{p}_{\text{loo}}\) is roughly proportional to \(J\) (each agent
contributes approximately one effective parameter for their individual
parameters), giving a ratio around \(50 / 6000 \approx 0.008\) in our
case. This is small --- but not zero, and the exact value matters when
the dataset is smaller or the model is more complex.

\begin{Shaded}
\begin{Highlighting}[]
\NormalTok{complexity\_table }\OtherTok{\textless{}{-}}\NormalTok{ pareto\_k\_table }\SpecialCharTok{|\textgreater{}}
\NormalTok{  dplyr}\SpecialCharTok{::}\FunctionTok{select}\NormalTok{(Combination, N, p\_hat\_loo, ratio) }\SpecialCharTok{|\textgreater{}}
  \FunctionTok{mutate}\NormalTok{(}\AttributeTok{visual\_ok =} \FunctionTok{ifelse}\NormalTok{(ratio }\SpecialCharTok{\textless{}} \FloatTok{0.05}\NormalTok{,}
    \StringTok{"Standard visual acceptable"}\NormalTok{,}
    \StringTok{"Formal test recommended (POT{-}C)"}\NormalTok{))}

\NormalTok{knitr}\SpecialCharTok{::}\FunctionTok{kable}\NormalTok{(complexity\_table, }\AttributeTok{caption =} \FunctionTok{paste}\NormalTok{(}
  \StringTok{"Model complexity ratios. For hierarchical models with J=50 agents,"}\NormalTok{,}
  \StringTok{"p\_hat\_loo grows with J: each agent contributes roughly one effective parameter."}\NormalTok{,}
  \StringTok{"The ratio p\_hat\_loo/N determines whether the standard visual is sufficient."}\NormalTok{))}
\end{Highlighting}
\end{Shaded}

\begin{longtable}[]{@{}lrrrl@{}}
\caption{Model complexity ratios. For hierarchical models with J=50
agents, p\_hat\_loo grows with J: each agent contributes roughly one
effective parameter. The ratio p\_hat\_loo/N determines whether the
standard visual is sufficient.}\tabularnewline
\toprule\noalign{}
Combination & N & p\_hat\_loo & ratio & visual\_ok \\
\midrule\noalign{}
\endfirsthead
\toprule\noalign{}
Combination & N & p\_hat\_loo & ratio & visual\_ok \\
\midrule\noalign{}
\endhead
\bottomrule\noalign{}
\endlastfoot
Biased-\textgreater Biased & 6000 & 41.1 & 0.007 & Standard visual
acceptable \\
Memory-\textgreater Biased & 6000 & 47.0 & 0.008 & Standard visual
acceptable \\
Biased-\textgreater Memory & 6000 & 35.6 & 0.006 & Standard visual
acceptable \\
Memory-\textgreater Memory & 6000 & 61.7 & 0.010 & Standard visual
acceptable \\
\end{longtable}

Examine the \texttt{visual\_ok} column. For most of our fits the ratio
is small enough that the standard visual is acceptable --- but the
formal test below provides additional confidence, and is essential
whenever the ratio is larger.

\subsection{Step 2: The Randomized ECDF Visual
Check}\label{step-2-the-randomized-ecdf-visual-check}

We always produce the visual check, regardless of whether the ratio is
small. The visual is indispensable for understanding the \emph{shape} of
any miscalibration: is the model overconfident? Biased in one direction?

Because our Matching Pennies data is binary (\(y \in \{0, 1\}\)), the
true predictive CDF is a discrete step function. Applying standard
continuous Density Overlays to this generates severe mathematical
artifacts. To evaluate calibration accurately, we rely on the
\textbf{Randomized Probability Integral Transform}. We mathematically
restore continuity by injecting uniform noise strictly between the
discrete CDF steps. We then plot the Empirical Cumulative Distribution
Function (ECDF) of these randomized values. If the model is
well-calibrated, the ECDF will closely track the diagonal and remain
completely inside the shaded confidence envelope.

\begin{Shaded}
\begin{Highlighting}[]
\FunctionTok{color\_scheme\_set}\NormalTok{(}\StringTok{"blue"}\NormalTok{)}

\CommentTok{\# 1. Define a rigorous function to compute Randomized PIT for Bernoulli data}
\NormalTok{compute\_randomized\_pit }\OtherTok{\textless{}{-}} \ControlFlowTok{function}\NormalTok{(y, loo\_obj, }\AttributeTok{seed =} \DecValTok{101}\NormalTok{) \{}
  \FunctionTok{set.seed}\NormalTok{(seed)}
\NormalTok{  p\_obs }\OtherTok{\textless{}{-}} \FunctionTok{exp}\NormalTok{(loo\_obj}\SpecialCharTok{$}\NormalTok{pointwise[, }\StringTok{"elpd\_loo"}\NormalTok{])}
  \CommentTok{\# Inject uniform noise between the discrete CDF steps}
\NormalTok{  pit\_vals }\OtherTok{\textless{}{-}} \FunctionTok{ifelse}\NormalTok{(y }\SpecialCharTok{==} \DecValTok{1}\NormalTok{, }\FunctionTok{runif}\NormalTok{(}\FunctionTok{length}\NormalTok{(y), }\DecValTok{1} \SpecialCharTok{{-}}\NormalTok{ p\_obs, }\DecValTok{1}\NormalTok{), }\FunctionTok{runif}\NormalTok{(}\FunctionTok{length}\NormalTok{(y), }\DecValTok{0}\NormalTok{, p\_obs))}
  \FunctionTok{return}\NormalTok{(pit\_vals)}
\NormalTok{\}}

\CommentTok{\# 2. Extract outcomes and compute randomized PIT values }
\NormalTok{pit\_b\_on\_b }\OtherTok{\textless{}{-}} \FunctionTok{compute\_randomized\_pit}\NormalTok{(}\FunctionTok{as.integer}\NormalTok{(sd\_biased\_on\_biased}\SpecialCharTok{$}\NormalTok{h), loo\_b\_on\_b)}
\NormalTok{pit\_b\_on\_m }\OtherTok{\textless{}{-}} \FunctionTok{compute\_randomized\_pit}\NormalTok{(}\FunctionTok{as.integer}\NormalTok{(sd\_biased\_on\_memory}\SpecialCharTok{$}\NormalTok{h), loo\_b\_on\_m)}

\CommentTok{\# 3. Helper function to generate the ECDF plots}
\NormalTok{plot\_pit\_ecdf }\OtherTok{\textless{}{-}} \ControlFlowTok{function}\NormalTok{(y, pit\_vals, plot\_title) \{}
  \FunctionTok{ppc\_pit\_ecdf}\NormalTok{(}\AttributeTok{y =}\NormalTok{ y, }\AttributeTok{pit =}\NormalTok{ pit\_vals, }\AttributeTok{pit\_bounds =} \FunctionTok{c}\NormalTok{(}\DecValTok{0}\NormalTok{, }\DecValTok{1}\NormalTok{)) }\SpecialCharTok{+}
    \FunctionTok{labs}\NormalTok{(}\AttributeTok{title =}\NormalTok{ plot\_title) }\SpecialCharTok{+} \FunctionTok{theme\_classic}\NormalTok{()}
\NormalTok{\}}

\CommentTok{\# 4. Generate plots for the true model and the miscalibrated model}
\NormalTok{p\_pit1 }\OtherTok{\textless{}{-}} \FunctionTok{plot\_pit\_ecdf}\NormalTok{(}\FunctionTok{as.integer}\NormalTok{(sd\_biased\_on\_biased}\SpecialCharTok{$}\NormalTok{h), pit\_b\_on\_b, }\StringTok{"Biased {-}\textgreater{} Biased (Calibrated)"}\NormalTok{)}
\NormalTok{p\_pit2 }\OtherTok{\textless{}{-}} \FunctionTok{plot\_pit\_ecdf}\NormalTok{(}\FunctionTok{as.integer}\NormalTok{(sd\_biased\_on\_memory}\SpecialCharTok{$}\NormalTok{h), pit\_b\_on\_m, }\StringTok{"Biased {-}\textgreater{} Memory (Miscalibrated)"}\NormalTok{)}

\NormalTok{p\_pit1 }\SpecialCharTok{|}\NormalTok{ p\_pit2 }\SpecialCharTok{+}
  \FunctionTok{plot\_annotation}\NormalTok{(}
    \AttributeTok{title =} \StringTok{"Randomized LOO{-}PIT ECDF Calibration Checks"}\NormalTok{,}
    \AttributeTok{subtitle =} \StringTok{"Lines inside the shaded envelope indicate uniformity (well{-}calibrated model)."}\NormalTok{)}
\end{Highlighting}
\end{Shaded}

\begin{center}
\includegraphics[width=0.8\linewidth,height=\textheight,keepaspectratio]{07-ModelComparison_files/figure-pdf/07-loo-pit-visual-1.pdf}
\end{center}

The flat, nearly-uniform overlay for the well-specified fits
(Biased-\textgreater Biased and Memory-\textgreater Memory) is a
reassuring sign. The misspecified fits could have shown characteristic
departures, but in our case it doesn't. So we can move forward.

\subsection{Step 3: The POT-C Formal Test -- Valid Under
Dependence}\label{step-3-the-pot-c-formal-test-valid-under-dependence}

The visual check is useful for identifying the \emph{shape} of
miscalibration, but it can miss subtle violations when LOO-PIT values
are correlated. The POT-C (Pointwise Order Test Combination) from Tesso
\& Vehtari (2026) provides a formal uniformity test that remains
statistically valid regardless of the dependence structure among the PIT
values.

\textbf{The intuition behind POT-C in three steps:}

\emph{Step 1 - Individual tests per order statistic.} Imagine sorting
all \(N\) PIT values from smallest to largest:
\(u_{(1)} \leq u_{(2)} \leq \ldots \leq u_{(N)}\). By the theory of
order statistics, under uniformity
\(u_{(i)} \sim \text{Beta}(i, N+1-i)\). This gives us a valid marginal
test for each order statistic: we compute the two-sided p-value for
whether the observed \(u_{(i)}\) is consistent with
\(\text{Beta}(i, N+1-i)\). We now have \(N\) individual p-values, one
for each position in the sorted sample, each marginally valid but
potentially correlated. p-values, you read correctly. We are not doing
null hypothesis testing, tho, we are trying to see whether we should
reject our actual model.

\emph{Step 2 - Combining correlated p-values via the Cauchy method.}
Standard methods for combining p-values, like Fisher's method, assume
independence and would give misleading results here. The Cauchy
combination method (Liu \& Xie, 2020) handles this by transforming each
p-value as \(t^{(i)} = \tan\left[(0.5 - p^{(i)}) \pi\right]\). A p-value
of exactly 0.5 (no evidence against uniformity) contributes zero. A very
small p-value contributes a large positive value. The test statistic is
the average: \[T = \frac{1}{N} \sum_{i=1}^{N} t^{(i)}\] The key
mathematical result (Liu \& Xie, 2020) is that the tail probabilities of
\(T\) approximate those of a standard Cauchy distribution even under
arbitrary correlation structures. The heavy tails of the Cauchy
distribution ``absorb'' the dependence without requiring us to model it
explicitly.

\emph{Step 3 - Shapley values for localization.} Beyond asking ``is the
model miscalibrated?'', we want to know ``where is it miscalibrated?''.
The Shapley value \(\phi_i\) for each observation measures how much that
observation's order statistic contributed to the overall test statistic
\(T\): how much it pushed the result toward rejection. By coloring the
ECDF difference plot with these values, we can visually highlight which
regions of the predictive distribution are problematic. This
localization connects the statistical diagnostic back to the cognitive
model.

\begin{Shaded}
\begin{Highlighting}[]
\CommentTok{\# {-}{-} POT{-}C implementation (Tesso \& Vehtari, 2026) {-}{-}{-}{-}{-}{-}{-}{-}{-}{-}{-}{-}{-}{-}{-}{-}{-}{-}{-}{-}{-}{-}{-}{-}{-}{-}{-}{-}{-}}

\NormalTok{pot\_c\_test }\OtherTok{\textless{}{-}} \ControlFlowTok{function}\NormalTok{(u\_vec, }\AttributeTok{alpha =} \FloatTok{0.05}\NormalTok{) \{}
\NormalTok{  u   }\OtherTok{\textless{}{-}} \FunctionTok{sort}\NormalTok{(u\_vec)}
\NormalTok{  n   }\OtherTok{\textless{}{-}} \FunctionTok{length}\NormalTok{(u)}

  \CommentTok{\# Step 1: Beta{-}based p{-}value for each order statistic}
\NormalTok{  p\_vals }\OtherTok{\textless{}{-}} \FunctionTok{map\_dbl}\NormalTok{(}\FunctionTok{seq\_len}\NormalTok{(n), }\ControlFlowTok{function}\NormalTok{(i) \{}
\NormalTok{    cdf\_val }\OtherTok{\textless{}{-}} \FunctionTok{pbeta}\NormalTok{(u[i], }\AttributeTok{shape1 =}\NormalTok{ i, }\AttributeTok{shape2 =}\NormalTok{ n }\SpecialCharTok{+} \DecValTok{1} \SpecialCharTok{{-}}\NormalTok{ i)}
    \DecValTok{2} \SpecialCharTok{*} \FunctionTok{min}\NormalTok{(cdf\_val, }\DecValTok{1} \SpecialCharTok{{-}}\NormalTok{ cdf\_val)}
\NormalTok{  \})}

  \CommentTok{\# Step 2: Cauchy transformation and global p{-}value}
\NormalTok{  t\_vals }\OtherTok{\textless{}{-}} \FunctionTok{tan}\NormalTok{((}\FloatTok{0.5} \SpecialCharTok{{-}}\NormalTok{ p\_vals) }\SpecialCharTok{*}\NormalTok{ pi)}
\NormalTok{  T\_stat }\OtherTok{\textless{}{-}} \FunctionTok{mean}\NormalTok{(t\_vals)}
\NormalTok{  p\_global }\OtherTok{\textless{}{-}} \DecValTok{1} \SpecialCharTok{{-}} \FunctionTok{pcauchy}\NormalTok{(T\_stat)}

  \CommentTok{\# Step 3: Shapley values {-}{-} each observation\textquotesingle{}s marginal contribution to T}
  \CommentTok{\# phi\_i = (1/n)*t\_i + (H\_n{-}1)/n * (t\_i {-} mean of other t\_j)}
\NormalTok{  H\_n      }\OtherTok{\textless{}{-}} \FunctionTok{sum}\NormalTok{(}\DecValTok{1} \SpecialCharTok{/} \FunctionTok{seq\_len}\NormalTok{(n))}
\NormalTok{  t\_sum    }\OtherTok{\textless{}{-}} \FunctionTok{sum}\NormalTok{(t\_vals)}
\NormalTok{  phi\_vals }\OtherTok{\textless{}{-}} \FunctionTok{map\_dbl}\NormalTok{(}\FunctionTok{seq\_len}\NormalTok{(n), }\ControlFlowTok{function}\NormalTok{(i) \{}
\NormalTok{    t\_i }\OtherTok{\textless{}{-}}\NormalTok{ t\_vals[i]}
\NormalTok{    (}\DecValTok{1}\SpecialCharTok{/}\NormalTok{n) }\SpecialCharTok{*}\NormalTok{ t\_i }\SpecialCharTok{+}\NormalTok{ ((H\_n }\SpecialCharTok{{-}} \DecValTok{1}\NormalTok{)}\SpecialCharTok{/}\NormalTok{n) }\SpecialCharTok{*}\NormalTok{ (t\_i }\SpecialCharTok{{-}}\NormalTok{ (t\_sum }\SpecialCharTok{{-}}\NormalTok{ t\_i) }\SpecialCharTok{/}\NormalTok{ (n }\SpecialCharTok{{-}} \DecValTok{1}\NormalTok{))}
\NormalTok{  \})}

  \FunctionTok{list}\NormalTok{(}\AttributeTok{p\_global=}\NormalTok{p\_global, }\AttributeTok{T\_stat=}\NormalTok{T\_stat, }\AttributeTok{p\_vals=}\NormalTok{p\_vals, }\AttributeTok{t\_vals=}\NormalTok{t\_vals,}
       \AttributeTok{phi=}\NormalTok{phi\_vals, }\AttributeTok{u\_sorted=}\NormalTok{u, }\AttributeTok{n=}\NormalTok{n)}
\NormalTok{\}}

\CommentTok{\# {-}{-} Randomized LOO{-}PIT for Bernoulli data via PSIS weights {-}{-}{-}{-}{-}{-}{-}{-}{-}{-}{-}{-}{-}{-}{-}{-}{-}{-}{-}}
\CommentTok{\# We approximate the LOO predictive P(h\_i=1|y\_\{{-}i\}) as a weighted mean}
\CommentTok{\# of the posterior predictive draws, using the PSIS log{-}weights.}
\CommentTok{\# The resulting u\_i values are then randomized per Czado et al. (2009).}

\NormalTok{compute\_loo\_pit\_bernoulli }\OtherTok{\textless{}{-}} \ControlFlowTok{function}\NormalTok{(y, yrep, psis\_lw, }\AttributeTok{seed =} \DecValTok{99}\NormalTok{) \{}
  \FunctionTok{set.seed}\NormalTok{(seed)}
\NormalTok{  w\_mat   }\OtherTok{\textless{}{-}} \FunctionTok{exp}\NormalTok{(psis\_lw)                           }\CommentTok{\# S x N probability weights}
\NormalTok{  p\_loo\_1 }\OtherTok{\textless{}{-}} \FunctionTok{colSums}\NormalTok{(w\_mat }\SpecialCharTok{*}\NormalTok{ yrep) }\SpecialCharTok{/} \FunctionTok{colSums}\NormalTok{(w\_mat) }\CommentTok{\# weighted mean per observation}
\NormalTok{  p\_loo\_1 }\OtherTok{\textless{}{-}} \FunctionTok{pmax}\NormalTok{(}\FloatTok{1e{-}6}\NormalTok{, }\FunctionTok{pmin}\NormalTok{(}\DecValTok{1} \SpecialCharTok{{-}} \FloatTok{1e{-}6}\NormalTok{, p\_loo\_1))}

  \CommentTok{\# Randomized PIT:}
  \CommentTok{\# h=1: u\_i \textasciitilde{} Uniform(1{-}p\_i, 1),   i.e. (1{-}p) + v*p}
  \CommentTok{\# h=0: u\_i \textasciitilde{} Uniform(0, 1{-}p\_i),   i.e. v*(1{-}p)}
\NormalTok{  V }\OtherTok{\textless{}{-}} \FunctionTok{runif}\NormalTok{(}\FunctionTok{length}\NormalTok{(y))}
  \FunctionTok{ifelse}\NormalTok{(y }\SpecialCharTok{==} \DecValTok{1}\NormalTok{,}
\NormalTok{         (}\DecValTok{1} \SpecialCharTok{{-}}\NormalTok{ p\_loo\_1) }\SpecialCharTok{+}\NormalTok{ V }\SpecialCharTok{*}\NormalTok{ p\_loo\_1,}
\NormalTok{         V }\SpecialCharTok{*}\NormalTok{ (}\DecValTok{1} \SpecialCharTok{{-}}\NormalTok{ p\_loo\_1))}
\NormalTok{\}}

\NormalTok{run\_pit\_analysis }\OtherTok{\textless{}{-}} \ControlFlowTok{function}\NormalTok{(fit\_obj, data\_list, loo\_obj, label) \{}
\NormalTok{  y      }\OtherTok{\textless{}{-}} \FunctionTok{as.integer}\NormalTok{(data\_list}\SpecialCharTok{$}\NormalTok{h)}
\NormalTok{  yrep   }\OtherTok{\textless{}{-}}\NormalTok{ fit\_obj}\SpecialCharTok{$}\FunctionTok{draws}\NormalTok{(}\StringTok{"h\_post\_rep"}\NormalTok{, }\AttributeTok{format =} \StringTok{"matrix"}\NormalTok{)}
\NormalTok{  lw     }\OtherTok{\textless{}{-}} \FunctionTok{weights}\NormalTok{(loo\_obj}\SpecialCharTok{$}\NormalTok{psis\_object)}
\NormalTok{  u\_vec  }\OtherTok{\textless{}{-}} \FunctionTok{compute\_loo\_pit\_bernoulli}\NormalTok{(y, yrep, lw)}
\NormalTok{  result }\OtherTok{\textless{}{-}} \FunctionTok{pot\_c\_test}\NormalTok{(u\_vec)}
  \FunctionTok{list}\NormalTok{(}\AttributeTok{label =}\NormalTok{ label, }\AttributeTok{u =}\NormalTok{ u\_vec, }\AttributeTok{result =}\NormalTok{ result)}
\NormalTok{\}}

\NormalTok{pit\_b\_on\_b }\OtherTok{\textless{}{-}} \FunctionTok{run\_pit\_analysis}\NormalTok{(fit\_b\_on\_b, sd\_biased\_on\_biased, loo\_b\_on\_b, }\StringTok{"Biased{-}\textgreater{}Biased"}\NormalTok{)}
\NormalTok{pit\_m\_on\_b }\OtherTok{\textless{}{-}} \FunctionTok{run\_pit\_analysis}\NormalTok{(fit\_m\_on\_b, sd\_memory\_on\_biased, loo\_m\_on\_b, }\StringTok{"Memory{-}\textgreater{}Biased"}\NormalTok{)}
\NormalTok{pit\_b\_on\_m }\OtherTok{\textless{}{-}} \FunctionTok{run\_pit\_analysis}\NormalTok{(fit\_b\_on\_m, sd\_biased\_on\_memory, loo\_b\_on\_m, }\StringTok{"Biased{-}\textgreater{}Memory"}\NormalTok{)}
\NormalTok{pit\_m\_on\_m }\OtherTok{\textless{}{-}} \FunctionTok{run\_pit\_analysis}\NormalTok{(fit\_m\_on\_m, sd\_memory\_on\_memory, loo\_m\_on\_m, }\StringTok{"Memory{-}\textgreater{}Memory"}\NormalTok{)}

\FunctionTok{map\_dfr}\NormalTok{(}\FunctionTok{list}\NormalTok{(pit\_b\_on\_b, pit\_m\_on\_b, pit\_b\_on\_m, pit\_m\_on\_m), }\ControlFlowTok{function}\NormalTok{(x) \{}
  \FunctionTok{tibble}\NormalTok{(}\AttributeTok{Fit =}\NormalTok{ x}\SpecialCharTok{$}\NormalTok{label, }\AttributeTok{p\_global =} \FunctionTok{round}\NormalTok{(x}\SpecialCharTok{$}\NormalTok{result}\SpecialCharTok{$}\NormalTok{p\_global,}\DecValTok{4}\NormalTok{),}
         \AttributeTok{T\_stat =} \FunctionTok{round}\NormalTok{(x}\SpecialCharTok{$}\NormalTok{result}\SpecialCharTok{$}\NormalTok{T\_stat,}\DecValTok{2}\NormalTok{),}
         \AttributeTok{Verdict =} \FunctionTok{ifelse}\NormalTok{(x}\SpecialCharTok{$}\NormalTok{result}\SpecialCharTok{$}\NormalTok{p\_global }\SpecialCharTok{\textless{}} \FloatTok{0.05}\NormalTok{,}
                        \StringTok{"REJECT {-}{-} miscalibrated"}\NormalTok{, }\StringTok{"Fail to reject {-}{-} acceptable"}\NormalTok{))}
\NormalTok{\}) }\SpecialCharTok{|\textgreater{}}
\NormalTok{  knitr}\SpecialCharTok{::}\FunctionTok{kable}\NormalTok{(}\AttributeTok{caption =} \FunctionTok{paste}\NormalTok{(}
    \StringTok{"POT{-}C formal uniformity test for LOO{-}PIT values."}\NormalTok{,}
    \StringTok{"Valid even when LOO{-}PIT values are correlated."}\NormalTok{,}
    \StringTok{"p \textless{} 0.05 = significant departure from calibration."}\NormalTok{))}
\end{Highlighting}
\end{Shaded}

\begin{longtable}[]{@{}lrrl@{}}
\caption{POT-C formal uniformity test for LOO-PIT values. Valid even
when LOO-PIT values are correlated. p \textless{} 0.05 = significant
departure from calibration.}\tabularnewline
\toprule\noalign{}
Fit & p\_global & T\_stat & Verdict \\
\midrule\noalign{}
\endfirsthead
\toprule\noalign{}
Fit & p\_global & T\_stat & Verdict \\
\midrule\noalign{}
\endhead
\bottomrule\noalign{}
\endlastfoot
Biased-\textgreater Biased & 0 & 2.096002e+09 & REJECT --
miscalibrated \\
Memory-\textgreater Biased & 0 & 2.242040e+09 & REJECT --
miscalibrated \\
Biased-\textgreater Memory & 0 & 1.243144e+16 & REJECT --
miscalibrated \\
Memory-\textgreater Memory & 0 & 1.206140e+16 & REJECT --
miscalibrated \\
\end{longtable}

\begin{Shaded}
\begin{Highlighting}[]
\CommentTok{\# {-}{-} Shapley{-}colored ECDF difference plots {-}{-}{-}{-}{-}{-}{-}{-}{-}{-}{-}{-}{-}{-}{-}{-}{-}{-}{-}{-}{-}{-}{-}{-}{-}{-}{-}{-}{-}{-}{-}{-}{-}{-}{-}{-}}
\CommentTok{\# The ECDF difference plot shows F\_hat(u) {-} u (deviation from perfect uniform).}
\CommentTok{\# Perfect calibration = horizontal line at y=0.}
\CommentTok{\# We color each point by its Shapley value:}
\CommentTok{\#   RED  = positive phi (pushes the test statistic toward rejection)}
\CommentTok{\#   BLUE = neutral or negative phi (consistent with uniformity)}

\NormalTok{make\_shapley\_ecdf }\OtherTok{\textless{}{-}} \ControlFlowTok{function}\NormalTok{(pit\_res) \{}
\NormalTok{  r   }\OtherTok{\textless{}{-}}\NormalTok{ pit\_res}\SpecialCharTok{$}\NormalTok{result}
\NormalTok{  n   }\OtherTok{\textless{}{-}}\NormalTok{ r}\SpecialCharTok{$}\NormalTok{n}
  
  \CommentTok{\# Corrected namespace assignment}
\NormalTok{  df  }\OtherTok{\textless{}{-}} \FunctionTok{tibble}\NormalTok{(}
    \AttributeTok{u\_sorted    =}\NormalTok{ r}\SpecialCharTok{$}\NormalTok{u\_sorted, }
    \AttributeTok{phi         =}\NormalTok{ r}\SpecialCharTok{$}\NormalTok{phi,}
    \AttributeTok{ecdf\_y      =} \FunctionTok{seq\_len}\NormalTok{(n) }\SpecialCharTok{/}\NormalTok{ n,}
    \AttributeTok{ecdf\_diff   =}\NormalTok{ ecdf\_y }\SpecialCharTok{{-}}\NormalTok{ u\_sorted,}
    \AttributeTok{influential =}\NormalTok{ phi }\SpecialCharTok{\textgreater{}} \DecValTok{0}
\NormalTok{  )}
  
\NormalTok{  p\_label }\OtherTok{\textless{}{-}} \FunctionTok{ifelse}\NormalTok{(r}\SpecialCharTok{$}\NormalTok{p\_global }\SpecialCharTok{\textless{}} \FloatTok{0.001}\NormalTok{, }\StringTok{"p \textless{} 0.001"}\NormalTok{,}
                    \FunctionTok{sprintf}\NormalTok{(}\StringTok{"p = \%.3f"}\NormalTok{, r}\SpecialCharTok{$}\NormalTok{p\_global))}
  
  \FunctionTok{ggplot}\NormalTok{(df, }\FunctionTok{aes}\NormalTok{(}\AttributeTok{x =}\NormalTok{ u\_sorted, }\AttributeTok{y =}\NormalTok{ ecdf\_diff)) }\SpecialCharTok{+}
    \FunctionTok{geom\_hline}\NormalTok{(}\AttributeTok{yintercept =} \DecValTok{0}\NormalTok{, }\AttributeTok{color =} \StringTok{"gray40"}\NormalTok{, }\AttributeTok{linewidth =} \FloatTok{0.5}\NormalTok{) }\SpecialCharTok{+}
    \FunctionTok{geom\_step}\NormalTok{(}\AttributeTok{color =} \StringTok{"gray30"}\NormalTok{, }\AttributeTok{linewidth =} \FloatTok{0.7}\NormalTok{) }\SpecialCharTok{+}
    \FunctionTok{geom\_point}\NormalTok{(}\FunctionTok{aes}\NormalTok{(}\AttributeTok{color =}\NormalTok{ influential), }\AttributeTok{size =} \FloatTok{1.0}\NormalTok{, }\AttributeTok{alpha =} \FloatTok{0.7}\NormalTok{) }\SpecialCharTok{+}
    \FunctionTok{scale\_color\_manual}\NormalTok{(}
      \AttributeTok{values =} \FunctionTok{c}\NormalTok{(}\StringTok{"FALSE"}\OtherTok{=}\StringTok{"steelblue"}\NormalTok{, }\StringTok{"TRUE"}\OtherTok{=}\StringTok{"firebrick"}\NormalTok{),}
      \AttributeTok{labels =} \FunctionTok{c}\NormalTok{(}\StringTok{"FALSE"}\OtherTok{=}\StringTok{"Neutral"}\NormalTok{, }\StringTok{"TRUE"}\OtherTok{=}\StringTok{"Drives rejection"}\NormalTok{),}
      \AttributeTok{name   =} \StringTok{"Shapley influence"}\NormalTok{) }\SpecialCharTok{+}
    \FunctionTok{annotate}\NormalTok{(}\StringTok{"text"}\NormalTok{, }\AttributeTok{x=}\FloatTok{0.02}\NormalTok{, }\AttributeTok{y=}\ConstantTok{Inf}\NormalTok{, }\AttributeTok{hjust=}\DecValTok{0}\NormalTok{, }\AttributeTok{vjust=}\FloatTok{1.5}\NormalTok{,}
             \AttributeTok{label=}\NormalTok{p\_label, }\AttributeTok{fontface=}\StringTok{"bold"}\NormalTok{, }\AttributeTok{size=}\FloatTok{3.5}\NormalTok{) }\SpecialCharTok{+}
    \FunctionTok{scale\_x\_continuous}\NormalTok{(}\AttributeTok{limits=}\FunctionTok{c}\NormalTok{(}\DecValTok{0}\NormalTok{,}\DecValTok{1}\NormalTok{)) }\SpecialCharTok{+}
    \FunctionTok{labs}\NormalTok{(}\AttributeTok{title    =}\NormalTok{ pit\_res}\SpecialCharTok{$}\NormalTok{label,}
         \AttributeTok{subtitle =} \StringTok{"ECDF(u) {-} u: zero everywhere = perfect calibration"}\NormalTok{,}
         \AttributeTok{x =} \StringTok{"LOO{-}PIT value u"}\NormalTok{, }\AttributeTok{y =} \StringTok{"ECDF(u) {-} u"}\NormalTok{) }\SpecialCharTok{+}
    \FunctionTok{theme}\NormalTok{(}\AttributeTok{legend.position =} \StringTok{"bottom"}\NormalTok{)}
\NormalTok{\}}

\CommentTok{\# Plot only the correctly specified and the explicitly miscalibrated model}
\FunctionTok{make\_shapley\_ecdf}\NormalTok{(pit\_b\_on\_b) }\SpecialCharTok{|} \FunctionTok{make\_shapley\_ecdf}\NormalTok{(pit\_b\_on\_m) }\SpecialCharTok{+}
  \FunctionTok{plot\_annotation}\NormalTok{(}\AttributeTok{title =} \StringTok{"Shapley{-}Colored ECDF Difference Plots (POT{-}C)"}\NormalTok{)}
\end{Highlighting}
\end{Shaded}

\begin{center}
\includegraphics[width=0.8\linewidth,height=\textheight,keepaspectratio]{07-ModelComparison_files/figure-pdf/07-pot-c-implementation-1.pdf}
\end{center}

The ECDF difference plot is a centered version of the usual ECDF:
instead of plotting \(\hat{F}(u)\) against \(u\), we plot
\(\hat{F}(u) - u\). Perfect calibration corresponds to a horizontal line
at \(y = 0\). Deviations reveal the shape of the problem:

\begin{itemize}
\item
  Red in the left tail (near \(u = 0\)): The model was confidently wrong
  when the agent chose ``left'' (\(y=0\)). The model predicted a
  near-certain probability of ``right'', forcing the randomized PIT
  value into the extreme lower bound.
\item
  Red in the right tail (near \(u = 1\)): The model was confidently
  wrong when the agent chose ``right'' (\(y=1\)). The model predicted a
  near-certain probability of ``left'', forcing the randomized PIT value
  into the extreme upper bound.
\item
  A global slope or S-curve: The model exhibits a systematic bias in
  base rates. It consistently over- or under-predicts the overall global
  frequency of one of the choices, rather than just failing on extreme
  boundary cases.
\end{itemize}

The Shapley coloring tells us exactly which observations are
responsible. Red points are the ``smoking guns'' --- the specific order
statistics mathematically pushing the POT-C test toward rejection. For
the Biased \(\rightarrow\) Memory combination, expect the red points to
concentrate heavily in the extreme tails. These correspond primarily to
the trials immediately following the block reversal, where the biased
model is rigidly and repeatedly surprised by behavioral changes it
structurally lacks the flexibility to predict.

\subsection{Connecting LOO-PIT to the Cognitive
Theory}\label{connecting-loo-pit-to-the-cognitive-theory}

Consider what happens when we fit the biased model to memory data. The
biased model says: ``Each agent has a fixed probability \(\psi_j\) of
choosing right, regardless of what the opponent has been doing.'' When
memory agents play against the reversal design, their trial-level choice
probabilities are anything but fixed. In the first block (trials 1--60),
as the running opponent rate climbs toward 0.8, memory agents
increasingly track it and choose right more often. At the reversal
(trial 61), the opponent flips to 20\% right: the memory signal inverts,
and memory agents start choosing left. The biased model cannot represent
either the rise or the reversal --- it is committed to a single
probability \(\psi_j\) that tries to average across both phases.

The predictive failure is asymmetric. In the first block the biased
model may do tolerably: its fixed \(\psi_j\) roughly matches the
early-session average. But immediately after the reversal, memory agents
start choosing left while the biased model is still confidently
predicting right. On those trials the biased model assigns high
probability to ``right'' and the agent chooses ``left'' --- exactly the
kind of systematic surprise that drives LOO-PIT values into the tails
and produces a large Pareto-\(\hat{k}\). This is the overconfidence
signature --- not overconfidence in the sense of arrogance, but in the
sense of structural rigidity: the model makes narrow, committed
predictions that the reversal repeatedly proves wrong.

\subsection{Agent-Level Calibration: A Transparent
Complement}\label{agent-level-calibration-a-transparent-complement}

LOO-PIT operates at the trial level, checking whether the model's
uncertainty about individual choices is well-calibrated. A complementary
check operates at the agent level: does the model get the average
behavior of each agent right?

\begin{Shaded}
\begin{Highlighting}[]
\NormalTok{agent\_calibration }\OtherTok{\textless{}{-}} \ControlFlowTok{function}\NormalTok{(fit\_obj, data\_df, label) \{}
\NormalTok{  obs\_rates }\OtherTok{\textless{}{-}}\NormalTok{ data\_df }\SpecialCharTok{|\textgreater{}}
    \FunctionTok{group\_by}\NormalTok{(agent\_id) }\SpecialCharTok{|\textgreater{}}
    \FunctionTok{summarise}\NormalTok{(}\AttributeTok{obs\_rate =} \FunctionTok{mean}\NormalTok{(choice), }\AttributeTok{.groups=}\StringTok{"drop"}\NormalTok{)}
\NormalTok{  yrep      }\OtherTok{\textless{}{-}}\NormalTok{ fit\_obj}\SpecialCharTok{$}\FunctionTok{draws}\NormalTok{(}\StringTok{"h\_post\_rep"}\NormalTok{, }\AttributeTok{format=}\StringTok{"matrix"}\NormalTok{)}
\NormalTok{  agent\_vec }\OtherTok{\textless{}{-}}\NormalTok{ data\_df}\SpecialCharTok{$}\NormalTok{agent\_id}
\NormalTok{  pred\_mean }\OtherTok{\textless{}{-}} \FunctionTok{map\_dbl}\NormalTok{(}\FunctionTok{sort}\NormalTok{(}\FunctionTok{unique}\NormalTok{(agent\_vec)), }\ControlFlowTok{function}\NormalTok{(j)}
    \FunctionTok{mean}\NormalTok{(yrep[, }\FunctionTok{which}\NormalTok{(agent\_vec }\SpecialCharTok{==}\NormalTok{ j)]))}
\NormalTok{  obs\_rates }\SpecialCharTok{|\textgreater{}} \FunctionTok{mutate}\NormalTok{(}\AttributeTok{pred\_rate=}\NormalTok{pred\_mean, }\AttributeTok{label=}\NormalTok{label)}
\NormalTok{\}}

\FunctionTok{bind\_rows}\NormalTok{(}
  \FunctionTok{agent\_calibration}\NormalTok{(fit\_b\_on\_b, d\_biased, }\StringTok{"Biased{-}\textgreater{}Biased"}\NormalTok{),}
  \FunctionTok{agent\_calibration}\NormalTok{(fit\_m\_on\_b, d\_biased, }\StringTok{"Memory{-}\textgreater{}Biased"}\NormalTok{),}
  \FunctionTok{agent\_calibration}\NormalTok{(fit\_b\_on\_m, d\_memory, }\StringTok{"Biased{-}\textgreater{}Memory"}\NormalTok{),}
  \FunctionTok{agent\_calibration}\NormalTok{(fit\_m\_on\_m, d\_memory, }\StringTok{"Memory{-}\textgreater{}Memory"}\NormalTok{)}
\NormalTok{) }\SpecialCharTok{|\textgreater{}}
  \FunctionTok{ggplot}\NormalTok{(}\FunctionTok{aes}\NormalTok{(}\AttributeTok{x=}\NormalTok{obs\_rate, }\AttributeTok{y=}\NormalTok{pred\_rate)) }\SpecialCharTok{+}
  \FunctionTok{geom\_abline}\NormalTok{(}\AttributeTok{intercept=}\DecValTok{0}\NormalTok{, }\AttributeTok{slope=}\DecValTok{1}\NormalTok{, }\AttributeTok{linetype=}\StringTok{"dashed"}\NormalTok{, }\AttributeTok{color=}\StringTok{"gray50"}\NormalTok{) }\SpecialCharTok{+}
  \FunctionTok{geom\_point}\NormalTok{(}\AttributeTok{color=}\StringTok{"steelblue"}\NormalTok{, }\AttributeTok{alpha=}\FloatTok{0.7}\NormalTok{, }\AttributeTok{size=}\DecValTok{2}\NormalTok{) }\SpecialCharTok{+}
  \FunctionTok{facet\_wrap}\NormalTok{(}\SpecialCharTok{\textasciitilde{}}\NormalTok{label, }\AttributeTok{ncol=}\DecValTok{2}\NormalTok{) }\SpecialCharTok{+}
  \FunctionTok{scale\_x\_continuous}\NormalTok{(}\AttributeTok{labels=}\NormalTok{scales}\SpecialCharTok{::}\NormalTok{percent) }\SpecialCharTok{+}
  \FunctionTok{scale\_y\_continuous}\NormalTok{(}\AttributeTok{labels=}\NormalTok{scales}\SpecialCharTok{::}\NormalTok{percent) }\SpecialCharTok{+}
  \FunctionTok{labs}\NormalTok{(}\AttributeTok{title    =} \StringTok{"Agent{-}Level Calibration: Observed vs. Predicted Choice Rates"}\NormalTok{,}
       \AttributeTok{subtitle =} \StringTok{"LOO{-}PIT checks trial{-}level calibration; this checks agent{-}level means. Both matter."}\NormalTok{,}
       \AttributeTok{x=}\StringTok{"Observed proportion \textquotesingle{}right\textquotesingle{}"}\NormalTok{, }\AttributeTok{y=}\StringTok{"Posterior mean predicted proportion \textquotesingle{}right\textquotesingle{}"}\NormalTok{)}
\end{Highlighting}
\end{Shaded}

\begin{center}
\includegraphics[width=0.8\linewidth,height=\textheight,keepaspectratio]{07-ModelComparison_files/figure-pdf/07-agent-calibration-1.pdf}
\end{center}

This plot answers a different question from LOO-PIT: not ``is the
model's predictive \emph{uncertainty} correct?'' but ``is the model
getting the agent-level \emph{means} right?'' A model can pass this
check and still fail LOO-PIT --- for example, the biased model might
accurately predict each agent's average choice rate (because it directly
estimates their \(\theta_j\)), while still being badly miscalibrated at
the trial level (because it cannot predict which specific trials will
deviate from that average). Both checks are necessary for a complete
picture.

\section{Formal Model Comparison}\label{formal-model-comparison}

\begin{quote}
\emph{Why ELPD/LOO rather than Bayes Factors?} Bayes Factors compare the
marginal likelihood of the data under each model --- an elegant idea,
but one that requires integrating over the \emph{entire prior}. For
hierarchical models, this integral is computationally intractable and
the result is exquisitely sensitive to the exact width of every prior.
PSIS-LOO avoids this by focusing entirely on predictive accuracy for
held-out observations --- a criterion that is robust to prior choices
and provides the Pareto-\(\hat{k}\) as a built-in reliability
diagnostic. The trade-off is philosophical: LOO measures pragmatic
predictive utility, not the probability that a model is true. For most
applied cognitive modeling questions, pragmatic predictive utility is
what we actually care about.
\end{quote}

\subsection{\texorpdfstring{Reading \texttt{loo\_compare()}
Correctly}{Reading loo\_compare() Correctly}}\label{reading-loo_compare-correctly}

Now we come to the main event: using PSIS-LOO to compare the two models
on each dataset. The \texttt{loo\_compare()} function ranks models by
their ELPD estimates and reports the differences. Understanding exactly
what those differences mean is important before trusting any conclusion.

The output has two key columns:

\begin{itemize}
\item
  \texttt{elpd\_diff}: The difference in ELPD between this model and the
  best-ranked model (which always appears in row 1, with
  \texttt{elpd\_diff\ =\ 0} by definition). For lower-ranked models,
  \texttt{elpd\_diff} is negative --- they predict worse. The magnitude
  matters: a difference of \(-5\) nats means the model is roughly
  \(e^5 \approx 148\) times worse per observation on average. A
  difference of \(-0.5\) is much less decisive.
\item
  \texttt{se\_diff}: The standard error of that difference. This is
  estimated directly from the \emph{pointwise} ELPD contributions,
  treating each observation as an independent draw from the distribution
  of possible pointwise differences. It accounts for the fact that
  observations are not equally predictable.
\end{itemize}

The ratio \(|\text{elpd\_diff}|/\text{se\_diff}\) behaves like a
Z-score: it tells us how many standard errors the difference is from
zero. As a rough guide:

\begin{itemize}
\tightlist
\item
  Below 2: inconclusive --- the models are not decisively separable at
  this data size.
\item
  Between 2 and 4: suggestive --- there is a preference but it may not
  be reliable.
\item
  Above 4: a practically meaningful difference that is unlikely to be
  sampling noise.
\end{itemize}

What if the result is inconclusive (\(< 2\))? An inconclusive result
means the current data cannot distinguish between your theories. In this
case, do not simply pick the model with the slightly higher ELPD.
Instead:

\begin{itemize}
\item
  Redesign the task: Create a more discriminative opponent (like our
  reversal design) to force the models apart.
\item
  Collect more data: Use your Model Recovery power analysis to see how
  many more participants and/or trials you need.
\item
  Check PPCs: One model might have a better ELPD but fail to capture a
  specific qualitative feature of the behavior.
\end{itemize}

\begin{Shaded}
\begin{Highlighting}[]
\FunctionTok{cat}\NormalTok{(}\StringTok{"=== BIASED DATA (truth = Biased) ===}\SpecialCharTok{\textbackslash{}n\textbackslash{}n}\StringTok{"}\NormalTok{)}
\end{Highlighting}
\end{Shaded}

\begin{verbatim}
=== BIASED DATA (truth = Biased) ===
\end{verbatim}

\begin{Shaded}
\begin{Highlighting}[]
\NormalTok{compare\_b }\OtherTok{\textless{}{-}} \FunctionTok{loo\_compare}\NormalTok{(loo\_b\_on\_b, loo\_m\_on\_b)}
\FunctionTok{print}\NormalTok{(compare\_b)}
\end{Highlighting}
\end{Shaded}

\begin{verbatim}
       elpd_diff se_diff
model1  0.0       0.0   
model2 -1.3       1.6   
\end{verbatim}

\begin{Shaded}
\begin{Highlighting}[]
\NormalTok{z\_b }\OtherTok{\textless{}{-}} \FunctionTok{abs}\NormalTok{(compare\_b[}\DecValTok{2}\NormalTok{,}\StringTok{"elpd\_diff"}\NormalTok{]}\SpecialCharTok{/}\NormalTok{compare\_b[}\DecValTok{2}\NormalTok{,}\StringTok{"se\_diff"}\NormalTok{])}
\FunctionTok{cat}\NormalTok{(}\FunctionTok{sprintf}\NormalTok{(}\StringTok{"}\SpecialCharTok{\textbackslash{}n}\StringTok{|elpd\_diff|/SE = \%.1f  {-}\textgreater{}  \%s}\SpecialCharTok{\textbackslash{}n\textbackslash{}n}\StringTok{"}\NormalTok{, z\_b,}
  \FunctionTok{ifelse}\NormalTok{(z\_b}\SpecialCharTok{\textgreater{}}\DecValTok{4}\NormalTok{,}\StringTok{"Clear preference."}\NormalTok{,}\FunctionTok{ifelse}\NormalTok{(z\_b}\SpecialCharTok{\textgreater{}}\DecValTok{2}\NormalTok{,}\StringTok{"Weak preference."}\NormalTok{,}\StringTok{"Inconclusive."}\NormalTok{))))}
\end{Highlighting}
\end{Shaded}

\begin{verbatim}

|elpd_diff|/SE = 0.8  ->  Inconclusive.
\end{verbatim}

\begin{Shaded}
\begin{Highlighting}[]
\FunctionTok{cat}\NormalTok{(}\StringTok{"=== MEMORY DATA (truth = Memory) ===}\SpecialCharTok{\textbackslash{}n\textbackslash{}n}\StringTok{"}\NormalTok{)}
\end{Highlighting}
\end{Shaded}

\begin{verbatim}
=== MEMORY DATA (truth = Memory) ===
\end{verbatim}

\begin{Shaded}
\begin{Highlighting}[]
\NormalTok{compare\_m }\OtherTok{\textless{}{-}} \FunctionTok{loo\_compare}\NormalTok{(loo\_m\_on\_m, loo\_b\_on\_m)}
\FunctionTok{print}\NormalTok{(compare\_m)}
\end{Highlighting}
\end{Shaded}

\begin{verbatim}
       elpd_diff se_diff
model1    0.0       0.0 
model2 -337.9      23.4 
\end{verbatim}

\begin{Shaded}
\begin{Highlighting}[]
\NormalTok{z\_m }\OtherTok{\textless{}{-}} \FunctionTok{abs}\NormalTok{(compare\_m[}\DecValTok{2}\NormalTok{,}\StringTok{"elpd\_diff"}\NormalTok{]}\SpecialCharTok{/}\NormalTok{compare\_m[}\DecValTok{2}\NormalTok{,}\StringTok{"se\_diff"}\NormalTok{])}
\FunctionTok{cat}\NormalTok{(}\FunctionTok{sprintf}\NormalTok{(}\StringTok{"}\SpecialCharTok{\textbackslash{}n}\StringTok{|elpd\_diff|/SE = \%.1f  {-}\textgreater{}  \%s}\SpecialCharTok{\textbackslash{}n}\StringTok{"}\NormalTok{, z\_m,}
  \FunctionTok{ifelse}\NormalTok{(z\_m}\SpecialCharTok{\textgreater{}}\DecValTok{4}\NormalTok{,}\StringTok{"Clear preference."}\NormalTok{,}\FunctionTok{ifelse}\NormalTok{(z\_m}\SpecialCharTok{\textgreater{}}\DecValTok{2}\NormalTok{,}\StringTok{"Weak preference."}\NormalTok{,}\StringTok{"Inconclusive."}\NormalTok{))))}
\end{Highlighting}
\end{Shaded}

\begin{verbatim}

|elpd_diff|/SE = 14.4  ->  Clear preference.
\end{verbatim}

Reading the output with the Z-score rule
(\(|\text{elpd\_diff}|/\text{se\_diff}\)):

\begin{itemize}
\tightlist
\item
  \textbf{Biased data:} The memory model is slightly worse, but the
  ratio is likely below 2 --- \emph{inconclusive}. This is the correct
  scientific answer: when data come from a purely biased agent, a memory
  model with \(\beta \approx 0\) produces indistinguishable predictions.
  The PSIS-LOO score is not broken; it is honestly reporting that the
  simpler model is sufficient here.
\item
  \textbf{Memory data:} The biased model is decisively worse, with a
  ratio well above 4 --- \emph{clear preference}. The reversal at trial
  61 produces behavioral changes that the biased model cannot track, and
  PSIS-LOO detects this.
\end{itemize}

This is \textbf{successful model recovery}: the method correctly
identifies the true generating model in the column where detection is
possible, and honestly reports uncertainty in the column where the
models overlap. Both outcomes are scientifically valid.

\subsection{Pointwise ELPD Differences: Where Does the Advantage Come
From?}\label{pointwise-elpd-differences-where-does-the-advantage-come-from}

An aggregate ELPD difference tells us \emph{that} one model predicts
better overall. But often the most scientifically interesting question
is \emph{where} that difference comes from: on which types of trials
does the advantage concentrate?

\begin{Shaded}
\begin{Highlighting}[]
\NormalTok{make\_pointwise\_df }\OtherTok{\textless{}{-}} \ControlFlowTok{function}\NormalTok{(loo\_a, loo\_b, data\_df, label)}
  \FunctionTok{tibble}\NormalTok{(}\AttributeTok{obs =} \FunctionTok{seq\_len}\NormalTok{(}\FunctionTok{nrow}\NormalTok{(data\_df)),}
         \AttributeTok{elpd\_diff=}\NormalTok{loo\_a}\SpecialCharTok{$}\NormalTok{pointwise[,}\StringTok{"elpd\_loo"}\NormalTok{]}\SpecialCharTok{{-}}\NormalTok{loo\_b}\SpecialCharTok{$}\NormalTok{pointwise[,}\StringTok{"elpd\_loo"}\NormalTok{],}
         \AttributeTok{trial =}\NormalTok{ data\_df}\SpecialCharTok{$}\NormalTok{trial, }\AttributeTok{agent\_id=}\NormalTok{data\_df}\SpecialCharTok{$}\NormalTok{agent\_id, }\AttributeTok{label=}\NormalTok{label)}

\NormalTok{pw\_all }\OtherTok{\textless{}{-}} \FunctionTok{bind\_rows}\NormalTok{(}
  \FunctionTok{make\_pointwise\_df}\NormalTok{(loo\_b\_on\_b, loo\_m\_on\_b, d\_biased, }\StringTok{"Biased data: Biased {-} Memory"}\NormalTok{),}
  \FunctionTok{make\_pointwise\_df}\NormalTok{(loo\_m\_on\_m, loo\_b\_on\_m, d\_memory, }\StringTok{"Memory data: Memory {-} Biased"}\NormalTok{))}

\FunctionTok{ggplot}\NormalTok{(pw\_all, }\FunctionTok{aes}\NormalTok{(}\AttributeTok{x=}\NormalTok{trial, }\AttributeTok{y=}\NormalTok{elpd\_diff)) }\SpecialCharTok{+}
  \FunctionTok{geom\_hline}\NormalTok{(}\AttributeTok{yintercept=}\DecValTok{0}\NormalTok{, }\AttributeTok{linetype=}\StringTok{"dashed"}\NormalTok{, }\AttributeTok{color=}\StringTok{"gray40"}\NormalTok{) }\SpecialCharTok{+}
  \FunctionTok{geom\_point}\NormalTok{(}\FunctionTok{aes}\NormalTok{(}\AttributeTok{color=}\NormalTok{elpd\_diff}\SpecialCharTok{\textgreater{}}\DecValTok{0}\NormalTok{), }\AttributeTok{size=}\FloatTok{0.7}\NormalTok{, }\AttributeTok{alpha=}\FloatTok{0.35}\NormalTok{) }\SpecialCharTok{+}
  \FunctionTok{geom\_smooth}\NormalTok{(}\AttributeTok{method=}\StringTok{"gam"}\NormalTok{, }\AttributeTok{formula=}\NormalTok{y}\SpecialCharTok{\textasciitilde{}}\FunctionTok{s}\NormalTok{(x,}\AttributeTok{bs=}\StringTok{"cs"}\NormalTok{),}
              \AttributeTok{color=}\StringTok{"midnightblue"}\NormalTok{, }\AttributeTok{linewidth=}\DecValTok{1}\NormalTok{, }\AttributeTok{se=}\ConstantTok{TRUE}\NormalTok{) }\SpecialCharTok{+}
  \FunctionTok{facet\_wrap}\NormalTok{(}\SpecialCharTok{\textasciitilde{}}\NormalTok{label) }\SpecialCharTok{+}
  \FunctionTok{scale\_color\_manual}\NormalTok{(}\AttributeTok{values=}\FunctionTok{c}\NormalTok{(}\StringTok{"FALSE"}\OtherTok{=}\StringTok{"firebrick"}\NormalTok{,}\StringTok{"TRUE"}\OtherTok{=}\StringTok{"steelblue"}\NormalTok{),}
    \AttributeTok{labels=}\FunctionTok{c}\NormalTok{(}\StringTok{"FALSE"}\OtherTok{=}\StringTok{"Favors alternative"}\NormalTok{,}\StringTok{"TRUE"}\OtherTok{=}\StringTok{"Favors true model"}\NormalTok{), }\AttributeTok{name=}\ConstantTok{NULL}\NormalTok{) }\SpecialCharTok{+}
  \FunctionTok{labs}\NormalTok{(}\AttributeTok{title =} \StringTok{"Pointwise ELPD Differences Across Trials"}\NormalTok{)}
\end{Highlighting}
\end{Shaded}

\begin{center}
\includegraphics[width=0.8\linewidth,height=\textheight,keepaspectratio]{07-ModelComparison_files/figure-pdf/07-pointwise-elpd-1.pdf}
\end{center}

The growing advantage of the memory model over trials is a direct
cognitive signature: early in the game the running opponent rate is near
0.5 (uninformative), so both models predict roughly equally. As the
opponent's behavior becomes established --- typically after 20-30 trials
--- the memory model increasingly exploits that signal while the biased
model cannot. Once the reversal happens, the advantage disappears, but
then picks up again. The advantage concentrates precisely where the
cognitive mechanism operates. This is model comparison doing what it is
supposed to do: not just telling us which model wins, but illuminating
\emph{why} and \emph{when}.

\subsection{Bayesian Stacking: Combining Models Into an
Ensemble}\label{bayesian-stacking-combining-models-into-an-ensemble}

\texttt{loo\_compare()} produces a ranking, but it does not tell us how
much to trust each model. \textbf{Bayesian stacking} (Yao et al., 2018)
goes further by finding the \emph{optimal mixture} of the competing
models for prediction.

The idea is to find weights \(w_1, w_2\) (summing to 1) such that the
weighted combination of the two models' LOO predictive distributions is
as good as possible. Unlike Bayes factors, stacking does not assume that
one of the models is the truth --- it simply finds the best predictive
combination. A weight of nearly 1.0 for one model means the other adds
nothing useful. A more balanced split means both models are capturing
something the other misses.

\begin{Shaded}
\begin{Highlighting}[]
\NormalTok{stack\_b }\OtherTok{\textless{}{-}} \FunctionTok{loo\_model\_weights}\NormalTok{(}\FunctionTok{list}\NormalTok{(}\StringTok{"Biased"} \OtherTok{=}\NormalTok{ loo\_b\_on\_b, }\StringTok{"Memory"} \OtherTok{=}\NormalTok{ loo\_m\_on\_b),}
                             \AttributeTok{method =} \StringTok{"stacking"}\NormalTok{)}
\NormalTok{stack\_m }\OtherTok{\textless{}{-}} \FunctionTok{loo\_model\_weights}\NormalTok{(}\FunctionTok{list}\NormalTok{(}\StringTok{"Memory"}\OtherTok{=}\NormalTok{loo\_m\_on\_m, }\StringTok{"Biased"}\OtherTok{=}\NormalTok{loo\_b\_on\_m),}
                             \AttributeTok{method=}\StringTok{"stacking"}\NormalTok{)}

\FunctionTok{tibble}\NormalTok{(}
  \AttributeTok{Dataset  =} \FunctionTok{c}\NormalTok{(}\StringTok{"Biased"}\NormalTok{,}\StringTok{"Biased"}\NormalTok{,}\StringTok{"Memory"}\NormalTok{,}\StringTok{"Memory"}\NormalTok{),}
  \AttributeTok{Model    =} \FunctionTok{c}\NormalTok{(}\StringTok{"Biased"}\NormalTok{,}\StringTok{"Memory"}\NormalTok{,}\StringTok{"Memory"}\NormalTok{,}\StringTok{"Biased"}\NormalTok{),}
  \AttributeTok{Stacking =} \FunctionTok{round}\NormalTok{(}\FunctionTok{c}\NormalTok{(stack\_b[}\StringTok{"Biased"}\NormalTok{],stack\_b[}\StringTok{"Memory"}\NormalTok{],}
\NormalTok{                     stack\_m[}\StringTok{"Memory"}\NormalTok{],stack\_m[}\StringTok{"Biased"}\NormalTok{]), }\DecValTok{3}\NormalTok{),}
  \AttributeTok{True\_model =} \FunctionTok{c}\NormalTok{(}\ConstantTok{TRUE}\NormalTok{,}\ConstantTok{FALSE}\NormalTok{,}\ConstantTok{TRUE}\NormalTok{,}\ConstantTok{FALSE}\NormalTok{)}
\NormalTok{) }\SpecialCharTok{|\textgreater{}}
  \FunctionTok{ggplot}\NormalTok{(}\FunctionTok{aes}\NormalTok{(}\AttributeTok{x=}\NormalTok{Model, }\AttributeTok{y=}\NormalTok{Stacking, }\AttributeTok{fill=}\NormalTok{True\_model)) }\SpecialCharTok{+}
  \FunctionTok{geom\_col}\NormalTok{() }\SpecialCharTok{+}
  \CommentTok{\# Add labels, offset slightly above the bar}
  \FunctionTok{geom\_text}\NormalTok{(}\FunctionTok{aes}\NormalTok{(}\AttributeTok{label =}\NormalTok{ scales}\SpecialCharTok{::}\FunctionTok{percent}\NormalTok{(Stacking, }\AttributeTok{accuracy =} \FloatTok{0.1}\NormalTok{)), }
            \AttributeTok{vjust =} \SpecialCharTok{{-}}\FloatTok{0.5}\NormalTok{, }\AttributeTok{size =} \DecValTok{4}\NormalTok{) }\SpecialCharTok{+}
  \FunctionTok{facet\_wrap}\NormalTok{(}\SpecialCharTok{\textasciitilde{}}\NormalTok{Dataset, }\AttributeTok{labeller=}\NormalTok{label\_both) }\SpecialCharTok{+}
  \FunctionTok{scale\_fill\_manual}\NormalTok{(}\AttributeTok{values=}\FunctionTok{c}\NormalTok{(}\StringTok{"FALSE"}\OtherTok{=}\StringTok{"gray70"}\NormalTok{,}\StringTok{"TRUE"}\OtherTok{=}\StringTok{"steelblue"}\NormalTok{),}
    \AttributeTok{labels=}\FunctionTok{c}\NormalTok{(}\StringTok{"FALSE"}\OtherTok{=}\StringTok{"Alternative"}\NormalTok{,}\StringTok{"TRUE"}\OtherTok{=}\StringTok{"True model"}\NormalTok{), }\AttributeTok{name=}\ConstantTok{NULL}\NormalTok{) }\SpecialCharTok{+}
  \CommentTok{\# Expand the upper y{-}limit to 1.1 so the 100\% labels are not clipped}
  \FunctionTok{scale\_y\_continuous}\NormalTok{(}\AttributeTok{labels=}\NormalTok{scales}\SpecialCharTok{::}\NormalTok{percent, }\AttributeTok{limits=}\FunctionTok{c}\NormalTok{(}\DecValTok{0}\NormalTok{, }\FloatTok{1.1}\NormalTok{)) }\SpecialCharTok{+}
  \FunctionTok{labs}\NormalTok{(}\AttributeTok{title    =} \StringTok{"Bayesian Stacking Weights"}\NormalTok{,}
       \AttributeTok{subtitle =} \StringTok{"Successful recovery: true model receives near{-}100\% weight."}\NormalTok{,}
       \AttributeTok{x=}\ConstantTok{NULL}\NormalTok{, }\AttributeTok{y=}\StringTok{"Model weight"}\NormalTok{) }\SpecialCharTok{+}
  \FunctionTok{theme}\NormalTok{(}\AttributeTok{legend.position=}\StringTok{"bottom"}\NormalTok{)}
\end{Highlighting}
\end{Shaded}

\begin{center}
\includegraphics[width=0.8\linewidth,height=\textheight,keepaspectratio]{07-ModelComparison_files/figure-pdf/07-stacking-1.pdf}
\end{center}

\section{Evaluating Our Pipeline: Full Model
Recovery}\label{evaluating-our-pipeline-full-model-recovery}

Let us now pull together everything into the model recovery matrix ---
the central test of whether PSIS-LOO is doing its job. For each of the
two datasets, we check whether the method identifies the correct
generating model.

\begin{Shaded}
\begin{Highlighting}[]
\NormalTok{elpd\_matrix }\OtherTok{\textless{}{-}} \FunctionTok{tribble}\NormalTok{(}
  \SpecialCharTok{\textasciitilde{}}\NormalTok{Model,   }\SpecialCharTok{\textasciitilde{}}\NormalTok{Data,    }\SpecialCharTok{\textasciitilde{}}\NormalTok{ELPD,       }\SpecialCharTok{\textasciitilde{}}\NormalTok{SE,}
  \StringTok{"Biased"}\NormalTok{, }\StringTok{"Biased"}\NormalTok{, loo\_b\_on\_b}\SpecialCharTok{$}\NormalTok{estimates[}\StringTok{"elpd\_loo"}\NormalTok{,}\StringTok{"Estimate"}\NormalTok{], loo\_b\_on\_b}\SpecialCharTok{$}\NormalTok{estimates[}\StringTok{"elpd\_loo"}\NormalTok{,}\StringTok{"SE"}\NormalTok{],}
  \StringTok{"Memory"}\NormalTok{, }\StringTok{"Biased"}\NormalTok{, loo\_m\_on\_b}\SpecialCharTok{$}\NormalTok{estimates[}\StringTok{"elpd\_loo"}\NormalTok{,}\StringTok{"Estimate"}\NormalTok{], loo\_m\_on\_b}\SpecialCharTok{$}\NormalTok{estimates[}\StringTok{"elpd\_loo"}\NormalTok{,}\StringTok{"SE"}\NormalTok{],}
  \StringTok{"Biased"}\NormalTok{, }\StringTok{"Memory"}\NormalTok{, loo\_b\_on\_m}\SpecialCharTok{$}\NormalTok{estimates[}\StringTok{"elpd\_loo"}\NormalTok{,}\StringTok{"Estimate"}\NormalTok{], loo\_b\_on\_m}\SpecialCharTok{$}\NormalTok{estimates[}\StringTok{"elpd\_loo"}\NormalTok{,}\StringTok{"SE"}\NormalTok{],}
  \StringTok{"Memory"}\NormalTok{, }\StringTok{"Memory"}\NormalTok{, loo\_m\_on\_m}\SpecialCharTok{$}\NormalTok{estimates[}\StringTok{"elpd\_loo"}\NormalTok{,}\StringTok{"Estimate"}\NormalTok{], loo\_m\_on\_m}\SpecialCharTok{$}\NormalTok{estimates[}\StringTok{"elpd\_loo"}\NormalTok{,}\StringTok{"SE"}\NormalTok{]}
\NormalTok{) }\SpecialCharTok{|\textgreater{}}
  \FunctionTok{mutate}\NormalTok{(}\AttributeTok{is\_true=}\NormalTok{(Model}\SpecialCharTok{==}\StringTok{"Biased"}\SpecialCharTok{\&}\NormalTok{Data}\SpecialCharTok{==}\StringTok{"Biased"}\NormalTok{)}\SpecialCharTok{|}\NormalTok{(Model}\SpecialCharTok{==}\StringTok{"Memory"}\SpecialCharTok{\&}\NormalTok{Data}\SpecialCharTok{==}\StringTok{"Memory"}\NormalTok{),}
         \AttributeTok{label=}\FunctionTok{sprintf}\NormalTok{(}\StringTok{"\%.1f}\SpecialCharTok{\textbackslash{}n}\StringTok{(SE=\%.1f)"}\NormalTok{, ELPD, SE))}

\FunctionTok{ggplot}\NormalTok{(elpd\_matrix, }\FunctionTok{aes}\NormalTok{(}\AttributeTok{x=}\NormalTok{Data, }\AttributeTok{y=}\NormalTok{Model, }\AttributeTok{fill=}\NormalTok{ELPD)) }\SpecialCharTok{+}
  \FunctionTok{geom\_tile}\NormalTok{(}\AttributeTok{color=}\StringTok{"white"}\NormalTok{, }\AttributeTok{linewidth=}\FloatTok{1.5}\NormalTok{) }\SpecialCharTok{+}
  \FunctionTok{geom\_text}\NormalTok{(}\FunctionTok{aes}\NormalTok{(}\AttributeTok{label=}\NormalTok{label), }\AttributeTok{size=}\DecValTok{4}\NormalTok{, }\AttributeTok{fontface=}\StringTok{"bold"}\NormalTok{) }\SpecialCharTok{+}
  \FunctionTok{geom\_tile}\NormalTok{(}\AttributeTok{data=}\FunctionTok{filter}\NormalTok{(elpd\_matrix,is\_true),}
            \AttributeTok{color=}\StringTok{"midnightblue"}\NormalTok{, }\AttributeTok{fill=}\ConstantTok{NA}\NormalTok{, }\AttributeTok{linewidth=}\DecValTok{2}\NormalTok{) }\SpecialCharTok{+}
  \FunctionTok{scale\_fill\_viridis\_c}\NormalTok{(}\AttributeTok{option=}\StringTok{"D"}\NormalTok{, }\AttributeTok{name=}\StringTok{"ELPD"}\NormalTok{,}
    \AttributeTok{guide=}\FunctionTok{guide\_colorbar}\NormalTok{(}\AttributeTok{barwidth=}\DecValTok{10}\NormalTok{, }\AttributeTok{barheight=}\FloatTok{0.5}\NormalTok{)) }\SpecialCharTok{+}
  \FunctionTok{labs}\NormalTok{(}\AttributeTok{title    =} \StringTok{"Model Recovery Matrix"}\NormalTok{,}
       \AttributeTok{subtitle =} \StringTok{"Blue borders = true generating model. Higher ELPD = better prediction."}\NormalTok{,}
       \AttributeTok{x=}\StringTok{"Data{-}generating model"}\NormalTok{, }\AttributeTok{y=}\StringTok{"Fitted model"}\NormalTok{) }\SpecialCharTok{+}
  \FunctionTok{theme}\NormalTok{(}\AttributeTok{legend.position=}\StringTok{"bottom"}\NormalTok{)}
\end{Highlighting}
\end{Shaded}

\begin{center}
\includegraphics[width=0.8\linewidth,height=\textheight,keepaspectratio]{07-ModelComparison_files/figure-pdf/07-recovery-matrix-1.pdf}
\end{center}

Each column of this matrix answers a different scientific question. The
left column (biased data) asks: can PSIS-LOO detect that the memory
model is unnecessarily complex when a simpler model suffices? The right
column (memory data) asks: can PSIS-LOO detect that the biased model is
missing a real cognitive mechanism? Within each column, the model with
the higher ELPD should be the one that generated the data --- that is
the meaning of successful model recovery.

For context, the Ch. 7 memory model (fitted to the full 50-agent
reversal dataset with 80--240 trials per agent) had ELPD = -4233 (SE =
44.3). The Ch. 8 memory-on-memory fit uses a shorter, fixed-length
session (120 trials vs.~80--240), so its absolute ELPD is not directly
comparable --- but the sign and direction of all comparisons should be
consistent.

\begin{Shaded}
\begin{Highlighting}[]
\NormalTok{final\_table }\OtherTok{\textless{}{-}} \FunctionTok{tibble}\NormalTok{(}
  \StringTok{\textasciigrave{}}\AttributeTok{Fit}\StringTok{\textasciigrave{}}           \OtherTok{=} \FunctionTok{c}\NormalTok{(}\StringTok{"Biased{-}\textgreater{}Biased"}\NormalTok{,}\StringTok{"Memory{-}\textgreater{}Biased"}\NormalTok{,}\StringTok{"Biased{-}\textgreater{}Memory"}\NormalTok{,}\StringTok{"Memory{-}\textgreater{}Memory"}\NormalTok{),}
  \StringTok{\textasciigrave{}}\AttributeTok{ELPD}\StringTok{\textasciigrave{}}          \OtherTok{=}\NormalTok{ pareto\_k\_table}\SpecialCharTok{$}\NormalTok{ELPD,}
  \StringTok{\textasciigrave{}}\AttributeTok{SE}\StringTok{\textasciigrave{}}            \OtherTok{=}\NormalTok{ pareto\_k\_table}\SpecialCharTok{$}\NormalTok{SE,}
  \StringTok{\textasciigrave{}}\AttributeTok{p\_loo/N}\StringTok{\textasciigrave{}}       \OtherTok{=}\NormalTok{ pareto\_k\_table}\SpecialCharTok{$}\NormalTok{ratio,}
  \StringTok{\textasciigrave{}}\AttributeTok{POT{-}C p}\StringTok{\textasciigrave{}}       \OtherTok{=} \FunctionTok{c}\NormalTok{(}\FunctionTok{round}\NormalTok{(pit\_b\_on\_b}\SpecialCharTok{$}\NormalTok{result}\SpecialCharTok{$}\NormalTok{p\_global,}\DecValTok{3}\NormalTok{),}
                       \FunctionTok{round}\NormalTok{(pit\_m\_on\_b}\SpecialCharTok{$}\NormalTok{result}\SpecialCharTok{$}\NormalTok{p\_global,}\DecValTok{3}\NormalTok{),}
                       \FunctionTok{round}\NormalTok{(pit\_b\_on\_m}\SpecialCharTok{$}\NormalTok{result}\SpecialCharTok{$}\NormalTok{p\_global,}\DecValTok{3}\NormalTok{),}
                       \FunctionTok{round}\NormalTok{(pit\_m\_on\_m}\SpecialCharTok{$}\NormalTok{result}\SpecialCharTok{$}\NormalTok{p\_global,}\DecValTok{3}\NormalTok{)),}
  \StringTok{\textasciigrave{}}\AttributeTok{True model?}\StringTok{\textasciigrave{}}   \OtherTok{=} \FunctionTok{c}\NormalTok{(}\StringTok{"YES"}\NormalTok{,}\StringTok{"no"}\NormalTok{,}\StringTok{"no"}\NormalTok{,}\StringTok{"YES"}\NormalTok{),}
  \StringTok{\textasciigrave{}}\AttributeTok{LOO preferred?}\StringTok{\textasciigrave{}}\OtherTok{=} \FunctionTok{c}\NormalTok{(}
    \CommentTok{\# Row 1: Does Biased beat Memory (on Biased data)?}
    \FunctionTok{ifelse}\NormalTok{(loo\_b\_on\_b}\SpecialCharTok{$}\NormalTok{estimates[}\StringTok{"elpd\_loo"}\NormalTok{,}\StringTok{"Estimate"}\NormalTok{] }\SpecialCharTok{\textgreater{}}
\NormalTok{           loo\_m\_on\_b}\SpecialCharTok{$}\NormalTok{estimates[}\StringTok{"elpd\_loo"}\NormalTok{,}\StringTok{"Estimate"}\NormalTok{],}\StringTok{"YES"}\NormalTok{,}\StringTok{"no"}\NormalTok{),}
    
    \CommentTok{\# Row 2: Does Memory beat Biased (on Biased data)?}
    \FunctionTok{ifelse}\NormalTok{(loo\_m\_on\_b}\SpecialCharTok{$}\NormalTok{estimates[}\StringTok{"elpd\_loo"}\NormalTok{,}\StringTok{"Estimate"}\NormalTok{] }\SpecialCharTok{\textgreater{}}
\NormalTok{           loo\_b\_on\_b}\SpecialCharTok{$}\NormalTok{estimates[}\StringTok{"elpd\_loo"}\NormalTok{,}\StringTok{"Estimate"}\NormalTok{],}\StringTok{"YES"}\NormalTok{,}\StringTok{"no"}\NormalTok{),}
    
    \CommentTok{\# Row 3: Does Biased beat Memory (on Memory data)?  \textless{}{-}{-} FIXED THIS ONE}
    \FunctionTok{ifelse}\NormalTok{(loo\_b\_on\_m}\SpecialCharTok{$}\NormalTok{estimates[}\StringTok{"elpd\_loo"}\NormalTok{,}\StringTok{"Estimate"}\NormalTok{] }\SpecialCharTok{\textgreater{}}
\NormalTok{           loo\_m\_on\_m}\SpecialCharTok{$}\NormalTok{estimates[}\StringTok{"elpd\_loo"}\NormalTok{,}\StringTok{"Estimate"}\NormalTok{],}\StringTok{"YES"}\NormalTok{,}\StringTok{"no"}\NormalTok{),}
    
    \CommentTok{\# Row 4: Does Memory beat Biased (on Memory data)?}
    \FunctionTok{ifelse}\NormalTok{(loo\_m\_on\_m}\SpecialCharTok{$}\NormalTok{estimates[}\StringTok{"elpd\_loo"}\NormalTok{,}\StringTok{"Estimate"}\NormalTok{] }\SpecialCharTok{\textgreater{}}
\NormalTok{           loo\_b\_on\_m}\SpecialCharTok{$}\NormalTok{estimates[}\StringTok{"elpd\_loo"}\NormalTok{,}\StringTok{"Estimate"}\NormalTok{],}\StringTok{"YES"}\NormalTok{,}\StringTok{"no"}\NormalTok{))}
\NormalTok{)}

\NormalTok{knitr}\SpecialCharTok{::}\FunctionTok{kable}\NormalTok{(final\_table, }\AttributeTok{caption=}\FunctionTok{paste}\NormalTok{(}
  \StringTok{"Full comparison summary. p\_loo/N: use POT{-}C if \textgreater{} 0.05."}\NormalTok{,}
  \StringTok{"POT{-}C p: formal uniformity test, valid under dependence."}\NormalTok{,}
  \StringTok{"Recovery: \textquotesingle{}True model?\textquotesingle{} should match \textquotesingle{}LOO preferred?\textquotesingle{}."}\NormalTok{))}
\end{Highlighting}
\end{Shaded}

\begin{longtable}[]{@{}
  >{\raggedright\arraybackslash}p{(\linewidth - 12\tabcolsep) * \real{0.2113}}
  >{\raggedleft\arraybackslash}p{(\linewidth - 12\tabcolsep) * \real{0.1127}}
  >{\raggedleft\arraybackslash}p{(\linewidth - 12\tabcolsep) * \real{0.0704}}
  >{\raggedleft\arraybackslash}p{(\linewidth - 12\tabcolsep) * \real{0.1127}}
  >{\raggedleft\arraybackslash}p{(\linewidth - 12\tabcolsep) * \real{0.1127}}
  >{\raggedright\arraybackslash}p{(\linewidth - 12\tabcolsep) * \real{0.1690}}
  >{\raggedright\arraybackslash}p{(\linewidth - 12\tabcolsep) * \real{0.2113}}@{}}
\caption{Full comparison summary. p\_loo/N: use POT-C if \textgreater{}
0.05. POT-C p: formal uniformity test, valid under dependence. Recovery:
`True model?' should match `LOO preferred?'.}\tabularnewline
\toprule\noalign{}
\begin{minipage}[b]{\linewidth}\raggedright
Fit
\end{minipage} & \begin{minipage}[b]{\linewidth}\raggedleft
ELPD
\end{minipage} & \begin{minipage}[b]{\linewidth}\raggedleft
SE
\end{minipage} & \begin{minipage}[b]{\linewidth}\raggedleft
p\_loo/N
\end{minipage} & \begin{minipage}[b]{\linewidth}\raggedleft
POT-C p
\end{minipage} & \begin{minipage}[b]{\linewidth}\raggedright
True model?
\end{minipage} & \begin{minipage}[b]{\linewidth}\raggedright
LOO preferred?
\end{minipage} \\
\midrule\noalign{}
\endfirsthead
\toprule\noalign{}
\begin{minipage}[b]{\linewidth}\raggedright
Fit
\end{minipage} & \begin{minipage}[b]{\linewidth}\raggedleft
ELPD
\end{minipage} & \begin{minipage}[b]{\linewidth}\raggedleft
SE
\end{minipage} & \begin{minipage}[b]{\linewidth}\raggedleft
p\_loo/N
\end{minipage} & \begin{minipage}[b]{\linewidth}\raggedleft
POT-C p
\end{minipage} & \begin{minipage}[b]{\linewidth}\raggedright
True model?
\end{minipage} & \begin{minipage}[b]{\linewidth}\raggedright
LOO preferred?
\end{minipage} \\
\midrule\noalign{}
\endhead
\bottomrule\noalign{}
\endlastfoot
Biased-\textgreater Biased & -3393.8 & 36.0 & 0.007 & 0 & YES & YES \\
Memory-\textgreater Biased & -3395.1 & 36.0 & 0.008 & 0 & no & no \\
Biased-\textgreater Memory & -4038.3 & 15.4 & 0.006 & 0 & no & no \\
Memory-\textgreater Memory & -3700.3 & 27.3 & 0.010 & 0 & YES & YES \\
\end{longtable}

\subsection{Power Analysis and Experimental
Design}\label{power-analysis-and-experimental-design}

e just saw a successful single instance of model recovery: for our
specific simulated dataset, loo\_compare() correctly preferred the true
generating model in both cases. But one lucky simulation does not
guarantee our experimental design is robust.

If a participant truly uses a simple Biased strategy, is our
experimental design rich enough for LOO cross-validation to correctly
prefer the Biased model over the Memory model? Conversely, if they truly
use the Memory model, do we have enough trials to confidently detect
that complexity, or will the simpler Biased model win by default due to
its lower complexity penalty?

Testing this is called a Model Recovery Analysis, and it serves as the
Bayesian cognitive modeling equivalent of a statistical power analysis.

\subsection{Measuring Power: The F1
Score}\label{measuring-power-the-f1-score}

To run a model recovery study, we simulate hundreds of datasets from
both models across a variety of experimental designs (varying the number
of participants and trials) and cross-fit both models to see which
wins.Because model comparison penalizes complexity, the errors are
rarely symmetrical. A complex model (like Memory) might be hard to
detect if the effect is weak (low Recall), but if LOO does select it, it
is almost certainly the true model (high Precision). Conversely, the
simpler model (Biased) acts as a default fallback: it is easily
recovered when true, but is often falsely selected when the true Memory
mechanism is too weak to detect.To capture this asymmetry in a single
metric, we calculate the F1 Score for each model at every experimental
design point. The F1 score is the harmonic mean of Precision and
Recall:\[F1 = 2 \times \frac{\text{Precision} \times \text{Recall}}{\text{Precision} + \text{Recall}}\]
An F1 score near 1.0 means we perfectly and reliably identify the model.
An F1 score below 0.80 suggests our experimental design is underpowered
for that specific cognitive mechanism.

\subsection{Designing the Experiment: Trials, Participants, and Task
Structure}\label{designing-the-experiment-trials-participants-and-task-structure}

In hierarchical cognitive modeling, we have two distinct sample sizes to
worry about: the number of trials per participant (\(N\)), and the
number of participants (\(J\)).These two dimensions govern different
types of statistical power:

\begin{itemize}
\item
  More Trials (\(N\)): Increases our ability to accurately estimate
  individual-level parameters. If we only have 30 trials, we cannot
  confidently detect the subtle influence of the Memory parameter within
  a single person's noisy choices.
\item
  More Participants (\(J\)): Increases our ability to accurately
  estimate the population-level distribution.
\end{itemize}

To conduct a formal power analysis for model comparison, we
systematically repeat our recovery pipeline across a grid of design
choices. We then plot the F1 score as a 2D heatmap to find the optimal
experimental threshold where both models can be reliably identified.

\begin{Shaded}
\begin{Highlighting}[]
\CommentTok{\# {-}{-}{-} Conceptual Power Trade{-}off Matrix (F1 Scores) {-}{-}{-}}
\NormalTok{power\_grid }\OtherTok{\textless{}{-}} \FunctionTok{expand\_grid}\NormalTok{(}
  \AttributeTok{N\_Trials =} \FunctionTok{c}\NormalTok{(}\DecValTok{30}\NormalTok{, }\DecValTok{60}\NormalTok{, }\DecValTok{90}\NormalTok{, }\DecValTok{120}\NormalTok{, }\DecValTok{150}\NormalTok{),}
  \AttributeTok{N\_Participants =} \FunctionTok{c}\NormalTok{(}\DecValTok{20}\NormalTok{, }\DecValTok{40}\NormalTok{, }\DecValTok{60}\NormalTok{, }\DecValTok{80}\NormalTok{, }\DecValTok{100}\NormalTok{)}
\NormalTok{) }\SpecialCharTok{|\textgreater{}}
  \FunctionTok{mutate}\NormalTok{(}
    \CommentTok{\# Simulate Recall (True Positive Rate)}
    \CommentTok{\# Memory is hard to detect with few trials.}
    \AttributeTok{Recall\_Memory =} \DecValTok{1} \SpecialCharTok{{-}} \FunctionTok{exp}\NormalTok{(}\SpecialCharTok{{-}}\FloatTok{0.025} \SpecialCharTok{*}\NormalTok{ N\_Trials) }\SpecialCharTok{*} \FunctionTok{exp}\NormalTok{(}\SpecialCharTok{{-}}\FloatTok{0.02} \SpecialCharTok{*}\NormalTok{ N\_Participants),}
    \CommentTok{\# Biased is easy to detect when true.}
    \AttributeTok{Recall\_Biased =} \DecValTok{1} \SpecialCharTok{{-}} \FunctionTok{exp}\NormalTok{(}\SpecialCharTok{{-}}\FloatTok{0.01} \SpecialCharTok{*}\NormalTok{ N\_Trials) }\SpecialCharTok{*} \FunctionTok{exp}\NormalTok{(}\SpecialCharTok{{-}}\FloatTok{0.01} \SpecialCharTok{*}\NormalTok{ N\_Participants), }

    \CommentTok{\# Simulate Precision (Positive Predictive Value)}
    \CommentTok{\# If LOO picks Memory, it\textquotesingle{}s almost certainly right (complex models rarely win by accident).}
    \AttributeTok{Precision\_Memory =} \FunctionTok{pmin}\NormalTok{(}\FloatTok{0.99}\NormalTok{, }\FloatTok{0.80} \SpecialCharTok{+} \FloatTok{0.001} \SpecialCharTok{*}\NormalTok{ N\_Trials }\SpecialCharTok{+} \FloatTok{0.001} \SpecialCharTok{*}\NormalTok{ N\_Participants),}
    \CommentTok{\# If LOO picks Biased, it might be a weak Memory model falling back (lower precision).}
    \AttributeTok{Precision\_Biased =} \DecValTok{1} \SpecialCharTok{{-}} \FunctionTok{exp}\NormalTok{(}\SpecialCharTok{{-}}\FloatTok{0.015} \SpecialCharTok{*}\NormalTok{ N\_Trials) }\SpecialCharTok{*} \FunctionTok{exp}\NormalTok{(}\SpecialCharTok{{-}}\FloatTok{0.02} \SpecialCharTok{*}\NormalTok{ N\_Participants)}
\NormalTok{  ) }\SpecialCharTok{|\textgreater{}}
  \FunctionTok{mutate}\NormalTok{(}
    \CommentTok{\# Calculate F1 Score}
    \AttributeTok{F1\_Memory =} \DecValTok{2} \SpecialCharTok{*}\NormalTok{ (Precision\_Memory }\SpecialCharTok{*}\NormalTok{ Recall\_Memory) }\SpecialCharTok{/}\NormalTok{ (Precision\_Memory }\SpecialCharTok{+}\NormalTok{ Recall\_Memory),}
    \AttributeTok{F1\_Biased =} \DecValTok{2} \SpecialCharTok{*}\NormalTok{ (Precision\_Biased }\SpecialCharTok{*}\NormalTok{ Recall\_Biased) }\SpecialCharTok{/}\NormalTok{ (Precision\_Biased }\SpecialCharTok{+}\NormalTok{ Recall\_Biased)}
\NormalTok{  ) }\SpecialCharTok{|\textgreater{}}
  \FunctionTok{mutate}\NormalTok{(}\FunctionTok{across}\NormalTok{(}\FunctionTok{starts\_with}\NormalTok{(}\StringTok{"F1"}\NormalTok{), }\SpecialCharTok{\textasciitilde{}}\FunctionTok{pmin}\NormalTok{(}\FloatTok{0.99}\NormalTok{, }\FunctionTok{round}\NormalTok{(.x, }\DecValTok{2}\NormalTok{)))) }\SpecialCharTok{|\textgreater{}}
\NormalTok{  dplyr}\SpecialCharTok{::}\FunctionTok{select}\NormalTok{(N\_Trials, N\_Participants, F1\_Memory, F1\_Biased) }\SpecialCharTok{|\textgreater{}}
  \FunctionTok{pivot\_longer}\NormalTok{(}\AttributeTok{cols =} \FunctionTok{starts\_with}\NormalTok{(}\StringTok{"F1"}\NormalTok{), }\AttributeTok{names\_to =} \StringTok{"Model"}\NormalTok{, }
               \AttributeTok{names\_prefix =} \StringTok{"F1\_"}\NormalTok{, }\AttributeTok{values\_to =} \StringTok{"F1\_Score"}\NormalTok{)}

\CommentTok{\# Plot the surfaces side{-}by{-}side}
\FunctionTok{ggplot}\NormalTok{(power\_grid, }\FunctionTok{aes}\NormalTok{(}\AttributeTok{x =} \FunctionTok{factor}\NormalTok{(N\_Trials), }\AttributeTok{y =} \FunctionTok{factor}\NormalTok{(N\_Participants), }\AttributeTok{fill =}\NormalTok{ F1\_Score)) }\SpecialCharTok{+}
  \FunctionTok{geom\_tile}\NormalTok{(}\AttributeTok{color =} \StringTok{"white"}\NormalTok{, }\AttributeTok{linewidth =} \DecValTok{1}\NormalTok{) }\SpecialCharTok{+}
  \FunctionTok{geom\_text}\NormalTok{(}\FunctionTok{aes}\NormalTok{(}\AttributeTok{label =}\NormalTok{ scales}\SpecialCharTok{::}\FunctionTok{percent}\NormalTok{(F1\_Score, }\AttributeTok{accuracy =} \DecValTok{1}\NormalTok{)), }
            \AttributeTok{size =} \DecValTok{4}\NormalTok{, }\AttributeTok{fontface =} \StringTok{"bold"}\NormalTok{, }
            \AttributeTok{color =} \FunctionTok{ifelse}\NormalTok{(power\_grid}\SpecialCharTok{$}\NormalTok{F1\_Score }\SpecialCharTok{\textgreater{}} \FloatTok{0.80}\NormalTok{, }\StringTok{"white"}\NormalTok{, }\StringTok{"black"}\NormalTok{)) }\SpecialCharTok{+}
  \FunctionTok{facet\_wrap}\NormalTok{(}\SpecialCharTok{\textasciitilde{}}\NormalTok{Model, }\AttributeTok{labeller =} \FunctionTok{as\_labeller}\NormalTok{(}\FunctionTok{c}\NormalTok{(}\AttributeTok{Biased =} \StringTok{"Biased Model Recovery (F1)"}\NormalTok{, }
                                              \AttributeTok{Memory =} \StringTok{"Memory Model Recovery (F1)"}\NormalTok{))) }\SpecialCharTok{+}
  \FunctionTok{scale\_fill\_viridis\_c}\NormalTok{(}\AttributeTok{option =} \StringTok{"mako"}\NormalTok{, }\AttributeTok{direction =} \SpecialCharTok{{-}}\DecValTok{1}\NormalTok{, }\AttributeTok{limits =} \FunctionTok{c}\NormalTok{(}\DecValTok{0}\NormalTok{, }\DecValTok{1}\NormalTok{), }
                       \AttributeTok{labels =}\NormalTok{ scales}\SpecialCharTok{::}\NormalTok{percent) }\SpecialCharTok{+}
       \FunctionTok{labs}\NormalTok{(}\AttributeTok{title =} \StringTok{"Model Recovery Power Analysis: Trials × Participants"}\NormalTok{,}
       \AttributeTok{subtitle =} \FunctionTok{paste0}\NormalTok{(}\StringTok{"F1 \textgreater{} 80\% (white cells) = reliably discriminable design. "}\NormalTok{,}
                         \StringTok{"Memory is harder to detect than Biased — it needs both more trials AND more participants."}\NormalTok{),}
       \AttributeTok{x =} \StringTok{"Number of Trials per Participant (N)"}\NormalTok{,}
       \AttributeTok{y =} \StringTok{"Number of Participants (J)"}\NormalTok{,}
       \AttributeTok{fill =} \StringTok{"F1 Score"}\NormalTok{) }\SpecialCharTok{+}
  \FunctionTok{theme\_minimal}\NormalTok{() }\SpecialCharTok{+}
  \FunctionTok{theme}\NormalTok{(}\AttributeTok{panel.grid =} \FunctionTok{element\_blank}\NormalTok{(),}
        \AttributeTok{axis.text =} \FunctionTok{element\_text}\NormalTok{(}\AttributeTok{size =} \DecValTok{10}\NormalTok{),}
        \AttributeTok{axis.title =} \FunctionTok{element\_text}\NormalTok{(}\AttributeTok{size =} \DecValTok{12}\NormalTok{, }\AttributeTok{face =} \StringTok{"bold"}\NormalTok{),}
        \AttributeTok{strip.text =} \FunctionTok{element\_text}\NormalTok{(}\AttributeTok{size =} \DecValTok{12}\NormalTok{, }\AttributeTok{face =} \StringTok{"bold"}\NormalTok{))}
\end{Highlighting}
\end{Shaded}

\begin{center}
\includegraphics[width=0.8\linewidth,height=\textheight,keepaspectratio]{07-ModelComparison_files/figure-pdf/07-model-recovery-matrix-1.pdf}
\end{center}

Looking at this trade-off matrix, we can definitively make a decision:
``To reliably distinguish between Biased and Memory strategies (e.g., F1
\textgreater{} 80\% for both models), we cannot rely on large
participant counts alone. Running 100 participants for only 30 trials
yields inadequate power for recovering the Memory mechanism (F1 = 66\%).
We must ensure sufficient within-subject data, requiring a minimum
design of roughly 60 participants completing 90 trials each to achieve
reliable discriminability across the board.''

\subsection{The Impact of Task Design: Better Data vs.~More
Data}\label{the-impact-of-task-design-better-data-vs.-more-data}

But power in cognitive modeling isn't just about collecting more data;
it is often about collecting better data. We can vary the experimental
setup itself.Consider why we explicitly forced the opponent to switch
from an 80\%-right bias to a 20\%-right bias halfway through our
Matching Pennies simulation. What would happen if we just let the
opponent play a constant 80\%-right bias for the entire game?If the
opponent's behavior is constant, the Memory agent's internal running
average (opp\_rate\_prev) quickly settles near 0.8 and stays there.
Behaviorally, an agent responding to a flat memory signal is
mathematically indistinguishable from a Biased agent with a fixed
probability. Because the independent variable (the memory signal) lacks
variance, the two models become collinear. When models are collinear,
PSIS-LOO will heavily penalize the Memory model for its extra \(\beta\)
parameter, and the simpler Biased model will falsely ``win'' the
comparison.We can formalize this by adding ``Task Design'' as an
explicit dimension in our power analysis. Below, we simulate the F1
recovery score for the Memory model under two different experimental
designs: a Constant 80\% Bias and an 80\% \(\to\) 20\% Reversal.

\begin{Shaded}
\begin{Highlighting}[]
\CommentTok{\# {-}{-}{-} Conceptual Power Trade{-}off: Task Design Impact {-}{-}{-}}
\NormalTok{design\_grid }\OtherTok{\textless{}{-}} \FunctionTok{expand\_grid}\NormalTok{(}
  \AttributeTok{N\_Trials =} \FunctionTok{c}\NormalTok{(}\DecValTok{30}\NormalTok{, }\DecValTok{60}\NormalTok{, }\DecValTok{90}\NormalTok{, }\DecValTok{120}\NormalTok{, }\DecValTok{150}\NormalTok{),}
  \AttributeTok{N\_Participants =} \FunctionTok{c}\NormalTok{(}\DecValTok{20}\NormalTok{, }\DecValTok{40}\NormalTok{, }\DecValTok{60}\NormalTok{, }\DecValTok{80}\NormalTok{, }\DecValTok{100}\NormalTok{),}
  \AttributeTok{Task\_Design =} \FunctionTok{c}\NormalTok{(}\StringTok{"Constant Bias (80\%)"}\NormalTok{, }\StringTok{"Reversal (80\% {-}\textgreater{} 20\%)"}\NormalTok{)}
\NormalTok{) }\SpecialCharTok{|\textgreater{}}
  \FunctionTok{mutate}\NormalTok{(}
    \CommentTok{\# Simulate Recall: Constant bias severely caps the ability to detect memory}
    \AttributeTok{Recall =} \FunctionTok{ifelse}\NormalTok{(Task\_Design }\SpecialCharTok{==} \StringTok{"Reversal (80\% {-}\textgreater{} 20\%)"}\NormalTok{,}
                    \DecValTok{1} \SpecialCharTok{{-}} \FunctionTok{exp}\NormalTok{(}\SpecialCharTok{{-}}\FloatTok{0.025} \SpecialCharTok{*}\NormalTok{ N\_Trials) }\SpecialCharTok{*} \FunctionTok{exp}\NormalTok{(}\SpecialCharTok{{-}}\FloatTok{0.02} \SpecialCharTok{*}\NormalTok{ N\_Participants),}
                    \FloatTok{0.60} \SpecialCharTok{{-}} \FunctionTok{exp}\NormalTok{(}\SpecialCharTok{{-}}\FloatTok{0.01} \SpecialCharTok{*}\NormalTok{ N\_Trials) }\SpecialCharTok{*} \FunctionTok{exp}\NormalTok{(}\SpecialCharTok{{-}}\FloatTok{0.01} \SpecialCharTok{*}\NormalTok{ N\_Participants)),}
    
    \CommentTok{\# Simulate Precision: Constant bias makes the model uncertain even when it wins}
    \AttributeTok{Precision =} \FunctionTok{ifelse}\NormalTok{(Task\_Design }\SpecialCharTok{==} \StringTok{"Reversal (80\% {-}\textgreater{} 20\%)"}\NormalTok{,}
                       \FunctionTok{pmin}\NormalTok{(}\FloatTok{0.99}\NormalTok{, }\FloatTok{0.80} \SpecialCharTok{+} \FloatTok{0.001} \SpecialCharTok{*}\NormalTok{ N\_Trials }\SpecialCharTok{+} \FloatTok{0.001} \SpecialCharTok{*}\NormalTok{ N\_Participants),}
                       \FloatTok{0.70} \SpecialCharTok{+} \FloatTok{0.0005} \SpecialCharTok{*}\NormalTok{ N\_Trials }\SpecialCharTok{+} \FloatTok{0.0005} \SpecialCharTok{*}\NormalTok{ N\_Participants)}
\NormalTok{  ) }\SpecialCharTok{|\textgreater{}}
  \FunctionTok{mutate}\NormalTok{(}
    \AttributeTok{F1\_Score =} \DecValTok{2} \SpecialCharTok{*}\NormalTok{ (Precision }\SpecialCharTok{*}\NormalTok{ Recall) }\SpecialCharTok{/}\NormalTok{ (Precision }\SpecialCharTok{+}\NormalTok{ Recall),}
    \AttributeTok{F1\_Score =} \FunctionTok{pmax}\NormalTok{(}\DecValTok{0}\NormalTok{, }\FunctionTok{pmin}\NormalTok{(}\FloatTok{0.99}\NormalTok{, }\FunctionTok{round}\NormalTok{(F1\_Score, }\DecValTok{2}\NormalTok{))) }\CommentTok{\# Bound between 0 and 0.99}
\NormalTok{  )}

\CommentTok{\# Plot the impact of Task Design}
\FunctionTok{ggplot}\NormalTok{(design\_grid, }\FunctionTok{aes}\NormalTok{(}\AttributeTok{x =} \FunctionTok{factor}\NormalTok{(N\_Trials), }\AttributeTok{y =} \FunctionTok{factor}\NormalTok{(N\_Participants), }\AttributeTok{fill =}\NormalTok{ F1\_Score)) }\SpecialCharTok{+}
  \FunctionTok{geom\_tile}\NormalTok{(}\AttributeTok{color =} \StringTok{"white"}\NormalTok{, }\AttributeTok{linewidth =} \DecValTok{1}\NormalTok{) }\SpecialCharTok{+}
  \FunctionTok{geom\_text}\NormalTok{(}\FunctionTok{aes}\NormalTok{(}\AttributeTok{label =}\NormalTok{ scales}\SpecialCharTok{::}\FunctionTok{percent}\NormalTok{(F1\_Score, }\AttributeTok{accuracy =} \DecValTok{1}\NormalTok{)), }
            \AttributeTok{size =} \DecValTok{4}\NormalTok{, }\AttributeTok{fontface =} \StringTok{"bold"}\NormalTok{, }
            \AttributeTok{color =} \FunctionTok{ifelse}\NormalTok{(design\_grid}\SpecialCharTok{$}\NormalTok{F1\_Score }\SpecialCharTok{\textgreater{}} \FloatTok{0.75}\NormalTok{, }\StringTok{"white"}\NormalTok{, }\StringTok{"black"}\NormalTok{)) }\SpecialCharTok{+}
  \FunctionTok{facet\_wrap}\NormalTok{(}\SpecialCharTok{\textasciitilde{}}\NormalTok{Task\_Design) }\SpecialCharTok{+}
  \FunctionTok{scale\_fill\_viridis\_c}\NormalTok{(}\AttributeTok{option =} \StringTok{"mako"}\NormalTok{, }\AttributeTok{direction =} \SpecialCharTok{{-}}\DecValTok{1}\NormalTok{, }\AttributeTok{limits =} \FunctionTok{c}\NormalTok{(}\DecValTok{0}\NormalTok{, }\DecValTok{1}\NormalTok{), }
                       \AttributeTok{labels =}\NormalTok{ scales}\SpecialCharTok{::}\NormalTok{percent) }\SpecialCharTok{+}
  \FunctionTok{labs}\NormalTok{(}\AttributeTok{title =} \StringTok{"The Impact of Task Design on Model Recovery"}\NormalTok{,}
       \AttributeTok{subtitle =} \StringTok{"F1 Score for recovering the Memory Model. A poor design cannot be saved by sample size."}\NormalTok{,}
       \AttributeTok{x =} \StringTok{"Number of Trials per Participant (N)"}\NormalTok{,}
       \AttributeTok{y =} \StringTok{"Number of Participants (J)"}\NormalTok{,}
       \AttributeTok{fill =} \StringTok{"F1 Score"}\NormalTok{) }\SpecialCharTok{+}
  \FunctionTok{theme\_minimal}\NormalTok{() }\SpecialCharTok{+}
  \FunctionTok{theme}\NormalTok{(}\AttributeTok{panel.grid =} \FunctionTok{element\_blank}\NormalTok{(),}
        \AttributeTok{axis.text =} \FunctionTok{element\_text}\NormalTok{(}\AttributeTok{size =} \DecValTok{10}\NormalTok{),}
        \AttributeTok{axis.title =} \FunctionTok{element\_text}\NormalTok{(}\AttributeTok{size =} \DecValTok{12}\NormalTok{, }\AttributeTok{face =} \StringTok{"bold"}\NormalTok{),}
        \AttributeTok{strip.text =} \FunctionTok{element\_text}\NormalTok{(}\AttributeTok{size =} \DecValTok{12}\NormalTok{, }\AttributeTok{face =} \StringTok{"bold"}\NormalTok{))}
\end{Highlighting}
\end{Shaded}

\begin{center}
\includegraphics[width=0.8\linewidth,height=\textheight,keepaspectratio]{07-ModelComparison_files/figure-pdf/ch08_power_task_design_viz-1.pdf}
\end{center}

This visualization delivers a stark warning. If you use a Constant Bias
design, even running a massive, expensive study with 100 participants
and 150 trials yields a dismal F1 score of roughly 65\%. The models are
structurally overlapping, and throwing more data at the problem barely
moves the needle.By simply changing the code of the task to include a
Reversal (costing exactly \(0\) extra dollars and \(0\) extra minutes of
participant time), a much smaller study of just 40 participants and 90
trials achieves a highly reliable 84\% F1 score.

\subsection{Beyond Binary Recovery: Simulation-Based Calibration Checks
for
Models}\label{beyond-binary-recovery-simulation-based-calibration-checks-for-models}

\subsection{From Binary Wins to Probabilistic
Power}\label{from-binary-wins-to-probabilistic-power}

Traditional model recovery treats model comparison as a binary,
``winner-takes-all'' election. If ELPD\_A \textgreater{} ELPD\_B by even
0.1 nats, Model A gets a ``1'' for accuracy, and Model B gets a ``0''.
This binary approach is mathematically impoverished because it hides our
uncertainty.

In a fully Bayesian workflow, we prefer a probabilistic approach.
Instead of asking ``Did the true model win?'', we can ask ``How much
stacking weight does the true model receive?''

This transforms our understanding of statistical power. Instead of a
binary accuracy curve, we can plot the average stacking weight assigned
to the true model across simulations. If the Biased model generated the
data, but its average stacking weight is only 0.55, our experimental
design is highly uncertain and under-powered, even if the Biased model
technically ``wins'' most of the time. Using continuous stacking weights
introduces a natural, continuous gradient of uncertainty into our model
identification. By using stacking weights, we move from a discrete
`Correct/Incorrect' recovery metric to a continuous measure of
probabilistic power. This allows us to quantify exactly how much
uncertainty remains in our model identification even when the true model
`wins' the comparison.

This realization leads us directly to the ultimate check of our entire
model comparison pipeline: Model Simulation-Based Calibration Checks
(Model SBC).

\subsection{\texorpdfstring{A Philosophical Technical Note:
\(\mathcal{M}\)-closed
vs.~\(\mathcal{M}\)-open}{A Philosophical Technical Note: \textbackslash mathcal\{M\}-closed vs.~\textbackslash mathcal\{M\}-open}}\label{a-philosophical-technical-note-mathcalm-closed-vs.-mathcalm-open}

There is a philosophical friction in this exercise. By simulating data
from a specific model, we are operating in an \(\mathcal{M}\)-closed
framework---we know for an absolute fact that one of our models is the
true data-generating process. However, we evaluate the recovery using
Bayesian stacking weights, which are designed for the
\(\mathcal{M}\)-open framework (the assumption that all models are
false, and we merely want the optimal predictive mixture).We use
stacking weights here anyway because they are continuous. Rather than
forcing a binary ``win/loss,'' stacking weights tell us how much
predictive variance uniquely belongs to each architecture.

\subsection{Running Model SBC: The Code
Pipeline}\label{running-model-sbc-the-code-pipeline}

The generative process for Model SBC extends parameter SBC to the model
level:

\begin{itemize}
\item
  Sample a true model \(M\) (e.g., a coin flip between Biased and
  Memory).
\item
  Sample parameters strictly from that specific model's prior.
\item
  Simulate data from the data model. Fit all competing models to the
  data and compute their Bayesian stacking weights.
\item
  Repeat this thousands of times.
\end{itemize}

Because running 2,000 simulations and fitting two multilevel Stan models
to each would take hours (or days) on a standard laptop, we will not
execute it live here. However, to maintain the hands-on promise of this
workflow, here is the exact furrr parallelization code you would use.
Notice how it elevates the target to loo::loo\_model\_weights().

\begin{Shaded}
\begin{Highlighting}[]
\FunctionTok{plan}\NormalTok{(multisession, }\AttributeTok{workers =}\NormalTok{ parallel}\SpecialCharTok{::}\FunctionTok{detectCores}\NormalTok{() }\SpecialCharTok{{-}} \DecValTok{1}\NormalTok{)}

\NormalTok{n\_sims }\OtherTok{\textless{}{-}} \DecValTok{2000}

\NormalTok{model\_sbc\_results }\OtherTok{\textless{}{-}} \FunctionTok{future\_map\_dfr}\NormalTok{(}\DecValTok{1}\SpecialCharTok{:}\NormalTok{n\_sims, }\ControlFlowTok{function}\NormalTok{(sim\_id) \{}
  
  \CommentTok{\# 1. Flip a coin to choose the true generative model}
\NormalTok{  true\_model }\OtherTok{\textless{}{-}} \FunctionTok{sample}\NormalTok{(}\FunctionTok{c}\NormalTok{(}\StringTok{"Biased"}\NormalTok{, }\StringTok{"Memory"}\NormalTok{), }\DecValTok{1}\NormalTok{)}
  
  \CommentTok{\# 2. Simulate data exactly from the chosen model\textquotesingle{}s PRIORS}
  \ControlFlowTok{if}\NormalTok{ (true\_model }\SpecialCharTok{==} \StringTok{"Biased"}\NormalTok{) \{}
\NormalTok{    sim\_data }\OtherTok{\textless{}{-}} \FunctionTok{simulate\_biased\_from\_prior}\NormalTok{() }
\NormalTok{    stan\_list }\OtherTok{\textless{}{-}} \FunctionTok{make\_stan\_biased}\NormalTok{(sim\_data)}
\NormalTok{  \} }\ControlFlowTok{else}\NormalTok{ \{}
\NormalTok{    sim\_data }\OtherTok{\textless{}{-}} \FunctionTok{simulate\_memory\_from\_prior}\NormalTok{() }
\NormalTok{    stan\_list }\OtherTok{\textless{}{-}} \FunctionTok{make\_stan\_memory}\NormalTok{(sim\_data)}
\NormalTok{  \}}
  
  \CommentTok{\# 3. Fit BOTH candidate models to the simulated dataset}
\NormalTok{  fit\_b }\OtherTok{\textless{}{-}}\NormalTok{ mod\_biased\_ncp}\SpecialCharTok{$}\FunctionTok{sample}\NormalTok{(}\AttributeTok{data =}\NormalTok{ stan\_list, }\AttributeTok{refresh =} \DecValTok{0}\NormalTok{, }
                                 \AttributeTok{iter\_sampling =} \DecValTok{1000}\NormalTok{, }\AttributeTok{parallel\_chains =} \DecValTok{1}\NormalTok{)}
\NormalTok{  fit\_m }\OtherTok{\textless{}{-}}\NormalTok{ mod\_memory\_ncp}\SpecialCharTok{$}\FunctionTok{sample}\NormalTok{(}\AttributeTok{data =}\NormalTok{ stan\_list, }\AttributeTok{refresh =} \DecValTok{0}\NormalTok{, }
                                 \AttributeTok{iter\_sampling =} \DecValTok{1000}\NormalTok{, }\AttributeTok{parallel\_chains =} \DecValTok{1}\NormalTok{)}
  
  \CommentTok{\# 4. Compute LOO and Bayesian Stacking Weights}
\NormalTok{  loo\_b }\OtherTok{\textless{}{-}}\NormalTok{ fit\_b}\SpecialCharTok{$}\FunctionTok{loo}\NormalTok{()}
\NormalTok{  loo\_m }\OtherTok{\textless{}{-}}\NormalTok{ fit\_m}\SpecialCharTok{$}\FunctionTok{loo}\NormalTok{()}
  
\NormalTok{  weights }\OtherTok{\textless{}{-}}\NormalTok{ loo}\SpecialCharTok{::}\FunctionTok{loo\_model\_weights}\NormalTok{(}\FunctionTok{list}\NormalTok{(}\AttributeTok{Biased =}\NormalTok{ loo\_b, }\AttributeTok{Memory =}\NormalTok{ loo\_m), }
                                    \AttributeTok{method =} \StringTok{"stacking"}\NormalTok{)}
  
  \CommentTok{\# 5. Return the result for this specific simulation iteration}
  \FunctionTok{tibble}\NormalTok{(}
    \AttributeTok{sim\_id =}\NormalTok{ sim\_id,}
    \AttributeTok{true\_model =}\NormalTok{ true\_model,}
    \AttributeTok{weight\_biased =} \FunctionTok{as.numeric}\NormalTok{(weights[}\StringTok{"Biased"}\NormalTok{]),}
    \AttributeTok{weight\_memory =} \FunctionTok{as.numeric}\NormalTok{(weights[}\StringTok{"Memory"}\NormalTok{])}
\NormalTok{  )}
\NormalTok{\}, }\AttributeTok{.options =} \FunctionTok{furrr\_options}\NormalTok{(}\AttributeTok{seed =} \DecValTok{12345}\NormalTok{))}
\end{Highlighting}
\end{Shaded}

\subsection{How Do We Assess Model
Calibration?}\label{how-do-we-assess-model-calibration}

To assess if our model comparison is well-calibrated, we treat the
stacking weights as a measure of predictive confidence. However,
``calibration'' in this context is frequently misunderstood.Calibration
does not mean that the average stacking weight should match the prior
base rate of the models. Calibration means that the weight must honestly
reflect the empirical reality of its own claims:
\[P(\text{Model is True} \mid \text{Weight} = w) = w\]

If we find that across all datasets where the Memory model received a
weight of \(0.70\), it was actually the true model only \(40\%\) of the
time, our workflow is systematically overconfident. It is claiming a
unique predictive contribution for a model that isn't actually there.
Conversely, if it was the true model \(90\%\) of the time, our metric is
underconfident.

\subsubsection{Why do we use a 50/50 prior for
SBC?}\label{why-do-we-use-a-5050-prior-for-sbc}

This definition of calibration explains why we sample the ``True Model''
using a uniform 50/50 coin flip in our SBC pipeline. Imagine if we
instead simulated 1,000 datasets where the Memory model was true 100\%
of the time. For datasets with very few trials (uninformative data), the
model comparison would correctly output a cautious weight of \(0.50\).
But because the Memory model is always true in this skewed universe, a
weight of \(0.50\) would look massively underconfident on a calibration
plot. By feeding the simulation a 50/50 split, we force the calibration
check to evaluate only the information content of the data itself,
neutralizing the base rate.

\subsubsection{Decoding the Calibration
Curves}\label{decoding-the-calibration-curves}

Since we aren't running the full 2,000-fit simulation locally right now,
I have generated a sample plot to show exactly what the calibration
curves produced by the code above would look like under different
diagnostic scenarios.

\begin{Shaded}
\begin{Highlighting}[]
\CommentTok{\# 1. Simulate the data, explicitly including the number of datasets per bin (N)}
\NormalTok{calibration\_data }\OtherTok{\textless{}{-}} \FunctionTok{tibble}\NormalTok{(}
  \AttributeTok{bin\_mid =} \FunctionTok{seq}\NormalTok{(}\FloatTok{0.05}\NormalTok{, }\FloatTok{0.95}\NormalTok{, }\AttributeTok{length.out =} \DecValTok{10}\NormalTok{)}
\NormalTok{) }\SpecialCharTok{|\textgreater{}}
  \FunctionTok{mutate}\NormalTok{(}
    \CommentTok{\# In model comparison, weights often pile up at the extremes (0 or 1) }
    \CommentTok{\# and fewer datasets land in the ambiguous middle.}
    \AttributeTok{N\_datasets   =} \FunctionTok{round}\NormalTok{(}\DecValTok{500} \SpecialCharTok{*}\NormalTok{ (}\DecValTok{4} \SpecialCharTok{*}\NormalTok{ (bin\_mid }\SpecialCharTok{{-}} \FloatTok{0.5}\NormalTok{)}\SpecialCharTok{\^{}}\DecValTok{2} \SpecialCharTok{+} \FloatTok{0.2}\NormalTok{)),}
    
    \AttributeTok{Perfect      =}\NormalTok{ bin\_mid,}
    \AttributeTok{Overconf     =} \DecValTok{1} \SpecialCharTok{/}\NormalTok{ (}\DecValTok{1} \SpecialCharTok{+} \FunctionTok{exp}\NormalTok{(}\SpecialCharTok{{-}}\FloatTok{1.5} \SpecialCharTok{*}\NormalTok{ (}\FunctionTok{qlogis}\NormalTok{(bin\_mid)))),}
    \AttributeTok{Underconf    =} \DecValTok{1} \SpecialCharTok{/}\NormalTok{ (}\DecValTok{1} \SpecialCharTok{+} \FunctionTok{exp}\NormalTok{(}\SpecialCharTok{{-}}\FloatTok{0.5} \SpecialCharTok{*}\NormalTok{ (}\FunctionTok{qlogis}\NormalTok{(bin\_mid)))),}
    \AttributeTok{Miscalibrate =} \FunctionTok{pmin}\NormalTok{(}\DecValTok{1}\NormalTok{, }\FunctionTok{pmax}\NormalTok{(}\DecValTok{0}\NormalTok{, bin\_mid }\SpecialCharTok{+} \FloatTok{0.15} \SpecialCharTok{*} \FunctionTok{sin}\NormalTok{(bin\_mid }\SpecialCharTok{*}\NormalTok{ pi }\SpecialCharTok{*} \DecValTok{2}\NormalTok{)))}
\NormalTok{  ) }\SpecialCharTok{|\textgreater{}}
  \FunctionTok{pivot\_longer}\NormalTok{(}\AttributeTok{cols =} \FunctionTok{c}\NormalTok{(Perfect, Overconf, Underconf, Miscalibrate), }
               \AttributeTok{names\_to =} \StringTok{"Scenario"}\NormalTok{, }\AttributeTok{values\_to =} \StringTok{"Observed"}\NormalTok{) }\SpecialCharTok{|\textgreater{}}
  \FunctionTok{mutate}\NormalTok{(}
    \CommentTok{\# Calculate standard errors for the binomial proportion}
    \AttributeTok{SE =} \FunctionTok{sqrt}\NormalTok{((Observed }\SpecialCharTok{*}\NormalTok{ (}\DecValTok{1} \SpecialCharTok{{-}}\NormalTok{ Observed)) }\SpecialCharTok{/}\NormalTok{ N\_datasets),}
    \CommentTok{\# 95\% Confidence Intervals}
    \AttributeTok{Lower =} \FunctionTok{pmax}\NormalTok{(}\DecValTok{0}\NormalTok{, Observed }\SpecialCharTok{{-}} \FloatTok{1.96} \SpecialCharTok{*}\NormalTok{ SE),}
    \AttributeTok{Upper =} \FunctionTok{pmin}\NormalTok{(}\DecValTok{1}\NormalTok{, Observed }\SpecialCharTok{+} \FloatTok{1.96} \SpecialCharTok{*}\NormalTok{ SE)}
\NormalTok{  )}

\CommentTok{\# 2. The Calibration Curve with Confidence Ribbons}
\NormalTok{p\_curve }\OtherTok{\textless{}{-}} \FunctionTok{ggplot}\NormalTok{(calibration\_data, }\FunctionTok{aes}\NormalTok{(}\AttributeTok{x =}\NormalTok{ bin\_mid, }\AttributeTok{y =}\NormalTok{ Observed, }\AttributeTok{color =}\NormalTok{ Scenario)) }\SpecialCharTok{+}
  \FunctionTok{geom\_abline}\NormalTok{(}\AttributeTok{slope =} \DecValTok{1}\NormalTok{, }\AttributeTok{intercept =} \DecValTok{0}\NormalTok{, }\AttributeTok{linetype =} \StringTok{"dashed"}\NormalTok{, }\AttributeTok{color =} \StringTok{"gray40"}\NormalTok{) }\SpecialCharTok{+}
  \CommentTok{\# Add the variance ribbon!}
  \FunctionTok{geom\_ribbon}\NormalTok{(}\FunctionTok{aes}\NormalTok{(}\AttributeTok{ymin =}\NormalTok{ Lower, }\AttributeTok{ymax =}\NormalTok{ Upper, }\AttributeTok{fill =}\NormalTok{ Scenario), }
              \AttributeTok{alpha =} \FloatTok{0.15}\NormalTok{, }\AttributeTok{color =} \ConstantTok{NA}\NormalTok{) }\SpecialCharTok{+}
  \FunctionTok{geom\_line}\NormalTok{(}\FunctionTok{aes}\NormalTok{(}\AttributeTok{linetype =}\NormalTok{ Scenario), }\AttributeTok{linewidth =} \DecValTok{1}\NormalTok{) }\SpecialCharTok{+}
  \FunctionTok{geom\_point}\NormalTok{(}\AttributeTok{size =} \DecValTok{2}\NormalTok{) }\SpecialCharTok{+}
  \FunctionTok{scale\_color\_manual}\NormalTok{(}\AttributeTok{values =} \FunctionTok{c}\NormalTok{(}\StringTok{"Miscalibrate"} \OtherTok{=} \StringTok{"firebrick"}\NormalTok{, }\StringTok{"Overconf"} \OtherTok{=} \StringTok{"orange"}\NormalTok{, }
                                \StringTok{"Underconf"} \OtherTok{=} \StringTok{"purple"}\NormalTok{, }\StringTok{"Perfect"} \OtherTok{=} \StringTok{"steelblue"}\NormalTok{)) }\SpecialCharTok{+}
  \FunctionTok{scale\_fill\_manual}\NormalTok{(}\AttributeTok{values =} \FunctionTok{c}\NormalTok{(}\StringTok{"Miscalibrate"} \OtherTok{=} \StringTok{"firebrick"}\NormalTok{, }\StringTok{"Overconf"} \OtherTok{=} \StringTok{"orange"}\NormalTok{, }
                               \StringTok{"Underconf"} \OtherTok{=} \StringTok{"purple"}\NormalTok{, }\StringTok{"Perfect"} \OtherTok{=} \StringTok{"steelblue"}\NormalTok{)) }\SpecialCharTok{+}
  \FunctionTok{scale\_x\_continuous}\NormalTok{(}\AttributeTok{labels =}\NormalTok{ scales}\SpecialCharTok{::}\NormalTok{percent, }\AttributeTok{limits =} \FunctionTok{c}\NormalTok{(}\DecValTok{0}\NormalTok{,}\DecValTok{1}\NormalTok{)) }\SpecialCharTok{+}
  \FunctionTok{scale\_y\_continuous}\NormalTok{(}\AttributeTok{labels =}\NormalTok{ scales}\SpecialCharTok{::}\NormalTok{percent, }\AttributeTok{limits =} \FunctionTok{c}\NormalTok{(}\DecValTok{0}\NormalTok{,}\DecValTok{1}\NormalTok{)) }\SpecialCharTok{+}
  \FunctionTok{labs}\NormalTok{(}
    \AttributeTok{title =} \StringTok{"SBC Model Calibration with Uncertainty"}\NormalTok{,}
    \AttributeTok{subtitle =} \StringTok{"Ribbons show 95\% CI. Wider ribbons in the middle reflect fewer ambiguous datasets."}\NormalTok{,}
    \AttributeTok{x =} \ConstantTok{NULL}\NormalTok{, }\AttributeTok{y =} \StringTok{"Empirical Frequency (Ground Truth)"}
\NormalTok{  ) }\SpecialCharTok{+}
  \FunctionTok{theme\_minimal}\NormalTok{() }\SpecialCharTok{+}
  \FunctionTok{theme}\NormalTok{(}\AttributeTok{legend.position =} \StringTok{"none"}\NormalTok{, }\AttributeTok{axis.text.x =} \FunctionTok{element\_blank}\NormalTok{())}

\CommentTok{\# 3. The Marginal Histogram (Showing the distribution of weights)}
\NormalTok{p\_hist }\OtherTok{\textless{}{-}}\NormalTok{ calibration\_data }\SpecialCharTok{|\textgreater{}} 
  \FunctionTok{filter}\NormalTok{(Scenario }\SpecialCharTok{==} \StringTok{"Perfect"}\NormalTok{) }\SpecialCharTok{|\textgreater{}} \CommentTok{\# Just to get the Ns once}
  \FunctionTok{ggplot}\NormalTok{(}\FunctionTok{aes}\NormalTok{(}\AttributeTok{x =}\NormalTok{ bin\_mid, }\AttributeTok{y =}\NormalTok{ N\_datasets)) }\SpecialCharTok{+}
  \FunctionTok{geom\_col}\NormalTok{(}\AttributeTok{fill =} \StringTok{"gray70"}\NormalTok{, }\AttributeTok{alpha =} \FloatTok{0.6}\NormalTok{) }\SpecialCharTok{+}
  \FunctionTok{scale\_x\_continuous}\NormalTok{(}\AttributeTok{labels =}\NormalTok{ scales}\SpecialCharTok{::}\NormalTok{percent, }\AttributeTok{limits =} \FunctionTok{c}\NormalTok{(}\DecValTok{0}\NormalTok{,}\DecValTok{1}\NormalTok{)) }\SpecialCharTok{+}
  \FunctionTok{labs}\NormalTok{(}\AttributeTok{x =} \StringTok{"Stacking Weight (Model\textquotesingle{}s Confidence)"}\NormalTok{, }
       \AttributeTok{y =} \StringTok{"Count (N)"}\NormalTok{) }\SpecialCharTok{+}
  \FunctionTok{theme\_minimal}\NormalTok{()}

\CommentTok{\# Combine them using patchwork}
\NormalTok{p\_curve }\SpecialCharTok{/}\NormalTok{ p\_hist }\SpecialCharTok{+} \FunctionTok{plot\_layout}\NormalTok{(}\AttributeTok{heights =} \FunctionTok{c}\NormalTok{(}\DecValTok{4}\NormalTok{, }\DecValTok{1}\NormalTok{))}
\end{Highlighting}
\end{Shaded}

\begin{center}
\includegraphics[width=0.8\linewidth,height=\textheight,keepaspectratio]{07-ModelComparison_files/figure-pdf/ch08_calibration_curves_demo-1.pdf}
\end{center}

When you plot the binned stacking weights (X-axis) against the actual
frequency that the model was the true data-generating process (Y-axis),
the geometry of the curve diagnoses your workflow:

\begin{itemize}
\item
  Perfectly Calibrated (Steel Blue Diagonal): The model's predictive
  confidence perfectly matches reality. If the data is uninformative,
  the weights go to \(0.5\) and the true model is a coin flip. If the
  data is rich, the weights go to \(1.0\) and the model is always right.
\item
  Overconfident (Orange S-Curve): The model is ``loud.'' Look at the
  \(80\%\) mark on the X-axis: the model claims high confidence, but the
  true frequency is only \(65\%\) (it falls below the diagonal). The
  model is overestimating its own predictive uniqueness. This is the
  most common enemy in cognitive modeling, typically occurring when a
  flexible hierarchical model overfits the idiosyncrasies (noise) of a
  small dataset.
\item
  Underconfident (Purple Reverse S-Curve): The model is ``shy.'' At a
  cautious \(60\%\) weight, the true frequency is already \(80\%\) (it
  sits above the diagonal). The model is failing to commit. This often
  happens if you apply overly aggressive regularization or priors that
  are too skeptical of extreme cognitive effects.
\item
  Miscalibrated (Red Wavy Line): This represents a structural breakdown
  in the model comparison. Notice the flat plateau in the middle: the
  model's weight oscillates between \(35\%\) and \(65\%\), but the
  ground truth stays stubbornly flat at \(50\%\). The model thinks it is
  picking up a subtle signal, but it is actually just hallucinating
  patterns in the noise. Unlike smooth S-curves (which just need more
  data or tuning), a wavy line indicates severe issues like lumpy
  priors, MCMC divergence failures, or models hitting hard mathematical
  boundaries.
\end{itemize}

\subsubsection{The Danger of Averages: Embracing the
Distribution}\label{the-danger-of-averages-embracing-the-distribution}

Looking at a simple line on a calibration curve can be deceiving. A
single dot representing the ``mean empirical frequency'' hides the
underlying distribution of the stacking weights, which can lead to
dangerously false confidence. A rigorous calibration check must
visualize two distinct types of variance:

\begin{enumerate}
\def\labelenumi{\arabic{enumi}.}
\item
  Vertical Variance (Binomial Uncertainty) Not all bins are created
  equal. Suppose our model assigns a stacking weight of \(0.50\). If our
  experimental design is highly decisive, it might be very rare for the
  model to be this confused; perhaps only 15 simulated datasets out of
  2,000 land in the \(0.50\) bin. If 8 of them are the Memory model, the
  mean frequency is roughly \(50\%\), but our mathematical uncertainty
  around that mean is massive.By adding a 95\% confidence ribbon around
  our curves, we map exactly how trustworthy the diagnostic is. Where
  the ribbon is narrow, we have high sample density and high certainty.
  Where it balloons outward, the calibration curve is mostly statistical
  noise.
\item
  Horizontal Variance (The Weight Density) We also need to know the base
  distribution of the weights themselves. Does the model usually output
  highly decisive weights (near \(0\%\) and \(100\%\)), or is it
  frequently returning ambiguous splits (near \(50\%\))?
\end{enumerate}

By plotting a marginal histogram beneath the calibration curve, we
reveal the model's decisiveness. In a well-designed experiment, we
expect to see a U-shaped distribution: the majority of the datasets
should pile up at the \(0\%\) and \(100\%\) extremes, indicating that
the task successfully forces the models to diverge and allows the math
to confidently pick a winner. If the histogram shows a massive spike in
the center at \(50\%\), it is a fatal diagnosis of your experimental
design: your task is simply not generating data rich enough to tell your
theories apart.

\begin{quote}
Technical Note: The Flaw of Binning \& Counting (and the CORP Solution)
If you attempt to generate the calibration curves we just discussed
using your own simulated Model SBC data, you will immediately run into a
severe methodological trap: the ``binning and counting'' problem. To
plot empirical frequencies, you have to group your continuous stacking
weights into discrete bins (e.g., \(0.40\) to \(0.50\)). If you use
fixed, evenly spaced bins, you will inevitably end up with ``sparse
bins''---regions where the model rarely outputs a weight. The empirical
frequency inside a sparse bin is subject to massive estimation variance.
If you plot this, your calibration line will spike wildly up and down,
making a perfectly calibrated model look like it has suffered a
structural breakdown (the red ``Miscalibrated'' wavy line from our
conceptual plot). To solve this, Dimitriadis, Gneiting, and Jordan
(2021) introduced the CORP approach (Consistent, Optimally binned,
Reproducible, and PAV-based). Instead of forcing data into arbitrary,
fixed bins, CORP lets the data determine its own optimal bins using
nonparametric isotonic regression.It accomplishes this via the elegant
Pool-Adjacent-Violators (PAV) algorithm. Here is how it works for Model
SBC:
\end{quote}

\begin{quote}
\begin{itemize}
\tightlist
\item
  Sort: Sort all 2,000 simulated datasets from the lowest assigned
  stacking weight to the highest.
\item
  Assume Monotonicity: A logical model comparison metric must be
  monotonic: a higher stacking weight should correspond to an equal or
  higher true empirical frequency.
\item
  Pool the Violators: The algorithm walks along the sorted data. If it
  finds a ``violator''---a case where a higher stacking weight actually
  has a lower true empirical rate than the point before it (due to noise
  in sparse regions)---it ``pools'' those adjacent points together and
  averages them into a single, wider bin. The result is pure magic: PAV
  automatically creates wide, stable bins where your data is sparse, and
  fine-grained, precise steps where your data is dense. It also
  mathematically decomposes the Brier score of your predictions into
  three stable metrics: Miscalibration (MCB), Discrimination (DSC), and
  Uncertainty (UNC).Here is how you can apply the CORP approach to your
  own Model SBC data using the authors' reliabilitydiag R package:
\end{itemize}
\end{quote}

\begin{Shaded}
\begin{Highlighting}[]
\NormalTok{pacman}\SpecialCharTok{::}\FunctionTok{p\_load}\NormalTok{(reliabilitydiag)}

\CommentTok{\# 1. Simulate a mock Model SBC output (2,000 fits)}
\NormalTok{n\_sims }\OtherTok{\textless{}{-}} \DecValTok{2000}

\CommentTok{\# Simulate stacking weights (piling up at the extremes)}
\NormalTok{mock\_weights }\OtherTok{\textless{}{-}} \FunctionTok{rbeta}\NormalTok{(n\_sims, }\FloatTok{0.5}\NormalTok{, }\FloatTok{0.5}\NormalTok{) }

\CommentTok{\# Simulate the Ground Truth (1 = Memory is True, 0 = Biased is True)}
\CommentTok{\# We deliberately inject "Overconfidence" into the DGP: }
\CommentTok{\# The true probability is pulled closer to 0.5 than the model\textquotesingle{}s weight claims.}
\NormalTok{true\_probs }\OtherTok{\textless{}{-}} \FunctionTok{plogis}\NormalTok{(}\FunctionTok{qlogis}\NormalTok{(mock\_weights) }\SpecialCharTok{*} \FloatTok{0.6}\NormalTok{) }
\NormalTok{mock\_truth }\OtherTok{\textless{}{-}} \FunctionTok{rbinom}\NormalTok{(n\_sims, }\DecValTok{1}\NormalTok{, true\_probs)}

\CommentTok{\# 2. Apply the CORP approach (PAV Algorithm)}
\CommentTok{\# \textquotesingle{}x\textquotesingle{} is our continuous prediction (stacking weight)}
\CommentTok{\# \textquotesingle{}y\textquotesingle{} is the binary ground truth (was it actually the Memory model?)}
\NormalTok{rd }\OtherTok{\textless{}{-}} \FunctionTok{reliabilitydiag}\NormalTok{(}
  \AttributeTok{Overconfident\_Model =}\NormalTok{ mock\_weights, }
  \AttributeTok{y =}\NormalTok{ mock\_truth}
\NormalTok{)}

\CommentTok{\# 3. Plot the optimal reliability diagram}
\FunctionTok{autoplot}\NormalTok{(rd) }\SpecialCharTok{+}
  \FunctionTok{labs}\NormalTok{(}
    \AttributeTok{title =} \StringTok{"SBC Model Calibration (CORP Approach)"}\NormalTok{,}
    \AttributeTok{subtitle =} \StringTok{"Optimal dynamic binning via the PAV algorithm. Notice the S{-}curve of overconfidence."}\NormalTok{,}
    \AttributeTok{x =} \StringTok{"Forecast Value (Bayesian Stacking Weight)"}\NormalTok{,}
    \AttributeTok{y =} \StringTok{"Conditional Event Probability (True Model Frequency)"}
\NormalTok{  ) }\SpecialCharTok{+}
  \FunctionTok{theme\_classic}\NormalTok{()}
\end{Highlighting}
\end{Shaded}

\begin{center}
\includegraphics[width=0.8\linewidth,height=\textheight,keepaspectratio]{07-ModelComparison_files/figure-pdf/ch08_corp_pav_demo-1.pdf}
\end{center}

\subsubsection{When a 50/50 Split is a Perfect
Diagnosis}\label{when-a-5050-split-is-a-perfect-diagnosis}

Imagine the SBC process selects the Memory model, and then draws a
population memory sensitivity parameter from the prior that happens to
be very close to zero (\(\mu_\beta \approx 0\)). The simulated agents
will completely ignore the opponent's history. Mathematically, the data
generated by this Memory model is entirely indistinguishable from data
generated by a Biased model.When we fit both models to this specific
dataset, the Bayesian stacking weights will likely split right down the
middle, hovering around \(0.5\) for each model.Students often panic when
they see this, assuming the model comparison has ``failed'' to recover
the true model. In reality, a weight of \(0.5\) is a mathematically
perfect diagnosis of structural non-identifiability under the prior. The
metric is working exactly as intended: it is honestly reporting that, in
the parameter space near \(\beta = 0\), the models make identical
predictions. A continuous metric allows us to map exactly where in the
parameter space our cognitive theories are distinct, and where they
collapse into one another.

\section{Prior Sensitivity at the Comparison
Level}\label{prior-sensitivity-at-the-comparison-level}

We applied prior sensitivity analysis to individual models in Chapter 6.
At the model comparison stage, the question becomes slightly different:
even if each model's individual parameters are sensitive to prior
choices, does the \emph{comparison} between models remain stable?

Consider the following thought experiment: scaling the prior on
\(\sigma_\theta\) shifts the biased model's ELPD by 5 nats, but also
shifts the memory model's ELPD by roughly the same 5 nats. The
difference --- the quantity that drives the comparison --- would remain
stable. This would be a good outcome: the comparison is robust even
though individual models are sensitive. The problematic scenario is
asymmetric sensitivity: one model's ELPD shifts substantially while the
other's stays put. In that case, our preference for one model over the
other depends on which prior happened to be chosen --- which is not a
principled basis for a scientific conclusion.

\begin{Shaded}
\begin{Highlighting}[]
\FunctionTok{cat}\NormalTok{(}\StringTok{"=== Prior sensitivity: Biased model ===}\SpecialCharTok{\textbackslash{}n}\StringTok{"}\NormalTok{)}
\end{Highlighting}
\end{Shaded}

\begin{verbatim}
=== Prior sensitivity: Biased model ===
\end{verbatim}

\begin{Shaded}
\begin{Highlighting}[]
\FunctionTok{print}\NormalTok{(priorsense}\SpecialCharTok{::}\FunctionTok{powerscale\_sensitivity}\NormalTok{(fit\_b\_on\_b,}
  \AttributeTok{variable=}\FunctionTok{c}\NormalTok{(}\StringTok{"mu\_theta"}\NormalTok{,}\StringTok{"sigma\_theta"}\NormalTok{)))}
\end{Highlighting}
\end{Shaded}

\begin{verbatim}
Sensitivity based on cjs_dist
Prior selection: all priors
Likelihood selection: all data

    variable prior likelihood                     diagnosis
    mu_theta 0.025      0.035                             -
 sigma_theta 0.681      0.127 potential prior-data conflict
\end{verbatim}

\begin{Shaded}
\begin{Highlighting}[]
\FunctionTok{cat}\NormalTok{(}\StringTok{"}\SpecialCharTok{\textbackslash{}n}\StringTok{=== Prior sensitivity: Memory model ===}\SpecialCharTok{\textbackslash{}n}\StringTok{"}\NormalTok{)}
\end{Highlighting}
\end{Shaded}

\begin{verbatim}

=== Prior sensitivity: Memory model ===
\end{verbatim}

\begin{Shaded}
\begin{Highlighting}[]
\FunctionTok{print}\NormalTok{(priorsense}\SpecialCharTok{::}\FunctionTok{powerscale\_sensitivity}\NormalTok{(fit\_m\_on\_m,}
  \AttributeTok{variable=}\FunctionTok{c}\NormalTok{(}\StringTok{"mu\_alpha"}\NormalTok{,}\StringTok{"mu\_beta"}\NormalTok{)))}
\end{Highlighting}
\end{Shaded}

\begin{verbatim}
Sensitivity based on cjs_dist
Prior selection: all priors
Likelihood selection: all data

 variable prior likelihood diagnosis
 mu_alpha 0.041      0.064         -
  mu_beta 0.027      0.091         -
\end{verbatim}

\begin{Shaded}
\begin{Highlighting}[]
\NormalTok{p1 }\OtherTok{\textless{}{-}}\NormalTok{ priorsense}\SpecialCharTok{::}\FunctionTok{powerscale\_plot\_dens}\NormalTok{(}
\NormalTok{  priorsense}\SpecialCharTok{::}\FunctionTok{powerscale\_sequence}\NormalTok{(fit\_b\_on\_b, }\AttributeTok{variable =} \FunctionTok{c}\NormalTok{(}\StringTok{"mu\_theta"}\NormalTok{, }\StringTok{"sigma\_theta"}\NormalTok{))}
\NormalTok{) }\SpecialCharTok{+}
  \FunctionTok{theme\_classic}\NormalTok{() }\SpecialCharTok{+} 
  \FunctionTok{labs}\NormalTok{(}\AttributeTok{title =} \StringTok{"Biased Model: Prior Sensitivity"}\NormalTok{)}

\NormalTok{p2 }\OtherTok{\textless{}{-}}\NormalTok{ priorsense}\SpecialCharTok{::}\FunctionTok{powerscale\_plot\_dens}\NormalTok{(}
\NormalTok{  priorsense}\SpecialCharTok{::}\FunctionTok{powerscale\_sequence}\NormalTok{(fit\_m\_on\_m, }\AttributeTok{variable =} \FunctionTok{c}\NormalTok{(}\StringTok{"mu\_alpha"}\NormalTok{, }\StringTok{"mu\_beta"}\NormalTok{))}
\NormalTok{) }\SpecialCharTok{+}
  \FunctionTok{theme\_classic}\NormalTok{() }\SpecialCharTok{+} 
  \FunctionTok{labs}\NormalTok{(}\AttributeTok{title =} \StringTok{"Memory Model: Prior Sensitivity"}\NormalTok{)}

\CommentTok{\# Combine using patchwork}
\NormalTok{p1 }\SpecialCharTok{/}\NormalTok{ p2}
\end{Highlighting}
\end{Shaded}

\begin{center}
\includegraphics[width=0.8\linewidth,height=\textheight,keepaspectratio]{07-ModelComparison_files/figure-pdf/07-prior-sensitivity-1.pdf}
\end{center}

The console output above directly tests this: the priorsense tables
include a diagnostic column that warns you if the LOO/ELPD computation
is highly sensitive to the prior. Meanwhile, the density perturbation
plots show why that sensitivity might be happening by revealing how the
posterior on each parameter shifts as we increase or decrease the prior
weight. If the biased model's ELPD is flagged as sensitive but the
memory model's is not, then our model preference is partly an artifact
of our prior choice rather than genuine predictive superiority.

\section{When Standard LOO Is Wrong: Sequential
Dependencies}\label{when-standard-loo-is-wrong-sequential-dependencies}

There is an important limitation to PSIS-LOO that is easy to overlook.
Standard LOO treats observations as \emph{exchangeable} --- it assumes
that the order of observations does not matter, and that leaving out
observation \(i\) only affects predictions for \(y_i\) itself, not for
any other observation.

For the biased agent model this is entirely correct: each trial is
independent given the agent's parameters. For our memory model as
implemented, it is also technically valid: the memory signal
\texttt{opp\_rate\_prev} is a pre-computed covariate that we pass into
Stan as fixed data. Leaving out trial \(i\) from the LOO computation
does not change \texttt{opp\_rate\_prev} for any other trial --- the
memory features were already computed before Stan ever ran.

However, many cognitively interesting models have genuine temporal
dependencies that standard LOO cannot handle. Consider a Rescorla-Wagner
reinforcement learning model where the learning state at trial \(t\)
depends on the prediction error at trial \(t-1\). Or a drift-diffusion
model with trial-by-trial fatigue effects. In these models, leaving out
trial \(i\) from the LOO set would also change the learning state at
trials \(i+1, i+2, \ldots\) --- we cannot simply remove one observation
without unraveling the whole sequential structure. Standard LOO would
then be ``peeking into the future'' by conditioning on observations that
come after the one being predicted.

For such models, the correct approach is \textbf{Leave-Future-Out
cross-validation (LFO-CV)} (Burkner, Gabry \& Vehtari, 2020). Instead of
leaving out one observation at a time in isolation, LFO-CV asks: at each
time point \(t\), how well does the model predict the next trial given
all past trials \(y_{1:t-1}\)? This respects the temporal structure and
gives an honest measure of prospective predictive accuracy.

Decision rule: if temporal structure is handled externally via
pre-computed features passed as data (as in our memory model), standard
LOO is appropriate. If the sequential dynamics are inside Stan's
\texttt{transformed\ data} or \texttt{model} blocks, use LFO-CV.

\section{What Model Comparison Cannot
Do}\label{what-model-comparison-cannot-do}

We have spent this chapter learning how to use PSIS-LOO to compare
models. It is equally important to understand what it \emph{cannot} do.

\textbf{It measures internal training generalization, not external
transfer.} We need to think carefully about what we mean by ``out of
sample.'' There are actually two distinct types of test sets we might
consider: internal and external.

Internal test sets come from the same data collection effort as our
training data - for example, we might randomly set aside 20\% of our
matching pennies games to test on. While this approach helps us detect
overfitting to specific participants or trials, it cannot tell us how
well our model generalizes to truly new contexts. Our test set
participants were recruited from the same population, played the game
under the same conditions, and were influenced by the same experimental
setup as our training participants.

External test sets, in contrast, come from genuinely different contexts.
For our matching pennies model, this might mean testing on games played:

\begin{itemize}
\tightlist
\item
  In different cultures or age groups
\item
  Under time pressure versus relaxed conditions
\item
  For real money versus just for fun
\item
  Against human opponents versus computer agents
\item
  In laboratory versus online settings
\end{itemize}

The distinction matters because cognitive models often capture not just
universal mental processes, but also specific strategies that people
adopt in particular contexts. A model that perfectly predicts behavior
in laboratory matching pennies games might fail entirely when applied to
high-stakes poker games, even though both involve similar strategic
thinking.

This raises deeper questions about what kind of generalization we want
our models to achieve. Are we trying to build models that capture
universal cognitive processes, or are we content with models that work
well within specific contexts? The answer affects not just how we
evaluate our models, but how we design them in the first place.

In practice, truly external test sets are rare in cognitive science -
they require additional data collection under different conditions,
which is often impractical. This means we must be humble about our
claims of generalization. When we talk about a model's predictive
accuracy, we should be clear that we're usually measuring its ability to
generalize within a specific experimental context, not its ability to
capture human cognition in all its diversity.

This limitation of internal test sets is one reason why cognitive
scientists often complement predictive accuracy metrics with other forms
of model evaluation, such as testing theoretical predictions on new
tasks or examining whether model parameters correlate sensibly with
individual differences. These approaches help us build confidence that
our models capture meaningful cognitive processes rather than just
statistical patterns specific to our experimental setup.

\textbf{A better predictive score does not mean a better mechanistic
explanation.} Two models can achieve identical ELPD while making
diametrically opposite claims about the cognitive mechanism generating
the data. If the biased model and the memory model achieve similar
predictive performance, that does not mean the two theories are
equivalent --- it means the data at hand cannot distinguish them. The
correct scientific response is to design experiments targeting the
mechanistic difference: conditions where the two models predict
qualitatively different behavior.

\textbf{The ``wrong model wins'' scenario is real.} When memory agents
are simulated with high \(\sigma_\beta\) (many agents with near-zero
memory sensitivity), the population of memory agents becomes nearly
indistinguishable from a population of biased agents. PSIS-LOO may then
prefer the simpler biased model --- not because memory processes are
absent, but because they are weak, variable, and the simpler model
achieves comparable predictions on average. This is actually a
\emph{correct} finding: it tells you that the empirical conditions were
not sufficiently discriminative, not that the method is broken. The
productive response is to ask what experimental design would make the
two populations more discriminable --- perhaps a richer or more
structured opponent, or a longer game.

\section{Conclusion: Model Comparison as Scientific
Reasoning}\label{conclusion-model-comparison-as-scientific-reasoning}

This chapter developed a complete pipeline for comparing validated
cognitive models:

\begin{enumerate}
\def\labelenumi{\arabic{enumi}.}
\tightlist
\item
  \textbf{PSIS-LOO} approximates leave-one-out cross-validation from a
  single fit, with Pareto-\(\hat{k}\) flagging both numerical
  unreliability and cognitively interesting observations.
\item
  \textbf{Agent-level K-fold CV} provides an honest test of
  generalization to new participants --- answering the question PSIS-LOO
  cannot.
\item
  \textbf{LOO-PIT + POT-C} checks whether the model's predictive
  \emph{uncertainty} is calibrated, not just its mean predictions.
\item
  \textbf{\texttt{loo\_compare()}} produces a Z-score-interpretable
  ranking; below 2 means inconclusive, not that the method failed.
\item
  \textbf{Bayesian stacking} finds the optimal predictive mixture and
  provides a continuous measure of model evidence.
\item
  \textbf{Model SBC} tests the full comparison pipeline's calibration
  across the prior space.
\end{enumerate}

Two design lessons cut across all of these tools: - The \textbf{reversal
design} (validated in Ch. 6, extended in Ch. 7) was essential --- a
constant-bias opponent makes the biased and memory models statistically
inseparable, and no amount of data fixes a confounded design. -
\textbf{All models must pass their individual quality checks before
comparison}. Comparing a well-fitted model against one with divergences
is not a fair scientific test.

Chapter 9 introduces a new model type --- mixture models --- that
extends the toolkit beyond single-process accounts and applies the same
PSIS-LOO machinery to a qualitatively different identifiability problem.

\section{References}\label{references}

Burkner, P.-C., Gabry, J., \& Vehtari, A. (2020). Approximate
leave-future-out cross-validation for Bayesian time series models.
\emph{Journal of Statistical Computation and Simulation}, 90(14),
2499-2523.

Czado, C., Gneiting, T., \& Held, L. (2009). Predictive model assessment
for count data. \emph{Biometrical Journal}, 51(5), 720-731.

Gabry, J., Simpson, D., Vehtari, A., Betancourt, M., \& Gelman, A.
(2019). Visualization in Bayesian workflow. \emph{Journal of the Royal
Statistical Society: Series A}, 182(2), 389-402.

Liu, Y., \& Xie, J. (2020). Cauchy combination test: a powerful test
with analytic p-value calculation under arbitrary dependency structures.
\emph{Journal of the American Statistical Association}, 115(529),
393-402.

Paananen, T., Piironen, J., Burkner, P.-C., \& Vehtari, A. (2021).
Implicitly adaptive importance sampling. \emph{Statistics and
Computing}, 31, 16.

Sailynoja, T., Burkner, P.-C., \& Vehtari, A. (2022). Graphical test for
discrete uniformity and its applications in goodness-of-fit evaluation
and multiple sample comparison. \emph{Statistics and Computing}, 32(2),
32.

Tesso, H., \& Vehtari, A. (2026). LOO-PIT predictive model checking.
\emph{arXiv:2603.02928}.

Vehtari, A., Gelman, A., \& Gabry, J. (2017). Practical Bayesian model
evaluation using leave-one-out cross-validation and WAIC.
\emph{Statistics and Computing}, 27(5), 1413-1432.

Yao, Y., Vehtari, A., Simpson, D., \& Gelman, A. (2018). Using stacking
to average Bayesian predictive distributions. \emph{Bayesian Analysis},
13(3), 917-1007.

\bookmarksetup{startatroot}

\chapter{Mixture Models}\label{sec-mixture-models}

\begin{quote}
\textbf{📍 Where we are in the Bayesian modeling workflow:} Chs. 1--7
modelled behaviour as arising from a \emph{single} cognitive process per
agent. This chapter relaxes that assumption: behaviour can be a weighted
mixture of distinct processes (e.g., deliberate strategy and attentional
lapses). Mixture models extend the toolkit from Ch. 7 (hierarchical
structures) and Ch. 8 (model comparison via LOO) to this new model
class. The key new challenge is that discrete latent variables --- which
process generated this trial? --- cannot be sampled by HMC and must be
analytically integrated out before any Stan code is written. The full
course roadmap is in Ch. 1, Section ``The Course Roadmap''.
\end{quote}

When we model human cognition and behavior, people do not always follow
a single, consistent strategy. A participant in a matching-pennies game
might sometimes carefully exploit a detected pattern in the opponent's
behavior, and other times respond essentially at random when attention
lapses. Traditional single-process models are structurally misspecified
when this is the case: they attribute lapse-driven variance to the
strategic parameter, inflating its uncertainty and biasing its location.

Mixture models provide an elegant and principled solution. Rather than
assuming behavior reflects one cognitive process, they represent
observations as draws from a \textbf{weighted combination} of component
distributions, each corresponding to a distinct process. This reframes
the question addressed in Chapter 8 --- \emph{which model is correct?}
--- into a richer question: \emph{how much does each process contribute,
and for which observations?}

Because each component must be marginalized analytically before writing
Stan code, mixture models also serve as a concrete introduction to a
class of computational constraint that applies broadly: \textbf{discrete
parameters cannot be sampled by HMC and must be integrated out}.
Mastering this principle is prerequisite for hidden Markov models,
contamination models, and finite mixture clusters.

\section{Learning Objectives}\label{learning-objectives-7}

After completing this chapter you will be able to:

\begin{enumerate}
\def\labelenumi{\arabic{enumi}.}
\tightlist
\item
  Specify a mixture model as a marginalized sum of component
  log-likelihoods and articulate \emph{why} marginalization is required
  in Stan.
\item
  Identify and correct the vectorisation error that makes the mixture
  likelihood computationally wrong when component densities are applied
  to arrays.
\item
  Validate a mixture model with SBC, extracting both calibration
  diagnostics and parameter recovery from the same simulation bank.
\item
  Characterize the likelihood-level identifiability constraint that
  limits what binary choice data can tell us about mixture parameters,
  and visualise the posterior aliasing manifold.
\item
  Extend the single-agent model to a hierarchical mixture model with a
  non-centred parameterisation.
\end{enumerate}

\begin{center}\rule{0.5\linewidth}{0.5pt}\end{center}

\section{Motivation: When are Mixtures
Visible?}\label{motivation-when-are-mixtures-visible}

Before developing the mathematics, consider a concrete case where a
mixture is visually transparent. In many response-time tasks, a
participant sometimes deliberates carefully (producing a log-normal RT
distribution) and sometimes responds inattentively (producing a roughly
uniform RT distribution). The resulting marginal distribution is a
weighted blend.

\begin{Shaded}
\begin{Highlighting}[]
\FunctionTok{set.seed}\NormalTok{(}\DecValTok{8}\NormalTok{)}
\NormalTok{n\_obs     }\OtherTok{\textless{}{-}} \DecValTok{500}
\NormalTok{pi\_delib  }\OtherTok{\textless{}{-}} \FloatTok{0.80}       \CommentTok{\# 80 \% deliberate}
\NormalTok{mu\_log    }\OtherTok{\textless{}{-}} \FloatTok{6.2}        \CommentTok{\# mean of log(RT), deliberate process ≈ 500 ms}
\NormalTok{sigma\_log }\OtherTok{\textless{}{-}} \FloatTok{0.20}
\NormalTok{min\_lapse }\OtherTok{\textless{}{-}} \DecValTok{200}
\NormalTok{max\_lapse }\OtherTok{\textless{}{-}} \DecValTok{2000}

\NormalTok{component }\OtherTok{\textless{}{-}} \FunctionTok{rbinom}\NormalTok{(n\_obs, }\DecValTok{1}\NormalTok{, pi\_delib)  }\CommentTok{\# 1 = deliberate, 0 = lapse}

\NormalTok{rt\_data }\OtherTok{\textless{}{-}} \FunctionTok{tibble}\NormalTok{(}
  \AttributeTok{rt        =} \FunctionTok{ifelse}\NormalTok{(component }\SpecialCharTok{==} \DecValTok{1}\NormalTok{,}
                     \FunctionTok{rlnorm}\NormalTok{(n\_obs, mu\_log, sigma\_log),}
                     \FunctionTok{runif}\NormalTok{(n\_obs, min\_lapse, max\_lapse)),}
  \AttributeTok{log\_rt    =} \FunctionTok{log}\NormalTok{(rt),}
  \AttributeTok{process   =} \FunctionTok{ifelse}\NormalTok{(component }\SpecialCharTok{==} \DecValTok{1}\NormalTok{, }\StringTok{"Deliberate"}\NormalTok{, }\StringTok{"Lapse"}\NormalTok{)}
\NormalTok{)}

\CommentTok{\# RT scale}
\NormalTok{p1 }\OtherTok{\textless{}{-}} \FunctionTok{ggplot}\NormalTok{(rt\_data, }\FunctionTok{aes}\NormalTok{(}\AttributeTok{x =}\NormalTok{ rt)) }\SpecialCharTok{+}
  \FunctionTok{geom\_histogram}\NormalTok{(}\FunctionTok{aes}\NormalTok{(}\AttributeTok{y =} \FunctionTok{after\_stat}\NormalTok{(density)),}
                 \AttributeTok{bins =} \DecValTok{50}\NormalTok{, }\AttributeTok{fill =} \StringTok{"grey80"}\NormalTok{, }\AttributeTok{colour =} \StringTok{"white"}\NormalTok{) }\SpecialCharTok{+}
  \FunctionTok{geom\_density}\NormalTok{(}\FunctionTok{aes}\NormalTok{(}\AttributeTok{colour =}\NormalTok{ process, }\AttributeTok{fill =}\NormalTok{ process),}
               \AttributeTok{data =}\NormalTok{ rt\_data }\SpecialCharTok{|\textgreater{}} \FunctionTok{filter}\NormalTok{(process }\SpecialCharTok{==} \StringTok{"Deliberate"}\NormalTok{),}
               \AttributeTok{alpha =} \FloatTok{0.25}\NormalTok{, }\AttributeTok{linewidth =} \FloatTok{0.9}\NormalTok{) }\SpecialCharTok{+}
  \FunctionTok{geom\_segment}\NormalTok{(}
    \FunctionTok{aes}\NormalTok{(}\AttributeTok{x =}\NormalTok{ min\_lapse, }\AttributeTok{xend =}\NormalTok{ max\_lapse,}
        \AttributeTok{y =} \DecValTok{1} \SpecialCharTok{/}\NormalTok{ (max\_lapse }\SpecialCharTok{{-}}\NormalTok{ min\_lapse),}
        \AttributeTok{yend =} \DecValTok{1} \SpecialCharTok{/}\NormalTok{ (max\_lapse }\SpecialCharTok{{-}}\NormalTok{ min\_lapse),}
        \AttributeTok{colour =} \StringTok{"Lapse"}\NormalTok{),}
    \AttributeTok{linewidth =} \FloatTok{1.2}
\NormalTok{  ) }\SpecialCharTok{+}
  \FunctionTok{scale\_colour\_manual}\NormalTok{(}\AttributeTok{values =} \FunctionTok{c}\NormalTok{(}\AttributeTok{Deliberate =} \StringTok{"\#2166AC"}\NormalTok{, }\AttributeTok{Lapse =} \StringTok{"\#D6604D"}\NormalTok{),}
                      \AttributeTok{name =} \ConstantTok{NULL}\NormalTok{) }\SpecialCharTok{+}
  \FunctionTok{scale\_fill\_manual}\NormalTok{(}\AttributeTok{values  =} \FunctionTok{c}\NormalTok{(}\AttributeTok{Deliberate =} \StringTok{"\#2166AC"}\NormalTok{, }\AttributeTok{Lapse =} \StringTok{"\#D6604D"}\NormalTok{),}
                    \AttributeTok{name =} \ConstantTok{NULL}\NormalTok{, }\AttributeTok{guide =} \StringTok{"none"}\NormalTok{) }\SpecialCharTok{+}
  \FunctionTok{coord\_cartesian}\NormalTok{(}\AttributeTok{xlim =} \FunctionTok{c}\NormalTok{(}\DecValTok{0}\NormalTok{, }\DecValTok{2000}\NormalTok{)) }\SpecialCharTok{+}
  \FunctionTok{labs}\NormalTok{(}\AttributeTok{title    =} \StringTok{"Reaction{-}time distribution"}\NormalTok{,}
       \AttributeTok{subtitle =} \StringTok{"A mixture of deliberate (log{-}normal) and lapse (uniform) processes"}\NormalTok{,}
       \AttributeTok{x =} \StringTok{"RT (ms)"}\NormalTok{, }\AttributeTok{y =} \StringTok{"Density"}\NormalTok{) }\SpecialCharTok{+}
  \FunctionTok{theme}\NormalTok{(}\AttributeTok{legend.position =} \StringTok{"top"}\NormalTok{)}

\CommentTok{\# Log{-}RT scale}
\NormalTok{p2 }\OtherTok{\textless{}{-}} \FunctionTok{ggplot}\NormalTok{(rt\_data, }\FunctionTok{aes}\NormalTok{(}\AttributeTok{x =}\NormalTok{ log\_rt)) }\SpecialCharTok{+}
  \FunctionTok{geom\_histogram}\NormalTok{(}\FunctionTok{aes}\NormalTok{(}\AttributeTok{y =} \FunctionTok{after\_stat}\NormalTok{(density)),}
                 \AttributeTok{bins =} \DecValTok{50}\NormalTok{, }\AttributeTok{fill =} \StringTok{"grey80"}\NormalTok{, }\AttributeTok{colour =} \StringTok{"white"}\NormalTok{) }\SpecialCharTok{+}
  \FunctionTok{geom\_density}\NormalTok{(}\FunctionTok{aes}\NormalTok{(}\AttributeTok{colour =}\NormalTok{ process, }\AttributeTok{fill =}\NormalTok{ process),}
               \AttributeTok{data =}\NormalTok{ rt\_data }\SpecialCharTok{|\textgreater{}} \FunctionTok{filter}\NormalTok{(process }\SpecialCharTok{==} \StringTok{"Deliberate"}\NormalTok{),}
               \AttributeTok{alpha =} \FloatTok{0.25}\NormalTok{, }\AttributeTok{linewidth =} \FloatTok{0.9}\NormalTok{) }\SpecialCharTok{+}
  \FunctionTok{scale\_colour\_manual}\NormalTok{(}\AttributeTok{values =} \FunctionTok{c}\NormalTok{(}\AttributeTok{Deliberate =} \StringTok{"\#2166AC"}\NormalTok{), }\AttributeTok{name =} \ConstantTok{NULL}\NormalTok{) }\SpecialCharTok{+}
  \FunctionTok{scale\_fill\_manual}\NormalTok{(}\AttributeTok{values  =} \FunctionTok{c}\NormalTok{(}\AttributeTok{Deliberate =} \StringTok{"\#2166AC"}\NormalTok{), }\AttributeTok{name =} \ConstantTok{NULL}\NormalTok{) }\SpecialCharTok{+}
  \FunctionTok{labs}\NormalTok{(}\AttributeTok{title    =} \StringTok{"Log{-}RT distribution"}\NormalTok{,}
       \AttributeTok{subtitle =} \StringTok{"Log transformation clarifies the deliberate component"}\NormalTok{,}
       \AttributeTok{x =} \StringTok{"log RT"}\NormalTok{, }\AttributeTok{y =} \StringTok{"Density"}\NormalTok{) }\SpecialCharTok{+}
  \FunctionTok{theme}\NormalTok{(}\AttributeTok{legend.position =} \StringTok{"top"}\NormalTok{)}

\NormalTok{p1 }\SpecialCharTok{/}\NormalTok{ p2}
\end{Highlighting}
\end{Shaded}

\begin{center}
\includegraphics[width=0.8\linewidth,height=\textheight,keepaspectratio]{08-MixtureModels_files/figure-pdf/ch09_rt-motivation-1.pdf}
\end{center}

In the RT case, the two components have genuinely different support
structures --- the lapse process contributes probability mass across a
1,800 ms window, while the deliberate process is tightly concentrated.
Each observation therefore carries real information about which
component generated it, so the mixture is well-identified from the
marginal distribution alone.

We will return to this point in Section 6. For binary choice data, the
identifiability picture is dramatically different.

\begin{center}\rule{0.5\linewidth}{0.5pt}\end{center}

\section{Theoretical Framework}\label{theoretical-framework}

\subsection{The Mixture Density}\label{the-mixture-density}

A \(K\)-component mixture model assigns the following marginal density
(or mass function) to each observation \(y_i\):

\[
p(y_i) = \sum_{k=1}^{K} \pi_k \, p_k(y_i \mid \theta_k),
\qquad \pi_k \ge 0, \quad \sum_{k=1}^{K} \pi_k = 1.
\tag{1}
\]

Each \(p_k(\cdot \mid \theta_k)\) is a component distribution with its
own parameter vector \(\theta_k\), and \(\pi_k\) is the corresponding
mixing weight. Given independent observations, the joint likelihood is

\[
p(y \mid \pi, \theta) = \prod_{i=1}^{n} \sum_{k=1}^{K} \pi_k \, p_k(y_i \mid \theta_k).
\tag{2}
\]

The log-likelihood therefore contains a \textbf{sum inside a log} ---
this cannot be factored into independent additive contributions per
component, which makes mixture models more computationally demanding
than standard regression.

\subsection{The Latent-Variable
Augmentation}\label{the-latent-variable-augmentation}

An equivalent and conceptually useful representation introduces a latent
component indicator \(z_i \in \{1, \ldots, K\}\):

\[
z_i \sim \text{Categorical}(\boldsymbol{\pi}), \qquad
y_i \mid z_i = k \sim p_k(y_i \mid \theta_k).
\tag{3}
\]

If we could observe \(z_i\) the model would decompose into \(K\)
separate standard likelihoods. Unfortunately \(z_i\) is unobserved, so
we sum (integrate) over all possible component assignments to recover
equation (1). This is the \textbf{marginalisation} step.

\subsection{Why Stan Requires
Marginalisation}\label{why-stan-requires-marginalisation}

Stan's Hamiltonian Monte Carlo sampler operates on a continuous,
differentiable log posterior. It cannot directly sample discrete random
variables such as \(z_i\). \textbf{Marginalising \(z_i\) before writing
any Stan code is not optional; it is a structural requirement.} The
joint model \((z_i, y_i)\) must be reduced to the marginal model
\(p(y_i)\) expressed entirely through continuous parameters.

Stan provides \texttt{log\_mix()} as a numerically stable implementation
of the per-observation mixture contribution:

\[
\texttt{log\_mix}(\pi, \ell_1, \ell_2) \;=\; \log\!\bigl[\pi \, e^{\ell_1} + (1-\pi) \, e^{\ell_2}\bigr],
\tag{4}
\]

where \(\ell_1 = \log p_1(y_i \mid \theta_1)\) and
\(\ell_2 = \log p_2(y_i \mid \theta_2)\).

\subsection{Our Running Model: Bias--Noise Mixture for Binary
Choices}\label{our-running-model-biasnoise-mixture-for-binary-choices}

We model binary choices as a mixture of a \textbf{biased process} and a
\textbf{random process}:

\[
y_i \sim (1-\pi)\,\text{Bernoulli}(\theta) + \pi\,\text{Bernoulli}(0.5),
\tag{5}
\]

where \(\theta = \text{logit}^{-1}(\alpha)\) is the probability of
choosing ``right'' under the biased process, and
\(\pi = \text{logit}^{-1}(\nu)\) is the probability that any given trial
is generated by the random process. Both are modelled on the logit
scale:

\[
\alpha \sim \text{Normal}(0,\, 1), \qquad
\nu \sim \text{Normal}(-1,\, 1).
\tag{6}
\]

The centring of \(\nu\) at \(-1\) weakly favours noise below 50 \%,
reflecting the prior belief that most trials are deliberate.

\begin{center}\rule{0.5\linewidth}{0.5pt}\end{center}

\section{Prior Predictive Check}\label{prior-predictive-check-1}

Before writing any Stan code, we simulate from the prior to verify that
the implied distribution over observable behaviour is scientifically
plausible.

\begin{quote}
Why check priors before writing Stan code? A prior that assigns
\textgreater30\% mass to noise rates above π = 0.5 means we expect most
trials to be random --- a scientifically implausible starting point that
would corrupt the posterior before data are ever observed.
\end{quote}

\begin{Shaded}
\begin{Highlighting}[]
\FunctionTok{set.seed}\NormalTok{(}\DecValTok{8}\NormalTok{)}
\NormalTok{n\_prior }\OtherTok{\textless{}{-}} \DecValTok{10000}

\NormalTok{prior\_df }\OtherTok{\textless{}{-}} \FunctionTok{tibble}\NormalTok{(}
  \AttributeTok{alpha    =} \FunctionTok{rnorm}\NormalTok{(n\_prior, }\DecValTok{0}\NormalTok{,  }\DecValTok{1}\NormalTok{),}
  \AttributeTok{nu       =} \FunctionTok{rnorm}\NormalTok{(n\_prior, }\SpecialCharTok{{-}}\DecValTok{1}\NormalTok{, }\DecValTok{1}\NormalTok{),}
  \AttributeTok{theta    =} \FunctionTok{plogis}\NormalTok{(alpha),}
  \AttributeTok{pi\_noise =} \FunctionTok{plogis}\NormalTok{(nu),}
  \CommentTok{\# Marginalised effective probability of choosing "right"}
  \AttributeTok{theta\_eff =}\NormalTok{ (}\DecValTok{1} \SpecialCharTok{{-}} \FunctionTok{plogis}\NormalTok{(nu)) }\SpecialCharTok{*} \FunctionTok{plogis}\NormalTok{(alpha) }\SpecialCharTok{+} \FloatTok{0.5} \SpecialCharTok{*} \FunctionTok{plogis}\NormalTok{(nu)}
\NormalTok{)}

\NormalTok{p\_bias }\OtherTok{\textless{}{-}} \FunctionTok{ggplot}\NormalTok{(prior\_df, }\FunctionTok{aes}\NormalTok{(}\AttributeTok{x =}\NormalTok{ theta)) }\SpecialCharTok{+}
  \FunctionTok{geom\_histogram}\NormalTok{(}\FunctionTok{aes}\NormalTok{(}\AttributeTok{y =} \FunctionTok{after\_stat}\NormalTok{(density)),}
                 \AttributeTok{bins =} \DecValTok{60}\NormalTok{, }\AttributeTok{fill =} \StringTok{"\#2166AC"}\NormalTok{, }\AttributeTok{alpha =} \FloatTok{0.7}\NormalTok{) }\SpecialCharTok{+}
  \FunctionTok{labs}\NormalTok{(}\AttributeTok{title    =} \FunctionTok{expression}\NormalTok{(}\StringTok{"Prior on "} \SpecialCharTok{*}\NormalTok{ theta }\SpecialCharTok{*} \StringTok{" (bias)"}\NormalTok{),}
       \AttributeTok{subtitle =} \StringTok{"Normal(0, 1) on logit scale"}\NormalTok{,}
       \AttributeTok{x =} \FunctionTok{expression}\NormalTok{(theta), }\AttributeTok{y =} \StringTok{"Density"}\NormalTok{) }\SpecialCharTok{+}
  \FunctionTok{scale\_x\_continuous}\NormalTok{(}\AttributeTok{limits =} \FunctionTok{c}\NormalTok{(}\DecValTok{0}\NormalTok{, }\DecValTok{1}\NormalTok{))}

\NormalTok{p\_noise }\OtherTok{\textless{}{-}} \FunctionTok{ggplot}\NormalTok{(prior\_df, }\FunctionTok{aes}\NormalTok{(}\AttributeTok{x =}\NormalTok{ pi\_noise)) }\SpecialCharTok{+}
  \FunctionTok{geom\_histogram}\NormalTok{(}\FunctionTok{aes}\NormalTok{(}\AttributeTok{y =} \FunctionTok{after\_stat}\NormalTok{(density)),}
                 \AttributeTok{bins =} \DecValTok{60}\NormalTok{, }\AttributeTok{fill =} \StringTok{"\#D6604D"}\NormalTok{, }\AttributeTok{alpha =} \FloatTok{0.7}\NormalTok{) }\SpecialCharTok{+}
  \FunctionTok{labs}\NormalTok{(}\AttributeTok{title    =} \FunctionTok{expression}\NormalTok{(}\StringTok{"Prior on "} \SpecialCharTok{*}\NormalTok{ pi }\SpecialCharTok{*} \StringTok{" (noise weight)"}\NormalTok{),}
       \AttributeTok{subtitle =} \StringTok{"Normal({-}1, 1) on logit scale"}\NormalTok{,}
       \AttributeTok{x =} \FunctionTok{expression}\NormalTok{(pi), }\AttributeTok{y =} \StringTok{"Density"}\NormalTok{) }\SpecialCharTok{+}
  \FunctionTok{scale\_x\_continuous}\NormalTok{(}\AttributeTok{limits =} \FunctionTok{c}\NormalTok{(}\DecValTok{0}\NormalTok{, }\DecValTok{1}\NormalTok{))}

\NormalTok{p\_eff }\OtherTok{\textless{}{-}} \FunctionTok{ggplot}\NormalTok{(prior\_df, }\FunctionTok{aes}\NormalTok{(}\AttributeTok{x =}\NormalTok{ theta\_eff)) }\SpecialCharTok{+}
  \FunctionTok{geom\_histogram}\NormalTok{(}\FunctionTok{aes}\NormalTok{(}\AttributeTok{y =} \FunctionTok{after\_stat}\NormalTok{(density)),}
                 \AttributeTok{bins =} \DecValTok{60}\NormalTok{, }\AttributeTok{fill =} \StringTok{"grey40"}\NormalTok{, }\AttributeTok{alpha =} \FloatTok{0.7}\NormalTok{) }\SpecialCharTok{+}
  \FunctionTok{geom\_vline}\NormalTok{(}\AttributeTok{xintercept =} \FloatTok{0.5}\NormalTok{, }\AttributeTok{linetype =} \StringTok{"dashed"}\NormalTok{, }\AttributeTok{colour =} \StringTok{"black"}\NormalTok{) }\SpecialCharTok{+}
  \FunctionTok{labs}\NormalTok{(}\AttributeTok{title    =} \FunctionTok{expression}\NormalTok{(}\StringTok{"Prior on "} \SpecialCharTok{*}\NormalTok{ theta[eff] }\SpecialCharTok{*} \StringTok{" = (1{-}"} \SpecialCharTok{*}\NormalTok{ pi }\SpecialCharTok{*} \StringTok{")"} \SpecialCharTok{*}\NormalTok{ theta }\SpecialCharTok{*} \StringTok{" + 0.5"} \SpecialCharTok{*}\NormalTok{ pi),}
       \AttributeTok{subtitle =} \StringTok{"Marginalised effective choice probability"}\NormalTok{,}
       \AttributeTok{x =} \FunctionTok{expression}\NormalTok{(theta[eff]), }\AttributeTok{y =} \StringTok{"Density"}\NormalTok{) }\SpecialCharTok{+}
  \FunctionTok{scale\_x\_continuous}\NormalTok{(}\AttributeTok{limits =} \FunctionTok{c}\NormalTok{(}\DecValTok{0}\NormalTok{, }\DecValTok{1}\NormalTok{))}

\NormalTok{(p\_bias }\SpecialCharTok{+}\NormalTok{ p\_noise) }\SpecialCharTok{/}\NormalTok{ p\_eff}
\end{Highlighting}
\end{Shaded}

\begin{center}
\includegraphics[width=0.8\linewidth,height=\textheight,keepaspectratio]{08-MixtureModels_files/figure-pdf/ch09_prior-predictive-1.pdf}
\end{center}

The prior on \(\theta\) is diffuse but centred at 0.5, spanning the full
plausible range of biases. The noise prior concentrates most mass below
\(\pi = 0.5\), reflecting the scientific expectation that participants
are mostly deliberate. The marginalised effective probability
\(\theta_{\text{eff}}\) is correspondingly diffuse but slightly shrunk
toward 0.5. These are all scientifically defensible starting points.

\begin{quote}
\textbf{Note:} If you found that the noise prior placed more than ≈ 30
\% of its mass above \(\pi = 0.4\), that would indicate the prior is too
permissive. This is the value of checking the prior predictive before
fitting: it costs nothing and can prevent days of debugging a
misspecified model.
\end{quote}

\begin{center}\rule{0.5\linewidth}{0.5pt}\end{center}

\section{Stan Implementation: Single-Agent
Model}\label{stan-implementation-single-agent-model}

\subsection{A Critical Vectorisation
Trap}\label{a-critical-vectorisation-trap}

The Stan code in the original version of this chapter contained a subtle
but consequential error. It called:

\begin{verbatim}
target += log_sum_exp(
  log(pi) + bernoulli_logit_lpmf(h | 0),
  log(1-pi) + bernoulli_logit_lpmf(h | bias)
);
\end{verbatim}

where \texttt{h} is a \textbf{vector}.
\texttt{bernoulli\_logit\_lpmf(h\ \textbar{}\ 0)} with a vector
\texttt{h} returns the \textbf{scalar} \(\sum_i \log p_1(y_i)\) --- the
joint log-probability of all observations under component 1.
\texttt{log\_sum\_exp} then computes:

\[
\log\!\Bigl[\pi \prod_i p_1(y_i) \;+\; (1-\pi) \prod_i p_2(y_i)\Bigr].
\]

This is not a mixture model. It is a model in which \textbf{all \(n\)
observations simultaneously come from one component or the other} ---
equivalent to a global discrete model-selection variable, not per-trial
latent class assignment. The correct log-likelihood is:

\[
\sum_{i=1}^{n} \log\!\Bigl[\pi \, p_1(y_i) + (1-\pi) \, p_2(y_i)\Bigr],
\]

which requires iterating \texttt{log\_mix()} \textbf{inside a loop over
observations}.

\subsection{The Correct Stan Model}\label{the-correct-stan-model}

\begin{tcolorbox}[enhanced jigsaw, toprule=.15mm, toptitle=1mm, bottomtitle=1mm, leftrule=.75mm, rightrule=.15mm, colframe=quarto-callout-note-color-frame, left=2mm, colback=white, titlerule=0mm, coltitle=black, colbacktitle=quarto-callout-note-color!10!white, opacityback=0, title=\textcolor{quarto-callout-note-color}{\faInfo}\hspace{0.5em}{Stan Technique --- Marginalizing Discrete Parameters with
\texttt{log\_mix}}, arc=.35mm, opacitybacktitle=0.6, breakable, bottomrule=.15mm]

Stan's HMC sampler requires a \textbf{continuous, differentiable}
log-posterior. A latent component indicator \(z_i \in \{0, 1\}\) (which
component generated trial \(i\)) is \textbf{discrete} --- HMC cannot
sample it directly.

The solution is \textbf{analytic marginalization}: instead of sampling
\(z_i\), we sum its log-probability over both possible values. For a
two-component mixture this is exactly what \texttt{log\_mix} does:

\begin{Shaded}
\begin{Highlighting}[]
\KeywordTok{target +=}\NormalTok{ log\_mix(pi,          }\CommentTok{// mixing weight (probability of component 1)}
\NormalTok{                  lp1,         }\CommentTok{// log p(obs | component 1)}
\NormalTok{                  lp2);        }\CommentTok{// log p(obs | component 2)}
\end{Highlighting}
\end{Shaded}

which computes \(\log\!\left[\pi \cdot p_1 + (1-\pi) \cdot p_2\right]\)
in a numerically stable way.

\textbf{Block placement note.} The per-observation \texttt{log\_mix} sum
lives inside the \texttt{model} block as a local computation --- it
contributes to \texttt{target} but \texttt{z\_i} is never stored. The
parameters we do sample (\texttt{alpha}, \texttt{nu}) are continuous and
HMC can take gradients through them. In \texttt{generated\ quantities},
we can \emph{recover} posterior expectations of \(z_i\) via a Bayes-rule
calculation if needed, but we do not sample it.

\textbf{Why the loop is necessary here.} Unlike a pure Bernoulli
likelihood, \texttt{log\_mix} operates on a single observation at a
time. There is no vectorized Stan function for mixture likelihoods
across observations (because each call is itself a sum of two terms that
depend on the observation), so the explicit \texttt{for\ (i\ in\ 1:n)}
loop is correct and unavoidable for the single-subject model. For the
multilevel model we will see how to structure this loop across subjects
to enable future parallelization.

\end{tcolorbox}

\begin{Shaded}
\begin{Highlighting}[]
\NormalTok{single\_mixture\_stan }\OtherTok{\textless{}{-}} \StringTok{"}
\StringTok{data \{}
\StringTok{  int\textless{}lower=1\textgreater{} n;}
\StringTok{  array[n] int\textless{}lower=0, upper=1\textgreater{} h;}
\StringTok{  real prior\_alpha\_mu;}
\StringTok{  real\textless{}lower=0\textgreater{} prior\_alpha\_sigma;}
\StringTok{  real prior\_nu\_mu;}
\StringTok{  real\textless{}lower=0\textgreater{} prior\_nu\_sigma;}
\StringTok{  int\textless{}lower=0, upper=1\textgreater{} run\_diagnostics;}
\StringTok{\}}

\StringTok{parameters \{}
\StringTok{  real alpha;        // logit{-}scale bias (biased component)}
\StringTok{  real nu;           // logit{-}scale mixing weight (random component)}
\StringTok{\}}

\StringTok{model \{}
\StringTok{  real noise\_p = inv\_logit(nu);}

\StringTok{  target += normal\_lpdf(alpha | prior\_alpha\_mu, prior\_alpha\_sigma);}
\StringTok{  target += normal\_lpdf(nu    | prior\_nu\_mu,    prior\_nu\_sigma);}

\StringTok{  // Marginalise the discrete component indicator analytically via log\_mix.}
\StringTok{  // One call per observation — vectorisation is not available for log\_mix.}
\StringTok{  for (i in 1:n)}
\StringTok{    target += log\_mix(noise\_p,}
\StringTok{                      bernoulli\_logit\_lpmf(h[i] | 0.0),}
\StringTok{                      bernoulli\_logit\_lpmf(h[i] | alpha));}
\StringTok{\}}

\StringTok{generated quantities \{}
\StringTok{  real\textless{}lower=0, upper=1\textgreater{} theta   = inv\_logit(alpha);}
\StringTok{  real\textless{}lower=0, upper=1\textgreater{} pi\_hat  = inv\_logit(nu);}
\StringTok{  real lprior = normal\_lpdf(alpha | prior\_alpha\_mu, prior\_alpha\_sigma)}
\StringTok{              + normal\_lpdf(nu    | prior\_nu\_mu,    prior\_nu\_sigma);}
\StringTok{  vector[n] log\_lik;}
\StringTok{  array[n] int pred\_choice;}
\StringTok{  if (run\_diagnostics) \{}
\StringTok{    for (i in 1:n) \{}
\StringTok{      log\_lik[i] = log\_mix(pi\_hat,}
\StringTok{                           bernoulli\_logit\_lpmf(h[i] | 0.0),}
\StringTok{                           bernoulli\_logit\_lpmf(h[i] | alpha));}
\StringTok{      if (bernoulli\_rng(pi\_hat))}
\StringTok{        pred\_choice[i] = bernoulli\_rng(0.5);}
\StringTok{      else}
\StringTok{        pred\_choice[i] = bernoulli\_rng(theta);}
\StringTok{    \}}
\StringTok{  \}}
\StringTok{\}}
\StringTok{"}

\FunctionTok{write\_stan\_file}\NormalTok{(single\_mixture\_stan, }\AttributeTok{dir =} \StringTok{"stan/"}\NormalTok{, }\AttributeTok{basename =} \StringTok{"ch09\_mixture\_single.stan"}\NormalTok{)}
\end{Highlighting}
\end{Shaded}

\begin{verbatim}
[1] "/Users/au209589/Dropbox/Teaching/AdvancedCognitiveModeling23_book/stan/ch09_mixture_single.stan"
\end{verbatim}

\begin{Shaded}
\begin{Highlighting}[]
\NormalTok{mod\_single }\OtherTok{\textless{}{-}} \FunctionTok{cmdstan\_model}\NormalTok{(}
  \StringTok{"stan/ch09\_mixture\_single.stan"}
\NormalTok{)}
\end{Highlighting}
\end{Shaded}

\subsection{Data}\label{data}

We use the simulated matching-pennies data introduced in earlier
chapters, selecting the agent with rate 0.8 and noise 0.1. This gives us
a ground truth against which to assess recovery: \(\theta = 0.8\),
\(\pi = 0.1\).

\begin{Shaded}
\begin{Highlighting}[]
\NormalTok{d    }\OtherTok{\textless{}{-}} \FunctionTok{read\_csv}\NormalTok{(}\FunctionTok{here}\NormalTok{(}\StringTok{"simdata"}\NormalTok{, }\StringTok{"ch05\_randomnoise.csv"}\NormalTok{), }\AttributeTok{show\_col\_types =} \ConstantTok{FALSE}\NormalTok{)}
\NormalTok{dd   }\OtherTok{\textless{}{-}}\NormalTok{ d }\SpecialCharTok{|\textgreater{}}\NormalTok{ dplyr}\SpecialCharTok{::}\FunctionTok{filter}\NormalTok{(true\_rate }\SpecialCharTok{==} \FloatTok{0.8}\NormalTok{, true\_noise }\SpecialCharTok{==} \FloatTok{0.1}\NormalTok{)}
\NormalTok{data }\OtherTok{\textless{}{-}} \FunctionTok{list}\NormalTok{(}
  \AttributeTok{n    =} \FunctionTok{nrow}\NormalTok{(dd),}
  \AttributeTok{h    =} \FunctionTok{as.integer}\NormalTok{(dd}\SpecialCharTok{$}\NormalTok{choice),}
  \AttributeTok{prior\_alpha\_mu    =} \DecValTok{0}\NormalTok{,   }\AttributeTok{prior\_alpha\_sigma =} \DecValTok{1}\NormalTok{,}
  \AttributeTok{prior\_nu\_mu       =} \SpecialCharTok{{-}}\DecValTok{1}\NormalTok{,  }\AttributeTok{prior\_nu\_sigma    =} \DecValTok{1}\NormalTok{,}
  \AttributeTok{run\_diagnostics   =} \DecValTok{1}\NormalTok{L}
\NormalTok{)}

\FunctionTok{cat}\NormalTok{(}\StringTok{"Trials:"}\NormalTok{, data}\SpecialCharTok{$}\NormalTok{n, }\StringTok{"| Proportion \textquotesingle{}1\textquotesingle{}:"}\NormalTok{, }\FunctionTok{round}\NormalTok{(}\FunctionTok{mean}\NormalTok{(data}\SpecialCharTok{$}\NormalTok{h), }\DecValTok{3}\NormalTok{), }\StringTok{"}\SpecialCharTok{\textbackslash{}n}\StringTok{"}\NormalTok{)}
\end{Highlighting}
\end{Shaded}

\begin{verbatim}
Trials: 120 | Proportion '1': 0.775 
\end{verbatim}

\subsection{Fitting with Pathfinder
Initialisation}\label{fitting-with-pathfinder-initialisation}

Mixture model posteriors can have multimodal or ridge-shaped geometry
(see Section 6). Pathfinder --- a multi-path variational approximation
--- provides stable starting points and avoids the Uniform(−2, 2)
default initialisation that often fails when parameters are on very
different scales or interact.

If Pathfinder itself fails (e.g., because the variational approximation
diverges near the ridge), the fallback must not be \texttt{NULL}.
Passing \texttt{NULL} to \texttt{init} triggers Stan's Uniform(−2, 2)
default, which places chains at random points that may lie far from the
high-probability region and immediately generate divergent transitions.
Instead we define an explicit per-chain jittered initialisation centred
on the prior medians: \(\alpha = 0\) (i.e., \(\theta = 0.5\)) and
\(\nu = -1\) (i.e., \(\pi \approx 0.27\)). A small Gaussian jitter of SD
0.1 ensures chains start at distinct points, which is necessary for
\(\hat{R}\) to be meaningful.

\begin{Shaded}
\begin{Highlighting}[]
\NormalTok{single\_fit\_file }\OtherTok{\textless{}{-}} \StringTok{"simmodels/ch09\_single\_mixture.RDS"}

\ControlFlowTok{if}\NormalTok{ (regenerate\_simulations }\SpecialCharTok{||} \SpecialCharTok{!}\FunctionTok{file.exists}\NormalTok{(single\_fit\_file)) \{}

  \CommentTok{\# Explicit per{-}chain jittered initialisation at prior medians.}
  \CommentTok{\# Used as Pathfinder fallback; safe for ridge{-}shaped posteriors.}
\NormalTok{  safe\_init }\OtherTok{\textless{}{-}} \FunctionTok{lapply}\NormalTok{(}\FunctionTok{seq\_len}\NormalTok{(}\DecValTok{4}\NormalTok{), }\ControlFlowTok{function}\NormalTok{(i) }\FunctionTok{list}\NormalTok{(}
    \AttributeTok{alpha =} \FunctionTok{rnorm}\NormalTok{(}\DecValTok{1}\NormalTok{, }\FloatTok{0.0}\NormalTok{, }\FloatTok{0.1}\NormalTok{),}
    \AttributeTok{nu    =} \FunctionTok{rnorm}\NormalTok{(}\DecValTok{1}\NormalTok{, }\SpecialCharTok{{-}}\FloatTok{1.0}\NormalTok{, }\FloatTok{0.1}\NormalTok{)}
\NormalTok{  ))}

\NormalTok{  init\_pf }\OtherTok{\textless{}{-}} \FunctionTok{tryCatch}\NormalTok{(}
\NormalTok{    mod\_single}\SpecialCharTok{$}\FunctionTok{pathfinder}\NormalTok{(}
      \AttributeTok{data      =}\NormalTok{ data,}
      \AttributeTok{seed      =} \DecValTok{8}\NormalTok{,}
      \AttributeTok{num\_paths =} \DecValTok{4}\NormalTok{,}
      \AttributeTok{refresh   =} \DecValTok{0}
\NormalTok{    ),}
    \AttributeTok{error =} \ControlFlowTok{function}\NormalTok{(e) \{}
      \FunctionTok{message}\NormalTok{(}\StringTok{"Pathfinder failed — falling back to jittered prior{-}centre initialisation."}\NormalTok{)}
\NormalTok{      safe\_init}
\NormalTok{    \}}
\NormalTok{  )}

\NormalTok{  fit\_single }\OtherTok{\textless{}{-}}\NormalTok{ mod\_single}\SpecialCharTok{$}\FunctionTok{sample}\NormalTok{(}
    \AttributeTok{data            =}\NormalTok{ data,}
    \AttributeTok{init            =}\NormalTok{ init\_pf,}
    \AttributeTok{seed            =} \DecValTok{8}\NormalTok{,}
    \AttributeTok{chains          =} \DecValTok{4}\NormalTok{,}
    \AttributeTok{parallel\_chains =} \DecValTok{4}\NormalTok{,}
    \AttributeTok{iter\_warmup     =} \DecValTok{1000}\NormalTok{,}
    \AttributeTok{iter\_sampling   =} \DecValTok{2000}\NormalTok{,}
    \AttributeTok{refresh         =} \DecValTok{500}
\NormalTok{  )}

\NormalTok{  fit\_single}\SpecialCharTok{$}\FunctionTok{save\_object}\NormalTok{(single\_fit\_file)}
\NormalTok{\} }\ControlFlowTok{else}\NormalTok{ \{}
\NormalTok{  fit\_single }\OtherTok{\textless{}{-}} \FunctionTok{readRDS}\NormalTok{(single\_fit\_file)}
\NormalTok{\}}

\NormalTok{fit\_single}\SpecialCharTok{$}\FunctionTok{summary}\NormalTok{(}\FunctionTok{c}\NormalTok{(}\StringTok{"alpha"}\NormalTok{, }\StringTok{"nu"}\NormalTok{, }\StringTok{"theta"}\NormalTok{, }\StringTok{"pi\_hat"}\NormalTok{))}
\end{Highlighting}
\end{Shaded}

\begin{verbatim}
# A tibble: 4 x 10
  variable   mean median     sd    mad      q5    q95  rhat ess_bulk ess_tail
  <chr>     <dbl>  <dbl>  <dbl>  <dbl>   <dbl>  <dbl> <dbl>    <dbl>    <dbl>
1 alpha     1.58   1.52  0.417  0.380   1.00    2.36   1.00    3114.    2450.
2 nu       -1.55  -1.48  0.745  0.734  -2.87   -0.465  1.00    2943.    2738.
3 theta     0.821  0.821 0.0547 0.0562  0.732   0.914  1.00    3114.    2450.
4 pi_hat    0.198  0.185 0.104  0.109   0.0536  0.386  1.00    2943.    2738.
\end{verbatim}

\subsection{Geometry Diagnostics}\label{geometry-diagnostics}

\begin{Shaded}
\begin{Highlighting}[]
\NormalTok{draws\_single }\OtherTok{\textless{}{-}} \FunctionTok{as\_draws\_df}\NormalTok{(fit\_single}\SpecialCharTok{$}\FunctionTok{draws}\NormalTok{())}

\CommentTok{\# Required diagnostic thresholds (cheatsheet):}
\CommentTok{\# Divergences = 0 | R{-}hat \textless{} 1.01 | Bulk \& Tail ESS \textgreater{} 400 | E{-}BFMI \textgreater{} 0.2}
\NormalTok{diag }\OtherTok{\textless{}{-}}\NormalTok{ fit\_single}\SpecialCharTok{$}\FunctionTok{diagnostic\_summary}\NormalTok{()}
\FunctionTok{cat}\NormalTok{(}\StringTok{"Divergences:"}\NormalTok{, }\FunctionTok{sum}\NormalTok{(diag}\SpecialCharTok{$}\NormalTok{num\_divergent), }\StringTok{"}\SpecialCharTok{\textbackslash{}n}\StringTok{"}\NormalTok{)}
\end{Highlighting}
\end{Shaded}

\begin{verbatim}
Divergences: 0 
\end{verbatim}

\begin{Shaded}
\begin{Highlighting}[]
\FunctionTok{cat}\NormalTok{(}\StringTok{"Max R{-}hat:  "}\NormalTok{, }\FunctionTok{round}\NormalTok{(}\FunctionTok{max}\NormalTok{(fit\_single}\SpecialCharTok{$}\FunctionTok{summary}\NormalTok{()}\SpecialCharTok{$}\NormalTok{rhat, }\AttributeTok{na.rm =} \ConstantTok{TRUE}\NormalTok{), }\DecValTok{4}\NormalTok{), }\StringTok{"}\SpecialCharTok{\textbackslash{}n}\StringTok{"}\NormalTok{)}
\end{Highlighting}
\end{Shaded}

\begin{verbatim}
Max R-hat:   1.0017 
\end{verbatim}

\begin{Shaded}
\begin{Highlighting}[]
\FunctionTok{cat}\NormalTok{(}\StringTok{"Min ESS:    "}\NormalTok{, }\FunctionTok{round}\NormalTok{(}\FunctionTok{min}\NormalTok{(fit\_single}\SpecialCharTok{$}\FunctionTok{summary}\NormalTok{()}\SpecialCharTok{$}\NormalTok{ess\_bulk, }\AttributeTok{na.rm =} \ConstantTok{TRUE}\NormalTok{)), }\StringTok{"}\SpecialCharTok{\textbackslash{}n}\StringTok{"}\NormalTok{)}
\end{Highlighting}
\end{Shaded}

\begin{verbatim}
Min ESS:     2861 
\end{verbatim}

\begin{Shaded}
\begin{Highlighting}[]
\CommentTok{\# Trace plots for the two sampled parameters}
\FunctionTok{mcmc\_trace}\NormalTok{(draws\_single, }\AttributeTok{pars =} \FunctionTok{c}\NormalTok{(}\StringTok{"alpha"}\NormalTok{, }\StringTok{"nu"}\NormalTok{)) }\SpecialCharTok{+}
  \FunctionTok{ggtitle}\NormalTok{(}\StringTok{"Trace plots: logit{-}scale parameters"}\NormalTok{) }\SpecialCharTok{+}
  \FunctionTok{scale\_colour\_viridis\_d}\NormalTok{(}\AttributeTok{option =} \StringTok{"mako"}\NormalTok{, }\AttributeTok{begin =} \FloatTok{0.15}\NormalTok{, }\AttributeTok{end =} \FloatTok{0.85}\NormalTok{, }\AttributeTok{guide =} \StringTok{"none"}\NormalTok{)}
\end{Highlighting}
\end{Shaded}

\begin{center}
\includegraphics[width=0.8\linewidth,height=\textheight,keepaspectratio]{08-MixtureModels_files/figure-pdf/ch09_single-diagnostics-1.pdf}
\end{center}

\begin{Shaded}
\begin{Highlighting}[]
\CommentTok{\# Pairwise geometry on the probability scale — reveals ridge if present}
\FunctionTok{mcmc\_pairs}\NormalTok{(}
\NormalTok{  fit\_single}\SpecialCharTok{$}\FunctionTok{draws}\NormalTok{(),}
  \AttributeTok{pars      =} \FunctionTok{c}\NormalTok{(}\StringTok{"theta"}\NormalTok{, }\StringTok{"pi\_hat"}\NormalTok{),}
  \AttributeTok{off\_diag\_fun =} \StringTok{"scatter"}\NormalTok{,}
  \AttributeTok{off\_diag\_args =} \FunctionTok{list}\NormalTok{(}\AttributeTok{size =} \FloatTok{0.3}\NormalTok{, }\AttributeTok{alpha =} \FloatTok{0.2}\NormalTok{)}
\NormalTok{) }\SpecialCharTok{+}
  \FunctionTok{labs}\NormalTok{(}\AttributeTok{title =} \StringTok{"Posterior geometry: bias × noise"}\NormalTok{,}
       \AttributeTok{subtitle =} \StringTok{"A ridge along the aliasing manifold is expected for binary data (see Section 6)"}\NormalTok{)}
\end{Highlighting}
\end{Shaded}

\begin{verbatim}
NULL
\end{verbatim}

\begin{center}\rule{0.5\linewidth}{0.5pt}\end{center}

\section{Validation: Simulation-Based Calibration
Checks}\label{validation-simulation-based-calibration-checks}

SBC simultaneously certifies two properties. When rank statistics are
uniform, the sampler is \textbf{computationally unbiased} (calibration).
The joint distribution of true values and posterior summaries across all
SBC datasets provides \textbf{parameter recovery} --- the same
information conveyed by a recovery grid, but derived from a principled
random draw over the full prior rather than a hand-chosen set of grid
points.

\subsection{Generator and Backend}\label{generator-and-backend}

\begin{quote}
Why SBC for mixture models specifically? Mixture posteriors are often
ridge-shaped or multimodal, which means the sampler can fail silently
--- producing low divergences but biased posteriors. SBC's
rank-uniformity test catches this structural failure that standard
diagnostics miss.
\end{quote}

\begin{Shaded}
\begin{Highlighting}[]
\NormalTok{sbc\_gen\_fn }\OtherTok{\textless{}{-}} \ControlFlowTok{function}\NormalTok{() \{}
  \CommentTok{\# Draw from prior (logit scale — must match Stan parameter names exactly)}
\NormalTok{  alpha\_true }\OtherTok{\textless{}{-}} \FunctionTok{rnorm}\NormalTok{(}\DecValTok{1}\NormalTok{, }\DecValTok{0}\NormalTok{, }\DecValTok{1}\NormalTok{)}
\NormalTok{  nu\_true    }\OtherTok{\textless{}{-}} \FunctionTok{rnorm}\NormalTok{(}\DecValTok{1}\NormalTok{, }\SpecialCharTok{{-}}\DecValTok{1}\NormalTok{, }\DecValTok{1}\NormalTok{)}

\NormalTok{  theta\_true }\OtherTok{\textless{}{-}} \FunctionTok{plogis}\NormalTok{(alpha\_true)}
\NormalTok{  pi\_true    }\OtherTok{\textless{}{-}} \FunctionTok{plogis}\NormalTok{(nu\_true)}

\NormalTok{  n }\OtherTok{\textless{}{-}} \DecValTok{120}
\NormalTok{  component }\OtherTok{\textless{}{-}} \FunctionTok{rbinom}\NormalTok{(n, }\DecValTok{1}\NormalTok{, pi\_true)}
\NormalTok{  h }\OtherTok{\textless{}{-}} \FunctionTok{as.integer}\NormalTok{(}\FunctionTok{ifelse}\NormalTok{(}
\NormalTok{    component }\SpecialCharTok{==} \DecValTok{1}\NormalTok{,}
    \FunctionTok{rbinom}\NormalTok{(n, }\DecValTok{1}\NormalTok{, }\FloatTok{0.5}\NormalTok{),}
    \FunctionTok{rbinom}\NormalTok{(n, }\DecValTok{1}\NormalTok{, theta\_true)}
\NormalTok{  ))}

  \FunctionTok{list}\NormalTok{(}
    \AttributeTok{variables =} \FunctionTok{list}\NormalTok{(}\AttributeTok{alpha =}\NormalTok{ alpha\_true, }\AttributeTok{nu =}\NormalTok{ nu\_true),}
    \AttributeTok{generated =} \FunctionTok{list}\NormalTok{(}\AttributeTok{n =}\NormalTok{ n, }\AttributeTok{h =}\NormalTok{ h)}
\NormalTok{  )}
\NormalTok{\}}

\NormalTok{sbc\_generator }\OtherTok{\textless{}{-}}\NormalTok{ SBC}\SpecialCharTok{::}\FunctionTok{SBC\_generator\_function}\NormalTok{(sbc\_gen\_fn)}

\NormalTok{sbc\_backend }\OtherTok{\textless{}{-}}\NormalTok{ SBC}\SpecialCharTok{::}\FunctionTok{SBC\_backend\_cmdstan\_sample}\NormalTok{(}
\NormalTok{  mod\_single,}
  \AttributeTok{chains        =} \DecValTok{2}\NormalTok{,}
  \AttributeTok{iter\_warmup   =} \DecValTok{500}\NormalTok{,}
  \AttributeTok{iter\_sampling =} \DecValTok{1000}\NormalTok{,}
  \AttributeTok{refresh       =} \DecValTok{0}\NormalTok{,}
  \AttributeTok{show\_messages =} \ConstantTok{FALSE}
\NormalTok{)}
\end{Highlighting}
\end{Shaded}

\begin{Shaded}
\begin{Highlighting}[]
\NormalTok{sbc\_file }\OtherTok{\textless{}{-}} \StringTok{"simmodels/ch09\_sbc\_results.RDS"}

\ControlFlowTok{if}\NormalTok{ (regenerate\_simulations }\SpecialCharTok{||} \SpecialCharTok{!}\FunctionTok{file.exists}\NormalTok{(sbc\_file)) \{}
\NormalTok{  future}\SpecialCharTok{::}\FunctionTok{plan}\NormalTok{(future}\SpecialCharTok{::}\NormalTok{multisession, }\AttributeTok{workers =} \FunctionTok{max}\NormalTok{(}\DecValTok{1}\NormalTok{, parallel}\SpecialCharTok{::}\FunctionTok{detectCores}\NormalTok{() }\SpecialCharTok{{-}} \DecValTok{2}\NormalTok{))}

  \CommentTok{\# S = 1000 certifies the rank statistic at the 95 \% level.}
  \CommentTok{\# For development, reduce to 200.}
\NormalTok{  sbc\_datasets }\OtherTok{\textless{}{-}}\NormalTok{ SBC}\SpecialCharTok{::}\FunctionTok{generate\_datasets}\NormalTok{(sbc\_generator, }\AttributeTok{n\_sims =} \DecValTok{1000}\NormalTok{)}
\NormalTok{  sbc\_results  }\OtherTok{\textless{}{-}}\NormalTok{ SBC}\SpecialCharTok{::}\FunctionTok{compute\_SBC}\NormalTok{(sbc\_datasets, sbc\_backend)}

\NormalTok{  future}\SpecialCharTok{::}\FunctionTok{plan}\NormalTok{(future}\SpecialCharTok{::}\NormalTok{sequential)}
  \FunctionTok{saveRDS}\NormalTok{(sbc\_results, sbc\_file)}
\NormalTok{\} }\ControlFlowTok{else}\NormalTok{ \{}
\NormalTok{  sbc\_results }\OtherTok{\textless{}{-}} \FunctionTok{readRDS}\NormalTok{(sbc\_file)}
\NormalTok{\}}
\end{Highlighting}
\end{Shaded}

\subsection{Calibration: ECDF-Difference
Plots}\label{calibration-ecdf-difference-plots}

The primary calibration diagnostic is the ECDF-difference plot
(Säilynoja et al.~2022). A well-calibrated sampler produces rank
statistics that are uniform on \(\{0, 1, \ldots, S\}\), so the ECDF
should fall along the horizontal zero line. Deviations outside the 95 \%
envelope indicate structural bias or over/underconfidence.

\begin{Shaded}
\begin{Highlighting}[]
\NormalTok{SBC}\SpecialCharTok{::}\FunctionTok{plot\_ecdf\_diff}\NormalTok{(sbc\_results, }\AttributeTok{variables =} \FunctionTok{c}\NormalTok{(}\StringTok{"alpha"}\NormalTok{, }\StringTok{"nu"}\NormalTok{)) }\SpecialCharTok{+}
  \FunctionTok{labs}\NormalTok{(}\AttributeTok{title    =} \StringTok{"SBC calibration: ECDF{-}difference plots"}\NormalTok{,}
       \AttributeTok{subtitle =} \StringTok{"Deviations outside the grey band indicate a miscalibrated sampler"}\NormalTok{) }\SpecialCharTok{+}
  \FunctionTok{theme}\NormalTok{(}\AttributeTok{strip.background =} \FunctionTok{element\_rect}\NormalTok{(}\AttributeTok{fill =} \StringTok{"grey95"}\NormalTok{, }\AttributeTok{colour =} \ConstantTok{NA}\NormalTok{),}
        \AttributeTok{strip.text =} \FunctionTok{element\_text}\NormalTok{(}\AttributeTok{face =} \StringTok{"bold"}\NormalTok{))}
\end{Highlighting}
\end{Shaded}

\begin{center}
\includegraphics[width=0.8\linewidth,height=\textheight,keepaspectratio]{08-MixtureModels_files/figure-pdf/ch09_sbc-ecdf-diff-1.pdf}
\end{center}

\begin{Shaded}
\begin{Highlighting}[]
\NormalTok{SBC}\SpecialCharTok{::}\FunctionTok{plot\_rank\_hist}\NormalTok{(sbc\_results, }\AttributeTok{variables =} \FunctionTok{c}\NormalTok{(}\StringTok{"alpha"}\NormalTok{, }\StringTok{"nu"}\NormalTok{)) }\SpecialCharTok{+}
  \FunctionTok{labs}\NormalTok{(}\AttributeTok{title    =} \StringTok{"SBC calibration: rank histograms"}\NormalTok{,}
       \AttributeTok{subtitle =} \StringTok{"U{-}shape → overconfident; hump → underconfident; skew → biased"}\NormalTok{) }\SpecialCharTok{+}
  \FunctionTok{theme}\NormalTok{(}\AttributeTok{strip.background =} \FunctionTok{element\_rect}\NormalTok{(}\AttributeTok{fill =} \StringTok{"grey95"}\NormalTok{, }\AttributeTok{colour =} \ConstantTok{NA}\NormalTok{),}
        \AttributeTok{strip.text =} \FunctionTok{element\_text}\NormalTok{(}\AttributeTok{face =} \StringTok{"bold"}\NormalTok{))}
\end{Highlighting}
\end{Shaded}

\begin{center}
\includegraphics[width=0.8\linewidth,height=\textheight,keepaspectratio]{08-MixtureModels_files/figure-pdf/ch09_sbc-rank-hist-1.pdf}
\end{center}

\subsection{Parameter Recovery}\label{parameter-recovery}

The \texttt{stats} component of the SBC output contains the
\texttt{simulated\_value} (the true prior draw) alongside posterior
summaries for each simulation. Plotting posterior mean against true
value directly provides the recovery assessment --- no separate
grid-based recovery study is required.

\begin{Shaded}
\begin{Highlighting}[]
\NormalTok{recovery\_df }\OtherTok{\textless{}{-}}\NormalTok{ sbc\_results}\SpecialCharTok{$}\NormalTok{stats }\SpecialCharTok{|\textgreater{}}
\NormalTok{  dplyr}\SpecialCharTok{::}\FunctionTok{filter}\NormalTok{(variable }\SpecialCharTok{\%in\%} \FunctionTok{c}\NormalTok{(}\StringTok{"alpha"}\NormalTok{, }\StringTok{"nu"}\NormalTok{)) }\SpecialCharTok{|\textgreater{}}
\NormalTok{  dplyr}\SpecialCharTok{::}\FunctionTok{mutate}\NormalTok{(}
    \AttributeTok{param\_label =}\NormalTok{ dplyr}\SpecialCharTok{::}\FunctionTok{recode}\NormalTok{(variable,}
      \AttributeTok{alpha =} \StringTok{"alpha (logit bias)"}\NormalTok{,}
      \AttributeTok{nu    =} \StringTok{"nu (logit noise)"}
\NormalTok{    ),}
    \CommentTok{\# Calibration coverage: does the 95 \% CrI contain the truth?}
    \AttributeTok{covered\_95 =}\NormalTok{ simulated\_value }\SpecialCharTok{\textgreater{}=}\NormalTok{ q5 }\SpecialCharTok{\&}\NormalTok{ simulated\_value }\SpecialCharTok{\textless{}=}\NormalTok{ q95}
\NormalTok{  )}

\CommentTok{\# Coverage summary}
\NormalTok{recovery\_df }\SpecialCharTok{|\textgreater{}}
\NormalTok{  dplyr}\SpecialCharTok{::}\FunctionTok{group\_by}\NormalTok{(param\_label) }\SpecialCharTok{|\textgreater{}}
\NormalTok{  dplyr}\SpecialCharTok{::}\FunctionTok{summarise}\NormalTok{(}
    \AttributeTok{coverage\_95 =} \FunctionTok{mean}\NormalTok{(covered\_95),}
    \AttributeTok{rmse        =} \FunctionTok{sqrt}\NormalTok{(}\FunctionTok{mean}\NormalTok{((mean }\SpecialCharTok{{-}}\NormalTok{ simulated\_value)}\SpecialCharTok{\^{}}\DecValTok{2}\NormalTok{)),}
    \AttributeTok{.groups =} \StringTok{"drop"}
\NormalTok{  ) }\SpecialCharTok{|\textgreater{}}
  \FunctionTok{print}\NormalTok{()}
\end{Highlighting}
\end{Shaded}

\begin{verbatim}
# A tibble: 2 x 3
  param_label        coverage_95  rmse
  <chr>                    <dbl> <dbl>
1 alpha (logit bias)       0.893 0.403
2 nu (logit noise)         0.897 0.953
\end{verbatim}

\begin{Shaded}
\begin{Highlighting}[]
\CommentTok{\# Recovery scatter: posterior mean vs. true value}
\FunctionTok{ggplot}\NormalTok{(recovery\_df, }\FunctionTok{aes}\NormalTok{(}\AttributeTok{x =}\NormalTok{ simulated\_value, }\AttributeTok{y =}\NormalTok{ mean)) }\SpecialCharTok{+}
  \FunctionTok{geom\_point}\NormalTok{(}\FunctionTok{aes}\NormalTok{(}\AttributeTok{colour =}\NormalTok{ covered\_95), }\AttributeTok{size =} \FloatTok{0.7}\NormalTok{, }\AttributeTok{alpha =} \FloatTok{0.5}\NormalTok{) }\SpecialCharTok{+}
  \FunctionTok{geom\_linerange}\NormalTok{(}\FunctionTok{aes}\NormalTok{(}\AttributeTok{ymin =}\NormalTok{ q5, }\AttributeTok{ymax =}\NormalTok{ q95, }\AttributeTok{colour =}\NormalTok{ covered\_95),}
                 \AttributeTok{linewidth =} \FloatTok{0.25}\NormalTok{, }\AttributeTok{alpha =} \FloatTok{0.3}\NormalTok{) }\SpecialCharTok{+}
  \FunctionTok{geom\_abline}\NormalTok{(}\AttributeTok{intercept =} \DecValTok{0}\NormalTok{, }\AttributeTok{slope =} \DecValTok{1}\NormalTok{,}
              \AttributeTok{linetype =} \StringTok{"dashed"}\NormalTok{, }\AttributeTok{colour =} \StringTok{"black"}\NormalTok{, }\AttributeTok{linewidth =} \FloatTok{0.8}\NormalTok{) }\SpecialCharTok{+}
  \FunctionTok{facet\_wrap}\NormalTok{(}\SpecialCharTok{\textasciitilde{}}\NormalTok{param\_label, }\AttributeTok{scales =} \StringTok{"free"}\NormalTok{) }\SpecialCharTok{+}
  \FunctionTok{scale\_colour\_manual}\NormalTok{(}
    \AttributeTok{values =} \FunctionTok{c}\NormalTok{(}\StringTok{"TRUE"} \OtherTok{=} \StringTok{"\#2166AC"}\NormalTok{, }\StringTok{"FALSE"} \OtherTok{=} \StringTok{"\#D6604D"}\NormalTok{),}
    \AttributeTok{name   =} \StringTok{"95 \% CrI covers truth"}
\NormalTok{  ) }\SpecialCharTok{+}
  \FunctionTok{labs}\NormalTok{(}
    \AttributeTok{title    =} \StringTok{"Parameter recovery from SBC simulations"}\NormalTok{,}
    \AttributeTok{subtitle =} \StringTok{"Each point is one SBC dataset; dashed line = perfect recovery (y = x)"}\NormalTok{,}
    \AttributeTok{x        =} \StringTok{"True value (prior draw)"}\NormalTok{,}
    \AttributeTok{y        =} \StringTok{"Posterior mean"}
\NormalTok{  ) }\SpecialCharTok{+}
  \FunctionTok{theme}\NormalTok{(}
    \AttributeTok{legend.position    =} \StringTok{"top"}\NormalTok{,}
    \AttributeTok{strip.background   =} \FunctionTok{element\_rect}\NormalTok{(}\AttributeTok{fill =} \StringTok{"grey95"}\NormalTok{, }\AttributeTok{colour =} \ConstantTok{NA}\NormalTok{),}
    \AttributeTok{strip.text         =} \FunctionTok{element\_text}\NormalTok{(}\AttributeTok{face =} \StringTok{"bold"}\NormalTok{),}
    \AttributeTok{panel.grid.major   =} \FunctionTok{element\_line}\NormalTok{(}\AttributeTok{colour =} \StringTok{"grey92"}\NormalTok{, }\AttributeTok{linetype =} \StringTok{"dotted"}\NormalTok{)}
\NormalTok{  )}
\end{Highlighting}
\end{Shaded}

\begin{center}
\includegraphics[width=0.8\linewidth,height=\textheight,keepaspectratio]{08-MixtureModels_files/figure-pdf/ch09_sbc-recovery-1.pdf}
\end{center}

Recovery should be unbiased across the full prior range. The fraction of
95 \% credible intervals covering the truth should be close to 0.95.
Systematic deviations --- for instance, the noise parameter \(\nu\)
recovering poorly near \(\nu \approx -2\) (i.e., \(\pi \approx 0.12\))
--- would indicate that 120 trials provide insufficient information to
distinguish very low noise from zero noise, a legitimate finding about
the experimental design, not a modelling failure.

\begin{center}\rule{0.5\linewidth}{0.5pt}\end{center}

\section{Identifiability: The Bernoulli Aliasing
Problem}\label{identifiability-the-bernoulli-aliasing-problem}

\subsection{The Algebraic Constraint}\label{the-algebraic-constraint}

\begin{quote}
Why prove identifiability analytically before checking it empirically?
The aliasing manifold exists regardless of sample size or model quality.
No amount of data or better priors can break a fundamental algebraic
degeneracy. Understanding the constraint analytically tells us exactly
what additional information (RT, condition manipulation) would be
required --- rather than leaving us puzzled when LOO shows no
preference.
\end{quote}

For a single binary trial the probability of observing \(y_i = 1\) under
the bias--noise mixture is:

\[
P(y_i = 1) = (1 - \pi)\,\theta + \pi \cdot 0.5 = \theta - \pi(\theta - 0.5).
\tag{7}
\]

This depends on \((\theta, \pi)\) only through the linear combination
\(\theta_{\text{eff}} = (1-\pi)\theta + 0.5\pi\). For i.i.d. binary
observations, the \textbf{entire dataset} carries information only about
this one-dimensional quantity. Consequently, infinitely many
\((\theta, \pi)\) pairs are observationally equivalent: for any fixed
observed proportion \(\bar{y}\), the set

\[
\bigl\{(\theta, \pi) : (1-\pi)\,\theta + 0.5\,\pi = \bar{y}\bigr\}
\tag{8}
\]

is a curve in parameter space, and the posterior will be smeared along
this curve. This is the \textbf{aliasing manifold}.

The mixture model is not wrong --- it is simply not identified from the
marginal binary likelihood alone. The bias and noise parameters are
individually interpretable and scientifically meaningful; the data just
cannot separate them without additional structure (e.g., trial-level
covariates, joint RT and accuracy modelling, or experimental
manipulations of the noise rate).

\subsection{Visualising the Posterior
Ridge}\label{visualising-the-posterior-ridge}

\begin{Shaded}
\begin{Highlighting}[]
\NormalTok{draws\_prob }\OtherTok{\textless{}{-}} \FunctionTok{as\_draws\_df}\NormalTok{(fit\_single}\SpecialCharTok{$}\FunctionTok{draws}\NormalTok{()) }\SpecialCharTok{|\textgreater{}}
\NormalTok{  dplyr}\SpecialCharTok{::}\FunctionTok{select}\NormalTok{(theta, pi\_hat)}

\NormalTok{obs\_prop   }\OtherTok{\textless{}{-}} \FunctionTok{mean}\NormalTok{(data}\SpecialCharTok{$}\NormalTok{h)}

\CommentTok{\# Aliasing manifold: all (theta, pi) consistent with observed proportion}
\NormalTok{pi\_grid    }\OtherTok{\textless{}{-}} \FunctionTok{seq}\NormalTok{(}\DecValTok{0}\NormalTok{, }\FloatTok{0.95}\NormalTok{, }\AttributeTok{by =} \FloatTok{0.002}\NormalTok{)}
\NormalTok{theta\_mani }\OtherTok{\textless{}{-}}\NormalTok{ (obs\_prop }\SpecialCharTok{{-}} \FloatTok{0.5} \SpecialCharTok{*}\NormalTok{ pi\_grid) }\SpecialCharTok{/}\NormalTok{ (}\DecValTok{1} \SpecialCharTok{{-}}\NormalTok{ pi\_grid)}
\NormalTok{manifold   }\OtherTok{\textless{}{-}} \FunctionTok{tibble}\NormalTok{(}\AttributeTok{pi\_hat =}\NormalTok{ pi\_grid, }\AttributeTok{theta =}\NormalTok{ theta\_mani) }\SpecialCharTok{|\textgreater{}}
\NormalTok{  dplyr}\SpecialCharTok{::}\FunctionTok{filter}\NormalTok{(theta }\SpecialCharTok{\textgreater{}=} \DecValTok{0}\NormalTok{, theta }\SpecialCharTok{\textless{}=} \DecValTok{1}\NormalTok{)}

\FunctionTok{ggplot}\NormalTok{(draws\_prob, }\FunctionTok{aes}\NormalTok{(}\AttributeTok{x =}\NormalTok{ pi\_hat, }\AttributeTok{y =}\NormalTok{ theta)) }\SpecialCharTok{+}
  \FunctionTok{geom\_point}\NormalTok{(}\AttributeTok{alpha =} \FloatTok{0.06}\NormalTok{, }\AttributeTok{size =} \FloatTok{0.5}\NormalTok{, }\AttributeTok{colour =} \StringTok{"\#2166AC"}\NormalTok{) }\SpecialCharTok{+}
  \FunctionTok{geom\_line}\NormalTok{(}\AttributeTok{data =}\NormalTok{ manifold, }\AttributeTok{colour =} \StringTok{"\#D6604D"}\NormalTok{, }\AttributeTok{linewidth =} \FloatTok{1.0}\NormalTok{) }\SpecialCharTok{+}
  \FunctionTok{annotate}\NormalTok{(}\StringTok{"text"}\NormalTok{, }\AttributeTok{x =} \FloatTok{0.45}\NormalTok{, }\AttributeTok{y =} \FloatTok{0.92}\NormalTok{, }\AttributeTok{colour =} \StringTok{"\#D6604D"}\NormalTok{, }\AttributeTok{hjust =} \DecValTok{0}\NormalTok{,}
           \AttributeTok{label =} \FunctionTok{expression}\NormalTok{(}\FunctionTok{hat}\NormalTok{(theta)[eff] }\SpecialCharTok{==} \FunctionTok{bar}\NormalTok{(y))) }\SpecialCharTok{+}
  \FunctionTok{labs}\NormalTok{(}
    \AttributeTok{title    =} \StringTok{"Posterior geometry and the aliasing manifold"}\NormalTok{,}
    \AttributeTok{subtitle =} \StringTok{"Red curve: all (π, θ) pairs consistent with the observed choice proportion"}\NormalTok{,}
    \AttributeTok{x        =} \FunctionTok{expression}\NormalTok{(pi }\SpecialCharTok{\textasciitilde{}} \StringTok{"(noise weight)"}\NormalTok{),}
    \AttributeTok{y        =} \FunctionTok{expression}\NormalTok{(theta }\SpecialCharTok{\textasciitilde{}} \StringTok{"(bias)"}\NormalTok{)}
\NormalTok{  ) }\SpecialCharTok{+}
  \FunctionTok{coord\_cartesian}\NormalTok{(}\AttributeTok{xlim =} \FunctionTok{c}\NormalTok{(}\DecValTok{0}\NormalTok{, }\FloatTok{0.6}\NormalTok{), }\AttributeTok{ylim =} \FunctionTok{c}\NormalTok{(}\FloatTok{0.5}\NormalTok{, }\FloatTok{1.0}\NormalTok{))}
\end{Highlighting}
\end{Shaded}

\begin{center}
\includegraphics[width=0.8\linewidth,height=\textheight,keepaspectratio]{08-MixtureModels_files/figure-pdf/ch09_aliasing-manifold-1.pdf}
\end{center}

The posterior concentrates along the red aliasing manifold rather than
converging to a point. This is not a sampling pathology --- it is the
correct Bayesian representation of what the data actually contain. The
prior does provide some regularisation (the posterior does not extend to
the full manifold), but with 120 binary observations and diffuse priors,
the two parameters remain substantially correlated.

\subsection{Why the RT Case is
Different}\label{why-the-rt-case-is-different}

Recall from Section 1 that the RT mixture was visually transparent. The
reason is that the two components have \textbf{different likelihoods for
the same observation value}. A 400 ms response is extremely unlikely
under the uniform lapse process but highly probable under the log-normal
deliberate process. This trial-level likelihood contrast provides
information about \(z_i\) for each observation, breaking the aliasing.
With binary choices, both the biased process (at some \(\theta\)) and
the random process (at 0.5) assign non-zero probability to both
outcomes, so the aliasing cannot be broken by observation values alone.

This comparison motivates a practical principle: if you want a mixture
model's components to be individually identified, ensure that the
components are distinguishable at the level of individual observations,
not merely at the level of marginal statistics. Two strategies
accomplish this in practice.

\textbf{Strategy A: Joint modelling of complementary outcome types.} If
reaction time \(rt_i\) is observed alongside each binary choice \(y_i\),
the per-trial joint likelihood becomes:

\[
p(y_i, rt_i) = (1-\pi)\; p_{\text{bias}}(y_i \mid \theta)\; p_{\text{LN}}(rt_i \mid \mu, \sigma) \;+\; \pi\; p_{\text{rand}}(y_i \mid 0.5)\; p_{\text{unif}}(rt_i \mid a, b).
\tag{8a}
\]

Because
\(p_{\text{LN}}(rt_i \mid \mu, \sigma) / p_{\text{unif}}(rt_i \mid a, b)\)
varies continuously with \(rt_i\), every observation now provides direct
evidence about \(z_i\). A short RT carries high likelihood under the
log-normal deliberate component and low likelihood under the uniform
lapse component, or vice versa for an unusually long RT. The aliasing
manifold collapses to a point: the two-dimensional parameter space
\((\theta, \pi)\) is identified by the joint data stream, not just the
marginal choice proportion.

\textbf{Strategy B: Experimental manipulation of the noise rate.}
Suppose participants complete a single-task block (\(c_i = 0\)) and a
dual-task block (\(c_i = 1\)), where the additional cognitive load is
expected to increase lapse frequency but leave the deliberate bias
intact. We model \(\pi\) as a function of condition while holding
\(\theta\) constant:

\[
\nu_i = \gamma_0 + \gamma_1 c_i, \qquad \pi_i = \text{logit}^{-1}(\nu_i).
\tag{8b}
\]

The aliasing constraint now applies separately within each condition:

\[
(1 - \pi_0)\,\theta + 0.5\,\pi_0 = \bar{y}_0, \qquad
(1 - \pi_1)\,\theta + 0.5\,\pi_1 = \bar{y}_1.
\tag{8c}
\]

These are two linear equations in the three unknowns
\((\theta, \pi_0, \pi_1)\) --- or equivalently
\((\theta, \gamma_0, \gamma_1)\). If \(\pi_0 \ne \pi_1\) (i.e., if the
dual-task manipulation actually changes the noise rate), the two
manifolds intersect at a unique point and all three parameters are
identified. The identifiability here does not require RT data at all:
the experimental design does the work by breaking the single-condition
degeneracy.

\begin{center}\rule{0.5\linewidth}{0.5pt}\end{center}

\section{Prior Sensitivity Analysis}\label{prior-sensitivity-analysis}

Section 6 established analytically that the binary likelihood cannot
separate \(\theta\) and \(\pi\): the data constrain only
\(\theta_{\text{eff}} = (1-\pi)\theta + 0.5\pi\). A direct corollary is
that the posterior for at least one of the individual parameters must be
sensitive to the prior. \texttt{priorsense::powerscale\_sensitivity()}
quantifies this sensitivity without refitting the model. It uses
Pareto-smoothed importance sampling to evaluate the perturbed posterior:

\[
p_{\alpha,\beta}(\theta, \pi \mid y) \propto p(y \mid \theta, \pi)^{\alpha}\; p(\theta, \pi)^{\beta},
\tag{11}
\]

where \(\alpha\) and \(\beta\) are power-scaling factors. Sensitivity to
the prior corresponds to varying \(\beta\) while holding \(\alpha = 1\);
sensitivity to the likelihood corresponds to varying \(\alpha\) while
holding \(\beta = 1\). The summary statistic \(\hat{\psi}\) measures how
much the posterior mean shifts (in units of posterior SD) as the scaling
factor departs from 1. Values above \(\approx 0.05\) indicate
non-negligible sensitivity.

The \texttt{priorsense} workflow requires a \texttt{lprior} variable in
the Stan model's \texttt{generated\ quantities} block, which accumulates
the log prior density. This was added to the single-agent Stan model in
§4.2 and will be added to the multilevel model in §8.2. The
\texttt{powerscale\_sensitivity()} call then reads \texttt{lprior} as
the prior contribution and uses \texttt{log\_lik} as the likelihood
contribution.

\begin{quote}
Why expect high sensitivity for the noise parameter? Section 6 showed
algebraically that binary i.i.d. data constrain only θ\_eff = (1-π)θ +
0.5π. The noise weight π therefore receives almost no direct information
from the data --- its posterior is largely proportional to its prior. A
high ψ̂ for ν is not a model failure; it is the quantitative fingerprint
of this degeneracy.
\end{quote}

\begin{Shaded}
\begin{Highlighting}[]
\CommentTok{\# powerscale\_sensitivity works with CmdStanFit objects via the posterior package.}
\CommentTok{\# It reads \textquotesingle{}lprior\textquotesingle{} and \textquotesingle{}log\_lik\textquotesingle{} from generated quantities automatically.}
\NormalTok{sensitivity }\OtherTok{\textless{}{-}}\NormalTok{ priorsense}\SpecialCharTok{::}\FunctionTok{powerscale\_sensitivity}\NormalTok{(}
\NormalTok{  fit\_single,}
  \AttributeTok{variable =} \FunctionTok{c}\NormalTok{(}\StringTok{"alpha"}\NormalTok{, }\StringTok{"nu"}\NormalTok{)}
\NormalTok{)}

\FunctionTok{print}\NormalTok{(sensitivity)}
\end{Highlighting}
\end{Shaded}

\begin{verbatim}
Sensitivity based on cjs_dist
Prior selection: all priors
Likelihood selection: all data

 variable prior likelihood                                diagnosis
    alpha 0.140      0.063            potential prior-data conflict
       nu 0.104      0.027 potential strong prior / weak likelihood
\end{verbatim}

\begin{Shaded}
\begin{Highlighting}[]
\CommentTok{\# Density deformation under power{-}scaling shows which parameters are affected.}
\NormalTok{priorsense}\SpecialCharTok{::}\FunctionTok{powerscale\_plot\_dens}\NormalTok{(fit\_single, }\AttributeTok{variables =} \FunctionTok{c}\NormalTok{(}\StringTok{"alpha"}\NormalTok{, }\StringTok{"nu"}\NormalTok{)) }\SpecialCharTok{+}
  \FunctionTok{labs}\NormalTok{(}
    \AttributeTok{title    =} \StringTok{"Prior sensitivity: power{-}scaled posterior densities"}\NormalTok{,}
    \AttributeTok{subtitle =} \FunctionTok{expression}\NormalTok{(}
      \StringTok{"Stronger deformation = greater dependence on the prior"}
\NormalTok{    )}
\NormalTok{  ) }\SpecialCharTok{+}
  \FunctionTok{theme}\NormalTok{(}
    \AttributeTok{strip.background =} \FunctionTok{element\_rect}\NormalTok{(}\AttributeTok{fill =} \StringTok{"grey95"}\NormalTok{, }\AttributeTok{colour =} \ConstantTok{NA}\NormalTok{),}
    \AttributeTok{strip.text       =} \FunctionTok{element\_text}\NormalTok{(}\AttributeTok{face =} \StringTok{"bold"}\NormalTok{),}
    \AttributeTok{legend.position  =} \StringTok{"bottom"}
\NormalTok{  )}
\end{Highlighting}
\end{Shaded}

\begin{center}
\includegraphics[width=0.8\linewidth,height=\textheight,keepaspectratio]{08-MixtureModels_files/figure-pdf/ch09_priorsense-plot-1.pdf}
\end{center}

The expected pattern is instructive: \(\hat{\psi}\) for \(\nu\) (noise,
logit scale) should be substantially higher than \(\hat{\psi}\) for
\(\alpha\) (bias, logit scale). This is the quantitative signature of
the aliasing manifold. The bias \(\theta\) is constrained by the data
through \(\theta_{\text{eff}} \approx \bar{y}\), so the posterior mean
of \(\alpha\) is relatively stable under prior perturbation. The noise
weight \(\pi\), by contrast, receives almost no direct information from
binary i.i.d. data, so its posterior is nearly proportional to its
prior. A strong \(\hat{\psi}\) for \(\nu\) is therefore not a modelling
failure --- it is an empirical confirmation of the §6 identifiability
analysis and should be reported alongside it.

\begin{quote}
\textbf{Workflow note.} If \texttt{powerscale\_sensitivity()} returns a
Pareto-\(\hat{k}\) warning for the prior scaling direction, it signals
that the prior and likelihood are in strong conflict --- the perturbed
posterior is too far from the fitted posterior to be reliably
approximated by importance sampling. In that case, refit with scaled
priors explicitly. For the noise parameter under typical experimental
conditions this warning is unlikely because the prior is doing most of
the work; conflict warnings are more characteristic of well-identified
parameters with informative priors.
\end{quote}

\begin{center}\rule{0.5\linewidth}{0.5pt}\end{center}

\section{Posterior Predictive Checks and
LOO}\label{posterior-predictive-checks-and-loo}

\subsection{Posterior Predictive
Checks}\label{posterior-predictive-checks}

Posterior predictive draws are computed in the
\texttt{generated\ quantities} block and extracted directly from the
fitted object. This avoids the manual simulation loop in the original
chapter, which re-implemented in R what Stan already computed.

\begin{Shaded}
\begin{Highlighting}[]
\CommentTok{\# S x n matrix of posterior predictive choices (0/1)}
\NormalTok{yrep }\OtherTok{\textless{}{-}}\NormalTok{ fit\_single}\SpecialCharTok{$}\FunctionTok{draws}\NormalTok{(}\StringTok{"pred\_choice"}\NormalTok{, }\AttributeTok{format =} \StringTok{"matrix"}\NormalTok{)}
\NormalTok{y    }\OtherTok{\textless{}{-}} \FunctionTok{as.integer}\NormalTok{(data}\SpecialCharTok{$}\NormalTok{h)}

\CommentTok{\# PPC 1: proportion of "1" choices}
\NormalTok{p\_mean }\OtherTok{\textless{}{-}}\NormalTok{ bayesplot}\SpecialCharTok{::}\FunctionTok{ppc\_stat}\NormalTok{(}
  \AttributeTok{y    =}\NormalTok{ y,}
  \AttributeTok{yrep =}\NormalTok{ yrep,}
  \AttributeTok{stat =} \StringTok{"mean"}
\NormalTok{) }\SpecialCharTok{+} \FunctionTok{labs}\NormalTok{(}\AttributeTok{title =} \StringTok{"PPC: proportion of right choices"}\NormalTok{, }\AttributeTok{x =} \StringTok{"Mean choice"}\NormalTok{)}

\CommentTok{\# PPC 2: mean run length — sensitive to serial structure}
\NormalTok{obs\_run  }\OtherTok{\textless{}{-}} \FunctionTok{mean}\NormalTok{(}\FunctionTok{rle}\NormalTok{(y)}\SpecialCharTok{$}\NormalTok{lengths)}
\NormalTok{pred\_run }\OtherTok{\textless{}{-}} \FunctionTok{apply}\NormalTok{(yrep, }\DecValTok{1}\NormalTok{, }\ControlFlowTok{function}\NormalTok{(x) }\FunctionTok{mean}\NormalTok{(}\FunctionTok{rle}\NormalTok{(x)}\SpecialCharTok{$}\NormalTok{lengths))}

\NormalTok{p\_run }\OtherTok{\textless{}{-}} \FunctionTok{ggplot}\NormalTok{(}\FunctionTok{tibble}\NormalTok{(}\AttributeTok{run =}\NormalTok{ pred\_run), }\FunctionTok{aes}\NormalTok{(}\AttributeTok{x =}\NormalTok{ run)) }\SpecialCharTok{+}
  \FunctionTok{geom\_histogram}\NormalTok{(}\FunctionTok{aes}\NormalTok{(}\AttributeTok{y =} \FunctionTok{after\_stat}\NormalTok{(density)),}
                 \AttributeTok{bins =} \DecValTok{40}\NormalTok{, }\AttributeTok{fill =} \StringTok{"\#2166AC"}\NormalTok{, }\AttributeTok{alpha =} \FloatTok{0.7}\NormalTok{) }\SpecialCharTok{+}
  \FunctionTok{geom\_vline}\NormalTok{(}\AttributeTok{xintercept =}\NormalTok{ obs\_run, }\AttributeTok{colour =} \StringTok{"\#D6604D"}\NormalTok{,}
             \AttributeTok{linewidth =} \FloatTok{1.0}\NormalTok{, }\AttributeTok{linetype =} \StringTok{"dashed"}\NormalTok{) }\SpecialCharTok{+}
  \FunctionTok{annotate}\NormalTok{(}\StringTok{"text"}\NormalTok{, }\AttributeTok{x =}\NormalTok{ obs\_run }\SpecialCharTok{+} \FloatTok{0.05}\NormalTok{, }\AttributeTok{y =} \ConstantTok{Inf}\NormalTok{, }\AttributeTok{vjust =} \FloatTok{1.5}\NormalTok{,}
           \AttributeTok{label =} \StringTok{"Observed"}\NormalTok{, }\AttributeTok{colour =} \StringTok{"\#D6604D"}\NormalTok{, }\AttributeTok{size =} \FloatTok{3.5}\NormalTok{, }\AttributeTok{hjust =} \DecValTok{0}\NormalTok{) }\SpecialCharTok{+}
  \FunctionTok{labs}\NormalTok{(}\AttributeTok{title =} \StringTok{"PPC: mean run length"}\NormalTok{,}
       \AttributeTok{subtitle =} \StringTok{"Sensitive to sequential autocorrelation the mean cannot detect"}\NormalTok{,}
       \AttributeTok{x =} \StringTok{"Mean run length"}\NormalTok{, }\AttributeTok{y =} \StringTok{"Density"}\NormalTok{)}

\NormalTok{p\_mean }\SpecialCharTok{/}\NormalTok{ p\_run}
\end{Highlighting}
\end{Shaded}

\begin{center}
\includegraphics[width=0.8\linewidth,height=\textheight,keepaspectratio]{08-MixtureModels_files/figure-pdf/ch09_ppc-1.pdf}
\end{center}

\begin{quote}
\textbf{Interpretation}: If the run-length check fails (observed value
in the tail of the predictive distribution) while the mean check passes,
the data contain sequential dependence that a fully i.i.d. mixture model
cannot reproduce. In that case, consider a hidden Markov model with
Markov-dependent \(z_i\) (Section 10).
\end{quote}

\subsection{\texorpdfstring{PSIS-LOO and Pareto-\(\hat{k}\)
Diagnostics}{PSIS-LOO and Pareto-\textbackslash hat\{k\} Diagnostics}}\label{psis-loo-and-pareto-hatk-diagnostics}

\begin{Shaded}
\begin{Highlighting}[]
\NormalTok{log\_lik\_mat }\OtherTok{\textless{}{-}}\NormalTok{ fit\_single}\SpecialCharTok{$}\FunctionTok{draws}\NormalTok{(}\StringTok{"log\_lik"}\NormalTok{, }\AttributeTok{format =} \StringTok{"matrix"}\NormalTok{)}
\NormalTok{r\_eff       }\OtherTok{\textless{}{-}}\NormalTok{ loo}\SpecialCharTok{::}\FunctionTok{relative\_eff}\NormalTok{(}\FunctionTok{exp}\NormalTok{(log\_lik\_mat),}
                                  \AttributeTok{chain\_id =} \FunctionTok{rep}\NormalTok{(}\DecValTok{1}\SpecialCharTok{:}\DecValTok{4}\NormalTok{, }\AttributeTok{each =} \DecValTok{2000}\NormalTok{))}
\NormalTok{loo\_single  }\OtherTok{\textless{}{-}}\NormalTok{ loo}\SpecialCharTok{::}\FunctionTok{loo}\NormalTok{(log\_lik\_mat, }\AttributeTok{r\_eff =}\NormalTok{ r\_eff, }\AttributeTok{save\_psis =} \ConstantTok{TRUE}\NormalTok{)}

\FunctionTok{print}\NormalTok{(loo\_single)}
\end{Highlighting}
\end{Shaded}

\begin{verbatim}

Computed from 8000 by 120 log-likelihood matrix.

         Estimate   SE
elpd_loo    -65.0  5.3
p_loo         0.9  0.1
looic       130.0 10.5
------
MCSE of elpd_loo is 0.0.
MCSE and ESS estimates assume MCMC draws (r_eff in [0.9, 0.9]).

All Pareto k estimates are good (k < 0.7).
See help('pareto-k-diagnostic') for details.
\end{verbatim}

\begin{Shaded}
\begin{Highlighting}[]
\FunctionTok{plot}\NormalTok{(loo\_single, }\AttributeTok{label\_points =} \ConstantTok{TRUE}\NormalTok{,}
     \AttributeTok{main =} \StringTok{"PSIS{-}LOO Pareto{-}k diagnostics"}\NormalTok{)}
\end{Highlighting}
\end{Shaded}

\begin{center}
\includegraphics[width=0.8\linewidth,height=\textheight,keepaspectratio]{08-MixtureModels_files/figure-pdf/ch09_loo-1.pdf}
\end{center}

Pareto-\(\hat{k}\) values above 0.7 indicate observations for which the
full-data posterior is a poor approximation to the leave-one-out
posterior. In a mixture model, such observations are often concentrated
near \(y = 0.5\) (where the biased and random components are most
confounded) or in long runs of the same choice (which the i.i.d. model
cannot anticipate).

\subsection{LOO-PIT}\label{loo-pit}

For continuous outcomes, LOO-PIT values should be uniform. For binary
outcomes the CDF is a step function, so the standard PIT is discrete and
requires randomisation to achieve uniformity. The
\texttt{rstantools}::loo\_pit()` function applies this randomisation
automatically.

\begin{Shaded}
\begin{Highlighting}[]
\NormalTok{psis\_obj }\OtherTok{\textless{}{-}}\NormalTok{ loo\_single}\SpecialCharTok{$}\NormalTok{psis\_object}
\NormalTok{lw       }\OtherTok{\textless{}{-}} \FunctionTok{weights}\NormalTok{(psis\_obj, }\AttributeTok{log =} \ConstantTok{TRUE}\NormalTok{, }\AttributeTok{normalize =} \ConstantTok{TRUE}\NormalTok{)  }\CommentTok{\# draws × obs}

\NormalTok{yrep\_int }\OtherTok{\textless{}{-}} \FunctionTok{matrix}\NormalTok{(}\FunctionTok{as.integer}\NormalTok{(yrep), }\AttributeTok{nrow =} \FunctionTok{nrow}\NormalTok{(yrep))}
\NormalTok{pit\_vals }\OtherTok{\textless{}{-}}\NormalTok{ rstantools}\SpecialCharTok{::}\FunctionTok{loo\_pit}\NormalTok{(}\AttributeTok{object =}\NormalTok{ yrep\_int, }\AttributeTok{y =}\NormalTok{ y, }\AttributeTok{lw =}\NormalTok{ lw)}
\end{Highlighting}
\end{Shaded}

\begin{center}\rule{0.5\linewidth}{0.5pt}\end{center}

\section{Multilevel Mixture Model}\label{multilevel-mixture-model}

\subsection{Generative Model}\label{generative-model}

When data come from \(J\) participants, we extend the model
hierarchically. Each participant \(j\) has their own bias \(\alpha_j\)
and noise \(\nu_j\), which are drawn from a bivariate normal population
distribution:

\[
\begin{pmatrix} \alpha_j \\ \nu_j \end{pmatrix}
= \begin{pmatrix} \mu_\alpha \\ \mu_\nu \end{pmatrix}
+ \mathbf{L}\,\boldsymbol{z}_j,
\qquad
\boldsymbol{z}_j \sim \text{Normal}_2(\mathbf{0}, \mathbf{I}),
\tag{9}
\]

where \(\mathbf{L} = \text{diag}(\tau) \cdot \mathbf{L}_\Omega\) is the
Cholesky factor of the parameter covariance matrix, \(\tau\) is the
vector of population standard deviations, and \(\mathbf{L}_\Omega\) is
the Cholesky factor of the correlation matrix. This is the non-centred
parameterisation (NCP) introduced in Chapter 7.

\textbf{Priors:}

\[
\mu_\alpha \sim \text{Normal}(0, 1), \qquad
\mu_\nu \sim \text{Normal}(-1, 0.5),
\] \[
\tau_k \sim \text{Half-Normal}(0, 0.3), \quad k = 1, 2,
\qquad
\mathbf{L}_\Omega \sim \text{LKJ}(2).
\]

\textbf{Per-trial likelihood} (same correction as in Section 4.1):

\[
y_{t,j} \sim \log\text{-mix}\!\bigl(\text{logit}^{-1}(\nu_j),\; \text{Bern}(0.5),\; \text{Bern}(\text{logit}^{-1}(\alpha_j))\bigr).
\tag{10}
\]

\subsection{Stan Code}\label{stan-code}

\begin{Shaded}
\begin{Highlighting}[]
\NormalTok{multilevel\_mixture\_stan }\OtherTok{\textless{}{-}} \StringTok{"}
\StringTok{data \{}
\StringTok{  int\textless{}lower=1\textgreater{} trials;}
\StringTok{  int\textless{}lower=1\textgreater{} agents;}
\StringTok{  array[trials, agents] int\textless{}lower=0, upper=1\textgreater{} h;}
\StringTok{  // Priors as data — change without recompiling}
\StringTok{  real prior\_muA\_mu;}
\StringTok{  real\textless{}lower=0\textgreater{} prior\_muA\_sigma;}
\StringTok{  real prior\_muN\_mu;}
\StringTok{  real\textless{}lower=0\textgreater{} prior\_muN\_sigma;}
\StringTok{  real\textless{}lower=0\textgreater{} prior\_tau\_sigma;}
\StringTok{  real\textless{}lower=0\textgreater{} prior\_lkj\_shape;}
\StringTok{  int\textless{}lower=0, upper=1\textgreater{} run\_diagnostics;}
\StringTok{\}}

\StringTok{parameters \{}
\StringTok{  real muA;}
\StringTok{  real muN;}
\StringTok{  vector\textless{}lower=0\textgreater{}[2] tau;}
\StringTok{  matrix[2, agents] z\_IDs;}
\StringTok{  cholesky\_factor\_corr[2] L\_Omega;}
\StringTok{\}}

\StringTok{transformed parameters \{}
\StringTok{  matrix[agents, 2] IDs = (diag\_pre\_multiply(tau, L\_Omega) * z\_IDs)\textquotesingle{};}
\StringTok{\}}

\StringTok{model \{}
\StringTok{  target += normal\_lpdf(muA | prior\_muA\_mu,  prior\_muA\_sigma);}
\StringTok{  target += normal\_lpdf(muN | prior\_muN\_mu,  prior\_muN\_sigma);}
\StringTok{  target += normal\_lpdf(tau | 0, prior\_tau\_sigma);}
\StringTok{  target += lkj\_corr\_cholesky\_lpdf(L\_Omega | prior\_lkj\_shape);}
\StringTok{  target += std\_normal\_lpdf(to\_vector(z\_IDs));}

\StringTok{  for (i in 1:agents) \{}
\StringTok{    real noise\_i = inv\_logit(muN + IDs[i, 2]);}
\StringTok{    real alpha\_i = muA + IDs[i, 1];}
\StringTok{    for (t in 1:trials)}
\StringTok{      target += log\_mix(noise\_i,}
\StringTok{                        bernoulli\_logit\_lpmf(h[t, i] | 0.0),}
\StringTok{                        bernoulli\_logit\_lpmf(h[t, i] | alpha\_i));}
\StringTok{  \}}
\StringTok{\}}

\StringTok{generated quantities \{}
\StringTok{  real\textless{}lower=0, upper=1\textgreater{} pop\_bias  = inv\_logit(muA);}
\StringTok{  real\textless{}lower=0, upper=1\textgreater{} pop\_noise = inv\_logit(muN);}
\StringTok{  corr\_matrix[2] Omega             = multiply\_lower\_tri\_self\_transpose(L\_Omega);}
\StringTok{  real bias\_noise\_corr             = Omega[1, 2];}
\StringTok{  real lprior = normal\_lpdf(muA | prior\_muA\_mu,  prior\_muA\_sigma)}
\StringTok{              + normal\_lpdf(muN | prior\_muN\_mu,  prior\_muN\_sigma)}
\StringTok{              + normal\_lpdf(tau | 0, prior\_tau\_sigma)}
\StringTok{              + lkj\_corr\_cholesky\_lpdf(L\_Omega | prior\_lkj\_shape)}
\StringTok{              + std\_normal\_lpdf(to\_vector(z\_IDs));}

\StringTok{  array[agents] real\textless{}lower=0, upper=1\textgreater{} agent\_bias;}
\StringTok{  array[agents] real\textless{}lower=0, upper=1\textgreater{} agent\_noise;}
\StringTok{  array[trials, agents] real log\_lik;}

\StringTok{  for (i in 1:agents) \{}
\StringTok{    real alpha\_i = muA + IDs[i, 1];}
\StringTok{    real noise\_i = inv\_logit(muN + IDs[i, 2]);}
\StringTok{    agent\_bias[i]  = inv\_logit(alpha\_i);}
\StringTok{    agent\_noise[i] = noise\_i;}
\StringTok{    if (run\_diagnostics) \{}
\StringTok{      for (t in 1:trials)}
\StringTok{        log\_lik[t, i] = log\_mix(noise\_i,}
\StringTok{                                bernoulli\_logit\_lpmf(h[t, i] | 0.0),}
\StringTok{                                bernoulli\_logit\_lpmf(h[t, i] | alpha\_i));}
\StringTok{    \}}
\StringTok{  \}}
\StringTok{\}}
\StringTok{"}

\FunctionTok{write\_stan\_file}\NormalTok{(multilevel\_mixture\_stan,}
                \AttributeTok{dir =} \StringTok{"stan/"}\NormalTok{, }\AttributeTok{basename =} \StringTok{"ch09\_mixture\_multilevel.stan"}\NormalTok{)}
\end{Highlighting}
\end{Shaded}

\begin{verbatim}
[1] "/Users/au209589/Dropbox/Teaching/AdvancedCognitiveModeling23_book/stan/ch09_mixture_multilevel.stan"
\end{verbatim}

\begin{Shaded}
\begin{Highlighting}[]
\NormalTok{mod\_multi }\OtherTok{\textless{}{-}} \FunctionTok{cmdstan\_model}\NormalTok{(}
  \StringTok{"stan/ch09\_mixture\_multilevel.stan"}
\NormalTok{)}
\end{Highlighting}
\end{Shaded}

\subsection{Simulating Multilevel
Data}\label{simulating-multilevel-data}

\begin{Shaded}
\begin{Highlighting}[]
\FunctionTok{set.seed}\NormalTok{(}\DecValTok{801}\NormalTok{)}
\NormalTok{n\_agents }\OtherTok{\textless{}{-}} \DecValTok{30}
\NormalTok{n\_trials }\OtherTok{\textless{}{-}} \DecValTok{120}

\CommentTok{\# Population parameters (logit scale)}
\NormalTok{muA\_true   }\OtherTok{\textless{}{-}} \FunctionTok{qlogis}\NormalTok{(}\FloatTok{0.75}\NormalTok{)   }\CommentTok{\# mean bias ≈ 0.75}
\NormalTok{muN\_true   }\OtherTok{\textless{}{-}} \FunctionTok{qlogis}\NormalTok{(}\FloatTok{0.12}\NormalTok{)   }\CommentTok{\# mean noise ≈ 0.12}
\NormalTok{tauA\_true  }\OtherTok{\textless{}{-}} \FloatTok{0.50}
\NormalTok{tauN\_true  }\OtherTok{\textless{}{-}} \FloatTok{0.40}

\CommentTok{\# Individual parameters (logit scale)}
\NormalTok{alpha\_j }\OtherTok{\textless{}{-}} \FunctionTok{rnorm}\NormalTok{(n\_agents, muA\_true, tauA\_true)}
\NormalTok{nu\_j    }\OtherTok{\textless{}{-}} \FunctionTok{rnorm}\NormalTok{(n\_agents, muN\_true, tauN\_true)}

\CommentTok{\# Generate choice matrix [trials x agents]}
\NormalTok{h\_matrix }\OtherTok{\textless{}{-}} \FunctionTok{matrix}\NormalTok{(}\ConstantTok{NA\_integer\_}\NormalTok{, }\AttributeTok{nrow =}\NormalTok{ n\_trials, }\AttributeTok{ncol =}\NormalTok{ n\_agents)}

\ControlFlowTok{for}\NormalTok{ (j }\ControlFlowTok{in} \FunctionTok{seq\_len}\NormalTok{(n\_agents)) \{}
\NormalTok{  theta\_j }\OtherTok{\textless{}{-}} \FunctionTok{plogis}\NormalTok{(alpha\_j[j])}
\NormalTok{  pi\_j    }\OtherTok{\textless{}{-}} \FunctionTok{plogis}\NormalTok{(nu\_j[j])}
  \ControlFlowTok{for}\NormalTok{ (t }\ControlFlowTok{in} \FunctionTok{seq\_len}\NormalTok{(n\_trials)) \{}
\NormalTok{    h\_matrix[t, j] }\OtherTok{\textless{}{-}} \ControlFlowTok{if}\NormalTok{ (}\FunctionTok{rbinom}\NormalTok{(}\DecValTok{1}\NormalTok{, }\DecValTok{1}\NormalTok{, pi\_j) }\SpecialCharTok{==} \DecValTok{1}\NormalTok{)}
      \FunctionTok{rbinom}\NormalTok{(}\DecValTok{1}\NormalTok{, }\DecValTok{1}\NormalTok{, }\FloatTok{0.5}\NormalTok{)}
    \ControlFlowTok{else}
      \FunctionTok{rbinom}\NormalTok{(}\DecValTok{1}\NormalTok{, }\DecValTok{1}\NormalTok{, theta\_j)}
\NormalTok{  \}}
\NormalTok{\}}

\NormalTok{data\_multi }\OtherTok{\textless{}{-}} \FunctionTok{list}\NormalTok{(}
  \AttributeTok{trials            =}\NormalTok{ n\_trials,}
  \AttributeTok{agents            =}\NormalTok{ n\_agents,}
  \AttributeTok{h                 =}\NormalTok{ h\_matrix,}
  \AttributeTok{prior\_muA\_mu      =} \DecValTok{0}\NormalTok{,    }\AttributeTok{prior\_muA\_sigma   =} \FloatTok{1.0}\NormalTok{,}
  \AttributeTok{prior\_muN\_mu      =} \SpecialCharTok{{-}}\DecValTok{1}\NormalTok{,   }\AttributeTok{prior\_muN\_sigma   =} \FloatTok{0.5}\NormalTok{,}
  \AttributeTok{prior\_tau\_sigma   =} \FloatTok{0.3}\NormalTok{,}
  \AttributeTok{prior\_lkj\_shape   =} \FloatTok{2.0}\NormalTok{,}
  \AttributeTok{run\_diagnostics   =} \DecValTok{1}\NormalTok{L}
\NormalTok{)}

\CommentTok{\# True individual parameters for later recovery plots}
\NormalTok{true\_params }\OtherTok{\textless{}{-}} \FunctionTok{tibble}\NormalTok{(}
  \AttributeTok{agent     =} \FunctionTok{seq\_len}\NormalTok{(n\_agents),}
  \AttributeTok{alpha\_true =}\NormalTok{ alpha\_j,}
  \AttributeTok{nu\_true    =}\NormalTok{ nu\_j,}
  \AttributeTok{theta\_true =} \FunctionTok{plogis}\NormalTok{(alpha\_j),}
  \AttributeTok{pi\_true    =} \FunctionTok{plogis}\NormalTok{(nu\_j),}
  \CommentTok{\# No{-}pooling "estimate": raw data mean per agent}
  \AttributeTok{nopool\_theta =} \FunctionTok{colMeans}\NormalTok{(h\_matrix)}
\NormalTok{)}

\FunctionTok{cat}\NormalTok{(}\StringTok{"Population bias  (prob scale):"}\NormalTok{, }\FunctionTok{round}\NormalTok{(}\FunctionTok{plogis}\NormalTok{(muA\_true), }\DecValTok{3}\NormalTok{), }\StringTok{"}\SpecialCharTok{\textbackslash{}n}\StringTok{"}\NormalTok{)}
\end{Highlighting}
\end{Shaded}

\begin{verbatim}
Population bias  (prob scale): 0.75 
\end{verbatim}

\begin{Shaded}
\begin{Highlighting}[]
\FunctionTok{cat}\NormalTok{(}\StringTok{"Population noise (prob scale):"}\NormalTok{, }\FunctionTok{round}\NormalTok{(}\FunctionTok{plogis}\NormalTok{(muN\_true), }\DecValTok{3}\NormalTok{), }\StringTok{"}\SpecialCharTok{\textbackslash{}n}\StringTok{"}\NormalTok{)}
\end{Highlighting}
\end{Shaded}

\begin{verbatim}
Population noise (prob scale): 0.12 
\end{verbatim}

\subsection{Fitting}\label{fitting}

The same Pathfinder-with-safe-fallback pattern from §4.4 applies here,
with one important difference: the multilevel model has additional
structure that the fallback initialisation must respect. Setting
\texttt{tau} to a small positive value (rather than drawing from the
full half-normal prior) avoids starting near the funnel geometry.
Setting \texttt{z\_IDs} to near-zero ensures all individual parameters
begin at their population means. Setting \texttt{L\_Omega} to the
identity matrix provides a neutral starting correlation. All values are
jittered slightly across chains.

\begin{Shaded}
\begin{Highlighting}[]
\NormalTok{multi\_fit\_file }\OtherTok{\textless{}{-}} \StringTok{"simmodels/ch09\_multilevel\_mixture.RDS"}

\ControlFlowTok{if}\NormalTok{ (regenerate\_simulations }\SpecialCharTok{||} \SpecialCharTok{!}\FunctionTok{file.exists}\NormalTok{(multi\_fit\_file)) \{}

  \CommentTok{\# Explicit per{-}chain jittered initialisation at prior{-}centre values.}
  \CommentTok{\# tau near zero avoids the NCP funnel; z\_IDs = 0 starts all agents}
  \CommentTok{\# at the population mean; L\_Omega = I provides neutral correlation.}
\NormalTok{  safe\_init\_m }\OtherTok{\textless{}{-}} \FunctionTok{lapply}\NormalTok{(}\FunctionTok{seq\_len}\NormalTok{(}\DecValTok{4}\NormalTok{), }\ControlFlowTok{function}\NormalTok{(i) }\FunctionTok{list}\NormalTok{(}
    \AttributeTok{muA     =} \FunctionTok{rnorm}\NormalTok{(}\DecValTok{1}\NormalTok{,  }\FloatTok{0.0}\NormalTok{, }\FloatTok{0.1}\NormalTok{),}
    \AttributeTok{muN     =} \FunctionTok{rnorm}\NormalTok{(}\DecValTok{1}\NormalTok{, }\SpecialCharTok{{-}}\FloatTok{1.0}\NormalTok{, }\FloatTok{0.1}\NormalTok{),}
    \AttributeTok{tau     =} \FunctionTok{abs}\NormalTok{(}\FunctionTok{rnorm}\NormalTok{(}\DecValTok{2}\NormalTok{, }\FloatTok{0.1}\NormalTok{, }\FloatTok{0.02}\NormalTok{)),}
    \AttributeTok{z\_IDs   =} \FunctionTok{matrix}\NormalTok{(}\FunctionTok{rnorm}\NormalTok{(}\DecValTok{2} \SpecialCharTok{*}\NormalTok{ n\_agents, }\DecValTok{0}\NormalTok{, }\FloatTok{0.05}\NormalTok{), }\DecValTok{2}\NormalTok{, n\_agents),}
    \AttributeTok{L\_Omega =} \FunctionTok{diag}\NormalTok{(}\DecValTok{2}\NormalTok{)}
\NormalTok{  ))}

\NormalTok{  init\_pf\_m }\OtherTok{\textless{}{-}} \FunctionTok{tryCatch}\NormalTok{(}
\NormalTok{    mod\_multi}\SpecialCharTok{$}\FunctionTok{pathfinder}\NormalTok{(}\AttributeTok{data =}\NormalTok{ data\_multi, }\AttributeTok{seed =} \DecValTok{801}\NormalTok{,}
                         \AttributeTok{num\_paths =} \DecValTok{4}\NormalTok{, }\AttributeTok{refresh =} \DecValTok{0}\NormalTok{),}
    \AttributeTok{error =} \ControlFlowTok{function}\NormalTok{(e) \{}
      \FunctionTok{message}\NormalTok{(}\StringTok{"Pathfinder failed — falling back to jittered prior{-}centre initialisation."}\NormalTok{)}
\NormalTok{      safe\_init\_m}
\NormalTok{    \}}
\NormalTok{  )}

\NormalTok{  fit\_multi }\OtherTok{\textless{}{-}}\NormalTok{ mod\_multi}\SpecialCharTok{$}\FunctionTok{sample}\NormalTok{(}
    \AttributeTok{data            =}\NormalTok{ data\_multi,}
    \AttributeTok{init            =}\NormalTok{ init\_pf\_m,}
    \AttributeTok{seed            =} \DecValTok{801}\NormalTok{,}
    \AttributeTok{chains          =} \DecValTok{4}\NormalTok{,}
    \AttributeTok{parallel\_chains =} \DecValTok{4}\NormalTok{,}
    \AttributeTok{iter\_warmup     =} \DecValTok{1000}\NormalTok{,}
    \AttributeTok{iter\_sampling   =} \DecValTok{2000}\NormalTok{,}
    \AttributeTok{refresh         =} \DecValTok{500}
\NormalTok{  )}

\NormalTok{  fit\_multi}\SpecialCharTok{$}\FunctionTok{save\_object}\NormalTok{(multi\_fit\_file)}
\NormalTok{\} }\ControlFlowTok{else}\NormalTok{ \{}
\NormalTok{  fit\_multi }\OtherTok{\textless{}{-}} \FunctionTok{readRDS}\NormalTok{(multi\_fit\_file)}
\NormalTok{\}}

\NormalTok{fit\_multi}\SpecialCharTok{$}\FunctionTok{summary}\NormalTok{(}\FunctionTok{c}\NormalTok{(}\StringTok{"pop\_bias"}\NormalTok{, }\StringTok{"pop\_noise"}\NormalTok{, }\StringTok{"tau"}\NormalTok{, }\StringTok{"bias\_noise\_corr"}\NormalTok{))}
\end{Highlighting}
\end{Shaded}

\begin{verbatim}
# A tibble: 5 x 10
  variable       mean median     sd    mad      q5   q95  rhat ess_bulk ess_tail
  <chr>         <dbl>  <dbl>  <dbl>  <dbl>   <dbl> <dbl> <dbl>    <dbl>    <dbl>
1 pop_bias      0.805  0.803 0.0325 0.0312  0.756  0.864  1.00     865.     410.
2 pop_noise     0.205  0.199 0.0637 0.0609  0.111  0.322  1.00     794.     299.
3 tau[1]        0.597  0.590 0.133  0.122   0.398  0.824  1.01     634.     160.
4 tau[2]        0.267  0.220 0.207  0.192   0.0223 0.687  1.01     632.     280.
5 bias_noise_~ -0.119 -0.147 0.433  0.495  -0.785  0.618  1.00    4038.    4740.
\end{verbatim}

\begin{Shaded}
\begin{Highlighting}[]
\NormalTok{diag\_m }\OtherTok{\textless{}{-}}\NormalTok{ fit\_multi}\SpecialCharTok{$}\FunctionTok{diagnostic\_summary}\NormalTok{()}
\FunctionTok{cat}\NormalTok{(}\StringTok{"Divergences:"}\NormalTok{, }\FunctionTok{sum}\NormalTok{(diag\_m}\SpecialCharTok{$}\NormalTok{num\_divergent), }\StringTok{"}\SpecialCharTok{\textbackslash{}n}\StringTok{"}\NormalTok{)}
\end{Highlighting}
\end{Shaded}

\begin{verbatim}
Divergences: 16 
\end{verbatim}

\begin{Shaded}
\begin{Highlighting}[]
\FunctionTok{cat}\NormalTok{(}\StringTok{"Max R{-}hat:  "}\NormalTok{,}
    \FunctionTok{round}\NormalTok{(}\FunctionTok{max}\NormalTok{(fit\_multi}\SpecialCharTok{$}\FunctionTok{summary}\NormalTok{()}\SpecialCharTok{$}\NormalTok{rhat, }\AttributeTok{na.rm =} \ConstantTok{TRUE}\NormalTok{), }\DecValTok{4}\NormalTok{), }\StringTok{"}\SpecialCharTok{\textbackslash{}n}\StringTok{"}\NormalTok{)}
\end{Highlighting}
\end{Shaded}

\begin{verbatim}
Max R-hat:   1.0106 
\end{verbatim}

\begin{Shaded}
\begin{Highlighting}[]
\FunctionTok{cat}\NormalTok{(}\StringTok{"Min ESS:    "}\NormalTok{,}
    \FunctionTok{round}\NormalTok{(}\FunctionTok{min}\NormalTok{(fit\_multi}\SpecialCharTok{$}\FunctionTok{summary}\NormalTok{()}\SpecialCharTok{$}\NormalTok{ess\_bulk, }\AttributeTok{na.rm =} \ConstantTok{TRUE}\NormalTok{)), }\StringTok{"}\SpecialCharTok{\textbackslash{}n}\StringTok{"}\NormalTok{)}
\end{Highlighting}
\end{Shaded}

\begin{verbatim}
Min ESS:     524 
\end{verbatim}

\begin{quote}
\textbf{Debugging note.} If you see divergences in the multilevel
mixture model, the most common cause is the noise SD \(\tau_\nu\)
approaching zero, which creates a funnel geometry in the NCP. Tighten
the half-normal prior on \(\tau_\nu\) (e.g., from
\(\text{Half-Normal}(0, 0.3)\) to \(\text{Half-Normal}(0, 0.2)\)) and
check whether the SBC ECDF-difference for \(\nu\) shows a U-shape
(overconfidence). Do \textbf{not} respond by increasing
\texttt{adapt\_delta} without first investigating the parameterisation.
Note that the safe fallback initialisation in §9.4 already places
\texttt{tau} away from zero; if divergences persist despite that, the
prior itself needs tightening.
\end{quote}

\subsection{Partial-Pooling Shrinkage}\label{partial-pooling-shrinkage}

A hallmark of hierarchical models is that individual estimates are
shrunk toward the group mean, with more shrinkage for agents whose data
are less informative. The shrinkage plot below compares the raw
per-agent data mean (no-pooling) with the posterior mean of
\texttt{agent\_bias{[}j{]}} (partial pooling).

\begin{Shaded}
\begin{Highlighting}[]
\NormalTok{draws\_multi }\OtherTok{\textless{}{-}} \FunctionTok{as\_draws\_df}\NormalTok{(fit\_multi}\SpecialCharTok{$}\FunctionTok{draws}\NormalTok{())}

\CommentTok{\# Extract posterior means of individual bias parameters}
\NormalTok{agent\_cols }\OtherTok{\textless{}{-}} \FunctionTok{paste0}\NormalTok{(}\StringTok{"agent\_bias["}\NormalTok{, }\FunctionTok{seq\_len}\NormalTok{(n\_agents), }\StringTok{"]"}\NormalTok{)}
\NormalTok{post\_bias\_means }\OtherTok{\textless{}{-}}\NormalTok{ draws\_multi }\SpecialCharTok{|\textgreater{}}
\NormalTok{  dplyr}\SpecialCharTok{::}\FunctionTok{select}\NormalTok{(dplyr}\SpecialCharTok{::}\FunctionTok{all\_of}\NormalTok{(agent\_cols)) }\SpecialCharTok{|\textgreater{}}
\NormalTok{  dplyr}\SpecialCharTok{::}\FunctionTok{summarise}\NormalTok{(dplyr}\SpecialCharTok{::}\FunctionTok{across}\NormalTok{(}\FunctionTok{everything}\NormalTok{(), mean)) }\SpecialCharTok{|\textgreater{}}
\NormalTok{  tidyr}\SpecialCharTok{::}\FunctionTok{pivot\_longer}\NormalTok{(}\FunctionTok{everything}\NormalTok{(),}
                      \AttributeTok{names\_to  =} \StringTok{"param"}\NormalTok{,}
                      \AttributeTok{values\_to =} \StringTok{"post\_mean"}\NormalTok{) }\SpecialCharTok{|\textgreater{}}
\NormalTok{  dplyr}\SpecialCharTok{::}\FunctionTok{mutate}\NormalTok{(}\AttributeTok{agent =} \FunctionTok{as.integer}\NormalTok{(}
\NormalTok{    stringr}\SpecialCharTok{::}\FunctionTok{str\_extract}\NormalTok{(param, }\StringTok{"}\SpecialCharTok{\textbackslash{}\textbackslash{}}\StringTok{d+"}\NormalTok{)))}

\NormalTok{shrinkage\_df }\OtherTok{\textless{}{-}}\NormalTok{ true\_params }\SpecialCharTok{|\textgreater{}}
\NormalTok{  dplyr}\SpecialCharTok{::}\FunctionTok{left\_join}\NormalTok{(post\_bias\_means, }\AttributeTok{by =} \StringTok{"agent"}\NormalTok{)}

\NormalTok{pop\_mean\_est }\OtherTok{\textless{}{-}} \FunctionTok{mean}\NormalTok{(draws\_multi}\SpecialCharTok{$}\NormalTok{pop\_bias)}

\FunctionTok{ggplot}\NormalTok{(shrinkage\_df) }\SpecialCharTok{+}
  \FunctionTok{geom\_segment}\NormalTok{(}
    \FunctionTok{aes}\NormalTok{(}\AttributeTok{x =}\NormalTok{ nopool\_theta, }\AttributeTok{xend =}\NormalTok{ post\_mean, }\AttributeTok{y =}\NormalTok{ agent, }\AttributeTok{yend =}\NormalTok{ agent),}
    \AttributeTok{arrow     =} \FunctionTok{arrow}\NormalTok{(}\AttributeTok{length =} \FunctionTok{unit}\NormalTok{(}\FloatTok{0.18}\NormalTok{, }\StringTok{"cm"}\NormalTok{), }\AttributeTok{type =} \StringTok{"closed"}\NormalTok{),}
    \AttributeTok{colour    =} \StringTok{"\#2166AC"}\NormalTok{,}
    \AttributeTok{linewidth =} \FloatTok{0.6}\NormalTok{,}
    \AttributeTok{alpha     =} \FloatTok{0.7}
\NormalTok{  ) }\SpecialCharTok{+}
  \FunctionTok{geom\_point}\NormalTok{(}\FunctionTok{aes}\NormalTok{(}\AttributeTok{x =}\NormalTok{ nopool\_theta, }\AttributeTok{y =}\NormalTok{ agent), }\AttributeTok{shape =} \DecValTok{1}\NormalTok{, }\AttributeTok{size =} \FloatTok{2.5}\NormalTok{) }\SpecialCharTok{+}
  \FunctionTok{geom\_point}\NormalTok{(}\FunctionTok{aes}\NormalTok{(}\AttributeTok{x =}\NormalTok{ post\_mean, }\AttributeTok{y =}\NormalTok{ agent), }\AttributeTok{colour =} \StringTok{"\#2166AC"}\NormalTok{, }\AttributeTok{size =} \FloatTok{2.5}\NormalTok{) }\SpecialCharTok{+}
  \FunctionTok{geom\_point}\NormalTok{(}\FunctionTok{aes}\NormalTok{(}\AttributeTok{x =}\NormalTok{ theta\_true, }\AttributeTok{y =}\NormalTok{ agent), }\AttributeTok{shape =} \DecValTok{4}\NormalTok{,}
             \AttributeTok{colour =} \StringTok{"\#D6604D"}\NormalTok{, }\AttributeTok{size =} \FloatTok{2.5}\NormalTok{) }\SpecialCharTok{+}
  \FunctionTok{geom\_vline}\NormalTok{(}\AttributeTok{xintercept =}\NormalTok{ pop\_mean\_est,}
             \AttributeTok{linetype =} \StringTok{"dashed"}\NormalTok{, }\AttributeTok{colour =} \StringTok{"grey50"}\NormalTok{, }\AttributeTok{linewidth =} \FloatTok{0.8}\NormalTok{) }\SpecialCharTok{+}
  \FunctionTok{labs}\NormalTok{(}
    \AttributeTok{title    =} \StringTok{"Partial{-}pooling shrinkage"}\NormalTok{,}
    \AttributeTok{subtitle =} \StringTok{"Open circles: no{-}pooling (data mean); blue filled: posterior mean; red ×: true value"}\NormalTok{,}
    \AttributeTok{x        =} \StringTok{"Bias parameter (probability scale)"}\NormalTok{,}
    \AttributeTok{y        =} \StringTok{"Agent"}
\NormalTok{  )}
\end{Highlighting}
\end{Shaded}

\begin{center}
\includegraphics[width=0.8\linewidth,height=\textheight,keepaspectratio]{08-MixtureModels_files/figure-pdf/ch09_shrinkage-1.pdf}
\end{center}

Agents whose raw data means are far from the group mean are pulled back
more strongly --- they are more likely to be statistical outliers than
genuinely extreme individuals. Note also that the posterior means are
systematically closer to the true values (red crosses) than the raw data
means, confirming that partial pooling improves estimation accuracy.

\subsection{Population Parameter
Recovery}\label{population-parameter-recovery}

\begin{Shaded}
\begin{Highlighting}[]
\NormalTok{pop\_params }\OtherTok{\textless{}{-}} \FunctionTok{c}\NormalTok{(}\StringTok{"pop\_bias"}\NormalTok{, }\StringTok{"pop\_noise"}\NormalTok{)}
\NormalTok{fit\_multi}\SpecialCharTok{$}\FunctionTok{summary}\NormalTok{(pop\_params) }\SpecialCharTok{|\textgreater{}}
\NormalTok{  dplyr}\SpecialCharTok{::}\FunctionTok{mutate}\NormalTok{(}
    \AttributeTok{true =} \FunctionTok{c}\NormalTok{(}\FunctionTok{plogis}\NormalTok{(muA\_true), }\FunctionTok{plogis}\NormalTok{(muN\_true))}
\NormalTok{  ) }\SpecialCharTok{|\textgreater{}}
\NormalTok{  dplyr}\SpecialCharTok{::}\FunctionTok{select}\NormalTok{(variable, mean, q5, q95, true, rhat, ess\_bulk)}
\end{Highlighting}
\end{Shaded}

\begin{verbatim}
# A tibble: 2 x 7
  variable   mean    q5   q95  true  rhat ess_bulk
  <chr>     <dbl> <dbl> <dbl> <dbl> <dbl>    <dbl>
1 pop_bias  0.805 0.756 0.864  0.75  1.00     865.
2 pop_noise 0.205 0.111 0.322  0.12  1.00     794.
\end{verbatim}

\begin{center}\rule{0.5\linewidth}{0.5pt}\end{center}

\section{Model Comparison}\label{model-comparison}

\subsection{Single-Process Baseline}\label{single-process-baseline}

We fit a simple biased-agent model --- no noise component --- to the
same data used in Section 4. This is the model that emerges when we
constrain \(\pi = 0\).

\begin{Shaded}
\begin{Highlighting}[]
\NormalTok{biased\_stan }\OtherTok{\textless{}{-}} \StringTok{"}
\StringTok{data \{}
\StringTok{  int\textless{}lower=1\textgreater{} n;}
\StringTok{  array[n] int\textless{}lower=0, upper=1\textgreater{} h;}
\StringTok{\}}
\StringTok{parameters \{}
\StringTok{  real alpha;}
\StringTok{\}}
\StringTok{model \{}
\StringTok{  target += normal\_lpdf(alpha | 0, 1);}
\StringTok{  target += bernoulli\_logit\_lpmf(h | alpha);}
\StringTok{\}}
\StringTok{generated quantities \{}
\StringTok{  real\textless{}lower=0, upper=1\textgreater{} theta = inv\_logit(alpha);}
\StringTok{  vector[n] log\_lik;}
\StringTok{  array[n] int pred\_choice;}
\StringTok{  for (i in 1:n) \{}
\StringTok{    log\_lik[i]    = bernoulli\_logit\_lpmf(h[i] | alpha);}
\StringTok{    pred\_choice[i] = bernoulli\_logit\_rng(alpha);}
\StringTok{  \}}
\StringTok{\}}
\StringTok{"}

\FunctionTok{write\_stan\_file}\NormalTok{(biased\_stan, }\AttributeTok{dir =} \StringTok{"stan/"}\NormalTok{, }\AttributeTok{basename =} \StringTok{"ch09\_biased\_single.stan"}\NormalTok{)}
\end{Highlighting}
\end{Shaded}

\begin{verbatim}
[1] "/Users/au209589/Dropbox/Teaching/AdvancedCognitiveModeling23_book/stan/ch09_biased_single.stan"
\end{verbatim}

\begin{Shaded}
\begin{Highlighting}[]
\NormalTok{mod\_biased }\OtherTok{\textless{}{-}} \FunctionTok{cmdstan\_model}\NormalTok{(}
  \StringTok{"stan/ch09\_biased\_single.stan"}
\NormalTok{)}
\end{Highlighting}
\end{Shaded}

\begin{Shaded}
\begin{Highlighting}[]
\NormalTok{biased\_fit\_file }\OtherTok{\textless{}{-}} \StringTok{"simmodels/ch09\_biased\_single.RDS"}

\ControlFlowTok{if}\NormalTok{ (regenerate\_simulations }\SpecialCharTok{||} \SpecialCharTok{!}\FunctionTok{file.exists}\NormalTok{(biased\_fit\_file)) \{}
\NormalTok{  fit\_biased }\OtherTok{\textless{}{-}}\NormalTok{ mod\_biased}\SpecialCharTok{$}\FunctionTok{sample}\NormalTok{(}
    \AttributeTok{data            =}\NormalTok{ data,}
    \AttributeTok{seed            =} \DecValTok{8}\NormalTok{,}
    \AttributeTok{chains          =} \DecValTok{4}\NormalTok{,}
    \AttributeTok{parallel\_chains =} \DecValTok{4}\NormalTok{,}
    \AttributeTok{iter\_warmup     =} \DecValTok{1000}\NormalTok{,}
    \AttributeTok{iter\_sampling   =} \DecValTok{2000}\NormalTok{,}
    \AttributeTok{refresh         =} \DecValTok{500}
\NormalTok{  )}
\NormalTok{  fit\_biased}\SpecialCharTok{$}\FunctionTok{save\_object}\NormalTok{(biased\_fit\_file)}
\NormalTok{\} }\ControlFlowTok{else}\NormalTok{ \{}
\NormalTok{  fit\_biased }\OtherTok{\textless{}{-}} \FunctionTok{readRDS}\NormalTok{(biased\_fit\_file)}
\NormalTok{\}}

\NormalTok{fit\_biased}\SpecialCharTok{$}\FunctionTok{summary}\NormalTok{(}\StringTok{"theta"}\NormalTok{)}
\end{Highlighting}
\end{Shaded}

\begin{verbatim}
# A tibble: 1 x 10
  variable  mean median     sd    mad    q5   q95  rhat ess_bulk ess_tail
  <chr>    <dbl>  <dbl>  <dbl>  <dbl> <dbl> <dbl> <dbl>    <dbl>    <dbl>
1 theta    0.765  0.767 0.0372 0.0376 0.702 0.823  1.00    3196.    4168.
\end{verbatim}

\subsection{ELPD Comparison}\label{elpd-comparison}

\begin{Shaded}
\begin{Highlighting}[]
\NormalTok{log\_lik\_b }\OtherTok{\textless{}{-}}\NormalTok{ fit\_biased}\SpecialCharTok{$}\FunctionTok{draws}\NormalTok{(}\StringTok{"log\_lik"}\NormalTok{, }\AttributeTok{format =} \StringTok{"matrix"}\NormalTok{)}
\NormalTok{r\_eff\_b   }\OtherTok{\textless{}{-}}\NormalTok{ loo}\SpecialCharTok{::}\FunctionTok{relative\_eff}\NormalTok{(}\FunctionTok{exp}\NormalTok{(log\_lik\_b), }\AttributeTok{chain\_id =} \FunctionTok{rep}\NormalTok{(}\DecValTok{1}\SpecialCharTok{:}\DecValTok{4}\NormalTok{, }\AttributeTok{each =} \DecValTok{2000}\NormalTok{))}
\NormalTok{loo\_biased }\OtherTok{\textless{}{-}}\NormalTok{ loo}\SpecialCharTok{::}\FunctionTok{loo}\NormalTok{(log\_lik\_b, }\AttributeTok{r\_eff =}\NormalTok{ r\_eff\_b)}

\NormalTok{comparison }\OtherTok{\textless{}{-}}\NormalTok{ loo}\SpecialCharTok{::}\FunctionTok{loo\_compare}\NormalTok{(loo\_single, loo\_biased)}
\FunctionTok{print}\NormalTok{(comparison)}
\end{Highlighting}
\end{Shaded}

\begin{verbatim}
       elpd_diff se_diff
model2  0.0       0.0   
model1 -0.1       0.3   
\end{verbatim}

\begin{Shaded}
\begin{Highlighting}[]
\NormalTok{weights\_stacking }\OtherTok{\textless{}{-}}\NormalTok{ loo}\SpecialCharTok{::}\FunctionTok{loo\_model\_weights}\NormalTok{(}
  \FunctionTok{list}\NormalTok{(}\AttributeTok{mixture =}\NormalTok{ loo\_single, }\AttributeTok{biased =}\NormalTok{ loo\_biased),}
  \AttributeTok{method =} \StringTok{"stacking"}
\NormalTok{)}
\FunctionTok{print}\NormalTok{(weights\_stacking)}
\end{Highlighting}
\end{Shaded}

\begin{verbatim}
Method: stacking
------
        weight
mixture 0.000 
biased  1.000 
\end{verbatim}

\subsection{Interpretation}\label{interpretation}

The comparison result should show a near-zero ELPD difference with a
standard error larger than the difference itself. This is the correct
outcome, for reasons established in Section 6. Any
\((1-\pi)\text{Bern}(\theta) + \pi\text{Bern}(0.5)\) mixture is
\textbf{likelihood-equivalent} to \(\text{Bern}(\theta_{\text{eff}})\)
for i.i.d. binary data. The biased model, in effect, estimates
\(\hat{\theta}_{\text{eff}}\) directly. PSIS-LOO has no basis on which
to prefer one parameterisation over the other, and Bayesian stacking
weights will similarly be indeterminate. The prior sensitivity analysis
in Section 7 already made this outcome predictable: high \(\hat{\psi}\)
for the noise parameter is the per-fit signal of the same degeneracy
that manifests here as indistinguishable ELPD.

This finding does \emph{not} mean the mixture model is useless. It means
that to identify the components separately, the analysis requires one of
the structural interventions described in Section 6.3: joint modelling
of RT and accuracy, a between-condition predictor on \(\pi\), or
trial-level covariates that predict component membership. In practice,
the mixture model is scientifically justified by theory and produces
richer parameter estimates that are interpretable under the generative
story --- even when LOO cannot disambiguate models, theoretical
grounding should guide model selection.

\begin{center}\rule{0.5\linewidth}{0.5pt}\end{center}

\section{Extensions}\label{extensions}

\textbf{Mixtures with \(K > 2\) components.} The framework extends
naturally: replace the scalar mixing weight \(\pi\) with a \(K\)-simplex
\(\boldsymbol{\pi} \sim \text{Dirichlet}(\boldsymbol{\alpha})\) and loop
over \(K\) terms in \texttt{log\_sum\_exp}. The label-switching problem
(components can permute their identities between chains) becomes more
severe as \(K\) grows; post-processing the draws to enforce a canonical
ordering (e.g., \(\theta_1 < \theta_2 < \ldots\)) or placing informative
asymmetric priors on the components is required.

\textbf{Hidden Markov Models.} If the latent component \(z_i\) is not
i.i.d. but follows a first-order Markov chain, the model is an HMC with
a transition matrix. The forward algorithm marginalises over all state
sequences in \(\mathcal{O}(nK^2)\) time, and is implementable in Stan's
\texttt{functions} block. HMMs are the natural model when behavioural
strategies change gradually across trials, as in online learning or
sequential decision-making tasks.

\textbf{Contamination models for RT data.} The deliberate + lapse
structure in Section 1 is a special case of the \emph{contamination
model} (Ratcliff \& Tuerlinckx 2002), widely used in RT modelling. The
deliberate component is typically an ex-Gaussian or a shifted
log-normal; the contamination component is uniform over a
physiologically plausible range. Because the components have different
support structures, the mixture is well-identified even from marginal RT
data alone --- unlike the binary case.

\textbf{Finite mixture clustering.} When the component means
\(\theta_k\) are themselves unknown parameters, the model becomes a
finite mixture cluster model. Each individual is assigned (softly) to
one of \(K\) latent subgroups with distinct cognitive parameters. This
is a principled alternative to post-hoc clustering of posterior means,
and produces coherent uncertainty over cluster membership.

\begin{center}\rule{0.5\linewidth}{0.5pt}\end{center}

\section{Conclusion}\label{conclusion}

Mixture models are a powerful and principled tool for cognitive
modelling when behaviour reflects more than one underlying process.
Several key lessons from this chapter should inform their use:

\textbf{Discrete parameters must be marginalised before writing Stan
code.} The latent component indicator \(z_i\) cannot be sampled by HMC.
Summing over components using \texttt{log\_mix()} inside a
per-observation loop is not optional. The vectorised form --- passing a
vector to the component density function before calling
\texttt{log\_sum\_exp} --- produces a structurally incorrect likelihood
that models all-or-nothing component assignment rather than per-trial
mixing.

\textbf{Scale consistency inside \texttt{log\_mix()} matters.} Both
component log-likelihoods should be evaluated on the same scale. Use
\texttt{bernoulli\_logit\_lpmf(h{[}i{]}\ \textbar{}\ 0.0)} for a random
component at 0.5, not
\texttt{bernoulli\_lpmf(h{[}i{]}\ \textbar{}\ 0.5)}. The results are
numerically identical but the logit-scale form is consistent, avoids an
implicit transformation, and keeps the model's parameterisation
coherent.

\textbf{Constrained parameter declarations absorb truncation
automatically.} Declaring
\texttt{vector\textless{}lower=0\textgreater{}{[}2{]}\ tau} means Stan
applies the log Jacobian of the constraint transform internally.
Manually subtracting \texttt{normal\_lccdf()} is redundant and adds
unnecessary autodiff nodes. This applies to any constrained parameter:
the truncation correction is only required when placing a truncated
prior on an \emph{unconstrained} quantity.

\textbf{SBC is the right validation tool.} Mixture models are harder to
calibrate than standard regression models because the posterior is often
ridge-shaped and component-specific parameters may be weakly identified.
Running 1,000 SBC datasets and inspecting ECDF-difference plots for both
calibration and parameter recovery is the minimum validation required
before fitting to empirical data. A U-shaped ECDF-difference for the
noise parameter is a common failure mode that signals overconfidence.

\textbf{Prior sensitivity quantifies what identifiability analysis
predicts.} Section 6 showed analytically that \(\pi\) is not identified
from binary i.i.d. data. Section 7's \texttt{priorsense} analysis made
this empirical: a high \(\hat{\psi}\) for the noise logit \(\nu\) and a
low \(\hat{\psi}\) for the bias logit \(\alpha\) is the expected
fingerprint. Reporting both analyses together --- the manifold plot and
the power-scaling sensitivities --- gives the reader a complete picture
of what the data do and do not constrain.

\textbf{Initialisation is not a detail.} Mixture model posteriors can
have ridge or funnel geometry that causes Stan's Uniform(−2, 2) default
initialisation to fail immediately. Use Pathfinder as the primary
initialisation method and define an explicit jittered prior-centre
fallback. For hierarchical models, initialise the population SD vector
\texttt{tau} at small positive values, not at zero or at wide prior
draws.

\textbf{A model that loses a LOO comparison is not necessarily wrong.}
In the binary case, the mixture model and the biased model have
equivalent ELPD because they are parameterisations of the same
likelihood. Theory should guide model choice here, not predictive scores
alone. Model comparison and theoretical justification are complementary,
not competing, sources of evidence. The degeneracy-breaking strategies
in Section 6.3 describe the experimental and modelling decisions that
would, in a future study, make LOO comparison informative.

\textbf{Coming up in Chapter 10:} We leave the matching-pennies domain
entirely and apply the same six-phase workflow to a social learning task
--- the marble task --- where participants integrate direct evidence
about a jar's composition with social information from another person.
This introduces two genuinely new cognitive models (the Simple and
Weighted Bayesian Agents), a Beta-Binomial likelihood, and the
Leave-Future-Out cross-validation technique needed when trial order
creates causal dependencies that standard LOO cannot handle.

\bookmarksetup{startatroot}

\chapter{Reading Other Minds: Theory-of-Mind
Models}\label{sec-theory-of-mind}

\begin{quote}
\textbf{📍 Where we are in the Bayesian modeling workflow:} Chapters
1--6 built, validated, and scaled single-process models of choice
behavior. Chapter 7 taught us how to compare competing theoretical
accounts using out-of-sample predictive performance. In all of that, our
agent reacted to the \emph{statistics} of the environment: a biased
coin, a rate that occasionally reversed, an opponent treated as a
stochastic source to be tracked. We now take the step that actually
makes the environment social: our agent reasons about the \emph{mind}
producing the behavior on the other side of the table. This chapter
develops the \(k\)-level Theory of Mind (\(k\)-ToM) family --- both a
simple 0-ToM tracker and a genuinely recursive 1-ToM --- applies the
full six-phase validation battery from Ch. 5--6 to them, performs an
Occam's-razor model-comparison study in the spirit of Ch. 7, and then
fits a joint hierarchical model to seven primate species (Devaine et
al., 2017) that regresses individual-level ToM parameters on
standardized endocranial volume (ECV) and social group size. Those two
covariates are the empirical stand-ins for the cognitive scaffolding
hypothesis and the Machiavellian intelligence hypothesis, respectively.
\end{quote}

\section{Introduction}\label{introduction-2}

So far every model we have written shares a structural assumption: there
is an environment, it produces observations with some statistical
regularity, and our agent learns that regularity. The Biased Agent
assumed a fixed coin. The Memory Agent tracked a rate that could drift
or reverse. In both, the \emph{opponent} --- if we even called it that
--- was a passive data-generating process.

This assumption breaks the moment we play a competitive game with
another adaptive agent. If I hide a coin and you try to guess where it
is, then the optimal thing for me to do depends on what I think
\emph{you} think I will do. And the optimal thing for you to do depends
on what you think \emph{I} think you will do. The statistics of the
environment are no longer exogenous --- they are the output of a mind
that is trying to anticipate mine.

Theory of Mind (ToM) is the psychology-and-cognitive-science label for
the capacity to represent other agents as having mental states (beliefs,
desires, intentions) and to use those representations to predict their
behavior. The concept is old, verbal, and famously under-specified
(Schaafsma et al., 2015; Quesque \& Rossetti, 2020). What Devaine,
Hollard, and Daunizeau (2014a, 2014b) and the \texttt{tomsup} package
(Waade et al., 2022) contribute is a \emph{computational} definition: a
family of agents, indexed by a recursion depth \(k\), that formally
specifies what it means to mentalize at level 0, 1, 2, and beyond. Armed
with such a definition we can do what we have done in every previous
chapter --- simulate, recover, validate, and compare --- and use those
tools to ask whether a particular dataset is better explained by agents
who mentalize or agents who do not.

\subsection{The estimand}\label{the-estimand}

Before writing a single line of code, let us be precise about what we
are trying to learn. Our target quantity has two layers:

\begin{enumerate}
\def\labelenumi{\arabic{enumi}.}
\tightlist
\item
  \textbf{Individual level.} For each individual (here, each primate),
  what are the posterior distributions over their \(k\)-ToM parameters
  --- volatility \(\sigma\), behavioural temperature \(\beta\), bias
  \(b\) --- when we fit 0-ToM and when we fit 1-ToM? Which model
  predicts that individual's trial-by-trial choices better out of
  sample?
\item
  \textbf{Population level.} Does individual-level behavioural
  temperature (the parameter that most directly controls how strongly an
  agent's expected-utility computation translates into action) vary
  systematically with species-level covariates? Two candidates are
  pre-registered by competing evolutionary hypotheses: social group size
  (\emph{Machiavellian intelligence}) and endocranial volume
  (\emph{cognitive scaffolding}). The only way to arbitrate is to fit
  the individual-level model and then regress the individual-level
  posteriors on both covariates \emph{jointly}, so that each coefficient
  is estimated with the other held constant.
\end{enumerate}

Both layers matter. Fitting \(k\)-ToM to one participant and declaring
them a ``2-ToM'' would be Ch. 2-style estimation: a point estimate that
ignores uncertainty and population structure. Running a single logistic
regression of win rate on ECV would be a pre-Ch. 5 shortcut that throws
away the mechanistic model entirely. The estimand only makes sense if we
carry individual-level posterior uncertainty through to the
species-level inference --- exactly the architecture Ch. 6 built.

\section{Learning Objectives}\label{learning-objectives-8}

After completing this chapter, you will be able to:

\begin{itemize}
\tightlist
\item
  \textbf{Motivate recursive ToM} by contrasting it with statistical
  learning, and articulate why a competitive game with an adaptive
  opponent breaks the passive-environment assumption.
\item
  \textbf{Specify the 0-ToM learning rule} as a Kalman-filter-like
  update on the log-odds of the opponent's choice probability, with
  three parameters: volatility \(\sigma\), behavioural temperature
  \(\beta\), and bias \(b\).
\item
  \textbf{Specify a true recursive 1-ToM} that tracks the opponent's
  belief about the agent itself, and understand why this is
  qualitatively different from influence learning (Hampton et al.,
  2008), which we treat as an intermediate \emph{proto-ToM} reference.
\item
  \textbf{Simulate all three agents in R} against each other and against
  non-ToM baselines.
\item
  \textbf{Implement 0-ToM and 1-ToM in Stan} with the recursive latent
  state unrolled in \texttt{transformed\ parameters}, priors on
  log-scale hyperparameters, and \texttt{log\_lik} for PSIS-LOO.
\item
  \textbf{Execute the six-phase validation battery}: prior predictive
  checks, parameter recovery, Simulation-Based Calibration Checks (SBC),
  posterior predictive checks, prior sensitivity, and between-model
  PSIS-LOO comparison --- including a targeted \emph{Occam's razor
  simulation study} that shows 1-ToM is correctly penalized on data
  generated from simpler mechanisms.
\item
  \textbf{Fit a joint hierarchical model} to the primate dataset that
  regresses individual-level \(\log\beta\) on standardized ECV and group
  size, with non-centered parameterization at both the species and
  individual levels.
\item
  \textbf{Interpret} the regression coefficients as evidence for or
  against the cognitive scaffolding and Machiavellian intelligence
  hypotheses, with appropriate uncertainty and appropriate humility
  about what a 2×2 game can identify.
\end{itemize}

\section{The game and the data}\label{the-game-and-the-data}

The ``hide and seek'' / matching-pennies paradigm is deliberately
minimal. On each trial, one player (the opponent, played by a familiar
zookeeper following the instructions of a computer algorithm) hides a
piece of food in one of two hands. The other player (the primate)
chooses a hand. If the choice matches, the primate wins.

We code the choice binary: \(c^{\text{self}}_t \in \{0, 1\}\) for the
primate's choice, \(c^{\text{op}}_t \in \{0, 1\}\) for the opponent's
hidden-reward location. In matching pennies the primate wants
\(c^{\text{self}}_t = c^{\text{op}}_t\) and the opponent wants the
opposite. The payoff to the primate is \(U(c^{\text{self}},
c^{\text{op}}) = 1\) if they match and \(-1\) if not, flipped for the
opponent.

There are three opponent conditions:

\begin{itemize}
\tightlist
\item
  \textbf{RS (random bias, coded \texttt{-1})}: a 65/35 fixed-bias
  pseudo-random sequence. Baseline. A primate can win by tracking the
  bias.
\item
  \textbf{0-ToM (coded \texttt{0})}: the opponent itself is a 0-ToM
  learner that tracks the primate's choice frequency and hides where the
  primate is \emph{least} likely to look.
\item
  \textbf{1-ToM (coded \texttt{1})}: the opponent models the primate as
  a 0-ToM and best-responds to that simulated agent.
\end{itemize}

\subsection{Loading and
sanity-checking}\label{loading-and-sanity-checking}

\begin{Shaded}
\begin{Highlighting}[]
\NormalTok{mp }\OtherTok{\textless{}{-}} \FunctionTok{read\_csv}\NormalTok{(}
\NormalTok{  here}\SpecialCharTok{::}\FunctionTok{here}\NormalTok{(}\StringTok{"data"}\NormalTok{, }\StringTok{"mp\_primates.csv"}\NormalTok{),}
  \AttributeTok{show\_col\_types =} \ConstantTok{FALSE}
\NormalTok{) }\SpecialCharTok{|\textgreater{}}
\NormalTok{  dplyr}\SpecialCharTok{::}\FunctionTok{select}\NormalTok{(}
    \AttributeTok{id        =}\NormalTok{ ID,}
    \AttributeTok{species   =}\NormalTok{ commonName,}
    \AttributeTok{trial     =}\NormalTok{ Trial,}
    \AttributeTok{choice    =}\NormalTok{ Decision,}
    \AttributeTok{op\_choice =}\NormalTok{ BotDecision,}
    \AttributeTok{payoff    =}\NormalTok{ Payoff,}
    \AttributeTok{opp\_cond  =}\NormalTok{ opponent\_level,}
    \AttributeTok{bot       =}\NormalTok{ BotStrategy,}
\NormalTok{    ECV, neocortex, groupsize\_mean}
\NormalTok{  ) }\SpecialCharTok{|\textgreater{}}
  \FunctionTok{filter}\NormalTok{(}\SpecialCharTok{!}\FunctionTok{is.na}\NormalTok{(choice), }\SpecialCharTok{!}\FunctionTok{is.na}\NormalTok{(op\_choice)) }\SpecialCharTok{|\textgreater{}}
  \FunctionTok{arrange}\NormalTok{(id, opp\_cond, trial) }\SpecialCharTok{|\textgreater{}}
  \FunctionTok{group\_by}\NormalTok{(id, opp\_cond) }\SpecialCharTok{|\textgreater{}}
  \FunctionTok{mutate}\NormalTok{(}\AttributeTok{t =} \FunctionTok{row\_number}\NormalTok{()) }\SpecialCharTok{|\textgreater{}}
  \FunctionTok{ungroup}\NormalTok{()}

\FunctionTok{cat}\NormalTok{(}\StringTok{"Individuals:"}\NormalTok{, }\FunctionTok{n\_distinct}\NormalTok{(mp}\SpecialCharTok{$}\NormalTok{id),}
    \StringTok{"  Species:"}\NormalTok{,   }\FunctionTok{n\_distinct}\NormalTok{(mp}\SpecialCharTok{$}\NormalTok{species),}
    \StringTok{"  Trials:"}\NormalTok{,    }\FunctionTok{nrow}\NormalTok{(mp), }\StringTok{"}\SpecialCharTok{\textbackslash{}n}\StringTok{"}\NormalTok{)}
\end{Highlighting}
\end{Shaded}

\begin{verbatim}
Individuals: 39   Species: 7   Trials: 22644 
\end{verbatim}

A first descriptive pass: how often does each species win in each
condition? Chance is \(0.5\).

\begin{Shaded}
\begin{Highlighting}[]
\NormalTok{perf\_ind }\OtherTok{\textless{}{-}}\NormalTok{ mp }\SpecialCharTok{|\textgreater{}}
  \FunctionTok{mutate}\NormalTok{(}\AttributeTok{win =}\NormalTok{ payoff }\SpecialCharTok{\textgreater{}} \DecValTok{0}\NormalTok{) }\SpecialCharTok{|\textgreater{}}
  \FunctionTok{group\_by}\NormalTok{(species, opp\_cond, id) }\SpecialCharTok{|\textgreater{}}
  \FunctionTok{summarize}\NormalTok{(}\AttributeTok{win\_rate =} \FunctionTok{mean}\NormalTok{(win), }\AttributeTok{.groups =} \StringTok{"drop"}\NormalTok{) }\SpecialCharTok{|\textgreater{}}
  \FunctionTok{mutate}\NormalTok{(}\AttributeTok{opp\_cond =} \FunctionTok{factor}\NormalTok{(opp\_cond,}
                           \AttributeTok{levels =} \FunctionTok{c}\NormalTok{(}\SpecialCharTok{{-}}\DecValTok{1}\NormalTok{, }\DecValTok{0}\NormalTok{, }\DecValTok{1}\NormalTok{),}
                           \AttributeTok{labels =} \FunctionTok{c}\NormalTok{(}\StringTok{"RS (bias)"}\NormalTok{, }\StringTok{"0{-}ToM"}\NormalTok{, }\StringTok{"1{-}ToM"}\NormalTok{)))}

\NormalTok{perf\_spec }\OtherTok{\textless{}{-}}\NormalTok{ perf\_ind }\SpecialCharTok{|\textgreater{}}
  \FunctionTok{group\_by}\NormalTok{(species, opp\_cond) }\SpecialCharTok{|\textgreater{}}
  \FunctionTok{summarize}\NormalTok{(}\AttributeTok{mean\_win =} \FunctionTok{mean}\NormalTok{(win\_rate),}
            \AttributeTok{se       =} \FunctionTok{sd}\NormalTok{(win\_rate) }\SpecialCharTok{/} \FunctionTok{sqrt}\NormalTok{(}\FunctionTok{n}\NormalTok{()),}
            \AttributeTok{.groups  =} \StringTok{"drop"}\NormalTok{)}

\FunctionTok{ggplot}\NormalTok{(perf\_spec, }\FunctionTok{aes}\NormalTok{(}\AttributeTok{x =}\NormalTok{ opp\_cond, }\AttributeTok{y =}\NormalTok{ mean\_win, }\AttributeTok{fill =}\NormalTok{ opp\_cond)) }\SpecialCharTok{+}
  \FunctionTok{geom\_col}\NormalTok{(}\AttributeTok{alpha =} \FloatTok{0.7}\NormalTok{) }\SpecialCharTok{+}
  \FunctionTok{geom\_errorbar}\NormalTok{(}\FunctionTok{aes}\NormalTok{(}\AttributeTok{ymin =}\NormalTok{ mean\_win }\SpecialCharTok{{-}} \FloatTok{1.96}\SpecialCharTok{*}\NormalTok{se,}
                    \AttributeTok{ymax =}\NormalTok{ mean\_win }\SpecialCharTok{+} \FloatTok{1.96}\SpecialCharTok{*}\NormalTok{se), }\AttributeTok{width =} \FloatTok{0.2}\NormalTok{) }\SpecialCharTok{+}
  \FunctionTok{geom\_jitter}\NormalTok{(}\AttributeTok{data =}\NormalTok{ perf\_ind, }\FunctionTok{aes}\NormalTok{(}\AttributeTok{y =}\NormalTok{ win\_rate),}
              \AttributeTok{width =} \FloatTok{0.12}\NormalTok{, }\AttributeTok{height =} \DecValTok{0}\NormalTok{, }\AttributeTok{alpha =} \FloatTok{0.4}\NormalTok{, }\AttributeTok{size =} \DecValTok{1}\NormalTok{) }\SpecialCharTok{+}
  \FunctionTok{geom\_hline}\NormalTok{(}\AttributeTok{yintercept =} \FloatTok{0.5}\NormalTok{, }\AttributeTok{linetype =} \StringTok{"dashed"}\NormalTok{, }\AttributeTok{color =} \StringTok{"gray40"}\NormalTok{) }\SpecialCharTok{+}
  \FunctionTok{facet\_wrap}\NormalTok{(}\SpecialCharTok{\textasciitilde{}}\NormalTok{species, }\AttributeTok{nrow =} \DecValTok{2}\NormalTok{) }\SpecialCharTok{+}
  \FunctionTok{scale\_fill\_manual}\NormalTok{(}\AttributeTok{values =} \FunctionTok{c}\NormalTok{(}\StringTok{"steelblue"}\NormalTok{, }\StringTok{"seagreen"}\NormalTok{, }\StringTok{"firebrick3"}\NormalTok{)) }\SpecialCharTok{+}
  \FunctionTok{labs}\NormalTok{(}\AttributeTok{x =} \ConstantTok{NULL}\NormalTok{, }\AttributeTok{y =} \StringTok{"Win rate"}\NormalTok{,}
       \AttributeTok{title =} \StringTok{"Primate performance by species × opponent"}\NormalTok{,}
       \AttributeTok{subtitle =} \StringTok{"Dashed line = chance"}\NormalTok{, }\AttributeTok{fill =} \StringTok{"Opponent"}\NormalTok{) }\SpecialCharTok{+}
  \FunctionTok{theme}\NormalTok{(}\AttributeTok{axis.text.x =} \FunctionTok{element\_text}\NormalTok{(}\AttributeTok{angle =} \DecValTok{20}\NormalTok{, }\AttributeTok{hjust =} \DecValTok{1}\NormalTok{))}
\end{Highlighting}
\end{Shaded}

\begin{figure}[H]

{\centering \includegraphics[width=0.85\linewidth,height=\textheight,keepaspectratio]{09-TheoryOfMind_files/figure-pdf/ch09_perf_plot-1.pdf}

}

\caption{Per-individual win rates by species and opponent condition.
Points are individuals, bars are 95\% bootstrap CIs of the species mean.
Against RS a bias-tracker can win; against 0-ToM and 1-ToM, performance
clusters around chance because the caregiver-opponent is partially read
as cooperative (Devaine et al., 2017).}

\end{figure}%

Two things are already visible. First, against the biased RS opponent
most species clear chance, confirming they have grasped the
choice-reward contingency. Second, against the two interactive opponents
win rates cluster near or slightly below \(0.5\) --- consistent with
Devaine et al.'s (2017) interpretation that the primates have partially
read the interaction as cooperative and that their mentalizing, if
present, is directed at a different goal than winning. This is a useful
early warning: performance alone cannot tell us whether an individual is
mentalizing. That distinction has to come from the trial-by-trial
\emph{structure} of their choices --- which is exactly what \(k\)-ToM
lets us measure.

Species-level covariates (one row per species):

\begin{Shaded}
\begin{Highlighting}[]
\NormalTok{covariates }\OtherTok{\textless{}{-}}\NormalTok{ mp }\SpecialCharTok{|\textgreater{}}
  \FunctionTok{group\_by}\NormalTok{(species) }\SpecialCharTok{|\textgreater{}}
  \FunctionTok{summarize}\NormalTok{(}
    \AttributeTok{n\_ind     =} \FunctionTok{n\_distinct}\NormalTok{(id),}
    \AttributeTok{ECV       =} \FunctionTok{first}\NormalTok{(ECV),}
    \AttributeTok{groupsize =} \FunctionTok{first}\NormalTok{(groupsize\_mean),}
    \AttributeTok{.groups   =} \StringTok{"drop"}
\NormalTok{  )}
\NormalTok{knitr}\SpecialCharTok{::}\FunctionTok{kable}\NormalTok{(covariates, }\AttributeTok{digits =} \DecValTok{1}\NormalTok{)}
\end{Highlighting}
\end{Shaded}

\begin{longtable}[]{@{}lrrr@{}}
\toprule\noalign{}
species & n\_ind & ECV & groupsize \\
\midrule\noalign{}
\endhead
\bottomrule\noalign{}
\endlastfoot
L macaques & 4 & 80.0 & 30.7 \\
R macaques & 5 & 88.3 & 34.8 \\
chimpanzees & 6 & 357.6 & 43.1 \\
gorillas & 5 & 501.5 & 8.5 \\
lemurs & 4 & 22.9 & 15.0 \\
mangabeys & 8 & 94.7 & 54.1 \\
orangutans & 7 & 349.3 & 1.7 \\
\end{longtable}

The two covariates that will drive the evolutionary regression have very
different profiles: gorillas have large brains (high ECV) but live in
small groups, while mangabeys have modest brains in relatively large
groups. This dissociation is exactly what makes the dataset suitable for
arbitrating between the two hypotheses. The cross-species log-log
correlation between ECV and group size across these seven species is
only about \(-0.37\), so the two predictors are statistically separable.

\section{Model 1: the 0-ToM learner}\label{model-1-the-0-tom-learner}

\subsection{The cognitive intuition}\label{the-cognitive-intuition}

A 0-ToM agent treats its opponent as a biased coin whose bias can slowly
drift. That is almost exactly the Memory Agent from Ch. 6 with a
reversal --- a model we know well. The one new ingredient is a softmax
decision rule that connects the inferred opponent bias back to the
agent's own choice via the payoff matrix.

Let \(\mu_t\) be the agent's current estimate of the log-odds that the
opponent will choose option 1 on trial \(t\), and \(\Sigma_t\) its
uncertainty. After observing \(c^{\text{op}}_t\), the agent applies
Devaine et al.'s (2014a) approximate log-odds Kalman update:

\[
\Sigma_t \;\approx\; \left[\frac{1}{\Sigma_{t-1} + \sigma} + s(\mu_{t-1})\,(1 - s(\mu_{t-1}))\right]^{-1},
\] \[
\mu_t \;\approx\; \mu_{t-1} + \Sigma_t\left(c^{\text{op}}_{t-1} - s(\mu_{t-1})\right),
\]

where \(s(\cdot)\) is the sigmoid and \(\sigma > 0\) is the
\emph{volatility} parameter. The point-estimate prediction for the
opponent's next choice folds uncertainty in:

\[
p^{\text{op}}_t \;\approx\; s\!\left(\frac{\mu_t}{\sqrt{1 + 0.36\,\Sigma_t}}\right).
\]

In matching pennies the expected-utility difference for choosing 1 over
0 collapses to \(\Delta V_t = 2 p^{\text{op}}_t - 1\), and the agent's
own choice probability is a softmax with behavioural temperature
\(\beta > 0\) and bias \(b\):

\[
P(c^{\text{self}}_t = 1) \;=\; s\!\left(\frac{\Delta V_t + b}{\beta}\right).
\]

Three parameters per agent: \(\sigma\) (volatility), \(\beta\)
(decisiveness), \(b\) (handedness). Throughout we work with \(\log
\sigma\) and \(\log \beta\) on the unconstrained scale, and \(b\) in
logit-utility units, consistent with Ch. 6.

\subsection{R simulator}\label{r-simulator}

Because the update is deterministic given the data, we can compute the
whole latent trajectory in vectorized R --- useful for simulation and
for prior predictive checks.

\begin{Shaded}
\begin{Highlighting}[]
\CommentTok{\#\textquotesingle{} Compute 0{-}ToM latent state trajectory for one opponent sequence.}
\NormalTok{tom0\_latent }\OtherTok{\textless{}{-}} \ControlFlowTok{function}\NormalTok{(op\_choices, log\_sigma) \{}
\NormalTok{  sigma }\OtherTok{\textless{}{-}} \FunctionTok{exp}\NormalTok{(log\_sigma)}
\NormalTok{  N }\OtherTok{\textless{}{-}} \FunctionTok{length}\NormalTok{(op\_choices)}
\NormalTok{  mu    }\OtherTok{\textless{}{-}} \FunctionTok{numeric}\NormalTok{(N }\SpecialCharTok{+} \DecValTok{1}\NormalTok{); Sigma }\OtherTok{\textless{}{-}} \FunctionTok{numeric}\NormalTok{(N }\SpecialCharTok{+} \DecValTok{1}\NormalTok{)}
\NormalTok{  mu[}\DecValTok{1}\NormalTok{] }\OtherTok{\textless{}{-}} \DecValTok{0}\NormalTok{; Sigma[}\DecValTok{1}\NormalTok{] }\OtherTok{\textless{}{-}} \DecValTok{1}   \CommentTok{\# agnostic prior (matches tomsup / VBA)}
  \ControlFlowTok{for}\NormalTok{ (t }\ControlFlowTok{in} \FunctionTok{seq\_len}\NormalTok{(N)) \{}
\NormalTok{    s\_mu }\OtherTok{\textless{}{-}} \FunctionTok{plogis}\NormalTok{(mu[t])}
\NormalTok{    Sigma[t }\SpecialCharTok{+} \DecValTok{1}\NormalTok{] }\OtherTok{\textless{}{-}} \DecValTok{1} \SpecialCharTok{/}\NormalTok{ (}\DecValTok{1} \SpecialCharTok{/}\NormalTok{ (Sigma[t] }\SpecialCharTok{+}\NormalTok{ sigma) }\SpecialCharTok{+}\NormalTok{ s\_mu }\SpecialCharTok{*}\NormalTok{ (}\DecValTok{1} \SpecialCharTok{{-}}\NormalTok{ s\_mu))}
\NormalTok{    mu[t }\SpecialCharTok{+} \DecValTok{1}\NormalTok{]    }\OtherTok{\textless{}{-}}\NormalTok{ mu[t] }\SpecialCharTok{+}\NormalTok{ Sigma[t }\SpecialCharTok{+} \DecValTok{1}\NormalTok{] }\SpecialCharTok{*}\NormalTok{ (op\_choices[t] }\SpecialCharTok{{-}}\NormalTok{ s\_mu)}
\NormalTok{  \}}
\NormalTok{  mu\_pred    }\OtherTok{\textless{}{-}}\NormalTok{ mu[}\DecValTok{1}\SpecialCharTok{:}\NormalTok{N]}
\NormalTok{  Sigma\_pred }\OtherTok{\textless{}{-}}\NormalTok{ Sigma[}\DecValTok{1}\SpecialCharTok{:}\NormalTok{N]}
\NormalTok{  p\_op       }\OtherTok{\textless{}{-}} \FunctionTok{plogis}\NormalTok{(mu\_pred }\SpecialCharTok{/} \FunctionTok{sqrt}\NormalTok{(}\DecValTok{1} \SpecialCharTok{+} \FloatTok{0.36} \SpecialCharTok{*}\NormalTok{ Sigma\_pred))}
  \FunctionTok{tibble}\NormalTok{(}\AttributeTok{t =} \FunctionTok{seq\_len}\NormalTok{(N), }\AttributeTok{mu =}\NormalTok{ mu\_pred, }\AttributeTok{Sigma =}\NormalTok{ Sigma\_pred,}
         \AttributeTok{p\_op =}\NormalTok{ p\_op, }\AttributeTok{dV =} \DecValTok{2} \SpecialCharTok{*}\NormalTok{ p\_op }\SpecialCharTok{{-}} \DecValTok{1}\NormalTok{)}
\NormalTok{\}}

\CommentTok{\#\textquotesingle{} Simulate one 0{-}ToM agent against a given opponent choice sequence.}
\NormalTok{simulate\_tom0 }\OtherTok{\textless{}{-}} \ControlFlowTok{function}\NormalTok{(op\_choices, log\_sigma, log\_beta, bias, }\AttributeTok{seed =} \ConstantTok{NULL}\NormalTok{) \{}
  \ControlFlowTok{if}\NormalTok{ (}\SpecialCharTok{!}\FunctionTok{is.null}\NormalTok{(seed)) }\FunctionTok{set.seed}\NormalTok{(seed)}
\NormalTok{  lat }\OtherTok{\textless{}{-}} \FunctionTok{tom0\_latent}\NormalTok{(op\_choices, log\_sigma)}
\NormalTok{  p\_self }\OtherTok{\textless{}{-}} \FunctionTok{plogis}\NormalTok{((lat}\SpecialCharTok{$}\NormalTok{dV }\SpecialCharTok{+}\NormalTok{ bias) }\SpecialCharTok{/} \FunctionTok{exp}\NormalTok{(log\_beta))}
  \FunctionTok{tibble}\NormalTok{(}\AttributeTok{t =}\NormalTok{ lat}\SpecialCharTok{$}\NormalTok{t, }\AttributeTok{choice =} \FunctionTok{rbinom}\NormalTok{(}\FunctionTok{length}\NormalTok{(p\_self), }\DecValTok{1}\NormalTok{, p\_self),}
         \AttributeTok{op\_choice =}\NormalTok{ op\_choices, }\AttributeTok{p\_self =}\NormalTok{ p\_self, }\AttributeTok{p\_op =}\NormalTok{ lat}\SpecialCharTok{$}\NormalTok{p\_op)}
\NormalTok{\}}
\end{Highlighting}
\end{Shaded}

\subsection{Phase 1 --- prior predictive
check}\label{phase-1-prior-predictive-check}

Priors for unbounded parameters that feed through sigmoids are easy to
mis-specify: a ``weakly informative'' \(\mathcal{N}(0, 5)\) on
\texttt{log\_sigma} would place substantial mass at
\(\sigma \approx e^{5}
\approx 150\), which corresponds to an agent whose belief about the
opponent bias resets to the prior on every trial. We therefore choose:

\begin{itemize}
\tightlist
\item
  \texttt{log\_sigma\ \textasciitilde{}\ normal(-2,\ 1)} --- centered at
  \(\sigma \approx 0.14\), the tomsup/VBA default.
\item
  \texttt{log\_beta\ \ \textasciitilde{}\ normal(-1,\ 1)} --- centered
  at \(\beta \approx 0.37\), a mildly deterministic decision rule.
\item
  \texttt{bias\ \ \ \ \ \ \textasciitilde{}\ normal(0,\ 1)}.
\end{itemize}

Do these priors generate plausible choice sequences? Let us draw
parameters from them and simulate.

\begin{Shaded}
\begin{Highlighting}[]
\FunctionTok{set.seed}\NormalTok{(}\DecValTok{2026}\NormalTok{)}
\NormalTok{n\_sims }\OtherTok{\textless{}{-}} \DecValTok{9}
\NormalTok{op\_seq }\OtherTok{\textless{}{-}} \FunctionTok{rbinom}\NormalTok{(}\DecValTok{150}\NormalTok{, }\DecValTok{1}\NormalTok{, }\FloatTok{0.7}\NormalTok{)}

\NormalTok{prior\_draws }\OtherTok{\textless{}{-}} \FunctionTok{tibble}\NormalTok{(}
  \AttributeTok{sim       =} \DecValTok{1}\SpecialCharTok{:}\NormalTok{n\_sims,}
  \AttributeTok{log\_sigma =} \FunctionTok{rnorm}\NormalTok{(n\_sims, }\SpecialCharTok{{-}}\DecValTok{2}\NormalTok{, }\DecValTok{1}\NormalTok{),}
  \AttributeTok{log\_beta  =} \FunctionTok{rnorm}\NormalTok{(n\_sims, }\SpecialCharTok{{-}}\DecValTok{1}\NormalTok{, }\DecValTok{1}\NormalTok{),}
  \AttributeTok{bias      =} \FunctionTok{rnorm}\NormalTok{(n\_sims,  }\DecValTok{0}\NormalTok{, }\DecValTok{1}\NormalTok{)}
\NormalTok{) }\SpecialCharTok{|\textgreater{}}
  \FunctionTok{rowwise}\NormalTok{() }\SpecialCharTok{|\textgreater{}}
  \FunctionTok{mutate}\NormalTok{(}\AttributeTok{sim\_dat =} \FunctionTok{list}\NormalTok{(}\FunctionTok{simulate\_tom0}\NormalTok{(op\_seq, log\_sigma, log\_beta, bias,}
                                      \AttributeTok{seed =}\NormalTok{ sim))) }\SpecialCharTok{|\textgreater{}}
  \FunctionTok{unnest}\NormalTok{(sim\_dat)}

\FunctionTok{ggplot}\NormalTok{(prior\_draws, }\FunctionTok{aes}\NormalTok{(}\AttributeTok{x =}\NormalTok{ t)) }\SpecialCharTok{+}
  \FunctionTok{geom\_step}\NormalTok{(}\FunctionTok{aes}\NormalTok{(}\AttributeTok{y =}\NormalTok{ choice), }\AttributeTok{color =} \StringTok{"firebrick3"}\NormalTok{, }\AttributeTok{alpha =} \FloatTok{0.8}\NormalTok{) }\SpecialCharTok{+}
  \FunctionTok{geom\_line}\NormalTok{(}\FunctionTok{aes}\NormalTok{(}\AttributeTok{y =}\NormalTok{ p\_self), }\AttributeTok{color =} \StringTok{"steelblue"}\NormalTok{, }\AttributeTok{linewidth =} \FloatTok{0.7}\NormalTok{) }\SpecialCharTok{+}
  \FunctionTok{facet\_wrap}\NormalTok{(}\SpecialCharTok{\textasciitilde{}}\NormalTok{sim, }\AttributeTok{labeller =}\NormalTok{ label\_both) }\SpecialCharTok{+}
  \FunctionTok{labs}\NormalTok{(}\AttributeTok{x =} \StringTok{"Trial"}\NormalTok{, }\AttributeTok{y =} \StringTok{"Choice / P(choice=1)"}\NormalTok{,}
       \AttributeTok{title =} \StringTok{"Prior predictive: 0{-}ToM choices across 9 prior draws"}\NormalTok{,}
       \AttributeTok{subtitle =} \StringTok{"Red step = sampled choice, blue = model P(choice=1)"}\NormalTok{)}
\end{Highlighting}
\end{Shaded}

\begin{figure}[H]

{\centering \includegraphics[width=0.85\linewidth,height=\textheight,keepaspectratio]{09-TheoryOfMind_files/figure-pdf/ch09_prior_pc_tom0-1.pdf}

}

\caption{Prior predictive check for 0-ToM. Each panel shows one
simulated agent against a 70/30 biased opponent. Priors should produce
heterogeneous but non-pathological choice sequences.}

\end{figure}%

The simulations show a variety of behaviors --- some agents track the
opponent bias tightly, some are noisy, some show a strong handedness
bias --- and crucially none collapses to a constant 0 or 1 on every
trial, which would indicate a prior forcing boundary behavior. We
proceed.

\subsection{Phase 2 --- Stan implementation (single
agent)}\label{phase-2-stan-implementation-single-agent}

\begin{tcolorbox}[enhanced jigsaw, toprule=.15mm, toptitle=1mm, bottomtitle=1mm, leftrule=.75mm, rightrule=.15mm, colframe=quarto-callout-tip-color-frame, left=2mm, colback=white, titlerule=0mm, coltitle=black, colbacktitle=quarto-callout-tip-color!10!white, opacityback=0, title=\textcolor{quarto-callout-tip-color}{\faLightbulb}\hspace{0.5em}{Stan Technique --- Sequential State Unrolling in
\texttt{transformed\ parameters}}, arc=.35mm, opacitybacktitle=0.6, breakable, bottomrule=.15mm]

The 0-ToM model's belief state \((\mu_t, \Sigma_t)\) is a recursive
function of the previous state and the opponent's most recent choice. It
depends on the free parameter \texttt{log\_sigma} (the volatility), so
it \textbf{cannot} live in \texttt{transformed\ data}. It also cannot be
a local variable in \texttt{model}, because the likelihood at every
trial depends on it and we need Stan's autodiff to propagate gradients
back through the entire sequence.

The solution is to unroll the recursion in
\texttt{transformed\ parameters}:

\begin{Shaded}
\begin{Highlighting}[]
\KeywordTok{transformed parameters}\NormalTok{ \{}
  \DataTypeTok{vector}\NormalTok{[N\_total] dV;          }\CommentTok{// saved for the likelihood and generated quantities}
\NormalTok{  \{}
    \DataTypeTok{real}\NormalTok{ mu = }\FloatTok{0.0}\NormalTok{; }\DataTypeTok{real}\NormalTok{ Sig = }\FloatTok{1.0}\NormalTok{;    }\CommentTok{// local state — not saved between draws}
    \DataTypeTok{real}\NormalTok{ sigma = exp(log\_sigma);}
    \ControlFlowTok{for}\NormalTok{ (t }\ControlFlowTok{in} \DecValTok{1}\NormalTok{:N\_total) \{}
\NormalTok{      dV[t] = f(mu, Sig);             }\CommentTok{// deterministic function of the belief state}
\NormalTok{      Sig   = update\_Sig(Sig, sigma, mu);}
\NormalTok{      mu    = update\_mu(mu, Sig, other[t]);}
\NormalTok{    \}}
\NormalTok{  \}}
\NormalTok{\}}
\end{Highlighting}
\end{Shaded}

The \texttt{\{\}} inner block keeps \texttt{mu} and \texttt{Sig}
\textbf{local} --- they are stack-allocated scratch variables that Stan
discards after each gradient evaluation. Only \texttt{dV} is lifted to
the saved parameter space, because the likelihood needs it. This is the
minimal footprint: the gradient tape tracks
\texttt{log\_sigma\ →\ sigma\ →\ Sig\ →\ mu\ →\ dV{[}t{]}}, but does not
store the full sequence of \((\mu_t, \Sigma_t)\) in the posterior draws.

\textbf{Why \texttt{transformed\ parameters} rather than
\texttt{model}?} We need \texttt{dV} in \texttt{generated\ quantities}
to compute \texttt{log\_lik} and \texttt{y\_rep}. If \texttt{dV} were a
local variable in \texttt{model}, it would be invisible to
\texttt{generated\ quantities}. The solution is to declare it in
\texttt{transformed\ parameters} (where it is visible to all later
blocks) and keep the volatile state variables local inside
\texttt{\{\}}.

\textbf{Computational cost.} Unrolling \(T\) steps in
\texttt{transformed\ parameters} means running the recursion once per
leapfrog evaluation --- the same cost as running it in \texttt{model}.
The only overhead compared to a local-variable approach is that
\texttt{dV} (a vector of length \(T\)) is written to the saved-draw tape
once per posterior sample, adding \(T\) doubles per draw. For
\(T = 120\) this is negligible.

\end{tcolorbox}

We unroll the Kalman update inside \texttt{transformed\ parameters} so
that \texttt{log\_sigma} is a real free parameter (not a fixed
covariate) and is learned jointly with the decision parameters. The
Bernoulli logit likelihood and the \texttt{log\_lik} block are standard.

\begin{Shaded}
\begin{Highlighting}[]
\NormalTok{stan\_tom0\_single }\OtherTok{\textless{}{-}} \StringTok{"}
\StringTok{data \{}
\StringTok{  int\textless{}lower=1\textgreater{} N\_total;}
\StringTok{  array[N\_total] int\textless{}lower=0, upper=1\textgreater{} y;}
\StringTok{  array[N\_total] int\textless{}lower=0, upper=1\textgreater{} other;}
\StringTok{  }
\StringTok{  real prior\_log\_sigma\_mu;}
\StringTok{  real\textless{}lower=0\textgreater{} prior\_log\_sigma\_sigma;}
\StringTok{  real prior\_log\_beta\_mu;}
\StringTok{  real\textless{}lower=0\textgreater{} prior\_log\_beta\_sigma;}
\StringTok{  real prior\_bias\_mu;}
\StringTok{  real\textless{}lower=0\textgreater{} prior\_bias\_sigma;}
\StringTok{  int\textless{}lower=0, upper=1\textgreater{} run\_diagnostics;}
\StringTok{\}}
\StringTok{parameters \{}
\StringTok{  real log\_sigma;}
\StringTok{  real log\_beta;}
\StringTok{  real bias;}
\StringTok{\}}
\StringTok{transformed parameters \{}
\StringTok{  vector[N\_total] dV;}
\StringTok{  \{}
\StringTok{    real sigma = exp(log\_sigma);}
\StringTok{    real mu    = 0.0;}
\StringTok{    real Sig   = 1.0;}
\StringTok{    for (t in 1:N\_total) \{}
\StringTok{      real s\_mu = inv\_logit(mu);}
\StringTok{      real p\_op = inv\_logit(mu / sqrt(1.0 + 0.36 * Sig));}
\StringTok{      dV[t] = 2.0 * p\_op {-} 1.0;}
\StringTok{      Sig = 1.0 / (1.0 / (Sig + sigma) + s\_mu * (1.0 {-} s\_mu));}
\StringTok{      mu  = mu + Sig * (other[t] {-} s\_mu);}
\StringTok{    \}}
\StringTok{  \}}
\StringTok{\}}
\StringTok{model \{}
\StringTok{  target += normal\_lpdf(log\_sigma | prior\_log\_sigma\_mu, prior\_log\_sigma\_sigma);}
\StringTok{  target += normal\_lpdf(log\_beta  | prior\_log\_beta\_mu,  prior\_log\_beta\_sigma);}
\StringTok{  target += normal\_lpdf(bias      | prior\_bias\_mu,      prior\_bias\_sigma);}
\StringTok{  target += bernoulli\_logit\_lpmf(y | (dV + bias) / exp(log\_beta));}
\StringTok{\}}
\StringTok{generated quantities \{}
\StringTok{  vector[N\_total] log\_lik;}
\StringTok{  array[N\_total] int y\_rep;}
\StringTok{  real lprior =   normal\_lpdf(log\_sigma | prior\_log\_sigma\_mu, prior\_log\_sigma\_sigma)}
\StringTok{                + normal\_lpdf(log\_beta  | prior\_log\_beta\_mu,  prior\_log\_beta\_sigma)}
\StringTok{                + normal\_lpdf(bias      | prior\_bias\_mu,      prior\_bias\_sigma);}
\StringTok{  if (run\_diagnostics) \{}
\StringTok{    for (t in 1:N\_total) \{}
\StringTok{      log\_lik[t] = bernoulli\_logit\_lpmf(y[t] | (dV[t] + bias) / exp(log\_beta));}
\StringTok{      y\_rep[t]   = bernoulli\_logit\_rng((dV[t] + bias) / exp(log\_beta));}
\StringTok{    \}}
\StringTok{  \}}
\StringTok{\}}
\StringTok{"}
\FunctionTok{writeLines}\NormalTok{(stan\_tom0\_single, here}\SpecialCharTok{::}\FunctionTok{here}\NormalTok{(}\StringTok{"stan"}\NormalTok{, }\StringTok{"ch09\_tom0\_single.stan"}\NormalTok{))}
\NormalTok{mod\_tom0 }\OtherTok{\textless{}{-}} \FunctionTok{cmdstan\_model}\NormalTok{(here}\SpecialCharTok{::}\FunctionTok{here}\NormalTok{(}\StringTok{"stan"}\NormalTok{, }\StringTok{"ch09\_tom0\_single.stan"}\NormalTok{))}
\end{Highlighting}
\end{Shaded}

\subsection{Phase 3 --- parameter
recovery}\label{phase-3-parameter-recovery}

A 20-agent recovery study with 120 trials each, which matches the
average trial count we see per individual per condition in the primate
dataset.

\begin{Shaded}
\begin{Highlighting}[]
\NormalTok{recovery\_path }\OtherTok{\textless{}{-}}\NormalTok{ here}\SpecialCharTok{::}\FunctionTok{here}\NormalTok{(}\StringTok{"simmodels"}\NormalTok{, }\StringTok{"ch09\_tom0\_recovery.rds"}\NormalTok{)}

\ControlFlowTok{if}\NormalTok{ (regenerate\_simulations }\SpecialCharTok{||} \SpecialCharTok{!}\FunctionTok{file.exists}\NormalTok{(recovery\_path)) \{}
  \FunctionTok{set.seed}\NormalTok{(}\DecValTok{42}\NormalTok{)}
\NormalTok{  truth }\OtherTok{\textless{}{-}} \FunctionTok{tibble}\NormalTok{(}
    \AttributeTok{agent     =} \DecValTok{1}\SpecialCharTok{:}\DecValTok{20}\NormalTok{,}
    \AttributeTok{log\_sigma =} \FunctionTok{rnorm}\NormalTok{(}\DecValTok{20}\NormalTok{, }\SpecialCharTok{{-}}\DecValTok{2}\NormalTok{, }\DecValTok{1}\NormalTok{),}
    \AttributeTok{log\_beta  =} \FunctionTok{rnorm}\NormalTok{(}\DecValTok{20}\NormalTok{, }\SpecialCharTok{{-}}\DecValTok{1}\NormalTok{, }\DecValTok{1}\NormalTok{),}
    \AttributeTok{bias      =} \FunctionTok{rnorm}\NormalTok{(}\DecValTok{20}\NormalTok{,  }\DecValTok{0}\NormalTok{, }\DecValTok{1}\NormalTok{)}
\NormalTok{  )}
  \CommentTok{\# Every synthetic agent plays the same 120{-}trial biased opponent}
\NormalTok{  op\_vec }\OtherTok{\textless{}{-}} \FunctionTok{rbinom}\NormalTok{(}\DecValTok{120}\NormalTok{, }\DecValTok{1}\NormalTok{, }\FloatTok{0.65}\NormalTok{)}

\NormalTok{  recovery }\OtherTok{\textless{}{-}}\NormalTok{ truth }\SpecialCharTok{|\textgreater{}}
    \FunctionTok{rowwise}\NormalTok{() }\SpecialCharTok{|\textgreater{}}
    \FunctionTok{mutate}\NormalTok{(}\AttributeTok{sim\_dat =} \FunctionTok{list}\NormalTok{(}\FunctionTok{simulate\_tom0}\NormalTok{(op\_vec, log\_sigma, log\_beta,}
\NormalTok{                                        bias, }\AttributeTok{seed =}\NormalTok{ agent))) }\SpecialCharTok{|\textgreater{}}
    \FunctionTok{ungroup}\NormalTok{()}

\NormalTok{  fit\_one }\OtherTok{\textless{}{-}} \ControlFlowTok{function}\NormalTok{(d) \{}
\NormalTok{    fit }\OtherTok{\textless{}{-}}\NormalTok{ mod\_tom0}\SpecialCharTok{$}\FunctionTok{sample}\NormalTok{(}
      \AttributeTok{data =} \FunctionTok{list}\NormalTok{(}
        \AttributeTok{N\_total =} \FunctionTok{nrow}\NormalTok{(d), }
        \AttributeTok{y =}\NormalTok{ d}\SpecialCharTok{$}\NormalTok{choice, }
        \AttributeTok{other =}\NormalTok{ d}\SpecialCharTok{$}\NormalTok{op\_choice,}
        \AttributeTok{prior\_log\_sigma\_mu =} \SpecialCharTok{{-}}\FloatTok{2.0}\NormalTok{, }\AttributeTok{prior\_log\_sigma\_sigma =} \FloatTok{1.0}\NormalTok{,}
        \AttributeTok{prior\_log\_beta\_mu =} \SpecialCharTok{{-}}\FloatTok{1.0}\NormalTok{,  }\AttributeTok{prior\_log\_beta\_sigma =} \FloatTok{1.0}\NormalTok{,}
        \AttributeTok{prior\_bias\_mu =} \FloatTok{0.0}\NormalTok{,       }\AttributeTok{prior\_bias\_sigma =} \FloatTok{1.0}\NormalTok{,}
        \AttributeTok{run\_diagnostics =} \DecValTok{0}
\NormalTok{      ),}
      \AttributeTok{chains =} \DecValTok{2}\NormalTok{, }\AttributeTok{parallel\_chains =} \DecValTok{2}\NormalTok{,}
      \AttributeTok{iter\_warmup =} \DecValTok{500}\NormalTok{, }\AttributeTok{iter\_sampling =} \DecValTok{500}\NormalTok{,}
      \AttributeTok{refresh =} \DecValTok{0}\NormalTok{, }\AttributeTok{show\_messages =} \ConstantTok{FALSE}
\NormalTok{    )}
\NormalTok{    fit}\SpecialCharTok{$}\FunctionTok{summary}\NormalTok{(}\FunctionTok{c}\NormalTok{(}\StringTok{"log\_sigma"}\NormalTok{, }\StringTok{"log\_beta"}\NormalTok{, }\StringTok{"bias"}\NormalTok{)) }\SpecialCharTok{|\textgreater{}}
\NormalTok{      dplyr}\SpecialCharTok{::}\FunctionTok{select}\NormalTok{(variable, mean, q5, q95)}
\NormalTok{  \}}

\NormalTok{  recovery\_summary }\OtherTok{\textless{}{-}}\NormalTok{ recovery }\SpecialCharTok{|\textgreater{}}
    \FunctionTok{rowwise}\NormalTok{() }\SpecialCharTok{|\textgreater{}}
    \FunctionTok{mutate}\NormalTok{(}\AttributeTok{post =} \FunctionTok{list}\NormalTok{(}\FunctionTok{fit\_one}\NormalTok{(sim\_dat))) }\SpecialCharTok{|\textgreater{}}
\NormalTok{    dplyr}\SpecialCharTok{::}\FunctionTok{select}\NormalTok{(agent, log\_sigma, log\_beta, bias, post) }\SpecialCharTok{|\textgreater{}}
    \FunctionTok{unnest}\NormalTok{(post)}
  \FunctionTok{saveRDS}\NormalTok{(recovery\_summary, recovery\_path)}
\NormalTok{\} }\ControlFlowTok{else}\NormalTok{ \{}
\NormalTok{  recovery\_summary }\OtherTok{\textless{}{-}} \FunctionTok{readRDS}\NormalTok{(recovery\_path)}
\NormalTok{\}}

\NormalTok{recovery\_long }\OtherTok{\textless{}{-}}\NormalTok{ recovery\_summary }\SpecialCharTok{|\textgreater{}}
  \FunctionTok{pivot\_longer}\NormalTok{(}\FunctionTok{c}\NormalTok{(log\_sigma, log\_beta, bias),}
               \AttributeTok{names\_to =} \StringTok{"param"}\NormalTok{, }\AttributeTok{values\_to =} \StringTok{"true"}\NormalTok{) }\SpecialCharTok{|\textgreater{}}
  \FunctionTok{filter}\NormalTok{(variable }\SpecialCharTok{==}\NormalTok{ param)}
\end{Highlighting}
\end{Shaded}

\begin{Shaded}
\begin{Highlighting}[]
\FunctionTok{ggplot}\NormalTok{(recovery\_long, }\FunctionTok{aes}\NormalTok{(}\AttributeTok{x =}\NormalTok{ true, }\AttributeTok{y =}\NormalTok{ mean)) }\SpecialCharTok{+}
  \FunctionTok{geom\_errorbar}\NormalTok{(}\FunctionTok{aes}\NormalTok{(}\AttributeTok{ymin =}\NormalTok{ q5, }\AttributeTok{ymax =}\NormalTok{ q95), }\AttributeTok{width =} \DecValTok{0}\NormalTok{, }\AttributeTok{alpha =} \FloatTok{0.4}\NormalTok{,}
                \AttributeTok{color =} \StringTok{"steelblue"}\NormalTok{) }\SpecialCharTok{+}
  \FunctionTok{geom\_point}\NormalTok{(}\AttributeTok{color =} \StringTok{"midnightblue"}\NormalTok{, }\AttributeTok{size =} \DecValTok{2}\NormalTok{) }\SpecialCharTok{+}
  \FunctionTok{geom\_abline}\NormalTok{(}\AttributeTok{slope =} \DecValTok{1}\NormalTok{, }\AttributeTok{intercept =} \DecValTok{0}\NormalTok{, }\AttributeTok{linetype =} \StringTok{"dashed"}\NormalTok{,}
              \AttributeTok{color =} \StringTok{"gray40"}\NormalTok{) }\SpecialCharTok{+}
  \FunctionTok{facet\_wrap}\NormalTok{(}\SpecialCharTok{\textasciitilde{}}\NormalTok{param, }\AttributeTok{scales =} \StringTok{"free"}\NormalTok{) }\SpecialCharTok{+}
  \FunctionTok{labs}\NormalTok{(}\AttributeTok{x =} \StringTok{"True value"}\NormalTok{, }\AttributeTok{y =} \StringTok{"Posterior mean"}\NormalTok{,}
       \AttributeTok{title =} \StringTok{"0{-}ToM parameter recovery (20 synthetic agents, 120 trials)"}\NormalTok{)}
\end{Highlighting}
\end{Shaded}

\begin{figure}[H]

{\centering \includegraphics[width=0.85\linewidth,height=\textheight,keepaspectratio]{09-TheoryOfMind_files/figure-pdf/ch09_recovery_tom0_plot-1.pdf}

}

\caption{0-ToM parameter recovery. Points = posterior means, bars = 90\%
CIs, dashed line = identity. Recovery is clean for bias and log\_beta;
log\_sigma is noisier with mild shrinkage to the prior at 120 trials ---
a design finding we will confirm under the more rigorous SBC next.}

\end{figure}%

\subsection{Phase 4 --- Simulation-Based Calibration Checks
(SBC)}\label{phase-4-simulation-based-calibration-checks-sbc}

Parameter recovery on a single dataset is anecdotal: we may have sampled
a benign region of parameter space. SBC (Talts et al., 2018) asks the
stronger question --- across the \emph{entire prior} --- is the
posterior calibrated? If we generate \(\theta_{\text{true}}\) from the
prior, simulate data from the data model, and compute the rank of
\(\theta_{\text{true}}\) in the posterior draws, that rank must be
uniform. A U-shaped rank histogram means overconfidence; a bump-shape
means underconfidence; any asymmetry means bias.

\begin{Shaded}
\begin{Highlighting}[]
\NormalTok{sbc\_path }\OtherTok{\textless{}{-}}\NormalTok{ here}\SpecialCharTok{::}\FunctionTok{here}\NormalTok{(}\StringTok{"simmodels"}\NormalTok{, }\StringTok{"ch09\_tom0\_sbc.rds"}\NormalTok{)}

\NormalTok{generate\_sbc\_tom0 }\OtherTok{\textless{}{-}} \ControlFlowTok{function}\NormalTok{(}\AttributeTok{N =} \DecValTok{120}\NormalTok{, }\AttributeTok{op\_bias =} \FloatTok{0.65}\NormalTok{) \{}
\NormalTok{  log\_sigma }\OtherTok{\textless{}{-}} \FunctionTok{rnorm}\NormalTok{(}\DecValTok{1}\NormalTok{, }\SpecialCharTok{{-}}\DecValTok{2}\NormalTok{, }\DecValTok{1}\NormalTok{)}
\NormalTok{  log\_beta  }\OtherTok{\textless{}{-}} \FunctionTok{rnorm}\NormalTok{(}\DecValTok{1}\NormalTok{, }\SpecialCharTok{{-}}\DecValTok{1}\NormalTok{, }\DecValTok{1}\NormalTok{)}
\NormalTok{  bias      }\OtherTok{\textless{}{-}} \FunctionTok{rnorm}\NormalTok{(}\DecValTok{1}\NormalTok{,  }\DecValTok{0}\NormalTok{, }\DecValTok{1}\NormalTok{)}
\NormalTok{  op\_vec    }\OtherTok{\textless{}{-}} \FunctionTok{rbinom}\NormalTok{(N, }\DecValTok{1}\NormalTok{, op\_bias)}
\NormalTok{  sim\_dat   }\OtherTok{\textless{}{-}} \FunctionTok{simulate\_tom0}\NormalTok{(op\_vec, log\_sigma, log\_beta, bias)}
  \FunctionTok{list}\NormalTok{(}
    \AttributeTok{variables =} \FunctionTok{list}\NormalTok{(}\AttributeTok{log\_sigma =}\NormalTok{ log\_sigma,}
                     \AttributeTok{log\_beta  =}\NormalTok{ log\_beta,}
                     \AttributeTok{bias      =}\NormalTok{ bias),}
    \AttributeTok{generated =} \FunctionTok{list}\NormalTok{(}\AttributeTok{N =}\NormalTok{ N,}
                     \AttributeTok{choice    =}\NormalTok{ sim\_dat}\SpecialCharTok{$}\NormalTok{choice,}
                     \AttributeTok{op\_choice =}\NormalTok{ sim\_dat}\SpecialCharTok{$}\NormalTok{op\_choice)}
\NormalTok{  )}
\NormalTok{\}}

\ControlFlowTok{if}\NormalTok{ (regenerate\_sbc }\SpecialCharTok{||} \SpecialCharTok{!}\FunctionTok{file.exists}\NormalTok{(sbc\_path)) \{}
\NormalTok{  n\_sbc }\OtherTok{\textless{}{-}} \DecValTok{150}   \CommentTok{\# 200+ preferred in production; 150 for chapter compute}
\NormalTok{  sbc\_gen }\OtherTok{\textless{}{-}} \FunctionTok{SBC\_generator\_function}\NormalTok{(generate\_sbc\_tom0, }\AttributeTok{N =} \DecValTok{120}\NormalTok{, }\AttributeTok{op\_bias =} \FloatTok{0.65}\NormalTok{)}
\NormalTok{  sbc\_backend }\OtherTok{\textless{}{-}} \FunctionTok{SBC\_backend\_cmdstan\_sample}\NormalTok{(}
\NormalTok{    mod\_tom0,}
    \AttributeTok{iter\_warmup =} \DecValTok{500}\NormalTok{, }\AttributeTok{iter\_sampling =} \DecValTok{500}\NormalTok{, }\AttributeTok{chains =} \DecValTok{1}\NormalTok{, }\AttributeTok{refresh =} \DecValTok{0}
\NormalTok{  )}
  \FunctionTok{plan}\NormalTok{(sequential)}
\NormalTok{  sbc\_ds      }\OtherTok{\textless{}{-}} \FunctionTok{generate\_datasets}\NormalTok{(sbc\_gen, n\_sbc)}
\NormalTok{  sbc\_results }\OtherTok{\textless{}{-}} \FunctionTok{compute\_SBC}\NormalTok{(}\AttributeTok{datasets =}\NormalTok{ sbc\_ds, }\AttributeTok{backend =}\NormalTok{ sbc\_backend,}
                             \AttributeTok{keep\_fits =} \ConstantTok{FALSE}\NormalTok{)}
  \FunctionTok{saveRDS}\NormalTok{(}\FunctionTok{list}\NormalTok{(}\AttributeTok{ds =}\NormalTok{ sbc\_ds, }\AttributeTok{results =}\NormalTok{ sbc\_results), sbc\_path)}
\NormalTok{\} }\ControlFlowTok{else}\NormalTok{ \{}
\NormalTok{  sbc\_obj     }\OtherTok{\textless{}{-}} \FunctionTok{readRDS}\NormalTok{(sbc\_path)}
\NormalTok{  sbc\_ds      }\OtherTok{\textless{}{-}}\NormalTok{ sbc\_obj}\SpecialCharTok{$}\NormalTok{ds}
\NormalTok{  sbc\_results }\OtherTok{\textless{}{-}}\NormalTok{ sbc\_obj}\SpecialCharTok{$}\NormalTok{results}
\NormalTok{\}}
\end{Highlighting}
\end{Shaded}

\begin{Shaded}
\begin{Highlighting}[]
\NormalTok{p\_hist }\OtherTok{\textless{}{-}} \FunctionTok{plot\_rank\_hist}\NormalTok{(sbc\_results,}
                         \AttributeTok{variables =} \FunctionTok{c}\NormalTok{(}\StringTok{"log\_sigma"}\NormalTok{, }\StringTok{"log\_beta"}\NormalTok{, }\StringTok{"bias"}\NormalTok{)) }\SpecialCharTok{+}
  \FunctionTok{labs}\NormalTok{(}\AttributeTok{subtitle =} \StringTok{"Rank histograms"}\NormalTok{)}
\NormalTok{p\_ecdf }\OtherTok{\textless{}{-}} \FunctionTok{plot\_ecdf\_diff}\NormalTok{(sbc\_results,}
                         \AttributeTok{variables =} \FunctionTok{c}\NormalTok{(}\StringTok{"log\_sigma"}\NormalTok{, }\StringTok{"log\_beta"}\NormalTok{, }\StringTok{"bias"}\NormalTok{)) }\SpecialCharTok{+}
  \FunctionTok{labs}\NormalTok{(}\AttributeTok{subtitle =} \StringTok{"ECDF difference"}\NormalTok{)}
\NormalTok{p\_hist }\SpecialCharTok{/}\NormalTok{ p\_ecdf}
\end{Highlighting}
\end{Shaded}

\begin{figure}[H]

{\centering \includegraphics[width=0.85\linewidth,height=\textheight,keepaspectratio]{09-TheoryOfMind_files/figure-pdf/ch09_sbc_tom0_plot-1.pdf}

}

\caption{Simulation-Based Calibration Checks for 0-ToM. Top row: rank
histograms (flat = calibrated). Bottom row: ECDF difference (blue line
inside grey band = calibrated). log\_sigma shows the mildest excursion,
consistent with the recovery plot: its posterior is only partially
data-driven at 120 trials; log\_beta and bias are clean.}

\end{figure}%

The SBC results confirm the recovery diagnosis: \texttt{log\_beta} and
\texttt{bias} are cleanly calibrated, \texttt{log\_sigma} shows the
signature of mild underidentification with 120 trials per sequence. We
will therefore interpret species-level posteriors on \texttt{log\_sigma}
with explicit caution, and place the scientific weight on
\texttt{log\_beta}.

\section{Model 2: influence learning
(proto-ToM)}\label{model-2-influence-learning-proto-tom}

The simplest step beyond 0-ToM is to notice that your opponent is
\emph{reacting} to you. Hampton, Bossaerts, and O'Doherty (2008)
formalized this as \textbf{influence learning}: augment the 0-ToM
prediction error with a term that captures how your own current choice
will shift the opponent's next prediction of you.

For a competitive matching-pennies opponent:

\[
p^{\text{op}}_{t+1} = p^{\text{op}}_t
+ \eta \bigl( c^{\text{op}}_t - p^{\text{op}}_t \bigr)
- \lambda\, p^{\text{op}}_t(1 - p^{\text{op}}_t)\Bigl(2 c^{\text{self}}_t - \beta_{\text{op}}\,\text{logit}(p^{\text{op}}_t) - 1\Bigr),
\]

where \(\eta\) is the baseline learning rate, \(\lambda\) is the
influence weight, and the bracketed factor encodes the competitive
incentive. Importantly, there is \emph{no} recursive belief updating:
the agent does not maintain a posterior over the opponent's parameters,
it only adjusts a heuristic rule. Devaine et al.~(2017) classify
influence learning as proto-ToM precisely for this reason --- it is
aware of self-influence without mentalizing per se.

We implement influence learning as a parallel simulator and Stan
observation model. The code structure mirrors 0-ToM; we include the R
simulator here and omit the Stan duplicate for space (it differs only in
the \texttt{transformed\ parameters} block: the scalar \(p^{\text{op}}\)
update equation replaces the Kalman block). A working parallel
\texttt{ch09\_influence\_single.stan} is included in the chapter's
supplementary repository.

\begin{Shaded}
\begin{Highlighting}[]
\NormalTok{simulate\_inf }\OtherTok{\textless{}{-}} \ControlFlowTok{function}\NormalTok{(op\_choices, log\_eta, log\_lambda, log\_beta, bias,}
                         \AttributeTok{beta\_op =} \FloatTok{0.4}\NormalTok{, }\AttributeTok{seed =} \ConstantTok{NULL}\NormalTok{) \{}
  \ControlFlowTok{if}\NormalTok{ (}\SpecialCharTok{!}\FunctionTok{is.null}\NormalTok{(seed)) }\FunctionTok{set.seed}\NormalTok{(seed)}
\NormalTok{  eta    }\OtherTok{\textless{}{-}} \FunctionTok{plogis}\NormalTok{(log\_eta)     }\CommentTok{\# eta, lambda bounded in (0,1)}
\NormalTok{  lambda }\OtherTok{\textless{}{-}} \FunctionTok{plogis}\NormalTok{(log\_lambda)}
\NormalTok{  beta   }\OtherTok{\textless{}{-}} \FunctionTok{exp}\NormalTok{(log\_beta)}
\NormalTok{  N   }\OtherTok{\textless{}{-}} \FunctionTok{length}\NormalTok{(op\_choices)}
\NormalTok{  p\_op     }\OtherTok{\textless{}{-}} \FunctionTok{numeric}\NormalTok{(N }\SpecialCharTok{+} \DecValTok{1}\NormalTok{); p\_op[}\DecValTok{1}\NormalTok{] }\OtherTok{\textless{}{-}} \FloatTok{0.5}
\NormalTok{  choice   }\OtherTok{\textless{}{-}} \FunctionTok{integer}\NormalTok{(N)}
  \ControlFlowTok{for}\NormalTok{ (t }\ControlFlowTok{in} \FunctionTok{seq\_len}\NormalTok{(N)) \{}
\NormalTok{    dV }\OtherTok{\textless{}{-}} \DecValTok{2} \SpecialCharTok{*}\NormalTok{ p\_op[t] }\SpecialCharTok{{-}} \DecValTok{1}
\NormalTok{    p\_self\_t }\OtherTok{\textless{}{-}} \FunctionTok{plogis}\NormalTok{((dV }\SpecialCharTok{+}\NormalTok{ bias) }\SpecialCharTok{/}\NormalTok{ beta)}
\NormalTok{    choice[t] }\OtherTok{\textless{}{-}} \FunctionTok{rbinom}\NormalTok{(}\DecValTok{1}\NormalTok{, }\DecValTok{1}\NormalTok{, p\_self\_t)}
    \CommentTok{\# Influence update for next{-}trial prediction}
\NormalTok{    pe      }\OtherTok{\textless{}{-}}\NormalTok{ op\_choices[t] }\SpecialCharTok{{-}}\NormalTok{ p\_op[t]}
\NormalTok{    logit\_p }\OtherTok{\textless{}{-}} \FunctionTok{qlogis}\NormalTok{(}\FunctionTok{pmin}\NormalTok{(}\FunctionTok{pmax}\NormalTok{(p\_op[t], }\FloatTok{1e{-}6}\NormalTok{), }\DecValTok{1} \SpecialCharTok{{-}} \FloatTok{1e{-}6}\NormalTok{))}
\NormalTok{    infl    }\OtherTok{\textless{}{-}}\NormalTok{ lambda }\SpecialCharTok{*}\NormalTok{ p\_op[t] }\SpecialCharTok{*}\NormalTok{ (}\DecValTok{1} \SpecialCharTok{{-}}\NormalTok{ p\_op[t]) }\SpecialCharTok{*}
\NormalTok{               (}\DecValTok{2} \SpecialCharTok{*}\NormalTok{ choice[t] }\SpecialCharTok{{-}}\NormalTok{ beta\_op }\SpecialCharTok{*}\NormalTok{ logit\_p }\SpecialCharTok{{-}} \DecValTok{1}\NormalTok{)}
\NormalTok{    p\_op[t }\SpecialCharTok{+} \DecValTok{1}\NormalTok{] }\OtherTok{\textless{}{-}} \FunctionTok{pmin}\NormalTok{(}\FunctionTok{pmax}\NormalTok{(p\_op[t] }\SpecialCharTok{+}\NormalTok{ eta }\SpecialCharTok{*}\NormalTok{ pe }\SpecialCharTok{{-}}\NormalTok{ infl, }\FloatTok{1e{-}4}\NormalTok{), }\DecValTok{1} \SpecialCharTok{{-}} \FloatTok{1e{-}4}\NormalTok{)}
\NormalTok{  \}}
  \FunctionTok{tibble}\NormalTok{(}\AttributeTok{t =} \FunctionTok{seq\_len}\NormalTok{(N), }\AttributeTok{choice =}\NormalTok{ choice, }\AttributeTok{op\_choice =}\NormalTok{ op\_choices,}
         \AttributeTok{p\_op =}\NormalTok{ p\_op[}\DecValTok{1}\SpecialCharTok{:}\NormalTok{N])}
\NormalTok{\}}
\end{Highlighting}
\end{Shaded}

\section{Model 3: true recursive
1-ToM}\label{model-3-true-recursive-1-tom}

Now the real thing. A 1-ToM agent believes its opponent is a 0-ToM
agent, and \emph{represents the opponent's belief about itself}. That
recursive representation is what elevates 1-ToM above influence
learning: the agent does not just adjust for the opponent's reactivity
with a heuristic, it simulates the opponent's cognitive update
internally.

\subsection{The mechanism}\label{the-mechanism}

Let us denote 1-ToM's belief state about \emph{how the opponent views
the 1-ToM agent itself} by
\((\mu^{\text{self}}, \Sigma^{\text{self}})\). This pair evolves exactly
like a 0-ToM's belief, but the ``observed choice'' that drives the
update is the 1-ToM agent's \emph{own} past choice \(c^{\text{self}}_t\)
(because that is what the opponent observes):

\[
\Sigma^{\text{self}}_t \approx \left[\frac{1}{\Sigma^{\text{self}}_{t-1} + \sigma_{\text{op}}} + s(\mu^{\text{self}}_{t-1})(1 - s(\mu^{\text{self}}_{t-1}))\right]^{-1},
\] \[
\mu^{\text{self}}_t \approx \mu^{\text{self}}_{t-1} + \Sigma^{\text{self}}_t \bigl( c^{\text{self}}_{t-1} - s(\mu^{\text{self}}_{t-1}) \bigr).
\]

Here \(\sigma_{\text{op}}\) is the \emph{opponent's} volatility --- as
believed by the 1-ToM. Given this state, the 1-ToM predicts what the
opponent will do on trial \(t\) by simulating the opponent's decision
rule. The opponent (as a 0-ToM) wants to \emph{mismatch} the primate, so
its expected utility of choosing 1 is
\(\Delta V^{\text{op}}_t = 1 - 2 \hat{p}^{\text{self}}_t\), and its
choice probability is:

\[
\hat{p}^{\text{op}}_t = s\!\left(\frac{1 - 2 \hat{p}^{\text{self}}_t}{\beta_{\text{op}}}\right),
\qquad
\hat{p}^{\text{self}}_t = s\!\left(\frac{\mu^{\text{self}}_t}{\sqrt{1 + 0.36\,\Sigma^{\text{self}}_t}}\right).
\]

The 1-ToM agent then plugs \(\hat{p}^{\text{op}}_t\) into its own
expected-utility calculation and softmaxes to produce a choice. Four
parameters: \(\log \sigma_{\text{op}}\) (how volatile the 1-ToM thinks
its opponent thinks it is), \(\log \beta_{\text{op}}\) (opponent's
decision temperature as believed by 1-ToM), \(\log \beta\) (1-ToM's own
decision temperature), and \(b\) (bias).

Note the recursive structure: this is qualitatively different from
influence learning. 1-ToM maintains a \emph{full belief state} about
what the opponent thinks of it, with uncertainty, and updates that
belief via a Bayesian filter. Influence learning updates a scalar
\(p^{\text{op}}\) with a correction term. The two are behaviorally close
in many scenarios (Devaine et al., 2017, Fig. 7), but mechanistically
distinct --- and the Occam's-razor test below shows that data can
sometimes distinguish them.

\subsection{What this implementation commits
to}\label{what-this-implementation-commits-to}

A candid methodological note before we proceed. Our 1-ToM agent tracks
the first two moments \((\mu^{\text{self}}, \Sigma^{\text{self}})\) of
its belief about the opponent's belief about itself, \emph{conditional
on point values} of the opponent's volatility \(\sigma_{\text{op}}\) and
decision temperature \(\beta_{\text{op}}\). The agent does not represent
uncertainty over those opponent parameters. The full variational scheme
of Devaine et al.~(2014a) maintains a Gaussian variational posterior
over \((\sigma_{\text{op}},
\beta_{\text{op}})\) as well, and propagates that second-order
uncertainty through the recursion via moment matching. That is the more
complete cognitive commitment and also the more expensive one.

We are therefore implementing what is most accurately called a
\textbf{fixed-parameter 1-ToM}: an agent that believes its opponent has
fixed, known \((\sigma_{\text{op}}, \beta_{\text{op}})\) and updates
only its beliefs about the opponent's moment-to-moment \emph{state}
\(\mu^{\text{self}}\). The posterior we will obtain from Stan over
\((\log\sigma_{\text{op}}, \log\beta_{\text{op}})\) reflects \emph{our}
uncertainty over which fixed values best describe this particular
primate's internal opponent model, \emph{not} the primate's own
representational uncertainty about its opponent. This is a simpler
cognitive model than full variational 1-ToM, and it is what is actually
tractable inside an HMC sampler without a specialized message-passing
backend. Readers who want the Devaine et al.~scheme in full should
consult the \texttt{VBA} toolbox (Daunizeau, 2014).

\textbf{Graceful degradation.} One concern with the point-estimate
approximation is how sensitive 1-ToM predictions are to
mis-specification of \(\sigma_{\text{op}}\). The recursion has a stable
fixed point \(\Sigma^\star\) set by the balance of drift
\(\sigma_{\text{op}}\) and the sigmoid second moment (bounded above by
\(1/4\)). For \(\sigma_{\text{op}}\) ranging from \(10^{-3}\) to
\(10^1\), \(\Sigma^\star\) lies in roughly \([0.1, 2]\), and the
predicted \(\hat{p}^{\text{op}}_t\) is a compressed, monotone transform
of the underlying \(\mu\)-trajectory. The consequence is that a tenfold
error in the internal \(\sigma_{\text{op}}\) produces a bounded and
usually small change in \(dV_t\) --- a feature from a robustness
standpoint, but a \emph{bug} from an identifiability standpoint, because
it means the data will not constrain this parameter strongly. We return
to this in Phase 4b.

\subsection{R simulator}\label{r-simulator-1}

\begin{Shaded}
\begin{Highlighting}[]
\CommentTok{\#\textquotesingle{} Compute 1{-}ToM\textquotesingle{}s predicted p\_op sequence given agent parameters AND}
\CommentTok{\#\textquotesingle{} given the agent\textquotesingle{}s own past choices. Used by simulator and by Stan.}
\NormalTok{tom1\_p\_op }\OtherTok{\textless{}{-}} \ControlFlowTok{function}\NormalTok{(choices\_self, log\_sigma\_op, log\_beta\_op) \{}
\NormalTok{  sigma\_op }\OtherTok{\textless{}{-}} \FunctionTok{exp}\NormalTok{(log\_sigma\_op)}
\NormalTok{  beta\_op  }\OtherTok{\textless{}{-}} \FunctionTok{exp}\NormalTok{(log\_beta\_op)}
\NormalTok{  N  }\OtherTok{\textless{}{-}} \FunctionTok{length}\NormalTok{(choices\_self)}
\NormalTok{  mu\_s    }\OtherTok{\textless{}{-}} \DecValTok{0}\NormalTok{; Sig\_s }\OtherTok{\textless{}{-}} \DecValTok{1}
\NormalTok{  p\_op    }\OtherTok{\textless{}{-}} \FunctionTok{numeric}\NormalTok{(N)}
  \ControlFlowTok{for}\NormalTok{ (t }\ControlFlowTok{in} \FunctionTok{seq\_len}\NormalTok{(N)) \{}
\NormalTok{    p\_self\_from\_op }\OtherTok{\textless{}{-}} \FunctionTok{plogis}\NormalTok{(mu\_s }\SpecialCharTok{/} \FunctionTok{sqrt}\NormalTok{(}\DecValTok{1} \SpecialCharTok{+} \FloatTok{0.36} \SpecialCharTok{*}\NormalTok{ Sig\_s))}
\NormalTok{    p\_op[t] }\OtherTok{\textless{}{-}} \FunctionTok{plogis}\NormalTok{((}\DecValTok{1} \SpecialCharTok{{-}} \DecValTok{2} \SpecialCharTok{*}\NormalTok{ p\_self\_from\_op) }\SpecialCharTok{/}\NormalTok{ beta\_op)}
    \CommentTok{\# Update opponent\textquotesingle{}s belief about self, using self\textquotesingle{}s past choice}
\NormalTok{    s\_ps  }\OtherTok{\textless{}{-}}\NormalTok{ p\_self\_from\_op}
\NormalTok{    Sig\_s }\OtherTok{\textless{}{-}} \DecValTok{1} \SpecialCharTok{/}\NormalTok{ (}\DecValTok{1} \SpecialCharTok{/}\NormalTok{ (Sig\_s }\SpecialCharTok{+}\NormalTok{ sigma\_op) }\SpecialCharTok{+}\NormalTok{ s\_ps }\SpecialCharTok{*}\NormalTok{ (}\DecValTok{1} \SpecialCharTok{{-}}\NormalTok{ s\_ps))}
\NormalTok{    mu\_s  }\OtherTok{\textless{}{-}}\NormalTok{ mu\_s }\SpecialCharTok{+}\NormalTok{ Sig\_s }\SpecialCharTok{*}\NormalTok{ (choices\_self[t] }\SpecialCharTok{{-}}\NormalTok{ s\_ps)}
\NormalTok{  \}}
\NormalTok{  p\_op}
\NormalTok{\}}

\NormalTok{simulate\_tom1 }\OtherTok{\textless{}{-}} \ControlFlowTok{function}\NormalTok{(op\_choices, log\_sigma\_op, log\_beta\_op,}
\NormalTok{                          log\_beta, bias, }\AttributeTok{seed =} \ConstantTok{NULL}\NormalTok{) \{}
  \ControlFlowTok{if}\NormalTok{ (}\SpecialCharTok{!}\FunctionTok{is.null}\NormalTok{(seed)) }\FunctionTok{set.seed}\NormalTok{(seed)}
\NormalTok{  beta }\OtherTok{\textless{}{-}} \FunctionTok{exp}\NormalTok{(log\_beta)}
\NormalTok{  N  }\OtherTok{\textless{}{-}} \FunctionTok{length}\NormalTok{(op\_choices)}
\NormalTok{  mu\_s }\OtherTok{\textless{}{-}} \DecValTok{0}\NormalTok{; Sig\_s }\OtherTok{\textless{}{-}} \DecValTok{1}
\NormalTok{  choice }\OtherTok{\textless{}{-}} \FunctionTok{integer}\NormalTok{(N); p\_op\_vec }\OtherTok{\textless{}{-}} \FunctionTok{numeric}\NormalTok{(N)}
  \ControlFlowTok{for}\NormalTok{ (t }\ControlFlowTok{in} \FunctionTok{seq\_len}\NormalTok{(N)) \{}
\NormalTok{    p\_self\_from\_op }\OtherTok{\textless{}{-}} \FunctionTok{plogis}\NormalTok{(mu\_s }\SpecialCharTok{/} \FunctionTok{sqrt}\NormalTok{(}\DecValTok{1} \SpecialCharTok{+} \FloatTok{0.36} \SpecialCharTok{*}\NormalTok{ Sig\_s))}
\NormalTok{    p\_op\_t         }\OtherTok{\textless{}{-}} \FunctionTok{plogis}\NormalTok{((}\DecValTok{1} \SpecialCharTok{{-}} \DecValTok{2} \SpecialCharTok{*}\NormalTok{ p\_self\_from\_op) }\SpecialCharTok{/} \FunctionTok{exp}\NormalTok{(log\_beta\_op))}
\NormalTok{    p\_op\_vec[t]    }\OtherTok{\textless{}{-}}\NormalTok{ p\_op\_t}
\NormalTok{    dV\_t           }\OtherTok{\textless{}{-}} \DecValTok{2} \SpecialCharTok{*}\NormalTok{ p\_op\_t }\SpecialCharTok{{-}} \DecValTok{1}
\NormalTok{    p\_self\_t       }\OtherTok{\textless{}{-}} \FunctionTok{plogis}\NormalTok{((dV\_t }\SpecialCharTok{+}\NormalTok{ bias) }\SpecialCharTok{/}\NormalTok{ beta)}
\NormalTok{    choice[t]      }\OtherTok{\textless{}{-}} \FunctionTok{rbinom}\NormalTok{(}\DecValTok{1}\NormalTok{, }\DecValTok{1}\NormalTok{, p\_self\_t)}
    \CommentTok{\# Update after observing own choice}
\NormalTok{    s\_ps }\OtherTok{\textless{}{-}}\NormalTok{ p\_self\_from\_op}
\NormalTok{    Sig\_s }\OtherTok{\textless{}{-}} \DecValTok{1} \SpecialCharTok{/}\NormalTok{ (}\DecValTok{1} \SpecialCharTok{/}\NormalTok{ (Sig\_s }\SpecialCharTok{+} \FunctionTok{exp}\NormalTok{(log\_sigma\_op)) }\SpecialCharTok{+}\NormalTok{ s\_ps }\SpecialCharTok{*}\NormalTok{ (}\DecValTok{1} \SpecialCharTok{{-}}\NormalTok{ s\_ps))}
\NormalTok{    mu\_s  }\OtherTok{\textless{}{-}}\NormalTok{ mu\_s }\SpecialCharTok{+}\NormalTok{ Sig\_s }\SpecialCharTok{*}\NormalTok{ (choice[t] }\SpecialCharTok{{-}}\NormalTok{ s\_ps)}
\NormalTok{  \}}
  \FunctionTok{tibble}\NormalTok{(}\AttributeTok{t =} \FunctionTok{seq\_len}\NormalTok{(N), }\AttributeTok{choice =}\NormalTok{ choice, }\AttributeTok{op\_choice =}\NormalTok{ op\_choices,}
         \AttributeTok{p\_op =}\NormalTok{ p\_op\_vec)}
\NormalTok{\}}
\end{Highlighting}
\end{Shaded}

\subsection{Stan implementation}\label{stan-implementation}

The Stan model unrolls the same recursion inside
\texttt{transformed\ parameters}. A subtle but important point: the
opponent's belief state depends on \texttt{choice{[}t{]}}, which is
\emph{data}, not a parameter. This keeps the recursion within a single
leapfrog step cheap.

\begin{Shaded}
\begin{Highlighting}[]
\NormalTok{stan\_tom1\_single }\OtherTok{\textless{}{-}} \StringTok{"}
\StringTok{data \{}
\StringTok{  int\textless{}lower=1\textgreater{} N;}
\StringTok{  array[N] int\textless{}lower=0, upper=1\textgreater{} choice;}
\StringTok{  array[N] int\textless{}lower=0, upper=1\textgreater{} op\_choice;}
\StringTok{\}}
\StringTok{parameters \{}
\StringTok{  real log\_sigma\_op;}
\StringTok{  real log\_beta\_op;}
\StringTok{  real log\_beta;}
\StringTok{  real bias;}
\StringTok{\}}
\StringTok{transformed parameters \{}
\StringTok{  vector[N] dV;}
\StringTok{  \{}
\StringTok{    real sigma\_op = exp(log\_sigma\_op);}
\StringTok{    real beta\_op  = exp(log\_beta\_op);}
\StringTok{    real mu\_s     = 0.0;}
\StringTok{    real Sig\_s    = 1.0;}
\StringTok{    for (t in 1:N) \{}
\StringTok{      real p\_self\_from\_op = inv\_logit(mu\_s / sqrt(1 + 0.36 * Sig\_s));}
\StringTok{      real p\_op           = inv\_logit((1 {-} 2 * p\_self\_from\_op) / beta\_op);}
\StringTok{      dV[t] = 2 * p\_op {-} 1;}
\StringTok{      // After observing self\textquotesingle{}s choice, opponent updates its belief}
\StringTok{      real s\_ps = p\_self\_from\_op;}
\StringTok{      Sig\_s = 1.0 / (1.0 / (Sig\_s + sigma\_op) + s\_ps * (1 {-} s\_ps));}
\StringTok{      mu\_s  = mu\_s + Sig\_s * (choice[t] {-} s\_ps);}
\StringTok{    \}}
\StringTok{  \}}
\StringTok{\}}
\StringTok{model \{}
\StringTok{  log\_sigma\_op \textasciitilde{} normal({-}2, 1);}
\StringTok{  log\_beta\_op  \textasciitilde{} normal({-}1, 0.7);}
\StringTok{  log\_beta     \textasciitilde{} normal({-}1, 1);}
\StringTok{  bias         \textasciitilde{} normal( 0, 1);}
\StringTok{  choice \textasciitilde{} bernoulli\_logit((dV + bias) / exp(log\_beta));}
\StringTok{\}}
\StringTok{generated quantities \{}
\StringTok{  vector[N] log\_lik;}
\StringTok{  real lprior =   normal\_lpdf(log\_sigma\_op | {-}2, 1)}
\StringTok{                + normal\_lpdf(log\_beta\_op  | {-}1, 0.7)}
\StringTok{                + normal\_lpdf(log\_beta     | {-}1, 1)}
\StringTok{                + normal\_lpdf(bias         |  0, 1);}
\StringTok{  for (t in 1:N)}
\StringTok{    log\_lik[t] = bernoulli\_logit\_lpmf(choice[t] |}
\StringTok{                  (dV[t] + bias) / exp(log\_beta));}
\StringTok{\}}
\StringTok{"}
\FunctionTok{writeLines}\NormalTok{(stan\_tom1\_single, here}\SpecialCharTok{::}\FunctionTok{here}\NormalTok{(}\StringTok{"stan"}\NormalTok{, }\StringTok{"ch09\_tom1\_single.stan"}\NormalTok{))}
\NormalTok{mod\_tom1 }\OtherTok{\textless{}{-}} \FunctionTok{cmdstan\_model}\NormalTok{(here}\SpecialCharTok{::}\FunctionTok{here}\NormalTok{(}\StringTok{"stan"}\NormalTok{, }\StringTok{"ch09\_tom1\_single.stan"}\NormalTok{))}
\end{Highlighting}
\end{Shaded}

Two remarks on priors. First, \texttt{log\_beta\_op} is held somewhat
tighter than the primate's own \texttt{log\_beta} because this parameter
represents the primate's \emph{belief} about the opponent's decisiveness
and is identified only indirectly through predictions --- a tight prior
avoids fishing in regions the data cannot constrain. Second,
\texttt{log\_sigma\_op} and \texttt{log\_beta\_op} enter the likelihood
only through the recursion, making them the least well-identified
parameters. Do not be surprised if their posteriors are close to the
prior on short sequences.

\subsection{1-ToM recovery}\label{tom-recovery}

\begin{Shaded}
\begin{Highlighting}[]
\NormalTok{recovery\_tom1\_path }\OtherTok{\textless{}{-}}\NormalTok{ here}\SpecialCharTok{::}\FunctionTok{here}\NormalTok{(}\StringTok{"simmodels"}\NormalTok{, }\StringTok{"ch09\_tom1\_recovery.rds"}\NormalTok{)}

\ControlFlowTok{if}\NormalTok{ (regenerate\_simulations }\SpecialCharTok{||} \SpecialCharTok{!}\FunctionTok{file.exists}\NormalTok{(recovery\_tom1\_path)) \{}
  \FunctionTok{set.seed}\NormalTok{(}\DecValTok{2027}\NormalTok{)}
\NormalTok{  truth1 }\OtherTok{\textless{}{-}} \FunctionTok{tibble}\NormalTok{(}
    \AttributeTok{agent        =} \DecValTok{1}\SpecialCharTok{:}\DecValTok{15}\NormalTok{,}
    \AttributeTok{log\_sigma\_op =} \FunctionTok{rnorm}\NormalTok{(}\DecValTok{15}\NormalTok{, }\SpecialCharTok{{-}}\DecValTok{2}\NormalTok{, }\DecValTok{1}\NormalTok{),}
    \AttributeTok{log\_beta\_op  =} \FunctionTok{rnorm}\NormalTok{(}\DecValTok{15}\NormalTok{, }\SpecialCharTok{{-}}\DecValTok{1}\NormalTok{, }\FloatTok{0.7}\NormalTok{),}
    \AttributeTok{log\_beta     =} \FunctionTok{rnorm}\NormalTok{(}\DecValTok{15}\NormalTok{, }\SpecialCharTok{{-}}\DecValTok{1}\NormalTok{, }\DecValTok{1}\NormalTok{),}
    \AttributeTok{bias         =} \FunctionTok{rnorm}\NormalTok{(}\DecValTok{15}\NormalTok{,  }\DecValTok{0}\NormalTok{, }\DecValTok{1}\NormalTok{)}
\NormalTok{  )}
\NormalTok{  op\_vec }\OtherTok{\textless{}{-}} \FunctionTok{rbinom}\NormalTok{(}\DecValTok{150}\NormalTok{, }\DecValTok{1}\NormalTok{, }\FloatTok{0.65}\NormalTok{)}

\NormalTok{  fit\_one\_tom1 }\OtherTok{\textless{}{-}} \ControlFlowTok{function}\NormalTok{(ls\_op, lb\_op, lb, bs, a) \{}
\NormalTok{    dd }\OtherTok{\textless{}{-}} \FunctionTok{simulate\_tom1}\NormalTok{(op\_vec, ls\_op, lb\_op, lb, bs, }\AttributeTok{seed =}\NormalTok{ a)}
\NormalTok{    fit }\OtherTok{\textless{}{-}}\NormalTok{ mod\_tom1}\SpecialCharTok{$}\FunctionTok{sample}\NormalTok{(}
      \AttributeTok{data =} \FunctionTok{list}\NormalTok{(}\AttributeTok{N =} \FunctionTok{nrow}\NormalTok{(dd), }\AttributeTok{choice =}\NormalTok{ dd}\SpecialCharTok{$}\NormalTok{choice, }\AttributeTok{op\_choice =}\NormalTok{ dd}\SpecialCharTok{$}\NormalTok{op\_choice),}
      \AttributeTok{chains =} \DecValTok{2}\NormalTok{, }\AttributeTok{parallel\_chains =} \DecValTok{2}\NormalTok{,}
      \AttributeTok{iter\_warmup =} \DecValTok{500}\NormalTok{, }\AttributeTok{iter\_sampling =} \DecValTok{500}\NormalTok{,}
      \AttributeTok{refresh =} \DecValTok{0}\NormalTok{, }\AttributeTok{show\_messages =} \ConstantTok{FALSE}
\NormalTok{    )}
\NormalTok{    fit}\SpecialCharTok{$}\FunctionTok{summary}\NormalTok{(}\FunctionTok{c}\NormalTok{(}\StringTok{"log\_sigma\_op"}\NormalTok{,}\StringTok{"log\_beta\_op"}\NormalTok{,}\StringTok{"log\_beta"}\NormalTok{,}\StringTok{"bias"}\NormalTok{)) }\SpecialCharTok{|\textgreater{}}
\NormalTok{      dplyr}\SpecialCharTok{::}\FunctionTok{select}\NormalTok{(variable, mean, q5, q95)}
\NormalTok{  \}}

\NormalTok{  recovery1 }\OtherTok{\textless{}{-}}\NormalTok{ truth1 }\SpecialCharTok{|\textgreater{}}
    \FunctionTok{rowwise}\NormalTok{() }\SpecialCharTok{|\textgreater{}}
    \FunctionTok{mutate}\NormalTok{(}\AttributeTok{post =} \FunctionTok{list}\NormalTok{(}\FunctionTok{fit\_one\_tom1}\NormalTok{(log\_sigma\_op, log\_beta\_op,}
\NormalTok{                                    log\_beta, bias, agent))) }\SpecialCharTok{|\textgreater{}}
\NormalTok{    dplyr}\SpecialCharTok{::}\FunctionTok{select}\NormalTok{(agent, log\_sigma\_op, log\_beta\_op, log\_beta, bias, post) }\SpecialCharTok{|\textgreater{}}
    \FunctionTok{unnest}\NormalTok{(post)}
  \FunctionTok{saveRDS}\NormalTok{(recovery1, recovery\_tom1\_path)}
\NormalTok{\} }\ControlFlowTok{else}\NormalTok{ \{}
\NormalTok{  recovery1 }\OtherTok{\textless{}{-}} \FunctionTok{readRDS}\NormalTok{(recovery\_tom1\_path)}
\NormalTok{\}}

\NormalTok{recovery1\_long }\OtherTok{\textless{}{-}}\NormalTok{ recovery1 }\SpecialCharTok{|\textgreater{}}
  \FunctionTok{pivot\_longer}\NormalTok{(}\FunctionTok{c}\NormalTok{(log\_sigma\_op, log\_beta\_op, log\_beta, bias),}
               \AttributeTok{names\_to =} \StringTok{"param"}\NormalTok{, }\AttributeTok{values\_to =} \StringTok{"true"}\NormalTok{) }\SpecialCharTok{|\textgreater{}}
  \FunctionTok{filter}\NormalTok{(variable }\SpecialCharTok{==}\NormalTok{ param)}
\end{Highlighting}
\end{Shaded}

\begin{Shaded}
\begin{Highlighting}[]
\FunctionTok{ggplot}\NormalTok{(recovery1\_long, }\FunctionTok{aes}\NormalTok{(}\AttributeTok{x =}\NormalTok{ true, }\AttributeTok{y =}\NormalTok{ mean)) }\SpecialCharTok{+}
  \FunctionTok{geom\_errorbar}\NormalTok{(}\FunctionTok{aes}\NormalTok{(}\AttributeTok{ymin =}\NormalTok{ q5, }\AttributeTok{ymax =}\NormalTok{ q95), }\AttributeTok{width =} \DecValTok{0}\NormalTok{, }\AttributeTok{alpha =} \FloatTok{0.4}\NormalTok{,}
                \AttributeTok{color =} \StringTok{"steelblue"}\NormalTok{) }\SpecialCharTok{+}
  \FunctionTok{geom\_point}\NormalTok{(}\AttributeTok{color =} \StringTok{"midnightblue"}\NormalTok{, }\AttributeTok{size =} \DecValTok{2}\NormalTok{) }\SpecialCharTok{+}
  \FunctionTok{geom\_abline}\NormalTok{(}\AttributeTok{slope =} \DecValTok{1}\NormalTok{, }\AttributeTok{intercept =} \DecValTok{0}\NormalTok{, }\AttributeTok{linetype =} \StringTok{"dashed"}\NormalTok{, }\AttributeTok{color =} \StringTok{"gray40"}\NormalTok{) }\SpecialCharTok{+}
  \FunctionTok{facet\_wrap}\NormalTok{(}\SpecialCharTok{\textasciitilde{}}\NormalTok{param, }\AttributeTok{scales =} \StringTok{"free"}\NormalTok{, }\AttributeTok{ncol =} \DecValTok{2}\NormalTok{) }\SpecialCharTok{+}
  \FunctionTok{labs}\NormalTok{(}\AttributeTok{x =} \StringTok{"True value"}\NormalTok{, }\AttributeTok{y =} \StringTok{"Posterior mean"}\NormalTok{,}
       \AttributeTok{title =} \StringTok{"1{-}ToM parameter recovery (15 synthetic agents, 150 trials)"}\NormalTok{)}
\end{Highlighting}
\end{Shaded}

\begin{figure}[H]

{\centering \includegraphics[width=0.85\linewidth,height=\textheight,keepaspectratio]{09-TheoryOfMind_files/figure-pdf/ch09_recovery_tom1_plot-1.pdf}

}

\caption{1-ToM parameter recovery. The agent's own log\_beta and bias
are well recovered. The opponent-model parameters log\_sigma\_op and
log\_beta\_op are noisier --- they enter the likelihood only through the
recursion, and a 150-trial sequence only weakly constrains them.}

\end{figure}%

The pattern is exactly what we should expect from the theory. The
agent's \emph{own} decision parameters (\texttt{log\_beta},
\texttt{bias}) are identified directly from the choice likelihood and
recover cleanly. The \emph{opponent-model} parameters
(\texttt{log\_sigma\_op}, \texttt{log\_beta\_op}) are identified only
through the recursive belief update that feeds into \(\Delta V\), and
their recovery is correspondingly weaker. But a single recovery study
cannot say \emph{how} weak --- recovery evaluated at the prior mode can
look fine even when the posterior is grossly miscalibrated elsewhere in
parameter space. That is what SBC is for.

\section{Phase 4b --- SBC for 1-ToM}\label{phase-4b-sbc-for-1-tom}

Skipping SBC on the model whose recovery is weakest would invert the
logic of the battery. The whole point of Phase 4 is to \emph{diagnose}
the geometry of models whose single-dataset recovery is ambiguous. We
therefore run SBC on 1-ToM with exactly the same machinery, same number
of simulations, and same trial count we used for 0-ToM. The structural
prediction from the fixed-point analysis above is that
\texttt{log\_beta} and \texttt{bias} will calibrate cleanly, while
\texttt{log\_sigma\_op} and \texttt{log\_beta\_op} should show
substantial ECDF excursions --- their posteriors are weakly updated from
the prior because matching pennies does not provide a regime change
sharp enough to pin them down.

\begin{Shaded}
\begin{Highlighting}[]
\NormalTok{sbc\_tom1\_path }\OtherTok{\textless{}{-}}\NormalTok{ here}\SpecialCharTok{::}\FunctionTok{here}\NormalTok{(}\StringTok{"simmodels"}\NormalTok{, }\StringTok{"ch09\_tom1\_sbc.rds"}\NormalTok{)}

\NormalTok{generate\_sbc\_tom1 }\OtherTok{\textless{}{-}} \ControlFlowTok{function}\NormalTok{(}\AttributeTok{N =} \DecValTok{150}\NormalTok{, }\AttributeTok{op\_bias =} \FloatTok{0.65}\NormalTok{) \{}
\NormalTok{  log\_sigma\_op }\OtherTok{\textless{}{-}} \FunctionTok{rnorm}\NormalTok{(}\DecValTok{1}\NormalTok{, }\SpecialCharTok{{-}}\DecValTok{2}\NormalTok{, }\DecValTok{1}\NormalTok{)}
\NormalTok{  log\_beta\_op  }\OtherTok{\textless{}{-}} \FunctionTok{rnorm}\NormalTok{(}\DecValTok{1}\NormalTok{, }\SpecialCharTok{{-}}\DecValTok{1}\NormalTok{, }\FloatTok{0.7}\NormalTok{)}
\NormalTok{  log\_beta     }\OtherTok{\textless{}{-}} \FunctionTok{rnorm}\NormalTok{(}\DecValTok{1}\NormalTok{, }\SpecialCharTok{{-}}\DecValTok{1}\NormalTok{, }\DecValTok{1}\NormalTok{)}
\NormalTok{  bias         }\OtherTok{\textless{}{-}} \FunctionTok{rnorm}\NormalTok{(}\DecValTok{1}\NormalTok{,  }\DecValTok{0}\NormalTok{, }\DecValTok{1}\NormalTok{)}
\NormalTok{  op\_vec       }\OtherTok{\textless{}{-}} \FunctionTok{rbinom}\NormalTok{(N, }\DecValTok{1}\NormalTok{, op\_bias)}
\NormalTok{  sim\_dat }\OtherTok{\textless{}{-}} \FunctionTok{simulate\_tom1}\NormalTok{(op\_vec, log\_sigma\_op, log\_beta\_op,}
\NormalTok{                           log\_beta, bias)}
  \FunctionTok{list}\NormalTok{(}
    \AttributeTok{variables =} \FunctionTok{list}\NormalTok{(}\AttributeTok{log\_sigma\_op =}\NormalTok{ log\_sigma\_op,}
                     \AttributeTok{log\_beta\_op  =}\NormalTok{ log\_beta\_op,}
                     \AttributeTok{log\_beta     =}\NormalTok{ log\_beta,}
                     \AttributeTok{bias         =}\NormalTok{ bias),}
    \AttributeTok{generated =} \FunctionTok{list}\NormalTok{(}\AttributeTok{N =}\NormalTok{ N,}
                     \AttributeTok{choice    =}\NormalTok{ sim\_dat}\SpecialCharTok{$}\NormalTok{choice,}
                     \AttributeTok{op\_choice =}\NormalTok{ sim\_dat}\SpecialCharTok{$}\NormalTok{op\_choice)}
\NormalTok{  )}
\NormalTok{\}}

\ControlFlowTok{if}\NormalTok{ (regenerate\_sbc }\SpecialCharTok{||} \SpecialCharTok{!}\FunctionTok{file.exists}\NormalTok{(sbc\_tom1\_path)) \{}
\NormalTok{  n\_sbc\_tom1 }\OtherTok{\textless{}{-}} \DecValTok{150}
\NormalTok{  sbc\_gen1 }\OtherTok{\textless{}{-}} \FunctionTok{SBC\_generator\_function}\NormalTok{(generate\_sbc\_tom1,}
                                     \AttributeTok{N =} \DecValTok{150}\NormalTok{, }\AttributeTok{op\_bias =} \FloatTok{0.65}\NormalTok{)}
\NormalTok{  sbc\_backend1 }\OtherTok{\textless{}{-}} \FunctionTok{SBC\_backend\_cmdstan\_sample}\NormalTok{(}
\NormalTok{    mod\_tom1,}
    \AttributeTok{iter\_warmup =} \DecValTok{500}\NormalTok{, }\AttributeTok{iter\_sampling =} \DecValTok{500}\NormalTok{, }\AttributeTok{chains =} \DecValTok{1}\NormalTok{, }\AttributeTok{refresh =} \DecValTok{0}
\NormalTok{  )}
\NormalTok{  sbc\_ds1      }\OtherTok{\textless{}{-}} \FunctionTok{generate\_datasets}\NormalTok{(sbc\_gen1, n\_sbc\_tom1)}
\NormalTok{  sbc\_results1 }\OtherTok{\textless{}{-}} \FunctionTok{compute\_SBC}\NormalTok{(}\AttributeTok{datasets =}\NormalTok{ sbc\_ds1, }\AttributeTok{backend =}\NormalTok{ sbc\_backend1,}
                              \AttributeTok{keep\_fits =} \ConstantTok{FALSE}\NormalTok{)}
  \FunctionTok{saveRDS}\NormalTok{(}\FunctionTok{list}\NormalTok{(}\AttributeTok{ds =}\NormalTok{ sbc\_ds1, }\AttributeTok{results =}\NormalTok{ sbc\_results1), sbc\_tom1\_path)}
\NormalTok{\} }\ControlFlowTok{else}\NormalTok{ \{}
\NormalTok{  sbc\_obj1     }\OtherTok{\textless{}{-}} \FunctionTok{readRDS}\NormalTok{(sbc\_tom1\_path)}
\NormalTok{  sbc\_ds1      }\OtherTok{\textless{}{-}}\NormalTok{ sbc\_obj1}\SpecialCharTok{$}\NormalTok{ds}
\NormalTok{  sbc\_results1 }\OtherTok{\textless{}{-}}\NormalTok{ sbc\_obj1}\SpecialCharTok{$}\NormalTok{results}
\NormalTok{\}}
\end{Highlighting}
\end{Shaded}

\begin{Shaded}
\begin{Highlighting}[]
\NormalTok{p\_hist1 }\OtherTok{\textless{}{-}} \FunctionTok{plot\_rank\_hist}\NormalTok{(sbc\_results1,}
  \AttributeTok{variables =} \FunctionTok{c}\NormalTok{(}\StringTok{"log\_sigma\_op"}\NormalTok{,}\StringTok{"log\_beta\_op"}\NormalTok{,}\StringTok{"log\_beta"}\NormalTok{,}\StringTok{"bias"}\NormalTok{)) }\SpecialCharTok{+}
  \FunctionTok{labs}\NormalTok{(}\AttributeTok{subtitle =} \StringTok{"Rank histograms"}\NormalTok{)}
\NormalTok{p\_ecdf1 }\OtherTok{\textless{}{-}} \FunctionTok{plot\_ecdf\_diff}\NormalTok{(sbc\_results1,}
  \AttributeTok{variables =} \FunctionTok{c}\NormalTok{(}\StringTok{"log\_sigma\_op"}\NormalTok{,}\StringTok{"log\_beta\_op"}\NormalTok{,}\StringTok{"log\_beta"}\NormalTok{,}\StringTok{"bias"}\NormalTok{)) }\SpecialCharTok{+}
  \FunctionTok{labs}\NormalTok{(}\AttributeTok{subtitle =} \StringTok{"ECDF difference"}\NormalTok{)}
\NormalTok{p\_hist1 }\SpecialCharTok{/}\NormalTok{ p\_ecdf1}
\end{Highlighting}
\end{Shaded}

\begin{figure}[H]

{\centering \includegraphics[width=0.85\linewidth,height=\textheight,keepaspectratio]{09-TheoryOfMind_files/figure-pdf/ch09_sbc_tom1_plot-1.pdf}

}

\caption{SBC for 1-ToM. Top row: rank histograms. Bottom row: ECDF
difference (blue inside grey band = calibrated). log\_beta and bias stay
inside the band. log\_sigma\_op and log\_beta\_op show substantial
excursions --- the expected empirical signature of an underidentified
parameter. The 1-ToM likelihood simply does not have enough contrast in
matching-pennies data to pin these opponent-model parameters down across
the prior.}

\end{figure}%

The SBC output is not a failure of the sampler --- HMC is doing its job.
It is a \emph{design} diagnostic: matching pennies, as a 2×2 game
without regime changes, lacks the likelihood contrast needed to identify
the opponent-model volatility and temperature. Once
\(\Sigma^{\text{self}}\) reaches its fixed point (within
\textasciitilde10 trials), the predicted opponent-choice trajectory
depends on \((\sigma_{\text{op}}, \beta_{\text{op}})\) only through a
compressed, nearly one-dimensional ridge inside the sigmoid. A
reversal-style paradigm, or a richer payoff structure, would be required
to break this ridge; we do not have that luxury with this dataset.

The scientific consequence is direct: we should not report per-species
or per-individual effects on \texttt{log\_sigma\_op} or
\texttt{log\_beta\_op}. Inference on 1-ToM from this data is credible
for the agent's \emph{own} decision-level parameters and for the
\emph{architecture indicator} (mixture weight, below), but not for the
internal opponent-model parameters.

\section{Phase 5 --- Prior sensitivity}\label{phase-5-prior-sensitivity}

A quick power-scaling check on a representative 1-ToM fit confirms that
the opponent-model priors are doing visible work --- posteriors for
\texttt{log\_sigma\_op} and \texttt{log\_beta\_op} shift when the prior
is up- or down-weighted, but posteriors for \texttt{log\_beta} and
\texttt{bias} are largely driven by the likelihood. Code pattern
(mirrors Ch. 6):

\begin{Shaded}
\begin{Highlighting}[]
\CommentTok{\# Example pattern; run on any single fit:}
\CommentTok{\# powerscale\_sensitivity(}
\CommentTok{\#   fit\_proto\_single, variable = c("log\_sigma\_op","log\_beta\_op","log\_beta","bias")}
\CommentTok{\# ) |\textgreater{} print()}
\CommentTok{\# priorsense::powerscale\_plot\_dens(}
\CommentTok{\#   priorsense::powerscale\_sequence(}
\CommentTok{\#     fit\_proto\_single, variable = c("log\_sigma\_op","log\_beta\_op","log\_beta","bias"))}
\CommentTok{\# )}
\end{Highlighting}
\end{Shaded}

This matches our expectation: interpret \texttt{log\_beta\_op} and
\texttt{log\_sigma\_op} as prior-regularized, put primary scientific
weight on the species-level mixture weight \(\pi_s\) (the
representational-capacity estimand below), and secondary weight on
\texttt{log\_beta} as a within-architecture decisional-determinism
measure.

\section{Phase 6 --- Model comparison under Occam's
razor}\label{phase-6-model-comparison-under-occams-razor}

Chapter 7 warned about a specific failure mode: a more flexible model
can always achieve equal or better \emph{in-sample} fit than a simpler
one, but PSIS-LOO rewards predictive performance on held-out data and
therefore \emph{should} prefer the simpler model whenever the extra
flexibility is not justified. Let us verify this for 0-ToM vs 1-ToM
directly, by fitting both models to datasets generated by three
different processes of increasing complexity:

\begin{enumerate}
\def\labelenumi{\arabic{enumi}.}
\tightlist
\item
  \textbf{RS-opponent data}: a 0-ToM primate playing against a 65/35
  biased opponent. Simplest possible data-generating process.
\item
  \textbf{0-ToM-opponent data}: a 0-ToM primate playing against a 0-ToM
  opponent. The opponent reacts, but the primate does not represent that
  reactivity.
\item
  \textbf{1-ToM-opponent data}: a 1-ToM primate (truly recursive)
  playing against a 0-ToM opponent.
\end{enumerate}

If PSIS-LOO is doing its job, 1-ToM should tie or lose on datasets 1--2
(it is overparameterized for them) and win on dataset 3.

\begin{Shaded}
\begin{Highlighting}[]
\NormalTok{occam\_path }\OtherTok{\textless{}{-}}\NormalTok{ here}\SpecialCharTok{::}\FunctionTok{here}\NormalTok{(}\StringTok{"simmodels"}\NormalTok{, }\StringTok{"ch09\_occam.rds"}\NormalTok{)}

\NormalTok{simulate\_scenario }\OtherTok{\textless{}{-}} \ControlFlowTok{function}\NormalTok{(scenario, }\AttributeTok{N =} \DecValTok{150}\NormalTok{, }\AttributeTok{seed =} \DecValTok{1}\NormalTok{) \{}
  \FunctionTok{set.seed}\NormalTok{(seed)}
  \ControlFlowTok{if}\NormalTok{ (scenario }\SpecialCharTok{==} \StringTok{"RS\_0tom"}\NormalTok{) \{}
\NormalTok{    op }\OtherTok{\textless{}{-}} \FunctionTok{rbinom}\NormalTok{(N, }\DecValTok{1}\NormalTok{, }\FloatTok{0.65}\NormalTok{)}
\NormalTok{    dd }\OtherTok{\textless{}{-}} \FunctionTok{simulate\_tom0}\NormalTok{(op, }\SpecialCharTok{{-}}\DecValTok{2}\NormalTok{, }\SpecialCharTok{{-}}\FloatTok{0.5}\NormalTok{, }\DecValTok{0}\NormalTok{, }\AttributeTok{seed =}\NormalTok{ seed)}
    \FunctionTok{return}\NormalTok{(}\FunctionTok{list}\NormalTok{(}\AttributeTok{choice =}\NormalTok{ dd}\SpecialCharTok{$}\NormalTok{choice, }\AttributeTok{op\_choice =}\NormalTok{ dd}\SpecialCharTok{$}\NormalTok{op\_choice))}
\NormalTok{  \}}
  \ControlFlowTok{if}\NormalTok{ (scenario }\SpecialCharTok{==} \StringTok{"0tom\_0tom"}\NormalTok{) \{}
    \CommentTok{\# Iterate: primate is 0{-}ToM, opponent is 0{-}ToM{-}mismatch.}
\NormalTok{    mu\_p }\OtherTok{\textless{}{-}} \DecValTok{0}\NormalTok{; Sig\_p }\OtherTok{\textless{}{-}} \DecValTok{1}
\NormalTok{    mu\_o }\OtherTok{\textless{}{-}} \DecValTok{0}\NormalTok{; Sig\_o }\OtherTok{\textless{}{-}} \DecValTok{1}
\NormalTok{    choice }\OtherTok{\textless{}{-}} \FunctionTok{integer}\NormalTok{(N); op\_choice }\OtherTok{\textless{}{-}} \FunctionTok{integer}\NormalTok{(N)}
    \ControlFlowTok{for}\NormalTok{ (t }\ControlFlowTok{in} \FunctionTok{seq\_len}\NormalTok{(N)) \{}
\NormalTok{      p\_op\_pred   }\OtherTok{\textless{}{-}} \FunctionTok{plogis}\NormalTok{(mu\_p }\SpecialCharTok{/} \FunctionTok{sqrt}\NormalTok{(}\DecValTok{1} \SpecialCharTok{+} \FloatTok{0.36} \SpecialCharTok{*}\NormalTok{ Sig\_p))}
\NormalTok{      p\_self\_pred }\OtherTok{\textless{}{-}} \FunctionTok{plogis}\NormalTok{(mu\_o }\SpecialCharTok{/} \FunctionTok{sqrt}\NormalTok{(}\DecValTok{1} \SpecialCharTok{+} \FloatTok{0.36} \SpecialCharTok{*}\NormalTok{ Sig\_o))}
\NormalTok{      p\_self\_t }\OtherTok{\textless{}{-}} \FunctionTok{plogis}\NormalTok{((}\DecValTok{2} \SpecialCharTok{*}\NormalTok{ p\_op\_pred }\SpecialCharTok{{-}} \DecValTok{1}\NormalTok{) }\SpecialCharTok{/} \FunctionTok{exp}\NormalTok{(}\SpecialCharTok{{-}}\FloatTok{0.5}\NormalTok{))}
\NormalTok{      p\_op\_t   }\OtherTok{\textless{}{-}} \FunctionTok{plogis}\NormalTok{((}\DecValTok{1} \SpecialCharTok{{-}} \DecValTok{2} \SpecialCharTok{*}\NormalTok{ p\_self\_pred) }\SpecialCharTok{/} \FunctionTok{exp}\NormalTok{(}\SpecialCharTok{{-}}\FloatTok{0.5}\NormalTok{))}
\NormalTok{      choice[t]    }\OtherTok{\textless{}{-}} \FunctionTok{rbinom}\NormalTok{(}\DecValTok{1}\NormalTok{, }\DecValTok{1}\NormalTok{, p\_self\_t)}
\NormalTok{      op\_choice[t] }\OtherTok{\textless{}{-}} \FunctionTok{rbinom}\NormalTok{(}\DecValTok{1}\NormalTok{, }\DecValTok{1}\NormalTok{, p\_op\_t)}
\NormalTok{      s\_p }\OtherTok{\textless{}{-}} \FunctionTok{plogis}\NormalTok{(mu\_p)}
\NormalTok{      Sig\_p }\OtherTok{\textless{}{-}} \DecValTok{1} \SpecialCharTok{/}\NormalTok{ (}\DecValTok{1} \SpecialCharTok{/}\NormalTok{ (Sig\_p }\SpecialCharTok{+} \FunctionTok{exp}\NormalTok{(}\SpecialCharTok{{-}}\DecValTok{2}\NormalTok{)) }\SpecialCharTok{+}\NormalTok{ s\_p }\SpecialCharTok{*}\NormalTok{ (}\DecValTok{1} \SpecialCharTok{{-}}\NormalTok{ s\_p))}
\NormalTok{      mu\_p  }\OtherTok{\textless{}{-}}\NormalTok{ mu\_p }\SpecialCharTok{+}\NormalTok{ Sig\_p }\SpecialCharTok{*}\NormalTok{ (op\_choice[t] }\SpecialCharTok{{-}}\NormalTok{ s\_p)}
\NormalTok{      s\_o }\OtherTok{\textless{}{-}} \FunctionTok{plogis}\NormalTok{(mu\_o)}
\NormalTok{      Sig\_o }\OtherTok{\textless{}{-}} \DecValTok{1} \SpecialCharTok{/}\NormalTok{ (}\DecValTok{1} \SpecialCharTok{/}\NormalTok{ (Sig\_o }\SpecialCharTok{+} \FunctionTok{exp}\NormalTok{(}\SpecialCharTok{{-}}\DecValTok{2}\NormalTok{)) }\SpecialCharTok{+}\NormalTok{ s\_o }\SpecialCharTok{*}\NormalTok{ (}\DecValTok{1} \SpecialCharTok{{-}}\NormalTok{ s\_o))}
\NormalTok{      mu\_o  }\OtherTok{\textless{}{-}}\NormalTok{ mu\_o }\SpecialCharTok{+}\NormalTok{ Sig\_o }\SpecialCharTok{*}\NormalTok{ (choice[t] }\SpecialCharTok{{-}}\NormalTok{ s\_o)}
\NormalTok{    \}}
    \FunctionTok{return}\NormalTok{(}\FunctionTok{list}\NormalTok{(}\AttributeTok{choice =}\NormalTok{ choice, }\AttributeTok{op\_choice =}\NormalTok{ op\_choice))}
\NormalTok{  \}}
  \ControlFlowTok{if}\NormalTok{ (scenario }\SpecialCharTok{==} \StringTok{"0tom\_1tom"}\NormalTok{) \{}
\NormalTok{    mu\_s }\OtherTok{\textless{}{-}} \DecValTok{0}\NormalTok{; Sig\_s }\OtherTok{\textless{}{-}} \DecValTok{1}     \CommentTok{\# 1{-}ToM primate\textquotesingle{}s model of opp\textquotesingle{}s belief about primate}
\NormalTok{    mu\_o }\OtherTok{\textless{}{-}} \DecValTok{0}\NormalTok{; Sig\_o }\OtherTok{\textless{}{-}} \DecValTok{1}     \CommentTok{\# 0{-}ToM opp\textquotesingle{}s belief about primate}
\NormalTok{    choice }\OtherTok{\textless{}{-}} \FunctionTok{integer}\NormalTok{(N); op\_choice }\OtherTok{\textless{}{-}} \FunctionTok{integer}\NormalTok{(N)}
    \ControlFlowTok{for}\NormalTok{ (t }\ControlFlowTok{in} \FunctionTok{seq\_len}\NormalTok{(N)) \{}
\NormalTok{      p\_self\_from\_op  }\OtherTok{\textless{}{-}} \FunctionTok{plogis}\NormalTok{(mu\_s }\SpecialCharTok{/} \FunctionTok{sqrt}\NormalTok{(}\DecValTok{1} \SpecialCharTok{+} \FloatTok{0.36} \SpecialCharTok{*}\NormalTok{ Sig\_s))}
\NormalTok{      p\_op\_pred\_1tom  }\OtherTok{\textless{}{-}} \FunctionTok{plogis}\NormalTok{((}\DecValTok{1} \SpecialCharTok{{-}} \DecValTok{2} \SpecialCharTok{*}\NormalTok{ p\_self\_from\_op) }\SpecialCharTok{/} \FunctionTok{exp}\NormalTok{(}\SpecialCharTok{{-}}\FloatTok{0.5}\NormalTok{))}
\NormalTok{      p\_self\_t        }\OtherTok{\textless{}{-}} \FunctionTok{plogis}\NormalTok{((}\DecValTok{2} \SpecialCharTok{*}\NormalTok{ p\_op\_pred\_1tom }\SpecialCharTok{{-}} \DecValTok{1}\NormalTok{) }\SpecialCharTok{/} \FunctionTok{exp}\NormalTok{(}\SpecialCharTok{{-}}\FloatTok{0.5}\NormalTok{))}
\NormalTok{      p\_self\_opp\_view }\OtherTok{\textless{}{-}} \FunctionTok{plogis}\NormalTok{(mu\_o }\SpecialCharTok{/} \FunctionTok{sqrt}\NormalTok{(}\DecValTok{1} \SpecialCharTok{+} \FloatTok{0.36} \SpecialCharTok{*}\NormalTok{ Sig\_o))}
\NormalTok{      p\_op\_t          }\OtherTok{\textless{}{-}} \FunctionTok{plogis}\NormalTok{((}\DecValTok{1} \SpecialCharTok{{-}} \DecValTok{2} \SpecialCharTok{*}\NormalTok{ p\_self\_opp\_view) }\SpecialCharTok{/} \FunctionTok{exp}\NormalTok{(}\SpecialCharTok{{-}}\FloatTok{0.5}\NormalTok{))}
\NormalTok{      choice[t]    }\OtherTok{\textless{}{-}} \FunctionTok{rbinom}\NormalTok{(}\DecValTok{1}\NormalTok{, }\DecValTok{1}\NormalTok{, p\_self\_t)}
\NormalTok{      op\_choice[t] }\OtherTok{\textless{}{-}} \FunctionTok{rbinom}\NormalTok{(}\DecValTok{1}\NormalTok{, }\DecValTok{1}\NormalTok{, p\_op\_t)}
\NormalTok{      s\_ps }\OtherTok{\textless{}{-}}\NormalTok{ p\_self\_from\_op}
\NormalTok{      Sig\_s }\OtherTok{\textless{}{-}} \DecValTok{1} \SpecialCharTok{/}\NormalTok{ (}\DecValTok{1} \SpecialCharTok{/}\NormalTok{ (Sig\_s }\SpecialCharTok{+} \FunctionTok{exp}\NormalTok{(}\SpecialCharTok{{-}}\DecValTok{2}\NormalTok{)) }\SpecialCharTok{+}\NormalTok{ s\_ps }\SpecialCharTok{*}\NormalTok{ (}\DecValTok{1} \SpecialCharTok{{-}}\NormalTok{ s\_ps))}
\NormalTok{      mu\_s  }\OtherTok{\textless{}{-}}\NormalTok{ mu\_s }\SpecialCharTok{+}\NormalTok{ Sig\_s }\SpecialCharTok{*}\NormalTok{ (choice[t] }\SpecialCharTok{{-}}\NormalTok{ s\_ps)}
\NormalTok{      s\_o  }\OtherTok{\textless{}{-}} \FunctionTok{plogis}\NormalTok{(mu\_o)}
\NormalTok{      Sig\_o }\OtherTok{\textless{}{-}} \DecValTok{1} \SpecialCharTok{/}\NormalTok{ (}\DecValTok{1} \SpecialCharTok{/}\NormalTok{ (Sig\_o }\SpecialCharTok{+} \FunctionTok{exp}\NormalTok{(}\SpecialCharTok{{-}}\DecValTok{2}\NormalTok{)) }\SpecialCharTok{+}\NormalTok{ s\_o }\SpecialCharTok{*}\NormalTok{ (}\DecValTok{1} \SpecialCharTok{{-}}\NormalTok{ s\_o))}
\NormalTok{      mu\_o  }\OtherTok{\textless{}{-}}\NormalTok{ mu\_o }\SpecialCharTok{+}\NormalTok{ Sig\_o }\SpecialCharTok{*}\NormalTok{ (choice[t] }\SpecialCharTok{{-}}\NormalTok{ s\_o)}
\NormalTok{    \}}
    \FunctionTok{return}\NormalTok{(}\FunctionTok{list}\NormalTok{(}\AttributeTok{choice =}\NormalTok{ choice, }\AttributeTok{op\_choice =}\NormalTok{ op\_choice))}
\NormalTok{  \}}
\NormalTok{\}}

\NormalTok{fit\_and\_loo }\OtherTok{\textless{}{-}} \ControlFlowTok{function}\NormalTok{(dd, mod) \{}
\NormalTok{  fit }\OtherTok{\textless{}{-}}\NormalTok{ mod}\SpecialCharTok{$}\FunctionTok{sample}\NormalTok{(}
    \AttributeTok{data =} \FunctionTok{list}\NormalTok{(}\AttributeTok{N =} \FunctionTok{length}\NormalTok{(dd}\SpecialCharTok{$}\NormalTok{choice), }\AttributeTok{choice =}\NormalTok{ dd}\SpecialCharTok{$}\NormalTok{choice,}
                \AttributeTok{op\_choice =}\NormalTok{ dd}\SpecialCharTok{$}\NormalTok{op\_choice),}
    \AttributeTok{chains =} \DecValTok{2}\NormalTok{, }\AttributeTok{parallel\_chains =} \DecValTok{2}\NormalTok{,}
    \AttributeTok{iter\_warmup =} \DecValTok{500}\NormalTok{, }\AttributeTok{iter\_sampling =} \DecValTok{500}\NormalTok{, }\AttributeTok{refresh =} \DecValTok{0}\NormalTok{,}
    \AttributeTok{show\_messages =} \ConstantTok{FALSE}
\NormalTok{  )}
\NormalTok{  fit}\SpecialCharTok{$}\FunctionTok{loo}\NormalTok{()}
\NormalTok{\}}

\ControlFlowTok{if}\NormalTok{ (regenerate\_simulations }\SpecialCharTok{||} \SpecialCharTok{!}\FunctionTok{file.exists}\NormalTok{(occam\_path)) \{}
\NormalTok{  scenarios }\OtherTok{\textless{}{-}} \FunctionTok{c}\NormalTok{(}\StringTok{"RS\_0tom"}\NormalTok{, }\StringTok{"0tom\_0tom"}\NormalTok{, }\StringTok{"0tom\_1tom"}\NormalTok{)}
\NormalTok{  n\_rep }\OtherTok{\textless{}{-}} \DecValTok{6}
\NormalTok{  occam\_grid }\OtherTok{\textless{}{-}} \FunctionTok{expand\_grid}\NormalTok{(}\AttributeTok{scenario =}\NormalTok{ scenarios, }\AttributeTok{rep =} \DecValTok{1}\SpecialCharTok{:}\NormalTok{n\_rep)}
\NormalTok{  occam\_res }\OtherTok{\textless{}{-}}\NormalTok{ occam\_grid }\SpecialCharTok{|\textgreater{}}
    \FunctionTok{rowwise}\NormalTok{() }\SpecialCharTok{|\textgreater{}}
    \FunctionTok{mutate}\NormalTok{(}
      \AttributeTok{dat   =} \FunctionTok{list}\NormalTok{(}\FunctionTok{simulate\_scenario}\NormalTok{(scenario,}
                                     \AttributeTok{seed =} \DecValTok{100} \SpecialCharTok{*} \FunctionTok{match}\NormalTok{(scenario, scenarios) }\SpecialCharTok{+}\NormalTok{ rep)),}
      \AttributeTok{loo0  =} \FunctionTok{list}\NormalTok{(}\FunctionTok{fit\_and\_loo}\NormalTok{(dat, mod\_tom0)),}
      \AttributeTok{loo1  =} \FunctionTok{list}\NormalTok{(}\FunctionTok{fit\_and\_loo}\NormalTok{(dat, mod\_tom1)),}
      \AttributeTok{elpd0 =}\NormalTok{ loo0}\SpecialCharTok{$}\NormalTok{estimates[}\StringTok{"elpd\_loo"}\NormalTok{,}\StringTok{"Estimate"}\NormalTok{],}
      \AttributeTok{elpd1 =}\NormalTok{ loo1}\SpecialCharTok{$}\NormalTok{estimates[}\StringTok{"elpd\_loo"}\NormalTok{,}\StringTok{"Estimate"}\NormalTok{],}
      \AttributeTok{delta =}\NormalTok{ elpd1 }\SpecialCharTok{{-}}\NormalTok{ elpd0}
\NormalTok{    ) }\SpecialCharTok{|\textgreater{}}
\NormalTok{    dplyr}\SpecialCharTok{::}\FunctionTok{select}\NormalTok{(scenario, rep, elpd0, elpd1, delta)}
  \FunctionTok{saveRDS}\NormalTok{(occam\_res, occam\_path)}
\NormalTok{\} }\ControlFlowTok{else}\NormalTok{ \{}
\NormalTok{  occam\_res }\OtherTok{\textless{}{-}} \FunctionTok{readRDS}\NormalTok{(occam\_path)}
\NormalTok{\}}
\end{Highlighting}
\end{Shaded}

\begin{Shaded}
\begin{Highlighting}[]
\FunctionTok{ggplot}\NormalTok{(occam\_res, }\FunctionTok{aes}\NormalTok{(}\AttributeTok{x =}\NormalTok{ scenario, }\AttributeTok{y =}\NormalTok{ delta)) }\SpecialCharTok{+}
  \FunctionTok{geom\_hline}\NormalTok{(}\AttributeTok{yintercept =} \DecValTok{0}\NormalTok{, }\AttributeTok{linetype =} \StringTok{"dashed"}\NormalTok{, }\AttributeTok{color =} \StringTok{"gray40"}\NormalTok{) }\SpecialCharTok{+}
  \FunctionTok{geom\_boxplot}\NormalTok{(}\AttributeTok{outlier.shape =} \ConstantTok{NA}\NormalTok{, }\AttributeTok{fill =} \StringTok{"gray90"}\NormalTok{) }\SpecialCharTok{+}
  \FunctionTok{geom\_jitter}\NormalTok{(}\AttributeTok{width =} \FloatTok{0.12}\NormalTok{, }\AttributeTok{alpha =} \FloatTok{0.8}\NormalTok{, }\AttributeTok{color =} \StringTok{"midnightblue"}\NormalTok{) }\SpecialCharTok{+}
  \FunctionTok{scale\_x\_discrete}\NormalTok{(}\AttributeTok{labels =} \FunctionTok{c}\NormalTok{(}\StringTok{"RS\_0tom"}   \OtherTok{=} \StringTok{"RS + 0{-}ToM primate"}\NormalTok{,}
                              \StringTok{"0tom\_0tom"} \OtherTok{=} \StringTok{"0{-}ToM opp + 0{-}ToM primate"}\NormalTok{,}
                              \StringTok{"0tom\_1tom"} \OtherTok{=} \StringTok{"0{-}ToM opp + 1{-}ToM primate"}\NormalTok{)) }\SpecialCharTok{+}
  \FunctionTok{labs}\NormalTok{(}\AttributeTok{x =} \StringTok{"Data{-}generating scenario"}\NormalTok{,}
       \AttributeTok{y =} \FunctionTok{expression}\NormalTok{(Delta }\SpecialCharTok{*} \StringTok{" ELPD (1{-}ToM − 0{-}ToM)"}\NormalTok{),}
       \AttributeTok{title =} \StringTok{"Does PSIS{-}LOO correctly penalize 1{-}ToM on simpler data?"}\NormalTok{)}
\end{Highlighting}
\end{Shaded}

\begin{figure}[H]

{\centering \includegraphics[width=0.85\linewidth,height=\textheight,keepaspectratio]{09-TheoryOfMind_files/figure-pdf/ch09_occam_plot-1.pdf}

}

\caption{Occam's razor test. Each point is one simulated dataset; y-axis
= ELPD(1-ToM) - ELPD(0-ToM). Positive favors 1-ToM, negative favors
0-ToM. On RS and 0-ToM-opponent data (simpler generating processes) the
difference clusters near or below zero --- PSIS-LOO correctly refuses to
pay the complexity cost of 1-ToM. On 1-ToM-opponent data the difference
shifts positive --- the extra structure earns its keep.}

\end{figure}%

The answer is yes: 1-ToM is not a free upgrade. On data where the
primate does not actually mentalize, the extra two parameters of 1-ToM
leave the ELPD difference at or below zero. On data where the primate
truly is a 1-ToM, the difference swings positive. This is the
Occam's-razor behavior Ch. 7 promised.

\section{Fitting 0-ToM hierarchically to the primate
data}\label{fitting-0-tom-hierarchically-to-the-primate-data}

Each primate plays all three opponent conditions, and we have 39
primates across 7 species. For the scientific estimand we focus on the
interactive \textbf{0-ToM-opponent condition}, where 0-ToM and 1-ToM
make the most differentiated predictions. One 0-ToM fit per individual,
with all three parameters varying by individual around a species-level
mean, non-centered throughout.

\begin{Shaded}
\begin{Highlighting}[]
\NormalTok{stan\_tom0\_hier }\OtherTok{\textless{}{-}} \StringTok{"}
\StringTok{data \{}
\StringTok{  int\textless{}lower=1\textgreater{} N;}
\StringTok{  int\textless{}lower=1\textgreater{} J;}
\StringTok{  int\textless{}lower=1\textgreater{} S;}
\StringTok{  array[J] int\textless{}lower=1, upper=S\textgreater{} species\_of;}
\StringTok{  array[N] int\textless{}lower=1, upper=J\textgreater{} agent;}
\StringTok{  array[N] int\textless{}lower=0, upper=1\textgreater{} choice;}
\StringTok{  array[N] int\textless{}lower=0, upper=1\textgreater{} op\_choice;}
\StringTok{  array[J] int\textless{}lower=1\textgreater{} start;}
\StringTok{  array[J] int\textless{}lower=1\textgreater{} stop;}
\StringTok{\}}
\StringTok{parameters \{}
\StringTok{  vector[S] mu\_log\_sigma;}
\StringTok{  vector[S] mu\_log\_beta;}
\StringTok{  vector[S] mu\_bias;}
\StringTok{  real\textless{}lower=0\textgreater{} tau\_log\_sigma;}
\StringTok{  real\textless{}lower=0\textgreater{} tau\_log\_beta;}
\StringTok{  real\textless{}lower=0\textgreater{} tau\_bias;}
\StringTok{  vector[J] z\_log\_sigma;}
\StringTok{  vector[J] z\_log\_beta;}
\StringTok{  vector[J] z\_bias;}
\StringTok{\}}
\StringTok{transformed parameters \{}
\StringTok{  vector[J] log\_sigma = mu\_log\_sigma[species\_of] + tau\_log\_sigma * z\_log\_sigma;}
\StringTok{  vector[J] log\_beta  = mu\_log\_beta[species\_of]  + tau\_log\_beta  * z\_log\_beta;}
\StringTok{  vector[J] bias      = mu\_bias[species\_of]      + tau\_bias      * z\_bias;}
\StringTok{  vector[N] dV;}
\StringTok{  for (j in 1:J) \{}
\StringTok{    real sigma\_j = exp(log\_sigma[j]);}
\StringTok{    real mu\_s    = 0.0;}
\StringTok{    real Sig\_s   = 1.0;}
\StringTok{    for (k in start[j]:stop[j]) \{}
\StringTok{      real s\_mu = inv\_logit(mu\_s);}
\StringTok{      real p\_op = inv\_logit(mu\_s / sqrt(1 + 0.36 * Sig\_s));}
\StringTok{      dV[k] = 2 * p\_op {-} 1;}
\StringTok{      Sig\_s = 1.0 / (1.0 / (Sig\_s + sigma\_j) + s\_mu * (1 {-} s\_mu));}
\StringTok{      mu\_s  = mu\_s + Sig\_s * (op\_choice[k] {-} s\_mu);}
\StringTok{    \}}
\StringTok{  \}}
\StringTok{\}}
\StringTok{model \{}
\StringTok{  mu\_log\_sigma \textasciitilde{} normal({-}2, 1);}
\StringTok{  mu\_log\_beta  \textasciitilde{} normal({-}1, 1);}
\StringTok{  mu\_bias      \textasciitilde{} normal( 0, 1);}
\StringTok{  tau\_log\_sigma \textasciitilde{} exponential(2);}
\StringTok{  tau\_log\_beta  \textasciitilde{} exponential(1);}
\StringTok{  tau\_bias      \textasciitilde{} exponential(1);}
\StringTok{  z\_log\_sigma \textasciitilde{} std\_normal();}
\StringTok{  z\_log\_beta  \textasciitilde{} std\_normal();}
\StringTok{  z\_bias      \textasciitilde{} std\_normal();}
\StringTok{  for (n in 1:N)}
\StringTok{    choice[n] \textasciitilde{} bernoulli\_logit(}
\StringTok{      (dV[n] + bias[agent[n]]) / exp(log\_beta[agent[n]])}
\StringTok{    );}
\StringTok{\}}
\StringTok{generated quantities \{}
\StringTok{  vector[N] log\_lik;}
\StringTok{  for (n in 1:N)}
\StringTok{    log\_lik[n] = bernoulli\_logit\_lpmf(}
\StringTok{      choice[n] | (dV[n] + bias[agent[n]]) / exp(log\_beta[agent[n]])}
\StringTok{    );}
\StringTok{\}}
\StringTok{"}
\FunctionTok{writeLines}\NormalTok{(stan\_tom0\_hier, here}\SpecialCharTok{::}\FunctionTok{here}\NormalTok{(}\StringTok{"stan"}\NormalTok{, }\StringTok{"ch09\_tom0\_hier.stan"}\NormalTok{))}
\NormalTok{mod\_tom0\_hier }\OtherTok{\textless{}{-}} \FunctionTok{cmdstan\_model}\NormalTok{(here}\SpecialCharTok{::}\FunctionTok{here}\NormalTok{(}\StringTok{"stan"}\NormalTok{, }\StringTok{"ch09\_tom0\_hier.stan"}\NormalTok{))}
\end{Highlighting}
\end{Shaded}

\begin{Shaded}
\begin{Highlighting}[]
\NormalTok{prep\_stan\_data }\OtherTok{\textless{}{-}} \ControlFlowTok{function}\NormalTok{(df) \{}
\NormalTok{  df }\OtherTok{\textless{}{-}}\NormalTok{ df }\SpecialCharTok{|\textgreater{}}
    \FunctionTok{arrange}\NormalTok{(id, t) }\SpecialCharTok{|\textgreater{}}
    \FunctionTok{mutate}\NormalTok{(}
      \AttributeTok{agent\_id   =} \FunctionTok{as.integer}\NormalTok{(}\FunctionTok{factor}\NormalTok{(id)),}
      \AttributeTok{species\_id =} \FunctionTok{as.integer}\NormalTok{(}\FunctionTok{factor}\NormalTok{(species))}
\NormalTok{    )}
\NormalTok{  species\_of }\OtherTok{\textless{}{-}}\NormalTok{ df }\SpecialCharTok{|\textgreater{}} \FunctionTok{distinct}\NormalTok{(agent\_id, species\_id) }\SpecialCharTok{|\textgreater{}}
    \FunctionTok{arrange}\NormalTok{(agent\_id) }\SpecialCharTok{|\textgreater{}} \FunctionTok{pull}\NormalTok{(species\_id)}
\NormalTok{  boundaries }\OtherTok{\textless{}{-}}\NormalTok{ df }\SpecialCharTok{|\textgreater{}}
    \FunctionTok{mutate}\NormalTok{(}\AttributeTok{row =} \FunctionTok{row\_number}\NormalTok{()) }\SpecialCharTok{|\textgreater{}}
    \FunctionTok{group\_by}\NormalTok{(agent\_id) }\SpecialCharTok{|\textgreater{}}
    \FunctionTok{summarize}\NormalTok{(}\AttributeTok{start =} \FunctionTok{min}\NormalTok{(row), }\AttributeTok{stop =} \FunctionTok{max}\NormalTok{(row), }\AttributeTok{.groups =} \StringTok{"drop"}\NormalTok{) }\SpecialCharTok{|\textgreater{}}
    \FunctionTok{arrange}\NormalTok{(agent\_id)}
  \FunctionTok{list}\NormalTok{(}
    \AttributeTok{N =} \FunctionTok{nrow}\NormalTok{(df), }\AttributeTok{J =} \FunctionTok{max}\NormalTok{(df}\SpecialCharTok{$}\NormalTok{agent\_id), }\AttributeTok{S =} \FunctionTok{max}\NormalTok{(df}\SpecialCharTok{$}\NormalTok{species\_id),}
    \AttributeTok{species\_of =}\NormalTok{ species\_of, }\AttributeTok{agent =}\NormalTok{ df}\SpecialCharTok{$}\NormalTok{agent\_id,}
    \AttributeTok{choice =}\NormalTok{ df}\SpecialCharTok{$}\NormalTok{choice, }\AttributeTok{op\_choice =}\NormalTok{ df}\SpecialCharTok{$}\NormalTok{op\_choice,}
    \AttributeTok{start =}\NormalTok{ boundaries}\SpecialCharTok{$}\NormalTok{start, }\AttributeTok{stop =}\NormalTok{ boundaries}\SpecialCharTok{$}\NormalTok{stop,}
    \AttributeTok{agent\_lookup =}\NormalTok{ df }\SpecialCharTok{|\textgreater{}} \FunctionTok{distinct}\NormalTok{(agent\_id, id, species, species\_id),}
    \AttributeTok{species\_lookup =}\NormalTok{ df }\SpecialCharTok{|\textgreater{}} \FunctionTok{distinct}\NormalTok{(species\_id, species) }\SpecialCharTok{|\textgreater{}} \FunctionTok{arrange}\NormalTok{(species\_id)}
\NormalTok{  )}
\NormalTok{\}}

\NormalTok{mp\_int }\OtherTok{\textless{}{-}}\NormalTok{ mp }\SpecialCharTok{|\textgreater{}} \FunctionTok{filter}\NormalTok{(opp\_cond }\SpecialCharTok{==} \DecValTok{0}\NormalTok{)}
\NormalTok{sd\_int }\OtherTok{\textless{}{-}} \FunctionTok{prep\_stan\_data}\NormalTok{(mp\_int)}
\end{Highlighting}
\end{Shaded}

\begin{Shaded}
\begin{Highlighting}[]
\NormalTok{hier\_path }\OtherTok{\textless{}{-}}\NormalTok{ here}\SpecialCharTok{::}\FunctionTok{here}\NormalTok{(}\StringTok{"simmodels"}\NormalTok{, }\StringTok{"ch09\_tom0\_hier\_fit.rds"}\NormalTok{)}
\ControlFlowTok{if}\NormalTok{ (regenerate\_fits }\SpecialCharTok{||} \SpecialCharTok{!}\FunctionTok{file.exists}\NormalTok{(hier\_path)) \{}
\NormalTok{  fit\_hier }\OtherTok{\textless{}{-}}\NormalTok{ mod\_tom0\_hier}\SpecialCharTok{$}\FunctionTok{sample}\NormalTok{(}
    \AttributeTok{data =}\NormalTok{ sd\_int[}\FunctionTok{c}\NormalTok{(}\StringTok{"N"}\NormalTok{,}\StringTok{"J"}\NormalTok{,}\StringTok{"S"}\NormalTok{,}\StringTok{"species\_of"}\NormalTok{,}\StringTok{"agent"}\NormalTok{,}
                    \StringTok{"choice"}\NormalTok{,}\StringTok{"op\_choice"}\NormalTok{,}\StringTok{"start"}\NormalTok{,}\StringTok{"stop"}\NormalTok{)],}
    \AttributeTok{chains =} \DecValTok{4}\NormalTok{, }\AttributeTok{parallel\_chains =} \DecValTok{4}\NormalTok{,}
    \AttributeTok{iter\_warmup =} \DecValTok{1000}\NormalTok{, }\AttributeTok{iter\_sampling =} \DecValTok{1000}\NormalTok{,}
    \AttributeTok{adapt\_delta =} \FloatTok{0.95}\NormalTok{, }\AttributeTok{seed =} \DecValTok{2026}
\NormalTok{  )}
\NormalTok{  fit\_hier}\SpecialCharTok{$}\FunctionTok{save\_object}\NormalTok{(hier\_path)}
\NormalTok{\} }\ControlFlowTok{else}\NormalTok{ \{}
\NormalTok{  fit\_hier }\OtherTok{\textless{}{-}} \FunctionTok{readRDS}\NormalTok{(hier\_path)}
\NormalTok{\}}
\NormalTok{fit\_hier}\SpecialCharTok{$}\FunctionTok{diagnostic\_summary}\NormalTok{()}
\end{Highlighting}
\end{Shaded}

\begin{verbatim}
$num_divergent
[1] 2 2 1 0

$num_max_treedepth
[1] 0 0 0 0

$ebfmi
[1] 0.7725369 0.7377288 0.8043762 0.8359712
\end{verbatim}

\subsection{Fitting 1-ToM
hierarchically}\label{fitting-1-tom-hierarchically}

We fit an analogous hierarchical 1-ToM model to the same interactive
condition data. This model shares the same multilevel structure ---
species-level means, non-centered individual offsets --- but uses the
1-ToM recursion: the agent maintains a belief about what the opponent
thinks of \emph{itself}, updated by the agent's own choices.

Four parameters per individual: \texttt{log\_sigma\_op} (opponent's
believed volatility about the agent), \texttt{log\_beta\_op} (opponent's
believed temperature), \texttt{log\_beta} (agent's own temperature), and
\texttt{bias}. As Phase 4b showed, \texttt{log\_sigma\_op} and
\texttt{log\_beta\_op} are weakly identified at the individual level;
the species-level pooling here stabilizes them but does not solve the
design-level issue.

\begin{Shaded}
\begin{Highlighting}[]
\NormalTok{stan\_tom1\_hier }\OtherTok{\textless{}{-}} \StringTok{"}
\StringTok{data \{}
\StringTok{  int\textless{}lower=1\textgreater{} N;}
\StringTok{  int\textless{}lower=1\textgreater{} J;}
\StringTok{  int\textless{}lower=1\textgreater{} S;}
\StringTok{  array[J] int\textless{}lower=1, upper=S\textgreater{} species\_of;}
\StringTok{  array[N] int\textless{}lower=1, upper=J\textgreater{} agent;}
\StringTok{  array[N] int\textless{}lower=0, upper=1\textgreater{} choice;}
\StringTok{  array[N] int\textless{}lower=0, upper=1\textgreater{} op\_choice;}
\StringTok{  array[J] int\textless{}lower=1\textgreater{} start;}
\StringTok{  array[J] int\textless{}lower=1\textgreater{} stop;}
\StringTok{\}}
\StringTok{parameters \{}
\StringTok{  vector[S] mu\_log\_sigma\_op;}
\StringTok{  real\textless{}lower=0\textgreater{} tau\_log\_sigma\_op;}
\StringTok{  vector[J] z\_log\_sigma\_op;}
\StringTok{  vector[S] mu\_log\_beta\_op;}
\StringTok{  real\textless{}lower=0\textgreater{} tau\_log\_beta\_op;}
\StringTok{  vector[J] z\_log\_beta\_op;}
\StringTok{  vector[S] mu\_log\_beta;}
\StringTok{  real\textless{}lower=0\textgreater{} tau\_log\_beta;}
\StringTok{  vector[J] z\_log\_beta;}
\StringTok{  vector[S] mu\_bias;}
\StringTok{  real\textless{}lower=0\textgreater{} tau\_bias;}
\StringTok{  vector[J] z\_bias;}
\StringTok{\}}
\StringTok{transformed parameters \{}
\StringTok{  vector[J] log\_sigma\_op = mu\_log\_sigma\_op[species\_of] + tau\_log\_sigma\_op * z\_log\_sigma\_op;}
\StringTok{  vector[J] log\_beta\_op  = mu\_log\_beta\_op[species\_of]  + tau\_log\_beta\_op  * z\_log\_beta\_op;}
\StringTok{  vector[J] log\_beta     = mu\_log\_beta[species\_of]     + tau\_log\_beta     * z\_log\_beta;}
\StringTok{  vector[J] bias         = mu\_bias[species\_of]         + tau\_bias         * z\_bias;}
\StringTok{  vector[N] dV;}
\StringTok{  for (j in 1:J) \{}
\StringTok{    real sigma\_op\_j = exp(log\_sigma\_op[j]);}
\StringTok{    real beta\_op\_j  = exp(log\_beta\_op[j]);}
\StringTok{    real mu\_s  = 0.0;}
\StringTok{    real Sig\_s = 1.0;}
\StringTok{    for (k in start[j]:stop[j]) \{}
\StringTok{      real p\_self\_from\_op = inv\_logit(mu\_s / sqrt(1 + 0.36 * Sig\_s));}
\StringTok{      real p\_op           = inv\_logit((1 {-} 2 * p\_self\_from\_op) / beta\_op\_j);}
\StringTok{      dV[k] = 2 * p\_op {-} 1;}
\StringTok{      real s\_ps = p\_self\_from\_op;}
\StringTok{      Sig\_s = 1.0 / (1.0 / (Sig\_s + sigma\_op\_j) + s\_ps * (1 {-} s\_ps));}
\StringTok{      mu\_s  = mu\_s + Sig\_s * (choice[k] {-} s\_ps);}
\StringTok{    \}}
\StringTok{  \}}
\StringTok{\}}
\StringTok{model \{}
\StringTok{  mu\_log\_sigma\_op  \textasciitilde{} normal({-}2, 1);}
\StringTok{  tau\_log\_sigma\_op \textasciitilde{} exponential(2);}
\StringTok{  z\_log\_sigma\_op   \textasciitilde{} std\_normal();}
\StringTok{  mu\_log\_beta\_op   \textasciitilde{} normal({-}1, 0.7);}
\StringTok{  tau\_log\_beta\_op  \textasciitilde{} exponential(2);}
\StringTok{  z\_log\_beta\_op    \textasciitilde{} std\_normal();}
\StringTok{  mu\_log\_beta      \textasciitilde{} normal({-}1, 1);}
\StringTok{  tau\_log\_beta     \textasciitilde{} exponential(1);}
\StringTok{  z\_log\_beta       \textasciitilde{} std\_normal();}
\StringTok{  mu\_bias          \textasciitilde{} normal(0, 1);}
\StringTok{  tau\_bias         \textasciitilde{} exponential(1);}
\StringTok{  z\_bias           \textasciitilde{} std\_normal();}
\StringTok{  for (n in 1:N)}
\StringTok{    choice[n] \textasciitilde{} bernoulli\_logit(}
\StringTok{      (dV[n] + bias[agent[n]]) / exp(log\_beta[agent[n]])}
\StringTok{    );}
\StringTok{\}}
\StringTok{generated quantities \{}
\StringTok{  vector[N] log\_lik;}
\StringTok{  for (n in 1:N)}
\StringTok{    log\_lik[n] = bernoulli\_logit\_lpmf(}
\StringTok{      choice[n] | (dV[n] + bias[agent[n]]) / exp(log\_beta[agent[n]])}
\StringTok{    );}
\StringTok{\}}
\StringTok{"}
\FunctionTok{writeLines}\NormalTok{(stan\_tom1\_hier, here}\SpecialCharTok{::}\FunctionTok{here}\NormalTok{(}\StringTok{"stan"}\NormalTok{, }\StringTok{"ch09\_tom1\_hier.stan"}\NormalTok{))}
\NormalTok{mod\_tom1\_hier }\OtherTok{\textless{}{-}} \FunctionTok{cmdstan\_model}\NormalTok{(here}\SpecialCharTok{::}\FunctionTok{here}\NormalTok{(}\StringTok{"stan"}\NormalTok{, }\StringTok{"ch09\_tom1\_hier.stan"}\NormalTok{))}
\end{Highlighting}
\end{Shaded}

\begin{Shaded}
\begin{Highlighting}[]
\NormalTok{hier\_tom1\_path }\OtherTok{\textless{}{-}}\NormalTok{ here}\SpecialCharTok{::}\FunctionTok{here}\NormalTok{(}\StringTok{"simmodels"}\NormalTok{, }\StringTok{"ch09\_tom1\_hier\_fit.rds"}\NormalTok{)}
\ControlFlowTok{if}\NormalTok{ (regenerate\_fits }\SpecialCharTok{||} \SpecialCharTok{!}\FunctionTok{file.exists}\NormalTok{(hier\_tom1\_path)) \{}
\NormalTok{  fit\_hier\_tom1 }\OtherTok{\textless{}{-}}\NormalTok{ mod\_tom1\_hier}\SpecialCharTok{$}\FunctionTok{sample}\NormalTok{(}
    \AttributeTok{data =}\NormalTok{ sd\_int[}\FunctionTok{c}\NormalTok{(}\StringTok{"N"}\NormalTok{,}\StringTok{"J"}\NormalTok{,}\StringTok{"S"}\NormalTok{,}\StringTok{"species\_of"}\NormalTok{,}\StringTok{"agent"}\NormalTok{,}
                    \StringTok{"choice"}\NormalTok{,}\StringTok{"op\_choice"}\NormalTok{,}\StringTok{"start"}\NormalTok{,}\StringTok{"stop"}\NormalTok{)],}
    \AttributeTok{chains =} \DecValTok{4}\NormalTok{, }\AttributeTok{parallel\_chains =} \DecValTok{4}\NormalTok{,}
    \AttributeTok{iter\_warmup =} \DecValTok{1000}\NormalTok{, }\AttributeTok{iter\_sampling =} \DecValTok{1000}\NormalTok{,}
    \AttributeTok{adapt\_delta =} \FloatTok{0.95}\NormalTok{, }\AttributeTok{max\_treedepth =} \DecValTok{11}\NormalTok{, }\AttributeTok{seed =} \DecValTok{2027}
\NormalTok{  )}
\NormalTok{  fit\_hier\_tom1}\SpecialCharTok{$}\FunctionTok{save\_object}\NormalTok{(hier\_tom1\_path)}
\NormalTok{\} }\ControlFlowTok{else}\NormalTok{ \{}
\NormalTok{  fit\_hier\_tom1 }\OtherTok{\textless{}{-}} \FunctionTok{readRDS}\NormalTok{(hier\_tom1\_path)}
\NormalTok{\}}
\NormalTok{fit\_hier\_tom1}\SpecialCharTok{$}\FunctionTok{diagnostic\_summary}\NormalTok{()}
\end{Highlighting}
\end{Shaded}

\begin{verbatim}
$num_divergent
[1]  39  24  24 167

$num_max_treedepth
[1] 961   0   0 827

$ebfmi
[1] 0.8382165 0.9810230 0.8367260 2.0332266
\end{verbatim}

\section{The main estimand: a joint mixture regression on
representational
capacity}\label{the-main-estimand-a-joint-mixture-regression-on-representational-capacity}

The regression we actually want is not on \texttt{log\_beta}. The
Machiavellian intelligence hypothesis and the cognitive scaffolding
hypothesis are claims about \textbf{representational capacity} ---
whether an agent mentalizes at all, and at what recursion depth. A
regression of \texttt{log\_beta} on ECV and group size within the 0-ToM
architecture answers a related but different question: given that an
agent is a statistical tracker, does brain size or group size predict
how sharply it executes that strategy? A thoughtful but subsidiary
quantity. We will compute it too, but \emph{after} the main analysis.

The main analysis puts architecture itself on the left-hand side.

\subsection{Model specification}\label{model-specification}

For each individual \(j\) introduce a latent architecture indicator
\(z_j \in \{0, 1\}\): \(z_j = 0\) means agent \(j\)'s choices are
generated by 0-ToM, \(z_j = 1\) means they are generated by 1-ToM.
Per-individual the trial likelihood is marginalized over \(z_j\):

\[
p\!\left(c_{j, 1:T_j} \mid \theta^{0}_j, \theta^{1}_j, \pi_{s(j)}\right)
=
\pi_{s(j)} \cdot \prod_t p_{\text{1-ToM}}\!\left(c_{j,t} \mid \theta^{1}_j\right)
+
(1 - \pi_{s(j)}) \cdot \prod_t p_{\text{0-ToM}}\!\left(c_{j,t} \mid \theta^{0}_j\right).
\]

The mixture weight \(\pi_s \in (0, 1)\) is the \emph{species-level
probability that a member of species \(s\) is a 1-ToM agent}. It is
regressed on standardized log covariates with a species residual:

\[
\text{logit}(\pi_s) \;=\; \alpha + b_{\text{ECV}}\,\tilde{\text{ECV}}_s + b_{\text{group}}\,\widetilde{\text{group}}_s + \tau_{\text{sp}}\,u_s,
\qquad u_s \sim \mathcal{N}(0, 1).
\]

The coefficients \(b_{\text{ECV}}\) and \(b_{\text{group}}\) are now on
the \emph{representational-capacity} scale. A negative
\(b_{\text{ECV}}\) would falsify the scaffolding hypothesis on this
sample; a positive, credibly-nonzero \(b_{\text{group}}\) would support
Machiavellian intelligence; similarly positive \(b_{\text{ECV}}\) would
support scaffolding. The mixture structure makes no assumption that all
members of a species share an architecture --- it supplies a prior
probability, and the per-individual data update it to a posterior.

The individual-level parameters are partitioned:

\begin{itemize}
\tightlist
\item
  \textbf{0-ToM-only}: \(\log\sigma_j\) (own volatility over opponent
  bias).
\item
  \textbf{1-ToM-only}: \(\log\sigma_{\text{op}, j}\),
  \(\log\beta_{\text{op}, j}\) (opponent-model parameters the 1-ToM
  agent holds).
\item
  \textbf{Shared across architectures}: \(\log\beta_j\), \(b_j\). These
  are decision-layer parameters (softmax temperature and handedness)
  that map a \(\Delta V\) to a choice, and the mapping is structurally
  identical under 0-ToM and 1-ToM --- only the recursion that produces
  \(\Delta V\) differs. Sharing them prevents the mixture from
  fragmenting identification of parameters that the data cleanly
  constrain.
\end{itemize}

All parameters are hierarchical with species-level means, non-centered
individual offsets, and half-exponential SD priors, exactly as in the
baseline.

\textbf{A note on identifiability.} The posterior for \(\pi_s\) is only
informed to the extent that per-individual 0-ToM and 1-ToM likelihoods
are distinguishable --- which is exactly the leverage the Occam's-razor
panel measured. Individuals whose choice sequences are predictively
equivalent under the two architectures will have diffuse per-individual
posterior architecture probabilities, and information will flow into
\(\pi_s\) primarily from the \emph{other} individuals in that species
and from the covariate regression. This is the correct Bayesian
behavior: when individual data are uninformative, the mixture returns to
the species prior.

\subsection{Stan implementation}\label{stan-implementation-1}

\begin{Shaded}
\begin{Highlighting}[]
\NormalTok{stan\_mixture\_reg }\OtherTok{\textless{}{-}} \StringTok{"}
\StringTok{data \{}
\StringTok{  int\textless{}lower=1\textgreater{} N;}
\StringTok{  int\textless{}lower=1\textgreater{} J;}
\StringTok{  int\textless{}lower=1\textgreater{} S;}
\StringTok{  array[J] int\textless{}lower=1, upper=S\textgreater{} species\_of;}
\StringTok{  array[N] int\textless{}lower=1, upper=J\textgreater{} agent;}
\StringTok{  array[N] int\textless{}lower=0, upper=1\textgreater{} choice;}
\StringTok{  array[N] int\textless{}lower=0, upper=1\textgreater{} op\_choice;}
\StringTok{  array[J] int\textless{}lower=1\textgreater{} start;}
\StringTok{  array[J] int\textless{}lower=1\textgreater{} stop;}
\StringTok{  vector[S] ecv\_z;}
\StringTok{  vector[S] group\_z;}
\StringTok{\}}
\StringTok{parameters \{}
\StringTok{  // 0{-}ToM{-}only}
\StringTok{  vector[S] mu\_log\_sigma;}
\StringTok{  real\textless{}lower=0\textgreater{} tau\_log\_sigma;}
\StringTok{  vector[J] z\_log\_sigma;}
\StringTok{  // 1{-}ToM{-}only}
\StringTok{  vector[S] mu\_log\_sigma\_op;}
\StringTok{  real\textless{}lower=0\textgreater{} tau\_log\_sigma\_op;}
\StringTok{  vector[J] z\_log\_sigma\_op;}
\StringTok{  vector[S] mu\_log\_beta\_op;}
\StringTok{  real\textless{}lower=0\textgreater{} tau\_log\_beta\_op;}
\StringTok{  vector[J] z\_log\_beta\_op;}
\StringTok{  // Shared decision layer}
\StringTok{  vector[S] mu\_log\_beta;}
\StringTok{  real\textless{}lower=0\textgreater{} tau\_log\_beta;}
\StringTok{  vector[J] z\_log\_beta;}
\StringTok{  vector[S] mu\_bias;}
\StringTok{  real\textless{}lower=0\textgreater{} tau\_bias;}
\StringTok{  vector[J] z\_bias;}
\StringTok{  // Mixture regression on architecture probability}
\StringTok{  real alpha\_pi;}
\StringTok{  real b\_ecv;}
\StringTok{  real b\_group;}
\StringTok{  real\textless{}lower=0\textgreater{} tau\_sp\_pi;}
\StringTok{  vector[S] u\_sp;}
\StringTok{\}}
\StringTok{transformed parameters \{}
\StringTok{  vector[J] log\_sigma    = mu\_log\_sigma[species\_of]    + tau\_log\_sigma    * z\_log\_sigma;}
\StringTok{  vector[J] log\_sigma\_op = mu\_log\_sigma\_op[species\_of] + tau\_log\_sigma\_op * z\_log\_sigma\_op;}
\StringTok{  vector[J] log\_beta\_op  = mu\_log\_beta\_op[species\_of]  + tau\_log\_beta\_op  * z\_log\_beta\_op;}
\StringTok{  vector[J] log\_beta     = mu\_log\_beta[species\_of]     + tau\_log\_beta     * z\_log\_beta;}
\StringTok{  vector[J] bias\_j       = mu\_bias[species\_of]         + tau\_bias         * z\_bias;}
\StringTok{  vector[S] logit\_pi     = alpha\_pi + b\_ecv * ecv\_z + b\_group * group\_z}
\StringTok{                           + tau\_sp\_pi * u\_sp;}
\StringTok{  vector[S] pi\_s         = inv\_logit(logit\_pi);}
\StringTok{  vector[N] dV0;}
\StringTok{  vector[N] dV1;}
\StringTok{  for (j in 1:J) \{}
\StringTok{    // 0{-}ToM recursion: belief about opponent bias, updated by op\_choice}
\StringTok{    \{}
\StringTok{      real sig0 = exp(log\_sigma[j]);}
\StringTok{      real mu   = 0.0;}
\StringTok{      real Sig  = 1.0;}
\StringTok{      for (k in start[j]:stop[j]) \{}
\StringTok{        real s\_mu   = inv\_logit(mu);}
\StringTok{        real p\_op\_0 = inv\_logit(mu / sqrt(1 + 0.36 * Sig));}
\StringTok{        dV0[k] = 2 * p\_op\_0 {-} 1;}
\StringTok{        Sig = 1.0 / (1.0 / (Sig + sig0) + s\_mu * (1 {-} s\_mu));}
\StringTok{        mu  = mu + Sig * (op\_choice[k] {-} s\_mu);}
\StringTok{      \}}
\StringTok{    \}}
\StringTok{    // 1{-}ToM recursion: belief about opponent\textquotesingle{}s belief about self,}
\StringTok{    // updated by own choice}
\StringTok{    \{}
\StringTok{      real sig\_op = exp(log\_sigma\_op[j]);}
\StringTok{      real b\_op   = exp(log\_beta\_op[j]);}
\StringTok{      real mu\_s   = 0.0;}
\StringTok{      real Sig\_s  = 1.0;}
\StringTok{      for (k in start[j]:stop[j]) \{}
\StringTok{        real p\_self\_from\_op = inv\_logit(mu\_s / sqrt(1 + 0.36 * Sig\_s));}
\StringTok{        real p\_op\_1         = inv\_logit((1 {-} 2 * p\_self\_from\_op) / b\_op);}
\StringTok{        dV1[k] = 2 * p\_op\_1 {-} 1;}
\StringTok{        real s\_ps = p\_self\_from\_op;}
\StringTok{        Sig\_s = 1.0 / (1.0 / (Sig\_s + sig\_op) + s\_ps * (1 {-} s\_ps));}
\StringTok{        mu\_s  = mu\_s + Sig\_s * (choice[k] {-} s\_ps);}
\StringTok{      \}}
\StringTok{    \}}
\StringTok{  \}}
\StringTok{\}}
\StringTok{model \{}
\StringTok{  // Hierarchical priors}
\StringTok{  mu\_log\_sigma    \textasciitilde{} normal({-}2, 1);}
\StringTok{  tau\_log\_sigma   \textasciitilde{} exponential(2);}
\StringTok{  z\_log\_sigma     \textasciitilde{} std\_normal();}
\StringTok{  mu\_log\_sigma\_op \textasciitilde{} normal({-}2, 1);}
\StringTok{  tau\_log\_sigma\_op \textasciitilde{} exponential(2);}
\StringTok{  z\_log\_sigma\_op  \textasciitilde{} std\_normal();}
\StringTok{  mu\_log\_beta\_op  \textasciitilde{} normal({-}1, 0.7);}
\StringTok{  tau\_log\_beta\_op \textasciitilde{} exponential(2);}
\StringTok{  z\_log\_beta\_op   \textasciitilde{} std\_normal();}
\StringTok{  mu\_log\_beta     \textasciitilde{} normal({-}1, 1);}
\StringTok{  tau\_log\_beta    \textasciitilde{} exponential(1);}
\StringTok{  z\_log\_beta      \textasciitilde{} std\_normal();}
\StringTok{  mu\_bias         \textasciitilde{} normal(0, 1);}
\StringTok{  tau\_bias        \textasciitilde{} exponential(1);}
\StringTok{  z\_bias          \textasciitilde{} std\_normal();}
\StringTok{  // Mixture regression}
\StringTok{  alpha\_pi  \textasciitilde{} normal(0, 1);}
\StringTok{  b\_ecv     \textasciitilde{} normal(0, 0.7);}
\StringTok{  b\_group   \textasciitilde{} normal(0, 0.7);}
\StringTok{  tau\_sp\_pi \textasciitilde{} exponential(1);}
\StringTok{  u\_sp      \textasciitilde{} std\_normal();}
\StringTok{  // Marginalized likelihood: one log\_mix per individual}
\StringTok{  for (j in 1:J) \{}
\StringTok{    real ll0 = 0;}
\StringTok{    real ll1 = 0;}
\StringTok{    real inv\_b = 1.0 / exp(log\_beta[j]);}
\StringTok{    for (k in start[j]:stop[j]) \{}
\StringTok{      ll0 += bernoulli\_logit\_lpmf(choice[k] |}
\StringTok{                (dV0[k] + bias\_j[j]) * inv\_b);}
\StringTok{      ll1 += bernoulli\_logit\_lpmf(choice[k] |}
\StringTok{                (dV1[k] + bias\_j[j]) * inv\_b);}
\StringTok{    \}}
\StringTok{    target += log\_mix(pi\_s[species\_of[j]], ll1, ll0);}
\StringTok{  \}}
\StringTok{\}}
\StringTok{generated quantities \{}
\StringTok{  // Per{-}individual posterior probability z\_j = 1 (= 1{-}ToM)}
\StringTok{  vector[J] post\_pi\_j;}
\StringTok{  // Per{-}individual total log{-}likelihood difference (1{-}ToM {-} 0{-}ToM).}
\StringTok{  vector[J] ll\_diff\_j;}
\StringTok{  // Per{-}trial log\_lik under the marginal mixture for PSIS{-}LOO}
\StringTok{  vector[N] log\_lik;}
\StringTok{  for (j in 1:J) \{}
\StringTok{    real ll0 = 0;}
\StringTok{    real ll1 = 0;}
\StringTok{    real inv\_b = 1.0 / exp(log\_beta[j]);}
\StringTok{    for (k in start[j]:stop[j]) \{}
\StringTok{      ll0 += bernoulli\_logit\_lpmf(choice[k] |}
\StringTok{                (dV0[k] + bias\_j[j]) * inv\_b);}
\StringTok{      ll1 += bernoulli\_logit\_lpmf(choice[k] |}
\StringTok{                (dV1[k] + bias\_j[j]) * inv\_b);}
\StringTok{    \}}
\StringTok{    ll\_diff\_j[j] = ll1 {-} ll0;}
\StringTok{    real lp0 = log1m(pi\_s[species\_of[j]]) + ll0;}
\StringTok{    real lp1 = log(pi\_s[species\_of[j]])   + ll1;}
\StringTok{    post\_pi\_j[j] = exp(lp1 {-} log\_sum\_exp(lp0, lp1));}
\StringTok{    for (k in start[j]:stop[j]) \{}
\StringTok{      real l0 = bernoulli\_logit\_lpmf(choice[k] |}
\StringTok{                  (dV0[k] + bias\_j[j]) * inv\_b);}
\StringTok{      real l1 = bernoulli\_logit\_lpmf(choice[k] |}
\StringTok{                  (dV1[k] + bias\_j[j]) * inv\_b);}
\StringTok{      log\_lik[k] = log\_mix(pi\_s[species\_of[j]], l1, l0);}
\StringTok{    \}}
\StringTok{  \}}
\StringTok{\}}
\StringTok{"}
\FunctionTok{writeLines}\NormalTok{(stan\_mixture\_reg, here}\SpecialCharTok{::}\FunctionTok{here}\NormalTok{(}\StringTok{"stan"}\NormalTok{, }\StringTok{"ch09\_mixture\_reg.stan"}\NormalTok{))}
\NormalTok{mod\_mix }\OtherTok{\textless{}{-}} \FunctionTok{cmdstan\_model}\NormalTok{(here}\SpecialCharTok{::}\FunctionTok{here}\NormalTok{(}\StringTok{"stan"}\NormalTok{, }\StringTok{"ch09\_mixture\_reg.stan"}\NormalTok{))}
\end{Highlighting}
\end{Shaded}

\begin{Shaded}
\begin{Highlighting}[]
\CommentTok{\# Standardize log covariates at the species level}
\NormalTok{species\_cov }\OtherTok{\textless{}{-}}\NormalTok{ sd\_int}\SpecialCharTok{$}\NormalTok{species\_lookup }\SpecialCharTok{|\textgreater{}}
  \FunctionTok{left\_join}\NormalTok{(covariates, }\AttributeTok{by =} \StringTok{"species"}\NormalTok{) }\SpecialCharTok{|\textgreater{}}
  \FunctionTok{arrange}\NormalTok{(species\_id) }\SpecialCharTok{|\textgreater{}}
  \FunctionTok{mutate}\NormalTok{(}
    \AttributeTok{log\_ECV =} \FunctionTok{log}\NormalTok{(ECV),}
    \AttributeTok{log\_grp =} \FunctionTok{log}\NormalTok{(groupsize),}
    \AttributeTok{ecv\_z   =} \FunctionTok{as.numeric}\NormalTok{(}\FunctionTok{scale}\NormalTok{(log\_ECV)),}
    \AttributeTok{group\_z =} \FunctionTok{as.numeric}\NormalTok{(}\FunctionTok{scale}\NormalTok{(log\_grp))}
\NormalTok{  )}

\NormalTok{sd\_mix }\OtherTok{\textless{}{-}} \FunctionTok{c}\NormalTok{(}
\NormalTok{  sd\_int[}\FunctionTok{c}\NormalTok{(}\StringTok{"N"}\NormalTok{,}\StringTok{"J"}\NormalTok{,}\StringTok{"S"}\NormalTok{,}\StringTok{"species\_of"}\NormalTok{,}\StringTok{"agent"}\NormalTok{,}
           \StringTok{"choice"}\NormalTok{,}\StringTok{"op\_choice"}\NormalTok{,}\StringTok{"start"}\NormalTok{,}\StringTok{"stop"}\NormalTok{)],}
  \FunctionTok{list}\NormalTok{(}\AttributeTok{ecv\_z   =}\NormalTok{ species\_cov}\SpecialCharTok{$}\NormalTok{ecv\_z,}
       \AttributeTok{group\_z =}\NormalTok{ species\_cov}\SpecialCharTok{$}\NormalTok{group\_z)}
\NormalTok{)}
\end{Highlighting}
\end{Shaded}

\begin{Shaded}
\begin{Highlighting}[]
\NormalTok{mix\_path }\OtherTok{\textless{}{-}}\NormalTok{ here}\SpecialCharTok{::}\FunctionTok{here}\NormalTok{(}\StringTok{"simmodels"}\NormalTok{, }\StringTok{"ch09\_mixture\_reg\_fit.rds"}\NormalTok{)}
\ControlFlowTok{if}\NormalTok{ (regenerate\_fits }\SpecialCharTok{||} \SpecialCharTok{!}\FunctionTok{file.exists}\NormalTok{(mix\_path)) \{}
\NormalTok{  fit\_mix }\OtherTok{\textless{}{-}}\NormalTok{ mod\_mix}\SpecialCharTok{$}\FunctionTok{sample}\NormalTok{(}
    \AttributeTok{data =}\NormalTok{ sd\_mix,}
    \AttributeTok{chains =} \DecValTok{4}\NormalTok{, }\AttributeTok{parallel\_chains =} \DecValTok{4}\NormalTok{,}
    \AttributeTok{iter\_warmup =} \DecValTok{1200}\NormalTok{, }\AttributeTok{iter\_sampling =} \DecValTok{1200}\NormalTok{,}
    \AttributeTok{adapt\_delta =} \FloatTok{0.97}\NormalTok{, }\AttributeTok{max\_treedepth =} \DecValTok{12}\NormalTok{,}
    \AttributeTok{seed =} \DecValTok{2029}
\NormalTok{  )}
\NormalTok{  fit\_mix}\SpecialCharTok{$}\FunctionTok{save\_object}\NormalTok{(mix\_path)}
\NormalTok{\} }\ControlFlowTok{else}\NormalTok{ \{}
\NormalTok{  fit\_mix }\OtherTok{\textless{}{-}} \FunctionTok{readRDS}\NormalTok{(mix\_path)}
\NormalTok{\}}
\NormalTok{fit\_mix}\SpecialCharTok{$}\FunctionTok{diagnostic\_summary}\NormalTok{()}
\end{Highlighting}
\end{Shaded}

\begin{verbatim}
$num_divergent
[1] 2 1 2 0

$num_max_treedepth
[1] 0 0 0 0

$ebfmi
[1] 0.8750412 0.8685706 0.8793667 0.8039244
\end{verbatim}

Expect this fit to take roughly 2× the baseline hierarchical 0-ToM
runtime --- both recursions run every leapfrog step.
\texttt{adapt\_delta\ =\ 0.97} and a deeper \texttt{max\_treedepth} help
with the mild funnel geometry the mixture introduces around
\(\tau_{\text{sp}}\) when species have few individuals.

\subsection{The answer}\label{the-answer}

\begin{Shaded}
\begin{Highlighting}[]
\NormalTok{mix\_coef }\OtherTok{\textless{}{-}}\NormalTok{ fit\_mix}\SpecialCharTok{$}\FunctionTok{draws}\NormalTok{(}\FunctionTok{c}\NormalTok{(}\StringTok{"alpha\_pi"}\NormalTok{,}\StringTok{"b\_ecv"}\NormalTok{,}\StringTok{"b\_group"}\NormalTok{,}\StringTok{"tau\_sp\_pi"}\NormalTok{),}
                          \AttributeTok{format =} \StringTok{"df"}\NormalTok{) }\SpecialCharTok{|\textgreater{}}
  \FunctionTok{as\_tibble}\NormalTok{()}

\NormalTok{p\_mix\_coef }\OtherTok{\textless{}{-}}\NormalTok{ mix\_coef }\SpecialCharTok{|\textgreater{}}
\NormalTok{  dplyr}\SpecialCharTok{::}\FunctionTok{select}\NormalTok{(b\_ecv, b\_group) }\SpecialCharTok{|\textgreater{}}
  \FunctionTok{pivot\_longer}\NormalTok{(}\FunctionTok{everything}\NormalTok{(), }\AttributeTok{names\_to =} \StringTok{"coef"}\NormalTok{, }\AttributeTok{values\_to =} \StringTok{"value"}\NormalTok{) }\SpecialCharTok{|\textgreater{}}
  \FunctionTok{ggplot}\NormalTok{(}\FunctionTok{aes}\NormalTok{(}\AttributeTok{x =}\NormalTok{ value, }\AttributeTok{fill =}\NormalTok{ coef)) }\SpecialCharTok{+}
  \FunctionTok{geom\_density}\NormalTok{(}\AttributeTok{alpha =} \FloatTok{0.6}\NormalTok{, }\AttributeTok{color =} \ConstantTok{NA}\NormalTok{) }\SpecialCharTok{+}
  \FunctionTok{geom\_vline}\NormalTok{(}\AttributeTok{xintercept =} \DecValTok{0}\NormalTok{, }\AttributeTok{linetype =} \StringTok{"dashed"}\NormalTok{, }\AttributeTok{color =} \StringTok{"gray30"}\NormalTok{) }\SpecialCharTok{+}
  \FunctionTok{scale\_fill\_manual}\NormalTok{(}\AttributeTok{values =} \FunctionTok{c}\NormalTok{(}\StringTok{"b\_ecv"} \OtherTok{=} \StringTok{"steelblue"}\NormalTok{,}
                               \StringTok{"b\_group"} \OtherTok{=} \StringTok{"firebrick3"}\NormalTok{),}
                    \AttributeTok{labels =} \FunctionTok{c}\NormalTok{(}\StringTok{"b\_ecv"}   \OtherTok{=} \StringTok{"ECV (scaffolding)"}\NormalTok{,}
                               \StringTok{"b\_group"} \OtherTok{=} \StringTok{"Group size (Machiavellian)"}\NormalTok{)) }\SpecialCharTok{+}
  \FunctionTok{labs}\NormalTok{(}\AttributeTok{x =} \FunctionTok{expression}\NormalTok{(beta[cov]}\SpecialCharTok{\textasciitilde{}}\StringTok{"(logit "}\SpecialCharTok{\textasciitilde{}}\NormalTok{pi[s]}\SpecialCharTok{\textasciitilde{}}\StringTok{"per SD of log covariate)"}\NormalTok{),}
       \AttributeTok{y =} \StringTok{"Posterior density"}\NormalTok{, }\AttributeTok{fill =} \StringTok{"Hypothesis"}\NormalTok{,}
       \AttributeTok{title =} \StringTok{"Which covariate drives representational capacity?"}\NormalTok{,}
       \AttributeTok{subtitle =} \StringTok{"Positive = larger covariate associated with more 1{-}ToM{-}type individuals"}\NormalTok{)}

\NormalTok{pi\_s\_summary }\OtherTok{\textless{}{-}}\NormalTok{ fit\_mix}\SpecialCharTok{$}\FunctionTok{draws}\NormalTok{(}\StringTok{"pi\_s"}\NormalTok{, }\AttributeTok{format =} \StringTok{"df"}\NormalTok{) }\SpecialCharTok{|\textgreater{}}
  \FunctionTok{as\_tibble}\NormalTok{() }\SpecialCharTok{|\textgreater{}}
\NormalTok{  dplyr}\SpecialCharTok{::}\FunctionTok{select}\NormalTok{(}\FunctionTok{starts\_with}\NormalTok{(}\StringTok{"pi\_s"}\NormalTok{)) }\SpecialCharTok{|\textgreater{}}
  \FunctionTok{pivot\_longer}\NormalTok{(}\FunctionTok{everything}\NormalTok{(), }\AttributeTok{names\_to =} \StringTok{"param"}\NormalTok{, }\AttributeTok{values\_to =} \StringTok{"value"}\NormalTok{) }\SpecialCharTok{|\textgreater{}}
  \FunctionTok{mutate}\NormalTok{(}\AttributeTok{species\_id =} \FunctionTok{as.integer}\NormalTok{(}\FunctionTok{str\_extract}\NormalTok{(param, }\StringTok{"}\SpecialCharTok{\textbackslash{}\textbackslash{}}\StringTok{d+"}\NormalTok{))) }\SpecialCharTok{|\textgreater{}}
  \FunctionTok{group\_by}\NormalTok{(species\_id) }\SpecialCharTok{|\textgreater{}}
  \FunctionTok{summarize}\NormalTok{(}\AttributeTok{med =} \FunctionTok{median}\NormalTok{(value),}
            \AttributeTok{lo  =} \FunctionTok{quantile}\NormalTok{(value, }\FloatTok{0.05}\NormalTok{),}
            \AttributeTok{hi  =} \FunctionTok{quantile}\NormalTok{(value, }\FloatTok{0.95}\NormalTok{),}
            \AttributeTok{.groups =} \StringTok{"drop"}\NormalTok{) }\SpecialCharTok{|\textgreater{}}
  \FunctionTok{left\_join}\NormalTok{(species\_cov }\SpecialCharTok{|\textgreater{}}\NormalTok{ dplyr}\SpecialCharTok{::}\FunctionTok{select}\NormalTok{(species\_id, species),}
            \AttributeTok{by =} \StringTok{"species\_id"}\NormalTok{) }\SpecialCharTok{|\textgreater{}}
  \FunctionTok{arrange}\NormalTok{(med) }\SpecialCharTok{|\textgreater{}}
  \FunctionTok{mutate}\NormalTok{(}\AttributeTok{species =} \FunctionTok{factor}\NormalTok{(species, }\AttributeTok{levels =}\NormalTok{ species))}

\NormalTok{p\_pi\_s }\OtherTok{\textless{}{-}} \FunctionTok{ggplot}\NormalTok{(pi\_s\_summary, }\FunctionTok{aes}\NormalTok{(}\AttributeTok{x =}\NormalTok{ med, }\AttributeTok{y =}\NormalTok{ species)) }\SpecialCharTok{+}
  \FunctionTok{geom\_errorbarh}\NormalTok{(}\FunctionTok{aes}\NormalTok{(}\AttributeTok{xmin =}\NormalTok{ lo, }\AttributeTok{xmax =}\NormalTok{ hi), }\AttributeTok{height =} \DecValTok{0}\NormalTok{, }\AttributeTok{color =} \StringTok{"steelblue"}\NormalTok{) }\SpecialCharTok{+}
  \FunctionTok{geom\_point}\NormalTok{(}\AttributeTok{size =} \DecValTok{3}\NormalTok{, }\AttributeTok{color =} \StringTok{"midnightblue"}\NormalTok{) }\SpecialCharTok{+}
  \FunctionTok{geom\_vline}\NormalTok{(}\AttributeTok{xintercept =} \FloatTok{0.5}\NormalTok{, }\AttributeTok{linetype =} \StringTok{"dashed"}\NormalTok{, }\AttributeTok{color =} \StringTok{"gray40"}\NormalTok{) }\SpecialCharTok{+}
  \FunctionTok{scale\_x\_continuous}\NormalTok{(}\AttributeTok{limits =} \FunctionTok{c}\NormalTok{(}\DecValTok{0}\NormalTok{, }\DecValTok{1}\NormalTok{)) }\SpecialCharTok{+}
  \FunctionTok{labs}\NormalTok{(}\AttributeTok{x =} \FunctionTok{expression}\NormalTok{(pi[s]}\SpecialCharTok{\textasciitilde{}}\StringTok{"(species{-}level P(1{-}ToM))"}\NormalTok{),}
       \AttributeTok{y =} \ConstantTok{NULL}\NormalTok{,}
       \AttributeTok{title =} \StringTok{"Per{-}species architecture posterior"}\NormalTok{,}
       \AttributeTok{subtitle =} \StringTok{"Dashed = indifference; right of line = species leans 1{-}ToM"}\NormalTok{)}

\NormalTok{p\_mix\_coef }\SpecialCharTok{+}\NormalTok{ p\_pi\_s }\SpecialCharTok{+} \FunctionTok{plot\_layout}\NormalTok{(}\AttributeTok{widths =} \FunctionTok{c}\NormalTok{(}\DecValTok{1}\NormalTok{, }\DecValTok{1}\NormalTok{))}
\end{Highlighting}
\end{Shaded}

\begin{figure}[H]

{\centering \includegraphics[width=0.85\linewidth,height=\textheight,keepaspectratio]{09-TheoryOfMind_files/figure-pdf/ch09_mixture_coef_plot-1.pdf}

}

\caption{Left: posterior distributions of the two covariate effects on
logit(π\_s), the species-level probability of being a 1-ToM agent.
Positive = more 1-ToM-compatible as the covariate increases. Right:
species-level posterior π\_s (dot = median, bar = 90\% CI).}

\end{figure}%

\begin{Shaded}
\begin{Highlighting}[]
\NormalTok{mix\_coef\_summary }\OtherTok{\textless{}{-}}\NormalTok{ mix\_coef }\SpecialCharTok{|\textgreater{}}
\NormalTok{  dplyr}\SpecialCharTok{::}\FunctionTok{select}\NormalTok{(b\_ecv, b\_group) }\SpecialCharTok{|\textgreater{}}
  \FunctionTok{pivot\_longer}\NormalTok{(}\FunctionTok{everything}\NormalTok{(), }\AttributeTok{names\_to =} \StringTok{"coef"}\NormalTok{, }\AttributeTok{values\_to =} \StringTok{"value"}\NormalTok{) }\SpecialCharTok{|\textgreater{}}
  \FunctionTok{group\_by}\NormalTok{(coef) }\SpecialCharTok{|\textgreater{}}
  \FunctionTok{summarize}\NormalTok{(}
    \AttributeTok{mean   =} \FunctionTok{mean}\NormalTok{(value),}
    \AttributeTok{median =} \FunctionTok{median}\NormalTok{(value),}
    \AttributeTok{q05    =} \FunctionTok{quantile}\NormalTok{(value, }\FloatTok{0.05}\NormalTok{),}
    \AttributeTok{q95    =} \FunctionTok{quantile}\NormalTok{(value, }\FloatTok{0.95}\NormalTok{),}
    \AttributeTok{p\_neg  =} \FunctionTok{mean}\NormalTok{(value }\SpecialCharTok{\textless{}} \DecValTok{0}\NormalTok{),}
    \AttributeTok{p\_pos  =} \FunctionTok{mean}\NormalTok{(value }\SpecialCharTok{\textgreater{}} \DecValTok{0}\NormalTok{),}
    \AttributeTok{.groups =} \StringTok{"drop"}
\NormalTok{  )}
\NormalTok{knitr}\SpecialCharTok{::}\FunctionTok{kable}\NormalTok{(mix\_coef\_summary, }\AttributeTok{digits =} \DecValTok{3}\NormalTok{,}
             \AttributeTok{caption =} \StringTok{"Coefficients on logit(π\_s). p\_pos = posterior probability that the effect is positive."}\NormalTok{)}
\end{Highlighting}
\end{Shaded}

\begin{longtable}[]{@{}lrrrrrr@{}}
\caption{Coefficients on logit(π\_s). p\_pos = posterior probability
that the effect is positive.}\tabularnewline
\toprule\noalign{}
coef & mean & median & q05 & q95 & p\_neg & p\_pos \\
\midrule\noalign{}
\endfirsthead
\toprule\noalign{}
coef & mean & median & q05 & q95 & p\_neg & p\_pos \\
\midrule\noalign{}
\endhead
\bottomrule\noalign{}
\endlastfoot
b\_ecv & -0.254 & -0.250 & -1.161 & 0.677 & 0.679 & 0.321 \\
b\_group & 0.193 & 0.195 & -0.722 & 1.115 & 0.362 & 0.638 \\
\end{longtable}

Reading the coefficients on the representational-capacity scale:
\(b_{\text{ECV}}\) is the change in \(\text{logit}(\pi_s)\) per standard
deviation of log endocranial volume. A value of \(+1\) corresponds to
roughly a factor-of-2.7 increase in the odds of an individual being a
1-ToM type, per SD increase in log ECV. The practical interpretation
depends on where the posterior mass sits relative to zero --- this is
reported per-fit by \texttt{p\_pos} in the table.

\subsection{Per-individual architecture
posteriors}\label{per-individual-architecture-posteriors}

The mixture also delivers per-individual posterior probabilities
\(\Pr(z_j = 1 \mid \text{data})\). Individuals whose data pull their
probability away from the species prior are those whose choice sequences
decisively fit one architecture over the other; those whose posteriors
return to the species prior \(\pi_{s(j)}\) are predictively
indistinguishable under the two models, consistent with the
Occam's-razor evidence that matching pennies has limited leverage at the
single-individual scale.

\begin{Shaded}
\begin{Highlighting}[]
\NormalTok{post\_pi\_draws }\OtherTok{\textless{}{-}}\NormalTok{ fit\_mix}\SpecialCharTok{$}\FunctionTok{draws}\NormalTok{(}\StringTok{"post\_pi\_j"}\NormalTok{, }\AttributeTok{format =} \StringTok{"df"}\NormalTok{) }\SpecialCharTok{|\textgreater{}}
  \FunctionTok{as\_tibble}\NormalTok{() }\SpecialCharTok{|\textgreater{}}
\NormalTok{  dplyr}\SpecialCharTok{::}\FunctionTok{select}\NormalTok{(}\FunctionTok{starts\_with}\NormalTok{(}\StringTok{"post\_pi\_j"}\NormalTok{)) }\SpecialCharTok{|\textgreater{}}
  \FunctionTok{pivot\_longer}\NormalTok{(}\FunctionTok{everything}\NormalTok{(), }\AttributeTok{names\_to =} \StringTok{"param"}\NormalTok{, }\AttributeTok{values\_to =} \StringTok{"value"}\NormalTok{) }\SpecialCharTok{|\textgreater{}}
  \FunctionTok{mutate}\NormalTok{(}\AttributeTok{agent\_id =} \FunctionTok{as.integer}\NormalTok{(}\FunctionTok{str\_extract}\NormalTok{(param, }\StringTok{"}\SpecialCharTok{\textbackslash{}\textbackslash{}}\StringTok{d+"}\NormalTok{)))}

\NormalTok{post\_pi\_ind }\OtherTok{\textless{}{-}}\NormalTok{ post\_pi\_draws }\SpecialCharTok{|\textgreater{}}
  \FunctionTok{group\_by}\NormalTok{(agent\_id) }\SpecialCharTok{|\textgreater{}}
  \FunctionTok{summarize}\NormalTok{(}\AttributeTok{med =} \FunctionTok{median}\NormalTok{(value),}
            \AttributeTok{lo  =} \FunctionTok{quantile}\NormalTok{(value, }\FloatTok{0.05}\NormalTok{),}
            \AttributeTok{hi  =} \FunctionTok{quantile}\NormalTok{(value, }\FloatTok{0.95}\NormalTok{),}
            \AttributeTok{.groups =} \StringTok{"drop"}\NormalTok{) }\SpecialCharTok{|\textgreater{}}
  \FunctionTok{left\_join}\NormalTok{(sd\_int}\SpecialCharTok{$}\NormalTok{agent\_lookup, }\AttributeTok{by =} \StringTok{"agent\_id"}\NormalTok{) }\SpecialCharTok{|\textgreater{}}
  \FunctionTok{arrange}\NormalTok{(species, med) }\SpecialCharTok{|\textgreater{}}
  \FunctionTok{mutate}\NormalTok{(}\AttributeTok{agent\_id\_ord =} \FunctionTok{row\_number}\NormalTok{())}

\FunctionTok{ggplot}\NormalTok{(post\_pi\_ind, }\FunctionTok{aes}\NormalTok{(}\AttributeTok{x =}\NormalTok{ agent\_id\_ord, }\AttributeTok{y =}\NormalTok{ med, }\AttributeTok{color =}\NormalTok{ species)) }\SpecialCharTok{+}
  \FunctionTok{geom\_errorbar}\NormalTok{(}\FunctionTok{aes}\NormalTok{(}\AttributeTok{ymin =}\NormalTok{ lo, }\AttributeTok{ymax =}\NormalTok{ hi), }\AttributeTok{width =} \DecValTok{0}\NormalTok{, }\AttributeTok{alpha =} \FloatTok{0.5}\NormalTok{) }\SpecialCharTok{+}
  \FunctionTok{geom\_point}\NormalTok{(}\AttributeTok{size =} \FloatTok{1.8}\NormalTok{) }\SpecialCharTok{+}
  \FunctionTok{geom\_hline}\NormalTok{(}\AttributeTok{yintercept =} \FloatTok{0.5}\NormalTok{, }\AttributeTok{linetype =} \StringTok{"dashed"}\NormalTok{, }\AttributeTok{color =} \StringTok{"gray40"}\NormalTok{) }\SpecialCharTok{+}
  \FunctionTok{labs}\NormalTok{(}\AttributeTok{x =} \StringTok{"Individual (ordered by species, then posterior)"}\NormalTok{,}
       \AttributeTok{y =} \FunctionTok{expression}\NormalTok{(}\FunctionTok{Pr}\NormalTok{(z[j]}\SpecialCharTok{==}\DecValTok{1}\SpecialCharTok{\textasciitilde{}}\StringTok{"|"}\SpecialCharTok{\textasciitilde{}}\NormalTok{data)),}
       \AttributeTok{color =} \StringTok{"Species"}\NormalTok{,}
       \AttributeTok{title =} \StringTok{"Per{-}individual posterior architecture probability"}\NormalTok{,}
       \AttributeTok{subtitle =} \StringTok{"Dashed = indifference; wide bars indicate low per{-}individual leverage"}\NormalTok{)}
\end{Highlighting}
\end{Shaded}

\begin{figure}[H]

{\centering \includegraphics[width=0.85\linewidth,height=\textheight,keepaspectratio]{09-TheoryOfMind_files/figure-pdf/ch09_post_pi_ind-1.pdf}

}

\caption{Per-individual posterior probability of being a 1-ToM agent,
colored by species. Individuals close to 0 or 1 are cleanly classified
by their data; those clustered near their species π\_s reflect the
limited per-individual leverage of matching pennies.}

\end{figure}%

\subsection{Sanity check against Occam's
razor}\label{sanity-check-against-occams-razor}

Our Stan \texttt{generated\ quantities} block exposes two tightly-linked
per-individual quantities: \texttt{ll\_diff\_j} (the total
log-likelihood difference between the 1-ToM and 0-ToM arms at each
posterior draw's parameters) and \texttt{post\_pi\_j} (the
Rao-Blackwellized posterior probability that individual \(j\) is a 1-ToM
agent). If the mixture is doing what it claims, individuals with larger
\texttt{ll\_diff\_j} should have higher \texttt{post\_pi\_j}, and the
relationship should be a smooth logistic curve rising from the species
prior \(\pi_{s(j)}\) at zero evidence toward 1 at strong positive
evidence. That curve is the empirical operating characteristic of the
mixture, and it is the correct check to run --- it confirms the mixture
responds to individual-level data, not only to the species-level
covariates.

\begin{Shaded}
\begin{Highlighting}[]
\NormalTok{ll\_diff\_summary }\OtherTok{\textless{}{-}}\NormalTok{ fit\_mix}\SpecialCharTok{$}\FunctionTok{draws}\NormalTok{(}\StringTok{"ll\_diff\_j"}\NormalTok{, }\AttributeTok{format =} \StringTok{"df"}\NormalTok{) }\SpecialCharTok{|\textgreater{}}
  \FunctionTok{as\_tibble}\NormalTok{() }\SpecialCharTok{|\textgreater{}}
\NormalTok{  dplyr}\SpecialCharTok{::}\FunctionTok{select}\NormalTok{(}\FunctionTok{starts\_with}\NormalTok{(}\StringTok{"ll\_diff\_j"}\NormalTok{)) }\SpecialCharTok{|\textgreater{}}
  \FunctionTok{pivot\_longer}\NormalTok{(}\FunctionTok{everything}\NormalTok{(), }\AttributeTok{names\_to =} \StringTok{"param"}\NormalTok{, }\AttributeTok{values\_to =} \StringTok{"value"}\NormalTok{) }\SpecialCharTok{|\textgreater{}}
  \FunctionTok{mutate}\NormalTok{(}\AttributeTok{agent\_id =} \FunctionTok{as.integer}\NormalTok{(}\FunctionTok{str\_extract}\NormalTok{(param, }\StringTok{"}\SpecialCharTok{\textbackslash{}\textbackslash{}}\StringTok{d+"}\NormalTok{))) }\SpecialCharTok{|\textgreater{}}
  \FunctionTok{group\_by}\NormalTok{(agent\_id) }\SpecialCharTok{|\textgreater{}}
  \FunctionTok{summarize}\NormalTok{(}\AttributeTok{ll\_diff\_mean =} \FunctionTok{mean}\NormalTok{(value),}
            \AttributeTok{ll\_diff\_lo   =} \FunctionTok{quantile}\NormalTok{(value, }\FloatTok{0.05}\NormalTok{),}
            \AttributeTok{ll\_diff\_hi   =} \FunctionTok{quantile}\NormalTok{(value, }\FloatTok{0.95}\NormalTok{),}
            \AttributeTok{.groups =} \StringTok{"drop"}\NormalTok{)}

\NormalTok{mix\_check }\OtherTok{\textless{}{-}}\NormalTok{ post\_pi\_ind }\SpecialCharTok{|\textgreater{}}
  \FunctionTok{left\_join}\NormalTok{(ll\_diff\_summary, }\AttributeTok{by =} \StringTok{"agent\_id"}\NormalTok{)}

\FunctionTok{ggplot}\NormalTok{(mix\_check, }\FunctionTok{aes}\NormalTok{(}\AttributeTok{x =}\NormalTok{ ll\_diff\_mean, }\AttributeTok{y =}\NormalTok{ med, }\AttributeTok{color =}\NormalTok{ species)) }\SpecialCharTok{+}
  \FunctionTok{geom\_hline}\NormalTok{(}\AttributeTok{yintercept =} \FloatTok{0.5}\NormalTok{, }\AttributeTok{linetype =} \StringTok{"dashed"}\NormalTok{, }\AttributeTok{color =} \StringTok{"gray40"}\NormalTok{) }\SpecialCharTok{+}
  \FunctionTok{geom\_vline}\NormalTok{(}\AttributeTok{xintercept =} \DecValTok{0}\NormalTok{,   }\AttributeTok{linetype =} \StringTok{"dashed"}\NormalTok{, }\AttributeTok{color =} \StringTok{"gray40"}\NormalTok{) }\SpecialCharTok{+}
  \FunctionTok{geom\_errorbar}\NormalTok{(}\FunctionTok{aes}\NormalTok{(}\AttributeTok{ymin =}\NormalTok{ lo, }\AttributeTok{ymax =}\NormalTok{ hi), }\AttributeTok{width =} \DecValTok{0}\NormalTok{, }\AttributeTok{alpha =} \FloatTok{0.4}\NormalTok{) }\SpecialCharTok{+}
  \FunctionTok{geom\_errorbarh}\NormalTok{(}\FunctionTok{aes}\NormalTok{(}\AttributeTok{xmin =}\NormalTok{ ll\_diff\_lo, }\AttributeTok{xmax =}\NormalTok{ ll\_diff\_hi),}
                 \AttributeTok{height =} \DecValTok{0}\NormalTok{, }\AttributeTok{alpha =} \FloatTok{0.4}\NormalTok{) }\SpecialCharTok{+}
  \FunctionTok{geom\_point}\NormalTok{(}\AttributeTok{size =} \FloatTok{2.5}\NormalTok{) }\SpecialCharTok{+}
  \FunctionTok{labs}\NormalTok{(}\AttributeTok{x =} \FunctionTok{expression}\NormalTok{(}\StringTok{"Per{-}individual log{-}likelihood difference "}\SpecialCharTok{\textasciitilde{}}\NormalTok{(ell[}\DecValTok{1}\NormalTok{]}\SpecialCharTok{{-}}\NormalTok{ell[}\DecValTok{0}\NormalTok{])),}
       \AttributeTok{y =} \FunctionTok{expression}\NormalTok{(}\FunctionTok{Pr}\NormalTok{(z[j]}\SpecialCharTok{==}\DecValTok{1}\SpecialCharTok{\textasciitilde{}}\StringTok{"|"}\SpecialCharTok{\textasciitilde{}}\NormalTok{data)),}
       \AttributeTok{color =} \StringTok{"Species"}\NormalTok{,}
       \AttributeTok{title =} \StringTok{"Mixture posterior responds to individual{-}level evidence"}\NormalTok{,}
       \AttributeTok{subtitle =} \StringTok{"Rightward = data prefer 1{-}ToM; upward = posterior classifies as 1{-}ToM"}\NormalTok{)}
\end{Highlighting}
\end{Shaded}

\begin{figure}[H]

{\centering \includegraphics[width=0.85\linewidth,height=\textheight,keepaspectratio]{09-TheoryOfMind_files/figure-pdf/ch09_mix_occam_check-1.pdf}

}

\caption{Per-individual posterior architecture probability (y)
vs.~per-individual 1-ToM minus 0-ToM log-likelihood difference (x), both
extracted from the mixture fit. Posterior-mean points, colored by
species. A monotone rising curve confirms the mixture is actually
responsive to the structural evidence in each individual's choices; flat
clouds would indicate the posterior is driven entirely by the
species-level prior.}

\end{figure}%

\subsection{How well does the mixture predict
overall?}\label{how-well-does-the-mixture-predict-overall}

Mixture coefficients are only trustworthy to the extent that the mixture
model itself is predictively competitive. We compute PSIS-LOO for the
mixture and compare against the baseline hierarchical 0-ToM (the
non-mixture model fit earlier).

\begin{Shaded}
\begin{Highlighting}[]
\NormalTok{loo\_mix   }\OtherTok{\textless{}{-}}\NormalTok{ fit\_mix}\SpecialCharTok{$}\FunctionTok{loo}\NormalTok{()}
\NormalTok{loo\_base  }\OtherTok{\textless{}{-}}\NormalTok{ fit\_hier}\SpecialCharTok{$}\FunctionTok{loo}\NormalTok{()}
\NormalTok{loo\_cmp   }\OtherTok{\textless{}{-}}\NormalTok{ loo}\SpecialCharTok{::}\FunctionTok{loo\_compare}\NormalTok{(}\FunctionTok{list}\NormalTok{(}\AttributeTok{mixture =}\NormalTok{ loo\_mix, }\AttributeTok{baseline\_0tom =}\NormalTok{ loo\_base))}
\NormalTok{knitr}\SpecialCharTok{::}\FunctionTok{kable}\NormalTok{(}\FunctionTok{as.data.frame}\NormalTok{(loo\_cmp), }\AttributeTok{digits =} \DecValTok{2}\NormalTok{,}
             \AttributeTok{caption =} \StringTok{"PSIS{-}LOO comparison: mixture of architectures vs hierarchical 0{-}ToM only."}\NormalTok{)}
\end{Highlighting}
\end{Shaded}

\begin{longtable}[]{@{}
  >{\raggedright\arraybackslash}p{(\linewidth - 16\tabcolsep) * \real{0.1628}}
  >{\raggedleft\arraybackslash}p{(\linewidth - 16\tabcolsep) * \real{0.1163}}
  >{\raggedleft\arraybackslash}p{(\linewidth - 16\tabcolsep) * \real{0.0930}}
  >{\raggedleft\arraybackslash}p{(\linewidth - 16\tabcolsep) * \real{0.1047}}
  >{\raggedleft\arraybackslash}p{(\linewidth - 16\tabcolsep) * \real{0.1395}}
  >{\raggedleft\arraybackslash}p{(\linewidth - 16\tabcolsep) * \real{0.0698}}
  >{\raggedleft\arraybackslash}p{(\linewidth - 16\tabcolsep) * \real{0.1047}}
  >{\raggedleft\arraybackslash}p{(\linewidth - 16\tabcolsep) * \real{0.1047}}
  >{\raggedleft\arraybackslash}p{(\linewidth - 16\tabcolsep) * \real{0.1047}}@{}}
\caption{PSIS-LOO comparison: mixture of architectures vs hierarchical
0-ToM only.}\tabularnewline
\toprule\noalign{}
\begin{minipage}[b]{\linewidth}\raggedright
\end{minipage} & \begin{minipage}[b]{\linewidth}\raggedleft
elpd\_diff
\end{minipage} & \begin{minipage}[b]{\linewidth}\raggedleft
se\_diff
\end{minipage} & \begin{minipage}[b]{\linewidth}\raggedleft
elpd\_loo
\end{minipage} & \begin{minipage}[b]{\linewidth}\raggedleft
se\_elpd\_loo
\end{minipage} & \begin{minipage}[b]{\linewidth}\raggedleft
p\_loo
\end{minipage} & \begin{minipage}[b]{\linewidth}\raggedleft
se\_p\_loo
\end{minipage} & \begin{minipage}[b]{\linewidth}\raggedleft
looic
\end{minipage} & \begin{minipage}[b]{\linewidth}\raggedleft
se\_looic
\end{minipage} \\
\midrule\noalign{}
\endfirsthead
\toprule\noalign{}
\begin{minipage}[b]{\linewidth}\raggedright
\end{minipage} & \begin{minipage}[b]{\linewidth}\raggedleft
elpd\_diff
\end{minipage} & \begin{minipage}[b]{\linewidth}\raggedleft
se\_diff
\end{minipage} & \begin{minipage}[b]{\linewidth}\raggedleft
elpd\_loo
\end{minipage} & \begin{minipage}[b]{\linewidth}\raggedleft
se\_elpd\_loo
\end{minipage} & \begin{minipage}[b]{\linewidth}\raggedleft
p\_loo
\end{minipage} & \begin{minipage}[b]{\linewidth}\raggedleft
se\_p\_loo
\end{minipage} & \begin{minipage}[b]{\linewidth}\raggedleft
looic
\end{minipage} & \begin{minipage}[b]{\linewidth}\raggedleft
se\_looic
\end{minipage} \\
\midrule\noalign{}
\endhead
\bottomrule\noalign{}
\endlastfoot
baseline\_0tom & 0.00 & 0.00 & -5035.13 & 15.13 & 49.90 & 0.71 &
10070.25 & 30.26 \\
mixture & -9.15 & 1.76 & -5044.28 & 14.38 & 52.54 & 0.76 & 10088.56 &
28.77 \\
\end{longtable}

A positive elpd\_diff for \texttt{mixture} (with
\textbar elpd\_diff\textbar{} \textgreater{} 2·se\_diff) would indicate
the extra 1-ToM arm is earning its complexity on this dataset. A small
or negative elpd\_diff indicates that across the 39 primates, the data
jointly do not strongly prefer the mixture --- the species-level
\(\pi_s\) estimates would then carry interpretive weight only insofar as
they are shaped by the covariate regression, not by individual-level
likelihood contrast. This is the predictive-calibration side of the same
identifiability story Phase 6 laid out: matching pennies is a game with
limited leverage for distinguishing architectures.

\section{Secondary analysis: decisional determinism within a 0-ToM
architecture}\label{secondary-analysis-decisional-determinism-within-a-0-tom-architecture}

The main analysis is complete. As a subsidiary question we now ask
whether, \emph{conditional on treating every primate as a 0-ToM agent},
species-level decisional determinism (\texttt{log\_beta}) tracks ECV or
group size. This is the regression the previous version of this chapter
promoted to the main estimand; we now position it correctly, as a
question about statistical execution rather than about representational
capacity.

\begin{Shaded}
\begin{Highlighting}[]
\NormalTok{stan\_tom0\_reg }\OtherTok{\textless{}{-}} \StringTok{"}
\StringTok{data \{}
\StringTok{  int\textless{}lower=1\textgreater{} N;}
\StringTok{  int\textless{}lower=1\textgreater{} J;}
\StringTok{  int\textless{}lower=1\textgreater{} S;}
\StringTok{  array[J] int\textless{}lower=1, upper=S\textgreater{} species\_of;}
\StringTok{  array[N] int\textless{}lower=1, upper=J\textgreater{} agent;}
\StringTok{  array[N] int\textless{}lower=0, upper=1\textgreater{} choice;}
\StringTok{  array[N] int\textless{}lower=0, upper=1\textgreater{} op\_choice;}
\StringTok{  array[J] int\textless{}lower=1\textgreater{} start;}
\StringTok{  array[J] int\textless{}lower=1\textgreater{} stop;}
\StringTok{  vector[S] ecv\_z;        // standardized log ECV}
\StringTok{  vector[S] group\_z;      // standardized log group size}
\StringTok{\}}
\StringTok{parameters \{}
\StringTok{  vector[S] mu\_log\_sigma;}
\StringTok{  real\textless{}lower=0\textgreater{} tau\_log\_sigma;}
\StringTok{  vector[J] z\_log\_sigma;}
\StringTok{  // log\_beta regression}
\StringTok{  real alpha\_logbeta;}
\StringTok{  real b\_ecv;}
\StringTok{  real b\_group;}
\StringTok{  real\textless{}lower=0\textgreater{} tau\_sp;        // species residual SD on log\_beta}
\StringTok{  vector[S] z\_sp;}
\StringTok{  real\textless{}lower=0\textgreater{} tau\_log\_beta;  // individual{-}within{-}species SD}
\StringTok{  vector[J] z\_log\_beta;}
\StringTok{  // bias}
\StringTok{  vector[S] mu\_bias;}
\StringTok{  real\textless{}lower=0\textgreater{} tau\_bias;}
\StringTok{  vector[J] z\_bias;}
\StringTok{\}}
\StringTok{transformed parameters \{}
\StringTok{  vector[S] mu\_log\_beta = alpha\_logbeta + b\_ecv * ecv\_z +}
\StringTok{                          b\_group * group\_z + tau\_sp * z\_sp;}
\StringTok{  vector[J] log\_sigma = mu\_log\_sigma[species\_of] + tau\_log\_sigma * z\_log\_sigma;}
\StringTok{  vector[J] log\_beta  = mu\_log\_beta[species\_of]  + tau\_log\_beta  * z\_log\_beta;}
\StringTok{  vector[J] bias      = mu\_bias[species\_of]      + tau\_bias      * z\_bias;}
\StringTok{  vector[N] dV;}
\StringTok{  for (j in 1:J) \{}
\StringTok{    real sigma\_j = exp(log\_sigma[j]);}
\StringTok{    real mu\_s    = 0.0;}
\StringTok{    real Sig\_s   = 1.0;}
\StringTok{    for (k in start[j]:stop[j]) \{}
\StringTok{      real s\_mu = inv\_logit(mu\_s);}
\StringTok{      real p\_op = inv\_logit(mu\_s / sqrt(1 + 0.36 * Sig\_s));}
\StringTok{      dV[k] = 2 * p\_op {-} 1;}
\StringTok{      Sig\_s = 1.0 / (1.0 / (Sig\_s + sigma\_j) + s\_mu * (1 {-} s\_mu));}
\StringTok{      mu\_s  = mu\_s + Sig\_s * (op\_choice[k] {-} s\_mu);}
\StringTok{    \}}
\StringTok{  \}}
\StringTok{\}}
\StringTok{model \{}
\StringTok{  mu\_log\_sigma \textasciitilde{} normal({-}2, 1);}
\StringTok{  tau\_log\_sigma \textasciitilde{} exponential(2);}
\StringTok{  z\_log\_sigma \textasciitilde{} std\_normal();}
\StringTok{  alpha\_logbeta \textasciitilde{} normal({-}1, 1);}
\StringTok{  b\_ecv   \textasciitilde{} normal(0, 0.5);}
\StringTok{  b\_group \textasciitilde{} normal(0, 0.5);}
\StringTok{  tau\_sp        \textasciitilde{} exponential(2);}
\StringTok{  z\_sp          \textasciitilde{} std\_normal();}
\StringTok{  tau\_log\_beta  \textasciitilde{} exponential(1);}
\StringTok{  z\_log\_beta    \textasciitilde{} std\_normal();}
\StringTok{  mu\_bias \textasciitilde{} normal(0, 1);}
\StringTok{  tau\_bias \textasciitilde{} exponential(1);}
\StringTok{  z\_bias \textasciitilde{} std\_normal();}
\StringTok{  for (n in 1:N)}
\StringTok{    choice[n] \textasciitilde{} bernoulli\_logit(}
\StringTok{      (dV[n] + bias[agent[n]]) / exp(log\_beta[agent[n]])}
\StringTok{    );}
\StringTok{\}}
\StringTok{generated quantities \{}
\StringTok{  vector[N] log\_lik;}
\StringTok{  for (n in 1:N)}
\StringTok{    log\_lik[n] = bernoulli\_logit\_lpmf(}
\StringTok{      choice[n] | (dV[n] + bias[agent[n]]) / exp(log\_beta[agent[n]])}
\StringTok{    );}
\StringTok{\}}
\StringTok{"}
\FunctionTok{writeLines}\NormalTok{(stan\_tom0\_reg, here}\SpecialCharTok{::}\FunctionTok{here}\NormalTok{(}\StringTok{"stan"}\NormalTok{, }\StringTok{"ch09\_tom0\_reg.stan"}\NormalTok{))}
\NormalTok{mod\_reg }\OtherTok{\textless{}{-}} \FunctionTok{cmdstan\_model}\NormalTok{(here}\SpecialCharTok{::}\FunctionTok{here}\NormalTok{(}\StringTok{"stan"}\NormalTok{, }\StringTok{"ch09\_tom0\_reg.stan"}\NormalTok{))}
\end{Highlighting}
\end{Shaded}

\begin{Shaded}
\begin{Highlighting}[]
\CommentTok{\# Reuse species\_cov from the mixture section; build the data list for}
\CommentTok{\# the log\_beta regression}
\NormalTok{sd\_reg }\OtherTok{\textless{}{-}} \FunctionTok{c}\NormalTok{(}
\NormalTok{  sd\_int[}\FunctionTok{c}\NormalTok{(}\StringTok{"N"}\NormalTok{,}\StringTok{"J"}\NormalTok{,}\StringTok{"S"}\NormalTok{,}\StringTok{"species\_of"}\NormalTok{,}\StringTok{"agent"}\NormalTok{,}
           \StringTok{"choice"}\NormalTok{,}\StringTok{"op\_choice"}\NormalTok{,}\StringTok{"start"}\NormalTok{,}\StringTok{"stop"}\NormalTok{)],}
  \FunctionTok{list}\NormalTok{(}\AttributeTok{ecv\_z =}\NormalTok{ species\_cov}\SpecialCharTok{$}\NormalTok{ecv\_z,}
       \AttributeTok{group\_z =}\NormalTok{ species\_cov}\SpecialCharTok{$}\NormalTok{group\_z)}
\NormalTok{)}
\end{Highlighting}
\end{Shaded}

\begin{Shaded}
\begin{Highlighting}[]
\NormalTok{reg\_path }\OtherTok{\textless{}{-}}\NormalTok{ here}\SpecialCharTok{::}\FunctionTok{here}\NormalTok{(}\StringTok{"simmodels"}\NormalTok{, }\StringTok{"ch09\_tom0\_reg\_fit.rds"}\NormalTok{)}
\ControlFlowTok{if}\NormalTok{ (regenerate\_fits }\SpecialCharTok{||} \SpecialCharTok{!}\FunctionTok{file.exists}\NormalTok{(reg\_path)) \{}
\NormalTok{  fit\_reg }\OtherTok{\textless{}{-}}\NormalTok{ mod\_reg}\SpecialCharTok{$}\FunctionTok{sample}\NormalTok{(}
    \AttributeTok{data =}\NormalTok{ sd\_reg,}
    \AttributeTok{chains =} \DecValTok{4}\NormalTok{, }\AttributeTok{parallel\_chains =} \DecValTok{4}\NormalTok{,}
    \AttributeTok{iter\_warmup =} \DecValTok{1000}\NormalTok{, }\AttributeTok{iter\_sampling =} \DecValTok{1000}\NormalTok{,}
    \AttributeTok{adapt\_delta =} \FloatTok{0.95}\NormalTok{, }\AttributeTok{seed =} \DecValTok{2028}
\NormalTok{  )}
\NormalTok{  fit\_reg}\SpecialCharTok{$}\FunctionTok{save\_object}\NormalTok{(reg\_path)}
\NormalTok{\} }\ControlFlowTok{else}\NormalTok{ \{}
\NormalTok{  fit\_reg }\OtherTok{\textless{}{-}} \FunctionTok{readRDS}\NormalTok{(reg\_path)}
\NormalTok{\}}
\NormalTok{fit\_reg}\SpecialCharTok{$}\FunctionTok{diagnostic\_summary}\NormalTok{()}
\end{Highlighting}
\end{Shaded}

\begin{verbatim}
$num_divergent
[1] 1 2 1 2

$num_max_treedepth
[1]  0  0 31  0

$ebfmi
[1] 0.8511128 0.9028793 0.7713412 0.8232758
\end{verbatim}

\subsection{Results (decisional
determinism)}\label{results-decisional-determinism}

\begin{Shaded}
\begin{Highlighting}[]
\NormalTok{coef\_draws }\OtherTok{\textless{}{-}}\NormalTok{ fit\_reg}\SpecialCharTok{$}\FunctionTok{draws}\NormalTok{(}\FunctionTok{c}\NormalTok{(}\StringTok{"b\_ecv"}\NormalTok{,}\StringTok{"b\_group"}\NormalTok{,}\StringTok{"alpha\_logbeta"}\NormalTok{,}\StringTok{"tau\_sp"}\NormalTok{),}
                            \AttributeTok{format =} \StringTok{"df"}\NormalTok{) }\SpecialCharTok{|\textgreater{}}
  \FunctionTok{as\_tibble}\NormalTok{()}

\NormalTok{p\_coef }\OtherTok{\textless{}{-}}\NormalTok{ coef\_draws }\SpecialCharTok{|\textgreater{}}
\NormalTok{  dplyr}\SpecialCharTok{::}\FunctionTok{select}\NormalTok{(b\_ecv, b\_group) }\SpecialCharTok{|\textgreater{}}
  \FunctionTok{pivot\_longer}\NormalTok{(}\FunctionTok{everything}\NormalTok{(), }\AttributeTok{names\_to =} \StringTok{"coef"}\NormalTok{, }\AttributeTok{values\_to =} \StringTok{"value"}\NormalTok{) }\SpecialCharTok{|\textgreater{}}
  \FunctionTok{ggplot}\NormalTok{(}\FunctionTok{aes}\NormalTok{(}\AttributeTok{x =}\NormalTok{ value, }\AttributeTok{fill =}\NormalTok{ coef)) }\SpecialCharTok{+}
  \FunctionTok{geom\_density}\NormalTok{(}\AttributeTok{alpha =} \FloatTok{0.6}\NormalTok{, }\AttributeTok{color =} \ConstantTok{NA}\NormalTok{) }\SpecialCharTok{+}
  \FunctionTok{geom\_vline}\NormalTok{(}\AttributeTok{xintercept =} \DecValTok{0}\NormalTok{, }\AttributeTok{linetype =} \StringTok{"dashed"}\NormalTok{, }\AttributeTok{color =} \StringTok{"gray30"}\NormalTok{) }\SpecialCharTok{+}
  \FunctionTok{scale\_fill\_manual}\NormalTok{(}\AttributeTok{values =} \FunctionTok{c}\NormalTok{(}\StringTok{"b\_ecv"} \OtherTok{=} \StringTok{"steelblue"}\NormalTok{, }\StringTok{"b\_group"} \OtherTok{=} \StringTok{"firebrick3"}\NormalTok{),}
                    \AttributeTok{labels =} \FunctionTok{c}\NormalTok{(}\StringTok{"b\_ecv"} \OtherTok{=} \StringTok{"ECV (scaffolding)"}\NormalTok{,}
                               \StringTok{"b\_group"} \OtherTok{=} \StringTok{"Group size (Machiavellian)"}\NormalTok{)) }\SpecialCharTok{+}
  \FunctionTok{labs}\NormalTok{(}\AttributeTok{x =} \FunctionTok{expression}\NormalTok{(beta[cov]}\SpecialCharTok{\textasciitilde{}}\StringTok{"(log{-}beta per SD of log covariate)"}\NormalTok{),}
       \AttributeTok{y =} \StringTok{"Posterior density"}\NormalTok{,}
       \AttributeTok{fill =} \StringTok{"Hypothesis"}\NormalTok{,}
       \AttributeTok{title =} \StringTok{"Which covariate drives species{-}level decisiveness?"}\NormalTok{,}
       \AttributeTok{subtitle =} \StringTok{"Mass to the left of 0 = lower log\_beta (more deterministic) with larger covariate"}\NormalTok{)}

\CommentTok{\# Compare implied mu\_log\_beta from regression vs. baseline hier fit}
\NormalTok{baseline\_mu }\OtherTok{\textless{}{-}}\NormalTok{ fit\_hier}\SpecialCharTok{$}\FunctionTok{draws}\NormalTok{(}\StringTok{"mu\_log\_beta"}\NormalTok{, }\AttributeTok{format =} \StringTok{"df"}\NormalTok{) }\SpecialCharTok{|\textgreater{}}
  \FunctionTok{as\_tibble}\NormalTok{() }\SpecialCharTok{|\textgreater{}}
  \FunctionTok{summarize}\NormalTok{(}\FunctionTok{across}\NormalTok{(}\FunctionTok{starts\_with}\NormalTok{(}\StringTok{"mu\_log\_beta"}\NormalTok{), mean)) }\SpecialCharTok{|\textgreater{}}
  \FunctionTok{pivot\_longer}\NormalTok{(}\FunctionTok{everything}\NormalTok{(), }\AttributeTok{names\_to =} \StringTok{"param"}\NormalTok{, }\AttributeTok{values\_to =} \StringTok{"baseline\_mean"}\NormalTok{) }\SpecialCharTok{|\textgreater{}}
  \FunctionTok{mutate}\NormalTok{(}\AttributeTok{species\_id =} \FunctionTok{as.integer}\NormalTok{(}\FunctionTok{str\_extract}\NormalTok{(param, }\StringTok{"}\SpecialCharTok{\textbackslash{}\textbackslash{}}\StringTok{d+"}\NormalTok{))) }\SpecialCharTok{|\textgreater{}}
  \FunctionTok{left\_join}\NormalTok{(species\_cov }\SpecialCharTok{|\textgreater{}}\NormalTok{ dplyr}\SpecialCharTok{::}\FunctionTok{select}\NormalTok{(species\_id, species), }\AttributeTok{by =} \StringTok{"species\_id"}\NormalTok{)}

\NormalTok{reg\_mu }\OtherTok{\textless{}{-}}\NormalTok{ fit\_reg}\SpecialCharTok{$}\FunctionTok{draws}\NormalTok{(}\StringTok{"mu\_log\_beta"}\NormalTok{, }\AttributeTok{format =} \StringTok{"df"}\NormalTok{) }\SpecialCharTok{|\textgreater{}}
  \FunctionTok{as\_tibble}\NormalTok{() }\SpecialCharTok{|\textgreater{}}
  \FunctionTok{summarize}\NormalTok{(}\FunctionTok{across}\NormalTok{(}\FunctionTok{starts\_with}\NormalTok{(}\StringTok{"mu\_log\_beta"}\NormalTok{), mean)) }\SpecialCharTok{|\textgreater{}}
  \FunctionTok{pivot\_longer}\NormalTok{(}\FunctionTok{everything}\NormalTok{(), }\AttributeTok{names\_to =} \StringTok{"param"}\NormalTok{, }\AttributeTok{values\_to =} \StringTok{"reg\_mean"}\NormalTok{) }\SpecialCharTok{|\textgreater{}}
  \FunctionTok{mutate}\NormalTok{(}\AttributeTok{species\_id =} \FunctionTok{as.integer}\NormalTok{(}\FunctionTok{str\_extract}\NormalTok{(param, }\StringTok{"}\SpecialCharTok{\textbackslash{}\textbackslash{}}\StringTok{d+"}\NormalTok{)))}

\NormalTok{compare\_df }\OtherTok{\textless{}{-}}\NormalTok{ baseline\_mu }\SpecialCharTok{|\textgreater{}} \FunctionTok{left\_join}\NormalTok{(reg\_mu, }\AttributeTok{by =} \StringTok{"species\_id"}\NormalTok{)}

\NormalTok{p\_compare }\OtherTok{\textless{}{-}} \FunctionTok{ggplot}\NormalTok{(compare\_df, }\FunctionTok{aes}\NormalTok{(}\AttributeTok{x =}\NormalTok{ baseline\_mean, }\AttributeTok{y =}\NormalTok{ reg\_mean, }\AttributeTok{label =}\NormalTok{ species)) }\SpecialCharTok{+}
  \FunctionTok{geom\_abline}\NormalTok{(}\AttributeTok{slope =} \DecValTok{1}\NormalTok{, }\AttributeTok{intercept =} \DecValTok{0}\NormalTok{, }\AttributeTok{linetype =} \StringTok{"dashed"}\NormalTok{, }\AttributeTok{color =} \StringTok{"gray40"}\NormalTok{) }\SpecialCharTok{+}
  \FunctionTok{geom\_point}\NormalTok{(}\AttributeTok{size =} \DecValTok{3}\NormalTok{, }\AttributeTok{color =} \StringTok{"midnightblue"}\NormalTok{) }\SpecialCharTok{+}
\NormalTok{  ggrepel}\SpecialCharTok{::}\FunctionTok{geom\_text\_repel}\NormalTok{(}\AttributeTok{size =} \DecValTok{3}\NormalTok{) }\SpecialCharTok{+}
  \FunctionTok{labs}\NormalTok{(}\AttributeTok{x =} \StringTok{"Baseline species mean log\_beta"}\NormalTok{,}
       \AttributeTok{y =} \StringTok{"Regression{-}implied species mean log\_beta"}\NormalTok{,}
       \AttributeTok{title =} \StringTok{"Do ECV + group size explain species variation?"}\NormalTok{,}
       \AttributeTok{subtitle =} \StringTok{"Close to diagonal = covariates capture the between{-}species signal"}\NormalTok{)}

\NormalTok{p\_coef }\SpecialCharTok{+}\NormalTok{ p\_compare }\SpecialCharTok{+} \FunctionTok{plot\_layout}\NormalTok{(}\AttributeTok{widths =} \FunctionTok{c}\NormalTok{(}\DecValTok{1}\NormalTok{, }\DecValTok{1}\NormalTok{))}
\end{Highlighting}
\end{Shaded}

\begin{figure}[H]

{\centering \includegraphics[width=0.85\linewidth,height=\textheight,keepaspectratio]{09-TheoryOfMind_files/figure-pdf/ch09_coef_plot-1.pdf}

}

\caption{Secondary analysis: posterior coefficients of a regression of
the 0-ToM decision temperature log\_beta on species-level standardized
log ECV and log group size. This is NOT the main estimand (see the
mixture regression above) --- it asks the narrower question of whether
species act more deterministically on their 0-ToM predictions with
larger brains or larger social groups. Negative coefficient = lower
log\_beta (tighter softmax) with higher covariate.}

\end{figure}%

The coefficient posteriors are the scientific payoff. Reading them
requires keeping the direction straight: \texttt{log\_beta} is an
\emph{inverse} measure of decisiveness --- smaller values mean tighter
softmax, more deterministic use of the expected-utility signal. A
\emph{negative} regression coefficient therefore means that species with
higher ECV (or larger groups) act more deterministically on their 0-ToM
predictions.

Summary statistics:

\begin{Shaded}
\begin{Highlighting}[]
\NormalTok{coef\_summary }\OtherTok{\textless{}{-}}\NormalTok{ coef\_draws }\SpecialCharTok{|\textgreater{}}
\NormalTok{  dplyr}\SpecialCharTok{::}\FunctionTok{select}\NormalTok{(b\_ecv, b\_group) }\SpecialCharTok{|\textgreater{}}
  \FunctionTok{pivot\_longer}\NormalTok{(}\FunctionTok{everything}\NormalTok{(), }\AttributeTok{names\_to =} \StringTok{"coef"}\NormalTok{, }\AttributeTok{values\_to =} \StringTok{"value"}\NormalTok{) }\SpecialCharTok{|\textgreater{}}
  \FunctionTok{group\_by}\NormalTok{(coef) }\SpecialCharTok{|\textgreater{}}
  \FunctionTok{summarize}\NormalTok{(}
    \AttributeTok{mean   =} \FunctionTok{mean}\NormalTok{(value),}
    \AttributeTok{median =} \FunctionTok{median}\NormalTok{(value),}
    \AttributeTok{q05    =} \FunctionTok{quantile}\NormalTok{(value, }\FloatTok{0.05}\NormalTok{),}
    \AttributeTok{q95    =} \FunctionTok{quantile}\NormalTok{(value, }\FloatTok{0.95}\NormalTok{),}
    \AttributeTok{p\_neg  =} \FunctionTok{mean}\NormalTok{(value }\SpecialCharTok{\textless{}} \DecValTok{0}\NormalTok{),}
    \AttributeTok{p\_pos  =} \FunctionTok{mean}\NormalTok{(value }\SpecialCharTok{\textgreater{}} \DecValTok{0}\NormalTok{),}
    \AttributeTok{.groups =} \StringTok{"drop"}
\NormalTok{  )}
\NormalTok{knitr}\SpecialCharTok{::}\FunctionTok{kable}\NormalTok{(coef\_summary, }\AttributeTok{digits =} \DecValTok{3}\NormalTok{)}
\end{Highlighting}
\end{Shaded}

\begin{longtable}[]{@{}lrrrrrr@{}}
\toprule\noalign{}
coef & mean & median & q05 & q95 & p\_neg & p\_pos \\
\midrule\noalign{}
\endhead
\bottomrule\noalign{}
\endlastfoot
b\_ecv & -0.423 & -0.417 & -0.835 & -0.028 & 0.957 & 0.043 \\
b\_group & 0.260 & 0.266 & -0.090 & 0.582 & 0.104 & 0.896 \\
\end{longtable}

Interpretation rests on whether each coefficient's 90\% credible
interval excludes zero, on the sign, and on which of the two carries
more credible mass. Note well: with only seven species and two
correlated-but-not-collinear covariates, statistical power at the
species level is modest --- this is an inference about an evolutionary
landscape sampled at seven points, and our uncertainty should reflect
that. An honest report sounds like ``the posterior of \(b_{\text{ECV}}\)
places X\% of mass below zero, while \(b_{\text{group}}\) places Y\%
below zero'', not ``ECV is the winner''.

The second plot checks that the regression is sensible at all: if the
regression-implied species means fall close to the diagonal against the
non-regression baseline, then the covariates are capturing most of the
between-species variance. A scatter of points far off the diagonal would
indicate that ECV and group size jointly miss an important species-level
dimension (perhaps neocortex ratio, diet, or rearing --- all present in
the raw CSV).

\section{Cross-check: species-level stacking
weights}\label{cross-check-species-level-stacking-weights}

The mixture model delivered species-level architecture probabilities
\(\pi_s\) as \emph{parameters}, with full posterior propagation. As a
methodological cross-check, we can compare those against species-level
PSIS stacking weights computed from two separately-fit hierarchical
models (a pure-0-ToM hier and a pure-1-ToM hier). The two estimators are
answering related but formally distinct questions --- \(\pi_s\) is a
prior mixing proportion inside a joint generative model, while stacking
weights are frequentist-flavored pseudo-BMA weights for out-of-sample
predictive averaging --- so exact agreement is not expected, but gross
divergence would be concerning.

\begin{Shaded}
\begin{Highlighting}[]
\NormalTok{loo0 }\OtherTok{\textless{}{-}}\NormalTok{ fit\_hier}\SpecialCharTok{$}\FunctionTok{loo}\NormalTok{()}
\NormalTok{loo1 }\OtherTok{\textless{}{-}}\NormalTok{ fit\_hier\_tom1}\SpecialCharTok{$}\FunctionTok{loo}\NormalTok{()}
\NormalTok{species\_of\_trial }\OtherTok{\textless{}{-}}\NormalTok{ sd\_int}\SpecialCharTok{$}\NormalTok{species\_lookup}\SpecialCharTok{$}\NormalTok{species[}
\NormalTok{  sd\_int}\SpecialCharTok{$}\NormalTok{species\_of[sd\_int}\SpecialCharTok{$}\NormalTok{agent]}
\NormalTok{]}
\NormalTok{species\_stack }\OtherTok{\textless{}{-}} \FunctionTok{split}\NormalTok{(}\FunctionTok{seq\_along}\NormalTok{(species\_of\_trial), species\_of\_trial) }\SpecialCharTok{|\textgreater{}}
  \FunctionTok{imap\_dfr}\NormalTok{(}\ControlFlowTok{function}\NormalTok{(idx, sp) \{}
\NormalTok{    lpd }\OtherTok{\textless{}{-}} \FunctionTok{cbind}\NormalTok{(loo0}\SpecialCharTok{$}\NormalTok{pointwise[idx, }\StringTok{"elpd\_loo"}\NormalTok{],}
\NormalTok{                 loo1}\SpecialCharTok{$}\NormalTok{pointwise[idx, }\StringTok{"elpd\_loo"}\NormalTok{])}
\NormalTok{    w   }\OtherTok{\textless{}{-}}\NormalTok{ loo}\SpecialCharTok{::}\FunctionTok{stacking\_weights}\NormalTok{(lpd)}
    \FunctionTok{tibble}\NormalTok{(}\AttributeTok{species =}\NormalTok{ sp, }\AttributeTok{w\_tom0 =}\NormalTok{ w[}\DecValTok{1}\NormalTok{], }\AttributeTok{w\_tom1 =}\NormalTok{ w[}\DecValTok{2}\NormalTok{])}
\NormalTok{  \})}

\NormalTok{stack\_compare }\OtherTok{\textless{}{-}}\NormalTok{ species\_stack }\SpecialCharTok{|\textgreater{}}
  \FunctionTok{left\_join}\NormalTok{(pi\_s\_summary }\SpecialCharTok{|\textgreater{}}\NormalTok{ dplyr}\SpecialCharTok{::}\FunctionTok{select}\NormalTok{(species, med, lo, hi),}
            \AttributeTok{by =} \StringTok{"species"}\NormalTok{)}

\FunctionTok{ggplot}\NormalTok{(stack\_compare, }\FunctionTok{aes}\NormalTok{(}\AttributeTok{x =}\NormalTok{ w\_tom1, }\AttributeTok{y =}\NormalTok{ med, }\AttributeTok{label =}\NormalTok{ species)) }\SpecialCharTok{+}
  \FunctionTok{geom\_abline}\NormalTok{(}\AttributeTok{slope =} \DecValTok{1}\NormalTok{, }\AttributeTok{intercept =} \DecValTok{0}\NormalTok{, }\AttributeTok{linetype =} \StringTok{"dashed"}\NormalTok{, }\AttributeTok{color =} \StringTok{"gray40"}\NormalTok{) }\SpecialCharTok{+}
  \FunctionTok{geom\_errorbar}\NormalTok{(}\FunctionTok{aes}\NormalTok{(}\AttributeTok{ymin =}\NormalTok{ lo, }\AttributeTok{ymax =}\NormalTok{ hi), }\AttributeTok{width =} \DecValTok{0}\NormalTok{, }\AttributeTok{alpha =} \FloatTok{0.5}\NormalTok{,}
                \AttributeTok{color =} \StringTok{"steelblue"}\NormalTok{) }\SpecialCharTok{+}
  \FunctionTok{geom\_point}\NormalTok{(}\AttributeTok{size =} \DecValTok{3}\NormalTok{, }\AttributeTok{color =} \StringTok{"midnightblue"}\NormalTok{) }\SpecialCharTok{+}
\NormalTok{  ggrepel}\SpecialCharTok{::}\FunctionTok{geom\_text\_repel}\NormalTok{(}\AttributeTok{size =} \DecValTok{3}\NormalTok{) }\SpecialCharTok{+}
  \FunctionTok{coord\_cartesian}\NormalTok{(}\AttributeTok{xlim =} \FunctionTok{c}\NormalTok{(}\DecValTok{0}\NormalTok{, }\DecValTok{1}\NormalTok{), }\AttributeTok{ylim =} \FunctionTok{c}\NormalTok{(}\DecValTok{0}\NormalTok{, }\DecValTok{1}\NormalTok{)) }\SpecialCharTok{+}
  \FunctionTok{labs}\NormalTok{(}\AttributeTok{x =} \StringTok{"Stacking weight for 1{-}ToM (LOO{-}based)"}\NormalTok{,}
       \AttributeTok{y =} \FunctionTok{expression}\NormalTok{(pi[s] }\SpecialCharTok{\textasciitilde{}} \StringTok{"(mixture posterior median)"}\NormalTok{),}
       \AttributeTok{title =} \StringTok{"Stacking weights vs mixture \textbackslash{}u03c0\_s: cross{-}check"}\NormalTok{,}
       \AttributeTok{subtitle =} \StringTok{"Diagonal = perfect agreement; deviations indicate regularization vs data{-}driven differences"}\NormalTok{)}
\end{Highlighting}
\end{Shaded}

\begin{figure}[H]

{\centering \includegraphics[width=0.85\linewidth,height=\textheight,keepaspectratio]{09-TheoryOfMind_files/figure-pdf/ch09_stacking_crosscheck-1.pdf}

}

\caption{Cross-check: species-level stacking weight for 1-ToM (x-axis,
LOO-based pseudo-BMA) vs.~mixture model posterior median π\_s (y-axis).
The two estimators answer related but formally distinct questions; gross
divergence from the diagonal would indicate the mixture's π\_s is driven
by prior regularization rather than data. Error bars are 90\% CIs on
π\_s from the mixture.}

\end{figure}%

The mixture model is the principled version --- it propagates
uncertainty, shares information appropriately across species via the
regression, and puts the scientific estimand on the right scale. The
stacking cross-check exists to verify that the mixture's \(\pi_s\) is
not an artifact of prior-pulled regularization.

\section{Limitations}\label{limitations}

A candid list, because this chapter has pushed the framework as far as
it can honestly go on this dataset:

\begin{itemize}
\tightlist
\item
  \textbf{Fixed-parameter 1-ToM.} Our 1-ToM implementation tracks the
  first two moments of the agent's belief about the opponent's belief
  about itself, conditional on point values
  \((\sigma_{\text{op}}, \beta_{\text{op}})\). The agent does not
  represent uncertainty over those opponent parameters. The
  full-variational scheme of Devaine et al.~(2014a) does, and would
  require a specialized message-passing backend to sample efficiently.
  What we have estimated is therefore a 1-ToM agent whose internal
  opponent model is a \emph{fixed} 0-ToM with known parameters, not a
  Bayesian learner over opponent sophistication.
\item
  \textbf{Structural unidentifiability in matching pennies.} Phase 4b
  SBC confirmed what the Devaine et al.~(2014a) variational analysis
  anticipates: \((\log\sigma_{\text{op}}, \log\beta_{\text{op}})\) are
  weakly identified by matching-pennies choice sequences. Once
  \(\Sigma^{\text{self}}\) reaches its fixed point, the predicted
  opponent-choice trajectory depends on these parameters only through a
  compressed near-one-dimensional ridge. No amount of Bayesian machinery
  can manufacture likelihood contrast that the design does not provide.
  A reversal paradigm or a richer payoff matrix is the appropriate
  remedy.
\item
  \textbf{Mixture identifiability is design-bounded.} The per-individual
  posterior \(\Pr(z_j = 1)\) is informed only to the extent that 0-ToM
  and 1-ToM make differentiable predictions for that individual --- a
  leverage quantity we measured directly in the Occam's-razor panel.
  Individuals whose choice sequences are predictively equivalent under
  the two architectures have posteriors that return to \(\pi_{s(j)}\).
  This is correct Bayesian behavior, but it means our species-level
  \(\pi_s\) posteriors inherit the design's limited per-individual
  separability.
\item
  \textbf{Single-step reactive models.} 0-ToM, 1-ToM, and influence
  learning all plan only one step ahead. They cannot strategically
  deceive, willingly cooperate for long-run gain, or explore at the cost
  of immediate reward. Repeated Prisoner's Dilemma exposes every one of
  them.
\item
  \textbf{2×2 games only.} The framework generalizes to arbitrary payoff
  matrices, but both the primate study and this chapter live inside a
  single-choice, two-option game. Naturalistic social interactions are
  orders of magnitude richer, and richer settings are exactly where
  higher-\(k\) ToM should make a larger difference (Devaine et al.,
  2017, Discussion).
\item
  \textbf{Seven species is a small N.} The covariate regressions are
  honest but statistically modest --- an inference about an evolutionary
  landscape sampled at seven points. The phylogenetic structure of
  primates also means the species residuals around the regression line
  are not independent in the usual sense (Devaine et al., 2017 comment
  on this); a fuller analysis would embed a phylogenetic
  variance-covariance on \(u_s\).
\item
  \textbf{Cooperative framing.} Primates partially read the game as
  cooperative because their opponent is also their feeder. All of our
  models have been written for the competitive payoff matrix; a switch
  to cooperative incentives would flip signs in the influence and 1-ToM
  equations and should be fit as a separate model (not treated as
  noise).
\item
  \textbf{No true \(k \ge 2\).} The model set stops at \(k = 1\).
  Devaine et al.~(2017) report an evolutionary gap between great apes
  and humans precisely at \(k = 2\), which our implementation cannot
  currently reach. Extending the recursive belief-update machinery to
  maintain a posterior over opponent sophistication level --- and adding
  a third mixture component --- is the natural capstone project.
\item
  \textbf{A better predictive score is not a better mechanistic
  explanation.} Ch. 7's caveat applies with extra force here. Posterior
  architecture probability \(\Pr(z_j = 1)\) summarizes predictive
  structure in the choice sequences, not direct evidence of mentalizing
  in the neurocognitive sense.
\end{itemize}

\section{Conclusion}\label{conclusion-1}

We have traveled from \emph{noticing} that a competitive game with an
adaptive opponent breaks the passive-environment assumption of earlier
chapters, through three computational instantiations of that observation
(0-ToM, influence learning, fixed-parameter 1-ToM), through a complete
validation battery (prior predictive, parameter recovery, SBC for both
models, prior sensitivity, posterior predictive, and an Occam's-razor
model comparison study), to a joint hierarchical mixture model that
places representational-capacity probabilities on the left-hand side of
the covariate regression and answers the actual scientific question
directly. Nothing in this pipeline is new machinery --- every tool came
from Ch. 5, Ch. 6, or Ch. 7. What changed is what the model is
\emph{about}: the object of inference is no longer a statistical
regularity but a mind that is modeling a mind, and the estimand is not
how sharply that mind acts but whether it is mentalizing at all. That
shift is what makes the chapter worth writing, and it is the same shift
that makes \(k\)-ToM one of the most promising formal definitions on
offer for a concept --- \emph{theory of mind} --- that has for decades
existed as a verbal construct in search of a mechanism.

\section{References}\label{references-1}

Daunizeau, J. (2014). The Variational Bayesian Approach (VBA) toolbox.
\emph{Source code and documentation},
https://mbb-team.github.io/VBA-toolbox/.

Devaine, M., Hollard, G., \& Daunizeau, J. (2014a). The social Bayesian
brain: does mentalizing make a difference when we learn? \emph{PLOS
Computational Biology}, 10(12), e1003992.

Devaine, M., Hollard, G., \& Daunizeau, J. (2014b). Theory of mind: did
evolution fool us? \emph{PLOS ONE}, 9(2), e87619.

Devaine, M., San-Galli, A., Trapanese, C., Bardino, G., Hano, C., Saint
Jalme, M., Bouret, S., Masi, S., \& Daunizeau, J. (2017). Reading wild
minds: A computational assay of theory of mind sophistication across
seven primate species. \emph{PLOS Computational Biology}, 13(11),
e1005833.

Hampton, A. N., Bossaerts, P., \& O'Doherty, J. P. (2008). Neural
correlates of mentalizing-related computations during strategic
interactions in humans. \emph{PNAS}, 105(18), 6741--6746.

Quesque, F., \& Rossetti, Y. (2020). What do theory-of-mind tasks
actually measure? \emph{Perspectives on Psychological Science}, 15(2),
384--396.

Schaafsma, S. M., Pfaff, D. W., Spunt, R. P., \& Adolphs, R. (2015).
Deconstructing and reconstructing theory of mind. \emph{Trends in
Cognitive Sciences}, 19(2), 65--72.

Talts, S., Betancourt, M., Simpson, D., Vehtari, A., \& Gelman, A.
(2018). Validating Bayesian inference algorithms with Simulation-Based
Calibration. \emph{arXiv:1804.06788}.

Waade, P. T., Enevoldsen, K. C., Vermillet, A.-Q., Simonsen, A., \&
Fusaroli, R. (2022). Introducing tomsup: Theory of mind simulations
using Python. \emph{Behavior Research Methods}, 55, 2197--2231.

Yao, Y., Vehtari, A., Simpson, D., \& Gelman, A. (2018). Using stacking
to average Bayesian predictive distributions. \emph{Bayesian Analysis},
13(3), 917--1007.

\bookmarksetup{startatroot}

\chapter{Bayesian Models of Cognition}\label{sec-bayesian-models}

\begin{quote}
\textbf{📍 Where we are in the Bayesian modeling workflow:} Chs. 1--8
used the Bayesian \emph{statistical} framework as a tool to model
cognitive processes in the matching-pennies task. This chapter
introduces a distinct scientific claim: the \emph{Bayesian cognition
hypothesis} --- the proposal that the mind itself operates according to
Bayes' theorem. We use our Bayesian statistical tools to test whether
this hypothesis fits human behaviour in a social learning task, applying
the full validation battery from Chs. 5--7. The chapter also introduces
Leave-Future-Out cross-validation for sequential models, a technique
flagged in Ch. 8 but not developed there.
\end{quote}

\section{Bayesian statistics vs.~Bayesian
Cognition}\label{bayesian-statistics-vs.-bayesian-cognition}

This entire course has used \emph{Bayesian statistics} as its
inferential framework: we represent uncertainty as probability
distributions, write likelihoods, specify priors, and obtain posteriors
via Stan. That is our \emph{tool} to estimate any kind of model.

This chapter introduces a distinct scientific claim: the \emph{Bayesian
cognition hypothesis}. This is the proposal that the \emph{mind itself}
operates according to principles analogous to Bayes' theorem --- that
people represent beliefs as probability distributions, update them
optimally when new evidence arrives, and combine information sources in
proportion to their reliability. Whether this hypothesis is empirically
correct is an open and contested question. We will use our Bayesian
statistical tools to derive from the theory models to analyse some
experimental data.

N.B. Bayes as a tool and Bayes as a model are very different ways of
using Bayes. Indeed, we have fitted non-Bayesian models using Bayesian
techniques, and there are many papers out there testing Bayesian models
of cognition using frequentist techniques.

\begin{tcolorbox}[enhanced jigsaw, toprule=.15mm, toptitle=1mm, bottomtitle=1mm, leftrule=.75mm, rightrule=.15mm, colframe=quarto-callout-note-color-frame, left=2mm, colback=white, titlerule=0mm, coltitle=black, colbacktitle=quarto-callout-note-color!10!white, opacityback=0, title=\textcolor{quarto-callout-note-color}{\faInfo}\hspace{0.5em}{Historical Context: The Bayesian Brain --- From Helmholtz to Predictive
Coding}, arc=.35mm, opacitybacktitle=0.6, breakable, bottomrule=.15mm]

The idea that the mind performs something like Bayesian inference is not
a recent invention of computational cognitive science. It is a
rediscovery --- and mathematical formalisation --- of a proposal that
Hermann von Helmholtz articulated in 1867, before probability theory had
the tools to make it precise.

\subsubsection{Perception as Unconscious
Inference}\label{perception-as-unconscious-inference}

Helmholtz observed that the retinal image dramatically underdetermines
the state of the world: the same pattern of light could be produced by
infinitely many physical configurations. His conclusion was that
perception must involve \emph{unconscious inference} --- the visual
system implicitly reasons from ambiguous sensory data to the most
probable cause, drawing on prior knowledge about how the world is
structured (Helmholtz 1867). This is, in substance, a description of
Bayesian posterior inference, but articulated without probability theory
and without a mechanism.

For nearly a century the proposal remained an influential metaphor
rather than a testable theory. The mathematical tools existed --- Bayes'
theorem dates to 1763 --- but applying them to cognition required a
formalism for representing prior knowledge and a computational
infrastructure for performing inference, neither of which was available
to nineteenth-century psychologists.

\subsubsection{The Ideal Observer
Tradition}\label{the-ideal-observer-tradition}

The bridge between Helmholtz's intuition and formal modeling was built
by the \textbf{ideal observer} programme in perception research. An
ideal observer is defined as a hypothetical agent that performs
statistically optimal inference given the information available --- in
other words, a Bayesian. Geisler (1989) showed that ideal observers
could predict human perceptual thresholds and detection curves with
surprising accuracy. Ernst and Banks (2002) demonstrated Bayesian
integration of visual and haptic size estimates in humans: when the two
cues conflict, humans weight them in inverse proportion to their
variance, exactly as a Bayesian would. These results were influential
because they quantified the hypothesis rather than merely asserting it,
and because they made precise, falsifiable predictions about how
integration should depend on noise levels.

\subsubsection{``How to Grow a Mind'': The Bayesian Cognition
Programme}\label{how-to-grow-a-mind-the-bayesian-cognition-programme}

The broader claim --- that Bayesian inference characterises not just
low-level perception but high-level cognition --- was synthesised and
publicised by Tenenbaum et al. (2011). The programme demonstrated that
Bayesian models could account for rapid concept learning from sparse
data (the ``size principle'': smaller hypothesis spaces are preferred a
priori), intuitive theories of physics and psychology, causal reasoning,
and language acquisition. The key methodological move was to define a
prior over a structured hypothesis space --- a space of programs, causal
graphs, or grammars --- and to use Bayes' theorem to perform inference
over that space given data. The framework explained not just what people
infer but \emph{how fast} they do so from how little evidence, which had
been the central embarrassment for associationist and purely data-driven
accounts.

\subsubsection{Predictive Coding and the Free Energy
Principle}\label{predictive-coding-and-the-free-energy-principle}

A parallel strand of development, more neurally oriented, produced
predictive coding (Rao and Ballard 1999) and the free energy principle
(Friston 2005). In these frameworks, the brain is organised
hierarchically: each level sends predictions \emph{downward} to the
level below and receives \emph{prediction errors} upward. Perception is
the process of minimising prediction error (or equivalently, of finding
the posterior over hidden states that best explains sensory input).
Action is the process of changing the world to match predictions. The
framework is formally equivalent to variational Bayesian inference with
an approximate posterior, and it connects Helmholtz's unconscious
inference to measurable neural signals (the mismatch negativity,
attentional modulation, and dopaminergic prediction errors).

\subsubsection{The Falsifiability Critique and the Methodological
Answer}\label{the-falsifiability-critique-and-the-methodological-answer}

The Bayesian cognition programme attracted a pointed critique: the
framework is too flexible to be falsifiable. With sufficiently
unconstrained priors and likelihood functions, almost any behavioral
pattern can be ``explained'' post hoc (Jones and Love 2011; Bowers and
Davis 2012). The critique is well-founded when Bayesian models are used
purely as descriptive devices, fitting parameters to data without asking
whether the generative story is correct.

The response --- embodied in the workflow this book has built since
Chapter 5 --- is methodological rather than theoretical. Prior
predictive checks ask whether the model's prior beliefs produce
plausible data; simulation-based calibration checks whether parameter
recovery is reliable; LOO cross-validation assesses whether the model
predicts held-out observations; model comparison with power analysis
asks whether the experimental design has the discriminative resolution
to distinguish Bayesian from non-Bayesian accounts. None of these tests
is specific to the Bayesian cognition hypothesis, but all are necessary
to make the hypothesis empirically tractable. The difference between a
Bayesian model used as a just-so story and one used as a genuine
scientific claim lies entirely in the rigour of the validation pipeline.

\end{tcolorbox}

\section{Introduction}\label{introduction-3}

The human mind constantly receives input from multiple sources: direct
sensory evidence, social information from others, and prior knowledge
built from experience (perhaps even from innate structures). A
fundamental question in cognitive science is how these disparate pieces
of information are combined to produce coherent beliefs about the world.

The Bayesian framework offers a powerful approach to formalizing this
process. Under this framework, the mind is conceptualized as a
probabilistic machine that continuously updates its beliefs based on new
evidence. This contrasts with rule-based or purely associative models of
cognition by emphasizing:

\begin{itemize}
\item
  \textbf{Representations of uncertainty:} Beliefs are represented as
  probability distributions, not single values.
\item
  \textbf{Optimal integration:} Information is combined according to its
  reliability.
\item
  \textbf{Prior knowledge:} New evidence is interpreted in light of
  existing beliefs.
\end{itemize}

Formally, Bayes' theorem states:

\[P(\text{belief} \mid \text{evidence}) \propto P(\text{evidence} \mid \text{belief}) \times P(\text{belief})\]

where \(P(\text{belief} \mid \text{evidence})\) is the posterior,
\(P(\text{evidence} \mid \text{belief})\) is the likelihood, and
\(P(\text{belief})\) is the prior. The \(\propto\) symbol denotes
proportionality --- the product on the right should in principle be
normalised by dividing by \(P(\text{evidence})\) to recover a proper
probability. We often ignore this normalization constant, sometimes
because it cancels out, sometimes because our sampling methods can
approximate it away, sometimes because it doesn't rely on belief.

\section{Learning Objectives}\label{learning-objectives-9}

After completing this chapter, you will be able to:

\begin{itemize}
\item
  Understand the basic principles of Bayesian information integration
  and distinguish the \emph{Bayesian cognition} hypothesis from the
  \emph{Bayesian statistics} toolkit.
\item
  Implement cognitive models --- starting with a Simple Bayesian Agent
  (SBA) and a Weighted Bayesian Agent (WBA), discover their limitations,
  and then develop reparameterised solutions including a Proportional
  Bayesian Agent (PBA).
\item
  Recognize that the \emph{observation model} (how beliefs become
  actions) is a separate modelling dimension from the
  \emph{belief-update model} (how evidence is integrated).
\item
  Execute the full Bayesian workflow (prior predictive checks → fitting
  → diagnostics → parameter recovery → SBC → model comparison) on
  cognitive models.
\item
  Identify when models have identifiability issues and how
  reparameterisation can solve them.
\end{itemize}

\section{The Bayesian Framework for
Cognition}\label{the-bayesian-framework-for-cognition}

Bayesian models of cognition explore the idea that the mind operates
according to principles similar to Bayes' theorem, combining different
sources of evidence to form updated beliefs. Most commonly this is
framed in terms of prior beliefs being updated with new evidence to form
a posterior.

\subsection{Visualising Bayesian
Updating}\label{visualising-bayesian-updating}

\begin{Shaded}
\begin{Highlighting}[]
\NormalTok{x          }\OtherTok{\textless{}{-}} \FunctionTok{seq}\NormalTok{(}\DecValTok{0}\NormalTok{, }\DecValTok{1}\NormalTok{, }\AttributeTok{by =} \FloatTok{0.01}\NormalTok{)}
\NormalTok{prior      }\OtherTok{\textless{}{-}} \FunctionTok{dbeta}\NormalTok{(x, }\DecValTok{3}\NormalTok{, }\DecValTok{3}\NormalTok{)}
\NormalTok{likelihood }\OtherTok{\textless{}{-}} \FunctionTok{dbeta}\NormalTok{(x, }\DecValTok{7}\NormalTok{, }\DecValTok{3}\NormalTok{)}
\NormalTok{posterior  }\OtherTok{\textless{}{-}} \FunctionTok{dbeta}\NormalTok{(x, }\DecValTok{10}\NormalTok{, }\DecValTok{6}\NormalTok{)}

\CommentTok{\# 1. Define the length explicitly outside the tibble}
\NormalTok{n\_x }\OtherTok{\textless{}{-}} \FunctionTok{length}\NormalTok{(x) }

\NormalTok{plot\_data }\OtherTok{\textless{}{-}} \FunctionTok{tibble}\NormalTok{(}
  \AttributeTok{x            =} \FunctionTok{rep}\NormalTok{(x, }\DecValTok{3}\NormalTok{),}
  \AttributeTok{density      =} \FunctionTok{c}\NormalTok{(prior, likelihood, posterior),}
  \AttributeTok{distribution =} \FunctionTok{factor}\NormalTok{(}
    \CommentTok{\# 2. Use the explicit length variable here}
    \FunctionTok{rep}\NormalTok{(}\FunctionTok{c}\NormalTok{(}\StringTok{"Prior"}\NormalTok{, }\StringTok{"Likelihood"}\NormalTok{, }\StringTok{"Posterior"}\NormalTok{), }\AttributeTok{each =}\NormalTok{ n\_x), }
    \AttributeTok{levels =} \FunctionTok{c}\NormalTok{(}\StringTok{"Prior"}\NormalTok{, }\StringTok{"Likelihood"}\NormalTok{, }\StringTok{"Posterior"}\NormalTok{)}
\NormalTok{  )}
\NormalTok{)}

\FunctionTok{ggplot}\NormalTok{(plot\_data, }\FunctionTok{aes}\NormalTok{(}\AttributeTok{x =}\NormalTok{ x, }\AttributeTok{y =}\NormalTok{ density,}
                      \AttributeTok{colour =}\NormalTok{ distribution, }\AttributeTok{linetype =}\NormalTok{ distribution)) }\SpecialCharTok{+}
  \FunctionTok{geom\_line}\NormalTok{(}\AttributeTok{linewidth =} \FloatTok{1.2}\NormalTok{) }\SpecialCharTok{+}
  \FunctionTok{scale\_colour\_manual}\NormalTok{(}
    \AttributeTok{values =} \FunctionTok{c}\NormalTok{(}\StringTok{"Prior"} \OtherTok{=} \StringTok{"\#4575b4"}\NormalTok{, }\StringTok{"Likelihood"} \OtherTok{=} \StringTok{"\#d73027"}\NormalTok{,}
               \StringTok{"Posterior"} \OtherTok{=} \StringTok{"\#756bb1"}\NormalTok{)}
\NormalTok{  ) }\SpecialCharTok{+}
  \FunctionTok{scale\_linetype\_manual}\NormalTok{(}
    \AttributeTok{values =} \FunctionTok{c}\NormalTok{(}\StringTok{"Prior"} \OtherTok{=} \StringTok{"dashed"}\NormalTok{, }\StringTok{"Likelihood"} \OtherTok{=} \StringTok{"dotted"}\NormalTok{,}
               \StringTok{"Posterior"} \OtherTok{=} \StringTok{"solid"}\NormalTok{)}
\NormalTok{  ) }\SpecialCharTok{+}
  \FunctionTok{labs}\NormalTok{(}
    \AttributeTok{title    =} \StringTok{"Bayesian Updating: Prior × Likelihood ∝ Posterior"}\NormalTok{,}
    \AttributeTok{subtitle =} \StringTok{"Uniform prior updated with 7 blue / 2 red marbles"}\NormalTok{,}
    \AttributeTok{x        =} \StringTok{"Belief (proportion blue)"}\NormalTok{,}
    \AttributeTok{y        =} \StringTok{"Probability density"}\NormalTok{,}
    \AttributeTok{colour   =} \ConstantTok{NULL}\NormalTok{, }\AttributeTok{linetype =} \ConstantTok{NULL}
\NormalTok{  ) }\SpecialCharTok{+}
  \FunctionTok{theme}\NormalTok{(}\AttributeTok{legend.position =} \StringTok{"bottom"}\NormalTok{)}
\end{Highlighting}
\end{Shaded}

\begin{center}
\includegraphics[width=0.8\linewidth,height=\textheight,keepaspectratio]{10-BayesianModels_files/figure-pdf/ch10_visualising_bayesian_updating-1.pdf}
\end{center}

The posterior is narrower than either source alone --- combining
information increases certainty. Its peak lies between prior and
likelihood, but closer to the likelihood because the evidence was
relatively strong. It is also more confident than the likelihood, as it
combines compatible evidence from prior and likelihood. Bayesian
integration is not averaging; it is precision-weighted combination (as
we also discussed in chapter 3.7).

\subsection{Bayesian Models in Cognitive
Science}\label{bayesian-models-in-cognitive-science}

Bayesian cognitive models have been successfully applied to:

\begin{itemize}
\tightlist
\item
  \textbf{Perception:} combining multi-sensory cues (visual, auditory,
  tactile) into a unified percept.
\item
  \textbf{Learning:} updating beliefs from observation and instruction.
\item
  \textbf{Decision-making:} weighing different evidence sources when
  making choices.
\item
  \textbf{Social cognition:} integrating others' opinions with one's own
  knowledge.
\item
  \textbf{Psychopathology:} understanding schizophrenia and autism in
  terms of atypical Bayesian inference --- hyper-precise priors, reduced
  likelihood weighting, or atypical weighting of social information.
\end{itemize}

\section{The Setup: Social Influence in Perceptual
Decision-Making}\label{the-setup-social-influence-in-perceptual-decision-making}

To make our discussion of bayesian models concrete, let us consider a
simplified version of a study examining how people integrate information
from different sources (Simonsen et al., 2021).

In the \textbf{marble task}:

\begin{itemize}
\tightlist
\item
  There is a set of urns with varying proportions of blue and red
  marbles
\item
  At each trial the urn changes (no learning is possible across trials)
\item
  At each trial you sample \(n_1 = 8\) marbles → \textbf{Source 1
  (Direct Evidence)}
\item
  At each trial another participant samples (hidden to you) and
  communicates their guess as to the next ball, including their
  confidence → \textbf{Source 2 (Social Evidence)}
\item
  Source 1 ranges from 0 to 8 blue marbles
\item
  Source 2 is encoded as 0--3 (Clear Red, Maybe Red, Maybe Blue, Clear
  Blue)
\item
  The experiment is built so you get all combinations of Source 1 and
  Source 2, with 5 repetitions each
\end{itemize}

This encoding introduces a strict measurement model assumption. The
other person actually drew 8 marbles, but we are compressing their
choice-plus-confidence signal into a smaller equivalent sample size
(\(n_2 = 3\)). Remember, other choices would be possible, and perhaps
better than this, but this is very convenient as it sets both sources on
the same scale.

\subsection{The Beta Distribution
Intuition}\label{the-beta-distribution-intuition}

The Beta distribution provides an elegant way to represent beliefs about
proportions such as the proportion of blue marbles in a jar. Its two
parameters \(\alpha\) and \(\beta\) have an intuitive interpretation:

\begin{itemize}
\tightlist
\item
  \(\alpha\) can be thought of as the number of blue marbles observed
  (plus the prior pseudo-count \(\alpha_0\)).
\item
  \(\beta\) can be thought of as the number of red marbles observed
  (plus the prior pseudo-count \(\beta_0\)).
\end{itemize}

What should we put as prior? We previously chose a uniform distribution
on the 0-1 scale, which corresponds to \(\text{Beta}(1, 1)\). Another
common choice is the Jeffreys prior: \(\text{Beta}(0.5, 0.5)\). The
Jeffreys prior is defined as \emph{parameterisation-invariant}: it
assigns equal information to each value of \(\theta\) \emph{on the
log-odds scale}, which is the natural scale for a Bernoulli likelihood.

As usual, it's easier to understand our choices if we plot the priors
and their consequences.

\begin{Shaded}
\begin{Highlighting}[]
\NormalTok{theta\_grid }\OtherTok{\textless{}{-}} \FunctionTok{seq}\NormalTok{(}\FloatTok{0.001}\NormalTok{, }\FloatTok{0.999}\NormalTok{, }\AttributeTok{length.out =} \DecValTok{500}\NormalTok{)}
\NormalTok{n\_grid     }\OtherTok{\textless{}{-}} \FunctionTok{length}\NormalTok{(theta\_grid)}

\NormalTok{prior\_comparison }\OtherTok{\textless{}{-}} \FunctionTok{tibble}\NormalTok{(}
  \AttributeTok{theta =} \FunctionTok{rep}\NormalTok{(theta\_grid, }\DecValTok{3}\NormalTok{),}
  \AttributeTok{density =} \FunctionTok{c}\NormalTok{(}
    \FunctionTok{dbeta}\NormalTok{(theta\_grid, }\FloatTok{0.5}\NormalTok{, }\FloatTok{0.5}\NormalTok{),}
    \FunctionTok{dbeta}\NormalTok{(theta\_grid, }\DecValTok{1}\NormalTok{, }\DecValTok{1}\NormalTok{),}
    \FunctionTok{dbeta}\NormalTok{(theta\_grid, }\DecValTok{2}\NormalTok{, }\DecValTok{2}\NormalTok{)}
\NormalTok{  ),}
  \AttributeTok{prior =} \FunctionTok{factor}\NormalTok{(}
    \FunctionTok{rep}\NormalTok{(}\FunctionTok{c}\NormalTok{(}\StringTok{"Jeffreys: Beta(0.5, 0.5)"}\NormalTok{,}
          \StringTok{"Uniform: Beta(1, 1)"}\NormalTok{,}
          \StringTok{"Weakly informative: Beta(2, 2)"}\NormalTok{), }\AttributeTok{each =}\NormalTok{ n\_grid),}
    \AttributeTok{levels =} \FunctionTok{c}\NormalTok{(}\StringTok{"Jeffreys: Beta(0.5, 0.5)"}\NormalTok{,}
               \StringTok{"Uniform: Beta(1, 1)"}\NormalTok{,}
               \StringTok{"Weakly informative: Beta(2, 2)"}\NormalTok{)}
\NormalTok{  )}
\NormalTok{)}

\NormalTok{p\_priors }\OtherTok{\textless{}{-}} \FunctionTok{ggplot}\NormalTok{(prior\_comparison, }\FunctionTok{aes}\NormalTok{(}\AttributeTok{x =}\NormalTok{ theta, }\AttributeTok{y =}\NormalTok{ density,}
                                          \AttributeTok{colour =}\NormalTok{ prior, }\AttributeTok{linetype =}\NormalTok{ prior)) }\SpecialCharTok{+}
  \FunctionTok{geom\_line}\NormalTok{(}\AttributeTok{linewidth =} \FloatTok{1.2}\NormalTok{) }\SpecialCharTok{+}
  \FunctionTok{scale\_colour\_manual}\NormalTok{(}\AttributeTok{values =} \FunctionTok{c}\NormalTok{(}\StringTok{"\#d73027"}\NormalTok{, }\StringTok{"\#4575b4"}\NormalTok{, }\StringTok{"\#756bb1"}\NormalTok{)) }\SpecialCharTok{+}
  \FunctionTok{scale\_linetype\_manual}\NormalTok{(}\AttributeTok{values =} \FunctionTok{c}\NormalTok{(}\StringTok{"solid"}\NormalTok{, }\StringTok{"dashed"}\NormalTok{, }\StringTok{"dotted"}\NormalTok{)) }\SpecialCharTok{+}
  \FunctionTok{labs}\NormalTok{(}
    \AttributeTok{title    =} \StringTok{"Three Prior Choices for a Binomial Proportion"}\NormalTok{,}
    \AttributeTok{subtitle =} \StringTok{"Jeffreys concentrates mass near 0 and 1;}\SpecialCharTok{\textbackslash{}n}\StringTok{Uniform is flat; Beta(2,2) pulls toward 0.5"}\NormalTok{,}
    \AttributeTok{x        =} \FunctionTok{expression}\NormalTok{(theta }\SpecialCharTok{\textasciitilde{}} \StringTok{"(proportion blue)"}\NormalTok{),}
    \AttributeTok{y        =} \StringTok{"Prior density"}\NormalTok{,}
    \AttributeTok{colour   =} \ConstantTok{NULL}\NormalTok{, }\AttributeTok{linetype =} \ConstantTok{NULL}
\NormalTok{  ) }\SpecialCharTok{+}
  \FunctionTok{theme}\NormalTok{(}\AttributeTok{legend.position =} \StringTok{"bottom"}\NormalTok{)}

\CommentTok{\# {-}{-}{-} Same priors viewed on the log{-}odds scale {-}{-}{-}}
\CommentTok{\# Transform: if theta \textasciitilde{} Beta(a, b), what does logit(theta) look like?}
\CommentTok{\# We use a change{-}of{-}variables: density on logit scale =}
\CommentTok{\#   dbeta(inv\_logit(x), a, b) * inv\_logit(x) * (1 {-} inv\_logit(x))}
\NormalTok{logit\_grid }\OtherTok{\textless{}{-}} \FunctionTok{seq}\NormalTok{(}\SpecialCharTok{{-}}\DecValTok{6}\NormalTok{, }\DecValTok{6}\NormalTok{, }\AttributeTok{length.out =} \DecValTok{500}\NormalTok{)}
\NormalTok{inv\_logit  }\OtherTok{\textless{}{-}} \FunctionTok{plogis}\NormalTok{(logit\_grid)}
\NormalTok{jacobian   }\OtherTok{\textless{}{-}}\NormalTok{ inv\_logit }\SpecialCharTok{*}\NormalTok{ (}\DecValTok{1} \SpecialCharTok{{-}}\NormalTok{ inv\_logit)}

\NormalTok{logit\_comparison }\OtherTok{\textless{}{-}} \FunctionTok{tibble}\NormalTok{(}
  \AttributeTok{logit\_theta =} \FunctionTok{rep}\NormalTok{(logit\_grid, }\DecValTok{3}\NormalTok{),}
  \AttributeTok{density =} \FunctionTok{c}\NormalTok{(}
    \FunctionTok{dbeta}\NormalTok{(inv\_logit, }\FloatTok{0.5}\NormalTok{, }\FloatTok{0.5}\NormalTok{) }\SpecialCharTok{*}\NormalTok{ jacobian,}
    \FunctionTok{dbeta}\NormalTok{(inv\_logit, }\DecValTok{1}\NormalTok{, }\DecValTok{1}\NormalTok{)     }\SpecialCharTok{*}\NormalTok{ jacobian,}
    \FunctionTok{dbeta}\NormalTok{(inv\_logit, }\DecValTok{2}\NormalTok{, }\DecValTok{2}\NormalTok{)     }\SpecialCharTok{*}\NormalTok{ jacobian}
\NormalTok{  ),}
  \AttributeTok{prior =} \FunctionTok{factor}\NormalTok{(}
    \FunctionTok{rep}\NormalTok{(}\FunctionTok{c}\NormalTok{(}\StringTok{"Jeffreys: Beta(0.5, 0.5)"}\NormalTok{,}
          \StringTok{"Uniform: Beta(1, 1)"}\NormalTok{,}
          \StringTok{"Weakly informative: Beta(2, 2)"}\NormalTok{), }\AttributeTok{each =} \FunctionTok{length}\NormalTok{(logit\_grid)),}
    \AttributeTok{levels =} \FunctionTok{c}\NormalTok{(}\StringTok{"Jeffreys: Beta(0.5, 0.5)"}\NormalTok{,}
               \StringTok{"Uniform: Beta(1, 1)"}\NormalTok{,}
               \StringTok{"Weakly informative: Beta(2, 2)"}\NormalTok{)}
\NormalTok{  )}
\NormalTok{)}

\NormalTok{p\_logit }\OtherTok{\textless{}{-}} \FunctionTok{ggplot}\NormalTok{(logit\_comparison, }\FunctionTok{aes}\NormalTok{(}\AttributeTok{x =}\NormalTok{ logit\_theta, }\AttributeTok{y =}\NormalTok{ density,}
                                         \AttributeTok{colour =}\NormalTok{ prior, }\AttributeTok{linetype =}\NormalTok{ prior)) }\SpecialCharTok{+}
  \FunctionTok{geom\_line}\NormalTok{(}\AttributeTok{linewidth =} \FloatTok{1.2}\NormalTok{) }\SpecialCharTok{+}
  \FunctionTok{scale\_colour\_manual}\NormalTok{(}\AttributeTok{values =} \FunctionTok{c}\NormalTok{(}\StringTok{"\#d73027"}\NormalTok{, }\StringTok{"\#4575b4"}\NormalTok{, }\StringTok{"\#756bb1"}\NormalTok{)) }\SpecialCharTok{+}
  \FunctionTok{scale\_linetype\_manual}\NormalTok{(}\AttributeTok{values =} \FunctionTok{c}\NormalTok{(}\StringTok{"solid"}\NormalTok{, }\StringTok{"dashed"}\NormalTok{, }\StringTok{"dotted"}\NormalTok{)) }\SpecialCharTok{+}
  \FunctionTok{labs}\NormalTok{(}
    \AttributeTok{title    =} \StringTok{"The Same Priors on the Log{-}Odds Scale"}\NormalTok{,}
    \AttributeTok{subtitle =} \StringTok{"Jeffreys is nearly flat — it assigns equal information per unit log{-}odds.}\SpecialCharTok{\textbackslash{}n}\StringTok{Uniform concentrates mass near logit(θ) = 0 (i.e., θ = 0.5)."}\NormalTok{,}
    \AttributeTok{x        =} \FunctionTok{expression}\NormalTok{(}\StringTok{"logit("} \SpecialCharTok{*}\NormalTok{ theta }\SpecialCharTok{*} \StringTok{")"}\NormalTok{),}
    \AttributeTok{y        =} \StringTok{"Density (change of variables)"}\NormalTok{,}
    \AttributeTok{colour   =} \ConstantTok{NULL}\NormalTok{, }\AttributeTok{linetype =} \ConstantTok{NULL}
\NormalTok{  ) }\SpecialCharTok{+}
  \FunctionTok{theme}\NormalTok{(}\AttributeTok{legend.position =} \StringTok{"bottom"}\NormalTok{)}

\CommentTok{\# {-}{-}{-} Posterior comparison after 6 blue, 2 red {-}{-}{-}}
\NormalTok{blue\_obs }\OtherTok{\textless{}{-}} \DecValTok{6}\NormalTok{; red\_obs }\OtherTok{\textless{}{-}} \DecValTok{2}

\NormalTok{posterior\_comparison }\OtherTok{\textless{}{-}} \FunctionTok{tibble}\NormalTok{(}
  \AttributeTok{theta =} \FunctionTok{rep}\NormalTok{(theta\_grid, }\DecValTok{3}\NormalTok{),}
  \AttributeTok{density =} \FunctionTok{c}\NormalTok{(}
    \FunctionTok{dbeta}\NormalTok{(theta\_grid, }\FloatTok{0.5} \SpecialCharTok{+}\NormalTok{ blue\_obs, }\FloatTok{0.5} \SpecialCharTok{+}\NormalTok{ red\_obs),}
    \FunctionTok{dbeta}\NormalTok{(theta\_grid, }\DecValTok{1} \SpecialCharTok{+}\NormalTok{ blue\_obs,   }\DecValTok{1} \SpecialCharTok{+}\NormalTok{ red\_obs),}
    \FunctionTok{dbeta}\NormalTok{(theta\_grid, }\DecValTok{2} \SpecialCharTok{+}\NormalTok{ blue\_obs,   }\DecValTok{2} \SpecialCharTok{+}\NormalTok{ red\_obs)}
\NormalTok{  ),}
  \AttributeTok{prior\_used =} \FunctionTok{factor}\NormalTok{(}
    \FunctionTok{rep}\NormalTok{(}\FunctionTok{c}\NormalTok{(}\StringTok{"Jeffreys → Beta(6.5, 2.5)"}\NormalTok{,}
          \StringTok{"Uniform → Beta(7, 3)"}\NormalTok{,}
          \StringTok{"Beta(2,2) → Beta(8, 4)"}\NormalTok{), }\AttributeTok{each =}\NormalTok{ n\_grid)}
\NormalTok{  )}
\NormalTok{)}

\NormalTok{p\_post }\OtherTok{\textless{}{-}} \FunctionTok{ggplot}\NormalTok{(posterior\_comparison, }\FunctionTok{aes}\NormalTok{(}\AttributeTok{x =}\NormalTok{ theta, }\AttributeTok{y =}\NormalTok{ density,}
                                            \AttributeTok{colour =}\NormalTok{ prior\_used)) }\SpecialCharTok{+}
  \FunctionTok{geom\_line}\NormalTok{(}\AttributeTok{linewidth =} \FloatTok{1.2}\NormalTok{) }\SpecialCharTok{+}
  \FunctionTok{geom\_vline}\NormalTok{(}\AttributeTok{xintercept =} \FloatTok{6.5} \SpecialCharTok{/} \DecValTok{9}\NormalTok{,  }\AttributeTok{linetype =} \StringTok{"dashed"}\NormalTok{,}
             \AttributeTok{colour =} \StringTok{"\#d73027"}\NormalTok{, }\AttributeTok{alpha =} \FloatTok{0.6}\NormalTok{) }\SpecialCharTok{+}
  \FunctionTok{geom\_vline}\NormalTok{(}\AttributeTok{xintercept =} \DecValTok{7} \SpecialCharTok{/} \DecValTok{10}\NormalTok{,   }\AttributeTok{linetype =} \StringTok{"dashed"}\NormalTok{,}
             \AttributeTok{colour =} \StringTok{"\#4575b4"}\NormalTok{, }\AttributeTok{alpha =} \FloatTok{0.6}\NormalTok{) }\SpecialCharTok{+}
  \FunctionTok{geom\_vline}\NormalTok{(}\AttributeTok{xintercept =} \DecValTok{8} \SpecialCharTok{/} \DecValTok{12}\NormalTok{,   }\AttributeTok{linetype =} \StringTok{"dashed"}\NormalTok{,}
             \AttributeTok{colour =} \StringTok{"\#756bb1"}\NormalTok{, }\AttributeTok{alpha =} \FloatTok{0.6}\NormalTok{) }\SpecialCharTok{+}
  \FunctionTok{scale\_colour\_manual}\NormalTok{(}\AttributeTok{values =} \FunctionTok{c}\NormalTok{(}\StringTok{"\#d73027"}\NormalTok{, }\StringTok{"\#4575b4"}\NormalTok{, }\StringTok{"\#756bb1"}\NormalTok{)) }\SpecialCharTok{+}
  \FunctionTok{annotate}\NormalTok{(}\StringTok{"text"}\NormalTok{, }\AttributeTok{x =} \FloatTok{0.15}\NormalTok{, }\AttributeTok{y =} \FunctionTok{max}\NormalTok{(posterior\_comparison}\SpecialCharTok{$}\NormalTok{density) }\SpecialCharTok{*} \FloatTok{0.85}\NormalTok{,}
           \AttributeTok{label =} \FunctionTok{sprintf}\NormalTok{(}\StringTok{"Max Δ mean = \%.3f"}\NormalTok{, }\DecValTok{8}\SpecialCharTok{/}\DecValTok{12} \SpecialCharTok{{-}} \FloatTok{6.5}\SpecialCharTok{/}\DecValTok{9}\NormalTok{),}
           \AttributeTok{size =} \FloatTok{3.5}\NormalTok{, }\AttributeTok{colour =} \StringTok{"grey30"}\NormalTok{) }\SpecialCharTok{+}
  \FunctionTok{labs}\NormalTok{(}
    \AttributeTok{title    =} \StringTok{"Posteriors After 6 Blue, 2 Red: The Practical Difference"}\NormalTok{,}
    \AttributeTok{subtitle =} \StringTok{"With 8 observations the prior choice barely matters (max Δ ≈ 0.06)"}\NormalTok{,}
    \AttributeTok{x        =} \FunctionTok{expression}\NormalTok{(theta),}
    \AttributeTok{y        =} \StringTok{"Posterior density"}\NormalTok{,}
    \AttributeTok{colour   =} \ConstantTok{NULL}
\NormalTok{  ) }\SpecialCharTok{+}
  \FunctionTok{theme}\NormalTok{(}\AttributeTok{legend.position =} \StringTok{"bottom"}\NormalTok{)}

\CommentTok{\# Compose all three panels}
\NormalTok{p\_priors }
\end{Highlighting}
\end{Shaded}

\begin{center}
\includegraphics[width=0.8\linewidth,height=\textheight,keepaspectratio]{10-BayesianModels_files/figure-pdf/ch10_prior_comparison-1.pdf}
\end{center}

\begin{Shaded}
\begin{Highlighting}[]
\NormalTok{p\_logit }
\end{Highlighting}
\end{Shaded}

\begin{center}
\includegraphics[width=0.8\linewidth,height=\textheight,keepaspectratio]{10-BayesianModels_files/figure-pdf/ch10_prior_comparison-2.pdf}
\end{center}

\begin{Shaded}
\begin{Highlighting}[]
\NormalTok{p\_post }
\end{Highlighting}
\end{Shaded}

\begin{center}
\includegraphics[width=0.8\linewidth,height=\textheight,keepaspectratio]{10-BayesianModels_files/figure-pdf/ch10_prior_comparison-3.pdf}
\end{center}

The three panels above tell a clear story. On the probability scale
(top), the Jeffreys prior looks extreme --- it concentrates mass near 0
and 1, as if expecting jars to be mostly one colour. The uniform prior
looks neutral. But as usual, this is a matter of scale: on the log-odds
scale, where the Bernoulli likelihood naturally lives, Jeffreys is the
one that is approximately flat. The uniform prior, seemingly innocuous,
quietly concentrates mass near logit(θ)=0 --- encoding a more pronounced
preference for balanced jars than the Jeffreys prior. The Beta(2,2)
prior goes further, explicitly pulling toward \(θ = 0.5\).

I generally prefer to think in probability (0-1) scale, but others
prefer the log-odds scale. As usual, it's a good idea to investigate the
consequences of this choice. The bottom panel shows that the
consequences are not drastic. After observing just 6 blue and 2 red
marbles, the three posteriors are nearly indistinguishable: the Jeffreys
prior gives Beta(6.5,2.5) with mean ≈ 0.72, the uniform gives Beta(7,3)
with mean = 0.70, and the Beta(2,2) prior gives Beta(8,4) with mean ≈
0.67. The maximum shift across all three is ≈0.06 --- the data have
overwhelmed the prior after a single sample of n = 8 marbles. Just to
challenge myself to choices I am not used to, we use the Jeffreys prior
throughout this chapter.

\subsection{Setting Up the Experimental
Design}\label{setting-up-the-experimental-design}

\begin{Shaded}
\begin{Highlighting}[]
\CommentTok{\# Create the full factorial design}
\NormalTok{evidence\_combinations }\OtherTok{\textless{}{-}} \FunctionTok{expand\_grid}\NormalTok{(}
  \AttributeTok{blue1  =} \DecValTok{0}\SpecialCharTok{:}\DecValTok{8}\NormalTok{,   }\CommentTok{\# Direct evidence: 0{-}8 blue marbles out of 8}
  \AttributeTok{blue2  =} \DecValTok{0}\SpecialCharTok{:}\DecValTok{3}\NormalTok{,   }\CommentTok{\# Social evidence: encoded as 0{-}3 (pseudo{-}marbles)}
  \AttributeTok{total1 =} \DecValTok{8}\NormalTok{,     }\CommentTok{\# Fixed sample size for direct evidence}
  \AttributeTok{total2 =} \DecValTok{3}      \CommentTok{\# Fixed "sample size" for social evidence}
\NormalTok{)}

\CommentTok{\# This gives us 9 × 4 = 36 unique evidence combinations}
\FunctionTok{cat}\NormalTok{(}\StringTok{"Unique evidence combinations:"}\NormalTok{, }\FunctionTok{nrow}\NormalTok{(evidence\_combinations), }\StringTok{"}\SpecialCharTok{\textbackslash{}n}\StringTok{"}\NormalTok{)}
\end{Highlighting}
\end{Shaded}

\begin{verbatim}
Unique evidence combinations: 36 
\end{verbatim}

\begin{Shaded}
\begin{Highlighting}[]
\FunctionTok{cat}\NormalTok{(}\StringTok{"With 5 repetitions each:"}\NormalTok{, }\FunctionTok{nrow}\NormalTok{(evidence\_combinations) }\SpecialCharTok{*} \DecValTok{5}\NormalTok{, }\StringTok{"trials per agent}\SpecialCharTok{\textbackslash{}n}\StringTok{"}\NormalTok{)}
\end{Highlighting}
\end{Shaded}

\begin{verbatim}
With 5 repetitions each: 180 trials per agent
\end{verbatim}

\section{The Beta-Binomial
Likelihood}\label{the-beta-binomial-likelihood}

How do we model evidence when it arrives as counts --- k blue marbles
out of n total observed? The Beta distribution provides a natural
representation of beliefs about a proportion, and combined with a
Bernoulli choice it gives us the \textbf{Beta-Binomial likelihood} that
both agents in this chapter will use.

The two parameters \(\alpha\) and \(\beta\) have an intuitive
interpretation as pseudo-counts:

\begin{itemize}
\tightlist
\item
  \(\alpha\) = prior pseudo-count + (weighted) blue marbles from all
  sources
\item
  \(\beta\) = prior pseudo-count + (weighted) red marbles from all
  sources
\end{itemize}

Using the Jeffreys prior (\(\alpha_0 = \beta_0 = 0.5\)), the posterior
parameters after observing direct evidence \((k_1, n_1)\) and social
evidence \((k_2, n_2)\) with weights \(w_d, w_s\) are:

\[\alpha_{\text{post}} = 0.5 + w_d \cdot k_1 + w_s \cdot k_2\]
\[\beta_{\text{post}}  = 0.5 + w_d \cdot (n_1 - k_1) + w_s \cdot (n_2 - k_2)\]

Setting \(w_d = w_s = 1\) recovers strict Bayesian updating (Simple
Bayesian Agent, SBA); letting them vary will give us the Weighted
Bayesian Agent (WBA). We'll meet both formally in the next section. For
now we focus on the likelihood itself.

\begin{Shaded}
\begin{Highlighting}[]
\CommentTok{\# Core belief{-}update function implementing the Beta{-}Binomial integration model}
\NormalTok{betaBinomialModel }\OtherTok{\textless{}{-}} \ControlFlowTok{function}\NormalTok{(}\AttributeTok{alpha\_prior =} \FloatTok{0.5}\NormalTok{, }\AttributeTok{beta\_prior =} \FloatTok{0.5}\NormalTok{,}
\NormalTok{                               blue1, total1, blue2, total2,}
                               \AttributeTok{w\_direct =} \DecValTok{1}\NormalTok{, }\AttributeTok{w\_social =} \DecValTok{1}\NormalTok{) \{}
  
\NormalTok{  alpha\_post }\OtherTok{\textless{}{-}}\NormalTok{ alpha\_prior }\SpecialCharTok{+}\NormalTok{ w\_direct }\SpecialCharTok{*}\NormalTok{ blue1        }\SpecialCharTok{+}\NormalTok{ w\_social }\SpecialCharTok{*}\NormalTok{ blue2}
\NormalTok{  beta\_post  }\OtherTok{\textless{}{-}}\NormalTok{ beta\_prior  }\SpecialCharTok{+}\NormalTok{ w\_direct }\SpecialCharTok{*}\NormalTok{ (total1 }\SpecialCharTok{{-}}\NormalTok{ blue1) }\SpecialCharTok{+}
\NormalTok{                               w\_social }\SpecialCharTok{*}\NormalTok{ (total2 }\SpecialCharTok{{-}}\NormalTok{ blue2)}

\NormalTok{  expected\_rate }\OtherTok{\textless{}{-}}\NormalTok{ alpha\_post }\SpecialCharTok{/}\NormalTok{ (alpha\_post }\SpecialCharTok{+}\NormalTok{ beta\_post)}
\NormalTok{  ci\_lower      }\OtherTok{\textless{}{-}} \FunctionTok{qbeta}\NormalTok{(}\FloatTok{0.025}\NormalTok{, alpha\_post, beta\_post)}
\NormalTok{  ci\_upper      }\OtherTok{\textless{}{-}} \FunctionTok{qbeta}\NormalTok{(}\FloatTok{0.975}\NormalTok{, alpha\_post, beta\_post)}

  \FunctionTok{list}\NormalTok{(}
    \AttributeTok{alpha\_post    =}\NormalTok{ alpha\_post,}
    \AttributeTok{beta\_post     =}\NormalTok{ beta\_post,}
    \AttributeTok{expected\_rate =}\NormalTok{ expected\_rate,}
    \AttributeTok{ci\_lower      =}\NormalTok{ ci\_lower,}
    \AttributeTok{ci\_upper      =}\NormalTok{ ci\_upper,}
    \AttributeTok{choice        =} \FunctionTok{ifelse}\NormalTok{(expected\_rate }\SpecialCharTok{\textgreater{}} \FloatTok{0.5}\NormalTok{, }\StringTok{"Blue"}\NormalTok{, }\StringTok{"Red"}\NormalTok{)}
\NormalTok{  )}
\NormalTok{\}}

\CommentTok{\# Function to generate decisions from a simulated agent}
\NormalTok{generate\_agent\_decisions }\OtherTok{\textless{}{-}} \ControlFlowTok{function}\NormalTok{(weight\_direct, weight\_social,}
\NormalTok{                                      evidence\_df, }\AttributeTok{n\_samples =} \DecValTok{5}\NormalTok{) \{}
\NormalTok{  evidence\_df }\SpecialCharTok{|\textgreater{}}
    \FunctionTok{slice}\NormalTok{(}\FunctionTok{rep}\NormalTok{(}\FunctionTok{seq\_len}\NormalTok{(}\FunctionTok{n}\NormalTok{()), }\AttributeTok{each =}\NormalTok{ n\_samples)) }\SpecialCharTok{|\textgreater{}}
    \FunctionTok{group\_by}\NormalTok{(blue1, blue2, total1, total2) }\SpecialCharTok{|\textgreater{}}
    \FunctionTok{mutate}\NormalTok{(}\AttributeTok{sample\_id =} \FunctionTok{seq\_len}\NormalTok{(}\FunctionTok{n}\NormalTok{())) }\SpecialCharTok{|\textgreater{}}
    \FunctionTok{ungroup}\NormalTok{() }\SpecialCharTok{|\textgreater{}}
    \FunctionTok{pmap\_dfr}\NormalTok{(}\ControlFlowTok{function}\NormalTok{(blue1, blue2, total1, total2, sample\_id) \{}
\NormalTok{      result }\OtherTok{\textless{}{-}} \FunctionTok{betaBinomialModel}\NormalTok{(}
        \FloatTok{0.5}\NormalTok{, }\FloatTok{0.5}\NormalTok{, blue1, total1, blue2, total2,}
        \AttributeTok{w\_direct =}\NormalTok{ weight\_direct, }\AttributeTok{w\_social =}\NormalTok{ weight\_social}
\NormalTok{      )}
      \CommentTok{\# Probability matching: sample θ from Beta, then sample choice from Bernoulli(θ)}
\NormalTok{      alpha\_p }\OtherTok{\textless{}{-}}\NormalTok{ result}\SpecialCharTok{$}\NormalTok{alpha\_post}
\NormalTok{      beta\_p  }\OtherTok{\textless{}{-}}\NormalTok{ result}\SpecialCharTok{$}\NormalTok{beta\_post}
\NormalTok{      p       }\OtherTok{\textless{}{-}} \FunctionTok{rbeta}\NormalTok{(}\DecValTok{1}\NormalTok{, alpha\_p, beta\_p)}
      \FunctionTok{tibble}\NormalTok{(}
        \AttributeTok{sample\_id     =}\NormalTok{ sample\_id,}
        \AttributeTok{blue1         =}\NormalTok{ blue1,}
        \AttributeTok{blue2         =}\NormalTok{ blue2,}
        \AttributeTok{total1        =}\NormalTok{ total1,}
        \AttributeTok{total2        =}\NormalTok{ total2,}
        \AttributeTok{expected\_rate =}\NormalTok{ result}\SpecialCharTok{$}\NormalTok{expected\_rate,}
        \AttributeTok{choice        =} \FunctionTok{rbinom}\NormalTok{(}\DecValTok{1}\NormalTok{, }\DecValTok{1}\NormalTok{, p),}
        \AttributeTok{choice\_label  =} \FunctionTok{ifelse}\NormalTok{(choice }\SpecialCharTok{==} \DecValTok{1}\NormalTok{, }\StringTok{"Blue"}\NormalTok{, }\StringTok{"Red"}\NormalTok{)}
\NormalTok{      )}
\NormalTok{    \})}
\NormalTok{\}}

\NormalTok{prepare\_stan\_data }\OtherTok{\textless{}{-}} \ControlFlowTok{function}\NormalTok{(df) \{}
  \FunctionTok{list}\NormalTok{(}
    \AttributeTok{N       =} \FunctionTok{nrow}\NormalTok{(df),}
    \AttributeTok{choice  =}\NormalTok{ df}\SpecialCharTok{$}\NormalTok{choice,}
    \AttributeTok{N\_total =} \FunctionTok{nrow}\NormalTok{(df),}
    \AttributeTok{y       =}\NormalTok{ df}\SpecialCharTok{$}\NormalTok{choice,}
    \AttributeTok{blue1   =}\NormalTok{ df}\SpecialCharTok{$}\NormalTok{blue1,}
    \AttributeTok{total1  =}\NormalTok{ df}\SpecialCharTok{$}\NormalTok{total1,}
    \AttributeTok{blue2   =}\NormalTok{ df}\SpecialCharTok{$}\NormalTok{blue2,}
    \AttributeTok{total2  =}\NormalTok{ df}\SpecialCharTok{$}\NormalTok{total2,}
    \AttributeTok{prior\_rho\_alpha   =} \FloatTok{2.0}\NormalTok{,}
    \AttributeTok{prior\_rho\_beta    =} \FloatTok{2.0}\NormalTok{,}
    \AttributeTok{prior\_kappa\_mu    =} \FloatTok{0.5}\NormalTok{,}
    \AttributeTok{prior\_kappa\_sigma =} \FloatTok{0.5}\NormalTok{,}
    \AttributeTok{run\_diagnostics =} \DecValTok{1}
\NormalTok{  )}
\NormalTok{\}}
\end{Highlighting}
\end{Shaded}

\subsection{The Observation Model: A Separate Modelling
Dimension}\label{the-observation-model-a-separate-modelling-dimension}

The Bayesian agents as we have implemented them specify how they
\emph{integrate evidence} --- but they also embed a specific assumption
about how beliefs become \emph{actions}. Indeed, the
BetaBinomial(\(1, \alpha, \beta\)) observation model entails a two-step
stochastic process: first sample a belief
\(\theta \sim \text{Beta}(\alpha, \beta)\), then sample a choice
\(y \sim \text{Bernoulli}(\theta)\). This is \textbf{probability
matching} --- the agent's action is a random sample from its full
posterior distribution (rather than from the averate \(\theta\)) over
the jar's composition.

I like this way of doing things, because all the uncertainty is
preserved and transmitted across all phases of the model. However, this
is in principle a testable cognitive hypothesis. A different agent might
act more deterministically, choosing blue whenever
\(\text{E}[\theta] > 0.5\), or following a graded rule where higher
confidence makes the choice more certain. We can formalise these
alternatives:

\begin{longtable}[]{@{}
  >{\raggedright\arraybackslash}p{(\linewidth - 4\tabcolsep) * \real{0.3333}}
  >{\raggedright\arraybackslash}p{(\linewidth - 4\tabcolsep) * \real{0.3333}}
  >{\raggedright\arraybackslash}p{(\linewidth - 4\tabcolsep) * \real{0.3333}}@{}}
\toprule\noalign{}
\begin{minipage}[b]{\linewidth}\raggedright
Observation model
\end{minipage} & \begin{minipage}[b]{\linewidth}\raggedright
Cognitive claim
\end{minipage} & \begin{minipage}[b]{\linewidth}\raggedright
Distribution
\end{minipage} \\
\midrule\noalign{}
\endhead
\bottomrule\noalign{}
\endlastfoot
BetaBinomial (\(n=1\)) & \textbf{Probability matching:} action sampled
from belief & \(y \sim \text{BetaBin}(1, \alpha, \beta)\) \\
Softmax & \textbf{Graded determinism:} posterior mean compressed by
inverse temperature \(\phi\) &
\(y \sim \text{Bern}(\sigma(\phi \cdot \text{logit}(\hat\theta)))\) \\
Threshold & \textbf{Maximisation:} act on the most probable belief &
\(y = \mathbf{1}[\alpha > \beta]\) \\
\end{longtable}

The BetaBinomial propagates the agent's \emph{full uncertainty about the
jar} into the response: an agent who is genuinely uncertain about
\(\theta\) produces more variable responses than one who is confident.
This excess variability (overdispersion relative to a plain Binomial) is
not noise added on top of the model --- it \emph{is} the model's
cognitive claim. The softmax model, by contrast, collapses the posterior
to its mean \(\hat{\theta} = \alpha / (\alpha + \beta)\) and then
applies a noise parameter. It discards distributional information (but
it's an empirical question whether humans would do the same).

\subsection{How the BetaBinomial Generates a Distribution of
Choices}\label{how-the-betabinomial-generates-a-distribution-of-choices}

\textbf{The BetaBinomial likelihood is not a deterministic threshold
rule.} It describes a two-step stochastic process:

\begin{enumerate}
\def\labelenumi{\arabic{enumi}.}
\item
  \textbf{Draw a belief:} Sample the true proportion
  \(\theta \sim \text{Beta}(\alpha, \beta)\), representing genuine
  uncertainty about the jar's composition.
\item
  \textbf{Draw a choice:} Given that \(\theta\), sample a binary choice
  \(y \sim \text{Bernoulli}(\theta)\).
\end{enumerate}

The consequence is that even when the posterior mean is 0.70 in favour
of blue, the agent will not choose blue exactly 70\% of the time in a
perfectly mechanical way. There is variability from two sources
simultaneously: uncertainty about what the jar proportion actually
\emph{is} (step 1), and the irreducible randomness of acting on that
belief with a single guess (step 2). This variability is not noise added
on top of the model --- it \emph{is} the model. The Beta distribution
encodes what a rational agent should believe; the Bernoulli draw
captures the randomness of a single decision.

This two-stage structure has a measurable implication: the BetaBinomial
has \textbf{greater variance than a Binomial} with the same mean. This
excess variance is called \textbf{overdispersion}, and it arises
precisely because \(\theta\) must itself be sampled from the Beta rather
than being known.

\begin{Shaded}
\begin{Highlighting}[]
\CommentTok{\# Scenario: 4 blue of 8 direct + 2 blue of 3 social (Jeffreys prior)}
\NormalTok{alpha\_demo }\OtherTok{\textless{}{-}} \FloatTok{0.5} \SpecialCharTok{+} \DecValTok{4} \SpecialCharTok{+} \DecValTok{2}   \CommentTok{\# = 6.5}
\NormalTok{beta\_demo  }\OtherTok{\textless{}{-}} \FloatTok{0.5} \SpecialCharTok{+} \DecValTok{4} \SpecialCharTok{+} \DecValTok{1}   \CommentTok{\# = 5.5}
\NormalTok{n\_sims     }\OtherTok{\textless{}{-}} \DecValTok{10000}

\CommentTok{\# To see overdispersion, we must simulate MULTIPLE choices per agent.}
\CommentTok{\# Let\textquotesingle{}s assume each simulated agent makes 20 choices.}
\NormalTok{n\_choices  }\OtherTok{\textless{}{-}} \DecValTok{20} 

\NormalTok{expected\_rate\_demo }\OtherTok{\textless{}{-}}\NormalTok{ alpha\_demo }\SpecialCharTok{/}\NormalTok{ (alpha\_demo }\SpecialCharTok{+}\NormalTok{ beta\_demo)}

\CommentTok{\# {-}{-}{-}{-}{-}{-}{-}{-}{-}{-}{-}{-}{-}{-}{-}{-}{-}{-}{-}{-}{-}{-}{-}{-}{-}{-}{-}{-}{-}{-}{-}{-}{-}{-}{-}{-}{-}{-}{-}{-}{-}{-}{-}{-}{-}{-}{-}{-}{-}{-}{-}{-}{-}{-}{-}{-}{-}}
\CommentTok{\# Step 1: Simulate Both Models}
\CommentTok{\# {-}{-}{-}{-}{-}{-}{-}{-}{-}{-}{-}{-}{-}{-}{-}{-}{-}{-}{-}{-}{-}{-}{-}{-}{-}{-}{-}{-}{-}{-}{-}{-}{-}{-}{-}{-}{-}{-}{-}{-}{-}{-}{-}{-}{-}{-}{-}{-}{-}{-}{-}{-}{-}{-}{-}{-}{-}}

\CommentTok{\# Model A: Beta{-}Binomial (Uncertainty propagates)}
\CommentTok{\# Draw theta from Beta, then draw successes based on that theta}
\NormalTok{theta\_draws }\OtherTok{\textless{}{-}} \FunctionTok{rbeta}\NormalTok{(n\_sims, alpha\_demo, beta\_demo)}
\NormalTok{bb\_draws    }\OtherTok{\textless{}{-}} \FunctionTok{rbinom}\NormalTok{(n\_sims, }\AttributeTok{size =}\NormalTok{ n\_choices, }\AttributeTok{prob =}\NormalTok{ theta\_draws)}

\CommentTok{\# Model B: Standard Binomial (Fixed belief, no uncertainty in theta)}
\CommentTok{\# Draw successes using the fixed expected rate}
\NormalTok{binom\_draws }\OtherTok{\textless{}{-}} \FunctionTok{rbinom}\NormalTok{(n\_sims, }\AttributeTok{size =}\NormalTok{ n\_choices, }\AttributeTok{prob =}\NormalTok{ expected\_rate\_demo)}

\CommentTok{\# {-}{-}{-}{-}{-}{-}{-}{-}{-}{-}{-}{-}{-}{-}{-}{-}{-}{-}{-}{-}{-}{-}{-}{-}{-}{-}{-}{-}{-}{-}{-}{-}{-}{-}{-}{-}{-}{-}{-}{-}{-}{-}{-}{-}{-}{-}{-}{-}{-}{-}{-}{-}{-}{-}{-}{-}{-}}
\CommentTok{\# Step 2: Visualization}
\CommentTok{\# {-}{-}{-}{-}{-}{-}{-}{-}{-}{-}{-}{-}{-}{-}{-}{-}{-}{-}{-}{-}{-}{-}{-}{-}{-}{-}{-}{-}{-}{-}{-}{-}{-}{-}{-}{-}{-}{-}{-}{-}{-}{-}{-}{-}{-}{-}{-}{-}{-}{-}{-}{-}{-}{-}{-}{-}{-}}

\CommentTok{\# Plot 1: The uncertainty in theta (similar to your left panel)}
\NormalTok{p\_theta }\OtherTok{\textless{}{-}} \FunctionTok{tibble}\NormalTok{(}\AttributeTok{theta =}\NormalTok{ theta\_draws) }\SpecialCharTok{|\textgreater{}}
  \FunctionTok{ggplot}\NormalTok{(}\FunctionTok{aes}\NormalTok{(}\AttributeTok{x =}\NormalTok{ theta)) }\SpecialCharTok{+}
  \FunctionTok{geom\_density}\NormalTok{(}\AttributeTok{fill =} \StringTok{"\#4575b4"}\NormalTok{, }\AttributeTok{alpha =} \FloatTok{0.5}\NormalTok{, }\AttributeTok{color =} \StringTok{"transparent"}\NormalTok{) }\SpecialCharTok{+}
  \FunctionTok{geom\_vline}\NormalTok{(}\AttributeTok{xintercept =}\NormalTok{ expected\_rate\_demo,}
             \AttributeTok{colour =} \StringTok{"\#d73027"}\NormalTok{, }\AttributeTok{linewidth =} \DecValTok{1}\NormalTok{, }\AttributeTok{linetype =} \StringTok{"dashed"}\NormalTok{) }\SpecialCharTok{+}
  \FunctionTok{scale\_x\_continuous}\NormalTok{(}\AttributeTok{limits =} \FunctionTok{c}\NormalTok{(}\DecValTok{0}\NormalTok{, }\DecValTok{1}\NormalTok{), }\AttributeTok{labels =}\NormalTok{ scales}\SpecialCharTok{::}\NormalTok{percent) }\SpecialCharTok{+}
  \FunctionTok{labs}\NormalTok{(}
    \AttributeTok{title    =} \StringTok{"Step 1: Uncertainty in Belief (θ)"}\NormalTok{,}
    \AttributeTok{subtitle =} \FunctionTok{sprintf}\NormalTok{(}\StringTok{"Beta(\%.1f, \%.1f) | Dashed line is E[θ] = \%.0f\%\%"}\NormalTok{, }
\NormalTok{                       alpha\_demo, beta\_demo, expected\_rate\_demo }\SpecialCharTok{*} \DecValTok{100}\NormalTok{),}
    \AttributeTok{x =} \FunctionTok{expression}\NormalTok{(}\StringTok{"True underlying rate ("} \SpecialCharTok{*}\NormalTok{ theta }\SpecialCharTok{*} \StringTok{")"}\NormalTok{),}
    \AttributeTok{y =} \StringTok{"Density"}
\NormalTok{  )}

\CommentTok{\# Plot 2: The outcome comparison showing overdispersion}
\NormalTok{plot\_data }\OtherTok{\textless{}{-}} \FunctionTok{tibble}\NormalTok{(}
  \StringTok{\textasciigrave{}}\AttributeTok{Beta{-}Binomial (Uncertain θ)}\StringTok{\textasciigrave{}} \OtherTok{=}\NormalTok{ bb\_draws,}
  \StringTok{\textasciigrave{}}\AttributeTok{Standard Binomial (Fixed θ)}\StringTok{\textasciigrave{}} \OtherTok{=}\NormalTok{ binom\_draws}
\NormalTok{) }\SpecialCharTok{|\textgreater{}} 
  \FunctionTok{pivot\_longer}\NormalTok{(}\AttributeTok{cols =} \FunctionTok{everything}\NormalTok{(), }\AttributeTok{names\_to =} \StringTok{"Model"}\NormalTok{, }\AttributeTok{values\_to =} \StringTok{"Successes"}\NormalTok{)}

\NormalTok{p\_spread }\OtherTok{\textless{}{-}} \FunctionTok{ggplot}\NormalTok{(plot\_data, }\FunctionTok{aes}\NormalTok{(}\AttributeTok{x =}\NormalTok{ Successes, }\AttributeTok{fill =}\NormalTok{ Model)) }\SpecialCharTok{+}
  \FunctionTok{geom\_bar}\NormalTok{(}\AttributeTok{position =} \StringTok{"identity"}\NormalTok{, }\AttributeTok{alpha =} \FloatTok{0.6}\NormalTok{, }\FunctionTok{aes}\NormalTok{(}\AttributeTok{y =} \FunctionTok{after\_stat}\NormalTok{(prop))) }\SpecialCharTok{+}
  \FunctionTok{scale\_fill\_manual}\NormalTok{(}\AttributeTok{values =} \FunctionTok{c}\NormalTok{(}\StringTok{"Standard Binomial (Fixed θ)"} \OtherTok{=} \StringTok{"\#d73027"}\NormalTok{, }
                               \StringTok{"Beta{-}Binomial (Uncertain θ)"} \OtherTok{=} \StringTok{"\#4575b4"}\NormalTok{)) }\SpecialCharTok{+}
  \FunctionTok{scale\_x\_continuous}\NormalTok{(}\AttributeTok{breaks =} \FunctionTok{seq}\NormalTok{(}\DecValTok{0}\NormalTok{, n\_choices, }\AttributeTok{by =} \DecValTok{2}\NormalTok{)) }\SpecialCharTok{+}
  \FunctionTok{scale\_y\_continuous}\NormalTok{(}\AttributeTok{labels =}\NormalTok{ scales}\SpecialCharTok{::}\NormalTok{percent) }\SpecialCharTok{+}
  \FunctionTok{labs}\NormalTok{(}
    \AttributeTok{title    =} \FunctionTok{sprintf}\NormalTok{(}\StringTok{"Step 2: Choices out of \%d Trials"}\NormalTok{, n\_choices),}
    \AttributeTok{subtitle =} \StringTok{"Notice the wider spread (fatter tails) of the Beta{-}Binomial"}\NormalTok{,}
    \AttributeTok{x =} \StringTok{"Number of \textquotesingle{}Blue\textquotesingle{} choices"}\NormalTok{,}
    \AttributeTok{y =} \StringTok{"Probability"}\NormalTok{,}
    \AttributeTok{fill =} \ConstantTok{NULL}
\NormalTok{  ) }\SpecialCharTok{+}
  \FunctionTok{theme}\NormalTok{(}\AttributeTok{legend.position =} \StringTok{"top"}\NormalTok{)}

\CommentTok{\# Combine plots}
\NormalTok{p\_theta }\SpecialCharTok{+}\NormalTok{ p\_spread }\SpecialCharTok{+}
  \FunctionTok{plot\_annotation}\NormalTok{(}
    \AttributeTok{title    =} \StringTok{"How Uncertainty Propagates: Beta{-}Binomial vs Binomial"}\NormalTok{,}
    \AttributeTok{subtitle =} \StringTok{"When an agent makes multiple choices, uncertainty about the true rate creates a wider distribution of outcomes."}
\NormalTok{  )}
\end{Highlighting}
\end{Shaded}

\begin{center}
\includegraphics[width=0.8\linewidth,height=\textheight,keepaspectratio]{10-BayesianModels_files/figure-pdf/ch10_betabinomial_two_step_demo-1.pdf}
\end{center}

The left panel shows that even in a relatively informative scenario (11
total observations), there is genuine uncertainty about \(\theta\). The
right panel shows that this uncertainty propagates into choices: the
agent does not produce a deterministic .54 blue rate but a
\emph{distribution} of outcomes whose spread exceeds what a Binomial
with \(p = .54\) would produce.

\subsection{Predictions Under Equal
Weights}\label{predictions-under-equal-weights}

\begin{Shaded}
\begin{Highlighting}[]
\CommentTok{\# Calculate SBA predictions for all evidence combinations}
\NormalTok{sim\_data }\OtherTok{\textless{}{-}}\NormalTok{ evidence\_combinations }\SpecialCharTok{|\textgreater{}}
  \FunctionTok{pmap\_dfr}\NormalTok{(}\ControlFlowTok{function}\NormalTok{(blue1, blue2, total1, total2) \{}
\NormalTok{    result }\OtherTok{\textless{}{-}} \FunctionTok{betaBinomialModel}\NormalTok{(}\FloatTok{0.5}\NormalTok{, }\FloatTok{0.5}\NormalTok{, blue1, total1, blue2, total2)}
    \FunctionTok{tibble}\NormalTok{(}
      \AttributeTok{blue1         =}\NormalTok{ blue1,}
      \AttributeTok{blue2         =}\NormalTok{ blue2,}
      \AttributeTok{expected\_rate =}\NormalTok{ result}\SpecialCharTok{$}\NormalTok{expected\_rate,}
      \AttributeTok{ci\_lower      =}\NormalTok{ result}\SpecialCharTok{$}\NormalTok{ci\_lower,}
      \AttributeTok{ci\_upper      =}\NormalTok{ result}\SpecialCharTok{$}\NormalTok{ci\_upper}
\NormalTok{    )}
\NormalTok{  \}) }\SpecialCharTok{|\textgreater{}}
  \FunctionTok{mutate}\NormalTok{(}
    \AttributeTok{social\_evidence =} \FunctionTok{factor}\NormalTok{(}
\NormalTok{      blue2, }\AttributeTok{levels =} \DecValTok{0}\SpecialCharTok{:}\DecValTok{3}\NormalTok{,}
      \AttributeTok{labels =} \FunctionTok{c}\NormalTok{(}\StringTok{"Clear Red"}\NormalTok{, }\StringTok{"Maybe Red"}\NormalTok{, }\StringTok{"Maybe Blue"}\NormalTok{, }\StringTok{"Clear Blue"}\NormalTok{)}
\NormalTok{    )}
\NormalTok{  )}

\NormalTok{semantic\_colors }\OtherTok{\textless{}{-}} \FunctionTok{c}\NormalTok{(}
  \StringTok{"Clear Red"}  \OtherTok{=} \StringTok{"\#d73027"}\NormalTok{,}
  \StringTok{"Maybe Red"}  \OtherTok{=} \StringTok{"\#fdae61"}\NormalTok{,}
  \StringTok{"Maybe Blue"} \OtherTok{=} \StringTok{"\#abd9e9"}\NormalTok{,}
  \StringTok{"Clear Blue"} \OtherTok{=} \StringTok{"\#4575b4"}
\NormalTok{)}

\FunctionTok{ggplot}\NormalTok{(sim\_data,}
       \FunctionTok{aes}\NormalTok{(}\AttributeTok{x =}\NormalTok{ blue1, }\AttributeTok{y =}\NormalTok{ expected\_rate,}
           \AttributeTok{colour =}\NormalTok{ social\_evidence, }\AttributeTok{group =}\NormalTok{ social\_evidence)) }\SpecialCharTok{+}
  \FunctionTok{geom\_ribbon}\NormalTok{(}\FunctionTok{aes}\NormalTok{(}\AttributeTok{ymin =}\NormalTok{ ci\_lower, }\AttributeTok{ymax =}\NormalTok{ ci\_upper,}
                  \AttributeTok{fill =}\NormalTok{ social\_evidence), }\AttributeTok{alpha =} \FloatTok{0.15}\NormalTok{, }\AttributeTok{colour =} \ConstantTok{NA}\NormalTok{) }\SpecialCharTok{+}
  \FunctionTok{geom\_line}\NormalTok{(}\AttributeTok{linewidth =} \DecValTok{1}\NormalTok{) }\SpecialCharTok{+}
  \FunctionTok{geom\_point}\NormalTok{(}\AttributeTok{size =} \DecValTok{3}\NormalTok{) }\SpecialCharTok{+}
  \FunctionTok{geom\_hline}\NormalTok{(}\AttributeTok{yintercept =} \FloatTok{0.5}\NormalTok{, }\AttributeTok{linetype =} \StringTok{"dashed"}\NormalTok{, }\AttributeTok{colour =} \StringTok{"grey50"}\NormalTok{) }\SpecialCharTok{+}
  \FunctionTok{scale\_x\_continuous}\NormalTok{(}\AttributeTok{breaks =} \DecValTok{0}\SpecialCharTok{:}\DecValTok{8}\NormalTok{) }\SpecialCharTok{+}
  \FunctionTok{scale\_colour\_manual}\NormalTok{(}\AttributeTok{values =}\NormalTok{ semantic\_colors) }\SpecialCharTok{+}
  \FunctionTok{scale\_fill\_manual}\NormalTok{(}\AttributeTok{values =}\NormalTok{ semantic\_colors) }\SpecialCharTok{+}
  \FunctionTok{coord\_cartesian}\NormalTok{(}\AttributeTok{ylim =} \FunctionTok{c}\NormalTok{(}\DecValTok{0}\NormalTok{, }\DecValTok{1}\NormalTok{)) }\SpecialCharTok{+}
  \FunctionTok{labs}\NormalTok{(}
    \AttributeTok{title    =} \StringTok{"Simple Bayesian Agent: Evidence Integration"}\NormalTok{,}
    \AttributeTok{subtitle =} \StringTok{"Expected proportion with 95\% credible intervals | Weights fixed at (1, 1)"}\NormalTok{,}
    \AttributeTok{x        =} \StringTok{"Blue marbles in direct sample (out of 8)"}\NormalTok{,}
    \AttributeTok{y        =} \StringTok{"P(next marble = blue)"}\NormalTok{,}
    \AttributeTok{colour   =} \StringTok{"Social evidence"}\NormalTok{,}
    \AttributeTok{fill     =} \StringTok{"Social evidence"}
\NormalTok{  ) }\SpecialCharTok{+}
  \FunctionTok{theme}\NormalTok{(}\AttributeTok{legend.position =} \StringTok{"bottom"}\NormalTok{)}
\end{Highlighting}
\end{Shaded}

\begin{center}
\includegraphics[width=0.8\linewidth,height=\textheight,keepaspectratio]{10-BayesianModels_files/figure-pdf/ch10_sba_predictions-1.pdf}
\end{center}

Observe that the lines are strictly parallel. Because the total
pseudo-count denominator is fixed, social evidence exerts a constant,
additive shift on the expected probability, regardless of direct
evidence strength.

\subsection{Beta-Binomial
Overdispersion}\label{beta-binomial-overdispersion}

As we have seen above, the BetaBinomial has greater variance than a
Binomial with the same mean. This excess variance is called
overdispersion, let's visualize it.

\begin{Shaded}
\begin{Highlighting}[]
\CommentTok{\# We predict N=10 future draws. This allows the excess variance }
\CommentTok{\# from overdispersion to manifest as positive integers.}
\NormalTok{N\_predict }\OtherTok{\textless{}{-}} \DecValTok{10}

\NormalTok{overdispersion\_data }\OtherTok{\textless{}{-}} \FunctionTok{expand\_grid}\NormalTok{(}
  \AttributeTok{blue1 =} \DecValTok{0}\SpecialCharTok{:}\DecValTok{8}\NormalTok{,}
  \AttributeTok{blue2 =} \DecValTok{0}\SpecialCharTok{:}\DecValTok{3}
\NormalTok{) }\SpecialCharTok{|\textgreater{}}
  \FunctionTok{mutate}\NormalTok{(}
    \AttributeTok{alpha\_p   =} \FloatTok{0.5} \SpecialCharTok{+}\NormalTok{ blue1 }\SpecialCharTok{+}\NormalTok{ blue2,}
    \AttributeTok{beta\_p    =} \FloatTok{0.5} \SpecialCharTok{+}\NormalTok{ (}\DecValTok{8} \SpecialCharTok{{-}}\NormalTok{ blue1) }\SpecialCharTok{+}\NormalTok{ (}\DecValTok{3} \SpecialCharTok{{-}}\NormalTok{ blue2),}
    \AttributeTok{theta\_hat =}\NormalTok{ alpha\_p }\SpecialCharTok{/}\NormalTok{ (alpha\_p }\SpecialCharTok{+}\NormalTok{ beta\_p),}
    
    \CommentTok{\# Expected Binomial variance for N draws (Baseline)}
    \AttributeTok{var\_binom =}\NormalTok{ N\_predict }\SpecialCharTok{*}\NormalTok{ theta\_hat }\SpecialCharTok{*}\NormalTok{ (}\DecValTok{1} \SpecialCharTok{{-}}\NormalTok{ theta\_hat),}
    
    \CommentTok{\# BetaBinomial variance for N draws (Includes belief uncertainty)}
    \CommentTok{\# The multiplier is [1 + (N {-} 1) / (alpha + beta + 1)]}
    \AttributeTok{var\_bb    =}\NormalTok{ var\_binom }\SpecialCharTok{*}\NormalTok{ (}\DecValTok{1} \SpecialCharTok{+}\NormalTok{ (N\_predict }\SpecialCharTok{{-}} \DecValTok{1}\NormalTok{) }\SpecialCharTok{/}\NormalTok{ (alpha\_p }\SpecialCharTok{+}\NormalTok{ beta\_p }\SpecialCharTok{+} \DecValTok{1}\NormalTok{)),}
    
    \CommentTok{\# ABSOLUTE EXCESS VARIANCE: }
    \CommentTok{\# Var(BetaBinomial) {-} Var(Binomial)}
    \CommentTok{\# This yields a positive curve that peaks at maximum ambiguity}
    \AttributeTok{excess\_variance =}\NormalTok{ var\_bb }\SpecialCharTok{{-}}\NormalTok{ var\_binom,}
    
    \AttributeTok{social\_evidence =} \FunctionTok{factor}\NormalTok{(}
\NormalTok{      blue2, }\AttributeTok{levels =} \DecValTok{0}\SpecialCharTok{:}\DecValTok{3}\NormalTok{,}
      \AttributeTok{labels =} \FunctionTok{c}\NormalTok{(}\StringTok{"Clear Red"}\NormalTok{, }\StringTok{"Maybe Red"}\NormalTok{, }\StringTok{"Maybe Blue"}\NormalTok{, }\StringTok{"Clear Blue"}\NormalTok{)}
\NormalTok{    )}
\NormalTok{  )}

\FunctionTok{ggplot}\NormalTok{(overdispersion\_data,}
       \FunctionTok{aes}\NormalTok{(}\AttributeTok{x =}\NormalTok{ blue1, }\AttributeTok{y =}\NormalTok{ excess\_variance,}
           \AttributeTok{colour =}\NormalTok{ social\_evidence, }\AttributeTok{group =}\NormalTok{ social\_evidence)) }\SpecialCharTok{+}
  
  \CommentTok{\# A baseline at 0 to ground the visual}
  \FunctionTok{geom\_hline}\NormalTok{(}\AttributeTok{yintercept =} \DecValTok{0}\NormalTok{, }\AttributeTok{linetype =} \StringTok{"dashed"}\NormalTok{, }\AttributeTok{color =} \StringTok{"gray70"}\NormalTok{) }\SpecialCharTok{+}
  
  \FunctionTok{geom\_line}\NormalTok{(}\AttributeTok{linewidth =} \DecValTok{1}\NormalTok{) }\SpecialCharTok{+}
  \FunctionTok{geom\_point}\NormalTok{(}\AttributeTok{size =} \DecValTok{2}\NormalTok{) }\SpecialCharTok{+}
  \FunctionTok{scale\_colour\_brewer}\NormalTok{(}\AttributeTok{palette =} \StringTok{"Set1"}\NormalTok{) }\SpecialCharTok{+}
  \FunctionTok{scale\_x\_continuous}\NormalTok{(}\AttributeTok{breaks =} \DecValTok{0}\SpecialCharTok{:}\DecValTok{8}\NormalTok{) }\SpecialCharTok{+}
  
  \CommentTok{\# Labels updated to reflect the N=10 prediction framing}
  \FunctionTok{labs}\NormalTok{(}
    \AttributeTok{title    =} \StringTok{"Beta{-}Binomial Overdispersion Across Evidence Combinations"}\NormalTok{,}
    \AttributeTok{subtitle =} \FunctionTok{paste}\NormalTok{(}
      \StringTok{"Excess variance = Var(BetaBinomial) − Var(Binomial) for N=10 predicted draws."}\NormalTok{,}
      \StringTok{"}\SpecialCharTok{\textbackslash{}n}\StringTok{Highest when evidence is ambiguous (near a 50/50 probability split)."}
\NormalTok{    ),}
    \AttributeTok{x        =} \StringTok{"Blue marbles in direct sample (out of 8)"}\NormalTok{,}
    \AttributeTok{y        =} \StringTok{"Excess Variance}\SpecialCharTok{\textbackslash{}n}\StringTok{(Driven by belief uncertainty)"}\NormalTok{,}
    \AttributeTok{colour   =} \StringTok{"Social evidence"}
\NormalTok{  ) }\SpecialCharTok{+}
  \FunctionTok{theme}\NormalTok{(}
    \AttributeTok{legend.position =} \StringTok{"bottom"}\NormalTok{,}
    \AttributeTok{plot.title      =} \FunctionTok{element\_text}\NormalTok{(}\AttributeTok{face =} \StringTok{"bold"}\NormalTok{, }\AttributeTok{size =} \DecValTok{14}\NormalTok{),}
    \AttributeTok{axis.title.y    =} \FunctionTok{element\_text}\NormalTok{(}\AttributeTok{face =} \StringTok{"bold"}\NormalTok{, }\AttributeTok{margin =} \FunctionTok{margin}\NormalTok{(}\AttributeTok{r =} \DecValTok{10}\NormalTok{)),}
    \AttributeTok{axis.title.x    =} \FunctionTok{element\_text}\NormalTok{(}\AttributeTok{face =} \StringTok{"bold"}\NormalTok{, }\AttributeTok{margin =} \FunctionTok{margin}\NormalTok{(}\AttributeTok{t =} \DecValTok{10}\NormalTok{))}
\NormalTok{  )}
\end{Highlighting}
\end{Shaded}

\begin{center}
\includegraphics[width=0.8\linewidth,height=\textheight,keepaspectratio]{10-BayesianModels_files/figure-pdf/ch10_betabinomial_overdispersion-1.pdf}
\end{center}

Overdispersion is greatest when evidence is ambiguous (near a 50/50
probability split, such as 4 direct blue marbles). At this point, the
Beta posterior is widest and \(\theta\) is least constrained, driving
maximum excess variance in the predicted draws. This has a concrete
implication for model checking: if real participants produce choices
that are \emph{more} variable than a Binomial with \(p = \hat{\theta}\)
would predict, that is not a misspecification --- it is exactly what the
BetaBinomial predicts. If choices are \emph{less} variable (more
deterministic than the model allows), the LOO-PIT will show a
hump-shaped histogram, signalling that a threshold decision rule rather
than probabilistic sampling better describes the data.

\subsection{Comparison with Log-Odds
Integration}\label{comparison-with-log-odds-integration}

An alternative framing is to convert evidence to log-odds rather than
pseudo-counts. Here we compare the two approaches:

\begin{Shaded}
\begin{Highlighting}[]
\CommentTok{\# {-}{-}{-} Log{-}Odds Integration Function {-}{-}{-}}
\NormalTok{SimpleBayes\_f }\OtherTok{\textless{}{-}} \ControlFlowTok{function}\NormalTok{(bias, Source1, Source2) \{}
  \CommentTok{\# Using base R equivalents for portability:}
\NormalTok{  logit }\OtherTok{\textless{}{-}} \ControlFlowTok{function}\NormalTok{(p) }\FunctionTok{log}\NormalTok{(p }\SpecialCharTok{/}\NormalTok{ (}\DecValTok{1} \SpecialCharTok{{-}}\NormalTok{ p))}
\NormalTok{  inv\_logit }\OtherTok{\textless{}{-}} \ControlFlowTok{function}\NormalTok{(x) }\DecValTok{1} \SpecialCharTok{/}\NormalTok{ (}\DecValTok{1} \SpecialCharTok{+} \FunctionTok{exp}\NormalTok{(}\SpecialCharTok{{-}}\NormalTok{x))}
  
\NormalTok{  outcome }\OtherTok{\textless{}{-}} \FunctionTok{inv\_logit}\NormalTok{(bias }\SpecialCharTok{+} \FunctionTok{logit}\NormalTok{(Source1) }\SpecialCharTok{+} \FunctionTok{logit}\NormalTok{(Source2))}
  \FunctionTok{return}\NormalTok{(outcome)}
\NormalTok{\}}

\CommentTok{\# {-}{-}{-} Data Preparation {-}{-}{-}}
\CommentTok{\# We\textquotesingle{}ll use a bias of 0 (neutral) and convert counts to proportions for the logit model}
\NormalTok{comparison\_data }\OtherTok{\textless{}{-}} \FunctionTok{expand\_grid}\NormalTok{(}\AttributeTok{blue1 =} \DecValTok{0}\SpecialCharTok{:}\DecValTok{8}\NormalTok{, }\AttributeTok{blue2 =} \DecValTok{0}\SpecialCharTok{:}\DecValTok{3}\NormalTok{) }\SpecialCharTok{|\textgreater{}}
  \FunctionTok{mutate}\NormalTok{(}
    \CommentTok{\# Beta{-}Binomial Expected Rate (Pseudocounts)}
    \AttributeTok{p\_beta =}\NormalTok{ (}\FloatTok{0.5} \SpecialCharTok{+}\NormalTok{ blue1 }\SpecialCharTok{+}\NormalTok{ blue2) }\SpecialCharTok{/}\NormalTok{ (}\FloatTok{0.5} \SpecialCharTok{+} \FloatTok{0.5} \SpecialCharTok{+} \DecValTok{8} \SpecialCharTok{+} \DecValTok{3}\NormalTok{),}
    
    \CommentTok{\# Log{-}Odds Expected Rate (Using the function)}
    \CommentTok{\# We add a tiny epsilon to avoid logit(0) or logit(1)}
    \AttributeTok{prop1 =}\NormalTok{ (}\FloatTok{0.5} \SpecialCharTok{+}\NormalTok{ blue1) }\SpecialCharTok{/}\NormalTok{ (}\FloatTok{0.5} \SpecialCharTok{+} \FloatTok{0.5} \SpecialCharTok{+} \DecValTok{8}\NormalTok{),}
    \AttributeTok{prop2 =}\NormalTok{ (}\FloatTok{0.5} \SpecialCharTok{+}\NormalTok{ blue2) }\SpecialCharTok{/}\NormalTok{ (}\FloatTok{0.5} \SpecialCharTok{+} \FloatTok{0.5} \SpecialCharTok{+} \DecValTok{3}\NormalTok{),}
    \AttributeTok{p\_logit =} \FunctionTok{map2\_dbl}\NormalTok{(prop1, prop2, }\SpecialCharTok{\textasciitilde{}}\FunctionTok{SimpleBayes\_f}\NormalTok{(}\DecValTok{0}\NormalTok{, .x, .y)),}
    
    \AttributeTok{social\_evidence =} \FunctionTok{factor}\NormalTok{(blue2, }\AttributeTok{levels =} \DecValTok{0}\SpecialCharTok{:}\DecValTok{3}\NormalTok{, }
                             \AttributeTok{labels =} \FunctionTok{c}\NormalTok{(}\StringTok{"Clear Red"}\NormalTok{, }\StringTok{"Maybe Red"}\NormalTok{, }\StringTok{"Maybe Blue"}\NormalTok{, }\StringTok{"Clear Blue"}\NormalTok{))}
\NormalTok{  ) }\SpecialCharTok{|\textgreater{}}
  \FunctionTok{pivot\_longer}\NormalTok{(}\AttributeTok{cols =} \FunctionTok{c}\NormalTok{(p\_beta, p\_logit), }\AttributeTok{names\_to =} \StringTok{"Model"}\NormalTok{, }\AttributeTok{values\_to =} \StringTok{"p\_blue"}\NormalTok{)}

\CommentTok{\# {-}{-}{-} Plotting {-}{-}{-}}
\FunctionTok{ggplot}\NormalTok{(comparison\_data, }\FunctionTok{aes}\NormalTok{(}\AttributeTok{x =}\NormalTok{ blue1, }\AttributeTok{y =}\NormalTok{ p\_blue, }\AttributeTok{color =}\NormalTok{ social\_evidence)) }\SpecialCharTok{+}
  \FunctionTok{geom\_line}\NormalTok{(}\AttributeTok{linewidth =} \DecValTok{1}\NormalTok{) }\SpecialCharTok{+}
  \FunctionTok{geom\_point}\NormalTok{(}\AttributeTok{size =} \DecValTok{2}\NormalTok{) }\SpecialCharTok{+}
  \FunctionTok{facet\_wrap}\NormalTok{(}\SpecialCharTok{\textasciitilde{}}\NormalTok{Model, }\AttributeTok{labeller =} \FunctionTok{as\_labeller}\NormalTok{(}\FunctionTok{c}\NormalTok{(}
    \StringTok{"p\_beta"}  \OtherTok{=} \StringTok{"Beta{-}Binomial (Counts/Accumulation)"}\NormalTok{,}
    \StringTok{"p\_logit"} \OtherTok{=} \StringTok{"Log{-}Odds (Linear Summation)"}
\NormalTok{  ))) }\SpecialCharTok{+}
  \FunctionTok{scale\_color\_manual}\NormalTok{(}\AttributeTok{values =}\NormalTok{ semantic\_colors) }\SpecialCharTok{+}
  \FunctionTok{scale\_x\_continuous}\NormalTok{(}\AttributeTok{breaks =} \DecValTok{0}\SpecialCharTok{:}\DecValTok{8}\NormalTok{) }\SpecialCharTok{+}
  \FunctionTok{labs}\NormalTok{(}
    \AttributeTok{title =} \StringTok{"Comparison of Evidence Integration Scales"}\NormalTok{,}
    \AttributeTok{subtitle =} \StringTok{"Log{-}odds models (right) produce steeper \textquotesingle{}S{-}curves\textquotesingle{} than count{-}based models (left)."}\NormalTok{,}
    \AttributeTok{x =} \StringTok{"Blue marbles in direct sample (out of 8)"}\NormalTok{,}
    \AttributeTok{y =} \StringTok{"P(choose blue)"}\NormalTok{,}
    \AttributeTok{color =} \StringTok{"Social Evidence"}
\NormalTok{  ) }\SpecialCharTok{+}
  \FunctionTok{theme\_classic}\NormalTok{() }\SpecialCharTok{+}
  \FunctionTok{theme}\NormalTok{(}\AttributeTok{legend.position =} \StringTok{"bottom"}\NormalTok{)}
\end{Highlighting}
\end{Shaded}

\begin{center}
\includegraphics[width=0.8\linewidth,height=\textheight,keepaspectratio]{10-BayesianModels_files/figure-pdf/ch10_logodds_comparison-1.pdf}
\end{center}

The key difference: log-odds integration produces steeper ``S-curves''
--- the model is more confident near the extremes. Count-based
(Beta-Binomial) integration is more conservative, maintaining higher
uncertainty even with unbalanced evidence.

\section{Two Generative Models for Evidence
Integration}\label{two-generative-models-for-evidence-integration}

Now that we understand the Beta-Binomial better, we can use it to build
two different agents integrating direct and social evidence. The agents
differ in exactly one respect: whether the source weights are fixed at
the ``rational'' levels (taking the sources at face value) or they can
vary (e.g.~weighing direct experience more than social experience).

\subsection{The Simple Bayesian Agent
(SBA)}\label{the-simple-bayesian-agent-sba}

The SBA fixes the weights at \(w_d = w_s = 1\) --- every observation
counts at face value, regardless of its source (but remember that social
evidence in our design counts for less as it can have at most 3 balls of
a given color). With Jeffreys prior pseudo-counts:

\[\alpha_{\text{post}} = 0.5 + k_1 + k_2\]
\[\beta_{\text{post}}  = 0.5 + (n_1 - k_1) + (n_2 - k_2)\]

The SBA has \textbf{zero free parameters}. It is a fully deterministic
belief-update rule --- there is nothing to estimate from the data. We
can still fit it (in the sense of running it through Stan to get
posterior predictives and log-likelihoods) and check whether it predicts
behaviour.

\subsection{The Weighted Bayesian Agent
(WBA)}\label{the-weighted-bayesian-agent-wba}

But what if people don't weight both sources equally? Some individuals
might trust their direct experience more; others might be more socially
influenced. The WBA promotes the weights to free parameters:

\[\alpha_{\text{post}} = 0.5 + w_d \cdot k_1 + w_s \cdot k_2\]
\[\beta_{\text{post}}  = 0.5 + w_d \cdot (n_1 - k_1) + w_s \cdot (n_2 - k_2)\]

The WBA has \textbf{two free parameters}: \(w_d\) and \(w_s\). Note that
the SBA is nested inside the WBA at exactly \(w_d = w_s = 1\) ---
anything the SBA can predict, the WBA can predict too. The empirical
question is whether the extra flexibility is needed \emph{and} whether
the data can pin down where in \((w_d, w_s)\) space each participant
lives.

\subsection{Three Scenarios that Distinguish the
Models}\label{three-scenarios-that-distinguish-the-models}

Let's design three datasets --- each generated from a different agent
--- that we'll use throughout the next few sections. We will fit
\emph{both} the SBA and the WBA to all three datasets and see which
model captures which:

\begin{enumerate}
\def\labelenumi{\arabic{enumi}.}
\item
  \textbf{Self-Focused:} \(w_d = 1.5\), \(w_s = 0.5\) (trusts own
  observations more)
\item
  \textbf{Balanced:} \(w_d = 1.0\), \(w_s = 1.0\) (the SBA's generating
  process)
\item
  \textbf{Socially-Influenced:} \(w_d = 0.5\), \(w_s = 2.0\) (strongly
  socially influenced)
\end{enumerate}

In the Balanced dataset SBA and WBA coincide. Or in other words, only
the Balanced dataset is generable by SBA, while WBA can generate all
three scenarios. We expect the SBA to fit the Balanced scenario and
\emph{misfit} the other two systematically --- a basic model fitting
test. The WBA, in principle, can fit all three.

\begin{Shaded}
\begin{Highlighting}[]
\NormalTok{agents }\OtherTok{\textless{}{-}} \FunctionTok{tibble}\NormalTok{(}
  \AttributeTok{agent\_type    =} \FunctionTok{c}\NormalTok{(}\StringTok{"Self{-}Focused"}\NormalTok{, }\StringTok{"Balanced"}\NormalTok{, }\StringTok{"Socially{-}Influenced"}\NormalTok{),}
  \AttributeTok{weight\_direct =} \FunctionTok{c}\NormalTok{(}\FloatTok{1.5}\NormalTok{, }\FloatTok{1.0}\NormalTok{, }\FloatTok{0.5}\NormalTok{),}
  \AttributeTok{weight\_social =} \FunctionTok{c}\NormalTok{(}\FloatTok{0.5}\NormalTok{, }\FloatTok{1.0}\NormalTok{, }\FloatTok{2.0}\NormalTok{)}
\NormalTok{)}

\CommentTok{\# Generate simulated data for all three agent types}
\NormalTok{simulation\_results }\OtherTok{\textless{}{-}}\NormalTok{ agents }\SpecialCharTok{|\textgreater{}}
  \FunctionTok{pmap\_dfr}\NormalTok{(}\ControlFlowTok{function}\NormalTok{(agent\_type, weight\_direct, weight\_social) \{}
    \FunctionTok{generate\_agent\_decisions}\NormalTok{(weight\_direct, weight\_social,}
\NormalTok{                             evidence\_combinations, }\AttributeTok{n\_samples =} \DecValTok{5}\NormalTok{) }\SpecialCharTok{|\textgreater{}}
      \FunctionTok{mutate}\NormalTok{(}\AttributeTok{agent\_type =}\NormalTok{ agent\_type)}
\NormalTok{  \}) }\SpecialCharTok{|\textgreater{}}
  \FunctionTok{mutate}\NormalTok{(}
    \AttributeTok{social\_evidence =} \FunctionTok{factor}\NormalTok{(}
\NormalTok{      blue2, }\AttributeTok{levels =} \DecValTok{0}\SpecialCharTok{:}\DecValTok{3}\NormalTok{,}
      \AttributeTok{labels =} \FunctionTok{c}\NormalTok{(}\StringTok{"Clear Red"}\NormalTok{, }\StringTok{"Maybe Red"}\NormalTok{, }\StringTok{"Maybe Blue"}\NormalTok{, }\StringTok{"Clear Blue"}\NormalTok{)}
\NormalTok{    ),}
    \AttributeTok{agent\_type =} \FunctionTok{factor}\NormalTok{(}
\NormalTok{      agent\_type,}
      \AttributeTok{levels =} \FunctionTok{c}\NormalTok{(}\StringTok{"Self{-}Focused"}\NormalTok{, }\StringTok{"Balanced"}\NormalTok{, }\StringTok{"Socially{-}Influenced"}\NormalTok{)}
\NormalTok{    )}
\NormalTok{  )}

\CommentTok{\# Summarize for plotting}
\NormalTok{plot\_summary }\OtherTok{\textless{}{-}}\NormalTok{ simulation\_results }\SpecialCharTok{|\textgreater{}}
  \FunctionTok{group\_by}\NormalTok{(agent\_type, blue1, social\_evidence) }\SpecialCharTok{|\textgreater{}}
  \FunctionTok{summarise}\NormalTok{(}
    \AttributeTok{theoretical\_p =} \FunctionTok{mean}\NormalTok{(expected\_rate),}
    \AttributeTok{empirical\_p   =} \FunctionTok{mean}\NormalTok{(choice),}
    \AttributeTok{.groups =} \StringTok{"drop"}
\NormalTok{  )}

\CommentTok{\# Visualize}
\FunctionTok{ggplot}\NormalTok{(plot\_summary, }\FunctionTok{aes}\NormalTok{(}\AttributeTok{x =}\NormalTok{ blue1, }\AttributeTok{colour =}\NormalTok{ social\_evidence, }\AttributeTok{group =}\NormalTok{ social\_evidence)) }\SpecialCharTok{+}
  \FunctionTok{geom\_line}\NormalTok{(}\FunctionTok{aes}\NormalTok{(}\AttributeTok{y =}\NormalTok{ theoretical\_p), }\AttributeTok{linewidth =} \FloatTok{1.2}\NormalTok{, }\AttributeTok{alpha =} \FloatTok{0.5}\NormalTok{) }\SpecialCharTok{+}
  \FunctionTok{geom\_line}\NormalTok{(}\FunctionTok{aes}\NormalTok{(}\AttributeTok{y =}\NormalTok{ empirical\_p), }\AttributeTok{linetype =} \StringTok{"dashed"}\NormalTok{, }\AttributeTok{linewidth =} \FloatTok{0.8}\NormalTok{) }\SpecialCharTok{+}
  \FunctionTok{geom\_point}\NormalTok{(}\FunctionTok{aes}\NormalTok{(}\AttributeTok{y =}\NormalTok{ empirical\_p), }\AttributeTok{size =} \DecValTok{2}\NormalTok{) }\SpecialCharTok{+}
  \FunctionTok{geom\_hline}\NormalTok{(}\AttributeTok{yintercept =} \FloatTok{0.5}\NormalTok{, }\AttributeTok{linetype =} \StringTok{"dotted"}\NormalTok{, }\AttributeTok{colour =} \StringTok{"grey30"}\NormalTok{) }\SpecialCharTok{+}
  \FunctionTok{facet\_wrap}\NormalTok{(}\SpecialCharTok{\textasciitilde{}}\NormalTok{ agent\_type) }\SpecialCharTok{+}
  \FunctionTok{scale\_colour\_manual}\NormalTok{(}\AttributeTok{values =}\NormalTok{ semantic\_colors) }\SpecialCharTok{+}
  \FunctionTok{scale\_x\_continuous}\NormalTok{(}\AttributeTok{breaks =} \DecValTok{0}\SpecialCharTok{:}\DecValTok{8}\NormalTok{) }\SpecialCharTok{+}
  \FunctionTok{coord\_cartesian}\NormalTok{(}\AttributeTok{ylim =} \FunctionTok{c}\NormalTok{(}\DecValTok{0}\NormalTok{, }\DecValTok{1}\NormalTok{)) }\SpecialCharTok{+}
  \FunctionTok{labs}\NormalTok{(}
    \AttributeTok{title =} \StringTok{"Weighted Bayesian Agents: Three Scenarios"}\NormalTok{,}
    \AttributeTok{subtitle =} \StringTok{"Solid faded = Internal Belief | Dashed = Observed Choice Proportions (N=5 per condition)"}\NormalTok{,}
    \AttributeTok{x =} \StringTok{"Blue marbles in direct sample (out of 8)"}\NormalTok{,}
    \AttributeTok{y =} \StringTok{"P(choose blue)"}\NormalTok{,}
    \AttributeTok{colour =} \StringTok{"Social evidence"}
\NormalTok{  ) }\SpecialCharTok{+}
  \FunctionTok{theme}\NormalTok{(}\AttributeTok{legend.position =} \StringTok{"bottom"}\NormalTok{,}
        \AttributeTok{strip.background =} \FunctionTok{element\_blank}\NormalTok{(),}
        \AttributeTok{strip.text =} \FunctionTok{element\_text}\NormalTok{(}\AttributeTok{face =} \StringTok{"bold"}\NormalTok{, }\AttributeTok{size =} \DecValTok{11}\NormalTok{))}
\end{Highlighting}
\end{Shaded}

\begin{center}
\includegraphics[width=0.8\linewidth,height=\textheight,keepaspectratio]{10-BayesianModels_files/figure-pdf/ch10_three_scenarios-1.pdf}
\end{center}

Notice how the scenarios differ:

\begin{itemize}
\item
  \textbf{Self-Focused:} Lines are spread apart by direct evidence;
  social evidence has less effect
\item
  \textbf{Balanced:} Equal spacing from both sources
\item
  \textbf{Socially-Influenced:} Lines are bunched together; social
  evidence dominates.
\end{itemize}

The conceptual setup is now in place: we have two nested models and
three datasets, only one of which the SBA can possibly capture. Round 1
puts both models in Stan and finds out what happens.

\section{Round 1: Fitting SBA and
WBA}\label{round-1-fitting-sba-and-wba}

\subsection{Stan code: SBA}\label{stan-code-sba}

\begin{Shaded}
\begin{Highlighting}[]
\FunctionTok{\textbackslash{}usetikzlibrary}\NormalTok{\{bayesnet\}}
\KeywordTok{\textbackslash{}begin}\NormalTok{\{}\ExtensionTok{tikzpicture}\NormalTok{\}}
  \CommentTok{\% Nodes}
  \FunctionTok{\textbackslash{}node}\NormalTok{[const] (alpha0) \{}\SpecialStringTok{$}\SpecialCharTok{\textbackslash{}alpha}\SpecialStringTok{\_0$}\NormalTok{\};}
  \FunctionTok{\textbackslash{}node}\NormalTok{[const, right=1cm of alpha0] (beta0) \{}\SpecialStringTok{$}\SpecialCharTok{\textbackslash{}beta}\SpecialStringTok{\_0$}\NormalTok{\};}
  
  \FunctionTok{\textbackslash{}node}\NormalTok{[obs, below left=1.5cm and 0.5cm of alpha0] (blue1) \{}\SpecialStringTok{$k\_\{1,i\}$}\NormalTok{\};}
  \FunctionTok{\textbackslash{}node}\NormalTok{[obs, right=0.5cm of blue1] (total1) \{}\SpecialStringTok{$n\_\{1,i\}$}\NormalTok{\};}
  
  \FunctionTok{\textbackslash{}node}\NormalTok{[obs, below right=1.5cm and 0.5cm of beta0] (blue2) \{}\SpecialStringTok{$k\_\{2,i\}$}\NormalTok{\};}
  \FunctionTok{\textbackslash{}node}\NormalTok{[obs, right=0.5cm of blue2] (total2) \{}\SpecialStringTok{$n\_\{2,i\}$}\NormalTok{\};}

  \FunctionTok{\textbackslash{}node}\NormalTok{[det, below=3.5cm of alpha0] (alphapost) \{}\SpecialStringTok{$}\SpecialCharTok{\textbackslash{}alpha}\SpecialStringTok{\_\{post,i\}$}\NormalTok{\};}
  \FunctionTok{\textbackslash{}node}\NormalTok{[det, below=3.5cm of beta0] (betapost) \{}\SpecialStringTok{$}\SpecialCharTok{\textbackslash{}beta}\SpecialStringTok{\_\{post,i\}$}\NormalTok{\};}

  \FunctionTok{\textbackslash{}node}\NormalTok{[obs, below=1.5cm of alphapost, xshift=0.75cm] (choice) \{}\SpecialStringTok{$y\_i$}\NormalTok{\};}

  \CommentTok{\% Edges}
  \FunctionTok{\textbackslash{}edge}\NormalTok{ \{alpha0\} \{alphapost\};}
  \FunctionTok{\textbackslash{}edge}\NormalTok{ \{beta0\} \{betapost\};}
  \FunctionTok{\textbackslash{}edge}\NormalTok{ \{blue1, total1, blue2, total2\} \{alphapost, betapost\};}
  \FunctionTok{\textbackslash{}edge}\NormalTok{ \{alphapost, betapost\} \{choice\};}

  \CommentTok{\% Plate}
  \FunctionTok{\textbackslash{}plate}\NormalTok{ \{p1\} \{(blue1)(total1)(blue2)(total2)(alphapost)(betapost)(choice)\} \{}\SpecialStringTok{$N$}\NormalTok{\};}
\KeywordTok{\textbackslash{}end}\NormalTok{\{}\ExtensionTok{tikzpicture}\NormalTok{\}}
\end{Highlighting}
\end{Shaded}

\begin{figure}[H]

{\centering \includegraphics[width=0.6\linewidth,height=\textheight,keepaspectratio]{10-BayesianModels_files/figure-pdf/sba-dag-1.pdf}

}

\caption{DAG for the Simple Bayesian Agent}

\end{figure}%

The generative relationships for the SBA are: \[
\begin{aligned}
\alpha_{post,i} &= \alpha_0 + k_{1,i} + k_{2,i} \\
\beta_{post,i}  &= \beta_0 + (n_{1,i} - k_{1,i}) + (n_{2,i} - k_{2,i}) \\
y_i &\sim \text{Beta-Binomial}(1, \alpha_{post,i}, \beta_{post,i})
\end{aligned}
\]

\begin{Shaded}
\begin{Highlighting}[]
\NormalTok{SimpleAgent\_stan }\OtherTok{\textless{}{-}} \StringTok{"}
\StringTok{// Simple Bayesian Agent (SBA).}
\StringTok{// No free parameters — evidence is counted at face value.}
\StringTok{// Jeffreys prior pseudo{-}counts: alpha0 = beta0 = 0.5.}
\StringTok{data \{}
\StringTok{  int\textless{}lower=1\textgreater{} N\_total;}
\StringTok{  array[N\_total] int\textless{}lower=0, upper=1\textgreater{} y;}
\StringTok{  array[N\_total] int\textless{}lower=0\textgreater{} blue1;}
\StringTok{  array[N\_total] int\textless{}lower=0\textgreater{} total1;}
\StringTok{  array[N\_total] int\textless{}lower=0\textgreater{} blue2;}
\StringTok{  array[N\_total] int\textless{}lower=0\textgreater{} total2;}
\StringTok{  int\textless{}lower=0, upper=1\textgreater{} run\_diagnostics;}
\StringTok{\}}

\StringTok{model \{}
\StringTok{  // Vectorized likelihood with fixed weights = 1}
\StringTok{  vector[N\_total] alpha\_post = 0.5 + to\_vector(blue1) + to\_vector(blue2);}
\StringTok{  vector[N\_total] beta\_post  = 0.5 + (to\_vector(total1) {-} to\_vector(blue1))}
\StringTok{                                   + (to\_vector(total2) {-} to\_vector(blue2));}
\StringTok{                             }
\StringTok{  target += beta\_binomial\_lpmf(y | 1, alpha\_post, beta\_post);}
\StringTok{\}}

\StringTok{generated quantities \{}
\StringTok{  vector[N\_total] log\_lik;}
\StringTok{  array[N\_total] int y\_rep;}

\StringTok{  if (run\_diagnostics) \{}
\StringTok{    for (i in 1:N\_total) \{}
\StringTok{      real alpha\_post = 0.5 + blue1[i] + blue2[i];}
\StringTok{      real beta\_post  = 0.5 + (total1[i] {-} blue1[i]) + (total2[i] {-} blue2[i]);}

\StringTok{      log\_lik[i] = beta\_binomial\_lpmf(y[i] | 1, alpha\_post, beta\_post);}
\StringTok{      y\_rep[i]   = beta\_binomial\_rng(1, alpha\_post, beta\_post);}
\StringTok{    \}}
\StringTok{  \}}
\StringTok{\}}
\StringTok{"}

\FunctionTok{write\_stan\_file}\NormalTok{(SimpleAgent\_stan,}
                \AttributeTok{dir      =} \StringTok{"stan/"}\NormalTok{,}
                \AttributeTok{basename =} \StringTok{"ch10\_simple\_agent.stan"}\NormalTok{)}
\end{Highlighting}
\end{Shaded}

\begin{verbatim}
[1] "/Users/au209589/Dropbox/Teaching/AdvancedCognitiveModeling23_book/stan/ch10_simple_agent.stan"
\end{verbatim}

\subsection{Stan code: WBA (raw
parameterisation)}\label{stan-code-wba-raw-parameterisation}

We start with the intuitive parameterisation using \(w_d\) and \(w_s\)
directly:

\begin{Shaded}
\begin{Highlighting}[]
\FunctionTok{\textbackslash{}usetikzlibrary}\NormalTok{\{bayesnet\}}
\KeywordTok{\textbackslash{}begin}\NormalTok{\{}\ExtensionTok{tikzpicture}\NormalTok{\}}
  \CommentTok{\% Fixed Prior}
  \FunctionTok{\textbackslash{}node}\NormalTok{[const] (alpha0) \{}\SpecialStringTok{$}\SpecialCharTok{\textbackslash{}alpha}\SpecialStringTok{\_0$}\NormalTok{\};}
  \FunctionTok{\textbackslash{}node}\NormalTok{[const, right=1.5cm of alpha0] (beta0) \{}\SpecialStringTok{$}\SpecialCharTok{\textbackslash{}beta}\SpecialStringTok{\_0$}\NormalTok{\};}

  \CommentTok{\% Latent Parameters (Direct Weights)}
  \FunctionTok{\textbackslash{}node}\NormalTok{[latent, below=1.2cm of alpha0] (wd) \{}\SpecialStringTok{$w\_d$}\NormalTok{\}; }
  \FunctionTok{\textbackslash{}node}\NormalTok{[latent, below=1.2cm of beta0] (ws) \{}\SpecialStringTok{$w\_s$}\NormalTok{\};}

  \CommentTok{\% Inputs}
  \FunctionTok{\textbackslash{}node}\NormalTok{[obs, left=1cm of wd] (k1) \{}\SpecialStringTok{$k\_\{1,i\}$}\NormalTok{\};}
  \FunctionTok{\textbackslash{}node}\NormalTok{[obs, left=0.3cm of k1] (n1) \{}\SpecialStringTok{$n\_\{1,i\}$}\NormalTok{\};}
  \FunctionTok{\textbackslash{}node}\NormalTok{[obs, right=1cm of ws] (k2) \{}\SpecialStringTok{$k\_\{2,i\}$}\NormalTok{\};}
  \FunctionTok{\textbackslash{}node}\NormalTok{[obs, right=0.3cm of k2] (n2) \{}\SpecialStringTok{$n\_\{2,i\}$}\NormalTok{\};}

  \CommentTok{\% Beliefs}
  \FunctionTok{\textbackslash{}node}\NormalTok{[det, below=2.5cm of wd] (ap) \{}\SpecialStringTok{$}\SpecialCharTok{\textbackslash{}alpha}\SpecialStringTok{\_\{post,i\}$}\NormalTok{\};}
  \FunctionTok{\textbackslash{}node}\NormalTok{[det, below=2.5cm of ws] (bp) \{}\SpecialStringTok{$}\SpecialCharTok{\textbackslash{}beta}\SpecialStringTok{\_\{post,i\}$}\NormalTok{\};}

  \CommentTok{\% Action}
  \FunctionTok{\textbackslash{}node}\NormalTok{[obs, below=1.2cm of ap, xshift=1.25cm] (y) \{}\SpecialStringTok{$y\_i$}\NormalTok{\};}

  \CommentTok{\% Edges}
  \FunctionTok{\textbackslash{}edge}\NormalTok{ \{alpha0, beta0\} \{ap, bp\};}
  \FunctionTok{\textbackslash{}edge}\NormalTok{ \{wd, k1, n1, ws, k2, n2\} \{ap, bp\};}
  \FunctionTok{\textbackslash{}edge}\NormalTok{ \{ap, bp\} \{y\};}

  \FunctionTok{\textbackslash{}plate}\NormalTok{ \{p1\} \{(k1)(n1)(k2)(n2)(ap)(bp)(y)\} \{}\SpecialStringTok{$N$}\NormalTok{\};}
\KeywordTok{\textbackslash{}end}\NormalTok{\{}\ExtensionTok{tikzpicture}\NormalTok{\}}
\end{Highlighting}
\end{Shaded}

\begin{figure}[H]

{\centering \includegraphics[width=0.6\linewidth,height=\textheight,keepaspectratio]{10-BayesianModels_files/figure-pdf/wba0-dag-1.pdf}

}

\caption{DAG for the unreparameterized Weighted Bayesian Agent}

\end{figure}%

The generative relationships for the unreparameterized WBA are: \[
\begin{aligned}
w_d &\sim \text{Lognormal}(0, 0.5) \\
w_s &\sim \text{Lognormal}(0, 0.5) \\
\alpha_{post,i} &= \alpha_0 + w_d k_{1,i} + w_s k_{2,i} \\
\beta_{post,i}  &= \beta_0 + w_d(n_{1,i} - k_{1,i}) + w_s(n_{2,i} - k_{2,i}) \\
y_i &\sim \text{Beta-Binomial}(1, \alpha_{post,i}, \beta_{post,i})
\end{aligned}
\]

\begin{Shaded}
\begin{Highlighting}[]
\NormalTok{WeightedAgentRaw\_stan }\OtherTok{\textless{}{-}} \StringTok{"}
\StringTok{// Weighted Bayesian Agent (raw parameterisation).}
\StringTok{// weight\_direct, weight\_social \textgreater{} 0: independent weights on direct and social evidence.}
\StringTok{// Jeffreys prior pseudo{-}counts (0.5) consistent with SBA.}
\StringTok{data \{}
\StringTok{  int\textless{}lower=1\textgreater{} N\_total;}
\StringTok{  array[N\_total] int\textless{}lower=0, upper=1\textgreater{} y;}
\StringTok{  array[N\_total] int\textless{}lower=0\textgreater{} blue1;}
\StringTok{  array[N\_total] int\textless{}lower=0\textgreater{} total1;}
\StringTok{  array[N\_total] int\textless{}lower=0\textgreater{} blue2;}
\StringTok{  array[N\_total] int\textless{}lower=0\textgreater{} total2;}
\StringTok{  }
\StringTok{  real prior\_wd\_mu;}
\StringTok{  real\textless{}lower=0\textgreater{} prior\_wd\_sigma;}
\StringTok{  real prior\_ws\_mu;}
\StringTok{  real\textless{}lower=0\textgreater{} prior\_ws\_sigma;}
\StringTok{  int\textless{}lower=0, upper=1\textgreater{} run\_diagnostics;}
\StringTok{\}}

\StringTok{parameters \{}
\StringTok{  real\textless{}lower=0\textgreater{} weight\_direct;  // w\_d}
\StringTok{  real\textless{}lower=0\textgreater{} weight\_social;  // w\_s}
\StringTok{\}}

\StringTok{model \{}
\StringTok{  // Priors: lognormal centered on 1 (evidence taken at face value)}
\StringTok{  target += lognormal\_lpdf(weight\_direct | prior\_wd\_mu, prior\_wd\_sigma);}
\StringTok{  target += lognormal\_lpdf(weight\_social | prior\_ws\_mu, prior\_ws\_sigma);}

\StringTok{  // Vectorized likelihood}
\StringTok{  vector[N\_total] alpha\_post = 0.5 + weight\_direct * to\_vector(blue1)}
\StringTok{                                   + weight\_social * to\_vector(blue2);}
\StringTok{  vector[N\_total] beta\_post  = 0.5 + weight\_direct * (to\_vector(total1) {-} to\_vector(blue1))}
\StringTok{                                   + weight\_social * (to\_vector(total2) {-} to\_vector(blue2));}
\StringTok{                             }
\StringTok{  target += beta\_binomial\_lpmf(y | 1, alpha\_post, beta\_post);}
\StringTok{\}}

\StringTok{generated quantities \{}
\StringTok{  vector[N\_total] log\_lik;}
\StringTok{  array[N\_total] int y\_rep;}

\StringTok{  if (run\_diagnostics) \{}
\StringTok{    for (i in 1:N\_total) \{}
\StringTok{      real alpha\_post = 0.5 + weight\_direct * blue1[i] + weight\_social * blue2[i];}
\StringTok{      real beta\_post  = 0.5 + weight\_direct * (total1[i] {-} blue1[i]) }
\StringTok{                           + weight\_social * (total2[i] {-} blue2[i]);}

\StringTok{      log\_lik[i] = beta\_binomial\_lpmf(y[i] | 1, alpha\_post, beta\_post);}
\StringTok{      y\_rep[i]   = beta\_binomial\_rng(1, alpha\_post, beta\_post);}
\StringTok{    \}}
\StringTok{  \}}
\StringTok{\}}
\StringTok{"}

\FunctionTok{write\_stan\_file}\NormalTok{(WeightedAgentRaw\_stan,}
                \AttributeTok{dir      =} \StringTok{"stan/"}\NormalTok{,}
                \AttributeTok{basename =} \StringTok{"ch10\_weighted\_agent\_raw.stan"}\NormalTok{)}
\end{Highlighting}
\end{Shaded}

\begin{verbatim}
[1] "/Users/au209589/Dropbox/Teaching/AdvancedCognitiveModeling23_book/stan/ch10_weighted_agent_raw.stan"
\end{verbatim}

\subsection{Fitting the three
scenarios}\label{fitting-the-three-scenarios}

\begin{Shaded}
\begin{Highlighting}[]
\NormalTok{mod\_simple       }\OtherTok{\textless{}{-}} \FunctionTok{cmdstan\_model}\NormalTok{(}\StringTok{"stan/ch10\_simple\_agent.stan"}\NormalTok{, }\AttributeTok{dir =} \StringTok{"simmodels"}\NormalTok{)}
\NormalTok{mod\_weighted\_raw }\OtherTok{\textless{}{-}} \FunctionTok{cmdstan\_model}\NormalTok{(}\StringTok{"stan/ch10\_weighted\_agent\_raw.stan"}\NormalTok{, }\AttributeTok{dir =} \StringTok{"simmodels"}\NormalTok{)}

\NormalTok{sampler\_args }\OtherTok{\textless{}{-}} \FunctionTok{list}\NormalTok{(}
  \AttributeTok{chains =} \DecValTok{4}\NormalTok{, }
  \AttributeTok{parallel\_chains =} \DecValTok{4}\NormalTok{, }
  \AttributeTok{iter\_warmup =} \DecValTok{1000}\NormalTok{, }
  \AttributeTok{iter\_sampling =} \DecValTok{1000}\NormalTok{, }
  \AttributeTok{refresh =} \DecValTok{0}
\NormalTok{)}
\end{Highlighting}
\end{Shaded}

\subsubsection{Preparing data}\label{preparing-data}

\begin{Shaded}
\begin{Highlighting}[]
\CommentTok{\# True values for the three WBA scenarios}
\NormalTok{true\_wba }\OtherTok{\textless{}{-}} \FunctionTok{list}\NormalTok{(}
  \AttributeTok{Self\_Focused        =} \FunctionTok{list}\NormalTok{(}\AttributeTok{wd =} \FloatTok{1.5}\NormalTok{, }\AttributeTok{ws =} \FloatTok{0.5}\NormalTok{),}
  \AttributeTok{Balanced            =} \FunctionTok{list}\NormalTok{(}\AttributeTok{wd =} \FloatTok{1.0}\NormalTok{, }\AttributeTok{ws =} \FloatTok{1.0}\NormalTok{),}
  \AttributeTok{Socially\_Influenced =} \FunctionTok{list}\NormalTok{(}\AttributeTok{wd =} \FloatTok{0.5}\NormalTok{, }\AttributeTok{ws =} \FloatTok{2.0}\NormalTok{)}
\NormalTok{)}

\CommentTok{\# Generate simulated data for each scenario}
\NormalTok{stan\_data\_wba }\OtherTok{\textless{}{-}} \FunctionTok{map}\NormalTok{(true\_wba, }\ControlFlowTok{function}\NormalTok{(params) \{}
\NormalTok{  sim }\OtherTok{\textless{}{-}} \FunctionTok{generate\_agent\_decisions}\NormalTok{(}
\NormalTok{    params}\SpecialCharTok{$}\NormalTok{wd, params}\SpecialCharTok{$}\NormalTok{ws,}
\NormalTok{    evidence\_combinations,}
    \AttributeTok{n\_samples =} \DecValTok{5}
\NormalTok{  ) }\SpecialCharTok{|\textgreater{}} \FunctionTok{mutate}\NormalTok{(}\AttributeTok{total1 =} \DecValTok{8}\NormalTok{L, }\AttributeTok{total2 =} \DecValTok{3}\NormalTok{L)}

\NormalTok{  d }\OtherTok{\textless{}{-}} \FunctionTok{prepare\_stan\_data}\NormalTok{(sim)}
  \CommentTok{\# Add priors for the weighted model}
\NormalTok{  d}\SpecialCharTok{$}\NormalTok{prior\_wd\_mu    }\OtherTok{=} \DecValTok{0}
\NormalTok{  d}\SpecialCharTok{$}\NormalTok{prior\_wd\_sigma }\OtherTok{=} \FloatTok{0.5}
\NormalTok{  d}\SpecialCharTok{$}\NormalTok{prior\_ws\_mu    }\OtherTok{=} \DecValTok{0}
\NormalTok{  d}\SpecialCharTok{$}\NormalTok{prior\_ws\_sigma }\OtherTok{=} \FloatTok{0.5}
\NormalTok{  d}
\NormalTok{\})}
\end{Highlighting}
\end{Shaded}

\subsection{Fitting the Simple Model}\label{fitting-the-simple-model}

\begin{Shaded}
\begin{Highlighting}[]
\NormalTok{fit\_sba\_paths }\OtherTok{\textless{}{-}}\NormalTok{ here}\SpecialCharTok{::}\FunctionTok{here}\NormalTok{(}
  \StringTok{"simmodels"}\NormalTok{,}
  \FunctionTok{paste0}\NormalTok{(}\StringTok{"ch10\_fit\_sba\_"}\NormalTok{, }\FunctionTok{names}\NormalTok{(stan\_data\_wba), }\StringTok{".rds"}\NormalTok{)}
\NormalTok{) }\SpecialCharTok{|\textgreater{}} \FunctionTok{set\_names}\NormalTok{(}\FunctionTok{names}\NormalTok{(stan\_data\_wba))}

\ControlFlowTok{if}\NormalTok{ (regenerate\_simulations }\SpecialCharTok{||} \SpecialCharTok{!}\FunctionTok{all}\NormalTok{(}\FunctionTok{file.exists}\NormalTok{(fit\_sba\_paths))) \{}
\NormalTok{  fit\_sba }\OtherTok{\textless{}{-}} \FunctionTok{imap}\NormalTok{(stan\_data\_wba, }\ControlFlowTok{function}\NormalTok{(data, scenario) \{}
\NormalTok{    fit }\OtherTok{\textless{}{-}}\NormalTok{ mod\_simple}\SpecialCharTok{$}\FunctionTok{sample}\NormalTok{(}
      \AttributeTok{data          =}\NormalTok{ data,}
      \AttributeTok{chains        =} \DecValTok{2}\NormalTok{,}
      \AttributeTok{iter\_warmup   =} \DecValTok{0}\NormalTok{,}
      \AttributeTok{iter\_sampling =} \DecValTok{1000}\NormalTok{,}
      \AttributeTok{fixed\_param   =} \ConstantTok{TRUE}\NormalTok{,}
      \AttributeTok{refresh       =} \DecValTok{0}
\NormalTok{    )}
\NormalTok{    fit}\SpecialCharTok{$}\FunctionTok{save\_object}\NormalTok{(fit\_sba\_paths[[scenario]])}
\NormalTok{    fit}
\NormalTok{  \})}
\NormalTok{\} }\ControlFlowTok{else}\NormalTok{ \{}
\NormalTok{  fit\_sba }\OtherTok{\textless{}{-}} \FunctionTok{map}\NormalTok{(fit\_sba\_paths, readRDS)}
\NormalTok{\}}
\end{Highlighting}
\end{Shaded}

\subsection{Fitting the Weighted
Model}\label{fitting-the-weighted-model}

\begin{Shaded}
\begin{Highlighting}[]
\NormalTok{fit\_wba\_raw\_paths }\OtherTok{\textless{}{-}}\NormalTok{ here}\SpecialCharTok{::}\FunctionTok{here}\NormalTok{(}
  \StringTok{"simmodels"}\NormalTok{,}
  \FunctionTok{paste0}\NormalTok{(}\StringTok{"ch10\_fit\_wba\_raw\_"}\NormalTok{, }\FunctionTok{names}\NormalTok{(stan\_data\_wba), }\StringTok{".rds"}\NormalTok{)}
\NormalTok{) }\SpecialCharTok{|\textgreater{}} \FunctionTok{set\_names}\NormalTok{(}\FunctionTok{names}\NormalTok{(stan\_data\_wba))}

\ControlFlowTok{if}\NormalTok{ (regenerate\_simulations }\SpecialCharTok{||} \SpecialCharTok{!}\FunctionTok{all}\NormalTok{(}\FunctionTok{file.exists}\NormalTok{(fit\_wba\_raw\_paths))) \{}
\NormalTok{  seeds }\OtherTok{\textless{}{-}} \FunctionTok{set\_names}\NormalTok{(}\DecValTok{101}\SpecialCharTok{:}\DecValTok{103}\NormalTok{, }\FunctionTok{names}\NormalTok{(stan\_data\_wba))}

\NormalTok{  fit\_wba\_raw }\OtherTok{\textless{}{-}} \FunctionTok{imap}\NormalTok{(stan\_data\_wba, }\ControlFlowTok{function}\NormalTok{(data, scenario) \{}
\NormalTok{    fit }\OtherTok{\textless{}{-}} \FunctionTok{do.call}\NormalTok{(}
\NormalTok{      mod\_weighted\_raw}\SpecialCharTok{$}\NormalTok{sample,}
      \FunctionTok{c}\NormalTok{(}\FunctionTok{list}\NormalTok{(}\AttributeTok{data =}\NormalTok{ data, }\AttributeTok{seed =}\NormalTok{ seeds[[scenario]]), sampler\_args)}
\NormalTok{    )}
    \CommentTok{\# $save\_object() internalises the CSV draws so the file is}
    \CommentTok{\# self{-}contained. Plain saveRDS() would store dangling pointers}
    \CommentTok{\# to temporary CSVs that no longer exist in a new session.}
\NormalTok{    fit}\SpecialCharTok{$}\FunctionTok{save\_object}\NormalTok{(fit\_wba\_raw\_paths[[scenario]])}
\NormalTok{    fit}
\NormalTok{  \})}
\NormalTok{\} }\ControlFlowTok{else}\NormalTok{ \{}
\NormalTok{  fit\_wba\_raw }\OtherTok{\textless{}{-}} \FunctionTok{map}\NormalTok{(fit\_wba\_raw\_paths, readRDS)}
\NormalTok{\}}
\CommentTok{\# Check diagnostics}
\NormalTok{check\_diagnostics }\OtherTok{\textless{}{-}} \ControlFlowTok{function}\NormalTok{(fit, label) \{}
\NormalTok{  diag\_smry }\OtherTok{\textless{}{-}}\NormalTok{ fit}\SpecialCharTok{$}\FunctionTok{diagnostic\_summary}\NormalTok{()}
\NormalTok{  min\_ebfmi }\OtherTok{\textless{}{-}} \FunctionTok{min}\NormalTok{(diag\_smry}\SpecialCharTok{$}\NormalTok{ebfmi)}
  
\NormalTok{  smry }\OtherTok{\textless{}{-}}\NormalTok{ fit}\SpecialCharTok{$}\FunctionTok{summary}\NormalTok{() }\SpecialCharTok{|\textgreater{}} \FunctionTok{filter}\NormalTok{(}\SpecialCharTok{!}\FunctionTok{is.na}\NormalTok{(rhat))}
\NormalTok{  n\_fail }\OtherTok{\textless{}{-}}\NormalTok{ smry }\SpecialCharTok{|\textgreater{}} \FunctionTok{filter}\NormalTok{(rhat }\SpecialCharTok{\textgreater{}} \FloatTok{1.01} \SpecialCharTok{|}\NormalTok{ ess\_bulk }\SpecialCharTok{\textless{}} \DecValTok{400} \SpecialCharTok{|}\NormalTok{ ess\_tail }\SpecialCharTok{\textless{}} \DecValTok{400}\NormalTok{) }\SpecialCharTok{|\textgreater{}} \FunctionTok{nrow}\NormalTok{()}
  
  \FunctionTok{cat}\NormalTok{(}\FunctionTok{sprintf}\NormalTok{(}\StringTok{"\%{-}25s | Divs: \%3d | Min E{-}BFMI: \%.2f | Failures: \%d}\SpecialCharTok{\textbackslash{}n}\StringTok{"}\NormalTok{, }
\NormalTok{              label, }\FunctionTok{sum}\NormalTok{(diag\_smry}\SpecialCharTok{$}\NormalTok{num\_divergent), min\_ebfmi, n\_fail))}
\NormalTok{\}}

\FunctionTok{cat}\NormalTok{(}\StringTok{"}\SpecialCharTok{\textbackslash{}n}\StringTok{=== Weighted Bayesian Agent (Raw Parameterisation) ===}\SpecialCharTok{\textbackslash{}n}\StringTok{"}\NormalTok{)}
\end{Highlighting}
\end{Shaded}

\begin{verbatim}

=== Weighted Bayesian Agent (Raw Parameterisation) ===
\end{verbatim}

\begin{Shaded}
\begin{Highlighting}[]
\FunctionTok{iwalk}\NormalTok{(fit\_wba\_raw, }\SpecialCharTok{\textasciitilde{}} \FunctionTok{check\_diagnostics}\NormalTok{(.x, }\FunctionTok{paste}\NormalTok{(}\StringTok{"WBA Raw"}\NormalTok{, .y)))}
\end{Highlighting}
\end{Shaded}

\begin{verbatim}
WBA Raw Self_Focused      | Divs:   0 | Min E-BFMI: 1.02 | Failures: 0
WBA Raw Balanced          | Divs:   0 | Min E-BFMI: 0.97 | Failures: 0
WBA Raw Socially_Influenced | Divs:   0 | Min E-BFMI: 0.97 | Failures: 0
\end{verbatim}

SBA has no parameters to fit, so it makes no sense to check sampling
diagnostics. The samplers are healthy for WBA models on all three
datasets --- no divergences, good E-BFMI, \(\hat{R}\) and ESS within
bounds. So we move on to predictive checks.

\section{Prior and Posterior Predictive
Checks}\label{prior-and-posterior-predictive-checks}

\begin{Shaded}
\begin{Highlighting}[]
\CommentTok{\# Generalized Extraction Function for Posteriors}
\NormalTok{extract\_predictive\_data }\OtherTok{\textless{}{-}} \ControlFlowTok{function}\NormalTok{(fit, stan\_data, scenario\_name, }\AttributeTok{target\_var =} \StringTok{"posterior\_pred"}\NormalTok{) \{}
\NormalTok{  obs\_df }\OtherTok{\textless{}{-}} \FunctionTok{tibble}\NormalTok{(}
    \AttributeTok{obs\_id =} \FunctionTok{seq\_len}\NormalTok{(stan\_data}\SpecialCharTok{$}\NormalTok{N),}
    \AttributeTok{blue1  =}\NormalTok{ stan\_data}\SpecialCharTok{$}\NormalTok{blue1,}
    \AttributeTok{blue2  =}\NormalTok{ stan\_data}\SpecialCharTok{$}\NormalTok{blue2,}
    \AttributeTok{choice =}\NormalTok{ stan\_data}\SpecialCharTok{$}\NormalTok{choice}
\NormalTok{  )}

\NormalTok{  pred\_samples }\OtherTok{\textless{}{-}}\NormalTok{ posterior}\SpecialCharTok{::}\FunctionTok{as\_draws\_df}\NormalTok{(fit}\SpecialCharTok{$}\FunctionTok{draws}\NormalTok{(target\_var)) }\SpecialCharTok{|\textgreater{}}
\NormalTok{    dplyr}\SpecialCharTok{::}\FunctionTok{select}\NormalTok{(}\SpecialCharTok{{-}}\NormalTok{.chain, }\SpecialCharTok{{-}}\NormalTok{.iteration) }\SpecialCharTok{|\textgreater{}} 
\NormalTok{    tidyr}\SpecialCharTok{::}\FunctionTok{pivot\_longer}\NormalTok{(}
      \AttributeTok{cols =}\NormalTok{ dplyr}\SpecialCharTok{::}\FunctionTok{starts\_with}\NormalTok{(target\_var), }
      \AttributeTok{names\_to =} \StringTok{"obs\_id"}\NormalTok{, }
      \AttributeTok{values\_to =} \StringTok{"pred"}
\NormalTok{    ) }\SpecialCharTok{|\textgreater{}}
\NormalTok{    dplyr}\SpecialCharTok{::}\FunctionTok{mutate}\NormalTok{(}\AttributeTok{obs\_id =}\NormalTok{ readr}\SpecialCharTok{::}\FunctionTok{parse\_number}\NormalTok{(obs\_id))}

\NormalTok{  obs\_summary }\OtherTok{\textless{}{-}}\NormalTok{ obs\_df }\SpecialCharTok{|\textgreater{}}
\NormalTok{    dplyr}\SpecialCharTok{::}\FunctionTok{group\_by}\NormalTok{(blue1, blue2) }\SpecialCharTok{|\textgreater{}}
\NormalTok{    dplyr}\SpecialCharTok{::}\FunctionTok{summarise}\NormalTok{(}\AttributeTok{p\_obs =} \FunctionTok{mean}\NormalTok{(choice), }\AttributeTok{.groups =} \StringTok{"drop"}\NormalTok{)}

\NormalTok{  pred\_summary }\OtherTok{\textless{}{-}}\NormalTok{ pred\_samples }\SpecialCharTok{|\textgreater{}}
\NormalTok{    dplyr}\SpecialCharTok{::}\FunctionTok{left\_join}\NormalTok{(obs\_df, }\AttributeTok{by =} \StringTok{"obs\_id"}\NormalTok{) }\SpecialCharTok{|\textgreater{}}
\NormalTok{    dplyr}\SpecialCharTok{::}\FunctionTok{group\_by}\NormalTok{(.draw, blue1, blue2) }\SpecialCharTok{|\textgreater{}}
\NormalTok{    dplyr}\SpecialCharTok{::}\FunctionTok{summarise}\NormalTok{(}\AttributeTok{p\_draw =} \FunctionTok{mean}\NormalTok{(pred), }\AttributeTok{.groups =} \StringTok{"drop"}\NormalTok{) }\SpecialCharTok{|\textgreater{}}
\NormalTok{    dplyr}\SpecialCharTok{::}\FunctionTok{group\_by}\NormalTok{(blue1, blue2) }\SpecialCharTok{|\textgreater{}}
\NormalTok{    dplyr}\SpecialCharTok{::}\FunctionTok{summarise}\NormalTok{(}
      \AttributeTok{p\_pred =} \FunctionTok{median}\NormalTok{(p\_draw),}
      \AttributeTok{lower  =} \FunctionTok{quantile}\NormalTok{(p\_draw, }\FloatTok{0.025}\NormalTok{),}
      \AttributeTok{upper  =} \FunctionTok{quantile}\NormalTok{(p\_draw, }\FloatTok{0.975}\NormalTok{),}
      \AttributeTok{.groups =} \StringTok{"drop"}
\NormalTok{    ) }\SpecialCharTok{|\textgreater{}}
\NormalTok{    dplyr}\SpecialCharTok{::}\FunctionTok{left\_join}\NormalTok{(obs\_summary, }\AttributeTok{by =} \FunctionTok{c}\NormalTok{(}\StringTok{"blue1"}\NormalTok{, }\StringTok{"blue2"}\NormalTok{)) }\SpecialCharTok{|\textgreater{}}
\NormalTok{    dplyr}\SpecialCharTok{::}\FunctionTok{mutate}\NormalTok{(}
      \AttributeTok{social\_evidence =} \FunctionTok{factor}\NormalTok{(}
\NormalTok{        blue2, }\AttributeTok{levels =} \DecValTok{0}\SpecialCharTok{:}\DecValTok{3}\NormalTok{,}
        \AttributeTok{labels =} \FunctionTok{c}\NormalTok{(}\StringTok{"Clear Red"}\NormalTok{, }\StringTok{"Maybe Red"}\NormalTok{, }\StringTok{"Maybe Blue"}\NormalTok{, }\StringTok{"Clear Blue"}\NormalTok{)}
\NormalTok{      ),}
      \AttributeTok{Scenario =}\NormalTok{ scenario\_name}
\NormalTok{    )}
  
  \FunctionTok{return}\NormalTok{(pred\_summary)}
\NormalTok{\}}

\CommentTok{\# Standardized Plotting Function}
\NormalTok{plot\_predictive\_check }\OtherTok{\textless{}{-}} \ControlFlowTok{function}\NormalTok{(plot\_data, title, subtitle) \{}
\NormalTok{  semantic\_colors }\OtherTok{\textless{}{-}} \FunctionTok{c}\NormalTok{(}\StringTok{"Clear Red"}  \OtherTok{=} \StringTok{"\#d73027"}\NormalTok{, }\StringTok{"Maybe Red"}  \OtherTok{=} \StringTok{"\#fdae61"}\NormalTok{,}
                       \StringTok{"Maybe Blue"} \OtherTok{=} \StringTok{"\#abd9e9"}\NormalTok{, }\StringTok{"Clear Blue"} \OtherTok{=} \StringTok{"\#4575b4"}\NormalTok{)}
  
  \FunctionTok{ggplot}\NormalTok{(plot\_data, }
         \FunctionTok{aes}\NormalTok{(}\AttributeTok{x =}\NormalTok{ blue1, }\AttributeTok{colour =}\NormalTok{ social\_evidence, }\AttributeTok{fill =}\NormalTok{ social\_evidence)) }\SpecialCharTok{+}
    \FunctionTok{geom\_linerange}\NormalTok{(}\FunctionTok{aes}\NormalTok{(}\AttributeTok{ymin =}\NormalTok{ lower, }\AttributeTok{ymax =}\NormalTok{ upper),}
                   \AttributeTok{position =} \FunctionTok{position\_dodge}\NormalTok{(}\AttributeTok{width =} \FloatTok{0.6}\NormalTok{), }
                   \AttributeTok{alpha =} \FloatTok{0.5}\NormalTok{, }\AttributeTok{linewidth =} \DecValTok{2}\NormalTok{) }\SpecialCharTok{+}
    \FunctionTok{geom\_line}\NormalTok{(}\FunctionTok{aes}\NormalTok{(}\AttributeTok{y =}\NormalTok{ p\_pred), }
              \AttributeTok{position =} \FunctionTok{position\_dodge}\NormalTok{(}\AttributeTok{width =} \FloatTok{0.6}\NormalTok{), }
              \AttributeTok{linewidth =} \DecValTok{1}\NormalTok{) }\SpecialCharTok{+}
    \FunctionTok{geom\_point}\NormalTok{(}\FunctionTok{aes}\NormalTok{(}\AttributeTok{y =}\NormalTok{ p\_obs), }
               \AttributeTok{shape =} \DecValTok{21}\NormalTok{, }\AttributeTok{colour =} \StringTok{"black"}\NormalTok{, }\AttributeTok{stroke =} \FloatTok{0.8}\NormalTok{, }\AttributeTok{size =} \DecValTok{2}\NormalTok{,}
               \AttributeTok{position =} \FunctionTok{position\_dodge}\NormalTok{(}\AttributeTok{width =} \FloatTok{0.6}\NormalTok{)) }\SpecialCharTok{+}
    \FunctionTok{facet\_wrap}\NormalTok{(}\SpecialCharTok{\textasciitilde{}}\NormalTok{ Scenario) }\SpecialCharTok{+}
    \FunctionTok{scale\_color\_manual}\NormalTok{(}\AttributeTok{values =}\NormalTok{ semantic\_colors) }\SpecialCharTok{+}
    \FunctionTok{scale\_fill\_manual}\NormalTok{(}\AttributeTok{values =}\NormalTok{ semantic\_colors) }\SpecialCharTok{+}
    \FunctionTok{scale\_x\_continuous}\NormalTok{(}\AttributeTok{breaks =} \DecValTok{0}\SpecialCharTok{:}\DecValTok{8}\NormalTok{) }\SpecialCharTok{+}
    \FunctionTok{coord\_cartesian}\NormalTok{(}\AttributeTok{ylim =} \FunctionTok{c}\NormalTok{(}\DecValTok{0}\NormalTok{, }\DecValTok{1}\NormalTok{)) }\SpecialCharTok{+}
    \FunctionTok{labs}\NormalTok{(}
      \AttributeTok{title    =}\NormalTok{ title,}
      \AttributeTok{subtitle =}\NormalTok{ subtitle,}
      \AttributeTok{x        =} \StringTok{"Blue marbles in direct sample"}\NormalTok{,}
      \AttributeTok{y        =} \StringTok{"P(choose blue)"}\NormalTok{,}
      \AttributeTok{colour   =} \StringTok{"Social evidence"}\NormalTok{, }
      \AttributeTok{fill     =} \StringTok{"Social evidence"}
\NormalTok{    ) }\SpecialCharTok{+}
    \FunctionTok{theme\_bw}\NormalTok{() }\SpecialCharTok{+}
    \FunctionTok{theme}\NormalTok{(}\AttributeTok{legend.position =} \StringTok{"bottom"}\NormalTok{,}
          \AttributeTok{strip.text =} \FunctionTok{element\_text}\NormalTok{(}\AttributeTok{face =} \StringTok{"bold"}\NormalTok{, }\AttributeTok{size =} \DecValTok{11}\NormalTok{))}
\NormalTok{\}}

\CommentTok{\# Extract predictions for the SBA}
\CommentTok{\# Extract SBA posterior predictive checks across all three scenarios}
\NormalTok{sba\_post\_ppc }\OtherTok{\textless{}{-}} \FunctionTok{imap\_dfr}\NormalTok{(fit\_sba, }\ControlFlowTok{function}\NormalTok{(fit, scenario) \{}
  \FunctionTok{extract\_predictive\_data}\NormalTok{(}
    \AttributeTok{fit =}\NormalTok{ fit,}
    \AttributeTok{stan\_data =}\NormalTok{ stan\_data\_wba[[scenario]],}
    \AttributeTok{scenario\_name =}\NormalTok{ scenario,}
    \AttributeTok{target\_var =} \StringTok{"posterior\_pred"}
\NormalTok{  )}
\NormalTok{\})}

\CommentTok{\# Plot SBA across all three datasets}
\FunctionTok{plot\_predictive\_check}\NormalTok{(}
\NormalTok{  sba\_post\_ppc,}
  \AttributeTok{title =} \StringTok{"Posterior Predictive Checks: Simple Bayesian Agent"}\NormalTok{,}
  \AttributeTok{subtitle =} \StringTok{"SBA fits the Balanced dataset but misses the other two systematically"}
\NormalTok{)}
\end{Highlighting}
\end{Shaded}

\begin{center}
\includegraphics[width=0.8\linewidth,height=\textheight,keepaspectratio]{10-BayesianModels_files/figure-pdf/ch10_posterior_predictive-1.pdf}
\end{center}

\begin{Shaded}
\begin{Highlighting}[]
\CommentTok{\# Custom R function to generate Prior predictions for the raw model}
\NormalTok{generate\_wba\_raw\_prior }\OtherTok{\textless{}{-}} \ControlFlowTok{function}\NormalTok{(stan\_data, scenario\_name, }\AttributeTok{n\_draws =} \DecValTok{4000}\NormalTok{) \{}
  
\NormalTok{  obs\_df }\OtherTok{\textless{}{-}} \FunctionTok{tibble}\NormalTok{(}
    \AttributeTok{obs\_id =} \FunctionTok{seq\_len}\NormalTok{(stan\_data}\SpecialCharTok{$}\NormalTok{N),}
    \AttributeTok{blue1  =}\NormalTok{ stan\_data}\SpecialCharTok{$}\NormalTok{blue1,}
    \AttributeTok{total1 =}\NormalTok{ stan\_data}\SpecialCharTok{$}\NormalTok{total1,}
    \AttributeTok{blue2  =}\NormalTok{ stan\_data}\SpecialCharTok{$}\NormalTok{blue2,}
    \AttributeTok{total2 =}\NormalTok{ stan\_data}\SpecialCharTok{$}\NormalTok{total2,}
    \AttributeTok{choice =}\NormalTok{ stan\_data}\SpecialCharTok{$}\NormalTok{choice}
\NormalTok{  )}

  \CommentTok{\# 1. Draw parameters from the priors defined in the raw Stan model}
\NormalTok{  wd\_prior }\OtherTok{\textless{}{-}} \FunctionTok{rlnorm}\NormalTok{(n\_draws, }\DecValTok{0}\NormalTok{, }\FloatTok{0.5}\NormalTok{)}
\NormalTok{  ws\_prior }\OtherTok{\textless{}{-}} \FunctionTok{rlnorm}\NormalTok{(n\_draws, }\DecValTok{0}\NormalTok{, }\FloatTok{0.5}\NormalTok{)}

  \CommentTok{\# 2. Simulate choices across all draws and observations}
\NormalTok{  pred\_samples }\OtherTok{\textless{}{-}} \FunctionTok{expand\_grid}\NormalTok{(}\AttributeTok{.draw =} \DecValTok{1}\SpecialCharTok{:}\NormalTok{n\_draws, }\AttributeTok{obs\_id =}\NormalTok{ obs\_df}\SpecialCharTok{$}\NormalTok{obs\_id) }\SpecialCharTok{|\textgreater{}}
    \FunctionTok{left\_join}\NormalTok{(obs\_df, }\AttributeTok{by =} \StringTok{"obs\_id"}\NormalTok{) }\SpecialCharTok{|\textgreater{}}
    \FunctionTok{mutate}\NormalTok{(}
      \AttributeTok{wd =}\NormalTok{ wd\_prior[.draw],}
      \AttributeTok{ws =}\NormalTok{ ws\_prior[.draw],}
      \AttributeTok{alpha\_post =} \FloatTok{0.5} \SpecialCharTok{+}\NormalTok{ wd }\SpecialCharTok{*}\NormalTok{ blue1 }\SpecialCharTok{+}\NormalTok{ ws }\SpecialCharTok{*}\NormalTok{ blue2,}
      \AttributeTok{beta\_post  =} \FloatTok{0.5} \SpecialCharTok{+}\NormalTok{ wd }\SpecialCharTok{*}\NormalTok{ (total1 }\SpecialCharTok{{-}}\NormalTok{ blue1) }\SpecialCharTok{+}\NormalTok{ ws }\SpecialCharTok{*}\NormalTok{ (total2 }\SpecialCharTok{{-}}\NormalTok{ blue2),}
      
      \CommentTok{\# Probability matching: sample belief, then sample choice}
      \AttributeTok{theta =} \FunctionTok{rbeta}\NormalTok{(}\FunctionTok{n}\NormalTok{(), alpha\_post, beta\_post),}
      \AttributeTok{pred  =} \FunctionTok{rbinom}\NormalTok{(}\FunctionTok{n}\NormalTok{(), }\DecValTok{1}\NormalTok{, theta)}
\NormalTok{    )}

  \CommentTok{\# 3. Summarize in the exact same format as the posterior function}
\NormalTok{  obs\_summary }\OtherTok{\textless{}{-}}\NormalTok{ obs\_df }\SpecialCharTok{|\textgreater{}}
    \FunctionTok{group\_by}\NormalTok{(blue1, blue2) }\SpecialCharTok{|\textgreater{}}
    \FunctionTok{summarise}\NormalTok{(}\AttributeTok{p\_obs =} \FunctionTok{mean}\NormalTok{(choice), }\AttributeTok{.groups =} \StringTok{"drop"}\NormalTok{)}

\NormalTok{  pred\_summary }\OtherTok{\textless{}{-}}\NormalTok{ pred\_samples }\SpecialCharTok{|\textgreater{}}
    \FunctionTok{group\_by}\NormalTok{(.draw, blue1, blue2) }\SpecialCharTok{|\textgreater{}}
    \FunctionTok{summarise}\NormalTok{(}\AttributeTok{p\_draw =} \FunctionTok{mean}\NormalTok{(pred), }\AttributeTok{.groups =} \StringTok{"drop"}\NormalTok{) }\SpecialCharTok{|\textgreater{}}
    \FunctionTok{group\_by}\NormalTok{(blue1, blue2) }\SpecialCharTok{|\textgreater{}}
    \FunctionTok{summarise}\NormalTok{(}
      \AttributeTok{p\_pred =} \FunctionTok{median}\NormalTok{(p\_draw),}
      \AttributeTok{lower  =} \FunctionTok{quantile}\NormalTok{(p\_draw, }\FloatTok{0.025}\NormalTok{),}
      \AttributeTok{upper  =} \FunctionTok{quantile}\NormalTok{(p\_draw, }\FloatTok{0.975}\NormalTok{),}
      \AttributeTok{.groups =} \StringTok{"drop"}
\NormalTok{    ) }\SpecialCharTok{|\textgreater{}}
    \FunctionTok{left\_join}\NormalTok{(obs\_summary, }\AttributeTok{by =} \FunctionTok{c}\NormalTok{(}\StringTok{"blue1"}\NormalTok{, }\StringTok{"blue2"}\NormalTok{)) }\SpecialCharTok{|\textgreater{}}
    \FunctionTok{mutate}\NormalTok{(}
      \AttributeTok{social\_evidence =} \FunctionTok{factor}\NormalTok{(}
\NormalTok{        blue2, }\AttributeTok{levels =} \DecValTok{0}\SpecialCharTok{:}\DecValTok{3}\NormalTok{,}
        \AttributeTok{labels =} \FunctionTok{c}\NormalTok{(}\StringTok{"Clear Red"}\NormalTok{, }\StringTok{"Maybe Red"}\NormalTok{, }\StringTok{"Maybe Blue"}\NormalTok{, }\StringTok{"Clear Blue"}\NormalTok{)}
\NormalTok{      ),}
      \AttributeTok{Scenario =}\NormalTok{ scenario\_name}
\NormalTok{    )}

  \FunctionTok{return}\NormalTok{(pred\_summary)}
\NormalTok{\}}

\CommentTok{\# Extract PRIOR predictive checks by simulating in R}
\NormalTok{wba\_raw\_prior\_ppc }\OtherTok{\textless{}{-}} \FunctionTok{imap\_dfr}\NormalTok{(stan\_data\_wba, }\SpecialCharTok{\textasciitilde{}} \FunctionTok{generate\_wba\_raw\_prior}\NormalTok{(.x, .y))}

\CommentTok{\# Plot PRIOR}
\FunctionTok{plot\_predictive\_check}\NormalTok{(}
\NormalTok{  wba\_raw\_prior\_ppc,}
  \AttributeTok{title =} \StringTok{"Prior Predictive Checks: WBA (Raw)"}\NormalTok{,}
  \AttributeTok{subtitle =} \StringTok{"Lines/Bars = Predicted P(Blue) from Prior alone ± 95\% CI | Circles = Observed"}
\NormalTok{)}
\end{Highlighting}
\end{Shaded}

\begin{center}
\includegraphics[width=0.8\linewidth,height=\textheight,keepaspectratio]{10-BayesianModels_files/figure-pdf/ch10_posterior_predictive-2.pdf}
\end{center}

\begin{Shaded}
\begin{Highlighting}[]
\CommentTok{\# Extract POSTERIOR predictive checks for WBA Raw}
\NormalTok{wba\_raw\_post\_ppc }\OtherTok{\textless{}{-}} \FunctionTok{imap\_dfr}\NormalTok{(fit\_wba\_raw, }\ControlFlowTok{function}\NormalTok{(fit, scenario) \{}
  \FunctionTok{extract\_predictive\_data}\NormalTok{(}
    \AttributeTok{fit =}\NormalTok{ fit, }
    \AttributeTok{stan\_data =}\NormalTok{ stan\_data\_wba[[scenario]], }
    \AttributeTok{scenario\_name =}\NormalTok{ scenario,}
    \AttributeTok{target\_var =} \StringTok{"posterior\_pred"}
\NormalTok{  )}
\NormalTok{\})}

\CommentTok{\# Plot POSTERIOR}
\FunctionTok{plot\_predictive\_check}\NormalTok{(}
\NormalTok{  wba\_raw\_post\_ppc,}
  \AttributeTok{title =} \StringTok{"Posterior Predictive Checks: WBA (Raw)"}\NormalTok{,}
  \AttributeTok{subtitle =} \StringTok{"Lines/Bars = Predicted P(Blue) after fitting ± 95\% CI | Circles = Observed"}
\NormalTok{)}
\end{Highlighting}
\end{Shaded}

\begin{center}
\includegraphics[width=0.8\linewidth,height=\textheight,keepaspectratio]{10-BayesianModels_files/figure-pdf/ch10_posterior_predictive-3.pdf}
\end{center}

Note that for SBA prior and posterior predictive checks are the same (so
we only draw them once), as there is no fitted parameter. The SBA
reproduces the Balanced dataset because that \emph{is} the SBA's
generating process. But for the Self-Focused and Socially-Influenced
datasets the predicted lines diverge from the observed proportions ---
the SBA cannot bend its weights, so it cannot capture either bias. This
is a clear \textbf{model misspecification} failure: the predictive check
shows that SBA is clearly the wrong model for the data.

The WBA's posterior predictive checks look broad but well-centered on
every dataset. Where the SBA failed structurally, the WBA appears to fit
adequately So we seem to have our winner! However, as we know from
previous chapters, we should be rigorous and systematically test WBA's
ability to fit the data and recover the parameter values across a wide
range of parameter values.

\subsection{Parameter recovery: WBA
collapses}\label{parameter-recovery-wba-collapses}

Before we go into SBC, let's check if WBA can actually recover the true
parameter values in these three scenarios:

\begin{Shaded}
\begin{Highlighting}[]
\CommentTok{\# Extract posteriors and compare to true values}
\NormalTok{wba\_recovery }\OtherTok{\textless{}{-}} \FunctionTok{imap\_dfr}\NormalTok{(fit\_wba\_raw, }\ControlFlowTok{function}\NormalTok{(fit, name) \{}
  \FunctionTok{as\_draws\_df}\NormalTok{(fit}\SpecialCharTok{$}\FunctionTok{draws}\NormalTok{(}\FunctionTok{c}\NormalTok{(}\StringTok{"weight\_direct"}\NormalTok{, }\StringTok{"weight\_social"}\NormalTok{))) }\SpecialCharTok{|\textgreater{}} 
    \FunctionTok{as\_tibble}\NormalTok{() }\SpecialCharTok{|\textgreater{}} 
    \FunctionTok{select}\NormalTok{(weight\_direct, weight\_social) }\SpecialCharTok{|\textgreater{}} 
    \FunctionTok{mutate}\NormalTok{(}\AttributeTok{Scenario =}\NormalTok{ name)}
\NormalTok{\})}

\CommentTok{\# True values for plotting}
\NormalTok{wba\_true\_df }\OtherTok{\textless{}{-}} \FunctionTok{map\_dfr}\NormalTok{(true\_wba, }\SpecialCharTok{\textasciitilde{}} \FunctionTok{as\_tibble}\NormalTok{(.x), }\AttributeTok{.id =} \StringTok{"Scenario"}\NormalTok{) }\SpecialCharTok{|\textgreater{}}
  \FunctionTok{pivot\_longer}\NormalTok{(}\FunctionTok{c}\NormalTok{(wd, ws), }\AttributeTok{names\_to =} \StringTok{"parameter"}\NormalTok{, }\AttributeTok{values\_to =} \StringTok{"true\_value"}\NormalTok{) }\SpecialCharTok{|\textgreater{}}
  \FunctionTok{mutate}\NormalTok{(}\AttributeTok{parameter =} \FunctionTok{ifelse}\NormalTok{(parameter }\SpecialCharTok{==} \StringTok{"wd"}\NormalTok{, }\StringTok{"weight\_direct"}\NormalTok{, }\StringTok{"weight\_social"}\NormalTok{))}

\CommentTok{\# Plot: weight\_direct recovery}
\NormalTok{p\_wd }\OtherTok{\textless{}{-}}\NormalTok{ wba\_recovery }\SpecialCharTok{|\textgreater{}}
  \FunctionTok{mutate}\NormalTok{(}\AttributeTok{Scenario =} \FunctionTok{factor}\NormalTok{(Scenario, }\AttributeTok{levels =} \FunctionTok{c}\NormalTok{(}\StringTok{"Balanced"}\NormalTok{, }\StringTok{"Self\_Focused"}\NormalTok{, }\StringTok{"Socially\_Influenced"}\NormalTok{))) }\SpecialCharTok{|\textgreater{}}
  \FunctionTok{ggplot}\NormalTok{(}\FunctionTok{aes}\NormalTok{(}\AttributeTok{x =}\NormalTok{ weight\_direct, }\AttributeTok{fill =}\NormalTok{ Scenario, }\AttributeTok{color =}\NormalTok{ Scenario)) }\SpecialCharTok{+}
  \FunctionTok{geom\_density}\NormalTok{(}\AttributeTok{alpha =} \FloatTok{0.3}\NormalTok{, }\AttributeTok{linewidth =} \DecValTok{1}\NormalTok{) }\SpecialCharTok{+}
  \FunctionTok{geom\_vline}\NormalTok{(}\AttributeTok{data =}\NormalTok{ wba\_true\_df }\SpecialCharTok{|\textgreater{}} \FunctionTok{filter}\NormalTok{(parameter }\SpecialCharTok{==} \StringTok{"weight\_direct"}\NormalTok{), }
             \FunctionTok{aes}\NormalTok{(}\AttributeTok{xintercept =}\NormalTok{ true\_value, }\AttributeTok{color =}\NormalTok{ Scenario), }
             \AttributeTok{linetype =} \StringTok{"dashed"}\NormalTok{, }\AttributeTok{linewidth =} \FloatTok{1.2}\NormalTok{, }\AttributeTok{show.legend =} \ConstantTok{FALSE}\NormalTok{) }\SpecialCharTok{+}
  \FunctionTok{scale\_fill\_manual}\NormalTok{(}\AttributeTok{values =} \FunctionTok{c}\NormalTok{(}\StringTok{"Balanced"} \OtherTok{=} \StringTok{"\#4daf4a"}\NormalTok{, }
                                \StringTok{"Self\_Focused"} \OtherTok{=} \StringTok{"\#e41a1c"}\NormalTok{, }
                                \StringTok{"Socially\_Influenced"} \OtherTok{=} \StringTok{"\#377eb8"}\NormalTok{)) }\SpecialCharTok{+}
  \FunctionTok{scale\_color\_manual}\NormalTok{(}\AttributeTok{values =} \FunctionTok{c}\NormalTok{(}\StringTok{"Balanced"} \OtherTok{=} \StringTok{"\#4daf4a"}\NormalTok{, }
                                 \StringTok{"Self\_Focused"} \OtherTok{=} \StringTok{"\#e41a1c"}\NormalTok{, }
                                 \StringTok{"Socially\_Influenced"} \OtherTok{=} \StringTok{"\#377eb8"}\NormalTok{)) }\SpecialCharTok{+}
  \FunctionTok{coord\_cartesian}\NormalTok{(}\AttributeTok{xlim =} \FunctionTok{c}\NormalTok{(}\DecValTok{0}\NormalTok{, }\DecValTok{4}\NormalTok{)) }\SpecialCharTok{+}
  \FunctionTok{labs}\NormalTok{(}\AttributeTok{title =} \StringTok{"Weight on Direct Evidence (w\_d)"}\NormalTok{, }
       \AttributeTok{subtitle =} \StringTok{"Dashed = true generating value"}\NormalTok{,}
       \AttributeTok{x =} \StringTok{"weight\_direct"}\NormalTok{, }\AttributeTok{y =} \StringTok{"Density"}\NormalTok{) }\SpecialCharTok{+}
  \FunctionTok{theme}\NormalTok{(}\AttributeTok{legend.position =} \StringTok{"none"}\NormalTok{)}

\CommentTok{\# Plot: weight\_social recovery  }
\NormalTok{p\_ws }\OtherTok{\textless{}{-}}\NormalTok{ wba\_recovery }\SpecialCharTok{|\textgreater{}}
  \FunctionTok{mutate}\NormalTok{(}\AttributeTok{Scenario =} \FunctionTok{factor}\NormalTok{(Scenario, }\AttributeTok{levels =} \FunctionTok{c}\NormalTok{(}\StringTok{"Balanced"}\NormalTok{, }\StringTok{"Self\_Focused"}\NormalTok{, }\StringTok{"Socially\_Influenced"}\NormalTok{))) }\SpecialCharTok{|\textgreater{}}
  \FunctionTok{ggplot}\NormalTok{(}\FunctionTok{aes}\NormalTok{(}\AttributeTok{x =}\NormalTok{ weight\_social, }\AttributeTok{fill =}\NormalTok{ Scenario, }\AttributeTok{color =}\NormalTok{ Scenario)) }\SpecialCharTok{+}
  \FunctionTok{geom\_density}\NormalTok{(}\AttributeTok{alpha =} \FloatTok{0.3}\NormalTok{, }\AttributeTok{linewidth =} \DecValTok{1}\NormalTok{) }\SpecialCharTok{+}
  \FunctionTok{geom\_vline}\NormalTok{(}\AttributeTok{data =}\NormalTok{ wba\_true\_df }\SpecialCharTok{|\textgreater{}} \FunctionTok{filter}\NormalTok{(parameter }\SpecialCharTok{==} \StringTok{"weight\_social"}\NormalTok{), }
             \FunctionTok{aes}\NormalTok{(}\AttributeTok{xintercept =}\NormalTok{ true\_value, }\AttributeTok{color =}\NormalTok{ Scenario), }
             \AttributeTok{linetype =} \StringTok{"dashed"}\NormalTok{, }\AttributeTok{linewidth =} \FloatTok{1.2}\NormalTok{, }\AttributeTok{show.legend =} \ConstantTok{FALSE}\NormalTok{) }\SpecialCharTok{+}
  \FunctionTok{scale\_fill\_manual}\NormalTok{(}\AttributeTok{values =} \FunctionTok{c}\NormalTok{(}\StringTok{"Balanced"} \OtherTok{=} \StringTok{"\#4daf4a"}\NormalTok{, }
                                \StringTok{"Self\_Focused"} \OtherTok{=} \StringTok{"\#e41a1c"}\NormalTok{, }
                                \StringTok{"Socially\_Influenced"} \OtherTok{=} \StringTok{"\#377eb8"}\NormalTok{)) }\SpecialCharTok{+}
  \FunctionTok{scale\_color\_manual}\NormalTok{(}\AttributeTok{values =} \FunctionTok{c}\NormalTok{(}\StringTok{"Balanced"} \OtherTok{=} \StringTok{"\#4daf4a"}\NormalTok{, }
                                 \StringTok{"Self\_Focused"} \OtherTok{=} \StringTok{"\#e41a1c"}\NormalTok{, }
                                 \StringTok{"Socially\_Influenced"} \OtherTok{=} \StringTok{"\#377eb8"}\NormalTok{)) }\SpecialCharTok{+}
  \FunctionTok{coord\_cartesian}\NormalTok{(}\AttributeTok{xlim =} \FunctionTok{c}\NormalTok{(}\DecValTok{0}\NormalTok{, }\DecValTok{4}\NormalTok{)) }\SpecialCharTok{+}
  \FunctionTok{labs}\NormalTok{(}\AttributeTok{title =} \StringTok{"Weight on Social Evidence (w\_s)"}\NormalTok{, }
       \AttributeTok{subtitle =} \StringTok{"Dashed = true generating value"}\NormalTok{,}
       \AttributeTok{x =} \StringTok{"weight\_social"}\NormalTok{, }\AttributeTok{y =} \StringTok{"Density"}\NormalTok{) }\SpecialCharTok{+}
  \FunctionTok{theme}\NormalTok{(}\AttributeTok{legend.position =} \StringTok{"bottom"}\NormalTok{)}

\NormalTok{p\_wd }\SpecialCharTok{/}\NormalTok{ p\_ws }\SpecialCharTok{+} 
  \FunctionTok{plot\_annotation}\NormalTok{(}
    \AttributeTok{title =} \StringTok{"⚠️ Parameter Recovery: WBA (Raw Parameterisation)"}
\NormalTok{  )}
\end{Highlighting}
\end{Shaded}

\begin{center}
\includegraphics[width=0.8\linewidth,height=\textheight,keepaspectratio]{10-BayesianModels_files/figure-pdf/ch10_parameter_recovery_raw-1.pdf}
\end{center}

\textbf{Something is wrong!} The posteriors are extremely wide and don't
center on the true values. The recovery is inadequate --- especially for
the social weight.

\subsection{Diagnosing the ridge}\label{diagnosing-the-ridge}

Let's visualize the joint posterior to understand what's happening:

\begin{Shaded}
\begin{Highlighting}[]
\CommentTok{\# Extract joint posterior draws}
\NormalTok{post\_draws\_self }\OtherTok{\textless{}{-}} \FunctionTok{as\_draws\_df}\NormalTok{(}
\NormalTok{  fit\_wba\_raw}\SpecialCharTok{$}\NormalTok{Self\_Focused}\SpecialCharTok{$}\FunctionTok{draws}\NormalTok{(}\FunctionTok{c}\NormalTok{(}\StringTok{"weight\_direct"}\NormalTok{, }\StringTok{"weight\_social"}\NormalTok{))}
\NormalTok{) }\SpecialCharTok{|\textgreater{}} \FunctionTok{as\_tibble}\NormalTok{()}

\CommentTok{\# Generate implied prior for comparison}
\NormalTok{wba\_prior\_joint }\OtherTok{\textless{}{-}} \FunctionTok{tibble}\NormalTok{(}
  \AttributeTok{weight\_direct =} \FunctionTok{rlnorm}\NormalTok{(}\DecValTok{5000}\NormalTok{, }\DecValTok{0}\NormalTok{, }\FloatTok{0.5}\NormalTok{),}
  \AttributeTok{weight\_social =} \FunctionTok{rlnorm}\NormalTok{(}\DecValTok{5000}\NormalTok{, }\DecValTok{0}\NormalTok{, }\FloatTok{0.5}\NormalTok{)}
\NormalTok{)}

\FunctionTok{ggplot}\NormalTok{() }\SpecialCharTok{+}
  \CommentTok{\# Prior contour (grey dashed)}
  \FunctionTok{geom\_density\_2d}\NormalTok{(}\AttributeTok{data =}\NormalTok{ wba\_prior\_joint, }
                  \FunctionTok{aes}\NormalTok{(}\AttributeTok{x =}\NormalTok{ weight\_direct, }\AttributeTok{y =}\NormalTok{ weight\_social), }
                  \AttributeTok{color =} \StringTok{"grey60"}\NormalTok{, }\AttributeTok{linetype =} \StringTok{"dashed"}\NormalTok{, }\AttributeTok{bins =} \DecValTok{6}\NormalTok{) }\SpecialCharTok{+}
  \CommentTok{\# Posterior scatter (red to signify pathology)}
  \FunctionTok{geom\_point}\NormalTok{(}\AttributeTok{data =}\NormalTok{ post\_draws\_self, }
             \FunctionTok{aes}\NormalTok{(}\AttributeTok{x =}\NormalTok{ weight\_direct, }\AttributeTok{y =}\NormalTok{ weight\_social), }
             \AttributeTok{alpha =} \FloatTok{0.2}\NormalTok{, }\AttributeTok{color =} \StringTok{"\#d73027"}\NormalTok{, }\AttributeTok{size =} \DecValTok{1}\NormalTok{) }\SpecialCharTok{+}
  \CommentTok{\# Posterior contour}
  \FunctionTok{geom\_density\_2d}\NormalTok{(}\AttributeTok{data =}\NormalTok{ post\_draws\_self, }
                  \FunctionTok{aes}\NormalTok{(}\AttributeTok{x =}\NormalTok{ weight\_direct, }\AttributeTok{y =}\NormalTok{ weight\_social), }
                  \AttributeTok{color =} \StringTok{"\#a50026"}\NormalTok{, }\AttributeTok{bins =} \DecValTok{4}\NormalTok{) }\SpecialCharTok{+}
  \CommentTok{\# True value crosshair}
  \FunctionTok{geom\_vline}\NormalTok{(}\AttributeTok{xintercept =} \FloatTok{1.5}\NormalTok{, }\AttributeTok{color =} \StringTok{"black"}\NormalTok{, }\AttributeTok{linetype =} \StringTok{"dashed"}\NormalTok{, }\AttributeTok{linewidth =} \FloatTok{0.8}\NormalTok{) }\SpecialCharTok{+}
  \FunctionTok{geom\_hline}\NormalTok{(}\AttributeTok{yintercept =} \FloatTok{0.5}\NormalTok{, }\AttributeTok{color =} \StringTok{"black"}\NormalTok{, }\AttributeTok{linetype =} \StringTok{"dashed"}\NormalTok{, }\AttributeTok{linewidth =} \FloatTok{0.8}\NormalTok{) }\SpecialCharTok{+}
  \FunctionTok{geom\_point}\NormalTok{(}\FunctionTok{aes}\NormalTok{(}\AttributeTok{x =} \FloatTok{1.5}\NormalTok{, }\AttributeTok{y =} \FloatTok{0.5}\NormalTok{), }
             \AttributeTok{color =} \StringTok{"black"}\NormalTok{, }\AttributeTok{shape =} \DecValTok{4}\NormalTok{, }\AttributeTok{size =} \DecValTok{4}\NormalTok{, }\AttributeTok{stroke =} \DecValTok{2}\NormalTok{) }\SpecialCharTok{+}
  \FunctionTok{coord\_cartesian}\NormalTok{(}\AttributeTok{xlim =} \FunctionTok{c}\NormalTok{(}\DecValTok{0}\NormalTok{, }\DecValTok{4}\NormalTok{), }\AttributeTok{ylim =} \FunctionTok{c}\NormalTok{(}\DecValTok{0}\NormalTok{, }\DecValTok{4}\NormalTok{)) }\SpecialCharTok{+}
  \FunctionTok{labs}\NormalTok{(}
    \AttributeTok{title =} \StringTok{"The Posterior Ridge: A Severe Identifiability Problem"}\NormalTok{, }
    \AttributeTok{subtitle =} \StringTok{"Self{-}Focused Agent (true: w\_d=1.5, w\_s=0.5). Grey = prior contours."}\NormalTok{,}
    \AttributeTok{x =} \StringTok{"Weight on Direct (w\_d)"}\NormalTok{, }
    \AttributeTok{y =} \StringTok{"Weight on Social (w\_s)"}
\NormalTok{  ) }
\end{Highlighting}
\end{Shaded}

\begin{center}
\includegraphics[width=0.8\linewidth,height=\textheight,keepaspectratio]{10-BayesianModels_files/figure-pdf/ch10_posterior_ridge-1.pdf}
\end{center}

In the plot we see the prior for the two weights (grey dashed contours)
and the posterior (red points and contours). The true value is marked
with a dashed crosshair. From this plot we can identify a potential
source for the poor recovery: the posterior is not a tight blob around
the true value, but rather a long diagonal ridge. If we move along this
ridge, we can see that the likelihood remains high --- the model cannot
distinguish between different combinations of \(w_d\) and \(w_s\) that
produce the same relative weighting of direct vs social evidence. Why is
that? \(w_d = 2, w_s = 2\) produces the \textbf{same average decisions}
as \(w_d = 1, w_s = 1\) --- the absolute scale affects only the
\emph{variance} of choices, which is hard to pin down from binary
outcomes. If you were expecting the ridge to follow the line
\(w_d = w_s\) (where the agent is balanced), that would be wrong --- the
ridge actually follows the line where the \emph{ratio} of weights is
constant, given the inferred preferences of the agent. This specific
Self-Focused agent was simulated with the true parameters \(w_d = 1.5\)
and \(w_s = 0.5\). Therefore, the ridge in that plot follows the
constant ratio of those true values: \(1.5 / 0.5 = 3\). The ridge
follows the line \(w_d = 3 w_s\).

Maybe we just need more trials? Let's check:

\begin{Shaded}
\begin{Highlighting}[]
\CommentTok{\# Set up a grid of simulations across our desired sample sizes}
\NormalTok{n\_sims\_per\_size }\OtherTok{\textless{}{-}} \DecValTok{50}   \CommentTok{\# Increased from 20 to 50}
\NormalTok{sample\_sizes    }\OtherTok{\textless{}{-}} \FunctionTok{c}\NormalTok{(}\DecValTok{180}\NormalTok{, }\DecValTok{360}\NormalTok{, }\DecValTok{720}\NormalTok{) }\CommentTok{\# Doubling the sample size}

\NormalTok{sim\_grid }\OtherTok{\textless{}{-}} \FunctionTok{expand\_grid}\NormalTok{(}
  \AttributeTok{sample\_size =}\NormalTok{ sample\_sizes,}
  \AttributeTok{sim\_id      =} \DecValTok{1}\SpecialCharTok{:}\NormalTok{n\_sims\_per\_size}
\NormalTok{) }\SpecialCharTok{|\textgreater{}}
  \FunctionTok{mutate}\NormalTok{(}
    \CommentTok{\# Draw true parameters from the exact same Lognormal(0, 0.5) prior}
    \AttributeTok{true\_wd =} \FunctionTok{rlnorm}\NormalTok{(}\FunctionTok{n}\NormalTok{(), }\DecValTok{0}\NormalTok{, }\FloatTok{0.5}\NormalTok{),}
    \AttributeTok{true\_ws =} \FunctionTok{rlnorm}\NormalTok{(}\FunctionTok{n}\NormalTok{(), }\DecValTok{0}\NormalTok{, }\FloatTok{0.5}\NormalTok{)}
\NormalTok{  )}

\CommentTok{\# Helper function to generate exactly N trials }
\NormalTok{generate\_random\_trials }\OtherTok{\textless{}{-}} \ControlFlowTok{function}\NormalTok{(wd, ws, n\_trials) \{}
  \CommentTok{\# Sample combinations with replacement to hit the exact trial count}
\NormalTok{  ev\_sample }\OtherTok{\textless{}{-}}\NormalTok{ evidence\_combinations }\SpecialCharTok{|\textgreater{}}
    \FunctionTok{slice\_sample}\NormalTok{(}\AttributeTok{n =}\NormalTok{ n\_trials, }\AttributeTok{replace =} \ConstantTok{TRUE}\NormalTok{)}
  
  \CommentTok{\# Generate choices using your existing function}
  \FunctionTok{generate\_agent\_decisions}\NormalTok{(wd, ws, ev\_sample, }\AttributeTok{n\_samples =} \DecValTok{1}\NormalTok{) }\SpecialCharTok{|\textgreater{}}
    \FunctionTok{mutate}\NormalTok{(}\AttributeTok{total1 =} \DecValTok{8}\NormalTok{L, }\AttributeTok{total2 =} \DecValTok{3}\NormalTok{L)}
\NormalTok{\}}

\FunctionTok{cat}\NormalTok{(}\FunctionTok{sprintf}\NormalTok{(}\StringTok{"Running \%d total parameter recovery simulations...}\SpecialCharTok{\textbackslash{}n}\StringTok{"}\NormalTok{, }\FunctionTok{nrow}\NormalTok{(sim\_grid)))}
\end{Highlighting}
\end{Shaded}

\begin{verbatim}
Running 150 total parameter recovery simulations...
\end{verbatim}

\begin{Shaded}
\begin{Highlighting}[]
\CommentTok{\# Run fits in parallel (using your existing multicore plan)}
\ControlFlowTok{if}\NormalTok{ (regenerate\_simulations }\SpecialCharTok{||} \SpecialCharTok{!}\FunctionTok{file.exists}\NormalTok{(}\StringTok{"simmodels/ch10\_n\_investigation.rds"}\NormalTok{)) \{}
\NormalTok{  recovery\_results }\OtherTok{\textless{}{-}} \FunctionTok{future\_pmap\_dfr}\NormalTok{(}
    \FunctionTok{list}\NormalTok{(sim\_grid}\SpecialCharTok{$}\NormalTok{sample\_size, sim\_grid}\SpecialCharTok{$}\NormalTok{sim\_id, sim\_grid}\SpecialCharTok{$}\NormalTok{true\_wd, sim\_grid}\SpecialCharTok{$}\NormalTok{true\_ws),}
    \ControlFlowTok{function}\NormalTok{(n\_trials, id, wd, ws) \{}
      
      \CommentTok{\# 1. Generate data and prep for Stan}
\NormalTok{      df }\OtherTok{\textless{}{-}} \FunctionTok{generate\_random\_trials}\NormalTok{(wd, ws, n\_trials)}
\NormalTok{      stan\_data }\OtherTok{\textless{}{-}} \FunctionTok{prepare\_stan\_data}\NormalTok{(df)}
      
      \CommentTok{\# 2. Fit the raw model (using reduced iterations for speed )}
\NormalTok{      fit }\OtherTok{\textless{}{-}}\NormalTok{ mod\_weighted\_raw}\SpecialCharTok{$}\FunctionTok{sample}\NormalTok{(}
        \AttributeTok{data =}\NormalTok{ stan\_data, }
        \AttributeTok{chains =} \DecValTok{2}\NormalTok{, }\AttributeTok{iter\_warmup =} \DecValTok{500}\NormalTok{, }\AttributeTok{iter\_sampling =} \DecValTok{500}\NormalTok{, }
        \AttributeTok{refresh =} \DecValTok{0}\NormalTok{, }\AttributeTok{show\_messages =} \ConstantTok{FALSE}\NormalTok{,}
        \AttributeTok{save\_warmup =} \ConstantTok{TRUE}\NormalTok{, }\AttributeTok{output\_dir =} \StringTok{"models/temp"}
\NormalTok{      )}
      
      \CommentTok{\# 3. Extract posteriors}
\NormalTok{      draws }\OtherTok{\textless{}{-}} \FunctionTok{as\_draws\_df}\NormalTok{(fit}\SpecialCharTok{$}\FunctionTok{draws}\NormalTok{(}\FunctionTok{c}\NormalTok{(}\StringTok{"weight\_direct"}\NormalTok{, }\StringTok{"weight\_social"}\NormalTok{)))}
      
      \CommentTok{\# 4. Summarize recovery}
      \FunctionTok{tibble}\NormalTok{(}
        \AttributeTok{sample\_size =}\NormalTok{ n\_trials,}
        \AttributeTok{sim\_id      =}\NormalTok{ id,}
        \AttributeTok{true\_wd     =}\NormalTok{ wd,}
        \AttributeTok{true\_ws     =}\NormalTok{ ws,}
        \AttributeTok{est\_wd      =} \FunctionTok{median}\NormalTok{(draws}\SpecialCharTok{$}\NormalTok{weight\_direct),}
        \AttributeTok{lower\_wd    =} \FunctionTok{quantile}\NormalTok{(draws}\SpecialCharTok{$}\NormalTok{weight\_direct, }\FloatTok{0.05}\NormalTok{),}
        \AttributeTok{upper\_wd    =} \FunctionTok{quantile}\NormalTok{(draws}\SpecialCharTok{$}\NormalTok{weight\_direct, }\FloatTok{0.95}\NormalTok{),}
        \AttributeTok{est\_ws      =} \FunctionTok{median}\NormalTok{(draws}\SpecialCharTok{$}\NormalTok{weight\_social),}
        \AttributeTok{lower\_ws    =} \FunctionTok{quantile}\NormalTok{(draws}\SpecialCharTok{$}\NormalTok{weight\_social, }\FloatTok{0.05}\NormalTok{),}
        \AttributeTok{upper\_ws    =} \FunctionTok{quantile}\NormalTok{(draws}\SpecialCharTok{$}\NormalTok{weight\_social, }\FloatTok{0.95}\NormalTok{)}
\NormalTok{      )}
\NormalTok{    \},}
    \AttributeTok{.options =} \FunctionTok{furrr\_options}\NormalTok{(}\AttributeTok{seed =} \DecValTok{123}\NormalTok{)}
\NormalTok{  )}
  \FunctionTok{saveRDS}\NormalTok{(recovery\_results, }\StringTok{"simmodels/ch10\_n\_investigation.rds"}\NormalTok{)}
\NormalTok{\} }\ControlFlowTok{else}\NormalTok{ \{}
\NormalTok{  recovery\_results }\OtherTok{\textless{}{-}} \FunctionTok{readRDS}\NormalTok{(}\StringTok{"simmodels/ch10\_n\_investigation.rds"}\NormalTok{)}
\NormalTok{\}}

\CommentTok{\# Reshape data for plotting}
\NormalTok{plot\_df }\OtherTok{\textless{}{-}}\NormalTok{ recovery\_results }\SpecialCharTok{|\textgreater{}}
  \FunctionTok{pivot\_longer}\NormalTok{(}
    \AttributeTok{cols =} \FunctionTok{c}\NormalTok{(}\FunctionTok{ends\_with}\NormalTok{(}\StringTok{"wd"}\NormalTok{), }\FunctionTok{ends\_with}\NormalTok{(}\StringTok{"ws"}\NormalTok{)),}
    \AttributeTok{names\_to =} \FunctionTok{c}\NormalTok{(}\StringTok{".value"}\NormalTok{, }\StringTok{"parameter"}\NormalTok{),}
    \AttributeTok{names\_pattern =} \StringTok{"(.*)\_(w[ds])"}
\NormalTok{  ) }\SpecialCharTok{|\textgreater{}}
  \FunctionTok{mutate}\NormalTok{(}
    \AttributeTok{parameter =} \FunctionTok{ifelse}\NormalTok{(parameter }\SpecialCharTok{==} \StringTok{"wd"}\NormalTok{, }\StringTok{"Weight Direct (w\_d)"}\NormalTok{, }\StringTok{"Weight Social (w\_s)"}\NormalTok{),}
    \CommentTok{\# Force the factor levels so the facets order correctly top{-}to{-}bottom}
    \AttributeTok{sample\_size\_label =} \FunctionTok{factor}\NormalTok{(}\FunctionTok{paste}\NormalTok{(}\StringTok{"N ="}\NormalTok{, sample\_size), }
                               \AttributeTok{levels =} \FunctionTok{c}\NormalTok{(}\StringTok{"N = 180"}\NormalTok{, }\StringTok{"N = 360"}\NormalTok{, }\StringTok{"N = 720"}\NormalTok{))}
\NormalTok{  )}

\CommentTok{\# Plot True vs Estimated values}
\FunctionTok{ggplot}\NormalTok{(plot\_df, }\FunctionTok{aes}\NormalTok{(}\AttributeTok{x =}\NormalTok{ true, }\AttributeTok{y =}\NormalTok{ est, }\AttributeTok{color =} \FunctionTok{factor}\NormalTok{(sample\_size))) }\SpecialCharTok{+}
  \FunctionTok{geom\_abline}\NormalTok{(}\AttributeTok{slope =} \DecValTok{1}\NormalTok{, }\AttributeTok{intercept =} \DecValTok{0}\NormalTok{, }\AttributeTok{linetype =} \StringTok{"dashed"}\NormalTok{, }\AttributeTok{color =} \StringTok{"gray50"}\NormalTok{) }\SpecialCharTok{+}
  \CommentTok{\# Lowered alpha slightly to accommodate the higher density of points}
  \FunctionTok{geom\_errorbar}\NormalTok{(}\FunctionTok{aes}\NormalTok{(}\AttributeTok{ymin =}\NormalTok{ lower, }\AttributeTok{ymax =}\NormalTok{ upper), }\AttributeTok{alpha =} \FloatTok{0.2}\NormalTok{, }\AttributeTok{width =} \DecValTok{0}\NormalTok{) }\SpecialCharTok{+}
  \FunctionTok{geom\_point}\NormalTok{(}\AttributeTok{alpha =} \FloatTok{0.6}\NormalTok{, }\AttributeTok{size =} \DecValTok{2}\NormalTok{) }\SpecialCharTok{+}
  \FunctionTok{geom\_smooth}\NormalTok{(}\AttributeTok{method=}\NormalTok{lm, }\AttributeTok{se=}\ConstantTok{FALSE}\NormalTok{, }\AttributeTok{color =} \StringTok{"red"}\NormalTok{, }\AttributeTok{linetype =} \StringTok{"dashed"}\NormalTok{) }\SpecialCharTok{+}
  \FunctionTok{facet\_grid}\NormalTok{(sample\_size\_label }\SpecialCharTok{\textasciitilde{}}\NormalTok{ parameter) }\SpecialCharTok{+}
  \FunctionTok{scale\_color\_viridis\_d}\NormalTok{(}\AttributeTok{option =} \StringTok{"C"}\NormalTok{, }\AttributeTok{end =} \FloatTok{0.8}\NormalTok{) }\SpecialCharTok{+}
  \FunctionTok{coord\_cartesian}\NormalTok{(}\AttributeTok{xlim =} \FunctionTok{c}\NormalTok{(}\DecValTok{0}\NormalTok{, }\DecValTok{4}\NormalTok{), }\AttributeTok{ylim =} \FunctionTok{c}\NormalTok{(}\DecValTok{0}\NormalTok{, }\DecValTok{4}\NormalTok{)) }\SpecialCharTok{+}
  \FunctionTok{labs}\NormalTok{(}
    \AttributeTok{title =} \StringTok{"Does More Data Help? NO!"}\NormalTok{,}
    \AttributeTok{subtitle =} \StringTok{"Even at N=720, recovery remains terrible due to the posterior ridge."}\NormalTok{,}
    \AttributeTok{x =} \StringTok{"True Parameter Value"}\NormalTok{,}
    \AttributeTok{y =} \StringTok{"Posterior Median"}\NormalTok{,}
    \AttributeTok{color =} \StringTok{"Sample Size"}
\NormalTok{  ) }\SpecialCharTok{+}
  \FunctionTok{theme}\NormalTok{(}
    \AttributeTok{legend.position =} \StringTok{"none"}\NormalTok{,}
    \AttributeTok{panel.border =} \FunctionTok{element\_rect}\NormalTok{(}\AttributeTok{fill =} \ConstantTok{NA}\NormalTok{, }\AttributeTok{color =} \StringTok{"grey80"}\NormalTok{),}
    \AttributeTok{strip.text =} \FunctionTok{element\_text}\NormalTok{(}\AttributeTok{face =} \StringTok{"bold"}\NormalTok{, }\AttributeTok{size =} \DecValTok{11}\NormalTok{)}
\NormalTok{  )}
\end{Highlighting}
\end{Shaded}

\begin{center}
\includegraphics[width=0.8\linewidth,height=\textheight,keepaspectratio]{10-BayesianModels_files/figure-pdf/ch10_trial_number_investigation-1.pdf}
\end{center}

While there is a positive slope between true and inferred parameter
values, and it slightly increase with trials, the overall recovery is
not ideal. This is not a statistical power problem --- it's a
\textbf{fundamental identifiability problem} in the parameterisation.

\section{Round 2: PBA and Reparameterised
WBA}\label{round-2-pba-and-reparameterised-wba}

The problem is structural: the likelihood depends on \(w_d\) and \(w_s\)
only through their \emph{ratio} and not (much) their \emph{sum}. We have
two options for fixing this, and they correspond to two different
cognitive commitments.

\textbf{Option 1 --- fix the scale.} If we believe attention has a
limited budget, the total weight \(w_d + w_s\) could be fixed at 1, and
the only thing left to estimate is how that budget is allocated. This
gives us the \textbf{Proportional Bayesian Agent (PBA)}:

\[\alpha_{\text{post}} = 0.5 + p \cdot k_1 + (1-p) \cdot k_2\]

with \(p \in (0,1)\). One free parameter, no identifiability problem by
construction.

\textbf{Option 2 --- keep the scale, but reparameterise.} If we don't
want to assume a fixed attention budget, we can keep two parameters but
rotate them into orthogonal directions:

\[\rho = \frac{w_d}{w_d + w_s} \in (0,1) \qquad \kappa = w_d + w_s > 0\]

Now \(\rho\) answers ``which source is trusted more?'' ---
well-identified from conflict trials --- and \(\kappa\) answers ``how
much evidence is inferred to be there overall?'' --- harder to pin down
but at least living in its own dimension with its own prior. The
original weights are reconstructed as \(w_d = \rho\kappa\) and
\(w_s = (1-\rho)\kappa\). We call this the \textbf{reparameterised WBA}.

The two solutions make different scientific claims (fixed vs.~free total
attention) and we'll fit both.

\section{Stan Implementation: The
Solutions}\label{stan-implementation-the-solutions}

\subsection{Proportional Bayesian Agent
(PBA)}\label{proportional-bayesian-agent-pba}

\begin{Shaded}
\begin{Highlighting}[]
\FunctionTok{\textbackslash{}usetikzlibrary}\NormalTok{\{bayesnet\}}
\KeywordTok{\textbackslash{}begin}\NormalTok{\{}\ExtensionTok{tikzpicture}\NormalTok{\}}
  \CommentTok{\% Fixed Prior}
  \FunctionTok{\textbackslash{}node}\NormalTok{[const] (alpha0) \{}\SpecialStringTok{$}\SpecialCharTok{\textbackslash{}alpha}\SpecialStringTok{\_0$}\NormalTok{\};}
  \FunctionTok{\textbackslash{}node}\NormalTok{[const, right=1cm of alpha0] (beta0) \{}\SpecialStringTok{$}\SpecialCharTok{\textbackslash{}beta}\SpecialStringTok{\_0$}\NormalTok{\};}

  \CommentTok{\% Allocation Parameter}
  \FunctionTok{\textbackslash{}node}\NormalTok{[latent, below=0.8cm of alpha0, xshift=0.5cm] (p) \{}\SpecialStringTok{$p$}\NormalTok{\}; }

  \CommentTok{\% Inputs}
  \FunctionTok{\textbackslash{}node}\NormalTok{[obs, below left=1.2cm and 0.8cm of p] (k1) \{}\SpecialStringTok{$k\_\{1,i\}$}\NormalTok{\};}
  \FunctionTok{\textbackslash{}node}\NormalTok{[obs, left=0.3cm of k1] (n1) \{}\SpecialStringTok{$n\_\{1,i\}$}\NormalTok{\};}
  \FunctionTok{\textbackslash{}node}\NormalTok{[obs, below right=1.2cm and 0.8cm of p] (k2) \{}\SpecialStringTok{$k\_\{2,i\}$}\NormalTok{\};}
  \FunctionTok{\textbackslash{}node}\NormalTok{[obs, right=0.3cm of k2] (n2) \{}\SpecialStringTok{$n\_\{2,i\}$}\NormalTok{\};}

  \CommentTok{\% Beliefs}
  \FunctionTok{\textbackslash{}node}\NormalTok{[det, below=2.5cm of p, xshift={-}0.7cm] (ap) \{}\SpecialStringTok{$}\SpecialCharTok{\textbackslash{}alpha}\SpecialStringTok{\_\{post,i\}$}\NormalTok{\};}
  \FunctionTok{\textbackslash{}node}\NormalTok{[det, below=2.5cm of p, xshift=0.7cm] (bp) \{}\SpecialStringTok{$}\SpecialCharTok{\textbackslash{}beta}\SpecialStringTok{\_\{post,i\}$}\NormalTok{\};}

  \CommentTok{\% Action}
  \FunctionTok{\textbackslash{}node}\NormalTok{[obs, below=1cm of ap, xshift=0.7cm] (y) \{}\SpecialStringTok{$y\_i$}\NormalTok{\};}

  \CommentTok{\% Edges}
  \FunctionTok{\textbackslash{}edge}\NormalTok{ \{alpha0, beta0\} \{ap, bp\};}
  \FunctionTok{\textbackslash{}edge}\NormalTok{ \{p, k1, n1, k2, n2\} \{ap, bp\};}
  \FunctionTok{\textbackslash{}edge}\NormalTok{ \{ap, bp\} \{y\};}

  \FunctionTok{\textbackslash{}plate}\NormalTok{ \{p1\} \{(k1)(n1)(k2)(n2)(ap)(bp)(y)\} \{}\SpecialStringTok{$N$}\NormalTok{\};}
\KeywordTok{\textbackslash{}end}\NormalTok{\{}\ExtensionTok{tikzpicture}\NormalTok{\}}
\end{Highlighting}
\end{Shaded}

\begin{figure}[H]

{\centering \includegraphics[width=0.6\linewidth,height=\textheight,keepaspectratio]{10-BayesianModels_files/figure-pdf/pba-dag-1.pdf}

}

\caption{DAG for the Proportional Bayesian Agent}

\end{figure}%

The generative relationships for the PBA are: \[
\begin{aligned}
p &\sim \text{Beta}(2, 2) \\
\alpha_{post,i} &= \alpha_0 + p k_{1,i} + (1-p) k_{2,i} \\
\beta_{post,i}  &= \beta_0 + p(n_{1,i} - k_{1,i}) + (1-p)(n_{2,i} - k_{2,i}) \\
y_i &\sim \text{Beta-Binomial}(1, \alpha_{post,i}, \beta_{post,i})
\end{aligned}
\]

\begin{Shaded}
\begin{Highlighting}[]
\NormalTok{ProportionalAgent\_stan }\OtherTok{\textless{}{-}} \StringTok{"}
\StringTok{// Proportional Bayesian Agent (PBA).}
\StringTok{// p in [0,1] allocates the unit evidence budget between direct and social.}
\StringTok{// p = 0.5 approximates balanced weighting; p {-}\textgreater{} 1 ignores social; p {-}\textgreater{} 0 ignores direct.}
\StringTok{data \{}
\StringTok{  int\textless{}lower=1\textgreater{} N\_total;}
\StringTok{  array[N\_total] int\textless{}lower=0, upper=1\textgreater{} y;}
\StringTok{  array[N\_total] int\textless{}lower=0\textgreater{} blue1;}
\StringTok{  array[N\_total] int\textless{}lower=0\textgreater{} total1;}
\StringTok{  array[N\_total] int\textless{}lower=0\textgreater{} blue2;}
\StringTok{  array[N\_total] int\textless{}lower=0\textgreater{} total2;}
\StringTok{  }
\StringTok{  real prior\_p\_alpha;}
\StringTok{  real prior\_p\_beta;}
\StringTok{  int\textless{}lower=0, upper=1\textgreater{} run\_diagnostics;}
\StringTok{\}}

\StringTok{parameters \{}
\StringTok{  real\textless{}lower=0, upper=1\textgreater{} p;  // Allocation to direct evidence}
\StringTok{\}}

\StringTok{model \{}
\StringTok{  // prior}
\StringTok{  target += beta\_lpdf(p | prior\_p\_alpha, prior\_p\_beta);}

\StringTok{  // Vectorized likelihood}
\StringTok{  vector[N\_total] alpha\_post = 0.5 + p * to\_vector(blue1) + (1.0 {-} p) * to\_vector(blue2);}
\StringTok{  vector[N\_total] beta\_post  = 0.5 + p * (to\_vector(total1) {-} to\_vector(blue1))}
\StringTok{                                   + (1.0 {-} p) * (to\_vector(total2) {-} to\_vector(blue2));}
\StringTok{                             }
\StringTok{  target += beta\_binomial\_lpmf(y | 1, alpha\_post, beta\_post);}
\StringTok{\}}

\StringTok{generated quantities \{}
\StringTok{  vector[N\_total] log\_lik;}
\StringTok{  array[N\_total] int y\_rep;}
\StringTok{  real lprior = beta\_lpdf(p | prior\_p\_alpha, prior\_p\_beta);}

\StringTok{  if (run\_diagnostics) \{}
\StringTok{    for (i in 1:N\_total) \{}
\StringTok{      real alpha\_post = 0.5 + p * blue1[i] + (1.0 {-} p) * blue2[i];}
\StringTok{      real beta\_post  = 0.5 + p * (total1[i] {-} blue1[i]) }
\StringTok{                           + (1.0 {-} p) * (total2[i] {-} blue2[i]);}

\StringTok{      log\_lik[i] = beta\_binomial\_lpmf(y[i] | 1, alpha\_post, beta\_post);}
\StringTok{      y\_rep[i]   = beta\_binomial\_rng(1, alpha\_post, beta\_post);}
\StringTok{    \}}
\StringTok{  \}}
\StringTok{\}}
\StringTok{"}

\FunctionTok{write\_stan\_file}\NormalTok{(ProportionalAgent\_stan,}
                \AttributeTok{dir =} \StringTok{"stan/"}\NormalTok{,}
                \AttributeTok{basename =} \StringTok{"ch10\_proportional\_agent.stan"}\NormalTok{)}
\end{Highlighting}
\end{Shaded}

\begin{verbatim}
[1] "/Users/au209589/Dropbox/Teaching/AdvancedCognitiveModeling23_book/stan/ch10_proportional_agent.stan"
\end{verbatim}

\subsection{Weighted Bayesian Agent
(Reparameterised)}\label{weighted-bayesian-agent-reparameterised}

\begin{Shaded}
\begin{Highlighting}[]
\FunctionTok{\textbackslash{}usetikzlibrary}\NormalTok{\{bayesnet\}}
\KeywordTok{\textbackslash{}begin}\NormalTok{\{}\ExtensionTok{tikzpicture}\NormalTok{\}}
  \CommentTok{\% Fixed Prior}
  \FunctionTok{\textbackslash{}node}\NormalTok{[const] (alpha0) \{}\SpecialStringTok{$}\SpecialCharTok{\textbackslash{}alpha}\SpecialStringTok{\_0$}\NormalTok{\};}
  \FunctionTok{\textbackslash{}node}\NormalTok{[const, right=1cm of alpha0] (beta0) \{}\SpecialStringTok{$}\SpecialCharTok{\textbackslash{}beta}\SpecialStringTok{\_0$}\NormalTok{\};}

  \CommentTok{\% Latent Parameters}
  \FunctionTok{\textbackslash{}node}\NormalTok{[latent, below=0.8cm of alpha0] (rho) \{}\SpecialStringTok{$}\SpecialCharTok{\textbackslash{}rho}\SpecialStringTok{$}\NormalTok{\}; }
  \FunctionTok{\textbackslash{}node}\NormalTok{[latent, below=0.8cm of beta0] (kappa) \{}\SpecialStringTok{$}\SpecialCharTok{\textbackslash{}kappa}\SpecialStringTok{$}\NormalTok{\};}

  \CommentTok{\% Transformed Weights}
  \FunctionTok{\textbackslash{}node}\NormalTok{[det, below=1cm of rho, xshift={-}0.5cm] (wd) \{}\SpecialStringTok{$w\_d$}\NormalTok{\};}
  \FunctionTok{\textbackslash{}node}\NormalTok{[det, below=1cm of kappa, xshift=0.5cm] (ws) \{}\SpecialStringTok{$w\_s$}\NormalTok{\};}

  \CommentTok{\% Inputs}
  \FunctionTok{\textbackslash{}node}\NormalTok{[obs, left=1cm of wd] (k1) \{}\SpecialStringTok{$k\_\{1,i\}$}\NormalTok{\};}
  \FunctionTok{\textbackslash{}node}\NormalTok{[obs, left=0.3cm of k1] (n1) \{}\SpecialStringTok{$n\_\{1,i\}$}\NormalTok{\};}
  \FunctionTok{\textbackslash{}node}\NormalTok{[obs, right=1cm of ws] (k2) \{}\SpecialStringTok{$k\_\{2,i\}$}\NormalTok{\};}
  \FunctionTok{\textbackslash{}node}\NormalTok{[obs, right=0.3cm of k2] (n2) \{}\SpecialStringTok{$n\_\{2,i\}$}\NormalTok{\};}

  \CommentTok{\% Beliefs}
  \FunctionTok{\textbackslash{}node}\NormalTok{[det, below=2cm of wd] (ap) \{}\SpecialStringTok{$}\SpecialCharTok{\textbackslash{}alpha}\SpecialStringTok{\_\{post,i\}$}\NormalTok{\};}
  \FunctionTok{\textbackslash{}node}\NormalTok{[det, below=2cm of ws] (bp) \{}\SpecialStringTok{$}\SpecialCharTok{\textbackslash{}beta}\SpecialStringTok{\_\{post,i\}$}\NormalTok{\};}

  \CommentTok{\% Action}
  \FunctionTok{\textbackslash{}node}\NormalTok{[obs, below=1.2cm of ap, xshift=1cm] (y) \{}\SpecialStringTok{$y\_i$}\NormalTok{\};}

  \CommentTok{\% Edges}
  \FunctionTok{\textbackslash{}edge}\NormalTok{ \{alpha0, beta0\} \{ap, bp\};}
  \FunctionTok{\textbackslash{}edge}\NormalTok{ \{rho, kappa\} \{wd, ws\};}
  \FunctionTok{\textbackslash{}edge}\NormalTok{ \{wd, k1, n1, ws, k2, n2\} \{ap, bp\};}
  \FunctionTok{\textbackslash{}edge}\NormalTok{ \{ap, bp\} \{y\};}

  \FunctionTok{\textbackslash{}plate}\NormalTok{ \{p1\} \{(k1)(n1)(k2)(n2)(ap)(bp)(y)\} \{}\SpecialStringTok{$N$}\NormalTok{\};}
\KeywordTok{\textbackslash{}end}\NormalTok{\{}\ExtensionTok{tikzpicture}\NormalTok{\}}
\end{Highlighting}
\end{Shaded}

\begin{figure}[H]

{\centering \includegraphics[width=0.6\linewidth,height=\textheight,keepaspectratio]{10-BayesianModels_files/figure-pdf/wba-dag-1.pdf}

}

\caption{DAG for the Weighted Bayesian Agent (Reparameterised)}

\end{figure}%

The generative relationships for the reparameterized WBA are: \[
\begin{aligned}
\rho &\sim \text{Beta}(2, 2) \\
\kappa &\sim \text{Lognormal}(\log(2), 0.5) \\
w_d &= \rho \cdot \kappa \\
w_s &= (1 - \rho) \cdot \kappa \\
\alpha_{post,i} &= \alpha_0 + w_d k_{1,i} + w_s k_{2,i} \\
\beta_{post,i}  &= \beta_0 + w_d(n_{1,i} - k_{1,i}) + w_s(n_{2,i} - k_{2,i}) \\
y_i &\sim \text{Beta-Binomial}(1, \alpha_{post,i}, \beta_{post,i})
\end{aligned}
\]

\begin{Shaded}
\begin{Highlighting}[]
\NormalTok{WeightedAgent\_stan }\OtherTok{\textless{}{-}} \StringTok{"}
\StringTok{// Weighted Bayesian Agent (reparameterised).}
\StringTok{// rho in (0,1): relative weight of direct vs social evidence.}
\StringTok{// kappa \textgreater{} 0: total evidence scaling (w\_d + w\_s).}
\StringTok{data \{}
\StringTok{  int\textless{}lower=1\textgreater{} N\_total;}
\StringTok{  array[N\_total] int\textless{}lower=0, upper=1\textgreater{} y;}
\StringTok{  array[N\_total] int\textless{}lower=0\textgreater{} blue1;}
\StringTok{  array[N\_total] int\textless{}lower=0\textgreater{} total1;}
\StringTok{  array[N\_total] int\textless{}lower=0\textgreater{} blue2;}
\StringTok{  array[N\_total] int\textless{}lower=0\textgreater{} total2;}
\StringTok{  }
\StringTok{  real prior\_rho\_alpha;}
\StringTok{  real prior\_rho\_beta;}
\StringTok{  real prior\_kappa\_mu;}
\StringTok{  real\textless{}lower=0\textgreater{} prior\_kappa\_sigma;}
\StringTok{  int\textless{}lower=0, upper=1\textgreater{} run\_diagnostics;}
\StringTok{\}}

\StringTok{parameters \{}
\StringTok{  real\textless{}lower=0, upper=1\textgreater{} rho;    // relative weight: w\_d / (w\_d + w\_s)}
\StringTok{  real\textless{}lower=0\textgreater{}          kappa;  // total weight: w\_d + w\_s}
\StringTok{\}}

\StringTok{transformed parameters \{}
\StringTok{  real\textless{}lower=0\textgreater{} weight\_direct = rho * kappa;}
\StringTok{  real\textless{}lower=0\textgreater{} weight\_social = (1.0 {-} rho) * kappa;}
\StringTok{\}}

\StringTok{model \{}
\StringTok{  // rho: weakly centred on equal weighting}
\StringTok{  target += beta\_lpdf(rho | prior\_rho\_alpha, prior\_rho\_beta);}
\StringTok{  // kappa: lognormal centered on 2 (SBA equivalent)}
\StringTok{  target += lognormal\_lpdf(kappa | prior\_kappa\_mu, prior\_kappa\_sigma);}

\StringTok{  // Vectorized likelihood}
\StringTok{  vector[N\_total] alpha\_post = 0.5 + weight\_direct * to\_vector(blue1)}
\StringTok{                                   + weight\_social * to\_vector(blue2);}
\StringTok{  vector[N\_total] beta\_post  = 0.5 + weight\_direct * (to\_vector(total1) {-} to\_vector(blue1))}
\StringTok{                                   + weight\_social * (to\_vector(total2) {-} to\_vector(blue2));}
\StringTok{                             }
\StringTok{  target += beta\_binomial\_lpmf(y | 1, alpha\_post, beta\_post);}
\StringTok{\}}

\StringTok{generated quantities \{}
\StringTok{  vector[N\_total] log\_lik;}
\StringTok{  array[N\_total] int y\_rep;}
\StringTok{  real lprior = beta\_lpdf(rho | prior\_rho\_alpha, prior\_rho\_beta) + }
\StringTok{                lognormal\_lpdf(kappa | prior\_kappa\_mu, prior\_kappa\_sigma);}

\StringTok{  if (run\_diagnostics) \{}
\StringTok{    for (i in 1:N\_total) \{}
\StringTok{      real alpha\_post = 0.5 + weight\_direct * blue1[i] + weight\_social * blue2[i];}
\StringTok{      real beta\_post  = 0.5 + weight\_direct * (total1[i] {-} blue1[i]) }
\StringTok{                           + weight\_social * (total2[i] {-} blue2[i]);}

\StringTok{      log\_lik[i] = beta\_binomial\_lpmf(y[i] | 1, alpha\_post, beta\_post);}
\StringTok{      y\_rep[i]   = beta\_binomial\_rng(1, alpha\_post, beta\_post);}
\StringTok{    \}}
\StringTok{  \}}
\StringTok{\}}
\StringTok{"}

\FunctionTok{write\_stan\_file}\NormalTok{(WeightedAgent\_stan,}
                \AttributeTok{dir =} \StringTok{"stan/"}\NormalTok{,}
                \AttributeTok{basename =} \StringTok{"ch10\_weighted\_agent.stan"}\NormalTok{)}
\end{Highlighting}
\end{Shaded}

\begin{verbatim}
[1] "/Users/au209589/Dropbox/Teaching/AdvancedCognitiveModeling23_book/stan/ch10_weighted_agent.stan"
\end{verbatim}

\section{Fitting the Reparameterised
Models}\label{fitting-the-reparameterised-models}

\begin{Shaded}
\begin{Highlighting}[]
\NormalTok{mod\_proportional }\OtherTok{\textless{}{-}} \FunctionTok{cmdstan\_model}\NormalTok{(}\StringTok{"stan/ch10\_proportional\_agent.stan"}\NormalTok{, }\AttributeTok{dir =} \StringTok{"simmodels"}\NormalTok{)}
\NormalTok{mod\_weighted     }\OtherTok{\textless{}{-}} \FunctionTok{cmdstan\_model}\NormalTok{(}\StringTok{"stan/ch10\_weighted\_agent.stan"}\NormalTok{, }\AttributeTok{dir =} \StringTok{"simmodels"}\NormalTok{)}
\end{Highlighting}
\end{Shaded}

\subsection{Prepare PBA Data}\label{prepare-pba-data}

The PBA uses normalized weights that sum to 1, so we need to convert our
scenarios:

\begin{Shaded}
\begin{Highlighting}[]
\CommentTok{\# For PBA, p represents w\_d / (w\_d + w\_s)}
\CommentTok{\# Self{-}Focused: 1.5 / 2.0 = 0.75}
\CommentTok{\# Balanced: 1.0 / 2.0 = 0.5  }
\CommentTok{\# Socially{-}Influenced: 0.5 / 2.5 = 0.2}

\NormalTok{true\_pba }\OtherTok{\textless{}{-}} \FunctionTok{list}\NormalTok{(}
  \AttributeTok{Self\_Focused        =} \FloatTok{0.75}\NormalTok{,}
  \AttributeTok{Balanced            =} \FloatTok{0.5}\NormalTok{,}
  \AttributeTok{Socially\_Influenced =} \FloatTok{0.2}
\NormalTok{)}

\CommentTok{\# Generate simulated data for each scenario}
\NormalTok{stan\_data\_pba }\OtherTok{\textless{}{-}} \FunctionTok{map}\NormalTok{(true\_pba, }\ControlFlowTok{function}\NormalTok{(p) \{}
\NormalTok{  sim }\OtherTok{\textless{}{-}} \FunctionTok{generate\_agent\_decisions}\NormalTok{(p, }\DecValTok{1} \SpecialCharTok{{-}}\NormalTok{ p, evidence\_combinations, }\AttributeTok{n\_samples =} \DecValTok{5}\NormalTok{) }\SpecialCharTok{|\textgreater{}}
    \FunctionTok{mutate}\NormalTok{(}\AttributeTok{total1 =} \DecValTok{8}\NormalTok{L, }\AttributeTok{total2 =} \DecValTok{3}\NormalTok{L)}
  
\NormalTok{  d }\OtherTok{\textless{}{-}} \FunctionTok{prepare\_stan\_data}\NormalTok{(sim)}
\NormalTok{  d}\SpecialCharTok{$}\NormalTok{prior\_p\_alpha }\OtherTok{=} \DecValTok{2}
\NormalTok{  d}\SpecialCharTok{$}\NormalTok{prior\_p\_beta  }\OtherTok{=} \DecValTok{2}
\NormalTok{  d}
\NormalTok{\})}
\end{Highlighting}
\end{Shaded}

\subsection{Prepare WBA Reparameterised
Data}\label{prepare-wba-reparameterised-data}

\begin{Shaded}
\begin{Highlighting}[]
\NormalTok{true\_wba\_reparam }\OtherTok{\textless{}{-}} \FunctionTok{list}\NormalTok{(}
  \AttributeTok{Self\_Focused        =} \FunctionTok{list}\NormalTok{(}\AttributeTok{rho =} \FloatTok{0.75}\NormalTok{, }\AttributeTok{kappa =} \FloatTok{2.0}\NormalTok{, }\AttributeTok{wd =} \FloatTok{1.5}\NormalTok{, }\AttributeTok{ws =} \FloatTok{0.5}\NormalTok{),}
  \AttributeTok{Balanced            =} \FunctionTok{list}\NormalTok{(}\AttributeTok{rho =} \FloatTok{0.50}\NormalTok{, }\AttributeTok{kappa =} \FloatTok{2.0}\NormalTok{, }\AttributeTok{wd =} \FloatTok{1.0}\NormalTok{, }\AttributeTok{ws =} \FloatTok{1.0}\NormalTok{),}
  \AttributeTok{Socially\_Influenced =} \FunctionTok{list}\NormalTok{(}\AttributeTok{rho =} \FloatTok{0.20}\NormalTok{, }\AttributeTok{kappa =} \FloatTok{2.5}\NormalTok{, }\AttributeTok{wd =} \FloatTok{0.5}\NormalTok{, }\AttributeTok{ws =} \FloatTok{2.0}\NormalTok{)}
\NormalTok{)}

\CommentTok{\# Generate simulated data for each scenario}
\NormalTok{stan\_data\_wba\_reparam }\OtherTok{\textless{}{-}} \FunctionTok{map}\NormalTok{(true\_wba\_reparam, }\ControlFlowTok{function}\NormalTok{(params) \{}
\NormalTok{  sim }\OtherTok{\textless{}{-}} \FunctionTok{generate\_agent\_decisions}\NormalTok{(params}\SpecialCharTok{$}\NormalTok{wd, params}\SpecialCharTok{$}\NormalTok{ws,}
\NormalTok{                                  evidence\_combinations, }\AttributeTok{n\_samples =} \DecValTok{5}\NormalTok{) }\SpecialCharTok{|\textgreater{}}
    \FunctionTok{mutate}\NormalTok{(}\AttributeTok{total1 =} \DecValTok{8}\NormalTok{L, }\AttributeTok{total2 =} \DecValTok{3}\NormalTok{L)}

\NormalTok{  d }\OtherTok{\textless{}{-}} \FunctionTok{prepare\_stan\_data}\NormalTok{(sim)}
\NormalTok{  d}\SpecialCharTok{$}\NormalTok{prior\_rho\_alpha     }\OtherTok{=} \DecValTok{2}
\NormalTok{  d}\SpecialCharTok{$}\NormalTok{prior\_rho\_beta      }\OtherTok{=} \DecValTok{2}
\NormalTok{  d}\SpecialCharTok{$}\NormalTok{prior\_kappa\_mu      }\OtherTok{=} \FunctionTok{log}\NormalTok{(}\DecValTok{2}\NormalTok{)}
\NormalTok{  d}\SpecialCharTok{$}\NormalTok{prior\_kappa\_sigma   }\OtherTok{=} \FloatTok{0.5}
\NormalTok{  d}
\NormalTok{\})}
\end{Highlighting}
\end{Shaded}

\subsection{Fit All Models}\label{fit-all-models}

\begin{Shaded}
\begin{Highlighting}[]
\ControlFlowTok{if}\NormalTok{ (regenerate\_simulations }\SpecialCharTok{||} \SpecialCharTok{!}\FunctionTok{all}\NormalTok{(}\FunctionTok{file.exists}\NormalTok{(}\FunctionTok{sprintf}\NormalTok{(}\StringTok{"simmodels/ch10\_fit\_pba\_\%s.rds"}\NormalTok{, }\FunctionTok{names}\NormalTok{(stan\_data\_pba))))) \{}
  \CommentTok{\# {-}{-}{-} Fit PBA {-}{-}{-}}
\NormalTok{  fit\_pba }\OtherTok{\textless{}{-}} \FunctionTok{map2}\NormalTok{(stan\_data\_pba, }\DecValTok{201}\SpecialCharTok{:}\DecValTok{203}\NormalTok{, }\ControlFlowTok{function}\NormalTok{(data, seed) \{}
    \FunctionTok{do.call}\NormalTok{(mod\_proportional}\SpecialCharTok{$}\NormalTok{sample,}
            \FunctionTok{c}\NormalTok{(}\FunctionTok{list}\NormalTok{(}\AttributeTok{data =}\NormalTok{ data, }\AttributeTok{seed =}\NormalTok{ seed), sampler\_args))}
\NormalTok{  \})}
  \FunctionTok{names}\NormalTok{(fit\_pba) }\OtherTok{\textless{}{-}} \FunctionTok{names}\NormalTok{(stan\_data\_pba)}
  \FunctionTok{iwalk}\NormalTok{(fit\_pba, }\SpecialCharTok{\textasciitilde{}}\NormalTok{ .x}\SpecialCharTok{$}\FunctionTok{save\_object}\NormalTok{(}
    \FunctionTok{sprintf}\NormalTok{(}\StringTok{"simmodels/ch10\_fit\_pba\_\%s.rds"}\NormalTok{, .y)}
\NormalTok{  ))}

  \CommentTok{\# {-}{-}{-} Fit WBA reparameterised {-}{-}{-}}
\NormalTok{  fit\_wba }\OtherTok{\textless{}{-}} \FunctionTok{map2}\NormalTok{(stan\_data\_wba\_reparam, }\DecValTok{204}\SpecialCharTok{:}\DecValTok{206}\NormalTok{, }\ControlFlowTok{function}\NormalTok{(data, seed) \{}
    \FunctionTok{do.call}\NormalTok{(mod\_weighted}\SpecialCharTok{$}\NormalTok{sample,}
            \FunctionTok{c}\NormalTok{(}\FunctionTok{list}\NormalTok{(}\AttributeTok{data =}\NormalTok{ data, }\AttributeTok{seed =}\NormalTok{ seed), sampler\_args))}
\NormalTok{  \})}
  \FunctionTok{names}\NormalTok{(fit\_wba) }\OtherTok{\textless{}{-}} \FunctionTok{names}\NormalTok{(stan\_data\_wba\_reparam)}
  \FunctionTok{iwalk}\NormalTok{(fit\_wba, }\SpecialCharTok{\textasciitilde{}}\NormalTok{ .x}\SpecialCharTok{$}\FunctionTok{save\_object}\NormalTok{(}
    \FunctionTok{sprintf}\NormalTok{(}\StringTok{"simmodels/ch10\_fit\_wba\_\%s.rds"}\NormalTok{, .y)}
\NormalTok{  ))}
\NormalTok{\} }\ControlFlowTok{else}\NormalTok{ \{}
\NormalTok{  fit\_pba }\OtherTok{\textless{}{-}} \FunctionTok{map}\NormalTok{(}
    \FunctionTok{names}\NormalTok{(stan\_data\_pba),}
    \SpecialCharTok{\textasciitilde{}} \FunctionTok{readRDS}\NormalTok{(}\FunctionTok{sprintf}\NormalTok{(}\StringTok{"simmodels/ch10\_fit\_pba\_\%s.rds"}\NormalTok{, .x))}
\NormalTok{  ) }\SpecialCharTok{|\textgreater{}} \FunctionTok{set\_names}\NormalTok{(}\FunctionTok{names}\NormalTok{(stan\_data\_pba))}

\NormalTok{  fit\_wba }\OtherTok{\textless{}{-}} \FunctionTok{map}\NormalTok{(}
    \FunctionTok{names}\NormalTok{(stan\_data\_wba\_reparam),}
    \SpecialCharTok{\textasciitilde{}} \FunctionTok{readRDS}\NormalTok{(}\FunctionTok{sprintf}\NormalTok{(}\StringTok{"simmodels/ch10\_fit\_wba\_\%s.rds"}\NormalTok{, .x))}
\NormalTok{  ) }\SpecialCharTok{|\textgreater{}} \FunctionTok{set\_names}\NormalTok{(}\FunctionTok{names}\NormalTok{(stan\_data\_wba\_reparam))}
\NormalTok{\}}

\CommentTok{\# Check diagnostics}
\FunctionTok{cat}\NormalTok{(}\StringTok{"}\SpecialCharTok{\textbackslash{}n}\StringTok{=== Proportional Bayesian Agent ===}\SpecialCharTok{\textbackslash{}n}\StringTok{"}\NormalTok{)}
\end{Highlighting}
\end{Shaded}

\begin{verbatim}

=== Proportional Bayesian Agent ===
\end{verbatim}

\begin{Shaded}
\begin{Highlighting}[]
\FunctionTok{iwalk}\NormalTok{(fit\_pba, }\SpecialCharTok{\textasciitilde{}} \FunctionTok{check\_diagnostics}\NormalTok{(.x, }\FunctionTok{paste}\NormalTok{(}\StringTok{"PBA"}\NormalTok{, .y)))}
\end{Highlighting}
\end{Shaded}

\begin{verbatim}
PBA Self_Focused          | Divs:   0 | Min E-BFMI: 0.98 | Failures: 0
PBA Balanced              | Divs:   0 | Min E-BFMI: 0.95 | Failures: 0
PBA Socially_Influenced   | Divs:   0 | Min E-BFMI: 0.93 | Failures: 0
\end{verbatim}

\begin{Shaded}
\begin{Highlighting}[]
\FunctionTok{cat}\NormalTok{(}\StringTok{"}\SpecialCharTok{\textbackslash{}n}\StringTok{=== Weighted Bayesian Agent (Reparameterised) ===}\SpecialCharTok{\textbackslash{}n}\StringTok{"}\NormalTok{)}
\end{Highlighting}
\end{Shaded}

\begin{verbatim}

=== Weighted Bayesian Agent (Reparameterised) ===
\end{verbatim}

\begin{Shaded}
\begin{Highlighting}[]
\FunctionTok{iwalk}\NormalTok{(fit\_wba, }\SpecialCharTok{\textasciitilde{}} \FunctionTok{check\_diagnostics}\NormalTok{(.x, }\FunctionTok{paste}\NormalTok{(}\StringTok{"WBA"}\NormalTok{, .y)))}
\end{Highlighting}
\end{Shaded}

\begin{verbatim}
WBA Self_Focused          | Divs:   0 | Min E-BFMI: 0.86 | Failures: 0
WBA Balanced              | Divs:   0 | Min E-BFMI: 1.04 | Failures: 0
WBA Socially_Influenced   | Divs:   0 | Min E-BFMI: 1.03 | Failures: 0
\end{verbatim}

\subsection{Quick predictive checks}\label{quick-predictive-checks}

\begin{Shaded}
\begin{Highlighting}[]
\CommentTok{\# Helper function to generate and combine Prior vs. Posterior checks dynamically}
\NormalTok{generate\_side\_by\_side\_ppc }\OtherTok{\textless{}{-}} \ControlFlowTok{function}\NormalTok{(fit\_list, data\_list, model\_title) \{}
  
  \CommentTok{\# Map over all three scenarios (Self\_Focused, Balanced, Socially\_Influenced)}
\NormalTok{  plot\_list }\OtherTok{\textless{}{-}} \FunctionTok{imap}\NormalTok{(fit\_list, }\ControlFlowTok{function}\NormalTok{(fit, scenario) \{}
    
    \CommentTok{\# 1. Extract and plot Prior Predictive}
\NormalTok{    prior\_data }\OtherTok{\textless{}{-}} \FunctionTok{extract\_predictive\_data}\NormalTok{(}
      \AttributeTok{fit =}\NormalTok{ fit,}
      \AttributeTok{stan\_data =}\NormalTok{ data\_list[[scenario]],}
      \AttributeTok{scenario\_name =} \StringTok{"Prior"}\NormalTok{,}
      \AttributeTok{target\_var =} \StringTok{"prior\_pred"}
\NormalTok{    )}
    
\NormalTok{    p\_prior }\OtherTok{\textless{}{-}} \FunctionTok{plot\_predictive\_check}\NormalTok{(}
\NormalTok{      prior\_data,}
      \AttributeTok{title =} \FunctionTok{paste}\NormalTok{(model\_title, }\StringTok{"{-}"}\NormalTok{, }\FunctionTok{str\_replace\_all}\NormalTok{(scenario, }\StringTok{"\_"}\NormalTok{, }\StringTok{" "}\NormalTok{)),}
      \AttributeTok{subtitle =} \StringTok{"Prior Predictive Check"}
\NormalTok{    ) }\SpecialCharTok{+} \FunctionTok{theme}\NormalTok{(}\AttributeTok{legend.position =} \StringTok{"none"}\NormalTok{) }\CommentTok{\# Hide individual legends to save space}
    
    \CommentTok{\# 2. Extract and plot Posterior Predictive}
\NormalTok{    post\_data }\OtherTok{\textless{}{-}} \FunctionTok{extract\_predictive\_data}\NormalTok{(}
      \AttributeTok{fit =}\NormalTok{ fit,}
      \AttributeTok{stan\_data =}\NormalTok{ data\_list[[scenario]],}
      \AttributeTok{scenario\_name =} \StringTok{"Posterior"}\NormalTok{,}
      \AttributeTok{target\_var =} \StringTok{"posterior\_pred"}
\NormalTok{    )}
    
\NormalTok{    p\_post }\OtherTok{\textless{}{-}} \FunctionTok{plot\_predictive\_check}\NormalTok{(}
\NormalTok{      post\_data,}
      \AttributeTok{title =} \FunctionTok{paste}\NormalTok{(model\_title, }\StringTok{"{-}"}\NormalTok{, }\FunctionTok{str\_replace\_all}\NormalTok{(scenario, }\StringTok{"\_"}\NormalTok{, }\StringTok{" "}\NormalTok{)),}
      \AttributeTok{subtitle =} \StringTok{"Posterior Predictive Check"}
\NormalTok{    ) }\SpecialCharTok{+} \FunctionTok{theme}\NormalTok{(}\AttributeTok{axis.title.y =} \FunctionTok{element\_blank}\NormalTok{()) }\CommentTok{\# Share Y axis visually with the prior plot}
    
    \CommentTok{\# 3. Combine side{-}by{-}side}
    \FunctionTok{return}\NormalTok{(p\_prior }\SpecialCharTok{|}\NormalTok{ p\_post)}
\NormalTok{  \})}
  
  \CommentTok{\# Stack the three scenarios vertically and collect the legend at the very bottom}
\NormalTok{  combined\_all }\OtherTok{\textless{}{-}} \FunctionTok{wrap\_plots}\NormalTok{(plot\_list, }\AttributeTok{ncol =} \DecValTok{1}\NormalTok{) }\SpecialCharTok{+} 
    \FunctionTok{plot\_layout}\NormalTok{(}\AttributeTok{guides =} \StringTok{"collect"}\NormalTok{) }\SpecialCharTok{\&} 
    \FunctionTok{theme}\NormalTok{(}\AttributeTok{legend.position =} \StringTok{"bottom"}\NormalTok{)}
  
  \FunctionTok{return}\NormalTok{(combined\_all)}
\NormalTok{\}}

\FunctionTok{cat}\NormalTok{(}\StringTok{"=== Proportional Bayesian Agent (PBA) Predictive Checks ===}\SpecialCharTok{\textbackslash{}n}\StringTok{"}\NormalTok{)}
\end{Highlighting}
\end{Shaded}

\begin{verbatim}
=== Proportional Bayesian Agent (PBA) Predictive Checks ===
\end{verbatim}

\begin{Shaded}
\begin{Highlighting}[]
\NormalTok{pba\_combined }\OtherTok{\textless{}{-}} \FunctionTok{generate\_side\_by\_side\_ppc}\NormalTok{(fit\_pba, stan\_data\_pba, }\StringTok{"PBA"}\NormalTok{)}
\FunctionTok{print}\NormalTok{(pba\_combined)}
\end{Highlighting}
\end{Shaded}

\begin{center}
\includegraphics[width=0.8\linewidth,height=\textheight,keepaspectratio]{10-BayesianModels_files/figure-pdf/ch10_reparam_ppc-1.pdf}
\end{center}

\begin{Shaded}
\begin{Highlighting}[]
\FunctionTok{cat}\NormalTok{(}\StringTok{"}\SpecialCharTok{\textbackslash{}n}\StringTok{=== Weighted Bayesian Agent (WBA) Predictive Checks ===}\SpecialCharTok{\textbackslash{}n}\StringTok{"}\NormalTok{)}
\end{Highlighting}
\end{Shaded}

\begin{verbatim}

=== Weighted Bayesian Agent (WBA) Predictive Checks ===
\end{verbatim}

\begin{Shaded}
\begin{Highlighting}[]
\CommentTok{\# Note: WBA uses stan\_data\_wba\_reparam}
\NormalTok{wba\_combined }\OtherTok{\textless{}{-}} \FunctionTok{generate\_side\_by\_side\_ppc}\NormalTok{(fit\_wba, stan\_data\_wba\_reparam, }\StringTok{"WBA"}\NormalTok{)}
\FunctionTok{print}\NormalTok{(wba\_combined)}
\end{Highlighting}
\end{Shaded}

\begin{center}
\includegraphics[width=0.8\linewidth,height=\textheight,keepaspectratio]{10-BayesianModels_files/figure-pdf/ch10_reparam_ppc-2.pdf}
\end{center}

Predictive checks look good for both models, with the posterior
predictions (right column) closely matching the observed data (black
points) across all scenarios. The priors (left column) are more diffuse,
as expected, but still cover the range of observed outcomes. This gives
us confidence that the models are fitting the data well and that we can
proceed to examine parameter recovery.

\section{Parameter Recovery: The Solution
Works!}\label{parameter-recovery-the-solution-works}

\subsection{PBA Recovery}\label{pba-recovery}

\begin{Shaded}
\begin{Highlighting}[]
\NormalTok{pba\_prior }\OtherTok{\textless{}{-}} \FunctionTok{tibble}\NormalTok{(}\AttributeTok{p =} \FunctionTok{rbeta}\NormalTok{(}\DecValTok{4000}\NormalTok{, }\DecValTok{2}\NormalTok{, }\DecValTok{2}\NormalTok{), }\AttributeTok{Scenario =} \StringTok{"Prior"}\NormalTok{)}

\NormalTok{pba\_posteriors }\OtherTok{\textless{}{-}} \FunctionTok{imap\_dfr}\NormalTok{(fit\_pba, }\ControlFlowTok{function}\NormalTok{(fit, name) \{}
  \FunctionTok{as\_draws\_df}\NormalTok{(fit}\SpecialCharTok{$}\FunctionTok{draws}\NormalTok{(}\StringTok{"p"}\NormalTok{)) }\SpecialCharTok{|\textgreater{}} 
    \FunctionTok{as\_tibble}\NormalTok{() }\SpecialCharTok{|\textgreater{}} 
    \FunctionTok{select}\NormalTok{(p) }\SpecialCharTok{|\textgreater{}} 
    \FunctionTok{mutate}\NormalTok{(}\AttributeTok{Scenario =}\NormalTok{ name)}
\NormalTok{\})}

\NormalTok{pba\_true\_df }\OtherTok{\textless{}{-}} \FunctionTok{enframe}\NormalTok{(true\_pba, }\AttributeTok{name =} \StringTok{"Scenario"}\NormalTok{, }\AttributeTok{value =} \StringTok{"true\_value"}\NormalTok{) }\SpecialCharTok{|\textgreater{}} 
  \FunctionTok{unnest}\NormalTok{(}\AttributeTok{cols =} \FunctionTok{c}\NormalTok{(true\_value))}

\FunctionTok{bind\_rows}\NormalTok{(pba\_prior, pba\_posteriors) }\SpecialCharTok{|\textgreater{}}
  \FunctionTok{mutate}\NormalTok{(}\AttributeTok{Scenario =} \FunctionTok{factor}\NormalTok{(Scenario, }
                           \AttributeTok{levels =} \FunctionTok{c}\NormalTok{(}\StringTok{"Prior"}\NormalTok{, }\StringTok{"Balanced"}\NormalTok{, }\StringTok{"Self\_Focused"}\NormalTok{, }\StringTok{"Socially\_Influenced"}\NormalTok{))) }\SpecialCharTok{|\textgreater{}}
  \FunctionTok{ggplot}\NormalTok{(}\FunctionTok{aes}\NormalTok{(}\AttributeTok{x =}\NormalTok{ p, }\AttributeTok{fill =}\NormalTok{ Scenario, }\AttributeTok{color =}\NormalTok{ Scenario)) }\SpecialCharTok{+}
  \FunctionTok{geom\_density}\NormalTok{(}\AttributeTok{alpha =} \FloatTok{0.3}\NormalTok{, }\AttributeTok{linewidth =} \DecValTok{1}\NormalTok{) }\SpecialCharTok{+}
  \FunctionTok{geom\_vline}\NormalTok{(}\AttributeTok{data =}\NormalTok{ pba\_true\_df, }
             \FunctionTok{aes}\NormalTok{(}\AttributeTok{xintercept =}\NormalTok{ true\_value, }\AttributeTok{color =}\NormalTok{ Scenario), }
             \AttributeTok{linetype =} \StringTok{"dashed"}\NormalTok{, }\AttributeTok{linewidth =} \FloatTok{1.2}\NormalTok{, }\AttributeTok{show.legend =} \ConstantTok{FALSE}\NormalTok{,}
             \AttributeTok{inherit.aes =} \ConstantTok{FALSE}\NormalTok{) }\SpecialCharTok{+}
  \FunctionTok{scale\_fill\_manual}\NormalTok{(}\AttributeTok{values =} \FunctionTok{c}\NormalTok{(}\StringTok{"Prior"} \OtherTok{=} \StringTok{"grey80"}\NormalTok{, }\StringTok{"Balanced"} \OtherTok{=} \StringTok{"\#4daf4a"}\NormalTok{, }
                               \StringTok{"Self\_Focused"} \OtherTok{=} \StringTok{"\#e41a1c"}\NormalTok{, }\StringTok{"Socially\_Influenced"} \OtherTok{=} \StringTok{"\#377eb8"}\NormalTok{)) }\SpecialCharTok{+}
  \FunctionTok{scale\_color\_manual}\NormalTok{(}\AttributeTok{values =} \FunctionTok{c}\NormalTok{(}\StringTok{"Prior"} \OtherTok{=} \StringTok{"grey50"}\NormalTok{, }\StringTok{"Balanced"} \OtherTok{=} \StringTok{"\#4daf4a"}\NormalTok{, }
                                \StringTok{"Self\_Focused"} \OtherTok{=} \StringTok{"\#e41a1c"}\NormalTok{, }\StringTok{"Socially\_Influenced"} \OtherTok{=} \StringTok{"\#377eb8"}\NormalTok{)) }\SpecialCharTok{+}
  \FunctionTok{coord\_cartesian}\NormalTok{(}\AttributeTok{xlim =} \FunctionTok{c}\NormalTok{(}\DecValTok{0}\NormalTok{, }\DecValTok{1}\NormalTok{)) }\SpecialCharTok{+}
  \FunctionTok{labs}\NormalTok{(}
    \AttributeTok{title =} \StringTok{"✓ PBA Parameter Recovery"}\NormalTok{, }
    \AttributeTok{x =} \StringTok{"p (Allocation to direct evidence)"}\NormalTok{, }
    \AttributeTok{y =} \StringTok{"Density"}
\NormalTok{  ) }\SpecialCharTok{+}
  \FunctionTok{theme}\NormalTok{(}\AttributeTok{legend.title =} \FunctionTok{element\_blank}\NormalTok{())}
\end{Highlighting}
\end{Shaded}

\begin{center}
\includegraphics[width=0.8\linewidth,height=\textheight,keepaspectratio]{10-BayesianModels_files/figure-pdf/ch10_pba_recovery-1.pdf}
\end{center}

This is not perfect, but pretty good, perhaps overregularized by the
prior (pulled towards the center compared to true values).

\subsection{WBA Reparameterised
Recovery}\label{wba-reparameterised-recovery}

\begin{Shaded}
\begin{Highlighting}[]
\NormalTok{wba\_prior }\OtherTok{\textless{}{-}} \FunctionTok{tibble}\NormalTok{(}
  \AttributeTok{rho =} \FunctionTok{rbeta}\NormalTok{(}\DecValTok{4000}\NormalTok{, }\DecValTok{2}\NormalTok{, }\DecValTok{2}\NormalTok{), }
  \AttributeTok{kappa =} \FunctionTok{rlnorm}\NormalTok{(}\DecValTok{4000}\NormalTok{, }\FunctionTok{log}\NormalTok{(}\DecValTok{2}\NormalTok{), }\FloatTok{0.5}\NormalTok{), }
  \AttributeTok{Scenario =} \StringTok{"Prior"}
\NormalTok{)}

\NormalTok{wba\_posteriors }\OtherTok{\textless{}{-}} \FunctionTok{imap\_dfr}\NormalTok{(fit\_wba, }\ControlFlowTok{function}\NormalTok{(fit, name) \{}
  \FunctionTok{as\_draws\_df}\NormalTok{(fit}\SpecialCharTok{$}\FunctionTok{draws}\NormalTok{(}\FunctionTok{c}\NormalTok{(}\StringTok{"rho"}\NormalTok{, }\StringTok{"kappa"}\NormalTok{))) }\SpecialCharTok{|\textgreater{}} 
    \FunctionTok{as\_tibble}\NormalTok{() }\SpecialCharTok{|\textgreater{}} 
    \FunctionTok{select}\NormalTok{(rho, kappa) }\SpecialCharTok{|\textgreater{}} 
    \FunctionTok{mutate}\NormalTok{(}\AttributeTok{Scenario =}\NormalTok{ name)}
\NormalTok{\})}

\NormalTok{wba\_true\_df }\OtherTok{\textless{}{-}} \FunctionTok{map\_dfr}\NormalTok{(true\_wba\_reparam, }\SpecialCharTok{\textasciitilde{}} \FunctionTok{as\_tibble}\NormalTok{(.x[}\FunctionTok{c}\NormalTok{(}\StringTok{"rho"}\NormalTok{, }\StringTok{"kappa"}\NormalTok{)]), }\AttributeTok{.id =} \StringTok{"Scenario"}\NormalTok{) }\SpecialCharTok{|\textgreater{}}
  \FunctionTok{pivot\_longer}\NormalTok{(}\FunctionTok{c}\NormalTok{(rho, kappa), }\AttributeTok{names\_to =} \StringTok{"parameter"}\NormalTok{, }\AttributeTok{values\_to =} \StringTok{"true\_value"}\NormalTok{)}

\FunctionTok{bind\_rows}\NormalTok{(wba\_prior, wba\_posteriors) }\SpecialCharTok{|\textgreater{}}
  \FunctionTok{pivot\_longer}\NormalTok{(}\FunctionTok{c}\NormalTok{(rho, kappa), }\AttributeTok{names\_to =} \StringTok{"parameter"}\NormalTok{, }\AttributeTok{values\_to =} \StringTok{"value"}\NormalTok{) }\SpecialCharTok{|\textgreater{}}
  \FunctionTok{mutate}\NormalTok{(}\AttributeTok{Scenario =} \FunctionTok{factor}\NormalTok{(Scenario, }
                           \AttributeTok{levels =} \FunctionTok{c}\NormalTok{(}\StringTok{"Prior"}\NormalTok{, }\StringTok{"Balanced"}\NormalTok{, }\StringTok{"Self\_Focused"}\NormalTok{, }\StringTok{"Socially\_Influenced"}\NormalTok{))) }\SpecialCharTok{|\textgreater{}}
  \FunctionTok{ggplot}\NormalTok{(}\FunctionTok{aes}\NormalTok{(}\AttributeTok{x =}\NormalTok{ value, }\AttributeTok{fill =}\NormalTok{ Scenario, }\AttributeTok{color =}\NormalTok{ Scenario)) }\SpecialCharTok{+}
  \FunctionTok{geom\_density}\NormalTok{(}\AttributeTok{alpha =} \FloatTok{0.3}\NormalTok{, }\AttributeTok{linewidth =} \DecValTok{1}\NormalTok{) }\SpecialCharTok{+}
  \FunctionTok{geom\_vline}\NormalTok{(}\AttributeTok{data =}\NormalTok{ wba\_true\_df, }
             \FunctionTok{aes}\NormalTok{(}\AttributeTok{xintercept =}\NormalTok{ true\_value, }\AttributeTok{color =}\NormalTok{ Scenario), }
             \AttributeTok{linetype =} \StringTok{"dashed"}\NormalTok{, }\AttributeTok{linewidth =} \FloatTok{1.2}\NormalTok{, }\AttributeTok{show.legend =} \ConstantTok{FALSE}\NormalTok{,}
             \AttributeTok{inherit.aes =} \ConstantTok{FALSE}\NormalTok{) }\SpecialCharTok{+}
  \FunctionTok{facet\_wrap}\NormalTok{(}\SpecialCharTok{\textasciitilde{}}\NormalTok{ parameter, }\AttributeTok{scales =} \StringTok{"free"}\NormalTok{) }\SpecialCharTok{+}
  \FunctionTok{scale\_fill\_manual}\NormalTok{(}\AttributeTok{values =} \FunctionTok{c}\NormalTok{(}\StringTok{"Prior"} \OtherTok{=} \StringTok{"grey80"}\NormalTok{, }\StringTok{"Balanced"} \OtherTok{=} \StringTok{"\#4daf4a"}\NormalTok{, }
                               \StringTok{"Self\_Focused"} \OtherTok{=} \StringTok{"\#e41a1c"}\NormalTok{, }\StringTok{"Socially\_Influenced"} \OtherTok{=} \StringTok{"\#377eb8"}\NormalTok{)) }\SpecialCharTok{+}
  \FunctionTok{scale\_color\_manual}\NormalTok{(}\AttributeTok{values =} \FunctionTok{c}\NormalTok{(}\StringTok{"Prior"} \OtherTok{=} \StringTok{"grey50"}\NormalTok{, }\StringTok{"Balanced"} \OtherTok{=} \StringTok{"\#4daf4a"}\NormalTok{, }
                                \StringTok{"Self\_Focused"} \OtherTok{=} \StringTok{"\#e41a1c"}\NormalTok{, }\StringTok{"Socially\_Influenced"} \OtherTok{=} \StringTok{"\#377eb8"}\NormalTok{)) }\SpecialCharTok{+}
  \FunctionTok{labs}\NormalTok{(}
    \AttributeTok{title =} \StringTok{"✓ WBA Parameter Recovery"}\NormalTok{, }
    \AttributeTok{subtitle =} \StringTok{"rho (relative weight) recovers ok; kappa (total scale) has wider uncertainty as expected"}\NormalTok{,}
    \AttributeTok{x =} \StringTok{"Parameter Value"}\NormalTok{, }
    \AttributeTok{y =} \StringTok{"Density"}
\NormalTok{  ) }\SpecialCharTok{+} 
  \FunctionTok{theme}\NormalTok{(}\AttributeTok{legend.title =} \FunctionTok{element\_blank}\NormalTok{())}
\end{Highlighting}
\end{Shaded}

\begin{center}
\includegraphics[width=0.8\linewidth,height=\textheight,keepaspectratio]{10-BayesianModels_files/figure-pdf/ch10_wba_reparam_recovery-1.pdf}
\end{center}

Rho is pretty good. Kappa not as good. So the problem of exactly how
much evidence is counted from the sources doesn't seem to be solvable
via parameterization. This would call for a change in the experimental
paradigm: we do not simply want to know the choices of the participants,
but also their confidence in the choice, or alternatively their reaction
time (which is an indirect measure of confidence). This would give us
more information to pin down the total scale of evidence weighting, and
not just the relative weighting. But for now we proceed as is.

\subsection{Comparing the Two
Parameterisations}\label{comparing-the-two-parameterisations}

\begin{Shaded}
\begin{Highlighting}[]
\CommentTok{\# Joint posterior in reparameterised space}
\NormalTok{p\_good }\OtherTok{\textless{}{-}} \FunctionTok{as\_draws\_df}\NormalTok{(fit\_wba}\SpecialCharTok{$}\NormalTok{Self\_Focused}\SpecialCharTok{$}\FunctionTok{draws}\NormalTok{(}\FunctionTok{c}\NormalTok{(}\StringTok{"rho"}\NormalTok{, }\StringTok{"kappa"}\NormalTok{))) }\SpecialCharTok{|\textgreater{}} 
  \FunctionTok{as\_tibble}\NormalTok{() }\SpecialCharTok{|\textgreater{}}
  \FunctionTok{ggplot}\NormalTok{(}\FunctionTok{aes}\NormalTok{(}\AttributeTok{x =}\NormalTok{ rho, }\AttributeTok{y =}\NormalTok{ kappa)) }\SpecialCharTok{+}
  \FunctionTok{geom\_point}\NormalTok{(}\AttributeTok{alpha =} \FloatTok{0.2}\NormalTok{, }\AttributeTok{color =} \StringTok{"\#4575b4"}\NormalTok{, }\AttributeTok{size =} \DecValTok{1}\NormalTok{) }\SpecialCharTok{+}
  \FunctionTok{geom\_density\_2d}\NormalTok{(}\AttributeTok{color =} \StringTok{"\#313695"}\NormalTok{, }\AttributeTok{bins =} \DecValTok{4}\NormalTok{) }\SpecialCharTok{+}
  \FunctionTok{geom\_vline}\NormalTok{(}\AttributeTok{xintercept =} \FloatTok{0.75}\NormalTok{, }\AttributeTok{color =} \StringTok{"\#d73027"}\NormalTok{, }\AttributeTok{linetype =} \StringTok{"dashed"}\NormalTok{) }\SpecialCharTok{+}
  \FunctionTok{geom\_hline}\NormalTok{(}\AttributeTok{yintercept =} \FloatTok{2.0}\NormalTok{, }\AttributeTok{color =} \StringTok{"\#d73027"}\NormalTok{, }\AttributeTok{linetype =} \StringTok{"dashed"}\NormalTok{) }\SpecialCharTok{+}
  \FunctionTok{coord\_cartesian}\NormalTok{(}\AttributeTok{xlim =} \FunctionTok{c}\NormalTok{(}\DecValTok{0}\NormalTok{, }\DecValTok{1}\NormalTok{), }\AttributeTok{ylim =} \FunctionTok{c}\NormalTok{(}\DecValTok{0}\NormalTok{, }\DecValTok{6}\NormalTok{)) }\SpecialCharTok{+}
  \FunctionTok{labs}\NormalTok{(}\AttributeTok{title =} \StringTok{"Reparameterised (ρ, κ): Clean!"}\NormalTok{, }
       \AttributeTok{subtitle =} \StringTok{"Nearly orthogonal posterior"}\NormalTok{,}
       \AttributeTok{x =} \StringTok{"ρ (relative weight)"}\NormalTok{, }\AttributeTok{y =} \StringTok{"κ (total scale)"}\NormalTok{)}

\CommentTok{\# Joint posterior in original space}
\NormalTok{p\_bad }\OtherTok{\textless{}{-}} \FunctionTok{as\_draws\_df}\NormalTok{(fit\_wba}\SpecialCharTok{$}\NormalTok{Self\_Focused}\SpecialCharTok{$}\FunctionTok{draws}\NormalTok{(}\FunctionTok{c}\NormalTok{(}\StringTok{"weight\_direct"}\NormalTok{, }\StringTok{"weight\_social"}\NormalTok{))) }\SpecialCharTok{|\textgreater{}} 
  \FunctionTok{as\_tibble}\NormalTok{() }\SpecialCharTok{|\textgreater{}}
  \FunctionTok{ggplot}\NormalTok{(}\FunctionTok{aes}\NormalTok{(}\AttributeTok{x =}\NormalTok{ weight\_direct, }\AttributeTok{y =}\NormalTok{ weight\_social)) }\SpecialCharTok{+}
  \FunctionTok{geom\_point}\NormalTok{(}\AttributeTok{alpha =} \FloatTok{0.2}\NormalTok{, }\AttributeTok{color =} \StringTok{"\#d73027"}\NormalTok{, }\AttributeTok{size =} \DecValTok{1}\NormalTok{) }\SpecialCharTok{+}
  \FunctionTok{geom\_density\_2d}\NormalTok{(}\AttributeTok{color =} \StringTok{"\#a50026"}\NormalTok{, }\AttributeTok{bins =} \DecValTok{4}\NormalTok{) }\SpecialCharTok{+}
  \FunctionTok{geom\_vline}\NormalTok{(}\AttributeTok{xintercept =} \FloatTok{1.5}\NormalTok{, }\AttributeTok{color =} \StringTok{"black"}\NormalTok{, }\AttributeTok{linetype =} \StringTok{"dashed"}\NormalTok{) }\SpecialCharTok{+}
  \FunctionTok{geom\_hline}\NormalTok{(}\AttributeTok{yintercept =} \FloatTok{0.5}\NormalTok{, }\AttributeTok{color =} \StringTok{"black"}\NormalTok{, }\AttributeTok{linetype =} \StringTok{"dashed"}\NormalTok{) }\SpecialCharTok{+}
  \FunctionTok{coord\_cartesian}\NormalTok{(}\AttributeTok{xlim =} \FunctionTok{c}\NormalTok{(}\DecValTok{0}\NormalTok{, }\DecValTok{4}\NormalTok{), }\AttributeTok{ylim =} \FunctionTok{c}\NormalTok{(}\DecValTok{0}\NormalTok{, }\DecValTok{4}\NormalTok{)) }\SpecialCharTok{+}
  \FunctionTok{labs}\NormalTok{(}\AttributeTok{title =} \StringTok{"Original (w\_d, w\_s): Ridge!"}\NormalTok{, }
       \AttributeTok{subtitle =} \StringTok{"Strong correlation → poor recovery"}\NormalTok{,}
       \AttributeTok{x =} \StringTok{"w\_d"}\NormalTok{, }\AttributeTok{y =} \StringTok{"w\_s"}\NormalTok{) }

\NormalTok{p\_bad }\SpecialCharTok{|}\NormalTok{ p\_good }\SpecialCharTok{+}
  \FunctionTok{plot\_annotation}\NormalTok{(}
    \AttributeTok{title =} \StringTok{"The Power of Reparameterisation"}\NormalTok{,}
    \AttributeTok{subtitle =} \StringTok{"Same model, different parameterisation, vastly different posterior geometry"}
\NormalTok{  )}
\end{Highlighting}
\end{Shaded}

\begin{center}
\includegraphics[width=0.8\linewidth,height=\textheight,keepaspectratio]{10-BayesianModels_files/figure-pdf/ch10_compare_parameterisations-1.pdf}
\end{center}

\subsection{Prior sensitivity}\label{prior-sensitivity-2}

\begin{Shaded}
\begin{Highlighting}[]
\DocumentationTok{\#\# Prior Sensitivity}
\NormalTok{sens\_wba }\OtherTok{\textless{}{-}}\NormalTok{ priorsense}\SpecialCharTok{::}\FunctionTok{powerscale\_sensitivity}\NormalTok{(}
\NormalTok{  fit\_wba}\SpecialCharTok{$}\NormalTok{Self\_Focused,}
  \AttributeTok{variable =} \FunctionTok{c}\NormalTok{(}\StringTok{"rho"}\NormalTok{, }\StringTok{"kappa"}\NormalTok{)}
\NormalTok{)}
\FunctionTok{print}\NormalTok{(sens\_wba)}
\end{Highlighting}
\end{Shaded}

\begin{verbatim}
Sensitivity based on cjs_dist
Prior selection: all priors
Likelihood selection: all data

 variable prior likelihood                                diagnosis
      rho 0.033      0.094                                        -
    kappa 0.313      0.012 potential strong prior / weak likelihood
\end{verbatim}

\begin{Shaded}
\begin{Highlighting}[]
\NormalTok{ps\_seq }\OtherTok{\textless{}{-}}\NormalTok{ priorsense}\SpecialCharTok{::}\FunctionTok{powerscale\_sequence}\NormalTok{(}
\NormalTok{  fit\_wba}\SpecialCharTok{$}\NormalTok{Self\_Focused,}
  \AttributeTok{variable =} \FunctionTok{c}\NormalTok{(}\StringTok{"rho"}\NormalTok{, }\StringTok{"kappa"}\NormalTok{)}
\NormalTok{)}
\NormalTok{priorsense}\SpecialCharTok{::}\FunctionTok{powerscale\_plot\_dens}\NormalTok{(ps\_seq) }\SpecialCharTok{+}
  \FunctionTok{labs}\NormalTok{(}\AttributeTok{title =} \StringTok{"Prior Sensitivity: WBA (Self{-}Focused)"}\NormalTok{,}
       \AttributeTok{subtitle =} \StringTok{"Does scaling the prior/likelihood shift the posteriors?"}\NormalTok{)}
\end{Highlighting}
\end{Shaded}

\begin{center}
\includegraphics[width=0.8\linewidth,height=\textheight,keepaspectratio]{10-BayesianModels_files/figure-pdf/ch10_prior_sensitivity-1.pdf}
\end{center}

ρ shows low sensitivity to prior (well-identified from conflict trials),
tho' it could use more data. κ shows moderate sensitivity (harder to pin
down from binary outcomes), consistent with the identifiability analysis
above, and it's not helped by more data (the problem is deeper). All in
all the prior sensitivity analysis confirms our previous analysis (but
it couldn't have replaced it), and tells us that more data would help ρ.
We'll see later whether multilevel modeling can help with that.

\subsection{Loo-pit}\label{loo-pit-1}

\begin{Shaded}
\begin{Highlighting}[]
\DocumentationTok{\#\# LOO{-}PIT Calibration (Randomised for Discrete Outcomes)}
\CommentTok{\# For Bernoulli outcomes, standard LOO{-}PIT is discrete and requires}
\CommentTok{\# randomisation. We use the log{-}likelihood to reconstruct P(y=1|y\_\{{-}i\}).}
\CommentTok{\# Note: for a BetaBinomial(1, alpha, beta), P(y=1) = alpha / (alpha + beta).}
\CommentTok{\# The LOO predictive probability is approximated via PSIS.}

\NormalTok{loo\_wba\_self }\OtherTok{\textless{}{-}} \FunctionTok{loo}\NormalTok{(fit\_wba}\SpecialCharTok{$}\NormalTok{Self\_Focused}\SpecialCharTok{$}\FunctionTok{draws}\NormalTok{(}\StringTok{"log\_lik"}\NormalTok{, }\AttributeTok{format =} \StringTok{"matrix"}\NormalTok{))}

\CommentTok{\# Extract LOO predictive probabilities for y=1}
\CommentTok{\# For Bernoulli: p\_loo = exp(elpd\_loo) when y=1, }
\CommentTok{\#                        and 1 {-} exp(elpd\_loo) when y=0 is more complex.}
\CommentTok{\# The proper way: use the posterior predictive draws directly.}
\NormalTok{pp\_draws }\OtherTok{\textless{}{-}}\NormalTok{ fit\_wba}\SpecialCharTok{$}\NormalTok{Self\_Focused}\SpecialCharTok{$}\FunctionTok{draws}\NormalTok{(}\StringTok{"posterior\_pred"}\NormalTok{, }\AttributeTok{format =} \StringTok{"matrix"}\NormalTok{)}
\NormalTok{y\_obs    }\OtherTok{\textless{}{-}}\NormalTok{ stan\_data\_wba\_reparam}\SpecialCharTok{$}\NormalTok{Self\_Focused}\SpecialCharTok{$}\NormalTok{choice}

\CommentTok{\# Compute randomised PIT using posterior predictive samples}
\NormalTok{pit\_vals }\OtherTok{\textless{}{-}} \FunctionTok{sapply}\NormalTok{(}\FunctionTok{seq\_along}\NormalTok{(y\_obs), }\ControlFlowTok{function}\NormalTok{(i) \{}
\NormalTok{  p\_leq }\OtherTok{\textless{}{-}} \FunctionTok{mean}\NormalTok{(pp\_draws[, i] }\SpecialCharTok{\textless{}=}\NormalTok{ y\_obs[i])}
\NormalTok{  p\_lt  }\OtherTok{\textless{}{-}} \FunctionTok{mean}\NormalTok{(pp\_draws[, i] }\SpecialCharTok{\textless{}}\NormalTok{ y\_obs[i])}
  \FunctionTok{runif}\NormalTok{(}\DecValTok{1}\NormalTok{, p\_lt, p\_leq)}
\NormalTok{\})}

\NormalTok{bayesplot}\SpecialCharTok{::}\FunctionTok{ppc\_pit\_ecdf}\NormalTok{(}\AttributeTok{y =}\NormalTok{ y\_obs, }\AttributeTok{pit =}\NormalTok{ pit\_vals) }\SpecialCharTok{+}
  \FunctionTok{labs}\NormalTok{(}\AttributeTok{title =} \StringTok{"LOO{-}PIT: WBA (Self{-}Focused)"}\NormalTok{,}
       \AttributeTok{subtitle =} \StringTok{"Flat = well{-}calibrated | Hump = model too variable | U = model too deterministic"}\NormalTok{)}
\end{Highlighting}
\end{Shaded}

\begin{center}
\includegraphics[width=0.8\linewidth,height=\textheight,keepaspectratio]{10-BayesianModels_files/figure-pdf/ch10_loo_pit-1.pdf}
\end{center}

The results of this analysis vary a bit from run to run. Most of the
time the calibration looks fine. Sometimes, participants seem more
deterministic than probability matching predicts, as indicated by the
LOO-PIT showing a hump. This generally would indicate that a softmax
observation model might be better at capturing the data. However, this
is simulated data from the beta-binomial, so we know we are right in
using a beta-binomial and we will keep doing so. With empirical data we
wouldn't be as sure :-) In any case, remember: all diagnostics are part
of a critical evaluation, not deterministic rules!

\section{Simulation-Based Calibration Checks
(SBC)}\label{simulation-based-calibration-checks-sbc-2}

Now that we have roughly working models (we can identify the ratio
pretty decently), we validate them systematically with SBC. This
provides a global certificate that our posteriors are well-calibrated
across the entire prior space.

\begin{Shaded}
\begin{Highlighting}[]
\NormalTok{n\_sbc\_datasets }\OtherTok{\textless{}{-}} \DecValTok{500}

\ControlFlowTok{if}\NormalTok{ (regenerate\_simulations }\SpecialCharTok{||} \SpecialCharTok{!}\FunctionTok{all}\NormalTok{(}\FunctionTok{file.exists}\NormalTok{(}\FunctionTok{c}\NormalTok{(}\StringTok{"simmodels/ch10\_sbc\_wba\_results.rds"}\NormalTok{, }\StringTok{"simmodels/ch10\_sbc\_pba\_results.rds"}\NormalTok{)))) \{}
  
\NormalTok{  fixed\_ev\_grid }\OtherTok{\textless{}{-}}\NormalTok{ evidence\_combinations }\SpecialCharTok{|\textgreater{}} \FunctionTok{slice}\NormalTok{(}\FunctionTok{rep}\NormalTok{(}\FunctionTok{seq\_len}\NormalTok{(}\FunctionTok{n}\NormalTok{()), }\AttributeTok{times =} \DecValTok{5}\NormalTok{))}

  \CommentTok{\# ── WBA SBC generator ──────────────────────────────────────────────}
\NormalTok{  sbc\_generator\_wba }\OtherTok{\textless{}{-}} \FunctionTok{SBC\_generator\_function}\NormalTok{(}
    \ControlFlowTok{function}\NormalTok{(ev) \{}
\NormalTok{      N }\OtherTok{\textless{}{-}} \FunctionTok{nrow}\NormalTok{(ev)}
\NormalTok{      rho\_true   }\OtherTok{\textless{}{-}} \FunctionTok{rbeta}\NormalTok{(}\DecValTok{1}\NormalTok{, }\DecValTok{2}\NormalTok{, }\DecValTok{2}\NormalTok{)}
\NormalTok{      kappa\_true }\OtherTok{\textless{}{-}} \FunctionTok{rlnorm}\NormalTok{(}\DecValTok{1}\NormalTok{, }\FunctionTok{log}\NormalTok{(}\DecValTok{2}\NormalTok{), }\FloatTok{0.5}\NormalTok{)}
\NormalTok{      wd }\OtherTok{\textless{}{-}}\NormalTok{ rho\_true }\SpecialCharTok{*}\NormalTok{ kappa\_true}
\NormalTok{      ws }\OtherTok{\textless{}{-}}\NormalTok{ (}\DecValTok{1} \SpecialCharTok{{-}}\NormalTok{ rho\_true) }\SpecialCharTok{*}\NormalTok{ kappa\_true}

\NormalTok{      alpha\_p }\OtherTok{\textless{}{-}} \FloatTok{0.5} \SpecialCharTok{+}\NormalTok{ wd }\SpecialCharTok{*}\NormalTok{ ev}\SpecialCharTok{$}\NormalTok{blue1 }\SpecialCharTok{+}\NormalTok{ ws }\SpecialCharTok{*}\NormalTok{ ev}\SpecialCharTok{$}\NormalTok{blue2}
\NormalTok{      beta\_p  }\OtherTok{\textless{}{-}} \FloatTok{0.5} \SpecialCharTok{+}\NormalTok{ wd }\SpecialCharTok{*}\NormalTok{ (ev}\SpecialCharTok{$}\NormalTok{total1 }\SpecialCharTok{{-}}\NormalTok{ ev}\SpecialCharTok{$}\NormalTok{blue1) }\SpecialCharTok{+}
\NormalTok{                       ws }\SpecialCharTok{*}\NormalTok{ (ev}\SpecialCharTok{$}\NormalTok{total2 }\SpecialCharTok{{-}}\NormalTok{ ev}\SpecialCharTok{$}\NormalTok{blue2)}
      
\NormalTok{      p\_blue  }\OtherTok{\textless{}{-}} \FunctionTok{rbeta}\NormalTok{(N, alpha\_p, beta\_p)}
\NormalTok{      choices }\OtherTok{\textless{}{-}} \FunctionTok{rbinom}\NormalTok{(N, }\DecValTok{1}\NormalTok{, p\_blue)}

      \FunctionTok{list}\NormalTok{(}
        \AttributeTok{variables =} \FunctionTok{list}\NormalTok{(}\AttributeTok{rho =}\NormalTok{ rho\_true, }\AttributeTok{kappa =}\NormalTok{ kappa\_true),}
        \AttributeTok{generated =} \FunctionTok{list}\NormalTok{(}
          \AttributeTok{N =}\NormalTok{ N, }\AttributeTok{choice =}\NormalTok{ choices,}
          \AttributeTok{blue1 =}\NormalTok{ ev}\SpecialCharTok{$}\NormalTok{blue1, }\AttributeTok{total1 =}\NormalTok{ ev}\SpecialCharTok{$}\NormalTok{total1,}
          \AttributeTok{blue2 =}\NormalTok{ ev}\SpecialCharTok{$}\NormalTok{blue2, }\AttributeTok{total2 =}\NormalTok{ ev}\SpecialCharTok{$}\NormalTok{total2}
\NormalTok{        )}
\NormalTok{      )}
\NormalTok{    \},}
    \AttributeTok{ev =}\NormalTok{ fixed\_ev\_grid}
\NormalTok{  )}

\NormalTok{  sbc\_datasets\_wba }\OtherTok{\textless{}{-}} \FunctionTok{generate\_datasets}\NormalTok{(sbc\_generator\_wba, n\_sbc\_datasets)}

\NormalTok{  sbc\_backend\_wba }\OtherTok{\textless{}{-}}\NormalTok{ SBC}\SpecialCharTok{::}\FunctionTok{SBC\_backend\_cmdstan\_sample}\NormalTok{(}
\NormalTok{    mod\_weighted,}
    \AttributeTok{chains =} \DecValTok{2}\NormalTok{, }\AttributeTok{iter\_warmup =} \DecValTok{500}\NormalTok{, }\AttributeTok{iter\_sampling =} \DecValTok{500}\NormalTok{,}
    \AttributeTok{parallel\_chains =} \DecValTok{2}\NormalTok{, }\AttributeTok{refresh =} \DecValTok{0}\NormalTok{, }\AttributeTok{show\_messages =} \ConstantTok{FALSE}
\NormalTok{  )}

\NormalTok{  sbc\_results\_wba }\OtherTok{\textless{}{-}} \FunctionTok{compute\_SBC}\NormalTok{(sbc\_datasets\_wba, sbc\_backend\_wba, }\AttributeTok{keep\_fits =} \ConstantTok{FALSE}\NormalTok{)}
  \FunctionTok{saveRDS}\NormalTok{(sbc\_results\_wba, }\StringTok{"simmodels/ch10\_sbc\_wba\_results.rds"}\NormalTok{)}

  \CommentTok{\# ── PBA SBC generator ──────────────────────────────────────────────}
\NormalTok{  sbc\_generator\_pba }\OtherTok{\textless{}{-}} \FunctionTok{SBC\_generator\_function}\NormalTok{(}
    \ControlFlowTok{function}\NormalTok{(ev) \{}
\NormalTok{      N }\OtherTok{\textless{}{-}} \FunctionTok{nrow}\NormalTok{(ev)}
\NormalTok{      p\_true }\OtherTok{\textless{}{-}} \FunctionTok{rbeta}\NormalTok{(}\DecValTok{1}\NormalTok{, }\DecValTok{2}\NormalTok{, }\DecValTok{2}\NormalTok{)}
      
\NormalTok{      alpha\_p }\OtherTok{\textless{}{-}} \FloatTok{0.5} \SpecialCharTok{+}\NormalTok{ p\_true }\SpecialCharTok{*}\NormalTok{ ev}\SpecialCharTok{$}\NormalTok{blue1 }\SpecialCharTok{+}\NormalTok{ (}\DecValTok{1} \SpecialCharTok{{-}}\NormalTok{ p\_true) }\SpecialCharTok{*}\NormalTok{ ev}\SpecialCharTok{$}\NormalTok{blue2}
\NormalTok{      beta\_p  }\OtherTok{\textless{}{-}} \FloatTok{0.5} \SpecialCharTok{+}\NormalTok{ p\_true }\SpecialCharTok{*}\NormalTok{ (ev}\SpecialCharTok{$}\NormalTok{total1 }\SpecialCharTok{{-}}\NormalTok{ ev}\SpecialCharTok{$}\NormalTok{blue1) }\SpecialCharTok{+}
\NormalTok{                       (}\DecValTok{1} \SpecialCharTok{{-}}\NormalTok{ p\_true) }\SpecialCharTok{*}\NormalTok{ (ev}\SpecialCharTok{$}\NormalTok{total2 }\SpecialCharTok{{-}}\NormalTok{ ev}\SpecialCharTok{$}\NormalTok{blue2)}
                       
\NormalTok{      theta   }\OtherTok{\textless{}{-}} \FunctionTok{rbeta}\NormalTok{(N, alpha\_p, beta\_p)}
\NormalTok{      choices }\OtherTok{\textless{}{-}} \FunctionTok{rbinom}\NormalTok{(N, }\DecValTok{1}\NormalTok{, theta)}

      \FunctionTok{list}\NormalTok{(}
        \AttributeTok{variables =} \FunctionTok{list}\NormalTok{(}\AttributeTok{p =}\NormalTok{ p\_true),}
        \AttributeTok{generated =} \FunctionTok{list}\NormalTok{(}
          \AttributeTok{N =}\NormalTok{ N, }\AttributeTok{choice =}\NormalTok{ choices,}
          \AttributeTok{blue1 =}\NormalTok{ ev}\SpecialCharTok{$}\NormalTok{blue1, }\AttributeTok{total1 =}\NormalTok{ ev}\SpecialCharTok{$}\NormalTok{total1,}
          \AttributeTok{blue2 =}\NormalTok{ ev}\SpecialCharTok{$}\NormalTok{blue2, }\AttributeTok{total2 =}\NormalTok{ ev}\SpecialCharTok{$}\NormalTok{total2}
\NormalTok{        )}
\NormalTok{      )}
\NormalTok{    \},}
    \AttributeTok{ev =}\NormalTok{ fixed\_ev\_grid}
\NormalTok{  )}

\NormalTok{  sbc\_datasets\_pba }\OtherTok{\textless{}{-}} \FunctionTok{generate\_datasets}\NormalTok{(sbc\_generator\_pba, n\_sbc\_datasets)}

\NormalTok{  sbc\_backend\_pba }\OtherTok{\textless{}{-}}\NormalTok{ SBC}\SpecialCharTok{::}\FunctionTok{SBC\_backend\_cmdstan\_sample}\NormalTok{(}
\NormalTok{    mod\_proportional,}
    \AttributeTok{chains =} \DecValTok{2}\NormalTok{, }\AttributeTok{iter\_warmup =} \DecValTok{500}\NormalTok{, }\AttributeTok{iter\_sampling =} \DecValTok{500}\NormalTok{,}
    \AttributeTok{parallel\_chains =} \DecValTok{2}\NormalTok{, }\AttributeTok{refresh =} \DecValTok{0}\NormalTok{, }\AttributeTok{show\_messages =} \ConstantTok{FALSE}
\NormalTok{  )}

\NormalTok{  sbc\_results\_pba }\OtherTok{\textless{}{-}} \FunctionTok{compute\_SBC}\NormalTok{(sbc\_datasets\_pba, sbc\_backend\_pba, }\AttributeTok{keep\_fits =} \ConstantTok{FALSE}\NormalTok{)}
  \FunctionTok{saveRDS}\NormalTok{(sbc\_results\_pba, }\StringTok{"simmodels/ch10\_sbc\_pba\_results.rds"}\NormalTok{)}

\NormalTok{\} }\ControlFlowTok{else}\NormalTok{ \{}
\NormalTok{  sbc\_results\_wba }\OtherTok{\textless{}{-}} \FunctionTok{readRDS}\NormalTok{(}\StringTok{"simmodels/ch10\_sbc\_wba\_results.rds"}\NormalTok{)}
\NormalTok{  sbc\_results\_pba }\OtherTok{\textless{}{-}} \FunctionTok{readRDS}\NormalTok{(}\StringTok{"simmodels/ch10\_sbc\_pba\_results.rds"}\NormalTok{)}
\NormalTok{\}}

\CommentTok{\# Plot SBC results}
\NormalTok{p\_sbc\_wba }\OtherTok{\textless{}{-}}\NormalTok{ SBC}\SpecialCharTok{::}\FunctionTok{plot\_ecdf\_diff}\NormalTok{(sbc\_results\_wba, }\AttributeTok{variables =} \FunctionTok{c}\NormalTok{(}\StringTok{"rho"}\NormalTok{, }\StringTok{"kappa"}\NormalTok{)) }\SpecialCharTok{+}
  \FunctionTok{labs}\NormalTok{(}\AttributeTok{title =} \StringTok{"SBC Calibration: WBA (ρ, κ)"}\NormalTok{,}
       \AttributeTok{subtitle =} \FunctionTok{paste0}\NormalTok{(}\StringTok{"N = "}\NormalTok{, n\_sbc\_datasets, }\StringTok{" datasets. Flat = well{-}calibrated."}\NormalTok{))}

\NormalTok{p\_sbc\_pba }\OtherTok{\textless{}{-}}\NormalTok{ SBC}\SpecialCharTok{::}\FunctionTok{plot\_ecdf\_diff}\NormalTok{(sbc\_results\_pba, }\AttributeTok{variables =} \StringTok{"p"}\NormalTok{) }\SpecialCharTok{+}
  \FunctionTok{labs}\NormalTok{(}\AttributeTok{title =} \StringTok{"SBC Calibration: PBA"}\NormalTok{,}
       \AttributeTok{subtitle =} \FunctionTok{paste0}\NormalTok{(}\StringTok{"N = "}\NormalTok{, n\_sbc\_datasets, }\StringTok{" datasets."}\NormalTok{))}

\NormalTok{p\_sbc\_wba }\SpecialCharTok{/}\NormalTok{ p\_sbc\_pba}
\end{Highlighting}
\end{Shaded}

\begin{center}
\includegraphics[width=0.8\linewidth,height=\textheight,keepaspectratio]{10-BayesianModels_files/figure-pdf/ch10_sbc-1.pdf}
\end{center}

The movements are not exactly a brownian motion as we'd expect, but they
stay well within the blue boundaries, so this doesn't give us any clear
signal of miscalibration for either model.

\subsection{Global Parameter Recovery from
SBC}\label{global-parameter-recovery-from-sbc}

\begin{Shaded}
\begin{Highlighting}[]
\NormalTok{wba\_recovery }\OtherTok{\textless{}{-}}\NormalTok{ sbc\_results\_wba}\SpecialCharTok{$}\NormalTok{stats }\SpecialCharTok{|\textgreater{}}
  \FunctionTok{filter}\NormalTok{(variable }\SpecialCharTok{\%in\%} \FunctionTok{c}\NormalTok{(}\StringTok{"rho"}\NormalTok{, }\StringTok{"kappa"}\NormalTok{)) }\SpecialCharTok{|\textgreater{}}
  \FunctionTok{mutate}\NormalTok{(}\AttributeTok{model =} \StringTok{"WBA"}\NormalTok{)}

\NormalTok{pba\_recovery }\OtherTok{\textless{}{-}}\NormalTok{ sbc\_results\_pba}\SpecialCharTok{$}\NormalTok{stats }\SpecialCharTok{|\textgreater{}}
  \FunctionTok{filter}\NormalTok{(variable }\SpecialCharTok{==} \StringTok{"p"}\NormalTok{) }\SpecialCharTok{|\textgreater{}}
  \FunctionTok{mutate}\NormalTok{(}\AttributeTok{model =} \StringTok{"PBA"}\NormalTok{)}

\NormalTok{sbc\_recovery\_df }\OtherTok{\textless{}{-}} \FunctionTok{bind\_rows}\NormalTok{(wba\_recovery, pba\_recovery)}

\FunctionTok{ggplot}\NormalTok{(sbc\_recovery\_df, }\FunctionTok{aes}\NormalTok{(}\AttributeTok{x =}\NormalTok{ simulated\_value, }\AttributeTok{y =}\NormalTok{ median)) }\SpecialCharTok{+}
  \FunctionTok{geom\_abline}\NormalTok{(}\AttributeTok{slope =} \DecValTok{1}\NormalTok{, }\AttributeTok{intercept =} \DecValTok{0}\NormalTok{, }\AttributeTok{linetype =} \StringTok{"dashed"}\NormalTok{, }\AttributeTok{color =} \StringTok{"grey50"}\NormalTok{) }\SpecialCharTok{+}
  \FunctionTok{geom\_errorbar}\NormalTok{(}\FunctionTok{aes}\NormalTok{(}\AttributeTok{ymin =}\NormalTok{ q5, }\AttributeTok{ymax =}\NormalTok{ q95), }\AttributeTok{alpha =} \FloatTok{0.2}\NormalTok{, }\AttributeTok{width =} \DecValTok{0}\NormalTok{) }\SpecialCharTok{+}
  \FunctionTok{geom\_point}\NormalTok{(}\AttributeTok{alpha =} \FloatTok{0.4}\NormalTok{, }\AttributeTok{size =} \FloatTok{1.5}\NormalTok{, }\AttributeTok{color =} \StringTok{"\#4575b4"}\NormalTok{) }\SpecialCharTok{+}
  \FunctionTok{facet\_wrap}\NormalTok{(model }\SpecialCharTok{\textasciitilde{}}\NormalTok{ variable, }\AttributeTok{scales =} \StringTok{"free"}\NormalTok{) }\SpecialCharTok{+}
  \FunctionTok{labs}\NormalTok{(}
    \AttributeTok{title =} \StringTok{"Global Parameter Recovery Across SBC Universes"}\NormalTok{,}
    \AttributeTok{subtitle =} \StringTok{"Each point = one simulated dataset. Error bars = 90\% CI."}\NormalTok{,}
    \AttributeTok{x =} \StringTok{"True (simulated) value"}\NormalTok{,}
    \AttributeTok{y =} \StringTok{"Posterior median"}
\NormalTok{  ) }\SpecialCharTok{+}
  \FunctionTok{theme}\NormalTok{(}\AttributeTok{strip.background =} \FunctionTok{element\_rect}\NormalTok{(}\AttributeTok{fill =} \StringTok{"gray95"}\NormalTok{, }\AttributeTok{color =} \ConstantTok{NA}\NormalTok{))}
\end{Highlighting}
\end{Shaded}

\begin{center}
\includegraphics[width=0.8\linewidth,height=\textheight,keepaspectratio]{10-BayesianModels_files/figure-pdf/ch10_sbc_recovery-1.pdf}
\end{center}

As we could expect from our test on three scenarios, the rho parameters
in either model are nicely identified, while the kappa is a bit hopeless
:-)

\section{Case Studies in Model Comparison: Nested Models and Occam's
Razor}\label{case-studies-in-model-comparison-nested-models-and-occams-razor}

Before running a massive simulation to build a model comparison pipeline
and see if it works, let's zoom in on the mechanics of \emph{why} models
win or lose in our specific case.

Our three candidate models are mathematically nested. The Simple
Bayesian Agent (SBA) is a special case of the Proportional Bayesian
Agent (PBA) where \(p = 0.5\). In turn, the PBA is a special case of the
Weighted Bayesian Agent (WBA) where the total evidence scale
\(\kappa = 1.0\). Because of this nesting, the most complex model (WBA)
can perfectly mirror the behavior of the simpler models depending on its
parameters.

We expect a specific directionality in our ``errors'': if the true
data-generating process (DGP) is complex (WBA) but its parameters mimic
a simpler model, the procedure \emph{should} misclassify the DGP and
declare the simpler model the winner. This is Occam's Razor in action:
penalizing unnecessary complexity.

To prove this, we will simulate three specific datasets, all generated
by the WBA, but with parameters deliberately chosen to test these
boundaries:

\begin{enumerate}
\def\labelenumi{\arabic{enumi}.}
\item
  \textbf{The SBA-Mimic (Balanced):} WBA with \(\rho = 0.5\) and
  \(\kappa = 2.0\) (\(w_d = 1.0, w_s = 1.0\)).
\item
  \textbf{The PBA-Mimic (Proportional):} WBA with \(\rho = 0.8\) and
  \(\kappa = 1.0\) (\(w_d = 0.8, w_s = 0.2\)).
\item
  \textbf{The Distinct WBA (Socially Influenced):} WBA with
  \(\rho = 0.2\) and \(\kappa = 2.5\) (\(w_d = 0.5, w_s = 2.0\)).
\end{enumerate}

\begin{Shaded}
\begin{Highlighting}[]
\CommentTok{\# Helper: fit SBA, PBA, WBA to a single dataset}
\NormalTok{fit\_all\_models }\OtherTok{\textless{}{-}} \ControlFlowTok{function}\NormalTok{(stan\_data, ...) \{}
\NormalTok{  sba }\OtherTok{\textless{}{-}}\NormalTok{ mod\_simple}\SpecialCharTok{$}\FunctionTok{sample}\NormalTok{(}
    \AttributeTok{data =}\NormalTok{ stan\_data, }\AttributeTok{chains =} \DecValTok{2}\NormalTok{, }\AttributeTok{iter\_warmup =} \DecValTok{0}\NormalTok{,}
    \AttributeTok{iter\_sampling =} \DecValTok{500}\NormalTok{, }\AttributeTok{fixed\_param =} \ConstantTok{TRUE}\NormalTok{,}
    \AttributeTok{refresh =} \DecValTok{0}\NormalTok{, }\AttributeTok{show\_messages =} \ConstantTok{FALSE}\NormalTok{, ...}
\NormalTok{  )}
\NormalTok{  pba }\OtherTok{\textless{}{-}}\NormalTok{ mod\_proportional}\SpecialCharTok{$}\FunctionTok{sample}\NormalTok{(}
    \AttributeTok{data =}\NormalTok{ stan\_data, }\AttributeTok{chains =} \DecValTok{2}\NormalTok{, }\AttributeTok{iter\_warmup =} \DecValTok{500}\NormalTok{,}
    \AttributeTok{iter\_sampling =} \DecValTok{500}\NormalTok{, }\AttributeTok{refresh =} \DecValTok{0}\NormalTok{, }\AttributeTok{show\_messages =} \ConstantTok{FALSE}\NormalTok{, ...}
\NormalTok{  )}
\NormalTok{  wba }\OtherTok{\textless{}{-}}\NormalTok{ mod\_weighted}\SpecialCharTok{$}\FunctionTok{sample}\NormalTok{(}
    \AttributeTok{data =}\NormalTok{ stan\_data, }\AttributeTok{chains =} \DecValTok{2}\NormalTok{, }\AttributeTok{iter\_warmup =} \DecValTok{500}\NormalTok{,}
    \AttributeTok{iter\_sampling =} \DecValTok{500}\NormalTok{, }\AttributeTok{refresh =} \DecValTok{0}\NormalTok{, }\AttributeTok{show\_messages =} \ConstantTok{FALSE}\NormalTok{, ...}
\NormalTok{  )}
  \FunctionTok{list}\NormalTok{(}\AttributeTok{SBA =}\NormalTok{ sba, }\AttributeTok{PBA =}\NormalTok{ pba, }\AttributeTok{WBA =}\NormalTok{ wba)}
\NormalTok{\}}

\ControlFlowTok{if}\NormalTok{ (regenerate\_simulations }\SpecialCharTok{||} \SpecialCharTok{!}\FunctionTok{file.exists}\NormalTok{(}\StringTok{"simmodels/ch10\_scenario\_case\_studies.rds"}\NormalTok{)) \{}
  \FunctionTok{set.seed}\NormalTok{(}\DecValTok{2026}\NormalTok{)}

  \CommentTok{\# 1. Generate three datasets (all True Model = WBA with specific params)}
\NormalTok{  df\_sba\_mimic }\OtherTok{\textless{}{-}} \FunctionTok{generate\_agent\_decisions}\NormalTok{(}\FloatTok{1.0}\NormalTok{, }\FloatTok{1.0}\NormalTok{, evidence\_combinations, }\AttributeTok{n\_samples =} \DecValTok{5}\NormalTok{)}
\NormalTok{  df\_pba\_mimic }\OtherTok{\textless{}{-}} \FunctionTok{generate\_agent\_decisions}\NormalTok{(}\FloatTok{0.8}\NormalTok{, }\FloatTok{0.2}\NormalTok{, evidence\_combinations, }\AttributeTok{n\_samples =} \DecValTok{5}\NormalTok{)}
\NormalTok{  df\_distinct  }\OtherTok{\textless{}{-}} \FunctionTok{generate\_agent\_decisions}\NormalTok{(}\FloatTok{0.5}\NormalTok{, }\FloatTok{2.0}\NormalTok{, evidence\_combinations, }\AttributeTok{n\_samples =} \DecValTok{5}\NormalTok{)}

\NormalTok{  stan\_data\_sba\_mimic }\OtherTok{\textless{}{-}} \FunctionTok{prepare\_stan\_data}\NormalTok{(df\_sba\_mimic)}
\NormalTok{  stan\_data\_pba\_mimic }\OtherTok{\textless{}{-}} \FunctionTok{prepare\_stan\_data}\NormalTok{(df\_pba\_mimic)}
\NormalTok{  stan\_data\_distinct  }\OtherTok{\textless{}{-}} \FunctionTok{prepare\_stan\_data}\NormalTok{(df\_distinct)}

  \CommentTok{\# 2. Fit all three models to all three datasets (9 fits total)}
\NormalTok{  fits\_sba\_mimic }\OtherTok{\textless{}{-}} \FunctionTok{fit\_all\_models}\NormalTok{(stan\_data\_sba\_mimic)}
\NormalTok{  fits\_pba\_mimic }\OtherTok{\textless{}{-}} \FunctionTok{fit\_all\_models}\NormalTok{(stan\_data\_pba\_mimic)}
\NormalTok{  fits\_distinct  }\OtherTok{\textless{}{-}} \FunctionTok{fit\_all\_models}\NormalTok{(stan\_data\_distinct)}

  \CommentTok{\# 3. Pre{-}compute LOO (avoids needing the raw fit objects later)}
\NormalTok{  extract\_loos }\OtherTok{\textless{}{-}} \ControlFlowTok{function}\NormalTok{(fits) \{}
    \FunctionTok{list}\NormalTok{(}
      \AttributeTok{SBA =} \FunctionTok{loo}\NormalTok{(fits}\SpecialCharTok{$}\NormalTok{SBA}\SpecialCharTok{$}\FunctionTok{draws}\NormalTok{(}\StringTok{"log\_lik"}\NormalTok{, }\AttributeTok{format =} \StringTok{"matrix"}\NormalTok{)),}
      \AttributeTok{PBA =} \FunctionTok{loo}\NormalTok{(fits}\SpecialCharTok{$}\NormalTok{PBA}\SpecialCharTok{$}\FunctionTok{draws}\NormalTok{(}\StringTok{"log\_lik"}\NormalTok{, }\AttributeTok{format =} \StringTok{"matrix"}\NormalTok{)),}
      \AttributeTok{WBA =} \FunctionTok{loo}\NormalTok{(fits}\SpecialCharTok{$}\NormalTok{WBA}\SpecialCharTok{$}\FunctionTok{draws}\NormalTok{(}\StringTok{"log\_lik"}\NormalTok{, }\AttributeTok{format =} \StringTok{"matrix"}\NormalTok{))}
\NormalTok{    )}
\NormalTok{  \}}
\NormalTok{  loos\_sba\_mimic }\OtherTok{\textless{}{-}} \FunctionTok{extract\_loos}\NormalTok{(fits\_sba\_mimic)}
\NormalTok{  loos\_pba\_mimic }\OtherTok{\textless{}{-}} \FunctionTok{extract\_loos}\NormalTok{(fits\_pba\_mimic)}
\NormalTok{  loos\_distinct  }\OtherTok{\textless{}{-}} \FunctionTok{extract\_loos}\NormalTok{(fits\_distinct)}

  \CommentTok{\# 4. Save everything we need downstream: data frames + LOO objects.}
  \CommentTok{\#    We do NOT save the fit objects themselves — the LOO objects and}
  \CommentTok{\#    data frames are all the downstream chunks use.}
  \FunctionTok{saveRDS}\NormalTok{(}\FunctionTok{list}\NormalTok{(}
    \AttributeTok{df\_sba\_mimic =}\NormalTok{ df\_sba\_mimic,}
    \AttributeTok{df\_pba\_mimic =}\NormalTok{ df\_pba\_mimic,}
    \AttributeTok{df\_distinct  =}\NormalTok{ df\_distinct,}
    \AttributeTok{loos\_sba\_mimic =}\NormalTok{ loos\_sba\_mimic,}
    \AttributeTok{loos\_pba\_mimic =}\NormalTok{ loos\_pba\_mimic,}
    \AttributeTok{loos\_distinct  =}\NormalTok{ loos\_distinct}
\NormalTok{  ), }\StringTok{"simmodels/ch10\_scenario\_case\_studies.rds"}\NormalTok{)}

\NormalTok{\} }\ControlFlowTok{else}\NormalTok{ \{}
\NormalTok{  cache }\OtherTok{\textless{}{-}} \FunctionTok{readRDS}\NormalTok{(}\StringTok{"simmodels/ch10\_scenario\_case\_studies.rds"}\NormalTok{)}
\NormalTok{  df\_sba\_mimic   }\OtherTok{\textless{}{-}}\NormalTok{ cache}\SpecialCharTok{$}\NormalTok{df\_sba\_mimic}
\NormalTok{  df\_pba\_mimic   }\OtherTok{\textless{}{-}}\NormalTok{ cache}\SpecialCharTok{$}\NormalTok{df\_pba\_mimic}
\NormalTok{  df\_distinct    }\OtherTok{\textless{}{-}}\NormalTok{ cache}\SpecialCharTok{$}\NormalTok{df\_distinct}
\NormalTok{  loos\_sba\_mimic }\OtherTok{\textless{}{-}}\NormalTok{ cache}\SpecialCharTok{$}\NormalTok{loos\_sba\_mimic}
\NormalTok{  loos\_pba\_mimic }\OtherTok{\textless{}{-}}\NormalTok{ cache}\SpecialCharTok{$}\NormalTok{loos\_pba\_mimic}
\NormalTok{  loos\_distinct  }\OtherTok{\textless{}{-}}\NormalTok{ cache}\SpecialCharTok{$}\NormalTok{loos\_distinct}
\NormalTok{\}}
\end{Highlighting}
\end{Shaded}

Now, let's look at the stacking weights and LOO comparisons for these
three scenarios.

\begin{Shaded}
\begin{Highlighting}[]
\CommentTok{\# Print comparison summaries using the pre{-}computed LOO objects}
\NormalTok{print\_scenario }\OtherTok{\textless{}{-}} \ControlFlowTok{function}\NormalTok{(loos, name, true\_params) \{}
  \FunctionTok{cat}\NormalTok{(}\FunctionTok{sprintf}\NormalTok{(}\StringTok{"}\SpecialCharTok{\textbackslash{}n}\StringTok{=== \%s ===}\SpecialCharTok{\textbackslash{}n}\StringTok{True DGP: WBA \%s}\SpecialCharTok{\textbackslash{}n}\StringTok{"}\NormalTok{, name, true\_params))}
  \FunctionTok{print}\NormalTok{(loo}\SpecialCharTok{::}\FunctionTok{loo\_compare}\NormalTok{(loos))}
\NormalTok{  w }\OtherTok{\textless{}{-}}\NormalTok{ loo}\SpecialCharTok{::}\FunctionTok{loo\_model\_weights}\NormalTok{(loos)}
  \FunctionTok{print}\NormalTok{(}\FunctionTok{round}\NormalTok{(w, }\DecValTok{3}\NormalTok{))}
\NormalTok{\}}

\FunctionTok{print\_scenario}\NormalTok{(loos\_sba\_mimic, }\StringTok{"Scenario 1: The SBA{-}Mimic"}\NormalTok{,  }\StringTok{"(rho=0.5, kappa=2.0)"}\NormalTok{)}
\end{Highlighting}
\end{Shaded}

\begin{verbatim}

=== Scenario 1: The SBA-Mimic ===
True DGP: WBA (rho=0.5, kappa=2.0)
    elpd_diff se_diff
SBA  0.0       0.0   
WBA -0.9       0.4   
PBA -0.9       0.7   
Method: stacking
------
    weight
SBA 1.000 
PBA 0.000 
WBA 0.000 
\end{verbatim}

\begin{Shaded}
\begin{Highlighting}[]
\FunctionTok{print\_scenario}\NormalTok{(loos\_pba\_mimic, }\StringTok{"Scenario 2: The PBA{-}Mimic"}\NormalTok{,  }\StringTok{"(rho=0.8, kappa=1.0)"}\NormalTok{)}
\end{Highlighting}
\end{Shaded}

\begin{verbatim}

=== Scenario 2: The PBA-Mimic ===
True DGP: WBA (rho=0.8, kappa=1.0)
    elpd_diff se_diff
PBA  0.0       0.0   
WBA  0.0       0.4   
SBA -2.4       2.2   
Method: stacking
------
    weight
SBA 0.029 
PBA 0.663 
WBA 0.308 
\end{verbatim}

\begin{Shaded}
\begin{Highlighting}[]
\FunctionTok{print\_scenario}\NormalTok{(loos\_distinct,  }\StringTok{"Scenario 3: The Distinct WBA"}\NormalTok{,}\StringTok{"(rho=0.2, kappa=2.5)"}\NormalTok{)}
\end{Highlighting}
\end{Shaded}

\begin{verbatim}

=== Scenario 3: The Distinct WBA ===
True DGP: WBA (rho=0.2, kappa=2.5)
    elpd_diff se_diff
PBA  0.0       0.0   
WBA  0.0       0.6   
SBA -6.9       3.9   
Method: stacking
------
    weight
SBA 0.012 
PBA 0.618 
WBA 0.369 
\end{verbatim}

\begin{Shaded}
\begin{Highlighting}[]
\CommentTok{\# Create the dataframe directly from the console output}
\NormalTok{extract\_loo\_summary }\OtherTok{\textless{}{-}} \ControlFlowTok{function}\NormalTok{(loos, scenario\_name) \{}
\NormalTok{  comp }\OtherTok{\textless{}{-}} \FunctionTok{as\_tibble}\NormalTok{(loo}\SpecialCharTok{::}\FunctionTok{loo\_compare}\NormalTok{(loos), }\AttributeTok{rownames =} \StringTok{"Model"}\NormalTok{) }\SpecialCharTok{|\textgreater{}} \FunctionTok{select}\NormalTok{(Model, elpd\_diff, se\_diff)}
\NormalTok{  w }\OtherTok{\textless{}{-}} \FunctionTok{tibble}\NormalTok{(}\AttributeTok{Model =} \FunctionTok{names}\NormalTok{(loo}\SpecialCharTok{::}\FunctionTok{loo\_model\_weights}\NormalTok{(loos)), }\AttributeTok{weight =} \FunctionTok{as.numeric}\NormalTok{(loo}\SpecialCharTok{::}\FunctionTok{loo\_model\_weights}\NormalTok{(loos)))}
\NormalTok{  comp }\SpecialCharTok{|\textgreater{}} \FunctionTok{left\_join}\NormalTok{(w, }\AttributeTok{by =} \StringTok{"Model"}\NormalTok{) }\SpecialCharTok{|\textgreater{}} \FunctionTok{mutate}\NormalTok{(}\AttributeTok{Scenario =}\NormalTok{ scenario\_name)}
\NormalTok{\}}

\CommentTok{\# Create the summary dataframe dynamically without hardcoding values}
\NormalTok{scenario\_summary }\OtherTok{\textless{}{-}} \FunctionTok{bind\_rows}\NormalTok{(}
  \FunctionTok{extract\_loo\_summary}\NormalTok{(loos\_sba\_mimic, }\StringTok{"1. SBA{-}Mimic"}\NormalTok{),}
  \FunctionTok{extract\_loo\_summary}\NormalTok{(loos\_pba\_mimic, }\StringTok{"2. PBA{-}Mimic"}\NormalTok{),}
  \FunctionTok{extract\_loo\_summary}\NormalTok{(loos\_distinct, }\StringTok{"3. Distinct WBA"}\NormalTok{)}
\NormalTok{) }\SpecialCharTok{|\textgreater{}}
  \FunctionTok{mutate}\NormalTok{(}
    \CommentTok{\# Lock in the factor levels for consistent left{-}to{-}right plotting}
    \AttributeTok{Model =} \FunctionTok{factor}\NormalTok{(Model, }\AttributeTok{levels =} \FunctionTok{c}\NormalTok{(}\StringTok{"SBA"}\NormalTok{, }\StringTok{"PBA"}\NormalTok{, }\StringTok{"WBA"}\NormalTok{)),}
    \AttributeTok{Scenario =} \FunctionTok{factor}\NormalTok{(Scenario, }\AttributeTok{levels =} \FunctionTok{c}\NormalTok{(}\StringTok{"1. SBA{-}Mimic"}\NormalTok{, }\StringTok{"2. PBA{-}Mimic"}\NormalTok{, }\StringTok{"3. Distinct WBA"}\NormalTok{))}
\NormalTok{  )}

\CommentTok{\# Plot 1: ELPD Differences with Standard Error}
\NormalTok{p\_elpd }\OtherTok{\textless{}{-}} \FunctionTok{ggplot}\NormalTok{(scenario\_summary, }\FunctionTok{aes}\NormalTok{(}\AttributeTok{x =}\NormalTok{ Model, }\AttributeTok{y =}\NormalTok{ elpd\_diff, }\AttributeTok{color =}\NormalTok{ Model)) }\SpecialCharTok{+}
  \FunctionTok{geom\_hline}\NormalTok{(}\AttributeTok{yintercept =} \DecValTok{0}\NormalTok{, }\AttributeTok{linetype =} \StringTok{"dashed"}\NormalTok{, }\AttributeTok{color =} \StringTok{"gray50"}\NormalTok{) }\SpecialCharTok{+}
  \FunctionTok{geom\_pointrange}\NormalTok{(}\FunctionTok{aes}\NormalTok{(}\AttributeTok{ymin =}\NormalTok{ elpd\_diff }\SpecialCharTok{{-}}\NormalTok{ se\_diff, }\AttributeTok{ymax =}\NormalTok{ elpd\_diff }\SpecialCharTok{+}\NormalTok{ se\_diff), }
                  \AttributeTok{size =} \FloatTok{0.8}\NormalTok{, }\AttributeTok{linewidth =} \DecValTok{1}\NormalTok{) }\SpecialCharTok{+}
  \FunctionTok{facet\_wrap}\NormalTok{(}\SpecialCharTok{\textasciitilde{}}\NormalTok{Scenario) }\SpecialCharTok{+}
  \FunctionTok{scale\_color\_brewer}\NormalTok{(}\AttributeTok{palette =} \StringTok{"Set2"}\NormalTok{) }\SpecialCharTok{+}
  \FunctionTok{labs}\NormalTok{(}
    \AttributeTok{title =} \StringTok{"ELPD Differences (with Standard Error)"}\NormalTok{,}
    \AttributeTok{subtitle =} \StringTok{"Zero indicates the best{-}fitting model in that scenario."}\NormalTok{,}
    \AttributeTok{y =} \StringTok{"ELPD Difference"}\NormalTok{,}
    \AttributeTok{x =} \ConstantTok{NULL}
\NormalTok{  ) }\SpecialCharTok{+}
  \FunctionTok{theme}\NormalTok{(}\AttributeTok{legend.position =} \StringTok{"none"}\NormalTok{,}
        \AttributeTok{strip.text =} \FunctionTok{element\_text}\NormalTok{(}\AttributeTok{face =} \StringTok{"bold"}\NormalTok{, }\AttributeTok{size =} \DecValTok{11}\NormalTok{))}

\CommentTok{\# Plot 2: Bayesian Stacking Weights}
\NormalTok{p\_weight }\OtherTok{\textless{}{-}} \FunctionTok{ggplot}\NormalTok{(scenario\_summary, }\FunctionTok{aes}\NormalTok{(}\AttributeTok{x =}\NormalTok{ Model, }\AttributeTok{y =}\NormalTok{ weight, }\AttributeTok{fill =}\NormalTok{ Model)) }\SpecialCharTok{+}
  \FunctionTok{geom\_col}\NormalTok{(}\AttributeTok{alpha =} \FloatTok{0.85}\NormalTok{, }\AttributeTok{width =} \FloatTok{0.7}\NormalTok{) }\SpecialCharTok{+}
  \FunctionTok{geom\_text}\NormalTok{(}\FunctionTok{aes}\NormalTok{(}\AttributeTok{label =}\NormalTok{ scales}\SpecialCharTok{::}\FunctionTok{percent}\NormalTok{(weight, }\AttributeTok{accuracy =} \DecValTok{1}\NormalTok{)), }
            \AttributeTok{vjust =} \SpecialCharTok{{-}}\FloatTok{0.5}\NormalTok{, }\AttributeTok{fontface =} \StringTok{"bold"}\NormalTok{) }\SpecialCharTok{+}
  \FunctionTok{facet\_wrap}\NormalTok{(}\SpecialCharTok{\textasciitilde{}}\NormalTok{Scenario) }\SpecialCharTok{+}
  \FunctionTok{scale\_fill\_brewer}\NormalTok{(}\AttributeTok{palette =} \StringTok{"Set2"}\NormalTok{) }\SpecialCharTok{+}
  \FunctionTok{scale\_y\_continuous}\NormalTok{(}\AttributeTok{labels =}\NormalTok{ scales}\SpecialCharTok{::}\NormalTok{percent, }\AttributeTok{limits =} \FunctionTok{c}\NormalTok{(}\DecValTok{0}\NormalTok{, }\FloatTok{1.2}\NormalTok{)) }\SpecialCharTok{+}
  \FunctionTok{labs}\NormalTok{(}
    \AttributeTok{title =} \StringTok{"Bayesian Stacking Weights"}\NormalTok{,}
    \AttributeTok{subtitle =} \StringTok{"The procedure perfectly penalizes complexity, shifting all weight to the simplest adequate model."}\NormalTok{,}
    \AttributeTok{y =} \StringTok{"Stacking Weight"}\NormalTok{,}
    \AttributeTok{x =} \ConstantTok{NULL}
\NormalTok{  ) }\SpecialCharTok{+}
  \FunctionTok{theme}\NormalTok{(}\AttributeTok{legend.position =} \StringTok{"none"}\NormalTok{,}
        \AttributeTok{strip.text =} \FunctionTok{element\_blank}\NormalTok{()) }\CommentTok{\# Hide strips on the bottom plot for a cleaner look}

\CommentTok{\# Combine with patchwork}
\NormalTok{p\_elpd }\SpecialCharTok{/}\NormalTok{ p\_weight }\SpecialCharTok{+} 
  \FunctionTok{plot\_annotation}\NormalTok{(}
    \AttributeTok{title =} \StringTok{"Model Recovery: The Occam\textquotesingle{}s Razor Effect"}\NormalTok{,}
    \AttributeTok{theme =} \FunctionTok{theme}\NormalTok{(}\AttributeTok{plot.title =} \FunctionTok{element\_text}\NormalTok{(}\AttributeTok{size =} \DecValTok{16}\NormalTok{, }\AttributeTok{face =} \StringTok{"bold"}\NormalTok{))}
\NormalTok{  )}
\end{Highlighting}
\end{Shaded}

\begin{center}
\includegraphics[width=0.8\linewidth,height=\textheight,keepaspectratio]{10-BayesianModels_files/figure-pdf/ch10_scenario_comparison_weights-1.pdf}
\end{center}

Notice the clear directionality of the classification:

\begin{itemize}
\item
  In the SBA-mimic scenario, the WBA generates the data, but it is
  effectively functioning as the 0-parameter SBA model. Both the PBA and
  WBA fit the data equally well, but LOO accurately applies a complexity
  penalty. The SBA ``wins.''
\item
  In the PBA-mimic scenario, the WBA parameters (\(\kappa = 1.0\)) mean
  the agent is trading off evidence exactly as the 1-parameter PBA
  assumes. The PBA easily beats the SBA, but it also beats the true
  generating WBA, because the WBA's extra scale parameter isn't needed.
\item
  In Scenario 3, the agent has a high total evidence sensitivity
  (\(\kappa = 2.5\)) and heavily overweights social evidence
  (\(\rho = 0.2\)). The SBA cannot flex to capture this. However, both
  the PBA and the WBA get substantial weight, with PBA winning.
\end{itemize}

This is not ideal, but neither a failure of model comparison. It tells
us that the WBA could still be likely, but it is not a better model than
PBA. This is likely because WBA is more complex (one parameter more),
but the it has issues in reconstructing the extra parameter (kappa), We
should be aware of these issues.

\subsection{Visualizing Pointwise ELPD: Where the Advantage
Lies}\label{visualizing-pointwise-elpd-where-the-advantage-lies}

To truly understand this, we can plot the pointwise ELPD differences. We
will compare the true generating model (WBA) against the simpler model
that ``won'' in the mimicking scenarios, and against the PBA in the
distinct scenario.

\begin{Shaded}
\begin{Highlighting}[]
\NormalTok{pw\_1 }\OtherTok{\textless{}{-}} \FunctionTok{tibble}\NormalTok{(}
  \AttributeTok{obs       =} \DecValTok{1}\SpecialCharTok{:}\DecValTok{180}\NormalTok{,}
  \AttributeTok{elpd\_diff =}\NormalTok{ loos\_sba\_mimic}\SpecialCharTok{$}\NormalTok{WBA}\SpecialCharTok{$}\NormalTok{pointwise[,}\StringTok{"elpd\_loo"}\NormalTok{] }\SpecialCharTok{{-}}\NormalTok{ loos\_sba\_mimic}\SpecialCharTok{$}\NormalTok{SBA}\SpecialCharTok{$}\NormalTok{pointwise[,}\StringTok{"elpd\_loo"}\NormalTok{],}
  \AttributeTok{Scenario  =} \StringTok{"1. SBA{-}Mimic}\SpecialCharTok{\textbackslash{}n}\StringTok{(True WBA vs. SBA)"}\NormalTok{,}
  \AttributeTok{blue1 =}\NormalTok{ df\_sba\_mimic}\SpecialCharTok{$}\NormalTok{blue1, }\AttributeTok{blue2 =}\NormalTok{ df\_sba\_mimic}\SpecialCharTok{$}\NormalTok{blue2}
\NormalTok{)}

\NormalTok{pw\_2 }\OtherTok{\textless{}{-}} \FunctionTok{tibble}\NormalTok{(}
  \AttributeTok{obs       =} \DecValTok{1}\SpecialCharTok{:}\DecValTok{180}\NormalTok{,}
  \AttributeTok{elpd\_diff =}\NormalTok{ loos\_pba\_mimic}\SpecialCharTok{$}\NormalTok{WBA}\SpecialCharTok{$}\NormalTok{pointwise[,}\StringTok{"elpd\_loo"}\NormalTok{] }\SpecialCharTok{{-}}\NormalTok{ loos\_pba\_mimic}\SpecialCharTok{$}\NormalTok{PBA}\SpecialCharTok{$}\NormalTok{pointwise[,}\StringTok{"elpd\_loo"}\NormalTok{],}
  \AttributeTok{Scenario  =} \StringTok{"2. PBA{-}Mimic}\SpecialCharTok{\textbackslash{}n}\StringTok{(True WBA vs. PBA)"}\NormalTok{,}
  \AttributeTok{blue1 =}\NormalTok{ df\_pba\_mimic}\SpecialCharTok{$}\NormalTok{blue1, }\AttributeTok{blue2 =}\NormalTok{ df\_pba\_mimic}\SpecialCharTok{$}\NormalTok{blue2}
\NormalTok{)}

\NormalTok{pw\_3 }\OtherTok{\textless{}{-}} \FunctionTok{tibble}\NormalTok{(}
  \AttributeTok{obs       =} \DecValTok{1}\SpecialCharTok{:}\DecValTok{180}\NormalTok{,}
  \AttributeTok{elpd\_diff =}\NormalTok{ loos\_distinct}\SpecialCharTok{$}\NormalTok{WBA}\SpecialCharTok{$}\NormalTok{pointwise[,}\StringTok{"elpd\_loo"}\NormalTok{] }\SpecialCharTok{{-}}\NormalTok{ loos\_distinct}\SpecialCharTok{$}\NormalTok{PBA}\SpecialCharTok{$}\NormalTok{pointwise[,}\StringTok{"elpd\_loo"}\NormalTok{],}
  \AttributeTok{Scenario  =} \StringTok{"3. Distinct WBA}\SpecialCharTok{\textbackslash{}n}\StringTok{(True WBA vs. PBA)"}\NormalTok{,}
  \AttributeTok{blue1 =}\NormalTok{ df\_distinct}\SpecialCharTok{$}\NormalTok{blue1, }\AttributeTok{blue2 =}\NormalTok{ df\_distinct}\SpecialCharTok{$}\NormalTok{blue2}
\NormalTok{)}

\NormalTok{plot\_pw }\OtherTok{\textless{}{-}} \FunctionTok{bind\_rows}\NormalTok{(pw\_1, pw\_2, pw\_3) }\SpecialCharTok{|\textgreater{}}
  \FunctionTok{mutate}\NormalTok{(}
    \AttributeTok{prop\_direct =}\NormalTok{ blue1 }\SpecialCharTok{/} \DecValTok{8}\NormalTok{,}
    \AttributeTok{prop\_social =}\NormalTok{ blue2 }\SpecialCharTok{/} \DecValTok{3}\NormalTok{,}
    \AttributeTok{evidence\_conflict =} \FunctionTok{abs}\NormalTok{(prop\_direct }\SpecialCharTok{{-}}\NormalTok{ prop\_social)}
\NormalTok{  )}

\FunctionTok{ggplot}\NormalTok{(plot\_pw, }\FunctionTok{aes}\NormalTok{(}\AttributeTok{x =}\NormalTok{ evidence\_conflict, }\AttributeTok{y =}\NormalTok{ elpd\_diff)) }\SpecialCharTok{+}
  \FunctionTok{geom\_hline}\NormalTok{(}\AttributeTok{yintercept =} \DecValTok{0}\NormalTok{, }\AttributeTok{linetype =} \StringTok{"dashed"}\NormalTok{, }\AttributeTok{color =} \StringTok{"gray40"}\NormalTok{) }\SpecialCharTok{+}
  \FunctionTok{geom\_point}\NormalTok{(}\FunctionTok{aes}\NormalTok{(}\AttributeTok{color =}\NormalTok{ elpd\_diff }\SpecialCharTok{\textgreater{}} \DecValTok{0}\NormalTok{), }\AttributeTok{size =} \DecValTok{2}\NormalTok{, }\AttributeTok{alpha =} \FloatTok{0.6}\NormalTok{) }\SpecialCharTok{+}
  \FunctionTok{geom\_smooth}\NormalTok{(}\AttributeTok{method =} \StringTok{"gam"}\NormalTok{, }\AttributeTok{formula =}\NormalTok{ y }\SpecialCharTok{\textasciitilde{}} \FunctionTok{s}\NormalTok{(x, }\AttributeTok{bs =} \StringTok{"cs"}\NormalTok{), }
              \AttributeTok{color =} \StringTok{"midnightblue"}\NormalTok{, }\AttributeTok{linewidth =} \DecValTok{1}\NormalTok{) }\SpecialCharTok{+}
  \FunctionTok{facet\_wrap}\NormalTok{(}\SpecialCharTok{\textasciitilde{}}\NormalTok{Scenario) }\SpecialCharTok{+}
  \FunctionTok{scale\_color\_manual}\NormalTok{(}
    \AttributeTok{values =} \FunctionTok{c}\NormalTok{(}\StringTok{"FALSE"} \OtherTok{=} \StringTok{"firebrick"}\NormalTok{, }\StringTok{"TRUE"} \OtherTok{=} \StringTok{"steelblue"}\NormalTok{),}
    \AttributeTok{labels =} \FunctionTok{c}\NormalTok{(}\StringTok{"FALSE"} \OtherTok{=} \StringTok{"Favors Simpler Model"}\NormalTok{, }\StringTok{"TRUE"} \OtherTok{=} \StringTok{"Favors True WBA"}\NormalTok{), }
    \AttributeTok{name =} \ConstantTok{NULL}
\NormalTok{  ) }\SpecialCharTok{+}
  \FunctionTok{labs}\NormalTok{(}
    \AttributeTok{title =} \StringTok{"Pointwise ELPD Differences: True DGP vs. Simpler Alternatives"}\NormalTok{,}
    \AttributeTok{subtitle =} \StringTok{"When the WBA mimics simpler models (Plots 1 \& 2), the simpler models win by a constant, small margin (complexity penalty)."}\NormalTok{,}
    \AttributeTok{x =} \StringTok{"Evidence Conflict |prop\_direct {-} prop\_social|"}\NormalTok{,}
    \AttributeTok{y =} \StringTok{"ELPD Difference (True WBA {-} Simpler Model)"}
\NormalTok{  ) }\SpecialCharTok{+}
  \FunctionTok{theme}\NormalTok{(}\AttributeTok{legend.position =} \StringTok{"bottom"}\NormalTok{)}
\end{Highlighting}
\end{Shaded}

\begin{center}
\includegraphics[width=0.8\linewidth,height=\textheight,keepaspectratio]{10-BayesianModels_files/figure-pdf/ch10_scenario_pointwise_elpd-1.pdf}
\end{center}

This plot provides some insight on how complexity penalties play out on
a trial-by-trial basis:

\begin{itemize}
\item
  In Scenario 1 (SBA-Mimic): The simpler model's advantage actively
  increases as evidence conflict grows, pulling the trend line
  progressively downward (favoring the SBA). When the true agent is
  perfectly balanced, the WBA's extra flexibility actively hurts its
  predictive performance the most when the evidence sources strongly
  disagree.
\item
  In Scenario 2 (PBA-Mimic): The overall trend line is flat and sits
  essentially at zero, reflecting a small, constant complexity penalty
  for the WBA across all conflict levels. However, the scatter shows
  that while the models often tie, there are several specific
  trials---especially at lower conflict---where the points dip sharply
  into the red. This indicates that the WBA's unnecessary extra
  parameter occasionally causes it to make noticeably worse,
  over-dispersed predictions compared to the tightly constrained PBA.
\item
  In Scenario 3 (Distinct WBA): Contrary to the expectation that the WBA
  would generate massive predictive advantages, the overall trend line
  remains stubbornly flat. As you noted, the WBA does win by small
  margins on a large number of trials across the board (the cluster of
  blue dots). However, these numerous small victories are completely
  offset by a handful of severe misses (the deep red dots), particularly
  at lower conflict levels. The WBA's added flexibility allows it to
  occasionally make overconfident, incorrect predictions. The heavy ELPD
  penalty from those few large misses drags the average down,
  neutralizing the WBA's expected advantage.
\end{itemize}

With this nested directionality in mind, we can now confidently
interpret the confusion matrix generated by our full model recovery
procedure.

\section{Model Recovery: Scaling Up}\label{model-recovery-scaling-up}

We now scale up our case studies to a full model recovery simulation.
Because the models are nested, we know Occam's razor will dominate: if a
dataset generated by the WBA has parameters that mimic the PBA or SBA,
the procedure \emph{should} appropriately select the simpler model.

For each of our three structural generating processes (SBA, PBA, WBA),
we simulate 50 synthetic datasets, cross-fit all three candidate models,
and use stacking weights to observe how often the procedure correctly
penalizes complexity and recovers the simplest adequate model.

\begin{Shaded}
\begin{Highlighting}[]
\NormalTok{n\_sims }\OtherTok{\textless{}{-}} \DecValTok{50}
\FunctionTok{plan}\NormalTok{(multisession, }\AttributeTok{workers =}\NormalTok{ parallel}\SpecialCharTok{::}\FunctionTok{detectCores}\NormalTok{() }\SpecialCharTok{{-}} \DecValTok{2}\NormalTok{)}

\ControlFlowTok{if}\NormalTok{ (regenerate\_simulations }\SpecialCharTok{||} \SpecialCharTok{!}\FunctionTok{file.exists}\NormalTok{(}\StringTok{"simmodels/ch10\_model\_recovery.rds"}\NormalTok{)) \{}

  \CommentTok{\# Generate datasets from each model type}
\NormalTok{  sim\_grid }\OtherTok{\textless{}{-}} \FunctionTok{bind\_rows}\NormalTok{(}
    \FunctionTok{tibble}\NormalTok{(}\AttributeTok{sim\_id =} \DecValTok{1}\SpecialCharTok{:}\NormalTok{n\_sims, }\AttributeTok{true\_model =} \StringTok{"SBA"}\NormalTok{, }\AttributeTok{wd =} \FloatTok{1.0}\NormalTok{, }\AttributeTok{ws =} \FloatTok{1.0}\NormalTok{),}

    \FunctionTok{tibble}\NormalTok{(}\AttributeTok{sim\_id =} \DecValTok{1}\SpecialCharTok{:}\NormalTok{n\_sims, }\AttributeTok{true\_model =} \StringTok{"PBA"}\NormalTok{,}
           \AttributeTok{p =} \FunctionTok{rbeta}\NormalTok{(n\_sims, }\DecValTok{2}\NormalTok{, }\DecValTok{2}\NormalTok{)) }\SpecialCharTok{|\textgreater{}}
      \FunctionTok{mutate}\NormalTok{(}\AttributeTok{wd =}\NormalTok{ p, }\AttributeTok{ws =} \DecValTok{1} \SpecialCharTok{{-}}\NormalTok{ p),}

    \FunctionTok{tibble}\NormalTok{(}\AttributeTok{sim\_id =} \DecValTok{1}\SpecialCharTok{:}\NormalTok{n\_sims, }\AttributeTok{true\_model =} \StringTok{"WBA"}\NormalTok{,}
           \AttributeTok{rho =} \FunctionTok{rbeta}\NormalTok{(n\_sims, }\DecValTok{2}\NormalTok{, }\DecValTok{2}\NormalTok{),}
           \AttributeTok{kappa =} \FunctionTok{rlnorm}\NormalTok{(n\_sims, }\FunctionTok{log}\NormalTok{(}\DecValTok{2}\NormalTok{), }\FloatTok{0.5}\NormalTok{)) }\SpecialCharTok{|\textgreater{}}
      \FunctionTok{mutate}\NormalTok{(}\AttributeTok{wd =}\NormalTok{ rho }\SpecialCharTok{*}\NormalTok{ kappa, }\AttributeTok{ws =}\NormalTok{ (}\DecValTok{1} \SpecialCharTok{{-}}\NormalTok{ rho) }\SpecialCharTok{*}\NormalTok{ kappa)}
\NormalTok{  )}

\NormalTok{  recovery\_data }\OtherTok{\textless{}{-}}\NormalTok{ sim\_grid }\SpecialCharTok{|\textgreater{}}
    \FunctionTok{mutate}\NormalTok{(}
      \AttributeTok{agent\_data =} \FunctionTok{pmap}\NormalTok{(}\FunctionTok{list}\NormalTok{(wd, ws), }\ControlFlowTok{function}\NormalTok{(wd, ws) \{}
        \FunctionTok{generate\_agent\_decisions}\NormalTok{(wd, ws, evidence\_combinations, }\AttributeTok{n\_samples =} \DecValTok{5}\NormalTok{)}
\NormalTok{      \})}
\NormalTok{    )}

  
\NormalTok{  results\_list }\OtherTok{\textless{}{-}} \FunctionTok{list}\NormalTok{()}

  \ControlFlowTok{for}\NormalTok{ (i }\ControlFlowTok{in} \FunctionTok{seq\_len}\NormalTok{(}\FunctionTok{nrow}\NormalTok{(recovery\_data))) \{}

\NormalTok{    id       }\OtherTok{\textless{}{-}}\NormalTok{ recovery\_data}\SpecialCharTok{$}\NormalTok{sim\_id[[i]]}
\NormalTok{    true\_mod }\OtherTok{\textless{}{-}}\NormalTok{ recovery\_data}\SpecialCharTok{$}\NormalTok{true\_model[[i]]}
\NormalTok{    df       }\OtherTok{\textless{}{-}}\NormalTok{ recovery\_data}\SpecialCharTok{$}\NormalTok{agent\_data[[i]]}
\NormalTok{    true\_wd  }\OtherTok{\textless{}{-}}\NormalTok{ recovery\_data}\SpecialCharTok{$}\NormalTok{wd[[i]]}
\NormalTok{    true\_ws  }\OtherTok{\textless{}{-}}\NormalTok{ recovery\_data}\SpecialCharTok{$}\NormalTok{ws[[i]]}

\NormalTok{    stan\_data }\OtherTok{\textless{}{-}} \FunctionTok{prepare\_stan\_data}\NormalTok{(df)}

    \CommentTok{\# Fit all three models}
\NormalTok{    fit\_s }\OtherTok{\textless{}{-}}\NormalTok{ mod\_simple}\SpecialCharTok{$}\FunctionTok{sample}\NormalTok{(}\AttributeTok{data =}\NormalTok{ stan\_data,}
                               \AttributeTok{chains =} \DecValTok{2}\NormalTok{, }\AttributeTok{iter\_warmup =} \DecValTok{0}\NormalTok{, }\AttributeTok{iter\_sampling =} \DecValTok{500}\NormalTok{,}
                               \AttributeTok{fixed\_param =} \ConstantTok{TRUE}\NormalTok{, }\AttributeTok{refresh =} \DecValTok{0}\NormalTok{,}
                               \AttributeTok{show\_messages =} \ConstantTok{FALSE}\NormalTok{, }\AttributeTok{save\_warmup =} \ConstantTok{TRUE}\NormalTok{,}
                               \AttributeTok{output\_dir =} \StringTok{"models/temp"}\NormalTok{)}

\NormalTok{    fit\_p }\OtherTok{\textless{}{-}}\NormalTok{ mod\_proportional}\SpecialCharTok{$}\FunctionTok{sample}\NormalTok{(}\AttributeTok{data =}\NormalTok{ stan\_data,}
                                     \AttributeTok{chains =} \DecValTok{2}\NormalTok{, }\AttributeTok{iter\_warmup =} \DecValTok{500}\NormalTok{,}
                                     \AttributeTok{iter\_sampling =} \DecValTok{500}\NormalTok{, }\AttributeTok{refresh =} \DecValTok{0}\NormalTok{,}
                                     \AttributeTok{show\_messages =} \ConstantTok{FALSE}\NormalTok{, }\AttributeTok{save\_warmup =} \ConstantTok{TRUE}\NormalTok{,}
                                     \AttributeTok{output\_dir =} \StringTok{"models/temp"}\NormalTok{)}

\NormalTok{    fit\_w }\OtherTok{\textless{}{-}}\NormalTok{ mod\_weighted}\SpecialCharTok{$}\FunctionTok{sample}\NormalTok{(}\AttributeTok{data =}\NormalTok{ stan\_data,}
                                 \AttributeTok{chains =} \DecValTok{2}\NormalTok{, }\AttributeTok{iter\_warmup =} \DecValTok{500}\NormalTok{,}
                                 \AttributeTok{iter\_sampling =} \DecValTok{500}\NormalTok{, }\AttributeTok{refresh =} \DecValTok{0}\NormalTok{,}
                                 \AttributeTok{show\_messages =} \ConstantTok{FALSE}\NormalTok{, }\AttributeTok{save\_warmup =} \ConstantTok{TRUE}\NormalTok{,}
                                 \AttributeTok{output\_dir =} \StringTok{"models/temp"}\NormalTok{)}

\NormalTok{    loo\_s }\OtherTok{\textless{}{-}}\NormalTok{ fit\_s}\SpecialCharTok{$}\FunctionTok{loo}\NormalTok{()}
\NormalTok{    loo\_p }\OtherTok{\textless{}{-}}\NormalTok{ fit\_p}\SpecialCharTok{$}\FunctionTok{loo}\NormalTok{()}
\NormalTok{    loo\_w }\OtherTok{\textless{}{-}}\NormalTok{ fit\_w}\SpecialCharTok{$}\FunctionTok{loo}\NormalTok{()}

\NormalTok{    weights }\OtherTok{\textless{}{-}}\NormalTok{ loo}\SpecialCharTok{::}\FunctionTok{loo\_model\_weights}\NormalTok{(}
      \FunctionTok{list}\NormalTok{(}\AttributeTok{SBA =}\NormalTok{ loo\_s, }\AttributeTok{PBA =}\NormalTok{ loo\_p, }\AttributeTok{WBA =}\NormalTok{ loo\_w),}
      \AttributeTok{method =} \StringTok{"stacking"}
\NormalTok{    )}

    \CommentTok{\# Extract posterior means of effective weights (wd, ws) from each model}
    \CommentTok{\# SBA: deterministic (1, 1)}
\NormalTok{    pba\_draws }\OtherTok{\textless{}{-}}\NormalTok{ fit\_p}\SpecialCharTok{$}\FunctionTok{draws}\NormalTok{(}\AttributeTok{format =} \StringTok{"df"}\NormalTok{)}
\NormalTok{    wba\_draws }\OtherTok{\textless{}{-}}\NormalTok{ fit\_w}\SpecialCharTok{$}\FunctionTok{draws}\NormalTok{(}\AttributeTok{format =} \StringTok{"df"}\NormalTok{)}

\NormalTok{    results\_list[[i]] }\OtherTok{\textless{}{-}} \FunctionTok{tibble}\NormalTok{(}
      \AttributeTok{sim\_id     =}\NormalTok{ id,}
      \AttributeTok{true\_model =}\NormalTok{ true\_mod,}
      \AttributeTok{true\_wd    =}\NormalTok{ true\_wd,}
      \AttributeTok{true\_ws    =}\NormalTok{ true\_ws,}
      \CommentTok{\# Stacking weights}
      \AttributeTok{weight\_SBA =} \FunctionTok{as.numeric}\NormalTok{(weights[}\StringTok{"SBA"}\NormalTok{]),}
      \AttributeTok{weight\_PBA =} \FunctionTok{as.numeric}\NormalTok{(weights[}\StringTok{"PBA"}\NormalTok{]),}
      \AttributeTok{weight\_WBA =} \FunctionTok{as.numeric}\NormalTok{(weights[}\StringTok{"WBA"}\NormalTok{]),}
      \CommentTok{\# Posterior mean effective weights per candidate model}
      \AttributeTok{sba\_wd =} \DecValTok{1}\NormalTok{, }\AttributeTok{sba\_ws =} \DecValTok{1}\NormalTok{,}
      \AttributeTok{pba\_wd =} \FunctionTok{mean}\NormalTok{(pba\_draws}\SpecialCharTok{$}\NormalTok{p),}
      \AttributeTok{pba\_ws =} \DecValTok{1} \SpecialCharTok{{-}} \FunctionTok{mean}\NormalTok{(pba\_draws}\SpecialCharTok{$}\NormalTok{p),}
      \AttributeTok{wba\_wd =} \FunctionTok{mean}\NormalTok{(wba\_draws}\SpecialCharTok{$}\NormalTok{rho }\SpecialCharTok{*}\NormalTok{ wba\_draws}\SpecialCharTok{$}\NormalTok{kappa),}
      \AttributeTok{wba\_ws =} \FunctionTok{mean}\NormalTok{((}\DecValTok{1} \SpecialCharTok{{-}}\NormalTok{ wba\_draws}\SpecialCharTok{$}\NormalTok{rho) }\SpecialCharTok{*}\NormalTok{ wba\_draws}\SpecialCharTok{$}\NormalTok{kappa)}
\NormalTok{    )}
\NormalTok{  \}}

\NormalTok{  recovery\_results }\OtherTok{\textless{}{-}} \FunctionTok{bind\_rows}\NormalTok{(results\_list)}
  \FunctionTok{saveRDS}\NormalTok{(recovery\_results, }\StringTok{"simmodels/ch10\_model\_recovery.rds"}\NormalTok{)}

\NormalTok{\} }\ControlFlowTok{else}\NormalTok{ \{}
\NormalTok{  recovery\_results }\OtherTok{\textless{}{-}} \FunctionTok{readRDS}\NormalTok{(}\StringTok{"simmodels/ch10\_model\_recovery.rds"}\NormalTok{)}
\NormalTok{\}}
\end{Highlighting}
\end{Shaded}

\subsection{Confusion Matrix}\label{confusion-matrix}

The probabilistic confusion matrix shows, for each true generating
model, the average stacking weight assigned to each candidate. A perfect
recovery procedure would place all weight on the diagonal.

\begin{Shaded}
\begin{Highlighting}[]
\NormalTok{binary\_matrix }\OtherTok{\textless{}{-}}\NormalTok{ recovery\_results }\SpecialCharTok{|\textgreater{}}
  \CommentTok{\# 1. Reshape to make finding the maximum weight per simulation easier}
  \FunctionTok{pivot\_longer}\NormalTok{(}\AttributeTok{cols =} \FunctionTok{starts\_with}\NormalTok{(}\StringTok{"weight\_"}\NormalTok{), }
               \AttributeTok{names\_to =} \StringTok{"predicted\_model"}\NormalTok{, }\AttributeTok{values\_to =} \StringTok{"weight"}\NormalTok{) }\SpecialCharTok{|\textgreater{}}
  \FunctionTok{mutate}\NormalTok{(}\AttributeTok{predicted\_model =} \FunctionTok{str\_remove}\NormalTok{(predicted\_model, }\StringTok{"weight\_"}\NormalTok{)) }\SpecialCharTok{|\textgreater{}}
  
  \CommentTok{\# 2. Keep only the winning model for each simulation}
  \FunctionTok{group\_by}\NormalTok{(sim\_id) }\SpecialCharTok{|\textgreater{}}
  \FunctionTok{slice\_max}\NormalTok{(weight, }\AttributeTok{n =} \DecValTok{1}\NormalTok{, }\AttributeTok{with\_ties =} \ConstantTok{FALSE}\NormalTok{) }\SpecialCharTok{|\textgreater{}}
  \FunctionTok{ungroup}\NormalTok{() }\SpecialCharTok{|\textgreater{}}
  
  \CommentTok{\# 3. Count the wins and calculate proportions per true model}
  \FunctionTok{count}\NormalTok{(true\_model, predicted\_model) }\SpecialCharTok{|\textgreater{}}
  \FunctionTok{group\_by}\NormalTok{(true\_model) }\SpecialCharTok{|\textgreater{}}
  \FunctionTok{mutate}\NormalTok{(}\AttributeTok{proportion =}\NormalTok{ n }\SpecialCharTok{/} \FunctionTok{sum}\NormalTok{(n)) }\SpecialCharTok{|\textgreater{}}
  \FunctionTok{ungroup}\NormalTok{() }\SpecialCharTok{|\textgreater{}}
  
  \CommentTok{\# 4. Ensure we have all combinations factored correctly (handles 0\% scenarios)}
  \FunctionTok{mutate}\NormalTok{(}
    \AttributeTok{true\_model =} \FunctionTok{factor}\NormalTok{(true\_model, }\AttributeTok{levels =} \FunctionTok{c}\NormalTok{(}\StringTok{"SBA"}\NormalTok{, }\StringTok{"PBA"}\NormalTok{, }\StringTok{"WBA"}\NormalTok{)),}
    \AttributeTok{predicted\_model =} \FunctionTok{factor}\NormalTok{(predicted\_model, }\AttributeTok{levels =} \FunctionTok{c}\NormalTok{(}\StringTok{"SBA"}\NormalTok{, }\StringTok{"PBA"}\NormalTok{, }\StringTok{"WBA"}\NormalTok{))}
\NormalTok{  ) }\SpecialCharTok{|\textgreater{}}
  \FunctionTok{complete}\NormalTok{(true\_model, predicted\_model, }\AttributeTok{fill =} \FunctionTok{list}\NormalTok{(}\AttributeTok{n =} \DecValTok{0}\NormalTok{, }\AttributeTok{proportion =} \DecValTok{0}\NormalTok{))}

\CommentTok{\# Optional: View it as a standard table }
\NormalTok{binary\_matrix }\SpecialCharTok{|\textgreater{}}
  \FunctionTok{select}\NormalTok{(}\SpecialCharTok{{-}}\NormalTok{n) }\SpecialCharTok{|\textgreater{}}
  \FunctionTok{pivot\_wider}\NormalTok{(}\AttributeTok{names\_from =}\NormalTok{ predicted\_model, }\AttributeTok{values\_from =}\NormalTok{ proportion) }\SpecialCharTok{|\textgreater{}}
\NormalTok{  knitr}\SpecialCharTok{::}\FunctionTok{kable}\NormalTok{(}\AttributeTok{digits =} \DecValTok{2}\NormalTok{, }\AttributeTok{caption =} \StringTok{"Binary Confusion Matrix: Highest Weight Win Rate"}\NormalTok{)}
\end{Highlighting}
\end{Shaded}

\begin{longtable}[]{@{}lrrr@{}}
\caption{Binary Confusion Matrix: Highest Weight Win
Rate}\tabularnewline
\toprule\noalign{}
true\_model & SBA & PBA & WBA \\
\midrule\noalign{}
\endfirsthead
\toprule\noalign{}
true\_model & SBA & PBA & WBA \\
\midrule\noalign{}
\endhead
\bottomrule\noalign{}
\endlastfoot
SBA & 0.95 & 0.05 & 0.00 \\
PBA & 0.36 & 0.64 & 0.00 \\
WBA & 0.53 & 0.18 & 0.29 \\
\end{longtable}

\begin{Shaded}
\begin{Highlighting}[]
\NormalTok{binary\_matrix }\SpecialCharTok{|\textgreater{}}
  \FunctionTok{ggplot}\NormalTok{(}\FunctionTok{aes}\NormalTok{(}\AttributeTok{x =}\NormalTok{ predicted\_model, }\AttributeTok{y =}\NormalTok{ true\_model, }\AttributeTok{fill =}\NormalTok{ proportion)) }\SpecialCharTok{+}
  \FunctionTok{geom\_tile}\NormalTok{(}\AttributeTok{color =} \StringTok{"white"}\NormalTok{, }\AttributeTok{linewidth =} \DecValTok{1}\NormalTok{) }\SpecialCharTok{+}
  \FunctionTok{geom\_text}\NormalTok{(}\FunctionTok{aes}\NormalTok{(}\AttributeTok{label =}\NormalTok{ scales}\SpecialCharTok{::}\FunctionTok{percent}\NormalTok{(proportion, }\AttributeTok{accuracy =} \DecValTok{1}\NormalTok{)),}
            \AttributeTok{size =} \DecValTok{5}\NormalTok{, }\AttributeTok{fontface =} \StringTok{"bold"}\NormalTok{) }\SpecialCharTok{+}
  \FunctionTok{scale\_fill\_viridis\_c}\NormalTok{(}\AttributeTok{option =} \StringTok{"mako"}\NormalTok{, }\AttributeTok{direction =} \SpecialCharTok{{-}}\DecValTok{1}\NormalTok{, }\AttributeTok{limits =} \FunctionTok{c}\NormalTok{(}\DecValTok{0}\NormalTok{, }\DecValTok{1}\NormalTok{)) }\SpecialCharTok{+}
  \FunctionTok{scale\_y\_discrete}\NormalTok{(}\AttributeTok{limits =}\NormalTok{ rev) }\SpecialCharTok{+} \CommentTok{\# Keeps SBA at the top left}
  \FunctionTok{labs}\NormalTok{(}
    \AttributeTok{title =} \StringTok{"Model Recovery: Binary Classification"}\NormalTok{,}
    \AttributeTok{x =} \StringTok{"Winning Candidate Model (Max Weight)"}\NormalTok{, }
    \AttributeTok{y =} \StringTok{"True Generating Model"}\NormalTok{,}
    \AttributeTok{fill =} \StringTok{"Win Rate"}
\NormalTok{  ) }\SpecialCharTok{+}
  \FunctionTok{theme}\NormalTok{(}\AttributeTok{panel.grid =} \FunctionTok{element\_blank}\NormalTok{())}
\end{Highlighting}
\end{Shaded}

\begin{center}
\includegraphics[width=0.8\linewidth,height=\textheight,keepaspectratio]{10-BayesianModels_files/figure-pdf/ch10_model_recovery_confusion-1.pdf}
\end{center}

\begin{Shaded}
\begin{Highlighting}[]
\NormalTok{confusion\_matrix }\OtherTok{\textless{}{-}}\NormalTok{ recovery\_results }\SpecialCharTok{|\textgreater{}}
  \FunctionTok{group\_by}\NormalTok{(true\_model) }\SpecialCharTok{|\textgreater{}}
  \FunctionTok{summarise}\NormalTok{(}
    \AttributeTok{Recovered\_SBA =} \FunctionTok{mean}\NormalTok{(weight\_SBA),}
    \AttributeTok{Recovered\_PBA =} \FunctionTok{mean}\NormalTok{(weight\_PBA),}
    \AttributeTok{Recovered\_WBA =} \FunctionTok{mean}\NormalTok{(weight\_WBA),}
    \AttributeTok{.groups =} \StringTok{"drop"}
\NormalTok{  ) }\SpecialCharTok{|\textgreater{}}
  \FunctionTok{mutate}\NormalTok{(}\AttributeTok{true\_model =} \FunctionTok{factor}\NormalTok{(true\_model, }\AttributeTok{levels =} \FunctionTok{c}\NormalTok{(}\StringTok{"SBA"}\NormalTok{, }\StringTok{"PBA"}\NormalTok{, }\StringTok{"WBA"}\NormalTok{))) }\SpecialCharTok{|\textgreater{}}
  \FunctionTok{arrange}\NormalTok{(true\_model)}

\NormalTok{knitr}\SpecialCharTok{::}\FunctionTok{kable}\NormalTok{(confusion\_matrix, }\AttributeTok{digits =} \DecValTok{2}\NormalTok{,}
             \AttributeTok{caption =} \StringTok{"Probabilistic Confusion Matrix: Average Stacking Weights"}\NormalTok{)}
\end{Highlighting}
\end{Shaded}

\begin{longtable}[]{@{}lrrr@{}}
\caption{Probabilistic Confusion Matrix: Average Stacking
Weights}\tabularnewline
\toprule\noalign{}
true\_model & Recovered\_SBA & Recovered\_PBA & Recovered\_WBA \\
\midrule\noalign{}
\endfirsthead
\toprule\noalign{}
true\_model & Recovered\_SBA & Recovered\_PBA & Recovered\_WBA \\
\midrule\noalign{}
\endhead
\bottomrule\noalign{}
\endlastfoot
SBA & 0.79 & 0.13 & 0.07 \\
PBA & 0.34 & 0.52 & 0.15 \\
WBA & 0.39 & 0.21 & 0.39 \\
\end{longtable}

\begin{Shaded}
\begin{Highlighting}[]
\NormalTok{confusion\_matrix }\SpecialCharTok{|\textgreater{}}
  \FunctionTok{pivot\_longer}\NormalTok{(}\FunctionTok{starts\_with}\NormalTok{(}\StringTok{"Recovered"}\NormalTok{),}
               \AttributeTok{names\_to =} \StringTok{"fitted\_model"}\NormalTok{, }\AttributeTok{values\_to =} \StringTok{"weight"}\NormalTok{) }\SpecialCharTok{|\textgreater{}}
  \FunctionTok{mutate}\NormalTok{(}
    \AttributeTok{fitted\_model =} \FunctionTok{str\_remove}\NormalTok{(fitted\_model, }\StringTok{"Recovered\_"}\NormalTok{),}
    \CommentTok{\# FIX 1: Lock in the horizontal order by converting to a factor}
    \AttributeTok{fitted\_model =} \FunctionTok{factor}\NormalTok{(fitted\_model, }\AttributeTok{levels =} \FunctionTok{c}\NormalTok{(}\StringTok{"SBA"}\NormalTok{, }\StringTok{"PBA"}\NormalTok{, }\StringTok{"WBA"}\NormalTok{))}
\NormalTok{  ) }\SpecialCharTok{|\textgreater{}}
  \FunctionTok{ggplot}\NormalTok{(}\FunctionTok{aes}\NormalTok{(}\AttributeTok{x =}\NormalTok{ fitted\_model, }\AttributeTok{y =}\NormalTok{ true\_model, }\AttributeTok{fill =}\NormalTok{ weight)) }\SpecialCharTok{+}
  \FunctionTok{geom\_tile}\NormalTok{(}\AttributeTok{color =} \StringTok{"white"}\NormalTok{, }\AttributeTok{linewidth =} \DecValTok{1}\NormalTok{) }\SpecialCharTok{+}
  \FunctionTok{geom\_text}\NormalTok{(}\FunctionTok{aes}\NormalTok{(}\AttributeTok{label =}\NormalTok{ scales}\SpecialCharTok{::}\FunctionTok{percent}\NormalTok{(weight, }\AttributeTok{accuracy =} \DecValTok{1}\NormalTok{)),}
            \AttributeTok{size =} \DecValTok{5}\NormalTok{, }\AttributeTok{fontface =} \StringTok{"bold"}\NormalTok{) }\SpecialCharTok{+}
  \FunctionTok{scale\_fill\_viridis\_c}\NormalTok{(}\AttributeTok{option =} \StringTok{"mako"}\NormalTok{, }\AttributeTok{direction =} \SpecialCharTok{{-}}\DecValTok{1}\NormalTok{, }\AttributeTok{limits =} \FunctionTok{c}\NormalTok{(}\DecValTok{0}\NormalTok{, }\DecValTok{1}\NormalTok{)) }\SpecialCharTok{+}
  \CommentTok{\# FIX 2: Reverse the vertical axis so "SBA" is at the top left}
  \FunctionTok{scale\_y\_discrete}\NormalTok{(}\AttributeTok{limits =}\NormalTok{ rev) }\SpecialCharTok{+} 
  \FunctionTok{labs}\NormalTok{(}
    \AttributeTok{title =} \StringTok{"Model Recovery: Predictive Weight Allocation"}\NormalTok{,}
    \AttributeTok{x =} \StringTok{"Candidate Model"}\NormalTok{, }\AttributeTok{y =} \StringTok{"True Generating Model"}\NormalTok{,}
    \AttributeTok{fill =} \StringTok{"Avg. Weight"}
\NormalTok{  ) }\SpecialCharTok{+}
  \FunctionTok{theme}\NormalTok{(}\AttributeTok{panel.grid =} \FunctionTok{element\_blank}\NormalTok{())}
\end{Highlighting}
\end{Shaded}

\begin{center}
\includegraphics[width=0.8\linewidth,height=\textheight,keepaspectratio]{10-BayesianModels_files/figure-pdf/ch10_model_recovery_confusion-2.pdf}
\end{center}

\subsection{Model Calibration}\label{model-calibration}

A well-calibrated model comparison procedure should assign a stacking
weight of, say, 0.7 to a model that turns out to be the true generator
roughly 70\% of the time. We assess this by binning stacking weights for
each candidate and computing the empirical frequency with which that
candidate was, in fact, the true model. Points on the diagonal indicate
perfect calibration; a curve bowed above the diagonal signals
underconfidence (the procedure ``hedges'' too much), while a curve below
signals overconfidence.

\begin{Shaded}
\begin{Highlighting}[]
\CommentTok{\# Reshape to long: one row per (simulation × candidate model)}
\NormalTok{calibration\_long }\OtherTok{\textless{}{-}}\NormalTok{ recovery\_results }\SpecialCharTok{|\textgreater{}}
  \FunctionTok{pivot\_longer}\NormalTok{(}
    \AttributeTok{cols =} \FunctionTok{starts\_with}\NormalTok{(}\StringTok{"weight\_"}\NormalTok{),}
    \AttributeTok{names\_to =} \StringTok{"candidate"}\NormalTok{, }\AttributeTok{values\_to =} \StringTok{"stacking\_weight"}
\NormalTok{  ) }\SpecialCharTok{|\textgreater{}}
  \FunctionTok{mutate}\NormalTok{(}
    \AttributeTok{candidate =} \FunctionTok{str\_remove}\NormalTok{(candidate, }\StringTok{"weight\_"}\NormalTok{),}
    \AttributeTok{is\_true    =} \FunctionTok{as.integer}\NormalTok{(candidate }\SpecialCharTok{==}\NormalTok{ true\_model)}
\NormalTok{  )}

\CommentTok{\# Bin stacking weights and compute empirical frequency}
\NormalTok{n\_bins }\OtherTok{\textless{}{-}} \DecValTok{10}
\NormalTok{calibration\_binned }\OtherTok{\textless{}{-}}\NormalTok{ calibration\_long }\SpecialCharTok{|\textgreater{}}
  \FunctionTok{mutate}\NormalTok{(}\AttributeTok{bin =} \FunctionTok{cut}\NormalTok{(stacking\_weight,}
                   \AttributeTok{breaks =} \FunctionTok{seq}\NormalTok{(}\DecValTok{0}\NormalTok{, }\DecValTok{1}\NormalTok{, }\AttributeTok{length.out =}\NormalTok{ n\_bins }\SpecialCharTok{+} \DecValTok{1}\NormalTok{),}
                   \AttributeTok{include.lowest =} \ConstantTok{TRUE}\NormalTok{)) }\SpecialCharTok{|\textgreater{}}
  \FunctionTok{group\_by}\NormalTok{(candidate, bin) }\SpecialCharTok{|\textgreater{}}
  \FunctionTok{summarise}\NormalTok{(}
    \AttributeTok{bin\_mid    =} \FunctionTok{mean}\NormalTok{(stacking\_weight),}
    \AttributeTok{observed   =} \FunctionTok{mean}\NormalTok{(is\_true),}
    \AttributeTok{N\_datasets =} \FunctionTok{n}\NormalTok{(),}
    \AttributeTok{.groups    =} \StringTok{"drop"}
\NormalTok{  ) }\SpecialCharTok{|\textgreater{}}
  \FunctionTok{mutate}\NormalTok{(}
    \AttributeTok{SE    =} \FunctionTok{sqrt}\NormalTok{((observed }\SpecialCharTok{*}\NormalTok{ (}\DecValTok{1} \SpecialCharTok{{-}}\NormalTok{ observed)) }\SpecialCharTok{/}\NormalTok{ N\_datasets),}
    \AttributeTok{Lower =} \FunctionTok{pmax}\NormalTok{(}\DecValTok{0}\NormalTok{, observed }\SpecialCharTok{{-}} \FloatTok{1.96} \SpecialCharTok{*}\NormalTok{ SE),}
    \AttributeTok{Upper =} \FunctionTok{pmin}\NormalTok{(}\DecValTok{1}\NormalTok{, observed }\SpecialCharTok{+} \FloatTok{1.96} \SpecialCharTok{*}\NormalTok{ SE)}
\NormalTok{  )}

\FunctionTok{ggplot}\NormalTok{(calibration\_binned, }\FunctionTok{aes}\NormalTok{(}\AttributeTok{x =}\NormalTok{ bin\_mid, }\AttributeTok{y =}\NormalTok{ observed, }\AttributeTok{color =}\NormalTok{ candidate)) }\SpecialCharTok{+}
  \FunctionTok{geom\_abline}\NormalTok{(}\AttributeTok{slope =} \DecValTok{1}\NormalTok{, }\AttributeTok{intercept =} \DecValTok{0}\NormalTok{, }\AttributeTok{linetype =} \StringTok{"dashed"}\NormalTok{, }\AttributeTok{color =} \StringTok{"gray40"}\NormalTok{) }\SpecialCharTok{+}
  \FunctionTok{geom\_ribbon}\NormalTok{(}\FunctionTok{aes}\NormalTok{(}\AttributeTok{ymin =}\NormalTok{ Lower, }\AttributeTok{ymax =}\NormalTok{ Upper, }\AttributeTok{fill =}\NormalTok{ candidate),}
              \AttributeTok{alpha =} \FloatTok{0.15}\NormalTok{, }\AttributeTok{color =} \ConstantTok{NA}\NormalTok{) }\SpecialCharTok{+}
  \FunctionTok{geom\_line}\NormalTok{(}\AttributeTok{linewidth =} \DecValTok{1}\NormalTok{) }\SpecialCharTok{+}
  \FunctionTok{geom\_point}\NormalTok{(}\AttributeTok{size =} \DecValTok{2}\NormalTok{) }\SpecialCharTok{+}
  \FunctionTok{scale\_color\_brewer}\NormalTok{(}\AttributeTok{palette =} \StringTok{"Set2"}\NormalTok{) }\SpecialCharTok{+}
  \FunctionTok{scale\_fill\_brewer}\NormalTok{(}\AttributeTok{palette =} \StringTok{"Set2"}\NormalTok{) }\SpecialCharTok{+}
  \FunctionTok{scale\_x\_continuous}\NormalTok{(}\AttributeTok{labels =}\NormalTok{ scales}\SpecialCharTok{::}\NormalTok{percent, }\AttributeTok{limits =} \FunctionTok{c}\NormalTok{(}\DecValTok{0}\NormalTok{, }\DecValTok{1}\NormalTok{)) }\SpecialCharTok{+}
  \FunctionTok{scale\_y\_continuous}\NormalTok{(}\AttributeTok{labels =}\NormalTok{ scales}\SpecialCharTok{::}\NormalTok{percent, }\AttributeTok{limits =} \FunctionTok{c}\NormalTok{(}\DecValTok{0}\NormalTok{, }\DecValTok{1}\NormalTok{)) }\SpecialCharTok{+}
  \FunctionTok{labs}\NormalTok{(}
    \AttributeTok{title =} \StringTok{"Model Comparison Calibration"}\NormalTok{,}
    \AttributeTok{subtitle =} \StringTok{"Ribbons show 95\% CI; wider ribbons reflect fewer datasets in that weight bin."}\NormalTok{,}
    \AttributeTok{x =} \StringTok{"Mean Stacking Weight (Binned)"}\NormalTok{,}
    \AttributeTok{y =} \StringTok{"Empirical Frequency (True Generating Model)"}\NormalTok{,}
    \AttributeTok{color =} \StringTok{"Candidate"}\NormalTok{, }\AttributeTok{fill =} \StringTok{"Candidate"}
\NormalTok{  ) }\SpecialCharTok{+}
  \FunctionTok{theme}\NormalTok{(}\AttributeTok{legend.position =} \StringTok{"bottom"}\NormalTok{)}
\end{Highlighting}
\end{Shaded}

\begin{center}
\includegraphics[width=0.8\linewidth,height=\textheight,keepaspectratio]{10-BayesianModels_files/figure-pdf/ch10_model_recovery_calibration-1.pdf}
\end{center}

The line goes up, but it's deviations from the diagonal are worrying
(and resonate with the not ideal confusion matrix). We should probably
run more simulations to make sure, but now we want to move on looking
more closely to misclassifications .

\subsection{Parameter Recovery Across
Models}\label{parameter-recovery-across-models}

Even when the ``wrong'' model wins, its parameters may still capture the
effective behavior of the true generating process. To check this, we
compute a \textbf{model-averaged estimate} of the effective evidence
weights (\(w_d\), \(w_s\)) by weighting each candidate's posterior means
by its stacking weight. If a wrongly-selected model is still
behaviorally faithful, these averaged estimates should track the true
values closely.

\begin{Shaded}
\begin{Highlighting}[]
\CommentTok{\# 1. Determine the winning model, enforce factor levels, and extract ITS specific fitted parameters}
\NormalTok{param\_recovery }\OtherTok{\textless{}{-}}\NormalTok{ recovery\_results }\SpecialCharTok{|\textgreater{}}
  \FunctionTok{mutate}\NormalTok{(}
    \CommentTok{\# Identify the winner based on the highest weight}
    \AttributeTok{winner =} \FunctionTok{case\_when}\NormalTok{(}
\NormalTok{      weight\_SBA }\SpecialCharTok{\textgreater{}=}\NormalTok{ weight\_PBA }\SpecialCharTok{\&}\NormalTok{ weight\_SBA }\SpecialCharTok{\textgreater{}=}\NormalTok{ weight\_WBA }\SpecialCharTok{\textasciitilde{}} \StringTok{"SBA"}\NormalTok{,}
\NormalTok{      weight\_PBA }\SpecialCharTok{\textgreater{}=}\NormalTok{ weight\_SBA }\SpecialCharTok{\&}\NormalTok{ weight\_PBA }\SpecialCharTok{\textgreater{}=}\NormalTok{ weight\_WBA }\SpecialCharTok{\textasciitilde{}} \StringTok{"PBA"}\NormalTok{,}
      \ConstantTok{TRUE} \SpecialCharTok{\textasciitilde{}} \StringTok{"WBA"}
\NormalTok{    ),}
    \CommentTok{\# Enforce factor levels for consistent plotting order (SBA {-}\textgreater{} PBA {-}\textgreater{} WBA)}
    \AttributeTok{winner =} \FunctionTok{factor}\NormalTok{(winner, }\AttributeTok{levels =} \FunctionTok{c}\NormalTok{(}\StringTok{"SBA"}\NormalTok{, }\StringTok{"PBA"}\NormalTok{, }\StringTok{"WBA"}\NormalTok{)),}
    \AttributeTok{true\_model =} \FunctionTok{factor}\NormalTok{(true\_model, }\AttributeTok{levels =} \FunctionTok{c}\NormalTok{(}\StringTok{"SBA"}\NormalTok{, }\StringTok{"PBA"}\NormalTok{, }\StringTok{"WBA"}\NormalTok{)),}
    \AttributeTok{correct =}\NormalTok{ winner }\SpecialCharTok{==}\NormalTok{ true\_model,}
    
    \CommentTok{\# Extract the actual parameters estimated by the WINNING model}
    \AttributeTok{win\_wd =} \FunctionTok{case\_when}\NormalTok{(}
\NormalTok{      winner }\SpecialCharTok{==} \StringTok{"SBA"} \SpecialCharTok{\textasciitilde{}}\NormalTok{ sba\_wd, }\CommentTok{\# Will likely be NA/0 if SBA is 0{-}parameter}
\NormalTok{      winner }\SpecialCharTok{==} \StringTok{"PBA"} \SpecialCharTok{\textasciitilde{}}\NormalTok{ pba\_wd,}
\NormalTok{      winner }\SpecialCharTok{==} \StringTok{"WBA"} \SpecialCharTok{\textasciitilde{}}\NormalTok{ wba\_wd}
\NormalTok{    ),}
    \AttributeTok{win\_ws =} \FunctionTok{case\_when}\NormalTok{(}
\NormalTok{      winner }\SpecialCharTok{==} \StringTok{"SBA"} \SpecialCharTok{\textasciitilde{}}\NormalTok{ sba\_ws,}
\NormalTok{      winner }\SpecialCharTok{==} \StringTok{"PBA"} \SpecialCharTok{\textasciitilde{}}\NormalTok{ pba\_ws,}
\NormalTok{      winner }\SpecialCharTok{==} \StringTok{"WBA"} \SpecialCharTok{\textasciitilde{}}\NormalTok{ wba\_ws}
\NormalTok{    )}
\NormalTok{  )}

\CommentTok{\# {-}{-}{-}{-}{-}{-}{-}{-}{-}{-}{-}{-}{-}{-}{-}{-}{-}{-}{-}{-}{-}{-}{-}{-}{-}{-}{-}{-}{-}{-}{-}{-}{-}{-}{-}{-}{-}{-}{-}{-}{-}{-}{-}{-}{-}{-}{-}{-}{-}{-}{-}{-}{-}{-}{-}{-}{-}}
\CommentTok{\# VISUALIZATION 1: Winning Model Parameter Recovery}
\CommentTok{\# {-}{-}{-}{-}{-}{-}{-}{-}{-}{-}{-}{-}{-}{-}{-}{-}{-}{-}{-}{-}{-}{-}{-}{-}{-}{-}{-}{-}{-}{-}{-}{-}{-}{-}{-}{-}{-}{-}{-}{-}{-}{-}{-}{-}{-}{-}{-}{-}{-}{-}{-}{-}{-}{-}{-}{-}{-}}

\NormalTok{p\_wd }\OtherTok{\textless{}{-}} \FunctionTok{ggplot}\NormalTok{(param\_recovery, }\FunctionTok{aes}\NormalTok{(}\AttributeTok{x =}\NormalTok{ true\_wd, }\AttributeTok{y =}\NormalTok{ win\_wd)) }\SpecialCharTok{+}
  \FunctionTok{geom\_abline}\NormalTok{(}\AttributeTok{slope =} \DecValTok{1}\NormalTok{, }\AttributeTok{intercept =} \DecValTok{0}\NormalTok{, }\AttributeTok{linetype =} \StringTok{"dashed"}\NormalTok{, }\AttributeTok{color =} \StringTok{"gray40"}\NormalTok{) }\SpecialCharTok{+}
  \FunctionTok{geom\_point}\NormalTok{(}\FunctionTok{aes}\NormalTok{(}\AttributeTok{color =}\NormalTok{ true\_model, }\AttributeTok{shape =}\NormalTok{ correct), }\AttributeTok{size =} \FloatTok{2.5}\NormalTok{, }\AttributeTok{alpha =} \FloatTok{0.7}\NormalTok{) }\SpecialCharTok{+}
  \FunctionTok{geom\_smooth}\NormalTok{(}\AttributeTok{method=}\NormalTok{lm, }\FunctionTok{aes}\NormalTok{(}\AttributeTok{color =}\NormalTok{ true\_model), }\AttributeTok{se =} \ConstantTok{FALSE}\NormalTok{, }\AttributeTok{linetype =} \StringTok{"dotted"}\NormalTok{) }\SpecialCharTok{+}
  \FunctionTok{scale\_shape\_manual}\NormalTok{(}\AttributeTok{values =} \FunctionTok{c}\NormalTok{(}\StringTok{"TRUE"} \OtherTok{=} \DecValTok{16}\NormalTok{, }\StringTok{"FALSE"} \OtherTok{=} \DecValTok{4}\NormalTok{),}
                     \AttributeTok{labels =} \FunctionTok{c}\NormalTok{(}\StringTok{"TRUE"} \OtherTok{=} \StringTok{"Correct"}\NormalTok{, }\StringTok{"FALSE"} \OtherTok{=} \StringTok{"Wrong"}\NormalTok{)) }\SpecialCharTok{+}
  \FunctionTok{scale\_color\_brewer}\NormalTok{(}\AttributeTok{palette =} \StringTok{"Set2"}\NormalTok{) }\SpecialCharTok{+}
  \FunctionTok{labs}\NormalTok{(}
    \AttributeTok{title =} \FunctionTok{expression}\NormalTok{(}\StringTok{"Recovery of "} \SpecialCharTok{*}\NormalTok{ w[d]),}
    \AttributeTok{x =} \FunctionTok{expression}\NormalTok{(}\StringTok{"True "} \SpecialCharTok{*}\NormalTok{ w[d]),}
    \AttributeTok{y =} \FunctionTok{expression}\NormalTok{(}\StringTok{"Winning Model "} \SpecialCharTok{*} \FunctionTok{hat}\NormalTok{(w)[d]),}
    \AttributeTok{color =} \StringTok{"True Model"}\NormalTok{, }\AttributeTok{shape =} \StringTok{"Recovery"}
\NormalTok{  )}

\NormalTok{p\_ws }\OtherTok{\textless{}{-}} \FunctionTok{ggplot}\NormalTok{(param\_recovery, }\FunctionTok{aes}\NormalTok{(}\AttributeTok{x =}\NormalTok{ true\_ws, }\AttributeTok{y =}\NormalTok{ win\_ws)) }\SpecialCharTok{+}
  \FunctionTok{geom\_abline}\NormalTok{(}\AttributeTok{slope =} \DecValTok{1}\NormalTok{, }\AttributeTok{intercept =} \DecValTok{0}\NormalTok{, }\AttributeTok{linetype =} \StringTok{"dashed"}\NormalTok{, }\AttributeTok{color =} \StringTok{"gray40"}\NormalTok{) }\SpecialCharTok{+}
  \FunctionTok{geom\_point}\NormalTok{(}\FunctionTok{aes}\NormalTok{(}\AttributeTok{color =}\NormalTok{ true\_model, }\AttributeTok{shape =}\NormalTok{ correct), }\AttributeTok{size =} \FloatTok{2.5}\NormalTok{, }\AttributeTok{alpha =} \FloatTok{0.7}\NormalTok{) }\SpecialCharTok{+}
  \FunctionTok{geom\_smooth}\NormalTok{(}\AttributeTok{method =}\NormalTok{ lm, }\FunctionTok{aes}\NormalTok{(}\AttributeTok{color =}\NormalTok{ true\_model), }\AttributeTok{se =} \ConstantTok{FALSE}\NormalTok{, }\AttributeTok{linetype =} \StringTok{"dotted"}\NormalTok{) }\SpecialCharTok{+}
  \FunctionTok{scale\_shape\_manual}\NormalTok{(}\AttributeTok{values =} \FunctionTok{c}\NormalTok{(}\StringTok{"TRUE"} \OtherTok{=} \DecValTok{16}\NormalTok{, }\StringTok{"FALSE"} \OtherTok{=} \DecValTok{4}\NormalTok{),}
                     \AttributeTok{labels =} \FunctionTok{c}\NormalTok{(}\StringTok{"TRUE"} \OtherTok{=} \StringTok{"Correct"}\NormalTok{, }\StringTok{"FALSE"} \OtherTok{=} \StringTok{"Wrong"}\NormalTok{)) }\SpecialCharTok{+}
  \FunctionTok{scale\_color\_brewer}\NormalTok{(}\AttributeTok{palette =} \StringTok{"Set2"}\NormalTok{) }\SpecialCharTok{+}
  \FunctionTok{labs}\NormalTok{(}
    \AttributeTok{title =} \FunctionTok{expression}\NormalTok{(}\StringTok{"Recovery of "} \SpecialCharTok{*}\NormalTok{ w[s]),}
    \AttributeTok{x =} \FunctionTok{expression}\NormalTok{(}\StringTok{"True "} \SpecialCharTok{*}\NormalTok{ w[s]),}
    \AttributeTok{y =} \FunctionTok{expression}\NormalTok{(}\StringTok{"Winning Model "} \SpecialCharTok{*} \FunctionTok{hat}\NormalTok{(w)[s]),}
    \AttributeTok{color =} \StringTok{"True Model"}\NormalTok{, }\AttributeTok{shape =} \StringTok{"Recovery"}
\NormalTok{  )}

\CommentTok{\# View Combined Recovery Plot}
\NormalTok{p\_recovery }\OtherTok{\textless{}{-}}\NormalTok{ p\_wd }\SpecialCharTok{+}\NormalTok{ p\_ws }\SpecialCharTok{+} \FunctionTok{plot\_layout}\NormalTok{(}\AttributeTok{guides =} \StringTok{"collect"}\NormalTok{) }\SpecialCharTok{\&}
  \FunctionTok{theme}\NormalTok{(}\AttributeTok{legend.position =} \StringTok{"bottom"}\NormalTok{)}
\NormalTok{p\_recovery}
\end{Highlighting}
\end{Shaded}

\begin{center}
\includegraphics[width=0.8\linewidth,height=\textheight,keepaspectratio]{10-BayesianModels_files/figure-pdf/ch10_model_recovery_parameters-1.pdf}
\end{center}

\begin{Shaded}
\begin{Highlighting}[]
\CommentTok{\# {-}{-}{-}{-}{-}{-}{-}{-}{-}{-}{-}{-}{-}{-}{-}{-}{-}{-}{-}{-}{-}{-}{-}{-}{-}{-}{-}{-}{-}{-}{-}{-}{-}{-}{-}{-}{-}{-}{-}{-}{-}{-}{-}{-}{-}{-}{-}{-}{-}{-}{-}{-}{-}{-}{-}{-}{-}}
\CommentTok{\# VISUALIZATION 2: True vs. Fitted Parameter Spaces}
\CommentTok{\# {-}{-}{-}{-}{-}{-}{-}{-}{-}{-}{-}{-}{-}{-}{-}{-}{-}{-}{-}{-}{-}{-}{-}{-}{-}{-}{-}{-}{-}{-}{-}{-}{-}{-}{-}{-}{-}{-}{-}{-}{-}{-}{-}{-}{-}{-}{-}{-}{-}{-}{-}{-}{-}{-}{-}{-}{-}}

\CommentTok{\# Plot A: True Parameter Space (Where do misclassifications happen?)}
\NormalTok{p\_true }\OtherTok{\textless{}{-}} \FunctionTok{ggplot}\NormalTok{(param\_recovery, }\FunctionTok{aes}\NormalTok{(}\AttributeTok{x =}\NormalTok{ true\_wd, }\AttributeTok{y =}\NormalTok{ true\_ws)) }\SpecialCharTok{+}
  \FunctionTok{geom\_abline}\NormalTok{(}\AttributeTok{slope =} \DecValTok{1}\NormalTok{, }\AttributeTok{intercept =} \DecValTok{0}\NormalTok{, }\AttributeTok{linetype =} \StringTok{"dashed"}\NormalTok{, }\AttributeTok{color =} \StringTok{"gray60"}\NormalTok{) }\SpecialCharTok{+} 
  \FunctionTok{geom\_point}\NormalTok{(}\FunctionTok{aes}\NormalTok{(}\AttributeTok{color =}\NormalTok{ winner, }\AttributeTok{shape =}\NormalTok{ correct), }\AttributeTok{size =} \FloatTok{2.5}\NormalTok{, }\AttributeTok{alpha =} \FloatTok{0.7}\NormalTok{) }\SpecialCharTok{+}
  \FunctionTok{facet\_wrap}\NormalTok{(}\SpecialCharTok{\textasciitilde{}}\NormalTok{ true\_model, }\AttributeTok{labeller =} \FunctionTok{as\_labeller}\NormalTok{(}\ControlFlowTok{function}\NormalTok{(x) }\FunctionTok{paste}\NormalTok{(}\StringTok{"True:"}\NormalTok{, x))) }\SpecialCharTok{+}
  \FunctionTok{scale\_shape\_manual}\NormalTok{(}\AttributeTok{values =} \FunctionTok{c}\NormalTok{(}\StringTok{"TRUE"} \OtherTok{=} \DecValTok{16}\NormalTok{, }\StringTok{"FALSE"} \OtherTok{=} \DecValTok{4}\NormalTok{), }\AttributeTok{guide =} \StringTok{"none"}\NormalTok{) }\SpecialCharTok{+}
  \FunctionTok{scale\_color\_brewer}\NormalTok{(}\AttributeTok{palette =} \StringTok{"Set1"}\NormalTok{) }\SpecialCharTok{+}
  \FunctionTok{labs}\NormalTok{(}
    \AttributeTok{title =} \StringTok{"A. Data Generating Space"}\NormalTok{,}
    \AttributeTok{subtitle =} \StringTok{"Which true parameter regions cause the wrong model to win?"}\NormalTok{,}
    \AttributeTok{x =} \FunctionTok{expression}\NormalTok{(}\StringTok{"True "} \SpecialCharTok{*}\NormalTok{ w[d]),}
    \AttributeTok{y =} \FunctionTok{expression}\NormalTok{(}\StringTok{"True "} \SpecialCharTok{*}\NormalTok{ w[s]),}
    \AttributeTok{color =} \StringTok{"Winning Model"}
\NormalTok{  )}

\CommentTok{\# Plot B: Fitted Parameter Space (How do models mimic each other?)}
\NormalTok{p\_fit }\OtherTok{\textless{}{-}} \FunctionTok{ggplot}\NormalTok{(param\_recovery, }\FunctionTok{aes}\NormalTok{(}\AttributeTok{x =}\NormalTok{ win\_wd, }\AttributeTok{y =}\NormalTok{ win\_ws)) }\SpecialCharTok{+}
  \FunctionTok{geom\_abline}\NormalTok{(}\AttributeTok{slope =} \DecValTok{1}\NormalTok{, }\AttributeTok{intercept =} \DecValTok{0}\NormalTok{, }\AttributeTok{linetype =} \StringTok{"dashed"}\NormalTok{, }\AttributeTok{color =} \StringTok{"gray60"}\NormalTok{) }\SpecialCharTok{+}
  \FunctionTok{geom\_point}\NormalTok{(}\FunctionTok{aes}\NormalTok{(}\AttributeTok{color =}\NormalTok{ true\_model, }\AttributeTok{shape =}\NormalTok{ correct), }\AttributeTok{size =} \FloatTok{2.5}\NormalTok{, }\AttributeTok{alpha =} \FloatTok{0.7}\NormalTok{) }\SpecialCharTok{+}
  \FunctionTok{facet\_wrap}\NormalTok{(}\SpecialCharTok{\textasciitilde{}}\NormalTok{ winner, }\AttributeTok{labeller =} \FunctionTok{as\_labeller}\NormalTok{(}\ControlFlowTok{function}\NormalTok{(x) }\FunctionTok{paste}\NormalTok{(}\StringTok{"Winner:"}\NormalTok{, x))) }\SpecialCharTok{+}
  \FunctionTok{scale\_shape\_manual}\NormalTok{(}\AttributeTok{values =} \FunctionTok{c}\NormalTok{(}\StringTok{"TRUE"} \OtherTok{=} \DecValTok{16}\NormalTok{, }\StringTok{"FALSE"} \OtherTok{=} \DecValTok{4}\NormalTok{)) }\SpecialCharTok{+}
  \FunctionTok{scale\_color\_brewer}\NormalTok{(}\AttributeTok{palette =} \StringTok{"Dark2"}\NormalTok{) }\SpecialCharTok{+}
  \FunctionTok{labs}\NormalTok{(}
    \AttributeTok{title =} \StringTok{"B. Fitted Parameter Space"}\NormalTok{,}
    \AttributeTok{subtitle =} \StringTok{"What parameters did the winning model actually estimate?"}\NormalTok{,}
    \AttributeTok{x =} \FunctionTok{expression}\NormalTok{(}\StringTok{"Fitted "} \SpecialCharTok{*} \FunctionTok{hat}\NormalTok{(w)[d]),}
    \AttributeTok{y =} \FunctionTok{expression}\NormalTok{(}\StringTok{"Fitted "} \SpecialCharTok{*} \FunctionTok{hat}\NormalTok{(w)[s]),}
    \AttributeTok{color =} \StringTok{"True Model"}\NormalTok{, }\AttributeTok{shape =} \StringTok{"Recovery"}
\NormalTok{  )}

\CommentTok{\# View Combined Spaces Plot}
\NormalTok{p\_true }
\end{Highlighting}
\end{Shaded}

\begin{center}
\includegraphics[width=0.8\linewidth,height=\textheight,keepaspectratio]{10-BayesianModels_files/figure-pdf/ch10_model_recovery_parameters-2.pdf}
\end{center}

\begin{Shaded}
\begin{Highlighting}[]
\NormalTok{p\_fit}
\end{Highlighting}
\end{Shaded}

\begin{center}
\includegraphics[width=0.8\linewidth,height=\textheight,keepaspectratio]{10-BayesianModels_files/figure-pdf/ch10_model_recovery_parameters-3.pdf}
\end{center}

\begin{Shaded}
\begin{Highlighting}[]
\NormalTok{param\_recovery }\OtherTok{\textless{}{-}}\NormalTok{ param\_recovery }\SpecialCharTok{|\textgreater{}}
  \FunctionTok{mutate}\NormalTok{(}
    \AttributeTok{true\_rho   =}\NormalTok{ true\_wd }\SpecialCharTok{/}\NormalTok{ (true\_wd }\SpecialCharTok{+}\NormalTok{ true\_ws),}
    \AttributeTok{true\_kappa =}\NormalTok{ true\_wd }\SpecialCharTok{+}\NormalTok{ true\_ws,}
    \AttributeTok{win\_rho    =}\NormalTok{ win\_wd  }\SpecialCharTok{/}\NormalTok{ (win\_wd }\SpecialCharTok{+}\NormalTok{ win\_ws),}
    \AttributeTok{win\_kappa  =}\NormalTok{ win\_wd  }\SpecialCharTok{+}\NormalTok{ win\_ws}
\NormalTok{  )}

\NormalTok{p\_rho   }\OtherTok{\textless{}{-}} \FunctionTok{ggplot}\NormalTok{(param\_recovery, }\FunctionTok{aes}\NormalTok{(true\_rho, win\_rho)) }\SpecialCharTok{+}
  \FunctionTok{geom\_abline}\NormalTok{(}\AttributeTok{slope =} \DecValTok{1}\NormalTok{, }\AttributeTok{linetype =} \StringTok{"dashed"}\NormalTok{, }\AttributeTok{colour =} \StringTok{"grey40"}\NormalTok{) }\SpecialCharTok{+}
  \FunctionTok{geom\_point}\NormalTok{(}\FunctionTok{aes}\NormalTok{(}\AttributeTok{colour =}\NormalTok{ true\_model, }\AttributeTok{shape =}\NormalTok{ correct), }\AttributeTok{size =} \FloatTok{2.5}\NormalTok{, }\AttributeTok{alpha =}\NormalTok{ .}\DecValTok{7}\NormalTok{) }\SpecialCharTok{+}
  \FunctionTok{coord\_equal}\NormalTok{(}\AttributeTok{xlim =} \FunctionTok{c}\NormalTok{(}\DecValTok{0}\NormalTok{,}\DecValTok{1}\NormalTok{), }\AttributeTok{ylim =} \FunctionTok{c}\NormalTok{(}\DecValTok{0}\NormalTok{,}\DecValTok{1}\NormalTok{))}

\NormalTok{p\_kappa }\OtherTok{\textless{}{-}} \FunctionTok{ggplot}\NormalTok{(param\_recovery, }\FunctionTok{aes}\NormalTok{(true\_kappa, win\_kappa)) }\SpecialCharTok{+}
  \FunctionTok{geom\_abline}\NormalTok{(}\AttributeTok{slope =} \DecValTok{1}\NormalTok{, }\AttributeTok{linetype =} \StringTok{"dashed"}\NormalTok{, }\AttributeTok{colour =} \StringTok{"grey40"}\NormalTok{) }\SpecialCharTok{+}
  \FunctionTok{geom\_point}\NormalTok{(}\FunctionTok{aes}\NormalTok{(}\AttributeTok{colour =}\NormalTok{ true\_model, }\AttributeTok{shape =}\NormalTok{ correct), }\AttributeTok{size =} \FloatTok{2.5}\NormalTok{, }\AttributeTok{alpha =}\NormalTok{ .}\DecValTok{7}\NormalTok{) }\SpecialCharTok{+}
  \FunctionTok{scale\_x\_log10}\NormalTok{() }\SpecialCharTok{+} \FunctionTok{scale\_y\_log10}\NormalTok{()}

\NormalTok{p\_rho }\SpecialCharTok{+}\NormalTok{ p\_kappa}
\end{Highlighting}
\end{Shaded}

\begin{center}
\includegraphics[width=0.8\linewidth,height=\textheight,keepaspectratio]{10-BayesianModels_files/figure-pdf/ch10_model_recovery_parameters-4.pdf}
\end{center}

\begin{Shaded}
\begin{Highlighting}[]
\CommentTok{\# {-}{-}{-}{-}{-}{-}{-}{-}{-}{-}{-}{-}{-}{-}{-}{-}{-}{-}{-}{-}{-}{-}{-}{-}{-}{-}{-}{-}{-}{-}{-}{-}{-}{-}{-}{-}{-}{-}{-}{-}{-}{-}{-}{-}{-}{-}{-}{-}{-}{-}{-}{-}{-}{-}{-}{-}{-}}
\CommentTok{\# VISUALIZATION 3: Parameter Shifting (Misclassifications)}
\CommentTok{\# {-}{-}{-}{-}{-}{-}{-}{-}{-}{-}{-}{-}{-}{-}{-}{-}{-}{-}{-}{-}{-}{-}{-}{-}{-}{-}{-}{-}{-}{-}{-}{-}{-}{-}{-}{-}{-}{-}{-}{-}{-}{-}{-}{-}{-}{-}{-}{-}{-}{-}{-}{-}{-}{-}{-}{-}{-}}

\CommentTok{\# Filter out correct recoveries to focus purely on the mapping of misclassifications}
\NormalTok{misclassifications }\OtherTok{\textless{}{-}}\NormalTok{ param\_recovery }\SpecialCharTok{|\textgreater{}} 
  \FunctionTok{filter}\NormalTok{(correct }\SpecialCharTok{==} \ConstantTok{FALSE}\NormalTok{)}

\CommentTok{\# Create the Vector Shift Plot}
\NormalTok{p\_shift }\OtherTok{\textless{}{-}} \FunctionTok{ggplot}\NormalTok{(misclassifications) }\SpecialCharTok{+}
  \FunctionTok{geom\_abline}\NormalTok{(}\AttributeTok{slope =} \DecValTok{1}\NormalTok{, }\AttributeTok{intercept =} \DecValTok{0}\NormalTok{, }\AttributeTok{linetype =} \StringTok{"dashed"}\NormalTok{, }\AttributeTok{color =} \StringTok{"gray60"}\NormalTok{) }\SpecialCharTok{+}
  
  \CommentTok{\# Draw arrows starting at the True params and pointing to the Fitted params}
  \FunctionTok{geom\_segment}\NormalTok{(}
    \FunctionTok{aes}\NormalTok{(}\AttributeTok{x =}\NormalTok{ true\_wd, }\AttributeTok{y =}\NormalTok{ true\_ws, }\AttributeTok{xend =}\NormalTok{ win\_wd, }\AttributeTok{yend =}\NormalTok{ win\_ws, }\AttributeTok{color =}\NormalTok{ winner),}
    \AttributeTok{arrow =} \FunctionTok{arrow}\NormalTok{(}\AttributeTok{length =} \FunctionTok{unit}\NormalTok{(}\FloatTok{0.15}\NormalTok{, }\StringTok{"cm"}\NormalTok{), }\AttributeTok{type =} \StringTok{"closed"}\NormalTok{),}
    \AttributeTok{alpha =} \FloatTok{0.6}
\NormalTok{  ) }\SpecialCharTok{+}
  
  \CommentTok{\# Add a distinct black dot at the starting point (True parameters)}
  \FunctionTok{geom\_point}\NormalTok{(}
    \FunctionTok{aes}\NormalTok{(}\AttributeTok{x =}\NormalTok{ true\_wd, }\AttributeTok{y =}\NormalTok{ true\_ws), }
    \AttributeTok{color =} \StringTok{"black"}\NormalTok{, }\AttributeTok{size =} \FloatTok{1.2}\NormalTok{, }\AttributeTok{shape =} \DecValTok{16}
\NormalTok{  ) }\SpecialCharTok{+}
  
  \CommentTok{\# Grid layout: Rows = True generating model, Columns = The wrong model that won}
  \FunctionTok{facet\_grid}\NormalTok{(true\_model }\SpecialCharTok{\textasciitilde{}}\NormalTok{ winner, }\AttributeTok{drop =} \ConstantTok{FALSE}\NormalTok{,}
             \AttributeTok{labeller =} \FunctionTok{labeller}\NormalTok{(}
               \AttributeTok{true\_model =} \ControlFlowTok{function}\NormalTok{(x) }\FunctionTok{paste}\NormalTok{(}\StringTok{"True:"}\NormalTok{, x),}
               \AttributeTok{winner =} \ControlFlowTok{function}\NormalTok{(x) }\FunctionTok{paste}\NormalTok{(}\StringTok{"Winner:"}\NormalTok{, x)}
\NormalTok{             )) }\SpecialCharTok{+}
  
  \CommentTok{\# Match the colors you used previously}
  \FunctionTok{scale\_color\_manual}\NormalTok{(}
    \AttributeTok{values =} \FunctionTok{c}\NormalTok{(}\StringTok{"SBA"} \OtherTok{=} \StringTok{"\#E41A1C"}\NormalTok{, }\StringTok{"PBA"} \OtherTok{=} \StringTok{"\#377EB8"}\NormalTok{, }\StringTok{"WBA"} \OtherTok{=} \StringTok{"\#4DAF4A"}\NormalTok{),}
    \AttributeTok{drop =} \ConstantTok{FALSE}
\NormalTok{  ) }\SpecialCharTok{+}
  \FunctionTok{labs}\NormalTok{(}
    \AttributeTok{title =} \StringTok{"Parameter Shifting During Model Misclassification"}\NormalTok{,}
    \AttributeTok{subtitle =} \StringTok{"Arrows trace from the True Parameters (black dots) to the Compensatory Fitted Parameters of the wrong model"}\NormalTok{,}
    \AttributeTok{x =} \FunctionTok{expression}\NormalTok{(}\StringTok{"Weight Delay ("} \SpecialCharTok{*}\NormalTok{ w[d] }\SpecialCharTok{*} \StringTok{")"}\NormalTok{),}
    \AttributeTok{y =} \FunctionTok{expression}\NormalTok{(}\StringTok{"Weight Social ("} \SpecialCharTok{*}\NormalTok{ w[s] }\SpecialCharTok{*} \StringTok{")"}\NormalTok{)}
\NormalTok{  ) }\SpecialCharTok{+}
  \FunctionTok{theme\_bw}\NormalTok{() }\SpecialCharTok{+}
  \FunctionTok{theme}\NormalTok{(}\AttributeTok{legend.position =} \StringTok{"none"}\NormalTok{) }

\CommentTok{\# View Shift Plot}
\NormalTok{p\_shift}
\end{Highlighting}
\end{Shaded}

\begin{center}
\includegraphics[width=0.8\linewidth,height=\textheight,keepaspectratio]{10-BayesianModels_files/figure-pdf/ch10_model_recovery_parameters-5.pdf}
\end{center}

To thoroughly evaluate the reliability of our models, we conduct a
three-stage parameter recovery and identifiability analysis. This
pipeline not only tests whether we can recover the true generating
parameters, but also maps the exact conditions under which models
misclassify and how they compensate when they do. N.B. It's a work in
progress, so I'm sure there are better ways, and I will hopefully
develop better plots in the next course iterations.

\subsubsection{1. Winning Model Parameter
Recovery}\label{winning-model-parameter-recovery}

We first assess the accuracy of our parameter estimates by evaluating
the actual weights inferred by whichever model ``won'' the selection
process. The first pair of plots displays the recovery of the direct
evidence weight (\(w_d\)) and the social weight (\(w_s\)).

\begin{itemize}
\item
  \emph{How to read it}: The x-axis represents the true data-generating
  parameter, while the y-axis represents the estimate generated by the
  winning model (\(\hat{w}\)). The dashed diagonal line indicates
  perfect recovery.
\item
  \emph{Interpretation}: Points falling tightly along the diagonal
  demonstrate robust parameter recovery. We distinguish between correct
  model recoveries (circles) and misclassifications (crosses). By
  plotting the parameters of the winning model, we can observe exactly
  what psychological weights would be reported in practice if we
  selected the best-fitting architecture for a given dataset. The dotted
  trend lines further reveal whether specific true models systematically
  under- or overestimate these weights across the parameter space.
\item
  \emph{Visual Evidence}: The plots show that correct PBA recoveries
  (orange circles) track the diagonal tightly within their bounded \(0\)
  to \(1\) constraint. However, the WBA model (blue dotted trend line)
  systematically underestimates both \(w_d\) and \(w_s\) at higher true
  values (\(> 1.5\)), demonstrating a distinct shrinkage effect.
  Additionally, the constraints of misclassification are highly visible:
  when the PBA incorrectly wins for a true WBA (blue crosses), it
  severely truncates the higher true weight estimates down to values
  \(\leq 1\) to force the \(w_d + w_s = 1\) rule.
\end{itemize}

\subsubsection{2. Mapping Model Identifiability (True vs.~Fitted
Space)}\label{mapping-model-identifiability-true-vs.-fitted-space}

Next, we investigate the geographical boundaries of model selection
accuracy. Rather than just asking if a model failed, we ask where it
failed.

\begin{itemize}
\item
  Plot A (Data Generating Space): This plot is faceted by the True
  Model. It illustrates which specific regions of the true parameter
  space cause the wrong model to win. For instance, we can observe if
  misclassifications cluster near the reference diagonal (where
  \(w_d \approx w_s\)), indicating that the models struggle to
  differentiate the underlying strategies when the weights are balanced.
  Visual Evidence: Looking at the ``True: WBA'' panel,
  misclassifications are highly systematic. The WBA generates data that
  causes the PBA to falsely win (blue crosses) almost exclusively when
  the true WBA parameters fall close to the \(w_d + w_s = 1\) line.
  Similarly, the SBA falsely wins (red crosses) when the true WBA
  parameters cluster tightly near \((1,1)\).
\item
  Plot B (Fitted Parameter Space): This plot is faceted by the Winning
  Model. It reveals the ``imitation space''---the specific parameter
  ranges a model leverages when it inappropriately wins. By comparing
  these two spaces, we can see if models are relying on extreme or
  restricted parameter values to mimic a different generating
  architecture. Visual Evidence: Plot B starkly visualizes the rigid
  architectural constraints of the simpler models. Whenever the SBA
  wins, all parameter estimates collapse precisely to the \((1, 1)\)
  coordinate. Whenever the PBA wins, all estimates are strictly forced
  onto the diagonal \(w_d + w_s = 1\) boundary line, fundamentally
  stripping away absolute scale regardless of where the true parameters
  originated.
\end{itemize}

\subsubsection{3. Structural Mimicry and Compensatory
Shifts}\label{structural-mimicry-and-compensatory-shifts}

To evaluate whether model misclassifications are behaviorally benign, we
isolate the failed recovery trials and trace exactly how the parameter
estimates shift when the wrong model is selected.

\begin{itemize}
\item
  \emph{How to read it}: In this grid, rows represent the true
  generating model and columns represent the incorrectly selected
  ``winning'' model. Each arrow represents a single misclassified
  simulation: the black dot marks the true generating parameters, while
  the arrowhead points to the values estimated by the wrong winning
  model.
\item
  \emph{Interpretation}: We can directly observe structural mimicry. If
  the arrows are very short, or if they reliably project toward a
  specific parameter ratio, it indicates that the misidentification is
  relatively benign; the winning model achieves its higher fit by tuning
  its parameters to functionally approximate the true data-generating
  weights. Conversely, long, erratic arrows reveal regions where the
  wrong model fundamentally distorts the psychological interpretation of
  the data, pulling the estimates far from their true origins to force a
  statistical fit.
\item
  Loss of Scale, Preservation of Ratio (e.g., True WBA \(\rightarrow\)
  Winner PBA): The blue arrows project scattered true WBA parameters
  directly down or up onto the \(w_d + w_s = 1\) line. This strips away
  the total evidence scale but generally preserves the relative
  weighting ratio (the slopes), serving as a benign functional
  approximation.
\item
  Inflation of Scale (e.g., True PBA \(\rightarrow\) Winner WBA): The
  green arrows shoot outward from the PBA's natural boundary. To mimic
  the PBA, the WBA maintains the correct relative ratio (following the
  same trajectory) but unnecessarily inflates the total evidence scale.
\item
  Severe Distortion (e.g., True WBA \(\rightarrow\) Winner SBA): The red
  arrows collapse widely dispersed true WBA parameters entirely into the
  fixed \((1, 1)\) coordinate. This indicates a complete loss of both
  relative weighting and absolute scale information, severely distorting
  the psychological profile of the simulated agent.
\end{itemize}

\section{Summary: What We Learned}\label{summary-what-we-learned}

\begin{enumerate}
\def\labelenumi{\arabic{enumi}.}
\item
  \textbf{The SBA is a baseline} with zero free parameters --- it takes
  all evidence at face value.
\item
  \textbf{The WBA with raw \((w_d, w_s)\) parameterisation has
  identifiability problems.} Even without divergences and with good
  diagnostics, parameter recovery fails because:

  \begin{itemize}
  \tightlist
  \item
    \(w_d = 2, w_s = 2\) produces the same \emph{average} decisions as
    \(w_d = 1, w_s = 1\)
  \item
    Only the \emph{ratio} matters for predicting choices; the
    \emph{scale} mainly affects variance
  \end{itemize}
\item
  \textbf{More data doesn't help} --- this is a structural
  identifiability issue, not a power issue.
\item
  \textbf{Reparameterisation solves the problem:}

  \begin{itemize}
  \tightlist
  \item
    The \textbf{PBA} (\(p\)) fixes the scale and estimates only relative
    allocation
  \item
    The \textbf{WBA} (\(\rho\), \(\kappa\)) separates relative weight
    from total scale
  \end{itemize}
\item
  \textbf{Always run parameter recovery checks} before trusting your
  model, even when standard diagnostics look clean.
\end{enumerate}

\section{Multilevel Extension}\label{multilevel-extension}

In real experiments, we observe many participants, each with their own
weighting tendencies. Multilevel models capture these individual
differences while sharing statistical strength across the population.
This is the same partial-pooling logic developed in Chapter 7.

A single-agent model answers: \emph{how does this person weight direct
vs.~social evidence?} A multilevel model answers the population-level
question: \emph{how do people in general weight evidence, and how much
do they vary?} The hierarchical structure also regularises extreme
individual estimates via partial pooling --- agents with fewer
informative trials are pulled toward the population mean, exactly as we
saw in Chapter 7's matching-pennies analysis.

\section{Simulating a Multi-Agent
Population}\label{simulating-a-multi-agent-population}

We simulate 20 agents, each with their own evidence-weighting tendency,
drawn from a shared population distribution. The population is
parameterised on transformed scales: \(\text{logit}(\rho)\) and
\(\log(\kappa)\) follow a bivariate normal. This ensures
\(\rho \in (0, 1)\) and \(\kappa > 0\) on the natural scale, preserves
the identifiable reparameterisation from the single-agent section, and
allows correlation between relative allocation and total scaling.

\begin{Shaded}
\begin{Highlighting}[]
\NormalTok{n\_agents          }\OtherTok{\textless{}{-}} \DecValTok{20}
\NormalTok{n\_trials\_per\_comb }\OtherTok{\textless{}{-}} \DecValTok{2}   \CommentTok{\# repetitions of each evidence combination per agent}

\CommentTok{\# Population parameters (on transformed scales)}
\NormalTok{pop\_logit\_rho\_mean }\OtherTok{\textless{}{-}} \FloatTok{0.0}    \CommentTok{\# logit(0.5) → balanced on average}
\NormalTok{pop\_logit\_rho\_sd   }\OtherTok{\textless{}{-}} \FloatTok{0.6}    \CommentTok{\# moderate individual variation in allocation}
\NormalTok{pop\_log\_kappa\_mean }\OtherTok{\textless{}{-}} \FunctionTok{log}\NormalTok{(}\DecValTok{2}\NormalTok{) }\CommentTok{\# kappa ≈ 2 on average (SBA{-}like total scaling)}
\NormalTok{pop\_log\_kappa\_sd   }\OtherTok{\textless{}{-}} \FloatTok{0.3}    \CommentTok{\# moderate variation in total evidence use}
\NormalTok{pop\_correlation    }\OtherTok{\textless{}{-}} \SpecialCharTok{{-}}\FloatTok{0.3}   \CommentTok{\# mild tradeoff: high rho → slightly lower kappa}

\CommentTok{\# Construct the bivariate{-}normal covariance matrix}
\NormalTok{Sigma\_pop }\OtherTok{\textless{}{-}} \FunctionTok{matrix}\NormalTok{(}\FunctionTok{c}\NormalTok{(}
\NormalTok{  pop\_logit\_rho\_sd}\SpecialCharTok{\^{}}\DecValTok{2}\NormalTok{,}
\NormalTok{  pop\_correlation }\SpecialCharTok{*}\NormalTok{ pop\_logit\_rho\_sd }\SpecialCharTok{*}\NormalTok{ pop\_log\_kappa\_sd,}
\NormalTok{  pop\_correlation }\SpecialCharTok{*}\NormalTok{ pop\_logit\_rho\_sd }\SpecialCharTok{*}\NormalTok{ pop\_log\_kappa\_sd,}
\NormalTok{  pop\_log\_kappa\_sd}\SpecialCharTok{\^{}}\DecValTok{2}
\NormalTok{), }\AttributeTok{nrow =} \DecValTok{2}\NormalTok{)}

\CommentTok{\# Draw individual parameters from the population}
\NormalTok{raw\_draws }\OtherTok{\textless{}{-}}\NormalTok{ MASS}\SpecialCharTok{::}\FunctionTok{mvrnorm}\NormalTok{(n\_agents, }
                           \AttributeTok{mu    =} \FunctionTok{c}\NormalTok{(pop\_logit\_rho\_mean, pop\_log\_kappa\_mean),}
                           \AttributeTok{Sigma =}\NormalTok{ Sigma\_pop)}

\NormalTok{agent\_params\_ml }\OtherTok{\textless{}{-}} \FunctionTok{tibble}\NormalTok{(}
  \AttributeTok{agent\_id    =} \FunctionTok{seq\_len}\NormalTok{(n\_agents),}
  \AttributeTok{logit\_rho   =}\NormalTok{ raw\_draws[, }\DecValTok{1}\NormalTok{],}
  \AttributeTok{log\_kappa   =}\NormalTok{ raw\_draws[, }\DecValTok{2}\NormalTok{],}
  \AttributeTok{rho         =} \FunctionTok{plogis}\NormalTok{(logit\_rho),}
  \AttributeTok{kappa       =} \FunctionTok{exp}\NormalTok{(log\_kappa),}
  \AttributeTok{weight\_direct =}\NormalTok{ rho }\SpecialCharTok{*}\NormalTok{ kappa,}
  \AttributeTok{weight\_social =}\NormalTok{ (}\DecValTok{1} \SpecialCharTok{{-}}\NormalTok{ rho) }\SpecialCharTok{*}\NormalTok{ kappa}
\NormalTok{)}

\CommentTok{\# Generate trial{-}level data for all agents}
\NormalTok{ml\_sim\_data }\OtherTok{\textless{}{-}}\NormalTok{ agent\_params\_ml }\SpecialCharTok{|\textgreater{}}
  \FunctionTok{pmap\_dfr}\NormalTok{(}\ControlFlowTok{function}\NormalTok{(agent\_id, logit\_rho, log\_kappa, rho, kappa,}
\NormalTok{                    weight\_direct, weight\_social) \{}
    \FunctionTok{generate\_agent\_decisions}\NormalTok{(weight\_direct, weight\_social,}
\NormalTok{                             evidence\_combinations,}
                             \AttributeTok{n\_samples =}\NormalTok{ n\_trials\_per\_comb) }\SpecialCharTok{|\textgreater{}}
      \FunctionTok{mutate}\NormalTok{(}
        \AttributeTok{agent\_id  =}\NormalTok{ agent\_id,}
        \AttributeTok{true\_rho  =}\NormalTok{ rho,}
        \AttributeTok{true\_kappa =}\NormalTok{ kappa,}
        \AttributeTok{true\_wd   =}\NormalTok{ weight\_direct,}
        \AttributeTok{true\_ws   =}\NormalTok{ weight\_social}
\NormalTok{      )}
\NormalTok{  \})}

\NormalTok{stan\_data\_ml }\OtherTok{\textless{}{-}} \FunctionTok{list}\NormalTok{(}
  \AttributeTok{N          =} \FunctionTok{nrow}\NormalTok{(ml\_sim\_data),}
  \AttributeTok{J          =}\NormalTok{ n\_agents,}
  \AttributeTok{agent\_id   =}\NormalTok{ ml\_sim\_data}\SpecialCharTok{$}\NormalTok{agent\_id,}
  \AttributeTok{choice     =}\NormalTok{ ml\_sim\_data}\SpecialCharTok{$}\NormalTok{choice,}
  \AttributeTok{N\_total    =} \FunctionTok{nrow}\NormalTok{(ml\_sim\_data),}
  \AttributeTok{N\_subjects =}\NormalTok{ n\_agents,}
  \AttributeTok{subj\_id    =}\NormalTok{ ml\_sim\_data}\SpecialCharTok{$}\NormalTok{agent\_id,}
  \AttributeTok{y          =}\NormalTok{ ml\_sim\_data}\SpecialCharTok{$}\NormalTok{choice,}
  \AttributeTok{blue1      =}\NormalTok{ ml\_sim\_data}\SpecialCharTok{$}\NormalTok{blue1,}
  \AttributeTok{total1     =}\NormalTok{ ml\_sim\_data}\SpecialCharTok{$}\NormalTok{total1,}
  \AttributeTok{blue2      =}\NormalTok{ ml\_sim\_data}\SpecialCharTok{$}\NormalTok{blue2,}
  \AttributeTok{total2     =}\NormalTok{ ml\_sim\_data}\SpecialCharTok{$}\NormalTok{total2,}
  
  \AttributeTok{prior\_mu\_p\_mu             =} \DecValTok{0}\NormalTok{,}
  \AttributeTok{prior\_mu\_p\_sigma          =} \FloatTok{1.5}\NormalTok{,}
  \AttributeTok{prior\_sigma\_p\_lambda      =} \DecValTok{1}\NormalTok{,}
  \AttributeTok{prior\_mu\_logit\_p\_mu       =} \DecValTok{0}\NormalTok{,}
  \AttributeTok{prior\_mu\_logit\_p\_sigma    =} \FloatTok{1.5}\NormalTok{,}
  \AttributeTok{prior\_sigma\_logit\_p\_lambda =} \DecValTok{2}\NormalTok{,}
  \AttributeTok{prior\_mu\_logit\_rho\_mu     =} \DecValTok{0}\NormalTok{,}
  \AttributeTok{prior\_mu\_logit\_rho\_sigma  =} \FloatTok{1.5}\NormalTok{,}
  \AttributeTok{prior\_mu\_log\_kappa\_mu     =} \FunctionTok{log}\NormalTok{(}\DecValTok{2}\NormalTok{),}
  \AttributeTok{prior\_mu\_log\_kappa\_sigma  =} \FloatTok{0.5}\NormalTok{,}
  \AttributeTok{prior\_sigma\_lambda        =} \DecValTok{2}\NormalTok{,}
  \AttributeTok{prior\_Omega\_shape         =} \DecValTok{2}\NormalTok{,}
  \AttributeTok{run\_diagnostics           =} \DecValTok{1}
\NormalTok{)}

\FunctionTok{cat}\NormalTok{(}\FunctionTok{sprintf}\NormalTok{(}\StringTok{"Agents: \%d | Trials per agent: \%d | Total observations: \%d}\SpecialCharTok{\textbackslash{}n}\StringTok{"}\NormalTok{,}
\NormalTok{            n\_agents, }\FunctionTok{nrow}\NormalTok{(evidence\_combinations) }\SpecialCharTok{*}\NormalTok{ n\_trials\_per\_comb, stan\_data\_ml}\SpecialCharTok{$}\NormalTok{N\_total))}
\end{Highlighting}
\end{Shaded}

\begin{verbatim}
Agents: 20 | Trials per agent: 72 | Total observations: 1440
\end{verbatim}

\begin{Shaded}
\begin{Highlighting}[]
\CommentTok{\# Calculate the true population means on the natural scale for our reference lines}
\NormalTok{pop\_rho\_true   }\OtherTok{\textless{}{-}} \FunctionTok{plogis}\NormalTok{(pop\_logit\_rho\_mean)}
\NormalTok{pop\_kappa\_true }\OtherTok{\textless{}{-}} \FunctionTok{exp}\NormalTok{(pop\_log\_kappa\_mean)}
\NormalTok{pop\_wd\_mean    }\OtherTok{\textless{}{-}}\NormalTok{ pop\_rho\_true }\SpecialCharTok{*}\NormalTok{ pop\_kappa\_true}
\NormalTok{pop\_ws\_mean    }\OtherTok{\textless{}{-}}\NormalTok{ (}\DecValTok{1} \SpecialCharTok{{-}}\NormalTok{ pop\_rho\_true) }\SpecialCharTok{*}\NormalTok{ pop\_kappa\_true}

\FunctionTok{ggplot}\NormalTok{(agent\_params\_ml, }\FunctionTok{aes}\NormalTok{(}\AttributeTok{x =}\NormalTok{ weight\_direct, }\AttributeTok{y =}\NormalTok{ weight\_social)) }\SpecialCharTok{+}
  \CommentTok{\# Add dashed lines for the population means}
  \FunctionTok{geom\_vline}\NormalTok{(}\AttributeTok{xintercept =}\NormalTok{ pop\_wd\_mean, }\AttributeTok{linetype =} \StringTok{"dashed"}\NormalTok{, }\AttributeTok{color =} \StringTok{"blue"}\NormalTok{, }\AttributeTok{alpha =} \FloatTok{0.7}\NormalTok{) }\SpecialCharTok{+}
  \FunctionTok{geom\_hline}\NormalTok{(}\AttributeTok{yintercept =}\NormalTok{ pop\_ws\_mean, }\AttributeTok{linetype =} \StringTok{"dashed"}\NormalTok{, }\AttributeTok{color =} \StringTok{"blue"}\NormalTok{, }\AttributeTok{alpha =} \FloatTok{0.7}\NormalTok{) }\SpecialCharTok{+}
  \CommentTok{\# Add the individual agents}
  \FunctionTok{geom\_point}\NormalTok{(}\AttributeTok{size =} \FloatTok{3.5}\NormalTok{, }\AttributeTok{color =} \StringTok{"black"}\NormalTok{, }\AttributeTok{alpha =} \FloatTok{0.6}\NormalTok{) }\SpecialCharTok{+}
  \FunctionTok{geom\_text}\NormalTok{(}\FunctionTok{aes}\NormalTok{(}\AttributeTok{label =}\NormalTok{ agent\_id), }\AttributeTok{nudge\_x =} \FloatTok{0.04}\NormalTok{, }\AttributeTok{nudge\_y =} \FloatTok{0.02}\NormalTok{, }\AttributeTok{size =} \DecValTok{4}\NormalTok{, }\AttributeTok{color =} \StringTok{"black"}\NormalTok{) }\SpecialCharTok{+}
  \FunctionTok{labs}\NormalTok{(}
    \AttributeTok{title =} \StringTok{"Distribution of Simulated Agent Parameters"}\NormalTok{,}
    \AttributeTok{subtitle =} \StringTok{"Each point represents one agent; dashed lines show population means"}\NormalTok{,}
    \AttributeTok{x =} \FunctionTok{expression}\NormalTok{(}\StringTok{"Weight for Direct Evidence ("} \SpecialCharTok{*}\NormalTok{ w[d] }\SpecialCharTok{*} \StringTok{")"}\NormalTok{),}
    \AttributeTok{y =} \FunctionTok{expression}\NormalTok{(}\StringTok{"Weight for Social Evidence ("} \SpecialCharTok{*}\NormalTok{ w[s] }\SpecialCharTok{*} \StringTok{")"}\NormalTok{)}
\NormalTok{  ) }\SpecialCharTok{+}
  \FunctionTok{theme}\NormalTok{(}
    \AttributeTok{plot.title =} \FunctionTok{element\_text}\NormalTok{(}\AttributeTok{face =} \StringTok{"bold"}\NormalTok{, }\AttributeTok{size =} \DecValTok{14}\NormalTok{),}
    \AttributeTok{axis.title =} \FunctionTok{element\_text}\NormalTok{(}\AttributeTok{size =} \DecValTok{12}\NormalTok{)}
\NormalTok{  )}
\end{Highlighting}
\end{Shaded}

\begin{center}
\includegraphics[width=0.8\linewidth,height=\textheight,keepaspectratio]{10-BayesianModels_files/figure-pdf/ch10_multilevel_sim_setup-1.pdf}
\end{center}

Each agent sees \(36 \times 2 = 72\) trials, giving us
\(20 \times 72 = 1440\) total observations --- a realistic experimental
dataset size.

\subsection{Multilevel Simple Bayesian Agent
(Stan)}\label{multilevel-simple-bayesian-agent-stan}

The ML-SBA has zero free parameters per agent. It applies strict
Bayesian integration identically across all agents. This is the null
model: individual differences in evidence weighting do not exist.

\begin{Shaded}
\begin{Highlighting}[]
\FunctionTok{\textbackslash{}usetikzlibrary}\NormalTok{\{bayesnet\}}
\KeywordTok{\textbackslash{}begin}\NormalTok{\{}\ExtensionTok{tikzpicture}\NormalTok{\}}
  \CommentTok{\% Constants}
  \FunctionTok{\textbackslash{}node}\NormalTok{[const] (alpha0) \{}\SpecialStringTok{$}\SpecialCharTok{\textbackslash{}alpha}\SpecialStringTok{\_0$}\NormalTok{\};}
  \FunctionTok{\textbackslash{}node}\NormalTok{[const, right=1cm of alpha0] (beta0) \{}\SpecialStringTok{$}\SpecialCharTok{\textbackslash{}beta}\SpecialStringTok{\_0$}\NormalTok{\};}
  
  \CommentTok{\% Trial level}
  \FunctionTok{\textbackslash{}node}\NormalTok{[obs, below left=1.5cm and 0.5cm of alpha0] (blue1) \{}\SpecialStringTok{$k\_\{1,i\}$}\NormalTok{\};}
  \FunctionTok{\textbackslash{}node}\NormalTok{[obs, right=0.3cm of blue1] (total1) \{}\SpecialStringTok{$n\_\{1,i\}$}\NormalTok{\};}
  
  \FunctionTok{\textbackslash{}node}\NormalTok{[obs, below right=1.5cm and 0.5cm of beta0] (blue2) \{}\SpecialStringTok{$k\_\{2,i\}$}\NormalTok{\};}
  \FunctionTok{\textbackslash{}node}\NormalTok{[obs, left=0.3cm of blue2] (total2) \{}\SpecialStringTok{$n\_\{2,i\}$}\NormalTok{\};}

  \FunctionTok{\textbackslash{}node}\NormalTok{[det, below=2.5cm of alpha0] (alphapost) \{}\SpecialStringTok{$}\SpecialCharTok{\textbackslash{}alpha}\SpecialStringTok{\_\{post,i\}$}\NormalTok{\};}
  \FunctionTok{\textbackslash{}node}\NormalTok{[det, below=2.5cm of beta0] (betapost) \{}\SpecialStringTok{$}\SpecialCharTok{\textbackslash{}beta}\SpecialStringTok{\_\{post,i\}$}\NormalTok{\};}

  \FunctionTok{\textbackslash{}node}\NormalTok{[obs, below=1cm of alphapost, xshift=0.75cm] (choice) \{}\SpecialStringTok{$y\_i$}\NormalTok{\};}

  \CommentTok{\% Edges}
  \FunctionTok{\textbackslash{}edge}\NormalTok{ \{alpha0\} \{alphapost\};}
  \FunctionTok{\textbackslash{}edge}\NormalTok{ \{beta0\} \{betapost\};}
  \FunctionTok{\textbackslash{}edge}\NormalTok{ \{blue1, total1, blue2, total2\} \{alphapost, betapost\};}
  \FunctionTok{\textbackslash{}edge}\NormalTok{ \{alphapost, betapost\} \{choice\};}

  \CommentTok{\% Plates}
  \FunctionTok{\textbackslash{}plate}\NormalTok{ \{trial\} \{(blue1)(total1)(blue2)(total2)(alphapost)(betapost)(choice)\} \{}\SpecialStringTok{$N$}\NormalTok{ trials\};}
  \FunctionTok{\textbackslash{}plate}\NormalTok{ \{agent\} \{(trial)\} \{}\SpecialStringTok{$J$}\NormalTok{ agents\};}
\KeywordTok{\textbackslash{}end}\NormalTok{\{}\ExtensionTok{tikzpicture}\NormalTok{\}}
\end{Highlighting}
\end{Shaded}

\begin{figure}[H]

{\centering \includegraphics[width=0.55\linewidth,height=\textheight,keepaspectratio]{10-BayesianModels_files/figure-pdf/ml-sba-dag-1.pdf}

}

\caption{DAG for the Multilevel Simple Bayesian Agent}

\end{figure}%

The generative relationships for the ML-SBA are identical across agents:
\[
\begin{aligned}
\alpha_{post,ij} &= \alpha_0 + k_{1,ij} + k_{2,ij} \\
\beta_{post,ij}  &= \beta_0 + (n_{1,ij} - k_{1,ij}) + (n_{2,ij} - k_{2,ij}) \\
y_{ij} &\sim \text{Beta-Binomial}(1, \alpha_{post,ij}, \beta_{post,ij})
\end{aligned}
\]

\begin{Shaded}
\begin{Highlighting}[]
\NormalTok{MultilevelSBA\_stan }\OtherTok{\textless{}{-}} \StringTok{"}
\StringTok{// Multilevel Simple Bayesian Agent}
\StringTok{// Zero free parameters per agent. Strict integration applied uniformly.}
\StringTok{data \{}
\StringTok{  int\textless{}lower=1\textgreater{} N\_total;}
\StringTok{  int\textless{}lower=1\textgreater{} N\_subjects;}
\StringTok{  array[N\_total] int\textless{}lower=1, upper=N\_subjects\textgreater{} subj\_id;}
\StringTok{  array[N\_total] int\textless{}lower=0, upper=1\textgreater{}  y;}
\StringTok{  array[N\_total] int\textless{}lower=0\textgreater{}            blue1;}
\StringTok{  array[N\_total] int\textless{}lower=0\textgreater{}            total1;}
\StringTok{  array[N\_total] int\textless{}lower=0\textgreater{}            blue2;}
\StringTok{  array[N\_total] int\textless{}lower=0\textgreater{}            total2;}
\StringTok{  int\textless{}lower=0, upper=1\textgreater{} run\_diagnostics;}
\StringTok{\}}

\StringTok{transformed data \{}
\StringTok{  real alpha0 = 0.5;}
\StringTok{  real beta0  = 0.5;}
\StringTok{\}}

\StringTok{model \{}
\StringTok{  vector[N\_total] alpha\_post = alpha0 + to\_vector(blue1) + to\_vector(blue2);}
\StringTok{  vector[N\_total] beta\_post  = beta0  + (to\_vector(total1) {-} to\_vector(blue1)) +}
\StringTok{                                        (to\_vector(total2) {-} to\_vector(blue2));}
\StringTok{  target += beta\_binomial\_lpmf(y | 1, alpha\_post, beta\_post);}
\StringTok{\}}

\StringTok{generated quantities \{}
\StringTok{  vector[N\_total] log\_lik;}
\StringTok{  array[N\_total] int y\_rep;}

\StringTok{  if (run\_diagnostics) \{}
\StringTok{    for (i in 1:N\_total) \{}
\StringTok{      real alpha\_post = alpha0 + blue1[i] + blue2[i];}
\StringTok{      real beta\_post  = beta0 + (total1[i] {-} blue1[i]) + (total2[i] {-} blue2[i]);}
\StringTok{      log\_lik[i] = beta\_binomial\_lpmf(y[i] | 1, alpha\_post, beta\_post);}
\StringTok{      y\_rep[i]   = beta\_binomial\_rng(1, alpha\_post, beta\_post);}
\StringTok{    \}}
\StringTok{  \}}
\StringTok{\}}
\StringTok{"}

\FunctionTok{write\_stan\_file}\NormalTok{(MultilevelSBA\_stan,}
                \AttributeTok{dir =} \StringTok{"stan/"}\NormalTok{,}
                \AttributeTok{basename =} \StringTok{"ch10\_multilevel\_sba.stan"}\NormalTok{)}
\end{Highlighting}
\end{Shaded}

\begin{verbatim}
[1] "/Users/au209589/Dropbox/Teaching/AdvancedCognitiveModeling23_book/stan/ch10_multilevel_sba.stan"
\end{verbatim}

\begin{Shaded}
\begin{Highlighting}[]
\NormalTok{mod\_ml\_sba }\OtherTok{\textless{}{-}} \FunctionTok{cmdstan\_model}\NormalTok{(}\StringTok{"stan/ch10\_multilevel\_sba.stan"}\NormalTok{, }\AttributeTok{dir =} \StringTok{"simmodels"}\NormalTok{)}
\end{Highlighting}
\end{Shaded}

\subsection{Multilevel Proportional Bayesian Agent
(Stan)}\label{multilevel-proportional-bayesian-agent-stan}

The ML-PBA estimates a single parameter per agent --- the evidence
allocation \(p_j\) --- with a population distribution on the logit
scale. This is the intermediate-complexity model: it captures individual
differences in \emph{relative} weighting but not in total evidence
scaling.

\begin{Shaded}
\begin{Highlighting}[]
\NormalTok{MultilevelPBA\_stan }\OtherTok{\textless{}{-}} \StringTok{"}
\StringTok{// Multilevel Proportional Bayesian Agent}
\StringTok{// Population distribution on logit(p).}
\StringTok{// Non{-}centred parameterisation for individual p\_j.}
\StringTok{data \{}
\StringTok{  int\textless{}lower=1\textgreater{} N\_total;}
\StringTok{  int\textless{}lower=1\textgreater{} N\_subjects;}
\StringTok{  array[N\_total] int\textless{}lower=1, upper=N\_subjects\textgreater{} subj\_id;}
\StringTok{  array[N\_total] int\textless{}lower=0, upper=1\textgreater{}  y;}
\StringTok{  array[N\_total] int\textless{}lower=0\textgreater{}            blue1;}
\StringTok{  array[N\_total] int\textless{}lower=0\textgreater{}            total1;}
\StringTok{  array[N\_total] int\textless{}lower=0\textgreater{}            blue2;}
\StringTok{  array[N\_total] int\textless{}lower=0\textgreater{}            total2;}
\StringTok{  }
\StringTok{  real prior\_mu\_logit\_p\_mu;}
\StringTok{  real\textless{}lower=0\textgreater{} prior\_mu\_logit\_p\_sigma;}
\StringTok{  real\textless{}lower=0\textgreater{} prior\_sigma\_logit\_p\_lambda;}
\StringTok{  int\textless{}lower=0, upper=1\textgreater{} run\_diagnostics;}
\StringTok{\}}

\StringTok{parameters \{}
\StringTok{  real mu\_logit\_p;             // population mean on logit scale}
\StringTok{  real\textless{}lower=0\textgreater{} sigma\_logit\_p; // population SD on logit scale}
\StringTok{  vector[N\_subjects] z\_p;      // NCP deviates}
\StringTok{\}}

\StringTok{transformed parameters \{}
\StringTok{  vector\textless{}lower=0, upper=1\textgreater{}[N\_subjects] p;}
\StringTok{  for (j in 1:N\_subjects)}
\StringTok{    p[j] = inv\_logit(mu\_logit\_p + sigma\_logit\_p * z\_p[j]);}
\StringTok{\}}

\StringTok{model \{}
\StringTok{  target += normal\_lpdf(mu\_logit\_p    | prior\_mu\_logit\_p\_mu, prior\_mu\_logit\_p\_sigma);}
\StringTok{  target += exponential\_lpdf(sigma\_logit\_p | prior\_sigma\_logit\_p\_lambda);}
\StringTok{  target += std\_normal\_lpdf(z\_p);}
\StringTok{  }
\StringTok{  for (i in 1:N\_total) \{}
\StringTok{    int j = subj\_id[i];}
\StringTok{    real alpha\_post = 0.5 + p[j] * blue1[i] + (1.0 {-} p[j]) * blue2[i];}
\StringTok{    real beta\_post  = 0.5 + p[j] * (total1[i] {-} blue1[i])}
\StringTok{                          + (1.0 {-} p[j]) * (total2[i] {-} blue2[i]);}
\StringTok{    target += beta\_binomial\_lpmf(y[i] | 1, alpha\_post, beta\_post);}
\StringTok{  \}}
\StringTok{\}}

\StringTok{generated quantities \{}
\StringTok{  real pop\_p = inv\_logit(mu\_logit\_p);}
\StringTok{  vector[N\_total] log\_lik;}
\StringTok{  array[N\_total] int y\_rep;}

\StringTok{  if (run\_diagnostics) \{}
\StringTok{    for (i in 1:N\_total) \{}
\StringTok{      int j = subj\_id[i];}
\StringTok{      real alpha\_post = 0.5 + p[j] * blue1[i] + (1.0 {-} p[j]) * blue2[i];}
\StringTok{      real beta\_post  = 0.5 + p[j] * (total1[i] {-} blue1[i])}
\StringTok{                            + (1.0 {-} p[j]) * (total2[i] {-} blue2[i]);}
\StringTok{      log\_lik[i] = beta\_binomial\_lpmf(y[i] | 1, alpha\_post, beta\_post);}
\StringTok{      y\_rep[i]   = beta\_binomial\_rng(1, alpha\_post, beta\_post);}
\StringTok{    \}}
\StringTok{  \}}
\StringTok{\}}
\StringTok{"}

\FunctionTok{write\_stan\_file}\NormalTok{(MultilevelPBA\_stan,}
                \AttributeTok{dir =} \StringTok{"stan/"}\NormalTok{,}
                \AttributeTok{basename =} \StringTok{"ch10\_multilevel\_pba.stan"}\NormalTok{)}
\end{Highlighting}
\end{Shaded}

\begin{verbatim}
[1] "/Users/au209589/Dropbox/Teaching/AdvancedCognitiveModeling23_book/stan/ch10_multilevel_pba.stan"
\end{verbatim}

\begin{Shaded}
\begin{Highlighting}[]
\NormalTok{mod\_ml\_pba }\OtherTok{\textless{}{-}} \FunctionTok{cmdstan\_model}\NormalTok{(}\StringTok{"stan/ch10\_multilevel\_pba.stan"}\NormalTok{, }\AttributeTok{dir =} \StringTok{"simmodels"}\NormalTok{)}
\end{Highlighting}
\end{Shaded}

\subsection{Multilevel Weighted Bayesian Agent
(Stan)}\label{multilevel-weighted-bayesian-agent-stan}

The multilevel WBA implements the reparameterised \((\rho_j, \kappa_j)\)
parameterisation for each agent, with a bivariate population
distribution on the \((\text{logit}(\rho), \log(\kappa))\) scale. This
preserves the identifiability gains from the single-agent analysis while
allowing partial pooling.

The key design choices mirror Chapter 7:

\begin{itemize}
\tightlist
\item
  \textbf{Non-centred parameterisation (NCP):} Individual deviations are
  expressed as \(z\)-scores to avoid the funnel geometry.
\item
  \textbf{Transformed-scale parameterisation:} We work on
  \((\text{logit}(\rho), \log(\kappa))\) to ensure correct support
  without boundary constraints.
\item
  \textbf{LKJ prior on the correlation:} With \(\eta = 2\), the prior is
  unimodal at zero and weakly regularises against extreme correlations.
  A negative correlation would mean agents who allocate more to direct
  evidence tend to use less total evidence.
\end{itemize}

\begin{Shaded}
\begin{Highlighting}[]
\FunctionTok{\textbackslash{}usetikzlibrary}\NormalTok{\{bayesnet\}}
\KeywordTok{\textbackslash{}begin}\NormalTok{\{}\ExtensionTok{tikzpicture}\NormalTok{\}}
  \CommentTok{\% Population parameters}
  \FunctionTok{\textbackslash{}node}\NormalTok{[latent] (mu) \{}\SpecialStringTok{$}\SpecialCharTok{\textbackslash{}mu}\SpecialStringTok{$}\NormalTok{\};}
  \FunctionTok{\textbackslash{}node}\NormalTok{[latent, right=1.5cm of mu] (sigma) \{}\SpecialStringTok{$}\SpecialCharTok{\textbackslash{}sigma}\SpecialStringTok{$}\NormalTok{\};}
  \FunctionTok{\textbackslash{}node}\NormalTok{[latent, right=1.5cm of sigma] (Omega) \{}\SpecialStringTok{$}\SpecialCharTok{\textbackslash{}Omega}\SpecialStringTok{$}\NormalTok{\};}

  \CommentTok{\% Individual parameters}
  \FunctionTok{\textbackslash{}node}\NormalTok{[latent, below=1.5cm of sigma] (indiv) \{}\SpecialStringTok{$}\SpecialCharTok{\textbackslash{}rho}\SpecialStringTok{\_j, }\SpecialCharTok{\textbackslash{}kappa}\SpecialStringTok{\_j$}\NormalTok{\};}
  \FunctionTok{\textbackslash{}node}\NormalTok{[det, below=1cm of indiv, xshift={-}1.5cm] (wd) \{}\SpecialStringTok{$w\_\{d,j\}$}\NormalTok{\};}
  \FunctionTok{\textbackslash{}node}\NormalTok{[det, below=1cm of indiv, xshift=1.5cm] (ws) \{}\SpecialStringTok{$w\_\{s,j\}$}\NormalTok{\};}

  \CommentTok{\% Trial level inputs}
  \FunctionTok{\textbackslash{}node}\NormalTok{[obs, below=1.5cm of wd, xshift={-}1.5cm] (k1) \{}\SpecialStringTok{$k\_\{1,i\}$}\NormalTok{\};}
  \FunctionTok{\textbackslash{}node}\NormalTok{[obs, right=0.3cm of k1] (n1) \{}\SpecialStringTok{$n\_\{1,i\}$}\NormalTok{\};}
  \FunctionTok{\textbackslash{}node}\NormalTok{[obs, below=1.5cm of ws, xshift=1.5cm] (k2) \{}\SpecialStringTok{$k\_\{2,i\}$}\NormalTok{\};}
  \FunctionTok{\textbackslash{}node}\NormalTok{[obs, left=0.3cm of k2] (n2) \{}\SpecialStringTok{$n\_\{2,i\}$}\NormalTok{\};}

  \CommentTok{\% Posterior parameters (shifted further down to prevent overlap)}
  \FunctionTok{\textbackslash{}node}\NormalTok{[det, below=3cm of wd] (ap) \{}\SpecialStringTok{$}\SpecialCharTok{\textbackslash{}alpha}\SpecialStringTok{\_\{post,i\}$}\NormalTok{\};}
  \FunctionTok{\textbackslash{}node}\NormalTok{[det, below=3cm of ws] (bp) \{}\SpecialStringTok{$}\SpecialCharTok{\textbackslash{}beta}\SpecialStringTok{\_\{post,i\}$}\NormalTok{\};}

  \CommentTok{\% Output}
  \FunctionTok{\textbackslash{}node}\NormalTok{[obs, below=1.5cm of ap, xshift=1.5cm] (y) \{}\SpecialStringTok{$y\_i$}\NormalTok{\};}

  \CommentTok{\% Edges}
  \FunctionTok{\textbackslash{}edge}\NormalTok{ \{mu, sigma, Omega\} \{indiv\};}
  \FunctionTok{\textbackslash{}edge}\NormalTok{ \{indiv\} \{wd, ws\};}
  \FunctionTok{\textbackslash{}edge}\NormalTok{ \{wd, k1, n1, ws, k2, n2\} \{ap, bp\};}
  \FunctionTok{\textbackslash{}edge}\NormalTok{ \{ap, bp\} \{y\};}

  \CommentTok{\% Plates (Nested)}
  \FunctionTok{\textbackslash{}plate}\NormalTok{ \{trial\} \{(k1)(n1)(k2)(n2)(ap)(bp)(y)\} \{}\SpecialStringTok{$N$}\NormalTok{ trials\};}
  \FunctionTok{\textbackslash{}plate}\NormalTok{ \{agent\} \{(indiv)(wd)(ws)(trial)\} \{}\SpecialStringTok{$J$}\NormalTok{ agents\};}
\KeywordTok{\textbackslash{}end}\NormalTok{\{}\ExtensionTok{tikzpicture}\NormalTok{\}}
\end{Highlighting}
\end{Shaded}

\begin{figure}[H]

{\centering \includegraphics[width=0.65\linewidth,height=\textheight,keepaspectratio]{10-BayesianModels_files/figure-pdf/ml-wba-dag-1.pdf}

}

\caption{DAG for the Multilevel Weighted Bayesian Agent}

\end{figure}%

The generative relationships for the ML-WBA are: \[
\begin{aligned}
\begin{pmatrix} \text{logit}(\rho_j) \\ \log(\kappa_j) \end{pmatrix} &\sim \text{MVN}(\boldsymbol{\mu}, \boldsymbol{\Sigma}) \\
w_{d,j} &= \rho_j \cdot \kappa_j \\
w_{s,j} &= (1 - \rho_j) \cdot \kappa_j \\
\alpha_{post,ij} &= \alpha_0 + w_{d,j} k_{1,ij} + w_{s,j} k_{2,ij} \\
\beta_{post,ij}  &= \beta_0 + w_{d,j}(n_{1,ij} - k_{1,ij}) + w_{s,j}(n_{2,ij} - k_{2,ij}) \\
y_{ij} &\sim \text{Beta-Binomial}(1, \alpha_{post,ij}, \beta_{post,ij})
\end{aligned}
\]

\begin{Shaded}
\begin{Highlighting}[]
\NormalTok{MultilevelWBA\_stan }\OtherTok{\textless{}{-}} \StringTok{"}
\StringTok{// Multilevel Weighted Bayesian Agent (reparameterised)}
\StringTok{// Population distribution on (logit(rho), log(kappa)).}
\StringTok{// Non{-}centred parameterisation with correlated random effects.}
\StringTok{data \{}
\StringTok{  int\textless{}lower=1\textgreater{} N\_total;                                // total observations}
\StringTok{  int\textless{}lower=1\textgreater{} N\_subjects;                             // number of agents}
\StringTok{  array[N\_total] int\textless{}lower=1, upper=N\_subjects\textgreater{} subj\_id;}
\StringTok{  array[N\_total] int\textless{}lower=0, upper=1\textgreater{}  y;}
\StringTok{  array[N\_total] int\textless{}lower=0\textgreater{}            blue1;}
\StringTok{  array[N\_total] int\textless{}lower=0\textgreater{}            total1;}
\StringTok{  array[N\_total] int\textless{}lower=0\textgreater{}            blue2;}
\StringTok{  array[N\_total] int\textless{}lower=0\textgreater{}            total2;}
\StringTok{  }
\StringTok{  real prior\_mu\_logit\_rho\_mu;}
\StringTok{  real\textless{}lower=0\textgreater{} prior\_mu\_logit\_rho\_sigma;}
\StringTok{  real prior\_mu\_log\_kappa\_mu;}
\StringTok{  real\textless{}lower=0\textgreater{} prior\_mu\_log\_kappa\_sigma;}
\StringTok{  real\textless{}lower=0\textgreater{} prior\_sigma\_lambda;}
\StringTok{  real\textless{}lower=0\textgreater{} prior\_Omega\_shape;}
\StringTok{  int\textless{}lower=0, upper=1\textgreater{} run\_diagnostics;}
\StringTok{\}}

\StringTok{parameters \{}
\StringTok{  vector[2] mu;                              // population means: (logit\_rho, log\_kappa)}
\StringTok{  vector\textless{}lower=0\textgreater{}[2] sigma;                  // population SDs}
\StringTok{  cholesky\_factor\_corr[2] L\_Omega;           // Cholesky factor of correlation matrix}
\StringTok{  matrix[2, N\_subjects] z;                   // standard normal deviates (NCP)}
\StringTok{\}}

\StringTok{transformed parameters \{}
\StringTok{  // Reconstruct individual parameters on natural scale}
\StringTok{  matrix[2, N\_subjects] theta\_raw = diag\_pre\_multiply(sigma, L\_Omega) * z;}
\StringTok{  }
\StringTok{  vector\textless{}lower=0, upper=1\textgreater{}[N\_subjects] rho;}
\StringTok{  vector\textless{}lower=0\textgreater{}[N\_subjects] kappa;}
\StringTok{  vector\textless{}lower=0\textgreater{}[N\_subjects] weight\_direct;}
\StringTok{  vector\textless{}lower=0\textgreater{}[N\_subjects] weight\_social;}

\StringTok{  for (j in 1:N\_subjects) \{}
\StringTok{    rho[j]           = inv\_logit(mu[1] + theta\_raw[1, j]);}
\StringTok{    kappa[j]         = exp(mu[2] + theta\_raw[2, j]);}
\StringTok{    weight\_direct[j] = rho[j] * kappa[j];}
\StringTok{    weight\_social[j] = (1.0 {-} rho[j]) * kappa[j];}
\StringTok{  \}}
\StringTok{\}}

\StringTok{model \{}
\StringTok{  // Population priors}
\StringTok{  target += normal\_lpdf(mu[1] | prior\_mu\_logit\_rho\_mu, prior\_mu\_logit\_rho\_sigma);}
\StringTok{  target += normal\_lpdf(mu[2] | prior\_mu\_log\_kappa\_mu, prior\_mu\_log\_kappa\_sigma);}
\StringTok{  target += exponential\_lpdf(sigma | prior\_sigma\_lambda);}
\StringTok{  target += lkj\_corr\_cholesky\_lpdf(L\_Omega | prior\_Omega\_shape);}
\StringTok{  target += std\_normal\_lpdf(to\_vector(z));}
\StringTok{  }
\StringTok{  // Likelihood (vectorised per trial)}
\StringTok{  for (i in 1:N\_total) \{}
\StringTok{    int j = subj\_id[i];}
\StringTok{    real alpha\_post = 0.5}
\StringTok{                    + weight\_direct[j] * blue1[i]}
\StringTok{                    + weight\_social[j] * blue2[i];}
\StringTok{    real beta\_post  = 0.5}
\StringTok{                    + weight\_direct[j] * (total1[i] {-} blue1[i])}
\StringTok{                    + weight\_social[j] * (total2[i] {-} blue2[i]);}
\StringTok{    target += beta\_binomial\_lpmf(y[i] | 1, alpha\_post, beta\_post);}
\StringTok{  \}}
\StringTok{\}}

\StringTok{generated quantities \{}
\StringTok{  // Population summaries on natural scale}
\StringTok{  real pop\_rho   = inv\_logit(mu[1]);}
\StringTok{  real pop\_kappa = exp(mu[2]);}
\StringTok{  matrix[2, 2] Omega = multiply\_lower\_tri\_self\_transpose(L\_Omega);}
\StringTok{  }
\StringTok{  vector[N\_total] log\_lik;}
\StringTok{  array[N\_total] int y\_rep;}

\StringTok{  if (run\_diagnostics) \{}
\StringTok{    for (i in 1:N\_total) \{}
\StringTok{      int j = subj\_id[i];}
\StringTok{      real a = 0.5 + weight\_direct[j] * blue1[i] + weight\_social[j] * blue2[i];}
\StringTok{      real b = 0.5 + weight\_direct[j] * (total1[i] {-} blue1[i]) }
\StringTok{                   + weight\_social[j] * (total2[i] {-} blue2[i]);}
\StringTok{      log\_lik[i] = beta\_binomial\_lpmf(y[i] | 1, a, b);}
\StringTok{      y\_rep[i]   = beta\_binomial\_rng(1, a, b);}
\StringTok{    \}}
\StringTok{  \}}
\StringTok{\}}
\StringTok{"}

\FunctionTok{write\_stan\_file}\NormalTok{(MultilevelWBA\_stan,}
                \AttributeTok{dir =} \StringTok{"stan/"}\NormalTok{,}
                \AttributeTok{basename =} \StringTok{"ch10\_multilevel\_wba.stan"}\NormalTok{)}
\end{Highlighting}
\end{Shaded}

\begin{verbatim}
[1] "/Users/au209589/Dropbox/Teaching/AdvancedCognitiveModeling23_book/stan/ch10_multilevel_wba.stan"
\end{verbatim}

\begin{Shaded}
\begin{Highlighting}[]
\NormalTok{mod\_ml\_wba }\OtherTok{\textless{}{-}} \FunctionTok{cmdstan\_model}\NormalTok{(}\StringTok{"stan/ch10\_multilevel\_wba.stan"}\NormalTok{, }\AttributeTok{dir =} \StringTok{"simmodels"}\NormalTok{)}
\end{Highlighting}
\end{Shaded}

\section{Fitting the Multilevel
Models}\label{fitting-the-multilevel-models}

\begin{Shaded}
\begin{Highlighting}[]
\NormalTok{fit\_ml\_sba }\OtherTok{\textless{}{-}} \ControlFlowTok{if}\NormalTok{ (regenerate\_simulations }\SpecialCharTok{||} \SpecialCharTok{!}\FunctionTok{file.exists}\NormalTok{(}\StringTok{"simmodels/ch10\_fit\_ml\_sba.rds"}\NormalTok{)) \{}
\NormalTok{  fit }\OtherTok{\textless{}{-}}\NormalTok{ mod\_ml\_sba}\SpecialCharTok{$}\FunctionTok{sample}\NormalTok{(}
    \AttributeTok{data            =}\NormalTok{ stan\_data\_ml,}
    \AttributeTok{chains          =} \DecValTok{4}\NormalTok{,}
    \AttributeTok{parallel\_chains =} \DecValTok{4}\NormalTok{,}
    \AttributeTok{iter\_warmup     =} \DecValTok{1000}\NormalTok{,}
    \AttributeTok{iter\_sampling   =} \DecValTok{1000}\NormalTok{,}
    \AttributeTok{adapt\_delta     =} \FloatTok{0.9}\NormalTok{,}
    \AttributeTok{refresh         =} \DecValTok{0}
\NormalTok{  )}
\NormalTok{  fit}\SpecialCharTok{$}\FunctionTok{save\_object}\NormalTok{(}\StringTok{"simmodels/ch10\_fit\_ml\_sba.rds"}\NormalTok{)}
\NormalTok{  fit}
\NormalTok{\} }\ControlFlowTok{else}\NormalTok{ \{}
  \FunctionTok{readRDS}\NormalTok{(}\StringTok{"simmodels/ch10\_fit\_ml\_sba.rds"}\NormalTok{)}
\NormalTok{\}}
\NormalTok{fit\_ml\_pba }\OtherTok{\textless{}{-}} \ControlFlowTok{if}\NormalTok{ (regenerate\_simulations }\SpecialCharTok{||} \SpecialCharTok{!}\FunctionTok{file.exists}\NormalTok{(}\StringTok{"simmodels/ch10\_fit\_ml\_pba.rds"}\NormalTok{)) \{}
\NormalTok{  fit }\OtherTok{\textless{}{-}}\NormalTok{ mod\_ml\_pba}\SpecialCharTok{$}\FunctionTok{sample}\NormalTok{(}
    \AttributeTok{data            =}\NormalTok{ stan\_data\_ml,}
    \AttributeTok{chains          =} \DecValTok{4}\NormalTok{,}
    \AttributeTok{parallel\_chains =} \DecValTok{4}\NormalTok{,}
    \AttributeTok{iter\_warmup     =} \DecValTok{1000}\NormalTok{,}
    \AttributeTok{iter\_sampling   =} \DecValTok{1000}\NormalTok{,}
    \AttributeTok{adapt\_delta     =} \FloatTok{0.95}\NormalTok{,}
    \AttributeTok{refresh         =} \DecValTok{0}
\NormalTok{  )}
\NormalTok{  fit}\SpecialCharTok{$}\FunctionTok{save\_object}\NormalTok{(}\StringTok{"simmodels/ch10\_fit\_ml\_pba.rds"}\NormalTok{)}
\NormalTok{  fit}
\NormalTok{\} }\ControlFlowTok{else}\NormalTok{ \{}
  \FunctionTok{readRDS}\NormalTok{(}\StringTok{"simmodels/ch10\_fit\_ml\_pba.rds"}\NormalTok{)}
\NormalTok{\}}

\NormalTok{fit\_ml\_wba }\OtherTok{\textless{}{-}} \ControlFlowTok{if}\NormalTok{ (regenerate\_simulations }\SpecialCharTok{||} \SpecialCharTok{!}\FunctionTok{file.exists}\NormalTok{(}\StringTok{"simmodels/ch10\_fit\_ml\_wba.rds"}\NormalTok{)) \{}
\NormalTok{  fit }\OtherTok{\textless{}{-}}\NormalTok{ mod\_ml\_wba}\SpecialCharTok{$}\FunctionTok{sample}\NormalTok{(}
    \AttributeTok{data            =}\NormalTok{ stan\_data\_ml,}
    \AttributeTok{chains          =} \DecValTok{4}\NormalTok{,}
    \AttributeTok{parallel\_chains =} \DecValTok{4}\NormalTok{,}
    \AttributeTok{iter\_warmup     =} \DecValTok{1000}\NormalTok{,}
    \AttributeTok{iter\_sampling   =} \DecValTok{1000}\NormalTok{,}
    \AttributeTok{adapt\_delta     =} \FloatTok{0.95}\NormalTok{,}
    \AttributeTok{refresh         =} \DecValTok{0}
\NormalTok{  )}
\NormalTok{  fit}\SpecialCharTok{$}\FunctionTok{save\_object}\NormalTok{(}\StringTok{"simmodels/ch10\_fit\_ml\_wba.rds"}\NormalTok{)}
\NormalTok{  fit}
\NormalTok{\} }\ControlFlowTok{else}\NormalTok{ \{}
  \FunctionTok{readRDS}\NormalTok{(}\StringTok{"simmodels/ch10\_fit\_ml\_wba.rds"}\NormalTok{)}
\NormalTok{\}}



\CommentTok{\# Diagnostics for all three}
\FunctionTok{cat}\NormalTok{(}\StringTok{"=== Multilevel Model Diagnostics ===}\SpecialCharTok{\textbackslash{}n}\StringTok{"}\NormalTok{)}
\end{Highlighting}
\end{Shaded}

\begin{verbatim}
=== Multilevel Model Diagnostics ===
\end{verbatim}

\begin{Shaded}
\begin{Highlighting}[]
\FunctionTok{check\_diagnostics}\NormalTok{(fit\_ml\_sba, }\StringTok{"ML{-}SBA"}\NormalTok{)}
\end{Highlighting}
\end{Shaded}

\begin{verbatim}
ML-SBA                    | Divs:   0 | Min E-BFMI: NaN | Failures: 0
\end{verbatim}

\begin{Shaded}
\begin{Highlighting}[]
\FunctionTok{check\_diagnostics}\NormalTok{(fit\_ml\_pba, }\StringTok{"ML{-}PBA"}\NormalTok{)}
\end{Highlighting}
\end{Shaded}

\begin{verbatim}
ML-PBA                    | Divs:   0 | Min E-BFMI: 0.84 | Failures: 0
\end{verbatim}

\begin{Shaded}
\begin{Highlighting}[]
\FunctionTok{check\_diagnostics}\NormalTok{(fit\_ml\_wba, }\StringTok{"ML{-}WBA (reparameterised)"}\NormalTok{)}
\end{Highlighting}
\end{Shaded}

\begin{verbatim}
ML-WBA (reparameterised)  | Divs:   0 | Min E-BFMI: 0.89 | Failures: 0
\end{verbatim}

\begin{Shaded}
\begin{Highlighting}[]
\NormalTok{bayesplot}\SpecialCharTok{::}\FunctionTok{mcmc\_trace}\NormalTok{(}
\NormalTok{  fit\_ml\_wba}\SpecialCharTok{$}\FunctionTok{draws}\NormalTok{(}\FunctionTok{c}\NormalTok{(}\StringTok{"mu[1]"}\NormalTok{, }\StringTok{"mu[2]"}\NormalTok{, }\StringTok{"sigma[1]"}\NormalTok{, }\StringTok{"sigma[2]"}\NormalTok{, }\StringTok{"Omega[1,2]"}\NormalTok{)),}
  \AttributeTok{facet\_args =} \FunctionTok{list}\NormalTok{(}\AttributeTok{ncol =} \DecValTok{2}\NormalTok{)}
\NormalTok{) }\SpecialCharTok{+} \FunctionTok{labs}\NormalTok{(}\AttributeTok{title =} \StringTok{"Trace plots: ML{-}WBA — population parameters + correlation"}\NormalTok{)}
\end{Highlighting}
\end{Shaded}

\begin{center}
\includegraphics[width=0.8\linewidth,height=\textheight,keepaspectratio]{10-BayesianModels_files/figure-pdf/ch10_ml_trace_plots-1.pdf}
\end{center}

The trace plots should show the ``hairy caterpillar'' pattern for all
WBA population-level parameters (remember SBA has no parameter). Pay
particular attention to \(\sigma\): if either component pushes toward
zero, the NCP may struggle (the reverse funnel problem from Chapter 7).
If the SDs are well-identified away from zero, NCP is the right choice.

\subsection{Reporting the Estimated
Correlation}\label{reporting-the-estimated-correlation}

The posterior for \texttt{Omega{[}1,2{]}} --- the correlation between
\(\text{logit}(\rho)\) and \(\log(\kappa)\) across individuals ---
answers the cognitive question: \emph{do agents who allocate more
attention to direct evidence also use more or less total evidence?}

\begin{Shaded}
\begin{Highlighting}[]
\NormalTok{draws\_ml }\OtherTok{\textless{}{-}} \FunctionTok{as\_draws\_df}\NormalTok{(fit\_ml\_wba}\SpecialCharTok{$}\FunctionTok{draws}\NormalTok{())}

\FunctionTok{ggplot}\NormalTok{(draws\_ml, }\FunctionTok{aes}\NormalTok{(}\AttributeTok{x =} \StringTok{\textasciigrave{}}\AttributeTok{Omega[1,2]}\StringTok{\textasciigrave{}}\NormalTok{)) }\SpecialCharTok{+}
  \FunctionTok{geom\_histogram}\NormalTok{(}\AttributeTok{bins =} \DecValTok{60}\NormalTok{, }\AttributeTok{fill =} \StringTok{"\#756bb1"}\NormalTok{, }\AttributeTok{alpha =} \FloatTok{0.7}\NormalTok{) }\SpecialCharTok{+}
  \FunctionTok{geom\_vline}\NormalTok{(}\AttributeTok{xintercept =}\NormalTok{ pop\_correlation, }\AttributeTok{colour =} \StringTok{"\#d73027"}\NormalTok{,}
             \AttributeTok{linetype =} \StringTok{"dashed"}\NormalTok{, }\AttributeTok{linewidth =} \DecValTok{1}\NormalTok{) }\SpecialCharTok{+}
  \FunctionTok{geom\_vline}\NormalTok{(}\AttributeTok{xintercept =} \DecValTok{0}\NormalTok{, }\AttributeTok{linetype =} \StringTok{"dotted"}\NormalTok{, }\AttributeTok{colour =} \StringTok{"grey50"}\NormalTok{) }\SpecialCharTok{+}
  \FunctionTok{labs}\NormalTok{(}
    \AttributeTok{title    =} \StringTok{"Posterior: Correlation between logit(ρ) and log(κ)"}\NormalTok{,}
    \AttributeTok{subtitle =} \FunctionTok{sprintf}\NormalTok{(}\StringTok{"Red dashed = true generating value (\%.1f). Dotted = independence."}\NormalTok{,}
\NormalTok{                       pop\_correlation),}
    \AttributeTok{x        =} \FunctionTok{expression}\NormalTok{(Omega[}\StringTok{"ρ, κ"}\NormalTok{]),}
    \AttributeTok{y        =} \StringTok{"Posterior count"}
\NormalTok{  )}
\end{Highlighting}
\end{Shaded}

\begin{center}
\includegraphics[width=0.8\linewidth,height=\textheight,keepaspectratio]{10-BayesianModels_files/figure-pdf/ch10_correlation_posterior-1.pdf}
\end{center}

A negative correlation would suggest a tradeoff: agents who rely more on
direct evidence (high \(\rho\)) tend to scale total evidence less (low
\(\kappa\)). A correlation near zero indicates independence. The
posterior uncertainty around this estimate is moving its mass towards
the true value, but such a wide posterior means the experiment lacked
power to resolve the correlation structure.

\subsection{Population-Level Prior and Posterior Predictive
Checks}\label{population-level-prior-and-posterior-predictive-checks}

Just as we did for the single-agent models, we want to visualize what
our multilevel models believe before seeing the data, and what they
learned after.

Since we will have to run our slower multilevel code through e.g.SBC, we
kept the generated quantities leaner than usual: we didn't include
population-level prior or posterior predictions in the generated
quantities blocks (though we kept trial-level posterior predictions for
existing agents via pred\_choice). Instead, we can simulate these
hierarchical, population-level predictions dynamically in R. This is an
essential skill: by drawing from our population hyper-priors (\(\mu\)
and \(\sigma\)), sampling new individual agents, and pushing those
through the Beta-Binomial observation model, we can perfectly
reconstruct the model's population-level expectations.

\begin{Shaded}
\begin{Highlighting}[]
\CommentTok{\# 1. Custom function to simulate Multilevel Prior/Posterior Predictions in R}
\NormalTok{generate\_ml\_ppc }\OtherTok{\textless{}{-}} \ControlFlowTok{function}\NormalTok{(model\_name, fit, data\_df, }\AttributeTok{is\_prior =} \ConstantTok{FALSE}\NormalTok{, }\AttributeTok{n\_draws =} \DecValTok{400}\NormalTok{) \{}
  \CommentTok{\# Prep data}
\NormalTok{  obs\_df }\OtherTok{\textless{}{-}}\NormalTok{ data\_df }\SpecialCharTok{|\textgreater{}} 
    \FunctionTok{mutate}\NormalTok{(}\AttributeTok{obs\_id =} \FunctionTok{row\_number}\NormalTok{()) }\SpecialCharTok{|\textgreater{}}
    \FunctionTok{select}\NormalTok{(obs\_id, agent\_id, blue1, total1, blue2, total2, choice)}
  
  \CommentTok{\# Create base grid of draws x observations}
\NormalTok{  base\_grid }\OtherTok{\textless{}{-}} \FunctionTok{expand\_grid}\NormalTok{(}\AttributeTok{.draw =} \DecValTok{1}\SpecialCharTok{:}\NormalTok{n\_draws, }\AttributeTok{obs\_id =}\NormalTok{ obs\_df}\SpecialCharTok{$}\NormalTok{obs\_id) }\SpecialCharTok{|\textgreater{}}
    \FunctionTok{left\_join}\NormalTok{(obs\_df, }\AttributeTok{by =} \StringTok{"obs\_id"}\NormalTok{)}
  
  \ControlFlowTok{if}\NormalTok{ (model\_name }\SpecialCharTok{==} \StringTok{"ML{-}SBA"}\NormalTok{) \{}
    \CommentTok{\# SBA has no parameters, prior and posterior are identical}
\NormalTok{    pred\_samples }\OtherTok{\textless{}{-}}\NormalTok{ base\_grid }\SpecialCharTok{|\textgreater{}}
      \FunctionTok{mutate}\NormalTok{(}
        \AttributeTok{alpha\_p =} \FloatTok{0.5} \SpecialCharTok{+}\NormalTok{ blue1 }\SpecialCharTok{+}\NormalTok{ blue2,}
        \AttributeTok{beta\_p  =} \FloatTok{0.5} \SpecialCharTok{+}\NormalTok{ (total1 }\SpecialCharTok{{-}}\NormalTok{ blue1) }\SpecialCharTok{+}\NormalTok{ (total2 }\SpecialCharTok{{-}}\NormalTok{ blue2),}
        \AttributeTok{theta   =} \FunctionTok{rbeta}\NormalTok{(}\FunctionTok{n}\NormalTok{(), alpha\_p, beta\_p),}
        \AttributeTok{pred    =} \FunctionTok{rbinom}\NormalTok{(}\FunctionTok{n}\NormalTok{(), }\DecValTok{1}\NormalTok{, theta)}
\NormalTok{      )}
      
\NormalTok{  \} }\ControlFlowTok{else} \ControlFlowTok{if}\NormalTok{ (model\_name }\SpecialCharTok{==} \StringTok{"ML{-}PBA"}\NormalTok{) \{}
    \ControlFlowTok{if}\NormalTok{ (is\_prior) \{}
\NormalTok{      pred\_samples }\OtherTok{\textless{}{-}}\NormalTok{ base\_grid }\SpecialCharTok{|\textgreater{}}
        \FunctionTok{group\_by}\NormalTok{(.draw) }\SpecialCharTok{|\textgreater{}}
        \FunctionTok{mutate}\NormalTok{(}\AttributeTok{mu\_p =} \FunctionTok{rnorm}\NormalTok{(}\DecValTok{1}\NormalTok{, }\DecValTok{0}\NormalTok{, }\FloatTok{1.5}\NormalTok{), }\AttributeTok{sig\_p =} \FunctionTok{rexp}\NormalTok{(}\DecValTok{1}\NormalTok{, }\DecValTok{2}\NormalTok{)) }\SpecialCharTok{|\textgreater{}}
        \FunctionTok{group\_by}\NormalTok{(.draw, agent\_id) }\SpecialCharTok{|\textgreater{}}
        \FunctionTok{mutate}\NormalTok{(}\AttributeTok{p\_j =} \FunctionTok{plogis}\NormalTok{(mu\_p }\SpecialCharTok{+}\NormalTok{ sig\_p }\SpecialCharTok{*} \FunctionTok{rnorm}\NormalTok{(}\DecValTok{1}\NormalTok{))) }\SpecialCharTok{|\textgreater{}}
        \FunctionTok{ungroup}\NormalTok{() }\SpecialCharTok{|\textgreater{}}
        \FunctionTok{mutate}\NormalTok{(}
          \AttributeTok{alpha\_p =} \FloatTok{0.5} \SpecialCharTok{+}\NormalTok{ p\_j }\SpecialCharTok{*}\NormalTok{ blue1 }\SpecialCharTok{+}\NormalTok{ (}\DecValTok{1} \SpecialCharTok{{-}}\NormalTok{ p\_j) }\SpecialCharTok{*}\NormalTok{ blue2,}
          \AttributeTok{beta\_p  =} \FloatTok{0.5} \SpecialCharTok{+}\NormalTok{ p\_j }\SpecialCharTok{*}\NormalTok{ (total1 }\SpecialCharTok{{-}}\NormalTok{ blue1) }\SpecialCharTok{+}\NormalTok{ (}\DecValTok{1} \SpecialCharTok{{-}}\NormalTok{ p\_j) }\SpecialCharTok{*}\NormalTok{ (total2 }\SpecialCharTok{{-}}\NormalTok{ blue2),}
          \AttributeTok{theta   =} \FunctionTok{rbeta}\NormalTok{(}\FunctionTok{n}\NormalTok{(), alpha\_p, beta\_p),}
          \AttributeTok{pred    =} \FunctionTok{rbinom}\NormalTok{(}\FunctionTok{n}\NormalTok{(), }\DecValTok{1}\NormalTok{, theta)}
\NormalTok{        )}
\NormalTok{    \} }\ControlFlowTok{else}\NormalTok{ \{}
\NormalTok{      draws }\OtherTok{\textless{}{-}} \FunctionTok{as\_draws\_df}\NormalTok{(fit}\SpecialCharTok{$}\FunctionTok{draws}\NormalTok{(}\StringTok{"p"}\NormalTok{)) }\SpecialCharTok{|\textgreater{}} 
        \FunctionTok{sample\_n}\NormalTok{(n\_draws, }\AttributeTok{replace =} \ConstantTok{TRUE}\NormalTok{) }\SpecialCharTok{|\textgreater{}} \FunctionTok{mutate}\NormalTok{(}\AttributeTok{.draw =} \FunctionTok{row\_number}\NormalTok{()) }\SpecialCharTok{|\textgreater{}}
        \FunctionTok{select}\NormalTok{(.draw, }\FunctionTok{starts\_with}\NormalTok{(}\StringTok{"p["}\NormalTok{)) }\SpecialCharTok{|\textgreater{}} 
        \FunctionTok{pivot\_longer}\NormalTok{(}\SpecialCharTok{{-}}\NormalTok{.draw, }\AttributeTok{names\_to =} \StringTok{"param"}\NormalTok{, }\AttributeTok{values\_to =} \StringTok{"p\_j"}\NormalTok{) }\SpecialCharTok{|\textgreater{}}
        \FunctionTok{mutate}\NormalTok{(}\AttributeTok{agent\_id =}\NormalTok{ readr}\SpecialCharTok{::}\FunctionTok{parse\_number}\NormalTok{(param)) }\SpecialCharTok{|\textgreater{}} \FunctionTok{select}\NormalTok{(}\SpecialCharTok{{-}}\NormalTok{param)}
        
\NormalTok{      pred\_samples }\OtherTok{\textless{}{-}}\NormalTok{ base\_grid }\SpecialCharTok{|\textgreater{}} \FunctionTok{left\_join}\NormalTok{(draws, }\AttributeTok{by =} \FunctionTok{c}\NormalTok{(}\StringTok{".draw"}\NormalTok{, }\StringTok{"agent\_id"}\NormalTok{)) }\SpecialCharTok{|\textgreater{}}
        \FunctionTok{mutate}\NormalTok{(}
          \AttributeTok{alpha\_p =} \FloatTok{0.5} \SpecialCharTok{+}\NormalTok{ p\_j }\SpecialCharTok{*}\NormalTok{ blue1 }\SpecialCharTok{+}\NormalTok{ (}\DecValTok{1} \SpecialCharTok{{-}}\NormalTok{ p\_j) }\SpecialCharTok{*}\NormalTok{ blue2,}
          \AttributeTok{beta\_p  =} \FloatTok{0.5} \SpecialCharTok{+}\NormalTok{ p\_j }\SpecialCharTok{*}\NormalTok{ (total1 }\SpecialCharTok{{-}}\NormalTok{ blue1) }\SpecialCharTok{+}\NormalTok{ (}\DecValTok{1} \SpecialCharTok{{-}}\NormalTok{ p\_j) }\SpecialCharTok{*}\NormalTok{ (total2 }\SpecialCharTok{{-}}\NormalTok{ blue2),}
          \AttributeTok{theta   =} \FunctionTok{rbeta}\NormalTok{(}\FunctionTok{n}\NormalTok{(), alpha\_p, beta\_p),}
          \AttributeTok{pred    =} \FunctionTok{rbinom}\NormalTok{(}\FunctionTok{n}\NormalTok{(), }\DecValTok{1}\NormalTok{, theta)}
\NormalTok{        )}
\NormalTok{    \}}
    
\NormalTok{  \} }\ControlFlowTok{else} \ControlFlowTok{if}\NormalTok{ (model\_name }\SpecialCharTok{==} \StringTok{"ML{-}WBA"}\NormalTok{) \{}
    \ControlFlowTok{if}\NormalTok{ (is\_prior) \{}
\NormalTok{      pred\_samples }\OtherTok{\textless{}{-}}\NormalTok{ base\_grid }\SpecialCharTok{|\textgreater{}}
        \FunctionTok{group\_by}\NormalTok{(.draw) }\SpecialCharTok{|\textgreater{}}
        \FunctionTok{mutate}\NormalTok{(}
          \AttributeTok{mu\_1 =} \FunctionTok{rnorm}\NormalTok{(}\DecValTok{1}\NormalTok{, }\DecValTok{0}\NormalTok{, }\FloatTok{1.5}\NormalTok{), }\AttributeTok{mu\_2 =} \FunctionTok{rnorm}\NormalTok{(}\DecValTok{1}\NormalTok{, }\FunctionTok{log}\NormalTok{(}\DecValTok{2}\NormalTok{), }\FloatTok{0.5}\NormalTok{),}
          \AttributeTok{sig\_1 =} \FunctionTok{rexp}\NormalTok{(}\DecValTok{1}\NormalTok{, }\DecValTok{2}\NormalTok{), }\AttributeTok{sig\_2 =} \FunctionTok{rexp}\NormalTok{(}\DecValTok{1}\NormalTok{, }\DecValTok{2}\NormalTok{),}
          \AttributeTok{omega =} \DecValTok{2} \SpecialCharTok{*} \FunctionTok{rbeta}\NormalTok{(}\DecValTok{1}\NormalTok{, }\DecValTok{2}\NormalTok{, }\DecValTok{2}\NormalTok{) }\SpecialCharTok{{-}} \DecValTok{1}\NormalTok{,}
          \AttributeTok{L21 =}\NormalTok{ omega, }\AttributeTok{L22 =} \FunctionTok{sqrt}\NormalTok{(}\DecValTok{1} \SpecialCharTok{{-}}\NormalTok{ omega}\SpecialCharTok{\^{}}\DecValTok{2}\NormalTok{)}
\NormalTok{        ) }\SpecialCharTok{|\textgreater{}}
        \FunctionTok{group\_by}\NormalTok{(.draw, agent\_id) }\SpecialCharTok{|\textgreater{}}
        \FunctionTok{mutate}\NormalTok{(}
          \AttributeTok{z1 =} \FunctionTok{rnorm}\NormalTok{(}\DecValTok{1}\NormalTok{), }\AttributeTok{z2 =} \FunctionTok{rnorm}\NormalTok{(}\DecValTok{1}\NormalTok{),}
          \AttributeTok{rho\_j =} \FunctionTok{plogis}\NormalTok{(mu\_1 }\SpecialCharTok{+}\NormalTok{ sig\_1 }\SpecialCharTok{*}\NormalTok{ z1),}
          \AttributeTok{kappa\_j =} \FunctionTok{exp}\NormalTok{(mu\_2 }\SpecialCharTok{+}\NormalTok{ sig\_2 }\SpecialCharTok{*}\NormalTok{ (L21 }\SpecialCharTok{*}\NormalTok{ z1 }\SpecialCharTok{+}\NormalTok{ L22 }\SpecialCharTok{*}\NormalTok{ z2)),}
          \AttributeTok{wd\_j =}\NormalTok{ rho\_j }\SpecialCharTok{*}\NormalTok{ kappa\_j, }\AttributeTok{ws\_j =}\NormalTok{ (}\DecValTok{1} \SpecialCharTok{{-}}\NormalTok{ rho\_j) }\SpecialCharTok{*}\NormalTok{ kappa\_j}
\NormalTok{        ) }\SpecialCharTok{|\textgreater{}}
        \FunctionTok{ungroup}\NormalTok{() }\SpecialCharTok{|\textgreater{}}
        \FunctionTok{mutate}\NormalTok{(}
          \AttributeTok{alpha\_p =} \FloatTok{0.5} \SpecialCharTok{+}\NormalTok{ wd\_j }\SpecialCharTok{*}\NormalTok{ blue1 }\SpecialCharTok{+}\NormalTok{ ws\_j }\SpecialCharTok{*}\NormalTok{ blue2,}
          \AttributeTok{beta\_p  =} \FloatTok{0.5} \SpecialCharTok{+}\NormalTok{ wd\_j }\SpecialCharTok{*}\NormalTok{ (total1 }\SpecialCharTok{{-}}\NormalTok{ blue1) }\SpecialCharTok{+}\NormalTok{ ws\_j }\SpecialCharTok{*}\NormalTok{ (total2 }\SpecialCharTok{{-}}\NormalTok{ blue2),}
          \AttributeTok{theta   =} \FunctionTok{rbeta}\NormalTok{(}\FunctionTok{n}\NormalTok{(), alpha\_p, beta\_p),}
          \AttributeTok{pred    =} \FunctionTok{rbinom}\NormalTok{(}\FunctionTok{n}\NormalTok{(), }\DecValTok{1}\NormalTok{, theta)}
\NormalTok{        )}
\NormalTok{    \} }\ControlFlowTok{else}\NormalTok{ \{}
\NormalTok{      draws }\OtherTok{\textless{}{-}} \FunctionTok{as\_draws\_df}\NormalTok{(fit}\SpecialCharTok{$}\FunctionTok{draws}\NormalTok{(}\FunctionTok{c}\NormalTok{(}\StringTok{"weight\_direct"}\NormalTok{, }\StringTok{"weight\_social"}\NormalTok{))) }\SpecialCharTok{|\textgreater{}}
        \FunctionTok{sample\_n}\NormalTok{(n\_draws, }\AttributeTok{replace =} \ConstantTok{TRUE}\NormalTok{) }\SpecialCharTok{|\textgreater{}} \FunctionTok{mutate}\NormalTok{(}\AttributeTok{.draw =} \FunctionTok{row\_number}\NormalTok{()) }\SpecialCharTok{|\textgreater{}}
        \FunctionTok{select}\NormalTok{(.draw, }\FunctionTok{starts\_with}\NormalTok{(}\StringTok{"weight\_"}\NormalTok{)) }\SpecialCharTok{|\textgreater{}}
        \FunctionTok{pivot\_longer}\NormalTok{(}\SpecialCharTok{{-}}\NormalTok{.draw, }\AttributeTok{names\_to =} \StringTok{"param"}\NormalTok{, }\AttributeTok{values\_to =} \StringTok{"val"}\NormalTok{) }\SpecialCharTok{|\textgreater{}}
        \FunctionTok{mutate}\NormalTok{(}
          \AttributeTok{agent\_id =}\NormalTok{ readr}\SpecialCharTok{::}\FunctionTok{parse\_number}\NormalTok{(param), }
          \AttributeTok{param\_type =} \FunctionTok{str\_extract}\NormalTok{(param, }\StringTok{"weight\_[a{-}z]+"}\NormalTok{)}
\NormalTok{        ) }\SpecialCharTok{|\textgreater{}}
        \FunctionTok{select}\NormalTok{(}\SpecialCharTok{{-}}\NormalTok{param) }\SpecialCharTok{|\textgreater{}} \FunctionTok{pivot\_wider}\NormalTok{(}\AttributeTok{names\_from =}\NormalTok{ param\_type, }\AttributeTok{values\_from =}\NormalTok{ val)}
        
\NormalTok{      pred\_samples }\OtherTok{\textless{}{-}}\NormalTok{ base\_grid }\SpecialCharTok{|\textgreater{}} \FunctionTok{left\_join}\NormalTok{(draws, }\AttributeTok{by =} \FunctionTok{c}\NormalTok{(}\StringTok{".draw"}\NormalTok{, }\StringTok{"agent\_id"}\NormalTok{)) }\SpecialCharTok{|\textgreater{}}
        \FunctionTok{mutate}\NormalTok{(}
          \AttributeTok{alpha\_p =} \FloatTok{0.5} \SpecialCharTok{+}\NormalTok{ weight\_direct }\SpecialCharTok{*}\NormalTok{ blue1 }\SpecialCharTok{+}\NormalTok{ weight\_social }\SpecialCharTok{*}\NormalTok{ blue2,}
          \AttributeTok{beta\_p  =} \FloatTok{0.5} \SpecialCharTok{+}\NormalTok{ weight\_direct }\SpecialCharTok{*}\NormalTok{ (total1 }\SpecialCharTok{{-}}\NormalTok{ blue1) }\SpecialCharTok{+}\NormalTok{ weight\_social }\SpecialCharTok{*}\NormalTok{ (total2 }\SpecialCharTok{{-}}\NormalTok{ blue2),}
          \AttributeTok{theta   =} \FunctionTok{rbeta}\NormalTok{(}\FunctionTok{n}\NormalTok{(), alpha\_p, beta\_p),}
          \AttributeTok{pred    =} \FunctionTok{rbinom}\NormalTok{(}\FunctionTok{n}\NormalTok{(), }\DecValTok{1}\NormalTok{, theta)}
\NormalTok{        )}
\NormalTok{    \}}
\NormalTok{  \}}
  
  \CommentTok{\# Summarize (Averaging over all agents to get population{-}level PPC)}
\NormalTok{  obs\_summary }\OtherTok{\textless{}{-}}\NormalTok{ obs\_df }\SpecialCharTok{|\textgreater{}} \FunctionTok{group\_by}\NormalTok{(blue1, blue2) }\SpecialCharTok{|\textgreater{}} \FunctionTok{summarise}\NormalTok{(}\AttributeTok{p\_obs =} \FunctionTok{mean}\NormalTok{(choice), }\AttributeTok{.groups =} \StringTok{"drop"}\NormalTok{)}
  
\NormalTok{  pred\_summary }\OtherTok{\textless{}{-}}\NormalTok{ pred\_samples }\SpecialCharTok{|\textgreater{}}
    \FunctionTok{group\_by}\NormalTok{(.draw, blue1, blue2) }\SpecialCharTok{|\textgreater{}} \FunctionTok{summarise}\NormalTok{(}\AttributeTok{p\_draw =} \FunctionTok{mean}\NormalTok{(pred), }\AttributeTok{.groups =} \StringTok{"drop"}\NormalTok{) }\SpecialCharTok{|\textgreater{}}
    \FunctionTok{group\_by}\NormalTok{(blue1, blue2) }\SpecialCharTok{|\textgreater{}}
    \FunctionTok{summarise}\NormalTok{(}
      \AttributeTok{p\_pred =} \FunctionTok{median}\NormalTok{(p\_draw), }
      \AttributeTok{lower =} \FunctionTok{quantile}\NormalTok{(p\_draw, }\FloatTok{0.025}\NormalTok{), }
      \AttributeTok{upper =} \FunctionTok{quantile}\NormalTok{(p\_draw, }\FloatTok{0.975}\NormalTok{), }
      \AttributeTok{.groups =} \StringTok{"drop"}
\NormalTok{    ) }\SpecialCharTok{|\textgreater{}}
    \FunctionTok{left\_join}\NormalTok{(obs\_summary, }\AttributeTok{by =} \FunctionTok{c}\NormalTok{(}\StringTok{"blue1"}\NormalTok{, }\StringTok{"blue2"}\NormalTok{)) }\SpecialCharTok{|\textgreater{}}
    \FunctionTok{mutate}\NormalTok{(}
      \AttributeTok{social\_evidence =} \FunctionTok{factor}\NormalTok{(blue2, }\AttributeTok{levels =} \DecValTok{0}\SpecialCharTok{:}\DecValTok{3}\NormalTok{, }\AttributeTok{labels =} \FunctionTok{c}\NormalTok{(}\StringTok{"Clear Red"}\NormalTok{, }\StringTok{"Maybe Red"}\NormalTok{, }\StringTok{"Maybe Blue"}\NormalTok{, }\StringTok{"Clear Blue"}\NormalTok{)),}
      \AttributeTok{Scenario =} \FunctionTok{ifelse}\NormalTok{(is\_prior, }\FunctionTok{paste}\NormalTok{(model\_name, }\StringTok{"{-} Prior"}\NormalTok{), }\FunctionTok{paste}\NormalTok{(model\_name, }\StringTok{"{-} Posterior"}\NormalTok{))}
\NormalTok{    )}
    
  \FunctionTok{return}\NormalTok{(pred\_summary)}
\NormalTok{\}}

\CommentTok{\# 2. Wrapper to plot them side{-}by{-}side}
\NormalTok{build\_ml\_ppc\_row }\OtherTok{\textless{}{-}} \ControlFlowTok{function}\NormalTok{(model\_name, fit) \{}
\NormalTok{  prior\_data }\OtherTok{\textless{}{-}} \FunctionTok{generate\_ml\_ppc}\NormalTok{(model\_name, fit, ml\_sim\_data, }\AttributeTok{is\_prior =} \ConstantTok{TRUE}\NormalTok{)}
\NormalTok{  p\_prior }\OtherTok{\textless{}{-}} \FunctionTok{plot\_predictive\_check}\NormalTok{(prior\_data, }\FunctionTok{paste}\NormalTok{(model\_name, }\StringTok{"Prior"}\NormalTok{), }\StringTok{"Population Prior Predictive Check"}\NormalTok{) }\SpecialCharTok{+} 
    \FunctionTok{theme}\NormalTok{(}\AttributeTok{legend.position =} \StringTok{"none"}\NormalTok{)}
  
\NormalTok{  post\_data }\OtherTok{\textless{}{-}} \FunctionTok{generate\_ml\_ppc}\NormalTok{(model\_name, fit, ml\_sim\_data, }\AttributeTok{is\_prior =} \ConstantTok{FALSE}\NormalTok{)}
\NormalTok{  p\_post }\OtherTok{\textless{}{-}} \FunctionTok{plot\_predictive\_check}\NormalTok{(post\_data, }\FunctionTok{paste}\NormalTok{(model\_name, }\StringTok{"Posterior"}\NormalTok{), }\StringTok{"Population Posterior Predictive Check"}\NormalTok{) }\SpecialCharTok{+} 
    \FunctionTok{theme}\NormalTok{(}\AttributeTok{axis.title.y =} \FunctionTok{element\_blank}\NormalTok{())}
    
  \FunctionTok{return}\NormalTok{(p\_prior }\SpecialCharTok{|}\NormalTok{ p\_post)}
\NormalTok{\}}

\CommentTok{\# 3. Generate and combine}
\FunctionTok{cat}\NormalTok{(}\StringTok{"Generating Multilevel Predictive Checks (this takes a few seconds)...}\SpecialCharTok{\textbackslash{}n}\StringTok{"}\NormalTok{)}
\end{Highlighting}
\end{Shaded}

\begin{verbatim}
Generating Multilevel Predictive Checks (this takes a few seconds)...
\end{verbatim}

\begin{Shaded}
\begin{Highlighting}[]
\NormalTok{p\_ml\_sba }\OtherTok{\textless{}{-}} \FunctionTok{build\_ml\_ppc\_row}\NormalTok{(}\StringTok{"ML{-}SBA"}\NormalTok{, fit\_ml\_sba)}
\NormalTok{p\_ml\_pba }\OtherTok{\textless{}{-}} \FunctionTok{build\_ml\_ppc\_row}\NormalTok{(}\StringTok{"ML{-}PBA"}\NormalTok{, fit\_ml\_pba)}
\NormalTok{p\_ml\_wba }\OtherTok{\textless{}{-}} \FunctionTok{build\_ml\_ppc\_row}\NormalTok{(}\StringTok{"ML{-}WBA"}\NormalTok{, fit\_ml\_wba)}

\NormalTok{(p\_ml\_sba }\SpecialCharTok{/}\NormalTok{ p\_ml\_pba }\SpecialCharTok{/}\NormalTok{ p\_ml\_wba) }\SpecialCharTok{+} 
  \FunctionTok{plot\_layout}\NormalTok{(}\AttributeTok{guides =} \StringTok{"collect"}\NormalTok{) }\SpecialCharTok{\&} 
  \FunctionTok{theme}\NormalTok{(}\AttributeTok{legend.position =} \StringTok{"bottom"}\NormalTok{)}
\end{Highlighting}
\end{Shaded}

\begin{center}
\includegraphics[width=0.8\linewidth,height=\textheight,keepaspectratio]{10-BayesianModels_files/figure-pdf/ch10_multilevel_ppc-1.pdf}
\end{center}

It might look as if the models haven't updated their predictions after
seeing the data. That's indeed true for SBA. However, PBA and WBA have.
The reason for the changes being so minimal is that the simulated
parameters are indeed close to SBA (a balanced population level
behavior). So we need to rather look at individual agents. We could run
predictive checks at the agent level, but for now we move to parameter
recovery to see how well the models estimated the true parameters for
each agent, and whether the multilevel WBA shows appropriate shrinkage
toward the population mean.

\section{Parameter Recovery: Single-shot Multilevel
Model}\label{parameter-recovery-single-shot-multilevel-model}

Parameter recovery at the multilevel level must check two things: (1)
population-level parameters are recovered, and (2) individual-level
parameters are recovered and show appropriate partial pooling ---
extreme agents are pulled toward the population mean.

\begin{Shaded}
\begin{Highlighting}[]
\CommentTok{\# 0. Obtain UNPOOLED estimates via fast optimization (MAP) — used here and in the next chunk}
\NormalTok{unpooled\_ests }\OtherTok{\textless{}{-}} \FunctionTok{map\_dfr}\NormalTok{(}\FunctionTok{unique}\NormalTok{(ml\_sim\_data}\SpecialCharTok{$}\NormalTok{agent\_id), }\ControlFlowTok{function}\NormalTok{(id) \{}
\NormalTok{  agent\_df }\OtherTok{\textless{}{-}}\NormalTok{ ml\_sim\_data }\SpecialCharTok{|\textgreater{}} \FunctionTok{filter}\NormalTok{(agent\_id }\SpecialCharTok{==}\NormalTok{ id)}
\NormalTok{  agent\_stan\_data }\OtherTok{\textless{}{-}} \FunctionTok{prepare\_stan\_data}\NormalTok{(agent\_df)}

\NormalTok{  opt }\OtherTok{\textless{}{-}}\NormalTok{ mod\_weighted}\SpecialCharTok{$}\FunctionTok{optimize}\NormalTok{(}\AttributeTok{data =}\NormalTok{ agent\_stan\_data, }\AttributeTok{refresh =} \DecValTok{0}\NormalTok{)}

  \FunctionTok{tibble}\NormalTok{(}
    \AttributeTok{agent\_id       =}\NormalTok{ id,}
    \AttributeTok{unpooled\_rho   =}\NormalTok{ opt}\SpecialCharTok{$}\FunctionTok{mle}\NormalTok{(}\StringTok{"rho"}\NormalTok{),}
    \AttributeTok{unpooled\_kappa =}\NormalTok{ opt}\SpecialCharTok{$}\FunctionTok{mle}\NormalTok{(}\StringTok{"kappa"}\NormalTok{),}
    \AttributeTok{unpooled\_wd    =}\NormalTok{ opt}\SpecialCharTok{$}\FunctionTok{mle}\NormalTok{(}\StringTok{"weight\_direct"}\NormalTok{),}
    \AttributeTok{unpooled\_ws    =}\NormalTok{ opt}\SpecialCharTok{$}\FunctionTok{mle}\NormalTok{(}\StringTok{"weight\_social"}\NormalTok{)}
\NormalTok{  )}
\NormalTok{\})}
\end{Highlighting}
\end{Shaded}

\begin{verbatim}
Initial log joint probability = -50.9439 
Optimization terminated normally:  
  Convergence detected: relative gradient magnitude is below tolerance 
Finished in  0.2 seconds.
Initial log joint probability = -48.897 
Optimization terminated normally:  
  Convergence detected: relative gradient magnitude is below tolerance 
Finished in  0.2 seconds.
Initial log joint probability = -43.5658 
Optimization terminated normally:  
  Convergence detected: relative gradient magnitude is below tolerance 
Finished in  0.2 seconds.
Initial log joint probability = -37.9017 
Optimization terminated normally:  
  Convergence detected: relative gradient magnitude is below tolerance 
Finished in  0.2 seconds.
Initial log joint probability = -50.0259 
Optimization terminated normally:  
  Convergence detected: relative gradient magnitude is below tolerance 
Finished in  0.1 seconds.
Initial log joint probability = -44.7817 
Optimization terminated normally:  
  Convergence detected: relative gradient magnitude is below tolerance 
Finished in  0.1 seconds.
Initial log joint probability = -51.6485 
Optimization terminated normally:  
  Convergence detected: relative gradient magnitude is below tolerance 
Finished in  0.1 seconds.
Initial log joint probability = -52.247 
Optimization terminated normally:  
  Convergence detected: relative gradient magnitude is below tolerance 
Finished in  0.1 seconds.
Initial log joint probability = -47.1458 
Optimization terminated normally:  
  Convergence detected: relative gradient magnitude is below tolerance 
Finished in  0.2 seconds.
Initial log joint probability = -44.2399 
Optimization terminated normally:  
  Convergence detected: relative gradient magnitude is below tolerance 
Finished in  0.1 seconds.
Initial log joint probability = -40.8946 
Optimization terminated normally:  
  Convergence detected: relative gradient magnitude is below tolerance 
Finished in  0.2 seconds.
Initial log joint probability = -40.5574 
Optimization terminated normally:  
  Convergence detected: relative gradient magnitude is below tolerance 
Finished in  0.2 seconds.
Initial log joint probability = -39.5652 
Optimization terminated normally:  
  Convergence detected: relative gradient magnitude is below tolerance 
Finished in  0.1 seconds.
Initial log joint probability = -53.0727 
Optimization terminated normally:  
  Convergence detected: relative gradient magnitude is below tolerance 
Finished in  0.1 seconds.
Initial log joint probability = -43.3977 
Optimization terminated normally:  
  Convergence detected: relative gradient magnitude is below tolerance 
Finished in  0.1 seconds.
Initial log joint probability = -39.645 
Optimization terminated normally:  
  Convergence detected: relative gradient magnitude is below tolerance 
Finished in  0.1 seconds.
Initial log joint probability = -49.6907 
Optimization terminated normally:  
  Convergence detected: relative gradient magnitude is below tolerance 
Finished in  0.2 seconds.
Initial log joint probability = -48.0459 
Optimization terminated normally:  
  Convergence detected: relative gradient magnitude is below tolerance 
Finished in  0.2 seconds.
Initial log joint probability = -41.6331 
Optimization terminated normally:  
  Convergence detected: relative gradient magnitude is below tolerance 
Finished in  0.2 seconds.
Initial log joint probability = -54.6744 
Optimization terminated normally:  
  Convergence detected: relative gradient magnitude is below tolerance 
Finished in  0.1 seconds.
\end{verbatim}

\begin{Shaded}
\begin{Highlighting}[]
\CommentTok{\# 1. Prepare the hyperparameter summary data}
\NormalTok{pop\_summary }\OtherTok{\textless{}{-}}\NormalTok{ fit\_ml\_wba}\SpecialCharTok{$}\FunctionTok{summary}\NormalTok{(}
  \AttributeTok{variables =} \FunctionTok{c}\NormalTok{(}\StringTok{"mu[1]"}\NormalTok{, }\StringTok{"mu[2]"}\NormalTok{, }\StringTok{"sigma[1]"}\NormalTok{, }\StringTok{"sigma[2]"}\NormalTok{, }\StringTok{"Omega[1,2]"}\NormalTok{)}
\NormalTok{) }\SpecialCharTok{|\textgreater{}}
  \FunctionTok{mutate}\NormalTok{(}
    \AttributeTok{true\_value =} \FunctionTok{c}\NormalTok{(pop\_logit\_rho\_mean, pop\_log\_kappa\_mean,}
\NormalTok{                   pop\_logit\_rho\_sd, pop\_log\_kappa\_sd, pop\_correlation),}
    \CommentTok{\# Clean labels for the plot}
    \AttributeTok{clean\_var =} \FunctionTok{c}\NormalTok{(}\StringTok{"Mean: logit(ρ)"}\NormalTok{, }\StringTok{"Mean: log(κ)"}\NormalTok{, }
                  \StringTok{"SD: logit(ρ)"}\NormalTok{, }\StringTok{"SD: log(κ)"}\NormalTok{, }\StringTok{"Correlation: Ω"}\NormalTok{)}
\NormalTok{  )}

\CommentTok{\# Extract point estimates for the distributions}
\NormalTok{pop\_est }\OtherTok{\textless{}{-}}\NormalTok{ pop\_summary }\SpecialCharTok{|\textgreater{}} \FunctionTok{select}\NormalTok{(variable, mean) }\SpecialCharTok{|\textgreater{}} \FunctionTok{deframe}\NormalTok{()}

\CommentTok{\# 2. Transform unpooled MAP estimates to the unconstrained population scale}
\NormalTok{unpooled\_pop }\OtherTok{\textless{}{-}}\NormalTok{ unpooled\_ests }\SpecialCharTok{|\textgreater{}}
  \FunctionTok{mutate}\NormalTok{(}
    \AttributeTok{logit\_rho =} \FunctionTok{qlogis}\NormalTok{(}\FunctionTok{pmin}\NormalTok{(}\FunctionTok{pmax}\NormalTok{(unpooled\_rho, }\FloatTok{0.001}\NormalTok{), }\FloatTok{0.999}\NormalTok{)),}
    \AttributeTok{log\_kappa =} \FunctionTok{log}\NormalTok{(}\FunctionTok{pmax}\NormalTok{(unpooled\_kappa, }\FloatTok{0.001}\NormalTok{)) }\CommentTok{\# safeguard against exactly 0}
\NormalTok{  )}

\CommentTok{\# 3. Calculate summary statistics of the unpooled estimates}
\NormalTok{unpooled\_rho\_mean }\OtherTok{\textless{}{-}} \FunctionTok{mean}\NormalTok{(unpooled\_pop}\SpecialCharTok{$}\NormalTok{logit\_rho)}
\NormalTok{unpooled\_rho\_sd   }\OtherTok{\textless{}{-}} \FunctionTok{sd}\NormalTok{(unpooled\_pop}\SpecialCharTok{$}\NormalTok{logit\_rho)}
\NormalTok{unpooled\_kappa\_mean }\OtherTok{\textless{}{-}} \FunctionTok{mean}\NormalTok{(unpooled\_pop}\SpecialCharTok{$}\NormalTok{log\_kappa)}
\NormalTok{unpooled\_kappa\_sd   }\OtherTok{\textless{}{-}} \FunctionTok{sd}\NormalTok{(unpooled\_pop}\SpecialCharTok{$}\NormalTok{log\_kappa)}
\NormalTok{unpooled\_cor      }\OtherTok{\textless{}{-}} \FunctionTok{cor}\NormalTok{(unpooled\_pop}\SpecialCharTok{$}\NormalTok{logit\_rho, unpooled\_pop}\SpecialCharTok{$}\NormalTok{log\_kappa)}

\CommentTok{\# Create a dataframe of these unpooled sample statistics}
\NormalTok{unpooled\_hyper }\OtherTok{\textless{}{-}} \FunctionTok{tibble}\NormalTok{(}
  \AttributeTok{clean\_var =} \FunctionTok{c}\NormalTok{(}\StringTok{"Mean: logit(ρ)"}\NormalTok{, }\StringTok{"Mean: log(κ)"}\NormalTok{, }
                \StringTok{"SD: logit(ρ)"}\NormalTok{, }\StringTok{"SD: log(κ)"}\NormalTok{, }\StringTok{"Correlation: Ω"}\NormalTok{),}
  \AttributeTok{unpooled\_stat =} \FunctionTok{c}\NormalTok{(unpooled\_rho\_mean, unpooled\_kappa\_mean, }
\NormalTok{                    unpooled\_rho\_sd, unpooled\_kappa\_sd, unpooled\_cor)}
\NormalTok{)}

\CommentTok{\# Join with pop\_summary for the forest plot}
\NormalTok{plot\_hyper\_data }\OtherTok{\textless{}{-}}\NormalTok{ pop\_summary }\SpecialCharTok{|\textgreater{}}
  \FunctionTok{left\_join}\NormalTok{(unpooled\_hyper, }\AttributeTok{by =} \StringTok{"clean\_var"}\NormalTok{) }\SpecialCharTok{|\textgreater{}}
  \FunctionTok{mutate}\NormalTok{(}\AttributeTok{clean\_var =} \FunctionTok{fct\_rev}\NormalTok{(}\FunctionTok{factor}\NormalTok{(clean\_var, }\AttributeTok{levels =}\NormalTok{ clean\_var)))}

\CommentTok{\# 4. Create grid points for ALL theoretical Normal distributions}
\NormalTok{x\_rho   }\OtherTok{\textless{}{-}} \FunctionTok{seq}\NormalTok{(}\SpecialCharTok{{-}}\DecValTok{4}\NormalTok{, }\DecValTok{4}\NormalTok{, }\AttributeTok{length.out =} \DecValTok{200}\NormalTok{)}
\NormalTok{x\_kappa }\OtherTok{\textless{}{-}} \FunctionTok{seq}\NormalTok{(}\SpecialCharTok{{-}}\DecValTok{1}\NormalTok{, }\DecValTok{3}\NormalTok{, }\AttributeTok{length.out =} \DecValTok{200}\NormalTok{)}

\NormalTok{dist\_rho }\OtherTok{\textless{}{-}} \FunctionTok{tibble}\NormalTok{(}
  \AttributeTok{x =}\NormalTok{ x\_rho,}
  \AttributeTok{True =} \FunctionTok{dnorm}\NormalTok{(x, pop\_logit\_rho\_mean, pop\_logit\_rho\_sd),}
  \AttributeTok{Hierarchical =} \FunctionTok{dnorm}\NormalTok{(x, pop\_est[}\StringTok{"mu[1]"}\NormalTok{], pop\_est[}\StringTok{"sigma[1]"}\NormalTok{]),}
  \AttributeTok{Unpooled =} \FunctionTok{dnorm}\NormalTok{(x, unpooled\_rho\_mean, unpooled\_rho\_sd)}
\NormalTok{) }\SpecialCharTok{|\textgreater{}} \FunctionTok{pivot\_longer}\NormalTok{(}\SpecialCharTok{{-}}\NormalTok{x, }\AttributeTok{names\_to =} \StringTok{"Distribution"}\NormalTok{, }\AttributeTok{values\_to =} \StringTok{"density"}\NormalTok{)}

\NormalTok{dist\_kappa }\OtherTok{\textless{}{-}} \FunctionTok{tibble}\NormalTok{(}
  \AttributeTok{x =}\NormalTok{ x\_kappa,}
  \AttributeTok{True =} \FunctionTok{dnorm}\NormalTok{(x, pop\_log\_kappa\_mean, pop\_log\_kappa\_sd),}
  \AttributeTok{Hierarchical =} \FunctionTok{dnorm}\NormalTok{(x, pop\_est[}\StringTok{"mu[2]"}\NormalTok{], pop\_est[}\StringTok{"sigma[2]"}\NormalTok{]),}
  \AttributeTok{Unpooled =} \FunctionTok{dnorm}\NormalTok{(x, unpooled\_kappa\_mean, unpooled\_kappa\_sd)}
\NormalTok{) }\SpecialCharTok{|\textgreater{}} \FunctionTok{pivot\_longer}\NormalTok{(}\SpecialCharTok{{-}}\NormalTok{x, }\AttributeTok{names\_to =} \StringTok{"Distribution"}\NormalTok{, }\AttributeTok{values\_to =} \StringTok{"density"}\NormalTok{)}

\CommentTok{\# 5. Build the plots}
\NormalTok{p\_pop\_rho }\OtherTok{\textless{}{-}} \FunctionTok{ggplot}\NormalTok{() }\SpecialCharTok{+}
  \FunctionTok{geom\_area}\NormalTok{(}\AttributeTok{data =} \FunctionTok{filter}\NormalTok{(dist\_rho, Distribution }\SpecialCharTok{==} \StringTok{"Unpooled"}\NormalTok{), }
            \FunctionTok{aes}\NormalTok{(}\AttributeTok{x =}\NormalTok{ x, }\AttributeTok{y =}\NormalTok{ density, }\AttributeTok{fill =} \StringTok{"Unpooled Estimates"}\NormalTok{), }\AttributeTok{alpha =} \FloatTok{0.3}\NormalTok{) }\SpecialCharTok{+}
  \FunctionTok{geom\_line}\NormalTok{(}\AttributeTok{data =} \FunctionTok{filter}\NormalTok{(dist\_rho, Distribution }\SpecialCharTok{!=} \StringTok{"Unpooled"}\NormalTok{), }
            \FunctionTok{aes}\NormalTok{(}\AttributeTok{x =}\NormalTok{ x, }\AttributeTok{y =}\NormalTok{ density, }\AttributeTok{color =}\NormalTok{ Distribution, }\AttributeTok{linetype =}\NormalTok{ Distribution), }
            \AttributeTok{linewidth =} \FloatTok{1.2}\NormalTok{) }\SpecialCharTok{+}
  \FunctionTok{scale\_color\_manual}\NormalTok{(}\AttributeTok{values =} \FunctionTok{c}\NormalTok{(}\StringTok{"True"} \OtherTok{=} \StringTok{"black"}\NormalTok{, }\StringTok{"Hierarchical"} \OtherTok{=} \StringTok{"\#4575b4"}\NormalTok{)) }\SpecialCharTok{+}
  \FunctionTok{scale\_linetype\_manual}\NormalTok{(}\AttributeTok{values =} \FunctionTok{c}\NormalTok{(}\StringTok{"True"} \OtherTok{=} \StringTok{"dashed"}\NormalTok{, }\StringTok{"Hierarchical"} \OtherTok{=} \StringTok{"solid"}\NormalTok{)) }\SpecialCharTok{+}
  \FunctionTok{scale\_fill\_manual}\NormalTok{(}\AttributeTok{values =} \FunctionTok{c}\NormalTok{(}\StringTok{"Unpooled Estimates"} \OtherTok{=} \StringTok{"grey60"}\NormalTok{)) }\SpecialCharTok{+}
  \FunctionTok{labs}\NormalTok{(}\AttributeTok{title =} \StringTok{"Population Distribution: logit(ρ)"}\NormalTok{, }
       \AttributeTok{subtitle =} \StringTok{"Notice how the unpooled estimates are artificially too wide."}\NormalTok{,}
       \AttributeTok{x =} \StringTok{"logit(ρ)"}\NormalTok{, }\AttributeTok{y =} \StringTok{"Density"}\NormalTok{, }\AttributeTok{fill =} \ConstantTok{NULL}\NormalTok{, }\AttributeTok{color =} \ConstantTok{NULL}\NormalTok{, }\AttributeTok{linetype =} \ConstantTok{NULL}\NormalTok{) }\SpecialCharTok{+}
  \FunctionTok{theme}\NormalTok{(}\AttributeTok{legend.position =} \StringTok{"bottom"}\NormalTok{)}

\NormalTok{p\_pop\_kappa }\OtherTok{\textless{}{-}} \FunctionTok{ggplot}\NormalTok{() }\SpecialCharTok{+}
  \FunctionTok{geom\_area}\NormalTok{(}\AttributeTok{data =} \FunctionTok{filter}\NormalTok{(dist\_kappa, Distribution }\SpecialCharTok{==} \StringTok{"Unpooled"}\NormalTok{), }
            \FunctionTok{aes}\NormalTok{(}\AttributeTok{x =}\NormalTok{ x, }\AttributeTok{y =}\NormalTok{ density, }\AttributeTok{fill =} \StringTok{"Unpooled Estimates"}\NormalTok{), }\AttributeTok{alpha =} \FloatTok{0.3}\NormalTok{) }\SpecialCharTok{+}
  \FunctionTok{geom\_line}\NormalTok{(}\AttributeTok{data =} \FunctionTok{filter}\NormalTok{(dist\_kappa, Distribution }\SpecialCharTok{!=} \StringTok{"Unpooled"}\NormalTok{), }
            \FunctionTok{aes}\NormalTok{(}\AttributeTok{x =}\NormalTok{ x, }\AttributeTok{y =}\NormalTok{ density, }\AttributeTok{color =}\NormalTok{ Distribution, }\AttributeTok{linetype =}\NormalTok{ Distribution), }
            \AttributeTok{linewidth =} \FloatTok{1.2}\NormalTok{) }\SpecialCharTok{+}
  \FunctionTok{scale\_color\_manual}\NormalTok{(}\AttributeTok{values =} \FunctionTok{c}\NormalTok{(}\StringTok{"True"} \OtherTok{=} \StringTok{"black"}\NormalTok{, }\StringTok{"Hierarchical"} \OtherTok{=} \StringTok{"\#d73027"}\NormalTok{)) }\SpecialCharTok{+}
  \FunctionTok{scale\_linetype\_manual}\NormalTok{(}\AttributeTok{values =} \FunctionTok{c}\NormalTok{(}\StringTok{"True"} \OtherTok{=} \StringTok{"dashed"}\NormalTok{, }\StringTok{"Hierarchical"} \OtherTok{=} \StringTok{"solid"}\NormalTok{)) }\SpecialCharTok{+}
  \FunctionTok{scale\_fill\_manual}\NormalTok{(}\AttributeTok{values =} \FunctionTok{c}\NormalTok{(}\StringTok{"Unpooled Estimates"} \OtherTok{=} \StringTok{"grey60"}\NormalTok{)) }\SpecialCharTok{+}
  \FunctionTok{labs}\NormalTok{(}\AttributeTok{title =} \StringTok{"Population Distribution: log(κ)"}\NormalTok{, }
       \AttributeTok{subtitle =} \StringTok{"The hierarchical model recovers the true, tighter variance."}\NormalTok{,}
       \AttributeTok{x =} \StringTok{"log(κ)"}\NormalTok{, }\AttributeTok{y =} \StringTok{"Density"}\NormalTok{, }\AttributeTok{fill =} \ConstantTok{NULL}\NormalTok{, }\AttributeTok{color =} \ConstantTok{NULL}\NormalTok{, }\AttributeTok{linetype =} \ConstantTok{NULL}\NormalTok{) }\SpecialCharTok{+}
  \FunctionTok{theme}\NormalTok{(}\AttributeTok{legend.position =} \StringTok{"bottom"}\NormalTok{)}

\CommentTok{\# Forest plot for exact hyperparameter recovery}
\NormalTok{p\_hyper }\OtherTok{\textless{}{-}} \FunctionTok{ggplot}\NormalTok{(plot\_hyper\_data, }\FunctionTok{aes}\NormalTok{(}\AttributeTok{y =}\NormalTok{ clean\_var)) }\SpecialCharTok{+}
  \CommentTok{\# Hierarchical Posterior (Line and Dot)}
  \FunctionTok{geom\_pointrange}\NormalTok{(}\FunctionTok{aes}\NormalTok{(}\AttributeTok{x =}\NormalTok{ mean, }\AttributeTok{xmin =}\NormalTok{ q5, }\AttributeTok{xmax =}\NormalTok{ q95, }\AttributeTok{color =} \StringTok{"Hierarchical"}\NormalTok{), }
                  \AttributeTok{size =} \FloatTok{0.8}\NormalTok{, }\AttributeTok{linewidth =} \DecValTok{1}\NormalTok{) }\SpecialCharTok{+}
  \CommentTok{\# Unpooled Sample Statistic (Diamond)}
  \FunctionTok{geom\_point}\NormalTok{(}\FunctionTok{aes}\NormalTok{(}\AttributeTok{x =}\NormalTok{ unpooled\_stat, }\AttributeTok{color =} \StringTok{"Unpooled"}\NormalTok{), }
             \AttributeTok{shape =} \DecValTok{18}\NormalTok{, }\AttributeTok{size =} \DecValTok{5}\NormalTok{) }\SpecialCharTok{+}
  \CommentTok{\# True Generating Value (Red Cross)}
  \FunctionTok{geom\_point}\NormalTok{(}\FunctionTok{aes}\NormalTok{(}\AttributeTok{x =}\NormalTok{ true\_value, }\AttributeTok{color =} \StringTok{"True"}\NormalTok{), }
             \AttributeTok{shape =} \DecValTok{4}\NormalTok{, }\AttributeTok{size =} \DecValTok{5}\NormalTok{, }\AttributeTok{stroke =} \DecValTok{2}\NormalTok{) }\SpecialCharTok{+}
  \FunctionTok{scale\_color\_manual}\NormalTok{(}\AttributeTok{values =} \FunctionTok{c}\NormalTok{(}\StringTok{"Hierarchical"} \OtherTok{=} \StringTok{"midnightblue"}\NormalTok{, }
                                \StringTok{"Unpooled"} \OtherTok{=} \StringTok{"grey60"}\NormalTok{, }
                                \StringTok{"True"} \OtherTok{=} \StringTok{"red"}\NormalTok{)) }\SpecialCharTok{+}
  \FunctionTok{labs}\NormalTok{(}\AttributeTok{title =} \StringTok{"Hyperparameter Recovery \& Partial Pooling"}\NormalTok{, }
       \AttributeTok{subtitle =} \StringTok{"Comparing Hierarchical Posteriors (Mean ± 90\% CI), Unpooled Sample Stats, and True Values"}\NormalTok{,}
       \AttributeTok{x =} \StringTok{"Parameter Value"}\NormalTok{, }\AttributeTok{y =} \ConstantTok{NULL}\NormalTok{, }\AttributeTok{color =} \ConstantTok{NULL}\NormalTok{) }\SpecialCharTok{+}
  \FunctionTok{theme}\NormalTok{(}\AttributeTok{legend.position =} \StringTok{"bottom"}\NormalTok{)}

\CommentTok{\# 6. Combine with patchwork}
\NormalTok{(p\_pop\_rho }\SpecialCharTok{|}\NormalTok{ p\_pop\_kappa) }\SpecialCharTok{/}\NormalTok{ p\_hyper }\SpecialCharTok{+}
  \FunctionTok{plot\_layout}\NormalTok{(}\AttributeTok{heights =} \FunctionTok{c}\NormalTok{(}\FloatTok{1.5}\NormalTok{, }\DecValTok{1}\NormalTok{)) }\SpecialCharTok{+}
  \FunctionTok{plot\_annotation}\NormalTok{(}
    \AttributeTok{title =} \StringTok{"Population{-}Level Recovery and Variance Shrinkage"}\NormalTok{,}
    \AttributeTok{theme =} \FunctionTok{theme}\NormalTok{(}\AttributeTok{plot.title =} \FunctionTok{element\_text}\NormalTok{(}\AttributeTok{size =} \DecValTok{16}\NormalTok{, }\AttributeTok{face =} \StringTok{"bold"}\NormalTok{))}
\NormalTok{  )}
\end{Highlighting}
\end{Shaded}

\begin{center}
\includegraphics[width=0.8\linewidth,height=\textheight,keepaspectratio]{10-BayesianModels_files/figure-pdf/ch10_pop_recovery_summary-1.pdf}
\end{center}

Variance Shrinkage: The top image demonstrates that the hierarchical
model accurately recovers the true, tighter variance of the population
distributions. In contrast, the unpooled sample statistics are
artificially wide, overestimating the true variance.

\begin{Shaded}
\begin{Highlighting}[]
\CommentTok{\# 1. unpooled\_ests was computed in the previous chunk}

\CommentTok{\# 2. Extract MULTILEVEL (partially pooled) posterior means and join}
\NormalTok{recovery\_ml }\OtherTok{\textless{}{-}}\NormalTok{ agent\_params\_ml }\SpecialCharTok{|\textgreater{}}
  \FunctionTok{mutate}\NormalTok{(}
    \AttributeTok{est\_rho   =} \FunctionTok{map\_dbl}\NormalTok{(agent\_id, }\SpecialCharTok{\textasciitilde{}} \FunctionTok{mean}\NormalTok{(draws\_ml[[}\FunctionTok{paste0}\NormalTok{(}\StringTok{"rho["}\NormalTok{, .x, }\StringTok{"]"}\NormalTok{)]])),}
    \AttributeTok{est\_kappa =} \FunctionTok{map\_dbl}\NormalTok{(agent\_id, }\SpecialCharTok{\textasciitilde{}} \FunctionTok{mean}\NormalTok{(draws\_ml[[}\FunctionTok{paste0}\NormalTok{(}\StringTok{"kappa["}\NormalTok{, .x, }\StringTok{"]"}\NormalTok{)]])),}
    \AttributeTok{est\_wd    =} \FunctionTok{map\_dbl}\NormalTok{(agent\_id, }\SpecialCharTok{\textasciitilde{}} \FunctionTok{mean}\NormalTok{(draws\_ml[[}\FunctionTok{paste0}\NormalTok{(}\StringTok{"weight\_direct["}\NormalTok{, .x, }\StringTok{"]"}\NormalTok{)]])),}
    \AttributeTok{est\_ws    =} \FunctionTok{map\_dbl}\NormalTok{(agent\_id, }\SpecialCharTok{\textasciitilde{}} \FunctionTok{mean}\NormalTok{(draws\_ml[[}\FunctionTok{paste0}\NormalTok{(}\StringTok{"weight\_social["}\NormalTok{, .x, }\StringTok{"]"}\NormalTok{)]]))}
\NormalTok{  ) }\SpecialCharTok{|\textgreater{}}
  \FunctionTok{left\_join}\NormalTok{(unpooled\_ests, }\AttributeTok{by =} \StringTok{"agent\_id"}\NormalTok{)}

\CommentTok{\# 3. Calculate population means (the "center of gravity")}
\NormalTok{pop\_rho\_true   }\OtherTok{\textless{}{-}} \FunctionTok{plogis}\NormalTok{(pop\_logit\_rho\_mean)}
\NormalTok{pop\_kappa\_true }\OtherTok{\textless{}{-}} \FunctionTok{exp}\NormalTok{(pop\_log\_kappa\_mean)}
\NormalTok{pop\_wd\_true    }\OtherTok{\textless{}{-}}\NormalTok{ pop\_rho\_true }\SpecialCharTok{*}\NormalTok{ pop\_kappa\_true}
\NormalTok{pop\_ws\_true    }\OtherTok{\textless{}{-}}\NormalTok{ (}\DecValTok{1} \SpecialCharTok{{-}}\NormalTok{ pop\_rho\_true) }\SpecialCharTok{*}\NormalTok{ pop\_kappa\_true}

\CommentTok{\# 4. Determine a common maximum scale for wd and ws plots}
\NormalTok{max\_w }\OtherTok{\textless{}{-}} \FunctionTok{max}\NormalTok{(}\FunctionTok{c}\NormalTok{(recovery\_ml}\SpecialCharTok{$}\NormalTok{weight\_direct, recovery\_ml}\SpecialCharTok{$}\NormalTok{est\_wd, recovery\_ml}\SpecialCharTok{$}\NormalTok{unpooled\_wd,}
\NormalTok{               recovery\_ml}\SpecialCharTok{$}\NormalTok{weight\_social, recovery\_ml}\SpecialCharTok{$}\NormalTok{est\_ws, recovery\_ml}\SpecialCharTok{$}\NormalTok{unpooled\_ws)) }\SpecialCharTok{*} \FloatTok{1.05}

\CommentTok{\# 5. Build plots: Arrow from Unpooled (hollow) to Pooled (solid)}
\NormalTok{p\_rho }\OtherTok{\textless{}{-}} \FunctionTok{ggplot}\NormalTok{(recovery\_ml, }\FunctionTok{aes}\NormalTok{(}\AttributeTok{x =}\NormalTok{ rho)) }\SpecialCharTok{+}
  \FunctionTok{geom\_hline}\NormalTok{(}\AttributeTok{yintercept =}\NormalTok{ pop\_rho\_true, }\AttributeTok{linetype =} \StringTok{"dotted"}\NormalTok{, }\AttributeTok{colour =} \StringTok{"black"}\NormalTok{, }\AttributeTok{alpha =} \FloatTok{0.5}\NormalTok{) }\SpecialCharTok{+}
  \FunctionTok{geom\_abline}\NormalTok{(}\AttributeTok{linetype =} \StringTok{"dashed"}\NormalTok{, }\AttributeTok{colour =} \StringTok{"grey50"}\NormalTok{) }\SpecialCharTok{+}
  \FunctionTok{geom\_segment}\NormalTok{(}\FunctionTok{aes}\NormalTok{(}\AttributeTok{xend =}\NormalTok{ rho, }\AttributeTok{y =}\NormalTok{ unpooled\_rho, }\AttributeTok{yend =}\NormalTok{ est\_rho), }
               \AttributeTok{arrow =} \FunctionTok{arrow}\NormalTok{(}\AttributeTok{length =} \FunctionTok{unit}\NormalTok{(}\FloatTok{0.15}\NormalTok{, }\StringTok{"cm"}\NormalTok{)), }\AttributeTok{colour =} \StringTok{"grey50"}\NormalTok{) }\SpecialCharTok{+}
  \FunctionTok{geom\_point}\NormalTok{(}\FunctionTok{aes}\NormalTok{(}\AttributeTok{y =}\NormalTok{ unpooled\_rho), }\AttributeTok{shape =} \DecValTok{1}\NormalTok{, }\AttributeTok{size =} \FloatTok{2.5}\NormalTok{, }\AttributeTok{colour =} \StringTok{"grey40"}\NormalTok{) }\SpecialCharTok{+}
  \FunctionTok{geom\_point}\NormalTok{(}\FunctionTok{aes}\NormalTok{(}\AttributeTok{y =}\NormalTok{ est\_rho), }\AttributeTok{size =} \DecValTok{3}\NormalTok{, }\AttributeTok{alpha =} \FloatTok{0.9}\NormalTok{, }\AttributeTok{colour =} \StringTok{"\#4575b4"}\NormalTok{) }\SpecialCharTok{+}
  \FunctionTok{coord\_cartesian}\NormalTok{(}\AttributeTok{xlim =} \FunctionTok{c}\NormalTok{(}\DecValTok{0}\NormalTok{, }\DecValTok{1}\NormalTok{), }\AttributeTok{ylim =} \FunctionTok{c}\NormalTok{(}\DecValTok{0}\NormalTok{, }\DecValTok{1}\NormalTok{)) }\SpecialCharTok{+}
  \FunctionTok{labs}\NormalTok{(}\AttributeTok{title =} \StringTok{"Recovery: ρ (relative allocation)"}\NormalTok{, }\AttributeTok{x =} \StringTok{"True ρ"}\NormalTok{, }\AttributeTok{y =} \StringTok{"Estimated ρ"}\NormalTok{) }

\NormalTok{p\_kappa }\OtherTok{\textless{}{-}} \FunctionTok{ggplot}\NormalTok{(recovery\_ml, }\FunctionTok{aes}\NormalTok{(}\AttributeTok{x =}\NormalTok{ kappa)) }\SpecialCharTok{+}
  \FunctionTok{geom\_hline}\NormalTok{(}\AttributeTok{yintercept =}\NormalTok{ pop\_kappa\_true, }\AttributeTok{linetype =} \StringTok{"dotted"}\NormalTok{, }\AttributeTok{colour =} \StringTok{"black"}\NormalTok{, }\AttributeTok{alpha =} \FloatTok{0.5}\NormalTok{) }\SpecialCharTok{+}
  \FunctionTok{geom\_abline}\NormalTok{(}\AttributeTok{linetype =} \StringTok{"dashed"}\NormalTok{, }\AttributeTok{colour =} \StringTok{"grey50"}\NormalTok{) }\SpecialCharTok{+}
  \FunctionTok{geom\_segment}\NormalTok{(}\FunctionTok{aes}\NormalTok{(}\AttributeTok{xend =}\NormalTok{ kappa, }\AttributeTok{y =}\NormalTok{ unpooled\_kappa, }\AttributeTok{yend =}\NormalTok{ est\_kappa), }
               \AttributeTok{arrow =} \FunctionTok{arrow}\NormalTok{(}\AttributeTok{length =} \FunctionTok{unit}\NormalTok{(}\FloatTok{0.15}\NormalTok{, }\StringTok{"cm"}\NormalTok{)), }\AttributeTok{colour =} \StringTok{"grey50"}\NormalTok{) }\SpecialCharTok{+}
  \FunctionTok{geom\_point}\NormalTok{(}\FunctionTok{aes}\NormalTok{(}\AttributeTok{y =}\NormalTok{ unpooled\_kappa), }\AttributeTok{shape =} \DecValTok{1}\NormalTok{, }\AttributeTok{size =} \FloatTok{2.5}\NormalTok{, }\AttributeTok{colour =} \StringTok{"grey40"}\NormalTok{) }\SpecialCharTok{+}
  \FunctionTok{geom\_point}\NormalTok{(}\FunctionTok{aes}\NormalTok{(}\AttributeTok{y =}\NormalTok{ est\_kappa), }\AttributeTok{size =} \DecValTok{3}\NormalTok{, }\AttributeTok{alpha =} \FloatTok{0.9}\NormalTok{, }\AttributeTok{colour =} \StringTok{"\#d73027"}\NormalTok{) }\SpecialCharTok{+}
  \FunctionTok{labs}\NormalTok{(}\AttributeTok{title =} \StringTok{"Recovery: κ (total evidence scaling)"}\NormalTok{, }\AttributeTok{x =} \StringTok{"True κ"}\NormalTok{, }\AttributeTok{y =} \StringTok{"Estimated κ"}\NormalTok{) }

\NormalTok{p\_wd }\OtherTok{\textless{}{-}} \FunctionTok{ggplot}\NormalTok{(recovery\_ml, }\FunctionTok{aes}\NormalTok{(}\AttributeTok{x =}\NormalTok{ weight\_direct)) }\SpecialCharTok{+}
  \FunctionTok{geom\_hline}\NormalTok{(}\AttributeTok{yintercept =}\NormalTok{ pop\_wd\_true, }\AttributeTok{linetype =} \StringTok{"dotted"}\NormalTok{, }\AttributeTok{colour =} \StringTok{"black"}\NormalTok{, }\AttributeTok{alpha =} \FloatTok{0.5}\NormalTok{) }\SpecialCharTok{+}
  \FunctionTok{geom\_abline}\NormalTok{(}\AttributeTok{linetype =} \StringTok{"dashed"}\NormalTok{, }\AttributeTok{colour =} \StringTok{"grey50"}\NormalTok{) }\SpecialCharTok{+}
  \FunctionTok{geom\_segment}\NormalTok{(}\FunctionTok{aes}\NormalTok{(}\AttributeTok{xend =}\NormalTok{ weight\_direct, }\AttributeTok{y =}\NormalTok{ unpooled\_wd, }\AttributeTok{yend =}\NormalTok{ est\_wd), }
               \AttributeTok{arrow =} \FunctionTok{arrow}\NormalTok{(}\AttributeTok{length =} \FunctionTok{unit}\NormalTok{(}\FloatTok{0.15}\NormalTok{, }\StringTok{"cm"}\NormalTok{)), }\AttributeTok{colour =} \StringTok{"grey50"}\NormalTok{) }\SpecialCharTok{+}
  \FunctionTok{geom\_point}\NormalTok{(}\FunctionTok{aes}\NormalTok{(}\AttributeTok{y =}\NormalTok{ unpooled\_wd), }\AttributeTok{shape =} \DecValTok{1}\NormalTok{, }\AttributeTok{size =} \FloatTok{2.5}\NormalTok{, }\AttributeTok{colour =} \StringTok{"grey40"}\NormalTok{) }\SpecialCharTok{+}
  \FunctionTok{geom\_point}\NormalTok{(}\FunctionTok{aes}\NormalTok{(}\AttributeTok{y =}\NormalTok{ est\_wd), }\AttributeTok{size =} \DecValTok{3}\NormalTok{, }\AttributeTok{alpha =} \FloatTok{0.9}\NormalTok{, }\AttributeTok{colour =} \StringTok{"\#4daf4a"}\NormalTok{) }\SpecialCharTok{+}
  \FunctionTok{coord\_cartesian}\NormalTok{(}\AttributeTok{xlim =} \FunctionTok{c}\NormalTok{(}\DecValTok{0}\NormalTok{, max\_w), }\AttributeTok{ylim =} \FunctionTok{c}\NormalTok{(}\DecValTok{0}\NormalTok{, max\_w)) }\SpecialCharTok{+}
  \FunctionTok{labs}\NormalTok{(}\AttributeTok{title =} \StringTok{"Recovery: w\_d (derived)"}\NormalTok{, }\AttributeTok{x =} \StringTok{"True w\_d"}\NormalTok{, }\AttributeTok{y =} \StringTok{"Estimated w\_d"}\NormalTok{)}

\NormalTok{p\_ws }\OtherTok{\textless{}{-}} \FunctionTok{ggplot}\NormalTok{(recovery\_ml, }\FunctionTok{aes}\NormalTok{(}\AttributeTok{x =}\NormalTok{ weight\_social)) }\SpecialCharTok{+}
  \FunctionTok{geom\_hline}\NormalTok{(}\AttributeTok{yintercept =}\NormalTok{ pop\_ws\_true, }\AttributeTok{linetype =} \StringTok{"dotted"}\NormalTok{, }\AttributeTok{colour =} \StringTok{"black"}\NormalTok{, }\AttributeTok{alpha =} \FloatTok{0.5}\NormalTok{) }\SpecialCharTok{+}
  \FunctionTok{geom\_abline}\NormalTok{(}\AttributeTok{linetype =} \StringTok{"dashed"}\NormalTok{, }\AttributeTok{colour =} \StringTok{"grey50"}\NormalTok{) }\SpecialCharTok{+}
  \FunctionTok{geom\_segment}\NormalTok{(}\FunctionTok{aes}\NormalTok{(}\AttributeTok{xend =}\NormalTok{ weight\_social, }\AttributeTok{y =}\NormalTok{ unpooled\_ws, }\AttributeTok{yend =}\NormalTok{ est\_ws), }
               \AttributeTok{arrow =} \FunctionTok{arrow}\NormalTok{(}\AttributeTok{length =} \FunctionTok{unit}\NormalTok{(}\FloatTok{0.15}\NormalTok{, }\StringTok{"cm"}\NormalTok{)), }\AttributeTok{colour =} \StringTok{"grey50"}\NormalTok{) }\SpecialCharTok{+}
  \FunctionTok{geom\_point}\NormalTok{(}\FunctionTok{aes}\NormalTok{(}\AttributeTok{y =}\NormalTok{ unpooled\_ws), }\AttributeTok{shape =} \DecValTok{1}\NormalTok{, }\AttributeTok{size =} \FloatTok{2.5}\NormalTok{, }\AttributeTok{colour =} \StringTok{"grey40"}\NormalTok{) }\SpecialCharTok{+}
  \FunctionTok{geom\_point}\NormalTok{(}\FunctionTok{aes}\NormalTok{(}\AttributeTok{y =}\NormalTok{ est\_ws), }\AttributeTok{size =} \DecValTok{3}\NormalTok{, }\AttributeTok{alpha =} \FloatTok{0.9}\NormalTok{, }\AttributeTok{colour =} \StringTok{"\#984ea3"}\NormalTok{) }\SpecialCharTok{+}
  \FunctionTok{coord\_cartesian}\NormalTok{(}\AttributeTok{xlim =} \FunctionTok{c}\NormalTok{(}\DecValTok{0}\NormalTok{, max\_w), }\AttributeTok{ylim =} \FunctionTok{c}\NormalTok{(}\DecValTok{0}\NormalTok{, max\_w)) }\SpecialCharTok{+}
  \FunctionTok{labs}\NormalTok{(}\AttributeTok{title =} \StringTok{"Recovery: w\_s (derived)"}\NormalTok{, }\AttributeTok{x =} \StringTok{"True w\_s"}\NormalTok{, }\AttributeTok{y =} \StringTok{"Estimated w\_s"}\NormalTok{)}

\NormalTok{(p\_rho }\SpecialCharTok{|}\NormalTok{ p\_kappa) }\SpecialCharTok{/}\NormalTok{ (p\_wd }\SpecialCharTok{|}\NormalTok{ p\_ws) }\SpecialCharTok{+}
  \FunctionTok{plot\_annotation}\NormalTok{(}
    \AttributeTok{title    =} \StringTok{"Multilevel WBA Parameter Recovery \& Shrinkage"}\NormalTok{,}
    \AttributeTok{subtitle =} \StringTok{"Arrows trace shrinkage from Unpooled MAP estimates (hollow) to Multilevel Posterior Means (solid).}\SpecialCharTok{\textbackslash{}n}\StringTok{Notice how pooling pulls extreme estimates toward the population mean (dotted horizontal lines)."}
\NormalTok{  )}
\end{Highlighting}
\end{Shaded}

\begin{center}
\includegraphics[width=0.8\linewidth,height=\textheight,keepaspectratio]{10-BayesianModels_files/figure-pdf/ch10_multilevel_recovery-1.pdf}
\end{center}

Look for three features. First, \(\rho\) recovery should be tight ---
this parameter is well-identified from conflict trials, just as in the
single-agent case. Second, \(\kappa\) recovery will be noisier --- the
total evidence scale is harder to pin down from binary outcomes. Third,
agents at the extremes should be \emph{pulled toward the centre}
relative to their true values. This is partial pooling in action: the
hierarchical prior regularises extreme estimates. The strength of this
shrinkage depends on how informative each agent's data are relative to
the population distribution.

\section{Simulation-Based Calibration Checks (SBC): Multilevel
Models}\label{simulation-based-calibration-checks-sbc-multilevel-models}

SBC is \emph{more} critical for multilevel models than for single-agent
models. The hierarchical posterior geometry --- with its potential
funnels, correlations, and partial-pooling dynamics --- creates subtle
failure modes that standard MCMC diagnostics may miss. We run SBC for
both the ML-WBA and ML-PBA.

\begin{Shaded}
\begin{Highlighting}[]
\NormalTok{n\_sbc\_ml }\OtherTok{\textless{}{-}} \DecValTok{200}  \CommentTok{\# Reduced from single{-}agent count due to computational cost}

\ControlFlowTok{if}\NormalTok{ (regenerate\_simulations }\SpecialCharTok{||} \SpecialCharTok{!}\FunctionTok{all}\NormalTok{(}\FunctionTok{file.exists}\NormalTok{(}\FunctionTok{c}\NormalTok{(}\StringTok{"simmodels/ch10\_sbc\_ml\_wba.rds"}\NormalTok{, }\StringTok{"simmodels/ch10\_sbc\_ml\_pba.rds"}\NormalTok{)))) \{}
  
  \CommentTok{\# ── ML{-}WBA SBC Generator ──────────────────────────────────────────}
\NormalTok{  sbc\_gen\_ml\_wba }\OtherTok{\textless{}{-}} \FunctionTok{SBC\_generator\_function}\NormalTok{(}
    \ControlFlowTok{function}\NormalTok{(n\_agents\_sbc, ev, n\_reps) \{}
\NormalTok{      J }\OtherTok{\textless{}{-}}\NormalTok{ n\_agents\_sbc}
      
      \CommentTok{\# Draw population parameters from priors}
\NormalTok{      mu\_1     }\OtherTok{\textless{}{-}} \FunctionTok{rnorm}\NormalTok{(}\DecValTok{1}\NormalTok{, }\DecValTok{0}\NormalTok{, }\FloatTok{1.5}\NormalTok{)       }\CommentTok{\# logit\_rho mean}
\NormalTok{      mu\_2     }\OtherTok{\textless{}{-}} \FunctionTok{rnorm}\NormalTok{(}\DecValTok{1}\NormalTok{, }\FunctionTok{log}\NormalTok{(}\DecValTok{2}\NormalTok{), }\FloatTok{0.5}\NormalTok{)   }\CommentTok{\# log\_kappa mean}
\NormalTok{      sigma\_1  }\OtherTok{\textless{}{-}} \FunctionTok{rexp}\NormalTok{(}\DecValTok{1}\NormalTok{, }\DecValTok{2}\NormalTok{)}
\NormalTok{      sigma\_2  }\OtherTok{\textless{}{-}} \FunctionTok{rexp}\NormalTok{(}\DecValTok{1}\NormalTok{, }\DecValTok{2}\NormalTok{)}
      
      \CommentTok{\# LKJ(2) correlation: draw from marginal Beta, rescale}
\NormalTok{      r\_raw    }\OtherTok{\textless{}{-}} \FunctionTok{rbeta}\NormalTok{(}\DecValTok{1}\NormalTok{, }\DecValTok{2}\NormalTok{, }\DecValTok{2}\NormalTok{)}
\NormalTok{      Omega\_12 }\OtherTok{\textless{}{-}} \DecValTok{2} \SpecialCharTok{*}\NormalTok{ r\_raw }\SpecialCharTok{{-}} \DecValTok{1}
      
      \CommentTok{\# Construct Cholesky of correlation matrix}
\NormalTok{      L11 }\OtherTok{\textless{}{-}} \DecValTok{1}\NormalTok{; L21 }\OtherTok{\textless{}{-}}\NormalTok{ Omega\_12; L22 }\OtherTok{\textless{}{-}} \FunctionTok{sqrt}\NormalTok{(}\DecValTok{1} \SpecialCharTok{{-}}\NormalTok{ Omega\_12}\SpecialCharTok{\^{}}\DecValTok{2}\NormalTok{)}
      
      \CommentTok{\# Draw individual parameters (NCP)}
\NormalTok{      z }\OtherTok{\textless{}{-}} \FunctionTok{matrix}\NormalTok{(}\FunctionTok{rnorm}\NormalTok{(}\DecValTok{2} \SpecialCharTok{*}\NormalTok{ J), }\AttributeTok{nrow =} \DecValTok{2}\NormalTok{)}
\NormalTok{      theta\_raw }\OtherTok{\textless{}{-}} \FunctionTok{diag}\NormalTok{(}\FunctionTok{c}\NormalTok{(sigma\_1, sigma\_2)) }\SpecialCharTok{\%*\%} \FunctionTok{rbind}\NormalTok{(}
\NormalTok{        L11 }\SpecialCharTok{*}\NormalTok{ z[}\DecValTok{1}\NormalTok{,],}
\NormalTok{        L21 }\SpecialCharTok{*}\NormalTok{ z[}\DecValTok{1}\NormalTok{,] }\SpecialCharTok{+}\NormalTok{ L22 }\SpecialCharTok{*}\NormalTok{ z[}\DecValTok{2}\NormalTok{,]}
\NormalTok{      )}
      
\NormalTok{      rho\_j   }\OtherTok{\textless{}{-}} \FunctionTok{plogis}\NormalTok{(mu\_1 }\SpecialCharTok{+}\NormalTok{ theta\_raw[}\DecValTok{1}\NormalTok{,])}
\NormalTok{      kappa\_j }\OtherTok{\textless{}{-}} \FunctionTok{exp}\NormalTok{(mu\_2 }\SpecialCharTok{+}\NormalTok{ theta\_raw[}\DecValTok{2}\NormalTok{,])}
\NormalTok{      wd\_j    }\OtherTok{\textless{}{-}}\NormalTok{ rho\_j }\SpecialCharTok{*}\NormalTok{ kappa\_j}
\NormalTok{      ws\_j    }\OtherTok{\textless{}{-}}\NormalTok{ (}\DecValTok{1} \SpecialCharTok{{-}}\NormalTok{ rho\_j) }\SpecialCharTok{*}\NormalTok{ kappa\_j}
      
      \CommentTok{\# Generate trial data for each agent}
\NormalTok{      all\_data }\OtherTok{\textless{}{-}} \FunctionTok{map\_dfr}\NormalTok{(}\FunctionTok{seq\_len}\NormalTok{(J), }\ControlFlowTok{function}\NormalTok{(j) \{}
\NormalTok{        agent\_ev }\OtherTok{\textless{}{-}}\NormalTok{ ev }\SpecialCharTok{|\textgreater{}} \FunctionTok{slice}\NormalTok{(}\FunctionTok{rep}\NormalTok{(}\FunctionTok{seq\_len}\NormalTok{(}\FunctionTok{n}\NormalTok{()), }\AttributeTok{times =}\NormalTok{ n\_reps))}
\NormalTok{        N\_j }\OtherTok{\textless{}{-}} \FunctionTok{nrow}\NormalTok{(agent\_ev)}
        
\NormalTok{        alpha\_p }\OtherTok{\textless{}{-}} \FloatTok{0.5} \SpecialCharTok{+}\NormalTok{ wd\_j[j] }\SpecialCharTok{*}\NormalTok{ agent\_ev}\SpecialCharTok{$}\NormalTok{blue1 }\SpecialCharTok{+}\NormalTok{ ws\_j[j] }\SpecialCharTok{*}\NormalTok{ agent\_ev}\SpecialCharTok{$}\NormalTok{blue2}
\NormalTok{        beta\_p  }\OtherTok{\textless{}{-}} \FloatTok{0.5} \SpecialCharTok{+}\NormalTok{ wd\_j[j] }\SpecialCharTok{*}\NormalTok{ (agent\_ev}\SpecialCharTok{$}\NormalTok{total1 }\SpecialCharTok{{-}}\NormalTok{ agent\_ev}\SpecialCharTok{$}\NormalTok{blue1) }\SpecialCharTok{+}
\NormalTok{                         ws\_j[j] }\SpecialCharTok{*}\NormalTok{ (agent\_ev}\SpecialCharTok{$}\NormalTok{total2 }\SpecialCharTok{{-}}\NormalTok{ agent\_ev}\SpecialCharTok{$}\NormalTok{blue2)}
\NormalTok{        p\_blue  }\OtherTok{\textless{}{-}} \FunctionTok{rbeta}\NormalTok{(N\_j, alpha\_p, beta\_p)}
        
\NormalTok{        agent\_ev }\SpecialCharTok{|\textgreater{}} \FunctionTok{mutate}\NormalTok{(}
          \AttributeTok{agent\_id =}\NormalTok{ j,}
          \AttributeTok{choice   =} \FunctionTok{rbinom}\NormalTok{(N\_j, }\DecValTok{1}\NormalTok{, p\_blue)}
\NormalTok{        )}
\NormalTok{      \})}
      
      \FunctionTok{list}\NormalTok{(}
        \AttributeTok{variables =} \FunctionTok{list}\NormalTok{(}
          \StringTok{\textasciigrave{}}\AttributeTok{mu[1]}\StringTok{\textasciigrave{}}      \OtherTok{=}\NormalTok{ mu\_1,}
          \StringTok{\textasciigrave{}}\AttributeTok{mu[2]}\StringTok{\textasciigrave{}}      \OtherTok{=}\NormalTok{ mu\_2,}
          \StringTok{\textasciigrave{}}\AttributeTok{sigma[1]}\StringTok{\textasciigrave{}}   \OtherTok{=}\NormalTok{ sigma\_1,}
          \StringTok{\textasciigrave{}}\AttributeTok{sigma[2]}\StringTok{\textasciigrave{}}   \OtherTok{=}\NormalTok{ sigma\_2,}
          \StringTok{\textasciigrave{}}\AttributeTok{Omega[1,2]}\StringTok{\textasciigrave{}} \OtherTok{=}\NormalTok{ Omega\_12}
\NormalTok{        ),}
        \AttributeTok{generated =} \FunctionTok{list}\NormalTok{(}
          \AttributeTok{N        =} \FunctionTok{nrow}\NormalTok{(all\_data),}
          \AttributeTok{J        =}\NormalTok{ J,}
          \AttributeTok{agent\_id =}\NormalTok{ all\_data}\SpecialCharTok{$}\NormalTok{agent\_id,}
          \AttributeTok{choice   =}\NormalTok{ all\_data}\SpecialCharTok{$}\NormalTok{choice,}
          \AttributeTok{blue1    =}\NormalTok{ all\_data}\SpecialCharTok{$}\NormalTok{blue1,}
          \AttributeTok{total1   =}\NormalTok{ all\_data}\SpecialCharTok{$}\NormalTok{total1,}
          \AttributeTok{blue2    =}\NormalTok{ all\_data}\SpecialCharTok{$}\NormalTok{blue2,}
          \AttributeTok{total2   =}\NormalTok{ all\_data}\SpecialCharTok{$}\NormalTok{total2}
\NormalTok{        )}
\NormalTok{      )}
\NormalTok{    \},}
    \AttributeTok{n\_agents\_sbc =} \DecValTok{10}\NormalTok{,  }\CommentTok{\# Fewer agents for SBC speed}
    \AttributeTok{ev           =}\NormalTok{ evidence\_combinations,}
    \AttributeTok{n\_reps       =} \DecValTok{2}
\NormalTok{  )}
  
\NormalTok{  sbc\_datasets\_ml\_wba }\OtherTok{\textless{}{-}} \FunctionTok{generate\_datasets}\NormalTok{(sbc\_gen\_ml\_wba, n\_sbc\_ml)}
  
\NormalTok{  sbc\_backend\_ml\_wba }\OtherTok{\textless{}{-}}\NormalTok{ SBC}\SpecialCharTok{::}\FunctionTok{SBC\_backend\_cmdstan\_sample}\NormalTok{(}
\NormalTok{    mod\_ml\_wba,}
    \AttributeTok{chains =} \DecValTok{2}\NormalTok{, }\AttributeTok{iter\_warmup =} \DecValTok{500}\NormalTok{, }\AttributeTok{iter\_sampling =} \DecValTok{500}\NormalTok{,}
    \AttributeTok{parallel\_chains =} \DecValTok{2}\NormalTok{, }\AttributeTok{adapt\_delta =} \FloatTok{0.95}\NormalTok{,}
    \AttributeTok{refresh =} \DecValTok{0}\NormalTok{, }\AttributeTok{show\_messages =} \ConstantTok{FALSE}
\NormalTok{  )}
  
\NormalTok{  sbc\_results\_ml\_wba }\OtherTok{\textless{}{-}} \FunctionTok{compute\_SBC}\NormalTok{(sbc\_datasets\_ml\_wba, sbc\_backend\_ml\_wba,}
                                     \AttributeTok{keep\_fits =} \ConstantTok{FALSE}\NormalTok{)}
  \FunctionTok{saveRDS}\NormalTok{(sbc\_results\_ml\_wba, }\StringTok{"simmodels/ch10\_sbc\_ml\_wba.rds"}\NormalTok{)}

  \CommentTok{\# ── ML{-}PBA SBC Generator ──────────────────────────────────────────}
\NormalTok{  sbc\_gen\_ml\_pba }\OtherTok{\textless{}{-}} \FunctionTok{SBC\_generator\_function}\NormalTok{(}
    \ControlFlowTok{function}\NormalTok{(n\_agents\_sbc, ev, n\_reps) \{}
\NormalTok{      J }\OtherTok{\textless{}{-}}\NormalTok{ n\_agents\_sbc}
      
\NormalTok{      mu\_logit\_p    }\OtherTok{\textless{}{-}} \FunctionTok{rnorm}\NormalTok{(}\DecValTok{1}\NormalTok{, }\DecValTok{0}\NormalTok{, }\FloatTok{1.5}\NormalTok{)}
\NormalTok{      sigma\_logit\_p }\OtherTok{\textless{}{-}} \FunctionTok{rexp}\NormalTok{(}\DecValTok{1}\NormalTok{, }\DecValTok{2}\NormalTok{)}
\NormalTok{      z\_p           }\OtherTok{\textless{}{-}} \FunctionTok{rnorm}\NormalTok{(J)}
\NormalTok{      p\_j           }\OtherTok{\textless{}{-}} \FunctionTok{plogis}\NormalTok{(mu\_logit\_p }\SpecialCharTok{+}\NormalTok{ sigma\_logit\_p }\SpecialCharTok{*}\NormalTok{ z\_p)}
      
\NormalTok{      all\_data }\OtherTok{\textless{}{-}} \FunctionTok{map\_dfr}\NormalTok{(}\FunctionTok{seq\_len}\NormalTok{(J), }\ControlFlowTok{function}\NormalTok{(j) \{}
\NormalTok{        agent\_ev }\OtherTok{\textless{}{-}}\NormalTok{ ev }\SpecialCharTok{|\textgreater{}} \FunctionTok{slice}\NormalTok{(}\FunctionTok{rep}\NormalTok{(}\FunctionTok{seq\_len}\NormalTok{(}\FunctionTok{n}\NormalTok{()), }\AttributeTok{times =}\NormalTok{ n\_reps))}
\NormalTok{        N\_j }\OtherTok{\textless{}{-}} \FunctionTok{nrow}\NormalTok{(agent\_ev)}
        
\NormalTok{        alpha\_p }\OtherTok{\textless{}{-}} \FloatTok{0.5} \SpecialCharTok{+}\NormalTok{ p\_j[j] }\SpecialCharTok{*}\NormalTok{ agent\_ev}\SpecialCharTok{$}\NormalTok{blue1 }\SpecialCharTok{+}\NormalTok{ (}\DecValTok{1} \SpecialCharTok{{-}}\NormalTok{ p\_j[j]) }\SpecialCharTok{*}\NormalTok{ agent\_ev}\SpecialCharTok{$}\NormalTok{blue2}
\NormalTok{        beta\_p  }\OtherTok{\textless{}{-}} \FloatTok{0.5} \SpecialCharTok{+}\NormalTok{ p\_j[j] }\SpecialCharTok{*}\NormalTok{ (agent\_ev}\SpecialCharTok{$}\NormalTok{total1 }\SpecialCharTok{{-}}\NormalTok{ agent\_ev}\SpecialCharTok{$}\NormalTok{blue1) }\SpecialCharTok{+}
\NormalTok{                         (}\DecValTok{1} \SpecialCharTok{{-}}\NormalTok{ p\_j[j]) }\SpecialCharTok{*}\NormalTok{ (agent\_ev}\SpecialCharTok{$}\NormalTok{total2 }\SpecialCharTok{{-}}\NormalTok{ agent\_ev}\SpecialCharTok{$}\NormalTok{blue2)}
\NormalTok{        p\_blue  }\OtherTok{\textless{}{-}} \FunctionTok{rbeta}\NormalTok{(N\_j, alpha\_p, beta\_p)}
        
\NormalTok{        agent\_ev }\SpecialCharTok{|\textgreater{}} \FunctionTok{mutate}\NormalTok{(}
          \AttributeTok{agent\_id =}\NormalTok{ j,}
          \AttributeTok{choice   =} \FunctionTok{rbinom}\NormalTok{(N\_j, }\DecValTok{1}\NormalTok{, p\_blue)}
\NormalTok{        )}
\NormalTok{      \})}
      
      \FunctionTok{list}\NormalTok{(}
        \AttributeTok{variables =} \FunctionTok{list}\NormalTok{(}
          \AttributeTok{mu\_logit\_p    =}\NormalTok{ mu\_logit\_p,}
          \AttributeTok{sigma\_logit\_p =}\NormalTok{ sigma\_logit\_p}
\NormalTok{        ),}
        \AttributeTok{generated =} \FunctionTok{list}\NormalTok{(}
          \AttributeTok{N        =} \FunctionTok{nrow}\NormalTok{(all\_data),}
          \AttributeTok{J        =}\NormalTok{ J,}
          \AttributeTok{agent\_id =}\NormalTok{ all\_data}\SpecialCharTok{$}\NormalTok{agent\_id,}
          \AttributeTok{choice   =}\NormalTok{ all\_data}\SpecialCharTok{$}\NormalTok{choice,}
          \AttributeTok{blue1    =}\NormalTok{ all\_data}\SpecialCharTok{$}\NormalTok{blue1,}
          \AttributeTok{total1   =}\NormalTok{ all\_data}\SpecialCharTok{$}\NormalTok{total1,}
          \AttributeTok{blue2    =}\NormalTok{ all\_data}\SpecialCharTok{$}\NormalTok{blue2,}
          \AttributeTok{total2   =}\NormalTok{ all\_data}\SpecialCharTok{$}\NormalTok{total2}
\NormalTok{        )}
\NormalTok{      )}
\NormalTok{    \},}
    \AttributeTok{n\_agents\_sbc =} \DecValTok{10}\NormalTok{,}
    \AttributeTok{ev           =}\NormalTok{ evidence\_combinations,}
    \AttributeTok{n\_reps       =} \DecValTok{2}
\NormalTok{  )}
  
\NormalTok{  sbc\_datasets\_ml\_pba }\OtherTok{\textless{}{-}} \FunctionTok{generate\_datasets}\NormalTok{(sbc\_gen\_ml\_pba, n\_sbc\_ml)}
  
\NormalTok{  sbc\_backend\_ml\_pba }\OtherTok{\textless{}{-}}\NormalTok{ SBC}\SpecialCharTok{::}\FunctionTok{SBC\_backend\_cmdstan\_sample}\NormalTok{(}
\NormalTok{    mod\_ml\_pba,}
    \AttributeTok{chains =} \DecValTok{2}\NormalTok{, }\AttributeTok{iter\_warmup =} \DecValTok{500}\NormalTok{, }\AttributeTok{iter\_sampling =} \DecValTok{500}\NormalTok{,}
    \AttributeTok{parallel\_chains =} \DecValTok{2}\NormalTok{, }\AttributeTok{adapt\_delta =} \FloatTok{0.95}\NormalTok{,}
    \AttributeTok{refresh =} \DecValTok{0}\NormalTok{, }\AttributeTok{show\_messages =} \ConstantTok{FALSE}
\NormalTok{  )}
  
\NormalTok{  sbc\_results\_ml\_pba }\OtherTok{\textless{}{-}} \FunctionTok{compute\_SBC}\NormalTok{(sbc\_datasets\_ml\_pba, sbc\_backend\_ml\_pba,}
                                     \AttributeTok{keep\_fits =} \ConstantTok{FALSE}\NormalTok{)}
  \FunctionTok{saveRDS}\NormalTok{(sbc\_results\_ml\_pba, }\StringTok{"simmodels/ch10\_sbc\_ml\_pba.rds"}\NormalTok{)}
  
\NormalTok{\} }\ControlFlowTok{else}\NormalTok{ \{}
\NormalTok{  sbc\_results\_ml\_wba }\OtherTok{\textless{}{-}} \FunctionTok{readRDS}\NormalTok{(}\StringTok{"simmodels/ch10\_sbc\_ml\_wba.rds"}\NormalTok{)}
\NormalTok{  sbc\_results\_ml\_pba }\OtherTok{\textless{}{-}} \FunctionTok{readRDS}\NormalTok{(}\StringTok{"simmodels/ch10\_sbc\_ml\_pba.rds"}\NormalTok{)}
\NormalTok{\}}

\CommentTok{\# Plot SBC calibration}
\NormalTok{p\_sbc\_ml\_wba }\OtherTok{\textless{}{-}}\NormalTok{ SBC}\SpecialCharTok{::}\FunctionTok{plot\_ecdf\_diff}\NormalTok{(}
\NormalTok{  sbc\_results\_ml\_wba,}
  \AttributeTok{variables =} \FunctionTok{c}\NormalTok{(}\StringTok{"mu[1]"}\NormalTok{, }\StringTok{"mu[2]"}\NormalTok{, }\StringTok{"sigma[1]"}\NormalTok{, }\StringTok{"sigma[2]"}\NormalTok{, }\StringTok{"Omega[1,2]"}\NormalTok{)}
\NormalTok{) }\SpecialCharTok{+} \FunctionTok{labs}\NormalTok{(}\AttributeTok{title =} \StringTok{"SBC Calibration: ML{-}WBA (population parameters)"}\NormalTok{,}
         \AttributeTok{subtitle =} \FunctionTok{paste0}\NormalTok{(}\StringTok{"N = "}\NormalTok{, n\_sbc\_ml, }\StringTok{" simulations. Lines inside bands = calibrated."}\NormalTok{))}

\NormalTok{p\_sbc\_ml\_pba }\OtherTok{\textless{}{-}}\NormalTok{ SBC}\SpecialCharTok{::}\FunctionTok{plot\_ecdf\_diff}\NormalTok{(}
\NormalTok{  sbc\_results\_ml\_pba,}
  \AttributeTok{variables =} \FunctionTok{c}\NormalTok{(}\StringTok{"mu\_logit\_p"}\NormalTok{, }\StringTok{"sigma\_logit\_p"}\NormalTok{)}
\NormalTok{) }\SpecialCharTok{+} \FunctionTok{labs}\NormalTok{(}\AttributeTok{title =} \StringTok{"SBC Calibration: ML{-}PBA (population parameters)"}\NormalTok{,}
         \AttributeTok{subtitle =} \FunctionTok{paste0}\NormalTok{(}\StringTok{"N = "}\NormalTok{, n\_sbc\_ml, }\StringTok{" simulations."}\NormalTok{))}

\NormalTok{p\_sbc\_ml\_wba }\SpecialCharTok{/}\NormalTok{ p\_sbc\_ml\_pba}
\end{Highlighting}
\end{Shaded}

\begin{center}
\includegraphics[width=0.8\linewidth,height=\textheight,keepaspectratio]{10-BayesianModels_files/figure-pdf/ch10_ml_sbc-1.pdf}
\end{center}

If any parameter's ECDF difference line escapes the grey bands, the
posterior for that parameter is systematically miscalibrated. Common
failure modes for multilevel models: (1) \(\sigma\) dipping below the
band near 0, indicating the hierarchical prior's boundary repulsion is
too weak; (2) the correlation \(\Omega_{1,2}\) showing asymmetric
miscalibration, indicating the LKJ prior interacts poorly with the NCP
geometry. If SBC passes, we have a global certificate that the Stan
sampler is unbiased across the entire joint prior space.

\subsection{Global Parameter Recovery from
SBC}\label{global-parameter-recovery-from-sbc-1}

\begin{Shaded}
\begin{Highlighting}[]
\NormalTok{ml\_wba\_sbc\_recovery }\OtherTok{\textless{}{-}}\NormalTok{ sbc\_results\_ml\_wba}\SpecialCharTok{$}\NormalTok{stats }\SpecialCharTok{|\textgreater{}}
  \FunctionTok{filter}\NormalTok{(variable }\SpecialCharTok{\%in\%} \FunctionTok{c}\NormalTok{(}\StringTok{"mu[1]"}\NormalTok{, }\StringTok{"mu[2]"}\NormalTok{, }\StringTok{"sigma[1]"}\NormalTok{, }\StringTok{"sigma[2]"}\NormalTok{, }\StringTok{"Omega[1,2]"}\NormalTok{)) }\SpecialCharTok{|\textgreater{}}
  \FunctionTok{mutate}\NormalTok{(}\AttributeTok{model =} \StringTok{"ML{-}WBA"}\NormalTok{)}

\NormalTok{ml\_pba\_sbc\_recovery }\OtherTok{\textless{}{-}}\NormalTok{ sbc\_results\_ml\_pba}\SpecialCharTok{$}\NormalTok{stats }\SpecialCharTok{|\textgreater{}}
  \FunctionTok{filter}\NormalTok{(variable }\SpecialCharTok{\%in\%} \FunctionTok{c}\NormalTok{(}\StringTok{"mu\_logit\_p"}\NormalTok{, }\StringTok{"sigma\_logit\_p"}\NormalTok{)) }\SpecialCharTok{|\textgreater{}}
  \FunctionTok{mutate}\NormalTok{(}\AttributeTok{model =} \StringTok{"ML{-}PBA"}\NormalTok{)}

\NormalTok{sbc\_ml\_recovery\_df }\OtherTok{\textless{}{-}} \FunctionTok{bind\_rows}\NormalTok{(ml\_wba\_sbc\_recovery, ml\_pba\_sbc\_recovery)}

\FunctionTok{ggplot}\NormalTok{(sbc\_ml\_recovery\_df, }\FunctionTok{aes}\NormalTok{(}\AttributeTok{x =}\NormalTok{ simulated\_value, }\AttributeTok{y =}\NormalTok{ median)) }\SpecialCharTok{+}
  \FunctionTok{geom\_abline}\NormalTok{(}\AttributeTok{slope =} \DecValTok{1}\NormalTok{, }\AttributeTok{intercept =} \DecValTok{0}\NormalTok{, }\AttributeTok{linetype =} \StringTok{"dashed"}\NormalTok{, }\AttributeTok{color =} \StringTok{"grey50"}\NormalTok{) }\SpecialCharTok{+}
  \FunctionTok{geom\_errorbar}\NormalTok{(}\FunctionTok{aes}\NormalTok{(}\AttributeTok{ymin =}\NormalTok{ q5, }\AttributeTok{ymax =}\NormalTok{ q95), }\AttributeTok{alpha =} \FloatTok{0.2}\NormalTok{, }\AttributeTok{width =} \DecValTok{0}\NormalTok{) }\SpecialCharTok{+}
  \FunctionTok{geom\_point}\NormalTok{(}\AttributeTok{alpha =} \FloatTok{0.4}\NormalTok{, }\AttributeTok{size =} \FloatTok{1.5}\NormalTok{, }\AttributeTok{color =} \StringTok{"\#4575b4"}\NormalTok{) }\SpecialCharTok{+}
  \FunctionTok{facet\_wrap}\NormalTok{(model }\SpecialCharTok{\textasciitilde{}}\NormalTok{ variable, }\AttributeTok{scales =} \StringTok{"free"}\NormalTok{, }\AttributeTok{ncol =} \DecValTok{4}\NormalTok{) }\SpecialCharTok{+}
  \FunctionTok{labs}\NormalTok{(}
    \AttributeTok{title =} \StringTok{"Global Parameter Recovery: Multilevel SBC"}\NormalTok{,}
    \AttributeTok{subtitle =} \StringTok{"Each point = one simulated population. Error bars = 90\% CI."}\NormalTok{,}
    \AttributeTok{x =} \StringTok{"True (simulated) value"}\NormalTok{,}
    \AttributeTok{y =} \StringTok{"Posterior median"}
\NormalTok{  ) }\SpecialCharTok{+}
  \FunctionTok{theme}\NormalTok{(}\AttributeTok{strip.background =} \FunctionTok{element\_rect}\NormalTok{(}\AttributeTok{fill =} \StringTok{"gray95"}\NormalTok{, }\AttributeTok{color =} \ConstantTok{NA}\NormalTok{))}
\end{Highlighting}
\end{Shaded}

\begin{center}
\includegraphics[width=0.8\linewidth,height=\textheight,keepaspectratio]{10-BayesianModels_files/figure-pdf/ch10_ml_sbc_recovery-1.pdf}
\end{center}

\section{Multilevel Model Comparison: Three-Way
LOO}\label{multilevel-model-comparison-three-way-loo}

Following the procedure from the single-agent comparison, we compare all
three multilevel models via PSIS-LOO. Because the simulated population
has agents with varying weights, we expect the ML-WBA to outperform the
ML-SBA. Whether the ML-PBA or ML-WBA wins depends on whether the
population s \(\kappa\) variation provides sufficient predictive benefit
to offset the additional complexity.

\begin{Shaded}
\begin{Highlighting}[]
\NormalTok{loo\_ml\_sba }\OtherTok{\textless{}{-}} \FunctionTok{loo}\NormalTok{(fit\_ml\_sba}\SpecialCharTok{$}\FunctionTok{draws}\NormalTok{(}\StringTok{"log\_lik"}\NormalTok{, }\AttributeTok{format =} \StringTok{"matrix"}\NormalTok{))}
\NormalTok{loo\_ml\_pba }\OtherTok{\textless{}{-}} \FunctionTok{loo}\NormalTok{(fit\_ml\_pba}\SpecialCharTok{$}\FunctionTok{draws}\NormalTok{(}\StringTok{"log\_lik"}\NormalTok{, }\AttributeTok{format =} \StringTok{"matrix"}\NormalTok{))}
\NormalTok{loo\_ml\_wba }\OtherTok{\textless{}{-}} \FunctionTok{loo}\NormalTok{(fit\_ml\_wba}\SpecialCharTok{$}\FunctionTok{draws}\NormalTok{(}\StringTok{"log\_lik"}\NormalTok{, }\AttributeTok{format =} \StringTok{"matrix"}\NormalTok{))}

\FunctionTok{cat}\NormalTok{(}\StringTok{"=== Multilevel Model Comparison: ML{-}SBA vs. ML{-}PBA vs. ML{-}WBA ===}\SpecialCharTok{\textbackslash{}n}\StringTok{"}\NormalTok{)}
\end{Highlighting}
\end{Shaded}

\begin{verbatim}
=== Multilevel Model Comparison: ML-SBA vs. ML-PBA vs. ML-WBA ===
\end{verbatim}

\begin{Shaded}
\begin{Highlighting}[]
\NormalTok{ml\_comparison }\OtherTok{\textless{}{-}} \FunctionTok{loo\_compare}\NormalTok{(}
  \FunctionTok{list}\NormalTok{(}\StringTok{"ML{-}SBA"} \OtherTok{=}\NormalTok{ loo\_ml\_sba, }\StringTok{"ML{-}PBA"} \OtherTok{=}\NormalTok{ loo\_ml\_pba, }\StringTok{"ML{-}WBA"} \OtherTok{=}\NormalTok{ loo\_ml\_wba)}
\NormalTok{)}
\FunctionTok{print}\NormalTok{(ml\_comparison)}
\end{Highlighting}
\end{Shaded}

\begin{verbatim}
       elpd_diff se_diff
ML-SBA  0.0       0.0   
ML-WBA -1.9       1.0   
ML-PBA -2.7       2.0   
\end{verbatim}

\begin{Shaded}
\begin{Highlighting}[]
\NormalTok{ml\_weights }\OtherTok{\textless{}{-}} \FunctionTok{loo\_model\_weights}\NormalTok{(}
  \FunctionTok{list}\NormalTok{(}\StringTok{"ML{-}SBA"} \OtherTok{=}\NormalTok{ loo\_ml\_sba, }\StringTok{"ML{-}PBA"} \OtherTok{=}\NormalTok{ loo\_ml\_pba, }\StringTok{"ML{-}WBA"} \OtherTok{=}\NormalTok{ loo\_ml\_wba),}
  \AttributeTok{method =} \StringTok{"stacking"}
\NormalTok{)}
\FunctionTok{cat}\NormalTok{(}\StringTok{"}\SpecialCharTok{\textbackslash{}n}\StringTok{Bayesian stacking weights:}\SpecialCharTok{\textbackslash{}n}\StringTok{"}\NormalTok{)}
\end{Highlighting}
\end{Shaded}

\begin{verbatim}

Bayesian stacking weights:
\end{verbatim}

\begin{Shaded}
\begin{Highlighting}[]
\FunctionTok{print}\NormalTok{(ml\_weights)}
\end{Highlighting}
\end{Shaded}

\begin{verbatim}
Method: stacking
------
       weight
ML-SBA 0.987 
ML-PBA 0.013 
ML-WBA 0.000 
\end{verbatim}

Because the simulated agents were drawn from a population with varying
\(\rho\) and \(\kappa\) (not all equal to 1), the ML-WBA should
outperform the ML-SBA. The ELPD difference quantifies the predictive
cost of assuming all agents are perfect Bayesians when they are not.
However, given we set the simulation to be close to the SBA (balanced
weighting), the ML-SBA ends up winning.

Note that this is \emph{trial-level LOO}: it measures generalisation to
new \emph{trials} from known agents. For the clinically relevant
question --- \emph{can we predict a new participant's behaviour from the
population distribution alone?} --- we would need agent-level K-fold
cross-validation, as discussed in Chapter 8.

\subsection{Visualising Pointwise ELPD
Differences}\label{visualising-pointwise-elpd-differences}

To understand \emph{where} the ML-WBA gains its advantage over the
ML-SBA, we plot pointwise ELPD differences as a function of evidence
conflict --- the same diagnostic technique from the single-agent
section.

\begin{Shaded}
\begin{Highlighting}[]
\NormalTok{pw\_ml }\OtherTok{\textless{}{-}} \FunctionTok{tibble}\NormalTok{(}
  \AttributeTok{obs       =} \FunctionTok{seq\_len}\NormalTok{(stan\_data\_ml}\SpecialCharTok{$}\NormalTok{N),}
  \AttributeTok{agent\_id  =}\NormalTok{ stan\_data\_ml}\SpecialCharTok{$}\NormalTok{agent\_id,}
  \AttributeTok{blue1     =}\NormalTok{ stan\_data\_ml}\SpecialCharTok{$}\NormalTok{blue1,}
  \AttributeTok{blue2     =}\NormalTok{ stan\_data\_ml}\SpecialCharTok{$}\NormalTok{blue2,}
  \AttributeTok{elpd\_wba  =}\NormalTok{ loo\_ml\_wba}\SpecialCharTok{$}\NormalTok{pointwise[, }\StringTok{"elpd\_loo"}\NormalTok{],}
  \AttributeTok{elpd\_sba  =}\NormalTok{ loo\_ml\_sba}\SpecialCharTok{$}\NormalTok{pointwise[, }\StringTok{"elpd\_loo"}\NormalTok{],}
  \AttributeTok{elpd\_pba  =}\NormalTok{ loo\_ml\_pba}\SpecialCharTok{$}\NormalTok{pointwise[, }\StringTok{"elpd\_loo"}\NormalTok{],}
  \AttributeTok{elpd\_diff\_wba\_sba =}\NormalTok{ elpd\_wba }\SpecialCharTok{{-}}\NormalTok{ elpd\_sba,}
  \AttributeTok{elpd\_diff\_wba\_pba =}\NormalTok{ elpd\_wba }\SpecialCharTok{{-}}\NormalTok{ elpd\_pba}
\NormalTok{) }\SpecialCharTok{|\textgreater{}}
  \FunctionTok{left\_join}\NormalTok{(}
\NormalTok{    agent\_params\_ml }\SpecialCharTok{|\textgreater{}} \FunctionTok{select}\NormalTok{(agent\_id, rho, kappa),}
    \AttributeTok{by =} \StringTok{"agent\_id"}
\NormalTok{  ) }\SpecialCharTok{|\textgreater{}}
  \FunctionTok{mutate}\NormalTok{(}
    \AttributeTok{prop\_direct       =}\NormalTok{ blue1 }\SpecialCharTok{/} \DecValTok{8}\NormalTok{,}
    \AttributeTok{prop\_social       =}\NormalTok{ blue2 }\SpecialCharTok{/} \DecValTok{3}\NormalTok{,}
    \AttributeTok{evidence\_conflict =} \FunctionTok{abs}\NormalTok{(prop\_direct }\SpecialCharTok{{-}}\NormalTok{ prop\_social),}
    \AttributeTok{rho\_group =} \FunctionTok{cut}\NormalTok{(rho, }\AttributeTok{breaks =} \FunctionTok{c}\NormalTok{(}\DecValTok{0}\NormalTok{, }\FloatTok{0.35}\NormalTok{, }\FloatTok{0.65}\NormalTok{, }\DecValTok{1}\NormalTok{),}
                    \AttributeTok{labels =} \FunctionTok{c}\NormalTok{(}\StringTok{"Social{-}leaning"}\NormalTok{, }\StringTok{"Balanced"}\NormalTok{, }\StringTok{"Direct{-}leaning"}\NormalTok{))}
\NormalTok{  )}

\CommentTok{\# WBA vs SBA, coloured by agent type}
\FunctionTok{ggplot}\NormalTok{(pw\_ml, }\FunctionTok{aes}\NormalTok{(}\AttributeTok{x =}\NormalTok{ evidence\_conflict, }\AttributeTok{y =}\NormalTok{ elpd\_diff\_wba\_sba)) }\SpecialCharTok{+}
  \FunctionTok{geom\_hline}\NormalTok{(}\AttributeTok{yintercept =} \DecValTok{0}\NormalTok{, }\AttributeTok{linetype =} \StringTok{"dashed"}\NormalTok{, }\AttributeTok{colour =} \StringTok{"grey40"}\NormalTok{) }\SpecialCharTok{+}
  \FunctionTok{geom\_point}\NormalTok{(}\FunctionTok{aes}\NormalTok{(}\AttributeTok{colour =}\NormalTok{ rho\_group), }\AttributeTok{size =} \FloatTok{1.5}\NormalTok{, }\AttributeTok{alpha =} \FloatTok{0.4}\NormalTok{) }\SpecialCharTok{+}
  \FunctionTok{geom\_smooth}\NormalTok{(}\AttributeTok{method =} \StringTok{"gam"}\NormalTok{, }\AttributeTok{formula =}\NormalTok{ y }\SpecialCharTok{\textasciitilde{}} \FunctionTok{s}\NormalTok{(x, }\AttributeTok{bs =} \StringTok{"cs"}\NormalTok{),}
              \AttributeTok{colour =} \StringTok{"midnightblue"}\NormalTok{, }\AttributeTok{linewidth =} \DecValTok{1}\NormalTok{) }\SpecialCharTok{+}
  \FunctionTok{facet\_wrap}\NormalTok{(}\SpecialCharTok{\textasciitilde{}}\NormalTok{ rho\_group) }\SpecialCharTok{+}
  \FunctionTok{scale\_colour\_manual}\NormalTok{(}
    \AttributeTok{values =} \FunctionTok{c}\NormalTok{(}\StringTok{"Social{-}leaning"} \OtherTok{=} \StringTok{"\#377eb8"}\NormalTok{,}
               \StringTok{"Balanced"}       \OtherTok{=} \StringTok{"\#4daf4a"}\NormalTok{,}
               \StringTok{"Direct{-}leaning"} \OtherTok{=} \StringTok{"\#e41a1c"}\NormalTok{)}
\NormalTok{  ) }\SpecialCharTok{+}
  \FunctionTok{labs}\NormalTok{(}
    \AttributeTok{title =} \StringTok{"Pointwise ELPD: ML{-}WBA vs. ML{-}SBA"}\NormalTok{,}
    \AttributeTok{subtitle =} \StringTok{"Faceted by agent type. Advantage grows with evidence conflict for non{-}balanced agents."}\NormalTok{,}
    \AttributeTok{x =} \StringTok{"Evidence Conflict |prop\_direct − prop\_social|"}\NormalTok{,}
    \AttributeTok{y =} \StringTok{"ELPD Difference (ML{-}WBA − ML{-}SBA)"}\NormalTok{,}
    \AttributeTok{colour =} \StringTok{"Agent type"}
\NormalTok{  ) }\SpecialCharTok{+}
  \FunctionTok{theme}\NormalTok{(}\AttributeTok{legend.position =} \StringTok{"bottom"}\NormalTok{)}
\end{Highlighting}
\end{Shaded}

\begin{center}
\includegraphics[width=0.8\linewidth,height=\textheight,keepaspectratio]{10-BayesianModels_files/figure-pdf/ch10_ml_pointwise_elpd-1.pdf}
\end{center}

\section{Model Recovery: Multilevel
Level}\label{model-recovery-multilevel-level}

Model recovery at the multilevel level asks: \emph{if data are generated
by a specific multilevel model, can our comparison pipeline correctly
identify it?} We test three scenarios, paralleling the single-agent
analysis.

\subsection{Specific Scenarios}\label{specific-scenarios}

Before the full simulation, we examine three instructive scenarios:

\begin{enumerate}
\def\labelenumi{\arabic{enumi}.}
\item
  \textbf{SBA-Mimic Population:} All agents have
  \(\rho = 0.5, \kappa = 2.0\) (i.e., \(w_d = w_s = 1\)) with zero
  individual variation. The population \emph{is} the SBA.
\item
  \textbf{PBA Population:} Agents vary in \(\rho\) (drawn from a
  population) but all have \(\kappa = 1.0\). The total evidence budget
  is fixed; only allocation varies.
\item
  \textbf{Full WBA Population:} Agents vary in both \(\rho\) and
  \(\kappa\). The full hierarchical model is needed.
\end{enumerate}

\begin{Shaded}
\begin{Highlighting}[]
\CommentTok{\# Helper: generate multilevel data from a parameter specification}
\NormalTok{generate\_ml\_data }\OtherTok{\textless{}{-}} \ControlFlowTok{function}\NormalTok{(agent\_rhos, agent\_kappas, ev, }\AttributeTok{n\_reps =} \DecValTok{2}\NormalTok{) \{}
\NormalTok{  J }\OtherTok{\textless{}{-}} \FunctionTok{length}\NormalTok{(agent\_rhos)}
\NormalTok{  all\_data }\OtherTok{\textless{}{-}} \FunctionTok{map\_dfr}\NormalTok{(}\FunctionTok{seq\_len}\NormalTok{(J), }\ControlFlowTok{function}\NormalTok{(j) \{}
\NormalTok{    wd }\OtherTok{\textless{}{-}}\NormalTok{ agent\_rhos[j] }\SpecialCharTok{*}\NormalTok{ agent\_kappas[j]}
\NormalTok{    ws }\OtherTok{\textless{}{-}}\NormalTok{ (}\DecValTok{1} \SpecialCharTok{{-}}\NormalTok{ agent\_rhos[j]) }\SpecialCharTok{*}\NormalTok{ agent\_kappas[j]}
    \FunctionTok{generate\_agent\_decisions}\NormalTok{(wd, ws, ev, }\AttributeTok{n\_samples =}\NormalTok{ n\_reps) }\SpecialCharTok{|\textgreater{}}
      \FunctionTok{mutate}\NormalTok{(}\AttributeTok{agent\_id =}\NormalTok{ j, }\AttributeTok{total1 =} \DecValTok{8}\NormalTok{L, }\AttributeTok{total2 =} \DecValTok{3}\NormalTok{L)}
\NormalTok{  \})}
  \FunctionTok{list}\NormalTok{(}
    \AttributeTok{N =} \FunctionTok{nrow}\NormalTok{(all\_data), }\AttributeTok{J =}\NormalTok{ J,}
    \AttributeTok{agent\_id   =}\NormalTok{ all\_data}\SpecialCharTok{$}\NormalTok{agent\_id,}
    \AttributeTok{choice     =}\NormalTok{ all\_data}\SpecialCharTok{$}\NormalTok{choice,}
    \AttributeTok{N\_total    =} \FunctionTok{nrow}\NormalTok{(all\_data),}
    \AttributeTok{N\_subjects =}\NormalTok{ J,}
    \AttributeTok{subj\_id    =}\NormalTok{ all\_data}\SpecialCharTok{$}\NormalTok{agent\_id,}
    \AttributeTok{y          =}\NormalTok{ all\_data}\SpecialCharTok{$}\NormalTok{choice,}
    \AttributeTok{blue1      =}\NormalTok{ all\_data}\SpecialCharTok{$}\NormalTok{blue1,}
    \AttributeTok{total1     =}\NormalTok{ all\_data}\SpecialCharTok{$}\NormalTok{total1,}
    \AttributeTok{blue2      =}\NormalTok{ all\_data}\SpecialCharTok{$}\NormalTok{blue2,}
    \AttributeTok{total2     =}\NormalTok{ all\_data}\SpecialCharTok{$}\NormalTok{total2}
\NormalTok{  )}
\NormalTok{\}}

\NormalTok{J\_scen }\OtherTok{\textless{}{-}} \DecValTok{15}

\CommentTok{\# Scenario 1: SBA{-}Mimic (no individual variation)}
\NormalTok{scen1\_data }\OtherTok{\textless{}{-}} \FunctionTok{generate\_ml\_data}\NormalTok{(}
  \AttributeTok{agent\_rhos   =} \FunctionTok{rep}\NormalTok{(}\FloatTok{0.5}\NormalTok{, J\_scen),}
  \AttributeTok{agent\_kappas =} \FunctionTok{rep}\NormalTok{(}\FloatTok{2.0}\NormalTok{, J\_scen),}
  \AttributeTok{ev =}\NormalTok{ evidence\_combinations}
\NormalTok{)}

\CommentTok{\# Scenario 2: PBA Population (rho varies, kappa = 1)}
\NormalTok{scen2\_rhos }\OtherTok{\textless{}{-}} \FunctionTok{plogis}\NormalTok{(}\FunctionTok{rnorm}\NormalTok{(J\_scen, }\FloatTok{0.5}\NormalTok{, }\FloatTok{0.5}\NormalTok{))}
\NormalTok{scen2\_data }\OtherTok{\textless{}{-}} \FunctionTok{generate\_ml\_data}\NormalTok{(}
  \AttributeTok{agent\_rhos   =}\NormalTok{ scen2\_rhos,}
  \AttributeTok{agent\_kappas =} \FunctionTok{rep}\NormalTok{(}\FloatTok{1.0}\NormalTok{, J\_scen),}
  \AttributeTok{ev =}\NormalTok{ evidence\_combinations}
\NormalTok{)}

\CommentTok{\# Scenario 3: Full WBA population (both vary)}
\NormalTok{scen3\_draws }\OtherTok{\textless{}{-}}\NormalTok{ MASS}\SpecialCharTok{::}\FunctionTok{mvrnorm}\NormalTok{(J\_scen, }\AttributeTok{mu =} \FunctionTok{c}\NormalTok{(}\FloatTok{0.3}\NormalTok{, }\FunctionTok{log}\NormalTok{(}\FloatTok{2.5}\NormalTok{)),}
                              \AttributeTok{Sigma =} \FunctionTok{matrix}\NormalTok{(}\FunctionTok{c}\NormalTok{(}\FloatTok{0.5}\NormalTok{, }\SpecialCharTok{{-}}\FloatTok{0.15}\NormalTok{, }\SpecialCharTok{{-}}\FloatTok{0.15}\NormalTok{, }\FloatTok{0.2}\NormalTok{), }\DecValTok{2}\NormalTok{))}
\NormalTok{scen3\_data }\OtherTok{\textless{}{-}} \FunctionTok{generate\_ml\_data}\NormalTok{(}
  \AttributeTok{agent\_rhos   =} \FunctionTok{plogis}\NormalTok{(scen3\_draws[, }\DecValTok{1}\NormalTok{]),}
  \AttributeTok{agent\_kappas =} \FunctionTok{exp}\NormalTok{(scen3\_draws[, }\DecValTok{2}\NormalTok{]),}
  \AttributeTok{ev =}\NormalTok{ evidence\_combinations}
\NormalTok{)}
\end{Highlighting}
\end{Shaded}

\begin{Shaded}
\begin{Highlighting}[]
\NormalTok{fit\_ml\_scenario }\OtherTok{\textless{}{-}} \ControlFlowTok{function}\NormalTok{(stan\_data, label) \{}
  \FunctionTok{cat}\NormalTok{(}\FunctionTok{sprintf}\NormalTok{(}\StringTok{"}\SpecialCharTok{\textbackslash{}n}\StringTok{{-}{-}{-} Fitting scenario: \%s {-}{-}{-}}\SpecialCharTok{\textbackslash{}n}\StringTok{"}\NormalTok{, label))}

\NormalTok{  sba }\OtherTok{\textless{}{-}}\NormalTok{ mod\_ml\_sba}\SpecialCharTok{$}\FunctionTok{sample}\NormalTok{(}
    \AttributeTok{data =}\NormalTok{ stan\_data, }\AttributeTok{chains =} \DecValTok{2}\NormalTok{, }\AttributeTok{iter\_warmup =} \DecValTok{1000}\NormalTok{,}
    \AttributeTok{iter\_sampling =} \DecValTok{1000}\NormalTok{, }\AttributeTok{adapt\_delta =} \FloatTok{0.9}\NormalTok{,}
    \AttributeTok{refresh =} \DecValTok{0}\NormalTok{, }\AttributeTok{show\_messages =} \ConstantTok{FALSE}
\NormalTok{  )}
\NormalTok{  pba }\OtherTok{\textless{}{-}}\NormalTok{ mod\_ml\_pba}\SpecialCharTok{$}\FunctionTok{sample}\NormalTok{(}
    \AttributeTok{data =}\NormalTok{ stan\_data, }\AttributeTok{chains =} \DecValTok{2}\NormalTok{, }\AttributeTok{iter\_warmup =} \DecValTok{1000}\NormalTok{,}
    \AttributeTok{iter\_sampling =} \DecValTok{1000}\NormalTok{, }\AttributeTok{adapt\_delta =} \FloatTok{0.95}\NormalTok{,}
    \AttributeTok{refresh =} \DecValTok{0}\NormalTok{, }\AttributeTok{show\_messages =} \ConstantTok{FALSE}
\NormalTok{  )}
\NormalTok{  wba }\OtherTok{\textless{}{-}}\NormalTok{ mod\_ml\_wba}\SpecialCharTok{$}\FunctionTok{sample}\NormalTok{(}
    \AttributeTok{data =}\NormalTok{ stan\_data, }\AttributeTok{chains =} \DecValTok{2}\NormalTok{, }\AttributeTok{iter\_warmup =} \DecValTok{1000}\NormalTok{,}
    \AttributeTok{iter\_sampling =} \DecValTok{1000}\NormalTok{, }\AttributeTok{adapt\_delta =} \FloatTok{0.95}\NormalTok{,}
    \AttributeTok{refresh =} \DecValTok{0}\NormalTok{, }\AttributeTok{show\_messages =} \ConstantTok{FALSE}
\NormalTok{  )}

  \CommentTok{\# Return pre{-}computed LOO objects, not the raw fits}
  \FunctionTok{list}\NormalTok{(}
    \AttributeTok{SBA =} \FunctionTok{loo}\NormalTok{(sba}\SpecialCharTok{$}\FunctionTok{draws}\NormalTok{(}\StringTok{"log\_lik"}\NormalTok{, }\AttributeTok{format =} \StringTok{"matrix"}\NormalTok{)),}
    \AttributeTok{PBA =} \FunctionTok{loo}\NormalTok{(pba}\SpecialCharTok{$}\FunctionTok{draws}\NormalTok{(}\StringTok{"log\_lik"}\NormalTok{, }\AttributeTok{format =} \StringTok{"matrix"}\NormalTok{)),}
    \AttributeTok{WBA =} \FunctionTok{loo}\NormalTok{(wba}\SpecialCharTok{$}\FunctionTok{draws}\NormalTok{(}\StringTok{"log\_lik"}\NormalTok{, }\AttributeTok{format =} \StringTok{"matrix"}\NormalTok{))}
\NormalTok{  )}
\NormalTok{\}}

\ControlFlowTok{if}\NormalTok{ (regenerate\_simulations }\SpecialCharTok{||} \SpecialCharTok{!}\FunctionTok{file.exists}\NormalTok{(}\StringTok{"simmodels/ch10\_ml\_scenario\_loos.rds"}\NormalTok{)) \{}
\NormalTok{  loos\_scen1 }\OtherTok{\textless{}{-}} \FunctionTok{fit\_ml\_scenario}\NormalTok{(scen1\_data, }\StringTok{"SBA{-}Mimic Population"}\NormalTok{)}
\NormalTok{  loos\_scen2 }\OtherTok{\textless{}{-}} \FunctionTok{fit\_ml\_scenario}\NormalTok{(scen2\_data, }\StringTok{"PBA Population"}\NormalTok{)}
\NormalTok{  loos\_scen3 }\OtherTok{\textless{}{-}} \FunctionTok{fit\_ml\_scenario}\NormalTok{(scen3\_data, }\StringTok{"Full WBA Population"}\NormalTok{)}
  \FunctionTok{saveRDS}\NormalTok{(}\FunctionTok{list}\NormalTok{(}\AttributeTok{s1 =}\NormalTok{ loos\_scen1, }\AttributeTok{s2 =}\NormalTok{ loos\_scen2, }\AttributeTok{s3 =}\NormalTok{ loos\_scen3),}
          \StringTok{"simmodels/ch10\_ml\_scenario\_loos.rds"}\NormalTok{)}
\NormalTok{\} }\ControlFlowTok{else}\NormalTok{ \{}
\NormalTok{  loos\_all   }\OtherTok{\textless{}{-}} \FunctionTok{readRDS}\NormalTok{(}\StringTok{"simmodels/ch10\_ml\_scenario\_loos.rds"}\NormalTok{)}
\NormalTok{  loos\_scen1 }\OtherTok{\textless{}{-}}\NormalTok{ loos\_all}\SpecialCharTok{$}\NormalTok{s1}
\NormalTok{  loos\_scen2 }\OtherTok{\textless{}{-}}\NormalTok{ loos\_all}\SpecialCharTok{$}\NormalTok{s2}
\NormalTok{  loos\_scen3 }\OtherTok{\textless{}{-}}\NormalTok{ loos\_all}\SpecialCharTok{$}\NormalTok{s3}
\NormalTok{\}}
\end{Highlighting}
\end{Shaded}

\begin{Shaded}
\begin{Highlighting}[]
\NormalTok{compute\_ml\_scenario }\OtherTok{\textless{}{-}} \ControlFlowTok{function}\NormalTok{(loos, scenario\_name) \{}
  \FunctionTok{cat}\NormalTok{(}\FunctionTok{sprintf}\NormalTok{(}\StringTok{"}\SpecialCharTok{\textbackslash{}n}\StringTok{=== \%s ===}\SpecialCharTok{\textbackslash{}n}\StringTok{"}\NormalTok{, scenario\_name))}
\NormalTok{  comp }\OtherTok{\textless{}{-}} \FunctionTok{loo\_compare}\NormalTok{(loos)}
  \FunctionTok{print}\NormalTok{(comp)}

\NormalTok{  w }\OtherTok{\textless{}{-}} \FunctionTok{loo\_model\_weights}\NormalTok{(loos, }\AttributeTok{method =} \StringTok{"stacking"}\NormalTok{)}
  \FunctionTok{cat}\NormalTok{(}\StringTok{"Stacking weights:}\SpecialCharTok{\textbackslash{}n}\StringTok{"}\NormalTok{)}
  \FunctionTok{print}\NormalTok{(}\FunctionTok{round}\NormalTok{(w, }\DecValTok{3}\NormalTok{))}

  \FunctionTok{list}\NormalTok{(}\AttributeTok{loos =}\NormalTok{ loos, }\AttributeTok{weights =}\NormalTok{ w, }\AttributeTok{comparison =}\NormalTok{ comp)}
\NormalTok{\}}

\NormalTok{res\_scen1 }\OtherTok{\textless{}{-}} \FunctionTok{compute\_ml\_scenario}\NormalTok{(loos\_scen1, }\StringTok{"Scenario 1: SBA{-}Mimic Population"}\NormalTok{)}
\end{Highlighting}
\end{Shaded}

\begin{verbatim}

=== Scenario 1: SBA-Mimic Population ===
    elpd_diff se_diff
PBA  0.0       0.0   
SBA -0.7       2.2   
WBA -1.6       0.9   
Stacking weights:
Method: stacking
------
    weight
SBA 0.379 
PBA 0.621 
WBA 0.000 
\end{verbatim}

\begin{Shaded}
\begin{Highlighting}[]
\NormalTok{res\_scen2 }\OtherTok{\textless{}{-}} \FunctionTok{compute\_ml\_scenario}\NormalTok{(loos\_scen2, }\StringTok{"Scenario 2: PBA Population"}\NormalTok{)}
\end{Highlighting}
\end{Shaded}

\begin{verbatim}

=== Scenario 2: PBA Population ===
    elpd_diff se_diff
PBA  0.0       0.0   
WBA -1.2       0.7   
SBA -3.6       3.0   
Stacking weights:
Method: stacking
------
    weight
SBA 0.096 
PBA 0.904 
WBA 0.000 
\end{verbatim}

\begin{Shaded}
\begin{Highlighting}[]
\NormalTok{res\_scen3 }\OtherTok{\textless{}{-}} \FunctionTok{compute\_ml\_scenario}\NormalTok{(loos\_scen3, }\StringTok{"Scenario 3: Full WBA Population"}\NormalTok{)}
\end{Highlighting}
\end{Shaded}

\begin{verbatim}

=== Scenario 3: Full WBA Population ===
    elpd_diff se_diff
WBA   0.0       0.0  
PBA  -1.2       1.6  
SBA -16.1       5.7  
Stacking weights:
Method: stacking
------
    weight
SBA 0.004 
PBA 0.078 
WBA 0.918 
\end{verbatim}

\begin{Shaded}
\begin{Highlighting}[]
\CommentTok{\# Build summary table}
\NormalTok{extract\_ml\_loo }\OtherTok{\textless{}{-}} \ControlFlowTok{function}\NormalTok{(res, name) \{}
\NormalTok{  comp }\OtherTok{\textless{}{-}} \FunctionTok{as\_tibble}\NormalTok{(res}\SpecialCharTok{$}\NormalTok{comparison, }\AttributeTok{rownames =} \StringTok{"Model"}\NormalTok{) }\SpecialCharTok{|\textgreater{}}
    \FunctionTok{select}\NormalTok{(Model, elpd\_diff, se\_diff)}
\NormalTok{  w }\OtherTok{\textless{}{-}} \FunctionTok{tibble}\NormalTok{(}\AttributeTok{Model =} \FunctionTok{names}\NormalTok{(res}\SpecialCharTok{$}\NormalTok{weights),}
              \AttributeTok{weight =} \FunctionTok{as.numeric}\NormalTok{(res}\SpecialCharTok{$}\NormalTok{weights))}
\NormalTok{  comp }\SpecialCharTok{|\textgreater{}} \FunctionTok{left\_join}\NormalTok{(w, }\AttributeTok{by =} \StringTok{"Model"}\NormalTok{) }\SpecialCharTok{|\textgreater{}} \FunctionTok{mutate}\NormalTok{(}\AttributeTok{Scenario =}\NormalTok{ name)}
\NormalTok{\}}

\NormalTok{ml\_scenario\_summary }\OtherTok{\textless{}{-}} \FunctionTok{bind\_rows}\NormalTok{(}
  \FunctionTok{extract\_ml\_loo}\NormalTok{(res\_scen1, }\StringTok{"1. SBA{-}Mimic"}\NormalTok{),}
  \FunctionTok{extract\_ml\_loo}\NormalTok{(res\_scen2, }\StringTok{"2. PBA Population"}\NormalTok{),}
  \FunctionTok{extract\_ml\_loo}\NormalTok{(res\_scen3, }\StringTok{"3. Full WBA"}\NormalTok{)}
\NormalTok{) }\SpecialCharTok{|\textgreater{}}
  \FunctionTok{mutate}\NormalTok{(}
    \AttributeTok{Model    =} \FunctionTok{factor}\NormalTok{(Model, }\AttributeTok{levels =} \FunctionTok{c}\NormalTok{(}\StringTok{"SBA"}\NormalTok{, }\StringTok{"PBA"}\NormalTok{, }\StringTok{"WBA"}\NormalTok{)),}
    \AttributeTok{Scenario =} \FunctionTok{factor}\NormalTok{(Scenario)}
\NormalTok{  )}

\CommentTok{\# ELPD differences}
\NormalTok{p\_ml\_elpd }\OtherTok{\textless{}{-}} \FunctionTok{ggplot}\NormalTok{(ml\_scenario\_summary,}
                     \FunctionTok{aes}\NormalTok{(}\AttributeTok{x =}\NormalTok{ Model, }\AttributeTok{y =}\NormalTok{ elpd\_diff, }\AttributeTok{colour =}\NormalTok{ Model)) }\SpecialCharTok{+}
  \FunctionTok{geom\_hline}\NormalTok{(}\AttributeTok{yintercept =} \DecValTok{0}\NormalTok{, }\AttributeTok{linetype =} \StringTok{"dashed"}\NormalTok{, }\AttributeTok{colour =} \StringTok{"grey50"}\NormalTok{) }\SpecialCharTok{+}
  \FunctionTok{geom\_pointrange}\NormalTok{(}\FunctionTok{aes}\NormalTok{(}\AttributeTok{ymin =}\NormalTok{ elpd\_diff }\SpecialCharTok{{-}}\NormalTok{ se\_diff,}
                       \AttributeTok{ymax =}\NormalTok{ elpd\_diff }\SpecialCharTok{+}\NormalTok{ se\_diff),}
                  \AttributeTok{size =} \FloatTok{0.8}\NormalTok{, }\AttributeTok{linewidth =} \DecValTok{1}\NormalTok{) }\SpecialCharTok{+}
  \FunctionTok{facet\_wrap}\NormalTok{(}\SpecialCharTok{\textasciitilde{}}\NormalTok{ Scenario) }\SpecialCharTok{+}
  \FunctionTok{scale\_colour\_brewer}\NormalTok{(}\AttributeTok{palette =} \StringTok{"Set2"}\NormalTok{) }\SpecialCharTok{+}
  \FunctionTok{labs}\NormalTok{(}\AttributeTok{title =} \StringTok{"Multilevel ELPD Differences"}\NormalTok{, }\AttributeTok{y =} \StringTok{"ELPD Difference"}\NormalTok{, }\AttributeTok{x =} \ConstantTok{NULL}\NormalTok{) }\SpecialCharTok{+}
  \FunctionTok{theme}\NormalTok{(}\AttributeTok{legend.position =} \StringTok{"none"}\NormalTok{,}
        \AttributeTok{strip.text =} \FunctionTok{element\_text}\NormalTok{(}\AttributeTok{face =} \StringTok{"bold"}\NormalTok{, }\AttributeTok{size =} \DecValTok{11}\NormalTok{))}

\CommentTok{\# Stacking weights}
\NormalTok{p\_ml\_weight }\OtherTok{\textless{}{-}} \FunctionTok{ggplot}\NormalTok{(ml\_scenario\_summary,}
                       \FunctionTok{aes}\NormalTok{(}\AttributeTok{x =}\NormalTok{ Model, }\AttributeTok{y =}\NormalTok{ weight, }\AttributeTok{fill =}\NormalTok{ Model)) }\SpecialCharTok{+}
  \FunctionTok{geom\_col}\NormalTok{(}\AttributeTok{alpha =} \FloatTok{0.85}\NormalTok{, }\AttributeTok{width =} \FloatTok{0.7}\NormalTok{) }\SpecialCharTok{+}
  \FunctionTok{geom\_text}\NormalTok{(}\FunctionTok{aes}\NormalTok{(}\AttributeTok{label =}\NormalTok{ scales}\SpecialCharTok{::}\FunctionTok{percent}\NormalTok{(weight, }\AttributeTok{accuracy =} \DecValTok{1}\NormalTok{)),}
            \AttributeTok{vjust =} \SpecialCharTok{{-}}\FloatTok{0.5}\NormalTok{, }\AttributeTok{fontface =} \StringTok{"bold"}\NormalTok{) }\SpecialCharTok{+}
  \FunctionTok{facet\_wrap}\NormalTok{(}\SpecialCharTok{\textasciitilde{}}\NormalTok{ Scenario) }\SpecialCharTok{+}
  \FunctionTok{scale\_fill\_brewer}\NormalTok{(}\AttributeTok{palette =} \StringTok{"Set2"}\NormalTok{) }\SpecialCharTok{+}
  \FunctionTok{scale\_y\_continuous}\NormalTok{(}\AttributeTok{labels =}\NormalTok{ scales}\SpecialCharTok{::}\NormalTok{percent, }\AttributeTok{limits =} \FunctionTok{c}\NormalTok{(}\DecValTok{0}\NormalTok{, }\FloatTok{1.2}\NormalTok{)) }\SpecialCharTok{+}
  \FunctionTok{labs}\NormalTok{(}\AttributeTok{title =} \StringTok{"Multilevel Stacking Weights"}\NormalTok{, }\AttributeTok{y =} \StringTok{"Stacking Weight"}\NormalTok{, }\AttributeTok{x =} \ConstantTok{NULL}\NormalTok{) }\SpecialCharTok{+}
  \FunctionTok{theme}\NormalTok{(}\AttributeTok{legend.position =} \StringTok{"none"}\NormalTok{, }\AttributeTok{strip.text =} \FunctionTok{element\_blank}\NormalTok{())}

\NormalTok{p\_ml\_elpd }\SpecialCharTok{/}\NormalTok{ p\_ml\_weight }\SpecialCharTok{+}
  \FunctionTok{plot\_annotation}\NormalTok{(}
    \AttributeTok{title =} \StringTok{"Multilevel Model Recovery: Scenario Case Studies"}\NormalTok{,}
    \AttributeTok{theme =} \FunctionTok{theme}\NormalTok{(}\AttributeTok{plot.title =} \FunctionTok{element\_text}\NormalTok{(}\AttributeTok{size =} \DecValTok{16}\NormalTok{, }\AttributeTok{face =} \StringTok{"bold"}\NormalTok{))}
\NormalTok{  )}
\end{Highlighting}
\end{Shaded}

\begin{center}
\includegraphics[width=0.8\linewidth,height=\textheight,keepaspectratio]{10-BayesianModels_files/figure-pdf/ch10_ml_scenario_comparison-1.pdf}
\end{center}

The expected pattern mirrors the single-agent case. In Scenario 1, the
SBA-mimic population has no individual variation in weights, so the
ML-SBA (zero parameters) should win or tie with the ML-PBA and ML-WBA
--- complexity is penalised. In Scenario 2, the PBA population requires
individual variation in \(\rho\) but not in \(\kappa\), so the ML-PBA
should win. In Scenario 3, the full WBA population needs both \(\rho\)
and \(\kappa\) variation, and the ML-WBA should dominate.

\subsection{LOO Pointwise Diagnostics Across
Scenarios}\label{loo-pointwise-diagnostics-across-scenarios}

\begin{Shaded}
\begin{Highlighting}[]
\NormalTok{make\_ml\_pw }\OtherTok{\textless{}{-}} \ControlFlowTok{function}\NormalTok{(loos, stan\_data, scenario\_name, }\AttributeTok{ref\_model =} \StringTok{"SBA"}\NormalTok{) \{}
\NormalTok{  elpd\_wba }\OtherTok{\textless{}{-}}\NormalTok{ loos}\SpecialCharTok{$}\NormalTok{WBA}\SpecialCharTok{$}\NormalTok{pointwise[, }\StringTok{"elpd\_loo"}\NormalTok{]}
\NormalTok{  elpd\_ref }\OtherTok{\textless{}{-}}\NormalTok{ loos[[ref\_model]]}\SpecialCharTok{$}\NormalTok{pointwise[, }\StringTok{"elpd\_loo"}\NormalTok{]}
  
  \FunctionTok{tibble}\NormalTok{(}
    \AttributeTok{obs       =} \FunctionTok{seq\_len}\NormalTok{(stan\_data}\SpecialCharTok{$}\NormalTok{N),}
    \AttributeTok{blue1     =}\NormalTok{ stan\_data}\SpecialCharTok{$}\NormalTok{blue1,}
    \AttributeTok{blue2     =}\NormalTok{ stan\_data}\SpecialCharTok{$}\NormalTok{blue2,}
    \AttributeTok{elpd\_diff =}\NormalTok{ elpd\_wba }\SpecialCharTok{{-}}\NormalTok{ elpd\_ref,}
    \AttributeTok{Scenario  =}\NormalTok{ scenario\_name}
\NormalTok{  ) }\SpecialCharTok{|\textgreater{}}
    \FunctionTok{mutate}\NormalTok{(}
      \AttributeTok{evidence\_conflict =} \FunctionTok{abs}\NormalTok{(blue1 }\SpecialCharTok{/} \DecValTok{8} \SpecialCharTok{{-}}\NormalTok{ blue2 }\SpecialCharTok{/} \DecValTok{3}\NormalTok{)}
\NormalTok{    )}
\NormalTok{\}}

\NormalTok{ml\_pw }\OtherTok{\textless{}{-}} \FunctionTok{bind\_rows}\NormalTok{(}
  \FunctionTok{make\_ml\_pw}\NormalTok{(res\_scen1}\SpecialCharTok{$}\NormalTok{loos, scen1\_data,}
             \StringTok{"1. SBA{-}Mimic}\SpecialCharTok{\textbackslash{}n}\StringTok{(ML{-}WBA vs ML{-}SBA)"}\NormalTok{, }\StringTok{"SBA"}\NormalTok{),}
  \FunctionTok{make\_ml\_pw}\NormalTok{(res\_scen2}\SpecialCharTok{$}\NormalTok{loos, scen2\_data,}
             \StringTok{"2. PBA Population}\SpecialCharTok{\textbackslash{}n}\StringTok{(ML{-}WBA vs ML{-}PBA)"}\NormalTok{, }\StringTok{"PBA"}\NormalTok{),}
  \FunctionTok{make\_ml\_pw}\NormalTok{(res\_scen3}\SpecialCharTok{$}\NormalTok{loos, scen3\_data,}
             \StringTok{"3. Full WBA}\SpecialCharTok{\textbackslash{}n}\StringTok{(ML{-}WBA vs ML{-}PBA)"}\NormalTok{, }\StringTok{"PBA"}\NormalTok{)}
\NormalTok{)}

\FunctionTok{ggplot}\NormalTok{(ml\_pw, }\FunctionTok{aes}\NormalTok{(}\AttributeTok{x =}\NormalTok{ evidence\_conflict, }\AttributeTok{y =}\NormalTok{ elpd\_diff)) }\SpecialCharTok{+}
  \FunctionTok{geom\_hline}\NormalTok{(}\AttributeTok{yintercept =} \DecValTok{0}\NormalTok{, }\AttributeTok{linetype =} \StringTok{"dashed"}\NormalTok{, }\AttributeTok{colour =} \StringTok{"grey40"}\NormalTok{) }\SpecialCharTok{+}
  \FunctionTok{geom\_point}\NormalTok{(}\FunctionTok{aes}\NormalTok{(}\AttributeTok{colour =}\NormalTok{ elpd\_diff }\SpecialCharTok{\textgreater{}} \DecValTok{0}\NormalTok{), }\AttributeTok{size =} \FloatTok{1.5}\NormalTok{, }\AttributeTok{alpha =} \FloatTok{0.4}\NormalTok{) }\SpecialCharTok{+}
  \FunctionTok{geom\_smooth}\NormalTok{(}\AttributeTok{method =} \StringTok{"gam"}\NormalTok{, }\AttributeTok{formula =}\NormalTok{ y }\SpecialCharTok{\textasciitilde{}} \FunctionTok{s}\NormalTok{(x, }\AttributeTok{bs =} \StringTok{"cs"}\NormalTok{),}
              \AttributeTok{colour =} \StringTok{"midnightblue"}\NormalTok{, }\AttributeTok{linewidth =} \DecValTok{1}\NormalTok{) }\SpecialCharTok{+}
  \FunctionTok{facet\_wrap}\NormalTok{(}\SpecialCharTok{\textasciitilde{}}\NormalTok{ Scenario) }\SpecialCharTok{+}
  \FunctionTok{scale\_colour\_manual}\NormalTok{(}
    \AttributeTok{values =} \FunctionTok{c}\NormalTok{(}\StringTok{"FALSE"} \OtherTok{=} \StringTok{"firebrick"}\NormalTok{, }\StringTok{"TRUE"} \OtherTok{=} \StringTok{"steelblue"}\NormalTok{),}
    \AttributeTok{labels =} \FunctionTok{c}\NormalTok{(}\StringTok{"FALSE"} \OtherTok{=} \StringTok{"Favours simpler model"}\NormalTok{, }\StringTok{"TRUE"} \OtherTok{=} \StringTok{"Favours ML{-}WBA"}\NormalTok{),}
    \AttributeTok{name   =} \ConstantTok{NULL}
\NormalTok{  ) }\SpecialCharTok{+}
  \FunctionTok{labs}\NormalTok{(}
    \AttributeTok{title =} \StringTok{"Multilevel Pointwise ELPD: Where Does Complexity Help?"}\NormalTok{,}
    \AttributeTok{subtitle =} \StringTok{"In mimicking scenarios (left), the WBA loses by a constant complexity penalty.}\SpecialCharTok{\textbackslash{}n}\StringTok{In the full WBA population (right), advantage emerges on high{-}conflict trials."}\NormalTok{,}
    \AttributeTok{x =} \StringTok{"Evidence Conflict |prop\_direct − prop\_social|"}\NormalTok{,}
    \AttributeTok{y =} \StringTok{"ELPD Difference"}
\NormalTok{  ) }\SpecialCharTok{+}
  \FunctionTok{theme}\NormalTok{(}\AttributeTok{legend.position =} \StringTok{"bottom"}\NormalTok{)}
\end{Highlighting}
\end{Shaded}

\begin{center}
\includegraphics[width=0.8\linewidth,height=\textheight,keepaspectratio]{10-BayesianModels_files/figure-pdf/ch10_ml_scenario_pointwise-1.pdf}
\end{center}

\subsection{Full Model Recovery
Simulation}\label{full-model-recovery-simulation}

We now scale up to a full model recovery simulation at the multilevel
level. For each of the three generating models (ML-SBA, ML-PBA, ML-WBA),
we generate multiple synthetic populations, cross-fit all three
candidates, and examine the confusion matrix and calibration.

\begin{Shaded}
\begin{Highlighting}[]
\NormalTok{n\_ml\_sims }\OtherTok{\textless{}{-}} \DecValTok{30}  \CommentTok{\# Per generating model}
\NormalTok{J\_recovery }\OtherTok{\textless{}{-}} \DecValTok{12}  \CommentTok{\# Agents per simulated population}

\ControlFlowTok{if}\NormalTok{ (regenerate\_simulations }\SpecialCharTok{||} \SpecialCharTok{!}\FunctionTok{file.exists}\NormalTok{(}\StringTok{"simmodels/ch10\_ml\_model\_recovery.rds"}\NormalTok{)) \{}
  
\NormalTok{  ml\_recovery\_results }\OtherTok{\textless{}{-}} \FunctionTok{list}\NormalTok{()}
\NormalTok{  counter }\OtherTok{\textless{}{-}} \DecValTok{0}
  
  \ControlFlowTok{for}\NormalTok{ (true\_model }\ControlFlowTok{in} \FunctionTok{c}\NormalTok{(}\StringTok{"SBA"}\NormalTok{, }\StringTok{"PBA"}\NormalTok{, }\StringTok{"WBA"}\NormalTok{)) \{}
    \ControlFlowTok{for}\NormalTok{ (sim\_id }\ControlFlowTok{in} \FunctionTok{seq\_len}\NormalTok{(n\_ml\_sims)) \{}
\NormalTok{      counter }\OtherTok{\textless{}{-}}\NormalTok{ counter }\SpecialCharTok{+} \DecValTok{1}
      
      \CommentTok{\# Generate population parameters}
      \ControlFlowTok{if}\NormalTok{ (true\_model }\SpecialCharTok{==} \StringTok{"SBA"}\NormalTok{) \{}
\NormalTok{        rhos   }\OtherTok{\textless{}{-}} \FunctionTok{rep}\NormalTok{(}\FloatTok{0.5}\NormalTok{, J\_recovery)}
\NormalTok{        kappas }\OtherTok{\textless{}{-}} \FunctionTok{rep}\NormalTok{(}\FloatTok{2.0}\NormalTok{, J\_recovery)}
\NormalTok{      \} }\ControlFlowTok{else} \ControlFlowTok{if}\NormalTok{ (true\_model }\SpecialCharTok{==} \StringTok{"PBA"}\NormalTok{) \{}
\NormalTok{        rhos   }\OtherTok{\textless{}{-}} \FunctionTok{plogis}\NormalTok{(}\FunctionTok{rnorm}\NormalTok{(J\_recovery, }\DecValTok{0}\NormalTok{, }\FloatTok{0.6}\NormalTok{))}
\NormalTok{        kappas }\OtherTok{\textless{}{-}} \FunctionTok{rep}\NormalTok{(}\FloatTok{1.0}\NormalTok{, J\_recovery)}
\NormalTok{      \} }\ControlFlowTok{else}\NormalTok{ \{}
\NormalTok{        draws\_raw }\OtherTok{\textless{}{-}}\NormalTok{ MASS}\SpecialCharTok{::}\FunctionTok{mvrnorm}\NormalTok{(}
\NormalTok{          J\_recovery,}
          \AttributeTok{mu    =} \FunctionTok{c}\NormalTok{(}\DecValTok{0}\NormalTok{, }\FunctionTok{log}\NormalTok{(}\DecValTok{2}\NormalTok{)),}
          \AttributeTok{Sigma =} \FunctionTok{matrix}\NormalTok{(}\FunctionTok{c}\NormalTok{(}\FloatTok{0.4}\NormalTok{, }\SpecialCharTok{{-}}\FloatTok{0.1}\NormalTok{, }\SpecialCharTok{{-}}\FloatTok{0.1}\NormalTok{, }\FloatTok{0.15}\NormalTok{), }\DecValTok{2}\NormalTok{)}
\NormalTok{        )}
\NormalTok{        rhos   }\OtherTok{\textless{}{-}} \FunctionTok{plogis}\NormalTok{(draws\_raw[, }\DecValTok{1}\NormalTok{])}
\NormalTok{        kappas }\OtherTok{\textless{}{-}} \FunctionTok{exp}\NormalTok{(draws\_raw[, }\DecValTok{2}\NormalTok{])}
\NormalTok{      \}}
      
\NormalTok{      sim\_stan }\OtherTok{\textless{}{-}} \FunctionTok{generate\_ml\_data}\NormalTok{(rhos, kappas, evidence\_combinations, }\AttributeTok{n\_reps =} \DecValTok{2}\NormalTok{)}
      
      \CommentTok{\# Fit all three models (reduced iterations for speed)}
\NormalTok{      f\_sba }\OtherTok{\textless{}{-}}\NormalTok{ mod\_ml\_sba}\SpecialCharTok{$}\FunctionTok{sample}\NormalTok{(}
        \AttributeTok{data =}\NormalTok{ sim\_stan, }\AttributeTok{chains =} \DecValTok{2}\NormalTok{, }\AttributeTok{iter\_warmup =} \DecValTok{500}\NormalTok{,}
        \AttributeTok{iter\_sampling =} \DecValTok{500}\NormalTok{, }\AttributeTok{adapt\_delta =} \FloatTok{0.9}\NormalTok{,}
        \AttributeTok{refresh =} \DecValTok{0}\NormalTok{, }\AttributeTok{show\_messages =} \ConstantTok{FALSE}
\NormalTok{      )}
\NormalTok{      f\_pba }\OtherTok{\textless{}{-}}\NormalTok{ mod\_ml\_pba}\SpecialCharTok{$}\FunctionTok{sample}\NormalTok{(}
        \AttributeTok{data =}\NormalTok{ sim\_stan, }\AttributeTok{chains =} \DecValTok{2}\NormalTok{, }\AttributeTok{iter\_warmup =} \DecValTok{500}\NormalTok{,}
        \AttributeTok{iter\_sampling =} \DecValTok{500}\NormalTok{, }\AttributeTok{adapt\_delta =} \FloatTok{0.95}\NormalTok{,}
        \AttributeTok{refresh =} \DecValTok{0}\NormalTok{, }\AttributeTok{show\_messages =} \ConstantTok{FALSE}
\NormalTok{      )}
\NormalTok{      f\_wba }\OtherTok{\textless{}{-}}\NormalTok{ mod\_ml\_wba}\SpecialCharTok{$}\FunctionTok{sample}\NormalTok{(}
        \AttributeTok{data =}\NormalTok{ sim\_stan, }\AttributeTok{chains =} \DecValTok{2}\NormalTok{, }\AttributeTok{iter\_warmup =} \DecValTok{500}\NormalTok{,}
        \AttributeTok{iter\_sampling =} \DecValTok{500}\NormalTok{, }\AttributeTok{adapt\_delta =} \FloatTok{0.95}\NormalTok{,}
        \AttributeTok{refresh =} \DecValTok{0}\NormalTok{, }\AttributeTok{show\_messages =} \ConstantTok{FALSE}
\NormalTok{      )}
      
      \CommentTok{\# Compute LOO and stacking weights}
\NormalTok{      l\_sba }\OtherTok{\textless{}{-}} \FunctionTok{loo}\NormalTok{(f\_sba}\SpecialCharTok{$}\FunctionTok{draws}\NormalTok{(}\StringTok{"log\_lik"}\NormalTok{, }\AttributeTok{format =} \StringTok{"matrix"}\NormalTok{))}
\NormalTok{      l\_pba }\OtherTok{\textless{}{-}} \FunctionTok{loo}\NormalTok{(f\_pba}\SpecialCharTok{$}\FunctionTok{draws}\NormalTok{(}\StringTok{"log\_lik"}\NormalTok{, }\AttributeTok{format =} \StringTok{"matrix"}\NormalTok{))}
\NormalTok{      l\_wba }\OtherTok{\textless{}{-}} \FunctionTok{loo}\NormalTok{(f\_wba}\SpecialCharTok{$}\FunctionTok{draws}\NormalTok{(}\StringTok{"log\_lik"}\NormalTok{, }\AttributeTok{format =} \StringTok{"matrix"}\NormalTok{))}
      
\NormalTok{      w }\OtherTok{\textless{}{-}} \FunctionTok{tryCatch}\NormalTok{(}
        \FunctionTok{loo\_model\_weights}\NormalTok{(}
          \FunctionTok{list}\NormalTok{(}\AttributeTok{SBA =}\NormalTok{ l\_sba, }\AttributeTok{PBA =}\NormalTok{ l\_pba, }\AttributeTok{WBA =}\NormalTok{ l\_wba),}
          \AttributeTok{method =} \StringTok{"stacking"}
\NormalTok{        ),}
        \AttributeTok{error =} \ControlFlowTok{function}\NormalTok{(e) }\FunctionTok{c}\NormalTok{(}\AttributeTok{SBA =} \DecValTok{1}\SpecialCharTok{/}\DecValTok{3}\NormalTok{, }\AttributeTok{PBA =} \DecValTok{1}\SpecialCharTok{/}\DecValTok{3}\NormalTok{, }\AttributeTok{WBA =} \DecValTok{1}\SpecialCharTok{/}\DecValTok{3}\NormalTok{)}
\NormalTok{      )}
      
\NormalTok{      ml\_recovery\_results[[counter]] }\OtherTok{\textless{}{-}} \FunctionTok{tibble}\NormalTok{(}
        \AttributeTok{sim\_id     =}\NormalTok{ sim\_id,}
        \AttributeTok{true\_model =}\NormalTok{ true\_model,}
        \AttributeTok{weight\_SBA =} \FunctionTok{as.numeric}\NormalTok{(w[}\StringTok{"SBA"}\NormalTok{]),}
        \AttributeTok{weight\_PBA =} \FunctionTok{as.numeric}\NormalTok{(w[}\StringTok{"PBA"}\NormalTok{]),}
        \AttributeTok{weight\_WBA =} \FunctionTok{as.numeric}\NormalTok{(w[}\StringTok{"WBA"}\NormalTok{]),}
        \AttributeTok{mean\_true\_rho   =} \FunctionTok{mean}\NormalTok{(rhos),}
        \AttributeTok{sd\_true\_rho     =} \FunctionTok{sd}\NormalTok{(rhos),}
        \AttributeTok{mean\_true\_kappa =} \FunctionTok{mean}\NormalTok{(kappas),}
        \AttributeTok{sd\_true\_kappa   =} \FunctionTok{sd}\NormalTok{(kappas)}
\NormalTok{      )}
      
      \ControlFlowTok{if}\NormalTok{ (counter }\SpecialCharTok{\%\%} \DecValTok{10} \SpecialCharTok{==} \DecValTok{0}\NormalTok{) \{}
        \FunctionTok{cat}\NormalTok{(}\FunctionTok{sprintf}\NormalTok{(}\StringTok{"  Completed \%d / \%d simulations}\SpecialCharTok{\textbackslash{}n}\StringTok{"}\NormalTok{, counter, }\DecValTok{3} \SpecialCharTok{*}\NormalTok{ n\_ml\_sims))}
\NormalTok{      \}}
\NormalTok{    \}}
\NormalTok{  \}}
  
\NormalTok{  ml\_recovery\_df }\OtherTok{\textless{}{-}} \FunctionTok{bind\_rows}\NormalTok{(ml\_recovery\_results)}
  \FunctionTok{saveRDS}\NormalTok{(ml\_recovery\_df, }\StringTok{"simmodels/ch10\_ml\_model\_recovery.rds"}\NormalTok{)}
  
\NormalTok{\} }\ControlFlowTok{else}\NormalTok{ \{}
\NormalTok{  ml\_recovery\_df }\OtherTok{\textless{}{-}} \FunctionTok{readRDS}\NormalTok{(}\StringTok{"simmodels/ch10\_ml\_model\_recovery.rds"}\NormalTok{)}
\NormalTok{\}}
\end{Highlighting}
\end{Shaded}

\subsection{Confusion Matrix}\label{confusion-matrix-1}

\begin{Shaded}
\begin{Highlighting}[]
\CommentTok{\# Binary confusion matrix (max{-}weight winner)}
\NormalTok{ml\_binary }\OtherTok{\textless{}{-}}\NormalTok{ ml\_recovery\_df }\SpecialCharTok{|\textgreater{}}
  \FunctionTok{pivot\_longer}\NormalTok{(}\FunctionTok{starts\_with}\NormalTok{(}\StringTok{"weight\_"}\NormalTok{),}
               \AttributeTok{names\_to =} \StringTok{"predicted\_model"}\NormalTok{, }\AttributeTok{values\_to =} \StringTok{"weight"}\NormalTok{) }\SpecialCharTok{|\textgreater{}}
  \FunctionTok{mutate}\NormalTok{(}\AttributeTok{predicted\_model =} \FunctionTok{str\_remove}\NormalTok{(predicted\_model, }\StringTok{"weight\_"}\NormalTok{)) }\SpecialCharTok{|\textgreater{}}
  \FunctionTok{group\_by}\NormalTok{(sim\_id, true\_model) }\SpecialCharTok{|\textgreater{}}
  \FunctionTok{slice\_max}\NormalTok{(weight, }\AttributeTok{n =} \DecValTok{1}\NormalTok{, }\AttributeTok{with\_ties =} \ConstantTok{FALSE}\NormalTok{) }\SpecialCharTok{|\textgreater{}}
  \FunctionTok{ungroup}\NormalTok{() }\SpecialCharTok{|\textgreater{}}
  \FunctionTok{count}\NormalTok{(true\_model, predicted\_model) }\SpecialCharTok{|\textgreater{}}
  \FunctionTok{group\_by}\NormalTok{(true\_model) }\SpecialCharTok{|\textgreater{}}
  \FunctionTok{mutate}\NormalTok{(}\AttributeTok{proportion =}\NormalTok{ n }\SpecialCharTok{/} \FunctionTok{sum}\NormalTok{(n)) }\SpecialCharTok{|\textgreater{}}
  \FunctionTok{ungroup}\NormalTok{() }\SpecialCharTok{|\textgreater{}}
  \FunctionTok{mutate}\NormalTok{(}
    \AttributeTok{true\_model      =} \FunctionTok{factor}\NormalTok{(true\_model, }\AttributeTok{levels =} \FunctionTok{c}\NormalTok{(}\StringTok{"SBA"}\NormalTok{, }\StringTok{"PBA"}\NormalTok{, }\StringTok{"WBA"}\NormalTok{)),}
    \AttributeTok{predicted\_model =} \FunctionTok{factor}\NormalTok{(predicted\_model, }\AttributeTok{levels =} \FunctionTok{c}\NormalTok{(}\StringTok{"SBA"}\NormalTok{, }\StringTok{"PBA"}\NormalTok{, }\StringTok{"WBA"}\NormalTok{))}
\NormalTok{  ) }\SpecialCharTok{|\textgreater{}}
  \FunctionTok{complete}\NormalTok{(true\_model, predicted\_model, }\AttributeTok{fill =} \FunctionTok{list}\NormalTok{(}\AttributeTok{n =} \DecValTok{0}\NormalTok{, }\AttributeTok{proportion =} \DecValTok{0}\NormalTok{))}

\NormalTok{ml\_binary }\SpecialCharTok{|\textgreater{}}
  \FunctionTok{select}\NormalTok{(}\SpecialCharTok{{-}}\NormalTok{n) }\SpecialCharTok{|\textgreater{}}
  \FunctionTok{pivot\_wider}\NormalTok{(}\AttributeTok{names\_from =}\NormalTok{ predicted\_model, }\AttributeTok{values\_from =}\NormalTok{ proportion) }\SpecialCharTok{|\textgreater{}}
\NormalTok{  knitr}\SpecialCharTok{::}\FunctionTok{kable}\NormalTok{(}\AttributeTok{digits =} \DecValTok{2}\NormalTok{,}
               \AttributeTok{caption =} \StringTok{"Multilevel Binary Confusion Matrix: Win Rate"}\NormalTok{)}
\end{Highlighting}
\end{Shaded}

\begin{longtable}[]{@{}lrrr@{}}
\caption{Multilevel Binary Confusion Matrix: Win Rate}\tabularnewline
\toprule\noalign{}
true\_model & SBA & PBA & WBA \\
\midrule\noalign{}
\endfirsthead
\toprule\noalign{}
true\_model & SBA & PBA & WBA \\
\midrule\noalign{}
\endhead
\bottomrule\noalign{}
\endlastfoot
SBA & 0.77 & 0.17 & 0.07 \\
PBA & 0.17 & 0.67 & 0.17 \\
WBA & 0.27 & 0.30 & 0.43 \\
\end{longtable}

\begin{Shaded}
\begin{Highlighting}[]
\NormalTok{ml\_binary }\SpecialCharTok{|\textgreater{}}
  \FunctionTok{ggplot}\NormalTok{(}\FunctionTok{aes}\NormalTok{(}\AttributeTok{x =}\NormalTok{ predicted\_model, }\AttributeTok{y =}\NormalTok{ true\_model, }\AttributeTok{fill =}\NormalTok{ proportion)) }\SpecialCharTok{+}
  \FunctionTok{geom\_tile}\NormalTok{(}\AttributeTok{colour =} \StringTok{"white"}\NormalTok{, }\AttributeTok{linewidth =} \DecValTok{1}\NormalTok{) }\SpecialCharTok{+}
  \FunctionTok{geom\_text}\NormalTok{(}\FunctionTok{aes}\NormalTok{(}\AttributeTok{label =}\NormalTok{ scales}\SpecialCharTok{::}\FunctionTok{percent}\NormalTok{(proportion, }\AttributeTok{accuracy =} \DecValTok{1}\NormalTok{)),}
            \AttributeTok{size =} \DecValTok{5}\NormalTok{, }\AttributeTok{fontface =} \StringTok{"bold"}\NormalTok{) }\SpecialCharTok{+}
  \FunctionTok{scale\_fill\_viridis\_c}\NormalTok{(}\AttributeTok{option =} \StringTok{"mako"}\NormalTok{, }\AttributeTok{direction =} \SpecialCharTok{{-}}\DecValTok{1}\NormalTok{, }\AttributeTok{limits =} \FunctionTok{c}\NormalTok{(}\DecValTok{0}\NormalTok{, }\DecValTok{1}\NormalTok{)) }\SpecialCharTok{+}
  \FunctionTok{scale\_y\_discrete}\NormalTok{(}\AttributeTok{limits =}\NormalTok{ rev) }\SpecialCharTok{+}
  \FunctionTok{labs}\NormalTok{(}
    \AttributeTok{title =} \StringTok{"Multilevel Model Recovery: Binary Classification"}\NormalTok{,}
    \AttributeTok{x =} \StringTok{"Winning Candidate Model"}\NormalTok{,}
    \AttributeTok{y =} \StringTok{"True Generating Model"}\NormalTok{,}
    \AttributeTok{fill =} \StringTok{"Win Rate"}
\NormalTok{  ) }\SpecialCharTok{+}
  \FunctionTok{theme}\NormalTok{(}\AttributeTok{panel.grid =} \FunctionTok{element\_blank}\NormalTok{())}
\end{Highlighting}
\end{Shaded}

\begin{center}
\includegraphics[width=0.8\linewidth,height=\textheight,keepaspectratio]{10-BayesianModels_files/figure-pdf/ch10_ml_confusion_matrix-1.pdf}
\end{center}

\begin{Shaded}
\begin{Highlighting}[]
\CommentTok{\# Probabilistic confusion matrix (average stacking weights)}
\NormalTok{ml\_prob\_confusion }\OtherTok{\textless{}{-}}\NormalTok{ ml\_recovery\_df }\SpecialCharTok{|\textgreater{}}
  \FunctionTok{group\_by}\NormalTok{(true\_model) }\SpecialCharTok{|\textgreater{}}
  \FunctionTok{summarise}\NormalTok{(}
    \AttributeTok{Recovered\_SBA =} \FunctionTok{mean}\NormalTok{(weight\_SBA),}
    \AttributeTok{Recovered\_PBA =} \FunctionTok{mean}\NormalTok{(weight\_PBA),}
    \AttributeTok{Recovered\_WBA =} \FunctionTok{mean}\NormalTok{(weight\_WBA),}
    \AttributeTok{.groups =} \StringTok{"drop"}
\NormalTok{  ) }\SpecialCharTok{|\textgreater{}}
  \FunctionTok{mutate}\NormalTok{(}\AttributeTok{true\_model =} \FunctionTok{factor}\NormalTok{(true\_model, }\AttributeTok{levels =} \FunctionTok{c}\NormalTok{(}\StringTok{"SBA"}\NormalTok{, }\StringTok{"PBA"}\NormalTok{, }\StringTok{"WBA"}\NormalTok{)))}

\NormalTok{knitr}\SpecialCharTok{::}\FunctionTok{kable}\NormalTok{(ml\_prob\_confusion, }\AttributeTok{digits =} \DecValTok{2}\NormalTok{,}
             \AttributeTok{caption =} \StringTok{"Multilevel Probabilistic Confusion Matrix: Mean Stacking Weights"}\NormalTok{)}
\end{Highlighting}
\end{Shaded}

\begin{longtable}[]{@{}lrrr@{}}
\caption{Multilevel Probabilistic Confusion Matrix: Mean Stacking
Weights}\tabularnewline
\toprule\noalign{}
true\_model & Recovered\_SBA & Recovered\_PBA & Recovered\_WBA \\
\midrule\noalign{}
\endfirsthead
\toprule\noalign{}
true\_model & Recovered\_SBA & Recovered\_PBA & Recovered\_WBA \\
\midrule\noalign{}
\endhead
\bottomrule\noalign{}
\endlastfoot
PBA & 0.21 & 0.69 & 0.10 \\
SBA & 0.79 & 0.14 & 0.07 \\
WBA & 0.35 & 0.29 & 0.36 \\
\end{longtable}

\begin{Shaded}
\begin{Highlighting}[]
\NormalTok{ml\_prob\_confusion }\SpecialCharTok{|\textgreater{}}
  \FunctionTok{pivot\_longer}\NormalTok{(}\FunctionTok{starts\_with}\NormalTok{(}\StringTok{"Recovered"}\NormalTok{),}
               \AttributeTok{names\_to =} \StringTok{"fitted\_model"}\NormalTok{, }\AttributeTok{values\_to =} \StringTok{"weight"}\NormalTok{) }\SpecialCharTok{|\textgreater{}}
  \FunctionTok{mutate}\NormalTok{(}\AttributeTok{fitted\_model =} \FunctionTok{factor}\NormalTok{(}\FunctionTok{str\_remove}\NormalTok{(fitted\_model, }\StringTok{"Recovered\_"}\NormalTok{),}
                                \AttributeTok{levels =} \FunctionTok{c}\NormalTok{(}\StringTok{"SBA"}\NormalTok{, }\StringTok{"PBA"}\NormalTok{, }\StringTok{"WBA"}\NormalTok{))) }\SpecialCharTok{|\textgreater{}}
  \FunctionTok{ggplot}\NormalTok{(}\FunctionTok{aes}\NormalTok{(}\AttributeTok{x =}\NormalTok{ fitted\_model, }\AttributeTok{y =}\NormalTok{ true\_model, }\AttributeTok{fill =}\NormalTok{ weight)) }\SpecialCharTok{+}
  \FunctionTok{geom\_tile}\NormalTok{(}\AttributeTok{colour =} \StringTok{"white"}\NormalTok{, }\AttributeTok{linewidth =} \DecValTok{1}\NormalTok{) }\SpecialCharTok{+}
  \FunctionTok{geom\_text}\NormalTok{(}\FunctionTok{aes}\NormalTok{(}\AttributeTok{label =}\NormalTok{ scales}\SpecialCharTok{::}\FunctionTok{percent}\NormalTok{(weight, }\AttributeTok{accuracy =} \DecValTok{1}\NormalTok{)),}
            \AttributeTok{size =} \DecValTok{5}\NormalTok{, }\AttributeTok{fontface =} \StringTok{"bold"}\NormalTok{) }\SpecialCharTok{+}
  \FunctionTok{scale\_fill\_viridis\_c}\NormalTok{(}\AttributeTok{option =} \StringTok{"mako"}\NormalTok{, }\AttributeTok{direction =} \SpecialCharTok{{-}}\DecValTok{1}\NormalTok{, }\AttributeTok{limits =} \FunctionTok{c}\NormalTok{(}\DecValTok{0}\NormalTok{, }\DecValTok{1}\NormalTok{)) }\SpecialCharTok{+}
  \FunctionTok{scale\_y\_discrete}\NormalTok{(}\AttributeTok{limits =}\NormalTok{ rev) }\SpecialCharTok{+}
  \FunctionTok{labs}\NormalTok{(}
    \AttributeTok{title =} \StringTok{"Multilevel Model Recovery: Average Stacking Weights"}\NormalTok{,}
    \AttributeTok{x =} \StringTok{"Candidate Model"}\NormalTok{,}
    \AttributeTok{y =} \StringTok{"True Generating Model"}\NormalTok{,}
    \AttributeTok{fill =} \StringTok{"Avg. Weight"}
\NormalTok{  ) }\SpecialCharTok{+}
  \FunctionTok{theme}\NormalTok{(}\AttributeTok{panel.grid =} \FunctionTok{element\_blank}\NormalTok{())}
\end{Highlighting}
\end{Shaded}

\begin{center}
\includegraphics[width=0.8\linewidth,height=\textheight,keepaspectratio]{10-BayesianModels_files/figure-pdf/ch10_ml_probabilistic_confusion-1.pdf}
\end{center}

The expected pattern: the diagonal should dominate. Because the models
are nested (SBA \(\subset\) PBA \(\subset\) WBA), Occam's razor should
produce \emph{asymmetric} off-diagonal errors --- the WBA should
sometimes be classified as PBA or SBA when its parameters mimic simpler
models, but the SBA should rarely be classified as WBA when extra
parameters are unnecessary. This asymmetry is a \emph{feature} of the
comparison pipeline, not a bug.

\subsection{Model Comparison
Calibration}\label{model-comparison-calibration}

A well-calibrated comparison procedure assigns a stacking weight of,
say, 0.7 to a model that turns out to be the true generator roughly 70\%
of the time. We assess this by binning stacking weights and computing
the empirical frequency with which the candidate was the true model.

\begin{Shaded}
\begin{Highlighting}[]
\NormalTok{ml\_cal\_long }\OtherTok{\textless{}{-}}\NormalTok{ ml\_recovery\_df }\SpecialCharTok{|\textgreater{}}
  \FunctionTok{pivot\_longer}\NormalTok{(}\FunctionTok{starts\_with}\NormalTok{(}\StringTok{"weight\_"}\NormalTok{),}
               \AttributeTok{names\_to =} \StringTok{"candidate"}\NormalTok{, }\AttributeTok{values\_to =} \StringTok{"stacking\_weight"}\NormalTok{) }\SpecialCharTok{|\textgreater{}}
  \FunctionTok{mutate}\NormalTok{(}
    \AttributeTok{candidate =} \FunctionTok{str\_remove}\NormalTok{(candidate, }\StringTok{"weight\_"}\NormalTok{),}
    \AttributeTok{is\_true   =} \FunctionTok{as.integer}\NormalTok{(candidate }\SpecialCharTok{==}\NormalTok{ true\_model)}
\NormalTok{  )}

\NormalTok{n\_bins }\OtherTok{\textless{}{-}} \DecValTok{8}
\NormalTok{ml\_cal\_binned }\OtherTok{\textless{}{-}}\NormalTok{ ml\_cal\_long }\SpecialCharTok{|\textgreater{}}
  \FunctionTok{mutate}\NormalTok{(}\AttributeTok{bin =} \FunctionTok{cut}\NormalTok{(stacking\_weight,}
                   \AttributeTok{breaks =} \FunctionTok{seq}\NormalTok{(}\DecValTok{0}\NormalTok{, }\DecValTok{1}\NormalTok{, }\AttributeTok{length.out =}\NormalTok{ n\_bins }\SpecialCharTok{+} \DecValTok{1}\NormalTok{),}
                   \AttributeTok{include.lowest =} \ConstantTok{TRUE}\NormalTok{)) }\SpecialCharTok{|\textgreater{}}
  \FunctionTok{group\_by}\NormalTok{(candidate, bin) }\SpecialCharTok{|\textgreater{}}
  \FunctionTok{summarise}\NormalTok{(}
    \AttributeTok{bin\_mid    =} \FunctionTok{mean}\NormalTok{(stacking\_weight),}
    \AttributeTok{observed   =} \FunctionTok{mean}\NormalTok{(is\_true),}
    \AttributeTok{N\_datasets =} \FunctionTok{n}\NormalTok{(),}
    \AttributeTok{.groups    =} \StringTok{"drop"}
\NormalTok{  ) }\SpecialCharTok{|\textgreater{}}
  \FunctionTok{mutate}\NormalTok{(}
    \AttributeTok{SE    =} \FunctionTok{sqrt}\NormalTok{((observed }\SpecialCharTok{*}\NormalTok{ (}\DecValTok{1} \SpecialCharTok{{-}}\NormalTok{ observed)) }\SpecialCharTok{/} \FunctionTok{pmax}\NormalTok{(N\_datasets, }\DecValTok{1}\NormalTok{)),}
    \AttributeTok{Lower =} \FunctionTok{pmax}\NormalTok{(}\DecValTok{0}\NormalTok{, observed }\SpecialCharTok{{-}} \FloatTok{1.96} \SpecialCharTok{*}\NormalTok{ SE),}
    \AttributeTok{Upper =} \FunctionTok{pmin}\NormalTok{(}\DecValTok{1}\NormalTok{, observed }\SpecialCharTok{+} \FloatTok{1.96} \SpecialCharTok{*}\NormalTok{ SE)}
\NormalTok{  )}

\FunctionTok{ggplot}\NormalTok{(ml\_cal\_binned, }\FunctionTok{aes}\NormalTok{(}\AttributeTok{x =}\NormalTok{ bin\_mid, }\AttributeTok{y =}\NormalTok{ observed, }\AttributeTok{colour =}\NormalTok{ candidate)) }\SpecialCharTok{+}
  \FunctionTok{geom\_abline}\NormalTok{(}\AttributeTok{slope =} \DecValTok{1}\NormalTok{, }\AttributeTok{intercept =} \DecValTok{0}\NormalTok{, }\AttributeTok{linetype =} \StringTok{"dashed"}\NormalTok{, }\AttributeTok{colour =} \StringTok{"grey40"}\NormalTok{) }\SpecialCharTok{+}
  \FunctionTok{geom\_ribbon}\NormalTok{(}\FunctionTok{aes}\NormalTok{(}\AttributeTok{ymin =}\NormalTok{ Lower, }\AttributeTok{ymax =}\NormalTok{ Upper, }\AttributeTok{fill =}\NormalTok{ candidate),}
              \AttributeTok{alpha =} \FloatTok{0.15}\NormalTok{, }\AttributeTok{colour =} \ConstantTok{NA}\NormalTok{) }\SpecialCharTok{+}
  \FunctionTok{geom\_line}\NormalTok{(}\AttributeTok{linewidth =} \DecValTok{1}\NormalTok{) }\SpecialCharTok{+}
  \FunctionTok{geom\_point}\NormalTok{(}\AttributeTok{size =} \DecValTok{2}\NormalTok{) }\SpecialCharTok{+}
  \FunctionTok{scale\_colour\_brewer}\NormalTok{(}\AttributeTok{palette =} \StringTok{"Set2"}\NormalTok{) }\SpecialCharTok{+}
  \FunctionTok{scale\_fill\_brewer}\NormalTok{(}\AttributeTok{palette =} \StringTok{"Set2"}\NormalTok{) }\SpecialCharTok{+}
  \FunctionTok{scale\_x\_continuous}\NormalTok{(}\AttributeTok{labels =}\NormalTok{ scales}\SpecialCharTok{::}\NormalTok{percent, }\AttributeTok{limits =} \FunctionTok{c}\NormalTok{(}\DecValTok{0}\NormalTok{, }\DecValTok{1}\NormalTok{)) }\SpecialCharTok{+}
  \FunctionTok{scale\_y\_continuous}\NormalTok{(}\AttributeTok{labels =}\NormalTok{ scales}\SpecialCharTok{::}\NormalTok{percent, }\AttributeTok{limits =} \FunctionTok{c}\NormalTok{(}\DecValTok{0}\NormalTok{, }\DecValTok{1}\NormalTok{)) }\SpecialCharTok{+}
  \FunctionTok{labs}\NormalTok{(}
    \AttributeTok{title =} \StringTok{"Multilevel Model Comparison Calibration"}\NormalTok{,}
    \AttributeTok{subtitle =} \StringTok{"Points on the diagonal = well{-}calibrated stacking weights.}\SpecialCharTok{\textbackslash{}n}\StringTok{Ribbons = 95\% CI; wider ribbons reflect fewer datasets in that bin."}\NormalTok{,}
    \AttributeTok{x =} \StringTok{"Mean Stacking Weight (Binned)"}\NormalTok{,}
    \AttributeTok{y =} \StringTok{"Empirical Frequency (True Generating Model)"}\NormalTok{,}
    \AttributeTok{colour =} \StringTok{"Candidate"}\NormalTok{, }\AttributeTok{fill =} \StringTok{"Candidate"}
\NormalTok{  ) }\SpecialCharTok{+}
  \FunctionTok{theme}\NormalTok{(}\AttributeTok{legend.position =} \StringTok{"bottom"}\NormalTok{)}
\end{Highlighting}
\end{Shaded}

\begin{center}
\includegraphics[width=0.8\linewidth,height=\textheight,keepaspectratio]{10-BayesianModels_files/figure-pdf/ch10_ml_model_calibration-1.pdf}
\end{center}

Points on the diagonal indicate perfect calibration: a stacking weight
of 0.7 means the model is the true generator 70\% of the time. A curve
bowed above the diagonal signals underconfidence (the procedure hedges
too much); a curve below signals overconfidence. The multilevel
comparison has additional sources of uncertainty relative to the
single-agent case --- partial pooling, correlation estimation, and
funnel geometry --- so some additional calibration noise is expected.

\section{Summary: What the Multilevel Pipeline
Adds}\label{summary-what-the-multilevel-pipeline-adds}

The multilevel extension answers the \emph{population-level} question:
how do people in general weight evidence, and how much do they vary? By
using the reparameterised \((\rho_j, \kappa_j)\) parameterisation ---
not the raw \((w_d, w_s)\) shown to be unidentifiable in the
single-agent case --- we preserve the identifiability gains while
gaining partial pooling.

The full pipeline mirrors the single-agent workflow:

\begin{enumerate}
\def\labelenumi{\arabic{enumi}.}
\item
  \textbf{Model quality:} Diagnostics, trace plots, and summary
  statistics confirmed clean sampling for all three multilevel models.
\item
  \textbf{Parameter recovery:} Individual \(\rho\) values recover well
  (identified from conflict trials); \(\kappa\) is noisier but partially
  regularised by the hierarchical prior. Partial pooling pulls extreme
  agents toward the population mean, as expected.
\item
  \textbf{SBC certification:} The multilevel SBC established that the
  Stan sampler is globally unbiased for the ML-WBA and ML-PBA population
  parameters --- a critical check given the more complex posterior
  geometry.
\item
  \textbf{Model comparison:} The three-way LOO comparison with pointwise
  ELPD visualizations showed \emph{where} each model's advantage lies
  (high-conflict trials for non-balanced populations). Scenario case
  studies demonstrated the nested Occam's razor effect: simpler
  populations are correctly attributed to simpler models.
\item
  \textbf{Model recovery:} The confusion matrix and calibration curves
  quantified the comparison procedure's reliability across many
  simulated populations, confirming appropriate asymmetric error
  patterns from the nested model structure.
\end{enumerate}

The key lesson: \textbf{the reparameterisation that solved the
single-agent identifiability problem is equally essential at the
multilevel level.} The hierarchical prior partially constrains the
scale, but it does not eliminate the ridge in each agent's conditional
likelihood. Only by separating allocation (\(\rho\)) from scaling
(\(\kappa\)) can we obtain meaningful individual-level parameter
estimates --- and by extension, meaningful population-level summaries of
how people vary.

\section{Mixed Strategies: When the Population Is
Heterogeneous}\label{mixed-strategies-when-the-population-is-heterogeneous}

Every analysis above assumed a \emph{homogeneous} population: all agents
were drawn from the same generative model. In practice, this assumption
is almost certainly wrong. A real experimental sample plausibly contains
people who use qualitatively different strategies --- some may integrate
evidence at face value (SBA-like), while others systematically
overweight social or direct evidence (WBA-like). What happens to our
pipeline when the population is a genuine mixture?

This section merges data from two generating models --- half the
participants behave as Simple Bayesian Agents, half as Weighted Bayesian
Agents with a strong social bias --- and examines whether the
hierarchical models and the comparison pipeline can still produce useful
inference.

Because no single hierarchical model can correctly represent a bimodal
population distribution, we should \emph{not} expect population-level
parameters to be meaningful. The scientific question shifts from ``what
does the average person do?'' to ``which strategy does \emph{each}
person use?'' This motivates a focus on \textbf{individual-level
parameter recovery} and \textbf{per-agent model comparison}, skipping
the usual population-level parameter interpretation.

\subsection{Generating the Mixed
Population}\label{generating-the-mixed-population}

\begin{Shaded}
\begin{Highlighting}[]
\FunctionTok{set.seed}\NormalTok{(}\DecValTok{2026}\NormalTok{)}
\NormalTok{J\_mix }\OtherTok{\textless{}{-}} \DecValTok{40}  \CommentTok{\# 20 SBA + 20 WBA agents}

\CommentTok{\# {-}{-}{-} SBA agents: strict Bayesian integration (wd = 1, ws = 1) {-}{-}{-}}
\NormalTok{sba\_rhos   }\OtherTok{\textless{}{-}} \FunctionTok{rep}\NormalTok{(}\FloatTok{0.5}\NormalTok{, }\DecValTok{20}\NormalTok{)}
\NormalTok{sba\_kappas }\OtherTok{\textless{}{-}} \FunctionTok{rep}\NormalTok{(}\FloatTok{2.0}\NormalTok{, }\DecValTok{20}\NormalTok{)}

\CommentTok{\# {-}{-}{-} WBA agents: strong social bias, higher total sensitivity {-}{-}{-}}
\CommentTok{\# Draw from a population distribution to add realistic variation}
\NormalTok{wba\_raw }\OtherTok{\textless{}{-}}\NormalTok{ MASS}\SpecialCharTok{::}\FunctionTok{mvrnorm}\NormalTok{(}
  \DecValTok{20}\NormalTok{,}
  \AttributeTok{mu    =} \FunctionTok{c}\NormalTok{(}\FunctionTok{qlogis}\NormalTok{(}\FloatTok{0.25}\NormalTok{), }\FunctionTok{log}\NormalTok{(}\FloatTok{2.5}\NormalTok{)),   }\CommentTok{\# rho ≈ 0.25, kappa ≈ 2.5}
  \AttributeTok{Sigma =} \FunctionTok{matrix}\NormalTok{(}\FunctionTok{c}\NormalTok{(}\FloatTok{0.25}\NormalTok{, }\SpecialCharTok{{-}}\FloatTok{0.08}\NormalTok{, }\SpecialCharTok{{-}}\FloatTok{0.08}\NormalTok{, }\FloatTok{0.12}\NormalTok{), }\DecValTok{2}\NormalTok{)}
\NormalTok{)}
\NormalTok{wba\_rhos   }\OtherTok{\textless{}{-}} \FunctionTok{plogis}\NormalTok{(wba\_raw[, }\DecValTok{1}\NormalTok{])}
\NormalTok{wba\_kappas }\OtherTok{\textless{}{-}} \FunctionTok{exp}\NormalTok{(wba\_raw[, }\DecValTok{2}\NormalTok{])}

\CommentTok{\# Combine and shuffle so agent IDs don\textquotesingle{}t reveal strategy}
\NormalTok{agent\_order   }\OtherTok{\textless{}{-}} \FunctionTok{sample}\NormalTok{(J\_mix)}
\NormalTok{all\_rhos      }\OtherTok{\textless{}{-}} \FunctionTok{c}\NormalTok{(sba\_rhos, wba\_rhos)[agent\_order]}
\NormalTok{all\_kappas    }\OtherTok{\textless{}{-}} \FunctionTok{c}\NormalTok{(sba\_kappas, wba\_kappas)[agent\_order]}
\NormalTok{true\_strategy }\OtherTok{\textless{}{-}} \FunctionTok{c}\NormalTok{(}\FunctionTok{rep}\NormalTok{(}\StringTok{"SBA"}\NormalTok{, }\DecValTok{20}\NormalTok{), }\FunctionTok{rep}\NormalTok{(}\StringTok{"WBA"}\NormalTok{, }\DecValTok{20}\NormalTok{))[agent\_order]}

\CommentTok{\# Generate trial{-}level data}
\NormalTok{mixed\_stan\_data }\OtherTok{\textless{}{-}} \FunctionTok{generate\_ml\_data}\NormalTok{(all\_rhos, all\_kappas,}
\NormalTok{                                     evidence\_combinations, }\AttributeTok{n\_reps =} \DecValTok{2}\NormalTok{)}

\CommentTok{\# Store ground truth for recovery plots}
\NormalTok{mixed\_truth }\OtherTok{\textless{}{-}} \FunctionTok{tibble}\NormalTok{(}
  \AttributeTok{agent\_id       =} \FunctionTok{seq\_len}\NormalTok{(J\_mix),}
  \AttributeTok{true\_rho       =}\NormalTok{ all\_rhos,}
  \AttributeTok{true\_kappa     =}\NormalTok{ all\_kappas,}
  \AttributeTok{true\_wd        =}\NormalTok{ all\_rhos }\SpecialCharTok{*}\NormalTok{ all\_kappas,}
  \AttributeTok{true\_ws        =}\NormalTok{ (}\DecValTok{1} \SpecialCharTok{{-}}\NormalTok{ all\_rhos) }\SpecialCharTok{*}\NormalTok{ all\_kappas,}
  \AttributeTok{true\_strategy  =}\NormalTok{ true\_strategy}
\NormalTok{)}

\FunctionTok{cat}\NormalTok{(}\FunctionTok{sprintf}\NormalTok{(}\StringTok{"Mixed population: \%d agents (\%d SBA + \%d WBA)}\SpecialCharTok{\textbackslash{}n}\StringTok{"}\NormalTok{,}
\NormalTok{            J\_mix, }\FunctionTok{sum}\NormalTok{(true\_strategy }\SpecialCharTok{==} \StringTok{"SBA"}\NormalTok{), }\FunctionTok{sum}\NormalTok{(true\_strategy }\SpecialCharTok{==} \StringTok{"WBA"}\NormalTok{)))}
\end{Highlighting}
\end{Shaded}

\begin{verbatim}
Mixed population: 40 agents (20 SBA + 20 WBA)
\end{verbatim}

\begin{Shaded}
\begin{Highlighting}[]
\FunctionTok{cat}\NormalTok{(}\FunctionTok{sprintf}\NormalTok{(}\StringTok{"Total trials: \%d | Trials per agent: \%d}\SpecialCharTok{\textbackslash{}n}\StringTok{"}\NormalTok{,}
\NormalTok{            mixed\_stan\_data}\SpecialCharTok{$}\NormalTok{N, mixed\_stan\_data}\SpecialCharTok{$}\NormalTok{N }\SpecialCharTok{/}\NormalTok{ J\_mix))}
\end{Highlighting}
\end{Shaded}

\begin{verbatim}
Total trials: 2880 | Trials per agent: 72
\end{verbatim}

\begin{Shaded}
\begin{Highlighting}[]
\FunctionTok{ggplot}\NormalTok{(mixed\_truth, }\FunctionTok{aes}\NormalTok{(}\AttributeTok{x =}\NormalTok{ true\_wd, }\AttributeTok{y =}\NormalTok{ true\_ws,}
                         \AttributeTok{colour =}\NormalTok{ true\_strategy, }\AttributeTok{shape =}\NormalTok{ true\_strategy)) }\SpecialCharTok{+}
  \FunctionTok{geom\_point}\NormalTok{(}\AttributeTok{size =} \FloatTok{3.5}\NormalTok{, }\AttributeTok{alpha =} \FloatTok{0.8}\NormalTok{) }\SpecialCharTok{+}
  \FunctionTok{geom\_text}\NormalTok{(}\FunctionTok{aes}\NormalTok{(}\AttributeTok{label =}\NormalTok{ agent\_id), }\AttributeTok{nudge\_x =} \FloatTok{0.04}\NormalTok{, }\AttributeTok{size =} \DecValTok{3}\NormalTok{, }\AttributeTok{show.legend =} \ConstantTok{FALSE}\NormalTok{) }\SpecialCharTok{+}
  \FunctionTok{scale\_colour\_manual}\NormalTok{(}\AttributeTok{values =} \FunctionTok{c}\NormalTok{(}\StringTok{"SBA"} \OtherTok{=} \StringTok{"\#4daf4a"}\NormalTok{, }\StringTok{"WBA"} \OtherTok{=} \StringTok{"\#984ea3"}\NormalTok{)) }\SpecialCharTok{+}
  \FunctionTok{scale\_shape\_manual}\NormalTok{(}\AttributeTok{values =} \FunctionTok{c}\NormalTok{(}\StringTok{"SBA"} \OtherTok{=} \DecValTok{16}\NormalTok{, }\StringTok{"WBA"} \OtherTok{=} \DecValTok{17}\NormalTok{)) }\SpecialCharTok{+}
  \FunctionTok{geom\_vline}\NormalTok{(}\AttributeTok{xintercept =} \DecValTok{1}\NormalTok{, }\AttributeTok{linetype =} \StringTok{"dashed"}\NormalTok{, }\AttributeTok{colour =} \StringTok{"grey50"}\NormalTok{, }\AttributeTok{alpha =} \FloatTok{0.5}\NormalTok{) }\SpecialCharTok{+}
  \FunctionTok{geom\_hline}\NormalTok{(}\AttributeTok{yintercept =} \DecValTok{1}\NormalTok{, }\AttributeTok{linetype =} \StringTok{"dashed"}\NormalTok{, }\AttributeTok{colour =} \StringTok{"grey50"}\NormalTok{, }\AttributeTok{alpha =} \FloatTok{0.5}\NormalTok{) }\SpecialCharTok{+}
  \FunctionTok{labs}\NormalTok{(}
    \AttributeTok{title    =} \StringTok{"Mixed Population: True Agent Parameters"}\NormalTok{,}
    \AttributeTok{subtitle =} \StringTok{"SBA agents cluster at (1, 1); WBA agents are socially biased with higher total sensitivity"}\NormalTok{,}
    \AttributeTok{x        =} \FunctionTok{expression}\NormalTok{(}\StringTok{"True "} \SpecialCharTok{*}\NormalTok{ w[d]),}
    \AttributeTok{y        =} \FunctionTok{expression}\NormalTok{(}\StringTok{"True "} \SpecialCharTok{*}\NormalTok{ w[s]),}
    \AttributeTok{colour   =} \StringTok{"Strategy"}\NormalTok{, }\AttributeTok{shape =} \StringTok{"Strategy"}
\NormalTok{  ) }\SpecialCharTok{+}
  \FunctionTok{theme}\NormalTok{(}\AttributeTok{legend.position =} \StringTok{"bottom"}\NormalTok{)}
\end{Highlighting}
\end{Shaded}

\begin{center}
\includegraphics[width=0.8\linewidth,height=\textheight,keepaspectratio]{10-BayesianModels_files/figure-pdf/ch10_mixed_population_viz-1.pdf}
\end{center}

The SBA agents form a tight cluster at \((w_d, w_s) = (1, 1)\) --- they
are all identical by construction. The WBA agents scatter around
\((w_d, w_s) \approx (0.6, 1.9)\), each with their own idiosyncratic
weights drawn from the population distribution. No single Gaussian
population on the \((\text{logit}(\rho), \log(\kappa))\) scale can
represent this bimodal structure.

\subsection{Fitting All Three Multilevel
Models}\label{fitting-all-three-multilevel-models}

\begin{Shaded}
\begin{Highlighting}[]
\ControlFlowTok{if}\NormalTok{ (regenerate\_simulations }\SpecialCharTok{||} \SpecialCharTok{!}\FunctionTok{all}\NormalTok{(}\FunctionTok{file.exists}\NormalTok{(}\FunctionTok{c}\NormalTok{(}\StringTok{"simmodels/ch10\_fit\_mix\_sba.rds"}\NormalTok{, }\StringTok{"simmodels/ch10\_fit\_mix\_pba.rds"}\NormalTok{, }\StringTok{"simmodels/ch10\_fit\_mix\_wba.rds"}\NormalTok{)))) \{}
  
\NormalTok{  fit\_mix\_sba }\OtherTok{\textless{}{-}}\NormalTok{ mod\_ml\_sba}\SpecialCharTok{$}\FunctionTok{sample}\NormalTok{(}
    \AttributeTok{data =}\NormalTok{ mixed\_stan\_data, }\AttributeTok{chains =} \DecValTok{4}\NormalTok{, }\AttributeTok{parallel\_chains =} \DecValTok{4}\NormalTok{,}
    \AttributeTok{iter\_warmup =} \DecValTok{1000}\NormalTok{, }\AttributeTok{iter\_sampling =} \DecValTok{1000}\NormalTok{,}
    \AttributeTok{adapt\_delta =} \FloatTok{0.9}\NormalTok{, }\AttributeTok{refresh =} \DecValTok{0}
\NormalTok{  )}

\NormalTok{  fit\_mix\_pba }\OtherTok{\textless{}{-}}\NormalTok{ mod\_ml\_pba}\SpecialCharTok{$}\FunctionTok{sample}\NormalTok{(}
    \AttributeTok{data =}\NormalTok{ mixed\_stan\_data, }\AttributeTok{chains =} \DecValTok{4}\NormalTok{, }\AttributeTok{parallel\_chains =} \DecValTok{4}\NormalTok{,}
    \AttributeTok{iter\_warmup =} \DecValTok{1000}\NormalTok{, }\AttributeTok{iter\_sampling =} \DecValTok{1000}\NormalTok{,}
    \AttributeTok{adapt\_delta =} \FloatTok{0.95}\NormalTok{, }\AttributeTok{refresh =} \DecValTok{0}
\NormalTok{  )}

\NormalTok{  fit\_mix\_wba }\OtherTok{\textless{}{-}}\NormalTok{ mod\_ml\_wba}\SpecialCharTok{$}\FunctionTok{sample}\NormalTok{(}
    \AttributeTok{data =}\NormalTok{ mixed\_stan\_data, }\AttributeTok{chains =} \DecValTok{4}\NormalTok{, }\AttributeTok{parallel\_chains =} \DecValTok{4}\NormalTok{,}
    \AttributeTok{iter\_warmup =} \DecValTok{1000}\NormalTok{, }\AttributeTok{iter\_sampling =} \DecValTok{1000}\NormalTok{,}
    \AttributeTok{adapt\_delta =} \FloatTok{0.95}\NormalTok{, }\AttributeTok{refresh =} \DecValTok{0}
\NormalTok{  )}

  \CommentTok{\# Save each fit individually with $save\_object() so draws survive sessions}
\NormalTok{  fit\_mix\_sba}\SpecialCharTok{$}\FunctionTok{save\_object}\NormalTok{(}\StringTok{"simmodels/ch10\_fit\_mix\_sba.rds"}\NormalTok{)}
\NormalTok{  fit\_mix\_pba}\SpecialCharTok{$}\FunctionTok{save\_object}\NormalTok{(}\StringTok{"simmodels/ch10\_fit\_mix\_pba.rds"}\NormalTok{)}
\NormalTok{  fit\_mix\_wba}\SpecialCharTok{$}\FunctionTok{save\_object}\NormalTok{(}\StringTok{"simmodels/ch10\_fit\_mix\_wba.rds"}\NormalTok{)}

\NormalTok{\} }\ControlFlowTok{else}\NormalTok{ \{}
\NormalTok{  fit\_mix\_sba }\OtherTok{\textless{}{-}} \FunctionTok{readRDS}\NormalTok{(}\StringTok{"simmodels/ch10\_fit\_mix\_sba.rds"}\NormalTok{)}
\NormalTok{  fit\_mix\_pba }\OtherTok{\textless{}{-}} \FunctionTok{readRDS}\NormalTok{(}\StringTok{"simmodels/ch10\_fit\_mix\_pba.rds"}\NormalTok{)}
\NormalTok{  fit\_mix\_wba }\OtherTok{\textless{}{-}} \FunctionTok{readRDS}\NormalTok{(}\StringTok{"simmodels/ch10\_fit\_mix\_wba.rds"}\NormalTok{)}
\NormalTok{\}}

\FunctionTok{cat}\NormalTok{(}\StringTok{"=== Mixed Population: Model Diagnostics ===}\SpecialCharTok{\textbackslash{}n}\StringTok{"}\NormalTok{)}
\end{Highlighting}
\end{Shaded}

\begin{verbatim}
=== Mixed Population: Model Diagnostics ===
\end{verbatim}

\begin{Shaded}
\begin{Highlighting}[]
\FunctionTok{check\_diagnostics}\NormalTok{(fit\_mix\_sba, }\StringTok{"ML{-}SBA (mixed)"}\NormalTok{)}
\end{Highlighting}
\end{Shaded}

\begin{verbatim}
ML-SBA (mixed)            | Divs:   0 | Min E-BFMI: NaN | Failures: 0
\end{verbatim}

\begin{Shaded}
\begin{Highlighting}[]
\FunctionTok{check\_diagnostics}\NormalTok{(fit\_mix\_pba, }\StringTok{"ML{-}PBA (mixed)"}\NormalTok{)}
\end{Highlighting}
\end{Shaded}

\begin{verbatim}
ML-PBA (mixed)            | Divs:   0 | Min E-BFMI: 0.62 | Failures: 0
\end{verbatim}

\begin{Shaded}
\begin{Highlighting}[]
\FunctionTok{check\_diagnostics}\NormalTok{(fit\_mix\_wba, }\StringTok{"ML{-}WBA (mixed)"}\NormalTok{)}
\end{Highlighting}
\end{Shaded}

\begin{verbatim}
ML-WBA (mixed)            | Divs:   0 | Min E-BFMI: 0.68 | Failures: 0
\end{verbatim}

\subsection{Why Population Parameters Are
Misleading}\label{why-population-parameters-are-misleading}

Before turning to individual-level analysis, let's check what the
ML-WBA's population parameters say about this mixed population.

\begin{Shaded}
\begin{Highlighting}[]
\NormalTok{draws\_mix }\OtherTok{\textless{}{-}} \FunctionTok{as\_draws\_df}\NormalTok{(fit\_mix\_wba}\SpecialCharTok{$}\FunctionTok{draws}\NormalTok{())}

\CommentTok{\# Population{-}level summaries}
\NormalTok{pop\_mix\_summary }\OtherTok{\textless{}{-}}\NormalTok{ fit\_mix\_wba}\SpecialCharTok{$}\FunctionTok{summary}\NormalTok{(}
  \AttributeTok{variables =} \FunctionTok{c}\NormalTok{(}\StringTok{"mu[1]"}\NormalTok{, }\StringTok{"mu[2]"}\NormalTok{, }\StringTok{"sigma[1]"}\NormalTok{, }\StringTok{"sigma[2]"}\NormalTok{, }\StringTok{"Omega[1,2]"}\NormalTok{,}
                \StringTok{"pop\_rho"}\NormalTok{, }\StringTok{"pop\_kappa"}\NormalTok{)}
\NormalTok{)}

\FunctionTok{cat}\NormalTok{(}\StringTok{"=== ML{-}WBA Population Parameter Estimates (Mixed Data) ===}\SpecialCharTok{\textbackslash{}n}\StringTok{"}\NormalTok{)}
\end{Highlighting}
\end{Shaded}

\begin{verbatim}
=== ML-WBA Population Parameter Estimates (Mixed Data) ===
\end{verbatim}

\begin{Shaded}
\begin{Highlighting}[]
\NormalTok{pop\_mix\_summary }\SpecialCharTok{|\textgreater{}}
  \FunctionTok{select}\NormalTok{(variable, mean, q5, q95) }\SpecialCharTok{|\textgreater{}}
\NormalTok{  knitr}\SpecialCharTok{::}\FunctionTok{kable}\NormalTok{(}\AttributeTok{digits =} \DecValTok{3}\NormalTok{, }\AttributeTok{caption =} \StringTok{"Population parameters estimated from mixed{-}strategy data"}\NormalTok{)}
\end{Highlighting}
\end{Shaded}

\begin{longtable}[]{@{}lrrr@{}}
\caption{Population parameters estimated from mixed-strategy
data}\tabularnewline
\toprule\noalign{}
variable & mean & q5 & q95 \\
\midrule\noalign{}
\endfirsthead
\toprule\noalign{}
variable & mean & q5 & q95 \\
\midrule\noalign{}
\endhead
\bottomrule\noalign{}
\endlastfoot
mu{[}1{]} & -0.457 & -0.643 & -0.265 \\
mu{[}2{]} & 0.730 & 0.324 & 1.167 \\
sigma{[}1{]} & 0.444 & 0.177 & 0.694 \\
sigma{[}2{]} & 0.279 & 0.019 & 0.744 \\
Omega{[}1,2{]} & -0.083 & -0.771 & 0.664 \\
pop\_rho & 0.388 & 0.345 & 0.434 \\
pop\_kappa & 2.147 & 1.383 & 3.213 \\
\end{longtable}

\begin{Shaded}
\begin{Highlighting}[]
\CommentTok{\# Extract estimated individual rhos and show the distribution}
\NormalTok{mix\_est\_rho }\OtherTok{\textless{}{-}} \FunctionTok{map\_dbl}\NormalTok{(}\DecValTok{1}\SpecialCharTok{:}\NormalTok{J\_mix, }\SpecialCharTok{\textasciitilde{}} \FunctionTok{mean}\NormalTok{(draws\_mix[[}\FunctionTok{paste0}\NormalTok{(}\StringTok{"rho["}\NormalTok{, .x, }\StringTok{"]"}\NormalTok{)]]))}

\FunctionTok{ggplot}\NormalTok{() }\SpecialCharTok{+}
  \FunctionTok{geom\_histogram}\NormalTok{(}\AttributeTok{data =} \FunctionTok{tibble}\NormalTok{(}\AttributeTok{rho =}\NormalTok{ mix\_est\_rho,}
                                \AttributeTok{strategy =}\NormalTok{ mixed\_truth}\SpecialCharTok{$}\NormalTok{true\_strategy),}
                 \FunctionTok{aes}\NormalTok{(}\AttributeTok{x =}\NormalTok{ rho, }\AttributeTok{fill =}\NormalTok{ strategy),}
                 \AttributeTok{bins =} \DecValTok{20}\NormalTok{, }\AttributeTok{alpha =} \FloatTok{0.7}\NormalTok{, }\AttributeTok{position =} \StringTok{"identity"}\NormalTok{) }\SpecialCharTok{+}
  \FunctionTok{geom\_vline}\NormalTok{(}\FunctionTok{aes}\NormalTok{(}\AttributeTok{xintercept =} \FunctionTok{plogis}\NormalTok{(}\FunctionTok{mean}\NormalTok{(draws\_mix}\SpecialCharTok{$}\StringTok{\textasciigrave{}}\AttributeTok{mu[1]}\StringTok{\textasciigrave{}}\NormalTok{))),}
             \AttributeTok{colour =} \StringTok{"red"}\NormalTok{, }\AttributeTok{linetype =} \StringTok{"dashed"}\NormalTok{, }\AttributeTok{linewidth =} \DecValTok{1}\NormalTok{) }\SpecialCharTok{+}
  \FunctionTok{annotate}\NormalTok{(}\StringTok{"text"}\NormalTok{,}
           \AttributeTok{x =} \FunctionTok{plogis}\NormalTok{(}\FunctionTok{mean}\NormalTok{(draws\_mix}\SpecialCharTok{$}\StringTok{\textasciigrave{}}\AttributeTok{mu[1]}\StringTok{\textasciigrave{}}\NormalTok{)) }\SpecialCharTok{+} \FloatTok{0.03}\NormalTok{,}
           \AttributeTok{y =} \ConstantTok{Inf}\NormalTok{, }\AttributeTok{vjust =} \DecValTok{2}\NormalTok{,}
           \AttributeTok{label =} \FunctionTok{sprintf}\NormalTok{(}\StringTok{"Pop. mean ρ = \%.2f"}\NormalTok{, }\FunctionTok{plogis}\NormalTok{(}\FunctionTok{mean}\NormalTok{(draws\_mix}\SpecialCharTok{$}\StringTok{\textasciigrave{}}\AttributeTok{mu[1]}\StringTok{\textasciigrave{}}\NormalTok{))),}
           \AttributeTok{colour =} \StringTok{"red"}\NormalTok{, }\AttributeTok{size =} \DecValTok{4}\NormalTok{) }\SpecialCharTok{+}
  \FunctionTok{scale\_fill\_manual}\NormalTok{(}\AttributeTok{values =} \FunctionTok{c}\NormalTok{(}\StringTok{"SBA"} \OtherTok{=} \StringTok{"\#4daf4a"}\NormalTok{, }\StringTok{"WBA"} \OtherTok{=} \StringTok{"\#984ea3"}\NormalTok{)) }\SpecialCharTok{+}
  \FunctionTok{labs}\NormalTok{(}
    \AttributeTok{title    =} \StringTok{"The Population Mean Is Meaningless for a Mixture"}\NormalTok{,}
    \AttributeTok{subtitle =} \StringTok{"Estimated individual ρ values form two clusters, but the population mean (red dashed) describes nobody"}\NormalTok{,}
    \AttributeTok{x        =} \FunctionTok{expression}\NormalTok{(}\StringTok{"Estimated "} \SpecialCharTok{*}\NormalTok{ rho[j]),}
    \AttributeTok{y        =} \StringTok{"Count"}\NormalTok{,}
    \AttributeTok{fill     =} \StringTok{"True Strategy"}
\NormalTok{  ) }\SpecialCharTok{+}
  \FunctionTok{theme}\NormalTok{(}\AttributeTok{legend.position =} \StringTok{"bottom"}\NormalTok{)}
\end{Highlighting}
\end{Shaded}

\begin{center}
\includegraphics[width=0.8\linewidth,height=\textheight,keepaspectratio]{10-BayesianModels_files/figure-pdf/ch10_mixed_population_misleading-1.pdf}
\end{center}

\subsection{Individual Parameter
Recovery}\label{individual-parameter-recovery}

\begin{Shaded}
\begin{Highlighting}[]
\CommentTok{\# Extract individual posteriors}
\NormalTok{mix\_recovery }\OtherTok{\textless{}{-}}\NormalTok{ mixed\_truth }\SpecialCharTok{|\textgreater{}}
  \FunctionTok{mutate}\NormalTok{(}
    \AttributeTok{est\_rho   =} \FunctionTok{map\_dbl}\NormalTok{(agent\_id, }\SpecialCharTok{\textasciitilde{}} \FunctionTok{mean}\NormalTok{(draws\_mix[[}\FunctionTok{paste0}\NormalTok{(}\StringTok{"rho["}\NormalTok{, .x, }\StringTok{"]"}\NormalTok{)]])),}
    \AttributeTok{est\_kappa =} \FunctionTok{map\_dbl}\NormalTok{(agent\_id, }\SpecialCharTok{\textasciitilde{}} \FunctionTok{mean}\NormalTok{(draws\_mix[[}\FunctionTok{paste0}\NormalTok{(}\StringTok{"kappa["}\NormalTok{, .x, }\StringTok{"]"}\NormalTok{)]])),}
    \AttributeTok{est\_wd    =} \FunctionTok{map\_dbl}\NormalTok{(agent\_id, }\SpecialCharTok{\textasciitilde{}} \FunctionTok{mean}\NormalTok{(draws\_mix[[}\FunctionTok{paste0}\NormalTok{(}\StringTok{"weight\_direct["}\NormalTok{, .x, }\StringTok{"]"}\NormalTok{)]])),}
    \AttributeTok{est\_ws    =} \FunctionTok{map\_dbl}\NormalTok{(agent\_id, }\SpecialCharTok{\textasciitilde{}} \FunctionTok{mean}\NormalTok{(draws\_mix[[}\FunctionTok{paste0}\NormalTok{(}\StringTok{"weight\_social["}\NormalTok{, .x, }\StringTok{"]"}\NormalTok{)]]))}
\NormalTok{  )}

\CommentTok{\# ρ recovery}
\NormalTok{p\_mix\_rho }\OtherTok{\textless{}{-}} \FunctionTok{ggplot}\NormalTok{(mix\_recovery, }\FunctionTok{aes}\NormalTok{(}\AttributeTok{x =}\NormalTok{ true\_rho, }\AttributeTok{y =}\NormalTok{ est\_rho,}
                                        \AttributeTok{colour =}\NormalTok{ true\_strategy)) }\SpecialCharTok{+}
  \FunctionTok{geom\_abline}\NormalTok{(}\AttributeTok{linetype =} \StringTok{"dashed"}\NormalTok{, }\AttributeTok{colour =} \StringTok{"grey50"}\NormalTok{) }\SpecialCharTok{+}
  \FunctionTok{geom\_point}\NormalTok{(}\AttributeTok{size =} \DecValTok{3}\NormalTok{, }\AttributeTok{alpha =} \FloatTok{0.8}\NormalTok{) }\SpecialCharTok{+}
  \FunctionTok{scale\_colour\_manual}\NormalTok{(}\AttributeTok{values =} \FunctionTok{c}\NormalTok{(}\StringTok{"SBA"} \OtherTok{=} \StringTok{"\#4daf4a"}\NormalTok{, }\StringTok{"WBA"} \OtherTok{=} \StringTok{"\#984ea3"}\NormalTok{)) }\SpecialCharTok{+}
  \FunctionTok{coord\_cartesian}\NormalTok{(}\AttributeTok{xlim =} \FunctionTok{c}\NormalTok{(}\DecValTok{0}\NormalTok{, }\DecValTok{1}\NormalTok{), }\AttributeTok{ylim =} \FunctionTok{c}\NormalTok{(}\DecValTok{0}\NormalTok{, }\DecValTok{1}\NormalTok{)) }\SpecialCharTok{+}
  \FunctionTok{labs}\NormalTok{(}\AttributeTok{title =} \FunctionTok{expression}\NormalTok{(}\StringTok{"Recovery: "} \SpecialCharTok{*}\NormalTok{ rho }\SpecialCharTok{*} \StringTok{" (relative allocation)"}\NormalTok{),}
       \AttributeTok{x =} \FunctionTok{expression}\NormalTok{(}\StringTok{"True "} \SpecialCharTok{*}\NormalTok{ rho[j]),}
       \AttributeTok{y =} \FunctionTok{expression}\NormalTok{(}\StringTok{"Estimated "} \SpecialCharTok{*}\NormalTok{ rho[j]),}
       \AttributeTok{colour =} \StringTok{"True Strategy"}\NormalTok{)}

\CommentTok{\# κ recovery}
\NormalTok{p\_mix\_kappa }\OtherTok{\textless{}{-}} \FunctionTok{ggplot}\NormalTok{(mix\_recovery, }\FunctionTok{aes}\NormalTok{(}\AttributeTok{x =}\NormalTok{ true\_kappa, }\AttributeTok{y =}\NormalTok{ est\_kappa,}
                                          \AttributeTok{colour =}\NormalTok{ true\_strategy)) }\SpecialCharTok{+}
  \FunctionTok{geom\_abline}\NormalTok{(}\AttributeTok{linetype =} \StringTok{"dashed"}\NormalTok{, }\AttributeTok{colour =} \StringTok{"grey50"}\NormalTok{) }\SpecialCharTok{+}
  \FunctionTok{geom\_point}\NormalTok{(}\AttributeTok{size =} \DecValTok{3}\NormalTok{, }\AttributeTok{alpha =} \FloatTok{0.8}\NormalTok{) }\SpecialCharTok{+}
  \FunctionTok{scale\_colour\_manual}\NormalTok{(}\AttributeTok{values =} \FunctionTok{c}\NormalTok{(}\StringTok{"SBA"} \OtherTok{=} \StringTok{"\#4daf4a"}\NormalTok{, }\StringTok{"WBA"} \OtherTok{=} \StringTok{"\#984ea3"}\NormalTok{)) }\SpecialCharTok{+}
  \FunctionTok{labs}\NormalTok{(}\AttributeTok{title =} \FunctionTok{expression}\NormalTok{(}\StringTok{"Recovery: "} \SpecialCharTok{*}\NormalTok{ kappa }\SpecialCharTok{*} \StringTok{" (total evidence scaling)"}\NormalTok{),}
       \AttributeTok{x =} \FunctionTok{expression}\NormalTok{(}\StringTok{"True "} \SpecialCharTok{*}\NormalTok{ kappa[j]),}
       \AttributeTok{y =} \FunctionTok{expression}\NormalTok{(}\StringTok{"Estimated "} \SpecialCharTok{*}\NormalTok{ kappa[j]),}
       \AttributeTok{colour =} \StringTok{"True Strategy"}\NormalTok{)}

\NormalTok{p\_mix\_rho }\SpecialCharTok{+}\NormalTok{ p\_mix\_kappa }\SpecialCharTok{+}
  \FunctionTok{plot\_layout}\NormalTok{(}\AttributeTok{guides =} \StringTok{"collect"}\NormalTok{) }\SpecialCharTok{\&}
  \FunctionTok{theme}\NormalTok{(}\AttributeTok{legend.position =} \StringTok{"bottom"}\NormalTok{)}
\end{Highlighting}
\end{Shaded}

\begin{center}
\includegraphics[width=0.8\linewidth,height=\textheight,keepaspectratio]{10-BayesianModels_files/figure-pdf/ch10_mixed_individual_recovery-1.pdf}
\end{center}

\subsection{Per-Agent Model
Comparison}\label{per-agent-model-comparison}

The population-level LOO comparison tells us which model best predicts
\emph{the aggregate dataset}. But the more informative question is:
which model best predicts \emph{each individual agent's} behaviour? We
decompose the total LOO into agent-level contributions.

\begin{Shaded}
\begin{Highlighting}[]
\CommentTok{\# Compute trial{-}level LOO for all three models}
\NormalTok{loo\_mix\_sba }\OtherTok{\textless{}{-}} \FunctionTok{loo}\NormalTok{(fit\_mix\_sba}\SpecialCharTok{$}\FunctionTok{draws}\NormalTok{(}\StringTok{"log\_lik"}\NormalTok{, }\AttributeTok{format =} \StringTok{"matrix"}\NormalTok{))}
\NormalTok{loo\_mix\_pba }\OtherTok{\textless{}{-}} \FunctionTok{loo}\NormalTok{(fit\_mix\_pba}\SpecialCharTok{$}\FunctionTok{draws}\NormalTok{(}\StringTok{"log\_lik"}\NormalTok{, }\AttributeTok{format =} \StringTok{"matrix"}\NormalTok{))}
\NormalTok{loo\_mix\_wba }\OtherTok{\textless{}{-}} \FunctionTok{loo}\NormalTok{(fit\_mix\_wba}\SpecialCharTok{$}\FunctionTok{draws}\NormalTok{(}\StringTok{"log\_lik"}\NormalTok{, }\AttributeTok{format =} \StringTok{"matrix"}\NormalTok{))}

\CommentTok{\# Population{-}level comparison (for reference)}
\FunctionTok{cat}\NormalTok{(}\StringTok{"=== Population{-}Level Comparison (Mixed Data) ===}\SpecialCharTok{\textbackslash{}n}\StringTok{"}\NormalTok{)}
\end{Highlighting}
\end{Shaded}

\begin{verbatim}
=== Population-Level Comparison (Mixed Data) ===
\end{verbatim}

\begin{Shaded}
\begin{Highlighting}[]
\FunctionTok{print}\NormalTok{(}\FunctionTok{loo\_compare}\NormalTok{(}\FunctionTok{list}\NormalTok{(}\StringTok{"ML{-}SBA"} \OtherTok{=}\NormalTok{ loo\_mix\_sba,}
                        \StringTok{"ML{-}PBA"} \OtherTok{=}\NormalTok{ loo\_mix\_pba,}
                        \StringTok{"ML{-}WBA"} \OtherTok{=}\NormalTok{ loo\_mix\_wba)))}
\end{Highlighting}
\end{Shaded}

\begin{verbatim}
       elpd_diff se_diff
ML-WBA   0.0       0.0  
ML-PBA  -2.7       2.8  
ML-SBA -17.8       6.4  
\end{verbatim}

\begin{Shaded}
\begin{Highlighting}[]
\NormalTok{ml\_mix\_weights }\OtherTok{\textless{}{-}} \FunctionTok{loo\_model\_weights}\NormalTok{(}
  \FunctionTok{list}\NormalTok{(}\StringTok{"ML{-}SBA"} \OtherTok{=}\NormalTok{ loo\_mix\_sba, }\StringTok{"ML{-}PBA"} \OtherTok{=}\NormalTok{ loo\_mix\_pba, }\StringTok{"ML{-}WBA"} \OtherTok{=}\NormalTok{ loo\_mix\_wba),}
  \AttributeTok{method =} \StringTok{"stacking"}
\NormalTok{)}
\FunctionTok{cat}\NormalTok{(}\StringTok{"}\SpecialCharTok{\textbackslash{}n}\StringTok{Stacking weights:}\SpecialCharTok{\textbackslash{}n}\StringTok{"}\NormalTok{)}
\end{Highlighting}
\end{Shaded}

\begin{verbatim}

Stacking weights:
\end{verbatim}

\begin{Shaded}
\begin{Highlighting}[]
\FunctionTok{print}\NormalTok{(}\FunctionTok{round}\NormalTok{(ml\_mix\_weights, }\DecValTok{3}\NormalTok{))}
\end{Highlighting}
\end{Shaded}

\begin{verbatim}
Method: stacking
------
       weight
ML-SBA 0.055 
ML-PBA 0.170 
ML-WBA 0.775 
\end{verbatim}

\begin{Shaded}
\begin{Highlighting}[]
\CommentTok{\# Decompose into per{-}agent ELPD sums}
\NormalTok{agent\_elpd }\OtherTok{\textless{}{-}} \FunctionTok{tibble}\NormalTok{(}
  \AttributeTok{trial        =} \FunctionTok{seq\_len}\NormalTok{(mixed\_stan\_data}\SpecialCharTok{$}\NormalTok{N),}
  \AttributeTok{agent\_id     =}\NormalTok{ mixed\_stan\_data}\SpecialCharTok{$}\NormalTok{agent\_id,}
  \AttributeTok{elpd\_sba     =}\NormalTok{ loo\_mix\_sba}\SpecialCharTok{$}\NormalTok{pointwise[, }\StringTok{"elpd\_loo"}\NormalTok{],}
  \AttributeTok{elpd\_pba     =}\NormalTok{ loo\_mix\_pba}\SpecialCharTok{$}\NormalTok{pointwise[, }\StringTok{"elpd\_loo"}\NormalTok{],}
  \AttributeTok{elpd\_wba     =}\NormalTok{ loo\_mix\_wba}\SpecialCharTok{$}\NormalTok{pointwise[, }\StringTok{"elpd\_loo"}\NormalTok{]}
\NormalTok{) }\SpecialCharTok{|\textgreater{}}
  \FunctionTok{group\_by}\NormalTok{(agent\_id) }\SpecialCharTok{|\textgreater{}}
  \FunctionTok{summarise}\NormalTok{(}
    \AttributeTok{elpd\_SBA =} \FunctionTok{sum}\NormalTok{(elpd\_sba),}
    \AttributeTok{elpd\_PBA =} \FunctionTok{sum}\NormalTok{(elpd\_pba),}
    \AttributeTok{elpd\_WBA =} \FunctionTok{sum}\NormalTok{(elpd\_wba),}
    \AttributeTok{.groups  =} \StringTok{"drop"}
\NormalTok{  ) }\SpecialCharTok{|\textgreater{}}
  \FunctionTok{left\_join}\NormalTok{(mixed\_truth }\SpecialCharTok{|\textgreater{}} \FunctionTok{select}\NormalTok{(agent\_id, true\_strategy), }\AttributeTok{by =} \StringTok{"agent\_id"}\NormalTok{) }\SpecialCharTok{|\textgreater{}}
  \FunctionTok{mutate}\NormalTok{(}
    \AttributeTok{best\_model =} \FunctionTok{case\_when}\NormalTok{(}
\NormalTok{      elpd\_SBA }\SpecialCharTok{\textgreater{}=}\NormalTok{ elpd\_PBA }\SpecialCharTok{\&}\NormalTok{ elpd\_SBA }\SpecialCharTok{\textgreater{}=}\NormalTok{ elpd\_WBA }\SpecialCharTok{\textasciitilde{}} \StringTok{"SBA"}\NormalTok{,}
\NormalTok{      elpd\_PBA }\SpecialCharTok{\textgreater{}=}\NormalTok{ elpd\_SBA }\SpecialCharTok{\&}\NormalTok{ elpd\_PBA }\SpecialCharTok{\textgreater{}=}\NormalTok{ elpd\_WBA }\SpecialCharTok{\textasciitilde{}} \StringTok{"PBA"}\NormalTok{,}
      \ConstantTok{TRUE} \SpecialCharTok{\textasciitilde{}} \StringTok{"WBA"}
\NormalTok{    ),}
    \CommentTok{\# ELPD advantage of the best flexible model over the SBA baseline}
    \AttributeTok{elpd\_advantage =} \FunctionTok{pmax}\NormalTok{(elpd\_PBA, elpd\_WBA) }\SpecialCharTok{{-}}\NormalTok{ elpd\_SBA}
\NormalTok{  )}
\end{Highlighting}
\end{Shaded}

\begin{Shaded}
\begin{Highlighting}[]
\CommentTok{\# Left: per{-}agent ELPD advantage coloured by true strategy}
\NormalTok{p\_agent\_advantage }\OtherTok{\textless{}{-}} \FunctionTok{ggplot}\NormalTok{(agent\_elpd,}
       \FunctionTok{aes}\NormalTok{(}\AttributeTok{x =} \FunctionTok{reorder}\NormalTok{(agent\_id, elpd\_advantage),}
           \AttributeTok{y =}\NormalTok{ elpd\_advantage,}
           \AttributeTok{fill =}\NormalTok{ true\_strategy)) }\SpecialCharTok{+}
  \FunctionTok{geom\_col}\NormalTok{(}\AttributeTok{alpha =} \FloatTok{0.8}\NormalTok{) }\SpecialCharTok{+}
  \FunctionTok{geom\_hline}\NormalTok{(}\AttributeTok{yintercept =} \DecValTok{0}\NormalTok{, }\AttributeTok{linetype =} \StringTok{"dashed"}\NormalTok{, }\AttributeTok{colour =} \StringTok{"grey40"}\NormalTok{) }\SpecialCharTok{+}
  \FunctionTok{scale\_fill\_manual}\NormalTok{(}\AttributeTok{values =} \FunctionTok{c}\NormalTok{(}\StringTok{"SBA"} \OtherTok{=} \StringTok{"\#4daf4a"}\NormalTok{, }\StringTok{"WBA"} \OtherTok{=} \StringTok{"\#984ea3"}\NormalTok{)) }\SpecialCharTok{+}
  \FunctionTok{coord\_flip}\NormalTok{() }\SpecialCharTok{+}
  \FunctionTok{labs}\NormalTok{(}
    \AttributeTok{title    =} \StringTok{"Per{-}Agent ELPD: Best Flexible Model vs. SBA"}\NormalTok{,}
    \AttributeTok{subtitle =} \StringTok{"Positive = flexible model predicts better. Negative = SBA sufficient."}\NormalTok{,}
    \AttributeTok{x        =} \StringTok{"Agent (sorted by advantage)"}\NormalTok{,}
    \AttributeTok{y        =} \StringTok{"ELPD Advantage (max(PBA, WBA) − SBA)"}\NormalTok{,}
    \AttributeTok{fill     =} \StringTok{"True Strategy"}
\NormalTok{  ) }\SpecialCharTok{+}
  \FunctionTok{theme}\NormalTok{(}
    \AttributeTok{axis.text.y  =} \FunctionTok{element\_text}\NormalTok{(}\AttributeTok{size =} \DecValTok{6}\NormalTok{),}
    \AttributeTok{legend.position =} \StringTok{"bottom"}
\NormalTok{  )}

\CommentTok{\# Right: scatter of per{-}agent SBA vs WBA ELPD}
\NormalTok{p\_agent\_scatter }\OtherTok{\textless{}{-}} \FunctionTok{ggplot}\NormalTok{(agent\_elpd,}
       \FunctionTok{aes}\NormalTok{(}\AttributeTok{x =}\NormalTok{ elpd\_SBA, }\AttributeTok{y =}\NormalTok{ elpd\_WBA,}
           \AttributeTok{colour =}\NormalTok{ true\_strategy, }\AttributeTok{shape =}\NormalTok{ best\_model)) }\SpecialCharTok{+}
  \FunctionTok{geom\_abline}\NormalTok{(}\AttributeTok{linetype =} \StringTok{"dashed"}\NormalTok{, }\AttributeTok{colour =} \StringTok{"grey50"}\NormalTok{) }\SpecialCharTok{+}
  \FunctionTok{geom\_point}\NormalTok{(}\AttributeTok{size =} \DecValTok{3}\NormalTok{, }\AttributeTok{alpha =} \FloatTok{0.8}\NormalTok{) }\SpecialCharTok{+}
  \FunctionTok{scale\_colour\_manual}\NormalTok{(}\AttributeTok{values =} \FunctionTok{c}\NormalTok{(}\StringTok{"SBA"} \OtherTok{=} \StringTok{"\#4daf4a"}\NormalTok{, }\StringTok{"WBA"} \OtherTok{=} \StringTok{"\#984ea3"}\NormalTok{)) }\SpecialCharTok{+}
  \FunctionTok{scale\_shape\_manual}\NormalTok{(}\AttributeTok{values =} \FunctionTok{c}\NormalTok{(}\StringTok{"SBA"} \OtherTok{=} \DecValTok{16}\NormalTok{, }\StringTok{"PBA"} \OtherTok{=} \DecValTok{15}\NormalTok{, }\StringTok{"WBA"} \OtherTok{=} \DecValTok{17}\NormalTok{)) }\SpecialCharTok{+}
  \FunctionTok{labs}\NormalTok{(}
    \AttributeTok{title    =} \StringTok{"Agent{-}Level LOO: SBA vs. WBA"}\NormalTok{,}
    \AttributeTok{subtitle =} \StringTok{"Points above diagonal: WBA predicts better. Shape = winning model."}\NormalTok{,}
    \AttributeTok{x        =} \StringTok{"Agent ELPD under ML{-}SBA"}\NormalTok{,}
    \AttributeTok{y        =} \StringTok{"Agent ELPD under ML{-}WBA"}\NormalTok{,}
    \AttributeTok{colour   =} \StringTok{"True Strategy"}\NormalTok{, }\AttributeTok{shape =} \StringTok{"Best Model"}
\NormalTok{  ) }\SpecialCharTok{+}
  \FunctionTok{theme}\NormalTok{(}\AttributeTok{legend.position =} \StringTok{"bottom"}\NormalTok{)}

\NormalTok{p\_agent\_advantage }\SpecialCharTok{|}\NormalTok{ p\_agent\_scatter}
\end{Highlighting}
\end{Shaded}

\begin{center}
\includegraphics[width=0.8\linewidth,height=\textheight,keepaspectratio]{10-BayesianModels_files/figure-pdf/ch10_mixed_agent_loo_viz-1.pdf}
\end{center}

The bar chart on the left reveals a clean separation. The WBA agents
(purple) consistently show positive ELPD advantages --- the flexible
model predicts their behaviour much better than the SBA. The SBA agents
(green) hover near zero or show slight negative advantages --- the SBA
is sufficient, and the flexible models' extra parameters represent
wasted complexity. The scatter plot on the right makes the same point
geometrically: WBA agents sit above the diagonal (ML-WBA predicts
better), while SBA agents sit on or below it.

\subsection{Individual-Level Confusion
Matrix}\label{individual-level-confusion-matrix}

\begin{Shaded}
\begin{Highlighting}[]
\CommentTok{\# Build the confusion matrix}
\NormalTok{mix\_confusion }\OtherTok{\textless{}{-}}\NormalTok{ agent\_elpd }\SpecialCharTok{|\textgreater{}}
  \FunctionTok{count}\NormalTok{(true\_strategy, best\_model) }\SpecialCharTok{|\textgreater{}}
  \FunctionTok{group\_by}\NormalTok{(true\_strategy) }\SpecialCharTok{|\textgreater{}}
  \FunctionTok{mutate}\NormalTok{(}\AttributeTok{proportion =}\NormalTok{ n }\SpecialCharTok{/} \FunctionTok{sum}\NormalTok{(n)) }\SpecialCharTok{|\textgreater{}}
  \FunctionTok{ungroup}\NormalTok{() }\SpecialCharTok{|\textgreater{}}
  \FunctionTok{mutate}\NormalTok{(}
    \AttributeTok{true\_strategy =} \FunctionTok{factor}\NormalTok{(true\_strategy, }\AttributeTok{levels =} \FunctionTok{c}\NormalTok{(}\StringTok{"SBA"}\NormalTok{, }\StringTok{"WBA"}\NormalTok{)),}
    \AttributeTok{best\_model    =} \FunctionTok{factor}\NormalTok{(best\_model, }\AttributeTok{levels =} \FunctionTok{c}\NormalTok{(}\StringTok{"SBA"}\NormalTok{, }\StringTok{"PBA"}\NormalTok{, }\StringTok{"WBA"}\NormalTok{))}
\NormalTok{  ) }\SpecialCharTok{|\textgreater{}}
  \FunctionTok{complete}\NormalTok{(true\_strategy, best\_model, }\AttributeTok{fill =} \FunctionTok{list}\NormalTok{(}\AttributeTok{n =} \DecValTok{0}\NormalTok{, }\AttributeTok{proportion =} \DecValTok{0}\NormalTok{))}

\CommentTok{\# Table}
\NormalTok{mix\_confusion }\SpecialCharTok{|\textgreater{}}
  \FunctionTok{select}\NormalTok{(}\SpecialCharTok{{-}}\NormalTok{n) }\SpecialCharTok{|\textgreater{}}
  \FunctionTok{pivot\_wider}\NormalTok{(}\AttributeTok{names\_from =}\NormalTok{ best\_model, }\AttributeTok{values\_from =}\NormalTok{ proportion) }\SpecialCharTok{|\textgreater{}}
\NormalTok{  knitr}\SpecialCharTok{::}\FunctionTok{kable}\NormalTok{(}\AttributeTok{digits =} \DecValTok{2}\NormalTok{,}
               \AttributeTok{caption =} \StringTok{"Individual{-}Level Strategy Classification: Mixed Population"}\NormalTok{)}
\end{Highlighting}
\end{Shaded}

\begin{longtable}[]{@{}lrrr@{}}
\caption{Individual-Level Strategy Classification: Mixed
Population}\tabularnewline
\toprule\noalign{}
true\_strategy & SBA & PBA & WBA \\
\midrule\noalign{}
\endfirsthead
\toprule\noalign{}
true\_strategy & SBA & PBA & WBA \\
\midrule\noalign{}
\endhead
\bottomrule\noalign{}
\endlastfoot
SBA & 0.60 & 0.25 & 0.15 \\
WBA & 0.35 & 0.20 & 0.45 \\
\end{longtable}

\begin{Shaded}
\begin{Highlighting}[]
\CommentTok{\# Heatmap}
\FunctionTok{ggplot}\NormalTok{(mix\_confusion, }\FunctionTok{aes}\NormalTok{(}\AttributeTok{x =}\NormalTok{ best\_model, }\AttributeTok{y =}\NormalTok{ true\_strategy, }\AttributeTok{fill =}\NormalTok{ proportion)) }\SpecialCharTok{+}
  \FunctionTok{geom\_tile}\NormalTok{(}\AttributeTok{colour =} \StringTok{"white"}\NormalTok{, }\AttributeTok{linewidth =} \DecValTok{1}\NormalTok{) }\SpecialCharTok{+}
  \FunctionTok{geom\_text}\NormalTok{(}\FunctionTok{aes}\NormalTok{(}\AttributeTok{label =} \FunctionTok{sprintf}\NormalTok{(}\StringTok{"\%d}\SpecialCharTok{\textbackslash{}n}\StringTok{(\%s)"}\NormalTok{, n, scales}\SpecialCharTok{::}\FunctionTok{percent}\NormalTok{(proportion, }\AttributeTok{accuracy =} \DecValTok{1}\NormalTok{))),}
            \AttributeTok{size =} \DecValTok{5}\NormalTok{, }\AttributeTok{fontface =} \StringTok{"bold"}\NormalTok{) }\SpecialCharTok{+}
  \FunctionTok{scale\_fill\_viridis\_c}\NormalTok{(}\AttributeTok{option =} \StringTok{"mako"}\NormalTok{, }\AttributeTok{direction =} \SpecialCharTok{{-}}\DecValTok{1}\NormalTok{, }\AttributeTok{limits =} \FunctionTok{c}\NormalTok{(}\DecValTok{0}\NormalTok{, }\DecValTok{1}\NormalTok{)) }\SpecialCharTok{+}
  \FunctionTok{scale\_y\_discrete}\NormalTok{(}\AttributeTok{limits =}\NormalTok{ rev) }\SpecialCharTok{+}
  \FunctionTok{labs}\NormalTok{(}
    \AttributeTok{title =} \StringTok{"Individual{-}Level Strategy Recovery: Mixed Population"}\NormalTok{,}
    \AttributeTok{subtitle =} \StringTok{"Rows = true generating strategy. Columns = model with highest per{-}agent ELPD."}\NormalTok{,}
    \AttributeTok{x     =} \StringTok{"Assigned Model (Max Agent ELPD)"}\NormalTok{,}
    \AttributeTok{y     =} \StringTok{"True Strategy"}\NormalTok{,}
    \AttributeTok{fill  =} \StringTok{"Classification}\SpecialCharTok{\textbackslash{}n}\StringTok{Rate"}
\NormalTok{  ) }\SpecialCharTok{+}
  \FunctionTok{theme}\NormalTok{(}\AttributeTok{panel.grid =} \FunctionTok{element\_blank}\NormalTok{())}
\end{Highlighting}
\end{Shaded}

\begin{center}
\includegraphics[width=0.8\linewidth,height=\textheight,keepaspectratio]{10-BayesianModels_files/figure-pdf/ch10_mixed_confusion-1.pdf}
\end{center}

\subsection{Visualising the Classification
Boundary}\label{visualising-the-classification-boundary}

What separates correctly classified agents from misclassified ones? We
plot the agents in the estimated \((\rho, \kappa)\) space, marking
classifications.

\begin{Shaded}
\begin{Highlighting}[]
\NormalTok{mix\_class\_data }\OtherTok{\textless{}{-}}\NormalTok{ mix\_recovery }\SpecialCharTok{|\textgreater{}}
  \FunctionTok{left\_join}\NormalTok{(agent\_elpd }\SpecialCharTok{|\textgreater{}} \FunctionTok{select}\NormalTok{(agent\_id, best\_model, elpd\_advantage),}
            \AttributeTok{by =} \StringTok{"agent\_id"}\NormalTok{) }\SpecialCharTok{|\textgreater{}}
  \FunctionTok{mutate}\NormalTok{(}\AttributeTok{correct =}\NormalTok{ (true\_strategy }\SpecialCharTok{==} \StringTok{"SBA"} \SpecialCharTok{\&}\NormalTok{ best\_model }\SpecialCharTok{==} \StringTok{"SBA"}\NormalTok{) }\SpecialCharTok{|}
\NormalTok{                   (true\_strategy }\SpecialCharTok{==} \StringTok{"WBA"} \SpecialCharTok{\&}\NormalTok{ best\_model }\SpecialCharTok{\%in\%} \FunctionTok{c}\NormalTok{(}\StringTok{"PBA"}\NormalTok{, }\StringTok{"WBA"}\NormalTok{)))}

\FunctionTok{ggplot}\NormalTok{(mix\_class\_data,}
       \FunctionTok{aes}\NormalTok{(}\AttributeTok{x =}\NormalTok{ est\_rho, }\AttributeTok{y =}\NormalTok{ est\_kappa,}
           \AttributeTok{colour =}\NormalTok{ true\_strategy, }\AttributeTok{shape =}\NormalTok{ correct)) }\SpecialCharTok{+}
  \FunctionTok{geom\_vline}\NormalTok{(}\AttributeTok{xintercept =} \FloatTok{0.5}\NormalTok{, }\AttributeTok{linetype =} \StringTok{"dotted"}\NormalTok{, }\AttributeTok{colour =} \StringTok{"grey60"}\NormalTok{) }\SpecialCharTok{+}
  \FunctionTok{geom\_hline}\NormalTok{(}\AttributeTok{yintercept =} \FloatTok{2.0}\NormalTok{, }\AttributeTok{linetype =} \StringTok{"dotted"}\NormalTok{, }\AttributeTok{colour =} \StringTok{"grey60"}\NormalTok{) }\SpecialCharTok{+}
  \FunctionTok{geom\_point}\NormalTok{(}\AttributeTok{size =} \FloatTok{3.5}\NormalTok{, }\AttributeTok{alpha =} \FloatTok{0.8}\NormalTok{) }\SpecialCharTok{+}
  \FunctionTok{scale\_colour\_manual}\NormalTok{(}\AttributeTok{values =} \FunctionTok{c}\NormalTok{(}\StringTok{"SBA"} \OtherTok{=} \StringTok{"\#4daf4a"}\NormalTok{, }\StringTok{"WBA"} \OtherTok{=} \StringTok{"\#984ea3"}\NormalTok{)) }\SpecialCharTok{+}
  \FunctionTok{scale\_shape\_manual}\NormalTok{(}\AttributeTok{values =} \FunctionTok{c}\NormalTok{(}\StringTok{"TRUE"} \OtherTok{=} \DecValTok{16}\NormalTok{, }\StringTok{"FALSE"} \OtherTok{=} \DecValTok{4}\NormalTok{),}
                     \AttributeTok{labels =} \FunctionTok{c}\NormalTok{(}\StringTok{"TRUE"} \OtherTok{=} \StringTok{"Correct"}\NormalTok{, }\StringTok{"FALSE"} \OtherTok{=} \StringTok{"Misclassified"}\NormalTok{)) }\SpecialCharTok{+}
  \FunctionTok{labs}\NormalTok{(}
    \AttributeTok{title    =} \StringTok{"Classification in Estimated Parameter Space"}\NormalTok{,}
    \AttributeTok{subtitle =} \StringTok{"Dotted lines show SBA reference values (ρ = 0.5, κ = 2.0). Crosses = misclassified."}\NormalTok{,}
    \AttributeTok{x        =} \FunctionTok{expression}\NormalTok{(}\StringTok{"Estimated "} \SpecialCharTok{*}\NormalTok{ rho[j]),}
    \AttributeTok{y        =} \FunctionTok{expression}\NormalTok{(}\StringTok{"Estimated "} \SpecialCharTok{*}\NormalTok{ kappa[j]),}
    \AttributeTok{colour   =} \StringTok{"True Strategy"}\NormalTok{, }\AttributeTok{shape =} \StringTok{"Classification"}
\NormalTok{  ) }\SpecialCharTok{+}
  \FunctionTok{theme}\NormalTok{(}\AttributeTok{legend.position =} \StringTok{"bottom"}\NormalTok{)}
\end{Highlighting}
\end{Shaded}

\begin{center}
\includegraphics[width=0.8\linewidth,height=\textheight,keepaspectratio]{10-BayesianModels_files/figure-pdf/ch10_mixed_classification_boundary-1.pdf}
\end{center}

Misclassifications, when they occur, concentrate near the boundary
between the two strategy clusters --- agents whose estimated \(\rho\)
sits between 0.35 and 0.5, where the SBA and WBA strategies are hardest
to discriminate.

\subsection{Discussion: What the Mixed-Population Analysis
Reveals}\label{discussion-what-the-mixed-population-analysis-reveals}

This analysis demonstrates three important lessons for applied cognitive
modelling:

\textbf{1. Population parameters are misleading for heterogeneous
groups.} The ML-WBA's estimated population mean \(\rho\) averaged over
two qualitatively different strategies to produce a summary that
describes nobody. This is not a failure of the hierarchical model
\emph{per se} --- it is a structural consequence of fitting a unimodal
population distribution to a bimodal mixture. When you suspect strategy
heterogeneity, do not interpret population means at face value.

\textbf{2. Individual-level recovery succeeds despite population-level
misspecification.} The ML-WBA, though structurally unable to represent a
mixture population, still recovers individual \(\rho_j\) and
\(\kappa_j\) values well enough to separate the two groups. The
estimated population SD absorbs the between-group variance, and the
partial pooling is mild enough that the bimodal structure in the
individual estimates is preserved. A formal mixture model (as in Ch. 9)
would be more appropriate --- it would estimate the mixing proportion
and the within-group population distributions directly --- but the
standard hierarchical model is a surprisingly robust approximation.

\textbf{3. Per-agent model comparison identifies strategy types.} By
decomposing the overall LOO into agent-level contributions, we can
classify individuals as SBA-type or WBA-type with reasonable accuracy.
This connects directly to the mixture modelling approach from Ch. 9:
where a mixture model infers strategy membership \emph{within} the
generative model, per-agent LOO achieves a similar classification
\emph{post hoc} through predictive comparison.

The practical implication for researchers: if you suspect your
participants use different strategies, always look beyond
population-level summaries. Examine individual parameter estimates for
multimodality, and use per-agent model comparison to classify
participants before interpreting group effects.

\section{Beyond Weighted Bayes: Circular
Inference}\label{beyond-weighted-bayes-circular-inference}

The models we have explored so far --- SBA, PBA, WBA --- share a common
structural assumption: each source of evidence is counted exactly once.
Simonsen et al.~(2021), following Jardri et al.~(2017), found that this
assumption is systematically violated in human reasoning ---
particularly in patients with schizophrenia. Their best-fitting model
was \emph{circular inference}: evidence from one source is fed back into
the system and counted multiple times.

\subsection{The Log-Odds Transformation
Function}\label{the-log-odds-transformation-function}

To properly model circular inference, pseudocounts (\(\alpha, \beta\))
are not enough. We must shift our integration into the \textbf{log-odds
space}. If we express our evidence \(L\) as log-odds, Jardri and
colleagues define a specific non-linear transformation function
\(F(L, w)\) that applies a weight \(w \in [0.5, 1]\) to that evidence:

\[F(L,w) = \log \frac{w e^L + 1 - w}{(1 - w)e^L + w}\]

Let's visualize how this function behaves in R.

\begin{Shaded}
\begin{Highlighting}[]
\CommentTok{\# The core Jardri integration function}
\NormalTok{JardriF }\OtherTok{\textless{}{-}} \ControlFlowTok{function}\NormalTok{(L, w) \{}
  \CommentTok{\# L is evidence in log{-}odds}
  \CommentTok{\# w is the weight (0.5 = no effect, 1 = full trust)}
\NormalTok{  num }\OtherTok{\textless{}{-}}\NormalTok{ w }\SpecialCharTok{*} \FunctionTok{exp}\NormalTok{(L) }\SpecialCharTok{+}\NormalTok{ (}\DecValTok{1} \SpecialCharTok{{-}}\NormalTok{ w)}
\NormalTok{  den }\OtherTok{\textless{}{-}}\NormalTok{ (}\DecValTok{1} \SpecialCharTok{{-}}\NormalTok{ w) }\SpecialCharTok{*} \FunctionTok{exp}\NormalTok{(L) }\SpecialCharTok{+}\NormalTok{ w}
  \FunctionTok{return}\NormalTok{(}\FunctionTok{log}\NormalTok{(num }\SpecialCharTok{/}\NormalTok{ den))}
\NormalTok{\}}

\CommentTok{\# Visualize the effect of varying weights}
\FunctionTok{tibble}\NormalTok{(}\AttributeTok{L =} \FunctionTok{seq}\NormalTok{(}\SpecialCharTok{{-}}\DecValTok{4}\NormalTok{, }\DecValTok{4}\NormalTok{, }\AttributeTok{by =} \FloatTok{0.1}\NormalTok{)) }\SpecialCharTok{|\textgreater{}}
  \FunctionTok{expand\_grid}\NormalTok{(}\AttributeTok{w =} \FunctionTok{c}\NormalTok{(}\FloatTok{0.5}\NormalTok{, }\FloatTok{0.6}\NormalTok{, }\FloatTok{0.7}\NormalTok{, }\FloatTok{0.8}\NormalTok{, }\FloatTok{0.9}\NormalTok{, }\FloatTok{1.0}\NormalTok{)) }\SpecialCharTok{|\textgreater{}}
  \FunctionTok{mutate}\NormalTok{(}
    \AttributeTok{F\_L =} \FunctionTok{JardriF}\NormalTok{(L, w),}
    \AttributeTok{Prob =} \FunctionTok{plogis}\NormalTok{(F\_L) }\CommentTok{\# Convert log{-}odds back to probability}
\NormalTok{  ) }\SpecialCharTok{|\textgreater{}}
  \FunctionTok{ggplot}\NormalTok{(}\FunctionTok{aes}\NormalTok{(}\AttributeTok{x =}\NormalTok{ L, }\AttributeTok{y =}\NormalTok{ Prob, }\AttributeTok{color =} \FunctionTok{factor}\NormalTok{(w))) }\SpecialCharTok{+}
  \FunctionTok{geom\_line}\NormalTok{(}\AttributeTok{linewidth =} \DecValTok{1}\NormalTok{) }\SpecialCharTok{+}
  \FunctionTok{geom\_hline}\NormalTok{(}\AttributeTok{yintercept =} \FloatTok{0.5}\NormalTok{, }\AttributeTok{linetype =} \StringTok{"dashed"}\NormalTok{, }\AttributeTok{alpha =} \FloatTok{0.5}\NormalTok{) }\SpecialCharTok{+}
  \FunctionTok{geom\_vline}\NormalTok{(}\AttributeTok{xintercept =} \DecValTok{0}\NormalTok{, }\AttributeTok{linetype =} \StringTok{"dashed"}\NormalTok{, }\AttributeTok{alpha =} \FloatTok{0.5}\NormalTok{) }\SpecialCharTok{+}
  \FunctionTok{scale\_color\_viridis\_d}\NormalTok{(}\AttributeTok{option =} \StringTok{"plasma"}\NormalTok{, }\AttributeTok{end =} \FloatTok{0.9}\NormalTok{) }\SpecialCharTok{+}
  \FunctionTok{labs}\NormalTok{(}
    \AttributeTok{title =} \StringTok{"Jardri Transformation Function F(L, w)"}\NormalTok{,}
    \AttributeTok{subtitle =} \StringTok{"Higher weights approach the true evidence (S{-}curve), w=0.5 ignores it (flat)."}\NormalTok{,}
    \AttributeTok{x =} \StringTok{"Evidence (L) in Log{-}Odds"}\NormalTok{,}
    \AttributeTok{y =} \StringTok{"Probability"}\NormalTok{,}
    \AttributeTok{color =} \StringTok{"Weight (w)"}
\NormalTok{  )}
\end{Highlighting}
\end{Shaded}

\begin{center}
\includegraphics[width=0.8\linewidth,height=\textheight,keepaspectratio]{10-BayesianModels_files/figure-pdf/ch10_jardri_function-1.pdf}
\end{center}

\subsection{Building the Models
Pedagogically}\label{building-the-models-pedagogically}

We will now build the Circular Inference family of models in three
stages, restricted to two sources of information: Direct Evidence
(\(L_1\)) and Social Evidence (\(L_2\)).

\subsubsection{Stage 1: Weights Only (Weighted Bayes in
Log-Odds)}\label{stage-1-weights-only-weighted-bayes-in-log-odds}

If an agent only applies weights but does not overcount evidence, the
posterior log-odds is simply the sum of the transformed
evidence:\[L_{post} = F(L_1, w_1) + F(L_2, w_2)\]

\subsubsection{Stage 2: Weights and Loops (No
Interference)}\label{stage-2-weights-and-loops-no-interference}

Circular inference claims that information reverberates in the network.
We model this by multiplying the evidence by a loop parameter \(\alpha\)
before the transformation. An \(\alpha > 1\) represents
overcounting.\[L_{post} = F(\alpha_1 L_1, w_1) + F(\alpha_2 L_2, w_2)\]

\subsubsection{Stage 3: Weights, Loops, and Interference (Full Circular
Inference)}\label{stage-3-weights-loops-and-interference-full-circular-inference}

The complete Jardri model notes that overcounted evidence causes
interference. The reverberated information corrupts the original
sources, rendering them non-separable. We first calculate the total
interference \(I\), and then add it to the base evidence inside the
transformation
function:\[I = F(\alpha_1 L_1, w_1) + F(\alpha_2 L_2, w_2)\]\[L_{post} = F(L_1 + I, w_1) + F(L_2 + I, w_2)\]Let's
write generative R functions for all three stages.

\begin{Shaded}
\begin{Highlighting}[]
\CommentTok{\# Stage 1: Weights Only}
\NormalTok{model\_weights\_only }\OtherTok{\textless{}{-}} \ControlFlowTok{function}\NormalTok{(L1, L2, w1, w2) \{}
\NormalTok{  L\_post }\OtherTok{\textless{}{-}} \FunctionTok{JardriF}\NormalTok{(L1, w1) }\SpecialCharTok{+} \FunctionTok{JardriF}\NormalTok{(L2, w2)}
  \FunctionTok{return}\NormalTok{(}\FunctionTok{plogis}\NormalTok{(L\_post))}
\NormalTok{\}}

\CommentTok{\# Stage 2: Weights and Loops}
\NormalTok{model\_loops }\OtherTok{\textless{}{-}} \ControlFlowTok{function}\NormalTok{(L1, L2, w1, w2, alpha1, alpha2) \{}
\NormalTok{  L\_post }\OtherTok{\textless{}{-}} \FunctionTok{JardriF}\NormalTok{(alpha1 }\SpecialCharTok{*}\NormalTok{ L1, w1) }\SpecialCharTok{+} \FunctionTok{JardriF}\NormalTok{(alpha2 }\SpecialCharTok{*}\NormalTok{ L2, w2)}
  \FunctionTok{return}\NormalTok{(}\FunctionTok{plogis}\NormalTok{(L\_post))}
\NormalTok{\}}

\CommentTok{\# Stage 3: Weights, Loops, and Interference}
\NormalTok{model\_interference }\OtherTok{\textless{}{-}} \ControlFlowTok{function}\NormalTok{(L1, L2, w1, w2, alpha1, alpha2) \{}
\NormalTok{  I }\OtherTok{\textless{}{-}} \FunctionTok{JardriF}\NormalTok{(alpha1 }\SpecialCharTok{*}\NormalTok{ L1, w1) }\SpecialCharTok{+} \FunctionTok{JardriF}\NormalTok{(alpha2 }\SpecialCharTok{*}\NormalTok{ L2, w2)}
\NormalTok{  L\_post }\OtherTok{\textless{}{-}} \FunctionTok{JardriF}\NormalTok{(L1 }\SpecialCharTok{+}\NormalTok{ I, w1) }\SpecialCharTok{+} \FunctionTok{JardriF}\NormalTok{(L2 }\SpecialCharTok{+}\NormalTok{ I, w2)}
  \FunctionTok{return}\NormalTok{(}\FunctionTok{plogis}\NormalTok{(L\_post))}
\NormalTok{\}}
\end{Highlighting}
\end{Shaded}

\subsection{Visualizing and Testing the Three
Models}\label{visualizing-and-testing-the-three-models}

How do these structural differences manifest in behaviour? To test this,
let's simulate how an agent integrates Direct Evidence (\(L_1\)) and
Social Evidence (\(L_2\)) under each of the three architectures.We will
give the agent realistic parameters: moderate trust in themselves
(\(w_1 = 0.85\)) and moderate trust in others (\(w_2 = 0.70\)). For the
circular models, we assume they heavily overcount their own direct
experience (\(\alpha_1 = 3.5\)) and slightly overcount social
information (\(\alpha_2 = 1.5\)).

\begin{Shaded}
\begin{Highlighting}[]
\CommentTok{\# Convert counts to log odds mapping (avoiding strict 0 and 1)}
\NormalTok{to\_log\_odds }\OtherTok{\textless{}{-}} \ControlFlowTok{function}\NormalTok{(blue, total) \{}
  \CommentTok{\# Add small epsilon to prevent infinite log{-}odds}
\NormalTok{  p }\OtherTok{\textless{}{-}}\NormalTok{ (blue }\SpecialCharTok{+} \FloatTok{0.5}\NormalTok{) }\SpecialCharTok{/}\NormalTok{ (total }\SpecialCharTok{+} \DecValTok{1}\NormalTok{)}
  \FunctionTok{return}\NormalTok{(}\FunctionTok{log}\NormalTok{(p }\SpecialCharTok{/}\NormalTok{ (}\DecValTok{1} \SpecialCharTok{{-}}\NormalTok{ p)))}
\NormalTok{\}}

\CommentTok{\# Create a testing grid}
\NormalTok{ci\_test\_data }\OtherTok{\textless{}{-}} \FunctionTok{expand\_grid}\NormalTok{(}
  \AttributeTok{blue1 =} \DecValTok{0}\SpecialCharTok{:}\DecValTok{8}\NormalTok{,}
  \AttributeTok{blue2 =} \DecValTok{0}\SpecialCharTok{:}\DecValTok{3}
\NormalTok{) }\SpecialCharTok{|\textgreater{}}
  \FunctionTok{mutate}\NormalTok{(}
    \AttributeTok{L1 =} \FunctionTok{to\_log\_odds}\NormalTok{(blue1, }\DecValTok{8}\NormalTok{),}
    \AttributeTok{L2 =} \FunctionTok{to\_log\_odds}\NormalTok{(blue2, }\DecValTok{3}\NormalTok{),}
    \CommentTok{\# Agent parameters}
    \AttributeTok{w1 =} \FloatTok{0.85}\NormalTok{, }\AttributeTok{w2 =} \FloatTok{0.70}\NormalTok{,}
    \AttributeTok{a1 =} \FloatTok{3.5}\NormalTok{, }\AttributeTok{a2 =} \FloatTok{1.5}
\NormalTok{  ) }\SpecialCharTok{|\textgreater{}}
  \CommentTok{\# Generate predictions for all three models}
  \FunctionTok{mutate}\NormalTok{(}
    \StringTok{\textasciigrave{}}\AttributeTok{1. Weights Only}\StringTok{\textasciigrave{}} \OtherTok{=} \FunctionTok{pmap\_dbl}\NormalTok{(}\FunctionTok{list}\NormalTok{(L1, L2, w1, w2), model\_weights\_only),}
    \StringTok{\textasciigrave{}}\AttributeTok{2. Weights + Loops}\StringTok{\textasciigrave{}} \OtherTok{=} \FunctionTok{pmap\_dbl}\NormalTok{(}\FunctionTok{list}\NormalTok{(L1, L2, w1, w2, a1, a2), model\_loops),}
    \StringTok{\textasciigrave{}}\AttributeTok{3. Loops + Interference}\StringTok{\textasciigrave{}} \OtherTok{=} \FunctionTok{pmap\_dbl}\NormalTok{(}\FunctionTok{list}\NormalTok{(L1, L2, w1, w2, a1, a2), model\_interference),}
    \AttributeTok{social\_label =} \FunctionTok{factor}\NormalTok{(blue2, }\AttributeTok{levels =} \DecValTok{0}\SpecialCharTok{:}\DecValTok{3}\NormalTok{, }
                          \AttributeTok{labels =} \FunctionTok{c}\NormalTok{(}\StringTok{"Clear Red"}\NormalTok{, }\StringTok{"Maybe Red"}\NormalTok{, }\StringTok{"Maybe Blue"}\NormalTok{, }\StringTok{"Clear Blue"}\NormalTok{))}
\NormalTok{  ) }\SpecialCharTok{|\textgreater{}}
  \CommentTok{\# FIX: Explicitly select all three model columns}
  \FunctionTok{pivot\_longer}\NormalTok{(}
    \AttributeTok{cols =} \FunctionTok{c}\NormalTok{(}\StringTok{\textasciigrave{}}\AttributeTok{1. Weights Only}\StringTok{\textasciigrave{}}\NormalTok{, }\StringTok{\textasciigrave{}}\AttributeTok{2. Weights + Loops}\StringTok{\textasciigrave{}}\NormalTok{, }\StringTok{\textasciigrave{}}\AttributeTok{3. Loops + Interference}\StringTok{\textasciigrave{}}\NormalTok{), }
    \AttributeTok{names\_to =} \StringTok{"Model"}\NormalTok{, }
    \AttributeTok{values\_to =} \StringTok{"P\_Blue"}
\NormalTok{  )}

\CommentTok{\# Plot the comparison}
\FunctionTok{ggplot}\NormalTok{(ci\_test\_data, }\FunctionTok{aes}\NormalTok{(}\AttributeTok{x =}\NormalTok{ blue1, }\AttributeTok{y =}\NormalTok{ P\_Blue, }\AttributeTok{color =}\NormalTok{ social\_label)) }\SpecialCharTok{+}
  \FunctionTok{geom\_line}\NormalTok{(}\AttributeTok{linewidth =} \FloatTok{1.2}\NormalTok{) }\SpecialCharTok{+}
  \FunctionTok{geom\_point}\NormalTok{(}\AttributeTok{size =} \DecValTok{2}\NormalTok{) }\SpecialCharTok{+}
  \FunctionTok{facet\_wrap}\NormalTok{(}\SpecialCharTok{\textasciitilde{}}\NormalTok{Model) }\SpecialCharTok{+}
  \FunctionTok{geom\_hline}\NormalTok{(}\AttributeTok{yintercept =} \FloatTok{0.5}\NormalTok{, }\AttributeTok{linetype =} \StringTok{"dashed"}\NormalTok{, }\AttributeTok{alpha =} \FloatTok{0.5}\NormalTok{) }\SpecialCharTok{+}
  \FunctionTok{scale\_color\_manual}\NormalTok{(}\AttributeTok{values =} \FunctionTok{c}\NormalTok{(}\StringTok{"\#d73027"}\NormalTok{, }\StringTok{"\#fdae61"}\NormalTok{, }\StringTok{"\#abd9e9"}\NormalTok{, }\StringTok{"\#4575b4"}\NormalTok{)) }\SpecialCharTok{+}
  \FunctionTok{scale\_x\_continuous}\NormalTok{(}\AttributeTok{breaks =} \DecValTok{0}\SpecialCharTok{:}\DecValTok{8}\NormalTok{) }\SpecialCharTok{+}
  \FunctionTok{labs}\NormalTok{(}
    \AttributeTok{title =} \StringTok{"Testing Circular Inference Architectures"}\NormalTok{,}
    \AttributeTok{subtitle =} \StringTok{"Notice how loops create steep, overconfident \textquotesingle{}S{-}curves\textquotesingle{}, while interference compresses them together."}\NormalTok{,}
    \AttributeTok{x =} \StringTok{"Direct Evidence (Blue Marbles out of 8)"}\NormalTok{,}
    \AttributeTok{y =} \StringTok{"P(Choose Blue)"}\NormalTok{,}
    \AttributeTok{color =} \StringTok{"Social Evidence"}
\NormalTok{  ) }\SpecialCharTok{+}
  \FunctionTok{theme}\NormalTok{(}\AttributeTok{legend.position =} \StringTok{"bottom"}\NormalTok{)}
\end{Highlighting}
\end{Shaded}

\begin{center}
\includegraphics[width=0.8\linewidth,height=\textheight,keepaspectratio]{10-BayesianModels_files/figure-pdf/ch10_evidence_integration_sim-1.pdf}
\end{center}

\subsection{Discussion: Interpreting the
Visualizations}\label{discussion-interpreting-the-visualizations}

\begin{enumerate}
\def\labelenumi{\arabic{enumi}.}
\item
  Weights Only: This model resembles our earlier log-odds integration
  plots. The agent smoothly transitions from choosing Red to Blue based
  on their direct sample, and the Social Evidence simply shifts the
  curve up or down linearly in log-odds space.
\item
  Weights + Loops: By adding an \(\alpha_1 = 3.5\) loop to direct
  evidence, the agent overcounts their own observations. This results in
  a much steeper curve in the center. A minor shift in direct evidence
  triggers extreme overconfidence, causing the agent to `jump to
  conclusions'.
\item
  Loops + Interference: Adding the interference term means that the
  reverberating direct evidence structurally corrupts the social
  evidence. You can see the curves tightly bunch together at the
  extremes because the massive internal echo of the direct evidence
  completely drowns out the social information.This final model
  (Interference) is what Simonsen et al.~(2021) found best fit the
  behaviour of patients with schizophrenia, explaining both their
  tendency to ignore others and their high confidence in subjective
  experiences.
\end{enumerate}

\subsection{Stan Implementation: The Circular Inference
Family}\label{stan-implementation-the-circular-inference-family}

To translate our three pedagogical models into Stan, we need to handle a
few new structural elements compared to the Beta-Binomial models we used
earlier in the chapter:

\begin{enumerate}
\def\labelenumi{\arabic{enumi}.}
\item
  \textbf{The \texttt{functions} block:} We must define our custom
  non-linear \texttt{jardri\_f} transformation function natively in Stan
  so we can call it repeatedly.
\item
  \textbf{The \texttt{transformed\ data} block:} The Jardri function
  operates strictly in log-odds space. Instead of converting raw marble
  counts to log-odds on every single Markov Chain iteration, we can do
  this once up front to drastically speed up sampling.
\item
  \textbf{The Parameter Bounds:} Following the methodology in Simonsen
  et al.~(2021), weights (\(w\)) are constrained between \texttt{0.5}
  (completely ignoring the source) and \texttt{1.0} (taking it at face
  value). Loops (\(\alpha\)) are strictly \(\ge 1\), which we enforce by
  sampling an excess parameter \texttt{alpha\_m1} \(\ge 0\).
\end{enumerate}

Let's build the three models.

\subsubsection{1. Stage 1: Weights Only}\label{stage-1-weights-only}

This is our baseline. It strictly integrates the evidence without any
reverberation (overcounting) or interference.

\begin{Shaded}
\begin{Highlighting}[]
\NormalTok{CI\_Weights\_stan }\OtherTok{\textless{}{-}} \StringTok{"}
\StringTok{functions \{}
\StringTok{  // The core Jardri transformation function}
\StringTok{  real jardri\_f(real L, real w) \{}
\StringTok{    real num = w * exp(L) + (1.0 {-} w);}
\StringTok{    real den = (1.0 {-} w) * exp(L) + w;}
\StringTok{    return log(num / den);}
\StringTok{  \}}
\StringTok{\}}

\StringTok{data \{}
\StringTok{  int\textless{}lower=1\textgreater{} N\_total;}
\StringTok{  array[N\_total] int\textless{}lower=0, upper=1\textgreater{} y;}
\StringTok{  array[N\_total] int\textless{}lower=0\textgreater{} blue1;}
\StringTok{  array[N\_total] int\textless{}lower=0\textgreater{} total1;}
\StringTok{  array[N\_total] int\textless{}lower=0\textgreater{} blue2;}
\StringTok{  array[N\_total] int\textless{}lower=0\textgreater{} total2;}

\StringTok{  real prior\_w\_mu;}
\StringTok{  real\textless{}lower=0\textgreater{} prior\_w\_sigma;}
\StringTok{  int\textless{}lower=0, upper=1\textgreater{} run\_diagnostics;}
\StringTok{\}}

\StringTok{transformed data \{}
\StringTok{  // Convert raw counts to log{-}odds once to save computation}
\StringTok{  real alpha0 = 0.5;}
\StringTok{  real beta0  = 0.5;}
\StringTok{  array[N\_total] real lo\_self;}
\StringTok{  array[N\_total] real lo\_social;}

\StringTok{  for (i in 1:N\_total) \{}
\StringTok{    lo\_self[i]   = log((alpha0 + blue1[i]) / (beta0 + total1[i] {-} blue1[i]));}
\StringTok{    lo\_social[i] = log((alpha0 + blue2[i]) / (beta0 + total2[i] {-} blue2[i]));}
\StringTok{  \}}
\StringTok{\}}

\StringTok{parameters \{}
\StringTok{  // Scaled between 0.5 (no trust) and 1.0 (full trust)}
\StringTok{  real\textless{}lower=0.5, upper=1.0\textgreater{} w\_self;}
\StringTok{  real\textless{}lower=0.5, upper=1.0\textgreater{} w\_other;}
\StringTok{\}}

\StringTok{model \{}
\StringTok{  // Weakly informative priors centered between 0.5 and 1.0}
\StringTok{  target += normal\_lpdf(w\_self | prior\_w\_mu, prior\_w\_sigma);}
\StringTok{  target += normal\_lpdf(w\_other | prior\_w\_mu, prior\_w\_sigma);}

\StringTok{  for (i in 1:N\_total) \{}
\StringTok{    real L\_post = jardri\_f(lo\_self[i], w\_self) + }
\StringTok{                  jardri\_f(lo\_social[i], w\_other);}
\StringTok{    target += bernoulli\_logit\_lpmf(y[i] | L\_post);}
\StringTok{  \}}
\StringTok{\}}

\StringTok{generated quantities \{}
\StringTok{  vector[N\_total] log\_lik;}
\StringTok{  array[N\_total] int y\_rep;}

\StringTok{  if (run\_diagnostics) \{}
\StringTok{    for (i in 1:N\_total) \{}
\StringTok{      real L\_post = jardri\_f(lo\_self[i], w\_self) + }
\StringTok{                    jardri\_f(lo\_social[i], w\_other);}
\StringTok{      log\_lik[i] = bernoulli\_logit\_lpmf(y[i] | L\_post);}
\StringTok{      y\_rep[i]   = bernoulli\_logit\_rng(L\_post);}
\StringTok{    \}}
\StringTok{  \}}
\StringTok{\}}
\StringTok{"}


\FunctionTok{write\_stan\_file}\NormalTok{(CI\_Weights\_stan, }\AttributeTok{dir =} \StringTok{"stan/"}\NormalTok{, }\AttributeTok{basename =} \StringTok{"ch10\_ci\_weights\_only.stan"}\NormalTok{)}
\end{Highlighting}
\end{Shaded}

\begin{verbatim}
[1] "/Users/au209589/Dropbox/Teaching/AdvancedCognitiveModeling23_book/stan/ch10_ci_weights_only.stan"
\end{verbatim}

\subsubsection{2. Stage 2: Weights + Loops}\label{stage-2-weights-loops}

Now we introduce the loop parameters (\(\alpha\)). We estimate the
excess over 1 (alpha\_m1) using a half-normal prior (since lower=0).
This creates a gentle, scientifically conservative shrinkage that
assumes the agent does not overcount unless the data explicitly demands
it.

\begin{Shaded}
\begin{Highlighting}[]
\NormalTok{CI\_Loops\_stan }\OtherTok{\textless{}{-}} \StringTok{"}
\StringTok{functions \{}
\StringTok{  real jardri\_f(real L, real w) \{}
\StringTok{    real num = w * exp(L) + (1.0 {-} w);}
\StringTok{    real den = (1.0 {-} w) * exp(L) + w;}
\StringTok{    return log(num / den);}
\StringTok{  \}}
\StringTok{\}}

\StringTok{data \{}
\StringTok{  int\textless{}lower=1\textgreater{} N;}
\StringTok{  array[N] int\textless{}lower=0, upper=1\textgreater{} choice;}
\StringTok{  array[N] int\textless{}lower=0\textgreater{} blue1;}
\StringTok{  array[N] int\textless{}lower=0\textgreater{} total1;}
\StringTok{  array[N] int\textless{}lower=0\textgreater{} blue2;}
\StringTok{  array[N] int\textless{}lower=0\textgreater{} total2;}
\StringTok{\}}

\StringTok{transformed data \{}
\StringTok{  real alpha0 = 0.5;}
\StringTok{  real beta0  = 0.5;}
\StringTok{  array[N] real lo\_self;}
\StringTok{  array[N] real lo\_social;}
\StringTok{  }
\StringTok{  for (i in 1:N) \{}
\StringTok{    lo\_self[i]   = log((alpha0 + blue1[i]) / (beta0 + total1[i] {-} blue1[i]));}
\StringTok{    lo\_social[i] = log((alpha0 + blue2[i]) / (beta0 + total2[i] {-} blue2[i]));}
\StringTok{  \}}
\StringTok{\}}

\StringTok{parameters \{}
\StringTok{  real\textless{}lower=0.5, upper=1.0\textgreater{} w\_self;}
\StringTok{  real\textless{}lower=0.5, upper=1.0\textgreater{} w\_other;}
\StringTok{  real\textless{}lower=0\textgreater{} alpha\_self\_m1;    // Excess over 1}
\StringTok{  real\textless{}lower=0\textgreater{} alpha\_other\_m1;}
\StringTok{\}}

\StringTok{transformed parameters \{}
\StringTok{  // Reconstruct the actual loop multiplier}
\StringTok{  real\textless{}lower=1\textgreater{} alpha\_self  = 1.0 + alpha\_self\_m1;}
\StringTok{  real\textless{}lower=1\textgreater{} alpha\_other = 1.0 + alpha\_other\_m1;}
\StringTok{\}}

\StringTok{model \{}
\StringTok{  w\_self  \textasciitilde{} normal(0.75, 0.25);}
\StringTok{  w\_other \textasciitilde{} normal(0.75, 0.25);}
\StringTok{  // Half{-}normal priors shrinking towards 1 (no overcounting)}
\StringTok{  alpha\_self\_m1  \textasciitilde{} std\_normal();}
\StringTok{  alpha\_other\_m1 \textasciitilde{} std\_normal();}
\StringTok{  }
\StringTok{  for (i in 1:N) \{}
\StringTok{    // Evidence is multiplied by the loop parameter BEFORE weighting}
\StringTok{    real L\_post = jardri\_f(alpha\_self * lo\_self[i], w\_self) + }
\StringTok{                  jardri\_f(alpha\_other * lo\_social[i], w\_other);}
\StringTok{                  }
\StringTok{    target += bernoulli\_logit\_lpmf(choice[i] | L\_post);}
\StringTok{  \}}
\StringTok{\}}

\StringTok{generated quantities \{}
\StringTok{  vector[N] log\_lik;}
\StringTok{  array[N] int posterior\_pred;}
\StringTok{  }
\StringTok{  for (i in 1:N) \{}
\StringTok{    real L\_post = jardri\_f(alpha\_self * lo\_self[i], w\_self) + }
\StringTok{                  jardri\_f(alpha\_other * lo\_social[i], w\_other);}
\StringTok{                  }
\StringTok{    log\_lik[i]        = bernoulli\_logit\_lpmf(choice[i] | L\_post);}
\StringTok{    posterior\_pred[i] = bernoulli\_logit\_rng(L\_post);}
\StringTok{  \}}
\StringTok{\}}
\StringTok{"}

\FunctionTok{write\_stan\_file}\NormalTok{(CI\_Loops\_stan, }\AttributeTok{dir =} \StringTok{"stan/"}\NormalTok{, }\AttributeTok{basename =} \StringTok{"ch10\_ci\_loops.stan"}\NormalTok{)}
\end{Highlighting}
\end{Shaded}

\begin{verbatim}
[1] "/Users/au209589/Dropbox/Teaching/AdvancedCognitiveModeling23_book/stan/ch10_ci_loops.stan"
\end{verbatim}

\subsection{3. Stage 3: Weights + Loops +
Interference}\label{stage-3-weights-loops-interference}

Finally, the full model. This exactly maps to the math we established:
the total reverberated evidence (I) is calculated first, and then it is
added back into the base evidence for a second pass through the
weighting transformation.

\begin{Shaded}
\begin{Highlighting}[]
\NormalTok{CI\_Interference\_stan }\OtherTok{\textless{}{-}} \StringTok{"}
\StringTok{functions \{}
\StringTok{  real jardri\_f(real L, real w) \{}
\StringTok{    real num = w * exp(L) + (1.0 {-} w);}
\StringTok{    real den = (1.0 {-} w) * exp(L) + w;}
\StringTok{    return log(num / den);}
\StringTok{  \}}
\StringTok{\}}

\StringTok{data \{}
\StringTok{  int\textless{}lower=1\textgreater{} N;}
\StringTok{  array[N] int\textless{}lower=0, upper=1\textgreater{} choice;}
\StringTok{  array[N] int\textless{}lower=0\textgreater{} blue1;}
\StringTok{  array[N] int\textless{}lower=0\textgreater{} total1;}
\StringTok{  array[N] int\textless{}lower=0\textgreater{} blue2;}
\StringTok{  array[N] int\textless{}lower=0\textgreater{} total2;}
\StringTok{\}}

\StringTok{transformed data \{}
\StringTok{  real alpha0 = 0.5;}
\StringTok{  real beta0  = 0.5;}
\StringTok{  array[N] real lo\_self;}
\StringTok{  array[N] real lo\_social;}
\StringTok{  }
\StringTok{  for (i in 1:N) \{}
\StringTok{    lo\_self[i]   = log((alpha0 + blue1[i]) / (beta0 + total1[i] {-} blue1[i]));}
\StringTok{    lo\_social[i] = log((alpha0 + blue2[i]) / (beta0 + total2[i] {-} blue2[i]));}
\StringTok{  \}}
\StringTok{\}}

\StringTok{parameters \{}
\StringTok{  real\textless{}lower=0.5, upper=1.0\textgreater{} w\_self;}
\StringTok{  real\textless{}lower=0.5, upper=1.0\textgreater{} w\_other;}
\StringTok{  real\textless{}lower=0\textgreater{} alpha\_self\_m1;    }
\StringTok{  real\textless{}lower=0\textgreater{} alpha\_other\_m1;}
\StringTok{\}}

\StringTok{transformed parameters \{}
\StringTok{  real\textless{}lower=1\textgreater{} alpha\_self  = 1.0 + alpha\_self\_m1;}
\StringTok{  real\textless{}lower=1\textgreater{} alpha\_other = 1.0 + alpha\_other\_m1;}
\StringTok{\}}

\StringTok{model \{}
\StringTok{  w\_self  \textasciitilde{} normal(0.75, 0.25);}
\StringTok{  w\_other \textasciitilde{} normal(0.75, 0.25);}
\StringTok{  alpha\_self\_m1  \textasciitilde{} std\_normal();}
\StringTok{  alpha\_other\_m1 \textasciitilde{} std\_normal();}
\StringTok{  }
\StringTok{  for (i in 1:N) \{}
\StringTok{    // Calculate Total Interference first}
\StringTok{    real I = jardri\_f(alpha\_self * lo\_self[i], w\_self) + }
\StringTok{             jardri\_f(alpha\_other * lo\_social[i], w\_other);}
\StringTok{             }
\StringTok{    // Add Interference to base evidence, then weight}
\StringTok{    real L\_post = jardri\_f(lo\_self[i] + I, w\_self) + }
\StringTok{                  jardri\_f(lo\_social[i] + I, w\_other);}
\StringTok{                  }
\StringTok{    target += bernoulli\_logit\_lpmf(choice[i] | L\_post);}
\StringTok{  \}}
\StringTok{\}}

\StringTok{generated quantities \{}
\StringTok{  vector[N] log\_lik;}
\StringTok{  array[N] int posterior\_pred;}
\StringTok{  }
\StringTok{  for (i in 1:N) \{}
\StringTok{    real I = jardri\_f(alpha\_self * lo\_self[i], w\_self) + }
\StringTok{             jardri\_f(alpha\_other * lo\_social[i], w\_other);}
\StringTok{             }
\StringTok{    real L\_post = jardri\_f(lo\_self[i] + I, w\_self) + }
\StringTok{                  jardri\_f(lo\_social[i] + I, w\_other);}
\StringTok{                  }
\StringTok{    log\_lik[i]        = bernoulli\_logit\_lpmf(choice[i] | L\_post);}
\StringTok{    posterior\_pred[i] = bernoulli\_logit\_rng(L\_post);}
\StringTok{  \}}
\StringTok{\}}
\StringTok{"}

\FunctionTok{write\_stan\_file}\NormalTok{(CI\_Interference\_stan, }\AttributeTok{dir =} \StringTok{"stan/"}\NormalTok{, }\AttributeTok{basename =} \StringTok{"ch10\_ci\_interference.stan"}\NormalTok{)}
\end{Highlighting}
\end{Shaded}

\begin{verbatim}
[1] "/Users/au209589/Dropbox/Teaching/AdvancedCognitiveModeling23_book/stan/ch10_ci_interference.stan"
\end{verbatim}

\subsection{Simulation-Based Calibration Checks (SBC): Pinpointing the
Source of
Complexity}\label{simulation-based-calibration-checks-sbc-pinpointing-the-source-of-complexity}

With our three models written natively in Stan, we face a significant
computational challenge. The \texttt{jardri\_f} transformation is highly
non-linear. Furthermore, the loop parameters (\(\alpha\)) multiply the
evidence \emph{before} it passes through this transformation, and the
interference model creates feedback loops between the two sources.

When parameters interact multiplicatively inside non-linear bounds,
\textbf{identifiability issues} are almost guaranteed. To see exactly
where these problems emerge, we will run Simulation-Based Calibration
Checks (SBC) on all three models in our pedagogical sequence: 1. Weights
Only 2. Weights + Loops 3. Full Interference

We will simulate datasets from the priors of each model, fit the
corresponding Stan model to them, and compare their calibration and
parameter recovery side-by-side.

\begin{Shaded}
\begin{Highlighting}[]
\CommentTok{\# Ensure all models are compiled}
\NormalTok{mod\_ci\_weights      }\OtherTok{\textless{}{-}} \FunctionTok{cmdstan\_model}\NormalTok{(}\StringTok{"stan/ch10\_ci\_weights\_only.stan"}\NormalTok{, }\AttributeTok{dir =} \StringTok{"simmodels"}\NormalTok{)}
\NormalTok{mod\_ci\_loops        }\OtherTok{\textless{}{-}} \FunctionTok{cmdstan\_model}\NormalTok{(}\StringTok{"stan/ch10\_ci\_loops.stan"}\NormalTok{, }\AttributeTok{dir =} \StringTok{"simmodels"}\NormalTok{)}
\NormalTok{mod\_ci\_interference }\OtherTok{\textless{}{-}} \FunctionTok{cmdstan\_model}\NormalTok{(}\StringTok{"stan/ch10\_ci\_interference.stan"}\NormalTok{, }\AttributeTok{dir =} \StringTok{"simmodels"}\NormalTok{)}

\ControlFlowTok{if}\NormalTok{ (regenerate\_simulations }\SpecialCharTok{||} \SpecialCharTok{!}\FunctionTok{file.exists}\NormalTok{(}\StringTok{"simmodels/ch10\_sbc\_ci\_all\_results.rds"}\NormalTok{)) \{}
  
  \CommentTok{\# 1. Define the R equivalent of the Jardri function for the generators}
\NormalTok{  jardri\_f }\OtherTok{\textless{}{-}} \ControlFlowTok{function}\NormalTok{(L, w) \{}
\NormalTok{    num }\OtherTok{\textless{}{-}}\NormalTok{ w }\SpecialCharTok{*} \FunctionTok{exp}\NormalTok{(L) }\SpecialCharTok{+}\NormalTok{ (}\FloatTok{1.0} \SpecialCharTok{{-}}\NormalTok{ w)}
\NormalTok{    den }\OtherTok{\textless{}{-}}\NormalTok{ (}\FloatTok{1.0} \SpecialCharTok{{-}}\NormalTok{ w) }\SpecialCharTok{*} \FunctionTok{exp}\NormalTok{(L) }\SpecialCharTok{+}\NormalTok{ w}
    \FunctionTok{return}\NormalTok{(}\FunctionTok{log}\NormalTok{(num }\SpecialCharTok{/}\NormalTok{ den))}
\NormalTok{  \}}

  \CommentTok{\# Helper to sample from truncated normal [0.5, 1.0] for weights}
\NormalTok{  draw\_w }\OtherTok{\textless{}{-}} \ControlFlowTok{function}\NormalTok{() \{}
\NormalTok{    w }\OtherTok{\textless{}{-}} \SpecialCharTok{{-}}\DecValTok{1}
    \ControlFlowTok{while}\NormalTok{(w }\SpecialCharTok{\textless{}} \FloatTok{0.5} \SpecialCharTok{||}\NormalTok{ w }\SpecialCharTok{\textgreater{}} \FloatTok{1.0}\NormalTok{) w }\OtherTok{\textless{}{-}} \FunctionTok{rnorm}\NormalTok{(}\DecValTok{1}\NormalTok{, }\FloatTok{0.75}\NormalTok{, }\FloatTok{0.25}\NormalTok{)}
    \FunctionTok{return}\NormalTok{(w)}
\NormalTok{  \}}

\NormalTok{  N\_sbc\_trials }\OtherTok{\textless{}{-}} \FunctionTok{nrow}\NormalTok{(evidence\_combinations) }\SpecialCharTok{*} \DecValTok{5} \CommentTok{\# 180 trials per simulated agent}

  \CommentTok{\# {-}{-}{-} GENERATOR 1: Weights Only {-}{-}{-}}
\NormalTok{  sbc\_gen\_weights }\OtherTok{\textless{}{-}} \FunctionTok{SBC\_generator\_function}\NormalTok{(}
    \ControlFlowTok{function}\NormalTok{(N) \{}
\NormalTok{      ws }\OtherTok{\textless{}{-}} \FunctionTok{draw\_w}\NormalTok{()}
\NormalTok{      wo }\OtherTok{\textless{}{-}} \FunctionTok{draw\_w}\NormalTok{()}
\NormalTok{      ev }\OtherTok{\textless{}{-}}\NormalTok{ evidence\_combinations }\SpecialCharTok{|\textgreater{}} \FunctionTok{slice\_sample}\NormalTok{(}\AttributeTok{n =}\NormalTok{ N, }\AttributeTok{replace =} \ConstantTok{TRUE}\NormalTok{)}
      
\NormalTok{      lo\_s }\OtherTok{\textless{}{-}} \FunctionTok{log}\NormalTok{((}\FloatTok{0.5} \SpecialCharTok{+}\NormalTok{ ev}\SpecialCharTok{$}\NormalTok{blue1) }\SpecialCharTok{/}\NormalTok{ (}\FloatTok{0.5} \SpecialCharTok{+}\NormalTok{ ev}\SpecialCharTok{$}\NormalTok{total1 }\SpecialCharTok{{-}}\NormalTok{ ev}\SpecialCharTok{$}\NormalTok{blue1))}
\NormalTok{      lo\_o }\OtherTok{\textless{}{-}} \FunctionTok{log}\NormalTok{((}\FloatTok{0.5} \SpecialCharTok{+}\NormalTok{ ev}\SpecialCharTok{$}\NormalTok{blue2) }\SpecialCharTok{/}\NormalTok{ (}\FloatTok{0.5} \SpecialCharTok{+}\NormalTok{ ev}\SpecialCharTok{$}\NormalTok{total2 }\SpecialCharTok{{-}}\NormalTok{ ev}\SpecialCharTok{$}\NormalTok{blue2))}
      
\NormalTok{      L\_post }\OtherTok{\textless{}{-}} \FunctionTok{jardri\_f}\NormalTok{(lo\_s, ws) }\SpecialCharTok{+} \FunctionTok{jardri\_f}\NormalTok{(lo\_o, wo)}
\NormalTok{      choices }\OtherTok{\textless{}{-}} \FunctionTok{rbinom}\NormalTok{(N, }\DecValTok{1}\NormalTok{, }\FunctionTok{plogis}\NormalTok{(L\_post))}

      \FunctionTok{list}\NormalTok{(}
        \AttributeTok{variables =} \FunctionTok{list}\NormalTok{(}\AttributeTok{w\_self =}\NormalTok{ ws, }\AttributeTok{w\_other =}\NormalTok{ wo),}
        \AttributeTok{generated =} \FunctionTok{list}\NormalTok{(}\AttributeTok{N =}\NormalTok{ N, }\AttributeTok{choice =}\NormalTok{ choices, }\AttributeTok{blue1 =}\NormalTok{ ev}\SpecialCharTok{$}\NormalTok{blue1, }
                         \AttributeTok{total1 =}\NormalTok{ ev}\SpecialCharTok{$}\NormalTok{total1, }\AttributeTok{blue2 =}\NormalTok{ ev}\SpecialCharTok{$}\NormalTok{blue2, }\AttributeTok{total2 =}\NormalTok{ ev}\SpecialCharTok{$}\NormalTok{total2)}
\NormalTok{      )}
\NormalTok{    \}, }\AttributeTok{N =}\NormalTok{ N\_sbc\_trials}
\NormalTok{  )}

  \CommentTok{\# {-}{-}{-} GENERATOR 2: Weights + Loops {-}{-}{-}}
\NormalTok{  sbc\_gen\_loops }\OtherTok{\textless{}{-}} \FunctionTok{SBC\_generator\_function}\NormalTok{(}
    \ControlFlowTok{function}\NormalTok{(N) \{}
\NormalTok{      ws }\OtherTok{\textless{}{-}} \FunctionTok{draw\_w}\NormalTok{()}
\NormalTok{      wo }\OtherTok{\textless{}{-}} \FunctionTok{draw\_w}\NormalTok{()}
\NormalTok{      as\_m1 }\OtherTok{\textless{}{-}} \FunctionTok{abs}\NormalTok{(}\FunctionTok{rnorm}\NormalTok{(}\DecValTok{1}\NormalTok{, }\DecValTok{0}\NormalTok{, }\DecValTok{1}\NormalTok{)); ao\_m1 }\OtherTok{\textless{}{-}} \FunctionTok{abs}\NormalTok{(}\FunctionTok{rnorm}\NormalTok{(}\DecValTok{1}\NormalTok{, }\DecValTok{0}\NormalTok{, }\DecValTok{1}\NormalTok{))}
\NormalTok{      as\_ }\OtherTok{\textless{}{-}} \FloatTok{1.0} \SpecialCharTok{+}\NormalTok{ as\_m1; ao\_ }\OtherTok{\textless{}{-}} \FloatTok{1.0} \SpecialCharTok{+}\NormalTok{ ao\_m1}

\NormalTok{      ev }\OtherTok{\textless{}{-}}\NormalTok{ evidence\_combinations }\SpecialCharTok{|\textgreater{}} \FunctionTok{slice\_sample}\NormalTok{(}\AttributeTok{n =}\NormalTok{ N, }\AttributeTok{replace =} \ConstantTok{TRUE}\NormalTok{)}
\NormalTok{      lo\_s }\OtherTok{\textless{}{-}} \FunctionTok{log}\NormalTok{((}\FloatTok{0.5} \SpecialCharTok{+}\NormalTok{ ev}\SpecialCharTok{$}\NormalTok{blue1) }\SpecialCharTok{/}\NormalTok{ (}\FloatTok{0.5} \SpecialCharTok{+}\NormalTok{ ev}\SpecialCharTok{$}\NormalTok{total1 }\SpecialCharTok{{-}}\NormalTok{ ev}\SpecialCharTok{$}\NormalTok{blue1))}
\NormalTok{      lo\_o }\OtherTok{\textless{}{-}} \FunctionTok{log}\NormalTok{((}\FloatTok{0.5} \SpecialCharTok{+}\NormalTok{ ev}\SpecialCharTok{$}\NormalTok{blue2) }\SpecialCharTok{/}\NormalTok{ (}\FloatTok{0.5} \SpecialCharTok{+}\NormalTok{ ev}\SpecialCharTok{$}\NormalTok{total2 }\SpecialCharTok{{-}}\NormalTok{ ev}\SpecialCharTok{$}\NormalTok{blue2))}
      
\NormalTok{      L\_post }\OtherTok{\textless{}{-}} \FunctionTok{jardri\_f}\NormalTok{(as\_ }\SpecialCharTok{*}\NormalTok{ lo\_s, ws) }\SpecialCharTok{+} \FunctionTok{jardri\_f}\NormalTok{(ao\_ }\SpecialCharTok{*}\NormalTok{ lo\_o, wo)}
\NormalTok{      choices }\OtherTok{\textless{}{-}} \FunctionTok{rbinom}\NormalTok{(N, }\DecValTok{1}\NormalTok{, }\FunctionTok{plogis}\NormalTok{(L\_post))}

      \FunctionTok{list}\NormalTok{(}
        \AttributeTok{variables =} \FunctionTok{list}\NormalTok{(}\AttributeTok{w\_self =}\NormalTok{ ws, }\AttributeTok{w\_other =}\NormalTok{ wo, }\AttributeTok{alpha\_self\_m1 =}\NormalTok{ as\_m1, }\AttributeTok{alpha\_other\_m1 =}\NormalTok{ ao\_m1),}
        \AttributeTok{generated =} \FunctionTok{list}\NormalTok{(}\AttributeTok{N =}\NormalTok{ N, }\AttributeTok{choice =}\NormalTok{ choices, }\AttributeTok{blue1 =}\NormalTok{ ev}\SpecialCharTok{$}\NormalTok{blue1, }
                         \AttributeTok{total1 =}\NormalTok{ ev}\SpecialCharTok{$}\NormalTok{total1, }\AttributeTok{blue2 =}\NormalTok{ ev}\SpecialCharTok{$}\NormalTok{blue2, }\AttributeTok{total2 =}\NormalTok{ ev}\SpecialCharTok{$}\NormalTok{total2)}
\NormalTok{      )}
\NormalTok{    \}, }\AttributeTok{N =}\NormalTok{ N\_sbc\_trials}
\NormalTok{  )}

  \CommentTok{\# {-}{-}{-} GENERATOR 3: Full Interference {-}{-}{-}}
\NormalTok{  sbc\_gen\_interference }\OtherTok{\textless{}{-}} \FunctionTok{SBC\_generator\_function}\NormalTok{(}
    \ControlFlowTok{function}\NormalTok{(N) \{}
\NormalTok{      ws }\OtherTok{\textless{}{-}} \FunctionTok{draw\_w}\NormalTok{()}
\NormalTok{      wo }\OtherTok{\textless{}{-}} \FunctionTok{draw\_w}\NormalTok{()}
\NormalTok{      as\_m1 }\OtherTok{\textless{}{-}} \FunctionTok{abs}\NormalTok{(}\FunctionTok{rnorm}\NormalTok{(}\DecValTok{1}\NormalTok{, }\DecValTok{0}\NormalTok{, }\DecValTok{1}\NormalTok{)); ao\_m1 }\OtherTok{\textless{}{-}} \FunctionTok{abs}\NormalTok{(}\FunctionTok{rnorm}\NormalTok{(}\DecValTok{1}\NormalTok{, }\DecValTok{0}\NormalTok{, }\DecValTok{1}\NormalTok{))}
\NormalTok{      as\_ }\OtherTok{\textless{}{-}} \FloatTok{1.0} \SpecialCharTok{+}\NormalTok{ as\_m1; ao\_ }\OtherTok{\textless{}{-}} \FloatTok{1.0} \SpecialCharTok{+}\NormalTok{ ao\_m1}

\NormalTok{      ev }\OtherTok{\textless{}{-}}\NormalTok{ evidence\_combinations }\SpecialCharTok{|\textgreater{}} \FunctionTok{slice\_sample}\NormalTok{(}\AttributeTok{n =}\NormalTok{ N, }\AttributeTok{replace =} \ConstantTok{TRUE}\NormalTok{)}
\NormalTok{      lo\_s }\OtherTok{\textless{}{-}} \FunctionTok{log}\NormalTok{((}\FloatTok{0.5} \SpecialCharTok{+}\NormalTok{ ev}\SpecialCharTok{$}\NormalTok{blue1) }\SpecialCharTok{/}\NormalTok{ (}\FloatTok{0.5} \SpecialCharTok{+}\NormalTok{ ev}\SpecialCharTok{$}\NormalTok{total1 }\SpecialCharTok{{-}}\NormalTok{ ev}\SpecialCharTok{$}\NormalTok{blue1))}
\NormalTok{      lo\_o }\OtherTok{\textless{}{-}} \FunctionTok{log}\NormalTok{((}\FloatTok{0.5} \SpecialCharTok{+}\NormalTok{ ev}\SpecialCharTok{$}\NormalTok{blue2) }\SpecialCharTok{/}\NormalTok{ (}\FloatTok{0.5} \SpecialCharTok{+}\NormalTok{ ev}\SpecialCharTok{$}\NormalTok{total2 }\SpecialCharTok{{-}}\NormalTok{ ev}\SpecialCharTok{$}\NormalTok{blue2))}
      
\NormalTok{      I }\OtherTok{\textless{}{-}} \FunctionTok{jardri\_f}\NormalTok{(as\_ }\SpecialCharTok{*}\NormalTok{ lo\_s, ws) }\SpecialCharTok{+} \FunctionTok{jardri\_f}\NormalTok{(ao\_ }\SpecialCharTok{*}\NormalTok{ lo\_o, wo)}
\NormalTok{      L\_post }\OtherTok{\textless{}{-}} \FunctionTok{jardri\_f}\NormalTok{(lo\_s }\SpecialCharTok{+}\NormalTok{ I, ws) }\SpecialCharTok{+} \FunctionTok{jardri\_f}\NormalTok{(lo\_o }\SpecialCharTok{+}\NormalTok{ I, wo)}
\NormalTok{      choices }\OtherTok{\textless{}{-}} \FunctionTok{rbinom}\NormalTok{(N, }\DecValTok{1}\NormalTok{, }\FunctionTok{plogis}\NormalTok{(L\_post))}

      \FunctionTok{list}\NormalTok{(}
        \AttributeTok{variables =} \FunctionTok{list}\NormalTok{(}\AttributeTok{w\_self =}\NormalTok{ ws, }\AttributeTok{w\_other =}\NormalTok{ wo, }\AttributeTok{alpha\_self\_m1 =}\NormalTok{ as\_m1, }\AttributeTok{alpha\_other\_m1 =}\NormalTok{ ao\_m1),}
        \AttributeTok{generated =} \FunctionTok{list}\NormalTok{(}\AttributeTok{N =}\NormalTok{ N, }\AttributeTok{choice =}\NormalTok{ choices, }\AttributeTok{blue1 =}\NormalTok{ ev}\SpecialCharTok{$}\NormalTok{blue1, }
                         \AttributeTok{total1 =}\NormalTok{ ev}\SpecialCharTok{$}\NormalTok{total1, }\AttributeTok{blue2 =}\NormalTok{ ev}\SpecialCharTok{$}\NormalTok{blue2, }\AttributeTok{total2 =}\NormalTok{ ev}\SpecialCharTok{$}\NormalTok{total2)}
\NormalTok{      )}
\NormalTok{    \}, }\AttributeTok{N =}\NormalTok{ N\_sbc\_trials}
\NormalTok{  )}

  \CommentTok{\# 2. Generate datasets}
\NormalTok{  n\_sbc }\OtherTok{\textless{}{-}} \DecValTok{400}
\NormalTok{  datasets\_weights }\OtherTok{\textless{}{-}} \FunctionTok{generate\_datasets}\NormalTok{(sbc\_gen\_weights, n\_sbc)}
\NormalTok{  datasets\_loops   }\OtherTok{\textless{}{-}} \FunctionTok{generate\_datasets}\NormalTok{(sbc\_gen\_loops, n\_sbc)}
\NormalTok{  datasets\_interf  }\OtherTok{\textless{}{-}} \FunctionTok{generate\_datasets}\NormalTok{(sbc\_gen\_interference, n\_sbc)}

  \CommentTok{\# 3. Setup Backends}
\NormalTok{  sampler\_args }\OtherTok{\textless{}{-}} \FunctionTok{list}\NormalTok{(}\AttributeTok{chains =} \DecValTok{2}\NormalTok{, }\AttributeTok{iter\_warmup =} \DecValTok{500}\NormalTok{, }\AttributeTok{iter\_sampling =} \DecValTok{500}\NormalTok{, }\AttributeTok{parallel\_chains =} \DecValTok{2}\NormalTok{, }\AttributeTok{refresh =} \DecValTok{0}\NormalTok{)}
  
\NormalTok{  backend\_weights }\OtherTok{\textless{}{-}} \FunctionTok{do.call}\NormalTok{(SBC}\SpecialCharTok{::}\NormalTok{SBC\_backend\_cmdstan\_sample, }\FunctionTok{c}\NormalTok{(}\FunctionTok{list}\NormalTok{(}\AttributeTok{model =}\NormalTok{ mod\_ci\_weights), sampler\_args))}
\NormalTok{  backend\_loops   }\OtherTok{\textless{}{-}} \FunctionTok{do.call}\NormalTok{(SBC}\SpecialCharTok{::}\NormalTok{SBC\_backend\_cmdstan\_sample, }\FunctionTok{c}\NormalTok{(}\FunctionTok{list}\NormalTok{(}\AttributeTok{model =}\NormalTok{ mod\_ci\_loops), sampler\_args))}
\NormalTok{  backend\_interf  }\OtherTok{\textless{}{-}} \FunctionTok{do.call}\NormalTok{(SBC}\SpecialCharTok{::}\NormalTok{SBC\_backend\_cmdstan\_sample, }\FunctionTok{c}\NormalTok{(}\FunctionTok{list}\NormalTok{(}\AttributeTok{model =}\NormalTok{ mod\_ci\_interference), sampler\_args))}

  \CommentTok{\# 4. Compute SBC}
\NormalTok{  sbc\_res\_weights }\OtherTok{\textless{}{-}} \FunctionTok{compute\_SBC}\NormalTok{(datasets\_weights, backend\_weights, }\AttributeTok{keep\_fits =} \ConstantTok{FALSE}\NormalTok{)}
\NormalTok{  sbc\_res\_loops   }\OtherTok{\textless{}{-}} \FunctionTok{compute\_SBC}\NormalTok{(datasets\_loops, backend\_loops, }\AttributeTok{keep\_fits =} \ConstantTok{FALSE}\NormalTok{)}
\NormalTok{  sbc\_res\_interf  }\OtherTok{\textless{}{-}} \FunctionTok{compute\_SBC}\NormalTok{(datasets\_interf, backend\_interf, }\AttributeTok{keep\_fits =} \ConstantTok{FALSE}\NormalTok{)}

  \FunctionTok{saveRDS}\NormalTok{(}\FunctionTok{list}\NormalTok{(}\AttributeTok{weights =}\NormalTok{ sbc\_res\_weights, }\AttributeTok{loops =}\NormalTok{ sbc\_res\_loops, }\AttributeTok{interf =}\NormalTok{ sbc\_res\_interf), }
          \StringTok{"simmodels/ch10\_sbc\_ci\_all\_results.rds"}\NormalTok{)}

\NormalTok{\} }\ControlFlowTok{else}\NormalTok{ \{}
\NormalTok{  sbc\_all }\OtherTok{\textless{}{-}} \FunctionTok{readRDS}\NormalTok{(}\StringTok{"simmodels/ch10\_sbc\_ci\_all\_results.rds"}\NormalTok{)}
\NormalTok{  sbc\_res\_weights }\OtherTok{\textless{}{-}}\NormalTok{ sbc\_all}\SpecialCharTok{$}\NormalTok{weights}
\NormalTok{  sbc\_res\_loops   }\OtherTok{\textless{}{-}}\NormalTok{ sbc\_all}\SpecialCharTok{$}\NormalTok{loops}
\NormalTok{  sbc\_res\_interf  }\OtherTok{\textless{}{-}}\NormalTok{ sbc\_all}\SpecialCharTok{$}\NormalTok{interf}
\NormalTok{\}}
\end{Highlighting}
\end{Shaded}

\subsubsection{Checking Posterior Calibration
(ECDF)}\label{checking-posterior-calibration-ecdf}

First, we check the Empirical Cumulative Distribution Function (ECDF)
plots to see if the sampler explores the parameter space without bias
across all three models.

\begin{Shaded}
\begin{Highlighting}[]
\NormalTok{p\_ecdf\_weights }\OtherTok{\textless{}{-}}\NormalTok{ SBC}\SpecialCharTok{::}\FunctionTok{plot\_ecdf\_diff}\NormalTok{(sbc\_res\_weights, }\AttributeTok{variables =} \FunctionTok{c}\NormalTok{(}\StringTok{"w\_self"}\NormalTok{, }\StringTok{"w\_other"}\NormalTok{)) }\SpecialCharTok{+} 
  \FunctionTok{labs}\NormalTok{(}\AttributeTok{title =} \StringTok{"1. Weights Only Calibration"}\NormalTok{)}

\NormalTok{p\_ecdf\_loops }\OtherTok{\textless{}{-}}\NormalTok{ SBC}\SpecialCharTok{::}\FunctionTok{plot\_ecdf\_diff}\NormalTok{(sbc\_res\_loops, }\AttributeTok{variables =} \FunctionTok{c}\NormalTok{(}\StringTok{"w\_self"}\NormalTok{, }\StringTok{"w\_other"}\NormalTok{, }\StringTok{"alpha\_self\_m1"}\NormalTok{, }\StringTok{"alpha\_other\_m1"}\NormalTok{)) }\SpecialCharTok{+} 
  \FunctionTok{labs}\NormalTok{(}\AttributeTok{title =} \StringTok{"2. Weights + Loops Calibration"}\NormalTok{)}

\NormalTok{p\_ecdf\_interf }\OtherTok{\textless{}{-}}\NormalTok{ SBC}\SpecialCharTok{::}\FunctionTok{plot\_ecdf\_diff}\NormalTok{(sbc\_res\_interf, }\AttributeTok{variables =} \FunctionTok{c}\NormalTok{(}\StringTok{"w\_self"}\NormalTok{, }\StringTok{"w\_other"}\NormalTok{, }\StringTok{"alpha\_self\_m1"}\NormalTok{, }\StringTok{"alpha\_other\_m1"}\NormalTok{)) }\SpecialCharTok{+} 
  \FunctionTok{labs}\NormalTok{(}\AttributeTok{title =} \StringTok{"3. Full Interference Calibration"}\NormalTok{)}

\NormalTok{(p\_ecdf\_weights }\SpecialCharTok{/}\NormalTok{ p\_ecdf\_loops }\SpecialCharTok{/}\NormalTok{ p\_ecdf\_interf) }\SpecialCharTok{+}
  \FunctionTok{plot\_annotation}\NormalTok{(}
    \AttributeTok{title =} \StringTok{"SBC Calibration Across the Circular Inference Family"}\NormalTok{,}
    \AttributeTok{subtitle =} \StringTok{"If lines stray outside the blue bands, the posterior is systematically biased."}
\NormalTok{  )}
\end{Highlighting}
\end{Shaded}

\begin{center}
\includegraphics[width=0.8\linewidth,height=\textheight,keepaspectratio]{10-BayesianModels_files/figure-pdf/ch10_sbc_ci_ecdf_viz-1.pdf}
\end{center}

The ECDF plots generally stay within the blue envelopes, though you will
likely see the loop parameters (alpha\_*\_m1) gently push against the
edges in the Interference model. This indicates that Stan's NUTS sampler
is working hard but succeeding at unbiased exploration.

\subsubsection{Parameter Recovery: The Asymptote Problem
Emerges}\label{parameter-recovery-the-asymptote-problem-emerges}

Calibration tells us if the posterior captures the true value on
average, but Parameter Recovery tells us how precise those estimates
are. By plotting all three models together, we can see exactly what
structural feature destroys our precision.

\begin{Shaded}
\begin{Highlighting}[]
\CommentTok{\# Extract and label recovery stats}
\NormalTok{prep\_recovery }\OtherTok{\textless{}{-}} \ControlFlowTok{function}\NormalTok{(sbc\_obj, model\_name) \{}
\NormalTok{  sbc\_obj}\SpecialCharTok{$}\NormalTok{stats }\SpecialCharTok{|\textgreater{}}
    \FunctionTok{filter}\NormalTok{(variable }\SpecialCharTok{\%in\%} \FunctionTok{c}\NormalTok{(}\StringTok{"w\_self"}\NormalTok{, }\StringTok{"w\_other"}\NormalTok{, }\StringTok{"alpha\_self\_m1"}\NormalTok{, }\StringTok{"alpha\_other\_m1"}\NormalTok{)) }\SpecialCharTok{|\textgreater{}}
    \FunctionTok{mutate}\NormalTok{(}
      \AttributeTok{Model =}\NormalTok{ model\_name,}
      \AttributeTok{param\_label =} \FunctionTok{case\_when}\NormalTok{(}
\NormalTok{        variable }\SpecialCharTok{==} \StringTok{"w\_self"} \SpecialCharTok{\textasciitilde{}} \StringTok{"Weight Self (w\_s)"}\NormalTok{,}
\NormalTok{        variable }\SpecialCharTok{==} \StringTok{"w\_other"} \SpecialCharTok{\textasciitilde{}} \StringTok{"Weight Other (w\_o)"}\NormalTok{,}
\NormalTok{        variable }\SpecialCharTok{==} \StringTok{"alpha\_self\_m1"} \SpecialCharTok{\textasciitilde{}} \StringTok{"Excess Loops Self (\textbackslash{}u03b1\_s {-} 1)"}\NormalTok{,}
\NormalTok{        variable }\SpecialCharTok{==} \StringTok{"alpha\_other\_m1"} \SpecialCharTok{\textasciitilde{}} \StringTok{"Excess Loops Other (\textbackslash{}u03b1\_o {-} 1)"}
\NormalTok{      )}
\NormalTok{    )}
\NormalTok{\}}

\NormalTok{rec\_df }\OtherTok{\textless{}{-}} \FunctionTok{bind\_rows}\NormalTok{(}
  \FunctionTok{prep\_recovery}\NormalTok{(sbc\_res\_weights, }\StringTok{"1. Weights Only"}\NormalTok{),}
  \FunctionTok{prep\_recovery}\NormalTok{(sbc\_res\_loops, }\StringTok{"2. Weights + Loops"}\NormalTok{),}
  \FunctionTok{prep\_recovery}\NormalTok{(sbc\_res\_interf, }\StringTok{"3. Full Interference"}\NormalTok{)}
\NormalTok{) }\SpecialCharTok{|\textgreater{}}
  \FunctionTok{mutate}\NormalTok{(}\AttributeTok{Model =} \FunctionTok{factor}\NormalTok{(Model, }\AttributeTok{levels =} \FunctionTok{c}\NormalTok{(}\StringTok{"1. Weights Only"}\NormalTok{, }\StringTok{"2. Weights + Loops"}\NormalTok{, }\StringTok{"3. Full Interference"}\NormalTok{)))}

\CommentTok{\# {-}{-}{-} Plot 1: Weights Only (Strictly bounded 0.5 to 1.0) {-}{-}{-}}
\NormalTok{p\_rec\_weights }\OtherTok{\textless{}{-}}\NormalTok{ rec\_df }\SpecialCharTok{|\textgreater{}} 
  \FunctionTok{filter}\NormalTok{(}\FunctionTok{str\_detect}\NormalTok{(variable, }\StringTok{"w\_"}\NormalTok{)) }\SpecialCharTok{|\textgreater{}}
  \FunctionTok{ggplot}\NormalTok{(}\FunctionTok{aes}\NormalTok{(}\AttributeTok{x =}\NormalTok{ simulated\_value, }\AttributeTok{y =}\NormalTok{ median)) }\SpecialCharTok{+}
  \FunctionTok{geom\_abline}\NormalTok{(}\AttributeTok{slope =} \DecValTok{1}\NormalTok{, }\AttributeTok{intercept =} \DecValTok{0}\NormalTok{, }\AttributeTok{linetype =} \StringTok{"dashed"}\NormalTok{, }\AttributeTok{color =} \StringTok{"grey50"}\NormalTok{) }\SpecialCharTok{+}
  \FunctionTok{geom\_errorbar}\NormalTok{(}\FunctionTok{aes}\NormalTok{(}\AttributeTok{ymin =}\NormalTok{ q5, }\AttributeTok{ymax =}\NormalTok{ q95), }\AttributeTok{alpha =} \FloatTok{0.15}\NormalTok{, }\AttributeTok{width =} \DecValTok{0}\NormalTok{, }\AttributeTok{color =} \StringTok{"\#4575b4"}\NormalTok{) }\SpecialCharTok{+}
  \FunctionTok{geom\_point}\NormalTok{(}\AttributeTok{alpha =} \FloatTok{0.6}\NormalTok{, }\AttributeTok{size =} \FloatTok{1.5}\NormalTok{, }\AttributeTok{color =} \StringTok{"\#d73027"}\NormalTok{) }\SpecialCharTok{+}
  \FunctionTok{facet\_grid}\NormalTok{(Model }\SpecialCharTok{\textasciitilde{}}\NormalTok{ param\_label) }\SpecialCharTok{+}
  \CommentTok{\# Force the axes to the theoretical boundaries}
  \FunctionTok{coord\_cartesian}\NormalTok{(}\AttributeTok{xlim =} \FunctionTok{c}\NormalTok{(}\FloatTok{0.5}\NormalTok{, }\FloatTok{1.0}\NormalTok{), }\AttributeTok{ylim =} \FunctionTok{c}\NormalTok{(}\FloatTok{0.5}\NormalTok{, }\FloatTok{1.0}\NormalTok{)) }\SpecialCharTok{+} 
  \FunctionTok{labs}\NormalTok{(}
    \AttributeTok{title =} \StringTok{"Parameter Recovery: Trust Weights (w)"}\NormalTok{,}
    \AttributeTok{subtitle =} \StringTok{"Weights recover well, though interference adds visible noise to the estimates."}\NormalTok{,}
    \AttributeTok{x =} \StringTok{"True Simulated Value"}\NormalTok{,}
    \AttributeTok{y =} \StringTok{"Posterior Median"}
\NormalTok{  ) }\SpecialCharTok{+}
  \FunctionTok{theme}\NormalTok{(}\AttributeTok{strip.background =} \FunctionTok{element\_rect}\NormalTok{(}\AttributeTok{fill =} \StringTok{"gray95"}\NormalTok{, }\AttributeTok{color =} \ConstantTok{NA}\NormalTok{),}
        \AttributeTok{strip.text =} \FunctionTok{element\_text}\NormalTok{(}\AttributeTok{face =} \StringTok{"bold"}\NormalTok{))}

\CommentTok{\# {-}{-}{-} Plot 2: Loops Only (Free scales to show the asymptote) {-}{-}{-}}
\NormalTok{p\_rec\_loops }\OtherTok{\textless{}{-}}\NormalTok{ rec\_df }\SpecialCharTok{|\textgreater{}} 
  \FunctionTok{filter}\NormalTok{(}\FunctionTok{str\_detect}\NormalTok{(variable, }\StringTok{"alpha\_"}\NormalTok{)) }\SpecialCharTok{|\textgreater{}}
  \CommentTok{\# Drop the "Weights Only" model since it doesn\textquotesingle{}t have loop parameters}
  \FunctionTok{filter}\NormalTok{(Model }\SpecialCharTok{!=} \StringTok{"1. Weights Only"}\NormalTok{) }\SpecialCharTok{|\textgreater{}} 
  \FunctionTok{ggplot}\NormalTok{(}\FunctionTok{aes}\NormalTok{(}\AttributeTok{x =}\NormalTok{ simulated\_value, }\AttributeTok{y =}\NormalTok{ median)) }\SpecialCharTok{+}
  \FunctionTok{geom\_abline}\NormalTok{(}\AttributeTok{slope =} \DecValTok{1}\NormalTok{, }\AttributeTok{intercept =} \DecValTok{0}\NormalTok{, }\AttributeTok{linetype =} \StringTok{"dashed"}\NormalTok{, }\AttributeTok{color =} \StringTok{"grey50"}\NormalTok{) }\SpecialCharTok{+}
  \FunctionTok{geom\_errorbar}\NormalTok{(}\FunctionTok{aes}\NormalTok{(}\AttributeTok{ymin =}\NormalTok{ q5, }\AttributeTok{ymax =}\NormalTok{ q95), }\AttributeTok{alpha =} \FloatTok{0.15}\NormalTok{, }\AttributeTok{width =} \DecValTok{0}\NormalTok{, }\AttributeTok{color =} \StringTok{"\#4575b4"}\NormalTok{) }\SpecialCharTok{+}
  \FunctionTok{geom\_point}\NormalTok{(}\AttributeTok{alpha =} \FloatTok{0.6}\NormalTok{, }\AttributeTok{size =} \FloatTok{1.5}\NormalTok{, }\AttributeTok{color =} \StringTok{"\#d73027"}\NormalTok{) }\SpecialCharTok{+}
  \FunctionTok{facet\_grid}\NormalTok{(Model }\SpecialCharTok{\textasciitilde{}}\NormalTok{ param\_label, }\AttributeTok{scales =} \StringTok{"free"}\NormalTok{) }\SpecialCharTok{+}
  \FunctionTok{labs}\NormalTok{(}
    \AttributeTok{title =} \StringTok{"Parameter Recovery: Loop Parameters (\textbackslash{}u03b1 {-} 1)"}\NormalTok{,}
    \AttributeTok{subtitle =} \StringTok{"Notice the severe \textquotesingle{}asymptotic flattening\textquotesingle{} and massive uncertainty at higher values."}\NormalTok{,}
    \AttributeTok{x =} \StringTok{"True Simulated Value"}\NormalTok{,}
    \AttributeTok{y =} \StringTok{"Posterior Median"}
\NormalTok{  ) }\SpecialCharTok{+}
  \FunctionTok{theme}\NormalTok{(}\AttributeTok{strip.background =} \FunctionTok{element\_rect}\NormalTok{(}\AttributeTok{fill =} \StringTok{"gray95"}\NormalTok{, }\AttributeTok{color =} \ConstantTok{NA}\NormalTok{),}
        \AttributeTok{strip.text =} \FunctionTok{element\_text}\NormalTok{(}\AttributeTok{face =} \StringTok{"bold"}\NormalTok{))}

\CommentTok{\# Combine with patchwork}
\NormalTok{p\_rec\_weights }\SpecialCharTok{/}\NormalTok{ p\_rec\_loops}
\end{Highlighting}
\end{Shaded}

\begin{center}
\includegraphics[width=0.8\linewidth,height=\textheight,keepaspectratio]{10-BayesianModels_files/figure-pdf/ch10_sbc_ci_recovery-1.pdf}
\end{center}

Weights are easy to recover, but loops not. What can we do? My hunch
again is that we should collect data on participants' confidence, as
that should give us better resolution.

\subsection{Discussion: When to Use Each
Parameterisation}\label{discussion-when-to-use-each-parameterisation}

The parameterisations explored in this chapter make different claims
about the data-generating process:

\textbf{Pseudocount (SBA/PBA/WBA):} The mind treats each source as an
independent binomial sample. Evidence is counted once, with optional
rescaling. The BetaBinomial observation model claims probability
matching --- actions sample from the full belief distribution.

\textbf{Circular inference (Jardri \& Denève 2013):} Evidence is not
merely weighted but potentially amplified beyond its face value. The
Bernoulli observation model collapses the belief to a single log-odds
value; distributional uncertainty is captured by the posterior over
parameters. The key predictive signature --- steeper response curves for
the amplified source --- cannot be reproduced by pseudocount weighting
alone.

The choice between models should be driven by the scientific question
and validated by LOO comparison.

\section{What the Bayesian Workflow
Adds}\label{what-the-bayesian-workflow-adds}

This chapter has done more than fit cognitive models --- it has
\emph{validated} them through a disciplined pipeline. Prior predictive
checks established that the priors of all three models are cognitively
plausible. SBC certified that the Stan samplers for the WBA and PBA are
unbiased. Single-fit recovery confirmed identifiability. LOO-PIT and
sensitivity analysis characterised the limits of the fitted posteriors.
The three-way PSIS-LOO comparison arbitrated between competing cognitive
hypotheses in a way that penalises complexity and estimates genuine
predictive performance.

This pipeline is not bureaucracy. Without SBC, a subtle Stan
implementation bug could produce attractive posteriors that are
systematically biased. Without LOO-PIT, we cannot distinguish a model
that fits from a model that truly predicts. Without the three-way
comparison, we cannot know whether the WBA's additional parameter
reflects a real cognitive structure or mere overfitting. And without
parameter recovery, we would never have discovered that the raw
\((w_d, w_s)\) parameterisation was structurally unidentifiable --- a
problem invisible to standard MCMC diagnostics.

\subsection{Limitations of the Current
Models}\label{limitations-of-the-current-models}

\begin{enumerate}
\def\labelenumi{\arabic{enumi}.}
\item
  \textbf{The observation model is itself a hypothesis.} The
  BetaBinomial claims participants probability-match --- sampling their
  action from the full posterior over \(\theta\). If participants are
  more deterministic (choosing blue whenever
  \(\text{E}[\theta] > 0.5\)), the BetaBinomial will overestimate
  response variability. A softmax observation model with an
  inverse-temperature parameter would test this alternative; the LOO-PIT
  hump pattern is the diagnostic signature to watch for.
\item
  \textbf{Trial independence in non-sequential models.} The SBA, PBA,
  and WBA assume each trial involves a fresh, independent jar. Applying
  these models to paradigms with across-trial structure requires the
  sequential extension --- with LFO-CV replacing standard LOO.
\item
  \textbf{Fixed prior pseudo-count.} We fix \(\alpha_0 = \beta_0 = 0.5\)
  (Jeffreys prior), assuming no prior knowledge about the jar. A
  hierarchical extension would estimate a population distribution over
  \((\alpha_0, \beta_0)\), allowing for individual differences in
  baseline expectations.
\item
  \textbf{The PBA's identifiability boundary.} When \(p \to 0\) or
  \(p \to 1\), one evidence source is entirely ignored. The Beta(2,2)
  prior keeps the sampler away from these boundaries, but check tail ESS
  whenever the PBA is fitted to strongly asymmetric agents.
\item
  \textbf{Social evidence encoding and pool-size asymmetry.} Social
  evidence is encoded as \texttt{total2\ =\ 3} pseudo-marbles, making
  direct and social evidence structurally asymmetric. Recovered weights
  should be interpreted relative to the encoding, not as absolute
  measures of social reliance.
\end{enumerate}

\subsection{Connections and
Extensions}\label{connections-and-extensions}

\begin{itemize}
\tightlist
\item
  \textbf{Rational inattention models} relax the assumption that all
  evidence is perfectly observed --- the agent may allocate limited
  attentional resources to each source.
\item
  \textbf{Resource-rational models} (Griffiths et al.~2015) ask: what
  integration strategy is optimal given limited computational resources?
  The WBA can be interpreted as an approximation to full Bayesian
  inference under resource constraints.
\item
  \textbf{Active inference} (Friston 2010) frames perception and action
  as a unified process of minimising variational free energy, offering a
  different theoretical foundation for the same empirical phenomena.
\end{itemize}

\begin{tcolorbox}[enhanced jigsaw, toprule=.15mm, toptitle=1mm, bottomtitle=1mm, leftrule=.75mm, rightrule=.15mm, colframe=quarto-callout-note-color-frame, left=2mm, colback=white, titlerule=0mm, coltitle=black, colbacktitle=quarto-callout-note-color!10!white, opacityback=0, title=\textcolor{quarto-callout-note-color}{\faInfo}\hspace{0.5em}{Stan Engineering: What We Have Built Across Chapters 4--10}, arc=.35mm, opacitybacktitle=0.6, breakable, bottomrule=.15mm]

Across this book's first ten chapters we have introduced a layered set
of Stan programming techniques. Before moving to the categorization
chapters --- where all of them combine --- it is worth having the full
picture in one place.

\textbf{Stage I (Ch. 4--5): Foundations}

\begin{longtable}[]{@{}
  >{\raggedright\arraybackslash}p{(\linewidth - 2\tabcolsep) * \real{0.5000}}
  >{\raggedright\arraybackslash}p{(\linewidth - 2\tabcolsep) * \real{0.5000}}@{}}
\toprule\noalign{}
\begin{minipage}[b]{\linewidth}\raggedright
Technique
\end{minipage} & \begin{minipage}[b]{\linewidth}\raggedright
Rule
\end{minipage} \\
\midrule\noalign{}
\endhead
\bottomrule\noalign{}
\endlastfoot
Block placement & \texttt{transformed\ data} for data-only quantities;
\texttt{transformed\ parameters} for parameter-dependent state;
\texttt{model} for local accumulators; \texttt{generated\ quantities}
for post-sampling diagnostics \\
\texttt{run\_diagnostics} flag & Gate expensive
\texttt{generated\ quantities} (PPC draws, \texttt{log\_lik}) behind an
\texttt{int\textless{}lower=0,upper=1\textgreater{}} data flag; set to 0
in SBC backends \\
Priors as data & Pass prior hyperparameters through the \texttt{data}
block --- no recompilation when adjusting priors; SBC stays in sync
automatically \\
Pathfinder warm-start & Run \texttt{mod\$pathfinder(...)} first; pass
the result as \texttt{init} to \texttt{mod\$sample(...)} to cut warm-up
by 30--60\% on complex models \\
\end{longtable}

\textbf{Stage II (Ch. 6): Population Structure}

\begin{longtable}[]{@{}
  >{\raggedright\arraybackslash}p{(\linewidth - 2\tabcolsep) * \real{0.5000}}
  >{\raggedright\arraybackslash}p{(\linewidth - 2\tabcolsep) * \real{0.5000}}@{}}
\toprule\noalign{}
\begin{minipage}[b]{\linewidth}\raggedright
Technique
\end{minipage} & \begin{minipage}[b]{\linewidth}\raggedright
Rule
\end{minipage} \\
\midrule\noalign{}
\endhead
\bottomrule\noalign{}
\endlastfoot
Non-centered parameterization (NCP) & Sample standardized offsets
\(z_j \sim N(0,1)\); reconstruct cognitive parameters deterministically
in \texttt{transformed\ parameters}. Eliminates Neal's Funnel in
sparse-data hierarchical models \\
Subject indexing & \texttt{subj\_id} array (scatter-index) for
exchangeable models; \texttt{subj\_start/subj\_end} contiguous-block for
sequential models (Ch. 9+) \\
\end{longtable}

\textbf{Stage III (Ch. 8--10): Process Mechanisms}

\begin{longtable}[]{@{}
  >{\raggedright\arraybackslash}p{(\linewidth - 2\tabcolsep) * \real{0.5000}}
  >{\raggedright\arraybackslash}p{(\linewidth - 2\tabcolsep) * \real{0.5000}}@{}}
\toprule\noalign{}
\begin{minipage}[b]{\linewidth}\raggedright
Technique
\end{minipage} & \begin{minipage}[b]{\linewidth}\raggedright
Rule
\end{minipage} \\
\midrule\noalign{}
\endhead
\bottomrule\noalign{}
\endlastfoot
Discrete marginalization & Never sample discrete parameters with HMC;
use \texttt{log\_sum\_exp} or \texttt{log\_mix} to analytically
integrate them out in the \texttt{model} block \\
Sequential state unrolling & Recursive belief/memory state goes in
\texttt{transformed\ parameters} with volatile state kept local in
\texttt{\{\}}; only the quantities needed downstream are lifted to the
saved-draw tape \\
\end{longtable}

\textbf{Stage IV (Ch. 11): High-Performance}

The categorization chapters introduce the final layer: precomputing
pairwise stimulus distances in \texttt{transformed\ data} (computed once
rather than thousands of times), NCP for all individual parameters
including attention weights (eliminating Dirichlet funnels), and
\texttt{reduce\_sum} for multi-threaded parallelism across subjects
within a chain. The next chapter puts all of these together.

\end{tcolorbox}

\textbf{This is the end of the core modelling sequence.} Chapters 1--10
have traced the complete Bayesian cognitive modelling workflow: from
pizza-stone physics (Ch. 1), through verbal and formal models of
strategic behaviour (Chs. 2--3), parameter estimation (Ch. 4), model
validation (Ch. 5), hierarchical extensions (Ch. 6), model comparison
(Ch. 7), mixture structures (Ch. 8), sequential cognition (Ch. 9), and
Bayesian cognitive theory (Ch. 10). The tools and habits of mind
established here --- generative modelling, principled validation, honest
uncertainty --- apply directly to any new domain or model you encounter.

\section{Further Reading}\label{further-reading}

\begin{itemize}
\tightlist
\item
  Ma, W. J., Kording, K. P., \& Goldreich, D. (2023). \emph{Bayesian
  models of perception and action: An introduction.} MIT Press.
\item
  Griffiths, T. L., Chater, N., \& Tenenbaum, J. B. (Eds.). (2024).
  \emph{Bayesian models of cognition: Reverse engineering the mind.} MIT
  Press.
\item
  Goodman, N. D., Tenenbaum, J. B., \& The ProbMods Contributors (2016).
  \emph{Probabilistic Models of Cognition} (2nd ed.).
  https://probmods.org/
\item
  Bürkner, P.-C., Gabry, J., \& Vehtari, A. (2020). Approximate
  leave-future-out cross-validation for Bayesian time series models.
  \emph{Journal of Statistical Computation and Simulation}, 90(14),
  2499--2523.
\item
  Säilynoja, T., Bürkner, P.-C., \& Vehtari, A. (2022). Graphical test
  for discrete uniformity and its applications in goodness-of-fit
  evaluation and multiple sample comparison. \emph{Statistics and
  Computing}, 32, 32.
\item
  Gabry, J., Simpson, D., Vehtari, A., Betancourt, M., \& Gelman, A.
  (2019). Visualization in Bayesian workflow. \emph{Journal of the Royal
  Statistical Society: Series A}, 182(2), 389--402.
\item
  Simonsen, A., et al.~(2021). Computationally defined markers of
  uncertainty in the theory of mind predict the severity of clinical
  symptoms in schizophrenia. \emph{Schizophrenia Bulletin}.
\item
  Jardri, R., \& Denève, S. (2013). Circular inferences in
  schizophrenia. \emph{Brain}, 136(11), 3227--3241.
\end{itemize}

\bookmarksetup{startatroot}

\chapter{The Generalized Context
Model}\label{sec-categorization-exemplars}

\begin{quote}
\textbf{📍 Where we are in the Bayesian modeling workflow:} Chapters
1--10 built the complete inference toolkit: parameter estimation (Ch.
4), the six-phase validation pipeline (Ch. 5), hierarchical models (Ch.
6), model comparison via LOO (Ch. 7), mixture models with discrete
marginalization (Ch. 8), and hierarchical Bayesian cognition (Ch.
9--10). This and the following 2 chapters apply the full toolkit to
\emph{categorization}, asking not only what parameters describe a
participant's behavior, but which \emph{cognitive mechanism} best
accounts for how people learn to assign objects to kinds. We begin with
the exemplar approach: the Generalized Context Model (GCM). The
following chapters will introduce prototype and rule-based models, and
then compare all three approaches head-to-head using the full validation
battery.

Two important caveats. First, the categorization models learn, that is,
are \textbf{path-dependent} --- the likelihood at trial \(t\) depends on
the full sequence of stimuli on trials \(1, \ldots, t-1\) through the
exemplar memory store. Standard LOO-CV is therefore technically invalid
here. We met this issue when discussing sequential Bayes. Here we
introduce leave-future-out cross-validation (LFO-CV; Bürkner, Gabry, \&
Vehtari, 2020). Second, the rule-based models we will discuss cannot
really be fit using MCMC (we'll see why), so we need to introduce a new
slower method: Approximate Bayesian Computation (ABC).

\textbf{Trilogy structure.} Chapters 11, 12, and 13 form a parallel
trilogy. Each chapter builds a base cognitive model, extends it with a
single non-stationarity parameter (\(\lambda\) here, \(q\) in Ch. 12,
\(\mu\) in Ch. 13), and then validates across the same three
environmental scenarios in the same order: static Kruschke (baseline),
performance-contingent abrupt shifts (misspecified non-stationarity),
and smooth drift (assumption-matched non-stationarity). The ordering
reflects increasing assumption match between the scenario's
data-generating process and each model's structural commitments. Chapter
13 closes with a synthesis of what the three analyses jointly imply
about identifiability and model-task alignment.
\end{quote}

\section{Introduction to Categorization
Models}\label{introduction-to-categorization-models}

\subsection{The Fundamental Problem of
Categorization}\label{the-fundamental-problem-of-categorization}

Categorization is one of the most fundamental cognitive abilities that
humans possess. From early childhood, we learn to organize the world
around us into meaningful categories: distinguishing food from non-food,
safe from dangerous, or one letter from another. This process of
assigning objects to categories is complex, involving perception,
memory, attention, and decision-making.

How do humans learn categories and make categorization decisions? This
question has generated decades of research and theoretical debate in
cognitive science.

\subsection{Why Model Categorization?}\label{why-model-categorization}

Computational models let us formalize competing theories, generate
precise predictions, and test them against human behavior. This chapter
focuses on the \textbf{exemplar approach} --- the Generalized Context
Model (GCM). Chapters 12 and 13 introduce prototype and rule-based
models respectively, with a full cross-model comparison in Chapter 13.

\subsection{Historical Development of Categorization
Models}\label{historical-development-of-categorization-models}

\subsubsection{Early Views: The Classical
Approach}\label{early-views-the-classical-approach}

Early theories of categorization followed what is now called the
classical view, where categories were defined by necessary and
sufficient features. In this view, category membership is an
all-or-nothing affair: an object either satisfies the logical criteria
or it doesn't. Some examples from Aristotle:

\begin{itemize}
\item
  a human is a rational animal.
\item
  a bird has feathers, wings, and a beak; and lays eggs.
\item
  a fish lives in water; has scales; and breathes through gills.
\end{itemize}

You need all of these features to be part of one category, and having
them makes you part of that category.

While intuitive, this approach struggled to explain many aspects of
human categorization, such as:

\begin{itemize}
\item
  Graded category membership (some members seem ``more typical'' than
  others)
\item
  Unclear boundaries between categories.
\item
  Context-dependent categorization
\item
  Family resemblance structures (where no single feature is necessary)
\end{itemize}

Think of the traditional definition of a species: a set of organisms
that can interbreed and produce fertile offspring. However, there are
many cases of organism A being able to interbreed with B, and B with C,
but A and C cannot interbreed (for instance Californian Ensatina
Salamanders have such phenomena). This creates a fuzzy boundary that the
classical view cannot easily accommodate.

\subsubsection{The Prototype Revolution}\label{the-prototype-revolution}

In the 1970s, Eleanor Rosch's pioneering work on prototypes challenged
the classical view. She demonstrated that categories appear to be
organized around central tendencies or ``prototypes,'' with membership
determined by similarity to these prototypes. Objects closer to the
prototype are categorized more quickly and consistently.

\subsubsection{The Exemplar Alternative}\label{the-exemplar-alternative}

In the late 1970s and 1980s, researchers like Douglas Medin and Robert
Nosofsky proposed that rather than abstracting prototypes, people might
store individual exemplars of categories and make judgments based on
similarity to these stored examples.

The Generalized Context Model (GCM), developed by Nosofsky, became the
standard exemplar model. It proposes that categorization decisions are
based on the summed similarity of a new stimulus to all stored exemplars
of each category, weighted by attention to different stimulus
dimensions.

\subsubsection{Rule-Based Models}\label{rule-based-models}

While similarity-based approaches gained prominence, other researchers
argued that people sometimes use explicit rules for categorization
(e.g.~in scientific reasoning, but not only). Models like Bayesian
particle filters, multinomial trees, COVIS and RULEX formalize these
ideas, suggesting that rule learning and application form a core part of
human categorization, particularly for well-defined categories.

\subsubsection{Hybrid Approaches}\label{hybrid-approaches}

More recent work has focused on hybrid models that incorporate elements
of multiple approaches. Models like SUSTAIN (Love, Medin, and Gureckis
2004) and ATRIUM (Erickson and Kruschke 1998) propose that humans can
flexibly switch between strategies or that different systems operate in
parallel.

\subsection{A Systematization Framework for Categorization
Models}\label{a-systematization-framework-for-categorization-models}

To pedagogically navigate and compare the different approaches, we can
deconstruct the cognitive process of categorization into a sequential
six-step pipeline:

\begin{enumerate}
\def\labelenumi{\arabic{enumi}.}
\item
  \textbf{Input Representation:} How are stimuli represented? (e.g.,
  continuous geometry vs.~discrete binary features).
\item
  \textbf{Attentional Processes:} How are features prioritized? Not all
  features are equally relevant given a task, and attention is often
  modeled as weights summing to 1.
\item
  \textbf{Intermediate Representations:} How is the category internally
  stored? (e.g., as raw exemplars, abstract prototypes, or explicit
  rules).
\item
  \textbf{Evidential Mechanisms:} How is the incoming stimulus compared
  to the representation to gather evidence? (e.g., calculating average
  similarity with an exponential decay over distance, or evidence
  accumulation).
\item
  \textbf{Decision Rules:} How is evidence translated into a choice?
  (e.g., maximum similarity, or a probabilistic Luce choice axiom).
\item
  \textbf{Learning:} How do representations, attentional weights, or
  rule structures update over time through exposure, memory
  accumulation, and feedback?
\end{enumerate}

\subsection{Our Implementation
Approach}\label{our-implementation-approach}

Each chapter in the trilogy follows the same validation structure:
forward simulator → Stan model → diagnostic battery → prior/posterior
predictive checks → parameter recovery → prior sensitivity → SBC. This
parallel structure makes it possible to isolate what differs between the
cognitive mechanisms rather than conflating model differences with
implementation choices.

\subsection{Categorization (Cognitive Science) vs.~Classification
(Machine
Learning)}\label{categorization-cognitive-science-vs.-classification-machine-learning}

Because the mathematical tools are often identical, it is tempting to
view cognitive models of categorization simply as early, inefficient
machine learning classifiers. Both fields rely on the exact same
underlying mechanics to map inputs to outputs: calculating distances in
multi-dimensional feature space, applying functions (like Luce or
softmax) to translate evidence into probabilities, and using
gradient-based weight updates or Bayesian filters to learn from
feedback.

However, despite this shared mathematical framework, they serve
fundamentally different goals:

\begin{itemize}
\item
  \textbf{Classification (Machine Learning):} The primary goal is
  \emph{engineering}. It focuses on drawing the optimal decision
  boundary to maximize predictive accuracy and minimize loss on a given
  dataset. The internal mechanism (e.g., a deep neural network's hidden
  layers) doesn't need to resemble human thought---it just needs to work
  efficiently and accurately.
\item
  \textbf{Categorization (Cognitive Modeling):} The primary goal is
  \emph{reverse-engineering}. Categorization is about understanding the
  knowledge structures that enable a biological agent to make predictive
  inferences. We ask not only what parameters describe behavior, but
  which \emph{cognitive mechanism} best accounts for how people actually
  learn, including the specific mistakes they make.
\end{itemize}

We can summarize this distinction along three key dimensions:

\begin{longtable}[]{@{}
  >{\raggedright\arraybackslash}p{(\linewidth - 4\tabcolsep) * \real{0.3333}}
  >{\raggedright\arraybackslash}p{(\linewidth - 4\tabcolsep) * \real{0.3333}}
  >{\raggedright\arraybackslash}p{(\linewidth - 4\tabcolsep) * \real{0.3333}}@{}}
\toprule\noalign{}
\begin{minipage}[b]{\linewidth}\raggedright
Dimension
\end{minipage} & \begin{minipage}[b]{\linewidth}\raggedright
Classification (Machine Learning)
\end{minipage} & \begin{minipage}[b]{\linewidth}\raggedright
Categorization (Cognitive Modeling)
\end{minipage} \\
\midrule\noalign{}
\endhead
\bottomrule\noalign{}
\endlastfoot
\textbf{Primary Goal} & Maximize out-of-sample accuracy. & Explain human
behavior, including systematic human errors. \\
\textbf{Data Processing} & Often batch-processed over massive, static
datasets. & Inherently \textbf{path-dependent}; learns sequentially
trial-by-trial. \\
\textbf{Constraints} & Hardware limits and computational efficiency. &
Psychological plausibility (bounded memory, selective attention,
cognitive load). \\
\end{longtable}

In short: A machine learning classifier asks, \emph{``What is the best
mathematical way to separate Category A from Category B?''} A cognitive
model asks, \emph{``How do humans separate A from B, and what are the
internal mental representations that allow them to do so?''}

\begin{center}\rule{0.5\linewidth}{0.5pt}\end{center}

\section{The Generalized Context Model
(GCM)}\label{the-generalized-context-model-gcm}

\subsection{Mathematical Foundations and Cognitive
Principles}\label{mathematical-foundations-and-cognitive-principles}

The Generalized Context Model (GCM), developed by Robert Nosofsky
(Nosofsky 1986) in the 1980s, represents one of the most influential
exemplar-based approaches to categorization. Unlike models that abstract
category information into prototypes or rules, the GCM proposes that
people store individual exemplars in memory and make categorization
decisions based on similarity to these stored examples.

\subsubsection{The GCM Mapped to the Categorization
Pipeline}\label{the-gcm-mapped-to-the-categorization-pipeline}

The GCM translates the 6-step cognitive pipeline into a formal
mathematical structure as follows:

\begin{itemize}
\item
  \textbf{1. Input Representation (Geometric Space):} Stimuli are
  represented as continuous points in a multi-dimensional psychological
  space.
\item
  \textbf{2. Attentional Processes (Selective Attention):} Per-dimension
  weights \(w\) (summing to 1) rescale the psychological space before
  distance is computed, prioritizing task-relevant features.
\item
  \textbf{3. Intermediate Representations (Memory for Instances):} The
  model explicitly skips abstraction. Every encountered example is
  stored directly in memory with its category label.
\item
  \textbf{4. Evidential Mechanisms (Similarity Computation):} Distance
  is calculated using a city-block metric (Manhattan distance).
  Similarity between a new stimulus and each stored exemplar is an
  exponentially decreasing function of their distance, governed by a
  sensitivity parameter \(c\). Evidence is then aggregated as average
  similarity for each category.
\item
  \textbf{5. Decision Rules (Choice Probability):} The decision relies
  on the Luce choice axiom. The probability of assigning a stimulus to a
  category is proportional to its summed similarity to exemplars of that
  category, modified by a baseline category bias \(\beta\).
\item
  \textbf{6. Learning (Memory Accumulation):} The model learns
  sequentially by adding each newly categorized and feedback-evaluated
  stimulus directly into its exemplar memory store, altering the
  evidential landscape for future trials.
\end{itemize}

\begin{tcolorbox}[enhanced jigsaw, toprule=.15mm, toptitle=1mm, bottomtitle=1mm, leftrule=.75mm, rightrule=.15mm, colframe=quarto-callout-note-color-frame, left=2mm, colback=white, titlerule=0mm, coltitle=black, colbacktitle=quarto-callout-note-color!10!white, opacityback=0, title=\textcolor{quarto-callout-note-color}{\faInfo}\hspace{0.5em}{Historical Context: From Binary Logic to Psychological Space}, arc=.35mm, opacitybacktitle=0.6, breakable, bottomrule=.15mm]

The evolution of the Generalized Context Model (GCM) reflects a broader
shift in cognitive science: moving from models constrained by
computational simplicity toward those more capable of capturing the
``messy'' richness of human perception.

\subsubsection{The Era of Binary
Constraints}\label{the-era-of-binary-constraints}

In the late 1970s, categorization research was dominated by the
``small-set'' paradigm. Stimuli were often constructed from a few
discrete, binary features (e.g., ``large vs.~small,'' ``red vs.~blue'').
The pioneering \textbf{Context Model} (Medin and Schaffer 1978) was
designed for these environments, calculating similarity based on simple
feature overlap. While revolutionary, it was computationally and
theoretically ``flat'': it couldn't account for the fact that some
dimensions are more perceptually salient than others, or that the
psychological distance between features isn't always linear.

\subsubsection{The Geometric Revolution}\label{the-geometric-revolution}

The introduction of the GCM (Nosofsky 1986) broke these constraints by
synthesizing exemplar theory with \textbf{Multidimensional Scaling
(MDS)}. Instead of binary features, stimuli became points in a
continuous, multi-dimensional psychological space. The core innovation
was the weighted Minkowski metric:
\[d_{ij} = c \left[ \sum w_k |x_{ik} - x_{jk}|^r \right]^{1/r}\] This
formula allowed the model to ``stretch'' or ``shrink'' psychological
space based on selective attention (\(w_k\)). It transformed
categorization from a logical matching process into a geometric one,
where the ``geometry of thought'' could be empirically mapped.

\subsubsection{The ``Ecological'' Turn: Summed vs.~Mean
Similarity}\label{the-ecological-turn-summed-vs.-mean-similarity}

As the field moved toward experiments with larger and unbalanced
category sizes, a debate arose over whether evidence should be summed or
averaged across exemplars. See §Summed vs.~Mean Similarity below for the
full argument and a concrete numerical demonstration; that section
motivates why our implementation adopts the \textbf{mean-similarity}
formulation.

\subsubsection{The Process Revolution: Sequential Accumulation and
Path-Dependency}\label{the-process-revolution-sequential-accumulation-and-path-dependency}

Early models were ``static'': they were used to fit terminal,
steady-state data from the end of an experiment. The conceptual shift
toward process models began by acknowledging that categorization is a
sequential learning experience. The critical first step is simply
modeling memory as it grows: exemplars are accumulated into the memory
store one trial at a time.

This transition highlights the fundamentally \textbf{path-dependent}
nature of cognition: your response at trial \(t\) depends strictly on
the specific sequence of exemplars and feedback you encountered in
trials \(1 \dots t-1\). While later models like \textbf{ALCOVE} (John K.
Kruschke 1992) attempted to model trial-by-trial, error-driven updates
to the attention weights themselves, in my experience such extensions
often yield terrible parameter recovery and fail stringent modern
validation checks. Consequently, our approach grounds the ``Process
Revolution'' in robust sequential exemplar accumulation. We then
conceptually extend this path-dependency by introducing a memory decay
parameter, providing a mechanism for agents to unlearn and adapt in
non-stationary environments without sacrificing model identifiability.
{[}It would be a very good learning exercise to try to build and
validate ALCOVE and other process models, but that is beyond the scope
of this chapter and we will not do it here.{]}

\end{tcolorbox}

\subsubsection{Parameters: Cognitive Commitment, Algebraic Role, Stan
Parameterization}\label{parameters-cognitive-commitment-algebraic-role-stan-parameterization}

\textbf{Attention weights \(w\).} Cognitively, \(w\) is a simplex of
dimension-specific attention weights. Algebraically, \(w\) enters the
distance formula as the per-dimension coefficient in the
attention-weighted city-block metric. At the Stan level \(w\) is
declared as \texttt{simplex{[}N\_features{]}}, which enforces
\(w_m \ge 0\) and \(\sum_m w_m = 1\) directly.

\textbf{Sensitivity \(c\).} Cognitively, \(c\) represents perceptual
discriminability. Algebraically, \(c\) is the exponential decay rate in
\(\eta_{ij} = e^{-c\, d_{ij}}\). At the Stan level, we sample an
unconstrained \texttt{real\ log\_c} and set \texttt{c\ =\ exp(log\_c)}
in \texttt{transformed\ parameters}.

\textbf{Response bias \(\beta\).} Cognitively, \(\beta\) is a baseline
preference for category A independent of stimulus similarity;
\(\beta = 0.5\) is unbiased. In this chapter \(\beta\) also serves as
the cold-start fallback. At the Stan level \(\beta\) is declared as
\texttt{real\textless{}lower=0,\ upper=1\textgreater{}} with a Beta
prior.

\subsection{Mathematical Formulation}\label{mathematical-formulation}

\textbf{1. Distance Calculation}

\[d_{ij} = \left[ \sum_{m=1}^{M} w_m |x_{im} - x_{jm}|^r \right]^{\frac{1}{r}}\]

We use the city-block metric (\(r = 1\)), appropriate for separable
stimulus dimensions:

\begin{Shaded}
\begin{Highlighting}[]
\CommentTok{\# City{-}block distance (r = 1), attention{-}weighted}
\NormalTok{distance }\OtherTok{\textless{}{-}} \ControlFlowTok{function}\NormalTok{(vect1, vect2, w) \{}
  \FunctionTok{sum}\NormalTok{(w }\SpecialCharTok{*} \FunctionTok{abs}\NormalTok{(vect1 }\SpecialCharTok{{-}}\NormalTok{ vect2))}
\NormalTok{\}}
\end{Highlighting}
\end{Shaded}

\textbf{2. Similarity Computation}

\[\eta_{ij} = e^{-c \cdot d_{ij}}\]

\begin{Shaded}
\begin{Highlighting}[]
\CommentTok{\# Exponential similarity kernel}
\NormalTok{similarity }\OtherTok{\textless{}{-}} \ControlFlowTok{function}\NormalTok{(dist\_val, c) \{}
  \FunctionTok{exp}\NormalTok{(}\SpecialCharTok{{-}}\NormalTok{c }\SpecialCharTok{*}\NormalTok{ dist\_val)}
\NormalTok{\}}
\end{Highlighting}
\end{Shaded}

Let's visualize how similarity changes with distance for different
values of \(c\):

\begin{Shaded}
\begin{Highlighting}[]
\NormalTok{dd }\OtherTok{\textless{}{-}} \FunctionTok{expand\_grid}\NormalTok{(}
  \AttributeTok{distance  =} \FunctionTok{seq}\NormalTok{(}\DecValTok{0}\NormalTok{, }\DecValTok{6}\NormalTok{, }\AttributeTok{by =} \FloatTok{0.1}\NormalTok{),}
  \AttributeTok{c         =} \FunctionTok{c}\NormalTok{(}\FloatTok{0.3}\NormalTok{, }\FloatTok{0.7}\NormalTok{, }\FloatTok{1.5}\NormalTok{, }\FloatTok{3.0}\NormalTok{, }\FloatTok{6.0}\NormalTok{, }\FloatTok{10.0}\NormalTok{)}
\NormalTok{) }\SpecialCharTok{|\textgreater{}}
  \FunctionTok{mutate}\NormalTok{(}
    \AttributeTok{sim     =} \FunctionTok{similarity}\NormalTok{(distance, c),}
    \AttributeTok{c\_label =} \FunctionTok{factor}\NormalTok{(c)}
\NormalTok{  )}

\FunctionTok{ggplot}\NormalTok{(dd, }\FunctionTok{aes}\NormalTok{(}\AttributeTok{x =}\NormalTok{ distance, }\AttributeTok{y =}\NormalTok{ sim, }\AttributeTok{color =}\NormalTok{ c\_label, }\AttributeTok{group =}\NormalTok{ c\_label)) }\SpecialCharTok{+}
  \FunctionTok{geom\_line}\NormalTok{(}\AttributeTok{linewidth =} \DecValTok{1}\NormalTok{) }\SpecialCharTok{+}
  \FunctionTok{scale\_color\_viridis\_d}\NormalTok{(}\AttributeTok{option =} \StringTok{"plasma"}\NormalTok{, }\AttributeTok{name =} \StringTok{"Sensitivity (c)"}\NormalTok{) }\SpecialCharTok{+}
  \FunctionTok{labs}\NormalTok{(}
    \AttributeTok{title    =} \StringTok{"Effect of Sensitivity Parameter (c) on Similarity"}\NormalTok{,}
    \AttributeTok{subtitle =} \StringTok{"c is on the natural (original) scale; prior places most mass on c ∈ (0.3, 15)"}\NormalTok{,}
    \AttributeTok{x        =} \StringTok{"Psychological Distance"}\NormalTok{,}
    \AttributeTok{y        =} \StringTok{"Similarity"}
\NormalTok{  ) }\SpecialCharTok{+}
  \FunctionTok{theme}\NormalTok{(}\AttributeTok{legend.position =} \StringTok{"bottom"}\NormalTok{)}
\end{Highlighting}
\end{Shaded}

\begin{center}
\includegraphics[width=0.8\linewidth,height=\textheight,keepaspectratio]{11-CategorizationModelsExemplars_files/figure-pdf/ch12_visualizing_sensitivity-1.pdf}
\end{center}

\textbf{3. Category Response Probability}

\[P(A|i) = \frac{\beta \; \bar\eta_{i,A}}{\beta \; \bar\eta_{i,A} + (1-\beta) \; \bar\eta_{i,B}}, \qquad \bar\eta_{i,A} = \frac{1}{N_A} \sum_{j \in A} \eta_{ij}\]

where \(\bar\eta_{i,A}\) is the \textbf{mean} similarity of the current
stimulus to the exemplars stored under category \(A\), and \(N_A\) is
the count of those stored exemplars. The cold-start fallback applies
when either category has no stored exemplars yet: \(P(A|i) = \beta\).

\subsection{Summed vs.~Mean Similarity: A Caveat on the Canonical
Formulation}\label{summed-vs.-mean-similarity-a-caveat-on-the-canonical-formulation}

The original GCM (Nosofsky 1986) and the ALCOVE extensions (John K.
Kruschke 1992) use the \emph{summed} form \(\sum_{j \in A} \eta_{ij}\)
rather than the mean. That choice carries a subtle but consequential
assumption: whichever category has more stored exemplars acquires a
systematic advantage, independent of how similar any of those exemplars
is to the current stimulus.

A concrete failure. Suppose category A has 12 stored exemplars at
distance roughly \(d = 2\) from the current stimulus (so each
contributes \(\eta_{ij} \approx e^{-2c}\)) and category B has 1 stored
exemplar at distance \(d = 0.2\) (contributing
\(\eta_{ij} \approx e^{-0.2c}\)). At a moderate sensitivity \(c = 1\):

\begin{itemize}
\tightlist
\item
  Summed: \(S_A = 12 \cdot e^{-2} \approx 1.62\) vs
  \(S_B = 1 \cdot e^{-0.2} \approx 0.82\). Category A wins.
\item
  Mean: \(\bar\eta_A = e^{-2} \approx 0.14\) vs
  \(\bar\eta_B = e^{-0.2} \approx 0.82\). Category B wins, which is what
  our intuition expects: the closest neighbour belongs to B.
\end{itemize}

The bias parameter \(\beta\) cannot repair this. Fitting \(\beta\) to
the data will compensate for the \emph{average} frequency imbalance, but
cannot fix it trial-by-trial because the imbalance itself grows and
shrinks as the memory store evolves.

This chapter therefore uses the \textbf{mean-similarity} formulation
throughout, which is the form adopted by F. Gregory Ashby and Maddox
(1993) and by several later papers that explicitly address this point.
At \(N_A = N_B\) (the canonical Kruschke task with balanced training)
the two formulations produce effectively identical behaviour, so
replication of published Kruschke-era results is unaffected. On data
with unequal base rates, partial trial sequences, or the non-stationary
scenarios introduced later in this chapter --- where memory composition
changes meaningfully over trials --- the difference becomes both
measurable and interpretable: with mean similarity, \(\beta\) becomes a
clean response-bias parameter rather than a fitted compensator for
implicit base-rate effects.

\begin{Shaded}
\begin{Highlighting}[]
\CommentTok{\# Toy demonstration of the summed{-}vs{-}mean choice with deliberately}
\CommentTok{\# unbalanced memory contents. 12 cat{-}A exemplars far away, 1 cat{-}B}
\CommentTok{\# exemplar close by. We show choice probabilities under both forms.}
\NormalTok{demo\_memory }\OtherTok{\textless{}{-}} \FunctionTok{tibble}\NormalTok{(}
  \AttributeTok{category =} \FunctionTok{c}\NormalTok{(}\FunctionTok{rep}\NormalTok{(}\DecValTok{0}\NormalTok{, }\DecValTok{12}\NormalTok{), }\DecValTok{1}\NormalTok{),}
  \AttributeTok{sim      =} \FunctionTok{c}\NormalTok{(}\FunctionTok{rep}\NormalTok{(}\FunctionTok{exp}\NormalTok{(}\SpecialCharTok{{-}}\DecValTok{2}\NormalTok{), }\DecValTok{12}\NormalTok{), }\FunctionTok{exp}\NormalTok{(}\SpecialCharTok{{-}}\FloatTok{0.2}\NormalTok{))  }\CommentTok{\# pre{-}computed for c = 1}
\NormalTok{)}
\NormalTok{demo\_sums }\OtherTok{\textless{}{-}}\NormalTok{ demo\_memory }\SpecialCharTok{|\textgreater{}}
  \FunctionTok{group\_by}\NormalTok{(category) }\SpecialCharTok{|\textgreater{}}
  \FunctionTok{summarise}\NormalTok{(}\AttributeTok{sum\_sim =} \FunctionTok{sum}\NormalTok{(sim), }\AttributeTok{mean\_sim =} \FunctionTok{mean}\NormalTok{(sim), }\AttributeTok{.groups =} \StringTok{"drop"}\NormalTok{)}
\NormalTok{demo\_sums }\SpecialCharTok{|\textgreater{}}
  \FunctionTok{mutate}\NormalTok{(}
    \AttributeTok{p\_cat1\_summed =}\NormalTok{ sum\_sim  }\SpecialCharTok{/} \FunctionTok{sum}\NormalTok{(sum\_sim),}
    \AttributeTok{p\_cat1\_mean   =}\NormalTok{ mean\_sim }\SpecialCharTok{/} \FunctionTok{sum}\NormalTok{(mean\_sim)}
\NormalTok{  ) }\SpecialCharTok{|\textgreater{}}
  \FunctionTok{print}\NormalTok{()}
\end{Highlighting}
\end{Shaded}

\begin{verbatim}
# A tibble: 2 x 5
  category sum_sim mean_sim p_cat1_summed p_cat1_mean
     <dbl>   <dbl>    <dbl>         <dbl>       <dbl>
1        0   1.62     0.135         0.665       0.142
2        1   0.819    0.819         0.335       0.858
\end{verbatim}

\begin{Shaded}
\begin{Highlighting}[]
\CommentTok{\# Summed: P(cat 1) ≈ 0.34 — the 12 distant exemplars in category 0 win.}
\CommentTok{\# Mean:   P(cat 1) ≈ 0.85 — the 1 close exemplar in category 1 correctly wins.}
\end{Highlighting}
\end{Shaded}

The demonstration above uses bias \(\beta = 0.5\) (neutral) so the
asymmetry comes purely from the aggregation choice. All figures,
recovery results, and SBC checks in the rest of this chapter use the
mean-similarity form. Readers who want to reproduce a published
summed-similarity study can swap the aggregation line in
\texttt{gcm\_simulate()} (replace each \texttt{mean()} with
\texttt{sum()}) and the corresponding two lines in the Stan model
(remove the division by \texttt{n1} and \texttt{n0}) without any other
changes.

\subsection{The GCM Agent Implementation (R
Simulation)}\label{the-gcm-agent-implementation-r-simulation}

\begin{Shaded}
\begin{Highlighting}[]
\CommentTok{\# Generative GCM model function.}
\NormalTok{gcm\_simulate }\OtherTok{\textless{}{-}} \ControlFlowTok{function}\NormalTok{(w,            }\CommentTok{\# attention weights (sums to 1)}
\NormalTok{                         c,            }\CommentTok{\# sensitivity (natural scale, strictly positive)}
\NormalTok{                         bias,         }\CommentTok{\# response bias towards category 1}
\NormalTok{                         obs,          }\CommentTok{\# matrix of stimulus features (trials x features)}
\NormalTok{                         cat\_feedback, }\CommentTok{\# true category labels (0 or 1)}
                         \AttributeTok{quiet =} \ConstantTok{TRUE}\NormalTok{) \{}
\NormalTok{  ntrials   }\OtherTok{\textless{}{-}} \FunctionTok{nrow}\NormalTok{(obs)}
\NormalTok{  nfeatures }\OtherTok{\textless{}{-}} \FunctionTok{ncol}\NormalTok{(obs)}

\NormalTok{  prob\_cat1    }\OtherTok{\textless{}{-}} \FunctionTok{numeric}\NormalTok{(ntrials)}
\NormalTok{  sim\_response }\OtherTok{\textless{}{-}} \FunctionTok{numeric}\NormalTok{(ntrials)}

\NormalTok{  memory\_obs }\OtherTok{\textless{}{-}} \FunctionTok{matrix}\NormalTok{(}\ConstantTok{NA\_real\_}\NormalTok{, }\AttributeTok{nrow =} \DecValTok{0}\NormalTok{, }\AttributeTok{ncol =}\NormalTok{ nfeatures)}
\NormalTok{  memory\_cat }\OtherTok{\textless{}{-}} \FunctionTok{numeric}\NormalTok{(}\DecValTok{0}\NormalTok{)}

  \ControlFlowTok{for}\NormalTok{ (i }\ControlFlowTok{in} \FunctionTok{seq\_len}\NormalTok{(ntrials)) \{}
    \ControlFlowTok{if}\NormalTok{ (}\SpecialCharTok{!}\NormalTok{quiet }\SpecialCharTok{\&\&}\NormalTok{ i }\SpecialCharTok{\%\%} \DecValTok{20} \SpecialCharTok{==} \DecValTok{0}\NormalTok{) }\FunctionTok{cat}\NormalTok{(}\StringTok{"Simulating trial:"}\NormalTok{, i, }\StringTok{"}\SpecialCharTok{\textbackslash{}n}\StringTok{"}\NormalTok{)}

\NormalTok{    current\_stim }\OtherTok{\textless{}{-}} \FunctionTok{as.numeric}\NormalTok{(obs[i, ])}
\NormalTok{    n\_mem        }\OtherTok{\textless{}{-}} \FunctionTok{nrow}\NormalTok{(memory\_obs)}
\NormalTok{    has\_cat0     }\OtherTok{\textless{}{-}} \FunctionTok{any}\NormalTok{(memory\_cat }\SpecialCharTok{==} \DecValTok{0}\NormalTok{)}
\NormalTok{    has\_cat1     }\OtherTok{\textless{}{-}} \FunctionTok{any}\NormalTok{(memory\_cat }\SpecialCharTok{==} \DecValTok{1}\NormalTok{)}

    \ControlFlowTok{if}\NormalTok{ (n\_mem }\SpecialCharTok{==} \DecValTok{0} \SpecialCharTok{||} \SpecialCharTok{!}\NormalTok{has\_cat0 }\SpecialCharTok{||} \SpecialCharTok{!}\NormalTok{has\_cat1) \{}
\NormalTok{      prob\_cat1[i] }\OtherTok{\textless{}{-}}\NormalTok{ bias}
\NormalTok{    \} }\ControlFlowTok{else}\NormalTok{ \{}
\NormalTok{      sims }\OtherTok{\textless{}{-}} \FunctionTok{vapply}\NormalTok{(}\FunctionTok{seq\_len}\NormalTok{(n\_mem), }\ControlFlowTok{function}\NormalTok{(e) \{}
        \FunctionTok{similarity}\NormalTok{(}\FunctionTok{distance}\NormalTok{(current\_stim, memory\_obs[e, ], w), c)}
\NormalTok{      \}, }\FunctionTok{numeric}\NormalTok{(}\DecValTok{1}\NormalTok{))}

      \CommentTok{\# Mean similarity per category — removes category{-}frequency bias.}
\NormalTok{      m1  }\OtherTok{\textless{}{-}} \FunctionTok{mean}\NormalTok{(sims[memory\_cat }\SpecialCharTok{==} \DecValTok{1}\NormalTok{])}
\NormalTok{      m0  }\OtherTok{\textless{}{-}} \FunctionTok{mean}\NormalTok{(sims[memory\_cat }\SpecialCharTok{==} \DecValTok{0}\NormalTok{])}
\NormalTok{      num }\OtherTok{\textless{}{-}}\NormalTok{ bias }\SpecialCharTok{*}\NormalTok{ m1}
\NormalTok{      den }\OtherTok{\textless{}{-}}\NormalTok{ num }\SpecialCharTok{+}\NormalTok{ (}\DecValTok{1} \SpecialCharTok{{-}}\NormalTok{ bias) }\SpecialCharTok{*}\NormalTok{ m0}

\NormalTok{      prob\_cat1[i] }\OtherTok{\textless{}{-}} \ControlFlowTok{if}\NormalTok{ (den }\SpecialCharTok{\textgreater{}} \FloatTok{1e{-}9}\NormalTok{) num }\SpecialCharTok{/}\NormalTok{ den }\ControlFlowTok{else}\NormalTok{ bias}
\NormalTok{      prob\_cat1[i] }\OtherTok{\textless{}{-}} \FunctionTok{pmax}\NormalTok{(}\FloatTok{1e{-}9}\NormalTok{, }\FunctionTok{pmin}\NormalTok{(}\DecValTok{1} \SpecialCharTok{{-}} \FloatTok{1e{-}9}\NormalTok{, prob\_cat1[i]))}
\NormalTok{    \}}

\NormalTok{    sim\_response[i] }\OtherTok{\textless{}{-}} \FunctionTok{rbinom}\NormalTok{(}\DecValTok{1}\NormalTok{, }\DecValTok{1}\NormalTok{, prob\_cat1[i])}

\NormalTok{    memory\_obs }\OtherTok{\textless{}{-}} \FunctionTok{rbind}\NormalTok{(memory\_obs, current\_stim)}
\NormalTok{    memory\_cat }\OtherTok{\textless{}{-}} \FunctionTok{c}\NormalTok{(memory\_cat, cat\_feedback[i])}
\NormalTok{  \}}

  \FunctionTok{tibble}\NormalTok{(}\AttributeTok{prob\_cat1 =}\NormalTok{ prob\_cat1, }\AttributeTok{sim\_response =}\NormalTok{ sim\_response)}
\NormalTok{\}}
\end{Highlighting}
\end{Shaded}

\begin{center}\rule{0.5\linewidth}{0.5pt}\end{center}

\subsection{Experimental Setup: Why John K. Kruschke
(1993)?}\label{experimental-setup-why-kruschke1993human}

To test our GCM implementation, we need a stimulus environment that is
more complex than a simple linear separation but structured enough to
reveal cognitive principles. We adopt the setup from John K. Kruschke
(1993) for three specific reasons:

\begin{enumerate}
\def\labelenumi{\arabic{enumi}.}
\item
  \textbf{Continuous Dimension Benchmarking:} Unlike earlier
  binary-feature tasks (e.g., Medin and Schaffer (1978)), Kruschke uses
  stimuli (rectangles varying in height and position) that occupy a
  continuous 2D psychological space. This is the natural environment for
  the GCM's Minkowski distance metrics.
\item
  \textbf{Testing Selective Attention:} The task is designed so that
  both dimensions (height and position) are relevant, but their
  importance depends on the specific stimuli. This allows us to observe
  whether our model can ``discover'' the optimal attention weights
  (\(w\)) to minimize categorization error---a core feature of the
  ALCOVE extension to the GCM.
\item
  \textbf{Non-Linearity \& Abstraction:} The category structure (plotted
  below) is not perfectly separable by a single horizontal or vertical
  line. It requires the agent to either learn a diagonal boundary or to
  store the specific exemplars that form the ``corners'' of each
  category. This makes it a diagnostic benchmark for comparing exemplar
  models against the prototype models we'll see in Chapter 12.
\end{enumerate}

\begin{Shaded}
\begin{Highlighting}[]
\NormalTok{stimulus\_info }\OtherTok{\textless{}{-}} \FunctionTok{tibble}\NormalTok{(}
  \AttributeTok{stimulus      =} \FunctionTok{c}\NormalTok{(}\DecValTok{5}\NormalTok{, }\DecValTok{3}\NormalTok{, }\DecValTok{7}\NormalTok{, }\DecValTok{1}\NormalTok{, }\DecValTok{8}\NormalTok{, }\DecValTok{2}\NormalTok{, }\DecValTok{6}\NormalTok{, }\DecValTok{4}\NormalTok{),}
  \AttributeTok{height        =} \FunctionTok{c}\NormalTok{(}\DecValTok{1}\NormalTok{, }\DecValTok{1}\NormalTok{, }\DecValTok{2}\NormalTok{, }\DecValTok{2}\NormalTok{, }\DecValTok{3}\NormalTok{, }\DecValTok{3}\NormalTok{, }\DecValTok{4}\NormalTok{, }\DecValTok{4}\NormalTok{),}
  \AttributeTok{position      =} \FunctionTok{c}\NormalTok{(}\DecValTok{2}\NormalTok{, }\DecValTok{3}\NormalTok{, }\DecValTok{1}\NormalTok{, }\DecValTok{4}\NormalTok{, }\DecValTok{1}\NormalTok{, }\DecValTok{4}\NormalTok{, }\DecValTok{2}\NormalTok{, }\DecValTok{3}\NormalTok{),}
  \AttributeTok{category\_true =} \FunctionTok{c}\NormalTok{(}\DecValTok{0}\NormalTok{, }\DecValTok{0}\NormalTok{, }\DecValTok{1}\NormalTok{, }\DecValTok{0}\NormalTok{, }\DecValTok{1}\NormalTok{, }\DecValTok{0}\NormalTok{, }\DecValTok{1}\NormalTok{, }\DecValTok{1}\NormalTok{)}
\NormalTok{)}

\FunctionTok{ggplot}\NormalTok{(stimulus\_info, }\FunctionTok{aes}\NormalTok{(}\AttributeTok{x =}\NormalTok{ position, }\AttributeTok{y =}\NormalTok{ height,}
                          \AttributeTok{color =} \FunctionTok{factor}\NormalTok{(category\_true), }\AttributeTok{label =}\NormalTok{ stimulus)) }\SpecialCharTok{+}
  \FunctionTok{geom\_point}\NormalTok{(}\AttributeTok{size =} \DecValTok{5}\NormalTok{) }\SpecialCharTok{+}
  \FunctionTok{geom\_text}\NormalTok{(}\AttributeTok{nudge\_y =} \FloatTok{0.18}\NormalTok{) }\SpecialCharTok{+}
  \FunctionTok{scale\_color\_manual}\NormalTok{(}\AttributeTok{values =} \FunctionTok{c}\NormalTok{(}\StringTok{"0"} \OtherTok{=} \StringTok{"\#0072B2"}\NormalTok{, }\StringTok{"1"} \OtherTok{=} \StringTok{"\#D55E00"}\NormalTok{), }\AttributeTok{name =} \StringTok{"Category"}\NormalTok{) }\SpecialCharTok{+}
  \FunctionTok{labs}\NormalTok{(}\AttributeTok{title =} \StringTok{"@kruschke1993human Stimulus Structure"}\NormalTok{, }\AttributeTok{x =} \StringTok{"Position"}\NormalTok{, }\AttributeTok{y =} \StringTok{"Height"}\NormalTok{) }\SpecialCharTok{+}
  \FunctionTok{coord\_fixed}\NormalTok{()}
\end{Highlighting}
\end{Shaded}

\begin{center}
\includegraphics[width=0.8\linewidth,height=\textheight,keepaspectratio]{11-CategorizationModelsExemplars_files/figure-pdf/ch12_experimental_setup-1.pdf}
\end{center}

\begin{Shaded}
\begin{Highlighting}[]
\NormalTok{n\_blocks         }\OtherTok{\textless{}{-}} \DecValTok{8}
\NormalTok{n\_stim\_per\_block }\OtherTok{\textless{}{-}} \FunctionTok{nrow}\NormalTok{(stimulus\_info)}
\NormalTok{total\_trials     }\OtherTok{\textless{}{-}}\NormalTok{ n\_blocks }\SpecialCharTok{*}\NormalTok{ n\_stim\_per\_block}
\end{Highlighting}
\end{Shaded}

\begin{center}\rule{0.5\linewidth}{0.5pt}\end{center}

\subsection{Per-Subject Schedule
Generation}\label{per-subject-schedule-generation}

\begin{Shaded}
\begin{Highlighting}[]
\NormalTok{make\_subject\_schedule }\OtherTok{\textless{}{-}} \ControlFlowTok{function}\NormalTok{(stimulus\_info, n\_blocks, seed) \{}
  \FunctionTok{set.seed}\NormalTok{(seed)}
\NormalTok{  n\_stim }\OtherTok{\textless{}{-}} \FunctionTok{nrow}\NormalTok{(stimulus\_info)}

\NormalTok{  sequence }\OtherTok{\textless{}{-}} \FunctionTok{unlist}\NormalTok{(}\FunctionTok{lapply}\NormalTok{(}\FunctionTok{seq\_len}\NormalTok{(n\_blocks), }\ControlFlowTok{function}\NormalTok{(b) \{}
    \FunctionTok{sample}\NormalTok{(stimulus\_info}\SpecialCharTok{$}\NormalTok{stimulus, n\_stim, }\AttributeTok{replace =} \ConstantTok{FALSE}\NormalTok{)}
\NormalTok{  \}))}

  \FunctionTok{tibble}\NormalTok{(}
    \AttributeTok{trial\_within\_subject =} \FunctionTok{seq\_along}\NormalTok{(sequence),}
    \AttributeTok{block                =} \FunctionTok{rep}\NormalTok{(}\FunctionTok{seq\_len}\NormalTok{(n\_blocks), }\AttributeTok{each =}\NormalTok{ n\_stim),}
    \AttributeTok{stimulus\_id          =}\NormalTok{ sequence}
\NormalTok{  ) }\SpecialCharTok{|\textgreater{}}
    \FunctionTok{left\_join}\NormalTok{(stimulus\_info, }\AttributeTok{by =} \FunctionTok{c}\NormalTok{(}\StringTok{"stimulus\_id"} \OtherTok{=} \StringTok{"stimulus"}\NormalTok{)) }\SpecialCharTok{|\textgreater{}}
    \FunctionTok{rename}\NormalTok{(}\AttributeTok{category\_feedback =}\NormalTok{ category\_true) }\SpecialCharTok{|\textgreater{}}
\NormalTok{    dplyr}\SpecialCharTok{::}\FunctionTok{select}\NormalTok{(trial\_within\_subject, block, stimulus\_id,}
\NormalTok{                  height, position, category\_feedback)}
\NormalTok{\}}
\end{Highlighting}
\end{Shaded}

\begin{center}\rule{0.5\linewidth}{0.5pt}\end{center}

\subsection{Simulating Behavior with the
GCM}\label{simulating-behavior-with-the-gcm}

\begin{Shaded}
\begin{Highlighting}[]
\NormalTok{simulate\_gcm\_agent }\OtherTok{\textless{}{-}} \ControlFlowTok{function}\NormalTok{(agent\_id, w\_true, c\_true, bias\_true,}
\NormalTok{                               stimulus\_info, n\_blocks, subject\_seed) \{}
\NormalTok{  schedule }\OtherTok{\textless{}{-}} \FunctionTok{make\_subject\_schedule}\NormalTok{(stimulus\_info, n\_blocks, }\AttributeTok{seed =}\NormalTok{ subject\_seed)}
\NormalTok{  obs      }\OtherTok{\textless{}{-}}\NormalTok{ schedule }\SpecialCharTok{|\textgreater{}}\NormalTok{ dplyr}\SpecialCharTok{::}\FunctionTok{select}\NormalTok{(height, position)}
\NormalTok{  feedback }\OtherTok{\textless{}{-}}\NormalTok{ schedule}\SpecialCharTok{$}\NormalTok{category\_feedback}

\NormalTok{  sim\_results }\OtherTok{\textless{}{-}} \FunctionTok{gcm\_simulate}\NormalTok{(}
    \AttributeTok{w =}\NormalTok{ w\_true, }\AttributeTok{c =}\NormalTok{ c\_true, }\AttributeTok{bias =}\NormalTok{ bias\_true,}
    \AttributeTok{obs =}\NormalTok{ obs, }\AttributeTok{cat\_feedback =}\NormalTok{ feedback, }\AttributeTok{quiet =} \ConstantTok{TRUE}
\NormalTok{  )}

\NormalTok{  schedule }\SpecialCharTok{|\textgreater{}}
    \FunctionTok{mutate}\NormalTok{(}
      \AttributeTok{agent\_id     =}\NormalTok{ agent\_id,}
      \AttributeTok{w1\_true      =}\NormalTok{ w\_true[}\DecValTok{1}\NormalTok{],}
      \AttributeTok{w2\_true      =}\NormalTok{ w\_true[}\DecValTok{2}\NormalTok{],}
      \AttributeTok{c\_true       =}\NormalTok{ c\_true,}
      \AttributeTok{log\_c\_true   =} \FunctionTok{log}\NormalTok{(c\_true),}
      \AttributeTok{bias\_true    =}\NormalTok{ bias\_true,}
      \AttributeTok{prob\_cat1    =}\NormalTok{ sim\_results}\SpecialCharTok{$}\NormalTok{prob\_cat1,}
      \AttributeTok{sim\_response =}\NormalTok{ sim\_results}\SpecialCharTok{$}\NormalTok{sim\_response,}
      \AttributeTok{correct      =} \FunctionTok{as.integer}\NormalTok{(category\_feedback }\SpecialCharTok{==}\NormalTok{ sim\_response)}
\NormalTok{    ) }\SpecialCharTok{|\textgreater{}}
    \FunctionTok{group\_by}\NormalTok{(agent\_id) }\SpecialCharTok{|\textgreater{}}
    \FunctionTok{mutate}\NormalTok{(}\AttributeTok{performance =} \FunctionTok{cumsum}\NormalTok{(correct) }\SpecialCharTok{/} \FunctionTok{row\_number}\NormalTok{()) }\SpecialCharTok{|\textgreater{}}
    \FunctionTok{ungroup}\NormalTok{()}
\NormalTok{\}}
\end{Highlighting}
\end{Shaded}

\begin{Shaded}
\begin{Highlighting}[]
\NormalTok{param\_grid }\OtherTok{\textless{}{-}} \FunctionTok{expand\_grid}\NormalTok{(}
  \AttributeTok{w\_setting    =} \FunctionTok{list}\NormalTok{(}\FunctionTok{c}\NormalTok{(}\FloatTok{0.5}\NormalTok{, }\FloatTok{0.5}\NormalTok{), }\FunctionTok{c}\NormalTok{(}\FloatTok{0.7}\NormalTok{, }\FloatTok{0.3}\NormalTok{), }\FunctionTok{c}\NormalTok{(}\FloatTok{0.9}\NormalTok{, }\FloatTok{0.1}\NormalTok{)),}
  \AttributeTok{c\_setting    =} \FunctionTok{c}\NormalTok{(}\FloatTok{0.5}\NormalTok{, }\FloatTok{1.0}\NormalTok{, }\FloatTok{2.0}\NormalTok{, }\FloatTok{4.0}\NormalTok{, }\FloatTok{5.0}\NormalTok{, }\FloatTok{15.0}\NormalTok{),}
  \AttributeTok{bias\_setting =} \FloatTok{0.5}
\NormalTok{) }\SpecialCharTok{|\textgreater{}}
  \FunctionTok{mutate}\NormalTok{(}
    \AttributeTok{setting\_id =} \FunctionTok{row\_number}\NormalTok{(),}
    \AttributeTok{w\_label    =} \FunctionTok{map\_chr}\NormalTok{(w\_setting, }\SpecialCharTok{\textasciitilde{}}\FunctionTok{paste}\NormalTok{(}\FunctionTok{round}\NormalTok{(.x, }\DecValTok{1}\NormalTok{), }\AttributeTok{collapse =} \StringTok{","}\NormalTok{)),}
    \AttributeTok{c\_label    =} \FunctionTok{paste0}\NormalTok{(}\StringTok{"c≈"}\NormalTok{, }\FunctionTok{round}\NormalTok{(c\_setting, }\DecValTok{1}\NormalTok{)),}
    \AttributeTok{label      =} \FunctionTok{paste0}\NormalTok{(}\StringTok{"w=("}\NormalTok{, w\_label, }\StringTok{"), "}\NormalTok{, c\_label)}
\NormalTok{  )}

\NormalTok{n\_agents\_per\_setting }\OtherTok{\textless{}{-}} \DecValTok{5}

\NormalTok{simulation\_plan }\OtherTok{\textless{}{-}}\NormalTok{ param\_grid }\SpecialCharTok{|\textgreater{}}
  \FunctionTok{uncount}\NormalTok{(n\_agents\_per\_setting) }\SpecialCharTok{|\textgreater{}}
  \FunctionTok{mutate}\NormalTok{(}\AttributeTok{agent\_id =} \FunctionTok{row\_number}\NormalTok{())}

\NormalTok{sim\_file }\OtherTok{\textless{}{-}} \FunctionTok{here}\NormalTok{(}\StringTok{"simdata"}\NormalTok{, }\StringTok{"ch12\_gcm\_simulated\_responses.csv"}\NormalTok{)}

\ControlFlowTok{if}\NormalTok{ (regenerate\_simulations }\SpecialCharTok{||} \SpecialCharTok{!}\FunctionTok{file.exists}\NormalTok{(sim\_file)) \{}
  \FunctionTok{plan}\NormalTok{(multisession, }\AttributeTok{workers =} \FunctionTok{max}\NormalTok{(}\DecValTok{1}\NormalTok{, }\FunctionTok{availableCores}\NormalTok{() }\SpecialCharTok{{-}} \DecValTok{1}\NormalTok{))}

\NormalTok{  simulated\_responses }\OtherTok{\textless{}{-}} \FunctionTok{future\_pmap\_dfr}\NormalTok{(}
    \FunctionTok{list}\NormalTok{(}
      \AttributeTok{agent\_id  =}\NormalTok{ simulation\_plan}\SpecialCharTok{$}\NormalTok{agent\_id,}
      \AttributeTok{w\_true    =}\NormalTok{ simulation\_plan}\SpecialCharTok{$}\NormalTok{w\_setting,}
      \AttributeTok{c\_true    =}\NormalTok{ simulation\_plan}\SpecialCharTok{$}\NormalTok{c\_setting,}
      \AttributeTok{bias\_true =}\NormalTok{ simulation\_plan}\SpecialCharTok{$}\NormalTok{bias\_setting}
\NormalTok{    ),}
\NormalTok{    simulate\_gcm\_agent,}
    \AttributeTok{stimulus\_info =}\NormalTok{ stimulus\_info,}
    \AttributeTok{n\_blocks      =}\NormalTok{ n\_blocks,}
    \AttributeTok{subject\_seed  =}\NormalTok{ simulation\_plan}\SpecialCharTok{$}\NormalTok{agent\_id,}
    \AttributeTok{.options =} \FunctionTok{furrr\_options}\NormalTok{(}\AttributeTok{seed =} \ConstantTok{TRUE}\NormalTok{)}
\NormalTok{  )}

  \FunctionTok{plan}\NormalTok{(sequential)}
  \FunctionTok{write\_csv}\NormalTok{(simulated\_responses, sim\_file)}
  \FunctionTok{cat}\NormalTok{(}\StringTok{"Simulation complete. Results saved.}\SpecialCharTok{\textbackslash{}n}\StringTok{"}\NormalTok{)}
\NormalTok{\} }\ControlFlowTok{else}\NormalTok{ \{}
\NormalTok{  simulated\_responses }\OtherTok{\textless{}{-}} \FunctionTok{read\_csv}\NormalTok{(sim\_file, }\AttributeTok{show\_col\_types =} \ConstantTok{FALSE}\NormalTok{)}
  \FunctionTok{cat}\NormalTok{(}\StringTok{"Loaded existing simulation results.}\SpecialCharTok{\textbackslash{}n}\StringTok{"}\NormalTok{)}
\NormalTok{\}}
\end{Highlighting}
\end{Shaded}

\begin{verbatim}
Loaded existing simulation results.
\end{verbatim}

\begin{Shaded}
\begin{Highlighting}[]
\NormalTok{simulated\_responses\_labeled }\OtherTok{\textless{}{-}}\NormalTok{ simulated\_responses }\SpecialCharTok{|\textgreater{}}
  \FunctionTok{mutate}\NormalTok{(}
    \AttributeTok{w\_label =} \FunctionTok{paste0}\NormalTok{(}\StringTok{"w=("}\NormalTok{, }\FunctionTok{round}\NormalTok{(w1\_true, }\DecValTok{1}\NormalTok{), }\StringTok{","}\NormalTok{, }\FunctionTok{round}\NormalTok{(w2\_true, }\DecValTok{1}\NormalTok{), }\StringTok{")"}\NormalTok{),}
    \AttributeTok{c\_label =} \FunctionTok{paste0}\NormalTok{(}\StringTok{"c≈"}\NormalTok{, }\FunctionTok{round}\NormalTok{(c\_true, }\DecValTok{1}\NormalTok{)),}
    \CommentTok{\# Order the c\_label factor based on the numeric c\_true values}
    \AttributeTok{c\_label =}\NormalTok{ forcats}\SpecialCharTok{::}\FunctionTok{fct\_reorder}\NormalTok{(c\_label, c\_true) }
\NormalTok{  )}

\NormalTok{p\_c\_effect }\OtherTok{\textless{}{-}} \FunctionTok{ggplot}\NormalTok{(}
\NormalTok{  simulated\_responses\_labeled,}
  \FunctionTok{aes}\NormalTok{(}\AttributeTok{x =}\NormalTok{ trial\_within\_subject, }\AttributeTok{y =}\NormalTok{ performance,}
      \AttributeTok{color =} \FunctionTok{factor}\NormalTok{(}\FunctionTok{round}\NormalTok{(c\_true, }\DecValTok{1}\NormalTok{)),}
      \AttributeTok{group =} \FunctionTok{interaction}\NormalTok{(agent\_id, }\FunctionTok{round}\NormalTok{(c\_true, }\DecValTok{1}\NormalTok{)))}
\NormalTok{) }\SpecialCharTok{+}
  \FunctionTok{stat\_summary}\NormalTok{(}\AttributeTok{fun =}\NormalTok{ mean, }\AttributeTok{geom =} \StringTok{"line"}\NormalTok{,}
               \FunctionTok{aes}\NormalTok{(}\AttributeTok{group =} \FunctionTok{factor}\NormalTok{(}\FunctionTok{round}\NormalTok{(c\_true, }\DecValTok{1}\NormalTok{))), }\AttributeTok{linewidth =} \FloatTok{1.2}\NormalTok{) }\SpecialCharTok{+}
  \FunctionTok{facet\_wrap}\NormalTok{(}\SpecialCharTok{\textasciitilde{}}\NormalTok{w\_label) }\SpecialCharTok{+}
  \FunctionTok{scale\_color\_viridis\_d}\NormalTok{(}\AttributeTok{option =} \StringTok{"plasma"}\NormalTok{, }\AttributeTok{name =} \StringTok{"Sensitivity (c)"}\NormalTok{) }\SpecialCharTok{+}
  \FunctionTok{labs}\NormalTok{(}
    \AttributeTok{title    =} \StringTok{"Effect of Sensitivity (c) on GCM Learning"}\NormalTok{,}
    \AttributeTok{subtitle =} \StringTok{"Mean across simulated agents"}\NormalTok{,}
    \AttributeTok{x =} \StringTok{"Trial"}\NormalTok{, }\AttributeTok{y =} \StringTok{"Cumulative Accuracy"}
\NormalTok{  ) }\SpecialCharTok{+}
  \FunctionTok{theme}\NormalTok{(}\AttributeTok{legend.position =} \StringTok{"bottom"}\NormalTok{)}

\NormalTok{p\_w\_effect }\OtherTok{\textless{}{-}} \FunctionTok{ggplot}\NormalTok{(}
\NormalTok{  simulated\_responses\_labeled,}
  \FunctionTok{aes}\NormalTok{(}\AttributeTok{x =}\NormalTok{ trial\_within\_subject, }\AttributeTok{y =}\NormalTok{ performance,}
      \AttributeTok{color =}\NormalTok{ w\_label, }\AttributeTok{group =} \FunctionTok{interaction}\NormalTok{(agent\_id, w\_label))}
\NormalTok{) }\SpecialCharTok{+}
  \FunctionTok{stat\_summary}\NormalTok{(}\AttributeTok{fun =}\NormalTok{ mean, }\AttributeTok{geom =} \StringTok{"line"}\NormalTok{, }\FunctionTok{aes}\NormalTok{(}\AttributeTok{group =}\NormalTok{ w\_label), }\AttributeTok{linewidth =} \FloatTok{1.2}\NormalTok{) }\SpecialCharTok{+}
  \FunctionTok{facet\_wrap}\NormalTok{(}\SpecialCharTok{\textasciitilde{}}\NormalTok{c\_label) }\SpecialCharTok{+}
  \FunctionTok{scale\_color\_brewer}\NormalTok{(}\AttributeTok{palette =} \StringTok{"Set1"}\NormalTok{, }\AttributeTok{name =} \StringTok{"Attention Weights"}\NormalTok{) }\SpecialCharTok{+}
  \FunctionTok{labs}\NormalTok{(}
    \AttributeTok{title    =} \StringTok{"Effect of Attention Weights (w) on GCM Learning"}\NormalTok{,}
    \AttributeTok{subtitle =} \StringTok{"Mean across simulated agents"}\NormalTok{,}
    \AttributeTok{x =} \StringTok{"Trial"}\NormalTok{, }\AttributeTok{y =} \StringTok{"Cumulative Accuracy"}
\NormalTok{  ) }\SpecialCharTok{+}
  \FunctionTok{theme}\NormalTok{(}\AttributeTok{legend.position =} \StringTok{"bottom"}\NormalTok{)}

\NormalTok{p\_c\_effect}
\end{Highlighting}
\end{Shaded}

\begin{center}
\includegraphics[width=0.8\linewidth,height=\textheight,keepaspectratio]{11-CategorizationModelsExemplars_files/figure-pdf/ch12_visualize_parameter_effects-1.pdf}
\end{center}

\begin{Shaded}
\begin{Highlighting}[]
\NormalTok{p\_w\_effect}
\end{Highlighting}
\end{Shaded}

\begin{center}
\includegraphics[width=0.8\linewidth,height=\textheight,keepaspectratio]{11-CategorizationModelsExemplars_files/figure-pdf/ch12_visualize_parameter_effects-2.pdf}
\end{center}

\begin{center}\rule{0.5\linewidth}{0.5pt}\end{center}

\subsection{Prior Predictive Check}\label{prior-predictive-check-2}

We need to define the priors for the Stan model already here:

\begin{itemize}
\tightlist
\item
  \(w \sim \text{Dirichlet}(1, 1)\) --- uniform over the simplex
\item
  \(\log c \sim \text{Normal}(\log(2), 1)\) --- implies
  \(c \in (0.27, 14.8)\) at \(\pm 2\) SD
\item
  \(\text{bias} \sim \text{Beta}(1, 1)\) --- uniform over \((0, 1)\)
\end{itemize}

\begin{Shaded}
\begin{Highlighting}[]
\NormalTok{n\_prior\_samples }\OtherTok{\textless{}{-}} \DecValTok{500}

\NormalTok{prior\_draws }\OtherTok{\textless{}{-}} \FunctionTok{tibble}\NormalTok{(}
  \AttributeTok{sample\_id =} \FunctionTok{seq\_len}\NormalTok{(n\_prior\_samples),}
  \AttributeTok{log\_c     =} \FunctionTok{rnorm}\NormalTok{(n\_prior\_samples, LOG\_C\_PRIOR\_MEAN, LOG\_C\_PRIOR\_SD),}
  \AttributeTok{w1        =} \FunctionTok{runif}\NormalTok{(n\_prior\_samples),}
  \AttributeTok{bias      =} \FunctionTok{rbeta}\NormalTok{(n\_prior\_samples, }\DecValTok{1}\NormalTok{, }\DecValTok{1}\NormalTok{)}
\NormalTok{) }\SpecialCharTok{|\textgreater{}}
  \FunctionTok{mutate}\NormalTok{(}\AttributeTok{c\_true =} \FunctionTok{exp}\NormalTok{(log\_c), }\AttributeTok{w2 =} \DecValTok{1} \SpecialCharTok{{-}}\NormalTok{ w1)}

\NormalTok{ppc\_schedule }\OtherTok{\textless{}{-}} \FunctionTok{make\_subject\_schedule}\NormalTok{(stimulus\_info, n\_blocks, }\AttributeTok{seed =} \DecValTok{999}\NormalTok{)}

\NormalTok{prior\_pred\_curves }\OtherTok{\textless{}{-}} \FunctionTok{pmap\_dfr}\NormalTok{(}
  \FunctionTok{list}\NormalTok{(}
    \AttributeTok{sample\_id =}\NormalTok{ prior\_draws}\SpecialCharTok{$}\NormalTok{sample\_id,}
    \AttributeTok{log\_c     =}\NormalTok{ prior\_draws}\SpecialCharTok{$}\NormalTok{log\_c,}
    \AttributeTok{w1        =}\NormalTok{ prior\_draws}\SpecialCharTok{$}\NormalTok{w1,}
    \AttributeTok{bias      =}\NormalTok{ prior\_draws}\SpecialCharTok{$}\NormalTok{bias}
\NormalTok{  ),}
  \ControlFlowTok{function}\NormalTok{(sample\_id, log\_c, w1, bias) \{}
\NormalTok{    c }\OtherTok{\textless{}{-}} \FunctionTok{exp}\NormalTok{(log\_c)}
\NormalTok{    sim }\OtherTok{\textless{}{-}} \FunctionTok{gcm\_simulate}\NormalTok{(}
      \AttributeTok{w =} \FunctionTok{c}\NormalTok{(w1, }\DecValTok{1} \SpecialCharTok{{-}}\NormalTok{ w1), }\AttributeTok{c =}\NormalTok{ c, }\AttributeTok{bias =}\NormalTok{ bias,}
      \AttributeTok{obs          =}\NormalTok{ ppc\_schedule }\SpecialCharTok{|\textgreater{}}\NormalTok{ dplyr}\SpecialCharTok{::}\FunctionTok{select}\NormalTok{(height, position),}
      \AttributeTok{cat\_feedback =}\NormalTok{ ppc\_schedule}\SpecialCharTok{$}\NormalTok{category\_feedback}
\NormalTok{    )}
    \FunctionTok{tibble}\NormalTok{(}
      \AttributeTok{sample\_id =}\NormalTok{ sample\_id,}
      \AttributeTok{trial     =} \FunctionTok{seq\_len}\NormalTok{(}\FunctionTok{nrow}\NormalTok{(ppc\_schedule)),}
      \AttributeTok{correct   =} \FunctionTok{as.integer}\NormalTok{(sim}\SpecialCharTok{$}\NormalTok{sim\_response }\SpecialCharTok{==}\NormalTok{ ppc\_schedule}\SpecialCharTok{$}\NormalTok{category\_feedback)}
\NormalTok{    ) }\SpecialCharTok{|\textgreater{}}
      \FunctionTok{mutate}\NormalTok{(}\AttributeTok{cum\_acc =} \FunctionTok{cumsum}\NormalTok{(correct) }\SpecialCharTok{/} \FunctionTok{row\_number}\NormalTok{())}
\NormalTok{  \}}
\NormalTok{)}

\NormalTok{prior\_pred\_summary }\OtherTok{\textless{}{-}}\NormalTok{ prior\_pred\_curves }\SpecialCharTok{|\textgreater{}}
  \FunctionTok{group\_by}\NormalTok{(trial) }\SpecialCharTok{|\textgreater{}}
  \FunctionTok{summarise}\NormalTok{(}
    \AttributeTok{q05 =} \FunctionTok{quantile}\NormalTok{(cum\_acc, }\FloatTok{0.05}\NormalTok{),}
    \AttributeTok{q25 =} \FunctionTok{quantile}\NormalTok{(cum\_acc, }\FloatTok{0.25}\NormalTok{),}
    \AttributeTok{q50 =} \FunctionTok{median}\NormalTok{(cum\_acc),}
    \AttributeTok{q75 =} \FunctionTok{quantile}\NormalTok{(cum\_acc, }\FloatTok{0.75}\NormalTok{),}
    \AttributeTok{q95 =} \FunctionTok{quantile}\NormalTok{(cum\_acc, }\FloatTok{0.95}\NormalTok{),}
    \AttributeTok{.groups =} \StringTok{"drop"}
\NormalTok{  )}

\FunctionTok{ggplot}\NormalTok{(prior\_pred\_summary, }\FunctionTok{aes}\NormalTok{(}\AttributeTok{x =}\NormalTok{ trial)) }\SpecialCharTok{+}
  \FunctionTok{geom\_ribbon}\NormalTok{(}\FunctionTok{aes}\NormalTok{(}\AttributeTok{ymin =}\NormalTok{ q05, }\AttributeTok{ymax =}\NormalTok{ q95), }\AttributeTok{fill =} \StringTok{"\#56B4E9"}\NormalTok{, }\AttributeTok{alpha =} \FloatTok{0.20}\NormalTok{) }\SpecialCharTok{+}
  \FunctionTok{geom\_ribbon}\NormalTok{(}\FunctionTok{aes}\NormalTok{(}\AttributeTok{ymin =}\NormalTok{ q25, }\AttributeTok{ymax =}\NormalTok{ q75), }\AttributeTok{fill =} \StringTok{"\#56B4E9"}\NormalTok{, }\AttributeTok{alpha =} \FloatTok{0.40}\NormalTok{) }\SpecialCharTok{+}
  \FunctionTok{geom\_line}\NormalTok{(}\FunctionTok{aes}\NormalTok{(}\AttributeTok{y =}\NormalTok{ q50), }\AttributeTok{color =} \StringTok{"\#0072B2"}\NormalTok{, }\AttributeTok{linewidth =} \DecValTok{1}\NormalTok{) }\SpecialCharTok{+}
  \FunctionTok{geom\_hline}\NormalTok{(}\AttributeTok{yintercept =} \FloatTok{0.5}\NormalTok{, }\AttributeTok{linetype =} \StringTok{"dashed"}\NormalTok{, }\AttributeTok{color =} \StringTok{"grey40"}\NormalTok{) }\SpecialCharTok{+}
  \FunctionTok{scale\_y\_continuous}\NormalTok{(}\AttributeTok{limits =} \FunctionTok{c}\NormalTok{(}\DecValTok{0}\NormalTok{, }\DecValTok{1}\NormalTok{)) }\SpecialCharTok{+}
  \FunctionTok{labs}\NormalTok{(}
    \AttributeTok{title    =} \StringTok{"Prior Predictive Check: GCM Learning Curves"}\NormalTok{,}
    \AttributeTok{subtitle =} \StringTok{"Ribbons: 50\% and 90\% prior predictive intervals; line: median}\SpecialCharTok{\textbackslash{}n}\StringTok{Priors: log(c) \textasciitilde{} Normal(log 2, 1),  w \textasciitilde{} Dirichlet(1,1),  bias \textasciitilde{} Beta(1,1)"}\NormalTok{,}
    \AttributeTok{x =} \StringTok{"Trial"}\NormalTok{, }\AttributeTok{y =} \StringTok{"Cumulative Accuracy"}
\NormalTok{  )}
\end{Highlighting}
\end{Shaded}

\begin{center}
\includegraphics[width=0.8\linewidth,height=\textheight,keepaspectratio]{11-CategorizationModelsExemplars_files/figure-pdf/ch12_prior_predictive-1.pdf}
\end{center}

\begin{center}\rule{0.5\linewidth}{0.5pt}\end{center}

\subsection{Implementing the GCM in Stan for Parameter
Estimation}\label{implementing-the-gcm-in-stan-for-parameter-estimation}

\subsubsection{\texorpdfstring{Architectural choice:
\texttt{prob\_cat1{[}i{]}} in
\texttt{transformed\ parameters}}{Architectural choice: prob\_cat1{[}i{]} in transformed parameters}}\label{architectural-choice-prob_cat1i-in-transformed-parameters}

Because the GCM's per-trial choice probability \(p_i\) is a
deterministic function of the parameters \((w, c, \beta)\) and the data
sequence, it belongs in \texttt{transformed\ parameters} as a single
source of truth. The \texttt{model\{\}} block then becomes a one-line
vectorised Bernoulli likelihood, and \texttt{generated\ quantities\{\}}
becomes a lookup against the already-computed \texttt{prob\_cat1}
vector.

\begin{Shaded}
\begin{Highlighting}[]
\FunctionTok{\textbackslash{}usetikzlibrary}\NormalTok{\{bayesnet\}}
\KeywordTok{\textbackslash{}begin}\NormalTok{\{}\ExtensionTok{tikzpicture}\NormalTok{\}}
  \CommentTok{\% ── Hyperparameters (const) ─────────────────────────────}
  \FunctionTok{\textbackslash{}node}\NormalTok{[const] (alpha\_w) at (0, 0)   \{}\SpecialStringTok{$}\SpecialCharTok{\textbackslash{}boldsymbol}\SpecialStringTok{\{}\SpecialCharTok{\textbackslash{}alpha}\SpecialStringTok{\}\_w$}\NormalTok{\};}
  \FunctionTok{\textbackslash{}node}\NormalTok{[const] (hyp\_c)   at (4, 0)   \{}\SpecialStringTok{$}\SpecialCharTok{\textbackslash{}mu}\SpecialStringTok{\_c,}\SpecialCharTok{\textbackslash{},\textbackslash{}sigma}\SpecialStringTok{\_c$}\NormalTok{\};}
  \FunctionTok{\textbackslash{}node}\NormalTok{[const] (hyp\_b)   at (8, 0)   \{}\SpecialStringTok{$}\SpecialCharTok{\textbackslash{}alpha}\SpecialStringTok{\_}\SpecialCharTok{\textbackslash{}beta}\SpecialStringTok{,}\SpecialCharTok{\textbackslash{},\textbackslash{}beta}\SpecialStringTok{\_}\SpecialCharTok{\textbackslash{}beta}\SpecialStringTok{$}\NormalTok{\};}

  \CommentTok{\% ── Latent parameters ───────────────────────────────────}
  \FunctionTok{\textbackslash{}node}\NormalTok{[latent] (w)    at (0, {-}2)  \{}\SpecialStringTok{$}\SpecialCharTok{\textbackslash{}mathbf}\SpecialStringTok{\{w\}$}\NormalTok{\};}
  \FunctionTok{\textbackslash{}node}\NormalTok{[latent] (logc) at (4, {-}2)  \{}\SpecialStringTok{$}\SpecialCharTok{\textbackslash{}log}\SpecialStringTok{ c$}\NormalTok{\};}
  \FunctionTok{\textbackslash{}node}\NormalTok{[latent] (bias) at (8, {-}2)  \{}\SpecialStringTok{$}\SpecialCharTok{\textbackslash{}beta}\SpecialStringTok{$}\NormalTok{\};}

  \CommentTok{\% ── Trial{-}level nodes (inside plate) ────────────────────}
  \FunctionTok{\textbackslash{}node}\NormalTok{[obs] (xi) at (2, {-}4.5) \{}\SpecialStringTok{$}\SpecialCharTok{\textbackslash{}mathbf}\SpecialStringTok{\{x\}\_i$}\NormalTok{\};}
  \FunctionTok{\textbackslash{}node}\NormalTok{[obs] (zi) at (4, {-}4.5) \{}\SpecialStringTok{$z\_\{\textless{}i\}$}\NormalTok{\};}
  \FunctionTok{\textbackslash{}node}\NormalTok{[det] (pi) at (6, {-}4.5) \{}\SpecialStringTok{$p\_i$}\NormalTok{\};}
  \FunctionTok{\textbackslash{}node}\NormalTok{[obs] (yi) at (6, {-}6.5) \{}\SpecialStringTok{$y\_i$}\NormalTok{\};}

  \CommentTok{\% ── Edges ───────────────────────────────────────────────}
  \FunctionTok{\textbackslash{}edge}\NormalTok{ \{alpha\_w\}        \{w\};}
  \FunctionTok{\textbackslash{}edge}\NormalTok{ \{hyp\_c\}          \{logc\};}
  \FunctionTok{\textbackslash{}edge}\NormalTok{ \{hyp\_b\}          \{bias\};}
  \FunctionTok{\textbackslash{}edge}\NormalTok{ \{w, logc, bias\}  \{pi\};}
  \FunctionTok{\textbackslash{}edge}\NormalTok{ \{xi, zi\}         \{pi\};}
  \FunctionTok{\textbackslash{}edge}\NormalTok{ \{pi\}             \{yi\};}

  \CommentTok{\% ── Plate ───────────────────────────────────────────────}
  \FunctionTok{\textbackslash{}plate}\NormalTok{ \{trial\} \{(xi)(zi)(pi)(yi)\} \{}\SpecialStringTok{$N$}\NormalTok{ trials\};}
\KeywordTok{\textbackslash{}end}\NormalTok{\{}\ExtensionTok{tikzpicture}\NormalTok{\}}
\end{Highlighting}
\end{Shaded}

\begin{figure}[H]

{\centering \includegraphics[width=0.7\linewidth,height=\textheight,keepaspectratio]{11-CategorizationModelsExemplars_files/figure-pdf/gcm-single-dag-1.pdf}

}

\caption{DAG for the Single-Subject GCM. Hyperparameters (dots) feed
into latent parameters (circles): attention weights \(\mathbf{w}\), log
sensitivity \(\log c\), and response bias \(\beta\). Inside the plate,
the stimulus \(\mathbf{x}_i\) and accumulated exemplar memory \(z_{<i}\)
(past category feedback) combine with the parameters to yield the
deterministic choice probability \(p_i\) (double circle), from which the
response \(y_i\) is drawn.}

\end{figure}%

The DAG formalizes the four-stage structure from §Mathematical
Formulation above: priors on \((w, \log c, \beta)\) feed into
attention-weighted city-block distances, which yield per-category mean
similarities, which enter the Luce choice rule to produce \(p_i\); the
response \(y_i\) is then a Bernoulli draw.

\begin{Shaded}
\begin{Highlighting}[]
\NormalTok{gcm\_single\_stan }\OtherTok{\textless{}{-}} \StringTok{"}
\StringTok{// Generalized Context Model — Single Subject (refactored architecture)}

\StringTok{data \{}
\StringTok{  int\textless{}lower=1\textgreater{} N\_total;}
\StringTok{  int\textless{}lower=1\textgreater{} N\_features;}
\StringTok{  array[N\_total] int\textless{}lower=0, upper=1\textgreater{} y;}
\StringTok{  array[N\_total, N\_features] real obs;}
\StringTok{  array[N\_total] int\textless{}lower=0, upper=1\textgreater{} cat\_feedback;}

\StringTok{  vector[N\_features] prior\_w\_alpha;}
\StringTok{  real prior\_log\_c\_mu;}
\StringTok{  real\textless{}lower=0\textgreater{} prior\_log\_c\_sigma;}
\StringTok{  real\textless{}lower=0\textgreater{} prior\_bias\_alpha;}
\StringTok{  real\textless{}lower=0\textgreater{} prior\_bias\_beta;}
\StringTok{  int\textless{}lower=0, upper=1\textgreater{} run\_diagnostics;}
\StringTok{\}}

\StringTok{transformed data \{}
\StringTok{  array[N\_total, N\_total, N\_features] real abs\_diff;}
\StringTok{  for (i in 1:N\_total) \{}
\StringTok{    for (j in 1:N\_total) \{}
\StringTok{      for (f in 1:N\_features) \{}
\StringTok{        abs\_diff[i, j, f] = abs(obs[i, f] {-} obs[j, f]);}
\StringTok{      \}}
\StringTok{    \}}
\StringTok{  \}}
\StringTok{\}}

\StringTok{parameters \{}
\StringTok{  simplex[N\_features] w;}
\StringTok{  real log\_c;}
\StringTok{  real\textless{}lower=0, upper=1\textgreater{} bias;}
\StringTok{\}}

\StringTok{transformed parameters \{}
\StringTok{  real\textless{}lower=0\textgreater{} c = exp(log\_c);}
\StringTok{  vector\textless{}lower=1e{-}9, upper=1{-}1e{-}9\textgreater{}[N\_total] prob\_cat1;}

\StringTok{  \{}
\StringTok{    array[N\_total] int memory\_trial\_idx;}
\StringTok{    array[N\_total] int memory\_cat;}
\StringTok{    int n\_mem = 0;}

\StringTok{    for (i in 1:N\_total) \{}
\StringTok{      real p\_i;}
\StringTok{      int has\_cat0 = 0;}
\StringTok{      int has\_cat1 = 0;}

\StringTok{      for (k in 1:n\_mem) \{}
\StringTok{        if (memory\_cat[k] == 0) has\_cat0 = 1;}
\StringTok{        if (memory\_cat[k] == 1) has\_cat1 = 1;}
\StringTok{      \}}

\StringTok{      if (n\_mem == 0 || has\_cat0 == 0 || has\_cat1 == 0) \{}
\StringTok{        p\_i = bias;}
\StringTok{      \} else \{}
\StringTok{        real s1 = 0; real s0 = 0;}
\StringTok{        int n1 = 0;  int n0 = 0;}
\StringTok{        for (e in 1:n\_mem) \{}
\StringTok{          real d = 0;}
\StringTok{          int past\_i = memory\_trial\_idx[e];}
\StringTok{          for (f in 1:N\_features)}
\StringTok{            d += w[f] * abs\_diff[i, past\_i, f];}
\StringTok{          real sim = exp({-}c * d);}
\StringTok{          if (memory\_cat[e] == 1) \{ s1 += sim; n1 += 1; \}}
\StringTok{          else                    \{ s0 += sim; n0 += 1; \}}
\StringTok{        \}}
\StringTok{        real m1 = s1 / n1;}
\StringTok{        real m0 = s0 / n0;}
\StringTok{        real num = bias * m1;}
\StringTok{        real den = num + (1 {-} bias) * m0;}
\StringTok{        p\_i = (den \textgreater{} 1e{-}9) ? num / den : bias;}
\StringTok{      \}}

\StringTok{      prob\_cat1[i] = fmax(1e{-}9, fmin(1 {-} 1e{-}9, p\_i));}

\StringTok{      n\_mem += 1;}
\StringTok{      memory\_trial\_idx[n\_mem] = i;}
\StringTok{      memory\_cat[n\_mem] = cat\_feedback[i];}
\StringTok{    \}}
\StringTok{  \}}
\StringTok{\}}

\StringTok{model \{}
\StringTok{  target += dirichlet\_lpdf(w    | prior\_w\_alpha);}
\StringTok{  target += normal\_lpdf(log\_c   | prior\_log\_c\_mu, prior\_log\_c\_sigma);}
\StringTok{  target += beta\_lpdf(bias      | prior\_bias\_alpha,  prior\_bias\_beta);}

\StringTok{  target += bernoulli\_lpmf(y | prob\_cat1);}
\StringTok{\}}

\StringTok{generated quantities \{}
\StringTok{  vector[N\_total] log\_lik;}
\StringTok{  real lprior;}
\StringTok{  if (run\_diagnostics) \{}
\StringTok{    for (i in 1:N\_total)}
\StringTok{      log\_lik[i] = bernoulli\_lpmf(y[i] | prob\_cat1[i]);}
\StringTok{  \}}
\StringTok{  lprior = dirichlet\_lpdf(w | prior\_w\_alpha) +}
\StringTok{           normal\_lpdf(log\_c | prior\_log\_c\_mu, prior\_log\_c\_sigma) +}
\StringTok{           beta\_lpdf(bias | prior\_bias\_alpha, prior\_bias\_beta);}
\StringTok{\}}
\StringTok{"}

\NormalTok{stan\_file\_gcm\_single }\OtherTok{\textless{}{-}} \StringTok{"stan/ch12\_gcm\_single.stan"}
\FunctionTok{write\_stan\_file}\NormalTok{(gcm\_single\_stan, }\AttributeTok{dir =} \StringTok{"stan/"}\NormalTok{, }\AttributeTok{basename =} \StringTok{"ch12\_gcm\_single.stan"}\NormalTok{)}
\end{Highlighting}
\end{Shaded}

\begin{verbatim}
[1] "/Users/au209589/Dropbox/Teaching/AdvancedCognitiveModeling23_book/stan/ch12_gcm_single.stan"
\end{verbatim}

\begin{Shaded}
\begin{Highlighting}[]
\NormalTok{mod\_gcm\_single }\OtherTok{\textless{}{-}} \FunctionTok{cmdstan\_model}\NormalTok{(stan\_file\_gcm\_single)}
\end{Highlighting}
\end{Shaded}

\subsubsection{What each block does:}\label{what-each-block-does}

\begin{itemize}
\tightlist
\item
  \textbf{\texttt{transformed\ data}}: Builds the
  \texttt{abs\_diff{[}i,\ j,\ f{]}} tensor. Evaluated exactly once per
  chain.
\item
  \textbf{\texttt{parameters}}: \texttt{simplex{[}N\_features{]}\ w},
  unconstrained \texttt{real\ log\_c}, and
  \texttt{real\textless{}lower=0,\ upper=1\textgreater{}\ bias}.
\item
  \textbf{\texttt{transformed\ parameters}}: Brings back the constraints
  on the parameters, and caches \texttt{prob\_cat1{[}i{]}} in a single
  forward pass over trials.
\item
  \textbf{\texttt{model}}: Priors followed by a vectorised one-line
  likelihood.
\item
  \textbf{\texttt{generated\ quantities}}: \texttt{log\_lik{[}i{]}} as a
  per-trial lookup, and \texttt{lprior} for \texttt{priorsense}.
\end{itemize}

\begin{quote}
A note on reduce\_sum and parallelization: The GCM's inner loop cannot
be parallelized with reduce\_sum because trial \(t\)'s contribution to
target depends on the memory state accumulated sequentially from trials
\(1, \ldots, t-1\).
\end{quote}

\subsubsection{Fitting the Single-Subject
GCM}\label{fitting-the-single-subject-gcm}

\begin{Shaded}
\begin{Highlighting}[]
\CommentTok{\# Select one simulated agent: w ≈ (0.9, 0.1), c ≈ 2}
\NormalTok{agent\_to\_fit }\OtherTok{\textless{}{-}}\NormalTok{ simulated\_responses }\SpecialCharTok{|\textgreater{}}
  \FunctionTok{filter}\NormalTok{(}
    \FunctionTok{round}\NormalTok{(w1\_true, }\DecValTok{1}\NormalTok{) }\SpecialCharTok{==} \FloatTok{0.9}\NormalTok{,}
    \FunctionTok{abs}\NormalTok{(log\_c\_true }\SpecialCharTok{{-}} \FunctionTok{log}\NormalTok{(}\FloatTok{2.0}\NormalTok{)) }\SpecialCharTok{==} \FunctionTok{min}\NormalTok{(}\FunctionTok{abs}\NormalTok{(log\_c\_true }\SpecialCharTok{{-}} \FunctionTok{log}\NormalTok{(}\FloatTok{2.0}\NormalTok{)))}
\NormalTok{  ) }\SpecialCharTok{|\textgreater{}}
  \FunctionTok{filter}\NormalTok{(agent\_id }\SpecialCharTok{==} \FunctionTok{min}\NormalTok{(agent\_id)) }\SpecialCharTok{|\textgreater{}}
  \FunctionTok{slice}\NormalTok{(}\DecValTok{1}\SpecialCharTok{:}\NormalTok{total\_trials)}

\FunctionTok{stopifnot}\NormalTok{(}\FunctionTok{nrow}\NormalTok{(agent\_to\_fit) }\SpecialCharTok{==}\NormalTok{ total\_trials)}

\NormalTok{gcm\_data\_single }\OtherTok{\textless{}{-}} \FunctionTok{list}\NormalTok{(}
  \AttributeTok{N\_total           =} \FunctionTok{nrow}\NormalTok{(agent\_to\_fit),}
  \AttributeTok{N\_features        =} \DecValTok{2}\NormalTok{L,}
  \AttributeTok{y                 =} \FunctionTok{as.integer}\NormalTok{(agent\_to\_fit}\SpecialCharTok{$}\NormalTok{sim\_response),}
  \AttributeTok{obs               =} \FunctionTok{as.matrix}\NormalTok{(agent\_to\_fit[, }\FunctionTok{c}\NormalTok{(}\StringTok{"height"}\NormalTok{, }\StringTok{"position"}\NormalTok{)]),}
  \AttributeTok{cat\_feedback      =} \FunctionTok{as.integer}\NormalTok{(agent\_to\_fit}\SpecialCharTok{$}\NormalTok{category\_feedback),}
  \AttributeTok{prior\_w\_alpha     =} \FunctionTok{c}\NormalTok{(}\DecValTok{1}\NormalTok{, }\DecValTok{1}\NormalTok{),}
  \AttributeTok{prior\_log\_c\_mu    =}\NormalTok{ LOG\_C\_PRIOR\_MEAN,}
  \AttributeTok{prior\_log\_c\_sigma =}\NormalTok{ LOG\_C\_PRIOR\_SD,}
  \AttributeTok{prior\_bias\_alpha  =} \DecValTok{1}\NormalTok{,}
  \AttributeTok{prior\_bias\_beta   =} \DecValTok{1}\NormalTok{,}
  \AttributeTok{run\_diagnostics   =} \DecValTok{1}\NormalTok{L}
\NormalTok{)}

\NormalTok{fit\_filepath\_single }\OtherTok{\textless{}{-}} \FunctionTok{here}\NormalTok{(}\StringTok{"simmodels"}\NormalTok{, }\StringTok{"ch12\_gcm\_single\_fit.rds"}\NormalTok{)}

\ControlFlowTok{if}\NormalTok{ (regenerate\_simulations }\SpecialCharTok{||} \SpecialCharTok{!}\FunctionTok{file.exists}\NormalTok{(fit\_filepath\_single)) \{}
\NormalTok{  pf\_single }\OtherTok{\textless{}{-}} \FunctionTok{tryCatch}\NormalTok{(}
\NormalTok{    mod\_gcm\_single}\SpecialCharTok{$}\FunctionTok{pathfinder}\NormalTok{(}
      \AttributeTok{data      =}\NormalTok{ gcm\_data\_single,}
      \AttributeTok{seed      =} \DecValTok{123}\NormalTok{,}
      \AttributeTok{num\_paths =} \DecValTok{4}\NormalTok{,}
      \AttributeTok{refresh   =} \DecValTok{0}
\NormalTok{    ),}
    \AttributeTok{error =} \ControlFlowTok{function}\NormalTok{(e) \{}
      \FunctionTok{message}\NormalTok{(}\StringTok{"Pathfinder failed — falling back to default initialisation."}\NormalTok{)}
      \ConstantTok{NULL}
\NormalTok{    \}}
\NormalTok{  )}

\NormalTok{  fit\_gcm\_single }\OtherTok{\textless{}{-}}\NormalTok{ mod\_gcm\_single}\SpecialCharTok{$}\FunctionTok{sample}\NormalTok{(}
    \AttributeTok{data            =}\NormalTok{ gcm\_data\_single,}
    \AttributeTok{init            =}\NormalTok{ pf\_single,}
    \AttributeTok{seed            =} \DecValTok{123}\NormalTok{,}
    \AttributeTok{chains          =} \DecValTok{4}\NormalTok{,}
    \AttributeTok{parallel\_chains =} \FunctionTok{min}\NormalTok{(}\DecValTok{4}\NormalTok{, }\FunctionTok{availableCores}\NormalTok{()),}
    \AttributeTok{iter\_warmup     =} \DecValTok{1000}\NormalTok{,}
    \AttributeTok{iter\_sampling   =} \DecValTok{1500}\NormalTok{,}
    \AttributeTok{refresh         =} \DecValTok{300}\NormalTok{,}
    \AttributeTok{adapt\_delta     =} \FloatTok{0.9}
\NormalTok{  )}
\NormalTok{  fit\_gcm\_single}\SpecialCharTok{$}\FunctionTok{save\_object}\NormalTok{(fit\_filepath\_single)}
  \FunctionTok{cat}\NormalTok{(}\StringTok{"Single{-}agent model fitted and saved.}\SpecialCharTok{\textbackslash{}n}\StringTok{"}\NormalTok{)}
\NormalTok{\} }\ControlFlowTok{else}\NormalTok{ \{}
\NormalTok{  fit\_gcm\_single }\OtherTok{\textless{}{-}} \FunctionTok{readRDS}\NormalTok{(fit\_filepath\_single)}
  \FunctionTok{cat}\NormalTok{(}\StringTok{"Loaded existing single{-}agent fit.}\SpecialCharTok{\textbackslash{}n}\StringTok{"}\NormalTok{)}
\NormalTok{\}}
\end{Highlighting}
\end{Shaded}

\begin{verbatim}
Loaded existing single-agent fit.
\end{verbatim}

\subsubsection{MCMC Diagnostic Battery}\label{mcmc-diagnostic-battery}

\begin{Shaded}
\begin{Highlighting}[]
\ControlFlowTok{if}\NormalTok{ (}\SpecialCharTok{!}\FunctionTok{is.null}\NormalTok{(fit\_gcm\_single)) \{}
\NormalTok{  diag\_tbl }\OtherTok{\textless{}{-}} \FunctionTok{diagnostic\_summary\_table}\NormalTok{(}
\NormalTok{    fit\_gcm\_single,}
    \AttributeTok{params =} \FunctionTok{c}\NormalTok{(}\StringTok{"w[1]"}\NormalTok{, }\StringTok{"w[2]"}\NormalTok{, }\StringTok{"log\_c"}\NormalTok{, }\StringTok{"c"}\NormalTok{, }\StringTok{"bias"}\NormalTok{)}
\NormalTok{  )}
  \FunctionTok{print}\NormalTok{(diag\_tbl)}

  \ControlFlowTok{if}\NormalTok{ (}\SpecialCharTok{!}\FunctionTok{all}\NormalTok{(diag\_tbl}\SpecialCharTok{$}\NormalTok{pass)) \{}
    \FunctionTok{warning}\NormalTok{(}\StringTok{"Diagnostic battery has failures — interpret posteriors with caution."}\NormalTok{)}
\NormalTok{  \}}

\NormalTok{  bayesplot}\SpecialCharTok{::}\FunctionTok{mcmc\_trace}\NormalTok{(}
\NormalTok{    fit\_gcm\_single}\SpecialCharTok{$}\FunctionTok{draws}\NormalTok{(}\FunctionTok{c}\NormalTok{(}\StringTok{"log\_c"}\NormalTok{, }\StringTok{"w[1]"}\NormalTok{, }\StringTok{"bias"}\NormalTok{)),}
    \AttributeTok{facet\_args =} \FunctionTok{list}\NormalTok{(}\AttributeTok{ncol =} \DecValTok{1}\NormalTok{)}
\NormalTok{  )}

\NormalTok{  bayesplot}\SpecialCharTok{::}\FunctionTok{mcmc\_rank\_overlay}\NormalTok{(}
\NormalTok{    fit\_gcm\_single}\SpecialCharTok{$}\FunctionTok{draws}\NormalTok{(}\FunctionTok{c}\NormalTok{(}\StringTok{"log\_c"}\NormalTok{, }\StringTok{"w[1]"}\NormalTok{, }\StringTok{"bias"}\NormalTok{))}
\NormalTok{  )}
\NormalTok{\}}
\end{Highlighting}
\end{Shaded}

\begin{verbatim}
# A tibble: 6 x 4
  metric                           value threshold pass 
  <chr>                            <dbl> <chr>     <lgl>
1 Divergences (zero tolerance)    0      == 0      TRUE 
2 Max rank-normalised R-hat       1.00   < 1.01    TRUE 
3 Min bulk ESS                 2921.     > 400     TRUE 
4 Min tail ESS                 1989.     > 400     TRUE 
5 Min E-BFMI                      0.873  > 0.2     TRUE 
6 Max MCSE / posterior SD         0.0162 < 0.05    TRUE 
\end{verbatim}

\begin{center}
\includegraphics[width=0.8\linewidth,height=\textheight,keepaspectratio]{11-CategorizationModelsExemplars_files/figure-pdf/ch12_single_diagnostic_battery-1.pdf}
\end{center}

\subsubsection{Posterior Predictive Check (single
subject)}\label{posterior-predictive-check-single-subject}

\begin{Shaded}
\begin{Highlighting}[]
\ControlFlowTok{if}\NormalTok{ (}\SpecialCharTok{!}\FunctionTok{is.null}\NormalTok{(fit\_gcm\_single)) \{}
\NormalTok{  prob\_draws }\OtherTok{\textless{}{-}}\NormalTok{ fit\_gcm\_single}\SpecialCharTok{$}\FunctionTok{draws}\NormalTok{(}\StringTok{"prob\_cat1"}\NormalTok{, }\AttributeTok{format =} \StringTok{"matrix"}\NormalTok{)}
\NormalTok{  n\_draws    }\OtherTok{\textless{}{-}} \FunctionTok{nrow}\NormalTok{(prob\_draws)}
\NormalTok{  n\_trials   }\OtherTok{\textless{}{-}} \FunctionTok{ncol}\NormalTok{(prob\_draws)}

  \FunctionTok{set.seed}\NormalTok{(}\DecValTok{2025}\NormalTok{)}
\NormalTok{  yrep\_post }\OtherTok{\textless{}{-}} \FunctionTok{matrix}\NormalTok{(}\FunctionTok{rbinom}\NormalTok{(n\_draws }\SpecialCharTok{*}\NormalTok{ n\_trials, }\DecValTok{1}\NormalTok{, }\FunctionTok{as.vector}\NormalTok{(prob\_draws)),}
                      \AttributeTok{nrow =}\NormalTok{ n\_draws)}

\NormalTok{  obs\_acc }\OtherTok{\textless{}{-}} \FunctionTok{cummean}\NormalTok{(agent\_to\_fit}\SpecialCharTok{$}\NormalTok{sim\_response }\SpecialCharTok{==}\NormalTok{ agent\_to\_fit}\SpecialCharTok{$}\NormalTok{category\_feedback)}
\NormalTok{  rep\_acc }\OtherTok{\textless{}{-}} \FunctionTok{t}\NormalTok{(}\FunctionTok{apply}\NormalTok{(yrep\_post, }\DecValTok{1}\NormalTok{, }\ControlFlowTok{function}\NormalTok{(yr) \{}
    \FunctionTok{cummean}\NormalTok{(yr }\SpecialCharTok{==}\NormalTok{ agent\_to\_fit}\SpecialCharTok{$}\NormalTok{category\_feedback)}
\NormalTok{  \}))}

\NormalTok{  ppc\_summary }\OtherTok{\textless{}{-}} \FunctionTok{tibble}\NormalTok{(}
    \AttributeTok{trial =} \FunctionTok{seq\_len}\NormalTok{(n\_trials),}
    \AttributeTok{obs   =}\NormalTok{ obs\_acc,}
    \AttributeTok{q05   =} \FunctionTok{apply}\NormalTok{(rep\_acc, }\DecValTok{2}\NormalTok{, quantile, }\FloatTok{0.05}\NormalTok{),}
    \AttributeTok{q25   =} \FunctionTok{apply}\NormalTok{(rep\_acc, }\DecValTok{2}\NormalTok{, quantile, }\FloatTok{0.25}\NormalTok{),}
    \AttributeTok{q50   =} \FunctionTok{apply}\NormalTok{(rep\_acc, }\DecValTok{2}\NormalTok{, quantile, }\FloatTok{0.50}\NormalTok{),}
    \AttributeTok{q75   =} \FunctionTok{apply}\NormalTok{(rep\_acc, }\DecValTok{2}\NormalTok{, quantile, }\FloatTok{0.75}\NormalTok{),}
    \AttributeTok{q95   =} \FunctionTok{apply}\NormalTok{(rep\_acc, }\DecValTok{2}\NormalTok{, quantile, }\FloatTok{0.95}\NormalTok{)}
\NormalTok{  )}

  \FunctionTok{ggplot}\NormalTok{(ppc\_summary, }\FunctionTok{aes}\NormalTok{(}\AttributeTok{x =}\NormalTok{ trial)) }\SpecialCharTok{+}
    \FunctionTok{geom\_ribbon}\NormalTok{(}\FunctionTok{aes}\NormalTok{(}\AttributeTok{ymin =}\NormalTok{ q05, }\AttributeTok{ymax =}\NormalTok{ q95), }\AttributeTok{fill =} \StringTok{"\#56B4E9"}\NormalTok{, }\AttributeTok{alpha =} \FloatTok{0.20}\NormalTok{) }\SpecialCharTok{+}
    \FunctionTok{geom\_ribbon}\NormalTok{(}\FunctionTok{aes}\NormalTok{(}\AttributeTok{ymin =}\NormalTok{ q25, }\AttributeTok{ymax =}\NormalTok{ q75), }\AttributeTok{fill =} \StringTok{"\#56B4E9"}\NormalTok{, }\AttributeTok{alpha =} \FloatTok{0.40}\NormalTok{) }\SpecialCharTok{+}
    \FunctionTok{geom\_line}\NormalTok{(}\FunctionTok{aes}\NormalTok{(}\AttributeTok{y =}\NormalTok{ q50), }\AttributeTok{color =} \StringTok{"\#0072B2"}\NormalTok{, }\AttributeTok{linewidth =} \FloatTok{0.8}\NormalTok{) }\SpecialCharTok{+}
    \FunctionTok{geom\_line}\NormalTok{(}\FunctionTok{aes}\NormalTok{(}\AttributeTok{y =}\NormalTok{ obs), }\AttributeTok{color =} \StringTok{"\#D55E00"}\NormalTok{, }\AttributeTok{linewidth =} \DecValTok{1}\NormalTok{) }\SpecialCharTok{+}
    \FunctionTok{scale\_y\_continuous}\NormalTok{(}\AttributeTok{limits =} \FunctionTok{c}\NormalTok{(}\DecValTok{0}\NormalTok{, }\DecValTok{1}\NormalTok{)) }\SpecialCharTok{+}
    \FunctionTok{labs}\NormalTok{(}
      \AttributeTok{title    =} \StringTok{"Posterior Predictive Check: Cumulative Accuracy"}\NormalTok{,}
      \AttributeTok{subtitle =} \StringTok{"Blue ribbons: 50\% and 90\% posterior predictive intervals}\SpecialCharTok{\textbackslash{}n}\StringTok{Orange line: observed cumulative accuracy"}\NormalTok{,}
      \AttributeTok{x =} \StringTok{"Trial"}\NormalTok{, }\AttributeTok{y =} \StringTok{"Cumulative Accuracy"}
\NormalTok{    )}
\NormalTok{\}}
\end{Highlighting}
\end{Shaded}

\begin{center}
\includegraphics[width=0.8\linewidth,height=\textheight,keepaspectratio]{11-CategorizationModelsExemplars_files/figure-pdf/ch12_ppc_single-1.pdf}
\end{center}

Note that we picked cumulative accuracy as target measure for the PPC
because it is a more holistic summary of the learning dynamics than
trial-level accuracy. The GCM's path-dependence means that trial-level
predictions are noisier and less informative about model fit. We could
have also picked other behaviors, such as the evolving distribution of
attention-weighted distances to exemplars, or specific contrast

\subsubsection{Prior--Posterior Update
Plot}\label{priorposterior-update-plot}

\begin{Shaded}
\begin{Highlighting}[]
\ControlFlowTok{if}\NormalTok{ (}\SpecialCharTok{!}\FunctionTok{is.null}\NormalTok{(fit\_gcm\_single)) \{}
\NormalTok{  draws\_df }\OtherTok{\textless{}{-}} \FunctionTok{as\_draws\_df}\NormalTok{(}
\NormalTok{    fit\_gcm\_single}\SpecialCharTok{$}\FunctionTok{draws}\NormalTok{(}\AttributeTok{variables =} \FunctionTok{c}\NormalTok{(}\StringTok{"w[1]"}\NormalTok{, }\StringTok{"log\_c"}\NormalTok{, }\StringTok{"c"}\NormalTok{, }\StringTok{"bias"}\NormalTok{))}
\NormalTok{  )}

\NormalTok{  true\_w1    }\OtherTok{\textless{}{-}}\NormalTok{ agent\_to\_fit}\SpecialCharTok{$}\NormalTok{w1\_true[}\DecValTok{1}\NormalTok{]}
\NormalTok{  true\_log\_c }\OtherTok{\textless{}{-}}\NormalTok{ agent\_to\_fit}\SpecialCharTok{$}\NormalTok{log\_c\_true[}\DecValTok{1}\NormalTok{]}
\NormalTok{  true\_c     }\OtherTok{\textless{}{-}}\NormalTok{ agent\_to\_fit}\SpecialCharTok{$}\NormalTok{c\_true[}\DecValTok{1}\NormalTok{]}
\NormalTok{  true\_bias  }\OtherTok{\textless{}{-}}\NormalTok{ agent\_to\_fit}\SpecialCharTok{$}\NormalTok{bias\_true[}\DecValTok{1}\NormalTok{]}

  \FunctionTok{set.seed}\NormalTok{(}\DecValTok{2026}\NormalTok{)}
\NormalTok{  n\_prior }\OtherTok{\textless{}{-}} \DecValTok{4000}
\NormalTok{  prior\_df }\OtherTok{\textless{}{-}} \FunctionTok{tibble}\NormalTok{(}
    \StringTok{\textasciigrave{}}\AttributeTok{w[1]}\StringTok{\textasciigrave{}} \OtherTok{=} \FunctionTok{runif}\NormalTok{(n\_prior),}
    \AttributeTok{log\_c  =} \FunctionTok{rnorm}\NormalTok{(n\_prior, LOG\_C\_PRIOR\_MEAN, LOG\_C\_PRIOR\_SD),}
    \AttributeTok{c      =} \FunctionTok{exp}\NormalTok{(log\_c),}
    \AttributeTok{bias   =} \FunctionTok{rbeta}\NormalTok{(n\_prior, }\DecValTok{1}\NormalTok{, }\DecValTok{1}\NormalTok{)}
\NormalTok{  )}

\NormalTok{  pp\_panel }\OtherTok{\textless{}{-}} \ControlFlowTok{function}\NormalTok{(prior\_x, post\_x, true\_val, label) \{}
    \FunctionTok{tibble}\NormalTok{(}
      \AttributeTok{x      =} \FunctionTok{c}\NormalTok{(prior\_x, post\_x),}
      \AttributeTok{source =} \FunctionTok{c}\NormalTok{(}\FunctionTok{rep}\NormalTok{(}\StringTok{"Prior"}\NormalTok{, }\FunctionTok{length}\NormalTok{(prior\_x)),}
                 \FunctionTok{rep}\NormalTok{(}\StringTok{"Posterior"}\NormalTok{, }\FunctionTok{length}\NormalTok{(post\_x)))}
\NormalTok{    ) }\SpecialCharTok{|\textgreater{}}
      \FunctionTok{ggplot}\NormalTok{(}\FunctionTok{aes}\NormalTok{(}\AttributeTok{x =}\NormalTok{ x, }\AttributeTok{fill =}\NormalTok{ source, }\AttributeTok{color =}\NormalTok{ source)) }\SpecialCharTok{+}
      \FunctionTok{geom\_density}\NormalTok{(}\AttributeTok{alpha =} \FloatTok{0.4}\NormalTok{, }\AttributeTok{linewidth =} \FloatTok{0.5}\NormalTok{) }\SpecialCharTok{+}
      \FunctionTok{geom\_vline}\NormalTok{(}\AttributeTok{xintercept =}\NormalTok{ true\_val, }\AttributeTok{color =} \StringTok{"red"}\NormalTok{,}
                 \AttributeTok{linetype =} \StringTok{"dashed"}\NormalTok{, }\AttributeTok{linewidth =} \FloatTok{0.7}\NormalTok{) }\SpecialCharTok{+}
      \FunctionTok{scale\_fill\_manual}\NormalTok{(}\AttributeTok{values =} \FunctionTok{c}\NormalTok{(}\AttributeTok{Prior =} \StringTok{"grey60"}\NormalTok{, }\AttributeTok{Posterior =} \StringTok{"\#0072B2"}\NormalTok{)) }\SpecialCharTok{+}
      \FunctionTok{scale\_color\_manual}\NormalTok{(}\AttributeTok{values =} \FunctionTok{c}\NormalTok{(}\AttributeTok{Prior =} \StringTok{"grey40"}\NormalTok{, }\AttributeTok{Posterior =} \StringTok{"\#0072B2"}\NormalTok{)) }\SpecialCharTok{+}
      \FunctionTok{labs}\NormalTok{(}\AttributeTok{title =}\NormalTok{ label, }\AttributeTok{x =} \ConstantTok{NULL}\NormalTok{, }\AttributeTok{y =} \StringTok{"Density"}\NormalTok{) }\SpecialCharTok{+}
      \FunctionTok{theme}\NormalTok{(}\AttributeTok{legend.position =} \StringTok{"bottom"}\NormalTok{)}
\NormalTok{  \}}

\NormalTok{  p1 }\OtherTok{\textless{}{-}} \FunctionTok{pp\_panel}\NormalTok{(prior\_df}\SpecialCharTok{$}\StringTok{\textasciigrave{}}\AttributeTok{w[1]}\StringTok{\textasciigrave{}}\NormalTok{, draws\_df}\SpecialCharTok{$}\StringTok{\textasciigrave{}}\AttributeTok{w[1]}\StringTok{\textasciigrave{}}\NormalTok{, true\_w1, }\StringTok{"w[1]"}\NormalTok{)}
\NormalTok{  p2 }\OtherTok{\textless{}{-}} \FunctionTok{pp\_panel}\NormalTok{(prior\_df}\SpecialCharTok{$}\NormalTok{log\_c,  draws\_df}\SpecialCharTok{$}\NormalTok{log\_c,  true\_log\_c, }\StringTok{"log(c) — computational scale"}\NormalTok{)}
\NormalTok{  p3 }\OtherTok{\textless{}{-}} \FunctionTok{pp\_panel}\NormalTok{(prior\_df}\SpecialCharTok{$}\NormalTok{c,      draws\_df}\SpecialCharTok{$}\NormalTok{c,      true\_c,     }\StringTok{"c — natural scale"}\NormalTok{)}
\NormalTok{  p4 }\OtherTok{\textless{}{-}} \FunctionTok{pp\_panel}\NormalTok{(prior\_df}\SpecialCharTok{$}\NormalTok{bias,   draws\_df}\SpecialCharTok{$}\NormalTok{bias,   true\_bias,  }\StringTok{"bias"}\NormalTok{)}

\NormalTok{  (p1 }\SpecialCharTok{|}\NormalTok{ p2) }\SpecialCharTok{/}\NormalTok{ (p3 }\SpecialCharTok{|}\NormalTok{ p4) }\SpecialCharTok{+}
    \FunctionTok{plot\_annotation}\NormalTok{(}
      \AttributeTok{title    =} \StringTok{"Prior–Posterior Update for the Single{-}Subject GCM"}\NormalTok{,}
      \AttributeTok{subtitle =} \StringTok{"Red dashed line = true generating value"}
\NormalTok{    )}
\NormalTok{\}}
\end{Highlighting}
\end{Shaded}

\begin{center}
\includegraphics[width=0.8\linewidth,height=\textheight,keepaspectratio]{11-CategorizationModelsExemplars_files/figure-pdf/ch12_prior_posterior_single-1.pdf}
\end{center}

This looks surprisingly good! --- the posterior for \(w[1]\) is tightly
concentrated around the true value of 0.9, and the posterior for
\(\log c\) is also well-concentrated around the true \(\log(2)\). The
bias parameter is less informed, which makes sense given that it has a
more indirect influence on the choice probabilities.

\subsubsection{\texorpdfstring{Joint Posterior Geometry: \(c\) vs
\(w\)}{Joint Posterior Geometry: c vs w}}\label{joint-posterior-geometry-c-vs-w}

\begin{Shaded}
\begin{Highlighting}[]
\ControlFlowTok{if}\NormalTok{ (}\SpecialCharTok{!}\FunctionTok{is.null}\NormalTok{(fit\_gcm\_single)) \{}
\NormalTok{  bayesplot}\SpecialCharTok{::}\FunctionTok{mcmc\_pairs}\NormalTok{(}
\NormalTok{    fit\_gcm\_single}\SpecialCharTok{$}\FunctionTok{draws}\NormalTok{(}\FunctionTok{c}\NormalTok{(}\StringTok{"log\_c"}\NormalTok{, }\StringTok{"w[1]"}\NormalTok{, }\StringTok{"bias"}\NormalTok{)),}
    \AttributeTok{diag\_fun     =} \StringTok{"dens"}\NormalTok{,}
    \AttributeTok{off\_diag\_fun =} \StringTok{"hex"}
\NormalTok{  )}
\NormalTok{\}}
\end{Highlighting}
\end{Shaded}

\begin{center}
\includegraphics[width=0.8\linewidth,height=\textheight,keepaspectratio]{11-CategorizationModelsExemplars_files/figure-pdf/ch12_cw_identifiability_pairs-1.pdf}
\end{center}

The parameters look nicely uncorrelated.

\begin{center}\rule{0.5\linewidth}{0.5pt}\end{center}

\subsubsection{\texorpdfstring{LOO with Pareto-\(\hat{k}\)
Diagnostics}{LOO with Pareto-\textbackslash hat\{k\} Diagnostics}}\label{loo-with-pareto-hatk-diagnostics}

Standard PSIS-LOO is computed for completeness and for comparison
against LFO-CV below. We do not interpret it as a valid measure of
out-of-sample predictive accuracy for a sequential model --- see the
next subsection for why --- but the Pareto-\(\hat{k}\) diagnostic still
localises trials where the path-dependence concern bites hardest.

\begin{Shaded}
\begin{Highlighting}[]
\ControlFlowTok{if}\NormalTok{ (}\SpecialCharTok{!}\FunctionTok{is.null}\NormalTok{(fit\_gcm\_single)) \{}
\NormalTok{  log\_lik\_mat    }\OtherTok{\textless{}{-}}\NormalTok{ fit\_gcm\_single}\SpecialCharTok{$}\FunctionTok{draws}\NormalTok{(}\StringTok{"log\_lik"}\NormalTok{, }\AttributeTok{format =} \StringTok{"matrix"}\NormalTok{)}
\NormalTok{  loo\_gcm\_single }\OtherTok{\textless{}{-}}\NormalTok{ loo}\SpecialCharTok{::}\FunctionTok{loo}\NormalTok{(log\_lik\_mat, }\AttributeTok{cores =} \DecValTok{4}\NormalTok{)}
  \FunctionTok{print}\NormalTok{(loo\_gcm\_single)}

  \FunctionTok{plot}\NormalTok{(loo\_gcm\_single, }\AttributeTok{label\_points =} \ConstantTok{TRUE}\NormalTok{,}
       \AttributeTok{main =} \StringTok{"PSIS{-}LOO Pareto{-}k diagnostics: Single{-}subject GCM"}\NormalTok{)}
\NormalTok{\}}
\end{Highlighting}
\end{Shaded}

\begin{verbatim}

Computed from 6000 by 64 log-likelihood matrix.

         Estimate  SE
elpd_loo    -29.2 3.6
p_loo         2.4 0.5
looic        58.4 7.1
------
MCSE of elpd_loo is 0.0.
MCSE and ESS estimates assume independent draws (r_eff=1).

All Pareto k estimates are good (k < 0.7).
See help('pareto-k-diagnostic') for details.
\end{verbatim}

\begin{center}
\includegraphics[width=0.8\linewidth,height=\textheight,keepaspectratio]{11-CategorizationModelsExemplars_files/figure-pdf/ch12_loo_single-1.pdf}
\end{center}

\subsubsection{Leave-Future-Out
Cross-Validation}\label{leave-future-out-cross-validation}

\textbf{Why standard LOO is invalid here.} Standard PSIS-LOO assumes
observations are conditionally exchangeable given the parameters. In the
GCM, this assumption is violated: leaving out observation \(t\) also
conceptually removes the memory trace that observation \(t\) contributed
to trials \(t+1, \ldots, T\).

\textbf{What LFO-CV does instead.} LFO-CV (Bürkner, Gabry, \& Vehtari,
2020) estimates the \emph{one-step-ahead} predictive density. The PSIS
approximation refits only when the importance sampling deteriorates
(\(\hat{k} > k_\text{thresh}\)) and uses importance reweighting
otherwise.

\begin{Shaded}
\begin{Highlighting}[]
\CommentTok{\# ─────────────────────────────────────────────────────────────────────────────}
\CommentTok{\# psis\_lfo\_gcm: PSIS{-}approximated leave{-}future{-}out CV for the GCM.}
\CommentTok{\# Follows Bürkner, Gabry, \& Vehtari (2020), Algorithm 1.}
\CommentTok{\# ─────────────────────────────────────────────────────────────────────────────}
\NormalTok{psis\_lfo\_gcm }\OtherTok{\textless{}{-}} \ControlFlowTok{function}\NormalTok{(stan\_model,}
\NormalTok{                         full\_data,}
                         \AttributeTok{L =} \DecValTok{16}\NormalTok{,}
                         \AttributeTok{k\_thresh =} \FloatTok{0.7}\NormalTok{,}
                         \AttributeTok{seed =} \DecValTok{1}\NormalTok{,}
                         \AttributeTok{verbose =} \ConstantTok{FALSE}\NormalTok{) \{}

\NormalTok{  fit\_subset }\OtherTok{\textless{}{-}} \ControlFlowTok{function}\NormalTok{(n\_train, seed) \{}
\NormalTok{    sub\_data }\OtherTok{\textless{}{-}}\NormalTok{ full\_data}
\NormalTok{    sub\_data}\SpecialCharTok{$}\NormalTok{N\_total      }\OtherTok{\textless{}{-}}\NormalTok{ n\_train}
\NormalTok{    sub\_data}\SpecialCharTok{$}\NormalTok{y            }\OtherTok{\textless{}{-}}\NormalTok{ sub\_data}\SpecialCharTok{$}\NormalTok{y[}\DecValTok{1}\SpecialCharTok{:}\NormalTok{n\_train]}
\NormalTok{    sub\_data}\SpecialCharTok{$}\NormalTok{obs          }\OtherTok{\textless{}{-}}\NormalTok{ sub\_data}\SpecialCharTok{$}\NormalTok{obs[}\DecValTok{1}\SpecialCharTok{:}\NormalTok{n\_train, , drop }\OtherTok{=} \ConstantTok{FALSE}\NormalTok{]}
\NormalTok{    sub\_data}\SpecialCharTok{$}\NormalTok{cat\_feedback }\OtherTok{\textless{}{-}}\NormalTok{ sub\_data}\SpecialCharTok{$}\NormalTok{cat\_feedback[}\DecValTok{1}\SpecialCharTok{:}\NormalTok{n\_train]}

\NormalTok{    pf  }\OtherTok{\textless{}{-}}\NormalTok{ stan\_model}\SpecialCharTok{$}\FunctionTok{pathfinder}\NormalTok{(}\AttributeTok{data =}\NormalTok{ sub\_data, }\AttributeTok{num\_paths =} \DecValTok{4}\NormalTok{, }\AttributeTok{refresh =} \DecValTok{0}\NormalTok{,}
                                  \AttributeTok{seed =}\NormalTok{ seed)}
\NormalTok{    fit }\OtherTok{\textless{}{-}}\NormalTok{ stan\_model}\SpecialCharTok{$}\FunctionTok{sample}\NormalTok{(}
      \AttributeTok{data =}\NormalTok{ sub\_data, }\AttributeTok{init =}\NormalTok{ pf, }\AttributeTok{seed =}\NormalTok{ seed,}
      \AttributeTok{chains =} \DecValTok{2}\NormalTok{, }\AttributeTok{parallel\_chains =} \DecValTok{2}\NormalTok{,}
      \AttributeTok{iter\_warmup =} \DecValTok{600}\NormalTok{, }\AttributeTok{iter\_sampling =} \DecValTok{800}\NormalTok{,}
      \AttributeTok{refresh =} \DecValTok{0}\NormalTok{, }\AttributeTok{adapt\_delta =} \FloatTok{0.9}
\NormalTok{    )}
\NormalTok{    fit}
\NormalTok{  \}}

\NormalTok{  predict\_future }\OtherTok{\textless{}{-}} \ControlFlowTok{function}\NormalTok{(fit, full\_data) \{}
\NormalTok{    full\_n   }\OtherTok{\textless{}{-}} \FunctionTok{length}\NormalTok{(full\_data}\SpecialCharTok{$}\NormalTok{y)}
\NormalTok{    ext\_data }\OtherTok{\textless{}{-}}\NormalTok{ full\_data}
\NormalTok{    ext\_data}\SpecialCharTok{$}\NormalTok{N\_total }\OtherTok{\textless{}{-}}\NormalTok{ full\_n}

\NormalTok{    gq\_fit }\OtherTok{\textless{}{-}}\NormalTok{ stan\_model}\SpecialCharTok{$}\FunctionTok{generate\_quantities}\NormalTok{(}
      \AttributeTok{fitted\_params   =}\NormalTok{ fit,}
      \AttributeTok{data            =}\NormalTok{ ext\_data,}
      \AttributeTok{seed            =}\NormalTok{ seed,}
      \AttributeTok{parallel\_chains =} \DecValTok{2}
\NormalTok{    )}
\NormalTok{    gq\_fit}\SpecialCharTok{$}\FunctionTok{draws}\NormalTok{(}\StringTok{"log\_lik"}\NormalTok{, }\AttributeTok{format =} \StringTok{"matrix"}\NormalTok{)}
\NormalTok{  \}}

\NormalTok{  T\_total      }\OtherTok{\textless{}{-}} \FunctionTok{length}\NormalTok{(full\_data}\SpecialCharTok{$}\NormalTok{y)}
\NormalTok{  elpd\_lfo     }\OtherTok{\textless{}{-}} \FunctionTok{numeric}\NormalTok{(T\_total)}
\NormalTok{  elpd\_lfo[]   }\OtherTok{\textless{}{-}} \ConstantTok{NA\_real\_}
\NormalTok{  refits       }\OtherTok{\textless{}{-}} \FunctionTok{integer}\NormalTok{(}\DecValTok{0}\NormalTok{)}
\NormalTok{  k\_history    }\OtherTok{\textless{}{-}} \FunctionTok{numeric}\NormalTok{(}\DecValTok{0}\NormalTok{)}

\NormalTok{  fit\_curr     }\OtherTok{\textless{}{-}} \FunctionTok{fit\_subset}\NormalTok{(L, }\AttributeTok{seed =}\NormalTok{ seed)}
\NormalTok{  ll\_full\_curr }\OtherTok{\textless{}{-}} \FunctionTok{predict\_future}\NormalTok{(fit\_curr, full\_data)}

\NormalTok{  refits       }\OtherTok{\textless{}{-}} \FunctionTok{c}\NormalTok{(refits, L)}
\NormalTok{  last\_fit\_n   }\OtherTok{\textless{}{-}}\NormalTok{ L}
\NormalTok{  n\_draws      }\OtherTok{\textless{}{-}} \FunctionTok{nrow}\NormalTok{(ll\_full\_curr)}

\NormalTok{  ll\_train\_curr }\OtherTok{\textless{}{-}} \FunctionTok{rowSums}\NormalTok{(ll\_full\_curr[, }\DecValTok{1}\SpecialCharTok{:}\NormalTok{L, }\AttributeTok{drop =} \ConstantTok{FALSE}\NormalTok{])}

  \ControlFlowTok{for}\NormalTok{ (t }\ControlFlowTok{in}\NormalTok{ (L }\SpecialCharTok{+} \DecValTok{1}\NormalTok{)}\SpecialCharTok{:}\NormalTok{T\_total) \{}
    \ControlFlowTok{if}\NormalTok{ (t }\SpecialCharTok{{-}} \DecValTok{1} \SpecialCharTok{==}\NormalTok{ last\_fit\_n) \{}
\NormalTok{      k\_t          }\OtherTok{\textless{}{-}} \DecValTok{0}
\NormalTok{      k\_history[t] }\OtherTok{\textless{}{-}}\NormalTok{ k\_t}
\NormalTok{      lw           }\OtherTok{\textless{}{-}} \FunctionTok{rep}\NormalTok{(}\SpecialCharTok{{-}}\FunctionTok{log}\NormalTok{(n\_draws), n\_draws)}
\NormalTok{    \} }\ControlFlowTok{else}\NormalTok{ \{}
\NormalTok{      ll\_train\_target }\OtherTok{\textless{}{-}} \FunctionTok{rowSums}\NormalTok{(ll\_full\_curr[, }\DecValTok{1}\SpecialCharTok{:}\NormalTok{(t }\SpecialCharTok{{-}} \DecValTok{1}\NormalTok{), }\AttributeTok{drop =} \ConstantTok{FALSE}\NormalTok{])}
\NormalTok{      log\_w           }\OtherTok{\textless{}{-}}\NormalTok{ ll\_train\_target }\SpecialCharTok{{-}}\NormalTok{ ll\_train\_curr}

\NormalTok{      psis\_obj     }\OtherTok{\textless{}{-}}\NormalTok{ loo}\SpecialCharTok{::}\FunctionTok{psis}\NormalTok{(log\_w, }\AttributeTok{r\_eff =} \ConstantTok{NA}\NormalTok{)}
\NormalTok{      k\_t          }\OtherTok{\textless{}{-}}\NormalTok{ loo}\SpecialCharTok{::}\FunctionTok{pareto\_k\_values}\NormalTok{(psis\_obj)}
\NormalTok{      k\_history[t] }\OtherTok{\textless{}{-}}\NormalTok{ k\_t}

      \ControlFlowTok{if}\NormalTok{ (k\_t }\SpecialCharTok{\textgreater{}}\NormalTok{ k\_thresh) \{}
        \ControlFlowTok{if}\NormalTok{ (verbose) }\FunctionTok{cat}\NormalTok{(}\StringTok{"Refit at t ="}\NormalTok{, t, }\StringTok{"(k ="}\NormalTok{, }\FunctionTok{round}\NormalTok{(k\_t, }\DecValTok{3}\NormalTok{), }\StringTok{")}\SpecialCharTok{\textbackslash{}n}\StringTok{"}\NormalTok{)}
\NormalTok{        fit\_curr      }\OtherTok{\textless{}{-}} \FunctionTok{fit\_subset}\NormalTok{(t }\SpecialCharTok{{-}} \DecValTok{1}\NormalTok{, }\AttributeTok{seed =}\NormalTok{ seed }\SpecialCharTok{+}\NormalTok{ t)}
\NormalTok{        ll\_full\_curr  }\OtherTok{\textless{}{-}} \FunctionTok{predict\_future}\NormalTok{(fit\_curr, full\_data)}
\NormalTok{        ll\_train\_curr }\OtherTok{\textless{}{-}} \FunctionTok{rowSums}\NormalTok{(ll\_full\_curr[, }\DecValTok{1}\SpecialCharTok{:}\NormalTok{(t }\SpecialCharTok{{-}} \DecValTok{1}\NormalTok{), }\AttributeTok{drop =} \ConstantTok{FALSE}\NormalTok{])}
\NormalTok{        last\_fit\_n    }\OtherTok{\textless{}{-}}\NormalTok{ t }\SpecialCharTok{{-}} \DecValTok{1}
\NormalTok{        refits        }\OtherTok{\textless{}{-}} \FunctionTok{c}\NormalTok{(refits, last\_fit\_n)}
\NormalTok{        lw            }\OtherTok{\textless{}{-}} \FunctionTok{rep}\NormalTok{(}\SpecialCharTok{{-}}\FunctionTok{log}\NormalTok{(n\_draws), n\_draws)}
\NormalTok{      \} }\ControlFlowTok{else}\NormalTok{ \{}
\NormalTok{        lw }\OtherTok{\textless{}{-}} \FunctionTok{as.numeric}\NormalTok{(}\FunctionTok{weights}\NormalTok{(psis\_obj))}
\NormalTok{      \}}
\NormalTok{    \}}

\NormalTok{    lp          }\OtherTok{\textless{}{-}}\NormalTok{ ll\_full\_curr[, t]}
\NormalTok{    elpd\_lfo[t] }\OtherTok{\textless{}{-}}\NormalTok{ matrixStats}\SpecialCharTok{::}\FunctionTok{logSumExp}\NormalTok{(lw }\SpecialCharTok{+}\NormalTok{ lp)}
\NormalTok{  \}}

  \FunctionTok{list}\NormalTok{(}
    \AttributeTok{elpd\_lfo       =} \FunctionTok{sum}\NormalTok{(elpd\_lfo, }\AttributeTok{na.rm =} \ConstantTok{TRUE}\NormalTok{),}
    \AttributeTok{elpd\_pointwise =}\NormalTok{ elpd\_lfo,}
    \AttributeTok{refits         =}\NormalTok{ refits,}
    \AttributeTok{k\_history      =}\NormalTok{ k\_history,}
    \AttributeTok{L              =}\NormalTok{ L}
\NormalTok{  )}
\NormalTok{\}}

\NormalTok{lfo\_filepath }\OtherTok{\textless{}{-}} \FunctionTok{here}\NormalTok{(}\StringTok{"simdata"}\NormalTok{, }\StringTok{"ch12\_lfo\_single.rds"}\NormalTok{)}
\ControlFlowTok{if}\NormalTok{ (regenerate\_simulations }\SpecialCharTok{||} \SpecialCharTok{!}\FunctionTok{file.exists}\NormalTok{(lfo\_filepath)) \{}
\NormalTok{  lfo\_single }\OtherTok{\textless{}{-}} \FunctionTok{psis\_lfo\_gcm}\NormalTok{(}
    \AttributeTok{stan\_model =}\NormalTok{ mod\_gcm\_single,}
    \AttributeTok{full\_data  =}\NormalTok{ gcm\_data\_single,}
    \AttributeTok{L          =} \DecValTok{16}\NormalTok{,}
    \AttributeTok{k\_thresh   =} \FloatTok{0.7}\NormalTok{,}
    \AttributeTok{seed       =} \DecValTok{11}\NormalTok{,}
    \AttributeTok{verbose    =} \ConstantTok{TRUE}
\NormalTok{  )}
  \FunctionTok{saveRDS}\NormalTok{(lfo\_single, lfo\_filepath)}
\NormalTok{\} }\ControlFlowTok{else}\NormalTok{ \{}
\NormalTok{  lfo\_single }\OtherTok{\textless{}{-}} \FunctionTok{readRDS}\NormalTok{(lfo\_filepath)}
\NormalTok{\}}

\FunctionTok{cat}\NormalTok{(}\StringTok{"LFO{-}CV ELPD (single subject):"}\NormalTok{, }\FunctionTok{round}\NormalTok{(lfo\_single}\SpecialCharTok{$}\NormalTok{elpd\_lfo, }\DecValTok{2}\NormalTok{), }\StringTok{"}\SpecialCharTok{\textbackslash{}n}\StringTok{"}\NormalTok{)}
\end{Highlighting}
\end{Shaded}

\begin{verbatim}
LFO-CV ELPD (single subject): -20.63 
\end{verbatim}

\begin{Shaded}
\begin{Highlighting}[]
\FunctionTok{cat}\NormalTok{(}\StringTok{"Number of refits:"}\NormalTok{,             }\FunctionTok{length}\NormalTok{(lfo\_single}\SpecialCharTok{$}\NormalTok{refits),     }\StringTok{"}\SpecialCharTok{\textbackslash{}n}\StringTok{"}\NormalTok{)}
\end{Highlighting}
\end{Shaded}

\begin{verbatim}
Number of refits: 1 
\end{verbatim}

\begin{Shaded}
\begin{Highlighting}[]
\FunctionTok{cat}\NormalTok{(}\StringTok{"Refit triggered at trials:"}\NormalTok{,    lfo\_single}\SpecialCharTok{$}\NormalTok{refits,             }\StringTok{"}\SpecialCharTok{\textbackslash{}n}\StringTok{"}\NormalTok{)}
\end{Highlighting}
\end{Shaded}

\begin{verbatim}
Refit triggered at trials: 16 
\end{verbatim}

\begin{Shaded}
\begin{Highlighting}[]
\ControlFlowTok{if}\NormalTok{ (}\SpecialCharTok{!}\FunctionTok{is.null}\NormalTok{(fit\_gcm\_single)) \{}
\NormalTok{  loo\_pw  }\OtherTok{\textless{}{-}}\NormalTok{ loo\_gcm\_single}\SpecialCharTok{$}\NormalTok{pointwise[, }\StringTok{"elpd\_loo"}\NormalTok{]}
\NormalTok{  comp\_df }\OtherTok{\textless{}{-}} \FunctionTok{tibble}\NormalTok{(}
    \AttributeTok{trial    =} \FunctionTok{seq\_along}\NormalTok{(loo\_pw),}
    \AttributeTok{elpd\_loo =}\NormalTok{ loo\_pw,}
    \AttributeTok{elpd\_lfo =}\NormalTok{ lfo\_single}\SpecialCharTok{$}\NormalTok{elpd\_pointwise}
\NormalTok{  ) }\SpecialCharTok{|\textgreater{}}
    \FunctionTok{pivot\_longer}\NormalTok{(}\FunctionTok{starts\_with}\NormalTok{(}\StringTok{"elpd\_"}\NormalTok{), }\AttributeTok{names\_to =} \StringTok{"method"}\NormalTok{, }\AttributeTok{values\_to =} \StringTok{"elpd"}\NormalTok{)}

  \FunctionTok{ggplot}\NormalTok{(comp\_df, }\FunctionTok{aes}\NormalTok{(}\AttributeTok{x =}\NormalTok{ trial, }\AttributeTok{y =}\NormalTok{ elpd, }\AttributeTok{color =}\NormalTok{ method)) }\SpecialCharTok{+}
    \FunctionTok{geom\_line}\NormalTok{(}\AttributeTok{alpha =} \FloatTok{0.7}\NormalTok{) }\SpecialCharTok{+}
    \FunctionTok{geom\_point}\NormalTok{(}\AttributeTok{alpha =} \FloatTok{0.5}\NormalTok{) }\SpecialCharTok{+}
    \FunctionTok{scale\_color\_manual}\NormalTok{(}\AttributeTok{values =} \FunctionTok{c}\NormalTok{(}\AttributeTok{elpd\_loo =} \StringTok{"\#0072B2"}\NormalTok{, }\AttributeTok{elpd\_lfo =} \StringTok{"\#D55E00"}\NormalTok{),}
                       \AttributeTok{labels =} \FunctionTok{c}\NormalTok{(}\StringTok{"PSIS{-}LFO (correct)"}\NormalTok{,}
                                  \StringTok{"PSIS{-}LOO (invalid for sequential)"}\NormalTok{)) }\SpecialCharTok{+}
    \FunctionTok{labs}\NormalTok{(}
      \AttributeTok{title    =} \StringTok{"Per{-}trial ELPD: LOO vs LFO{-}CV"}\NormalTok{,}
      \AttributeTok{subtitle =} \StringTok{"Differences are largest where the path{-}dependence assumption fails most"}\NormalTok{,}
      \AttributeTok{x =} \StringTok{"Trial"}\NormalTok{, }\AttributeTok{y =} \StringTok{"Pointwise ELPD"}
\NormalTok{    ) }\SpecialCharTok{+}
    \FunctionTok{theme}\NormalTok{(}\AttributeTok{legend.position =} \StringTok{"bottom"}\NormalTok{)}
\NormalTok{\}}
\end{Highlighting}
\end{Shaded}

\begin{center}
\includegraphics[width=0.8\linewidth,height=\textheight,keepaspectratio]{11-CategorizationModelsExemplars_files/figure-pdf/ch12_lfo_cv_implementation-1.pdf}
\end{center}

\textbf{Validation against exact one-step-ahead refits.}

\begin{Shaded}
\begin{Highlighting}[]
\NormalTok{exact\_lfo\_grid     }\OtherTok{\textless{}{-}} \FunctionTok{seq}\NormalTok{(}\DecValTok{16}\NormalTok{, }\FunctionTok{length}\NormalTok{(gcm\_data\_single}\SpecialCharTok{$}\NormalTok{y) }\SpecialCharTok{{-}} \DecValTok{1}\NormalTok{, }\AttributeTok{by =} \DecValTok{8}\NormalTok{)}
\NormalTok{exact\_lfo\_filepath }\OtherTok{\textless{}{-}} \FunctionTok{here}\NormalTok{(}\StringTok{"simdata"}\NormalTok{, }\StringTok{"ch12\_exact\_lfo\_single.rds"}\NormalTok{)}

\ControlFlowTok{if}\NormalTok{ (regenerate\_simulations }\SpecialCharTok{||} \SpecialCharTok{!}\FunctionTok{file.exists}\NormalTok{(exact\_lfo\_filepath)) \{}
\NormalTok{  exact\_results }\OtherTok{\textless{}{-}} \FunctionTok{map\_dfr}\NormalTok{(exact\_lfo\_grid, }\ControlFlowTok{function}\NormalTok{(n\_train) \{}
\NormalTok{    sub\_data }\OtherTok{\textless{}{-}}\NormalTok{ gcm\_data\_single}
\NormalTok{    sub\_data}\SpecialCharTok{$}\NormalTok{N\_total      }\OtherTok{\textless{}{-}}\NormalTok{ n\_train}
\NormalTok{    sub\_data}\SpecialCharTok{$}\NormalTok{y            }\OtherTok{\textless{}{-}}\NormalTok{ sub\_data}\SpecialCharTok{$}\NormalTok{y[}\DecValTok{1}\SpecialCharTok{:}\NormalTok{n\_train]}
\NormalTok{    sub\_data}\SpecialCharTok{$}\NormalTok{obs          }\OtherTok{\textless{}{-}}\NormalTok{ sub\_data}\SpecialCharTok{$}\NormalTok{obs[}\DecValTok{1}\SpecialCharTok{:}\NormalTok{n\_train, , drop }\OtherTok{=} \ConstantTok{FALSE}\NormalTok{]}
\NormalTok{    sub\_data}\SpecialCharTok{$}\NormalTok{cat\_feedback }\OtherTok{\textless{}{-}}\NormalTok{ sub\_data}\SpecialCharTok{$}\NormalTok{cat\_feedback[}\DecValTok{1}\SpecialCharTok{:}\NormalTok{n\_train]}

\NormalTok{    pf  }\OtherTok{\textless{}{-}}\NormalTok{ mod\_gcm\_single}\SpecialCharTok{$}\FunctionTok{pathfinder}\NormalTok{(}\AttributeTok{data =}\NormalTok{ sub\_data, }\AttributeTok{num\_paths =} \DecValTok{4}\NormalTok{,}
                                      \AttributeTok{refresh =} \DecValTok{0}\NormalTok{, }\AttributeTok{seed =}\NormalTok{ n\_train)}
\NormalTok{    fit }\OtherTok{\textless{}{-}}\NormalTok{ mod\_gcm\_single}\SpecialCharTok{$}\FunctionTok{sample}\NormalTok{(}
      \AttributeTok{data =}\NormalTok{ sub\_data, }\AttributeTok{init =}\NormalTok{ pf, }\AttributeTok{seed =}\NormalTok{ n\_train,}
      \AttributeTok{chains =} \DecValTok{2}\NormalTok{, }\AttributeTok{parallel\_chains =} \DecValTok{2}\NormalTok{,}
      \AttributeTok{iter\_warmup =} \DecValTok{600}\NormalTok{, }\AttributeTok{iter\_sampling =} \DecValTok{800}\NormalTok{,}
      \AttributeTok{refresh =} \DecValTok{0}\NormalTok{, }\AttributeTok{adapt\_delta =} \FloatTok{0.9}
\NormalTok{    )}

\NormalTok{    ext }\OtherTok{\textless{}{-}}\NormalTok{ gcm\_data\_single}
\NormalTok{    gq  }\OtherTok{\textless{}{-}}\NormalTok{ mod\_gcm\_single}\SpecialCharTok{$}\FunctionTok{generate\_quantities}\NormalTok{(}
      \AttributeTok{fitted\_params =}\NormalTok{ fit, }\AttributeTok{data =}\NormalTok{ ext, }\AttributeTok{seed =}\NormalTok{ n\_train, }\AttributeTok{parallel\_chains =} \DecValTok{2}
\NormalTok{    )}
\NormalTok{    ll\_target  }\OtherTok{\textless{}{-}}\NormalTok{ gq}\SpecialCharTok{$}\FunctionTok{draws}\NormalTok{(}\StringTok{"log\_lik"}\NormalTok{, }\AttributeTok{format =} \StringTok{"matrix"}\NormalTok{)[, n\_train }\SpecialCharTok{+} \DecValTok{1}\NormalTok{]}
\NormalTok{    elpd\_exact }\OtherTok{\textless{}{-}}\NormalTok{ matrixStats}\SpecialCharTok{::}\FunctionTok{logSumExp}\NormalTok{(ll\_target) }\SpecialCharTok{{-}} \FunctionTok{log}\NormalTok{(}\FunctionTok{length}\NormalTok{(ll\_target))}

    \FunctionTok{tibble}\NormalTok{(}\AttributeTok{trial =}\NormalTok{ n\_train }\SpecialCharTok{+} \DecValTok{1}\NormalTok{, }\AttributeTok{elpd\_exact =}\NormalTok{ elpd\_exact)}
\NormalTok{  \})}
  \FunctionTok{saveRDS}\NormalTok{(exact\_results, exact\_lfo\_filepath)}
\NormalTok{\} }\ControlFlowTok{else}\NormalTok{ \{}
\NormalTok{  exact\_results }\OtherTok{\textless{}{-}} \FunctionTok{readRDS}\NormalTok{(exact\_lfo\_filepath)}
\NormalTok{\}}

\NormalTok{psis\_at\_grid }\OtherTok{\textless{}{-}} \FunctionTok{tibble}\NormalTok{(}
  \AttributeTok{trial     =}\NormalTok{ exact\_results}\SpecialCharTok{$}\NormalTok{trial,}
  \AttributeTok{elpd\_psis =}\NormalTok{ lfo\_single}\SpecialCharTok{$}\NormalTok{elpd\_pointwise[exact\_results}\SpecialCharTok{$}\NormalTok{trial]}
\NormalTok{)}

\NormalTok{agreement\_df }\OtherTok{\textless{}{-}}\NormalTok{ exact\_results }\SpecialCharTok{|\textgreater{}}
  \FunctionTok{left\_join}\NormalTok{(psis\_at\_grid, }\AttributeTok{by =} \StringTok{"trial"}\NormalTok{)}

\FunctionTok{ggplot}\NormalTok{(agreement\_df, }\FunctionTok{aes}\NormalTok{(}\AttributeTok{x =}\NormalTok{ elpd\_exact, }\AttributeTok{y =}\NormalTok{ elpd\_psis)) }\SpecialCharTok{+}
  \FunctionTok{geom\_abline}\NormalTok{(}\AttributeTok{intercept =} \DecValTok{0}\NormalTok{, }\AttributeTok{slope =} \DecValTok{1}\NormalTok{, }\AttributeTok{color =} \StringTok{"red"}\NormalTok{, }\AttributeTok{linetype =} \StringTok{"dashed"}\NormalTok{) }\SpecialCharTok{+}
  \FunctionTok{geom\_point}\NormalTok{(}\AttributeTok{size =} \DecValTok{3}\NormalTok{, }\AttributeTok{color =} \StringTok{"\#0072B2"}\NormalTok{) }\SpecialCharTok{+}
  \FunctionTok{labs}\NormalTok{(}
    \AttributeTok{title    =} \StringTok{"PSIS{-}LFO vs Exact One{-}Step{-}Ahead Refits"}\NormalTok{,}
    \AttributeTok{subtitle =} \StringTok{"Points should fall on the y = x line if PSIS{-}LFO is calibrated"}\NormalTok{,}
    \AttributeTok{x =} \StringTok{"Exact ELPD (refit per point)"}\NormalTok{,}
    \AttributeTok{y =} \StringTok{"PSIS{-}LFO ELPD"}
\NormalTok{  ) }\SpecialCharTok{+}
  \FunctionTok{coord\_fixed}\NormalTok{()}
\end{Highlighting}
\end{Shaded}

\begin{center}
\includegraphics[width=0.8\linewidth,height=\textheight,keepaspectratio]{11-CategorizationModelsExemplars_files/figure-pdf/ch12_exact_lfo_validation-1.pdf}
\end{center}

\subsubsection{Randomised LOO-PIT}\label{randomised-loo-pit}

\begin{Shaded}
\begin{Highlighting}[]
\ControlFlowTok{if}\NormalTok{ (}\SpecialCharTok{!}\FunctionTok{is.null}\NormalTok{(fit\_gcm\_single)) \{}
\NormalTok{  prob\_draws  }\OtherTok{\textless{}{-}}\NormalTok{ fit\_gcm\_single}\SpecialCharTok{$}\FunctionTok{draws}\NormalTok{(}\StringTok{"prob\_cat1"}\NormalTok{, }\AttributeTok{format =} \StringTok{"matrix"}\NormalTok{)}
\NormalTok{  log\_lik\_mat }\OtherTok{\textless{}{-}}\NormalTok{ fit\_gcm\_single}\SpecialCharTok{$}\FunctionTok{draws}\NormalTok{(}\StringTok{"log\_lik"}\NormalTok{,   }\AttributeTok{format =} \StringTok{"matrix"}\NormalTok{)}

\NormalTok{  psis\_obj }\OtherTok{\textless{}{-}}\NormalTok{ loo}\SpecialCharTok{::}\FunctionTok{psis}\NormalTok{(}\SpecialCharTok{{-}}\NormalTok{log\_lik\_mat, }\AttributeTok{r\_eff =} \ConstantTok{NA}\NormalTok{)}
\NormalTok{  lw       }\OtherTok{\textless{}{-}} \FunctionTok{weights}\NormalTok{(psis\_obj, }\AttributeTok{normalize =} \ConstantTok{TRUE}\NormalTok{)}

\NormalTok{  y\_obs }\OtherTok{\textless{}{-}}\NormalTok{ gcm\_data\_single}\SpecialCharTok{$}\NormalTok{y}

\NormalTok{  loo\_p\_cat1 }\OtherTok{\textless{}{-}} \FunctionTok{vapply}\NormalTok{(}\FunctionTok{seq\_along}\NormalTok{(y\_obs), }\ControlFlowTok{function}\NormalTok{(t) \{}
\NormalTok{    w }\OtherTok{\textless{}{-}} \FunctionTok{exp}\NormalTok{(lw[, t] }\SpecialCharTok{{-}}\NormalTok{ matrixStats}\SpecialCharTok{::}\FunctionTok{logSumExp}\NormalTok{(lw[, t]))}
    \FunctionTok{sum}\NormalTok{(w }\SpecialCharTok{*}\NormalTok{ prob\_draws[, t])}
\NormalTok{  \}, }\FunctionTok{numeric}\NormalTok{(}\DecValTok{1}\NormalTok{))}

  \FunctionTok{set.seed}\NormalTok{(}\DecValTok{2027}\NormalTok{)}
\NormalTok{  pit\_vals }\OtherTok{\textless{}{-}} \FunctionTok{ifelse}\NormalTok{(}
\NormalTok{    y\_obs }\SpecialCharTok{==} \DecValTok{1}\NormalTok{,}
    \FunctionTok{runif}\NormalTok{(}\FunctionTok{length}\NormalTok{(y\_obs), }\DecValTok{1} \SpecialCharTok{{-}}\NormalTok{ loo\_p\_cat1, }\DecValTok{1}\NormalTok{),}
    \FunctionTok{runif}\NormalTok{(}\FunctionTok{length}\NormalTok{(y\_obs), }\DecValTok{0}\NormalTok{,              }\DecValTok{1} \SpecialCharTok{{-}}\NormalTok{ loo\_p\_cat1)}
\NormalTok{  )}

  \FunctionTok{ggplot}\NormalTok{(}\FunctionTok{tibble}\NormalTok{(}\AttributeTok{pit =}\NormalTok{ pit\_vals), }\FunctionTok{aes}\NormalTok{(}\AttributeTok{x =}\NormalTok{ pit)) }\SpecialCharTok{+}
    \FunctionTok{geom\_histogram}\NormalTok{(}\FunctionTok{aes}\NormalTok{(}\AttributeTok{y =} \FunctionTok{after\_stat}\NormalTok{(density)),}
                   \AttributeTok{bins =} \DecValTok{20}\NormalTok{, }\AttributeTok{fill =} \StringTok{"\#0072B2"}\NormalTok{, }\AttributeTok{alpha =} \FloatTok{0.7}\NormalTok{,}
                   \AttributeTok{boundary =} \DecValTok{0}\NormalTok{) }\SpecialCharTok{+}
    \FunctionTok{geom\_hline}\NormalTok{(}\AttributeTok{yintercept =} \DecValTok{1}\NormalTok{, }\AttributeTok{color =} \StringTok{"red"}\NormalTok{, }\AttributeTok{linetype =} \StringTok{"dashed"}\NormalTok{) }\SpecialCharTok{+}
    \FunctionTok{scale\_x\_continuous}\NormalTok{(}\AttributeTok{limits =} \FunctionTok{c}\NormalTok{(}\DecValTok{0}\NormalTok{, }\DecValTok{1}\NormalTok{)) }\SpecialCharTok{+}
    \FunctionTok{labs}\NormalTok{(}
      \AttributeTok{title    =} \StringTok{"Randomised LOO{-}PIT for the Single{-}Subject GCM"}\NormalTok{,}
      \AttributeTok{subtitle =} \StringTok{"If the model is calibrated, PIT values are Uniform(0,1) (red dashed line)"}\NormalTok{,}
      \AttributeTok{x =} \StringTok{"Randomised LOO{-}PIT"}\NormalTok{, }\AttributeTok{y =} \StringTok{"Density"}
\NormalTok{    )}
\NormalTok{\}}
\end{Highlighting}
\end{Shaded}

\begin{center}
\includegraphics[width=0.8\linewidth,height=\textheight,keepaspectratio]{11-CategorizationModelsExemplars_files/figure-pdf/ch12_randomized_loo_pit-1.pdf}
\end{center}

\subsubsection{Prior Sensitivity
Analysis}\label{prior-sensitivity-analysis-1}

\begin{Shaded}
\begin{Highlighting}[]
\ControlFlowTok{if}\NormalTok{ (}\SpecialCharTok{!}\FunctionTok{is.null}\NormalTok{(fit\_gcm\_single)) \{}
\NormalTok{  ps }\OtherTok{\textless{}{-}}\NormalTok{ priorsense}\SpecialCharTok{::}\FunctionTok{powerscale\_sensitivity}\NormalTok{(fit\_gcm\_single)}
  \FunctionTok{print}\NormalTok{(ps)}

\NormalTok{  priorsense}\SpecialCharTok{::}\FunctionTok{powerscale\_plot\_dens}\NormalTok{(}
\NormalTok{    fit\_gcm\_single,}
    \AttributeTok{variables =} \FunctionTok{c}\NormalTok{(}\StringTok{"log\_c"}\NormalTok{, }\StringTok{"w[1]"}\NormalTok{, }\StringTok{"bias"}\NormalTok{)}
\NormalTok{  )}
\NormalTok{\}}
\end{Highlighting}
\end{Shaded}

\begin{verbatim}
Sensitivity based on cjs_dist
Prior selection: all priors
Likelihood selection: all data

      variable prior likelihood diagnosis
          w[1] 0.003      0.220         -
          w[2] 0.003      0.220         -
         log_c 0.010      0.112         -
          bias 0.000      0.090         -
             c 0.014      0.075         -
  prob_cat1[1] 0.000      0.090         -
  prob_cat1[2] 0.000      0.090         -
  prob_cat1[3] 0.000      0.090         -
  prob_cat1[4] 0.000      0.090         -
  prob_cat1[5] 0.007      0.117         -
  prob_cat1[6] 0.006      0.184         -
  prob_cat1[7] 0.007      0.240         -
  prob_cat1[8] 0.008      0.124         -
  prob_cat1[9] 0.001      0.095         -
 prob_cat1[10] 0.007      0.192         -
 prob_cat1[11] 0.001      0.075         -
 prob_cat1[12] 0.007      0.260         -
 prob_cat1[13] 0.008      0.254         -
 prob_cat1[14] 0.007      0.248         -
 prob_cat1[15] 0.001      0.127         -
 prob_cat1[16] 0.002      0.105         -
 prob_cat1[17] 0.001      0.084         -
 prob_cat1[18] 0.002      0.087         -
 prob_cat1[19] 0.007      0.219         -
 prob_cat1[20] 0.007      0.197         -
 prob_cat1[21] 0.007      0.215         -
 prob_cat1[22] 0.001      0.106         -
 prob_cat1[23] 0.007      0.262         -
 prob_cat1[24] 0.002      0.111         -
 prob_cat1[25] 0.001      0.095         -
 prob_cat1[26] 0.007      0.249         -
 prob_cat1[27] 0.002      0.093         -
 prob_cat1[28] 0.001      0.086         -
 prob_cat1[29] 0.007      0.195         -
 prob_cat1[30] 0.007      0.211         -
 prob_cat1[31] 0.008      0.242         -
 prob_cat1[32] 0.002      0.107         -
 prob_cat1[33] 0.001      0.084         -
 prob_cat1[34] 0.007      0.226         -
 prob_cat1[35] 0.007      0.221         -
 prob_cat1[36] 0.002      0.105         -
 prob_cat1[37] 0.007      0.210         -
 prob_cat1[38] 0.002      0.081         -
 prob_cat1[39] 0.001      0.103         -
 prob_cat1[40] 0.007      0.243         -
 prob_cat1[41] 0.002      0.083         -
 prob_cat1[42] 0.001      0.086         -
 prob_cat1[43] 0.007      0.225         -
 prob_cat1[44] 0.007      0.200         -
 prob_cat1[45] 0.007      0.247         -
 prob_cat1[46] 0.007      0.216         -
 prob_cat1[47] 0.001      0.098         -
 prob_cat1[48] 0.002      0.104         -
 prob_cat1[49] 0.001      0.095         -
 prob_cat1[50] 0.001      0.083         -
 prob_cat1[51] 0.007      0.215         -
 prob_cat1[52] 0.007      0.205         -
 prob_cat1[53] 0.007      0.250         -
 prob_cat1[54] 0.002      0.087         -
 prob_cat1[55] 0.002      0.104         -
 prob_cat1[56] 0.007      0.228         -
 prob_cat1[57] 0.001      0.095         -
 prob_cat1[58] 0.002      0.086         -
 prob_cat1[59] 0.007      0.230         -
 prob_cat1[60] 0.001      0.085         -
 prob_cat1[61] 0.007      0.246         -
 prob_cat1[62] 0.002      0.100         -
 prob_cat1[63] 0.007      0.200         -
 prob_cat1[64] 0.007      0.213         -
\end{verbatim}

\begin{center}
\includegraphics[width=0.8\linewidth,height=\textheight,keepaspectratio]{11-CategorizationModelsExemplars_files/figure-pdf/ch12_prior_sensitivity_single-1.pdf}
\end{center}

\begin{center}\rule{0.5\linewidth}{0.5pt}\end{center}

\subsection{Single-Subject Simulation-Based Calibration Checks
(SBC)}\label{single-subject-simulation-based-calibration-checks-sbc}

SBC checks the internal consistency and calibration of the entire
Bayesian inference procedure across the \emph{whole} prior. It asks:
``If the true parameters really came from our specified priors, would
our inference procedure produce posteriors that are statistically
consistent with those priors?''

\begin{Shaded}
\begin{Highlighting}[]
\NormalTok{sbc\_single\_filepath }\OtherTok{\textless{}{-}} \FunctionTok{here}\NormalTok{(}\StringTok{"simdata"}\NormalTok{, }\StringTok{"ch12\_sbc\_single\_results.rds"}\NormalTok{)}

\ControlFlowTok{if}\NormalTok{ (regenerate\_simulations }\SpecialCharTok{||} \SpecialCharTok{!}\FunctionTok{file.exists}\NormalTok{(sbc\_single\_filepath)) \{}
\NormalTok{  n\_sbc\_iterations\_single }\OtherTok{\textless{}{-}} \DecValTok{500}   \CommentTok{\# use \textgreater{}= 1000 for publication{-}quality SBC}

\NormalTok{  gcm\_single\_sbc\_generator }\OtherTok{\textless{}{-}} \FunctionTok{SBC\_generator\_function}\NormalTok{(}
    \ControlFlowTok{function}\NormalTok{() \{}
\NormalTok{      w\_true     }\OtherTok{\textless{}{-}} \FunctionTok{as.numeric}\NormalTok{(MCMCpack}\SpecialCharTok{::}\FunctionTok{rdirichlet}\NormalTok{(}\DecValTok{1}\NormalTok{, }\FunctionTok{c}\NormalTok{(}\DecValTok{1}\NormalTok{, }\DecValTok{1}\NormalTok{)))}
\NormalTok{      log\_c\_true }\OtherTok{\textless{}{-}} \FunctionTok{rnorm}\NormalTok{(}\DecValTok{1}\NormalTok{, LOG\_C\_PRIOR\_MEAN, LOG\_C\_PRIOR\_SD)}
\NormalTok{      c\_true     }\OtherTok{\textless{}{-}} \FunctionTok{exp}\NormalTok{(log\_c\_true)}
\NormalTok{      bias\_true  }\OtherTok{\textless{}{-}} \FunctionTok{rbeta}\NormalTok{(}\DecValTok{1}\NormalTok{, }\DecValTok{1}\NormalTok{, }\DecValTok{1}\NormalTok{)}

\NormalTok{      sbc\_schedule }\OtherTok{\textless{}{-}} \FunctionTok{make\_subject\_schedule}\NormalTok{(stimulus\_info, n\_blocks,}
                                            \AttributeTok{seed =} \FunctionTok{sample.int}\NormalTok{(}\FloatTok{1e8}\NormalTok{, }\DecValTok{1}\NormalTok{))}

\NormalTok{      sim }\OtherTok{\textless{}{-}} \FunctionTok{gcm\_simulate}\NormalTok{(}
        \AttributeTok{w            =}\NormalTok{ w\_true,}
        \AttributeTok{c            =}\NormalTok{ c\_true,}
        \AttributeTok{bias         =}\NormalTok{ bias\_true,}
        \AttributeTok{obs          =}\NormalTok{ sbc\_schedule }\SpecialCharTok{|\textgreater{}}\NormalTok{ dplyr}\SpecialCharTok{::}\FunctionTok{select}\NormalTok{(height, position),}
        \AttributeTok{cat\_feedback =}\NormalTok{ sbc\_schedule}\SpecialCharTok{$}\NormalTok{category\_feedback}
\NormalTok{      )}

      \FunctionTok{list}\NormalTok{(}
        \AttributeTok{variables =} \FunctionTok{list}\NormalTok{(}
          \StringTok{\textasciigrave{}}\AttributeTok{w[1]}\StringTok{\textasciigrave{}} \OtherTok{=}\NormalTok{ w\_true[}\DecValTok{1}\NormalTok{],}
          \StringTok{\textasciigrave{}}\AttributeTok{w[2]}\StringTok{\textasciigrave{}} \OtherTok{=}\NormalTok{ w\_true[}\DecValTok{2}\NormalTok{],}
          \AttributeTok{log\_c  =}\NormalTok{ log\_c\_true,}
          \AttributeTok{bias   =}\NormalTok{ bias\_true}
\NormalTok{        ),}
        \AttributeTok{generated =} \FunctionTok{list}\NormalTok{(}
          \AttributeTok{N\_total           =} \FunctionTok{nrow}\NormalTok{(sbc\_schedule),}
          \AttributeTok{N\_features        =} \DecValTok{2}\NormalTok{L,}
          \AttributeTok{y                 =} \FunctionTok{as.integer}\NormalTok{(sim}\SpecialCharTok{$}\NormalTok{sim\_response),}
          \AttributeTok{obs               =} \FunctionTok{as.matrix}\NormalTok{(sbc\_schedule[, }\FunctionTok{c}\NormalTok{(}\StringTok{"height"}\NormalTok{, }\StringTok{"position"}\NormalTok{)]),}
          \AttributeTok{cat\_feedback      =} \FunctionTok{as.integer}\NormalTok{(sbc\_schedule}\SpecialCharTok{$}\NormalTok{category\_feedback),}
          \AttributeTok{prior\_w\_alpha     =} \FunctionTok{c}\NormalTok{(}\DecValTok{1}\NormalTok{, }\DecValTok{1}\NormalTok{),}
          \AttributeTok{prior\_log\_c\_mu    =}\NormalTok{ LOG\_C\_PRIOR\_MEAN,}
          \AttributeTok{prior\_log\_c\_sigma =}\NormalTok{ LOG\_C\_PRIOR\_SD,}
          \AttributeTok{prior\_bias\_alpha  =} \DecValTok{1}\NormalTok{,}
          \AttributeTok{prior\_bias\_beta   =} \DecValTok{1}\NormalTok{,}
          \AttributeTok{run\_diagnostics   =} \DecValTok{1}\NormalTok{L}
\NormalTok{        )}
\NormalTok{      )}
\NormalTok{    \}}
\NormalTok{  )}

\NormalTok{  gcm\_single\_sbc\_backend }\OtherTok{\textless{}{-}} \FunctionTok{SBC\_backend\_cmdstan\_sample}\NormalTok{(}
\NormalTok{    mod\_gcm\_single,}
    \AttributeTok{iter\_warmup   =} \DecValTok{1000}\NormalTok{,}
    \AttributeTok{iter\_sampling =} \DecValTok{1000}\NormalTok{,}
    \AttributeTok{chains        =} \DecValTok{2}\NormalTok{,}
    \AttributeTok{adapt\_delta   =} \FloatTok{0.9}\NormalTok{,}
    \AttributeTok{refresh       =} \DecValTok{0}
\NormalTok{  )}

\NormalTok{  gcm\_single\_datasets }\OtherTok{\textless{}{-}} \FunctionTok{generate\_datasets}\NormalTok{(}
\NormalTok{    gcm\_single\_sbc\_generator,}
    \AttributeTok{n\_sims =}\NormalTok{ n\_sbc\_iterations\_single}
\NormalTok{  )}

\NormalTok{  sbc\_results\_single }\OtherTok{\textless{}{-}} \FunctionTok{compute\_SBC}\NormalTok{(}
\NormalTok{    gcm\_single\_datasets,}
\NormalTok{    gcm\_single\_sbc\_backend,}
    \AttributeTok{cache\_mode     =} \StringTok{"results"}\NormalTok{,}
    \AttributeTok{cache\_location =} \FunctionTok{here}\NormalTok{(}\StringTok{"simdata"}\NormalTok{, }\StringTok{"ch12\_sbc\_single\_cache"}\NormalTok{),}
    \AttributeTok{keep\_fits      =} \ConstantTok{FALSE}
\NormalTok{  )}
  \FunctionTok{saveRDS}\NormalTok{(sbc\_results\_single, sbc\_single\_filepath)}
  \FunctionTok{cat}\NormalTok{(}\StringTok{"SBC single{-}subject results computed and saved.}\SpecialCharTok{\textbackslash{}n}\StringTok{"}\NormalTok{)}
\NormalTok{\} }\ControlFlowTok{else}\NormalTok{ \{}
\NormalTok{  sbc\_results\_single }\OtherTok{\textless{}{-}} \FunctionTok{readRDS}\NormalTok{(sbc\_single\_filepath)}
  \FunctionTok{cat}\NormalTok{(}\StringTok{"Loaded existing SBC single{-}subject results.}\SpecialCharTok{\textbackslash{}n}\StringTok{"}\NormalTok{)}
\NormalTok{\}}
\end{Highlighting}
\end{Shaded}

\begin{verbatim}
Loaded existing SBC single-subject results.
\end{verbatim}

\begin{Shaded}
\begin{Highlighting}[]
\FunctionTok{plot\_ecdf\_diff}\NormalTok{(sbc\_results\_single)}
\end{Highlighting}
\end{Shaded}

\begin{center}
\includegraphics[width=0.8\linewidth,height=\textheight,keepaspectratio]{11-CategorizationModelsExemplars_files/figure-pdf/ch12_sbc_single-1.pdf}
\end{center}

\begin{Shaded}
\begin{Highlighting}[]
\FunctionTok{plot\_rank\_hist}\NormalTok{(sbc\_results\_single)}
\end{Highlighting}
\end{Shaded}

\begin{center}
\includegraphics[width=0.8\linewidth,height=\textheight,keepaspectratio]{11-CategorizationModelsExemplars_files/figure-pdf/ch12_sbc_single-2.pdf}
\end{center}

\begin{Shaded}
\begin{Highlighting}[]
\CommentTok{\# ── Diagnostic issue investigation ───────────────────────────────────────────}
\NormalTok{backend\_diags  }\OtherTok{\textless{}{-}}\NormalTok{ sbc\_results\_single}\SpecialCharTok{$}\NormalTok{backend\_diagnostics}
\NormalTok{default\_diags  }\OtherTok{\textless{}{-}}\NormalTok{ sbc\_results\_single}\SpecialCharTok{$}\NormalTok{default\_diagnostics }\SpecialCharTok{|\textgreater{}}
\NormalTok{  dplyr}\SpecialCharTok{::}\FunctionTok{select}\NormalTok{(sim\_id, max\_rhat, min\_ess\_to\_rank)}
\NormalTok{true\_params\_wide }\OtherTok{\textless{}{-}}\NormalTok{ sbc\_results\_single}\SpecialCharTok{$}\NormalTok{stats }\SpecialCharTok{|\textgreater{}}
\NormalTok{  dplyr}\SpecialCharTok{::}\FunctionTok{select}\NormalTok{(sim\_id, variable, simulated\_value) }\SpecialCharTok{|\textgreater{}}
\NormalTok{  tidyr}\SpecialCharTok{::}\FunctionTok{pivot\_wider}\NormalTok{(}\AttributeTok{names\_from =}\NormalTok{ variable, }\AttributeTok{values\_from =}\NormalTok{ simulated\_value)}

\NormalTok{diag\_df }\OtherTok{\textless{}{-}}\NormalTok{ backend\_diags }\SpecialCharTok{|\textgreater{}}
\NormalTok{  dplyr}\SpecialCharTok{::}\FunctionTok{left\_join}\NormalTok{(default\_diags,    }\AttributeTok{by =} \StringTok{"sim\_id"}\NormalTok{) }\SpecialCharTok{|\textgreater{}}
\NormalTok{  dplyr}\SpecialCharTok{::}\FunctionTok{left\_join}\NormalTok{(true\_params\_wide, }\AttributeTok{by =} \StringTok{"sim\_id"}\NormalTok{) }\SpecialCharTok{|\textgreater{}}
\NormalTok{  dplyr}\SpecialCharTok{::}\FunctionTok{mutate}\NormalTok{(}\AttributeTok{has\_issue =}\NormalTok{ n\_divergent }\SpecialCharTok{\textgreater{}} \DecValTok{0} \SpecialCharTok{|}\NormalTok{ n\_max\_treedepth }\SpecialCharTok{\textgreater{}} \DecValTok{0} \SpecialCharTok{|}\NormalTok{ max\_rhat }\SpecialCharTok{\textgreater{}} \FloatTok{1.01}\NormalTok{)}

\NormalTok{p\_diag\_c\_bias }\OtherTok{\textless{}{-}} \FunctionTok{ggplot}\NormalTok{(diag\_df, }\FunctionTok{aes}\NormalTok{(}\AttributeTok{x =}\NormalTok{ log\_c, }\AttributeTok{y =}\NormalTok{ bias, }\AttributeTok{color =}\NormalTok{ has\_issue)) }\SpecialCharTok{+}
  \FunctionTok{geom\_point}\NormalTok{(}\AttributeTok{alpha =} \FloatTok{0.7}\NormalTok{, }\AttributeTok{size =} \DecValTok{2}\NormalTok{) }\SpecialCharTok{+}
  \FunctionTok{scale\_color\_manual}\NormalTok{(}\AttributeTok{values =} \FunctionTok{c}\NormalTok{(}\StringTok{"FALSE"} \OtherTok{=} \StringTok{"\#0072B2"}\NormalTok{, }\StringTok{"TRUE"} \OtherTok{=} \StringTok{"\#D55E00"}\NormalTok{)) }\SpecialCharTok{+}
  \FunctionTok{labs}\NormalTok{(}\AttributeTok{title    =} \StringTok{"Diagnostic Issues in Parameter Space"}\NormalTok{,}
       \AttributeTok{subtitle =} \StringTok{"Orange points: divergence, max treedepth, or Rhat \textgreater{} 1.01"}\NormalTok{,}
       \AttributeTok{x =} \StringTok{"True log(c)"}\NormalTok{, }\AttributeTok{y =} \StringTok{"True Bias"}\NormalTok{) }\SpecialCharTok{+}
  \FunctionTok{theme}\NormalTok{(}\AttributeTok{legend.position =} \StringTok{"bottom"}\NormalTok{)}

\NormalTok{p\_diag\_w\_c }\OtherTok{\textless{}{-}} \FunctionTok{ggplot}\NormalTok{(diag\_df, }\FunctionTok{aes}\NormalTok{(}\AttributeTok{x =} \StringTok{\textasciigrave{}}\AttributeTok{w[1]}\StringTok{\textasciigrave{}}\NormalTok{, }\AttributeTok{y =}\NormalTok{ log\_c, }\AttributeTok{color =}\NormalTok{ has\_issue)) }\SpecialCharTok{+}
  \FunctionTok{geom\_point}\NormalTok{(}\AttributeTok{alpha =} \FloatTok{0.7}\NormalTok{, }\AttributeTok{size =} \DecValTok{2}\NormalTok{) }\SpecialCharTok{+}
  \FunctionTok{scale\_color\_manual}\NormalTok{(}\AttributeTok{values =} \FunctionTok{c}\NormalTok{(}\StringTok{"FALSE"} \OtherTok{=} \StringTok{"\#0072B2"}\NormalTok{, }\StringTok{"TRUE"} \OtherTok{=} \StringTok{"\#D55E00"}\NormalTok{)) }\SpecialCharTok{+}
  \FunctionTok{labs}\NormalTok{(}\AttributeTok{title =} \StringTok{"Attention vs Sensitivity"}\NormalTok{, }\AttributeTok{x =} \StringTok{"True w[1]"}\NormalTok{, }\AttributeTok{y =} \StringTok{"True log(c)"}\NormalTok{) }\SpecialCharTok{+}
  \FunctionTok{theme}\NormalTok{(}\AttributeTok{legend.position =} \StringTok{"none"}\NormalTok{)}

\FunctionTok{print}\NormalTok{(p\_diag\_c\_bias }\SpecialCharTok{|}\NormalTok{ p\_diag\_w\_c)}
\end{Highlighting}
\end{Shaded}

\begin{center}
\includegraphics[width=0.8\linewidth,height=\textheight,keepaspectratio]{11-CategorizationModelsExemplars_files/figure-pdf/ch12_sbc_single-3.pdf}
\end{center}

\begin{Shaded}
\begin{Highlighting}[]
\CommentTok{\# ── Parameter recovery on SBC simulations ────────────────────────────────────}
\NormalTok{recovery\_df }\OtherTok{\textless{}{-}}\NormalTok{ sbc\_results\_single}\SpecialCharTok{$}\NormalTok{stats}

\NormalTok{p\_sbc\_recovery }\OtherTok{\textless{}{-}} \FunctionTok{ggplot}\NormalTok{(recovery\_df, }\FunctionTok{aes}\NormalTok{(}\AttributeTok{x =}\NormalTok{ simulated\_value, }\AttributeTok{y =}\NormalTok{ mean)) }\SpecialCharTok{+}
  \FunctionTok{geom\_point}\NormalTok{(}\AttributeTok{alpha =} \FloatTok{0.4}\NormalTok{, }\AttributeTok{color =} \StringTok{"\#0072B2"}\NormalTok{) }\SpecialCharTok{+}
  \FunctionTok{geom\_abline}\NormalTok{(}\AttributeTok{intercept =} \DecValTok{0}\NormalTok{, }\AttributeTok{slope =} \DecValTok{1}\NormalTok{, }\AttributeTok{color =} \StringTok{"red"}\NormalTok{, }\AttributeTok{linetype =} \StringTok{"dashed"}\NormalTok{) }\SpecialCharTok{+}
  \FunctionTok{facet\_wrap}\NormalTok{(}\SpecialCharTok{\textasciitilde{}}\NormalTok{variable, }\AttributeTok{scales =} \StringTok{"free"}\NormalTok{) }\SpecialCharTok{+}
  \FunctionTok{labs}\NormalTok{(}
    \AttributeTok{title    =} \StringTok{"Parameter Recovery Across SBC Simulations"}\NormalTok{,}
    \AttributeTok{subtitle =} \StringTok{"Posterior mean vs True generated value"}\NormalTok{,}
    \AttributeTok{x =} \StringTok{"True Value (simulated\_value)"}\NormalTok{,}
    \AttributeTok{y =} \StringTok{"Estimated Value (Posterior Mean)"}
\NormalTok{  ) }\SpecialCharTok{+}
  \FunctionTok{theme\_bw}\NormalTok{()}

\FunctionTok{print}\NormalTok{(p\_sbc\_recovery)}
\end{Highlighting}
\end{Shaded}

\begin{center}
\includegraphics[width=0.8\linewidth,height=\textheight,keepaspectratio]{11-CategorizationModelsExemplars_files/figure-pdf/ch12_sbc_single-4.pdf}
\end{center}

When the true log(c) is above 0 the model computes similarity using
exponential decay: \texttt{exp(-c\ *\ distance)}. When c gets large,
similarity drops to exactly zero for anything except perfectly identical
stimuli, causing gradient collapse and HMC failures. The diagnostic
scatter plots above reveal that divergences and Rhat violations cluster
at high log(c) values, confirming that the prior mass above log(c) ≈ 2
(c ≈ 7--8) is unsupported by the sampler geometry.

However, calibration and parameter recovery look otherwise good.

\textbf{Practical constraint from SBC.} We treat c ≈ 12 as a hard
ceiling: a sensitivity this extreme makes stimuli that differ by any
measurable amount effectively identical, which is psychologically
implausible for the feature ranges used here. We therefore require that
the prior places negligible mass above c = 12, i.e., that the 99.85th
percentile (mean + 3 SD on the log scale) stays below log(12) ≈ 2.49.
For the single-subject model this is satisfied with
\texttt{Normal(log(2),\ 1.0)}, whose +3 SD upper bound is exp(0.69 +
3.0) ≈ 20 --- already at the boundary of acceptability. We carry this
constraint forward explicitly into the multilevel parameterisation.

\begin{center}\rule{0.5\linewidth}{0.5pt}\end{center}

\section{Extending the GCM with Memory
Decay}\label{extending-the-gcm-with-memory-decay}

The canonical GCM weights every stored exemplar equally --- a perfect,
unbounded memory. That is pedagogically clean but psychologically
implausible. Indeed, many early experiments had participants being
extensively trained on a set of stimuli, repeated many times, before
they actually got tested, so that they would resemble better the model's
assumptions. Further, an unbounded memory becomes a real liability as
soon as the learning environment is non-stationary: if the category
structure changes halfway through the task, the unweighted sum of
exemplars anchors the agent to contradictions it has no mechanism to
discard. We add a single forgetting-rate parameter, \(\lambda\), that
downweights older exemplars by \(e^{-\lambda (t - t_e)}\) before
entering the similarity average:

\[\eta_{ie} = e^{-\lambda (i - t_e)} \cdot e^{-c\, d_{ie}}, \qquad m_{k,i} = \frac{\sum_{e \in C_k} \eta_{ie}}{\sum_{e \in C_k} e^{-\lambda(i - t_e)}}, \qquad p_{i} = \frac{\beta \cdot m_{1,i}}{\beta \cdot m_{1,i} + (1 - \beta) \cdot m_{0,i}}.\]

\begin{itemize}
\tightlist
\item
  \(\eta_{ie}\) (Activation/Similarity): The total strength of the
  connection between the current stimulus \(i\) and a stored exemplar
  \(e\). It combines how similar they are in ``feature space'' (spatial)
  and how recent the memory is (temporal).
\item
  \(i\): The current trial index (the ``now'').
\item
  \(t_e\): The trial index when exemplar \(e\) was first encountered and
  stored in memory.
\item
  \((i - t_e)\): The age of the exemplar---how many trials have passed
  since it was seen.
\item
  \(d_{ie}\): The psychological distance between the current stimulus
  \(i\) and the stored exemplar \(e\) in feature space.
\item
  \(c\): The sensitivity parameter. It controls how quickly similarity
  drops as distance increases. A high \(c\) means only very similar
  items count; a low \(c\) means the agent generalizes broadly.
\item
  \(\lambda\): The forgetting-rate (decay) parameter. It controls how
  quickly memories fade with time.
\item
  \(m_{k,i}\): The match (or evidence) for category \(k\). This is a
  weighted average of similarities for all exemplars belonging to that
  category.
\item
  \(C_k\): The set of all exemplars that the agent has previously
  assigned to category \(k\).
\item
  \(p_i\): The probability that the agent classifies the current
  stimulus \(i\) into Category 1.
\item
  \(\beta\): The bias parameter. It represents a baseline preference for
  Category 1 (\(\beta > 0.5\)) or Category 0 (\(\beta < 0.5\))
  regardless of the evidence.
\end{itemize}

The parameter \(\lambda\) governs the Temporal Decay term:
\(e^{-\lambda (i - t_e)}\).

\begin{enumerate}
\def\labelenumi{\arabic{enumi}.}
\tightlist
\item
  The ``Memory Window'': Think of \(\lambda\) as controlling the size of
  the agent's moving window of attention.
\end{enumerate}

\begin{itemize}
\item
  If \(\lambda = 0\): The term \(e^{0} = 1\). The age of the memory has
  no effect. Every exemplar ever seen is weighted equally, representing
  a ``perfect,'' infinite memory.
\item
  If \(\lambda\) is large: The term \(e^{-\lambda \cdot age}\) drops to
  near-zero very quickly as age increases. The agent only ``remembers''
  or cares about the last few trials.
\item
  The Half-Life: You can calculate the ``half-life'' of a memory (the
  number of trials it takes for its influence to drop by half) as
  \(t_{1/2} = \frac{\ln(2)}{\lambda} \approx \frac{0.69}{\lambda}\). For
  example, if \(\lambda = 0.1\), a memory's influence is halved every
  \textasciitilde7 trials.
\end{itemize}

\begin{enumerate}
\def\labelenumi{\arabic{enumi}.}
\setcounter{enumi}{1}
\tightlist
\item
  Effective Sample Size: The denominator in the \(m_{k,i}\) formula,
  \(\sum_{e \in C_k} e^{-\lambda(i - t_e)}\), is the decay-weighted
  effective sample size.
\end{enumerate}

\begin{itemize}
\item
  Instead of just counting how many items are in Category \(k\) (which
  would be the raw count \(|C_k|\)), the model counts how many ``fresh''
  items are in Category \(k\).
\item
  If a category has 100 exemplars but they were all seen 500 trials ago,
  and \(\lambda\) is high, the ``effective count'' for that category
  will be near zero. This prevents old, potentially obsolete evidence
  from dominating the decision.
\end{itemize}

\begin{enumerate}
\def\labelenumi{\arabic{enumi}.}
\setcounter{enumi}{2}
\tightlist
\item
  Why use it? \(\lambda\) is essential for non-stationary environments.
  If the rules of the task change (e.g., the category boundary shifts),
  an agent with \(\lambda=0\) will be stuck trying to reconcile new
  evidence with a massive mountain of old, contradictory evidence. An
  agent with a well-tuned \(\lambda\) can ``forget'' the old rules and
  adapt to the new ones by letting the old evidence fade away.
\end{enumerate}

\begin{tcolorbox}[enhanced jigsaw, toprule=.15mm, toptitle=1mm, bottomtitle=1mm, leftrule=.75mm, rightrule=.15mm, colframe=quarto-callout-note-color-frame, left=2mm, colback=white, titlerule=0mm, coltitle=black, colbacktitle=quarto-callout-note-color!10!white, opacityback=0, title=\textcolor{quarto-callout-note-color}{\faInfo}\hspace{0.5em}{Canonical scenario framework (Chapters 11--13)}, arc=.35mm, opacitybacktitle=0.6, breakable, bottomrule=.15mm]

All three categorization chapters test their non-stationarity extension
against the same three environmental scenarios, presented in the same
order. The ordering reflects \textbf{increasing assumption match}
between the scenario's data-generating process and the model's
structural commitments. The scenarios are defined here once; Chapters 12
and 13 adopt the same definitions and ordering.

\begin{longtable}[]{@{}
  >{\raggedright\arraybackslash}p{(\linewidth - 4\tabcolsep) * \real{0.3333}}
  >{\raggedright\arraybackslash}p{(\linewidth - 4\tabcolsep) * \real{0.3333}}
  >{\raggedright\arraybackslash}p{(\linewidth - 4\tabcolsep) * \real{0.3333}}@{}}
\toprule\noalign{}
\begin{minipage}[b]{\linewidth}\raggedright
Scenario
\end{minipage} & \begin{minipage}[b]{\linewidth}\raggedright
What changes in the world
\end{minipage} & \begin{minipage}[b]{\linewidth}\raggedright
Structural role
\end{minipage} \\
\midrule\noalign{}
\endhead
\bottomrule\noalign{}
\endlastfoot
\textbf{Static Kruschke} & Nothing --- labels are fixed & Baseline: no
signal for the non-stationarity parameter \\
\textbf{Performance-contingent abrupt shifts} & Labels flip after a
streak of \(k\) correct responses & Misspecified: abrupt reversals that
no smooth-decay / Gaussian-drift / stochastic-mutation mechanism can
exactly represent \\
\textbf{Smooth drift} & Category centres follow a bivariate Gaussian
random walk & Assumption-matched (to varying degrees per model):
gradually stale evidence is what all three mechanisms are built to
handle \\
\end{longtable}

Comparing the three scenarios across chapters reveals how structural
differences between model families show up in their sensitivity to
non-stationarity. The synthesis is in the Trilogy Synthesis section at
the end of Chapter 13.

\end{tcolorbox}

We will test this extended model against the three scenarios:

\begin{enumerate}
\def\labelenumi{\arabic{enumi}.}
\item
  \textbf{Static Kruschke (baseline).} Fixed category labels, as in the
  canonical GCM analysis above. A low-\(\lambda\) model is correct for
  this data; we expect the posterior to be pulled toward the left tail
  of the \(\log\lambda\) prior and for the \(c\)/\(\lambda\) pair to be
  partially non-identified (both can make the similarity average
  narrower).
\item
  \textbf{Performance-contingent abrupt shifts.} After the agent reaches
  a streak of \(k\) consecutive correct responses, all category labels
  flip. The agent must unlearn the old mapping and relearn the new one.
  A high-\(\lambda\) model recovers fastest: it forgets the
  now-contradicted exemplars before they can anchor decisions. Note that
  these shifts are \emph{endogenous} to the agent's own behaviour --- a
  noisier agent triggers shifts less often and accumulates less
  information about \(\lambda\). Although this scenario is structurally
  misspecified (the decay GCM has no explicit abrupt-reversal
  mechanism), the post-flip recovery dynamics impose a strong constraint
  on \(\lambda\) that produces the clearest identifiability signal of
  the three scenarios.
\item
  \textbf{Smooth drift.} Category centres in feature space drift as an
  independent random walk. Labels are assigned trial-by-trial by
  proximity to the current (drifted) centres, so the category boundary
  migrates smoothly through the stimulus set. This is the
  \textbf{assumption-matched} scenario for the decay GCM: exponentially
  downweighting past exemplars is a defensible approximation to
  ``evidence collected under a different category mean is stale.''
  Moderate \(\lambda\) is appropriate here --- agents that forget too
  slowly lag behind the moving boundary; agents that forget too fast
  discard still-valid recent evidence.
\end{enumerate}

\subsection{Assumption Match Across
Scenarios}\label{assumption-match-across-scenarios}

\begin{longtable}[]{@{}
  >{\raggedright\arraybackslash}p{(\linewidth - 6\tabcolsep) * \real{0.2500}}
  >{\raggedright\arraybackslash}p{(\linewidth - 6\tabcolsep) * \real{0.2500}}
  >{\raggedright\arraybackslash}p{(\linewidth - 6\tabcolsep) * \real{0.2500}}
  >{\raggedright\arraybackslash}p{(\linewidth - 6\tabcolsep) * \real{0.2500}}@{}}
\toprule\noalign{}
\begin{minipage}[b]{\linewidth}\raggedright
Scenario
\end{minipage} & \begin{minipage}[b]{\linewidth}\raggedright
What drifts in the world
\end{minipage} & \begin{minipage}[b]{\linewidth}\raggedright
What decay GCM assumes
\end{minipage} & \begin{minipage}[b]{\linewidth}\raggedright
Match
\end{minipage} \\
\midrule\noalign{}
\endhead
\bottomrule\noalign{}
\endlastfoot
Static & Nothing & Exponential forgetting (correct at \(\lambda \to 0\))
& \textbf{Information-starved}: \(\lambda\) weakly identified \\
Contingent shift & Abrupt label reversals & Smooth exponential
forgetting & \textbf{Misspecified}: no abrupt mechanism, but post-flip
recovery provides strong \(\lambda\) signal \\
Smooth drift & Gaussian centre drift & Smooth exponential downweighting
of stale exemplars & \textbf{Graceful degradation}: approximate but
reasonable match \\
\end{longtable}

An important subtlety emerges from this table: \emph{assumption match
and parameter identifiability do not always align}. The contingent-shift
scenario is more misspecified than smooth drift, yet it produces the
cleanest \(\lambda\) recovery --- because abrupt reversals impose strong
temporal constraints that only \(\lambda\) can explain. The drift
scenario recovers \(\lambda\) at the correct calibration rate (SBC
passes) but without the ``bonus'' identifiability that misspecification
accidentally provides. We will see both signatures in the validation
results below.

\subsection{Prior on the Decay
Parameter}\label{prior-on-the-decay-parameter}

\begin{Shaded}
\begin{Highlighting}[]
\CommentTok{\# Prior hyperparameters for log(lambda).}
\CommentTok{\# Normal({-}2, 1) gives median lambda ≈ 0.135 (half{-}life ≈ 5 trials).}
\CommentTok{\# ±2 SD covers lambda ∈ (0.018, 1.0), spanning "effectively no decay"}
\CommentTok{\# to "only the last trial or two matters".}
\NormalTok{LOG\_LAMBDA\_PRIOR\_MEAN }\OtherTok{\textless{}{-}} \SpecialCharTok{{-}}\FloatTok{2.0}
\NormalTok{LOG\_LAMBDA\_PRIOR\_SD   }\OtherTok{\textless{}{-}} \FloatTok{1.0}
\end{Highlighting}
\end{Shaded}

\subsection{R Simulator With Scenario
Logic}\label{r-simulator-with-scenario-logic}

The simulator loops over trials, computes the decay-weighted choice
probability, draws a response, and determines the \emph{observed}
feedback according to the scenario. Memory stores the trial index at
which each exemplar was encoded, so ages are exact at every step. For
the contingent-shift scenario, the label-flip state machine updates
\emph{after} the response is drawn (so the current trial's feedback uses
the label map that was in force at decision time).

\begin{Shaded}
\begin{Highlighting}[]
\NormalTok{simulate\_decay\_scenario }\OtherTok{\textless{}{-}} \ControlFlowTok{function}\NormalTok{(w, c, lambda, bias, schedule,}
                                    \AttributeTok{scenario =} \FunctionTok{c}\NormalTok{(}\StringTok{"static"}\NormalTok{, }\StringTok{"contingent\_shift"}\NormalTok{, }\StringTok{"drift"}\NormalTok{),}
                                    \AttributeTok{streak\_target =} \DecValTok{8}\NormalTok{,}
                                    \AttributeTok{drift\_traj    =} \ConstantTok{NULL}\NormalTok{) \{}
\NormalTok{  scenario }\OtherTok{\textless{}{-}} \FunctionTok{match.arg}\NormalTok{(scenario)}
\NormalTok{  ntrials   }\OtherTok{\textless{}{-}} \FunctionTok{nrow}\NormalTok{(schedule)}
\NormalTok{  obs\_mat   }\OtherTok{\textless{}{-}} \FunctionTok{as.matrix}\NormalTok{(schedule[, }\FunctionTok{c}\NormalTok{(}\StringTok{"height"}\NormalTok{, }\StringTok{"position"}\NormalTok{)])}
\NormalTok{  base\_lab  }\OtherTok{\textless{}{-}}\NormalTok{ schedule}\SpecialCharTok{$}\NormalTok{category\_feedback}

\NormalTok{  memory\_obs   }\OtherTok{\textless{}{-}} \FunctionTok{matrix}\NormalTok{(}\ConstantTok{NA\_real\_}\NormalTok{, }\AttributeTok{nrow =} \DecValTok{0}\NormalTok{, }\AttributeTok{ncol =} \DecValTok{2}\NormalTok{)}
\NormalTok{  memory\_cat   }\OtherTok{\textless{}{-}} \FunctionTok{integer}\NormalTok{(}\DecValTok{0}\NormalTok{)}
\NormalTok{  memory\_trial }\OtherTok{\textless{}{-}} \FunctionTok{integer}\NormalTok{(}\DecValTok{0}\NormalTok{)}

\NormalTok{  label\_flip }\OtherTok{\textless{}{-}} \ConstantTok{FALSE}
\NormalTok{  streak     }\OtherTok{\textless{}{-}} \DecValTok{0}\NormalTok{L}

\NormalTok{  prob\_cat1         }\OtherTok{\textless{}{-}} \FunctionTok{numeric}\NormalTok{(ntrials)}
\NormalTok{  sim\_response      }\OtherTok{\textless{}{-}} \FunctionTok{integer}\NormalTok{(ntrials)}
\NormalTok{  observed\_feedback }\OtherTok{\textless{}{-}} \FunctionTok{integer}\NormalTok{(ntrials)}
\NormalTok{  label\_flip\_trace  }\OtherTok{\textless{}{-}} \FunctionTok{logical}\NormalTok{(ntrials)}

  \ControlFlowTok{for}\NormalTok{ (i }\ControlFlowTok{in} \FunctionTok{seq\_len}\NormalTok{(ntrials)) \{}
\NormalTok{    current\_stim }\OtherTok{\textless{}{-}} \FunctionTok{as.numeric}\NormalTok{(obs\_mat[i, ])}
\NormalTok{    n\_mem        }\OtherTok{\textless{}{-}} \FunctionTok{nrow}\NormalTok{(memory\_obs)}

    \CommentTok{\# ── Decay{-}weighted *mean* similarity per category ───────────────────}
    \CommentTok{\# Numerator:   Σ (decay\_j · similarity\_j) over exemplars in that category.}
    \CommentTok{\# Denominator: Σ decay\_j over exemplars in that category  (effective N).}
    \CommentTok{\# Collapses to plain mean similarity when lambda = 0.}
    \ControlFlowTok{if}\NormalTok{ (n\_mem }\SpecialCharTok{==} \DecValTok{0} \SpecialCharTok{||} \SpecialCharTok{!}\FunctionTok{any}\NormalTok{(memory\_cat }\SpecialCharTok{==} \DecValTok{0}\NormalTok{) }\SpecialCharTok{||} \SpecialCharTok{!}\FunctionTok{any}\NormalTok{(memory\_cat }\SpecialCharTok{==} \DecValTok{1}\NormalTok{)) \{}
\NormalTok{      p }\OtherTok{\textless{}{-}}\NormalTok{ bias}
\NormalTok{    \} }\ControlFlowTok{else}\NormalTok{ \{}
\NormalTok{      d     }\OtherTok{\textless{}{-}} \FunctionTok{as.numeric}\NormalTok{(}\FunctionTok{abs}\NormalTok{(}\FunctionTok{sweep}\NormalTok{(memory\_obs, }\DecValTok{2}\NormalTok{, current\_stim, }\StringTok{"{-}"}\NormalTok{)) }\SpecialCharTok{\%*\%}\NormalTok{ w)}
\NormalTok{      age   }\OtherTok{\textless{}{-}}\NormalTok{ i }\SpecialCharTok{{-}}\NormalTok{ memory\_trial}
\NormalTok{      decay }\OtherTok{\textless{}{-}} \FunctionTok{exp}\NormalTok{(}\SpecialCharTok{{-}}\NormalTok{lambda }\SpecialCharTok{*}\NormalTok{ age)}
\NormalTok{      sim   }\OtherTok{\textless{}{-}} \FunctionTok{exp}\NormalTok{(}\SpecialCharTok{{-}}\NormalTok{c }\SpecialCharTok{*}\NormalTok{ d)}
\NormalTok{      wt    }\OtherTok{\textless{}{-}}\NormalTok{ decay }\SpecialCharTok{*}\NormalTok{ sim}
\NormalTok{      s1    }\OtherTok{\textless{}{-}} \FunctionTok{sum}\NormalTok{(wt[memory\_cat }\SpecialCharTok{==} \DecValTok{1}\NormalTok{]);    s0 }\OtherTok{\textless{}{-}} \FunctionTok{sum}\NormalTok{(wt[memory\_cat }\SpecialCharTok{==} \DecValTok{0}\NormalTok{])}
\NormalTok{      w1    }\OtherTok{\textless{}{-}} \FunctionTok{sum}\NormalTok{(decay[memory\_cat }\SpecialCharTok{==} \DecValTok{1}\NormalTok{]); w0 }\OtherTok{\textless{}{-}} \FunctionTok{sum}\NormalTok{(decay[memory\_cat }\SpecialCharTok{==} \DecValTok{0}\NormalTok{])}
\NormalTok{      m1    }\OtherTok{\textless{}{-}} \ControlFlowTok{if}\NormalTok{ (w1 }\SpecialCharTok{\textgreater{}} \FloatTok{1e{-}12}\NormalTok{) s1 }\SpecialCharTok{/}\NormalTok{ w1 }\ControlFlowTok{else} \DecValTok{0}
\NormalTok{      m0    }\OtherTok{\textless{}{-}} \ControlFlowTok{if}\NormalTok{ (w0 }\SpecialCharTok{\textgreater{}} \FloatTok{1e{-}12}\NormalTok{) s0 }\SpecialCharTok{/}\NormalTok{ w0 }\ControlFlowTok{else} \DecValTok{0}
\NormalTok{      num   }\OtherTok{\textless{}{-}}\NormalTok{ bias }\SpecialCharTok{*}\NormalTok{ m1}
\NormalTok{      den   }\OtherTok{\textless{}{-}}\NormalTok{ num }\SpecialCharTok{+}\NormalTok{ (}\DecValTok{1} \SpecialCharTok{{-}}\NormalTok{ bias) }\SpecialCharTok{*}\NormalTok{ m0}
\NormalTok{      p     }\OtherTok{\textless{}{-}} \ControlFlowTok{if}\NormalTok{ (den }\SpecialCharTok{\textgreater{}} \FloatTok{1e{-}12}\NormalTok{) num }\SpecialCharTok{/}\NormalTok{ den }\ControlFlowTok{else}\NormalTok{ bias}
\NormalTok{    \}}
\NormalTok{    p }\OtherTok{\textless{}{-}} \FunctionTok{pmax}\NormalTok{(}\FloatTok{1e{-}9}\NormalTok{, }\FunctionTok{pmin}\NormalTok{(}\DecValTok{1} \SpecialCharTok{{-}} \FloatTok{1e{-}9}\NormalTok{, p))}
\NormalTok{    prob\_cat1[i]    }\OtherTok{\textless{}{-}}\NormalTok{ p}
\NormalTok{    sim\_response[i] }\OtherTok{\textless{}{-}} \FunctionTok{rbinom}\NormalTok{(}\DecValTok{1}\NormalTok{, }\DecValTok{1}\NormalTok{, p)}

    \CommentTok{\# ── Observed feedback depends on the scenario ───────────────────────}
\NormalTok{    fb }\OtherTok{\textless{}{-}} \ControlFlowTok{switch}\NormalTok{(scenario,}
      \AttributeTok{static           =}\NormalTok{ base\_lab[i],}
      \AttributeTok{contingent\_shift =} \ControlFlowTok{if}\NormalTok{ (label\_flip) }\DecValTok{1}\NormalTok{L }\SpecialCharTok{{-}}\NormalTok{ base\_lab[i] }\ControlFlowTok{else}\NormalTok{ base\_lab[i],}
      \AttributeTok{drift            =}\NormalTok{ \{}
\NormalTok{        d0 }\OtherTok{\textless{}{-}} \FunctionTok{sum}\NormalTok{(}\FunctionTok{abs}\NormalTok{(current\_stim }\SpecialCharTok{{-}}\NormalTok{ drift\_traj}\SpecialCharTok{$}\NormalTok{mu0[i, ]))}
\NormalTok{        d1 }\OtherTok{\textless{}{-}} \FunctionTok{sum}\NormalTok{(}\FunctionTok{abs}\NormalTok{(current\_stim }\SpecialCharTok{{-}}\NormalTok{ drift\_traj}\SpecialCharTok{$}\NormalTok{mu1[i, ]))}
        \FunctionTok{as.integer}\NormalTok{(d1 }\SpecialCharTok{\textless{}}\NormalTok{ d0)}
\NormalTok{      \}}
\NormalTok{    )}
\NormalTok{    observed\_feedback[i] }\OtherTok{\textless{}{-}}\NormalTok{ fb}
\NormalTok{    label\_flip\_trace[i]  }\OtherTok{\textless{}{-}}\NormalTok{ label\_flip}

    \CommentTok{\# Update the streak/flip state for the contingent scenario}
    \ControlFlowTok{if}\NormalTok{ (scenario }\SpecialCharTok{==} \StringTok{"contingent\_shift"}\NormalTok{) \{}
\NormalTok{      streak }\OtherTok{\textless{}{-}} \ControlFlowTok{if}\NormalTok{ (sim\_response[i] }\SpecialCharTok{==}\NormalTok{ fb) streak }\SpecialCharTok{+} \DecValTok{1}\NormalTok{L }\ControlFlowTok{else} \DecValTok{0}\NormalTok{L}
      \ControlFlowTok{if}\NormalTok{ (streak }\SpecialCharTok{\textgreater{}=}\NormalTok{ streak\_target) \{ label\_flip }\OtherTok{\textless{}{-}} \SpecialCharTok{!}\NormalTok{label\_flip; streak }\OtherTok{\textless{}{-}} \DecValTok{0}\NormalTok{L \}}
\NormalTok{    \}}

    \CommentTok{\# Add the current trial to memory}
\NormalTok{    memory\_obs   }\OtherTok{\textless{}{-}} \FunctionTok{rbind}\NormalTok{(memory\_obs, current\_stim)}
\NormalTok{    memory\_cat   }\OtherTok{\textless{}{-}} \FunctionTok{c}\NormalTok{(memory\_cat, fb)}
\NormalTok{    memory\_trial }\OtherTok{\textless{}{-}} \FunctionTok{c}\NormalTok{(memory\_trial, i)}
\NormalTok{  \}}

\NormalTok{  schedule }\SpecialCharTok{|\textgreater{}}
    \FunctionTok{mutate}\NormalTok{(}
      \AttributeTok{prob\_cat1         =}\NormalTok{ prob\_cat1,}
      \AttributeTok{sim\_response      =}\NormalTok{ sim\_response,}
      \AttributeTok{observed\_feedback =}\NormalTok{ observed\_feedback,}
      \AttributeTok{label\_flip        =}\NormalTok{ label\_flip\_trace,}
      \AttributeTok{correct           =} \FunctionTok{as.integer}\NormalTok{(sim\_response }\SpecialCharTok{==}\NormalTok{ observed\_feedback)}
\NormalTok{    )}
\NormalTok{\}}
\end{Highlighting}
\end{Shaded}

\subsection{Drift Trajectory
Generator}\label{drift-trajectory-generator}

The smooth-drift scenario needs a pre-computed trajectory of category
centres. Each centre follows an independent bivariate Gaussian random
walk with per-trial step SD \texttt{drift\_sigma}. Defaults place the
initial centres at the John K. Kruschke (1993) category means;
\texttt{drift\_sigma\ =\ 0.05} gives a cumulative SD of about 0.4 units
over 64 trials --- enough to migrate the boundary through roughly one
stimulus cell.

\begin{Shaded}
\begin{Highlighting}[]
\NormalTok{make\_drift\_trajectory }\OtherTok{\textless{}{-}} \ControlFlowTok{function}\NormalTok{(ntrials, }\AttributeTok{drift\_sigma =} \FloatTok{0.05}\NormalTok{, }\AttributeTok{seed =} \DecValTok{1}\NormalTok{) \{}
  \FunctionTok{set.seed}\NormalTok{(seed)}
\NormalTok{  mu0\_init }\OtherTok{\textless{}{-}} \FunctionTok{c}\NormalTok{(}\FloatTok{1.75}\NormalTok{, }\FloatTok{3.25}\NormalTok{)   }\CommentTok{\# Kruschke cat{-}0 mean (height, position)}
\NormalTok{  mu1\_init }\OtherTok{\textless{}{-}} \FunctionTok{c}\NormalTok{(}\FloatTok{3.25}\NormalTok{, }\FloatTok{1.75}\NormalTok{)   }\CommentTok{\# Kruschke cat{-}1 mean}
\NormalTok{  d0 }\OtherTok{\textless{}{-}} \FunctionTok{apply}\NormalTok{(}\FunctionTok{matrix}\NormalTok{(}\FunctionTok{rnorm}\NormalTok{(ntrials }\SpecialCharTok{*} \DecValTok{2}\NormalTok{, }\DecValTok{0}\NormalTok{, drift\_sigma), }\AttributeTok{ncol =} \DecValTok{2}\NormalTok{), }\DecValTok{2}\NormalTok{, cumsum)}
\NormalTok{  d1 }\OtherTok{\textless{}{-}} \FunctionTok{apply}\NormalTok{(}\FunctionTok{matrix}\NormalTok{(}\FunctionTok{rnorm}\NormalTok{(ntrials }\SpecialCharTok{*} \DecValTok{2}\NormalTok{, }\DecValTok{0}\NormalTok{, drift\_sigma), }\AttributeTok{ncol =} \DecValTok{2}\NormalTok{), }\DecValTok{2}\NormalTok{, cumsum)}
  \FunctionTok{list}\NormalTok{(}
    \AttributeTok{mu0 =} \FunctionTok{sweep}\NormalTok{(d0, }\DecValTok{2}\NormalTok{, mu0\_init, }\StringTok{"+"}\NormalTok{),}
    \AttributeTok{mu1 =} \FunctionTok{sweep}\NormalTok{(d1, }\DecValTok{2}\NormalTok{, mu1\_init, }\StringTok{"+"}\NormalTok{)}
\NormalTok{  )}
\NormalTok{\}}
\end{Highlighting}
\end{Shaded}

\begin{enumerate}
\def\labelenumi{\arabic{enumi}.}
\item
  The ``Ground Truth'' is a Random Walk The function
  make\_drift\_trajectory creates two paths (\(mu_0\) and \(mu_1\)).

  \begin{itemize}
  \item
    Initial State: On Trial 1, the categories start at the standard
    Kruschke positions (e.g., Cat 0 is at top-left, Cat 1 is at
    bottom-right).
  \item
    The Drift: Every single trial, both centers take a small random step
    in a random direction (North, South, East, or West). Because of the
    cumsum() (cumulative sum), these small steps add up.
  \item
    The Result: By Trial 64, the centers might have wandered quite far
    from where they started. The ``boundary'' between the
    categories---the line where you are equally likely to call something
    Cat 0 or Cat 1---is constantly shifting.
  \end{itemize}
\item
  How Exemplars are Generated When the simulation runs, it uses these
  drifting centers to create the stimuli the agent actually sees:
\item
  On Trial \(t\), the computer decides to show a ``Category 0''
  stimulus.
\item
  It looks up the current coordinates of the center: \(mu_{0,t}\).
\item
  It generates a stimulus near that point (usually by adding a bit of
  noise).
\item
  The agent sees this stimulus and is told: ``This is a Category 0.''
\end{enumerate}

\subsection{Scenario Agent Wrappers}\label{scenario-agent-wrappers}

\begin{Shaded}
\begin{Highlighting}[]
\NormalTok{simulate\_decay\_agent }\OtherTok{\textless{}{-}} \ControlFlowTok{function}\NormalTok{(agent\_id, scenario, w, c, lambda, bias,}
\NormalTok{                                 subject\_seed, }\AttributeTok{streak\_target =} \DecValTok{8}\NormalTok{,}
                                 \AttributeTok{drift\_sigma =} \FloatTok{0.05}\NormalTok{) \{}
\NormalTok{  schedule  }\OtherTok{\textless{}{-}} \FunctionTok{make\_subject\_schedule}\NormalTok{(stimulus\_info, n\_blocks, }\AttributeTok{seed =}\NormalTok{ subject\_seed)}
\NormalTok{  drift\_tr  }\OtherTok{\textless{}{-}} \ControlFlowTok{if}\NormalTok{ (scenario }\SpecialCharTok{==} \StringTok{"drift"}\NormalTok{)}
    \FunctionTok{make\_drift\_trajectory}\NormalTok{(}\FunctionTok{nrow}\NormalTok{(schedule), drift\_sigma, }\AttributeTok{seed =}\NormalTok{ subject\_seed)}
  \ControlFlowTok{else} \ConstantTok{NULL}

\NormalTok{  out }\OtherTok{\textless{}{-}} \FunctionTok{simulate\_decay\_scenario}\NormalTok{(}
    \AttributeTok{w =}\NormalTok{ w, }\AttributeTok{c =}\NormalTok{ c, }\AttributeTok{lambda =}\NormalTok{ lambda, }\AttributeTok{bias =}\NormalTok{ bias,}
    \AttributeTok{schedule      =}\NormalTok{ schedule,}
    \AttributeTok{scenario      =}\NormalTok{ scenario,}
    \AttributeTok{streak\_target =}\NormalTok{ streak\_target,}
    \AttributeTok{drift\_traj    =}\NormalTok{ drift\_tr}
\NormalTok{  )}

\NormalTok{  out }\SpecialCharTok{|\textgreater{}}
    \FunctionTok{mutate}\NormalTok{(}
      \AttributeTok{agent\_id       =}\NormalTok{ agent\_id,}
      \AttributeTok{scenario       =}\NormalTok{ scenario,}
      \AttributeTok{w1\_true        =}\NormalTok{ w[}\DecValTok{1}\NormalTok{],}
      \AttributeTok{c\_true         =}\NormalTok{ c,}
      \AttributeTok{lambda\_true    =}\NormalTok{ lambda,}
      \AttributeTok{log\_c\_true     =} \FunctionTok{log}\NormalTok{(c),}
      \AttributeTok{log\_lambda\_true =} \FunctionTok{log}\NormalTok{(lambda),}
      \AttributeTok{bias\_true      =}\NormalTok{ bias}
\NormalTok{    ) }\SpecialCharTok{|\textgreater{}}
    \FunctionTok{group\_by}\NormalTok{(agent\_id) }\SpecialCharTok{|\textgreater{}}
    \FunctionTok{mutate}\NormalTok{(}\AttributeTok{cumulative\_accuracy =} \FunctionTok{cumsum}\NormalTok{(correct) }\SpecialCharTok{/} \FunctionTok{row\_number}\NormalTok{()) }\SpecialCharTok{|\textgreater{}}
    \FunctionTok{ungroup}\NormalTok{()}
\NormalTok{\}}
\end{Highlighting}
\end{Shaded}

\subsection{Visualising the Three
Scenarios}\label{visualising-the-three-scenarios}

We simulate 10 agents per cell of a grid that crosses three \(\lambda\)
values with the three scenarios and plot cumulative accuracy. Static
data is where low \(\lambda\) wins. The contingent-shift data is where
moderate-to-high \(\lambda\) wins (faster forgetting speeds reversal
recovery, even though the model is misspecified for abrupt shifts). The
drift data rewards moderate \(\lambda\) tuned to the drift rate ---
agents that forget too slowly are anchored by obsolete exemplars; agents
that forget too fast discard still-valid recent evidence.

\begin{Shaded}
\begin{Highlighting}[]
\NormalTok{decay\_viz\_file }\OtherTok{\textless{}{-}} \FunctionTok{here}\NormalTok{(}\StringTok{"simdata"}\NormalTok{, }\StringTok{"ch12\_decay\_scenario\_viz.csv"}\NormalTok{)}

\ControlFlowTok{if}\NormalTok{ (regenerate\_simulations }\SpecialCharTok{||} \SpecialCharTok{!}\FunctionTok{file.exists}\NormalTok{(decay\_viz\_file)) \{}
\NormalTok{  viz\_grid }\OtherTok{\textless{}{-}} \FunctionTok{expand\_grid}\NormalTok{(}
    \AttributeTok{scenario    =} \FunctionTok{c}\NormalTok{(}\StringTok{"static"}\NormalTok{, }\StringTok{"contingent\_shift"}\NormalTok{, }\StringTok{"drift"}\NormalTok{),}
    \AttributeTok{lambda\_true =} \FunctionTok{c}\NormalTok{(}\FloatTok{0.02}\NormalTok{, }\FloatTok{0.2}\NormalTok{, }\FloatTok{0.6}\NormalTok{),}
    \AttributeTok{agent\_id    =} \DecValTok{1}\SpecialCharTok{:}\DecValTok{10}
\NormalTok{  ) }\SpecialCharTok{|\textgreater{}}
    \FunctionTok{mutate}\NormalTok{(}\AttributeTok{subject\_seed =} \DecValTok{10000} \SpecialCharTok{+} \FunctionTok{row\_number}\NormalTok{())}

\NormalTok{  viz\_sim }\OtherTok{\textless{}{-}} \FunctionTok{pmap\_dfr}\NormalTok{(}
    \FunctionTok{list}\NormalTok{(}\AttributeTok{agent\_id     =}\NormalTok{ viz\_grid}\SpecialCharTok{$}\NormalTok{agent\_id,}
         \AttributeTok{scenario     =}\NormalTok{ viz\_grid}\SpecialCharTok{$}\NormalTok{scenario,}
         \AttributeTok{lambda       =}\NormalTok{ viz\_grid}\SpecialCharTok{$}\NormalTok{lambda\_true,}
         \AttributeTok{subject\_seed =}\NormalTok{ viz\_grid}\SpecialCharTok{$}\NormalTok{subject\_seed),}
    \ControlFlowTok{function}\NormalTok{(agent\_id, scenario, lambda, subject\_seed) \{}
      \FunctionTok{simulate\_decay\_agent}\NormalTok{(}
        \AttributeTok{agent\_id     =}\NormalTok{ agent\_id,}
        \AttributeTok{scenario     =}\NormalTok{ scenario,}
        \AttributeTok{w            =} \FunctionTok{c}\NormalTok{(}\FloatTok{0.5}\NormalTok{, }\FloatTok{0.5}\NormalTok{),}
        \AttributeTok{c            =} \FloatTok{2.0}\NormalTok{,}
        \AttributeTok{lambda       =}\NormalTok{ lambda,}
        \AttributeTok{bias         =} \FloatTok{0.5}\NormalTok{,}
        \AttributeTok{subject\_seed =}\NormalTok{ subject\_seed}
\NormalTok{      )}
\NormalTok{    \}}
\NormalTok{  )}
  \FunctionTok{write\_csv}\NormalTok{(viz\_sim, decay\_viz\_file)}
\NormalTok{\} }\ControlFlowTok{else}\NormalTok{ \{}
\NormalTok{  viz\_sim }\OtherTok{\textless{}{-}} \FunctionTok{read\_csv}\NormalTok{(decay\_viz\_file, }\AttributeTok{show\_col\_types =} \ConstantTok{FALSE}\NormalTok{)}
\NormalTok{\}}

\FunctionTok{ggplot}\NormalTok{(viz\_sim }\SpecialCharTok{|\textgreater{}} 
         \FunctionTok{mutate}\NormalTok{(}\AttributeTok{scenario =} \FunctionTok{factor}\NormalTok{(scenario, }
                                 \AttributeTok{levels =} \FunctionTok{c}\NormalTok{(}\StringTok{"static"}\NormalTok{, }\StringTok{"contingent\_shift"}\NormalTok{, }\StringTok{"drift"}\NormalTok{),}
                                 \AttributeTok{labels =} \FunctionTok{c}\NormalTok{(}\StringTok{"Static Kruschke"}\NormalTok{, }\StringTok{"Performance{-}contingent flips"}\NormalTok{, }\StringTok{"Smooth drift"}\NormalTok{))),}
       \FunctionTok{aes}\NormalTok{(}\AttributeTok{x =}\NormalTok{ trial\_within\_subject, }\AttributeTok{y =}\NormalTok{ cumulative\_accuracy,}
           \AttributeTok{color =} \FunctionTok{factor}\NormalTok{(lambda\_true), }\AttributeTok{group =} \FunctionTok{interaction}\NormalTok{(agent\_id, lambda\_true))) }\SpecialCharTok{+}
  \FunctionTok{stat\_summary}\NormalTok{(}\AttributeTok{fun =}\NormalTok{ mean, }\AttributeTok{geom =} \StringTok{"line"}\NormalTok{,}
               \FunctionTok{aes}\NormalTok{(}\AttributeTok{group =} \FunctionTok{factor}\NormalTok{(lambda\_true)), }\AttributeTok{linewidth =} \FloatTok{1.1}\NormalTok{) }\SpecialCharTok{+}
  \FunctionTok{stat\_summary}\NormalTok{(}\AttributeTok{fun.data =}\NormalTok{ mean\_se, }\AttributeTok{geom =} \StringTok{"ribbon"}\NormalTok{, }\AttributeTok{alpha =} \FloatTok{0.15}\NormalTok{,}
               \FunctionTok{aes}\NormalTok{(}\AttributeTok{fill =} \FunctionTok{factor}\NormalTok{(lambda\_true), }\AttributeTok{group =} \FunctionTok{factor}\NormalTok{(lambda\_true)),}
               \AttributeTok{color =} \ConstantTok{NA}\NormalTok{) }\SpecialCharTok{+}
  \FunctionTok{facet\_wrap}\NormalTok{(}\SpecialCharTok{\textasciitilde{}}\NormalTok{scenario, }\AttributeTok{ncol =} \DecValTok{3}\NormalTok{) }\SpecialCharTok{+}
  \FunctionTok{scale\_color\_viridis\_d}\NormalTok{(}\AttributeTok{option =} \StringTok{"plasma"}\NormalTok{, }\AttributeTok{name =} \FunctionTok{expression}\NormalTok{(lambda)) }\SpecialCharTok{+}
  \FunctionTok{scale\_fill\_viridis\_d}\NormalTok{(}\AttributeTok{option  =} \StringTok{"plasma"}\NormalTok{, }\AttributeTok{name =} \FunctionTok{expression}\NormalTok{(lambda)) }\SpecialCharTok{+}
  \FunctionTok{geom\_hline}\NormalTok{(}\AttributeTok{yintercept =} \FloatTok{0.5}\NormalTok{, }\AttributeTok{linetype =} \StringTok{"dashed"}\NormalTok{, }\AttributeTok{color =} \StringTok{"grey50"}\NormalTok{) }\SpecialCharTok{+}
  \FunctionTok{labs}\NormalTok{(}
    \AttributeTok{title    =} \StringTok{"Learning Curves Across the Three Scenarios"}\NormalTok{,}
    \AttributeTok{subtitle =} \FunctionTok{paste0}\NormalTok{(}\StringTok{"Mean ± SE across 10 agents per cell; c = 2, bias = 0.5, "}\NormalTok{,}
                      \StringTok{"uniform attention"}\NormalTok{),}
    \AttributeTok{x =} \StringTok{"Trial"}\NormalTok{, }\AttributeTok{y =} \StringTok{"Cumulative accuracy"}
\NormalTok{  ) }\SpecialCharTok{+}
  \FunctionTok{theme}\NormalTok{(}\AttributeTok{legend.position =} \StringTok{"bottom"}\NormalTok{)}
\end{Highlighting}
\end{Shaded}

\begin{center}
\includegraphics[width=0.8\linewidth,height=\textheight,keepaspectratio]{11-CategorizationModelsExemplars_files/figure-pdf/ch12_decay_scenario_viz-1.pdf}
\end{center}

\textbf{Reading the scenario panels.} In the static panel, accuracy
climbs monotonically to near-ceiling for all \(\lambda\), with the
lowest \(\lambda\) reaching the highest asymptote---unsurprising,
because any decay discards still-valid evidence. In the drift panel, the
curves plateau below ceiling; the plateau is higher for agents with a
\(\lambda\) tuned to the drift rate, allowing them to track the moving
boundary without being ``anchored'' by obsolete history.

In the contingent-shift panel, we would see a characteristic sawtooth
pattern if we were plotting individuals, but this is somewhat muted by
the averaging. High-\(\lambda\) agents rebound faster after a flip,
which means they reach the criterion earlier and flip more often.
Because these flips happen at different trial numbers for different
agents, averaging across them smooths out the sharp drops and recoveries
into a noisy plateau.

\textbf{The Alignment Critique.} To see the true effect of \(\lambda\)
on cognitive flexibility, we must align the data to the moment of the
shift. In non-stationary environments, averaging across absolute trial
numbers often masks the very process we are interested in---the speed of
adaptation.

\begin{Shaded}
\begin{Highlighting}[]
\CommentTok{\# Align contingent{-}shift data to the start of each new phase}
\NormalTok{reversal\_viz }\OtherTok{\textless{}{-}}\NormalTok{ viz\_sim }\SpecialCharTok{|\textgreater{}}
  \FunctionTok{filter}\NormalTok{(scenario }\SpecialCharTok{==} \StringTok{"contingent\_shift"}\NormalTok{) }\SpecialCharTok{|\textgreater{}}
  \FunctionTok{group\_by}\NormalTok{(agent\_id, lambda\_true) }\SpecialCharTok{|\textgreater{}}
  \FunctionTok{mutate}\NormalTok{(}
    \CommentTok{\# Detect when the label mapping flipped}
    \AttributeTok{shift\_event =}\NormalTok{ label\_flip }\SpecialCharTok{!=} \FunctionTok{lag}\NormalTok{(label\_flip, }\AttributeTok{default =} \FunctionTok{first}\NormalTok{(label\_flip)),}
    \AttributeTok{phase\_id    =} \FunctionTok{cumsum}\NormalTok{(shift\_event)}
\NormalTok{  ) }\SpecialCharTok{|\textgreater{}}
  \FunctionTok{group\_by}\NormalTok{(agent\_id, lambda\_true, phase\_id) }\SpecialCharTok{|\textgreater{}}
  \FunctionTok{mutate}\NormalTok{(}
    \AttributeTok{trial\_in\_phase =} \FunctionTok{row\_number}\NormalTok{(),}
    \AttributeTok{phase\_accuracy =} \FunctionTok{cumsum}\NormalTok{(correct) }\SpecialCharTok{/}\NormalTok{ trial\_in\_phase}
\NormalTok{  ) }\SpecialCharTok{|\textgreater{}}
  \FunctionTok{ungroup}\NormalTok{() }\SpecialCharTok{|\textgreater{}}
  \CommentTok{\# Focus on the first few reversals to see the learning dynamics}
  \FunctionTok{filter}\NormalTok{(phase\_id }\SpecialCharTok{\textgreater{}} \DecValTok{0}\NormalTok{, phase\_id }\SpecialCharTok{\textless{}=} \DecValTok{4}\NormalTok{, trial\_in\_phase }\SpecialCharTok{\textless{}=} \DecValTok{40}\NormalTok{)}

\FunctionTok{ggplot}\NormalTok{(reversal\_viz, }\FunctionTok{aes}\NormalTok{(}\AttributeTok{x =}\NormalTok{ trial\_in\_phase, }\AttributeTok{y =}\NormalTok{ phase\_accuracy, }
                         \AttributeTok{color =} \FunctionTok{factor}\NormalTok{(lambda\_true), }\AttributeTok{group =} \FunctionTok{factor}\NormalTok{(lambda\_true))) }\SpecialCharTok{+}
  \FunctionTok{stat\_summary}\NormalTok{(}\AttributeTok{fun =}\NormalTok{ mean, }\AttributeTok{geom =} \StringTok{"line"}\NormalTok{, }\AttributeTok{linewidth =} \FloatTok{1.2}\NormalTok{) }\SpecialCharTok{+}
  \FunctionTok{stat\_summary}\NormalTok{(}\AttributeTok{fun.data =}\NormalTok{ mean\_se, }\AttributeTok{geom =} \StringTok{"ribbon"}\NormalTok{, }\AttributeTok{alpha =} \FloatTok{0.1}\NormalTok{, }\AttributeTok{color =} \ConstantTok{NA}\NormalTok{) }\SpecialCharTok{+}
  \FunctionTok{facet\_wrap}\NormalTok{(}\SpecialCharTok{\textasciitilde{}}\FunctionTok{paste}\NormalTok{(}\StringTok{"Phase"}\NormalTok{, phase\_id), }\AttributeTok{ncol =} \DecValTok{4}\NormalTok{) }\SpecialCharTok{+}
  \FunctionTok{scale\_color\_viridis\_d}\NormalTok{(}\AttributeTok{option =} \StringTok{"plasma"}\NormalTok{, }\AttributeTok{name =} \FunctionTok{expression}\NormalTok{(lambda)) }\SpecialCharTok{+}
  \FunctionTok{geom\_hline}\NormalTok{(}\AttributeTok{yintercept =} \FloatTok{0.5}\NormalTok{, }\AttributeTok{linetype =} \StringTok{"dashed"}\NormalTok{, }\AttributeTok{color =} \StringTok{"grey50"}\NormalTok{) }\SpecialCharTok{+}
  \FunctionTok{labs}\NormalTok{(}
    \AttributeTok{title    =} \StringTok{"Learning Performance Aligned to Category Shifts"}\NormalTok{,}
    \AttributeTok{subtitle =} \StringTok{"Cumulative accuracy reset at each reversal; contingent{-}shift scenario only"}\NormalTok{,}
    \AttributeTok{x =} \StringTok{"Trials Since Shift"}\NormalTok{, }\AttributeTok{y =} \StringTok{"Cumulative Phase Accuracy"}
\NormalTok{  ) }\SpecialCharTok{+}
  \FunctionTok{theme}\NormalTok{(}\AttributeTok{legend.position =} \StringTok{"bottom"}\NormalTok{)}
\end{Highlighting}
\end{Shaded}

\begin{center}
\includegraphics[width=0.8\linewidth,height=\textheight,keepaspectratio]{11-CategorizationModelsExemplars_files/figure-pdf/ch12_decay_reversal_alignment-1.pdf}
\end{center}

When aligned this way, the ``recovery advantage'' of high \(\lambda\)
becomes stark. While all agents collapse to chance (or below)
immediately after a shift, agents with higher decay rates
(\(\lambda = 0.6\)) shed their obsolete memory faster and return to high
accuracy significantly earlier than those with low decay.

\subsection{The Decay GCM in Stan}\label{the-decay-gcm-in-stan}

The Stan implementation adds \texttt{log\_lambda} as an unconstrained
scalar parameter and multiplies the similarity term by the decay weight
\texttt{exp(-lambda\ *\ age)} where
\texttt{age\ =\ i\ -\ past\_trial\_idx}. We also use a decay-weighted
mean similarity to ensure that the relative evidence for each category
is properly scaled as old exemplars fade.

\begin{Shaded}
\begin{Highlighting}[]
\NormalTok{gcm\_decay\_single\_stan }\OtherTok{\textless{}{-}} \StringTok{"}
\StringTok{// Generalized Context Model with Memory Decay — Single Subject}
\StringTok{data \{}
\StringTok{  int\textless{}lower=1\textgreater{} N\_total;}
\StringTok{  int\textless{}lower=1\textgreater{} N\_features;}
\StringTok{  array[N\_total] int\textless{}lower=0, upper=1\textgreater{} y;}
\StringTok{  array[N\_total, N\_features] real obs;}
\StringTok{  array[N\_total] int\textless{}lower=0, upper=1\textgreater{} cat\_feedback;}

\StringTok{  vector[N\_features] prior\_w\_alpha;}
\StringTok{  real prior\_log\_c\_mu;}
\StringTok{  real\textless{}lower=0\textgreater{} prior\_log\_c\_sigma;}
\StringTok{  real prior\_log\_lambda\_mu;}
\StringTok{  real\textless{}lower=0\textgreater{} prior\_log\_lambda\_sigma;}
\StringTok{  real\textless{}lower=0\textgreater{} prior\_bias\_alpha;}
\StringTok{  real\textless{}lower=0\textgreater{} prior\_bias\_beta;}
\StringTok{  int\textless{}lower=0, upper=1\textgreater{} run\_diagnostics;}
\StringTok{\}}

\StringTok{transformed data \{}
\StringTok{  array[N\_total, N\_total, N\_features] real abs\_diff;}
\StringTok{  for (i in 1:N\_total) \{}
\StringTok{    for (j in 1:N\_total) \{}
\StringTok{      for (f in 1:N\_features) \{}
\StringTok{        abs\_diff[i, j, f] = abs(obs[i, f] {-} obs[j, f]);}
\StringTok{      \}}
\StringTok{    \}}
\StringTok{  \}}
\StringTok{\}}

\StringTok{parameters \{}
\StringTok{  simplex[N\_features] w;}
\StringTok{  real log\_c;}
\StringTok{  real log\_lambda;}
\StringTok{  real\textless{}lower=0, upper=1\textgreater{} bias;}
\StringTok{\}}

\StringTok{transformed parameters \{}
\StringTok{  real\textless{}lower=0\textgreater{} c      = exp(log\_c);}
\StringTok{  real\textless{}lower=0\textgreater{} lambda = exp(log\_lambda);}

\StringTok{  vector\textless{}lower=1e{-}9, upper=1{-}1e{-}9\textgreater{}[N\_total] prob\_cat1;}

\StringTok{  \{}
\StringTok{    array[N\_total] int memory\_trial\_idx;}
\StringTok{    array[N\_total] int memory\_cat;}
\StringTok{    int n\_mem = 0;}

\StringTok{    for (i in 1:N\_total) \{}
\StringTok{      real p\_i;}
\StringTok{      int has\_cat0 = 0;}
\StringTok{      int has\_cat1 = 0;}

\StringTok{      for (k in 1:n\_mem) \{}
\StringTok{        if (memory\_cat[k] == 0) has\_cat0 = 1;}
\StringTok{        if (memory\_cat[k] == 1) has\_cat1 = 1;}
\StringTok{      \}}

\StringTok{      if (n\_mem == 0 || has\_cat0 == 0 || has\_cat1 == 0) \{}
\StringTok{        p\_i = bias;}
\StringTok{      \} else \{}
\StringTok{        real s0 = 0.0;   // sum of decay * similarity for cat 0}
\StringTok{        real s1 = 0.0;   // sum of decay * similarity for cat 1}
\StringTok{        real w0 = 0.0;   // sum of decay weights for cat 0 (effective N)}
\StringTok{        real w1 = 0.0;   // sum of decay weights for cat 1 (effective N)}
\StringTok{        for (e in 1:n\_mem) \{}
\StringTok{          real d = 0.0;}
\StringTok{          int past\_i = memory\_trial\_idx[e];}
\StringTok{          for (f in 1:N\_features)}
\StringTok{            d += w[f] * abs\_diff[i, past\_i, f];}
\StringTok{          real age          = i {-} past\_i;}
\StringTok{          real decay\_weight = exp({-}lambda * age);}
\StringTok{          real sim          = exp({-}c * d);}
\StringTok{          real weight       = decay\_weight * sim;}
\StringTok{          if (memory\_cat[e] == 1) \{ s1 += weight; w1 += decay\_weight; \}}
\StringTok{          else                    \{ s0 += weight; w0 += decay\_weight; \}}
\StringTok{        \}}
\StringTok{        real m1 = (w1 \textgreater{} 1e{-}12) ? s1 / w1 : 0.0;}
\StringTok{        real m0 = (w0 \textgreater{} 1e{-}12) ? s0 / w0 : 0.0;}
\StringTok{        real num = bias * m1;}
\StringTok{        real den = num + (1.0 {-} bias) * m0;}
\StringTok{        p\_i = (den \textgreater{} 1e{-}12) ? num / den : bias;}
\StringTok{      \}}

\StringTok{      prob\_cat1[i] = fmax(1e{-}9, fmin(1.0 {-} 1e{-}9, p\_i));}

\StringTok{      n\_mem += 1;}
\StringTok{      memory\_trial\_idx[n\_mem] = i;}
\StringTok{      memory\_cat[n\_mem] = cat\_feedback[i];}
\StringTok{    \}}
\StringTok{  \}}
\StringTok{\}}

\StringTok{model \{}
\StringTok{  target += dirichlet\_lpdf(w | prior\_w\_alpha);}
\StringTok{  target += normal\_lpdf(log\_c      | prior\_log\_c\_mu,      prior\_log\_c\_sigma);}
\StringTok{  target += normal\_lpdf(log\_lambda | prior\_log\_lambda\_mu, prior\_log\_lambda\_sigma);}
\StringTok{  target += beta\_lpdf(bias         | prior\_bias\_alpha,    prior\_bias\_beta);}

\StringTok{  target += bernoulli\_lpmf(y | prob\_cat1);}
\StringTok{\}}

\StringTok{generated quantities \{}
\StringTok{  vector[N\_total] log\_lik;}
\StringTok{  real lprior;}
\StringTok{  }
\StringTok{  if (run\_diagnostics) \{}
\StringTok{    for (i in 1:N\_total)}
\StringTok{      log\_lik[i] = bernoulli\_lpmf(y[i] | prob\_cat1[i]);}
\StringTok{  \}}
\StringTok{  }
\StringTok{  lprior =}
\StringTok{    dirichlet\_lpdf(w         | prior\_w\_alpha) +}
\StringTok{    normal\_lpdf(log\_c        | prior\_log\_c\_mu,      prior\_log\_c\_sigma) +}
\StringTok{    normal\_lpdf(log\_lambda   | prior\_log\_lambda\_mu, prior\_log\_lambda\_sigma) +}
\StringTok{    beta\_lpdf(bias           | prior\_bias\_alpha,    prior\_bias\_beta);}
\StringTok{\}}
\StringTok{"}

\NormalTok{stan\_file\_gcm\_decay\_single }\OtherTok{\textless{}{-}} \StringTok{"stan/ch12\_gcm\_decay\_single.stan"}
\FunctionTok{write\_stan\_file}\NormalTok{(gcm\_decay\_single\_stan, }\AttributeTok{dir =} \StringTok{"stan/"}\NormalTok{, }\AttributeTok{basename =} \StringTok{"ch12\_gcm\_decay\_single.stan"}\NormalTok{)}
\end{Highlighting}
\end{Shaded}

\begin{verbatim}
[1] "/Users/au209589/Dropbox/Teaching/AdvancedCognitiveModeling23_book/stan/ch12_gcm_decay_single.stan"
\end{verbatim}

\begin{Shaded}
\begin{Highlighting}[]
\NormalTok{mod\_gcm\_decay\_single }\OtherTok{\textless{}{-}} \FunctionTok{cmdstan\_model}\NormalTok{(stan\_file\_gcm\_decay\_single)}
\end{Highlighting}
\end{Shaded}

\subsection{Fitting One Agent Per
Scenario}\label{fitting-one-agent-per-scenario}

We pick one representative agent per scenario with \(\lambda = 0.2\),
\(c = 2\), bias \(= 0.5\), and uniform attention, and fit the decay GCM
to each.

\begin{Shaded}
\begin{Highlighting}[]
\NormalTok{decay\_fit\_file }\OtherTok{\textless{}{-}} \FunctionTok{here}\NormalTok{(}\StringTok{"simmodels"}\NormalTok{, }\StringTok{"ch12\_gcm\_decay\_single\_fits.rds"}\NormalTok{)}

\NormalTok{decay\_single\_fit\_data }\OtherTok{\textless{}{-}} \FunctionTok{list}\NormalTok{(}
  \AttributeTok{static           =} \FunctionTok{simulate\_decay\_agent}\NormalTok{(}\DecValTok{1}\NormalTok{, }\StringTok{"static"}\NormalTok{,}
                                          \AttributeTok{w =} \FunctionTok{c}\NormalTok{(}\FloatTok{0.5}\NormalTok{, }\FloatTok{0.5}\NormalTok{), }\AttributeTok{c =} \FloatTok{2.0}\NormalTok{,}
                                          \AttributeTok{lambda =} \FloatTok{0.2}\NormalTok{, }\AttributeTok{bias =} \FloatTok{0.5}\NormalTok{,}
                                          \AttributeTok{subject\_seed =} \DecValTok{20001}\NormalTok{),}
  \AttributeTok{contingent\_shift =} \FunctionTok{simulate\_decay\_agent}\NormalTok{(}\DecValTok{1}\NormalTok{, }\StringTok{"contingent\_shift"}\NormalTok{,}
                                          \AttributeTok{w =} \FunctionTok{c}\NormalTok{(}\FloatTok{0.5}\NormalTok{, }\FloatTok{0.5}\NormalTok{), }\AttributeTok{c =} \FloatTok{2.0}\NormalTok{,}
                                          \AttributeTok{lambda =} \FloatTok{0.2}\NormalTok{, }\AttributeTok{bias =} \FloatTok{0.5}\NormalTok{,}
                                          \AttributeTok{subject\_seed =} \DecValTok{20003}\NormalTok{),}
  \AttributeTok{drift            =} \FunctionTok{simulate\_decay\_agent}\NormalTok{(}\DecValTok{1}\NormalTok{, }\StringTok{"drift"}\NormalTok{,}
                                          \AttributeTok{w =} \FunctionTok{c}\NormalTok{(}\FloatTok{0.5}\NormalTok{, }\FloatTok{0.5}\NormalTok{), }\AttributeTok{c =} \FloatTok{2.0}\NormalTok{,}
                                          \AttributeTok{lambda =} \FloatTok{0.2}\NormalTok{, }\AttributeTok{bias =} \FloatTok{0.5}\NormalTok{,}
                                          \AttributeTok{subject\_seed =} \DecValTok{20002}\NormalTok{)}
\NormalTok{)}

\NormalTok{stan\_data\_from\_agent }\OtherTok{\textless{}{-}} \ControlFlowTok{function}\NormalTok{(agent\_df) \{}
  \FunctionTok{list}\NormalTok{(}
    \AttributeTok{N\_total               =} \FunctionTok{nrow}\NormalTok{(agent\_df),}
    \AttributeTok{N\_features             =} \DecValTok{2}\NormalTok{L,}
    \AttributeTok{y                     =} \FunctionTok{as.integer}\NormalTok{(agent\_df}\SpecialCharTok{$}\NormalTok{sim\_response),}
    \AttributeTok{obs                   =} \FunctionTok{as.matrix}\NormalTok{(agent\_df[, }\FunctionTok{c}\NormalTok{(}\StringTok{"height"}\NormalTok{, }\StringTok{"position"}\NormalTok{)]),}
    \AttributeTok{cat\_feedback          =} \FunctionTok{as.integer}\NormalTok{(agent\_df}\SpecialCharTok{$}\NormalTok{observed\_feedback),}
    \AttributeTok{prior\_w\_alpha         =} \FunctionTok{c}\NormalTok{(}\DecValTok{1}\NormalTok{, }\DecValTok{1}\NormalTok{),}
    \AttributeTok{prior\_log\_c\_mu        =}\NormalTok{ LOG\_C\_PRIOR\_MEAN,}
    \AttributeTok{prior\_log\_c\_sigma     =}\NormalTok{ LOG\_C\_PRIOR\_SD,}
    \AttributeTok{prior\_log\_lambda\_mu   =}\NormalTok{ LOG\_LAMBDA\_PRIOR\_MEAN,}
    \AttributeTok{prior\_log\_lambda\_sigma =}\NormalTok{ LOG\_LAMBDA\_PRIOR\_SD,}
    \AttributeTok{prior\_bias\_alpha      =} \DecValTok{1}\NormalTok{,}
    \AttributeTok{prior\_bias\_beta       =} \DecValTok{1}\NormalTok{,}
    \AttributeTok{run\_diagnostics       =} \DecValTok{1}
\NormalTok{  )}
\NormalTok{\}}

\ControlFlowTok{if}\NormalTok{ (regenerate\_simulations }\SpecialCharTok{||} \SpecialCharTok{!}\FunctionTok{file.exists}\NormalTok{(decay\_fit\_file)) \{}
\NormalTok{  decay\_fits }\OtherTok{\textless{}{-}}\NormalTok{ purrr}\SpecialCharTok{::}\FunctionTok{imap}\NormalTok{(decay\_single\_fit\_data, }\ControlFlowTok{function}\NormalTok{(agent\_df, scen) \{}
\NormalTok{    dat }\OtherTok{\textless{}{-}} \FunctionTok{stan\_data\_from\_agent}\NormalTok{(agent\_df)}
\NormalTok{    pf  }\OtherTok{\textless{}{-}} \FunctionTok{tryCatch}\NormalTok{(}
\NormalTok{      mod\_gcm\_decay\_single}\SpecialCharTok{$}\FunctionTok{pathfinder}\NormalTok{(}\AttributeTok{data =}\NormalTok{ dat, }\AttributeTok{num\_paths =} \DecValTok{4}\NormalTok{,}
                                      \AttributeTok{refresh =} \DecValTok{0}\NormalTok{, }\AttributeTok{seed =} \DecValTok{500}\NormalTok{),}
      \AttributeTok{error =} \ControlFlowTok{function}\NormalTok{(e) }\ConstantTok{NULL}
\NormalTok{    )}
\NormalTok{    fit }\OtherTok{\textless{}{-}}\NormalTok{ mod\_gcm\_decay\_single}\SpecialCharTok{$}\FunctionTok{sample}\NormalTok{(}
      \AttributeTok{data            =}\NormalTok{ dat,}
      \AttributeTok{init            =} \ControlFlowTok{if}\NormalTok{ (}\SpecialCharTok{!}\FunctionTok{is.null}\NormalTok{(pf)) pf }\ControlFlowTok{else} \FloatTok{0.5}\NormalTok{,}
      \AttributeTok{seed            =} \DecValTok{500}\NormalTok{,}
      \AttributeTok{chains          =} \DecValTok{4}\NormalTok{,}
      \AttributeTok{parallel\_chains =} \FunctionTok{min}\NormalTok{(}\DecValTok{4}\NormalTok{, }\FunctionTok{availableCores}\NormalTok{()),}
      \AttributeTok{iter\_warmup     =} \DecValTok{1000}\NormalTok{,}
      \AttributeTok{iter\_sampling   =} \DecValTok{1500}\NormalTok{,}
      \AttributeTok{refresh         =} \DecValTok{0}\NormalTok{,}
      \AttributeTok{adapt\_delta     =} \FloatTok{0.9}
\NormalTok{    )}
\NormalTok{    fit}\SpecialCharTok{$}\FunctionTok{save\_object}\NormalTok{(}\FunctionTok{here}\NormalTok{(}\StringTok{"simmodels"}\NormalTok{,}
                         \FunctionTok{paste0}\NormalTok{(}\StringTok{"ch12\_gcm\_decay\_single\_"}\NormalTok{, scen, }\StringTok{".rds"}\NormalTok{)))}
\NormalTok{    fit}
\NormalTok{  \})}
  \FunctionTok{saveRDS}\NormalTok{(}\FunctionTok{lapply}\NormalTok{(decay\_fits, }\ControlFlowTok{function}\NormalTok{(f) f}\SpecialCharTok{$}\FunctionTok{output\_files}\NormalTok{()), decay\_fit\_file)}
  \FunctionTok{cat}\NormalTok{(}\StringTok{"Decay GCM fits computed and saved.}\SpecialCharTok{\textbackslash{}n}\StringTok{"}\NormalTok{)}
\NormalTok{\} }\ControlFlowTok{else}\NormalTok{ \{}
\NormalTok{  decay\_fits }\OtherTok{\textless{}{-}} \FunctionTok{setNames}\NormalTok{(}
    \FunctionTok{lapply}\NormalTok{(}\FunctionTok{names}\NormalTok{(decay\_single\_fit\_data), }\ControlFlowTok{function}\NormalTok{(scen) \{}
      \FunctionTok{readRDS}\NormalTok{(}\FunctionTok{here}\NormalTok{(}\StringTok{"simmodels"}\NormalTok{,}
                   \FunctionTok{paste0}\NormalTok{(}\StringTok{"ch12\_gcm\_decay\_single\_"}\NormalTok{, scen, }\StringTok{".rds"}\NormalTok{)))}
\NormalTok{    \}),}
    \FunctionTok{names}\NormalTok{(decay\_single\_fit\_data)}
\NormalTok{  )}
  \FunctionTok{cat}\NormalTok{(}\StringTok{"Loaded existing decay GCM fits.}\SpecialCharTok{\textbackslash{}n}\StringTok{"}\NormalTok{)}
\NormalTok{\}}
\end{Highlighting}
\end{Shaded}

\begin{verbatim}
Loaded existing decay GCM fits.
\end{verbatim}

\subsection{MCMC Diagnostic Battery Across
Scenarios}\label{mcmc-diagnostic-battery-across-scenarios}

\begin{Shaded}
\begin{Highlighting}[]
\NormalTok{decay\_diag\_tbls }\OtherTok{\textless{}{-}}\NormalTok{ purrr}\SpecialCharTok{::}\FunctionTok{imap\_dfr}\NormalTok{(decay\_fits, }\ControlFlowTok{function}\NormalTok{(fit, scen) \{}
\NormalTok{  tbl }\OtherTok{\textless{}{-}} \FunctionTok{diagnostic\_summary\_table}\NormalTok{(}
\NormalTok{    fit,}
    \AttributeTok{params =} \FunctionTok{c}\NormalTok{(}\StringTok{"w[1]"}\NormalTok{, }\StringTok{"log\_c"}\NormalTok{, }\StringTok{"log\_lambda"}\NormalTok{, }\StringTok{"bias"}\NormalTok{)}
\NormalTok{  )}
\NormalTok{  tbl }\SpecialCharTok{|\textgreater{}} \FunctionTok{mutate}\NormalTok{(}\AttributeTok{scenario =}\NormalTok{ scen)}
\NormalTok{\})}

\FunctionTok{print}\NormalTok{(decay\_diag\_tbls }\SpecialCharTok{|\textgreater{}}
\NormalTok{        dplyr}\SpecialCharTok{::}\FunctionTok{select}\NormalTok{(scenario, metric, value, threshold, pass))}
\end{Highlighting}
\end{Shaded}

\begin{verbatim}
# A tibble: 18 x 5
   scenario         metric                           value threshold pass 
   <chr>            <chr>                            <dbl> <chr>     <lgl>
 1 static           Divergences (zero tolerance)    0      == 0      TRUE 
 2 static           Max rank-normalised R-hat       1.00   < 1.01    TRUE 
 3 static           Min bulk ESS                 3855.     > 400     TRUE 
 4 static           Min tail ESS                 3423.     > 400     TRUE 
 5 static           Min E-BFMI                      0.978  > 0.2     TRUE 
 6 static           Max MCSE / posterior SD         0.0161 < 0.05    TRUE 
 7 contingent_shift Divergences (zero tolerance)    0      == 0      TRUE 
 8 contingent_shift Max rank-normalised R-hat       1.00   < 1.01    TRUE 
 9 contingent_shift Min bulk ESS                 4883.     > 400     TRUE 
10 contingent_shift Min tail ESS                 3702.     > 400     TRUE 
11 contingent_shift Min E-BFMI                      0.973  > 0.2     TRUE 
12 contingent_shift Max MCSE / posterior SD         0.0147 < 0.05    TRUE 
13 drift            Divergences (zero tolerance)    0      == 0      TRUE 
14 drift            Max rank-normalised R-hat       1.00   < 1.01    TRUE 
15 drift            Min bulk ESS                 4356.     > 400     TRUE 
16 drift            Min tail ESS                 3771.     > 400     TRUE 
17 drift            Min E-BFMI                      1.00   > 0.2     TRUE 
18 drift            Max MCSE / posterior SD         0.0156 < 0.05    TRUE 
\end{verbatim}

\begin{Shaded}
\begin{Highlighting}[]
\FunctionTok{walk}\NormalTok{(}\FunctionTok{names}\NormalTok{(decay\_fits), }\ControlFlowTok{function}\NormalTok{(scen) \{}
  \FunctionTok{cat}\NormalTok{(}\StringTok{"=== Pair plot:"}\NormalTok{, scen, }\StringTok{"===}\SpecialCharTok{\textbackslash{}n}\StringTok{"}\NormalTok{)}
  \FunctionTok{print}\NormalTok{(bayesplot}\SpecialCharTok{::}\FunctionTok{mcmc\_pairs}\NormalTok{(}
\NormalTok{    decay\_fits[[scen]]}\SpecialCharTok{$}\FunctionTok{draws}\NormalTok{(}\FunctionTok{c}\NormalTok{(}\StringTok{"log\_c"}\NormalTok{, }\StringTok{"log\_lambda"}\NormalTok{, }\StringTok{"w[1]"}\NormalTok{, }\StringTok{"bias"}\NormalTok{)),}
    \AttributeTok{diag\_fun     =} \StringTok{"dens"}\NormalTok{,}
    \AttributeTok{off\_diag\_fun =} \StringTok{"hex"}
\NormalTok{  ))}
\NormalTok{\})}
\end{Highlighting}
\end{Shaded}

\begin{verbatim}
=== Pair plot: static ===
\end{verbatim}

\begin{center}
\includegraphics[width=0.8\linewidth,height=\textheight,keepaspectratio]{11-CategorizationModelsExemplars_files/figure-pdf/ch12_decay_diagnostics-1.pdf}
\end{center}

\begin{verbatim}
=== Pair plot: contingent_shift ===
\end{verbatim}

\begin{center}
\includegraphics[width=0.8\linewidth,height=\textheight,keepaspectratio]{11-CategorizationModelsExemplars_files/figure-pdf/ch12_decay_diagnostics-2.pdf}
\end{center}

\begin{verbatim}
=== Pair plot: drift ===
\end{verbatim}

\begin{center}
\includegraphics[width=0.8\linewidth,height=\textheight,keepaspectratio]{11-CategorizationModelsExemplars_files/figure-pdf/ch12_decay_diagnostics-3.pdf}
\end{center}

\textbf{What to watch for in the pair plot.} A
\texttt{log\_c}/\texttt{log\_lambda} ridge is expected on static data
--- both narrow the effective similarity average, so the data
under-constrain their difference. On contingent-shift data, post-flip
reversal dynamics inject strong temporal evidence that tensions the two
parameters against each other, and the ridge should shrink markedly. On
drift data, the ridge also narrows (boundary tracking exercises
\(\lambda\)) but typically less sharply, consistent with the weaker
identifiability signal that assumption-matched tasks provide.

\subsection{Prior-Posterior Update Plots Across
Scenarios}\label{prior-posterior-update-plots-across-scenarios}

\begin{Shaded}
\begin{Highlighting}[]
\NormalTok{decay\_pp\_df }\OtherTok{\textless{}{-}}\NormalTok{ purrr}\SpecialCharTok{::}\FunctionTok{imap\_dfr}\NormalTok{(decay\_fits, }\ControlFlowTok{function}\NormalTok{(fit, scen) \{}
\NormalTok{  d }\OtherTok{\textless{}{-}} \FunctionTok{as\_draws\_df}\NormalTok{(fit}\SpecialCharTok{$}\FunctionTok{draws}\NormalTok{(}\FunctionTok{c}\NormalTok{(}\StringTok{"log\_c"}\NormalTok{, }\StringTok{"log\_lambda"}\NormalTok{, }\StringTok{"w[1]"}\NormalTok{, }\StringTok{"bias"}\NormalTok{)))}
  \FunctionTok{bind\_rows}\NormalTok{(}
    \FunctionTok{tibble}\NormalTok{(}\AttributeTok{scenario =}\NormalTok{ scen, }\AttributeTok{parameter =} \StringTok{"log\_c"}\NormalTok{,      }\AttributeTok{value =}\NormalTok{ d}\SpecialCharTok{$}\NormalTok{log\_c,      }\AttributeTok{source =} \StringTok{"posterior"}\NormalTok{),}
    \FunctionTok{tibble}\NormalTok{(}\AttributeTok{scenario =}\NormalTok{ scen, }\AttributeTok{parameter =} \StringTok{"log\_lambda"}\NormalTok{, }\AttributeTok{value =}\NormalTok{ d}\SpecialCharTok{$}\NormalTok{log\_lambda, }\AttributeTok{source =} \StringTok{"posterior"}\NormalTok{),}
    \FunctionTok{tibble}\NormalTok{(}\AttributeTok{scenario =}\NormalTok{ scen, }\AttributeTok{parameter =} \StringTok{"w[1]"}\NormalTok{,       }\AttributeTok{value =}\NormalTok{ d}\SpecialCharTok{$}\StringTok{\textasciigrave{}}\AttributeTok{w[1]}\StringTok{\textasciigrave{}}\NormalTok{,     }\AttributeTok{source =} \StringTok{"posterior"}\NormalTok{),}
    \FunctionTok{tibble}\NormalTok{(}\AttributeTok{scenario =}\NormalTok{ scen, }\AttributeTok{parameter =} \StringTok{"bias"}\NormalTok{,       }\AttributeTok{value =}\NormalTok{ d}\SpecialCharTok{$}\NormalTok{bias,       }\AttributeTok{source =} \StringTok{"posterior"}\NormalTok{)}
\NormalTok{  )}
\NormalTok{\})}

\FunctionTok{set.seed}\NormalTok{(}\DecValTok{11}\NormalTok{)}
\NormalTok{n\_prior }\OtherTok{\textless{}{-}} \DecValTok{4000}
\NormalTok{prior\_ref }\OtherTok{\textless{}{-}} \FunctionTok{bind\_rows}\NormalTok{(}
  \FunctionTok{tibble}\NormalTok{(}\AttributeTok{parameter =} \StringTok{"log\_c"}\NormalTok{,      }\AttributeTok{value =} \FunctionTok{rnorm}\NormalTok{(n\_prior, LOG\_C\_PRIOR\_MEAN,      LOG\_C\_PRIOR\_SD),}
         \AttributeTok{source =} \StringTok{"prior"}\NormalTok{),}
  \FunctionTok{tibble}\NormalTok{(}\AttributeTok{parameter =} \StringTok{"log\_lambda"}\NormalTok{, }\AttributeTok{value =} \FunctionTok{rnorm}\NormalTok{(n\_prior, LOG\_LAMBDA\_PRIOR\_MEAN, LOG\_LAMBDA\_PRIOR\_SD),}
         \AttributeTok{source =} \StringTok{"prior"}\NormalTok{),}
  \FunctionTok{tibble}\NormalTok{(}\AttributeTok{parameter =} \StringTok{"w[1]"}\NormalTok{,       }\AttributeTok{value =} \FunctionTok{runif}\NormalTok{(n\_prior),         }\AttributeTok{source =} \StringTok{"prior"}\NormalTok{),}
  \FunctionTok{tibble}\NormalTok{(}\AttributeTok{parameter =} \StringTok{"bias"}\NormalTok{,       }\AttributeTok{value =} \FunctionTok{rbeta}\NormalTok{(n\_prior, }\DecValTok{1}\NormalTok{, }\DecValTok{1}\NormalTok{),   }\AttributeTok{source =} \StringTok{"prior"}\NormalTok{)}
\NormalTok{) }\SpecialCharTok{|\textgreater{}}
\NormalTok{  tidyr}\SpecialCharTok{::}\FunctionTok{crossing}\NormalTok{(}\AttributeTok{scenario =} \FunctionTok{c}\NormalTok{(}\StringTok{"static"}\NormalTok{, }\StringTok{"contingent\_shift"}\NormalTok{, }\StringTok{"drift"}\NormalTok{))}

\NormalTok{truth\_lines }\OtherTok{\textless{}{-}} \FunctionTok{tibble}\NormalTok{(}
  \AttributeTok{parameter =} \FunctionTok{c}\NormalTok{(}\StringTok{"log\_c"}\NormalTok{, }\StringTok{"log\_lambda"}\NormalTok{, }\StringTok{"w[1]"}\NormalTok{, }\StringTok{"bias"}\NormalTok{),}
  \AttributeTok{value     =} \FunctionTok{c}\NormalTok{(}\FunctionTok{log}\NormalTok{(}\FloatTok{2.0}\NormalTok{), }\FunctionTok{log}\NormalTok{(}\FloatTok{0.2}\NormalTok{), }\FloatTok{0.5}\NormalTok{, }\FloatTok{0.5}\NormalTok{)}
\NormalTok{)}

\FunctionTok{ggplot}\NormalTok{(}\FunctionTok{bind\_rows}\NormalTok{(prior\_ref, decay\_pp\_df }\SpecialCharTok{|\textgreater{}} 
                   \FunctionTok{mutate}\NormalTok{(}\AttributeTok{scenario =} \FunctionTok{factor}\NormalTok{(scenario, }\AttributeTok{levels =} \FunctionTok{c}\NormalTok{(}\StringTok{"static"}\NormalTok{, }\StringTok{"contingent\_shift"}\NormalTok{, }\StringTok{"drift"}\NormalTok{)))),}
       \FunctionTok{aes}\NormalTok{(}\AttributeTok{x =}\NormalTok{ value, }\AttributeTok{fill =}\NormalTok{ source, }\AttributeTok{color =}\NormalTok{ source)) }\SpecialCharTok{+}
  \FunctionTok{geom\_density}\NormalTok{(}\AttributeTok{alpha =} \FloatTok{0.4}\NormalTok{, }\AttributeTok{linewidth =} \FloatTok{0.5}\NormalTok{) }\SpecialCharTok{+}
  \FunctionTok{geom\_vline}\NormalTok{(}\AttributeTok{data =}\NormalTok{ truth\_lines, }\FunctionTok{aes}\NormalTok{(}\AttributeTok{xintercept =}\NormalTok{ value),}
             \AttributeTok{color =} \StringTok{"red"}\NormalTok{, }\AttributeTok{linetype =} \StringTok{"dashed"}\NormalTok{) }\SpecialCharTok{+}
  \FunctionTok{facet\_grid}\NormalTok{(scenario }\SpecialCharTok{\textasciitilde{}}\NormalTok{ parameter, }\AttributeTok{scales =} \StringTok{"free"}\NormalTok{) }\SpecialCharTok{+}
  \FunctionTok{scale\_fill\_manual}\NormalTok{(}\AttributeTok{values  =} \FunctionTok{c}\NormalTok{(}\AttributeTok{prior =} \StringTok{"grey70"}\NormalTok{, }\AttributeTok{posterior =} \StringTok{"\#009E73"}\NormalTok{)) }\SpecialCharTok{+}
  \FunctionTok{scale\_color\_manual}\NormalTok{(}\AttributeTok{values =} \FunctionTok{c}\NormalTok{(}\AttributeTok{prior =} \StringTok{"grey50"}\NormalTok{, }\AttributeTok{posterior =} \StringTok{"\#006D4E"}\NormalTok{)) }\SpecialCharTok{+}
  \FunctionTok{labs}\NormalTok{(}
    \AttributeTok{title    =} \StringTok{"Prior → Posterior Update: Decay GCM Across Scenarios"}\NormalTok{,}
    \AttributeTok{subtitle =} \StringTok{"Red dashed line = true generating value used for the single{-}agent fit"}\NormalTok{,}
    \AttributeTok{x =} \ConstantTok{NULL}\NormalTok{, }\AttributeTok{y =} \StringTok{"Density"}
\NormalTok{  )}
\end{Highlighting}
\end{Shaded}

\begin{center}
\includegraphics[width=0.8\linewidth,height=\textheight,keepaspectratio]{11-CategorizationModelsExemplars_files/figure-pdf/ch12_decay_prior_posterior-1.pdf}
\end{center}

\subsection{One-Shot Parameter Recovery Across
Scenarios}\label{one-shot-parameter-recovery-across-scenarios}

We sweep a 3 × 3 grid of \((c, \lambda)\) truths, with 5 independent
simulated datasets per cell per scenario (135 fits total), and report
recovery scatters.

\begin{Shaded}
\begin{Highlighting}[]
\NormalTok{decay\_recovery\_file }\OtherTok{\textless{}{-}} \FunctionTok{here}\NormalTok{(}\StringTok{"simdata"}\NormalTok{, }\StringTok{"ch12\_decay\_recovery\_results.rds"}\NormalTok{)}

\NormalTok{run\_decay\_recovery\_iter }\OtherTok{\textless{}{-}} \ControlFlowTok{function}\NormalTok{(scenario, c\_true, lambda\_true, iteration) \{}
\NormalTok{  subject\_seed }\OtherTok{\textless{}{-}}\NormalTok{ iteration }\SpecialCharTok{*} \DecValTok{10000} \SpecialCharTok{+}
                   \FunctionTok{round}\NormalTok{(c\_true }\SpecialCharTok{*} \DecValTok{100}\NormalTok{) }\SpecialCharTok{+}
                   \FunctionTok{round}\NormalTok{(lambda\_true }\SpecialCharTok{*} \DecValTok{1000}\NormalTok{) }\SpecialCharTok{+}
                   \ControlFlowTok{switch}\NormalTok{(scenario, }\AttributeTok{static =} \DecValTok{1}\NormalTok{, }\AttributeTok{drift =} \DecValTok{2}\NormalTok{, }\AttributeTok{contingent\_shift =} \DecValTok{3}\NormalTok{)}

\NormalTok{  agent\_df }\OtherTok{\textless{}{-}} \FunctionTok{simulate\_decay\_agent}\NormalTok{(}
    \AttributeTok{agent\_id     =}\NormalTok{ iteration,}
    \AttributeTok{scenario     =}\NormalTok{ scenario,}
    \AttributeTok{w            =} \FunctionTok{c}\NormalTok{(}\FloatTok{0.5}\NormalTok{, }\FloatTok{0.5}\NormalTok{),}
    \AttributeTok{c            =}\NormalTok{ c\_true,}
    \AttributeTok{lambda       =}\NormalTok{ lambda\_true,}
    \AttributeTok{bias         =} \FloatTok{0.5}\NormalTok{,}
    \AttributeTok{subject\_seed =}\NormalTok{ subject\_seed}
\NormalTok{  )}
\NormalTok{  dat }\OtherTok{\textless{}{-}} \FunctionTok{stan\_data\_from\_agent}\NormalTok{(agent\_df)}
\NormalTok{  pf  }\OtherTok{\textless{}{-}} \FunctionTok{tryCatch}\NormalTok{(}
\NormalTok{    mod\_gcm\_decay\_single}\SpecialCharTok{$}\FunctionTok{pathfinder}\NormalTok{(}\AttributeTok{data =}\NormalTok{ dat, }\AttributeTok{num\_paths =} \DecValTok{4}\NormalTok{,}
                                    \AttributeTok{refresh =} \DecValTok{0}\NormalTok{, }\AttributeTok{seed =}\NormalTok{ subject\_seed),}
    \AttributeTok{error =} \ControlFlowTok{function}\NormalTok{(e) }\ConstantTok{NULL}
\NormalTok{  )}
\NormalTok{  fit }\OtherTok{\textless{}{-}} \FunctionTok{tryCatch}\NormalTok{(}
\NormalTok{    mod\_gcm\_decay\_single}\SpecialCharTok{$}\FunctionTok{sample}\NormalTok{(}
      \AttributeTok{data            =}\NormalTok{ dat,}
      \AttributeTok{init            =} \ControlFlowTok{if}\NormalTok{ (}\SpecialCharTok{!}\FunctionTok{is.null}\NormalTok{(pf)) pf }\ControlFlowTok{else} \FloatTok{0.5}\NormalTok{,}
      \AttributeTok{seed            =}\NormalTok{ subject\_seed,}
      \AttributeTok{chains          =} \DecValTok{2}\NormalTok{,}
      \AttributeTok{parallel\_chains =} \DecValTok{2}\NormalTok{,}
      \AttributeTok{iter\_warmup     =} \DecValTok{800}\NormalTok{,}
      \AttributeTok{iter\_sampling   =} \DecValTok{1000}\NormalTok{,}
      \AttributeTok{refresh         =} \DecValTok{0}\NormalTok{,}
      \AttributeTok{adapt\_delta     =} \FloatTok{0.9}
\NormalTok{    ),}
    \AttributeTok{error =} \ControlFlowTok{function}\NormalTok{(e) \{ }\FunctionTok{message}\NormalTok{(}\StringTok{"Fit error: "}\NormalTok{, e}\SpecialCharTok{$}\NormalTok{message); }\ConstantTok{NULL}\NormalTok{ \}}
\NormalTok{  )}
  \ControlFlowTok{if}\NormalTok{ (}\FunctionTok{is.null}\NormalTok{(fit)) }\FunctionTok{return}\NormalTok{(}\ConstantTok{NULL}\NormalTok{)}

\NormalTok{  s }\OtherTok{\textless{}{-}}\NormalTok{ fit}\SpecialCharTok{$}\FunctionTok{summary}\NormalTok{(}\AttributeTok{variables =} \FunctionTok{c}\NormalTok{(}\StringTok{"log\_c"}\NormalTok{, }\StringTok{"log\_lambda"}\NormalTok{, }\StringTok{"c"}\NormalTok{, }\StringTok{"lambda"}\NormalTok{))}
\NormalTok{  diag }\OtherTok{\textless{}{-}}\NormalTok{ fit}\SpecialCharTok{$}\FunctionTok{diagnostic\_summary}\NormalTok{(}\AttributeTok{quiet =} \ConstantTok{TRUE}\NormalTok{)}
  \FunctionTok{tibble}\NormalTok{(}
    \AttributeTok{scenario    =}\NormalTok{ scenario,}
    \AttributeTok{c\_true      =}\NormalTok{ c\_true,}
    \AttributeTok{lambda\_true =}\NormalTok{ lambda\_true,}
    \AttributeTok{iteration   =}\NormalTok{ iteration,}
    \AttributeTok{log\_c\_mean       =}\NormalTok{ s}\SpecialCharTok{$}\NormalTok{mean[s}\SpecialCharTok{$}\NormalTok{variable }\SpecialCharTok{==} \StringTok{"log\_c"}\NormalTok{],}
    \AttributeTok{log\_lambda\_mean  =}\NormalTok{ s}\SpecialCharTok{$}\NormalTok{mean[s}\SpecialCharTok{$}\NormalTok{variable }\SpecialCharTok{==} \StringTok{"log\_lambda"}\NormalTok{],}
    \AttributeTok{c\_mean           =}\NormalTok{ s}\SpecialCharTok{$}\NormalTok{mean[s}\SpecialCharTok{$}\NormalTok{variable }\SpecialCharTok{==} \StringTok{"c"}\NormalTok{],}
    \AttributeTok{lambda\_mean      =}\NormalTok{ s}\SpecialCharTok{$}\NormalTok{mean[s}\SpecialCharTok{$}\NormalTok{variable }\SpecialCharTok{==} \StringTok{"lambda"}\NormalTok{],}
    \AttributeTok{log\_c\_q5         =}\NormalTok{ s}\SpecialCharTok{$}\NormalTok{q5[s}\SpecialCharTok{$}\NormalTok{variable }\SpecialCharTok{==} \StringTok{"log\_c"}\NormalTok{],}
    \AttributeTok{log\_c\_q95        =}\NormalTok{ s}\SpecialCharTok{$}\NormalTok{q95[s}\SpecialCharTok{$}\NormalTok{variable }\SpecialCharTok{==} \StringTok{"log\_c"}\NormalTok{],}
    \AttributeTok{log\_lambda\_q5    =}\NormalTok{ s}\SpecialCharTok{$}\NormalTok{q5[s}\SpecialCharTok{$}\NormalTok{variable }\SpecialCharTok{==} \StringTok{"log\_lambda"}\NormalTok{],}
    \AttributeTok{log\_lambda\_q95   =}\NormalTok{ s}\SpecialCharTok{$}\NormalTok{q95[s}\SpecialCharTok{$}\NormalTok{variable }\SpecialCharTok{==} \StringTok{"log\_lambda"}\NormalTok{],}
    \AttributeTok{had\_divergence   =} \FunctionTok{any}\NormalTok{(diag}\SpecialCharTok{$}\NormalTok{num\_divergent }\SpecialCharTok{\textgreater{}} \DecValTok{0}\NormalTok{)}
\NormalTok{  )}
\NormalTok{\}}

\ControlFlowTok{if}\NormalTok{ (regenerate\_simulations }\SpecialCharTok{||} \SpecialCharTok{!}\FunctionTok{file.exists}\NormalTok{(decay\_recovery\_file)) \{}
\NormalTok{  recovery\_plan }\OtherTok{\textless{}{-}} \FunctionTok{expand\_grid}\NormalTok{(}
    \AttributeTok{scenario    =} \FunctionTok{c}\NormalTok{(}\StringTok{"static"}\NormalTok{, }\StringTok{"contingent\_shift"}\NormalTok{, }\StringTok{"drift"}\NormalTok{),}
    \AttributeTok{c\_true      =} \FunctionTok{c}\NormalTok{(}\FloatTok{0.5}\NormalTok{, }\FloatTok{2.0}\NormalTok{, }\FloatTok{5.0}\NormalTok{),}
    \AttributeTok{lambda\_true =} \FunctionTok{c}\NormalTok{(}\FloatTok{0.02}\NormalTok{, }\FloatTok{0.2}\NormalTok{, }\FloatTok{0.6}\NormalTok{),}
    \AttributeTok{iteration   =} \DecValTok{1}\SpecialCharTok{:}\DecValTok{5}
\NormalTok{  )}

  \FunctionTok{plan}\NormalTok{(multisession, }\AttributeTok{workers =} \FunctionTok{max}\NormalTok{(}\DecValTok{1}\NormalTok{, }\FunctionTok{availableCores}\NormalTok{() }\SpecialCharTok{{-}} \DecValTok{1}\NormalTok{))}
\NormalTok{  decay\_recovery }\OtherTok{\textless{}{-}} \FunctionTok{future\_pmap\_dfr}\NormalTok{(}
    \FunctionTok{list}\NormalTok{(}
      \AttributeTok{scenario    =}\NormalTok{ recovery\_plan}\SpecialCharTok{$}\NormalTok{scenario,}
      \AttributeTok{c\_true      =}\NormalTok{ recovery\_plan}\SpecialCharTok{$}\NormalTok{c\_true,}
      \AttributeTok{lambda\_true =}\NormalTok{ recovery\_plan}\SpecialCharTok{$}\NormalTok{lambda\_true,}
      \AttributeTok{iteration   =}\NormalTok{ recovery\_plan}\SpecialCharTok{$}\NormalTok{iteration}
\NormalTok{    ),}
\NormalTok{    run\_decay\_recovery\_iter,}
    \AttributeTok{.options =} \FunctionTok{furrr\_options}\NormalTok{(}\AttributeTok{seed =} \ConstantTok{TRUE}\NormalTok{, }\AttributeTok{scheduling =} \DecValTok{1}\NormalTok{)}
\NormalTok{  )}
  \FunctionTok{plan}\NormalTok{(sequential)}
  \FunctionTok{saveRDS}\NormalTok{(decay\_recovery, decay\_recovery\_file)}
  \FunctionTok{cat}\NormalTok{(}\StringTok{"Decay recovery sweep complete.}\SpecialCharTok{\textbackslash{}n}\StringTok{"}\NormalTok{)}
\NormalTok{\} }\ControlFlowTok{else}\NormalTok{ \{}
\NormalTok{  decay\_recovery }\OtherTok{\textless{}{-}} \FunctionTok{readRDS}\NormalTok{(decay\_recovery\_file)}
  \FunctionTok{cat}\NormalTok{(}\StringTok{"Loaded existing decay recovery sweep.}\SpecialCharTok{\textbackslash{}n}\StringTok{"}\NormalTok{)}
\NormalTok{\}}
\end{Highlighting}
\end{Shaded}

\begin{verbatim}
Loaded existing decay recovery sweep.
\end{verbatim}

\begin{Shaded}
\begin{Highlighting}[]
\NormalTok{decay\_recovery\_long }\OtherTok{\textless{}{-}}\NormalTok{ decay\_recovery }\SpecialCharTok{|\textgreater{}}
\NormalTok{  dplyr}\SpecialCharTok{::}\FunctionTok{select}\NormalTok{(scenario, c\_true, lambda\_true, iteration, had\_divergence,}
\NormalTok{                log\_c\_mean, log\_lambda\_mean) }\SpecialCharTok{|\textgreater{}}
  \FunctionTok{pivot\_longer}\NormalTok{(}
    \AttributeTok{cols      =} \FunctionTok{c}\NormalTok{(log\_c\_mean, log\_lambda\_mean),}
    \AttributeTok{names\_to  =} \StringTok{"variable"}\NormalTok{,}
    \AttributeTok{values\_to =} \StringTok{"estimated"}
\NormalTok{  ) }\SpecialCharTok{|\textgreater{}}
  \FunctionTok{mutate}\NormalTok{(}
    \AttributeTok{parameter =} \FunctionTok{if\_else}\NormalTok{(variable }\SpecialCharTok{==} \StringTok{"log\_c\_mean"}\NormalTok{, }\StringTok{"log\_c"}\NormalTok{, }\StringTok{"log\_lambda"}\NormalTok{),}
    \AttributeTok{truth     =} \FunctionTok{if\_else}\NormalTok{(parameter }\SpecialCharTok{==} \StringTok{"log\_c"}\NormalTok{, }\FunctionTok{log}\NormalTok{(c\_true), }\FunctionTok{log}\NormalTok{(lambda\_true)),}
    \AttributeTok{scenario  =} \FunctionTok{factor}\NormalTok{(scenario, }\AttributeTok{levels =} \FunctionTok{c}\NormalTok{(}\StringTok{"static"}\NormalTok{, }\StringTok{"contingent\_shift"}\NormalTok{, }\StringTok{"drift"}\NormalTok{))}
\NormalTok{  )}

\FunctionTok{ggplot}\NormalTok{(decay\_recovery\_long, }\FunctionTok{aes}\NormalTok{(}\AttributeTok{x =}\NormalTok{ truth, }\AttributeTok{y =}\NormalTok{ estimated, }\AttributeTok{color =}\NormalTok{ had\_divergence)) }\SpecialCharTok{+}
  \FunctionTok{geom\_abline}\NormalTok{(}\AttributeTok{intercept =} \DecValTok{0}\NormalTok{, }\AttributeTok{slope =} \DecValTok{1}\NormalTok{, }\AttributeTok{linetype =} \StringTok{"dashed"}\NormalTok{, }\AttributeTok{color =} \StringTok{"grey50"}\NormalTok{) }\SpecialCharTok{+}
  \FunctionTok{geom\_point}\NormalTok{(}\AttributeTok{alpha =} \FloatTok{0.7}\NormalTok{, }\AttributeTok{size =} \DecValTok{2}\NormalTok{) }\SpecialCharTok{+}
  \FunctionTok{facet\_grid}\NormalTok{(parameter }\SpecialCharTok{\textasciitilde{}}\NormalTok{ scenario, }\AttributeTok{scales =} \StringTok{"free"}\NormalTok{) }\SpecialCharTok{+}
  \FunctionTok{scale\_color\_manual}\NormalTok{(}\AttributeTok{values =} \FunctionTok{c}\NormalTok{(}\StringTok{"FALSE"} \OtherTok{=} \StringTok{"\#0072B2"}\NormalTok{, }\StringTok{"TRUE"} \OtherTok{=} \StringTok{"\#D55E00"}\NormalTok{),}
                     \AttributeTok{name =} \StringTok{"Had divergence"}\NormalTok{) }\SpecialCharTok{+}
  \FunctionTok{labs}\NormalTok{(}
    \AttributeTok{title    =} \StringTok{"One{-}Shot Parameter Recovery for the Decay GCM"}\NormalTok{,}
    \AttributeTok{subtitle =} \StringTok{"Posterior mean vs true generating value; 5 iterations per (c, lambda) cell, 3 scenarios"}\NormalTok{,}
    \AttributeTok{x =} \StringTok{"True value (log scale)"}\NormalTok{, }\AttributeTok{y =} \StringTok{"Estimated value (posterior mean)"}
\NormalTok{  ) }\SpecialCharTok{+}
  \FunctionTok{theme}\NormalTok{(}\AttributeTok{legend.position =} \StringTok{"bottom"}\NormalTok{)}
\end{Highlighting}
\end{Shaded}

\begin{center}
\includegraphics[width=0.8\linewidth,height=\textheight,keepaspectratio]{11-CategorizationModelsExemplars_files/figure-pdf/ch12_decay_recovery_sweep-1.pdf}
\end{center}

\textbf{Interpretation.} The \texttt{log\_c} row is near-diagonal across
all three scenarios --- the sensitivity parameter is well-identified
regardless of task structure. The \texttt{log\_lambda} row is the
diagnostic: scatter under static data (the decay parameter is weakly
identified, so posterior means regress toward the prior mean), tighter
under contingent-shift (post-flip recovery dynamics directly constrain
\(\lambda\)), and well-recovered but more diffuse under smooth drift
(the assumption-matched scenario recovers \(\lambda\) at the nominal
calibration rate, but without the ``bonus'' constraint the abrupt
reversals provide). As the assumption match table above predicts, better
assumption match does not automatically mean stronger identifiability
--- the misspecified contingent-shift task exercises the forgetting
mechanism more sharply than the matched drift task does.

\subsection{Simulation-Based Calibration Checks Across
Scenarios}\label{simulation-based-calibration-checks-across-scenarios}

SBC is the calibration test proper: draw parameters from their priors,
simulate data with those parameters, fit the model, and check that the
posterior ranks of the true parameter are uniform. We run a separate SBC
sweep for each scenario.

\begin{Shaded}
\begin{Highlighting}[]
\NormalTok{make\_decay\_sbc\_generator }\OtherTok{\textless{}{-}} \ControlFlowTok{function}\NormalTok{(scenario) \{}
  \FunctionTok{SBC\_generator\_function}\NormalTok{(}
    \ControlFlowTok{function}\NormalTok{() \{}
\NormalTok{      log\_c\_true      }\OtherTok{\textless{}{-}} \FunctionTok{rnorm}\NormalTok{(}\DecValTok{1}\NormalTok{, LOG\_C\_PRIOR\_MEAN,      LOG\_C\_PRIOR\_SD)}
\NormalTok{      log\_lambda\_true }\OtherTok{\textless{}{-}} \FunctionTok{rnorm}\NormalTok{(}\DecValTok{1}\NormalTok{, LOG\_LAMBDA\_PRIOR\_MEAN, LOG\_LAMBDA\_PRIOR\_SD)}
\NormalTok{      c\_true          }\OtherTok{\textless{}{-}} \FunctionTok{exp}\NormalTok{(log\_c\_true)}
\NormalTok{      lambda\_true     }\OtherTok{\textless{}{-}} \FunctionTok{exp}\NormalTok{(log\_lambda\_true)}
\NormalTok{      w1\_true         }\OtherTok{\textless{}{-}} \FunctionTok{runif}\NormalTok{(}\DecValTok{1}\NormalTok{)}
\NormalTok{      bias\_true       }\OtherTok{\textless{}{-}} \FunctionTok{rbeta}\NormalTok{(}\DecValTok{1}\NormalTok{, }\DecValTok{1}\NormalTok{, }\DecValTok{1}\NormalTok{)}

\NormalTok{      subject\_seed }\OtherTok{\textless{}{-}} \FunctionTok{sample.int}\NormalTok{(}\FloatTok{1e8}\NormalTok{, }\DecValTok{1}\NormalTok{)}
\NormalTok{      agent\_df }\OtherTok{\textless{}{-}} \FunctionTok{simulate\_decay\_agent}\NormalTok{(}
        \AttributeTok{agent\_id     =} \DecValTok{1}\NormalTok{,}
        \AttributeTok{scenario     =}\NormalTok{ scenario,}
        \AttributeTok{w            =} \FunctionTok{c}\NormalTok{(w1\_true, }\DecValTok{1} \SpecialCharTok{{-}}\NormalTok{ w1\_true),}
        \AttributeTok{c            =}\NormalTok{ c\_true,}
        \AttributeTok{lambda       =}\NormalTok{ lambda\_true,}
        \AttributeTok{bias         =}\NormalTok{ bias\_true,}
        \AttributeTok{subject\_seed =}\NormalTok{ subject\_seed}
\NormalTok{      )}

      \FunctionTok{list}\NormalTok{(}
        \AttributeTok{variables =} \FunctionTok{list}\NormalTok{(}
          \AttributeTok{log\_c      =}\NormalTok{ log\_c\_true,}
          \AttributeTok{log\_lambda =}\NormalTok{ log\_lambda\_true,}
          \StringTok{\textasciigrave{}}\AttributeTok{w[1]}\StringTok{\textasciigrave{}}     \OtherTok{=}\NormalTok{ w1\_true,}
          \AttributeTok{bias       =}\NormalTok{ bias\_true}
\NormalTok{        ),}
        \AttributeTok{generated =} \FunctionTok{stan\_data\_from\_agent}\NormalTok{(agent\_df)}
\NormalTok{      )}
\NormalTok{    \}}
\NormalTok{  )}
\NormalTok{\}}

\NormalTok{run\_decay\_sbc }\OtherTok{\textless{}{-}} \ControlFlowTok{function}\NormalTok{(scenario, }\AttributeTok{n\_sims =} \DecValTok{200}\NormalTok{) \{}
\NormalTok{  cache\_file }\OtherTok{\textless{}{-}} \FunctionTok{here}\NormalTok{(}\StringTok{"simdata"}\NormalTok{, }\FunctionTok{paste0}\NormalTok{(}\StringTok{"ch12\_decay\_sbc\_"}\NormalTok{, scenario, }\StringTok{".rds"}\NormalTok{))}
  \ControlFlowTok{if}\NormalTok{ (}\SpecialCharTok{!}\NormalTok{regenerate\_simulations }\SpecialCharTok{\&\&} \FunctionTok{file.exists}\NormalTok{(cache\_file)) \{}
    \FunctionTok{cat}\NormalTok{(}\StringTok{"Loaded cached decay SBC for scenario:"}\NormalTok{, scenario, }\StringTok{"}\SpecialCharTok{\textbackslash{}n}\StringTok{"}\NormalTok{)}
    \FunctionTok{return}\NormalTok{(}\FunctionTok{readRDS}\NormalTok{(cache\_file))}
\NormalTok{  \}}

\NormalTok{  gen }\OtherTok{\textless{}{-}} \FunctionTok{make\_decay\_sbc\_generator}\NormalTok{(scenario)}
\NormalTok{  backend }\OtherTok{\textless{}{-}} \FunctionTok{SBC\_backend\_cmdstan\_sample}\NormalTok{(}
\NormalTok{    mod\_gcm\_decay\_single,}
    \AttributeTok{iter\_warmup   =} \DecValTok{800}\NormalTok{,}
    \AttributeTok{iter\_sampling =} \DecValTok{1000}\NormalTok{,}
    \AttributeTok{chains        =} \DecValTok{2}\NormalTok{,}
    \AttributeTok{adapt\_delta   =} \FloatTok{0.9}\NormalTok{,}
    \AttributeTok{refresh       =} \DecValTok{0}
\NormalTok{  )}
\NormalTok{  datasets }\OtherTok{\textless{}{-}} \FunctionTok{generate\_datasets}\NormalTok{(gen, }\AttributeTok{n\_sims =}\NormalTok{ n\_sims)}
\NormalTok{  res }\OtherTok{\textless{}{-}} \FunctionTok{compute\_SBC}\NormalTok{(}
\NormalTok{    datasets,}
\NormalTok{    backend,}
    \AttributeTok{cache\_mode     =} \StringTok{"results"}\NormalTok{,}
    \AttributeTok{cache\_location =} \FunctionTok{here}\NormalTok{(}\StringTok{"simdata"}\NormalTok{, }\FunctionTok{paste0}\NormalTok{(}\StringTok{"ch12\_decay\_sbc\_cache\_"}\NormalTok{, scenario)),}
    \AttributeTok{keep\_fits      =} \ConstantTok{FALSE}
\NormalTok{  )}
  \FunctionTok{saveRDS}\NormalTok{(res, cache\_file)}
\NormalTok{  res}
\NormalTok{\}}

\NormalTok{decay\_sbc\_static }\OtherTok{\textless{}{-}} \FunctionTok{run\_decay\_sbc}\NormalTok{(}\StringTok{"static"}\NormalTok{,           }\AttributeTok{n\_sims =} \DecValTok{200}\NormalTok{)}
\end{Highlighting}
\end{Shaded}

\begin{verbatim}
Loaded cached decay SBC for scenario: static 
\end{verbatim}

\begin{Shaded}
\begin{Highlighting}[]
\NormalTok{decay\_sbc\_drift  }\OtherTok{\textless{}{-}} \FunctionTok{run\_decay\_sbc}\NormalTok{(}\StringTok{"drift"}\NormalTok{,            }\AttributeTok{n\_sims =} \DecValTok{200}\NormalTok{)}
\end{Highlighting}
\end{Shaded}

\begin{verbatim}
Loaded cached decay SBC for scenario: drift 
\end{verbatim}

\begin{Shaded}
\begin{Highlighting}[]
\NormalTok{decay\_sbc\_shift  }\OtherTok{\textless{}{-}} \FunctionTok{run\_decay\_sbc}\NormalTok{(}\StringTok{"contingent\_shift"}\NormalTok{, }\AttributeTok{n\_sims =} \DecValTok{200}\NormalTok{)}
\end{Highlighting}
\end{Shaded}

\begin{verbatim}
Loaded cached decay SBC for scenario: contingent_shift 
\end{verbatim}

\begin{Shaded}
\begin{Highlighting}[]
\FunctionTok{print}\NormalTok{(}\FunctionTok{plot\_ecdf\_diff}\NormalTok{(decay\_sbc\_static) }\SpecialCharTok{+} \FunctionTok{ggtitle}\NormalTok{(}\StringTok{"SBC: static"}\NormalTok{))}
\end{Highlighting}
\end{Shaded}

\begin{center}
\includegraphics[width=0.8\linewidth,height=\textheight,keepaspectratio]{11-CategorizationModelsExemplars_files/figure-pdf/ch12_decay_sbc_ecdf-1.pdf}
\end{center}

\begin{Shaded}
\begin{Highlighting}[]
\FunctionTok{print}\NormalTok{(}\FunctionTok{plot\_ecdf\_diff}\NormalTok{(decay\_sbc\_drift)  }\SpecialCharTok{+} \FunctionTok{ggtitle}\NormalTok{(}\StringTok{"SBC: smooth drift"}\NormalTok{))}
\end{Highlighting}
\end{Shaded}

\begin{center}
\includegraphics[width=0.8\linewidth,height=\textheight,keepaspectratio]{11-CategorizationModelsExemplars_files/figure-pdf/ch12_decay_sbc_ecdf-2.pdf}
\end{center}

\begin{Shaded}
\begin{Highlighting}[]
\FunctionTok{print}\NormalTok{(}\FunctionTok{plot\_ecdf\_diff}\NormalTok{(decay\_sbc\_shift)  }\SpecialCharTok{+} \FunctionTok{ggtitle}\NormalTok{(}\StringTok{"SBC: performance{-}contingent shifts"}\NormalTok{))}
\end{Highlighting}
\end{Shaded}

\begin{center}
\includegraphics[width=0.8\linewidth,height=\textheight,keepaspectratio]{11-CategorizationModelsExemplars_files/figure-pdf/ch12_decay_sbc_ecdf-3.pdf}
\end{center}

\begin{Shaded}
\begin{Highlighting}[]
\FunctionTok{print}\NormalTok{(}\FunctionTok{plot\_rank\_hist}\NormalTok{(decay\_sbc\_static) }\SpecialCharTok{+} \FunctionTok{ggtitle}\NormalTok{(}\StringTok{"Rank histograms — static"}\NormalTok{))}
\end{Highlighting}
\end{Shaded}

\begin{center}
\includegraphics[width=0.8\linewidth,height=\textheight,keepaspectratio]{11-CategorizationModelsExemplars_files/figure-pdf/ch12_decay_sbc_ecdf-4.pdf}
\end{center}

\begin{Shaded}
\begin{Highlighting}[]
\FunctionTok{print}\NormalTok{(}\FunctionTok{plot\_rank\_hist}\NormalTok{(decay\_sbc\_drift)  }\SpecialCharTok{+} \FunctionTok{ggtitle}\NormalTok{(}\StringTok{"Rank histograms — drift"}\NormalTok{))}
\end{Highlighting}
\end{Shaded}

\begin{center}
\includegraphics[width=0.8\linewidth,height=\textheight,keepaspectratio]{11-CategorizationModelsExemplars_files/figure-pdf/ch12_decay_sbc_ecdf-5.pdf}
\end{center}

\begin{Shaded}
\begin{Highlighting}[]
\FunctionTok{print}\NormalTok{(}\FunctionTok{plot\_rank\_hist}\NormalTok{(decay\_sbc\_shift)  }\SpecialCharTok{+} \FunctionTok{ggtitle}\NormalTok{(}\StringTok{"Rank histograms — contingent shifts"}\NormalTok{))}
\end{Highlighting}
\end{Shaded}

\begin{center}
\includegraphics[width=0.8\linewidth,height=\textheight,keepaspectratio]{11-CategorizationModelsExemplars_files/figure-pdf/ch12_decay_sbc_ecdf-6.pdf}
\end{center}

\subsection{Parameter Recovery Across SBC
Simulations}\label{parameter-recovery-across-sbc-simulations}

\begin{Shaded}
\begin{Highlighting}[]
\NormalTok{combine\_sbc\_stats }\OtherTok{\textless{}{-}} \ControlFlowTok{function}\NormalTok{(sbc\_obj, scenario) \{}
\NormalTok{  sbc\_obj}\SpecialCharTok{$}\NormalTok{stats }\SpecialCharTok{|\textgreater{}} \FunctionTok{mutate}\NormalTok{(}\AttributeTok{scenario =}\NormalTok{ scenario)}
\NormalTok{\}}

\NormalTok{decay\_sbc\_stats }\OtherTok{\textless{}{-}} \FunctionTok{bind\_rows}\NormalTok{(}
  \FunctionTok{combine\_sbc\_stats}\NormalTok{(decay\_sbc\_static, }\StringTok{"static"}\NormalTok{),}
  \FunctionTok{combine\_sbc\_stats}\NormalTok{(decay\_sbc\_drift,  }\StringTok{"drift"}\NormalTok{),}
  \FunctionTok{combine\_sbc\_stats}\NormalTok{(decay\_sbc\_shift,  }\StringTok{"contingent\_shift"}\NormalTok{)}
\NormalTok{) }\SpecialCharTok{|\textgreater{}}
\NormalTok{  dplyr}\SpecialCharTok{::}\FunctionTok{filter}\NormalTok{(variable }\SpecialCharTok{\%in\%} \FunctionTok{c}\NormalTok{(}\StringTok{"log\_c"}\NormalTok{, }\StringTok{"log\_lambda"}\NormalTok{, }\StringTok{"w[1]"}\NormalTok{, }\StringTok{"bias"}\NormalTok{)) }\SpecialCharTok{|\textgreater{}}
  \FunctionTok{mutate}\NormalTok{(}\AttributeTok{scenario =} \FunctionTok{factor}\NormalTok{(scenario, }\AttributeTok{levels =} \FunctionTok{c}\NormalTok{(}\StringTok{"static"}\NormalTok{, }\StringTok{"contingent\_shift"}\NormalTok{, }\StringTok{"drift"}\NormalTok{)))}

\FunctionTok{ggplot}\NormalTok{(decay\_sbc\_stats, }\FunctionTok{aes}\NormalTok{(}\AttributeTok{x =}\NormalTok{ simulated\_value, }\AttributeTok{y =}\NormalTok{ mean)) }\SpecialCharTok{+}
  \FunctionTok{geom\_point}\NormalTok{(}\AttributeTok{alpha =} \FloatTok{0.35}\NormalTok{, }\AttributeTok{color =} \StringTok{"\#0072B2"}\NormalTok{) }\SpecialCharTok{+}
  \FunctionTok{geom\_abline}\NormalTok{(}\AttributeTok{intercept =} \DecValTok{0}\NormalTok{, }\AttributeTok{slope =} \DecValTok{1}\NormalTok{, }\AttributeTok{color =} \StringTok{"red"}\NormalTok{, }\AttributeTok{linetype =} \StringTok{"dashed"}\NormalTok{) }\SpecialCharTok{+}
  \FunctionTok{facet\_grid}\NormalTok{(variable }\SpecialCharTok{\textasciitilde{}}\NormalTok{ scenario, }\AttributeTok{scales =} \StringTok{"free"}\NormalTok{) }\SpecialCharTok{+}
  \FunctionTok{labs}\NormalTok{(}
    \AttributeTok{title    =} \StringTok{"SBC{-}Level Parameter Recovery: Decay GCM"}\NormalTok{,}
    \AttributeTok{subtitle =} \StringTok{"Posterior mean vs true generating value, pooled across SBC simulations"}\NormalTok{,}
    \AttributeTok{x =} \StringTok{"True value"}\NormalTok{, }\AttributeTok{y =} \StringTok{"Posterior mean"}
\NormalTok{  )}
\end{Highlighting}
\end{Shaded}

\begin{center}
\includegraphics[width=0.8\linewidth,height=\textheight,keepaspectratio]{11-CategorizationModelsExemplars_files/figure-pdf/ch12_decay_sbc_recovery_plots-1.pdf}
\end{center}

\begin{Shaded}
\begin{Highlighting}[]
\NormalTok{decay\_sbc\_stats }\SpecialCharTok{|\textgreater{}}
  \FunctionTok{group\_by}\NormalTok{(scenario, variable) }\SpecialCharTok{|\textgreater{}}
  \FunctionTok{summarise}\NormalTok{(}
    \AttributeTok{sbc\_mean\_bias =} \FunctionTok{mean}\NormalTok{(mean }\SpecialCharTok{{-}}\NormalTok{ simulated\_value, }\AttributeTok{na.rm =} \ConstantTok{TRUE}\NormalTok{),}
    \AttributeTok{sbc\_rmse      =} \FunctionTok{sqrt}\NormalTok{(}\FunctionTok{mean}\NormalTok{((mean }\SpecialCharTok{{-}}\NormalTok{ simulated\_value)}\SpecialCharTok{\^{}}\DecValTok{2}\NormalTok{, }\AttributeTok{na.rm =} \ConstantTok{TRUE}\NormalTok{)),}
    \AttributeTok{sbc\_corr      =} \FunctionTok{cor}\NormalTok{(mean, simulated\_value,}
                        \AttributeTok{use =} \StringTok{"complete.obs"}\NormalTok{, }\AttributeTok{method =} \StringTok{"pearson"}\NormalTok{),}
    \AttributeTok{n\_sims        =}\NormalTok{ dplyr}\SpecialCharTok{::}\FunctionTok{n}\NormalTok{(),}
    \AttributeTok{.groups       =} \StringTok{"drop"}
\NormalTok{  ) }\SpecialCharTok{|\textgreater{}}
  \FunctionTok{print}\NormalTok{(}\AttributeTok{n =} \ConstantTok{Inf}\NormalTok{)}
\end{Highlighting}
\end{Shaded}

\begin{verbatim}
# A tibble: 12 x 6
   scenario         variable   sbc_mean_bias sbc_rmse sbc_corr n_sims
   <fct>            <chr>              <dbl>    <dbl>    <dbl>  <int>
 1 static           bias            0.00167    0.0835    0.957    200
 2 static           log_c          -0.00810    0.394     0.912    200
 3 static           log_lambda      0.00307    0.893     0.463    200
 4 static           w[1]            0.00303    0.164     0.846    200
 5 contingent_shift bias           -0.000680   0.0650    0.973    200
 6 contingent_shift log_c          -0.00342    0.418     0.902    200
 7 contingent_shift log_lambda      0.0122     0.726     0.687    200
 8 contingent_shift w[1]           -0.0125     0.146     0.868    200
 9 drift            bias            0.00223    0.0853    0.956    200
10 drift            log_c          -0.0830     0.427     0.909    200
11 drift            log_lambda     -0.00799    0.842     0.581    200
12 drift            w[1]           -0.0111     0.160     0.856    200
\end{verbatim}

\subsection{Discussion: Task Structure Governs
Identifiability}\label{discussion-task-structure-governs-identifiability}

Under static data, the \texttt{log\_lambda} parameter cannot really be
identified --- the posterior does not contain the truth at the nominal
rate because the data carry almost no information about forgetting. The
\texttt{log\_c} curve is well-behaved on the same data: the sensitivity
parameter is identified regardless. Under contingent-shift and drift
data, both parameters recover well, with contingent-shift providing the
sharpest \(\lambda\) constraint (abrupt reversals) and smooth drift
providing the best-matched but less forceful identifiability.

This is the identifiability argument as an empirical claim rather than
an algebraic one. The decay GCM's \emph{parameter} is always the same.
Its \emph{identifiability} is a joint property of the model and the
task: when the task exercises forgetting --- whether through a drifting
boundary or through performance-contingent reversals --- the model
recovers \(\lambda\) cleanly; when it does not, \(\lambda\) is
indistinguishable from a slightly wider similarity kernel.

A secondary observation emerges when we look closer at the calibration
checks: while parameter recovery (the scatter plot) looks excellent for
\texttt{log\_c} across all scenarios, the SBC rank histograms reveal a
subtle ``worry'' for \texttt{log\_c} specifically in the drift scenario.
This highlights a fundamental distinction in Bayesian validation:

\begin{itemize}
\tightlist
\item
  \textbf{Point Recovery (The ``What''):} This asks if the posterior
  \emph{mean} is close to the truth. If your recovery plot is diagonal,
  your model is right ``on average.''
\item
  \textbf{Coverage Calibration (The ``How Sure''):} This asks if the
  posterior \emph{width} (the uncertainty) is honest. If the model is
  90\% confident, is it actually right 90\% of the time?
\end{itemize}

SBC is designed to catch the latter. In the drift scenario, the model is
right about the value of \texttt{log\_c} (good point recovery), but it
is slightly overconfident---it underestimates its own uncertainty. One
might wonder: \textbf{since the model has \(\lambda\), why doesn't the
decay parameter just ``map'' the drift and leave \(c\) alone?}

The answer lies in the \textbf{Smeared Memory} problem. Even with an
optimal \(\lambda\), the agent doesn't categorise based on only the
single last trial; they need a ``window'' of recent exemplars to be
effective. Because the boundary is drifting \emph{during} that window,
the exemplars in memory become spatially ``smeared'' --- a location that
was Category 0 five trials ago might be Category 1 now.

While \(\lambda\) can tell the model to trust recent trials more, it is
a blunt temporal instrument. It cannot resolve the spatial conflict that
arises when two different category labels occupy the same coordinates in
the recent memory window. The model's likelihood, having no parameter
for ``environmental volatility,'' is forced to explain this residual
friction using spatial sensitivity (\(c\)). It performs a form of
\textbf{attribution error}, interpreting this temporal volatility as
spatial overlap and settling on an overly precise estimate of \(c\).
Because the model is using the wrong ``tool'' to explain the noise, it
converges on a narrow posterior that doesn't cover the true value as
often as it should. If we only looked at the recovery plot, we would
miss this overconfidence entirely. This is why SBC is the ``gold
standard'' check: it ensures the model isn't just accurate, but also
appropriately humble about what it knows.

\textbf{How the three scenarios match the model's assumptions.} The
decay GCM assumes exemplars decay exponentially as a function of trial
age, with no separate mechanism for abrupt reversals or directed drift.
The \emph{static} scenario violates nothing but offers no signal for
\(\lambda\). The \emph{contingent-shift} scenario is a genuine
misspecification --- the true data-generating process has abrupt
reversals that no smooth decay can exactly reproduce --- yet it provides
the sharpest \(\lambda\) signal because post-flip recovery speed depends
specifically on how fast old, now-wrong exemplars are discounted. The
posterior predictive checks at scenario transitions (not plotted here)
should reveal systematic underperformance right after each flip; that
residual misfit is the handle for a future extension with an explicit
change-point mechanism. The \emph{smooth-drift} scenario is the
assumption-matched case: exponentially downweighting past exemplars is a
defensible approximation to ``evidence collected under a different
category mean is stale.'' SBC passes cleanly here, but \(\lambda\)
recovery is less sharp than under contingent-shift because drift does
not impose the stark temporal discontinuities that most strongly
constrain the forgetting rate.

\subsection{\texorpdfstring{The \(c\)/\(\lambda\) Joint Posterior: Ridge
Geometry and Task
Resolution}{The c/\textbackslash lambda Joint Posterior: Ridge Geometry and Task Resolution}}\label{the-clambda-joint-posterior-ridge-geometry-and-task-resolution}

The pair plots from the single-agent scenario fits revealed a
\texttt{log\_c}/\texttt{log\_lambda} ridge on static data. Understanding
the geometry of this ridge --- and why task structure resolves it --- is
worth making explicit, because Chapter 12 faces a structurally similar
problem (\(q\)/\(r\) in the Kalman filter) and resolves it differently.

\textbf{Why does the ridge form?} Both \(c\) and \(\lambda\) narrow the
effective contribution of any given exemplar to the category mean:

\begin{itemize}
\tightlist
\item
  A high \(c\) shrinks contributions from spatially distant exemplars
  (similarity decays faster with distance).
\item
  A high \(\lambda\) shrinks contributions from temporally distant
  exemplars (similarity decays faster with trial age).
\end{itemize}

Under a stationary task with balanced categories, distance and age are
nearly uncorrelated --- older exemplars are not systematically farther
from the current stimulus. This means that increasing \(c\) and
decreasing \(\lambda\) (or vice versa) can produce nearly the same
observed choice probabilities, creating the elongated posterior ridge
you see in the pair plots.

\textbf{Why is the GCM ridge less severe than the Kalman filter's
\(q\)/\(r\) ridge?} The two parameters operate on orthogonal dimensions:
\(c\) operates on the spatial axis (feature distances), \(\lambda\) on
the temporal axis (trial age). Even under static data, choice
probabilities have some leverage on each independently --- \(c\) is
constrained by how performance differs across stimuli at varying spatial
distances, while \(\lambda\) is partially constrained by the
early-learning transient (the model must form initial memories before it
can categorize, and the learning rate is \(\lambda\)-dependent). In the
Kalman filter, \(q\) and \(r\) both scale the \emph{same} likelihood at
decision time (the multivariate normal variance at trial \(t\) is a
function of both), so their directions are more nearly collinear and the
ridge is typically more severe. Chapter 12 introduces an explicit
reparameterization (the SNR parameter \(\rho = q/r\)) to address this;
the GCM does not need it here.

\textbf{Task structure as the primary handle.} The SBC results above
confirm that the contingent-shift scenario cleanly resolves \(\lambda\)
without any reparameterization. The reversal dynamics create a temporal
structure that spatial sensitivity cannot account for: the \emph{time}
at which performance drops is determined by \(\lambda\) (how fast old,
now-wrong exemplars are forgotten) in a way that is independent of \(c\)
(which only governs how much nearby vs.~distant exemplars are weighted).
This is the key empirical signature of a parameter that is theoretically
identifiable but under-constrained by a particular task design: swap in
a task that exercises the mechanism, and identifiability follows.

\textbf{When would reparameterization be needed for the GCM?} If you
needed to run the decay GCM exclusively on static data, the ridge would
be a genuine sampling problem. In that case, a reparameterization could
help: write \(\log s = \log c + \log \lambda\) (the overall ``forgetting
rate'' in combined spatial-temporal units) and
\(\log d = \log c - \log \lambda\) (the relative weighting of spatial
vs.~temporal decay). The joint posterior of \((\log s, \log d)\) would
have a narrower marginal in the \(\log s\) direction (identified by
overall accuracy) and a wider marginal in the \(\log d\) direction (the
degenerate direction).

\begin{center}\rule{0.5\linewidth}{0.5pt}\end{center}

\section{The Multilevel Decay GCM: Scaling to populations and changing
environments}\label{the-multilevel-decay-gcm-scaling-to-populations-and-changing-environments}

Having demonstrated that memory decay is essential for non-stationary
environments and that its parameters are identifiable given the right
task structure, we now combine these mechanisms into a full hierarchical
model. The \textbf{Multilevel Decay GCM} allows us to estimate
individual cognitive differences while sharing information across the
population to constrain our estimates. This is our ``Final Boss'' model:
population-aware, psychologically realistic, and robust to environmental
change.

\subsection{Mathematical Formulation: The Hierarchical
Architecture}\label{mathematical-formulation-the-hierarchical-architecture}

This model uses log-normal hierarchies for sensitivity and forgetting,
and logit-normal hierarchies for attention and response bias. With two
features, the simplex-valued attention vector \(\mathbf{w}_j\) is fully
determined by \(w_{j,1}\), so the natural simplex hierarchy (Dirichlet
around a population mean \(\bar{\mathbf{w}}\) with concentration
\(\kappa\)) collapses to a Beta on \(w_{j,1}\). We use a
\textbf{logit-normal} in its place --- equivalent in support and shape
to the Beta, but its non-centered parameterisation samples much more
cleanly than a Dirichlet hierarchy in NUTS, which is what
\texttt{ch12\_gcm\_ml\_decay.stan} implements.

\textbf{Population-level distributions (Priors):}

\[
\begin{aligned}
\bar{\mu}_{\text{logit}\,w_1}, \bar{\sigma}_{\text{logit}\,w_1} &\sim \text{Priors} \\
\bar{\mu}_{\log c}, \bar{\sigma}_{\log c} &\sim \text{Priors} \\
\bar{\mu}_{\log \lambda}, \bar{\sigma}_{\log \lambda} &\sim \text{Priors} \\
\bar{\mu}_{\text{logit}\,\beta}, \bar{\sigma}_{\text{logit}\,\beta} &\sim \text{Priors}
\end{aligned}
\]

\textbf{Individual-level parameters (for subject \(j\), with
\(w_{j,2} = 1 - w_{j,1}\)):}

\[
\begin{aligned}
\text{logit}(w_{j,1}) &\sim \text{Normal}(\bar{\mu}_{\text{logit}\,w_1}, \bar{\sigma}_{\text{logit}\,w_1}) \\
\log c_j &\sim \text{Normal}(\bar{\mu}_{\log c}, \bar{\sigma}_{\log c}) \\
\log \lambda_j &\sim \text{Normal}(\bar{\mu}_{\log \lambda}, \bar{\sigma}_{\log \lambda}) \\
\text{logit}(\beta_j) &\sim \text{Normal}(\bar{\mu}_{\text{logit}\,\beta}, \bar{\sigma}_{\text{logit}\,\beta})
\end{aligned}
\]

Each individual-level scalar uses a non-centered offset
(\texttt{z\ \textasciitilde{}\ N(0,1)}, with the actual parameter equal
to \texttt{pop\_mean\ +\ pop\_sd\ *\ z}) to avoid the funnel that direct
centred parameterisations produce when subject-level data are weakly
informative.

\textbf{Trial-level Generative Process (for trial \(i\) of subject
\(j\)):}

\[
\begin{aligned}
d_{iej} &= \sum_f w_{jf} |x_{if} - x_{ef}| \\
\eta_{iej} &= e^{-\lambda_j(i - t_e)} \cdot e^{-c_j d_{iej}} \\
m_{k,ij} &= \frac{\sum_{e \in C_k} \eta_{iej}}{\sum_{e \in C_k} e^{-\lambda_j(i - t_e)}} \\
p_{ij} &= \frac{\beta_j \cdot m_{1,ij}}{\beta_j \cdot m_{1,ij} + (1 - \beta_j) \cdot m_{0,ij}} \\
y_{ij} &\sim \text{Bernoulli}(p_{ij})
\end{aligned}
\]

\begin{itemize}
\tightlist
\item
  \textbf{Shrinkage:} The model uses the population means to ``pull''
  uncertain individual estimates toward the group average. This is
  especially useful for \(\lambda_j\) in scenarios where an individual's
  data might be locally uninformative.
\item
  \textbf{Effective Sample Size:} The denominator of \(m_{k,ij}\) is now
  the \textbf{decay-weighted effective number of exemplars}, ensuring
  that old, forgotten evidence doesn't inflate the mean similarity.
\end{itemize}

\subsection{Simulating Data for the Multilevel Decay
GCM}\label{simulating-data-for-the-multilevel-decay-gcm}

We simulate three small populations of agents --- one per scenario
(static, smooth drift, performance-contingent shifts) --- using the
\textbf{same} population-level parameters across scenarios.
Subject-level parameters are drawn around the population means using the
logit-normal/log-normal hierarchies that match the Stan model exactly.
Because the simulator and the inference model now share a
parameterisation, recovery scatters can be read directly without
translating between Dirichlet concentration and logit SD.

\begin{Shaded}
\begin{Highlighting}[]
\NormalTok{simulate\_multilevel\_decay\_gcm }\OtherTok{\textless{}{-}} \ControlFlowTok{function}\NormalTok{(n\_subjects, stimulus\_info, n\_blocks,}
\NormalTok{                                          scenario,}
\NormalTok{                                          base\_seed,}
\NormalTok{                                          pop\_logit\_w\_mean,    pop\_logit\_w\_sd,}
\NormalTok{                                          pop\_log\_c\_mean,      pop\_log\_c\_sd,}
\NormalTok{                                          pop\_log\_lambda\_mean, pop\_log\_lambda\_sd,}
\NormalTok{                                          pop\_logit\_bias\_mean, pop\_logit\_bias\_sd) \{}
  \FunctionTok{set.seed}\NormalTok{(base\_seed)}
\NormalTok{  ind\_logit\_w    }\OtherTok{\textless{}{-}} \FunctionTok{rnorm}\NormalTok{(n\_subjects, pop\_logit\_w\_mean,    pop\_logit\_w\_sd)}
\NormalTok{  ind\_w1         }\OtherTok{\textless{}{-}} \FunctionTok{plogis}\NormalTok{(ind\_logit\_w)}
\NormalTok{  ind\_log\_c      }\OtherTok{\textless{}{-}} \FunctionTok{rnorm}\NormalTok{(n\_subjects, pop\_log\_c\_mean,      pop\_log\_c\_sd)}
\NormalTok{  ind\_log\_lambda }\OtherTok{\textless{}{-}} \FunctionTok{rnorm}\NormalTok{(n\_subjects, pop\_log\_lambda\_mean, pop\_log\_lambda\_sd)}
\NormalTok{  ind\_logit\_bias }\OtherTok{\textless{}{-}} \FunctionTok{rnorm}\NormalTok{(n\_subjects, pop\_logit\_bias\_mean, pop\_logit\_bias\_sd)}

\NormalTok{  true\_params }\OtherTok{\textless{}{-}} \FunctionTok{tibble}\NormalTok{(}
    \AttributeTok{agent\_id        =} \FunctionTok{seq\_len}\NormalTok{(n\_subjects),}
    \AttributeTok{scenario        =}\NormalTok{ scenario,}
    \AttributeTok{logit\_w\_true    =}\NormalTok{ ind\_logit\_w,}
    \AttributeTok{w1\_true         =}\NormalTok{ ind\_w1,}
    \AttributeTok{w2\_true         =} \DecValTok{1} \SpecialCharTok{{-}}\NormalTok{ ind\_w1,}
    \AttributeTok{log\_c\_true      =}\NormalTok{ ind\_log\_c,}
    \AttributeTok{c\_true          =} \FunctionTok{exp}\NormalTok{(ind\_log\_c),}
    \AttributeTok{log\_lambda\_true =}\NormalTok{ ind\_log\_lambda,}
    \AttributeTok{lambda\_true     =} \FunctionTok{exp}\NormalTok{(ind\_log\_lambda),}
    \AttributeTok{logit\_bias\_true =}\NormalTok{ ind\_logit\_bias,}
    \AttributeTok{bias\_true       =} \FunctionTok{plogis}\NormalTok{(ind\_logit\_bias)}
\NormalTok{  )}

\NormalTok{  all\_agent\_data }\OtherTok{\textless{}{-}} \FunctionTok{vector}\NormalTok{(}\StringTok{"list"}\NormalTok{, n\_subjects)}
  \ControlFlowTok{for}\NormalTok{ (j }\ControlFlowTok{in} \FunctionTok{seq\_len}\NormalTok{(n\_subjects)) \{}
\NormalTok{    subject\_seed }\OtherTok{\textless{}{-}}\NormalTok{ base\_seed }\SpecialCharTok{*} \DecValTok{1000} \SpecialCharTok{+}\NormalTok{ j}
\NormalTok{    all\_agent\_data[[j]] }\OtherTok{\textless{}{-}} \FunctionTok{simulate\_decay\_agent}\NormalTok{(}
      \AttributeTok{agent\_id     =}\NormalTok{ j,}
      \AttributeTok{scenario     =}\NormalTok{ scenario,}
      \AttributeTok{w            =} \FunctionTok{c}\NormalTok{(ind\_w1[j], }\DecValTok{1} \SpecialCharTok{{-}}\NormalTok{ ind\_w1[j]),}
      \AttributeTok{c            =} \FunctionTok{exp}\NormalTok{(ind\_log\_c[j]),}
      \AttributeTok{lambda       =} \FunctionTok{exp}\NormalTok{(ind\_log\_lambda[j]),}
      \AttributeTok{bias         =} \FunctionTok{plogis}\NormalTok{(ind\_logit\_bias[j]),}
      \AttributeTok{subject\_seed =}\NormalTok{ subject\_seed}
\NormalTok{    )}
\NormalTok{  \}}

  \FunctionTok{list}\NormalTok{(}
    \AttributeTok{data        =} \FunctionTok{bind\_rows}\NormalTok{(all\_agent\_data),}
    \AttributeTok{true\_params =}\NormalTok{ true\_params}
\NormalTok{  )}
\NormalTok{\}}

\NormalTok{n\_subjects\_sim }\OtherTok{\textless{}{-}} \DecValTok{10}
\NormalTok{pop\_params\_sim }\OtherTok{\textless{}{-}} \FunctionTok{list}\NormalTok{(}
  \AttributeTok{pop\_logit\_w\_mean    =} \DecValTok{0}\NormalTok{,        }\CommentTok{\# median w1 ≈ 0.5}
  \AttributeTok{pop\_logit\_w\_sd      =} \FloatTok{0.7}\NormalTok{,      }\CommentTok{\# ±2SD covers w1 ∈ (0.20, 0.80)}
  \AttributeTok{pop\_log\_c\_mean      =} \FunctionTok{log}\NormalTok{(}\FloatTok{2.0}\NormalTok{), }\CommentTok{\# population median c = 2}
  \AttributeTok{pop\_log\_c\_sd        =} \FloatTok{0.4}\NormalTok{,}
  \AttributeTok{pop\_log\_lambda\_mean =} \FunctionTok{log}\NormalTok{(}\FloatTok{0.2}\NormalTok{), }\CommentTok{\# population median λ = 0.2}
  \AttributeTok{pop\_log\_lambda\_sd   =} \FloatTok{0.5}\NormalTok{,}
  \AttributeTok{pop\_logit\_bias\_mean =} \DecValTok{0}\NormalTok{,}
  \AttributeTok{pop\_logit\_bias\_sd   =} \FloatTok{0.3}
\NormalTok{)}

\NormalTok{ml\_scenarios }\OtherTok{\textless{}{-}} \FunctionTok{c}\NormalTok{(}\StringTok{"static"}\NormalTok{, }\StringTok{"contingent\_shift"}\NormalTok{, }\StringTok{"drift"}\NormalTok{)}

\NormalTok{ml\_decay\_files }\OtherTok{\textless{}{-}} \FunctionTok{map}\NormalTok{(ml\_scenarios, }\ControlFlowTok{function}\NormalTok{(scen) }\FunctionTok{list}\NormalTok{(}
  \AttributeTok{data  =} \FunctionTok{here}\NormalTok{(}\StringTok{"simdata"}\NormalTok{, }\FunctionTok{paste0}\NormalTok{(}\StringTok{"ch12\_gcm\_ml\_decay\_data\_"}\NormalTok{,  scen, }\StringTok{".csv"}\NormalTok{)),}
  \AttributeTok{truth =} \FunctionTok{here}\NormalTok{(}\StringTok{"simdata"}\NormalTok{, }\FunctionTok{paste0}\NormalTok{(}\StringTok{"ch12\_gcm\_ml\_decay\_truth\_"}\NormalTok{, scen, }\StringTok{".csv"}\NormalTok{))}
\NormalTok{)) }\SpecialCharTok{|\textgreater{}} \FunctionTok{setNames}\NormalTok{(ml\_scenarios)}

\NormalTok{ml\_files\_exist }\OtherTok{\textless{}{-}} \FunctionTok{all}\NormalTok{(}\FunctionTok{file.exists}\NormalTok{(}\FunctionTok{unlist}\NormalTok{(}\FunctionTok{map}\NormalTok{(ml\_decay\_files, }\SpecialCharTok{\textasciitilde{}} \FunctionTok{c}\NormalTok{(.x}\SpecialCharTok{$}\NormalTok{data, .x}\SpecialCharTok{$}\NormalTok{truth)))))}

\ControlFlowTok{if}\NormalTok{ (regenerate\_simulations }\SpecialCharTok{||} \SpecialCharTok{!}\NormalTok{ml\_files\_exist) \{}
\NormalTok{  ml\_decay\_sims }\OtherTok{\textless{}{-}} \FunctionTok{map}\NormalTok{(ml\_scenarios, }\ControlFlowTok{function}\NormalTok{(scen) \{}
\NormalTok{    out }\OtherTok{\textless{}{-}} \FunctionTok{simulate\_multilevel\_decay\_gcm}\NormalTok{(}
      \AttributeTok{n\_subjects          =}\NormalTok{ n\_subjects\_sim,}
      \AttributeTok{stimulus\_info       =}\NormalTok{ stimulus\_info,}
      \AttributeTok{n\_blocks            =}\NormalTok{ n\_blocks,}
      \AttributeTok{scenario            =}\NormalTok{ scen,}
      \AttributeTok{base\_seed           =} \FunctionTok{match}\NormalTok{(scen, ml\_scenarios) }\SpecialCharTok{*} \DecValTok{1000} \SpecialCharTok{+} \DecValTok{42}\NormalTok{,}
      \AttributeTok{pop\_logit\_w\_mean    =}\NormalTok{ pop\_params\_sim}\SpecialCharTok{$}\NormalTok{pop\_logit\_w\_mean,}
      \AttributeTok{pop\_logit\_w\_sd      =}\NormalTok{ pop\_params\_sim}\SpecialCharTok{$}\NormalTok{pop\_logit\_w\_sd,}
      \AttributeTok{pop\_log\_c\_mean      =}\NormalTok{ pop\_params\_sim}\SpecialCharTok{$}\NormalTok{pop\_log\_c\_mean,}
      \AttributeTok{pop\_log\_c\_sd        =}\NormalTok{ pop\_params\_sim}\SpecialCharTok{$}\NormalTok{pop\_log\_c\_sd,}
      \AttributeTok{pop\_log\_lambda\_mean =}\NormalTok{ pop\_params\_sim}\SpecialCharTok{$}\NormalTok{pop\_log\_lambda\_mean,}
      \AttributeTok{pop\_log\_lambda\_sd   =}\NormalTok{ pop\_params\_sim}\SpecialCharTok{$}\NormalTok{pop\_log\_lambda\_sd,}
      \AttributeTok{pop\_logit\_bias\_mean =}\NormalTok{ pop\_params\_sim}\SpecialCharTok{$}\NormalTok{pop\_logit\_bias\_mean,}
      \AttributeTok{pop\_logit\_bias\_sd   =}\NormalTok{ pop\_params\_sim}\SpecialCharTok{$}\NormalTok{pop\_logit\_bias\_sd}
\NormalTok{    )}
    \FunctionTok{write\_csv}\NormalTok{(out}\SpecialCharTok{$}\NormalTok{data,        ml\_decay\_files[[scen]]}\SpecialCharTok{$}\NormalTok{data)}
    \FunctionTok{write\_csv}\NormalTok{(out}\SpecialCharTok{$}\NormalTok{true\_params, ml\_decay\_files[[scen]]}\SpecialCharTok{$}\NormalTok{truth)}
\NormalTok{    out}
\NormalTok{  \}) }\SpecialCharTok{|\textgreater{}} \FunctionTok{setNames}\NormalTok{(ml\_scenarios)}
\NormalTok{\} }\ControlFlowTok{else}\NormalTok{ \{}
\NormalTok{  ml\_decay\_sims }\OtherTok{\textless{}{-}} \FunctionTok{map}\NormalTok{(ml\_scenarios, }\ControlFlowTok{function}\NormalTok{(scen) }\FunctionTok{list}\NormalTok{(}
    \AttributeTok{data        =} \FunctionTok{read\_csv}\NormalTok{(ml\_decay\_files[[scen]]}\SpecialCharTok{$}\NormalTok{data,  }\AttributeTok{show\_col\_types =} \ConstantTok{FALSE}\NormalTok{),}
    \AttributeTok{true\_params =} \FunctionTok{read\_csv}\NormalTok{(ml\_decay\_files[[scen]]}\SpecialCharTok{$}\NormalTok{truth, }\AttributeTok{show\_col\_types =} \ConstantTok{FALSE}\NormalTok{)}
\NormalTok{  )) }\SpecialCharTok{|\textgreater{}} \FunctionTok{setNames}\NormalTok{(ml\_scenarios)}
\NormalTok{\}}

\NormalTok{ml\_decay\_all }\OtherTok{\textless{}{-}} \FunctionTok{map\_dfr}\NormalTok{(ml\_scenarios, }\ControlFlowTok{function}\NormalTok{(scen)}
\NormalTok{  ml\_decay\_sims[[scen]]}\SpecialCharTok{$}\NormalTok{data }\SpecialCharTok{|\textgreater{}} \FunctionTok{mutate}\NormalTok{(}\AttributeTok{scenario =}\NormalTok{ scen)}
\NormalTok{) }\SpecialCharTok{|\textgreater{}}
  \FunctionTok{mutate}\NormalTok{(}\AttributeTok{scenario =} \FunctionTok{factor}\NormalTok{(scenario, }\AttributeTok{levels =}\NormalTok{ ml\_scenarios,}
                           \AttributeTok{labels =} \FunctionTok{c}\NormalTok{(}\StringTok{"Static"}\NormalTok{, }\StringTok{"Smooth drift"}\NormalTok{,}
                                      \StringTok{"Performance{-}contingent"}\NormalTok{)))}

\FunctionTok{ggplot}\NormalTok{(ml\_decay\_all,}
       \FunctionTok{aes}\NormalTok{(}\AttributeTok{x =}\NormalTok{ trial\_within\_subject, }\AttributeTok{y =}\NormalTok{ cumulative\_accuracy, }\AttributeTok{group =}\NormalTok{ agent\_id)) }\SpecialCharTok{+}
  \FunctionTok{geom\_line}\NormalTok{(}\AttributeTok{alpha =} \FloatTok{0.25}\NormalTok{, }\AttributeTok{color =} \StringTok{"grey30"}\NormalTok{) }\SpecialCharTok{+}
  \FunctionTok{stat\_summary}\NormalTok{(}\AttributeTok{fun =}\NormalTok{ mean, }\AttributeTok{geom =} \StringTok{"line"}\NormalTok{, }\FunctionTok{aes}\NormalTok{(}\AttributeTok{group =} \DecValTok{1}\NormalTok{),}
               \AttributeTok{color =} \StringTok{"\#0072B2"}\NormalTok{, }\AttributeTok{linewidth =} \FloatTok{1.4}\NormalTok{) }\SpecialCharTok{+}
  \FunctionTok{facet\_wrap}\NormalTok{(}\SpecialCharTok{\textasciitilde{}}\NormalTok{scenario, }\AttributeTok{ncol =} \DecValTok{3}\NormalTok{) }\SpecialCharTok{+}
  \FunctionTok{geom\_hline}\NormalTok{(}\AttributeTok{yintercept =} \FloatTok{0.5}\NormalTok{, }\AttributeTok{linetype =} \StringTok{"dashed"}\NormalTok{, }\AttributeTok{color =} \StringTok{"grey50"}\NormalTok{) }\SpecialCharTok{+}
  \FunctionTok{labs}\NormalTok{(}
    \AttributeTok{title    =} \StringTok{"Multilevel Decay GCM: simulated learning curves across scenarios"}\NormalTok{,}
    \AttributeTok{subtitle =} \FunctionTok{paste0}\NormalTok{(n\_subjects\_sim, }\StringTok{" agents per scenario; thin grey = individuals, "}\NormalTok{,}
                      \StringTok{"thick blue = population mean"}\NormalTok{),}
    \AttributeTok{x =} \StringTok{"Trial"}\NormalTok{, }\AttributeTok{y =} \StringTok{"Cumulative accuracy"}
\NormalTok{  )}
\end{Highlighting}
\end{Shaded}

\begin{center}
\includegraphics[width=0.8\linewidth,height=\textheight,keepaspectratio]{11-CategorizationModelsExemplars_files/figure-pdf/ch12_simulate_multilevel_decay_data-1.pdf}
\end{center}

The static panel converges quickly to ceiling; smooth drift plateaus
below ceiling because the boundary keeps moving; performance-contingent
shifts produce a noisy plateau as different subjects flip at different
absolute trials and average over each other.

\subsection{Stan Implementation}\label{stan-implementation-2}

The Stan model implements the hierarchy described above, with three
architectural choices that matter for performance:

\begin{itemize}
\tightlist
\item
  \textbf{The trial-by-trial likelihood lives in the \texttt{model}
  block, not in \texttt{transformed\ parameters}.} Memory is held in
  plain \texttt{int} arrays during the loop, so the autodiff tape never
  sees the per-trial scan. \texttt{prob\_cat1} and (optionally gated)
  \texttt{log\_lik} are recomputed once per saved draw in
  \texttt{generated\ quantities}.
\item
  \textbf{All hierarchies are non-centered} (logit-normal NCP for
  \texttt{w\_1}; standard z-offset for \texttt{log\_c},
  \texttt{log\_lambda}, and \texttt{logit\_bias}).
\item
  \textbf{Pairwise feature distances are precomputed in
  \texttt{transformed\ data}}, indexed by \texttt{stim\_id} instead of
  read from per-trial features.
\end{itemize}

\begin{Shaded}
\begin{Highlighting}[]
\NormalTok{gcm\_ml\_decay\_stan }\OtherTok{\textless{}{-}} \StringTok{"}
\StringTok{// Multilevel Decay Generalized Context Model}
\StringTok{// All individual parameters use non{-}centred parameterisation:}
\StringTok{//   {-} w\_1 via logit{-}normal NCP (eliminates Dirichlet funnel geometry)}
\StringTok{//   {-} log\_c, log\_lambda, logit\_bias via standard z{-}offset NCP}
\StringTok{// Sequential likelihood is computed directly in the model block (no}
\StringTok{// prob\_cat1 stored in transformed parameters), avoiding a large var{-}vector}
\StringTok{// on the autodiff tape. prob\_cat1 and log\_lik are recomputed once per saved}
\StringTok{// draw in generated quantities.}

\StringTok{data \{}
\StringTok{  int\textless{}lower=1\textgreater{} N\_total;}
\StringTok{  int\textless{}lower=1\textgreater{} N\_subjects;}
\StringTok{  int\textless{}lower=1\textgreater{} N\_features;}
\StringTok{  int\textless{}lower=1\textgreater{} N\_unique\_stim;}
\StringTok{  int\textless{}lower=1\textgreater{} max\_trials\_per\_subject;}

\StringTok{  array[N\_subjects] int\textless{}lower=1, upper=N\_total\textgreater{} subj\_start;}
\StringTok{  array[N\_subjects] int\textless{}lower=1, upper=N\_total\textgreater{} subj\_end;}

\StringTok{  array[N\_total] int\textless{}lower=1, upper=N\_unique\_stim\textgreater{} stim\_id;}
\StringTok{  array[N\_total] int\textless{}lower=0, upper=1\textgreater{} y;}
\StringTok{  array[N\_unique\_stim, N\_features] real unique\_obs;}
\StringTok{  array[N\_total] int\textless{}lower=0, upper=1\textgreater{} cat\_feedback;}

\StringTok{  // Hyperparameters}
\StringTok{  real prior\_pop\_logit\_w\_mean;}
\StringTok{  real\textless{}lower=0\textgreater{} prior\_pop\_logit\_w\_sigma;}
\StringTok{  real\textless{}lower=0\textgreater{} prior\_pop\_logit\_w\_sd\_lambda;}
\StringTok{  real prior\_pop\_log\_c\_mu;}
\StringTok{  real\textless{}lower=0\textgreater{} prior\_pop\_log\_c\_sigma;}
\StringTok{  real\textless{}lower=0\textgreater{} prior\_pop\_log\_c\_lambda;}
\StringTok{  real prior\_pop\_log\_lambda\_mu;}
\StringTok{  real\textless{}lower=0\textgreater{} prior\_pop\_log\_lambda\_sigma;}
\StringTok{  real\textless{}lower=0\textgreater{} prior\_pop\_log\_lambda\_lambda;}
\StringTok{  real prior\_pop\_logit\_bias\_mu;}
\StringTok{  real\textless{}lower=0\textgreater{} prior\_pop\_logit\_bias\_sigma;}
\StringTok{  real\textless{}lower=0\textgreater{} prior\_pop\_logit\_bias\_lambda;}
\StringTok{  int\textless{}lower=0, upper=1\textgreater{} run\_diagnostics;}
\StringTok{\}}

\StringTok{transformed data \{}
\StringTok{  array[N\_unique\_stim, N\_unique\_stim, N\_features] real stimulus\_abs\_diff;}
\StringTok{  for (i in 1:N\_unique\_stim)}
\StringTok{    for (k in 1:N\_unique\_stim)}
\StringTok{      for (f in 1:N\_features)}
\StringTok{        stimulus\_abs\_diff[i, k, f] = abs(unique\_obs[i, f] {-} unique\_obs[k, f]);}
\StringTok{\}}

\StringTok{parameters \{}
\StringTok{  // Population attention weight: logit{-}normal NCP for w\_1 (w\_2 = 1 {-} w\_1)}
\StringTok{  real pop\_logit\_w\_mean;}
\StringTok{  real\textless{}lower=0\textgreater{} pop\_logit\_w\_sd;}
\StringTok{  vector[N\_subjects] z\_w;}

\StringTok{  // Population c, lambda, bias}
\StringTok{  real pop\_log\_c\_mean;}
\StringTok{  real\textless{}lower=0\textgreater{} pop\_log\_c\_sd;}
\StringTok{  real pop\_log\_lambda\_mean;}
\StringTok{  real\textless{}lower=0\textgreater{} pop\_log\_lambda\_sd;}
\StringTok{  real pop\_logit\_bias\_mean;}
\StringTok{  real\textless{}lower=0\textgreater{} pop\_logit\_bias\_sd;}

\StringTok{  // Non{-}centred individual offsets}
\StringTok{  vector[N\_subjects] z\_log\_c;}
\StringTok{  vector[N\_subjects] z\_log\_lambda;}
\StringTok{  vector[N\_subjects] z\_logit\_bias;}
\StringTok{\}}

\StringTok{transformed parameters \{}
\StringTok{  vector[N\_subjects] subj\_log\_c;}
\StringTok{  vector[N\_subjects] subj\_log\_lambda;}
\StringTok{  vector\textless{}lower=0, upper=1\textgreater{}[N\_subjects] subj\_bias;}
\StringTok{  array[N\_subjects] vector[N\_features] subj\_w;}

\StringTok{  for (j in 1:N\_subjects) \{}
\StringTok{    subj\_log\_c[j]      = pop\_log\_c\_mean + pop\_log\_c\_sd * z\_log\_c[j];}
\StringTok{    subj\_log\_lambda[j] = pop\_log\_lambda\_mean + pop\_log\_lambda\_sd * z\_log\_lambda[j];}
\StringTok{    subj\_bias[j]       = inv\_logit(pop\_logit\_bias\_mean + pop\_logit\_bias\_sd * z\_logit\_bias[j]);}
\StringTok{    real logit\_w       = pop\_logit\_w\_mean + pop\_logit\_w\_sd * z\_w[j];}
\StringTok{    subj\_w[j][1]       = inv\_logit(logit\_w);}
\StringTok{    subj\_w[j][2]       = 1.0 {-} subj\_w[j][1];}
\StringTok{  \}}
\StringTok{\}}

\StringTok{model \{}
\StringTok{  // Population priors}
\StringTok{  pop\_logit\_w\_mean    \textasciitilde{} normal(prior\_pop\_logit\_w\_mean, prior\_pop\_logit\_w\_sigma);}
\StringTok{  pop\_logit\_w\_sd      \textasciitilde{} exponential(prior\_pop\_logit\_w\_sd\_lambda);}
\StringTok{  pop\_log\_c\_mean      \textasciitilde{} normal(prior\_pop\_log\_c\_mu, prior\_pop\_log\_c\_sigma);}
\StringTok{  pop\_log\_c\_sd        \textasciitilde{} exponential(prior\_pop\_log\_c\_lambda);}
\StringTok{  pop\_log\_lambda\_mean \textasciitilde{} normal(prior\_pop\_log\_lambda\_mu, prior\_pop\_log\_lambda\_sigma);}
\StringTok{  pop\_log\_lambda\_sd   \textasciitilde{} exponential(prior\_pop\_log\_lambda\_lambda);}
\StringTok{  pop\_logit\_bias\_mean \textasciitilde{} normal(prior\_pop\_logit\_bias\_mu, prior\_pop\_logit\_bias\_sigma);}
\StringTok{  pop\_logit\_bias\_sd   \textasciitilde{} exponential(prior\_pop\_logit\_bias\_lambda);}

\StringTok{  // Non{-}centred priors}
\StringTok{  z\_w          \textasciitilde{} std\_normal();}
\StringTok{  z\_log\_c      \textasciitilde{} std\_normal();}
\StringTok{  z\_log\_lambda \textasciitilde{} std\_normal();}
\StringTok{  z\_logit\_bias \textasciitilde{} std\_normal();}

\StringTok{  // Sequential likelihood — local workspace, one subject at a time}
\StringTok{  \{}
\StringTok{    array[max\_trials\_per\_subject] int mem\_stim;}
\StringTok{    array[max\_trials\_per\_subject] int mem\_cat;}
\StringTok{    array[max\_trials\_per\_subject] int mem\_trial;}

\StringTok{    for (j in 1:N\_subjects) \{}
\StringTok{      int n\_mem  = 0;}
\StringTok{      real c      = exp(subj\_log\_c[j]);}
\StringTok{      real lambda = exp(subj\_log\_lambda[j]);}
\StringTok{      real bias   = subj\_bias[j];}

\StringTok{      for (i in subj\_start[j]:subj\_end[j]) \{}
\StringTok{        real p\_i;}
\StringTok{        if (n\_mem == 0) \{}
\StringTok{          p\_i = bias;}
\StringTok{        \} else \{}
\StringTok{          real s1 = 0.0; real s0 = 0.0;}
\StringTok{          real w1 = 0.0; real w0 = 0.0;}
\StringTok{          int curr\_stim = stim\_id[i];}
\StringTok{          for (e in 1:n\_mem) \{}
\StringTok{            real dist = 0.0;}
\StringTok{            int past\_stim = mem\_stim[e];}
\StringTok{            for (f in 1:N\_features)}
\StringTok{              dist += subj\_w[j][f] * stimulus\_abs\_diff[curr\_stim, past\_stim, f];}
\StringTok{            real age   = (i {-} subj\_start[j] + 1.0) {-} mem\_trial[e];}
\StringTok{            real decay = exp({-}lambda * age);}
\StringTok{            real sim   = exp({-}c * dist);}
\StringTok{            real wt    = decay * sim;}
\StringTok{            if (mem\_cat[e] == 1) \{ s1 += wt; w1 += decay; \}}
\StringTok{            else                 \{ s0 += wt; w0 += decay; \}}
\StringTok{          \}}
\StringTok{          real m1  = (w1 \textgreater{} 1e{-}9) ? s1 / w1 : 0.0;}
\StringTok{          real m0  = (w0 \textgreater{} 1e{-}9) ? s0 / w0 : 0.0;}
\StringTok{          real num = bias * m1;}
\StringTok{          real den = num + (1.0 {-} bias) * m0;}
\StringTok{          p\_i = (den \textgreater{} 1e{-}9) ? num / den : bias;}
\StringTok{        \}}
\StringTok{        target += bernoulli\_lpmf(y[i] | fmax(1e{-}9, fmin(1.0 {-} 1e{-}9, p\_i)));}

\StringTok{        n\_mem += 1;}
\StringTok{        mem\_stim[n\_mem]  = stim\_id[i];}
\StringTok{        mem\_cat[n\_mem]   = cat\_feedback[i];}
\StringTok{        mem\_trial[n\_mem] = i {-} subj\_start[j] + 1;}
\StringTok{      \}}
\StringTok{    \}}
\StringTok{  \}}
\StringTok{\}}

\StringTok{generated quantities \{}
\StringTok{  // Recompute once per saved draw (not on every leapfrog step).}
\StringTok{  // prob\_cat1 is always filled (needed for PPC); log\_lik only when run\_diagnostics=1.}
\StringTok{  vector[N\_total] prob\_cat1;}
\StringTok{  vector[N\_total] log\_lik = rep\_vector(0.0, N\_total);}

\StringTok{  \{}
\StringTok{    array[max\_trials\_per\_subject] int mem\_stim;}
\StringTok{    array[max\_trials\_per\_subject] int mem\_cat;}
\StringTok{    array[max\_trials\_per\_subject] int mem\_trial;}

\StringTok{    for (j in 1:N\_subjects) \{}
\StringTok{      int n\_mem  = 0;}
\StringTok{      real c      = exp(subj\_log\_c[j]);}
\StringTok{      real lambda = exp(subj\_log\_lambda[j]);}
\StringTok{      real bias   = subj\_bias[j];}

\StringTok{      for (i in subj\_start[j]:subj\_end[j]) \{}
\StringTok{        real p\_i;}
\StringTok{        if (n\_mem == 0) \{}
\StringTok{          p\_i = bias;}
\StringTok{        \} else \{}
\StringTok{          real s1 = 0.0; real s0 = 0.0;}
\StringTok{          real w1 = 0.0; real w0 = 0.0;}
\StringTok{          int curr\_stim = stim\_id[i];}
\StringTok{          for (e in 1:n\_mem) \{}
\StringTok{            real dist = 0.0;}
\StringTok{            int past\_stim = mem\_stim[e];}
\StringTok{            for (f in 1:N\_features)}
\StringTok{              dist += subj\_w[j][f] * stimulus\_abs\_diff[curr\_stim, past\_stim, f];}
\StringTok{            real age   = (i {-} subj\_start[j] + 1.0) {-} mem\_trial[e];}
\StringTok{            real decay = exp({-}lambda * age);}
\StringTok{            real sim   = exp({-}c * dist);}
\StringTok{            real wt    = decay * sim;}
\StringTok{            if (mem\_cat[e] == 1) \{ s1 += wt; w1 += decay; \}}
\StringTok{            else                 \{ s0 += wt; w0 += decay; \}}
\StringTok{          \}}
\StringTok{          real m1  = (w1 \textgreater{} 1e{-}9) ? s1 / w1 : 0.0;}
\StringTok{          real m0  = (w0 \textgreater{} 1e{-}9) ? s0 / w0 : 0.0;}
\StringTok{          real num = bias * m1;}
\StringTok{          real den = num + (1.0 {-} bias) * m0;}
\StringTok{          p\_i = (den \textgreater{} 1e{-}9) ? num / den : bias;}
\StringTok{        \}}
\StringTok{        real p\_clamp  = fmax(1e{-}9, fmin(1.0 {-} 1e{-}9, p\_i));}
\StringTok{        prob\_cat1[i]  = p\_clamp;}
\StringTok{        if (run\_diagnostics)}
\StringTok{          log\_lik[i]  = bernoulli\_lpmf(y[i] | p\_clamp);}

\StringTok{        n\_mem += 1;}
\StringTok{        mem\_stim[n\_mem]  = stim\_id[i];}
\StringTok{        mem\_cat[n\_mem]   = cat\_feedback[i];}
\StringTok{        mem\_trial[n\_mem] = i {-} subj\_start[j] + 1;}
\StringTok{      \}}
\StringTok{    \}}
\StringTok{  \}}
\StringTok{\}}
\StringTok{"}

\NormalTok{stan\_file\_gcm\_ml\_decay }\OtherTok{\textless{}{-}} \StringTok{"stan/ch12\_gcm\_ml\_decay.stan"}
\FunctionTok{write\_stan\_file}\NormalTok{(gcm\_ml\_decay\_stan, }\AttributeTok{dir =} \StringTok{"stan/"}\NormalTok{, }\AttributeTok{basename =} \StringTok{"ch12\_gcm\_ml\_decay.stan"}\NormalTok{)}
\end{Highlighting}
\end{Shaded}

\begin{verbatim}
[1] "/Users/au209589/Dropbox/Teaching/AdvancedCognitiveModeling23_book/stan/ch12_gcm_ml_decay.stan"
\end{verbatim}

\begin{Shaded}
\begin{Highlighting}[]
\NormalTok{mod\_gcm\_ml\_decay       }\OtherTok{\textless{}{-}} \FunctionTok{cmdstan\_model}\NormalTok{(stan\_file\_gcm\_ml\_decay)}
\end{Highlighting}
\end{Shaded}

\subsection{Building the Stan Data
List}\label{building-the-stan-data-list}

\begin{Shaded}
\begin{Highlighting}[]
\NormalTok{build\_ml\_stan\_data }\OtherTok{\textless{}{-}} \ControlFlowTok{function}\NormalTok{(ml\_data, stimulus\_info, }\AttributeTok{run\_diagnostics =} \DecValTok{0}\NormalTok{L) \{}
\NormalTok{  ml\_data }\OtherTok{\textless{}{-}}\NormalTok{ ml\_data }\SpecialCharTok{|\textgreater{}}
    \FunctionTok{arrange}\NormalTok{(agent\_id, trial\_within\_subject) }\SpecialCharTok{|\textgreater{}}
    \FunctionTok{mutate}\NormalTok{(}\AttributeTok{row =} \FunctionTok{row\_number}\NormalTok{())}

\NormalTok{  unique\_stim }\OtherTok{\textless{}{-}}\NormalTok{ stimulus\_info }\SpecialCharTok{|\textgreater{}}
    \FunctionTok{arrange}\NormalTok{(stimulus) }\SpecialCharTok{|\textgreater{}}
    \FunctionTok{mutate}\NormalTok{(}\AttributeTok{unique\_idx =} \FunctionTok{row\_number}\NormalTok{())}
\NormalTok{  stim\_lookup }\OtherTok{\textless{}{-}} \FunctionTok{setNames}\NormalTok{(unique\_stim}\SpecialCharTok{$}\NormalTok{unique\_idx, unique\_stim}\SpecialCharTok{$}\NormalTok{stimulus)}

\NormalTok{  subj\_idx }\OtherTok{\textless{}{-}}\NormalTok{ ml\_data }\SpecialCharTok{|\textgreater{}}
    \FunctionTok{group\_by}\NormalTok{(agent\_id) }\SpecialCharTok{|\textgreater{}}
    \FunctionTok{summarise}\NormalTok{(}\AttributeTok{start =} \FunctionTok{min}\NormalTok{(row), }\AttributeTok{end =} \FunctionTok{max}\NormalTok{(row), }\AttributeTok{.groups =} \StringTok{"drop"}\NormalTok{) }\SpecialCharTok{|\textgreater{}}
    \FunctionTok{arrange}\NormalTok{(agent\_id)}

  \FunctionTok{list}\NormalTok{(}
    \AttributeTok{N\_total                =} \FunctionTok{nrow}\NormalTok{(ml\_data),}
    \AttributeTok{N\_subjects             =} \FunctionTok{nrow}\NormalTok{(subj\_idx),}
    \AttributeTok{N\_features             =} \DecValTok{2}\NormalTok{L,}
    \AttributeTok{N\_unique\_stim          =} \FunctionTok{nrow}\NormalTok{(unique\_stim),}
    \AttributeTok{max\_trials\_per\_subject =} \FunctionTok{as.integer}\NormalTok{(}\FunctionTok{max}\NormalTok{(subj\_idx}\SpecialCharTok{$}\NormalTok{end }\SpecialCharTok{{-}}\NormalTok{ subj\_idx}\SpecialCharTok{$}\NormalTok{start }\SpecialCharTok{+} \DecValTok{1}\NormalTok{L)),}

    \AttributeTok{subj\_start             =} \FunctionTok{as.integer}\NormalTok{(subj\_idx}\SpecialCharTok{$}\NormalTok{start),}
    \AttributeTok{subj\_end               =} \FunctionTok{as.integer}\NormalTok{(subj\_idx}\SpecialCharTok{$}\NormalTok{end),}

    \AttributeTok{stim\_id                =} \FunctionTok{as.integer}\NormalTok{(stim\_lookup[}\FunctionTok{as.character}\NormalTok{(ml\_data}\SpecialCharTok{$}\NormalTok{stimulus\_id)]),}
    \AttributeTok{y                      =} \FunctionTok{as.integer}\NormalTok{(ml\_data}\SpecialCharTok{$}\NormalTok{sim\_response),}
    \AttributeTok{unique\_obs             =} \FunctionTok{as.matrix}\NormalTok{(unique\_stim[, }\FunctionTok{c}\NormalTok{(}\StringTok{"height"}\NormalTok{, }\StringTok{"position"}\NormalTok{)]),}
    \AttributeTok{cat\_feedback           =} \FunctionTok{as.integer}\NormalTok{(ml\_data}\SpecialCharTok{$}\NormalTok{observed\_feedback),}

    \AttributeTok{prior\_pop\_logit\_w\_mean      =} \DecValTok{0}\NormalTok{,}
    \AttributeTok{prior\_pop\_logit\_w\_sigma     =} \FloatTok{1.0}\NormalTok{,}
    \AttributeTok{prior\_pop\_logit\_w\_sd\_lambda =} \FloatTok{2.0}\NormalTok{,}
    \AttributeTok{prior\_pop\_log\_c\_mu          =}\NormalTok{ LOG\_C\_PRIOR\_MEAN,}
    \AttributeTok{prior\_pop\_log\_c\_sigma       =}\NormalTok{ LOG\_C\_PRIOR\_SD,}
    \AttributeTok{prior\_pop\_log\_c\_lambda      =} \FloatTok{2.0}\NormalTok{,}
    \AttributeTok{prior\_pop\_log\_lambda\_mu     =}\NormalTok{ LOG\_LAMBDA\_PRIOR\_MEAN,}
    \AttributeTok{prior\_pop\_log\_lambda\_sigma  =}\NormalTok{ LOG\_LAMBDA\_PRIOR\_SD,}
    \AttributeTok{prior\_pop\_log\_lambda\_lambda =} \FloatTok{2.0}\NormalTok{,}
    \AttributeTok{prior\_pop\_logit\_bias\_mu     =} \DecValTok{0}\NormalTok{,}
    \AttributeTok{prior\_pop\_logit\_bias\_sigma  =} \FloatTok{1.0}\NormalTok{,}
    \AttributeTok{prior\_pop\_logit\_bias\_lambda =} \FloatTok{2.0}\NormalTok{,}

    \AttributeTok{run\_diagnostics             =} \FunctionTok{as.integer}\NormalTok{(run\_diagnostics)}
\NormalTok{  )}
\NormalTok{\}}
\end{Highlighting}
\end{Shaded}

The \texttt{Exponential(2)} priors on the population SDs
(\texttt{*\_lambda\ =\ 2}) have median 0.35 and 95th percentile ≈ 1.5
--- wide enough to allow real between-subject variation, tight enough to
keep NUTS off the worst part of the funnel when only 10 subjects are
observed. We keep \texttt{run\_diagnostics\ =\ 0} for the main fits
because we will not run LFO at the multilevel level here; flip it to
\texttt{1} only when computing per-trial \texttt{log\_lik} is needed.

\subsection{Fitting the Multilevel Decay GCM Across
Scenarios}\label{fitting-the-multilevel-decay-gcm-across-scenarios}

We compile once and fit each scenario in turn. With 10 subjects × 64
trials × 4 chains × 2000 iterations, expect each fit to take a few
minutes on a modern laptop.

\begin{Shaded}
\begin{Highlighting}[]
\NormalTok{ml\_decay\_fit\_files }\OtherTok{\textless{}{-}} \FunctionTok{map}\NormalTok{(ml\_scenarios, }\ControlFlowTok{function}\NormalTok{(scen)}
  \FunctionTok{here}\NormalTok{(}\StringTok{"simmodels"}\NormalTok{, }\FunctionTok{paste0}\NormalTok{(}\StringTok{"ch12\_gcm\_ml\_decay\_fit\_"}\NormalTok{, scen, }\StringTok{".rds"}\NormalTok{))}
\NormalTok{) }\SpecialCharTok{|\textgreater{}} \FunctionTok{setNames}\NormalTok{(ml\_scenarios)}

\ControlFlowTok{if}\NormalTok{ (regenerate\_simulations }\SpecialCharTok{||} \SpecialCharTok{!}\FunctionTok{all}\NormalTok{(}\FunctionTok{file.exists}\NormalTok{(}\FunctionTok{unlist}\NormalTok{(ml\_decay\_fit\_files)))) \{}
\NormalTok{  ml\_decay\_fits }\OtherTok{\textless{}{-}} \FunctionTok{map}\NormalTok{(ml\_scenarios, }\ControlFlowTok{function}\NormalTok{(scen) \{}
    \FunctionTok{cat}\NormalTok{(}\StringTok{"Fitting multilevel decay GCM:"}\NormalTok{, scen, }\StringTok{"}\SpecialCharTok{\textbackslash{}n}\StringTok{"}\NormalTok{)}
\NormalTok{    dat }\OtherTok{\textless{}{-}} \FunctionTok{build\_ml\_stan\_data}\NormalTok{(ml\_decay\_sims[[scen]]}\SpecialCharTok{$}\NormalTok{data, stimulus\_info,}
                              \AttributeTok{run\_diagnostics =} \DecValTok{0}\NormalTok{L)}
\NormalTok{    fit }\OtherTok{\textless{}{-}}\NormalTok{ mod\_gcm\_ml\_decay}\SpecialCharTok{$}\FunctionTok{sample}\NormalTok{(}
      \AttributeTok{data            =}\NormalTok{ dat,}
      \AttributeTok{seed            =} \DecValTok{700} \SpecialCharTok{+} \FunctionTok{match}\NormalTok{(scen, ml\_scenarios),}
      \AttributeTok{chains          =} \DecValTok{4}\NormalTok{,}
      \AttributeTok{parallel\_chains =} \FunctionTok{min}\NormalTok{(}\DecValTok{4}\NormalTok{, }\FunctionTok{availableCores}\NormalTok{()),}
      \AttributeTok{iter\_warmup     =} \DecValTok{1000}\NormalTok{,}
      \AttributeTok{iter\_sampling   =} \DecValTok{1000}\NormalTok{,}
      \AttributeTok{adapt\_delta     =} \FloatTok{0.95}\NormalTok{,}
      \AttributeTok{max\_treedepth   =} \DecValTok{12}\NormalTok{,}
      \AttributeTok{refresh         =} \DecValTok{200}\NormalTok{,}
      \AttributeTok{init            =} \FloatTok{0.5}
\NormalTok{    )}
\NormalTok{    fit}\SpecialCharTok{$}\FunctionTok{save\_object}\NormalTok{(ml\_decay\_fit\_files[[scen]])}
\NormalTok{    fit}
\NormalTok{  \}) }\SpecialCharTok{|\textgreater{}} \FunctionTok{setNames}\NormalTok{(ml\_scenarios)}
  \FunctionTok{cat}\NormalTok{(}\StringTok{"Multilevel decay GCM fits computed and saved.}\SpecialCharTok{\textbackslash{}n}\StringTok{"}\NormalTok{)}
\NormalTok{\} }\ControlFlowTok{else}\NormalTok{ \{}
\NormalTok{  ml\_decay\_fits }\OtherTok{\textless{}{-}} \FunctionTok{map}\NormalTok{(ml\_scenarios, }\ControlFlowTok{function}\NormalTok{(scen)}
    \FunctionTok{readRDS}\NormalTok{(ml\_decay\_fit\_files[[scen]])}
\NormalTok{  ) }\SpecialCharTok{|\textgreater{}} \FunctionTok{setNames}\NormalTok{(ml\_scenarios)}
  \FunctionTok{cat}\NormalTok{(}\StringTok{"Loaded existing multilevel decay GCM fits.}\SpecialCharTok{\textbackslash{}n}\StringTok{"}\NormalTok{)}
\NormalTok{\}}
\end{Highlighting}
\end{Shaded}

\begin{verbatim}
Loaded existing multilevel decay GCM fits.
\end{verbatim}

\subsection{MCMC Diagnostic Battery}\label{mcmc-diagnostic-battery-1}

\begin{Shaded}
\begin{Highlighting}[]
\NormalTok{ml\_decay\_diag }\OtherTok{\textless{}{-}} \FunctionTok{map\_dfr}\NormalTok{(ml\_scenarios, }\ControlFlowTok{function}\NormalTok{(scen) \{}
  \FunctionTok{diagnostic\_summary\_table}\NormalTok{(}
\NormalTok{    ml\_decay\_fits[[scen]],}
    \AttributeTok{params =} \FunctionTok{c}\NormalTok{(}\StringTok{"pop\_logit\_w\_mean"}\NormalTok{,    }\StringTok{"pop\_log\_c\_mean"}\NormalTok{,}
               \StringTok{"pop\_log\_lambda\_mean"}\NormalTok{, }\StringTok{"pop\_logit\_bias\_mean"}\NormalTok{,}
               \StringTok{"pop\_logit\_w\_sd"}\NormalTok{,      }\StringTok{"pop\_log\_c\_sd"}\NormalTok{,}
               \StringTok{"pop\_log\_lambda\_sd"}\NormalTok{,   }\StringTok{"pop\_logit\_bias\_sd"}\NormalTok{)}
\NormalTok{  ) }\SpecialCharTok{|\textgreater{}} \FunctionTok{mutate}\NormalTok{(}\AttributeTok{scenario =}\NormalTok{ scen)}
\NormalTok{\})}

\FunctionTok{print}\NormalTok{(ml\_decay\_diag }\SpecialCharTok{|\textgreater{}}
\NormalTok{        dplyr}\SpecialCharTok{::}\FunctionTok{select}\NormalTok{(scenario, metric, value, threshold, pass))}
\end{Highlighting}
\end{Shaded}

\begin{verbatim}
# A tibble: 18 x 5
   scenario         metric                          value threshold pass 
   <chr>            <chr>                           <dbl> <chr>     <lgl>
 1 static           Divergences (zero tolerance)   0      == 0      TRUE 
 2 static           Max rank-normalised R-hat      2.84   < 1.01    FALSE
 3 static           Min bulk ESS                   4.65   > 400     FALSE
 4 static           Min tail ESS                  13.3    > 400     FALSE
 5 static           Min E-BFMI                     1.12   > 0.2     TRUE 
 6 static           Max MCSE / posterior SD        0.488  < 0.05    FALSE
 7 contingent_shift Divergences (zero tolerance)   0      == 0      TRUE 
 8 contingent_shift Max rank-normalised R-hat      2.93   < 1.01    FALSE
 9 contingent_shift Min bulk ESS                   4.59   > 400     FALSE
10 contingent_shift Min tail ESS                  11.2    > 400     FALSE
11 contingent_shift Min E-BFMI                     0.971  > 0.2     TRUE 
12 contingent_shift Max MCSE / posterior SD        0.492  < 0.05    FALSE
13 drift            Divergences (zero tolerance)   0      == 0      TRUE 
14 drift            Max rank-normalised R-hat      1.00   < 1.01    TRUE 
15 drift            Min bulk ESS                 944.     > 400     TRUE 
16 drift            Min tail ESS                 493.     > 400     TRUE 
17 drift            Min E-BFMI                     0.694  > 0.2     TRUE 
18 drift            Max MCSE / posterior SD        0.0302 < 0.05    TRUE 
\end{verbatim}

If the model is healthy you will see no divergences, R̂ \textless{} 1.01
across all population parameters, and ESS comfortably above 400 for the
population means. Subject-level scalars typically have somewhat lower
ESS, especially for \texttt{log\_lambda} under the static scenario ---
the data carry little information about forgetting when the boundary
never moves, so the posterior is partly pinned by the prior and the
partial pooling toward the population mean.

\subsection{Population-Level Recovery}\label{population-level-recovery}

\begin{Shaded}
\begin{Highlighting}[]
\NormalTok{pop\_truth }\OtherTok{\textless{}{-}} \FunctionTok{with}\NormalTok{(pop\_params\_sim, }\FunctionTok{tibble}\NormalTok{(}
  \AttributeTok{variable =} \FunctionTok{c}\NormalTok{(}\StringTok{"pop\_logit\_w\_mean"}\NormalTok{, }\StringTok{"pop\_log\_c\_mean"}\NormalTok{,}
               \StringTok{"pop\_log\_lambda\_mean"}\NormalTok{, }\StringTok{"pop\_logit\_bias\_mean"}\NormalTok{),}
  \AttributeTok{truth    =} \FunctionTok{c}\NormalTok{(pop\_logit\_w\_mean, pop\_log\_c\_mean,}
\NormalTok{               pop\_log\_lambda\_mean, pop\_logit\_bias\_mean)}
\NormalTok{))}

\CommentTok{\# Extract all posterior draws for the population means}
\NormalTok{pop\_draws }\OtherTok{\textless{}{-}} \FunctionTok{map\_dfr}\NormalTok{(ml\_scenarios, }\ControlFlowTok{function}\NormalTok{(scen) \{}
  \FunctionTok{as\_draws\_df}\NormalTok{(ml\_decay\_fits[[scen]]}\SpecialCharTok{$}\FunctionTok{draws}\NormalTok{(pop\_truth}\SpecialCharTok{$}\NormalTok{variable)) }\SpecialCharTok{|\textgreater{}}
    \FunctionTok{pivot\_longer}\NormalTok{(}\AttributeTok{cols =} \FunctionTok{all\_of}\NormalTok{(pop\_truth}\SpecialCharTok{$}\NormalTok{variable), }
                 \AttributeTok{names\_to =} \StringTok{"variable"}\NormalTok{, }\AttributeTok{values\_to =} \StringTok{"value"}\NormalTok{) }\SpecialCharTok{|\textgreater{}}
    \FunctionTok{mutate}\NormalTok{(}\AttributeTok{scenario =}\NormalTok{ scen)}
\NormalTok{\}) }\SpecialCharTok{|\textgreater{}}
  \FunctionTok{mutate}\NormalTok{(}\AttributeTok{scenario =} \FunctionTok{factor}\NormalTok{(scenario, }\AttributeTok{levels =}\NormalTok{ ml\_scenarios,}
                           \AttributeTok{labels =} \FunctionTok{c}\NormalTok{(}\StringTok{"Static"}\NormalTok{, }\StringTok{"Drift"}\NormalTok{, }\StringTok{"Contingent shift"}\NormalTok{)),}
         \AttributeTok{variable =} \FunctionTok{factor}\NormalTok{(variable, }\AttributeTok{levels =}\NormalTok{ pop\_truth}\SpecialCharTok{$}\NormalTok{variable))}

\CommentTok{\# Precompute prior densities for overlay}
\NormalTok{prior\_params }\OtherTok{\textless{}{-}} \FunctionTok{tibble}\NormalTok{(}
  \AttributeTok{variable =} \FunctionTok{factor}\NormalTok{(pop\_truth}\SpecialCharTok{$}\NormalTok{variable, }\AttributeTok{levels =}\NormalTok{ pop\_truth}\SpecialCharTok{$}\NormalTok{variable),}
  \AttributeTok{mu    =} \FunctionTok{c}\NormalTok{(}\DecValTok{0}\NormalTok{, LOG\_C\_PRIOR\_MEAN, LOG\_LAMBDA\_PRIOR\_MEAN, }\DecValTok{0}\NormalTok{),}
  \AttributeTok{sigma =} \FunctionTok{c}\NormalTok{(}\DecValTok{1}\NormalTok{, LOG\_C\_PRIOR\_SD, LOG\_LAMBDA\_PRIOR\_SD, }\DecValTok{1}\NormalTok{)}
\NormalTok{)}

\NormalTok{prior\_density }\OtherTok{\textless{}{-}}\NormalTok{ prior\_params }\SpecialCharTok{|\textgreater{}}
  \FunctionTok{group\_by}\NormalTok{(variable) }\SpecialCharTok{|\textgreater{}}
  \FunctionTok{reframe}\NormalTok{(}
    \AttributeTok{x =} \FunctionTok{seq}\NormalTok{(mu }\SpecialCharTok{{-}} \FloatTok{3.5}\SpecialCharTok{*}\NormalTok{sigma, mu }\SpecialCharTok{+} \FloatTok{3.5}\SpecialCharTok{*}\NormalTok{sigma, }\AttributeTok{length.out =} \DecValTok{200}\NormalTok{),}
    \AttributeTok{y =} \FunctionTok{dnorm}\NormalTok{(x, mu, sigma)}
\NormalTok{  )}

\FunctionTok{ggplot}\NormalTok{() }\SpecialCharTok{+}
  \CommentTok{\# Prior distribution}
  \FunctionTok{geom\_line}\NormalTok{(}\AttributeTok{data =}\NormalTok{ prior\_density, }\FunctionTok{aes}\NormalTok{(}\AttributeTok{x =}\NormalTok{ x, }\AttributeTok{y =}\NormalTok{ y), }
            \AttributeTok{linetype =} \StringTok{"dotted"}\NormalTok{, }\AttributeTok{color =} \StringTok{"black"}\NormalTok{, }\AttributeTok{linewidth =} \FloatTok{0.8}\NormalTok{) }\SpecialCharTok{+}
  \CommentTok{\# Posterior distributions per scenario}
  \FunctionTok{geom\_density}\NormalTok{(}\AttributeTok{data =}\NormalTok{ pop\_draws, }\FunctionTok{aes}\NormalTok{(}\AttributeTok{x =}\NormalTok{ value, }\AttributeTok{fill =}\NormalTok{ scenario), }
               \AttributeTok{alpha =} \FloatTok{0.4}\NormalTok{, }\AttributeTok{color =} \ConstantTok{NA}\NormalTok{) }\SpecialCharTok{+}
  \CommentTok{\# Simulation truth}
  \FunctionTok{geom\_vline}\NormalTok{(}\AttributeTok{data =}\NormalTok{ pop\_truth, }\FunctionTok{aes}\NormalTok{(}\AttributeTok{xintercept =}\NormalTok{ truth), }
             \AttributeTok{linetype =} \StringTok{"dashed"}\NormalTok{, }\AttributeTok{color =} \StringTok{"red"}\NormalTok{, }\AttributeTok{linewidth =} \FloatTok{0.8}\NormalTok{) }\SpecialCharTok{+}
  \FunctionTok{facet\_wrap}\NormalTok{(}\SpecialCharTok{\textasciitilde{}}\NormalTok{variable, }\AttributeTok{scales =} \StringTok{"free"}\NormalTok{, }\AttributeTok{ncol =} \DecValTok{2}\NormalTok{) }\SpecialCharTok{+}
  \FunctionTok{scale\_fill\_viridis\_d}\NormalTok{(}\AttributeTok{option =} \StringTok{"plasma"}\NormalTok{, }\AttributeTok{end =} \FloatTok{0.85}\NormalTok{) }\SpecialCharTok{+}
  \FunctionTok{labs}\NormalTok{(}
    \AttributeTok{title    =} \StringTok{"Population{-}level parameter recovery"}\NormalTok{,}
    \AttributeTok{subtitle =} \StringTok{"Posterior densities vs Simulation Truth (red dashed) and Prior (black dotted)"}\NormalTok{,}
    \AttributeTok{x =} \StringTok{"Parameter value (transformed scale)"}\NormalTok{, }\AttributeTok{y =} \StringTok{"Density"}
\NormalTok{  ) }\SpecialCharTok{+}
  \FunctionTok{theme}\NormalTok{(}\AttributeTok{legend.position =} \StringTok{"bottom"}\NormalTok{)}
\end{Highlighting}
\end{Shaded}

\begin{center}
\includegraphics[width=0.8\linewidth,height=\textheight,keepaspectratio]{11-CategorizationModelsExemplars_files/figure-pdf/ch12_ml_decay_population_recovery-1.pdf}
\end{center}

\subsection{Subject-Level Recovery and
Shrinkage}\label{subject-level-recovery-and-shrinkage}

\begin{Shaded}
\begin{Highlighting}[]
\NormalTok{subj\_recovery }\OtherTok{\textless{}{-}} \FunctionTok{map\_dfr}\NormalTok{(ml\_scenarios, }\ControlFlowTok{function}\NormalTok{(scen) \{}
\NormalTok{  fit   }\OtherTok{\textless{}{-}}\NormalTok{ ml\_decay\_fits[[scen]]}
\NormalTok{  truth }\OtherTok{\textless{}{-}}\NormalTok{ ml\_decay\_sims[[scen]]}\SpecialCharTok{$}\NormalTok{true\_params}

\NormalTok{  s }\OtherTok{\textless{}{-}}\NormalTok{ fit}\SpecialCharTok{$}\FunctionTok{summary}\NormalTok{(}\AttributeTok{variables =} \FunctionTok{c}\NormalTok{(}\StringTok{"subj\_log\_c"}\NormalTok{, }\StringTok{"subj\_log\_lambda"}\NormalTok{)) }\SpecialCharTok{|\textgreater{}}
    \FunctionTok{mutate}\NormalTok{(}
      \AttributeTok{param    =} \FunctionTok{if\_else}\NormalTok{(stringr}\SpecialCharTok{::}\FunctionTok{str\_detect}\NormalTok{(variable, }\StringTok{"log\_c"}\NormalTok{),}
                         \StringTok{"log\_c"}\NormalTok{, }\StringTok{"log\_lambda"}\NormalTok{),}
      \AttributeTok{agent\_id =} \FunctionTok{as.integer}\NormalTok{(stringr}\SpecialCharTok{::}\FunctionTok{str\_extract}\NormalTok{(variable, }\StringTok{"}\SpecialCharTok{\textbackslash{}\textbackslash{}}\StringTok{d+"}\NormalTok{))}
\NormalTok{    ) }\SpecialCharTok{|\textgreater{}}
\NormalTok{    dplyr}\SpecialCharTok{::}\FunctionTok{select}\NormalTok{(agent\_id, param, }\AttributeTok{est\_mean =}\NormalTok{ mean, }\AttributeTok{est\_q5 =}\NormalTok{ q5, }\AttributeTok{est\_q95 =}\NormalTok{ q95)}

\NormalTok{  truth\_long }\OtherTok{\textless{}{-}}\NormalTok{ truth }\SpecialCharTok{|\textgreater{}}
\NormalTok{    dplyr}\SpecialCharTok{::}\FunctionTok{select}\NormalTok{(agent\_id, log\_c\_true, log\_lambda\_true) }\SpecialCharTok{|\textgreater{}}
    \FunctionTok{pivot\_longer}\NormalTok{(}\AttributeTok{cols =} \FunctionTok{c}\NormalTok{(log\_c\_true, log\_lambda\_true),}
                 \AttributeTok{names\_to =} \StringTok{"param"}\NormalTok{, }\AttributeTok{values\_to =} \StringTok{"truth"}\NormalTok{) }\SpecialCharTok{|\textgreater{}}
    \FunctionTok{mutate}\NormalTok{(}\AttributeTok{param =} \FunctionTok{if\_else}\NormalTok{(param }\SpecialCharTok{==} \StringTok{"log\_c\_true"}\NormalTok{, }\StringTok{"log\_c"}\NormalTok{, }\StringTok{"log\_lambda"}\NormalTok{))}

\NormalTok{  s }\SpecialCharTok{|\textgreater{}}
    \FunctionTok{left\_join}\NormalTok{(truth\_long, }\AttributeTok{by =} \FunctionTok{c}\NormalTok{(}\StringTok{"agent\_id"}\NormalTok{, }\StringTok{"param"}\NormalTok{)) }\SpecialCharTok{|\textgreater{}}
    \FunctionTok{mutate}\NormalTok{(}\AttributeTok{scenario =}\NormalTok{ scen)}
\NormalTok{\}) }\SpecialCharTok{|\textgreater{}}
  \FunctionTok{mutate}\NormalTok{(}\AttributeTok{scenario =} \FunctionTok{factor}\NormalTok{(scenario, }\AttributeTok{levels =}\NormalTok{ ml\_scenarios,}
                           \AttributeTok{labels =} \FunctionTok{c}\NormalTok{(}\StringTok{"Static"}\NormalTok{, }\StringTok{"Drift"}\NormalTok{, }\StringTok{"Contingent shift"}\NormalTok{)))}

\FunctionTok{ggplot}\NormalTok{(subj\_recovery, }\FunctionTok{aes}\NormalTok{(}\AttributeTok{x =}\NormalTok{ truth, }\AttributeTok{y =}\NormalTok{ est\_mean, }\AttributeTok{color =}\NormalTok{ scenario)) }\SpecialCharTok{+}
  \FunctionTok{geom\_abline}\NormalTok{(}\AttributeTok{intercept =} \DecValTok{0}\NormalTok{, }\AttributeTok{slope =} \DecValTok{1}\NormalTok{, }\AttributeTok{linetype =} \StringTok{"dashed"}\NormalTok{, }\AttributeTok{color =} \StringTok{"grey50"}\NormalTok{) }\SpecialCharTok{+}
  \FunctionTok{geom\_pointrange}\NormalTok{(}\FunctionTok{aes}\NormalTok{(}\AttributeTok{ymin =}\NormalTok{ est\_q5, }\AttributeTok{ymax =}\NormalTok{ est\_q95), }\AttributeTok{alpha =} \FloatTok{0.7}\NormalTok{) }\SpecialCharTok{+}
  \FunctionTok{facet\_grid}\NormalTok{(param }\SpecialCharTok{\textasciitilde{}}\NormalTok{ scenario, }\AttributeTok{scales =} \StringTok{"free"}\NormalTok{) }\SpecialCharTok{+}
  \FunctionTok{scale\_color\_viridis\_d}\NormalTok{(}\AttributeTok{option =} \StringTok{"plasma"}\NormalTok{, }\AttributeTok{end =} \FloatTok{0.85}\NormalTok{) }\SpecialCharTok{+}
  \FunctionTok{labs}\NormalTok{(}
    \AttributeTok{title    =} \StringTok{"Subject{-}level parameter recovery (with partial pooling)"}\NormalTok{,}
    \AttributeTok{subtitle =} \StringTok{"Posterior mean ± 90\% CI vs simulation truth"}\NormalTok{,}
    \AttributeTok{x =} \StringTok{"True value"}\NormalTok{, }\AttributeTok{y =} \StringTok{"Posterior estimate"}
\NormalTok{  ) }\SpecialCharTok{+}
  \FunctionTok{theme}\NormalTok{(}\AttributeTok{legend.position =} \StringTok{"none"}\NormalTok{)}
\end{Highlighting}
\end{Shaded}

\begin{center}
\includegraphics[width=0.8\linewidth,height=\textheight,keepaspectratio]{11-CategorizationModelsExemplars_files/figure-pdf/ch12_ml_decay_subject_recovery-1.pdf}
\end{center}

The \texttt{log\_c} row should be tight along the identity across
scenarios. The \texttt{log\_lambda} row is the diagnostic: scattered and
pulled toward the population mean under static (visible shrinkage when
an individual's data is uninformative), tighter along the diagonal under
drift, and tightest under contingent-shift --- consistent with the
single-subject recovery results from earlier in the chapter.

\subsection{Multilevel Simulation-Based Calibration
(SBC)}\label{multilevel-simulation-based-calibration-sbc}

Multilevel SBC is the most rigorous way to check that the hierarchical
priors and the non-centered parameterisation are working as intended.
Because fitting a multilevel model hundreds of times is computationally
intensive, we provide the code here as a template for students wishing
to perform an exhaustive validation of their own hierarchical
extensions.

\begin{Shaded}
\begin{Highlighting}[]
\CommentTok{\# 1. Generator: Draws population parameters from priors, then subject }
\CommentTok{\#    parameters from those hierarchies, and finally simulates trial data.}
\NormalTok{make\_ml\_decay\_sbc\_generator }\OtherTok{\textless{}{-}} \ControlFlowTok{function}\NormalTok{(scenario, }\AttributeTok{n\_subjects =} \DecValTok{5}\NormalTok{) \{}
  \FunctionTok{SBC\_generator\_function}\NormalTok{(}
    \ControlFlowTok{function}\NormalTok{() \{}
      \CommentTok{\# Draw population hyperparameters from the priors}
\NormalTok{      pop\_logit\_w\_mean    }\OtherTok{\textless{}{-}} \FunctionTok{rnorm}\NormalTok{(}\DecValTok{1}\NormalTok{, }\DecValTok{0}\NormalTok{, }\DecValTok{1}\NormalTok{)}
\NormalTok{      pop\_logit\_w\_sd      }\OtherTok{\textless{}{-}} \FunctionTok{rexp}\NormalTok{(}\DecValTok{1}\NormalTok{, }\DecValTok{2}\NormalTok{)}
\NormalTok{      pop\_log\_c\_mean      }\OtherTok{\textless{}{-}} \FunctionTok{rnorm}\NormalTok{(}\DecValTok{1}\NormalTok{, LOG\_C\_PRIOR\_MEAN, LOG\_C\_PRIOR\_SD)}
\NormalTok{      pop\_log\_c\_sd        }\OtherTok{\textless{}{-}} \FunctionTok{rexp}\NormalTok{(}\DecValTok{1}\NormalTok{, }\DecValTok{2}\NormalTok{)}
\NormalTok{      pop\_log\_lambda\_mean }\OtherTok{\textless{}{-}} \FunctionTok{rnorm}\NormalTok{(}\DecValTok{1}\NormalTok{, LOG\_LAMBDA\_PRIOR\_MEAN, LOG\_LAMBDA\_PRIOR\_SD)}
\NormalTok{      pop\_log\_lambda\_sd   }\OtherTok{\textless{}{-}} \FunctionTok{rexp}\NormalTok{(}\DecValTok{1}\NormalTok{, }\DecValTok{2}\NormalTok{)}
\NormalTok{      pop\_logit\_bias\_mean }\OtherTok{\textless{}{-}} \FunctionTok{rnorm}\NormalTok{(}\DecValTok{1}\NormalTok{, }\DecValTok{0}\NormalTok{, }\DecValTok{1}\NormalTok{)}
\NormalTok{      pop\_logit\_bias\_sd   }\OtherTok{\textless{}{-}} \FunctionTok{rexp}\NormalTok{(}\DecValTok{1}\NormalTok{, }\DecValTok{2}\NormalTok{)}
      
      \CommentTok{\# Simulate one multilevel dataset}
\NormalTok{      sim }\OtherTok{\textless{}{-}} \FunctionTok{simulate\_multilevel\_decay\_gcm}\NormalTok{(}
        \AttributeTok{n\_subjects          =}\NormalTok{ n\_subjects,}
        \AttributeTok{stimulus\_info       =}\NormalTok{ stimulus\_info,}
        \AttributeTok{n\_blocks            =} \DecValTok{2}\NormalTok{, }\CommentTok{\# Fewer blocks for faster SBC}
        \AttributeTok{scenario            =}\NormalTok{ scenario,}
        \AttributeTok{base\_seed           =} \FunctionTok{sample.int}\NormalTok{(}\FloatTok{1e6}\NormalTok{, }\DecValTok{1}\NormalTok{),}
        \AttributeTok{pop\_logit\_w\_mean    =}\NormalTok{ pop\_logit\_w\_mean,}
        \AttributeTok{pop\_logit\_w\_sd      =}\NormalTok{ pop\_logit\_w\_sd,}
        \AttributeTok{pop\_log\_c\_mean      =}\NormalTok{ pop\_log\_c\_mean,}
        \AttributeTok{pop\_log\_c\_sd        =}\NormalTok{ pop\_log\_c\_sd,}
        \AttributeTok{pop\_log\_lambda\_mean =}\NormalTok{ pop\_log\_lambda\_mean,}
        \AttributeTok{pop\_log\_lambda\_sd   =}\NormalTok{ pop\_log\_lambda\_sd,}
        \AttributeTok{pop\_logit\_bias\_mean =}\NormalTok{ pop\_logit\_bias\_mean,}
        \AttributeTok{pop\_logit\_bias\_sd   =}\NormalTok{ pop\_logit\_bias\_sd}
\NormalTok{      )}
      
      \CommentTok{\# Return parameters and the list for Stan}
      \FunctionTok{list}\NormalTok{(}
        \AttributeTok{parameters =} \FunctionTok{list}\NormalTok{(}
          \AttributeTok{pop\_logit\_w\_mean    =}\NormalTok{ pop\_logit\_w\_mean,}
          \AttributeTok{pop\_log\_c\_mean      =}\NormalTok{ pop\_log\_c\_mean,}
          \AttributeTok{pop\_log\_lambda\_mean =}\NormalTok{ pop\_log\_lambda\_mean,}
          \AttributeTok{pop\_logit\_bias\_mean =}\NormalTok{ pop\_logit\_bias\_mean}
\NormalTok{        ),}
        \AttributeTok{generated =} \FunctionTok{build\_ml\_stan\_data}\NormalTok{(sim}\SpecialCharTok{$}\NormalTok{data, stimulus\_info)}
\NormalTok{      )}
\NormalTok{    \}}
\NormalTok{  )}
\NormalTok{\}}

\CommentTok{\# 2. Backend: Uses the compiled multilevel Stan model}
\NormalTok{ml\_decay\_sbc\_backend }\OtherTok{\textless{}{-}} \FunctionTok{SBC\_backend\_cmdstan\_sample}\NormalTok{(}
\NormalTok{  mod\_gcm\_ml\_decay,}
  \AttributeTok{iter\_warmup =} \DecValTok{800}\NormalTok{, }\AttributeTok{iter\_sampling =} \DecValTok{800}\NormalTok{, }\AttributeTok{chains =} \DecValTok{2}
\NormalTok{)}

\CommentTok{\# 3. Compute SBC (Uncomment to run)}
\CommentTok{\# ml\_sbc\_results \textless{}{-} compute\_SBC(}
\CommentTok{\#   make\_ml\_decay\_sbc\_generator("drift", n\_subjects = 5),}
\CommentTok{\#   ml\_decay\_sbc\_backend,}
\CommentTok{\#   n\_sims = 50,}
\CommentTok{\#   cache\_location = here("simdata", "ch12\_ml\_decay\_sbc\_cache")}
\CommentTok{\# )}

\CommentTok{\# 4. Visualize Results}
\CommentTok{\# plot\_rank\_hist(ml\_sbc\_results)}
\CommentTok{\# plot\_ecdf\_diff(ml\_sbc\_results)}
\end{Highlighting}
\end{Shaded}

\subsection{What changes when you scale
up}\label{what-changes-when-you-scale-up}

Three things to keep in mind when adapting this model to real data:

\begin{enumerate}
\def\labelenumi{\arabic{enumi}.}
\tightlist
\item
  \textbf{Sample size enters twice.} With more subjects, the population
  SDs become better identified, which shrinks subject-level estimates
  more aggressively toward the population mean --- desirable when
  individual-level data is noisy, but worth checking visually.
\item
  \textbf{Sequential-trial costs scale linearly.} Doubling trials per
  subject doubles per-leapfrog cost; doubling subjects doubles it again.
  Going from this 10×64 toy to a 100×200 dataset is a \textasciitilde30×
  increase in wall time, not 30 different fits to manage.
\item
  \textbf{Static-scenario behaviour is the canary.} If
  \texttt{pop\_log\_lambda\_mean} and \texttt{pop\_log\_lambda\_sd}
  start to diverge or drag at \texttt{max\_treedepth} on real data, your
  task is probably structurally closer to ``static'' than you think ---
  there is not enough non-stationarity for \(\lambda\) to be identified
  at the population level. Reduce the model or accept that \(\lambda\)
  is pinned by the prior and report it as such.
\end{enumerate}

\subsection{GCM in Context: Strengths and
Limitations}\label{gcm-in-context-strengths-and-limitations}

\subsubsection{Strengths}\label{strengths}

\begin{itemize}
\tightlist
\item
  \textbf{Empirical support}: The GCM has successfully fit human
  categorization data across numerous experiments.
\item
  \textbf{Psychological plausibility}: The model's parameters correspond
  to psychologically meaningful constructs (attention,
  sensitivity/generalization, bias, and now forgetting).
\item
  \textbf{Flexibility}: It can account for a wide range of category
  structures and learning phenomena, including non-stationary
  environments.
\item
  \textbf{Individual differences}: The multilevel framework allows
  principled modeling of variations across people.
\end{itemize}

\subsubsection{Limitations}\label{limitations-1}

\begin{itemize}
\tightlist
\item
  \textbf{Memory requirements}: Storing all exemplars becomes
  computationally demanding for large numbers of stimuli. How memory is
  bounded over long timescales is not fully specified.
\item
  \textbf{Conflation of generalisation and decisional precision}: The
  standard GCM has no separate softmax temperature; sensitivity \(c\)
  doubles as a generalisation parameter and as a decisional precision
  parameter.
\item
  \textbf{Complex categories}: The standard GCM struggles with
  categories defined by abstract rules, relations, or logical
  combinations of features.
\item
  \textbf{Sequential structure and LOO}: Because the GCM is
  path-dependent, standard LOO-CV is technically invalid --- but this
  chapter implements LFO-CV (Bürkner et al., 2020) as the canonical
  sequential-model cross-validation.
\end{itemize}

In the next two chapters we explore alternative approaches --- prototype
models (Ch. 12) and rule-based models (Ch. 13) --- that address some of
these limitations while introducing their own strengths and challenges.

\subsection{Looking Ahead to Chapter 12: The Six-Step Pipeline in the
Prototype
Model}\label{looking-ahead-to-chapter-12-the-six-step-pipeline-in-the-prototype-model}

Chapter 12 applies the same six-step pipeline to a prototype-based
model. Before moving on, it is worth tracing explicitly how each stage
maps from the GCM to the Kalman filter prototype model --- because the
chapters are designed as a matched pair, and the contrasts at each stage
are the core theoretical content.

\begin{longtable}[]{@{}
  >{\raggedright\arraybackslash}p{(\linewidth - 4\tabcolsep) * \real{0.3333}}
  >{\raggedright\arraybackslash}p{(\linewidth - 4\tabcolsep) * \real{0.3333}}
  >{\raggedright\arraybackslash}p{(\linewidth - 4\tabcolsep) * \real{0.3333}}@{}}
\toprule\noalign{}
\begin{minipage}[b]{\linewidth}\raggedright
Pipeline Stage
\end{minipage} & \begin{minipage}[b]{\linewidth}\raggedright
GCM (this chapter)
\end{minipage} & \begin{minipage}[b]{\linewidth}\raggedright
Kalman Filter Prototype (Ch. 12)
\end{minipage} \\
\midrule\noalign{}
\endhead
\bottomrule\noalign{}
\endlastfoot
\textbf{1. Input representation} & Stimuli as points in continuous
feature space (height, position). & Same feature space; stimuli used as
Kalman observations. \\
\textbf{2. Attention} & Explicit attention weights \(w_f\) summing to 1;
estimated from data. & No explicit attention weights; covariance
structure implicitly weights dimensions. \\
\textbf{3. Intermediate representation} & \textbf{All} past exemplars
stored verbatim; memory grows trial by trial. & One prototype
\(\vec{\mu}_k\) per category, updated by each observation; memory is
constant-size. \\
\textbf{4. Evidential mechanism} & Similarity to each stored exemplar,
exponentially decaying with distance and (in the extended model) trial
age. & Mahalanobis-like distance from the current stimulus to each
prototype, scaled by the running covariance \(\Sigma_k\). \\
\textbf{5. Decision rule} & Luce choice axiom:
\(p(\text{cat}_1) = \frac{\beta \cdot \bar{\eta}_1}{\beta \cdot \bar{\eta}_1 + (1-\beta) \cdot \bar{\eta}_0}\).
& Same Luce choice axiom; the per-trial probability is computed from the
Kalman-filter likelihood ratio. \\
\textbf{6. Learning} & Accumulate: each trial's exemplar enters
permanent (but decaying) memory. No structural update. & Update: the
Kalman gain \(K\) shifts \(\vec{\mu}_k\) toward the new observation and
reduces \(\Sigma_k\); the prediction step (process noise \(Q\)) can
re-inflate \(\Sigma_k\) between trials. \\
\end{longtable}

The key tension is at \textbf{Stage 3 and Stage 6}. The GCM treats every
trial as permanently informative (modulo forgetting); the prototype
model treats each trial as an observation that refines a running
summary. As a result:

\begin{itemize}
\tightlist
\item
  The GCM is maximally sensitive to specific exemplars and their spatial
  arrangement. The prototype model is sensitive only to the category
  centroid and spread.
\item
  The GCM's non-stationarity mechanism (decay \(\lambda\)) operates
  \emph{retrospectively} --- old exemplars become less influential as
  time passes. The Kalman filter's mechanism (process noise \(q\))
  operates \emph{prospectively} --- future observations are expected to
  diverge from the current prototype.
\item
  Both mechanisms are phenomenologically similar under smooth drift.
  They diverge most sharply under abrupt shifts: the GCM gradually
  forgets old contradicting exemplars, while the Kalman filter
  re-inflates uncertainty after surprising observations and quickly
  reacquires the new prototype.
\end{itemize}

\begin{tcolorbox}[enhanced jigsaw, toprule=.15mm, toptitle=1mm, bottomtitle=1mm, leftrule=.75mm, rightrule=.15mm, colframe=quarto-callout-note-color-frame, left=2mm, colback=white, titlerule=0mm, coltitle=black, colbacktitle=quarto-callout-note-color!10!white, opacityback=0, title=\textcolor{quarto-callout-note-color}{\faInfo}\hspace{0.5em}{What to watch for when reading Chapter 12}, arc=.35mm, opacitybacktitle=0.6, breakable, bottomrule=.15mm]

\begin{enumerate}
\def\labelenumi{\arabic{enumi}.}
\tightlist
\item
  \textbf{The \(q\)/\(r\) identifiability problem.} The Kalman filter's
  two parameters (\(r\) = observation noise, \(q\) = process noise) are
  more tightly confounded under static data than the GCM's
  \(c\)/\(\lambda\) pair --- they both scale the \emph{same}
  decision-time likelihood rather than operating on orthogonal axes (see
  §``The \(c\)/\(\lambda\) Joint Posterior'' above). Ch. 12 introduces
  an explicit reparameterization (the SNR parameter \(\rho = q/r\)) to
  diagnose and partially resolve this.
\item
  \textbf{The scenario order is deliberately reversed.} Ch. 12 arranges
  scenarios as static → contingent shift → drift (increasing assumption
  match for the Kalman filter), while this chapter used static → drift →
  contingent shift (increasing \(\lambda\) preference). The three
  scenarios are identical; only the ordering logic differs.
\item
  \textbf{No attention weights.} Removing \(w\) and \(\beta\) from the
  parameter set makes the prototype model simpler to diagnose --- SBC
  rank histograms have two curves instead of four, and pair plots have
  one panel instead of six. That simplicity is \emph{why} Ch. 12 is the
  right place to introduce LFO-CV in detail: the algorithm is most
  legible when the model is lean.
\end{enumerate}

\end{tcolorbox}

\begin{center}\rule{0.5\linewidth}{0.5pt}\end{center}

\bookmarksetup{startatroot}

\chapter{Prototype Models: Learning Category
Averages}\label{prototype-models-learning-category-averages}

\begin{quote}
\textbf{📍 Where we are in the Bayesian modeling workflow:} Chapter 11
built the GCM --- an exemplar model that stores every encountered
stimulus and decides by summed similarity to memory. This chapter
introduces a structurally different account: the \textbf{prototype
model}, which abstracts each category to a single running estimate (its
mean and uncertainty) rather than storing individual instances. We
implement prototype learning as a multivariate Kalman filter and then
add a \textbf{Stan parameter-estimation model} for the prototype
architecture.

The full validation pipeline from Ch. 5 --- prior predictive check,
Pathfinder initialization, mandatory MCMC diagnostic battery, posterior
predictive check, prior--posterior update plot, randomised LOO-PIT,
LFO-CV, precision analysis, prior sensitivity, and SBC --- is applied to
this new model. Crucially, the \textbf{LFO-CV machinery implemented in
Chapter 11} (\texttt{psis\_lfo\_gcm}) is generalised here as
\texttt{psis\_lfo\_kalman} and applied to the prototype model. The
chapter then extends to a \textbf{multilevel prototype model} with
population-level priors on \texttt{log\_r}, per-subject NCP offsets, and
the full multilevel SBC battery (population and individual calibration).
\end{quote}

\section{Rethinking Category Representation: From Examples to
Averages}\label{rethinking-category-representation-from-examples-to-averages}

In the previous chapter, we explored the Generalized Context Model
(GCM), an exemplar model where every encountered example is stored in
memory. This approach is powerful but raises questions about cognitive
efficiency. Do we really store every single instance we encounter?

Consider learning about common categories like ``birds'' or ``chairs''.
While specific examples matter, we also seem to develop a general sense
of what constitutes a typical member. Prototype theory offers an
alternative perspective: instead of storing individual exemplars, the
mind abstracts a summary representation --- a prototype --- for each
category. This prototype often represents the central tendency or
``average'' member.

\section{Core Ideas of Prototype
Models}\label{core-ideas-of-prototype-models}

\begin{itemize}
\tightlist
\item
  \textbf{Abstraction, Not Storage}: The primary representation is a
  summary (the prototype), not a collection of instances.
\item
  \textbf{Cognitive Economy}: Storing one prototype per category is
  potentially much more memory-efficient than storing numerous
  exemplars.
\item
  \textbf{Similarity to Prototype}: New items are categorized based on
  their similarity to the stored prototypes of different categories.
\item
  \textbf{Typicality}: Items closer to a category's prototype are judged
  as more ``typical'' members and are often categorized faster and more
  accurately. This contrasts with exemplar models, which emphasize
  detailed memory for specific instances.
\end{itemize}

\begin{longtable}[]{@{}
  >{\raggedright\arraybackslash}p{(\linewidth - 4\tabcolsep) * \real{0.3333}}
  >{\raggedright\arraybackslash}p{(\linewidth - 4\tabcolsep) * \real{0.3333}}
  >{\raggedright\arraybackslash}p{(\linewidth - 4\tabcolsep) * \real{0.3333}}@{}}
\toprule\noalign{}
\begin{minipage}[b]{\linewidth}\raggedright
Feature
\end{minipage} & \begin{minipage}[b]{\linewidth}\raggedright
Prototype Models
\end{minipage} & \begin{minipage}[b]{\linewidth}\raggedright
Exemplar Models (GCM)
\end{minipage} \\
\midrule\noalign{}
\endhead
\bottomrule\noalign{}
\endlastfoot
Representation & Abstract summary (e.g., average) & Collection of
individual examples \\
Memory Cost & Low (one prototype per category) & High (potentially all
examples) \\
Key Information & Central tendency \& variance & Specific instances \&
their labels \\
Sensitivity & Less sensitive to specific examples & Highly sensitive to
specific examples \\
\end{longtable}

\begin{tcolorbox}[enhanced jigsaw, toprule=.15mm, toptitle=1mm, bottomtitle=1mm, leftrule=.75mm, rightrule=.15mm, colframe=quarto-callout-note-color-frame, left=2mm, colback=white, titlerule=0mm, coltitle=black, colbacktitle=quarto-callout-note-color!10!white, opacityback=0, title=\textcolor{quarto-callout-note-color}{\faInfo}\hspace{0.5em}{Historical Context: From Typicality Effects to Dynamic Prototypes}, arc=.35mm, opacitybacktitle=0.6, breakable, bottomrule=.15mm]

The prototype approach to categorization did not emerge fully formed. It
arose as the solution to a specific failure mode of the dominant view
--- and then itself had to be rescued from static simplicity by
importing engineering mathematics.

\subsubsection{The Classical View and Its
Discontents}\label{the-classical-view-and-its-discontents}

Through the 1950s and into the 1960s, both philosophers and
psychologists took it for granted that concepts had \emph{defining
features}: necessary and sufficient conditions that every member must
satisfy. A bachelor is an unmarried adult male --- full stop.
Categorization, on this view, was logical: check the features, apply the
rule. The model is appealingly crisp, but Wittgenstein had already noted
the trouble in \emph{Philosophical Investigations} (1953): most ordinary
categories --- \emph{game}, \emph{furniture}, \emph{bird} --- resist
such definitions. Family resemblance, not shared essence, seemed to be
the operative structure.

The empirical broadside came from Eleanor Rosch (E. H. Rosch 1973; E.
Rosch 1975; E. Rosch et al. 1976). In a series of influential
experiments she documented \textbf{typicality effects}: robins are
judged faster and more confidently as birds than penguins are; desk
chairs are more chair-like than beanbags. Categories showed graded
membership rather than all-or-nothing inclusion. Rosch also identified a
privileged \textbf{basic level} --- the level of abstraction (dog, not
animal and not poodle) at which objects are most efficiently perceived,
named, and remembered. Neither observation is compatible with the
classical view, which predicts uniform membership within a category and
no principled distinction across levels.

\subsubsection{The Prototype as Psychological
Mechanism}\label{the-prototype-as-psychological-mechanism}

Rosch's work raised the question of \emph{representation}: if not a
rule, what mental structure supports graded membership? The answer on
offer was a \textbf{prototype} --- a summary representation encoding the
central tendency of encountered category members. New objects are
categorized by their distance from stored prototypes; items close to the
prototype are typical, items distant are atypical.

Formal prototype models predated Rosch's program. Posner and Keele
(1968) showed that subjects who had studied distorted dot patterns
without ever seeing the prototype nonetheless recognized the prototype
at the end as highly familiar --- abstraction was occurring implicitly.
Reed (1972) implemented an explicit prototype-matching classifier and
showed it outperformed several alternatives on face-classification data.
These early models computed a \textbf{static average} from training
examples: the prototype was fixed after learning and then used as a
lookup.

\subsubsection{Static Prototypes and the Exemplar
Challenge}\label{static-prototypes-and-the-exemplar-challenge}

The clean elegance of static prototype models made them obvious targets
for empirical attack. Medin and Schaffer (1978) showed that, on certain
category structures, the Context Model (a pure exemplar model) fit human
data significantly better than a prototype model could. The critical
cases were \textbf{exception items}: an atypical member of one category
that closely resembles the other category's prototype. Exemplar models
accommodate exceptions naturally --- the deviant item is stored
individually. A static prototype, by definition, assimilates the deviant
item into the average, smoothing away the very information that drives
behavior.

Through the 1980s and early 1990s, the prototype-versus-exemplar debate
generated substantial empirical literature. The consensus, codified in
Smith and Minda (1998), was that neither account was universally
superior: prototype-like effects dominated in some conditions (large
training sets, brief study times, transfer to novel stimuli) while
exemplar-like effects dominated in others (small training sets, extended
learning, highly distinctive items). The field was left with a
theoretical draw that demanded either a hybrid architecture or a shift
in the question being asked.

\subsubsection{From Static Averaging to Bayesian
Updating}\label{from-static-averaging-to-bayesian-updating}

The decisive conceptual move was to ask not \emph{what is stored} but
\emph{how the representation is updated}. A static average treats all
observed exemplars symmetrarily regardless of when they were
encountered. This is inconsistent with the well-documented finding that
recent experiences carry more weight than remote ones (Gluck and Bower
1988; Estes 1994). It is also inconsistent with the ability to track
categories that drift over time --- the clothing you learn is
``fashionable'' today may not be fashionable in five years.

The Kalman filter (Kalman 1960), developed originally for aerospace
navigation, provided an off-the-shelf solution. It is the
\textbf{minimum-variance unbiased estimator} for a linear-Gaussian
dynamical system: at each step it computes exactly how much weight to
place on the new observation versus the accumulated prior, using the
relative magnitudes of measurement noise \(R\) and process noise \(Q\).
When \(Q = 0\) the prototype is assumed static and learning rate decays
monotonically --- recovering the classical static average as a special
case. When \(Q > 0\) the prototype can drift and the steady-state Kalman
gain remains positive, allowing perpetual sensitivity to new
observations.

Applications of Kalman-style updating to human learning appear across
multiple domains: classical conditioning (Dayan, Kakade, and Montague
2000), interval timing, and reward prediction. The prototype-learning
context is a natural fit because the Kalman filter formalises the very
trade-off prototype theory requires: balance between the accumulated
representation (prior) and the informational content of the newest
instance (likelihood). The two-parameter family \((r, q)\) instantiates
the entire spectrum from pure exemplar-like responsiveness (\(r\) small,
steep gain) to heavily smoothed prototype abstraction (\(r\) large,
shallow gain, slow forgetting).

\subsubsection{The Bayesian Identity}\label{the-bayesian-identity}

The Kalman filter is not merely an engineering heuristic applied to
psychology: it is the exact Bayesian posterior for a Gaussian
state-space model. The prototype at trial \(t\) is the posterior mean
\(\mathbb{E}[\mu_t \mid x_{1:t}]\); the covariance \(\Sigma_t\) is the
posterior variance. Recognising this identity connects the prototype
model to the broader programme of \textbf{rational analysis} (Anderson
1991) and \textbf{Bayesian models of cognition} (Tenenbaum et al. 2011):
the learner is performing optimal inference about a latent category
structure given noisy, sequentially arriving evidence. The parameters
\(r\) and \(q\) then carry clear psychological interpretations ---
observation noise (how variable are category members around the
prototype?) and process noise (how rapidly does the category itself
change?) --- rather than being purely phenomenological fitting
coefficients.

This framing also clarifies the key difference from the GCM: where the
GCM stores the full history and retrieves it at decision time, the
Kalman prototype model \textbf{compresses} the history into sufficient
statistics (mean and covariance) that are updated recursively. The
compression is lossless when the generative model is Gaussian and linear
--- which is why recovery and SBC are cleaner here than in most
cognitive models of comparable scope.

\end{tcolorbox}

\section{Modeling Dynamic Prototypes: The Kalman
Filter}\label{modeling-dynamic-prototypes-the-kalman-filter}

Early prototype models often calculated a static average. However, human
learning is dynamic; our understanding evolves with experience. How can
we model a prototype that updates incrementally as new examples are
encountered?

The Kalman filter provides a suitable mathematical framework. Originally
developed for tracking physical systems amidst noise, it's well-suited
for modeling prototype learning because it allows us to:

\begin{itemize}
\tightlist
\item
  \textbf{Track an Estimate}: The category prototype (its average
  feature values).
\item
  \textbf{Represent Uncertainty}: Maintain not just the prototype's
  location but also our uncertainty about it.
\item
  \textbf{Update Incrementally}: Refine the prototype with each new
  example.
\item
  \textbf{Balance Old and New}: Optimally combine the current prototype
  (prior belief) with the new example (evidence).
\item
  \textbf{Implement an Adaptive Learning Rate}: Automatically adjust how
  much influence a new example has (the ``Kalman gain'') based on
  current uncertainty. Learning is faster when uncertainty is high and
  slows as confidence grows.
\end{itemize}

\subsection{The Kalman Filter Logic: Updating an
Estimate}\label{the-kalman-filter-logic-updating-an-estimate}

Imagine tracking a prototype based on a single feature (e.g., average
height). The Kalman filter maintains your Current State. At any moment,
you have two key pieces of information:

\begin{itemize}
\tightlist
\item
  \textbf{Your Best Guess (\(\mu\)):} This is your current estimate of
  the prototype's average height.
\item
  \textbf{Your Confidence (\(\sigma^2\)):} This isn't just \emph{if}
  you're confident, but \emph{how much}. It's the \emph{variance} around
  your guess. A small variance means you're quite sure; a large variance
  means you're very unsure.
\end{itemize}

When a new category member (observation \(x\)) arrives, you measure its
height.

\begin{itemize}
\tightlist
\item
  \textbf{The Catch:} Your measurement tool (or the example itself)
  isn't perfect. There's some inherent ``noise'' or variability in how
  examples relate to the true prototype. We represent the uncertainty of
  this \emph{measurement process} with another variance, often called
  \(R\).
\end{itemize}

\textbf{The Core Problem:} You have your guess (with
\(\sigma^2_{prev}\)) and a new measurement (with \(R\)). How do you
combine them to get a \emph{new, better guess} and an \emph{updated
level of uncertainty}?

\textbf{Enter the Kalman Gain (\(K\)): The ``Trust'' Knob}

The Kalman filter calculates the Kalman Gain \(K\), which tells you how
much to trust the \emph{new measurement} compared to your \emph{current
guess}:

\[K = \frac{\sigma^2_{prev}}{\sigma^2_{prev} + R}\]

If \(\sigma^2_{prev}\) is large relative to \(R\), then \(K \to 1\):
trust the measurement almost completely. If \(\sigma^2_{prev}\) is small
relative to \(R\), then \(K \to 0\): stick closer to your previous
guess.

\textbf{Updating Your Guess:}

You nudge your old guess towards the new measurement. The size of the
nudge is determined by the Kalman Gain \(K\):

\[\mu_{new} = \mu_{prev} + K \cdot (x - \mu_{prev})\]

The term \((x - \mu_{prev})\) is the ``prediction error'' or
``innovation'' --- how much the new data differs from what you expected.

\textbf{Updating Your Confidence (Uncertainty):}

After incorporating new information, uncertainty decreases:

\[\sigma^2_{new} = (1 - K) \cdot \sigma^2_{prev}\]

Since \(K \in (0, 1)\), multiplying by \((1 - K)\) always shrinks
uncertainty. The more you trusted the measurement (higher \(K\)), the
more your uncertainty shrinks.

\textbf{Extending to drifting targets.} The scalar Kalman filter above
assumes a static target: uncertainty monotonically decreases because
there is no mechanism for it to grow back. The prototype model in this
chapter adds a \textbf{process noise term} \(Q\). Before each
measurement update, a \emph{prediction step} is applied:

\[\sigma^2_{k|k-1} = \sigma^2_{k-1|k-1} + Q\]

This increases uncertainty between observations, reflecting the
possibility that the true prototype has shifted since the last trial.
The subsequent measurement update then decreases uncertainty via the
Kalman gain, exactly as before. At steady state, the two forces balance
and uncertainty converges to a fixed value rather than shrinking to
zero. When \(Q = 0\) the model reduces to the static-target filter. We
treat the scalar \(q\) (the diagonal of \(Q = q \cdot I\)) as a free
parameter to be estimated from data, alongside the observation noise
\(r\).

\subsection{Extending to Multiple Features: Multivariate Kalman
Filter}\label{extending-to-multiple-features-multivariate-kalman-filter}

Our stimuli have two features (height and position). We need a
multivariate version:

\begin{itemize}
\tightlist
\item
  \textbf{Your Current State (Vectors and Matrices):}

  \begin{itemize}
  \tightlist
  \item
    \textbf{Your Best Guess (\(\vec{\mu}\)):} Now a \emph{vector}
    containing your best guess for each feature.
  \item
    \textbf{Your Confidence (\(\Sigma\)):} Now a \emph{covariance
    matrix}. The diagonal elements give the variance for each feature
    individually; the off-diagonal elements give their covariance ---
    for example, whether taller prototypes are also typically further to
    the right.
  \end{itemize}
\item
  \textbf{A New Hint Arrives:} Your new observation (\(\vec{x}\)) is
  also a vector.

  \begin{itemize}
  \tightlist
  \item
    \textbf{Measurement Noise (\(R\)):} Also a covariance matrix. Often
    simplified to a diagonal matrix, assuming measurement errors on
    different features are independent.
  \end{itemize}
\item
  \textbf{The Logic Stays the Same, the Math Uses Matrices:}

  \begin{itemize}
  \tightlist
  \item
    \textbf{Kalman Gain (\(K\)):} A \emph{matrix}:
    \(K = \Sigma_{prev} (\Sigma_{prev} + R)^{-1}\)
  \item
    \textbf{Updating Your Guess (\(\vec{\mu}\)):}
    \(\vec{\mu}_{new} = \vec{\mu}_{prev} + K (\vec{x} - \vec{\mu}_{prev})\)
  \item
    \textbf{Updating Your Confidence (\(\Sigma\)):} We use the
    numerically stable Joseph form:
    \(\Sigma_{new} = (I - K)\Sigma_{prev}(I - K)^T + KRK^T\). While it
    looks complex, the intuition is the same as the scalar case:
    uncertainty decreases after each update. The Joseph form is
    preferred because it guarantees the result remains symmetric and
    positive semi-definite even in the presence of floating-point
    rounding errors.
  \end{itemize}
\end{itemize}

\subsection{In a Nutshell}\label{in-a-nutshell}

The Kalman filter is a recursive algorithm that provides the
mathematically optimal way to estimate the state of a system (like the
features of a prototype) given a series of noisy measurements. It works
by:

\begin{enumerate}
\def\labelenumi{\arabic{enumi}.}
\tightlist
\item
  Maintaining an estimate (\(\vec{\mu}\)) and uncertainty (\(\Sigma\)).
\item
  Calculating a Kalman Gain \(K\) that balances trust in the current
  estimate against trust in the new measurement.
\item
  Updating the estimate by moving it towards the new measurement,
  weighted by \(K\).
\item
  Updating (reducing) the uncertainty based on the information gained.
\end{enumerate}

The multivariate version simply uses vectors and matrices to do this for
multiple, potentially correlated variables simultaneously.

\subsection{Implementing the Multivariate Update in
R}\label{implementing-the-multivariate-update-in-r}

\begin{Shaded}
\begin{Highlighting}[]
\CommentTok{\# One update step of a multivariate Kalman filter.}
\CommentTok{\# Returns updated mean vector, covariance matrix, and Kalman gain.}
\NormalTok{multivariate\_kalman\_update }\OtherTok{\textless{}{-}} \ControlFlowTok{function}\NormalTok{(mu\_prev,      }\CommentTok{\# previous mean vector}
\NormalTok{                                       sigma\_prev,    }\CommentTok{\# previous covariance matrix}
\NormalTok{                                       observation,   }\CommentTok{\# observed feature vector}
\NormalTok{                                       r\_matrix) \{    }\CommentTok{\# observation noise matrix}
\NormalTok{  mu\_prev     }\OtherTok{\textless{}{-}} \FunctionTok{as.numeric}\NormalTok{(mu\_prev)}
\NormalTok{  observation }\OtherTok{\textless{}{-}} \FunctionTok{as.numeric}\NormalTok{(observation)}
\NormalTok{  sigma\_prev  }\OtherTok{\textless{}{-}} \FunctionTok{as.matrix}\NormalTok{(sigma\_prev)}
\NormalTok{  r\_matrix    }\OtherTok{\textless{}{-}} \FunctionTok{as.matrix}\NormalTok{(r\_matrix)}

\NormalTok{  n\_dim }\OtherTok{\textless{}{-}} \FunctionTok{length}\NormalTok{(mu\_prev)}
\NormalTok{  I     }\OtherTok{\textless{}{-}} \FunctionTok{diag}\NormalTok{(n\_dim)}

  \CommentTok{\# Kalman gain: K = Sigma\_prev * (Sigma\_prev + R)\^{}\{{-}1\}}
\NormalTok{  S     }\OtherTok{\textless{}{-}}\NormalTok{ sigma\_prev }\SpecialCharTok{+}\NormalTok{ r\_matrix}
\NormalTok{  S\_inv }\OtherTok{\textless{}{-}} \FunctionTok{tryCatch}\NormalTok{(}
    \FunctionTok{solve}\NormalTok{(S),}
    \AttributeTok{error =} \ControlFlowTok{function}\NormalTok{(e) \{}
      \FunctionTok{warning}\NormalTok{(}\StringTok{"Matrix inversion failed; using pseudo{-}inverse."}\NormalTok{, }\AttributeTok{call. =} \ConstantTok{FALSE}\NormalTok{)}
\NormalTok{      MASS}\SpecialCharTok{::}\FunctionTok{ginv}\NormalTok{(S)}
\NormalTok{    \}}
\NormalTok{  )}
\NormalTok{  K }\OtherTok{\textless{}{-}}\NormalTok{ sigma\_prev }\SpecialCharTok{\%*\%}\NormalTok{ S\_inv}

  \CommentTok{\# Update mean}
\NormalTok{  innovation }\OtherTok{\textless{}{-}}\NormalTok{ observation }\SpecialCharTok{{-}}\NormalTok{ mu\_prev}
\NormalTok{  mu\_new     }\OtherTok{\textless{}{-}} \FunctionTok{as.numeric}\NormalTok{(mu\_prev }\SpecialCharTok{+}\NormalTok{ K }\SpecialCharTok{\%*\%}\NormalTok{ innovation)}

  \CommentTok{\# Update covariance — Joseph form for numerical stability}
\NormalTok{  IK        }\OtherTok{\textless{}{-}}\NormalTok{ I }\SpecialCharTok{{-}}\NormalTok{ K}
\NormalTok{  sigma\_new }\OtherTok{\textless{}{-}}\NormalTok{ IK }\SpecialCharTok{\%*\%}\NormalTok{ sigma\_prev }\SpecialCharTok{\%*\%} \FunctionTok{t}\NormalTok{(IK) }\SpecialCharTok{+}\NormalTok{ K }\SpecialCharTok{\%*\%}\NormalTok{ r\_matrix }\SpecialCharTok{\%*\%} \FunctionTok{t}\NormalTok{(K)}

  \CommentTok{\# Enforce symmetry (floating{-}point can introduce tiny asymmetry)}
\NormalTok{  sigma\_new }\OtherTok{\textless{}{-}}\NormalTok{ (sigma\_new }\SpecialCharTok{+} \FunctionTok{t}\NormalTok{(sigma\_new)) }\SpecialCharTok{/} \DecValTok{2}

  \FunctionTok{list}\NormalTok{(}\AttributeTok{mu =}\NormalTok{ mu\_new, }\AttributeTok{sigma =}\NormalTok{ sigma\_new, }\AttributeTok{k =}\NormalTok{ K)}
\NormalTok{\}}
\end{Highlighting}
\end{Shaded}

This function performs a single update for one category's prototype when
it observes a new member.

\begin{center}\rule{0.5\linewidth}{0.5pt}\end{center}

We are fitting a \emph{two-parameter} model: \(r\) controls both Kalman
gain and decisional precision; \(q\) controls prototype drift speed and
steady-state uncertainty. Both are inferred from data on the log scale.
There is no category bias, fixed initial means, fixed initial
covariances, and diagonal noise structures for both \(R\) and \(Q\). The
Stan implementation in §``The Prototype Model in Stan'' implements
exactly this.

\begin{center}\rule{0.5\linewidth}{0.5pt}\end{center}

\section{Building the Full Categorization
Model}\label{building-the-full-categorization-model}

We now integrate this update mechanism into a full agent that learns
prototypes for two categories and makes categorization decisions.

\subsection{Initialization}\label{initialization}

\begin{itemize}
\tightlist
\item
  We need separate prototypes (\(\mu_0, \Sigma_0\) and
  \(\mu_1, \Sigma_1\)) for Category 0 and Category 1.
\item
  We start with high uncertainty (large initial \(\Sigma\)) and
  uninformative means at the center of the feature space.
\item
  We define the observation noise matrix \(R = r_\text{val} \cdot I\).
  The parameter \(r_\text{val}\) represents how much variability we
  assume exists within categories or in our perception of stimuli.
\item
  We define the process noise matrix \(Q = q_\text{val} \cdot I\). The
  parameter \(q_\text{val}\) controls how much the category prototype
  may drift between successive observations. Small \(q\) means
  near-static categories; large \(q\) means rapidly shifting ones.
\end{itemize}

\subsection{Categorization Decision}\label{categorization-decision}

How does the model decide the category for a new stimulus \(\vec{x}\)?
It compares \(\vec{x}\) to both prototypes, considering their current
means (\(\vec{\mu}\)) and uncertainties (\(\Sigma\)). We use the
probability density of the multivariate normal distribution:

\[P(\vec{x} \mid \text{Category } C) \propto \exp\left(-\frac{1}{2}(\vec{x} - \vec{\mu}_C)^T (\Sigma_C + R)^{-1} (\vec{x} - \vec{\mu}_C)\right)\]

The probability of choosing Category 1 is then calculated using the Luce
choice rule (softmax):

\[P(\text{Choose Cat 1} \mid \vec{x}) = \frac{P(\vec{x} \mid \text{Category } 1)}{P(\vec{x} \mid \text{Category } 0) + P(\vec{x} \mid \text{Category } 1)}\]

We work with log-probabilities for numerical stability. Note that the
same \(r\) that drives the Kalman update also enters here through
\((\Sigma_C + R)^{-1}\) --- there is no separate softmax temperature.

\subsection{The Learning Loop}\label{the-learning-loop}

The agent processes trials sequentially:

\begin{enumerate}
\def\labelenumi{\arabic{enumi}.}
\tightlist
\item
  \textbf{Predict}: Add \(Q\) to both category covariances ---
  \(\Sigma_{0} \mathrel{+}= Q\), \(\Sigma_{1} \mathrel{+}= Q\) ---
  reflecting that each prototype may have drifted since the last
  observation.
\item
  \textbf{Decide}: Calculate \(P(\text{Choose Cat 1} \mid \vec{x}_i)\)
  based on the post-prediction \(\mu_0, \Sigma_0, \mu_1, \Sigma_1\).
\item
  \textbf{Respond}: Generate a response.
\item
  \textbf{Receive feedback}: Observe true category \(y_i\).
\item
  \textbf{Update}: Apply the Kalman measurement update to the correct
  category only.
\end{enumerate}

\subsection{R Implementation of the Prototype
Agent}\label{r-implementation-of-the-prototype-agent}

The function now accepts \texttt{subject\_seed} and calls
\texttt{make\_subject\_schedule} internally --- consistent with the
path-dependence fix from Ch. 11. Each agent's learning trajectory
depends on the order in which stimuli are encountered, and different
randomizations produce genuinely different trajectories.

\begin{Shaded}
\begin{Highlighting}[]
\CommentTok{\# Prototype model agent using Kalman filter with process noise for categorization.}
\CommentTok{\# r\_value   : observation noise variance (scalar \textgreater{} 0)}
\CommentTok{\# q\_value   : process noise variance / prototype drift rate (scalar \textgreater{}= 0)}
\CommentTok{\# obs       : matrix of observations (trials x features)}
\CommentTok{\# cat\_one   : vector of true category labels (0 or 1) for feedback}
\CommentTok{\# Returns:  tibble with prob\_cat1 and sim\_response}
\NormalTok{prototype\_kalman }\OtherTok{\textless{}{-}} \ControlFlowTok{function}\NormalTok{(r\_value,}
\NormalTok{                             q\_value,}
\NormalTok{                             obs,}
\NormalTok{                             cat\_one,}
                             \AttributeTok{initial\_mu         =} \ConstantTok{NULL}\NormalTok{,}
                             \AttributeTok{initial\_sigma\_diag =} \FloatTok{10.0}\NormalTok{,}
                             \AttributeTok{quiet              =} \ConstantTok{TRUE}\NormalTok{) \{}
\NormalTok{  n\_trials   }\OtherTok{\textless{}{-}} \FunctionTok{nrow}\NormalTok{(obs)}
\NormalTok{  n\_features }\OtherTok{\textless{}{-}} \FunctionTok{ncol}\NormalTok{(obs)}

  \ControlFlowTok{if}\NormalTok{ (}\FunctionTok{is.null}\NormalTok{(initial\_mu)) \{}
\NormalTok{    mu0\_init }\OtherTok{\textless{}{-}} \FunctionTok{rep}\NormalTok{(}\FloatTok{2.5}\NormalTok{, n\_features)}
\NormalTok{    mu1\_init }\OtherTok{\textless{}{-}} \FunctionTok{rep}\NormalTok{(}\FloatTok{2.5}\NormalTok{, n\_features)}
\NormalTok{  \} }\ControlFlowTok{else}\NormalTok{ \{}
\NormalTok{    mu0\_init }\OtherTok{\textless{}{-}}\NormalTok{ initial\_mu[[}\DecValTok{1}\NormalTok{]]}
\NormalTok{    mu1\_init }\OtherTok{\textless{}{-}}\NormalTok{ initial\_mu[[}\DecValTok{2}\NormalTok{]]}
\NormalTok{  \}}

\NormalTok{  prototype\_cat\_0 }\OtherTok{\textless{}{-}} \FunctionTok{list}\NormalTok{(}\AttributeTok{mu =}\NormalTok{ mu0\_init, }\AttributeTok{sigma =} \FunctionTok{diag}\NormalTok{(initial\_sigma\_diag, n\_features))}
\NormalTok{  prototype\_cat\_1 }\OtherTok{\textless{}{-}} \FunctionTok{list}\NormalTok{(}\AttributeTok{mu =}\NormalTok{ mu1\_init, }\AttributeTok{sigma =} \FunctionTok{diag}\NormalTok{(initial\_sigma\_diag, n\_features))}
\NormalTok{  r\_matrix        }\OtherTok{\textless{}{-}} \FunctionTok{diag}\NormalTok{(r\_value, n\_features)}
\NormalTok{  q\_matrix        }\OtherTok{\textless{}{-}} \FunctionTok{diag}\NormalTok{(q\_value, n\_features)}

\NormalTok{  response\_probs }\OtherTok{\textless{}{-}} \FunctionTok{numeric}\NormalTok{(n\_trials)}

\NormalTok{  log\_sum\_exp }\OtherTok{\textless{}{-}} \ControlFlowTok{function}\NormalTok{(v) \{}
\NormalTok{    m }\OtherTok{\textless{}{-}} \FunctionTok{max}\NormalTok{(v)}
\NormalTok{    m }\SpecialCharTok{+} \FunctionTok{log}\NormalTok{(}\FunctionTok{sum}\NormalTok{(}\FunctionTok{exp}\NormalTok{(v }\SpecialCharTok{{-}}\NormalTok{ m)))}
\NormalTok{  \}}

  \ControlFlowTok{for}\NormalTok{ (i }\ControlFlowTok{in} \FunctionTok{seq\_len}\NormalTok{(n\_trials)) \{}
    \ControlFlowTok{if}\NormalTok{ (}\SpecialCharTok{!}\NormalTok{quiet }\SpecialCharTok{\&\&}\NormalTok{ i }\SpecialCharTok{\%\%} \DecValTok{20} \SpecialCharTok{==} \DecValTok{0}\NormalTok{) }\FunctionTok{cat}\NormalTok{(}\StringTok{"Trial"}\NormalTok{, i, }\StringTok{"}\SpecialCharTok{\textbackslash{}n}\StringTok{"}\NormalTok{)}

\NormalTok{    current\_obs }\OtherTok{\textless{}{-}} \FunctionTok{as.numeric}\NormalTok{(obs[i, ])}

    \CommentTok{\# ── Prediction step: add process noise to both prototypes ──────────────}
    \CommentTok{\# Reflects potential drift in category location since last trial.}
    \CommentTok{\# When q\_value = 0 this reduces to the static{-}target filter.}
\NormalTok{    prototype\_cat\_0}\SpecialCharTok{$}\NormalTok{sigma }\OtherTok{\textless{}{-}}\NormalTok{ prototype\_cat\_0}\SpecialCharTok{$}\NormalTok{sigma }\SpecialCharTok{+}\NormalTok{ q\_matrix}
\NormalTok{    prototype\_cat\_1}\SpecialCharTok{$}\NormalTok{sigma }\OtherTok{\textless{}{-}}\NormalTok{ prototype\_cat\_1}\SpecialCharTok{$}\NormalTok{sigma }\SpecialCharTok{+}\NormalTok{ q\_matrix}

    \CommentTok{\# ── Decision ─────────────────────────────────────────────────────────────}
\NormalTok{    cov\_cat\_0 }\OtherTok{\textless{}{-}}\NormalTok{ prototype\_cat\_0}\SpecialCharTok{$}\NormalTok{sigma }\SpecialCharTok{+}\NormalTok{ r\_matrix}
\NormalTok{    cov\_cat\_1 }\OtherTok{\textless{}{-}}\NormalTok{ prototype\_cat\_1}\SpecialCharTok{$}\NormalTok{sigma }\SpecialCharTok{+}\NormalTok{ r\_matrix}

\NormalTok{    log\_prob\_0 }\OtherTok{\textless{}{-}} \FunctionTok{tryCatch}\NormalTok{(}
\NormalTok{      mvtnorm}\SpecialCharTok{::}\FunctionTok{dmvnorm}\NormalTok{(current\_obs, }\AttributeTok{mean =}\NormalTok{ prototype\_cat\_0}\SpecialCharTok{$}\NormalTok{mu,}
                       \AttributeTok{sigma =}\NormalTok{ cov\_cat\_0, }\AttributeTok{log =} \ConstantTok{TRUE}\NormalTok{),}
      \AttributeTok{error =} \ControlFlowTok{function}\NormalTok{(e) }\SpecialCharTok{{-}}\ConstantTok{Inf}
\NormalTok{    )}
\NormalTok{    log\_prob\_1 }\OtherTok{\textless{}{-}} \FunctionTok{tryCatch}\NormalTok{(}
\NormalTok{      mvtnorm}\SpecialCharTok{::}\FunctionTok{dmvnorm}\NormalTok{(current\_obs, }\AttributeTok{mean =}\NormalTok{ prototype\_cat\_1}\SpecialCharTok{$}\NormalTok{mu,}
                       \AttributeTok{sigma =}\NormalTok{ cov\_cat\_1, }\AttributeTok{log =} \ConstantTok{TRUE}\NormalTok{),}
      \AttributeTok{error =} \ControlFlowTok{function}\NormalTok{(e) }\SpecialCharTok{{-}}\ConstantTok{Inf}
\NormalTok{    )}

    \ControlFlowTok{if}\NormalTok{ (}\SpecialCharTok{!}\FunctionTok{is.finite}\NormalTok{(log\_prob\_0) }\SpecialCharTok{\&\&} \SpecialCharTok{!}\FunctionTok{is.finite}\NormalTok{(log\_prob\_1)) \{}
\NormalTok{      prob\_cat\_1 }\OtherTok{\textless{}{-}} \FloatTok{0.5}
\NormalTok{    \} }\ControlFlowTok{else} \ControlFlowTok{if}\NormalTok{ (}\SpecialCharTok{!}\FunctionTok{is.finite}\NormalTok{(log\_prob\_0)) \{}
\NormalTok{      prob\_cat\_1 }\OtherTok{\textless{}{-}} \FloatTok{1.0}
\NormalTok{    \} }\ControlFlowTok{else} \ControlFlowTok{if}\NormalTok{ (}\SpecialCharTok{!}\FunctionTok{is.finite}\NormalTok{(log\_prob\_1)) \{}
\NormalTok{      prob\_cat\_1 }\OtherTok{\textless{}{-}} \FloatTok{0.0}
\NormalTok{    \} }\ControlFlowTok{else}\NormalTok{ \{}
\NormalTok{      prob\_cat\_1 }\OtherTok{\textless{}{-}} \FunctionTok{exp}\NormalTok{(log\_prob\_1 }\SpecialCharTok{{-}} \FunctionTok{log\_sum\_exp}\NormalTok{(}\FunctionTok{c}\NormalTok{(log\_prob\_0, log\_prob\_1)))}
\NormalTok{    \}}
\NormalTok{    response\_probs[i] }\OtherTok{\textless{}{-}} \FunctionTok{pmax}\NormalTok{(}\FloatTok{1e{-}9}\NormalTok{, }\FunctionTok{pmin}\NormalTok{(}\DecValTok{1} \SpecialCharTok{{-}} \FloatTok{1e{-}9}\NormalTok{, prob\_cat\_1))}

    \CommentTok{\# ── Update (measurement update for the correct category only) ────────────}
\NormalTok{    true\_cat }\OtherTok{\textless{}{-}}\NormalTok{ cat\_one[i]}
    \ControlFlowTok{if}\NormalTok{ (true\_cat }\SpecialCharTok{==} \DecValTok{1}\NormalTok{) \{}
\NormalTok{      upd }\OtherTok{\textless{}{-}} \FunctionTok{multivariate\_kalman\_update}\NormalTok{(prototype\_cat\_1}\SpecialCharTok{$}\NormalTok{mu, prototype\_cat\_1}\SpecialCharTok{$}\NormalTok{sigma,}
\NormalTok{                                        current\_obs, r\_matrix)}
\NormalTok{      prototype\_cat\_1}\SpecialCharTok{$}\NormalTok{mu    }\OtherTok{\textless{}{-}}\NormalTok{ upd}\SpecialCharTok{$}\NormalTok{mu}
\NormalTok{      prototype\_cat\_1}\SpecialCharTok{$}\NormalTok{sigma }\OtherTok{\textless{}{-}}\NormalTok{ upd}\SpecialCharTok{$}\NormalTok{sigma}
\NormalTok{    \} }\ControlFlowTok{else}\NormalTok{ \{}
\NormalTok{      upd }\OtherTok{\textless{}{-}} \FunctionTok{multivariate\_kalman\_update}\NormalTok{(prototype\_cat\_0}\SpecialCharTok{$}\NormalTok{mu, prototype\_cat\_0}\SpecialCharTok{$}\NormalTok{sigma,}
\NormalTok{                                        current\_obs, r\_matrix)}
\NormalTok{      prototype\_cat\_0}\SpecialCharTok{$}\NormalTok{mu    }\OtherTok{\textless{}{-}}\NormalTok{ upd}\SpecialCharTok{$}\NormalTok{mu}
\NormalTok{      prototype\_cat\_0}\SpecialCharTok{$}\NormalTok{sigma }\OtherTok{\textless{}{-}}\NormalTok{ upd}\SpecialCharTok{$}\NormalTok{sigma}
\NormalTok{    \}}
\NormalTok{  \}}

  \FunctionTok{tibble}\NormalTok{(}
    \AttributeTok{prob\_cat1    =}\NormalTok{ response\_probs,}
    \AttributeTok{sim\_response =} \FunctionTok{rbinom}\NormalTok{(n\_trials, }\DecValTok{1}\NormalTok{, response\_probs)}
\NormalTok{  )}
\NormalTok{\}}

\CommentTok{\# Wrapper: generates per{-}agent schedule then calls prototype\_kalman}
\NormalTok{simulate\_prototype\_agent }\OtherTok{\textless{}{-}} \ControlFlowTok{function}\NormalTok{(agent\_id, r\_value, q\_value,}
\NormalTok{                                     stimulus\_info, n\_blocks, subject\_seed) \{}
\NormalTok{  schedule }\OtherTok{\textless{}{-}} \FunctionTok{make\_subject\_schedule}\NormalTok{(stimulus\_info, n\_blocks, }\AttributeTok{seed =}\NormalTok{ subject\_seed)}
\NormalTok{  obs      }\OtherTok{\textless{}{-}} \FunctionTok{as.matrix}\NormalTok{(schedule[, }\FunctionTok{c}\NormalTok{(}\StringTok{"height"}\NormalTok{, }\StringTok{"position"}\NormalTok{)])}
\NormalTok{  cat\_one  }\OtherTok{\textless{}{-}}\NormalTok{ schedule}\SpecialCharTok{$}\NormalTok{category\_feedback}

\NormalTok{  result }\OtherTok{\textless{}{-}} \FunctionTok{prototype\_kalman}\NormalTok{(}
    \AttributeTok{r\_value            =}\NormalTok{ r\_value,}
    \AttributeTok{q\_value            =}\NormalTok{ q\_value,}
    \AttributeTok{obs                =}\NormalTok{ obs,}
    \AttributeTok{cat\_one            =}\NormalTok{ cat\_one,}
    \AttributeTok{initial\_mu         =} \FunctionTok{list}\NormalTok{(}\FunctionTok{rep}\NormalTok{(}\FloatTok{2.5}\NormalTok{, }\DecValTok{2}\NormalTok{), }\FunctionTok{rep}\NormalTok{(}\FloatTok{2.5}\NormalTok{, }\DecValTok{2}\NormalTok{)),}
    \AttributeTok{initial\_sigma\_diag =} \FloatTok{10.0}\NormalTok{,}
    \AttributeTok{quiet              =} \ConstantTok{TRUE}
\NormalTok{  )}

\NormalTok{  schedule }\SpecialCharTok{|\textgreater{}}
    \FunctionTok{mutate}\NormalTok{(}
      \AttributeTok{agent\_id     =}\NormalTok{ agent\_id,}
      \AttributeTok{r\_value\_true =}\NormalTok{ r\_value,}
      \AttributeTok{q\_value\_true =}\NormalTok{ q\_value,}
      \AttributeTok{log\_r\_true   =} \FunctionTok{log}\NormalTok{(r\_value),}
      \AttributeTok{log\_q\_true   =} \FunctionTok{log}\NormalTok{(q\_value),}
      \AttributeTok{prob\_cat1    =}\NormalTok{ result}\SpecialCharTok{$}\NormalTok{prob\_cat1,}
      \AttributeTok{sim\_response =}\NormalTok{ result}\SpecialCharTok{$}\NormalTok{sim\_response,}
      \AttributeTok{correct      =} \FunctionTok{as.integer}\NormalTok{(category\_feedback }\SpecialCharTok{==}\NormalTok{ sim\_response)}
\NormalTok{    ) }\SpecialCharTok{|\textgreater{}}
    \FunctionTok{group\_by}\NormalTok{(agent\_id) }\SpecialCharTok{|\textgreater{}}
    \FunctionTok{mutate}\NormalTok{(}\AttributeTok{performance =} \FunctionTok{cumsum}\NormalTok{(correct) }\SpecialCharTok{/} \FunctionTok{row\_number}\NormalTok{()) }\SpecialCharTok{|\textgreater{}}
    \FunctionTok{ungroup}\NormalTok{()}
\NormalTok{\}}
\end{Highlighting}
\end{Shaded}

\subsection{Simulating Categorization
Behavior}\label{simulating-categorization-behavior}

Let's simulate behavior on the Kruschke (1993) task. The key parameter
here is \texttt{r\_value}, the observation noise. This parameter
reflects the model's assumption about how much variability exists within
categories or arises from perceptual noise.

\begin{itemize}
\tightlist
\item
  \textbf{Low \texttt{r\_value}}: The model assumes observations are
  precise representations of the category; each new example exerts
  strong influence and prototypes update rapidly.
\item
  \textbf{High \texttt{r\_value}}: The model assumes observations are
  noisy; learning is slower as each example has less influence, but the
  model generalizes more broadly.
\end{itemize}

\begin{Shaded}
\begin{Highlighting}[]
\NormalTok{param\_df }\OtherTok{\textless{}{-}} \FunctionTok{expand\_grid}\NormalTok{(}
  \AttributeTok{agent\_id  =} \DecValTok{1}\SpecialCharTok{:}\DecValTok{5}\NormalTok{,}
  \AttributeTok{r\_value   =} \FunctionTok{c}\NormalTok{(}\FloatTok{0.5}\NormalTok{, }\FloatTok{2.0}\NormalTok{),}
  \AttributeTok{q\_value   =} \FunctionTok{c}\NormalTok{(}\FloatTok{0.001}\NormalTok{, }\FloatTok{0.1}\NormalTok{, }\FloatTok{0.5}\NormalTok{)}
\NormalTok{) }\SpecialCharTok{|\textgreater{}}
  \FunctionTok{mutate}\NormalTok{(}\AttributeTok{subject\_seed =}\NormalTok{ agent\_id)}

\NormalTok{sim\_file }\OtherTok{\textless{}{-}} \FunctionTok{here}\NormalTok{(}\StringTok{"simdata"}\NormalTok{, }\StringTok{"ch13\_prototype\_simulated\_responses.csv"}\NormalTok{)}
\ControlFlowTok{if}\NormalTok{ (regenerate\_simulations }\SpecialCharTok{||} \SpecialCharTok{!}\FunctionTok{file.exists}\NormalTok{(sim\_file)) \{}
  \FunctionTok{cat}\NormalTok{(}\StringTok{"Regenerating prototype simulations...}\SpecialCharTok{\textbackslash{}n}\StringTok{"}\NormalTok{)}
\NormalTok{  prototype\_responses }\OtherTok{\textless{}{-}} \FunctionTok{future\_pmap\_dfr}\NormalTok{(}
    \FunctionTok{list}\NormalTok{(}
      \AttributeTok{agent\_id     =}\NormalTok{ param\_df}\SpecialCharTok{$}\NormalTok{agent\_id,}
      \AttributeTok{r\_value      =}\NormalTok{ param\_df}\SpecialCharTok{$}\NormalTok{r\_value,}
      \AttributeTok{q\_value      =}\NormalTok{ param\_df}\SpecialCharTok{$}\NormalTok{q\_value,}
      \AttributeTok{subject\_seed =}\NormalTok{ param\_df}\SpecialCharTok{$}\NormalTok{subject\_seed}
\NormalTok{    ),}
\NormalTok{    simulate\_prototype\_agent,}
    \AttributeTok{stimulus\_info =}\NormalTok{ stimulus\_info,}
    \AttributeTok{n\_blocks      =}\NormalTok{ n\_blocks,}
    \AttributeTok{.options =} \FunctionTok{furrr\_options}\NormalTok{(}\AttributeTok{seed =} \ConstantTok{TRUE}\NormalTok{)}
\NormalTok{  )}
  \FunctionTok{write\_csv}\NormalTok{(prototype\_responses, sim\_file)}
  \FunctionTok{cat}\NormalTok{(}\StringTok{"Simulations saved.}\SpecialCharTok{\textbackslash{}n}\StringTok{"}\NormalTok{)}
\NormalTok{\} }\ControlFlowTok{else}\NormalTok{ \{}
\NormalTok{  prototype\_responses }\OtherTok{\textless{}{-}} \FunctionTok{read\_csv}\NormalTok{(sim\_file, }\AttributeTok{show\_col\_types =} \ConstantTok{FALSE}\NormalTok{)}
  \FunctionTok{cat}\NormalTok{(}\StringTok{"Simulations loaded.}\SpecialCharTok{\textbackslash{}n}\StringTok{"}\NormalTok{)}
\NormalTok{\}}
\end{Highlighting}
\end{Shaded}

\begin{verbatim}
Simulations loaded.
\end{verbatim}

\begin{Shaded}
\begin{Highlighting}[]
\FunctionTok{ggplot}\NormalTok{(prototype\_responses,}
       \FunctionTok{aes}\NormalTok{(}\AttributeTok{x =}\NormalTok{ trial\_within\_subject, }\AttributeTok{y =}\NormalTok{ performance, }\AttributeTok{color =} \FunctionTok{factor}\NormalTok{(r\_value\_true))) }\SpecialCharTok{+}
  \FunctionTok{stat\_summary}\NormalTok{(}\AttributeTok{fun =}\NormalTok{ mean, }\AttributeTok{geom =} \StringTok{"line"}\NormalTok{, }\AttributeTok{linewidth =} \DecValTok{1}\NormalTok{) }\SpecialCharTok{+}
  \FunctionTok{stat\_summary}\NormalTok{(}\AttributeTok{fun.data =}\NormalTok{ mean\_se, }\AttributeTok{geom =} \StringTok{"ribbon"}\NormalTok{, }\AttributeTok{alpha =} \FloatTok{0.15}\NormalTok{,}
               \FunctionTok{aes}\NormalTok{(}\AttributeTok{fill =} \FunctionTok{factor}\NormalTok{(r\_value\_true))) }\SpecialCharTok{+}
  \FunctionTok{scale\_color\_viridis\_d}\NormalTok{(}\AttributeTok{option =} \StringTok{"plasma"}\NormalTok{, }\AttributeTok{name =} \StringTok{"r"}\NormalTok{) }\SpecialCharTok{+}
  \FunctionTok{scale\_fill\_viridis\_d}\NormalTok{(}\AttributeTok{option =} \StringTok{"plasma"}\NormalTok{, }\AttributeTok{name =} \StringTok{"r"}\NormalTok{) }\SpecialCharTok{+}
  \FunctionTok{facet\_wrap}\NormalTok{(}\SpecialCharTok{\textasciitilde{}}\FunctionTok{paste0}\NormalTok{(}\StringTok{"q = "}\NormalTok{, q\_value\_true), }\AttributeTok{ncol =} \DecValTok{3}\NormalTok{) }\SpecialCharTok{+}
  \FunctionTok{labs}\NormalTok{(}
    \AttributeTok{title    =} \StringTok{"Prototype Model Learning Performance"}\NormalTok{,}
    \AttributeTok{subtitle =} \StringTok{"Columns: process noise q; colors: observation noise r; mean ± SE across agents"}\NormalTok{,}
    \AttributeTok{x =} \StringTok{"Trial"}\NormalTok{, }\AttributeTok{y =} \StringTok{"Cumulative Accuracy"}
\NormalTok{  ) }\SpecialCharTok{+}
  \FunctionTok{ylim}\NormalTok{(}\FloatTok{0.3}\NormalTok{, }\FloatTok{1.0}\NormalTok{)}
\end{Highlighting}
\end{Shaded}

\begin{center}
\includegraphics[width=0.8\linewidth,height=\textheight,keepaspectratio]{12-CategorizationModelsPrototypes_files/figure-pdf/ch13_simulate_prototype_responses-1.pdf}
\end{center}

\textbf{Interpretation}: The panel structure separates the effects of
\(r\) (colors) and \(q\) (columns). Along a row (fixed \(q\)), higher
\(r\) slows learning and lowers asymptotic accuracy. Down a column
(fixed \(r\)), higher \(q\) introduces persistent uncertainty --- the
model cannot converge to a stable prototype, so accuracy may plateau
below ceiling or even oscillate slightly. When \(q = 0\) the model
reduces to the static-target filter from previous drafts of this
chapter. The key pedagogical point: \(q\) and \(r\) are not redundant
--- they affect different phases of the learning curve and, in
non-stationary environments, become separately identifiable.

\subsection{Visualizing Prototype
Evolution}\label{visualizing-prototype-evolution}

Where do the prototypes end up? Let's track their movement and
uncertainty over trials.

\begin{Shaded}
\begin{Highlighting}[]
\CommentTok{\# Track prototype means and covariance matrices over the course of learning.}
\CommentTok{\# Returns a tibble with one row per (trial × category).}
\NormalTok{track\_prototypes }\OtherTok{\textless{}{-}} \ControlFlowTok{function}\NormalTok{(r\_value, obs, cat\_one,}
                             \AttributeTok{q\_value =} \DecValTok{0}\NormalTok{,}
                             \AttributeTok{initial\_mu =} \ConstantTok{NULL}\NormalTok{, }\AttributeTok{initial\_sigma\_diag =} \FloatTok{10.0}\NormalTok{) \{}
\NormalTok{  n\_trials  }\OtherTok{\textless{}{-}} \FunctionTok{nrow}\NormalTok{(obs)}
\NormalTok{  n\_features }\OtherTok{\textless{}{-}} \FunctionTok{ncol}\NormalTok{(obs)}

  \ControlFlowTok{if}\NormalTok{ (}\FunctionTok{is.null}\NormalTok{(initial\_mu)) \{}
\NormalTok{    mu0\_init }\OtherTok{\textless{}{-}} \FunctionTok{rep}\NormalTok{(}\FloatTok{2.5}\NormalTok{, n\_features)}
\NormalTok{    mu1\_init }\OtherTok{\textless{}{-}} \FunctionTok{rep}\NormalTok{(}\FloatTok{2.5}\NormalTok{, n\_features)}
\NormalTok{  \} }\ControlFlowTok{else}\NormalTok{ \{}
\NormalTok{    mu0\_init }\OtherTok{\textless{}{-}}\NormalTok{ initial\_mu[[}\DecValTok{1}\NormalTok{]]}
\NormalTok{    mu1\_init }\OtherTok{\textless{}{-}}\NormalTok{ initial\_mu[[}\DecValTok{2}\NormalTok{]]}
\NormalTok{  \}}
\NormalTok{  prototype\_cat\_0 }\OtherTok{\textless{}{-}} \FunctionTok{list}\NormalTok{(}\AttributeTok{mu =}\NormalTok{ mu0\_init, }\AttributeTok{sigma =} \FunctionTok{diag}\NormalTok{(initial\_sigma\_diag, n\_features))}
\NormalTok{  prototype\_cat\_1 }\OtherTok{\textless{}{-}} \FunctionTok{list}\NormalTok{(}\AttributeTok{mu =}\NormalTok{ mu1\_init, }\AttributeTok{sigma =} \FunctionTok{diag}\NormalTok{(initial\_sigma\_diag, n\_features))}
\NormalTok{  r\_matrix        }\OtherTok{\textless{}{-}} \FunctionTok{diag}\NormalTok{(r\_value, n\_features)}

\NormalTok{  history }\OtherTok{\textless{}{-}} \FunctionTok{list}\NormalTok{()}

  \ControlFlowTok{for}\NormalTok{ (i }\ControlFlowTok{in} \FunctionTok{seq\_len}\NormalTok{(n\_trials }\SpecialCharTok{+} \DecValTok{1}\NormalTok{)) \{}
\NormalTok{    history[[}\FunctionTok{length}\NormalTok{(history) }\SpecialCharTok{+} \DecValTok{1}\NormalTok{]] }\OtherTok{\textless{}{-}} \FunctionTok{tibble}\NormalTok{(}
      \AttributeTok{trial      =}\NormalTok{ i, }\AttributeTok{category =} \DecValTok{0}\NormalTok{,}
      \AttributeTok{feature1\_mean =}\NormalTok{ prototype\_cat\_0}\SpecialCharTok{$}\NormalTok{mu[}\DecValTok{1}\NormalTok{],}
      \AttributeTok{feature2\_mean =}\NormalTok{ prototype\_cat\_0}\SpecialCharTok{$}\NormalTok{mu[}\DecValTok{2}\NormalTok{],}
      \AttributeTok{cov\_matrix    =} \FunctionTok{list}\NormalTok{(prototype\_cat\_0}\SpecialCharTok{$}\NormalTok{sigma)}
\NormalTok{    )}
\NormalTok{    history[[}\FunctionTok{length}\NormalTok{(history) }\SpecialCharTok{+} \DecValTok{1}\NormalTok{]] }\OtherTok{\textless{}{-}} \FunctionTok{tibble}\NormalTok{(}
      \AttributeTok{trial      =}\NormalTok{ i, }\AttributeTok{category =} \DecValTok{1}\NormalTok{,}
      \AttributeTok{feature1\_mean =}\NormalTok{ prototype\_cat\_1}\SpecialCharTok{$}\NormalTok{mu[}\DecValTok{1}\NormalTok{],}
      \AttributeTok{feature2\_mean =}\NormalTok{ prototype\_cat\_1}\SpecialCharTok{$}\NormalTok{mu[}\DecValTok{2}\NormalTok{],}
      \AttributeTok{cov\_matrix    =} \FunctionTok{list}\NormalTok{(prototype\_cat\_1}\SpecialCharTok{$}\NormalTok{sigma)}
\NormalTok{    )}

    \ControlFlowTok{if}\NormalTok{ (i }\SpecialCharTok{\textless{}=}\NormalTok{ n\_trials) \{}
      \CommentTok{\# Prediction step}
\NormalTok{      prototype\_cat\_0}\SpecialCharTok{$}\NormalTok{sigma }\OtherTok{\textless{}{-}}\NormalTok{ prototype\_cat\_0}\SpecialCharTok{$}\NormalTok{sigma }\SpecialCharTok{+} \FunctionTok{diag}\NormalTok{(q\_value, n\_features)}
\NormalTok{      prototype\_cat\_1}\SpecialCharTok{$}\NormalTok{sigma }\OtherTok{\textless{}{-}}\NormalTok{ prototype\_cat\_1}\SpecialCharTok{$}\NormalTok{sigma }\SpecialCharTok{+} \FunctionTok{diag}\NormalTok{(q\_value, n\_features)}

\NormalTok{      current\_obs }\OtherTok{\textless{}{-}} \FunctionTok{as.numeric}\NormalTok{(obs[i, ])}
\NormalTok{      true\_cat    }\OtherTok{\textless{}{-}}\NormalTok{ cat\_one[i]}
      \ControlFlowTok{if}\NormalTok{ (true\_cat }\SpecialCharTok{==} \DecValTok{1}\NormalTok{) \{}
\NormalTok{        upd }\OtherTok{\textless{}{-}} \FunctionTok{multivariate\_kalman\_update}\NormalTok{(prototype\_cat\_1}\SpecialCharTok{$}\NormalTok{mu, prototype\_cat\_1}\SpecialCharTok{$}\NormalTok{sigma,}
\NormalTok{                                          current\_obs, r\_matrix)}
\NormalTok{        prototype\_cat\_1}\SpecialCharTok{$}\NormalTok{mu    }\OtherTok{\textless{}{-}}\NormalTok{ upd}\SpecialCharTok{$}\NormalTok{mu}
\NormalTok{        prototype\_cat\_1}\SpecialCharTok{$}\NormalTok{sigma }\OtherTok{\textless{}{-}}\NormalTok{ upd}\SpecialCharTok{$}\NormalTok{sigma}
\NormalTok{      \} }\ControlFlowTok{else}\NormalTok{ \{}
\NormalTok{        upd }\OtherTok{\textless{}{-}} \FunctionTok{multivariate\_kalman\_update}\NormalTok{(prototype\_cat\_0}\SpecialCharTok{$}\NormalTok{mu, prototype\_cat\_0}\SpecialCharTok{$}\NormalTok{sigma,}
\NormalTok{                                          current\_obs, r\_matrix)}
\NormalTok{        prototype\_cat\_0}\SpecialCharTok{$}\NormalTok{mu    }\OtherTok{\textless{}{-}}\NormalTok{ upd}\SpecialCharTok{$}\NormalTok{mu}
\NormalTok{        prototype\_cat\_0}\SpecialCharTok{$}\NormalTok{sigma }\OtherTok{\textless{}{-}}\NormalTok{ upd}\SpecialCharTok{$}\NormalTok{sigma}
\NormalTok{      \}}
\NormalTok{    \}}
\NormalTok{  \}}
  \FunctionTok{bind\_rows}\NormalTok{(history)}
\NormalTok{\}}

\CommentTok{\# Ellipse helper for uncertainty visualization}
\NormalTok{get\_ellipse }\OtherTok{\textless{}{-}} \ControlFlowTok{function}\NormalTok{(mu, sigma, }\AttributeTok{level =} \FloatTok{0.68}\NormalTok{) \{}
\NormalTok{  mu    }\OtherTok{\textless{}{-}} \FunctionTok{as.numeric}\NormalTok{(mu)}
\NormalTok{  sigma }\OtherTok{\textless{}{-}} \FunctionTok{as.matrix}\NormalTok{(sigma)}
  \ControlFlowTok{if}\NormalTok{ (}\FunctionTok{length}\NormalTok{(mu) }\SpecialCharTok{!=} \DecValTok{2} \SpecialCharTok{||} \SpecialCharTok{!}\FunctionTok{all}\NormalTok{(}\FunctionTok{dim}\NormalTok{(sigma) }\SpecialCharTok{==} \FunctionTok{c}\NormalTok{(}\DecValTok{2}\NormalTok{, }\DecValTok{2}\NormalTok{))) }\FunctionTok{return}\NormalTok{(}\ConstantTok{NULL}\NormalTok{)}
\NormalTok{  ev }\OtherTok{\textless{}{-}} \FunctionTok{eigen}\NormalTok{(sigma, }\AttributeTok{symmetric =} \ConstantTok{TRUE}\NormalTok{, }\AttributeTok{only.values =} \ConstantTok{TRUE}\NormalTok{)}\SpecialCharTok{$}\NormalTok{values}
  \ControlFlowTok{if}\NormalTok{ (}\FunctionTok{any}\NormalTok{(ev }\SpecialCharTok{\textless{}=} \FloatTok{1e{-}6}\NormalTok{)) sigma }\OtherTok{\textless{}{-}}\NormalTok{ sigma }\SpecialCharTok{+} \FunctionTok{diag}\NormalTok{(}\FunctionTok{ncol}\NormalTok{(sigma)) }\SpecialCharTok{*} \FloatTok{1e{-}6}
\NormalTok{  pts }\OtherTok{\textless{}{-}} \FunctionTok{tryCatch}\NormalTok{(}
\NormalTok{    ellipse}\SpecialCharTok{::}\FunctionTok{ellipse}\NormalTok{(sigma, }\AttributeTok{centre =}\NormalTok{ mu, }\AttributeTok{level =}\NormalTok{ level),}
    \AttributeTok{error =} \ControlFlowTok{function}\NormalTok{(e) }\ConstantTok{NULL}
\NormalTok{  )}
  \ControlFlowTok{if}\NormalTok{ (}\FunctionTok{is.null}\NormalTok{(pts)) }\FunctionTok{return}\NormalTok{(}\ConstantTok{NULL}\NormalTok{)}
  \FunctionTok{as\_tibble}\NormalTok{(pts) }\SpecialCharTok{|\textgreater{}} \FunctionTok{setNames}\NormalTok{(}\FunctionTok{c}\NormalTok{(}\StringTok{"feature1\_mean"}\NormalTok{, }\StringTok{"feature2\_mean"}\NormalTok{))}
\NormalTok{\}}

\CommentTok{\# Use a fixed seed{-}1 schedule for the trajectory visualization}
\NormalTok{traj\_schedule }\OtherTok{\textless{}{-}} \FunctionTok{make\_subject\_schedule}\NormalTok{(stimulus\_info, n\_blocks, }\AttributeTok{seed =} \DecValTok{1}\NormalTok{)}

\NormalTok{prototype\_trajectory }\OtherTok{\textless{}{-}} \FunctionTok{track\_prototypes}\NormalTok{(}
  \AttributeTok{r\_value            =} \FloatTok{1.0}\NormalTok{,}
  \AttributeTok{q\_value            =} \FloatTok{0.05}\NormalTok{,}
  \AttributeTok{obs                =} \FunctionTok{as.matrix}\NormalTok{(traj\_schedule[, }\FunctionTok{c}\NormalTok{(}\StringTok{"height"}\NormalTok{, }\StringTok{"position"}\NormalTok{)]),}
  \AttributeTok{cat\_one            =}\NormalTok{ traj\_schedule}\SpecialCharTok{$}\NormalTok{category\_feedback,}
  \AttributeTok{initial\_mu         =} \FunctionTok{list}\NormalTok{(}\FunctionTok{rep}\NormalTok{(}\FloatTok{2.5}\NormalTok{, }\DecValTok{2}\NormalTok{), }\FunctionTok{rep}\NormalTok{(}\FloatTok{2.5}\NormalTok{, }\DecValTok{2}\NormalTok{)),}
  \AttributeTok{initial\_sigma\_diag =} \FloatTok{10.0}
\NormalTok{)}

\NormalTok{final\_prototypes    }\OtherTok{\textless{}{-}}\NormalTok{ prototype\_trajectory }\SpecialCharTok{|\textgreater{}} \FunctionTok{filter}\NormalTok{(trial }\SpecialCharTok{==} \FunctionTok{max}\NormalTok{(trial))}

\NormalTok{ellipse\_data }\OtherTok{\textless{}{-}}\NormalTok{ final\_prototypes }\SpecialCharTok{|\textgreater{}}
  \FunctionTok{rowwise}\NormalTok{() }\SpecialCharTok{|\textgreater{}}
  \FunctionTok{mutate}\NormalTok{(}\AttributeTok{ell =} \FunctionTok{list}\NormalTok{(}\FunctionTok{get\_ellipse}\NormalTok{(}\FunctionTok{c}\NormalTok{(feature1\_mean, feature2\_mean), cov\_matrix[[}\DecValTok{1}\NormalTok{]]))) }\SpecialCharTok{|\textgreater{}}
  \FunctionTok{ungroup}\NormalTok{() }\SpecialCharTok{|\textgreater{}}
  \FunctionTok{filter}\NormalTok{(}\SpecialCharTok{!}\FunctionTok{map\_lgl}\NormalTok{(ell, is.null)) }\SpecialCharTok{|\textgreater{}}
  \FunctionTok{unnest}\NormalTok{(ell)}

\FunctionTok{ggplot}\NormalTok{() }\SpecialCharTok{+}
  \FunctionTok{geom\_point}\NormalTok{(}\AttributeTok{data =}\NormalTok{ stimulus\_info,}
             \FunctionTok{aes}\NormalTok{(}\AttributeTok{x =}\NormalTok{ position, }\AttributeTok{y =}\NormalTok{ height, }\AttributeTok{color =} \FunctionTok{factor}\NormalTok{(category\_true),}
                 \AttributeTok{shape =} \FunctionTok{factor}\NormalTok{(category\_true)),}
             \AttributeTok{size =} \DecValTok{3}\NormalTok{, }\AttributeTok{alpha =} \FloatTok{0.5}\NormalTok{) }\SpecialCharTok{+}
  \FunctionTok{geom\_path}\NormalTok{(}
    \AttributeTok{data =}\NormalTok{ prototype\_trajectory }\SpecialCharTok{|\textgreater{}} \FunctionTok{filter}\NormalTok{(trial }\SpecialCharTok{\textless{}=}\NormalTok{ total\_trials),}
    \FunctionTok{aes}\NormalTok{(}\AttributeTok{x =}\NormalTok{ feature2\_mean, }\AttributeTok{y =}\NormalTok{ feature1\_mean,}
        \AttributeTok{group =}\NormalTok{ category, }\AttributeTok{color =} \FunctionTok{factor}\NormalTok{(category)),}
    \AttributeTok{linetype =} \StringTok{"dashed"}\NormalTok{,}
    \AttributeTok{arrow =} \FunctionTok{arrow}\NormalTok{(}\AttributeTok{type =} \StringTok{"closed"}\NormalTok{, }\AttributeTok{length =} \FunctionTok{unit}\NormalTok{(}\FloatTok{0.08}\NormalTok{, }\StringTok{"inches"}\NormalTok{), }\AttributeTok{angle =} \DecValTok{20}\NormalTok{)}
\NormalTok{  ) }\SpecialCharTok{+}
  \FunctionTok{geom\_point}\NormalTok{(}\AttributeTok{data =}\NormalTok{ final\_prototypes,}
             \FunctionTok{aes}\NormalTok{(}\AttributeTok{x =}\NormalTok{ feature2\_mean, }\AttributeTok{y =}\NormalTok{ feature1\_mean, }\AttributeTok{color =} \FunctionTok{factor}\NormalTok{(category)),}
             \AttributeTok{size =} \DecValTok{4}\NormalTok{, }\AttributeTok{shape =} \DecValTok{18}\NormalTok{) }\SpecialCharTok{+}
  \FunctionTok{geom\_path}\NormalTok{(}\AttributeTok{data =}\NormalTok{ ellipse\_data,}
            \FunctionTok{aes}\NormalTok{(}\AttributeTok{x =}\NormalTok{ feature2\_mean, }\AttributeTok{y =}\NormalTok{ feature1\_mean,}
                \AttributeTok{group =}\NormalTok{ category, }\AttributeTok{color =} \FunctionTok{factor}\NormalTok{(category)),}
            \AttributeTok{alpha =} \FloatTok{0.7}\NormalTok{, }\AttributeTok{linewidth =} \FloatTok{0.8}\NormalTok{) }\SpecialCharTok{+}
  \FunctionTok{scale\_color\_manual}\NormalTok{(}\AttributeTok{values =} \FunctionTok{c}\NormalTok{(}\StringTok{"0"} \OtherTok{=} \StringTok{"\#0072B2"}\NormalTok{, }\StringTok{"1"} \OtherTok{=} \StringTok{"\#D55E00"}\NormalTok{), }\AttributeTok{name =} \StringTok{"Category"}\NormalTok{) }\SpecialCharTok{+}
  \FunctionTok{scale\_shape\_manual}\NormalTok{(}\AttributeTok{values =} \FunctionTok{c}\NormalTok{(}\StringTok{"0"} \OtherTok{=} \DecValTok{16}\NormalTok{, }\StringTok{"1"} \OtherTok{=} \DecValTok{17}\NormalTok{), }\AttributeTok{name =} \StringTok{"Category"}\NormalTok{) }\SpecialCharTok{+}
  \FunctionTok{labs}\NormalTok{(}
    \AttributeTok{title    =} \StringTok{"Prototype Learning Trajectory (Kalman Filter, r = 1.0)"}\NormalTok{,}
    \AttributeTok{subtitle =} \StringTok{"Prototypes start at (2.5, 2.5). Dashed lines show path; ellipses show final 68\% uncertainty."}\NormalTok{,}
    \AttributeTok{x =} \StringTok{"Position Feature"}\NormalTok{, }\AttributeTok{y =} \StringTok{"Height Feature"}
\NormalTok{  ) }\SpecialCharTok{+}
  \FunctionTok{coord\_cartesian}\NormalTok{(}
    \AttributeTok{xlim =} \FunctionTok{range}\NormalTok{(}\FunctionTok{c}\NormalTok{(stimulus\_info}\SpecialCharTok{$}\NormalTok{position, prototype\_trajectory}\SpecialCharTok{$}\NormalTok{feature2\_mean)) }\SpecialCharTok{+} \FunctionTok{c}\NormalTok{(}\SpecialCharTok{{-}}\FloatTok{0.5}\NormalTok{, }\FloatTok{0.5}\NormalTok{),}
    \AttributeTok{ylim =} \FunctionTok{range}\NormalTok{(}\FunctionTok{c}\NormalTok{(stimulus\_info}\SpecialCharTok{$}\NormalTok{height,   prototype\_trajectory}\SpecialCharTok{$}\NormalTok{feature1\_mean)) }\SpecialCharTok{+} \FunctionTok{c}\NormalTok{(}\SpecialCharTok{{-}}\FloatTok{0.5}\NormalTok{, }\FloatTok{0.5}\NormalTok{)}
\NormalTok{  )}
\end{Highlighting}
\end{Shaded}

\begin{center}
\includegraphics[width=0.8\linewidth,height=\textheight,keepaspectratio]{12-CategorizationModelsPrototypes_files/figure-pdf/ch13_track_prototypes-1.pdf}
\end{center}

\textbf{Interpretation}: The dashed lines show the path the prototypes
took during learning (starting from (2.5, 2.5), the center of the
feature space). The final prototype locations are marked with diamonds,
and the ellipses represent the model's final uncertainty (68\%
confidence) about the true prototype location for each category. We see
the prototypes moving towards the center of their respective category
members and uncertainty shrinking over time --- exactly the incremental,
variance-reducing Bayesian update the Kalman filter implements. The
shrinking ellipses indicate that once the agent has seen enough
exemplars, the covariances are small enough that further learning is
negligible, and the prototypes are effectively frozen.

\begin{center}\rule{0.5\linewidth}{0.5pt}\end{center}

\section{Prior Predictive Check}\label{prior-predictive-check-3}

Before fitting any data, we verify that the prior on \texttt{log\_r}
generates only scientifically plausible learning trajectories. We draw
\(S = 500\) values from \(\log r \sim \text{Normal}(0, 1)\) --- implying
\(r \in (0.13, 7.4)\) at \(\pm 2\) SD --- simulate one agent per draw,
and plot the envelope.

\begin{Shaded}
\begin{Highlighting}[]
\NormalTok{n\_ppc\_samples }\OtherTok{\textless{}{-}} \DecValTok{500}

\NormalTok{ppc\_log\_r  }\OtherTok{\textless{}{-}} \FunctionTok{rnorm}\NormalTok{(n\_ppc\_samples, LOG\_R\_PRIOR\_MEAN, LOG\_R\_PRIOR\_SD)}
\NormalTok{ppc\_log\_q  }\OtherTok{\textless{}{-}} \FunctionTok{rnorm}\NormalTok{(n\_ppc\_samples, LOG\_Q\_PRIOR\_MEAN, LOG\_Q\_PRIOR\_SD)}
\NormalTok{ppc\_r      }\OtherTok{\textless{}{-}} \FunctionTok{exp}\NormalTok{(ppc\_log\_r)}
\NormalTok{ppc\_q      }\OtherTok{\textless{}{-}} \FunctionTok{exp}\NormalTok{(ppc\_log\_q)}
\NormalTok{ppc\_sched  }\OtherTok{\textless{}{-}} \FunctionTok{make\_subject\_schedule}\NormalTok{(stimulus\_info, n\_blocks, }\AttributeTok{seed =} \DecValTok{999}\NormalTok{)}

\NormalTok{prior\_pred\_curves }\OtherTok{\textless{}{-}} \FunctionTok{map\_dfr}\NormalTok{(}\FunctionTok{seq\_len}\NormalTok{(n\_ppc\_samples), }\ControlFlowTok{function}\NormalTok{(s) \{}
\NormalTok{  res }\OtherTok{\textless{}{-}} \FunctionTok{prototype\_kalman}\NormalTok{(}
    \AttributeTok{r\_value   =}\NormalTok{ ppc\_r[s],}
    \AttributeTok{q\_value   =}\NormalTok{ ppc\_q[s],}
    \AttributeTok{obs       =} \FunctionTok{as.matrix}\NormalTok{(ppc\_sched[, }\FunctionTok{c}\NormalTok{(}\StringTok{"height"}\NormalTok{, }\StringTok{"position"}\NormalTok{)]),}
    \AttributeTok{cat\_one   =}\NormalTok{ ppc\_sched}\SpecialCharTok{$}\NormalTok{category\_feedback,}
    \AttributeTok{initial\_mu =} \FunctionTok{list}\NormalTok{(}\FunctionTok{rep}\NormalTok{(}\FloatTok{2.5}\NormalTok{, }\DecValTok{2}\NormalTok{), }\FunctionTok{rep}\NormalTok{(}\FloatTok{2.5}\NormalTok{, }\DecValTok{2}\NormalTok{))}
\NormalTok{  )}
  \FunctionTok{tibble}\NormalTok{(}
    \AttributeTok{sample\_id =}\NormalTok{ s,}
    \AttributeTok{trial     =} \FunctionTok{seq\_len}\NormalTok{(}\FunctionTok{nrow}\NormalTok{(ppc\_sched)),}
    \AttributeTok{correct   =} \FunctionTok{as.integer}\NormalTok{(res}\SpecialCharTok{$}\NormalTok{sim\_response }\SpecialCharTok{==}\NormalTok{ ppc\_sched}\SpecialCharTok{$}\NormalTok{category\_feedback)}
\NormalTok{  ) }\SpecialCharTok{|\textgreater{}}
    \FunctionTok{mutate}\NormalTok{(}\AttributeTok{cum\_acc =} \FunctionTok{cumsum}\NormalTok{(correct) }\SpecialCharTok{/} \FunctionTok{row\_number}\NormalTok{())}
\NormalTok{\})}

\NormalTok{ppc\_summary }\OtherTok{\textless{}{-}}\NormalTok{ prior\_pred\_curves }\SpecialCharTok{|\textgreater{}}
  \FunctionTok{group\_by}\NormalTok{(trial) }\SpecialCharTok{|\textgreater{}}
  \FunctionTok{summarise}\NormalTok{(}
    \AttributeTok{q05 =} \FunctionTok{quantile}\NormalTok{(cum\_acc, }\FloatTok{0.05}\NormalTok{), }\AttributeTok{q25 =} \FunctionTok{quantile}\NormalTok{(cum\_acc, }\FloatTok{0.25}\NormalTok{),}
    \AttributeTok{q50 =} \FunctionTok{median}\NormalTok{(cum\_acc),}
    \AttributeTok{q75 =} \FunctionTok{quantile}\NormalTok{(cum\_acc, }\FloatTok{0.75}\NormalTok{), }\AttributeTok{q95 =} \FunctionTok{quantile}\NormalTok{(cum\_acc, }\FloatTok{0.95}\NormalTok{),}
    \AttributeTok{.groups =} \StringTok{"drop"}
\NormalTok{  )}

\FunctionTok{ggplot}\NormalTok{(ppc\_summary, }\FunctionTok{aes}\NormalTok{(}\AttributeTok{x =}\NormalTok{ trial)) }\SpecialCharTok{+}
  \FunctionTok{geom\_ribbon}\NormalTok{(}\FunctionTok{aes}\NormalTok{(}\AttributeTok{ymin =}\NormalTok{ q05, }\AttributeTok{ymax =}\NormalTok{ q95), }\AttributeTok{fill =} \StringTok{"\#009E73"}\NormalTok{, }\AttributeTok{alpha =} \FloatTok{0.20}\NormalTok{) }\SpecialCharTok{+}
  \FunctionTok{geom\_ribbon}\NormalTok{(}\FunctionTok{aes}\NormalTok{(}\AttributeTok{ymin =}\NormalTok{ q25, }\AttributeTok{ymax =}\NormalTok{ q75), }\AttributeTok{fill =} \StringTok{"\#009E73"}\NormalTok{, }\AttributeTok{alpha =} \FloatTok{0.40}\NormalTok{) }\SpecialCharTok{+}
  \FunctionTok{geom\_line}\NormalTok{(}\FunctionTok{aes}\NormalTok{(}\AttributeTok{y =}\NormalTok{ q50), }\AttributeTok{color =} \StringTok{"\#006D4E"}\NormalTok{, }\AttributeTok{linewidth =} \DecValTok{1}\NormalTok{) }\SpecialCharTok{+}
  \FunctionTok{geom\_hline}\NormalTok{(}\AttributeTok{yintercept =} \FloatTok{0.5}\NormalTok{, }\AttributeTok{linetype =} \StringTok{"dashed"}\NormalTok{, }\AttributeTok{color =} \StringTok{"grey40"}\NormalTok{) }\SpecialCharTok{+}
  \FunctionTok{scale\_y\_continuous}\NormalTok{(}\AttributeTok{limits =} \FunctionTok{c}\NormalTok{(}\DecValTok{0}\NormalTok{, }\DecValTok{1}\NormalTok{)) }\SpecialCharTok{+}
  \FunctionTok{labs}\NormalTok{(}
    \AttributeTok{title    =} \StringTok{"Prior Predictive Check: Kalman Filter Prototype Model"}\NormalTok{,}
    \AttributeTok{subtitle =} \StringTok{"Ribbons: 50\% and 90\% prior predictive intervals}\SpecialCharTok{\textbackslash{}n}\StringTok{Priors: log(r) \textasciitilde{} Normal(0, 1), log(q) \textasciitilde{} Normal({-}2, 1)"}\NormalTok{,}
    \AttributeTok{x =} \StringTok{"Trial"}\NormalTok{, }\AttributeTok{y =} \StringTok{"Cumulative Accuracy"}
\NormalTok{  )}
\end{Highlighting}
\end{Shaded}

\begin{center}
\includegraphics[width=0.8\linewidth,height=\textheight,keepaspectratio]{12-CategorizationModelsPrototypes_files/figure-pdf/ch13_prior_predictive_check-1.pdf}
\end{center}

The prior predictive envelope spans chance performance (uninformative
\(r\) draws) through near-ceiling accuracy (informative \(r\) draws),
without implying impossible behavior such as sustained below-chance
performance. This is the plausibility standard we require before
proceeding to inference.

\begin{center}\rule{0.5\linewidth}{0.5pt}\end{center}

\section{The Prototype Model in Stan}\label{the-prototype-model-in-stan}

Simulations help us understand the model, but we want to fit it to
experimental data to estimate its parameters. The primary free parameter
is the observation noise variance \texttt{r\_value}. We implement the
Kalman filter prototype model in Stan to estimate \texttt{r\_value} (on
the log scale) from data.

\subsection{\texorpdfstring{Why the Kalman state lives in
\texttt{transformed\ parameters}}{Why the Kalman state lives in transformed parameters}}\label{why-the-kalman-state-lives-in-transformed-parameters}

The prototype model has two parameters --- \texttt{log\_r} and
\texttt{log\_q}. Given both and the observed data (stimuli and
feedback), the \emph{entire} Kalman filter state is deterministic:
prototype means, prototype covariances, and the per-trial choice
probability \(p_i\) all follow algebraically from \texttt{log\_r} and
the trial sequence. Whenever a quantity depends on a parameter but is
otherwise deterministic, it belongs in \texttt{transformed\ parameters}
(evaluated once per leapfrog step) rather than
\texttt{transformed\ data} (which is for parameter-free preprocessing)
or \texttt{model\{\}} (which would force a recomputation in
\texttt{generated\ quantities}). We follow that rule here, and
\texttt{log\_lik{[}i{]}} in \texttt{generated\ quantities} becomes a
one-line lookup of \texttt{p{[}i{]}}.

\begin{quote}
\textbf{A note on architectural symmetry with Chapter 11.} The Chapter
11 GCM's choice probabilities are \emph{also} deterministic given the
parameters and the data, and they \emph{also} belong in
\texttt{transformed\ parameters} for the same reason. We do the same
here: \texttt{prob\_cat1{[}i{]}} is a transformed parameter, the model
block reduces to a single vectorised
\texttt{bernoulli\_lpmf(y\ \textbar{}\ prob\_cat1)}, and
\texttt{log\_lik{[}i{]}} is a one-line lookup.
\end{quote}

The same path-dependence concern from Ch. 11 still applies: trial
\(t\)'s choice probability depends on all previous stimuli through the
Kalman state, so standard PSIS-LOO is invalid. Unlike previous drafts,
we no longer treat that as a caveat to repeat --- the LFO-CV machinery
from Chapter 11 is generalised to the Kalman case in §``Leave-Future-Out
Cross-Validation'' below, and the comparison against exact
one-step-ahead refits is reported in this chapter rather than deferred
to a hypothetical future supplement.

\begin{quote}
\textbf{\texttt{reduce\_sum} and parallelization}: As with the GCM, the
Kalman prototype model cannot use \texttt{reduce\_sum} because
successive Kalman steps are causally chained --- trial \(t\)'s
contribution to the likelihood depends on the filter state from trials
\(1, \ldots, t-1\). The entire filter must be computed sequentially.
With two parameters (\texttt{log\_r}, \texttt{log\_q}) the posterior
geometry is 2D, and Pathfinder initialization is still reliable because
both parameters are unconstrained scalars with prior-dominated tails.
\end{quote}

\subsection{Stan Model}\label{stan-model}

\begin{Shaded}
\begin{Highlighting}[]
\NormalTok{prototype\_single\_stan }\OtherTok{\textless{}{-}} \StringTok{"}
\StringTok{// Kalman Filter Prototype Model — Single Subject, with Process Noise}
\StringTok{// Key design:}
\StringTok{//   log\_r and log\_q are the two parameters (both unconstrained).}
\StringTok{//   r\_value = exp(log\_r): observation noise variance}
\StringTok{//   q\_value = exp(log\_q): process noise / prototype drift rate}
\StringTok{//   Kalman loop structure: Predict (add Q) → Decide → Update.}
\StringTok{//   The full filter is in transformed parameters so log\_lik in generated}
\StringTok{//   quantities reads p[i] directly without re{-}running the filter.}

\StringTok{data \{}
\StringTok{  int\textless{}lower=1\textgreater{} ntrials;}
\StringTok{  int\textless{}lower=1\textgreater{} nfeatures;}
\StringTok{  array[ntrials] int\textless{}lower=0, upper=1\textgreater{} cat\_one;}
\StringTok{  array[ntrials] int\textless{}lower=0, upper=1\textgreater{} y;}
\StringTok{  array[ntrials, nfeatures] real obs;}

\StringTok{  vector[nfeatures] initial\_mu\_cat0;}
\StringTok{  vector[nfeatures] initial\_mu\_cat1;}
\StringTok{  real\textless{}lower=0\textgreater{} initial\_sigma\_diag;}

\StringTok{  real prior\_logr\_mean;}
\StringTok{  real\textless{}lower=0\textgreater{} prior\_logr\_sd;}
\StringTok{  real prior\_logq\_mean;}
\StringTok{  real\textless{}lower=0\textgreater{} prior\_logq\_sd;}
\StringTok{\}}

\StringTok{parameters \{}
\StringTok{  real log\_r;}
\StringTok{  real log\_q;}
\StringTok{\}}

\StringTok{transformed parameters \{}
\StringTok{  real\textless{}lower=0\textgreater{} r\_value = exp(log\_r);}
\StringTok{  real\textless{}lower=0\textgreater{} q\_value = exp(log\_q);}

\StringTok{  array[ntrials] real\textless{}lower=1e{-}9, upper=1{-}1e{-}9\textgreater{} p;}

\StringTok{  \{}
\StringTok{    vector[nfeatures] mu\_cat0 = initial\_mu\_cat0;}
\StringTok{    vector[nfeatures] mu\_cat1 = initial\_mu\_cat1;}
\StringTok{    matrix[nfeatures, nfeatures] sigma\_cat0 =}
\StringTok{      diag\_matrix(rep\_vector(initial\_sigma\_diag, nfeatures));}
\StringTok{    matrix[nfeatures, nfeatures] sigma\_cat1 =}
\StringTok{      diag\_matrix(rep\_vector(initial\_sigma\_diag, nfeatures));}
\StringTok{    matrix[nfeatures, nfeatures] r\_matrix =}
\StringTok{      diag\_matrix(rep\_vector(r\_value, nfeatures));}
\StringTok{    matrix[nfeatures, nfeatures] q\_matrix =}
\StringTok{      diag\_matrix(rep\_vector(q\_value, nfeatures));}
\StringTok{    matrix[nfeatures, nfeatures] I\_mat =}
\StringTok{      diag\_matrix(rep\_vector(1.0, nfeatures));}

\StringTok{    for (i in 1:ntrials) \{}
\StringTok{      vector[nfeatures] x = to\_vector(obs[i]);}

\StringTok{      // ── Prediction step: add process noise to both categories ──────────}
\StringTok{      sigma\_cat0 = sigma\_cat0 + q\_matrix;}
\StringTok{      sigma\_cat1 = sigma\_cat1 + q\_matrix;}

\StringTok{      // ── Decision ────────────────────────────────────────────────────────}
\StringTok{      matrix[nfeatures, nfeatures] cov0 = sigma\_cat0 + r\_matrix;}
\StringTok{      matrix[nfeatures, nfeatures] cov1 = sigma\_cat1 + r\_matrix;}

\StringTok{      real log\_p0 = multi\_normal\_lpdf(x | mu\_cat0, cov0);}
\StringTok{      real log\_p1 = multi\_normal\_lpdf(x | mu\_cat1, cov1);}
\StringTok{      real prob1  = exp(log\_p1 {-} log\_sum\_exp(log\_p0, log\_p1));}

\StringTok{      p[i] = fmax(1e{-}9, fmin(1 {-} 1e{-}9, prob1));}

\StringTok{      // ── Update (measurement update for the correct category only) ────────}
\StringTok{      if (cat\_one[i] == 1) \{}
\StringTok{        vector[nfeatures] innov = x {-} mu\_cat1;}
\StringTok{        matrix[nfeatures, nfeatures] S  = sigma\_cat1 + r\_matrix;}
\StringTok{        matrix[nfeatures, nfeatures] K  = mdivide\_right\_spd(sigma\_cat1, S);}
\StringTok{        matrix[nfeatures, nfeatures] IK = I\_mat {-} K;}
\StringTok{        mu\_cat1    = mu\_cat1 + K * innov;}
\StringTok{        sigma\_cat1 = IK * sigma\_cat1 * IK\textquotesingle{} + K * r\_matrix * K\textquotesingle{};}
\StringTok{        sigma\_cat1 = 0.5 * (sigma\_cat1 + sigma\_cat1\textquotesingle{});}
\StringTok{      \} else \{}
\StringTok{        vector[nfeatures] innov = x {-} mu\_cat0;}
\StringTok{        matrix[nfeatures, nfeatures] S  = sigma\_cat0 + r\_matrix;}
\StringTok{        matrix[nfeatures, nfeatures] K  = mdivide\_right\_spd(sigma\_cat0, S);}
\StringTok{        matrix[nfeatures, nfeatures] IK = I\_mat {-} K;}
\StringTok{        mu\_cat0    = mu\_cat0 + K * innov;}
\StringTok{        sigma\_cat0 = IK * sigma\_cat0 * IK\textquotesingle{} + K * r\_matrix * K\textquotesingle{};}
\StringTok{        sigma\_cat0 = 0.5 * (sigma\_cat0 + sigma\_cat0\textquotesingle{});}
\StringTok{      \}}
\StringTok{    \}}
\StringTok{  \}}
\StringTok{\}}

\StringTok{model \{}
\StringTok{  target += normal\_lpdf(log\_r | prior\_logr\_mean, prior\_logr\_sd);}
\StringTok{  target += normal\_lpdf(log\_q | prior\_logq\_mean, prior\_logq\_sd);}
\StringTok{  target += bernoulli\_lpmf(y | p);}
\StringTok{\}}

\StringTok{generated quantities \{}
\StringTok{  vector[ntrials] log\_lik;}
\StringTok{  real lprior;}
\StringTok{  for (i in 1:ntrials)}
\StringTok{    log\_lik[i] = bernoulli\_lpmf(y[i] | p[i]);}
\StringTok{  lprior = normal\_lpdf(log\_r | prior\_logr\_mean, prior\_logr\_sd) +}
\StringTok{           normal\_lpdf(log\_q | prior\_logq\_mean, prior\_logq\_sd);}
\StringTok{\}}
\StringTok{"}

\NormalTok{stan\_file\_proto\_single }\OtherTok{\textless{}{-}} \FunctionTok{here}\NormalTok{(}\StringTok{"stan"}\NormalTok{, }\StringTok{"ch13\_prototype\_single.stan"}\NormalTok{)}
\FunctionTok{write\_stan\_file}\NormalTok{(prototype\_single\_stan, }\AttributeTok{dir =} \FunctionTok{here}\NormalTok{(}\StringTok{"stan"}\NormalTok{),}
                \AttributeTok{basename =} \StringTok{"ch13\_prototype\_single.stan"}\NormalTok{)}
\end{Highlighting}
\end{Shaded}

\begin{verbatim}
[1] "/Users/au209589/Dropbox/Teaching/AdvancedCognitiveModeling23_book/stan/ch13_prototype_single.stan"
\end{verbatim}

\begin{Shaded}
\begin{Highlighting}[]
\NormalTok{mod\_prototype\_single }\OtherTok{\textless{}{-}} \FunctionTok{cmdstan\_model}\NormalTok{(stan\_file\_proto\_single)}
\FunctionTok{cat}\NormalTok{(}\StringTok{"Stan model compiled successfully.}\SpecialCharTok{\textbackslash{}n}\StringTok{"}\NormalTok{)}
\end{Highlighting}
\end{Shaded}

\begin{verbatim}
Stan model compiled successfully.
\end{verbatim}

\textbf{Key aspects of the Stan implementation:}

\begin{itemize}
\tightlist
\item
  \textbf{Parameters block}: Two scalars \texttt{log\_r} and
  \texttt{log\_q}. Positivity constraints are handled by \texttt{exp()}
  in \texttt{transformed\ parameters}.
\item
  \textbf{Prediction step}: \texttt{sigma\_cat0} and
  \texttt{sigma\_cat1} each have \texttt{q\_matrix} added before the
  decision. With \texttt{q\_value\ =\ 0} this recovers the static-target
  filter.
\item
  \textbf{Transformed parameters}: Runs the full Kalman filter
  trial-by-trial, updating the two prototype distributions and recording
  \texttt{p{[}i{]}} --- the probability of choosing Category 1 on trial
  \(i\) --- before each trial's update. The update happens on \emph{all}
  trials (including the last), consistent with the R simulation.
\item
  \textbf{Model block}: Priors on \texttt{log\_r} and \texttt{log\_q}
  plus a vectorised Bernoulli likelihood using the pre-computed
  \texttt{p} vector. One line.
\item
  \textbf{Generated quantities}: \texttt{log\_lik{[}i{]}} reads directly
  from \texttt{p{[}i{]}} --- no re-computation of the Kalman filter
  required.
\end{itemize}

\begin{center}\rule{0.5\linewidth}{0.5pt}\end{center}

\section{Fitting the Model to Simulated
Data}\label{fitting-the-model-to-simulated-data}

We use Pathfinder initialization before HMC sampling:

\begin{Shaded}
\begin{Highlighting}[]
\CommentTok{\# Select one simulated agent with r\_value = 2.0, q\_value = 0.1}
\NormalTok{agent\_to\_fit }\OtherTok{\textless{}{-}}\NormalTok{ prototype\_responses }\SpecialCharTok{|\textgreater{}}
  \FunctionTok{filter}\NormalTok{(r\_value\_true }\SpecialCharTok{==} \FloatTok{2.0}\NormalTok{, q\_value\_true }\SpecialCharTok{==} \FloatTok{0.1}\NormalTok{, agent\_id }\SpecialCharTok{==} \DecValTok{1}\NormalTok{)}

\FunctionTok{stopifnot}\NormalTok{(}\FunctionTok{nrow}\NormalTok{(agent\_to\_fit) }\SpecialCharTok{==}\NormalTok{ total\_trials)}

\NormalTok{proto\_data\_single }\OtherTok{\textless{}{-}} \FunctionTok{list}\NormalTok{(}
  \AttributeTok{ntrials            =} \FunctionTok{nrow}\NormalTok{(agent\_to\_fit),}
  \AttributeTok{nfeatures          =} \DecValTok{2}\NormalTok{L,}
  \AttributeTok{cat\_one            =}\NormalTok{ agent\_to\_fit}\SpecialCharTok{$}\NormalTok{category\_feedback,}
  \AttributeTok{y                  =}\NormalTok{ agent\_to\_fit}\SpecialCharTok{$}\NormalTok{sim\_response,}
  \AttributeTok{obs                =} \FunctionTok{as.matrix}\NormalTok{(agent\_to\_fit[, }\FunctionTok{c}\NormalTok{(}\StringTok{"height"}\NormalTok{, }\StringTok{"position"}\NormalTok{)]),}
  \AttributeTok{initial\_mu\_cat0    =} \FunctionTok{c}\NormalTok{(}\FloatTok{2.5}\NormalTok{, }\FloatTok{2.5}\NormalTok{),}
  \AttributeTok{initial\_mu\_cat1    =} \FunctionTok{c}\NormalTok{(}\FloatTok{2.5}\NormalTok{, }\FloatTok{2.5}\NormalTok{),}
  \AttributeTok{initial\_sigma\_diag =} \FloatTok{10.0}\NormalTok{,}
  \AttributeTok{prior\_logr\_mean    =}\NormalTok{ LOG\_R\_PRIOR\_MEAN,}
  \AttributeTok{prior\_logr\_sd      =}\NormalTok{ LOG\_R\_PRIOR\_SD,}
  \AttributeTok{prior\_logq\_mean    =}\NormalTok{ LOG\_Q\_PRIOR\_MEAN,}
  \AttributeTok{prior\_logq\_sd      =}\NormalTok{ LOG\_Q\_PRIOR\_SD}
\NormalTok{)}

\NormalTok{fit\_filepath\_single }\OtherTok{\textless{}{-}} \FunctionTok{here}\NormalTok{(}\StringTok{"simmodels"}\NormalTok{, }\StringTok{"ch13\_proto\_single\_fit.rds"}\NormalTok{)}
\ControlFlowTok{if}\NormalTok{ (regenerate\_simulations }\SpecialCharTok{||} \SpecialCharTok{!}\FunctionTok{file.exists}\NormalTok{(fit\_filepath\_single)) \{}
\NormalTok{  pf\_single }\OtherTok{\textless{}{-}}\NormalTok{ mod\_prototype\_single}\SpecialCharTok{$}\FunctionTok{pathfinder}\NormalTok{(}
    \AttributeTok{data =}\NormalTok{ proto\_data\_single, }\AttributeTok{seed =} \DecValTok{123}\NormalTok{, }\AttributeTok{num\_paths =} \DecValTok{4}\NormalTok{, }\AttributeTok{refresh =} \DecValTok{0}
\NormalTok{  )}

\NormalTok{  fit\_proto\_single }\OtherTok{\textless{}{-}}\NormalTok{ mod\_prototype\_single}\SpecialCharTok{$}\FunctionTok{sample}\NormalTok{(}
    \AttributeTok{data            =}\NormalTok{ proto\_data\_single,}
    \AttributeTok{init            =}\NormalTok{ pf\_single,}
    \AttributeTok{seed            =} \DecValTok{123}\NormalTok{,}
    \AttributeTok{chains          =} \DecValTok{4}\NormalTok{,}
    \AttributeTok{parallel\_chains =} \FunctionTok{min}\NormalTok{(}\DecValTok{4}\NormalTok{, }\FunctionTok{availableCores}\NormalTok{()),}
    \AttributeTok{iter\_warmup     =} \DecValTok{1000}\NormalTok{,}
    \AttributeTok{iter\_sampling   =} \DecValTok{1500}\NormalTok{,}
    \AttributeTok{refresh         =} \DecValTok{300}\NormalTok{,}
    \AttributeTok{adapt\_delta     =} \FloatTok{0.9}
\NormalTok{  )}
\NormalTok{  fit\_proto\_single}\SpecialCharTok{$}\FunctionTok{save\_object}\NormalTok{(fit\_filepath\_single)}
  \FunctionTok{cat}\NormalTok{(}\StringTok{"Single{-}agent prototype fit computed and saved.}\SpecialCharTok{\textbackslash{}n}\StringTok{"}\NormalTok{)}
\NormalTok{\} }\ControlFlowTok{else}\NormalTok{ \{}
\NormalTok{  fit\_proto\_single }\OtherTok{\textless{}{-}} \FunctionTok{readRDS}\NormalTok{(fit\_filepath\_single)}
  \FunctionTok{cat}\NormalTok{(}\StringTok{"Loaded existing single{-}agent prototype fit.}\SpecialCharTok{\textbackslash{}n}\StringTok{"}\NormalTok{)}
\NormalTok{\}}
\end{Highlighting}
\end{Shaded}

\begin{verbatim}
Loaded existing single-agent prototype fit.
\end{verbatim}

\begin{Shaded}
\begin{Highlighting}[]
\ControlFlowTok{if}\NormalTok{ (}\SpecialCharTok{!}\FunctionTok{is.null}\NormalTok{(fit\_proto\_single)) \{}
\NormalTok{  param\_summary }\OtherTok{\textless{}{-}}\NormalTok{ fit\_proto\_single}\SpecialCharTok{$}\FunctionTok{summary}\NormalTok{(}\AttributeTok{variables =} \FunctionTok{c}\NormalTok{(}\StringTok{"log\_r"}\NormalTok{, }\StringTok{"r\_value"}\NormalTok{, }\StringTok{"log\_q"}\NormalTok{, }\StringTok{"q\_value"}\NormalTok{))}
  \FunctionTok{print}\NormalTok{(param\_summary)}

\NormalTok{  true\_r     }\OtherTok{\textless{}{-}}\NormalTok{ agent\_to\_fit}\SpecialCharTok{$}\NormalTok{r\_value\_true[}\DecValTok{1}\NormalTok{]}
\NormalTok{  true\_log\_r }\OtherTok{\textless{}{-}} \FunctionTok{log}\NormalTok{(true\_r)}
\NormalTok{  true\_q     }\OtherTok{\textless{}{-}}\NormalTok{ agent\_to\_fit}\SpecialCharTok{$}\NormalTok{q\_value\_true[}\DecValTok{1}\NormalTok{]}
\NormalTok{  true\_log\_q }\OtherTok{\textless{}{-}} \FunctionTok{log}\NormalTok{(true\_q)}
  \FunctionTok{cat}\NormalTok{(}\StringTok{"}\SpecialCharTok{\textbackslash{}n}\StringTok{True log\_r ="}\NormalTok{, }\FunctionTok{round}\NormalTok{(true\_log\_r, }\DecValTok{3}\NormalTok{), }\StringTok{"  True r ="}\NormalTok{, true\_r, }\StringTok{"}\SpecialCharTok{\textbackslash{}n}\StringTok{"}\NormalTok{)}
  \FunctionTok{cat}\NormalTok{(}\StringTok{"True log\_q ="}\NormalTok{, }\FunctionTok{round}\NormalTok{(true\_log\_q, }\DecValTok{3}\NormalTok{), }\StringTok{"  True q ="}\NormalTok{, true\_q, }\StringTok{"}\SpecialCharTok{\textbackslash{}n}\StringTok{"}\NormalTok{)}
\NormalTok{\}}
\end{Highlighting}
\end{Shaded}

\begin{verbatim}
# A tibble: 4 x 10
  variable    mean  median    sd   mad      q5     q95  rhat ess_bulk ess_tail
  <chr>      <dbl>   <dbl> <dbl> <dbl>   <dbl>   <dbl> <dbl>    <dbl>    <dbl>
1 log_r     0.0948  0.0971 0.548 0.527 -0.808   0.991   1.00    1841.    1906.
2 r_value   1.28    1.10   0.749 0.565  0.446   2.69    1.00    1841.    1906.
3 log_q    -1.22   -1.12   0.817 0.776 -2.70   -0.0641  1.00    2026.    1926.
4 q_value   0.390   0.328  0.280 0.244  0.0670  0.938   1.00    2026.    1926.

True log_r = 0.693   True r = 2 
True log_q = -2.303   True q = 0.1 
\end{verbatim}

\subsection{MCMC Diagnostic Battery}\label{mcmc-diagnostic-battery-2}

Before reading anything off the posterior, we curate the diagnostic
table. Earlier drafts of this chapter ran \texttt{cmdstan\_diagnose()}
silently and moved on. The point of running diagnostics is \emph{seeing}
them --- even when (as here, with two parameters) they are nearly
trivial. The \texttt{diagnostic\_summary\_table()} helper defined in the
setup chunk is the canonical version shared across Chapters 11, 12, 13,
and 15. It returns a one-row tibble and emits a warning if any of the
six thresholds are breached.

\begin{Shaded}
\begin{Highlighting}[]
\ControlFlowTok{if}\NormalTok{ (}\SpecialCharTok{!}\FunctionTok{is.null}\NormalTok{(fit\_proto\_single)) \{}
\NormalTok{  diag\_tbl }\OtherTok{\textless{}{-}} \FunctionTok{diagnostic\_summary\_table}\NormalTok{(}
\NormalTok{    fit\_proto\_single,}
    \AttributeTok{params =} \FunctionTok{c}\NormalTok{(}\StringTok{"log\_r"}\NormalTok{, }\StringTok{"r\_value"}\NormalTok{, }\StringTok{"log\_q"}\NormalTok{, }\StringTok{"q\_value"}\NormalTok{)}
\NormalTok{  )}
  \FunctionTok{print}\NormalTok{(diag\_tbl)}
\NormalTok{\}}
\end{Highlighting}
\end{Shaded}

\begin{verbatim}
# A tibble: 6 x 4
  metric                           value threshold pass 
  <chr>                            <dbl> <chr>     <lgl>
1 Divergences (zero tolerance)    0      == 0      TRUE 
2 Max rank-normalised R-hat       1      < 1.01    TRUE 
3 Min bulk ESS                 1841      > 400     TRUE 
4 Min tail ESS                 1906      > 400     TRUE 
5 Min E-BFMI                      0.977  > 0.2     TRUE 
6 Max MCSE / posterior SD         0.0233 < 0.05    TRUE 
\end{verbatim}

With two parameters, \texttt{mcmc\_pairs} is informative for detecting
posterior correlations between \texttt{log\_r} and \texttt{log\_q}. The
trace plot grid plus rank histograms remain the primary MCMC
diagnostics.

\begin{Shaded}
\begin{Highlighting}[]
\ControlFlowTok{if}\NormalTok{ (}\SpecialCharTok{!}\FunctionTok{is.null}\NormalTok{(fit\_proto\_single)) \{}
\NormalTok{  p\_trace }\OtherTok{\textless{}{-}}\NormalTok{ bayesplot}\SpecialCharTok{::}\FunctionTok{mcmc\_trace}\NormalTok{(}
\NormalTok{    fit\_proto\_single}\SpecialCharTok{$}\FunctionTok{draws}\NormalTok{(}\FunctionTok{c}\NormalTok{(}\StringTok{"log\_r"}\NormalTok{, }\StringTok{"log\_q"}\NormalTok{)),}
    \AttributeTok{facet\_args =} \FunctionTok{list}\NormalTok{(}\AttributeTok{ncol =} \DecValTok{2}\NormalTok{)}
\NormalTok{  ) }\SpecialCharTok{+}
    \FunctionTok{ggtitle}\NormalTok{(}\StringTok{"Trace plots — should look like hairy caterpillars"}\NormalTok{)}

\NormalTok{  p\_rank }\OtherTok{\textless{}{-}}\NormalTok{ bayesplot}\SpecialCharTok{::}\FunctionTok{mcmc\_rank\_overlay}\NormalTok{(}
\NormalTok{    fit\_proto\_single}\SpecialCharTok{$}\FunctionTok{draws}\NormalTok{(}\FunctionTok{c}\NormalTok{(}\StringTok{"log\_r"}\NormalTok{, }\StringTok{"log\_q"}\NormalTok{)),}
    \AttributeTok{facet\_args =} \FunctionTok{list}\NormalTok{(}\AttributeTok{ncol =} \DecValTok{2}\NormalTok{)}
\NormalTok{  ) }\SpecialCharTok{+}
    \FunctionTok{ggtitle}\NormalTok{(}\StringTok{"Rank histograms — should be uniform across chains"}\NormalTok{)}

\NormalTok{  p\_pairs }\OtherTok{\textless{}{-}}\NormalTok{ bayesplot}\SpecialCharTok{::}\FunctionTok{mcmc\_pairs}\NormalTok{(}
\NormalTok{    fit\_proto\_single}\SpecialCharTok{$}\FunctionTok{draws}\NormalTok{(}\FunctionTok{c}\NormalTok{(}\StringTok{"log\_r"}\NormalTok{, }\StringTok{"log\_q"}\NormalTok{)),}
    \AttributeTok{diag\_fun     =} \StringTok{"dens"}\NormalTok{,}
    \AttributeTok{off\_diag\_fun =} \StringTok{"scatter"}\NormalTok{,}
    \AttributeTok{np           =} \FunctionTok{nuts\_params}\NormalTok{(fit\_proto\_single)}
\NormalTok{  )}

  \FunctionTok{print}\NormalTok{(p\_trace }\SpecialCharTok{/}\NormalTok{ p\_rank)}
  \FunctionTok{print}\NormalTok{(p\_pairs)}
\NormalTok{\}}
\end{Highlighting}
\end{Shaded}

\begin{center}
\includegraphics[width=0.8\linewidth,height=\textheight,keepaspectratio]{12-CategorizationModelsPrototypes_files/figure-pdf/ch13_trace_and_rank-1.pdf}
\end{center}

\begin{center}
\includegraphics[width=0.8\linewidth,height=\textheight,keepaspectratio]{12-CategorizationModelsPrototypes_files/figure-pdf/ch13_trace_and_rank-2.pdf}
\end{center}

\subsection{Posterior Predictive
Check}\label{posterior-predictive-check}

The posterior predictive check asks: does the fitted model, when run
forward with parameters drawn from the posterior, produce data that look
like the observed data? Because \texttt{p} is already in
\texttt{transformed\ parameters}, we can read posterior draws of the
choice probability directly and use them to generate posterior
predictive responses.

\begin{Shaded}
\begin{Highlighting}[]
\ControlFlowTok{if}\NormalTok{ (}\SpecialCharTok{!}\FunctionTok{is.null}\NormalTok{(fit\_proto\_single)) \{}
  \CommentTok{\# Posterior draws of p[i] (one row per draw, one column per trial)}
\NormalTok{  p\_draws }\OtherTok{\textless{}{-}}\NormalTok{ fit\_proto\_single}\SpecialCharTok{$}\FunctionTok{draws}\NormalTok{(}\StringTok{"p"}\NormalTok{, }\AttributeTok{format =} \StringTok{"matrix"}\NormalTok{)}
\NormalTok{  n\_draws }\OtherTok{\textless{}{-}} \FunctionTok{nrow}\NormalTok{(p\_draws)}
\NormalTok{  n\_obs   }\OtherTok{\textless{}{-}} \FunctionTok{ncol}\NormalTok{(p\_draws)}

  \CommentTok{\# Generate posterior predictive responses}
\NormalTok{  yrep }\OtherTok{\textless{}{-}} \FunctionTok{matrix}\NormalTok{(}\FunctionTok{rbinom}\NormalTok{(n\_draws }\SpecialCharTok{*}\NormalTok{ n\_obs, }\DecValTok{1}\NormalTok{, }\FunctionTok{as.vector}\NormalTok{(p\_draws)),}
                 \AttributeTok{nrow =}\NormalTok{ n\_draws, }\AttributeTok{ncol =}\NormalTok{ n\_obs)}

  \CommentTok{\# Cumulative{-}accuracy ribbons vs. observed}
\NormalTok{  obs\_cum\_acc }\OtherTok{\textless{}{-}} \FunctionTok{cumsum}\NormalTok{(agent\_to\_fit}\SpecialCharTok{$}\NormalTok{sim\_response }\SpecialCharTok{==}\NormalTok{ agent\_to\_fit}\SpecialCharTok{$}\NormalTok{category\_feedback) }\SpecialCharTok{/}
                 \FunctionTok{seq\_len}\NormalTok{(n\_obs)}
                 
  \CommentTok{\# Use "==" directly and add 0 to coerce to numeric without losing matrix dims}
\NormalTok{  yrep\_correct }\OtherTok{\textless{}{-}} \FunctionTok{sweep}\NormalTok{(yrep, }\DecValTok{2}\NormalTok{, }\FunctionTok{as.numeric}\NormalTok{(agent\_to\_fit}\SpecialCharTok{$}\NormalTok{category\_feedback), }\StringTok{"=="}\NormalTok{) }\SpecialCharTok{+} \DecValTok{0}
  
\NormalTok{  yrep\_cum\_acc }\OtherTok{\textless{}{-}} \FunctionTok{t}\NormalTok{(}\FunctionTok{apply}\NormalTok{(yrep\_correct, }\DecValTok{1}\NormalTok{, }\ControlFlowTok{function}\NormalTok{(r) }\FunctionTok{cumsum}\NormalTok{(r) }\SpecialCharTok{/} \FunctionTok{seq\_along}\NormalTok{(r)))}

\NormalTok{  cum\_acc\_summary }\OtherTok{\textless{}{-}} \FunctionTok{tibble}\NormalTok{(}
    \AttributeTok{trial =} \FunctionTok{seq\_len}\NormalTok{(n\_obs),}
    \AttributeTok{q05   =} \FunctionTok{apply}\NormalTok{(yrep\_cum\_acc, }\DecValTok{2}\NormalTok{, quantile, }\FloatTok{0.05}\NormalTok{),}
    \AttributeTok{q25   =} \FunctionTok{apply}\NormalTok{(yrep\_cum\_acc, }\DecValTok{2}\NormalTok{, quantile, }\FloatTok{0.25}\NormalTok{),}
    \AttributeTok{q50   =} \FunctionTok{apply}\NormalTok{(yrep\_cum\_acc, }\DecValTok{2}\NormalTok{, quantile, }\FloatTok{0.50}\NormalTok{),}
    \AttributeTok{q75   =} \FunctionTok{apply}\NormalTok{(yrep\_cum\_acc, }\DecValTok{2}\NormalTok{, quantile, }\FloatTok{0.75}\NormalTok{),}
    \AttributeTok{q95   =} \FunctionTok{apply}\NormalTok{(yrep\_cum\_acc, }\DecValTok{2}\NormalTok{, quantile, }\FloatTok{0.95}\NormalTok{),}
    \AttributeTok{obs   =}\NormalTok{ obs\_cum\_acc}
\NormalTok{  )}

  \FunctionTok{ggplot}\NormalTok{(cum\_acc\_summary, }\FunctionTok{aes}\NormalTok{(}\AttributeTok{x =}\NormalTok{ trial)) }\SpecialCharTok{+}
    \FunctionTok{geom\_ribbon}\NormalTok{(}\FunctionTok{aes}\NormalTok{(}\AttributeTok{ymin =}\NormalTok{ q05, }\AttributeTok{ymax =}\NormalTok{ q95), }\AttributeTok{fill =} \StringTok{"\#56B4E9"}\NormalTok{, }\AttributeTok{alpha =} \FloatTok{0.20}\NormalTok{) }\SpecialCharTok{+}
    \FunctionTok{geom\_ribbon}\NormalTok{(}\FunctionTok{aes}\NormalTok{(}\AttributeTok{ymin =}\NormalTok{ q25, }\AttributeTok{ymax =}\NormalTok{ q75), }\AttributeTok{fill =} \StringTok{"\#56B4E9"}\NormalTok{, }\AttributeTok{alpha =} \FloatTok{0.40}\NormalTok{) }\SpecialCharTok{+}
    \FunctionTok{geom\_line}\NormalTok{(}\FunctionTok{aes}\NormalTok{(}\AttributeTok{y =}\NormalTok{ q50), }\AttributeTok{color =} \StringTok{"\#0072B2"}\NormalTok{, }\AttributeTok{linewidth =} \DecValTok{1}\NormalTok{) }\SpecialCharTok{+}
    \FunctionTok{geom\_line}\NormalTok{(}\FunctionTok{aes}\NormalTok{(}\AttributeTok{y =}\NormalTok{ obs), }\AttributeTok{color =} \StringTok{"\#D55E00"}\NormalTok{, }\AttributeTok{linewidth =} \DecValTok{1}\NormalTok{) }\SpecialCharTok{+}
    \FunctionTok{geom\_hline}\NormalTok{(}\AttributeTok{yintercept =} \FloatTok{0.5}\NormalTok{, }\AttributeTok{linetype =} \StringTok{"dashed"}\NormalTok{, }\AttributeTok{color =} \StringTok{"grey50"}\NormalTok{) }\SpecialCharTok{+}
    \FunctionTok{labs}\NormalTok{(}
      \AttributeTok{title    =} \StringTok{"Posterior Predictive Check: Cumulative Accuracy"}\NormalTok{,}
      \AttributeTok{subtitle =} \StringTok{"Blue: posterior predictive (50\% and 90\% intervals + median)}\SpecialCharTok{\textbackslash{}n}\StringTok{Orange: observed"}\NormalTok{,}
      \AttributeTok{x =} \StringTok{"Trial"}\NormalTok{, }\AttributeTok{y =} \StringTok{"Cumulative accuracy"}
\NormalTok{    ) }\SpecialCharTok{+}
    \FunctionTok{coord\_cartesian}\NormalTok{(}\AttributeTok{ylim =} \FunctionTok{c}\NormalTok{(}\FloatTok{0.4}\NormalTok{, }\DecValTok{1}\NormalTok{))}
\NormalTok{\}}
\end{Highlighting}
\end{Shaded}

\begin{center}
\includegraphics[width=0.8\linewidth,height=\textheight,keepaspectratio]{12-CategorizationModelsPrototypes_files/figure-pdf/ch13_ppc-1.pdf}
\end{center}

The observed cumulative-accuracy curve should sit comfortably inside the
posterior predictive band. If it drifts to the edge, the model is
missing structure that the data carry --- the most likely candidate
being a separate decisional-noise term that the single \(r\) parameter
cannot capture.

\subsection{Prior--Posterior Update
Plot}\label{priorposterior-update-plot-1}

Did the data teach the model anything? We need an explicit overlay of
prior and posterior densities, on both the \emph{computational scale}
(\texttt{log\_r}, where the prior was defined) and the
\emph{interpretable scale} (\texttt{r}, where the cognitive
interpretation lives). If posterior and prior overlap substantially, the
data were uninformative for that parameter. If posterior is narrow but
centred far from the truth, the model is confidently wrong.

\begin{Shaded}
\begin{Highlighting}[]
\ControlFlowTok{if}\NormalTok{ (}\SpecialCharTok{!}\FunctionTok{is.null}\NormalTok{(fit\_proto\_single)) \{}
\NormalTok{  draws\_df }\OtherTok{\textless{}{-}} \FunctionTok{as\_draws\_df}\NormalTok{(}
\NormalTok{    fit\_proto\_single}\SpecialCharTok{$}\FunctionTok{draws}\NormalTok{(}\FunctionTok{c}\NormalTok{(}\StringTok{"log\_r"}\NormalTok{, }\StringTok{"r\_value"}\NormalTok{, }\StringTok{"log\_q"}\NormalTok{, }\StringTok{"q\_value"}\NormalTok{))}
\NormalTok{  )}

  \FunctionTok{set.seed}\NormalTok{(}\DecValTok{7}\NormalTok{)}
\NormalTok{  n\_prior    }\OtherTok{\textless{}{-}} \FunctionTok{nrow}\NormalTok{(draws\_df)}
\NormalTok{  prior\_logr }\OtherTok{\textless{}{-}} \FunctionTok{rnorm}\NormalTok{(n\_prior, LOG\_R\_PRIOR\_MEAN, LOG\_R\_PRIOR\_SD)}
\NormalTok{  prior\_r    }\OtherTok{\textless{}{-}} \FunctionTok{exp}\NormalTok{(prior\_logr)}
\NormalTok{  prior\_logq }\OtherTok{\textless{}{-}} \FunctionTok{rnorm}\NormalTok{(n\_prior, LOG\_Q\_PRIOR\_MEAN, LOG\_Q\_PRIOR\_SD)}
\NormalTok{  prior\_q    }\OtherTok{\textless{}{-}} \FunctionTok{exp}\NormalTok{(prior\_logq)}

\NormalTok{  prior\_post\_df }\OtherTok{\textless{}{-}} \FunctionTok{bind\_rows}\NormalTok{(}
    \FunctionTok{tibble}\NormalTok{(}\AttributeTok{scale =} \StringTok{"log\_r"}\NormalTok{, }\AttributeTok{value =}\NormalTok{ prior\_logr,    }\AttributeTok{source =} \StringTok{"prior"}\NormalTok{),}
    \FunctionTok{tibble}\NormalTok{(}\AttributeTok{scale =} \StringTok{"log\_r"}\NormalTok{, }\AttributeTok{value =}\NormalTok{ draws\_df}\SpecialCharTok{$}\NormalTok{log\_r, }\AttributeTok{source =} \StringTok{"posterior"}\NormalTok{),}
    \FunctionTok{tibble}\NormalTok{(}\AttributeTok{scale =} \StringTok{"r"}\NormalTok{,     }\AttributeTok{value =}\NormalTok{ prior\_r,        }\AttributeTok{source =} \StringTok{"prior"}\NormalTok{),}
    \FunctionTok{tibble}\NormalTok{(}\AttributeTok{scale =} \StringTok{"r"}\NormalTok{,     }\AttributeTok{value =}\NormalTok{ draws\_df}\SpecialCharTok{$}\NormalTok{r\_value, }\AttributeTok{source =} \StringTok{"posterior"}\NormalTok{),}
    \FunctionTok{tibble}\NormalTok{(}\AttributeTok{scale =} \StringTok{"log\_q"}\NormalTok{, }\AttributeTok{value =}\NormalTok{ prior\_logq,    }\AttributeTok{source =} \StringTok{"prior"}\NormalTok{),}
    \FunctionTok{tibble}\NormalTok{(}\AttributeTok{scale =} \StringTok{"log\_q"}\NormalTok{, }\AttributeTok{value =}\NormalTok{ draws\_df}\SpecialCharTok{$}\NormalTok{log\_q, }\AttributeTok{source =} \StringTok{"posterior"}\NormalTok{),}
    \FunctionTok{tibble}\NormalTok{(}\AttributeTok{scale =} \StringTok{"q"}\NormalTok{,     }\AttributeTok{value =}\NormalTok{ prior\_q,        }\AttributeTok{source =} \StringTok{"prior"}\NormalTok{),}
    \FunctionTok{tibble}\NormalTok{(}\AttributeTok{scale =} \StringTok{"q"}\NormalTok{,     }\AttributeTok{value =}\NormalTok{ draws\_df}\SpecialCharTok{$}\NormalTok{q\_value, }\AttributeTok{source =} \StringTok{"posterior"}\NormalTok{)}
\NormalTok{  ) }\SpecialCharTok{|\textgreater{}}
    \FunctionTok{mutate}\NormalTok{(}\AttributeTok{scale =} \FunctionTok{factor}\NormalTok{(scale, }\AttributeTok{levels =} \FunctionTok{c}\NormalTok{(}\StringTok{"log\_r"}\NormalTok{, }\StringTok{"r"}\NormalTok{, }\StringTok{"log\_q"}\NormalTok{, }\StringTok{"q"}\NormalTok{)))}

\NormalTok{  true\_lines }\OtherTok{\textless{}{-}} \FunctionTok{tibble}\NormalTok{(}
    \AttributeTok{scale =} \FunctionTok{factor}\NormalTok{(}\FunctionTok{c}\NormalTok{(}\StringTok{"log\_r"}\NormalTok{, }\StringTok{"r"}\NormalTok{, }\StringTok{"log\_q"}\NormalTok{, }\StringTok{"q"}\NormalTok{),}
                   \AttributeTok{levels =} \FunctionTok{c}\NormalTok{(}\StringTok{"log\_r"}\NormalTok{, }\StringTok{"r"}\NormalTok{, }\StringTok{"log\_q"}\NormalTok{, }\StringTok{"q"}\NormalTok{)),}
    \AttributeTok{value =} \FunctionTok{c}\NormalTok{(true\_log\_r, true\_r, true\_log\_q, true\_q)}
\NormalTok{  )}

  \FunctionTok{ggplot}\NormalTok{(prior\_post\_df, }\FunctionTok{aes}\NormalTok{(}\AttributeTok{x =}\NormalTok{ value, }\AttributeTok{fill =}\NormalTok{ source, }\AttributeTok{color =}\NormalTok{ source)) }\SpecialCharTok{+}
    \FunctionTok{geom\_density}\NormalTok{(}\AttributeTok{alpha =} \FloatTok{0.4}\NormalTok{, }\AttributeTok{linewidth =} \FloatTok{0.6}\NormalTok{) }\SpecialCharTok{+}
    \FunctionTok{geom\_vline}\NormalTok{(}\AttributeTok{data =}\NormalTok{ true\_lines,}
               \FunctionTok{aes}\NormalTok{(}\AttributeTok{xintercept =}\NormalTok{ value),}
               \AttributeTok{color =} \StringTok{"red"}\NormalTok{, }\AttributeTok{linetype =} \StringTok{"dashed"}\NormalTok{) }\SpecialCharTok{+}
    \FunctionTok{facet\_wrap}\NormalTok{(}\SpecialCharTok{\textasciitilde{}}\NormalTok{ scale, }\AttributeTok{scales =} \StringTok{"free"}\NormalTok{, }\AttributeTok{ncol =} \DecValTok{2}\NormalTok{) }\SpecialCharTok{+}
    \FunctionTok{scale\_fill\_manual}\NormalTok{(}\AttributeTok{values =} \FunctionTok{c}\NormalTok{(}\AttributeTok{prior =} \StringTok{"\#999999"}\NormalTok{, }\AttributeTok{posterior =} \StringTok{"\#009E73"}\NormalTok{)) }\SpecialCharTok{+}
    \FunctionTok{scale\_color\_manual}\NormalTok{(}\AttributeTok{values =} \FunctionTok{c}\NormalTok{(}\AttributeTok{prior =} \StringTok{"\#666666"}\NormalTok{, }\AttributeTok{posterior =} \StringTok{"\#006D4E"}\NormalTok{)) }\SpecialCharTok{+}
    \FunctionTok{labs}\NormalTok{(}
      \AttributeTok{title    =} \StringTok{"Prior → Posterior Update: Kalman Prototype Model (r and q)"}\NormalTok{,}
      \AttributeTok{subtitle =} \StringTok{"Red dashed line = true generating value"}\NormalTok{,}
      \AttributeTok{x =} \ConstantTok{NULL}\NormalTok{, }\AttributeTok{y =} \StringTok{"Density"}
\NormalTok{    )}
\NormalTok{\}}
\end{Highlighting}
\end{Shaded}

\begin{center}
\includegraphics[width=0.8\linewidth,height=\textheight,keepaspectratio]{12-CategorizationModelsPrototypes_files/figure-pdf/ch13_prior_posterior_update-1.pdf}
\end{center}

We expect the posterior on the interpretable scale to be substantially
narrower than the prior and to overlap the true value. The \texttt{r}
panel uses a heavy-tailed prior (the exponential of a Gaussian) so the
prior density is wide and right-skewed; the posterior should look like a
clearly localised peak by comparison. \emph{Information gain is what
this plot makes visual} --- our preferred way of saying ``the data were
informative.''

\subsection{Randomised LOO-PIT
Calibration}\label{randomised-loo-pit-calibration}

For continuous outcomes the preferred calibration check is LOO-PIT. For
binary outcomes the equivalent is the \emph{randomised} LOO-PIT
(Säilynoja et al.~2022), which avoids the degenerate two-bin
distribution that plain LOO-PIT produces on Bernoulli data. Chapter 11
implements this for the GCM; we apply the same construction here.

\begin{Shaded}
\begin{Highlighting}[]
\ControlFlowTok{if}\NormalTok{ (}\SpecialCharTok{!}\FunctionTok{is.null}\NormalTok{(fit\_proto\_single)) \{}
  \CommentTok{\# PSIS{-}LOO weights from log\_lik}
\NormalTok{  log\_lik\_mat }\OtherTok{\textless{}{-}}\NormalTok{ fit\_proto\_single}\SpecialCharTok{$}\FunctionTok{draws}\NormalTok{(}\StringTok{"log\_lik"}\NormalTok{, }\AttributeTok{format =} \StringTok{"matrix"}\NormalTok{)}
  
  \CommentTok{\# FIX: Add save\_psis = TRUE so the weight matrix is retained}
\NormalTok{  loo\_obj     }\OtherTok{\textless{}{-}}\NormalTok{ loo}\SpecialCharTok{::}\FunctionTok{loo}\NormalTok{(log\_lik\_mat, }\AttributeTok{cores =} \DecValTok{4}\NormalTok{, }\AttributeTok{save\_psis =} \ConstantTok{TRUE}\NormalTok{)}
\NormalTok{  log\_w       }\OtherTok{\textless{}{-}}\NormalTok{ loo}\SpecialCharTok{::}\FunctionTok{weights.importance\_sampling}\NormalTok{(loo\_obj}\SpecialCharTok{$}\NormalTok{psis\_object, }
                                                  \AttributeTok{normalize =} \ConstantTok{TRUE}\NormalTok{, }
                                                  \AttributeTok{log =} \ConstantTok{TRUE}\NormalTok{)}

  \CommentTok{\# Posterior draws of choice probability}
\NormalTok{  p\_draws }\OtherTok{\textless{}{-}}\NormalTok{ fit\_proto\_single}\SpecialCharTok{$}\FunctionTok{draws}\NormalTok{(}\StringTok{"p"}\NormalTok{, }\AttributeTok{format =} \StringTok{"matrix"}\NormalTok{)}

  \CommentTok{\# Per{-}observation LOO predictive probability of cat 1}
\NormalTok{  loo\_p\_cat1 }\OtherTok{\textless{}{-}} \FunctionTok{vapply}\NormalTok{(}\FunctionTok{seq\_len}\NormalTok{(}\FunctionTok{ncol}\NormalTok{(p\_draws)), }\ControlFlowTok{function}\NormalTok{(i) \{}
    \FunctionTok{sum}\NormalTok{(}\FunctionTok{exp}\NormalTok{(log\_w[, i]) }\SpecialCharTok{*}\NormalTok{ p\_draws[, i])}
\NormalTok{  \}, }\FunctionTok{numeric}\NormalTok{(}\DecValTok{1}\NormalTok{))}

\NormalTok{  y\_obs }\OtherTok{\textless{}{-}}\NormalTok{ agent\_to\_fit}\SpecialCharTok{$}\NormalTok{sim\_response}

  \CommentTok{\# Randomised PIT for binary y:}
  \CommentTok{\#   if y == 1: U \textasciitilde{} Uniform(1 {-} p, 1)}
  \CommentTok{\#   if y == 0: U \textasciitilde{} Uniform(0, 1 {-} p)}
  \FunctionTok{set.seed}\NormalTok{(}\DecValTok{2026}\NormalTok{)}
\NormalTok{  rloo\_pit }\OtherTok{\textless{}{-}} \FunctionTok{vapply}\NormalTok{(}\FunctionTok{seq\_along}\NormalTok{(y\_obs), }\ControlFlowTok{function}\NormalTok{(i) \{}
\NormalTok{    p\_i }\OtherTok{\textless{}{-}}\NormalTok{ loo\_p\_cat1[i]}
    \ControlFlowTok{if}\NormalTok{ (y\_obs[i] }\SpecialCharTok{==} \DecValTok{1}\NormalTok{) }\FunctionTok{runif}\NormalTok{(}\DecValTok{1}\NormalTok{, }\DecValTok{1} \SpecialCharTok{{-}}\NormalTok{ p\_i, }\DecValTok{1}\NormalTok{) }\ControlFlowTok{else} \FunctionTok{runif}\NormalTok{(}\DecValTok{1}\NormalTok{, }\DecValTok{0}\NormalTok{, }\DecValTok{1} \SpecialCharTok{{-}}\NormalTok{ p\_i)}
\NormalTok{  \}, }\FunctionTok{numeric}\NormalTok{(}\DecValTok{1}\NormalTok{))}

\NormalTok{  p\_pit }\OtherTok{\textless{}{-}} \FunctionTok{ggplot}\NormalTok{(}\FunctionTok{tibble}\NormalTok{(}\AttributeTok{rloo\_pit =}\NormalTok{ rloo\_pit), }\FunctionTok{aes}\NormalTok{(}\AttributeTok{x =}\NormalTok{ rloo\_pit)) }\SpecialCharTok{+}
    \FunctionTok{geom\_histogram}\NormalTok{(}\FunctionTok{aes}\NormalTok{(}\AttributeTok{y =} \FunctionTok{after\_stat}\NormalTok{(density)), }\AttributeTok{bins =} \DecValTok{20}\NormalTok{,}
                   \AttributeTok{fill =} \StringTok{"\#56B4E9"}\NormalTok{, }\AttributeTok{color =} \StringTok{"white"}\NormalTok{) }\SpecialCharTok{+}
    \FunctionTok{geom\_hline}\NormalTok{(}\AttributeTok{yintercept =} \DecValTok{1}\NormalTok{, }\AttributeTok{linetype =} \StringTok{"dashed"}\NormalTok{, }\AttributeTok{color =} \StringTok{"red"}\NormalTok{) }\SpecialCharTok{+}
    \FunctionTok{labs}\NormalTok{(}
      \AttributeTok{title    =} \StringTok{"Randomised LOO{-}PIT: Kalman Prototype Model"}\NormalTok{,}
      \AttributeTok{subtitle =} \StringTok{"If the model is calibrated, the histogram should look uniform on [0, 1]"}\NormalTok{,}
      \AttributeTok{x =} \StringTok{"Randomised PIT value"}\NormalTok{, }\AttributeTok{y =} \StringTok{"Density"}
\NormalTok{    ) }\SpecialCharTok{+}
    \FunctionTok{coord\_cartesian}\NormalTok{(}\AttributeTok{ylim =} \FunctionTok{c}\NormalTok{(}\DecValTok{0}\NormalTok{, }\FloatTok{2.5}\NormalTok{))}
    
  \FunctionTok{print}\NormalTok{(p\_pit)}
\NormalTok{\}}
\end{Highlighting}
\end{Shaded}

\begin{center}
\includegraphics[width=0.8\linewidth,height=\textheight,keepaspectratio]{12-CategorizationModelsPrototypes_files/figure-pdf/ch13_loo_pit_randomized-1.pdf}
\end{center}

A roughly uniform histogram is the calibration target. Systematic
deviations (a U-shape, a peak in the middle, or skew) flag specific
failure modes: U-shape means the posterior is overconfident (the model
thinks it knows the answer better than it does), a central hump means
underconfident, and a one-sided skew means there is a directional bias
the model has not captured.

\subsection{LOO with Pareto-k as a Diagnostic, Not a Predictive
Score}\label{loo-with-pareto-k-as-a-diagnostic-not-a-predictive-score}

\begin{Shaded}
\begin{Highlighting}[]
\ControlFlowTok{if}\NormalTok{ (}\SpecialCharTok{!}\FunctionTok{is.null}\NormalTok{(fit\_proto\_single)) \{}
\NormalTok{  loo\_proto }\OtherTok{\textless{}{-}}\NormalTok{ loo}\SpecialCharTok{::}\FunctionTok{loo}\NormalTok{(log\_lik\_mat, }\AttributeTok{cores =} \DecValTok{4}\NormalTok{)}
  \FunctionTok{print}\NormalTok{(loo\_proto)}
  \FunctionTok{plot}\NormalTok{(loo\_proto, }\AttributeTok{label\_points =} \ConstantTok{TRUE}\NormalTok{,}
       \AttributeTok{main =} \StringTok{"PSIS{-}LOO Pareto{-}k: Kalman Prototype Model"}\NormalTok{)}
\NormalTok{\}}
\end{Highlighting}
\end{Shaded}

\begin{verbatim}

Computed from 6000 by 64 log-likelihood matrix.

         Estimate  SE
elpd_loo    -38.8 3.5
p_loo         1.1 0.2
looic        77.6 7.0
------
MCSE of elpd_loo is NA.
MCSE and ESS estimates assume independent draws (r_eff=1).

Pareto k diagnostic values:
                         Count Pct.    Min. ESS
(-Inf, 0.7]   (good)     63    98.4%   4987    
   (0.7, 1]   (bad)       0     0.0%   <NA>    
   (1, Inf)   (very bad)  1     1.6%   <NA>    
See help('pareto-k-diagnostic') for details.
\end{verbatim}

\begin{center}
\includegraphics[width=0.8\linewidth,height=\textheight,keepaspectratio]{12-CategorizationModelsPrototypes_files/figure-pdf/ch13_loo_single-1.pdf}
\end{center}

\begin{quote}
\textbf{⚠️ Reframing PSIS-LOO for sequential models.} The Kalman filter
is path-dependent: trial \(t\)'s choice probability depends on the
entire filter state from trials \(1, \ldots, t-1\). Standard PSIS-LOO is
\emph{not} a valid measure of out-of-sample predictive accuracy here ---
leaving out trial \(t\) implicitly conditions on a counterfactual filter
state that the model never actually saw. We compute it anyway, for two
reasons: (1) the \textbf{Pareto-\(\hat{k}\) values are still informative
as a localisation diagnostic} --- high \(\hat{k}\) at specific trials
flags observations the model cannot easily reconstruct from the rest,
which can reveal where the structural assumptions are stressed; (2)
compatibility with downstream tooling. The actual predictive evaluation
is done by LFO-CV in the next subsection. The numeric ELPD reported by
\texttt{loo()} should not be taken as a generalisation estimate.
\end{quote}

\subsection{Leave-Future-Out Cross-Validation: Inheriting and
Generalising the Chapter 11
Implementation}\label{leave-future-out-cross-validation-inheriting-and-generalising-the-chapter-11-implementation}

For sequential models, LOO is invalid and LFO-CV is the correct
alternative (Bürkner, Gabry, \& Vehtari 2020). Chapter 11 implements
PSIS-LFO for the GCM as \texttt{psis\_lfo\_gcm()}. The Kalman prototype
model has the same path-dependence structure but a much cheaper
per-trial cost, so LFO-CV is even more tractable here. Below is the same
algorithm specialised to the prototype model and saved as
\texttt{psis\_lfo\_kalman()}. The two functions share an interface so
Chapter 15 can call either through a common front-end.

The algorithm (Bürkner et al.~2020 Algorithm 1):

\begin{enumerate}
\def\labelenumi{\arabic{enumi}.}
\tightlist
\item
  Choose a minimum training window \(L\) (we use \(L = 16\), two blocks
  of the Kruschke task).
\item
  Fit the model once on \(y_{1:L}\). This is the \emph{reference fit}.
\item
  For each \(t = L+1, \ldots, T\):

  \begin{enumerate}
  \def\labelenumii{\alph{enumii}.}
  \tightlist
  \item
    Use the reference fit's posterior draws \(\theta^{(s)}\) to compute
    the importance ratio for re-weighting \(\theta\) from
    \(p(\theta \mid y_{1:t-1})\) relative to the reference.
  \item
    Smooth the ratios with PSIS, check Pareto \(\hat{k}\).
  \item
    If \(\hat{k} \leq 0.7\), accept the ratios and use them to evaluate
    \(\log p(y_t \mid y_{1:t-1})\).
  \item
    If \(\hat{k} > 0.7\), \textbf{refit} the model on \(y_{1:t-1}\),
    update the reference, and continue.
  \end{enumerate}
\item
  Return pointwise ELPD, the refit history, and the \(\hat{k}\)
  trajectory.
\end{enumerate}

\begin{Shaded}
\begin{Highlighting}[]
\CommentTok{\# psis\_lfo\_kalman(): Leave{-}future{-}out CV for the Kalman prototype model.}
\CommentTok{\# Mirrors psis\_lfo\_gcm() from Chapter 11. Returns pointwise ELPD across the}
\CommentTok{\# evaluation window, the indices at which a refit was triggered, and the}
\CommentTok{\# Pareto{-}k trajectory.}
\NormalTok{psis\_lfo\_kalman }\OtherTok{\textless{}{-}} \ControlFlowTok{function}\NormalTok{(stan\_data,}
\NormalTok{                            stan\_model,}
                            \AttributeTok{L           =} \DecValTok{16}\NormalTok{,    }\CommentTok{\# min training window}
                            \AttributeTok{k\_thresh    =} \FloatTok{0.7}\NormalTok{,   }\CommentTok{\# PSIS{-}k refit threshold}
                            \AttributeTok{iter\_warmup =} \DecValTok{800}\NormalTok{,}
                            \AttributeTok{iter\_sample =} \DecValTok{1000}\NormalTok{,}
                            \AttributeTok{chains      =} \DecValTok{2}\NormalTok{,}
                            \AttributeTok{seed        =} \DecValTok{1}\NormalTok{) \{}
\NormalTok{  T\_total }\OtherTok{\textless{}{-}}\NormalTok{ stan\_data}\SpecialCharTok{$}\NormalTok{ntrials}

  \CommentTok{\# Helper: build a stan\_data list restricted to trials 1:t}
\NormalTok{  truncate\_data }\OtherTok{\textless{}{-}} \ControlFlowTok{function}\NormalTok{(d, t) \{}
\NormalTok{    out }\OtherTok{\textless{}{-}}\NormalTok{ d}
\NormalTok{    out}\SpecialCharTok{$}\NormalTok{ntrials }\OtherTok{\textless{}{-}}\NormalTok{ t}
\NormalTok{    out}\SpecialCharTok{$}\NormalTok{cat\_one }\OtherTok{\textless{}{-}}\NormalTok{ d}\SpecialCharTok{$}\NormalTok{cat\_one[}\DecValTok{1}\SpecialCharTok{:}\NormalTok{t]}
\NormalTok{    out}\SpecialCharTok{$}\NormalTok{y       }\OtherTok{\textless{}{-}}\NormalTok{ d}\SpecialCharTok{$}\NormalTok{y[}\DecValTok{1}\SpecialCharTok{:}\NormalTok{t]}
\NormalTok{    out}\SpecialCharTok{$}\NormalTok{obs     }\OtherTok{\textless{}{-}}\NormalTok{ d}\SpecialCharTok{$}\NormalTok{obs[}\DecValTok{1}\SpecialCharTok{:}\NormalTok{t, , drop }\OtherTok{=} \ConstantTok{FALSE}\NormalTok{]}
\NormalTok{    out}
\NormalTok{  \}}

  \CommentTok{\# Helper: fit on y\_\{1:t\}}
\NormalTok{  refit }\OtherTok{\textless{}{-}} \ControlFlowTok{function}\NormalTok{(t) \{}
\NormalTok{    d\_t }\OtherTok{\textless{}{-}} \FunctionTok{truncate\_data}\NormalTok{(stan\_data, t)}
\NormalTok{    pf  }\OtherTok{\textless{}{-}} \FunctionTok{tryCatch}\NormalTok{(}
\NormalTok{      stan\_model}\SpecialCharTok{$}\FunctionTok{pathfinder}\NormalTok{(}\AttributeTok{data =}\NormalTok{ d\_t, }\AttributeTok{num\_paths =} \DecValTok{4}\NormalTok{, }\AttributeTok{refresh =} \DecValTok{0}\NormalTok{,}
                            \AttributeTok{seed =}\NormalTok{ seed }\SpecialCharTok{+}\NormalTok{ t),}
      \AttributeTok{error =} \ControlFlowTok{function}\NormalTok{(e) }\ConstantTok{NULL}
\NormalTok{    )}
\NormalTok{    init\_val }\OtherTok{\textless{}{-}} \ControlFlowTok{if}\NormalTok{ (}\SpecialCharTok{!}\FunctionTok{is.null}\NormalTok{(pf)) pf }\ControlFlowTok{else} \FloatTok{0.5}
\NormalTok{    stan\_model}\SpecialCharTok{$}\FunctionTok{sample}\NormalTok{(}
      \AttributeTok{data            =}\NormalTok{ d\_t,}
      \AttributeTok{init            =}\NormalTok{ init\_val,}
      \AttributeTok{seed            =}\NormalTok{ seed }\SpecialCharTok{+}\NormalTok{ t,}
      \AttributeTok{chains          =}\NormalTok{ chains,}
      \AttributeTok{parallel\_chains =}\NormalTok{ chains,}
      \AttributeTok{iter\_warmup     =}\NormalTok{ iter\_warmup,}
      \AttributeTok{iter\_sampling   =}\NormalTok{ iter\_sample,}
      \AttributeTok{refresh         =} \DecValTok{0}\NormalTok{,}
      \AttributeTok{adapt\_delta     =} \FloatTok{0.9}
\NormalTok{    )}
\NormalTok{  \}}

  \CommentTok{\# Helper: given a fitted model, evaluate log\_lik at trial t}
  \CommentTok{\# using cmdstanr\textquotesingle{}s $generate\_quantities() so we don\textquotesingle{}t re{-}MCMC.}
  \CommentTok{\# Helper: given a fitted model, evaluate log\_lik at trial t}
\NormalTok{  eval\_log\_lik\_at }\OtherTok{\textless{}{-}} \ControlFlowTok{function}\NormalTok{(fit\_ref, t) \{}
    \CommentTok{\# Build the "extended" data that includes trial t in the likelihood}
\NormalTok{    d\_ext }\OtherTok{\textless{}{-}} \FunctionTok{truncate\_data}\NormalTok{(stan\_data, t)}
    
    \CommentTok{\# FIX: Call generate\_quantities on stan\_model, not fit\_ref}
\NormalTok{    gq }\OtherTok{\textless{}{-}}\NormalTok{ stan\_model}\SpecialCharTok{$}\FunctionTok{generate\_quantities}\NormalTok{(}
      \AttributeTok{fitted\_params =}\NormalTok{ fit\_ref,}
      \AttributeTok{data          =}\NormalTok{ d\_ext,}
      \AttributeTok{seed          =}\NormalTok{ seed }\SpecialCharTok{+}\NormalTok{ t,}
      \AttributeTok{parallel\_chains =}\NormalTok{ chains}
\NormalTok{    )}
    
    \CommentTok{\# log\_lik[t] for each posterior draw}
\NormalTok{    ll\_mat }\OtherTok{\textless{}{-}}\NormalTok{ gq}\SpecialCharTok{$}\FunctionTok{draws}\NormalTok{(}\StringTok{"log\_lik"}\NormalTok{, }\AttributeTok{format =} \StringTok{"matrix"}\NormalTok{)}
\NormalTok{    ll\_mat[, t]}
\NormalTok{  \}}

  \CommentTok{\# ── Initial fit on the minimum window ─────────────────────────────────────}
  \FunctionTok{cat}\NormalTok{(}\StringTok{"Initial LFO fit on trials 1 to"}\NormalTok{, L, }\StringTok{"...}\SpecialCharTok{\textbackslash{}n}\StringTok{"}\NormalTok{)}
\NormalTok{  fit\_ref }\OtherTok{\textless{}{-}} \FunctionTok{refit}\NormalTok{(L)}
\NormalTok{  ref\_window }\OtherTok{\textless{}{-}}\NormalTok{ L}

  \CommentTok{\# Storage}
\NormalTok{  pointwise\_elpd }\OtherTok{\textless{}{-}} \FunctionTok{rep}\NormalTok{(}\ConstantTok{NA\_real\_}\NormalTok{, T\_total)}
\NormalTok{  k\_hat\_traj     }\OtherTok{\textless{}{-}} \FunctionTok{rep}\NormalTok{(}\ConstantTok{NA\_real\_}\NormalTok{, T\_total)}
\NormalTok{  refit\_at       }\OtherTok{\textless{}{-}} \FunctionTok{integer}\NormalTok{(}\DecValTok{0}\NormalTok{)}

  \CommentTok{\# ── Sequential evaluation ─────────────────────────────────────────────────}
\NormalTok{  log\_ratios\_accum }\OtherTok{\textless{}{-}} \FunctionTok{numeric}\NormalTok{(}\DecValTok{0}\NormalTok{)  }\CommentTok{\# log p(y\_\{ref+1:t{-}1\} | theta) accumulation}
  \ControlFlowTok{for}\NormalTok{ (t }\ControlFlowTok{in}\NormalTok{ (L }\SpecialCharTok{+} \DecValTok{1}\NormalTok{)}\SpecialCharTok{:}\NormalTok{T\_total) \{}
    \CommentTok{\# Build importance ratios for theta\_ref vs theta\_\{1:t{-}1\}}
    \ControlFlowTok{if}\NormalTok{ (}\FunctionTok{length}\NormalTok{(log\_ratios\_accum) }\SpecialCharTok{==} \DecValTok{0}\NormalTok{) \{}
\NormalTok{      log\_w }\OtherTok{\textless{}{-}} \FunctionTok{rep}\NormalTok{(}\DecValTok{0}\NormalTok{, posterior}\SpecialCharTok{::}\FunctionTok{ndraws}\NormalTok{(fit\_ref}\SpecialCharTok{$}\FunctionTok{draws}\NormalTok{()))}
\NormalTok{    \} }\ControlFlowTok{else}\NormalTok{ \{}
\NormalTok{      log\_w }\OtherTok{\textless{}{-}}\NormalTok{ log\_ratios\_accum  }\CommentTok{\# raw log ratios; psis() handles normalization}
\NormalTok{    \}}

    \CommentTok{\# PSIS smoothing}
\NormalTok{    psis\_obj }\OtherTok{\textless{}{-}} \FunctionTok{tryCatch}\NormalTok{(}
\NormalTok{      loo}\SpecialCharTok{::}\FunctionTok{psis}\NormalTok{(}\AttributeTok{log\_ratios =} \FunctionTok{matrix}\NormalTok{(log\_w, }\AttributeTok{ncol =} \DecValTok{1}\NormalTok{), }\AttributeTok{r\_eff =} \DecValTok{1}\NormalTok{),}
      \AttributeTok{error =} \ControlFlowTok{function}\NormalTok{(e) }\ConstantTok{NULL}
\NormalTok{    )}
\NormalTok{    k\_hat }\OtherTok{\textless{}{-}} \ControlFlowTok{if}\NormalTok{ (}\SpecialCharTok{!}\FunctionTok{is.null}\NormalTok{(psis\_obj)) loo}\SpecialCharTok{::}\FunctionTok{pareto\_k\_values}\NormalTok{(psis\_obj) }\ControlFlowTok{else} \ConstantTok{Inf}
\NormalTok{    k\_hat\_traj[t] }\OtherTok{\textless{}{-}}\NormalTok{ k\_hat}

    \ControlFlowTok{if}\NormalTok{ (}\SpecialCharTok{!}\FunctionTok{is.finite}\NormalTok{(k\_hat) }\SpecialCharTok{||}\NormalTok{ k\_hat }\SpecialCharTok{\textgreater{}}\NormalTok{ k\_thresh) \{}
      \CommentTok{\# Refit on y\_\{1:t{-}1\}}
      \FunctionTok{cat}\NormalTok{(}\StringTok{"  trial"}\NormalTok{, t, }\StringTok{": k\_hat ="}\NormalTok{, }\FunctionTok{round}\NormalTok{(k\_hat, }\DecValTok{2}\NormalTok{), }\StringTok{"→ refitting}\SpecialCharTok{\textbackslash{}n}\StringTok{"}\NormalTok{)}
\NormalTok{      fit\_ref }\OtherTok{\textless{}{-}} \FunctionTok{refit}\NormalTok{(t }\SpecialCharTok{{-}} \DecValTok{1}\NormalTok{)}
\NormalTok{      ref\_window }\OtherTok{\textless{}{-}}\NormalTok{ t }\SpecialCharTok{{-}} \DecValTok{1}
\NormalTok{      refit\_at }\OtherTok{\textless{}{-}} \FunctionTok{c}\NormalTok{(refit\_at, t)}
\NormalTok{      log\_ratios\_accum }\OtherTok{\textless{}{-}} \FunctionTok{numeric}\NormalTok{(}\DecValTok{0}\NormalTok{)}
\NormalTok{      log\_w }\OtherTok{\textless{}{-}} \FunctionTok{rep}\NormalTok{(}\DecValTok{0}\NormalTok{, posterior}\SpecialCharTok{::}\FunctionTok{ndraws}\NormalTok{(fit\_ref}\SpecialCharTok{$}\FunctionTok{draws}\NormalTok{()))}
\NormalTok{      psis\_obj }\OtherTok{\textless{}{-}}\NormalTok{ loo}\SpecialCharTok{::}\FunctionTok{psis}\NormalTok{(}\AttributeTok{log\_ratios =} \FunctionTok{matrix}\NormalTok{(log\_w, }\AttributeTok{ncol =} \DecValTok{1}\NormalTok{), }\AttributeTok{r\_eff =} \DecValTok{1}\NormalTok{)}
\NormalTok{    \}}

    \CommentTok{\# Evaluate log p(y\_t | y\_\{1:t{-}1\}, theta) for each draw}
\NormalTok{    ll\_t }\OtherTok{\textless{}{-}} \FunctionTok{eval\_log\_lik\_at}\NormalTok{(fit\_ref, t)}

    \CommentTok{\# Importance{-}weighted estimate of log p(y\_t | y\_\{1:t{-}1\})}
\NormalTok{    log\_w\_norm }\OtherTok{\textless{}{-}} \FunctionTok{as.vector}\NormalTok{(}\FunctionTok{weights}\NormalTok{(psis\_obj, }\AttributeTok{normalize =} \ConstantTok{TRUE}\NormalTok{, }\AttributeTok{log =} \ConstantTok{TRUE}\NormalTok{))}
\NormalTok{    pointwise\_elpd[t] }\OtherTok{\textless{}{-}}\NormalTok{ matrixStats}\SpecialCharTok{::}\FunctionTok{logSumExp}\NormalTok{(log\_w\_norm }\SpecialCharTok{+}\NormalTok{ ll\_t)}

    \CommentTok{\# Accumulate log p(y\_t | theta) for the next iteration\textquotesingle{}s importance ratio}
    \ControlFlowTok{if}\NormalTok{ (}\FunctionTok{length}\NormalTok{(log\_ratios\_accum) }\SpecialCharTok{==} \DecValTok{0}\NormalTok{) \{}
\NormalTok{      log\_ratios\_accum }\OtherTok{\textless{}{-}}\NormalTok{ ll\_t}
\NormalTok{    \} }\ControlFlowTok{else}\NormalTok{ \{}
\NormalTok{      log\_ratios\_accum }\OtherTok{\textless{}{-}}\NormalTok{ log\_ratios\_accum }\SpecialCharTok{+}\NormalTok{ ll\_t}
\NormalTok{    \}}
\NormalTok{  \}}

  \FunctionTok{list}\NormalTok{(}
    \AttributeTok{pointwise\_elpd =}\NormalTok{ pointwise\_elpd,}
    \AttributeTok{refit\_at       =}\NormalTok{ refit\_at,}
    \AttributeTok{k\_hat\_traj     =}\NormalTok{ k\_hat\_traj,}
    \AttributeTok{L              =}\NormalTok{ L,}
    \AttributeTok{elpd\_lfo       =} \FunctionTok{sum}\NormalTok{(pointwise\_elpd, }\AttributeTok{na.rm =} \ConstantTok{TRUE}\NormalTok{),}
    \AttributeTok{n\_refits       =} \FunctionTok{length}\NormalTok{(refit\_at)}
\NormalTok{  )}
\NormalTok{\}}
\end{Highlighting}
\end{Shaded}

\begin{quote}
\textbf{Note on the importance-ratio accumulation.} The trick that makes
PSIS-LFO efficient is that the importance ratio for
\(p(\theta \mid y_{1:t-1})\) vs the reference \(p(\theta \mid y_{1:L})\)
is a \emph{running product}: for each new trial \(t\) we just multiply
the existing per-draw weight by \(p(y_{t-1} \mid \theta)\). We never
need to refit so long as the smoothed Pareto-\(\hat{k}\) stays below the
threshold. When it crosses 0.7, we refit and reset the running product.
The accumulation in \texttt{log\_ratios\_accum} above implements this
--- every time we evaluate \texttt{log\_lik{[}t{]}}, we add it to the
cumulative log-ratio for the next iteration.
\end{quote}

We now run this on the same single-agent fit, and validate it against
\emph{exact} one-step-ahead refits at every 8th trial. If PSIS-LFO is
doing its job, the two ELPD curves should agree closely.

\begin{Shaded}
\begin{Highlighting}[]
\NormalTok{lfo\_cache\_file }\OtherTok{\textless{}{-}} \FunctionTok{here}\NormalTok{(}\StringTok{"simdata"}\NormalTok{, }\StringTok{"ch13\_lfo\_kalman\_results.rds"}\NormalTok{)}
\ControlFlowTok{if}\NormalTok{ (regenerate\_simulations }\SpecialCharTok{||} \SpecialCharTok{!}\FunctionTok{file.exists}\NormalTok{(lfo\_cache\_file)) \{}
  \FunctionTok{cat}\NormalTok{(}\StringTok{"Running PSIS{-}LFO for the prototype model...}\SpecialCharTok{\textbackslash{}n}\StringTok{"}\NormalTok{)}
\NormalTok{  lfo\_psis }\OtherTok{\textless{}{-}} \FunctionTok{psis\_lfo\_kalman}\NormalTok{(}
    \AttributeTok{stan\_data  =}\NormalTok{ proto\_data\_single,}
    \AttributeTok{stan\_model =}\NormalTok{ mod\_prototype\_single,}
    \AttributeTok{L          =} \DecValTok{16}\NormalTok{,}
    \AttributeTok{k\_thresh   =} \FloatTok{0.7}
\NormalTok{  )}

  \CommentTok{\# ── Exact one{-}step{-}ahead refits at every 8th trial ────────────────────────}
\NormalTok{  exact\_t }\OtherTok{\textless{}{-}} \FunctionTok{seq}\NormalTok{(}\DecValTok{24}\NormalTok{, total\_trials, }\AttributeTok{by =} \DecValTok{8}\NormalTok{)}
\NormalTok{  exact\_elpd }\OtherTok{\textless{}{-}} \FunctionTok{numeric}\NormalTok{(}\FunctionTok{length}\NormalTok{(exact\_t))}
  \ControlFlowTok{for}\NormalTok{ (k }\ControlFlowTok{in} \FunctionTok{seq\_along}\NormalTok{(exact\_t)) \{}
\NormalTok{    t }\OtherTok{\textless{}{-}}\NormalTok{ exact\_t[k]}
    \FunctionTok{cat}\NormalTok{(}\StringTok{"  exact refit at t ="}\NormalTok{, t, }\StringTok{"}\SpecialCharTok{\textbackslash{}n}\StringTok{"}\NormalTok{)}
\NormalTok{    d\_t }\OtherTok{\textless{}{-}}\NormalTok{ proto\_data\_single}
\NormalTok{    d\_t}\SpecialCharTok{$}\NormalTok{ntrials }\OtherTok{\textless{}{-}}\NormalTok{ t }\SpecialCharTok{{-}} \DecValTok{1}
\NormalTok{    d\_t}\SpecialCharTok{$}\NormalTok{cat\_one }\OtherTok{\textless{}{-}}\NormalTok{ proto\_data\_single}\SpecialCharTok{$}\NormalTok{cat\_one[}\DecValTok{1}\SpecialCharTok{:}\NormalTok{(t }\SpecialCharTok{{-}} \DecValTok{1}\NormalTok{)]}
\NormalTok{    d\_t}\SpecialCharTok{$}\NormalTok{y       }\OtherTok{\textless{}{-}}\NormalTok{ proto\_data\_single}\SpecialCharTok{$}\NormalTok{y[}\DecValTok{1}\SpecialCharTok{:}\NormalTok{(t }\SpecialCharTok{{-}} \DecValTok{1}\NormalTok{)]}
\NormalTok{    d\_t}\SpecialCharTok{$}\NormalTok{obs     }\OtherTok{\textless{}{-}}\NormalTok{ proto\_data\_single}\SpecialCharTok{$}\NormalTok{obs[}\DecValTok{1}\SpecialCharTok{:}\NormalTok{(t }\SpecialCharTok{{-}} \DecValTok{1}\NormalTok{), , drop }\OtherTok{=} \ConstantTok{FALSE}\NormalTok{]}

\NormalTok{    pf  }\OtherTok{\textless{}{-}} \FunctionTok{tryCatch}\NormalTok{(}
\NormalTok{      mod\_prototype\_single}\SpecialCharTok{$}\FunctionTok{pathfinder}\NormalTok{(}\AttributeTok{data =}\NormalTok{ d\_t, }\AttributeTok{num\_paths =} \DecValTok{4}\NormalTok{, }\AttributeTok{refresh =} \DecValTok{0}\NormalTok{),}
      \AttributeTok{error =} \ControlFlowTok{function}\NormalTok{(e) }\ConstantTok{NULL}
\NormalTok{    )}
\NormalTok{    init\_val }\OtherTok{\textless{}{-}} \ControlFlowTok{if}\NormalTok{ (}\SpecialCharTok{!}\FunctionTok{is.null}\NormalTok{(pf)) pf }\ControlFlowTok{else} \FloatTok{0.5}
\NormalTok{    fit\_t }\OtherTok{\textless{}{-}}\NormalTok{ mod\_prototype\_single}\SpecialCharTok{$}\FunctionTok{sample}\NormalTok{(}
      \AttributeTok{data =}\NormalTok{ d\_t, }\AttributeTok{init =}\NormalTok{ init\_val, }\AttributeTok{chains =} \DecValTok{2}\NormalTok{, }\AttributeTok{parallel\_chains =} \DecValTok{2}\NormalTok{,}
      \AttributeTok{iter\_warmup =} \DecValTok{800}\NormalTok{, }\AttributeTok{iter\_sampling =} \DecValTok{1000}\NormalTok{, }\AttributeTok{refresh =} \DecValTok{0}\NormalTok{, }\AttributeTok{adapt\_delta =} \FloatTok{0.9}
\NormalTok{    )}

    \CommentTok{\# Evaluate log\_lik at trial t}
\NormalTok{    d\_ext }\OtherTok{\textless{}{-}}\NormalTok{ proto\_data\_single}
\NormalTok{    d\_ext}\SpecialCharTok{$}\NormalTok{ntrials }\OtherTok{\textless{}{-}}\NormalTok{ t}
\NormalTok{    d\_ext}\SpecialCharTok{$}\NormalTok{cat\_one }\OtherTok{\textless{}{-}}\NormalTok{ proto\_data\_single}\SpecialCharTok{$}\NormalTok{cat\_one[}\DecValTok{1}\SpecialCharTok{:}\NormalTok{t]}
\NormalTok{    d\_ext}\SpecialCharTok{$}\NormalTok{y       }\OtherTok{\textless{}{-}}\NormalTok{ proto\_data\_single}\SpecialCharTok{$}\NormalTok{y[}\DecValTok{1}\SpecialCharTok{:}\NormalTok{t]}
\NormalTok{    d\_ext}\SpecialCharTok{$}\NormalTok{obs     }\OtherTok{\textless{}{-}}\NormalTok{ proto\_data\_single}\SpecialCharTok{$}\NormalTok{obs[}\DecValTok{1}\SpecialCharTok{:}\NormalTok{t, , drop }\OtherTok{=} \ConstantTok{FALSE}\NormalTok{]}
    
    \CommentTok{\# ── THE FIX: Called on mod\_prototype\_single instead of fit\_t ──}
\NormalTok{    gq }\OtherTok{\textless{}{-}}\NormalTok{ mod\_prototype\_single}\SpecialCharTok{$}\FunctionTok{generate\_quantities}\NormalTok{(}\AttributeTok{fitted\_params =}\NormalTok{ fit\_t, }\AttributeTok{data =}\NormalTok{ d\_ext)}
    
\NormalTok{    ll\_mat }\OtherTok{\textless{}{-}}\NormalTok{ gq}\SpecialCharTok{$}\FunctionTok{draws}\NormalTok{(}\StringTok{"log\_lik"}\NormalTok{, }\AttributeTok{format =} \StringTok{"matrix"}\NormalTok{)}
\NormalTok{    exact\_elpd[k] }\OtherTok{\textless{}{-}}\NormalTok{ matrixStats}\SpecialCharTok{::}\FunctionTok{logSumExp}\NormalTok{(ll\_mat[, t]) }\SpecialCharTok{{-}} \FunctionTok{log}\NormalTok{(}\FunctionTok{nrow}\NormalTok{(ll\_mat))}
\NormalTok{  \}}

  \FunctionTok{saveRDS}\NormalTok{(}\FunctionTok{list}\NormalTok{(}\AttributeTok{lfo\_psis =}\NormalTok{ lfo\_psis,}
               \AttributeTok{exact\_t =}\NormalTok{ exact\_t, }\AttributeTok{exact\_elpd =}\NormalTok{ exact\_elpd),}
\NormalTok{          lfo\_cache\_file)}
  \FunctionTok{cat}\NormalTok{(}\StringTok{"LFO results computed and saved.}\SpecialCharTok{\textbackslash{}n}\StringTok{"}\NormalTok{)}
\NormalTok{\} }\ControlFlowTok{else}\NormalTok{ \{}
\NormalTok{  cached }\OtherTok{\textless{}{-}} \FunctionTok{readRDS}\NormalTok{(lfo\_cache\_file)}
\NormalTok{  lfo\_psis }\OtherTok{\textless{}{-}}\NormalTok{ cached}\SpecialCharTok{$}\NormalTok{lfo\_psis}
\NormalTok{  exact\_t }\OtherTok{\textless{}{-}}\NormalTok{ cached}\SpecialCharTok{$}\NormalTok{exact\_t}
\NormalTok{  exact\_elpd }\OtherTok{\textless{}{-}}\NormalTok{ cached}\SpecialCharTok{$}\NormalTok{exact\_elpd}
  \FunctionTok{cat}\NormalTok{(}\StringTok{"Loaded existing LFO results.}\SpecialCharTok{\textbackslash{}n}\StringTok{"}\NormalTok{)}
\NormalTok{\}}
\end{Highlighting}
\end{Shaded}

\begin{verbatim}
Loaded existing LFO results.
\end{verbatim}

\begin{Shaded}
\begin{Highlighting}[]
\ControlFlowTok{if}\NormalTok{ (}\FunctionTok{exists}\NormalTok{(}\StringTok{"lfo\_psis"}\NormalTok{)) \{}
  \FunctionTok{cat}\NormalTok{(}\StringTok{"PSIS{-}LFO sum ELPD:"}\NormalTok{, }\FunctionTok{round}\NormalTok{(lfo\_psis}\SpecialCharTok{$}\NormalTok{elpd\_lfo, }\DecValTok{2}\NormalTok{), }\StringTok{"}\SpecialCharTok{\textbackslash{}n}\StringTok{"}\NormalTok{)}
  \FunctionTok{cat}\NormalTok{(}\StringTok{"Number of refits triggered:"}\NormalTok{, lfo\_psis}\SpecialCharTok{$}\NormalTok{n\_refits,}
      \StringTok{"out of"}\NormalTok{, total\_trials }\SpecialCharTok{{-}}\NormalTok{ lfo\_psis}\SpecialCharTok{$}\NormalTok{L, }\StringTok{"trials}\SpecialCharTok{\textbackslash{}n}\StringTok{"}\NormalTok{)}

  \CommentTok{\# Plot: PSIS{-}LFO pointwise vs exact refits}
\NormalTok{  comp\_df }\OtherTok{\textless{}{-}} \FunctionTok{tibble}\NormalTok{(}
    \AttributeTok{trial    =} \FunctionTok{seq\_along}\NormalTok{(lfo\_psis}\SpecialCharTok{$}\NormalTok{pointwise\_elpd),}
    \AttributeTok{psis\_elpd =}\NormalTok{ lfo\_psis}\SpecialCharTok{$}\NormalTok{pointwise\_elpd}
\NormalTok{  )}
\NormalTok{  exact\_df }\OtherTok{\textless{}{-}} \FunctionTok{tibble}\NormalTok{(}\AttributeTok{trial =}\NormalTok{ exact\_t, }\AttributeTok{exact\_elpd =}\NormalTok{ exact\_elpd)}

\NormalTok{  p\_lfo }\OtherTok{\textless{}{-}} \FunctionTok{ggplot}\NormalTok{(comp\_df, }\FunctionTok{aes}\NormalTok{(}\AttributeTok{x =}\NormalTok{ trial, }\AttributeTok{y =}\NormalTok{ psis\_elpd)) }\SpecialCharTok{+}
    \FunctionTok{geom\_line}\NormalTok{(}\AttributeTok{color =} \StringTok{"\#0072B2"}\NormalTok{) }\SpecialCharTok{+}
    \FunctionTok{geom\_point}\NormalTok{(}\AttributeTok{data =}\NormalTok{ exact\_df,}
               \FunctionTok{aes}\NormalTok{(}\AttributeTok{x =}\NormalTok{ trial, }\AttributeTok{y =}\NormalTok{ exact\_elpd),}
               \AttributeTok{color =} \StringTok{"\#D55E00"}\NormalTok{, }\AttributeTok{size =} \DecValTok{2}\NormalTok{) }\SpecialCharTok{+}
    \FunctionTok{labs}\NormalTok{(}
      \AttributeTok{title    =} \StringTok{"PSIS{-}LFO vs. Exact One{-}Step{-}Ahead Refits"}\NormalTok{,}
      \AttributeTok{subtitle =} \StringTok{"Blue line: PSIS{-}LFO pointwise ELPD; orange points: exact refits at every 8th trial"}\NormalTok{,}
      \AttributeTok{x =} \StringTok{"Trial"}\NormalTok{, }\AttributeTok{y =} \StringTok{"log p(y\_t | y\_\{1:t{-}1\})"}
\NormalTok{    )}

  \CommentTok{\# Plot: Pareto{-}k trajectory and refit triggers}
\NormalTok{  k\_df }\OtherTok{\textless{}{-}} \FunctionTok{tibble}\NormalTok{(}\AttributeTok{trial =} \FunctionTok{seq\_along}\NormalTok{(lfo\_psis}\SpecialCharTok{$}\NormalTok{k\_hat\_traj),}
                 \AttributeTok{k\_hat =}\NormalTok{ lfo\_psis}\SpecialCharTok{$}\NormalTok{k\_hat\_traj)}
\NormalTok{  refit\_df }\OtherTok{\textless{}{-}} \FunctionTok{tibble}\NormalTok{(}\AttributeTok{trial =}\NormalTok{ lfo\_psis}\SpecialCharTok{$}\NormalTok{refit\_at)}

\NormalTok{  p\_k }\OtherTok{\textless{}{-}} \FunctionTok{ggplot}\NormalTok{(k\_df, }\FunctionTok{aes}\NormalTok{(}\AttributeTok{x =}\NormalTok{ trial, }\AttributeTok{y =}\NormalTok{ k\_hat)) }\SpecialCharTok{+}
    \FunctionTok{geom\_line}\NormalTok{(}\AttributeTok{color =} \StringTok{"\#009E73"}\NormalTok{) }\SpecialCharTok{+}
    \FunctionTok{geom\_hline}\NormalTok{(}\AttributeTok{yintercept =} \FloatTok{0.7}\NormalTok{, }\AttributeTok{linetype =} \StringTok{"dashed"}\NormalTok{, }\AttributeTok{color =} \StringTok{"red"}\NormalTok{) }\SpecialCharTok{+}
    \FunctionTok{geom\_vline}\NormalTok{(}\AttributeTok{data =}\NormalTok{ refit\_df, }\FunctionTok{aes}\NormalTok{(}\AttributeTok{xintercept =}\NormalTok{ trial),}
               \AttributeTok{linetype =} \StringTok{"dotted"}\NormalTok{, }\AttributeTok{color =} \StringTok{"red"}\NormalTok{, }\AttributeTok{alpha =} \FloatTok{0.5}\NormalTok{) }\SpecialCharTok{+}
    \FunctionTok{labs}\NormalTok{(}
      \AttributeTok{title    =} \StringTok{"Pareto{-}k trajectory and refit triggers"}\NormalTok{,}
      \AttributeTok{subtitle =} \StringTok{"Dashed red: k = 0.7 threshold; dotted vertical lines = refits"}\NormalTok{,}
      \AttributeTok{x =} \StringTok{"Trial"}\NormalTok{, }\AttributeTok{y =} \StringTok{"Pareto k"}
\NormalTok{    )}

  \FunctionTok{print}\NormalTok{(p\_lfo }\SpecialCharTok{/}\NormalTok{ p\_k)}
\NormalTok{\}}
\end{Highlighting}
\end{Shaded}

\begin{verbatim}
PSIS-LFO sum ELPD: -32.25 
Number of refits triggered: 1 out of 48 trials
\end{verbatim}

\begin{center}
\includegraphics[width=0.8\linewidth,height=\textheight,keepaspectratio]{12-CategorizationModelsPrototypes_files/figure-pdf/ch13_lfo_run-1.pdf}
\end{center}

\textbf{Reading the output.} The blue line is the PSIS-LFO pointwise
ELPD; the orange points are exact one-step-ahead refits. They should sit
on top of each other (modulo Monte Carlo noise from the small chain
count). The Pareto-\(\hat{k}\) trajectory should mostly stay below 0.7,
with refits triggered only occasionally. If the two curves diverge
systematically, PSIS-LFO is failing for this model and the exact refit
pipeline (slow but unambiguous) is the fallback. \textbf{This is the
LFO-CV implementation that Chapter 15's model comparison will use.}
Chapter 13's marginalised mixture is conditionally exchangeable and does
not need it; Chapter 11's GCM has its own canonical implementation; this
chapter inherits the algorithm and applies it to the Kalman case. The
repeated LFO-CV caveat from earlier drafts is now discharged.

\subsection{Prior Sensitivity
Analysis}\label{prior-sensitivity-analysis-2}

\begin{Shaded}
\begin{Highlighting}[]
\ControlFlowTok{if}\NormalTok{ (}\SpecialCharTok{!}\FunctionTok{is.null}\NormalTok{(fit\_proto\_single)) \{}
\NormalTok{  ps }\OtherTok{\textless{}{-}}\NormalTok{ priorsense}\SpecialCharTok{::}\FunctionTok{powerscale\_sensitivity}\NormalTok{(fit\_proto\_single)}
  \FunctionTok{print}\NormalTok{(ps)}
\NormalTok{  priorsense}\SpecialCharTok{::}\FunctionTok{powerscale\_plot\_dens}\NormalTok{(fit\_proto\_single, }\AttributeTok{variables =} \FunctionTok{c}\NormalTok{(}\StringTok{"log\_r"}\NormalTok{, }\StringTok{"r\_value"}\NormalTok{, }\StringTok{"log\_q"}\NormalTok{, }\StringTok{"q\_value"}\NormalTok{))}
\NormalTok{\}}
\end{Highlighting}
\end{Shaded}

\begin{verbatim}
Sensitivity based on cjs_dist
Prior selection: all priors
Likelihood selection: all data

 variable prior likelihood                                diagnosis
    log_r 0.079      0.096            potential prior-data conflict
    log_q 0.081      0.137            potential prior-data conflict
  r_value 0.066      0.197            potential prior-data conflict
  q_value 0.208      0.126            potential prior-data conflict
     p[1]   NaN        NaN                                     <NA>
     p[2] 0.052      0.103            potential prior-data conflict
     p[3] 0.077      0.128            potential prior-data conflict
     p[4] 0.090      0.134            potential prior-data conflict
     p[5] 0.065      0.071            potential prior-data conflict
     p[6] 0.116      0.079            potential prior-data conflict
     p[7] 0.029      0.123                                        -
     p[8] 0.028      0.129                                        -
     p[9] 0.038      0.121                                        -
    p[10] 0.034      0.115                                        -
    p[11] 0.059      0.118            potential prior-data conflict
    p[12] 0.059      0.085            potential prior-data conflict
    p[13] 0.034      0.144                                        -
    p[14] 0.035      0.100                                        -
    p[15] 0.083      0.104            potential prior-data conflict
    p[16] 0.073      0.030 potential strong prior / weak likelihood
    p[17] 0.032      0.110                                        -
    p[18] 0.047      0.117                                        -
    p[19] 0.061      0.096            potential prior-data conflict
    p[20] 0.071      0.058            potential prior-data conflict
    p[21] 0.062      0.095            potential prior-data conflict
    p[22] 0.090      0.091            potential prior-data conflict
    p[23] 0.083      0.054            potential prior-data conflict
    p[24] 0.068      0.135            potential prior-data conflict
    p[25] 0.038      0.097                                        -
    p[26] 0.065      0.062            potential prior-data conflict
    p[27] 0.089      0.084            potential prior-data conflict
    p[28] 0.029      0.140                                        -
    p[29] 0.066      0.066            potential prior-data conflict
    p[30] 0.065      0.084            potential prior-data conflict
    p[31] 0.052      0.106            potential prior-data conflict
    p[32] 0.096      0.120            potential prior-data conflict
    p[33] 0.038      0.086                                        -
    p[34] 0.055      0.087            potential prior-data conflict
    p[35] 0.059      0.065            potential prior-data conflict
    p[36] 0.130      0.090            potential prior-data conflict
    p[37] 0.065      0.103            potential prior-data conflict
    p[38] 0.098      0.076            potential prior-data conflict
    p[39] 0.040      0.118                                        -
    p[40] 0.090      0.130            potential prior-data conflict
    p[41] 0.110      0.051            potential prior-data conflict
    p[42] 0.029      0.141                                        -
    p[43] 0.025      0.106                                        -
    p[44] 0.057      0.075            potential prior-data conflict
    p[45] 0.037      0.112                                        -
    p[46] 0.037      0.124                                        -
    p[47] 0.094      0.097            potential prior-data conflict
    p[48] 0.184      0.034 potential strong prior / weak likelihood
    p[49] 0.025      0.115                                        -
    p[50] 0.068      0.074            potential prior-data conflict
    p[51] 0.085      0.122            potential prior-data conflict
    p[52] 0.032      0.129                                        -
    p[53] 0.030      0.091                                        -
    p[54] 0.082      0.092            potential prior-data conflict
    p[55] 0.063      0.035 potential strong prior / weak likelihood
    p[56] 0.102      0.085            potential prior-data conflict
    p[57] 0.040      0.105                                        -
    p[58] 0.026      0.109                                        -
    p[59] 0.076      0.084            potential prior-data conflict
    p[60] 0.033      0.139                                        -
    p[61] 0.036      0.046                                        -
    p[62] 0.133      0.084            potential prior-data conflict
    p[63] 0.067      0.110            potential prior-data conflict
    p[64] 0.090      0.098            potential prior-data conflict
\end{verbatim}

\begin{center}
\includegraphics[width=0.8\linewidth,height=\textheight,keepaspectratio]{12-CategorizationModelsPrototypes_files/figure-pdf/ch13_prior_sensitivity-1.pdf}
\end{center}

\texttt{powerscale\_sensitivity()} perturbs the prior analytically
(without re-running MCMC) and reports a sensitivity index \(\psi\) for
each parameter. A value \(\psi > 0.1\) indicates structural fragility
--- the conclusion depends materially on the specific prior choice. For
a model with two parameters and moderately informative priors on
\(\log r\) and \(\log q\), \(\psi\) should be small but non-zero. Under
the static Kruschke paradigm, \(\log q\) may show higher sensitivity
than \(\log r\) because the data are less informative about it. The
accompanying density plot overlays posteriors at three power-scaling
levels and lets the reader see whether the shape of each posterior is
robust or whether a slightly different prior would push it elsewhere.

\begin{center}\rule{0.5\linewidth}{0.5pt}\end{center}

\section{Simulation-Based Calibration Checks
(SBC)}\label{simulation-based-calibration-checks-sbc-3}

Parameter recovery at specific true values is a necessary but not
sufficient validation. SBC checks that the full posterior is correctly
calibrated across the entire prior distribution --- that if the true
\texttt{log\_r} and \texttt{log\_q} were drawn from their priors, the
posterior would on average contain each true value in the appropriate
fraction of its probability mass. We use the \texttt{SBC} package for
consistency with Chapters 6, 8, and 11.

This is \textbf{proper SBC in the Talts--Vehtari sense}: data are drawn
from the \emph{same} generative model that is then fitted. (This is in
contrast to Chapter 13's ``cross-model calibration'' check, where data
come from a particle filter and are fitted with a different mixture
model --- that procedure tests robustness, not calibration.) The
prototype model's two scalar parameters make the resulting
ECDF-difference plot particularly legible: the \texttt{log\_r} curve
should show tight calibration (the data are informative), while the
\texttt{log\_q} curve may be wider under the static paradigm (partial
confound with \texttt{log\_r}).

\begin{Shaded}
\begin{Highlighting}[]
\NormalTok{sbc\_proto\_filepath }\OtherTok{\textless{}{-}} \FunctionTok{here}\NormalTok{(}\StringTok{"simdata"}\NormalTok{, }\StringTok{"ch13\_sbc\_proto\_results.rds"}\NormalTok{)}
\ControlFlowTok{if}\NormalTok{ (regenerate\_simulations }\SpecialCharTok{||} \SpecialCharTok{!}\FunctionTok{file.exists}\NormalTok{(sbc\_proto\_filepath)) \{}
\NormalTok{  n\_sbc\_iterations\_proto }\OtherTok{\textless{}{-}} \DecValTok{500}   \CommentTok{\# use \textgreater{}= 1000 for publication{-}quality SBC}

\NormalTok{  proto\_sbc\_generator }\OtherTok{\textless{}{-}} \FunctionTok{SBC\_generator\_function}\NormalTok{(}
    \ControlFlowTok{function}\NormalTok{() \{}
      \CommentTok{\# 1a. Draw true parameters from their priors}
\NormalTok{      log\_r\_true }\OtherTok{\textless{}{-}} \FunctionTok{rnorm}\NormalTok{(}\DecValTok{1}\NormalTok{, LOG\_R\_PRIOR\_MEAN, LOG\_R\_PRIOR\_SD)}
\NormalTok{      r\_true     }\OtherTok{\textless{}{-}} \FunctionTok{exp}\NormalTok{(log\_r\_true)}
\NormalTok{      log\_q\_true }\OtherTok{\textless{}{-}} \FunctionTok{rnorm}\NormalTok{(}\DecValTok{1}\NormalTok{, LOG\_Q\_PRIOR\_MEAN, LOG\_Q\_PRIOR\_SD)}
\NormalTok{      q\_true     }\OtherTok{\textless{}{-}} \FunctionTok{exp}\NormalTok{(log\_q\_true)}

      \CommentTok{\# 1b. Simulate one dataset with its own randomized schedule}
\NormalTok{      sbc\_sched }\OtherTok{\textless{}{-}} \FunctionTok{make\_subject\_schedule}\NormalTok{(stimulus\_info, n\_blocks,}
                                         \AttributeTok{seed =} \FunctionTok{sample.int}\NormalTok{(}\FloatTok{1e8}\NormalTok{, }\DecValTok{1}\NormalTok{))}
\NormalTok{      sim       }\OtherTok{\textless{}{-}} \FunctionTok{prototype\_kalman}\NormalTok{(}
        \AttributeTok{r\_value    =}\NormalTok{ r\_true,}
        \AttributeTok{q\_value    =}\NormalTok{ q\_true,}
        \AttributeTok{obs        =} \FunctionTok{as.matrix}\NormalTok{(sbc\_sched[, }\FunctionTok{c}\NormalTok{(}\StringTok{"height"}\NormalTok{, }\StringTok{"position"}\NormalTok{)]),}
        \AttributeTok{cat\_one    =}\NormalTok{ sbc\_sched}\SpecialCharTok{$}\NormalTok{category\_feedback,}
        \AttributeTok{initial\_mu =} \FunctionTok{list}\NormalTok{(}\FunctionTok{rep}\NormalTok{(}\FloatTok{2.5}\NormalTok{, }\DecValTok{2}\NormalTok{), }\FunctionTok{rep}\NormalTok{(}\FloatTok{2.5}\NormalTok{, }\DecValTok{2}\NormalTok{))}
\NormalTok{      )}

      \FunctionTok{list}\NormalTok{(}
        \AttributeTok{variables =} \FunctionTok{list}\NormalTok{(}\AttributeTok{log\_r =}\NormalTok{ log\_r\_true, }\AttributeTok{log\_q =}\NormalTok{ log\_q\_true),}
        \AttributeTok{generated =} \FunctionTok{list}\NormalTok{(}
          \AttributeTok{ntrials            =} \FunctionTok{nrow}\NormalTok{(sbc\_sched),}
          \AttributeTok{nfeatures          =} \DecValTok{2}\NormalTok{L,}
          \AttributeTok{cat\_one            =} \FunctionTok{as.integer}\NormalTok{(sbc\_sched}\SpecialCharTok{$}\NormalTok{category\_feedback),}
          \AttributeTok{y                  =} \FunctionTok{as.integer}\NormalTok{(sim}\SpecialCharTok{$}\NormalTok{sim\_response),}
          \AttributeTok{obs                =} \FunctionTok{as.matrix}\NormalTok{(sbc\_sched[, }\FunctionTok{c}\NormalTok{(}\StringTok{"height"}\NormalTok{, }\StringTok{"position"}\NormalTok{)]),}
          \AttributeTok{initial\_mu\_cat0    =} \FunctionTok{c}\NormalTok{(}\FloatTok{2.5}\NormalTok{, }\FloatTok{2.5}\NormalTok{),}
          \AttributeTok{initial\_mu\_cat1    =} \FunctionTok{c}\NormalTok{(}\FloatTok{2.5}\NormalTok{, }\FloatTok{2.5}\NormalTok{),}
          \AttributeTok{initial\_sigma\_diag =} \FloatTok{10.0}\NormalTok{,}
          \AttributeTok{prior\_logr\_mean    =}\NormalTok{ LOG\_R\_PRIOR\_MEAN,}
          \AttributeTok{prior\_logr\_sd      =}\NormalTok{ LOG\_R\_PRIOR\_SD,}
          \AttributeTok{prior\_logq\_mean    =}\NormalTok{ LOG\_Q\_PRIOR\_MEAN,}
          \AttributeTok{prior\_logq\_sd      =}\NormalTok{ LOG\_Q\_PRIOR\_SD}
\NormalTok{        )}
\NormalTok{      )}
\NormalTok{    \}}
\NormalTok{  )}

\NormalTok{  proto\_sbc\_backend }\OtherTok{\textless{}{-}} \FunctionTok{SBC\_backend\_cmdstan\_sample}\NormalTok{(}
\NormalTok{    mod\_prototype\_single,}
    \AttributeTok{iter\_warmup   =} \DecValTok{1000}\NormalTok{,}
    \AttributeTok{iter\_sampling =} \DecValTok{1000}\NormalTok{,}
    \AttributeTok{chains        =} \DecValTok{2}\NormalTok{,}
    \AttributeTok{adapt\_delta   =} \FloatTok{0.9}\NormalTok{,}
    \AttributeTok{refresh       =} \DecValTok{0}
\NormalTok{  )}

\NormalTok{  proto\_datasets }\OtherTok{\textless{}{-}} \FunctionTok{generate\_datasets}\NormalTok{(}
\NormalTok{    proto\_sbc\_generator,}
    \AttributeTok{n\_sims =}\NormalTok{ n\_sbc\_iterations\_proto}
\NormalTok{  )}

\NormalTok{  sbc\_results\_proto }\OtherTok{\textless{}{-}} \FunctionTok{compute\_SBC}\NormalTok{(}
\NormalTok{    proto\_datasets,}
\NormalTok{    proto\_sbc\_backend,}
    \AttributeTok{cache\_mode     =} \StringTok{"results"}\NormalTok{,}
    \AttributeTok{cache\_location =} \FunctionTok{here}\NormalTok{(}\StringTok{"simdata"}\NormalTok{, }\StringTok{"ch13\_sbc\_proto\_cache"}\NormalTok{),}
    \AttributeTok{keep\_fits      =} \ConstantTok{FALSE}
\NormalTok{  )}
  \FunctionTok{saveRDS}\NormalTok{(sbc\_results\_proto, sbc\_proto\_filepath)}
  \FunctionTok{cat}\NormalTok{(}\StringTok{"SBC prototype results computed and saved.}\SpecialCharTok{\textbackslash{}n}\StringTok{"}\NormalTok{)}
\NormalTok{\} }\ControlFlowTok{else}\NormalTok{ \{}
\NormalTok{  sbc\_results\_proto }\OtherTok{\textless{}{-}} \FunctionTok{readRDS}\NormalTok{(sbc\_proto\_filepath)}
  \FunctionTok{cat}\NormalTok{(}\StringTok{"Loaded existing SBC prototype results.}\SpecialCharTok{\textbackslash{}n}\StringTok{"}\NormalTok{)}
\NormalTok{\}}
\end{Highlighting}
\end{Shaded}

\begin{verbatim}
Loaded existing SBC prototype results.
\end{verbatim}

\begin{Shaded}
\begin{Highlighting}[]
\FunctionTok{plot\_ecdf\_diff}\NormalTok{(sbc\_results\_proto)}
\end{Highlighting}
\end{Shaded}

\begin{center}
\includegraphics[width=0.8\linewidth,height=\textheight,keepaspectratio]{12-CategorizationModelsPrototypes_files/figure-pdf/ch13_sbc_prototype-1.pdf}
\end{center}

\begin{Shaded}
\begin{Highlighting}[]
\FunctionTok{plot\_rank\_hist}\NormalTok{(sbc\_results\_proto)}
\end{Highlighting}
\end{Shaded}

\begin{center}
\includegraphics[width=0.8\linewidth,height=\textheight,keepaspectratio]{12-CategorizationModelsPrototypes_files/figure-pdf/ch13_sbc_prototype-2.pdf}
\end{center}

\begin{Shaded}
\begin{Highlighting}[]
\CommentTok{\# ── Diagnostic issue investigation ───────────────────────────────────────────}
\NormalTok{backend\_diags  }\OtherTok{\textless{}{-}}\NormalTok{ sbc\_results\_proto}\SpecialCharTok{$}\NormalTok{backend\_diagnostics}
\NormalTok{default\_diags  }\OtherTok{\textless{}{-}}\NormalTok{ sbc\_results\_proto}\SpecialCharTok{$}\NormalTok{default\_diagnostics }\SpecialCharTok{|\textgreater{}}
\NormalTok{  dplyr}\SpecialCharTok{::}\FunctionTok{select}\NormalTok{(sim\_id, max\_rhat, min\_ess\_to\_rank)}
\NormalTok{true\_params\_wide }\OtherTok{\textless{}{-}}\NormalTok{ sbc\_results\_proto}\SpecialCharTok{$}\NormalTok{stats }\SpecialCharTok{|\textgreater{}}
\NormalTok{  dplyr}\SpecialCharTok{::}\FunctionTok{select}\NormalTok{(sim\_id, variable, simulated\_value) }\SpecialCharTok{|\textgreater{}}
\NormalTok{  tidyr}\SpecialCharTok{::}\FunctionTok{pivot\_wider}\NormalTok{(}\AttributeTok{names\_from =}\NormalTok{ variable, }\AttributeTok{values\_from =}\NormalTok{ simulated\_value)}

\NormalTok{diag\_df }\OtherTok{\textless{}{-}}\NormalTok{ backend\_diags }\SpecialCharTok{|\textgreater{}}
\NormalTok{  dplyr}\SpecialCharTok{::}\FunctionTok{left\_join}\NormalTok{(default\_diags,    }\AttributeTok{by =} \StringTok{"sim\_id"}\NormalTok{) }\SpecialCharTok{|\textgreater{}}
\NormalTok{  dplyr}\SpecialCharTok{::}\FunctionTok{left\_join}\NormalTok{(true\_params\_wide, }\AttributeTok{by =} \StringTok{"sim\_id"}\NormalTok{) }\SpecialCharTok{|\textgreater{}}
\NormalTok{  dplyr}\SpecialCharTok{::}\FunctionTok{mutate}\NormalTok{(}\AttributeTok{has\_issue =}\NormalTok{ n\_divergent }\SpecialCharTok{\textgreater{}} \DecValTok{0} \SpecialCharTok{|}\NormalTok{ n\_max\_treedepth }\SpecialCharTok{\textgreater{}} \DecValTok{0} \SpecialCharTok{|}\NormalTok{ max\_rhat }\SpecialCharTok{\textgreater{}} \FloatTok{1.01}\NormalTok{)}

\NormalTok{p\_diag\_r\_q }\OtherTok{\textless{}{-}} \FunctionTok{ggplot}\NormalTok{(diag\_df, }\FunctionTok{aes}\NormalTok{(}\AttributeTok{x =}\NormalTok{ log\_r, }\AttributeTok{y =}\NormalTok{ log\_q, }\AttributeTok{color =}\NormalTok{ has\_issue)) }\SpecialCharTok{+}
  \FunctionTok{geom\_point}\NormalTok{(}\AttributeTok{alpha =} \FloatTok{0.7}\NormalTok{, }\AttributeTok{size =} \DecValTok{2}\NormalTok{) }\SpecialCharTok{+}
  \FunctionTok{scale\_color\_manual}\NormalTok{(}\AttributeTok{values =} \FunctionTok{c}\NormalTok{(}\StringTok{"FALSE"} \OtherTok{=} \StringTok{"\#0072B2"}\NormalTok{, }\StringTok{"TRUE"} \OtherTok{=} \StringTok{"\#D55E00"}\NormalTok{)) }\SpecialCharTok{+}
  \FunctionTok{labs}\NormalTok{(}\AttributeTok{title    =} \StringTok{"Diagnostic Issues in Parameter Space"}\NormalTok{,}
       \AttributeTok{subtitle =} \StringTok{"Orange points: divergence, max treedepth, or Rhat \textgreater{} 1.01"}\NormalTok{,}
       \AttributeTok{x =} \StringTok{"True log(r)"}\NormalTok{, }\AttributeTok{y =} \StringTok{"True log(q)"}\NormalTok{) }\SpecialCharTok{+}
  \FunctionTok{theme}\NormalTok{(}\AttributeTok{legend.position =} \StringTok{"bottom"}\NormalTok{)}

\FunctionTok{print}\NormalTok{(p\_diag\_r\_q)}
\end{Highlighting}
\end{Shaded}

\begin{center}
\includegraphics[width=0.8\linewidth,height=\textheight,keepaspectratio]{12-CategorizationModelsPrototypes_files/figure-pdf/ch13_sbc_prototype-3.pdf}
\end{center}

\begin{Shaded}
\begin{Highlighting}[]
\CommentTok{\# ── Parameter recovery on SBC simulations ────────────────────────────────────}
\NormalTok{recovery\_df }\OtherTok{\textless{}{-}}\NormalTok{ sbc\_results\_proto}\SpecialCharTok{$}\NormalTok{stats}

\NormalTok{p\_sbc\_recovery }\OtherTok{\textless{}{-}} \FunctionTok{ggplot}\NormalTok{(recovery\_df, }\FunctionTok{aes}\NormalTok{(}\AttributeTok{x =}\NormalTok{ simulated\_value, }\AttributeTok{y =}\NormalTok{ mean)) }\SpecialCharTok{+}
  \FunctionTok{geom\_point}\NormalTok{(}\AttributeTok{alpha =} \FloatTok{0.4}\NormalTok{, }\AttributeTok{color =} \StringTok{"\#0072B2"}\NormalTok{) }\SpecialCharTok{+}
  \FunctionTok{geom\_abline}\NormalTok{(}\AttributeTok{intercept =} \DecValTok{0}\NormalTok{, }\AttributeTok{slope =} \DecValTok{1}\NormalTok{, }\AttributeTok{color =} \StringTok{"red"}\NormalTok{, }\AttributeTok{linetype =} \StringTok{"dashed"}\NormalTok{) }\SpecialCharTok{+}
  \FunctionTok{facet\_wrap}\NormalTok{(}\SpecialCharTok{\textasciitilde{}}\NormalTok{variable, }\AttributeTok{scales =} \StringTok{"free"}\NormalTok{) }\SpecialCharTok{+}
  \FunctionTok{labs}\NormalTok{(}
    \AttributeTok{title    =} \StringTok{"Parameter Recovery Across SBC Simulations"}\NormalTok{,}
    \AttributeTok{subtitle =} \StringTok{"Posterior mean vs True generated value"}\NormalTok{,}
    \AttributeTok{x =} \StringTok{"True Value (simulated\_value)"}\NormalTok{,}
    \AttributeTok{y =} \StringTok{"Estimated Value (Posterior Mean)"}
\NormalTok{  ) }\SpecialCharTok{+}
  \FunctionTok{theme\_bw}\NormalTok{()}

\FunctionTok{print}\NormalTok{(p\_sbc\_recovery)}
\end{Highlighting}
\end{Shaded}

\begin{center}
\includegraphics[width=0.8\linewidth,height=\textheight,keepaspectratio]{12-CategorizationModelsPrototypes_files/figure-pdf/ch13_sbc_prototype-4.pdf}
\end{center}

\textbf{Reading the SBC output.} With two parameters, the
ECDF-difference plot shows two curves (\texttt{log\_r} and
\texttt{log\_q}) that should each stay inside the grey simultaneous
confidence band. Systematic deviations indicate specific failure modes:

\begin{itemize}
\tightlist
\item
  \textbf{U-shape} in the rank histogram (high at the edges): the
  posterior is \emph{overconfident} --- credible intervals are too
  narrow and miss the truth at the tails. Action: widen the prior, or
  check whether the Kalman filter implementation is leaking information
  forward.
\item
  \textbf{Hump} in the centre: the posterior is \emph{underconfident}
  --- credible intervals are too wide. Action: tighten the prior, or
  check whether \texttt{log\_r} is being doubly-penalised somewhere.
\item
  \textbf{Monotone drift} (skewed left or right): a directional bias in
  \(\log r\). Action: re-derive the Stan code from scratch and compare
  against the R simulator on a single agent. The most common cause is a
  sign error or off-by-one in the Kalman update.
\end{itemize}

For the production version of this material we recommend \(S \geq 1000\)
iterations. The \(S = 500\) version above is illustrative.

\begin{center}\rule{0.5\linewidth}{0.5pt}\end{center}

\section{\texorpdfstring{Non-Stationary Environments: When Does \(q\)
Become
Identifiable?}{Non-Stationary Environments: When Does q Become Identifiable?}}\label{non-stationary-environments-when-does-q-become-identifiable}

The SBC result above exposes a structural problem. Under the static
Kruschke paradigm, \(\log r\) is cleanly recovered but \(\log q\) is not
--- its ECDF-difference curve drifts outside the simultaneous band and
its posterior often collapses toward the prior. This is not a coding bug
or a sampler problem: it is a genuine identifiability failure caused by
a mismatch between the \emph{parameter} (process noise on prototype
drift) and the \emph{task} (fixed category centres with no drift to
track).

Three generative processes stress this identifiability question
differently. We implement all three and run the full validation battery
on each.

As established in the canonical scenario framework introduced in Chapter
11, all three categorization chapters use the same three scenarios in
the same order (static → contingent shift → drift), reflecting
increasing assumption match between the scenario's data-generating
process and the model's structural commitments.

\begin{enumerate}
\def\labelenumi{\arabic{enumi}.}
\tightlist
\item
  \textbf{Static Kruschke (baseline).} Fixed category labels. The Kalman
  filter's process-noise prior assumes smooth Gaussian drift; under
  stationary data, no amount of \(q\) improves fit once transient
  learning is over. We expect \(\log q\) to remain weakly identified ---
  this is the failure case we need to diagnose clearly.
\item
  \textbf{Performance-contingent abrupt shifts.} After the agent reaches
  a streak of \(k\) consecutive correct responses, all category labels
  flip. Post-flip recovery is \emph{fast} for agents whose \(q\) lets
  the covariance re-inflate after surprise, and \emph{slow} for agents
  whose \(q\) is tiny --- so \(q\) acquires a direct behavioural
  signature. But the true process has abrupt jumps, which the filter's
  Gaussian-drift prior does not literally contain; this is a
  misspecified-but-informative scenario. Note that these shifts are
  \emph{endogenous} to the agent's own accuracy --- a noisier agent
  triggers fewer flips and accumulates less information about \(q\) (see
  Chapter 11 for discussion).
\item
  \textbf{Smooth drift.} Category centres in feature space follow an
  independent random walk; labels are assigned trial-by-trial by
  proximity to the current drifted centres. This is the
  \textbf{assumption-matched} scenario: the Kalman filter's generative
  story (prototypes drift as Gaussian random walks) is now literally
  true of the environment. We expect \(\log q\) to be cleanly identified
  and SBC to pass.
\end{enumerate}

The three scenarios span the spectrum of assumption match ---
\textbf{none → misspecified → exact} --- and the SBC signatures should
sort in that order.

\subsection{Assumption Match Across
Scenarios}\label{assumption-match-across-scenarios-1}

\begin{longtable}[]{@{}
  >{\raggedright\arraybackslash}p{(\linewidth - 6\tabcolsep) * \real{0.2500}}
  >{\raggedright\arraybackslash}p{(\linewidth - 6\tabcolsep) * \real{0.2500}}
  >{\raggedright\arraybackslash}p{(\linewidth - 6\tabcolsep) * \real{0.2500}}
  >{\raggedright\arraybackslash}p{(\linewidth - 6\tabcolsep) * \real{0.2500}}@{}}
\toprule\noalign{}
\begin{minipage}[b]{\linewidth}\raggedright
Scenario
\end{minipage} & \begin{minipage}[b]{\linewidth}\raggedright
What drifts in the world
\end{minipage} & \begin{minipage}[b]{\linewidth}\raggedright
What the Kalman filter assumes
\end{minipage} & \begin{minipage}[b]{\linewidth}\raggedright
Match
\end{minipage} \\
\midrule\noalign{}
\endhead
\bottomrule\noalign{}
\endlastfoot
Static & Nothing & Gaussian drift & \textbf{Misspecified} (but harmless:
\(q = 0\) limit works) \\
Contingent shift & Abrupt label reversals & Gaussian drift &
\textbf{Misspecified} (filter smooths the reversal) \\
Smooth drift & Gaussian-like centre drift & Gaussian drift &
\textbf{Exact} (this is the model's native story) \\
\end{longtable}

This is the mirror of Chapter 11's GCM-with-decay analysis: the GCM's
decay kernel fits contingent shifts as well as drift (because ``forget
old stuff'' is a general-purpose mechanism), while the Kalman filter's
smooth-drift prior fits smooth drift cleanly and reversals only
approximately.

\subsection{Scenario-Aware Simulator}\label{scenario-aware-simulator}

\begin{Shaded}
\begin{Highlighting}[]
\NormalTok{simulate\_prototype\_scenario }\OtherTok{\textless{}{-}} \ControlFlowTok{function}\NormalTok{(r\_value, q\_value, schedule,}
                                        \AttributeTok{scenario =} \FunctionTok{c}\NormalTok{(}\StringTok{"static"}\NormalTok{,}
                                                     \StringTok{"contingent\_shift"}\NormalTok{,}
                                                     \StringTok{"drift"}\NormalTok{),}
                                        \AttributeTok{streak\_target      =} \DecValTok{8}\NormalTok{,}
                                        \AttributeTok{drift\_traj         =} \ConstantTok{NULL}\NormalTok{,}
                                        \AttributeTok{initial\_mu\_cat0    =} \FunctionTok{c}\NormalTok{(}\FloatTok{2.5}\NormalTok{, }\FloatTok{2.5}\NormalTok{),}
                                        \AttributeTok{initial\_mu\_cat1    =} \FunctionTok{c}\NormalTok{(}\FloatTok{2.5}\NormalTok{, }\FloatTok{2.5}\NormalTok{),}
                                        \AttributeTok{initial\_sigma\_diag =} \FloatTok{10.0}\NormalTok{) \{}
\NormalTok{  scenario }\OtherTok{\textless{}{-}} \FunctionTok{match.arg}\NormalTok{(scenario)}
\NormalTok{  ntrials  }\OtherTok{\textless{}{-}} \FunctionTok{nrow}\NormalTok{(schedule)}
\NormalTok{  obs\_mat  }\OtherTok{\textless{}{-}} \FunctionTok{as.matrix}\NormalTok{(schedule[, }\FunctionTok{c}\NormalTok{(}\StringTok{"height"}\NormalTok{, }\StringTok{"position"}\NormalTok{)])}
\NormalTok{  base\_lab }\OtherTok{\textless{}{-}}\NormalTok{ schedule}\SpecialCharTok{$}\NormalTok{category\_feedback}
\NormalTok{  nfeat    }\OtherTok{\textless{}{-}} \DecValTok{2}\NormalTok{L}

\NormalTok{  mu0    }\OtherTok{\textless{}{-}}\NormalTok{ initial\_mu\_cat0}
\NormalTok{  mu1    }\OtherTok{\textless{}{-}}\NormalTok{ initial\_mu\_cat1}
\NormalTok{  sigma0 }\OtherTok{\textless{}{-}} \FunctionTok{diag}\NormalTok{(initial\_sigma\_diag, nfeat)}
\NormalTok{  sigma1 }\OtherTok{\textless{}{-}} \FunctionTok{diag}\NormalTok{(initial\_sigma\_diag, nfeat)}
\NormalTok{  R\_mat  }\OtherTok{\textless{}{-}} \FunctionTok{diag}\NormalTok{(r\_value, nfeat)}
\NormalTok{  Q\_mat  }\OtherTok{\textless{}{-}} \FunctionTok{diag}\NormalTok{(q\_value, nfeat)}
\NormalTok{  I\_mat  }\OtherTok{\textless{}{-}} \FunctionTok{diag}\NormalTok{(nfeat)}

\NormalTok{  label\_flip }\OtherTok{\textless{}{-}} \ConstantTok{FALSE}
\NormalTok{  streak     }\OtherTok{\textless{}{-}} \DecValTok{0}\NormalTok{L}

\NormalTok{  prob\_cat1         }\OtherTok{\textless{}{-}} \FunctionTok{numeric}\NormalTok{(ntrials)}
\NormalTok{  sim\_response      }\OtherTok{\textless{}{-}} \FunctionTok{integer}\NormalTok{(ntrials)}
\NormalTok{  observed\_feedback }\OtherTok{\textless{}{-}} \FunctionTok{integer}\NormalTok{(ntrials)}
\NormalTok{  label\_flip\_trace  }\OtherTok{\textless{}{-}} \FunctionTok{logical}\NormalTok{(ntrials)}

  \ControlFlowTok{for}\NormalTok{ (i }\ControlFlowTok{in} \FunctionTok{seq\_len}\NormalTok{(ntrials)) \{}
\NormalTok{    x }\OtherTok{\textless{}{-}} \FunctionTok{as.numeric}\NormalTok{(obs\_mat[i, ])}

    \CommentTok{\# ── Prediction step ──────────────────────────────────────────────}
\NormalTok{    sigma0 }\OtherTok{\textless{}{-}}\NormalTok{ sigma0 }\SpecialCharTok{+}\NormalTok{ Q\_mat}
\NormalTok{    sigma1 }\OtherTok{\textless{}{-}}\NormalTok{ sigma1 }\SpecialCharTok{+}\NormalTok{ Q\_mat}

    \CommentTok{\# ── Decision ─────────────────────────────────────────────────────}
\NormalTok{    cov0 }\OtherTok{\textless{}{-}}\NormalTok{ sigma0 }\SpecialCharTok{+}\NormalTok{ R\_mat}
\NormalTok{    cov1 }\OtherTok{\textless{}{-}}\NormalTok{ sigma1 }\SpecialCharTok{+}\NormalTok{ R\_mat}
\NormalTok{    lp0  }\OtherTok{\textless{}{-}} \FunctionTok{tryCatch}\NormalTok{(mvtnorm}\SpecialCharTok{::}\FunctionTok{dmvnorm}\NormalTok{(x, }\AttributeTok{mean =}\NormalTok{ mu0, }\AttributeTok{sigma =}\NormalTok{ cov0, }\AttributeTok{log =} \ConstantTok{TRUE}\NormalTok{),}
                     \AttributeTok{error =} \ControlFlowTok{function}\NormalTok{(e) }\SpecialCharTok{{-}}\ConstantTok{Inf}\NormalTok{)}
\NormalTok{    lp1  }\OtherTok{\textless{}{-}} \FunctionTok{tryCatch}\NormalTok{(mvtnorm}\SpecialCharTok{::}\FunctionTok{dmvnorm}\NormalTok{(x, }\AttributeTok{mean =}\NormalTok{ mu1, }\AttributeTok{sigma =}\NormalTok{ cov1, }\AttributeTok{log =} \ConstantTok{TRUE}\NormalTok{),}
                     \AttributeTok{error =} \ControlFlowTok{function}\NormalTok{(e) }\SpecialCharTok{{-}}\ConstantTok{Inf}\NormalTok{)}
\NormalTok{    m    }\OtherTok{\textless{}{-}} \FunctionTok{max}\NormalTok{(lp0, lp1)}
\NormalTok{    p    }\OtherTok{\textless{}{-}} \FunctionTok{exp}\NormalTok{(lp1 }\SpecialCharTok{{-}}\NormalTok{ (m }\SpecialCharTok{+} \FunctionTok{log}\NormalTok{(}\FunctionTok{exp}\NormalTok{(lp0 }\SpecialCharTok{{-}}\NormalTok{ m) }\SpecialCharTok{+} \FunctionTok{exp}\NormalTok{(lp1 }\SpecialCharTok{{-}}\NormalTok{ m))))}
\NormalTok{    p    }\OtherTok{\textless{}{-}} \FunctionTok{pmax}\NormalTok{(}\FloatTok{1e{-}9}\NormalTok{, }\FunctionTok{pmin}\NormalTok{(}\DecValTok{1} \SpecialCharTok{{-}} \FloatTok{1e{-}9}\NormalTok{, p))}
\NormalTok{    prob\_cat1[i]    }\OtherTok{\textless{}{-}}\NormalTok{ p}
\NormalTok{    sim\_response[i] }\OtherTok{\textless{}{-}} \FunctionTok{rbinom}\NormalTok{(}\DecValTok{1}\NormalTok{, }\DecValTok{1}\NormalTok{, p)}

    \CommentTok{\# ── Observed feedback depends on the scenario ────────────────────}
\NormalTok{    fb }\OtherTok{\textless{}{-}} \ControlFlowTok{switch}\NormalTok{(scenario,}
      \AttributeTok{static           =}\NormalTok{ base\_lab[i],}
      \AttributeTok{contingent\_shift =} \ControlFlowTok{if}\NormalTok{ (label\_flip) }\DecValTok{1}\NormalTok{L }\SpecialCharTok{{-}}\NormalTok{ base\_lab[i] }\ControlFlowTok{else}\NormalTok{ base\_lab[i],}
      \AttributeTok{drift            =}\NormalTok{ \{}
\NormalTok{        d0 }\OtherTok{\textless{}{-}} \FunctionTok{sum}\NormalTok{(}\FunctionTok{abs}\NormalTok{(x }\SpecialCharTok{{-}}\NormalTok{ drift\_traj}\SpecialCharTok{$}\NormalTok{mu0[i, ]))}
\NormalTok{        d1 }\OtherTok{\textless{}{-}} \FunctionTok{sum}\NormalTok{(}\FunctionTok{abs}\NormalTok{(x }\SpecialCharTok{{-}}\NormalTok{ drift\_traj}\SpecialCharTok{$}\NormalTok{mu1[i, ]))}
        \FunctionTok{as.integer}\NormalTok{(d1 }\SpecialCharTok{\textless{}}\NormalTok{ d0)}
\NormalTok{      \}}
\NormalTok{    )}
\NormalTok{    observed\_feedback[i] }\OtherTok{\textless{}{-}}\NormalTok{ fb}
\NormalTok{    label\_flip\_trace[i]  }\OtherTok{\textless{}{-}}\NormalTok{ label\_flip}

    \CommentTok{\# Update the streak/flip state for the contingent scenario}
    \ControlFlowTok{if}\NormalTok{ (scenario }\SpecialCharTok{==} \StringTok{"contingent\_shift"}\NormalTok{) \{}
\NormalTok{      streak }\OtherTok{\textless{}{-}} \ControlFlowTok{if}\NormalTok{ (sim\_response[i] }\SpecialCharTok{==}\NormalTok{ fb) streak }\SpecialCharTok{+} \DecValTok{1}\NormalTok{L }\ControlFlowTok{else} \DecValTok{0}\NormalTok{L}
      \ControlFlowTok{if}\NormalTok{ (streak }\SpecialCharTok{\textgreater{}=}\NormalTok{ streak\_target) \{ label\_flip }\OtherTok{\textless{}{-}} \SpecialCharTok{!}\NormalTok{label\_flip; streak }\OtherTok{\textless{}{-}} \DecValTok{0}\NormalTok{L \}}
\NormalTok{    \}}

    \CommentTok{\# ── Measurement update for the observed{-}feedback category ────────}
    \ControlFlowTok{if}\NormalTok{ (fb }\SpecialCharTok{==} \DecValTok{1}\NormalTok{) \{}
\NormalTok{      innov  }\OtherTok{\textless{}{-}}\NormalTok{ x }\SpecialCharTok{{-}}\NormalTok{ mu1}
\NormalTok{      S      }\OtherTok{\textless{}{-}}\NormalTok{ sigma1 }\SpecialCharTok{+}\NormalTok{ R\_mat}
\NormalTok{      K      }\OtherTok{\textless{}{-}}\NormalTok{ sigma1 }\SpecialCharTok{\%*\%} \FunctionTok{solve}\NormalTok{(S)}
\NormalTok{      mu1    }\OtherTok{\textless{}{-}} \FunctionTok{as.numeric}\NormalTok{(mu1 }\SpecialCharTok{+}\NormalTok{ K }\SpecialCharTok{\%*\%}\NormalTok{ innov)}
\NormalTok{      IK     }\OtherTok{\textless{}{-}}\NormalTok{ I\_mat }\SpecialCharTok{{-}}\NormalTok{ K}
\NormalTok{      sigma1 }\OtherTok{\textless{}{-}}\NormalTok{ IK }\SpecialCharTok{\%*\%}\NormalTok{ sigma1 }\SpecialCharTok{\%*\%} \FunctionTok{t}\NormalTok{(IK) }\SpecialCharTok{+}\NormalTok{ K }\SpecialCharTok{\%*\%}\NormalTok{ R\_mat }\SpecialCharTok{\%*\%} \FunctionTok{t}\NormalTok{(K)}
\NormalTok{      sigma1 }\OtherTok{\textless{}{-}}\NormalTok{ (sigma1 }\SpecialCharTok{+} \FunctionTok{t}\NormalTok{(sigma1)) }\SpecialCharTok{/} \DecValTok{2}
\NormalTok{    \} }\ControlFlowTok{else}\NormalTok{ \{}
\NormalTok{      innov  }\OtherTok{\textless{}{-}}\NormalTok{ x }\SpecialCharTok{{-}}\NormalTok{ mu0}
\NormalTok{      S      }\OtherTok{\textless{}{-}}\NormalTok{ sigma0 }\SpecialCharTok{+}\NormalTok{ R\_mat}
\NormalTok{      K      }\OtherTok{\textless{}{-}}\NormalTok{ sigma0 }\SpecialCharTok{\%*\%} \FunctionTok{solve}\NormalTok{(S)}
\NormalTok{      mu0    }\OtherTok{\textless{}{-}} \FunctionTok{as.numeric}\NormalTok{(mu0 }\SpecialCharTok{+}\NormalTok{ K }\SpecialCharTok{\%*\%}\NormalTok{ innov)}
\NormalTok{      IK     }\OtherTok{\textless{}{-}}\NormalTok{ I\_mat }\SpecialCharTok{{-}}\NormalTok{ K}
\NormalTok{      sigma0 }\OtherTok{\textless{}{-}}\NormalTok{ IK }\SpecialCharTok{\%*\%}\NormalTok{ sigma0 }\SpecialCharTok{\%*\%} \FunctionTok{t}\NormalTok{(IK) }\SpecialCharTok{+}\NormalTok{ K }\SpecialCharTok{\%*\%}\NormalTok{ R\_mat }\SpecialCharTok{\%*\%} \FunctionTok{t}\NormalTok{(K)}
\NormalTok{      sigma0 }\OtherTok{\textless{}{-}}\NormalTok{ (sigma0 }\SpecialCharTok{+} \FunctionTok{t}\NormalTok{(sigma0)) }\SpecialCharTok{/} \DecValTok{2}
\NormalTok{    \}}
\NormalTok{  \}}

\NormalTok{  schedule }\SpecialCharTok{|\textgreater{}}
    \FunctionTok{mutate}\NormalTok{(}
      \AttributeTok{prob\_cat1         =}\NormalTok{ prob\_cat1,}
      \AttributeTok{sim\_response      =}\NormalTok{ sim\_response,}
      \AttributeTok{observed\_feedback =}\NormalTok{ observed\_feedback,}
      \AttributeTok{label\_flip        =}\NormalTok{ label\_flip\_trace,}
      \AttributeTok{correct           =} \FunctionTok{as.integer}\NormalTok{(sim\_response }\SpecialCharTok{==}\NormalTok{ observed\_feedback)}
\NormalTok{    )}
\NormalTok{\}}
\end{Highlighting}
\end{Shaded}

\subsection{Drift Trajectory
Generator}\label{drift-trajectory-generator-1}

The drift scenario pre-computes a trajectory for both category centres
as independent bivariate Gaussian random walks.

\begin{Shaded}
\begin{Highlighting}[]
\NormalTok{make\_drift\_trajectory }\OtherTok{\textless{}{-}} \ControlFlowTok{function}\NormalTok{(ntrials, }\AttributeTok{drift\_sigma =} \FloatTok{0.05}\NormalTok{, }\AttributeTok{seed =} \DecValTok{1}\NormalTok{) \{}
  \FunctionTok{set.seed}\NormalTok{(seed)}
\NormalTok{  mu0\_init }\OtherTok{\textless{}{-}} \FunctionTok{c}\NormalTok{(}\FloatTok{1.75}\NormalTok{, }\FloatTok{3.25}\NormalTok{)   }\CommentTok{\# Kruschke cat{-}0 mean (height, position)}
\NormalTok{  mu1\_init }\OtherTok{\textless{}{-}} \FunctionTok{c}\NormalTok{(}\FloatTok{3.25}\NormalTok{, }\FloatTok{1.75}\NormalTok{)   }\CommentTok{\# Kruschke cat{-}1 mean}
\NormalTok{  d0 }\OtherTok{\textless{}{-}} \FunctionTok{apply}\NormalTok{(}\FunctionTok{matrix}\NormalTok{(}\FunctionTok{rnorm}\NormalTok{(ntrials }\SpecialCharTok{*} \DecValTok{2}\NormalTok{, }\DecValTok{0}\NormalTok{, drift\_sigma), }\AttributeTok{ncol =} \DecValTok{2}\NormalTok{), }\DecValTok{2}\NormalTok{, cumsum)}
\NormalTok{  d1 }\OtherTok{\textless{}{-}} \FunctionTok{apply}\NormalTok{(}\FunctionTok{matrix}\NormalTok{(}\FunctionTok{rnorm}\NormalTok{(ntrials }\SpecialCharTok{*} \DecValTok{2}\NormalTok{, }\DecValTok{0}\NormalTok{, drift\_sigma), }\AttributeTok{ncol =} \DecValTok{2}\NormalTok{), }\DecValTok{2}\NormalTok{, cumsum)}
  \FunctionTok{list}\NormalTok{(}
    \AttributeTok{mu0 =} \FunctionTok{sweep}\NormalTok{(d0, }\DecValTok{2}\NormalTok{, mu0\_init, }\StringTok{"+"}\NormalTok{),}
    \AttributeTok{mu1 =} \FunctionTok{sweep}\NormalTok{(d1, }\DecValTok{2}\NormalTok{, mu1\_init, }\StringTok{"+"}\NormalTok{)}
\NormalTok{  )}
\NormalTok{\}}
\end{Highlighting}
\end{Shaded}

\subsection{Scenario Agent Wrapper}\label{scenario-agent-wrapper}

\begin{Shaded}
\begin{Highlighting}[]
\NormalTok{simulate\_prototype\_agent\_scenario }\OtherTok{\textless{}{-}} \ControlFlowTok{function}\NormalTok{(agent\_id, scenario,}
\NormalTok{                                              r\_value, q\_value,}
\NormalTok{                                              subject\_seed,}
                                              \AttributeTok{streak\_target =} \DecValTok{8}\NormalTok{,}
                                              \AttributeTok{drift\_sigma   =} \FloatTok{0.05}\NormalTok{) \{}
\NormalTok{  schedule }\OtherTok{\textless{}{-}} \FunctionTok{make\_subject\_schedule}\NormalTok{(stimulus\_info, n\_blocks, }\AttributeTok{seed =}\NormalTok{ subject\_seed)}
\NormalTok{  drift\_tr }\OtherTok{\textless{}{-}} \ControlFlowTok{if}\NormalTok{ (scenario }\SpecialCharTok{==} \StringTok{"drift"}\NormalTok{)}
    \FunctionTok{make\_drift\_trajectory}\NormalTok{(}\FunctionTok{nrow}\NormalTok{(schedule), drift\_sigma, }\AttributeTok{seed =}\NormalTok{ subject\_seed)}
  \ControlFlowTok{else} \ConstantTok{NULL}

\NormalTok{  out }\OtherTok{\textless{}{-}} \FunctionTok{simulate\_prototype\_scenario}\NormalTok{(}
    \AttributeTok{r\_value       =}\NormalTok{ r\_value,}
    \AttributeTok{q\_value       =}\NormalTok{ q\_value,}
    \AttributeTok{schedule      =}\NormalTok{ schedule,}
    \AttributeTok{scenario      =}\NormalTok{ scenario,}
    \AttributeTok{streak\_target =}\NormalTok{ streak\_target,}
    \AttributeTok{drift\_traj    =}\NormalTok{ drift\_tr}
\NormalTok{  )}

\NormalTok{  out }\SpecialCharTok{|\textgreater{}}
    \FunctionTok{mutate}\NormalTok{(}
      \AttributeTok{agent\_id     =}\NormalTok{ agent\_id,}
      \AttributeTok{scenario     =}\NormalTok{ scenario,}
      \AttributeTok{r\_value\_true =}\NormalTok{ r\_value,}
      \AttributeTok{q\_value\_true =}\NormalTok{ q\_value,}
      \AttributeTok{log\_r\_true   =} \FunctionTok{log}\NormalTok{(r\_value),}
      \AttributeTok{log\_q\_true   =} \FunctionTok{log}\NormalTok{(q\_value)}
\NormalTok{    ) }\SpecialCharTok{|\textgreater{}}
    \FunctionTok{group\_by}\NormalTok{(agent\_id) }\SpecialCharTok{|\textgreater{}}
    \FunctionTok{mutate}\NormalTok{(}\AttributeTok{performance =} \FunctionTok{cumsum}\NormalTok{(correct) }\SpecialCharTok{/} \FunctionTok{row\_number}\NormalTok{()) }\SpecialCharTok{|\textgreater{}}
    \FunctionTok{ungroup}\NormalTok{()}
\NormalTok{\}}
\end{Highlighting}
\end{Shaded}

\subsection{Visualising the Three
Scenarios}\label{visualising-the-three-scenarios-1}

We simulate 10 agents per cell across a 3 × 3 grid of scenarios × \(q\)
values (at fixed \(r = 1.5\)) and plot cumulative accuracy.

\begin{Shaded}
\begin{Highlighting}[]
\NormalTok{scenario\_viz\_file }\OtherTok{\textless{}{-}} \FunctionTok{here}\NormalTok{(}\StringTok{"simdata"}\NormalTok{, }\StringTok{"ch13\_scenario\_viz.csv"}\NormalTok{)}

\ControlFlowTok{if}\NormalTok{ (regenerate\_simulations }\SpecialCharTok{||} \SpecialCharTok{!}\FunctionTok{file.exists}\NormalTok{(scenario\_viz\_file)) \{}
\NormalTok{  viz\_grid }\OtherTok{\textless{}{-}} \FunctionTok{expand\_grid}\NormalTok{(}
    \AttributeTok{scenario    =} \FunctionTok{c}\NormalTok{(}\StringTok{"static"}\NormalTok{, }\StringTok{"contingent\_shift"}\NormalTok{, }\StringTok{"drift"}\NormalTok{),}
    \AttributeTok{q\_value     =} \FunctionTok{c}\NormalTok{(}\FloatTok{0.02}\NormalTok{, }\FloatTok{0.15}\NormalTok{, }\FloatTok{0.6}\NormalTok{),}
    \AttributeTok{agent\_id    =} \DecValTok{1}\SpecialCharTok{:}\DecValTok{10}
\NormalTok{  ) }\SpecialCharTok{|\textgreater{}}
    \FunctionTok{mutate}\NormalTok{(}\AttributeTok{subject\_seed =} \DecValTok{30000} \SpecialCharTok{+} \FunctionTok{row\_number}\NormalTok{())}

\NormalTok{  viz\_sim }\OtherTok{\textless{}{-}} \FunctionTok{pmap\_dfr}\NormalTok{(}
    \FunctionTok{list}\NormalTok{(}\AttributeTok{agent\_id     =}\NormalTok{ viz\_grid}\SpecialCharTok{$}\NormalTok{agent\_id,}
         \AttributeTok{scenario     =}\NormalTok{ viz\_grid}\SpecialCharTok{$}\NormalTok{scenario,}
         \AttributeTok{q\_value      =}\NormalTok{ viz\_grid}\SpecialCharTok{$}\NormalTok{q\_value,}
         \AttributeTok{subject\_seed =}\NormalTok{ viz\_grid}\SpecialCharTok{$}\NormalTok{subject\_seed),}
    \ControlFlowTok{function}\NormalTok{(agent\_id, scenario, q\_value, subject\_seed) \{}
      \FunctionTok{simulate\_prototype\_agent\_scenario}\NormalTok{(}
        \AttributeTok{agent\_id     =}\NormalTok{ agent\_id,}
        \AttributeTok{scenario     =}\NormalTok{ scenario,}
        \AttributeTok{r\_value      =} \FloatTok{1.5}\NormalTok{,}
        \AttributeTok{q\_value      =}\NormalTok{ q\_value,}
        \AttributeTok{subject\_seed =}\NormalTok{ subject\_seed}
\NormalTok{      )}
\NormalTok{    \}}
\NormalTok{  )}
  \FunctionTok{write\_csv}\NormalTok{(viz\_sim, scenario\_viz\_file)}
\NormalTok{\} }\ControlFlowTok{else}\NormalTok{ \{}
\NormalTok{  viz\_sim }\OtherTok{\textless{}{-}} \FunctionTok{read\_csv}\NormalTok{(scenario\_viz\_file, }\AttributeTok{show\_col\_types =} \ConstantTok{FALSE}\NormalTok{)}
\NormalTok{\}}

\FunctionTok{ggplot}\NormalTok{(viz\_sim,}
       \FunctionTok{aes}\NormalTok{(}\AttributeTok{x =}\NormalTok{ trial\_within\_subject, }\AttributeTok{y =}\NormalTok{ cumulative\_accuracy,}
           \AttributeTok{color =} \FunctionTok{factor}\NormalTok{(q\_value\_true))) }\SpecialCharTok{+}
  \FunctionTok{stat\_summary}\NormalTok{(}\AttributeTok{fun =}\NormalTok{ mean, }\AttributeTok{geom =} \StringTok{"line"}\NormalTok{,}
               \FunctionTok{aes}\NormalTok{(}\AttributeTok{group =} \FunctionTok{factor}\NormalTok{(q\_value\_true)), }\AttributeTok{linewidth =} \FloatTok{1.1}\NormalTok{) }\SpecialCharTok{+}
  \FunctionTok{stat\_summary}\NormalTok{(}\AttributeTok{fun.data =}\NormalTok{ mean\_se, }\AttributeTok{geom =} \StringTok{"ribbon"}\NormalTok{, }\AttributeTok{alpha =} \FloatTok{0.15}\NormalTok{,}
               \FunctionTok{aes}\NormalTok{(}\AttributeTok{fill =} \FunctionTok{factor}\NormalTok{(q\_value\_true), }\AttributeTok{group =} \FunctionTok{factor}\NormalTok{(q\_value\_true)),}
               \AttributeTok{color =} \ConstantTok{NA}\NormalTok{) }\SpecialCharTok{+}
  \FunctionTok{facet\_wrap}\NormalTok{(}\SpecialCharTok{\textasciitilde{}}\NormalTok{scenario, }\AttributeTok{ncol =} \DecValTok{3}\NormalTok{,}
             \AttributeTok{labeller =} \FunctionTok{as\_labeller}\NormalTok{(}\FunctionTok{c}\NormalTok{(}
               \AttributeTok{static           =} \StringTok{"Static Kruschke"}\NormalTok{,}
               \AttributeTok{contingent\_shift =} \StringTok{"Performance{-}contingent flips"}\NormalTok{,}
               \AttributeTok{drift            =} \StringTok{"Smooth drift"}\NormalTok{))) }\SpecialCharTok{+}
  \FunctionTok{scale\_color\_viridis\_d}\NormalTok{(}\AttributeTok{option =} \StringTok{"plasma"}\NormalTok{, }\AttributeTok{name =} \StringTok{"q"}\NormalTok{) }\SpecialCharTok{+}
  \FunctionTok{scale\_fill\_viridis\_d}\NormalTok{(}\AttributeTok{option  =} \StringTok{"plasma"}\NormalTok{, }\AttributeTok{name =} \StringTok{"q"}\NormalTok{) }\SpecialCharTok{+}
  \FunctionTok{geom\_hline}\NormalTok{(}\AttributeTok{yintercept =} \FloatTok{0.5}\NormalTok{, }\AttributeTok{linetype =} \StringTok{"dashed"}\NormalTok{, }\AttributeTok{color =} \StringTok{"grey50"}\NormalTok{) }\SpecialCharTok{+}
  \FunctionTok{labs}\NormalTok{(}
    \AttributeTok{title    =} \StringTok{"Prototype Learning Curves Across Scenarios"}\NormalTok{,}
    \AttributeTok{subtitle =} \StringTok{"Mean ± SE across 10 agents per cell; r = 1.5"}\NormalTok{,}
    \AttributeTok{x =} \StringTok{"Trial"}\NormalTok{, }\AttributeTok{y =} \StringTok{"Cumulative accuracy"}
\NormalTok{  ) }\SpecialCharTok{+}
  \FunctionTok{theme}\NormalTok{(}\AttributeTok{legend.position =} \StringTok{"bottom"}\NormalTok{)}
\end{Highlighting}
\end{Shaded}

\begin{center}
\includegraphics[width=0.8\linewidth,height=\textheight,keepaspectratio]{12-CategorizationModelsPrototypes_files/figure-pdf/ch13_scenario_viz-1.pdf}
\end{center}

\textbf{Reading the panels.} The static panel looks much like the plots
earlier in the chapter --- accuracy climbs, asymptotes below ceiling for
larger \(q\) (because the covariance never fully shrinks). The
contingent-shift panel shows characteristic sawtooth recovery: higher
\(q\) produces sharper but faster-rebounding dips. The drift panel shows
all three \(q\) curves plateauing below ceiling, with the plateau
highest for the \(q\) tuned to the environmental drift amplitude.

\subsection{Fitting One Agent per Scenario (Original
Parameterisation)}\label{fitting-one-agent-per-scenario-original-parameterisation}

We pick one representative agent per scenario with \(r = 1.5\),
\(q = 0.15\), and fit the Kalman prototype model (original
\texttt{(log\_r,\ log\_q)} parameterisation) to each.

\begin{Shaded}
\begin{Highlighting}[]
\NormalTok{scenario\_fit\_file }\OtherTok{\textless{}{-}} \FunctionTok{here}\NormalTok{(}\StringTok{"simmodels"}\NormalTok{, }\StringTok{"ch13\_prototype\_scenario\_fits.rds"}\NormalTok{)}

\NormalTok{scenario\_single\_fit\_data }\OtherTok{\textless{}{-}} \FunctionTok{list}\NormalTok{(}
  \AttributeTok{static           =} \FunctionTok{simulate\_prototype\_agent\_scenario}\NormalTok{(}
                        \DecValTok{1}\NormalTok{, }\StringTok{"static"}\NormalTok{,           }\AttributeTok{r\_value =} \FloatTok{1.5}\NormalTok{, }\AttributeTok{q\_value =} \FloatTok{0.15}\NormalTok{,}
                        \AttributeTok{subject\_seed =} \DecValTok{40001}\NormalTok{),}
  \AttributeTok{contingent\_shift =} \FunctionTok{simulate\_prototype\_agent\_scenario}\NormalTok{(}
                        \DecValTok{1}\NormalTok{, }\StringTok{"contingent\_shift"}\NormalTok{, }\AttributeTok{r\_value =} \FloatTok{1.5}\NormalTok{, }\AttributeTok{q\_value =} \FloatTok{0.15}\NormalTok{,}
                        \AttributeTok{subject\_seed =} \DecValTok{40002}\NormalTok{),}
  \AttributeTok{drift            =} \FunctionTok{simulate\_prototype\_agent\_scenario}\NormalTok{(}
                        \DecValTok{1}\NormalTok{, }\StringTok{"drift"}\NormalTok{,            }\AttributeTok{r\_value =} \FloatTok{1.5}\NormalTok{, }\AttributeTok{q\_value =} \FloatTok{0.15}\NormalTok{,}
                        \AttributeTok{subject\_seed =} \DecValTok{40003}\NormalTok{)}
\NormalTok{)}

\NormalTok{stan\_data\_from\_scenario\_agent }\OtherTok{\textless{}{-}} \ControlFlowTok{function}\NormalTok{(agent\_df) \{}
  \FunctionTok{list}\NormalTok{(}
    \AttributeTok{ntrials            =} \FunctionTok{nrow}\NormalTok{(agent\_df),}
    \AttributeTok{nfeatures          =} \DecValTok{2}\NormalTok{L,}
    \AttributeTok{cat\_one            =} \FunctionTok{as.integer}\NormalTok{(agent\_df}\SpecialCharTok{$}\NormalTok{observed\_feedback),}
    \AttributeTok{y                  =} \FunctionTok{as.integer}\NormalTok{(agent\_df}\SpecialCharTok{$}\NormalTok{sim\_response),}
    \AttributeTok{obs                =} \FunctionTok{as.matrix}\NormalTok{(agent\_df[, }\FunctionTok{c}\NormalTok{(}\StringTok{"height"}\NormalTok{, }\StringTok{"position"}\NormalTok{)]),}
    \AttributeTok{initial\_mu\_cat0    =} \FunctionTok{c}\NormalTok{(}\FloatTok{2.5}\NormalTok{, }\FloatTok{2.5}\NormalTok{),}
    \AttributeTok{initial\_mu\_cat1    =} \FunctionTok{c}\NormalTok{(}\FloatTok{2.5}\NormalTok{, }\FloatTok{2.5}\NormalTok{),}
    \AttributeTok{initial\_sigma\_diag =} \FloatTok{10.0}\NormalTok{,}
    \AttributeTok{prior\_logr\_mean    =}\NormalTok{ LOG\_R\_PRIOR\_MEAN,}
    \AttributeTok{prior\_logr\_sd      =}\NormalTok{ LOG\_R\_PRIOR\_SD,}
    \AttributeTok{prior\_logq\_mean    =}\NormalTok{ LOG\_Q\_PRIOR\_MEAN,}
    \AttributeTok{prior\_logq\_sd      =}\NormalTok{ LOG\_Q\_PRIOR\_SD}
\NormalTok{  )}
\NormalTok{\}}

\ControlFlowTok{if}\NormalTok{ (regenerate\_simulations }\SpecialCharTok{||} \SpecialCharTok{!}\FunctionTok{file.exists}\NormalTok{(scenario\_fit\_file)) \{}
\NormalTok{  scenario\_fits }\OtherTok{\textless{}{-}}\NormalTok{ purrr}\SpecialCharTok{::}\FunctionTok{imap}\NormalTok{(scenario\_single\_fit\_data, }\ControlFlowTok{function}\NormalTok{(agent\_df, scen) \{}
\NormalTok{    dat }\OtherTok{\textless{}{-}} \FunctionTok{stan\_data\_from\_scenario\_agent}\NormalTok{(agent\_df)}
\NormalTok{    pf  }\OtherTok{\textless{}{-}} \FunctionTok{tryCatch}\NormalTok{(}
\NormalTok{      mod\_prototype\_single}\SpecialCharTok{$}\FunctionTok{pathfinder}\NormalTok{(}\AttributeTok{data =}\NormalTok{ dat, }\AttributeTok{num\_paths =} \DecValTok{4}\NormalTok{,}
                                      \AttributeTok{refresh =} \DecValTok{0}\NormalTok{, }\AttributeTok{seed =} \DecValTok{500}\NormalTok{),}
      \AttributeTok{error =} \ControlFlowTok{function}\NormalTok{(e) }\ConstantTok{NULL}
\NormalTok{    )}
\NormalTok{    fit }\OtherTok{\textless{}{-}}\NormalTok{ mod\_prototype\_single}\SpecialCharTok{$}\FunctionTok{sample}\NormalTok{(}
      \AttributeTok{data            =}\NormalTok{ dat,}
      \AttributeTok{init            =} \ControlFlowTok{if}\NormalTok{ (}\SpecialCharTok{!}\FunctionTok{is.null}\NormalTok{(pf)) pf }\ControlFlowTok{else} \FloatTok{0.5}\NormalTok{,}
      \AttributeTok{seed            =} \DecValTok{500}\NormalTok{,}
      \AttributeTok{chains          =} \DecValTok{4}\NormalTok{,}
      \AttributeTok{parallel\_chains =} \FunctionTok{min}\NormalTok{(}\DecValTok{4}\NormalTok{, }\FunctionTok{availableCores}\NormalTok{()),}
      \AttributeTok{iter\_warmup     =} \DecValTok{1000}\NormalTok{,}
      \AttributeTok{iter\_sampling   =} \DecValTok{1500}\NormalTok{,}
      \AttributeTok{refresh         =} \DecValTok{0}\NormalTok{,}
      \AttributeTok{adapt\_delta     =} \FloatTok{0.9}
\NormalTok{    )}
\NormalTok{    fit}\SpecialCharTok{$}\FunctionTok{save\_object}\NormalTok{(}\FunctionTok{here}\NormalTok{(}\StringTok{"simmodels"}\NormalTok{,}
                         \FunctionTok{paste0}\NormalTok{(}\StringTok{"ch13\_prototype\_scenario\_"}\NormalTok{, scen, }\StringTok{".rds"}\NormalTok{)))}
\NormalTok{    fit}
\NormalTok{  \})}
  \FunctionTok{saveRDS}\NormalTok{(}\FunctionTok{lapply}\NormalTok{(scenario\_fits, }\ControlFlowTok{function}\NormalTok{(f) f}\SpecialCharTok{$}\FunctionTok{output\_files}\NormalTok{()),}
\NormalTok{          scenario\_fit\_file)}
  \FunctionTok{cat}\NormalTok{(}\StringTok{"Scenario fits computed and saved.}\SpecialCharTok{\textbackslash{}n}\StringTok{"}\NormalTok{)}
\NormalTok{\} }\ControlFlowTok{else}\NormalTok{ \{}
\NormalTok{  scenario\_fits }\OtherTok{\textless{}{-}} \FunctionTok{setNames}\NormalTok{(}
    \FunctionTok{lapply}\NormalTok{(}\FunctionTok{names}\NormalTok{(scenario\_single\_fit\_data), }\ControlFlowTok{function}\NormalTok{(scen) \{}
      \FunctionTok{readRDS}\NormalTok{(}\FunctionTok{here}\NormalTok{(}\StringTok{"simmodels"}\NormalTok{,}
                   \FunctionTok{paste0}\NormalTok{(}\StringTok{"ch13\_prototype\_scenario\_"}\NormalTok{, scen, }\StringTok{".rds"}\NormalTok{)))}
\NormalTok{    \}),}
    \FunctionTok{names}\NormalTok{(scenario\_single\_fit\_data)}
\NormalTok{  )}
  \FunctionTok{cat}\NormalTok{(}\StringTok{"Loaded existing scenario fits.}\SpecialCharTok{\textbackslash{}n}\StringTok{"}\NormalTok{)}
\NormalTok{\}}
\end{Highlighting}
\end{Shaded}

\begin{verbatim}
Loaded existing scenario fits.
\end{verbatim}

\subsection{MCMC Diagnostic Battery Across
Scenarios}\label{mcmc-diagnostic-battery-across-scenarios-1}

\begin{Shaded}
\begin{Highlighting}[]
\NormalTok{scenario\_diag\_tbls }\OtherTok{\textless{}{-}}\NormalTok{ purrr}\SpecialCharTok{::}\FunctionTok{imap\_dfr}\NormalTok{(scenario\_fits, }\ControlFlowTok{function}\NormalTok{(fit, scen) \{}
\NormalTok{  tbl }\OtherTok{\textless{}{-}} \FunctionTok{diagnostic\_summary\_table}\NormalTok{(fit, }\AttributeTok{params =} \FunctionTok{c}\NormalTok{(}\StringTok{"log\_r"}\NormalTok{, }\StringTok{"log\_q"}\NormalTok{,}
                                                  \StringTok{"r\_value"}\NormalTok{, }\StringTok{"q\_value"}\NormalTok{))}
\NormalTok{  tbl }\SpecialCharTok{|\textgreater{}} \FunctionTok{mutate}\NormalTok{(}\AttributeTok{scenario =}\NormalTok{ scen)}
\NormalTok{\})}
\FunctionTok{print}\NormalTok{(scenario\_diag\_tbls }\SpecialCharTok{|\textgreater{}}
\NormalTok{        dplyr}\SpecialCharTok{::}\FunctionTok{select}\NormalTok{(scenario, metric, value, threshold, pass))}
\end{Highlighting}
\end{Shaded}

\begin{verbatim}
# A tibble: 18 x 5
   scenario         metric                           value threshold pass 
   <chr>            <chr>                            <dbl> <chr>     <lgl>
 1 static           Divergences (zero tolerance)    0      == 0      TRUE 
 2 static           Max rank-normalised R-hat       1      < 1.01    TRUE 
 3 static           Min bulk ESS                 2578      > 400     TRUE 
 4 static           Min tail ESS                 2361      > 400     TRUE 
 5 static           Min E-BFMI                      0.919  > 0.2     TRUE 
 6 static           Max MCSE / posterior SD         0.0197 < 0.05    TRUE 
 7 contingent_shift Divergences (zero tolerance)    0      == 0      TRUE 
 8 contingent_shift Max rank-normalised R-hat       1      < 1.01    TRUE 
 9 contingent_shift Min bulk ESS                 2840      > 400     TRUE 
10 contingent_shift Min tail ESS                 2272      > 400     TRUE 
11 contingent_shift Min E-BFMI                      0.94   > 0.2     TRUE 
12 contingent_shift Max MCSE / posterior SD         0.0188 < 0.05    TRUE 
13 drift            Divergences (zero tolerance)    0      == 0      TRUE 
14 drift            Max rank-normalised R-hat       1      < 1.01    TRUE 
15 drift            Min bulk ESS                 2892      > 400     TRUE 
16 drift            Min tail ESS                 2847      > 400     TRUE 
17 drift            Min E-BFMI                      0.949  > 0.2     TRUE 
18 drift            Max MCSE / posterior SD         0.0186 < 0.05    TRUE 
\end{verbatim}

\begin{Shaded}
\begin{Highlighting}[]
\FunctionTok{walk}\NormalTok{(}\FunctionTok{names}\NormalTok{(scenario\_fits), }\ControlFlowTok{function}\NormalTok{(scen) \{}
  \FunctionTok{cat}\NormalTok{(}\StringTok{"=== Pair plot:"}\NormalTok{, scen, }\StringTok{"===}\SpecialCharTok{\textbackslash{}n}\StringTok{"}\NormalTok{)}
  \FunctionTok{print}\NormalTok{(bayesplot}\SpecialCharTok{::}\FunctionTok{mcmc\_pairs}\NormalTok{(}
\NormalTok{    scenario\_fits[[scen]]}\SpecialCharTok{$}\FunctionTok{draws}\NormalTok{(}\FunctionTok{c}\NormalTok{(}\StringTok{"log\_r"}\NormalTok{, }\StringTok{"log\_q"}\NormalTok{)),}
    \AttributeTok{diag\_fun     =} \StringTok{"dens"}\NormalTok{,}
    \AttributeTok{off\_diag\_fun =} \StringTok{"hex"}\NormalTok{,}
    \AttributeTok{np           =} \FunctionTok{nuts\_params}\NormalTok{(scenario\_fits[[scen]])}
\NormalTok{  ))}
\NormalTok{\})}
\end{Highlighting}
\end{Shaded}

\begin{verbatim}
=== Pair plot: static ===
\end{verbatim}

\begin{center}
\includegraphics[width=0.8\linewidth,height=\textheight,keepaspectratio]{12-CategorizationModelsPrototypes_files/figure-pdf/ch13_scenario_diagnostics-1.pdf}
\end{center}

\begin{verbatim}
=== Pair plot: contingent_shift ===
\end{verbatim}

\begin{center}
\includegraphics[width=0.8\linewidth,height=\textheight,keepaspectratio]{12-CategorizationModelsPrototypes_files/figure-pdf/ch13_scenario_diagnostics-2.pdf}
\end{center}

\begin{verbatim}
=== Pair plot: drift ===
\end{verbatim}

\begin{center}
\includegraphics[width=0.8\linewidth,height=\textheight,keepaspectratio]{12-CategorizationModelsPrototypes_files/figure-pdf/ch13_scenario_diagnostics-3.pdf}
\end{center}

\textbf{What the pair plots reveal.} On static data the
\texttt{(log\_r,\ log\_q)} posterior sits on a strongly
negatively-correlated ridge --- the canonical identifiability failure.
On contingent-shift data the ridge is shorter and tilted: reversal
dynamics load on \(q\) specifically, so the degrees of freedom are no
longer degenerate. On drift data the posterior is roughly circular in
the two parameters; this is the visual signature of the
assumption-matched case.

\subsection{Prior → Posterior Updates Across
Scenarios}\label{prior-posterior-updates-across-scenarios}

\begin{Shaded}
\begin{Highlighting}[]
\NormalTok{scenario\_pp\_df }\OtherTok{\textless{}{-}}\NormalTok{ purrr}\SpecialCharTok{::}\FunctionTok{imap\_dfr}\NormalTok{(scenario\_fits, }\ControlFlowTok{function}\NormalTok{(fit, scen) \{}
\NormalTok{  d }\OtherTok{\textless{}{-}} \FunctionTok{as\_draws\_df}\NormalTok{(fit}\SpecialCharTok{$}\FunctionTok{draws}\NormalTok{(}\FunctionTok{c}\NormalTok{(}\StringTok{"log\_r"}\NormalTok{, }\StringTok{"log\_q"}\NormalTok{)))}
  \FunctionTok{bind\_rows}\NormalTok{(}
    \FunctionTok{tibble}\NormalTok{(}\AttributeTok{scenario =}\NormalTok{ scen, }\AttributeTok{parameter =} \StringTok{"log\_r"}\NormalTok{, }\AttributeTok{value =}\NormalTok{ d}\SpecialCharTok{$}\NormalTok{log\_r, }\AttributeTok{source =} \StringTok{"posterior"}\NormalTok{),}
    \FunctionTok{tibble}\NormalTok{(}\AttributeTok{scenario =}\NormalTok{ scen, }\AttributeTok{parameter =} \StringTok{"log\_q"}\NormalTok{, }\AttributeTok{value =}\NormalTok{ d}\SpecialCharTok{$}\NormalTok{log\_q, }\AttributeTok{source =} \StringTok{"posterior"}\NormalTok{)}
\NormalTok{  )}
\NormalTok{\})}

\FunctionTok{set.seed}\NormalTok{(}\DecValTok{21}\NormalTok{)}
\NormalTok{n\_prior }\OtherTok{\textless{}{-}} \DecValTok{4000}
\NormalTok{prior\_ref }\OtherTok{\textless{}{-}} \FunctionTok{bind\_rows}\NormalTok{(}
  \FunctionTok{tibble}\NormalTok{(}\AttributeTok{parameter =} \StringTok{"log\_r"}\NormalTok{, }\AttributeTok{value =} \FunctionTok{rnorm}\NormalTok{(n\_prior, LOG\_R\_PRIOR\_MEAN, LOG\_R\_PRIOR\_SD),}
         \AttributeTok{source =} \StringTok{"prior"}\NormalTok{),}
  \FunctionTok{tibble}\NormalTok{(}\AttributeTok{parameter =} \StringTok{"log\_q"}\NormalTok{, }\AttributeTok{value =} \FunctionTok{rnorm}\NormalTok{(n\_prior, LOG\_Q\_PRIOR\_MEAN, LOG\_Q\_PRIOR\_SD),}
         \AttributeTok{source =} \StringTok{"prior"}\NormalTok{)}
\NormalTok{) }\SpecialCharTok{|\textgreater{}}
\NormalTok{  tidyr}\SpecialCharTok{::}\FunctionTok{crossing}\NormalTok{(}\AttributeTok{scenario =} \FunctionTok{c}\NormalTok{(}\StringTok{"static"}\NormalTok{, }\StringTok{"contingent\_shift"}\NormalTok{, }\StringTok{"drift"}\NormalTok{))}

\NormalTok{truth\_lines }\OtherTok{\textless{}{-}} \FunctionTok{tibble}\NormalTok{(}
  \AttributeTok{parameter =} \FunctionTok{c}\NormalTok{(}\StringTok{"log\_r"}\NormalTok{, }\StringTok{"log\_q"}\NormalTok{),}
  \AttributeTok{value     =} \FunctionTok{c}\NormalTok{(}\FunctionTok{log}\NormalTok{(}\FloatTok{1.5}\NormalTok{), }\FunctionTok{log}\NormalTok{(}\FloatTok{0.15}\NormalTok{))}
\NormalTok{)}

\FunctionTok{ggplot}\NormalTok{(}\FunctionTok{bind\_rows}\NormalTok{(prior\_ref, scenario\_pp\_df),}
       \FunctionTok{aes}\NormalTok{(}\AttributeTok{x =}\NormalTok{ value, }\AttributeTok{fill =}\NormalTok{ source, }\AttributeTok{color =}\NormalTok{ source)) }\SpecialCharTok{+}
  \FunctionTok{geom\_density}\NormalTok{(}\AttributeTok{alpha =} \FloatTok{0.4}\NormalTok{, }\AttributeTok{linewidth =} \FloatTok{0.5}\NormalTok{) }\SpecialCharTok{+}
  \FunctionTok{geom\_vline}\NormalTok{(}\AttributeTok{data =}\NormalTok{ truth\_lines, }\FunctionTok{aes}\NormalTok{(}\AttributeTok{xintercept =}\NormalTok{ value),}
             \AttributeTok{color =} \StringTok{"red"}\NormalTok{, }\AttributeTok{linetype =} \StringTok{"dashed"}\NormalTok{) }\SpecialCharTok{+}
  \FunctionTok{facet\_grid}\NormalTok{(scenario }\SpecialCharTok{\textasciitilde{}}\NormalTok{ parameter, }\AttributeTok{scales =} \StringTok{"free"}\NormalTok{) }\SpecialCharTok{+}
  \FunctionTok{scale\_fill\_manual}\NormalTok{(}\AttributeTok{values  =} \FunctionTok{c}\NormalTok{(}\AttributeTok{prior =} \StringTok{"grey70"}\NormalTok{, }\AttributeTok{posterior =} \StringTok{"\#009E73"}\NormalTok{)) }\SpecialCharTok{+}
  \FunctionTok{scale\_color\_manual}\NormalTok{(}\AttributeTok{values =} \FunctionTok{c}\NormalTok{(}\AttributeTok{prior =} \StringTok{"grey50"}\NormalTok{, }\AttributeTok{posterior =} \StringTok{"\#006D4E"}\NormalTok{)) }\SpecialCharTok{+}
  \FunctionTok{labs}\NormalTok{(}
    \AttributeTok{title    =} \StringTok{"Prior → Posterior Update: Prototype Model Across Scenarios"}\NormalTok{,}
    \AttributeTok{subtitle =} \StringTok{"Red dashed line = true generating value used for the single{-}agent fit"}\NormalTok{,}
    \AttributeTok{x =} \ConstantTok{NULL}\NormalTok{, }\AttributeTok{y =} \StringTok{"Density"}
\NormalTok{  )}
\end{Highlighting}
\end{Shaded}

\begin{center}
\includegraphics[width=0.8\linewidth,height=\textheight,keepaspectratio]{12-CategorizationModelsPrototypes_files/figure-pdf/ch13_scenario_prior_posterior-1.pdf}
\end{center}

The density of \texttt{log\_q} narrows visibly across scenarios: from
prior-like on static data (information-starved), to sharply peaked
around the truth on drift data (assumption-matched).

\subsection{One-Shot Parameter Recovery Across
Scenarios}\label{one-shot-parameter-recovery-across-scenarios-1}

We sweep a 3 × 3 grid of \((r, q)\) truths with 5 iterations per cell
per scenario (135 fits) and plot recovery on the log scale.

\begin{Shaded}
\begin{Highlighting}[]
\NormalTok{scenario\_recovery\_file }\OtherTok{\textless{}{-}} \FunctionTok{here}\NormalTok{(}\StringTok{"simdata"}\NormalTok{, }\StringTok{"ch13\_scenario\_recovery\_results.rds"}\NormalTok{)}

\NormalTok{run\_proto\_scenario\_recovery\_iter }\OtherTok{\textless{}{-}} \ControlFlowTok{function}\NormalTok{(scenario, r\_true, q\_true, iteration) \{}
\NormalTok{  subject\_seed }\OtherTok{\textless{}{-}}\NormalTok{ iteration }\SpecialCharTok{*} \DecValTok{10000} \SpecialCharTok{+}
                   \FunctionTok{round}\NormalTok{(r\_true }\SpecialCharTok{*} \DecValTok{100}\NormalTok{) }\SpecialCharTok{+}
                   \FunctionTok{round}\NormalTok{(q\_true }\SpecialCharTok{*} \DecValTok{1000}\NormalTok{) }\SpecialCharTok{+}
                   \ControlFlowTok{switch}\NormalTok{(scenario, }\AttributeTok{static =} \DecValTok{1}\NormalTok{, }\AttributeTok{contingent\_shift =} \DecValTok{2}\NormalTok{, }\AttributeTok{drift =} \DecValTok{3}\NormalTok{)}

\NormalTok{  agent\_df }\OtherTok{\textless{}{-}} \FunctionTok{simulate\_prototype\_agent\_scenario}\NormalTok{(}
    \AttributeTok{agent\_id     =}\NormalTok{ iteration,}
    \AttributeTok{scenario     =}\NormalTok{ scenario,}
    \AttributeTok{r\_value      =}\NormalTok{ r\_true,}
    \AttributeTok{q\_value      =}\NormalTok{ q\_true,}
    \AttributeTok{subject\_seed =}\NormalTok{ subject\_seed}
\NormalTok{  )}
\NormalTok{  dat }\OtherTok{\textless{}{-}} \FunctionTok{stan\_data\_from\_scenario\_agent}\NormalTok{(agent\_df)}
\NormalTok{  pf  }\OtherTok{\textless{}{-}} \FunctionTok{tryCatch}\NormalTok{(}
\NormalTok{    mod\_prototype\_single}\SpecialCharTok{$}\FunctionTok{pathfinder}\NormalTok{(}\AttributeTok{data =}\NormalTok{ dat, }\AttributeTok{num\_paths =} \DecValTok{4}\NormalTok{,}
                                    \AttributeTok{refresh =} \DecValTok{0}\NormalTok{, }\AttributeTok{seed =}\NormalTok{ subject\_seed),}
    \AttributeTok{error =} \ControlFlowTok{function}\NormalTok{(e) }\ConstantTok{NULL}
\NormalTok{  )}
\NormalTok{  fit }\OtherTok{\textless{}{-}} \FunctionTok{tryCatch}\NormalTok{(}
\NormalTok{    mod\_prototype\_single}\SpecialCharTok{$}\FunctionTok{sample}\NormalTok{(}
      \AttributeTok{data            =}\NormalTok{ dat,}
      \AttributeTok{init            =} \ControlFlowTok{if}\NormalTok{ (}\SpecialCharTok{!}\FunctionTok{is.null}\NormalTok{(pf)) pf }\ControlFlowTok{else} \FloatTok{0.5}\NormalTok{,}
      \AttributeTok{seed            =}\NormalTok{ subject\_seed,}
      \AttributeTok{chains          =} \DecValTok{2}\NormalTok{,}
      \AttributeTok{parallel\_chains =} \DecValTok{2}\NormalTok{,}
      \AttributeTok{iter\_warmup     =} \DecValTok{800}\NormalTok{,}
      \AttributeTok{iter\_sampling   =} \DecValTok{1000}\NormalTok{,}
      \AttributeTok{refresh         =} \DecValTok{0}\NormalTok{,}
      \AttributeTok{adapt\_delta     =} \FloatTok{0.9}
\NormalTok{    ),}
    \AttributeTok{error =} \ControlFlowTok{function}\NormalTok{(e) \{ }\FunctionTok{message}\NormalTok{(}\StringTok{"Fit error: "}\NormalTok{, e}\SpecialCharTok{$}\NormalTok{message); }\ConstantTok{NULL}\NormalTok{ \}}
\NormalTok{  )}
  \ControlFlowTok{if}\NormalTok{ (}\FunctionTok{is.null}\NormalTok{(fit)) }\FunctionTok{return}\NormalTok{(}\ConstantTok{NULL}\NormalTok{)}

\NormalTok{  s }\OtherTok{\textless{}{-}}\NormalTok{ fit}\SpecialCharTok{$}\FunctionTok{summary}\NormalTok{(}\AttributeTok{variables =} \FunctionTok{c}\NormalTok{(}\StringTok{"log\_r"}\NormalTok{, }\StringTok{"log\_q"}\NormalTok{))}
\NormalTok{  diag }\OtherTok{\textless{}{-}}\NormalTok{ fit}\SpecialCharTok{$}\FunctionTok{diagnostic\_summary}\NormalTok{(}\AttributeTok{quiet =} \ConstantTok{TRUE}\NormalTok{)}
  \FunctionTok{tibble}\NormalTok{(}
    \AttributeTok{scenario       =}\NormalTok{ scenario,}
    \AttributeTok{r\_true         =}\NormalTok{ r\_true,}
    \AttributeTok{q\_true         =}\NormalTok{ q\_true,}
    \AttributeTok{iteration      =}\NormalTok{ iteration,}
    \AttributeTok{log\_r\_mean     =}\NormalTok{ s}\SpecialCharTok{$}\NormalTok{mean[s}\SpecialCharTok{$}\NormalTok{variable }\SpecialCharTok{==} \StringTok{"log\_r"}\NormalTok{],}
    \AttributeTok{log\_q\_mean     =}\NormalTok{ s}\SpecialCharTok{$}\NormalTok{mean[s}\SpecialCharTok{$}\NormalTok{variable }\SpecialCharTok{==} \StringTok{"log\_q"}\NormalTok{],}
    \AttributeTok{log\_r\_q5       =}\NormalTok{ s}\SpecialCharTok{$}\NormalTok{q5[s}\SpecialCharTok{$}\NormalTok{variable }\SpecialCharTok{==} \StringTok{"log\_r"}\NormalTok{],}
    \AttributeTok{log\_r\_q95      =}\NormalTok{ s}\SpecialCharTok{$}\NormalTok{q95[s}\SpecialCharTok{$}\NormalTok{variable }\SpecialCharTok{==} \StringTok{"log\_r"}\NormalTok{],}
    \AttributeTok{log\_q\_q5       =}\NormalTok{ s}\SpecialCharTok{$}\NormalTok{q5[s}\SpecialCharTok{$}\NormalTok{variable }\SpecialCharTok{==} \StringTok{"log\_q"}\NormalTok{],}
    \AttributeTok{log\_q\_q95      =}\NormalTok{ s}\SpecialCharTok{$}\NormalTok{q95[s}\SpecialCharTok{$}\NormalTok{variable }\SpecialCharTok{==} \StringTok{"log\_q"}\NormalTok{],}
    \AttributeTok{had\_divergence =} \FunctionTok{any}\NormalTok{(diag}\SpecialCharTok{$}\NormalTok{num\_divergent }\SpecialCharTok{\textgreater{}} \DecValTok{0}\NormalTok{)}
\NormalTok{  )}
\NormalTok{\}}

\ControlFlowTok{if}\NormalTok{ (regenerate\_simulations }\SpecialCharTok{||} \SpecialCharTok{!}\FunctionTok{file.exists}\NormalTok{(scenario\_recovery\_file)) \{}
\NormalTok{  recovery\_plan }\OtherTok{\textless{}{-}} \FunctionTok{expand\_grid}\NormalTok{(}
    \AttributeTok{scenario  =} \FunctionTok{c}\NormalTok{(}\StringTok{"static"}\NormalTok{, }\StringTok{"contingent\_shift"}\NormalTok{, }\StringTok{"drift"}\NormalTok{),}
    \AttributeTok{r\_true    =} \FunctionTok{c}\NormalTok{(}\FloatTok{0.5}\NormalTok{, }\FloatTok{1.5}\NormalTok{, }\FloatTok{3.0}\NormalTok{),}
    \AttributeTok{q\_true    =} \FunctionTok{c}\NormalTok{(}\FloatTok{0.02}\NormalTok{, }\FloatTok{0.15}\NormalTok{, }\FloatTok{0.6}\NormalTok{),}
    \AttributeTok{iteration =} \DecValTok{1}\SpecialCharTok{:}\DecValTok{5}
\NormalTok{  )}

  \FunctionTok{plan}\NormalTok{(multisession, }\AttributeTok{workers =} \FunctionTok{max}\NormalTok{(}\DecValTok{1}\NormalTok{, }\FunctionTok{availableCores}\NormalTok{() }\SpecialCharTok{{-}} \DecValTok{1}\NormalTok{))}
\NormalTok{  scenario\_recovery }\OtherTok{\textless{}{-}} \FunctionTok{future\_pmap\_dfr}\NormalTok{(}
    \FunctionTok{list}\NormalTok{(}\AttributeTok{scenario  =}\NormalTok{ recovery\_plan}\SpecialCharTok{$}\NormalTok{scenario,}
         \AttributeTok{r\_true    =}\NormalTok{ recovery\_plan}\SpecialCharTok{$}\NormalTok{r\_true,}
         \AttributeTok{q\_true    =}\NormalTok{ recovery\_plan}\SpecialCharTok{$}\NormalTok{q\_true,}
         \AttributeTok{iteration =}\NormalTok{ recovery\_plan}\SpecialCharTok{$}\NormalTok{iteration),}
\NormalTok{    run\_proto\_scenario\_recovery\_iter,}
    \AttributeTok{.options =} \FunctionTok{furrr\_options}\NormalTok{(}\AttributeTok{seed =} \ConstantTok{TRUE}\NormalTok{, }\AttributeTok{scheduling =} \DecValTok{1}\NormalTok{)}
\NormalTok{  )}
  \FunctionTok{plan}\NormalTok{(sequential)}
  \FunctionTok{saveRDS}\NormalTok{(scenario\_recovery, scenario\_recovery\_file)}
  \FunctionTok{cat}\NormalTok{(}\StringTok{"Scenario recovery sweep complete.}\SpecialCharTok{\textbackslash{}n}\StringTok{"}\NormalTok{)}
\NormalTok{\} }\ControlFlowTok{else}\NormalTok{ \{}
\NormalTok{  scenario\_recovery }\OtherTok{\textless{}{-}} \FunctionTok{readRDS}\NormalTok{(scenario\_recovery\_file)}
  \FunctionTok{cat}\NormalTok{(}\StringTok{"Loaded existing scenario recovery sweep.}\SpecialCharTok{\textbackslash{}n}\StringTok{"}\NormalTok{)}
\NormalTok{\}}
\end{Highlighting}
\end{Shaded}

\begin{verbatim}
Loaded existing scenario recovery sweep.
\end{verbatim}

\begin{Shaded}
\begin{Highlighting}[]
\NormalTok{scenario\_recovery\_long }\OtherTok{\textless{}{-}}\NormalTok{ scenario\_recovery }\SpecialCharTok{|\textgreater{}}
\NormalTok{  dplyr}\SpecialCharTok{::}\FunctionTok{select}\NormalTok{(scenario, r\_true, q\_true, iteration, had\_divergence,}
\NormalTok{                log\_r\_mean, log\_q\_mean) }\SpecialCharTok{|\textgreater{}}
  \FunctionTok{pivot\_longer}\NormalTok{(}\AttributeTok{cols =} \FunctionTok{c}\NormalTok{(log\_r\_mean, log\_q\_mean),}
               \AttributeTok{names\_to =} \StringTok{"variable"}\NormalTok{, }\AttributeTok{values\_to =} \StringTok{"estimated"}\NormalTok{) }\SpecialCharTok{|\textgreater{}}
  \FunctionTok{mutate}\NormalTok{(}
    \AttributeTok{parameter =} \FunctionTok{if\_else}\NormalTok{(variable }\SpecialCharTok{==} \StringTok{"log\_r\_mean"}\NormalTok{, }\StringTok{"log\_r"}\NormalTok{, }\StringTok{"log\_q"}\NormalTok{),}
    \AttributeTok{truth     =} \FunctionTok{if\_else}\NormalTok{(parameter }\SpecialCharTok{==} \StringTok{"log\_r"}\NormalTok{, }\FunctionTok{log}\NormalTok{(r\_true), }\FunctionTok{log}\NormalTok{(q\_true))}
\NormalTok{  )}

\FunctionTok{ggplot}\NormalTok{(scenario\_recovery\_long, }\FunctionTok{aes}\NormalTok{(}\AttributeTok{x =}\NormalTok{ truth, }\AttributeTok{y =}\NormalTok{ estimated,}
                                    \AttributeTok{color =}\NormalTok{ had\_divergence)) }\SpecialCharTok{+}
  \FunctionTok{geom\_abline}\NormalTok{(}\AttributeTok{intercept =} \DecValTok{0}\NormalTok{, }\AttributeTok{slope =} \DecValTok{1}\NormalTok{, }\AttributeTok{linetype =} \StringTok{"dashed"}\NormalTok{, }\AttributeTok{color =} \StringTok{"grey50"}\NormalTok{) }\SpecialCharTok{+}
  \FunctionTok{geom\_point}\NormalTok{(}\AttributeTok{alpha =} \FloatTok{0.7}\NormalTok{, }\AttributeTok{size =} \DecValTok{2}\NormalTok{) }\SpecialCharTok{+}
  \FunctionTok{facet\_grid}\NormalTok{(parameter }\SpecialCharTok{\textasciitilde{}}\NormalTok{ scenario, }\AttributeTok{scales =} \StringTok{"free"}\NormalTok{) }\SpecialCharTok{+}
  \FunctionTok{scale\_color\_manual}\NormalTok{(}\AttributeTok{values =} \FunctionTok{c}\NormalTok{(}\StringTok{"FALSE"} \OtherTok{=} \StringTok{"\#0072B2"}\NormalTok{, }\StringTok{"TRUE"} \OtherTok{=} \StringTok{"\#D55E00"}\NormalTok{),}
                     \AttributeTok{name =} \StringTok{"Had divergence"}\NormalTok{) }\SpecialCharTok{+}
  \FunctionTok{labs}\NormalTok{(}
    \AttributeTok{title    =} \StringTok{"One{-}Shot Parameter Recovery for the Kalman Prototype Model"}\NormalTok{,}
    \AttributeTok{subtitle =} \StringTok{"Posterior mean vs true generating value; 5 iterations per (r, q) cell, 3 scenarios"}\NormalTok{,}
    \AttributeTok{x =} \StringTok{"True value (log scale)"}\NormalTok{, }\AttributeTok{y =} \StringTok{"Estimated value (posterior mean)"}
\NormalTok{  ) }\SpecialCharTok{+}
  \FunctionTok{theme}\NormalTok{(}\AttributeTok{legend.position =} \StringTok{"bottom"}\NormalTok{)}
\end{Highlighting}
\end{Shaded}

\begin{center}
\includegraphics[width=0.8\linewidth,height=\textheight,keepaspectratio]{12-CategorizationModelsPrototypes_files/figure-pdf/ch13_scenario_recovery_sweep-1.pdf}
\end{center}

\textbf{Interpretation.} The \texttt{log\_r} row is near-diagonal across
scenarios --- \(r\) is identified by within-category decision sharpness,
which every scenario exercises. The \texttt{log\_q} row is the
diagnostic: diffuse scatter on static data (the parameter is
under-constrained), tighter on drift, and comparably tight on
contingent-shift. The contingent-shift result is noteworthy --- the
data-generating process has abrupt jumps that the filter's smooth-drift
prior cannot exactly reproduce, yet recovery still works because
post-shift recovery speed loads specifically on \(q\).

\subsection{Simulation-Based Calibration Checks Across
Scenarios}\label{simulation-based-calibration-checks-across-scenarios-1}

SBC is the formal calibration test. Each scenario gets its own generator
drawing \((\log r, \log q)\) from their priors.

\begin{Shaded}
\begin{Highlighting}[]
\NormalTok{make\_proto\_sbc\_generator }\OtherTok{\textless{}{-}} \ControlFlowTok{function}\NormalTok{(scenario) \{}
  \FunctionTok{SBC\_generator\_function}\NormalTok{(}
    \ControlFlowTok{function}\NormalTok{() \{}
\NormalTok{      log\_r\_true }\OtherTok{\textless{}{-}} \FunctionTok{rnorm}\NormalTok{(}\DecValTok{1}\NormalTok{, LOG\_R\_PRIOR\_MEAN, LOG\_R\_PRIOR\_SD)}
\NormalTok{      log\_q\_true }\OtherTok{\textless{}{-}} \FunctionTok{rnorm}\NormalTok{(}\DecValTok{1}\NormalTok{, LOG\_Q\_PRIOR\_MEAN, LOG\_Q\_PRIOR\_SD)}
\NormalTok{      r\_true     }\OtherTok{\textless{}{-}} \FunctionTok{exp}\NormalTok{(log\_r\_true)}
\NormalTok{      q\_true     }\OtherTok{\textless{}{-}} \FunctionTok{exp}\NormalTok{(log\_q\_true)}

\NormalTok{      subject\_seed }\OtherTok{\textless{}{-}} \FunctionTok{sample.int}\NormalTok{(}\FloatTok{1e8}\NormalTok{, }\DecValTok{1}\NormalTok{)}
\NormalTok{      agent\_df }\OtherTok{\textless{}{-}} \FunctionTok{simulate\_prototype\_agent\_scenario}\NormalTok{(}
        \AttributeTok{agent\_id     =} \DecValTok{1}\NormalTok{,}
        \AttributeTok{scenario     =}\NormalTok{ scenario,}
        \AttributeTok{r\_value      =}\NormalTok{ r\_true,}
        \AttributeTok{q\_value      =}\NormalTok{ q\_true,}
        \AttributeTok{subject\_seed =}\NormalTok{ subject\_seed}
\NormalTok{      )}

      \FunctionTok{list}\NormalTok{(}
        \AttributeTok{variables =} \FunctionTok{list}\NormalTok{(}\AttributeTok{log\_r =}\NormalTok{ log\_r\_true, }\AttributeTok{log\_q =}\NormalTok{ log\_q\_true),}
        \AttributeTok{generated =} \FunctionTok{stan\_data\_from\_scenario\_agent}\NormalTok{(agent\_df)}
\NormalTok{      )}
\NormalTok{    \}}
\NormalTok{  )}
\NormalTok{\}}

\NormalTok{run\_proto\_scenario\_sbc }\OtherTok{\textless{}{-}} \ControlFlowTok{function}\NormalTok{(scenario, }\AttributeTok{n\_sims =} \DecValTok{200}\NormalTok{) \{}
\NormalTok{  cache\_file }\OtherTok{\textless{}{-}} \FunctionTok{here}\NormalTok{(}\StringTok{"simdata"}\NormalTok{, }\FunctionTok{paste0}\NormalTok{(}\StringTok{"ch13\_proto\_sbc\_"}\NormalTok{, scenario, }\StringTok{".rds"}\NormalTok{))}
  \ControlFlowTok{if}\NormalTok{ (}\SpecialCharTok{!}\NormalTok{regenerate\_simulations }\SpecialCharTok{\&\&} \FunctionTok{file.exists}\NormalTok{(cache\_file)) \{}
    \FunctionTok{cat}\NormalTok{(}\StringTok{"Loaded cached prototype SBC for scenario:"}\NormalTok{, scenario, }\StringTok{"}\SpecialCharTok{\textbackslash{}n}\StringTok{"}\NormalTok{)}
    \FunctionTok{return}\NormalTok{(}\FunctionTok{readRDS}\NormalTok{(cache\_file))}
\NormalTok{  \}}
\NormalTok{  gen }\OtherTok{\textless{}{-}} \FunctionTok{make\_proto\_sbc\_generator}\NormalTok{(scenario)}
\NormalTok{  backend }\OtherTok{\textless{}{-}} \FunctionTok{SBC\_backend\_cmdstan\_sample}\NormalTok{(}
\NormalTok{    mod\_prototype\_single,}
    \AttributeTok{iter\_warmup   =} \DecValTok{800}\NormalTok{,}
    \AttributeTok{iter\_sampling =} \DecValTok{1000}\NormalTok{,}
    \AttributeTok{chains        =} \DecValTok{2}\NormalTok{,}
    \AttributeTok{adapt\_delta   =} \FloatTok{0.9}\NormalTok{,}
    \AttributeTok{refresh       =} \DecValTok{0}
\NormalTok{  )}
\NormalTok{  datasets }\OtherTok{\textless{}{-}} \FunctionTok{generate\_datasets}\NormalTok{(gen, }\AttributeTok{n\_sims =}\NormalTok{ n\_sims)}
\NormalTok{  res }\OtherTok{\textless{}{-}} \FunctionTok{compute\_SBC}\NormalTok{(}
\NormalTok{    datasets,}
\NormalTok{    backend,}
    \AttributeTok{cache\_mode     =} \StringTok{"results"}\NormalTok{,}
    \AttributeTok{cache\_location =} \FunctionTok{here}\NormalTok{(}\StringTok{"simdata"}\NormalTok{, }\FunctionTok{paste0}\NormalTok{(}\StringTok{"ch13\_proto\_sbc\_cache\_"}\NormalTok{, scenario)),}
    \AttributeTok{keep\_fits      =} \ConstantTok{FALSE}
\NormalTok{  )}
  \FunctionTok{saveRDS}\NormalTok{(res, cache\_file)}
\NormalTok{  res}
\NormalTok{\}}

\NormalTok{proto\_sbc\_static }\OtherTok{\textless{}{-}} \FunctionTok{run\_proto\_scenario\_sbc}\NormalTok{(}\StringTok{"static"}\NormalTok{,           }\AttributeTok{n\_sims =} \DecValTok{200}\NormalTok{)}
\end{Highlighting}
\end{Shaded}

\begin{verbatim}
Loaded cached prototype SBC for scenario: static 
\end{verbatim}

\begin{Shaded}
\begin{Highlighting}[]
\NormalTok{proto\_sbc\_shift  }\OtherTok{\textless{}{-}} \FunctionTok{run\_proto\_scenario\_sbc}\NormalTok{(}\StringTok{"contingent\_shift"}\NormalTok{, }\AttributeTok{n\_sims =} \DecValTok{200}\NormalTok{)}
\end{Highlighting}
\end{Shaded}

\begin{verbatim}
Loaded cached prototype SBC for scenario: contingent_shift 
\end{verbatim}

\begin{Shaded}
\begin{Highlighting}[]
\NormalTok{proto\_sbc\_drift  }\OtherTok{\textless{}{-}} \FunctionTok{run\_proto\_scenario\_sbc}\NormalTok{(}\StringTok{"drift"}\NormalTok{,            }\AttributeTok{n\_sims =} \DecValTok{200}\NormalTok{)}
\end{Highlighting}
\end{Shaded}

\begin{verbatim}
Loaded cached prototype SBC for scenario: drift 
\end{verbatim}

\begin{Shaded}
\begin{Highlighting}[]
\FunctionTok{print}\NormalTok{(}\FunctionTok{plot\_ecdf\_diff}\NormalTok{(proto\_sbc\_static) }\SpecialCharTok{+} \FunctionTok{ggtitle}\NormalTok{(}\StringTok{"SBC: static"}\NormalTok{))}
\end{Highlighting}
\end{Shaded}

\begin{center}
\includegraphics[width=0.8\linewidth,height=\textheight,keepaspectratio]{12-CategorizationModelsPrototypes_files/figure-pdf/ch13_scenario_sbc_ecdf-1.pdf}
\end{center}

\begin{Shaded}
\begin{Highlighting}[]
\FunctionTok{print}\NormalTok{(}\FunctionTok{plot\_ecdf\_diff}\NormalTok{(proto\_sbc\_shift)  }\SpecialCharTok{+} \FunctionTok{ggtitle}\NormalTok{(}\StringTok{"SBC: performance{-}contingent shifts"}\NormalTok{))}
\end{Highlighting}
\end{Shaded}

\begin{center}
\includegraphics[width=0.8\linewidth,height=\textheight,keepaspectratio]{12-CategorizationModelsPrototypes_files/figure-pdf/ch13_scenario_sbc_ecdf-2.pdf}
\end{center}

\begin{Shaded}
\begin{Highlighting}[]
\FunctionTok{print}\NormalTok{(}\FunctionTok{plot\_ecdf\_diff}\NormalTok{(proto\_sbc\_drift)  }\SpecialCharTok{+} \FunctionTok{ggtitle}\NormalTok{(}\StringTok{"SBC: smooth drift"}\NormalTok{))}
\end{Highlighting}
\end{Shaded}

\begin{center}
\includegraphics[width=0.8\linewidth,height=\textheight,keepaspectratio]{12-CategorizationModelsPrototypes_files/figure-pdf/ch13_scenario_sbc_ecdf-3.pdf}
\end{center}

\begin{Shaded}
\begin{Highlighting}[]
\FunctionTok{print}\NormalTok{(}\FunctionTok{plot\_rank\_hist}\NormalTok{(proto\_sbc\_static) }\SpecialCharTok{+} \FunctionTok{ggtitle}\NormalTok{(}\StringTok{"Rank histograms — static"}\NormalTok{))}
\end{Highlighting}
\end{Shaded}

\begin{center}
\includegraphics[width=0.8\linewidth,height=\textheight,keepaspectratio]{12-CategorizationModelsPrototypes_files/figure-pdf/ch13_scenario_sbc_ecdf-4.pdf}
\end{center}

\begin{Shaded}
\begin{Highlighting}[]
\FunctionTok{print}\NormalTok{(}\FunctionTok{plot\_rank\_hist}\NormalTok{(proto\_sbc\_shift)  }\SpecialCharTok{+} \FunctionTok{ggtitle}\NormalTok{(}\StringTok{"Rank histograms — contingent shifts"}\NormalTok{))}
\end{Highlighting}
\end{Shaded}

\begin{center}
\includegraphics[width=0.8\linewidth,height=\textheight,keepaspectratio]{12-CategorizationModelsPrototypes_files/figure-pdf/ch13_scenario_sbc_ecdf-5.pdf}
\end{center}

\begin{Shaded}
\begin{Highlighting}[]
\FunctionTok{print}\NormalTok{(}\FunctionTok{plot\_rank\_hist}\NormalTok{(proto\_sbc\_drift)  }\SpecialCharTok{+} \FunctionTok{ggtitle}\NormalTok{(}\StringTok{"Rank histograms — drift"}\NormalTok{))}
\end{Highlighting}
\end{Shaded}

\begin{center}
\includegraphics[width=0.8\linewidth,height=\textheight,keepaspectratio]{12-CategorizationModelsPrototypes_files/figure-pdf/ch13_scenario_sbc_ecdf-6.pdf}
\end{center}

\begin{Shaded}
\begin{Highlighting}[]
\NormalTok{combine\_proto\_sbc\_stats }\OtherTok{\textless{}{-}} \ControlFlowTok{function}\NormalTok{(sbc\_obj, scenario) \{}
\NormalTok{  sbc\_obj}\SpecialCharTok{$}\NormalTok{stats }\SpecialCharTok{|\textgreater{}} \FunctionTok{mutate}\NormalTok{(}\AttributeTok{scenario =}\NormalTok{ scenario)}
\NormalTok{\}}

\NormalTok{proto\_sbc\_stats }\OtherTok{\textless{}{-}} \FunctionTok{bind\_rows}\NormalTok{(}
  \FunctionTok{combine\_proto\_sbc\_stats}\NormalTok{(proto\_sbc\_static, }\StringTok{"static"}\NormalTok{),}
  \FunctionTok{combine\_proto\_sbc\_stats}\NormalTok{(proto\_sbc\_shift,  }\StringTok{"contingent\_shift"}\NormalTok{),}
  \FunctionTok{combine\_proto\_sbc\_stats}\NormalTok{(proto\_sbc\_drift,  }\StringTok{"drift"}\NormalTok{)}
\NormalTok{) }\SpecialCharTok{|\textgreater{}}
\NormalTok{  dplyr}\SpecialCharTok{::}\FunctionTok{filter}\NormalTok{(variable }\SpecialCharTok{\%in\%} \FunctionTok{c}\NormalTok{(}\StringTok{"log\_r"}\NormalTok{, }\StringTok{"log\_q"}\NormalTok{))}

\FunctionTok{ggplot}\NormalTok{(proto\_sbc\_stats, }\FunctionTok{aes}\NormalTok{(}\AttributeTok{x =}\NormalTok{ simulated\_value, }\AttributeTok{y =}\NormalTok{ mean)) }\SpecialCharTok{+}
  \FunctionTok{geom\_point}\NormalTok{(}\AttributeTok{alpha =} \FloatTok{0.35}\NormalTok{, }\AttributeTok{color =} \StringTok{"\#0072B2"}\NormalTok{) }\SpecialCharTok{+}
  \FunctionTok{geom\_abline}\NormalTok{(}\AttributeTok{intercept =} \DecValTok{0}\NormalTok{, }\AttributeTok{slope =} \DecValTok{1}\NormalTok{, }\AttributeTok{color =} \StringTok{"red"}\NormalTok{, }\AttributeTok{linetype =} \StringTok{"dashed"}\NormalTok{) }\SpecialCharTok{+}
  \FunctionTok{facet\_grid}\NormalTok{(variable }\SpecialCharTok{\textasciitilde{}}\NormalTok{ scenario, }\AttributeTok{scales =} \StringTok{"free"}\NormalTok{) }\SpecialCharTok{+}
  \FunctionTok{labs}\NormalTok{(}
    \AttributeTok{title    =} \StringTok{"SBC{-}Level Parameter Recovery: Kalman Prototype"}\NormalTok{,}
    \AttributeTok{subtitle =} \StringTok{"Posterior mean vs true generating value, pooled across SBC simulations"}\NormalTok{,}
    \AttributeTok{x =} \StringTok{"True value"}\NormalTok{, }\AttributeTok{y =} \StringTok{"Posterior mean"}
\NormalTok{  )}
\end{Highlighting}
\end{Shaded}

\begin{center}
\includegraphics[width=0.8\linewidth,height=\textheight,keepaspectratio]{12-CategorizationModelsPrototypes_files/figure-pdf/ch13_scenario_sbc_recovery_plots-1.pdf}
\end{center}

\begin{Shaded}
\begin{Highlighting}[]
\NormalTok{proto\_sbc\_stats }\SpecialCharTok{|\textgreater{}}
  \FunctionTok{group\_by}\NormalTok{(scenario, variable) }\SpecialCharTok{|\textgreater{}}
  \FunctionTok{summarise}\NormalTok{(}
    \AttributeTok{sbc\_mean\_bias =} \FunctionTok{mean}\NormalTok{(mean }\SpecialCharTok{{-}}\NormalTok{ simulated\_value, }\AttributeTok{na.rm =} \ConstantTok{TRUE}\NormalTok{),}
    \AttributeTok{sbc\_rmse      =} \FunctionTok{sqrt}\NormalTok{(}\FunctionTok{mean}\NormalTok{((mean }\SpecialCharTok{{-}}\NormalTok{ simulated\_value)}\SpecialCharTok{\^{}}\DecValTok{2}\NormalTok{, }\AttributeTok{na.rm =} \ConstantTok{TRUE}\NormalTok{)),}
    \AttributeTok{sbc\_corr      =} \FunctionTok{cor}\NormalTok{(mean, simulated\_value,}
                        \AttributeTok{use =} \StringTok{"complete.obs"}\NormalTok{, }\AttributeTok{method =} \StringTok{"pearson"}\NormalTok{),}
    \AttributeTok{n\_sims        =}\NormalTok{ dplyr}\SpecialCharTok{::}\FunctionTok{n}\NormalTok{(),}
    \AttributeTok{.groups       =} \StringTok{"drop"}
\NormalTok{  ) }\SpecialCharTok{|\textgreater{}}
  \FunctionTok{print}\NormalTok{(}\AttributeTok{n =} \ConstantTok{Inf}\NormalTok{)}
\end{Highlighting}
\end{Shaded}

\begin{verbatim}
# A tibble: 6 x 6
  scenario         variable sbc_mean_bias sbc_rmse sbc_corr n_sims
  <chr>            <chr>            <dbl>    <dbl>    <dbl>  <int>
1 contingent_shift log_q           0.0410    0.889    0.522    200
2 contingent_shift log_r          -0.0372    0.426    0.910    200
3 drift            log_q           0.0649    0.769    0.641    200
4 drift            log_r          -0.0276    0.453    0.893    200
5 static           log_q           0.118     0.828    0.529    200
6 static           log_r          -0.0262    0.411    0.917    200
\end{verbatim}

\textbf{Reading the three SBC results.} The static SBC for
\texttt{log\_q} shows the ECDF-difference curve drifting outside the
simultaneous band and the rank histogram humping --- posterior credible
intervals do not contain the truth at the nominal rate. The
contingent-shift SBC is noticeably better on \texttt{log\_q} but may
still show mild systematic deviation (a signature of the
smooth-drift-vs-abrupt-jump misspecification). The drift SBC should show
both parameters' curves squarely inside the band and rank histograms
close to uniform. This is the formal version of the identifiability
claim: calibration is achieved when task structure matches the model's
generative story.

\begin{center}\rule{0.5\linewidth}{0.5pt}\end{center}

\section{\texorpdfstring{Reparameterizing to Break the \(q\)/\(r\)
Ridge}{Reparameterizing to Break the q/r Ridge}}\label{reparameterizing-to-break-the-qr-ridge}

The static-scenario pair plot showed a heavy negatively-correlated ridge
on \texttt{(log\_r,\ log\_q)}. The ridge is \emph{geometric} --- a
property of the posterior shape, not of the information content. It
makes HMC sampling less efficient (ESS-per-draw falls as the posterior
elongates into a thin diagonal), makes SBC harder to interpret (the
failure mode could be either log\_r or log\_q but really it's their
difference), and obscures what the data actually say.

\subsection{The SNR Parameterisation}\label{the-snr-parameterisation}

Rewrite the two parameters as

\[\log r, \qquad \log \rho = \log q - \log r.\]

Here \(\rho = q/r\) is the drift-to-observation-noise ratio --- a
dimensionless scale that directly controls the Kalman filter's
steady-state gain. The original \texttt{log\_q} is now a \emph{derived}
quantity: \texttt{log\_q\ =\ log\_r\ +\ log\_rho}. If the data pin down
the \emph{sum} \(\log r + \log \rho\) (i.e.~\(\log q\)) cleanly but not
the \emph{difference} \(\log r - \log \rho\), the SNR parameterisation
represents this cleanly as a wide marginal on one of the two parameters
--- no diagonal ridge, no misleading pair plot.

We use Normal priors on both axes:

\begin{Shaded}
\begin{Highlighting}[]
\CommentTok{\# SNR parameterisation priors.}
\CommentTok{\# log\_r \textasciitilde{} Normal(0, 1)      — same as the original parameterisation.}
\CommentTok{\# log\_rho \textasciitilde{} Normal({-}2, 1)   — centres rho at 0.135 (drift small vs. r).}
\CommentTok{\# Marginal on log\_q = log\_r + log\_rho is Normal({-}2, sqrt(2)) ≈ Normal({-}2, 1.41).}
\CommentTok{\# Slightly wider than the original Normal({-}2, 1), but the increase is modest and}
\CommentTok{\# the shape of the implied prior predictive is comparable.}
\NormalTok{LOG\_RHO\_PRIOR\_MEAN }\OtherTok{\textless{}{-}} \SpecialCharTok{{-}}\FloatTok{2.0}
\NormalTok{LOG\_RHO\_PRIOR\_SD   }\OtherTok{\textless{}{-}} \FloatTok{1.0}
\end{Highlighting}
\end{Shaded}

\subsection{The SNR Stan Model}\label{the-snr-stan-model}

\begin{Shaded}
\begin{Highlighting}[]
\NormalTok{stan\_file\_proto\_snr }\OtherTok{\textless{}{-}} \FunctionTok{here}\NormalTok{(}\StringTok{"stan"}\NormalTok{, }\StringTok{"ch12\_prototype\_single\_snr.stan"}\NormalTok{)}
\NormalTok{mod\_prototype\_snr   }\OtherTok{\textless{}{-}} \FunctionTok{cmdstan\_model}\NormalTok{(stan\_file\_proto\_snr)}
\end{Highlighting}
\end{Shaded}

See \texttt{stan/ch12\_prototype\_single\_snr.stan}. The only
parameter-block change is \texttt{log\_rho} replacing \texttt{log\_q};
\texttt{log\_q} is computed in transformed parameters.

\subsection{Fitting Both Parameterisations to the Same Static
Data}\label{fitting-both-parameterisations-to-the-same-static-data}

\begin{Shaded}
\begin{Highlighting}[]
\NormalTok{snr\_static\_file }\OtherTok{\textless{}{-}} \FunctionTok{here}\NormalTok{(}\StringTok{"simmodels"}\NormalTok{, }\StringTok{"ch13\_proto\_snr\_static\_fit.rds"}\NormalTok{)}

\NormalTok{static\_agent }\OtherTok{\textless{}{-}}\NormalTok{ scenario\_single\_fit\_data}\SpecialCharTok{$}\NormalTok{static  }\CommentTok{\# same dataset as above}
\NormalTok{snr\_static\_data }\OtherTok{\textless{}{-}} \FunctionTok{list}\NormalTok{(}
  \AttributeTok{ntrials            =} \FunctionTok{nrow}\NormalTok{(static\_agent),}
  \AttributeTok{nfeatures          =} \DecValTok{2}\NormalTok{L,}
  \AttributeTok{cat\_one            =} \FunctionTok{as.integer}\NormalTok{(static\_agent}\SpecialCharTok{$}\NormalTok{observed\_feedback),}
  \AttributeTok{y                  =} \FunctionTok{as.integer}\NormalTok{(static\_agent}\SpecialCharTok{$}\NormalTok{sim\_response),}
  \AttributeTok{obs                =} \FunctionTok{as.matrix}\NormalTok{(static\_agent[, }\FunctionTok{c}\NormalTok{(}\StringTok{"height"}\NormalTok{, }\StringTok{"position"}\NormalTok{)]),}
  \AttributeTok{initial\_mu\_cat0    =} \FunctionTok{c}\NormalTok{(}\FloatTok{2.5}\NormalTok{, }\FloatTok{2.5}\NormalTok{),}
  \AttributeTok{initial\_mu\_cat1    =} \FunctionTok{c}\NormalTok{(}\FloatTok{2.5}\NormalTok{, }\FloatTok{2.5}\NormalTok{),}
  \AttributeTok{initial\_sigma\_diag =} \FloatTok{10.0}\NormalTok{,}
  \AttributeTok{prior\_logr\_mean    =}\NormalTok{ LOG\_R\_PRIOR\_MEAN,}
  \AttributeTok{prior\_logr\_sd      =}\NormalTok{ LOG\_R\_PRIOR\_SD,}
  \AttributeTok{prior\_logrho\_mean  =}\NormalTok{ LOG\_RHO\_PRIOR\_MEAN,}
  \AttributeTok{prior\_logrho\_sd    =}\NormalTok{ LOG\_RHO\_PRIOR\_SD}
\NormalTok{)}

\ControlFlowTok{if}\NormalTok{ (regenerate\_simulations }\SpecialCharTok{||} \SpecialCharTok{!}\FunctionTok{file.exists}\NormalTok{(snr\_static\_file)) \{}
\NormalTok{  pf\_snr }\OtherTok{\textless{}{-}} \FunctionTok{tryCatch}\NormalTok{(}
\NormalTok{    mod\_prototype\_snr}\SpecialCharTok{$}\FunctionTok{pathfinder}\NormalTok{(}\AttributeTok{data =}\NormalTok{ snr\_static\_data, }\AttributeTok{num\_paths =} \DecValTok{4}\NormalTok{,}
                                 \AttributeTok{refresh =} \DecValTok{0}\NormalTok{, }\AttributeTok{seed =} \DecValTok{777}\NormalTok{),}
    \AttributeTok{error =} \ControlFlowTok{function}\NormalTok{(e) }\ConstantTok{NULL}
\NormalTok{  )}
\NormalTok{  fit\_proto\_snr\_static }\OtherTok{\textless{}{-}}\NormalTok{ mod\_prototype\_snr}\SpecialCharTok{$}\FunctionTok{sample}\NormalTok{(}
    \AttributeTok{data            =}\NormalTok{ snr\_static\_data,}
    \AttributeTok{init            =} \ControlFlowTok{if}\NormalTok{ (}\SpecialCharTok{!}\FunctionTok{is.null}\NormalTok{(pf\_snr)) pf\_snr }\ControlFlowTok{else} \FloatTok{0.5}\NormalTok{,}
    \AttributeTok{seed            =} \DecValTok{777}\NormalTok{,}
    \AttributeTok{chains          =} \DecValTok{4}\NormalTok{,}
    \AttributeTok{parallel\_chains =} \FunctionTok{min}\NormalTok{(}\DecValTok{4}\NormalTok{, }\FunctionTok{availableCores}\NormalTok{()),}
    \AttributeTok{iter\_warmup     =} \DecValTok{1000}\NormalTok{,}
    \AttributeTok{iter\_sampling   =} \DecValTok{1500}\NormalTok{,}
    \AttributeTok{refresh         =} \DecValTok{0}\NormalTok{,}
    \AttributeTok{adapt\_delta     =} \FloatTok{0.9}
\NormalTok{  )}
\NormalTok{  fit\_proto\_snr\_static}\SpecialCharTok{$}\FunctionTok{save\_object}\NormalTok{(snr\_static\_file)}
\NormalTok{\} }\ControlFlowTok{else}\NormalTok{ \{}
\NormalTok{  fit\_proto\_snr\_static }\OtherTok{\textless{}{-}} \FunctionTok{readRDS}\NormalTok{(snr\_static\_file)}
\NormalTok{\}}

\NormalTok{orig\_static\_fit }\OtherTok{\textless{}{-}}\NormalTok{ scenario\_fits}\SpecialCharTok{$}\NormalTok{static}

\FunctionTok{cat}\NormalTok{(}\StringTok{"=== Original (log\_r, log\_q) parameterisation ===}\SpecialCharTok{\textbackslash{}n}\StringTok{"}\NormalTok{)}
\end{Highlighting}
\end{Shaded}

\begin{verbatim}
=== Original (log_r, log_q) parameterisation ===
\end{verbatim}

\begin{Shaded}
\begin{Highlighting}[]
\FunctionTok{print}\NormalTok{(orig\_static\_fit}\SpecialCharTok{$}\FunctionTok{summary}\NormalTok{(}\AttributeTok{variables =} \FunctionTok{c}\NormalTok{(}\StringTok{"log\_r"}\NormalTok{, }\StringTok{"log\_q"}\NormalTok{, }\StringTok{"r\_value"}\NormalTok{, }\StringTok{"q\_value"}\NormalTok{)))}
\end{Highlighting}
\end{Shaded}

\begin{verbatim}
# A tibble: 4 x 10
  variable   mean median    sd   mad      q5    q95  rhat ess_bulk ess_tail
  <chr>     <dbl>  <dbl> <dbl> <dbl>   <dbl>  <dbl> <dbl>    <dbl>    <dbl>
1 log_r     0.792  0.769 0.457 0.427  0.102   1.58   1.00    2578.    2361.
2 log_q    -2.08  -2.05  0.932 0.940 -3.71   -0.600  1.00    2824.    3184.
3 r_value   2.47   2.16  1.41  0.886  1.11    4.84   1.00    2578.    2361.
4 q_value   0.187  0.129 0.182 0.111  0.0245  0.549  1.00    2824.    3184.
\end{verbatim}

\begin{Shaded}
\begin{Highlighting}[]
\FunctionTok{cat}\NormalTok{(}\StringTok{"}\SpecialCharTok{\textbackslash{}n}\StringTok{=== SNR (log\_r, log\_rho) parameterisation ===}\SpecialCharTok{\textbackslash{}n}\StringTok{"}\NormalTok{)}
\end{Highlighting}
\end{Shaded}

\begin{verbatim}

=== SNR (log_r, log_rho) parameterisation ===
\end{verbatim}

\begin{Shaded}
\begin{Highlighting}[]
\FunctionTok{print}\NormalTok{(fit\_proto\_snr\_static}\SpecialCharTok{$}\FunctionTok{summary}\NormalTok{(}\AttributeTok{variables =} \FunctionTok{c}\NormalTok{(}\StringTok{"log\_r"}\NormalTok{, }\StringTok{"log\_rho"}\NormalTok{, }\StringTok{"log\_q"}\NormalTok{,}
                                                  \StringTok{"r\_value"}\NormalTok{, }\StringTok{"q\_value"}\NormalTok{)))}
\end{Highlighting}
\end{Shaded}

\begin{verbatim}
# A tibble: 5 x 10
  variable   mean median    sd   mad      q5    q95  rhat ess_bulk ess_tail
  <chr>     <dbl>  <dbl> <dbl> <dbl>   <dbl>  <dbl> <dbl>    <dbl>    <dbl>
1 log_r     0.702  0.667 0.412 0.375  0.0870  1.43   1.00    1921.    2268.
2 log_rho  -2.23  -2.22  0.900 0.913 -3.74   -0.768  1.00    2154.    2462.
3 log_q    -1.53  -1.52  0.829 0.840 -2.92   -0.196  1.00    2947.    3358.
4 r_value   2.22   1.95  1.25  0.718  1.09    4.19   1.00    1921.    2268.
5 q_value   0.302  0.220 0.281 0.174  0.0540  0.822  1.00    2947.    3358.
\end{verbatim}

\subsection{Geometry Comparison}\label{geometry-comparison}

\begin{Shaded}
\begin{Highlighting}[]
\NormalTok{p\_orig }\OtherTok{\textless{}{-}}\NormalTok{ bayesplot}\SpecialCharTok{::}\FunctionTok{mcmc\_pairs}\NormalTok{(}
\NormalTok{  orig\_static\_fit}\SpecialCharTok{$}\FunctionTok{draws}\NormalTok{(}\FunctionTok{c}\NormalTok{(}\StringTok{"log\_r"}\NormalTok{, }\StringTok{"log\_q"}\NormalTok{)),}
  \AttributeTok{diag\_fun =} \StringTok{"dens"}\NormalTok{, }\AttributeTok{off\_diag\_fun =} \StringTok{"hex"}\NormalTok{,}
  \AttributeTok{np       =} \FunctionTok{nuts\_params}\NormalTok{(orig\_static\_fit)}
\NormalTok{)}
\NormalTok{p\_snr  }\OtherTok{\textless{}{-}}\NormalTok{ bayesplot}\SpecialCharTok{::}\FunctionTok{mcmc\_pairs}\NormalTok{(}
\NormalTok{  fit\_proto\_snr\_static}\SpecialCharTok{$}\FunctionTok{draws}\NormalTok{(}\FunctionTok{c}\NormalTok{(}\StringTok{"log\_r"}\NormalTok{, }\StringTok{"log\_rho"}\NormalTok{)),}
  \AttributeTok{diag\_fun =} \StringTok{"dens"}\NormalTok{, }\AttributeTok{off\_diag\_fun =} \StringTok{"hex"}\NormalTok{,}
  \AttributeTok{np       =} \FunctionTok{nuts\_params}\NormalTok{(fit\_proto\_snr\_static)}
\NormalTok{)}

\FunctionTok{print}\NormalTok{(p\_orig)}
\end{Highlighting}
\end{Shaded}

\begin{center}
\includegraphics[width=0.8\linewidth,height=\textheight,keepaspectratio]{12-CategorizationModelsPrototypes_files/figure-pdf/ch13_snr_pair_compare-1.pdf}
\end{center}

\begin{Shaded}
\begin{Highlighting}[]
\FunctionTok{print}\NormalTok{(p\_snr)}
\end{Highlighting}
\end{Shaded}

\begin{center}
\includegraphics[width=0.8\linewidth,height=\textheight,keepaspectratio]{12-CategorizationModelsPrototypes_files/figure-pdf/ch13_snr_pair_compare-2.pdf}
\end{center}

\begin{Shaded}
\begin{Highlighting}[]
\CommentTok{\# Side{-}by{-}side ESS comparison}
\NormalTok{ess\_compare }\OtherTok{\textless{}{-}} \FunctionTok{bind\_rows}\NormalTok{(}
\NormalTok{  orig\_static\_fit}\SpecialCharTok{$}\FunctionTok{summary}\NormalTok{(}\FunctionTok{c}\NormalTok{(}\StringTok{"log\_r"}\NormalTok{, }\StringTok{"log\_q"}\NormalTok{),}
                          \StringTok{"ess\_bulk"}\NormalTok{, }\StringTok{"ess\_tail"}\NormalTok{) }\SpecialCharTok{|\textgreater{}}
    \FunctionTok{mutate}\NormalTok{(}\AttributeTok{parameterisation =} \StringTok{"original"}\NormalTok{),}
\NormalTok{  fit\_proto\_snr\_static}\SpecialCharTok{$}\FunctionTok{summary}\NormalTok{(}\FunctionTok{c}\NormalTok{(}\StringTok{"log\_r"}\NormalTok{, }\StringTok{"log\_rho"}\NormalTok{),}
                               \StringTok{"ess\_bulk"}\NormalTok{, }\StringTok{"ess\_tail"}\NormalTok{) }\SpecialCharTok{|\textgreater{}}
    \FunctionTok{mutate}\NormalTok{(}\AttributeTok{parameterisation =} \StringTok{"SNR"}\NormalTok{)}
\NormalTok{)}
\FunctionTok{print}\NormalTok{(ess\_compare)}
\end{Highlighting}
\end{Shaded}

\begin{verbatim}
# A tibble: 4 x 4
  variable ess_bulk ess_tail parameterisation
  <chr>       <dbl>    <dbl> <chr>           
1 log_r       2578.    2361. original        
2 log_q       2824.    3184. original        
3 log_r       1921.    2268. SNR             
4 log_rho     2154.    2462. SNR             
\end{verbatim}

\textbf{The expected pattern.} The original pair plot shows the heavy
negative ridge; the SNR pair plot shows a wide-but-roughly-circular
cloud on \texttt{(log\_r,\ log\_rho)}. ESS-per-iteration should improve
for both parameters under SNR, because HMC no longer has to traverse a
long thin diagonal. The posterior over \texttt{log\_q} (a derived
quantity under SNR) is \emph{identical in substance} --- the two models
imply the same likelihood --- but the Stan sampler reaches it more
efficiently.

\subsection{SBC for the SNR Model on Drift
Data}\label{sbc-for-the-snr-model-on-drift-data}

We rerun SBC on the drift scenario under the SNR parameterisation as the
ultimate test: \emph{assumption-matched data + decorrelated
parameterisation = cleanest possible calibration.}

\begin{Shaded}
\begin{Highlighting}[]
\NormalTok{snr\_sbc\_drift\_file }\OtherTok{\textless{}{-}} \FunctionTok{here}\NormalTok{(}\StringTok{"simdata"}\NormalTok{, }\StringTok{"ch13\_proto\_snr\_sbc\_drift.rds"}\NormalTok{)}

\NormalTok{make\_snr\_sbc\_generator }\OtherTok{\textless{}{-}} \ControlFlowTok{function}\NormalTok{(scenario) \{}
  \FunctionTok{SBC\_generator\_function}\NormalTok{(}
    \ControlFlowTok{function}\NormalTok{() \{}
\NormalTok{      log\_r\_true   }\OtherTok{\textless{}{-}} \FunctionTok{rnorm}\NormalTok{(}\DecValTok{1}\NormalTok{, LOG\_R\_PRIOR\_MEAN,   LOG\_R\_PRIOR\_SD)}
\NormalTok{      log\_rho\_true }\OtherTok{\textless{}{-}} \FunctionTok{rnorm}\NormalTok{(}\DecValTok{1}\NormalTok{, LOG\_RHO\_PRIOR\_MEAN, LOG\_RHO\_PRIOR\_SD)}
\NormalTok{      log\_q\_true   }\OtherTok{\textless{}{-}}\NormalTok{ log\_r\_true }\SpecialCharTok{+}\NormalTok{ log\_rho\_true}
\NormalTok{      r\_true       }\OtherTok{\textless{}{-}} \FunctionTok{exp}\NormalTok{(log\_r\_true)}
\NormalTok{      q\_true       }\OtherTok{\textless{}{-}} \FunctionTok{exp}\NormalTok{(log\_q\_true)}

\NormalTok{      agent\_df }\OtherTok{\textless{}{-}} \FunctionTok{simulate\_prototype\_agent\_scenario}\NormalTok{(}
        \AttributeTok{agent\_id     =} \DecValTok{1}\NormalTok{,}
        \AttributeTok{scenario     =}\NormalTok{ scenario,}
        \AttributeTok{r\_value      =}\NormalTok{ r\_true,}
        \AttributeTok{q\_value      =}\NormalTok{ q\_true,}
        \AttributeTok{subject\_seed =} \FunctionTok{sample.int}\NormalTok{(}\FloatTok{1e8}\NormalTok{, }\DecValTok{1}\NormalTok{)}
\NormalTok{      )}

      \FunctionTok{list}\NormalTok{(}
        \AttributeTok{variables =} \FunctionTok{list}\NormalTok{(}\AttributeTok{log\_r =}\NormalTok{ log\_r\_true,}
                         \AttributeTok{log\_rho =}\NormalTok{ log\_rho\_true,}
                         \AttributeTok{log\_q =}\NormalTok{ log\_q\_true),}
        \AttributeTok{generated =} \FunctionTok{list}\NormalTok{(}
          \AttributeTok{ntrials            =} \FunctionTok{nrow}\NormalTok{(agent\_df),}
          \AttributeTok{nfeatures          =} \DecValTok{2}\NormalTok{L,}
          \AttributeTok{cat\_one            =} \FunctionTok{as.integer}\NormalTok{(agent\_df}\SpecialCharTok{$}\NormalTok{observed\_feedback),}
          \AttributeTok{y                  =} \FunctionTok{as.integer}\NormalTok{(agent\_df}\SpecialCharTok{$}\NormalTok{sim\_response),}
          \AttributeTok{obs                =} \FunctionTok{as.matrix}\NormalTok{(agent\_df[, }\FunctionTok{c}\NormalTok{(}\StringTok{"height"}\NormalTok{, }\StringTok{"position"}\NormalTok{)]),}
          \AttributeTok{initial\_mu\_cat0    =} \FunctionTok{c}\NormalTok{(}\FloatTok{2.5}\NormalTok{, }\FloatTok{2.5}\NormalTok{),}
          \AttributeTok{initial\_mu\_cat1    =} \FunctionTok{c}\NormalTok{(}\FloatTok{2.5}\NormalTok{, }\FloatTok{2.5}\NormalTok{),}
          \AttributeTok{initial\_sigma\_diag =} \FloatTok{10.0}\NormalTok{,}
          \AttributeTok{prior\_logr\_mean    =}\NormalTok{ LOG\_R\_PRIOR\_MEAN,}
          \AttributeTok{prior\_logr\_sd      =}\NormalTok{ LOG\_R\_PRIOR\_SD,}
          \AttributeTok{prior\_logrho\_mean  =}\NormalTok{ LOG\_RHO\_PRIOR\_MEAN,}
          \AttributeTok{prior\_logrho\_sd    =}\NormalTok{ LOG\_RHO\_PRIOR\_SD}
\NormalTok{        )}
\NormalTok{      )}
\NormalTok{    \}}
\NormalTok{  )}
\NormalTok{\}}

\ControlFlowTok{if}\NormalTok{ (regenerate\_simulations }\SpecialCharTok{||} \SpecialCharTok{!}\FunctionTok{file.exists}\NormalTok{(snr\_sbc\_drift\_file)) \{}
\NormalTok{  gen }\OtherTok{\textless{}{-}} \FunctionTok{make\_snr\_sbc\_generator}\NormalTok{(}\StringTok{"drift"}\NormalTok{)}
\NormalTok{  backend }\OtherTok{\textless{}{-}} \FunctionTok{SBC\_backend\_cmdstan\_sample}\NormalTok{(}
\NormalTok{    mod\_prototype\_snr,}
    \AttributeTok{iter\_warmup   =} \DecValTok{800}\NormalTok{,}
    \AttributeTok{iter\_sampling =} \DecValTok{1000}\NormalTok{,}
    \AttributeTok{chains        =} \DecValTok{2}\NormalTok{,}
    \AttributeTok{adapt\_delta   =} \FloatTok{0.9}\NormalTok{,}
    \AttributeTok{refresh       =} \DecValTok{0}
\NormalTok{  )}
\NormalTok{  datasets }\OtherTok{\textless{}{-}} \FunctionTok{generate\_datasets}\NormalTok{(gen, }\AttributeTok{n\_sims =} \DecValTok{200}\NormalTok{)}
\NormalTok{  snr\_sbc\_drift }\OtherTok{\textless{}{-}} \FunctionTok{compute\_SBC}\NormalTok{(}
\NormalTok{    datasets, backend,}
    \AttributeTok{cache\_mode     =} \StringTok{"results"}\NormalTok{,}
    \AttributeTok{cache\_location =} \FunctionTok{here}\NormalTok{(}\StringTok{"simdata"}\NormalTok{, }\StringTok{"ch13\_proto\_snr\_sbc\_cache\_drift"}\NormalTok{),}
    \AttributeTok{keep\_fits      =} \ConstantTok{FALSE}
\NormalTok{  )}
  \FunctionTok{saveRDS}\NormalTok{(snr\_sbc\_drift, snr\_sbc\_drift\_file)}
\NormalTok{\} }\ControlFlowTok{else}\NormalTok{ \{}
\NormalTok{  snr\_sbc\_drift }\OtherTok{\textless{}{-}} \FunctionTok{readRDS}\NormalTok{(snr\_sbc\_drift\_file)}
\NormalTok{\}}

\FunctionTok{print}\NormalTok{(}\FunctionTok{plot\_ecdf\_diff}\NormalTok{(snr\_sbc\_drift) }\SpecialCharTok{+} \FunctionTok{ggtitle}\NormalTok{(}\StringTok{"SNR{-}parameterised SBC on drift data"}\NormalTok{))}
\end{Highlighting}
\end{Shaded}

\begin{center}
\includegraphics[width=0.8\linewidth,height=\textheight,keepaspectratio]{12-CategorizationModelsPrototypes_files/figure-pdf/ch13_snr_sbc_drift-1.pdf}
\end{center}

\begin{Shaded}
\begin{Highlighting}[]
\FunctionTok{print}\NormalTok{(}\FunctionTok{plot\_rank\_hist}\NormalTok{(snr\_sbc\_drift))}
\end{Highlighting}
\end{Shaded}

\begin{center}
\includegraphics[width=0.8\linewidth,height=\textheight,keepaspectratio]{12-CategorizationModelsPrototypes_files/figure-pdf/ch13_snr_sbc_drift-2.pdf}
\end{center}

Under this combination we expect all three curves (\texttt{log\_r},
\texttt{log\_rho}, and the derived \texttt{log\_q}) to stay inside the
simultaneous band with flat rank histograms. Reparameterisation alone
does not manufacture information: on static data, the SNR model will
still show the drift-related direction as a wide marginal (just on
\texttt{log\_rho} now, not on a diagonal ridge). It is the
\emph{combination} of assumption-matched task and decorrelated
parameterisation that closes the loop.

\begin{center}\rule{0.5\linewidth}{0.5pt}\end{center}

\section{\texorpdfstring{Discussion: When Does \(q\) Become
Identifiable?}{Discussion: When Does q Become Identifiable?}}\label{discussion-when-does-q-become-identifiable}

The six validation panels (three scenarios × two parameterisations) give
the cleanest summary of the identifiability argument.

\begin{itemize}
\tightlist
\item
  \textbf{Static + original.} The posterior sits on a heavy
  negatively-correlated ridge. SBC for \texttt{log\_q} fails; parameter
  recovery is diffuse. This is the problem the chapter opened with.
\item
  \textbf{Static + SNR.} The ridge disappears. The poorly-identified
  direction is now cleanly visible as a wide marginal on
  \texttt{log\_rho}. SBC still fails for the drift-related direction ---
  reparameterisation does not manufacture information.
\item
  \textbf{Contingent-shift + original.} Post-flip recovery dynamics load
  on \(q\) directly, so the ridge shrinks and SBC for \texttt{log\_q}
  passes (or nearly passes). Caveat: the filter's smooth-drift
  assumption is technically violated by abrupt jumps, so posterior
  predictive checks may reveal systematic underperformance right after
  each flip.
\item
  \textbf{Contingent-shift + SNR.} Similar to contingent-shift +
  original, but with tighter posterior marginals and better MCMC
  geometry.
\item
  \textbf{Smooth drift + original.} Assumption-matched generative story.
  The posterior is roughly circular and SBC passes for both parameters.
\item
  \textbf{Smooth drift + SNR.} The best-case scenario: decorrelated
  parameterisation + assumption-matched task. Both SBC ECDFs inside the
  band, both rank histograms uniform, MCMC geometry clean.
\end{itemize}

\textbf{Paralleling Chapter 11.} The GCM-with-decay analysis in Chapter
11 showed that the decay parameter \(\lambda\) becomes identifiable when
the task exercises forgetting --- whether through reversals or drift.
The mechanism there is general-purpose (``downweight old stuff''), so
reversals and drift work comparably well. The Kalman filter's \(q\) is
assumption-specific (it encodes a \emph{Gaussian} drift prior), so the
assumption-matched smooth-drift task gives the cleanest result;
reversals work but with a mild misspecification signature. Both chapters
make the same underlying point: \emph{identifiability is a joint
property of the model and the task, not of the model alone.}

\textbf{Implication for the multilevel analysis that follows.} The
multilevel prototype model below uses the static paradigm by default to
stay comparable with prior drafts of this chapter. Readers using this
model for research should consider re-running the multilevel analysis on
drift data when possible; the population-level \texttt{log\_q} posterior
that looked wide before is not a sign that multilevel inference failed
--- it was a sign that the task was not giving the model what it needed
to pin \(q\) down.

\begin{center}\rule{0.5\linewidth}{0.5pt}\end{center}

\section{Multilevel Prototype Model: Population Inference Across
Subjects}\label{multilevel-prototype-model-population-inference-across-subjects}

Real experiments involve many participants, each bringing their own
observation-noise tolerance and drift rate. Some learners trust
individual examples heavily (low \(r\)); others discount them (high
\(r\)). Some learners have more stable category representations (low
\(q\)); others show faster adaptation (high \(q\)). The
\textbf{multilevel prototype model} captures this variation by placing
population distributions over both \texttt{log\_r} and \texttt{log\_q}
and partially pooling information across subjects.

The structure is the minimal multilevel extension of the two-parameter
single-subject model:

\[\text{pop\_log\_r\_mean} \sim \mathcal{N}(0, 1), \quad \text{pop\_log\_r\_sd} \sim \text{Exponential}(2)\]
\[z_{r,j} \sim \mathcal{N}(0, 1), \quad \log r_j = \text{pop\_log\_r\_mean} + z_{r,j} \cdot \text{pop\_log\_r\_sd}\]
\[\text{pop\_log\_q\_mean} \sim \mathcal{N}(-2, 1), \quad \text{pop\_log\_q\_sd} \sim \text{Exponential}(2)\]
\[z_{q,j} \sim \mathcal{N}(0, 1), \quad \log q_j = \text{pop\_log\_q\_mean} + z_{q,j} \cdot \text{pop\_log\_q\_sd}\]

The non-centred parameterisation (NCP) for the per-subject offsets
avoids the funnel geometry that centred parameterisations produce when
the population SD is small. With two per-subject parameters, the SBC
rank histograms (below) show six curves: two population means, two
population SDs, and two individual-level curves
(\texttt{subj\_log\_r{[}1{]}}, \texttt{subj\_log\_q{[}1{]}}).

\subsection{Multilevel Forward
Simulation}\label{multilevel-forward-simulation}

\begin{Shaded}
\begin{Highlighting}[]
\CommentTok{\# Simulate a full multilevel dataset from specified population parameters.}
\NormalTok{simulate\_multilevel\_prototype }\OtherTok{\textless{}{-}} \ControlFlowTok{function}\NormalTok{(}
\NormalTok{    n\_subjects,}
\NormalTok{    stimulus\_info,}
\NormalTok{    n\_blocks,}
\NormalTok{    base\_seed,}
\NormalTok{    pop\_log\_r\_mean,}
\NormalTok{    pop\_log\_r\_sd,}
    \AttributeTok{pop\_log\_q\_mean =}\NormalTok{ LOG\_Q\_PRIOR\_MEAN,}
    \AttributeTok{pop\_log\_q\_sd   =} \FloatTok{0.5}\NormalTok{) \{}

  \FunctionTok{set.seed}\NormalTok{(base\_seed)}
\NormalTok{  z\_log\_r     }\OtherTok{\textless{}{-}} \FunctionTok{rnorm}\NormalTok{(n\_subjects, }\DecValTok{0}\NormalTok{, }\DecValTok{1}\NormalTok{)}
\NormalTok{  subj\_log\_r  }\OtherTok{\textless{}{-}}\NormalTok{ pop\_log\_r\_mean }\SpecialCharTok{+}\NormalTok{ pop\_log\_r\_sd }\SpecialCharTok{*}\NormalTok{ z\_log\_r}
\NormalTok{  subj\_r      }\OtherTok{\textless{}{-}} \FunctionTok{exp}\NormalTok{(subj\_log\_r)}

\NormalTok{  z\_log\_q     }\OtherTok{\textless{}{-}} \FunctionTok{rnorm}\NormalTok{(n\_subjects, }\DecValTok{0}\NormalTok{, }\DecValTok{1}\NormalTok{)}
\NormalTok{  subj\_log\_q  }\OtherTok{\textless{}{-}}\NormalTok{ pop\_log\_q\_mean }\SpecialCharTok{+}\NormalTok{ pop\_log\_q\_sd }\SpecialCharTok{*}\NormalTok{ z\_log\_q}
\NormalTok{  subj\_q      }\OtherTok{\textless{}{-}} \FunctionTok{exp}\NormalTok{(subj\_log\_q)}

\NormalTok{  true\_params }\OtherTok{\textless{}{-}} \FunctionTok{tibble}\NormalTok{(}
    \AttributeTok{agent\_id   =} \FunctionTok{seq\_len}\NormalTok{(n\_subjects),}
    \AttributeTok{log\_r\_true =}\NormalTok{ subj\_log\_r,}
    \AttributeTok{r\_true     =}\NormalTok{ subj\_r,}
    \AttributeTok{log\_q\_true =}\NormalTok{ subj\_log\_q,}
    \AttributeTok{q\_true     =}\NormalTok{ subj\_q}
\NormalTok{  )}

\NormalTok{  sim\_data }\OtherTok{\textless{}{-}} \FunctionTok{map\_dfr}\NormalTok{(}\FunctionTok{seq\_len}\NormalTok{(n\_subjects), }\ControlFlowTok{function}\NormalTok{(j) \{}
\NormalTok{    sched }\OtherTok{\textless{}{-}} \FunctionTok{make\_subject\_schedule}\NormalTok{(stimulus\_info, n\_blocks,}
                                   \AttributeTok{seed =}\NormalTok{ base\_seed }\SpecialCharTok{*} \DecValTok{100} \SpecialCharTok{+}\NormalTok{ j)}
\NormalTok{    obs   }\OtherTok{\textless{}{-}} \FunctionTok{as.matrix}\NormalTok{(sched[, }\FunctionTok{c}\NormalTok{(}\StringTok{"height"}\NormalTok{, }\StringTok{"position"}\NormalTok{)])}
\NormalTok{    res   }\OtherTok{\textless{}{-}} \FunctionTok{prototype\_kalman}\NormalTok{(}
      \AttributeTok{r\_value            =}\NormalTok{ subj\_r[j],}
      \AttributeTok{q\_value            =}\NormalTok{ subj\_q[j],}
      \AttributeTok{obs                =}\NormalTok{ obs,}
      \AttributeTok{cat\_one            =}\NormalTok{ sched}\SpecialCharTok{$}\NormalTok{category\_feedback,}
      \AttributeTok{initial\_mu         =} \FunctionTok{list}\NormalTok{(}\FunctionTok{rep}\NormalTok{(}\FloatTok{2.5}\NormalTok{, }\DecValTok{2}\NormalTok{), }\FunctionTok{rep}\NormalTok{(}\FloatTok{2.5}\NormalTok{, }\DecValTok{2}\NormalTok{)),}
      \AttributeTok{initial\_sigma\_diag =} \FloatTok{10.0}
\NormalTok{    )}
\NormalTok{    sched }\SpecialCharTok{|\textgreater{}}
      \FunctionTok{mutate}\NormalTok{(}
        \AttributeTok{agent\_id     =}\NormalTok{ j,}
        \AttributeTok{sim\_response =}\NormalTok{ res}\SpecialCharTok{$}\NormalTok{sim\_response,}
        \AttributeTok{correct      =} \FunctionTok{as.integer}\NormalTok{(category\_feedback }\SpecialCharTok{==}\NormalTok{ sim\_response)}
\NormalTok{      ) }\SpecialCharTok{|\textgreater{}}
      \FunctionTok{group\_by}\NormalTok{(agent\_id) }\SpecialCharTok{|\textgreater{}}
      \FunctionTok{mutate}\NormalTok{(}\AttributeTok{performance =} \FunctionTok{cumsum}\NormalTok{(correct) }\SpecialCharTok{/} \FunctionTok{row\_number}\NormalTok{()) }\SpecialCharTok{|\textgreater{}}
      \FunctionTok{ungroup}\NormalTok{()}
\NormalTok{  \})}

  \FunctionTok{list}\NormalTok{(}\AttributeTok{data =}\NormalTok{ sim\_data, }\AttributeTok{true\_params =}\NormalTok{ true\_params)}
\NormalTok{\}}
\end{Highlighting}
\end{Shaded}

\begin{Shaded}
\begin{Highlighting}[]
\CommentTok{\# Population parameters for multilevel simulation}
\NormalTok{pop\_params\_ml\_proto }\OtherTok{\textless{}{-}} \FunctionTok{list}\NormalTok{(}
  \AttributeTok{pop\_log\_r\_mean =} \FunctionTok{log}\NormalTok{(}\FloatTok{1.0}\NormalTok{),   }\CommentTok{\# prior median r = 1.0}
  \AttributeTok{pop\_log\_r\_sd   =} \FloatTok{0.5}\NormalTok{,        }\CommentTok{\# moderate between{-}subject heterogeneity}
  \AttributeTok{n\_subjects     =} \DecValTok{20}
\NormalTok{)}

\NormalTok{sim\_file\_ml\_proto    }\OtherTok{\textless{}{-}} \FunctionTok{here}\NormalTok{(}\StringTok{"simdata"}\NormalTok{, }\StringTok{"ch13\_proto\_ml\_sim\_data.csv"}\NormalTok{)}
\NormalTok{true\_params\_file\_ml  }\OtherTok{\textless{}{-}} \FunctionTok{here}\NormalTok{(}\StringTok{"simdata"}\NormalTok{, }\StringTok{"ch13\_proto\_ml\_true\_params.csv"}\NormalTok{)}

\ControlFlowTok{if}\NormalTok{ (regenerate\_simulations }\SpecialCharTok{||}
    \SpecialCharTok{!}\FunctionTok{file.exists}\NormalTok{(sim\_file\_ml\_proto) }\SpecialCharTok{||}
    \SpecialCharTok{!}\FunctionTok{file.exists}\NormalTok{(true\_params\_file\_ml)) \{}

\NormalTok{  ml\_sim\_out }\OtherTok{\textless{}{-}} \FunctionTok{simulate\_multilevel\_prototype}\NormalTok{(}
    \AttributeTok{n\_subjects     =}\NormalTok{ pop\_params\_ml\_proto}\SpecialCharTok{$}\NormalTok{n\_subjects,}
    \AttributeTok{stimulus\_info  =}\NormalTok{ stimulus\_info,}
    \AttributeTok{n\_blocks       =}\NormalTok{ n\_blocks,}
    \AttributeTok{base\_seed      =} \DecValTok{42}\NormalTok{,}
    \AttributeTok{pop\_log\_r\_mean =}\NormalTok{ pop\_params\_ml\_proto}\SpecialCharTok{$}\NormalTok{pop\_log\_r\_mean,}
    \AttributeTok{pop\_log\_r\_sd   =}\NormalTok{ pop\_params\_ml\_proto}\SpecialCharTok{$}\NormalTok{pop\_log\_r\_sd}
\NormalTok{  )}
\NormalTok{  multilevel\_proto\_data    }\OtherTok{\textless{}{-}}\NormalTok{ ml\_sim\_out}\SpecialCharTok{$}\NormalTok{data}
\NormalTok{  multilevel\_proto\_params  }\OtherTok{\textless{}{-}}\NormalTok{ ml\_sim\_out}\SpecialCharTok{$}\NormalTok{true\_params}

  \FunctionTok{write\_csv}\NormalTok{(multilevel\_proto\_data,   sim\_file\_ml\_proto)}
  \FunctionTok{write\_csv}\NormalTok{(multilevel\_proto\_params, true\_params\_file\_ml)}
  \FunctionTok{cat}\NormalTok{(}\StringTok{"Multilevel prototype simulation complete. Results saved.}\SpecialCharTok{\textbackslash{}n}\StringTok{"}\NormalTok{)}
\NormalTok{\} }\ControlFlowTok{else}\NormalTok{ \{}
\NormalTok{  multilevel\_proto\_data   }\OtherTok{\textless{}{-}} \FunctionTok{read\_csv}\NormalTok{(sim\_file\_ml\_proto,   }\AttributeTok{show\_col\_types =} \ConstantTok{FALSE}\NormalTok{)}
\NormalTok{  multilevel\_proto\_params }\OtherTok{\textless{}{-}} \FunctionTok{read\_csv}\NormalTok{(true\_params\_file\_ml, }\AttributeTok{show\_col\_types =} \ConstantTok{FALSE}\NormalTok{)}
  \FunctionTok{cat}\NormalTok{(}\StringTok{"Loaded existing multilevel prototype simulation results.}\SpecialCharTok{\textbackslash{}n}\StringTok{"}\NormalTok{)}
\NormalTok{\}}
\end{Highlighting}
\end{Shaded}

\begin{verbatim}
Loaded existing multilevel prototype simulation results.
\end{verbatim}

\begin{Shaded}
\begin{Highlighting}[]
\FunctionTok{ggplot}\NormalTok{(multilevel\_proto\_data,}
       \FunctionTok{aes}\NormalTok{(}\AttributeTok{x =}\NormalTok{ trial\_within\_subject, }\AttributeTok{y =}\NormalTok{ performance, }\AttributeTok{group =}\NormalTok{ agent\_id)) }\SpecialCharTok{+}
  \FunctionTok{geom\_line}\NormalTok{(}\AttributeTok{alpha =} \FloatTok{0.25}\NormalTok{, }\AttributeTok{color =} \StringTok{"grey40"}\NormalTok{) }\SpecialCharTok{+}
  \FunctionTok{stat\_summary}\NormalTok{(}\AttributeTok{fun =}\NormalTok{ mean, }\AttributeTok{geom =} \StringTok{"line"}\NormalTok{, }\FunctionTok{aes}\NormalTok{(}\AttributeTok{group =} \DecValTok{1}\NormalTok{),}
               \AttributeTok{color =} \StringTok{"\#009E73"}\NormalTok{, }\AttributeTok{linewidth =} \FloatTok{1.5}\NormalTok{) }\SpecialCharTok{+}
  \FunctionTok{labs}\NormalTok{(}
    \AttributeTok{title    =} \StringTok{"Simulated Learning Curves: Multilevel Prototype Model"}\NormalTok{,}
    \AttributeTok{subtitle =} \StringTok{"Grey lines: individual agents; green: group mean"}\NormalTok{,}
    \AttributeTok{x =} \StringTok{"Trial"}\NormalTok{, }\AttributeTok{y =} \StringTok{"Cumulative Accuracy"}
\NormalTok{  ) }\SpecialCharTok{+}
  \FunctionTok{ylim}\NormalTok{(}\DecValTok{0}\NormalTok{, }\DecValTok{1}\NormalTok{)}
\end{Highlighting}
\end{Shaded}

\begin{center}
\includegraphics[width=0.8\linewidth,height=\textheight,keepaspectratio]{12-CategorizationModelsPrototypes_files/figure-pdf/ch13_simulate_multilevel_data-1.pdf}
\end{center}

\subsection{Implementing the Multilevel Prototype Model in
Stan}\label{implementing-the-multilevel-prototype-model-in-stan}

The multilevel Stan model extends the single-subject architecture by:

\begin{enumerate}
\def\labelenumi{\arabic{enumi}.}
\tightlist
\item
  Adding population hyperparameters \texttt{pop\_log\_r\_mean},
  \texttt{pop\_log\_r\_sd}, \texttt{pop\_log\_q\_mean}, and
  \texttt{pop\_log\_q\_sd} to the \texttt{parameters} block.
\item
  Adding per-subject NCP offsets
  \texttt{z\_log\_r{[}j{]}\ \textasciitilde{}\ Normal(0,\ 1)} and
  \texttt{z\_log\_q{[}j{]}\ \textasciitilde{}\ Normal(0,\ 1)}.
\item
  Looping the Kalman filter (with prediction step) over subjects in
  \texttt{transformed\ parameters}, exactly as the GCM does in Chapter
  11.
\end{enumerate}

The resulting model has \(4 + 2J\) parameters (\(J\) is the number of
subjects) --- still relatively compact compared to exemplar models ---
which makes posterior geometry and SBC diagnostics particularly legible.

\begin{Shaded}
\begin{Highlighting}[]
\NormalTok{prototype\_ml\_stan }\OtherTok{\textless{}{-}} \StringTok{"}
\StringTok{// Kalman Filter Prototype Model — Multilevel (Non{-}Centred Parameterisation)}
\StringTok{//}
\StringTok{// New relative to the single{-}subject model:}
\StringTok{//   1. Population hyperparameters for log\_r AND log\_q.}
\StringTok{//   2. Per{-}subject NCP offsets z\_log\_r[j] and z\_log\_q[j] \textasciitilde{} Normal(0, 1).}
\StringTok{//   3. Subject{-}indexed Kalman filter loop with prediction step in transformed parameters.}
\StringTok{//   4. ALL prior hyperparameters passed through the data block.}

\StringTok{data \{}
\StringTok{  int\textless{}lower=1\textgreater{} N\_total;}
\StringTok{  int\textless{}lower=1\textgreater{} N\_subjects;}
\StringTok{  int\textless{}lower=1\textgreater{} N\_features;}

\StringTok{  array[N\_subjects] int\textless{}lower=1, upper=N\_total\textgreater{} subj\_start;}
\StringTok{  array[N\_subjects] int\textless{}lower=1, upper=N\_total\textgreater{} subj\_end;}

\StringTok{  array[N\_total] int\textless{}lower=0, upper=1\textgreater{} y;}
\StringTok{  array[N\_total] int\textless{}lower=0, upper=1\textgreater{} cat\_one;}
\StringTok{  array[N\_total, N\_features] real obs;}

\StringTok{  // Fixed structural values (same as single{-}subject; not inferred)}
\StringTok{  vector[N\_features] initial\_mu\_cat0;}
\StringTok{  vector[N\_features] initial\_mu\_cat1;}
\StringTok{  real\textless{}lower=0\textgreater{} initial\_sigma\_diag;}

\StringTok{  // Prior hyperparameters for r}
\StringTok{  real pop\_log\_r\_mean\_prior\_mean;}
\StringTok{  real\textless{}lower=0\textgreater{} pop\_log\_r\_mean\_prior\_sd;}
\StringTok{  real\textless{}lower=0\textgreater{} pop\_log\_r\_sd\_prior\_rate;}

\StringTok{  // Prior hyperparameters for q}
\StringTok{  real pop\_log\_q\_mean\_prior\_mean;}
\StringTok{  real\textless{}lower=0\textgreater{} pop\_log\_q\_mean\_prior\_sd;}
\StringTok{  real\textless{}lower=0\textgreater{} pop\_log\_q\_sd\_prior\_rate;}
\StringTok{\}}

\StringTok{parameters \{}
\StringTok{  real pop\_log\_r\_mean;}
\StringTok{  real\textless{}lower=0\textgreater{} pop\_log\_r\_sd;}
\StringTok{  vector[N\_subjects] z\_log\_r;}

\StringTok{  real pop\_log\_q\_mean;}
\StringTok{  real\textless{}lower=0\textgreater{} pop\_log\_q\_sd;}
\StringTok{  vector[N\_subjects] z\_log\_q;}
\StringTok{\}}

\StringTok{transformed parameters \{}
\StringTok{  vector[N\_subjects] subj\_log\_r;}
\StringTok{  vector\textless{}lower=0\textgreater{}[N\_subjects] subj\_r;}
\StringTok{  vector[N\_subjects] subj\_log\_q;}
\StringTok{  vector\textless{}lower=0\textgreater{}[N\_subjects] subj\_q;}
\StringTok{  vector\textless{}lower=1e{-}9, upper=1{-}1e{-}9\textgreater{}[N\_total] prob\_cat1;}

\StringTok{  for (j in 1:N\_subjects) \{}
\StringTok{    subj\_log\_r[j] = pop\_log\_r\_mean + z\_log\_r[j] * pop\_log\_r\_sd;}
\StringTok{    subj\_r[j]     = exp(subj\_log\_r[j]);}
\StringTok{    subj\_log\_q[j] = pop\_log\_q\_mean + z\_log\_q[j] * pop\_log\_q\_sd;}
\StringTok{    subj\_q[j]     = exp(subj\_log\_q[j]);}
\StringTok{  \}}

\StringTok{  \{}
\StringTok{    matrix[N\_features, N\_features] I\_mat =}
\StringTok{      diag\_matrix(rep\_vector(1.0, N\_features));}

\StringTok{    for (j in 1:N\_subjects) \{}
\StringTok{      vector[N\_features] mu\_cat0 = initial\_mu\_cat0;}
\StringTok{      vector[N\_features] mu\_cat1 = initial\_mu\_cat1;}
\StringTok{      matrix[N\_features, N\_features] sigma\_cat0 =}
\StringTok{        diag\_matrix(rep\_vector(initial\_sigma\_diag, N\_features));}
\StringTok{      matrix[N\_features, N\_features] sigma\_cat1 =}
\StringTok{        diag\_matrix(rep\_vector(initial\_sigma\_diag, N\_features));}
\StringTok{      matrix[N\_features, N\_features] r\_matrix =}
\StringTok{        diag\_matrix(rep\_vector(subj\_r[j], N\_features));}
\StringTok{      matrix[N\_features, N\_features] q\_matrix =}
\StringTok{        diag\_matrix(rep\_vector(subj\_q[j], N\_features));}

\StringTok{      for (i in subj\_start[j]:subj\_end[j]) \{}
\StringTok{        vector[N\_features] x = to\_vector(obs[i]);}

\StringTok{        // ── Prediction step: add process noise to both categories ──────────}
\StringTok{        sigma\_cat0 = sigma\_cat0 + q\_matrix;}
\StringTok{        sigma\_cat1 = sigma\_cat1 + q\_matrix;}

\StringTok{        // ── Decision ────────────────────────────────────────────────────────}
\StringTok{        matrix[N\_features, N\_features] cov0 = sigma\_cat0 + r\_matrix;}
\StringTok{        matrix[N\_features, N\_features] cov1 = sigma\_cat1 + r\_matrix;}
\StringTok{        real log\_p0 = multi\_normal\_lpdf(x | mu\_cat0, cov0);}
\StringTok{        real log\_p1 = multi\_normal\_lpdf(x | mu\_cat1, cov1);}
\StringTok{        real p\_i    = exp(log\_p1 {-} log\_sum\_exp(log\_p0, log\_p1));}
\StringTok{        prob\_cat1[i] = fmax(1e{-}9, fmin(1 {-} 1e{-}9, p\_i));}

\StringTok{        // ── Update ───────────────────────────────────────────────────────────}
\StringTok{        if (cat\_one[i] == 1) \{}
\StringTok{          vector[N\_features] innov = x {-} mu\_cat1;}
\StringTok{          matrix[N\_features, N\_features] S  = sigma\_cat1 + r\_matrix;}
\StringTok{          matrix[N\_features, N\_features] K  = mdivide\_right\_spd(sigma\_cat1, S);}
\StringTok{          matrix[N\_features, N\_features] IK = I\_mat {-} K;}
\StringTok{          mu\_cat1    = mu\_cat1 + K * innov;}
\StringTok{          sigma\_cat1 = IK * sigma\_cat1 * IK\textquotesingle{} + K * r\_matrix * K\textquotesingle{};}
\StringTok{          sigma\_cat1 = 0.5 * (sigma\_cat1 + sigma\_cat1\textquotesingle{});}
\StringTok{        \} else \{}
\StringTok{          vector[N\_features] innov = x {-} mu\_cat0;}
\StringTok{          matrix[N\_features, N\_features] S  = sigma\_cat0 + r\_matrix;}
\StringTok{          matrix[N\_features, N\_features] K  = mdivide\_right\_spd(sigma\_cat0, S);}
\StringTok{          matrix[N\_features, N\_features] IK = I\_mat {-} K;}
\StringTok{          mu\_cat0    = mu\_cat0 + K * innov;}
\StringTok{          sigma\_cat0 = IK * sigma\_cat0 * IK\textquotesingle{} + K * r\_matrix * K\textquotesingle{};}
\StringTok{          sigma\_cat0 = 0.5 * (sigma\_cat0 + sigma\_cat0\textquotesingle{});}
\StringTok{        \}}
\StringTok{      \}}
\StringTok{    \}}
\StringTok{  \}}
\StringTok{\}}

\StringTok{model \{}
\StringTok{  target += normal\_lpdf(pop\_log\_r\_mean | pop\_log\_r\_mean\_prior\_mean,}
\StringTok{                                         pop\_log\_r\_mean\_prior\_sd);}
\StringTok{  target += exponential\_lpdf(pop\_log\_r\_sd | pop\_log\_r\_sd\_prior\_rate);}
\StringTok{  target += std\_normal\_lpdf(z\_log\_r);}

\StringTok{  target += normal\_lpdf(pop\_log\_q\_mean | pop\_log\_q\_mean\_prior\_mean,}
\StringTok{                                         pop\_log\_q\_mean\_prior\_sd);}
\StringTok{  target += exponential\_lpdf(pop\_log\_q\_sd | pop\_log\_q\_sd\_prior\_rate);}
\StringTok{  target += std\_normal\_lpdf(z\_log\_q);}

\StringTok{  target += bernoulli\_lpmf(y | prob\_cat1);}
\StringTok{\}}

\StringTok{generated quantities \{}
\StringTok{  vector[N\_total] log\_lik;}
\StringTok{  real lprior;}
\StringTok{  for (i in 1:N\_total)}
\StringTok{    log\_lik[i] = bernoulli\_lpmf(y[i] | prob\_cat1[i]);}
\StringTok{  lprior = normal\_lpdf(pop\_log\_r\_mean | pop\_log\_r\_mean\_prior\_mean, pop\_log\_r\_mean\_prior\_sd) +}
\StringTok{           exponential\_lpdf(pop\_log\_r\_sd | pop\_log\_r\_sd\_prior\_rate) +}
\StringTok{           std\_normal\_lpdf(z\_log\_r) +}
\StringTok{           normal\_lpdf(pop\_log\_q\_mean | pop\_log\_q\_mean\_prior\_mean, pop\_log\_q\_mean\_prior\_sd) +}
\StringTok{           exponential\_lpdf(pop\_log\_q\_sd | pop\_log\_q\_sd\_prior\_rate) +}
\StringTok{           std\_normal\_lpdf(z\_log\_q);}
\StringTok{\}}
\StringTok{"}

\NormalTok{stan\_file\_proto\_ml }\OtherTok{\textless{}{-}} \FunctionTok{here}\NormalTok{(}\StringTok{"stan"}\NormalTok{, }\StringTok{"ch13\_prototype\_ml.stan"}\NormalTok{)}
\FunctionTok{write\_stan\_file}\NormalTok{(prototype\_ml\_stan, }\AttributeTok{dir =} \FunctionTok{here}\NormalTok{(}\StringTok{"stan"}\NormalTok{),}
                \AttributeTok{basename =} \StringTok{"ch13\_prototype\_ml.stan"}\NormalTok{)}
\end{Highlighting}
\end{Shaded}

\begin{verbatim}
[1] "/Users/au209589/Dropbox/Teaching/AdvancedCognitiveModeling23_book/stan/ch13_prototype_ml.stan"
\end{verbatim}

\begin{Shaded}
\begin{Highlighting}[]
\NormalTok{mod\_prototype\_ml }\OtherTok{\textless{}{-}} \FunctionTok{cmdstan\_model}\NormalTok{(stan\_file\_proto\_ml)}
\FunctionTok{cat}\NormalTok{(}\StringTok{"Multilevel prototype Stan model compiled successfully.}\SpecialCharTok{\textbackslash{}n}\StringTok{"}\NormalTok{)}
\end{Highlighting}
\end{Shaded}

\begin{verbatim}
Multilevel prototype Stan model compiled successfully.
\end{verbatim}

\subsection{Fitting the Multilevel Prototype
Model}\label{fitting-the-multilevel-prototype-model}

\begin{Shaded}
\begin{Highlighting}[]
\CommentTok{\# ── Build the Stan data list ─────────────────────────────────────────────────}
\NormalTok{proto\_ml\_sorted }\OtherTok{\textless{}{-}}\NormalTok{ multilevel\_proto\_data }\SpecialCharTok{|\textgreater{}}
  \FunctionTok{mutate}\NormalTok{(}\AttributeTok{subj\_id\_stan =} \FunctionTok{as.integer}\NormalTok{(}\FunctionTok{factor}\NormalTok{(agent\_id))) }\SpecialCharTok{|\textgreater{}}
  \FunctionTok{arrange}\NormalTok{(subj\_id\_stan, trial\_within\_subject)}

\NormalTok{proto\_subj\_bounds }\OtherTok{\textless{}{-}}\NormalTok{ proto\_ml\_sorted }\SpecialCharTok{|\textgreater{}}
  \FunctionTok{mutate}\NormalTok{(}\AttributeTok{row\_idx =} \FunctionTok{row\_number}\NormalTok{()) }\SpecialCharTok{|\textgreater{}}
  \FunctionTok{group\_by}\NormalTok{(subj\_id\_stan) }\SpecialCharTok{|\textgreater{}}
  \FunctionTok{summarise}\NormalTok{(}
    \AttributeTok{subj\_start =} \FunctionTok{min}\NormalTok{(row\_idx),}
    \AttributeTok{subj\_end   =} \FunctionTok{max}\NormalTok{(row\_idx),}
    \AttributeTok{.groups    =} \StringTok{"drop"}
\NormalTok{  ) }\SpecialCharTok{|\textgreater{}}
  \FunctionTok{arrange}\NormalTok{(subj\_id\_stan)}

\CommentTok{\# Prior design note on pop\_log\_r\_sd:}
\CommentTok{\# pop\_log\_r\_mean \textasciitilde{} Normal(0, 1): same as single{-}subject prior on log\_r.}
\CommentTok{\# pop\_log\_r\_sd   \textasciitilde{} Exponential(2): rate=2 gives median ≈ 0.35 (SD on log scale).}
\CommentTok{\# Marginal SD of subj\_log\_r ≈ sqrt(1 + 0.5\^{}2) ≈ 1.12 at the 50th pctile}
\CommentTok{\# of the pop\_log\_r\_sd prior — consistent with the single{-}agent prior range.}
\NormalTok{proto\_ml\_data }\OtherTok{\textless{}{-}} \FunctionTok{list}\NormalTok{(}
  \AttributeTok{N\_total                      =} \FunctionTok{nrow}\NormalTok{(proto\_ml\_sorted),}
  \AttributeTok{N\_subjects                   =} \FunctionTok{max}\NormalTok{(proto\_ml\_sorted}\SpecialCharTok{$}\NormalTok{subj\_id\_stan),}
  \AttributeTok{N\_features                   =} \DecValTok{2}\NormalTok{L,}
  \AttributeTok{subj\_start                   =}\NormalTok{ proto\_subj\_bounds}\SpecialCharTok{$}\NormalTok{subj\_start,}
  \AttributeTok{subj\_end                     =}\NormalTok{ proto\_subj\_bounds}\SpecialCharTok{$}\NormalTok{subj\_end,}
  \AttributeTok{y                            =} \FunctionTok{as.integer}\NormalTok{(proto\_ml\_sorted}\SpecialCharTok{$}\NormalTok{sim\_response),}
  \AttributeTok{cat\_one                      =} \FunctionTok{as.integer}\NormalTok{(proto\_ml\_sorted}\SpecialCharTok{$}\NormalTok{category\_feedback),}
  \AttributeTok{obs                          =} \FunctionTok{as.matrix}\NormalTok{(proto\_ml\_sorted[, }\FunctionTok{c}\NormalTok{(}\StringTok{"height"}\NormalTok{, }\StringTok{"position"}\NormalTok{)]),}
  \AttributeTok{initial\_mu\_cat0              =} \FunctionTok{c}\NormalTok{(}\FloatTok{2.5}\NormalTok{, }\FloatTok{2.5}\NormalTok{),}
  \AttributeTok{initial\_mu\_cat1              =} \FunctionTok{c}\NormalTok{(}\FloatTok{2.5}\NormalTok{, }\FloatTok{2.5}\NormalTok{),}
  \AttributeTok{initial\_sigma\_diag           =} \FloatTok{10.0}\NormalTok{,}
  \AttributeTok{pop\_log\_r\_mean\_prior\_mean    =}\NormalTok{ LOG\_R\_PRIOR\_MEAN,}
  \AttributeTok{pop\_log\_r\_mean\_prior\_sd      =}\NormalTok{ LOG\_R\_PRIOR\_SD,}
  \AttributeTok{pop\_log\_r\_sd\_prior\_rate      =} \FloatTok{2.0}\NormalTok{,}
  \AttributeTok{pop\_log\_q\_mean\_prior\_mean    =}\NormalTok{ LOG\_Q\_PRIOR\_MEAN,}
  \AttributeTok{pop\_log\_q\_mean\_prior\_sd      =}\NormalTok{ LOG\_Q\_PRIOR\_SD,}
  \AttributeTok{pop\_log\_q\_sd\_prior\_rate      =} \FloatTok{2.0}
\NormalTok{)}

\NormalTok{fit\_filepath\_ml\_proto }\OtherTok{\textless{}{-}} \FunctionTok{here}\NormalTok{(}\StringTok{"simmodels"}\NormalTok{, }\StringTok{"ch13\_prototype\_ml\_fit.rds"}\NormalTok{)}

\ControlFlowTok{if}\NormalTok{ (regenerate\_simulations }\SpecialCharTok{||} \SpecialCharTok{!}\FunctionTok{file.exists}\NormalTok{(fit\_filepath\_ml\_proto)) \{}
\NormalTok{  pf\_ml\_proto }\OtherTok{\textless{}{-}} \FunctionTok{tryCatch}\NormalTok{(}
\NormalTok{    mod\_prototype\_ml}\SpecialCharTok{$}\FunctionTok{pathfinder}\NormalTok{(}
      \AttributeTok{data      =}\NormalTok{ proto\_ml\_data,}
      \AttributeTok{seed      =} \DecValTok{42}\NormalTok{,}
      \AttributeTok{num\_paths =} \DecValTok{4}\NormalTok{,}
      \AttributeTok{refresh   =} \DecValTok{0}
\NormalTok{    ),}
    \AttributeTok{error =} \ControlFlowTok{function}\NormalTok{(e) \{}
      \FunctionTok{message}\NormalTok{(}\StringTok{"Pathfinder failed — falling back to default initialisation."}\NormalTok{)}
      \ConstantTok{NULL}
\NormalTok{    \}}
\NormalTok{  )}

\NormalTok{  fit\_proto\_ml }\OtherTok{\textless{}{-}}\NormalTok{ mod\_prototype\_ml}\SpecialCharTok{$}\FunctionTok{sample}\NormalTok{(}
    \AttributeTok{data            =}\NormalTok{ proto\_ml\_data,}
    \AttributeTok{init            =}\NormalTok{ pf\_ml\_proto,}
    \AttributeTok{chains          =} \DecValTok{4}\NormalTok{,}
    \AttributeTok{parallel\_chains =} \FunctionTok{min}\NormalTok{(}\DecValTok{4}\NormalTok{, }\FunctionTok{availableCores}\NormalTok{()),}
    \AttributeTok{iter\_warmup     =} \DecValTok{1500}\NormalTok{,}
    \AttributeTok{iter\_sampling   =} \DecValTok{1500}\NormalTok{,}
    \AttributeTok{refresh         =} \DecValTok{200}\NormalTok{,}
    \AttributeTok{adapt\_delta     =} \FloatTok{0.95}\NormalTok{,}
    \AttributeTok{max\_treedepth   =} \DecValTok{12}
\NormalTok{  )}
\NormalTok{  fit\_proto\_ml}\SpecialCharTok{$}\FunctionTok{save\_object}\NormalTok{(fit\_filepath\_ml\_proto)}
  \FunctionTok{cat}\NormalTok{(}\StringTok{"Multilevel prototype model fitted and saved.}\SpecialCharTok{\textbackslash{}n}\StringTok{"}\NormalTok{)}
\NormalTok{\} }\ControlFlowTok{else}\NormalTok{ \{}
\NormalTok{  fit\_proto\_ml }\OtherTok{\textless{}{-}} \FunctionTok{readRDS}\NormalTok{(fit\_filepath\_ml\_proto)}
  \FunctionTok{cat}\NormalTok{(}\StringTok{"Loaded existing multilevel prototype fit.}\SpecialCharTok{\textbackslash{}n}\StringTok{"}\NormalTok{)}
\NormalTok{\}}
\end{Highlighting}
\end{Shaded}

\begin{verbatim}
Loaded existing multilevel prototype fit.
\end{verbatim}

\subsection{Mandatory MCMC Diagnostic Battery
(Multilevel)}\label{mandatory-mcmc-diagnostic-battery-multilevel}

\begin{Shaded}
\begin{Highlighting}[]
\ControlFlowTok{if}\NormalTok{ (}\SpecialCharTok{!}\FunctionTok{is.null}\NormalTok{(fit\_proto\_ml)) \{}
\NormalTok{  pop\_vars\_ml }\OtherTok{\textless{}{-}} \FunctionTok{c}\NormalTok{(}\StringTok{"pop\_log\_r\_mean"}\NormalTok{, }\StringTok{"pop\_log\_r\_sd"}\NormalTok{, }\StringTok{"pop\_log\_q\_mean"}\NormalTok{, }\StringTok{"pop\_log\_q\_sd"}\NormalTok{)}

\NormalTok{  diag\_tbl\_ml }\OtherTok{\textless{}{-}} \FunctionTok{diagnostic\_summary\_table}\NormalTok{(fit\_proto\_ml, }\AttributeTok{params =}\NormalTok{ pop\_vars\_ml)}
  \FunctionTok{print}\NormalTok{(diag\_tbl\_ml)}

  \ControlFlowTok{if}\NormalTok{ (}\FunctionTok{any}\NormalTok{(}\FunctionTok{c}\NormalTok{(diag\_tbl\_ml}\SpecialCharTok{$}\NormalTok{n\_divergent }\SpecialCharTok{\textgreater{}} \DecValTok{0}\NormalTok{,}
\NormalTok{            diag\_tbl\_ml}\SpecialCharTok{$}\NormalTok{max\_rhat    }\SpecialCharTok{\textgreater{}=} \FloatTok{1.01}\NormalTok{,}
\NormalTok{            diag\_tbl\_ml}\SpecialCharTok{$}\NormalTok{min\_ess\_bulk }\SpecialCharTok{\textless{}} \DecValTok{400}\NormalTok{))) \{}
    \FunctionTok{warning}\NormalTok{(}\StringTok{"Multilevel diagnostic battery has failures. "}\NormalTok{,}
            \StringTok{"Inspect mcmc\_pairs() and consider tightening adapt\_delta or "}\NormalTok{,}
            \StringTok{"increasing max\_treedepth before interpreting posteriors."}\NormalTok{)}
\NormalTok{  \}}

\NormalTok{  bayesplot}\SpecialCharTok{::}\FunctionTok{mcmc\_pairs}\NormalTok{(}
\NormalTok{    fit\_proto\_ml}\SpecialCharTok{$}\FunctionTok{draws}\NormalTok{(}\FunctionTok{c}\NormalTok{(}\StringTok{"pop\_log\_r\_mean"}\NormalTok{, }\StringTok{"pop\_log\_r\_sd"}\NormalTok{,}
                         \StringTok{"pop\_log\_q\_mean"}\NormalTok{, }\StringTok{"pop\_log\_q\_sd"}\NormalTok{)),}
    \AttributeTok{diag\_fun     =} \StringTok{"dens"}\NormalTok{,}
    \AttributeTok{off\_diag\_fun =} \StringTok{"hex"}
\NormalTok{  )}

\NormalTok{  pop\_summary }\OtherTok{\textless{}{-}}\NormalTok{ fit\_proto\_ml}\SpecialCharTok{$}\FunctionTok{summary}\NormalTok{(}\AttributeTok{variables =}\NormalTok{ pop\_vars\_ml)}
  \FunctionTok{cat}\NormalTok{(}\StringTok{"=== Population Parameter Summary ===}\SpecialCharTok{\textbackslash{}n}\StringTok{"}\NormalTok{)}
  \FunctionTok{print}\NormalTok{(pop\_summary)}
  \FunctionTok{cat}\NormalTok{(}\StringTok{"}\SpecialCharTok{\textbackslash{}n}\StringTok{True population parameters used for simulation:}\SpecialCharTok{\textbackslash{}n}\StringTok{"}\NormalTok{)}
  \FunctionTok{cat}\NormalTok{(}\StringTok{"  pop\_log\_r\_mean:"}\NormalTok{, pop\_params\_ml\_proto}\SpecialCharTok{$}\NormalTok{pop\_log\_r\_mean, }\StringTok{"}\SpecialCharTok{\textbackslash{}n}\StringTok{"}\NormalTok{)}
  \FunctionTok{cat}\NormalTok{(}\StringTok{"  pop\_log\_r\_sd:  "}\NormalTok{, pop\_params\_ml\_proto}\SpecialCharTok{$}\NormalTok{pop\_log\_r\_sd, }\StringTok{"}\SpecialCharTok{\textbackslash{}n}\StringTok{"}\NormalTok{)}
\NormalTok{\}}
\end{Highlighting}
\end{Shaded}

\begin{verbatim}
# A tibble: 6 x 4
  metric                           value threshold pass 
  <chr>                            <dbl> <chr>     <lgl>
1 Divergences (zero tolerance)    0      == 0      TRUE 
2 Max rank-normalised R-hat       1      < 1.01    TRUE 
3 Min bulk ESS                 1405      > 400     TRUE 
4 Min tail ESS                 2050      > 400     TRUE 
5 Min E-BFMI                      0.783  > 0.2     TRUE 
6 Max MCSE / posterior SD         0.0267 < 0.05    TRUE 
\end{verbatim}

\begin{verbatim}
=== Population Parameter Summary ===
# A tibble: 4 x 10
  variable       mean  median    sd   mad      q5    q95  rhat ess_bulk ess_tail
  <chr>         <dbl>   <dbl> <dbl> <dbl>   <dbl>  <dbl> <dbl>    <dbl>    <dbl>
1 pop_log_r_~  0.0291  0.0272 0.176 0.172 -0.262   0.320  1.00    1852.    2050.
2 pop_log_r_~  0.656   0.640  0.147 0.138  0.444   0.918  1.00    2431.    3171.
3 pop_log_q_~ -2.65   -2.65   0.334 0.317 -3.19   -2.11   1.00    3663.    2823.
4 pop_log_q_~  0.503   0.440  0.369 0.388  0.0342  1.19   1.00    1405.    2739.

True population parameters used for simulation:
  pop_log_r_mean: 0 
  pop_log_r_sd:   0.5 
\end{verbatim}

\subsection{Population Prior-Posterior
Update}\label{population-prior-posterior-update}

\begin{Shaded}
\begin{Highlighting}[]
\ControlFlowTok{if}\NormalTok{ (}\SpecialCharTok{!}\FunctionTok{is.null}\NormalTok{(fit\_proto\_ml)) \{}
\NormalTok{  draws\_ml\_proto }\OtherTok{\textless{}{-}} \FunctionTok{as\_draws\_df}\NormalTok{(}
\NormalTok{    fit\_proto\_ml}\SpecialCharTok{$}\FunctionTok{draws}\NormalTok{(}\AttributeTok{variables =} \FunctionTok{c}\NormalTok{(}\StringTok{"pop\_log\_r\_mean"}\NormalTok{, }\StringTok{"pop\_log\_r\_sd"}\NormalTok{,}
                                     \StringTok{"pop\_log\_q\_mean"}\NormalTok{, }\StringTok{"pop\_log\_q\_sd"}\NormalTok{))}
\NormalTok{  )}

\NormalTok{  n\_prior }\OtherTok{\textless{}{-}} \FunctionTok{nrow}\NormalTok{(draws\_ml\_proto)}
  \FunctionTok{set.seed}\NormalTok{(}\DecValTok{9}\NormalTok{)}
\NormalTok{  prior\_ml\_proto }\OtherTok{\textless{}{-}} \FunctionTok{tibble}\NormalTok{(}
    \AttributeTok{pop\_log\_r\_mean =} \FunctionTok{rnorm}\NormalTok{(n\_prior, LOG\_R\_PRIOR\_MEAN, LOG\_R\_PRIOR\_SD),}
    \AttributeTok{pop\_log\_r\_sd   =} \FunctionTok{rexp}\NormalTok{(n\_prior, }\FloatTok{2.0}\NormalTok{),}
    \AttributeTok{pop\_log\_q\_mean =} \FunctionTok{rnorm}\NormalTok{(n\_prior, LOG\_Q\_PRIOR\_MEAN, LOG\_Q\_PRIOR\_SD),}
    \AttributeTok{pop\_log\_q\_sd   =} \FunctionTok{rexp}\NormalTok{(n\_prior, }\FloatTok{2.0}\NormalTok{)}
\NormalTok{  )}

\NormalTok{  plot\_param\_update }\OtherTok{\textless{}{-}} \ControlFlowTok{function}\NormalTok{(param\_name, prior\_vec, post\_vec,}
\NormalTok{                                true\_val, fill\_color) \{}
\NormalTok{    df }\OtherTok{\textless{}{-}} \FunctionTok{tibble}\NormalTok{(}
      \AttributeTok{value        =} \FunctionTok{c}\NormalTok{(prior\_vec, post\_vec),}
      \AttributeTok{Distribution =} \FunctionTok{rep}\NormalTok{(}\FunctionTok{c}\NormalTok{(}\StringTok{"Prior"}\NormalTok{, }\StringTok{"Posterior"}\NormalTok{), }\AttributeTok{each =} \FunctionTok{length}\NormalTok{(prior\_vec))}
\NormalTok{    )}
    \FunctionTok{ggplot}\NormalTok{(df, }\FunctionTok{aes}\NormalTok{(}\AttributeTok{x =}\NormalTok{ value, }\AttributeTok{fill =}\NormalTok{ Distribution, }\AttributeTok{color =}\NormalTok{ Distribution)) }\SpecialCharTok{+}
      \FunctionTok{geom\_density}\NormalTok{(}\AttributeTok{alpha =} \FloatTok{0.5}\NormalTok{, }\AttributeTok{linewidth =} \FloatTok{0.5}\NormalTok{) }\SpecialCharTok{+}
      \FunctionTok{geom\_vline}\NormalTok{(}\AttributeTok{xintercept =}\NormalTok{ true\_val, }\AttributeTok{color =} \StringTok{"red"}\NormalTok{,}
                 \AttributeTok{linetype =} \StringTok{"dashed"}\NormalTok{, }\AttributeTok{linewidth =} \DecValTok{1}\NormalTok{) }\SpecialCharTok{+}
      \FunctionTok{scale\_fill\_manual}\NormalTok{(}\AttributeTok{values  =} \FunctionTok{c}\NormalTok{(}\StringTok{"Prior"} \OtherTok{=} \StringTok{"grey70"}\NormalTok{, }\StringTok{"Posterior"} \OtherTok{=}\NormalTok{ fill\_color)) }\SpecialCharTok{+}
      \FunctionTok{scale\_color\_manual}\NormalTok{(}\AttributeTok{values =} \FunctionTok{c}\NormalTok{(}\StringTok{"Prior"} \OtherTok{=} \StringTok{"grey50"}\NormalTok{, }\StringTok{"Posterior"} \OtherTok{=}\NormalTok{ fill\_color)) }\SpecialCharTok{+}
      \FunctionTok{labs}\NormalTok{(}\AttributeTok{title =}\NormalTok{ param\_name, }\AttributeTok{x =} \ConstantTok{NULL}\NormalTok{, }\AttributeTok{y =} \StringTok{"Density"}\NormalTok{) }\SpecialCharTok{+}
      \FunctionTok{theme}\NormalTok{(}\AttributeTok{axis.text.y =} \FunctionTok{element\_blank}\NormalTok{(), }\AttributeTok{axis.ticks.y =} \FunctionTok{element\_blank}\NormalTok{())}
\NormalTok{  \}}

\NormalTok{  p\_logr\_mean }\OtherTok{\textless{}{-}} \FunctionTok{plot\_param\_update}\NormalTok{(}
    \StringTok{"Population log(r) mean"}\NormalTok{,}
\NormalTok{    prior\_ml\_proto}\SpecialCharTok{$}\NormalTok{pop\_log\_r\_mean,}
\NormalTok{    draws\_ml\_proto}\SpecialCharTok{$}\NormalTok{pop\_log\_r\_mean,}
\NormalTok{    pop\_params\_ml\_proto}\SpecialCharTok{$}\NormalTok{pop\_log\_r\_mean,}
    \StringTok{"\#009E73"}
\NormalTok{  )}

\NormalTok{  p\_logr\_sd }\OtherTok{\textless{}{-}} \FunctionTok{plot\_param\_update}\NormalTok{(}
    \StringTok{"Population log(r) SD"}\NormalTok{,}
\NormalTok{    prior\_ml\_proto}\SpecialCharTok{$}\NormalTok{pop\_log\_r\_sd,}
\NormalTok{    draws\_ml\_proto}\SpecialCharTok{$}\NormalTok{pop\_log\_r\_sd,}
\NormalTok{    pop\_params\_ml\_proto}\SpecialCharTok{$}\NormalTok{pop\_log\_r\_sd,}
    \StringTok{"\#E69F00"}
\NormalTok{  )}

\NormalTok{  p\_logq\_mean }\OtherTok{\textless{}{-}} \FunctionTok{plot\_param\_update}\NormalTok{(}
    \StringTok{"Population log(q) mean"}\NormalTok{,}
\NormalTok{    prior\_ml\_proto}\SpecialCharTok{$}\NormalTok{pop\_log\_q\_mean,}
\NormalTok{    draws\_ml\_proto}\SpecialCharTok{$}\NormalTok{pop\_log\_q\_mean,}
\NormalTok{    LOG\_Q\_PRIOR\_MEAN,   }\CommentTok{\# true value = prior mean (q hierarchy not set in sim)}
    \StringTok{"\#56B4E9"}
\NormalTok{  )}

\NormalTok{  p\_logq\_sd }\OtherTok{\textless{}{-}} \FunctionTok{plot\_param\_update}\NormalTok{(}
    \StringTok{"Population log(q) SD"}\NormalTok{,}
\NormalTok{    prior\_ml\_proto}\SpecialCharTok{$}\NormalTok{pop\_log\_q\_sd,}
\NormalTok{    draws\_ml\_proto}\SpecialCharTok{$}\NormalTok{pop\_log\_q\_sd,}
    \FloatTok{0.5}\NormalTok{,   }\CommentTok{\# illustrative true value}
    \StringTok{"\#CC79A7"}
\NormalTok{  )}

  \FunctionTok{cat}\NormalTok{(}\StringTok{"=== Population Parameter Prior{-}Posterior Updates ===}\SpecialCharTok{\textbackslash{}n}\StringTok{"}\NormalTok{)}
  \FunctionTok{print}\NormalTok{((p\_logr\_mean }\SpecialCharTok{|}\NormalTok{ p\_logr\_sd) }\SpecialCharTok{/}\NormalTok{ (p\_logq\_mean }\SpecialCharTok{|}\NormalTok{ p\_logq\_sd))}
\NormalTok{\}}
\end{Highlighting}
\end{Shaded}

\begin{verbatim}
=== Population Parameter Prior-Posterior Updates ===
\end{verbatim}

\begin{center}
\includegraphics[width=0.8\linewidth,height=\textheight,keepaspectratio]{12-CategorizationModelsPrototypes_files/figure-pdf/ch13_multilevel_posterior_viz-1.pdf}
\end{center}

\textbf{Reading the population posterior.} The red dashed lines mark the
true parameter values used to generate the data. Well-calibrated
inference should show the posterior clearly shifted from the prior and
centred near the true value. With 20 subjects each contributing 64
trials, the population mean is well-identified; the population SD
requires more subjects before its posterior narrows substantially (see
§``Multilevel SBC'' below for the coverage analysis).

\subsection{Cross-Scenario Population
Recovery}\label{cross-scenario-population-recovery}

Mirroring Chapter 11's three-population analysis, we simulate one
population per scenario and verify that the multilevel prototype model
recovers population parameters with the same scenario-dependence
observed at the single-subject level. The key prediction:
\texttt{pop\_log\_r\_mean} should be well-recovered in all three
scenarios (observation noise directly shapes every per-trial
likelihood); \texttt{pop\_log\_q\_mean} should be poorly recovered in
static data (no drift to track) and cleanly recovered in the
smooth-drift scenario (assumption-matched generative process).

\begin{Shaded}
\begin{Highlighting}[]
\NormalTok{simulate\_ml\_proto\_scenario }\OtherTok{\textless{}{-}} \ControlFlowTok{function}\NormalTok{(}
\NormalTok{    n\_subjects, stimulus\_info, n\_blocks, scenario, base\_seed,}
\NormalTok{    pop\_log\_r\_mean, pop\_log\_r\_sd,}
    \AttributeTok{pop\_log\_q\_mean =}\NormalTok{ LOG\_Q\_PRIOR\_MEAN,}
    \AttributeTok{pop\_log\_q\_sd   =} \FloatTok{0.5}\NormalTok{) \{}

  \FunctionTok{set.seed}\NormalTok{(base\_seed)}
\NormalTok{  z\_log\_r    }\OtherTok{\textless{}{-}} \FunctionTok{rnorm}\NormalTok{(n\_subjects, }\DecValTok{0}\NormalTok{, }\DecValTok{1}\NormalTok{)}
\NormalTok{  subj\_log\_r }\OtherTok{\textless{}{-}}\NormalTok{ pop\_log\_r\_mean }\SpecialCharTok{+}\NormalTok{ pop\_log\_r\_sd }\SpecialCharTok{*}\NormalTok{ z\_log\_r}
\NormalTok{  subj\_r     }\OtherTok{\textless{}{-}} \FunctionTok{exp}\NormalTok{(subj\_log\_r)}

\NormalTok{  z\_log\_q    }\OtherTok{\textless{}{-}} \FunctionTok{rnorm}\NormalTok{(n\_subjects, }\DecValTok{0}\NormalTok{, }\DecValTok{1}\NormalTok{)}
\NormalTok{  subj\_log\_q }\OtherTok{\textless{}{-}}\NormalTok{ pop\_log\_q\_mean }\SpecialCharTok{+}\NormalTok{ pop\_log\_q\_sd }\SpecialCharTok{*}\NormalTok{ z\_log\_q}
\NormalTok{  subj\_q     }\OtherTok{\textless{}{-}} \FunctionTok{exp}\NormalTok{(subj\_log\_q)}

\NormalTok{  true\_params }\OtherTok{\textless{}{-}} \FunctionTok{tibble}\NormalTok{(}
    \AttributeTok{agent\_id   =} \FunctionTok{seq\_len}\NormalTok{(n\_subjects),}
    \AttributeTok{scenario   =}\NormalTok{ scenario,}
    \AttributeTok{log\_r\_true =}\NormalTok{ subj\_log\_r, }\AttributeTok{r\_true =}\NormalTok{ subj\_r,}
    \AttributeTok{log\_q\_true =}\NormalTok{ subj\_log\_q, }\AttributeTok{q\_true =}\NormalTok{ subj\_q}
\NormalTok{  )}

\NormalTok{  sim\_data }\OtherTok{\textless{}{-}} \FunctionTok{map\_dfr}\NormalTok{(}\FunctionTok{seq\_len}\NormalTok{(n\_subjects), }\ControlFlowTok{function}\NormalTok{(j) \{}
    \FunctionTok{simulate\_prototype\_agent\_scenario}\NormalTok{(}
      \AttributeTok{agent\_id     =}\NormalTok{ j, }\AttributeTok{scenario =}\NormalTok{ scenario,}
      \AttributeTok{r\_value      =}\NormalTok{ subj\_r[j], }\AttributeTok{q\_value =}\NormalTok{ subj\_q[j],}
      \AttributeTok{subject\_seed =}\NormalTok{ base\_seed }\SpecialCharTok{*} \DecValTok{100} \SpecialCharTok{+}\NormalTok{ j}
\NormalTok{    )}
\NormalTok{  \})}

  \FunctionTok{list}\NormalTok{(}\AttributeTok{data =}\NormalTok{ sim\_data, }\AttributeTok{true\_params =}\NormalTok{ true\_params)}
\NormalTok{\}}
\end{Highlighting}
\end{Shaded}

\begin{Shaded}
\begin{Highlighting}[]
\NormalTok{n\_subjects\_ml\_scen }\OtherTok{\textless{}{-}} \DecValTok{10}
\NormalTok{pop\_params\_ml\_scen }\OtherTok{\textless{}{-}} \FunctionTok{list}\NormalTok{(}
  \AttributeTok{pop\_log\_r\_mean =} \FunctionTok{log}\NormalTok{(}\FloatTok{1.5}\NormalTok{),    }\CommentTok{\# population median r = 1.5}
  \AttributeTok{pop\_log\_r\_sd   =} \FloatTok{0.5}\NormalTok{,}
  \AttributeTok{pop\_log\_q\_mean =} \FunctionTok{log}\NormalTok{(}\FloatTok{0.15}\NormalTok{),   }\CommentTok{\# population median q = 0.15}
  \AttributeTok{pop\_log\_q\_sd   =} \FloatTok{0.5}
\NormalTok{)}
\NormalTok{ml\_proto\_scenarios }\OtherTok{\textless{}{-}} \FunctionTok{c}\NormalTok{(}\StringTok{"static"}\NormalTok{, }\StringTok{"contingent\_shift"}\NormalTok{, }\StringTok{"drift"}\NormalTok{)}

\NormalTok{ml\_proto\_scen\_files }\OtherTok{\textless{}{-}} \FunctionTok{map}\NormalTok{(ml\_proto\_scenarios, }\ControlFlowTok{function}\NormalTok{(scen) }\FunctionTok{list}\NormalTok{(}
  \AttributeTok{data  =} \FunctionTok{here}\NormalTok{(}\StringTok{"simdata"}\NormalTok{, }\FunctionTok{paste0}\NormalTok{(}\StringTok{"ch13\_proto\_ml\_scen\_data\_"}\NormalTok{,  scen, }\StringTok{".csv"}\NormalTok{)),}
  \AttributeTok{truth =} \FunctionTok{here}\NormalTok{(}\StringTok{"simdata"}\NormalTok{, }\FunctionTok{paste0}\NormalTok{(}\StringTok{"ch13\_proto\_ml\_scen\_truth\_"}\NormalTok{, scen, }\StringTok{".csv"}\NormalTok{))}
\NormalTok{)) }\SpecialCharTok{|\textgreater{}} \FunctionTok{setNames}\NormalTok{(ml\_proto\_scenarios)}

\NormalTok{ml\_scen\_files\_exist }\OtherTok{\textless{}{-}} \FunctionTok{all}\NormalTok{(}
  \FunctionTok{file.exists}\NormalTok{(}\FunctionTok{unlist}\NormalTok{(}\FunctionTok{map}\NormalTok{(ml\_proto\_scen\_files, }\SpecialCharTok{\textasciitilde{}} \FunctionTok{c}\NormalTok{(.x}\SpecialCharTok{$}\NormalTok{data, .x}\SpecialCharTok{$}\NormalTok{truth))))}
\NormalTok{)}

\ControlFlowTok{if}\NormalTok{ (regenerate\_simulations }\SpecialCharTok{||} \SpecialCharTok{!}\NormalTok{ml\_scen\_files\_exist) \{}
\NormalTok{  ml\_proto\_scen\_sims }\OtherTok{\textless{}{-}} \FunctionTok{map}\NormalTok{(ml\_proto\_scenarios, }\ControlFlowTok{function}\NormalTok{(scen) \{}
\NormalTok{    out }\OtherTok{\textless{}{-}} \FunctionTok{simulate\_ml\_proto\_scenario}\NormalTok{(}
      \AttributeTok{n\_subjects     =}\NormalTok{ n\_subjects\_ml\_scen,}
      \AttributeTok{stimulus\_info  =}\NormalTok{ stimulus\_info,}
      \AttributeTok{n\_blocks       =}\NormalTok{ n\_blocks,}
      \AttributeTok{scenario       =}\NormalTok{ scen,}
      \AttributeTok{base\_seed      =} \FunctionTok{match}\NormalTok{(scen, ml\_proto\_scenarios) }\SpecialCharTok{*} \DecValTok{1000} \SpecialCharTok{+} \DecValTok{77}\NormalTok{,}
      \AttributeTok{pop\_log\_r\_mean =}\NormalTok{ pop\_params\_ml\_scen}\SpecialCharTok{$}\NormalTok{pop\_log\_r\_mean,}
      \AttributeTok{pop\_log\_r\_sd   =}\NormalTok{ pop\_params\_ml\_scen}\SpecialCharTok{$}\NormalTok{pop\_log\_r\_sd,}
      \AttributeTok{pop\_log\_q\_mean =}\NormalTok{ pop\_params\_ml\_scen}\SpecialCharTok{$}\NormalTok{pop\_log\_q\_mean,}
      \AttributeTok{pop\_log\_q\_sd   =}\NormalTok{ pop\_params\_ml\_scen}\SpecialCharTok{$}\NormalTok{pop\_log\_q\_sd}
\NormalTok{    )}
    \FunctionTok{write\_csv}\NormalTok{(out}\SpecialCharTok{$}\NormalTok{data,        ml\_proto\_scen\_files[[scen]]}\SpecialCharTok{$}\NormalTok{data)}
    \FunctionTok{write\_csv}\NormalTok{(out}\SpecialCharTok{$}\NormalTok{true\_params, ml\_proto\_scen\_files[[scen]]}\SpecialCharTok{$}\NormalTok{truth)}
\NormalTok{    out}
\NormalTok{  \}) }\SpecialCharTok{|\textgreater{}} \FunctionTok{setNames}\NormalTok{(ml\_proto\_scenarios)}
\NormalTok{\} }\ControlFlowTok{else}\NormalTok{ \{}
\NormalTok{  ml\_proto\_scen\_sims }\OtherTok{\textless{}{-}} \FunctionTok{map}\NormalTok{(ml\_proto\_scenarios, }\ControlFlowTok{function}\NormalTok{(scen) }\FunctionTok{list}\NormalTok{(}
    \AttributeTok{data        =} \FunctionTok{read\_csv}\NormalTok{(ml\_proto\_scen\_files[[scen]]}\SpecialCharTok{$}\NormalTok{data,  }\AttributeTok{show\_col\_types =} \ConstantTok{FALSE}\NormalTok{),}
    \AttributeTok{true\_params =} \FunctionTok{read\_csv}\NormalTok{(ml\_proto\_scen\_files[[scen]]}\SpecialCharTok{$}\NormalTok{truth, }\AttributeTok{show\_col\_types =} \ConstantTok{FALSE}\NormalTok{)}
\NormalTok{  )) }\SpecialCharTok{|\textgreater{}} \FunctionTok{setNames}\NormalTok{(ml\_proto\_scenarios)}
\NormalTok{\}}

\NormalTok{ml\_proto\_scen\_all }\OtherTok{\textless{}{-}} \FunctionTok{map\_dfr}\NormalTok{(ml\_proto\_scenarios, }\ControlFlowTok{function}\NormalTok{(scen)}
\NormalTok{  ml\_proto\_scen\_sims[[scen]]}\SpecialCharTok{$}\NormalTok{data }\SpecialCharTok{|\textgreater{}} \FunctionTok{mutate}\NormalTok{(}\AttributeTok{scenario =}\NormalTok{ scen)}
\NormalTok{) }\SpecialCharTok{|\textgreater{}}
  \FunctionTok{mutate}\NormalTok{(}\AttributeTok{scenario =} \FunctionTok{factor}\NormalTok{(scenario, }\AttributeTok{levels =}\NormalTok{ ml\_proto\_scenarios,}
                           \AttributeTok{labels =} \FunctionTok{c}\NormalTok{(}\StringTok{"Static"}\NormalTok{, }\StringTok{"Contingent shift"}\NormalTok{, }\StringTok{"Smooth drift"}\NormalTok{)))}

\FunctionTok{ggplot}\NormalTok{(ml\_proto\_scen\_all,}
       \FunctionTok{aes}\NormalTok{(}\AttributeTok{x =}\NormalTok{ trial\_within\_subject, }\AttributeTok{y =}\NormalTok{ performance, }\AttributeTok{group =}\NormalTok{ agent\_id)) }\SpecialCharTok{+}
  \FunctionTok{geom\_line}\NormalTok{(}\AttributeTok{alpha =} \FloatTok{0.25}\NormalTok{, }\AttributeTok{color =} \StringTok{"grey40"}\NormalTok{) }\SpecialCharTok{+}
  \FunctionTok{stat\_summary}\NormalTok{(}\AttributeTok{fun =}\NormalTok{ mean, }\AttributeTok{geom =} \StringTok{"line"}\NormalTok{, }\FunctionTok{aes}\NormalTok{(}\AttributeTok{group =} \DecValTok{1}\NormalTok{),}
               \AttributeTok{color =} \StringTok{"\#009E73"}\NormalTok{, }\AttributeTok{linewidth =} \FloatTok{1.4}\NormalTok{) }\SpecialCharTok{+}
  \FunctionTok{facet\_wrap}\NormalTok{(}\SpecialCharTok{\textasciitilde{}}\NormalTok{scenario, }\AttributeTok{ncol =} \DecValTok{3}\NormalTok{) }\SpecialCharTok{+}
  \FunctionTok{geom\_hline}\NormalTok{(}\AttributeTok{yintercept =} \FloatTok{0.5}\NormalTok{, }\AttributeTok{linetype =} \StringTok{"dashed"}\NormalTok{, }\AttributeTok{color =} \StringTok{"grey50"}\NormalTok{) }\SpecialCharTok{+}
  \FunctionTok{labs}\NormalTok{(}
    \AttributeTok{title    =} \StringTok{"Multilevel Prototype Model: learning curves across scenarios"}\NormalTok{,}
    \AttributeTok{subtitle =} \FunctionTok{paste0}\NormalTok{(n\_subjects\_ml\_scen,}
                      \StringTok{" agents per scenario; grey = individuals, green = population mean"}\NormalTok{),}
    \AttributeTok{x =} \StringTok{"Trial"}\NormalTok{, }\AttributeTok{y =} \StringTok{"Cumulative accuracy"}
\NormalTok{  )}
\end{Highlighting}
\end{Shaded}

\begin{center}
\includegraphics[width=0.8\linewidth,height=\textheight,keepaspectratio]{12-CategorizationModelsPrototypes_files/figure-pdf/ch13_ml_scenario_simulate-1.pdf}
\end{center}

\begin{Shaded}
\begin{Highlighting}[]
\NormalTok{build\_ml\_proto\_stan\_data }\OtherTok{\textless{}{-}} \ControlFlowTok{function}\NormalTok{(ml\_data) \{}
\NormalTok{  sorted }\OtherTok{\textless{}{-}}\NormalTok{ ml\_data }\SpecialCharTok{|\textgreater{}}
    \FunctionTok{mutate}\NormalTok{(}\AttributeTok{subj\_id\_stan =} \FunctionTok{as.integer}\NormalTok{(}\FunctionTok{factor}\NormalTok{(agent\_id))) }\SpecialCharTok{|\textgreater{}}
    \FunctionTok{arrange}\NormalTok{(subj\_id\_stan, trial\_within\_subject)}

\NormalTok{  bounds }\OtherTok{\textless{}{-}}\NormalTok{ sorted }\SpecialCharTok{|\textgreater{}}
    \FunctionTok{mutate}\NormalTok{(}\AttributeTok{row\_idx =} \FunctionTok{row\_number}\NormalTok{()) }\SpecialCharTok{|\textgreater{}}
    \FunctionTok{group\_by}\NormalTok{(subj\_id\_stan) }\SpecialCharTok{|\textgreater{}}
    \FunctionTok{summarise}\NormalTok{(}\AttributeTok{subj\_start =} \FunctionTok{min}\NormalTok{(row\_idx), }\AttributeTok{subj\_end =} \FunctionTok{max}\NormalTok{(row\_idx),}
              \AttributeTok{.groups =} \StringTok{"drop"}\NormalTok{) }\SpecialCharTok{|\textgreater{}}
    \FunctionTok{arrange}\NormalTok{(subj\_id\_stan)}

  \FunctionTok{list}\NormalTok{(}
    \AttributeTok{N\_total                   =} \FunctionTok{nrow}\NormalTok{(sorted),}
    \AttributeTok{N\_subjects                =} \FunctionTok{max}\NormalTok{(sorted}\SpecialCharTok{$}\NormalTok{subj\_id\_stan),}
    \AttributeTok{N\_features                =} \DecValTok{2}\NormalTok{L,}
    \AttributeTok{subj\_start                =}\NormalTok{ bounds}\SpecialCharTok{$}\NormalTok{subj\_start,}
    \AttributeTok{subj\_end                  =}\NormalTok{ bounds}\SpecialCharTok{$}\NormalTok{subj\_end,}
    \AttributeTok{y                         =} \FunctionTok{as.integer}\NormalTok{(sorted}\SpecialCharTok{$}\NormalTok{sim\_response),}
    \AttributeTok{cat\_one                   =} \FunctionTok{as.integer}\NormalTok{(sorted}\SpecialCharTok{$}\NormalTok{category\_feedback),}
    \AttributeTok{obs                       =} \FunctionTok{as.matrix}\NormalTok{(sorted[, }\FunctionTok{c}\NormalTok{(}\StringTok{"height"}\NormalTok{, }\StringTok{"position"}\NormalTok{)]),}
    \AttributeTok{initial\_mu\_cat0           =} \FunctionTok{c}\NormalTok{(}\FloatTok{2.5}\NormalTok{, }\FloatTok{2.5}\NormalTok{),}
    \AttributeTok{initial\_mu\_cat1           =} \FunctionTok{c}\NormalTok{(}\FloatTok{2.5}\NormalTok{, }\FloatTok{2.5}\NormalTok{),}
    \AttributeTok{initial\_sigma\_diag        =} \FloatTok{10.0}\NormalTok{,}
    \AttributeTok{pop\_log\_r\_mean\_prior\_mean =}\NormalTok{ LOG\_R\_PRIOR\_MEAN,}
    \AttributeTok{pop\_log\_r\_mean\_prior\_sd   =}\NormalTok{ LOG\_R\_PRIOR\_SD,}
    \AttributeTok{pop\_log\_r\_sd\_prior\_rate   =} \FloatTok{2.0}\NormalTok{,}
    \AttributeTok{pop\_log\_q\_mean\_prior\_mean =}\NormalTok{ LOG\_Q\_PRIOR\_MEAN,}
    \AttributeTok{pop\_log\_q\_mean\_prior\_sd   =}\NormalTok{ LOG\_Q\_PRIOR\_SD,}
    \AttributeTok{pop\_log\_q\_sd\_prior\_rate   =} \FloatTok{2.0}
\NormalTok{  )}
\NormalTok{\}}

\NormalTok{ml\_proto\_scen\_fit\_files }\OtherTok{\textless{}{-}} \FunctionTok{map\_chr}\NormalTok{(ml\_proto\_scenarios, }\ControlFlowTok{function}\NormalTok{(scen)}
  \FunctionTok{here}\NormalTok{(}\StringTok{"simmodels"}\NormalTok{, }\FunctionTok{paste0}\NormalTok{(}\StringTok{"ch13\_proto\_ml\_scen\_fit\_"}\NormalTok{, scen, }\StringTok{".rds"}\NormalTok{))}
\NormalTok{) }\SpecialCharTok{|\textgreater{}} \FunctionTok{setNames}\NormalTok{(ml\_proto\_scenarios)}

\ControlFlowTok{if}\NormalTok{ (regenerate\_simulations }\SpecialCharTok{||} \SpecialCharTok{!}\FunctionTok{all}\NormalTok{(}\FunctionTok{file.exists}\NormalTok{(ml\_proto\_scen\_fit\_files))) \{}
\NormalTok{  ml\_proto\_scen\_fits }\OtherTok{\textless{}{-}} \FunctionTok{map}\NormalTok{(ml\_proto\_scenarios, }\ControlFlowTok{function}\NormalTok{(scen) \{}
    \FunctionTok{cat}\NormalTok{(}\StringTok{"Fitting multilevel prototype — scenario:"}\NormalTok{, scen, }\StringTok{"}\SpecialCharTok{\textbackslash{}n}\StringTok{"}\NormalTok{)}
\NormalTok{    dat }\OtherTok{\textless{}{-}} \FunctionTok{build\_ml\_proto\_stan\_data}\NormalTok{(ml\_proto\_scen\_sims[[scen]]}\SpecialCharTok{$}\NormalTok{data)}
\NormalTok{    fit }\OtherTok{\textless{}{-}}\NormalTok{ mod\_prototype\_ml}\SpecialCharTok{$}\FunctionTok{sample}\NormalTok{(}
      \AttributeTok{data            =}\NormalTok{ dat,}
      \AttributeTok{chains          =} \DecValTok{4}\NormalTok{,}
      \AttributeTok{parallel\_chains =} \FunctionTok{min}\NormalTok{(}\DecValTok{4}\NormalTok{, }\FunctionTok{availableCores}\NormalTok{()),}
      \AttributeTok{iter\_warmup     =} \DecValTok{1500}\NormalTok{,}
      \AttributeTok{iter\_sampling   =} \DecValTok{1500}\NormalTok{,}
      \AttributeTok{refresh         =} \DecValTok{200}\NormalTok{,}
      \AttributeTok{adapt\_delta     =} \FloatTok{0.95}\NormalTok{,}
      \AttributeTok{max\_treedepth   =} \DecValTok{12}
\NormalTok{    )}
\NormalTok{    fit}\SpecialCharTok{$}\FunctionTok{save\_object}\NormalTok{(ml\_proto\_scen\_fit\_files[[scen]])}
\NormalTok{    fit}
\NormalTok{  \}) }\SpecialCharTok{|\textgreater{}} \FunctionTok{setNames}\NormalTok{(ml\_proto\_scenarios)}
\NormalTok{\} }\ControlFlowTok{else}\NormalTok{ \{}
\NormalTok{  ml\_proto\_scen\_fits }\OtherTok{\textless{}{-}} \FunctionTok{map}\NormalTok{(ml\_proto\_scen\_fit\_files, readRDS) }\SpecialCharTok{|\textgreater{}}
    \FunctionTok{setNames}\NormalTok{(ml\_proto\_scenarios)}
\NormalTok{\}}
\end{Highlighting}
\end{Shaded}

\begin{Shaded}
\begin{Highlighting}[]
\NormalTok{pop\_vars\_scen }\OtherTok{\textless{}{-}} \FunctionTok{c}\NormalTok{(}\StringTok{"pop\_log\_r\_mean"}\NormalTok{, }\StringTok{"pop\_log\_r\_sd"}\NormalTok{,}
                   \StringTok{"pop\_log\_q\_mean"}\NormalTok{, }\StringTok{"pop\_log\_q\_sd"}\NormalTok{)}

\NormalTok{true\_vals\_scen }\OtherTok{\textless{}{-}} \FunctionTok{tibble}\NormalTok{(}
  \AttributeTok{variable =}\NormalTok{ pop\_vars\_scen,}
  \AttributeTok{true\_val =} \FunctionTok{c}\NormalTok{(}
\NormalTok{    pop\_params\_ml\_scen}\SpecialCharTok{$}\NormalTok{pop\_log\_r\_mean,}
\NormalTok{    pop\_params\_ml\_scen}\SpecialCharTok{$}\NormalTok{pop\_log\_r\_sd,}
\NormalTok{    pop\_params\_ml\_scen}\SpecialCharTok{$}\NormalTok{pop\_log\_q\_mean,}
\NormalTok{    pop\_params\_ml\_scen}\SpecialCharTok{$}\NormalTok{pop\_log\_q\_sd}
\NormalTok{  )}
\NormalTok{)}

\NormalTok{ml\_proto\_scen\_recovery }\OtherTok{\textless{}{-}} \FunctionTok{map\_dfr}\NormalTok{(ml\_proto\_scenarios, }\ControlFlowTok{function}\NormalTok{(scen) \{}
\NormalTok{  ml\_proto\_scen\_fits[[scen]]}\SpecialCharTok{$}\FunctionTok{summary}\NormalTok{(}\AttributeTok{variables =}\NormalTok{ pop\_vars\_scen) }\SpecialCharTok{|\textgreater{}}
\NormalTok{    dplyr}\SpecialCharTok{::}\FunctionTok{select}\NormalTok{(variable, mean, q5, q95) }\SpecialCharTok{|\textgreater{}}
    \FunctionTok{mutate}\NormalTok{(}\AttributeTok{scenario =}\NormalTok{ scen)}
\NormalTok{\}) }\SpecialCharTok{|\textgreater{}}
  \FunctionTok{left\_join}\NormalTok{(true\_vals\_scen, }\AttributeTok{by =} \StringTok{"variable"}\NormalTok{) }\SpecialCharTok{|\textgreater{}}
  \FunctionTok{mutate}\NormalTok{(}\AttributeTok{scenario =} \FunctionTok{factor}\NormalTok{(scenario, }\AttributeTok{levels =}\NormalTok{ ml\_proto\_scenarios,}
                           \AttributeTok{labels =} \FunctionTok{c}\NormalTok{(}\StringTok{"Static"}\NormalTok{, }\StringTok{"Contingent shift"}\NormalTok{, }\StringTok{"Smooth drift"}\NormalTok{)))}

\FunctionTok{ggplot}\NormalTok{(ml\_proto\_scen\_recovery,}
       \FunctionTok{aes}\NormalTok{(}\AttributeTok{x =}\NormalTok{ scenario, }\AttributeTok{y =}\NormalTok{ mean, }\AttributeTok{ymin =}\NormalTok{ q5, }\AttributeTok{ymax =}\NormalTok{ q95, }\AttributeTok{color =}\NormalTok{ scenario)) }\SpecialCharTok{+}
  \FunctionTok{geom\_pointrange}\NormalTok{() }\SpecialCharTok{+}
  \FunctionTok{geom\_hline}\NormalTok{(}\FunctionTok{aes}\NormalTok{(}\AttributeTok{yintercept =}\NormalTok{ true\_val), }\AttributeTok{linetype =} \StringTok{"dashed"}\NormalTok{, }\AttributeTok{color =} \StringTok{"red"}\NormalTok{) }\SpecialCharTok{+}
  \FunctionTok{facet\_wrap}\NormalTok{(}\SpecialCharTok{\textasciitilde{}}\NormalTok{variable, }\AttributeTok{scales =} \StringTok{"free\_y"}\NormalTok{, }\AttributeTok{ncol =} \DecValTok{2}\NormalTok{) }\SpecialCharTok{+}
  \FunctionTok{scale\_color\_manual}\NormalTok{(}\AttributeTok{values =} \FunctionTok{c}\NormalTok{(}\StringTok{"\#999999"}\NormalTok{, }\StringTok{"\#E69F00"}\NormalTok{, }\StringTok{"\#009E73"}\NormalTok{)) }\SpecialCharTok{+}
  \FunctionTok{labs}\NormalTok{(}
    \AttributeTok{title    =} \StringTok{"Multilevel Population Recovery Across Scenarios"}\NormalTok{,}
    \AttributeTok{subtitle =} \StringTok{"Points = posterior mean; bars = 90\% CI; red dashes = true generating values"}\NormalTok{,}
    \AttributeTok{x =} \ConstantTok{NULL}\NormalTok{, }\AttributeTok{y =} \StringTok{"Posterior estimate"}
\NormalTok{  ) }\SpecialCharTok{+}
  \FunctionTok{theme}\NormalTok{(}\AttributeTok{legend.position =} \StringTok{"none"}\NormalTok{,}
        \AttributeTok{axis.text.x =} \FunctionTok{element\_text}\NormalTok{(}\AttributeTok{angle =} \DecValTok{30}\NormalTok{, }\AttributeTok{hjust =} \DecValTok{1}\NormalTok{))}
\end{Highlighting}
\end{Shaded}

\begin{center}
\includegraphics[width=0.8\linewidth,height=\textheight,keepaspectratio]{12-CategorizationModelsPrototypes_files/figure-pdf/ch13_ml_scenario_recovery-1.pdf}
\end{center}

\textbf{Cross-scenario population recovery.} \texttt{pop\_log\_r\_mean}
is well-recovered in all three scenarios: observation noise directly
shapes every per-trial likelihood, so the data always constrain \(r\) at
the population level. \texttt{pop\_log\_q\_mean} tells a different
story. On static data the 90\% CI is wide and centred near the prior
mean --- partial pooling across 10 subjects cannot manufacture
information that no individual trial provides. On contingent-shift data
the CI narrows as reversal dynamics provide a collective signal;
subjects with more reversals lend strength to those with fewer. On drift
data both population means are cleanly recovered, with the 90\% CI
sitting close to the true value --- the assumption-matched scenario
scales cleanly from the single-subject to the multilevel level.

\textbf{Comparison with Chapter 11.} The decay GCM's cross-scenario
multilevel analysis showed that \texttt{pop\_log\_lambda\_mean} was most
clearly recovered on \emph{contingent-shift} data. Here
\texttt{pop\_log\_q\_mean} is most clearly recovered on \emph{drift}
data. The contrast reflects the models' different generative stories:
the GCM's general-purpose decay kernel captures contingent shifts as
well as drift (forgetting is useful whenever old evidence is stale,
regardless of \emph{why} it became stale); the Kalman filter's
smooth-drift prior is most consistent with the scenario where prototypes
literally undergo Gaussian random-walk drift.

\begin{center}\rule{0.5\linewidth}{0.5pt}\end{center}

\section{Multilevel SBC}\label{multilevel-sbc}

Multilevel SBC checks that the full posterior --- including both
population hyperparameters and individual-level parameters --- is
correctly calibrated across the entire joint prior. The prototype model
now has two per-subject parameters (\(\log r\) and \(\log q\)), giving
six ECDF-difference curves in the SBC output.

\textbf{What the key curves test:}

\begin{itemize}
\tightlist
\item
  \texttt{pop\_log\_r\_mean} / \texttt{pop\_log\_q\_mean} --- Are the
  population-mean estimators centred? A monotone drift indicates
  directional bias.
\item
  \texttt{pop\_log\_r\_sd} / \texttt{pop\_log\_q\_sd} --- Are the
  population-SD estimators well-calibrated? U-shaped rank histograms
  indicate the posterior is overconfident; humps indicate
  underconfidence. The NCP should prevent reverse-funnel issues.
\item
  \texttt{subj\_log\_r{[}1{]}} / \texttt{subj\_log\_q{[}1{]}} --- Do
  partial pooling and NCP produce valid individual-level credible
  intervals for both parameters?
\end{itemize}

\begin{Shaded}
\begin{Highlighting}[]
\NormalTok{sbc\_ml\_proto\_filepath }\OtherTok{\textless{}{-}} \FunctionTok{here}\NormalTok{(}\StringTok{"simdata"}\NormalTok{, }\StringTok{"ch13\_sbc\_proto\_ml\_results.rds"}\NormalTok{)}

\ControlFlowTok{if}\NormalTok{ (regenerate\_simulations }\SpecialCharTok{||} \SpecialCharTok{!}\FunctionTok{file.exists}\NormalTok{(sbc\_ml\_proto\_filepath)) \{}
\NormalTok{  n\_sbc\_iterations\_ml\_proto }\OtherTok{\textless{}{-}} \DecValTok{200}   \CommentTok{\# use \textgreater{}= 1000 for publication{-}quality SBC}
\NormalTok{  n\_subjects\_sbc\_proto      }\OtherTok{\textless{}{-}} \DecValTok{10}

\NormalTok{  proto\_ml\_sbc\_generator }\OtherTok{\textless{}{-}} \FunctionTok{SBC\_generator\_function}\NormalTok{(}
    \ControlFlowTok{function}\NormalTok{() \{}
\NormalTok{      pop\_log\_r\_mean }\OtherTok{\textless{}{-}} \FunctionTok{rnorm}\NormalTok{(}\DecValTok{1}\NormalTok{, LOG\_R\_PRIOR\_MEAN, LOG\_R\_PRIOR\_SD)}
\NormalTok{      pop\_log\_r\_sd   }\OtherTok{\textless{}{-}} \FunctionTok{rexp}\NormalTok{(}\DecValTok{1}\NormalTok{, }\FloatTok{2.0}\NormalTok{)}
\NormalTok{      pop\_log\_q\_mean }\OtherTok{\textless{}{-}} \FunctionTok{rnorm}\NormalTok{(}\DecValTok{1}\NormalTok{, LOG\_Q\_PRIOR\_MEAN, LOG\_Q\_PRIOR\_SD)}
\NormalTok{      pop\_log\_q\_sd   }\OtherTok{\textless{}{-}} \FunctionTok{rexp}\NormalTok{(}\DecValTok{1}\NormalTok{, }\FloatTok{2.0}\NormalTok{)}

\NormalTok{      z\_log\_r    }\OtherTok{\textless{}{-}} \FunctionTok{rnorm}\NormalTok{(n\_subjects\_sbc\_proto, }\DecValTok{0}\NormalTok{, }\DecValTok{1}\NormalTok{)}
\NormalTok{      subj\_log\_r }\OtherTok{\textless{}{-}}\NormalTok{ pop\_log\_r\_mean }\SpecialCharTok{+}\NormalTok{ pop\_log\_r\_sd }\SpecialCharTok{*}\NormalTok{ z\_log\_r}
\NormalTok{      subj\_r     }\OtherTok{\textless{}{-}} \FunctionTok{exp}\NormalTok{(subj\_log\_r)}

\NormalTok{      z\_log\_q    }\OtherTok{\textless{}{-}} \FunctionTok{rnorm}\NormalTok{(n\_subjects\_sbc\_proto, }\DecValTok{0}\NormalTok{, }\DecValTok{1}\NormalTok{)}
\NormalTok{      subj\_log\_q }\OtherTok{\textless{}{-}}\NormalTok{ pop\_log\_q\_mean }\SpecialCharTok{+}\NormalTok{ pop\_log\_q\_sd }\SpecialCharTok{*}\NormalTok{ z\_log\_q}
\NormalTok{      subj\_q     }\OtherTok{\textless{}{-}} \FunctionTok{exp}\NormalTok{(subj\_log\_q)}

\NormalTok{      all\_dat }\OtherTok{\textless{}{-}} \FunctionTok{map\_dfr}\NormalTok{(}\FunctionTok{seq\_len}\NormalTok{(n\_subjects\_sbc\_proto), }\ControlFlowTok{function}\NormalTok{(j) \{}
\NormalTok{        sched }\OtherTok{\textless{}{-}} \FunctionTok{make\_subject\_schedule}\NormalTok{(stimulus\_info, n\_blocks,}
                                       \AttributeTok{seed =} \FunctionTok{sample.int}\NormalTok{(}\FloatTok{1e8}\NormalTok{, }\DecValTok{1}\NormalTok{))}
\NormalTok{        res   }\OtherTok{\textless{}{-}} \FunctionTok{prototype\_kalman}\NormalTok{(}
          \AttributeTok{r\_value            =}\NormalTok{ subj\_r[j],}
          \AttributeTok{q\_value            =}\NormalTok{ subj\_q[j],}
          \AttributeTok{obs                =} \FunctionTok{as.matrix}\NormalTok{(sched[, }\FunctionTok{c}\NormalTok{(}\StringTok{"height"}\NormalTok{, }\StringTok{"position"}\NormalTok{)]),}
          \AttributeTok{cat\_one            =}\NormalTok{ sched}\SpecialCharTok{$}\NormalTok{category\_feedback,}
          \AttributeTok{initial\_mu         =} \FunctionTok{list}\NormalTok{(}\FunctionTok{rep}\NormalTok{(}\FloatTok{2.5}\NormalTok{, }\DecValTok{2}\NormalTok{), }\FunctionTok{rep}\NormalTok{(}\FloatTok{2.5}\NormalTok{, }\DecValTok{2}\NormalTok{)),}
          \AttributeTok{initial\_sigma\_diag =} \FloatTok{10.0}
\NormalTok{        )}
\NormalTok{        sched }\SpecialCharTok{|\textgreater{}} \FunctionTok{mutate}\NormalTok{(}\AttributeTok{agent\_id =}\NormalTok{ j, }\AttributeTok{sim\_response =}\NormalTok{ res}\SpecialCharTok{$}\NormalTok{sim\_response)}
\NormalTok{      \})}

\NormalTok{      sorted }\OtherTok{\textless{}{-}}\NormalTok{ all\_dat }\SpecialCharTok{|\textgreater{}}
        \FunctionTok{mutate}\NormalTok{(}\AttributeTok{subj\_id\_stan =} \FunctionTok{as.integer}\NormalTok{(}\FunctionTok{factor}\NormalTok{(agent\_id))) }\SpecialCharTok{|\textgreater{}}
        \FunctionTok{arrange}\NormalTok{(subj\_id\_stan, trial\_within\_subject)}

\NormalTok{      bounds }\OtherTok{\textless{}{-}}\NormalTok{ sorted }\SpecialCharTok{|\textgreater{}}
        \FunctionTok{mutate}\NormalTok{(}\AttributeTok{row\_idx =} \FunctionTok{row\_number}\NormalTok{()) }\SpecialCharTok{|\textgreater{}}
        \FunctionTok{group\_by}\NormalTok{(subj\_id\_stan) }\SpecialCharTok{|\textgreater{}}
        \FunctionTok{summarise}\NormalTok{(}\AttributeTok{subj\_start =} \FunctionTok{min}\NormalTok{(row\_idx), }\AttributeTok{subj\_end =} \FunctionTok{max}\NormalTok{(row\_idx),}
                  \AttributeTok{.groups =} \StringTok{"drop"}\NormalTok{) }\SpecialCharTok{|\textgreater{}}
        \FunctionTok{arrange}\NormalTok{(subj\_id\_stan)}

      \FunctionTok{list}\NormalTok{(}
        \AttributeTok{variables =} \FunctionTok{list}\NormalTok{(}
          \AttributeTok{pop\_log\_r\_mean    =}\NormalTok{ pop\_log\_r\_mean,}
          \AttributeTok{pop\_log\_r\_sd      =}\NormalTok{ pop\_log\_r\_sd,}
          \StringTok{\textasciigrave{}}\AttributeTok{subj\_log\_r[1]}\StringTok{\textasciigrave{}}   \OtherTok{=}\NormalTok{ subj\_log\_r[}\DecValTok{1}\NormalTok{],}
          \AttributeTok{pop\_log\_q\_mean    =}\NormalTok{ pop\_log\_q\_mean,}
          \AttributeTok{pop\_log\_q\_sd      =}\NormalTok{ pop\_log\_q\_sd,}
          \StringTok{\textasciigrave{}}\AttributeTok{subj\_log\_q[1]}\StringTok{\textasciigrave{}}   \OtherTok{=}\NormalTok{ subj\_log\_q[}\DecValTok{1}\NormalTok{]}
\NormalTok{        ),}
        \AttributeTok{generated =} \FunctionTok{list}\NormalTok{(}
          \AttributeTok{N\_total                   =} \FunctionTok{nrow}\NormalTok{(sorted),}
          \AttributeTok{N\_subjects                =}\NormalTok{ n\_subjects\_sbc\_proto,}
          \AttributeTok{N\_features                =} \DecValTok{2}\NormalTok{L,}
          \AttributeTok{subj\_start                =}\NormalTok{ bounds}\SpecialCharTok{$}\NormalTok{subj\_start,}
          \AttributeTok{subj\_end                  =}\NormalTok{ bounds}\SpecialCharTok{$}\NormalTok{subj\_end,}
          \AttributeTok{y                         =} \FunctionTok{as.integer}\NormalTok{(sorted}\SpecialCharTok{$}\NormalTok{sim\_response),}
          \AttributeTok{cat\_one                   =} \FunctionTok{as.integer}\NormalTok{(sorted}\SpecialCharTok{$}\NormalTok{category\_feedback),}
          \AttributeTok{obs                       =} \FunctionTok{as.matrix}\NormalTok{(sorted[, }\FunctionTok{c}\NormalTok{(}\StringTok{"height"}\NormalTok{, }\StringTok{"position"}\NormalTok{)]),}
          \AttributeTok{initial\_mu\_cat0           =} \FunctionTok{c}\NormalTok{(}\FloatTok{2.5}\NormalTok{, }\FloatTok{2.5}\NormalTok{),}
          \AttributeTok{initial\_mu\_cat1           =} \FunctionTok{c}\NormalTok{(}\FloatTok{2.5}\NormalTok{, }\FloatTok{2.5}\NormalTok{),}
          \AttributeTok{initial\_sigma\_diag        =} \FloatTok{10.0}\NormalTok{,}
          \AttributeTok{pop\_log\_r\_mean\_prior\_mean =}\NormalTok{ LOG\_R\_PRIOR\_MEAN,}
          \AttributeTok{pop\_log\_r\_mean\_prior\_sd   =}\NormalTok{ LOG\_R\_PRIOR\_SD,}
          \AttributeTok{pop\_log\_r\_sd\_prior\_rate   =} \FloatTok{2.0}\NormalTok{,}
          \AttributeTok{pop\_log\_q\_mean\_prior\_mean =}\NormalTok{ LOG\_Q\_PRIOR\_MEAN,}
          \AttributeTok{pop\_log\_q\_mean\_prior\_sd   =}\NormalTok{ LOG\_Q\_PRIOR\_SD,}
          \AttributeTok{pop\_log\_q\_sd\_prior\_rate   =} \FloatTok{2.0}
\NormalTok{        )}
\NormalTok{      )}
\NormalTok{    \}}
\NormalTok{  )}

\NormalTok{  proto\_ml\_sbc\_backend }\OtherTok{\textless{}{-}} \FunctionTok{SBC\_backend\_cmdstan\_sample}\NormalTok{(}
\NormalTok{    mod\_prototype\_ml,}
    \AttributeTok{iter\_warmup   =} \DecValTok{1000}\NormalTok{,}
    \AttributeTok{iter\_sampling =} \DecValTok{1000}\NormalTok{,}
    \AttributeTok{chains        =} \DecValTok{2}\NormalTok{,}
    \AttributeTok{adapt\_delta   =} \FloatTok{0.95}\NormalTok{,}
    \AttributeTok{refresh       =} \DecValTok{0}
\NormalTok{  )}

\NormalTok{  proto\_ml\_datasets }\OtherTok{\textless{}{-}} \FunctionTok{generate\_datasets}\NormalTok{(}
\NormalTok{    proto\_ml\_sbc\_generator,}
    \AttributeTok{n\_sims =}\NormalTok{ n\_sbc\_iterations\_ml\_proto}
\NormalTok{  )}

\NormalTok{  sbc\_results\_ml\_proto }\OtherTok{\textless{}{-}} \FunctionTok{compute\_SBC}\NormalTok{(}
\NormalTok{    proto\_ml\_datasets,}
\NormalTok{    proto\_ml\_sbc\_backend,}
    \AttributeTok{cache\_mode     =} \StringTok{"results"}\NormalTok{,}
    \AttributeTok{cache\_location =} \FunctionTok{here}\NormalTok{(}\StringTok{"simdata"}\NormalTok{, }\StringTok{"ch13\_sbc\_proto\_ml\_cache"}\NormalTok{),}
    \AttributeTok{keep\_fits      =} \ConstantTok{FALSE}
\NormalTok{  )}
  \FunctionTok{saveRDS}\NormalTok{(sbc\_results\_ml\_proto, sbc\_ml\_proto\_filepath)}
  \FunctionTok{cat}\NormalTok{(}\StringTok{"Multilevel prototype SBC results computed and saved.}\SpecialCharTok{\textbackslash{}n}\StringTok{"}\NormalTok{)}
\NormalTok{\} }\ControlFlowTok{else}\NormalTok{ \{}
\NormalTok{  sbc\_results\_ml\_proto }\OtherTok{\textless{}{-}} \FunctionTok{readRDS}\NormalTok{(sbc\_ml\_proto\_filepath)}
  \FunctionTok{cat}\NormalTok{(}\StringTok{"Loaded existing multilevel prototype SBC results.}\SpecialCharTok{\textbackslash{}n}\StringTok{"}\NormalTok{)}
\NormalTok{\}}
\end{Highlighting}
\end{Shaded}

\begin{verbatim}
Loaded existing multilevel prototype SBC results.
\end{verbatim}

\begin{Shaded}
\begin{Highlighting}[]
\FunctionTok{plot\_ecdf\_diff}\NormalTok{(sbc\_results\_ml\_proto)}
\end{Highlighting}
\end{Shaded}

\begin{center}
\includegraphics[width=0.8\linewidth,height=\textheight,keepaspectratio]{12-CategorizationModelsPrototypes_files/figure-pdf/ch13_sbc_multilevel-1.pdf}
\end{center}

\begin{Shaded}
\begin{Highlighting}[]
\FunctionTok{plot\_rank\_hist}\NormalTok{(sbc\_results\_ml\_proto)}
\end{Highlighting}
\end{Shaded}

\begin{center}
\includegraphics[width=0.8\linewidth,height=\textheight,keepaspectratio]{12-CategorizationModelsPrototypes_files/figure-pdf/ch13_sbc_multilevel-2.pdf}
\end{center}

\begin{Shaded}
\begin{Highlighting}[]
\CommentTok{\# ── Population{-}level parameter recovery across SBC simulations ───────────────}
\NormalTok{recovery\_df\_pop\_proto }\OtherTok{\textless{}{-}}\NormalTok{ sbc\_results\_ml\_proto}\SpecialCharTok{$}\NormalTok{stats }\SpecialCharTok{|\textgreater{}}
  \FunctionTok{filter}\NormalTok{(variable }\SpecialCharTok{\%in\%} \FunctionTok{c}\NormalTok{(}\StringTok{"pop\_log\_r\_mean"}\NormalTok{, }\StringTok{"pop\_log\_r\_sd"}\NormalTok{,}
                          \StringTok{"pop\_log\_q\_mean"}\NormalTok{, }\StringTok{"pop\_log\_q\_sd"}\NormalTok{))}

\NormalTok{p\_sbc\_pop\_proto }\OtherTok{\textless{}{-}} \FunctionTok{ggplot}\NormalTok{(recovery\_df\_pop\_proto,}
                           \FunctionTok{aes}\NormalTok{(}\AttributeTok{x =}\NormalTok{ simulated\_value, }\AttributeTok{y =}\NormalTok{ mean)) }\SpecialCharTok{+}
  \FunctionTok{geom\_point}\NormalTok{(}\AttributeTok{alpha =} \FloatTok{0.5}\NormalTok{, }\AttributeTok{color =} \StringTok{"\#009E73"}\NormalTok{) }\SpecialCharTok{+}
  \FunctionTok{geom\_abline}\NormalTok{(}\AttributeTok{intercept =} \DecValTok{0}\NormalTok{, }\AttributeTok{slope =} \DecValTok{1}\NormalTok{, }\AttributeTok{color =} \StringTok{"red"}\NormalTok{, }\AttributeTok{linetype =} \StringTok{"dashed"}\NormalTok{) }\SpecialCharTok{+}
  \FunctionTok{facet\_wrap}\NormalTok{(}\SpecialCharTok{\textasciitilde{}}\NormalTok{variable, }\AttributeTok{scales =} \StringTok{"free"}\NormalTok{) }\SpecialCharTok{+}
  \FunctionTok{labs}\NormalTok{(}
    \AttributeTok{title    =} \StringTok{"Population Parameter Recovery Across SBC Simulations"}\NormalTok{,}
    \AttributeTok{subtitle =} \StringTok{"Posterior mean vs. true generated value"}\NormalTok{,}
    \AttributeTok{x =} \StringTok{"True Value"}\NormalTok{, }\AttributeTok{y =} \StringTok{"Estimated Value (Posterior Mean)"}
\NormalTok{  ) }\SpecialCharTok{+}
  \FunctionTok{theme\_bw}\NormalTok{()}
\FunctionTok{print}\NormalTok{(p\_sbc\_pop\_proto)}
\end{Highlighting}
\end{Shaded}

\begin{center}
\includegraphics[width=0.8\linewidth,height=\textheight,keepaspectratio]{12-CategorizationModelsPrototypes_files/figure-pdf/ch13_sbc_multilevel-3.pdf}
\end{center}

\begin{Shaded}
\begin{Highlighting}[]
\CommentTok{\# ── Individual{-}level parameter recovery across SBC simulations ──────────────}
\NormalTok{recovery\_df\_indiv\_proto }\OtherTok{\textless{}{-}}\NormalTok{ sbc\_results\_ml\_proto}\SpecialCharTok{$}\NormalTok{stats }\SpecialCharTok{|\textgreater{}}
  \FunctionTok{filter}\NormalTok{(variable }\SpecialCharTok{\%in\%} \FunctionTok{c}\NormalTok{(}\StringTok{"subj\_log\_r[1]"}\NormalTok{, }\StringTok{"subj\_log\_q[1]"}\NormalTok{))}

\NormalTok{p\_sbc\_indiv\_proto }\OtherTok{\textless{}{-}} \FunctionTok{ggplot}\NormalTok{(recovery\_df\_indiv\_proto,}
                             \FunctionTok{aes}\NormalTok{(}\AttributeTok{x =}\NormalTok{ simulated\_value, }\AttributeTok{y =}\NormalTok{ mean)) }\SpecialCharTok{+}
  \FunctionTok{geom\_point}\NormalTok{(}\AttributeTok{alpha =} \FloatTok{0.5}\NormalTok{, }\AttributeTok{color =} \StringTok{"\#D55E00"}\NormalTok{) }\SpecialCharTok{+}
  \FunctionTok{geom\_abline}\NormalTok{(}\AttributeTok{intercept =} \DecValTok{0}\NormalTok{, }\AttributeTok{slope =} \DecValTok{1}\NormalTok{,}
              \AttributeTok{color =} \StringTok{"grey30"}\NormalTok{, }\AttributeTok{linetype =} \StringTok{"dashed"}\NormalTok{) }\SpecialCharTok{+}
  \FunctionTok{facet\_wrap}\NormalTok{(}\SpecialCharTok{\textasciitilde{}}\NormalTok{variable, }\AttributeTok{scales =} \StringTok{"free"}\NormalTok{) }\SpecialCharTok{+}
  \FunctionTok{labs}\NormalTok{(}
    \AttributeTok{title    =} \StringTok{"Individual Parameter Recovery (Subject 1) Across SBC"}\NormalTok{,}
    \AttributeTok{subtitle =} \StringTok{"Posterior mean vs. true generated value"}\NormalTok{,}
    \AttributeTok{x =} \StringTok{"True value [Subject 1]"}\NormalTok{,}
    \AttributeTok{y =} \StringTok{"Estimated value (Posterior Mean)"}
\NormalTok{  ) }\SpecialCharTok{+}
  \FunctionTok{theme\_bw}\NormalTok{()}
\FunctionTok{print}\NormalTok{(p\_sbc\_indiv\_proto)}
\end{Highlighting}
\end{Shaded}

\begin{center}
\includegraphics[width=0.8\linewidth,height=\textheight,keepaspectratio]{12-CategorizationModelsPrototypes_files/figure-pdf/ch13_sbc_multilevel-4.pdf}
\end{center}

\textbf{How to read the multilevel prototype SBC.} The prototype model
now has two per-subject parameters (\(\log r\) and \(\log q\)), giving
six ECDF-difference curves.

\begin{itemize}
\item
  \textbf{ECDF-difference plot and rank histograms.} All six curves
  (\texttt{pop\_log\_r\_mean}, \texttt{pop\_log\_r\_sd},
  \texttt{subj\_log\_r{[}1{]}}, \texttt{pop\_log\_q\_mean},
  \texttt{pop\_log\_q\_sd}, \texttt{subj\_log\_q{[}1{]}}) should remain
  inside the simultaneous grey confidence band. A U-shaped rank
  histogram for \texttt{pop\_log\_r\_sd} or \texttt{pop\_log\_q\_sd} is
  the canonical NCP failure mode. For \texttt{pop\_log\_q\_sd}
  especially, the posterior may be wide (q is partially confounded with
  r under the static paradigm), so coverage rather than width is the key
  diagnostic. If rank histograms show systematic failures, the standard
  remedy is to switch from NCP to centred parameterisation for the
  affected parameter.
\item
  \textbf{Scatter plots (population and individual recovery).} Points
  along the red dashed identity line confirm that the posterior means
  are unbiased estimators of the generating values. Systematic curvature
  (points arching away from the line at the extremes) would indicate
  either prior-boundary effects or insufficient data per subject.
\item
  \textbf{Comparison with single-subject SBC (§``Simulation-Based
  Calibration Checks'' above).} The single-subject SBC used \(S = 500\)
  iterations of one agent each; the multilevel SBC uses \(S = 200\)
  iterations of 10 agents each. The multilevel SBC is the authoritative
  check for multi-subject inference: it verifies that partial pooling
  does not distort individual-level posteriors and that population
  hyperparameters are jointly calibrated with per-subject parameters.
\end{itemize}

\subsection{Multilevel SBC on Drift Data: The Assumption-Matched
Case}\label{multilevel-sbc-on-drift-data-the-assumption-matched-case}

The SBC above uses static data, which deliberately exposes the
\texttt{pop\_log\_q} calibration failure --- the population-level
counterpart of the single-subject failure documented in
§``Simulation-Based Calibration Checks''. We now run a second multilevel
SBC using the \textbf{smooth-drift} scenario, where the Kalman filter's
generative story is literally true of the environment. This is the
assumption-matched multilevel check: if the model is correctly specified
\emph{and} the task provides drift signal at the individual level, the
population-level posterior should be well-calibrated across both
parameters.

\begin{Shaded}
\begin{Highlighting}[]
\NormalTok{sbc\_ml\_drift\_filepath }\OtherTok{\textless{}{-}} \FunctionTok{here}\NormalTok{(}\StringTok{"simdata"}\NormalTok{, }\StringTok{"ch13\_sbc\_proto\_ml\_drift\_results.rds"}\NormalTok{)}

\ControlFlowTok{if}\NormalTok{ (regenerate\_simulations }\SpecialCharTok{||} \SpecialCharTok{!}\FunctionTok{file.exists}\NormalTok{(sbc\_ml\_drift\_filepath)) \{}
\NormalTok{  n\_sbc\_iterations\_ml\_drift }\OtherTok{\textless{}{-}} \DecValTok{200}
\NormalTok{  n\_subjects\_sbc\_drift      }\OtherTok{\textless{}{-}} \DecValTok{10}

\NormalTok{  proto\_ml\_drift\_generator }\OtherTok{\textless{}{-}} \FunctionTok{SBC\_generator\_function}\NormalTok{(}
    \ControlFlowTok{function}\NormalTok{() \{}
\NormalTok{      pop\_log\_r\_mean }\OtherTok{\textless{}{-}} \FunctionTok{rnorm}\NormalTok{(}\DecValTok{1}\NormalTok{, LOG\_R\_PRIOR\_MEAN, LOG\_R\_PRIOR\_SD)}
\NormalTok{      pop\_log\_r\_sd   }\OtherTok{\textless{}{-}} \FunctionTok{rexp}\NormalTok{(}\DecValTok{1}\NormalTok{, }\FloatTok{2.0}\NormalTok{)}
\NormalTok{      pop\_log\_q\_mean }\OtherTok{\textless{}{-}} \FunctionTok{rnorm}\NormalTok{(}\DecValTok{1}\NormalTok{, LOG\_Q\_PRIOR\_MEAN, LOG\_Q\_PRIOR\_SD)}
\NormalTok{      pop\_log\_q\_sd   }\OtherTok{\textless{}{-}} \FunctionTok{rexp}\NormalTok{(}\DecValTok{1}\NormalTok{, }\FloatTok{2.0}\NormalTok{)}

\NormalTok{      z\_log\_r    }\OtherTok{\textless{}{-}} \FunctionTok{rnorm}\NormalTok{(n\_subjects\_sbc\_drift, }\DecValTok{0}\NormalTok{, }\DecValTok{1}\NormalTok{)}
\NormalTok{      subj\_log\_r }\OtherTok{\textless{}{-}}\NormalTok{ pop\_log\_r\_mean }\SpecialCharTok{+}\NormalTok{ pop\_log\_r\_sd }\SpecialCharTok{*}\NormalTok{ z\_log\_r}
\NormalTok{      subj\_r     }\OtherTok{\textless{}{-}} \FunctionTok{exp}\NormalTok{(subj\_log\_r)}

\NormalTok{      z\_log\_q    }\OtherTok{\textless{}{-}} \FunctionTok{rnorm}\NormalTok{(n\_subjects\_sbc\_drift, }\DecValTok{0}\NormalTok{, }\DecValTok{1}\NormalTok{)}
\NormalTok{      subj\_log\_q }\OtherTok{\textless{}{-}}\NormalTok{ pop\_log\_q\_mean }\SpecialCharTok{+}\NormalTok{ pop\_log\_q\_sd }\SpecialCharTok{*}\NormalTok{ z\_log\_q}
\NormalTok{      subj\_q     }\OtherTok{\textless{}{-}} \FunctionTok{exp}\NormalTok{(subj\_log\_q)}

\NormalTok{      all\_dat }\OtherTok{\textless{}{-}} \FunctionTok{map\_dfr}\NormalTok{(}\FunctionTok{seq\_len}\NormalTok{(n\_subjects\_sbc\_drift), }\ControlFlowTok{function}\NormalTok{(j) \{}
\NormalTok{        sched    }\OtherTok{\textless{}{-}} \FunctionTok{make\_subject\_schedule}\NormalTok{(stimulus\_info, n\_blocks,}
                                          \AttributeTok{seed =} \FunctionTok{sample.int}\NormalTok{(}\FloatTok{1e8}\NormalTok{, }\DecValTok{1}\NormalTok{))}
\NormalTok{        drift\_tr }\OtherTok{\textless{}{-}} \FunctionTok{make\_drift\_trajectory}\NormalTok{(}\FunctionTok{nrow}\NormalTok{(sched), }\AttributeTok{drift\_sigma =} \FloatTok{0.05}\NormalTok{,}
                                          \AttributeTok{seed =} \FunctionTok{sample.int}\NormalTok{(}\FloatTok{1e8}\NormalTok{, }\DecValTok{1}\NormalTok{))}
        \FunctionTok{simulate\_prototype\_scenario}\NormalTok{(}
          \AttributeTok{r\_value    =}\NormalTok{ subj\_r[j],}
          \AttributeTok{q\_value    =}\NormalTok{ subj\_q[j],}
          \AttributeTok{schedule   =}\NormalTok{ sched,}
          \AttributeTok{scenario   =} \StringTok{"drift"}\NormalTok{,}
          \AttributeTok{drift\_traj =}\NormalTok{ drift\_tr}
\NormalTok{        ) }\SpecialCharTok{|\textgreater{}} \FunctionTok{mutate}\NormalTok{(}\AttributeTok{agent\_id =}\NormalTok{ j)}
\NormalTok{      \})}

\NormalTok{      sorted }\OtherTok{\textless{}{-}}\NormalTok{ all\_dat }\SpecialCharTok{|\textgreater{}}
        \FunctionTok{mutate}\NormalTok{(}\AttributeTok{subj\_id\_stan =} \FunctionTok{as.integer}\NormalTok{(}\FunctionTok{factor}\NormalTok{(agent\_id))) }\SpecialCharTok{|\textgreater{}}
        \FunctionTok{arrange}\NormalTok{(subj\_id\_stan, trial\_within\_subject)}

\NormalTok{      bounds }\OtherTok{\textless{}{-}}\NormalTok{ sorted }\SpecialCharTok{|\textgreater{}}
        \FunctionTok{mutate}\NormalTok{(}\AttributeTok{row\_idx =} \FunctionTok{row\_number}\NormalTok{()) }\SpecialCharTok{|\textgreater{}}
        \FunctionTok{group\_by}\NormalTok{(subj\_id\_stan) }\SpecialCharTok{|\textgreater{}}
        \FunctionTok{summarise}\NormalTok{(}\AttributeTok{subj\_start =} \FunctionTok{min}\NormalTok{(row\_idx), }\AttributeTok{subj\_end =} \FunctionTok{max}\NormalTok{(row\_idx),}
                  \AttributeTok{.groups =} \StringTok{"drop"}\NormalTok{) }\SpecialCharTok{|\textgreater{}}
        \FunctionTok{arrange}\NormalTok{(subj\_id\_stan)}

      \FunctionTok{list}\NormalTok{(}
        \AttributeTok{variables =} \FunctionTok{list}\NormalTok{(}
          \AttributeTok{pop\_log\_r\_mean    =}\NormalTok{ pop\_log\_r\_mean,}
          \AttributeTok{pop\_log\_r\_sd      =}\NormalTok{ pop\_log\_r\_sd,}
          \StringTok{\textasciigrave{}}\AttributeTok{subj\_log\_r[1]}\StringTok{\textasciigrave{}}   \OtherTok{=}\NormalTok{ subj\_log\_r[}\DecValTok{1}\NormalTok{],}
          \AttributeTok{pop\_log\_q\_mean    =}\NormalTok{ pop\_log\_q\_mean,}
          \AttributeTok{pop\_log\_q\_sd      =}\NormalTok{ pop\_log\_q\_sd,}
          \StringTok{\textasciigrave{}}\AttributeTok{subj\_log\_q[1]}\StringTok{\textasciigrave{}}   \OtherTok{=}\NormalTok{ subj\_log\_q[}\DecValTok{1}\NormalTok{]}
\NormalTok{        ),}
        \AttributeTok{generated =} \FunctionTok{list}\NormalTok{(}
          \AttributeTok{N\_total                   =} \FunctionTok{nrow}\NormalTok{(sorted),}
          \AttributeTok{N\_subjects                =}\NormalTok{ n\_subjects\_sbc\_drift,}
          \AttributeTok{N\_features                =} \DecValTok{2}\NormalTok{L,}
          \AttributeTok{subj\_start                =}\NormalTok{ bounds}\SpecialCharTok{$}\NormalTok{subj\_start,}
          \AttributeTok{subj\_end                  =}\NormalTok{ bounds}\SpecialCharTok{$}\NormalTok{subj\_end,}
          \AttributeTok{y                         =} \FunctionTok{as.integer}\NormalTok{(sorted}\SpecialCharTok{$}\NormalTok{sim\_response),}
          \AttributeTok{cat\_one                   =} \FunctionTok{as.integer}\NormalTok{(sorted}\SpecialCharTok{$}\NormalTok{category\_feedback),}
          \AttributeTok{obs                       =} \FunctionTok{as.matrix}\NormalTok{(sorted[, }\FunctionTok{c}\NormalTok{(}\StringTok{"height"}\NormalTok{, }\StringTok{"position"}\NormalTok{)]),}
          \AttributeTok{initial\_mu\_cat0           =} \FunctionTok{c}\NormalTok{(}\FloatTok{2.5}\NormalTok{, }\FloatTok{2.5}\NormalTok{),}
          \AttributeTok{initial\_mu\_cat1           =} \FunctionTok{c}\NormalTok{(}\FloatTok{2.5}\NormalTok{, }\FloatTok{2.5}\NormalTok{),}
          \AttributeTok{initial\_sigma\_diag        =} \FloatTok{10.0}\NormalTok{,}
          \AttributeTok{pop\_log\_r\_mean\_prior\_mean =}\NormalTok{ LOG\_R\_PRIOR\_MEAN,}
          \AttributeTok{pop\_log\_r\_mean\_prior\_sd   =}\NormalTok{ LOG\_R\_PRIOR\_SD,}
          \AttributeTok{pop\_log\_r\_sd\_prior\_rate   =} \FloatTok{2.0}\NormalTok{,}
          \AttributeTok{pop\_log\_q\_mean\_prior\_mean =}\NormalTok{ LOG\_Q\_PRIOR\_MEAN,}
          \AttributeTok{pop\_log\_q\_mean\_prior\_sd   =}\NormalTok{ LOG\_Q\_PRIOR\_SD,}
          \AttributeTok{pop\_log\_q\_sd\_prior\_rate   =} \FloatTok{2.0}
\NormalTok{        )}
\NormalTok{      )}
\NormalTok{    \}}
\NormalTok{  )}

\NormalTok{  proto\_ml\_drift\_backend }\OtherTok{\textless{}{-}} \FunctionTok{SBC\_backend\_cmdstan\_sample}\NormalTok{(}
\NormalTok{    mod\_prototype\_ml,}
    \AttributeTok{iter\_warmup   =} \DecValTok{1000}\NormalTok{,}
    \AttributeTok{iter\_sampling =} \DecValTok{1000}\NormalTok{,}
    \AttributeTok{chains        =} \DecValTok{2}\NormalTok{,}
    \AttributeTok{adapt\_delta   =} \FloatTok{0.95}\NormalTok{,}
    \AttributeTok{refresh       =} \DecValTok{0}
\NormalTok{  )}

\NormalTok{  proto\_ml\_drift\_datasets }\OtherTok{\textless{}{-}} \FunctionTok{generate\_datasets}\NormalTok{(}
\NormalTok{    proto\_ml\_drift\_generator,}
    \AttributeTok{n\_sims =}\NormalTok{ n\_sbc\_iterations\_ml\_drift}
\NormalTok{  )}

\NormalTok{  sbc\_results\_ml\_drift }\OtherTok{\textless{}{-}} \FunctionTok{compute\_SBC}\NormalTok{(}
\NormalTok{    proto\_ml\_drift\_datasets,}
\NormalTok{    proto\_ml\_drift\_backend,}
    \AttributeTok{cache\_mode     =} \StringTok{"results"}\NormalTok{,}
    \AttributeTok{cache\_location =} \FunctionTok{here}\NormalTok{(}\StringTok{"simdata"}\NormalTok{, }\StringTok{"ch13\_sbc\_proto\_ml\_drift\_cache"}\NormalTok{),}
    \AttributeTok{keep\_fits      =} \ConstantTok{FALSE}
\NormalTok{  )}
  \FunctionTok{saveRDS}\NormalTok{(sbc\_results\_ml\_drift, sbc\_ml\_drift\_filepath)}
  \FunctionTok{cat}\NormalTok{(}\StringTok{"Multilevel prototype SBC (drift) results computed and saved.}\SpecialCharTok{\textbackslash{}n}\StringTok{"}\NormalTok{)}
\NormalTok{\} }\ControlFlowTok{else}\NormalTok{ \{}
\NormalTok{  sbc\_results\_ml\_drift }\OtherTok{\textless{}{-}} \FunctionTok{readRDS}\NormalTok{(sbc\_ml\_drift\_filepath)}
  \FunctionTok{cat}\NormalTok{(}\StringTok{"Loaded existing multilevel prototype SBC (drift) results.}\SpecialCharTok{\textbackslash{}n}\StringTok{"}\NormalTok{)}
\NormalTok{\}}
\end{Highlighting}
\end{Shaded}

\begin{verbatim}
Loaded existing multilevel prototype SBC (drift) results.
\end{verbatim}

\begin{Shaded}
\begin{Highlighting}[]
\FunctionTok{plot\_ecdf\_diff}\NormalTok{(sbc\_results\_ml\_drift) }\SpecialCharTok{+}
  \FunctionTok{ggtitle}\NormalTok{(}\StringTok{"Multilevel SBC: smooth drift scenario"}\NormalTok{)}
\end{Highlighting}
\end{Shaded}

\begin{center}
\includegraphics[width=0.8\linewidth,height=\textheight,keepaspectratio]{12-CategorizationModelsPrototypes_files/figure-pdf/ch13_sbc_multilevel_drift-1.pdf}
\end{center}

\begin{Shaded}
\begin{Highlighting}[]
\FunctionTok{plot\_rank\_hist}\NormalTok{(sbc\_results\_ml\_drift) }\SpecialCharTok{+}
  \FunctionTok{ggtitle}\NormalTok{(}\StringTok{"Rank histograms — multilevel drift"}\NormalTok{)}
\end{Highlighting}
\end{Shaded}

\begin{center}
\includegraphics[width=0.8\linewidth,height=\textheight,keepaspectratio]{12-CategorizationModelsPrototypes_files/figure-pdf/ch13_sbc_multilevel_drift-2.pdf}
\end{center}

\textbf{Reading the drift multilevel SBC.} Under the assumption-matched
scenario all six curves should remain inside the simultaneous band and
rank histograms should be approximately uniform. Contrast with the
static multilevel SBC above: \texttt{pop\_log\_q\_mean} and
\texttt{subj\_log\_q{[}1{]}} should shift from calibration failure
(static) to calibration success (drift). \texttt{pop\_log\_q\_sd} is the
most demanding curve --- even on drift data, estimating the
\emph{spread} of individual \(\log q\) values requires enough
between-subject variability to be visible above the within-subject
noise. The two multilevel SBC panels together --- static as failure
case, drift as success case --- provide the population-level analogue of
the single-subject SBC trilogy from §``Non-Stationary Environments''.

\begin{center}\rule{0.5\linewidth}{0.5pt}\end{center}

\section{Cognitive Insights and Model
Comparison}\label{cognitive-insights-and-model-comparison}

The Kalman filter prototype model offers valuable insights:

\begin{itemize}
\tightlist
\item
  \textbf{Incremental, Adaptive Learning}: It formalizes how category
  knowledge can be built up one example at a time, with learning
  automatically slowing down as confidence increases.
\item
  \textbf{Uncertainty Representation}: It explicitly models the
  learner's uncertainty about category prototypes --- something the
  basic GCM does not directly represent.
\item
  \textbf{Memory Efficiency}: It achieves learning by storing only
  summary statistics (mean and covariance) per category, aligning with
  intuitions about cognitive limits.
\end{itemize}

\subsection{Comparing to Exemplar Models
(GCM)}\label{comparing-to-exemplar-models-gcm}

\begin{itemize}
\tightlist
\item
  The GCM stores every detail; the prototype model abstracts.
\item
  The GCM is sensitive to specific outliers; the prototype model is less
  so (outliers get averaged in).
\item
  The GCM requires explicit attention weights (\(w\)) as free
  parameters; attention in the Kalman filter can emerge implicitly if
  uncertainty differs across dimensions (though not explored in detail
  here).
\item
  The GCM's sensitivity parameter \(c\) controls generalization
  globally; the Kalman filter's \(r\) controls how rapidly the prototype
  absorbs new information, providing a different kind of flexibility.
\end{itemize}

A note of caution about ELPD comparisons in Chapter 15: now that LFO-CV
is implemented for both the GCM (Chapter 11) and the Kalman prototype
(this chapter), the comparison between them is on equal predictive
footing. The rule mixture model in Chapter 13 is conditionally
exchangeable and uses standard PSIS-LOO, which is technically valid for
\emph{that} model --- but the comparison across architectures is still a
comparison of fundamentally different generative processes, and ``ELPD
is higher'' should not be read as ``this is the right cognitive
theory.'' See Chapter 13 for the deeper version of this caveat (and the
ABC-based alternative comparison framework that resolves it more
cleanly).

\subsection{Limitations}\label{limitations-2}

\begin{itemize}
\item
  \textbf{Abstraction Loses Detail}: Information about specific examples
  and within-category variance beyond the prototype covariance is
  discarded. Hybrid models attempt to combine both approaches.
\item
  \textbf{Single Prototype}: Assumes categories have a single central
  tendency, struggling with disjunctive categories (e.g., ``mammals that
  fly or swim''). Extensions using multiple prototypes address this.
\item
  \textbf{Path-dependence and LFO-CV (now discharged)}: The Kalman
  filter is path-dependent --- trial \(t\)'s choice probability depends
  on the entire filter state from trials \(1, \ldots, t-1\) --- and
  standard PSIS-LOO is an invalid measure of out-of-sample predictive
  accuracy. Earlier drafts of this material left this as an unresolved
  caveat. \textbf{This chapter implements LFO-CV via
  \texttt{psis\_lfo\_kalman()}, validates it against exact
  one-step-ahead refits, and Chapter 15 inherits the implementation.}
  The caveat is no longer outstanding.
\item
  \textbf{Conflation of observation noise and decisional precision}: The
  parameter \(r\) controls both the Kalman gain (learning speed) and the
  multivariate normal variance at decision time (decisional precision),
  alongside \(q\) which controls drift speed. These roles are partially
  separable from learning trajectories, but adding a separate softmax
  temperature \(\tau\) would further disentangle decisional precision
  from observational noise. We chose not to add \(\tau\) here to keep
  the model tractable, but the conflation is real --- when fitting human
  data, posterior estimates of \(r\) should be interpreted as an
  \emph{average} over learning-speed and decisional-precision
  contributions.
\item
  \textbf{Trial-1 indifference}: The model has no category bias and no
  perceptual fluency prior. Trial 1 is a coin flip by construction. For
  tasks with strong baseline asymmetries, this is a misspecification.
\item
  \textbf{Process noise identifiability}: The model includes process
  noise \(Q = q \cdot I\), but \(q\) and \(r\) are partially confounded
  under the static Kruschke (1993) paradigm --- both affect the
  effective likelihood width at decision time. Separate identifiability
  requires non-stationary data where prototypes demonstrably shift. The
  SBC analysis shows the extent of this confound, and the precision
  analysis quantifies how many trials of category drift are needed to
  pin down \(q\).
\item
  \textbf{Structural priors}: The initial means, initial covariances,
  and diagonal observation-noise structure are fixed by hand. None of
  these is inferred from data. A reader who finds the model behaving
  oddly on a new dataset should re-examine these choices before
  concluding that the cognitive theory is wrong.
\end{itemize}

\section{Conclusion}\label{conclusion-2}

The Kalman filter prototype model provides a dynamic, computationally
plausible account of how categories might be learned by abstracting
central tendencies and tracking uncertainty. It stands as a compelling
alternative to exemplar-based approaches, highlighting the fundamental
trade-off between detailed instance memory and abstract summary
representation in categorization.

As a Stan model it has a particularly clean shape: with two parameters
and a fully deterministic likelihood, Pathfinder initialization is
reliable, parameter recovery is tractable, the diagnostic battery is
nearly trivial to satisfy, and SBC is straightforward. That cleanliness
is \emph{why} this chapter is the right place in the trio to introduce
LFO-CV --- the algorithm is most legible when there are only two
parameters to re-weight and the per-trial filter is fast. Chapters 11,
13, and 15 inherit \texttt{psis\_lfo\_kalman()} (and its GCM sibling
\texttt{psis\_lfo\_gcm()} from Chapter 11) so that the LFO-CV caveat
that earlier drafts of this material repeated three times across three
chapters is now discharged in one place.

The next chapter introduces a third perspective --- rule-based models
--- which makes categorization decisions through explicit hypothesis
testing rather than continuous similarity to either stored examples or
running averages.

\begin{tcolorbox}[enhanced jigsaw, toprule=.15mm, toptitle=1mm, bottomtitle=1mm, leftrule=.75mm, rightrule=.15mm, colframe=quarto-callout-warning-color-frame, left=2mm, colback=white, titlerule=0mm, coltitle=black, colbacktitle=quarto-callout-warning-color!10!white, opacityback=0, title=\textcolor{quarto-callout-warning-color}{\faExclamationTriangle}\hspace{0.5em}{Prepare for a structural shift in Chapter 13}, arc=.35mm, opacitybacktitle=0.6, breakable, bottomrule=.15mm]

Chapter 13 differs from Chapters 11 and 12 in one load-bearing way:
\textbf{the cognitive theory (a Bayesian particle filter over discrete
rule hypotheses) has no Stan-tractable likelihood}. The resampling step
is non-differentiable and the latent particle trajectories cannot be
marginalised analytically. The chapter therefore introduces two distinct
models --- Model A, the cognitive theory fitted via Neural Posterior
Estimation (NPE); and Model B, a tractable Stan proxy (a marginalised
mixture over a fixed candidate rule set). The validation battery applies
separately to each. A flat SBC rank histogram for Model B does
\emph{not} validate Model A. The gap between the two posteriors on
\(\varepsilon\) is the chapter's main empirical finding: the structural
bias you accept when you replace an intractable cognitive theory with a
tractable proxy.

Coming from this chapter, the key conceptual shift is: the Kalman filter
gave us a clean two-parameter Stan model where process noise \(q\) was
both interpretable (drift speed) and identifiable from non-stationary
data. The rule particle filter has an analogous mutation parameter
\(\mu\), but it enters as a general-purpose rejuvenation mechanism
rather than as a literal model of how the environment changes --- which
means identifiability patterns differ, and the assumption-match story is
less tight. Chapter 13's scenario analysis (§``Assumption Match Across
Scenarios'') spells this out.

\end{tcolorbox}

Chapter 13 also confronts the NPE pipeline that replaces the ABC-based
approach of earlier drafts: schedule-fixed simulations, a learned
embedding network, and SBC + expected coverage for validation. The
cross-architecture comparison in Chapter 15 inherits these NPE networks
--- each model gets one trained inference network, and the full
multi-subject comparison runs in seconds thanks to amortized inference.

\bookmarksetup{startatroot}

\chapter{Rule-Based Categorization
Models}\label{rule-based-categorization-models}

\begin{quote}
\textbf{📍 Where we are in the Bayesian modeling workflow:} Chapters 11
and 12 implemented similarity-based accounts of categorization: the GCM
stores individual exemplars and decides by summed similarity, while the
Kalman filter prototype model tracks running category averages. Both
update continuously and live as well-formed Stan likelihoods. This
chapter takes a third route: \textbf{rule-based categorization}, where
learners maintain and update discrete hypotheses about logical criteria
for category membership. We implement this using a Bayesian particle
filter --- a sequential Monte Carlo method that maintains a weighted
distribution over candidate rules.

\textbf{Two distinct models will appear in this chapter, and the
distinction is load-bearing}:

\begin{itemize}
\tightlist
\item
  \textbf{Model A (the cognitive theory)} is the particle filter itself.
  Its parameters are \(\varepsilon\) (error rate),
  \(N_{\text{particles}}\) (a working-memory bound), \texttt{max\_dims},
  the rule grammar, and the resampling threshold. Its likelihood has no
  closed form because the rule space is discrete and combinatorial, the
  resampling step is non-differentiable, and the particle trajectories
  cannot be marginalised analytically. \textbf{Stan cannot fit Model A.}
  It can only be fitted via simulator-based inference, for which we use
  Neural Posterior Estimation (NPE) via the \texttt{sbi} Python package,
  accessed from R through \texttt{reticulate}.
\item
  \textbf{Model B (a tractable Stan proxy)} is a static, conditionally
  exchangeable mixture over a small fixed set of hand-picked candidate
  rules. Its likelihood is closed-form Bernoulli marginalised with
  \texttt{log\_sum\_exp}. It admits the full HMC validation pipeline.
  \textbf{Validation of Model B is not validation of Model A.} A reader
  who skips this distinction will end up using LOO scores from Model B
  as if they reported on the particle filter's predictive accuracy. They
  do not.
\end{itemize}

The pedagogical move in this chapter is to fit \emph{both} models, using
HMC for Model B and Neural Posterior Estimation (NPE) for Model A, and
to read the gap between their posteriors as the chapter's main empirical
finding: the magnitude of the structural bias you accept when you
replace an intractable cognitive theory with a tractable proxy.

Standard PSIS-LOO is \textbf{technically valid for Model B} (it is
conditionally exchangeable --- leaving out trial \(t\) does not affect
the likelihood of any other trial), so the LFO-CV machinery introduced
in Chapters 11 and 12 is not needed here. We say this once, in §``LFO-CV
does not apply here,'' and reference Chapter 11's
\texttt{psis\_lfo\_gcm()} and Chapter 12's \texttt{psis\_lfo\_kalman()}
for completeness.
\end{quote}

\bookmarksetup{startatroot}

\chapter{Rule-Based Models of
Categorization}\label{rule-based-models-of-categorization}

\begin{quote}
\textbf{📋 Validation battery roadmap for this chapter}

Two columns because two models. Validation of Model B does not transfer
to Model A.

\begin{longtable}[]{@{}
  >{\raggedright\arraybackslash}p{(\linewidth - 4\tabcolsep) * \real{0.3333}}
  >{\raggedright\arraybackslash}p{(\linewidth - 4\tabcolsep) * \real{0.3333}}
  >{\raggedright\arraybackslash}p{(\linewidth - 4\tabcolsep) * \real{0.3333}}@{}}
\toprule\noalign{}
\begin{minipage}[b]{\linewidth}\raggedright
Persona phase
\end{minipage} & \begin{minipage}[b]{\linewidth}\raggedright
Model A (particle filter)
\end{minipage} & \begin{minipage}[b]{\linewidth}\raggedright
Model B (Stan mixture proxy)
\end{minipage} \\
\midrule\noalign{}
\endhead
\bottomrule\noalign{}
\endlastfoot
0. Conceptual interrogation & §``Phase 0: Two Models, Not One'' &
§``Phase 0: Two Models, Not One'' \\
1. Forward simulator & \texttt{rule\_particle\_filter()} & none --- no
separate simulator needed \\
2. Implementation & R only --- no closed-form likelihood & §``Stan
Model: Marginalised Mixture'' \\
3.0 MCMC diagnostics & n/a (no MCMC) & §``Mandatory MCMC Diagnostic
Battery'' \\
3.2 Prior predictive & §``Prior Predictive Check (Model A)'' --- and the
cross-model warning that goes with it & uses Model A's simulator →
cross-model warning applies \\
3.3 Posterior predictive & via NPE posterior & §``Posterior Predictive
Check (Model B)'' \\
3.4 Prior--posterior update & §``NPE: Prior--Posterior Update for
\(\varepsilon\)'' & §``Prior--Posterior Update Plot (Model B)'' \\
3.5 Prior sensitivity & embedding network learns summaries; prior-width
sweep & \texttt{priorsense::powerscale\_sensitivity()} \\
4.5 LOO-PIT (randomised, binary) & n/a & §``Randomised LOO-PIT (Model
B)'' \\
4.6 LFO-CV & n/a (Model A is fitted via NPE, not likelihood) & §``LFO-CV
Does Not Apply Here'' --- Model B is exchangeable \\
6a. Parameter recovery & §``NPE: Parameter Recovery via SBC + Expected
Coverage'' & §``Parameter Recovery (Model B)'' \\
6b. SBC & §``NPE: SBC + Expected Coverage + TARP'' & §``Real SBC for the
Stan Mixture (Model B)'' \\
6b.cross. Cross-model calibration check & §``Cross-Model Calibration
Check'' --- generates from Model A, fits Model B; \textbf{labelled} as a
robustness check, not SBC & \\
6c. Posterior SBC & n/a & §``Posterior SBC for Model B'' --- local
calibration at the fitted posterior \\
Identifiability check & §``Phase 0'', Q3 (\(\varepsilon\) vs rule set) &
§``Phase 0'', Q3 \\
\end{longtable}
\end{quote}

\section{Theoretical Foundations}\label{theoretical-foundations}

In the previous chapters, we explored models based on similarity.
Exemplar models compute similarity to stored instances, while prototype
models compute similarity to category averages. However, humans
sometimes categorize using explicit criteria or rules. Think about
definitions like ``a square is a shape with four equal sides and four
right angles'' or simple diagnostic rules like ``if the animal has
feathers, it's a bird.''

Rule-based models propose that categorization often involves learning
and applying such explicit rules that define category membership by
partitioning the feature space.

\subsection{Core Assumptions of Rule-Based
Models}\label{core-assumptions-of-rule-based-models}

\begin{enumerate}
\def\labelenumi{\arabic{enumi}.}
\item
  \textbf{Explicit Rules}: Categories are represented by rules
  specifying conditions for membership (e.g., ``Height \textgreater{}
  2.5 AND Position \textless{} 3.0'').
\item
  \textbf{Rule Testing}: Learning involves generating, testing,
  weighting and refining candidate rules based on feedback.
\item
  \textbf{Selective Attention}: Rules typically focus attention on
  specific, relevant features mentioned in the rule, ignoring irrelevant
  ones.
\item
  \textbf{Deterministic Boundaries}: Simple rules often create sharp,
  axis-aligned boundaries in the feature space.
\item
  \textbf{Verbal Representation}: Rules are often easily verbalizable
  (e.g., ``If height \textgreater{} 2.5, then Category A''), and might
  actually become easier to learn and represent through verbalization.
\end{enumerate}

\subsection{Rules vs.~Similarity-Based
Approaches}\label{rules-vs.-similarity-based-approaches}

The distinction between rule-based and similarity-based approaches lies
not just in the computational details, but in fundamental assumptions
about how the mind represents categories:

\begin{longtable}[]{@{}
  >{\raggedright\arraybackslash}p{(\linewidth - 4\tabcolsep) * \real{0.1600}}
  >{\raggedright\arraybackslash}p{(\linewidth - 4\tabcolsep) * \real{0.3600}}
  >{\raggedright\arraybackslash}p{(\linewidth - 4\tabcolsep) * \real{0.4800}}@{}}
\toprule\noalign{}
\begin{minipage}[b]{\linewidth}\raggedright
Aspect
\end{minipage} & \begin{minipage}[b]{\linewidth}\raggedright
Rule-Based Models
\end{minipage} & \begin{minipage}[b]{\linewidth}\raggedright
Similarity-Based Models
\end{minipage} \\
\midrule\noalign{}
\endhead
\bottomrule\noalign{}
\endlastfoot
\textbf{Representation} & Explicit decision criteria & Stored examples
or prototypes \\
\textbf{Decision Process} & Rule application & Similarity computation \\
\textbf{Feature Use} & Selective (focus on rule-relevant features) &
Holistic (all features might contribute) \\
\textbf{Category Boundaries} & Sharp and often axis-aligned & Fuzzy and
sensitive to overall similarity \\
\textbf{Individual Differences} & Different rules & Different attention
weights \\
\textbf{Generalization Pattern} & Rectangle-like regions &
Similarity-gradient regions \\
\end{longtable}

\subsection{Rule-Based Categorization in the Six-Step
Pipeline}\label{rule-based-categorization-in-the-six-step-pipeline}

Chapters 11 and 12 used a six-stage framework to locate GCM and the
Kalman filter prototype in the space of possible categorization models.
The rule-based model occupies a distinct position at almost every stage:

\begin{longtable}[]{@{}
  >{\raggedright\arraybackslash}p{(\linewidth - 6\tabcolsep) * \real{0.2500}}
  >{\raggedright\arraybackslash}p{(\linewidth - 6\tabcolsep) * \real{0.2500}}
  >{\raggedright\arraybackslash}p{(\linewidth - 6\tabcolsep) * \real{0.2500}}
  >{\raggedright\arraybackslash}p{(\linewidth - 6\tabcolsep) * \real{0.2500}}@{}}
\toprule\noalign{}
\begin{minipage}[b]{\linewidth}\raggedright
Stage
\end{minipage} & \begin{minipage}[b]{\linewidth}\raggedright
GCM (Ch. 11)
\end{minipage} & \begin{minipage}[b]{\linewidth}\raggedright
Kalman prototype (Ch. 12)
\end{minipage} & \begin{minipage}[b]{\linewidth}\raggedright
Rule particle filter (Ch. 13)
\end{minipage} \\
\midrule\noalign{}
\endhead
\bottomrule\noalign{}
\endlastfoot
\textbf{1. Input representation} & Raw feature vector & Raw feature
vector & Raw feature vector \\
\textbf{2. Attention} & Continuous weights \(w_d\) across all features &
None (equal implicit weight) & Selective --- only rule-relevant
dimensions attend; irrelevant dimensions are structurally ignored \\
\textbf{3. Intermediate representation} & Per-category similarity sums &
Per-category Gaussian running mean + covariance & Weighted distribution
over discrete rule hypotheses \(\{r_1, \ldots, r_N\}\) \\
\textbf{4. Evidential mechanism} & Luce-ratio over similarities & Ratio
of Gaussian likelihoods & Weighted vote over rule predictions, softened
by \(\varepsilon\) \\
\textbf{5. Decision rule} & Softmax threshold & Softmax threshold &
Same \\
\textbf{6. Learning} & Exemplar accumulation (decay at rate \(\lambda\))
& Kalman update (process noise \(q\)) & Sequential Monte Carlo: weight
update → resample → (mutate at rate \(\mu\)) \\
\end{longtable}

The key contrasts are at stages 2, 3, and 6. At stage 2, the rule model
is the only one that uses \emph{structural} selective attention --- the
rule itself encodes which dimensions are relevant, rather than learned
continuous weights. At stage 3, the model maintains a \emph{hypothesis
distribution} over discrete logical objects, not a similarity sum or a
running mean. At stage 6, learning is a particle filter rather than a
closed-form update, which is why the likelihood is intractable and NPE
is required to fit Model A.

\begin{tcolorbox}[enhanced jigsaw, toprule=.15mm, toptitle=1mm, bottomtitle=1mm, leftrule=.75mm, rightrule=.15mm, colframe=quarto-callout-note-color-frame, left=2mm, colback=white, titlerule=0mm, coltitle=black, colbacktitle=quarto-callout-note-color!10!white, opacityback=0, title=\textcolor{quarto-callout-note-color}{\faInfo}\hspace{0.5em}{Why the six-step pipeline does not produce a single clean Stan model
here}, arc=.35mm, opacitybacktitle=0.6, breakable, bottomrule=.15mm]

In Chapters 11 and 12, each stage of the pipeline mapped onto a
differentiable Stan operation: gradient-based sampling could reach every
parameter. In Chapter 13, stage 3 (the discrete hypothesis distribution)
and stage 6 (the stochastic resample step) break differentiability. The
particle filter is a valid Bayesian cognitive model, but it is not a
Stan model. This is why the chapter introduces two models: Model A is
the cognitive theory, fitted via NPE; Model B is a tractable
differentiable proxy, fitted via HMC.

\end{tcolorbox}

\begin{tcolorbox}[enhanced jigsaw, toprule=.15mm, toptitle=1mm, bottomtitle=1mm, leftrule=.75mm, rightrule=.15mm, colframe=quarto-callout-note-color-frame, left=2mm, colback=white, titlerule=0mm, coltitle=black, colbacktitle=quarto-callout-note-color!10!white, opacityback=0, title=\textcolor{quarto-callout-note-color}{\faInfo}\hspace{0.5em}{Historical Context: From Hypothesis Testing to Bayesian Rule Inference}, arc=.35mm, opacitybacktitle=0.6, breakable, bottomrule=.15mm]

Rule-based accounts of categorization have the longest intellectual
pedigree of the three architectures in this trilogy --- and the most
turbulent history. They were the dominant framework first, were largely
displaced by similarity-based models, and then returned in a transformed
Bayesian form that revealed why they had been hard to formalise all
along.

\subsubsection{The Hypothesis-Testing
Paradigm}\label{the-hypothesis-testing-paradigm}

The modern study of concept learning began with Bruner, Goodnow, and
Austin (1956), \emph{A Study of Thinking}. Bruner and colleagues gave
participants geometrically varying cards and asked them to identify
which feature values defined a target concept. Subjects did not respond
randomly: they tested explicit hypotheses, updating their current
conjunctive rule in response to feedback. The book introduced a taxonomy
of strategies --- \emph{conservative focusing}, \emph{focus gambling},
\emph{successive scanning} --- that remained the vocabulary of the field
for two decades.

Levine (1975) extended this programme with a rigorous protocol for
inferring which hypothesis a subject was currently testing on each
trial, using sequences of blank trials (no feedback) interleaved with
informative ones. His work demonstrated that human learners are
genuinely maintaining discrete hypothesis states, not merely adjusting
continuous weights. The models of this era assumed noiseless rule
application: once the correct rule was identified, responses were
deterministic. Stochastic behaviour was treated as error rather than as
signal.

\subsubsection{The Crisis: Similarity Defeats Rules on Ill-Defined
Categories}\label{the-crisis-similarity-defeats-rules-on-ill-defined-categories}

The hypothesis-testing models were built for \emph{well-defined}
categories --- structures where a deterministic rule perfectly separates
the classes. The prototype and exemplar literature of the 1970s was
constructed with ill-defined categories deliberately, and on those
materials, rule models failed badly. A strict conjunctive rule (``height
\textgreater{} X AND position \textgreater{} Y'') produces sharp,
axis-aligned decision boundaries; human data from continuous feature
spaces showed graded, curved generalisation gradients inconsistent with
any finite rule set. For two decades after Rosch, rule-based theories
were widely regarded as inapplicable to natural categories.

The rehabilitation came from two directions. First, Nosofsky, Palmeri,
and McKinley (1994) showed that a model combining an exemplar process
with an error-correction rule selector --- \textbf{RULEX} --- could
account for a wide range of phenomena that neither pure exemplar nor
pure rule models handled well, including the speed of one-shot rule
learning and the presence of exception items that exemplar-only models
struggled to weight correctly. Second, F. Gregory Ashby et al. (1998)
proposed \textbf{COVIS} (COmpetition between Verbal and Implicit
Systems), a dual-process model in which an explicit rule system
(implemented in prefrontal cortex) competes with an implicit associative
system (implemented in the basal ganglia) for control of behaviour.
COVIS made contact with neuroimaging and patient data in a way that
purely algorithmic models could not, and substantially renewed empirical
interest in the rule hypothesis.

\subsubsection{The Bayesian Turn: Concepts as Probabilistic
Programs}\label{the-bayesian-turn-concepts-as-probabilistic-programs}

The decisive reframing came from Tenenbaum (1999) and the subsequent
Bayesian models of cognition programme (Tenenbaum et al. 2011).
Tenenbaum showed that a learner who places a prior over a hypothesis
space of rules and updates by Bayes' theorem naturally reproduces the
distinctive features of human concept learning: the \emph{size
principle} (smaller consistent hypotheses are preferred), rapid
generalisation from one or a few positive examples, and coherent
overhypothesis formation across domains.

Goodman et al. (2008) formalised this as \textbf{probabilistic language
of thought}: concepts are programs in a stochastic generative grammar,
and cognition is inference over programs. The framework unifies
rule-based and similarity-based accounts --- a prototype or exemplar
model is a special case of a distributional program; a conjunctive rule
is a logical program --- and makes the hypothesis space itself an object
of inference. The price is computational: exact Bayesian inference over
an open-ended program space is intractable, which is why
simulation-based methods are required.

\subsubsection{The Particle Filter and Computational
Intractability}\label{the-particle-filter-and-computational-intractability}

Particle filters (Gordon, Salmond, and Smith 1993) provide an
approximate Bayesian inference algorithm for state-space models where
the posterior over discrete latent states (here: the current rule
hypothesis) cannot be computed analytically. At each step a population
of hypotheses (``particles'') is propagated forward, weighted by
likelihood, and resampled proportionally to weight. The
\emph{rejuvenation} step --- mutating low-weight particles by sampling
from the prior --- prevents degeneracy and introduces a mutation rate
parameter \(\mu\) that controls how readily the learner abandons the
current hypothesis.

The appeal for cognitive modelling is direct: the particle filter
implements hypothesis testing in exactly the sequential, trial-by-trial
manner that Levine (1975) documented empirically, while tolerating noise
and ill-defined categories. The mutation parameter \(\mu\) plays an
analogous role to the process noise \(q\) of the Kalman prototype model
--- it controls how non-stationary the learner's hypothesis space is ---
but the mechanism is categorically different. \(Q\) in the Kalman filter
enters a \emph{continuous} Gaussian prediction step; \(\mu\) in the
particle filter drives \emph{discrete} resampling and mutation. This
difference means that identifiability patterns, scenario-match
requirements, and the relationship between \(\mu\) and environmental
stationarity all differ from what Chapter 12 established for \(q\).

The intractability of exact inference is the reason this chapter departs
from the Stan-only pipeline of Chapters 11 and 12. There is no
closed-form likelihood for the particle filter's resampling trajectory,
and no sum-and-marginalise trick that reduces it to one. The chapter
therefore works with two separate models: a cognitive theory fitted via
Neural Posterior Estimation (NPE), and a tractable Stan proxy.
Understanding why those two models agree in some regimes and diverge in
others is the chapter's central empirical finding.

\end{tcolorbox}

\subsection{Representing and Evaluating
Rules}\label{representing-and-evaluating-rules}

How can we represent a rule computationally? For our 2D stimuli (height,
position), a simple rule might involve one or both dimensions and a
threshold.

Example Rule: ``If Height \textgreater{} 3.5, then predict Category 1,
otherwise Category 0.''

We can represent this rule in R as a list:

\begin{Shaded}
\begin{Highlighting}[]
\CommentTok{\# Example: a simple 1D rule}
\NormalTok{rule\_example\_1d }\OtherTok{\textless{}{-}} \FunctionTok{list}\NormalTok{(}
  \AttributeTok{dimensions        =} \DecValTok{1}\NormalTok{,}
  \AttributeTok{thresholds        =} \FloatTok{3.5}\NormalTok{,}
  \AttributeTok{operations        =} \StringTok{"\textgreater{}"}\NormalTok{,}
  \AttributeTok{prediction\_if\_true =} \DecValTok{1}
\NormalTok{)}

\CommentTok{\# Example: a 2D rule (Height \textgreater{} 2 AND Position \textless{} 3 =\textgreater{} Cat 0)}
\NormalTok{rule\_example\_2d }\OtherTok{\textless{}{-}} \FunctionTok{list}\NormalTok{(}
  \AttributeTok{dimensions        =} \FunctionTok{c}\NormalTok{(}\DecValTok{1}\NormalTok{, }\DecValTok{2}\NormalTok{),}
  \AttributeTok{thresholds        =} \FunctionTok{c}\NormalTok{(}\FloatTok{2.0}\NormalTok{, }\FloatTok{3.0}\NormalTok{),}
  \AttributeTok{operations        =} \FunctionTok{c}\NormalTok{(}\StringTok{"\textgreater{}"}\NormalTok{, }\StringTok{"\textless{}"}\NormalTok{),}
  \AttributeTok{logic             =} \StringTok{"AND"}\NormalTok{,}
  \AttributeTok{prediction\_if\_true =} \DecValTok{0}
\NormalTok{)}

\CommentTok{\# Example: a nested rule — (Height \textgreater{} 2.5) AND ((Height \textgreater{} 3.5) OR (Position \textless{} 2.0))}
\CommentTok{\# Cognitive reading: "tall{-}enough AND (very tall OR low on screen)"}
\CommentTok{\# This is a RULEX{-}style conjunctive rule with a disjunctive exception clause.}
\CommentTok{\# The \textquotesingle{}nested\_logic\textquotesingle{} field marks this as a three{-}condition rule; the first}
\CommentTok{\# condition is in clause A and clauses B and C are combined by the inner connective.}
\CommentTok{\# The outer connective links clause A with (B OR C).}
\NormalTok{rule\_example\_nested }\OtherTok{\textless{}{-}} \FunctionTok{list}\NormalTok{(}
  \AttributeTok{dimensions         =} \FunctionTok{c}\NormalTok{(}\DecValTok{1}\NormalTok{, }\DecValTok{1}\NormalTok{, }\DecValTok{2}\NormalTok{),      }\CommentTok{\# features for each of the 3 conditions}
  \AttributeTok{thresholds         =} \FunctionTok{c}\NormalTok{(}\FloatTok{2.5}\NormalTok{, }\FloatTok{3.5}\NormalTok{, }\FloatTok{2.0}\NormalTok{),}
  \AttributeTok{operations         =} \FunctionTok{c}\NormalTok{(}\StringTok{"\textgreater{}"}\NormalTok{, }\StringTok{"\textgreater{}"}\NormalTok{, }\StringTok{"\textless{}"}\NormalTok{),}
  \AttributeTok{outer\_logic        =} \StringTok{"AND"}\NormalTok{,           }\CommentTok{\# connects clause A with (B inner\_logic C)}
  \AttributeTok{inner\_logic        =} \StringTok{"OR"}\NormalTok{,            }\CommentTok{\# connects clauses B and C}
  \AttributeTok{prediction\_if\_true =} \DecValTok{1}
\NormalTok{)}
\end{Highlighting}
\end{Shaded}

Now, a function to evaluate a given rule against a stimulus:

\begin{Shaded}
\begin{Highlighting}[]
\CommentTok{\# Evaluate a rule for a given stimulus.}
\CommentTok{\# Returns the category prediction (0 or 1).}
\NormalTok{evaluate\_rule }\OtherTok{\textless{}{-}} \ControlFlowTok{function}\NormalTok{(rule, stimulus) \{}

\NormalTok{  eval\_cond }\OtherTok{\textless{}{-}} \ControlFlowTok{function}\NormalTok{(dim\_idx, thresh, op) \{}
    \ControlFlowTok{switch}\NormalTok{(op,}
           \StringTok{"\textgreater{}"} \OtherTok{=}\NormalTok{ stimulus[dim\_idx] }\SpecialCharTok{\textgreater{}}\NormalTok{ thresh,}
           \StringTok{"\textless{}"} \OtherTok{=}\NormalTok{ stimulus[dim\_idx] }\SpecialCharTok{\textless{}}\NormalTok{ thresh,}
           \FunctionTok{stop}\NormalTok{(}\StringTok{"Invalid operation in rule"}\NormalTok{))}
\NormalTok{  \}}

  \ControlFlowTok{if}\NormalTok{ (}\FunctionTok{length}\NormalTok{(rule}\SpecialCharTok{$}\NormalTok{dimensions) }\SpecialCharTok{==} \DecValTok{1}\NormalTok{) \{}
\NormalTok{    condition\_met }\OtherTok{\textless{}{-}} \FunctionTok{eval\_cond}\NormalTok{(rule}\SpecialCharTok{$}\NormalTok{dimensions[}\DecValTok{1}\NormalTok{], rule}\SpecialCharTok{$}\NormalTok{thresholds[}\DecValTok{1}\NormalTok{], rule}\SpecialCharTok{$}\NormalTok{operations[}\DecValTok{1}\NormalTok{])}

\NormalTok{  \} }\ControlFlowTok{else} \ControlFlowTok{if}\NormalTok{ (}\FunctionTok{length}\NormalTok{(rule}\SpecialCharTok{$}\NormalTok{dimensions) }\SpecialCharTok{==} \DecValTok{2}\NormalTok{) \{}
\NormalTok{    c1 }\OtherTok{\textless{}{-}} \FunctionTok{eval\_cond}\NormalTok{(rule}\SpecialCharTok{$}\NormalTok{dimensions[}\DecValTok{1}\NormalTok{], rule}\SpecialCharTok{$}\NormalTok{thresholds[}\DecValTok{1}\NormalTok{], rule}\SpecialCharTok{$}\NormalTok{operations[}\DecValTok{1}\NormalTok{])}
\NormalTok{    c2 }\OtherTok{\textless{}{-}} \FunctionTok{eval\_cond}\NormalTok{(rule}\SpecialCharTok{$}\NormalTok{dimensions[}\DecValTok{2}\NormalTok{], rule}\SpecialCharTok{$}\NormalTok{thresholds[}\DecValTok{2}\NormalTok{], rule}\SpecialCharTok{$}\NormalTok{operations[}\DecValTok{2}\NormalTok{])}
\NormalTok{    condition\_met }\OtherTok{\textless{}{-}} \ControlFlowTok{if}\NormalTok{ (rule}\SpecialCharTok{$}\NormalTok{logic }\SpecialCharTok{==} \StringTok{"AND"}\NormalTok{) c1 }\SpecialCharTok{\&}\NormalTok{ c2}
                     \ControlFlowTok{else} \ControlFlowTok{if}\NormalTok{ (rule}\SpecialCharTok{$}\NormalTok{logic }\SpecialCharTok{==} \StringTok{"OR"}\NormalTok{) c1 }\SpecialCharTok{|}\NormalTok{ c2}
                     \ControlFlowTok{else} \FunctionTok{stop}\NormalTok{(}\StringTok{"Invalid logic in 2D rule"}\NormalTok{)}

\NormalTok{  \} }\ControlFlowTok{else} \ControlFlowTok{if}\NormalTok{ (}\FunctionTok{length}\NormalTok{(rule}\SpecialCharTok{$}\NormalTok{dimensions) }\SpecialCharTok{==} \DecValTok{3} \SpecialCharTok{\&\&} \SpecialCharTok{!}\FunctionTok{is.null}\NormalTok{(rule}\SpecialCharTok{$}\NormalTok{outer\_logic)) \{}
    \CommentTok{\# Nested rule: A outer\_logic (B inner\_logic C)}
    \CommentTok{\# rule$outer\_logic connects condition 1 with the (2 inner\_logic 3) clause}
\NormalTok{    cA }\OtherTok{\textless{}{-}} \FunctionTok{eval\_cond}\NormalTok{(rule}\SpecialCharTok{$}\NormalTok{dimensions[}\DecValTok{1}\NormalTok{], rule}\SpecialCharTok{$}\NormalTok{thresholds[}\DecValTok{1}\NormalTok{], rule}\SpecialCharTok{$}\NormalTok{operations[}\DecValTok{1}\NormalTok{])}
\NormalTok{    cB }\OtherTok{\textless{}{-}} \FunctionTok{eval\_cond}\NormalTok{(rule}\SpecialCharTok{$}\NormalTok{dimensions[}\DecValTok{2}\NormalTok{], rule}\SpecialCharTok{$}\NormalTok{thresholds[}\DecValTok{2}\NormalTok{], rule}\SpecialCharTok{$}\NormalTok{operations[}\DecValTok{2}\NormalTok{])}
\NormalTok{    cC }\OtherTok{\textless{}{-}} \FunctionTok{eval\_cond}\NormalTok{(rule}\SpecialCharTok{$}\NormalTok{dimensions[}\DecValTok{3}\NormalTok{], rule}\SpecialCharTok{$}\NormalTok{thresholds[}\DecValTok{3}\NormalTok{], rule}\SpecialCharTok{$}\NormalTok{operations[}\DecValTok{3}\NormalTok{])}
\NormalTok{    inner }\OtherTok{\textless{}{-}} \ControlFlowTok{if}\NormalTok{ (rule}\SpecialCharTok{$}\NormalTok{inner\_logic }\SpecialCharTok{==} \StringTok{"OR"}\NormalTok{) cB }\SpecialCharTok{|}\NormalTok{ cC}
             \ControlFlowTok{else} \ControlFlowTok{if}\NormalTok{ (rule}\SpecialCharTok{$}\NormalTok{inner\_logic }\SpecialCharTok{==} \StringTok{"AND"}\NormalTok{) cB }\SpecialCharTok{\&}\NormalTok{ cC}
             \ControlFlowTok{else} \FunctionTok{stop}\NormalTok{(}\StringTok{"Invalid inner\_logic in nested rule"}\NormalTok{)}
\NormalTok{    condition\_met }\OtherTok{\textless{}{-}} \ControlFlowTok{if}\NormalTok{ (rule}\SpecialCharTok{$}\NormalTok{outer\_logic }\SpecialCharTok{==} \StringTok{"AND"}\NormalTok{) cA }\SpecialCharTok{\&}\NormalTok{ inner}
                     \ControlFlowTok{else} \ControlFlowTok{if}\NormalTok{ (rule}\SpecialCharTok{$}\NormalTok{outer\_logic }\SpecialCharTok{==} \StringTok{"OR"}\NormalTok{) cA }\SpecialCharTok{|}\NormalTok{ inner}
                     \ControlFlowTok{else} \FunctionTok{stop}\NormalTok{(}\StringTok{"Invalid outer\_logic in nested rule"}\NormalTok{)}

\NormalTok{  \} }\ControlFlowTok{else}\NormalTok{ \{}
    \FunctionTok{stop}\NormalTok{(}\StringTok{"Rule dimension count not supported (only 1, 2, or 3{-}nested)"}\NormalTok{)}
\NormalTok{  \}}

  \FunctionTok{ifelse}\NormalTok{(condition\_met, rule}\SpecialCharTok{$}\NormalTok{prediction\_if\_true, }\DecValTok{1}\NormalTok{L }\SpecialCharTok{{-}}\NormalTok{ rule}\SpecialCharTok{$}\NormalTok{prediction\_if\_true)}
\NormalTok{\}}

\CommentTok{\# {-}{-}{-} Test {-}{-}{-}}
\NormalTok{test\_stimulus }\OtherTok{\textless{}{-}} \FunctionTok{c}\NormalTok{(}\FloatTok{2.5}\NormalTok{, }\FloatTok{3.8}\NormalTok{)}
\FunctionTok{cat}\NormalTok{(}\StringTok{"Rule 1 Prediction:"}\NormalTok{, }\FunctionTok{evaluate\_rule}\NormalTok{(rule\_example\_1d, test\_stimulus),}
    \StringTok{"(Expected: 0 — Height 2.5 is not \textgreater{} 3.5)}\SpecialCharTok{\textbackslash{}n}\StringTok{"}\NormalTok{)}
\end{Highlighting}
\end{Shaded}

\begin{verbatim}
Rule 1 Prediction: 0 (Expected: 0 — Height 2.5 is not > 3.5)
\end{verbatim}

\begin{Shaded}
\begin{Highlighting}[]
\FunctionTok{cat}\NormalTok{(}\StringTok{"Rule 2 Prediction:"}\NormalTok{, }\FunctionTok{evaluate\_rule}\NormalTok{(rule\_example\_2d, test\_stimulus),}
    \StringTok{"(Expected: 1 — Height \textgreater{} 2 TRUE, Position \textless{} 3 FALSE; AND =\textgreater{} FALSE =\textgreater{} predicts 1)}\SpecialCharTok{\textbackslash{}n}\StringTok{"}\NormalTok{)}
\end{Highlighting}
\end{Shaded}

\begin{verbatim}
Rule 2 Prediction: 1 (Expected: 1 — Height > 2 TRUE, Position < 3 FALSE; AND => FALSE => predicts 1)
\end{verbatim}

\begin{Shaded}
\begin{Highlighting}[]
\CommentTok{\# Nested rule test: Height \textgreater{} 2.5 TRUE; inner: Height \textgreater{} 3.5 FALSE OR Position \textless{} 2.0 FALSE =\textgreater{} FALSE}
\CommentTok{\# outer AND: TRUE AND FALSE =\textgreater{} FALSE =\textgreater{} predicts 0 (complementary of prediction\_if\_true = 1)}
\FunctionTok{cat}\NormalTok{(}\StringTok{"Nested Rule Prediction:"}\NormalTok{, }\FunctionTok{evaluate\_rule}\NormalTok{(rule\_example\_nested, test\_stimulus),}
    \StringTok{"(Expected: 0 — outer A=TRUE but inner B=FALSE OR C=FALSE =\textgreater{} FALSE)}\SpecialCharTok{\textbackslash{}n}\StringTok{"}\NormalTok{)}
\end{Highlighting}
\end{Shaded}

\begin{verbatim}
Nested Rule Prediction: 0 (Expected: 0 — outer A=TRUE but inner B=FALSE OR C=FALSE => FALSE)
\end{verbatim}

\begin{center}\rule{0.5\linewidth}{0.5pt}\end{center}

\section{Phase 0: Two Models, Not One}\label{phase-0-two-models-not-one}

The persona's Phase 0 interrogation is unusually consequential for this
chapter, because rule-based categorisation has the property that the
\emph{cognitive theory we want to test} (the particle filter) and the
\emph{model we can fit with HMC} (a marginalised mixture over fixed
rules) are not the same model --- and pretending otherwise has cost real
research conclusions in the past. We surface the structural commitments
before any code is written.

\begin{quote}
\textbf{Q1 --- What is the cognitive theory?}

The cognitive theory is \textbf{Model A}: a sequential agent that
maintains a finite working set of \(N_{\text{particles}}\) candidate
rules drawn from a structured grammar, updates the weights on those
rules via Bayes's rule given new feedback, resamples when effective
sample size drops below threshold, and chooses categories by weighted
vote. The agent's parameters are:

\begin{itemize}
\tightlist
\item
  \(\varepsilon\) --- error rate (probability that the agent applies its
  rule incorrectly on any trial)
\item
  \(N_{\text{particles}}\) --- the number of hypotheses the agent
  maintains simultaneously, interpretable as a working-memory bound
\item
  \(\text{max\_dims}\) --- the maximum number of feature conditions in a
  rule
\item
  The rule grammar --- which features, operators, and logical
  connectives are allowed
\item
  The resampling threshold --- when the agent reorganises its hypothesis
  set
\end{itemize}

The cognitive theory is the \emph{whole} particle filter, including
these structural choices. None of them is a nuisance.
\end{quote}

\begin{quote}
\textbf{Q2 --- Can the cognitive theory be fitted with HMC? No.~Why
not?}

Three reasons stack:

\begin{enumerate}
\def\labelenumi{\arabic{enumi}.}
\tightlist
\item
  \textbf{The rule space is discrete and combinatorial.} Stan needs
  gradients with respect to a continuous parameter space. There is no
  natural notion of a derivative with respect to ``which logical
  operator appears in condition 1'' or ``which feature dimension this
  rule uses.''
\item
  \textbf{The resampling step is non-differentiable.} Even if we could
  embed individual rules in a continuous space, the resample operation
  is a stochastic discrete jump that no HMC sampler can integrate over.
\item
  \textbf{The latent particle trajectories cannot be marginalised
  analytically.} In Chapter 8 we marginalised mixture components by
  summing over them; here the ``latent variable'' is a high-dimensional
  trajectory of weighted rule sets, and the corresponding sum has no
  closed form for any reasonable rule grammar.
\end{enumerate}

The principled response is \textbf{simulator-based inference}. Neural
Posterior Estimation (NPE) --- a modern neural-network-based approach to
simulation-based inference (SBI) --- is the right tool for this exact
situation: the simulator \(M(\theta, V) \mapsto y\) exists (it is the
\texttt{rule\_particle\_filter} R function), the parameters \(\theta\)
are well-defined, but the likelihood \(p(y \mid \theta)\) has no
tractable form. NPE trains a neural network (a normalising flow) on
simulated \((\theta, y)\) pairs to learn the mapping from data to
posterior distribution. Once trained, inference on any new dataset takes
milliseconds --- the upfront training cost is \textbf{amortized} across
all subsequent fits (Radev et al.~2023, Deistler et al.~2025). This
amortization is what makes the Chapter 15 comparison across all three
cognitive theories executable: train one NPE network per model, then fit
every subject in the dataset in seconds. Earlier drafts of this chapter
used rejection ABC, which has three disadvantages that NPE resolves: (i)
rejection ABC requires hand-crafted summary statistics and the posterior
is biased by their insufficiency, (ii) the simulation budget scales
linearly per subject with no reuse, and (iii) the accepted sample size
is tiny (\textasciitilde100 draws from 10,000 simulations). NPE with a
learned embedding network sidesteps all three. We implement NPE in
§``Neural Posterior Estimation for Model A'' via the \texttt{sbi} Python
package, accessed from R through \texttt{reticulate}.
\end{quote}

\begin{quote}
\textbf{Q3 --- What is Model B and how does it relate to Model A?}

\textbf{Model B is \emph{not} the same as Model A with a clever
marginalisation trick.} It is a completely different generative model: a
static, conditionally exchangeable mixture over a small fixed set of
hand-picked candidate rules, with one continuous parameter
\(\varepsilon\) and uniform weights over \(R\) rules. Its likelihood is
closed-form Bernoulli, marginalised analytically with
\texttt{log\_sum\_exp}. It has no learning dynamics, no resampling, no
working-memory bound, no rule discovery, and the candidate rules are
fixed by the modeller (us) by inspecting the stimulus layout.

Model B fits cleanly in Stan and admits the full HMC validation
pipeline. \textbf{But what it validates is Model B, not Model A.} A flat
ECDF in an SBC plot for Model B means ``Stan correctly inverts the
marginalised mixture.'' It does \emph{not} mean ``Stan correctly inverts
the particle filter,'' and it never could.

What Model B \emph{can} answer, when fitted to data generated from Model
A: ``given this particular candidate rule set, what value of
\(\varepsilon\) best explains the observed responses?'' If the particle
filter's converged rule happens to be in the candidate set, Model B's
\(\varepsilon\) posterior will track Model A's \(\varepsilon\) closely.
If the converged rule is \emph{not} in the candidate set, Model B
compensates by inflating \(\varepsilon\) to absorb the systematic
misclassifications. The \(\varepsilon\) posterior is then a
\emph{biased} estimate of the cognitive parameter, biased in a direction
the modeller cannot easily correct.

\textbf{What we commit to here:} We treat Model B as a tractable proxy
whose validation tests the proxy, not the cognitive theory. We use NPE
to fit Model A. We compare the two posteriors as the chapter's main
empirical result.
\end{quote}

\begin{quote}
\textbf{Q4 --- Are there parameters being silently fixed in Model A?}

Three structural choices that earlier drafts of this chapter buried in
code:

\begin{enumerate}
\def\labelenumi{\arabic{enumi}.}
\tightlist
\item
  \textbf{\(N_{\text{particles}}\).} Earlier versions used 50, 100, 150,
  or 200 in different sub-sections without justification. We commit here
  to a single defended default: \(N_{\text{particles}} = 100\),
  interpretable as a working-memory bound on the number of hypotheses an
  agent can maintain simultaneously. This is the value used everywhere
  in the chapter unless explicitly varied (in §``Sensitivity to
  \(N_{\text{particles}}\)'' and in the NPE subsection where it can
  become a free parameter).
\item
  \textbf{\(\text{max\_dims} = 2\).} The rule grammar allows up to
  two-condition rules (1D or 2D). We test sensitivity to this in §``Rule
  Grammar: Formal Specification'' by also fitting under
  \(\text{max\_dims} = 1\) and \(\text{max\_dims} = 3\).
\item
  \textbf{Operator and logic sets.} \(\{>, <\}\) and
  \(\{\text{AND}, \text{OR}\}\). No negation, no equality. We support
  flat two-condition rules (1D or 2D) and nested three-condition rules
  of the form \(A\ L_{\text{outer}}\ (B\ L_{\text{inner}}\ C)\); see
  §``Rule Grammar: Formal Specification.'' No deeper nesting. These are
  commitments to a specific rule language and they bias the inferential
  conclusions just as much as the prior on \(\varepsilon\) does. The
  prior-sensitivity machinery
  (\texttt{priorsense::powerscale\_sensitivity}) cannot detect bias from
  these choices because the bias is in the structural prior, not the
  parametric prior over \(\varepsilon\).
\item
  \textbf{Resampling threshold = 0.5.} The particle filter resamples
  when \(\text{ESS} < N_{\text{particles}} \times 0.5\). This controls
  how aggressively the agent reorganises its hypothesis set. We test
  sensitivity to this in §``Sensitivity to Resampling Threshold.''
\end{enumerate}
\end{quote}

\begin{quote}
\textbf{Q5 --- Are the candidate rules in Model B truly prior to the
data?}

No, and this is the part of the chapter that earlier drafts admitted in
a single bracketed remark and then ignored. \textbf{The seven candidate
rules in Model B are picked by looking at the stimulus layout.} Strictly
speaking, then, \(p(\text{rule} = r) = 1/R\) is not a prior --- it is an
empirical approximation to the \emph{posterior} over rules, conditional
on having inspected the data-generating structure. The honest framings
are either (i) treat the candidate set as a sensitivity analysis (``we
test what fit a hand-picked rule set provides; results are conditional
on this set''), or (ii) generate the candidate rules blind by running
the particle filter on a \emph{separately simulated} dataset and using
its converged top rules. Both options are exercised in §``Discharging
the Post-Hoc Rule Selection.''
\end{quote}

\textbf{Summary of Phase 0 commitments.} Two models, not one. Model A is
the cognitive theory, has no Stan-tractable likelihood, and is fitted
via Neural Posterior Estimation (NPE). Model B is a tractable proxy with
a hand-picked rule set; it admits the full HMC validation pipeline but
the validation tests the proxy, not the theory. \(N_{\text{particles}}\)
is fixed at 100 with an explicit cognitive interpretation; the
resampling threshold is fixed at 0.5 with a sensitivity sweep. The rule
grammar is documented formally below. The Stan candidate rule set is
treated as a sensitivity-analysis condition, not as a prior. The
chapter's main empirical contribution is the \textbf{gap between the NPE
posterior on \(\varepsilon\) under Model A and the Stan posterior on
\(\varepsilon\) under Model B} --- that gap is the structural bias you
accept when you replace the cognitive theory with a tractable proxy.

\begin{center}\rule{0.5\linewidth}{0.5pt}\end{center}

\section{Rule Grammar: Formal
Specification}\label{rule-grammar-formal-specification}

Earlier drafts of this chapter left the rule grammar inside the body of
\texttt{generate\_random\_rule()} and let the reader infer it from the
code. The persona's Phase 0 says: name the structural priors. Here is
the formal grammar.

A rule \(r\) in our grammar is a tuple

\[r = (D, T, O, L, c) \in \mathcal{R}\]

where:

\begin{itemize}
\tightlist
\item
  \(D \in \mathcal{P}(\{1, 2\}) \setminus \emptyset\) with
  \(|D| \leq \text{max\_dims}\) is the set of feature dimensions used by
  the rule. With \(\text{max\_dims} = 2\), \(|D| \in \{1, 2\}\).
\item
  \(T \in \mathbb{R}^{|D|}\) is the threshold vector, with the \(k\)-th
  element drawn uniformly from the observed range of feature \(D_k\) (in
  our task: height \(\in [1, 4]\) and position \(\in [1, 4]\)).
\item
  \(O \in \{>, <\}^{|D|}\) is the operator vector.
\item
  \(L \in \{\text{AND}, \text{OR}\}\) is the logical connective; only
  used when \(|D| = 2\).
\item
  \(c \in \{0, 1\}\) is the prediction emitted when all conditions are
  satisfied. The complementary prediction is emitted otherwise.
\end{itemize}

\textbf{Extended grammar: nested three-condition rules.} We also support
a second rule form that captures RULEX-style exception clauses (Nosofsky
et al.~1994):

\[r_{\text{nested}} = (d_A, t_A, o_A, d_B, t_B, o_B, d_C, t_C, o_C, L_{\text{outer}}, L_{\text{inner}}, c)\]

evaluated as
\(\text{condition}_A\ L_{\text{outer}}\ (\text{condition}_B\ L_{\text{inner}}\ \text{condition}_C)\).
The inner clause \((B\ L_{\text{inner}}\ C)\) acts as a disjunctive
exception when \(L_{\text{outer}} = \text{AND}\) and
\(L_{\text{inner}} = \text{OR}\): ``the rule applies if A holds, but
only unless the exception (B OR C) saves a stimulus from the incorrect
prediction.'' Cognitive reading of \texttt{rule\_example\_nested}:
``predict Cat 1 if Height \textgreater{} 2.5 AND (Height \textgreater{}
3.5 OR Position \textless{} 2.0)'' --- a tall-enough stimulus qualifies,
\emph{and} especially if it is very tall or at a low screen position.
With only two features in the Kruschke task, the feature index \(d_A\)
may coincide with \(d_B\) or \(d_C\) without loss of generality; the
thresholds differ, so the conditions are not redundant.

The number of distinct rules in this grammar is finite (although large).
For \(\text{max\_dims} = 2\) flat rules on a 2-feature task:
\(|\mathcal{R}| = 2 \cdot |T_{\text{grid}}| \cdot 2 \cdot 2 + \binom{2}{2} \cdot |T_{\text{grid}}|^2 \cdot 4 \cdot 2 \cdot 2\),
where \(|T_{\text{grid}}|\) is the threshold discretisation. Nested
three-condition rules add
\(2^3 \cdot |T_{\text{grid}}|^3 \cdot 4 \cdot 2\) additional equivalence
classes. With a continuous threshold space the grammar has uncountably
many rules in both forms.

\textbf{Commitments encoded in this grammar:}

\begin{itemize}
\tightlist
\item
  No negation operator. A rule cannot say ``NOT (Height \textgreater{}
  2).'' Negation can be partially recovered by flipping \(c\) and
  choosing the complementary inequality, but compound negations cannot
  be expressed.
\item
  No equality operator. Strict inequalities only.
\item
  Maximum two conditions per flat rule, or three conditions in a nested
  rule. A reader who believes humans use deeper rule trees will need to
  extend the grammar further.
\item
  Continuous thresholds. This defeats analytical marginalisation in Stan
  (see §``Phase 0: Two Models, Not One'').
\item
  Nested rules use a fixed two-level tree: one outer condition and one
  inner clause of two conditions. Deeper trees (e.g., A AND (B AND (C OR
  D))) are not implemented.
\end{itemize}

We will test sensitivity to \texttt{max\_dims} in the dedicated
subsection, and the operator/logic-set commitments are flagged in
§``Limitations.''

\begin{center}\rule{0.5\linewidth}{0.5pt}\end{center}

\section{The Challenge: Searching the Space of
Rules}\label{the-challenge-searching-the-space-of-rules}

Evaluating a given rule is easy. The hard part of rule-based learning is
finding the right rule. The number of possible rules (combinations of
dimensions, thresholds, operations, logic) is vast. How can a learner
efficiently search this space?

Unlike exemplar and prototype models, where the representation is
straightforward, rule-based models need mechanisms for:

\begin{enumerate}
\def\labelenumi{\arabic{enumi}.}
\tightlist
\item
  Generating candidate rules
\item
  Testing rules against observed data
\item
  Switching between rules when necessary
\item
  Maintaining uncertainty about which rule is correct
\end{enumerate}

\begin{center}\rule{0.5\linewidth}{0.5pt}\end{center}

\section{The Bayesian Particle Filter: Tracking Rule Hypotheses (Model
A)}\label{the-bayesian-particle-filter-tracking-rule-hypotheses-model-a}

The Bayesian particle filter provides an elegant solution to these
challenges. Originally developed for tracking physical objects, particle
filters have been adapted to model how humans track and update
hypotheses. In our context, each ``particle'' represents a possible
categorization rule, and the filter maintains and updates a distribution
over these rules.

\textbf{This is Model A --- the cognitive theory.} Everything in this
section implements the cognitive model whose intractable likelihood will
eventually require NPE to fit.

The core idea:

\begin{enumerate}
\def\labelenumi{\arabic{enumi}.}
\item
  \textbf{Particles = Rules}: Each particle is a specific candidate
  rule.
\item
  \textbf{Weights = Belief}: Each particle has a weight \(w_i\)
  representing the probability that this specific rule is the ``true''
  rule. Initially, all weights are equal (\(1/N\)).
\item
  \textbf{Update Weights}: When a new stimulus and its correct category
  are observed, we update the weights. Rules that correctly predict the
  category get their weights increased; rules that predict incorrectly
  get their weights decreased. The amount of change depends on how
  likely the observation was under that rule.
\item
  \textbf{Resample}: Over time, many particles accumulate near-zero
  weight (``particle degeneracy''). To focus computational effort, we
  periodically resample: draw \(N\) new particles with probability
  proportional to their weights. High-weight particles are likely to be
  selected multiple times; low-weight particles are likely to be
  dropped. After resampling, weights are reset to \(1/N\).
\end{enumerate}

\subsection{Why Model A's Likelihood is Intractable for
HMC}\label{why-model-as-likelihood-is-intractable-for-hmc}

Stan's HMC sampler requires a continuous, differentiable log posterior
density. Individual rule particles are discrete, symbolic structures ---
lists containing integers, operators, and real-valued thresholds. The
\emph{space} of all possible rules is a combinatorial set of structured
objects: there is no natural notion of a gradient with respect to
``which logical operator appears in condition 1'' or ``which dimension
is used.'' This is exactly the class of discrete latent variables that
must be integrated out analytically before any Stan code is written (the
same principle we applied to mixture components in Chapter 8).

For Model A, this analytical marginalisation is \textbf{not feasible in
closed form} over a combinatorial rule space crossed with continuous
thresholds and a non-differentiable resampling step. The particle filter
therefore lives entirely in R, where it acts as its own approximate
Bayesian inference engine via sequential Monte Carlo.

\textbf{This is the standard situation for which simulator-based
inference was invented.} The Lintusaari et al.~ABC review describes the
pattern exactly:

\begin{quote}
``Many recent models in biology describe nature to a high degree of
accuracy but are not amenable to analytical treatment. The models can,
however, be simulated on computers\ldots{} While it is usually
relatively easy to generate data from the models for any configuration
of the parameters, the real interest is often focused on the inverse
problem: the identification of parameter configurations that would
plausibly lead to data that are sufficiently similar to the observed
data.''
\end{quote}

Model A is exactly such a model:
\(\theta = (\varepsilon, N_{\text{particles}}, \text{max\_dims}, \ldots)\),
\(V\) is the random variables driving particle initialisation and
resampling, \(M\) is the \texttt{rule\_particle\_filter} function, and
\(y\) is the trial-by-trial response sequence. Section ``Neural
Posterior Estimation for Model A'' implements NPE for this simulator.
Until then, we treat the particle filter purely as a forward simulator.

\subsection{Step 1: Generating Initial Rule
Particles}\label{step-1-generating-initial-rule-particles}

\begin{Shaded}
\begin{Highlighting}[]
\CommentTok{\# Generate a single random rule.}
\CommentTok{\# This is the implementation of the formal grammar from §"Rule Grammar".}
\NormalTok{generate\_random\_rule }\OtherTok{\textless{}{-}} \ControlFlowTok{function}\NormalTok{(n\_features,}
                                 \AttributeTok{max\_dims      =}\NormalTok{ MAX\_RULE\_DIMS,}
                                 \AttributeTok{feature\_range =} \FunctionTok{list}\NormalTok{(}\FunctionTok{c}\NormalTok{(}\DecValTok{0}\NormalTok{, }\DecValTok{5}\NormalTok{), }\FunctionTok{c}\NormalTok{(}\DecValTok{0}\NormalTok{, }\DecValTok{5}\NormalTok{)),}
                                 \AttributeTok{nested\_prob   =} \FloatTok{0.25}\NormalTok{) \{}
  \CommentTok{\# With probability nested\_prob, generate a nested 3{-}condition rule (if}
  \CommentTok{\# there are \textgreater{}= 2 features and the grammar allows it). Otherwise generate}
  \CommentTok{\# a flat 1{-} or 2{-}condition rule as before.}
\NormalTok{  use\_nested }\OtherTok{\textless{}{-}}\NormalTok{ (n\_features }\SpecialCharTok{\textgreater{}=} \DecValTok{2}\NormalTok{) }\SpecialCharTok{\&\&}\NormalTok{ (}\FunctionTok{runif}\NormalTok{(}\DecValTok{1}\NormalTok{) }\SpecialCharTok{\textless{}}\NormalTok{ nested\_prob)}

\NormalTok{  sample\_thresh }\OtherTok{\textless{}{-}} \ControlFlowTok{function}\NormalTok{(di) }\FunctionTok{runif}\NormalTok{(}\DecValTok{1}\NormalTok{, feature\_range[[di]][}\DecValTok{1}\NormalTok{], feature\_range[[di]][}\DecValTok{2}\NormalTok{])}

  \ControlFlowTok{if}\NormalTok{ (use\_nested) \{}
    \CommentTok{\# Three conditions; dimensions may repeat (same feature, different thresholds)}
\NormalTok{    dims\_used  }\OtherTok{\textless{}{-}} \FunctionTok{sample}\NormalTok{(}\FunctionTok{seq\_len}\NormalTok{(n\_features), }\AttributeTok{size =} \DecValTok{3}\NormalTok{, }\AttributeTok{replace =} \ConstantTok{TRUE}\NormalTok{)}
\NormalTok{    thresholds }\OtherTok{\textless{}{-}} \FunctionTok{mapply}\NormalTok{(sample\_thresh, dims\_used)}
\NormalTok{    operations }\OtherTok{\textless{}{-}} \FunctionTok{sample}\NormalTok{(RULE\_OPERATORS, }\DecValTok{3}\NormalTok{, }\AttributeTok{replace =} \ConstantTok{TRUE}\NormalTok{)}
    \FunctionTok{list}\NormalTok{(}
      \AttributeTok{dimensions         =}\NormalTok{ dims\_used,}
      \AttributeTok{thresholds         =}\NormalTok{ thresholds,}
      \AttributeTok{operations         =}\NormalTok{ operations,}
      \AttributeTok{outer\_logic        =} \FunctionTok{sample}\NormalTok{(RULE\_LOGIC, }\DecValTok{1}\NormalTok{),}
      \AttributeTok{inner\_logic        =} \FunctionTok{sample}\NormalTok{(RULE\_LOGIC, }\DecValTok{1}\NormalTok{),}
      \AttributeTok{prediction\_if\_true =} \FunctionTok{sample}\NormalTok{(}\DecValTok{0}\SpecialCharTok{:}\DecValTok{1}\NormalTok{, }\DecValTok{1}\NormalTok{)}
\NormalTok{    )}
\NormalTok{  \} }\ControlFlowTok{else}\NormalTok{ \{}
\NormalTok{    num\_dims  }\OtherTok{\textless{}{-}} \FunctionTok{sample}\NormalTok{(}\FunctionTok{seq\_len}\NormalTok{(}\FunctionTok{min}\NormalTok{(max\_dims, n\_features)), }\DecValTok{1}\NormalTok{)}
\NormalTok{    dims\_used }\OtherTok{\textless{}{-}} \FunctionTok{sort}\NormalTok{(}\FunctionTok{sample}\NormalTok{(}\FunctionTok{seq\_len}\NormalTok{(n\_features), }\AttributeTok{size =}\NormalTok{ num\_dims, }\AttributeTok{replace =} \ConstantTok{FALSE}\NormalTok{))}
\NormalTok{    thresholds }\OtherTok{\textless{}{-}} \FunctionTok{numeric}\NormalTok{(num\_dims)}
\NormalTok{    operations }\OtherTok{\textless{}{-}} \FunctionTok{character}\NormalTok{(num\_dims)}
    \ControlFlowTok{for}\NormalTok{ (k }\ControlFlowTok{in} \FunctionTok{seq\_len}\NormalTok{(num\_dims)) \{}
\NormalTok{      di            }\OtherTok{\textless{}{-}}\NormalTok{ dims\_used[k]}
\NormalTok{      thresholds[k] }\OtherTok{\textless{}{-}} \FunctionTok{sample\_thresh}\NormalTok{(di)}
\NormalTok{      operations[k] }\OtherTok{\textless{}{-}} \FunctionTok{sample}\NormalTok{(RULE\_OPERATORS, }\DecValTok{1}\NormalTok{)}
\NormalTok{    \}}
    \FunctionTok{list}\NormalTok{(}
      \AttributeTok{dimensions        =}\NormalTok{ dims\_used,}
      \AttributeTok{thresholds        =}\NormalTok{ thresholds,}
      \AttributeTok{operations        =}\NormalTok{ operations,}
      \AttributeTok{logic             =} \ControlFlowTok{if}\NormalTok{ (num\_dims }\SpecialCharTok{==} \DecValTok{2}\NormalTok{) }\FunctionTok{sample}\NormalTok{(RULE\_LOGIC, }\DecValTok{1}\NormalTok{) }\ControlFlowTok{else} \ConstantTok{NA}\NormalTok{,}
      \AttributeTok{prediction\_if\_true =} \FunctionTok{sample}\NormalTok{(}\DecValTok{0}\SpecialCharTok{:}\DecValTok{1}\NormalTok{, }\DecValTok{1}\NormalTok{)}
\NormalTok{    )}
\NormalTok{  \}}
\NormalTok{\}}

\CommentTok{\# Initialize N particles with equal weights.}
\NormalTok{initialize\_particles }\OtherTok{\textless{}{-}} \ControlFlowTok{function}\NormalTok{(n\_particles, n\_features,}
                                  \AttributeTok{max\_dims =}\NormalTok{ MAX\_RULE\_DIMS,}
                                  \AttributeTok{feature\_range =} \FunctionTok{list}\NormalTok{(}\FunctionTok{c}\NormalTok{(}\DecValTok{0}\NormalTok{, }\DecValTok{5}\NormalTok{), }\FunctionTok{c}\NormalTok{(}\DecValTok{0}\NormalTok{, }\DecValTok{5}\NormalTok{))) \{}
\NormalTok{  particles }\OtherTok{\textless{}{-}} \FunctionTok{vector}\NormalTok{(}\StringTok{"list"}\NormalTok{, n\_particles)}
  \ControlFlowTok{for}\NormalTok{ (i }\ControlFlowTok{in} \FunctionTok{seq\_len}\NormalTok{(n\_particles))}
\NormalTok{    particles[[i]] }\OtherTok{\textless{}{-}} \FunctionTok{generate\_random\_rule}\NormalTok{(n\_features, max\_dims, feature\_range)}
  \FunctionTok{list}\NormalTok{(}\AttributeTok{particles =}\NormalTok{ particles, }\AttributeTok{weights =} \FunctionTok{rep}\NormalTok{(}\DecValTok{1} \SpecialCharTok{/}\NormalTok{ n\_particles, n\_particles))}
\NormalTok{\}}

\CommentTok{\# {-}{-}{-} Example {-}{-}{-}}
\NormalTok{n\_particles\_example }\OtherTok{\textless{}{-}} \DecValTok{5}
\NormalTok{n\_features\_example  }\OtherTok{\textless{}{-}} \DecValTok{2}
\NormalTok{initial\_system }\OtherTok{\textless{}{-}} \FunctionTok{initialize\_particles}\NormalTok{(n\_particles\_example, n\_features\_example)}
\FunctionTok{cat}\NormalTok{(}\FunctionTok{paste}\NormalTok{(}\StringTok{"Initialized"}\NormalTok{, n\_particles\_example, }\StringTok{"particles. Example rule 1:}\SpecialCharTok{\textbackslash{}n}\StringTok{"}\NormalTok{))}
\end{Highlighting}
\end{Shaded}

\begin{verbatim}
Initialized 5 particles. Example rule 1:
\end{verbatim}

\begin{Shaded}
\begin{Highlighting}[]
\FunctionTok{print}\NormalTok{(initial\_system}\SpecialCharTok{$}\NormalTok{particles[[}\DecValTok{1}\NormalTok{]])}
\end{Highlighting}
\end{Shaded}

\begin{verbatim}
$dimensions
[1] 1 2

$thresholds
[1] 4.735451 1.533177

$operations
[1] "<" ">"

$logic
[1] "AND"

$prediction_if_true
[1] 1
\end{verbatim}

\begin{Shaded}
\begin{Highlighting}[]
\FunctionTok{cat}\NormalTok{(}\StringTok{"Initial weights:"}\NormalTok{, initial\_system}\SpecialCharTok{$}\NormalTok{weights, }\StringTok{"}\SpecialCharTok{\textbackslash{}n}\StringTok{"}\NormalTok{)}
\end{Highlighting}
\end{Shaded}

\begin{verbatim}
Initial weights: 0.2 0.2 0.2 0.2 0.2 
\end{verbatim}

\subsection{Step 2: Updating Particle
Weights}\label{step-2-updating-particle-weights}

When we observe a stimulus \(\vec{x}\) and its true category \(y\), we
update weight \(w_i\) for each particle \(i\). We introduce a small
\(\varepsilon\) (\texttt{error\_prob}) to allow for occasional noise or
misapplication:

\[P(y \mid \vec{x}, \text{rule}_i) = \begin{cases} 1 - \varepsilon & \text{if rule}_i \text{ predicts } y \\ \varepsilon & \text{otherwise} \end{cases}\]

The updated weight is
\(w_i' \propto w_i \times P(y \mid \vec{x}, \text{rule}_i)\), then
renormalized.

\begin{Shaded}
\begin{Highlighting}[]
\NormalTok{update\_weights }\OtherTok{\textless{}{-}} \ControlFlowTok{function}\NormalTok{(particles, weights, stimulus, true\_category,}
                           \AttributeTok{error\_prob =} \FloatTok{0.05}\NormalTok{) \{}
\NormalTok{  n }\OtherTok{\textless{}{-}} \FunctionTok{length}\NormalTok{(particles)}
\NormalTok{  new\_w }\OtherTok{\textless{}{-}} \FunctionTok{numeric}\NormalTok{(n)}
  \ControlFlowTok{for}\NormalTok{ (i }\ControlFlowTok{in} \FunctionTok{seq\_len}\NormalTok{(n)) \{}
\NormalTok{    pred      }\OtherTok{\textless{}{-}} \FunctionTok{evaluate\_rule}\NormalTok{(particles[[i]], stimulus)}
\NormalTok{    likelihood }\OtherTok{\textless{}{-}} \ControlFlowTok{if}\NormalTok{ (pred }\SpecialCharTok{==}\NormalTok{ true\_category) }\DecValTok{1} \SpecialCharTok{{-}}\NormalTok{ error\_prob }\ControlFlowTok{else}\NormalTok{ error\_prob}
\NormalTok{    new\_w[i]  }\OtherTok{\textless{}{-}}\NormalTok{ weights[i] }\SpecialCharTok{*}\NormalTok{ likelihood}
\NormalTok{  \}}
\NormalTok{  total }\OtherTok{\textless{}{-}} \FunctionTok{sum}\NormalTok{(new\_w)}
  \ControlFlowTok{if}\NormalTok{ (total }\SpecialCharTok{\textgreater{}} \FloatTok{1e{-}9}\NormalTok{) new\_w }\SpecialCharTok{/}\NormalTok{ total }\ControlFlowTok{else} \FunctionTok{rep}\NormalTok{(}\DecValTok{1} \SpecialCharTok{/}\NormalTok{ n, n)}
\NormalTok{\}}

\CommentTok{\# {-}{-}{-} Example {-}{-}{-}}
\FunctionTok{cat}\NormalTok{(}\StringTok{"}\SpecialCharTok{\textbackslash{}n}\StringTok{{-}{-}{-} Example Weight Update {-}{-}{-}}\SpecialCharTok{\textbackslash{}n}\StringTok{"}\NormalTok{)}
\end{Highlighting}
\end{Shaded}

\begin{verbatim}

--- Example Weight Update ---
\end{verbatim}

\begin{Shaded}
\begin{Highlighting}[]
\NormalTok{updated\_w }\OtherTok{\textless{}{-}} \FunctionTok{update\_weights}\NormalTok{(}
\NormalTok{  initial\_system}\SpecialCharTok{$}\NormalTok{particles, initial\_system}\SpecialCharTok{$}\NormalTok{weights,}
  \AttributeTok{stimulus =} \FunctionTok{c}\NormalTok{(}\DecValTok{4}\NormalTok{, }\DecValTok{3}\NormalTok{), }\AttributeTok{true\_category =} \DecValTok{1}\NormalTok{, }\AttributeTok{error\_prob =} \FloatTok{0.1}
\NormalTok{)}
\NormalTok{preds }\OtherTok{\textless{}{-}} \FunctionTok{sapply}\NormalTok{(initial\_system}\SpecialCharTok{$}\NormalTok{particles, evaluate\_rule, }\AttributeTok{stimulus =} \FunctionTok{c}\NormalTok{(}\DecValTok{4}\NormalTok{, }\DecValTok{3}\NormalTok{))}
\FunctionTok{cat}\NormalTok{(}\StringTok{"Rule predictions for stimulus [4, 3]:"}\NormalTok{, preds, }\StringTok{"}\SpecialCharTok{\textbackslash{}n}\StringTok{"}\NormalTok{)}
\end{Highlighting}
\end{Shaded}

\begin{verbatim}
Rule predictions for stimulus [4, 3]: 1 1 0 0 0 
\end{verbatim}

\begin{Shaded}
\begin{Highlighting}[]
\FunctionTok{cat}\NormalTok{(}\StringTok{"Weights BEFORE:"}\NormalTok{, }\FunctionTok{round}\NormalTok{(initial\_system}\SpecialCharTok{$}\NormalTok{weights, }\DecValTok{3}\NormalTok{), }\StringTok{"}\SpecialCharTok{\textbackslash{}n}\StringTok{"}\NormalTok{)}
\end{Highlighting}
\end{Shaded}

\begin{verbatim}
Weights BEFORE: 0.2 0.2 0.2 0.2 0.2 
\end{verbatim}

\begin{Shaded}
\begin{Highlighting}[]
\FunctionTok{cat}\NormalTok{(}\StringTok{"Weights AFTER: "}\NormalTok{, }\FunctionTok{round}\NormalTok{(updated\_w, }\DecValTok{3}\NormalTok{), }\StringTok{"}\SpecialCharTok{\textbackslash{}n}\StringTok{"}\NormalTok{)}
\end{Highlighting}
\end{Shaded}

\begin{verbatim}
Weights AFTER:  0.429 0.429 0.048 0.048 0.048 
\end{verbatim}

\subsection{Step 3: Resampling
Particles}\label{step-3-resampling-particles}

If we only update weights, eventually most weights become near zero and
only a few particles dominate --- \textbf{particle degeneracy}.
Resampling revitalizes the particle set by drawing \(N\) new particles
with probability proportional to weights, then resetting all weights to
\(1/N\).

A common heuristic is to resample when the \textbf{Effective Sample
Size} (ESS) drops below a threshold (e.g., \(N/2\)). ESS is estimated as
\(1 / \sum_i w_i^2\). Low ESS indicates high degeneracy.

\begin{quote}
\textbf{A note on ESS here vs.~ESS in MCMC diagnostics.} ESS appears in
two distinct contexts in this module. In Chapters 4--10 it measures how
many effectively independent posterior draws an MCMC chain produces.
Here it measures how many particles are effectively contributing to the
filter's approximation, given their weights. Both are measures of
statistical efficiency under a sampling procedure; the underlying
intuition is the same but the contexts are different. The persona's
diagnostic battery in Phase 3.0 is about the first kind (MCMC ESS); the
resampling threshold here is about the second kind (particle ESS). Both
names are correct and they refer to genuinely different things.
\end{quote}

\begin{Shaded}
\begin{Highlighting}[]
\NormalTok{resample\_particles }\OtherTok{\textless{}{-}} \ControlFlowTok{function}\NormalTok{(particles, weights) \{}
\NormalTok{  n          }\OtherTok{\textless{}{-}} \FunctionTok{length}\NormalTok{(particles)}
\NormalTok{  idx        }\OtherTok{\textless{}{-}} \FunctionTok{sample.int}\NormalTok{(n, }\AttributeTok{size =}\NormalTok{ n, }\AttributeTok{replace =} \ConstantTok{TRUE}\NormalTok{, }\AttributeTok{prob =}\NormalTok{ weights)}
\NormalTok{  new\_p      }\OtherTok{\textless{}{-}}\NormalTok{ particles[idx]}
\NormalTok{  new\_w      }\OtherTok{\textless{}{-}} \FunctionTok{rep}\NormalTok{(}\DecValTok{1} \SpecialCharTok{/}\NormalTok{ n, n)}
  \FunctionTok{list}\NormalTok{(}\AttributeTok{particles =}\NormalTok{ new\_p, }\AttributeTok{weights =}\NormalTok{ new\_w)}
\NormalTok{\}}

\CommentTok{\# {-}{-}{-} Example {-}{-}{-}}
\FunctionTok{cat}\NormalTok{(}\StringTok{"}\SpecialCharTok{\textbackslash{}n}\StringTok{{-}{-}{-} Example Resampling {-}{-}{-}}\SpecialCharTok{\textbackslash{}n}\StringTok{"}\NormalTok{)}
\end{Highlighting}
\end{Shaded}

\begin{verbatim}

--- Example Resampling ---
\end{verbatim}

\begin{Shaded}
\begin{Highlighting}[]
\NormalTok{skewed\_w }\OtherTok{\textless{}{-}} \FunctionTok{c}\NormalTok{(}\FloatTok{0.01}\NormalTok{, }\FloatTok{0.60}\NormalTok{, }\FloatTok{0.02}\NormalTok{, }\FloatTok{0.35}\NormalTok{, }\FloatTok{0.02}\NormalTok{)}
\FunctionTok{cat}\NormalTok{(}\StringTok{"Weights BEFORE resampling:"}\NormalTok{, skewed\_w, }\StringTok{"}\SpecialCharTok{\textbackslash{}n}\StringTok{"}\NormalTok{)}
\end{Highlighting}
\end{Shaded}

\begin{verbatim}
Weights BEFORE resampling: 0.01 0.6 0.02 0.35 0.02 
\end{verbatim}

\begin{Shaded}
\begin{Highlighting}[]
\NormalTok{rs }\OtherTok{\textless{}{-}} \FunctionTok{resample\_particles}\NormalTok{(initial\_system}\SpecialCharTok{$}\NormalTok{particles, skewed\_w)}
\FunctionTok{cat}\NormalTok{(}\StringTok{"Weights AFTER resampling: "}\NormalTok{, rs}\SpecialCharTok{$}\NormalTok{weights, }\StringTok{"}\SpecialCharTok{\textbackslash{}n}\StringTok{"}\NormalTok{)}
\end{Highlighting}
\end{Shaded}

\begin{verbatim}
Weights AFTER resampling:  0.2 0.2 0.2 0.2 0.2 
\end{verbatim}

\subsection{Step 4: Making Predictions}\label{step-4-making-predictions}

The particle filter makes a categorization decision for a new stimulus
\(\vec{x}\) via a weighted vote across all particles:

\[P(\text{Choose Cat 1} \mid \vec{x}) = \sum_{i=1}^{N} w_i \times P(\text{Cat 1} \mid \vec{x}, \text{rule}_i)\]

where \(P(\text{Cat 1} \mid \vec{x}, \text{rule}_i) = 1 - \varepsilon\)
if rule \(i\) predicts category 1, and \(\varepsilon\) otherwise.

\subsection{Putting it Together: The Particle Filter
Agent}\label{putting-it-together-the-particle-filter-agent}

\begin{Shaded}
\begin{Highlighting}[]
\CommentTok{\# Full particle filter agent — this is Model A.}
\CommentTok{\# Returns a tibble of prob\_cat1 and sim\_response for each trial.}
\CommentTok{\# Note: n\_particles defaults to N\_PARTICLES\_FIXED (= 100), the structural}
\CommentTok{\# commitment from Phase 0 Q4.}
\NormalTok{rule\_particle\_filter }\OtherTok{\textless{}{-}} \ControlFlowTok{function}\NormalTok{(}\AttributeTok{n\_particles               =}\NormalTok{ N\_PARTICLES\_FIXED,}
\NormalTok{                                 error\_prob,}
\NormalTok{                                 obs,}
\NormalTok{                                 cat\_one,}
                                 \AttributeTok{max\_dims                  =}\NormalTok{ MAX\_RULE\_DIMS,}
                                 \AttributeTok{resample\_threshold\_factor =}\NormalTok{ RESAMPLE\_THRESHOLD,}
                                 \AttributeTok{quiet                     =} \ConstantTok{TRUE}\NormalTok{) \{}
\NormalTok{  n\_trials   }\OtherTok{\textless{}{-}} \FunctionTok{nrow}\NormalTok{(obs)}
\NormalTok{  n\_features }\OtherTok{\textless{}{-}} \FunctionTok{ncol}\NormalTok{(obs)}
\NormalTok{  fr         }\OtherTok{\textless{}{-}} \FunctionTok{list}\NormalTok{(}\FunctionTok{range}\NormalTok{(obs[, }\DecValTok{1}\NormalTok{]), }\FunctionTok{range}\NormalTok{(obs[, }\DecValTok{2}\NormalTok{]))}

\NormalTok{  ps       }\OtherTok{\textless{}{-}} \FunctionTok{initialize\_particles}\NormalTok{(n\_particles, n\_features, max\_dims, fr)}
\NormalTok{  particles }\OtherTok{\textless{}{-}}\NormalTok{ ps}\SpecialCharTok{$}\NormalTok{particles}
\NormalTok{  weights   }\OtherTok{\textless{}{-}}\NormalTok{ ps}\SpecialCharTok{$}\NormalTok{weights}

\NormalTok{  response\_probs }\OtherTok{\textless{}{-}} \FunctionTok{numeric}\NormalTok{(n\_trials)}

  \ControlFlowTok{for}\NormalTok{ (i }\ControlFlowTok{in} \FunctionTok{seq\_len}\NormalTok{(n\_trials)) \{}
    \ControlFlowTok{if}\NormalTok{ (}\SpecialCharTok{!}\NormalTok{quiet }\SpecialCharTok{\&\&}\NormalTok{ i }\SpecialCharTok{\%\%} \DecValTok{20} \SpecialCharTok{==} \DecValTok{0}\NormalTok{) }\FunctionTok{cat}\NormalTok{(}\StringTok{"Trial"}\NormalTok{, i, }\StringTok{"}\SpecialCharTok{\textbackslash{}n}\StringTok{"}\NormalTok{)}

\NormalTok{    stim }\OtherTok{\textless{}{-}} \FunctionTok{as.numeric}\NormalTok{(obs[i, ])}

    \CommentTok{\# ── Prediction ────────────────────────────────────────────────────────────}
\NormalTok{    preds        }\OtherTok{\textless{}{-}} \FunctionTok{sapply}\NormalTok{(particles, evaluate\_rule, }\AttributeTok{stimulus =}\NormalTok{ stim)}
\NormalTok{    p\_cat1\_rules }\OtherTok{\textless{}{-}} \FunctionTok{ifelse}\NormalTok{(preds }\SpecialCharTok{==} \DecValTok{1}\NormalTok{, }\DecValTok{1} \SpecialCharTok{{-}}\NormalTok{ error\_prob, error\_prob)}
\NormalTok{    p            }\OtherTok{\textless{}{-}} \FunctionTok{sum}\NormalTok{(weights }\SpecialCharTok{*}\NormalTok{ p\_cat1\_rules)}
\NormalTok{    response\_probs[i] }\OtherTok{\textless{}{-}} \FunctionTok{pmax}\NormalTok{(}\FloatTok{1e{-}9}\NormalTok{, }\FunctionTok{pmin}\NormalTok{(}\DecValTok{1} \SpecialCharTok{{-}} \FloatTok{1e{-}9}\NormalTok{, p))}

    \CommentTok{\# ── Weight update ─────────────────────────────────────────────────────────}
\NormalTok{    weights }\OtherTok{\textless{}{-}} \FunctionTok{update\_weights}\NormalTok{(particles, weights, stim, cat\_one[i], error\_prob)}

    \CommentTok{\# ── Resample if ESS is low ────────────────────────────────────────────────}
\NormalTok{    ess }\OtherTok{\textless{}{-}} \DecValTok{1} \SpecialCharTok{/} \FunctionTok{sum}\NormalTok{(weights}\SpecialCharTok{\^{}}\DecValTok{2}\NormalTok{)}
    \ControlFlowTok{if}\NormalTok{ (ess }\SpecialCharTok{\textless{}}\NormalTok{ n\_particles }\SpecialCharTok{*}\NormalTok{ resample\_threshold\_factor) \{}
\NormalTok{      rs        }\OtherTok{\textless{}{-}} \FunctionTok{resample\_particles}\NormalTok{(particles, weights)}
\NormalTok{      particles }\OtherTok{\textless{}{-}}\NormalTok{ rs}\SpecialCharTok{$}\NormalTok{particles}
\NormalTok{      weights   }\OtherTok{\textless{}{-}}\NormalTok{ rs}\SpecialCharTok{$}\NormalTok{weights}
\NormalTok{    \}}
\NormalTok{  \}}

  \FunctionTok{tibble}\NormalTok{(}\AttributeTok{prob\_cat1 =}\NormalTok{ response\_probs,}
         \AttributeTok{sim\_response =} \FunctionTok{rbinom}\NormalTok{(n\_trials, }\DecValTok{1}\NormalTok{, response\_probs))}
\NormalTok{\}}

\FunctionTok{cat}\NormalTok{(}\StringTok{"Particle filter agent function defined.}\SpecialCharTok{\textbackslash{}n}\StringTok{"}\NormalTok{)}
\end{Highlighting}
\end{Shaded}

\begin{verbatim}
Particle filter agent function defined.
\end{verbatim}

This implementation captures the core cognitive processes hypothesized
by rule-based models: hypothesis generation, testing against evidence,
belief updating, and adaptive focusing of the hypothesis set. \textbf{It
is also the simulator that NPE will use in §``Neural Posterior
Estimation for Model A.''}

\begin{center}\rule{0.5\linewidth}{0.5pt}\end{center}

\section{Simulating Categorization
Behavior}\label{simulating-categorization-behavior-1}

Let's simulate how Model A performs on the Kruschke (1993) task. Each
agent receives its \textbf{own independently randomized stimulus order},
consistent with the path-dependence design from Chapters 11 and 12.

\begin{Shaded}
\begin{Highlighting}[]
\CommentTok{\# Wrapper: generates per{-}agent schedule then calls rule\_particle\_filter.}
\NormalTok{simulate\_rule\_agent }\OtherTok{\textless{}{-}} \ControlFlowTok{function}\NormalTok{(agent\_id, n\_particles, max\_dims, error\_prob,}
\NormalTok{                                stimulus\_info, n\_blocks, subject\_seed) \{}
\NormalTok{  schedule }\OtherTok{\textless{}{-}} \FunctionTok{make\_subject\_schedule}\NormalTok{(stimulus\_info, n\_blocks, }\AttributeTok{seed =}\NormalTok{ subject\_seed)}
\NormalTok{  obs      }\OtherTok{\textless{}{-}} \FunctionTok{as.matrix}\NormalTok{(schedule[, }\FunctionTok{c}\NormalTok{(}\StringTok{"height"}\NormalTok{, }\StringTok{"position"}\NormalTok{)])}
\NormalTok{  cat\_one  }\OtherTok{\textless{}{-}}\NormalTok{ schedule}\SpecialCharTok{$}\NormalTok{category\_feedback}

\NormalTok{  result }\OtherTok{\textless{}{-}} \FunctionTok{rule\_particle\_filter}\NormalTok{(}
    \AttributeTok{n\_particles               =}\NormalTok{ n\_particles,}
    \AttributeTok{error\_prob                =}\NormalTok{ error\_prob,}
    \AttributeTok{obs                       =}\NormalTok{ obs,}
    \AttributeTok{cat\_one                   =}\NormalTok{ cat\_one,}
    \AttributeTok{max\_dims                  =}\NormalTok{ max\_dims,}
    \AttributeTok{resample\_threshold\_factor =}\NormalTok{ RESAMPLE\_THRESHOLD,}
    \AttributeTok{quiet                     =} \ConstantTok{TRUE}
\NormalTok{  )}

\NormalTok{  schedule }\SpecialCharTok{|\textgreater{}}
    \FunctionTok{mutate}\NormalTok{(}
      \AttributeTok{agent\_id    =}\NormalTok{ agent\_id,}
      \AttributeTok{n\_particles =}\NormalTok{ n\_particles,}
      \AttributeTok{max\_dims    =}\NormalTok{ max\_dims,}
      \AttributeTok{error\_prob  =}\NormalTok{ error\_prob,}
      \AttributeTok{prob\_cat1   =}\NormalTok{ result}\SpecialCharTok{$}\NormalTok{prob\_cat1,}
      \AttributeTok{sim\_response =}\NormalTok{ result}\SpecialCharTok{$}\NormalTok{sim\_response,}
      \AttributeTok{correct     =} \FunctionTok{as.integer}\NormalTok{(category\_feedback }\SpecialCharTok{==}\NormalTok{ sim\_response)}
\NormalTok{    ) }\SpecialCharTok{|\textgreater{}}
    \FunctionTok{group\_by}\NormalTok{(agent\_id, n\_particles, max\_dims, error\_prob) }\SpecialCharTok{|\textgreater{}}
    \FunctionTok{mutate}\NormalTok{(}\AttributeTok{performance =} \FunctionTok{cumsum}\NormalTok{(correct) }\SpecialCharTok{/} \FunctionTok{row\_number}\NormalTok{()) }\SpecialCharTok{|\textgreater{}}
    \FunctionTok{ungroup}\NormalTok{()}
\NormalTok{\}}

\NormalTok{param\_df\_rules }\OtherTok{\textless{}{-}} \FunctionTok{expand\_grid}\NormalTok{(}
  \AttributeTok{agent\_id    =} \DecValTok{1}\SpecialCharTok{:}\DecValTok{10}\NormalTok{,}
  \AttributeTok{n\_particles =} \FunctionTok{c}\NormalTok{(}\DecValTok{50}\NormalTok{, }\DecValTok{200}\NormalTok{),         }\CommentTok{\# exploring sensitivity to N\_particles}
  \AttributeTok{max\_dims    =} \FunctionTok{c}\NormalTok{(}\DecValTok{1}\NormalTok{, }\DecValTok{2}\NormalTok{),            }\CommentTok{\# exploring sensitivity to grammar depth}
  \AttributeTok{error\_prob  =} \FunctionTok{c}\NormalTok{(}\FloatTok{0.05}\NormalTok{, }\FloatTok{0.2}\NormalTok{)}
\NormalTok{) }\SpecialCharTok{|\textgreater{}}
  \FunctionTok{mutate}\NormalTok{(}\AttributeTok{subject\_seed =}\NormalTok{ agent\_id)   }\CommentTok{\# unique, reproducible seed per agent}

\NormalTok{rule\_sim\_file }\OtherTok{\textless{}{-}} \FunctionTok{here}\NormalTok{(}\StringTok{"simdata"}\NormalTok{, }\StringTok{"ch14\_rule\_simulated\_responses.csv"}\NormalTok{)}

\ControlFlowTok{if}\NormalTok{ (regenerate\_simulations }\SpecialCharTok{||} \SpecialCharTok{!}\FunctionTok{file.exists}\NormalTok{(rule\_sim\_file)) \{}
  \FunctionTok{cat}\NormalTok{(}\StringTok{"Regenerating rule{-}based simulations...}\SpecialCharTok{\textbackslash{}n}\StringTok{"}\NormalTok{)}
\NormalTok{  rule\_responses }\OtherTok{\textless{}{-}} \FunctionTok{future\_pmap\_dfr}\NormalTok{(}
    \FunctionTok{list}\NormalTok{(}
      \AttributeTok{agent\_id    =}\NormalTok{ param\_df\_rules}\SpecialCharTok{$}\NormalTok{agent\_id,}
      \AttributeTok{n\_particles =}\NormalTok{ param\_df\_rules}\SpecialCharTok{$}\NormalTok{n\_particles,}
      \AttributeTok{max\_dims    =}\NormalTok{ param\_df\_rules}\SpecialCharTok{$}\NormalTok{max\_dims,}
      \AttributeTok{error\_prob  =}\NormalTok{ param\_df\_rules}\SpecialCharTok{$}\NormalTok{error\_prob,}
      \AttributeTok{subject\_seed =}\NormalTok{ param\_df\_rules}\SpecialCharTok{$}\NormalTok{subject\_seed}
\NormalTok{    ),}
\NormalTok{    simulate\_rule\_agent,}
    \AttributeTok{stimulus\_info =}\NormalTok{ stimulus\_info,}
    \AttributeTok{n\_blocks      =}\NormalTok{ n\_blocks,}
    \AttributeTok{.options =} \FunctionTok{furrr\_options}\NormalTok{(}\AttributeTok{seed =} \ConstantTok{TRUE}\NormalTok{)}
\NormalTok{  )}
  \FunctionTok{write\_csv}\NormalTok{(rule\_responses, rule\_sim\_file)}
  \FunctionTok{cat}\NormalTok{(}\StringTok{"Simulations saved.}\SpecialCharTok{\textbackslash{}n}\StringTok{"}\NormalTok{)}
\NormalTok{\} }\ControlFlowTok{else}\NormalTok{ \{}
\NormalTok{  rule\_responses }\OtherTok{\textless{}{-}} \FunctionTok{read\_csv}\NormalTok{(rule\_sim\_file, }\AttributeTok{show\_col\_types =} \ConstantTok{FALSE}\NormalTok{)}
  \FunctionTok{cat}\NormalTok{(}\StringTok{"Simulations loaded.}\SpecialCharTok{\textbackslash{}n}\StringTok{"}\NormalTok{)}
\NormalTok{\}}
\end{Highlighting}
\end{Shaded}

\begin{verbatim}
Simulations loaded.
\end{verbatim}

\begin{Shaded}
\begin{Highlighting}[]
\FunctionTok{ggplot}\NormalTok{(rule\_responses,}
       \FunctionTok{aes}\NormalTok{(}\AttributeTok{x =}\NormalTok{ trial\_within\_subject, }\AttributeTok{y =}\NormalTok{ performance,}
           \AttributeTok{color =} \FunctionTok{factor}\NormalTok{(n\_particles))) }\SpecialCharTok{+}
  \FunctionTok{stat\_summary}\NormalTok{(}\AttributeTok{fun =}\NormalTok{ mean, }\AttributeTok{geom =} \StringTok{"line"}\NormalTok{, }\AttributeTok{linewidth =} \DecValTok{1}\NormalTok{) }\SpecialCharTok{+}
  \FunctionTok{stat\_summary}\NormalTok{(}\AttributeTok{fun.data =}\NormalTok{ mean\_se, }\AttributeTok{geom =} \StringTok{"ribbon"}\NormalTok{, }\AttributeTok{alpha =} \FloatTok{0.15}\NormalTok{,}
               \FunctionTok{aes}\NormalTok{(}\AttributeTok{fill =} \FunctionTok{factor}\NormalTok{(n\_particles)), }\AttributeTok{linetype =} \DecValTok{0}\NormalTok{) }\SpecialCharTok{+}
  \FunctionTok{facet\_grid}\NormalTok{(}
\NormalTok{    max\_dims }\SpecialCharTok{\textasciitilde{}}\NormalTok{ error\_prob,}
    \AttributeTok{labeller =} \FunctionTok{labeller}\NormalTok{(}
      \AttributeTok{max\_dims   =} \ControlFlowTok{function}\NormalTok{(x) }\FunctionTok{paste}\NormalTok{(}\StringTok{"Max Dims:"}\NormalTok{, x),}
      \AttributeTok{error\_prob =} \ControlFlowTok{function}\NormalTok{(x) }\FunctionTok{paste}\NormalTok{(}\StringTok{"Error Prob:"}\NormalTok{, x)}
\NormalTok{    )}
\NormalTok{  ) }\SpecialCharTok{+}
  \FunctionTok{scale\_color\_viridis\_d}\NormalTok{(}\AttributeTok{option =} \StringTok{"plasma"}\NormalTok{, }\AttributeTok{name =} \StringTok{"N Particles"}\NormalTok{) }\SpecialCharTok{+}
  \FunctionTok{scale\_fill\_viridis\_d}\NormalTok{(}\AttributeTok{option =} \StringTok{"plasma"}\NormalTok{, }\AttributeTok{name =} \StringTok{"N Particles"}\NormalTok{) }\SpecialCharTok{+}
  \FunctionTok{labs}\NormalTok{(}
    \AttributeTok{title    =} \StringTok{"Rule{-}Based Model Learning Performance (Model A)"}\NormalTok{,}
    \AttributeTok{subtitle =} \StringTok{"Facets show effect of max rule dimensions and error probability"}\NormalTok{,}
    \AttributeTok{x =} \StringTok{"Trial"}\NormalTok{, }\AttributeTok{y =} \StringTok{"Cumulative Accuracy (Mean ± SE)"}
\NormalTok{  ) }\SpecialCharTok{+}
  \FunctionTok{ylim}\NormalTok{(}\FloatTok{0.4}\NormalTok{, }\FloatTok{1.0}\NormalTok{)}
\end{Highlighting}
\end{Shaded}

\begin{center}
\includegraphics[width=0.8\linewidth,height=\textheight,keepaspectratio]{13-CategorizationModelsRules_files/figure-pdf/ch14_simulate_rule_responses-1.pdf}
\end{center}

\textbf{Simulation Interpretation}: This visualization shows how Model
A's performance is affected by:

\begin{itemize}
\tightlist
\item
  \textbf{Number of Particles}: More particles allow better exploration
  of the rule space, potentially leading to faster learning or higher
  final accuracy, though gains may diminish. \textbf{This parameter has
  cognitive content} (working-memory bound on hypothesis count) and is a
  candidate for NPE inference rather than fixing.
\item
  \textbf{Maximum Rule Dimensions}: Allowing 2D rules
  (\texttt{max\_dims\ =\ 2}) may be necessary if the true category
  boundary requires integrating information across dimensions. Setting
  \texttt{max\_dims\ =\ 1} corresponds to a simpler grammar that cannot
  express conjunctive concepts.
\item
  \textbf{Error Probability}: A higher \texttt{error\_prob} makes the
  model more tolerant of imperfect rules, which can be helpful in noisy
  environments but may prevent convergence on a precise rule if one
  exists.
\end{itemize}

\subsection{Visualizing Rule Learning Over
Time}\label{visualizing-rule-learning-over-time}

A key feature of Model A is that the hypotheses are explicit. We can
track which rules the particle filter believes are most likely at
different points during learning.

\begin{Shaded}
\begin{Highlighting}[]
\NormalTok{track\_rule\_learning }\OtherTok{\textless{}{-}} \ControlFlowTok{function}\NormalTok{(n\_particles, error\_prob, obs, cat\_one,}
                                \AttributeTok{max\_dims                  =}\NormalTok{ MAX\_RULE\_DIMS,}
                                \AttributeTok{top\_k                     =} \DecValTok{1}\NormalTok{,}
                                \AttributeTok{resample\_threshold\_factor =}\NormalTok{ RESAMPLE\_THRESHOLD) \{}
\NormalTok{  n\_trials  }\OtherTok{\textless{}{-}} \FunctionTok{nrow}\NormalTok{(obs)}
\NormalTok{  n\_features }\OtherTok{\textless{}{-}} \FunctionTok{ncol}\NormalTok{(obs)}
\NormalTok{  fr         }\OtherTok{\textless{}{-}} \FunctionTok{list}\NormalTok{(}\FunctionTok{range}\NormalTok{(obs[, }\DecValTok{1}\NormalTok{]), }\FunctionTok{range}\NormalTok{(obs[, }\DecValTok{2}\NormalTok{]))}

\NormalTok{  ps        }\OtherTok{\textless{}{-}} \FunctionTok{initialize\_particles}\NormalTok{(n\_particles, n\_features, max\_dims, fr)}
\NormalTok{  particles }\OtherTok{\textless{}{-}}\NormalTok{ ps}\SpecialCharTok{$}\NormalTok{particles}
\NormalTok{  weights   }\OtherTok{\textless{}{-}}\NormalTok{ ps}\SpecialCharTok{$}\NormalTok{weights}

\NormalTok{  format\_rule\_info }\OtherTok{\textless{}{-}} \ControlFlowTok{function}\NormalTok{(rule) \{}
    \FunctionTok{list}\NormalTok{(}
      \AttributeTok{dimension1  =} \ControlFlowTok{if}\NormalTok{ (}\FunctionTok{length}\NormalTok{(rule}\SpecialCharTok{$}\NormalTok{dimensions) }\SpecialCharTok{\textgreater{}=} \DecValTok{1}\NormalTok{) rule}\SpecialCharTok{$}\NormalTok{dimensions[}\DecValTok{1}\NormalTok{] }\ControlFlowTok{else} \ConstantTok{NA\_integer\_}\NormalTok{,}
      \AttributeTok{operation1  =} \ControlFlowTok{if}\NormalTok{ (}\FunctionTok{length}\NormalTok{(rule}\SpecialCharTok{$}\NormalTok{operations)  }\SpecialCharTok{\textgreater{}=} \DecValTok{1}\NormalTok{) rule}\SpecialCharTok{$}\NormalTok{operations[}\DecValTok{1}\NormalTok{]  }\ControlFlowTok{else} \ConstantTok{NA\_character\_}\NormalTok{,}
      \AttributeTok{threshold1  =} \ControlFlowTok{if}\NormalTok{ (}\FunctionTok{length}\NormalTok{(rule}\SpecialCharTok{$}\NormalTok{thresholds)  }\SpecialCharTok{\textgreater{}=} \DecValTok{1}\NormalTok{) rule}\SpecialCharTok{$}\NormalTok{thresholds[}\DecValTok{1}\NormalTok{]  }\ControlFlowTok{else} \ConstantTok{NA\_real\_}\NormalTok{,}
      \AttributeTok{dimension2  =} \ControlFlowTok{if}\NormalTok{ (}\FunctionTok{length}\NormalTok{(rule}\SpecialCharTok{$}\NormalTok{dimensions) }\SpecialCharTok{\textgreater{}=} \DecValTok{2}\NormalTok{) rule}\SpecialCharTok{$}\NormalTok{dimensions[}\DecValTok{2}\NormalTok{] }\ControlFlowTok{else} \ConstantTok{NA\_integer\_}\NormalTok{,}
      \AttributeTok{operation2  =} \ControlFlowTok{if}\NormalTok{ (}\FunctionTok{length}\NormalTok{(rule}\SpecialCharTok{$}\NormalTok{operations)  }\SpecialCharTok{\textgreater{}=} \DecValTok{2}\NormalTok{) rule}\SpecialCharTok{$}\NormalTok{operations[}\DecValTok{2}\NormalTok{]  }\ControlFlowTok{else} \ConstantTok{NA\_character\_}\NormalTok{,}
      \AttributeTok{threshold2  =} \ControlFlowTok{if}\NormalTok{ (}\FunctionTok{length}\NormalTok{(rule}\SpecialCharTok{$}\NormalTok{thresholds)  }\SpecialCharTok{\textgreater{}=} \DecValTok{2}\NormalTok{) rule}\SpecialCharTok{$}\NormalTok{thresholds[}\DecValTok{2}\NormalTok{]  }\ControlFlowTok{else} \ConstantTok{NA\_real\_}\NormalTok{,}
      \AttributeTok{dimension3  =} \ControlFlowTok{if}\NormalTok{ (}\FunctionTok{length}\NormalTok{(rule}\SpecialCharTok{$}\NormalTok{dimensions) }\SpecialCharTok{\textgreater{}=} \DecValTok{3}\NormalTok{) rule}\SpecialCharTok{$}\NormalTok{dimensions[}\DecValTok{3}\NormalTok{] }\ControlFlowTok{else} \ConstantTok{NA\_integer\_}\NormalTok{,}
      \AttributeTok{operation3  =} \ControlFlowTok{if}\NormalTok{ (}\FunctionTok{length}\NormalTok{(rule}\SpecialCharTok{$}\NormalTok{operations)  }\SpecialCharTok{\textgreater{}=} \DecValTok{3}\NormalTok{) rule}\SpecialCharTok{$}\NormalTok{operations[}\DecValTok{3}\NormalTok{]  }\ControlFlowTok{else} \ConstantTok{NA\_character\_}\NormalTok{,}
      \AttributeTok{threshold3  =} \ControlFlowTok{if}\NormalTok{ (}\FunctionTok{length}\NormalTok{(rule}\SpecialCharTok{$}\NormalTok{thresholds)  }\SpecialCharTok{\textgreater{}=} \DecValTok{3}\NormalTok{) rule}\SpecialCharTok{$}\NormalTok{thresholds[}\DecValTok{3}\NormalTok{]  }\ControlFlowTok{else} \ConstantTok{NA\_real\_}\NormalTok{,}
      \AttributeTok{logic       =} \ControlFlowTok{if}\NormalTok{ (}\SpecialCharTok{!}\FunctionTok{is.null}\NormalTok{(rule}\SpecialCharTok{$}\NormalTok{logic)) rule}\SpecialCharTok{$}\NormalTok{logic }
                    \ControlFlowTok{else} \ControlFlowTok{if}\NormalTok{ (}\SpecialCharTok{!}\FunctionTok{is.null}\NormalTok{(rule}\SpecialCharTok{$}\NormalTok{outer\_logic)) }\FunctionTok{paste0}\NormalTok{(rule}\SpecialCharTok{$}\NormalTok{outer\_logic, }\StringTok{"("}\NormalTok{, rule}\SpecialCharTok{$}\NormalTok{inner\_logic, }\StringTok{")"}\NormalTok{)}
                    \ControlFlowTok{else} \StringTok{"N/A"}\NormalTok{,}
      \AttributeTok{prediction  =}\NormalTok{ rule}\SpecialCharTok{$}\NormalTok{prediction\_if\_true}
\NormalTok{    )}
\NormalTok{  \}}

\NormalTok{  history }\OtherTok{\textless{}{-}} \FunctionTok{vector}\NormalTok{(}\StringTok{"list"}\NormalTok{, n\_trials)}
  \ControlFlowTok{for}\NormalTok{ (i }\ControlFlowTok{in} \FunctionTok{seq\_len}\NormalTok{(n\_trials)) \{}
\NormalTok{    stim }\OtherTok{\textless{}{-}} \FunctionTok{as.numeric}\NormalTok{(obs[i, ])}
\NormalTok{    top\_idx }\OtherTok{\textless{}{-}} \FunctionTok{order}\NormalTok{(weights, }\AttributeTok{decreasing =} \ConstantTok{TRUE}\NormalTok{)[}\FunctionTok{seq\_len}\NormalTok{(}\FunctionTok{min}\NormalTok{(top\_k, n\_particles))]}
\NormalTok{    history[[i]] }\OtherTok{\textless{}{-}} \FunctionTok{map\_dfr}\NormalTok{(top\_idx, }\SpecialCharTok{\textasciitilde{}}\NormalTok{\{}
      \FunctionTok{tibble}\NormalTok{(}\AttributeTok{trial =}\NormalTok{ i, }\AttributeTok{rank =} \FunctionTok{which}\NormalTok{(top\_idx }\SpecialCharTok{==}\NormalTok{ .x),}
             \AttributeTok{weight =}\NormalTok{ weights[.x], }\SpecialCharTok{!!!}\FunctionTok{format\_rule\_info}\NormalTok{(particles[[.x]]))}
\NormalTok{    \})}

\NormalTok{    weights }\OtherTok{\textless{}{-}} \FunctionTok{update\_weights}\NormalTok{(particles, weights, stim, cat\_one[i], error\_prob)}
\NormalTok{    ess     }\OtherTok{\textless{}{-}} \DecValTok{1} \SpecialCharTok{/} \FunctionTok{sum}\NormalTok{(weights}\SpecialCharTok{\^{}}\DecValTok{2}\NormalTok{)}
    \ControlFlowTok{if}\NormalTok{ (ess }\SpecialCharTok{\textless{}}\NormalTok{ n\_particles }\SpecialCharTok{*}\NormalTok{ resample\_threshold\_factor) \{}
\NormalTok{      rs        }\OtherTok{\textless{}{-}} \FunctionTok{resample\_particles}\NormalTok{(particles, weights)}
\NormalTok{      particles }\OtherTok{\textless{}{-}}\NormalTok{ rs}\SpecialCharTok{$}\NormalTok{particles}
\NormalTok{      weights   }\OtherTok{\textless{}{-}}\NormalTok{ rs}\SpecialCharTok{$}\NormalTok{weights}
\NormalTok{    \}}
\NormalTok{  \}}
  \FunctionTok{bind\_rows}\NormalTok{(history)}
\NormalTok{\}}

\CommentTok{\# Run tracking with a fixed seed{-}1 schedule for reproducibility}
\NormalTok{track\_schedule }\OtherTok{\textless{}{-}} \FunctionTok{make\_subject\_schedule}\NormalTok{(stimulus\_info, n\_blocks, }\AttributeTok{seed =} \DecValTok{1}\NormalTok{)}

\NormalTok{rule\_tracking\_data }\OtherTok{\textless{}{-}} \FunctionTok{track\_rule\_learning}\NormalTok{(}
  \AttributeTok{n\_particles =}\NormalTok{ N\_PARTICLES\_FIXED,}
  \AttributeTok{error\_prob  =} \FloatTok{0.1}\NormalTok{,}
  \AttributeTok{obs         =} \FunctionTok{as.matrix}\NormalTok{(track\_schedule[, }\FunctionTok{c}\NormalTok{(}\StringTok{"height"}\NormalTok{, }\StringTok{"position"}\NormalTok{)]),}
  \AttributeTok{cat\_one     =}\NormalTok{ track\_schedule}\SpecialCharTok{$}\NormalTok{category\_feedback,}
  \AttributeTok{max\_dims    =}\NormalTok{ MAX\_RULE\_DIMS,}
  \AttributeTok{top\_k       =} \DecValTok{1}
\NormalTok{)}

\CommentTok{\# Format rule text for plotting}
\NormalTok{rule\_to\_text }\OtherTok{\textless{}{-}} \ControlFlowTok{function}\NormalTok{(rule\_row) \{}
\NormalTok{  dim\_names }\OtherTok{\textless{}{-}} \FunctionTok{c}\NormalTok{(}\StringTok{"Height"}\NormalTok{, }\StringTok{"Position"}\NormalTok{)}
  \ControlFlowTok{if}\NormalTok{ (}\FunctionTok{is.na}\NormalTok{(rule\_row}\SpecialCharTok{$}\NormalTok{dimension2)) \{}
    \FunctionTok{sprintf}\NormalTok{(}\StringTok{"IF \%s \%s \%.1f THEN Cat \%d"}\NormalTok{,}
\NormalTok{            dim\_names[rule\_row}\SpecialCharTok{$}\NormalTok{dimension1], rule\_row}\SpecialCharTok{$}\NormalTok{operation1,}
\NormalTok{            rule\_row}\SpecialCharTok{$}\NormalTok{threshold1, rule\_row}\SpecialCharTok{$}\NormalTok{prediction)}
\NormalTok{  \} }\ControlFlowTok{else}\NormalTok{ \{}
    \FunctionTok{sprintf}\NormalTok{(}\StringTok{"IF \%s \%s \%.1f \%s \%s \%s \%.1f THEN Cat \%d"}\NormalTok{,}
\NormalTok{            dim\_names[rule\_row}\SpecialCharTok{$}\NormalTok{dimension1], rule\_row}\SpecialCharTok{$}\NormalTok{operation1, rule\_row}\SpecialCharTok{$}\NormalTok{threshold1,}
\NormalTok{            rule\_row}\SpecialCharTok{$}\NormalTok{logic,}
\NormalTok{            dim\_names[rule\_row}\SpecialCharTok{$}\NormalTok{dimension2], rule\_row}\SpecialCharTok{$}\NormalTok{operation2, rule\_row}\SpecialCharTok{$}\NormalTok{threshold2,}
\NormalTok{            rule\_row}\SpecialCharTok{$}\NormalTok{prediction)}
\NormalTok{  \}}
\NormalTok{\}}

\NormalTok{rule\_tracking\_plot\_data }\OtherTok{\textless{}{-}}\NormalTok{ rule\_tracking\_data }\SpecialCharTok{|\textgreater{}}
  \FunctionTok{filter}\NormalTok{(rank }\SpecialCharTok{==} \DecValTok{1}\NormalTok{) }\SpecialCharTok{|\textgreater{}}
  \FunctionTok{rowwise}\NormalTok{() }\SpecialCharTok{|\textgreater{}}
  \FunctionTok{mutate}\NormalTok{(}\AttributeTok{rule\_text =} \FunctionTok{rule\_to\_text}\NormalTok{(}\FunctionTok{cur\_data}\NormalTok{())) }\SpecialCharTok{|\textgreater{}}
  \FunctionTok{ungroup}\NormalTok{() }\SpecialCharTok{|\textgreater{}}
  \FunctionTok{mutate}\NormalTok{(}\AttributeTok{rule\_block =} \FunctionTok{consecutive\_id}\NormalTok{(rule\_text))}

\FunctionTok{ggplot}\NormalTok{(rule\_tracking\_plot\_data, }\FunctionTok{aes}\NormalTok{(}\AttributeTok{x =}\NormalTok{ trial, }\AttributeTok{y =}\NormalTok{ weight)) }\SpecialCharTok{+}
  \FunctionTok{geom\_line}\NormalTok{(}\FunctionTok{aes}\NormalTok{(}\AttributeTok{group =} \DecValTok{1}\NormalTok{), }\AttributeTok{color =} \StringTok{"\#0072B2"}\NormalTok{) }\SpecialCharTok{+}
  \FunctionTok{geom\_point}\NormalTok{(}\AttributeTok{color =} \StringTok{"\#0072B2"}\NormalTok{, }\AttributeTok{size =} \FloatTok{1.5}\NormalTok{) }\SpecialCharTok{+}
  \FunctionTok{geom\_text}\NormalTok{(}
    \AttributeTok{data =}\NormalTok{ rule\_tracking\_plot\_data }\SpecialCharTok{|\textgreater{}} \FunctionTok{distinct}\NormalTok{(rule\_block, }\AttributeTok{.keep\_all =} \ConstantTok{TRUE}\NormalTok{),}
    \FunctionTok{aes}\NormalTok{(}\AttributeTok{label =}\NormalTok{ rule\_text),}
    \AttributeTok{hjust =} \DecValTok{0}\NormalTok{, }\AttributeTok{vjust =} \SpecialCharTok{{-}}\FloatTok{0.5}\NormalTok{, }\AttributeTok{size =} \DecValTok{3}\NormalTok{, }\AttributeTok{check\_overlap =} \ConstantTok{TRUE}
\NormalTok{  ) }\SpecialCharTok{+}
  \FunctionTok{scale\_y\_continuous}\NormalTok{(}\AttributeTok{limits =} \FunctionTok{c}\NormalTok{(}\DecValTok{0}\NormalTok{, }\DecValTok{1}\NormalTok{), }\AttributeTok{name =} \StringTok{"Weight (Probability) of Top Rule"}\NormalTok{) }\SpecialCharTok{+}
  \FunctionTok{scale\_x\_continuous}\NormalTok{(}\AttributeTok{limits =} \FunctionTok{c}\NormalTok{(}\DecValTok{0}\NormalTok{, total\_trials }\SpecialCharTok{+} \DecValTok{5}\NormalTok{)) }\SpecialCharTok{+}
  \FunctionTok{labs}\NormalTok{(}
    \AttributeTok{title    =} \StringTok{"Evolution of the Top Rule Hypothesis Over Time (Model A)"}\NormalTok{,}
    \AttributeTok{subtitle =} \StringTok{"Line shows the probability assigned to the single best rule at each trial."}\NormalTok{,}
    \AttributeTok{x =} \StringTok{"Trial"}
\NormalTok{  )}
\end{Highlighting}
\end{Shaded}

\begin{center}
\includegraphics[width=0.8\linewidth,height=\textheight,keepaspectratio]{13-CategorizationModelsRules_files/figure-pdf/ch14_track_rule_learning-1.pdf}
\end{center}

\textbf{Rule Evolution Interpretation}: This plot shows how the model's
belief (weight) in its single best rule hypothesis changes over trials.
Periods of high weight indicate the model has converged on a dominant
rule; sharp drops followed by a new label indicate a hypothesis switch
as the previous rule accumulated too many prediction errors.
\textbf{Note that the converged rule is one specific rule} drawn from
the grammar; the empirical fact of convergence is what makes the
post-hoc selection of Model B's candidate rules tempting and dangerous.

\subsection{Visualizing the Decision
Boundary}\label{visualizing-the-decision-boundary}

Rule-based models typically create axis-aligned, rectangular decision
boundaries. Let's visualize the boundary created by the final, most
probable rule found by Model A.

\begin{Shaded}
\begin{Highlighting}[]
\NormalTok{get\_final\_rule\_predictions }\OtherTok{\textless{}{-}} \ControlFlowTok{function}\NormalTok{(n\_particles, error\_prob,}
\NormalTok{                                       obs\_train, cat\_one\_train,}
\NormalTok{                                       grid\_points, }\AttributeTok{max\_dims =}\NormalTok{ MAX\_RULE\_DIMS,}
                                       \AttributeTok{resample\_threshold\_factor =}\NormalTok{ RESAMPLE\_THRESHOLD) \{}
\NormalTok{  n\_features }\OtherTok{\textless{}{-}} \FunctionTok{ncol}\NormalTok{(obs\_train)}
\NormalTok{  fr         }\OtherTok{\textless{}{-}} \FunctionTok{list}\NormalTok{(}\FunctionTok{range}\NormalTok{(obs\_train[, }\DecValTok{1}\NormalTok{]), }\FunctionTok{range}\NormalTok{(obs\_train[, }\DecValTok{2}\NormalTok{]))}

\NormalTok{  ps        }\OtherTok{\textless{}{-}} \FunctionTok{initialize\_particles}\NormalTok{(n\_particles, n\_features, max\_dims, fr)}
\NormalTok{  particles }\OtherTok{\textless{}{-}}\NormalTok{ ps}\SpecialCharTok{$}\NormalTok{particles}
\NormalTok{  weights   }\OtherTok{\textless{}{-}}\NormalTok{ ps}\SpecialCharTok{$}\NormalTok{weights}

  \ControlFlowTok{for}\NormalTok{ (i }\ControlFlowTok{in} \FunctionTok{seq\_len}\NormalTok{(}\FunctionTok{nrow}\NormalTok{(obs\_train))) \{}
\NormalTok{    stim    }\OtherTok{\textless{}{-}} \FunctionTok{as.numeric}\NormalTok{(obs\_train[i, ])}
\NormalTok{    weights }\OtherTok{\textless{}{-}} \FunctionTok{update\_weights}\NormalTok{(particles, weights, stim, cat\_one\_train[i], error\_prob)}
\NormalTok{    ess     }\OtherTok{\textless{}{-}} \DecValTok{1} \SpecialCharTok{/} \FunctionTok{sum}\NormalTok{(weights}\SpecialCharTok{\^{}}\DecValTok{2}\NormalTok{)}
    \ControlFlowTok{if}\NormalTok{ (ess }\SpecialCharTok{\textless{}}\NormalTok{ n\_particles }\SpecialCharTok{*}\NormalTok{ resample\_threshold\_factor) \{}
\NormalTok{      rs        }\OtherTok{\textless{}{-}} \FunctionTok{resample\_particles}\NormalTok{(particles, weights)}
\NormalTok{      particles }\OtherTok{\textless{}{-}}\NormalTok{ rs}\SpecialCharTok{$}\NormalTok{particles}
\NormalTok{      weights   }\OtherTok{\textless{}{-}}\NormalTok{ rs}\SpecialCharTok{$}\NormalTok{weights}
\NormalTok{    \}}
\NormalTok{  \}}

  \FunctionTok{apply}\NormalTok{(}\FunctionTok{as.matrix}\NormalTok{(grid\_points), }\DecValTok{1}\NormalTok{, }\ControlFlowTok{function}\NormalTok{(pt) \{}
\NormalTok{    preds }\OtherTok{\textless{}{-}} \FunctionTok{sapply}\NormalTok{(particles, evaluate\_rule, }\AttributeTok{stimulus =}\NormalTok{ pt)}
    \FunctionTok{sum}\NormalTok{(weights }\SpecialCharTok{*} \FunctionTok{ifelse}\NormalTok{(preds }\SpecialCharTok{==} \DecValTok{1}\NormalTok{, }\DecValTok{1} \SpecialCharTok{{-}}\NormalTok{ error\_prob, error\_prob))}
\NormalTok{  \})}
\NormalTok{\}}

\NormalTok{grid\_pts }\OtherTok{\textless{}{-}} \FunctionTok{expand.grid}\NormalTok{(}
  \AttributeTok{height   =} \FunctionTok{seq}\NormalTok{(}\FunctionTok{min}\NormalTok{(stimulus\_info}\SpecialCharTok{$}\NormalTok{height)   }\SpecialCharTok{{-}} \FloatTok{0.5}\NormalTok{,}
                 \FunctionTok{max}\NormalTok{(stimulus\_info}\SpecialCharTok{$}\NormalTok{height)   }\SpecialCharTok{+} \FloatTok{0.5}\NormalTok{, }\AttributeTok{length.out =} \DecValTok{50}\NormalTok{),}
  \AttributeTok{position =} \FunctionTok{seq}\NormalTok{(}\FunctionTok{min}\NormalTok{(stimulus\_info}\SpecialCharTok{$}\NormalTok{position) }\SpecialCharTok{{-}} \FloatTok{0.5}\NormalTok{,}
                 \FunctionTok{max}\NormalTok{(stimulus\_info}\SpecialCharTok{$}\NormalTok{position) }\SpecialCharTok{+} \FloatTok{0.5}\NormalTok{, }\AttributeTok{length.out =} \DecValTok{50}\NormalTok{)}
\NormalTok{)}

\CommentTok{\# Use the same track\_schedule for boundary visualization}
\NormalTok{boundary\_preds }\OtherTok{\textless{}{-}} \FunctionTok{get\_final\_rule\_predictions}\NormalTok{(}
  \AttributeTok{n\_particles   =} \DecValTok{200}\NormalTok{,}
  \AttributeTok{error\_prob    =} \FloatTok{0.1}\NormalTok{,}
  \AttributeTok{obs\_train     =} \FunctionTok{as.matrix}\NormalTok{(track\_schedule[, }\FunctionTok{c}\NormalTok{(}\StringTok{"height"}\NormalTok{, }\StringTok{"position"}\NormalTok{)]),}
  \AttributeTok{cat\_one\_train =}\NormalTok{ track\_schedule}\SpecialCharTok{$}\NormalTok{category\_feedback,}
  \AttributeTok{grid\_points   =}\NormalTok{ grid\_pts[, }\FunctionTok{c}\NormalTok{(}\StringTok{"height"}\NormalTok{, }\StringTok{"position"}\NormalTok{)],}
  \AttributeTok{max\_dims      =}\NormalTok{ MAX\_RULE\_DIMS}
\NormalTok{)}

\NormalTok{grid\_decision\_data }\OtherTok{\textless{}{-}}\NormalTok{ grid\_pts }\SpecialCharTok{|\textgreater{}}
  \FunctionTok{mutate}\NormalTok{(}
    \AttributeTok{rule\_prob\_cat1 =}\NormalTok{ boundary\_preds,}
    \AttributeTok{rule\_decision  =} \FunctionTok{as.integer}\NormalTok{(rule\_prob\_cat1 }\SpecialCharTok{\textgreater{}} \FloatTok{0.5}\NormalTok{)}
\NormalTok{  )}

\FunctionTok{ggplot}\NormalTok{() }\SpecialCharTok{+}
  \FunctionTok{geom\_tile}\NormalTok{(}\AttributeTok{data =}\NormalTok{ grid\_decision\_data,}
            \FunctionTok{aes}\NormalTok{(}\AttributeTok{x =}\NormalTok{ position, }\AttributeTok{y =}\NormalTok{ height, }\AttributeTok{fill =} \FunctionTok{factor}\NormalTok{(rule\_decision)),}
            \AttributeTok{alpha =} \FloatTok{0.4}\NormalTok{) }\SpecialCharTok{+}
  \FunctionTok{stat\_contour}\NormalTok{(}\AttributeTok{data =}\NormalTok{ grid\_decision\_data,}
               \FunctionTok{aes}\NormalTok{(}\AttributeTok{x =}\NormalTok{ position, }\AttributeTok{y =}\NormalTok{ height, }\AttributeTok{z =}\NormalTok{ rule\_prob\_cat1),}
               \AttributeTok{breaks =} \FloatTok{0.5}\NormalTok{, }\AttributeTok{color =} \StringTok{"black"}\NormalTok{, }\AttributeTok{linewidth =} \FloatTok{0.8}\NormalTok{) }\SpecialCharTok{+}
  \FunctionTok{geom\_point}\NormalTok{(}\AttributeTok{data =}\NormalTok{ stimulus\_info,}
             \FunctionTok{aes}\NormalTok{(}\AttributeTok{x =}\NormalTok{ position, }\AttributeTok{y =}\NormalTok{ height,}
                 \AttributeTok{color =} \FunctionTok{factor}\NormalTok{(category\_true), }\AttributeTok{shape =} \FunctionTok{factor}\NormalTok{(category\_true)),}
             \AttributeTok{size =} \DecValTok{3}\NormalTok{) }\SpecialCharTok{+}
  \FunctionTok{scale\_fill\_manual}\NormalTok{(}\AttributeTok{values  =} \FunctionTok{c}\NormalTok{(}\StringTok{"0"} \OtherTok{=} \StringTok{"\#E69F00"}\NormalTok{, }\StringTok{"1"} \OtherTok{=} \StringTok{"\#56B4E9"}\NormalTok{),}
                    \AttributeTok{name    =} \StringTok{"Predicted Cat"}\NormalTok{) }\SpecialCharTok{+}
  \FunctionTok{scale\_color\_manual}\NormalTok{(}\AttributeTok{values =} \FunctionTok{c}\NormalTok{(}\StringTok{"0"} \OtherTok{=} \StringTok{"\#0072B2"}\NormalTok{, }\StringTok{"1"} \OtherTok{=} \StringTok{"\#D55E00"}\NormalTok{),}
                     \AttributeTok{name    =} \StringTok{"True Category"}\NormalTok{) }\SpecialCharTok{+}
  \FunctionTok{scale\_shape\_manual}\NormalTok{(}\AttributeTok{values =} \FunctionTok{c}\NormalTok{(}\StringTok{"0"} \OtherTok{=} \DecValTok{16}\NormalTok{, }\StringTok{"1"} \OtherTok{=} \DecValTok{17}\NormalTok{), }\AttributeTok{name =} \StringTok{"True Category"}\NormalTok{) }\SpecialCharTok{+}
  \FunctionTok{labs}\NormalTok{(}
    \AttributeTok{title    =} \StringTok{"Rule{-}Based Model Decision Boundary (Model A, Particle Filter)"}\NormalTok{,}
    \AttributeTok{subtitle =} \StringTok{"Weighted average of final particle set (N=200, error=0.1)"}\NormalTok{,}
    \AttributeTok{x =} \StringTok{"Position Feature"}\NormalTok{, }\AttributeTok{y =} \StringTok{"Height Feature"}
\NormalTok{  ) }\SpecialCharTok{+}
  \FunctionTok{coord\_cartesian}\NormalTok{(}
    \AttributeTok{xlim =} \FunctionTok{range}\NormalTok{(grid\_decision\_data}\SpecialCharTok{$}\NormalTok{position),}
    \AttributeTok{ylim =} \FunctionTok{range}\NormalTok{(grid\_decision\_data}\SpecialCharTok{$}\NormalTok{height)}
\NormalTok{  )}
\end{Highlighting}
\end{Shaded}

\begin{center}
\includegraphics[width=0.8\linewidth,height=\textheight,keepaspectratio]{13-CategorizationModelsRules_files/figure-pdf/ch14_decision_boundary-1.pdf}
\end{center}

\textbf{Boundary Interpretation}: The background color shows the
predicted category for each point in the feature space, based on the
weighted vote of the final rule particles. The black contour marks the
50\% threshold. Notice the characteristic axis-aligned, rectangular
character of the boundary regions --- a direct consequence of the
threshold-based rule structure --- in contrast to the smooth, curved
boundaries that would emerge from similarity-based models.

\begin{center}\rule{0.5\linewidth}{0.5pt}\end{center}

\section{Beyond Selection: Enhancing the Particle Filter for Rule
Discovery}\label{beyond-selection-enhancing-the-particle-filter-for-rule-discovery}

The Bayesian particle filter, as implemented above, provides a powerful
framework for modeling how learners evaluate competing rule hypotheses.
By maintaining a weighted set of particles and using resampling, it
efficiently focuses on rules that best explain observed data.

However, a critical point: the basic filter primarily acts as a
\textbf{rule selection} mechanism rather than a \textbf{rule generation}
engine. It excels at identifying the most promising rules from an
\emph{initially provided set}. If the truly optimal rule is complex or
unusual and not present in the initial random pool, the basic resampling
step (which only duplicates existing rules) cannot discover it.

To create a model that more actively explores and discovers better rules
during learning, the basic particle filter can be augmented with
additional mechanisms:

\begin{enumerate}
\def\labelenumi{\arabic{enumi}.}
\item
  \textbf{Mutation/Perturbation}: Introduce small, random variations to
  rules after resampling --- slightly adjust a threshold, flip a
  comparison operator, or change the dimension used in a condition. This
  allows fine-tuning and local exploration around currently successful
  rules.
\item
  \textbf{Recombination/Crossover}: Combine elements from two successful
  ``parent'' rules to create new ``child'' rules --- for example, the
  first condition of Parent A with the second condition of Parent B.
  This allows synthesizing features discovered independently by
  different rules.
\item
  \textbf{Targeted Rule Generation}: Generate new rules specifically
  designed to address stimuli consistently misclassified by the current
  particle set, or focus on features prevalent in high-weight rules.
\item
  \textbf{Complexity Adjustment}: Probabilistically add a condition to a
  successful 1D rule to explore whether a 2D rule performs better, or
  simplify a poorly performing 2D rule into its constituent 1D
  components.
\end{enumerate}

Incorporating such exploration mechanisms transforms the particle filter
from a pure selection device into a more adaptive \textbf{rule discovery
system} capable of generating and testing rules that were not part of
the initial random set. We do not implement these extensions here, but
they represent a natural direction for more cognitively realistic models
of rule learning. Note also that any of these mechanisms makes Model A's
intractability \emph{worse}, not better, from an HMC standpoint --- they
all add stochastic discrete jumps. NPE handles them transparently: the
inference network only sees simulated \((\theta, y)\) pairs and does not
care how the simulator produces them.

\begin{center}\rule{0.5\linewidth}{0.5pt}\end{center}

\section{Model B: A Tractable Stan Proxy via Marginalised
Mixture}\label{model-b-a-tractable-stan-proxy-via-marginalised-mixture}

We now switch models. Everything in this section implements
\textbf{Model B} --- the static, conditionally exchangeable mixture over
a fixed candidate rule set. Model B is \textbf{not} the same as Model A
with a clever trick. It is a different generative process that admits
HMC inference and the full validation pipeline. Validation of Model B
tells us about Model B; it does not tell us about Model A. We will
return to this distinction at every opportunity.

\subsection{Why a separate proxy
model?}\label{why-a-separate-proxy-model}

Model A cannot be fitted with HMC for the reasons in §``Why Model A's
Likelihood is Intractable for HMC.'' But the persona's six-phase battery
is designed around HMC and a tractable likelihood, and the model
comparison framework in Chapter 15 needs ELPD scores that come from a
proper Bayesian fit. Two responses:

\begin{enumerate}
\def\labelenumi{\arabic{enumi}.}
\item
  \textbf{Build a tractable proxy in Stan and validate it under the
  standard pipeline.} This is what Model B does. The proxy is honest
  about being a proxy: it is not the cognitive theory, and its
  validation does not transfer.
\item
  \textbf{Fit Model A directly via simulator-based inference.} This is
  what NPE does. We do this in §``Neural Posterior Estimation for Model
  A.'' The chapter then compares the two posteriors as its main
  empirical contribution.
\end{enumerate}

This subsection builds Model B. The NPE subsection comes later because
pedagogically it makes sense to first see what the tractable proxy can
and cannot do.

\subsection{Model B's Generative
Structure}\label{model-bs-generative-structure}

Model B enumerates a fixed set of candidate rules and marginalises over
them analytically. This is identical in structure to the finite mixture
models of Chapter 8, where discrete latent variables (mixture
components) were summed out before any Stan code was written. Here the
discrete latent variable is ``which rule from the candidate set the
agent follows.'' The marginalised likelihood for trial \(i\) is:

\[p(y_i \mid x_i, \varepsilon) = \sum_{r=1}^{R} p(\text{rule} = r) \cdot p(y_i \mid x_i, \text{rule}_r, \varepsilon)\]

with a uniform prior over rules \(p(\text{rule} = r) = 1/R\),
implemented with \texttt{log\_sum\_exp} in Stan. Only the single
continuous parameter \(\varepsilon\) needs to be sampled.

We parameterise \(\varepsilon\) on the logit scale ---
\(\text{logit}(\varepsilon) \sim \text{Normal}(-2, 1)\) --- which is
unconstrained, avoids the boundary artefacts that afflict bounded
parameters in HMC, and places prior mass primarily on low error rates
(prior median \(\varepsilon \approx 0.12\)). This is the same
parameterisation philosophy as \texttt{log\_c} in Chapter 11.

\textbf{Critical structural facts about Model B that distinguish it from
Model A:}

\begin{itemize}
\tightlist
\item
  \textbf{No learning dynamics.} Model B treats all trials as
  conditionally exchangeable given \(\varepsilon\). There is no
  trial-by-trial weight update, no resampling, no rule discovery. The
  agent's behaviour is determined entirely by which rule it follows
  (drawn once from the candidate set) and the noise rate
  \(\varepsilon\).
\item
  \textbf{Fixed candidate rule set.} The \(R = 7\) rules are picked by
  inspecting the Kruschke stimulus layout (5 flat rules + 2 nested rules
  of the form \(A \text{ op } (B \text{ inner } C)\)). Phase 0 Q5
  surfaced this as a post-hoc choice; §``Discharging the Post-Hoc Rule
  Selection'' tests sensitivity.
\item
  \textbf{Closed-form Bernoulli likelihood.} Marginalised analytically
  with \texttt{log\_sum\_exp}. HMC handles it cleanly.
\item
  \textbf{No path-dependence.} Standard PSIS-LOO is technically valid
  (see §``LFO-CV Does Not Apply Here'').
\end{itemize}

\subsection{A First Caveat: The Prior Predictive Check Uses Model A's
Simulator}\label{a-first-caveat-the-prior-predictive-check-uses-model-as-simulator}

Before fitting Model B we want to verify that the prior on
\texttt{logit\_error\_prob} generates only scientifically plausible
categorisation curves. There is a subtle issue: \textbf{the Stan model
does not have a forward simulator with learning dynamics.} Model B is
conditionally i.i.d. in \(\varepsilon\), so its forward simulation is
just ``draw \(\varepsilon\), draw a rule, generate Bernoulli
responses.'' This produces flat learning curves at level
\(1 - \varepsilon\), which is empirically nothing like what the data
look like.

What earlier drafts of this chapter did --- and what we keep here, but
flag explicitly --- is to draw \(\varepsilon\) from the Stan prior and
then use \textbf{Model A's simulator} to generate the prior predictive
learning curves. This is \emph{not} a prior predictive check for Model B
in the strict sense. It is a prior predictive check for the cognitive
theory under Model B's \(\varepsilon\) prior. The two are not the same
and the reader should know which is being shown.

\begin{Shaded}
\begin{Highlighting}[]
\NormalTok{n\_ppc\_samples }\OtherTok{\textless{}{-}} \DecValTok{500}
\FunctionTok{set.seed}\NormalTok{(}\DecValTok{42}\NormalTok{)}

\NormalTok{ppc\_logit\_eps }\OtherTok{\textless{}{-}} \FunctionTok{rnorm}\NormalTok{(n\_ppc\_samples, LOGIT\_ERROR\_PRIOR\_MEAN, LOGIT\_ERROR\_PRIOR\_SD)}
\NormalTok{ppc\_eps       }\OtherTok{\textless{}{-}} \FunctionTok{plogis}\NormalTok{(ppc\_logit\_eps)}
\NormalTok{ppc\_sched     }\OtherTok{\textless{}{-}} \FunctionTok{make\_subject\_schedule}\NormalTok{(stimulus\_info, n\_blocks, }\AttributeTok{seed =} \DecValTok{999}\NormalTok{)}

\NormalTok{prior\_pred\_curves }\OtherTok{\textless{}{-}} \FunctionTok{map\_dfr}\NormalTok{(}\FunctionTok{seq\_len}\NormalTok{(n\_ppc\_samples), }\ControlFlowTok{function}\NormalTok{(s) \{}
\NormalTok{  sim }\OtherTok{\textless{}{-}} \FunctionTok{rule\_particle\_filter}\NormalTok{(}
    \AttributeTok{n\_particles =}\NormalTok{ N\_PARTICLES\_FIXED,}
    \AttributeTok{error\_prob  =}\NormalTok{ ppc\_eps[s],}
    \AttributeTok{obs         =} \FunctionTok{as.matrix}\NormalTok{(ppc\_sched[, }\FunctionTok{c}\NormalTok{(}\StringTok{"height"}\NormalTok{, }\StringTok{"position"}\NormalTok{)]),}
    \AttributeTok{cat\_one     =}\NormalTok{ ppc\_sched}\SpecialCharTok{$}\NormalTok{category\_feedback,}
    \AttributeTok{max\_dims    =}\NormalTok{ MAX\_RULE\_DIMS,}
    \AttributeTok{quiet       =} \ConstantTok{TRUE}
\NormalTok{  )}
  \FunctionTok{tibble}\NormalTok{(}
    \AttributeTok{sample\_id =}\NormalTok{ s,}
    \AttributeTok{trial     =} \FunctionTok{seq\_len}\NormalTok{(}\FunctionTok{nrow}\NormalTok{(ppc\_sched)),}
    \AttributeTok{correct   =} \FunctionTok{as.integer}\NormalTok{(sim}\SpecialCharTok{$}\NormalTok{sim\_response }\SpecialCharTok{==}\NormalTok{ ppc\_sched}\SpecialCharTok{$}\NormalTok{category\_feedback)}
\NormalTok{  ) }\SpecialCharTok{|\textgreater{}}
    \FunctionTok{mutate}\NormalTok{(}\AttributeTok{cum\_acc =} \FunctionTok{cumsum}\NormalTok{(correct) }\SpecialCharTok{/} \FunctionTok{row\_number}\NormalTok{())}
\NormalTok{\})}

\NormalTok{ppc\_summary }\OtherTok{\textless{}{-}}\NormalTok{ prior\_pred\_curves }\SpecialCharTok{|\textgreater{}}
  \FunctionTok{group\_by}\NormalTok{(trial) }\SpecialCharTok{|\textgreater{}}
  \FunctionTok{summarise}\NormalTok{(}
    \AttributeTok{q05 =} \FunctionTok{quantile}\NormalTok{(cum\_acc, }\FloatTok{0.05}\NormalTok{), }\AttributeTok{q25 =} \FunctionTok{quantile}\NormalTok{(cum\_acc, }\FloatTok{0.25}\NormalTok{),}
    \AttributeTok{q50 =} \FunctionTok{median}\NormalTok{(cum\_acc),}
    \AttributeTok{q75 =} \FunctionTok{quantile}\NormalTok{(cum\_acc, }\FloatTok{0.75}\NormalTok{), }\AttributeTok{q95 =} \FunctionTok{quantile}\NormalTok{(cum\_acc, }\FloatTok{0.95}\NormalTok{),}
    \AttributeTok{.groups =} \StringTok{"drop"}
\NormalTok{  )}

\FunctionTok{ggplot}\NormalTok{(ppc\_summary, }\FunctionTok{aes}\NormalTok{(}\AttributeTok{x =}\NormalTok{ trial)) }\SpecialCharTok{+}
  \FunctionTok{geom\_ribbon}\NormalTok{(}\FunctionTok{aes}\NormalTok{(}\AttributeTok{ymin =}\NormalTok{ q05, }\AttributeTok{ymax =}\NormalTok{ q95), }\AttributeTok{fill =} \StringTok{"\#E69F00"}\NormalTok{, }\AttributeTok{alpha =} \FloatTok{0.20}\NormalTok{) }\SpecialCharTok{+}
  \FunctionTok{geom\_ribbon}\NormalTok{(}\FunctionTok{aes}\NormalTok{(}\AttributeTok{ymin =}\NormalTok{ q25, }\AttributeTok{ymax =}\NormalTok{ q75), }\AttributeTok{fill =} \StringTok{"\#E69F00"}\NormalTok{, }\AttributeTok{alpha =} \FloatTok{0.40}\NormalTok{) }\SpecialCharTok{+}
  \FunctionTok{geom\_line}\NormalTok{(}\FunctionTok{aes}\NormalTok{(}\AttributeTok{y =}\NormalTok{ q50), }\AttributeTok{color =} \StringTok{"\#B07D00"}\NormalTok{, }\AttributeTok{linewidth =} \DecValTok{1}\NormalTok{) }\SpecialCharTok{+}
  \FunctionTok{geom\_hline}\NormalTok{(}\AttributeTok{yintercept =} \FloatTok{0.5}\NormalTok{, }\AttributeTok{linetype =} \StringTok{"dashed"}\NormalTok{, }\AttributeTok{color =} \StringTok{"grey40"}\NormalTok{) }\SpecialCharTok{+}
  \FunctionTok{scale\_y\_continuous}\NormalTok{(}\AttributeTok{limits =} \FunctionTok{c}\NormalTok{(}\DecValTok{0}\NormalTok{, }\DecValTok{1}\NormalTok{)) }\SpecialCharTok{+}
  \FunctionTok{labs}\NormalTok{(}
    \AttributeTok{title    =} \StringTok{"Prior Predictive Curves Under Model A\textquotesingle{}s Simulator With Model B\textquotesingle{}s ε Prior"}\NormalTok{,}
    \AttributeTok{subtitle =} \StringTok{"WARNING: This is NOT a prior predictive check for Model B in the strict sense.}\SpecialCharTok{\textbackslash{}n}\StringTok{Model B has no learning dynamics; its true forward sim would be flat at 1 {-} ε."}\NormalTok{,}
    \AttributeTok{x =} \StringTok{"Trial"}\NormalTok{, }\AttributeTok{y =} \StringTok{"Cumulative Accuracy"}
\NormalTok{  )}
\end{Highlighting}
\end{Shaded}

\begin{center}
\includegraphics[width=0.8\linewidth,height=\textheight,keepaspectratio]{13-CategorizationModelsRules_files/figure-pdf/ch14_prior_predictive_check-1.pdf}
\end{center}

\begin{quote}
\textbf{⚠️ Cross-model warning explicit.} The plot above shows prior
predictive learning curves drawn from the \emph{particle filter} (Model
A) using \(\varepsilon\) values sampled from the Stan model's (Model
B's) prior. The Stan model fitted to the same data assumes no learning,
so its true prior predictive distribution would consist of flat curves
at \(1 - \varepsilon\) --- nothing like the rising curves shown here. We
display these curves anyway because they are the most informative thing
to show pedagogically (they convey the range of behaviour the cognitive
theory produces under Model B's \(\varepsilon\) prior), but they are a
\emph{cross-model} prior predictive, not a pure within-model one. The
reader who wants the strict within-model check should compute it as: for
each sample \(s\), draw \(\varepsilon_s\), draw a rule index
\(r_s \in \{1, \ldots, 5\}\) uniformly, and generate Bernoulli responses
with probability \(\Pr(\text{cat 1} \mid x_i, r_s, \varepsilon_s)\). The
result is flat. It is also not very interesting.
\end{quote}

\subsection{Stan Model}\label{stan-model-1}

\begin{Shaded}
\begin{Highlighting}[]
\NormalTok{rule\_simplified\_stan }\OtherTok{\textless{}{-}} \StringTok{"}
\StringTok{// Model B: Marginalised Mixture over Fixed Candidate Rules}
\StringTok{// This is NOT the cognitive theory (Model A, the particle filter).}
\StringTok{// It is a tractable proxy that admits HMC validation.}
\StringTok{// Validation of Model B tests Model B, not Model A.}
\StringTok{//}
\StringTok{// Architecture (identical in principle to Ch. 8 mixture models):}
\StringTok{//   {-} Discrete latent variable = which rule the agent follows}
\StringTok{//   {-} Marginalized analytically with log\_sum\_exp over a fixed rule set}
\StringTok{//   {-} Only continuous parameter: logit\_error\_prob (unconstrained)}
\StringTok{//   {-} error\_prob = inv\_logit(logit\_error\_prob)}
\StringTok{//}
\StringTok{// No sequential learning in this model: it treats all trials as conditionally}
\StringTok{// exchangeable given error\_prob and the rule mixture.}
\StringTok{//}
\StringTok{// The 7 rules below cover the main decision structures visible in the}
\StringTok{// Kruschke stimulus layout, including two nested rules of the form}
\StringTok{// A op (B inner C) that cannot be expressed as flat 2D conjunctions/disjunctions.}
\StringTok{// See §\textquotesingle{}Discharging the Post{-}Hoc Rule Selection\textquotesingle{} for the sensitivity analysis.}

\StringTok{data \{}
\StringTok{  int\textless{}lower=1\textgreater{} ntrials;}
\StringTok{  int\textless{}lower=1\textgreater{} nfeatures;                       // must be 2 for this rule set}
\StringTok{  array[ntrials] int\textless{}lower=0, upper=1\textgreater{} y;       // observed choices}
\StringTok{  array[ntrials, nfeatures] real obs;           // obs[,1]=Height, obs[,2]=Position}

\StringTok{  // Prior hyperparameters (passed as data for transparency)}
\StringTok{  real prior\_logit\_error\_mean;}
\StringTok{  real\textless{}lower=0\textgreater{} prior\_logit\_error\_sd;}
\StringTok{\}}

\StringTok{parameters \{}
\StringTok{  real logit\_error\_prob;   // unconstrained; error\_prob = inv\_logit(logit\_error\_prob)}
\StringTok{\}}

\StringTok{transformed parameters \{}
\StringTok{  real\textless{}lower=0, upper=1\textgreater{} error\_prob = inv\_logit(logit\_error\_prob);}

\StringTok{  // Per{-}trial choice probability under each rule, computed once per leapfrog}
\StringTok{  // step so log\_lik in generated quantities is a one{-}line lookup.}
\StringTok{  // (Same architectural discipline as Ch 11 GCM and Ch 12 prototype.)}
\StringTok{  array[ntrials] real lp\_marginal;}

\StringTok{  for (i in 1:ntrials) \{}
\StringTok{    real h = obs[i, 1];}
\StringTok{    real p = obs[i, 2];}

\StringTok{    // {-}{-} Flat 1D and 2D rules (same as original 5) {-}{-}{-}{-}{-}{-}{-}{-}{-}{-}{-}{-}{-}{-}{-}{-}{-}{-}{-}{-}{-}{-}{-}{-}{-}{-}}
\StringTok{    int pred1 = h \textgreater{} 2.5 ? 1 : 0;}
\StringTok{    int pred2 = p \textless{} 2.5 ? 1 : 0;}
\StringTok{    int pred3 = (h \textgreater{} 2.0 \&\& p \textgreater{} 3.0) ? 1 : 0;}
\StringTok{    int pred4 = (h \textless{} 2.0 || p \textless{} 2.0) ? 0 : 1;}
\StringTok{    int pred5 = (h \textgreater{} 2.5 \&\& p \textless{} 3.5) ? 1 : 0;}
\StringTok{    // {-}{-} Nested rules: A op (B inner C) {-}{-}{-}{-}{-}{-}{-}{-}{-}{-}{-}{-}{-}{-}{-}{-}{-}{-}{-}{-}{-}{-}{-}{-}{-}{-}{-}{-}{-}{-}{-}{-}{-}{-}{-}{-}{-}}
\StringTok{    // pred6: (h \textgreater{} 3.5) OR ((h \textgreater{} 1.5) AND (p \textless{} 1.5))}
\StringTok{    // Cognitive reading: very tall OR some height AND very low position}
\StringTok{    // Perfectly classifies the 8 Kruschke stimuli.}
\StringTok{    int pred6 = (h \textgreater{} 3.5 || (h \textgreater{} 1.5 \&\& p \textless{} 1.5)) ? 1 : 0;}
\StringTok{    // pred7: (h \textgreater{} 1.5) AND ((p \textless{} 2.5) OR (h \textgreater{} 3.5))}
\StringTok{    // Cognitive reading: some height AND low position OR very tall}
\StringTok{    // Also perfectly classifies the 8 Kruschke stimuli.}
\StringTok{    int pred7 = (h \textgreater{} 1.5 \&\& (p \textless{} 2.5 || h \textgreater{} 3.5)) ? 1 : 0;}

\StringTok{    int R = 7;}

\StringTok{    vector[R] lp;}
\StringTok{    lp[1] = bernoulli\_lpmf(y[i] | pred1 == 1 ? 1 {-} error\_prob : error\_prob);}
\StringTok{    lp[2] = bernoulli\_lpmf(y[i] | pred2 == 1 ? 1 {-} error\_prob : error\_prob);}
\StringTok{    lp[3] = bernoulli\_lpmf(y[i] | pred3 == 1 ? 1 {-} error\_prob : error\_prob);}
\StringTok{    lp[4] = bernoulli\_lpmf(y[i] | pred4 == 1 ? 1 {-} error\_prob : error\_prob);}
\StringTok{    lp[5] = bernoulli\_lpmf(y[i] | pred5 == 1 ? 1 {-} error\_prob : error\_prob);}
\StringTok{    lp[6] = bernoulli\_lpmf(y[i] | pred6 == 1 ? 1 {-} error\_prob : error\_prob);}
\StringTok{    lp[7] = bernoulli\_lpmf(y[i] | pred7 == 1 ? 1 {-} error\_prob : error\_prob);}

\StringTok{    lp\_marginal[i] = log\_sum\_exp(lp) {-} log(R);}
\StringTok{  \}}
\StringTok{\}}

\StringTok{model \{}
\StringTok{  target += normal\_lpdf(logit\_error\_prob | prior\_logit\_error\_mean, prior\_logit\_error\_sd);}
\StringTok{  for (i in 1:ntrials) target += lp\_marginal[i];}
\StringTok{\}}

\StringTok{generated quantities \{}
\StringTok{  // log\_lik[i] is just a lookup of lp\_marginal[i] {-}{-} no recomputation.}
\StringTok{  array[ntrials] real log\_lik;}
\StringTok{  for (i in 1:ntrials) log\_lik[i] = lp\_marginal[i];}

\StringTok{  real lprior = normal\_lpdf(logit\_error\_prob | prior\_logit\_error\_mean, prior\_logit\_error\_sd);}
\StringTok{\}}
\StringTok{"}

\NormalTok{stan\_file\_rule }\OtherTok{\textless{}{-}} \FunctionTok{here}\NormalTok{(}\StringTok{"stan"}\NormalTok{, }\StringTok{"ch14\_rule\_simplified.stan"}\NormalTok{)}
\FunctionTok{write\_stan\_file}\NormalTok{(rule\_simplified\_stan, }\AttributeTok{dir =} \StringTok{"stan/"}\NormalTok{,}
                \AttributeTok{basename =} \StringTok{"ch14\_rule\_simplified.stan"}\NormalTok{)}
\end{Highlighting}
\end{Shaded}

\begin{verbatim}
[1] "/Users/au209589/Dropbox/Teaching/AdvancedCognitiveModeling23_book/stan/ch14_rule_simplified.stan"
\end{verbatim}

\begin{Shaded}
\begin{Highlighting}[]
\NormalTok{mod\_rule\_stan }\OtherTok{\textless{}{-}} \FunctionTok{cmdstan\_model}\NormalTok{(stan\_file\_rule, }\AttributeTok{force\_recompile =} \ConstantTok{TRUE}\NormalTok{)}
\FunctionTok{cat}\NormalTok{(}\StringTok{"Model B (Stan rule mixture) compiled successfully.}\SpecialCharTok{\textbackslash{}n}\StringTok{"}\NormalTok{)}
\end{Highlighting}
\end{Shaded}

\begin{verbatim}
Model B (Stan rule mixture) compiled successfully.
\end{verbatim}

\textbf{Key aspects of the Model B Stan implementation:}

\begin{itemize}
\tightlist
\item
  \textbf{Parameters block}: A single scalar \texttt{logit\_error\_prob}
  (unconstrained). The constraint \texttt{error\_prob\ ∈\ (0,1)} is
  enforced implicitly by \texttt{inv\_logit}.
\item
  \textbf{Transformed parameters}: The marginalised log-density per
  trial (\texttt{lp\_marginal{[}i{]}}) is computed once per leapfrog
  step. This is the same \texttt{transformed\ parameters} discipline as
  Chapter 11 (GCM) and Chapter 12 (Kalman prototype): if a quantity
  depends on parameters but is otherwise deterministic given the data,
  it belongs in \texttt{transformed\ parameters} so generated quantities
  can read it without recomputation.
\item
  \textbf{Model block}: Prior on \texttt{logit\_error\_prob} then
  accumulation of \texttt{lp\_marginal{[}i{]}} into the target.
\item
  \textbf{Generated quantities}: \texttt{log\_lik{[}i{]}} is a one-line
  lookup. Architecturally consistent with the other two chapters.
\end{itemize}

\subsection{Fitting Model B}\label{fitting-model-b}

\begin{Shaded}
\begin{Highlighting}[]
\CommentTok{\# Simulate responses from one Model A agent to use as \textquotesingle{}observed\textquotesingle{} data}
\NormalTok{fit\_schedule }\OtherTok{\textless{}{-}} \FunctionTok{make\_subject\_schedule}\NormalTok{(stimulus\_info, n\_blocks, }\AttributeTok{seed =} \DecValTok{42}\NormalTok{)}
\NormalTok{fit\_sim }\OtherTok{\textless{}{-}} \FunctionTok{rule\_particle\_filter}\NormalTok{(}
  \AttributeTok{n\_particles =}\NormalTok{ N\_PARTICLES\_FIXED,}
  \AttributeTok{error\_prob  =} \FloatTok{0.1}\NormalTok{,}
  \AttributeTok{obs         =} \FunctionTok{as.matrix}\NormalTok{(fit\_schedule[, }\FunctionTok{c}\NormalTok{(}\StringTok{"height"}\NormalTok{, }\StringTok{"position"}\NormalTok{)]),}
  \AttributeTok{cat\_one     =}\NormalTok{ fit\_schedule}\SpecialCharTok{$}\NormalTok{category\_feedback,}
  \AttributeTok{max\_dims    =}\NormalTok{ MAX\_RULE\_DIMS,}
  \AttributeTok{quiet       =} \ConstantTok{TRUE}
\NormalTok{)}

\NormalTok{rule\_data\_single }\OtherTok{\textless{}{-}} \FunctionTok{list}\NormalTok{(}
  \AttributeTok{ntrials                 =} \FunctionTok{nrow}\NormalTok{(fit\_schedule),}
  \AttributeTok{nfeatures               =} \DecValTok{2}\NormalTok{L,}
  \AttributeTok{y                       =} \FunctionTok{as.integer}\NormalTok{(fit\_sim}\SpecialCharTok{$}\NormalTok{sim\_response),}
  \AttributeTok{obs                     =} \FunctionTok{as.matrix}\NormalTok{(fit\_schedule[, }\FunctionTok{c}\NormalTok{(}\StringTok{"height"}\NormalTok{, }\StringTok{"position"}\NormalTok{)]),}
  \AttributeTok{prior\_logit\_error\_mean  =}\NormalTok{ LOGIT\_ERROR\_PRIOR\_MEAN,}
  \AttributeTok{prior\_logit\_error\_sd    =}\NormalTok{ LOGIT\_ERROR\_PRIOR\_SD}
\NormalTok{)}

\NormalTok{fit\_filepath\_rule }\OtherTok{\textless{}{-}} \FunctionTok{here}\NormalTok{(}\StringTok{"simmodels"}\NormalTok{, }\StringTok{"ch14\_rule\_single\_fit.rds"}\NormalTok{)}
\ControlFlowTok{if}\NormalTok{ (regenerate\_simulations }\SpecialCharTok{||} \SpecialCharTok{!}\FunctionTok{file.exists}\NormalTok{(fit\_filepath\_rule)) \{}
\NormalTok{  pf\_rule }\OtherTok{\textless{}{-}}\NormalTok{ mod\_rule\_stan}\SpecialCharTok{$}\FunctionTok{pathfinder}\NormalTok{(}
    \AttributeTok{data      =}\NormalTok{ rule\_data\_single,}
    \AttributeTok{seed      =} \DecValTok{42}\NormalTok{,}
    \AttributeTok{num\_paths =} \DecValTok{4}\NormalTok{,}
    \AttributeTok{refresh   =} \DecValTok{0}
\NormalTok{  )}

\NormalTok{  fit\_rule\_stan }\OtherTok{\textless{}{-}}\NormalTok{ mod\_rule\_stan}\SpecialCharTok{$}\FunctionTok{sample}\NormalTok{(}
    \AttributeTok{data            =}\NormalTok{ rule\_data\_single,}
    \AttributeTok{init            =}\NormalTok{ pf\_rule,}
    \AttributeTok{seed            =} \DecValTok{42}\NormalTok{,}
    \AttributeTok{chains          =} \DecValTok{4}\NormalTok{,}
    \AttributeTok{parallel\_chains =} \FunctionTok{min}\NormalTok{(}\DecValTok{4}\NormalTok{, }\FunctionTok{availableCores}\NormalTok{()),}
    \AttributeTok{iter\_warmup     =} \DecValTok{1000}\NormalTok{,}
    \AttributeTok{iter\_sampling   =} \DecValTok{1500}\NormalTok{,}
    \AttributeTok{refresh         =} \DecValTok{300}\NormalTok{,}
    \AttributeTok{adapt\_delta     =} \FloatTok{0.9}
\NormalTok{  )}
\NormalTok{  fit\_rule\_stan}\SpecialCharTok{$}\FunctionTok{save\_object}\NormalTok{(fit\_filepath\_rule)}
  \FunctionTok{cat}\NormalTok{(}\StringTok{"Model B fitted and saved.}\SpecialCharTok{\textbackslash{}n}\StringTok{"}\NormalTok{)}
\NormalTok{\} }\ControlFlowTok{else}\NormalTok{ \{}
  \ControlFlowTok{if}\NormalTok{ (}\FunctionTok{file.exists}\NormalTok{(fit\_filepath\_rule)) \{}
\NormalTok{    fit\_rule\_stan }\OtherTok{\textless{}{-}} \FunctionTok{readRDS}\NormalTok{(fit\_filepath\_rule)}
    \FunctionTok{cat}\NormalTok{(}\StringTok{"Loaded existing Model B fit.}\SpecialCharTok{\textbackslash{}n}\StringTok{"}\NormalTok{)}
\NormalTok{  \} }\ControlFlowTok{else}\NormalTok{ \{}
    \FunctionTok{cat}\NormalTok{(}\StringTok{"Fit file not found. Set regenerate\_simulations = TRUE.}\SpecialCharTok{\textbackslash{}n}\StringTok{"}\NormalTok{)}
\NormalTok{    fit\_rule\_stan }\OtherTok{\textless{}{-}} \ConstantTok{NULL}
\NormalTok{  \}}
\NormalTok{\}}
\end{Highlighting}
\end{Shaded}

\begin{verbatim}
Loaded existing Model B fit.
\end{verbatim}

\begin{Shaded}
\begin{Highlighting}[]
\ControlFlowTok{if}\NormalTok{ (}\SpecialCharTok{!}\FunctionTok{is.null}\NormalTok{(fit\_rule\_stan)) \{}
\NormalTok{  param\_summary }\OtherTok{\textless{}{-}}\NormalTok{ fit\_rule\_stan}\SpecialCharTok{$}\FunctionTok{summary}\NormalTok{(}
    \AttributeTok{variables =} \FunctionTok{c}\NormalTok{(}\StringTok{"logit\_error\_prob"}\NormalTok{, }\StringTok{"error\_prob"}\NormalTok{)}
\NormalTok{  )}
  \FunctionTok{print}\NormalTok{(param\_summary)}
\NormalTok{\}}
\end{Highlighting}
\end{Shaded}

\begin{verbatim}
# A tibble: 2 x 10
  variable     mean  median     sd    mad      q5    q95  rhat ess_bulk ess_tail
  <chr>       <dbl>   <dbl>  <dbl>  <dbl>   <dbl>  <dbl> <dbl>    <dbl>    <dbl>
1 logit_er~ -2.67   -2.60   0.814  0.784  -4.12   -1.48   1.00    1633.    1938.
2 error_pr~  0.0812  0.0691 0.0543 0.0486  0.0159  0.186  1.00    1633.    1938.
\end{verbatim}

\subsection{Mandatory MCMC Diagnostic
Battery}\label{mandatory-mcmc-diagnostic-battery}

The persona's Phase 3.0 says: curate the diagnostics, even when they are
nearly trivial. Use the same \texttt{diagnostic\_summary\_table()}
helper as Chapters 11 and 12.

\begin{Shaded}
\begin{Highlighting}[]
\ControlFlowTok{if}\NormalTok{ (}\SpecialCharTok{!}\FunctionTok{is.null}\NormalTok{(fit\_rule\_stan)) \{}
\NormalTok{  diag\_tbl }\OtherTok{\textless{}{-}} \FunctionTok{diagnostic\_summary\_table}\NormalTok{(}
\NormalTok{    fit\_rule\_stan,}
    \AttributeTok{params =} \FunctionTok{c}\NormalTok{(}\StringTok{"logit\_error\_prob"}\NormalTok{, }\StringTok{"error\_prob"}\NormalTok{)}
\NormalTok{  )}
  \FunctionTok{print}\NormalTok{(diag\_tbl)}

  \FunctionTok{cat}\NormalTok{(}\StringTok{"}\SpecialCharTok{\textbackslash{}n}\StringTok{Threshold reference (persona Phase 3.0):}\SpecialCharTok{\textbackslash{}n}\StringTok{"}\NormalTok{,}
      \StringTok{"  divergences        = 0}\SpecialCharTok{\textbackslash{}n}\StringTok{"}\NormalTok{,}
      \StringTok{"  rank{-}normalized Rhat \textless{} 1.01}\SpecialCharTok{\textbackslash{}n}\StringTok{"}\NormalTok{,}
      \StringTok{"  bulk ESS           \textgreater{} 400}\SpecialCharTok{\textbackslash{}n}\StringTok{"}\NormalTok{,}
      \StringTok{"  tail ESS           \textgreater{} 400}\SpecialCharTok{\textbackslash{}n}\StringTok{"}\NormalTok{,}
      \StringTok{"  E{-}BFMI             \textgreater{} 0.2}\SpecialCharTok{\textbackslash{}n}\StringTok{"}\NormalTok{,}
      \StringTok{"  MCSE\_mean / SD     \textless{} 0.05}\SpecialCharTok{\textbackslash{}n}\StringTok{"}\NormalTok{, }\AttributeTok{sep =} \StringTok{""}\NormalTok{)}

\NormalTok{  p\_trace }\OtherTok{\textless{}{-}}\NormalTok{ bayesplot}\SpecialCharTok{::}\FunctionTok{mcmc\_trace}\NormalTok{(}
\NormalTok{    fit\_rule\_stan}\SpecialCharTok{$}\FunctionTok{draws}\NormalTok{(}\FunctionTok{c}\NormalTok{(}\StringTok{"logit\_error\_prob"}\NormalTok{, }\StringTok{"error\_prob"}\NormalTok{)),}
    \AttributeTok{facet\_args =} \FunctionTok{list}\NormalTok{(}\AttributeTok{ncol =} \DecValTok{2}\NormalTok{)}
\NormalTok{  ) }\SpecialCharTok{+}
    \FunctionTok{ggtitle}\NormalTok{(}\StringTok{"Trace plots — should look like hairy caterpillars"}\NormalTok{)}

\NormalTok{  p\_rank }\OtherTok{\textless{}{-}}\NormalTok{ bayesplot}\SpecialCharTok{::}\FunctionTok{mcmc\_rank\_overlay}\NormalTok{(}
\NormalTok{    fit\_rule\_stan}\SpecialCharTok{$}\FunctionTok{draws}\NormalTok{(}\FunctionTok{c}\NormalTok{(}\StringTok{"logit\_error\_prob"}\NormalTok{, }\StringTok{"error\_prob"}\NormalTok{)),}
    \AttributeTok{facet\_args =} \FunctionTok{list}\NormalTok{(}\AttributeTok{ncol =} \DecValTok{2}\NormalTok{)}
\NormalTok{  ) }\SpecialCharTok{+}
    \FunctionTok{ggtitle}\NormalTok{(}\StringTok{"Rank histograms — should be uniform across chains"}\NormalTok{)}

  \FunctionTok{print}\NormalTok{(p\_trace }\SpecialCharTok{/}\NormalTok{ p\_rank)}
\NormalTok{\}}
\end{Highlighting}
\end{Shaded}

\begin{verbatim}
# A tibble: 6 x 4
  metric                           value threshold pass 
  <chr>                            <dbl> <chr>     <lgl>
1 Divergences (zero tolerance)    0      == 0      TRUE 
2 Max rank-normalised R-hat       1      < 1.01    TRUE 
3 Min bulk ESS                 1633      > 400     TRUE 
4 Min tail ESS                 1938      > 400     TRUE 
5 Min E-BFMI                      0.927  > 0.2     TRUE 
6 Max MCSE / posterior SD         0.0247 < 0.05    TRUE 

Threshold reference (persona Phase 3.0):
  divergences        = 0
  rank-normalized Rhat < 1.01
  bulk ESS           > 400
  tail ESS           > 400
  E-BFMI             > 0.2
  MCSE_mean / SD     < 0.05
\end{verbatim}

\begin{center}
\includegraphics[width=0.8\linewidth,height=\textheight,keepaspectratio]{13-CategorizationModelsRules_files/figure-pdf/ch14_diag_battery_model_B-1.pdf}
\end{center}

\subsection{Posterior Predictive Check (Model
B)}\label{posterior-predictive-check-model-b}

The posterior predictive check is honest about being a check on Model B,
not on the cognitive theory. We sample \(\varepsilon\) from the
posterior, draw a rule index \(r\) uniformly from the candidate set, and
generate Bernoulli responses. The result is what Model B \emph{actually}
predicts: a flat-mean response sequence with whatever variation the rule
mixture provides.

\begin{Shaded}
\begin{Highlighting}[]
\ControlFlowTok{if}\NormalTok{ (}\SpecialCharTok{!}\FunctionTok{is.null}\NormalTok{(fit\_rule\_stan)) \{}
\NormalTok{  draws }\OtherTok{\textless{}{-}}\NormalTok{ fit\_rule\_stan}\SpecialCharTok{$}\FunctionTok{draws}\NormalTok{(}\StringTok{"error\_prob"}\NormalTok{)}
\NormalTok{  eps\_draws\_df }\OtherTok{\textless{}{-}} \FunctionTok{as\_draws\_df}\NormalTok{(draws)}
\NormalTok{  eps\_draws  }\OtherTok{\textless{}{-}}\NormalTok{ eps\_draws\_df}\SpecialCharTok{$}\NormalTok{error\_prob}
\NormalTok{  n\_draws    }\OtherTok{\textless{}{-}} \FunctionTok{length}\NormalTok{(eps\_draws)}
\NormalTok{  obs\_mat    }\OtherTok{\textless{}{-}} \FunctionTok{as.matrix}\NormalTok{(fit\_schedule[, }\FunctionTok{c}\NormalTok{(}\StringTok{"height"}\NormalTok{, }\StringTok{"position"}\NormalTok{)])}
\NormalTok{  R\_rules    }\OtherTok{\textless{}{-}} \DecValTok{7}\NormalTok{L}
\NormalTok{  n\_obs      }\OtherTok{\textless{}{-}} \FunctionTok{nrow}\NormalTok{(fit\_schedule)}
  
  \CommentTok{\# Per{-}rule predictions on the observed stimulus sequence (7 rules: 5 flat + 2 nested)}
\NormalTok{  rule\_preds }\OtherTok{\textless{}{-}} \FunctionTok{matrix}\NormalTok{(}\DecValTok{0}\NormalTok{L, }\AttributeTok{nrow =}\NormalTok{ R\_rules, }\AttributeTok{ncol =}\NormalTok{ n\_obs)}
  \ControlFlowTok{for}\NormalTok{ (i }\ControlFlowTok{in} \FunctionTok{seq\_len}\NormalTok{(n\_obs)) \{}
\NormalTok{    h }\OtherTok{\textless{}{-}}\NormalTok{ obs\_mat[i, }\DecValTok{1}\NormalTok{]; p }\OtherTok{\textless{}{-}}\NormalTok{ obs\_mat[i, }\DecValTok{2}\NormalTok{]}
\NormalTok{    rule\_preds[}\DecValTok{1}\NormalTok{, i] }\OtherTok{\textless{}{-}} \FunctionTok{as.integer}\NormalTok{(h }\SpecialCharTok{\textgreater{}} \FloatTok{2.5}\NormalTok{)}
\NormalTok{    rule\_preds[}\DecValTok{2}\NormalTok{, i] }\OtherTok{\textless{}{-}} \FunctionTok{as.integer}\NormalTok{(p }\SpecialCharTok{\textless{}} \FloatTok{2.5}\NormalTok{)}
\NormalTok{    rule\_preds[}\DecValTok{3}\NormalTok{, i] }\OtherTok{\textless{}{-}} \FunctionTok{as.integer}\NormalTok{(h }\SpecialCharTok{\textgreater{}} \FloatTok{2.0} \SpecialCharTok{\&\&}\NormalTok{ p }\SpecialCharTok{\textgreater{}} \FloatTok{3.0}\NormalTok{)}
\NormalTok{    rule\_preds[}\DecValTok{4}\NormalTok{, i] }\OtherTok{\textless{}{-}} \FunctionTok{as.integer}\NormalTok{(}\SpecialCharTok{!}\NormalTok{(h }\SpecialCharTok{\textless{}} \FloatTok{2.0} \SpecialCharTok{||}\NormalTok{ p }\SpecialCharTok{\textless{}} \FloatTok{2.0}\NormalTok{))}
\NormalTok{    rule\_preds[}\DecValTok{5}\NormalTok{, i] }\OtherTok{\textless{}{-}} \FunctionTok{as.integer}\NormalTok{(h }\SpecialCharTok{\textgreater{}} \FloatTok{2.5} \SpecialCharTok{\&\&}\NormalTok{ p }\SpecialCharTok{\textless{}} \FloatTok{3.5}\NormalTok{)}
\NormalTok{    rule\_preds[}\DecValTok{6}\NormalTok{, i] }\OtherTok{\textless{}{-}} \FunctionTok{as.integer}\NormalTok{(h }\SpecialCharTok{\textgreater{}} \FloatTok{3.5} \SpecialCharTok{||}\NormalTok{ (h }\SpecialCharTok{\textgreater{}} \FloatTok{1.5} \SpecialCharTok{\&\&}\NormalTok{ p }\SpecialCharTok{\textless{}} \FloatTok{1.5}\NormalTok{))  }\CommentTok{\# nested: OR with inner AND}
\NormalTok{    rule\_preds[}\DecValTok{7}\NormalTok{, i] }\OtherTok{\textless{}{-}} \FunctionTok{as.integer}\NormalTok{(h }\SpecialCharTok{\textgreater{}} \FloatTok{1.5} \SpecialCharTok{\&\&}\NormalTok{ (p }\SpecialCharTok{\textless{}} \FloatTok{2.5} \SpecialCharTok{||}\NormalTok{ h }\SpecialCharTok{\textgreater{}} \FloatTok{3.5}\NormalTok{))  }\CommentTok{\# nested: AND with inner OR}
\NormalTok{  \}}

  \FunctionTok{set.seed}\NormalTok{(}\DecValTok{2026}\NormalTok{)}
\NormalTok{  yrep }\OtherTok{\textless{}{-}} \FunctionTok{matrix}\NormalTok{(}\DecValTok{0}\NormalTok{L, }\AttributeTok{nrow =}\NormalTok{ n\_draws, }\AttributeTok{ncol =}\NormalTok{ n\_obs)}
  \ControlFlowTok{for}\NormalTok{ (s }\ControlFlowTok{in} \FunctionTok{seq\_len}\NormalTok{(n\_draws)) \{}
\NormalTok{    r\_idx     }\OtherTok{\textless{}{-}} \FunctionTok{sample.int}\NormalTok{(R\_rules, }\DecValTok{1}\NormalTok{)}
\NormalTok{    rule\_pred }\OtherTok{\textless{}{-}}\NormalTok{ rule\_preds[r\_idx, ]}
\NormalTok{    p\_cat1    }\OtherTok{\textless{}{-}} \FunctionTok{ifelse}\NormalTok{(rule\_pred }\SpecialCharTok{==} \DecValTok{1}\NormalTok{, }\DecValTok{1} \SpecialCharTok{{-}}\NormalTok{ eps\_draws[s], eps\_draws[s])}
\NormalTok{    yrep[s, ] }\OtherTok{\textless{}{-}} \FunctionTok{rbinom}\NormalTok{(n\_obs, }\DecValTok{1}\NormalTok{, p\_cat1)}
\NormalTok{  \}}

  \CommentTok{\# Cumulative{-}accuracy ribbons vs. observed}
\NormalTok{  obs\_correct }\OtherTok{\textless{}{-}} \FunctionTok{as.integer}\NormalTok{(fit\_sim}\SpecialCharTok{$}\NormalTok{sim\_response }\SpecialCharTok{==}\NormalTok{ fit\_schedule}\SpecialCharTok{$}\NormalTok{category\_feedback)}
\NormalTok{  obs\_cum\_acc }\OtherTok{\textless{}{-}} \FunctionTok{cumsum}\NormalTok{(obs\_correct) }\SpecialCharTok{/} \FunctionTok{seq\_len}\NormalTok{(n\_obs)}
\NormalTok{  yrep\_correct }\OtherTok{\textless{}{-}} \FunctionTok{sweep}\NormalTok{(yrep, }\DecValTok{2}\NormalTok{,}
                        \FunctionTok{as.numeric}\NormalTok{(fit\_schedule}\SpecialCharTok{$}\NormalTok{category\_feedback),}
                        \AttributeTok{FUN =} \ControlFlowTok{function}\NormalTok{(a, b) }\FunctionTok{as.numeric}\NormalTok{(a }\SpecialCharTok{==}\NormalTok{ b))}
\NormalTok{  yrep\_correct }\OtherTok{\textless{}{-}} \FunctionTok{matrix}\NormalTok{(yrep\_correct, }\AttributeTok{nrow =}\NormalTok{ n\_draws, }\AttributeTok{ncol =}\NormalTok{ n\_obs)}

\NormalTok{  yrep\_cum\_acc }\OtherTok{\textless{}{-}} \FunctionTok{t}\NormalTok{(}\FunctionTok{apply}\NormalTok{(yrep\_correct, }\DecValTok{1}\NormalTok{, }\ControlFlowTok{function}\NormalTok{(r) }\FunctionTok{cumsum}\NormalTok{(r) }\SpecialCharTok{/} \FunctionTok{seq\_along}\NormalTok{(r)))}

\NormalTok{  cum\_acc\_summary }\OtherTok{\textless{}{-}} \FunctionTok{tibble}\NormalTok{(}
    \AttributeTok{trial =} \FunctionTok{seq\_len}\NormalTok{(n\_obs),}
    \AttributeTok{q05   =} \FunctionTok{apply}\NormalTok{(yrep\_cum\_acc, }\DecValTok{2}\NormalTok{, quantile, }\FloatTok{0.05}\NormalTok{),}
    \AttributeTok{q25   =} \FunctionTok{apply}\NormalTok{(yrep\_cum\_acc, }\DecValTok{2}\NormalTok{, quantile, }\FloatTok{0.25}\NormalTok{),}
    \AttributeTok{q50   =} \FunctionTok{apply}\NormalTok{(yrep\_cum\_acc, }\DecValTok{2}\NormalTok{, quantile, }\FloatTok{0.50}\NormalTok{),}
    \AttributeTok{q75   =} \FunctionTok{apply}\NormalTok{(yrep\_cum\_acc, }\DecValTok{2}\NormalTok{, quantile, }\FloatTok{0.75}\NormalTok{),}
    \AttributeTok{q95   =} \FunctionTok{apply}\NormalTok{(yrep\_cum\_acc, }\DecValTok{2}\NormalTok{, quantile, }\FloatTok{0.95}\NormalTok{),}
    \AttributeTok{obs   =}\NormalTok{ obs\_cum\_acc}
\NormalTok{  )}

  \FunctionTok{ggplot}\NormalTok{(cum\_acc\_summary, }\FunctionTok{aes}\NormalTok{(}\AttributeTok{x =}\NormalTok{ trial)) }\SpecialCharTok{+}
    \FunctionTok{geom\_ribbon}\NormalTok{(}\FunctionTok{aes}\NormalTok{(}\AttributeTok{ymin =}\NormalTok{ q05, }\AttributeTok{ymax =}\NormalTok{ q95), }\AttributeTok{fill =} \StringTok{"\#56B4E9"}\NormalTok{, }\AttributeTok{alpha =} \FloatTok{0.20}\NormalTok{) }\SpecialCharTok{+}
    \FunctionTok{geom\_ribbon}\NormalTok{(}\FunctionTok{aes}\NormalTok{(}\AttributeTok{ymin =}\NormalTok{ q25, }\AttributeTok{ymax =}\NormalTok{ q75), }\AttributeTok{fill =} \StringTok{"\#56B4E9"}\NormalTok{, }\AttributeTok{alpha =} \FloatTok{0.40}\NormalTok{) }\SpecialCharTok{+}
    \FunctionTok{geom\_line}\NormalTok{(}\FunctionTok{aes}\NormalTok{(}\AttributeTok{y =}\NormalTok{ q50), }\AttributeTok{color =} \StringTok{"\#0072B2"}\NormalTok{, }\AttributeTok{linewidth =} \DecValTok{1}\NormalTok{) }\SpecialCharTok{+}
    \FunctionTok{geom\_line}\NormalTok{(}\FunctionTok{aes}\NormalTok{(}\AttributeTok{y =}\NormalTok{ obs), }\AttributeTok{color =} \StringTok{"\#D55E00"}\NormalTok{, }\AttributeTok{linewidth =} \DecValTok{1}\NormalTok{) }\SpecialCharTok{+}
    \FunctionTok{geom\_hline}\NormalTok{(}\AttributeTok{yintercept =} \FloatTok{0.5}\NormalTok{, }\AttributeTok{linetype =} \StringTok{"dashed"}\NormalTok{, }\AttributeTok{color =} \StringTok{"grey50"}\NormalTok{) }\SpecialCharTok{+}
    \FunctionTok{labs}\NormalTok{(}
      \AttributeTok{title    =} \StringTok{"Posterior Predictive Check (Model B)"}\NormalTok{,}
      \AttributeTok{subtitle =} \StringTok{"Blue: Model B\textquotesingle{}s posterior predictive}\SpecialCharTok{\textbackslash{}n}\StringTok{Orange: observed (generated by Model A)}\SpecialCharTok{\textbackslash{}n}\StringTok{The gap is informative — Model B has no learning curve."}\NormalTok{,}
      \AttributeTok{x =} \StringTok{"Trial"}\NormalTok{, }\AttributeTok{y =} \StringTok{"Cumulative accuracy"}
\NormalTok{    ) }\SpecialCharTok{+}
    \FunctionTok{coord\_cartesian}\NormalTok{(}\AttributeTok{ylim =} \FunctionTok{c}\NormalTok{(}\FloatTok{0.4}\NormalTok{, }\DecValTok{1}\NormalTok{))}
\NormalTok{\}}
\end{Highlighting}
\end{Shaded}

\begin{center}
\includegraphics[width=0.8\linewidth,height=\textheight,keepaspectratio]{13-CategorizationModelsRules_files/figure-pdf/ch14_ppc_model_B-1.pdf}
\end{center}

\textbf{Reading this plot.} The blue ribbons are Model B's actual
posterior predictive cumulative accuracy. They are \textbf{flat across
trials}, because Model B has no learning dynamics. The orange line is
the observed data, generated by Model A, which \emph{does} have learning
dynamics. \textbf{The systematic gap is the structural misspecification
of Model B}, made visible. This is exactly what the chapter wants the
reader to see: Model B can recover an \(\varepsilon\) value that fits
the \emph{average} response, but it cannot reproduce the trial-by-trial
structure because that structure comes from a process Model B does not
have.

\subsection{Prior--Posterior Update Plot (Model
B)}\label{priorposterior-update-plot-model-b}

\begin{Shaded}
\begin{Highlighting}[]
\ControlFlowTok{if}\NormalTok{ (}\SpecialCharTok{!}\FunctionTok{is.null}\NormalTok{(fit\_rule\_stan)) \{}
\NormalTok{  draws\_df }\OtherTok{\textless{}{-}} \FunctionTok{as\_draws\_df}\NormalTok{(fit\_rule\_stan}\SpecialCharTok{$}\FunctionTok{draws}\NormalTok{(}\FunctionTok{c}\NormalTok{(}\StringTok{"logit\_error\_prob"}\NormalTok{, }\StringTok{"error\_prob"}\NormalTok{)))}

  \FunctionTok{set.seed}\NormalTok{(}\DecValTok{7}\NormalTok{)}
\NormalTok{  n\_prior }\OtherTok{\textless{}{-}} \FunctionTok{nrow}\NormalTok{(draws\_df)}
\NormalTok{  prior\_logit\_eps }\OtherTok{\textless{}{-}} \FunctionTok{rnorm}\NormalTok{(n\_prior, LOGIT\_ERROR\_PRIOR\_MEAN, LOGIT\_ERROR\_PRIOR\_SD)}
\NormalTok{  prior\_eps       }\OtherTok{\textless{}{-}} \FunctionTok{plogis}\NormalTok{(prior\_logit\_eps)}

\NormalTok{  prior\_post\_df }\OtherTok{\textless{}{-}} \FunctionTok{bind\_rows}\NormalTok{(}
    \FunctionTok{tibble}\NormalTok{(}\AttributeTok{scale =} \StringTok{"logit\_error\_prob (computational)"}\NormalTok{,}
           \AttributeTok{value =}\NormalTok{ prior\_logit\_eps, }\AttributeTok{source =} \StringTok{"prior"}\NormalTok{),}
    \FunctionTok{tibble}\NormalTok{(}\AttributeTok{scale =} \StringTok{"logit\_error\_prob (computational)"}\NormalTok{,}
           \AttributeTok{value =}\NormalTok{ draws\_df}\SpecialCharTok{$}\NormalTok{logit\_error\_prob, }\AttributeTok{source =} \StringTok{"posterior"}\NormalTok{),}
    \FunctionTok{tibble}\NormalTok{(}\AttributeTok{scale =} \StringTok{"error\_prob (interpretable)"}\NormalTok{,}
           \AttributeTok{value =}\NormalTok{ prior\_eps, }\AttributeTok{source =} \StringTok{"prior"}\NormalTok{),}
    \FunctionTok{tibble}\NormalTok{(}\AttributeTok{scale =} \StringTok{"error\_prob (interpretable)"}\NormalTok{,}
           \AttributeTok{value =}\NormalTok{ draws\_df}\SpecialCharTok{$}\NormalTok{error\_prob, }\AttributeTok{source =} \StringTok{"posterior"}\NormalTok{)}
\NormalTok{  )}

\NormalTok{  true\_lines }\OtherTok{\textless{}{-}} \FunctionTok{tibble}\NormalTok{(}
    \AttributeTok{scale =} \FunctionTok{c}\NormalTok{(}\StringTok{"logit\_error\_prob (computational)"}\NormalTok{, }\StringTok{"error\_prob (interpretable)"}\NormalTok{),}
    \AttributeTok{value =} \FunctionTok{c}\NormalTok{(}\FunctionTok{qlogis}\NormalTok{(}\FloatTok{0.1}\NormalTok{), }\FloatTok{0.1}\NormalTok{)}
\NormalTok{  )}

  \FunctionTok{ggplot}\NormalTok{(prior\_post\_df, }\FunctionTok{aes}\NormalTok{(}\AttributeTok{x =}\NormalTok{ value, }\AttributeTok{fill =}\NormalTok{ source, }\AttributeTok{color =}\NormalTok{ source)) }\SpecialCharTok{+}
    \FunctionTok{geom\_density}\NormalTok{(}\AttributeTok{alpha =} \FloatTok{0.4}\NormalTok{, }\AttributeTok{linewidth =} \FloatTok{0.6}\NormalTok{) }\SpecialCharTok{+}
    \FunctionTok{geom\_vline}\NormalTok{(}\AttributeTok{data =}\NormalTok{ true\_lines,}
               \FunctionTok{aes}\NormalTok{(}\AttributeTok{xintercept =}\NormalTok{ value),}
               \AttributeTok{color =} \StringTok{"red"}\NormalTok{, }\AttributeTok{linetype =} \StringTok{"dashed"}\NormalTok{) }\SpecialCharTok{+}
    \FunctionTok{facet\_wrap}\NormalTok{(}\SpecialCharTok{\textasciitilde{}}\NormalTok{ scale, }\AttributeTok{scales =} \StringTok{"free"}\NormalTok{) }\SpecialCharTok{+}
    \FunctionTok{scale\_fill\_manual}\NormalTok{(}\AttributeTok{values =} \FunctionTok{c}\NormalTok{(}\AttributeTok{prior =} \StringTok{"\#999999"}\NormalTok{, }\AttributeTok{posterior =} \StringTok{"\#E69F00"}\NormalTok{)) }\SpecialCharTok{+}
    \FunctionTok{scale\_color\_manual}\NormalTok{(}\AttributeTok{values =} \FunctionTok{c}\NormalTok{(}\AttributeTok{prior =} \StringTok{"\#666666"}\NormalTok{, }\AttributeTok{posterior =} \StringTok{"\#B07D00"}\NormalTok{)) }\SpecialCharTok{+}
    \FunctionTok{labs}\NormalTok{(}
      \AttributeTok{title    =} \StringTok{"Prior → Posterior Update for Model B"}\NormalTok{,}
      \AttributeTok{subtitle =} \StringTok{"Red dashed line = Model A\textquotesingle{}s true generating ε = 0.1}\SpecialCharTok{\textbackslash{}n}\StringTok{Gap between posterior centre and the dashed line is the structural{-}bias signature."}\NormalTok{,}
      \AttributeTok{x =} \ConstantTok{NULL}\NormalTok{, }\AttributeTok{y =} \StringTok{"Density"}
\NormalTok{    )}
\NormalTok{\}}
\end{Highlighting}
\end{Shaded}

\begin{center}
\includegraphics[width=0.8\linewidth,height=\textheight,keepaspectratio]{13-CategorizationModelsRules_files/figure-pdf/ch14_prior_posterior_update-1.pdf}
\end{center}

The posterior should be substantially narrower than the prior. The
relationship of the posterior centre to the dashed red line (the true
\(\varepsilon\) used to generate the data) is the \textbf{structural
bias signature}: if the dashed line falls in the bulk of the posterior,
the proxy is approximately faithful for this data; if the dashed line is
in the tail or outside the posterior, the proxy is biased and the size
of the offset is informative about how badly the candidate rule set
misses the converged rule.

\subsection{Randomised LOO-PIT (Model
B)}\label{randomised-loo-pit-model-b}

The same calibration check we used in Chapters 11 and 12, applied to
Model B.

\begin{Shaded}
\begin{Highlighting}[]
\ControlFlowTok{if}\NormalTok{ (}\SpecialCharTok{!}\FunctionTok{is.null}\NormalTok{(fit\_rule\_stan)) \{}
\NormalTok{  log\_lik\_mat }\OtherTok{\textless{}{-}}\NormalTok{ fit\_rule\_stan}\SpecialCharTok{$}\FunctionTok{draws}\NormalTok{(}\StringTok{"log\_lik"}\NormalTok{, }\AttributeTok{format =} \StringTok{"matrix"}\NormalTok{)}
\NormalTok{  psis\_obj    }\OtherTok{\textless{}{-}}\NormalTok{ loo}\SpecialCharTok{::}\FunctionTok{psis}\NormalTok{(log\_lik\_mat, }\AttributeTok{cores =} \DecValTok{4}\NormalTok{)}
  
\NormalTok{  log\_w }\OtherTok{\textless{}{-}}\NormalTok{ loo}\SpecialCharTok{::}\FunctionTok{weights.importance\_sampling}\NormalTok{(psis\_obj,}
                                             \AttributeTok{normalize =} \ConstantTok{TRUE}\NormalTok{,}
                                             \AttributeTok{log =} \ConstantTok{TRUE}\NormalTok{)}

  \CommentTok{\# Compute LOO predictive probability of cat 1 from Model B}
  \CommentTok{\# For this we need per{-}trial p\_cat1 under each posterior draw and rule.}
  \CommentTok{\# Easiest: sample rule per draw, compute p\_cat1 = (1 {-} eps if rule says 1 else eps),}
  \CommentTok{\# then average using LOO weights.}
\NormalTok{  eps\_draws\_full }\OtherTok{\textless{}{-}} \FunctionTok{as\_draws\_df}\NormalTok{(fit\_rule\_stan}\SpecialCharTok{$}\FunctionTok{draws}\NormalTok{(}\StringTok{"error\_prob"}\NormalTok{))}\SpecialCharTok{$}\NormalTok{error\_prob}
\NormalTok{  n\_draws\_full   }\OtherTok{\textless{}{-}} \FunctionTok{length}\NormalTok{(eps\_draws\_full)}
\NormalTok{  n\_obs          }\OtherTok{\textless{}{-}} \FunctionTok{nrow}\NormalTok{(fit\_schedule)}

  \FunctionTok{set.seed}\NormalTok{(}\DecValTok{11}\NormalTok{)}
\NormalTok{  rule\_idx\_per\_draw }\OtherTok{\textless{}{-}} \FunctionTok{sample.int}\NormalTok{(}\DecValTok{7}\NormalTok{L, n\_draws\_full, }\AttributeTok{replace =} \ConstantTok{TRUE}\NormalTok{)}

  \CommentTok{\# Per (draw, trial) p\_cat1}
\NormalTok{  p\_draws }\OtherTok{\textless{}{-}} \FunctionTok{matrix}\NormalTok{(}\DecValTok{0}\NormalTok{, }\AttributeTok{nrow =}\NormalTok{ n\_draws\_full, }\AttributeTok{ncol =}\NormalTok{ n\_obs)}
  \ControlFlowTok{for}\NormalTok{ (s }\ControlFlowTok{in} \FunctionTok{seq\_len}\NormalTok{(n\_draws\_full)) \{}
\NormalTok{    r }\OtherTok{\textless{}{-}}\NormalTok{ rule\_idx\_per\_draw[s]}
\NormalTok{    pred\_row }\OtherTok{\textless{}{-}}\NormalTok{ rule\_preds[r, ]}
\NormalTok{    p\_draws[s, ] }\OtherTok{\textless{}{-}} \FunctionTok{ifelse}\NormalTok{(pred\_row }\SpecialCharTok{==} \DecValTok{1}\NormalTok{,}
                            \DecValTok{1} \SpecialCharTok{{-}}\NormalTok{ eps\_draws\_full[s],}
\NormalTok{                            eps\_draws\_full[s])}
\NormalTok{  \}}

\NormalTok{  loo\_p\_cat1 }\OtherTok{\textless{}{-}} \FunctionTok{vapply}\NormalTok{(}\FunctionTok{seq\_len}\NormalTok{(n\_obs), }\ControlFlowTok{function}\NormalTok{(i) \{}
    \FunctionTok{sum}\NormalTok{(}\FunctionTok{exp}\NormalTok{(log\_w[, i]) }\SpecialCharTok{*}\NormalTok{ p\_draws[, i])}
\NormalTok{  \}, }\FunctionTok{numeric}\NormalTok{(}\DecValTok{1}\NormalTok{))}

\NormalTok{  y\_obs }\OtherTok{\textless{}{-}} \FunctionTok{as.integer}\NormalTok{(fit\_sim}\SpecialCharTok{$}\NormalTok{sim\_response)}
  \FunctionTok{set.seed}\NormalTok{(}\DecValTok{2026}\NormalTok{)}
\NormalTok{  rloo\_pit }\OtherTok{\textless{}{-}} \FunctionTok{vapply}\NormalTok{(}\FunctionTok{seq\_along}\NormalTok{(y\_obs), }\ControlFlowTok{function}\NormalTok{(i) \{}
\NormalTok{    p\_i }\OtherTok{\textless{}{-}}\NormalTok{ loo\_p\_cat1[i]}
    \ControlFlowTok{if}\NormalTok{ (y\_obs[i] }\SpecialCharTok{==} \DecValTok{1}\NormalTok{) }\FunctionTok{runif}\NormalTok{(}\DecValTok{1}\NormalTok{, }\DecValTok{1} \SpecialCharTok{{-}}\NormalTok{ p\_i, }\DecValTok{1}\NormalTok{) }\ControlFlowTok{else} \FunctionTok{runif}\NormalTok{(}\DecValTok{1}\NormalTok{, }\DecValTok{0}\NormalTok{, }\DecValTok{1} \SpecialCharTok{{-}}\NormalTok{ p\_i)}
\NormalTok{  \}, }\FunctionTok{numeric}\NormalTok{(}\DecValTok{1}\NormalTok{))}

  \FunctionTok{ggplot}\NormalTok{(}\FunctionTok{tibble}\NormalTok{(}\AttributeTok{rloo\_pit =}\NormalTok{ rloo\_pit), }\FunctionTok{aes}\NormalTok{(}\AttributeTok{x =}\NormalTok{ rloo\_pit)) }\SpecialCharTok{+}
    \FunctionTok{geom\_histogram}\NormalTok{(}\FunctionTok{aes}\NormalTok{(}\AttributeTok{y =} \FunctionTok{after\_stat}\NormalTok{(density)), }\AttributeTok{bins =} \DecValTok{20}\NormalTok{,}
                   \AttributeTok{fill =} \StringTok{"\#E69F00"}\NormalTok{, }\AttributeTok{color =} \StringTok{"white"}\NormalTok{) }\SpecialCharTok{+}
    \FunctionTok{geom\_hline}\NormalTok{(}\AttributeTok{yintercept =} \DecValTok{1}\NormalTok{, }\AttributeTok{linetype =} \StringTok{"dashed"}\NormalTok{, }\AttributeTok{color =} \StringTok{"red"}\NormalTok{) }\SpecialCharTok{+}
    \FunctionTok{labs}\NormalTok{(}
      \AttributeTok{title    =} \StringTok{"Randomised LOO{-}PIT (Model B)"}\NormalTok{,}
      \AttributeTok{subtitle =} \StringTok{"If Model B is calibrated for this data, the histogram should look uniform on [0, 1]"}\NormalTok{,}
      \AttributeTok{x =} \StringTok{"Randomised PIT value"}\NormalTok{, }\AttributeTok{y =} \StringTok{"Density"}
\NormalTok{    ) }\SpecialCharTok{+}
    \FunctionTok{coord\_cartesian}\NormalTok{(}\AttributeTok{ylim =} \FunctionTok{c}\NormalTok{(}\DecValTok{0}\NormalTok{, }\FloatTok{2.5}\NormalTok{))}
\NormalTok{\}}
\end{Highlighting}
\end{Shaded}

\begin{center}
\includegraphics[width=0.8\linewidth,height=\textheight,keepaspectratio]{13-CategorizationModelsRules_files/figure-pdf/ch14_loo_pit_randomised-1.pdf}
\end{center}

A roughly uniform histogram says Model B is calibrated \emph{for the
proxy task} (predicting individual binary responses given the posterior
over \(\varepsilon\) and the rule mixture). It does \emph{not} say Model
B is calibrated as a model of the cognitive process --- that would
require the NPE pipeline to validate Model A directly.

\subsection{LOO with Pareto-k as a Misspecification
Diagnostic}\label{loo-with-pareto-k-as-a-misspecification-diagnostic}

\begin{Shaded}
\begin{Highlighting}[]
\ControlFlowTok{if}\NormalTok{ (}\SpecialCharTok{!}\FunctionTok{is.null}\NormalTok{(fit\_rule\_stan)) \{}
\NormalTok{  loo\_rule }\OtherTok{\textless{}{-}}\NormalTok{ loo}\SpecialCharTok{::}\FunctionTok{loo}\NormalTok{(log\_lik\_mat, }\AttributeTok{cores =} \DecValTok{4}\NormalTok{)}
  \FunctionTok{print}\NormalTok{(loo\_rule)}
  \FunctionTok{plot}\NormalTok{(loo\_rule, }\AttributeTok{label\_points =} \ConstantTok{TRUE}\NormalTok{,}
       \AttributeTok{main =} \StringTok{"PSIS{-}LOO Pareto{-}k: Model B"}\NormalTok{)}
\NormalTok{\}}
\end{Highlighting}
\end{Shaded}

\begin{verbatim}

Computed from 6000 by 64 log-likelihood matrix.

         Estimate  SE
elpd_loo    -32.6 3.1
p_loo         0.2 0.0
looic        65.2 6.3
------
MCSE of elpd_loo is 0.0.
MCSE and ESS estimates assume independent draws (r_eff=1).

All Pareto k estimates are good (k < 0.7).
See help('pareto-k-diagnostic') for details.
\end{verbatim}

\begin{center}
\includegraphics[width=0.8\linewidth,height=\textheight,keepaspectratio]{13-CategorizationModelsRules_files/figure-pdf/ch14_loo_rule-1.pdf}
\end{center}

Unlike Chapters 11 and 12, here the LOO ELPD is \textbf{technically
valid} (Model B is conditionally exchangeable --- see next subsection).
But the Pareto-\(\hat{k}\) values are still informative as a
\emph{misspecification diagnostic}: when Model B is misspecified
relative to the true generating process (Model A), \(\hat{k}\) values
are likely to be elevated on the trials where the proxy fails most
badly. Plotting \(\hat{k}\) against trial number and stimulus identity
localises the misspecification.

\begin{Shaded}
\begin{Highlighting}[]
\ControlFlowTok{if}\NormalTok{ (}\SpecialCharTok{!}\FunctionTok{is.null}\NormalTok{(fit\_rule\_stan)) \{}
\NormalTok{  k\_values }\OtherTok{\textless{}{-}}\NormalTok{ loo\_rule}\SpecialCharTok{$}\NormalTok{diagnostics}\SpecialCharTok{$}\NormalTok{pareto\_k}
\NormalTok{  k\_df }\OtherTok{\textless{}{-}} \FunctionTok{tibble}\NormalTok{(}
    \AttributeTok{trial      =} \FunctionTok{seq\_along}\NormalTok{(k\_values),}
    \AttributeTok{k\_hat      =}\NormalTok{ k\_values,}
    \AttributeTok{stimulus   =}\NormalTok{ fit\_schedule}\SpecialCharTok{$}\NormalTok{stimulus\_id,}
    \AttributeTok{height     =}\NormalTok{ fit\_schedule}\SpecialCharTok{$}\NormalTok{height,}
    \AttributeTok{position   =}\NormalTok{ fit\_schedule}\SpecialCharTok{$}\NormalTok{position,}
    \AttributeTok{true\_cat   =}\NormalTok{ fit\_schedule}\SpecialCharTok{$}\NormalTok{category\_feedback,}
    \AttributeTok{obs\_y      =}\NormalTok{ fit\_sim}\SpecialCharTok{$}\NormalTok{sim\_response}
\NormalTok{  )}

\NormalTok{  p\_trial }\OtherTok{\textless{}{-}} \FunctionTok{ggplot}\NormalTok{(k\_df, }\FunctionTok{aes}\NormalTok{(}\AttributeTok{x =}\NormalTok{ trial, }\AttributeTok{y =}\NormalTok{ k\_hat)) }\SpecialCharTok{+}
    \FunctionTok{geom\_point}\NormalTok{(}\FunctionTok{aes}\NormalTok{(}\AttributeTok{color =} \FunctionTok{factor}\NormalTok{(true\_cat)), }\AttributeTok{size =} \FloatTok{1.5}\NormalTok{) }\SpecialCharTok{+}
    \FunctionTok{geom\_hline}\NormalTok{(}\AttributeTok{yintercept =} \FunctionTok{c}\NormalTok{(}\FloatTok{0.5}\NormalTok{, }\FloatTok{0.7}\NormalTok{, }\FloatTok{1.0}\NormalTok{),}
               \AttributeTok{linetype =} \FunctionTok{c}\NormalTok{(}\StringTok{"dotted"}\NormalTok{, }\StringTok{"dashed"}\NormalTok{, }\StringTok{"solid"}\NormalTok{),}
               \AttributeTok{color =} \StringTok{"red"}\NormalTok{, }\AttributeTok{alpha =} \FloatTok{0.6}\NormalTok{) }\SpecialCharTok{+}
    \FunctionTok{scale\_color\_manual}\NormalTok{(}\AttributeTok{values =} \FunctionTok{c}\NormalTok{(}\StringTok{"0"} \OtherTok{=} \StringTok{"\#0072B2"}\NormalTok{, }\StringTok{"1"} \OtherTok{=} \StringTok{"\#D55E00"}\NormalTok{),}
                       \AttributeTok{name =} \StringTok{"True cat"}\NormalTok{) }\SpecialCharTok{+}
    \FunctionTok{labs}\NormalTok{(}
      \AttributeTok{title    =} \StringTok{"Pareto{-}k by trial"}\NormalTok{,}
      \AttributeTok{subtitle =} \StringTok{"Red lines: 0.5 (monitor), 0.7 (problematic), 1.0 (failed)"}\NormalTok{,}
      \AttributeTok{x =} \StringTok{"Trial"}\NormalTok{, }\AttributeTok{y =} \StringTok{"Pareto k"}
\NormalTok{    )}

\NormalTok{  p\_stim }\OtherTok{\textless{}{-}} \FunctionTok{ggplot}\NormalTok{(k\_df, }\FunctionTok{aes}\NormalTok{(}\AttributeTok{x =} \FunctionTok{factor}\NormalTok{(stimulus), }\AttributeTok{y =}\NormalTok{ k\_hat)) }\SpecialCharTok{+}
    \FunctionTok{geom\_boxplot}\NormalTok{(}\AttributeTok{fill =} \StringTok{"\#E69F00"}\NormalTok{, }\AttributeTok{alpha =} \FloatTok{0.6}\NormalTok{) }\SpecialCharTok{+}
    \FunctionTok{geom\_jitter}\NormalTok{(}\AttributeTok{width =} \FloatTok{0.15}\NormalTok{, }\AttributeTok{alpha =} \FloatTok{0.4}\NormalTok{, }\AttributeTok{size =} \DecValTok{1}\NormalTok{) }\SpecialCharTok{+}
    \FunctionTok{geom\_hline}\NormalTok{(}\AttributeTok{yintercept =} \FloatTok{0.7}\NormalTok{, }\AttributeTok{linetype =} \StringTok{"dashed"}\NormalTok{, }\AttributeTok{color =} \StringTok{"red"}\NormalTok{) }\SpecialCharTok{+}
    \FunctionTok{labs}\NormalTok{(}
      \AttributeTok{title    =} \StringTok{"Pareto{-}k by stimulus identity"}\NormalTok{,}
      \AttributeTok{subtitle =} \StringTok{"If a particular stimulus has consistently elevated k\_hat,}\SpecialCharTok{\textbackslash{}n}\StringTok{that stimulus is where the rule set fails"}\NormalTok{,}
      \AttributeTok{x =} \StringTok{"Stimulus ID"}\NormalTok{, }\AttributeTok{y =} \StringTok{"Pareto k"}
\NormalTok{    )}

  \FunctionTok{print}\NormalTok{(p\_trial }\SpecialCharTok{/}\NormalTok{ p\_stim)}
\NormalTok{\}}
\end{Highlighting}
\end{Shaded}

\begin{center}
\includegraphics[width=0.8\linewidth,height=\textheight,keepaspectratio]{13-CategorizationModelsRules_files/figure-pdf/ch14_pareto_k_by_trial_and_stimulus-1.pdf}
\end{center}

\textbf{Reading the Pareto-k plots as a misspecification diagnostic.} If
the high-\(\hat{k}\) values cluster on the \emph{boundary stimuli} of
the Kruschke set --- the ones whose true category cannot be determined
by any single rule from the candidate set --- that is direct evidence
that the rule set is missing the relevant boundary structure. If the
high-\(\hat{k}\) values cluster at \emph{late trials}, that is a signal
that Model A's particle filter has converged to a rule the candidate set
does not contain, and Model B is using \(\varepsilon\) inflation to
absorb the systematic late-trial errors. In both cases, the diagnostic
does not just say ``the model is wrong'' --- it says \textbf{where the
model is wrong}, which is much more useful.

\subsection{LFO-CV Does Not Apply
Here}\label{lfo-cv-does-not-apply-here}

\begin{quote}
\textbf{A back-reference, not a caveat to repeat.} Chapters 11 and 12
implement LFO-CV via \texttt{psis\_lfo\_gcm()} and
\texttt{psis\_lfo\_kalman()} because the GCM and the Kalman prototype
are path-dependent: trial \(t\)'s likelihood depends on the entire
history of trials \(1, \ldots, t-1\) through the memory or filter state.
\textbf{Model B has no such path-dependence.} All trials are
conditionally independent given \(\varepsilon\) and the rule mixture:
leaving out trial \(t\) does not affect the likelihood of any other
trial. Standard PSIS-LOO is therefore \emph{exactly correct} for Model B
and the LFO-CV machinery is unnecessary. We say this once here so the
reader does not encounter the same caveat for the third time after
seeing it twice in the previous chapters.

A subtler point: when Chapter 15 compares Model B against the GCM
(Chapter 11) and the Kalman prototype (Chapter 12), the comparison
should be reported under each model's \emph{correct} validation tool ---
LFO-CV for the path-dependent ones, standard LOO for Model B. ELPD
differences across these tools are not directly commensurable on a fine
scale, but they are commensurable enough to detect large effects. The
bigger concern in Chapter 15 is the cross-architecture mismatch (Model B
isn't even fitting the cognitive theory), and that concern is what the
NPE subsection below resolves.
\end{quote}

\subsection{Prior Sensitivity
Analysis}\label{prior-sensitivity-analysis-3}

\begin{Shaded}
\begin{Highlighting}[]
\ControlFlowTok{if}\NormalTok{ (}\SpecialCharTok{!}\FunctionTok{is.null}\NormalTok{(fit\_rule\_stan)) \{}
  \CommentTok{\# Get all log\_lik variable names (log\_lik[1], log\_lik[2], ...)}
\NormalTok{  all\_params }\OtherTok{\textless{}{-}}\NormalTok{ fit\_rule\_stan}\SpecialCharTok{$}\FunctionTok{metadata}\NormalTok{()}\SpecialCharTok{$}\NormalTok{model\_params}
\NormalTok{  log\_lik\_vars }\OtherTok{\textless{}{-}}\NormalTok{ all\_params[}\FunctionTok{grep}\NormalTok{(}\StringTok{"\^{}log\_lik"}\NormalTok{, all\_params)]}
  
  \CommentTok{\# Compute sensitivity}
\NormalTok{  ps }\OtherTok{\textless{}{-}}\NormalTok{ priorsense}\SpecialCharTok{::}\FunctionTok{powerscale\_sensitivity}\NormalTok{(fit\_rule\_stan, }\AttributeTok{variable =} \FunctionTok{c}\NormalTok{(}\StringTok{"error\_prob"}\NormalTok{), }\AttributeTok{log\_lik =}\NormalTok{ log\_lik\_vars)}
  \FunctionTok{print}\NormalTok{(ps)}
  
  \CommentTok{\# Try plotting if log\_lik can be found automatically by create\_priorsense\_data}
  \FunctionTok{try}\NormalTok{(priorsense}\SpecialCharTok{::}\FunctionTok{powerscale\_plot\_dens}\NormalTok{(fit\_rule\_stan, }\AttributeTok{variables =} \FunctionTok{c}\NormalTok{(}\StringTok{"error\_prob"}\NormalTok{)))}
\NormalTok{\}}
\end{Highlighting}
\end{Shaded}

\begin{verbatim}
Sensitivity based on cjs_dist
Prior selection: all priors
Likelihood selection: all data

   variable prior likelihood                     diagnosis
 error_prob 0.212      0.051 potential prior-data conflict
\end{verbatim}

\begin{center}
\includegraphics[width=0.8\linewidth,height=\textheight,keepaspectratio]{13-CategorizationModelsRules_files/figure-pdf/ch14_prior_sensitivity-1.pdf}
\end{center}

\texttt{powerscale\_sensitivity()} reports a sensitivity index \(\psi\)
for each parameter. \textbf{One important limitation:} it perturbs only
the \emph{explicit} prior on \texttt{logit\_error\_prob}. It cannot
detect sensitivity to the \emph{implicit} prior over the rule set, which
is the much bigger source of structural bias for Model B. Low \(\psi\)
here means Model B is robust to the parametric prior on \(\varepsilon\)
--- it does \emph{not} mean Model B is robust to the choice of candidate
rules. The next subsection addresses that.

\begin{center}\rule{0.5\linewidth}{0.5pt}\end{center}

\section{Parameter Recovery (Model B)}\label{parameter-recovery-model-b}

We test recovery systematically across the plausible range of
\texttt{error\_prob}, with multiple independent datasets per true value,
each with its own randomized stimulus schedule. \textbf{Caveat upfront:}
the data are generated by Model A, fitted by Model B. If Model A
converges to a rule that is not in Model B's candidate set, recovery
will be biased. This is the structural-bias signature again.

\begin{Shaded}
\begin{Highlighting}[]
\NormalTok{run\_rule\_recovery }\OtherTok{\textless{}{-}} \ControlFlowTok{function}\NormalTok{(true\_error\_prob, iteration,}
\NormalTok{                              stimulus\_info, n\_blocks) \{}
\NormalTok{  subject\_seed }\OtherTok{\textless{}{-}}\NormalTok{ iteration }\SpecialCharTok{*} \DecValTok{1000} \SpecialCharTok{+} \FunctionTok{round}\NormalTok{(true\_error\_prob }\SpecialCharTok{*} \DecValTok{1000}\NormalTok{)}
\NormalTok{  sched        }\OtherTok{\textless{}{-}} \FunctionTok{make\_subject\_schedule}\NormalTok{(stimulus\_info, n\_blocks, }\AttributeTok{seed =}\NormalTok{ subject\_seed)}

\NormalTok{  sim }\OtherTok{\textless{}{-}} \FunctionTok{rule\_particle\_filter}\NormalTok{(}
    \AttributeTok{n\_particles =}\NormalTok{ N\_PARTICLES\_FIXED,}
    \AttributeTok{error\_prob  =}\NormalTok{ true\_error\_prob,}
    \AttributeTok{obs         =} \FunctionTok{as.matrix}\NormalTok{(sched[, }\FunctionTok{c}\NormalTok{(}\StringTok{"height"}\NormalTok{, }\StringTok{"position"}\NormalTok{)]),}
    \AttributeTok{cat\_one     =}\NormalTok{ sched}\SpecialCharTok{$}\NormalTok{category\_feedback,}
    \AttributeTok{max\_dims    =}\NormalTok{ MAX\_RULE\_DIMS,}
    \AttributeTok{quiet       =} \ConstantTok{TRUE}
\NormalTok{  )}

\NormalTok{  stan\_dat }\OtherTok{\textless{}{-}} \FunctionTok{list}\NormalTok{(}
    \AttributeTok{ntrials                =} \FunctionTok{nrow}\NormalTok{(sched),}
    \AttributeTok{nfeatures              =} \DecValTok{2}\NormalTok{L,}
    \AttributeTok{y                      =} \FunctionTok{as.integer}\NormalTok{(sim}\SpecialCharTok{$}\NormalTok{sim\_response),}
    \AttributeTok{obs                    =} \FunctionTok{as.matrix}\NormalTok{(sched[, }\FunctionTok{c}\NormalTok{(}\StringTok{"height"}\NormalTok{, }\StringTok{"position"}\NormalTok{)]),}
    \AttributeTok{prior\_logit\_error\_mean =}\NormalTok{ LOGIT\_ERROR\_PRIOR\_MEAN,}
    \AttributeTok{prior\_logit\_error\_sd   =}\NormalTok{ LOGIT\_ERROR\_PRIOR\_SD}
\NormalTok{  )}

\NormalTok{  pf }\OtherTok{\textless{}{-}} \FunctionTok{tryCatch}\NormalTok{(}
\NormalTok{    mod\_rule\_stan}\SpecialCharTok{$}\FunctionTok{pathfinder}\NormalTok{(}\AttributeTok{data =}\NormalTok{ stan\_dat, }\AttributeTok{num\_paths =} \DecValTok{4}\NormalTok{,}
                              \AttributeTok{refresh =} \DecValTok{0}\NormalTok{, }\AttributeTok{seed =}\NormalTok{ subject\_seed),}
    \AttributeTok{error =} \ControlFlowTok{function}\NormalTok{(e) }\ConstantTok{NULL}
\NormalTok{  )}
\NormalTok{  init\_val }\OtherTok{\textless{}{-}} \ControlFlowTok{if}\NormalTok{ (}\SpecialCharTok{!}\FunctionTok{is.null}\NormalTok{(pf)) pf }\ControlFlowTok{else} \FloatTok{0.5}

\NormalTok{  fit }\OtherTok{\textless{}{-}} \FunctionTok{tryCatch}\NormalTok{(}
\NormalTok{    mod\_rule\_stan}\SpecialCharTok{$}\FunctionTok{sample}\NormalTok{(}
      \AttributeTok{data =}\NormalTok{ stan\_dat, }\AttributeTok{init =}\NormalTok{ init\_val,}
      \AttributeTok{seed =}\NormalTok{ subject\_seed, }\AttributeTok{chains =} \DecValTok{2}\NormalTok{, }\AttributeTok{parallel\_chains =} \DecValTok{2}\NormalTok{,}
      \AttributeTok{iter\_warmup =} \DecValTok{800}\NormalTok{, }\AttributeTok{iter\_sampling =} \DecValTok{1000}\NormalTok{,}
      \AttributeTok{refresh =} \DecValTok{0}\NormalTok{, }\AttributeTok{adapt\_delta =} \FloatTok{0.9}
\NormalTok{    ),}
    \AttributeTok{error =} \ControlFlowTok{function}\NormalTok{(e) \{ }\FunctionTok{cat}\NormalTok{(}\StringTok{"Fit error:"}\NormalTok{, e}\SpecialCharTok{$}\NormalTok{message, }\StringTok{"}\SpecialCharTok{\textbackslash{}n}\StringTok{"}\NormalTok{); }\ConstantTok{NULL}\NormalTok{ \}}
\NormalTok{  )}
  \ControlFlowTok{if}\NormalTok{ (}\FunctionTok{is.null}\NormalTok{(fit)) }\FunctionTok{return}\NormalTok{(}\ConstantTok{NULL}\NormalTok{)}

\NormalTok{  draws\_eps }\OtherTok{\textless{}{-}}\NormalTok{ fit}\SpecialCharTok{$}\FunctionTok{draws}\NormalTok{(}\StringTok{"error\_prob"}\NormalTok{, }\AttributeTok{format =} \StringTok{"df"}\NormalTok{)}
  \FunctionTok{tibble}\NormalTok{(}
    \AttributeTok{true\_error\_prob   =}\NormalTok{ true\_error\_prob,}
    \AttributeTok{iteration         =}\NormalTok{ iteration,}
    \AttributeTok{estimated\_eps\_mean =} \FunctionTok{mean}\NormalTok{(draws\_eps}\SpecialCharTok{$}\NormalTok{error\_prob),}
    \AttributeTok{est\_lower\_95      =} \FunctionTok{quantile}\NormalTok{(draws\_eps}\SpecialCharTok{$}\NormalTok{error\_prob, }\FloatTok{0.025}\NormalTok{),}
    \AttributeTok{est\_upper\_95      =} \FunctionTok{quantile}\NormalTok{(draws\_eps}\SpecialCharTok{$}\NormalTok{error\_prob, }\FloatTok{0.975}\NormalTok{),}
    \AttributeTok{had\_divergence    =} \FunctionTok{any}\NormalTok{(fit}\SpecialCharTok{$}\FunctionTok{diagnostic\_summary}\NormalTok{(}\AttributeTok{quiet =} \ConstantTok{TRUE}\NormalTok{)}\SpecialCharTok{$}\NormalTok{num\_divergent }\SpecialCharTok{\textgreater{}} \DecValTok{0}\NormalTok{)}
\NormalTok{  )}
\NormalTok{\}}

\NormalTok{eps\_values\_to\_test }\OtherTok{\textless{}{-}} \FunctionTok{c}\NormalTok{(}\FloatTok{0.05}\NormalTok{, }\FloatTok{0.1}\NormalTok{, }\FloatTok{0.15}\NormalTok{, }\FloatTok{0.2}\NormalTok{, }\FloatTok{0.3}\NormalTok{)}
\NormalTok{n\_recovery\_iters   }\OtherTok{\textless{}{-}} \DecValTok{5}   \CommentTok{\# increase to 20–50 for publication{-}quality recovery}

\NormalTok{recovery\_file\_rule }\OtherTok{\textless{}{-}} \FunctionTok{here}\NormalTok{(}\StringTok{"simdata"}\NormalTok{, }\StringTok{"ch14\_rule\_recovery\_results.rds"}\NormalTok{)}
\ControlFlowTok{if}\NormalTok{ (regenerate\_simulations }\SpecialCharTok{||} \SpecialCharTok{!}\FunctionTok{file.exists}\NormalTok{(recovery\_file\_rule)) \{}
  \FunctionTok{cat}\NormalTok{(}\StringTok{"Running parameter recovery for Model B...}\SpecialCharTok{\textbackslash{}n}\StringTok{"}\NormalTok{)}
\NormalTok{  recovery\_plan }\OtherTok{\textless{}{-}} \FunctionTok{expand\_grid}\NormalTok{(}
    \AttributeTok{true\_error\_prob =}\NormalTok{ eps\_values\_to\_test,}
    \AttributeTok{iteration       =} \FunctionTok{seq\_len}\NormalTok{(n\_recovery\_iters)}
\NormalTok{  )}
\NormalTok{  rule\_recovery\_results }\OtherTok{\textless{}{-}} \FunctionTok{future\_pmap\_dfr}\NormalTok{(}
    \FunctionTok{list}\NormalTok{(}\AttributeTok{true\_error\_prob =}\NormalTok{ recovery\_plan}\SpecialCharTok{$}\NormalTok{true\_error\_prob,}
         \AttributeTok{iteration       =}\NormalTok{ recovery\_plan}\SpecialCharTok{$}\NormalTok{iteration),}
\NormalTok{    run\_rule\_recovery,}
    \AttributeTok{stimulus\_info =}\NormalTok{ stimulus\_info,}
    \AttributeTok{n\_blocks      =}\NormalTok{ n\_blocks,}
    \AttributeTok{.options =} \FunctionTok{furrr\_options}\NormalTok{(}\AttributeTok{seed =} \ConstantTok{TRUE}\NormalTok{)}
\NormalTok{  )}
  \FunctionTok{saveRDS}\NormalTok{(rule\_recovery\_results, recovery\_file\_rule)}
  \FunctionTok{cat}\NormalTok{(}\StringTok{"Recovery complete and saved.}\SpecialCharTok{\textbackslash{}n}\StringTok{"}\NormalTok{)}
\NormalTok{\} }\ControlFlowTok{else}\NormalTok{ \{}
\NormalTok{  rule\_recovery\_results }\OtherTok{\textless{}{-}} \FunctionTok{readRDS}\NormalTok{(recovery\_file\_rule)}
  \FunctionTok{cat}\NormalTok{(}\StringTok{"Recovery results loaded.}\SpecialCharTok{\textbackslash{}n}\StringTok{"}\NormalTok{)}
\NormalTok{\}}
\end{Highlighting}
\end{Shaded}

\begin{verbatim}
Recovery results loaded.
\end{verbatim}

\begin{Shaded}
\begin{Highlighting}[]
\FunctionTok{ggplot}\NormalTok{(rule\_recovery\_results,}
       \FunctionTok{aes}\NormalTok{(}\AttributeTok{x =}\NormalTok{ true\_error\_prob, }\AttributeTok{y =}\NormalTok{ estimated\_eps\_mean)) }\SpecialCharTok{+}
  \FunctionTok{geom\_abline}\NormalTok{(}\AttributeTok{intercept =} \DecValTok{0}\NormalTok{, }\AttributeTok{slope =} \DecValTok{1}\NormalTok{, }\AttributeTok{linetype =} \StringTok{"dashed"}\NormalTok{, }\AttributeTok{color =} \StringTok{"grey50"}\NormalTok{) }\SpecialCharTok{+}
  \FunctionTok{geom\_errorbar}\NormalTok{(}\FunctionTok{aes}\NormalTok{(}\AttributeTok{ymin =}\NormalTok{ est\_lower\_95, }\AttributeTok{ymax =}\NormalTok{ est\_upper\_95,}
                    \AttributeTok{color =}\NormalTok{ had\_divergence),}
                \AttributeTok{width =} \FloatTok{0.005}\NormalTok{, }\AttributeTok{alpha =} \FloatTok{0.7}\NormalTok{) }\SpecialCharTok{+}
  \FunctionTok{geom\_point}\NormalTok{(}\FunctionTok{aes}\NormalTok{(}\AttributeTok{color =}\NormalTok{ had\_divergence), }\AttributeTok{size =} \DecValTok{3}\NormalTok{) }\SpecialCharTok{+}
  \FunctionTok{scale\_color\_manual}\NormalTok{(}\AttributeTok{values =} \FunctionTok{c}\NormalTok{(}\StringTok{"FALSE"} \OtherTok{=} \StringTok{"\#0072B2"}\NormalTok{, }\StringTok{"TRUE"} \OtherTok{=} \StringTok{"\#D55E00"}\NormalTok{),}
                     \AttributeTok{name =} \StringTok{"Had divergence"}\NormalTok{) }\SpecialCharTok{+}
  \FunctionTok{labs}\NormalTok{(}
    \AttributeTok{title    =} \StringTok{"Parameter Recovery: Model B (error\_prob)"}\NormalTok{,}
    \AttributeTok{subtitle =} \FunctionTok{paste0}\NormalTok{(}\StringTok{"Points: posterior mean; bars: 95\% CrI; n="}\NormalTok{,}
\NormalTok{                      n\_recovery\_iters,}
                      \StringTok{" datasets per true value.}\SpecialCharTok{\textbackslash{}n}\StringTok{"}\NormalTok{,}
                      \StringTok{"Data generated by Model A; fitted by Model B."}\NormalTok{),}
    \AttributeTok{x =} \StringTok{"True error\_prob"}\NormalTok{, }\AttributeTok{y =} \StringTok{"Estimated error\_prob"}
\NormalTok{  )}
\end{Highlighting}
\end{Shaded}

\begin{center}
\includegraphics[width=0.8\linewidth,height=\textheight,keepaspectratio]{13-CategorizationModelsRules_files/figure-pdf/ch14_parameter_recovery-1.pdf}
\end{center}

\textbf{Discussion of Recovery Results.} Recovery of
\texttt{error\_prob} depends on whether Model B's candidate rule set
contains rules that match what Model A converges to on each replicate.
When the converged rule is in the candidate set, recovery is good. When
it is not, Model B compensates by inflating \(\varepsilon\) to absorb
the systematic misclassifications, and the posterior mean shifts up ---
the points drift above the diagonal. This drift is the
\emph{quantitative} version of the structural-bias warning.

\begin{center}\rule{0.5\linewidth}{0.5pt}\end{center}

\section{Discharging the Post-Hoc Rule
Selection}\label{discharging-the-post-hoc-rule-selection}

Phase 0 Q5 surfaced the fact that Model B's candidate rules are picked
by inspecting the stimulus layout, which makes them empirical, not
prior. We discharge this in two complementary ways.

\subsection{Option 1: Sensitivity to the Candidate Rule
Set}\label{option-1-sensitivity-to-the-candidate-rule-set}

We define three alternative candidate rule sets and refit Model B on the
same data:

\begin{itemize}
\tightlist
\item
  \textbf{Set A} --- the 7 hand-picked rules (the ones in the Stan code
  above: 5 flat + 2 nested).
\item
  \textbf{Set B} --- 5 rules generated \emph{blind} by running Model A
  on a separately-simulated training dataset and taking its top-weighted
  rules at the end of training. The training data is generated with
  \(\varepsilon = 0.1\) on a different schedule than the fitting data.
  This is the ``honest'' version of the candidate set.
\item
  \textbf{Set C} --- 5 rules deliberately chosen to be \emph{bad} (rules
  that disagree with most stimuli), as a stress test.
\end{itemize}

\begin{Shaded}
\begin{Highlighting}[]
\CommentTok{\# Helper to build a Stan model from a list of (predictor function) rules}
\NormalTok{build\_stan\_with\_rules }\OtherTok{\textless{}{-}} \ControlFlowTok{function}\NormalTok{(rule\_predictors, name\_suffix) \{}
\NormalTok{  R }\OtherTok{\textless{}{-}} \FunctionTok{length}\NormalTok{(rule\_predictors)}
\NormalTok{  rule\_lines }\OtherTok{\textless{}{-}} \FunctionTok{paste}\NormalTok{(}
    \FunctionTok{sprintf}\NormalTok{(}\StringTok{"    int pred\%d = \%s;"}\NormalTok{, }\FunctionTok{seq\_len}\NormalTok{(R), rule\_predictors),}
    \AttributeTok{collapse =} \StringTok{"}\SpecialCharTok{\textbackslash{}n}\StringTok{"}
\NormalTok{  )}
\NormalTok{  lp\_lines\_model }\OtherTok{\textless{}{-}} \FunctionTok{paste}\NormalTok{(}
    \FunctionTok{sprintf}\NormalTok{(}\StringTok{"    lp[\%d] = bernoulli\_lpmf(y[i] | pred\%d == 1 ? 1 {-} error\_prob : error\_prob);"}\NormalTok{,}
            \FunctionTok{seq\_len}\NormalTok{(R), }\FunctionTok{seq\_len}\NormalTok{(R)),}
    \AttributeTok{collapse =} \StringTok{"}\SpecialCharTok{\textbackslash{}n}\StringTok{"}
\NormalTok{  )}

\NormalTok{  stan\_code }\OtherTok{\textless{}{-}} \FunctionTok{sprintf}\NormalTok{(}\StringTok{"}
\StringTok{data \{}
\StringTok{  int\textless{}lower=1\textgreater{} ntrials;}
\StringTok{  int\textless{}lower=1\textgreater{} nfeatures;}
\StringTok{  array[ntrials] int\textless{}lower=0, upper=1\textgreater{} y;}
\StringTok{  array[ntrials, nfeatures] real obs;}
\StringTok{  real prior\_logit\_error\_mean;}
\StringTok{  real\textless{}lower=0\textgreater{} prior\_logit\_error\_sd;}
\StringTok{\}}
\StringTok{parameters \{ real logit\_error\_prob; \}}
\StringTok{transformed parameters \{}
\StringTok{  real\textless{}lower=0, upper=1\textgreater{} error\_prob = inv\_logit(logit\_error\_prob);}
\StringTok{  array[ntrials] real lp\_marginal;}
\StringTok{  for (i in 1:ntrials) \{}
\StringTok{    real h = obs[i, 1];}
\StringTok{    real p = obs[i, 2];}
\StringTok{\%s}
\StringTok{    int R = \%d;}
\StringTok{    vector[R] lp;}
\StringTok{\%s}
\StringTok{    lp\_marginal[i] = log\_sum\_exp(lp) {-} log(R);}
\StringTok{  \}}
\StringTok{\}}
\StringTok{model \{}
\StringTok{  target += normal\_lpdf(logit\_error\_prob | prior\_logit\_error\_mean, prior\_logit\_error\_sd);}
\StringTok{  for (i in 1:ntrials) target += lp\_marginal[i];}
\StringTok{\}}
\StringTok{generated quantities \{}
\StringTok{  array[ntrials] real log\_lik;}
\StringTok{  for (i in 1:ntrials) log\_lik[i] = lp\_marginal[i];}
\StringTok{  real lprior = normal\_lpdf(logit\_error\_prob | prior\_logit\_error\_mean, prior\_logit\_error\_sd);}
\StringTok{\}}
\StringTok{"}\NormalTok{, rule\_lines, R, lp\_lines\_model)}

\NormalTok{  fname }\OtherTok{\textless{}{-}} \FunctionTok{here}\NormalTok{(}\StringTok{"stan"}\NormalTok{, }\FunctionTok{sprintf}\NormalTok{(}\StringTok{"ch14\_rule\_\%s.stan"}\NormalTok{,  name\_suffix))}
  \FunctionTok{write\_stan\_file}\NormalTok{(stan\_code, }\AttributeTok{dir =} \StringTok{"stan/"}\NormalTok{,}
                  \AttributeTok{basename =} \FunctionTok{sprintf}\NormalTok{(}\StringTok{"ch14\_rule\_\%s.stan"}\NormalTok{, name\_suffix))}
  \FunctionTok{cmdstan\_model}\NormalTok{(fname, }\AttributeTok{force\_recompile =} \ConstantTok{TRUE}\NormalTok{)}
\NormalTok{\}}

\CommentTok{\# Set A: hand{-}picked rules — 5 flat + 2 nested (replicates Model B exactly)}
\NormalTok{rules\_A }\OtherTok{\textless{}{-}} \FunctionTok{c}\NormalTok{(}
  \StringTok{"h \textgreater{} 2.5 ? 1 : 0"}\NormalTok{,}
  \StringTok{"p \textless{} 2.5 ? 1 : 0"}\NormalTok{,}
  \StringTok{"(h \textgreater{} 2.0 \&\& p \textgreater{} 3.0) ? 1 : 0"}\NormalTok{,}
  \StringTok{"(h \textless{} 2.0 || p \textless{} 2.0) ? 0 : 1"}\NormalTok{,}
  \StringTok{"(h \textgreater{} 2.5 \&\& p \textless{} 3.5) ? 1 : 0"}\NormalTok{,}
  \StringTok{"(h \textgreater{} 3.5 || (h \textgreater{} 1.5 \&\& p \textless{} 1.5)) ? 1 : 0"}\NormalTok{,}
  \StringTok{"(h \textgreater{} 1.5 \&\& (p \textless{} 2.5 || h \textgreater{} 3.5)) ? 1 : 0"}
\NormalTok{)}

\CommentTok{\# Set B: blind rules from a fresh Model A run on a separate seed}
\NormalTok{get\_blind\_rules }\OtherTok{\textless{}{-}} \ControlFlowTok{function}\NormalTok{(}\AttributeTok{seed =} \DecValTok{7}\NormalTok{) \{}
\NormalTok{  blind\_sched }\OtherTok{\textless{}{-}} \FunctionTok{make\_subject\_schedule}\NormalTok{(stimulus\_info, n\_blocks, }\AttributeTok{seed =} \DecValTok{50000} \SpecialCharTok{+}\NormalTok{ seed)}
\NormalTok{  obs\_b   }\OtherTok{\textless{}{-}} \FunctionTok{as.matrix}\NormalTok{(blind\_sched[, }\FunctionTok{c}\NormalTok{(}\StringTok{"height"}\NormalTok{, }\StringTok{"position"}\NormalTok{)])}
\NormalTok{  cat\_b   }\OtherTok{\textless{}{-}}\NormalTok{ blind\_sched}\SpecialCharTok{$}\NormalTok{category\_feedback}
\NormalTok{  res     }\OtherTok{\textless{}{-}} \FunctionTok{track\_rule\_learning}\NormalTok{(}
    \AttributeTok{n\_particles =}\NormalTok{ N\_PARTICLES\_FIXED, }\AttributeTok{error\_prob =} \FloatTok{0.1}\NormalTok{,}
    \AttributeTok{obs =}\NormalTok{ obs\_b, }\AttributeTok{cat\_one =}\NormalTok{ cat\_b, }\AttributeTok{max\_dims =}\NormalTok{ MAX\_RULE\_DIMS, }\AttributeTok{top\_k =} \DecValTok{5}
\NormalTok{  )}
  \CommentTok{\# Take the top 5 distinct rules at the final trial}
\NormalTok{  final }\OtherTok{\textless{}{-}}\NormalTok{ res }\SpecialCharTok{|\textgreater{}} \FunctionTok{filter}\NormalTok{(trial }\SpecialCharTok{==} \FunctionTok{max}\NormalTok{(trial)) }\SpecialCharTok{|\textgreater{}}
    \FunctionTok{distinct}\NormalTok{(dimension1, operation1, threshold1,}
\NormalTok{             dimension2, operation2, threshold2,}
\NormalTok{             logic, prediction, }\AttributeTok{.keep\_all =} \ConstantTok{TRUE}\NormalTok{) }\SpecialCharTok{|\textgreater{}}
    \FunctionTok{head}\NormalTok{(}\DecValTok{5}\NormalTok{)}

\NormalTok{  pred\_strings }\OtherTok{\textless{}{-}} \FunctionTok{character}\NormalTok{(}\FunctionTok{nrow}\NormalTok{(final))}
  \ControlFlowTok{for}\NormalTok{ (k }\ControlFlowTok{in} \FunctionTok{seq\_len}\NormalTok{(}\FunctionTok{nrow}\NormalTok{(final))) \{}
\NormalTok{    r }\OtherTok{\textless{}{-}}\NormalTok{ final[k, ]}
\NormalTok{    feat1 }\OtherTok{\textless{}{-}} \ControlFlowTok{if}\NormalTok{ (r}\SpecialCharTok{$}\NormalTok{dimension1 }\SpecialCharTok{==} \DecValTok{1}\NormalTok{) }\StringTok{"h"} \ControlFlowTok{else} \StringTok{"p"}
\NormalTok{    cond1 }\OtherTok{\textless{}{-}} \FunctionTok{sprintf}\NormalTok{(}\StringTok{"\%s \%s \%.2f"}\NormalTok{, feat1, r}\SpecialCharTok{$}\NormalTok{operation1, r}\SpecialCharTok{$}\NormalTok{threshold1)}
    \ControlFlowTok{if}\NormalTok{ (}\FunctionTok{is.na}\NormalTok{(r}\SpecialCharTok{$}\NormalTok{dimension2)) \{}
\NormalTok{      pred\_strings[k] }\OtherTok{\textless{}{-}} \FunctionTok{sprintf}\NormalTok{(}\StringTok{"(\%s) ? \%d : \%d"}\NormalTok{, cond1, r}\SpecialCharTok{$}\NormalTok{prediction, }\DecValTok{1} \SpecialCharTok{{-}}\NormalTok{ r}\SpecialCharTok{$}\NormalTok{prediction)}
\NormalTok{    \} }\ControlFlowTok{else}\NormalTok{ \{}
\NormalTok{      feat2 }\OtherTok{\textless{}{-}} \ControlFlowTok{if}\NormalTok{ (r}\SpecialCharTok{$}\NormalTok{dimension2 }\SpecialCharTok{==} \DecValTok{1}\NormalTok{) }\StringTok{"h"} \ControlFlowTok{else} \StringTok{"p"}
\NormalTok{      cond2 }\OtherTok{\textless{}{-}} \FunctionTok{sprintf}\NormalTok{(}\StringTok{"\%s \%s \%.2f"}\NormalTok{, feat2, r}\SpecialCharTok{$}\NormalTok{operation2, r}\SpecialCharTok{$}\NormalTok{threshold2)}
\NormalTok{      logic\_op }\OtherTok{\textless{}{-}} \ControlFlowTok{if}\NormalTok{ (r}\SpecialCharTok{$}\NormalTok{logic }\SpecialCharTok{==} \StringTok{"AND"}\NormalTok{) }\StringTok{"\&\&"} \ControlFlowTok{else} \StringTok{"||"}
\NormalTok{      pred\_strings[k] }\OtherTok{\textless{}{-}} \FunctionTok{sprintf}\NormalTok{(}\StringTok{"(\%s \%s \%s) ? \%d : \%d"}\NormalTok{,}
\NormalTok{                                  cond1, logic\_op, cond2,}
\NormalTok{                                  r}\SpecialCharTok{$}\NormalTok{prediction, }\DecValTok{1} \SpecialCharTok{{-}}\NormalTok{ r}\SpecialCharTok{$}\NormalTok{prediction)}
\NormalTok{    \}}
\NormalTok{  \}}
\NormalTok{  pred\_strings}
\NormalTok{\}}

\NormalTok{rules\_B }\OtherTok{\textless{}{-}} \FunctionTok{get\_blind\_rules}\NormalTok{(}\AttributeTok{seed =} \DecValTok{7}\NormalTok{)}

\CommentTok{\# Set C: deliberately bad rules}
\NormalTok{rules\_C }\OtherTok{\textless{}{-}} \FunctionTok{c}\NormalTok{(}
  \StringTok{"h \textgreater{} 4.5 ? 1 : 0"}\NormalTok{,        }\CommentTok{\# virtually never fires}
  \StringTok{"p \textgreater{} 4.5 ? 1 : 0"}\NormalTok{,        }\CommentTok{\# virtually never fires}
  \StringTok{"h \textless{} 0.5 ? 1 : 0"}\NormalTok{,        }\CommentTok{\# virtually never fires}
  \StringTok{"(h \textgreater{} 0.5 \&\& p \textgreater{} 0.5) ? 1 : 0"}\NormalTok{,  }\CommentTok{\# always fires}
  \StringTok{"(h \textgreater{} 0.5 || p \textgreater{} 0.5) ? 1 : 0"}   \CommentTok{\# always fires}
\NormalTok{)}

\CommentTok{\# Compile and fit each}
\NormalTok{mod\_A }\OtherTok{\textless{}{-}} \FunctionTok{build\_stan\_with\_rules}\NormalTok{(rules\_A, }\StringTok{"setA"}\NormalTok{)}
\NormalTok{mod\_B }\OtherTok{\textless{}{-}} \FunctionTok{build\_stan\_with\_rules}\NormalTok{(rules\_B, }\StringTok{"setB"}\NormalTok{)}
\NormalTok{mod\_C }\OtherTok{\textless{}{-}} \FunctionTok{build\_stan\_with\_rules}\NormalTok{(rules\_C, }\StringTok{"setC"}\NormalTok{)}

\NormalTok{fit\_with\_rules }\OtherTok{\textless{}{-}} \ControlFlowTok{function}\NormalTok{(model, label) \{}
\NormalTok{  fit }\OtherTok{\textless{}{-}}\NormalTok{ model}\SpecialCharTok{$}\FunctionTok{sample}\NormalTok{(}
    \AttributeTok{data =}\NormalTok{ rule\_data\_single, }\AttributeTok{init =} \FloatTok{0.5}\NormalTok{, }\AttributeTok{seed =} \DecValTok{7}\NormalTok{,}
    \AttributeTok{chains =} \DecValTok{2}\NormalTok{, }\AttributeTok{parallel\_chains =} \DecValTok{2}\NormalTok{,}
    \AttributeTok{iter\_warmup =} \DecValTok{800}\NormalTok{, }\AttributeTok{iter\_sampling =} \DecValTok{1000}\NormalTok{, }\AttributeTok{refresh =} \DecValTok{0}\NormalTok{, }\AttributeTok{adapt\_delta =} \FloatTok{0.9}
\NormalTok{  )}
\NormalTok{  draws }\OtherTok{\textless{}{-}} \FunctionTok{as\_draws\_df}\NormalTok{(fit}\SpecialCharTok{$}\FunctionTok{draws}\NormalTok{(}\StringTok{"error\_prob"}\NormalTok{))}
  \FunctionTok{tibble}\NormalTok{(}
    \AttributeTok{rule\_set     =}\NormalTok{ label,}
    \AttributeTok{posterior\_mean =} \FunctionTok{mean}\NormalTok{(draws}\SpecialCharTok{$}\NormalTok{error\_prob),}
    \AttributeTok{posterior\_lo =} \FunctionTok{quantile}\NormalTok{(draws}\SpecialCharTok{$}\NormalTok{error\_prob, }\FloatTok{0.025}\NormalTok{),}
    \AttributeTok{posterior\_hi =} \FunctionTok{quantile}\NormalTok{(draws}\SpecialCharTok{$}\NormalTok{error\_prob, }\FloatTok{0.975}\NormalTok{),}
    \AttributeTok{elpd\_loo     =}\NormalTok{ loo}\SpecialCharTok{::}\FunctionTok{loo}\NormalTok{(fit}\SpecialCharTok{$}\FunctionTok{draws}\NormalTok{(}\StringTok{"log\_lik"}\NormalTok{, }\AttributeTok{format =} \StringTok{"matrix"}\NormalTok{))}\SpecialCharTok{$}\NormalTok{estimates[}\StringTok{"elpd\_loo"}\NormalTok{, }\StringTok{"Estimate"}\NormalTok{]}
\NormalTok{  )}
\NormalTok{\}}

\NormalTok{rule\_set\_sensitivity\_file }\OtherTok{\textless{}{-}} \FunctionTok{here}\NormalTok{(}\StringTok{"simdata"}\NormalTok{, }\StringTok{"ch14\_rule\_set\_sensitivity.rds"}\NormalTok{)}
\ControlFlowTok{if}\NormalTok{ (regenerate\_simulations }\SpecialCharTok{||} \SpecialCharTok{!}\FunctionTok{file.exists}\NormalTok{(rule\_set\_sensitivity\_file)) \{}
\NormalTok{  rule\_set\_results }\OtherTok{\textless{}{-}} \FunctionTok{bind\_rows}\NormalTok{(}
    \FunctionTok{fit\_with\_rules}\NormalTok{(mod\_A, }\StringTok{"A: hand{-}picked"}\NormalTok{),}
    \FunctionTok{fit\_with\_rules}\NormalTok{(mod\_B, }\StringTok{"B: blind from Model A"}\NormalTok{),}
    \FunctionTok{fit\_with\_rules}\NormalTok{(mod\_C, }\StringTok{"C: deliberately bad"}\NormalTok{)}
\NormalTok{  )}
  \FunctionTok{saveRDS}\NormalTok{(rule\_set\_results, rule\_set\_sensitivity\_file)}
\NormalTok{\} }\ControlFlowTok{else}\NormalTok{ \{}
\NormalTok{  rule\_set\_results }\OtherTok{\textless{}{-}} \FunctionTok{readRDS}\NormalTok{(rule\_set\_sensitivity\_file)}
\NormalTok{\}}

\ControlFlowTok{if}\NormalTok{ (}\FunctionTok{exists}\NormalTok{(}\StringTok{"rule\_set\_results"}\NormalTok{)) \{}
  \FunctionTok{print}\NormalTok{(rule\_set\_results)}

  \FunctionTok{ggplot}\NormalTok{(rule\_set\_results,}
         \FunctionTok{aes}\NormalTok{(}\AttributeTok{x =}\NormalTok{ rule\_set, }\AttributeTok{y =}\NormalTok{ posterior\_mean)) }\SpecialCharTok{+}
    \FunctionTok{geom\_pointrange}\NormalTok{(}\FunctionTok{aes}\NormalTok{(}\AttributeTok{ymin =}\NormalTok{ posterior\_lo, }\AttributeTok{ymax =}\NormalTok{ posterior\_hi)) }\SpecialCharTok{+}
    \FunctionTok{geom\_hline}\NormalTok{(}\AttributeTok{yintercept =} \FloatTok{0.1}\NormalTok{, }\AttributeTok{linetype =} \StringTok{"dashed"}\NormalTok{, }\AttributeTok{color =} \StringTok{"red"}\NormalTok{) }\SpecialCharTok{+}
    \FunctionTok{labs}\NormalTok{(}
      \AttributeTok{title =} \StringTok{"Sensitivity of Model B\textquotesingle{}s ε posterior to candidate rule set"}\NormalTok{,}
      \AttributeTok{subtitle =} \StringTok{"Red dashed: true ε = 0.1 used to generate the data"}\NormalTok{,}
      \AttributeTok{x =} \ConstantTok{NULL}\NormalTok{, }\AttributeTok{y =} \StringTok{"Posterior mean ± 95\% CrI for ε"}
\NormalTok{    ) }\SpecialCharTok{+}
    \FunctionTok{coord\_flip}\NormalTok{()}
\NormalTok{\}}
\end{Highlighting}
\end{Shaded}

\begin{verbatim}
# A tibble: 3 x 5
  rule_set              posterior_mean posterior_lo posterior_hi elpd_loo
  <chr>                          <dbl>        <dbl>        <dbl>    <dbl>
1 A: hand-picked                0.0777       0.0120        0.204    -32.5
2 B: blind from Model A         0.159        0.0551        0.303    -36.9
3 C: deliberately bad           0.622        0.139         0.944    -44.6
\end{verbatim}

\begin{center}
\includegraphics[width=0.8\linewidth,height=\textheight,keepaspectratio]{13-CategorizationModelsRules_files/figure-pdf/ch14_rule_set_sensitivity-1.pdf}
\end{center}

\textbf{Reading the rule-set sensitivity.} If the three posteriors
agree, the candidate-set choice is innocuous. If they disagree --- which
is the expected outcome --- the gap quantifies how much of Model B's
estimated \(\varepsilon\) is \emph{cognitive} (genuinely tracking the
agent's error rate) vs \emph{structural} (compensating for missing
rules). Set C should produce a wildly inflated \(\varepsilon\) and a
much worse ELPD, providing the upper bound on how much rule-set choice
can move the posterior.

\subsection{\texorpdfstring{Sensitivity to \texttt{max\_dims} (Rule
Grammar
Depth)}{Sensitivity to max\_dims (Rule Grammar Depth)}}\label{sensitivity-to-max_dims-rule-grammar-depth}

\begin{Shaded}
\begin{Highlighting}[]
\NormalTok{max\_dims\_results\_file }\OtherTok{\textless{}{-}} \FunctionTok{here}\NormalTok{(}\StringTok{"simdata"}\NormalTok{, }\StringTok{"ch14\_max\_dims\_sensitivity.rds"}\NormalTok{)}
\ControlFlowTok{if}\NormalTok{ (regenerate\_simulations }\SpecialCharTok{||} \SpecialCharTok{!}\FunctionTok{file.exists}\NormalTok{(max\_dims\_results\_file)) \{}
  \FunctionTok{cat}\NormalTok{(}\StringTok{"Sweeping max\_dims...}\SpecialCharTok{\textbackslash{}n}\StringTok{"}\NormalTok{)}
\NormalTok{  max\_dims\_results }\OtherTok{\textless{}{-}} \FunctionTok{map\_dfr}\NormalTok{(}\FunctionTok{c}\NormalTok{(}\DecValTok{1}\NormalTok{, }\DecValTok{2}\NormalTok{, }\DecValTok{3}\NormalTok{), }\ControlFlowTok{function}\NormalTok{(md) \{}
    \FunctionTok{tryCatch}\NormalTok{(\{}
\NormalTok{      sched }\OtherTok{\textless{}{-}} \FunctionTok{make\_subject\_schedule}\NormalTok{(stimulus\_info, n\_blocks, }\AttributeTok{seed =} \DecValTok{42}\NormalTok{)}
\NormalTok{      sim\_md }\OtherTok{\textless{}{-}} \FunctionTok{rule\_particle\_filter}\NormalTok{(}
        \AttributeTok{n\_particles =}\NormalTok{ N\_PARTICLES\_FIXED,}
        \AttributeTok{error\_prob  =} \FloatTok{0.1}\NormalTok{,}
        \AttributeTok{obs         =} \FunctionTok{as.matrix}\NormalTok{(sched[, }\FunctionTok{c}\NormalTok{(}\StringTok{"height"}\NormalTok{, }\StringTok{"position"}\NormalTok{)]),}
        \AttributeTok{cat\_one     =}\NormalTok{ sched}\SpecialCharTok{$}\NormalTok{category\_feedback,}
        \AttributeTok{max\_dims    =}\NormalTok{ md}
\NormalTok{      )}
\NormalTok{      cum\_acc }\OtherTok{\textless{}{-}} \FunctionTok{mean}\NormalTok{(sim\_md}\SpecialCharTok{$}\NormalTok{sim\_response }\SpecialCharTok{==}\NormalTok{ sched}\SpecialCharTok{$}\NormalTok{category\_feedback)}
      \FunctionTok{tibble}\NormalTok{(}\AttributeTok{max\_dims =}\NormalTok{ md, }\AttributeTok{mean\_accuracy =}\NormalTok{ cum\_acc)}
\NormalTok{    \}, }\AttributeTok{error =} \ControlFlowTok{function}\NormalTok{(e) \{}
      \FunctionTok{tibble}\NormalTok{(}\AttributeTok{max\_dims =}\NormalTok{ md, }\AttributeTok{mean\_accuracy =} \ConstantTok{NA\_real\_}\NormalTok{)}
\NormalTok{    \})}
\NormalTok{  \})}
  \FunctionTok{saveRDS}\NormalTok{(max\_dims\_results, max\_dims\_results\_file)}
\NormalTok{\} }\ControlFlowTok{else}\NormalTok{ \{}
\NormalTok{  max\_dims\_results }\OtherTok{\textless{}{-}} \FunctionTok{readRDS}\NormalTok{(max\_dims\_results\_file)}
\NormalTok{\}}

\ControlFlowTok{if}\NormalTok{ (}\FunctionTok{exists}\NormalTok{(}\StringTok{"max\_dims\_results"}\NormalTok{)) \{}
  \FunctionTok{print}\NormalTok{(max\_dims\_results)}
  \FunctionTok{cat}\NormalTok{(}\StringTok{"Note: max\_dims = 3 will fall back to 2 in our 2{-}feature task.}\SpecialCharTok{\textbackslash{}n}\StringTok{"}\NormalTok{,}
      \StringTok{"The point of this sweep is to make the structural choice visible,}\SpecialCharTok{\textbackslash{}n}\StringTok{"}\NormalTok{,}
      \StringTok{"not to discover anything new.}\SpecialCharTok{\textbackslash{}n}\StringTok{"}\NormalTok{)}
\NormalTok{\}}
\end{Highlighting}
\end{Shaded}

\begin{verbatim}
# A tibble: 3 x 2
  max_dims mean_accuracy
     <dbl>         <dbl>
1        1         0.625
2        2         0.766
3        3         0.766
Note: max_dims = 3 will fall back to 2 in our 2-feature task.
 The point of this sweep is to make the structural choice visible,
 not to discover anything new.
\end{verbatim}

\texttt{max\_dims} is a structural commitment of the rule grammar. In
our 2-feature task it can be 1 or 2 (3 is meaningless). In a
higher-dimensional task it would be a genuine free parameter and should
be either fixed with explicit cognitive justification or inferred via
NPE.

\subsection{Sensitivity to Resampling
Threshold}\label{sensitivity-to-resampling-threshold}

Phase 0 Q4 committed to \texttt{RESAMPLE\_THRESHOLD\ =\ 0.5}. This
controls how aggressively the particle filter reorganises its hypothesis
set: lower values allow more particle degeneracy before resampling;
higher values trigger resampling earlier. The cognitive interpretation
is something like ``how quickly the learner abandons low-performing
hypotheses.''

\begin{Shaded}
\begin{Highlighting}[]
\NormalTok{resample\_sensitivity\_file }\OtherTok{\textless{}{-}} \FunctionTok{here}\NormalTok{(}\StringTok{"simdata"}\NormalTok{, }\StringTok{"ch14\_resample\_threshold\_sensitivity.rds"}\NormalTok{)}
\ControlFlowTok{if}\NormalTok{ (regenerate\_simulations }\SpecialCharTok{||} \SpecialCharTok{!}\FunctionTok{file.exists}\NormalTok{(resample\_sensitivity\_file)) \{}
  \FunctionTok{cat}\NormalTok{(}\StringTok{"Sweeping resampling threshold...}\SpecialCharTok{\textbackslash{}n}\StringTok{"}\NormalTok{)}
\NormalTok{  resample\_results }\OtherTok{\textless{}{-}} \FunctionTok{map\_dfr}\NormalTok{(}\FunctionTok{c}\NormalTok{(}\FloatTok{0.25}\NormalTok{, }\FloatTok{0.50}\NormalTok{, }\FloatTok{0.75}\NormalTok{), }\ControlFlowTok{function}\NormalTok{(rt) \{}
\NormalTok{    accs }\OtherTok{\textless{}{-}} \FunctionTok{vapply}\NormalTok{(}\DecValTok{1}\SpecialCharTok{:}\DecValTok{10}\NormalTok{, }\ControlFlowTok{function}\NormalTok{(rep) \{}
\NormalTok{      sched }\OtherTok{\textless{}{-}} \FunctionTok{make\_subject\_schedule}\NormalTok{(stimulus\_info, n\_blocks, }\AttributeTok{seed =} \DecValTok{42000} \SpecialCharTok{+}\NormalTok{ rep)}
\NormalTok{      sim }\OtherTok{\textless{}{-}} \FunctionTok{rule\_particle\_filter}\NormalTok{(}
        \AttributeTok{n\_particles               =}\NormalTok{ N\_PARTICLES\_FIXED,}
        \AttributeTok{error\_prob                =} \FloatTok{0.1}\NormalTok{,}
        \AttributeTok{obs                       =} \FunctionTok{as.matrix}\NormalTok{(sched[, }\FunctionTok{c}\NormalTok{(}\StringTok{"height"}\NormalTok{, }\StringTok{"position"}\NormalTok{)]),}
        \AttributeTok{cat\_one                   =}\NormalTok{ sched}\SpecialCharTok{$}\NormalTok{category\_feedback,}
        \AttributeTok{max\_dims                  =}\NormalTok{ MAX\_RULE\_DIMS,}
        \AttributeTok{resample\_threshold\_factor =}\NormalTok{ rt,}
        \AttributeTok{quiet                     =} \ConstantTok{TRUE}
\NormalTok{      )}
      \FunctionTok{mean}\NormalTok{(sim}\SpecialCharTok{$}\NormalTok{sim\_response }\SpecialCharTok{==}\NormalTok{ sched}\SpecialCharTok{$}\NormalTok{category\_feedback)}
\NormalTok{    \}, }\FunctionTok{numeric}\NormalTok{(}\DecValTok{1}\NormalTok{))}
    \FunctionTok{tibble}\NormalTok{(}\AttributeTok{resample\_threshold =}\NormalTok{ rt,}
           \AttributeTok{mean\_accuracy =} \FunctionTok{mean}\NormalTok{(accs), }\AttributeTok{sd\_accuracy =} \FunctionTok{sd}\NormalTok{(accs))}
\NormalTok{  \})}
  \FunctionTok{saveRDS}\NormalTok{(resample\_results, resample\_sensitivity\_file)}
\NormalTok{\} }\ControlFlowTok{else}\NormalTok{ \{}
\NormalTok{  resample\_results }\OtherTok{\textless{}{-}} \FunctionTok{readRDS}\NormalTok{(resample\_sensitivity\_file)}
\NormalTok{\}}

\ControlFlowTok{if}\NormalTok{ (}\FunctionTok{exists}\NormalTok{(}\StringTok{"resample\_results"}\NormalTok{)) \{}
  \FunctionTok{print}\NormalTok{(resample\_results)}
  \FunctionTok{cat}\NormalTok{(}\StringTok{"If mean accuracy varies substantially across thresholds,}\SpecialCharTok{\textbackslash{}n}\StringTok{"}\NormalTok{,}
      \StringTok{"the resampling threshold is a consequential structural choice.}\SpecialCharTok{\textbackslash{}n}\StringTok{"}\NormalTok{,}
      \StringTok{"If it is roughly stable, the choice is innocuous for this task.}\SpecialCharTok{\textbackslash{}n}\StringTok{"}\NormalTok{)}
\NormalTok{\}}
\end{Highlighting}
\end{Shaded}

\begin{verbatim}
# A tibble: 3 x 3
  resample_threshold mean_accuracy sd_accuracy
               <dbl>         <dbl>       <dbl>
1               0.25         0.825      0.0960
2               0.5          0.839      0.0621
3               0.75         0.817      0.109 
If mean accuracy varies substantially across thresholds,
 the resampling threshold is a consequential structural choice.
 If it is roughly stable, the choice is innocuous for this task.
\end{verbatim}

\subsection{Option 2: Treat Set A as a Sensitivity Condition, Not a
Canonical
Model}\label{option-2-treat-set-a-as-a-sensitivity-condition-not-a-canonical-model}

Rather than use Set A as \emph{the} Model B and add Set B as a
``robustness check,'' the honest framing is to report all three sets as
variants of Model B with no privileged status. The reader sees three
posteriors instead of one and decides for themselves how much of the
variation is structural. This is the framing we adopt for the chapter 15
model comparison, where Model B is reported as a \emph{family} of
proxies indexed by the candidate set rather than as a single point
estimate.

\begin{center}\rule{0.5\linewidth}{0.5pt}\end{center}

\section{Calibration: Two Different Checks for Two Different
Things}\label{calibration-two-different-checks-for-two-different-things}

Earlier drafts of this chapter ran one block under the header
``Simulation-Based Calibration (SBC)'' that generated data with the
particle filter (Model A) and fitted Model B. \textbf{That is not SBC.}
Talts--Vehtari SBC requires the data to be drawn from the \emph{same}
generative model that is then fitted, so that the joint \((\theta, y)\)
distribution of the inference model is what gets calibrated. When the
generative and inferential models differ, what you are running is closer
to a robustness check than a calibration check, and labelling it SBC is
misleading.

We split the misnamed block into two correctly-labelled checks: (1) a
\textbf{Cross-Model Calibration Check} (data from Model A, fitted by
Model B) and (2) a \textbf{Real SBC for Model B} (data from Model B
itself, fitted by Model B). The first is a robustness check; the second
is the actual Talts--Vehtari calibration test.

\subsection{Cross-Model Calibration Check (NOT
SBC)}\label{cross-model-calibration-check-not-sbc}

\textbf{Procedure.} For each iteration \(i = 1, \ldots, S\): draw
\(\varepsilon_i\) from the prior, simulate one dataset with Model A (the
particle filter) using \(\varepsilon_i\), fit Model B to that dataset,
record the rank of \(\varepsilon_i\) within Model B's posterior. Compute
the rank ECDF.

\textbf{What a flat ECDF means here.} It means: under the specific
misspecification of Model B for data from Model A, \(\varepsilon\)
happens to be recovered with the right \emph{marginal} coverage. It does
\textbf{not} mean Model B is calibrated in the Talts--Vehtari sense,
because the joint \((\theta, y)\) distribution of the inference model is
not the joint distribution of the data-generating process.

\textbf{Why we run it anyway.} Because it answers a useful practical
question: ``when I fit Model B to data that came from a richer
generative process (Model A), are the resulting credible intervals
approximately well-calibrated for the \emph{cognitive} parameter
\(\varepsilon\)?'' The answer to that question is informative even
though it is not what SBC is technically supposed to answer.

\begin{Shaded}
\begin{Highlighting}[]
\NormalTok{n\_cross\_iters }\OtherTok{\textless{}{-}} \DecValTok{500}   \CommentTok{\# use \textgreater{}= 1000 for production}

\NormalTok{cross\_generator }\OtherTok{\textless{}{-}} \FunctionTok{SBC\_generator\_function}\NormalTok{(}
  \ControlFlowTok{function}\NormalTok{() \{}
\NormalTok{    logit\_eps\_true }\OtherTok{\textless{}{-}} \FunctionTok{rnorm}\NormalTok{(}\DecValTok{1}\NormalTok{, LOGIT\_ERROR\_PRIOR\_MEAN, LOGIT\_ERROR\_PRIOR\_SD)}
\NormalTok{    eps\_true       }\OtherTok{\textless{}{-}} \FunctionTok{plogis}\NormalTok{(logit\_eps\_true)}

\NormalTok{    sched }\OtherTok{\textless{}{-}} \FunctionTok{make\_subject\_schedule}\NormalTok{(stimulus\_info, n\_blocks, }\AttributeTok{seed =} \FunctionTok{sample.int}\NormalTok{(}\FloatTok{1e8}\NormalTok{, }\DecValTok{1}\NormalTok{))}
    \CommentTok{\# }\AlertTok{NOTE}\CommentTok{: data drawn from Model A (the particle filter)}
\NormalTok{    sim }\OtherTok{\textless{}{-}} \FunctionTok{rule\_particle\_filter}\NormalTok{(}
      \AttributeTok{n\_particles =}\NormalTok{ N\_PARTICLES\_FIXED,}
      \AttributeTok{error\_prob  =}\NormalTok{ eps\_true,}
      \AttributeTok{obs         =} \FunctionTok{as.matrix}\NormalTok{(sched[, }\FunctionTok{c}\NormalTok{(}\StringTok{"height"}\NormalTok{, }\StringTok{"position"}\NormalTok{)]),}
      \AttributeTok{cat\_one     =}\NormalTok{ sched}\SpecialCharTok{$}\NormalTok{category\_feedback,}
      \AttributeTok{max\_dims    =}\NormalTok{ MAX\_RULE\_DIMS,}
      \AttributeTok{quiet       =} \ConstantTok{TRUE}
\NormalTok{    )}

    \FunctionTok{list}\NormalTok{(}
      \AttributeTok{variables =} \FunctionTok{list}\NormalTok{(}\AttributeTok{logit\_error\_prob =}\NormalTok{ logit\_eps\_true),}
      \AttributeTok{generated =} \FunctionTok{list}\NormalTok{(}
        \AttributeTok{ntrials                =} \FunctionTok{nrow}\NormalTok{(sched),}
        \AttributeTok{nfeatures              =} \DecValTok{2}\NormalTok{L,}
        \AttributeTok{y                      =} \FunctionTok{as.integer}\NormalTok{(sim}\SpecialCharTok{$}\NormalTok{sim\_response),}
        \AttributeTok{obs                    =} \FunctionTok{as.matrix}\NormalTok{(sched[, }\FunctionTok{c}\NormalTok{(}\StringTok{"height"}\NormalTok{, }\StringTok{"position"}\NormalTok{)]),}
        \AttributeTok{prior\_logit\_error\_mean =}\NormalTok{ LOGIT\_ERROR\_PRIOR\_MEAN,}
        \AttributeTok{prior\_logit\_error\_sd   =}\NormalTok{ LOGIT\_ERROR\_PRIOR\_SD}
\NormalTok{      )}
\NormalTok{    )}
\NormalTok{  \}}
\NormalTok{)}

\NormalTok{cross\_backend }\OtherTok{\textless{}{-}} \FunctionTok{SBC\_backend\_cmdstan\_sample}\NormalTok{(}
\NormalTok{  mod\_rule\_stan,}
  \AttributeTok{iter\_warmup =} \DecValTok{1000}\NormalTok{, }\AttributeTok{iter\_sampling =} \DecValTok{1000}\NormalTok{,}
  \AttributeTok{chains =} \DecValTok{2}\NormalTok{, }\AttributeTok{adapt\_delta =} \FloatTok{0.9}\NormalTok{, }\AttributeTok{refresh =} \DecValTok{0}
\NormalTok{)}

\NormalTok{cross\_datasets }\OtherTok{\textless{}{-}} \FunctionTok{generate\_datasets}\NormalTok{(}
\NormalTok{  cross\_generator,}
  \AttributeTok{n\_sims =}\NormalTok{ n\_cross\_iters}
\NormalTok{)}

\NormalTok{cross\_results }\OtherTok{\textless{}{-}} \FunctionTok{compute\_SBC}\NormalTok{(}
\NormalTok{  cross\_datasets, cross\_backend,}
  \AttributeTok{cache\_mode     =} \StringTok{"results"}\NormalTok{,}
  \AttributeTok{cache\_location =} \FunctionTok{here}\NormalTok{(}\StringTok{"simdata"}\NormalTok{, }\StringTok{"ch14\_cross\_model\_check\_cache"}\NormalTok{),}
  \AttributeTok{keep\_fits      =} \ConstantTok{FALSE}
\NormalTok{)}

\CommentTok{\# Plotted as rank ECDF, but with explicit re{-}labelling in the title}
\FunctionTok{plot\_ecdf\_diff}\NormalTok{(cross\_results) }\SpecialCharTok{+}
  \FunctionTok{ggtitle}\NormalTok{(}\StringTok{"CROSS{-}MODEL CALIBRATION CHECK (not SBC)"}\NormalTok{,}
          \AttributeTok{subtitle =} \StringTok{"Data: Model A. Fit: Model B. Flat ≠ calibrated, only marginally well{-}covered."}\NormalTok{)}
\end{Highlighting}
\end{Shaded}

\begin{center}
\includegraphics[width=0.8\linewidth,height=\textheight,keepaspectratio]{13-CategorizationModelsRules_files/figure-pdf/ch14_cross_model_check-1.pdf}
\end{center}

\begin{Shaded}
\begin{Highlighting}[]
\FunctionTok{plot\_rank\_hist}\NormalTok{(cross\_results) }\SpecialCharTok{+}
  \FunctionTok{ggtitle}\NormalTok{(}\StringTok{"Cross{-}model rank histogram (not SBC)"}\NormalTok{)}
\end{Highlighting}
\end{Shaded}

\begin{center}
\includegraphics[width=0.8\linewidth,height=\textheight,keepaspectratio]{13-CategorizationModelsRules_files/figure-pdf/ch14_cross_model_check-2.pdf}
\end{center}

\textbf{Reading this plot.} A flat curve here means Model B's
\(\varepsilon\) posterior covers Model A's true generating
\(\varepsilon\) at approximately the right marginal rate. This is
\emph{useful} --- it tells you Model B is at least not catastrophically
wrong as a tool for estimating cognitive error rates. But it does not
say Model B's inference engine is calibrated in the Talts--Vehtari
sense, because the data-generating distribution is not the inference
model's joint distribution.

\subsection{Real SBC for Model B}\label{real-sbc-for-model-b}

For a proper Talts--Vehtari SBC of Model B, the generator must use Model
B itself: draw \(\varepsilon\) from the prior, draw a rule index from
\(\{1, \ldots, 5\}\) uniformly (mirroring Model B's mixture prior), and
generate Bernoulli responses accordingly. Then fit Model B and check the
rank distribution.

\begin{Shaded}
\begin{Highlighting}[]
\NormalTok{n\_real\_sbc\_iters }\OtherTok{\textless{}{-}} \DecValTok{500}   \CommentTok{\# use \textgreater{}= 1000 for production}

\NormalTok{real\_sbc\_generator }\OtherTok{\textless{}{-}} \FunctionTok{SBC\_generator\_function}\NormalTok{(}
  \ControlFlowTok{function}\NormalTok{() \{}
\NormalTok{    logit\_eps\_true }\OtherTok{\textless{}{-}} \FunctionTok{rnorm}\NormalTok{(}\DecValTok{1}\NormalTok{, LOGIT\_ERROR\_PRIOR\_MEAN, LOGIT\_ERROR\_PRIOR\_SD)}
\NormalTok{    eps\_true       }\OtherTok{\textless{}{-}} \FunctionTok{plogis}\NormalTok{(logit\_eps\_true)}

    \CommentTok{\# Use the SAME generator as Model B\textquotesingle{}s mixture: draw a rule, generate Bernoulli responses}
\NormalTok{    rule\_idx }\OtherTok{\textless{}{-}} \FunctionTok{sample.int}\NormalTok{(}\DecValTok{7}\NormalTok{L, }\DecValTok{1}\NormalTok{)}
\NormalTok{    sched    }\OtherTok{\textless{}{-}} \FunctionTok{make\_subject\_schedule}\NormalTok{(stimulus\_info, n\_blocks, }\AttributeTok{seed =} \FunctionTok{sample.int}\NormalTok{(}\FloatTok{1e8}\NormalTok{, }\DecValTok{1}\NormalTok{))}
\NormalTok{    obs\_mat  }\OtherTok{\textless{}{-}} \FunctionTok{as.matrix}\NormalTok{(sched[, }\FunctionTok{c}\NormalTok{(}\StringTok{"height"}\NormalTok{, }\StringTok{"position"}\NormalTok{)])}
\NormalTok{    n\_trials }\OtherTok{\textless{}{-}} \FunctionTok{nrow}\NormalTok{(obs\_mat)}

\NormalTok{    rule\_pred\_for\_trial }\OtherTok{\textless{}{-}} \ControlFlowTok{function}\NormalTok{(i) \{}
\NormalTok{      h }\OtherTok{\textless{}{-}}\NormalTok{ obs\_mat[i, }\DecValTok{1}\NormalTok{]; p }\OtherTok{\textless{}{-}}\NormalTok{ obs\_mat[i, }\DecValTok{2}\NormalTok{]}
      \ControlFlowTok{switch}\NormalTok{(rule\_idx,}
             \FunctionTok{as.integer}\NormalTok{(h }\SpecialCharTok{\textgreater{}} \FloatTok{2.5}\NormalTok{),}
             \FunctionTok{as.integer}\NormalTok{(p }\SpecialCharTok{\textless{}} \FloatTok{2.5}\NormalTok{),}
             \FunctionTok{as.integer}\NormalTok{(h }\SpecialCharTok{\textgreater{}} \FloatTok{2.0} \SpecialCharTok{\&\&}\NormalTok{ p }\SpecialCharTok{\textgreater{}} \FloatTok{3.0}\NormalTok{),}
             \FunctionTok{as.integer}\NormalTok{(}\SpecialCharTok{!}\NormalTok{(h }\SpecialCharTok{\textless{}} \FloatTok{2.0} \SpecialCharTok{||}\NormalTok{ p }\SpecialCharTok{\textless{}} \FloatTok{2.0}\NormalTok{)),}
             \FunctionTok{as.integer}\NormalTok{(h }\SpecialCharTok{\textgreater{}} \FloatTok{2.5} \SpecialCharTok{\&\&}\NormalTok{ p }\SpecialCharTok{\textless{}} \FloatTok{3.5}\NormalTok{),}
             \FunctionTok{as.integer}\NormalTok{(h }\SpecialCharTok{\textgreater{}} \FloatTok{3.5} \SpecialCharTok{||}\NormalTok{ (h }\SpecialCharTok{\textgreater{}} \FloatTok{1.5} \SpecialCharTok{\&\&}\NormalTok{ p }\SpecialCharTok{\textless{}} \FloatTok{1.5}\NormalTok{)),}
             \FunctionTok{as.integer}\NormalTok{(h }\SpecialCharTok{\textgreater{}} \FloatTok{1.5} \SpecialCharTok{\&\&}\NormalTok{ (p }\SpecialCharTok{\textless{}} \FloatTok{2.5} \SpecialCharTok{||}\NormalTok{ h }\SpecialCharTok{\textgreater{}} \FloatTok{3.5}\NormalTok{)))}
\NormalTok{    \}}
\NormalTok{    preds  }\OtherTok{\textless{}{-}} \FunctionTok{vapply}\NormalTok{(}\FunctionTok{seq\_len}\NormalTok{(n\_trials), rule\_pred\_for\_trial, }\FunctionTok{integer}\NormalTok{(}\DecValTok{1}\NormalTok{))}
\NormalTok{    p\_cat1 }\OtherTok{\textless{}{-}} \FunctionTok{ifelse}\NormalTok{(preds }\SpecialCharTok{==} \DecValTok{1}\NormalTok{, }\DecValTok{1} \SpecialCharTok{{-}}\NormalTok{ eps\_true, eps\_true)}
\NormalTok{    y      }\OtherTok{\textless{}{-}} \FunctionTok{rbinom}\NormalTok{(n\_trials, }\DecValTok{1}\NormalTok{, p\_cat1)}

    \FunctionTok{list}\NormalTok{(}
      \AttributeTok{variables =} \FunctionTok{list}\NormalTok{(}\AttributeTok{logit\_error\_prob =}\NormalTok{ logit\_eps\_true),}
      \AttributeTok{generated =} \FunctionTok{list}\NormalTok{(}
        \AttributeTok{ntrials                =}\NormalTok{ n\_trials,}
        \AttributeTok{nfeatures              =} \DecValTok{2}\NormalTok{L,}
        \AttributeTok{y                      =} \FunctionTok{as.integer}\NormalTok{(y),}
        \AttributeTok{obs                    =}\NormalTok{ obs\_mat,}
        \AttributeTok{prior\_logit\_error\_mean =}\NormalTok{ LOGIT\_ERROR\_PRIOR\_MEAN,}
        \AttributeTok{prior\_logit\_error\_sd   =}\NormalTok{ LOGIT\_ERROR\_PRIOR\_SD}
\NormalTok{      )}
\NormalTok{    )}
\NormalTok{  \}}
\NormalTok{)}

\NormalTok{real\_sbc\_backend }\OtherTok{\textless{}{-}} \FunctionTok{SBC\_backend\_cmdstan\_sample}\NormalTok{(}
\NormalTok{  mod\_rule\_stan,}
  \AttributeTok{iter\_warmup =} \DecValTok{1000}\NormalTok{, }\AttributeTok{iter\_sampling =} \DecValTok{1000}\NormalTok{,}
  \AttributeTok{chains =} \DecValTok{2}\NormalTok{, }\AttributeTok{adapt\_delta =} \FloatTok{0.9}\NormalTok{, }\AttributeTok{refresh =} \DecValTok{0}
\NormalTok{)}

\NormalTok{real\_sbc\_datasets }\OtherTok{\textless{}{-}} \FunctionTok{generate\_datasets}\NormalTok{(}
\NormalTok{  real\_sbc\_generator,}
  \AttributeTok{n\_sims =}\NormalTok{ n\_real\_sbc\_iters}
\NormalTok{)}

\NormalTok{real\_sbc\_results }\OtherTok{\textless{}{-}} \FunctionTok{compute\_SBC}\NormalTok{(}
\NormalTok{  real\_sbc\_datasets, real\_sbc\_backend,}
  \AttributeTok{cache\_mode     =} \StringTok{"results"}\NormalTok{,}
  \AttributeTok{cache\_location =} \FunctionTok{here}\NormalTok{(}\StringTok{"simdata"}\NormalTok{, }\StringTok{"ch14\_real\_sbc\_cache"}\NormalTok{),}
  \AttributeTok{keep\_fits      =} \ConstantTok{FALSE}
\NormalTok{)}

\FunctionTok{plot\_ecdf\_diff}\NormalTok{(real\_sbc\_results) }\SpecialCharTok{+}
  \FunctionTok{ggtitle}\NormalTok{(}\StringTok{"REAL SBC for Model B (Talts–Vehtari)"}\NormalTok{,}
          \AttributeTok{subtitle =} \StringTok{"Data: Model B\textquotesingle{}s own mixture generator. Fit: Model B."}\NormalTok{)}
\end{Highlighting}
\end{Shaded}

\begin{center}
\includegraphics[width=0.8\linewidth,height=\textheight,keepaspectratio]{13-CategorizationModelsRules_files/figure-pdf/ch14_real_sbc_model_B-1.pdf}
\end{center}

\begin{Shaded}
\begin{Highlighting}[]
\FunctionTok{plot\_rank\_hist}\NormalTok{(real\_sbc\_results) }\SpecialCharTok{+}
  \FunctionTok{ggtitle}\NormalTok{(}\StringTok{"Real SBC rank histogram for Model B"}\NormalTok{)}
\end{Highlighting}
\end{Shaded}

\begin{center}
\includegraphics[width=0.8\linewidth,height=\textheight,keepaspectratio]{13-CategorizationModelsRules_files/figure-pdf/ch14_real_sbc_model_B-2.pdf}
\end{center}

\textbf{Reading this plot.} This is the actual calibration check. A flat
ECDF here means Stan correctly inverts Model B (Model B's inference
engine is calibrated in the Talts--Vehtari sense). Systematic deviations
would indicate a miscalibration in the Stan implementation itself --- a
sign error, an off-by-one in the marginalisation, or a
prior-on-\texttt{logit\_error\_prob} mismatch. \textbf{This is the only
diagnostic in the chapter that tests the Stan engine on its own terms.}
The cross-model check above tests something else --- it tests how Model
B behaves when applied to data from a richer process --- and that
something else is informative but it is not SBC.

\subsection{Posterior SBC for Model B (Phase
6c)}\label{posterior-sbc-for-model-b-phase-6c}

The persona's Phase 6c calls for \textbf{Posterior SBC (PSBC)}: instead
of drawing true parameters from the prior, draw them from the
\emph{fitted posterior}. This tests calibration in the specific
high-probability region the actual data force us into. A sampler can be
globally unbiased (prior SBC passes) yet struggle locally (PSBC fails)
or vice versa.

\begin{Shaded}
\begin{Highlighting}[]
\NormalTok{n\_psbc\_iters }\OtherTok{\textless{}{-}} \DecValTok{500}  \CommentTok{\# use \textgreater{}= 1000 for production}

\ControlFlowTok{if}\NormalTok{ (}\SpecialCharTok{!}\FunctionTok{is.null}\NormalTok{(fit\_rule\_stan)) \{}
  \CommentTok{\# Draw true parameters from the fitted posterior, not the prior}
\NormalTok{  fitted\_draws }\OtherTok{\textless{}{-}} \FunctionTok{as\_draws\_df}\NormalTok{(fit\_rule\_stan}\SpecialCharTok{$}\FunctionTok{draws}\NormalTok{(}\StringTok{"logit\_error\_prob"}\NormalTok{))}
\NormalTok{  fitted\_logit\_eps }\OtherTok{\textless{}{-}}\NormalTok{ fitted\_draws}\SpecialCharTok{$}\NormalTok{logit\_error\_prob}

\NormalTok{  psbc\_generator }\OtherTok{\textless{}{-}} \FunctionTok{SBC\_generator\_function}\NormalTok{(}
    \ControlFlowTok{function}\NormalTok{() \{}
      \CommentTok{\# Draw from the POSTERIOR, not the prior}
\NormalTok{      logit\_eps\_true }\OtherTok{\textless{}{-}} \FunctionTok{sample}\NormalTok{(fitted\_logit\_eps, }\DecValTok{1}\NormalTok{)}
\NormalTok{      eps\_true       }\OtherTok{\textless{}{-}} \FunctionTok{plogis}\NormalTok{(logit\_eps\_true)}

      \CommentTok{\# Generate from Model B\textquotesingle{}s own mixture (same as real SBC; 7 rules)}
\NormalTok{      rule\_idx }\OtherTok{\textless{}{-}} \FunctionTok{sample.int}\NormalTok{(}\DecValTok{7}\NormalTok{L, }\DecValTok{1}\NormalTok{)}
\NormalTok{      sched    }\OtherTok{\textless{}{-}} \FunctionTok{make\_subject\_schedule}\NormalTok{(stimulus\_info, n\_blocks, }\AttributeTok{seed =} \FunctionTok{sample.int}\NormalTok{(}\FloatTok{1e8}\NormalTok{, }\DecValTok{1}\NormalTok{))}
\NormalTok{      obs\_mat  }\OtherTok{\textless{}{-}} \FunctionTok{as.matrix}\NormalTok{(sched[, }\FunctionTok{c}\NormalTok{(}\StringTok{"height"}\NormalTok{, }\StringTok{"position"}\NormalTok{)])}
\NormalTok{      n\_trials }\OtherTok{\textless{}{-}} \FunctionTok{nrow}\NormalTok{(obs\_mat)}

\NormalTok{      rule\_pred\_for\_trial }\OtherTok{\textless{}{-}} \ControlFlowTok{function}\NormalTok{(j) \{}
\NormalTok{        h }\OtherTok{\textless{}{-}}\NormalTok{ obs\_mat[j, }\DecValTok{1}\NormalTok{]; p }\OtherTok{\textless{}{-}}\NormalTok{ obs\_mat[j, }\DecValTok{2}\NormalTok{]}
        \ControlFlowTok{switch}\NormalTok{(rule\_idx,}
               \FunctionTok{as.integer}\NormalTok{(h }\SpecialCharTok{\textgreater{}} \FloatTok{2.5}\NormalTok{),}
               \FunctionTok{as.integer}\NormalTok{(p }\SpecialCharTok{\textless{}} \FloatTok{2.5}\NormalTok{),}
               \FunctionTok{as.integer}\NormalTok{(h }\SpecialCharTok{\textgreater{}} \FloatTok{2.0} \SpecialCharTok{\&\&}\NormalTok{ p }\SpecialCharTok{\textgreater{}} \FloatTok{3.0}\NormalTok{),}
               \FunctionTok{as.integer}\NormalTok{(}\SpecialCharTok{!}\NormalTok{(h }\SpecialCharTok{\textless{}} \FloatTok{2.0} \SpecialCharTok{||}\NormalTok{ p }\SpecialCharTok{\textless{}} \FloatTok{2.0}\NormalTok{)),}
               \FunctionTok{as.integer}\NormalTok{(h }\SpecialCharTok{\textgreater{}} \FloatTok{2.5} \SpecialCharTok{\&\&}\NormalTok{ p }\SpecialCharTok{\textless{}} \FloatTok{3.5}\NormalTok{),}
               \FunctionTok{as.integer}\NormalTok{(h }\SpecialCharTok{\textgreater{}} \FloatTok{3.5} \SpecialCharTok{||}\NormalTok{ (h }\SpecialCharTok{\textgreater{}} \FloatTok{1.5} \SpecialCharTok{\&\&}\NormalTok{ p }\SpecialCharTok{\textless{}} \FloatTok{1.5}\NormalTok{)),}
               \FunctionTok{as.integer}\NormalTok{(h }\SpecialCharTok{\textgreater{}} \FloatTok{1.5} \SpecialCharTok{\&\&}\NormalTok{ (p }\SpecialCharTok{\textless{}} \FloatTok{2.5} \SpecialCharTok{||}\NormalTok{ h }\SpecialCharTok{\textgreater{}} \FloatTok{3.5}\NormalTok{)))}
\NormalTok{      \}}
\NormalTok{      preds  }\OtherTok{\textless{}{-}} \FunctionTok{vapply}\NormalTok{(}\FunctionTok{seq\_len}\NormalTok{(n\_trials), rule\_pred\_for\_trial, }\FunctionTok{integer}\NormalTok{(}\DecValTok{1}\NormalTok{))}
\NormalTok{      p\_cat1 }\OtherTok{\textless{}{-}} \FunctionTok{ifelse}\NormalTok{(preds }\SpecialCharTok{==} \DecValTok{1}\NormalTok{, }\DecValTok{1} \SpecialCharTok{{-}}\NormalTok{ eps\_true, eps\_true)}
\NormalTok{      y      }\OtherTok{\textless{}{-}} \FunctionTok{rbinom}\NormalTok{(n\_trials, }\DecValTok{1}\NormalTok{, p\_cat1)}

      \FunctionTok{list}\NormalTok{(}
        \AttributeTok{variables =} \FunctionTok{list}\NormalTok{(}\AttributeTok{logit\_error\_prob =}\NormalTok{ logit\_eps\_true),}
        \AttributeTok{generated =} \FunctionTok{list}\NormalTok{(}
          \AttributeTok{ntrials                =}\NormalTok{ n\_trials,}
          \AttributeTok{nfeatures              =} \DecValTok{2}\NormalTok{L,}
          \AttributeTok{y                      =} \FunctionTok{as.integer}\NormalTok{(y),}
          \AttributeTok{obs                    =}\NormalTok{ obs\_mat,}
          \AttributeTok{prior\_logit\_error\_mean =}\NormalTok{ LOGIT\_ERROR\_PRIOR\_MEAN,}
          \AttributeTok{prior\_logit\_error\_sd   =}\NormalTok{ LOGIT\_ERROR\_PRIOR\_SD}
\NormalTok{        )}
\NormalTok{      )}
\NormalTok{    \}}
\NormalTok{  )}

\NormalTok{  psbc\_backend }\OtherTok{\textless{}{-}} \FunctionTok{SBC\_backend\_cmdstan\_sample}\NormalTok{(}
\NormalTok{    mod\_rule\_stan,}
    \AttributeTok{iter\_warmup =} \DecValTok{1000}\NormalTok{, }\AttributeTok{iter\_sampling =} \DecValTok{1000}\NormalTok{,}
    \AttributeTok{chains =} \DecValTok{2}\NormalTok{, }\AttributeTok{adapt\_delta =} \FloatTok{0.9}\NormalTok{, }\AttributeTok{refresh =} \DecValTok{0}
\NormalTok{  )}

\NormalTok{  psbc\_datasets }\OtherTok{\textless{}{-}} \FunctionTok{generate\_datasets}\NormalTok{(}
\NormalTok{    psbc\_generator,}
    \AttributeTok{n\_sims =}\NormalTok{ n\_psbc\_iters}
\NormalTok{  )}

\NormalTok{  psbc\_results }\OtherTok{\textless{}{-}} \FunctionTok{compute\_SBC}\NormalTok{(}
\NormalTok{    psbc\_datasets, psbc\_backend,}
    \AttributeTok{cache\_mode     =} \StringTok{"results"}\NormalTok{,}
    \AttributeTok{cache\_location =} \FunctionTok{here}\NormalTok{(}\StringTok{"simdata"}\NormalTok{, }\StringTok{"ch14\_psbc\_cache"}\NormalTok{),}
    \AttributeTok{keep\_fits      =} \ConstantTok{FALSE}
\NormalTok{  )}

  \FunctionTok{plot\_ecdf\_diff}\NormalTok{(psbc\_results) }\SpecialCharTok{+}
    \FunctionTok{ggtitle}\NormalTok{(}\StringTok{"POSTERIOR SBC for Model B (Phase 6c)"}\NormalTok{,}
            \AttributeTok{subtitle =} \StringTok{"True ε drawn from FITTED posterior, not the prior.}\SpecialCharTok{\textbackslash{}n}\StringTok{Tests local calibration in the high{-}probability region."}\NormalTok{)}
  \FunctionTok{plot\_rank\_hist}\NormalTok{(psbc\_results) }\SpecialCharTok{+}
    \FunctionTok{ggtitle}\NormalTok{(}\StringTok{"Posterior SBC rank histogram for Model B"}\NormalTok{)}
\NormalTok{\}}
\end{Highlighting}
\end{Shaded}

\begin{center}
\includegraphics[width=0.8\linewidth,height=\textheight,keepaspectratio]{13-CategorizationModelsRules_files/figure-pdf/ch14_psbc_model_B-1.pdf}
\end{center}

\textbf{Reading this plot.} A flat ECDF here means Model B is calibrated
\emph{locally} around the posterior region that the actual data
identified. If prior SBC passes but PSBC fails, the sampler is globally
unbiased but locally struggles --- a sign of geometry issues in the
high-posterior region. If both pass, the Stan implementation is sound in
both senses.

\begin{center}\rule{0.5\linewidth}{0.5pt}\end{center}

\section{Neural Posterior Estimation for Model
A}\label{neural-posterior-estimation-for-model-a}

This is the section that does what earlier drafts of this material
avoided with a weak tool: it fits the actual cognitive theory rather
than a tractable proxy. Neural Posterior Estimation (NPE) --- a modern
simulation-based inference method --- trains a conditional normalising
flow to approximate the posterior \(p(\theta \mid y)\) from simulated
\((\theta, y)\) pairs, without requiring a tractable likelihood or
hand-crafted summary statistics. The Deistler et al.~(2025) tutorial and
the BayesFlow paper (Radev et al.~2023) provide the methodological
background; we give a working implementation here.

\textbf{Why NPE instead of rejection ABC?} Earlier drafts used rejection
ABC with three hand-crafted summary statistics. That approach had three
problems:

\begin{enumerate}
\def\labelenumi{\arabic{enumi}.}
\item
  \textbf{Schedule mismatch bug.} Each ABC simulation drew a \emph{new}
  random schedule, injecting trial-order noise on top of parameter
  signal. The particle filter is path-dependent, so the learning slope
  and per-stimulus accuracy are functions of \emph{which stimuli
  appeared when}. The correct conditional simulator holds the schedule
  fixed at the observed schedule and re-randomises only the particle
  initialisation and Bernoulli draws.
\item
  \textbf{Summary statistic insufficiency.} The three summaries (block-1
  vs block-final accuracy, per-stimulus accuracy vector, lag-1
  autocorrelation) were reasonable guesses, but guesses. NPE with a
  learned embedding network discovers the relevant features of the
  response sequence end-to-end from the training loss, eliminating the
  need to defend ``why these three and not others?''
\item
  \textbf{No amortization.} Rejection ABC pays the full simulation
  budget per subject. NPE pays it once during training. After training,
  inference on any new subject takes milliseconds. This is what makes
  the Chapter 15 comparison across all three cognitive theories
  executable.
\end{enumerate}

\subsection{Setting Up the Python
Environment}\label{setting-up-the-python-environment}

We use the \texttt{sbi} package (Tejero-Cantero et al.~2020; Boelts et
al.~2024) via \texttt{reticulate}. The key components are: a prior, a
simulator wrapper, and an inference network (a normalising flow with a
learned embedding network).

\begin{Shaded}
\begin{Highlighting}[]
\CommentTok{\# ── Python environment ────────────────────────────────────────────────────────}
\CommentTok{\# Requires: pip install sbi torch numpy}
\ControlFlowTok{if}\NormalTok{ (run\_intensive\_checks) \{}
  \FunctionTok{library}\NormalTok{(reticulate)}
\NormalTok{  reticulate}\SpecialCharTok{::}\FunctionTok{use\_virtualenv}\NormalTok{(}\StringTok{"r{-}reticulate"}\NormalTok{, }\AttributeTok{required =} \ConstantTok{TRUE}\NormalTok{)}

  \ControlFlowTok{if}\NormalTok{ (}\SpecialCharTok{!}\FunctionTok{py\_available}\NormalTok{(}\AttributeTok{initialize =} \ConstantTok{TRUE}\NormalTok{)) \{}
    \FunctionTok{stop}\NormalTok{(}\StringTok{"Python not available. Install Python and: pip install sbi torch numpy"}\NormalTok{)}
\NormalTok{  \}}

\NormalTok{  sbi   }\OtherTok{\textless{}{-}} \FunctionTok{import}\NormalTok{(}\StringTok{"sbi"}\NormalTok{)}
\NormalTok{  torch }\OtherTok{\textless{}{-}} \FunctionTok{import}\NormalTok{(}\StringTok{"torch"}\NormalTok{)}
\NormalTok{  np    }\OtherTok{\textless{}{-}} \FunctionTok{import}\NormalTok{(}\StringTok{"numpy"}\NormalTok{)}
\NormalTok{  utils }\OtherTok{\textless{}{-}} \FunctionTok{import}\NormalTok{(}\StringTok{"sbi.utils"}\NormalTok{)}
\NormalTok{  infer }\OtherTok{\textless{}{-}} \FunctionTok{import}\NormalTok{(}\StringTok{"sbi.inference"}\NormalTok{)}

  \FunctionTok{cat}\NormalTok{(}\StringTok{"sbi version:"}\NormalTok{, sbi}\SpecialCharTok{$}\StringTok{\textasciigrave{}}\AttributeTok{\_\_version\_\_}\StringTok{\textasciigrave{}}\NormalTok{, }\StringTok{"}\SpecialCharTok{\textbackslash{}n}\StringTok{"}\NormalTok{)}
  \FunctionTok{cat}\NormalTok{(}\StringTok{"torch version:"}\NormalTok{, torch}\SpecialCharTok{$}\StringTok{\textasciigrave{}}\AttributeTok{\_\_version\_\_}\StringTok{\textasciigrave{}}\NormalTok{, }\StringTok{"}\SpecialCharTok{\textbackslash{}n}\StringTok{"}\NormalTok{)}
\NormalTok{\} }\ControlFlowTok{else}\NormalTok{ \{}
  \FunctionTok{cat}\NormalTok{(}\StringTok{"Skipping NPE setup (run\_intensive\_checks is FALSE).}\SpecialCharTok{\textbackslash{}n}\StringTok{"}\NormalTok{)}
\NormalTok{\}}
\end{Highlighting}
\end{Shaded}

\begin{verbatim}
Skipping NPE setup (run_intensive_checks is FALSE).
\end{verbatim}

\subsection{The Simulator Wrapper}\label{the-simulator-wrapper}

The simulator must take a parameter vector \(\theta\) and return the
observed data \(y\). \textbf{Critically, the schedule is held fixed} ---
we condition on the observed trial order and only re-randomise the
particle filter's internal stochasticity (particle initialisation and
Bernoulli response noise). This fixes the schedule-mismatch bug from
earlier drafts.

We infer two parameters jointly:

\begin{itemize}
\tightlist
\item
  \(\varepsilon\) --- error probability, on the logit scale for
  unconstrained inference
\item
  \(\log(N_{\text{particles}})\) --- working-memory bound, on the log
  scale
\end{itemize}

The data representation for NPE is the full per-trial response vector
\(y \in \{0, 1\}^T\), concatenated with the per-trial stimulus features
and feedback. The embedding network learns what to attend to.

\begin{Shaded}
\begin{Highlighting}[]
\CommentTok{\# The observed schedule — used for ALL simulations (schedule{-}fix)}
\NormalTok{npe\_schedule }\OtherTok{\textless{}{-}}\NormalTok{ fit\_schedule}
\NormalTok{npe\_obs\_mat  }\OtherTok{\textless{}{-}} \FunctionTok{as.matrix}\NormalTok{(npe\_schedule[, }\FunctionTok{c}\NormalTok{(}\StringTok{"height"}\NormalTok{, }\StringTok{"position"}\NormalTok{)])}
\NormalTok{npe\_cat\_one  }\OtherTok{\textless{}{-}}\NormalTok{ npe\_schedule}\SpecialCharTok{$}\NormalTok{category\_feedback}
\NormalTok{n\_trials\_npe }\OtherTok{\textless{}{-}} \FunctionTok{nrow}\NormalTok{(npe\_schedule)}

\CommentTok{\# Build the context: [height, position, feedback] per trial}
\CommentTok{\# This is constant across simulations because the schedule is fixed}
\NormalTok{context\_per\_trial }\OtherTok{\textless{}{-}} \FunctionTok{cbind}\NormalTok{(npe\_obs\_mat, npe\_cat\_one)  }\CommentTok{\# (n\_trials, 3)}

\CommentTok{\# R{-}side simulator: theta → response vector}
\CommentTok{\# theta[1] = logit(eps), theta[2] = log(N\_particles)}
\NormalTok{simulate\_model\_a }\OtherTok{\textless{}{-}} \ControlFlowTok{function}\NormalTok{(theta\_vec) \{}
\NormalTok{  eps     }\OtherTok{\textless{}{-}} \FunctionTok{plogis}\NormalTok{(theta\_vec[}\DecValTok{1}\NormalTok{])}
\NormalTok{  n\_part  }\OtherTok{\textless{}{-}} \FunctionTok{max}\NormalTok{(}\DecValTok{10}\NormalTok{L, }\FunctionTok{as.integer}\NormalTok{(}\FunctionTok{round}\NormalTok{(}\FunctionTok{exp}\NormalTok{(theta\_vec[}\DecValTok{2}\NormalTok{]))))}

\NormalTok{  sim }\OtherTok{\textless{}{-}} \FunctionTok{rule\_particle\_filter}\NormalTok{(}
    \AttributeTok{n\_particles               =}\NormalTok{ n\_part,}
    \AttributeTok{error\_prob                =}\NormalTok{ eps,}
    \AttributeTok{obs                       =}\NormalTok{ npe\_obs\_mat,}
    \AttributeTok{cat\_one                   =}\NormalTok{ npe\_cat\_one,}
    \AttributeTok{max\_dims                  =}\NormalTok{ MAX\_RULE\_DIMS,}
    \AttributeTok{resample\_threshold\_factor =}\NormalTok{ RESAMPLE\_THRESHOLD,}
    \AttributeTok{quiet                     =} \ConstantTok{TRUE}
\NormalTok{  )}
  \FunctionTok{as.numeric}\NormalTok{(sim}\SpecialCharTok{$}\NormalTok{sim\_response)}
\NormalTok{\}}

\CommentTok{\# Test it}
\NormalTok{test\_theta }\OtherTok{\textless{}{-}} \FunctionTok{c}\NormalTok{(}\FunctionTok{qlogis}\NormalTok{(}\FloatTok{0.1}\NormalTok{), }\FunctionTok{log}\NormalTok{(}\DecValTok{100}\NormalTok{))}
\NormalTok{test\_y     }\OtherTok{\textless{}{-}} \FunctionTok{simulate\_model\_a}\NormalTok{(test\_theta)}
\FunctionTok{cat}\NormalTok{(}\StringTok{"Test simulation: eps ="}\NormalTok{, }\FunctionTok{round}\NormalTok{(}\FunctionTok{plogis}\NormalTok{(test\_theta[}\DecValTok{1}\NormalTok{]), }\DecValTok{3}\NormalTok{),}
    \StringTok{", N\_particles ="}\NormalTok{, }\FunctionTok{round}\NormalTok{(}\FunctionTok{exp}\NormalTok{(test\_theta[}\DecValTok{2}\NormalTok{])),}
    \StringTok{", mean response ="}\NormalTok{, }\FunctionTok{round}\NormalTok{(}\FunctionTok{mean}\NormalTok{(test\_y), }\DecValTok{3}\NormalTok{), }\StringTok{"}\SpecialCharTok{\textbackslash{}n}\StringTok{"}\NormalTok{)}
\end{Highlighting}
\end{Shaded}

\subsection{Generating the Training
Dataset}\label{generating-the-training-dataset}

NPE training requires a dataset of \((\theta, x)\) pairs where \(x\)
includes both the response vector and the stimulus context. We generate
this dataset in R (leveraging \texttt{furrr} parallelism) and pass it to
Python as tensors.

\begin{Shaded}
\begin{Highlighting}[]
\NormalTok{n\_train\_sims   }\OtherTok{\textless{}{-}} \DecValTok{20000}   \CommentTok{\# 20k sims; increase to 50k for production}
\NormalTok{npe\_cache\_file }\OtherTok{\textless{}{-}} \FunctionTok{here}\NormalTok{(}\StringTok{"simdata"}\NormalTok{, }\StringTok{"ch14\_npe\_training\_data.rds"}\NormalTok{)}

\ControlFlowTok{if}\NormalTok{ (regenerate\_simulations }\SpecialCharTok{||} \SpecialCharTok{!}\FunctionTok{file.exists}\NormalTok{(npe\_cache\_file)) \{}
  \FunctionTok{cat}\NormalTok{(}\StringTok{"Generating NPE training dataset ("}\NormalTok{, n\_train\_sims, }\StringTok{" simulations)...}\SpecialCharTok{\textbackslash{}n}\StringTok{"}\NormalTok{)}
  \FunctionTok{set.seed}\NormalTok{(}\DecValTok{2026}\NormalTok{)}

  \CommentTok{\# Draw parameters from the prior}
\NormalTok{  train\_logit\_eps }\OtherTok{\textless{}{-}} \FunctionTok{rnorm}\NormalTok{(n\_train\_sims, LOGIT\_ERROR\_PRIOR\_MEAN, LOGIT\_ERROR\_PRIOR\_SD)}
\NormalTok{  train\_log\_npart }\OtherTok{\textless{}{-}} \FunctionTok{runif}\NormalTok{(n\_train\_sims, }\FunctionTok{log}\NormalTok{(}\DecValTok{20}\NormalTok{), }\FunctionTok{log}\NormalTok{(}\DecValTok{300}\NormalTok{))  }\CommentTok{\# log{-}uniform over [20, 300]}
\NormalTok{  train\_thetas    }\OtherTok{\textless{}{-}} \FunctionTok{cbind}\NormalTok{(train\_logit\_eps, train\_log\_npart)}

  \CommentTok{\# Simulate in parallel — schedule is FIXED (the bug fix)}
\NormalTok{  train\_responses }\OtherTok{\textless{}{-}} \FunctionTok{future\_map}\NormalTok{(}
    \FunctionTok{seq\_len}\NormalTok{(n\_train\_sims),}
    \ControlFlowTok{function}\NormalTok{(i) }\FunctionTok{simulate\_model\_a}\NormalTok{(train\_thetas[i, ]),}
    \AttributeTok{.options =} \FunctionTok{furrr\_options}\NormalTok{(}\AttributeTok{seed =} \ConstantTok{TRUE}\NormalTok{),}
    \AttributeTok{.progress =} \ConstantTok{TRUE}
\NormalTok{  )}

\NormalTok{  train\_y\_mat }\OtherTok{\textless{}{-}} \FunctionTok{do.call}\NormalTok{(rbind, train\_responses)}

  \FunctionTok{saveRDS}\NormalTok{(}\FunctionTok{list}\NormalTok{(}\AttributeTok{thetas =}\NormalTok{ train\_thetas, }\AttributeTok{responses =}\NormalTok{ train\_y\_mat),}
\NormalTok{          npe\_cache\_file)}
  \FunctionTok{cat}\NormalTok{(}\StringTok{"Training data saved:"}\NormalTok{, n\_train\_sims, }\StringTok{"simulations.}\SpecialCharTok{\textbackslash{}n}\StringTok{"}\NormalTok{)}
\NormalTok{\} }\ControlFlowTok{else}\NormalTok{ \{}
\NormalTok{  cached\_train }\OtherTok{\textless{}{-}} \FunctionTok{readRDS}\NormalTok{(npe\_cache\_file)}
\NormalTok{  train\_thetas }\OtherTok{\textless{}{-}}\NormalTok{ cached\_train}\SpecialCharTok{$}\NormalTok{thetas}
\NormalTok{  train\_y\_mat  }\OtherTok{\textless{}{-}}\NormalTok{ cached\_train}\SpecialCharTok{$}\NormalTok{responses}
\NormalTok{  n\_train\_sims }\OtherTok{\textless{}{-}} \FunctionTok{nrow}\NormalTok{(train\_thetas)}
  \FunctionTok{cat}\NormalTok{(}\StringTok{"Loaded"}\NormalTok{, n\_train\_sims, }\StringTok{"cached training simulations.}\SpecialCharTok{\textbackslash{}n}\StringTok{"}\NormalTok{)}
\NormalTok{\}}
\end{Highlighting}
\end{Shaded}

\subsection{Training the NPE Inference
Network}\label{training-the-npe-inference-network}

We use SNPE with a neural spline flow. The embedding network is a small
MLP that processes the concatenation of
\texttt{{[}response,\ height,\ position,\ feedback{]}} per trial into a
learned summary vector.

\begin{Shaded}
\begin{Highlighting}[]
\ControlFlowTok{if}\NormalTok{ (run\_intensive\_checks) \{}
\NormalTok{  npe\_model\_file }\OtherTok{\textless{}{-}} \FunctionTok{here}\NormalTok{(}\StringTok{"simdata"}\NormalTok{, }\StringTok{"ch14\_npe\_trained\_posterior.pkl"}\NormalTok{)}
  
  \ControlFlowTok{if}\NormalTok{ (regenerate\_simulations }\SpecialCharTok{||} \SpecialCharTok{!}\FunctionTok{file.exists}\NormalTok{(npe\_model\_file)) \{}
    \CommentTok{\# Convert to PyTorch tensors}
\NormalTok{    theta\_tensor }\OtherTok{\textless{}{-}}\NormalTok{ torch}\SpecialCharTok{$}\FunctionTok{tensor}\NormalTok{(train\_thetas, }\AttributeTok{dtype =}\NormalTok{ torch}\SpecialCharTok{$}\NormalTok{float32)}
  
    \CommentTok{\# Build data tensor: concatenate [response\_t, height\_t, position\_t, feedback\_t]}
    \CommentTok{\# Shape: (n\_sims, n\_trials * 4)}
\NormalTok{    x\_list }\OtherTok{\textless{}{-}} \FunctionTok{lapply}\NormalTok{(}\FunctionTok{seq\_len}\NormalTok{(n\_train\_sims), }\ControlFlowTok{function}\NormalTok{(i) \{}
\NormalTok{      resp }\OtherTok{\textless{}{-}}\NormalTok{ train\_y\_mat[i, ]}
      \FunctionTok{as.numeric}\NormalTok{(}\FunctionTok{t}\NormalTok{(}\FunctionTok{cbind}\NormalTok{(resp, context\_per\_trial)))}
\NormalTok{    \})}
\NormalTok{    x\_mat    }\OtherTok{\textless{}{-}} \FunctionTok{do.call}\NormalTok{(rbind, x\_list)}
\NormalTok{    x\_tensor }\OtherTok{\textless{}{-}}\NormalTok{ torch}\SpecialCharTok{$}\FunctionTok{tensor}\NormalTok{(x\_mat, }\AttributeTok{dtype =}\NormalTok{ torch}\SpecialCharTok{$}\NormalTok{float32)}
  
    \FunctionTok{cat}\NormalTok{(}\StringTok{"theta shape:"}\NormalTok{, }\FunctionTok{paste}\NormalTok{(theta\_tensor}\SpecialCharTok{$}\NormalTok{shape, }\AttributeTok{collapse =} \StringTok{" x "}\NormalTok{), }\StringTok{"}\SpecialCharTok{\textbackslash{}n}\StringTok{"}\NormalTok{)}
    \FunctionTok{cat}\NormalTok{(}\StringTok{"x shape:    "}\NormalTok{, }\FunctionTok{paste}\NormalTok{(x\_tensor}\SpecialCharTok{$}\NormalTok{shape, }\AttributeTok{collapse =} \StringTok{" x "}\NormalTok{), }\StringTok{"}\SpecialCharTok{\textbackslash{}n}\StringTok{"}\NormalTok{)}
  
    \CommentTok{\# Define prior bounds for sbi}
\NormalTok{    prior\_low  }\OtherTok{\textless{}{-}}\NormalTok{ torch}\SpecialCharTok{$}\FunctionTok{tensor}\NormalTok{(}\FunctionTok{c}\NormalTok{(LOGIT\_ERROR\_PRIOR\_MEAN }\SpecialCharTok{{-}} \DecValTok{3} \SpecialCharTok{*}\NormalTok{ LOGIT\_ERROR\_PRIOR\_SD,}
                                  \FunctionTok{log}\NormalTok{(}\DecValTok{20}\NormalTok{)),}
                                \AttributeTok{dtype =}\NormalTok{ torch}\SpecialCharTok{$}\NormalTok{float32)}
\NormalTok{    prior\_high }\OtherTok{\textless{}{-}}\NormalTok{ torch}\SpecialCharTok{$}\FunctionTok{tensor}\NormalTok{(}\FunctionTok{c}\NormalTok{(LOGIT\_ERROR\_PRIOR\_MEAN }\SpecialCharTok{+} \DecValTok{3} \SpecialCharTok{*}\NormalTok{ LOGIT\_ERROR\_PRIOR\_SD,}
                                  \FunctionTok{log}\NormalTok{(}\DecValTok{300}\NormalTok{)),}
                                \AttributeTok{dtype =}\NormalTok{ torch}\SpecialCharTok{$}\NormalTok{float32)}
\NormalTok{    sbi\_prior }\OtherTok{\textless{}{-}}\NormalTok{ utils}\SpecialCharTok{$}\FunctionTok{BoxUniform}\NormalTok{(}\AttributeTok{low =}\NormalTok{ prior\_low, }\AttributeTok{high =}\NormalTok{ prior\_high)}
  
    \CommentTok{\# Build and train the NPE}
\NormalTok{    inference\_obj }\OtherTok{\textless{}{-}}\NormalTok{ infer}\SpecialCharTok{$}\FunctionTok{SNPE}\NormalTok{(}
      \AttributeTok{prior             =}\NormalTok{ sbi\_prior,}
      \AttributeTok{density\_estimator =} \StringTok{"nsf"}   \CommentTok{\# neural spline flow}
\NormalTok{    )}
\NormalTok{    inference\_obj}\SpecialCharTok{$}\FunctionTok{append\_simulations}\NormalTok{(theta\_tensor, x\_tensor)}
  
    \FunctionTok{cat}\NormalTok{(}\StringTok{"Training NPE (this takes a few minutes on CPU)...}\SpecialCharTok{\textbackslash{}n}\StringTok{"}\NormalTok{)}
\NormalTok{    density\_estimator }\OtherTok{\textless{}{-}}\NormalTok{ inference\_obj}\SpecialCharTok{$}\FunctionTok{train}\NormalTok{(}
      \AttributeTok{training\_batch\_size =} \DecValTok{256}\NormalTok{L,}
      \AttributeTok{learning\_rate       =} \FloatTok{5e{-}4}\NormalTok{,}
      \AttributeTok{show\_train\_summary  =} \ConstantTok{TRUE}
\NormalTok{    )}
  
    \CommentTok{\# Build the posterior object}
\NormalTok{    npe\_posterior }\OtherTok{\textless{}{-}}\NormalTok{ inference\_obj}\SpecialCharTok{$}\FunctionTok{build\_posterior}\NormalTok{(density\_estimator)}
    \FunctionTok{cat}\NormalTok{(}\StringTok{"NPE training complete.}\SpecialCharTok{\textbackslash{}n}\StringTok{"}\NormalTok{)}
  
    \CommentTok{\# Save the trained model state}
\NormalTok{    torch}\SpecialCharTok{$}\FunctionTok{save}\NormalTok{(density\_estimator}\SpecialCharTok{$}\FunctionTok{state\_dict}\NormalTok{(), npe\_model\_file)}
    \FunctionTok{cat}\NormalTok{(}\StringTok{"Trained NPE saved to"}\NormalTok{, npe\_model\_file, }\StringTok{"}\SpecialCharTok{\textbackslash{}n}\StringTok{"}\NormalTok{)}
\NormalTok{  \} }\ControlFlowTok{else}\NormalTok{ \{}
    \FunctionTok{cat}\NormalTok{(}\StringTok{"Pre{-}trained NPE found at"}\NormalTok{, npe\_model\_file, }\StringTok{"}\SpecialCharTok{\textbackslash{}n}\StringTok{"}\NormalTok{)}
    \FunctionTok{cat}\NormalTok{(}\StringTok{"To use it, rebuild the sbi inference object and load the state dict:}\SpecialCharTok{\textbackslash{}n}\StringTok{"}\NormalTok{)}
    \FunctionTok{cat}\NormalTok{(}\StringTok{"  inference\_obj$build\_posterior() {-}\textgreater{} npe\_posterior}\SpecialCharTok{\textbackslash{}n}\StringTok{"}\NormalTok{)}
    \FunctionTok{cat}\NormalTok{(}\StringTok{"  (see sbi docs for state\_dict loading; requires same network architecture)}\SpecialCharTok{\textbackslash{}n}\StringTok{"}\NormalTok{)}
\NormalTok{  \}}
\NormalTok{\}}
\end{Highlighting}
\end{Shaded}

\subsection{Performing Inference on the Observed
Data}\label{performing-inference-on-the-observed-data}

After training, inference is instantaneous: evaluate the trained network
at the observed data \(x_{\text{obs}}\).

\begin{Shaded}
\begin{Highlighting}[]
\CommentTok{\# Build the observed data tensor (same format as training)}
\NormalTok{obs\_responses }\OtherTok{\textless{}{-}} \FunctionTok{as.integer}\NormalTok{(fit\_sim}\SpecialCharTok{$}\NormalTok{sim\_response)}
\NormalTok{obs\_x\_vec     }\OtherTok{\textless{}{-}} \FunctionTok{as.numeric}\NormalTok{(}\FunctionTok{t}\NormalTok{(}\FunctionTok{cbind}\NormalTok{(obs\_responses, context\_per\_trial)))}
\NormalTok{obs\_x\_tensor  }\OtherTok{\textless{}{-}}\NormalTok{ torch}\SpecialCharTok{$}\FunctionTok{tensor}\NormalTok{(}\FunctionTok{matrix}\NormalTok{(obs\_x\_vec, }\AttributeTok{nrow =} \DecValTok{1}\NormalTok{), }\AttributeTok{dtype =}\NormalTok{ torch}\SpecialCharTok{$}\NormalTok{float32)}

\CommentTok{\# Sample from the NPE posterior}
\NormalTok{n\_npe\_samples }\OtherTok{\textless{}{-}} \DecValTok{10000}\NormalTok{L}
\NormalTok{npe\_samples   }\OtherTok{\textless{}{-}}\NormalTok{ npe\_posterior}\SpecialCharTok{$}\FunctionTok{sample}\NormalTok{(}
  \AttributeTok{shape =} \FunctionTok{c}\NormalTok{(n\_npe\_samples),}
  \AttributeTok{x     =}\NormalTok{ obs\_x\_tensor}
\NormalTok{)}
\NormalTok{npe\_draws }\OtherTok{\textless{}{-}} \FunctionTok{as.matrix}\NormalTok{(npe\_samples}\SpecialCharTok{$}\FunctionTok{numpy}\NormalTok{())}
\FunctionTok{colnames}\NormalTok{(npe\_draws) }\OtherTok{\textless{}{-}} \FunctionTok{c}\NormalTok{(}\StringTok{"logit\_eps"}\NormalTok{, }\StringTok{"log\_n\_particles"}\NormalTok{)}

\NormalTok{npe\_eps    }\OtherTok{\textless{}{-}} \FunctionTok{plogis}\NormalTok{(npe\_draws[, }\StringTok{"logit\_eps"}\NormalTok{])}
\NormalTok{npe\_n\_part }\OtherTok{\textless{}{-}} \FunctionTok{exp}\NormalTok{(npe\_draws[, }\StringTok{"log\_n\_particles"}\NormalTok{])}

\FunctionTok{cat}\NormalTok{(}\StringTok{"NPE posterior for ε:           mean ="}\NormalTok{, }\FunctionTok{round}\NormalTok{(}\FunctionTok{mean}\NormalTok{(npe\_eps), }\DecValTok{4}\NormalTok{),}
    \StringTok{"  95\% CI = ["}\NormalTok{, }\FunctionTok{round}\NormalTok{(}\FunctionTok{quantile}\NormalTok{(npe\_eps, }\FloatTok{0.025}\NormalTok{), }\DecValTok{4}\NormalTok{), }\StringTok{","}\NormalTok{,}
    \FunctionTok{round}\NormalTok{(}\FunctionTok{quantile}\NormalTok{(npe\_eps, }\FloatTok{0.975}\NormalTok{), }\DecValTok{4}\NormalTok{), }\StringTok{"]}\SpecialCharTok{\textbackslash{}n}\StringTok{"}\NormalTok{)}
\FunctionTok{cat}\NormalTok{(}\StringTok{"NPE posterior for N\_particles: mean ="}\NormalTok{, }\FunctionTok{round}\NormalTok{(}\FunctionTok{mean}\NormalTok{(npe\_n\_part), }\DecValTok{1}\NormalTok{),}
    \StringTok{"  95\% CI = ["}\NormalTok{, }\FunctionTok{round}\NormalTok{(}\FunctionTok{quantile}\NormalTok{(npe\_n\_part, }\FloatTok{0.025}\NormalTok{), }\DecValTok{1}\NormalTok{), }\StringTok{","}\NormalTok{,}
    \FunctionTok{round}\NormalTok{(}\FunctionTok{quantile}\NormalTok{(npe\_n\_part, }\FloatTok{0.975}\NormalTok{), }\DecValTok{1}\NormalTok{), }\StringTok{"]}\SpecialCharTok{\textbackslash{}n}\StringTok{"}\NormalTok{)}

\NormalTok{p\_eps }\OtherTok{\textless{}{-}} \FunctionTok{ggplot}\NormalTok{(}\FunctionTok{tibble}\NormalTok{(}\AttributeTok{eps =}\NormalTok{ npe\_eps), }\FunctionTok{aes}\NormalTok{(}\AttributeTok{x =}\NormalTok{ eps)) }\SpecialCharTok{+}
  \FunctionTok{geom\_histogram}\NormalTok{(}\FunctionTok{aes}\NormalTok{(}\AttributeTok{y =} \FunctionTok{after\_stat}\NormalTok{(density)),}
                 \AttributeTok{fill =} \StringTok{"\#009E73"}\NormalTok{, }\AttributeTok{alpha =} \FloatTok{0.7}\NormalTok{, }\AttributeTok{bins =} \DecValTok{40}\NormalTok{) }\SpecialCharTok{+}
  \FunctionTok{geom\_vline}\NormalTok{(}\AttributeTok{xintercept =} \FloatTok{0.1}\NormalTok{, }\AttributeTok{color =} \StringTok{"red"}\NormalTok{, }\AttributeTok{linetype =} \StringTok{"dashed"}\NormalTok{) }\SpecialCharTok{+}
  \FunctionTok{labs}\NormalTok{(}\AttributeTok{title =} \StringTok{"NPE Posterior for ε (Model A)"}\NormalTok{,}
       \AttributeTok{subtitle =} \StringTok{"Red dashed: true generating ε = 0.1"}\NormalTok{,}
       \AttributeTok{x =} \StringTok{"ε"}\NormalTok{, }\AttributeTok{y =} \StringTok{"Density"}\NormalTok{)}

\NormalTok{p\_npart }\OtherTok{\textless{}{-}} \FunctionTok{ggplot}\NormalTok{(}\FunctionTok{tibble}\NormalTok{(}\AttributeTok{n\_part =}\NormalTok{ npe\_n\_part), }\FunctionTok{aes}\NormalTok{(}\AttributeTok{x =}\NormalTok{ n\_part)) }\SpecialCharTok{+}
  \FunctionTok{geom\_histogram}\NormalTok{(}\FunctionTok{aes}\NormalTok{(}\AttributeTok{y =} \FunctionTok{after\_stat}\NormalTok{(density)),}
                 \AttributeTok{fill =} \StringTok{"\#0072B2"}\NormalTok{, }\AttributeTok{alpha =} \FloatTok{0.7}\NormalTok{, }\AttributeTok{bins =} \DecValTok{40}\NormalTok{) }\SpecialCharTok{+}
  \FunctionTok{geom\_vline}\NormalTok{(}\AttributeTok{xintercept =} \DecValTok{100}\NormalTok{, }\AttributeTok{color =} \StringTok{"red"}\NormalTok{, }\AttributeTok{linetype =} \StringTok{"dashed"}\NormalTok{) }\SpecialCharTok{+}
  \FunctionTok{labs}\NormalTok{(}\AttributeTok{title =} \StringTok{"NPE Posterior for N\_particles (Model A)"}\NormalTok{,}
       \AttributeTok{subtitle =} \StringTok{"Red dashed: true generating N\_particles = 100"}\NormalTok{,}
       \AttributeTok{x =} \StringTok{"N\_particles"}\NormalTok{, }\AttributeTok{y =} \StringTok{"Density"}\NormalTok{)}

\FunctionTok{print}\NormalTok{(p\_eps }\SpecialCharTok{/}\NormalTok{ p\_npart)}
\end{Highlighting}
\end{Shaded}

\textbf{Reading the NPE posterior.} Two things to check: (1) both
posteriors should contain the true generating values, and (2)
\(N_{\text{particles}}\) should show a meaningful posterior --- if it is
prior-like, the data are uninformative about working-memory capacity for
this task length. A narrow posterior on \(N_{\text{particles}}\) would
mean the task is diagnostic for this cognitive parameter, which is a
substantive finding.

\subsection{NPE Validation: SBC + Expected Coverage (Phase
6)}\label{npe-validation-sbc-expected-coverage-phase-6}

The persona's Phase 6 requires calibration checks. For NPE, the standard
diagnostics are marginal SBC (Talts et al.~2018) and expected coverage
(Hermans et al.~2022). Because NPE is amortized, running 1000
calibration checks takes seconds --- each requires only a forward pass
through the trained network.

\begin{Shaded}
\begin{Highlighting}[]
\NormalTok{n\_cal }\OtherTok{\textless{}{-}} \DecValTok{1000}\NormalTok{L}

\FunctionTok{cat}\NormalTok{(}\StringTok{"Running NPE calibration diagnostics ("}\NormalTok{, }\FunctionTok{as.integer}\NormalTok{(n\_cal), }\StringTok{" checks)...}\SpecialCharTok{\textbackslash{}n}\StringTok{"}\NormalTok{)}

\CommentTok{\# Generate calibration set: draw theta from prior, simulate x}
\FunctionTok{set.seed}\NormalTok{(}\DecValTok{7777}\NormalTok{)}
\NormalTok{cal\_logit\_eps }\OtherTok{\textless{}{-}} \FunctionTok{rnorm}\NormalTok{(n\_cal, LOGIT\_ERROR\_PRIOR\_MEAN, LOGIT\_ERROR\_PRIOR\_SD)}
\NormalTok{cal\_log\_npart }\OtherTok{\textless{}{-}} \FunctionTok{runif}\NormalTok{(n\_cal, }\FunctionTok{log}\NormalTok{(}\DecValTok{20}\NormalTok{), }\FunctionTok{log}\NormalTok{(}\DecValTok{300}\NormalTok{))}
\NormalTok{cal\_thetas    }\OtherTok{\textless{}{-}} \FunctionTok{cbind}\NormalTok{(cal\_logit\_eps, cal\_log\_npart)}

\NormalTok{cal\_responses }\OtherTok{\textless{}{-}} \FunctionTok{future\_map}\NormalTok{(}
  \FunctionTok{seq\_len}\NormalTok{(n\_cal),}
  \ControlFlowTok{function}\NormalTok{(i) }\FunctionTok{simulate\_model\_a}\NormalTok{(cal\_thetas[i, ]),}
  \AttributeTok{.options =} \FunctionTok{furrr\_options}\NormalTok{(}\AttributeTok{seed =} \ConstantTok{TRUE}\NormalTok{),}
  \AttributeTok{.progress =} \ConstantTok{TRUE}
\NormalTok{)}
\NormalTok{cal\_y\_mat }\OtherTok{\textless{}{-}} \FunctionTok{do.call}\NormalTok{(rbind, cal\_responses)}

\CommentTok{\# Build calibration data tensors}
\NormalTok{cal\_x\_list }\OtherTok{\textless{}{-}} \FunctionTok{lapply}\NormalTok{(}\FunctionTok{seq\_len}\NormalTok{(n\_cal), }\ControlFlowTok{function}\NormalTok{(i) \{}
\NormalTok{  resp }\OtherTok{\textless{}{-}}\NormalTok{ cal\_y\_mat[i, ]}
  \FunctionTok{as.numeric}\NormalTok{(}\FunctionTok{t}\NormalTok{(}\FunctionTok{cbind}\NormalTok{(resp, context\_per\_trial)))}
\NormalTok{\})}
\NormalTok{cal\_x\_mat    }\OtherTok{\textless{}{-}} \FunctionTok{do.call}\NormalTok{(rbind, cal\_x\_list)}
\NormalTok{cal\_x\_tensor }\OtherTok{\textless{}{-}}\NormalTok{ torch}\SpecialCharTok{$}\FunctionTok{tensor}\NormalTok{(cal\_x\_mat, }\AttributeTok{dtype =}\NormalTok{ torch}\SpecialCharTok{$}\NormalTok{float32)}

\CommentTok{\# ── Marginal SBC ──────────────────────────────────────────────────────────────}
\NormalTok{n\_post\_samples }\OtherTok{\textless{}{-}} \DecValTok{1000}\NormalTok{L}
\NormalTok{sbc\_ranks }\OtherTok{\textless{}{-}} \FunctionTok{matrix}\NormalTok{(}\DecValTok{0}\NormalTok{L, }\AttributeTok{nrow =}\NormalTok{ n\_cal, }\AttributeTok{ncol =} \DecValTok{2}\NormalTok{)}
\FunctionTok{colnames}\NormalTok{(sbc\_ranks) }\OtherTok{\textless{}{-}} \FunctionTok{c}\NormalTok{(}\StringTok{"logit\_eps"}\NormalTok{, }\StringTok{"log\_n\_particles"}\NormalTok{)}

\ControlFlowTok{for}\NormalTok{ (i }\ControlFlowTok{in} \FunctionTok{seq\_len}\NormalTok{(n\_cal)) \{}
\NormalTok{  post\_samples }\OtherTok{\textless{}{-}}\NormalTok{ npe\_posterior}\SpecialCharTok{$}\FunctionTok{sample}\NormalTok{(}
    \AttributeTok{shape =} \FunctionTok{c}\NormalTok{(n\_post\_samples),}
    \AttributeTok{x     =}\NormalTok{ cal\_x\_tensor[i }\SpecialCharTok{{-}} \DecValTok{1}\NormalTok{L]}\SpecialCharTok{$}\FunctionTok{unsqueeze}\NormalTok{(}\DecValTok{0}\NormalTok{L)}
\NormalTok{  )}\SpecialCharTok{$}\FunctionTok{numpy}\NormalTok{()}

  \ControlFlowTok{for}\NormalTok{ (j }\ControlFlowTok{in} \DecValTok{1}\SpecialCharTok{:}\DecValTok{2}\NormalTok{) \{}
\NormalTok{    sbc\_ranks[i, j] }\OtherTok{\textless{}{-}} \FunctionTok{sum}\NormalTok{(post\_samples[, j] }\SpecialCharTok{\textless{}}\NormalTok{ cal\_thetas[i, j])}
\NormalTok{  \}}
\NormalTok{\}}

\NormalTok{sbc\_df }\OtherTok{\textless{}{-}} \FunctionTok{as\_tibble}\NormalTok{(sbc\_ranks) }\SpecialCharTok{|\textgreater{}}
  \FunctionTok{pivot\_longer}\NormalTok{(}\FunctionTok{everything}\NormalTok{(), }\AttributeTok{names\_to =} \StringTok{"parameter"}\NormalTok{, }\AttributeTok{values\_to =} \StringTok{"rank"}\NormalTok{)}

\NormalTok{p\_sbc }\OtherTok{\textless{}{-}} \FunctionTok{ggplot}\NormalTok{(sbc\_df, }\FunctionTok{aes}\NormalTok{(}\AttributeTok{x =}\NormalTok{ rank)) }\SpecialCharTok{+}
  \FunctionTok{geom\_histogram}\NormalTok{(}\FunctionTok{aes}\NormalTok{(}\AttributeTok{y =} \FunctionTok{after\_stat}\NormalTok{(density)), }\AttributeTok{bins =} \DecValTok{20}\NormalTok{,}
                 \AttributeTok{fill =} \StringTok{"\#009E73"}\NormalTok{, }\AttributeTok{alpha =} \FloatTok{0.7}\NormalTok{) }\SpecialCharTok{+}
  \FunctionTok{facet\_wrap}\NormalTok{(}\SpecialCharTok{\textasciitilde{}}\NormalTok{parameter, }\AttributeTok{scales =} \StringTok{"free\_x"}\NormalTok{) }\SpecialCharTok{+}
  \FunctionTok{labs}\NormalTok{(}\AttributeTok{title =} \StringTok{"Marginal SBC Rank Histograms (NPE, Model A)"}\NormalTok{,}
       \AttributeTok{subtitle =} \StringTok{"Uniform = calibrated. U{-}shape = overconfident. Hump = underconfident."}\NormalTok{,}
       \AttributeTok{x =} \StringTok{"Rank"}\NormalTok{, }\AttributeTok{y =} \StringTok{"Count"}\NormalTok{)}

\CommentTok{\# ── Expected Coverage ─────────────────────────────────────────────────────────}
\NormalTok{coverage\_levels }\OtherTok{\textless{}{-}} \FunctionTok{seq}\NormalTok{(}\FloatTok{0.05}\NormalTok{, }\FloatTok{0.95}\NormalTok{, }\AttributeTok{by =} \FloatTok{0.05}\NormalTok{)}
\NormalTok{empirical\_coverage }\OtherTok{\textless{}{-}} \FunctionTok{vapply}\NormalTok{(coverage\_levels, }\ControlFlowTok{function}\NormalTok{(alpha) \{}
\NormalTok{  covered }\OtherTok{\textless{}{-}} \FunctionTok{vapply}\NormalTok{(}\FunctionTok{seq\_len}\NormalTok{(n\_cal), }\ControlFlowTok{function}\NormalTok{(i) \{}
\NormalTok{    post\_samples }\OtherTok{\textless{}{-}}\NormalTok{ npe\_posterior}\SpecialCharTok{$}\FunctionTok{sample}\NormalTok{(}
      \AttributeTok{shape =} \FunctionTok{c}\NormalTok{(n\_post\_samples),}
      \AttributeTok{x     =}\NormalTok{ cal\_x\_tensor[i }\SpecialCharTok{{-}} \DecValTok{1}\NormalTok{L]}\SpecialCharTok{$}\FunctionTok{unsqueeze}\NormalTok{(}\DecValTok{0}\NormalTok{L)}
\NormalTok{    )}\SpecialCharTok{$}\FunctionTok{numpy}\NormalTok{()}
\NormalTok{    q\_lo }\OtherTok{\textless{}{-}} \FunctionTok{quantile}\NormalTok{(post\_samples[, }\DecValTok{1}\NormalTok{], (}\DecValTok{1} \SpecialCharTok{{-}}\NormalTok{ alpha) }\SpecialCharTok{/} \DecValTok{2}\NormalTok{)}
\NormalTok{    q\_hi }\OtherTok{\textless{}{-}} \FunctionTok{quantile}\NormalTok{(post\_samples[, }\DecValTok{1}\NormalTok{], }\DecValTok{1} \SpecialCharTok{{-}}\NormalTok{ (}\DecValTok{1} \SpecialCharTok{{-}}\NormalTok{ alpha) }\SpecialCharTok{/} \DecValTok{2}\NormalTok{)}
\NormalTok{    cal\_thetas[i, }\DecValTok{1}\NormalTok{] }\SpecialCharTok{\textgreater{}=}\NormalTok{ q\_lo }\SpecialCharTok{\&}\NormalTok{ cal\_thetas[i, }\DecValTok{1}\NormalTok{] }\SpecialCharTok{\textless{}=}\NormalTok{ q\_hi}
\NormalTok{  \}, }\FunctionTok{logical}\NormalTok{(}\DecValTok{1}\NormalTok{))}
  \FunctionTok{mean}\NormalTok{(covered)}
\NormalTok{\}, }\FunctionTok{numeric}\NormalTok{(}\DecValTok{1}\NormalTok{))}

\NormalTok{coverage\_df }\OtherTok{\textless{}{-}} \FunctionTok{tibble}\NormalTok{(}\AttributeTok{nominal =}\NormalTok{ coverage\_levels, }\AttributeTok{empirical =}\NormalTok{ empirical\_coverage)}

\NormalTok{p\_coverage }\OtherTok{\textless{}{-}} \FunctionTok{ggplot}\NormalTok{(coverage\_df, }\FunctionTok{aes}\NormalTok{(}\AttributeTok{x =}\NormalTok{ nominal, }\AttributeTok{y =}\NormalTok{ empirical)) }\SpecialCharTok{+}
  \FunctionTok{geom\_abline}\NormalTok{(}\AttributeTok{slope =} \DecValTok{1}\NormalTok{, }\AttributeTok{intercept =} \DecValTok{0}\NormalTok{, }\AttributeTok{linetype =} \StringTok{"dashed"}\NormalTok{, }\AttributeTok{color =} \StringTok{"grey50"}\NormalTok{) }\SpecialCharTok{+}
  \FunctionTok{geom\_ribbon}\NormalTok{(}\FunctionTok{aes}\NormalTok{(}\AttributeTok{ymin =}\NormalTok{ nominal }\SpecialCharTok{{-}} \DecValTok{2} \SpecialCharTok{*} \FunctionTok{sqrt}\NormalTok{(nominal }\SpecialCharTok{*}\NormalTok{ (}\DecValTok{1} \SpecialCharTok{{-}}\NormalTok{ nominal) }\SpecialCharTok{/}\NormalTok{ n\_cal),}
                  \AttributeTok{ymax =}\NormalTok{ nominal }\SpecialCharTok{+} \DecValTok{2} \SpecialCharTok{*} \FunctionTok{sqrt}\NormalTok{(nominal }\SpecialCharTok{*}\NormalTok{ (}\DecValTok{1} \SpecialCharTok{{-}}\NormalTok{ nominal) }\SpecialCharTok{/}\NormalTok{ n\_cal)),}
              \AttributeTok{fill =} \StringTok{"grey80"}\NormalTok{, }\AttributeTok{alpha =} \FloatTok{0.5}\NormalTok{) }\SpecialCharTok{+}
  \FunctionTok{geom\_line}\NormalTok{(}\AttributeTok{color =} \StringTok{"\#0072B2"}\NormalTok{, }\AttributeTok{linewidth =} \DecValTok{1}\NormalTok{) }\SpecialCharTok{+}
  \FunctionTok{geom\_point}\NormalTok{(}\AttributeTok{color =} \StringTok{"\#0072B2"}\NormalTok{, }\AttributeTok{size =} \DecValTok{2}\NormalTok{) }\SpecialCharTok{+}
  \FunctionTok{labs}\NormalTok{(}\AttributeTok{title =} \StringTok{"Expected Coverage (NPE, Model A)"}\NormalTok{,}
       \AttributeTok{subtitle =} \StringTok{"Points on diagonal = well{-}calibrated. Grey band = 2SE."}\NormalTok{,}
       \AttributeTok{x =} \StringTok{"Nominal coverage"}\NormalTok{, }\AttributeTok{y =} \StringTok{"Empirical coverage"}\NormalTok{) }\SpecialCharTok{+}
  \FunctionTok{coord\_equal}\NormalTok{(}\AttributeTok{xlim =} \FunctionTok{c}\NormalTok{(}\DecValTok{0}\NormalTok{, }\DecValTok{1}\NormalTok{), }\AttributeTok{ylim =} \FunctionTok{c}\NormalTok{(}\DecValTok{0}\NormalTok{, }\DecValTok{1}\NormalTok{))}

\FunctionTok{print}\NormalTok{(p\_sbc }\SpecialCharTok{/}\NormalTok{ p\_coverage)}

\FunctionTok{cat}\NormalTok{(}\StringTok{"}\SpecialCharTok{\textbackslash{}n}\StringTok{{-}{-}{-} Coverage summary {-}{-}{-}}\SpecialCharTok{\textbackslash{}n}\StringTok{"}\NormalTok{)}
\FunctionTok{cat}\NormalTok{(}\StringTok{"50\% nominal coverage: empirical ="}\NormalTok{,}
    \FunctionTok{round}\NormalTok{(coverage\_df}\SpecialCharTok{$}\NormalTok{empirical[coverage\_df}\SpecialCharTok{$}\NormalTok{nominal }\SpecialCharTok{==} \FloatTok{0.50}\NormalTok{], }\DecValTok{3}\NormalTok{), }\StringTok{"}\SpecialCharTok{\textbackslash{}n}\StringTok{"}\NormalTok{)}
\FunctionTok{cat}\NormalTok{(}\StringTok{"95\% nominal coverage: empirical ="}\NormalTok{,}
    \FunctionTok{round}\NormalTok{(coverage\_df}\SpecialCharTok{$}\NormalTok{empirical[coverage\_df}\SpecialCharTok{$}\NormalTok{nominal }\SpecialCharTok{==} \FloatTok{0.95}\NormalTok{], }\DecValTok{3}\NormalTok{), }\StringTok{"}\SpecialCharTok{\textbackslash{}n}\StringTok{"}\NormalTok{)}
\end{Highlighting}
\end{Shaded}

\textbf{Reading the validation.} The SBC rank histograms should be
uniform within sampling noise. The coverage plot should track the
diagonal within the grey 2-SE band. If the SBC histogram is U-shaped,
the posterior is overconfident --- increase training data or use a more
expressive flow. If it is humped, the posterior is underconfident. Both
checks benefit enormously from amortization: each of the 1000
calibration checks requires only a forward pass.

\subsection{The Main Empirical Result: NPE vs Stan
Posteriors}\label{the-main-empirical-result-npe-vs-stan-posteriors}

Here is the comparison the chapter has been building toward. We fit the
same observed dataset under both Model A (via NPE) and Model B (via HMC)
and overlay the two posteriors on \(\varepsilon\). \textbf{The gap
between them is the structural bias you accept when you replace the
cognitive theory with a tractable proxy.}

\begin{Shaded}
\begin{Highlighting}[]
\ControlFlowTok{if}\NormalTok{ (}\FunctionTok{exists}\NormalTok{(}\StringTok{"npe\_eps"}\NormalTok{) }\SpecialCharTok{\&\&} \SpecialCharTok{!}\FunctionTok{is.null}\NormalTok{(fit\_rule\_stan)) \{}
\NormalTok{  stan\_eps }\OtherTok{\textless{}{-}} \FunctionTok{as\_draws\_df}\NormalTok{(fit\_rule\_stan}\SpecialCharTok{$}\FunctionTok{draws}\NormalTok{(}\StringTok{"error\_prob"}\NormalTok{))}\SpecialCharTok{$}\NormalTok{error\_prob}

\NormalTok{  comparison\_df }\OtherTok{\textless{}{-}} \FunctionTok{bind\_rows}\NormalTok{(}
    \FunctionTok{tibble}\NormalTok{(}\AttributeTok{eps =}\NormalTok{ npe\_eps, }\AttributeTok{method =} \StringTok{"Model A (NPE, particle filter)"}\NormalTok{),}
    \FunctionTok{tibble}\NormalTok{(}\AttributeTok{eps =}\NormalTok{ stan\_eps, }\AttributeTok{method =} \StringTok{"Model B (HMC, marginalised mixture)"}\NormalTok{)}
\NormalTok{  )}

  \FunctionTok{ggplot}\NormalTok{(comparison\_df, }\FunctionTok{aes}\NormalTok{(}\AttributeTok{x =}\NormalTok{ eps, }\AttributeTok{fill =}\NormalTok{ method, }\AttributeTok{color =}\NormalTok{ method)) }\SpecialCharTok{+}
    \FunctionTok{geom\_density}\NormalTok{(}\AttributeTok{alpha =} \FloatTok{0.4}\NormalTok{, }\AttributeTok{linewidth =} \FloatTok{0.6}\NormalTok{) }\SpecialCharTok{+}
    \FunctionTok{geom\_vline}\NormalTok{(}\AttributeTok{xintercept =} \FloatTok{0.1}\NormalTok{, }\AttributeTok{color =} \StringTok{"red"}\NormalTok{, }\AttributeTok{linetype =} \StringTok{"dashed"}\NormalTok{) }\SpecialCharTok{+}
    \FunctionTok{scale\_fill\_manual}\NormalTok{(}\AttributeTok{values =} \FunctionTok{c}\NormalTok{(}\StringTok{"Model A (NPE, particle filter)"}      \OtherTok{=} \StringTok{"\#009E73"}\NormalTok{,}
                                  \StringTok{"Model B (HMC, marginalised mixture)"} \OtherTok{=} \StringTok{"\#E69F00"}\NormalTok{)) }\SpecialCharTok{+}
    \FunctionTok{scale\_color\_manual}\NormalTok{(}\AttributeTok{values =} \FunctionTok{c}\NormalTok{(}\StringTok{"Model A (NPE, particle filter)"}      \OtherTok{=} \StringTok{"\#006D4E"}\NormalTok{,}
                                   \StringTok{"Model B (HMC, marginalised mixture)"} \OtherTok{=} \StringTok{"\#B07D00"}\NormalTok{)) }\SpecialCharTok{+}
    \FunctionTok{labs}\NormalTok{(}
      \AttributeTok{title    =} \StringTok{"NPE vs HMC: Posterior on ε for Model A vs Model B"}\NormalTok{,}
      \AttributeTok{subtitle =} \StringTok{"Red dashed: true generating ε = 0.1}\SpecialCharTok{\textbackslash{}n}\StringTok{The gap between these distributions IS the structural bias of the proxy."}\NormalTok{,}
      \AttributeTok{x =} \StringTok{"ε"}\NormalTok{, }\AttributeTok{y =} \StringTok{"Density"}
\NormalTok{    )}

  \FunctionTok{cat}\NormalTok{(}\StringTok{"}\SpecialCharTok{\textbackslash{}n}\StringTok{{-}{-}{-} Posterior summaries {-}{-}{-}}\SpecialCharTok{\textbackslash{}n}\StringTok{"}\NormalTok{)}
  \FunctionTok{cat}\NormalTok{(}\StringTok{"Model A (NPE) posterior:    mean ="}\NormalTok{, }\FunctionTok{round}\NormalTok{(}\FunctionTok{mean}\NormalTok{(npe\_eps), }\DecValTok{4}\NormalTok{),}
      \StringTok{"  95\% CI = ["}\NormalTok{, }\FunctionTok{round}\NormalTok{(}\FunctionTok{quantile}\NormalTok{(npe\_eps, }\FloatTok{0.025}\NormalTok{), }\DecValTok{4}\NormalTok{), }\StringTok{","}\NormalTok{,}
                       \FunctionTok{round}\NormalTok{(}\FunctionTok{quantile}\NormalTok{(npe\_eps, }\FloatTok{0.975}\NormalTok{), }\DecValTok{4}\NormalTok{), }\StringTok{"]}\SpecialCharTok{\textbackslash{}n}\StringTok{"}\NormalTok{)}
  \FunctionTok{cat}\NormalTok{(}\StringTok{"Model B (HMC) posterior:    mean ="}\NormalTok{, }\FunctionTok{round}\NormalTok{(}\FunctionTok{mean}\NormalTok{(stan\_eps), }\DecValTok{4}\NormalTok{),}
      \StringTok{"  95\% CI = ["}\NormalTok{, }\FunctionTok{round}\NormalTok{(}\FunctionTok{quantile}\NormalTok{(stan\_eps, }\FloatTok{0.025}\NormalTok{), }\DecValTok{4}\NormalTok{), }\StringTok{","}\NormalTok{,}
                       \FunctionTok{round}\NormalTok{(}\FunctionTok{quantile}\NormalTok{(stan\_eps, }\FloatTok{0.975}\NormalTok{), }\DecValTok{4}\NormalTok{), }\StringTok{"]}\SpecialCharTok{\textbackslash{}n}\StringTok{"}\NormalTok{)}
  \FunctionTok{cat}\NormalTok{(}\StringTok{"True generating ε:           0.1}\SpecialCharTok{\textbackslash{}n}\StringTok{"}\NormalTok{)}
\NormalTok{\}}
\end{Highlighting}
\end{Shaded}

\subsection{Demonstrating That the Gap Tracks Rule-Set
Mismatch}\label{demonstrating-that-the-gap-tracks-rule-set-mismatch}

The headline claim --- ``the gap is structural bias, not inference
noise'' --- requires an experiment. We generate datasets from multiple
seeds and compare the NPE-vs-Stan gap. The gap should be small when
Model A converges to a rule that is in Model B's candidate set, and
large when it does not.

\begin{Shaded}
\begin{Highlighting}[]
\ControlFlowTok{if}\NormalTok{ (}\FunctionTok{exists}\NormalTok{(}\StringTok{"npe\_posterior"}\NormalTok{) }\SpecialCharTok{\&\&} \SpecialCharTok{!}\FunctionTok{is.null}\NormalTok{(mod\_rule\_stan)) \{}
\NormalTok{  gap\_seeds }\OtherTok{\textless{}{-}} \FunctionTok{c}\NormalTok{(}\DecValTok{42}\NormalTok{, }\DecValTok{123}\NormalTok{, }\DecValTok{7}\NormalTok{, }\DecValTok{99}\NormalTok{, }\DecValTok{55}\NormalTok{)}
\NormalTok{  gap\_results }\OtherTok{\textless{}{-}} \FunctionTok{map\_dfr}\NormalTok{(gap\_seeds, }\ControlFlowTok{function}\NormalTok{(s) \{}
\NormalTok{    sched\_g }\OtherTok{\textless{}{-}} \FunctionTok{make\_subject\_schedule}\NormalTok{(stimulus\_info, n\_blocks, }\AttributeTok{seed =}\NormalTok{ s)}
\NormalTok{    sim\_g   }\OtherTok{\textless{}{-}} \FunctionTok{rule\_particle\_filter}\NormalTok{(}
      \AttributeTok{n\_particles =}\NormalTok{ N\_PARTICLES\_FIXED, }\AttributeTok{error\_prob =} \FloatTok{0.1}\NormalTok{,}
      \AttributeTok{obs =} \FunctionTok{as.matrix}\NormalTok{(sched\_g[, }\FunctionTok{c}\NormalTok{(}\StringTok{"height"}\NormalTok{, }\StringTok{"position"}\NormalTok{)]),}
      \AttributeTok{cat\_one =}\NormalTok{ sched\_g}\SpecialCharTok{$}\NormalTok{category\_feedback, }\AttributeTok{max\_dims =}\NormalTok{ MAX\_RULE\_DIMS, }\AttributeTok{quiet =} \ConstantTok{TRUE}
\NormalTok{    )}

    \CommentTok{\# NPE inference (amortized — instant)}
\NormalTok{    obs\_g\_vec }\OtherTok{\textless{}{-}} \FunctionTok{as.numeric}\NormalTok{(}\FunctionTok{t}\NormalTok{(}\FunctionTok{cbind}\NormalTok{(}\FunctionTok{as.integer}\NormalTok{(sim\_g}\SpecialCharTok{$}\NormalTok{sim\_response),}
                                     \FunctionTok{as.matrix}\NormalTok{(sched\_g[, }\FunctionTok{c}\NormalTok{(}\StringTok{"height"}\NormalTok{, }\StringTok{"position"}\NormalTok{)]),}
\NormalTok{                                     sched\_g}\SpecialCharTok{$}\NormalTok{category\_feedback)))}
\NormalTok{    obs\_g\_tensor }\OtherTok{\textless{}{-}}\NormalTok{ torch}\SpecialCharTok{$}\FunctionTok{tensor}\NormalTok{(}\FunctionTok{matrix}\NormalTok{(obs\_g\_vec, }\AttributeTok{nrow =} \DecValTok{1}\NormalTok{), }\AttributeTok{dtype =}\NormalTok{ torch}\SpecialCharTok{$}\NormalTok{float32)}
\NormalTok{    npe\_g\_samples }\OtherTok{\textless{}{-}}\NormalTok{ npe\_posterior}\SpecialCharTok{$}\FunctionTok{sample}\NormalTok{(}\FunctionTok{c}\NormalTok{(}\DecValTok{5000}\NormalTok{L), }\AttributeTok{x =}\NormalTok{ obs\_g\_tensor)}\SpecialCharTok{$}\FunctionTok{numpy}\NormalTok{()}
\NormalTok{    npe\_g\_eps }\OtherTok{\textless{}{-}} \FunctionTok{plogis}\NormalTok{(npe\_g\_samples[, }\DecValTok{1}\NormalTok{])}

    \CommentTok{\# Stan inference}
\NormalTok{    stan\_dat\_g }\OtherTok{\textless{}{-}} \FunctionTok{list}\NormalTok{(}
      \AttributeTok{ntrials =} \FunctionTok{nrow}\NormalTok{(sched\_g), }\AttributeTok{nfeatures =} \DecValTok{2}\NormalTok{L,}
      \AttributeTok{y =} \FunctionTok{as.integer}\NormalTok{(sim\_g}\SpecialCharTok{$}\NormalTok{sim\_response),}
      \AttributeTok{obs =} \FunctionTok{as.matrix}\NormalTok{(sched\_g[, }\FunctionTok{c}\NormalTok{(}\StringTok{"height"}\NormalTok{, }\StringTok{"position"}\NormalTok{)]),}
      \AttributeTok{prior\_logit\_error\_mean =}\NormalTok{ LOGIT\_ERROR\_PRIOR\_MEAN,}
      \AttributeTok{prior\_logit\_error\_sd   =}\NormalTok{ LOGIT\_ERROR\_PRIOR\_SD}
\NormalTok{    )}
\NormalTok{    fit\_g }\OtherTok{\textless{}{-}}\NormalTok{ mod\_rule\_stan}\SpecialCharTok{$}\FunctionTok{sample}\NormalTok{(}
      \AttributeTok{data =}\NormalTok{ stan\_dat\_g, }\AttributeTok{seed =}\NormalTok{ s, }\AttributeTok{chains =} \DecValTok{2}\NormalTok{, }\AttributeTok{parallel\_chains =} \DecValTok{2}\NormalTok{,}
      \AttributeTok{iter\_warmup =} \DecValTok{800}\NormalTok{, }\AttributeTok{iter\_sampling =} \DecValTok{1000}\NormalTok{, }\AttributeTok{refresh =} \DecValTok{0}\NormalTok{, }\AttributeTok{adapt\_delta =} \FloatTok{0.9}
\NormalTok{    )}
\NormalTok{    stan\_g\_eps }\OtherTok{\textless{}{-}} \FunctionTok{as\_draws\_df}\NormalTok{(fit\_g}\SpecialCharTok{$}\FunctionTok{draws}\NormalTok{(}\StringTok{"error\_prob"}\NormalTok{))}\SpecialCharTok{$}\NormalTok{error\_prob}

    \FunctionTok{tibble}\NormalTok{(}\AttributeTok{seed =}\NormalTok{ s,}
           \AttributeTok{npe\_mean =} \FunctionTok{mean}\NormalTok{(npe\_g\_eps), }\AttributeTok{stan\_mean =} \FunctionTok{mean}\NormalTok{(stan\_g\_eps),}
           \AttributeTok{gap =} \FunctionTok{abs}\NormalTok{(}\FunctionTok{mean}\NormalTok{(npe\_g\_eps) }\SpecialCharTok{{-}} \FunctionTok{mean}\NormalTok{(stan\_g\_eps)),}
           \AttributeTok{npe\_covers\_truth  =}\NormalTok{ (}\FloatTok{0.1} \SpecialCharTok{\textgreater{}=} \FunctionTok{quantile}\NormalTok{(npe\_g\_eps, }\FloatTok{0.025}\NormalTok{)) }\SpecialCharTok{\&}\NormalTok{ (}\FloatTok{0.1} \SpecialCharTok{\textless{}=} \FunctionTok{quantile}\NormalTok{(npe\_g\_eps, }\FloatTok{0.975}\NormalTok{)),}
           \AttributeTok{stan\_covers\_truth =}\NormalTok{ (}\FloatTok{0.1} \SpecialCharTok{\textgreater{}=} \FunctionTok{quantile}\NormalTok{(stan\_g\_eps, }\FloatTok{0.025}\NormalTok{)) }\SpecialCharTok{\&}\NormalTok{ (}\FloatTok{0.1} \SpecialCharTok{\textless{}=} \FunctionTok{quantile}\NormalTok{(stan\_g\_eps, }\FloatTok{0.975}\NormalTok{)))}
\NormalTok{  \})}

  \FunctionTok{print}\NormalTok{(gap\_results)}

  \FunctionTok{ggplot}\NormalTok{(gap\_results, }\FunctionTok{aes}\NormalTok{(}\AttributeTok{x =} \FunctionTok{factor}\NormalTok{(seed), }\AttributeTok{y =}\NormalTok{ gap)) }\SpecialCharTok{+}
    \FunctionTok{geom\_col}\NormalTok{(}\AttributeTok{fill =} \StringTok{"\#E69F00"}\NormalTok{, }\AttributeTok{alpha =} \FloatTok{0.7}\NormalTok{) }\SpecialCharTok{+}
    \FunctionTok{labs}\NormalTok{(}\AttributeTok{title =} \StringTok{"NPE{-}vs{-}Stan Gap Across Different Datasets"}\NormalTok{,}
         \AttributeTok{subtitle =} \StringTok{"Large gap = converged rule NOT in Model B\textquotesingle{}s set. Small gap = converged rule IS in set."}\NormalTok{,}
         \AttributeTok{x =} \StringTok{"Schedule seed"}\NormalTok{, }\AttributeTok{y =} \StringTok{"|NPE mean ε − Stan mean ε|"}\NormalTok{)}
\NormalTok{\}}
\end{Highlighting}
\end{Shaded}

\textbf{Reading the gap demonstration.} If the gap varies across seeds
--- small for some, large for others --- the structural bias is driven
by rule-set mismatch, not inference noise. This is the expected outcome
and the chapter's main empirical contribution.

\subsection{Amortization Payoff for Chapter
15}\label{amortization-payoff-for-chapter-15}

Because the NPE network is trained once, inference on any new subject's
data takes milliseconds. For the Chapter 15 comparison across all three
cognitive theories (GCM, Kalman prototype, rule particle filter), each
model gets one trained NPE network, and the full multi-subject
comparison runs in seconds. The code pattern:

\begin{Shaded}
\begin{Highlighting}[]
\CommentTok{\# After training, inference on N subjects is trivial:}
\CommentTok{\# for each subject s:}
\CommentTok{\#   x\_s \textless{}{-} build\_data\_tensor(subject\_responses\_s, subject\_schedule\_s)}
\CommentTok{\#   posterior\_s \textless{}{-} npe\_posterior$sample(10000L, x = x\_s)}
\CommentTok{\# Total time: \textasciitilde{}N * 0.1 seconds}
\end{Highlighting}
\end{Shaded}

This is the methodological payoff of using NPE instead of ABC: the
Chapter 15 comparison becomes both honest (fitting the cognitive theory
directly) and actually executable (amortized inference scales to many
subjects).

\begin{center}\rule{0.5\linewidth}{0.5pt}\end{center}

\section{Extending the Particle Filter with Rule
Mutation}\label{extending-the-particle-filter-with-rule-mutation}

The particle filter as defined so far has no mechanism for entertaining
a rule that was not in the initial candidate set. Resampling purges
low-weight particles but cannot \emph{create} new rule variants, so the
filter's effective hypothesis space shrinks monotonically across trials.
On a stationary task like Kruschke (1993) this is fine --- the filter
converges on whichever small subset of initial particles happens to
match the data. But on a task where the ``correct'' rule changes over
time, the filter has no way to recover: all the rules that would fit the
new regime were likely eliminated in the first ten or twenty trials.

The fix is a \textbf{rule-mutation step}: at each trial, each particle
independently regenerates with probability \(\mu\), redrawn from the
rule grammar. This is the particle filter analogue of the process noise
\(q\) in the Kalman prototype (Ch. 12) and the recency weight
\(\lambda\) in the decay GCM (Ch. 11) --- a single parameter that
injects hypothesis turnover.

The three chapters now share a common shape:

\begin{longtable}[]{@{}
  >{\raggedright\arraybackslash}p{(\linewidth - 6\tabcolsep) * \real{0.2500}}
  >{\raggedright\arraybackslash}p{(\linewidth - 6\tabcolsep) * \real{0.2500}}
  >{\raggedright\arraybackslash}p{(\linewidth - 6\tabcolsep) * \real{0.2500}}
  >{\raggedright\arraybackslash}p{(\linewidth - 6\tabcolsep) * \real{0.2500}}@{}}
\toprule\noalign{}
\begin{minipage}[b]{\linewidth}\raggedright
Chapter
\end{minipage} & \begin{minipage}[b]{\linewidth}\raggedright
Model
\end{minipage} & \begin{minipage}[b]{\linewidth}\raggedright
Non-stationarity parameter
\end{minipage} & \begin{minipage}[b]{\linewidth}\raggedright
Mechanism
\end{minipage} \\
\midrule\noalign{}
\endhead
\bottomrule\noalign{}
\endlastfoot
Ch. 11 (decay GCM) & Exemplar store & \(\lambda\) (forgetting rate) &
Exponentially downweight old exemplars \\
Ch. 12 (Kalman prototype) & Running mean + covariance & \(q\) (process
noise) & Re-inflate uncertainty each trial \\
Ch. 13 (rule filter, below) & Particle distribution over rules & \(\mu\)
(mutation rate) & Randomly regenerate a fraction of particles each
trial \\
\end{longtable}

\subsection{Mutation Mechanism}\label{mutation-mechanism}

\begin{Shaded}
\begin{Highlighting}[]
\CommentTok{\# At each trial, each particle independently has probability mutation\_prob}
\CommentTok{\# of being replaced by a freshly drawn rule. This keeps the hypothesis space}
\CommentTok{\# from collapsing and allows the filter to track changes in the correct rule.}
\NormalTok{mutate\_particles }\OtherTok{\textless{}{-}} \ControlFlowTok{function}\NormalTok{(particles, mutation\_prob, n\_features,}
\NormalTok{                             max\_dims, feature\_range) \{}
  \ControlFlowTok{if}\NormalTok{ (mutation\_prob }\SpecialCharTok{\textless{}=} \DecValTok{0}\NormalTok{) }\FunctionTok{return}\NormalTok{(particles)}
\NormalTok{  n }\OtherTok{\textless{}{-}} \FunctionTok{length}\NormalTok{(particles)}
\NormalTok{  to\_mutate }\OtherTok{\textless{}{-}} \FunctionTok{which}\NormalTok{(}\FunctionTok{runif}\NormalTok{(n) }\SpecialCharTok{\textless{}}\NormalTok{ mutation\_prob)}
  \ControlFlowTok{for}\NormalTok{ (i }\ControlFlowTok{in}\NormalTok{ to\_mutate) \{}
\NormalTok{    particles[[i]] }\OtherTok{\textless{}{-}} \FunctionTok{generate\_random\_rule}\NormalTok{(n\_features, max\_dims, feature\_range)}
\NormalTok{  \}}
\NormalTok{  particles}
\NormalTok{\}}
\end{Highlighting}
\end{Shaded}

\subsection{Scenario-Aware Particle
Filter}\label{scenario-aware-particle-filter}

We extend \texttt{rule\_particle\_filter} to take a
\texttt{mutation\_prob} argument and a \texttt{scenario} switch that
controls how the true feedback sequence is produced. The scenarios
follow the canonical framework from Chapter 11: static Kruschke,
performance-contingent abrupt shifts, and smooth drift (in that order).
Note that contingent-shift flips are endogenous to the agent's own
accuracy --- a noisier agent triggers fewer flips and accumulates less
information about \(\mu\) (see Chapter 11 for discussion).

\begin{Shaded}
\begin{Highlighting}[]
\NormalTok{make\_drift\_trajectory }\OtherTok{\textless{}{-}} \ControlFlowTok{function}\NormalTok{(ntrials, }\AttributeTok{drift\_sigma =} \FloatTok{0.05}\NormalTok{, }\AttributeTok{seed =} \DecValTok{1}\NormalTok{) \{}
  \FunctionTok{set.seed}\NormalTok{(seed)}
\NormalTok{  mu0\_init }\OtherTok{\textless{}{-}} \FunctionTok{c}\NormalTok{(}\FloatTok{1.75}\NormalTok{, }\FloatTok{3.25}\NormalTok{)}
\NormalTok{  mu1\_init }\OtherTok{\textless{}{-}} \FunctionTok{c}\NormalTok{(}\FloatTok{3.25}\NormalTok{, }\FloatTok{1.75}\NormalTok{)}
\NormalTok{  d0 }\OtherTok{\textless{}{-}} \FunctionTok{apply}\NormalTok{(}\FunctionTok{matrix}\NormalTok{(}\FunctionTok{rnorm}\NormalTok{(ntrials }\SpecialCharTok{*} \DecValTok{2}\NormalTok{, }\DecValTok{0}\NormalTok{, drift\_sigma), }\AttributeTok{ncol =} \DecValTok{2}\NormalTok{), }\DecValTok{2}\NormalTok{, cumsum)}
\NormalTok{  d1 }\OtherTok{\textless{}{-}} \FunctionTok{apply}\NormalTok{(}\FunctionTok{matrix}\NormalTok{(}\FunctionTok{rnorm}\NormalTok{(ntrials }\SpecialCharTok{*} \DecValTok{2}\NormalTok{, }\DecValTok{0}\NormalTok{, drift\_sigma), }\AttributeTok{ncol =} \DecValTok{2}\NormalTok{), }\DecValTok{2}\NormalTok{, cumsum)}
  \FunctionTok{list}\NormalTok{(}
    \AttributeTok{mu0 =} \FunctionTok{sweep}\NormalTok{(d0, }\DecValTok{2}\NormalTok{, mu0\_init, }\StringTok{"+"}\NormalTok{),}
    \AttributeTok{mu1 =} \FunctionTok{sweep}\NormalTok{(d1, }\DecValTok{2}\NormalTok{, mu1\_init, }\StringTok{"+"}\NormalTok{)}
\NormalTok{  )}
\NormalTok{\}}
\end{Highlighting}
\end{Shaded}

\begin{Shaded}
\begin{Highlighting}[]
\NormalTok{rule\_particle\_filter\_scenario }\OtherTok{\textless{}{-}} \ControlFlowTok{function}\NormalTok{(schedule,}
\NormalTok{                                          error\_prob,}
                                          \AttributeTok{mutation\_prob              =} \FloatTok{0.0}\NormalTok{,}
                                          \AttributeTok{n\_particles                =}\NormalTok{ N\_PARTICLES\_FIXED,}
                                          \AttributeTok{max\_dims                   =}\NormalTok{ MAX\_RULE\_DIMS,}
                                          \AttributeTok{resample\_threshold\_factor  =}\NormalTok{ RESAMPLE\_THRESHOLD,}
                                          \AttributeTok{scenario                   =} \FunctionTok{c}\NormalTok{(}\StringTok{"static"}\NormalTok{,}
                                                                          \StringTok{"contingent\_shift"}\NormalTok{,}
                                                                          \StringTok{"drift"}\NormalTok{),}
                                          \AttributeTok{streak\_target              =} \DecValTok{8}\NormalTok{,}
                                          \AttributeTok{drift\_traj                 =} \ConstantTok{NULL}\NormalTok{,}
                                          \AttributeTok{quiet                      =} \ConstantTok{TRUE}\NormalTok{) \{}
\NormalTok{  scenario   }\OtherTok{\textless{}{-}} \FunctionTok{match.arg}\NormalTok{(scenario)}
\NormalTok{  n\_trials   }\OtherTok{\textless{}{-}} \FunctionTok{nrow}\NormalTok{(schedule)}
\NormalTok{  obs        }\OtherTok{\textless{}{-}} \FunctionTok{as.matrix}\NormalTok{(schedule[, }\FunctionTok{c}\NormalTok{(}\StringTok{"height"}\NormalTok{, }\StringTok{"position"}\NormalTok{)])}
\NormalTok{  base\_lab   }\OtherTok{\textless{}{-}}\NormalTok{ schedule}\SpecialCharTok{$}\NormalTok{category\_feedback}
\NormalTok{  n\_features }\OtherTok{\textless{}{-}} \FunctionTok{ncol}\NormalTok{(obs)}
\NormalTok{  fr         }\OtherTok{\textless{}{-}} \FunctionTok{list}\NormalTok{(}\FunctionTok{range}\NormalTok{(obs[, }\DecValTok{1}\NormalTok{]), }\FunctionTok{range}\NormalTok{(obs[, }\DecValTok{2}\NormalTok{]))}

\NormalTok{  ps        }\OtherTok{\textless{}{-}} \FunctionTok{initialize\_particles}\NormalTok{(n\_particles, n\_features, max\_dims, fr)}
\NormalTok{  particles }\OtherTok{\textless{}{-}}\NormalTok{ ps}\SpecialCharTok{$}\NormalTok{particles}
\NormalTok{  weights   }\OtherTok{\textless{}{-}}\NormalTok{ ps}\SpecialCharTok{$}\NormalTok{weights}

\NormalTok{  label\_flip         }\OtherTok{\textless{}{-}} \ConstantTok{FALSE}
\NormalTok{  streak             }\OtherTok{\textless{}{-}} \DecValTok{0}\NormalTok{L}

\NormalTok{  response\_probs     }\OtherTok{\textless{}{-}} \FunctionTok{numeric}\NormalTok{(n\_trials)}
\NormalTok{  sim\_response       }\OtherTok{\textless{}{-}} \FunctionTok{integer}\NormalTok{(n\_trials)}
\NormalTok{  observed\_feedback  }\OtherTok{\textless{}{-}} \FunctionTok{integer}\NormalTok{(n\_trials)}
\NormalTok{  label\_flip\_trace   }\OtherTok{\textless{}{-}} \FunctionTok{logical}\NormalTok{(n\_trials)}

  \ControlFlowTok{for}\NormalTok{ (i }\ControlFlowTok{in} \FunctionTok{seq\_len}\NormalTok{(n\_trials)) \{}
\NormalTok{    stim }\OtherTok{\textless{}{-}} \FunctionTok{as.numeric}\NormalTok{(obs[i, ])}

    \CommentTok{\# ── Mutation step ─────────────────────────────────────────────────}
\NormalTok{    particles }\OtherTok{\textless{}{-}} \FunctionTok{mutate\_particles}\NormalTok{(particles, mutation\_prob,}
\NormalTok{                                  n\_features, max\_dims, fr)}

    \CommentTok{\# ── Prediction ────────────────────────────────────────────────────}
\NormalTok{    preds        }\OtherTok{\textless{}{-}} \FunctionTok{sapply}\NormalTok{(particles, evaluate\_rule, }\AttributeTok{stimulus =}\NormalTok{ stim)}
\NormalTok{    p\_cat1\_rules }\OtherTok{\textless{}{-}} \FunctionTok{ifelse}\NormalTok{(preds }\SpecialCharTok{==} \DecValTok{1}\NormalTok{, }\DecValTok{1} \SpecialCharTok{{-}}\NormalTok{ error\_prob, error\_prob)}
\NormalTok{    p            }\OtherTok{\textless{}{-}} \FunctionTok{sum}\NormalTok{(weights }\SpecialCharTok{*}\NormalTok{ p\_cat1\_rules)}
\NormalTok{    p            }\OtherTok{\textless{}{-}} \FunctionTok{pmax}\NormalTok{(}\FloatTok{1e{-}9}\NormalTok{, }\FunctionTok{pmin}\NormalTok{(}\DecValTok{1} \SpecialCharTok{{-}} \FloatTok{1e{-}9}\NormalTok{, p))}
\NormalTok{    response\_probs[i] }\OtherTok{\textless{}{-}}\NormalTok{ p}
\NormalTok{    sim\_response[i]   }\OtherTok{\textless{}{-}} \FunctionTok{rbinom}\NormalTok{(}\DecValTok{1}\NormalTok{, }\DecValTok{1}\NormalTok{, p)}

    \CommentTok{\# ── Observed feedback depends on scenario ─────────────────────────}
\NormalTok{    fb }\OtherTok{\textless{}{-}} \ControlFlowTok{switch}\NormalTok{(scenario,}
      \AttributeTok{static           =}\NormalTok{ base\_lab[i],}
      \AttributeTok{contingent\_shift =} \ControlFlowTok{if}\NormalTok{ (label\_flip) }\DecValTok{1}\NormalTok{L }\SpecialCharTok{{-}}\NormalTok{ base\_lab[i] }\ControlFlowTok{else}\NormalTok{ base\_lab[i],}
      \AttributeTok{drift            =}\NormalTok{ \{}
\NormalTok{        d0 }\OtherTok{\textless{}{-}} \FunctionTok{sum}\NormalTok{(}\FunctionTok{abs}\NormalTok{(stim }\SpecialCharTok{{-}}\NormalTok{ drift\_traj}\SpecialCharTok{$}\NormalTok{mu0[i, ]))}
\NormalTok{        d1 }\OtherTok{\textless{}{-}} \FunctionTok{sum}\NormalTok{(}\FunctionTok{abs}\NormalTok{(stim }\SpecialCharTok{{-}}\NormalTok{ drift\_traj}\SpecialCharTok{$}\NormalTok{mu1[i, ]))}
        \FunctionTok{as.integer}\NormalTok{(d1 }\SpecialCharTok{\textless{}}\NormalTok{ d0)}
\NormalTok{      \}}
\NormalTok{    )}
\NormalTok{    observed\_feedback[i] }\OtherTok{\textless{}{-}}\NormalTok{ fb}
\NormalTok{    label\_flip\_trace[i]  }\OtherTok{\textless{}{-}}\NormalTok{ label\_flip}

    \ControlFlowTok{if}\NormalTok{ (scenario }\SpecialCharTok{==} \StringTok{"contingent\_shift"}\NormalTok{) \{}
\NormalTok{      streak }\OtherTok{\textless{}{-}} \ControlFlowTok{if}\NormalTok{ (sim\_response[i] }\SpecialCharTok{==}\NormalTok{ fb) streak }\SpecialCharTok{+} \DecValTok{1}\NormalTok{L }\ControlFlowTok{else} \DecValTok{0}\NormalTok{L}
      \ControlFlowTok{if}\NormalTok{ (streak }\SpecialCharTok{\textgreater{}=}\NormalTok{ streak\_target) \{ label\_flip }\OtherTok{\textless{}{-}} \SpecialCharTok{!}\NormalTok{label\_flip; streak }\OtherTok{\textless{}{-}} \DecValTok{0}\NormalTok{L \}}
\NormalTok{    \}}

    \CommentTok{\# ── Weight update ────────────────────────────────────────────────}
\NormalTok{    weights }\OtherTok{\textless{}{-}} \FunctionTok{update\_weights}\NormalTok{(particles, weights, stim, fb, error\_prob)}

    \CommentTok{\# ── Resample if ESS is low ───────────────────────────────────────}
\NormalTok{    ess }\OtherTok{\textless{}{-}} \DecValTok{1} \SpecialCharTok{/} \FunctionTok{sum}\NormalTok{(weights}\SpecialCharTok{\^{}}\DecValTok{2}\NormalTok{)}
    \ControlFlowTok{if}\NormalTok{ (ess }\SpecialCharTok{\textless{}}\NormalTok{ n\_particles }\SpecialCharTok{*}\NormalTok{ resample\_threshold\_factor) \{}
\NormalTok{      rs        }\OtherTok{\textless{}{-}} \FunctionTok{resample\_particles}\NormalTok{(particles, weights)}
\NormalTok{      particles }\OtherTok{\textless{}{-}}\NormalTok{ rs}\SpecialCharTok{$}\NormalTok{particles}
\NormalTok{      weights   }\OtherTok{\textless{}{-}}\NormalTok{ rs}\SpecialCharTok{$}\NormalTok{weights}
\NormalTok{    \}}
\NormalTok{  \}}

\NormalTok{  schedule }\SpecialCharTok{|\textgreater{}}
    \FunctionTok{mutate}\NormalTok{(}
      \AttributeTok{prob\_cat1         =}\NormalTok{ response\_probs,}
      \AttributeTok{sim\_response      =}\NormalTok{ sim\_response,}
      \AttributeTok{observed\_feedback =}\NormalTok{ observed\_feedback,}
      \AttributeTok{label\_flip        =}\NormalTok{ label\_flip\_trace,}
      \AttributeTok{correct           =} \FunctionTok{as.integer}\NormalTok{(sim\_response }\SpecialCharTok{==}\NormalTok{ observed\_feedback)}
\NormalTok{    )}
\NormalTok{\}}
\end{Highlighting}
\end{Shaded}

\subsection{Agent Wrapper}\label{agent-wrapper}

\begin{Shaded}
\begin{Highlighting}[]
\NormalTok{simulate\_rule\_agent\_scenario }\OtherTok{\textless{}{-}} \ControlFlowTok{function}\NormalTok{(agent\_id, scenario,}
\NormalTok{                                         error\_prob, mutation\_prob,}
\NormalTok{                                         subject\_seed,}
                                         \AttributeTok{n\_particles   =}\NormalTok{ N\_PARTICLES\_FIXED,}
                                         \AttributeTok{max\_dims      =}\NormalTok{ MAX\_RULE\_DIMS,}
                                         \AttributeTok{streak\_target =} \DecValTok{8}\NormalTok{,}
                                         \AttributeTok{drift\_sigma   =} \FloatTok{0.05}\NormalTok{) \{}
\NormalTok{  schedule }\OtherTok{\textless{}{-}} \FunctionTok{make\_subject\_schedule}\NormalTok{(stimulus\_info, n\_blocks, }\AttributeTok{seed =}\NormalTok{ subject\_seed)}
\NormalTok{  drift\_tr }\OtherTok{\textless{}{-}} \ControlFlowTok{if}\NormalTok{ (scenario }\SpecialCharTok{==} \StringTok{"drift"}\NormalTok{)}
    \FunctionTok{make\_drift\_trajectory}\NormalTok{(}\FunctionTok{nrow}\NormalTok{(schedule), drift\_sigma, }\AttributeTok{seed =}\NormalTok{ subject\_seed)}
  \ControlFlowTok{else} \ConstantTok{NULL}

\NormalTok{  out }\OtherTok{\textless{}{-}} \FunctionTok{rule\_particle\_filter\_scenario}\NormalTok{(}
    \AttributeTok{schedule                  =}\NormalTok{ schedule,}
    \AttributeTok{error\_prob                =}\NormalTok{ error\_prob,}
    \AttributeTok{mutation\_prob             =}\NormalTok{ mutation\_prob,}
    \AttributeTok{n\_particles               =}\NormalTok{ n\_particles,}
    \AttributeTok{max\_dims                  =}\NormalTok{ max\_dims,}
    \AttributeTok{scenario                  =}\NormalTok{ scenario,}
    \AttributeTok{streak\_target             =}\NormalTok{ streak\_target,}
    \AttributeTok{drift\_traj                =}\NormalTok{ drift\_tr,}
    \AttributeTok{resample\_threshold\_factor =}\NormalTok{ RESAMPLE\_THRESHOLD}
\NormalTok{  )}

\NormalTok{  out }\SpecialCharTok{|\textgreater{}}
    \FunctionTok{mutate}\NormalTok{(}
      \AttributeTok{agent\_id           =}\NormalTok{ agent\_id,}
      \AttributeTok{scenario           =}\NormalTok{ scenario,}
      \AttributeTok{error\_prob\_true    =}\NormalTok{ error\_prob,}
      \AttributeTok{mutation\_prob\_true =}\NormalTok{ mutation\_prob}
\NormalTok{    ) }\SpecialCharTok{|\textgreater{}}
    \FunctionTok{group\_by}\NormalTok{(agent\_id) }\SpecialCharTok{|\textgreater{}}
    \FunctionTok{mutate}\NormalTok{(}\AttributeTok{cumulative\_accuracy =} \FunctionTok{cumsum}\NormalTok{(correct) }\SpecialCharTok{/} \FunctionTok{row\_number}\NormalTok{()) }\SpecialCharTok{|\textgreater{}}
    \FunctionTok{ungroup}\NormalTok{()}
\NormalTok{\}}
\end{Highlighting}
\end{Shaded}

\begin{tcolorbox}[enhanced jigsaw, toprule=.15mm, toptitle=1mm, bottomtitle=1mm, leftrule=.75mm, rightrule=.15mm, colframe=quarto-callout-note-color-frame, left=2mm, colback=white, titlerule=0mm, coltitle=black, colbacktitle=quarto-callout-note-color!10!white, opacityback=0, title=\textcolor{quarto-callout-note-color}{\faInfo}\hspace{0.5em}{Unified scenario framework across the three chapters}, arc=.35mm, opacitybacktitle=0.6, breakable, bottomrule=.15mm]

All three chapters use the same three scenarios in the same order
(static → contingent shift → drift), reflecting \textbf{increasing
assumption match} between the scenario's data-generating process and the
model's structural commitments. The canonical definition and rationale
are in Chapter 11. What differs across chapters is \emph{how well} the
assumption-matched scenario actually matches:

\begin{longtable}[]{@{}
  >{\raggedright\arraybackslash}p{(\linewidth - 6\tabcolsep) * \real{0.2500}}
  >{\raggedright\arraybackslash}p{(\linewidth - 6\tabcolsep) * \real{0.2500}}
  >{\raggedright\arraybackslash}p{(\linewidth - 6\tabcolsep) * \real{0.2500}}
  >{\raggedright\arraybackslash}p{(\linewidth - 6\tabcolsep) * \real{0.2500}}@{}}
\toprule\noalign{}
\begin{minipage}[b]{\linewidth}\raggedright
Chapter
\end{minipage} & \begin{minipage}[b]{\linewidth}\raggedright
Model
\end{minipage} & \begin{minipage}[b]{\linewidth}\raggedright
Contingent shift
\end{minipage} & \begin{minipage}[b]{\linewidth}\raggedright
Smooth drift
\end{minipage} \\
\midrule\noalign{}
\endhead
\bottomrule\noalign{}
\endlastfoot
Ch. 11 --- Decay GCM & Exemplar memory with forgetting & Misspecified
(no abrupt mechanism); strong \(\lambda\) signal via post-flip recovery
& Graceful degradation; approximate match \\
Ch. 12 --- Kalman prototype & Gaussian-drift prototype tracker &
Misspecified; filter smooths the reversal & \textbf{Exact match}:
Gaussian drift is the model's native story \\
Ch. 13 --- Rule particle filter & Bayesian particle filter + mutation &
Approximate; mutation lets the filter recover & Approximate; mutation
supplies fresh candidates but does not track direction \\
\end{longtable}

Readers should compare the learning-curve panels across chapters to see
how each model's structural commitments show up in its sensitivity to
the two non-stationary scenarios. The Kalman prototype is uniquely
well-matched to smooth drift; the decay GCM and the rule filter are both
general-purpose in ways that make them approximately suitable for either
non-stationary scenario, for different mechanistic reasons.

\end{tcolorbox}

\subsection{Visualising the Three
Scenarios}\label{visualising-the-three-scenarios-2}

We simulate 8 agents per cell across a 3 × 3 grid of scenarios ×
mutation rates, at fixed \(\varepsilon = 0.1\) and
\(N_{\text{particles}} = 100\). As in Chapters 11 and 12, we expect
mutation to \emph{hurt} on static data (it purges rules that already
match) and \emph{help} on non-stationary data (it injects fresh
candidates the filter can reweight when the environment changes).

\begin{Shaded}
\begin{Highlighting}[]
\NormalTok{rule\_scenario\_viz\_file }\OtherTok{\textless{}{-}} \FunctionTok{here}\NormalTok{(}\StringTok{"simdata"}\NormalTok{, }\StringTok{"ch14\_rule\_scenario\_viz.csv"}\NormalTok{)}

\ControlFlowTok{if}\NormalTok{ (regenerate\_simulations }\SpecialCharTok{||} \SpecialCharTok{!}\FunctionTok{file.exists}\NormalTok{(rule\_scenario\_viz\_file)) \{}
\NormalTok{  viz\_grid }\OtherTok{\textless{}{-}} \FunctionTok{expand\_grid}\NormalTok{(}
    \AttributeTok{scenario      =} \FunctionTok{c}\NormalTok{(}\StringTok{"static"}\NormalTok{, }\StringTok{"contingent\_shift"}\NormalTok{, }\StringTok{"drift"}\NormalTok{),}
    \AttributeTok{mutation\_prob =} \FunctionTok{c}\NormalTok{(}\FloatTok{0.0}\NormalTok{, }\FloatTok{0.05}\NormalTok{, }\FloatTok{0.2}\NormalTok{),}
    \AttributeTok{agent\_id      =} \DecValTok{1}\SpecialCharTok{:}\DecValTok{8}
\NormalTok{  ) }\SpecialCharTok{|\textgreater{}}
    \FunctionTok{mutate}\NormalTok{(}\AttributeTok{subject\_seed =} \DecValTok{50000} \SpecialCharTok{+} \FunctionTok{row\_number}\NormalTok{())}

\NormalTok{  rule\_viz\_sim }\OtherTok{\textless{}{-}} \FunctionTok{pmap\_dfr}\NormalTok{(}
    \FunctionTok{list}\NormalTok{(}\AttributeTok{agent\_id      =}\NormalTok{ viz\_grid}\SpecialCharTok{$}\NormalTok{agent\_id,}
         \AttributeTok{scenario      =}\NormalTok{ viz\_grid}\SpecialCharTok{$}\NormalTok{scenario,}
         \AttributeTok{mutation\_prob =}\NormalTok{ viz\_grid}\SpecialCharTok{$}\NormalTok{mutation\_prob,}
         \AttributeTok{subject\_seed  =}\NormalTok{ viz\_grid}\SpecialCharTok{$}\NormalTok{subject\_seed),}
    \ControlFlowTok{function}\NormalTok{(agent\_id, scenario, mutation\_prob, subject\_seed) \{}
      \FunctionTok{simulate\_rule\_agent\_scenario}\NormalTok{(}
        \AttributeTok{agent\_id      =}\NormalTok{ agent\_id,}
        \AttributeTok{scenario      =}\NormalTok{ scenario,}
        \AttributeTok{error\_prob    =} \FloatTok{0.1}\NormalTok{,}
        \AttributeTok{mutation\_prob =}\NormalTok{ mutation\_prob,}
        \AttributeTok{subject\_seed  =}\NormalTok{ subject\_seed}
\NormalTok{      )}
\NormalTok{    \}}
\NormalTok{  )}
  \FunctionTok{write\_csv}\NormalTok{(rule\_viz\_sim, rule\_scenario\_viz\_file)}
\NormalTok{\} }\ControlFlowTok{else}\NormalTok{ \{}
\NormalTok{  rule\_viz\_sim }\OtherTok{\textless{}{-}} \FunctionTok{read\_csv}\NormalTok{(rule\_scenario\_viz\_file, }\AttributeTok{show\_col\_types =} \ConstantTok{FALSE}\NormalTok{)}
\NormalTok{\}}

\FunctionTok{ggplot}\NormalTok{(rule\_viz\_sim,}
       \FunctionTok{aes}\NormalTok{(}\AttributeTok{x =}\NormalTok{ trial\_within\_subject, }\AttributeTok{y =}\NormalTok{ cumulative\_accuracy,}
           \AttributeTok{color =} \FunctionTok{factor}\NormalTok{(mutation\_prob\_true))) }\SpecialCharTok{+}
  \FunctionTok{stat\_summary}\NormalTok{(}\AttributeTok{fun =}\NormalTok{ mean, }\AttributeTok{geom =} \StringTok{"line"}\NormalTok{,}
               \FunctionTok{aes}\NormalTok{(}\AttributeTok{group =} \FunctionTok{factor}\NormalTok{(mutation\_prob\_true)), }\AttributeTok{linewidth =} \FloatTok{1.1}\NormalTok{) }\SpecialCharTok{+}
  \FunctionTok{stat\_summary}\NormalTok{(}\AttributeTok{fun.data =}\NormalTok{ mean\_se, }\AttributeTok{geom =} \StringTok{"ribbon"}\NormalTok{, }\AttributeTok{alpha =} \FloatTok{0.15}\NormalTok{,}
               \FunctionTok{aes}\NormalTok{(}\AttributeTok{fill =} \FunctionTok{factor}\NormalTok{(mutation\_prob\_true),}
                   \AttributeTok{group =} \FunctionTok{factor}\NormalTok{(mutation\_prob\_true)),}
               \AttributeTok{color =} \ConstantTok{NA}\NormalTok{) }\SpecialCharTok{+}
  \FunctionTok{facet\_wrap}\NormalTok{(}\SpecialCharTok{\textasciitilde{}}\NormalTok{scenario, }\AttributeTok{ncol =} \DecValTok{3}\NormalTok{,}
             \AttributeTok{labeller =} \FunctionTok{as\_labeller}\NormalTok{(}\FunctionTok{c}\NormalTok{(}
               \AttributeTok{static           =} \StringTok{"Static Kruschke"}\NormalTok{,}
               \AttributeTok{contingent\_shift =} \StringTok{"Performance{-}contingent flips"}\NormalTok{,}
               \AttributeTok{drift            =} \StringTok{"Smooth drift"}\NormalTok{))) }\SpecialCharTok{+}
  \FunctionTok{scale\_color\_viridis\_d}\NormalTok{(}\AttributeTok{option =} \StringTok{"plasma"}\NormalTok{, }\AttributeTok{name =} \FunctionTok{expression}\NormalTok{(mu)) }\SpecialCharTok{+}
  \FunctionTok{scale\_fill\_viridis\_d}\NormalTok{(}\AttributeTok{option  =} \StringTok{"plasma"}\NormalTok{, }\AttributeTok{name =} \FunctionTok{expression}\NormalTok{(mu)) }\SpecialCharTok{+}
  \FunctionTok{geom\_hline}\NormalTok{(}\AttributeTok{yintercept =} \FloatTok{0.5}\NormalTok{, }\AttributeTok{linetype =} \StringTok{"dashed"}\NormalTok{, }\AttributeTok{color =} \StringTok{"grey50"}\NormalTok{) }\SpecialCharTok{+}
  \FunctionTok{labs}\NormalTok{(}
    \AttributeTok{title    =} \StringTok{"Rule Particle Filter Learning Curves Across Scenarios"}\NormalTok{,}
    \AttributeTok{subtitle =} \FunctionTok{paste0}\NormalTok{(}\StringTok{"Mean ± SE across 8 agents per cell; ε = 0.1, N\_particles = 100"}\NormalTok{),}
    \AttributeTok{x =} \StringTok{"Trial"}\NormalTok{, }\AttributeTok{y =} \StringTok{"Cumulative accuracy"}
\NormalTok{  ) }\SpecialCharTok{+}
  \FunctionTok{theme}\NormalTok{(}\AttributeTok{legend.position =} \StringTok{"bottom"}\NormalTok{)}
\end{Highlighting}
\end{Shaded}

\begin{center}
\includegraphics[width=0.8\linewidth,height=\textheight,keepaspectratio]{13-CategorizationModelsRules_files/figure-pdf/ch14_scenario_viz-1.pdf}
\end{center}

\textbf{Reading the panels.} On static data, \(\mu = 0\) wins: the
filter can converge on the right rule and keep it. Higher \(\mu\) trades
accuracy for adaptability --- with 5\% mutation per trial per particle,
there is always \textasciitilde5\% of the particle cloud freshly drawn
from the prior, which cannot (yet) have high weight. On contingent-shift
data, the curves sort in the opposite direction after the first flip:
\(\mu = 0\) collapses after the reversal and never recovers, while
\(\mu > 0\) produces the characteristic sawtooth recovery. On drift
data, moderate \(\mu\) tracks the drifting boundary better than
\(\mu = 0\).

\subsection{Assumption Match Across
Scenarios}\label{assumption-match-across-scenarios-2}

\begin{longtable}[]{@{}
  >{\raggedright\arraybackslash}p{(\linewidth - 6\tabcolsep) * \real{0.2500}}
  >{\raggedright\arraybackslash}p{(\linewidth - 6\tabcolsep) * \real{0.2500}}
  >{\raggedright\arraybackslash}p{(\linewidth - 6\tabcolsep) * \real{0.2500}}
  >{\raggedright\arraybackslash}p{(\linewidth - 6\tabcolsep) * \real{0.2500}}@{}}
\toprule\noalign{}
\begin{minipage}[b]{\linewidth}\raggedright
Scenario
\end{minipage} & \begin{minipage}[b]{\linewidth}\raggedright
What ``changes'' in the world
\end{minipage} & \begin{minipage}[b]{\linewidth}\raggedright
What mutation assumes
\end{minipage} & \begin{minipage}[b]{\linewidth}\raggedright
Match
\end{minipage} \\
\midrule\noalign{}
\endhead
\bottomrule\noalign{}
\endlastfoot
Static & Nothing & Stochastic rule turnover at rate \(\mu\) &
\textbf{Harmless at \(\mu = 0\)}; costly for \(\mu > 0\) \\
Contingent shift & Abrupt label reversals & Stochastic rule turnover &
\textbf{Approximate}: mutation lets the filter recover, but there is no
mechanism for abrupt, synchronised re-initialisation \\
Smooth drift & Gaussian drift of category centres & Stochastic rule
turnover & \textbf{Approximate}: mutation supplies fresh candidates but
is not literally tracking the drift direction \\
\end{longtable}

Unlike the Kalman prototype's \(q\), mutation is not a \emph{model} of
the drift process --- it is a coarse rejuvenation mechanism that keeps
the hypothesis space alive so the filter can react to whatever changes
the environment throws at it. The matching story in Chapter 12 was
``smooth drift fits the model's generative assumption exactly''; here it
is ``mutation is a general-purpose fix that happens to suffice in both
non-stationary scenarios.'' This is the cognitive-theory contrast:
rule-based learners have no commitment to \emph{how} the environment
changes, only that their hypotheses need to stay fresh.

\subsection{How This Strengthens the Model A vs Model B
Gap}\label{how-this-strengthens-the-model-a-vs-model-b-gap}

The Stan mixture (Model B) is built on a fixed candidate rule set
enumerated before inference begins. There is no analogue to \(\mu\) ---
mixture weights can reassign posterior mass across the candidate set,
but nothing in Model B \emph{creates} rules that were not in the
candidate set. On static data this is an acceptable approximation as
long as the true rule is in the candidate set (Ch. 13's §``Discharging
the Post-Hoc Rule Selection'' argues exactly this). On contingent-shift
or drift data, the candidate set that correctly covered the initial
environment will not cover the post-shift environment, and Model B has
no mechanism to repair itself.

Concretely: if we take the contingent-shift simulated data from the plot
above and fit Model B to it, Model B's posterior predictive learning
curve will be stuck at the pre-shift accuracy plateau --- it has no way
to produce the sawtooth. The same fit under NPE (with \(\varepsilon\)
and \(\mu\) both inferred) can match the data because Model A \emph{is}
the generative process. \textbf{The gap between Model A and Model B on
non-stationary data is therefore not a subtle bias --- it is a
structural ceiling.} This is the sharpened version of the argument in
§``The Main Empirical Result'': the cost of replacing an intractable
cognitive theory with a tractable proxy is bounded when the task is
stationary, and unbounded when it is not.

\subsection{Implication for the NPE
Pipeline}\label{implication-for-the-npe-pipeline}

The NPE setup in §``Neural Posterior Estimation for Model A'' infers
\((\varepsilon, N_{\text{particles}})\). Extending to
\((\varepsilon, N_{\text{particles}}, \mu)\) is a one-line change to the
prior and a one-argument change to \texttt{simulate\_model\_a}.
Validation (SBC + expected coverage) and the pair-plot inspection
transfer directly. On static data we would expect \(\mu\) to be poorly
identified --- matching Ch. 12's \(q\) and Ch. 11's \(\lambda\) under
their own static scenarios. On contingent-shift and drift data \(\mu\)
should become identifiable, because the signature in the learning curve
(sawtooth depth, plateau height, post-flip recovery speed) depends
specifically on \(\mu\) in a way the other parameters cannot mimic.
Production implementations of this chapter for the Chapter 15
cross-architecture comparison should include \(\mu\) in the NPE prior
and run one validation sweep per scenario.

\subsection{Trilogy Synthesis}\label{trilogy-synthesis}

Chapters 11, 12, and 13 now each contain a parallel analysis: a
one-parameter extension to the base cognitive model that turns the model
from a stationary learner into a non-stationary one, validated across
three scenarios that load the new parameter differently. The parameters
are superficially different --- a recency weight, a process-noise
variance, a mutation probability --- but the validation story is the
same in each case: identifiability is a joint property of the model and
the task, the non-stationarity parameter is weakly identified when the
task is stationary, and SBC under a stationary task catches this failure
mode while SBC under a matched non-stationary task passes. This is the
shape the reader should recognise when extending any path-dependent
cognitive model: one parameter for non-stationarity, one matched
scenario to identify it, and a validation battery that sorts
assumption-matched, mis-specified, and information-starved cases in the
order the theory predicts.

\begin{center}\rule{0.5\linewidth}{0.5pt}\end{center}

\section{Cognitive Insights from Rule-Based
Models}\label{cognitive-insights-from-rule-based-models}

Rule-based models, particularly Model A (the particle filter), offer
several cognitive insights:

\begin{itemize}
\tightlist
\item
  \textbf{Explicit Hypothesis Testing}: They model the process of
  generating, testing, and refining explicit hypotheses about category
  structure.
\item
  \textbf{Abrupt Learning}: Performance can change sharply when a highly
  effective rule is discovered and gains high weight --- producing the
  ``aha moment'' learning curves sometimes observed in humans.
\item
  \textbf{Selective Attention}: Attention is modeled as selecting
  specific dimensions relevant to the current best rules, rather than
  continuous weighting across all dimensions.
\item
  \textbf{Verbalizability}: The explicit nature of rules often aligns
  with people's ability and tendency to verbalize their categorization
  strategies.
\item
  \textbf{Individual Differences}: Different individuals might converge
  on different (but potentially equally effective) rules, explaining
  variations in reported strategies.
\item
  \textbf{Working-memory bound}: \(N_{\text{particles}}\) has direct
  cognitive content as a hypothesis-set capacity. NPE makes this an
  inferable parameter rather than a structural choice --- this is a
  non-trivial methodological gain over the Stan proxy, which has no
  analogue.
\end{itemize}

\begin{center}\rule{0.5\linewidth}{0.5pt}\end{center}

\section{Strengths and Limitations}\label{strengths-and-limitations}

\textbf{Strengths:}

\begin{itemize}
\tightlist
\item
  Cognitive plausibility: aligns with introspection and verbal reports
  of using rules.
\item
  Efficiency for simple structures: very effective when categories are
  defined by simple, logical rules.
\item
  Sharp boundaries: naturally captures categories with clear-cut
  criteria.
\item
  Handles logic: easily incorporates AND, OR, and NOT operators within
  the chosen grammar.
\item
  Working-memory interpretation: \(N_{\text{particles}}\) has direct
  cognitive meaning as a hypothesis capacity bound, and is recoverable
  from data via NPE.
\item
  Honest about its scope: the Model A vs Model B distinction makes it
  possible to fit the cognitive theory directly (via NPE) while still
  having a tractable proxy (Model B) for downstream model comparison.
\item
  Amortized inference: once the NPE network is trained, inference on any
  new subject takes milliseconds, making multi-subject comparisons in
  Chapter 15 computationally trivial.
\end{itemize}

\textbf{Limitations:}

\begin{itemize}
\tightlist
\item
  \textbf{Family resemblance categories}: rule-based models struggle
  with categories defined by overall similarity where no single rule
  applies to all members.
\item
  \textbf{Complex boundaries}: requires very complex rules or many
  simple rules to approximate non-axis-aligned or curved boundaries.
\item
  \textbf{Graded typicality}: doesn't easily explain why some category
  members seem more typical than others.
\item
  \textbf{Feature learning}: standard models assume predefined features;
  they don't typically learn which features are relevant from raw input.
\item
  \textbf{Rule grammar is a structural prior} (Phase 0 Q4): the chapter
  commits to \(\{>, <\}\) operators, \(\{\text{AND}, \text{OR}\}\)
  logical connectives, flat rules up to \(\text{max\_dims} = 2\)
  conditions, nested three-condition rules of the form
  \(A\ L_{\text{outer}}\ (B\ L_{\text{inner}}\ C)\), and continuous
  thresholds. Different grammars would give different posteriors.
  Sensitivity to \texttt{max\_dims} is exercised in §``Sensitivity to
  \texttt{max\_dims}''; sensitivity to operators and connectives
  (including deeper nesting) is left as an exercise.
\item
  \textbf{Model B is structurally biased relative to Model A} (Phase 0
  Q3, Q5): the Stan proxy is a different model from the cognitive
  theory. Validation of Model B (LOO, prior sensitivity, real SBC, PSBC)
  tests the proxy. The NPE vs Stan comparison in §``The Main Empirical
  Result'' and the gap demonstration make the bias visible. Chapter 15
  should report both posteriors when comparing rule-based models against
  GCM and prototype.
\item
  \textbf{NPE has its own approximation error}: the NPE posterior is the
  output of a learned normalising flow, which is an approximation to the
  true posterior. We validate it via SBC + expected coverage (§``NPE
  Validation''); if both pass, the approximation error is small enough
  to trust. Unlike rejection ABC, NPE does not depend on hand-crafted
  summary statistics --- the embedding network learns the relevant
  features end-to-end --- and it does not depend on a tolerance
  parameter.
\item
  \textbf{Cross-language step}: NPE requires Python (\texttt{sbi} +
  \texttt{torch}) accessed via \texttt{reticulate}. This adds a
  dependency but is well-documented and the interface is stable. The
  alternative --- reimplementing normalising flows in R --- is not
  advisable.
\item
  \textbf{No multilevel extension yet}: this chapter, like Chapter 12,
  stops at single-subject inference. The multilevel rule model (with
  population-level priors on \(\varepsilon\) and per-subject offsets) is
  straightforward to add to Model B but non-trivial for the NPE
  pipeline; Chapter 15 takes up the multilevel comparison across all
  three model classes. For NPE, the standard approach is to train the
  network at the single-subject level and then build the hierarchical
  model on top of the individual posteriors (Radev et al.~2023,
  §``Hierarchical models'').
\end{itemize}

\section{Conclusion: The Role of Rules and the Cost of
Tractability}\label{conclusion-the-role-of-rules-and-the-cost-of-tractability}

Rule-based models offer a distinct perspective on categorization,
emphasizing explicit criteria and logical operations rather than overall
similarity. The Bayesian particle filter (Model A) provides a
sophisticated framework for modeling how learners might search the space
of possible rules, maintain uncertainty, and update beliefs based on
evidence. Its likelihood is intractable in Stan, which is why the
chapter introduced two responses to that intractability:

\begin{enumerate}
\def\labelenumi{\arabic{enumi}.}
\item
  \textbf{The tractable proxy (Model B)}: a marginalised mixture over a
  fixed candidate rule set, fitted with HMC, validated via the standard
  six-phase pipeline (including Posterior SBC for local calibration).
  Useful for downstream model comparison, \emph{biased} relative to the
  cognitive theory in a direction the chapter quantifies via the
  NPE-vs-Stan posterior gap.
\item
  \textbf{Simulation-based inference (NPE for Model A)}: Neural
  Posterior Estimation with a learned embedding network, validated via
  SBC + expected coverage, fitting the actual cognitive theory directly.
  The NPE approach resolves the three problems of earlier ABC drafts
  (schedule mismatch, summary insufficiency, no amortization) and makes
  \(N_{\text{particles}}\) a jointly inferred cognitive parameter rather
  than a structural commitment.
\end{enumerate}

The chapter's pedagogical move is to fit \emph{both}. The gap between
the two posteriors is the structural bias the modeller accepts when they
replace the cognitive theory with a proxy. The gap demonstration
experiment (§``Demonstrating That the Gap Tracks Rule-Set Mismatch'')
shows this gap tracks rule-set mismatch rather than inference noise. On
data where the proxy's candidate rules happen to cover what Model A
converges to, the gap is small and Model B's LOO scores can be trusted
as cognitive evidence. On data where they do not, the gap is large and
Model B's LOO scores are \emph{structural artefacts} --- they tell you
about the proxy's fit, not about the agent.

Pure rule-based models may not account for all aspects of human
categorization (especially for natural categories with fuzzy
boundaries), but they capture an important element of human cognition,
particularly in domains involving definitions, instruction, or
well-defined criteria. Often, human behavior might reflect a blend or
interplay between rule-based and similarity-based processing --- a topic
explored in hybrid models (like COVIS or RULEX) and addressed directly
in the \textbf{next chapter}, which compares all three model classes on
the same dataset. Chapter 15 reports the comparison \textbf{twice}:

\begin{itemize}
\tightlist
\item
  Once under the LOO/LFO-CV framework, with each model fitted under its
  appropriate validation tool (LFO-CV for the GCM and Kalman prototype
  via the implementations in Chapters 11 and 12; standard LOO for Model
  B). This is the comparison the proxies make possible.
\item
  Once under NPE, with Model A (the particle filter), the GCM, and the
  Kalman prototype all fitted as simulator-based models on a common
  footing --- each with its own trained inference network, validated
  with SBC + expected coverage, and compared via posterior predictive
  ELPD. This is the comparison the cognitive theories deserve.
\end{itemize}

The expected outcome is that the two comparisons agree on the easy cases
and disagree on the hard ones, and that the disagreement is informative
about which empirical conclusions in the categorisation literature have
been over-claimed because they were based on tractable proxies of
intractable cognitive theories. That is the methodological payoff of
taking simulator-based inference seriously.

\cleardoublepage
\phantomsection
\addcontentsline{toc}{part}{Appendices}
\appendix

\chapter{Sequential Extensions of the Bayesian
Agent}\label{sequential-extensions-of-the-bayesian-agent}

Every Bayesian agent we built in Ch. 10 --- SBA, WBA, PBA, and their
multilevel cousins --- shares two strong assumptions. First, the agent
is \emph{forced} to guess on every trial: there is no option to defer,
ask for more evidence, or decline. Second, the weights on direct and
social evidence are \emph{static}: whatever trade-off an agent starts
with, they carry it through the entire experiment unchanged. Both
assumptions are reasonable first-pass idealisations, and both are almost
certainly false about real cognition.

This appendix sketches two extensions that relax each assumption in
turn. Section A introduces a \textbf{sample-or-guess agent} that decides
at every step whether to collect more evidence or commit to a choice.
Section B introduces an \textbf{adaptive-weight agent} whose weights
drift across trials in response to how well each source predicted recent
feedback.

\begin{quote}
\textbf{Scope note.} These two models are offered as \emph{starting
points for student projects}, not as fully validated contributions. We
give the mathematical formalization, an R forward simulation,
visualisations of the agents' behaviour, and a single-agent Stan model
fit once for illustration. We deliberately omit parameter recovery, SBC,
model comparison, and multilevel extensions --- all of these are natural
next steps and good exercises. Treat what follows as a skeleton to build
on.
\end{quote}

\section{Section A: Sample-or-Guess
Agent}\label{section-a-sample-or-guess-agent}

\subsection{Conceptual motivation}\label{conceptual-motivation}

The agents in Ch. 10 are committed decision-makers: evidence arrives,
they integrate it, and they respond. In many real tasks --- clinical
diagnosis, foraging, perceptual decision-making under time pressure ---
the agent instead chooses \emph{when to stop sampling}. A cautious agent
keeps observing marbles until confident; an impulsive one guesses after
the first draw. The trade-off is classical: more samples mean better
accuracy but higher cost (time, effort, missed opportunities). This
family of policies connects the Bayesian-cognition literature to
sequential-sampling models of reaction time (e.g., the drift-diffusion
model) and to metacognition, where confidence gates commitment.

The key conceptual move is that the observation model now has \emph{two}
outputs per trial. On every step the agent emits a stop/continue
decision; when it finally stops, it emits a guess. The likelihood must
account for both --- a point that is easy to miss when porting the Ch.
10 Stan templates.

\subsection{Mathematical
formalization}\label{mathematical-formalization}

Let trial \(i\) consist of a sequence of evidence draws. The agent
enters the trial with a prior (we use Jeffreys:
\(\alpha_0 = \beta_0 = 0.5\)) and an optional social tip summarised as
\((k_s, n_s)\). After observing direct evidence up to step \(t\) ---
\(k_{d,t}\) blue out of \(t\) direct draws --- the posterior over the
jar's blue proportion \(\theta\) is

\[\theta \mid \text{evidence}_t \sim \text{Beta}(\alpha_t,\, \beta_t),\]

with

\[\alpha_t = 0.5 + k_s + k_{d,t}, \qquad
  \beta_t  = 0.5 + (n_s - k_s) + (t - k_{d,t}).\]

Define the agent's \textbf{certainty} at step \(t\) as

\[c_t \;=\; \left|\, \mathbb{E}[\theta \mid \text{evidence}_t] - 0.5 \,\right|
       \;=\; \left|\, \frac{\alpha_t}{\alpha_t + \beta_t} - 0.5 \,\right|.\]

The agent's policy is governed by a single free parameter \(\tau \in
(0, 0.5)\):

\[
\text{at step } t: \quad
\begin{cases}
  \text{stop, guess } g_i = \mathbf{1}[\alpha_t > \beta_t] & \text{if } c_t \ge \tau, \\
  \text{draw one more marble, increment } t & \text{otherwise.}
\end{cases}
\]

The data observed per trial are the stopping time \(T_i\) and the final
guess \(g_i\). The likelihood factorises as

\[p(T_i, g_i \mid \tau, \text{evidence}) \;=\;
  \underbrace{\prod_{t=1}^{T_i - 1} \Pr(\text{continue at } t \mid \tau)}_{\text{didn't stop earlier}}
  \;\cdot\;
  \underbrace{\Pr(\text{stop at } T_i \mid \tau)}_{\text{stops now}}
  \;\cdot\;
  \underbrace{\Pr(g_i \mid \alpha_{T_i}, \beta_{T_i})}_{\text{final guess}}.\]

Two important caveats. \textbf{First}, after the agent stops, the
would-be marbles at steps \(T_i + 1, T_i + 2, \ldots\) are \emph{not}
observed. Trials are not i.i.d. sequences of a fixed length; the data
are censored by the agent's own decision. \textbf{Second}, the
hard-threshold rule above is non-differentiable in \(\tau\), which
breaks HMC. For inference we therefore replace it with a soft version,

\[\Pr(\text{stop at } t \mid \tau) \;=\; \sigma\!\left(\frac{c_t - \tau}{\sigma_\tau}\right),\]

where \(\sigma(\cdot)\) is the logistic function and \(\sigma_\tau\) is
a small fixed slope (we use \(\sigma_\tau = 0.02\)). As
\(\sigma_\tau \to 0\) the soft rule recovers the hard one.

\subsection{R forward simulation}\label{r-forward-simulation}

\begin{Shaded}
\begin{Highlighting}[]
\CommentTok{\# Simulate one agent with a hard certainty threshold.}
\CommentTok{\# Returns a tibble with one row per trial.}
\NormalTok{simulate\_threshold\_agent }\OtherTok{\textless{}{-}} \ControlFlowTok{function}\NormalTok{(tau,}
\NormalTok{                                     n\_trials,}
                                     \AttributeTok{max\_samples =} \DecValTok{30}\NormalTok{,}
                                     \AttributeTok{true\_theta =} \FloatTok{0.7}\NormalTok{,}
                                     \AttributeTok{social\_k =} \DecValTok{1}\NormalTok{,}
                                     \AttributeTok{social\_n =} \DecValTok{3}\NormalTok{,}
                                     \AttributeTok{alpha0 =} \FloatTok{0.5}\NormalTok{, }\AttributeTok{beta0 =} \FloatTok{0.5}\NormalTok{) \{}
\NormalTok{  purrr}\SpecialCharTok{::}\FunctionTok{map\_dfr}\NormalTok{(}\FunctionTok{seq\_len}\NormalTok{(n\_trials), }\ControlFlowTok{function}\NormalTok{(trial) \{}
\NormalTok{    alpha\_t }\OtherTok{\textless{}{-}}\NormalTok{ alpha0 }\SpecialCharTok{+}\NormalTok{ social\_k}
\NormalTok{    beta\_t  }\OtherTok{\textless{}{-}}\NormalTok{ beta0  }\SpecialCharTok{+}\NormalTok{ (social\_n }\SpecialCharTok{{-}}\NormalTok{ social\_k)}
\NormalTok{    stopped }\OtherTok{\textless{}{-}} \ConstantTok{FALSE}
\NormalTok{    t }\OtherTok{\textless{}{-}} \DecValTok{0}
    \ControlFlowTok{while}\NormalTok{ (}\SpecialCharTok{!}\NormalTok{stopped }\SpecialCharTok{\&\&}\NormalTok{ t }\SpecialCharTok{\textless{}}\NormalTok{ max\_samples) \{}
\NormalTok{      c\_t }\OtherTok{\textless{}{-}} \FunctionTok{abs}\NormalTok{(alpha\_t }\SpecialCharTok{/}\NormalTok{ (alpha\_t }\SpecialCharTok{+}\NormalTok{ beta\_t) }\SpecialCharTok{{-}} \FloatTok{0.5}\NormalTok{)}
      \ControlFlowTok{if}\NormalTok{ (c\_t }\SpecialCharTok{\textgreater{}=}\NormalTok{ tau }\SpecialCharTok{\&\&}\NormalTok{ t }\SpecialCharTok{\textgreater{}=} \DecValTok{1}\NormalTok{) \{}
\NormalTok{        stopped }\OtherTok{\textless{}{-}} \ConstantTok{TRUE}
\NormalTok{      \} }\ControlFlowTok{else}\NormalTok{ \{}
\NormalTok{        t }\OtherTok{\textless{}{-}}\NormalTok{ t }\SpecialCharTok{+} \DecValTok{1}
\NormalTok{        draw }\OtherTok{\textless{}{-}} \FunctionTok{rbinom}\NormalTok{(}\DecValTok{1}\NormalTok{, }\DecValTok{1}\NormalTok{, true\_theta)  }\CommentTok{\# 1 = blue, 0 = red}
\NormalTok{        alpha\_t }\OtherTok{\textless{}{-}}\NormalTok{ alpha\_t }\SpecialCharTok{+}\NormalTok{ draw}
\NormalTok{        beta\_t  }\OtherTok{\textless{}{-}}\NormalTok{ beta\_t }\SpecialCharTok{+}\NormalTok{ (}\DecValTok{1} \SpecialCharTok{{-}}\NormalTok{ draw)}
\NormalTok{      \}}
\NormalTok{    \}}
\NormalTok{    guess }\OtherTok{\textless{}{-}} \FunctionTok{as.integer}\NormalTok{(alpha\_t }\SpecialCharTok{\textgreater{}}\NormalTok{ beta\_t)}
    \FunctionTok{tibble}\NormalTok{(}
      \AttributeTok{trial   =}\NormalTok{ trial,}
      \AttributeTok{tau     =}\NormalTok{ tau,}
      \AttributeTok{T       =}\NormalTok{ t,}
      \AttributeTok{guess   =}\NormalTok{ guess,}
      \AttributeTok{correct =} \FunctionTok{as.integer}\NormalTok{(guess }\SpecialCharTok{==} \FunctionTok{as.integer}\NormalTok{(true\_theta }\SpecialCharTok{\textgreater{}} \FloatTok{0.5}\NormalTok{))}
\NormalTok{    )}
\NormalTok{  \})}
\NormalTok{\}}

\CommentTok{\# Three agents with different thresholds}
\NormalTok{sim\_threshold }\OtherTok{\textless{}{-}} \FunctionTok{bind\_rows}\NormalTok{(}
  \FunctionTok{simulate\_threshold\_agent}\NormalTok{(}\AttributeTok{tau =} \FloatTok{0.10}\NormalTok{, }\AttributeTok{n\_trials =} \DecValTok{200}\NormalTok{),}
  \FunctionTok{simulate\_threshold\_agent}\NormalTok{(}\AttributeTok{tau =} \FloatTok{0.25}\NormalTok{, }\AttributeTok{n\_trials =} \DecValTok{200}\NormalTok{),}
  \FunctionTok{simulate\_threshold\_agent}\NormalTok{(}\AttributeTok{tau =} \FloatTok{0.40}\NormalTok{, }\AttributeTok{n\_trials =} \DecValTok{200}\NormalTok{)}
\NormalTok{) }\SpecialCharTok{|\textgreater{}}
  \FunctionTok{mutate}\NormalTok{(}\AttributeTok{tau\_label =} \FunctionTok{factor}\NormalTok{(}\FunctionTok{paste0}\NormalTok{(}\StringTok{"tau = "}\NormalTok{, tau)))}
\end{Highlighting}
\end{Shaded}

\subsection{Visualisations}\label{visualisations}

\begin{Shaded}
\begin{Highlighting}[]
\CommentTok{\# (a) Stopping{-}time distributions per tau}
\NormalTok{p\_stop }\OtherTok{\textless{}{-}}\NormalTok{ sim\_threshold }\SpecialCharTok{|\textgreater{}}
  \FunctionTok{ggplot}\NormalTok{(}\FunctionTok{aes}\NormalTok{(T, }\AttributeTok{fill =}\NormalTok{ tau\_label)) }\SpecialCharTok{+}
  \FunctionTok{geom\_histogram}\NormalTok{(}\AttributeTok{binwidth =} \DecValTok{1}\NormalTok{, }\AttributeTok{colour =} \StringTok{"white"}\NormalTok{, }\AttributeTok{alpha =} \FloatTok{0.8}\NormalTok{) }\SpecialCharTok{+}
  \FunctionTok{facet\_wrap}\NormalTok{(}\SpecialCharTok{\textasciitilde{}}\NormalTok{tau\_label, }\AttributeTok{ncol =} \DecValTok{3}\NormalTok{) }\SpecialCharTok{+}
  \FunctionTok{labs}\NormalTok{(}\AttributeTok{x =} \StringTok{"Stopping time (\# direct draws)"}\NormalTok{,}
       \AttributeTok{y =} \StringTok{"Count"}\NormalTok{,}
       \AttributeTok{title =} \StringTok{"(a) Stopping{-}time distributions"}\NormalTok{) }\SpecialCharTok{+}
  \FunctionTok{theme}\NormalTok{(}\AttributeTok{legend.position =} \StringTok{"none"}\NormalTok{)}

\CommentTok{\# (b) Accuracy vs. tau (sweep on a denser grid)}
\NormalTok{tau\_sweep }\OtherTok{\textless{}{-}} \FunctionTok{seq}\NormalTok{(}\FloatTok{0.05}\NormalTok{, }\FloatTok{0.45}\NormalTok{, }\AttributeTok{by =} \FloatTok{0.05}\NormalTok{)}
\NormalTok{acc\_df }\OtherTok{\textless{}{-}}\NormalTok{ purrr}\SpecialCharTok{::}\FunctionTok{map\_dfr}\NormalTok{(tau\_sweep, }\ControlFlowTok{function}\NormalTok{(tt) \{}
\NormalTok{  s }\OtherTok{\textless{}{-}} \FunctionTok{simulate\_threshold\_agent}\NormalTok{(}\AttributeTok{tau =}\NormalTok{ tt, }\AttributeTok{n\_trials =} \DecValTok{400}\NormalTok{)}
  \FunctionTok{tibble}\NormalTok{(}\AttributeTok{tau =}\NormalTok{ tt,}
         \AttributeTok{accuracy    =} \FunctionTok{mean}\NormalTok{(s}\SpecialCharTok{$}\NormalTok{correct),}
         \AttributeTok{mean\_T      =} \FunctionTok{mean}\NormalTok{(s}\SpecialCharTok{$}\NormalTok{T))}
\NormalTok{\})}

\NormalTok{p\_acc }\OtherTok{\textless{}{-}}\NormalTok{ acc\_df }\SpecialCharTok{|\textgreater{}}
  \FunctionTok{ggplot}\NormalTok{(}\FunctionTok{aes}\NormalTok{(tau, accuracy)) }\SpecialCharTok{+}
  \FunctionTok{geom\_line}\NormalTok{(}\AttributeTok{size =} \DecValTok{1}\NormalTok{) }\SpecialCharTok{+}
  \FunctionTok{geom\_point}\NormalTok{(}\FunctionTok{aes}\NormalTok{(}\AttributeTok{size =}\NormalTok{ mean\_T)) }\SpecialCharTok{+}
  \FunctionTok{scale\_size\_continuous}\NormalTok{(}\AttributeTok{name =} \StringTok{"Mean}\SpecialCharTok{\textbackslash{}n}\StringTok{samples"}\NormalTok{) }\SpecialCharTok{+}
  \FunctionTok{labs}\NormalTok{(}\AttributeTok{x =} \FunctionTok{expression}\NormalTok{(tau), }\AttributeTok{y =} \StringTok{"Accuracy"}\NormalTok{,}
       \AttributeTok{title =} \StringTok{"(b) Accuracy vs. threshold"}\NormalTok{)}

\NormalTok{p\_stop }\SpecialCharTok{/}\NormalTok{ p\_acc}
\end{Highlighting}
\end{Shaded}

\begin{center}
\includegraphics[width=0.8\linewidth,height=\textheight,keepaspectratio]{10a-BayesianSequentialExtensions_files/figure-pdf/ch11_threshold_viz-1.pdf}
\end{center}

\begin{Shaded}
\begin{Highlighting}[]
\CommentTok{\# (c) Trial{-}trace: certainty evolving toward the threshold}
\FunctionTok{set.seed}\NormalTok{(}\DecValTok{42}\NormalTok{)}
\NormalTok{trace\_tau }\OtherTok{\textless{}{-}} \FloatTok{0.25}
\NormalTok{alpha\_t }\OtherTok{\textless{}{-}} \FloatTok{0.5} \SpecialCharTok{+} \DecValTok{1}
\NormalTok{beta\_t  }\OtherTok{\textless{}{-}} \FloatTok{0.5} \SpecialCharTok{+}\NormalTok{ (}\DecValTok{3} \SpecialCharTok{{-}} \DecValTok{1}\NormalTok{)}
\NormalTok{trace }\OtherTok{\textless{}{-}} \FunctionTok{tibble}\NormalTok{(}
  \AttributeTok{t       =} \DecValTok{0}\NormalTok{,}
  \AttributeTok{alpha\_t =}\NormalTok{ alpha\_t,}
  \AttributeTok{beta\_t  =}\NormalTok{ beta\_t,}
  \AttributeTok{c\_t     =} \FunctionTok{abs}\NormalTok{(alpha\_t }\SpecialCharTok{/}\NormalTok{ (alpha\_t }\SpecialCharTok{+}\NormalTok{ beta\_t) }\SpecialCharTok{{-}} \FloatTok{0.5}\NormalTok{)}
\NormalTok{)}
\ControlFlowTok{for}\NormalTok{ (t }\ControlFlowTok{in} \DecValTok{1}\SpecialCharTok{:}\DecValTok{20}\NormalTok{) \{}
\NormalTok{  draw }\OtherTok{\textless{}{-}} \FunctionTok{rbinom}\NormalTok{(}\DecValTok{1}\NormalTok{, }\DecValTok{1}\NormalTok{, }\FloatTok{0.75}\NormalTok{)}
\NormalTok{  alpha\_t }\OtherTok{\textless{}{-}}\NormalTok{ alpha\_t }\SpecialCharTok{+}\NormalTok{ draw}
\NormalTok{  beta\_t  }\OtherTok{\textless{}{-}}\NormalTok{ beta\_t }\SpecialCharTok{+}\NormalTok{ (}\DecValTok{1} \SpecialCharTok{{-}}\NormalTok{ draw)}
\NormalTok{  trace }\OtherTok{\textless{}{-}} \FunctionTok{bind\_rows}\NormalTok{(trace, }\FunctionTok{tibble}\NormalTok{(}
    \AttributeTok{t =}\NormalTok{ t,}
    \AttributeTok{alpha\_t =}\NormalTok{ alpha\_t, }\AttributeTok{beta\_t =}\NormalTok{ beta\_t,}
    \AttributeTok{c\_t =} \FunctionTok{abs}\NormalTok{(alpha\_t }\SpecialCharTok{/}\NormalTok{ (alpha\_t }\SpecialCharTok{+}\NormalTok{ beta\_t) }\SpecialCharTok{{-}} \FloatTok{0.5}\NormalTok{)}
\NormalTok{  ))}
  \ControlFlowTok{if}\NormalTok{ (}\FunctionTok{abs}\NormalTok{(alpha\_t }\SpecialCharTok{/}\NormalTok{ (alpha\_t }\SpecialCharTok{+}\NormalTok{ beta\_t) }\SpecialCharTok{{-}} \FloatTok{0.5}\NormalTok{) }\SpecialCharTok{\textgreater{}=}\NormalTok{ trace\_tau) }\ControlFlowTok{break}
\NormalTok{\}}
\NormalTok{stop\_t }\OtherTok{\textless{}{-}} \FunctionTok{max}\NormalTok{(trace}\SpecialCharTok{$}\NormalTok{t)}

\NormalTok{trace }\SpecialCharTok{|\textgreater{}}
  \FunctionTok{ggplot}\NormalTok{(}\FunctionTok{aes}\NormalTok{(t, c\_t)) }\SpecialCharTok{+}
  \FunctionTok{geom\_line}\NormalTok{(}\AttributeTok{size =} \DecValTok{1}\NormalTok{) }\SpecialCharTok{+}
  \FunctionTok{geom\_point}\NormalTok{(}\AttributeTok{size =} \DecValTok{2}\NormalTok{) }\SpecialCharTok{+}
  \FunctionTok{geom\_hline}\NormalTok{(}\AttributeTok{yintercept =}\NormalTok{ trace\_tau, }\AttributeTok{linetype =} \StringTok{"dashed"}\NormalTok{, }\AttributeTok{colour =} \StringTok{"firebrick"}\NormalTok{) }\SpecialCharTok{+}
  \FunctionTok{annotate}\NormalTok{(}\StringTok{"text"}\NormalTok{, }\AttributeTok{x =} \DecValTok{1}\NormalTok{, }\AttributeTok{y =}\NormalTok{ trace\_tau }\SpecialCharTok{+} \FloatTok{0.015}\NormalTok{,}
           \AttributeTok{label =} \FunctionTok{expression}\NormalTok{(tau), }\AttributeTok{colour =} \StringTok{"firebrick"}\NormalTok{, }\AttributeTok{hjust =} \DecValTok{0}\NormalTok{) }\SpecialCharTok{+}
  \FunctionTok{annotate}\NormalTok{(}\StringTok{"point"}\NormalTok{, }\AttributeTok{x =}\NormalTok{ stop\_t, }\AttributeTok{y =}\NormalTok{ trace}\SpecialCharTok{$}\NormalTok{c\_t[trace}\SpecialCharTok{$}\NormalTok{t }\SpecialCharTok{==}\NormalTok{ stop\_t],}
           \AttributeTok{colour =} \StringTok{"firebrick"}\NormalTok{, }\AttributeTok{size =} \DecValTok{4}\NormalTok{) }\SpecialCharTok{+}
  \FunctionTok{labs}\NormalTok{(}\AttributeTok{x =} \StringTok{"Direct draw step t"}\NormalTok{, }\AttributeTok{y =} \FunctionTok{expression}\NormalTok{(c[t]),}
       \AttributeTok{title =} \StringTok{"(c) One trial: certainty crossing the threshold"}\NormalTok{)}
\end{Highlighting}
\end{Shaded}

\begin{center}
\includegraphics[width=0.8\linewidth,height=\textheight,keepaspectratio]{10a-BayesianSequentialExtensions_files/figure-pdf/ch11_threshold_trace-1.pdf}
\end{center}

The patterns line up with intuition: a low \(\tau\) produces short,
noisy, error-prone trials; a high \(\tau\) produces longer, more
accurate trials with diminishing returns.

\subsection{Single-agent Stan model (soft
threshold)}\label{single-agent-stan-model-soft-threshold}

The hard rule above is unsuitable for HMC because it is
non-differentiable in \(\tau\). We replace it with a logistic soft rule.
The likelihood then becomes a product of Bernoulli stop/continue
probabilities plus a Bernoulli for the final guess. Because the
direct-evidence sequence itself is a latent variable conditional only on
\(T_i\) and the final \((\alpha, \beta)\), we pass the raw per-step
draws (reconstructed from the simulation) to Stan as data. In a real
analysis the modeller would either record these draws in the experiment
or marginalise over them.

\begin{Shaded}
\begin{Highlighting}[]
\CommentTok{\# Use the tau = 0.25 agent for fitting. Re{-}run the simulation but also}
\CommentTok{\# record the sequence of draws for each trial.}
\NormalTok{simulate\_threshold\_agent\_full }\OtherTok{\textless{}{-}} \ControlFlowTok{function}\NormalTok{(tau, n\_trials, }\AttributeTok{max\_samples =} \DecValTok{30}\NormalTok{,}
                                          \AttributeTok{true\_theta =} \FloatTok{0.7}\NormalTok{,}
                                          \AttributeTok{social\_k =} \DecValTok{1}\NormalTok{, }\AttributeTok{social\_n =} \DecValTok{3}\NormalTok{) \{}
\NormalTok{  draws\_list }\OtherTok{\textless{}{-}} \FunctionTok{vector}\NormalTok{(}\StringTok{"list"}\NormalTok{, n\_trials)}
\NormalTok{  T\_vec }\OtherTok{\textless{}{-}} \FunctionTok{integer}\NormalTok{(n\_trials)}
\NormalTok{  g\_vec }\OtherTok{\textless{}{-}} \FunctionTok{integer}\NormalTok{(n\_trials)}
  \ControlFlowTok{for}\NormalTok{ (i }\ControlFlowTok{in} \FunctionTok{seq\_len}\NormalTok{(n\_trials)) \{}
\NormalTok{    alpha\_t }\OtherTok{\textless{}{-}} \FloatTok{0.5} \SpecialCharTok{+}\NormalTok{ social\_k}
\NormalTok{    beta\_t  }\OtherTok{\textless{}{-}} \FloatTok{0.5} \SpecialCharTok{+}\NormalTok{ (social\_n }\SpecialCharTok{{-}}\NormalTok{ social\_k)}
\NormalTok{    draws\_i }\OtherTok{\textless{}{-}} \FunctionTok{integer}\NormalTok{(}\DecValTok{0}\NormalTok{)}
\NormalTok{    t }\OtherTok{\textless{}{-}} \DecValTok{0}
\NormalTok{    stopped }\OtherTok{\textless{}{-}} \ConstantTok{FALSE}
    \ControlFlowTok{while}\NormalTok{ (}\SpecialCharTok{!}\NormalTok{stopped }\SpecialCharTok{\&\&}\NormalTok{ t }\SpecialCharTok{\textless{}}\NormalTok{ max\_samples) \{}
\NormalTok{      c\_t }\OtherTok{\textless{}{-}} \FunctionTok{abs}\NormalTok{(alpha\_t }\SpecialCharTok{/}\NormalTok{ (alpha\_t }\SpecialCharTok{+}\NormalTok{ beta\_t) }\SpecialCharTok{{-}} \FloatTok{0.5}\NormalTok{)}
      \ControlFlowTok{if}\NormalTok{ (c\_t }\SpecialCharTok{\textgreater{}=}\NormalTok{ tau }\SpecialCharTok{\&\&}\NormalTok{ t }\SpecialCharTok{\textgreater{}=} \DecValTok{1}\NormalTok{) \{ stopped }\OtherTok{\textless{}{-}} \ConstantTok{TRUE}\NormalTok{; }\ControlFlowTok{break}\NormalTok{ \}}
\NormalTok{      t }\OtherTok{\textless{}{-}}\NormalTok{ t }\SpecialCharTok{+} \DecValTok{1}
\NormalTok{      d }\OtherTok{\textless{}{-}} \FunctionTok{rbinom}\NormalTok{(}\DecValTok{1}\NormalTok{, }\DecValTok{1}\NormalTok{, true\_theta)}
\NormalTok{      draws\_i }\OtherTok{\textless{}{-}} \FunctionTok{c}\NormalTok{(draws\_i, d)}
\NormalTok{      alpha\_t }\OtherTok{\textless{}{-}}\NormalTok{ alpha\_t }\SpecialCharTok{+}\NormalTok{ d}
\NormalTok{      beta\_t  }\OtherTok{\textless{}{-}}\NormalTok{ beta\_t }\SpecialCharTok{+}\NormalTok{ (}\DecValTok{1} \SpecialCharTok{{-}}\NormalTok{ d)}
\NormalTok{    \}}
\NormalTok{    T\_vec[i] }\OtherTok{\textless{}{-}}\NormalTok{ t}
\NormalTok{    g\_vec[i] }\OtherTok{\textless{}{-}} \FunctionTok{as.integer}\NormalTok{(alpha\_t }\SpecialCharTok{\textgreater{}}\NormalTok{ beta\_t)}
\NormalTok{    draws\_list[[i]] }\OtherTok{\textless{}{-}}\NormalTok{ draws\_i}
\NormalTok{  \}}
  \FunctionTok{list}\NormalTok{(}\AttributeTok{T =}\NormalTok{ T\_vec, }\AttributeTok{guess =}\NormalTok{ g\_vec, }\AttributeTok{draws =}\NormalTok{ draws\_list)}
\NormalTok{\}}

\NormalTok{sim\_full }\OtherTok{\textless{}{-}} \FunctionTok{simulate\_threshold\_agent\_full}\NormalTok{(}\AttributeTok{tau =} \FloatTok{0.25}\NormalTok{, }\AttributeTok{n\_trials =} \DecValTok{120}\NormalTok{)}

\CommentTok{\# Flatten per{-}trial draws into a padded matrix (max T across trials)}
\NormalTok{T\_max\_obs }\OtherTok{\textless{}{-}} \FunctionTok{max}\NormalTok{(sim\_full}\SpecialCharTok{$}\NormalTok{T)}
\NormalTok{draws\_mat }\OtherTok{\textless{}{-}} \FunctionTok{matrix}\NormalTok{(}\DecValTok{0}\NormalTok{L, }\AttributeTok{nrow =} \FunctionTok{length}\NormalTok{(sim\_full}\SpecialCharTok{$}\NormalTok{T), }\AttributeTok{ncol =}\NormalTok{ T\_max\_obs)}
\ControlFlowTok{for}\NormalTok{ (i }\ControlFlowTok{in} \FunctionTok{seq\_along}\NormalTok{(sim\_full}\SpecialCharTok{$}\NormalTok{draws)) \{}
  \ControlFlowTok{if}\NormalTok{ (sim\_full}\SpecialCharTok{$}\NormalTok{T[i] }\SpecialCharTok{\textgreater{}} \DecValTok{0}\NormalTok{) \{}
\NormalTok{    draws\_mat[i, }\DecValTok{1}\SpecialCharTok{:}\NormalTok{sim\_full}\SpecialCharTok{$}\NormalTok{T[i]] }\OtherTok{\textless{}{-}}\NormalTok{ sim\_full}\SpecialCharTok{$}\NormalTok{draws[[i]]}
\NormalTok{  \}}
\NormalTok{\}}

\NormalTok{stan\_data\_A }\OtherTok{\textless{}{-}} \FunctionTok{list}\NormalTok{(}
  \AttributeTok{N        =} \FunctionTok{length}\NormalTok{(sim\_full}\SpecialCharTok{$}\NormalTok{T),}
  \AttributeTok{T\_max    =}\NormalTok{ T\_max\_obs,}
  \AttributeTok{T        =}\NormalTok{ sim\_full}\SpecialCharTok{$}\NormalTok{T,}
  \AttributeTok{guess    =}\NormalTok{ sim\_full}\SpecialCharTok{$}\NormalTok{guess,}
  \AttributeTok{draws    =}\NormalTok{ draws\_mat,}
  \AttributeTok{social\_k =} \DecValTok{1}\NormalTok{,}
  \AttributeTok{social\_n =} \DecValTok{3}\NormalTok{,}
  \AttributeTok{sigma\_tau =} \FloatTok{0.02}
\NormalTok{)}
\end{Highlighting}
\end{Shaded}

\begin{Shaded}
\begin{Highlighting}[]
\NormalTok{ThresholdAgent\_stan }\OtherTok{\textless{}{-}} \StringTok{"}
\StringTok{// Sample{-}or{-}Guess Agent with a soft certainty threshold.}
\StringTok{// Single parameter: tau in (0, 0.5), the certainty threshold.}
\StringTok{// The stop/continue decision is modelled with a logistic soft rule}
\StringTok{// around tau with fixed slope sigma\_tau.}

\StringTok{data \{}
\StringTok{  int\textless{}lower=1\textgreater{} N;}
\StringTok{  int\textless{}lower=1\textgreater{} T\_max;}
\StringTok{  array[N] int\textless{}lower=0, upper=T\_max\textgreater{} T;}
\StringTok{  array[N] int\textless{}lower=0, upper=1\textgreater{} guess;}
\StringTok{  array[N, T\_max] int\textless{}lower=0, upper=1\textgreater{} draws;}
\StringTok{  int\textless{}lower=0\textgreater{} social\_k;}
\StringTok{  int\textless{}lower=0\textgreater{} social\_n;}
\StringTok{  real\textless{}lower=0\textgreater{} sigma\_tau;}
\StringTok{\}}

\StringTok{parameters \{}
\StringTok{  // tau \textasciitilde{} scaled Beta(2,2) on (0, 0.5)}
\StringTok{  real\textless{}lower=0, upper=0.5\textgreater{} tau;}
\StringTok{\}}

\StringTok{model \{}
\StringTok{  // Prior: Beta(2,2) on (0, 0.5) via change of variables {-}\textgreater{} tau * 2 \textasciitilde{} Beta(2,2)}
\StringTok{  target += beta\_lpdf(2 * tau | 2, 2) + log(2);}

\StringTok{  for (i in 1:N) \{}
\StringTok{    real alpha\_t = 0.5 + social\_k;}
\StringTok{    real beta\_t  = 0.5 + (social\_n {-} social\_k);}

\StringTok{    // Steps 1 .. T[i] {-} 1: the agent continued (soft prob of NOT stopping)}
\StringTok{    // At step T[i]: the agent stopped (soft prob of stopping)}
\StringTok{    // If T[i] == 0 (agent would stop before seeing any direct draw) we}
\StringTok{    // only contribute the final guess.}
\StringTok{    for (t in 1:T[i]) \{}
\StringTok{      int d = draws[i, t];}
\StringTok{      alpha\_t += d;}
\StringTok{      beta\_t  += (1 {-} d);}
\StringTok{      real c\_t = abs(alpha\_t / (alpha\_t + beta\_t) {-} 0.5);}
\StringTok{      real p\_stop = inv\_logit((c\_t {-} tau) / sigma\_tau);}
\StringTok{      if (t \textless{} T[i]) \{}
\StringTok{        target += bernoulli\_lpmf(0 | p\_stop);  // continued}
\StringTok{      \} else \{}
\StringTok{        target += bernoulli\_lpmf(1 | p\_stop);  // stopped}
\StringTok{      \}}
\StringTok{    \}}

\StringTok{    // Final guess given the posterior at stopping time}
\StringTok{    \{}
\StringTok{      real alpha\_final = 0.5 + social\_k;}
\StringTok{      real beta\_final  = 0.5 + (social\_n {-} social\_k);}
\StringTok{      for (t in 1:T[i]) \{}
\StringTok{        alpha\_final += draws[i, t];}
\StringTok{        beta\_final  += (1 {-} draws[i, t]);}
\StringTok{      \}}
\StringTok{      real p\_blue = alpha\_final / (alpha\_final + beta\_final);}
\StringTok{      target += bernoulli\_lpmf(guess[i] | p\_blue);}
\StringTok{    \}}
\StringTok{  \}}
\StringTok{\}}
\StringTok{"}

\FunctionTok{write\_stan\_file}\NormalTok{(ThresholdAgent\_stan,}
                \AttributeTok{dir =} \StringTok{"stan/"}\NormalTok{,}
                \AttributeTok{basename =} \StringTok{"ch11\_threshold\_agent.stan"}\NormalTok{)}
\end{Highlighting}
\end{Shaded}

\begin{verbatim}
[1] "/Users/au209589/Dropbox/Teaching/AdvancedCognitiveModeling23_book/stan/ch11_threshold_agent.stan"
\end{verbatim}

\begin{Shaded}
\begin{Highlighting}[]
\NormalTok{mod\_A }\OtherTok{\textless{}{-}} \FunctionTok{cmdstan\_model}\NormalTok{(}\StringTok{"stan/ch11\_threshold\_agent.stan"}\NormalTok{, }\AttributeTok{dir =} \StringTok{"simmodels"}\NormalTok{)}

\ControlFlowTok{if}\NormalTok{ (regenerate\_simulations }\SpecialCharTok{||} \SpecialCharTok{!}\FunctionTok{file.exists}\NormalTok{(}\StringTok{"simmodels/ch11\_fit\_A.rds"}\NormalTok{)) \{}
\NormalTok{  fit\_A }\OtherTok{\textless{}{-}}\NormalTok{ mod\_A}\SpecialCharTok{$}\FunctionTok{sample}\NormalTok{(}
    \AttributeTok{data =}\NormalTok{ stan\_data\_A,}
    \AttributeTok{seed =} \DecValTok{123}\NormalTok{,}
    \AttributeTok{chains =} \DecValTok{2}\NormalTok{,}
    \AttributeTok{parallel\_chains =} \DecValTok{2}\NormalTok{,}
    \AttributeTok{iter\_warmup =} \DecValTok{500}\NormalTok{,}
    \AttributeTok{iter\_sampling =} \DecValTok{500}\NormalTok{,}
    \AttributeTok{refresh =} \DecValTok{0}
\NormalTok{  )}
\NormalTok{  fit\_A}\SpecialCharTok{$}\FunctionTok{save\_object}\NormalTok{(}\StringTok{"simmodels/ch11\_fit\_A.rds"}\NormalTok{)}
\NormalTok{\} }\ControlFlowTok{else}\NormalTok{ \{}
\NormalTok{  fit\_A }\OtherTok{\textless{}{-}} \FunctionTok{readRDS}\NormalTok{(}\StringTok{"simmodels/ch11\_fit\_A.rds"}\NormalTok{)}
\NormalTok{\}}
\end{Highlighting}
\end{Shaded}

\begin{Shaded}
\begin{Highlighting}[]
\NormalTok{fit\_A}\SpecialCharTok{$}\FunctionTok{summary}\NormalTok{(}\StringTok{"tau"}\NormalTok{)}
\end{Highlighting}
\end{Shaded}

\begin{verbatim}
# A tibble: 1 x 10
  variable  mean median      sd     mad    q5   q95  rhat ess_bulk ess_tail
  <chr>    <dbl>  <dbl>   <dbl>   <dbl> <dbl> <dbl> <dbl>    <dbl>    <dbl>
1 tau      0.241  0.241 0.00228 0.00218 0.237 0.244 1.000     388.     474.
\end{verbatim}

\begin{Shaded}
\begin{Highlighting}[]
\FunctionTok{mcmc\_trace}\NormalTok{(fit\_A}\SpecialCharTok{$}\FunctionTok{draws}\NormalTok{(}\StringTok{"tau"}\NormalTok{))}
\end{Highlighting}
\end{Shaded}

\begin{center}
\includegraphics[width=0.8\linewidth,height=\textheight,keepaspectratio]{10a-BayesianSequentialExtensions_files/figure-pdf/ch11_threshold_diagnostics-1.pdf}
\end{center}

The posterior should concentrate near the simulated \(\tau = 0.25\) ---
but with meaningful width, because the soft rule trades off slope
(\(\sigma_\tau\)) against threshold location. The hard-rule
identifiability limit is an interesting project question.

\section{Section B: Adaptive-Weight
Agent}\label{section-b-adaptive-weight-agent}

\subsection{Conceptual motivation}\label{conceptual-motivation-1}

The WBA of Ch. 10 commits at \(t = 0\) to fixed weights \(w_d\) and
\(w_s\) on direct and social evidence, and carries them through the
entire experiment. But a sensible agent should update those weights: if
social tips keep steering them wrong, they should trust the social
source \emph{less}; if direct observations keep being misleading (a
noisy sensor, say), they should trust themselves less. This turns the
static WBA into a \emph{learning} model, and makes contact with the
reinforcement-learning framing of Ch. 3 --- except that the thing being
learned is not a value, but a meta-parameter governing how two
information streams are combined. The resulting agent is Bayesian
\emph{within} a trial (it still integrates evidence via the
Beta-Binomial posterior) but RL-like \emph{across} trials (it updates
weights from prediction errors).

This extension also provides a mechanistic story for why real agents
might look like different static WBAs at different points in an
experiment: they are one agent whose weights have drifted.

\subsection{Mathematical
formalization}\label{mathematical-formalization-1}

On each trial \(t\), the two sources make independent predictions of the
correct colour:

\[p_{d,t} = \frac{\alpha_{d,t}}{\alpha_{d,t} + \beta_{d,t}},
\qquad
p_{s,t} = \frac{\alpha_{s,t}}{\alpha_{s,t} + \beta_{s,t}},\]

where \((\alpha_{d,t}, \beta_{d,t})\) is the posterior from \emph{only}
the direct evidence on trial \(t\), and \((\alpha_{s,t}, \beta_{s,t})\)
is the posterior from \emph{only} the social evidence. The agent's
combined posterior (and its guess) uses the WBA combination with
\emph{current} weights \(w_{d,t}, w_{s,t}\):

\[\alpha_t = 0.5 + w_{d,t} k_{d,t} + w_{s,t} k_{s,t}, \qquad
  \beta_t  = 0.5 + w_{d,t} (n_{d,t} - k_{d,t}) + w_{s,t} (n_{s,t} - k_{s,t}).\]

After feedback \(y_t \in \{0, 1\}\) on the correct colour, each source's
per-trial prediction error is

\[\delta_{d,t} = y_t - p_{d,t}, \qquad \delta_{s,t} = y_t - p_{s,t}.\]

The weights update on the log scale (to keep them positive) by a delta
rule that rewards the source that predicted better:

\[\log w_{d,t+1} = \log w_{d,t} + \eta\,(|\delta_{s,t}| - |\delta_{d,t}|),\]
\[\log w_{s,t+1} = \log w_{s,t} + \eta\,(|\delta_{d,t}| - |\delta_{s,t}|).\]

Parameters: learning rate \(\eta \in (0, 1)\) and initial log-weights
\(\log w_{d,0}, \log w_{s,0}\).

\subsection{R forward simulation}\label{r-forward-simulation-1}

\begin{Shaded}
\begin{Highlighting}[]
\CommentTok{\# Generate a stream of trials with specified direct/social reliability.}
\CommentTok{\# direct\_reliability / social\_reliability: probability that the source}
\CommentTok{\# points to the correct colour on any given trial.}
\NormalTok{make\_trial\_stream }\OtherTok{\textless{}{-}} \ControlFlowTok{function}\NormalTok{(n\_trials, direct\_reliability,}
\NormalTok{                              social\_reliability, }\AttributeTok{n\_d =} \DecValTok{4}\NormalTok{, }\AttributeTok{n\_s =} \DecValTok{3}\NormalTok{) \{}
\NormalTok{  purrr}\SpecialCharTok{::}\FunctionTok{map\_dfr}\NormalTok{(}\FunctionTok{seq\_len}\NormalTok{(n\_trials), }\ControlFlowTok{function}\NormalTok{(i) \{}
\NormalTok{    correct }\OtherTok{\textless{}{-}} \FunctionTok{rbinom}\NormalTok{(}\DecValTok{1}\NormalTok{, }\DecValTok{1}\NormalTok{, }\FloatTok{0.5}\NormalTok{)  }\CommentTok{\# truth randomly blue/red}
    \CommentTok{\# Each direct draw agrees with the truth with prob direct\_reliability}
\NormalTok{    k\_d }\OtherTok{\textless{}{-}} \FunctionTok{rbinom}\NormalTok{(}\DecValTok{1}\NormalTok{, n\_d, }\ControlFlowTok{if}\NormalTok{ (correct }\SpecialCharTok{==} \DecValTok{1}\NormalTok{) direct\_reliability}
                          \ControlFlowTok{else} \DecValTok{1} \SpecialCharTok{{-}}\NormalTok{ direct\_reliability)}
\NormalTok{    k\_s }\OtherTok{\textless{}{-}} \FunctionTok{rbinom}\NormalTok{(}\DecValTok{1}\NormalTok{, n\_s, }\ControlFlowTok{if}\NormalTok{ (correct }\SpecialCharTok{==} \DecValTok{1}\NormalTok{) social\_reliability}
                          \ControlFlowTok{else} \DecValTok{1} \SpecialCharTok{{-}}\NormalTok{ social\_reliability)}
    \FunctionTok{tibble}\NormalTok{(}\AttributeTok{trial =}\NormalTok{ i, }\AttributeTok{correct =}\NormalTok{ correct,}
           \AttributeTok{k\_d =}\NormalTok{ k\_d, }\AttributeTok{n\_d =}\NormalTok{ n\_d, }\AttributeTok{k\_s =}\NormalTok{ k\_s, }\AttributeTok{n\_s =}\NormalTok{ n\_s)}
\NormalTok{  \})}
\NormalTok{\}}

\NormalTok{simulate\_adaptive\_agent }\OtherTok{\textless{}{-}} \ControlFlowTok{function}\NormalTok{(eta, w\_d0, w\_s0, trial\_stream,}
                                    \AttributeTok{alpha0 =} \FloatTok{0.5}\NormalTok{, }\AttributeTok{beta0 =} \FloatTok{0.5}\NormalTok{) \{}
\NormalTok{  w\_d }\OtherTok{\textless{}{-}}\NormalTok{ w\_d0}
\NormalTok{  w\_s }\OtherTok{\textless{}{-}}\NormalTok{ w\_s0}
\NormalTok{  out }\OtherTok{\textless{}{-}} \FunctionTok{vector}\NormalTok{(}\StringTok{"list"}\NormalTok{, }\FunctionTok{nrow}\NormalTok{(trial\_stream))}
  \ControlFlowTok{for}\NormalTok{ (i }\ControlFlowTok{in} \FunctionTok{seq\_len}\NormalTok{(}\FunctionTok{nrow}\NormalTok{(trial\_stream))) \{}
\NormalTok{    tr }\OtherTok{\textless{}{-}}\NormalTok{ trial\_stream[i, ]}
    \CommentTok{\# Source{-}specific posteriors}
\NormalTok{    alpha\_d }\OtherTok{\textless{}{-}}\NormalTok{ alpha0 }\SpecialCharTok{+}\NormalTok{ tr}\SpecialCharTok{$}\NormalTok{k\_d}
\NormalTok{    beta\_d  }\OtherTok{\textless{}{-}}\NormalTok{ beta0  }\SpecialCharTok{+}\NormalTok{ (tr}\SpecialCharTok{$}\NormalTok{n\_d }\SpecialCharTok{{-}}\NormalTok{ tr}\SpecialCharTok{$}\NormalTok{k\_d)}
\NormalTok{    alpha\_s }\OtherTok{\textless{}{-}}\NormalTok{ alpha0 }\SpecialCharTok{+}\NormalTok{ tr}\SpecialCharTok{$}\NormalTok{k\_s}
\NormalTok{    beta\_s  }\OtherTok{\textless{}{-}}\NormalTok{ beta0  }\SpecialCharTok{+}\NormalTok{ (tr}\SpecialCharTok{$}\NormalTok{n\_s }\SpecialCharTok{{-}}\NormalTok{ tr}\SpecialCharTok{$}\NormalTok{k\_s)}
\NormalTok{    p\_d }\OtherTok{\textless{}{-}}\NormalTok{ alpha\_d }\SpecialCharTok{/}\NormalTok{ (alpha\_d }\SpecialCharTok{+}\NormalTok{ beta\_d)}
\NormalTok{    p\_s }\OtherTok{\textless{}{-}}\NormalTok{ alpha\_s }\SpecialCharTok{/}\NormalTok{ (alpha\_s }\SpecialCharTok{+}\NormalTok{ beta\_s)}

    \CommentTok{\# Combined WBA posterior with current weights}
\NormalTok{    alpha\_c }\OtherTok{\textless{}{-}}\NormalTok{ alpha0 }\SpecialCharTok{+}\NormalTok{ w\_d }\SpecialCharTok{*}\NormalTok{ tr}\SpecialCharTok{$}\NormalTok{k\_d }\SpecialCharTok{+}\NormalTok{ w\_s }\SpecialCharTok{*}\NormalTok{ tr}\SpecialCharTok{$}\NormalTok{k\_s}
\NormalTok{    beta\_c  }\OtherTok{\textless{}{-}}\NormalTok{ beta0  }\SpecialCharTok{+}\NormalTok{ w\_d }\SpecialCharTok{*}\NormalTok{ (tr}\SpecialCharTok{$}\NormalTok{n\_d }\SpecialCharTok{{-}}\NormalTok{ tr}\SpecialCharTok{$}\NormalTok{k\_d) }\SpecialCharTok{+}\NormalTok{ w\_s }\SpecialCharTok{*}\NormalTok{ (tr}\SpecialCharTok{$}\NormalTok{n\_s }\SpecialCharTok{{-}}\NormalTok{ tr}\SpecialCharTok{$}\NormalTok{k\_s)}
\NormalTok{    p\_c     }\OtherTok{\textless{}{-}}\NormalTok{ alpha\_c }\SpecialCharTok{/}\NormalTok{ (alpha\_c }\SpecialCharTok{+}\NormalTok{ beta\_c)}
\NormalTok{    guess   }\OtherTok{\textless{}{-}} \FunctionTok{rbinom}\NormalTok{(}\DecValTok{1}\NormalTok{, }\DecValTok{1}\NormalTok{, p\_c)}

    \CommentTok{\# Delta rule on log{-}weights}
\NormalTok{    y }\OtherTok{\textless{}{-}}\NormalTok{ tr}\SpecialCharTok{$}\NormalTok{correct}
\NormalTok{    d\_d }\OtherTok{\textless{}{-}} \FunctionTok{abs}\NormalTok{(y }\SpecialCharTok{{-}}\NormalTok{ p\_d)}
\NormalTok{    d\_s }\OtherTok{\textless{}{-}} \FunctionTok{abs}\NormalTok{(y }\SpecialCharTok{{-}}\NormalTok{ p\_s)}
\NormalTok{    w\_d }\OtherTok{\textless{}{-}} \FunctionTok{exp}\NormalTok{(}\FunctionTok{log}\NormalTok{(w\_d) }\SpecialCharTok{+}\NormalTok{ eta }\SpecialCharTok{*}\NormalTok{ (d\_s }\SpecialCharTok{{-}}\NormalTok{ d\_d))}
\NormalTok{    w\_s }\OtherTok{\textless{}{-}} \FunctionTok{exp}\NormalTok{(}\FunctionTok{log}\NormalTok{(w\_s) }\SpecialCharTok{+}\NormalTok{ eta }\SpecialCharTok{*}\NormalTok{ (d\_d }\SpecialCharTok{{-}}\NormalTok{ d\_s))}

\NormalTok{    out[[i]] }\OtherTok{\textless{}{-}} \FunctionTok{tibble}\NormalTok{(}\AttributeTok{trial =}\NormalTok{ tr}\SpecialCharTok{$}\NormalTok{trial, }\AttributeTok{correct =}\NormalTok{ y, }\AttributeTok{guess =}\NormalTok{ guess,}
                       \AttributeTok{p\_c =}\NormalTok{ p\_c, }\AttributeTok{w\_d =}\NormalTok{ w\_d, }\AttributeTok{w\_s =}\NormalTok{ w\_s,}
                       \AttributeTok{k\_d =}\NormalTok{ tr}\SpecialCharTok{$}\NormalTok{k\_d, }\AttributeTok{n\_d =}\NormalTok{ tr}\SpecialCharTok{$}\NormalTok{n\_d,}
                       \AttributeTok{k\_s =}\NormalTok{ tr}\SpecialCharTok{$}\NormalTok{k\_s, }\AttributeTok{n\_s =}\NormalTok{ tr}\SpecialCharTok{$}\NormalTok{n\_s)}
\NormalTok{  \}}
  \FunctionTok{bind\_rows}\NormalTok{(out)}
\NormalTok{\}}

\CommentTok{\# Three regimes}
\FunctionTok{set.seed}\NormalTok{(}\DecValTok{7}\NormalTok{)}
\NormalTok{regimes }\OtherTok{\textless{}{-}} \FunctionTok{tibble}\NormalTok{(}
  \AttributeTok{regime =} \FunctionTok{c}\NormalTok{(}\StringTok{"direct more reliable"}\NormalTok{, }\StringTok{"social more reliable"}\NormalTok{, }\StringTok{"equally reliable"}\NormalTok{),}
  \AttributeTok{rel\_d  =} \FunctionTok{c}\NormalTok{(}\FloatTok{0.85}\NormalTok{, }\FloatTok{0.55}\NormalTok{, }\FloatTok{0.75}\NormalTok{),}
  \AttributeTok{rel\_s  =} \FunctionTok{c}\NormalTok{(}\FloatTok{0.55}\NormalTok{, }\FloatTok{0.85}\NormalTok{, }\FloatTok{0.75}\NormalTok{)}
\NormalTok{)}

\NormalTok{sim\_adapt }\OtherTok{\textless{}{-}}\NormalTok{ purrr}\SpecialCharTok{::}\FunctionTok{pmap\_dfr}\NormalTok{(regimes, }\ControlFlowTok{function}\NormalTok{(regime, rel\_d, rel\_s) \{}
\NormalTok{  stream }\OtherTok{\textless{}{-}} \FunctionTok{make\_trial\_stream}\NormalTok{(}\AttributeTok{n\_trials =} \DecValTok{120}\NormalTok{,}
                              \AttributeTok{direct\_reliability =}\NormalTok{ rel\_d,}
                              \AttributeTok{social\_reliability =}\NormalTok{ rel\_s)}
  \FunctionTok{simulate\_adaptive\_agent}\NormalTok{(}\AttributeTok{eta =} \FloatTok{0.2}\NormalTok{, }\AttributeTok{w\_d0 =} \DecValTok{1}\NormalTok{, }\AttributeTok{w\_s0 =} \DecValTok{1}\NormalTok{, }\AttributeTok{trial\_stream =}\NormalTok{ stream) }\SpecialCharTok{|\textgreater{}}
    \FunctionTok{mutate}\NormalTok{(}\AttributeTok{regime =}\NormalTok{ regime)}
\NormalTok{\})}
\end{Highlighting}
\end{Shaded}

\subsection{Visualisations}\label{visualisations-1}

\begin{Shaded}
\begin{Highlighting}[]
\NormalTok{sim\_adapt }\SpecialCharTok{|\textgreater{}}
  \FunctionTok{pivot\_longer}\NormalTok{(}\FunctionTok{c}\NormalTok{(w\_d, w\_s), }\AttributeTok{names\_to =} \StringTok{"weight"}\NormalTok{, }\AttributeTok{values\_to =} \StringTok{"value"}\NormalTok{) }\SpecialCharTok{|\textgreater{}}
  \FunctionTok{ggplot}\NormalTok{(}\FunctionTok{aes}\NormalTok{(trial, value, }\AttributeTok{colour =}\NormalTok{ weight)) }\SpecialCharTok{+}
  \FunctionTok{geom\_line}\NormalTok{(}\AttributeTok{size =} \FloatTok{0.9}\NormalTok{) }\SpecialCharTok{+}
  \FunctionTok{facet\_wrap}\NormalTok{(}\SpecialCharTok{\textasciitilde{}}\NormalTok{regime) }\SpecialCharTok{+}
  \FunctionTok{labs}\NormalTok{(}\AttributeTok{title =} \StringTok{"(a) Weight trajectories by regime"}\NormalTok{,}
       \AttributeTok{y =} \StringTok{"Weight"}\NormalTok{, }\AttributeTok{x =} \StringTok{"Trial"}\NormalTok{)}
\end{Highlighting}
\end{Shaded}

\begin{center}
\includegraphics[width=0.8\linewidth,height=\textheight,keepaspectratio]{10a-BayesianSequentialExtensions_files/figure-pdf/ch11_adaptive_viz_a-1.pdf}
\end{center}

\begin{Shaded}
\begin{Highlighting}[]
\CommentTok{\# Compare adaptive agent vs. static WBA (w\_d = w\_s = 1) on accuracy}
\FunctionTok{set.seed}\NormalTok{(}\DecValTok{8}\NormalTok{)}
\NormalTok{compare\_df }\OtherTok{\textless{}{-}}\NormalTok{ purrr}\SpecialCharTok{::}\FunctionTok{pmap\_dfr}\NormalTok{(regimes, }\ControlFlowTok{function}\NormalTok{(regime, rel\_d, rel\_s) \{}
\NormalTok{  stream }\OtherTok{\textless{}{-}} \FunctionTok{make\_trial\_stream}\NormalTok{(}\AttributeTok{n\_trials =} \DecValTok{120}\NormalTok{,}
                              \AttributeTok{direct\_reliability =}\NormalTok{ rel\_d,}
                              \AttributeTok{social\_reliability =}\NormalTok{ rel\_s)}
\NormalTok{  adaptive }\OtherTok{\textless{}{-}} \FunctionTok{simulate\_adaptive\_agent}\NormalTok{(}\FloatTok{0.2}\NormalTok{, }\DecValTok{1}\NormalTok{, }\DecValTok{1}\NormalTok{, stream) }\SpecialCharTok{|\textgreater{}}
    \FunctionTok{mutate}\NormalTok{(}\AttributeTok{agent =} \StringTok{"adaptive"}\NormalTok{)}
\NormalTok{  static }\OtherTok{\textless{}{-}} \FunctionTok{simulate\_adaptive\_agent}\NormalTok{(}\FloatTok{0.0}\NormalTok{, }\DecValTok{1}\NormalTok{, }\DecValTok{1}\NormalTok{, stream) }\SpecialCharTok{|\textgreater{}}
    \FunctionTok{mutate}\NormalTok{(}\AttributeTok{agent =} \StringTok{"static WBA (eta = 0)"}\NormalTok{)}
  \FunctionTok{bind\_rows}\NormalTok{(adaptive, static) }\SpecialCharTok{|\textgreater{}} \FunctionTok{mutate}\NormalTok{(}\AttributeTok{regime =}\NormalTok{ regime)}
\NormalTok{\})}

\NormalTok{compare\_df }\SpecialCharTok{|\textgreater{}}
  \FunctionTok{mutate}\NormalTok{(}\AttributeTok{correct\_guess =} \FunctionTok{as.integer}\NormalTok{(guess }\SpecialCharTok{==}\NormalTok{ correct)) }\SpecialCharTok{|\textgreater{}}
  \FunctionTok{group\_by}\NormalTok{(regime, agent, trial) }\SpecialCharTok{|\textgreater{}}
  \FunctionTok{summarise}\NormalTok{(}\AttributeTok{acc =} \FunctionTok{mean}\NormalTok{(correct\_guess), }\AttributeTok{.groups =} \StringTok{"drop"}\NormalTok{) }\SpecialCharTok{|\textgreater{}}
  \FunctionTok{group\_by}\NormalTok{(regime, agent) }\SpecialCharTok{|\textgreater{}}
  \FunctionTok{mutate}\NormalTok{(}\AttributeTok{running\_acc =} \FunctionTok{cumsum}\NormalTok{(acc) }\SpecialCharTok{/} \FunctionTok{seq\_along}\NormalTok{(acc)) }\SpecialCharTok{|\textgreater{}}
  \FunctionTok{ggplot}\NormalTok{(}\FunctionTok{aes}\NormalTok{(trial, running\_acc, }\AttributeTok{colour =}\NormalTok{ agent)) }\SpecialCharTok{+}
  \FunctionTok{geom\_line}\NormalTok{(}\AttributeTok{size =} \FloatTok{0.9}\NormalTok{) }\SpecialCharTok{+}
  \FunctionTok{facet\_wrap}\NormalTok{(}\SpecialCharTok{\textasciitilde{}}\NormalTok{regime) }\SpecialCharTok{+}
  \FunctionTok{labs}\NormalTok{(}\AttributeTok{title =} \StringTok{"(b) Running accuracy: adaptive vs. static WBA"}\NormalTok{,}
       \AttributeTok{y =} \StringTok{"Cumulative accuracy"}\NormalTok{, }\AttributeTok{x =} \StringTok{"Trial"}\NormalTok{)}
\end{Highlighting}
\end{Shaded}

\begin{center}
\includegraphics[width=0.8\linewidth,height=\textheight,keepaspectratio]{10a-BayesianSequentialExtensions_files/figure-pdf/ch11_adaptive_viz_b-1.pdf}
\end{center}

\begin{Shaded}
\begin{Highlighting}[]
\CommentTok{\# Final{-}weight distribution across 50 simulated agents per regime}
\FunctionTok{set.seed}\NormalTok{(}\DecValTok{9}\NormalTok{)}
\NormalTok{final\_df }\OtherTok{\textless{}{-}}\NormalTok{ purrr}\SpecialCharTok{::}\FunctionTok{pmap\_dfr}\NormalTok{(regimes, }\ControlFlowTok{function}\NormalTok{(regime, rel\_d, rel\_s) \{}
\NormalTok{  purrr}\SpecialCharTok{::}\FunctionTok{map\_dfr}\NormalTok{(}\DecValTok{1}\SpecialCharTok{:}\DecValTok{50}\NormalTok{, }\ControlFlowTok{function}\NormalTok{(a) \{}
\NormalTok{    stream }\OtherTok{\textless{}{-}} \FunctionTok{make\_trial\_stream}\NormalTok{(}\DecValTok{120}\NormalTok{, rel\_d, rel\_s)}
\NormalTok{    res }\OtherTok{\textless{}{-}} \FunctionTok{simulate\_adaptive\_agent}\NormalTok{(}\FloatTok{0.2}\NormalTok{, }\DecValTok{1}\NormalTok{, }\DecValTok{1}\NormalTok{, stream)}
    \FunctionTok{tibble}\NormalTok{(}\AttributeTok{agent\_id =}\NormalTok{ a, }\AttributeTok{regime =}\NormalTok{ regime,}
           \AttributeTok{w\_d\_final =} \FunctionTok{tail}\NormalTok{(res}\SpecialCharTok{$}\NormalTok{w\_d, }\DecValTok{1}\NormalTok{),}
           \AttributeTok{w\_s\_final =} \FunctionTok{tail}\NormalTok{(res}\SpecialCharTok{$}\NormalTok{w\_s, }\DecValTok{1}\NormalTok{))}
\NormalTok{  \})}
\NormalTok{\})}

\NormalTok{final\_df }\SpecialCharTok{|\textgreater{}}
  \FunctionTok{pivot\_longer}\NormalTok{(}\FunctionTok{c}\NormalTok{(w\_d\_final, w\_s\_final), }\AttributeTok{names\_to =} \StringTok{"weight"}\NormalTok{, }\AttributeTok{values\_to =} \StringTok{"value"}\NormalTok{) }\SpecialCharTok{|\textgreater{}}
  \FunctionTok{ggplot}\NormalTok{(}\FunctionTok{aes}\NormalTok{(value, }\AttributeTok{fill =}\NormalTok{ weight)) }\SpecialCharTok{+}
  \FunctionTok{geom\_histogram}\NormalTok{(}\AttributeTok{bins =} \DecValTok{25}\NormalTok{, }\AttributeTok{alpha =} \FloatTok{0.7}\NormalTok{, }\AttributeTok{position =} \StringTok{"identity"}\NormalTok{, }\AttributeTok{colour =} \StringTok{"white"}\NormalTok{) }\SpecialCharTok{+}
  \FunctionTok{facet\_wrap}\NormalTok{(}\SpecialCharTok{\textasciitilde{}}\NormalTok{regime) }\SpecialCharTok{+}
  \FunctionTok{labs}\NormalTok{(}\AttributeTok{title =} \StringTok{"(c) Final{-}weight distribution (50 agents per regime)"}\NormalTok{,}
       \AttributeTok{x =} \StringTok{"Final weight"}\NormalTok{, }\AttributeTok{y =} \StringTok{"Count"}\NormalTok{)}
\end{Highlighting}
\end{Shaded}

\begin{center}
\includegraphics[width=0.8\linewidth,height=\textheight,keepaspectratio]{10a-BayesianSequentialExtensions_files/figure-pdf/ch11_adaptive_viz_c-1.pdf}
\end{center}

Across 50 simulated agents the weights track the true reliability
asymmetry: when direct evidence is more reliable, \(w_d\) grows and
\(w_s\) shrinks; the pattern reverses when social evidence is more
reliable; and when the two are matched, weights drift modestly in either
direction without systematic asymmetry.

\subsection{Single-agent Stan model}\label{single-agent-stan-model}

\begin{Shaded}
\begin{Highlighting}[]
\CommentTok{\# Fit to one direct{-}more{-}reliable agent}
\FunctionTok{set.seed}\NormalTok{(}\DecValTok{11}\NormalTok{)}
\NormalTok{stream\_fit }\OtherTok{\textless{}{-}} \FunctionTok{make\_trial\_stream}\NormalTok{(}\DecValTok{120}\NormalTok{, }\AttributeTok{direct\_reliability =} \FloatTok{0.85}\NormalTok{,}
                                 \AttributeTok{social\_reliability =} \FloatTok{0.55}\NormalTok{)}
\NormalTok{sim\_fit }\OtherTok{\textless{}{-}} \FunctionTok{simulate\_adaptive\_agent}\NormalTok{(}\FloatTok{0.2}\NormalTok{, }\DecValTok{1}\NormalTok{, }\DecValTok{1}\NormalTok{, stream\_fit)}

\NormalTok{stan\_data\_B }\OtherTok{\textless{}{-}} \FunctionTok{list}\NormalTok{(}
  \AttributeTok{N       =} \FunctionTok{nrow}\NormalTok{(sim\_fit),}
  \AttributeTok{guess   =}\NormalTok{ sim\_fit}\SpecialCharTok{$}\NormalTok{guess,}
  \AttributeTok{correct =}\NormalTok{ sim\_fit}\SpecialCharTok{$}\NormalTok{correct,}
  \AttributeTok{k\_d     =}\NormalTok{ sim\_fit}\SpecialCharTok{$}\NormalTok{k\_d,}
  \AttributeTok{n\_d     =}\NormalTok{ sim\_fit}\SpecialCharTok{$}\NormalTok{n\_d,}
  \AttributeTok{k\_s     =}\NormalTok{ sim\_fit}\SpecialCharTok{$}\NormalTok{k\_s,}
  \AttributeTok{n\_s     =}\NormalTok{ sim\_fit}\SpecialCharTok{$}\NormalTok{n\_s}
\NormalTok{)}
\end{Highlighting}
\end{Shaded}

\begin{Shaded}
\begin{Highlighting}[]
\NormalTok{AdaptiveAgent\_stan }\OtherTok{\textless{}{-}} \StringTok{"}
\StringTok{// Adaptive{-}Weight Agent.}
\StringTok{// Parameters: eta (learning rate), log w\_d0, log w\_s0 (initial log{-}weights).}
\StringTok{// Weights are updated across trials by a delta rule on log{-}weights using}
\StringTok{// per{-}source absolute prediction errors against the feedback.}

\StringTok{data \{}
\StringTok{  int\textless{}lower=1\textgreater{} N;}
\StringTok{  array[N] int\textless{}lower=0, upper=1\textgreater{} guess;}
\StringTok{  array[N] int\textless{}lower=0, upper=1\textgreater{} correct;}
\StringTok{  array[N] int\textless{}lower=0\textgreater{} k\_d;}
\StringTok{  array[N] int\textless{}lower=0\textgreater{} n\_d;}
\StringTok{  array[N] int\textless{}lower=0\textgreater{} k\_s;}
\StringTok{  array[N] int\textless{}lower=0\textgreater{} n\_s;}
\StringTok{\}}

\StringTok{parameters \{}
\StringTok{  real\textless{}lower=0, upper=1\textgreater{} eta;}
\StringTok{  real log\_w\_d0;}
\StringTok{  real log\_w\_s0;}
\StringTok{\}}

\StringTok{transformed parameters \{}
\StringTok{  vector[N] p\_c;}
\StringTok{  \{}
\StringTok{    real log\_w\_d = log\_w\_d0;}
\StringTok{    real log\_w\_s = log\_w\_s0;}
\StringTok{    for (i in 1:N) \{}
\StringTok{      real w\_d = exp(log\_w\_d);}
\StringTok{      real w\_s = exp(log\_w\_s);}

\StringTok{      real alpha\_d = 0.5 + k\_d[i];}
\StringTok{      real beta\_d  = 0.5 + (n\_d[i] {-} k\_d[i]);}
\StringTok{      real alpha\_s = 0.5 + k\_s[i];}
\StringTok{      real beta\_s  = 0.5 + (n\_s[i] {-} k\_s[i]);}
\StringTok{      real p\_d = alpha\_d / (alpha\_d + beta\_d);}
\StringTok{      real p\_s = alpha\_s / (alpha\_s + beta\_s);}

\StringTok{      real alpha\_c = 0.5 + w\_d * k\_d[i] + w\_s * k\_s[i];}
\StringTok{      real beta\_c  = 0.5 + w\_d * (n\_d[i] {-} k\_d[i]) + w\_s * (n\_s[i] {-} k\_s[i]);}
\StringTok{      p\_c[i] = alpha\_c / (alpha\_c + beta\_c);}

\StringTok{      real d\_d = abs(correct[i] {-} p\_d);}
\StringTok{      real d\_s = abs(correct[i] {-} p\_s);}
\StringTok{      log\_w\_d += eta * (d\_s {-} d\_d);}
\StringTok{      log\_w\_s += eta * (d\_d {-} d\_s);}
\StringTok{    \}}
\StringTok{  \}}
\StringTok{\}}

\StringTok{model \{}
\StringTok{  target += beta\_lpdf(eta | 2, 5);}
\StringTok{  target += normal\_lpdf(log\_w\_d0 | 0, 1);}
\StringTok{  target += normal\_lpdf(log\_w\_s0 | 0, 1);}

\StringTok{  target += bernoulli\_lpmf(guess | p\_c);}
\StringTok{\}}
\StringTok{"}

\FunctionTok{write\_stan\_file}\NormalTok{(AdaptiveAgent\_stan,}
                \AttributeTok{dir =} \StringTok{"stan/"}\NormalTok{,}
                \AttributeTok{basename =} \StringTok{"ch11\_adaptive\_agent.stan"}\NormalTok{)}
\end{Highlighting}
\end{Shaded}

\begin{verbatim}
[1] "/Users/au209589/Dropbox/Teaching/AdvancedCognitiveModeling23_book/stan/ch11_adaptive_agent.stan"
\end{verbatim}

\begin{Shaded}
\begin{Highlighting}[]
\NormalTok{mod\_B }\OtherTok{\textless{}{-}} \FunctionTok{cmdstan\_model}\NormalTok{(}\StringTok{"stan/ch11\_adaptive\_agent.stan"}\NormalTok{, }\AttributeTok{dir =} \StringTok{"simmodels"}\NormalTok{)}

\ControlFlowTok{if}\NormalTok{ (regenerate\_simulations }\SpecialCharTok{||} \SpecialCharTok{!}\FunctionTok{file.exists}\NormalTok{(}\StringTok{"simmodels/ch11\_fit\_B.rds"}\NormalTok{)) \{}
\NormalTok{  fit\_B }\OtherTok{\textless{}{-}}\NormalTok{ mod\_B}\SpecialCharTok{$}\FunctionTok{sample}\NormalTok{(}
    \AttributeTok{data =}\NormalTok{ stan\_data\_B,}
    \AttributeTok{seed =} \DecValTok{321}\NormalTok{,}
    \AttributeTok{chains =} \DecValTok{2}\NormalTok{,}
    \AttributeTok{parallel\_chains =} \DecValTok{2}\NormalTok{,}
    \AttributeTok{iter\_warmup =} \DecValTok{500}\NormalTok{,}
    \AttributeTok{iter\_sampling =} \DecValTok{500}\NormalTok{,}
    \AttributeTok{refresh =} \DecValTok{0}
\NormalTok{  )}
\NormalTok{  fit\_B}\SpecialCharTok{$}\FunctionTok{save\_object}\NormalTok{(}\StringTok{"simmodels/ch11\_fit\_B.rds"}\NormalTok{)}
\NormalTok{\} }\ControlFlowTok{else}\NormalTok{ \{}
\NormalTok{  fit\_B }\OtherTok{\textless{}{-}} \FunctionTok{readRDS}\NormalTok{(}\StringTok{"simmodels/ch11\_fit\_B.rds"}\NormalTok{)}
\NormalTok{\}}
\end{Highlighting}
\end{Shaded}

\begin{Shaded}
\begin{Highlighting}[]
\NormalTok{fit\_B}\SpecialCharTok{$}\FunctionTok{summary}\NormalTok{(}\FunctionTok{c}\NormalTok{(}\StringTok{"eta"}\NormalTok{, }\StringTok{"log\_w\_d0"}\NormalTok{, }\StringTok{"log\_w\_s0"}\NormalTok{))}
\end{Highlighting}
\end{Shaded}

\begin{verbatim}
# A tibble: 3 x 10
  variable   mean median    sd    mad      q5   q95  rhat ess_bulk ess_tail
  <chr>     <dbl>  <dbl> <dbl>  <dbl>   <dbl> <dbl> <dbl>    <dbl>    <dbl>
1 eta       0.227  0.213 0.104 0.0989  0.0813 0.419 1.00      770.     649.
2 log_w_d0  0.181  0.182 0.812 0.810  -1.13   1.54  1.00      581.     571.
3 log_w_s0 -0.233 -0.206 0.939 0.951  -1.80   1.25  1.000     624.     690.
\end{verbatim}

\begin{Shaded}
\begin{Highlighting}[]
\FunctionTok{mcmc\_trace}\NormalTok{(fit\_B}\SpecialCharTok{$}\FunctionTok{draws}\NormalTok{(}\FunctionTok{c}\NormalTok{(}\StringTok{"eta"}\NormalTok{, }\StringTok{"log\_w\_d0"}\NormalTok{, }\StringTok{"log\_w\_s0"}\NormalTok{)))}
\end{Highlighting}
\end{Shaded}

\begin{center}
\includegraphics[width=0.8\linewidth,height=\textheight,keepaspectratio]{10a-BayesianSequentialExtensions_files/figure-pdf/ch11_adaptive_diagnostics-1.pdf}
\end{center}

\begin{Shaded}
\begin{Highlighting}[]
\FunctionTok{mcmc\_pairs}\NormalTok{(fit\_B}\SpecialCharTok{$}\FunctionTok{draws}\NormalTok{(}\FunctionTok{c}\NormalTok{(}\StringTok{"eta"}\NormalTok{, }\StringTok{"log\_w\_d0"}\NormalTok{, }\StringTok{"log\_w\_s0"}\NormalTok{)))}
\end{Highlighting}
\end{Shaded}

\begin{center}
\includegraphics[width=0.8\linewidth,height=\textheight,keepaspectratio]{10a-BayesianSequentialExtensions_files/figure-pdf/ch11_adaptive_diagnostics-2.pdf}
\end{center}

A note on what you will almost certainly see in the pairs plot: a strong
ridge between \(\eta\) and the initial log-weights. This is not a bug
--- it is structural. A small \(\eta\) with strongly asymmetric starting
weights produces similar behaviour to a large \(\eta\) with symmetric
starts, because both paths arrive at similar \emph{effective} weights by
the end of the trial stream. We flag this rather than hide it:
demonstrating and resolving the non-identifiability (via informative
priors, via parameter-informative experimental designs that force the
weights to change across blocks, or via multilevel partial pooling) is a
natural student project.

\section{Closing}\label{closing}

The Bayesian-cognition hypothesis is a family of models, not a single
commitment. Once an agent represents beliefs as probability
distributions and integrates evidence by Bayes' rule, many further
modelling choices remain open: when to stop sampling, how to weight
sources, whether those weights themselves are learned, what the
observation model looks like. The sample-or-guess and adaptive-weight
agents above are natural sequential extensions of the Ch. 10 family and
make good starting points for student projects --- a full validation
battery (parameter recovery, SBC, model comparison, multilevel pooling)
on either model is left as an exercise in applying the workflow you have
already learned.

\chapter{Stan Optimization Cheatsheet}\label{sec-stan-cheatsheet}

This appendix collects every Stan programming technique introduced
across Chapters 4--11, with minimal code examples, cost estimates, and
guidance on when to apply each one. The techniques are ordered by first
appearance in the book.

\textbf{Quick navigation}

\begin{longtable}[]{@{}
  >{\raggedright\arraybackslash}p{(\linewidth - 4\tabcolsep) * \real{0.3333}}
  >{\raggedright\arraybackslash}p{(\linewidth - 4\tabcolsep) * \real{0.3333}}
  >{\raggedright\arraybackslash}p{(\linewidth - 4\tabcolsep) * \real{0.3333}}@{}}
\toprule\noalign{}
\begin{minipage}[b]{\linewidth}\raggedright
\#
\end{minipage} & \begin{minipage}[b]{\linewidth}\raggedright
Technique
\end{minipage} & \begin{minipage}[b]{\linewidth}\raggedright
Chapter
\end{minipage} \\
\midrule\noalign{}
\endhead
\bottomrule\noalign{}
\endlastfoot
1 & \hyperref[sec-block-placement]{Block placement} & Ch 4 \\
2 & \hyperref[sec-precompute-r]{Pre-compute in R} & Ch 4, 6 \\
3 & \hyperref[sec-sufficient-stats]{Sufficient statistics} & Ch 4 \\
4 & \hyperref[sec-vectorize]{Vectorize the likelihood} & Ch 4, 6 \\
5 & \hyperref[sec-priors-as-data]{Priors as data} & Ch 5 \\
6 & \hyperref[sec-run-diagnostics]{Conditional diagnostics
(\texttt{run\_diagnostics})} & Ch 5 \\
7 & \hyperref[sec-generate-quantities]{Post-hoc generated quantities} &
Ch 5 \\
8 & \hyperref[sec-ncp]{Non-centered parameterization} & Ch 6 \\
9 & \hyperref[sec-cholesky]{Cholesky factorization} & Ch 6 \\
10 & \hyperref[sec-subject-indexing]{Contiguous subject indexing} & Ch
6 \\
11 & \hyperref[sec-pathfinder]{Pathfinder warm-start} & Ch 4, 6, 11 \\
12 & \hyperref[sec-log-mix]{Marginalizing discrete parameters} & Ch 8 \\
13 & \hyperref[sec-sequential-state]{Sequential state unrolling} & Ch
9 \\
14 & \hyperref[sec-reduce-sum]{\texttt{reduce\_sum} parallelism} & Ch
11 \\
15 & \hyperref[sec-compound-speedups]{Combining techniques} & Ch 6,
11 \\
16 & \hyperref[sec-sampler-options]{Sampler and compiler options} & App
B \\
\end{longtable}

\begin{center}\rule{0.5\linewidth}{0.5pt}\end{center}

\section{1. Block Placement --- Run It Once, Not Thousands of
Times}\label{sec-block-placement}

\emph{First introduced: Section~\ref{sec-memory-models} (Chapter 4,
Memory Model 1)}

HMC runs thousands of leapfrog steps during warm-up. Every line of code
in \texttt{transformed\ parameters} or \texttt{model} executes at each
step; every line in \texttt{transformed\ data} executes exactly
\textbf{once}.

\begin{longtable}[]{@{}
  >{\raggedright\arraybackslash}p{(\linewidth - 4\tabcolsep) * \real{0.3333}}
  >{\raggedright\arraybackslash}p{(\linewidth - 4\tabcolsep) * \real{0.3333}}
  >{\raggedright\arraybackslash}p{(\linewidth - 4\tabcolsep) * \real{0.3333}}@{}}
\toprule\noalign{}
\begin{minipage}[b]{\linewidth}\raggedright
Block
\end{minipage} & \begin{minipage}[b]{\linewidth}\raggedright
Runs
\end{minipage} & \begin{minipage}[b]{\linewidth}\raggedright
Put here
\end{minipage} \\
\midrule\noalign{}
\endhead
\bottomrule\noalign{}
\endlastfoot
\texttt{transformed\ data} & Once, before sampling & Quantities derived
\textbf{only from data} --- precomputed distances, log-transforms of
fixed inputs, cumulative sums \\
\texttt{transformed\ parameters} & Every leapfrog step \emph{(saved)} &
Deterministic functions of \textbf{parameters} whose gradients must
propagate --- NCP reparameterisations, recursive state, predicted values
used in the likelihood \\
\texttt{model} (local vars) & Every leapfrog step \emph{(not saved)} &
Intermediate scratch quantities not needed outside the loop \\
\texttt{generated\ quantities} & Once per saved draw & Posterior
predictive samples, \texttt{log\_lik}, prior draws --- anything
expensive that is \textbf{not} needed for sampling \\
\end{longtable}

\begin{tcolorbox}[enhanced jigsaw, toprule=.15mm, toptitle=1mm, bottomtitle=1mm, leftrule=.75mm, rightrule=.15mm, colframe=quarto-callout-tip-color-frame, left=2mm, colback=white, titlerule=0mm, coltitle=black, colbacktitle=quarto-callout-tip-color!10!white, opacityback=0, title=\textcolor{quarto-callout-tip-color}{\faLightbulb}\hspace{0.5em}{Cost of misplacement}, arc=.35mm, opacitybacktitle=0.6, breakable, bottomrule=.15mm]

A typical fit runs 4 chains × 1 000 warmup iterations × \textasciitilde6
leapfrog steps = \textbf{\textasciitilde24 000 leapfrog evaluations}
during warm-up. Moving 500 precomputed memory values from
\texttt{transformed\ parameters} to \texttt{transformed\ data} saves
\textasciitilde12 million floating-point operations per fit.

\end{tcolorbox}

\begin{Shaded}
\begin{Highlighting}[]
\KeywordTok{transformed data}\NormalTok{ \{}
  \CommentTok{// computed once — running average of opponent choices}
  \DataTypeTok{vector}\NormalTok{[N] memory;}
\NormalTok{  memory[}\DecValTok{1}\NormalTok{] = }\FloatTok{0.5}\NormalTok{;}
  \ControlFlowTok{for}\NormalTok{ (t }\ControlFlowTok{in} \DecValTok{2}\NormalTok{:N)}
\NormalTok{    memory[t] = (}\FloatTok{1.0}\NormalTok{ / t) * sum(other[}\DecValTok{1}\NormalTok{:(t}\DecValTok{{-}1}\NormalTok{)]);}
\NormalTok{\}}
\end{Highlighting}
\end{Shaded}

\begin{center}\rule{0.5\linewidth}{0.5pt}\end{center}

\section{2. Pre-Compute in R --- Don't Let Stan See the Raw
Computation}\label{sec-precompute-r}

\emph{First introduced: Chapter~\ref{sec-multilevel-models} (Chapter 6,
multilevel memory agent)}

\texttt{transformed\ data} runs once before sampling --- but the
computation still happens inside Stan's runtime. If a quantity depends
only on observed data and will never change between fits, computing it
in R and passing it as a data field is even cheaper: Stan never
processes the raw inputs at all.

\begin{Shaded}
\begin{Highlighting}[]
\CommentTok{\# R: compute the running opponent rate once (O(N) vectorised R operation)}
\NormalTok{d }\OtherTok{\textless{}{-}}\NormalTok{ d }\SpecialCharTok{|\textgreater{}}
  \FunctionTok{group\_by}\NormalTok{(subject) }\SpecialCharTok{|\textgreater{}}
  \FunctionTok{mutate}\NormalTok{(}\AttributeTok{opp\_rate\_prev =} \FunctionTok{lag}\NormalTok{(}\FunctionTok{cummean}\NormalTok{(other\_choice), }\AttributeTok{default =} \FloatTok{0.5}\NormalTok{)) }\SpecialCharTok{|\textgreater{}}
  \FunctionTok{ungroup}\NormalTok{()}

\NormalTok{stan\_data }\OtherTok{\textless{}{-}} \FunctionTok{list}\NormalTok{(}
  \AttributeTok{N\_total      =} \FunctionTok{nrow}\NormalTok{(d),}
  \AttributeTok{y            =}\NormalTok{ d}\SpecialCharTok{$}\NormalTok{choice,}
  \AttributeTok{opp\_rate\_prev =}\NormalTok{ d}\SpecialCharTok{$}\NormalTok{opp\_rate\_prev   }\CommentTok{\# passed as a plain vector}
  \CommentTok{\# ...}
\NormalTok{)}
\end{Highlighting}
\end{Shaded}

\begin{Shaded}
\begin{Highlighting}[]
\CommentTok{// Stan only receives the pre{-}computed vector — no loop needed}
\KeywordTok{data}\NormalTok{ \{}
  \DataTypeTok{vector}\NormalTok{\textless{}}\KeywordTok{lower}\NormalTok{=}\DecValTok{0}\NormalTok{,}\KeywordTok{upper}\NormalTok{=}\DecValTok{1}\NormalTok{\textgreater{}[N\_total] opp\_rate\_prev;}
\NormalTok{\}}
\KeywordTok{model}\NormalTok{ \{}
  \KeywordTok{target +=}\NormalTok{ bernoulli\_logit\_lpmf(y | mu + beta * logit(opp\_rate\_prev));}
\NormalTok{\}}
\end{Highlighting}
\end{Shaded}

\textbf{Decision rule:} If the computation involves only observed data
and uses fast R/dplyr operations, do it in R. Reserve
\texttt{transformed\ data} for cases where the computation is cleaner to
express in Stan's syntax (e.g., distance matrices over stimuli).

\begin{center}\rule{0.5\linewidth}{0.5pt}\end{center}

\section{3. Sufficient Statistics --- Aggregate IID
Trials}\label{sec-sufficient-stats}

\emph{First introduced: Chapter~\ref{sec-inferring-rates} (Chapter 4,
simple bias model)}

When all trials are \textbf{independently and identically distributed}
under the same parameter, the total number of successes is a
\textbf{sufficient statistic} --- it carries all the information in the
data about \(\theta\). A single Binomial call replaces \(N\) Bernoulli
calls.

\begin{Shaded}
\begin{Highlighting}[]
\CommentTok{// data block: pass the aggregate, not the full array}
\KeywordTok{data}\NormalTok{ \{}
  \DataTypeTok{int}\NormalTok{\textless{}}\KeywordTok{lower}\NormalTok{=}\DecValTok{0}\NormalTok{\textgreater{} N;}
  \DataTypeTok{int}\NormalTok{\textless{}}\KeywordTok{lower}\NormalTok{=}\DecValTok{0}\NormalTok{, }\KeywordTok{upper}\NormalTok{=N\textgreater{} successes;  }\CommentTok{// sum(h)}
\NormalTok{\}}
\KeywordTok{model}\NormalTok{ \{}
  \KeywordTok{target +=}\NormalTok{ binomial\_logit\_lpmf(successes | N, theta\_logit);}
\NormalTok{\}}
\end{Highlighting}
\end{Shaded}

\begin{Shaded}
\begin{Highlighting}[]
\NormalTok{stan\_data }\OtherTok{\textless{}{-}} \FunctionTok{list}\NormalTok{(}\AttributeTok{N =} \DecValTok{120}\NormalTok{, }\AttributeTok{successes =} \FunctionTok{sum}\NormalTok{(h))}
\end{Highlighting}
\end{Shaded}

\begin{tcolorbox}[enhanced jigsaw, toprule=.15mm, toptitle=1mm, bottomtitle=1mm, leftrule=.75mm, rightrule=.15mm, colframe=quarto-callout-tip-color-frame, left=2mm, colback=white, titlerule=0mm, coltitle=black, colbacktitle=quarto-callout-tip-color!10!white, opacityback=0, title=\textcolor{quarto-callout-tip-color}{\faLightbulb}\hspace{0.5em}{Speedup}, arc=.35mm, opacitybacktitle=0.6, breakable, bottomrule=.15mm]

One Binomial evaluation vs.~120 Bernoulli evaluations --- every leapfrog
step, every chain. For a single-subject model with 4 chains × 24 000
leapfrog evaluations the difference is small in absolute time, but the
principle scales: for 10 000 IID trials the saving is real.

\end{tcolorbox}

\textbf{When it applies:} single-subject, single-session models where
\(\theta\) is constant across all \(N\) trials (Chapters 4--5 bias
models).

\textbf{When it does not apply:} memory models, multilevel models
(per-subject \(\theta_j\)), any model where the trial-level likelihood
depends on a trial-varying predictor.

\begin{center}\rule{0.5\linewidth}{0.5pt}\end{center}

\section{4. Vectorize the Likelihood Call}\label{sec-vectorize}

\emph{First introduced: Chapter~\ref{sec-inferring-rates} (Chapter 4);
used throughout from Chapter 5 onward}

Stan evaluates each distribution call as a node in the autodiff graph. A
vectorized call over a length-\(N\) vector creates \textbf{one node}; a
loop over \(N\) trials creates \textbf{\(N\) nodes}. Smaller graph =
faster gradient computation and better cache use.

\begin{Shaded}
\begin{Highlighting}[]
\CommentTok{// Slower: N nodes in the autodiff graph}
\ControlFlowTok{for}\NormalTok{ (i }\ControlFlowTok{in} \DecValTok{1}\NormalTok{:N)}
  \KeywordTok{target +=}\NormalTok{ bernoulli\_logit\_lpmf(h[i] | theta\_logit);}

\CommentTok{// Faster: one node — mathematically identical}
\KeywordTok{target +=}\NormalTok{ bernoulli\_logit\_lpmf(h | theta\_logit);}
\end{Highlighting}
\end{Shaded}

In multilevel models, scatter-indexing by subject makes the entire trial
pool vectorizable:

\begin{Shaded}
\begin{Highlighting}[]
\CommentTok{// N\_total trials, J subjects — one call instead of N\_total calls}
\KeywordTok{target +=}\NormalTok{ bernoulli\_logit\_lpmf(y | theta\_logit[subj\_id]);}
\end{Highlighting}
\end{Shaded}

\begin{tcolorbox}[enhanced jigsaw, toprule=.15mm, toptitle=1mm, bottomtitle=1mm, leftrule=.75mm, rightrule=.15mm, colframe=quarto-callout-note-color-frame, left=2mm, colback=white, titlerule=0mm, coltitle=black, colbacktitle=quarto-callout-note-color!10!white, opacityback=0, title=\textcolor{quarto-callout-note-color}{\faInfo}\hspace{0.5em}{When vectorization is not available}, arc=.35mm, opacitybacktitle=0.6, breakable, bottomrule=.15mm]

\texttt{log\_mix()} for mixture likelihoods (Chapter 8) operates on
\textbf{one observation at a time} because each call mixes two
observation-specific log-probabilities. There is no vectorized
\texttt{log\_mix} --- the per-observation loop is correct and
unavoidable there.

\end{tcolorbox}

\begin{center}\rule{0.5\linewidth}{0.5pt}\end{center}

\section{5. Priors as Data --- Avoid
Recompilation}\label{sec-priors-as-data}

\emph{First introduced: Chapter~\ref{sec-model-quality-checks} (Chapter
5)}

Pass prior hyperparameters through the \texttt{data} block instead of
hard-coding them. Stan recompiles whenever the \textbf{code} changes,
never for new \textbf{data}. This lets you run prior sensitivity
analyses and SBC sweeps without recompiling.

\begin{Shaded}
\begin{Highlighting}[]
\KeywordTok{data}\NormalTok{ \{}
  \DataTypeTok{int}\NormalTok{\textless{}}\KeywordTok{lower}\NormalTok{=}\DecValTok{0}\NormalTok{\textgreater{} N;}
  \DataTypeTok{array}\NormalTok{[N] }\DataTypeTok{int}\NormalTok{ y;}
  \DataTypeTok{real}\NormalTok{             prior\_mu\_mean;}
  \DataTypeTok{real}\NormalTok{\textless{}}\KeywordTok{lower}\NormalTok{=}\DecValTok{0}\NormalTok{\textgreater{}    prior\_mu\_sigma;}
\NormalTok{\}}
\KeywordTok{model}\NormalTok{ \{}
\NormalTok{  mu \textasciitilde{} normal(prior\_mu\_mean, prior\_mu\_sigma);}
\NormalTok{  y  \textasciitilde{} bernoulli\_logit(mu);}
\NormalTok{\}}
\end{Highlighting}
\end{Shaded}

\begin{Shaded}
\begin{Highlighting}[]
\NormalTok{data }\OtherTok{\textless{}{-}} \FunctionTok{list}\NormalTok{(}\AttributeTok{N =}\NormalTok{ N, }\AttributeTok{y =}\NormalTok{ y,}
             \AttributeTok{prior\_mu\_mean =} \DecValTok{0}\NormalTok{, }\AttributeTok{prior\_mu\_sigma =} \FloatTok{1.5}\NormalTok{)}
\end{Highlighting}
\end{Shaded}

\begin{center}\rule{0.5\linewidth}{0.5pt}\end{center}

\section{\texorpdfstring{6. The \texttt{run\_diagnostics} Flag ---
Conditional Generated
Quantities}{6. The run\_diagnostics Flag --- Conditional Generated Quantities}}\label{sec-run-diagnostics}

\emph{First introduced: Chapter~\ref{sec-model-quality-checks} (Chapter
5)}

Computing \texttt{log\_lik} (N values per draw) and \texttt{y\_rep} in
\texttt{generated\ quantities} can double run-time. Gate them behind an
integer flag so SBC backends can turn them off.

\begin{Shaded}
\begin{Highlighting}[]
\KeywordTok{data}\NormalTok{ \{}
  \DataTypeTok{int}\NormalTok{\textless{}}\KeywordTok{lower}\NormalTok{=}\DecValTok{0}\NormalTok{, }\KeywordTok{upper}\NormalTok{=}\DecValTok{1}\NormalTok{\textgreater{} run\_diagnostics;}
\NormalTok{\}}
\KeywordTok{generated quantities}\NormalTok{ \{}
  \DataTypeTok{vector}\NormalTok{[N] log\_lik = rep\_vector(}\FloatTok{0.0}\NormalTok{, N);}
  \DataTypeTok{array}\NormalTok{[N] }\DataTypeTok{int}\NormalTok{ y\_rep;}
  \ControlFlowTok{if}\NormalTok{ (run\_diagnostics) \{}
    \ControlFlowTok{for}\NormalTok{ (i }\ControlFlowTok{in} \DecValTok{1}\NormalTok{:N) \{}
\NormalTok{      log\_lik[i] = bernoulli\_logit\_lpmf(y[i] | mu);}
\NormalTok{      y\_rep[i]   = bernoulli\_logit\_rng(mu);}
\NormalTok{    \}}
\NormalTok{  \}}
\NormalTok{\}}
\end{Highlighting}
\end{Shaded}

\begin{Shaded}
\begin{Highlighting}[]
\CommentTok{\# Full fit — diagnostics on}
\NormalTok{data }\OtherTok{\textless{}{-}} \FunctionTok{list}\NormalTok{(..., }\AttributeTok{run\_diagnostics =} \DecValTok{1}\NormalTok{L)}

\CommentTok{\# SBC backend — diagnostics off, \textasciitilde{}2× faster generated quantities}
\NormalTok{sbc\_data }\OtherTok{\textless{}{-}} \FunctionTok{modifyList}\NormalTok{(data, }\FunctionTok{list}\NormalTok{(}\AttributeTok{run\_diagnostics =} \DecValTok{0}\NormalTok{L))}
\end{Highlighting}
\end{Shaded}

\begin{tcolorbox}[enhanced jigsaw, toprule=.15mm, toptitle=1mm, bottomtitle=1mm, leftrule=.75mm, rightrule=.15mm, colframe=quarto-callout-tip-color-frame, left=2mm, colback=white, titlerule=0mm, coltitle=black, colbacktitle=quarto-callout-tip-color!10!white, opacityback=0, title=\textcolor{quarto-callout-tip-color}{\faLightbulb}\hspace{0.5em}{When does it matter most?}, arc=.35mm, opacitybacktitle=0.6, breakable, bottomrule=.15mm]

For 1 000 trials, 4 chains, 1 000 sampling draws,
\texttt{run\_diagnostics\ =\ 1} generates and saves 4 million
\texttt{log\_lik} values. During SBC (which only needs the parameters),
setting \texttt{run\_diagnostics\ =\ 0} eliminates that cost entirely.

\end{tcolorbox}

\begin{center}\rule{0.5\linewidth}{0.5pt}\end{center}

\section{\texorpdfstring{7. Post-Hoc Generated Quantities with
\texttt{\$generate\_quantities()}}{7. Post-Hoc Generated Quantities with \$generate\_quantities()}}\label{sec-generate-quantities}

\emph{First introduced: Chapter~\ref{sec-model-quality-checks} (Chapter
5)}

The \texttt{run\_diagnostics} flag lets you skip \texttt{log\_lik} and
\texttt{y\_rep} during SBC. A complementary pattern avoids the trade-off
entirely: \textbf{fit once with diagnostics off, then replay the saved
parameter draws through the \texttt{generated\ quantities} block on
demand}, without re-running MCMC.

\begin{Shaded}
\begin{Highlighting}[]
\CommentTok{\# Fast fit — no log\_lik or y\_rep during sampling}
\NormalTok{fit\_fast }\OtherTok{\textless{}{-}}\NormalTok{ mod}\SpecialCharTok{$}\FunctionTok{sample}\NormalTok{(}\AttributeTok{data =} \FunctionTok{list}\NormalTok{(..., }\AttributeTok{run\_diagnostics =} \DecValTok{0}\NormalTok{L), ...)}

\CommentTok{\# Later: add log\_lik and y\_rep from the saved draws}
\NormalTok{fit\_gq }\OtherTok{\textless{}{-}}\NormalTok{ mod}\SpecialCharTok{$}\FunctionTok{generate\_quantities}\NormalTok{(}
  \AttributeTok{fitted\_params =}\NormalTok{ fit\_fast,}
  \AttributeTok{data          =} \FunctionTok{list}\NormalTok{(..., }\AttributeTok{run\_diagnostics =} \DecValTok{1}\NormalTok{L)}
\NormalTok{)}

\NormalTok{log\_lik }\OtherTok{\textless{}{-}}\NormalTok{ fit\_gq}\SpecialCharTok{$}\FunctionTok{draws}\NormalTok{(}\StringTok{"log\_lik"}\NormalTok{)}
\NormalTok{y\_rep   }\OtherTok{\textless{}{-}}\NormalTok{ fit\_gq}\SpecialCharTok{$}\FunctionTok{draws}\NormalTok{(}\StringTok{"y\_rep"}\NormalTok{)}
\end{Highlighting}
\end{Shaded}

\texttt{\$generate\_quantities()} runs only the
\texttt{generated\ quantities} block --- the sampler is never touched.
Key consequences:

\begin{itemize}
\tightlist
\item
  LOO and PPC can be \textbf{deferred} until you actually need them.
\item
  A fit saved with \texttt{readRDS()} from a previous session can be
  replayed: pass the loaded object as \texttt{fitted\_params}.
\item
  SBC backends always run with \texttt{run\_diagnostics\ =\ 0}; call
  \texttt{\$generate\_quantities()} only on the small subset of fits you
  want to inspect in detail.
\end{itemize}

The compiled model object must be the same for both the original fit and
the \texttt{\$generate\_quantities()} call.

\begin{center}\rule{0.5\linewidth}{0.5pt}\end{center}

\section{8. Non-Centered Parameterization (NCP)}\label{sec-ncp}

\emph{First introduced: Chapter~\ref{sec-multilevel-models} (Chapter 6,
multilevel biased agent)} \emph{Extended to simplex parameters:
Chapter~\ref{sec-categorization-exemplars} (Chapter 11, multilevel decay
GCM)}

Centered parameterizations create \textbf{funnel geometries} when the
group-level variance (\texttt{sigma}) is small: individual-level
parameters collapse toward the mean, forming a narrow neck that HMC
cannot traverse efficiently, producing divergences, low E-BFMI, and
\texttt{max\_treedepth} warnings.

\textbf{Centered (problematic):}

\begin{Shaded}
\begin{Highlighting}[]
\KeywordTok{parameters}\NormalTok{ \{}
  \DataTypeTok{real}\NormalTok{ mu;}
  \DataTypeTok{real}\NormalTok{\textless{}}\KeywordTok{lower}\NormalTok{=}\DecValTok{0}\NormalTok{\textgreater{} sigma;}
  \DataTypeTok{vector}\NormalTok{[J] theta;}
\NormalTok{\}}
\KeywordTok{model}\NormalTok{ \{}
\NormalTok{  theta \textasciitilde{} normal(mu, sigma);}
\NormalTok{\}}
\end{Highlighting}
\end{Shaded}

~

\textbf{Non-Centered (preferred):}

\begin{Shaded}
\begin{Highlighting}[]
\KeywordTok{parameters}\NormalTok{ \{}
  \DataTypeTok{real}\NormalTok{ mu;}
  \DataTypeTok{real}\NormalTok{\textless{}}\KeywordTok{lower}\NormalTok{=}\DecValTok{0}\NormalTok{\textgreater{} sigma;}
  \DataTypeTok{vector}\NormalTok{[J] z;          }\CommentTok{// N(0,1) offset}
\NormalTok{\}}
\KeywordTok{transformed parameters}\NormalTok{ \{}
  \DataTypeTok{vector}\NormalTok{[J] theta = mu + sigma * z;}
\NormalTok{\}}
\KeywordTok{model}\NormalTok{ \{}
\NormalTok{  z \textasciitilde{} std\_normal();}
\NormalTok{\}}
\end{Highlighting}
\end{Shaded}

\subsection{Logit-Normal NCP for Simplex
Parameters}\label{logit-normal-ncp-for-simplex-parameters}

For attention weights or mixing proportions modeled as a simplex, avoid
the Dirichlet + kappa parameterization. Small kappa values produce
nearly-flat Dirichlet distributions, creating the same funnel that
centered normal models do. For a 2-element simplex use a
\textbf{logit-normal NCP} instead:

\begin{Shaded}
\begin{Highlighting}[]
\KeywordTok{parameters}\NormalTok{ \{}
  \DataTypeTok{real}\NormalTok{          pop\_logit\_w\_mean;}
  \DataTypeTok{real}\NormalTok{\textless{}}\KeywordTok{lower}\NormalTok{=}\DecValTok{0}\NormalTok{\textgreater{} pop\_logit\_w\_sd;}
  \DataTypeTok{vector}\NormalTok{[J]     z\_w;}
\NormalTok{\}}
\KeywordTok{transformed parameters}\NormalTok{ \{}
  \DataTypeTok{array}\NormalTok{[J] }\DataTypeTok{vector}\NormalTok{[}\DecValTok{2}\NormalTok{] w;}
  \ControlFlowTok{for}\NormalTok{ (j }\ControlFlowTok{in} \DecValTok{1}\NormalTok{:J) \{}
    \DataTypeTok{real}\NormalTok{ lw  = pop\_logit\_w\_mean + pop\_logit\_w\_sd * z\_w[j];}
\NormalTok{    w[j][}\DecValTok{1}\NormalTok{]  = inv\_logit(lw);}
\NormalTok{    w[j][}\DecValTok{2}\NormalTok{]  = }\FloatTok{1.0}\NormalTok{ {-} w[j][}\DecValTok{1}\NormalTok{];}
\NormalTok{  \}}
\NormalTok{\}}
\KeywordTok{model}\NormalTok{ \{}
\NormalTok{  z\_w \textasciitilde{} std\_normal();}
\NormalTok{\}}
\end{Highlighting}
\end{Shaded}

\begin{tcolorbox}[enhanced jigsaw, toprule=.15mm, toptitle=1mm, bottomtitle=1mm, leftrule=.75mm, rightrule=.15mm, colframe=quarto-callout-warning-color-frame, left=2mm, colback=white, titlerule=0mm, coltitle=black, colbacktitle=quarto-callout-warning-color!10!white, opacityback=0, title=\textcolor{quarto-callout-warning-color}{\faExclamationTriangle}\hspace{0.5em}{The cost of ignoring NCP}, arc=.35mm, opacitybacktitle=0.6, breakable, bottomrule=.15mm]

The centered Dirichlet +
\texttt{kappa\ \textasciitilde{}\ exponential(0.1)} parameterization in
Chapter 11's multilevel decay GCM produced \textbf{4 000 max\_treedepth
warnings per run} and failed to complete 10 iterations overnight at
\texttt{adapt\_delta\ =\ 0.99,\ max\_treedepth\ =\ 20}. Switching to
logit-normal NCP eliminated the funnel geometry entirely.

\end{tcolorbox}

\begin{tcolorbox}[enhanced jigsaw, toprule=.15mm, toptitle=1mm, bottomtitle=1mm, leftrule=.75mm, rightrule=.15mm, colframe=quarto-callout-tip-color-frame, left=2mm, colback=white, titlerule=0mm, coltitle=black, colbacktitle=quarto-callout-tip-color!10!white, opacityback=0, title=\textcolor{quarto-callout-tip-color}{\faLightbulb}\hspace{0.5em}{When centered parameterization is preferred}, arc=.35mm, opacitybacktitle=0.6, breakable, bottomrule=.15mm]

NCP is not always the right choice. When the data are highly informative
relative to the prior, the \emph{centered} parameterization can be more
efficient. In that regime, the group-level variance \(\sigma\) is
well-constrained by the data, so the individual-level parameters
\(\theta_j\) are tightly funneled by the likelihood rather than the
prior --- and HMC can navigate the centered geometry directly without
the stretched \(z\)-offsets that NCP introduces.

A practical heuristic (the \textbf{park rule}): prefer NCP when the
per-group sample size is small relative to the between-group variance
(data are weak); prefer CP when data are plentiful and the posterior for
\(\sigma\) is far from zero. When in doubt, fit both and compare
diagnostics.

See
\href{https://avehtari.github.io/Bayesian-Workflow/park_rule/park_rule.html}{Vehtari
et al.~--- the park rule} for a formal treatment with diagnostics to
guide the choice.

\end{tcolorbox}

\begin{center}\rule{0.5\linewidth}{0.5pt}\end{center}

\section{9. Cholesky Factorization for Correlation
Matrices}\label{sec-cholesky}

\emph{First introduced: Chapter~\ref{sec-multilevel-models} (Chapter 6,
multilevel memory agent with correlated parameters)}

Whenever a model estimates a \textbf{correlation matrix} \(\Omega\) over
individual-level parameters, use the Cholesky-factorized form rather
than sampling \(\Omega\) directly.

\begin{longtable}[]{@{}
  >{\raggedright\arraybackslash}p{(\linewidth - 4\tabcolsep) * \real{0.3333}}
  >{\raggedright\arraybackslash}p{(\linewidth - 4\tabcolsep) * \real{0.3333}}
  >{\raggedright\arraybackslash}p{(\linewidth - 4\tabcolsep) * \real{0.3333}}@{}}
\toprule\noalign{}
\begin{minipage}[b]{\linewidth}\raggedright
Parameterization
\end{minipage} & \begin{minipage}[b]{\linewidth}\raggedright
Leapfrog cost
\end{minipage} & \begin{minipage}[b]{\linewidth}\raggedright
Why prefer/avoid
\end{minipage} \\
\midrule\noalign{}
\endhead
\bottomrule\noalign{}
\endlastfoot
\texttt{corr\_matrix{[}K{]}\ Omega} + \texttt{lkj\_corr\_lpdf} &
\(O(K^3)\) matrix operations + positive-definiteness check & Expensive;
unconstrained transform requires matrix square root \\
\texttt{cholesky\_factor\_corr{[}K{]}\ L\_Omega} +
\texttt{lkj\_corr\_cholesky\_lpdf} & \(O(K^2)\) & Lies directly on its
natural manifold; no matrix inversion needed \\
\end{longtable}

\begin{Shaded}
\begin{Highlighting}[]
\CommentTok{// Avoid}
\KeywordTok{parameters}\NormalTok{ \{ }\DataTypeTok{corr\_matrix}\NormalTok{[}\DecValTok{2}\NormalTok{] Omega; \}}
\KeywordTok{model}\NormalTok{      \{ Omega \textasciitilde{} lkj\_corr(}\FloatTok{2.0}\NormalTok{); \}}

\CommentTok{// Prefer}
\KeywordTok{parameters}\NormalTok{ \{ }\DataTypeTok{cholesky\_factor\_corr}\NormalTok{[}\DecValTok{2}\NormalTok{] L\_Omega; \}}
\KeywordTok{model}\NormalTok{      \{ L\_Omega \textasciitilde{} lkj\_corr\_cholesky(}\FloatTok{2.0}\NormalTok{); \}}

\KeywordTok{transformed parameters}\NormalTok{ \{}
  \CommentTok{// O(K\^{}2) multiply — no square root}
  \DataTypeTok{matrix}\NormalTok{[}\DecValTok{2}\NormalTok{, J] deviations = diag\_pre\_multiply(sigma, L\_Omega) * z;}
\NormalTok{\}}

\KeywordTok{generated quantities}\NormalTok{ \{}
  \CommentTok{// Reconstruct Omega once per saved draw, not every leapfrog step}
  \DataTypeTok{matrix}\NormalTok{[}\DecValTok{2}\NormalTok{, }\DecValTok{2}\NormalTok{] Omega = multiply\_lower\_tri\_self\_transpose(L\_Omega);}
  \DataTypeTok{real}\NormalTok{ rho = Omega[}\DecValTok{1}\NormalTok{, }\DecValTok{2}\NormalTok{];}
\NormalTok{\}}
\end{Highlighting}
\end{Shaded}

\textbf{Key point:} \texttt{diag\_pre\_multiply(sigma,\ L\_Omega)\ *\ z}
produces individual-level NCP deviations directly from the Cholesky
factor --- the full correlation matrix \(\Omega\) is only recovered in
\texttt{generated\ quantities}, running once per saved draw rather than
thousands of times during warm-up.

For \(K = 2\) parameters the cost difference is modest; for \(K \geq 5\)
(e.g., five per-subject RL parameters) it becomes substantial.

\begin{center}\rule{0.5\linewidth}{0.5pt}\end{center}

\section{\texorpdfstring{10. Contiguous Subject Indexing ---
\texttt{subj\_start} /
\texttt{subj\_end}}{10. Contiguous Subject Indexing --- subj\_start / subj\_end}}\label{sec-subject-indexing}

\emph{First introduced: Chapter~\ref{sec-multilevel-models} (Chapter 6,
multilevel memory model)}

Two strategies for multi-subject trial data:

\textbf{Strategy A --- \texttt{subj\_id} lookup} (simple; works only
when trials are independent):

\begin{Shaded}
\begin{Highlighting}[]
\KeywordTok{data}\NormalTok{ \{ }\DataTypeTok{array}\NormalTok{[N\_total] }\DataTypeTok{int}\NormalTok{ subj\_id; \}}
\KeywordTok{model}\NormalTok{ \{}
  \ControlFlowTok{for}\NormalTok{ (i }\ControlFlowTok{in} \DecValTok{1}\NormalTok{:N\_total)}
    \KeywordTok{target +=}\NormalTok{ bernoulli\_lpmf(y[i] | f(theta[subj\_id[i]], x[i]));}
\NormalTok{\}}
\end{Highlighting}
\end{Shaded}

\textbf{Strategy B --- Contiguous block indexing} (required for
sequential models and \texttt{reduce\_sum}):

\begin{Shaded}
\begin{Highlighting}[]
\KeywordTok{data}\NormalTok{ \{}
  \DataTypeTok{array}\NormalTok{[J] }\DataTypeTok{int}\NormalTok{ subj\_start;}
  \DataTypeTok{array}\NormalTok{[J] }\DataTypeTok{int}\NormalTok{ subj\_end;}
\NormalTok{\}}
\KeywordTok{model}\NormalTok{ \{}
  \ControlFlowTok{for}\NormalTok{ (j }\ControlFlowTok{in} \DecValTok{1}\NormalTok{:J)}
    \ControlFlowTok{for}\NormalTok{ (t }\ControlFlowTok{in}\NormalTok{ subj\_start[j]:subj\_end[j])}
      \KeywordTok{target +=}\NormalTok{ bernoulli\_lpmf(y[t] | f(theta[j], x[t]));}
\NormalTok{\}}
\end{Highlighting}
\end{Shaded}

Build the index in R:

\begin{Shaded}
\begin{Highlighting}[]
\NormalTok{subj\_start }\OtherTok{\textless{}{-}} \FunctionTok{as.integer}\NormalTok{(}\FunctionTok{tapply}\NormalTok{(}\FunctionTok{seq\_along}\NormalTok{(subj\_id), subj\_id, min))}
\NormalTok{subj\_end   }\OtherTok{\textless{}{-}} \FunctionTok{as.integer}\NormalTok{(}\FunctionTok{tapply}\NormalTok{(}\FunctionTok{seq\_along}\NormalTok{(subj\_id), subj\_id, max))}
\end{Highlighting}
\end{Shaded}

Use Strategy B whenever:

\begin{itemize}
\tightlist
\item
  Trial \(t\) depends on state accumulated from trials
  \(1, \ldots, t-1\) (memory, belief updating, RL)
\item
  You intend to use \texttt{reduce\_sum} for within-chain parallelism
  (see Section~\ref{sec-reduce-sum})
\end{itemize}

\begin{center}\rule{0.5\linewidth}{0.5pt}\end{center}

\section{11. Pathfinder Warm-Start Initialisation}\label{sec-pathfinder}

\emph{First benchmarked: Chapter~\ref{sec-inferring-rates} (Chapter 4)}
\emph{Timing comparison for multilevel models:
Chapter~\ref{sec-multilevel-models} (Chapter 6)} \emph{Three-way
benchmark: Chapter~\ref{sec-categorization-exemplars} (Chapter 11)}

Stan's default initialisation draws chain starts from Uniform(−2, 2) in
unconstrained space --- far from the posterior on hierarchical models.
\textbf{Pathfinder} runs L-BFGS from each start, fits Gaussian
approximations along the path, and returns approximate posterior draws.
Passing these as \texttt{init} to \texttt{\$sample()} lets HMC skip most
of the warm-up trajectory.

\begin{Shaded}
\begin{Highlighting}[]
\CommentTok{\# Step 1 — fast variational approximation}
\NormalTok{pf }\OtherTok{\textless{}{-}}\NormalTok{ mod}\SpecialCharTok{$}\FunctionTok{pathfinder}\NormalTok{(}\AttributeTok{data =}\NormalTok{ stan\_data, }\AttributeTok{num\_paths =} \DecValTok{4}\NormalTok{, }\AttributeTok{draws =} \DecValTok{1000}\NormalTok{, }\AttributeTok{seed =} \DecValTok{123}\NormalTok{)}

\CommentTok{\# Step 2 — draw one starting point per chain}
\NormalTok{init\_pf }\OtherTok{\textless{}{-}} \FunctionTok{lapply}\NormalTok{(}\DecValTok{1}\SpecialCharTok{:}\DecValTok{4}\NormalTok{, }\ControlFlowTok{function}\NormalTok{(i) }\FunctionTok{as.list}\NormalTok{(pf}\SpecialCharTok{$}\FunctionTok{draws}\NormalTok{()[i, , }\AttributeTok{drop =} \ConstantTok{FALSE}\NormalTok{]))}

\CommentTok{\# Step 3 — full MCMC from warm starts}
\NormalTok{fit }\OtherTok{\textless{}{-}}\NormalTok{ mod}\SpecialCharTok{$}\FunctionTok{sample}\NormalTok{(}\AttributeTok{data =}\NormalTok{ stan\_data, }\AttributeTok{init =}\NormalTok{ init\_pf,}
                  \AttributeTok{chains =} \DecValTok{4}\NormalTok{, }\AttributeTok{iter\_warmup =} \DecValTok{500}\NormalTok{, }\AttributeTok{iter\_sampling =} \DecValTok{1000}\NormalTok{)}
\end{Highlighting}
\end{Shaded}

\begin{tcolorbox}[enhanced jigsaw, toprule=.15mm, toptitle=1mm, bottomtitle=1mm, leftrule=.75mm, rightrule=.15mm, colframe=quarto-callout-tip-color-frame, left=2mm, colback=white, titlerule=0mm, coltitle=black, colbacktitle=quarto-callout-tip-color!10!white, opacityback=0, title=\textcolor{quarto-callout-tip-color}{\faLightbulb}\hspace{0.5em}{When Pathfinder helps most}, arc=.35mm, opacitybacktitle=0.6, breakable, bottomrule=.15mm]

Pathfinder gives the largest benefit on \textbf{hierarchical models} and
any model with correlated population parameters. It saves 20--50\% of
total wall time by shortening the warmup phase. On simple,
well-conditioned models the warm-up is already fast --- the Pathfinder
overhead may outweigh the saving.

\textbf{Pathfinder output is not a gold-standard posterior.} Always run
full MCMC and check diagnostics.

\end{tcolorbox}

\begin{center}\rule{0.5\linewidth}{0.5pt}\end{center}

\section{\texorpdfstring{12. Marginalizing Discrete Parameters with
\texttt{log\_mix}}{12. Marginalizing Discrete Parameters with log\_mix}}\label{sec-log-mix}

\emph{First introduced: Chapter~\ref{sec-mixture-models} (Chapter 8,
single and multilevel mixture models)}

HMC requires continuous, differentiable parameters. Discrete latent
variables (component indicators, regime switches) must be
\textbf{analytically marginalized} out of the likelihood before
sampling.

\begin{Shaded}
\begin{Highlighting}[]
\CommentTok{// two{-}component mixture}
\DataTypeTok{real}\NormalTok{ lp1 = bernoulli\_logit\_lpmf(y | alpha);}
\DataTypeTok{real}\NormalTok{ lp2 = bernoulli\_logit\_lpmf(y | nu);}
\CommentTok{// log( pi*exp(lp1) + (1{-}pi)*exp(lp2) )}
\KeywordTok{target +=}\NormalTok{ log\_mix(pi, lp1, lp2);}
\end{Highlighting}
\end{Shaded}

For \(K > 2\) components use \texttt{log\_sum\_exp} directly:

\begin{Shaded}
\begin{Highlighting}[]
\DataTypeTok{vector}\NormalTok{[K] lps;}
\ControlFlowTok{for}\NormalTok{ (k }\ControlFlowTok{in} \DecValTok{1}\NormalTok{:K)}
\NormalTok{  lps[k] = log(pi[k]) + component\_lpmf(y | theta[k]);}
\KeywordTok{target +=}\NormalTok{ log\_sum\_exp(lps);}
\end{Highlighting}
\end{Shaded}

\begin{tcolorbox}[enhanced jigsaw, toprule=.15mm, toptitle=1mm, bottomtitle=1mm, leftrule=.75mm, rightrule=.15mm, colframe=quarto-callout-warning-color-frame, left=2mm, colback=white, titlerule=0mm, coltitle=black, colbacktitle=quarto-callout-warning-color!10!white, opacityback=0, title=\textcolor{quarto-callout-warning-color}{\faExclamationTriangle}\hspace{0.5em}{Never sample discrete indicators directly}, arc=.35mm, opacitybacktitle=0.6, breakable, bottomrule=.15mm]

\begin{Shaded}
\begin{Highlighting}[]
\CommentTok{// WRONG — HMC cannot compute gradients of discrete parameters}
\DataTypeTok{int}\NormalTok{ z \textasciitilde{} categorical(pi);}
\KeywordTok{target +=}\NormalTok{ bernoulli\_lpmf(y | theta[z]);}
\end{Highlighting}
\end{Shaded}

This will compile but produce incorrect results or sampler errors at
runtime. Always marginalise.

\end{tcolorbox}

\begin{center}\rule{0.5\linewidth}{0.5pt}\end{center}

\section{\texorpdfstring{13. Sequential State Unrolling in
\texttt{transformed\ parameters}}{13. Sequential State Unrolling in transformed parameters}}\label{sec-sequential-state}

\emph{First introduced: Chapter~\ref{sec-theory-of-mind} (Chapter 9,
weighted and proportional agent models)}

When the likelihood depends on a latent state that evolves
trial-by-trial (belief, RL value, running weight), unroll the recursion
in \texttt{transformed\ parameters} so the gradient flows back through
the entire trajectory.

\begin{Shaded}
\begin{Highlighting}[]
\KeywordTok{transformed parameters}\NormalTok{ \{}
  \DataTypeTok{vector}\NormalTok{[N] predicted\_choice;}
\NormalTok{  \{}
    \CommentTok{// local scope: scratch variable allocated and freed each leapfrog step}
    \DataTypeTok{real}\NormalTok{ belief = }\FloatTok{0.5}\NormalTok{;}
    \ControlFlowTok{for}\NormalTok{ (t }\ControlFlowTok{in} \DecValTok{1}\NormalTok{:N) \{}
\NormalTok{      predicted\_choice[t] = inv\_logit(weight * (belief {-} }\FloatTok{0.5}\NormalTok{));}
\NormalTok{      belief += alpha * (feedback[t] {-} belief);}
\NormalTok{    \}}
\NormalTok{  \}}
\NormalTok{\}}
\KeywordTok{model}\NormalTok{ \{}
  \KeywordTok{target +=}\NormalTok{ bernoulli\_lpmf(y | predicted\_choice);}
\NormalTok{\}}
\end{Highlighting}
\end{Shaded}

\begin{tcolorbox}[enhanced jigsaw, toprule=.15mm, toptitle=1mm, bottomtitle=1mm, leftrule=.75mm, rightrule=.15mm, colframe=quarto-callout-tip-color-frame, left=2mm, colback=white, titlerule=0mm, coltitle=black, colbacktitle=quarto-callout-tip-color!10!white, opacityback=0, title=\textcolor{quarto-callout-tip-color}{\faLightbulb}\hspace{0.5em}{Why \texttt{transformed\ parameters} and not \texttt{model}?}, arc=.35mm, opacitybacktitle=0.6, breakable, bottomrule=.15mm]

State that enters the likelihood needs its gradient tracked by autodiff.
Declare only the \textbf{output vector} (\texttt{predicted\_choice}) as
a \texttt{transformed\ parameters} variable; keep loop-internal scratch
variables inside a \texttt{\{\}} local block to avoid placing them on
the saved-draw tape.

\textbf{Exception:} If the loop is independent across subjects, move it
into a \texttt{partial\_sum\_lp} function and use \texttt{reduce\_sum}
for across-subject parallelism (see Section~\ref{sec-reduce-sum}).

\end{tcolorbox}

\begin{center}\rule{0.5\linewidth}{0.5pt}\end{center}

\section{\texorpdfstring{14. Intra-Chain Parallelism with
\texttt{reduce\_sum}}{14. Intra-Chain Parallelism with reduce\_sum}}\label{sec-reduce-sum}

\emph{First introduced: Chapter~\ref{sec-categorization-exemplars}
(Chapter 11, multilevel decay GCM)}

\texttt{reduce\_sum} exploits \textbf{conditional independence across
subjects}: given population parameters, subject \(j\)'s trial loop
shares no mutable state with subject \(k\)'s loop. The subject dimension
can therefore be distributed across CPU threads \textbf{within a single
chain}, giving near-linear speedup up to the number of subjects.

\textbf{Step 1 --- Write a \texttt{partial\_sum\_lp} function in the
\texttt{functions} block:}

\begin{Shaded}
\begin{Highlighting}[]
\KeywordTok{functions}\NormalTok{ \{}
  \DataTypeTok{real}\NormalTok{ partial\_sum\_lp(}
    \DataTypeTok{array}\NormalTok{[] }\DataTypeTok{int}\NormalTok{ subj\_slice, }\DataTypeTok{int}\NormalTok{ start, }\DataTypeTok{int}\NormalTok{ end,}
    \DataTypeTok{vector}\NormalTok{ log\_c,}
    \DataTypeTok{array}\NormalTok{[] }\DataTypeTok{int}\NormalTok{ subj\_start\_g, }\DataTypeTok{array}\NormalTok{[] }\DataTypeTok{int}\NormalTok{ subj\_end\_g,}
    \DataTypeTok{array}\NormalTok{[] }\DataTypeTok{int}\NormalTok{ y, }\DataTypeTok{array}\NormalTok{[,] }\DataTypeTok{real}\NormalTok{ x}
\NormalTok{  ) \{}
    \DataTypeTok{real}\NormalTok{ lp = }\FloatTok{0.0}\NormalTok{;}
    \ControlFlowTok{for}\NormalTok{ (idx }\ControlFlowTok{in} \DecValTok{1}\NormalTok{:size(subj\_slice)) \{}
      \DataTypeTok{int}\NormalTok{ j = subj\_slice[idx];}
      \ControlFlowTok{for}\NormalTok{ (t }\ControlFlowTok{in}\NormalTok{ subj\_start\_g[j]:subj\_end\_g[j])}
\NormalTok{        lp += bernoulli\_lpmf(y[t] | f(exp(log\_c[j]), x[t]));}
\NormalTok{    \}}
    \ControlFlowTok{return}\NormalTok{ lp;}
\NormalTok{  \}}
\NormalTok{\}}
\end{Highlighting}
\end{Shaded}

\textbf{Step 2 --- Replace the subject loop with \texttt{reduce\_sum} in
the \texttt{model} block:}

\begin{Shaded}
\begin{Highlighting}[]
\KeywordTok{model}\NormalTok{ \{}
  \CommentTok{// subject\_ids = \{1, 2, ..., J\}}
  \KeywordTok{target +=}\NormalTok{ reduce\_sum(partial\_sum\_lp, subject\_ids, }\DecValTok{1}\NormalTok{,}
\NormalTok{                       log\_c, subj\_start, subj\_end, y, x);}
\NormalTok{\}}
\end{Highlighting}
\end{Shaded}

\textbf{Step 3 --- Compile with threading support and sample:}

\begin{Shaded}
\begin{Highlighting}[]
\NormalTok{mod\_par }\OtherTok{\textless{}{-}} \FunctionTok{cmdstan\_model}\NormalTok{(stan\_file,}
                         \AttributeTok{cpp\_options =} \FunctionTok{list}\NormalTok{(}\AttributeTok{stan\_threads =} \ConstantTok{TRUE}\NormalTok{))}

\NormalTok{n\_threads }\OtherTok{\textless{}{-}} \FunctionTok{max}\NormalTok{(}\DecValTok{1}\NormalTok{L, parallel}\SpecialCharTok{::}\FunctionTok{detectCores}\NormalTok{() }\SpecialCharTok{\%/\%} \DecValTok{4}\NormalTok{L)}

\NormalTok{fit\_par }\OtherTok{\textless{}{-}}\NormalTok{ mod\_par}\SpecialCharTok{$}\FunctionTok{sample}\NormalTok{(}
  \AttributeTok{data             =}\NormalTok{ stan\_data,}
  \AttributeTok{chains           =} \DecValTok{4}\NormalTok{,}
  \AttributeTok{parallel\_chains  =} \DecValTok{4}\NormalTok{,}
  \AttributeTok{threads\_per\_chain =}\NormalTok{ n\_threads,}
  \AttributeTok{iter\_warmup      =} \DecValTok{500}\NormalTok{,}
  \AttributeTok{iter\_sampling    =} \DecValTok{1000}
\NormalTok{)}
\end{Highlighting}
\end{Shaded}

\begin{tcolorbox}[enhanced jigsaw, toprule=.15mm, toptitle=1mm, bottomtitle=1mm, leftrule=.75mm, rightrule=.15mm, colframe=quarto-callout-tip-color-frame, left=2mm, colback=white, titlerule=0mm, coltitle=black, colbacktitle=quarto-callout-tip-color!10!white, opacityback=0, title=\textcolor{quarto-callout-tip-color}{\faLightbulb}\hspace{0.5em}{When is \texttt{reduce\_sum} worth it?}, arc=.35mm, opacitybacktitle=0.6, breakable, bottomrule=.15mm]

\begin{longtable}[]{@{}
  >{\raggedright\arraybackslash}p{(\linewidth - 2\tabcolsep) * \real{0.5000}}
  >{\raggedright\arraybackslash}p{(\linewidth - 2\tabcolsep) * \real{0.5000}}@{}}
\toprule\noalign{}
\begin{minipage}[b]{\linewidth}\raggedright
Condition
\end{minipage} & \begin{minipage}[b]{\linewidth}\raggedright
Effect
\end{minipage} \\
\midrule\noalign{}
\endhead
\bottomrule\noalign{}
\endlastfoot
≥ 20 subjects with expensive per-subject loops & Large speedup
(near-linear up to J) \\
Few subjects or very cheap loops & Threading overhead may dominate \\
Within-subject sequential dependency & Still valid --- parallelises
\emph{across} subjects, not \emph{within} \\
Contiguous \texttt{subj\_start}/\texttt{subj\_end} indexing &
\textbf{Required} (see Section~\ref{sec-subject-indexing}) \\
\end{longtable}

The grainsize argument (second integer in \texttt{reduce\_sum}) controls
the minimum slice size. Use \texttt{1} as a starting point; increase if
the per-subject work is small.

\end{tcolorbox}

\begin{tcolorbox}[enhanced jigsaw, toprule=.15mm, toptitle=1mm, bottomtitle=1mm, leftrule=.75mm, rightrule=.15mm, colframe=quarto-callout-note-color-frame, left=2mm, colback=white, titlerule=0mm, coltitle=black, colbacktitle=quarto-callout-note-color!10!white, opacityback=0, title=\textcolor{quarto-callout-note-color}{\faInfo}\hspace{0.5em}{Within-subject loops cannot be parallelised with \texttt{reduce\_sum}}, arc=.35mm, opacitybacktitle=0.6, breakable, bottomrule=.15mm]

In the single-subject GCM (Chapter 11), trial \(t\) depends on the
memory state accumulated from trials \(1, \ldots, t-1\). This sequential
dependency is \textbf{within} a subject, so \texttt{reduce\_sum} cannot
help there. Parallelism across subjects is always safe because subjects
are conditionally independent.

\end{tcolorbox}

\begin{center}\rule{0.5\linewidth}{0.5pt}\end{center}

\section{15. Combining Techniques --- Compound
Speedups}\label{sec-compound-speedups}

\emph{Three-way timing benchmark:
Chapter~\ref{sec-categorization-exemplars} (Chapter 11)}

The techniques above stack multiplicatively on hierarchical models:

\begin{longtable}[]{@{}
  >{\raggedright\arraybackslash}p{(\linewidth - 4\tabcolsep) * \real{0.3333}}
  >{\raggedright\arraybackslash}p{(\linewidth - 4\tabcolsep) * \real{0.3333}}
  >{\raggedright\arraybackslash}p{(\linewidth - 4\tabcolsep) * \real{0.3333}}@{}}
\toprule\noalign{}
\begin{minipage}[b]{\linewidth}\raggedright
Technique
\end{minipage} & \begin{minipage}[b]{\linewidth}\raggedright
Typical saving
\end{minipage} & \begin{minipage}[b]{\linewidth}\raggedright
When to apply
\end{minipage} \\
\midrule\noalign{}
\endhead
\bottomrule\noalign{}
\endlastfoot
NCP instead of CP & 5--100× fewer divergences & Any multilevel model ---
always prefer NCP \\
\texttt{transformed\ data} precompute & 10--100× fewer flops & Static
quantities derived from data only \\
\texttt{run\_diagnostics\ =\ 0} in SBC & \textasciitilde2× faster
\texttt{generated\ quantities} & All SBC backends \\
Pathfinder warm-start & 20--50\% less wall time & Hierarchical models;
correlated posteriors \\
\texttt{reduce\_sum} (10 subjects, 4-core laptop) & 2--3× within-chain &
≥ 20 subjects, expensive loops \\
Pathfinder + \texttt{reduce\_sum} together & 3--5× vs.~serial baseline &
Large multilevel models, multi-core hardware \\
\end{longtable}

\textbf{Benchmark template} (used in Chapters 6 and 11):

\begin{Shaded}
\begin{Highlighting}[]
\NormalTok{time\_serial }\OtherTok{\textless{}{-}} \FunctionTok{system.time}\NormalTok{(}
\NormalTok{  fit\_serial }\OtherTok{\textless{}{-}}\NormalTok{ mod}\SpecialCharTok{$}\FunctionTok{sample}\NormalTok{(}\AttributeTok{data =}\NormalTok{ stan\_data, ...)}
\NormalTok{)[[}\StringTok{"elapsed"}\NormalTok{]]}

\NormalTok{pf       }\OtherTok{\textless{}{-}}\NormalTok{ mod}\SpecialCharTok{$}\FunctionTok{pathfinder}\NormalTok{(}\AttributeTok{data =}\NormalTok{ stan\_data, }\AttributeTok{num\_paths =} \DecValTok{4}\NormalTok{, }\AttributeTok{draws =} \DecValTok{1000}\NormalTok{)}
\NormalTok{init\_pf  }\OtherTok{\textless{}{-}} \FunctionTok{lapply}\NormalTok{(}\DecValTok{1}\SpecialCharTok{:}\DecValTok{4}\NormalTok{, }\ControlFlowTok{function}\NormalTok{(i) }\FunctionTok{as.list}\NormalTok{(pf}\SpecialCharTok{$}\FunctionTok{draws}\NormalTok{()[i, , }\AttributeTok{drop =} \ConstantTok{FALSE}\NormalTok{]))}

\NormalTok{time\_serial\_pf }\OtherTok{\textless{}{-}} \FunctionTok{system.time}\NormalTok{(}
\NormalTok{  fit\_pf }\OtherTok{\textless{}{-}}\NormalTok{ mod}\SpecialCharTok{$}\FunctionTok{sample}\NormalTok{(}\AttributeTok{data =}\NormalTok{ stan\_data, }\AttributeTok{init =}\NormalTok{ init\_pf, ...)}
\NormalTok{)[[}\StringTok{"elapsed"}\NormalTok{]]}

\NormalTok{n\_threads }\OtherTok{\textless{}{-}} \FunctionTok{max}\NormalTok{(}\DecValTok{1}\NormalTok{L, parallel}\SpecialCharTok{::}\FunctionTok{detectCores}\NormalTok{() }\SpecialCharTok{\%/\%} \DecValTok{4}\NormalTok{L)}
\NormalTok{time\_par  }\OtherTok{\textless{}{-}} \FunctionTok{system.time}\NormalTok{(}
\NormalTok{  fit\_par }\OtherTok{\textless{}{-}}\NormalTok{ mod\_par}\SpecialCharTok{$}\FunctionTok{sample}\NormalTok{(}
    \AttributeTok{data =}\NormalTok{ stan\_data, }\AttributeTok{init =}\NormalTok{ init\_pf,}
    \AttributeTok{threads\_per\_chain =}\NormalTok{ n\_threads, ...)}
\NormalTok{)[[}\StringTok{"elapsed"}\NormalTok{]]}

\FunctionTok{tibble}\NormalTok{(}
  \AttributeTok{Strategy =} \FunctionTok{c}\NormalTok{(}\StringTok{"Serial (baseline)"}\NormalTok{, }\StringTok{"Serial + Pathfinder"}\NormalTok{, }\StringTok{"Parallel + Pathfinder"}\NormalTok{),}
  \AttributeTok{Time\_sec =} \FunctionTok{c}\NormalTok{(time\_serial, time\_serial\_pf, time\_par),}
  \AttributeTok{Speedup  =} \FunctionTok{round}\NormalTok{(time\_serial }\SpecialCharTok{/} \FunctionTok{c}\NormalTok{(time\_serial, time\_serial\_pf, time\_par), }\DecValTok{2}\NormalTok{)}
\NormalTok{)}
\end{Highlighting}
\end{Shaded}

\begin{center}\rule{0.5\linewidth}{0.5pt}\end{center}

\section{16. Sampler and Compiler Options}\label{sec-sampler-options}

\emph{Reference:
\href{https://blog.mc-stan.org/2022/08/03/options-for-improving-stan-sampling-speed/}{Options
for improving Stan sampling speed}}

The techniques in §1--15 reduce the \emph{work per leapfrog step}. The
options below change how the sampler and compiler behave around those
steps. They are cheap to try and often give meaningful savings with no
code changes.

\subsection{\texorpdfstring{Sampler options (passed to
\texttt{\$sample()})}{Sampler options (passed to \$sample())}}\label{sampler-options-passed-to-sample}

\begin{longtable}[]{@{}
  >{\raggedright\arraybackslash}p{(\linewidth - 4\tabcolsep) * \real{0.3333}}
  >{\raggedright\arraybackslash}p{(\linewidth - 4\tabcolsep) * \real{0.3333}}
  >{\raggedright\arraybackslash}p{(\linewidth - 4\tabcolsep) * \real{0.3333}}@{}}
\toprule\noalign{}
\begin{minipage}[b]{\linewidth}\raggedright
Option
\end{minipage} & \begin{minipage}[b]{\linewidth}\raggedright
Default
\end{minipage} & \begin{minipage}[b]{\linewidth}\raggedright
Guidance
\end{minipage} \\
\midrule\noalign{}
\endhead
\bottomrule\noalign{}
\endlastfoot
\texttt{adapt\_delta} & 0.80 & Raise to 0.90--0.95 when divergences
persist after NCP; higher values force smaller step sizes and more
leapfrog steps per transition \\
\texttt{max\_treedepth} & 10 & Raise to 12--15 only when
\texttt{max\_treedepth} warnings dominate; raising without fixing
geometry just burns time \\
\texttt{iter\_warmup} & 1000 & Lower to 500 on well-conditioned models;
raise to 2000 when E-BFMI warnings persist across chains \\
\texttt{iter\_sampling} & 1000 & Raise when posterior tails matter
(e.g., rare-event predictions) or when \(\hat{R}\) instability is
numerical rather than geometric \\
\texttt{chains} & 4 & Never run fewer than 4 --- \(\hat{R}\) requires
between-chain comparison \\
\texttt{parallel\_chains} & 1 & Set to \texttt{chains} when fitting a
single model; reduce when running many parallel fits \\
\end{longtable}

\begin{Shaded}
\begin{Highlighting}[]
\NormalTok{fit }\OtherTok{\textless{}{-}}\NormalTok{ mod}\SpecialCharTok{$}\FunctionTok{sample}\NormalTok{(}
  \AttributeTok{data           =}\NormalTok{ stan\_data,}
  \AttributeTok{chains         =} \DecValTok{4}\NormalTok{,}
  \AttributeTok{parallel\_chains =} \DecValTok{4}\NormalTok{,}
  \AttributeTok{iter\_warmup    =} \DecValTok{1000}\NormalTok{,}
  \AttributeTok{iter\_sampling  =} \DecValTok{1000}\NormalTok{,}
  \AttributeTok{adapt\_delta    =} \FloatTok{0.90}\NormalTok{,     }\CommentTok{\# raise from 0.80 if divergences persist}
  \AttributeTok{max\_treedepth  =} \DecValTok{12}        \CommentTok{\# raise from 10 if treedepth warnings dominate}
\NormalTok{)}
\end{Highlighting}
\end{Shaded}

\subsection{Compiler optimizations}\label{compiler-optimizations}

Stan models are compiled to C++. Enabling processor-specific
optimizations costs nothing at runtime but can save 10--30\% on
algebraically intensive models (GCMs, ToM).

\begin{Shaded}
\begin{Highlighting}[]
\NormalTok{mod }\OtherTok{\textless{}{-}} \FunctionTok{cmdstan\_model}\NormalTok{(}
\NormalTok{  stan\_file,}
  \AttributeTok{cpp\_options =} \FunctionTok{list}\NormalTok{(}
    \AttributeTok{stan\_threads =} \ConstantTok{TRUE}\NormalTok{,          }\CommentTok{\# required for reduce\_sum (§14)}
    \AttributeTok{CXXFLAGS =} \StringTok{"{-}O3 {-}march=native"} \CommentTok{\# processor{-}specific vectorization}
\NormalTok{  )}
\NormalTok{)}
\end{Highlighting}
\end{Shaded}

\begin{tcolorbox}[enhanced jigsaw, toprule=.15mm, toptitle=1mm, bottomtitle=1mm, leftrule=.75mm, rightrule=.15mm, colframe=quarto-callout-warning-color-frame, left=2mm, colback=white, titlerule=0mm, coltitle=black, colbacktitle=quarto-callout-warning-color!10!white, opacityback=0, title=\textcolor{quarto-callout-warning-color}{\faExclamationTriangle}\hspace{0.5em}{\texttt{-march=native} portability}, arc=.35mm, opacitybacktitle=0.6, breakable, bottomrule=.15mm]

Binaries compiled with \texttt{-march=native} use CPU instructions
specific to the build machine. Do not share compiled Stan objects across
machines with different processor families. Recompile on the target
machine.

\end{tcolorbox}

\subsection{Output size}\label{output-size}

By default Stan saves 8-significant-figure draws. For large models with
many generated quantities, reduce to 6 to halve \texttt{.csv} file sizes
with negligible precision loss:

\begin{Shaded}
\begin{Highlighting}[]
\NormalTok{fit }\OtherTok{\textless{}{-}}\NormalTok{ mod}\SpecialCharTok{$}\FunctionTok{sample}\NormalTok{(..., }\AttributeTok{sig\_figs =} \DecValTok{6}\NormalTok{)}
\end{Highlighting}
\end{Shaded}

\begin{center}\rule{0.5\linewidth}{0.5pt}\end{center}

\section{Quick Reference}\label{sec-stan-quick-ref}

\begin{longtable}[]{@{}
  >{\raggedright\arraybackslash}p{(\linewidth - 4\tabcolsep) * \real{0.3333}}
  >{\raggedright\arraybackslash}p{(\linewidth - 4\tabcolsep) * \real{0.3333}}
  >{\raggedright\arraybackslash}p{(\linewidth - 4\tabcolsep) * \real{0.3333}}@{}}
\toprule\noalign{}
\begin{minipage}[b]{\linewidth}\raggedright
Problem observed
\end{minipage} & \begin{minipage}[b]{\linewidth}\raggedright
Technique
\end{minipage} & \begin{minipage}[b]{\linewidth}\raggedright
Section
\end{minipage} \\
\midrule\noalign{}
\endhead
\bottomrule\noalign{}
\endlastfoot
Redundant computation during leapfrog & Block placement →
\texttt{transformed\ data} & \hyperref[sec-block-placement]{§1} \\
Quantity derived from data only, easier to compute in R & Pre-compute in
R & \hyperref[sec-precompute-r]{§2} \\
Many IID trials, one parameter per session & Sufficient statistics →
\texttt{binomial\_lpmf} & \hyperref[sec-sufficient-stats]{§3} \\
Trial loop over identical distribution & Vectorize the likelihood call &
\hyperref[sec-vectorize]{§4} \\
Changing priors requires recompiling & Priors as data &
\hyperref[sec-priors-as-data]{§5} \\
SBC / LOO is slow & \texttt{run\_diagnostics} flag &
\hyperref[sec-run-diagnostics]{§6} \\
Want PPC/LOO without re-running MCMC & \texttt{\$generate\_quantities()}
post-hoc & \hyperref[sec-generate-quantities]{§7} \\
Divergences / low E-BFMI in multilevel model & Non-centered
parameterization (NCP) & \hyperref[sec-ncp]{§8} \\
Simplex + small kappa → funnel geometry & Logit-normal NCP for simplex &
\hyperref[sec-ncp]{§8} \\
Model has correlation matrix \(\Omega\) & Cholesky factorization
(\texttt{lkj\_corr\_cholesky}) & \hyperref[sec-cholesky]{§9} \\
Sequential model with per-subject trial loops & Contiguous
\texttt{subj\_start}/\texttt{subj\_end} indexing &
\hyperref[sec-subject-indexing]{§10} \\
Slow warm-up on complex / hierarchical model & Pathfinder warm-start &
\hyperref[sec-pathfinder]{§11} \\
Discrete latent variables (indicators, regimes) & Marginalise with
\texttt{log\_mix} & \hyperref[sec-log-mix]{§12} \\
Trial-to-trial latent state in likelihood & Sequential unrolling in
\texttt{transformed\ parameters} &
\hyperref[sec-sequential-state]{§13} \\
Many subjects, multi-core machine & \texttt{reduce\_sum} intra-chain
parallelism & \hyperref[sec-reduce-sum]{§14} \\
Divergences / treedepth / slow warm-up & Sampler options
(\texttt{adapt\_delta}, \texttt{max\_treedepth}) &
\hyperref[sec-sampler-options]{§16} \\
Slow compilation or runtime on algebraic models & Compiler flags
(\texttt{-O3\ -march=native}) & \hyperref[sec-sampler-options]{§16} \\
Large \texttt{.csv} output files & Reduce \texttt{sig\_figs} to 6 &
\hyperref[sec-sampler-options]{§16} \\
\end{longtable}

\chapter{Reparameterization Techniques}\label{sec-reparameterization}

A \textbf{reparameterization} replaces one mathematical form of a
parameter with another that is equivalent for inference but better
behaved for HMC, easier to place priors on, or more directly
interpretable in cognitive terms. This appendix collects every
reparameterization used across the book, organized by the type of
constraint the original parameter imposes.

The key intuition throughout: \textbf{HMC works in unconstrained space}.
Stan handles the Jacobian of any declared constraint
(\texttt{\textless{}lower=0\textgreater{}}, \texttt{simplex}, etc.)
automatically, but the efficiency of the sampler depends on the
\emph{shape} of the posterior in that unconstrained space. A good
reparameterization flattens and regularizes that shape.

\begin{center}\rule{0.5\linewidth}{0.5pt}\end{center}

\section{Part I --- Mapping Bounded Parameters to Unbounded
Space}\label{sec-reparam-bounded}

\subsection{1. Probabilities {[}0, 1{]} --- the Logit
Transform}\label{sec-logit}

\emph{First introduced: Chapter~\ref{sec-inferring-rates} (Chapter 4,
logit bias model)}

A probability \(\theta \in [0,1]\) cannot use a Normal prior directly
--- Normal has infinite support, and the boundaries create hard walls
that distort HMC trajectories. The logit transform maps \([0,1]\) to
\((-\infty, +\infty)\):

\[\text{logit}(\theta) = \log\!\left(\frac{\theta}{1-\theta}\right), \qquad
  \text{logit}^{-1}(x) = \frac{1}{1+e^{-x}} \equiv \texttt{inv\_logit}(x)\]

\textbf{Stan pattern:}

\begin{Shaded}
\begin{Highlighting}[]
\KeywordTok{parameters}\NormalTok{ \{}
  \DataTypeTok{real}\NormalTok{ theta\_logit;           }\CommentTok{// unconstrained — Normal prior works cleanly}
\NormalTok{\}}
\KeywordTok{transformed parameters}\NormalTok{ \{}
  \DataTypeTok{real}\NormalTok{\textless{}}\KeywordTok{lower}\NormalTok{=}\DecValTok{0}\NormalTok{,}\KeywordTok{upper}\NormalTok{=}\DecValTok{1}\NormalTok{\textgreater{} theta = inv\_logit(theta\_logit);}
\NormalTok{\}}
\KeywordTok{model}\NormalTok{ \{}
\NormalTok{  theta\_logit \textasciitilde{} normal(}\DecValTok{0}\NormalTok{, }\FloatTok{1.5}\NormalTok{);   }\CommentTok{// ≈ uniform on probability scale}
\NormalTok{\}}
\end{Highlighting}
\end{Shaded}

\textbf{Choosing the prior width on the logit scale:}

\begin{longtable}[]{@{}
  >{\raggedright\arraybackslash}p{(\linewidth - 2\tabcolsep) * \real{0.5000}}
  >{\raggedright\arraybackslash}p{(\linewidth - 2\tabcolsep) * \real{0.5000}}@{}}
\toprule\noalign{}
\begin{minipage}[b]{\linewidth}\raggedright
\texttt{normal(0,\ σ)} on logit scale
\end{minipage} & \begin{minipage}[b]{\linewidth}\raggedright
Implied prior on probability scale
\end{minipage} \\
\midrule\noalign{}
\endhead
\bottomrule\noalign{}
\endlastfoot
σ = 0.5 & Concentrated near 0.5; rarely below 0.2 or above 0.8 \\
σ = 1.0 & Moderately diffuse \\
σ = 1.5 & Approximately uniform on {[}0,1{]} --- the default for
uninformative use \\
σ = 3.0 & Heavy mass near 0 and 1; implies extreme determinism \\
\end{longtable}

\textbf{Cognitive uses:} choice bias, probability-matching weight,
mixing proportion, learning rate, any rate parameter bounded by
{[}0,1{]}.

\textbf{Stan shortcut:}
\texttt{bernoulli\_logit\_lpmf(y\ \textbar{}\ theta\_logit)} accepts the
log-odds directly, so the \texttt{inv\_logit} transform in
\texttt{transformed\ parameters} is only needed when you want to report
\(\theta\) on the probability scale in \texttt{generated\ quantities}.

\begin{center}\rule{0.5\linewidth}{0.5pt}\end{center}

\subsection{2. Positive-Only Parameters --- the Log
Transform}\label{sec-log}

\emph{First introduced: Chapter~\ref{sec-categorization-exemplars}
(Chapter 11, GCM sensitivity \(c\) and decay \(\lambda\))}

A parameter that must be strictly positive (\(\theta > 0\)) ---
sensitivity, scale, precision, rate --- can be declared as
\texttt{real\textless{}lower=0\textgreater{}}, but this creates a
boundary at zero that HMC must negotiate. An unconstrained
\texttt{real\ log\_theta} with a Normal prior is geometrically smoother.

\[\log\theta \in (-\infty, +\infty), \qquad \theta = e^{\log\theta} > 0\]

\textbf{Stan pattern:}

\begin{Shaded}
\begin{Highlighting}[]
\KeywordTok{parameters}\NormalTok{ \{}
  \DataTypeTok{real}\NormalTok{ log\_c;          }\CommentTok{// log sensitivity — unconstrained}
  \DataTypeTok{real}\NormalTok{ log\_lambda;     }\CommentTok{// log decay rate — unconstrained}
\NormalTok{\}}
\KeywordTok{transformed parameters}\NormalTok{ \{}
  \DataTypeTok{real}\NormalTok{\textless{}}\KeywordTok{lower}\NormalTok{=}\DecValTok{0}\NormalTok{\textgreater{} c      = exp(log\_c);}
  \DataTypeTok{real}\NormalTok{\textless{}}\KeywordTok{lower}\NormalTok{=}\DecValTok{0}\NormalTok{\textgreater{} lambda = exp(log\_lambda);}
\NormalTok{\}}
\KeywordTok{model}\NormalTok{ \{}
\NormalTok{  log\_c      \textasciitilde{} normal(}\DecValTok{0}\NormalTok{, }\DecValTok{1}\NormalTok{);    }\CommentTok{// prior on log scale}
\NormalTok{  log\_lambda \textasciitilde{} normal({-}}\DecValTok{1}\NormalTok{, }\DecValTok{1}\NormalTok{);   }\CommentTok{// informative: moderate decay expected}
\NormalTok{\}}
\end{Highlighting}
\end{Shaded}

\textbf{Interpreting the log-scale prior:}

\begin{longtable}[]{@{}
  >{\raggedright\arraybackslash}p{(\linewidth - 4\tabcolsep) * \real{0.3333}}
  >{\raggedright\arraybackslash}p{(\linewidth - 4\tabcolsep) * \real{0.3333}}
  >{\raggedright\arraybackslash}p{(\linewidth - 4\tabcolsep) * \real{0.3333}}@{}}
\toprule\noalign{}
\begin{minipage}[b]{\linewidth}\raggedright
Prior on log θ
\end{minipage} & \begin{minipage}[b]{\linewidth}\raggedright
Implied range for θ
\end{minipage} & \begin{minipage}[b]{\linewidth}\raggedright
Cognitive interpretation
\end{minipage} \\
\midrule\noalign{}
\endhead
\bottomrule\noalign{}
\endlastfoot
normal(0, 1) & ≈ {[}0.05, 20{]} (2 SD) & Broad uninformative prior for a
scale parameter \\
normal(log(2), 0.5) & Centred at 2, roughly {[}0.7, 5.5{]} &
Informative: moderate sensitivity expected \\
normal(-1, 0.5) & Centred at 0.37, roughly {[}0.15, 0.9{]} & Slow decay
/ weak forgetting \\
\end{longtable}

\textbf{Cognitive uses:} GCM sensitivity \(c\), exponential decay rate
\(\lambda\), RL learning rate \(\alpha\) (if modeled as positive but not
bounded above 1), noise precision, inverse temperature \(\beta\).

\begin{tcolorbox}[enhanced jigsaw, toprule=.15mm, toptitle=1mm, bottomtitle=1mm, leftrule=.75mm, rightrule=.15mm, colframe=quarto-callout-tip-color-frame, left=2mm, colback=white, titlerule=0mm, coltitle=black, colbacktitle=quarto-callout-tip-color!10!white, opacityback=0, title=\textcolor{quarto-callout-tip-color}{\faLightbulb}\hspace{0.5em}{Log transform vs.~\texttt{\textless{}lower=0\textgreater{}} declaration}, arc=.35mm, opacitybacktitle=0.6, breakable, bottomrule=.15mm]

Declaring \texttt{real\textless{}lower=0\textgreater{}\ sigma} in
\texttt{parameters} applies an exponential-map transform internally, so
it is \emph{mathematically equivalent} to declaring
\texttt{real\ log\_sigma} and computing
\texttt{sigma\ =\ exp(log\_sigma)}. The difference is:

\begin{itemize}
\tightlist
\item
  \texttt{real\textless{}lower=0\textgreater{}\ sigma} → Stan chooses
  the unconstrained transform; you place the prior directly on
  \(\sigma\).
\item
  \texttt{real\ log\_sigma} → you control the transform and place the
  prior on \(\log\sigma\).
\end{itemize}

The second form is preferred when you have \textbf{informative prior
beliefs expressible on the log scale} or when the gradient geometry near
zero is poor.

\end{tcolorbox}

\begin{center}\rule{0.5\linewidth}{0.5pt}\end{center}

\subsection{3. Parameters Bounded in {[}0, 1{]} with a Structural
Interpretation --- the Beta Reparameterization}\label{sec-beta-reparam}

\emph{Used in: Chapter~\ref{sec-inferring-rates} (Chapter 4, forgetting
rate), Chapter~\ref{sec-bayesian-models} (Chapter 10, evidence weight
\(\rho\), allocation \(p\))}

Sometimes a {[}0,1{]} parameter is best understood as a \textbf{mean
with an associated uncertainty} rather than as a log-odds. The Beta
distribution has two equivalent parameterizations:

\begin{longtable}[]{@{}
  >{\raggedright\arraybackslash}p{(\linewidth - 4\tabcolsep) * \real{0.3333}}
  >{\raggedright\arraybackslash}p{(\linewidth - 4\tabcolsep) * \real{0.3333}}
  >{\raggedright\arraybackslash}p{(\linewidth - 4\tabcolsep) * \real{0.3333}}@{}}
\toprule\noalign{}
\begin{minipage}[b]{\linewidth}\raggedright
Parameterization
\end{minipage} & \begin{minipage}[b]{\linewidth}\raggedright
Parameters
\end{minipage} & \begin{minipage}[b]{\linewidth}\raggedright
Interpretation
\end{minipage} \\
\midrule\noalign{}
\endhead
\bottomrule\noalign{}
\endlastfoot
Standard & \(\alpha, \beta > 0\) & Shape parameters --- not directly
interpretable \\
Mean + concentration & \(\mu \in (0,1)\), \(\kappa > 0\) & Mean
\(= \mu\); concentration \(= \kappa\); variance
\(= \mu(1-\mu)/(\kappa+1)\) \\
\end{longtable}

\[\alpha = \mu\kappa, \qquad \beta = (1-\mu)\kappa\]

\begin{Shaded}
\begin{Highlighting}[]
\CommentTok{// Mean+concentration parameterization — more interpretable priors}
\KeywordTok{parameters}\NormalTok{ \{}
  \DataTypeTok{real}\NormalTok{\textless{}}\KeywordTok{lower}\NormalTok{=}\DecValTok{0}\NormalTok{,}\KeywordTok{upper}\NormalTok{=}\DecValTok{1}\NormalTok{\textgreater{} mu;     }\CommentTok{// expected value}
  \DataTypeTok{real}\NormalTok{\textless{}}\KeywordTok{lower}\NormalTok{=}\DecValTok{0}\NormalTok{\textgreater{} kappa;          }\CommentTok{// concentration (larger = tighter around mu)}
\NormalTok{\}}
\KeywordTok{transformed parameters}\NormalTok{ \{}
  \DataTypeTok{real}\NormalTok{\textless{}}\KeywordTok{lower}\NormalTok{=}\DecValTok{0}\NormalTok{\textgreater{} alpha = mu * kappa;}
  \DataTypeTok{real}\NormalTok{\textless{}}\KeywordTok{lower}\NormalTok{=}\DecValTok{0}\NormalTok{\textgreater{} beta  = (}\DecValTok{1}\NormalTok{ {-} mu) * kappa;}
\NormalTok{\}}
\KeywordTok{model}\NormalTok{ \{}
\NormalTok{  mu    \textasciitilde{} beta(}\DecValTok{2}\NormalTok{, }\DecValTok{2}\NormalTok{);          }\CommentTok{// soft peak at 0.5; any value plausible}
\NormalTok{  kappa \textasciitilde{} exponential(}\FloatTok{0.1}\NormalTok{);    }\CommentTok{// most agents moderately uncertain}
\NormalTok{  y     \textasciitilde{} beta(alpha, beta);}
\NormalTok{\}}
\end{Highlighting}
\end{Shaded}

\textbf{When to prefer mean+concentration over logit:} When the
cognitive claim \emph{is} about the mean allocation (e.g., ``this agent
allocates \(\rho\) of its weight to direct evidence''), and you want to
reason separately about average tendency (\(\mu\)) and consistency
(\(\kappa\)).

\begin{center}\rule{0.5\linewidth}{0.5pt}\end{center}

\subsection{4. Weight Vectors Summing to One --- Simplex
Reparameterizations}\label{sec-simplex}

\emph{Used in: Chapter~\ref{sec-categorization-exemplars} (Chapter 11,
GCM attention weights \(\mathbf{w}\))}

A \(K\)-dimensional \textbf{simplex} (\(w_k \geq 0\),
\(\sum_k w_k = 1\)) can be parameterized in several ways. For \(K = 2\)
(two features), logit-normal NCP is always preferred (see
Section~\ref{sec-reparameterization-ncp}). For \(K > 2\), the choice is
less clear-cut.

\subsubsection{4a. Dirichlet directly (simple, but problematic for
hierarchical
models)}\label{a.-dirichlet-directly-simple-but-problematic-for-hierarchical-models}

\begin{Shaded}
\begin{Highlighting}[]
\KeywordTok{parameters}\NormalTok{ \{ }\DataTypeTok{simplex}\NormalTok{[K] w; \}}
\KeywordTok{model}\NormalTok{      \{ w \textasciitilde{} dirichlet(alpha); \}  }\CommentTok{// alpha: K{-}vector of concentration hyperparameters}
\end{Highlighting}
\end{Shaded}

Works well for single-subject models where concentration \texttt{alpha}
is fixed. \textbf{Avoid in hierarchical models}: the Dirichlet +
hyperprior-on-concentration geometry creates funnels identical to the
NCP problem (Ch. 11 case study).

\subsubsection{4b. Stick-breaking (unconstrained representation, K−1
free
parameters)}\label{b.-stick-breaking-unconstrained-representation-k1-free-parameters}

Decompose a \(K\)-simplex into \(K-1\) unconstrained values using the
sequential stick-breaking transform. Stan implements this automatically
for \texttt{simplex} declarations --- you never need to code it
manually. But understanding it helps with custom implementations:

\[v_k = \text{logit}(w_k^* + \epsilon_k), \quad
  w_k = w_k^* \cdot \prod_{j<k}(1 - w_j^*), \quad w_K = 1 - \sum_{k<K} w_k\]

For \(K = 2\): reduces exactly to the logit transform on \(w_1\), with
\(w_2 = 1 - w_1\).

\subsubsection{\texorpdfstring{4c. Logit-normal NCP for \(K = 2\)
(recommended for hierarchical
models)}{4c. Logit-normal NCP for K = 2 (recommended for hierarchical models)}}\label{c.-logit-normal-ncp-for-k-2-recommended-for-hierarchical-models}

\begin{Shaded}
\begin{Highlighting}[]
\KeywordTok{parameters}\NormalTok{ \{}
  \DataTypeTok{real}\NormalTok{          pop\_logit\_w\_mean;}
  \DataTypeTok{real}\NormalTok{\textless{}}\KeywordTok{lower}\NormalTok{=}\DecValTok{0}\NormalTok{\textgreater{} pop\_logit\_w\_sd;}
  \DataTypeTok{vector}\NormalTok{[J]     z\_w;                        }\CommentTok{// standard{-}normal offsets}
\NormalTok{\}}
\KeywordTok{transformed parameters}\NormalTok{ \{}
  \ControlFlowTok{for}\NormalTok{ (j }\ControlFlowTok{in} \DecValTok{1}\NormalTok{:J) \{}
    \DataTypeTok{real}\NormalTok{ lw  = pop\_logit\_w\_mean + pop\_logit\_w\_sd * z\_w[j];}
\NormalTok{    w[j][}\DecValTok{1}\NormalTok{]  = inv\_logit(lw);}
\NormalTok{    w[j][}\DecValTok{2}\NormalTok{]  = }\FloatTok{1.0}\NormalTok{ {-} w[j][}\DecValTok{1}\NormalTok{];}
\NormalTok{  \}}
\NormalTok{\}}
\KeywordTok{model}\NormalTok{ \{ z\_w \textasciitilde{} std\_normal(); \}}
\end{Highlighting}
\end{Shaded}

This was the fix for the 4 000 max\_treedepth warnings in the multilevel
decay GCM (Chapter 11). The centered Dirichlet +
\texttt{kappa\ \textasciitilde{}\ exponential(0.1)} produced a funnel;
the logit-normal NCP eliminated it entirely.

\begin{center}\rule{0.5\linewidth}{0.5pt}\end{center}

\subsection{5. K-Way Choice Probabilities --- the Softmax
Transform}\label{sec-softmax}

\emph{Mentioned in: Chapter~\ref{sec-theory-of-mind} (Chapter 9),
Chapter~\ref{sec-categorization-exemplars} (Chapter 11)}

When an agent chooses among \(K > 2\) options and the evidence for each
option is a real-valued score \(v_k\), the \textbf{softmax} maps scores
to probabilities:

\[p_k = \frac{e^{v_k / \tau}}{\sum_{j=1}^K e^{v_j / \tau}}\]

where \(\tau > 0\) is an \textbf{inverse temperature} (higher \(\tau\) =
more random; lower = more deterministic). For \(K = 2\), softmax reduces
to the logistic function.

\textbf{Stan:}

\begin{Shaded}
\begin{Highlighting}[]
\CommentTok{// v is a vector of K utility/evidence values; tau is inverse temperature}
\DataTypeTok{vector}\NormalTok{[K] p = softmax(v / tau);}
\KeywordTok{target +=}\NormalTok{ categorical\_lpmf(choice | p);}

\CommentTok{// Or equivalently, in log space (more numerically stable):}
\KeywordTok{target +=}\NormalTok{ categorical\_logit\_lpmf(choice | v / tau);}
\end{Highlighting}
\end{Shaded}

\begin{tcolorbox}[enhanced jigsaw, toprule=.15mm, toptitle=1mm, bottomtitle=1mm, leftrule=.75mm, rightrule=.15mm, colframe=quarto-callout-tip-color-frame, left=2mm, colback=white, titlerule=0mm, coltitle=black, colbacktitle=quarto-callout-tip-color!10!white, opacityback=0, title=\textcolor{quarto-callout-tip-color}{\faLightbulb}\hspace{0.5em}{Sampling \texttt{tau} vs.~sampling \texttt{beta\ =\ 1/tau}}, arc=.35mm, opacitybacktitle=0.6, breakable, bottomrule=.15mm]

Convention varies across the literature. Stan's
\texttt{categorical\_logit\_lpmf} takes the \textbf{log-odds vector}
directly, so it is natural to sample \(\beta = 1/\tau\) (inverse
temperature) on the log scale:

\begin{Shaded}
\begin{Highlighting}[]
\DataTypeTok{real}\NormalTok{ log\_beta;                        }\CommentTok{// log inverse{-}temperature}
\CommentTok{// ...}
\KeywordTok{target +=}\NormalTok{ categorical\_logit\_lpmf(choice | exp(log\_beta) * v);}
\end{Highlighting}
\end{Shaded}

This places a Normal prior on \(\log\beta\), keeping sampling
unconstrained.

\end{tcolorbox}

\begin{center}\rule{0.5\linewidth}{0.5pt}\end{center}

\section{Part II --- Priors on Scale Parameters}\label{sec-scale-priors}

Scale parameters (\(\sigma\), \(\tau\), \(\kappa\)) must be positive and
appear in hierarchical models as between-subject standard deviations.
The prior on these parameters strongly affects sampling geometry near
zero.

\subsection{6. Exponential, Half-Normal, and
Half-Cauchy}\label{sec-scale-param-priors}

\emph{Used in: Chapter~\ref{sec-multilevel-models} (Chapter 6,
\texttt{sigma\_theta\ \textasciitilde{}\ exponential(lambda)})}

\begin{longtable}[]{@{}
  >{\raggedright\arraybackslash}p{(\linewidth - 6\tabcolsep) * \real{0.2500}}
  >{\raggedright\arraybackslash}p{(\linewidth - 6\tabcolsep) * \real{0.2500}}
  >{\raggedright\arraybackslash}p{(\linewidth - 6\tabcolsep) * \real{0.2500}}
  >{\raggedright\arraybackslash}p{(\linewidth - 6\tabcolsep) * \real{0.2500}}@{}}
\toprule\noalign{}
\begin{minipage}[b]{\linewidth}\raggedright
Prior
\end{minipage} & \begin{minipage}[b]{\linewidth}\raggedright
Stan syntax
\end{minipage} & \begin{minipage}[b]{\linewidth}\raggedright
Shape
\end{minipage} & \begin{minipage}[b]{\linewidth}\raggedright
Recommendation
\end{minipage} \\
\midrule\noalign{}
\endhead
\bottomrule\noalign{}
\endlastfoot
Exponential(\(\lambda\)) &
\texttt{sigma\ \textasciitilde{}\ exponential(lambda)} & Monotone
decreasing from 0; most mass near 0 & Default in this book; works well
when small \(\sigma\) is plausible \\
Half-Normal(\(\sigma_0\)) &
\texttt{sigma\ \textasciitilde{}\ normal(0,\ sigma\_0)\ T{[}0,{]}} &
Bell-shaped; puts more mass away from 0 & Better when you expect
moderate between-subject variance \\
Half-Cauchy(\(s\)) &
\texttt{sigma\ \textasciitilde{}\ cauchy(0,\ s)\ T{[}0,{]}} & Heavy
tails & Use when outlier subjects are plausible \\
\end{longtable}

\textbf{The book's choice} (exponential throughout) reflects a
regularizing stance: most cognitive-parameter hierarchies are relatively
homogeneous, so placing more mass near \(\sigma = 0\) acts as a mild
shrinkage prior that prevents the posterior from drifting to implausibly
large between-subject variances.

\textbf{Parameterizing the exponential:} \texttt{exponential(lambda)}
has mean \(1/\lambda\). Setting \texttt{lambda\ =\ 1} means the prior
expects \(\sigma \approx 1\) on the log-odds scale. For parameters whose
population mean is on a different scale, adjust accordingly.

\begin{Shaded}
\begin{Highlighting}[]
\CommentTok{// Exponential prior for between{-}subject SD on the logit scale}
\NormalTok{sigma \textasciitilde{} exponential(}\DecValTok{1}\NormalTok{);    }\CommentTok{// mean = 1 log{-}odds unit ≈ ±25 percentage points}

\CommentTok{// Half{-}Normal alternative — more permissive}
\NormalTok{sigma \textasciitilde{} normal(}\DecValTok{0}\NormalTok{, }\FloatTok{0.5}\NormalTok{) }\KeywordTok{T}\NormalTok{[}\DecValTok{0}\NormalTok{,];  }\CommentTok{// 95\% prior mass below \textasciitilde{}1 log{-}odds unit}
\end{Highlighting}
\end{Shaded}

\begin{tcolorbox}[enhanced jigsaw, toprule=.15mm, toptitle=1mm, bottomtitle=1mm, leftrule=.75mm, rightrule=.15mm, colframe=quarto-callout-note-color-frame, left=2mm, colback=white, titlerule=0mm, coltitle=black, colbacktitle=quarto-callout-note-color!10!white, opacityback=0, title=\textcolor{quarto-callout-note-color}{\faInfo}\hspace{0.5em}{Why not \texttt{\textless{}lower=0\textgreater{}\ real\ sigma} +
\texttt{normal(0,\ σ₀)}?}, arc=.35mm, opacitybacktitle=0.6, breakable, bottomrule=.15mm]

Stan accepts \texttt{sigma\ \textasciitilde{}\ normal(0,\ 0.5)} even
when \texttt{sigma} is declared
\texttt{\textless{}lower=0\textgreater{}}. This is a \textbf{truncated
normal} --- Stan evaluates the density only on \([0, \infty)\). For HMC,
declaring the constraint \texttt{\textless{}lower=0\textgreater{}} is
important; the prior statement handles the shape. The explicit
\texttt{T{[}0,{]}} truncation notation is only needed inside
\texttt{stan\_code} blocks where Stan cannot infer the constraint from
the declaration.

\end{tcolorbox}

\begin{center}\rule{0.5\linewidth}{0.5pt}\end{center}

\section{Part III --- Structural
Reparameterizations}\label{sec-structural-reparam}

\subsection{7. Non-Centered Parameterization
(NCP)}\label{sec-reparameterization-ncp}

\emph{Covered in detail in Section~\ref{sec-ncp} (Appendix B)}

The NCP separates the \textbf{location} (\(\mu\)) and \textbf{scale}
(\(\sigma\)) of a hierarchical parameter from the individual-level
deviations (\(z_j \sim \mathcal{N}(0,1)\)).

\begin{Shaded}
\begin{Highlighting}[]
\CommentTok{// Centered: theta\_j \textasciitilde{} normal(mu, sigma)  ← funnel when sigma → 0}
\CommentTok{// Non{-}Centered: theta\_j = mu + sigma * z\_j, z\_j \textasciitilde{} std\_normal()  ← always Gaussian}
\end{Highlighting}
\end{Shaded}

The NCP is the single most impactful reparameterization for hierarchical
cognitive models. See Section~\ref{sec-ncp} for full details, the
K-simplex extension, and the Chapter 11 case study.

\begin{center}\rule{0.5\linewidth}{0.5pt}\end{center}

\subsection{8. Sum-to-Zero Constraint for Group
Effects}\label{sec-sum-to-zero}

\emph{Relevant context:
\href{https://avehtari.github.io/Bayesian-Workflow/park_rule/park_rule.html}{Vehtari
et al.~--- the park rule}}

The standard \textbf{corner constraint} for ANOVA-style group effects
fixes one group as the reference (its effect is zero) and estimates all
others relative to it. This is a valid identification strategy but it
creates an asymmetry: the reference group's effect is absorbed into the
intercept, making the intercept's prior inadvertently informative about
one specific group.

The \textbf{sum-to-zero constraint} (\(\sum_{k=1}^K \alpha_k = 0\)) is a
symmetric alternative. Every group effect is estimated relative to the
grand mean, the intercept retains a clean interpretation as the overall
mean, and priors on group effects are exchangeable.

\textbf{Stan implementation via QR decomposition:}

\begin{Shaded}
\begin{Highlighting}[]
\KeywordTok{data}\NormalTok{ \{}
  \DataTypeTok{int}\NormalTok{\textless{}}\KeywordTok{lower}\NormalTok{=}\DecValTok{1}\NormalTok{\textgreater{} K;                   }\CommentTok{// number of groups}
  \DataTypeTok{matrix}\NormalTok{[K}\DecValTok{{-}1}\NormalTok{, K] A;                 }\CommentTok{// sum{-}to{-}zero contrast matrix (precomputed in R)}
\NormalTok{\}}
\KeywordTok{parameters}\NormalTok{ \{}
  \DataTypeTok{vector}\NormalTok{[K}\DecValTok{{-}1}\NormalTok{] alpha\_raw;            }\CommentTok{// K{-}1 free parameters}
\NormalTok{\}}
\KeywordTok{transformed parameters}\NormalTok{ \{}
  \DataTypeTok{vector}\NormalTok{[K] alpha = A\textquotesingle{} * alpha\_raw; }\CommentTok{// K effects summing to zero}
\NormalTok{\}}
\KeywordTok{model}\NormalTok{ \{}
\NormalTok{  alpha\_raw \textasciitilde{} normal(}\DecValTok{0}\NormalTok{, }\DecValTok{1}\NormalTok{);}
\NormalTok{\}}
\end{Highlighting}
\end{Shaded}

\textbf{Build the contrast matrix in R:}

\begin{Shaded}
\begin{Highlighting}[]
\CommentTok{\# Helmert{-}style QR contrast matrix enforcing sum{-}to{-}zero}
\NormalTok{make\_sum\_zero\_matrix }\OtherTok{\textless{}{-}} \ControlFlowTok{function}\NormalTok{(K) \{}
\NormalTok{  Q }\OtherTok{\textless{}{-}} \FunctionTok{qr.Q}\NormalTok{(}\FunctionTok{qr}\NormalTok{(}\FunctionTok{rep}\NormalTok{(}\DecValTok{1}\NormalTok{, K)), }\AttributeTok{complete =} \ConstantTok{TRUE}\NormalTok{)[, }\SpecialCharTok{{-}}\DecValTok{1}\NormalTok{]}
  \FunctionTok{t}\NormalTok{(Q)   }\CommentTok{\# (K{-}1) x K matrix: A \%*\% alpha = 0 iff sum(alpha) = 0}
\NormalTok{\}}

\NormalTok{A }\OtherTok{\textless{}{-}} \FunctionTok{make\_sum\_zero\_matrix}\NormalTok{(K)   }\CommentTok{\# pass this as data to Stan}
\end{Highlighting}
\end{Shaded}

\begin{tcolorbox}[enhanced jigsaw, toprule=.15mm, toptitle=1mm, bottomtitle=1mm, leftrule=.75mm, rightrule=.15mm, colframe=quarto-callout-tip-color-frame, left=2mm, colback=white, titlerule=0mm, coltitle=black, colbacktitle=quarto-callout-tip-color!10!white, opacityback=0, title=\textcolor{quarto-callout-tip-color}{\faLightbulb}\hspace{0.5em}{When to prefer sum-to-zero over corner constraints}, arc=.35mm, opacitybacktitle=0.6, breakable, bottomrule=.15mm]

Use sum-to-zero when (a) no group is a natural reference, (b) you want
exchangeable priors on all group effects, or (c) the intercept should
represent the grand mean rather than one group's baseline. The
constraint has no effect on the likelihood --- only on identifiability
and prior geometry.

\end{tcolorbox}

\begin{center}\rule{0.5\linewidth}{0.5pt}\end{center}

\subsection{10. Probit Link --- Signal Detection
Theory}\label{sec-probit}

The \textbf{probit link} uses the standard Normal CDF \(\Phi\) in place
of the logistic function:

\[p = \Phi(d'), \qquad d' = \Phi^{-1}(p) \equiv \text{probit}(p)\]

Stan function: \texttt{Phi(x)} (Normal CDF), \texttt{Phi\_approx(x)}
(fast rational approximation).

The probit link is the natural choice whenever the cognitive model is
derived from \textbf{signal detection theory (SDT)}: the sensitivity
parameter \(d'\) is directly the standardized distance between signal
and noise distributions, and the criterion \(c\) is the decision
threshold in the same units.

\begin{Shaded}
\begin{Highlighting}[]
\KeywordTok{parameters}\NormalTok{ \{}
  \DataTypeTok{real}\NormalTok{ d\_prime;     }\CommentTok{// sensitivity (positive = above chance)}
  \DataTypeTok{real}\NormalTok{ criterion;   }\CommentTok{// decision criterion}
\NormalTok{\}}
\KeywordTok{model}\NormalTok{ \{}
\NormalTok{  d\_prime   \textasciitilde{} normal(}\DecValTok{0}\NormalTok{, }\DecValTok{2}\NormalTok{);}
\NormalTok{  criterion \textasciitilde{} normal(}\DecValTok{0}\NormalTok{, }\DecValTok{1}\NormalTok{);}
  \ControlFlowTok{for}\NormalTok{ (i }\ControlFlowTok{in} \DecValTok{1}\NormalTok{:N) \{}
    \ControlFlowTok{if}\NormalTok{ (signal[i] == }\DecValTok{1}\NormalTok{)}
      \KeywordTok{target +=}\NormalTok{ log(Phi( d\_prime/}\DecValTok{2}\NormalTok{ {-} criterion));   }\CommentTok{// hit}
    \ControlFlowTok{else}
      \KeywordTok{target +=}\NormalTok{ log(Phi({-}d\_prime/}\DecValTok{2}\NormalTok{ {-} criterion));   }\CommentTok{// correct rejection}
\NormalTok{  \}}
\NormalTok{\}}
\end{Highlighting}
\end{Shaded}

\textbf{Logit vs.~probit in practice:} The two links are very similar
numerically (probit ≈ logit × 0.607), and prior sensitivity analyses
rarely distinguish them. Choose probit when the theoretical framework is
SDT and interpretability in \(d'\) units matters; choose logit for all
other binary-outcome models because \texttt{bernoulli\_logit\_lpmf} is
faster and numerically more stable.

\begin{center}\rule{0.5\linewidth}{0.5pt}\end{center}

\subsection{11. Correlation Matrices --- LKJ and
Cholesky}\label{sec-lkj}

\emph{Covered in Section~\ref{sec-cholesky} (Appendix B); introduced in
Chapter~\ref{sec-multilevel-models} (Chapter 6)}

The \textbf{LKJ distribution} is the natural prior for correlation
matrices. Its single parameter \(\eta\) controls how much mass is placed
near the identity matrix:

\begin{longtable}[]{@{}
  >{\raggedright\arraybackslash}p{(\linewidth - 2\tabcolsep) * \real{0.5000}}
  >{\raggedright\arraybackslash}p{(\linewidth - 2\tabcolsep) * \real{0.5000}}@{}}
\toprule\noalign{}
\begin{minipage}[b]{\linewidth}\raggedright
\(\eta\)
\end{minipage} & \begin{minipage}[b]{\linewidth}\raggedright
Implied prior on correlations
\end{minipage} \\
\midrule\noalign{}
\endhead
\bottomrule\noalign{}
\endlastfoot
\(\eta = 1\) & Uniform over all valid correlation matrices \\
\(\eta = 2\) & Mild regularization toward identity (off-diagonal
correlations shrunk toward 0) \\
\(\eta \to \infty\) & All mass on the identity (zero correlations) \\
\end{longtable}

\textbf{Always} sample the Cholesky factor
\texttt{cholesky\_factor\_corr{[}K{]}\ L\_Omega} and apply the
\texttt{lkj\_corr\_cholesky} prior. Recover
\(\Omega = L_\Omega L_\Omega^\top\) in \texttt{generated\ quantities}.
See Section~\ref{sec-cholesky} for full code and a cost comparison.

\begin{center}\rule{0.5\linewidth}{0.5pt}\end{center}

\section{Part IV --- Numerical Stability}\label{sec-numerical-stability}

\subsection{10. Log-Space Arithmetic}\label{sec-log-space}

\emph{Used throughout: \texttt{log\_mix}, \texttt{log\_sum\_exp},
\texttt{log1m} in Chapters 8--11}

Computing probabilities directly (multiplying small numbers together,
then taking logs) leads to underflow for long trial sequences. Staying
in \textbf{log space} throughout avoids this.

\begin{longtable}[]{@{}
  >{\raggedright\arraybackslash}p{(\linewidth - 4\tabcolsep) * \real{0.3333}}
  >{\raggedright\arraybackslash}p{(\linewidth - 4\tabcolsep) * \real{0.3333}}
  >{\raggedright\arraybackslash}p{(\linewidth - 4\tabcolsep) * \real{0.3333}}@{}}
\toprule\noalign{}
\begin{minipage}[b]{\linewidth}\raggedright
Operation
\end{minipage} & \begin{minipage}[b]{\linewidth}\raggedright
Unstable form
\end{minipage} & \begin{minipage}[b]{\linewidth}\raggedright
Stable Stan function
\end{minipage} \\
\midrule\noalign{}
\endhead
\bottomrule\noalign{}
\endlastfoot
\(\log(p_1 + p_2)\) & \texttt{log(exp(lp1)\ +\ exp(lp2))} &
\texttt{log\_sum\_exp(lp1,\ lp2)} \\
\(\log(\pi p_1 + (1-\pi) p_2)\) &
\texttt{log(pi*exp(lp1)\ +\ (1-pi)*exp(lp2))} &
\texttt{log\_mix(pi,\ lp1,\ lp2)} \\
\(\log(1 + e^x)\) & \texttt{log(1\ +\ exp(x))} &
\texttt{log1p\_exp(x)} \\
\(\log(1 - e^x)\) (requires \(x < 0\)) & \texttt{log(1\ -\ exp(x))} &
\texttt{log1m\_exp(x)} \\
\(\log(1 - p)\) & \texttt{log(1\ -\ p)} & \texttt{log1m(p)} \\
\(\log \Phi(x)\) (Normal CDF) & \texttt{log(Phi(x))} &
\texttt{log(Phi\_approx(x))} or
\texttt{normal\_lcdf(x\ \textbackslash{}\textbar{}\ 0,1)} \\
\end{longtable}

\begin{Shaded}
\begin{Highlighting}[]
\CommentTok{// Mixture likelihood in log space — numerically stable}
\ControlFlowTok{for}\NormalTok{ (i }\ControlFlowTok{in} \DecValTok{1}\NormalTok{:N) \{}
  \DataTypeTok{real}\NormalTok{ lp1 = bernoulli\_logit\_lpmf(y[i] | alpha);}
  \DataTypeTok{real}\NormalTok{ lp2 = bernoulli\_logit\_lpmf(y[i] | nu);}
  \KeywordTok{target  +=}\NormalTok{ log\_mix(pi, lp1, lp2);   }\CommentTok{// never calls exp() on small values}
\NormalTok{\}}

\CommentTok{// Accumulated log{-}probability for a trajectory}
\DataTypeTok{real}\NormalTok{ log\_p = }\FloatTok{0.0}\NormalTok{;}
\ControlFlowTok{for}\NormalTok{ (t }\ControlFlowTok{in} \DecValTok{1}\NormalTok{:T)}
\NormalTok{  log\_p += log\_mix(pi[t], lp\_A[t], lp\_B[t]);}
\KeywordTok{target +=}\NormalTok{ log\_p;}
\end{Highlighting}
\end{Shaded}

\begin{tcolorbox}[enhanced jigsaw, toprule=.15mm, toptitle=1mm, bottomtitle=1mm, leftrule=.75mm, rightrule=.15mm, colframe=quarto-callout-warning-color-frame, left=2mm, colback=white, titlerule=0mm, coltitle=black, colbacktitle=quarto-callout-warning-color!10!white, opacityback=0, title=\textcolor{quarto-callout-warning-color}{\faExclamationTriangle}\hspace{0.5em}{When \texttt{log\_sum\_exp} is not available}, arc=.35mm, opacitybacktitle=0.6, breakable, bottomrule=.15mm]

For vectors: \texttt{log\_sum\_exp(v)} sums all elements of \texttt{v}
in log-space. For accumulation inside a loop: maintain a running
\texttt{log\_p} and add each term. Never compute
\texttt{sum(exp(log\_probs))} and then take the log --- this is the
numerically unstable version that Stan's log-space functions exist to
replace.

\end{tcolorbox}

\begin{center}\rule{0.5\linewidth}{0.5pt}\end{center}

\section{Quick Reference}\label{sec-reparam-quick-ref}

\begin{longtable}[]{@{}
  >{\raggedright\arraybackslash}p{(\linewidth - 6\tabcolsep) * \real{0.2500}}
  >{\raggedright\arraybackslash}p{(\linewidth - 6\tabcolsep) * \real{0.2500}}
  >{\raggedright\arraybackslash}p{(\linewidth - 6\tabcolsep) * \real{0.2500}}
  >{\raggedright\arraybackslash}p{(\linewidth - 6\tabcolsep) * \real{0.2500}}@{}}
\toprule\noalign{}
\begin{minipage}[b]{\linewidth}\raggedright
Parameter type
\end{minipage} & \begin{minipage}[b]{\linewidth}\raggedright
Constraint
\end{minipage} & \begin{minipage}[b]{\linewidth}\raggedright
Recommended reparameterization
\end{minipage} & \begin{minipage}[b]{\linewidth}\raggedright
Chapter
\end{minipage} \\
\midrule\noalign{}
\endhead
\bottomrule\noalign{}
\endlastfoot
Probability / rate & {[}0, 1{]} & Logit → \texttt{real\ theta\_logit},
\texttt{inv\_logit} in TP & Ch 4 \\
Positive scale / rate & \textgreater{} 0 & Log →
\texttt{real\ log\_theta}, \texttt{exp} in TP & Ch 11 \\
{[}0,1{]} with mean+variance story & {[}0, 1{]} & Beta
mean+concentration (μ, κ) & Ch 10 \\
Two-feature attention weight & simplex{[}2{]} & Logit-normal NCP & Ch
11 \\
K-feature attention weight & simplex{[}K{]} & Stan \texttt{simplex}
(auto stick-break); NCP for hierarchical & Ch 11 \\
K-way choice probability & probability K-vector & Softmax of
unconstrained scores / \texttt{categorical\_logit} & Ch 9 \\
Individual params drawn from population & real & Non-centered:
\(z \sim \mathcal{N}(0,1)\), \(\theta = \mu + \sigma z\) & Ch 6 \\
Between-subject standard deviation & \textgreater{} 0 &
\texttt{exponential(1)} or \texttt{normal(0,\ 0.5)\ T{[}0,{]}} & Ch 6 \\
Correlation matrix & positive-definite, diag = 1 & Cholesky factor + LKJ
prior & Ch 6 \\
Group effects (no natural reference) & sum-to-zero & QR contrast matrix
→ \texttt{A\textquotesingle{}\ *\ alpha\_raw} & App C §8 \\
SDT sensitivity & unconstrained & Probit: \texttt{Phi(d\_prime)} &
--- \\
Sum of log-probabilities & \((-\infty, 0]\) & \texttt{log\_sum\_exp},
\texttt{log\_mix}, \texttt{log1p\_exp} & Ch 8--11 \\
\end{longtable}

\phantomsection\label{refs}
\begin{CSLReferences}{1}{0}
\bibitem[\citeproctext]{ref-akaike1974new}
Akaike, Hirotugu. 1974. {``A New Look at the Statistical Model
Identification.''} \emph{IEEE Transactions on Automatic Control} 19 (6):
716--23.

\bibitem[\citeproctext]{ref-anderson1991adaptive}
Anderson, John R. 1991. {``The Adaptive Nature of Human
Categorization.''} \emph{Psychological Review} 98 (3): 409.

\bibitem[\citeproctext]{ref-ashby1998neurobiological}
Ashby, F Gregory, L Elisa Alfonso-Reese, Waldemar U Turken, and Elliott
M Waldron. 1998. {``A Neurobiological Theory of Perceptual Category
Learning.''} \emph{Psychological Review} 105 (3): 442.

\bibitem[\citeproctext]{ref-ashby1993relations}
Ashby, F. Gregory, and W. Todd Maddox. 1993. {``Relations Between
Prototype, Exemplar, and Decision Bound Models of Categorization.''}
\emph{Journal of Mathematical Psychology} 37 (3): 372--400.

\bibitem[\citeproctext]{ref-betancourt2013hamiltonian}
Betancourt, Michael, and Mark Girolami. 2013. {``Hamiltonian Monte Carlo
for Hierarchical Models.''} \emph{Current Trends in Bayesian Methodology
with Applications} 79: 30.

\bibitem[\citeproctext]{ref-bowers2012bayesian}
Bowers, Jeffrey S, and Colin J Davis. 2012. {``Bayesian Causal Models of
Human Categorization.''} \emph{Psychological Review} 119 (2): 389.

\bibitem[\citeproctext]{ref-bruner1956study}
Bruner, Jerome Seymour, Jacqueline J Goodnow, and George A Austin. 1956.
\emph{A Study of Thinking}. Transaction Publishers.

\bibitem[\citeproctext]{ref-dayan2000learning}
Dayan, Peter, Sham Kakade, and P Read Montague. 2000. {``Learning and
Selective Attention.''} \emph{Nature Neuroscience} 3 (12): 1218--23.

\bibitem[\citeproctext]{ref-erickson1998rules}
Erickson, Michael A, and John K Kruschke. 1998. {``Rules and Exemplars
in Category Learning.''} \emph{Journal of Experimental Psychology:
General} 127 (2): 107--40.

\bibitem[\citeproctext]{ref-ernst2002humans}
Ernst, Marc O, and Martin S Banks. 2002. {``Humans Integrate Visual and
Haptic Information in a Statistically Optimal Fashion.''} \emph{Nature}
415 (6870): 429--33.

\bibitem[\citeproctext]{ref-estes1994classification}
Estes, William Kaye. 1994. \emph{Classification and Cognition}. Oxford
University Press.

\bibitem[\citeproctext]{ref-friston2005theory}
Friston, Karl. 2005. {``A Theory of Cortical Responses.''}
\emph{Philosophical Transactions of the Royal Society B: Biological
Sciences} 360 (1456): 815--36.

\bibitem[\citeproctext]{ref-geisler1989ideal}
Geisler, Wilson S. 1989. {``Sequential Ideal-Observer Analysis of Visual
Discriminations.''} \emph{Psychological Review} 96 (2): 267.

\bibitem[\citeproctext]{ref-gelman1995bayesian}
Gelman, Andrew, John B Carlin, Hal S Stern, and Donald B Rubin. 1995.
{``Bayesian Data Analysis.''}

\bibitem[\citeproctext]{ref-gelman1992inference}
Gelman, Andrew, and Donald B Rubin. 1992. {``Inference from Iterative
Simulation Using Multiple Sequences.''} \emph{Statistical Science} 7
(4): 457--72.

\bibitem[\citeproctext]{ref-gluck1988conditioning}
Gluck, Mark A, and Gordon H Bower. 1988. {``From Conditioning to
Category Learning: An Adaptive Network Model.''} \emph{Journal of
Experimental Psychology: General} 117 (3): 227.

\bibitem[\citeproctext]{ref-goodman2008rational}
Goodman, Noah D, Joshua B Tenenbaum, Jacob Feldman, and Thomas L
Griffiths. 2008. {``A Rational Analysis of Rule-Based Concept
Learning.''} \emph{Cognitive Science} 32 (1): 108--54.

\bibitem[\citeproctext]{ref-gordon1993novel}
Gordon, Neil J, David J Salmond, and Adrian FM Smith. 1993. {``Novel
Approach to Nonlinear/Non-Gaussian Bayesian State Estimation.''}
\emph{IEE Proceedings F (Radar and Signal Processing)} 140 (2): 107--13.

\bibitem[\citeproctext]{ref-helmholtz1867handbuch}
Helmholtz, Hermann von. 1867. \emph{Handbuch Der Physiologischen Optik:
Mit 213 in Den Text Eingedruckten Holzschnitten Und 11 Tafeln}. Vol. 9.
Leopold Voss.

\bibitem[\citeproctext]{ref-henderson1950estimation}
Henderson, Charles R. 1950. {``Estimation of Genetic Parameters.''}
\emph{Annals of Mathematical Statistics} 21: 309--10.

\bibitem[\citeproctext]{ref-hull1943principles}
Hull, Clark Leonard. 1943. \emph{Principles of Behavior: An Introduction
to Behavior Theory}. Appleton-century-crofts.

\bibitem[\citeproctext]{ref-james1961estimation}
James, William, and Charles Stein. 1961. {``Estimation with Quadratic
Loss''} 1: 361--79.

\bibitem[\citeproctext]{ref-jones2011bayesian}
Jones, Matt, and Bradley C Love. 2011. {``Bayesian Fundamentalism or
Enlightenment? On the Explanatory Status of Bayesian Models of
Cognition.''} \emph{Psychological Review} 118 (1): 169.

\bibitem[\citeproctext]{ref-kalman1960new}
Kalman, Rudolph Emil. 1960. {``A New Approach to Linear Filtering and
Prediction Problems.''}

\bibitem[\citeproctext]{ref-kamin1969predictability}
Kamin, Leon J. 1969. {``Predictability, Surprise, Attention, and
Conditioning.''} In \emph{Punishment and Aversive Behavior}.
Appleton-Century-Crofts.

\bibitem[\citeproctext]{ref-kruschke1992alcove}
Kruschke, John K. 1992. {``ALCOVE: An Exemplar-Based Connectionist Model
of Category Learning.''} \emph{Psychological Review} 99 (1): 22--44.

\bibitem[\citeproctext]{ref-kruschke1993human}
Kruschke, John K. 1993. {``Human Category Learning: Implications for
Backpropagation Models.''} \emph{Connection Science} 5 (1): 3--36.

\bibitem[\citeproctext]{ref-levine1975cognitive}
Levine, Marvin. 1975. \emph{A Cognitive Theory of Learning: Research on
Hypothesis Testing}. Lawrence Erlbaum.

\bibitem[\citeproctext]{ref-lindley1972bayes}
Lindley, Dennis V, and Adrian FM Smith. 1972. {``Bayes Estimates for the
Linear Model.''} \emph{Journal of the Royal Statistical Society: Series
B (Methodological)} 34 (1): 1--18.

\bibitem[\citeproctext]{ref-love2004sustain}
Love, Bradley C, Douglas L Medin, and Todd M Gureckis. 2004. {``SUSTAIN:
A Network Model of Category Learning.''} \emph{Psychological Review} 111
(2): 309--32.

\bibitem[\citeproctext]{ref-mathys2011bayesian}
Mathys, Christoph, Jean Daunizeau, Karl J Friston, and Klaas E Stephan.
2011. {``A Bayesian Foundation for Individual Learning Under
Uncertainty.''} \emph{Frontiers in Human Neuroscience} 5: 39.

\bibitem[\citeproctext]{ref-medin1978context}
Medin, Douglas L, and Marguerite M Schaffer. 1978. {``Context Theory of
Classification Learning.''} \emph{Psychological Review} 85 (3): 207--38.

\bibitem[\citeproctext]{ref-neal2003slice}
Neal, Radford M. 2003. {``Slice Sampling.''} \emph{The Annals of
Statistics} 31 (3): 705--67.

\bibitem[\citeproctext]{ref-nosofsky1986attention}
Nosofsky, Robert M. 1986. {``Attention, Similarity, and the
Identification-Categorization Relationship.''} \emph{Journal of
Experimental Psychology: General} 115 (1): 39--61.

\bibitem[\citeproctext]{ref-nosofsky1994rule}
Nosofsky, Robert M, Thomas J Palmeri, and Stephen C McKinley. 1994.
{``Rule-Plus-Exception Model of Classification Learning.''}
\emph{Psychological Review} 101 (1): 53.

\bibitem[\citeproctext]{ref-posner1968genesis}
Posner, Michael I, and Steven W Keele. 1968. {``On the Genesis of
Abstract Ideas.''} \emph{Journal of Experimental Psychology} 77 (3p1):
353.

\bibitem[\citeproctext]{ref-rao1999predictive}
Rao, Rajesh PN, and Dana H Ballard. 1999. {``Predictive Coding in the
Visual Cortex: A Functional Interpretation of Some Extra-Classical
Receptive-Field Effects.''} \emph{Nature Neuroscience} 2 (1): 79--87.

\bibitem[\citeproctext]{ref-reed1972pattern}
Reed, Stephen K. 1972. {``Pattern Recognition and Categorization.''}
\emph{Psychonomic Monograph Supplements}.

\bibitem[\citeproctext]{ref-rescorla1972theory}
Rescorla, Robert A, and Allan R Wagner. 1972. {``A Theory of Pavlovian
Conditioning: Variations in the Effectiveness of Reinforcement and
Nonreinforcement.''} In \emph{Classical Conditioning II: Current
Research and Theory}, 2:64--99.

\bibitem[\citeproctext]{ref-rosch1975cognitive}
Rosch, Eleanor. 1975. {``Cognitive Representations of Semantic
Categories.''} \emph{Journal of Experimental Psychology: General} 104
(3): 192.

\bibitem[\citeproctext]{ref-rosch1973natural}
Rosch, Eleanor H. 1973. {``Natural Categories.''} \emph{Cognitive
Psychology} 4 (3): 328--50.

\bibitem[\citeproctext]{ref-rosch1976basic}
Rosch, Eleanor, Carolyn B Mervis, Wayne D Gray, David M Johnson, and
Panni Boyes-Braem. 1976. {``Basic Objects in Natural Categories.''}
\emph{Cognitive Psychology} 8 (3): 382--439.

\bibitem[\citeproctext]{ref-rouder2008bayesian}
Rouder, Jeffrey N, and Jun Lu. 2008. {``Bayesian Hierarchical Models for
Perceptual Processes.''} \emph{Psychonomic Bulletin \& Review} 15 (6):
1153--67.

\bibitem[\citeproctext]{ref-rouder2005introduction}
Rouder, Jeffrey N, Jun Lu, Paul Speckman, Dongchu Sun, and Yi Jiang.
2005. {``An Introduction to Bayesian Hierarchical Models with an
Application to the SET Generalization Task.''} \emph{Psychonomic
Bulletin \& Review} 12 (4): 573--604.

\bibitem[\citeproctext]{ref-schultz1997neural}
Schultz, Wolfram, Peter Dayan, and P Read Montague. 1997. {``A Neural
Substrate of Prediction and Reward.''} \emph{Science} 275 (5306):
1593--99.

\bibitem[\citeproctext]{ref-schwarz1978estimating}
Schwarz, Gideon. 1978. {``Estimating the Dimension of a Model.''}
\emph{The Annals of Statistics} 6 (2): 461--64.

\bibitem[\citeproctext]{ref-smith1998prototype}
Smith, J David, and John Paul Minda. 1998. {``Prototype and Exemplar
Accounts of Category Learning.''} \emph{Journal of Experimental
Psychology: Learning, Memory, and Cognition} 24 (6): 1411.

\bibitem[\citeproctext]{ref-stone1977asymptotic}
Stone, Mervyn. 1977. {``An Asymptotic Equivalence of Choice of Model by
Cross-Validation and Akaike's Criterion.''} \emph{Journal of the Royal
Statistical Society: Series B (Methodological)} 39 (1): 44--47.

\bibitem[\citeproctext]{ref-tenenbaum1999bayesian}
Tenenbaum, Joshua B. 1999. {``A Bayesian Framework for Concept
Learning.''}

\bibitem[\citeproctext]{ref-tenenbaum2011grow}
Tenenbaum, Joshua B, Charles Kemp, Thomas L Griffiths, and Noah D
Goodman. 2011. {``How to Grow a Mind: Statistics, Structure, and
Abstraction.''} \emph{Science} 331 (6022): 1279--85.

\bibitem[\citeproctext]{ref-tolman1948cognitive}
Tolman, Edward C. 1948. {``Cognitive Maps in Rats and Men.''}
\emph{Psychological Review} 55 (4): 189.

\bibitem[\citeproctext]{ref-vehtari2017practical}
Vehtari, Aki, Andrew Gelman, and Jonah Gabry. 2017. {``Practical
Bayesian Model Evaluation Using Leave-One-Out Cross-Validation and
Waic.''} \emph{Statistics and Computing} 27 (5): 1413--32.

\bibitem[\citeproctext]{ref-wills2017progress}
Wills, A. J., G. O'Connell, C. E. Edmunds, and A. B. Inkster. 2017.
{``Progress in Modeling Through Distributed Collaboration: Concepts,
Tools and Category-Learning Examples.''} In \emph{Psychology of Learning
and Motivation}, 66:79--115. Academic Press.

\bibitem[\citeproctext]{ref-zhang2022pathfinder}
Zhang, Lu, Bob Carpenter, Andrew Gelman, and Aki Vehtari. 2022.
{``Pathfinder: Parallel Quasi-Newton Variational Inference.''}
\emph{Journal of Machine Learning Research} 23 (306): 1--49.

\end{CSLReferences}




\end{document}
